% SPDX-FileCopyrightText: 2024-2026 Fabian Januszewski
% SPDX-License-Identifier: CC-BY-4.0

%======================================================================
% gl2_contract.tex --- THE UNIFIED BINDING MATHEMATICAL CONTRACT
%
%   VERSION 1.26.4 (slot 26, 2026-11-10: SIGNED BY THE OWNER 2026-11-10 ("I sign those title pages."); the text of 1.26.4-pre.1, unchanged: header, date, release record and changelog only -- no statement, label, number, tag, status, marker (229 on 88 labels), written word or value moves; the title page names no sign-off; the 1.26.4-pre.1 and earlier history below is kept)
%   VERSION 1.26.4-pre.1 (slot 26, 2026-11-10, PRE-RELEASE for owner signature, NOT signed): title page only, no change of mathematical content (public release HACS 0.1.1): the \date line names only the version, the date and the revision history; the words 'binding document; signed by the owner 2026-11-09' leave the title page (the sign-off record stays in the Release record and the changelog of Appendix app:changelog). No statement, label, marker (229 on 88), number, tag or value moves.
%   VERSION 1.26.3 (slot 26, 2026-11-09: SIGNED BY THE OWNER 2026-11-09 ("I resign everything", effective on the referee's ACCEPT, RELFIX_signed_referee.md); the text of 1.26.3-pre.1, unchanged: header, date, release record and changelog only -- no statement, label, number, tag, status, marker (229 on 88 labels), written word or value moves; the 1.26.3-pre.1 and earlier history below is kept)
%   VERSION 1.26.3-pre.1 (slot 26, 2026-11-09, PRE-RELEASE for owner signature, NOT signed): text-only: private path removed, no change of mathematical content (release audit REL-AUDIT S1): one scratch-script path reduced to its file name; two report file names that carried a cluster name replaced by a description. No statement, label, marker (229 on 88), number, tag or value moves.
%   VERSION 1.26.2 (slot 26, 2026-11-09: SIGNED BY THE OWNER 2026-11-09 ("I sign."); the text of 1.26.2-pre.1, unchanged: header, date, release record and changelog only -- no statement, label, number, tag, status, marker (229 on 88 labels), written word or value moves; referee LAYOUT_signed_referee.md ACCEPT; the 1.26.2-pre.1 and earlier history below is kept)
%   VERSION 1.26.2-pre.1 (slot 26, 2026-11-09, PRE-RELEASE for owner re-signature, NOT signed): paths only, no change of content (repository layout, owner ruling 2026-11-09): the code moved under src/ (old -> new: tools/layout/mapping.tsv); 12 lines name the new paths (11 long-form file paths rewritten, one grep census sentence names src/). No statement, label, marker (229 on 88), number, tag or value moves.
%   VERSION 1.26.1 (slot 26, 2026-11-09: SIGNED BY THE OWNER 2026-11-09; the text of 1.26.1-pre.1, unchanged: header, date, release record and changelog only -- no statement, label, number, tag, status, marker (229 on 88 labels), written word or value moves; the 1.26.1-pre.1 and 1.26.0 history below is kept)
%   VERSION 1.26.1-pre.1 (slot 26, 2026-11-09, PRE-RELEASE for owner signature, NOT signed): no change of mathematical content: machine names and private paths removed for the public release (owner ruling 2026-11-09). Text only; no statement, label, marker (229 on 88), number, tag or value moves.
%
%   VERSION 1.26.0 (slot 26, 2026-11-09: SIGNED BY THE OWNER 2026-11-09; the text of 1.26.0-pre.11, unchanged: header, date, release record and changelog only -- no statement, label, number, tag, status, marker (229 on 88 labels), written word or value moves; item 8 (d')(vii): both values of G10.contract stay (1.26.0-pre.5, 1.26.0-pre.9); the 1.26.0-pre.1 ... pre.11 history below is kept)
%   VERSION 1.26.0-pre.11 (slot 26, 2026-11-09, the condition C1 and should-fix W-A, W-B of the confirm-referee docs/implementation/slot26/C126pre10_confirm_referee.md (ACCEPT-with-conditions)): a PRE-RELEASE, NOT owner-released;
%   TEXT ONLY; the status sites written before UP-Z, worded to hold before and after Z2's merge (thm:heckeZcert 'what is reachable', thm:descentcert(5), prop:descentpairhecke, def:heckeZ, sec:heckeZ 'Not here', prop:heckeZcong(5), 'What is not available at 2 and 3', prop:descentN's [R-Q4] pointer); item 8 (d')(ii): W-A (no row: 77 of 85, (11,0)), W-B; (d')(vii): pre.11 moves no written word, both values of G10.contract stay (1.26.0-pre.5, 1.26.0-pre.9).
%   No written word, golden entry, level-1 or anemic byte, hash input, label or marker (229 on 88) moves; no code change, no rerun.  Disposition docs/implementation/slot26/C126pre11_disposition.md.
%   VERSION 1.26.0-pre.10 (slot 26, 2026-11-08, the condition E1 of the referee docs/implementation/slot26/C126pre9_referee.md (ACCEPT-with-conditions)): a PRE-RELEASE, NOT owner-released;
%   TEXT ONLY; item 8 (d')(ii), The G3 words from Z2 on: E1 (the reason for G3.C_T.lower = not_run: d >= 3 eta_L of a pair undefined, d = 2 defined but no route performs the read, the non-reduced
%   branch by the Definition's scope; "prescribes a read of a lower bound"), W-1 (the words presuppose that the tags ran), N-1, N-2, N-3; (d')(vii): pre.10 moves no written word, both values of
%   G10.contract stay (1.26.0-pre.5, 1.26.0-pre.9).  No written word, golden entry, level-1 or anemic byte, hash input, label or marker (229 on 88) moves; no code change, no rerun.
%   Disposition docs/implementation/slot26/C126pre10_disposition.md.
%   VERSION 1.26.0-pre.9 (slot 26, 2026-11-08, Z2's G3 words, contract first): a PRE-RELEASE, NOT owner-released, NOT refereed; TEXT ONLY; item 8 (d')(ii): new paragraph "The G3 words
%   from Z2 on" (G3.C_T[U_p].status, G3.C_T.status, G3.C_T.lower final, not_run where no bound or no lower is read on the U_p_tags pairs (wording of pre.10); open(sandwich) at ell in {2,3} on an
%   open(cap) record; amendment A3: not_run in the closed list utag), the pre-Z2 rules kept for a binary without Z2's route; (d')(v) pointer; (d')(vii): G10.contract = 1.26.0-pre.9
%   for a binary that writes these words, 1.26.0-pre.5 for one that does not.  No golden entry, level-1 or anemic byte, label or marker (229 on 88) moves; no code change, no rerun.
%   Disposition docs/implementation/slot26/C126pre9_disposition.md.
%   VERSION 1.26.0-pre.8 (slot 26, 2026-11-08, the conditions of the confirm-referee docs/implementation/slot26/C126pre7_confirm_referee.md (ACCEPT-with-conditions)): a
%   PRE-RELEASE, NOT owner-released; TEXT ONLY; item 8 (c): D1 (two readings apart, not one: the derived maximal_ideals = not_run and the aggregate G3.C_T.status while a C_T read or
%   G3.C_T.lower it depends on is pending); (d')(vii): D2 (naming for the first time the word of a field (d') had left open moves a written word for a binary that writes it; the
%   manifest's assertion is of a record carrying that value).  G10.contract = 1.26.0-pre.5 stays correct.  No record word, hash input, golden entry, label or marker (229 on 88)
%   moves; no code change, no rerun.  Disposition docs/implementation/slot26/C126pre8_disposition.md.
%   VERSION 1.26.0-pre.7 (slot 26, 2026-11-08, the conditions of the confirm-referee docs/implementation/slot26/C126pre6_confirm_referee.md (ACCEPT-with-conditions)): a
%   PRE-RELEASE, NOT owner-released; TEXT ONLY; item 8 (d')(ii): C1 (not golden-eligible after Z2 alone only for a record carrying a field that waits on HZ-EISUP: 67 of Z1's 85,
%   not the 18), C3 (the derived not_run on any record, also 4 anemic export_word and 2 full N = 1 k1_route fields; pointers in (c), (d)), W-D (HZ-EISUP has its card in plan
%   v0.6); (d')(vii): C2 (<version> a decidable test; 'the code' = the binary that wrote the record; the manifest applies the test; G10.contract = 1.26.0-pre.5 stays correct); (c):
%   W-A (HZ-EISUP named as mover); (d')(ix): W-B ('exactly the status fields'; G0.operators and G1.S_cand.* named); N-a, N-b.  No record word, hash input, golden entry, label or
%   marker (229 on 88) moves; no code change, no rerun.  Disposition docs/implementation/slot26/C126pre7_disposition.md.
%   VERSION 1.26.0-pre.6 (slot 26, 2026-11-08, the conditions of the referee-light docs/implementation/slot26/C126pre5_referee_light.md (ACCEPT-with-conditions)): a
%   PRE-RELEASE, NOT owner-released; TEXT ONLY; item 8 (d')(ii): R1 (maximal_ideals = not_run beside a pending generation.status is a DERIVED word, rewritten when generation.status
%   moves; pointers in (c), (d), (v), (ix), the last covering edition 1's not_run -> blocked(U_p)), W1 (the Eisenstein U_p systems: no card of plan v0.5, working name HZ-EISUP, Z2 not
%   sufficient), W2 (G1.components = not_run at open(cap) stated), N1; (d')(v): R2 (G9.Up.* written by the package that serves U_p inside the record build, no Z1 record carries it,
%   a dump run does not create it); (d')(vii): R3 (<version> = the last version that moved a written word, a written value never rewritten, the manifest states both; the value
%   gl2_contract 1.26.0-pre.5 stays correct); N2, N3.  No record word, hash input, golden entry, label or marker (229 on 88) moves; no code change, no rerun.
%   Disposition docs/implementation/slot26/C126pre6_disposition.md.
%   VERSION 1.26.0-pre.5 (slot 26, 2026-11-08, the text aligned with the audited Z1 code, should-fix S1 (and S2's key) of docs/implementation/slot26/Z1_code_audit.md):
%   a PRE-RELEASE, NOT owner-released, referee-light-checked afterwards (docs/implementation/slot26/C126pre5_referee_light.md, ACCEPT-with-conditions, discharged at pre.6); item 8 (d')(ii): where the search ran and (C1) is certified, G1.components and
%   G2.ell=<l>.maximal_ideals (generation certified) stay blocked(U_p) (pending) until the Eisenstein U_p systems are built (Z2 or a named successor) -- replaces pre.4's
%   'no longer written' for G1.components; (d')(v): G9.Up.p=<p>.mat_sha256 unhashed, present only on full N > 1 records where U_p over Z was computed; (d')(vii): the key
%   is exactly G10.contract, value 'gl2_contract <version>, item 8 (d')'.  No final word introduced; only full_N words; no label added, renamed or removed; no marker
%   (229 on 88) moves; no record_fnv1a64, golden line or level-1 / anemic byte moves.  Disposition docs/implementation/slot26/C126pre5_disposition.md.
%   VERSION 1.26.0-pre.4 (slot 26, 2026-11-08, the should-fix items of the confirm-referee docs/implementation/slot26/C126pre3_confirm.md (ACCEPT)): a PRE-RELEASE, NOT
%   owner-released; item 8 (d')(ii): W3 (R1(b)'s fields enumerated: G3.C_T.lower = blocked(U_p); G3.C_T.status = upper_bound_certified where every U_p tag certifies a bound,
%   else blocked(U_p), sharpened to certified only together with .lower, by Z2; G3.C_T.exact absent), W4 (G1.components the fifth ell-independent field of R2), N1 (G3 on the
%   non-reduced branch: C_T.status = blocked(U_p), never upper_bound_certified; the final words of C_T.status / .lower there left explicitly open for Z2, which must name them);
%   item 8 (d), G1: W5 (the reading sentence of partial(unfactored): no ed.-1 record reaches that code path).  W6, N2-N4 are Z1 / K3 items (disposition).  No label added,
%   renamed or removed; no marker (229 on 88) moves; no value, record or hash moves (only full_N words; no record carries them yet).
%   VERSION 1.26.0-pre.3 (slot 26, 2026-11-08, the conditions of the adversarial re-referee docs/implementation/slot26/C126pre2_rereferee.md (ACCEPT-with-conditions)): a
%   PRE-RELEASE, NOT owner-released; item 8 (d')(ii): R1 (the non-reduced branch's G2 words written out -- open(sandwich) at ell >= 5, blocked(U_p) at ell in {2,3} until Z2;
%   the rule: G3.C_T.lower = blocked(U_p) on every full N > 1 record until Z2 lands or Z2-D's negative is recorded and audited, so no full N > 1 record is final before Z2),
%   R2 (the 'otherwise' branch's pending exit: where Z1's word-prime U_p route did not run the four ell-independent fields keep blocked(U_p) / blocked(export)), W1 (lead
%   decision: accurate words on the non-reduced branch -- G1.C1.status = refused(non_reduced) (final, new refusal reason), G1.theta = not_run, G1.C2_status = not_run --
%   amendment A2 after the freeze, hash-neutral; the reading of partial(unfactored) beside a non-certified C1), W2 (item 8 (d), C2_status as built).  No label added, renamed
%   or removed; no marker (229 on 88) moves; no value, record or hash moves (no record carries the new words yet).
%   VERSION 1.26.0-pre.2 (slot 26, 2026-11-08, the fixes of the adversarial referee docs/implementation/slot26/C126pre1_referee.md (ACCEPT-with-conditions)): a PRE-RELEASE,
%   NOT owner-released; conditions S1 (proof of prop:strictliftN(9)(c): descentpairhecke(1) reads nothing of N), S4 (item 8 (d')(ii): the non-reduced full cells --
%   no theta search, G1.C1.status = open(cap) final, theta-dependent fields named), should-fix S2 ('What is used'), S3 ((d')(iv): rmk:stabconvention and [U]'s orientation
%   remark cited and reconciled), E1-E9 editorial (E10: nothing to do).  No label added, renamed or removed; no marker (229 on 88) moves; no value, record or hash moves.
%   VERSION 1.26.0-pre.1 (slot 26, 2026-11-08, UP-Z contract-first text, condition C5 of docs/implementation/slot26/UPZ_design_audit.md (ACCEPT route (i) with C1-C7)): a PRE-RELEASE,
%   NOT owner-released, NOT yet refereed; on top of the signed 1.25.0.  prop:strictliftN gains item (9) (p | N: U_p on the free model through the certified level-1 lift and the
%   index-p transversal; items (2)-(6) verbatim, (1), (7), (8) not asserted) with its proof, the title / hypothesis / notation sentence amended; sec:recordschema item 8 gains
%   (d') (the full-N reading of the caps text, frozen with v1: operator list = G1.P's T_p then U_p, P_rule's (C1) = anemic (C1), theta_coeff_max = binom(d,2) with the full d,
%   the adjoint U_p* convention named; G1.theta_ops and G10.contract required on full N > 1 records).  No label added, renamed or removed; no marker (229 on 88) moves.
%   VERSION 1.25.0 (slot 25, 2026-11-07: SIGNED BY THE OWNER 2026-11-07; the text of 1.25.0-pre.8, unchanged: header, date and changelog only -- no statement, label, number, tag, status, marker (229 on 88 labels) or value moves; the 1.25.0-pre.1 ... pre.8 history below is kept)
%   VERSION 1.25.0-pre.8 (slot 25, 2026-11-07, LaTeX polish stage 5c, docs/hacs/LATEX_polish_plan.md; owner rulings Q1, Q3, Q5 of 2026-11-06): a PRE-RELEASE,
%   NOT owner-released.  The one binding change, rule (V1) amended (ruling Q1): the supersession record is the first subsection of Appendix E
%   (app:changelog).  Otherwise no content change: one-line date, uniform author block, generic abstract, the old date field and abstract moved
%   verbatim into Appendix E, two \part* headings (VIII, IX).  No label, number, statement, tag set or marker (229 on 88) moves.
%   VERSION 1.25.0-pre.7 (slot 25, 2026-11-06, the HZ-RAT-PERF fix round after its audit (docs/implementation/slot25/HZRATPERF_audit.md, PASS_WITH_MINORS, F2 MAJOR + F6): a
%   PRE-RELEASE, NOT owner-released; BINDING CHANGES NOT YET SEEN BY THE OWNER.  Item 8 only: (c) the pending word refused(word_prime_cap) in the closed list (the word-prime
%   list of the rational step or of a U_p tag would exceed its count cap 2^21; what a refused record contains); G1 C2_status's closed list as built (+ blocked(U_p),
%   refused(word_prime_cap)); (i) the cap in the unhashed word_primes= provenance key.  pre.7 incl. the re-audit's text fixes N2/N3
%   (docs/implementation/slot25/HZRATPERF_reaudit.md: "predicted (HZ-KIT)" need at (1,702); the envelope's non-field refusal text).  No label, tag set or marker (229 on 88) moves; no record, golden line or digest moves
%   VERSION 1.25.0-pre.6 (slot 25, 2026-11-06, the H1 merge after its audit (docs/implementation/slot25/H1_audit.md, PASS_WITH_MINORS): a PRE-RELEASE, NOT owner-released;
%   BINDING CHANGES NOT YET SEEN BY THE OWNER.  Item 8 only: (h) a G8 paragraph beside G5-G7 (heckez.G8 = hida only; the paths, closed lists, meanings, caps; G9 K9 and G10 hida);
%   (i) a pointer to it; (k) golden only under the default settings (F1).  No label, tag set or marker (229 on 88) moves; no default record, golden line or digest moves
%   VERSION 1.25.0-pre.5 (slot 25, 2026-11-06, the HZ-REC schema re-audit fix round (docs/implementation/slot25/HZREC_schema_reaudit.md, FAIL narrow: R1 MAJOR, R2 MEDIUM,
%   R3-R7 MINOR, R8-R12 NOTE; lead decisions, final for the round): a PRE-RELEASE, NOT owner-released; BINDING CHANGES NOT YET SEEN BY THE OWNER.  Item 8 only:
%   (c) [Lift] refused(no_lift) only for a certificate code (R6); (d) the Eisenstein multiplicity e_k (R1), the cusp block iff d >= 2 (R3), cusp mult not_run at N > 1 (R4),
%   the irreducibility cap 256 (R8); (e) the theta search over P u {T_ell} (R9), components_through / pair_kinds.kind not_run (R4); (f) the second route's word caps and
%   not_read (R2); (h) G7.pair, the theory tag (R4, R11); (i) the (C1) probe primes (R10); (j) heckez.second_route, the widened lines.  No label, tag set or marker (229 on 88) moves
%   VERSION 1.25.0-pre.4 (slot 25, 2026-11-06, the HZ-REC schema-audit fix round (docs/implementation/slot25/HZREC_schema_audit.md, FAIL as a short fix round:
%   2 MAJOR, 4 MEDIUM, 8 MINOR, 2 NOTE; lead decisions on F1-F3 and the Sec. 8 rulings): a PRE-RELEASE, NOT owner-released; BINDING CHANGES NOT YET SEEN BY THE OWNER
%   (he signs at the 1.25.0 release).  Item 8 only: (c) no pending word beyond the four, the export's refusals by name (F7, F8); (d) the components as written (F5,
%   F10); (e) the maximal-ideal rule, degree_attr not_run (F6, F12); (f) the partition list, agrees = not_read, eta_L.sign (F10, O5, O6); (g) torsion points, lower = 0
%   (F3, F12); (h) the new final value split_on_P, per-point G5 / G6, generation_tors and q = 2 not_run (F1, F3, F12, F13); (i) no G8 in default records, word primes
%   provenance (F2, F11); (j) the caps text, defaults never change (F2, F11, F16); (k) split_on_P.  No label, tag set, marker (229 on 88) or computed value moves)
%   VERSION 1.25.0-pre.3 (slot 25, 2026-11-06, the Z0 re-check (docs/implementation/slot25/Z0_recheck.md): a PRE-RELEASE, NOT owner-released; text only.  Item 8 (c):
%   export_word also covers K above EXPORT's K_max below the word bound (the precision EXPORT's ring loses to l-torsion), which pre.2 left without a record word;
%   open(precision) is the settings' step cap only.  No label, tag set, marker (229 on 88), value, record, golden line or digest moves)
%   VERSION 1.25.0-pre.2 (slot 25, 2026-11-06, the Z0 merge after its referee (docs/implementation/slot25/Z0_referee.md, REVISE: 3 MAJOR, 4 MINOR, 4 NOTE): a PRE-RELEASE,
%   NOT owner-released; text only.  F2: the counting criterion's binding sentence narrowed to N a power of p (a statement narrowed).  F3/O2: export_word is PENDING
%   (EXPORT-Z step 2 widens the ring), not final.  F4: the as-built refusal rule (word_modulus first; export_word iff l^K < 2^63 and K > EXPORT's K_max at an EXPORT-served read (pre.3); as built
%   at l >= 2^31 for every K), open(precision) added to the closed list.  O1 (lead): (C3) at K = 1, level 1, l in [2^31.5, 2^63) is the pending refused(k1_route).
%   F5: the (C2) / g_P cofactor route as built (gcd over the primitive theta's, residue-checked; upper_bound in (d) and (f)).  F6/O3: t_ell's closed vocabulary
%   (in_P, added, omitted(...)), omitted(export:<Name>), the 23170 cap in P_rule's default text.  F7: App. Part VII referee item.  No label, tag set, marker
%   (229 on 88), value, record, golden line or digest moves)
%   VERSION 1.25.0-pre.1 (slot 25, 2026-11-06, the Z0 merge: a PRE-RELEASE, NOT owner-released; on top of the signed 1.24.0.  The refereed HeckeZ text of the branch
%   s24-wip-heckez-z0 (1.25.0-pre.0c, A7 REVISE -> A7b ACCEPT) enters the main line by owner ruling Q8 (2026-11-04), as Z0 defines it: sec:heckeZ (23 statements, no number
%   moves), sec:recordschema item 8, the thm:descentcert(7)(a) pointer, 23 Appendix-B lines, 3 bib entries, the pre.0/0b/0c history; no LEAN marker (229 on 88).
%   Reconciled with 04-record-census.md (incl. its slot-25 amendment 1.5) and the wave-1/2 code (EXPORT, HZ-RAT/UTAG, HZ-LOC), each a dated note: item 8 (c) gains the
%   refusal reason export_word, (d) cofactor_test gains upper_bound, (j) the read set at l (P plus T_l for l not dividing N, l <= 23170; field t_ell), the factor budget as
%   built, the code map; sec:heckeZ: the word-modulus paragraph as built, the counting criterion's scope (N a power of p), the C2 frontier as built.  No label added,
%   renamed or removed beyond Z0's 23; no tag set moves; no value, record, golden line or digest moves.  Report: docs/implementation/slot25/Z0_merge_report.md)
%   VERSION 1.24.0 (slot 25, 2026-11-05: SIGNED BY THE OWNER 2026-11-05; the text of 1.24.0-rc.2, unchanged: header, date and changelog only -- no statement, label, number, tag, status, marker (229 on 88 labels) or value moves; the 1.24.0-pre.1, rc.1, rc.2 history below is kept)
%   VERSION 1.24.0-rc.2 (slot 25, 2026-11-05, RC124 rc.2: the RELEASE CANDIDATE of 1.24.0, NOT owner-released; rc.1 refereed ACCEPT_WITH_MINORS
%   (docs/implementation/slot25/RC124_referee.md), its F1-F4 and F6-F9 applied as text: rmk:bocksteinexp(4) the CCS grading (deg x = deg y = 2) and
%   the scope sentence; sec:status the marker rule amended by one sentence (beside PROVED or STD; on definitions / displayed equations, clause-scoped;
%   moves no marker and no tag), the rc.1 sentence shortened to a pointer; rmk:levelNcurvefamily the filtration dimensions dim m^i W = 255 - 2i, and
%   the two computational sentences reworded as attributed records not asserted here; prop:ltwothreescalar the Hatada attribution; two bib notes.
%   No label added, renamed or removed; no tag set or marker moves (229 on 88); no value, record, golden line or digest moves)
%   VERSION 1.24.0-rc.1 (slot 25, 2026-11-05, RC124: the RELEASE CANDIDATE of 1.24.0, NOT owner-released, NOT yet refereed (its referee is
%   docs/implementation/slot25/RC124_referee.md); = pre.1 (FMTMP item 7 (l), unchanged) + (a) the Q4 D5 Lean markers: 3 on thm:classical,
%   4 on eq:corconjres (slot24/LEAN_M6_report.md section 4 as audited; 222 on 86 -> 229 on 88 labels, 215 LEAN + 14 LEANlit) and the sec:status
%   marker list; (b) the phenomenology binding-text edits (docs/phenomenology/FINDINGS_SUMMARY.md section 3, catalogue section 11; NOT yet seen by
%   the owner): rmk:bocksteinexp (4) = Callegaro-Cohen-Salvetti Thm 4.5 (STD, VERIFYloc; the old (4) 'not asserted' narrowed to [L]'s estimate),
%   (3) a confirmation of it; prop:ltwothreescalar (F1) its origin (Calegari-Emerton Lemma 3.1, Hatada); Hi11 a route to a proof, tag unchanged;
%   rmk:levelNcurvefamily (21929,2) is CE's own example, the CE section-6 sentence on 2143 false; three bib entries (ccs13,
%   calegari_emerton_largeindex, rustom_primepowers).  No label added, renamed or removed; no tag set moves (rmk:bocksteinexp keeps PROVED, COMP^a, STD); no value, record, golden line or digest moves)
%   VERSION 1.25.0-pre.0c (slot 24, 2026-11-04, HeckeZ Z0: the A7b re-audit's MINOR (ACCEPT: 0 BLOCKER, 0 MAJOR, 1 MINOR, 5 NOTE; docs/implementation/slot24/heckez/A7b-Z0-reaudit.md): still a BRANCH PRE-RELEASE on s24-wip-heckez-z0, NOT owner-released; thm:heckeZcert(C2) states the drop rule through which the other theta_j enter (A7b-m1), the measured (C2) frontier re-measured on one D_{j0}; NOTEs N1, N4 applied; no label, number, tag set or marker moves; see the changelog);
%   VERSION 1.25.0-pre.0b (slot 24, 2026-11-04, HeckeZ Z0 fix round after the adversarial referee A7 (REVISE: 1 MAJOR, 7 MINOR, 10 NOTE; docs/implementation/slot24/heckez/A7-Z0-referee.md): still a BRANCH PRE-RELEASE on s24-wip-heckez-z0, NOT owner-released; prop:heckeZeiscert(3) gains the valuation check (A7-M1, counterexample (61,5,37) to the good-prime-only statement), the A7 MINORs and cheap NOTEs applied; no label, number, tag set or marker moves; see the changelog);
%   VERSION 1.25.0-pre.0 (slot 24, 2026-11-04, HeckeZ package Z0: a BRANCH PRE-RELEASE on s24-wip-heckez-z0, NOT owner-released, NOT refereed; it enters the main line as 1.25.0-pre.1 only after the owner has released 1.24.0 (ruling Q8 of 2026-11-04); on top of 1.24.0-pre.1.  NEW subsection sec:heckeZ at the end of sec:isomodelell (after sec:descentpair and sec:treepair), 23 statements appended, no number moves: def:heckeZ, prop:heckeZspec, prop:heckeZrat, lem:heckeZbridge, prop:heckeZmax, prop:heckeZgen, prop:heckeZetale, prop:heckeZmono, prop:heckeZlocal, prop:heckeZthrough, prop:heckeZuct, prop:heckeZsplit, prop:heckeZlevelone23, prop:heckeZglue, rmk:heckeZcbone, rmk:heckeZcbtwo, rmk:heckeZcbthree (CONJ), thm:heckeZdescent, thm:heckeZcert, rmk:heckeZcaveat, def:heckeZcong, prop:heckeZcong, prop:heckeZeiscert; sec:recordschema item 8 integral_hecke_structure/v1, (a)-(l); thm:descentcert(7)(a) points to prop:heckeZcong(5) for the read at a root; no marker moves (222 on 86); source docs/implementation/slot24/heckez/, signed 2026-11-04)
%   VERSION 1.24.0-pre.1 (slot 24, 2026-11-03, FMTMP: a pre-release, NOT owner-released, NOT refereed; on top of the signed 1.23.0; sec:recordschema item 7 gains the item (l) `The faithful coefficients' (every char_Z coefficient beyond 2^63 also as its exact value, canonical signed hex, the JSON list `big`, one hashed body line each), with pointers in (d) and (h); a record without such a coefficient keeps its bytes and record_fnv1a64; no golden line; text only, no label, number, tag set or marker moves)
%   VERSION 1.23.0 (slot 24, 2026-11-03: SIGNED BY THE OWNER 2026-11-03 ("I sign off 1.23.0"); the text of 1.23.0-pre.40, unchanged: header, date and changelog only -- no statement, label, number, tag, status, marker (222 on 86 labels) or value moves; the 1.23.0-pre.1 ... pre.40 history below is kept)
%   VERSION 1.23.0-pre.40 (slot 24, 2026-11-02, ICENSUS-P2: a pre-release, NOT owner-released, NOT refereed; sec:recordschema item 7 gains the item (k) `Parts not run' (spec key parts, hecke=not_run, id tail _partsG/_partsGH); text only; on top of pre.39 = ICENSUS-P1)
%   VERSION 1.23.0-pre.39 (slot 24, 2026-11-02, ICENSUS-P1: a pre-release, NOT owner-released, NOT refereed; sec:recordschema item 7 (e), (j): the ker/coker read's width ladder k_0 = min(K_max, 2 a_max + 3), then K_max(l) on a
%   precision shortfall (30 at 2, 19 at 3; every l); lifts the l = 2 [KerCokerPrecision] refusals at level 1; a failed safeguard refuses [NotEndomorphism]; moves no hash of an ok entry and no golden line; text only)
%   VERSION 1.23.0-pre.38 (hecke-extensions slot 24, 2026-10-28: a pre-release, NOT owner-released, NOT re-refereed; pre.37 refereed PASS-with-fixes (CON P7 (a),
%   docs/implementation/slot24/CON_pre37_referee.md, 2 M / 10 m / 5 N / 2 T); the fixes, text only (docs/implementation/slot24/CON_pre38_fixes.md).  M1 item 7 (d): the record's
%   coefficients run from the leading 1 down (c_i = coefficient of x^(r-i)), an exception to item 2, stated.  M2: item 7's code map, (f), (j) and the pointers restated as
%   built, audited (D6_code_audit.md), fixed (19a969de), merged (2aa61397), golden R1 landed (6958a5d2); console bound sum p <= 100, <= 6 primes <= 97; arm (B).  m1 (a) the
%   two refusal mechanisms (CoefficientRingK at construction, GateVerbsOnField at registration) and the residual risk; m2 (b) record level vs pair.bulk; m3 (e) indexed by
%   torsion_primes; m4 (e) the ker/coker conditions; m5 (g) groups' refusal is the record's, refusal text on non-refused entries; m6 (h) exhaustive partition, body() normative;
%   m7 (d) the digest's byte layout; m8 (h) the (1, 12) hashes illustrative, not pinned; m9 sec:isomodeltree row + sentence; m10 def:ringsemantics(3) at level N, (4) built;
%   N1-N4, T1, T2.  No label, number, tag set or marker moves (222 on 86).
%   (previous:) VERSION 1.23.0-pre.37 (hecke-extensions slot 24, 2026-10-26: a pre-release, NOT owner-released, NOT refereed; on top of pre.36.  The integral descent record gets its
%   contract home: sec:recordschema item 7, integral_descent_record/v1 (D6, branch s24-d6impl tip 538cf881, pending audit and merge; golden event staged, pending the
%   owner's go): ring Z and object group; q_max = 3 derived from thm:descentcert(6), truncated, periodic_from = 2; groups, bad primes, K per prime (provenance);
%   char_Z digests (FNV over c mod 2^64 and mod 3^40, residues when undetermined); torsion Hecke invariants; the pair, each complex q_max = 4 (compact's list of
%   names missing in v1: code item); refusals (a refused record is never an empty answer); hashed / unhashed split justified by thm:descentcert's choice-independence,
%   rmk:descentcertscope and the read-only rollbacks HECKE_DESCENT_TREE(_PAIR)=0.  (1, 200), l = 2: [NotEndomorphism] (30 < 34) is a precision shortfall of the
%   ker/coker read at 2^30 (30 < 34 counts pivots); the code name is to be [KerCokerPrecision] (code item).  thm:descentcert(4), (5) not recorded.
%   def:ringsemantics(3): Z un-reserved for the descent only; every other reader refuses it by name; Z/<m> stays reserved.  No label, number, tag set or marker
%   moves (222 on 86).
%   (previous:) VERSION 1.23.0-pre.36 (hecke-extensions slot 24, 2026-10-26: a pre-release, NOT owner-released, NOT re-refereed; pre.35 refereed ACCEPT after fixes (C-D8b,
%   docs/implementation/slot24/C-D8b_referee.md, 0 M / 1 m / 7 N); the fixes, and D5c MERGED at 4064e164.  m1 code map: the certificates restated for the merged
%   code -- certify_transport checks the blocks of the APPLIED operator (T'_T, eta_T, T'_d, strict block) against separately assembled ones and the eta identity on the
%   same values, entrywise for D <= 1300, above on the generators + 64 random vectors uniform in Z/l^K seeded by (l, N, n, j), miss <= l^-64 per check; what is
%   deterministic, what probabilistic, why admissible (no value read from a check; the read's proof is deterministic); "not new coverage" gone.  N1 (4)(e)(iv) "gives",
%   not "is", the coboundary half; above D = 1300 no deterministic cochain-level check of it.  N2 (4)(e) admits a truncation of T (the pair lift to degree 6).  N3 proof
%   of (e)(i) for projective sources.  N4 (e)(ii) refined (bulk in j+1, boundary and Y in j) and extended to H^q, H^q(d).  N5 the identity on the generators' alpha-parts
%   is the second component of (ii) modulo l^(K - sigma_max).  N6 def:treepair(2) "computed from these cochains".  N7 obligation (c) every 2- and 3-adic pair piece,
%   every N, incl. 3 in extra_primes with n odd; eta_skipped + arm (Q5); (g) met by transport_bytes.  Code map: merged, audit f7af1744 (m1-m4, N1, N2, N4, N6 fixed;
%   N5 deferred), reach unchanged; "pending" remarks updated (sec:isomodeltree, sec:descenthecke + D5c's additions to D3's code).  No label, number, tag set or
%   marker moves (222 on 86).
%   (previous:) VERSION 1.23.0-pre.35 (hecke-extensions slot 24, 2026-10-26: a pre-release, NOT owner-released, NOT refereed; sec:treepair brought in line with the D5c code
%   (branch s24-d5cimpl, tip e720f11c, pending audit).  def:treepair(2): S vs S^{-1} at l in {0,2} is immaterial on cochains where c acts trivially -- rho_{M_N}(c) = (-1)^n,
%   so the same cochains for n even at every l; the code map's control restated (chain level, or n odd mod 2^K, K >= 2; asserted clean for n even).  NEW prop:treepair(4)(e)
%   (truncation): with the comparisons in chain degrees <= m (homotopies <= m-1) the read of H^q_c, q <= m, is that of every extension; the identity of (4)(a) in cone degree j
%   needs degree j+2, formable iff j <= m-2, not used by the read; m = 5, so q_max = 5 STAYS (refusal q_max not in {4,5} kept).  Code map: built, default 2-/3-adic pair route,
%   rollbacks HECKE_DESCENT_TREE_PAIR=0 / HECKE_DESCENT_TREE=0, files, validations, reach (D ~ 1060-1240 -> refusal ~ 6340, ~ 8000 at 40 GiB, next wall D2's dense Smith
%   with transforms); certificates (i)-(iii) restated, "on the generators" = deterministic, partial, not probabilistic, admissible (no value read from it); N7 closed
%   (D3Refusal::HeldOutMismatch removed); eps^(l) = the code's.  cd8_check.py extended (arms M3, Q).  No label, number, tag set or marker moves (222 on 86).
%   (previous:) VERSION 1.23.0-pre.34 (hecke-extensions slot 24, 2026-10-26: a pre-release, NOT owner-released; the C-D8 referee's fixes m1-m4, N1-N8,
%   docs/implementation/slot24/C-D8_referee.md (pre.33 refereed ACCEPT after fixes, 0 M / 4 m / 8 N); NOT re-refereed.  sec:treepair (and one code-map note in sec:descenthecke).
%   m1 prop:treepair(4)(a) / rmk:treepairscope(iii): on the pair model (Cyl, B) the transported datum is strict (eta = 0, the plan's sentence holds there); on Cone(a^*)
%   T^c ~ varpi T^Cyl varsigma' acquires eta_T; "eta != 0 / diagonal not a cochain map" are OBSERVED in the checks (theta^*, the referee's T_eps), not proved; no "forced",
%   no "no strict pair datum".  m2 (3)(b): h_p, (L), (L') for any column augmentation of a complex filtered by finitely many free periodic columns, over every R after
%   (x) R; eps_B and cmpI|Isob listed; the wrong "stated for any such p" gone.  m3 code map (f) / check paragraph: S for S^{-1} only at l in {0,2}; at l = 3 drop the
%   S^{-1}e_U v^U_1 term; eta dropped only for non-fibre operators, not at l = 3 with n odd, detection count printed.  m4 def:treepair(4): eps^(l)_T, eps^(l)_B defined;
%   (1)(b) with them; the four comparisons lie over the identity.  N1 lem:treepaircyl f.g. projective, zero in negative degrees; N2 path orientation; N3 (3)(b)
%   parenthesis and scalars 1; N4 Cone(rres_Cyl) not admissible, rank C^0_{T,l}; N5 eta's cost in (e), (g); N6 rollback preserves the record's invariants; N7 open
%   code item (D3Refusal::HeldOutMismatch vestigial) in sec:descenthecke's code map; N8 S vs S^{-1} immaterial at l >= 3.  No label, number, tag set or marker moves (222 on 86).
%   (previous:) VERSION 1.23.0-pre.33 (hecke-extensions slot 24, 2026-10-26: a pre-release draft, NOT owner-released, NOT refereed; plan milestone D5c, clause C-D8:
%   NEW sec:treepair (after sec:descentpair): def:treepair (the cusp model B^(l)_T = cone of r_{(1-T)e} on the edge column, the axis map a^(l): B -> Iso^(l)_T
%   in closed form, the tree cone Cone(a^*), the geometric tables on the line complex C(L)), lem:treepaircyl (the cylinder is a pair model, its cone deformed
%   onto Cone(a^*) strictly compatibly with coneinc / coneproj, signs of prop:descentpair(1)(c)), prop:treepair (what it computes, closed forms, effective
%   comparisons Phi^Z_d, Psi^Z_d by (L), square homotopies by (L'), cone maps j_Phi, j_Psi mutually inverse on H through comparisons of pairs, T_p by
%   transport with a non-zero determined off-diagonal eta_T, precision rule, refusal / failure / scope), rmk:treepairscope (the plan's section 5 corrected).
%   sec:descenthecke code map: D3 merged (dc8261a3); l >= 5 cokernel read refuses RssCeiling before building; HeldOutMismatch and ResidueRank throw
%   HeckeInternalError (failures).  Check docs/implementation/slot24/cd8_check.py.  Four statements added (407 -> 411); no marker moves (222 on 86).
%   (previous:) VERSION 1.23.0-pre.32 (hecke-extensions slot 24, 2026-10-26: a pre-release, NOT owner-released; the C-D3h referee's fixes m1-m3, N1-N6,
%   docs/implementation/slot24/C-D3h_referee.md (pre.31 refereed ACCEPT after fixes, 0 M / 3 m / 6 N); NOT re-refereed.  sec:descenthecke only.
%   m1 prop:strictliftN(2) independence of the lift: d_I unsigned (lem:cosettransport(2)), H = Theta_{h_Iso} in every degree, no re-signing.
%   m2 thm:descentcert(4) at ell >= 5: K > sigma_max(Delta) + v_ell(chi(c)) suffices (k = K'), with proof; sigma_max = 0 when no ell-torsion.
%   m3 prop:torsionhecke(3'): any exact lifts x_j = x~_j mod ell^K form a basis (P = I + ell^K M), proved.  N1 presentation independence per
%   certifying read; N2 U = ST; N3 K' sharp for the free part; N4 'constructible' H (with m1); N5 code map: certify() at q <= 2 suffices,
%   Rel_N-stability a theorem; N6 'What is used' completed.  No label, number, tag set or marker moves (222 on 86).
%   (previous:) VERSION 1.23.0-pre.31 (hecke-extensions slot 24, 2026-10-26: a pre-release, NOT owner-released, NOT refereed; sec:descenthecke only: three clauses
%   closing D3's contract notes (DESCENT3_report.md section 6, branch s24-d3impl).  (1) prop:torsionhecke(3') the read on a presentation coker(Delta) with
%   F(im Delta) in im Delta, the free part with k = K', presentation-independence; prop:crtcharpolyN(2) and thm:descentcert(3), (4) admit [H_{N,p}] on
%   M_N/Rel_N at ell >= 5 / at the residue primes.  (2) prop:strictliftN(2): obligation (b) met through the lift (tau_Iso's certificate with the strict identity
%   on every basis element makes the table a datum over Z).  (3) prop:strictliftN(2): every lift with (T1)-(T3) and the strict identity is admissible (the
%   pair lift's bulk component is the one read); independence of the lift (an explicit cochain homotopy).  Code map: D3 built on s24-d3impl, pending merge
%   and audit; thm:descentcert(4) at ell in {2,3} not built.  No label moves (the new sub-item (3') is unlabelled); no number, tag set or marker moves (222 on 86).
%   (previous:) VERSION 1.23.0-pre.30 (hecke-extensions slot 24, 2026-10-26: a pre-release, NOT owner-released; NON-MATHEMATICAL code-map update, no statement changes:
%   sec:isomodeltree code map and its table row record D5 BUILT and MERGED (5eb7e9de: tree = default route for the dense ell-local bulk pieces, rollback
%   HECKE_DESCENT_TREE=0; TreeCochains, IsotropyModelTree, DescentTree, TreeComparisonArms; HeckeIsotropyModelTree(Slow)Validation; Hecke by transport = hook,
%   real arm pending D3; tree cusp pair D5c not built).  Code map (f) of sec:descentpair: hecke::HeckeInternalError (bc128cb7) -- internal errors propagate,
%   only a real overflow refuses as IntegerCap; the pre.29 "known width" is resolved.  No label, number, tag set or marker moves (222 on 86).
%   (previous:) VERSION 1.23.0-pre.29 (hecke-extensions slot 24, 2026-10-25: a pre-release, NOT owner-released; the C-D7p referee's fixes M1, m1-m4, N1,
%   docs/implementation/slot24/C-D7p_referee.md (pre.28 refereed ACCEPT-with-fixes); NOT re-refereed).  sec:descentpair only; no label / number / marker moves.
%   M1 prop:descentpairhecke(1), (5), code map (f): "refused exactly when" = the code's eleven PairRefusal codes (NotPrime, PDividesN, BulkOnly,
%   DegreeRange, RssCeiling, IntegerCap, DomainCap, InputTableCertificate, LiftCertificate, IsobClosure, StrictIdentity); the cap = int64 products,
%   int32 tau coefficients, the uint32 coset domain, a failed input-table certification (lift construction AND operator assembly); RSS ceiling 4 GiB
%   before building; every other exception a failure, not a refusal; IntegerCap's known width (HeckeOverflow from internal guards) recorded.
%   m1 proof of (1): the forced component map (i)-(iii) from (T1)-(T3), Isob-preservation and def:descentpair's augmentation; the eps_I route kept as
%   a second proof.  m2 q_max <= 5, q_top > 6 refused, higher degrees by (7).  m3 overflow in the assembly refuses.  m4 code map (i): the homotopy-fault
%   controls (StrictLift::with_homotopy).  N1 sizes model-dependent (x3 per two degrees on another model).  Code map: D4 MERGED at 3991a2f1.
%   222 markers on 86 labels; app:status counts stand.
%
%   VERSION 1.23.0-pre.28 (hecke-extensions slot 24, 2026-10-25: a pre-release, NOT owner-released, NOT refereed; the owner's scope addition of
%   2026-10-25).  sec:descentpair: prop:descentpairhecke(1) the pair lift built ON DEMAND for every prime p (degrees <= 6), certified at run time by
%   rmk:strictliftbuilt(2) incl. the boundary items, cached; refused on a failed certificate or the cap; the certificate suffices on its own (component
%   map forced); proof: boundary step spelled out, no step verified only at p in {2,3,5}.  (5) certified at every l, N for every prime p !| N whose lift
%   passed its certificate (run time); refused only for p not prime, p | N, failed certificate, cap.  rmk:descentpairscope(iii); code map BUILT on
%   s24-d4impl (d82c7761, not merged), (f) restated; rmk:strictliftbuilt, rmk:integralscope(2) historical notes added.  222 markers on 86 labels.
%
%   VERSION 1.23.0-pre.27 (hecke-extensions slot 24, 2026-10-25: a pre-release, NOT owner-released; the C-D7 referee's fixes M1-M3, m1-m5,
%   docs/implementation/slot24/C-D7_referee.md, ACCEPT-with-fixes; NOT re-refereed).  sec:descentpair only; no label / number / marker moves.
%   M1 the check paragraph: H^2_c does have l >= 5 torsion (Z/4+Z/60 at (1,8), Z/2+Z/12+Z/1260 at (1,12)); the false sentence and its reading replaced.
%   M2 prop:descentpair new (5)(e): H^2_c (x) Z_(l) = H_0(Gamma; M) (x) Z_(l) for l >= 5 (norms + Manin resolution + generation by S, U), for M_N = V_A,
%   zero for l > n+1 (Borel coinvariants, Nakayama); (6)(c) D^0 by the Betti recursion OR a prime l' >= 5, l' > n+1 (as prop:descentN(7)); a-priori
%   cap on D^0 at l >= 5; rmk:descentpairscope(ii), code map (c).  M3 prop:descentpairhecke(5) WEAKENED: boundary / cone T_p certified exactly at the p
%   with a built pair lift (p in {2,3,5}), and at l >= 5, level 1, where the shipped cone operator is built; level N at other p (l >= 5 included):
%   exists, determined, not computed, refused (code map (f): not a D4 obligation).  m1 locators: VII.6 withdrawn, (4) periodicity by the tensor
%   resolution, (5)(d) by the coefficient sequence (n1) with the Coind Shapiro-cup compatibility; dangling lem:bulkperiodic(4) removed; III.8-9.
%   m2 the cone unbounded at 3 exactly when H^{q>=2}(A;V)_(3) != 0.  m3 order n! for n even, 2^{floor(n/2)+1} for n odd.  m4 (1)(c) signs pinned:
%   coneinc_* = j^*, coneproj_* = -delta (proof by a graded splitting and an explicit resolution E).  m5 exponent <= 12 observed.  cd7_pair_check.py
%   arm R (308 checks): 2528 checks, 0 failing; the referee's check rerun, 2457, 0 failing.  222 markers on 86 labels; app:status counts stand.
%
%   VERSION 1.23.0-pre.26 (hecke-extensions slot 24, 2026-10-25: a pre-release draft, NOT owner-released, NOT refereed; plan milestone D4, clause C-D7
%   of DESCENT_plan.md v0.4).  NEW subsection sec:descentpair (after sec:descenthecke, before sec:mpmodel, so no statement number moves):
%   def:descentpair (pair models (P, P^d): Iso/Isob over Z, Iso^(l)/Isob^(l) over O_l, the cellular pair over Z[1/6]; comparisons and homotopies of pairs),
%   lem:pairlifting (PROVED, STD Weibel 2.2.6: comparisons of pairs exist boundary first and are unique up to homotopy of pairs; the cone maps and the
%   cone homotopy), prop:descentpair (PROVED, STD Brown III.2, III.5-6, VI.9, VII.6, Weibel: what a pair model computes over every ring; H^q_c =
%   Ext^{q-1}(St, M), H^1_c torsion-free; Isob^(2) free, Isob^(l) for l >= 3 the circle model (Bredon boundary on Iso^(3)); the glued LES over Z with the
%   pair map cmpI the only glueing datum, exactness local; level N by Shapiro, block-diagonal by cusps; the boundary in closed form from the cusp widths,
%   |tors (Sym^n)_{T^w}| = n! w^n; rational ranks h^1_c = b_1, h^2_c = b_0(M^v); 48 H^q_c = 0 (q >= 3), the l-local shapes; periodicity, q_max >= 4 for the
%   pair; the certified exponent on boundary and cone with the admissible rank sources and caps -- D-delta6), prop:descentpairhecke (PROVED, STD Brown,
%   Weibel: the pair lift for every p produced, built for p in {2,3,5}; the level-N operator on Isob, strict with rres, eta-hat = 0; what it induces on
%   H_c = Ext(St, V) and on the boundary, choice-free over Z; the LES T_p-equivariant over Z, Z_(2), Z_(3); the pieces; what is certified at which p;
%   D2's maps on the pair; periodicity Hecke-equivariant), rmk:descentpairscope (PROVED: where the plan's D4 / C-D7 wording is corrected).
%   rmk:integralscope (4), (5) and thm:descentcert(7)(c) gain pointers.  Check outside the build: docs/implementation/slot24/cd7_pair_check.py, 24 cells,
%   2220 checks, 0 failing.  app:status 402 -> 407 (def 67 -> 68, lem 61 -> 62, prop 98 -> 100, rmk 113 -> 114, PROVED 236 -> 240, STD 102 -> 105,
%   tagged 271 -> 275, untagged 131 -> 132).  222 markers on 86 labels (none added).
%
%   VERSION 1.23.0-pre.25 (hecke-extensions slot 24, 2026-10-25: a pre-release, NOT owner-released; the C-D4 referee's fixes M1, M2, m1-m5, n1,
%   t1-t3, docs/implementation/slot24/C-D4_referee.md, ACCEPT-with-fixes; NOT re-refereed).  sec:descenthecke, and one sentence each in
%   prop:descentmaps(2)(g)(ii) and rmk:bredonhecke; no number / label / marker moves.  M1 prop:torsionhecke new (1)(e): the layer-j block's char
%   poly is choice-free mod l^{e_j}, e_j = min(j, d_j), d_j the distance to the nearest other layer (proof by paths through the layers); (2) "only
%   mod l" replaced, with sharpness pairs for every 1 <= d < j; rmk:descentcertscope(iii) amended.  M2 resolution (a): the l = 3 piece of
%   thm:descentcert(3) is the transport T'^(3)_p (or T'_T,p) for every N; prop:strictliftN(6)(b) holds for an integral datum certified over Z only
%   and is read by no statement; the shipped Bredon operator is never read at l in {2,3}; code obligation (a) in the code map.  m1 (3): 2 sigma_max
%   <= K necessary for the module (minimal example), P_j and Fitting types at every certified K; not the conservative BuildPres policy of
%   (2)(g)(ii).  m2 Garner needs l' < 2^32, Horner mod Q > 2^32 a word mulmod, the negative case.  m3 tau_X any datum over Z (canonical tau once
%   certified; no boundary datum then); uniform in p.  m4 tau_Iso built once, independent of r_t and N.  m5 (7)(a) "not asserted" in place of
%   "open mathematics".  n1 l' = p admissible.  t1 sharpness iff r_N > 0.  t2 prop:symconv(2).  t3 levelconehecke(4) with tau_Iso.  Referee's
%   check extended by an arm M (sharpness, minimal example, block polys on real reads).  222 markers on 86 labels; app:status counts stand.
%
%   VERSION 1.23.0-pre.24 (hecke-extensions slot 24, 2026-10-25: a pre-release draft, NOT owner-released, NOT refereed; plan milestone D3, clauses
%   C-D4 and C-D5 of DESCENT_plan.md v0.4).  NEW subsection sec:descenthecke (after sec:isomodeltree, before sec:mpmodel, so no statement number
%   moves): prop:strictliftN (PROVED, STD Brown III.8-10, Weibel 2.2.6: AgA = A diag(p,1) A for p not dividing N; a bulk strict lift tau_Iso for
%   every p, constructed by (L) without a boundary step; the level-N operator eq:levelcorr on Iso over every ring; cmpI^* intertwines it on the
%   nose; transports to Iso^(l) and the tree (the route at 2); cheap routes at l >= 5 and at 3 with nu_3(N) = 0; strict square with alpha and the
%   level LES equivariant over Z, Z_(2), Z_(3); commutativity), prop:crtcharpolyN (PROVED, STD Diamond-Im, Deligne, Diamond-Shurman: char_Z of
%   T_p on H^1/tors at level N from O-piece residues only, the bound 1 + p^{n+1} as the input, recovery in 64-bit words by Garner's digits, the
%   output rule: integer / digits + bound / residues + bound, never a lifted integer), prop:torsionhecke (PROVED: choice-free invariants of an
%   endomorphism of a finite l-group -- layer polynomials over F_l, Fitting components, ker/coker types; per-Smith-summand char polys are NOT
%   choice-free; the exact read from a local run at K >= 2 sigma_max; lattice invariant factors of T - c at non-roots), thm:descentcert (PROVED,
%   C-D5 as restricted and corrected), rmk:descentcertscope (PROVED: where the plan's C-D4/C-D5 wording is corrected).  rmk:integralscope (2), (4)
%   and prop:descentmaps(3)(c) gain pointers.  Bibliography: deligne_weil1.  Check outside the build:
%   docs/implementation/slot24/cd4_hecke_check.py, 21 cells, 931 checks, 0 failing.  app:status 397 -> 402 (prop 95 -> 98, thm 19 -> 20, rmk
%   112 -> 113, PROVED 231 -> 236, STD 100 -> 102, tagged 266 -> 271).  222 markers on 86 labels (none added).
%
%   VERSION 1.23.0-pre.23 (hecke-extensions slot 24, 2026-10-25: a pre-release, NOT owner-released; the C-D3g referee's fixes M1, m1-m5, t1, t2,
%   docs/implementation/slot24/C-D3g_referee.md, ACCEPT-with-fixes; NOT re-refereed).  prop:descentmaps and, editorially, prop:isomodeltree; no number /
%   label / marker moves: M1 (2)(g)(i) "moreover": the first s rows of L, the last r-s columns of L^-1 and the last c-s rows of R^-1 are exact mod
%   ell^K (proof from the Schur-complement factorisations); (ii) lattice coordinates D^-1 (L y)_{<=s} mod ell^{K - sigma_i}, (L y)_{>s} = 0 mod ell^K
%   (LatticeQuotient); m1 torsion coordinates exact for every sigma_i < K, 2 sigma_max <= K a conservative policy (BuildPres), not a necessity;
%   m2 (iii) exact division, torsion-row reduction, mixed precisions; m3 (iii) the certifying ranks of derived forms, with the rational reason;
%   m4 pivots u_i ell^{sigma_i}, u_i units; m5 (2)(b) delta's convention, equivariance unaffected by the sign; t1 isomodeltree(6)(b) Phi_l / Psi_T
%   references; t2 isomodeltree(4)(b) all five decompose() callers.  222 markers on 86 labels.
%
%   VERSION 1.23.0-pre.22 (hecke-extensions slot 24, 2026-10-25: a pre-release, NOT owner-released, refereed ACCEPT-with-fixes (C-D3g_referee.md); the D2 code audit's F1 and F8,
%   docs/implementation/slot24/D2_code_audit.md).  prop:descentmaps only, no number / label / marker moves: NEW item (2)(g), the precision rule
%   (generators, kernel bases, coordinates and induced maps read from a CERTIFIED Smith form L A R = Sigma over Z/ell^K are exact modulo
%   ell^{K - sigma_max}; composition; refusal), with proof (the audit's Appendix B, made integral over Z_(ell)); (2)(b) the signs of the cone's two
%   maps ((c, c') |-> c induces -delta); the code map records D2 as built.  prop:isomodeltree, editorial (the C-D6 re-referee of pre.21, ACCEPT,
%   C-D6_referee.md): N1 the representative convention is PeriodicColumn.h decompose() (signed lex order, j rel. the fixed generator); N2 (L), (L')
%   finite per degree, not periodic (coefficients grow linearly), built to the degree in use; N3 p' needs no surjectivity; N4 (6)(b) names
%   theta = Phi_l o Psi_T.  222 markers on 86 labels.
%
%   VERSION 1.23.0-pre.21 (hecke-extensions slot 24, 2026-10-25: a pre-release, NOT owner-released; the C-D6 referee's fixes M1, m1-m5,
%   docs/implementation/slot24/C-D6_referee.md, ACCEPT-with-fixes; NOT re-refereed).  prop:isomodeltree only, no number / label / marker moves:
%   M1 (4) made EFFECTIVE: (4)(a) the tree as a Gamma-subcomplex of Xbulk (iota), an explicit retraction r (19 values) and homotopy G with
%   id - iota r = dG + Gd; (4)(b) the column contractions with fixed coset representatives, the solvers h_p for ker cmpI and ker eps_T, the
%   lifting formula (L) (lem:strictlifting made effective) and its homotopy form (L'); (4)(c) Phi^Z := (L) of r cmpI, Psi^Z := (L) of iota eps_T,
%   explicit homotopies H_T, H_X, H_l; (4)(d) Phi_T, Psi_T with explicit homotopies K, K'; (5) and the code map say so.  m1 isomodelell(4), (5)
%   hold over every O_l-algebra (stated in (4)(d), (6)(b)).  m2 "strict deformation retract" -> a retract, Psi*Phi* = id, Phi*Psi* ~ id by H_l.
%   m3 (1)(d)(iv) e_x central.  m4 (5) other maps enter through Phi_T, Psi_T.  m5 (6) stated for every cochain operator; the sign nit in
%   (1)(b)'s proof.  Companion check docs/implementation/slot24/cd6_iso_check.py (builds Iso): 12870 checks, 0 mismatches.  222 markers on 86 labels.
%
%   VERSION 1.23.0-pre.20 (hecke-extensions slot 24, 2026-10-25: a pre-release draft, NOT owner-released, NOT refereed; plan milestone D5, clause C-D6,
%   brought forward by the owner's ruling at gate G-T of 2026-10-25, DESCENT_plan.md v0.4).  NEW subsection sec:isomodeltree (after the last
%   statement of sec:descentN, before sec:mpmodel, so no statement number moves): def:isomodeltree (the free and the ell-local tree models: cones of
%   lifts of the tree boundary between ell-local columns; canonical lifts at ell = 2, 3, >= 5) and prop:isomodeltree (PROVED, STD Weibel 2.2.6,
%   Brown III.6.2): (1) projective resolutions computing H^*(Gamma; M) in every degree for every M (Ind modules, level N), one integral form over
%   O_ell, prop:certifiedeell(2)-(6) on it; (2) cochains in closed form, periodicity; (3) closed-form comparisons with the free tree, Phi o Psi = id
%   strictly (strict deformation retract); (4) comparisons Phi_T, Psi_T with Iso over Z (lem:strictlifting), not over C(Xbulk); (5) (T1)-(T3) of
%   prop:descentmaps(5) hold, so fibre maps and the level cone hold on the tree; (6) Hecke by transport T'_T = Psi_T^* T Phi_T^* (p not dividing N;
%   implemented with D3).  Bulk only (C-D8 not stated).  Check outside the build: docs/implementation/slot24/cd6_tree_check.py, 145 cells, 18 pairs,
%   49303 checks, 0 mismatches.  app:status 395 -> 397 (def 66 -> 67, prop 94 -> 95, PROVED 230 -> 231, STD 99 -> 100, tagged 265 -> 266);
%   labelindex three labels; codemap row.  222 markers on 86 labels (none added).
%
%   VERSION 1.23.0-pre.19 (hecke-extensions slot 24, 2026-10-25: a pre-release, NOT owner-released; the C-D3 referee's minors m1-m5,
%   docs/implementation/slot24/C-D3_referee.md, ACCEPT-with-fixes; NOT re-refereed).  prop:descentmaps only, no number / label / marker moves:
%   m1 (3)(a)/(b): [AgA], g = diag(1,p), is this document's U_p = the adjoint U_p^* on p-old spaces (levelconehecke(3)(d)); the res-square failure
%   is that convention's (the classical U_p commutes at chi = 1); NEW item (3)(b'): p | N => cor o [BgB] = [AgA] o cor on H^q, every q, every M,
%   every ring, over Z (Mackey with A' = A cap g^-1 A g, B cap A' = B', A = A'B via T^j in B): the Hida tower's cor is U_p-equivariant; fails for
%   p not dividing N, p | M.  m2 (3)(a) cites levelconehecke(1) alone.  m3 (3)(c) needs p not dividing M.  m4 (2)(d)(i) reduce the integral
%   (O_ell-) form mod ell'.  m5 (4)(c) wording: H^q(Gamma_0(p^r)) = H^{q-1}(G; H^1(Gamma_1(p^r))), q >= 2, all 2-, 3-primary.  Referee's script on
%   10 more cells (cor square: 8 commute with p | N, 2 fail with p not dividing N).  app:status note (no count moves).  222 markers on 86 labels.
%
%   VERSION 1.23.0-pre.18 (hecke-extensions slot 24, 2026-10-25: a pre-release, NOT owner-released, refereed ACCEPT-with-fixes (C-D3_referee.md); plan milestone D2, clause C-D3).
%   NEW prop:descentmaps (PROVED, STD Brown III.6.2, III.8-10, VII.6, Weibel 2.2.6), last statement of sec:descentN, so no statement number moves:
%   (1) fibre maps (res, cor, R_sigma, diamonds) commute strictly with every comparison (cmpI^*, Phi_ell^*, Psi_ell^*, the tree models of C-D6) and
%   glue to maps over Z; Shapiro gives res, cor, conjugation in every degree over every ring; (2) the level cone on every model: exact, computes
%   H^{*-1}(A; Q_{A/B}), the long exact sequence over Z in every degree (not six-term at 2, 3), glued by prop:cohglue; certified K on the cone with
%   the Betti recursion beta_j = b_j(B) - b_j(A); bad primes of the cone; tree / sparse routes at ell >= 5 and at ell = 3 when nu_3(N) = 0;
%   (3) res and cor commute with [g], g = diag(1,p), p not dividing M, every degree (T_p, p not dividing MN; U_p, p | N); level raising fails;
%   the relative term at 2, 3 deferred to C-D4; (4) normal case: five-term sequence over Z; torsion-free B (Gamma_1(N) < Gamma_0(N), the Hida
%   tower): H^q(A) = H^{q-1}(G; H^1(B)), n > 0; (5) what the tree must supply (T1)-(T3).  Deferred list.  Referee m1-m4 of C-D1_three_referee.md.
%   Check outside the build: docs/implementation/slot24/cd3_descentmaps_check.py, 74 + 12 cells, 888 oracle rows, 0 mismatches.  app:status
%   394 -> 395 (prop 93 -> 94, PROVED 229 -> 230, STD 98 -> 99, tagged 264 -> 265); labelindex two labels.  222 markers on 86 labels.
%
%   VERSION 1.23.0-pre.17 (hecke-extensions slot 24, 2026-10-25: a pre-release, NOT owner-released, refereed ACCEPT (C-D1_three_referee.md); the D1 code audit's
%   finding M1, docs/implementation/slot24/D1_code_audit.md).  NEW prop:descentNthree (PROVED, STD Brown III.6.2, Weibel 2.2.6), last statement of
%   sec:descentN, so no statement number moves: when nu_3(N) = 0 (Gamma_0(N) has no element of order 3; iff 9 | N or some p = 2 mod 3 divides N),
%   every stabiliser of Gamma_0(N) on the tree and on Xbulk has order 1, 2 or 4, so the tree complex T(M_N (x) Z_(3)) and Ceq compute
%   H^*(Gamma_0(N); Sym^n) (x) Z_(3) (cmpI^* (x) Z_(3) a quasi-isomorphism), H^{>=2} has no 3-torsion, and prop:descentN(7)(b)'s sparse cokernel
%   route and its certificate hold verbatim at ell = 3 (no U-fixed points).  Necessity: nu_3 > 0 cells with 3-torsion in H^{>=2}.  Pointer in
%   prop:descentN(7); code-map paragraph and a sec:codemap row.  app:status one line (393 -> 394), prop 92 -> 93, PROVED 228 -> 229, STD 97 -> 98,
%   tagged 263 -> 264; app:labelindex one label.  222 markers on 86 labels.
%
%   VERSION 1.23.0-pre.16 (hecke-extensions slot 24, 2026-10-25: a pre-release, NOT owner-released; the Dickson-proof referee's fixes
%   (docs/implementation/slot24/C-D1_dickson_referee.md, ACCEPT-with-fixes, 0 Blocker, 1 Minor, 4 nits), not re-refereed).  D-1: the proof of
%   cl:dimH1(2) becomes a lemma lem:dicksonSL2 (every prime ell, SL_2 and GL_2), cited by cl:dimH1(2), prop:oddnltwo(2), prop:dicksonH0(1)/(2);
%   "first non-zero at n = ell+1" restricted to ell >= 3 (delta_2(2) = 1).  D-2: prop:oddnltwo(2)'s \STD tag removed.  D-3: prop:dicksonH0(2)
%   rests on the SL_2 half only; the GL_2 half proved (det acts on L by itself, Q invariant).  D-4: Step 3 says why theta has a minimal
%   polynomial, and which five coefficients are compared.  D-5: b_1 = 1 exactly at N in {2,3,5,7,13}.  app:status: one line (392 -> 393),
%   PROVED 227 -> 228, lem 60 -> 61; app:labelindex one label.  Statement numbers after cl:dimH1 in its section shift by one.  222 markers on 86 labels.
%
%   VERSION 1.23.0-pre.15 (hecke-extensions slot 24, 2026-10-25: a pre-release, NOT owner-released; the C-D1 / C-D2 referee's fixes
%   (docs/implementation/slot24/C-D1_referee.md, ACCEPT-with-fixes, 0 Blocker, 1 Major, 4 minor, 2 nits), not re-refereed).  F1 (Major):
%   cl:dimH1(2) -- the Dickson invariants of SL_2(F_ell) -- PROVED in place (L, Q = L_2/L from the Moore determinant, integrality, the
%   degree bound [K:k] <= |G| against the trivial stabiliser of x + y^2, factorial => integrally closed), so the unverified \cite{wilkerson}
%   pointer no longer carries anything; prop:badprimesN(2)'s good-prime rule ell > n made independent of it (Claim A + the lower unipotent).
%   F2: prop:descentN(4) b_1 >= 1 for every N > 1 (6 b_1 = 6 + m_N - 3 nu_2 - 4 nu_3, CRT + Hensel), equality at N = 2, 3, 5, 7, 13.
%   F3: the code map's SparseCokernel sentence corrected (averaging-idempotent generators; the blocker is the field gate and RREF).  F4: (7)(b)'s
%   generator certificate (fixed by a generator, spanning mod ell; Nakayama).  F5: (5) at chain level (Hom(E, M) induces e cup) and U_p's
%   defining label.  F6, F7: nits.  Pointers to cl:dimH1(2)'s unverified citation updated in prop:oddnltwo(2), prop:badprimes' proof,
%   prop:dicksonH0(1) (its SL_2 half).  app:status: cl:dimH1's tag set gains PROVED, loses \cite{wilkerson} (PROVED 226 -> 227); prop:descentN's
%   citation list gains Brown Ch. V.3.  No label, statement number or marker moves (222 on 86 labels).
%
%   VERSION 1.23.0-pre.14 (hecke-extensions slot 24, 2026-10-25: a pre-release DRAFT, NOT owner-released, NOT refereed; milestone D1 of
%   docs/implementation/slot24/DESCENT_plan.md v0.3, clauses C-D1 and C-D2, audit findings M8, m3, m4, m7 of DESCENT_plan_audit.md; contract first,
%   rule 21).  NEW subsection sec:descentN after the last statement of sec:isomodelell, so that no statement number moves: prop:descentN (PROVED)
%   -- prop:descent / rmk:truncdevice for every Z[Gamma]-module M finitely generated and free over Z (M_N = Ind Sym^n, M_N(chi) over O_chi =
%   Z[zeta_m]); Hecke determined locally by prop:cohglue's morphism clause; prop:homvanish(2) holds for M iff b_0 b_1 = 0 (so for M_N(chi) iff
%   n > 0 or chi != 1; fails at n = 0, chi = 1, N = 2); periodicity above the vcd by res(e) cup, equivariant for every [AgA]; the good prime
%   (M_N: every ell' > n+1, ell' not dividing N not needed; chi != 1: ell' not dividing m_chi N needed); the certified K on the direct Smith of
%   C_eq (x) Z/ell^K and on the level-N sparse cokernel over Z/ell^K at ell >= 5 (rank_Q d^0 = rank C^0 - b_0); prop:certifiedeell(4) at level
%   N; the nebentypus module mod ell^K as a sum of W_K(F_q) character models.  prop:badprimesN (PROVED) -- for every such M and ell >= 5,
%   torsion only in H^1 with s_ell = dim H^0(M/ell) - b_0 summands; at level N, s_ell = delta_ell(n) (ell not dividing N, the level-1 Dickson
%   count) or beta_ell(n) (ell | N, the Borel count, reaching ell = n+1); the closed-form list B(n,N); nebentypus at ell not dividing 6 m_chi.
%   Code map (not built: the direct Smith route, the sparse cokernel over Z/ell^K, the level-N assembly, M_N(chi) over O_chi).  rmk:integralscope(4)
%   one parenthesis (C-D7, the level half); the sec:codemap row one sentence; app:status two lines (390 -> 392); app:labelindex three labels.
%   No existing statement, label, number, tag set, marker or value moves (222 on 86 labels; the new statements carry none).
%
%   VERSION 1.23.0-pre.13 (hecke-extensions slot 24, 2026-10-25: a pre-release, NOT owner-released; the C-D0 referee's fixes
%   (docs/implementation/slot24/C-D0_referee.md, ACCEPT-with-fixes, 0 Blocker, 2 Major, 6 minor), not re-refereed).  prop:certifiedeell:
%   (4) gains (a) q >= 1: 24 H^j(Gamma;M) = 0 for j >= 2 and EVERY Gamma-module M (Mayer-Vietoris of cl:amalgam(1) with Shapiro), so every
%   ell' >= 5 qualifies, q = 1 included, and b_j = 0 for j >= 2; (b) q = 0: ell' qualifies iff dim H^0(M/ell') = b_0 (universal coefficients
%   over O_ell), rank_Q d^0 = rank C^0 - b_0, b_0 = 0 (n > 0) / 1 (n = 0) for Sym^n and Ind_{Gamma_0(N)} Sym^n, proved; (c) the two admissible
%   sources of rank_Q d^q (this complex mod ell', or certified Betti numbers), census / shipped C_7 / Iso^(ell') ranks NOT admissible (F1, F2).
%   (6): the valuations a list with multiplicity, count s (F3); a-priori exponents a_t <= v_ell(24) for q >= 1, the stopping rule (F4).
%   Code map: the denominator check and the ell()-guard as D0 obligations (F5, F2); W <= 61 reachable now through ModCtx<uint64_t>, W in
%   [62,64] not built, wording harmonised with the sec:codemap row, which said "Z/2^64" (F6); the DenseBuf -> Ring2k / Ring3k hand-off and
%   q_max (F7); the plan row amended (F8, DESCENT_plan.md v0.3).  app:status: the proof's new STD input (Brown III.10) named on the line, no
%   count moves.  No label, number, tag set or marker moves (222 on 86 labels).
%
%   VERSION 1.23.0-pre.12 (hecke-extensions slot 24, 2026-10-25: a pre-release DRAFT, NOT owner-released, NOT refereed; milestone D0 of
%   docs/implementation/slot24/DESCENT_plan.md v0.2, clause C-D0, audit finding M3 of DESCENT_plan_audit.md; contract first, rule 21).  NEW
%   prop:certifiedeell (PROVED) at the end of sec:isomodelell's statements, so that no statement number moves: Iso^(ell) built once over
%   O_ell = Z[1/q_ell], q_ell = lcm |Q_sigma| (q_2 = 3, q_3 = 4: |Q_sigma| in {1,3} resp. {2,4}; the plan's "{1,2} at ell = 3" corrected), its base
%   changes to Z/ell^K and F_ell' (ell' not dividing q_ell) the complexes the code builds; prop:certifiede(1)-(3) and rmk:certifiedeperdiff over
%   Z_(ell), rank_Q read at ell'; the full layout (delta^q, P_q) of the code; the Ring2k / Ring3k Smith-valuation count = r_K.  Its code map
%   (built: one complex over Q[Gamma] with denominators |Q_sigma| only; not built: the Ring2k / Ring3k read-out, the reduction to Z/2^64,
%   Ring3k's K <= 19).  sec:isomodelell's fixes list (v), notation sentence, scope sentence, formalization note; the sec:codemap row;
%   app:status one line (389 -> 390); app:labelindex one label.  No existing statement, label, number, tag set, marker or value moves
%   (222 on 86 labels; the new statement carries none).
%
%   VERSION 1.23.0-pre.11 (hecke-extensions slot 24, 2026-10-21: a pre-release, NOT owner-released, not refereed; code-map text only,
%   from the audited packages G2-H (docs/implementation/slot24/G2H_report.md sec. 9, G2H_audit.md F6) and G6 (G6_report.md, G6_audit.md)).
%   sec:obligations' code map: G2-H (on the stream route under RunModuleT's lean retention, obl:welldef(iii) certified without forming
%   Theta_eq by unit_pivot_residual_product, exactly the read-off's residual, so discharged as written; route 2 subtracted in place,
%   obl:agreement(i) as Theta - route 2 = 0; early release of X_q, Zc; howell_rows adopts its store; rollback HECKE_MODULE_LEAN_PEAK=0).
%   The sec:cone code-map row: G6 (ConeCertificate's span inclusions C2 / C3 / C5 by the unit-row route, fallback to the elimination; the
%   shared Z; the census cone release; rollback HECKE_FIELD_DIET=0, which does not cover tape::factor_basis's two memory-only re-blockings).
%   No statement, label, number, tag set, marker or value moves (222 on 86 labels).
%
%   VERSION 1.23.0-pre.10 (hecke-extensions slot 24, 2026-10-20: a pre-release, NOT owner-released, not refereed; code-map text only,
%   from the audited packages FNTT (docs/implementation/slot24/FNTT_report.md, FNTT_audit.md) and MT-1 (MT1_report.md, MT1_audit.md)).
%   sec:saturation's code map: the forms cross-route's transform-domain route (cached spectra, one pass, delayed reduction in
%   HeckeOnCusp; rollback HECKE_FORMS_NTT_CACHE=0) and the forms ell-range restated (FormsSpaceSupported: every 5 <= ell < 2^16; for
%   2^16 <= ell < 2^27 only N < 1024 or the two-prime NTT exact, N(ell-1)^2 < P1 P2, 2N-1 <= 2^25; refused before the module run
%   otherwise; ell = 100003 admitted to n = 464, largest ell at n = 1000 is 77702).  sec:obligations' code map: MT-1 (the pivot-search
%   exit at the first unit at every thread count; further pool splits; obl:welldef, obl:agreement(i) discharged as written; rollback
%   HECKE_MODULE_MT=0).  L23-H1 (orbifold_cell/v2) touches no text here.  No statement, label, number, tag set, marker or value moves
%   (222 on 86 labels).
%
%   VERSION 1.23.0-pre.9 (hecke-extensions slot 24, 2026-10-20: a pre-release DRAFT, NOT owner-released, NOT refereed; package
%   ELL-W, docs/implementation/slot24/ELLW_report.md; plan owner question Q5).  AMENDS A PROVED LEMMA: lem:bounds gains clause (7)
%   (the wide tier: products of two stored-class matrices in 128 bits under 2 N B <= 2^63 - 1) with its proof; thm:backend gains
%   consequence (4) (the stored wide tier) and one proof sentence; obl:backend(ii) names the wide tier's assertion (code 52);
%   rmk:btabsmall gains the measured B_tab past 23170; sec:heckedatum's reach paragraph and the backend code map record the lift
%   (ell <= 23170 narrow, ell <= 1048573 by clause (e) of the tau certification).  No label is added, renamed or removed; no existing
%   clause number moves (clause (7) is appended); no tag set or LEAN marker moves (222 on 86 labels); every value at ell <= 23170 is
%   unchanged (the 64-bit evaluation runs first; golden and digests unchanged).
%
%   VERSION 1.23.0-pre.8 (hecke-extensions slot 23, 2026-10-18: a pre-release, NOT owner-released, not refereed; audited G2-T2,
%   docs/implementation/slot23/G2T2_audit.md, PASS, and G2-P, G2P_audit.md, PASS, fixes F1-F4 applied).  Code-map text only:
%   sec:obligations' G2 code-map paragraph adds G2-T2 (the SIMD rho fill, lever HECKE_RHO_SIMD; route 2's memo shared across chunks,
%   lever HECKE_ROUTE2_SHARED_MEMO) and G2-P (the module route on the consumer pool, lever HECKE_MODULE_PARALLEL; thread-count
%   independence).  Every lever value-neutral.  No statement, label, number, tag set, marker or value moves (222 on 86 labels).
%
%   VERSION 1.23.0-pre.7 (hecke-extensions slot 23, 2026-10-18: a pre-release, NOT owner-released, not refereed; audited G2-T,
%   docs/implementation/slot23/G2T_audit.md, PASS-WITH-FIXES).  Code-map text only: sec:obligations' G2 code-map paragraph adds G2-T, the
%   module route's term evaluation (ComputeSymmetricPowerRowsRing, factor_low_rank fast_scan) and its value-neutral rollback lever
%   HECKE_MODULE_EVAL=0.  No statement, label, number, tag set, marker or value moves (222 on 86 labels).
%
%   VERSION 1.23.0-pre.6 (hecke-extensions slot 23, 2026-10-18: a repair-only pre-release, NOT owner-released; the fixes F1-F6 of
%   docs/implementation/slot23/C1230_pre3_referee.md (Opus, adversarial; pre.1-pre.3 ACCEPT-WITH-FIXES, MINOR 2, NIT 4); on top of v1.22.0, signed by
%   the owner 2026-10-16).  F1: the proof of lem:precisionreduction states the tau-homotopy over C_eq(Z_ell) (prop:coneindep(2) with lem:tauexists(2),
%   6 a unit; prop:coneindep(3),(4) for rechosen s_j / g), not over C_eq(Z).  F2: the lemma and its formalization note name thm:relations(1) (the T_p
%   commute on M_K; its proof uses thm:classical).  F3: K_cls attributed correctly (e_max of prop:torsionprepass); KMAX (a run setting) and K_stab
%   defined.  F4: (3) covers readings across levels via (1); obl:precisionreuse(iv) names the computed LOEWY_RING fields.  F5 is code only
%   (HeckeLevelConeValidation's tag); F6 this header's and \date's sign-off record.  Beyond the one clause added to lem:precisionreduction(3) (a
%   consequence of its (1)) no statement's mathematics, label, number, tag set, marker or value moves (222 on 86 labels).
%
%   VERSION 1.23.0-pre.5 (hecke-extensions slot 23, 2026-10-18: a pre-release, NOT owner-released, not refereed; audited fix F1 of
%   docs/implementation/slot23/G2F_audit.md).  Code-map text only: obl:tau's code map and its status row state that a
%   TauTableContents::ChainsOnly table (the module route's, RunModuleT on the stream route) carries no consumer table, so (vi) is vacuous there
%   and (i)-(v), (vii) run unchanged; the module route evaluates rho_n mod ell^k from exact gamma.  No statement, label, number, tag set, marker
%   or value moves (222 on 86 labels).
%
%   VERSION 1.23.0-pre.4 (hecke-extensions slot 23, 2026-10-18: a mid-slot pre-release, NOT owner-released, not a release candidate -- the X.Y.0
%   candidate is cut at the slot close; refereed ACCEPT-WITH-FIXES, docs/implementation/slot23/C1230_codemap_referee.md, fixes F1-F10 applied;
%   pre.3 stays audited, not refereed).  Code-map text only, the
%   follow-up of docs/implementation/slot23/HANDOVER_1018.md: the code map of sec:satprecisionsec flipped from planned to built (G1, G1_impl_audit F8:
%   ringalg::LevelCache, ReduceToLevel / CheckDerivedModule / CheckRunModule, RunLevelPolicy, rollback HECKE_PRECISION_REUSE=0); the code map of
%   sec:obligations gains the slot-23 G2 module-route re-implementation (G2-A/B/C/D/E/G: Theta_eq, X_q and H_class_coordinates read off the unit
%   pivots of C, d o d by GEMM, Omega streamed per coset, route 2 in sigma-row chunks, the product read-off, lean retention for RunModuleT, lean
%   work buffers; rollback HECKE_MODULE_STREAM=0); the code map of sec:saturation the G2-X forms product (NTT + CRT; rollback HECKE_FORMS_FASTMUL=0);
%   the sec:codemap row of sec:saturation a pointer sentence.  Every rollback lever is value-neutral.  No statement, label, number, tag set, marker
%   or value moves (222 on 86 labels).
%
%   VERSION 1.23.0-pre.3 (hecke-extensions slot 23, G1, 2026-10-16; a pre-release, NOT owner-released; audited, ACCEPT-WITH-FIXES, docs/implementation/slot23/G1_clause_audit.md, fixes F1-F12 applied; on top of the signed v1.22.0,
%   owner sign-off 2026-10-16, through pre.1 / pre.2): new sec:satprecisionsec at
%   the end of sec:saturation -- lem:precisionreduction (PROVED: for ell >= 5 at level 1 the module route's H^1 over Z/ell^K is the reduction of its
%   H^1 over Z/ell^K', K <= K', Hecke-equivariantly, because H^2(C_eq(Z_ell)) = 0; downward only; the hypothesis is needed, with the H^0 / Dickson
%   counter-case) and obl:precisionreuse (one module run per cell, the per-K route kept as oracle and rollback arm, byte-identical output).  Owner
%   decision 2026-10-16 approving lever G1 of docs/implementation/slot23/ELL_scout.md sec. 3; note docs/implementation/slot23/G1_clause_note.md.
%   Additive: two new statements and one subsection label; no existing statement, number, tag set, marker or value moves (222 on 86 labels).
%
%   VERSION 1.23.0-pre.2 (hecke-extensions slot 23, LES-2 repairs, 2026-10-12; a pre-release, NOT owner-released): the repairs of
%   docs/implementation/slot23/F-C123pre1_referee.md (Fable, ACCEPT, MINOR 2, NIT 5): MINOR-1 (changelog bullet "Not edited here" reworded as the audit's
%   proposal), MINOR-2 ((ix) now says that the two-transversal comparison on C_eq is GATED, the flat one recorded, neither a negative control), NIT-1..-5
%   (wording; chapter comment before sec:nebles).  Repair-only: no statement, label, number, tag set, marker or value moves (222 on 86 labels).
%
%   VERSION 1.23.0-pre.1 (hecke-extensions slot 23, LES-2, 2026-10-12; a pre-release, NOT owner-released, refereed by F-C123pre1, ACCEPT): prop:levelconehecke(4) is
%   REPLACED -- by (T1) of def:tau the level-correspondence operator is independent of the transversal on equivariant cochains (so "independently fixed
%   transversals" is no counter-example and the strict square holds for every pair of transversals under the criterion (1)); the implemented operator is
%   moreover transversal-independent on the whole flat cochain space (it uses r_t only through its coset).  obl:levelconeexact(ix): the transversal-mismatch
%   negative control (which can never fire) is replaced by the criterion-fails control.  Sources: docs/implementation/slot23/LES2_audit.md sec. 3.2 (with
%   the flat-cochain sentence of the proposal corrected, R1 post-run note of LES2_predictions.md sec. 7), LES2_report.md sec. 3.  No label, number, tag set,
%   marker or value moves (222 on 86 labels).
%
%   VERSION 1.22.0 (hecke-extensions slot 22, CLOSE-T T3, 2026-10-12: the RELEASE CANDIDATE from v1.22.0-pre.17 -- every pre.N text of
%   1.22.0-pre.1 ... pre.17 has a referee record or is listed as unrefereed in its changelog entry; docs/implementation/slot22/
%   release_referee_coverage.md, F-T1_release_referee.md, slot22_split_close_plan.md sec. 1.1).  NOT YET OWNER-SIGNED: owner sign-off is
%   CLOSE-T step T8.  Header, \date and changelog only: no statement's mathematics, label, number, tag, status or marker changes (222 on
%   86 labels); a release upgrades no referee status.  Companions [E] 1.22.0, [U] 0.6.0, [L37] 0.5.2, [RC] 0.2.0, [MPT] 0.3.1, [CLI] 1.4.0, [P] 3.12.0.
%
%   VERSION 1.22.0-pre.17 (hecke-extensions slot 22, CLOSE-T T1 repairs, 2026-10-12; a pre-release, NOT owner-released): the repairs of
%   docs/implementation/slot22/F-T1_release_referee.md (Fable, ACCEPT on every gap, MAJOR 0, MINOR 2, NIT 15).  MINOR-1: the missing
%   changelog subsection of v1.22.0-pre.9 added.  NIT-1, -2, -3, -4, -8, -9, -10 applied (wording / punctuation / comment order);
%   NIT-5, -6, -7 not applied (they touch the scope of a statement), NIT-11..-15 concern [E] / [U] (see there).  No statement, label,
%   number, marker, tag set or value moves.
%
%   VERSION 1.22.0-pre.16 (hecke-extensions slot 22, LR-3, 2026-10-12; a pre-release, NOT owner-released, NOT refereed as [C] text --
%   the text is LRY_phase5_report.md sec. D as repaired by LRY5_audit.md (Fable, ACCEPT, MINOR-1 applied: "and the identifications of
%   thm:mp-mv-terms" in the thm:mp-mv-res and thm:mp-mv-les scopes, written as a [GT] reference, not a second marker); applied on the owner's
%   approval of LR-3, 2026-10-12): LR-Y phase 5.  NEW \LEAN markers: 4 on 0 new labels -- cl:amalgam (1) thm:mp-treeses-rep;
%   cl:amalgam (2) thm:mp-mv-res, thm:mp-mv-les, thm:mp-mv-natural; 222 markers on 86 labels in all (208 \LEAN, 14 \LEANlit).  The scope
%   of thm:mp-model7-tree-orbits loses "the sequence is not assembled as one object of Rep k Gamma"; the "Not in Lean" sentence of
%   cl:amalgam (2) is removed (its three items are now the three new markers).  No statement, label, number, tag set or value moves.
%   Report cpre16_report.md.
%
%   VERSION 1.22.0-pre.15 (hecke-extensions slot 22, F-pre14 referee repairs, 2026-10-11; a pre-release, NOT owner-released, NOT refereed as [C]
%   text -- the wording of docs/implementation/slot22/F-pre14_referee.md (commit 98b9ad8b, Fable, ACCEPT-WITH-REPAIRS) sec. 6, applied verbatim:
%   MINOR-1, NIT-1, NIT-2, NIT-3, NIT-5.  NIT-4 (label-base convention) and NIT-6 (app:labelindex, slot close) not applied.  Wording only:
%   no \LEAN marker, label, statement, number, tag or value moves; marker count stays 218 on 86 labels.  Report cpre15_report.md.
%
%   VERSION 1.22.0-pre.14 (hecke-extensions slot 22, LV markers, 2026-10-11; a pre-release, NOT owner-released, NOT refereed as [C]
%   text -- worded within LV2_audit.md sec. 8 (Fable, 3d7b868d, ACCEPT-WITH-REPAIRS) and LVprov_audit.md (a2ea84b5)): the Lean
%   dimension results at ell = 2, 3 (LV-0..LV-2, [GT] Part MP sections LV-1, LV-2).  NEW \LEAN markers: 12 on 4 labels, none of which
%   carried a marker -- lem:centralminusone (last clause, FIELD case only) thm:lv3-tree-negone; prop:oddnltwo (2) chk:lv-cells (the
%   value of h^0 at the odd n <= 63 and n = 101, 129, 255 only); prop:ellthreeclosed (4) thm:lv3-tree-negone (field case only); NEW
%   rmk:lvdimsltwothree (PROVED, at the end of sec:oddnltwo) thm:lv-tree-f2-dims, thm:lv-f1, thm:lv-v2-dims, thm:lv3-tree-triv,
%   thm:lv3-rankf3, thm:lv3-f1, thm:lv3-v3-dims, chk:lv-cells, chk:lv3-cells; 218 markers on 86 labels in all (204 \LEAN, 14 \LEANlit).
%   The L23Cells cells (n > 64) are cited as data through docs/implementation/slot22/lv/, never by a marker.  No marker on
%   prop:ltwothreescalar (Lean proves nothing about T_p).  The marker-history paragraph before sec:labelpolicy brought from pre.8 to
%   pre.14.  app:status: one line (388 -> 389).  No existing statement text, label, number before the new remark, tag set or value
%   moves; statement numbers after rmk:lvdimsltwothree within sec:six shift.  Report cpre14_report.md.
%
%   VERSION 1.22.0-pre.13 (hecke-extensions slot 22, LRY4 repairs, 2026-10-11; a pre-release, NOT owner-released, NOT refereed as [C]
%   text -- the text is LRY4_audit.md sec. 6 (Fable, 9feb57e2, ACCEPT-WITH-REPAIRS), verbatim, which carries MINOR-1, -2, -3 and NIT-1):
%   LR-Y phase 4 (docs/implementation/slot22/LRY_phase4_report.md).  NEW \LEAN markers: 4 on 0 new labels -- cl:amalgam (2)
%   thm:mp-mv-terms, thm:mp-mv-resres, thm:mp-mv-low; lem:bulkperiodic (2) thm:mp-periodic-natural; 206 markers on 82 labels in all
%   (192 \LEAN, 14 \LEANlit).  Rewritten marker scopes (no statement text): cl:amalgam (2) thm:mp-tree-mv ("replaced by" -> "built on";
%   the comparison named as mathlib's groupCohomologyIso through Ext; the "Not in Lean" list now: the map H^q(Gamma) -> H^q(<S>)+H^q(<U>)
%   is restriction, the connecting map in groupCohomology terms, naturality in V); lem:bulkperiodic (2) thm:mp-tree-periodic (only the
%   existence of an isomorphism, natural on the explicit complex; "not naturality on H^q(Gamma;V)" dropped, now covered by the new
%   marker, which is NOT shown to be the same isomorphism; LRY4_audit MINOR-1, -2).  No statement, label, number, tag set or value
%   moves.  Report cpre13_report.md.
%
%   VERSION 1.22.0-pre.12 (hecke-extensions slot 22, LR-3, 2026-10-11; a pre-release, NOT owner-released, NOT refereed as [C] text -- the
%   proposed text was audited by Fable, docs/implementation/slot22/LRY3_audit.md, c3bbe552, ACCEPT-WITH-REPAIRS, all MINOR-1..6 and the [C]
%   NITs applied): LR-Y phases 1-3 (docs/implementation/slot22/LRY_phase1_report.md, LRY_phase2_report.md, LRY_phase3_report.md sec. C).
%   NEW \LEAN markers: 7 on 3 labels -- cl:amalgam (1) thm:mp-treeses-exact, thm:mp-model7-tree-orbits; cl:amalgam (2) thm:mp-tree-mv;
%   lem:bulkperiodic (2) thm:mp-tree-periodic; rmk:treeconeres thm:mp-model7-cone-res, thm:mp-tree-mp-comparison, thm:mp-tree-cochains;
%   202 markers on 82 labels in all (188 \LEAN, 14 \LEANlit).  NEW rmk:treeconeres after cl:amalgam (the tree resolution P_. over every
%   ring, its orientation vs the tree Y of sec:levelN, the explicit cochain complex V^2 -> V^3 -> V^3 -> ...).  cl:amalgam (1): the tree's
%   vertex / edge description added after eq:treeses, and -- OWNER-VISIBLE, the one change to an existing clause's wording -- the last
%   sentence of its \STD pointer ("A reader who prefers not to rely on it may take eq:treeses itself as the hypothesis; ...") replaced by
%   "eq:treeses is proved independently of this pointer (the markers below), so nothing in this document rests on it; the amalgam
%   decomposition itself is used nowhere else" (LRY3_audit MINOR-4).  No statement, label, number of an existing statement before
%   rmk:treeconeres, tag set or value moves; statement numbers after rmk:treeconeres within sec:six shift.  Report cpre12_report.md.
%
%   VERSION 1.22.0-pre.11 (hecke-extensions slot 22, pre11 repairs, 2026-10-10; a pre-release, NOT owner-released, NOT refereed):
%   F-L23 referee (docs/implementation/slot22/F-L23_referee.md, d3ef8b2d, ACCEPT) repairs R1 (prop:oddnhecke(4): the "equivalently"
%   replaced by the one-way implication), R2 (prop:ltwothreescalar: the algebra generated by all T_p), R3 ((I1)'s representative-
%   independence proved in place), R4 ((F1)'s module identification derived from the cusp count), R5 (Step 0(b) wording).  F-MPT referee
%   (docs/implementation/slot22/mp_triangulations_referee.md, 54a55ec0) sec. 3: NEW def:mpqd (Q-domination (QD)) and prop:mpqreduced (the
%   Q-reduced nerve D = F'/~ is a projective resolution; PROVED, COMP^a = fmptri_spotcheck.py), closing Gap G1 (rmk:mphybrid(3) rewritten).
%   Two new statements; no existing tag set moves; no \LEAN marker.
%
%   VERSION 1.22.0-pre.10 (hecke-extensions slot 22, L23-0 extended, 2026-10-10; a pre-release, NOT owner-released, NOT refereed): new
%   subsection sec:oddnltwo in sec:six -- odd n at ell = 2 (lem:centralminusone, prop:oddnltwo, lem:oddnsplit, prop:oddnhecke), the all-ell
%   degree-0 Bockstein (lem:bocksteinallell), cyclic restriction (lem:cyclicres), T_ell rank <= 2 (lem:tellrank), T_p = 1+p at ell in {2,3}
%   (prop:ltwothreescalar, given (F1), (I1)), the c_p rule (lem:cprule), the standard twist (rmk:standardtwist), the [R] B(6) erratum
%   (rmk:oddnltwoerr).  Eleven new statements; no existing statement, tag or marker moves; no \LEAN marker.
%
%   VERSION 1.22.0-pre.9 (hecke-extensions slot 22, C-pre9, 2026-10-10; a pre-release, NOT owner-released, NOT refereed): sec:mpmodel
%   code-vs-chapter wording repairs: rmk:mprulesread COMP (four oracle rows, 12 cells; R_zeta row ungated), good-fails labels
%   (R_Z^+ at m only, R_Z at [1,2,7]), four rows in good_forms_agree, MPC3 planted-stabiliser control (only the H^2 change gated),
%   ex:mpteardrop control (second non-inverse mu pair). No statement, number or marker moves.
%
%   VERSION 1.22.0-pre.8 (hecke-extensions slot 22, C-pre8, 2026-10-10; a pre-release, NOT owner-released, NOT refereed): sec:mpmodel
%   pre.7 referee (docs/implementation/slot22/C_pre7_referee.md, Fable, ACCEPT) minor repairs F1-F4: prop:mpbredon proof "clause (2)
%   below"; prop:mprescor(1) "any free Gamma-set resolving Z"; def:mprules(3) "under (IO)"; thm:mpnatural(4) proof "Y connected by (AC)".
%   MP-4 (merged eac3f3b3, audit F-MP4 sec. 3): rmk:mphecke replaced (the twisted chain map tau on the MP resolution itself, unique up to
%   twisted homotopy over every ring since N Z[F'] is free over Z[Gamma''] -- no 6 in R^x; the label factor rho_n(g); G3 not needed to
%   compute), its COMP^r record (OrbifoldMPHeckeValidation, 2059 checks, non-vacuous counts, H^{<=2}, not gated: H^3, T_{p^nu}, Z/l^k
%   charpolys, level N, iota-sign), tag commentary -> COMP^r; def:mpmorphism(3) the g-twist clause; rmk:mpcheckable MPC12.  MP-3b (merged
%   ec982f93, audit F-MP3b R1-R2): MPC11 537 -> 723 checks with the Z/4, Z/9, Z/25, R_Z and -I not in A pairs, MPC10 331 -> 332; the
%   "not gated" clauses reduced.  app:status COMP 56 -> 57 (registered 28 -> 29).  Report C_pre8_report.md.
%   VERSION 1.22.0-pre.7 (hecke-extensions slot 22, C-pre7, 2026-10-10; a pre-release, NOT owner-released; NOT refereed at pre.7 -- refereed
%   at pre.8 by the Fable referee C_pre7_referee.md, ACCEPT, minor repairs F1-F4 applied there): sec:mpmodel
%   NEW \LEAN markers per docs/implementation/slot22/MPL3_report.md v1.0 sec. 5 (Lean MP-L3 merged at c5617bc7, [GT] Part G12): 25 on
%   def:mpmorphism, thm:mpnatural (and its COMP record), prop:mpbredon, def:mpset, def:mpfree, thm:mpcomparison; 195 markers on 79 labels
%   in all (181 \LEAN, 14 \LEANlit); thm:mp-dd's scope note "not the normalised complex" dropped.  Unnormalised clauses written into
%   def:mpmorphism(3), thm:mpnatural(1)-(3) (with proofs: u commutes with degeneracies; C in C_un a homotopy equivalence), prop:mpbredon(1)
%   (any lift Phi of the identity: Phi^* ~ epsilon^MP followed by C in C_un; proof); def:mpmorphism(2) and thm:mpnatural(1)'s proof: an
%   equivariant vertex map certifies a label table, only "non-empty to non-empty" is used.  MP-3 (F-MP3 sec. 10): prop:mprescor(1) on
%   F'_Gamma (either nerve serves), prop:mpboundary(2) the cone indexing, COMP^r records after both (tags PROVED -> PROVED COMP^r),
%   rmk:mpcheckable MPC10, MPC11 and what is not gated (Z/l^k level pairs, R_Z pairs, -I not in A); code map.  Report C_pre7_report.md.
%   VERSION 1.22.0-pre.6 (hecke-extensions slot 22, C-pre6, 2026-10-10; a pre-release, NOT owner-released): sec:mpmodel refereed by
%   F-MP2 (contract part, docs/implementation/slot22/F-MP2_contract_referee.md, ACCEPT WITH REPAIRS), repairs R1-R8 applied: R1
%   thm:mpnatural(4) first sentence (false as stated: labels +-I do not make f o u a cochain map unless they act by one and the same
%   central automorphism) and one clause of its proof; R2 prop:mpboundary(2) proof (Z[j] over the augmentation Z[Gamma/Gamma_bd] -> Z,
%   Ext(eps, M), Shapiro = res; Brown III.9, VERIFYloc) and its app:status locator; R3 (PD) as purity; R4 def:mpphi's choices; R5
%   rmk:mphybrid(1) "an instance of"; R6 rmk:mpwhat(3) the coset sentence and which count is which; R7 D declared; R8 eight hygiene
%   items.  def:mprules(3): the strong form of goodness is "T(v) in the stabiliser Gamma_p of one vertex p of St(v)", the form checked by
%   C++ and Lean (Lean goodVtx_discrepancy: "T(v) <= Gamma_v" fails at every vertex orbit of M7); rmk:mprulesread(3) accordingly.  The
%   unnormalised MP complex C_un named in def:mpset, thm:mpcomparison(3),(4) (the form Lean proves).  NEW \LEAN markers on sec:mpmodel
%   per docs/implementation/slot22/MPL1_report.md v1.1 sec. 5: 30 on 11 labels ([GT] Part G12, gt_part_MP.tex); 170 markers on 76
%   labels in all (156 \LEAN, 14 \LEANlit).  No statement number, label, tag set, status or value moves.  Report C_pre6_report.md.
%   VERSION 1.22.0-pre.5 (hecke-extensions slot 22, MP-CT, 2026-10-09; a pre-release, NOT owner-released; NOT refereed at pre.5 --
%   refereed at pre.6 by F-MP2 (contract part), ACCEPT WITH REPAIRS, repairs R1-R8 applied there; owner rulings Q1-Q7 of the Moerdijk-Pronk plans taken at the recommendations, 2026-10-09): NEW chapter sec:mpmodel after
%   sec:isomodelell, before the appendices (so no chapter or statement number of pre.4 moves; the implementation plan had named the place
%   after sec:orbifoldmodel): orbifold triangulations as descriptors and MP's rules as decidable predicates (def:mptri, def:mprules,
%   def:mpell, def:mpRZ, tab:mprules, rmk:mprulesread), the choice-free MP set (def:mpset, rmk:mpwhat), the comparison theorem in every
%   degree over every ring, -I included (def:mpfree, lem:mprows, lem:mpcols, lem:mpassembly, thm:mpcomparison), the l-local hybrid
%   = prop:isomodelell (rmk:mphybrid), all degrees with its exact scope (def:mpphi, prop:mpalldegrees: bulk, Gamma = SL_2(Z), D = 3,
%   24 H^{>=2} = 0, n-range certified), morphisms with the label factor rho(h_{sigma_0}) (def:mpmorphism, thm:mpnatural), the Bredon
%   comparison (prop:mpbredon: quasi-iso if 6 in R^x, Z[1/2] odd-n witness), the boundary (prop:mpboundary), res/cor and the pair of
%   groups (prop:mprescor), Hecke (rmk:mphecke, commentary), the teardrop (ex:mpteardrop), rmk:mpcheckable ((MPC1)-(MPC9)); addenda to
%   rmk:mpcaveat, rmk:orbifoldrefine, rmk:isomodellmp(3); a sec:codemap row.  23 statements (app:status 351 -> 374), app:labelindex
%   +40 labels.  No shipped value, default or golden object changes (Q4: F_2 / F_3 keep their Bredon semantics); no existing number,
%   label, tag set or marker moves (140 on 65; the new labels carry none).  Sources docs/implementation/slot22/MP_math_plan.md v0.2.1,
%   MP_impl_plan.md v0.2.1, F-MP2_code_audit.md sec. 1.3, MP1a/MP1b/MPT/MP2 reports; report MPCT_report.md.
%   VERSION 1.22.0-pre.4 (hecke-extensions slot 22, MP-0, 2026-10-09; a pre-release, NOT owner-released; owner-approved erratum, MP question Q6):
%   rmk:negI, second bullet: the words "under the additional hypothesis 6 in R^x" are deleted --- the group-side vanishing argument
%   (-I central, acting trivially on H^*(Gamma;-) and by (-1)^n = -1 on the coefficients, so 2 H^* = 0) needs 2 in R^x only, which is
%   automatic at ell != 2 and is the bullet's own case.  Wording only: no statement, label, tag, marker (140 on 65) or value moves.
%   Source docs/implementation/slot22/MP_math_plan.md sec. 4.2; companion docs/implementation/slot22/MP_predictions.md.
%   VERSION 1.22.0-pre.3 (hecke-extensions slot 22, LES-C, 2026-10-09; a pre-release, NOT owner-released, NOT yet refereed: the Fable referee
%   F-LES reads it before any LES code, rule 22; owner questions Q1-Q7 of docs/implementation/slot22/LES0_spec.md taken at the spec's
%   recommendations by the lead pending owner objection): NEW last subsection sec:nebles of sec:nebentypus (after sec:nebwildsec, so no
%   statement number of pre.2 moves): the level-change cone of B <= A with compatible characters on the 7-vertex model, its six-term
%   sequence at every ring and every ell (lem:levelconeexact), the split case, the Hecke squares (T_p and same-level U_p strict with the
%   smaller group's transversal; level-raising U_p OPENSTAT), the normal case = inflation-restriction plus two pieces (prop:levelinfres),
%   the cyclic Z/p step with its Jordan types (prop:cycliccone), Gamma_0(MN) normal in Gamma_0(N) iff M | gcd(N,24) (prop:gamma0normal),
%   the steps of the p-tower (prop:hidanormalsteps), level one (rmk:normalN1), the model tags at ell in {2,3} (prop:modeltorsion),
%   obl:levelconeexact, rmk:lescheckable ((LES1)-(LES9)), rmk:oldspacecone (deferred).  21 statements (app:status 330 -> 351),
%   app:labelindex +25 labels.  No existing statement number, label, tag set or marker moves (140 on 65; the new labels carry none).
%   Predictions docs/implementation/slot22/LES_predictions.md, numbers docs/implementation/slot22/les0/out_predict.txt.
%   F-LES referee repairs (same pre-release, 2026-10-09; docs/implementation/slot22/LES-C_referee.md): R1 prop:levelconehecke(3)
%   OPENSTAT -> PROVED with negative content (U^B_p o alpha = p [AgB] on cochains; the square fails at (Gamma_0(25) < Gamma_0(5), ell=7,
%   n=0); [C]'s U_p is U_p^* on old spaces), propagated to obl:levelconeexact(vi), LES9, app:status (OPENSTAT 8 -> 6, COMP 45 -> 46);
%   R3, R4 wording; R5 bib entries adamson1954, bieri_eckmann1978; R6 Gamma(M)/Gamma^2/Gamma^3 T_p instance, Gamma(3) torsion-free.
%
%   VERSION 1.22.0-pre.2 (hecke-extensions slot 22, DOC22, 2026-10-08; a pre-release, NOT owner-released): rmk:levelNcurvefamily, item (1),
%   reworded: Mazur's good-prime condition is stated as SUFFICIENT for generation (not 'iff'), and not necessary at p = 2 ((2473,2): eta_2
%   generates m T_m / 2 with 2 a quartic residue), with the p-odd, ell = 2 case; source docs/implementation/slot22/LN23F_report.md sec. 6
%   item 1.  Nothing else in [C] changed: no number, label, tag, marker or status moves.
%
%   VERSION 1.22.0-pre.1 (hecke-extensions slot 22, P10-C, 2026-10-06; a pre-release, NOT owner-released; decisions taken as recommended):
%   NEW last subsection sec:nebwildsec of sec:nebentypus (so no statement number of v1.21.0 moves), the text of
%   docs/implementation/slot22/P10C_contract_text.tex.txt v0.2 part (A) verbatim, as refereed by the Fable referee F-U5
%   (P10C_F4_referee_report.md, ACCEPT WITH REPAIRS, 0 B / 3 M / 13 m, all applied in the text): def:nebwild, lem:nebwild, def:nebgammaH,
%   prop:nebgammaH, prop:nebgammaHheilbronn, lem:nebshapiro (the M1-corrected statement: Gamma' acts on P through the inverse of the quotient
%   map), lem:nebwildtors, lem:nebchainrref, rmk:nebwildcheckable ((NB7)--(NB10); NB8's rank clause PROVED given [U] prop:u-towerorient, COMP at
%   r = 1, gated), rmk:nebcosetlevel, the code map (planned, slot-22 P10); one sentence appended to rmk:nebfield(2); the refusal box's item (ii)
%   narrowed (wild characters stay refused over F_q and W_K(F_q)); an unnumbered addendum after thm:hidacontrol (P10C_addenda.md (1): Hida's
%   Lambda-structure vs the natural diamond action, verified on Hida pp. 589/593 by F-U5; the theorem's statement unchanged).  app:status +10
%   lines (320 -> 330), app:labelindex +11 labels.  Macros \Zl, \Ql, \VERIFYloc provided (\providecommand).  No existing statement number,
%   label, tag set or marker moves (140 on 65; the new labels carry none).  Companion [U] 0.6.0-pre.1 (part 5, sec:u-tower).
%
%   VERSION 1.21.0 (hecke-extensions slot 21, the close, 2026-10-06: the RELEASE from v1.21.0-pre.4 under the owner's decision D16
%   (release of the pre-release documents at the close, accepted as recommended, 2026-10-06; docs/implementation/slot21/
%   slot21_status_and_close_plan.md Sec. 5).  Header, \date and changelog only: no statement's mathematics, label, number, tag,
%   status or marker changes (140 on 65 labels); a release upgrades no referee status.  Companions [E] 1.21.0, [U] 0.5.0,
%   [L37] 0.5.1, [RC] 0.1.0, [LG] 0.1.5, [EX] 1.4.1, [CLI] 1.3.0.
%
%   VERSION 1.21.0-pre.4 (hecke-extensions slot 21/22, LN23, 2026-10-06; a pre-release, NOT owner-released; owner decision O6 taken as
%   recommended): NEW last subsection sec:levelNcurve of sec:levelN (so no number of pre.3 moves) with prop:levelNcurve (PROVED: the n = 0
%   cokernel presentation with orbit-sum fixed spaces is H^1 of the quotient graph Gamma_0(N)\T = H^1(Y_0(N)) over Z and over F_p for every p,
%   p in {2,3} included; free of rank 2g + c - 1; H_{N,q} induces T_q over Z; one STD step (the spine, Serre Trees I.4.2, locator not
%   verified) with a prop:bredon(2) alternative) and rmk:levelNcurvefamily (STD Mazur II.16-17 / Calegari--Emerton / Sturm, CONJ, COMP^a:
%   generation by a Mazur good prime vs separation), both as delivered by the slot-22 F4 referee of LN23
%   (docs/implementation/slot22/LN23_C_proposition.tex.txt, LN23_F4_referee_report.md Sec. 11); def:levelNmodel(c) admits p in {2,3} at
%   n = 0 on the curve object; (LN6) gains the F2 caveat for the extension of (LN4)/(LN6) to p in {2,3}; the sec:levelN scope paragraph
%   names the exception; app:status +2 lines (318 -> 320), app:labelindex +3 labels; gl2_refs.bib gains calegari_emerton.  No existing
%   statement number, label, tag set or marker moves (140 on 65; the new labels carry none).  Companions unchanged.
%   VERSION 1.21.0-pre.3 (hecke-extensions slot 21, 2026-10-06; a pre-release, NOT owner-released): one addendum sentence in (Hi5):
%   [U] (U17)(a) is equivalent to [E] cl:lveis(2) on the ordinary m_Eis-block for n < ell - 1 (the reading after [E]'s proof of
%   cl:lveis; slot-21 F4 on [U], M3).  No statement, label, tag, status or marker count (140 on 65) changes.  Companions [E]
%   1.21.0-pre.5, [U] 0.5.0-pre.6.
%   VERSION 1.21.0-pre.2 (hecke-extensions slot 21, 2026-10-05; a pre-release, NOT owner-released): addenda only --- rmk:stabconvention(3)
%   records [U] lem:u-latticeord (ordinary parts are lattice-independent; Hi8 (a) = (b), (c) a theorem; slot-21 F4 on [U] part 4);
%   thm:hidacontrol(2) gains the normalisation of Hida's P_n (= P_{n+2}, Invent. p. 545, read by F4, M2); rmk:thetaclassical(1), (4) and
%   rmk:nebcheckable (NB1) gain the locators read in the slot-21 literature pass (Katz Thm (2) p. 55, Jochnowitz local Lemma 4.1(b) p. 257 and
%   Sec. 7, Cohen--Oesterle Thm 1); gl2_refs.bib notes katz1977, jochnowitz_local, cohen_oesterle.  No clause's mathematics, label, tag,
%   status or marker count (140 on 65) changes.  Companions [E] 1.21.0-pre.1, [U] 0.5.0-pre.3.
%   VERSION 1.21.0-pre.1 (hecke-extensions slot 21, R20-L, 2026-10-05; a pre-release, NOT owner-released): the Shimura / Brown /
%   Ash--Stevens locators of the nine Shimura \LEANlit markers (thm:relations x4, obl:algebra x3, thm:nebrelations x2) verified on the
%   owner's scan (Shimura Ch. 3 / Ch. 8 page images; Brown GTM 87; Ash--Stevens on GDZ), docs/implementation/slot21/R20-L_report.md;
%   "locator unverified (R20-L)" replaced by the verified locators (hcop = Prop. 3.16, hTp = Thm. 3.24(5) + Prop. 3.17, hcomm = Prop. 3.8
%   with p. 56, hUchi = Thm. 3.34(3) + Prop. 3.33, hTpchi0/hTpchi = Thm. 3.24(5) via Prop. 3.31); further STD locators repaired
%   (prop:cosets, thm:relations proof (T(p,p) not inverted; coprime = Prop. 3.16), err:e1 evidence (Thm. 3.24, not 3.34), def:nebdata,
%   Lemma 3.29, lem:heckerankone's secondary pointers, Thm. 8.4 sites, Brown 9.5(ii)).  No clause's mathematics, label, tag, marker
%   count (140 on 65) or status changes.  Companions unchanged.
%
%   VERSION 1.20.0 (hecke-extensions slot 20, the close, WP-18, 2026-10-05: the RELEASE from v1.20.0-pre.2 under the owner's ruling
%   O6 of 2026-10-04 (releases after the Fable windows), on F1 (WP-3) and F1' ACCEPT (code-map text) and the Lean audit F3 of M4 (L20-6,
%   ACCEPT WITH REPAIRS, repaired).  New over pre.2: the milestone-M4 \LEAN / \LEANlit markers clause by clause per
%   docs/implementation/slot20/L20-6_F3_audit.md Sec. 8 as amended by L20-6_repairs.md (MAJOR 1: no marker on thm:classical or on
%   eq:corconjres, not even on its maps (stricter than Sec. 8, by the lead at the release); MINOR 2: none on eq:fullop): 42 new (33 \LEAN, 9 \LEANlit), 140 markers on 65 labels in all
%   (126 \LEAN, 14 \LEANlit; 98 on 48 at v1.19.0); every \LEANlit of slot 20 cites Shimura with "locator unverified (R20-L)";
%   sec:status names the list.  No statement's mathematics, label, tag or count changes.  Companions [E] 1.20.0, [U] 0.4.0,
%   [L37] 0.5.0, [EX] 1.4.0.
%
%   VERSION 1.20.0-pre.2 (hecke-extensions slot 20, the close delta for the Fable referee F1', 2026-10-04; a pre-release, NOT
%   owner-released).  Code-map text only: the sec:notimplemented box and the sec:involution code-map row now say "built" (WP-MEAS
%   f2897f74, NB-4 6218ebc6, WP-LV3 c1fda051, S1/S2 af3bf133 merged; iota over Z/ell^k in Hecke/HeckeInvolution and the console,
%   schema hecke_involution_ring/v1).  No statement, tag, label or count change.  Companions [E] 1.20.0-pre.6, [U] 0.4.0-pre.3,
%   [L37] 0.5.0-pre.5, [EX] 1.4.0.  List: docs/implementation/slot20/close_delta.md.
%
%   VERSION 1.20.0-pre.1 (hecke-extensions slot 20, WP-2, 2026-10-04; a pre-release, NOT owner-released; for the slot-20 Fable referee F1
%   (WP-3)).  Code-map text only: a v1.20.0-pre.1 box in sec:notimplemented (slot-20 WP-LV3, WP-MEAS, NB-4 planned, contract first in
%   [E] 1.20.0-pre.1 sec:lvmeasure / sec:lvdepth2 and [P]; WP-JET's Jetson patches value-neutral), and a "Planned at v1.20.0-pre.1"
%   sentence in the sec:involution code-map row (iota over Z/ell^k in Hecke/HeckeInvolution and the console; prop:iotasplit already
%   covers it).  Not in this pre-release: the milestone-M4 \LEAN markers (they follow slot-20 L20-6, the Lean audit F3).  No statement,
%   tag, label or count change.  Companions [E] 1.20.0-pre.1, [U] 0.4.0-pre.1, [L37] 0.5.0-pre.1, [EX] 1.4.0.
%
%   VERSION 1.19.0 (hecke-extensions slot 19, the close, WP-18, 2026-10-04: the RELEASE from v1.19.0-pre.6 under the owner's ruling O3,
%   on F1 / F1-delta ACCEPT WITH REPAIRS (repairs at pre.2 and at 166ae476)).  Over pre.6: the milestone-M3 \LEAN / \LEANlit markers,
%   clause by clause per docs/implementation/slot19/L19-6_F3_audit.md Sec. 8 as amended by L19-6_repairs.md Sec. 3 and L19-4_report.md
%   Sec. 3 (17 new markers, one \LEANlit -> \LEAN at cor:stablecert(2), the prop:model7h4(ii) marker rescoped to prop:model7h4-ii;
%   98 on 48 labels in all, 93 \LEAN + 5 \LEANlit); sec:status names the v1.19.0 list.  The Shimura locators stay marked unverified.
%   No statement's mathematics, label or app:status count changes.  Companions [E] 1.19.0, [U] 0.3.0, [L37] 0.4.0, [EX] 1.4.0.
%
%   VERSION 1.19.0-pre.6 (F1-delta m-C1, 2026-10-04; NOT owner-released): one sentence in err:e6, textual only.
%
%   VERSION 1.19.0-pre.5 (hecke-extensions slot 19, WP-2 continued, 2026-10-04; NOT owner-released): two corrections of fact and one
%   code-map box -- (i) Shimura's notation in sec:primfull, def:operators, thm:relations, sec:relationsmodel and err:e1: the
%   full operator is Shimura's T(p^nu), his T(1,p^nu) the primitive coset (docs/implementation/slot19/shimura_convention_check.md;
%   locators from memory, the book is not in the library); (ii) prop:model7h4's H4(ii) clause: acyclicity is proved in [GT] since L19-3
%   (model7_acyclic), not "stated" -- no LEAN marker moves (tagging at the close); (iii) a v1.19.0-pre.5 box in sec:notimplemented
%   (WP-BIL, WP-LIN built).  No statement's mathematics, tag, label or count changes.  Companions [E] 1.19.0-pre.6, [U] 0.3.0-pre.4,
%   [L37] 0.4.0-pre.3.
%
%   VERSION 1.19.0-pre.4 (hecke-extensions slot 19, WP-2 continued, 2026-10-04; NOT owner-released): code-map text only -- WP-LV2
%   (the twisted Eisenstein route on the sec:nebentypus character model) and WP-HL (the truncated Lambda-adic arm, HeckeLambdaValidation)
%   BUILT: the sec:nebentypus and sec:hidafull rows, the sec:saturation row (WP-LV2's deep rungs), a v1.19.0-pre.4 box in
%   sec:notimplemented.  No statement, tag, label or count change.  Companions [E] 1.19.0-pre.5, [U] 0.3.0-pre.4, [L37] 0.4.0-pre.2.
%
%   VERSION 1.19.0-pre.3 (hecke-extensions slot 19, WP-2 continued, 2026-10-04; NOT owner-released): code-map text only -- the slot-19
%   packages as built: the sec:saturation row (WP-LC and WP-LV untwisted, built and merged), the sec:nebentypus row (WP-LV twisted: not
%   built), the sec:hidafull row (WP-HL: not built; R19-2's prototype), and a v1.19.0-pre.3 box in sec:notimplemented (built: WP-LC,
%   WP-LV untwisted, WP-NF, NB-2).  No statement, tag, label or count change.  Companions [E] 1.19.0-pre.3, [U] 0.3.0-pre.3,
%   [L37] 0.4.0-pre.1, [EX] 1.4.0.
%
%   VERSION 1.19.0-pre.2 (hecke-extensions slot 19, WP-3 repairs, 2026-10-03; NOT owner-released): the slot-19 F1's ACCEPT with one
%   typographic note m-C1 -- the three "Planned at v1.19.0-pre.1" code-map sentences: no space before the full stop, each row's text
%   ends with a period; the WP-LV row's "stable block" reads "saturated block (stable at K=1)", following [E] 1.19.0-pre.2's M-1.
%   No statement, tag, label or count change.  Companions [E] 1.19.0-pre.2, [U] 0.3.0-pre.2, [EX] 1.4.0.
%
%   VERSION 1.19.0-pre.1 (hecke-extensions slot 19, WP-2, 2026-10-03; a pre-release, NOT owner-released; for the slot-19 Fable referee F1
%   (WP-3); release at the slot-19 close on the owner's yes).  Code-map text only (plan docs/implementation/slot19/00-slot19-plan.md S2
%   "WP-2" item (4)): a sec:notimplemented box for slot 19 (WP-NF, WP-LC, WP-LV, WP-HL planned, contract first in [E] 1.19.0-pre.1 /
%   [U] 0.3.0-pre.1; NB-2's --serve, spec [P] v3.10 by NB-1), the "Planned at v1.19.0-pre.1" sentences in the sec:codemap rows of
%   sec:hidafull, sec:nebentypus, sec:saturation.  No statement, tag, label or app:status count change; the M3 \LEAN markers wait for
%   slot-19 L19-6 / F3 (not in this pre-release).  Companions [E] 1.19.0-pre.1, [U] 0.3.0-pre.1, [EX] 1.4.0.
%
%   VERSION 1.18.0 (hecke-extensions slot 18, the close, WP-18, 2026-10-03: the RELEASE from v1.18.0-pre.2 under the owner's ruling O3 of
%   2026-10-03, on F1's ACCEPT WITH REPAIRS).  Over pre.2: the milestone-M2 \LEAN / \LEANlit markers, clause-scoped, exactly per
%   docs/implementation/slot18/L18-6_F3_audit.md Sec. 8 + L18-6_repairs.md (39 new: 36 \LEAN, 3 \LEANlit; 81 on 42 labels in all); the
%   thm:welldef(c) ERRATUM confirmed by F3 with a counterexample (Gamma_0(p), gamma_1 = S) and by the owner (2026-10-03, with the
%   slot-19 rulings); sec:status names the v1.18.0 list.  No statement, tag, label or app:status count change.  Companions [E] 1.18.0, [EX] 1.4.0, [U] 0.2.1.
%
%   VERSION 1.18.0-pre.2 (hecke-extensions slot 18, WP-3 repairs, 2026-10-03; a pre-release, NOT owner-released; release at the slot-18
%   close on the owner's yes, ruling O3).  Slot-18 referee F1 ACCEPT WITH REPAIRS (0 blocking): M-1 (IL5's route: prop:isomodelell(4),(5)
%   in place of rmk:borelfree, the hypothesis "Iso^(l)(Y') is an l-local isotropy model over Y'", (T1) over Gamma''), m-C1 (rmk:levelNscale(4)
%   figures per cell), m-C2 (IL5 at l = 2 a code-path check).  Folded in: WP-5 (generator order built; the structure-aware fill model FAILED
%   its x1.5 rule, hook_rss/v1 stays; windows pending), WP-6 (sec:nebentypus (iv)-(vi) built), WP-7 (sec:isomodelell (v), (vi) run; IL5
%   COMP^a line); TAG MOVE IL3, IL4: COMP^a partial -> COMP^a (F1 C-4 on WP-7's pass).  ERRATUM for owner confirmation: thm:welldef(c)'s
%   proof, representatives {gamma_1 s_j}, not {s_j gamma_1^{-1}} (L18-2, Lean thetaG_transported); conclusion unchanged.  def:pair states
%   aug o j = aug_bd (L18-4, eta_exists's hjaug); no statement change.  Companions: [E] 1.18.0-pre.2, [EX] 1.4.0-pre.2, [U] 0.2.1.
%   (WP-17, still pre.2, wording only: the erratum parenthesis names the cosets Gamma''_1 s_j gamma_1^{-1} = gamma_1(Gamma''
%   gamma_1^{-1} s_j)gamma_1^{-1} -- F3-18 MINOR 9.)
%
%   VERSION 1.18.0-pre.1 (hecke-extensions slot 18, package WP-2, 2026-10-03; a pre-release, NOT owner-released; for the slot-18 Fable
%   referee F1 = WP-3; release at the slot-18 close only on F1's accept and the owner's yes, ruling O3 of 2026-10-03).  Contract first
%   (rule 21) for slot-18 WP-5 (sec:levelN code map; rmk:levelNscale(4) "the structure-aware fill model", a cost statement), WP-6
%   (sec:nebentypus code map (iv)-(vi): W_K(F_q) at d >= 2; the model-route T_{p^2} = [diag(1,p^2)]_chi + chi(p)p^n against
%   thm:nebrelations(2); the WP-6 audit minors; rmk:nebcost(4) pointer) and WP-7 (sec:isomodelell code map (v), (vi): the strict lift on
%   a synthetic refined descriptor with its MP requirements; IL3 over F_9 and Z/9); NEW checkable item IL5 (rmk:isomodellcheckable,
%   PROVED given (T1)-(T3) of the built map; before the code, rule 20); sec:codemap rows; sec:notimplemented box.  F1 m-C4 of slot 17:
%   lem:heckerankone's Eichler--Shimura input = Diamond--Im Thm 12.2.2 p. 108 + Prop 12.4.10 p. 117 (VERIFIED, L17-5 repairs); Shimura
%   Ch. 8 kept, not verified.  No statement content, tag, label rename/removal or app:status count changes; one label added (IL5).
%   Companions: [E] 1.18.0-pre.1, [EX] 1.4.0-pre.1, [U] 0.2.1 (unchanged).
%
%   VERSION 1.17.0 (hecke-extensions slot 17, the close, WP-18, 2026-10-02: the RELEASE from v1.17.0-pre.5 under the owner's ruling O3
%   of 2026-10-01 ("release [C] v1.17.0, [E] 1.17.0, [U] 0.2.1, [EX] 1.3.0 on F1 accept, incl. the \LEAN tags").  The Fable referee F1
%   (docs/implementation/slot17/WP-3_referee_F1.md) was ACCEPT WITH REPAIRS for [C] (0 blocking, 1 major M-1, 5 minor); the repairs were
%   applied at pre.3 / pre.4 (with the owner's rulings of 2026-10-02 on M-1 and B-1).  Over pre.5 the only change is the FIRST USE of the
%   \LEAN / \LEANlit marker (Lean milestone M1): 42 clause-scoped markers on 20 labels, exactly per
%   docs/implementation/slot17/lean_taggable_labels.md (the F3 audit's list after its repairs; census decls=440 proved=383
%   proved_mod_lit=20 hyp_bad=0 verdict=OK), \LEANlit at cor:stablecert (modulo lem:heckerankone) and prop:hidaalgconst(1) (modulo
%   Hida's control theorem); sec:status says a definition may carry the marker.  No statement, tag, label or count change.
%   Companions: [E] 1.17.0, [U] 0.2.1, [EX] 1.3.0.
%
%   VERSION 1.17.0-pre.5 (hecke-extensions slot 17, package WP-2, 2026-10-02; a pre-release, NOT owner-released): pointers only.
%   (1) Slot-17 R17-4's verified locators: (NB4)(c)'s STD input is Ash--Stevens Lemma 1.1.1 p. 194 (LES; not 1.1.3) and Def. 1.1.2 /
%   Lemma 1.1.4 pp. 194-195 (Shapiro), VERIFIED -- the "not verified" flag is dropped, the PROVED tag stands, the COMP^r fallback is
%   not triggered; compatibility (b) stated in A-S's orientation; prop:barformula's "in exactly this form" and eq:corconjres's "in this
%   exact form" narrowed (A-S: right module, homogeneous cochains); rmk:thetaclassical (3) narrowed to A-S Lemma 2.1 p. 204 for the
%   inclusion only, (2) Jochnowitz Thm 4.1 p. 276; lem:heckerankone's eigenbasis from Diamond--Im S3.5-3.6 pp. 54-55 (verified; new bib
%   key diamond_im), Shimura / Diamond--Shurman kept as unverified secondaries; LN5 from Merel Theoreme 2 p. 72 (verified) with its
%   hypotheses q > 3, q | (N-1)/gcd(N-1,12); the generic A-S S1 sites near thm:relations. (2) sec:isomodelell: IL3 / IL4 status
%   sentences and the code map after slot-17 WP-7 (run, audited ACCEPT; IL3N against the pipeline over F_3 only; arm REF a descriptor
%   round trip); the IL3 / IL4 tags stay COMP^a partial (the move is a proposal for the next referee). (3) sec:status: cor:stablecert
%   will take \LEANlit, not \LEAN (the list in docs/implementation/slot17/lean_taggable_labels.md); no \LEAN tag is used.  No
%   statement, tag, label or count change.  Companion: [E] 1.17.0-pre.5.
%
%   VERSION 1.17.0-pre.4 (hecke-extensions slot 17, package WP-2, 2026-10-02; a pre-release, NOT owner-released): two owner rulings
%   of 2026-10-02 on the slot-17 referee F1.  (1) M-1: (NB4)(c)'s sharp form at n = 0 tagged PROVED (given the one STD input, its
%   locator Ash--Stevens S1.1 flagged NOT verified -- the lem:heckerankone pattern); rmk:nebcheckable's tag set unchanged, no count
%   moves.  (2) B-1 ([E]): (R) restated on congruence-free classes; the two dependent sentences here (LN6's route-R3 sentence and
%   sec:saturation's scope paragraph) take their final wording.  No label added, renamed or removed.  Companion: [E] 1.17.0-pre.4.
%
%   VERSION 1.17.0-pre.3 (hecke-extensions slot 17, package WP-2 resumed, 2026-10-02; a pre-release, NOT owner-released): the
%   repairs of the slot-17 Fable referee F1 (docs/implementation/slot17/WP-3_referee_F1.md; [C] ACCEPT WITH REPAIRS, 0 blocking,
%   1 major M-1, 5 minor; WP-6's code-map text M-2 + m-W1..m-W3).  M-1: (NB4)(c) gains the orientation proof in this document's
%   conventions (0 -> Sym^{l-1-i} (x) det^i -> Ind -> Sym^i -> 0) and the locator [S1.1] on the one remaining Ash--Stevens STD
%   input; that locator is NOT verified (no copy in the library), so the sharp form at n = 0 is re-tagged CONJ -> COMP^r (the
%   referee's alternative), PROVED on verification.  M-2: sec:nebentypus's code map records WP-6 as built (merged, audited), with
%   T_{p^nu} from Mer_{p^nu} TESTED, NOT PROVED, and the model-route double coset not built; sec:codemap row, sec:notimplemented box.
%   WP-5: sec:levelN's code map records the int64-then-int128 coset products (merged, audited; WP-5 audit MINOR-1's wording on the
%   predictions), rmk:levelNscale(4) gains the N = 99330 data point (an open measurement, a failed prediction).  Minors m-C1, m-C3
%   (two sites citing (R), and the pre.2 changelog sentence corrected), m-C4 (Diamond--Shurman Thm 5.5.4, not verified), m-C5 (the
%   -I -> -1 control in sec:isomodelell (iv)); m-C2 a carry for a code package.  No label added, renamed or removed; no tag set of a
%   numbered statement changes, no count moves.  Companion pointer: [E] 1.17.0-pre.3 in the tree.
%
%   VERSION 1.17.0-pre.2 (hecke-extensions slot 17, package WP-2, second pass, 2026-10-02; a pre-release, NOT owner-released): over
%   pre.1 ONE citation repair only -- lem:heckerankone's proof cites the eigenbasis from Shimura Thm 3.41 or Diamond--Shurman S5.5
%   (both NOT verified: neither book is in the library) and Serre, Cours, Ch. VII S5 for multiplicity one only (verified locators);
%   the app:status row follows; bib entry diamond_shurman added.  No statement, tag, label or count changes.  Companion pointer:
%   [E] 1.17.0-pre.2 in the tree (its pr:latticesplit (R) restated again, owner ruling O-R17-2a of 2026-10-02).  For the same
%   referee F1 (WP-3).
%
%   VERSION 1.17.0-pre.1 (hecke-extensions slot 17, package WP-2, drafted 2026-10-01 PRE-START during the slot-16 close, NOT
%   owner-released; plan docs/implementation/slot17/00-slot17-plan.md S2 "WP-2"; owner ruling O3 of 2026-10-01: release at the
%   slot-17 close only on the Fable referee F1 (WP-3) accepting and the owner's yes).  The deferred items of
%   docs/implementation/slot16/contract_release_report.md S4, drafted once for one referee:
%   (a) NEW rmk:sattor in a NEW last subsection sec:sattorsec of sec:saturation (so no v1.16.0 number moves; slot-16 WP-5's Tor identity: I = Tor_1^{Z_l}(T^sat/T, F_l) = (T^sat/T)[l]
%   = (T cap l T^sat)/l T = ker(T/l -> End(L/l)); dim I = eps_1; the converse of prop:saturation(2)'s sufficient condition fails,
%   with an example); PROVED.  (b) rmk:nebcheckable (NB4): the n = 0 sharp form recorded as COMP^r (132/132, l <= 43, with the
%   orientation it pins: Sym^i untwisted), the n > 0 reading (Clebsch--Gordan, the H^0 correction at i = n; 240/240, reported), and
%   a PROPOSED proof of the n = 0 sharp form by the Jordan--Hoelder long exact sequence -- the referee decides between PROVED (from
%   the named STD inputs) and COMP^r only; the general (NB4) statement stays CONJ.  (c) rmk:levelNscale(4): slot-16 WP-13 F-3, the
%   fitted fill law underestimates final_nnz at composite N, k > 2, by 1.5--6.8x (COMP^a; a cost statement only).  (d) Code maps,
%   contract first for slot 17 (rule 21): sec:levelN (the cocycle width beyond N ~ 5.5e4, WP-5), sec:nebentypus (the sparse route
%   with a character and T_{p^nu}, nu >= 2, WP-6: scope only, already in thm:nebrelations(2) / rmk:nebcost), sec:isomodelell (what
%   WP-7 will run; what the MP constraint requires of a synthetic refined descriptor), sec:saturation (slot-16 WP-5 built).
%   (e) The \LEAN / \LEANlit macros (Lean programme Q9, owner-ruled 2026-10-01), DEFINED BUT UNUSED until the slot-17 close.
%   (f) lem:heckerankone's proof: the eigenbasis of M_k cited from Petersson / Shimura Ch. 3 (VERIFY), Serre Cours VII S5 only for
%   multiplicity one, with verified page locators (slot-17 L17-5's finding, relayed by the lead); no statement or tag changes.
%   One numbered statement added (rmk:sattor), none removed or renamed; tag sets: rmk:nebcheckable gains COMP (registered),
%   rmk:levelNscale gains COMP (ad hoc).  Companion pointer: [E] v1.16.3 (1.17.0-pre.1 in the tree).  Report
%   docs/implementation/slot17/WP-2_report.md; referee brief docs/implementation/slot17/WP-3_referee_brief.md.
%
%   VERSION 1.16.0 (hecke-extensions slot 16, the close, 2026-10-01: the RELEASE from v1.16.0-pre.3 under the owner's ruling O2
%   of 2026-10-01 ("release [C] v1.16.0 and [U] 0.2.0 since WP-14 accepted and the repairs are confirmed").  The Fable referee
%   docs/implementation/slot16/WP-14_referee.md was ACCEPT WITH REPAIRS (0 blocking); every repair applied at pre.3 and CONFIRMED
%   by docs/implementation/slot16/WP-14_repair_confirmation.md (16 / 16).  Over pre.3, no statement and no tag changes:
%   (a) locators LN5, LN6 verified on the arXiv text of Wake--Wang-Erickson (slot-16 WP-8; the CONJ tags stay); wwe's bib note;
%   (b) sec:isomodelell: the proof of prop:isomodelell(1) says why the e_sigma average is a no-op on d (stabilisers {+-I},
%   central) and where it is load-bearing (Psi_ell); (IL3) and (IL4) marked COMP^a partial (level N; Z/3^e and the pair/cone
%   half not run); code maps "built": sec:isomodelell at ell = 3 (WP-16), sec:levelNdegsec (WP-13), sec:nebentypus (WP-10,
%   WP-11), with the not-implemented boxes; (c) L16-2 audit note 5 / referee C-9 verified as applied at pre.3.  No label added,
%   renamed or removed; app:status unchanged.  Companion pointer -> [E] v1.16.2.  Report
%   docs/implementation/slot16/contract_release_report.md.
%
%   VERSION 1.16.0-pre.3 (hecke-extensions slot 16, WP-12 resumed, 2026-10-01, NOT owner-released): the repairs answering the
%   Fable referee docs/implementation/slot16/WP-14_referee.md (ACCEPT WITH REPAIRS; 0 blocking, 1 major, 8 minor), every one applied:
%   C-1 prop:isomodelell (1),(4) lifts averaged by e_sigma (proof, code map, (IL1)); C-2 kappa^mod on a refinement; C-3 the
%   identification of prop:lhcompliant(4) named; C-4 (LN8)/(LN9) wording, the torsion-freeness behind p > k+1; C-5 the admissibility
%   count 2 m_N/m_d; C-6 SparseCokernel's final_nnz; C-7 "|P_sigma| > 2" at ell = 2; C-8 H (+) H^anti in (NB5), the char-0 -> F_q and
%   d(chi) = 1 clauses; C-9 (NB5)'s dictionary re-tagged PROVED modulo (L1)-(L5) citing [GT, bridge:nebentypus-dictionary] (owner).
%   No label added, renamed or removed; no tag set of a numbered statement changes, so no count moves.
%
%   VERSION 1.16.0-pre.2 (hecke-extensions slot 16, WP-12 follow-up, 2026-10-01, NOT owner-released): rmk:nebcheckable (NB5)
%   repaired per the slot-16 WP-9 audit (docs/implementation/slot16/WP-9_audit.md S3): the cuspidal factor on the chi-part is
%   prod_{f in S_k(N,chi)} (X - a_p(f))^2, not (X - a_p)(X - conj a_p) (wrong for non-real chi: 20 / 24 cases); proof via the
%   involution of epsilon = diag(1,-1) swapping the Eichler--Shimura halves; the reduction rule psi(g) = s (least primitive
%   root) stated.  No label added, renamed or removed; no count moves.
%
%   VERSION 1.16.0-pre.1 (hecke-extensions slot 16, package WP-12, 2026-10-01: the pre-release of an X bump, the contract
%   first; plan docs/implementation/slot16/00-slot16-plan.md S2 "WP-12"; owner rulings of 2026-10-01 (O1-O4)).  NOT owner-released;
%   the Fable referee is slot-16 WP-14.  (a) NEW subsection sec:levelNdegsec at the end of sec:levelN (no number moves):
%   prop:kappamodel (kappa_N read on the fixed model: f -> -f(e_46), the pinned model -> presentation map, naturality for fibre
%   maps, the entrywise operator comparison), prop:degenpres (alpha_d = [Gamma_0(M) diag(1,d) Gamma_0(N)], every d with dM | N, on
%   the cokernel presentation by Merel's set of determinant d, eq:degenpres, PROVED equal to prop:levelcorrespondence(4)(c) on H^1;
%   d = 1 is alpha^0; composite d; the trace; with a character), rmk:degenpresconv (traps, COMP^a evidence: wp12/degen.py,
%   negctl.py), rmk:levelNscale ((c): scope at k > 2 and composite N, the memory model), rmk:levelNcheckabletwo ((LN7)-(LN9)).
%   v1.15.0-pre.3's unrefereed scope sentence is superseded (moved into prop:kappamodel(3)).  (b) NEW chapter sec:isomodelell
%   after sec:saturation, before the appendices: def:isomodelell, prop:isomodelell, prop:ellthreeclosed (the closed form at
%   ell = 3 and the exact comparison with the shipped Bredon route), rmk:isomodellmp (the MP constraint), rmk:isomodellcheckable
%   ((IL1)-(IL4)).  Ten numbered statements added, none removed or renamed, no tag moved.  gl2_refs.bib: billerey_menares (not
%   verified).  Companion pointer stays at [E] v1.16.1.  Report docs/implementation/slot16/WP-12_report.md.
%
%   VERSION 1.15.0 (hecke-extensions slot 15, WP-18, 2026-10-01: the RELEASE at the slot-15 close from v1.15.0-pre.3,
%   under the owner's ruling Q2 of 2026-09-30 ("release [C] v1.15.0 at the close if WP-4 accepts").  The Fable referee
%   docs/implementation/slot15/WP-4_referee.md: sec:nebentypus and cor:stablecert ACCEPT WITH REPAIRS, P15/P17/P18/P20 and
%   prop:cohforms ACCEPT, 0 blocking, all repairs applied at pre.2.  sec:nebentypus released unbuilt (code: slot 16, Q3).
%   pre.3's one sentence in sec:levelN's scope (slot-15 WP-17, proposed by WP-8: fibre maps on the cokernel presentation by
%   naturality; alpha_d, d > 1, on the model only) is released UNREFEREED and listed as such in the changelog.  Code map of
%   sec:saturation: the certified stable verdict and the digit route built (slot-15 WP-5).  No statement added or removed, no
%   label renamed, no tag moved.  Companion pointer stays at [E] v1.15.0 ([E] 1.15.1-pre.6 / [EX] 1.2.0-pre.13 not released:
%   pending the owner's ruling on the pr:latticesplit edit).
%
%   VERSION 1.15.0-pre.3 (hecke-extensions slot 15, WP-17, 2026-09-30: the one sentence slot-15 WP-8 proposed
%   (docs/implementation/slot15/WP-8_report.md sec. 7 item 3): in sec:levelN's scope paragraph, a Gamma-module map between
%   induced modules induces the map on the cokernel presentations (naturality of prop:levelorbit (2) in the coefficient
%   module), so the presentation-level alpha^0 and omega rest on a [C] statement, and alpha_d (d > 1) is asserted on the
%   model only.  Code map: the sparse cokernel elimination and the composite-N degeneracy maps (on the model) built
%   (slot-15 WP-8, Hecke/HeckeLevelNSparse); R3 at scale (WP-9, Hecke/HeckeLevelNR3).  No statement added or removed, no
%   label renamed, no tag moved.
%
%   VERSION 1.15.0-pre.2 (hecke-extensions slot 15, WP-3 resumed, 2026-09-30: the repairs answering the Fable referee
%   docs/implementation/slot15/WP-4_referee.md -- sec:nebentypus ACCEPT WITH REPAIRS, P15/P17/P18/P20/prop:cohforms ACCEPT,
%   cor:stablecert ACCEPT WITH REPAIRS, 0 blocking.  A1 prop:nebcorr(4) proved in-document (the bijection
%   A_1''\Gamma_1(N) -> A_0''\Gamma_0(N)); A2 parity clause in rmk:nebheilbronn; A4 (NB4)'s sharp form with the two
%   constituents Sym^i, Sym^{l-1-i} (x) det^i; A5 typo; A9 Galois-ring remark; N1 per-scalar eigenspace cost; B1
%   lem:heckerankone for any idempotent, cor:stablecert for a maximal ideal of T' = Z_l[T_p : p in S]; B4 scope sentences
%   (cl:mkfaithful, pr:latticesplit refuted at l = 37).  Code map: H_{N,p} built (slot-15 WP-7), HeckeHidaComponents built
%   (WP-6).  No statement added or removed, no label renamed, no tag moved.
%
%   VERSION 1.15.0-pre.1 (hecke-extensions slot 15, package WP-3, 2026-09-30: the pre-release of an X bump, the
%   contract first; plan docs/implementation/slot15/00-slot15-plan.md S2 "WP-3"; owner rulings Q2/Q3/Q4 of 2026-09-30).
%   NEW chapter sec:nebentypus (Gamma_1(N) and (Gamma_0(N),chi) for every N, chi:(Z/N)^x -> F_q^x, q = ell^d, and the
%   Teichmueller lift over W_K(F_q); the character model is the monomial level-N model with a root of unity per block;
%   diamonds as fibre maps; T_p, U_p, T(p,p) on the level-1 tau-table; orbit-wise fixed spaces, parity, cokernel,
%   Euler characteristic; the Heilbronn action with a character; predictions (NB1)-(NB6); NOT BUILT, slot 16).
%   NEW chapter sec:saturation (P15 = [M1] Prop. 1: prop:saturation, lem:freepart, cor:stablecert -- the stopping rule
%   dim A_K/ell = rank T is equivalent to A_K/ell = T/ell; P17: prop:satdefect; lem:heckerankone; prop:cohforms, the
%   torsion-aware cohomology -> forms comparison; P18: prop:hidaalgconst; P20: rmk:thetaforms).  Twenty-three numbered
%   statements added, none removed or renamed, no tag moved; thm:hidaweight and sec:levelN's scope gain one pointer each.
%   gl2_refs.bib: cohen_oesterle, ribet77 (not verified).  Report docs/implementation/slot15/WP-3_report.md.
%
%   VERSION 1.14.2 (hecke-extensions slot 14, 2026-09-29: Z-bump, owner ruling Q-F: the "not built -> built" flip of
%   sec:levelN's code maps (WP-12, 7a5158d; Fable review 64e6b9c, WP-12_code_review.md sec.5): Hecke/HeckeLevelN built,
%   O3 closed, O2 open; (LN3) tau-route half and (LN6) second-route COMP^a sentences; the r_t alternative in the box
%   after lem:cosetrepsN; nu_2, nu_3 checked.  No statement changed, no label added/renamed/removed, no tag moved.)
%
%   VERSION 1.14.1 (hecke-extensions slot 14, 2026-09-29: owner rulings on the Fable audit
%   docs/implementation/slot14/AUDIT_contract_currency_2026-09-29.md.  Q-B: the second-cluster-run data (job 3467349,
%   e85e25a) signed off -> WP-L1's L1-L4 as COMP^a sentences: Hi11 held at n = 4150, Hi12 held at (4830,2411),
%   (2420,2411) recomputed with the levers (19.2 GiB), probe peaks and sampler values stated apart, the admission
%   predictor's refit (549200e) named, levelO5's pre-lever fit not restated.  Q-H: NEW lem:rhorank (audit C4) after
%   sec:hidafull's code map -- rho_n(g), l | det g, has <= min(k, n+1) non-zero invariant factors over Z/l^k, rank-one
%   closed form over F_l; the claim HeckeLowRank.h ENGINE_CHECKs per block.  Q-F: WP-12's "not built -> built" flip
%   NOT made (made at v1.14.2).  No existing statement changed, no label renamed or removed, no tag moved.
%
%   VERSION 1.14.0 (hecke-extensions slot 14, package WP-13, 2026-09-28: the OWNER-SIGNED RELEASE of sec:levelN
%   from v1.14.0-pre.2, owner decision 2026-09-28 "Go for publication and release" ([S14] Q2 and Q3 at the recommended
%   options).  The WP-5 referee was ACCEPT WITH REPAIRS, 0 blocking, all 11 findings applied at pre.2.  sec:levelN is
%   released unbuilt (contract first; the code maps say "not built at v1.14.0").  Q3's routing of [S14] D1 to this
%   document: P1 (Hi12 held at (2420,2411)), P2 + P8 (Hi11 at n = 1390, 2080, 2770; (700,691) at m = 2) and P3
%   (rmk:bocksteinexp(3), the Z/l^2 depth window) re-inserted verbatim from pre.1 (cc97dd7), each COMP^a sentence
%   citing the now-published docs/matrix-engine/results/noctua2_run1/ (Q2); P7 (493 GiB) back to its pre.1 form.
%   The sec:codemap row and the v1.13.0 scope box say Hi12 ran and held.  Later proposals deferred to slot 15 as new
%   statements: P14, P15 (= P-WP6-1/2), P17, P18 (docs/implementation/slot14/WP-13_release_note.md).  No statement
%   added, no label renamed or removed, no tag moved.  Companion pointer -> [E] v1.15.0.
%
%   VERSION 1.14.0-pre.2 (hecke-extensions slot 14, package WP-5, 2026-09-28: the repairs answering the Fable
%   referee docs/implementation/slot14/WP-5_referee.md, ACCEPT WITH REPAIRS 0 blocking / 3 should-fix / 8 minor,
%   all applied).  F1 rmk:heilbronncost(3),(4),(5): a Hecke MATRIX on H^1 costs O(|Mer_p| d (dim H^1)^2) through the
%   D_N x dim H^1 normal form of Rel_N (quadratic in m_N); O1 loses the N_q^2 factor, not the quadratic order.  F2
%   (LN4)/(LN6): W_{N,p} computed jointly for q <= 7 only (contains LN4's W; equality predicted), e read off dim W.
%   F3 def:levelellmodel -> def:cosetrepsN.  F4 N = 1 in def:cosetrepsN(3); F5 case (iii) a bounded loop; F6 the new
%   symbols listed (rule A5); F7 the Brown pointer definitional; F8 (LN3)'s tau-route half is WP-12's, the referee's
%   seven cells added; F9 sum x_0 = p^2+1 spelled out, |Mer_p|/(p log p) increasing; F10 checked, no change.  F11:
%   [S14] Q2's default stands -- P1, P2, P3, P8 taken back out (Hi11, Hi12, rmk:bocksteinexp(3), the sec:codemap
%   row, sec:hidafull's code map); P7 kept.  No statement added, no label renamed or removed, no tag moved.
%   Report docs/implementation/slot14/WP-4_report.md "pre.2 repairs".
%
%   VERSION 1.14.0-pre.1 (hecke-extensions slot 14, package WP-4, 2026-09-28: the pre-release of an X bump,
%   the contract first; plan docs/implementation/slot14/00-slot14-plan.md S2 "WP-4").  NEW chapter sec:levelN after
%   sec:hidafull (no chapter number moves): def:cosetrepsN + lem:cosetrepsN (O3: the binding coset table for every N,
%   Definition cosetreps verbatim at prime powers, the B''\B representatives read off the level-N' table), def:levelNmodel
%   (every N, every ring; ell | N, ell not | N, the R3 setting F_p, p not | N), lem:treeres, lem:manin (Manin's exact
%   sequence, proved via the Farey complex), def:kappaN (the cokernel isomorphism of prop:levelorbit(2) pinned),
%   lem:levelduality (cokernel <-> cusp-divisor coinvariants, natural in res, cor, conj), lem:merelpath (Merel's set and
%   its path condition), prop:heilbronn (O2: T_p, p not | N, and U_p, p | N, on the cokernel presentation by Merel's
%   matrices in this document's conventions -- rho_n(adj h_j(eps)) on block j -> j.eps -- PROVED equal to
%   prop:levelcorrespondence(4)(d) on H^1), rmk:heilbronnconv, rmk:heilbronncost, rmk:levelNcheckable ((LN1)-(LN6)).
%   Twelve numbered statements added, none removed or renamed.  rmk:levelscaling: O2, O3 gain a status sentence.
%   Data remarks of [S14] D1 (P1, P2, P3, P7, P8), each marked "publication pending the owner's sign-off, [S14] Q2":
%   Hi11 and Hi12 (COMP^a sentences), rmk:bocksteinexp (3) (the Z/l^2 depth window), sec:hidafull's code map (the
%   493 GiB provenance).  Companion pointer -> [E] v1.14.0.  gl2_refs.bib: merel96, wwe added (not verified).
%   Report docs/implementation/slot14/WP-4_report.md.
%
%   VERSION 1.13.0 (hecke-extensions slot 13, the close, package WP-15, 2026-09-27: the OWNER-SIGNED
%   RELEASE of sec:hidafull from v1.13.0-pre.2, owner decision Q14: released because the WP-3 referee
%   was ACCEPT with 0 blocking (0B/2M/10m, all applied) and slot-13 WP-6's arms held (eq:Ulevell
%   against the independent group-cocycle route, prop:oldspace, prop:stabilisation, (Hi10) 92/92).
%   prop:stabilisation is released with its per-item tags as at pre.2.  One normative sentence
%   amended, as the WP-13 code audit (A3) required: def:hidalattice's variant (H) no longer says
%   "nothing else changes" -- on Hida's lattice the module route's Howell pivots of the (free)
%   cochain module may be non-units, and the route then takes a Nakayama basis (slot-13 WP-5),
%   value-preserving wherever it computed before.  Code maps: sec:hidafull and sec:codemap now
%   read "built" (HeckeLevelEll, HeckeHidaTower part 2, HeckeHidaHeavyCellProbe); (Hi11) HELD,
%   (Hi12) not run (deferred to the cluster roadmap).  No statement added, no label renamed or
%   removed, no tag moved.  Companion pointer unchanged ([E] v1.13.1; its v1.14.0-pre.1..pre.3 are
%   slot-13 pre-releases, unreleased).  Report docs/implementation/slot13/WP-15_report.md.)
%
%   VERSION 1.13.0-pre.2 (hecke-extensions slot 13, package WP-3, 2026-09-26: the repairs answering
%   the Fable referee docs/implementation/slot13/WP-3_referee.md, ACCEPT WITH REPAIRS 0B/2M/10m, all
%   twelve applied; response docs/implementation/slot13/WP-3_response.md).  M1: Hida's lattice L_n is
%   rho_n conjugated by D = diag(binom(n,j)) (rho^H_n(A) = rho_n(A^T)^T = D^-1 rho_n(A) D), not Sym^n up
%   to basis order -- rmk:stabconvention(3); glosses in thm:hidacontrol / thm:hidaweight; def:hidalattice
%   names its lattice and gains variant (H) on Hida's lattice; (Hi8) gets arms (a) V, (b) L_n, (c) the
%   D_* comparison, and its source is restated as an extrapolation.  M2: rmk:levelcost counts the
%   cubic projector on the built dense route; NEW rmk:levelscaling -- the two regimes, the obstacles
%   (O1)-(O6), the Manin--Merel Heilbronn-matrix action on the cokernel named as the specified-by-
%   pointer, unbuilt scalable route.  Minors: per-item tags in prop:stabilisation (m1); composite-N
%   cosets and composite degeneracies (m2); exact nu_2, nu_3 (m3); (Hi8) a decision procedure (m4);
%   Hida's p >= 5, N = 1 (m5); COMP^a scope (m6, m10); Scope at ell in {2,3} (m7); the U^* old-space
%   matrix and no trace inverse in Hida's convention, rmk:stabconvention(4) (m8); (Hi10)'s
%   commutativity STD (m9).  One numbered statement added (rmk:levelscaling), none removed or renamed.
%
%   VERSION 1.13.0-pre.1 (hecke-extensions slot 13, package WP-2, 2026-09-26: the pre-release of an
%   X bump, the contract first).  NEW chapter sec:hidafull after sec:hida (no chapter number moves):
%   level structure by Shapiro on the fixed model (def:levelellmodel, prop:levelcorrespondence --
%   every level-N Hecke operator, degeneracy map and level-change fibre map as one block formula on
%   the induced fibre reusing the level-1 tau-table; prop:levelellU -- U_ell at level ell as the block
%   row varpi, no new coset or tau table; prop:oldspace -- the old-space relations; prop:levelorbit --
%   Mackey orbit-wise fixed points and H^1 as a cokernel; rmk:levelcost -- the cost model), the
%   stabilisation (lem:levelelltorsion, prop:stabilisation, rmk:stabconvention), Hida's control and
%   weight independence as literature with VERIFIED locators (thm:hidacontrol, thm:hidaweight), the
%   finite-level lattice datum (def:hidalattice) and the predictions (Hi7)--(Hi12)
%   (rmk:hidacheckable2).  Thirteen numbered statements added, none removed or renamed; (Hi3)/(Hi4)
%   gain a locator note, their tags unchanged.  gl2_refs.bib: atkin_lehner added, hida_ens /
%   hida_invent notes record the verified locators.  Companion pointer unchanged ([E] v1.13.1).
%   Report docs/implementation/slot13/WP-2_report.md.)
%
%   VERSION 1.12.1 (hecke-extensions slot 12, post-close, package S12-census2, 2026-09-26: a
%   NON-NORMATIVE Y bump -- code maps only.  sec:hida's code map records what slot-12 WP-4 / WP-5
%   built (Hecke/HeckeOrdinaryPart, HeckeExt/HeckeHidaTower, their validations) in place of "not
%   built"; sec:codemap's sec:hida row likewise; the sec:boundarydatum and sec:dicksontheta rows
%   name HeckeExt/HeckeFinalityRoutes, the finality certificates (d0) "dickson" / (d1) "boundary"
%   of [E] v1.13.x sec:certgen that rest on prop:heckeH0 and prop:bdrytranscription /
%   lem:bdrytorus; the companion pointer moves to [E] v1.13.1.  No statement, label, tag or
%   hypothesis changes.  Report docs/implementation/slot12/census2_report.md.)
%
%   VERSION 1.12.0 (hecke-extensions slot 12, post-close, the owner-signed release, 2026-09-26,
%   from v1.12.0-pre.3: publishes sec:hida (WP-2, refereed WP-3) and the Shapiro transcription
%   in degree 1 (prop:bdrytranscription, lem:bdrytorus; WP-16, refereed WP-16_referee.md), both
%   accepted after repairs, all applied, none outstanding.  The one status change this version
%   makes: open:ordprecision is ANSWERED at ell >= 5, n >= 1 by cor:ordprecision (with
%   prop:torsionprepass and rmk:ordtorsionfree), so its tag moves OPENSTAT -> PROVED and the item
%   is retitled "-- ANSWERED 2026-09-26", exactly as open:ell2 lost its [open] tag at v1.3.0.
%   No statement added, no label renamed or removed, one tag moves; the question text and the
%   boxed statement of open:ordprecision are unchanged (rules L1, L2).  app:status: OPENSTAT
%   7 -> 6, PROVED 129 -> 130.  Companion pointer unchanged ([E] v1.12.1).  Report: this release
%   note (owner sign-off 2026-09-26; docs/implementation/slot12/WP-2_report.md, WP-3_referee.md,
%   WP-16_report.md and WP-16_referee.md record the refereed content).)
%
%   VERSION 1.12.0-pre.3 (hecke-extensions slot 12, package WP-16, post-close, owner-approved
%   2026-09-26: the WP-13 referee's F1 -- "Proposition T" proved).  NEW prop:bdrytranscription
%   (PROVED -- the Shapiro transcription in degree 1: Phi an isomorphism over every ring at
%   rho(-I) = id; at every nu and every p, Phi o Theta^bdry is a sum over the left cosets
%   B_N/Gamma_bdry (eq:bdrygroupformnu); at nu = 1 it is eq:bdrygroupform.  No p in R^x needed:
%   the correction terms telescope under j -> j.T, so p = ell is included) and lem:bdrytorus
%   (PROVED -- "Lemma B'": away from ell the boundary operators factor through the torus,
%   diagonalisable, commuting, p mod ell only; weight rule chi_k in degree 0, chi_{n-k} in degree 1).
%   Status notes: prop:bdrychar(3),(4), obl:taubdry(v),(v.b) (corroboration of a theorem),
%   cl:chi(2) (semisimplicity proved, multiset CONJ), rmk:bdrynugeneral, the OP1 decision box.
%   No label renamed or removed, no tag moved; no [E] finality certificate added.
%   Report docs/implementation/slot12/WP-16_report.md.)
%
%   VERSION 1.12.0-pre.2 (hecke-extensions slot 12, package WP-3, second half: the Fable referee's
%   (docs/implementation/slot12/WP-3_referee.md; 0 B / 0 M / 10 m, R1-R10, all applied) repairs.
%   R1: prop:torsionprepass(6) -- H^1(C_eq(Z)) agrees with group cohomology only after inverting 6;
%   R2: (2)'s ell >= 5 hedge explained (fails at (2,2)); R3: cor:ordprecision(1) pins eps as Hida's
%   projector lim U^{j!}, and e = e' (not max(1,e')); R4: h = 2d_k - 1 known before the run; R5:
%   rmk:ordexponentuse(1) gains COMP^a (exponent of GL_g(Z/ell^K) = N^lcm on eleven cases); R6: the
%   F'-block sentence reworded (e_ord(U_F') exists, canonical only mod ell^e'); R7: (Hi2) on named
%   modules; R8: (Hi3)'s weight-2 step explicit, classical COMP^a evidence; R9: (Hi5) simultaneous
%   eigenline; R10: (Hi6) Bezout step named, tau(2411) re-derived.  No label renamed or removed, no
%   statement added; one COMP^a tag added.  Report docs/implementation/slot12/WP-2_report.md sec. 6.)
%
%   VERSION 1.12.0-pre.1 (hecke-extensions slot 12, package WP-2: the pre-release of the slot's
%   X bump, "the contract first" -- written before any slot-12 package commits code for it
%   (slot-12 plan docs/implementation/slot12/00-slot12-plan.md S2 "WP-2", rules 20 and 21).
%   NEW sec:hida (Part VII, after sec:ordinary, so no chapter number moves): prop:torsionprepass
%   (PROVED -- H^0(Gamma; V/ell^e) = H^1(Gamma; V)[ell^e] by the e-fold Bockstein, at EVERY prime
%   for group cohomology (the plan said ell >= 5; only the use as a model precision needs it);
%   the lengths, the stopping rule for e_max, the l-adic elementary divisors of the stacked
%   [rho(S)-1; rho(T)-1] ARE the torsion exponents, Hecke equivariance, the integral torsion);
%   cor:ordprecision (PROVED -- at ell >= 5, n >= 1: the ORDINARY part needs no precision beyond
%   the target e' (torsion is non-ordinary), the free quotient needs K = e' + e_max, known before
%   the run; WP-18c's free-block reading justified); lem:ordexponent (PROVED -- U^N = e_ord for
%   N = N_ord(g, K) = ell^{K-1+ceil(log_ell g)} prod_{f<=g}(ell^f - 1), and for every multiple of
%   the lcm variant; g = generator count, the plan's r); rmk:ordexponentuse (PROVED, sharpness and
%   application); def:hidatower; rmk:hidacheckable (the predictions (Hi1)-(Hi6) = the plan's
%   (H1)-(H6), committed before WP-4 / WP-5's code; Hi3, Hi4 CONJ from Hida's theorems, locators
%   unverified).  open:ordprecision: a pre.1 status note -- answered at ell >= 5, n >= 1, the
%   OPENSTAT macro kept until the owner-signed release (plan S5 L3).  Three bibliography entries
%   (hida_ens, hida_invent, hida_book).  app:status: 249 -> 255 lines.  Companion pointer
%   unchanged ([E] v1.12.1).  Report docs/implementation/slot12/WP-2_report.md.)
%
%   VERSION 1.11.1 (hecke-extensions slot 11, package WP-17, the slot close: a Y bump, the
%   WP-15 Fable referee's repairs (docs/implementation/slot11/WP-15_referee.md; 0 B / 2 M / 6 m).
%   M-1: rmk:thetacheckable's COMP^r provenance sentence restated to the reports (WP-5's rule-20
%   deviation, WP-6's C5 / C7 re-alignment after the first runs).  F-4 / m-2: a new remark
%   rmk:ordtorsionfree (PROVED) -- at ell >= 5, n >= 1 the torsion of H^1 (x) Z_ell is
%   non-ordinary, so the ordinary part is torsion-free at level 1 and equals the free quotient's;
%   open:ordprecision narrowed by a status note to the per-cell bound (stays OPEN).  m-3: Part VII's
%   code map sentence marked "(through v1.11.0-pre.2)".  m-4: the companion pointer's history
%   names [E] v1.11.1 (4999a21).  F-2: the code map records the ell in {2,3} refusal as the
%   precondition [E] sec:thetaobj (h) specifies (exit 3).  sec:notimplemented v1.11.1 box with the
%   counts re-derived from --list (141 registered, 82 Hecke*).  app:status: 248 -> 249 lines.
%   v1.11.0's commit (a8e0afb) back-filled.  Companion pointer [E] v1.12.1.  Report
%   docs/implementation/slot11/WP-17_report.md.)
%
%   VERSION 1.11.0 (hecke-extensions slot 11, package WP-15: the slot's X bump published,
%   from v1.11.0-pre.2, with what slot 11 built; a short Fable referee of the delta since
%   pre.2 follows.  Normative: obl:backend(iv)'s scope sentence worded to the digit route
%   as built (ENGINE_API.md 13.11.2 / 13.11.8: single words and one fixed double-word
%   product, not "w-bit operations of the backend of (v)"; K fixed per cell from the
%   certified bound K = k + e_max and capped by [E]'s grammar; the digit route is outside the
%   integer-tier ladder, which is unchanged).  rmk:thetacheckable (C9): the indexing of
%   rank delta said -- source n' = n - 6 (sources 30, 38, targets 36, 44) versus the cell
%   itself (38, 42) -- slot-11 WP-6's finding 1; the list gains COMP^r (WP-5 / WP-6's
%   registered validations; every checkable item held, C7's open alternative decided at
%   n = 14, C9 decided cell by cell).  Status: sec:notimplemented's v1.11.0 box; Part VII's
%   code map (built); sec:codemap gains Part VII's row and sec:integral's (g);
%   prop:relationsmodel's "asserted to exist, not to be built" answered by slot-11 WP-11
%   (lifts for diag(1,p^2), diag(1,p^3), pI built; X7 on H^q, q <= 5); rmk:integralevidence
%   nine validations; open:ordprecision stays OPEN, with what WP-18c built (the free-part
%   bound and the ordinary idempotent) and its pointer corrected (ENGINE_API.md 13.11.11).
%   app:status: COMP 23 -> 24 (registered 16 -> 17); no statement added, no label renamed or
%   removed.  Commits of pre.2 (fb9dcbf) and v1.10.1 (bee8afb) back-filled.  Companion
%   pointer [E] v1.12.0.  Report docs/implementation/slot11/WP-15_report.md.)
%
%   VERSION 1.11.0-pre.2 (hecke-extensions slot 11, package WP-2: the Fable referee's
%   (WP-3, docs/implementation/slot11/WP-3_referee.md; 0 B / 1 M / 9 m, prop:heckeH0
%   confirmed) repairs.  M-1: obl:backend(iv) gains the owner-ruled scope sentence -- the
%   fixed-width l-adic digit-plane route is not arbitrary precision.  m-1: prop:heckeH0's
%   hypotheses per clause and its (4) at ell = 2 (U_2 D^3 = L^2 at m = 6); m-2: C5 records
%   that p+p^2 has multiplicity 2 in H^1(M_22) at ell = 11; m-3: thm:thetales' ad hoc
%   paragraph is a re-derivation note, not a tag (app:status unchanged); m-4:
%   rmk:thetasmallell per prime; m-5: Coleman-Voloch beside Gross, the endpoint k = ell-1;
%   m-7: a proof clause; m-8: pre.1's commit 3b223ef back-filled; m-9: C7's answer (n = 14
%   at ell = 13) recorded.  No label renamed or removed, no tag moved; one bib entry.)
%
%   VERSION 1.11.0-pre.1 (hecke-extensions slot 11, package WP-2: the pre-release of the
%   slot's X bump, "the contract first" -- written before any slot-11 package commits a
%   claim, on the owner's rulings of 2026-09-25 (slot-11 plan S10, all at recommendation;
%   decision 2 (A): no Hida section).  NEW Part VII (placed last, so no chapter number
%   moves): sec:dicksontheta -- def:dicksoninv, lem:dicksontwist (the det twist for all of
%   M_2(F_ell), singular included), prop:dicksonH0 (H^0 = F_ell[L,D], the -I parity),
%   prop:heckeH0 (T_p = (1+p)p^a; U_ell = v(1,0) G_m of rank <= 1 and square zero on H^0 --
%   NOT zero on all of H^0: U_5 D^6 = -L^20 at (120,5); the slot-11 plan's audit M-1 is
%   corrected, rmk:heckeH0plan), lem:heckefibre, def:Qn (Q_n with the FIXED complement basis
%   u_0..u_{ell-1}, u_{n+}), lem:Qnfunctions (dim H^0(Q_n) = [(ell-1)|(n+2)]), def:theta,
%   prop:thetatwist (T_p theta = p theta T_p at cochain level, the direction p pinned;
%   U_ell theta = 0 on cochains), thm:thetales (the LES, H^2 = 0, ker theta = im delta of
%   rank [(ell-1)|(n+2)] - [ell(ell-1)|(n+ell+1)], the Eisenstein line, the closed forms),
%   rmk:thetasmallell, rmk:thetaclassical (Serre, Swinnerton-Dyer, Katz, Jochnowitz,
%   Ash-Stevens, Edixhoven, Gross -- every locator flagged not verified), rmk:thetacheckable
%   (the prediction list C1-C9 of slot-11 WP-6); sec:ordinary -- def:ordinary, prop:ordinary
%   (Fitting, canonical, functorial, exact, compatible with reduction),
%   cor:thetanonordinary, rmk:ordinaryeis, open:ordprecision (OPEN: certified precision).
%   In Part VI: sec:relationsmodel, prop:relationsmodel -- thm:relations' recursion on
%   H^q(C_I(R)) for EVERY ring, ell in {2,3} included (the precondition of line 2's X7).
%   Twenty statements added, no label renamed or removed, no existing tag moved.  Eight
%   bibliography entries added.  Report docs/implementation/slot11/WP-2_report.md.)
%
%   VERSION 1.10.1 (hecke-extensions slot 10, package WP-15, phase 2: a Y bump -- the
%   phase-1 Fable referee's M-1, M-2, m-1 ... m-4 as wording
%   (docs/implementation/slot10/WP-15_referee_phase1.md).  No normative predicate, no
%   hypothesis, no displayed formula, no tag and no label moves; app:status unchanged.
%   M-1: rmk:bredonhecke's "unchanged since v1.0.0" is false of tau and of (a) -- the
%   canonical tau is alg:tau's since v1.1.0 (rmk:iotatausym); its operator (b) equals
%   the legacy tau's at every ell in {2,3} golden point, rmk:tauexistsscope(1), an
%   observation, not a theorem; "compares nothing there" qualified to the operators
%   named.  M-2: rmk:strictliftbuilt(3) 349 -> 359 cells.  m-1: the v1.10.0 header
%   comment's "narrowed" -> "gains a pointer".  m-2: "eight rings among those with a
%   shipped Theta^s_g".  m-3: why the canonical tau is integral (rmk:bredonhecke (a)).
%   m-4: the 224-element check is chi = zeta_+ (rmk:integralevidence).  Status:
%   rmk:integralevidence's counts after slot-10 WP-8 (801, 2483); the v1.10.0 commit
%   back-filled (ca6c9b9); companion pointer [E] v1.11.0.  Report
%   docs/implementation/slot10/WP-15_report.md.)
%
%   VERSION 1.10.0 (hecke-extensions slot 10, package WP-15: an X bump, on the owner's
%   rulings of 2026-09-25 on the slot-10 decisions, all at recommendation (A): 12 (A)
%   this X with a Fable referee, 3 (A) obl:coneexact(vii) over Z/ell^e, 4 (A) the Hecke
%   action on the Bredon object at ell in {2,3} declined.  Normative: lem:bulkperiodic's
%   Step 0 fixes the sign of the Bockstein delta (the snake-lemma connecting map), so the
%   class e of (1) is pinned, not only up to sign (slot-10 WP-14 review m-6); NEW
%   rmk:strictliftbuilt (COMP, registered): prop:strictlift (1), (3), (4) as built for
%   g = diag(1,p), p in {2,3,5}, degrees <= q_max + 1 = 6, with the certificates and the
%   integer sizes measured -- "existence only" stays for every other g and degree; NEW
%   rmk:bredonhecke (commentary): beside prop:strictlift(4), the operator on cohomology of
%   the Bredon object at ell in {2,3} is NOT SHOWN independent of tau (homotopy
%   uniqueness needs 6 in R^x) and [C] compares nothing there -- not a claim of
%   dependence.  Status: obl:coneexact(vii) over Z/ell^e, ell >= 5: sentences 1-2
%   discharged at ell^e < 2^32, every E8 point (slot-10 WP-3), sentence 3 (the PD
%   adjointness gate) recorded unavailable over the ring (the pairing and the
%   contragredient refuse a composite ring; slot-10 WP-14 review M-1); sec:integral:
%   tau_Iso and Theta_Iso,g built (L2-1b (1), WP-10), e and E_q built (L2-1b (2), WP-13),
%   the integral Hecke action read basis-free (M1(a), WP-11); rmk:integralevidence gains
%   three validations; rmk:integralscope(2) gains a pointer; sec:codemap,
%   sec:notimplemented, the obligation roster, the abstract.  app:status +2 lines (228).
%   No label renamed or removed, no existing tag moves.  Validation counts left to the
%   slot close.  Companion pointer [E] v1.11.0.  Report
%   docs/implementation/slot10/WP-15_report.md.)
%
%   VERSION 1.9.1 (hecke-extensions slot 9, package WP-17, the slot close: a Y bump --
%   the WP-15 Fable referee's M-1, M-2, m-1 ... m-4 as wording, and status.  No
%   normative predicate, no hypothesis, no displayed formula, no tag and no label
%   moves; app:status unchanged.  M-1: the slot's independent review (WP-14, pass 1)
%   is now in the tree, docs/implementation/slot9/WP-14_review.md (reconstructed from
%   the reviewer's returns after the 2026-09-25 disk-full loss of the original), and
%   rmk:certifiedeperdiff cites it; M-2: obl:coneexact(iv)'s as-built sentence says
%   arm (4b-C) covers five of the seven points and the two word-regime cells, (60,5)
%   at both k in HeckeReductionCoverageSlowValidation since slot-9 WP-4t; m-1: the
%   commits of v1.9.0 (1d7599a) and v1.8.1 (b7391f7) back-filled; m-2: the
%   non-freeness of C_eq(Z/ell^e) at ell in {2,3} cited to the computation it rests
%   on ([E] sec:ring (c1')), not to prop:ringinforce(2),(3); m-3: obl:coneexact(vii)'s
%   as-built sentence's grammar; m-4: the sec:integral code-map row's lead.  Status:
%   the Hecke* count of the sec:obligations row 65 -> 71 (from --list: 125 registered,
%   52 fast / 63 standard / 10 slow); sec:notimplemented a v1.9.1 sentence; companion
%   pointer [E] v1.10.1.  Close report docs/implementation/slot9/WP-17_report.md.)
%
%   VERSION 1.9.0 (hecke-extensions slot 9, package WP-15: an X bump, on the owner's
%   rulings of 2026-09-24 on the slot-9 decisions, all at recommendation (A): 10 (A)
%   this X with a Fable referee, 3 (iii) (A) obl:coneexact(vii) over Z/ell^e by route 1
%   only, 4 (A) the lift at ell in {2,3} for the cohomology object only.  Normative:
%   rmk:bocksteinexp(2) stated at every prime (its proof never used ell >= 5; slot-8
%   WP-9 finding 6), with the ell in {2,3} instances Z/2^j at n = 2^j and Z/9 at
%   n = 12, 18; NEW rmk:certifiedeperdiff (PROVED): prop:certifiede(3) checked
%   differential by differential (slot-9 WP-11 finding 2, confirmed by the WP-14
%   review); def:ringsemantics(2)'s stale "the implementation of this build refuses it
%   there" replaced by the as-built statement (lengths and invariant factors of the
%   Bredon object at ell in {2,3} by a flat-coordinate route; K-mode, Hecke and
%   reductions refused), (3)(a)-(b) pointed; obl:coneexact(iv)'s "until it exists ...
%   recorded as unchecked" restated: over Z/ell^e, ell >= 5, the non-field half is
%   checked (ConeModuleCertificate item 14, slot-9 WP-4).  Status: obl:coneexact(vi)
%   partial over the ring (commutator asserted, nu = 2 recursion not computed), (vii)
%   route 1 only, composite half open; phase E8 complete as specified (items 1-4b, 4b
%   by route 1); sec:integral partly built (the scoped lift at 2 and 3, WP-3; Iso, Isob,
%   cmp, WP-10 / L2-1a; the integral groups n <= 20, WP-11 / M0); rmk:integralevidence
%   gains three validations.  app:status +1 line (226).  No label renamed or removed,
%   no existing tag moves.  Validation counts left to the slot close.  Companion
%   pointer [E] v1.10.0.  Report docs/implementation/slot9/WP-15_report.md.)
%
%   VERSION 1.8.1 (hecke-extensions slot 8, package WP-14, the slot close: a Y bump --
%   the WP-9 referee's m-1 ... m-5 as wording, and status.  No normative predicate, no
%   hypothesis, no displayed formula, no tag and no label moves; app:status unchanged.
%   Wording (m-1 ... m-5 of the Fable referee of v1.8.0): lem:strictlifting's proof no
%   longer reads "for q = 0, z = 0" as if E_{-1} = 0 were assumed; the notation row of
%   F^sigma_bullet calls it the analogue of def:borelmodel's F_bullet for the stabiliser,
%   not an instance of it (N2); thm:specialisation's v1.8.0 sentence says what v1.8.0 adds
%   for (Xbulk, Xbdry) -- a finite-rank free model with cmp and a Hecke datum, as
%   existence -- instead of reading as if the unrefined case had been deferred;
%   prop:bredon's family D gains the parenthesis that it is the family of all finite
%   subgroups (unused here); rmk:truncdevice cites lem:bulkperiodic(2),(3).  Pre-existing
%   gaps closed: the abstract's stale "two chapters state objects nothing builds"
%   (sec:cone is implemented since v1.7.3; sec:integral is the second unbuilt chapter);
%   changelog subsections for v1.7.2, v1.7.3 and v1.7.4, reconstructed from their date
%   lines, header comments and commits; every resolvable "Commit: [hash]" back-filled;
%   CL5's list of Y bumps completed; app:changelog-gaps' census sentence.  Status:
%   phase E8 items 2, 3 and 4a are implemented (slot-8 WP-1, WP-3, WP-4) -- the complex
%   over the ring in force at Z/ell^k, ell >= 5, its reductions and Bockstein; item 4b
%   (boundary and cone over Z/ell^k, obl:coneexact(iv)'s non-field half) is slot 9's;
%   sec:codemap and sec:notimplemented record it; counts 108 -> 114 registered and
%   61 -> 65 Hecke*, from --list.  Companion pointer [E] v1.9.1.  NOT decided here: the
%   referee's recommendation that rmk:bocksteinexp(2) be stated at every ell (owner:
%   at the next X).  Close report docs/implementation/slot8/WP-14_report.md.)
%
%   VERSION 1.8.0 (hecke-extensions slot 8, package WP-9: the integral line, "L2-C" -- an
%   X bump.  New Part VI, sec:integral (placed last so no chapter number moves): the three
%   complexes on the model and what the pipeline computes over the ring in force
%   (def:ringinforce, prop:ringinforce -- not a quasi-isomorphism at (2, F_2), proved);
%   the shipped complex is Bredon cohomology over the ring in force, triangulation-free
%   (prop:bredon, via the orbit category of the family D and cl:amalgam(1)); the
%   isotropy-resolution model on the UNREFINED tessellation with its boundary subcomplex
%   and one integral comparison map cmp^* (def:isomodel, prop:isomodel; ranks 5, 14, 19,
%   correcting the design brief's 7, 18, 23); C-A5 re-scoped onto it and OD-T1 answered
%   "not needed" (rmk:cafive; owner ruling 6c of 2026-09-24); the fracture square,
%   cohomology-level glueing (flatness only), the descent realised by cmp^*, Hom^{-1} = 0
%   for the bulk (prop:fracture, prop:cohglue, prop:descent, prop:homvanish,
%   rmk:truncdevice); the strict lift of the strict pair datum, existence only
%   (lem:strictlifting, prop:strictlift); the bad primes, the certified-exponent rule with
%   its scope clause, the characteristic polynomial over Z by residues (prop:badprimes,
%   prop:certifiede, prop:crtcharpoly); what each coefficient ring names, the key being
%   [E] sec:ring's (def:ringsemantics); evidence and scope (rmk:integralevidence,
%   rmk:integralscope).  The F-A remark rmk:bocksteinexp beside cl:bockstein (owner
%   ruling 10): the Bockstein counts the exponent-one summands and the Frobenius lower
%   bound, proved; the probe's exponents, COMP^a; cl:bockstein unchanged.  H3 amended
%   ("no comparison map" -> "no isomorphism", with pointers); thm:specialisation's closing
%   paragraph pointed; notation subsection sec:notation-integral and seven preamble
%   macros.  Twenty statements added, all new; no label renamed or removed, no existing
%   tag moves; app:status re-counted mechanically (225 lines).  NOT done: the
%   arbitrary-precision ban is unchanged (OD-3 (c)); no record-schema rule for Z.  Report
%   docs/implementation/slot8/WP-9_report.md.)
%
%   VERSION 1.7.4 (hecke-extensions slot 7, package WP-14, the slot close: STATUS ONLY.
%   No normative predicate, no hypothesis, no displayed formula, no tag and no label moves.
%   Phase E7 (sec:contragredient, cl:pd) is ACCEPTED: cl:pd(1)-(2) are tested on all 54
%   cells of the E2 case list at ell >= 5 (slot-7 WP-1; HeckeDualityValidation arm A16 and
%   HeckeDualityCoverageSlowValidation), failing on none -- cl:pd stays OPEN: a test, not a
%   proof -- and its golden objects are emitted add-only.  The extension record of
%   sec:recordschema item 6 is hecke_ext/v3 since slot-7 stage 3 (nine keys more than v2,
%   the field inventory [E]'s, so no rule of item 6 moves); rule (a) holds as written, and
%   the golden texts carry the settings block reduced to its fingerprint (slot-7 stage 1).
%   Ten mathematical keys are registered.  sec:codemap and sec:notimplemented record it;
%   counts 106 -> 108 registered and 60 -> 61 Hecke*, from the registry (re-derived by
%   --list at the close's second commit).  Companion pointer [E] v1.8.2.
%   NOT decided here: nothing.  Close report docs/implementation/slot7/WP-14_report.md.)
%
%   VERSION 1.7.3 (hecke-extensions slot 6, package WP-14, the slot close: STATUS ONLY.
%   No normative predicate, no hypothesis, no displayed formula, no tag and no label moves.
%   sec:contragredient is IMPLEMENTED (slot-6 WP-1: Mat22::adj, rho_n^vee =
%   rho_n(adj .)^T, the contragredient relation list, Phi_n, the coefficient pairing,
%   H^*(D_n) and A(D_n)(T_p) in Hecke/HeckeContragredient and HeckeOperatorOnK's
%   FibreRepresentation argument; lem:symdual's intertwining asserted UNTWISTED, the
%   p^{+-n}-twisted relations asserted to fail).  obl:coneexact(vii) is IMPLEMENTED in all
%   three sentences: the strict-pair route (WP-4, HeckeConeStrictPairRoute, aborts at 6 in
%   R^x, records otherwise), route 3 (WP-5, HeckeConeRelativeRoute, aborts at every ell),
%   and the PD adjointness gate for X_3 at ell >= 5 (WP-3, Hecke/HeckePairing), which
%   moves from "recorded as unavailable" to asserted -- cl:pd stays OPEN: the gate TESTS
%   it on the computed cells, it does not prove it.  sec:codemap names HeckeContragredient,
%   HeckePairing, HeckePresentation (WP-2) and HeckeAtlas (WP-8), and identity-certificate
%   items 6 (WP-9) and 7 (WP-3).  sec:notimplemented: one chapter unimplemented
%   (sec:orbifoldmodel), none partial except obl:indep's clause (iii).  Counts 100 -> 106
%   registered and 55 -> 60 Hecke*, from --list.  Companion pointer [E] v1.7.2.
%   NOT decided here: nothing.  Close report docs/implementation/slot6/WP-14_report.md.)
%   VERSION 1.7.2 (hecke-extensions slot 5, package WP-13, the slot close: STATUS ONLY.
%   No normative predicate, no hypothesis, no displayed formula, no tag and no label moves.
%   sec:codemap records what slot 5 shipped -- the schema is hecke_ext/v2 and the
%   boundary / pair / cone complexes are emitted objects (WP-1); the cone's Z/B/H/P and
%   ConeCertificate's span inclusions run on engine::tape and ConeCertificate item 11
%   evaluates eq:conedims(1) without consulting H2_zero (WP-4, WP-5), both
%   value-preserving; obl:coneexact(iii)'s r-half is asserted at every ell in the shipped
%   code (WP-4) and ConeActionCertificate gains item 12, record honesty, so
%   restriction_equivariance_asserted reports what the run asserts (WP-12a, owner ruling
%   Q2); obl:iota(ii)'s rider gets the emitted record iota_exact_comparison with verify()
%   code 23 (WP-12a, owner ruling Q3); sec:recordschema item 5's refusal is verify() code
%   22 (WP-11a); Hecke/RhoMemo.h, HeckeExt/HeckeDicksonMultiplication (which discharges the
%   det-twisted identity rho_n(A) L = det(A) L rho_m(A) as code) and
%   HeckeExt/HeckeCanonicalField are named.  Counts 94 -> 100 registered and 49 -> 55
%   Hecke*, both from --list.  The owner's reading of obl:taubdry(vii) left open by v1.7.1
%   is DECIDED: Reading A, a distinct weaker clause, so v1.7.1's text stands as drafted.
%   NOT decided here: nothing.  Close report
%   docs/implementation/slot5/WP-13_report.md.)
%   VERSION 1.7.1 (hecke-extensions slot 5, package WP-12: the independent referee pass on
%   v1.7.0, a Y bump -- POINTERS AND SELF-CONTRADICTIONS ONLY.  No normative predicate, no
%   hypothesis, no displayed formula, no tag and no label moves.  Six repairs, five of them
%   of text v1.7.0 itself introduced: sec:recordschema item 6(f) said the item binds
%   "(a)--(e) and (g)--(h)" while the item's own preamble binds "the nine below", so rule
%   (i) was bound and unbound at once -- now (g)--(i); rmk:conerestrictsmallell's opening
%   said "load-bearing for no clause" and its own next clause says it still discharges
%   obl:taubdry(vii) -- now "no clause of sec:cone", as the changelog already had it;
%   obl:taubdry(vii)'s asymmetry paragraph asserted that "the data this chapter stores are
%   tau_bdry and Theta^bdry alone", which (vii)'s own statement falsifies (it reads the bulk
%   T_g and the restriction r) -- the missing datum is eta/h_derived, and the paragraph now
%   says so; (vii)'s opening "belongs to the pair and is stated in full at obl:coneexact"
%   is aligned with v1.7.0's own new "the pair's form of" in obl:coneexact(iii), the two
%   having become distinct clauses when v1.7.0 strengthened (iii) alone -- WITH AN OPEN
%   POINT FOR THE OWNER, since the other resolution (that (vii) is a pointer and is
%   strengthened with (iii)) is equally available and is a normative choice;
%   obl:coneexact(iii)'s proof now names z in Z^q_eq (WP-11's own residual risk 5, the "one
%   more word"); and sec:recordschema item 5's implementation note records that WP-11a
%   landed the refusal (verify() code 22).  Open for the owner: the (vii)-versus-(iii)
%   reading, and three further findings that are not this pass's to fix.  Report
%   docs/implementation/referee_report_2026-09-22_contract_v1_7_0.md.)
%   VERSION 1.7.0 (hecke-extensions slot 5, package WP-11: the slot's one X revision.
%   obl:coneexact(iii)'s ell in {2,3} carve-out on the r-half is DROPPED and the identity
%   asserted at EVERY ell, WITH THE PROOF WRITTEN IN -- clause (i) is unconditional in
%   ell, so on a cocycle z the difference is d(h z), a coboundary (owner ruling [OD]
%   slot-4 close-out 3; slot-5 WP-4 re-checked the hypothesis in the code) -- together
%   with what the clause does NOT say (the computed complexes, not group cohomology at
%   ell in {2,3}) and the ASYMMETRY: obl:taubdry(vii) keeps its carve-out, having neither
%   eta nor h, and gains a pointer; rmk:conerestrictsmallell is demoted to the
%   measurement that preceded the proof, load-bearing for no clause (still COMP).
%   obl:coneexact(v)'s last sentence names BOTH the certificate and the validation arm,
%   the shipped ConeCertificate now carrying the unconditional first eq:conedims identity
%   under its own guard (OI-2).  sec:recordschema: item 5 makes the derived-complex NAMES
%   the record's own derived flag, so "set" on a derived complex is refusable BY A
%   VERIFIER (WP-1 finding 3, decided here); item 6 becomes hecke_ext/v2 with nine rules
%   -- (b) a run-level verdict at record level with ONE resumable and a count of
%   DECISIONS, (c) a not-computed scalar or flag is null and never a default, (g) the
%   maps between complexes a record carries, (h) the +/- split as two further records
%   whose dimensions SUM (asserted), (i) a matching reported with its witness -- and
%   (a)/(d) answer the two halves of v1.6.0's deferred item 3 (the engine family is
%   excluded by rule, not enumerated; the closed certificate-name list is [E] S7's).
%   obl:iota(ii) gains the rider WP-3 asked for: the fibre half always, the exact half on
%   every row whose gamma the complex stores, recorded either way.  Deferred again:
%   iota-equivariance of h, no accepted package having measured it.  No label added,
%   renamed or removed; no tag moves.  Status sentences are the slot close's except where
%   this version's own text falsifies them.  Report
%   docs/implementation/slot5/WP-11_report.md.)
%   VERSION 1.6.2 (hecke-extensions slot 4, package WP-11: the slot close, a Y bump --
%   STATUS ONLY, no normative statement, hypothesis, displayed formula, asserted
%   predicate, tag or label moved.  The production involution EXISTS: sec:involution's
%   2026-09-16 census, its specification paragraph, sec:codemap's row and
%   sec:notimplemented are rewritten from the pointers v1.6.1 left, naming
%   Hecke/HeckeInvolution, HeckeExt/HeckeSignSplit and HeckeInvolutionValidation (WP-3,
%   phase E3).  sec:recordschema item 6's status -> "written, bulk complex", naming
%   HeckeExt/HeckeExtDriver (WP-1, WP-2, phase E2).  sec:codemap gains
%   HeckeExt/HeckeExtRecord + HeckeExtReport + HeckeExtFinal, HeckePredictions,
%   GloverFactors, the settings keys commutant_budget and final_mode, and -- in the
%   sec:cone row -- WP-4's three VALUE-PRESERVING certificate re-implementations, with
%   the statement that no clause was sampled, relaxed or moved out of compute().
%   rmk:labelcensus records the content-keyed allow-list (WP-6) and the slot-4 run.
%   Counts 89 -> 94 registered and 45 -> 49 Hecke*, both from --list.  NOT decided here
%   (owner's, and the first two are slot-5 CODE): ConeCertificate C11 vs the strengthened
%   clause (v); obl:coneexact(iii)'s ell in {2,3} carve-out, which clause (i) already
%   proves away; whether WP-9's B1 pointer repair belonged to that pass.  Report
%   docs/implementation/slot4/WP-11_report.md.)
%   VERSION 1.6.1 (hecke-extensions slot 4, package WP-9: the independent referee pass on
%   v1.6.0 -- a Z bump, no hypothesis, displayed formula, asserted predicate, tag or label
%   moves, and app:status's counts are unchanged.  Two recorded measurements corrected
%   (rmk:iotapairing's "thirteen cells" -> eight, (10,691) having been counted twice and
%   the four n = 0 cells folded in where obl:iota(v) predicts nothing; sec:coneaction's
%   commutator recorded at (10,691) alone and not on the 15-cell grid); three status
%   sentences that v1.6.0's own new text falsifies given a minimal pointer, the rewrite
%   left to the slot close; the title page's companion pointer moved to [E] v1.5.1.
%   Companion repairs at [E] v1.5.1: cl:harder's new predicate paragraph (false at
%   ell in {2,3} and incomplete), cl:glover(4)'s factor bound, and two proof defects in
%   cl:sharpbound.  Report
%   docs/implementation/referee_report_2026-09-22_contract_v1_6_0.md.)
%   VERSION 1.6.0 (hecke-extensions slot 4, package WP-8: the eight items deferred at
%   v1.5.2, plus the normative findings of the slot-4 packages accepted at the time
%   (WP-1, WP-3) -- an X bump.  obl:coneexact (i) "as matrices" is on the stored C_eq
%   basis with the flat identities RECORDED, (iii) gains the ell in {2,3} carve-out on
%   the r-half, (iv) the explicit im = ker check read degree-generically, (v) the first
%   identity of eq:conedims asserted with no H^2_eq = 0 hypothesis, (vi) where the FULL
%   T_{p^nu} comes from, the preamble "six maps" -> the three map families;
%   lem:conerelative(5) restructured with its proof in two halves; prop:coneambiguity(2)
%   read inside C_eq; lem:etaexists(V)'s "single edge path"; sec:recordschema gains
%   item 6, the extension record hecke_ext/v1 and its six rules (settings + fingerprint,
%   the resumable sidecar's injective name and relative path, the "unavailable"
%   placeholder, certificate_failed:<code or certificate name> and phase_reached per item
%   vs high-water, the two recorded cone flags, and -- from the E2b driver -- that a
%   RESUMED run writes the settings block a fresh run at the same effective values would
%   write, and that the incomplete-item count counts record ENTRIES); rmk:iotapairing's
%   mirror-split prediction becomes a COMP^r record on 13 cells; a new closing paragraph of
%   sec:coneaction states that [iota_c, A_c] = 0 follows from nothing here.  WP-4's E5(d)
%   findings are [X]'s, not this document's: it states no cost bound anywhere.  Owner
%   confirmations of 2026-09-21: item 1(a) yes, 1(b) yes, 1(c) per mode (an [E] item).
%   No label added, renamed or removed; one tag moves (rmk:iotapairing), and
%   app:status's counts with it: COMP 19 = 14 + 5, tagged/untagged 117/88.  Status
%   sentences stay with the slot close, except sec:recordschema item 5's half-sentence
%   that this version's own item 6 falsifies.  v1.5.2's commit back-filled.)
%   VERSION 1.5.2 (hecke-extensions slot 3, package WP-11: the slot close, a Y bump --
%   status only, no normative statement, tag or label moved.  sec:codemap's level (b),(c)
%   row no longer says strict locality throws (WP-3 implemented it -- the one outright-
%   false status sentence of v1.5.1); the cone chapter implemented except obl:coneexact
%   (vii): sec:cone's code map (EquivariantMappingCone, Hecke/SL2ZBoundaryPair,
%   Hecke/HeckeConeAction, the two vertex policies), the obligation roster (obl:eta
%   partial -> discharged, (v) by HeckeConeActionValidation; obl:coneexact unimplemented
%   -> partial, (vii) the cross-routes of slot 6), sec:notimplemented, app:formalization's
%   oracle item; the U_ell repair in the code maps (QuadPointT<I>, escalate_on_overflow,
%   HeckeExt/HeckeUEllBudget, the three U_ell validations); HeckeExt/HeckeSturmSweep;
%   sec:machinery item 2 and rmk:labelcensus (label_census.sh registered as
%   ContractCensusValidation, 0 unresolved); rmk:conerestrictsmallell re-recorded through
%   the cone (424/424); the validation count 80 -> 89 and Hecke* 39 -> 45; v1.5.1's commit
%   back-filled; the deferred list for the next X carries WP-6's four normative findings
%   and the not-applied points of WP-8 / WP-9.)
%   VERSION 1.5.1 (hecke-extensions slot 3, package WP-9: the independent referee pass
%   on v1.5.0 -- a Y bump, no normative statement, tag or label moved.  All ten
%   resolutions accepted; L1 and app:status's counts re-verified mechanically (205 rows,
%   COMP 18 = 13 + 5, PROVED 93, BSTD 34, CONJ 7, OPENSTAT 6, 116/89).  Six defects
%   repaired, all inside the U_ell item: thm:backend(1)'s "every N <= 2^15" is
%   4N^2 <= 2^31-1, i.e. N <= 23170 -- the reach sec:heckedatum already recorded; its
%   "N <~ 4580" is the MEASURED overflow, not a consequence of 7 B_red^2 < 2^63 (which
%   gives ~2700); item (iv)'s "B_red <= 7 C^2 N^2 by lem:bounds(6)" holds for the ENTRY
%   point only, the rest being measured at rmk:btabsmall; lem:bounds(6)'s proof gains the
%   two shapes alg:tau step 1 applies (S, and the nearest-integer T-shift, where the
%   generic 7 |gamma|^2 |z|^2 is not 7 |z|^2) and names w', not v', as the coordinate
%   attaining C^2 N^2; "144.8 N^2 = C^2 N^2" is an approximation, not an equality; and
%   obl:oracle's "off the Kummer loci E is a line" is measured, not proved; and rule A4's
%   worked example still quoted obl:eta(ii) in the form v1.5.0 corrected (the other five
%   "shared (T1) helper" sites are about the tau-table alone and are right).  The v1.5.0
%   lem:bounds changelog line gains PNEW under CL4 with app:status's row kept at
%   [repaired]; v1.5.0's commit back-filled.)
%   VERSION 1.5.0 (hecke-extensions slot 3, package WP-8: the seven items v1.4.2
%   deferred, plus the normative findings of the accepted slot-3 packages WP-1 and
%   WP-5 -- an X bump.  rmk:coneambiguity691's arithmetic rewritten from the measured
%   values (rank r^1 = 1, dim H^1_! = 2, dim X_3 = 3 at (10,691); conclusion unchanged)
%   and the mod-ell reading of "the Eisenstein class" as the joint eigenspace E stated
%   there and at obl:oracle; obl:eta (ii) each table's own (T1) helper on the one (P1)
%   split, (iv) an implementation MAY omit the mod-m regime, preamble the three table
%   hashes (identifying, never certifying); lem:etaexists(V) names the pairwise path
%   routing as the choice made; sec:recordschema item 5 gains the rule for a MAP (the
%   source complex's fields, the shift a schema constant) and for a DERIVED complex;
%   obl:coneexact(iv) gains the per-complex meaning of "reported"; and the U_ell item:
%   lem:bounds gains clause (6) with proof (the alg:tau step-1 reduction point, which
%   (1)-(5) never bounded), thm:backend gains closure item (iv) and its clause (1) is
%   corrected -- "covers every N < 2^15 outright" was false -- with the new measured
%   reach ell <= 23170, refused by obl:tau's own certification past it.  No label
%   added, renamed or removed; one tag change (rmk:coneambiguity691 gains COMP^r);
%   status sentences stay with WP-11 except item 7's, which IS the finding.
%   v1.4.2's commit back-filled.)
%   VERSION 1.4.2 (hecke-extensions slot 2, package WP-13: the slot close, a Y bump --
%   status only, no normative statement changed: sec:recordschema item 5 is
%   "implemented at hecke_report/v2" (WP-2, the stage's one golden re-baseline, exactly
%   the 384 predicted report lines); rmk:conestrict(1)'s COMP counts re-read on the
%   canonical tau (10 of 12 vertex and 5 of 9 edge slots leave the boundary; 11 of 12
%   and 6 of 9 carry a non-zero eta); the code maps of sec:boundarydatum, sec:cone,
%   sec:heckedatum, sec:obligations and sec:codemap carry the slot-2 units
%   (HeckeSettings, HeckeModuleDecomposition's certified loop, HeckeBoundaryDiagnostics,
%   HeckeExt/HeckeBoundary, HeckePairHomotopy, HeckeIntegerSolve, HeckeExt/HeckeLocalAlgebra,
%   HeckeExt/HeckeGenerationCertificates); obl:eta (i)-(iv) discharged by
%   HeckePairHomotopyValidation, (v) awaiting the cone; sec:notimplemented: chapter 15
%   partly implemented; rmk:coneambiguity691 and the U_ell reach record each gain a
%   status sentence (r^1 has rank 1 at (10,691); the reach repair, slot-2 WP-3, did not
%   run -- slot 3); the validation count 77 -> 80; v1.4.1's commit back-filled.)
%   VERSION 1.4.1 (hecke-extensions slot 2, package WP-10: the referee pass on v1.4.0,
%   a Y bump -- all nine resolutions accepted, no normative statement changed; repairs:
%   rmk:taubdrynoiota(3)'s (2,1) exhibit moved to the slots it occurs on (g^{-1} acts
%   on H as z -> N z, so the tie at -1/3 is on the vertex v_1 of K_bdry itself, not on
%   S v_1), (2)'s "consequently" given its attainment step (x(v_0) = -N/2 at p = 2) and
%   (4) made consistent with it; obl:taubdry (vii)'s H^1 control named as the scalar
%   chi_1(p), (v.c)'s unproved cochain-homotopy aside removed; prop:circlemodel(4)'s
%   proof names face closure and the standing hypothesis it does not need; alg:taubdry
%   step 3 wording; Appendix B's tagged/untagged split (obl:taubdry carries COMP^r);
%   v1.4.0's commit back-filled.)
%   VERSION 1.4.0 (hecke-extensions slot 2, package WP-9: the nine open points that
%   v1.3.3 recorded against chapter 14 are resolved -- obl:taubdry (v) split into (v.a)-(v.c)
%   and narrowed to degree 0 at nu >= 2, (vii) qualified; prop:circlemodel(4) identifies
%   the circle model with Rel_bdry INCLUDING the -I self-relations for every n;
%   prop:bdrychar (2) says what the right-coset sum is, (3) states the H^1 coordinate;
%   alg:taubdry (the canonical tau_bdry) and rmk:taubdrynoiota are new; the record
%   fields of sec:recordschema item 5 are per complex; rmk:conerestrictsmallell records
%   (vii) at ell in {2,3} -- an X bump.  v1.3.3 = slot 1 close: chapter 14 implemented,
%   status and code maps only; v1.3.2 = OQ3/OC-1 coherence pass, prose only; v1.3.1 = referee pass on
%   lem:bulkperiodic; v1.3.0 = OC-1
%   decided; v1.2.0 = OQ3 / chapter 9 was prepared in parallel
%   with it -- see Appendix E).  Every chapter and appendix planned by
%   docs/latex/00-unified-contract-architecture.md is written: chapters 1-22
%   and Appendices A-E, so rule (V5)(i) is met; and (V5)(ii) is met as well --
%   hecke_action.tex and monodromy_v5.tex are frozen and banner-ed, and
%   hecke_extensions.tex v1.3 cites the nine labels of chapter 16 instead of
%   defining them, so no label of this document is defined by a live companion.
%   Chapter 17 still states objects nothing implements and chapter 15 is
%   implemented in part (the eta-table since v1.4.2, the cone not) -- chapter 14,
%   the boundary datum, is implemented and its obligation discharged since
%   v1.3.3 -- and a complete document is not a finished programme.  What
%   changed in each version, statement by statement, is Appendix E.
%
%   THE LABEL NAMESPACE IS GATED BY THE BUILD.  All 68 preserved labels (67 from
%   hecke_action.tex + thm:explicit_complex_equiv from monodromy_v5.tex) are
%   declared below, attached to the statement that carries them, and every one of
%   them is dereferenced in Appendix D, so a missing or duplicated label is a
%   build error.  Appendix E dereferences every label it names in the same way.
%
%   BUILD: ./build_gl2_contract.sh   (pdflatex, bibtex, pdflatex, pdflatex)
%======================================================================
\documentclass[11pt]{article}
% SPDX-FileCopyrightText: 2024-2026 Fabian Januszewski
% SPDX-License-Identifier: CC-BY-4.0

%======================================================================
% gl2_preamble.tex --- shared preamble and macro file
%
%   The SINGLE source of macros for the unified contract (gl2_contract.tex)
%   and for every companion document that cites it.  Documents \input this
%   file; they do NOT copy it.  (The pre-2026-09 corpus copied macros between
%   hecke_action.tex, hecke_extensions.tex and monodromy_v5.tex; that is what
%   let three incompatible meanings of \De, \tau and \iota coexist.)
%
%   Usage, after \documentclass and before \begin{document}:
%       % SPDX-FileCopyrightText: 2024-2026 Fabian Januszewski
% SPDX-License-Identifier: CC-BY-4.0

%======================================================================
% gl2_preamble.tex --- shared preamble and macro file
%
%   The SINGLE source of macros for the unified contract (gl2_contract.tex)
%   and for every companion document that cites it.  Documents \input this
%   file; they do NOT copy it.  (The pre-2026-09 corpus copied macros between
%   hecke_action.tex, hecke_extensions.tex and monodromy_v5.tex; that is what
%   let three incompatible meanings of \De, \tau and \iota coexist.)
%
%   Usage, after \documentclass and before \begin{document}:
%       % SPDX-FileCopyrightText: 2024-2026 Fabian Januszewski
% SPDX-License-Identifier: CC-BY-4.0

%======================================================================
% gl2_preamble.tex --- shared preamble and macro file
%
%   The SINGLE source of macros for the unified contract (gl2_contract.tex)
%   and for every companion document that cites it.  Documents \input this
%   file; they do NOT copy it.  (The pre-2026-09 corpus copied macros between
%   hecke_action.tex, hecke_extensions.tex and monodromy_v5.tex; that is what
%   let three incompatible meanings of \De, \tau and \iota coexist.)
%
%   Usage, after \documentclass and before \begin{document}:
%       \input{gl2_preamble}
%
%   Every macro below is normative: the notation chapter of the contract
%   (gl2_notation.tex) is the prose form of this file, and the two must agree.
%   A symbol that is not declared here is not part of the fixed notation.
%
%   RETIRED MACROS.  Symbols retired by the notation chapter are redefined
%   here to raise a LaTeX error rather than silently typeset, so that text
%   re-homed from the old documents cannot carry an old meaning in unnoticed.
%   See the block "Retired symbols" at the end.
%======================================================================

\usepackage{amsmath, amssymb, amsthm}
\usepackage{mathtools}
\usepackage{geometry}
\usepackage{booktabs}
\usepackage{array}
\usepackage{longtable}
\usepackage{enumitem}
\usepackage{tikz}
\usepackage{tikz-cd}
\usepackage[colorlinks=true,linkcolor=blue,citecolor=blue,urlcolor=blue]{hyperref}
\setlist{nosep}
\emergencystretch=3em

%----------------------------------------------------------------------
% Theorem environments.  One counter, sectioned, so every statement has a
% unique number of the form <section>.<n> that can be cited from code.
%----------------------------------------------------------------------
\newtheorem{theorem}{Theorem}[section]
\newtheorem{lemma}[theorem]{Lemma}
\newtheorem{proposition}[theorem]{Proposition}
\newtheorem{corollary}[theorem]{Corollary}
\newtheorem{claim}[theorem]{Claim}

\theoremstyle{definition}
\newtheorem{definition}[theorem]{Definition}
\newtheorem{algorithm}[theorem]{Algorithm}
\newtheorem{obligation}[theorem]{Obligation}
\newtheorem{certificate}[theorem]{Certificate}

\theoremstyle{remark}
\newtheorem{remark}[theorem]{Remark}
\newtheorem{assumption}[theorem]{Assumption}
\newtheorem{example}[theorem]{Example}
\newtheorem{prediction}[theorem]{Prediction}
\newtheorem{openitem}[theorem]{Open construction item}

%----------------------------------------------------------------------
% Normative boxes.
%----------------------------------------------------------------------
\newenvironment{contractbox}
  {\par\medskip\noindent\begin{quote}\textbf{CONTRACT.}\ }
  {\end{quote}\medskip}
\newenvironment{decisionbox}
  {\par\medskip\noindent\begin{quote}\textbf{OWNER DECISION.}\ }
  {\end{quote}\medskip}
\newenvironment{erratumbox}
  {\par\medskip\noindent\begin{quote}\textbf{ERRATUM.}\ }
  {\end{quote}\medskip}

%----------------------------------------------------------------------
% STATUS tags.  Every non-definitional assertion carries exactly one.
%----------------------------------------------------------------------
\newcommand{\PROVED}{\textsf{[proved here]}}
\newcommand{\STD}[1]{\textsf{[standard: #1]}}   % MUST carry a precise citation
\newcommand{\CONJ}{\textsf{[conjectural]}}
\newcommand{\COMP}{\textsf{[verified computationally]}}
\newcommand{\OPENSTAT}{\textsf{[open]}}

%----------------------------------------------------------------------
% Number systems, rings, groups
%----------------------------------------------------------------------
\newcommand{\Z}{\mathbb{Z}}
\newcommand{\Q}{\mathbb{Q}}
\newcommand{\Rl}{\mathbb{R}}              % the real numbers (NOT \R: see below)
\newcommand{\Cx}{\mathbb{C}}              % the complex numbers (NOT \C)
\newcommand{\F}{\mathbb{F}}
\newcommand{\Fl}{\F_\ell}
\newcommand{\Flf}{\F_{\ell^{f}}}          % residue field, residue degree f
\newcommand{\Hyp}{\mathcal{H}}            % the upper half plane
\newcommand{\GL}{\mathrm{GL}}
\newcommand{\SL}{\mathrm{SL}}
\newcommand{\Mat}{\mathrm{M}}             % matrix algebra: \Mat_2(\Z)
\newcommand{\Ga}{\Gamma}
\newcommand{\Gbdry}{\Ga_\infty}           % the boundary group <T> x <-I>
\newcommand{\De}{\Delta}                  % the monoid {g in M_2(Z) : det g > 0}
\newcommand{\Deprim}{\De^{\mathrm{prim}}}
\newcommand{\cring}{R}                    % the coefficient ring: F_ell, Z/ell^e, Z/m
\newcommand{\Rel}{\mathcal{R}}            % the relation set of the model (K,\Rel)

%----------------------------------------------------------------------
% Operators
%----------------------------------------------------------------------
\newcommand{\Hom}{\mathrm{Hom}}
\newcommand{\End}{\mathrm{End}}
\newcommand{\Ind}{\mathrm{Ind}}
\newcommand{\Sym}{\mathrm{Sym}}
\newcommand{\id}{\mathrm{id}}
\DeclareMathOperator{\res}{res}
\DeclareMathOperator{\cor}{cor}
\DeclareMathOperator{\conj}{conj}
\DeclareMathOperator{\adjg}{adj}
\DeclareMathOperator{\tr}{tr}
\DeclareMathOperator{\rank}{rank}
\DeclareMathOperator{\sgn}{sign}
\DeclareMathOperator{\fl}{fl}
\DeclareMathOperator{\im}{im}
\DeclareMathOperator{\coker}{coker}
\DeclareMathOperator{\soc}{soc}
\DeclareMathOperator{\rad}{rad}
\DeclareMathOperator{\Ext}{Ext}
\DeclareMathOperator{\Ann}{Ann}
\DeclareMathOperator{\MaxSpec}{MaxSpec}
\DeclareMathOperator{\Frob}{Frob}
\DeclareMathOperator{\supp}{supp}
\newcommand{\aug}{\varepsilon}            % augmentation C_0 -> Z

%----------------------------------------------------------------------
% The simplicial framework
%----------------------------------------------------------------------
\newcommand{\Ku}{K^{u}}                   % universal cover complex
\newcommand{\KR}{K/\Rel}                  % quotient complex
\newcommand{\prR}{\mathrm{pr}_{\Rel}}     % projection K -> K/\Rel
\newcommand{\pru}{\mathrm{pr}^{u}}        % projection K^u -> K/\Rel
\newcommand{\piOne}{\pi_1}
\newcommand{\LCS}{\mathbf{LCS}}
\newcommand{\Rep}{\mathbf{Rep}}
\newcommand{\LS}{\mathbf{LS}}
\newcommand{\cat}[1]{\mathbf{#1}}
\newcommand{\orient}[1]{o_{#1}}           % oriented q-simplices: \orient{q}(K)
% The distinguished lift of a simplex of K into the universal cover K^u, for the
% strict fundamental domain K \subseteq K^u of (H4)(iii).  The code spells this
% i(sigma); the document cannot, because i is sqrt(-1) (notation chapter, C15/C16).
\newcommand{\lift}[1]{\widetilde{#1}}
\newcommand{\Ceq}{C^{\bullet}_{\mathrm{eq}}}
\newcommand{\Cq}[1]{C^{#1}_{\mathrm{eq}}}
\newcommand{\Zq}[1]{Z^{#1}_{\mathrm{eq}}}
\newcommand{\Bq}[1]{B^{#1}_{\mathrm{eq}}}
\newcommand{\Heq}[1]{H^{#1}_{\mathrm{eq}}}

%----------------------------------------------------------------------
% Coefficients
%----------------------------------------------------------------------
\newcommand{\rhon}{\rho_n}                % the symmetric power representation
\newcommand{\rhodual}{\rho_n^{\vee}}      % the contragredient
\newcommand{\Vcoef}{V}                    % V := Sym^n(Z^2)
\newcommand{\VR}{V_{\cring}}              % V \otimes \cring
\newcommand{\Mn}{M_n}                     % Sym^n of F_ell^2 (synonym of \VR at R = F_ell)
\newcommand{\Dn}{D_n}                     % the divided-power / contragredient module
\newcommand{\cf}{f_{12}}                  % the normalised weight-12 cusp form (Ramanujan's Delta)

%----------------------------------------------------------------------
% The Hecke datum
%----------------------------------------------------------------------
\newcommand{\Th}{\Theta}                  % the fused cochain map Theta_g
\newcommand{\Thb}{\Theta^{\partial}}      % the boundary cochain map
\newcommand{\Xig}{\Xi_g}
\newcommand{\taub}{\tau_{\partial}}       % the boundary g-twisted chain datum
\newcommand{\Phig}{\Phi_g}                % the cone cochain map
\newcommand{\Amat}[1]{A^{(#1)}}           % cohomology-level Hecke matrix
\newcommand{\Acone}[1]{A_c^{(#1)}}
\newcommand{\Abdry}[1]{A_{\partial}^{(#1)}}
\newcommand{\Tdc}[1]{T_{\Ga #1 \Ga}}      % T_{Gamma g Gamma}
\newcommand{\diam}{T(p,p)}                % the diamond operator

%----------------------------------------------------------------------
% The pair, the cone, compact supports
%----------------------------------------------------------------------
\newcommand{\Xbulk}{X}                    % the Gamma-tessellation of the truncated half plane
\newcommand{\Xbdry}{\partial X}
\newcommand{\Kbdry}{K_{\partial}}         % the circle model 0-1-2-3
\newcommand{\rres}{r}                     % the restriction cochain map r = j^*
\newcommand{\Cone}[1]{C(#1)}              % the mapping cone
\newcommand{\Dcone}{D}                    % the cone differential
\newcommand{\coneinc}{\mathrm{inc}_c}     % C_eq(K_bdry) -> C(r)
\newcommand{\coneproj}{\mathrm{pr}_c}     % C(r) -> C_eq(K)[1]
\newcommand{\Hcpt}[1]{H^{#1}_{c}}
\newcommand{\etapair}{\eta}               % the pair homotopy
\newcommand{\hder}{h}                     % the derived cochain homotopy h = \tilde\eta
\newcommand{\Lam}{\Lambda}                % the unique C_0(bdry X) -> C_2(X) of eta-uniqueness

%----------------------------------------------------------------------
% The integral line (sec:integral, since [C] v1.8.0; notation
% sec:notation-integral)
%----------------------------------------------------------------------
\newcommand{\Iso}{\mathcal{I}}            % the isotropy-resolution model on X
\newcommand{\Isob}{\mathcal{I}^{\partial}} % its boundary subcomplex
\newcommand{\cmpI}{\mathrm{cmp}}          % the comparison map I_* -> C_*(X)
\newcommand{\CI}{C^{\bullet}_{\Iso}}      % Hom_{Z[Gamma]}(I_*, V_R)
\newcommand{\Hfr}[1]{H^{#1}_{\mathrm{fr}}} % H^q(Gamma;V) modulo torsion
\newcommand{\Zsix}{\widehat{\Z}_{6}}      % Z_2 x Z_3
\newcommand{\Qsix}{\widehat{\Q}_{6}}      % Q_2 x Q_3

%----------------------------------------------------------------------
% The Dickson line and the ordinary projector (sec:dicksontheta,
% sec:ordinary, since [C] v1.11.0-pre.1; notation sec:notation-dickson).
% Sans-serif L, D: the plain letters are taken (L the local system, D_n the
% contragredient module), rule A5.
%----------------------------------------------------------------------
\newcommand{\DkL}{\mathsf{L}}             % the first Dickson invariant x^ell y - x y^ell (the code's L)
\newcommand{\DkD}{\mathsf{D}}             % the second Dickson invariant, degree ell^2 - ell
\newcommand{\thL}{\theta_{\DkL}}          % multiplication by L on cochains and cohomology
\newcommand{\nplus}{n^{+}}               % n + ell + 1 (N is the Hecke level only)
\newcommand{\Qn}{Q_{n}}                   % the quotient Sym^{n+ell+1} / L Sym^n
\newcommand{\eord}{e_{\mathrm{ord}}}      % the ordinary projector of def:ordinary (e alone is lem:bulkperiodic's class)

%----------------------------------------------------------------------
% The involution
%----------------------------------------------------------------------
\newcommand{\inv}{\iota}                  % the involution diag(-1,1) / z -> -zbar
\newcommand{\invK}{\iota_K}               % its simplicial realisation on the 7-vertex model

%----------------------------------------------------------------------
% The Hecke algebra and its local pieces
%----------------------------------------------------------------------
\newcommand{\Talg}{\mathbb{T}}
\newcommand{\midl}{\mathfrak{m}}
\newcommand{\Htot}{H_{\mathrm{tot}}}

%----------------------------------------------------------------------
% Misc
%----------------------------------------------------------------------
\newcommand{\abs}[1]{\left| #1 \right|}
\newcommand{\smat}[4]{\begin{pmatrix} #1 & #2 \\ #3 & #4 \end{pmatrix}}
\newcommand{\code}[1]{\texttt{#1}}
\newcommand{\file}[1]{\texttt{#1}}

%----------------------------------------------------------------------
% Cross-document citation.  \HA{} survives for documents that still cite the
% OLD hecke_action.tex; inside the unified contract, use \ref{} instead.
%----------------------------------------------------------------------
\newcommand{\HA}[1]{[C,~\texttt{#1}]}
\newcommand{\vfive}[1]{[v5,~\texttt{#1}]}
\newcommand{\Espec}[1]{[E,~\texttt{#1}]}

%----------------------------------------------------------------------
% RETIRED SYMBOLS.  These macros existed in the pre-consolidation corpus with
% a meaning the notation chapter has retired.  Using one is an error, not a
% typographic choice.  (\K, \R, \C, \pR, \pu, \image were in the copied macro
% blocks of hecke_action.tex and monodromy_v5.tex.)
%----------------------------------------------------------------------
\newcommand{\gltwoRETIRED}[2]{%
  \PackageError{gl2_preamble}{Retired macro \string#1 -- #2}{%
    See the notation chapter (sec:notation) for the replacement.}}
\newcommand{\K}{\gltwoRETIRED{K}{ambiguous complex name; write K directly}}
\newcommand{\R}{\gltwoRETIRED{R}{R is the coefficient ring; write \string\Rl\space for the reals}}
\newcommand{\C}{\gltwoRETIRED{C}{write \string\Cx\space for the complex numbers}}
\newcommand{\pR}{\gltwoRETIRED{pR}{write \string\prR}}
\newcommand{\pu}{\gltwoRETIRED{pu}{write \string\pru}}
\newcommand{\image}{\gltwoRETIRED{image}{write \string\im}}
\newcommand{\M}{\gltwoRETIRED{M}{write \string\Mat\space for the matrix algebra}}

%
%   Every macro below is normative: the notation chapter of the contract
%   (gl2_notation.tex) is the prose form of this file, and the two must agree.
%   A symbol that is not declared here is not part of the fixed notation.
%
%   RETIRED MACROS.  Symbols retired by the notation chapter are redefined
%   here to raise a LaTeX error rather than silently typeset, so that text
%   re-homed from the old documents cannot carry an old meaning in unnoticed.
%   See the block "Retired symbols" at the end.
%======================================================================

\usepackage{amsmath, amssymb, amsthm}
\usepackage{mathtools}
\usepackage{geometry}
\usepackage{booktabs}
\usepackage{array}
\usepackage{longtable}
\usepackage{enumitem}
\usepackage{tikz}
\usepackage{tikz-cd}
\usepackage[colorlinks=true,linkcolor=blue,citecolor=blue,urlcolor=blue]{hyperref}
\setlist{nosep}
\emergencystretch=3em

%----------------------------------------------------------------------
% Theorem environments.  One counter, sectioned, so every statement has a
% unique number of the form <section>.<n> that can be cited from code.
%----------------------------------------------------------------------
\newtheorem{theorem}{Theorem}[section]
\newtheorem{lemma}[theorem]{Lemma}
\newtheorem{proposition}[theorem]{Proposition}
\newtheorem{corollary}[theorem]{Corollary}
\newtheorem{claim}[theorem]{Claim}

\theoremstyle{definition}
\newtheorem{definition}[theorem]{Definition}
\newtheorem{algorithm}[theorem]{Algorithm}
\newtheorem{obligation}[theorem]{Obligation}
\newtheorem{certificate}[theorem]{Certificate}

\theoremstyle{remark}
\newtheorem{remark}[theorem]{Remark}
\newtheorem{assumption}[theorem]{Assumption}
\newtheorem{example}[theorem]{Example}
\newtheorem{prediction}[theorem]{Prediction}
\newtheorem{openitem}[theorem]{Open construction item}

%----------------------------------------------------------------------
% Normative boxes.
%----------------------------------------------------------------------
\newenvironment{contractbox}
  {\par\medskip\noindent\begin{quote}\textbf{CONTRACT.}\ }
  {\end{quote}\medskip}
\newenvironment{decisionbox}
  {\par\medskip\noindent\begin{quote}\textbf{OWNER DECISION.}\ }
  {\end{quote}\medskip}
\newenvironment{erratumbox}
  {\par\medskip\noindent\begin{quote}\textbf{ERRATUM.}\ }
  {\end{quote}\medskip}

%----------------------------------------------------------------------
% STATUS tags.  Every non-definitional assertion carries exactly one.
%----------------------------------------------------------------------
\newcommand{\PROVED}{\textsf{[proved here]}}
\newcommand{\STD}[1]{\textsf{[standard: #1]}}   % MUST carry a precise citation
\newcommand{\CONJ}{\textsf{[conjectural]}}
\newcommand{\COMP}{\textsf{[verified computationally]}}
\newcommand{\OPENSTAT}{\textsf{[open]}}

%----------------------------------------------------------------------
% Number systems, rings, groups
%----------------------------------------------------------------------
\newcommand{\Z}{\mathbb{Z}}
\newcommand{\Q}{\mathbb{Q}}
\newcommand{\Rl}{\mathbb{R}}              % the real numbers (NOT \R: see below)
\newcommand{\Cx}{\mathbb{C}}              % the complex numbers (NOT \C)
\newcommand{\F}{\mathbb{F}}
\newcommand{\Fl}{\F_\ell}
\newcommand{\Flf}{\F_{\ell^{f}}}          % residue field, residue degree f
\newcommand{\Hyp}{\mathcal{H}}            % the upper half plane
\newcommand{\GL}{\mathrm{GL}}
\newcommand{\SL}{\mathrm{SL}}
\newcommand{\Mat}{\mathrm{M}}             % matrix algebra: \Mat_2(\Z)
\newcommand{\Ga}{\Gamma}
\newcommand{\Gbdry}{\Ga_\infty}           % the boundary group <T> x <-I>
\newcommand{\De}{\Delta}                  % the monoid {g in M_2(Z) : det g > 0}
\newcommand{\Deprim}{\De^{\mathrm{prim}}}
\newcommand{\cring}{R}                    % the coefficient ring: F_ell, Z/ell^e, Z/m
\newcommand{\Rel}{\mathcal{R}}            % the relation set of the model (K,\Rel)

%----------------------------------------------------------------------
% Operators
%----------------------------------------------------------------------
\newcommand{\Hom}{\mathrm{Hom}}
\newcommand{\End}{\mathrm{End}}
\newcommand{\Ind}{\mathrm{Ind}}
\newcommand{\Sym}{\mathrm{Sym}}
\newcommand{\id}{\mathrm{id}}
\DeclareMathOperator{\res}{res}
\DeclareMathOperator{\cor}{cor}
\DeclareMathOperator{\conj}{conj}
\DeclareMathOperator{\adjg}{adj}
\DeclareMathOperator{\tr}{tr}
\DeclareMathOperator{\rank}{rank}
\DeclareMathOperator{\sgn}{sign}
\DeclareMathOperator{\fl}{fl}
\DeclareMathOperator{\im}{im}
\DeclareMathOperator{\coker}{coker}
\DeclareMathOperator{\soc}{soc}
\DeclareMathOperator{\rad}{rad}
\DeclareMathOperator{\Ext}{Ext}
\DeclareMathOperator{\Ann}{Ann}
\DeclareMathOperator{\MaxSpec}{MaxSpec}
\DeclareMathOperator{\Frob}{Frob}
\DeclareMathOperator{\supp}{supp}
\newcommand{\aug}{\varepsilon}            % augmentation C_0 -> Z

%----------------------------------------------------------------------
% The simplicial framework
%----------------------------------------------------------------------
\newcommand{\Ku}{K^{u}}                   % universal cover complex
\newcommand{\KR}{K/\Rel}                  % quotient complex
\newcommand{\prR}{\mathrm{pr}_{\Rel}}     % projection K -> K/\Rel
\newcommand{\pru}{\mathrm{pr}^{u}}        % projection K^u -> K/\Rel
\newcommand{\piOne}{\pi_1}
\newcommand{\LCS}{\mathbf{LCS}}
\newcommand{\Rep}{\mathbf{Rep}}
\newcommand{\LS}{\mathbf{LS}}
\newcommand{\cat}[1]{\mathbf{#1}}
\newcommand{\orient}[1]{o_{#1}}           % oriented q-simplices: \orient{q}(K)
% The distinguished lift of a simplex of K into the universal cover K^u, for the
% strict fundamental domain K \subseteq K^u of (H4)(iii).  The code spells this
% i(sigma); the document cannot, because i is sqrt(-1) (notation chapter, C15/C16).
\newcommand{\lift}[1]{\widetilde{#1}}
\newcommand{\Ceq}{C^{\bullet}_{\mathrm{eq}}}
\newcommand{\Cq}[1]{C^{#1}_{\mathrm{eq}}}
\newcommand{\Zq}[1]{Z^{#1}_{\mathrm{eq}}}
\newcommand{\Bq}[1]{B^{#1}_{\mathrm{eq}}}
\newcommand{\Heq}[1]{H^{#1}_{\mathrm{eq}}}

%----------------------------------------------------------------------
% Coefficients
%----------------------------------------------------------------------
\newcommand{\rhon}{\rho_n}                % the symmetric power representation
\newcommand{\rhodual}{\rho_n^{\vee}}      % the contragredient
\newcommand{\Vcoef}{V}                    % V := Sym^n(Z^2)
\newcommand{\VR}{V_{\cring}}              % V \otimes \cring
\newcommand{\Mn}{M_n}                     % Sym^n of F_ell^2 (synonym of \VR at R = F_ell)
\newcommand{\Dn}{D_n}                     % the divided-power / contragredient module
\newcommand{\cf}{f_{12}}                  % the normalised weight-12 cusp form (Ramanujan's Delta)

%----------------------------------------------------------------------
% The Hecke datum
%----------------------------------------------------------------------
\newcommand{\Th}{\Theta}                  % the fused cochain map Theta_g
\newcommand{\Thb}{\Theta^{\partial}}      % the boundary cochain map
\newcommand{\Xig}{\Xi_g}
\newcommand{\taub}{\tau_{\partial}}       % the boundary g-twisted chain datum
\newcommand{\Phig}{\Phi_g}                % the cone cochain map
\newcommand{\Amat}[1]{A^{(#1)}}           % cohomology-level Hecke matrix
\newcommand{\Acone}[1]{A_c^{(#1)}}
\newcommand{\Abdry}[1]{A_{\partial}^{(#1)}}
\newcommand{\Tdc}[1]{T_{\Ga #1 \Ga}}      % T_{Gamma g Gamma}
\newcommand{\diam}{T(p,p)}                % the diamond operator

%----------------------------------------------------------------------
% The pair, the cone, compact supports
%----------------------------------------------------------------------
\newcommand{\Xbulk}{X}                    % the Gamma-tessellation of the truncated half plane
\newcommand{\Xbdry}{\partial X}
\newcommand{\Kbdry}{K_{\partial}}         % the circle model 0-1-2-3
\newcommand{\rres}{r}                     % the restriction cochain map r = j^*
\newcommand{\Cone}[1]{C(#1)}              % the mapping cone
\newcommand{\Dcone}{D}                    % the cone differential
\newcommand{\coneinc}{\mathrm{inc}_c}     % C_eq(K_bdry) -> C(r)
\newcommand{\coneproj}{\mathrm{pr}_c}     % C(r) -> C_eq(K)[1]
\newcommand{\Hcpt}[1]{H^{#1}_{c}}
\newcommand{\etapair}{\eta}               % the pair homotopy
\newcommand{\hder}{h}                     % the derived cochain homotopy h = \tilde\eta
\newcommand{\Lam}{\Lambda}                % the unique C_0(bdry X) -> C_2(X) of eta-uniqueness

%----------------------------------------------------------------------
% The integral line (sec:integral, since [C] v1.8.0; notation
% sec:notation-integral)
%----------------------------------------------------------------------
\newcommand{\Iso}{\mathcal{I}}            % the isotropy-resolution model on X
\newcommand{\Isob}{\mathcal{I}^{\partial}} % its boundary subcomplex
\newcommand{\cmpI}{\mathrm{cmp}}          % the comparison map I_* -> C_*(X)
\newcommand{\CI}{C^{\bullet}_{\Iso}}      % Hom_{Z[Gamma]}(I_*, V_R)
\newcommand{\Hfr}[1]{H^{#1}_{\mathrm{fr}}} % H^q(Gamma;V) modulo torsion
\newcommand{\Zsix}{\widehat{\Z}_{6}}      % Z_2 x Z_3
\newcommand{\Qsix}{\widehat{\Q}_{6}}      % Q_2 x Q_3

%----------------------------------------------------------------------
% The Dickson line and the ordinary projector (sec:dicksontheta,
% sec:ordinary, since [C] v1.11.0-pre.1; notation sec:notation-dickson).
% Sans-serif L, D: the plain letters are taken (L the local system, D_n the
% contragredient module), rule A5.
%----------------------------------------------------------------------
\newcommand{\DkL}{\mathsf{L}}             % the first Dickson invariant x^ell y - x y^ell (the code's L)
\newcommand{\DkD}{\mathsf{D}}             % the second Dickson invariant, degree ell^2 - ell
\newcommand{\thL}{\theta_{\DkL}}          % multiplication by L on cochains and cohomology
\newcommand{\nplus}{n^{+}}               % n + ell + 1 (N is the Hecke level only)
\newcommand{\Qn}{Q_{n}}                   % the quotient Sym^{n+ell+1} / L Sym^n
\newcommand{\eord}{e_{\mathrm{ord}}}      % the ordinary projector of def:ordinary (e alone is lem:bulkperiodic's class)

%----------------------------------------------------------------------
% The involution
%----------------------------------------------------------------------
\newcommand{\inv}{\iota}                  % the involution diag(-1,1) / z -> -zbar
\newcommand{\invK}{\iota_K}               % its simplicial realisation on the 7-vertex model

%----------------------------------------------------------------------
% The Hecke algebra and its local pieces
%----------------------------------------------------------------------
\newcommand{\Talg}{\mathbb{T}}
\newcommand{\midl}{\mathfrak{m}}
\newcommand{\Htot}{H_{\mathrm{tot}}}

%----------------------------------------------------------------------
% Misc
%----------------------------------------------------------------------
\newcommand{\abs}[1]{\left| #1 \right|}
\newcommand{\smat}[4]{\begin{pmatrix} #1 & #2 \\ #3 & #4 \end{pmatrix}}
\newcommand{\code}[1]{\texttt{#1}}
\newcommand{\file}[1]{\texttt{#1}}

%----------------------------------------------------------------------
% Cross-document citation.  \HA{} survives for documents that still cite the
% OLD hecke_action.tex; inside the unified contract, use \ref{} instead.
%----------------------------------------------------------------------
\newcommand{\HA}[1]{[C,~\texttt{#1}]}
\newcommand{\vfive}[1]{[v5,~\texttt{#1}]}
\newcommand{\Espec}[1]{[E,~\texttt{#1}]}

%----------------------------------------------------------------------
% RETIRED SYMBOLS.  These macros existed in the pre-consolidation corpus with
% a meaning the notation chapter has retired.  Using one is an error, not a
% typographic choice.  (\K, \R, \C, \pR, \pu, \image were in the copied macro
% blocks of hecke_action.tex and monodromy_v5.tex.)
%----------------------------------------------------------------------
\newcommand{\gltwoRETIRED}[2]{%
  \PackageError{gl2_preamble}{Retired macro \string#1 -- #2}{%
    See the notation chapter (sec:notation) for the replacement.}}
\newcommand{\K}{\gltwoRETIRED{K}{ambiguous complex name; write K directly}}
\newcommand{\R}{\gltwoRETIRED{R}{R is the coefficient ring; write \string\Rl\space for the reals}}
\newcommand{\C}{\gltwoRETIRED{C}{write \string\Cx\space for the complex numbers}}
\newcommand{\pR}{\gltwoRETIRED{pR}{write \string\prR}}
\newcommand{\pu}{\gltwoRETIRED{pu}{write \string\pru}}
\newcommand{\image}{\gltwoRETIRED{image}{write \string\im}}
\newcommand{\M}{\gltwoRETIRED{M}{write \string\Mat\space for the matrix algebra}}

%
%   Every macro below is normative: the notation chapter of the contract
%   (gl2_notation.tex) is the prose form of this file, and the two must agree.
%   A symbol that is not declared here is not part of the fixed notation.
%
%   RETIRED MACROS.  Symbols retired by the notation chapter are redefined
%   here to raise a LaTeX error rather than silently typeset, so that text
%   re-homed from the old documents cannot carry an old meaning in unnoticed.
%   See the block "Retired symbols" at the end.
%======================================================================

\usepackage{amsmath, amssymb, amsthm}
\usepackage{mathtools}
\usepackage{geometry}
\usepackage{booktabs}
\usepackage{array}
\usepackage{longtable}
\usepackage{enumitem}
\usepackage{tikz}
\usepackage{tikz-cd}
\usepackage[colorlinks=true,linkcolor=blue,citecolor=blue,urlcolor=blue]{hyperref}
\setlist{nosep}
\emergencystretch=3em

%----------------------------------------------------------------------
% Theorem environments.  One counter, sectioned, so every statement has a
% unique number of the form <section>.<n> that can be cited from code.
%----------------------------------------------------------------------
\newtheorem{theorem}{Theorem}[section]
\newtheorem{lemma}[theorem]{Lemma}
\newtheorem{proposition}[theorem]{Proposition}
\newtheorem{corollary}[theorem]{Corollary}
\newtheorem{claim}[theorem]{Claim}

\theoremstyle{definition}
\newtheorem{definition}[theorem]{Definition}
\newtheorem{algorithm}[theorem]{Algorithm}
\newtheorem{obligation}[theorem]{Obligation}
\newtheorem{certificate}[theorem]{Certificate}

\theoremstyle{remark}
\newtheorem{remark}[theorem]{Remark}
\newtheorem{assumption}[theorem]{Assumption}
\newtheorem{example}[theorem]{Example}
\newtheorem{prediction}[theorem]{Prediction}
\newtheorem{openitem}[theorem]{Open construction item}

%----------------------------------------------------------------------
% Normative boxes.
%----------------------------------------------------------------------
\newenvironment{contractbox}
  {\par\medskip\noindent\begin{quote}\textbf{CONTRACT.}\ }
  {\end{quote}\medskip}
\newenvironment{decisionbox}
  {\par\medskip\noindent\begin{quote}\textbf{OWNER DECISION.}\ }
  {\end{quote}\medskip}
\newenvironment{erratumbox}
  {\par\medskip\noindent\begin{quote}\textbf{ERRATUM.}\ }
  {\end{quote}\medskip}

%----------------------------------------------------------------------
% STATUS tags.  Every non-definitional assertion carries exactly one.
%----------------------------------------------------------------------
\newcommand{\PROVED}{\textsf{[proved here]}}
\newcommand{\STD}[1]{\textsf{[standard: #1]}}   % MUST carry a precise citation
\newcommand{\CONJ}{\textsf{[conjectural]}}
\newcommand{\COMP}{\textsf{[verified computationally]}}
\newcommand{\OPENSTAT}{\textsf{[open]}}

%----------------------------------------------------------------------
% Number systems, rings, groups
%----------------------------------------------------------------------
\newcommand{\Z}{\mathbb{Z}}
\newcommand{\Q}{\mathbb{Q}}
\newcommand{\Rl}{\mathbb{R}}              % the real numbers (NOT \R: see below)
\newcommand{\Cx}{\mathbb{C}}              % the complex numbers (NOT \C)
\newcommand{\F}{\mathbb{F}}
\newcommand{\Fl}{\F_\ell}
\newcommand{\Flf}{\F_{\ell^{f}}}          % residue field, residue degree f
\newcommand{\Hyp}{\mathcal{H}}            % the upper half plane
\newcommand{\GL}{\mathrm{GL}}
\newcommand{\SL}{\mathrm{SL}}
\newcommand{\Mat}{\mathrm{M}}             % matrix algebra: \Mat_2(\Z)
\newcommand{\Ga}{\Gamma}
\newcommand{\Gbdry}{\Ga_\infty}           % the boundary group <T> x <-I>
\newcommand{\De}{\Delta}                  % the monoid {g in M_2(Z) : det g > 0}
\newcommand{\Deprim}{\De^{\mathrm{prim}}}
\newcommand{\cring}{R}                    % the coefficient ring: F_ell, Z/ell^e, Z/m
\newcommand{\Rel}{\mathcal{R}}            % the relation set of the model (K,\Rel)

%----------------------------------------------------------------------
% Operators
%----------------------------------------------------------------------
\newcommand{\Hom}{\mathrm{Hom}}
\newcommand{\End}{\mathrm{End}}
\newcommand{\Ind}{\mathrm{Ind}}
\newcommand{\Sym}{\mathrm{Sym}}
\newcommand{\id}{\mathrm{id}}
\DeclareMathOperator{\res}{res}
\DeclareMathOperator{\cor}{cor}
\DeclareMathOperator{\conj}{conj}
\DeclareMathOperator{\adjg}{adj}
\DeclareMathOperator{\tr}{tr}
\DeclareMathOperator{\rank}{rank}
\DeclareMathOperator{\sgn}{sign}
\DeclareMathOperator{\fl}{fl}
\DeclareMathOperator{\im}{im}
\DeclareMathOperator{\coker}{coker}
\DeclareMathOperator{\soc}{soc}
\DeclareMathOperator{\rad}{rad}
\DeclareMathOperator{\Ext}{Ext}
\DeclareMathOperator{\Ann}{Ann}
\DeclareMathOperator{\MaxSpec}{MaxSpec}
\DeclareMathOperator{\Frob}{Frob}
\DeclareMathOperator{\supp}{supp}
\newcommand{\aug}{\varepsilon}            % augmentation C_0 -> Z

%----------------------------------------------------------------------
% The simplicial framework
%----------------------------------------------------------------------
\newcommand{\Ku}{K^{u}}                   % universal cover complex
\newcommand{\KR}{K/\Rel}                  % quotient complex
\newcommand{\prR}{\mathrm{pr}_{\Rel}}     % projection K -> K/\Rel
\newcommand{\pru}{\mathrm{pr}^{u}}        % projection K^u -> K/\Rel
\newcommand{\piOne}{\pi_1}
\newcommand{\LCS}{\mathbf{LCS}}
\newcommand{\Rep}{\mathbf{Rep}}
\newcommand{\LS}{\mathbf{LS}}
\newcommand{\cat}[1]{\mathbf{#1}}
\newcommand{\orient}[1]{o_{#1}}           % oriented q-simplices: \orient{q}(K)
% The distinguished lift of a simplex of K into the universal cover K^u, for the
% strict fundamental domain K \subseteq K^u of (H4)(iii).  The code spells this
% i(sigma); the document cannot, because i is sqrt(-1) (notation chapter, C15/C16).
\newcommand{\lift}[1]{\widetilde{#1}}
\newcommand{\Ceq}{C^{\bullet}_{\mathrm{eq}}}
\newcommand{\Cq}[1]{C^{#1}_{\mathrm{eq}}}
\newcommand{\Zq}[1]{Z^{#1}_{\mathrm{eq}}}
\newcommand{\Bq}[1]{B^{#1}_{\mathrm{eq}}}
\newcommand{\Heq}[1]{H^{#1}_{\mathrm{eq}}}

%----------------------------------------------------------------------
% Coefficients
%----------------------------------------------------------------------
\newcommand{\rhon}{\rho_n}                % the symmetric power representation
\newcommand{\rhodual}{\rho_n^{\vee}}      % the contragredient
\newcommand{\Vcoef}{V}                    % V := Sym^n(Z^2)
\newcommand{\VR}{V_{\cring}}              % V \otimes \cring
\newcommand{\Mn}{M_n}                     % Sym^n of F_ell^2 (synonym of \VR at R = F_ell)
\newcommand{\Dn}{D_n}                     % the divided-power / contragredient module
\newcommand{\cf}{f_{12}}                  % the normalised weight-12 cusp form (Ramanujan's Delta)

%----------------------------------------------------------------------
% The Hecke datum
%----------------------------------------------------------------------
\newcommand{\Th}{\Theta}                  % the fused cochain map Theta_g
\newcommand{\Thb}{\Theta^{\partial}}      % the boundary cochain map
\newcommand{\Xig}{\Xi_g}
\newcommand{\taub}{\tau_{\partial}}       % the boundary g-twisted chain datum
\newcommand{\Phig}{\Phi_g}                % the cone cochain map
\newcommand{\Amat}[1]{A^{(#1)}}           % cohomology-level Hecke matrix
\newcommand{\Acone}[1]{A_c^{(#1)}}
\newcommand{\Abdry}[1]{A_{\partial}^{(#1)}}
\newcommand{\Tdc}[1]{T_{\Ga #1 \Ga}}      % T_{Gamma g Gamma}
\newcommand{\diam}{T(p,p)}                % the diamond operator

%----------------------------------------------------------------------
% The pair, the cone, compact supports
%----------------------------------------------------------------------
\newcommand{\Xbulk}{X}                    % the Gamma-tessellation of the truncated half plane
\newcommand{\Xbdry}{\partial X}
\newcommand{\Kbdry}{K_{\partial}}         % the circle model 0-1-2-3
\newcommand{\rres}{r}                     % the restriction cochain map r = j^*
\newcommand{\Cone}[1]{C(#1)}              % the mapping cone
\newcommand{\Dcone}{D}                    % the cone differential
\newcommand{\coneinc}{\mathrm{inc}_c}     % C_eq(K_bdry) -> C(r)
\newcommand{\coneproj}{\mathrm{pr}_c}     % C(r) -> C_eq(K)[1]
\newcommand{\Hcpt}[1]{H^{#1}_{c}}
\newcommand{\etapair}{\eta}               % the pair homotopy
\newcommand{\hder}{h}                     % the derived cochain homotopy h = \tilde\eta
\newcommand{\Lam}{\Lambda}                % the unique C_0(bdry X) -> C_2(X) of eta-uniqueness

%----------------------------------------------------------------------
% The integral line (sec:integral, since [C] v1.8.0; notation
% sec:notation-integral)
%----------------------------------------------------------------------
\newcommand{\Iso}{\mathcal{I}}            % the isotropy-resolution model on X
\newcommand{\Isob}{\mathcal{I}^{\partial}} % its boundary subcomplex
\newcommand{\cmpI}{\mathrm{cmp}}          % the comparison map I_* -> C_*(X)
\newcommand{\CI}{C^{\bullet}_{\Iso}}      % Hom_{Z[Gamma]}(I_*, V_R)
\newcommand{\Hfr}[1]{H^{#1}_{\mathrm{fr}}} % H^q(Gamma;V) modulo torsion
\newcommand{\Zsix}{\widehat{\Z}_{6}}      % Z_2 x Z_3
\newcommand{\Qsix}{\widehat{\Q}_{6}}      % Q_2 x Q_3

%----------------------------------------------------------------------
% The Dickson line and the ordinary projector (sec:dicksontheta,
% sec:ordinary, since [C] v1.11.0-pre.1; notation sec:notation-dickson).
% Sans-serif L, D: the plain letters are taken (L the local system, D_n the
% contragredient module), rule A5.
%----------------------------------------------------------------------
\newcommand{\DkL}{\mathsf{L}}             % the first Dickson invariant x^ell y - x y^ell (the code's L)
\newcommand{\DkD}{\mathsf{D}}             % the second Dickson invariant, degree ell^2 - ell
\newcommand{\thL}{\theta_{\DkL}}          % multiplication by L on cochains and cohomology
\newcommand{\nplus}{n^{+}}               % n + ell + 1 (N is the Hecke level only)
\newcommand{\Qn}{Q_{n}}                   % the quotient Sym^{n+ell+1} / L Sym^n
\newcommand{\eord}{e_{\mathrm{ord}}}      % the ordinary projector of def:ordinary (e alone is lem:bulkperiodic's class)

%----------------------------------------------------------------------
% The involution
%----------------------------------------------------------------------
\newcommand{\inv}{\iota}                  % the involution diag(-1,1) / z -> -zbar
\newcommand{\invK}{\iota_K}               % its simplicial realisation on the 7-vertex model

%----------------------------------------------------------------------
% The Hecke algebra and its local pieces
%----------------------------------------------------------------------
\newcommand{\Talg}{\mathbb{T}}
\newcommand{\midl}{\mathfrak{m}}
\newcommand{\Htot}{H_{\mathrm{tot}}}

%----------------------------------------------------------------------
% Misc
%----------------------------------------------------------------------
\newcommand{\abs}[1]{\left| #1 \right|}
\newcommand{\smat}[4]{\begin{pmatrix} #1 & #2 \\ #3 & #4 \end{pmatrix}}
\newcommand{\code}[1]{\texttt{#1}}
\newcommand{\file}[1]{\texttt{#1}}

%----------------------------------------------------------------------
% Cross-document citation.  \HA{} survives for documents that still cite the
% OLD hecke_action.tex; inside the unified contract, use \ref{} instead.
%----------------------------------------------------------------------
\newcommand{\HA}[1]{[C,~\texttt{#1}]}
\newcommand{\vfive}[1]{[v5,~\texttt{#1}]}
\newcommand{\Espec}[1]{[E,~\texttt{#1}]}

%----------------------------------------------------------------------
% RETIRED SYMBOLS.  These macros existed in the pre-consolidation corpus with
% a meaning the notation chapter has retired.  Using one is an error, not a
% typographic choice.  (\K, \R, \C, \pR, \pu, \image were in the copied macro
% blocks of hecke_action.tex and monodromy_v5.tex.)
%----------------------------------------------------------------------
\newcommand{\gltwoRETIRED}[2]{%
  \PackageError{gl2_preamble}{Retired macro \string#1 -- #2}{%
    See the notation chapter (sec:notation) for the replacement.}}
\newcommand{\K}{\gltwoRETIRED{K}{ambiguous complex name; write K directly}}
\newcommand{\R}{\gltwoRETIRED{R}{R is the coefficient ring; write \string\Rl\space for the reals}}
\newcommand{\C}{\gltwoRETIRED{C}{write \string\Cx\space for the complex numbers}}
\newcommand{\pR}{\gltwoRETIRED{pR}{write \string\prR}}
\newcommand{\pu}{\gltwoRETIRED{pu}{write \string\pru}}
\newcommand{\image}{\gltwoRETIRED{image}{write \string\im}}
\newcommand{\M}{\gltwoRETIRED{M}{write \string\Mat\space for the matrix algebra}}

\geometry{a4paper, margin=1in}

% v1.17.0-pre.1 (Lean programme [LP] Q9, owner-ruled 2026-10-01): the formal-verification tag.  First used at v1.17.0 (the slot-17
% close, milestone M1), clause-scoped, beside \PROVED, never in place of it (sec:status).
\newcommand{\LEAN}[1]{\textsf{[formally verified: [GT,~\texttt{\detokenize{#1}}]]}}
\newcommand{\LEANlit}[2]{\textsf{[formally verified modulo #2: [GT,~\texttt{\detokenize{#1}}]]}}

% v1.22.0-pre.1 (sec:nebwildsec, P10-C): three macros of [U]'s local block used by the refereed text, provided here so that it is inserted verbatim.
\providecommand{\Zl}{\Z_\ell}
\providecommand{\Ql}{\Q_\ell}
\providecommand{\VERIFYloc}{\textsf{VERIFY}}          % locator not checked against the text (rule A2)

\title{Local Systems, Equivariant Cohomology and the Integral Hecke Action\\
for $\GL(2)/\SL_2(\Z)$\\[1ex]
\large The Binding Mathematical Contract}
\author{Fabian Januszewski\\ Universit\"at Paderborn\\ \texttt{fabian.januszewski@math.uni-paderborn.de}}
\date{Version 1.26.4 --- 2026-11-10 (revision history: Appendix~\ref{app:changelog})}

\begin{document}
\maketitle

\begin{abstract}
This is the binding mathematical contract for the implementation: it fixes the objects it constructs
and everything asserted about them. It sets out the simplicial model $(K,\Rel)$ of
$\SL_2(\Z)\backslash\Hyp$, the matrix local systems on it, above all the symmetric powers $\Sym^n$ of the
standard representation, and the equivariant cochain complex and its cohomology. It defines the integral
Hecke datum and the operators $T_p$, $T_{p^{\nu}}$, $T_N$, $\diam$ and
$U_\ell$ in all degrees; the involution $\inv$; and the boundary datum and the mapping cone, which carry
the Hecke action to the boundary and to compact supports, with the six groups and their long exact
sequence. It states how these are realised and certified: a combinatorial reduction primitive with
integer bounds, a certified floating-point filter, numbered correctness obligations and a worked example.
Later parts treat the integral line and descent over $\Z$; the Dickson invariant, $\thL$, the ordinary
projector and Hida theory at level one; level $N$ and nebentypus; the $\Z/\ell^{K}$ Hecke algebra and
its saturation; the $\ell$-local isotropy model, descent at level $N$ and the Hecke algebra over $\Z$;
and the Moerdijk--Pronk model. Code maps say what is implemented. Every assertion carries a status tag;
Appendix~\ref{app:status} indexes every numbered statement, and Appendix~\ref{app:changelog} records what each version changed.
\end{abstract}

\tableofcontents

%======================================================================
\section{Scope, status, and how to read this document}\label{sec:scope}
%======================================================================
% SOURCE:  hecke_action.tex:99-160 (RH) + 00-unified-contract-architecture.md
%          SS1.3, SS2, SS7, SS8.5 (NEW)
% NOTATION: rows Rel/cring, nu, T_N, sigma of sec:notation-crosswalk applied
% ERRATA:  E16, E17 (App. A)

\paragraph{What this chapter fixes.}
What this document is normative about and what it is not; the vocabulary in which
every later statement declares its epistemic status; the rule that makes a
\code{\textbackslash label} of this document a stable API for the C++ source that
cites it; the rules its own authors are bound by; and the version convention under
which a statement of the form ``as of [C] v$X.Y$'' is verifiable.

\subsection*{What this document is}

This is the \emph{binding mathematical contract} for the implementation in this
repository: the objects it constructs, and everything asserted about them. It fixes
the simplicial model $(K,\Rel)$ of $\SL_2(\Z)\backslash\Hyp$, the matrix local systems
on it, the equivariant cochain complex and its cohomology $H^q(K,\Rel;L)$, the
integral Hecke datum and the operators $T_p$, $T_{p^{\nu}}$, $T_N$, $\diam$ and
$U_\ell$ in \emph{all} cohomological degrees, the boundary and cone constructions that
carry that action to $H^q(\partial)$ and $H^q_c$, the involution $\inv$, the
combinatorial reduction primitive with its integer bounds, the certified
floating-point filter, and the numbered correctness obligations that a conforming
implementation must discharge.

It is cited as \textbf{[C]}. That short key is not decorative: $53$ C++ source files
already cite statements of this document in the form \code{[C, alg:tau]}, at $519$
sites across $71$ files. The key and every cited label are therefore preserved
verbatim (\S\ref{sec:labelpolicy}, Appendix~\ref{app:labelindex}).

\subsection*{What is normative here, and what is not}

The corpus this document consolidates was split along a line that put predictions and
plain facts about the objects on the same side. The line is redrawn here:

\begin{contractbox}
\textbf{[C]} --- this document --- contains the objects, their construction, and
everything proved about them. Every statement in it is normative. \\[0.5ex]
\textbf{[E]} --- the companion \file{hecke\_extensions.tex} --- is the
\emph{programme specification}: what to compute with those objects and how to certify
it. It is versioned, expected to churn, and cited from no source file.
\end{contractbox}

\noindent
Consequently the following moved \emph{into} this document and are no longer sought in
[E]: the six groups $X_1,\dots,X_6$ and the long exact sequence relating them, the
boundary groups and their characters $\chi_b$, Poincar\'e duality, the amalgam
presentation of $\Ga$ and the resulting formula for $\dim H^1$, and the
Dickson/Bockstein statement (\S\ref{sec:six}). They are facts about the objects and
they are the numerical oracles that Obligation~\ref{obl:oracle} reads; they are not a
programme. What remains in [E]: the Hecke algebra $\Talg$ and its local pieces
$\Talg_{\midl}$, the separation and generation bounds, the classification data and the
staged algorithms, the scalar extension to $\Flf$, the output schema, and the
predictions. Appendix~\ref{app:errata}, entry~\ref{err:e18}, records the citation
changes this move causes.

Two things this document deliberately does \emph{not} do. It does not assert anything
about a coefficient ring beyond the hypotheses written into the individual statements
that use one --- $\cring=\Fl$, $\Z/\ell^{e}$ and $\Z/m$ are distinguished at the point
of use, never globally (rule~\ref{A1}). And it fixes the $\ell=2$ truncation policy
only as a \emph{policy}: \ref{open:ell2} (OC-1), decided by the owner on 2026-09-16 as
\ref{oc1a} and stated in \S\ref{sec:ell2} with its alternatives kept reachable, rather
than silently assumed away; the stabilisation theorem \ref{oc1c} that would replace the
policy by a proved bound is half proved there (Lemma~\ref{lem:bulkperiodic}) and half
open (Remark~\ref{rmk:bulkperiodicscope}).

\subsection*{The superseded documents}

\begin{itemize}
\item \file{hecke\_action.tex} (\textbf{[C]} as it stood; PDF of 2026-07-03) is
  superseded \emph{by this document}, which keeps its key, its $67$ labels and its
  statements, with the corrections of Appendix~\ref{app:errata} applied.
\item \file{monodromy\_v5.tex} (\textbf{[v5]}) is superseded and frozen. Its framework
  --- the pair $(K,\Rel)$, the universal cover $\Ku$, the three equivalent categories
  and the explicit equivariant complex --- is re-homed here in full
  (\S\S\ref{sec:framework}--\ref{sec:systematiccomputation}), and its labels come with
  it. This is not tidiness: five C++ source files cite
  \code{[C, thm:explicit\_complex\_equiv]}, and until this document existed that
  citation dangled, because the label lived in [v5]
  (Appendix~\ref{app:errata}, entry~\ref{err:e7}). What stays in the frozen [v5] is its
  \S6 M\"obius pedagogy and its \S12, the latter superseded in full
  (entries~\ref{err:e10}, \ref{err:e19}, \ref{err:e20}).
\item The earlier document cited [v5] \emph{by statement number} --- Prop.~24, Thm.~26,
  Def.~28/29, Thm.~31/33, Thm.~35, Def.~39/40, Thm.~44. Most of those statements carry
  no label and [v5] fixes no version, so those references were not stable. Every one of
  them is now an internal \code{\textbackslash ref} of this document.
\end{itemize}

\subsection*{How to read this document}

Part~I (\S\S\ref{sec:framework}--\ref{sec:systematiccomputation}) is the framework: the
model $(K,\Rel)$, local systems, the equivariant complex, and the isomorphism $\Psi$
through which every later formula is transported onto the fixed complex. Part~II
(\S\S\ref{sec:conventions}--\ref{sec:involution}) fixes the arithmetic conventions, the
coefficient module $\Sym^n$ with its contragredient, and the involution. Part~III
(\S\S\ref{sec:heckedatum}--\ref{sec:equivalence}) is the Hecke datum and the primary
definition of the operators. Part~IV
(\S\S\ref{sec:boundarydatum}--\ref{sec:orbifoldmodel}) carries that action to the
boundary, the cone and compact supports. Part~V
(\S\S\ref{sec:reduction}--\ref{sec:machinery}) is realisation and certification: the
reduction primitive, the filter, the obligations, the worked example and the code map.
Part~VI (\S\ref{sec:integral}, since v1.8.0) is the integral line: what each coefficient
ring names, the free model on the unrefined tessellation, and the descent over $\Z$; it
is placed last so that no earlier chapter number moves, and continues
\S\ref{sec:orbifoldmodel} in substance.
Part~VII (\S\S\ref{sec:dicksontheta}--\ref{sec:ordinary}, since v1.11.0-pre.1, \S\ref{sec:hida}, since
v1.12.0-pre.1, and \S\ref{sec:hidafull}, since v1.13.0-pre.1) is multiplication by
the Dickson invariant, the operator $\thL$ and its long exact sequence, the ordinary projector, the foundation of Hida theory at
level~$1$, and its continuation: level structure by Shapiro on the fixed model, the stabilisation to level~$\ell$, control
and the lattice datum; it is placed after Part~VI for the same reason.
Part~VIII (\S\ref{sec:levelN}, since v1.14.0-pre.1, and \S\ref{sec:nebentypus}) is level $N$ for every $N$ --- the composite coset table, the cokernel presentation and its Hecke
action by Merel's matrices --- and nebentypus; Part~IX (\S\S\ref{sec:saturation}--\ref{sec:mpmodel}) is the $\Z/\ell^{K}$ Hecke algebra, the
$\ell$-local isotropy model with the descent at level $N$ and the Hecke algebra over $\Z$, and the Moerdijk--Pronk model.
Chapter~\ref{sec:notation} is \emph{normative and exhaustive} and binds all nine parts;
a reader who skips it will mis-read formulas, because several symbols here mean
something different from what they meant in the superseded documents.

\subsection{Status vocabulary}\label{sec:status}

Every assertion that is not a definition carries \emph{exactly one} status tag, and
the tag is part of the statement.

\begin{center}
\begin{tabular}{@{}lll@{}}
\toprule
tag & macro & what it claims \\
\midrule
\PROVED & \code{\textbackslash PROVED} & proved in this document, from the hypotheses in its own statement \\
\STD{\ldots} & \code{\textbackslash STD\{\ldots\}} & standard; discharged by the \emph{precise} citation in the argument \\
\CONJ & \code{\textbackslash CONJ} & believed, not proved here and not discharged by a citation \\
\COMP & \code{\textbackslash COMP} & verified computationally, over the stated range, by a named validation \\
\OPENSTAT & \code{\textbackslash OPENSTAT} & open: the statement is not settled and the document says so \\
\bottomrule
\end{tabular}
\end{center}

\begin{contractbox}
\code{\textbackslash STD} takes a \textbf{mandatory} argument, and that argument must
be a precise citation --- author, work, and the section or theorem number --- not a
bare pointer to a book. \code{[standard]} with no pointer does not compile. Where no
precise citation exists, the tag is \CONJ{} and the reason is written out.
\end{contractbox}

\noindent
\COMP{} never upgrades to \PROVED. A computation over a finite grid is evidence about
that grid; the range it covers is stated with the tag, and the validation that produces
it is named, so the claim can be re-run. Appendix~\ref{app:status} indexes every
numbered statement with its tag.

\noindent
\textbf{A sixth marker, not a status tag} (v1.17.0-pre.1; the Lean programme's question Q9, ruled by the owner on 2026-10-01).
\code{\textbackslash LEAN\{decl\}} and \code{\textbackslash LEANlit\{decl\}\{hypotheses\}} record that a statement already tagged \PROVED{} is
formally verified in the Lean ground truth [GT] (\file{lean/blueprint/src/gl2\_groundtruth.tex}) --- with no literature hypothesis, or modulo the
named literature hypotheses of [GT]. The marker stands \emph{beside} \PROVED{}, never in place of it, and never beside \CONJ{}, \COMP{} or
\OPENSTAT{} alone; it does not change the exactly-one-tag rule above, and Appendix~\ref{app:status} does not count it. \emph{Amended at v1.24.0-rc.2} (the rule as already
applied, e.g.\ at Remark~\ref{rmk:orbifold} and Definition~\ref{def:hecke}; the referee \file{docs/implementation/slot25/RC124\_referee.md} F2): the marker stands
beside \PROVED{} or \STD{} --- on a \STD{} statement it records that [GT] proves the clause; it neither turns \STD{} into \PROVED{} nor replaces the
citation --- and, clause-scoped, on a definition or a displayed equation; never beside \CONJ{}, \COMP{} or \OPENSTAT{} alone. It was defined at
v1.17.0-pre.1; the first markers enter at v1.17.0, the release that records the Lean milestone M1.
\textbf{The marker is clause-scoped.} It stands at the clause, or the sentence, that [GT] formalises, never at the head of a statement of which
only some clauses are formalised; a parenthesis after it narrows it further where [GT] states a clause in a weaker or different form (for
instance ``as the injectivity of the canonical surjection''). A clause of a tagged statement that carries no marker is \emph{not} claimed
formally verified, whatever its \PROVED{} tag. A \emph{definition} carries no status tag; on a definition the marker records only that [GT]
defines the same objects (with the stated differences of indexing). The list of v1.17.0 is
\file{docs/implementation/slot17/lean\_taggable\_labels.md} (the slot-17 Lean audit F3's, after its repairs): $42$ markers on $20$ labels;
Corollary~\ref{cor:stablecert} takes \code{\textbackslash LEANlit}, its [GT] proof carrying Lemma~\ref{lem:heckerankone} as a hypothesis
(Lemma~\ref{lem:heckerankone} itself is not tagged), and Proposition~\ref{prop:hidaalgconst}(1) takes \code{\textbackslash LEANlit}
modulo Hida's control theorem. The list of v1.18.0 adds the slot-18 Lean audit F3's \S8
(\file{docs/implementation/slot18/L18-6\_F3\_audit.md}, milestone M2) with its repairs (\file{docs/implementation/slot18/L18-6\_repairs.md}:
Lemma~\ref{lem:model7coincidences}(2), Definition~\ref{def:model7rel}'s rows, \ref{FD1}): $81$ markers on $42$ labels in all ($75$
\code{\textbackslash LEAN}, $6$ \code{\textbackslash LEANlit}); Lemma~\ref{lem:heckerankone} and Proposition~\ref{prop:cohforms}(3), (4)'s second surjection take
\code{\textbackslash LEANlit} modulo Eichler--Shimura at level one (the [GT] hypotheses \texttt{hES}, \texttt{hF}); Corollary~\ref{cor:stablecert}'s
\code{\textbackslash LEANlit} stays (its [GT] proof still carries the rank hypothesis, not yet discharged from Lemma~\ref{lem:heckerankone}).
The list of v1.19.0 adds the slot-19 Lean audit F3's \S8 (\file{docs/implementation/slot19/L19-6\_F3\_audit.md}, milestone M3) as amended by its
repairs (\file{docs/implementation/slot19/L19-6\_repairs.md} \S3) and by the pair-hypothesis results of slot-19 L19-4
(\file{docs/implementation/slot19/L19-4\_report.md} \S3): $98$ markers on $48$ labels in all ($93$ \code{\textbackslash LEAN}, $5$
\code{\textbackslash LEANlit}; the boundary hypotheses \ref{B1}--\ref{B3} and the standing hypotheses of \S\ref{sec:cone} counted under their sections).
\ref{H4}(ii) is now tagged in the form this document uses it, the acyclicity \eqref{eq:contractibleuse}; Remark~\ref{rmk:orbifold},
Lemma~\ref{lem:tauexists}(2), Theorem~\ref{thm:welldef}(b) and the whole of (c), and the cone chapter's hypotheses on the model are tagged;
Corollary~\ref{cor:stablecert}(1) takes \code{\textbackslash LEANlit} modulo Eichler--Shimura at level one (the [GT] hypothesis \texttt{hES}, the rank
hypothesis now discharged), its (2) \code{\textbackslash LEAN}.
The list of v1.20.0 adds the slot-20 Lean audit F3's \S8 (\file{docs/implementation/slot20/L20-6\_F3\_audit.md}, milestone M4) as amended by its
repairs (\file{docs/implementation/slot20/L20-6\_repairs.md}): $140$ markers on $65$ labels in all ($126$ \code{\textbackslash LEAN}, $14$
\code{\textbackslash LEANlit}). The new \code{\textbackslash LEANlit} markers (Theorem~\ref{thm:relations}, Obligation~\ref{obl:algebra},
Theorem~\ref{thm:nebrelations}) are modulo the Hecke-ring identities of \cite[Prop.~3.8, Prop.~3.16, Prop.~3.17, Thm.~3.24, Prop.~3.31, Prop.~3.33, Thm.~3.34, \S3.2--3.3]{shimura}
as the [GT] hypotheses \texttt{hTp}, \texttt{hcomm}, \texttt{hcop}, \texttt{hTpchi0}, \texttt{hTpchi}, \texttt{hUchi}. \textbf{Since v1.21.0-pre.1 their locators are
verified} on the owner's scan of Shimura Ch.~3 (step R20-L, \file{docs/implementation/slot21/R20-L\_report.md}), with the transport of
\cite[\S1.1]{ash_stevens} read on the GDZ page images; R20-L restated the locators, not the clauses: \texttt{hcop} is Prop.~3.16 (Thm.~3.24(3) is
the full-operator form), \texttt{hTp} is Thm.~3.24(5) with Prop.~3.17 (Prop.~3.30, the level-$N$-to-level-$1$ homomorphism, is not used),
\texttt{hcomm} is Prop.~3.8 as applied on p.~56 (Thm.~3.24 does not state it), and the nebentypus identities are Shimura's character-free
identities in $R(\Ga_0(N),\De_0(N))$ (Prop.~3.31, Thm.~3.34(3) with Prop.~3.33), $\chi$ entering only through $T(p,p)$. Shimura's notation
($T(p^{\nu})$ the full sum, $T(1,p^{\nu})$ the primitive coset) is confirmed by his Thm.~3.24(1), (2). The full operator is formalised only in the
assembled form \eqref{eq:fullassembly}: \eqref{eq:fullop} carries no marker. Theorem~\ref{thm:classical} and the equation \eqref{eq:corconjres}
carry none either --- [GT] proves only a weak form, with $\cor$ characterised by $\cor\circ\res=[\Ga:\Ga'']$ alone; the maps $\res$, $(\conj_g)_{\ast}$
and their compatible pair are formalised (\texttt{def:m4-classical}), but \eqref{eq:corconjres} carries no marker, not even on them. The level-$N$ and nebentypus
markers are on the model and, for the operators, for a fixed table and $\tau$ (their independence of both is not formalised).
The list of v1.22.0-pre.6 adds \S\ref{sec:mpmodel}'s markers, clause-scoped per \file{docs/implementation/slot22/MPL1\_report.md} v1.1 \S5 ([GT] Part~G12,
the slot-22 Lean packages MP-L1 / MP-L2): $30$ markers on $11$ labels, all \code{\textbackslash LEAN}, so $170$ markers on $76$ labels in all ($156$
\code{\textbackslash LEAN}, $14$ \code{\textbackslash LEANlit}). They cover the descriptor predicates, the MP set, the \emph{unnormalised} MP complex
$C^{\bullet}_{\mathrm{un}}$ and its comparison theorem, discharged in full for $\mathsf M_7$ and $\mathsf R_{\Z}$; none is on Theorem~\ref{thm:mpnatural},
Propositions~\ref{prop:mpalldegrees}, \ref{prop:mpbredon}, \ref{prop:mpboundary}, \ref{prop:mprescor} or Remark~\ref{rmk:mphybrid} (not formalised), nor on the
normalised complex.
The list of v1.22.0-pre.7 adds $25$ markers of \S\ref{sec:mpmodel}, clause-scoped per \file{docs/implementation/slot22/MPL3\_report.md} v1.0 \S5 (the slot-22 Lean package
MP-L3), all \code{\textbackslash LEAN}, on Definitions~\ref{def:mpmorphism} ($4$), \ref{def:mpset} ($1$), \ref{def:mpfree} ($2$), Theorems~\ref{thm:mpnatural} ($10$, and $3$ on its
\COMP\textsuperscript{r} record), \ref{thm:mpcomparison} ($2$) and Proposition~\ref{prop:mpbredon} ($3$), so $195$ markers on $79$ labels in all ($181$ \code{\textbackslash LEAN},
$14$ \code{\textbackslash LEANlit}). They cover the morphisms with the label factor, naturality and refinement invariance for $C^{\bullet}_{\mathrm{un}}$ (label independence on
cohomology), $\varepsilon^{\mathsf{MP}}$ with the ``if'' of Proposition~\ref{prop:mpbredon}(2), and the normalised complex with its equivalence to $C^{\bullet}_{\mathrm{un}}$; none is on
Propositions~\ref{prop:mpalldegrees}, \ref{prop:mpboundary}, \ref{prop:mprescor} or Remark~\ref{rmk:mphybrid}.
v1.22.0-pre.8 adds no marker ($195$ on $79$ labels); Remark~\ref{rmk:mphecke} (the MP-model Hecke operators) carries none.
v1.22.0-pre.9 to pre.11 add none. v1.22.0-pre.12 adds the LR-Y markers (\file{docs/implementation/slot22/LRY\_phase3\_report.md} \S C as repaired by
\file{LRY3\_audit.md}): $7$ on Claim~\ref{cl:amalgam}(1),(2), Lemma~\ref{lem:bulkperiodic}(2) and the new Remark~\ref{rmk:treeconeres}, so $202$ markers on $82$
labels ($188$ \code{\textbackslash LEAN}, $14$ \code{\textbackslash LEANlit}). v1.22.0-pre.13 adds $4$ (LR-Y phase 4, \file{LRY4\_audit.md} \S6) on
Claim~\ref{cl:amalgam}(2) and Lemma~\ref{lem:bulkperiodic}(2), labels that already carried markers, so $206$ on $82$ ($192$ \code{\textbackslash LEAN}). v1.22.0-pre.14 adds
$12$ (the LV dimension results, worded within \file{LV2\_audit.md} \S8) on Lemma~\ref{lem:centralminusone} (its last clause, for fields only),
Proposition~\ref{prop:oddnltwo}(2) (the value of $h^{0}$ at the certified cells), Proposition~\ref{prop:ellthreeclosed}(4) (for fields only) and the new
Remark~\ref{rmk:lvdimsltwothree} ($9$), so $218$ markers on $86$ labels in all ($204$ \code{\textbackslash LEAN}, $14$ \code{\textbackslash LEANlit}); the Lean cells
outside [GT]'s census (\texttt{L23Cells}, $n>64$) are cited as data, never by a marker, and Proposition~\ref{prop:ltwothreescalar} carries none.
v1.22.0-pre.15 adds none. v1.22.0-pre.16 adds $4$ (LR-Y phase 5, \file{LRY\_phase5\_report.md} \S D as repaired by \file{LRY5\_audit.md}) on
Claim~\ref{cl:amalgam}(1),(2), a label that already carried markers, so $222$ markers on $86$ labels in all ($208$ \code{\textbackslash LEAN}, $14$ \code{\textbackslash LEANlit}).
v1.24.0-rc.1 adds $7$ (the slot-24 Lean package M6's proposal D5, \file{docs/implementation/slot24/LEAN\_M6\_report.md} \S4 as corrected by
\file{LEAN\_M6\_audit.md} F1, F3; the owner's question Q4, approved in principle on 2026-10-18): $3$ on Theorem~\ref{thm:classical} and $4$ on
\eqref{eq:corconjres}, so $229$ markers on $88$ labels in all ($215$ \code{\textbackslash LEAN}, $14$ \code{\textbackslash LEANlit}). [GT] Part~M6 now proves the
group-cohomological form on the model ($\Theta$ is the canonical comparison, $\cor$ is Brown's transfer, the operator is independent of the resolution and of $\tau$), so
the sentence of v1.20.0 above (``carry none either'') describes that version. Both statements are tagged \STD{\cite[Ch.~III.9--10]{brown}}; the markers stand
beside \STD{} under the rule as amended at v1.24.0-rc.2 above, clause by clause in their parentheses, and move no tag. The third marker of that proposal,
on the witness of Proposition~\ref{prop:mpbredon}(2), is not applied (it is not part of Q4).

\subsection{Label stability is part of the contract}\label{sec:labelpolicy}

A \code{\textbackslash label} of this document is an API, consumed by source code that
this document does not control.

\begin{enumerate}[label=(L\arabic*)]
\item\label{L1} A label is never renamed and never deleted.
\item\label{L2} A label is never retargeted to a statement that asserts something
  different. If what a statement says changes, the label stays with the statement whose
  claim it named, and the change is recorded in Appendix~\ref{app:errata}.
\item\label{L3} A retired statement keeps its label, attached to the erratum that
  retires it. Removal is not a way to fix a wrong statement; correction is.
\item\label{L4} A new label never reuses a retired string.
\item\label{L5} New labels use the scheme \code{kind:stem} with
  \code{kind} one of \code{sec}, \code{def}, \code{thm}, \code{lem}, \code{prop},
  \code{cor}, \code{rmk}, \code{alg}, \code{obl}, \code{cert}, \code{cl}, \code{eq},
  \code{app}, \code{open}, \code{err}. No global prefix: a prefix would break every
  citing source file at once, for nothing.
\item\label{L6} The policy is machine-checked, not merely declared
  (\S\ref{sec:machinery}): every label cited from source or from a project document is
  dereferenced in Appendix~\ref{app:labelindex}, so a missing or duplicated label fails
  the build, and a census script fails on a \code{[C, \ldots]} citation this document
  does not define.
\end{enumerate}

\subsection{Authoring rules}\label{sec:authoring}

These bind the authors of this document. They are cheap now and expensive later; the
cautionary examples behind each are real and are recorded in Appendix~\ref{app:errata}.

\begin{enumerate}[label=(A\arabic*)]
\item\label{A1} \textbf{Every hypothesis at the point of use.} A statement carries every
  hypothesis it needs \emph{in its own statement}, even when that hypothesis was
  declared for the chapter. Chapter-level standing hypotheses may still be declared; a
  statement that uses one repeats it. (Motivating failure: ``$\Xbulk$ is contractible
  and two-dimensional'' was load-bearing for the uniqueness of $\Lam$ and for the
  existence of the fills, and was unwritten --- so a reader could not tell the argument
  from a wish. A second: entry~\ref{err:e3}, a realisation theorem proved over a field
  and then applied over $\Z/\ell^{e}$.)
\item\label{A2} \textbf{No claim rests on ``standard''.} See \S\ref{sec:status}: every
  \STD{} carries a precise citation, and the tag will not compile without one.
\item\label{A3} \textbf{Definitions are constructive where the code is constructive.}
  Where the implementation builds an object by a terminating procedure on explicit data
  --- as Algorithm~\ref{alg:tau} builds $\tau$ --- the definition \emph{is} that
  procedure together with the properties it guarantees, and existence and uniqueness are
  then theorems \emph{about the construction}. ``Choose any $\tau$ satisfying
  (T1)--(T3)'' with existence deferred to an abstract lemma is not the definition this
  document gives, because it is not the object the code produces.
\item\label{A4} \textbf{Every proof obligation is a checkable predicate on emitted
  data.} An obligation states: the data it reads (which stored table, which matrix, in
  which basis), the predicate, the ring the predicate is evaluated over, and what
  happens on failure (abort / record / diagnostic). ``The tables must be consistent'' is
  not an obligation; ``$\partial\etapair+\etapair\partial = j\taub - \tau j$ on every
  stored representative, both sides evaluated through each table's own (T1) helper, over $\Z$;
  a failing table aborts'' is.
\item\label{A5} \textbf{One symbol, one meaning, document-wide} ---
  \S\ref{sec:notation}, enforced by the retired macros of \file{gl2\_preamble.tex},
  which raise a \LaTeX{} error rather than typeset. This is not a style rule: three
  incompatible meanings of $\De$, $\tau$ and $\inv$ coexisted in the corpus, two of them
  inside a single document.
\item\label{A6} \textbf{Diagnostics are never branched on.} Where the implementation
  records a quantity that could be mistaken for a decision procedure ---
  \code{cochain\_commutes}, \code{eta\_zero}, \code{solve\_regime} --- this document says
  so explicitly. The entrywise identity $\rres\,\Th_g = \Thb_g\,\rres$ holding on the
  stored data does \emph{not} license $\hder=0$, and the text must make that impossible
  to misread.
\end{enumerate}

\subsection{Version marker and changelog}\label{sec:versioning}

Neither superseded document carried a version marker: both used
\code{\textbackslash date\{\textbackslash today\}}, so every rebuild silently re-dated
the PDF and no ``as of'' statement about either was verifiable
(Appendix~\ref{app:errata}, entry~\ref{err:e17}). This document fixes a version:

\begin{enumerate}[label=(V\arabic*)]
\item\label{V1} The date line is \code{Version $X.Y$ --- YYYY-MM-DD} and a status clause, never
  \code{\textbackslash today}. The supersession record --- one entry per version --- is the first
  subsection of Appendix~\ref{app:changelog}.
\item\label{V2} \textbf{$X$ increments on any change to a normative statement}; $Y$
  increments on prose, on recording an erratum, and on new non-normative material.
\item\label{V3} \S\emph{Changelog} --- Appendix~\ref{app:changelog} --- carries one
  entry per version and, within it, \textbf{one line per changed statement, naming its
  label}. ``Editorial pass'' is not a changelog entry: a contract whose statements move
  must say which ones moved. The rules under which that appendix is kept, including what
  counts as one entry and how a statement's provenance is marked, are stated at its
  head.
\item\label{V4} \file{hecke\_action.tex} and \file{monodromy\_v5.tex} each receive a
  frozen version marker and a superseded-by banner when this document ships.
  \textbf{Done at v1.0.0}: \file{hecke\_action.tex} is frozen at Version~1.0 (2026-07-02)
  and \file{monodromy\_v5.tex} at Version~5.0 (2026-07-03), each carrying a banner that
  names this document as the successor, says that a \code{[C,~<label>]} citation resolves
  here and not there, and directs the reader to Appendix~\ref{app:errata}. The content of
  neither was otherwise touched: they are the historical record. [v5]'s banner also says
  where \code{thm:explicit\_complex\_equiv} now lives, that being the label its five C++
  citing sites use (entry~\ref{err:e7}), and that \S12's $\beta$ is not re-homed
  (entry~\ref{err:e19}).
\item\label{V5} \textbf{Before completion the version is $0.X.Y$.} While this document
  is incomplete in the sense below, the leading field is $0$ and marks that fact; $X$ and
  $Y$ are the fields of \ref{V1} and carry \ref{V2}'s bump rules unchanged. The leading
  field becomes $1$, and $0.X.Y$ then reads $1.X.Y$ with no renumbering of the history,
  when \textbf{both} of the following hold:
  \begin{enumerate}[label=(\roman*)]
  \item\label{V5i} \textbf{no chapter of this document is a stub}; and
  \item\label{V5ii} \textbf{the label namespace is unambiguous} --- \ref{V4}'s two
    freezes have been done, and no label defined here is also defined by a live
    companion document.
  \end{enumerate}
  \textbf{Criterion \ref{V5i} was met at v0.19.0} and \textbf{\ref{V5ii} at v1.0.0}
  (Appendix~\ref{app:changelog}). \ref{V5ii}'s obstruction was concrete rather than
  ceremonial and was discharged concretely. Nine labels of \S\ref{sec:six} ---
  \code{sec:six}, \code{cl:amalgam}, \code{cl:dimH1}, \code{cl:bockstein},
  \code{cl:boundary}, \code{cl:les}, \code{eq:les}, \code{cl:chi}, \code{cl:pd} --- were
  defined \emph{both} here, by the scope cut of entry~\ref{err:e18}, \emph{and} in the
  live companion \file{hecke\_extensions.tex}, so that a reference \code{[C,~cl:pd]} did
  not resolve to a unique statement and \ref{L6} --- the policy is machine-checked, not
  merely declared --- could not be honoured across the corpus. \file{hecke\_extensions.tex}
  v1.3 removes its nine copies and cites this document instead, and \ref{V4}'s two freezes
  are done, so the criterion holds; the evidence is
  Remark~\ref{rmk:labelcensus}. The refinement of \ref{V1}'s form that the leading field
  represents is kept in the numbering, not dropped: ``as of [C] v$0.X.Y$'' says which
  normative revision is meant \emph{and} that the document was incomplete when it was
  made, and $0.19.0$ reads as $0.19.0$ still. A leading $1$ says only that this document
  is complete and its namespace unambiguous. It does \emph{not} say the programme is
  finished --- two chapters state objects nothing implements
  (\S\ref{sec:notimplemented}; three until v1.3.3), and \ref{open:ell2}, decided at
  v1.3.0, governs a construction none of them has --- and \S\ref{app:changelog-gaps}
  keeps that list.
\end{enumerate}

% SPDX-FileCopyrightText: 2024-2026 Fabian Januszewski
% SPDX-License-Identifier: CC-BY-4.0

%======================================================================
% gl2_notation.tex --- Chapter 2 of the unified contract.
% DRAFTED IN FULL (design pass, 2026-09-15).  Binds every later chapter.
%======================================================================
\section{Notation, fixed once}\label{sec:notation}

This chapter is \emph{normative and exhaustive}. A symbol used anywhere in this
document has exactly the meaning fixed here, in every chapter, in every proof, and
in every displayed formula. A symbol not listed here is either local to the
statement that introduces it --- in which case that statement introduces it
explicitly --- or it is a mistake.

The machine-readable form of this chapter is \file{gl2\_preamble.tex}: every macro
named below is defined there and nowhere else, and companion documents
\code{\textbackslash input} that file rather than copying it. The pre-consolidation
corpus copied macro blocks between \file{hecke\_action.tex},
\file{hecke\_extensions.tex} and \file{monodromy\_v5.tex}, which is precisely how
three incompatible meanings of $\De$, $\tau$ and $\inv$ came to coexist; \S\ref{sec:notation-collisions}
lists the nineteen resulting collisions and their resolution.

\subsection{Reading rules}\label{sec:notation-rules}

\begin{enumerate}[label=(N\arabic*)]
\item\label{N1} \textbf{One symbol, one meaning, document-wide.} Not
  ``one meaning per chapter'': a reader who opens the document at a proof must be
  able to read every symbol in it without knowing which chapter it is in.
\item\label{N2} \textbf{Decoration does not create a new symbol.} $\pi_1$ and a map
  called $\pi$ are a collision, not a disambiguation; so are $\tau$ and $\tau'$ with
  different meanings. Where two objects are genuinely two instances of one notion
  (the transporters $s_1,\dots,s_m$; the vertices $v_0,\dots,v_6$) indices are of
  course used.
\item\label{N3} \textbf{Single letters are a scarce resource; operators get names.}
  Where a map has no traditional single-letter name, it is written as a named
  operator in upright type ($\coneinc$, $\coneproj$, $\prR$, $\adjg$) rather than
  given a Greek letter. This costs three characters and buys a permanently free letter.
\item\label{N4} \textbf{Indices are letters, not primes.} $p$, $\ell$ and $q$ are
  reserved (\S\ref{sec:notation-primes}); a running index is $t$, $j$ or $s$. $i$ is
  never an index: it is reserved for $\sqrt{-1}$ and the elliptic point $v_6=i$
  (\S\ref{sec:notation-framework}, C15).
\item\label{N5} \textbf{Retired symbols are errors, not deprecations.} The macros
  \code{\textbackslash K}, \code{\textbackslash R}, \code{\textbackslash C},
  \code{\textbackslash M}, \code{\textbackslash pR}, \code{\textbackslash pu},
  \code{\textbackslash image} are defined in \file{gl2\_preamble.tex} to raise a
  \LaTeX{} error. Text re-homed from the old documents therefore cannot carry an old
  meaning in unnoticed: it fails to build.
\end{enumerate}

\subsection{The three arithmetic parameters}\label{sec:notation-primes}

This is the single most error-prone point in the corpus, and the one the owner has
already decided (OQ1, \file{hecke\_extensions\_owner\_decisions.md} \S OQ1:
\emph{adopt the contract's} \code{ell} / \code{n\_sym} / \code{p\_hecke}, weight
$k=n+2$). It is restated here as the binding rule, extended to the cohomological degree.

\begin{contractbox}
\begin{center}
\begin{tabular}{@{}llll@{}}
\toprule
role & symbol & code identifier & range \\
\midrule
\textbf{coefficient prime} (characteristic of $\cring$)      & $\ell$ & \code{ell}      & any prime \\
\textbf{Hecke prime} (index of the operator $T_p$)           & $p$    & \code{p\_hecke} & any prime, $p=\ell$ allowed \\
\textbf{symmetric-power degree}                              & $n$    & \code{n\_sym}   & $n\ge 0$ \\
\textbf{cohomological degree}                                & $q$    & \code{q}        & $q\ge 0$ \\
modular weight (derived, never independent)                  & $k:=n+2$ & ---           & \\
prime-power exponent in $T_{p^{\nu}}$, $N=p^{\nu}$           & $\nu$  & \code{r} (code) & $\nu\ge1$ \\
coefficient-ring exponent in $\cring=\Z/\ell^{e}$            & $e$    & ---             & $e\ge1$ \\
residue degree of $\Flf=\Talg_{\midl}/\midl$                 & $f$    & \code{r} (code) & $f\ge1$ \\
\bottomrule
\end{tabular}
\end{center}
\end{contractbox}

\noindent
Consequences, each of which changes text somewhere:
\begin{itemize}
\item $q$ is \textbf{the cohomological degree and nothing else}. The spec
  \file{hecke\_extensions.tex} used $q$ for the Hecke prime while also writing $H^q$
  and $d_q=\dim H^q$; every $T_q$, $B_q$, $\chi_b(q)$, $\Ga_0(q)$ there becomes
  $T_p$, $B_p$, $\chi_b(p)$, $\Ga_0(p)$ here.
\item $p$ is \textbf{the Hecke prime and nothing else}. The owner's notes
  \file{hecke\_extensions\_mathematics.md} and
  \file{hecke\_extensions\_computational\_strategy.md} use $p$ for the
  \emph{coefficient} prime --- the exact opposite assignment. A formula copied from
  them without applying the dictionary is not merely mislabelled, it is wrong.
\item $n$ is \textbf{the symmetric-power degree and nothing else}. Two other uses
  are retired: the cone degree (now $q$, \S\ref{sec:notation-cone}) and the composite
  Hecke index $T_n$ (now $T_N$ with $N\in\Z_{>0}$, which
  \file{hecke\_action.tex} already used at one site and not at two others).
\item $k$ is \textbf{the modular weight and nothing else}, and it is a
  \emph{derived} quantity: the document writes $k$ only where the classical
  weight-$k$ statement is being invoked, and always with $k=n+2$ visible nearby. The
  three other uses are retired: $k:=\#S$ (the number of Hecke operators; now $s$),
  the basis index $u_k=x^{n-k}y^k$ (now $u_j$), and the chain-term index
  $n^{(i,\sigma)}_k$ (now $n^{(j,\sigma)}_t$, C15).
\item The exponent of a coefficient ring $\Z/\ell^{e}$ is $e$, \textbf{not} $k$.
  This departs from the letter of the owner's OQ18 wording (``$\Z/\ell^k$''), which
  was written before $k$ was pinned to the weight; the two symbols would otherwise
  collide in the Sturm chapter, where $\Z/\ell^{e}$-coefficients and weight-$k$ forms
  appear in one sentence. \textbf{Owner confirmation requested} (\S\ref{sec:notation-open}, O1).
\end{itemize}

\subsection{The framework: complexes, the model, the pair}\label{sec:notation-framework}

\begin{center}
\begin{longtable}{@{}lll@{}}
\toprule
symbol & macro & meaning \\
\midrule
\endhead
$K$ & --- & the finite simplicial complex of the model (the 7-vertex complex) \\
$\Rel$ & \code{\textbackslash Rel} & the \emph{relation set} on $K$; the model is the pair $(K,\Rel)$ \\
$\KR$ & \code{\textbackslash KR} & the quotient complex \\
$\Ku$ & \code{\textbackslash Ku} & the universal cover complex \\
$\prR$, $\pru$ & \code{\textbackslash prR}, \code{\textbackslash pru} & the projections $K\to\KR$, $\Ku\to\KR$ \\
$\piOne$ & \code{\textbackslash piOne} & the fundamental group (only with the subscript $1$) \\
$\sigma$ & --- & an \emph{oriented simplex}; $\orient{q}(K)$ is the set of oriented $q$-simplices \\
$v_0,\dots,v_6$ & --- & the seven vertices of the model \\
$\Ga$ & \code{\textbackslash Ga} & $\SL_2(\Z)$ \\
$\Ga''$ & --- & $\Ga_0(N)$, $N=p^{\nu}$, the group of the correspondence \\
$\Gbdry$ & \code{\textbackslash Gbdry} & the boundary group $\langle T\rangle\times\langle -I\rangle$ \\
$\De$ & \code{\textbackslash De} & the monoid $\{g\in\Mat_2(\Z):\det g>0\}$ \\
$\Deprim_{p^{\nu}}$ & \code{\textbackslash Deprim} & the \emph{primitive} double coset \\
$\Xbulk$ & \code{\textbackslash Xbulk} & the $\Ga$-tessellation of the truncated upper half plane \\
$\Xbdry$ & \code{\textbackslash Xbdry} & $\Ga\cdot i(\Kbdry)\subset\Xbulk$, the $\Ga$-subcomplex of horocycle lines \\
$\Kbdry$ & \code{\textbackslash Kbdry} & the circle model: the path $v_0v_1v_2v_3$ with $v_3=Tv_0$ \\
$j$ & --- & the inclusion $\Xbdry\hookrightarrow\Xbulk$ (of \emph{spaces/chains}) \\
$\rres$ & \code{\textbackslash rres} & the restriction cochain map $\rres:=j^{\ast}$ \\
$\zeta$ & --- & $e^{\pi i/3}$; the elliptic vertices are $v_4=\zeta^2$, $v_5=\zeta$, $v_6=i$ \\
$T,S$ & --- & the standard generators $\smat1101$, $\smat0{-1}10$ of $\Ga$ \\
$\inv$, $\invK$ & \code{\textbackslash inv}, \code{\textbackslash invK} & the involution $\mathrm{diag}(-1,1)$ and its simplicial realisation \\
$\lift{\sigma}$ & \code{\textbackslash lift} & the distinguished lift of $\sigma\in\orient{q}(K)$ into $\orient{q}(\Ku)$, for the strict fundamental domain $K\subseteq\Ku$; the code spells this \code{i(sigma)}, which this document cannot (C16); the tilde's other readings are C19 \\
\bottomrule
\end{longtable}
\end{center}

\noindent\textbf{Note on $\Rel$ and $\cring$.} The old contract wrote the model as
$(K,R)$ and the cohomology as $H^q(K,R;L)$, with $R$ the \emph{relation}; the cone
and coefficient-ring material writes $\cring$ for the \emph{coefficient ring}
($6\in\cring^{\times}$, $\VR$, $\Z/\ell^{e}$). Both are indispensable and they cannot
share a letter. Here $\cring$ is the coefficient ring and the relation set is
$\Rel$, so the model is $(K,\Rel)$ and its cohomology $H^q(K,\Rel;L)$. Individual
relations are written as triples $(\sigma_{\mathrm{lhs}},\sigma_{\mathrm{rhs}},\gamma)
\in\Rel$, $\gamma\in\Ga$ (C17), and are given no single-letter name (the old
$r\in R_+$ is retired: $r$ is the restriction map). Code identifiers are unaffected:
\code{SimplexRelation} and
\code{EquivariantCohomologyOnK} name the same objects.

\subsection{Cochains, coefficients, bases}\label{sec:notation-cochains}

\begin{center}
\begin{longtable}{@{}lll@{}}
\toprule
symbol & macro & meaning \\
\midrule
\endhead
$\cring$ & \code{\textbackslash cring} & the coefficient ring: $\Fl$, $\Z/\ell^{e}$, or $\Z/m$ \\
$\Vcoef$ & \code{\textbackslash Vcoef} & $\Sym^n(\Z^2)$, the integral coefficient module \\
$\VR$ & \code{\textbackslash VR} & $\Vcoef\otimes\cring$ \\
$\Mn$, $\Dn$ & \code{\textbackslash Mn}, \code{\textbackslash Dn} & $\Sym^n\Fl^2$ and its contragredient (divided powers) \\
$\rhon$, $\rhodual$ & \code{\textbackslash rhon}, \code{\textbackslash rhodual} & the symmetric-power action and its contragredient \\
$u_j=x^{n-j}y^j$ & --- & the ordered basis of $\Sym^n$, $0\le j\le n$ \\
$d_V$ & --- & $\dim\Vcoef=n+1$ \\
$M_\gamma$ & --- & the monodromy matrix $\rhon(\gamma)$ of a relation \\
$\Cq{q}$, $\Zq{q}$, $\Bq{q}$, $\Heq{q}$ & \code{\textbackslash Cq}\ldots & equivariant cochains, cocycles, coboundaries, cohomology \\
$c^{(q)}_j$, $z^{(q)}_j$, $b^{(q)}_j$ & --- & the stored bases of $\Cq{q}$, $\Zq{q}$, $\Bq{q}$ \\
$d$ & --- & the \emph{coboundary} $\Cq{q}\to\Cq{q+1}$ \\
$\partial$ & --- & the \emph{chain boundary} $C_q\to C_{q-1}$ \\
$\aug$ & \code{\textbackslash aug} & the augmentation $C_0\to\Z$ \\
$P_q$ & --- & the projector $\Zq{q}\to\Heq{q}$ \\
$\cf$ & \code{\textbackslash cf} & the normalised weight-$12$ cusp form (classically ``$\Delta$'') \\
\bottomrule
\end{longtable}
\end{center}

\noindent The basis-name scheme $c/z/b$ replaces \file{monodromy\_v5.tex}'s
$\beta_j^{(q)}$, $\delta_s^{(q)}$, $\eta_p^{(q)}$, each of which collided with a
different reserved symbol --- and the last of which was additionally indexed by $p$.

\subsection{The Hecke datum}\label{sec:notation-hecke}

\begin{center}
\begin{longtable}{@{}lll@{}}
\toprule
symbol & macro & meaning \\
\midrule
\endhead
$\tau$ & --- & the $g$-twisted chain datum on $\Xbulk$ (\emph{only} this) \\
$\taub$ & \code{\textbackslash taub} & the $g$-twisted chain datum on $\Xbdry$ \\
$s_1,\dots,s_m$ & --- & the transporters, $\Ga=\bigsqcup_j\Ga''s_j$ \\
$\delta_j$ & --- & the coset elements $\delta_j=s_j^{-1}g$ \\
$\beta_j$, $\beta_\infty$ & --- & the upper-triangular \emph{right}-coset representatives \\
$\delta_j$, $\delta_\infty$ & --- & the upper-triangular \emph{left}-coset representatives (the normative ones) \\
$\Xig$, $\Th_g$ & \code{\textbackslash Xig}, \code{\textbackslash Th} & the induced-system map and the fused cochain map $\Th_g=\omega_\bullet\circ\Xig$ \\
$\alpha$, $\omega$ & --- & the restriction unit and the corestriction counit \\
$\mu$ & --- & the would-be middle map of the induced sequence (\emph{provably does not exist}) \\
$\Thb_g$ & \code{\textbackslash Thb} & the boundary cochain map \\
$\Amat{q}$, $\Abdry{q}$, $\Acone{q}$ & \code{\textbackslash Amat}\ldots & the cohomology-level matrices on $H^q$, $H^q(\partial)$, $H^q_c$ \\
$T_p$, $T_{p^{\nu}}$, $T_N$, $\diam$, $U_\ell$ & --- & the operators; $N\in\Z_{>0}$ a composite index \\
$\Talg$, $\midl$ & \code{\textbackslash Talg}, \code{\textbackslash midl} & the Hecke algebra and a maximal ideal of it \\
$\chi_b$ & --- & a boundary character \\
$a_N(\cf)$ & --- & the $N$-th Fourier coefficient of $\cf$ (classically ``$\tau(N)$'') \\
\bottomrule
\end{longtable}
\end{center}

\subsection{The pair, the cone, compact supports}\label{sec:notation-cone}

\begin{center}
\begin{longtable}{@{}lll@{}}
\toprule
symbol & macro & meaning \\
\midrule
\endhead
$\etapair$ & \code{\textbackslash etapair} & the \emph{pair homotopy} $C_\ast(\Xbdry)\to C_{\ast+1}(\Xbulk)$ \\
$\hder$ & \code{\textbackslash hder} & the \emph{derived} cochain homotopy $\hder:=\widetilde{\etapair}$ (never freely chosen) \\
$\Lam$ & \code{\textbackslash Lam} & the unique $C_0(\Xbdry)\to C_2(\Xbulk)$ with $\etapair'-\etapair=\partial\Lam-\Lam\partial$ \\
$\Cone{\rres}$ & \code{\textbackslash Cone} & the mapping cone, $\Cone{\rres}^q=\Cq{q+1}(K)\oplus\Cq{q}(\Kbdry)$ \\
$(c,c_\partial)$ & --- & a cone cochain (bulk component, boundary component) \\
$\Dcone$ & \code{\textbackslash Dcone} & the cone differential $\Dcone(c,c_\partial)=(-dc,\;dc_\partial+\rres c)$ \\
$\Phig$ & \code{\textbackslash Phig} & the cone cochain map $\Phig(c,c_\partial)=(\Th_g c,\;\Thb_g c_\partial+\hder c)$ \\
$\coneinc$, $\coneproj$ & \code{\textbackslash coneinc}, \code{\textbackslash coneproj} & the cone inclusion and projection \\
$\Hcpt{q}$ & \code{\textbackslash Hcpt} & $\Hcpt{q}:=H^{q-1}(\Cone{\rres})$ \\
$X_1,\dots,X_6$ & --- & the six groups of the long exact sequence \\
\bottomrule
\end{longtable}
\end{center}

\noindent Three points where this differs from the sources it is re-homed from:
\begin{itemize}
\item \textbf{The cone degree is $q$.} The OQ2 verdict and
  \file{monodromy\_v5.tex} \S8 write $\Cone{\rres}^n=A^{n+1}\oplus B^n$ and
  $H^n_c:=H^{n-1}(\Cone{\rres})$ --- with $n$ also the symmetric power, in the same
  notation block. Here: $\Cone{\rres}^q$, $\Hcpt{q}=H^{q-1}(\Cone{\rres})$.
\item \textbf{The cone inclusion is $\coneinc$, not $\inv$.} Both the verdict
  ($\inv_\ast\circ k_\ast\circ\pi_\ast$) and the implementation plan
  ($\Phig\,\inv=\inv\,\Thb_g$) wrote $\inv$ for the cone inclusion, $130$ lines after
  the same plan fixed $\inv$ as the involution --- both inside the text proposed for
  the contract. Here $\inv$ is the involution; the cone maps are named operators.
\item \textbf{The cone element is $(c,c_\partial)$, not $(a,b)$.} $a$ is a Fourier
  coefficient and an exponent in $\Z/\ell^a$; $b$ indexes the boundary characters
  $\chi_b$. The component names also say which component is which, which the
  formalization will want.
\end{itemize}

\subsection{The nineteen collisions and their resolution}\label{sec:notation-collisions}

Eleven collisions were found by the corpus survey of 2026-09-15; three more
(C12, C13, C14) were found during this design pass, a fifteenth (C15) during
the Chapters~7--8 drafting pass, a sixteenth (C16) during the Chapter~11
drafting pass, and a seventeenth (C17) during the Chapter~13 slot~4 follow-up
pass, all 2026-09-15. Two further ones were found while Chapters~14 and~15 were
written (2026-09-16) and resolved \emph{locally} there, each with a note saying it was
not yet listed here: they are C18 and C19, added 2026-09-16. Each
row says which document's usage \emph{changes}.

\begin{center}
\begin{longtable}{@{}p{0.055\textwidth}p{0.40\textwidth}p{0.47\textwidth}@{}}
\toprule
\# & the two (or more) live meanings & resolution \\
\midrule
\endhead
\textbf{C3} &
$\inv=\mathrm{diag}(-1,1)$ (spec (A3); owner decision OQ3; record field
\code{iota\_sign}) \emph{vs} $\inv=$ the cone inclusion $B\to\Cone{\rres}$
(\file{monodromy\_v5.tex}:1226; verdict \S0, \S2; plan \S0.4 C-A3) &
$\inv$ \textbf{is the involution}. The cone inclusion becomes $\coneinc$ and the
projection $\coneproj$. \emph{Changes:} \file{monodromy\_v5.tex} \S8, the verdict's
and the plan's cone text. \\
\addlinespace
\textbf{C4} &
$\tau=$ the $g$-twisted chain datum (contract, $36$ code sites via \code{alg:tau})
\emph{vs} $\tau=$ an oriented simplex (\file{hecke\_action.tex}:1012--1015,
\emph{inside the proof of} \code{thm:xi}; \file{monodromy\_v5.tex}:619,628--630)
\emph{vs} $\tau(N)=$ Ramanujan's function (spec, six sites) &
$\tau$ \textbf{is the chain datum}. Oriented simplices are $\sigma$ (already the
contract's own convention elsewhere). Ramanujan's function is written
$a_N(\cf)$ and never appears as $\tau$. \emph{Changes:} the proof of
\code{thm:xi}, \file{monodromy\_v5.tex} Thm~44's proof, the spec's $\ell=691$ section. \\
\addlinespace
\textbf{C5} &
$\De=\{g\in\Mat_2(\Z):\det g>0\}$ (both documents; ``note $\inv\notin\De$'')
\emph{vs} $\Delta=$ Ramanujan's cusp form (spec:162, 424--425, 1037, 1327, 1334) ---
\emph{same macro, same document} &
$\De$ \textbf{is the monoid}. The weight-$12$ cusp form is $\cf$, with
$a_N(\cf)$ its coefficients. \emph{Changes:} the spec. \\
\addlinespace
\textbf{C6} &
$\etapair=$ the pair homotopy (verdict \S2.3ff; spec:779; plan:346) \emph{vs}
$\eta_p^{(q)}=$ a coboundary basis vector (\file{monodromy\_v5.tex}:1182, 1208,
1211, 1325, 1328 --- \emph{inside \S8, the very cone section}), itself indexed by $p$ &
$\etapair$ \textbf{is the pair homotopy}. Basis vectors are $c^{(q)}_j$ / $z^{(q)}_j$
/ $b^{(q)}_j$ for $\Cq{q}$ / $\Zq{q}$ / $\Bq{q}$, indexed by $j$.
\emph{Changes:} \file{monodromy\_v5.tex} \S7--\S8. \\
\addlinespace
\textbf{C7} &
$\delta_j=s_j^{-1}g$ (contract, \code{alg:assembly}) \emph{vs} $\delta$ the
coboundary $(\delta\varphi)(x)=\varphi(\partial x)$ (\file{hecke\_action.tex}:690,
1014 --- same proof as the coset elements) \emph{vs} $\delta_s^{(q)}$ a cocycle basis
vector (\file{monodromy\_v5.tex}:1098--1109) &
$\delta$ \textbf{is the coset element}. The coboundary is $d$ throughout
($\partial$ for chains, $d$ for cochains, $\Dcone$ for the cone). Basis vectors as in
C6. \emph{Changes:} \file{hecke\_action.tex} \S\S4--5, \file{monodromy\_v5.tex} \S7. \\
\addlinespace
\textbf{C8} &
$\beta_j=\smat1j0p$ right-coset representatives (contract \code{prop:cosets})
\emph{vs} $\beta$ the middle map of the induced sequence
(\file{monodromy\_v5.tex}:1519, the map proved not to exist) \emph{vs}
$\beta_j^{(q)}$ a basis vector (\file{monodromy\_v5.tex}:1098) &
$\beta_j$ \textbf{is the right-coset representative}. The non-existent middle map is
$\mu$, with a parenthetical ``called $\beta$ in [v5] \S12''. Basis vectors as in C6.
\emph{Changes:} \file{monodromy\_v5.tex} \S7, \S12; the no-go remark keeps its label
\code{rmk:nogo}. \\
\addlinespace
\textbf{C9} &
$\rho=e^{2\pi i/3}$, $\rho^2$, ``the elliptic points $i,\rho,\rho^2$''
(\file{hecke\_action.tex}:194, 222--223, 243--244) \emph{vs} $\rhon$, $\rho_L$ the
representation (\emph{same subsection}, :245) &
$\rho$ \textbf{is the representation}. The elliptic points are $\zeta:=e^{\pi i/3}$,
$\zeta^2$ and $i$; the vertices are $v_4=\zeta^2$, $v_5=\zeta$, $v_6=i$. (The old
naming was non-standard as well as colliding: $e^{\pi i/3}$ is conventionally
$-\bar\rho$.) \emph{Changes:} \file{hecke\_action.tex} \S2.2. \\
\addlinespace
\textbf{C2} &
$n=$ symmetric power (everywhere) \emph{vs} $n=$ the cone degree (verdict:32; plan
C-A3(ii.4); \file{monodromy\_v5.tex}:1152) \emph{vs} $T_n$ the composite Hecke index
(\file{hecke\_action.tex}:80, 1439, against $T_N$ at :518) &
$n$ \textbf{is the symmetric power}. Cone degree $\to q$; composite index $\to T_N$.
\emph{Changes:} the verdict's and the plan's cone text, \file{monodromy\_v5.tex} \S8,
two sites in \file{hecke\_action.tex}. \\
\addlinespace
\textbf{C10} &
$k=n+2$ the weight (spec:136, 206, 323, 419) \emph{vs} $k:=\#S$ (spec:746,
\emph{same document}) \emph{vs} $u_k=x^{n-k}y^k$ (contract:294) \emph{vs}
$n_k^{(i,\sigma)}$ (contract:977) \emph{vs} $k=$ the symmetric power in [M]/[S] &
$k$ \textbf{is the weight}, always derived and always with $k=n+2$ visible.
$\#S\to s$; $u_k\to u_j$; $n_k^{(i,\sigma)}\to n_t^{(j,\sigma)}$ (the transporter index by C15); [M]/[S]'s $k\to n$.
\emph{Changes:} spec \S7, contract \S2.3 and \S5, both owner notes. \\
\addlinespace
\textbf{C11} &
$r$ the exponent in $T_{p^r}$ (contract:496) \emph{vs} $r$ the residue degree of
$\F_{\ell^r}$ (spec:40, 345) \emph{vs} $r^0,r^1$ the restriction maps (spec:238;
verdict; plan) \emph{vs} $r\in R_+$ a relation (contract:107, 121) ---
two of them as superscripts on one line at spec:280 &
$r$ \textbf{is the restriction cochain map} $\rres=j^{\ast}$ (it is the symbol the
cone notation is built on: $\Cone{\rres}$, $\ker\rres^1$, $\coker\rres^0$). Exponent
$\to\nu$ ($T_{p^{\nu}}$, $N=p^{\nu}$); residue degree $\to f$ ($\Flf$) --- the
standard number-theoretic letter; relations get no letter (see
\S\ref{sec:notation-framework}). \emph{Changes:} contract \S3, \S2.1; spec \S6. \\
\addlinespace
\textbf{C1} &
$q$ the Hecke prime (spec:134, and $\sim60$ occurrences of $T_q$) \emph{vs} $q$ the
cohomological degree (spec:262, 746, \emph{same Throughout-line}; the whole contract;
the code) &
$q$ \textbf{is the cohomological degree}; the Hecke prime is $p$
(\S\ref{sec:notation-primes}, owner decision OQ1). \emph{Changes:} the spec throughout. \\
\addlinespace
\label{coll:p}D3 &
$p$ the Hecke prime (contract, code) \emph{vs} $p$ the coefficient prime
([M], [S], and the owner's own wording in the decision record) &
$p$ \textbf{is the Hecke prime}; the coefficient prime is $\ell$
(owner decision OQ1). \emph{Changes:} [M], [S]. \\
\addlinespace
\textbf{C12} &
$M_\gamma:=\rhon(\gamma)$ a monodromy matrix (contract:245; the code's
\code{SimplexRelation::M}) \emph{vs} $\Mn=\Sym^n\Fl^2$ and $\Dn$ (spec) \emph{vs}
$M:C_0(\Xbdry)\to C_2(\Xbulk)$ the unique map of the $\etapair$-uniqueness lemma
(verdict Lemma 2.4.1; plan C-A3(ii.3)) \emph{vs} $\Mat_2(\Z)$ (macro
\code{\textbackslash M}) \emph{vs} $M^T$, $M_T$ boundary (co)invariants &
The coefficient module is $\Vcoef$ / $\VR$, with $\Mn$ retained as a synonym
\emph{defined once} for the classification literature; monodromy matrices keep
$M_\gamma$ (always subscripted); the unique degree-$(+2)$ map is $\Lam$; the matrix
algebra is $\Mat_2(\Z)$ (\code{\textbackslash Mat}; \code{\textbackslash M} is
retired to an error); boundary (co)invariants are $\VR^{T}$, $(\VR)_{T}$.
\emph{Changes:} the verdict's and plan's cone text, the spec. \\
\addlinespace
\textbf{C13} &
$R$ the relation set, in $(K,R)$ and $H^q(K,R;L)$ --- \emph{including the title of
the contract} --- \emph{vs} $R$ the coefficient ring ($6\in R^{\times}$, $V_R$,
$R=\Z/\ell^e$) in all of the cone and integrality material &
$\cring$ \textbf{is the coefficient ring}; the relation set is $\Rel$, so the model
is $(K,\Rel)$ and its cohomology $H^q(K,\Rel;L)$. \emph{Changes:} the contract's
title and \S\S1--2, \file{monodromy\_v5.tex} throughout (typographic only; no label
and no code identifier moves). \\
\addlinespace
\textbf{C14} &
$s_i$ the transporters $\Ga=\bigsqcup_i\Ga''s_i$ (contract) \emph{vs} $s$ an
(un)oriented simplex (\file{monodromy\_v5.tex}; the code's flat index $s\cdot d_V+k$)
\emph{vs} $S=\smat0{-1}10$ the generator \emph{vs} $S$ the set of computed operators
(spec \S7) &
$s_j$ \textbf{is a transporter} (index $j$, not $i$: C15); simplices are $\sigma$; $S$ is the generator; the
operator set is $\mathcal{S}$ with $s:=\#\mathcal{S}$. \emph{Changes:}
\file{monodromy\_v5.tex}, spec \S7. \\
\addlinespace
\textbf{C15} &
\ref{N4}'s reading rule, ``a running index is $i$, $j$, $t$ or $s$''
\emph{vs} $v_6=i=\sqrt{-1}$, the elliptic point
(\S\ref{sec:notation-framework}) --- the same chapter, the reading rules and the
framework table a page apart &
$i$ \textbf{is $\sqrt{-1}$ (and the elliptic point $v_6$) only}. Running indices
are $t$, $j$ or $s$; \ref{N4} is amended accordingly. The transporters and coset
elements are therefore $s_j$, $\delta_j=s_j^{-1}g$ and the chain-term coefficients
$n^{(j,\sigma)}_t$, which is what Chapters~10--11 write. \emph{Changes:} this
chapter's own \ref{N4} and its tables (\S\ref{sec:notation-hecke},
C7, C10, C14); Chapters~7--8 already avoided $i$ as an
index (found during their drafting); \file{hecke\_action.tex} throughout, where the
transporter index is $i$ (crosswalk, \S\ref{sec:notation-crosswalk}). Owner decision
recorded \S\ref{sec:notation-open}, 2026-09-15. \\
\addlinespace
\textbf{C16} &
$i(\sigma)$ the distinguished lift of a simplex $\sigma\in K$ into the universal cover
$\Ku$ (the code's \code{i(sigma)}; \file{monodromy\_v5.tex}'s inclusion
$i:K\to\Ku$ and \code{rmk:injectivityofi}; \file{hecke\_action.tex}, the $\tau$-datum
and \code{alg:tau}) \emph{vs} $i=\sqrt{-1}$ and the elliptic point $v_6$ (C15) ---
and, in the same formulas, $i$ as the transporter index &
$i$ \textbf{is $\sqrt{-1}$} (C15), so the lift gets no letter at all: it is the
decoration $\lift{\sigma}$ (\code{\textbackslash lift}), rendered
$\widetilde{\sigma}$, for the strict fundamental domain $K\subseteq\Ku$ of
\ref{H4}(iii). It is a \emph{map}, not an index, so \ref{N3} does not force an
upright operator name. \emph{Changes:} \file{monodromy\_v5.tex} \S\S1--2 and
\file{hecke\_action.tex} \S\S5--6 wherever $i(\sigma)$ appears; the code identifier
\code{i(sigma)} is unaffected and does \emph{not} move. \\
\addlinespace
\textbf{C17} &
$g$ the relation element of a triple $(\sigma_{\mathrm{lhs}},\sigma_{\mathrm{rhs}},g)
\in\Rel$ (\code{def:model7rel}, the seven-vertex model's relation table) \emph{vs}
$g=\smat{1}{0}{0}{N}$ the Hecke element fixed throughout the Hecke chapters
(\S\ref{sec:notation-hecke}) --- the two meet in one formula, \code{thm:xi}(1) &
$g$ \textbf{is the Hecke element only}. Relation elements are $\gamma\in\Ga$, so a
relation is a triple $(\sigma_{\mathrm{lhs}},\sigma_{\mathrm{rhs}},\gamma)\in\Rel$;
this model's relation elements ($T$, $S$, $ST$, $TS$, $-I$) all lie in $\Ga$, so the
rename costs nothing. \emph{Changes:} \code{def:model7rel} and its contract box
(both instances of $g$), and the one further site where the relation triple's third
component was written $g$ (the contragredient local system, \code{sec:contragredient}
in the contract's numbering). \\
\addlinespace
\textbf{C18} &
$N=p^{\nu}$ the Hecke level, fixed throughout the Hecke chapters
(\S\ref{sec:notation-primes}; $\Ga''=\Ga_0(N)$, $g=\mathrm{diag}(1,N)$, $T_N$,
$a_N(\cf)$) \emph{vs} $N:=\prod_{c\in\F_\ell}(y+cx)=y^{\ell}-x^{\ell-1}y$, the
degree-$\ell$ boundary invariant of \code{cl:boundary} and \code{cl:chi}
(\file{hecke\_extensions.tex}; the spelling the boundary chapter inherited) --- the two
meet in one statement, \code{prop:bdrychar}, whose eigenvalues are indexed by the
\emph{level} $p$ and whose basis is built from the \emph{invariant} &
$N$ \textbf{is the Hecke level only}. The invariant is $\mathrm{Nm}$ --- it is the norm,
i.e.\ the $\langle T\rangle$-orbit product, of $y$ --- and it gets an upright operator
name, not a letter (\ref{N3}). \emph{Changes:} \file{hecke\_extensions.tex}
\code{cl:boundary}, \code{cl:chi} and every occurrence of the invariant in their
neighbourhood; in this document the spelling is already $\mathrm{Nm}$ everywhere
(\code{prop:bdrychar} and the $\VR^{T}$ basis $x^{\,n-b\ell}\mathrm{Nm}^{\,b}$). Found
while Chapter~14 was written; resolved locally there on 2026-09-16 and listed here on
the same day. \\
\addlinespace
\textbf{C19} &
The \textbf{tilde} as a decoration, with four live readings: $\lift{\sigma}$, the
distinguished lift of a simplex (C16, \code{\textbackslash lift}); $\widetilde F$, the
coset transport of a (T1)-map (\code{def:cosettransport} --- $\widetilde\tau$,
$\widetilde H$, $\widetilde\Lam$, $\widetilde{N_\partial}$, $\hder=\widetilde{\etapair}$);
$\widetilde T$, the \emph{cochain-level} operator of \code{def:hecke}; and, in Part~I
only, $\widetilde K$ the saturation and $\widetilde X$ the universal cover &
\textbf{The decoration is read from what it decorates}, and no rename is made: on a
\emph{simplex or chain} it is the lift, on a \emph{map with (T1)} it is coset transport.
The two never meet in one formula and are typographically distinct (a decorated symbol
versus a decorated operator); the point of definition of each says so
(\S\ref{sec:notation-framework}, \code{def:cosettransport}). $\widetilde T$ is
\emph{consistent} with the second reading only because $\widetilde T=\widetilde\tau$
(\code{def:cosettransport}), $T$ itself being no chain map; it is kept because $53$ code
sites and both superseded documents spell it so. Part~I's two are out of scope of the
other three --- no lift and no coset transport occurs there --- and are not renamed.
\emph{Rule:} no \textbf{new} tilde decoration may be introduced, in this document or in
the code, without a row here. \emph{Changes:} none to any document; this row records the
reading rule. Found while Chapter~15 was written; noted there on 2026-09-16 and listed
here on the same day. \\
\bottomrule
\end{longtable}
\end{center}

\subsection{Crosswalk: reading the old documents}\label{sec:notation-crosswalk}

A reader holding a pre-consolidation document applies this table. It is also the
checklist for the re-homing work: a chapter is not done until every row that touches
its source has been applied.

\begin{center}
\begin{longtable}{@{}llll@{}}
\toprule
here & \file{hecke\_action.tex} & \file{hecke\_extensions.tex} & \file{monodromy\_v5.tex} / [M], [S], [V] \\
\midrule
\endhead
$\ell$ & $\ell$ & $\ell$ & $p$ \quad([M], [S]) \\
$p$ & $p$ & $q$ & $q$ \quad([M], [S]); $p$ ([V]) \\
$n$ & $n$ & $n$ & $k$ \quad([M], [S]) \\
$q$ (degree) & $q$ & $q$ \emph{and} the prime & $k$, $n$ \\
$k=n+2$ & --- & $k$, \emph{also} $\#S$ & $k+2$ \quad([M]) \\
$\nu$ & $r$ (in $T_{p^r}$) & $r$ & --- \\
$f$ & --- & $r$ (in $\F_{\ell^r}$) & --- \\
$e$ & --- & --- & $k$ (in $\Z/\ell^k$) \\
$s$ & --- & $k:=\#S$ & --- \\
$t,j,s$ (running index) & $i$ & $i$ & $i$ \\
$s_j$, $\delta_j$, $n^{(j,\sigma)}_t$ & $s_i$, $\delta_i$, $n^{(i,\sigma)}_k$ & --- & --- \\
$\lift{\sigma}$ & $i(\sigma)$ & --- & $i(\sigma)$, $i:K\to\Ku$ \\
$\rres$ & --- & $r^0,r^1$ & $r$ \\
$\Rel$ & $R$ & --- & $R$ \\
$\gamma$ (relation element) & $g_r$ & --- & $g_r$ \\
$\cring$ & --- & $R$ & $R$ \quad([V], [X]) \\
$\sigma$ & $\sigma$, \emph{also} $\tau$ & --- & $\tau$, $s$ \\
$\zeta$ & $\rho$ & --- & --- \\
$u_j$ & $u_k$ & --- & --- \\
$d$ (coboundary) & $\delta$ & --- & $\delta_{\mathrm{simp}}$ \\
$c^{(q)}_j,z^{(q)}_j,b^{(q)}_j$ & --- & --- & $\beta_j^{(q)},\delta_s^{(q)},\eta_p^{(q)}$ \\
$\coneinc,\coneproj$ & --- & --- & $\inv$, $\pi$ \\
$(c,c_\partial)$ & --- & --- & $(a,b)$ \\
$\Lam$ & --- & --- & $M$ \quad([V]) \\
$\mu$ & $\beta$ (in \code{rmk:nogo}) & --- & $\beta$ \\
$\cf$, $a_N(\cf)$ & --- & $\Delta$, $\tau(q)$ & --- \\
$\Mat_2(\Z)$ & $\mathrm{M}_2(\Z)$ & $\Mat_2(\Z)$ & --- \\
$T_N$ & $T_n$ \emph{and} $T_N$ & --- & --- \\
$\mathrm{Nm}$ (the invariant) & --- & $N$ \quad(C18) & --- \\
the tilde & \multicolumn{3}{l}{read by what it decorates; no rename (C19)} \\
\bottomrule
\end{longtable}
\end{center}

\subsection{What the owner must confirm}\label{sec:notation-open}

OQ1 fixed the three arithmetic parameters and nothing else. The following go beyond
it and are decisions taken in this design pass; each is defensible and each is
reversible at zero cost \emph{now} and at the cost of a full re-read \emph{later}.
One item that would otherwise belong here does not: C15 ($i$ reserved for
$\sqrt{-1}$/the elliptic point, never a running index) was already decided by the
owner on 2026-09-15 and needs no confirmation; it is recorded, settled, at its row
in \S\ref{sec:notation-collisions} and in \ref{N4}.

\begin{enumerate}[label=(O\arabic*)]
\item\label{O1} $\cring=\Z/\ell^{e}$ rather than $\Z/\ell^{k}$
  (\S\ref{sec:notation-primes}). Departs from the letter of OQ18's wording; needed
  because $k$ is the weight.
\item\label{O2} $\cring$ for the coefficient ring and $\Rel$ for the relation set
  (C13) --- which changes the \emph{title} of the binding contract from
  $H^{\ast}(K,R;L)$ to $H^{\ast}(K,\Rel;L)$. The alternative (keep $R$ for the
  relation, rename the coefficient ring) was rejected: $6\in R^{\times}$, $V_R$ and
  $\Z/\ell^e$ appear in far more statements than $(K,R)$ does, and no other letter for
  a commutative ring reads naturally.
\item\label{O3} $\nu$ for the prime-power exponent and $f$ for the residue degree
  (C11), so that $r$ can be the restriction map. The alternative (keep
  $T_{p^r}$, rename the restriction to $\res^{q}$) was rejected because $\res$ is
  already the group-cohomology restriction and $\Cone{\rres}$, $\ker\rres^1$,
  $\coker\rres^0$ are the load-bearing notation of the newest chapter.
\item\label{O4} The elliptic points as $\zeta=e^{\pi i/3}$, $\zeta^2$
  (C9), rather than keeping $\rho$ and renaming the representation.
\item\label{O5} Ramanujan's $\Delta$ and $\tau(N)$ written $\cf$ and $a_N(\cf)$
  (C5, C4). This is the one resolution that departs from
  universal mathematical custom; it is taken because $\De$ and $\tau$ are both
  normative \emph{implementation} objects with code citations, and the cusp form
  appears only in the classification literature.
\item\label{O6} The basis-name scheme $c^{(q)}_j$ / $z^{(q)}_j$ / $b^{(q)}_j$
  (C6, C7, C8) replacing three colliding
  Greek letters at once.
\item\label{O7} Cone maps as named operators $\coneinc$, $\coneproj$ rather than
  letters (\ref{N3}, C3).
\end{enumerate}

\subsection{The integral line}\label{sec:notation-integral}

Added at v1.8.0 with \S\ref{sec:integral}; every symbol below is used only there and in
the statements that point at it.

\begin{center}
\begin{longtable}{@{}p{0.2\textwidth}p{0.17\textwidth}p{0.55\textwidth}@{}}
\toprule
symbol & macro & meaning \\
\midrule
\endhead
$\Ceq(\cring)$ & --- & $\Ceq(K,\Rel;L)$ with fibre $\VR$ and the relations imposed \emph{over $\cring$}: the complex the implementation builds for the ring in force. \textbf{Never} the base change $\Ceq(\Z)\otimes\cring$, which is written so; the two agree when $6\in\cring^{\times}$ (\ref{prop:ringinforce}) \\
$\Iso_{\bullet}$, $\Isob_{\bullet}$ & \code{\textbackslash Iso}, \code{\textbackslash Isob} & the isotropy-resolution model on the unrefined $(\Xbulk,\Xbdry)$ and its boundary subcomplex (\ref{def:isomodel}) \\
$F^{\sigma}_{\bullet}$ & --- & the periodic free resolution of $\Z$ over $\Z[\Ga_\sigma]$; the analogue, for the stabiliser, of the $F_{\bullet}$ of \ref{def:borelmodel} over $\Z[\Ga]$ (\ref{N2}: the same kind of object for a subgroup, not the same object; worded so at v1.8.1) \\
$\Sigma_s$, $\Sigma^{\partial}_{s}$ & --- & the orbit representatives of $s$-simplices of $\Xbulk$, $\Xbdry$, as in the proof of \ref{thm:specialisation} \\
$\cmpI$ & \code{\textbackslash cmpI} & the comparison map $\Iso_{\bullet}\to C_{\bullet}(\Xbulk)$, a named operator (\ref{N3}); $\cmpI^{\ast}$ its cochain map \\
$\CI(\cring)$ & \code{\textbackslash CI} & $\Hom_{\Z[\Ga]}(\Iso_{\bullet},\VR)$ \\
$\tau_{\Iso}$, $\tau^{\partial}_{\Iso}$ & --- & the $g$-twisted chain datum on $\Iso_{\bullet}$, $\Isob_{\bullet}$ (\ref{prop:strictlift}); an instance of \ref{def:tau}'s notion, not a new meaning of $\tau$ \\
$\Zsix$, $\Qsix$ & \code{\textbackslash Zsix}, \code{\textbackslash Qsix} & $\Z_2\times\Z_3$ and $\Q_2\times\Q_3$ \\
$\Z_\ell$, $\Z_{(\ell)}$ & --- & the $\ell$-adic integers; the localisation of $\Z$ at $\ell$ \\
$\Hfr{q}$ & \code{\textbackslash Hfr} & $H^{q}(\Ga;\Vcoef)$ modulo its torsion \\
$e$ & --- & the exponent of the ring $\Z/\ell^{e}$ in force (\S\ref{sec:notation-primes}); the design notes' ``$K$'' for a certified level is this $e$, since $K$ is the model complex \\
\bottomrule
\end{longtable}
\end{center}

\noindent\emph{Reading the design notes.} The integral-line brief
(\file{docs/implementation/integral\_line\_design\_brief.md}) and the $\ell$-adic
investigation write $R$ for the free resolution (here $\Iso_{\bullet}$: $\cring$ is the
coefficient ring), $\varepsilon$ for the comparison map (here $\cmpI$: $\aug$ is the
augmentation), $\widehat C$ for $\Hom_{\Z[\Ga]}(R,V)$ (here $\CI$), $\hat\tau$ for the lifted
datum (here $\tau_{\Iso}$), $R_\partial$ (here $\Isob_{\bullet}$), $\Lambda^{q}$ for the free
quotient (here $\Hfr{q}$: $\Lam$ is the $\etapair$-uniqueness map), $C_7\otimes\Z/\ell^{K}$
for the pipeline object (here $\Ceq(\Z/\ell^{e})$), $\Z/\ell^{k}$ and $K$ for rings and
certified levels (here $\Z/\ell^{e}$), and $e_i$ for torsion exponents (here $a_t$, local to
the statement that introduces it).

\subsection{The Dickson line and the ordinary projector}\label{sec:notation-dickson}

Added at v1.11.0-pre.1 with \S\ref{sec:dicksontheta} and \S\ref{sec:ordinary}. The plan and
the code write $L$, $D$ and $\theta$; this document cannot, because $L$ is the matrix local
system (\S\ref{sec:notation-framework}), $\Dn$ the contragredient module and $\theta_g$ the
translation of \S\ref{sec:correspondence} (rule~\ref{A5}).

\begin{center}
\begin{longtable}{@{}p{0.2\textwidth}p{0.17\textwidth}p{0.55\textwidth}@{}}
\toprule
symbol & macro & meaning \\
\midrule
\endhead
$\DkL$ & \code{\textbackslash DkL} & the first Dickson invariant $x^{\ell}y-xy^{\ell}\in\Sym^{\ell+1}(\Fl^{2})$ (\ref{def:dicksoninv}); the code's and the plan's $L$ \\
$\DkD$ & \code{\textbackslash DkD} & the second Dickson invariant, of degree $\ell^{2}-\ell$ (\ref{def:dicksoninv}); the plan's $D$ \\
$\thL$ & \code{\textbackslash thL} & multiplication by $\DkL$, on fibres, cochains and cohomology (\ref{def:theta}); the literature's and the plan's $\theta$ \\
$n^{+}$ & \code{\textbackslash nplus} & $n+\ell+1$, the target degree of $\thL$ (the plan's and the code's $m$; $N$ is the Hecke level only, $m$ the coset count) \\
$\Qn$, $\rho_{\Qn}$ & \code{\textbackslash Qn} & the quotient $\Sym^{n^{+}}/\DkL\cdot\Sym^{n}$ over $\Fl$ and its monoid action (\ref{def:Qn}) \\
$\bar u_t$ & --- & the class of $u_t=x^{n^{+}-t}y^{t}$ in $\Qn$; the complement basis is $(\bar u_0,\ldots,\bar u_{\ell-1},\bar u_{n^{+}})$ \\
$\mathbf{1}_n$ & --- & the class $\bar u_0-\bar u_{\ell-1}+\bar u_{n^{+}}$, the invariant of $\Qn$ when $(\ell-1)\mid(n+2)$ (\ref{lem:Qnfunctions}) \\
$h^{q}(m)$ & --- & $\dim_{\Fl}H^{q}(\Ga;M_m)$ \\
$\eord$ & \code{\textbackslash eord} & the ordinary projector of an endomorphism of a finite-length module (\ref{def:ordinary}); \emph{not} the class $e$ of \ref{lem:bulkperiodic} nor the exponent $e$ of $\Z/\ell^{e}$ \\
\midrule
\multicolumn{3}{@{}l}{\emph{Added at v1.12.0-pre.1 with \S\ref{sec:hida}.}}\\
$J_n$ & --- & $\rhon(S)-1$ stacked over $\rhon(T)-1$, a $2(n+1)\times(n+1)$ integer matrix (\ref{prop:torsionprepass}) \\
$\lambda_e$ & --- & $\lambda(H^{0}(\Ga;\Vcoef\otimes\Z/\ell^{e}))$, the length (\ref{prop:torsionprepass}) \\
$e_{\max}$ & --- & the largest exponent of the $\ell$-primary torsion of $H^{1}(\Ga;\Vcoef)$ ($0$ if none); the code's \code{e\_max} \\
$e'$ & --- & a target precision; [E]'s and the code's $k$ in $K=k+e_{\max}$, which this document cannot write ($k$ is the weight) \\
$\mathfrak g$ & --- & a number of generators (\ref{lem:ordexponent}); the plan's $r$, which is the restriction map here \\
$N_{\mathrm{ord}}(\mathfrak g,K)$ & --- & the ordinary exponent $\ell^{K-1+\lceil\log_\ell\mathfrak g\rceil}\prod_{f\le\mathfrak g}(\ell^{f}-1)$ (\ref{lem:ordexponent}); $N$ alone is the Hecke level \\
$h^{1}_{\mathrm{ord}}(n)$, $\chi^{\mathrm{ord}}_{n,p}$ & --- & the ordinary rank and the characteristic polynomial of $T_p$ on the ordinary part (\ref{def:hidatower}) \\
$\mathrm{Tw}_m(\ell;n_0)$ & --- & the weight tower of level $m$ through $n_0$ (\ref{def:hidatower}) \\
(Hi1)--(Hi6) & --- & the predictions of \ref{rmk:hidacheckable}; the plan's (H1)--(H6), which are the standing hypotheses here \\
\midrule
\multicolumn{3}{@{}l}{\emph{Added at v1.13.0-pre.1 with \S\ref{sec:hidafull}.}}\\
$m_N$ & --- & $[\Ga:\Ga_0(N)]=\abs{\mathbb{P}^1(\Z/N)}$ (\ref{def:levelellmodel}); subscripted, as $m$ alone is a coset count or a tower level \\
$[A\,g\,B]$ & --- & $\cor^{B}_{B\cap gAg^{-1}}\circ(\conj_g)_{\ast}\circ\res^{A}_{A\cap g^{-1}Bg}$ between two groups; $[\Ga g\Ga]=\Tdc{g}$ \\
$\widetilde{[A\,g\,B]}$ & --- & its block formula on the induced fibre \eqref{eq:levelcorr}, reusing the level-$1$ $\tau$-table (\ref{prop:levelcorrespondence}) \\
$U^{(\ell)}_\ell$ & --- & $[\Ga_0(\ell)\,g\,\Ga_0(\ell)]$, $g=\mathrm{diag}(1,\ell)$ (\ref{prop:levelellU}); the plan's $U_\ell^{(\ell)}$; $U_\ell=T_\ell$ stays the level-$1$ operator \\
$U^{\ast}$ & --- & $[\Ga_0(\ell)\,\mathrm{diag}(\ell,1)\,\Ga_0(\ell)]$, the Atkin--Lehner conjugate, Hida's $T(p)$ at level $p$ (\ref{rmk:stabconvention}) \\
$\mathrm{dg}_1$, $\mathrm{dg}_\ell$ & --- & the degeneracy maps $\res^{\Ga}_{\Ga_0(\ell)}$ and $[\Ga\,g\,\Ga_0(\ell)]$; the plan's $\alpha_1$, $\alpha_\ell$ ($\alpha$ is \ref{prop:alpha}'s restriction unit) \\
$\alpha^{0}_{B/A}$, $\omega_{A/B}$ & --- & the level-change fibre maps, projection and trace (\ref{prop:levelcorrespondence}(4)(b)); $\omega_{\Ga_0(N)/\Ga}=\omega$ of \ref{prop:omega} \\
$\varpi$ & --- & the block row $(\rhon(gs_1^{-1}),\dots,\rhon(gs_\ell^{-1}),0)$ of \eqref{eq:Ulevell} \\
$s^{(2)}_b$ & --- & $\smat{1}{0}{\ell b}{1}$, representatives of $\Ga_0(\ell^{2})\backslash\Ga_0(\ell)$ (\ref{prop:levelellU}) \\
$\nu_2(N)$, $\nu_3(N)$ & --- & the numbers of fixed points of $S$, $ST$ on $\mathbb{P}^1(\Z/N)$ (\ref{prop:levelorbit}) \\
$e_1$, $e_\ell$ & --- & the $\ell$-adic ordinary projectors of $T_\ell$ at level $1$ and of $U^{(\ell)}_\ell$ at level $\ell$ (\ref{prop:stabilisation}) \\
$c_{\mathrm{st}}$, $u_{\mathrm{st}}$ & --- & $\omega\,e_\ell\,\mathrm{dg}_\ell$ on $e_1H(1)$, and the unit root of $u^{2}-T_\ell u+\ell^{n+1}$ (\ref{prop:stabilisation}) \\
$M_m(n)$, $\Talg^{\mathrm{ord}}_{m,\mathcal P}(n)$ & --- & the lattice datum and its Hecke algebra (\ref{def:hidalattice}) \\
(Hi7)--(Hi12) & --- & the predictions of \ref{rmk:hidacheckable2} \\
\bottomrule
\end{longtable}
\end{center}


%======================================================================
\part*{Part I --- The framework}
\addcontentsline{toc}{section}{Part I --- The framework}
%======================================================================

%----------------------------------------------------------------------
\section{Simplicial complexes with relations, and the universal cover}\label{sec:framework}
%----------------------------------------------------------------------
% SOURCE:  monodromy_v5.tex:51-363 (RH)
% NOTATION: rows \Rel (was R), \cring, \sigma (was s/tau), \lift{\sigma} (was
%           i(sigma)), \gamma (was g_r), t/j/s for running indices (was i) of
%           sec:notation-crosswalk applied.  The acting group of Part I is
%           written G (see the note on G below); in Parts II--V it is
%           instantiated as G = \Ga = \SL_2(\Z).
% ERRATA:  E8 (App. A) is referenced, not restated: it lands in sec:conventions.

\paragraph{What this chapter fixes.}
The combinatorial input of the whole document: a finite simplicial complex $K$, a
\emph{relation set} $\Rel$ on it, the quotient $\KR$, and the complex $\Ku$ on which a
group $G$ acts with $K$ as a strict fundamental domain. It fixes what it \emph{means}
for the pair $(K,\Rel)$ to present that quotient (Assumption~\ref{ass:fund_domain}) ---
the property \ref{H4}(iii) of \S\ref{sec:conventions} points at --- and it proves that
the elements occurring in $\Rel$ generate $G$
(Theorem~\ref{thm:fund_group_generators}). Everything here is stated for an arbitrary
finite $(K,\Rel)$: no triangulation is fixed, and a subdivision of one admissible input
is another admissible input.

\paragraph{Standing hypotheses.}
Assumptions \ref{ass:K_contractible}, \ref{ass:edge_non_ident},
\ref{ass:open_interior} and \ref{ass:fund_domain} are the hypotheses of this Part. Per
rule~\ref{A1} each statement below repeats the ones it uses; no statement of this
document may be read as holding without them merely because they were declared here.

\paragraph{A note on the letter $G$.}
Part~I is a framework for an arbitrary group acting on an arbitrary complex, and that
group is written $G$. From \S\ref{sec:conventions} onwards it is instantiated once and
for all as $G=\Ga=\SL_2(\Z)$ acting on $\Ku\cong\Xbulk$ (\ref{H4}), and the letter $G$
is not used again. $G$ is the one symbol of these four chapters that
\S\ref{sec:notation} does not tabulate; it is used here in exactly the sense [v5] used
it, and nowhere else in this document.

\subsection*{Complexes, orientations, and the quotient by a relation}

A finite \emph{simplicial complex} $K$ on an ordered vertex set $V_0$ is a set of
non-empty subsets of $V_0$ containing every singleton and closed under passing to
non-empty subsets; $\dim\sigma=\abs{\sigma}-1$ and $\dim K=\max_\sigma\dim\sigma$;
$\abs{K}$ is the geometric realisation. $\orient{q}(K)$ denotes the set of
\emph{oriented} $q$-simplices, and throughout this document the orientation of a simplex
is the one induced by the vertex order --- the stored orientation of
\S\ref{sec:conventions}, which is where the storage conventions are fixed and which this
chapter does not restate. $\partial_t\sigma$ is the $t$-th face of $\sigma$
($0\le t\le\dim\sigma$), with its induced orientation; $t$, $j$, $s$ are the running
indices of this document (\ref{N4}).

\begin{assumption}[Contractibility of the fundamental domain]\label{ass:K_contractible}
$K$ is a finite simplicial complex whose geometric realisation $\abs{K}$ is connected
and contractible.
\end{assumption}

\begin{definition}[Vertex relation and quotient complex]\label{def:quotientcomplex}
Let $V_0$ be the vertex set of $K$ and let $\sim$ be an equivalence relation on $V_0$,
written $[v]$ for the class of $v$. The \emph{quotient complex} $\KR$ is the set of
subsets $\{[v_0],\dots,[v_q]\}$ of $V_0/\!\sim$ for which some
$\{u_0,\dots,u_q\}\in K$ has $[u_t]=[v_t]$ for all $t$.
\end{definition}

\begin{assumption}[Non-identification on edges]\label{ass:edge_non_ident}
If $v_0,v_1\in V_0$ with $v_0\sim v_1$ and $v_0\neq v_1$, then $\{v_0,v_1\}$ is not a
$1$-simplex of $K$.
\end{assumption}

\begin{lemma}\label{lemma:pR_simplicial}
\PROVED{} Assume \ref{ass:edge_non_ident}. Then
\begin{enumerate}
\item for every $\sigma=\{v_0,\dots,v_q\}\in K$ the classes $[v_0],\dots,[v_q]$ are
  pairwise distinct;
\item $\KR$ is a simplicial complex;
\item $\prR\colon K\to\KR$, $\prR(\{v_0,\dots,v_q\})=\{[v_0],\dots,[v_q]\}$, is a
  surjective simplicial map.
\end{enumerate}
\end{lemma}

\begin{proof}
(1) If $[v_t]=[v_{t'}]$ with $t\neq t'$ then $v_t\sim v_{t'}$ and $v_t\neq v_{t'}$,
while $\{v_t,v_{t'}\}$ is a face of $\sigma$ and hence a $1$-simplex of $K$ ---
contradicting \ref{ass:edge_non_ident}. (2) Every singleton $\{[v]\}$ lies in $\KR$
because $\{v\}\in K$; and if $\sigma'=\{[v_0],\dots,[v_q]\}\in\KR$ comes from
$\sigma=\{u_0,\dots,u_q\}\in K$, any non-empty subset of $\sigma'$ comes from the
corresponding face of $\sigma$, which lies in $K$. (3) Immediate from (1) and
Definition~\ref{def:quotientcomplex}; surjectivity is that definition.
\end{proof}

\begin{definition}[Boundary and interior]\label{def:bdryint}
Assume $K$ is \emph{pure} of dimension $d$, i.e.\ every maximal simplex has dimension
$d$. The \emph{boundary} $\partial K$ is the subcomplex generated by the non-maximal
simplices of $K$ that are contained in a unique maximal simplex. Write
$K^{\circ}:=K\setminus\partial K$ for the \emph{interior}; $K^{\circ}$ is in general not
a subcomplex, and we use the convention $\abs{K^{\circ}}=\abs{K}^{\circ}$.
\end{definition}

\begin{assumption}[Open interior]\label{ass:open_interior}
$K$ is pure, and for all $\sigma,\sigma'\in K$,
\[
  \sigma\neq\sigma'\ \wedge\ \prR(\sigma)=\prR(\sigma')
  \quad\Longrightarrow\quad \sigma,\sigma'\in\partial K .
\]
\end{assumption}

\noindent
\ref{ass:open_interior} says that $\KR$ is obtained by glueing $K$ to itself
\emph{along its boundary only}: no two distinct $d$-simplices are identified, and the
relation is trivial on $K^{\circ}$.

\begin{example}[Simplicial $S^1$]\label{ex:S1}
$K=\{\{0\},\{1\},\{2\},\{3\},\{0,1\},\{1,2\},\{2,3\}\}$ is the interval in three
segments, $\partial K=\{\{0\},\{3\}\}$, and $\sim$ is generated by $0\sim3$. Then
\ref{ass:K_contractible}, \ref{ass:edge_non_ident} and \ref{ass:open_interior} all hold
and $\abs{\KR}\cong S^1$. This is the shape of the boundary model $\Kbdry$ of
\S\ref{sec:boundarydatum}; it is an example, not a definition, and nothing below
specialises to it.
\end{example}

\begin{proposition}\label{prop:injectivity}
\PROVED{} Assume \ref{ass:edge_non_ident} and \ref{ass:open_interior}. Then
$\partial K/\!\sim$ is a simplicial subcomplex of $\KR$, set-theoretically
$\abs{K^{\circ}}\sqcup\abs{\partial K/\!\sim}=\abs{\KR}$, the map
$\abs{\prR}\colon\abs{K}\to\abs{\KR}$ is injective on $\abs{K}^{\circ}$, and $\prR$ is
injective on $d$-simplices of $K$.
\end{proposition}

\begin{proof}
Purity (\ref{ass:open_interior}) makes $\partial K$ well defined, and
\ref{ass:edge_non_ident} is inherited by $(\partial K,\sim)$, so
Lemma~\ref{lemma:pR_simplicial} applies to it and exhibits $\partial K/\!\sim$ as a
subcomplex. The displayed implication of \ref{ass:open_interior} says exactly that the
only identifications happen inside $\partial K$, whence the set-theoretic disjoint
decomposition and the injectivity on the interior; injectivity on $d$-simplices is the
case $\sigma,\sigma'\notin\partial K$.
\end{proof}

\begin{definition}[Saturated complexes]\label{def:saturated}
$K$ is \emph{saturated} for $\sim$ if for every $\sigma\in K$ and every subset
$\sigma'\subseteq V_0$ such that each vertex of $\sigma$ is equivalent to exactly one
vertex of $\sigma'$ and conversely, one has $\sigma'\in K$: replacing any number of
vertices of a simplex by other representatives of their classes again yields a simplex.
\end{definition}

\begin{proposition}\label{prop:saturated_complexes}
\PROVED{} Let $K$ be a finite simplicial complex with an equivalence relation $\sim$ on
its vertices. There is a saturated pair $(\widetilde K,\widetilde\sim)$ together with a
simplicial map $\iota_{\mathrm{sat}}\colon K\to\widetilde K$ descending to
$\KR\to\widetilde K/\widetilde\sim$, universal among simplicial maps $f\colon K\to L$
that descend to a map $\KR\to L/{\sim_L}$: every such $f$ factors as
$f=\widetilde f\circ\iota_{\mathrm{sat}}$ for a unique simplicial
$\widetilde f\colon\widetilde K\to L$ descending to $\widetilde K/\widetilde\sim\;\to
L/{\sim_L}$. Moreover
\begin{enumerate}
\item[(a)] for any such $(\widetilde K,\widetilde\sim)$ the induced map
  $\KR\to\widetilde K/\widetilde\sim$ is an isomorphism of simplicial complexes;
\item[(b)] $\widetilde K$ may be chosen with the same vertex set as $K$ and
  $\iota_{\mathrm{sat}}$ induced by the identity on vertices, in which case
  $\widetilde\sim\;=\;\sim$;
\item[(c)] if $(K,\sim)$ satisfies \ref{ass:K_contractible} (resp.\
  \ref{ass:edge_non_ident}) then so does $(\widetilde K,\widetilde\sim)$.
\end{enumerate}
\end{proposition}

\begin{proof}
Take $\widetilde K$ on the vertex set $V_0$ to consist of all $\sigma'\subseteq V_0$
obtained from some $\sigma\in K$ by replacing vertices with equivalent ones as in
Definition~\ref{def:saturated}, and $\iota_{\mathrm{sat}}$ the identity on vertices.
This is a simplicial complex (a subset of such a $\sigma'$ is obtained the same way
from the corresponding face of $\sigma$), it is saturated by construction, and it
satisfies (b). Any $f$ as in the statement sends equivalent vertices to equivalent
vertices, hence sends each replacement $\sigma'$ of $\sigma$ to a simplex of $L$,
which forces and defines $\widetilde f$ uniquely on vertices; this is the universal
property. Since $\KR=\widetilde K/\!\sim$ canonically --- the two complexes have the
same simplices by Definition~\ref{def:quotientcomplex} --- (a) follows, and (c) is
immediate from (b), because $\abs{\widetilde K}$ deformation retracts onto $\abs{K}$
along the identity on vertices and $\sim$ is unchanged.
\end{proof}

\subsection*{The universal cover complex}

\begin{definition}[Universal cover complex]\label{def:universal_cover}
Assume \ref{ass:K_contractible}, \ref{ass:edge_non_ident} and \ref{ass:open_interior}.
Put $X:=\abs{\KR}$, let $p\colon\widetilde X\to X$ be a universal topological cover and
let $G:=\piOne(X,x_0)$ act on $\widetilde X$ by deck transformations. A
\emph{simplicial universal cover} of $(K,\Rel)$ is a (possibly infinite) simplicial
complex $\Ku$ with:
\begin{enumerate}
\item[(a)] $\abs{\Ku}=\widetilde X$;
\item[(b)] $p$ is the realisation of a simplicial map $\pru\colon\Ku\to\KR$;
\item[(c)] there is a simplicial map $K\to\Ku$, $\sigma\mapsto\lift{\sigma}$, with
  $\prR=\pru\circ\lift{(-)}$;
\item[(d)] $G$ acts simplicially on $\Ku$, realising the deck action, and $\pru$ is
  $G$-invariant;
\item[(e)] $\Ku$ is finest with these properties: any $G$-invariant simplicial complex
  $\Ku{}'$ covering $\KR$ admits a $G$-equivariant simplicial map $\Ku{}'\to\Ku$ over
  $\KR$.
\end{enumerate}
The map of (c) is written $\sigma\mapsto\lift{\sigma}$ throughout. [v5] and the code
write $i(\sigma)$; this document cannot, because $i$ is $\sqrt{-1}$ and the elliptic
vertex $v_6$ (\S\ref{sec:notation-collisions}, C15--C16). The code identifier
\code{i(sigma)} is unaffected.
\end{definition}

\begin{proposition}\label{prop:universalcover}
\STD{\cite[\S1.3, Props.~1.33--1.34]{hatcher}} Assume \ref{ass:K_contractible},
\ref{ass:edge_non_ident} and \ref{ass:open_interior}. Then a simplicial universal cover
$(\Ku,\pru,\lift{(-)})$ exists and is unique up to unique isomorphism.
\end{proposition}

\begin{proof}
Uniqueness is the universal property (e). For existence: $\KR$ is a simplicial complex
by Lemma~\ref{lemma:pR_simplicial}, so $X=\abs{\KR}$ is a connected, locally
path-connected, semilocally simply connected CW complex and a universal cover
$p\colon\widetilde X\to X$ exists \STD{\cite[\S1.3]{hatcher}}. Take as vertices of $\Ku$
the fibre $p^{-1}(V_0/\!\sim)$, and declare
$\lift\sigma=\{\lift w_0,\dots,\lift w_q\}$ a $q$-simplex whenever
$p(\lift\sigma)$ is the vertex set of a $q$-simplex $\sigma'$ of $\KR$ \emph{and} the
characteristic map $\abs{\sigma'}\to X$ lifts to $\abs{\sigma'}\to\widetilde X$ carrying
those vertices to $\lift w_0,\dots,\lift w_q$; since $\abs{\sigma'}$ is contractible,
such a lift exists and is unique once the image of one vertex is prescribed
\STD{\cite[Props.~1.33--1.34]{hatcher}}. This gives (a), (b). For (c): $\abs{K}$ is
contractible (\ref{ass:K_contractible}), so $\abs{\prR}$ lifts uniquely to
$\abs{K}\to\widetilde X$ once a basepoint lift is chosen, and that lift carries vertices
to vertices of $\Ku$ and simplices to simplices, i.e.\ it is the realisation of a
simplicial $\sigma\mapsto\lift\sigma$ with $\pru\circ\lift{(-)}=\prR$. (d) holds because
deck transformations permute $p^{-1}(V_0/\!\sim)$ and preserve lifted characteristic
maps. (e) holds because any $G$-invariant simplicial structure on $\widetilde X$ over
$\KR$ has the lifted characteristic maps among its cells, so $\Ku$ refines it.
\end{proof}

\begin{remark}\label{rmk:injectivityofi}
Assume \ref{ass:K_contractible}, \ref{ass:edge_non_ident} and \ref{ass:open_interior}.
Then $\sigma\mapsto\lift\sigma$ is injective. Indeed, suppose
$\lift{y_1}=\lift{y_2}$ for $y_1,y_2\in\abs{K}$, so $\prR(y_1)=\prR(y_2)$. If both lie
in $\abs{K}^{\circ}$, Proposition~\ref{prop:injectivity} gives $y_1=y_2$. Otherwise, if
$y_1\neq y_2$, a path $c$ in $\abs{K}$ from $y_1$ to $y_2$ --- which exists by
\ref{ass:K_contractible} --- projects to a loop $\prR(c)$ in $X$ whose lift is the loop
$\lift{c}$ in the simply connected $\widetilde X$; hence $\prR(c)$ is nullhomotopic, and
the identification of $y_1$ with $y_2$ contributes no non-trivial element of $G$.
Injectivity on simplices follows, since a simplex is determined by its vertices. This is
the hypothesis Theorem~\ref{thm:explicit_complex_equiv} needs and the reason
\ref{ass:open_interior} is assumed.
\end{remark}

\begin{definition}[Ordered relation]\label{def:orderedrelation}
With $V_0$ ordered, put
$\Rel_{+}:=\{(u,v)\in V_0\times V_0: u\sim v,\ u<v\}$ and
$\Rel_{-}:=\{(u,v): u\sim v,\ u>v\}$; together with the diagonal these generate $\sim$.
Set
$V^{\pm}_{\mathrm{cl}}:=\{u\in V_0:\exists v,\ (u,v)\in \Rel_{\pm}\}$ and let
$K^{\pm}_{\mathrm{cl}}$ be the subcomplexes of simplices all of whose vertices lie in
$V^{\pm}_{\mathrm{cl}}$. Write $K_{\mathrm{cl}}:=K^{-}_{\mathrm{cl}}\cup
K^{+}_{\mathrm{cl}}$. Under \ref{ass:open_interior},
$K^{\pm}_{\mathrm{cl}}\subseteq\partial K$: the relation identifies simplices of
$K^{-}_{\mathrm{cl}}$ with distinct simplices of $K^{+}_{\mathrm{cl}}$, and it is
trivial on the interior.
\end{definition}

\subsection*{The relation datum, and what it means to present the quotient}

The pair $(K,\Rel)$ of this document carries more than a vertex relation: each
identification comes with the group element realising it, and \emph{simplices may be
identified with themselves} by a non-trivial element --- exactly the case that occurs at
the elliptic vertices and, centrally, everywhere (\S\ref{sec:model},
Remark~\ref{rmk:negI}). The following definition is the form of the relation datum the
rest of the document uses.

\begin{definition}[Relations, transporters, the coincidence groupoid]\label{def:relationdatum}
Let a group $G$ act simplicially and \emph{without inversions} on a simplicial complex
$\Ku$ (i.e.\ if $\gamma$ maps a simplex to itself it fixes it pointwise, hence preserves
its orientation), and let $K\subseteq\Ku$ be a finite subcomplex, with
$\lift\sigma\in\orient{q}(\Ku)$ the image of $\sigma\in\orient{q}(K)$. For
$\sigma,\sigma'\in\orient{q}(K)$ put
\[
  \mathrm{Tr}(\sigma,\sigma')\;:=\;\{\gamma\in G:\ \gamma\cdot\lift{\sigma'}=\lift{\sigma}
  \ \text{in }\orient{q}(\Ku)\}.
\]
A \emph{relation} is a triple $\varrho=(\sigma_{\mathrm{lhs}},\sigma_{\mathrm{rhs}},\gamma)$
with $\sigma_{\mathrm{lhs}},\sigma_{\mathrm{rhs}}\in\orient{q}(K)$ for a common $q$
(written $\deg\varrho=q$) and $\gamma\in\mathrm{Tr}(\sigma_{\mathrm{lhs}},\sigma_{\mathrm{rhs}})$;
$\gamma$ is its \emph{relation element} (C17: relation elements are $\gamma\in G$, never
$g$, which is the Hecke element). The set of all relations is the \emph{coincidence
groupoid} $\mathcal{C}(K,\Ku,G)$: it contains $(\sigma,\sigma,1)$ for every $\sigma$ and
is closed under
\[
 (\sigma,\sigma',\gamma),(\sigma',\sigma'',\gamma')\ \longmapsto\ (\sigma,\sigma'',\gamma\gamma'),
 \qquad
 (\sigma,\sigma',\gamma)\ \longmapsto\ (\sigma',\sigma,\gamma^{-1}).
\]
A set $\Rel$ of relations \emph{generates} $\mathcal{C}(K,\Ku,G)$ if the smallest subset
of $\mathcal{C}(K,\Ku,G)$ containing $\Rel$ and all $(\sigma,\sigma,1)$ and closed under
the two operations above is $\mathcal{C}(K,\Ku,G)$ itself. Note
$\mathrm{Tr}(\sigma,\sigma)=\mathrm{Stab}_G(\lift\sigma)$ and that
$\mathrm{Tr}(\sigma,\sigma')$, when non-empty, is a coset of that stabiliser.
\end{definition}

\begin{definition}[The ineffective kernel $N$]\label{def:kernelN}
Let $G$ act simplicially and without inversions on $\Ku$ as in
Definition~\ref{def:relationdatum}. Write
\[
  N\;:=\;\ker\bigl(G\longrightarrow\mathrm{Aut}(\Ku)\bigr)
  \;=\;\{\gamma\in G:\ \gamma\cdot\lift\tau=\lift\tau\ \text{for every oriented simplex
  $\lift\tau$ of }\Ku\}
\]
for the subgroup acting \emph{trivially} on $\Ku$; the action is \emph{effective} when
$N=1$. $N$ is normal in $G$, is contained in $\mathrm{Stab}_G(\lift\sigma)$ for every
$\sigma$, and every non-empty transporter $\mathrm{Tr}(\sigma,\sigma')$ is a union of
cosets of $N$. Both cases occur in this document: $N=1$ in the free instance of
Remark~\ref{rmk:twosettings}, while for $G=\Ga=\SL_2(\Z)$ acting on $\Ku\cong\Xbulk$ one
has $N=\{\pm I\}$, because $-I$ acts trivially on $\Hyp$ and, by \ref{H4}(iv), no other
element of $\Ga$ fixes every simplex. A \emph{non-effective} action is therefore an
explicit, supported case of this Part, not an unstated exception.
\end{definition}

\begin{assumption}[Fundamental domain; presenting the quotient]\label{ass:fund_domain}
Let $G$, $\Ku$, $K$, $\Rel$ be as in Definition~\ref{def:relationdatum}, with
$\abs{\Ku}$ connected, and let $N\trianglelefteq G$ be the ineffective kernel of
Definition~\ref{def:kernelN}. We say that \textbf{$K$ is a strict fundamental domain and
$(K,\Rel)$ presents the quotient} when all four of the following hold. Write
$\widetilde D:=\abs{K}\setminus\abs{K_{\mathrm{cl}}}$, where $K_{\mathrm{cl}}\subseteq K$
is the subcomplex generated by the simplices that a coincidence \emph{moves}: those
$\sigma\in\orient{\bullet}(K)$ admitting a relation
$(\sigma,\sigma',\gamma)\in\mathcal{C}(K,\Ku,G)$ with
$\gamma\cdot\lift\sigma\neq\lift\sigma$ --- equivalently, with $\sigma'\neq\sigma$, the
two conditions agreeing because $\gamma\lift{\sigma'}=\lift\sigma$ and $\gamma$ is
injective on simplices. A simplex merely \emph{fixed} by a non-trivial element of $G$ is
therefore \emph{not} in $K_{\mathrm{cl}}$: its isotropy is the business of \ref{FD4}, not
of $K_{\mathrm{cl}}$ (Remark~\ref{rmk:fd13centre}). In the free instance of
Remark~\ref{rmk:twosettings} every stabiliser is trivial and this $K_{\mathrm{cl}}$ is the
one of Definition~\ref{def:orderedrelation}.
\begin{enumerate}[label=(FD\arabic*)]
\item\label{FD1} $\abs{K}$ is connected, $\widetilde D$ is connected and non-empty, and
  $\widetilde D$ is dense in $\abs{K}$: $\overline{\widetilde D}=\abs{K}$.
\item\label{FD2} $G\cdot K=\Ku$: every simplex of $\Ku$ is $\gamma\cdot\lift\sigma$ for
  some $\gamma\in G$ and $\sigma\in K$.
\item\label{FD3} $\gamma\widetilde D\cap\widetilde D=\emptyset$ for every
  $\gamma\in G\setminus N$, i.e.\ for every $\gamma$ acting non-trivially on $\Ku$.
  Since $\gamma\widetilde D$ depends only on the coset $\gamma N$, this says exactly that
  the translates $\gamma\widetilde D$ are pairwise disjoint as $\gamma$ ranges over
  $G/N$. (What the condition \emph{buys} --- the decomposition of $\abs{\Ku}$, and that
  all identifications occur on $\abs{K_{\mathrm{cl}}}$ --- is
  Remark~\ref{rmk:fdconsequences}; it is a consequence, not part of the hypothesis.)
\item\label{FD4} \textbf{Completeness.} $\Rel$ generates the coincidence groupoid
  $\mathcal{C}(K,\Ku,G)$ in the sense of Definition~\ref{def:relationdatum}.
\end{enumerate}
\end{assumption}

\begin{remark}[Why $N$ is quantified over and never divided out]\label{rmk:fd13centre}
\ref{FD1} and \ref{FD3} are the two conditions of
Assumption~\ref{ass:fund_domain} that speak about \emph{motion}, and both are stated
modulo $N$. Read without $N$ they degenerate the moment the action is non-effective, and
the seven-vertex model is exactly such a case: every one of its simplices admits the
coincidence $(\sigma,\sigma,-I)$, so a $K_{\mathrm{cl}}$ built from ``$\sigma$ admits
$(\sigma,\sigma',\gamma)$ with $(\sigma',\gamma)\neq(\sigma,1)$'' would be all of $K$,
leaving $\widetilde D=\emptyset$ and \ref{FD1} unsatisfiable; and
``$\gamma\widetilde D\cap\widetilde D=\emptyset$ for $\gamma\neq1$'' would fail at
$\gamma=-I$ whatever $\widetilde D$ is. Quantifying modulo $N$, and building
$K_{\mathrm{cl}}$ from the simplices actually moved, is what makes both conditions say
what they were meant to say.

\textbf{What must not be done instead is to replace $G$ by $G/N$.} The action on $\Ku$
factors through $G/N$; \textbf{the action on the coefficients does not.} For the model,
$-I$ is the identity of $\mathrm{Aut}(\Ku)$ but acts on $\Vcoef=\Sym^n$ as
$\rhon(-I)=(-1)^{n}\id$ (Definition~\ref{def:rhon}), so for odd $n$ it acts as $-\id$;
the constraint $f(\sigma)=\rhon(-I)f(\sigma)$ that \ref{FD4} therefore imposes on
\emph{every} simplex is precisely the odd-$n$ vanishing mechanism
$\Heq{\bullet}(K,\Rel;L)=0$ of Remark~\ref{rmk:negI}, and passing to $G/N$ would delete
it. Hence $G$ stays $\SL_2(\Z)$, $N$ is quantified over in \ref{FD1} and \ref{FD3}, and
the whole of the isotropy --- $N$ included --- is carried by \ref{FD4}, which is where
isotropy belongs (Lemma~\ref{lem:fd4criterion}\ref{FD4a}).
\end{remark}

\begin{remark}[What \ref{FD1}--\ref{FD3} buy: the decomposition of $\abs{\Ku}$]
\label{rmk:fdconsequences}
\PROVED{} Assume \ref{FD1}--\ref{FD3}. Then no translate of $\widetilde D$ meets
$\abs{K_{\mathrm{cl}}}$, and
\begin{equation}\label{eq:fddecomp}
  \abs{\Ku}\;=\;\Bigl(\ \bigsqcup_{\gamma N\in G/N}\gamma\widetilde D\ \Bigr)
  \ \sqcup\ \bigcup_{\gamma N\in G/N}\gamma\abs{K_{\mathrm{cl}}}\,,
\end{equation}
where the first union is disjoint and the second in general is not. Moreover
$\overline{\widetilde D}=\abs{K}$ maps onto $\abs{\KR}$ under $\abs{\pru}$, and any two
distinct points of $\abs{K}$ with the same image both lie in $\abs{K_{\mathrm{cl}}}$.

\textbf{The form ``hence $\abs{\Ku}=\bigsqcup_{\gamma\in G}\gamma\widetilde D$'' would be
false}, and does not follow from the disjointness in \ref{FD3} --- not even for an
effective action: the translates of $\widetilde D$ omit the translates of
$\abs{K_{\mathrm{cl}}}$, which is the whole identified boundary. They can be made to
cover only by adding back one representative of each identified boundary simplex, i.e.\
by passing to a \emph{half-open} fundamental domain. Nothing in this document needs one,
and none is constructed.
\end{remark}

\begin{proof}[Proof of Remark~\ref{rmk:fdconsequences}]
\emph{No translate of $\widetilde D$ meets $\abs{K_{\mathrm{cl}}}$.} Since
$K_{\mathrm{cl}}$ is a subcomplex, $\abs{K_{\mathrm{cl}}}$ is the union of the relative
interiors of the $\lift\sigma$, $\sigma\in K_{\mathrm{cl}}$; so let $z$ lie in the
relative interior of such a $\lift\sigma$ and suppose $z\in\gamma\widetilde D$. Then
$\gamma^{-1}z\in\widetilde D\subseteq\abs{K}$. The relative interiors of the simplices of
$\Ku$ partition $\abs{\Ku}$ and $\gamma^{-1}$ is simplicial, so $\gamma^{-1}z$ lies in
the relative interior of $\gamma^{-1}\lift\sigma$; as it lies in $\abs{K}$ and $K$ is a
subcomplex of $\Ku$, that simplex is $\lift\tau$ for some $\tau\in K$. If
$\tau\neq\sigma$ then $(\sigma,\tau,\gamma)\in\mathcal{C}(K,\Ku,G)$ moves $\tau$, so
$\tau\in K_{\mathrm{cl}}$; if $\tau=\sigma$ then $\tau\in K_{\mathrm{cl}}$ by hypothesis
on $\sigma$. Either way $\gamma^{-1}z\in\abs{K_{\mathrm{cl}}}$, contradicting
$\gamma^{-1}z\in\widetilde D=\abs{K}\setminus\abs{K_{\mathrm{cl}}}$. Hence
$\gamma\widetilde D\cap\abs{K_{\mathrm{cl}}}=\emptyset$ for every $\gamma\in G$, and
therefore $\gamma'\widetilde D\cap\gamma\abs{K_{\mathrm{cl}}}=\emptyset$ for all
$\gamma,\gamma'\in G$.

\emph{Equation~\eqref{eq:fddecomp}.} The $\gamma\widetilde D$ are pairwise disjoint over
$G/N$ by \ref{FD3}. Given $z\in\abs{\Ku}$, \ref{FD2} supplies $\gamma$ with
$\gamma^{-1}z\in\abs{K}=\widetilde D\sqcup\abs{K_{\mathrm{cl}}}$, so $z$ lies in
$\gamma\widetilde D$ or in $\gamma\abs{K_{\mathrm{cl}}}$; the two families are disjoint
by the previous paragraph.

\emph{Identifications occur on $\abs{K_{\mathrm{cl}}}$.} Surjectivity of
$\overline{\widetilde D}=\abs{K}\to\abs{\KR}$ is \ref{FD1} together with the definition
of $\KR$. If $z\neq z'$ in $\abs{K}$ have the same image then $z'=\gamma z$ for some
$\gamma\in G$, necessarily $\gamma\notin N$. Both in $\widetilde D$ contradicts
\ref{FD3}; one in $\widetilde D$ and the other in $\abs{K_{\mathrm{cl}}}$ contradicts the
first paragraph. So both lie in $\abs{K_{\mathrm{cl}}}$.
\end{proof}

\begin{remark}[The two instances that occur in this document]\label{rmk:twosettings}
\ref{FD4} is the substantive condition, and it is the one \ref{H4}(iii) of
\S\ref{sec:conventions} asserts for the seven-vertex model. Two instances occur:
\begin{itemize}
\item \emph{The free instance ([v5]'s setting).} $K$ satisfies
  \ref{ass:K_contractible}, \ref{ass:edge_non_ident}, \ref{ass:open_interior};
  $G=\piOne(\abs{\KR},x_0)$; $\Ku$ is the simplicial universal cover of
  Definition~\ref{def:universal_cover}. Deck transformations of a universal cover act
  \emph{freely}, so every stabiliser $\mathrm{Tr}(\sigma,\sigma)$ is trivial --- in
  particular the ineffective kernel is $N=1$ (Definition~\ref{def:kernelN}) and
  \ref{FD3} is the usual condition ``for $\gamma\neq1$'' --- $\Rel$
  contains no relation $(\sigma,\sigma,\gamma)$ with $\gamma\neq1$, and $\Rel$ is the
  datum $\{(\sigma^-,\sigma^+,\gamma_r)\}$ indexed by $\Rel_+$ of
  Definition~\ref{def:orderedrelation}. \ref{FD1}--\ref{FD3} then hold with
  $K_{\mathrm{cl}}$ as in Definition~\ref{def:orderedrelation}, provided $\widetilde D$
  is connected, dense in $\abs{K}$ and simply connected.
\item \emph{The instance this document computes in.} $G=\Ga=\SL_2(\Z)$ acting on
  $\Ku\cong\Xbulk$, which has \emph{finite non-trivial} simplex stabilisers --- every
  one of them containing $-I$ (\ref{H4}(iv)), so the action is \emph{not} effective and
  $N=\{\pm I\}$ (Definition~\ref{def:kernelN}, Remark~\ref{rmk:fd13centre}). Here
  $\Rel$ necessarily contains
  self-relations $(\sigma,\sigma,\gamma)$, $\gamma\neq1$: the three elliptic stabiliser
  rows and the $\abs{\orient{\bullet}(K)}$ central rows of
  Definition~\ref{def:model7rel}. \ref{FD4} is then a statement with content, and
  omitting the central rows makes it false --- which is precisely
  Remark~\ref{rmk:negI}, and the reason [v5]'s ``issues with torsion, which we ignore''
  could not be carried over (Appendix~\ref{app:errata}, entry~\ref{err:e8}).
\end{itemize}
\textbf{Nothing in \S\S\ref{sec:framework}--\ref{sec:systematiccomputation} fixes a
triangulation.} A subdivision of an admissible $(K,\Rel)$, equivariant for the $G$-action,
is again an admissible input; what is triangulation-indexed ($N_q$, $\abs{\Rel}$, the
stored simplex order) is recomputed per input and appears only from
\S\ref{sec:conventions} onwards.
\end{remark}

\begin{proposition}[Relation elements]\label{prop:fund_group_action}
\PROVED{} Let $G$, $\Ku$, $K$, $\Rel$ satisfy \ref{ass:fund_domain}.
\begin{enumerate}
\item For $\sigma,\sigma'\in\orient{q}(K)$ the transporter
  $\mathrm{Tr}(\sigma,\sigma')$ is either empty or a left coset
  $\gamma\cdot\mathrm{Stab}_G(\lift{\sigma'})$, and
  $\mathrm{Tr}(\sigma,\sigma)=\mathrm{Stab}_G(\lift\sigma)$.
\item In the free instance of Remark~\ref{rmk:twosettings} --- $G=\piOne(\abs{\KR},x_0)$
  acting on the simplicial universal cover under \ref{ass:K_contractible},
  \ref{ass:edge_non_ident}, \ref{ass:open_interior} --- every non-empty transporter is a
  \emph{singleton}. In particular each pair $(u,v)\in\Rel_+$ of
  Definition~\ref{def:orderedrelation}, with $\sigma^+=\{u\}\in K^{+}_{\mathrm{cl}}$ and
  $\sigma^-=\{v\}\in K^{-}_{\mathrm{cl}}$, determines a unique
  $\gamma_r\in G$ with $\gamma_r\cdot\lift{\sigma^+}=\lift{\sigma^-}$, and likewise in
  every degree for the higher-dimensional simplices identified by $r$.
\end{enumerate}
\end{proposition}

\begin{proof}
(1) If $\gamma,\gamma'\in\mathrm{Tr}(\sigma,\sigma')$ then
$\gamma^{-1}\gamma'\cdot\lift{\sigma'}=\lift{\sigma'}$, so
$\gamma'\in\gamma\,\mathrm{Stab}_G(\lift{\sigma'})$; conversely each such element
transports. The second claim is the definition. (2) Deck transformations of a universal
cover act freely on $\widetilde X$ \STD{\cite[\S1.3]{hatcher}}, so every stabiliser is
trivial and (1) makes non-empty transporters singletons. That the pair
$(\sigma^-,\sigma^+)$ has non-empty transporter is \ref{FD2} together with
$\prR(\sigma^-)=\prR(\sigma^+)$ and $\pru\circ\lift{(-)}=\prR$.
\end{proof}


\begin{lemma}[Criterion for \ref{FD4}: stabilisers and orbits]\label{lem:fd4criterion}
\PROVED{} Let $G$, $\Ku$, $K$ be as in Definition~\ref{def:relationdatum} ($G$ acting
simplicially and without inversions, $K\subseteq\Ku$ a finite subcomplex) and let
$\Rel\subseteq\mathcal{C}(K,\Ku,G)$ be a set of relations. Write $\langle\Rel\rangle$ for the
sub-groupoid of $\mathcal{C}(K,\Ku,G)$ generated by $\Rel$ in the sense of
Definition~\ref{def:relationdatum}. Then $\Rel$ satisfies \ref{FD4} --- that is,
$\langle\Rel\rangle=\mathcal{C}(K,\Ku,G)$ --- \emph{if and only if} both of the following hold.
\begin{enumerate}[label=(\alph*)]
\item\label{FD4a} \textbf{Every stabiliser is reached.} For every
  $\sigma\in\orient{\bullet}(K)$,
  \[
    \{\gamma\in G:\ (\sigma,\sigma,\gamma)\in\langle\Rel\rangle\}
    \;=\;\mathrm{Stab}_G(\lift\sigma).
  \]
\item\label{FD4b} \textbf{Every coincidence of simplices is reached.} For every pair
  $\sigma,\sigma'\in\orient{q}(K)$ with $\mathrm{Tr}(\sigma,\sigma')\neq\emptyset$ there
  is \emph{at least one} $\gamma$ with $(\sigma,\sigma',\gamma)\in\langle\Rel\rangle$.
\end{enumerate}
Moreover it suffices, for \ref{FD4a}, that for each $\sigma$ the relation elements of the
\emph{self-relations} $(\sigma,\sigma,\gamma)\in\Rel$ generate
$\mathrm{Stab}_G(\lift\sigma)$; and for \ref{FD4b}, that each such pair
$\{\sigma,\sigma'\}$ be joined by a path of members of $\Rel$, traversed in either
direction.
\end{lemma}

\begin{remark}[Why the criterion is the right shape]\label{rmk:fd4criterion}
\ref{FD4a} and \ref{FD4b} split \ref{FD4} into its two independent contents: a
\emph{local} one --- $\Rel$ must see each simplex's full isotropy --- and a \emph{global}
one --- $\Rel$ must connect each $G$-orbit of simplices of $K$. Neither implies the
other, and the split is exactly the one the seven-vertex model needs: its eight listed
rows discharge \ref{FD4b} and the three elliptic stabilisers, while
\ref{FD4a} at the remaining twenty simplices is discharged by, and only by, the central
rows $(\sigma,\sigma,-I)$ of Definition~\ref{def:model7rel} (Remark~\ref{rmk:negI},
Proposition~\ref{prop:model7h4}).

Both conditions are finite and decidable for finite $K$ with finite stabilisers, which is
what makes \ref{FD4} a formalization target: \ref{FD4a} is one equality of subgroups of a
finite group per simplex of $K$, and \ref{FD4b} is a reachability question in the finite
graph on $\orient{\bullet}(K)$ whose edges are the members of $\Rel$, compared against the
partition of $\orient{\bullet}(K)$ into $G$-orbits. No enumeration of $\Ku$, and no
enumeration of $G$, is involved in either.
\end{remark}

\begin{proof}[Proof of Lemma~\ref{lem:fd4criterion}]
\emph{Necessity.} Both conditions are special cases of
$\langle\Rel\rangle=\mathcal{C}(K,\Ku,G)$: the members of $\mathcal{C}(K,\Ku,G)$ of the form
$(\sigma,\sigma,-)$ are precisely the elements of $\mathrm{Stab}_G(\lift\sigma)$, and
$\mathrm{Tr}(\sigma,\sigma')\neq\emptyset$ says that some triple $(\sigma,\sigma',\gamma)$
lies in $\mathcal{C}(K,\Ku,G)$.

\emph{Sufficiency.} Let $(\sigma,\sigma',\gamma)\in\mathcal{C}(K,\Ku,G)$. Then
$\gamma\in\mathrm{Tr}(\sigma,\sigma')$, which is therefore non-empty, so \ref{FD4b}
supplies a $\gamma_0$ with $(\sigma,\sigma',\gamma_0)\in\langle\Rel\rangle$. By
Proposition~\ref{prop:fund_group_action}(1) the transporter is the left coset
$\gamma_0\,\mathrm{Stab}_G(\lift{\sigma'})$, so
$h:=\gamma_0^{-1}\gamma\in\mathrm{Stab}_G(\lift{\sigma'})$, and \ref{FD4a} at $\sigma'$
gives $(\sigma',\sigma',h)\in\langle\Rel\rangle$. Composing the two,
\[
  (\sigma,\sigma',\gamma_0)\cdot(\sigma',\sigma',h)
  \;=\;(\sigma,\sigma',\gamma_0h)\;=\;(\sigma,\sigma',\gamma)\ \in\ \langle\Rel\rangle.
\]

\emph{The two sufficient conditions.} $\langle\Rel\rangle$ is closed under composition and
inversion, so for a fixed $\sigma$ the set $\{\gamma:(\sigma,\sigma,\gamma)\in\langle\Rel\rangle\}$
is a \emph{subgroup} of $\mathrm{Stab}_G(\lift\sigma)$ --- it is contained in the
stabiliser because every member of $\mathcal{C}(K,\Ku,G)$ is --- and it contains the
relation element of every self-relation of $\Rel$ at $\sigma$; if those generate the
stabiliser, the subgroup is all of it. Likewise a path
$\sigma=\sigma_0,\sigma_1,\dots,\sigma_m=\sigma'$ whose consecutive terms are the two
simplices of a member of $\Rel$ composes, after inverting the members traversed
backwards, to a member of $\langle\Rel\rangle$ over the pair $(\sigma,\sigma')$.
\end{proof}
\begin{theorem}[Generation by the relation elements]\label{thm:fund_group_generators}
\PROVED{} Let $G$, $\Ku$, $K$, $\Rel$ satisfy \ref{ass:fund_domain}, with $\abs{\Ku}$
connected and $K$ finite. Then $G$ is generated by the relation elements
\[
  \mathcal{E}(\Rel)\;:=\;\{\gamma_\varrho:\varrho\in\Rel\}\subseteq G .
\]
If moreover $K$ satisfies \ref{ass:K_contractible}, \ref{ass:edge_non_ident} and
\ref{ass:open_interior} and $G=\piOne(X,x_0)$, $X=\abs{\KR}$, acts on the simplicial
universal cover of Definition~\ref{def:universal_cover}, then for $r=(u,v)\in\Rel_{+}$
and any path $c$ in $\abs{K}$ from $u$ to $v$ whose interior lies in $\abs{K}^{\circ}$
and passes through a point $x_K$ with $\prR(x_K)=x_0$, the loop $\prR(c)$ represents
$\gamma_r\in\piOne(X,x_0)$, independently of the choice of $c$.
\end{theorem}

\begin{proof}
Put $\Sigma:=\{\gamma\in G:\gamma K\cap K\neq\emptyset\ \text{as sets of simplices of }\Ku\}$
and $H:=\langle\Sigma\rangle$.

\emph{$\Sigma\subseteq\langle\mathcal E(\Rel)\rangle$.} Let $\gamma\in\Sigma$ and let
$\lift{\sigma}=\gamma\cdot\lift{\sigma'}$ with $\sigma,\sigma'\in K$; as $G$ acts without
inversions the equality holds for the stored orientations, so
$(\sigma,\sigma',\gamma)\in\mathcal{C}(K,\Ku,G)$. By \ref{FD4} that triple is obtained
from members of $\Rel$ and identity triples by composition and inversion, so $\gamma$ is
a word in $\mathcal E(\Rel)^{\pm1}$.

\emph{$H=G$.} Since $\abs{\Ku}$ is connected, its $1$-skeleton is connected; fix a
vertex $\lift{v_\ast}$ with $v_\ast\in K$. We show by induction on the edge-path
distance from $\lift{v_\ast}$ that every vertex $\lift w$ of $\Ku$ lies in $\eta K$ for
some $\eta\in H$. For distance $0$ take $\eta=1$. If $\lift w$ is adjacent to a vertex
$\lift u\in\eta K$, $\eta\in H$, then by \ref{FD2} the edge $\{\lift u,\lift w\}$ lies
in $\delta K$ for some $\delta\in G$; then $\lift u\in\eta K\cap\delta K$, so
$\eta^{-1}\delta\in\Sigma\subseteq H$, so $\delta\in H$, and $\lift w\in\delta K$. Now
let $\gamma\in G$ be arbitrary and pick a vertex $v\in K$. By the above,
$\gamma\lift v\in\delta K$ for some $\delta\in H$, i.e.\ $\delta^{-1}\gamma\lift v$ is a
vertex of $K$; hence $\delta^{-1}\gamma\in\Sigma\subseteq H$ and $\gamma\in H$. Combining
with the previous paragraph, $G=H=\langle\mathcal E(\Rel)\rangle$.

\emph{The loop description.} Assume the further hypotheses. The lift of $\prR(c)$
starting at $\lift u$ is $\lift c$, which ends at $\lift v$; by definition of the deck
action the loop $\prR(c)$, based at $\prR(u)=\prR(v)$, represents the unique
$\gamma\in G$ with $\gamma\lift u=\lift v$, which is $\gamma_r$
(Proposition~\ref{prop:fund_group_action}(2)). Changing the basepoint from $\prR(u)$ to
$x_0$ conjugates by the path $\prR(c|_{[u,x_K]})$, whose lift $\lift{c}|_{[u,x_K]}$ lies
in $\abs{K}=\overline{\widetilde D}$; carrying out the conjugation, the lift of the
based loop $\prR(c|_{[x_K,u]})\cdot\prR(c)\cdot\prR(c|_{[u,x_K]})$ starting at
$\lift{x_K}$ ends at $\gamma_r\lift{x_K}$, so that loop represents $\gamma_r$ in
$\piOne(X,x_0)$. Independence of $c$: any two admissible paths are homotopic rel
endpoints in the contractible $\abs{K}$ (\ref{ass:K_contractible}), and the homotopy
projects.
\end{proof}

\paragraph{Code map.}
Definition~\ref{def:relationdatum}: \code{SimplexRelation}
(\file{LocalSystemSimplicialComplex.h}) is the triple
$(\sigma_{\mathrm{lhs}},\sigma_{\mathrm{rhs}},\gamma)$, stored with its
$\gamma\in\Mat_2(\Z)$ and the derived matrix $\rhon(\gamma)$; the degree is
\code{q\_dim} and the two simplices are \code{s\_lhs\_idx\_in\_Kq},
\code{s\_rhs\_idx\_in\_Kq}. \ref{FD4} for the model:
\code{SL2ZRelationSeeds\_7VertexModel} plus
\code{GenerateSL2ZRelationsWithGroupData} (\file{LocalSystemSL2Z.cpp}), checked by
\code{HeckeTessellationValidation} and \code{HeckeTessellationSmokeValidation}.
Definition~\ref{def:universal_cover}(c): \code{i(sigma)} in the $\tau$-table sources
(\S\ref{sec:reduction}).

\paragraph{Formalization note.}
Definition~\ref{def:relationdatum} and \ref{FD4} are the formalization-friendly core:
for a finite $K$ and a decidable equality on $\orient{\bullet}(\Ku)$, the coincidence
groupoid is a decidable relation and ``$\Rel$ generates it'' is a reachability
statement in a finite graph once the relevant simplices are enumerated. What needs
genuine topology is only Proposition~\ref{prop:universalcover} (existence of the
universal cover) and the free-instance half of
Proposition~\ref{prop:fund_group_action}; the instance actually computed in
(Remark~\ref{rmk:twosettings}, second bullet) takes $\Ku\cong\Xbulk$ and its
$\Ga$-action as given by Proposition~\ref{prop:model7h4} and never constructs a
universal cover. Theorem~\ref{thm:fund_group_generators}'s first half is a finite
induction over an edge path and is the natural first target.

%----------------------------------------------------------------------
\section{Matrix local systems and the three equivalent categories}\label{sec:localsystems}
%----------------------------------------------------------------------
% SOURCE:  monodromy_v5.tex:364-453 (RH)
% NOTATION: rows \Rel, \gamma (relation element, was g_r), \cring, t/j/s of
%           sec:notation-crosswalk applied; the acting group is G (see sec:framework).
% ERRATA:  none land here.

\paragraph{What this chapter fixes.}
The coefficient side of the datum: a \emph{matrix local system} is a rule assigning an
invertible matrix to each relation, subject to a cocycle condition, and this is the same
thing as a representation of $G$ (Proposition~\ref{prop:matrix_rep_equiv}) and --- over
a field, in the free instance --- the same thing as a locally constant sheaf on
$\abs{\KR}$ (Theorems~\ref{thm:lcs_rep_equiv}, \ref{thm:rep_ls_equiv}). The chapter
exists so that the computational datum $(K,\Rel,\{M_\varrho\})$ may be used with a clear
conscience: it is not an ad hoc encoding but an object of an equivalent category.

\paragraph{Standing hypotheses.}
$\cring$ is a commutative ring with $1$ and $\Vcoef$ a free $\cring$-module of finite
rank $d_V$; $G$, $\Ku$, $K$, $\Rel$ satisfy \ref{ass:fund_domain}. Where a statement
needs $\cring$ to be a \emph{field} it says so: [v5] worked over a field $\F$
throughout and did not mark where that was used, and $\cring=\Z/\ell^{e}$ is in scope
here (\S\ref{sec:notation-primes}).

\begin{definition}[Matrix local system]\label{def:matrix_ls}
A \emph{matrix local system} $L=(\Vcoef,\{M_\varrho\}_{\varrho\in\Rel})$ of rank $d_V$
on $(K,\Rel)$ consists of a free $\cring$-module $\Vcoef$ of rank $d_V$ and, for each
relation $\varrho\in\Rel$, an invertible $M_\varrho\in\GL(\Vcoef)$, subject to the
\emph{cocycle condition}: for every word
$\gamma_{\varrho_1}^{\epsilon_1}\cdots\gamma_{\varrho_m}^{\epsilon_m}=1$ in $G$ with
$\varrho_t\in\Rel$ and $\epsilon_t=\pm1$, one has
$M_{\varrho_1}^{\epsilon_1}\cdots M_{\varrho_m}^{\epsilon_m}=\id_{\Vcoef}$; and
$M_\varrho=M_{\varrho'}$ whenever $\gamma_\varrho=\gamma_{\varrho'}$. A
\emph{morphism} $L\to L'=(\Vcoef',\{M'_\varrho\})$ is an $\cring$-linear
$\phi\colon\Vcoef\to\Vcoef'$ with $\phi\circ M_\varrho=M'_\varrho\circ\phi$ for every
$\varrho\in\Rel$. The resulting category is $\LS=\LS(K,\Rel)$.
\end{definition}

\begin{remark}[The matrices are derived, never entered]\label{rmk:derivedmatrices}
In this document $\Vcoef=\Sym^n(\Z^2)\otimes\cring$ and
$M_\varrho:=\rhon(\gamma_\varrho)$ for the representation $\rhon$ of
Definition~\ref{def:rhon}; the cocycle condition is then automatic, because
$\rhon$ is a homomorphism on $\SL_2(\Z)$. That is the contract box of
\S\ref{sec:model}: the list of relations is the single source of truth for the pairs
$(\text{relation},\gamma)$, and the matrix is derived. Definition~\ref{def:matrix_ls} is
stated in the general form because the cocycle condition is what makes the general
framework work, and because an implementation that entered matrices by hand would need
it as a check.
\end{remark}

\begin{proposition}\label{prop:matrix_rep_equiv}
\PROVED{} Let $G,\Ku,K,\Rel$ satisfy \ref{ass:fund_domain} and let $\Vcoef$ be a free
$\cring$-module of finite rank. A family $\{M_\varrho\}_{\varrho\in\Rel}$ of invertible
matrices is a matrix local system in the sense of Definition~\ref{def:matrix_ls} if and
only if the assignment $\gamma_\varrho\mapsto M_\varrho$ extends to a (necessarily
unique) group homomorphism $\rho_L\colon G\to\GL(\Vcoef)$. This assignment is a
bijection, functorial in $\Vcoef$.
\end{proposition}

\begin{proof}
By Theorem~\ref{thm:fund_group_generators} (whose hypotheses are
\ref{ass:fund_domain} with $\abs{\Ku}$ connected and $K$ finite), $G$ is generated by
$\mathcal E(\Rel)=\{\gamma_\varrho\}$. A map defined on a generating set extends to a
homomorphism exactly when the images satisfy every relation the generators satisfy,
which is the cocycle condition; the extension is then unique because the
$\gamma_\varrho$ generate. The consistency requirement
$\gamma_\varrho=\gamma_{\varrho'}\Rightarrow M_\varrho=M_{\varrho'}$ of
Definition~\ref{def:matrix_ls} is what makes $\gamma_\varrho\mapsto M_\varrho$ a
well-defined map on $\mathcal E(\Rel)$ in the first place. Functoriality: a morphism
$\phi$ intertwines the $M_\varrho$ iff it intertwines $\rho_L$ on generators iff it
intertwines $\rho_L$.
\end{proof}

\begin{theorem}[Monodromy equivalence]\label{thm:lcs_rep_equiv}
\STD{\cite[\S1.3, Thm.~1.38, and \S3.H]{hatcher}} Let $X$ be a connected, locally
path-connected and semilocally simply connected topological space with basepoint $x_0$,
let $G=\piOne(X,x_0)$, and let $\cring$ be a commutative ring. The category $\LCS(X)$ of
locally constant sheaves of finitely generated free $\cring$-modules on $X$ is
equivalent to the category $\Rep(G)$ of finitely generated free $\cring$-modules with a
$G$-action. The equivalence sends a sheaf to its stalk at $x_0$ with the monodromy
action by parallel transport, with quasi-inverse
$(\Vcoef,\rho)\mapsto(\widetilde X\times \Vcoef)/G$.
\end{theorem}

\noindent
$\abs{\KR}$ satisfies those hypotheses whenever \ref{ass:edge_non_ident} holds
(Lemma~\ref{lemma:pR_simplicial} makes it the realisation of a simplicial complex, hence
a CW complex) and $\abs{K}$ is connected.

\begin{theorem}\label{thm:rep_ls_equiv}
\PROVED{} Let $G,\Ku,K,\Rel$ satisfy \ref{ass:fund_domain} with $\abs{\Ku}$ connected
and $K$ finite, and let $\cring$ be a commutative ring. Then $\Rep(G)$ and
$\LS(K,\Rel)$ are equivalent, by mutually quasi-inverse functors
\[
  S\colon(\Vcoef,\rho)\longmapsto(\Vcoef,\{\rho(\gamma_\varrho)\}_{\varrho\in\Rel}),
  \qquad
  T\colon(\Vcoef,\{M_\varrho\})\longmapsto(\Vcoef,\rho_L),
\]
each the identity on underlying modules and on morphisms.
\end{theorem}

\begin{proof}
$S$ lands in $\LS$: the cocycle condition holds because $\rho$ is a homomorphism, and
$\gamma_\varrho=\gamma_{\varrho'}$ forces equal matrices. $T$ is
Proposition~\ref{prop:matrix_rep_equiv}. Both are the identity on underlying modules,
and a linear map intertwines all $\rho(\gamma_\varrho)$ iff it intertwines $\rho$ on the
generating set $\mathcal E(\Rel)$ (Theorem~\ref{thm:fund_group_generators}) iff it
intertwines $\rho$; so both are fully faithful, and $TS=\id$, $ST=\id$ on the nose.
\end{proof}

\begin{corollary}\label{cor:threecategories}
\PROVED{} In the free instance of Remark~\ref{rmk:twosettings} --- $K$ satisfying
\ref{ass:K_contractible}, \ref{ass:edge_non_ident}, \ref{ass:open_interior},
$X=\abs{\KR}$, $G=\piOne(X,x_0)$, $\Ku$ the simplicial universal cover --- the three
categories $\LCS(X)$, $\Rep(G)$ and $\LS(K,\Rel)$ are equivalent, over any commutative
ring $\cring$.
\end{corollary}

\begin{proof}
Compose Theorems~\ref{thm:lcs_rep_equiv} and \ref{thm:rep_ls_equiv}; the hypotheses of
the former hold for $X=\abs{\KR}$ by the remark after it, those of the latter by
Remark~\ref{rmk:twosettings}.
\end{proof}

\begin{remark}[What the equivalence does \emph{not} say in the instance computed in]
\label{rmk:lcsnotfree}
Corollary~\ref{cor:threecategories} is stated for the \emph{free} instance only.
For $G=\Ga=\SL_2(\Z)$ acting on $\Xbulk$ the stabilisers are non-trivial, the quotient
is an orbifold rather than a space with fundamental group $\Ga$, and $\LCS(\abs{\KR})$
is \emph{not} equivalent to $\Rep(\Ga)$. What survives without any freeness hypothesis
is Theorem~\ref{thm:rep_ls_equiv}, $\Rep(G)\simeq\LS(K,\Rel)$ --- which is all the rest
of this document uses, since the coefficient datum enters only through $\rhon$
(Remark~\ref{rmk:derivedmatrices}) --- together with the isomorphism $\Psi$ of
Theorem~\ref{thm:explicit_complex_equiv}, which likewise needs no freeness. The
comparison with $H^{\bullet}(\Ga;\VR)$ is a separate statement carrying the hypothesis
$6\in\cring^{\times}$ at its own point of use (Remark~\ref{rmk:orbifold}, \ref{H3});
[v5]'s ``up to issues with torsion, which we ignore'' is retired
(Appendix~\ref{app:errata}, entry~\ref{err:e8}).
\end{remark}

\paragraph{Code map.}
Definition~\ref{def:matrix_ls}: \code{LocalSystemSimplicialComplexT<Elt>} carries
$\Vcoef$ as \code{representation\_dim}~$=d_V$ and the family $\{M_\varrho\}$ as the
\code{SimplexRelation::M} fields; $\rho_L$ is $\rhon$
(\code{ComputeSymmetricPowerMatrixFM}, \file{FlatMatrix.cpp}).
Remark~\ref{rmk:derivedmatrices}: \code{GenerateSL2ZRelationsWithGroupData} derives
every \code{M} from the stored $\gamma$ and never accepts one by hand;
\code{HeckeTessellationValidation} checks that.

\paragraph{Formalization note.}
Proposition~\ref{prop:matrix_rep_equiv} and Theorem~\ref{thm:rep_ls_equiv} are pure
algebra over a generating set and are formalizable as stated; over
$\cring=\Fl$ they are statements about \code{Matrix.GeneralLinearGroup} and a
\code{MonoidHom}. Theorem~\ref{thm:lcs_rep_equiv} is the only genuinely topological
input of this chapter, is used nowhere downstream (Remark~\ref{rmk:lcsnotfree}), and
would be the one item to import rather than prove.

%----------------------------------------------------------------------
\section{The explicit equivariant cochain complex}\label{sec:explicitcochaincomplex}
%----------------------------------------------------------------------
% SOURCE:  monodromy_v5.tex:454-702 (RH)
% NOTATION: rows \Rel, \sigma (was tau/s), d (coboundary; was delta/delta_simp),
%           \lift{\sigma} (was i(sigma)), \gamma, t/j/s of
%           sec:notation-crosswalk applied.
% ERRATA:  E7 (App. A) closes here: this section carries
%          thm:explicit_complex_equiv under exactly that label string.
%          E8 and E10 are referenced, not restated.

\paragraph{What this chapter fixes.}
The complex this document computes with, and the isomorphism through which every later
formula is transported onto it. Three descriptions of the same cohomology are given ---
the cochain complex of $\KR$ with local coefficients, the $G$-equivariant complex on
$\Ku$, and the \emph{explicit} complex $\Ceq(K,\Rel;L)$ of functions on
$\orient{\bullet}(K)$ constrained by $\Rel$ --- and the last two are identified by
$\Psi$ (Theorem~\ref{thm:explicit_complex_equiv}). $\Psi$ is what makes a cochain a
finite list of vectors indexed by the simplices of the \emph{finite} complex $K$, and it
is the statement six C++ source files cite
(Appendix~\ref{app:errata}, entry~\ref{err:e7}).

\paragraph{Standing hypotheses.}
$\cring$ is a commutative ring, $L=(\Vcoef,\{M_\varrho\})$ a matrix local system on
$(K,\Rel)$ with representation $\rho_L\colon G\to\GL(\Vcoef)$
(Proposition~\ref{prop:matrix_rep_equiv}), and $G,\Ku,K,\Rel$ satisfy
\ref{ass:fund_domain}. Statements that additionally need a \emph{free} action, a
\emph{contractible} $\abs{\Ku}$, or $\cring$ a field, say so in their own statement
(rule~\ref{A1}); Theorem~\ref{thm:explicit_complex_equiv}(1) needs none of the three.

\subsection*{Cochains on the quotient, and on the universal cover}

Let $L$ be a matrix local system with representation $\rho_L$. Fix a total order on the
vertices of $\KR$, orienting its simplices, and for each $\sigma'\in\orient{q}(\KR)$
choose a lift $\lift{\sigma'}\in\orient{q}(\Ku)$ with
$\pru(\lift{\sigma'})=\sigma'$. A $q$-cochain of $\KR$ with coefficients in $L$ is a
function $f\colon\orient{q}(\KR)\to\Vcoef$, so
$C^q(\KR;L)=\bigoplus_{\sigma'\in\orient{q}(\KR)}\Vcoef$, and the coboundary is
\begin{equation}\label{eq:quotientcoboundary}
  (d^{q}f)(\vartheta)\;=\;\sum_{t=0}^{q+1}(-1)^{t}\,
  \rho_L\bigl(\gamma_t(\vartheta)\bigr)\,f(\partial_t\vartheta),
  \qquad \vartheta\in\orient{q+1}(\KR),
\end{equation}
where $\gamma_t(\vartheta)\in G$ is the unique element with
$\partial_t\lift{\vartheta}=\gamma_t(\vartheta)\cdot\widetilde{(\partial_t\vartheta)}$
when the action is free, and any element of the corresponding transporter otherwise ---
in which case \eqref{eq:quotientcoboundary} is well defined only if $\rho_L$ is trivial
on the relevant stabilisers, which is one more reason the equivariant description below
is the primary one. $H^{\bullet}(\KR;L)$ denotes the cohomology of
\eqref{eq:quotientcoboundary}. The coboundary is written $d$ throughout;
$\delta$ is a coset element (\S\ref{sec:notation-collisions}, C7).

\begin{remark}[Equivariant description]\label{rmk:universalcovercohomology}
\PROVED{} For a free $G$-action on $\Ku$ and any choice of lifts as above there is a
natural isomorphism of cochain complexes
\[
  C^q(\KR;L)\;\cong\;\Hom_{\Z[G]}\bigl(C_q(\Ku),\Vcoef\bigr),
  \qquad
  f_\phi(\sigma'):=\phi(\lift{\sigma'}),
\]
where $C_{\bullet}(\Ku)$ is the integral simplicial chain complex of $\Ku$ and
$\Hom_{\Z[G]}$ means $\phi(\gamma\cdot\lift\sigma)=\rho_L(\gamma)\phi(\lift\sigma)$.
Under it the coboundary \eqref{eq:quotientcoboundary} becomes
$(d\phi)(\lift{\vartheta})=\phi(\partial\lift{\vartheta})$, with
$\partial=\sum_t(-1)^t\partial_t$ the chain boundary: a $G$-equivariant $\phi$ is
determined by its values on one lift per orbit, and the $\rho_L(\gamma_t)$ of
\eqref{eq:quotientcoboundary} are exactly the correction factors moving
$\partial_t\lift{\vartheta}$ to the chosen lift. \textbf{The right-hand side is
intrinsic: it involves no choice of lifts and no freeness.} It is the definition this
document uses, and $H^{\bullet}(\KR;L)$ is to be read through it.
\end{remark}

\begin{theorem}[Equivalence of cohomologies]\label{thm:cohom_equiv}
Let $G,\Ku,K,\Rel$ satisfy \ref{ass:fund_domain}, let $L$ be a matrix local system with
representation $\rho_L\colon G\to\GL(\Vcoef)$, and assume in addition that
\textbf{$\abs{\Ku}$ is contractible and the $G$-action on $\Ku$ is free}. Put
$X:=\abs{\KR}$ and let $\mathcal F_L$ be the locally constant sheaf on $X$
corresponding to $\rho_L$ under Theorem~\ref{thm:lcs_rep_equiv}. Then there are natural
isomorphisms
\[
  H^{\bullet}(\KR;L)\;\cong\;H^{\bullet}(G;\Vcoef)\;\cong\;H^{\bullet}(X;\mathcal F_L).
\]
\end{theorem}

\begin{proof}
\PROVED{} from the following two standard inputs. \emph{(1)} $\Ku$ is a $G$-equivariant
triangulation of the contractible space $\abs{\Ku}$ on which $G$ acts freely, so
$C_{\bullet}(\Ku;\Z)$ is a complex of free $\Z[G]$-modules --- freely permuted basis,
trivial stabilisers --- and is acyclic with $H_0=\Z$ via the augmentation
$\aug$; i.e.\ it is a free $\Z[G]$-resolution of $\Z$
\STD{\cite[Ch.~I.4]{brown}, where a free action on a complex with contractible
realisation is shown to give a free resolution of $\Z$ over $\Z[G]$}. Group cohomology
may be computed from any such resolution
\STD{\cite[Ch.~I.7.4]{brown} for the independence of the resolution,
\cite[Ch.~III.1]{brown} for the definition with coefficients}, so
$H^{\bullet}(G;\Vcoef)=H^{\bullet}\bigl(\Hom_{\Z[G]}(C_{\bullet}(\Ku),\Vcoef)\bigr)$,
which is $H^{\bullet}(\KR;L)$ by Remark~\ref{rmk:universalcovercohomology}.
\emph{(2)} $X=\abs{\Ku}/G$ with $\abs{\Ku}$ contractible and $G$ acting freely and
properly discontinuously, so $X$ is a $K(G,1)$ and $\abs{\Ku}$ is its universal cover;
cohomology of $X$ with coefficients in the local system $\mathcal F_L$ is by definition
the cohomology of $\Hom_{\Z[\piOne(X)]}(C_{\bullet}(\widetilde X),\Vcoef)$
\STD{\cite[\S3.H]{hatcher}}, which is the same complex.
Naturality in $L$ is naturality of both identifications in $\Vcoef$.
\end{proof}

\noindent
\textbf{The freeness hypothesis is not decorative.} For $G=\Ga=\SL_2(\Z)$ on
$\Xbulk$ the chain groups $C_q(\Xbulk)$ are sums of $\Z[\Ga/\Ga_\sigma]$ with
$\Ga_\sigma\neq1$, hence not free, and step~(1) fails; what replaces it, under the
hypothesis $6\in\cring^{\times}$, is the relative-projectivity argument of
Remark~\ref{rmk:orbifold}, and at $\ell\in\{2,3\}$ nothing replaces it and no comparison
is asserted (\ref{H3}). Theorem~\ref{thm:explicit_complex_equiv} below is the part of
this chapter that survives without freeness, and it is the part the implementation uses.

\subsection*{The explicit complex on $K$}

\begin{definition}[Naive cochains on $K$]\label{def:naivecochains}
$C^q(K;\Vcoef)$ is the $\cring$-module of functions $f\colon\orient{q}(K)\to\Vcoef$,
i.e.\ $\bigoplus_{\sigma\in\orient{q}(K)}\Vcoef$, with the standard coboundary
\begin{equation}\label{eq:naivecoboundary}
  (d^{q}f)(\sigma)\;=\;\sum_{t=0}^{q+1}(-1)^{t}f(\partial_t\sigma),
  \qquad\sigma\in\orient{q+1}(K),
\end{equation}
so $d^{q}=d^{q}_{1}\otimes\id_{\Vcoef}$ with $d^{q}_{1}$ the coboundary of $K$ with
trivial coefficients. The stored orientation and the flat coordinates in which
\eqref{eq:naivecoboundary} is realised are fixed once in \S\ref{sec:conventions} and are
not restated here.
\end{definition}

\noindent
$(C^{\bullet}(K;\Vcoef),d)$ computes $H^{\bullet}(K;\Vcoef)$, which under
\ref{ass:K_contractible} is $\Vcoef$ in degree $0$ and $0$ above: it carries no
information about $L$. The local system enters as a \emph{constraint}.

\begin{definition}[Equivariant cochains]\label{def:eq_cochains}
\LEAN{def:relcx} ([GT] imposes the identity over the full transporter groupoid, Lemma~\ref{lem:eqgroupoid} under \ref{FD4}; the coboundary carries $\rho_L(h_j)$, $=1$ on the model)
$\Cq{q}(K,\Rel;L)\subseteq C^q(K;\Vcoef)$ is the submodule of those $f$ satisfying, for
every relation $\varrho=(\sigma_{\mathrm{lhs}},\sigma_{\mathrm{rhs}},\gamma)\in\Rel$ of
degree $q$,
\begin{equation}\label{eq:equivariance}
  f(\sigma_{\mathrm{lhs}})\;=\;M_\varrho\,f(\sigma_{\mathrm{rhs}})
  \;=\;\rho_L(\gamma)\,f(\sigma_{\mathrm{rhs}}).
\end{equation}
Relations with $\sigma_{\mathrm{lhs}}=\sigma_{\mathrm{rhs}}$ are included and are not
vacuous: they read $\bigl(\id-\rho_L(\gamma)\bigr)f(\sigma)=0$
(Remark~\ref{rmk:negI}).
\end{definition}

\begin{lemma}[Consequences of $\Rel$ are imposed too]\label{lem:eqgroupoid}
\PROVED{} \LEAN{def:relcx} Let $f\in\Cq{q}(K,\Rel;L)$ and let $(\sigma,\sigma',\gamma)$ lie in the
sub-groupoid generated by the degree-$q$ members of $\Rel$ in the sense of
Definition~\ref{def:relationdatum}. Then $f(\sigma)=\rho_L(\gamma)f(\sigma')$. In
particular, under \ref{FD4} the identity holds for \emph{every} $\gamma$ with
$\gamma\lift{\sigma'}=\lift{\sigma}$.
\end{lemma}

\begin{proof}
The set of triples $(\sigma,\sigma',\gamma)$ for which the identity holds for all
$f\in\Cq{q}$ contains $\Rel$ and all $(\sigma,\sigma,1)$, and is closed under
composition ($f(\sigma)=\rho_L(\gamma)f(\sigma')=\rho_L(\gamma)\rho_L(\gamma')f(\sigma'')
=\rho_L(\gamma\gamma')f(\sigma'')$) and inversion ($\rho_L(\gamma)$ is invertible). So it
contains the generated sub-groupoid. Under \ref{FD4} that is all of
$\mathcal{C}(K,\Ku,G)$ in degree $q$, whose degree-$q$ members are exactly the
$(\sigma,\sigma',\gamma)$ with $\gamma\lift{\sigma'}=\lift{\sigma}$.
\end{proof}

\begin{assumption}[Consistent face identification]\label{ass:consistent_face_id}
$\Rel$ is \emph{closed under faces}: for every
$(\sigma_{\mathrm{lhs}},\sigma_{\mathrm{rhs}},\gamma)\in\Rel$ of degree $q+1$ and every
$0\le t\le q+1$, the triple
$(\partial_t\sigma_{\mathrm{lhs}},\partial_t\sigma_{\mathrm{rhs}},\gamma)$ lies in the
sub-groupoid generated by the degree-$q$ members of $\Rel$.
\end{assumption}

\begin{remark}\label{rmk:facesautomatic}
\PROVED{} Under \ref{ass:fund_domain}, \ref{ass:consistent_face_id} holds automatically
as soon as $\Rel$ contains \emph{all} of $\mathcal{C}(K,\Ku,G)$ in each degree, and for
a given $\Rel$ it is a finite check. The geometric reason is one line: $K\subseteq\Ku$ is
a subcomplex and $G$ acts simplicially, so $\lift{\partial_t\sigma}=\partial_t\lift\sigma$
and
\[
  \lift{\partial_t\sigma_{\mathrm{lhs}}}=\partial_t\lift{\sigma_{\mathrm{lhs}}}
  =\partial_t\bigl(\gamma\cdot\lift{\sigma_{\mathrm{rhs}}}\bigr)
  =\gamma\cdot\partial_t\lift{\sigma_{\mathrm{rhs}}}
  =\gamma\cdot\lift{\partial_t\sigma_{\mathrm{rhs}}} ,
\]
so the face triple is a coincidence, and \ref{FD4} puts it in the generated groupoid.
The seven-vertex model's $\Rel$ is closed under faces literally, row by row: the
$T$-edge row has the two $T$-vertex rows as faces, the $S$-edge row the $S$-vertex rows,
and each central row $(\sigma,\sigma,-I)$ has the central rows of its faces
(Definition~\ref{def:model7rel}).
\end{remark}

\begin{lemma}\label{lemma:d_maps_C_eq}
\PROVED{} \LEAN{def:relcx} Assume \ref{ass:consistent_face_id}. Then
$d^{q}\bigl(\Cq{q}(K,\Rel;L)\bigr)\subseteq\Cq{q+1}(K,\Rel;L)$, so
$\bigl(\Ceq(K,\Rel;L),d\bigr)$ is a subcomplex of $(C^{\bullet}(K;\Vcoef),d)$.
\end{lemma}

\begin{proof}
Let $f\in\Cq{q}$ and
$\varrho=(\sigma_{\mathrm{lhs}},\sigma_{\mathrm{rhs}},\gamma)\in\Rel$ of degree $q+1$.
By \ref{ass:consistent_face_id} and Lemma~\ref{lem:eqgroupoid},
$f(\partial_t\sigma_{\mathrm{lhs}})=\rho_L(\gamma)f(\partial_t\sigma_{\mathrm{rhs}})$ for
every $t$. Hence
\[
  (d^{q}f)(\sigma_{\mathrm{lhs}})=\sum_{t}(-1)^{t}f(\partial_t\sigma_{\mathrm{lhs}})
  =\rho_L(\gamma)\sum_{t}(-1)^{t}f(\partial_t\sigma_{\mathrm{rhs}})
  =\rho_L(\gamma)\,(d^{q}f)(\sigma_{\mathrm{rhs}}).\qedhere
\]
\end{proof}

\begin{definition}[Equivariant cohomology of $(K,\Rel;L)$]\label{def:eqcohomology}
$\Zq{q}(K,\Rel;L):=\ker\bigl(d^{q}|_{\Cq{q}}\bigr)$,
$\Bq{q}(K,\Rel;L):=d^{q-1}\bigl(\Cq{q-1}\bigr)$ (with $\Bq{0}=0$) and
\[
  \Heq{q}(K,\Rel;L)\;:=\;\Zq{q}(K,\Rel;L)\big/\Bq{q}(K,\Rel;L),
\]
also written $H^q(K,\Rel;L)$. The inclusion $\Bq{q}\subseteq\Zq{q}$ is
Lemma~\ref{lemma:d_maps_C_eq} together with $d\circ d=0$.
\end{definition}

\begin{theorem}[The explicit equivariant complex computes $H^{\bullet}(K,\Rel;L)$]
\label{thm:explicit_complex_equiv}
Let $G$ act simplicially and without inversions on a simplicial complex $\Ku$, let
$K\subseteq\Ku$ be a finite subcomplex, let $\Rel$ be a relation set such that
$(K,\Rel)$ presents the quotient (\ref{ass:fund_domain}) and is closed under faces
(\ref{ass:consistent_face_id}), and let $L=(\Vcoef,\{M_\varrho\})$ be a matrix local
system over a commutative ring $\cring$ with representation
$\rho_L\colon G\to\GL(\Vcoef)$ (Proposition~\ref{prop:matrix_rep_equiv}). Then:
\begin{enumerate}
\item \PROVED{} \LEAN{thm:psi} (with \eqref{eq:psiinverse} and the naturality in $L$; not (2)) The map
\[
  \Psi\colon\;\Hom_{\Z[G]}\bigl(C_{q}(\Ku),\Vcoef\bigr)\longrightarrow\Cq{q}(K,\Rel;L),
  \qquad (\Psi\phi)(\sigma):=\phi(\lift\sigma),
\]
is an isomorphism of cochain complexes, natural in $L$, with inverse
$\Phi(f)=\phi_f$ determined by
\begin{equation}\label{eq:psiinverse}
  \phi_f\bigl(\gamma\cdot\lift\sigma\bigr)\;=\;\rho_L(\gamma)\,f(\sigma),
  \qquad\gamma\in G,\ \sigma\in\orient{q}(K).
\end{equation}
No freeness, contractibility or field hypothesis is used.
\item \PROVED{} If in addition $\abs{\Ku}$ is contractible and the $G$-action is free,
then $H^{\bullet}(\KR;L)\cong\Heq{\bullet}(K,\Rel;L)$, and by
Theorem~\ref{thm:cohom_equiv} both are $H^{\bullet}(G;\Vcoef)$ and
$H^{\bullet}(\abs{\KR};\mathcal F_L)$.
\end{enumerate}
\end{theorem}

\begin{proof}
(1) \emph{$\Psi$ lands in $\Cq{q}$.} Let $\varrho=(\sigma_{\mathrm{lhs}},\sigma_{\mathrm{rhs}},\gamma)\in\Rel$
of degree $q$, so $\lift{\sigma_{\mathrm{lhs}}}=\gamma\cdot\lift{\sigma_{\mathrm{rhs}}}$
by Definition~\ref{def:relationdatum}. Then
$(\Psi\phi)(\sigma_{\mathrm{lhs}})=\phi(\gamma\cdot\lift{\sigma_{\mathrm{rhs}}})
=\rho_L(\gamma)\phi(\lift{\sigma_{\mathrm{rhs}}})=M_\varrho(\Psi\phi)(\sigma_{\mathrm{rhs}})$.

\emph{$\Psi$ is a chain map.} $K\subseteq\Ku$ is a subcomplex, so
$\lift{\partial_t\sigma}=\partial_t\lift\sigma$ and
\[
  (d\,\Psi\phi)(\sigma)=\sum_t(-1)^t\phi(\lift{\partial_t\sigma})
  =\phi\Bigl(\sum_t(-1)^t\partial_t\lift\sigma\Bigr)=\phi(\partial\lift\sigma)
  =(d\phi)(\lift\sigma)=\bigl(\Psi(d\phi)\bigr)(\sigma),
\]
the fourth equality by Remark~\ref{rmk:universalcovercohomology}.

\emph{$\Phi$ is well defined.} Every $\lift{\vartheta}\in\orient{q}(\Ku)$ is
$\gamma\cdot\lift\sigma$ for some $\gamma\in G$, $\sigma\in\orient{q}(K)$, by \ref{FD2}.
Suppose $\gamma_1\lift{\sigma_1}=\gamma_2\lift{\sigma_2}$. Then
$\gamma_1^{-1}\gamma_2\cdot\lift{\sigma_2}=\lift{\sigma_1}$, i.e.\
$(\sigma_1,\sigma_2,\gamma_1^{-1}\gamma_2)\in\mathcal{C}(K,\Ku,G)$, so by
Lemma~\ref{lem:eqgroupoid} (using \ref{FD4})
$f(\sigma_1)=\rho_L(\gamma_1^{-1}\gamma_2)f(\sigma_2)$ and therefore
$\rho_L(\gamma_1)f(\sigma_1)=\rho_L(\gamma_2)f(\sigma_2)$: the right-hand side of
\eqref{eq:psiinverse} does not depend on the presentation. Extending
$\Z$-linearly to $C_q(\Ku)$ gives $\phi_f$, and
$\phi_f(\gamma'\lift{\vartheta})=\rho_L(\gamma')\phi_f(\lift{\vartheta})$ is immediate
from \eqref{eq:psiinverse}, so $\phi_f\in\Hom_{\Z[G]}(C_q(\Ku),\Vcoef)$. \emph{This step
is the whole content of the theorem, and \ref{FD4} is exactly what it needs}: the case
$\sigma_1=\sigma_2$, $\gamma_1^{-1}\gamma_2\in\mathrm{Stab}_G(\lift{\sigma_1})\neq1$ is
the case of a non-trivial stabiliser, and it is handled by the same line --- provided
$\Rel$ carries the corresponding self-relations. If it does not, $\phi_f$ is not well
defined and $\Psi$ is not surjective; that is the content of
Remark~\ref{rmk:negI} and of \ref{H4}(iv).

\emph{$\Psi$ and $\Phi$ are mutually inverse.} $(\Psi\phi_f)(\sigma)=\phi_f(\lift\sigma)
=\rho_L(1)f(\sigma)=f(\sigma)$, so $\Psi\Phi=\id$. Conversely, with
$f_\phi:=\Psi\phi$, $\phi_{f_\phi}(\gamma\lift\sigma)=\rho_L(\gamma)\phi(\lift\sigma)
=\phi(\gamma\lift\sigma)$ by $G$-equivariance of $\phi$, so $\Phi\Psi=\id$.
Naturality in $L$ is clear from the formulas. Both maps are $\cring$-linear and
commute with $d$, so $\Psi$ is an isomorphism of cochain complexes.

(2) Under the extra hypotheses, Remark~\ref{rmk:universalcovercohomology} identifies
$C^{\bullet}(\KR;L)$ with $\Hom_{\Z[G]}(C_{\bullet}(\Ku),\Vcoef)$, and (1) identifies
that with $\Ceq(K,\Rel;L)$; pass to cohomology and apply
Theorem~\ref{thm:cohom_equiv}.
\end{proof}

\begin{remark}[What the theorem is used for]\label{rmk:psiuse}
Part (1) is used, in this form and without part (2), at four places: it is the
identification $\Ceq(K,\Rel;L)\cong\Hom_{\Z[\Ga]}(C_{\bullet}(\Ku),\VR)$ of
Remark~\ref{rmk:orbifold}; it forces the central relation \eqref{eq:negI} on every
simplex (Remark~\ref{rmk:negI}); it is the passage $c\mapsto\phi_c$ of
\eqref{eq:psiinverse} on which $\Xig$ and $\Th_g$ are defined
(Definition~\ref{def:xi}, Theorem~\ref{thm:xi}); and it is what lets the induced local
system of Definition~\ref{def:inducedls} be compared with the group-cohomological
$\Ind$. [v5] \S12's ``Simplicial interpretation'' was reaching for part (1) and
garbled it; it is not re-homed (Appendix~\ref{app:errata}, entry~\ref{err:e10}).
\end{remark}

\begin{remark}[Injectivity of the lift]\label{rmk:psiinjectivelift}
Theorem~\ref{thm:explicit_complex_equiv}(1) presupposes that
$\sigma\mapsto\lift\sigma$ is injective on $\orient{q}(K)$, which is automatic here
because $K\subseteq\Ku$ is a subcomplex and the lift is the inclusion. In the free
instance, where $K$ is given abstractly and mapped into $\Ku$ by
Definition~\ref{def:universal_cover}(c), injectivity is
Remark~\ref{rmk:injectivityofi} and rests on \ref{ass:open_interior}. Without it a
single $\lift\vartheta$ could be the image of several simplices of $K$ and $\Phi$
would need a choice; [v5] flagged this and this document removes the issue by taking
$K\subseteq\Ku$.
\end{remark}

\paragraph{Code map.}
Definition~\ref{def:eq_cochains}: the rows of $A_{C^q}$
(\S\ref{sec:systematiccomputation}, \code{assemble\_A\_Cq} in
\file{EquivariantCohomologyOnK.cpp}) are exactly \eqref{eq:equivariance}, one per
relation. Definition~\ref{def:naivecochains}: \code{coboundary\_operator\_elt} is
$d^{q}_{1}$. Definition~\ref{def:eqcohomology}: \code{compute\_cohomology()} and the
accessors \code{C\_eq\_block(q)}, \code{Z\_eq\_block(q)}, \code{B\_eq\_block(q)},
\code{H\_eq\_block(q)}. Theorem~\ref{thm:explicit_complex_equiv} is cited as
\code{[C, thm:explicit\_complex\_equiv]} from \file{LocalSystemSL2Z.h},
\file{LocalSystemSL2Z.cpp}, \file{Hecke/InducedLocalSystem.cpp},
\file{Hecke/HeckeTessellationValidation.h}, \file{Hecke/HeckeTessellationValidation.cpp}
and \file{Hecke/HeckeTessellationSmokeValidation.cpp}; \eqref{eq:psiinverse} is the map
$c\mapsto\phi_c$ of \code{HeckeOperatorOnK} and \code{InducedLocalSystem}.

\paragraph{Formalization note.}
Theorem~\ref{thm:explicit_complex_equiv}(1) is finite and constructive: $\Psi$ is
restriction along a finite inclusion and $\Phi$ is defined by \eqref{eq:psiinverse},
whose well-definedness is the single lemma~\ref{lem:eqgroupoid} plus \ref{FD4}. It needs
no topology, no field, and no freeness, and is the natural formalization target of
Part~I --- the more so because it is the statement six source files cite. Part (2) and
Theorem~\ref{thm:cohom_equiv} rest on \cite[Ch.~I.4]{brown} and \cite[\S3.H]{hatcher}
and would be imported.

%----------------------------------------------------------------------
\section{Systematic computation: bases, projectors, and the stored objects}\label{sec:systematiccomputation}
%----------------------------------------------------------------------
% SOURCE:  monodromy_v5.tex:1083-1147 (RH), with the B^q step replaced by the
%          shipped one (prop:Bqimage, proved here) and the basis-name scheme of
%          sec:notation-cochains applied.
% NOTATION: rows c/z/b (was beta/delta/eta -- E11), \Rel, \cring, t/j/s.
% ERRATA:  E11 (App. A) lands here.

\paragraph{What this chapter fixes.}
\emph{Which} objects the implementation computes and stores for each degree $q$, in
which order, and by which linear algebra --- because every matrix of Parts~III--V is
written in the bases fixed here, and a Hecke matrix $\Amat{q}$ is meaningless without
them. Two linear systems, $A_{C^q}$ and $A_{Z^q}$, and three basis-extraction steps are
all that is needed.

\paragraph{Standing hypotheses.}
$\cring$ is a field --- the elimination-based basis extractions below use division, and
the implementation runs over $\cring=\Fl$; the case $\cring=\Z/\ell^{e}$ is not claimed
here. $(K,\Rel)$ satisfies \ref{ass:fund_domain} and \ref{ass:consistent_face_id}, $L$
is a matrix local system of rank $d_V$, and $N_q:=\abs{\orient{q}(K)}$, so
$\dim_{\cring}C^q(K;\Vcoef)=N_qd_V$. Cochains are written in the flat coordinates of
\eqref{eq:flatindex}; that convention is fixed in \S\ref{sec:conventions} and is not
restated.

\subsection*{The two constraint matrices}

Write $\Rel^{(q)}$ for the relations of degree $q$ and $N_{\Rel,q}:=\abs{\Rel^{(q)}}$.

\begin{proposition}[The cochain system $A_{C^q}$]\label{prop:ACq}
\PROVED{} Let $A_{C^q}$ be the $N_{\Rel,q}\times N_q$ block matrix over $\cring$, with
blocks of size $d_V\times d_V$, whose row $\varrho=(\sigma_{\mathrm{lhs}},\sigma_{\mathrm{rhs}},\gamma)$
is
\[
  \begin{cases}
    \ \id\ \text{at column }\sigma_{\mathrm{lhs}},\quad -M_\varrho\ \text{at column }\sigma_{\mathrm{rhs}},
      & \sigma_{\mathrm{lhs}}\neq\sigma_{\mathrm{rhs}},\\[2pt]
    \ \id-M_\varrho\ \text{at column }\sigma_{\mathrm{lhs}}, & \sigma_{\mathrm{lhs}}=\sigma_{\mathrm{rhs}},
  \end{cases}
\]
and zero elsewhere. Then $\Cq{q}(K,\Rel;L)=\ker A_{C^q}$ inside $\cring^{N_qd_V}$.
Moreover a relation with $\sigma_{\mathrm{lhs}}=\sigma_{\mathrm{rhs}}$ and
$M_\varrho=\id$ contributes the zero row and may be dropped; a relation with
$M_\varrho=\alpha\,\id$, $\alpha\neq1$, contributes the non-zero scalar block
$(1-\alpha)\id$ and may \emph{not}.
\end{proposition}

\begin{proof}
Row $\varrho$ of $A_{C^q}f=0$ is \eqref{eq:equivariance} verbatim, and
Definition~\ref{def:eq_cochains} imposes exactly one such equation per relation. The
two special cases are the two displayed forms evaluated at $M_\varrho=\id$ and
$M_\varrho=\alpha\id$.
\end{proof}

\noindent
The second sentence is the \emph{vacuous-relation short-circuit}: for even $n$ the
central relations of Definition~\ref{def:model7rel} have $\rhon(-I)=\id$ and drop out,
for odd $n$ and $\ell\neq2$ they have $\rhon(-I)=-\id$ and force $f=0$
(Remark~\ref{rmk:negI}). Dropping a zero row changes no value.

\begin{proposition}[The cocycle system $A_{Z^q}$]\label{prop:AZq}
\PROVED{} Let $Z^q_{1}:=\ker d^{q}_{1}\subseteq\cring^{N_q}$ be the trivial-coefficient
cocycle space, with a fixed ordered basis $z^{(q)}_{1,1},\dots,z^{(q)}_{1,k_q}$,
$k_q=\dim Z^q_1$; write $z^{(q)}_{1,j}(\sigma)\in\cring$ for its component at
$\sigma\in\orient{q}(K)$, and
\[
  E^{(q)}\colon\;\cring^{k_qd_V}\longrightarrow\cring^{N_qd_V},
  \qquad (v_1,\dots,v_{k_q})\longmapsto\sum_{j=1}^{k_q}z^{(q)}_{1,j}\otimes v_j ,
\]
the $N_q\times k_q$ block matrix with scalar block $z^{(q)}_{1,j}(\sigma)\,\id$ at
$(\sigma,j)$. Let $A_{Z^q}$ be the $N_{\Rel,q}\times k_q$ block matrix whose
$(\varrho,j)$ block is the pencil
\[
  z^{(q)}_{1,j}(\sigma_{\mathrm{lhs}})\,\id\;-\;z^{(q)}_{1,j}(\sigma_{\mathrm{rhs}})\,M_\varrho .
\]
Then $E^{(q)}$ is injective and
\[
  \Zq{q}(K,\Rel;L)\;=\;E^{(q)}\bigl(\ker A_{Z^q}\bigr).
\]
\end{proposition}

\begin{proof}
Since $d^{q}=d^{q}_{1}\otimes\id_{\Vcoef}$ (Definition~\ref{def:naivecochains}),
$\ker d^{q}=Z^q_1\otimes\Vcoef$, and $\Zq{q}=(\ker d^{q})\cap\Cq{q}$ by
Lemma~\ref{lemma:d_maps_C_eq}. Every $f\in Z^q_1\otimes\Vcoef$ is uniquely
$f=\sum_j z^{(q)}_{1,j}\otimes v_j$ --- uniqueness, hence injectivity of $E^{(q)}$,
because the $z^{(q)}_{1,j}$ are linearly independent and $\Vcoef$ is free. Evaluating at
$\sigma$ gives $f(\sigma)=\sum_j z^{(q)}_{1,j}(\sigma)\,v_j$, so
\eqref{eq:equivariance} for $\varrho$ reads
$\sum_j\bigl(z^{(q)}_{1,j}(\sigma_{\mathrm{lhs}})\id
-z^{(q)}_{1,j}(\sigma_{\mathrm{rhs}})M_\varrho\bigr)v_j=0$, which is row $\varrho$ of
$A_{Z^q}v=0$.
\end{proof}

\noindent
The point of $A_{Z^q}$ is size: it has $k_qd_V$ columns where $A_{C^q}$ has $N_qd_V$,
and $k_q\le N_q$ is typically much smaller. A block with both scalars zero is
structurally absent; a block with $z^{(q)}_{1,j}(\sigma_{\mathrm{rhs}})=0$, or with
$M_\varrho$ scalar, is itself scalar and carries no data.

\subsection*{The stored bases, degree by degree}

\begin{definition}[The stored objects]\label{def:storedbases}
For each degree $q$ the implementation computes and stores
\begin{center}
\begin{tabular}{@{}lll@{}}
\toprule
object & what it is & how it is obtained \\
\midrule
$c^{(q)}_1,\dots,c^{(q)}_{m_q}$ & a basis of $\Cq{q}$, $m_q=\dim\Cq{q}$ & kernel basis of $A_{C^q}$ \\
$z^{(q)}_1,\dots,z^{(q)}_{d_{Z,q}}$ & a basis of $\Zq{q}$ & $E^{(q)}$ applied to a kernel basis of $A_{Z^q}$ \\
$b^{(q)}_1,\dots,b^{(q)}_{d_{B,q}}$ & a basis of $\Bq{q}$ & image basis of $d^{q-1}$ on $\Cq{q-1}$ (Proposition~\ref{prop:Bqimage}) \\
$w^{(q)}_1,\dots,w^{(q)}_{m_q-d_{Z,q}}$ & a basis of $W^q:=\Cq{q}/\Zq{q}$ & the $c^{(q)}_j$ at the pivots of $[\,z^{(q)}\,|\,c^{(q)}\,]$ \\
$\mathfrak{h}_q\subseteq\{1,\dots,d_{Z,q}\}$ & the indices of a basis of $\Heq{q}$ & the pivots of $[\,b^{(q)}\,|\,z^{(q)}\,]$ \\
$P_q$ & the projector of Proposition~\ref{prop:projector} & from the same elimination \\
\bottomrule
\end{tabular}
\end{center}
all in the flat coordinates \eqref{eq:flatindex}. The basis names $c/z/b$ replace [v5]'s
$\beta^{(q)}_j$, $\delta^{(q)}_s$, $\eta^{(q)}_p$, each of which collided with a reserved
symbol (Appendix~\ref{app:errata}, entry~\ref{err:e11}); $w^{(q)}_j$ keeps [v5]'s letter
for the quotient $W^q$. The stored basis of $\Heq{q}$ is the sub-family
$\{[z^{(q)}_j]:j\in\mathfrak{h}_q\}$, re-indexed as $[z^{(q)}_1],\dots,[z^{(q)}_{d_q}]$,
$d_q:=\dim\Heq{q}=\abs{\mathfrak{h}_q}$, wherever a later chapter writes the stored
cohomology basis (Algorithm~\ref{alg:assembly}).
\end{definition}

\begin{contractbox}
\textbf{The basis choice is part of the datum.} $\Amat{q}$, the projectors and every
matrix of Parts~III--V are written in the stored bases of
Definition~\ref{def:storedbases}, and those are fixed by a deterministic elimination
rule: leftmost pivot column, topmost non-zero pivot row, leading $1$, elimination above
and below; kernel bases are the reduced free-column vectors in ascending order of free
column; image bases are the \emph{original} columns at the pivots; a factor basis is the
greedy leftmost-independent family of generators modulo the relations. Changing the rule
changes the stored bases and hence the matrices, without changing any of the invariants
of \S\ref{sec:obligations}. Cross-checks that compare matrices entry by entry are
comparing this convention as well.
\end{contractbox}

\begin{proposition}[Coboundaries without $W^{q-1}$]\label{prop:Bqimage}
\PROVED{} Assume \ref{ass:consistent_face_id} and $q\ge1$. Let $C^{(q-1)}$ be the
$N_{q-1}d_V\times m_{q-1}$ matrix whose columns are $c^{(q-1)}_1,\dots,c^{(q-1)}_{m_{q-1}}$.
Then the column $j$ of $d^{q-1}C^{(q-1)}$ is a pivot column of $d^{q-1}C^{(q-1)}$ if and
only if $c^{(q-1)}_j$ is a pivot column of $[\,z^{(q-1)}\,|\,c^{(q-1)}\,]$; consequently
\[
  \bigl\{\text{pivot columns of }d^{q-1}C^{(q-1)}\bigr\}
  \;=\;\bigl\{d^{q-1}w^{(q-1)}_1,\dots,d^{q-1}w^{(q-1)}_{m_{q-1}-d_{Z,q-1}}\bigr\}
\]
as ordered families, and either side is a basis of $\Bq{q}$.
\end{proposition}

\begin{proof}
With the leftmost-pivot rule, column $j$ of a matrix is a pivot column iff it is not in
the span of the columns before it. So column $j$ of $d^{q-1}C^{(q-1)}$ is a pivot column
iff $d^{q-1}c^{(q-1)}_j\notin\mathrm{span}\{d^{q-1}c^{(q-1)}_t:t<j\}$, i.e.\ iff
$c^{(q-1)}_j\notin\mathrm{span}\{c^{(q-1)}_t:t<j\}+\ker\bigl(d^{q-1}|_{\Cq{q-1}}\bigr)$.
By Lemma~\ref{lemma:d_maps_C_eq} and Definition~\ref{def:eqcohomology} that kernel is
$\Zq{q-1}$, so the condition is
$c^{(q-1)}_j\notin\mathrm{span}\bigl(\Zq{q-1}\cup\{c^{(q-1)}_t:t<j\}\bigr)$, which is
exactly the pivot condition for column $j$ of the $c$-block of
$[\,z^{(q-1)}\,|\,c^{(q-1)}\,]$, i.e.\ the definition of $w^{(q-1)}$ in
Definition~\ref{def:storedbases}. Both families therefore consist of the same vectors in
the same order. They span $\Bq{q}=d^{q-1}(\Cq{q-1})$, and they are independent because
they are pivot columns.
\end{proof}

\noindent
This is the shipped route, and it is why $\Zq{q-1}$ need not be held while $\Bq{q}$ is
computed: the elimination is $N_qd_V\times m_{q-1}$ instead of
$N_{q-1}d_V\times(d_{Z,q-1}+m_{q-1})$. [v5]'s step $(q-1).4$ computed $W^{q-1}$ first
and applied $d^{q-1}$ to it; Proposition~\ref{prop:Bqimage} says the two agree
\emph{column for column}, which is what makes the change value-preserving.

\begin{proposition}[The stored projector]\label{prop:projector}
\PROVED{} Let $\mathfrak{h}_q$ be as in Definition~\ref{def:storedbases}. There is a
unique $\cring$-linear $P_q\colon\cring^{N_qd_V}\to\cring^{d_q}$ with
\[
  P_q\,z^{(q)}_{j}=e_{\iota(j)}\ (j\in\mathfrak{h}_q),\qquad
  P_q\,\Bq{q}=0,\qquad
  P_q\,z^{(q)}_{j}=\text{the coordinates of }[z^{(q)}_{j}]\ (j\notin\mathfrak{h}_q),
\]
$\iota$ the increasing enumeration of $\mathfrak{h}_q$, and $P_q$ vanishes on a fixed
complement of $\Zq{q}$. Equivalently, $P_q$ reads off coordinates in the stored basis of
$\Heq{q}$ modulo $\Bq{q}$. Its defining certificates are
\begin{equation}\label{eq:projcert}
  P_q\,H^{(q)}=\id_{d_q},\qquad P_q\,B^{(q)}=0,
\end{equation}
where $H^{(q)}$, $B^{(q)}$ are the matrices whose columns are the stored bases of
$\Heq{q}$ and $\Bq{q}$.
\end{proposition}

\begin{proof}
$\Zq{q}=\Bq{q}\oplus\mathrm{span}\{z^{(q)}_j:j\in\mathfrak h_q\}$ by the pivot property
of $\mathfrak h_q$ in $[\,b^{(q)}\,|\,z^{(q)}\,]$, so the three displayed requirements
prescribe $P_q$ on $\Zq{q}$ and on the chosen complement, i.e.\ on all of
$\cring^{N_qd_V}$; linearity then determines it uniquely. \eqref{eq:projcert} restates
the first two requirements.
\end{proof}

\begin{algorithm}[The cohomology of $(K,\Rel;L)$ in all degrees]\label{alg:cohomology}
\textbf{Input.} $(K,\Rel)$ satisfying \ref{ass:fund_domain} and
\ref{ass:consistent_face_id}; a matrix local system $L$ of rank $d_V$ over a field
$\cring$; the trivial-coefficient data $d^{q}_1$ and a basis of $Z^q_1$ for each $q$
(independent of $L$ and of $d_V$, so precomputed once).
\textbf{Steps.} For $q=0,\dots,\dim K$:
\begin{enumerate}
\item assemble $A_{C^q}$ (Proposition~\ref{prop:ACq}) and take
  $c^{(q)}_{\bullet}:=$ its kernel basis; $m_q=\dim\Cq{q}$;
\item assemble $A_{Z^q}$ (Proposition~\ref{prop:AZq}), take a kernel basis, and apply
  $E^{(q)}$: $z^{(q)}_{\bullet}$, $d_{Z,q}=\dim\Zq{q}$;
\item for $q\ge1$, $b^{(q)}_{\bullet}:=$ the image basis of $d^{q-1}C^{(q-1)}$
  (Proposition~\ref{prop:Bqimage}); $b^{(0)}_{\bullet}$ is empty;
\item $\mathfrak h_q$, the stored basis of $\Heq{q}$ and $P_q$ from the factor basis of
  the generators $z^{(q)}_{\bullet}$ modulo the relations $b^{(q)}_{\bullet}$
  (Proposition~\ref{prop:projector}); $d_q=\dim\Heq{q}=d_{Z,q}-d_{B,q}$.
\end{enumerate}
\textbf{Dependencies.} Steps 1--2 depend only on $q$; step 3 in degree $q$ reads only
$c^{(q-1)}_{\bullet}$; step 4 reads only $z^{(q)}_{\bullet}$ and $b^{(q)}_{\bullet}$. So
the degrees of steps 1--2 are mutually independent, and the chains
$3\to4$ are mutually independent across $q$ once step 1 has run --- which is how the
implementation schedules them. \textbf{No step depends on the order in which the degrees
are run}, so the stored objects are independent of that schedule; this is asserted, and
checked, as a determinism property of the implementation.
\end{algorithm}

\begin{remark}[Cost]\label{rmk:cost}
With $d_V=n+1$ the dominant costs are the two kernel extractions: $A_{C^q}$ is
$N_{\Rel,q}d_V\times N_qd_V$ and $A_{Z^q}$ is $N_{\Rel,q}d_V\times k_qd_V$, so both are
cubic in $d_V$ with a triangulation-dependent constant, and step 4 eliminates a
$(d_{Z,q}+d_{B,q})\times N_qd_V$ system. The $d_V$-independent data --- $d^{q}_1$ and
the basis of $Z^q_1$ --- are computed once per $(K,\Rel)$; that separation is the point
of Proposition~\ref{prop:AZq}, and it is the only structural optimisation this chapter
asserts. The measured behaviour of the shipped implementation is not part of this
contract.
\end{remark}

\paragraph{Code map.}
Proposition~\ref{prop:ACq}: \code{assemble\_A\_Cq}, and the vacuous-relation
short-circuit is \code{relation\_is\_vacuous} / \code{live\_relations}, gated by
\code{VacuousRelationValidation}. Proposition~\ref{prop:AZq}: \code{assemble\_A\_Zq}
and \code{assemble\_beta\_expansion} ($E^{(q)}$), with $Z^q_1$ from
\code{get\_trivial\_cocycle\_basis\_Z1\_elt}. Definition~\ref{def:storedbases}: the
stored \code{C\_eq\_block(q)}, \code{Z\_eq\_block(q)}, \code{B\_eq\_block(q)},
\code{H\_eq\_block(q)}, \code{P\_block(q)}, each a multi-vector whose column $t$ is one
cochain in the layout \eqref{eq:flatindex}. Proposition~\ref{prop:Bqimage}:
\code{image\_basis} applied to \code{cohom.DC}. Proposition~\ref{prop:projector}:
\code{factor\_basis}, whose \code{pivots} are $\mathfrak h_q$ and whose
\code{projector} is $P_q$; \eqref{eq:projcert} is checked in process by
\code{CohomologyBlockModelValidation}. The contract box's elimination rule is the
engine's documented basis-choice contract, pinned entrywise by
\code{BlockEliminationValidation} and \code{GoldenRegressionValidation}.

\paragraph{Formalization note.}
Everything in this chapter is finite linear algebra over a field, and
Propositions~\ref{prop:ACq}, \ref{prop:AZq}, \ref{prop:Bqimage} and
\ref{prop:projector} are statements about kernels, images and quotients of explicit
matrices --- the most directly formalizable material in the document, and the natural
place to state the stored bases as \emph{definitions by a deterministic elimination}
rather than as ``choose a basis''. The one thing a formalization must not skip is the
contract box: the basis is part of the specification, not an implementation detail, and
a proof that only fixes the spans would not pin $\Amat{q}$.

%======================================================================
\part*{Part II --- The model and its coefficients}
\addcontentsline{toc}{section}{Part II --- The model and its coefficients}
%======================================================================

%----------------------------------------------------------------------
\section{Conventions, the two primes, and the 7-vertex model}\label{sec:conventions}
%----------------------------------------------------------------------
% SOURCE:  hecke_action.tex:161-286, 373-383 (RH)
% NOTATION: rows Rel, cring, sigma, zeta, d, T_N of sec:notation-crosswalk applied
% ERRATA:  E8 (App. A)

\paragraph{What this chapter fixes.}
The two roles a prime plays in this document, and the standing hypotheses
\ref{H1}--\ref{H4} under which every later statement is read; the concrete arithmetic
datum --- the group $\Ga=\SL_2(\Z)$, its generators, the seven-vertex complex $K$ and
its relation set $\Rel$ --- on which every later chapter computes; the central $-I$
relation, which is a modelling \emph{requirement} and not a convenience; and the
cochain coordinates in which every matrix of this document is written. Nothing here is
about the Hecke action: this chapter fixes the object the Hecke action will act on.

\paragraph{Standing hypotheses.}
\ref{H1}--\ref{H4} below are declared once for Parts~II--V and, per rule~\ref{A1}, are
\emph{repeated in every statement that uses one}. No statement of this document may be
read as holding for all $\ell$ merely because \ref{H1} was declared in this chapter.

\subsection*{The two primes $\ell$ and $p$}

Two distinct roles of primes occur throughout and must never be conflated. The symbols
are fixed in \S\ref{sec:notation-primes}; the box below is the operational form of that
fixing, and it is the one the code base implements.

\begin{contractbox}
$\ell$ is the \emph{residue characteristic of the coefficients}: the local system has
fibre $\VR$ over a coefficient ring $\cring$ of characteristic $\ell$, and
$\dim_{\Fl}\Mn=n+1$ for $\cring=\Fl$. $p$ is the \emph{Hecke prime}: $T_p$ is attached
to the double coset $\Ga\,\mathrm{diag}(1,p)\,\Ga$ (\S\ref{sec:heckedatum}). New code
uses the identifiers \code{ell} (coefficients) and \code{p\_hecke} (Hecke prime); the
legacy identifier \code{p} of the arithmetic layer ---
\code{LocalSystemSimplicialComplex::prime}, the \code{p\_mod} arguments --- is to be
read as $\ell$, never as the Hecke prime.
\end{contractbox}

\noindent
The composite-index operator is written $T_N$ with $N\in\Z_{>0}$, never $T_n$: $n$ is
the symmetric-power degree and nothing else (\S\ref{sec:notation-primes};
Appendix~\ref{app:errata}, entry~\ref{err:e16}).

\subsection*{The standing hypotheses (H1)--(H4)}

\begin{enumerate}[label=(H\arabic*)]
\item\label{H1} \textbf{Good case, $p\neq\ell$.} The generic situation. Every statement
  of \S\S\ref{sec:heckedatum}--\ref{sec:example} applies without restriction.
\item\label{H2} \textbf{Bad case, $p=\ell$.} The operator defined for $p=\ell$ is
  $U_\ell$. Its \emph{definition} is identical --- the construction of
  \S\S\ref{sec:chainlevel}--\ref{sec:induced} never inverts an element of determinant
  $p$ --- but $\rhon(g)$ for $\det g=\ell$ is \emph{singular} modulo $\ell$ for
  $n\ge1$ (Lemma~\ref{lem:detsym}), and the Hecke recursion degenerates:
  $T_{\ell^{\nu}}=U_\ell^{\,\nu}$ over $\cring=\Fl$ (Corollary~\ref{cor:ulrecursion}).
\item\label{H3} \textbf{Torsion-of-stabilisers caveat, $\ell\in\{2,3\}$.} The
  stabilisers in $\Ga=\SL_2(\Z)$ of the elliptic points $\zeta^2,\zeta,i$ have orders
  $6,6,4$, and every simplex stabiliser contains the centre $\{\pm I\}$; the orders
  occurring therefore divide $12$. If $6\in\cring^{\times}$ --- equivalently
  $\ell\ge5$, for $\cring$ of characteristic $\ell$ --- the complex $\Ceq(K,\Rel;L)$
  computes $H^{\bullet}(\Ga;\VR)$ (Remark~\ref{rmk:orbifold}). If $\ell\in\{2,3\}$ it
  computes the cohomology of the explicit complex itself, which we call the
  \textbf{orbifold model cohomology} of $(K,\Rel;L)$; \textbf{this may differ from
  $H^{\bullet}(\Ga;\VR)$, and no isomorphism between the two is asserted there, in
  either direction.} (Since v1.8.0 the orbifold model cohomology is identified, over the
  ring in force, as a Bredon cohomology independent of the triangulation
  (Proposition~\ref{prop:bredon}), and one integral comparison map to group cohomology is
  constructed, $\cmpI^{\ast}$ of Proposition~\ref{prop:isomodel}(4); it is a
  quasi-isomorphism only when $6\in\cring^{\times}$. Until v1.7.4 this hypothesis said that
  no comparison \emph{map} was asserted.) All \emph{constructions} and all \emph{well-definedness}
  statements of this document hold for every $\ell$; the \emph{comparison} statements
  with $H^{\bullet}(\Ga;\VR)$, the Hecke-algebra relations of
  Theorem~\ref{thm:relations} and the eigenvalue oracles of
  Obligation~\ref{obl:oracle} carry the hypothesis $6\in\cring^{\times}$ in their own
  statements.
\item\label{H4} \textbf{Tessellation hypothesis.} Fix a truncation height $Y>1$ and let
  $\Xbulk$ be the $\Ga$-tessellation of the truncated upper half plane: the
  $\Ga$-invariant simplicial complex whose underlying space is $\Hyp$ with the
  $\Ga$-translates of the open horoball $\{\mathrm{Im}\,z>Y\}$ removed, triangulated by
  the $\Ga$-translates of the five $2$-simplices of Definition~\ref{def:model7}. We use,
  and therefore state, all four of:
  \begin{enumerate}[label=(\roman*)]
  \item $\Ga$ acts simplicially on $\Xbulk$, without inversions, with \emph{finite}
    stabilisers, and $\Ku\cong\Xbulk$ as $\Ga$-complexes;
  \item \textbf{$\Xbulk$ is contractible and $2$-dimensional};
  \item $K$ is a strict fundamental domain for that action, and $(K,\Rel)$ presents the
    quotient in the sense of \S\ref{sec:framework}
    (Assumption~\ref{ass:fund_domain});
  \item the stabiliser of every simplex of $\Xbulk$ contains $-I$; the stabilisers of
    the elliptic vertices are $\langle ST\rangle$, $\langle TS\rangle$,
    $\langle S\rangle$, of orders $6,6,4$, at $v_4,v_5,v_6$ respectively; every other
    simplex stabiliser is exactly $\{\pm I\}$.
  \end{enumerate}
\end{enumerate}

\noindent
\textbf{Why \ref{H4}(ii) is written down.} It is load-bearing three chapters later ---
the uniqueness of $\Lam$ and the existence of the fills in \S\ref{sec:cone} both use it
--- and in the superseded corpus it was \emph{never written}, so a reader could not tell
the argument from a wish. Rule~\ref{A1} exists for exactly this statement.

\subsection{The arithmetic group datum and the 7-vertex model}\label{sec:model}

$\Ga:=\SL_2(\Z)$ acts on the upper half plane $\Hyp$ by $z\mapsto\frac{az+b}{cz+d}$.
The generators used by the model are
\[
T:=\smat{1}{1}{0}{1},\qquad
S:=\smat{0}{-1}{1}{0},\qquad
ST=\smat{0}{-1}{1}{1},\qquad
TS=\smat{1}{-1}{1}{0},
\]
with $S^2=(ST)^3=(TS)^3=-I$ and $S^4=(ST)^6=(TS)^6=I$. Write
$\mathcal{G}:=\{T,S,ST,TS\}$ for this generator multiset. The elliptic points are
$\zeta:=e^{\pi i/3}$, $\zeta^2$ and $i$; the letter $\rho$ is reserved
for the representation (\S\ref{sec:notation-collisions}, C9).

\begin{definition}[The seven-vertex model]\label{def:model7}
\LEAN{def:model7-lean} (the simplex list: the seven vertices, the five triangles, the eleven edges, increasing vertex order, lexicographic within each dimension)
$K$ is the finite $2$-dimensional simplicial complex with
\begin{itemize}
\item \textbf{vertices} $v_0,\dots,v_6$: $v_0,v_1,v_2,v_3$ on the horocycle
  $\{\mathrm{Im}\,z=Y\}$, with $\mathrm{Re}$ running from $-\tfrac12$ at $v_0$ to
  $+\tfrac12$ at $v_3$; $v_4=\zeta^2$ (on $\mathrm{Re}\,z=-\tfrac12$), $v_5=\zeta$ (on
  $\mathrm{Re}\,z=+\tfrac12$), $v_6=i$;
\item \textbf{$2$-simplices} $\{v_0v_1v_4\}$, $\{v_1v_2v_6\}$, $\{v_2v_3v_5\}$,
  $\{v_1v_4v_6\}$, $\{v_2v_5v_6\}$;
\item \textbf{$1$-simplices} the eleven edges of those five triangles: $v_0v_1$,
  $v_0v_4$, $v_1v_2$, $v_1v_4$, $v_1v_6$, $v_2v_3$, $v_2v_5$, $v_2v_6$, $v_3v_5$,
  $v_4v_6$, $v_5v_6$.
\end{itemize}
So $N_0=7$, $N_1=11$, $N_2=5$ and $\chi(K)=1$. Every simplex is stored with increasing
vertex index, and the stored order within each dimension is lexicographic; that order is
part of the datum, because every matrix of this document is written in it. Code:
\code{CreateSL2ZFundamentalDomainTriangulation} (\file{LocalSystemSL2Z.cpp}), which
derives the edge list from the face list exactly as above.
\end{definition}

\begin{definition}[The relation set $\Rel$ of the model]\label{def:model7rel}
\LEAN{prop:model7coincidences-2} (the eight displayed rows: each $\gamma$ lies in the transporter of its row and carries the stored orientation to the stored orientation)
$\Rel$ consists of the triples $(\sigma_{\mathrm{lhs}},\sigma_{\mathrm{rhs}},\gamma)$
below, each read as the equivariance constraint
$f(\sigma_{\mathrm{lhs}})=\rhon(\gamma)\,f(\sigma_{\mathrm{rhs}})$ --- the
\code{SimplexRelation} format $f(s_{\mathrm{lhs}})=M\,f(s_{\mathrm{rhs}})$ with
$M=\rhon(\gamma)$:
\begin{center}
\begin{tabular}{@{}lllll@{}}
\toprule
$q$ & $\sigma_{\mathrm{lhs}}$ & $\sigma_{\mathrm{rhs}}$ & $\gamma$ & code tag \\
\midrule
$0$ & $\{v_3\}$      & $\{v_0\}$      & $T$  & \code{T\_v3\_v0} \\
$0$ & $\{v_5\}$      & $\{v_4\}$      & $T$  & \code{T\_v5\_v4} \\
$1$ & $\{v_3,v_5\}$  & $\{v_0,v_4\}$  & $T$  & \code{T\_e(0,4)\_e(3,5)} \\
$0$ & $\{v_4\}$      & $\{v_5\}$      & $S$  & \code{S\_v4\_v5} \\
$0$ & $\{v_6\}$      & $\{v_6\}$      & $S$  & \code{S\_v6\_stabilizer} \\
$1$ & $\{v_4,v_6\}$  & $\{v_5,v_6\}$  & $S$  & \code{S\_e(5,6)\_e(4,6)} \\
$0$ & $\{v_4\}$      & $\{v_4\}$      & $ST$ & \code{ST\_v4\_stabilizer} \\
$0$ & $\{v_5\}$      & $\{v_5\}$      & $TS$ & \code{TS\_v5\_stabilizer} \\
\midrule
\multicolumn{5}{@{}p{0.80\textwidth}@{}}{and, for \emph{every} simplex $\sigma$ of $K$
in every dimension, the central relation $(\sigma,\sigma,-I)$ --- Remark~\ref{rmk:negI},
code tags \code{negI\_q0\_0}, \code{negI\_q1\_0}, \dots} \\
\bottomrule
\end{tabular}
\end{center}
The eight listed rows plus the $7+11+5=23$ central relations give $\abs{\Rel}=31$.
Code: \code{SL2ZRelationSeeds\_7VertexModel} and
\code{GenerateSL2ZRelationsWithGroupData} (\file{LocalSystemSL2Z.cpp}), in this order.
\end{definition}

\begin{contractbox}
The list of \code{SimplexRelation}s is the single source of truth for the pairs
$(\text{relation},\,\gamma)$. Any relation added to the model carries its $\gamma\in\Ga$
as exact integer $2\times2$ data; its local-system matrix is then \emph{derived} as
$\rhon(\gamma)$ (Definition~\ref{def:rhon}), never entered by hand. The superseded
dictionary
``$r\in R_+\leftrightarrow g_r$'' of [v5], Prop.~24 is the fourth column of the table
above; individual relations carry no single-letter name here, because $r$ is the
restriction cochain map (\S\ref{sec:notation-framework}).
\end{contractbox}

\begin{lemma}[The coincidences of the seven-vertex model]\label{lem:model7coincidences}
\PROVED{} Let $K$ be as in Definition~\ref{def:model7}, let $Y>1$, let $\Xbulk$ be as in
\ref{H4}, and let $\lift{(-)}$ be the inclusion of $K$ into $\Xbulk$ as the truncated
fundamental domain
$D_Y:=\{z\in\Hyp:\abs{z}\ge1,\ \abs{\mathrm{Re}\,z}\le\tfrac12,\ \mathrm{Im}\,z\le Y\}$.
Then:
\begin{enumerate}
\item \textbf{Stabilisers.} \LEAN{prop:model7h4-iv} $\mathrm{Stab}_{\Ga}(\lift\sigma)=\{\pm I\}$ for every
  simplex $\sigma$ of $K$ other than the three elliptic vertices, where
  \[
    \mathrm{Stab}_{\Ga}(\lift{v_4})=\langle ST\rangle,\quad
    \mathrm{Stab}_{\Ga}(\lift{v_5})=\langle TS\rangle,\quad
    \mathrm{Stab}_{\Ga}(\lift{v_6})=\langle S\rangle,
  \]
  of orders $6$, $6$, $4$. In particular every simplex stabiliser is finite, contains
  $-I$, and has order dividing $12$; and no $\gamma\in\Ga$ permutes the vertices of a
  simplex of $K$ non-trivially, so $\Ga$ acts without inversions.
\item \textbf{Coincidences between distinct simplices.} \LEAN{prop:model7coincidences-2} For
  $\sigma\neq\sigma'$ in $\orient{\bullet}(K)$ the transporter
  $\mathrm{Tr}(\sigma,\sigma')$ is non-empty for exactly four unordered pairs
  $\{\sigma,\sigma'\}$:
  \[
    \{v_0,v_3\},\qquad \{v_4,v_5\},\qquad
    \bigl\{\,\{v_0,v_4\},\ \{v_3,v_5\}\,\bigr\},\qquad
    \bigl\{\,\{v_4,v_6\},\ \{v_5,v_6\}\,\bigr\},
  \]
  realised by $T\in\mathrm{Tr}(v_3,v_0)$, $T\in\mathrm{Tr}(v_5,v_4)$,
  $T\in\mathrm{Tr}(\{v_3,v_5\},\{v_0,v_4\})$ and
  $S\in\mathrm{Tr}(\{v_4,v_6\},\{v_5,v_6\})$, each carrying the stored (increasing
  vertex order) orientation to the stored orientation. All five $2$-simplices, the three
  vertices $v_1,v_2,v_6$ and the seven edges other than the four listed are alone in
  their $\Ga$-orbit.
\end{enumerate}
\end{lemma}

\begin{proof}
Write $D:=\{z\in\Hyp:\abs{z}\ge1,\ \abs{\mathrm{Re}\,z}\le\tfrac12\}$, so
$\abs{K}=D_Y=D\cap\{\mathrm{Im}\,z\le Y\}$. Take as given the classical description of
the $\Ga$-action on $\Hyp$ \STD{\cite[Ch.~VII, \S1.2, Thms.~1--2]{serre_cours}}, in the
two parts used here:
\begin{enumerate}[label=(\Alph*)]
\item\label{serreA} \emph{Identifications.} If $z,z'\in D$ with $z\neq z'$ and
  $z=\gamma z'$ for some $\gamma\in\Ga$, then either $\mathrm{Re}\,z'=\mp\tfrac12$ and
  $z=z'\pm1$, with $\gamma=\pm T^{\pm1}$; or $\abs{z'}=1$ and $z=-1/z'$, with
  $\gamma=\pm S$.
\item\label{serreB} \emph{Stabilisers of points.} For $z\in D$, $\mathrm{Stab}_{\Ga}(z)$
  is $\{\pm I\}$ except at the three points $z=i$, $z=\zeta$, $z=\zeta^2$, where it is
  $\langle S\rangle$ of order $4$, $\langle TS\rangle$ of order $6$ and
  $\langle ST\rangle$ of order $6$ respectively. ($S$ fixes $i$; $ST:z\mapsto-1/(z+1)$
  fixes the roots of $z^2+z+1$, so $\zeta^2$; $TS:z\mapsto(z-1)/z$ fixes the roots of
  $z^2-z+1$, so $\zeta$. Serre's statement is in $\mathrm{PSL}_2(\Z)$, where the three
  stabilisers have orders $2,3,3$; the orders above are those of their preimages in
  $\Ga=\SL_2(\Z)$, the central $-I$ acting trivially on $\Hyp$.)
\end{enumerate}

\emph{Part (2).} Let $\sigma\neq\sigma'$ in $\orient{\bullet}(K)$ and
$\gamma\in\mathrm{Tr}(\sigma,\sigma')$, so $\gamma$ maps $\abs{\sigma'}$ onto
$\abs{\sigma}$ and hence the relative interior of $\abs{\sigma'}$ onto that of
$\abs{\sigma}$. Distinct simplices of a simplicial complex have disjoint relative
interiors, so $\gamma z'\neq z'$ for every $z'$ in the relative interior of
$\abs{\sigma'}$, while both $z'$ and $\gamma z'$ lie in $D$. By \ref{serreA} that
relative interior is contained in $L_-\cup L_+\cup C$, where
$L_{\mp}:=\{\mathrm{Re}\,z=\mp\tfrac12\}$ and $C:=\{\abs{z}=1\}$ --- three geodesics of
$\Hyp$. A $2$-simplex has $2$-dimensional relative interior and is therefore excluded, so
$\dim\sigma'\le1$; and if $\dim\sigma'=1$ its relative interior is an infinite subset of
a geodesic segment, so it meets one of the three geodesics in at least two points and
therefore lies in that one, two distinct geodesics of $\Hyp$ meeting in at most one
point. Hence $\abs{\sigma'}$, and likewise $\abs{\sigma}$, lies in $L_-$, in $L_+$ or in
$C$.

By Definition~\ref{def:model7} the simplices of $K$ so contained are exactly:
in $L_-\cap D_Y$ (the vertical segment from $v_4=\zeta^2$ to $v_0=-\tfrac12+iY$) the
simplices $\{v_0\}$, $\{v_4\}$, $\{v_0,v_4\}$; in $L_+\cap D_Y$ (from $v_5=\zeta$ to
$v_3=\tfrac12+iY$) the simplices $\{v_3\}$, $\{v_5\}$, $\{v_3,v_5\}$; and in $C\cap D_Y$
(the unit arc from $\zeta^2$ through $i$ to $\zeta$) the simplices $\{v_4\}$, $\{v_5\}$,
$\{v_6\}$, $\{v_4,v_6\}$, $\{v_5,v_6\}$ --- there are no other vertices of $K$ on the two
walls or on the arc, and no other edges of $K$ joining two of these vertices along them.
Now \ref{serreA} pins $\gamma$ and hence $\sigma$: if $\abs{\sigma'}\subseteq L_-$ then
$\gamma=\pm T$ and $\sigma=T\sigma'$, giving $v_0\mapsto v_3$, $v_4\mapsto v_5$,
$\{v_0,v_4\}\mapsto\{v_3,v_5\}$; if $\abs{\sigma'}\subseteq L_+$ then $\gamma=\pm T^{-1}$
and $\sigma=T^{-1}\sigma'$, the three inverse assignments; if $\abs{\sigma'}\subseteq C$
then $\gamma=\pm S$ and $\sigma=S\sigma'$, giving $v_4\leftrightarrow v_5$,
$\{v_4,v_6\}\leftrightarrow\{v_5,v_6\}$ and $v_6\mapsto v_6$, the last excluded by
$\sigma\neq\sigma'$. Collecting the unordered pairs gives exactly the four listed, with
the stated transporter elements. Each of $T$ and $S$ maps the smaller vertex of the
source edge to the smaller vertex of the target edge ($0\mapsto3$, $4\mapsto5$ for $T$;
$5\mapsto4$, $6\mapsto6$ for $S$), so the stored orientations correspond.

\emph{Part (1).} For a vertex $v$ of $K$ the claim is \ref{serreB}: $v_4=\zeta^2$,
$v_5=\zeta$, $v_6=i$ are the three exceptional points, and $v_0,\dots,v_3$ lie on
$\mathrm{Im}\,z=Y>1$, so are none of them. Let now $\dim\sigma\ge1$ and
$\gamma\in\mathrm{Stab}_{\Ga}(\lift\sigma)$; then $\gamma$ permutes the vertices of
$\sigma$. If that permutation were non-trivial it would move some vertex $v$ of $\sigma$
to another vertex $w$ of $\sigma$, so $\{v,w\}$ would be one of the two vertex pairs of
part~(2), namely $\{v_0,v_3\}$ or $\{v_4,v_5\}$ --- but neither is an edge of $K$
(Definition~\ref{def:model7}), so neither pair is contained in a common simplex. Hence
$\gamma$ fixes every vertex of $\sigma$, i.e.\ at least two distinct points of $\Hyp$;
a non-identity element of $\mathrm{PSL}_2(\Rl)$ fixes at most one point of $\Hyp$, so
$\gamma=\pm I$. This also proves the last sentence of part~(1): no $\gamma$ permutes the
vertices of a simplex of $K$ non-trivially, and since every simplex of $\Xbulk$ is a
$\Ga$-translate of a simplex of $K$ (that is how $\Xbulk$ is defined in \ref{H4}), the
same holds on $\Xbulk$: $\Ga$ acts without inversions, and
$\mathrm{Stab}_{\Ga}(\gamma\lift\sigma)=\gamma\,\mathrm{Stab}_{\Ga}(\lift\sigma)\gamma^{-1}$
is finite of order dividing $12$.
\end{proof}

\begin{proposition}[The model satisfies \ref{H4}]\label{prop:model7h4}
\PROVED{} \LEAN{prop:model7h4-i} (\ref{H4}(i) without the identification $\Ku\cong\Xbulk$, which is the trust boundary) \LEAN{prop:model7h4-ii} (\ref{H4}(ii) in the form this document uses it, since v1.19.0: the acyclicity of $C_\bullet(\Xbulk)$ over $\Z$ in the three forms of \eqref{eq:contractibleuse} --- $H_0(\Xbulk)\cong\Z$ by the augmentation, $Z_1=B_1$, $Z_2=0$ --- and $\dim\Xbulk=2$ ([GT] row \texttt{prop:model7h4-ii-dim}); proved combinatorially in [GT] (slot-19 L19-3, \texttt{model7\_acyclic}), not by the horoball retraction; contractibility as a topological statement is not formalised) \LEAN{prop:model7h4-iii} (\ref{H4}(iii): \ref{FD1} and \ref{FD2} only) \LEAN{prop:model7h4-iv} (\ref{H4}(iv) in full) Let $K$ and $\Rel$ be as in Definitions~\ref{def:model7} and
\ref{def:model7rel}, let $Y>1$, and let $\Xbulk$ be as in \ref{H4}. Then
\ref{H4}(i)--(iv) hold. The ineffective kernel of Definition~\ref{def:kernelN} is
$N=\{\pm I\}$ here, so \ref{FD1} and \ref{FD3} of \ref{H4}(iii) are the conditions
modulo $\{\pm I\}$ (Remark~\ref{rmk:fd13centre}).
\end{proposition}

\begin{proof}
\emph{(i) and (iv).} $\Ga$ permutes the simplices of $\Xbulk$ because $\Xbulk$ is by
definition the $\Ga$-translate tessellation of the five $2$-simplices of
Definition~\ref{def:model7} and their faces, and $\Ku\cong\Xbulk$ as $\Ga$-complexes is
the identification fixed in Remark~\ref{rmk:twosettings} (second bullet) --- it is not a
claim proved here but the instantiation of the framework of \S\ref{sec:framework}. That
the action is \emph{without inversions}, that every stabiliser is \emph{finite}, and the
whole of (iv) --- the stabiliser inventory, and $-I$ in every stabiliser --- are
Lemma~\ref{lem:model7coincidences}(1).

\emph{(ii).} $\dim K=2$ by Definition~\ref{def:model7}, and $\Xbulk$ is covered by
$\Ga$-translates of $\abs{K}$, so $\dim\Xbulk=2$. For contractibility: since $Y>1$, the
open horoballs $\gamma^{-1}\{\mathrm{Im}\,z>Y\}$, $\gamma\in\Ga$, are pairwise disjoint
or equal; the map $z\mapsto\mathrm{Re}\,z+i\min(\mathrm{Im}\,z,Y)$ is a deformation
retraction of $\{\mathrm{Im}\,z>Y\}$ onto the horocycle $\{\mathrm{Im}\,z=Y\}$, fixing
that horocycle pointwise; transporting it by $\gamma^{-1}$ on each horoball and taking
the identity elsewhere is well defined by disjointness, is $\Ga$-equivariant, and is a
deformation retraction of $\Hyp$ onto $\Xbulk$. As $\Hyp$ is contractible, so is
$\Xbulk$.

\emph{(iii): \ref{FD2}.} $\Ga\cdot K=\Xbulk$ holds by the definition of $\Xbulk$ just
recalled.

\emph{(iii): \ref{FD4}, the substantive clause.} By Lemma~\ref{lem:fd4criterion} it
suffices to verify its two conditions, and for each of them the sufficient criterion
given there. The stabilisers and the coincidences they refer to are supplied by
Lemma~\ref{lem:model7coincidences}; neither is read off a table of $\Rel$.

Condition~\ref{FD4a} asks that, at each of the $7+11+5=23$ simplices $\sigma$ of $K$, the relation
elements of the \emph{self}-relations of $\Rel$ at $\sigma$ generate
$\mathrm{Stab}_{\Ga}(\lift\sigma)$. Reading off Definition~\ref{def:model7rel} against
Lemma~\ref{lem:model7coincidences}(1):
\begin{itemize}
\item $\sigma=\{v_4\}$: the row $(\{v_4\},\{v_4\},ST)$ generates
  $\langle ST\rangle=\mathrm{Stab}_{\Ga}(\lift{v_4})$;
\item $\sigma=\{v_5\}$: the row $(\{v_5\},\{v_5\},TS)$ generates $\langle TS\rangle$;
\item $\sigma=\{v_6\}$: the row $(\{v_6\},\{v_6\},S)$ generates $\langle S\rangle$;
\item each of the remaining $20$ simplices --- $v_0,v_1,v_2,v_3$, all $11$ edges and all
  $5$ triangles --- has stabiliser $\{\pm I\}$, generated by the relation element of its
  central row $(\sigma,\sigma,-I)$.
\end{itemize}
\textbf{This is the entire content of the $23$ central rows, and it is why omitting them
is not an option.} Apart from them $\Rel$ has no self-relation at any of those twenty
simplices, so the subgroup reached at each would be trivial while the stabiliser is
$\{\pm I\}$: \ref{FD4a} would fail, \ref{FD4} with it, and with \ref{FD4} the
surjectivity of $\Psi$ in Theorem~\ref{thm:explicit_complex_equiv}. That is
Remark~\ref{rmk:negI} and Appendix~\ref{app:errata}, entry~\ref{err:e8}, seen as the
failure of a stated criterion
rather than as a discrepancy in a computed dimension.

Condition~\ref{FD4b} asks that each pair of distinct simplices with non-empty transporter be joined
by members of $\Rel$. By Lemma~\ref{lem:model7coincidences}(2) there are exactly four
such pairs, and $\Rel$ contains one row for each: \code{T\_v3\_v0} for $\{v_0,v_3\}$,
\code{T\_v5\_v4} for $\{v_4,v_5\}$, \code{T\_e(0,4)\_e(3,5)} for
$\{\{v_0,v_4\},\{v_3,v_5\}\}$ and \code{S\_e(5,6)\_e(4,6)} for
$\{\{v_4,v_6\},\{v_5,v_6\}\}$. This exhausts the four pairs, so \ref{FD4b} holds.

Seven of the eight listed rows have now been used --- three for \ref{FD4a} at the
elliptic vertices, four for \ref{FD4b}. The eighth, $(\{v_4\},\{v_5\},S)$
(\code{S\_v4\_v5}), is \emph{redundant} for \ref{FD4}: inverting \code{T\_v5\_v4} gives
$(\{v_4\},\{v_5\},T^{-1})$, composing it with the self-relation
$(\{v_5\},\{v_5\},TS)$ gives $(\{v_4\},\{v_5\},T^{-1}\cdot TS)=(\{v_4\},\{v_5\},S)$, and
both factors lie in $\langle\Rel\rangle$. It is retained because it is the $S$-wall's own
identification and the $\tau$-table path templates of \S\ref{sec:reduction} consume it.

\emph{(iii): $N$, \ref{FD1} and \ref{FD3}.} First $N=\{\pm I\}$: $-I$ acts trivially on
$\Hyp$, hence on $\Xbulk$; conversely $N\subseteq\mathrm{Stab}_{\Ga}(\lift{v_1})=\{\pm I\}$
by Lemma~\ref{lem:model7coincidences}(1).

By Assumption~\ref{ass:fund_domain}, $K_{\mathrm{cl}}$ is generated by the simplices a
coincidence moves, which by Lemma~\ref{lem:model7coincidences}(2) are the eight simplices
$v_0,v_3,v_4,v_5$, $\{v_0,v_4\},\{v_3,v_5\},\{v_4,v_6\},\{v_5,v_6\}$; the subcomplex they
generate is the two vertical walls and the unit arc together with their vertices
($v_6$ included, as a face). Hence
$\widetilde D=D_Y\setminus(L_-\cup L_+\cup C)$ --- the notation $L_{\mp}$, $C$ for the two
vertical walls and the unit circle being that of the proof just given --- is the interior
of $D_Y$ together with the open horocycle segment from $v_0$ to $v_3$: non-empty,
connected and dense in $\abs{K}=D_Y$, which with the connectedness of $D_Y$ is
\ref{FD1}. Note that the three elliptic vertices are \emph{not} put into
$K_{\mathrm{cl}}$ by their stabilisers --- $v_4,v_5$ are there because $T$ moves them and
$v_6$ only as a face --- which is the point of Remark~\ref{rmk:fd13centre}.

For \ref{FD3}, let $\gamma\in\Ga$ with $\gamma\widetilde D\cap\widetilde D\neq\emptyset$,
say $z,\gamma z\in\widetilde D$. Then $z$ lies on none of $L_-$, $L_+$, $C$, so
\ref{serreA} of Lemma~\ref{lem:model7coincidences}'s proof forces $\gamma z=z$, and
\ref{serreB} forces $\gamma\in\{\pm I\}=N$. So $\gamma\widetilde D\cap\widetilde
D=\emptyset$ for every $\gamma\in\Ga\setminus N$, which is \ref{FD3}; the decomposition
\eqref{eq:fddecomp} of $\abs{\Xbulk}$ and the statement that all identifications occur on
$\abs{K_{\mathrm{cl}}}$ then follow from Remark~\ref{rmk:fdconsequences}.
\end{proof}

\begin{remark}[\ref{FD4} for the model is re-verified as a finite computation]
\label{rmk:fd4computed}
\COMP{} The two conditions of Lemma~\ref{lem:fd4criterion} are precisely what
\code{SL2ZTessellationTables} (\file{Hecke/SL2ZTessellationTables.cpp}) computes, so the
discharge above is re-run on every construction of that class and by the validation named
below; it does not rest on the printed argument alone. The class seeds the transporter
sets $\mathrm{Tr}(\sigma,\sigma')$ with exactly a generating set of $\langle\Rel\rangle$
--- $\{\pm I\}$ at every simplex (the $23$ central rows), the cyclic groups
$\langle ST\rangle$, $\langle TS\rangle$, $\langle S\rangle$ at $v_4,v_5,v_6$ (the three
stabiliser rows), and one element per row of \code{SL2ZRelationSeeds\_7VertexModel} (the
eight listed rows) --- and closes them under composition, inverse and multiplication by
$-I$ up to a fixed point, which it verifies is one. Its constructor-time
\code{validate\_closure}, and independently \code{HeckeTessellationValidation} against a
hand-derived table committed before the implementation was read, then assert:
\begin{itemize}
\item self-transporter orders $2,2,2,2,6,6,4$ at $v_0,\dots,v_6$ and $2$ at each of the
  $11$ edges and $5$ triangles. Since
  $\{\gamma:(\sigma,\sigma,\gamma)\in\langle\Rel\rangle\}$ is a subgroup of
  $\mathrm{Stab}_{\Ga}(\lift\sigma)$ and the orders match those of
  Lemma~\ref{lem:model7coincidences}(1), the two groups are equal: this is
  \ref{FD4a}, all $23$ simplices of it.
\item the orbits of the generated groupoid --- $\{v_0,v_3\}$, $\{v_1\}$, $\{v_2\}$,
  $\{v_4,v_5\}$, $\{v_6\}$ on vertices; $\{\{v_0,v_4\},\{v_3,v_5\}\}$ and
  $\{\{v_4,v_6\},\{v_5,v_6\}\}$ on edges with the other seven singletons; all five
  triangles singletons --- together with
  $\mathrm{Tr}(v_3,v_0)=\{\pm T\}$,
  $\mathrm{Tr}(\{v_3,v_5\},\{v_0,v_4\})=\{\pm T\}$ and
  $\mathrm{Tr}(\{v_4,v_6\},\{v_5,v_6\})=\{\pm S\}$ by exact set equality. Matched against
  Lemma~\ref{lem:model7coincidences}(2), this is \ref{FD4b}.
\item that $\Rel$ has exactly $8+23=31$ rows: the eight listed ones first, in the order of
  Definition~\ref{def:model7rel} and with relation elements $T,T,T,S,S,S,ST,TS$, then one
  central row $(\sigma,\sigma,-I)$ per simplex of every dimension, each with
  $M_\varrho=\rhon(-I)$ recomputed independently.
\end{itemize}
What the computation does \emph{not} supply is Lemma~\ref{lem:model7coincidences} itself:
the \emph{true} stabilisers and the \emph{true} orbits, against which the generated ones
are compared. Those come from \cite[Ch.~VII, \S1.2]{serre_cours} through the printed
proof, and a hand-derived expected table is what
\code{HeckeTessellationValidation} encodes. Neither half alone discharges \ref{FD4}.
\end{remark}

\begin{remark}[The orbifold caveat]\label{rmk:orbifold}
Assume \ref{H4}. Besides the boundary identifications for $T$ and $S$, the model imposes
the \emph{stabiliser relations} $f(v_4)=\rhon(ST)f(v_4)$, $f(v_5)=\rhon(TS)f(v_5)$,
$f(v_6)=\rhon(S)f(v_6)$ (Definition~\ref{def:model7rel}) and --- by
Remark~\ref{rmk:negI} --- the central relation $f(\sigma)=\rhon(-I)f(\sigma)$ on every
simplex. Consequently
\[
  \Ceq(K,\Rel;L)\;\cong\;\Hom_{\Z[\Ga]}\bigl(C_{\bullet}(\Ku),\VR\bigr),
  \qquad \Ku\cong\Xbulk ,
\]
by Theorem~\ref{thm:explicit_complex_equiv}. By \ref{H4}(i),(iv) each chain group
$C_q(\Xbulk)$ is a direct sum of modules $\Z[\Ga/\Ga_\sigma]$ with
$\abs{\Ga_\sigma}\mid12$ \LEAN{lem:xbulk-orient} (the direct-sum decomposition, for every $\Xbulk$-descriptor, the orientation of each orbit fixed by \ref{H4}(i)); in particular $C_{\bullet}(\Xbulk)$ is \emph{not} free over
$\Z[\Ga]$. If moreover $6\in\cring^{\times}$, then $\abs{\Ga_\sigma}\in\cring^{\times}$
for every $\sigma$, each $\cring[\Ga/\Ga_\sigma]$ is projective relative to the trivial
subgroup, and the standard comparison arguments --- the fundamental lemma of relative
homological algebra, Shapiro's lemma, transfer --- apply verbatim, giving
$\Heq{\bullet}(K,\Rel;L)\cong H^{\bullet}(\Ga;\VR)$
\STD{\cite[\S3.6]{benson} for relative projectivity and transfer;
\cite[Ch.~III.6]{brown} for Shapiro's lemma; \cite[Ch.~VII, \S7]{brown} for the isotropy
spectral sequence of a group acting on a complex, whose rows $E_1^{\ast q}$, $q>0$,
vanish precisely because $\abs{\Ga_\sigma}\in\cring^{\times}$}. \LEAN{thm:xbulk-repaug} (the ingredient that
$C_{\bullet}(\Xbulk)\otimes\cring\to\cring$ is a resolution, for every acyclic $\Xbulk$-descriptor with a vertex and every $\cring$; projective on the model
when $6\in\cring^{\times}$) \LEAN{thm:model7-rmkorbifold} (the isomorphism $\Heq{n}(K,\Rel;L)\cong H^{n}(\Ga;\VR)$ for every $n\ge0$ under
$6\in\cring^{\times}$, for every $\Xbulk$-descriptor with \ref{H4}(ii) (as acyclicity), a vertex, and stabiliser orders dividing $12$ (\ref{H4}(iv)),
in particular the model; $L$ a local system with representation $\rho_L$ as in Proposition~\ref{prop:matrix_rep_equiv}, read on $\Ga$ (in [GT] a $\Delta$-module,
$\Ga\subseteq\Delta\subseteq G$, through $\rho_L=\rho\circ e^{-1}$ for $e\colon\Ga\cong\SL_2(\Z)$, not only a representation of $\SL_2(\Z)$); the statement is tagged, not this
argument: [GT] replaces relative projectivity, Shapiro's lemma, transfer and the isotropy spectral sequence by mathlib's \texttt{groupCohomologyIso} on the
projective resolution; a $K$ outside the descriptor framework is not claimed) This is the precise
content of \ref{H3}.

For $\ell\in\{2,3\}$ that vanishing fails and the isomorphism is \emph{not} asserted:
the left-hand side computes the orbifold model cohomology of \ref{H3}, a well-defined
invariant of $(K,\Rel;L)$ and the object this implementation emits at those primes. Two
further points, neither of them settled here. First, $\piOne$ of the quotient
\emph{orbifold} is $\Ga$ only in the orbifold sense: the torsion may not be ignored, and
Remark~\ref{rmk:negI} is the place where ignoring it changes an answer
(Appendix~\ref{app:errata}, entry~\ref{err:e8}). Second, the systematic treatment of the
regime $6\notin\cring^{\times}$, by a free $\Ga$-model of which the present one is a
specialisation, is \S\ref{sec:orbifoldmodel}. The separate question of what happens at
$\ell=2$ \emph{on the boundary}, where $H^{\bullet}(\Gbdry;\VR)$ is unbounded above for
the boundary group $\Gbdry$ of \S\ref{sec:boundarydatum}
(Lemma~\ref{lem:ell2unbounded}), is construction item~\ref{open:ell2} (OC-1): the
mathematics is \S\ref{sec:ell2} --- including the bulk periodicity of
Lemma~\ref{lem:bulkperiodic} --- the \emph{policy} there was decided by the owner on
2026-09-16 (\ref{oc1a}), and neither is this chapter's.

\emph{A design direction, recorded here and not restated elsewhere in this document.}
The model of Definitions~\ref{def:model7} and \ref{def:model7rel} is fixed for the
regime $6\in\cring^{\times}$. The owner's stated expectation (2026-09-15;
\file{docs/implementation/hecke\_extensions\_owner\_decisions.md}) is that a correct
orbifold treatment at $\ell\in\{2,3\}$, and over $\Z$, will require a \emph{refinement}
of this triangulation --- a subdivision of the present complex, not a different model.
This is an expectation about how \S\ref{sec:orbifoldmodel} and \ref{open:ell2} will be
constructed, \textbf{not} a theorem and not a claim that the present model is
insufficient there; nothing stated in this document depends on it. Its force is on the
authors of \S\ref{sec:orbifoldmodel} and \S\ref{sec:boundarydatum}: the definitions
there must be phrased so as to admit a refinement of the $7$-vertex complex rather
than hard-wiring its seven vertices, eleven edges and five triangles.
\emph{Answered at v1.8.0} (owner ruling 6c, 2026-09-24): no refinement is needed. The
free model at $\ell\in\{2,3\}$ and over $\Z$ is built on the unrefined tessellation, the
missing degrees coming from the stabiliser resolutions (\S\ref{sec:integral},
Remark~\ref{rmk:cafive}); and what the present model computes at those primes is
identified there (Proposition~\ref{prop:bredon}).
\end{remark}

\begin{remark}[Modelling requirement: the central $-I$ relation on \emph{every} simplex]
\label{rmk:negI}
Assume \ref{H4}. Because $-I\in\Ga$ is central and acts trivially on $\Hyp$, it is the
identity deck transformation of $\Xbulk\cong\Ku$: it fixes the image in $\Xbulk$ of
\emph{every} oriented simplex $\sigma\in\orient{q}(K)$, in every degree $q$ --- this is
\ref{H4}(iv). The isomorphism of Theorem~\ref{thm:explicit_complex_equiv} therefore
forces the equivariance relation
\begin{equation}\label{eq:negI}
  f(\sigma)\;=\;\rhon(-I)\,f(\sigma)
  \qquad\text{for \emph{all} }\sigma\in\orient{q}(K),\ \text{all }q,
\end{equation}
and not merely at the elliptic fixed points $v_4,v_5,v_6$, where the stabiliser rows of
Definition~\ref{def:model7rel} already supply $-I=S^2=(ST)^3=(TS)^3$ implicitly.
Imposing \eqref{eq:negI} on every free (non-elliptic) simplex as well is a
\textbf{binding modelling requirement}, stated here because the eight listed rows do not
imply it. Its effect, for $\cring$ of characteristic $\ell$, is \PROVED{} as follows.
\begin{itemize}
\item \emph{$n$ even.} $\rhon(-I)=+I$ and \eqref{eq:negI} is vacuous; omitting it is a
  no-op. The implementation emits it anyway and short-circuits it
  (\code{VacuousRelationValidation}).
\item \emph{$n$ odd, $\ell\neq2$.} $\rhon(-I)=-I$, so \eqref{eq:negI} reads
  $2f(\sigma)=0$, i.e.\ $f(\sigma)=0$ on every simplex, and the complex vanishes:
  $\Heq{\bullet}(K,\Rel;L)=0$. \emph{Omitting} \eqref{eq:negI} on the free simplices
  leaves $H^{\bullet}$ spuriously non-zero --- concretely $\dim H^1=n+1$ instead of $0$
  --- whereas the correct value is $0$. This is correct also on the group side:
  $-I$ is central, hence acts
  trivially on $H^{\bullet}(\Ga;-)$, while it acts on the coefficients by
  $\rhon(-I)=(-1)^{n}=-1$; so $2H^{\bullet}(\Ga;\VR)=0$ and, $2$ being invertible,
  $H^{\bullet}(\Ga;\VR)=0$. This is the odd-$n$ oracle of
  Obligation~\ref{obl:oracle} and \S\ref{sec:example}.
\item \emph{$\ell=2$.} $\rhon(-I)=+I$ for every $n$, since $-1=1$ in $\cring$;
  \eqref{eq:negI} is again vacuous, and the odd-$n$ vanishing argument above does not
  apply --- it used $2\in\cring^{\times}$.
\end{itemize}
The implementation imposes \eqref{eq:negI} globally, as the \code{FIX-negI} block of
\code{GenerateSL2ZRelationsWithGroupData}.
\end{remark}

\subsection*{Cochain conventions and flat coordinates}

\emph{Orientations.} Every $q$-simplex of $K$ is stored with increasing vertex order,
and that stored orientation is the one all signs are taken relative to.

\emph{Coboundary.} The cochain differential is
\[
  (d^{q}f)(\sigma)\;=\;\sum_{t=0}^{q+1}(-1)^{t}\,f(\partial_t\sigma),
  \qquad d^{q}=d^{q}_{1}\otimes\id_{\VR},
\]
with $\partial_t$ the $t$-th face map. It is written $d$ for cochains throughout;
$\partial$ is the chain boundary and $\Dcone$ the cone differential
(\S\ref{sec:notation-cochains}). The superseded documents wrote $\delta$ for this map;
$\delta$ now denotes a coset element (\S\ref{sec:notation-collisions}, C7).

\emph{Flat coordinates.} A cochain $c\in C^q(K;\VR)$ is the column
$\bigl(c(\sigma_0)^{\top},\dots,c(\sigma_{N_q-1})^{\top}\bigr)^{\top}
\in\cring^{\,N_q d_V}$ in the stored simplex order, $d_V=n+1$: the flat index of the
component of $c(\sigma_s)$ along the basis vector $u_j$ of
Definition~\ref{def:rhon} is
\begin{equation}\label{eq:flatindex}
  s\cdot d_V + j,\qquad 0\le s<N_q,\quad 0\le j\le n .
\end{equation}
Every matrix in this document --- $A_{C^q}$, $A_{Z^q}$, $\Th^{(q)}_{\mathrm{flat}}$, the
stored bases $c^{(q)}_j,z^{(q)}_j,b^{(q)}_j$ and the projectors $P_q$ --- is written in
these coordinates. In the implementation \eqref{eq:flatindex} is
\code{engine::BlockMat::to\_flat} / \code{from\_flat} under the one-block-column
contract, which is its single source of truth, and
\code{H\_eq\_projectors\_from\_C\_eq\_flat} is $P_q$ in it. The stored bases and
projectors themselves are the subject of \S\ref{sec:systematiccomputation}.

\paragraph{Code map.}
\ref{H1}, \ref{H2}: \code{ell} / \code{p\_hecke} throughout \file{Hecke/}; the $p=\ell$
branch is \code{HeckeOperatorFamily}'s $U_\ell$. \ref{H3}: gated at every comparison
site, e.g.\ \file{Hecke/BarCocycleH1.h} (``the denominators $4,6$ must be invertible''),
\file{Hecke/HeckeRelationsValidation.h}, \file{ValidationRegistry.cpp}. \ref{H4}:
\file{Hecke/SL2ZTessellationTables.cpp}. Definition~\ref{def:model7}:
\code{CreateSL2ZFundamentalDomainTriangulation}. Definition~\ref{def:model7rel}:
\code{SL2ZRelationSeeds\_7VertexModel}, \code{GenerateSL2ZRelationsWithGroupData}.
Remark~\ref{rmk:negI}: the \code{FIX-negI} block of \file{LocalSystemSL2Z.cpp}, checked
by \code{VacuousRelationValidation} and by the odd-$n$ oracle of \S\ref{sec:example}.
Equation~\eqref{eq:flatindex}: \code{engine::BlockMat::to\_flat}, gated by
\code{CohomologyBlockModelValidation}.

\paragraph{Formalization note.}
Definitions~\ref{def:model7} and \ref{def:model7rel} are finite explicit data --- a list
of vertex lists per dimension, and a list of
$(\text{degree},\text{index},\text{index},\Mat_2(\Z))$ tuples --- so both are
\code{decidable} to state and to check. Proposition~\ref{prop:model7h4}(i),(iii),(iv) is
a finite check once the classical fundamental-domain theorem is available;
\ref{H4}(ii) is the one hypothesis of this chapter that needs a genuine topological
argument, and the proof given reduces it to contractibility of $\Hyp$ plus disjointness
of horoballs at $Y>1$. Remark~\ref{rmk:negI}'s odd-$n$ statement is two lines from
centrality and $2\in\cring^{\times}$, and is the natural first formalization target of
Part~II.

%----------------------------------------------------------------------
\section{The $\Sym^n$ convention, its semigroup extension, and the contragredient}\label{sec:symconv}
%----------------------------------------------------------------------
% SOURCE:  hecke_action.tex:287-372 (RH);
%          sec:contragredient NEW, from [X] SS0.4 C-A4 and [OD] SSOQ4
% NOTATION: rows u_j, sigma, Mat_2(Z), cring of sec:notation-crosswalk applied
% ERRATA:  none land here

\paragraph{What this chapter fixes.}
The coefficient module and the exact convention by which a $2\times2$ integer matrix
acts on it: $\rhon$, defined as the \emph{monoid} homomorphism on all of
$\Mat_2(\Z)$ --- not only on $\Ga$ --- because the Hecke coset representatives have
determinant $p$ and are not invertible; the extension of that action to the monoid
$\De$; the determinant of $\rhon$, which is where the bad prime $p=\ell$ enters; and the
contragredient module $\Dn$ with its adjugate action and the perfect pairing
$\VR\times\Dn\to\cring$ against which $T_p$ is self-adjoint.

\paragraph{Standing hypotheses.}
$\cring$ is a commutative ring with $1$; where a characteristic or an invertibility
hypothesis is needed it is stated in the statement that needs it (rule~\ref{A1}).
\emph{No} hypothesis on $\det$ is used anywhere in this chapter: every definition and
every statement below is valid for singular matrices, which is what makes the bad case
\ref{H2} a case of the same construction rather than a different one.

\subsection*{The symmetric power and its action}

\begin{definition}[Integral symmetric power action]\label{def:rhon}
Let $W:=\Z^2$ with ordered basis $(x,y)$, and let $A=\smat{a}{b}{c}{d}\in\Mat_2(\Z)$ act
on $W$ in the \emph{column convention}: $A\cdot x=ax+cy$ and $A\cdot y=bx+dy$, i.e.\
$A\cdot e_j=\sum_t A_{tj}e_t$ with $x=e_1$, $y=e_2$. Let
\[
  \Vcoef:=\Sym^n W,\qquad
  u_j:=x^{\,n-j}y^{\,j}\ \ (0\le j\le n),\qquad d_V:=\dim\Vcoef=n+1 ,
\]
with $(u_0,\dots,u_n)$ as ordered basis. Then $\rhon(A):=\Sym^n(A)\in\End_\Z(\Vcoef)$ is
the multiplicative extension of the above to degree-$n$ monomials; concretely, column
$j$ of the matrix of $\rhon(A)$ in the basis $(u_t)$ holds the coefficients of
\begin{equation}\label{eq:rhondef}
  A_{\ast}(u_j)\;=\;(a\,x+c\,y)^{\,n-j}\,(b\,x+d\,y)^{\,j},
\end{equation}
that is, dehomogenising by $x=1$,
\[
  \rhon(A)_{t,j}\;=\;\bigl[y^{t}\bigr]\,(a+cy)^{\,n-j}(b+dy)^{\,j}
  \;=\;\sum_{s}\binom{n-j}{s}a^{\,n-j-s}c^{\,s}\binom{j}{t-s}b^{\,j-t+s}d^{\,t-s}.
\]
For a coefficient ring $\cring$ put $\VR:=\Vcoef\otimes_\Z\cring$ and
$\rhon^{\cring}:=\rhon\bmod\cring$; we drop the superscript. For $\cring=\Fl$ we also
write $\Mn=\VR$, the notation of the classification literature.
\end{definition}

\noindent
This is a constructive definition in the sense of rule~\ref{A3}: \eqref{eq:rhondef} is
literally the object the implementation builds
(\code{ComputeSymmetricPowerMatrixFM}, \file{ModArithmetic/FlatMatrix.cpp}), and
$\rhon(A)$ is therefore a mathematically unique matrix that any correct evaluation route
must reproduce entry for entry --- which is what
\code{SymmetricPowerValidation} asserts of the eight branches that route actually takes.

\begin{proposition}[Code convention $=$ standard action; multiplicativity]\label{prop:symconv}
\PROVED{} For every $A\in\Mat_2(\Z)$ and every $n\ge0$:
\begin{enumerate}
\item $\rhon(A)\bmod\ell=\code{ComputeSymmetricPowerMatrixFM}(A\bmod\ell,\,n)$;
\item $A\mapsto\rhon(A)$ is a homomorphism of \emph{monoids}
  $(\Mat_2(\Z),\cdot)\to(\End_\Z(\Vcoef),\cdot)$: $\rhon(A_1A_2)=\rhon(A_1)\rhon(A_2)$
  and $\rhon(I)=I$. In particular no transpose, inverse or contragredient is involved;
  $\rhon$ restricted to $\Ga$ is a group homomorphism, and its values on $T,S,ST,TS$ are
  exactly the matrices attached to the eight rows of
  Definition~\ref{def:model7rel}; \LEAN{prop:symconv}
\item $\rhon$ is integral, $\rhon(\Mat_2(\Z))\subseteq\Mat_{n+1}(\Z)$, and every entry of
  $\rhon(A)$ is bounded in absolute value by $(2\max_{t,j}\abs{A_{tj}})^{\,n}$. \LEAN{prop:symconv}
\end{enumerate}
\end{proposition}

\begin{proof}
(1) The substitution $x\mapsto ax+cy$, $y\mapsto bx+dy$ is an algebra endomorphism
$\mathrm{subst}_A$ of $\Z[x,y]$; its matrix on $\mathrm{span}(x,y)$ in the basis
$(x,y)$, images as columns, is $A$ itself, i.e.\ it is the action of
Definition~\ref{def:rhon}; and the implementation tabulates its matrix on the degree-$n$
part in the basis $(u_j)$, columns being images, which is $\Sym^n(A)$.

(2) $\mathrm{subst}_{A_1}\circ\mathrm{subst}_{A_2}=\mathrm{subst}_{A_1A_2}$: on $x$,
$\mathrm{subst}_{A_1}(a_2x+c_2y)=a_2(a_1x+c_1y)+c_2(b_1x+d_1y)
=(a_1a_2+b_1c_2)x+(c_1a_2+d_1c_2)y$, and $(A_1A_2)_{11}=a_1a_2+b_1c_2$,
$(A_1A_2)_{21}=c_1a_2+d_1c_2$; likewise on $y$; extend multiplicatively to degree $n$.
Composition of matrices in the columns-are-images convention matches composition of
maps, and $\mathrm{subst}_I=\id$.

(3) By \eqref{eq:rhondef} the sum of the absolute values of the coefficients of
$A_\ast(u_j)$ is at most $(\abs a+\abs c)^{n-j}(\abs b+\abs d)^{j}\le(2M)^n$ with
$M=\max_{t,j}\abs{A_{tj}}$; each single coefficient is bounded by that sum.
\end{proof}

\begin{contractbox}
The action of \emph{any} $g\in\Mat_2(\Z)$ on $\Sym^n$ coefficients --- Hecke coset
representatives, non-invertible modulo $\ell$, included --- is $\rhon(g)$ of
Definition~\ref{def:rhon}, computed as
\code{ComputeSymmetricPowerMatrixFM}$(g\bmod\ell,\,n)$. It acts on cochain values by
\emph{left} matrix multiplication, matching the \code{SimplexRelation} usage
$f(\sigma_{\mathrm{lhs}})=\rhon(g)f(\sigma_{\mathrm{rhs}})$ of
Definition~\ref{def:model7rel}. \textbf{No determinant twist is applied}: the
coefficient module is plain $\Sym^n(\Z^2)$, not $\Sym^n\otimes\det^{m}$. This is the
\emph{integral normalisation}, and it is what fixes the eigenvalue normalisation of the
oracles in \S\ref{sec:example}.
\end{contractbox}

\begin{definition}[Semigroup extension]\label{def:semigroup}
Let
\[
  \De:=\Mat_2(\Z)\cap\GL_2(\Q)^{+}=\{g\in\Mat_2(\Z):\det g>0\},
\]
a multiplicative submonoid of $\GL_2(\Q)$ containing $\Ga$. ``The representation extends
to the sub-semigroup of $\GL_2(\Q)$ generated by $\Ga$ and the chosen coset
representatives'' means precisely: there is a monoid homomorphism
$\rho\colon\De\to\End_\Z(\Vcoef)$ with $\rho|_{\Ga}=\rhon|_{\Ga}$. For $\Vcoef=\Sym^n$
this holds with $\rho=\rhon$, by Proposition~\ref{prop:symconv}(2), since
$\De\subseteq\Mat_2(\Z)$ and all coset representatives of
Proposition~\ref{prop:cosets} lie in $\De$.
\end{definition}

\noindent
The sentence just quoted is [v5], \S12.1. That section is superseded in full
(Appendix~\ref{app:errata}, entries~\ref{err:e19}, \ref{err:e20}), but the requirement
it states is not among its defects and is met as shown.

\begin{lemma}[Determinant of the symmetric power]\label{lem:detsym}
\PROVED{} \LEAN{lem:detsym} (the second sentence for $\cring$ a field of characteristic~$\ell$) $\det\rhon(A)=(\det A)^{\,n(n+1)/2}$ for every $A\in\Mat_2(\Z)$. Hence for
$g\in\De$ with $\det g=p$ and $\cring$ of characteristic $\ell$: $\rhon(g)$ is
invertible over $\cring$ if and only if $\ell\neq p$ or $n=0$.
\end{lemma}

\begin{proof}
Both sides are polynomials in $a,b,c,d$ with integer coefficients, so it suffices to
prove the identity on a Zariski-dense subset of $\Mat_2(\Cx)$; the matrices with two
distinct eigenvalues form such a subset. For $A=P\,\mathrm{diag}(\alpha,\beta)\,P^{-1}$
with $P$ invertible, Proposition~\ref{prop:symconv}(2) gives
$\rhon(A)=\rhon(P)\rhon(\mathrm{diag}(\alpha,\beta))\rhon(P)^{-1}$ (with
$\rhon(P)^{-1}=\rhon(P^{-1})$), so
$\det\rhon(A)=\det\rhon(\mathrm{diag}(\alpha,\beta))$. The latter matrix is diagonal
with entries $\alpha^{\,n-j}\beta^{\,j}$, whence
$\det=\alpha^{\sum_j(n-j)}\beta^{\sum_j j}=(\alpha\beta)^{n(n+1)/2}=(\det A)^{n(n+1)/2}$.
For the second claim: $n(n+1)/2=0$ exactly for $n=0$, and $\ell\mid p^{\,n(n+1)/2}$ with
$n\ge1$ exactly when $\ell=p$.
\end{proof}

\noindent
Lemma~\ref{lem:detsym} is the whole content of hypothesis \ref{H2}: at $p=\ell$ and
$n\ge1$ every coset representative acts by a singular matrix, which is why $U_\ell$ can
be nilpotent on $H^{\bullet}$ and why the recursion of
Corollary~\ref{cor:ulrecursion} degenerates --- \emph{not} a defect of the
construction, which never inverts $\rhon(g)$.

\subsection{The contragredient local system and the pairing}\label{sec:contragredient}

New in this document; owner decision OQ4 (\file{hecke\_extensions\_owner\_decisions.md}
\S OQ4) makes $\Dn$ a first-class local system \emph{together with} its pairing map,
rather than the source of duality predictions only, and requires (architecture rule R2)
that both be defined over $\Z$ and the $\cring$-versions be their base changes. That is
how they are defined here.

Recall the adjugate on $2\times2$ matrices,
\[
  \adjg\smat{a}{b}{c}{d}:=\smat{d}{-b}{-c}{a},
\]
which satisfies, for all $g,h\in\Mat_2(\Z)$ and by direct computation,
\begin{equation}\label{eq:adjrules}
  g\,\adjg g=\adjg g\,g=(\det g)I,\qquad
  \adjg(gh)=(\adjg h)(\adjg g),\qquad
  \adjg(\adjg g)=g,\qquad \adjg I=I .
\end{equation}
So $\adjg$ is an anti-automorphism of the monoid $\Mat_2(\Z)$, restricting to $\De$, and
$\adjg\gamma=\gamma^{-1}$ for $\gamma\in\Ga$.

\begin{definition}[The contragredient action $\rhodual(g)=\rhon(\adjg g)^{\top}$]\label{def:contragredient}
Let $\cring$ be a commutative ring. Put
\[
  \Dn:=\Hom_\Z(\Vcoef,\Z),\qquad \Dn\otimes_\Z\cring=\Hom_{\cring}(\VR,\cring),
\]
free of rank $d_V=n+1$ with the basis $(u_0^{\vee},\dots,u_n^{\vee})$ dual to
$(u_0,\dots,u_n)$ (the second equality because $\Vcoef$ is free of finite rank). Define
\[
  \rhodual(g):=\rhon(\adjg g)^{\top}\in\End_\Z(\Dn)
  \qquad\text{in the dual basis, }g\in\Mat_2(\Z),
\]
equivalently, basis-free,
\begin{equation}\label{eq:dualaction}
  \bigl(\rhodual(g)\varphi\bigr)(v)\;=\;\varphi\bigl(\rhon(\adjg g)\,v\bigr),
  \qquad \varphi\in\Dn,\ v\in\Vcoef .
\end{equation}
\LEAN{def:contragredient} (the action \eqref{eq:dualaction})
The \emph{contragredient local system} on $(K,\Rel)$ is the local system whose fibre is
$\Dn\otimes\cring$ and whose relation matrix for the relation
$(\sigma_{\mathrm{lhs}},\sigma_{\mathrm{rhs}},\gamma)\in\Rel$ of
Definition~\ref{def:model7rel} is $\rhodual(\gamma)$. Both $\Dn$ and $\rhodual$ are defined
over $\Z$; every coefficient ring is a base change of that datum, and the reduction maps
$\Z/\ell^{\,e+1}\to\Z/\ell^{\,e}$ commute with $\rhodual$ entrywise.
\end{definition}

\begin{proposition}[The perfect pairing $\Vcoef\times\Dn\to\cring$ and self-adjointness of $T_p$]\label{prop:pairing}
\PROVED{} \LEAN{prop:pairing} (clauses (1)--(5); (2) as the reflexivity of the evaluation, (5) as two matrix identities) Let $\cring$ be a commutative ring, $n\ge0$, and let
\[
  \langle\,\cdot\,,\,\cdot\,\rangle\colon \VR\times(\Dn\otimes\cring)\longrightarrow\cring,
  \qquad \langle v,\varphi\rangle:=\varphi(v)
\]
be the evaluation pairing. Then:
\begin{enumerate}
\item $\rhodual\colon\Mat_2(\Z)\to\End_\Z(\Dn)$ is a homomorphism of \emph{monoids}, and
  $\rhodual(\gamma)=\rhon(\gamma^{-1})^{\top}$ for $\gamma\in\Ga$, i.e.\ $\rhodual|_\Ga$
  is the contragredient of $\rhon|_\Ga$ in the usual sense.
\item The pairing is \emph{perfect} over every $\cring$: its Gram matrix in the bases
  $(u_j)$, $(u_j^{\vee})$ is the identity, so the induced map
  $\Dn\otimes\cring\to\Hom_{\cring}(\VR,\cring)$ is an isomorphism. It is defined over
  $\Z$ and the $\cring$-pairing is its base change; in particular it commutes with the
  reduction maps $\Z/\ell^{\,e+1}\to\Z/\ell^{\,e}$ on both arguments.
\item \textbf{Adjointness.} For all $g\in\Mat_2(\Z)$, $v\in\VR$,
  $\varphi\in\Dn\otimes\cring$,
  \begin{equation}\label{eq:pairadjoint}
    \bigl\langle \rhon(g)\,v,\ \varphi\bigr\rangle
    \;=\;\bigl\langle v,\ \rhodual(\adjg g)\,\varphi\bigr\rangle .
  \end{equation}
\item \textbf{Twisted invariance.} For all $g$,
  $\langle\rhon(g)v,\rhodual(g)\varphi\rangle=(\det g)^{\,n}\langle v,\varphi\rangle$. In
  particular the pairing is $\Ga$-invariant, so it is a pairing of local systems on
  $(K,\Rel)$ in the sense of Definition~\ref{def:contragredient}; on $\De$ it is
  invariant only up to the character $\det^{\,n}$.
\item \textbf{The double coset is self-adjugate.} $\adjg(\Ga g\Ga)=\Ga\,(\adjg g)\,\Ga$
  for every $g$, and for $g=\mathrm{diag}(1,p)$ one has
  $\adjg g=\mathrm{diag}(p,1)=SgS^{-1}$, hence $\Ga(\adjg g)\Ga=\Ga g\Ga$.
\end{enumerate}
\end{proposition}

\begin{proof}
(1) By the identities \eqref{eq:adjrules} and Proposition~\ref{prop:symconv}(2),
$\rhodual(gh)=\rhon(\adjg(gh))^{\top}=\bigl(\rhon(\adjg h)\rhon(\adjg g)\bigr)^{\top}
=\rhon(\adjg g)^{\top}\rhon(\adjg h)^{\top}=\rhodual(g)\rhodual(h)$, and
$\rhodual(I)=I$. For $\gamma\in\Ga$, $\adjg\gamma=(\det\gamma)\gamma^{-1}=\gamma^{-1}$.

(2) $\langle u_t,u_j^{\vee}\rangle=\delta_{tj}$ by definition of the dual basis, so the
Gram matrix is $I$, which is invertible over every $\cring$; everything in sight is the
base change of the corresponding $\Z$-datum because $\Vcoef$ is free of finite rank.

(3) By \eqref{eq:dualaction} applied to $\adjg g$ and $\adjg(\adjg g)=g$:
$\langle v,\rhodual(\adjg g)\varphi\rangle=\bigl(\rhodual(\adjg g)\varphi\bigr)(v)
=\varphi\bigl(\rhon(\adjg\adjg g)v\bigr)=\varphi(\rhon(g)v)=\langle\rhon(g)v,\varphi\rangle$.

(4) By \eqref{eq:dualaction}, Proposition~\ref{prop:symconv}(2) and the identities \eqref{eq:adjrules},
$\langle\rhon(g)v,\rhodual(g)\varphi\rangle
=\varphi\bigl(\rhon(\adjg g)\rhon(g)v\bigr)
=\varphi\bigl(\rhon((\det g)I)\,v\bigr)=(\det g)^{\,n}\varphi(v)$, using
$\rhon(\lambda I)=\lambda^{n}I$, which is Definition~\ref{def:rhon} for a scalar matrix. For
$\gamma\in\Ga$, $\det\gamma=1$.

(5) $\adjg(\gamma_1g\gamma_2)=(\adjg\gamma_2)(\adjg g)(\adjg\gamma_1)
=\gamma_2^{-1}(\adjg g)\gamma_1^{-1}$, and $\gamma\mapsto\gamma^{-1}$ is a bijection of
$\Ga$; so $\adjg(\Ga g\Ga)=\Ga(\adjg g)\Ga$. For $g=\mathrm{diag}(1,p)$,
$\adjg g=\mathrm{diag}(p,1)$, and
$S\,\mathrm{diag}(1,p)\,S^{-1}=\smat{0}{-1}{1}{0}\smat{1}{0}{0}{p}\smat{0}{1}{-1}{0}
=\mathrm{diag}(p,1)$.
\end{proof}

\begin{remark}[What self-adjointness of $T_p$ means, and where it is proved]\label{rmk:selfadjoint}
Equation~\eqref{eq:pairadjoint} is the coefficient-level input to the classical
statement that the Hecke operators are self-adjoint for Poincar\'e duality, and it is
exactly what makes the adjoint operator computable: if
$\Ga g\Ga=\bigsqcup_j\Ga g_j$ (Proposition~\ref{prop:cosets}), then the operator adjoint
to $T_{\Ga g\Ga}$ on the $\Dn$-side is the one built from the representatives
$\adjg g_j$, i.e.\ $T_{\Ga(\adjg g)\Ga}$; and by
Proposition~\ref{prop:pairing}(5) that is again $T_{\Ga g\Ga}$ for $g=\mathrm{diag}(1,p)$
--- \emph{no} determinant factor appears, because \eqref{eq:pairadjoint} carries none.
The cohomological statement --- perfect Hecke-compatible pairings
$H^{t}_{c}(\Ga;\Mn)\times H^{2-t}(\Ga;\Dn)\to\Fl$ and
$H^{t}(\partial;\Mn)\times H^{1-t}(\partial;\Dn)\to\Fl$, under the hypothesis
$6\in\cring^{\times}$ --- is Claim~\ref{cl:pd} of \S\ref{sec:six}, where the compact
supports and the boundary it refers to are constructed. This chapter proves the
coefficient-level half and nothing more: Proposition~\ref{prop:pairing} says nothing
about cohomology, and \eqref{eq:pairadjoint} must not be read as if it did.
\end{remark}

\begin{lemma}[$\Vcoef$ versus $\Dn$: the integral invariant form]\label{lem:symdual}
\PROVED{} \LEAN{lem:symdual} (the displayed identities, the equivariance of $\Phi_n$ and the invertibility criterion; [GT] defines $B_n$ by its Gram matrix \eqref{eq:BnGram}, so the permutation-sum formula and its well-definedness are not covered) Define $B_n\colon\Vcoef\times\Vcoef\to\Z$ on products of vectors of $W=\Z^2$ by
\[
  B_n(v_1\cdots v_n,\,w_1\cdots w_n)
  \;:=\;\sum_{\pi\in\mathfrak{S}_n}\ \prod_{t=1}^{n}\det\bigl[\,v_t\ \big|\ w_{\pi(t)}\,\bigr] .
\]
Then $B_n$ is a well-defined symmetric-in-each-argument bilinear form with
\begin{equation}\label{eq:BnGram}
  B_n(u_t,u_j)=(-1)^{t}\,t!\,(n-t)!\ \delta_{t+j,n},
\end{equation}
and for all $g\in\Mat_2(\Z)$ and $v,w\in\Vcoef$
\[
  B_n\bigl(\rhon(g)v,\rhon(g)w\bigr)=(\det g)^{\,n}B_n(v,w),
  \qquad
  B_n\bigl(\rhon(g)v,\,w\bigr)=B_n\bigl(v,\,\rhon(\adjg g)w\bigr).
\]
Consequently $\Phi_n\colon\Vcoef\to\Dn$, $\Phi_n(v):=B_n(v,\cdot)$, satisfies
$\Phi_n\circ\rhon(g)=\rhodual(g)\circ\Phi_n$ for every $g\in\Mat_2(\Z)$ --- it is a
morphism of $\De$-modules, hence of local systems on $(K,\Rel)$ --- its matrix is
antidiagonal with entries \eqref{eq:BnGram}, and
\[
  \Phi_n\otimes\cring\ \text{is an isomorphism}
  \iff \textstyle\prod_{t=0}^{n}t!\,(n-t)!\in\cring^{\times}
  \iff n!\in\cring^{\times},
\]
which for $\cring$ a field of characteristic $\ell$ means $n\le\ell-1$.
\end{lemma}

\begin{proof}
$\det[\,v\mid w\,]$ is the standard symplectic form on $W$, and
$\det[\,gv\mid gw\,]=(\det g)\det[\,v\mid w\,]$. The displayed sum is unchanged under a
permutation of the $v_t$ and under a permutation of the $w_t$ (reindex $\pi$), so it
descends to $\Sym^n W\times\Sym^n W$ and is bilinear. Multiplying each of the $n$ factors
by $\det g$ gives the first equivariance. For \eqref{eq:BnGram}, take
$v_1\cdots v_n=x^{n-t}y^{t}$ and $w_1\cdots w_n=x^{n-j}y^{j}$ and use
$\det[x\mid x]=\det[y\mid y]=0$, $\det[x\mid y]=1=-\det[y\mid x]$: a permutation
contributes only if it matches each of the $n-t$ letters $x$ on the left to a letter $y$
on the right and each of the $t$ letters $y$ on the left to a letter $x$ on the right,
which forces $n-t=j$ and $t=n-j$, i.e.\ $t+j=n$; there are then $(n-t)!\,t!$ such
permutations, each contributing $1^{\,n-t}(-1)^{t}$.

For the second equivariance, substitute $w\mapsto\rhon(\adjg g)w$ in the first and use
$\rhon(g)\rhon(\adjg g)=\rhon((\det g)I)=(\det g)^{n}\id$:
$(\det g)^{n}B_n(\rhon(g)v,w)=(\det g)^{n}B_n(v,\rhon(\adjg g)w)$. Both sides are
polynomial in the entries of $g$ with integer coefficients, and $\Z[a,b,c,d]$ is a
domain in which $(\det g)^{n}=(ad-bc)^{n}\neq0$, so the factor cancels there; the
identity over any $\cring$ follows by specialising the resulting polynomial identity.
Then, for $w\in\Vcoef$,
$\bigl(\Phi_n(\rhon(g)v)\bigr)(w)=B_n(\rhon(g)v,w)=B_n(v,\rhon(\adjg g)w)
=\bigl(\Phi_n(v)\bigr)(\rhon(\adjg g)w)=\bigl(\rhodual(g)\Phi_n(v)\bigr)(w)$ by
\eqref{eq:dualaction}.

Finally $\Phi_n(u_t)=(-1)^{t}t!(n-t)!\,u_{n-t}^{\vee}$, so the matrix of $\Phi_n$ is
antidiagonal with those entries and $\Phi_n\otimes\cring$ is bijective iff all of them
are units. Since $n!=\binom{n}{t}t!(n-t)!$, each $t!(n-t)!$ divides $n!$, so $n!$ being a
unit makes all of them units; conversely the factor at $t=0$ is $0!\,n!=n!$. Over a
field of characteristic $\ell$, $n!\neq0$ iff $n<\ell$.
\end{proof}

\begin{remark}[Divided powers, and what is \emph{not} claimed]\label{rmk:divided}
The classification literature calls $\Dn$ the \emph{divided-power} module and records
that $\Mn$ fails to be self-dual for $n\ge\ell$ (\Espec{sec:symdiv}). What
Lemma~\ref{lem:symdual} proves is the sharper and weaker statement that the canonical
integral comparison map $\Phi_n$ fails to be an isomorphism exactly when
$n!\notin\cring^{\times}$. Whether $\VR$ and $\Dn\otimes\cring$ are \emph{abstractly}
isomorphic as $\Ga$-modules in that range is a different question; it is not used
anywhere in this document and is not asserted here.
\end{remark}

\paragraph{Code map.}
Definition~\ref{def:rhon}, Proposition~\ref{prop:symconv}:
\code{ComputeSymmetricPowerMatrixFM} (\file{ModArithmetic/FlatMatrix.cpp}), whose
branch table realises \eqref{eq:rhondef} in $O(d_V^2)$, against the double-sum reference
\code{ComputeSymmetricPowerMatrixFM\_reference}; gated by
\code{SymmetricPowerValidation} (and its slow-tier twin) entry for entry, and by
$\rhon(A)\rhon(B)=\rhon(AB)$ as a structural check. The deprecated
\code{ComputeSymmetricPowerMatrix} returns the same matrix in the retired \code{Matrix}
type. Lemma~\ref{lem:detsym}: the singular-matrix regime is exercised by the $p=\ell$
arms of \code{HeckeBadPrimeValidation}. Definition~\ref{def:semigroup}: the coset
representatives of \file{Hecke/HeckeCosetTables.cpp}, all of determinant $p^{\nu}>0$.
\textbf{Since v1.7.3 (hecke-extensions slot-6 WP-1, phase E7a) \S\ref{sec:contragredient} is
implemented}: \code{Mat22::adj} (\file{Hecke/ExactInt}, \eqref{eq:adjrules}),
$\rhodual=\rhon(\adjg\,\cdot)^{T}$ as \code{ComputeContragredientPowerMatrixFM} with the
dispatcher \code{ComputeFibreMatrixFM} (\code{LocalSystemUtils}), the contragredient
relation list \code{GenerateSL2ZContragredientRelationsWithGroupData}
(\code{LocalSystemSL2Z}), $\Phi_n$ of Lemma~\ref{lem:symdual} and the pairing of
Proposition~\ref{prop:pairing} in \file{Hecke/HeckeContragredient}, and $A(\Dn)(T_p)^{q}$
through \code{HeckeOperatorOnK}'s \code{FibreRepresentation} argument (default \code{SymN},
byte for byte the pre-slot-6 operator). \code{HeckeDualityValidation} asserts the
intertwining \emph{untwisted} (exponent $0$, as Lemma~\ref{lem:symdual} gives it) and the
$p^{\pm n}$-twisted relations to fail. The cohomology-level pairings are
\file{Hecke/HeckePairing} (slot-6 WP-3; \S\ref{sec:codemap}). The v1.7.2 status follows,
kept as history. \S\ref{sec:contragredient} \textbf{had no implementation at v1.7.2}: $\adjg$ occurs in the
code only as the $\adjg(g)/N$ projective trick of \file{Hecke/HeckeTauTable.cpp} and
\file{Hecke/BarCocycleH1.cpp} (each spelling the $2\times2$ adjugate out inline; there is
no adjugate operation on \code{Mat22}), and neither $\rhodual$ nor the pairing is built. They are phase E7 of the extension plan, made unconditional by
owner decision OQ4; the perfectness certificate that
Proposition~\ref{prop:pairing}(2) licenses --- Gram matrix $=I$ in the stored bases ---
is the cheapest possible one.

\paragraph{Formalization note.}
Definition~\ref{def:rhon} is a closed-form matrix over $\Z$ and
Proposition~\ref{prop:symconv}(2) is the statement that $\Sym^n$ is a functor on the
multiplicative monoid, available in \code{mathlib4} as the algebra map induced on
\code{MvPolynomial}; the bound in Proposition~\ref{prop:symconv}(3) is a decidable inequality. Lemma~\ref{lem:detsym}
is the only statement in this chapter proved by Zariski density; a formalization would
more comfortably prove it by induction on $n$ from the closed form \eqref{eq:rhondef} for
triangular $A$ together with a Bruhat-style factorisation, or simply
accept the polynomial-identity argument over $\Z[a,b,c,d]$ used in
Lemma~\ref{lem:symdual}. Definition~\ref{def:contragredient} and the whole of
Proposition~\ref{prop:pairing} are finite linear algebra over a commutative ring and
carry no analytic content; \eqref{eq:pairadjoint} is the natural statement to formalise
first, as everything else about the duality rests on it.

%----------------------------------------------------------------------
\section{The involution $\inv$}\label{sec:involution}
%----------------------------------------------------------------------
% SOURCE:  NEW.  plan SS0.4 C-A1 (hecke_extensions_implementation_plan.md:287-305);
%          spec (A3) (hecke_extensions.tex:882-889); owner decision OQ3.
%          No re-homed LaTeX: no superseded document states any of this.
% NOTATION: C3 (\inv is the involution and never the cone inclusion), C9 (zeta for
%          the elliptic points), C12 (M is not free), C15 (i is sqrt(-1), never an
%          index), C16 (\lift{sigma} for the distinguished lift), C17 (relation
%          elements are gamma; g is the Hecke element).
% ERRATA:  none land here.

\paragraph{What this chapter fixes.}
The element $\inv=\mathrm{diag}(-1,1)\in\GL_2(\Z)$, its simplicial realisation
$\invK$ with the orientation sign of every simplex, the involution $\inv^{\ast}$ it
induces on $\Ceq(K,\Rel;L)$, the exact hypothesis under which $\inv^{\ast}$ splits
$\Heq{\bullet}$ into $\pm$ eigenspaces --- it is $2\in\cring^{\times}$, and what
replaces the splitting at $\ell=2$ is stated here and not deferred --- and the
commutation $[\inv^{\ast},\Tdc{g}]=0$, with the level at which it holds and the
hypothesis each level needs. \textbf{Since v1.1.1 every clause of
Obligation~\ref{obl:iota} is implemented and has been run}: (i) and (iv) since
v0.13.0 --- and since v1.1.0 clause (iv) \emph{holds} on the canonical $\tau$ and
gates it (Remark~\ref{rmk:iotatausym}) --- and (ii), (iii), (v) since v1.1.1, when
$\inv^{\ast}$ was first built as a map of cochain complexes. \textbf{Since v1.2.0 the
sign of $\inv^{\ast}$ is fixed by owner decision OQ3 (2026-09-16)}:
Definition~\ref{def:iota}(iv) defines $\inv^{\ast}:=-\inv^{\ast}_{\mathrm{geo}}$, the
negative of the geometric pullback, so that ``$+$'' \emph{is by definition} the
eigenspace containing the Eisenstein classes. The v1.1.1 run had found those classes in
the $(-1)$-eigenspace of the geometric involution on every cell of a $77$-cell grid, by
two independent routes; the owner resolved it by negating the definition, not at a call
site, and Remark~\ref{rmk:iotasign} says what the record field \code{iota\_sign} now
means. Every predicate of the five clauses is gating and passes, the former
``recorded half'' of clause (v) included, which is now a definition-consistency oracle;
the honest cost of the choice --- constants in $H^{0}$ are odd, $\inv^{\ast}$ is the
geometric involution twisted by the sign character $\det$ of $\GL_2(\Z)/\SL_2(\Z)$ rather
than a pullback --- is stated at Definition~\ref{def:iota}(iv) and in
\file{docs/implementation/oq3\_resolution.md}. The code map at the end says which
validation discharges what.

\paragraph{Standing hypotheses.}
\ref{H1}--\ref{H4} of \S\ref{sec:conventions}, each repeated where it is used
(rule~\ref{A1}), together with one hypothesis proper to this chapter. It is stated for
an arbitrary $\Ga$-tessellation rather than for the seven vertices of
Definition~\ref{def:model7}, so that a refinement of that complex is admissible input
(the design direction recorded in Remark~\ref{rmk:orbifold}).

\begin{enumerate}[label=(I\arabic*)]
\item\label{I1} \textbf{$\inv$-stable tessellation.} The map $z\mapsto-\bar z$ of
  $\Hyp$ carries the complex $\Xbulk$ of \ref{H4} to itself simplicially --- each
  simplex of $\Xbulk$ onto a simplex of $\Xbulk$ --- and carries the strict fundamental
  domain $K$ of \ref{H4}(iii) onto itself.
\end{enumerate}

\noindent
\ref{I1} is a property of the triangulation alone. It is verified for the seven-vertex
model in Lemma~\ref{lem:iotamodel}, and an $\inv$-stable refinement of that model
satisfies it verbatim.

\subsection*{The element, its symmetric power, and its realisation}

\begin{definition}[The involution $\inv$, its simplicial realisation $\invK$, and $\rhon(\inv)$]\label{def:iota}
\begin{enumerate}
\item \textbf{The element.} $\inv:=\mathrm{diag}(-1,1)\in\GL_2(\Z)$, with $\inv^2=I$,
  $\det\inv=-1$ and therefore $\inv\notin\De$: it is \emph{not} a Hecke element and
  $\Tdc{\inv}$ is not defined. Conjugation
  $\conj_{\inv}\colon\gamma\mapsto\inv\gamma\inv^{-1}=\inv\gamma\inv$ is an involutive
  automorphism of $\Ga=\SL_2(\Z)$ and of the monoid $\Mat_2(\Z)$, and on the generator
  multiset $\mathcal{G}=\{T,S,ST,TS\}$ of \S\ref{sec:model} and on the centre it is
  \begin{equation}\label{eq:iotaconj}
    \conj_{\inv}(T)=T^{-1},\quad \conj_{\inv}(S)=S^{-1},\quad
    \conj_{\inv}(ST)=(TS)^{-1},\quad \conj_{\inv}(TS)=(ST)^{-1},\quad
    \conj_{\inv}(-I)=-I .
  \end{equation}
  On $\Hyp$, $\inv$ acts by $z\mapsto-\bar z$ (the $\det<0$ convention
  $z\mapsto(a\bar z+b)/(c\bar z+d)$), and $\inv(\gamma z)=\conj_{\inv}(\gamma)(\inv z)$.
\item \textbf{On coefficients.} $\rhon(\inv)\in\GL(\Vcoef)$ is the value at $\inv$ of
  the monoid homomorphism $\rhon$ on all of $\Mat_2(\Z)$ (Definition~\ref{def:rhon},
  Proposition~\ref{prop:symconv}(2)) --- which is why that chapter defined $\rhon$ on
  $\Mat_2(\Z)$ and not on $\De$. Explicitly, from \eqref{eq:rhondef} with
  $(a,b,c,d)=(-1,0,0,1)$,
  \begin{equation}\label{eq:rhoniota}
    \rhon(\inv)\,u_j=(-1)^{\,n-j}u_j,\qquad
    \rhon(\inv)=\mathrm{diag}\bigl((-1)^{\,n-j}\bigr)_{0\le j\le n},\qquad
    \rhon(\inv)^2=\rhon(I)=I .
  \end{equation}
  In characteristic $2$, $\rhon(\inv)=I$ for every $n$.
\item \textbf{On the complex.} Assume \ref{H4} and \ref{I1}. Write $\inv_{\Xbulk}$ for
  the simplicial automorphism of $\Xbulk$ induced by $z\mapsto-\bar z$ and $\invK$ for
  its restriction to $K$, both extended $\Z$-linearly to simplicial chains, where they
  carry the stored orientations of \S\ref{sec:conventions} to stored orientations up to
  sign; $\invK$ and $\inv_{\Xbulk}$ are two realisations of one notion and are indexed
  accordingly (\ref{N2}). They satisfy
  $\inv_{\Xbulk}(\gamma\cdot x)=\conj_{\inv}(\gamma)\cdot\inv_{\Xbulk}(x)$,
  $\inv_{\Xbulk}\circ\partial=\partial\circ\inv_{\Xbulk}$,
  $\aug\circ\inv_{\Xbulk}=\aug$, $\inv_{\Xbulk}^2=\id$, and
  $\inv_{\Xbulk}(\lift\sigma)=\widetilde{\invK\sigma}$ for $\sigma\in\orient{q}(K)$.
\item \textbf{On cochains.} For $f\in C^q(K;\VR)$, read as a homomorphism on $C_q(K)$,
  write $\inv^{\ast}_{\mathrm{geo}}$ for the \emph{geometric} pullback along
  $(\invK,\rhon(\inv))$ and define $\inv^{\ast}$ as its negative:
  \begin{equation}\label{eq:iotacochain}
    (\inv^{\ast}_{\mathrm{geo}}f)(\sigma)\;:=\;\rhon(\inv)\,f\bigl(\invK\sigma\bigr),
    \qquad
    \inv^{\ast}\;:=\;-\,\inv^{\ast}_{\mathrm{geo}},\quad\text{i.e.}\quad
    (\inv^{\ast}f)(\sigma)\;=\;-\rhon(\inv)\,f\bigl(\invK\sigma\bigr),
    \qquad\sigma\in\orient{q}(K),
  \end{equation}
  in every degree $q$ and for every coefficient module $\VR$ this document uses ($\Dn$
  of Definition~\ref{def:contragredient} included). \textbf{$\inv^{\ast}$, not
  $\inv^{\ast}_{\mathrm{geo}}$, is the involution of this document}: $H^{q}_{\mathrm{eq},\pm}$,
  \code{iota\_sign}, Obligation~\ref{obl:iota} and every $\pm$ statement below refer to
  it. The sign is the owner decision OQ3 of 2026-09-16 (\textbf{v1.2.0}; previously
  $\inv^{\ast}=\inv^{\ast}_{\mathrm{geo}}$): it is chosen so that ``$+$'' is \emph{by
  definition} the eigenspace containing the Eisenstein classes, Remark~\ref{rmk:iotasign}.
  What the sign is, honestly: $-\rhon(\inv)=\det(\inv)\,\rhon(\inv)=(\det\otimes\rhon)(\inv)$
  with $\det|_{\Ga}=1$ --- so on every relation element $\gamma\in\Ga$ the twisted fibre
  map $(\det\otimes\rhon)(\gamma)$ equals $\rhon(\gamma)$ and the relations of $\Rel$ are
  preserved --- so $\inv^{\ast}$ is the pullback along
  $(\invK,\det(\inv)\rhon(\inv))$ --- an automorphism of the local system $(K,\Rel,L)$ in the
  sense of Proposition~\ref{prop:iotacommutes}(1)--(2), since the intertwining relation is
  linear in the fibre map --- but, for $n$ even, \emph{not} the pullback along any lift of
  $z\mapsto-\bar z$ to $\GL_2(\Z)$: the two lifts $\pm\inv$ give
  $\rhon(\pm\inv)=(\pm1)^{n}\rhon(\inv)$, the same matrix for $n$ even. (For $n$ odd it
  \emph{is} the pullback along the lift $-\inv$; but for $\ell\neq2$ and $n$ odd
  $\Heq{\bullet}=0$, Remark~\ref{rmk:negI}, and at $\ell=2$ the sign is invisible, so on
  every group this document computes with a non-zero answer the twist is not a lift.) It is the
  geometric involution twisted by the sign character of $\GL_2(\Z)/\SL_2(\Z)$, the
  mechanism by which the classical modular-symbol ``$+$'' (the side of the boundary symbol
  \STD{\cite[Ch.~2]{cremona}}) differs from the naive pullback. Consequences:
  $\inv^{\ast}=-\id$ on $H^{0}(\Ga;\Fl)$ at $n=0$ (constants are \emph{odd}); at $\ell=2$
  the sign is invisible, $\inv^{\ast}=\inv^{\ast}_{\mathrm{geo}}$. In the flat coordinates
  \eqref{eq:flatindex} $\inv^{\ast}$ is a signed block permutation: one block
  $-\varepsilon_{\sigma}\rhon(\inv)$ per simplex $\sigma$, at block position
  $(\sigma,\abs{\invK\sigma})$, where $\varepsilon_{\sigma}\in\{\pm1\}$ is the sign of
  $\invK\sigma$ in the table of Lemma~\ref{lem:iotamodel} (the form of the code map below
  and of \code{HeckeIotaSignValidation}; $\invK$ being an involution with
  $\varepsilon_{\invK\sigma}=\varepsilon_{\sigma}$ and $\rhon(\inv)$ diagonal, the matrix
  equals its transpose), and no data of its own beyond the single $d_V\times d_V$
  diagonal matrix \eqref{eq:rhoniota}. A third party reproduces any $\pm$ statement of a
  record from $\inv^{\ast}=-\inv^{\ast}_{\mathrm{geo}}$ and nothing else.
\end{enumerate}
\end{definition}

\begin{lemma}[The seven-vertex model is $\inv$-stable, with its sign table]\label{lem:iotamodel}
\PROVED{} Let $K$ be as in Definition~\ref{def:model7}, with the four horocycle vertices
placed symmetrically about the imaginary axis, i.e.\
$\mathrm{Re}\,v_1=-\mathrm{Re}\,v_2$. (Definition~\ref{def:model7} fixes
$\mathrm{Re}\,v_0=-\tfrac12$ and $\mathrm{Re}\,v_3=+\tfrac12$ and leaves
$\mathrm{Re}\,v_1,\mathrm{Re}\,v_2$ free, subject only to
$\abs{\mathrm{Re}}<\tfrac12$; the symmetric placement is what makes $\inv$ simplicial
and it is the only use made of it in this document.) Then \ref{I1} holds, and $\invK$ is
the vertex permutation
\[
  0\leftrightarrow3,\qquad 1\leftrightarrow2,\qquad 4\leftrightarrow5,\qquad
  6\ \text{fixed},
\]
an involutive simplicial automorphism of $K$ whose orientation sign is $-1$ on exactly
six of the $23$ simplices:
\begin{center}
\begin{tabular}{@{}lll@{}}
\toprule
$q$ & $\invK\sigma=-\sigma'$ (sign $-1$) & $\invK\sigma=+\sigma'$ (sign $+1$) \\
\midrule
$0$ & --- & all seven vertices \\
$1$ & $\{v_0v_1\}\!\leftrightarrow\!\{v_2v_3\}$;\ \ $\{v_1v_2\}$ to itself
    & the other eight edges \\
$2$ & $\{v_0v_1v_4\}\!\leftrightarrow\!\{v_2v_3v_5\}$;\ \ $\{v_1v_2v_6\}$ to itself
    & $\{v_1v_4v_6\}\!\leftrightarrow\!\{v_2v_5v_6\}$ \\
\bottomrule
\end{tabular}
\end{center}
\end{lemma}

\begin{proof}
$z\mapsto-\bar z$ preserves $\mathrm{Im}$ and $\abs{z}$ and negates $\mathrm{Re}$, so it
fixes setwise the horocycle $\{\mathrm{Im}\,z=Y\}$, the unit circle and the region
$D_Y=\abs{K}$ of Lemma~\ref{lem:model7coincidences}. On vertices it exchanges
$v_0\leftrightarrow v_3$ and, by the symmetric placement, $v_1\leftrightarrow v_2$; and
$v_4=\zeta^2=-\tfrac12+\tfrac{\sqrt3}{2}i\mapsto\zeta=v_5$, $v_6=i\mapsto v_6$. It maps
the five faces of Definition~\ref{def:model7} to $\{v_2v_3v_5\}$, $\{v_1v_2v_6\}$,
$\{v_0v_1v_4\}$, $\{v_2v_5v_6\}$, $\{v_1v_4v_6\}$ --- the same five --- hence the eleven
edges to the eleven edges; so $\invK$ is a simplicial automorphism of $K$ with
$\invK^2=\id$. That $\inv$ is simplicial on all of $\Xbulk$ follows from \ref{H4}'s
definition of $\Xbulk$ as the tessellation by $\Ga$-translates of those five simplices,
together with $\inv(\gamma\cdot x)=\conj_{\inv}(\gamma)\cdot\inv(x)$ and
$\conj_{\inv}(\Ga)=\Ga$; and $\inv_{\Xbulk}(K)=K$ is what was just shown. This is
\ref{I1}.

The signs are the signs of the permutations sorting the images back into increasing
vertex order. Edges: $(0,1)\mapsto(3,2)$, $(1,2)\mapsto(2,1)$ and $(2,3)\mapsto(1,0)$
each need one transposition; the other eight map increasingly ---
$(0,4)\mapsto(3,5)$, $(1,4)\mapsto(2,5)$, $(1,6)\mapsto(2,6)$, $(2,5)\mapsto(1,4)$,
$(2,6)\mapsto(1,6)$, $(3,5)\mapsto(0,4)$, $(4,6)\mapsto(5,6)$, $(5,6)\mapsto(4,6)$.
Triangles: $(0,1,4)\mapsto(3,2,5)$, $(1,2,6)\mapsto(2,1,6)$ and $(2,3,5)\mapsto(1,0,4)$
each need one transposition, while $(1,4,6)\mapsto(2,5,6)$ and $(2,5,6)\mapsto(1,4,6)$
are already increasing.
\end{proof}

\begin{remark}[What the sign table is, and what it is not]\label{rmk:iotasigns}
The six minus signs are an artefact of the \emph{stored} orientation --- increasing
vertex index, \S\ref{sec:conventions} --- which is not a coherent orientation of the
surface. They are not a convention and nothing may be chosen about them: they are forced
by Definition~\ref{def:model7}'s storage rule, and $\inv^{\ast}$ is simply wrong without
them. They are printed because a reader cannot recompute them from the vertex
permutation alone without redoing the sort, and because a refinement of $K$ will have a
different table, computed the same way.
\end{remark}

\subsection*{Automorphism of the local system, and commutation with the Hecke action}

\begin{proposition}[$(\invK,\rhon(\inv))$ is an automorphism of $(K,\Rel,L)$, and $\inv$ commutes with $\Tdc{g}$]\label{prop:iotacommutes}
Assume \ref{H4} and \ref{I1}, and let $\cring$ be a commutative ring.
\begin{enumerate}
\item \PROVED{} \textbf{(Automorphism of the relation datum; no hypothesis on $\ell$.)}
  For every relation $(\sigma_{\mathrm{lhs}},\sigma_{\mathrm{rhs}},\gamma)\in\Rel$ of
  Definition~\ref{def:model7rel}, the triple
  $\bigl(\invK\sigma_{\mathrm{lhs}},\invK\sigma_{\mathrm{rhs}},\conj_{\inv}(\gamma)\bigr)$
  lies in the sub-groupoid $\langle\Rel\rangle$ of
  Definition~\ref{def:relationdatum}. It is in fact the \emph{inverse} of a member of
  the list: each of the six $T$- and $S$-rows goes to the inverse of itself, the two
  elliptic stabiliser rows $(\{v_4\},\{v_4\},ST)$ and $(\{v_5\},\{v_5\},TS)$ are
  exchanged, and each central row $(\sigma,\sigma,-I)$ goes to the central row at
  $\abs{\invK\sigma}$. \textbf{Up to inversion is the strongest statement available}:
  $\Rel$ as a \emph{list} is not preserved, its generated groupoid is. Part of the
  claim is that in every row the two sides carry the \emph{same} orientation sign, so
  that the common sign cancels and the image triple is a triple of \emph{stored}
  simplices, which is what membership in $\langle\Rel\rangle$ asserts.
\item \PROVED{} \textbf{(Involution of the cochain complex; no hypothesis on $\ell$.)}
  $\inv^{\ast}$ of \eqref{eq:iotacochain} maps $\Cq{q}(K,\Rel;L)$ into itself, commutes
  with $d$, and satisfies $(\inv^{\ast})^2=\id$; so does $\inv^{\ast}_{\mathrm{geo}}$, and
  the two differ by the global sign $-\id$, which commutes with everything in sight.
  Hence $\inv^{\ast}$ preserves $\Zq{q}$ and $\Bq{q}$ and induces an involution of
  $\Heq{q}(K,\Rel;L)$ in every degree $q$. Under $\Psi$ of
  Theorem~\ref{thm:explicit_complex_equiv} it corresponds to
  $\inv^{\#}:=-\inv^{\#}_{\mathrm{geo}}$, where
  $(\inv^{\#}_{\mathrm{geo}}\phi)(x):=\rhon(\inv)\,\phi(\inv_{\Xbulk}x)$ on
  $\Hom_{\Z[\Ga]}\bigl(C_{\bullet}(\Ku),\VR\bigr)$, i.e.\
  $\Psi\circ\inv^{\#}=\inv^{\ast}\circ\Psi$ (and likewise with the subscript
  $\mathrm{geo}$ on both sides).
\item \textbf{(Hecke.)} Let $g\in\De$ with $\conj_{\inv}(g)=g$ --- which holds for
  \emph{every} operator of \S\ref{sec:heckedatum}, all of whose elements
  $\mathrm{diag}(1,p^{\nu})$ and $pI$ are diagonal --- let $s_1,\dots,s_m$,
  $\delta_j=s_j^{-1}g$ and $\tau$ be as in Definition~\ref{def:hecke}, and put
  $s^{\inv}_j:=\conj_{\inv}(s_j)$ and
  $\tau^{\inv}:=\inv_{\Xbulk}\circ\tau\circ\inv_{\Xbulk}$. Then:
  \begin{enumerate}
  \item[(3a)] \PROVED{} \textbf{no hypothesis on $\ell$.} $\{s^{\inv}_j\}$ represents
    $\Ga''\backslash\Ga$, $\tau^{\inv}$ is a $g$-twisted chain datum over $\cring$
    (Definition~\ref{def:tau}), and at \emph{cochain} level
    \[
      \inv^{\#}\circ\widetilde T\bigl[g,\{s_j\},\tau\bigr]
      \;=\;\widetilde T\bigl[g,\{s^{\inv}_j\},\tau^{\inv}\bigr]\circ\inv^{\#},
    \]
    where $\widetilde T[\,\cdot\,]$ is \eqref{eq:heckechain} built from the displayed
    datum.
  \item[(3b)] \PROVED{} \textbf{no hypothesis on $\ell$.} If moreover
    $\tau^{\inv}=\tau$, then $\inv^{\#}\circ\widetilde T=\widetilde T\circ\inv^{\#}$ on
    the nose, hence $\inv^{\ast}\Th_g=\Th_g\inv^{\ast}$ as cochain maps.
  \item[(3c)] \PROVED{} \textbf{under $6\in\cring^{\times}$} (hypothesis \ref{H3}; for
    $\cring$ of characteristic $\ell$, $\ell\ge5$). $[\inv^{\ast},\Tdc{g}]=0$ on
    $\Heq{q}(K,\Rel;L)$ for every $q$; equivalently $\inv^{\ast}$ and $\Amat{q}$ commute
    as matrices in the stored basis of $\Heq{q}$. The same then holds for $T_p$,
    $T_{p^{\nu}}$, $T_N$, $\diam$ and $U_\ell$, and for the primitive operator of
    \S\ref{sec:primfull}.
  \end{enumerate}
\end{enumerate}
\end{proposition}

\begin{proof}
\emph{(1).} Read Definition~\ref{def:model7rel} against \eqref{eq:iotaconj} and
Lemma~\ref{lem:iotamodel}. $(\{v_3\},\{v_0\},T)\mapsto(\{v_0\},\{v_3\},T^{-1})$,
$(\{v_5\},\{v_4\},T)\mapsto(\{v_4\},\{v_5\},T^{-1})$,
$(\{v_4\},\{v_5\},S)\mapsto(\{v_5\},\{v_4\},S^{-1})$ and
$(\{v_6\},\{v_6\},S)\mapsto(\{v_6\},\{v_6\},S^{-1})$ are the inverses of themselves. The
two edge rows likewise, and with sign $+1$ on \emph{both} sides by
Lemma~\ref{lem:iotamodel}, so no sign survives:
$(\{v_3,v_5\},\{v_0,v_4\},T)\mapsto(\{v_0,v_4\},\{v_3,v_5\},T^{-1})$ and
$(\{v_4,v_6\},\{v_5,v_6\},S)\mapsto(\{v_5,v_6\},\{v_4,v_6\},S^{-1})$. The elliptic
stabiliser rows are exchanged:
$(\{v_4\},\{v_4\},ST)\mapsto(\{v_5\},\{v_5\},(TS)^{-1})$, the inverse of
$(\{v_5\},\{v_5\},TS)$, and symmetrically. A central row $(\sigma,\sigma,-I)$ goes to
$(\invK\sigma,\invK\sigma,-I)$, whose two sides carry the same sign, hence to the
central row at $\abs{\invK\sigma}$. $\langle\Rel\rangle$ is closed under inversion, so
every image lies in it.

\emph{(2).} The four properties are proved for $\inv^{\ast}_{\mathrm{geo}}$; they pass
to $\inv^{\ast}=-\inv^{\ast}_{\mathrm{geo}}$ because the relations of $\Rel$ and $d$ are
linear, $(-\id)^2=\id$, and $\Psi$ is linear, so $\Psi\circ(-\inv^{\#}_{\mathrm{geo}})
=-\inv^{\ast}_{\mathrm{geo}}\circ\Psi$. (Equivalently: $-\rhon(\inv)$ satisfies the
same intertwining relation as $\rhon(\inv)$ below, so the computation may be read with
$-\rhon(\inv)$ throughout.)
\emph{Equivariance.} Let $f\in\Cq{q}$ and
$(\sigma_{\mathrm{lhs}},\sigma_{\mathrm{rhs}},\gamma)\in\Rel$. From
$\inv\,\conj_{\inv}(\gamma)=\gamma\,\inv$ in $\Mat_2(\Z)$ and
Proposition~\ref{prop:symconv}(2),
$\rhon(\inv)\rhon(\conj_{\inv}(\gamma))=\rhon(\gamma)\rhon(\inv)$; so by (1) and
Lemma~\ref{lem:eqgroupoid},
\[
  (\inv^{\ast}_{\mathrm{geo}}f)(\sigma_{\mathrm{lhs}})
  =\rhon(\inv)f(\invK\sigma_{\mathrm{lhs}})
  =\rhon(\inv)\rhon(\conj_{\inv}(\gamma))f(\invK\sigma_{\mathrm{rhs}})
  =\rhon(\gamma)\,(\inv^{\ast}_{\mathrm{geo}}f)(\sigma_{\mathrm{rhs}}).
\]
\emph{Commutation with $d$.} $\invK$ is simplicial, so $\invK\partial=\partial\invK$ on
chains, and $\rhon(\inv)$ is a constant matrix; hence
$(\inv^{\ast}_{\mathrm{geo}}(df))(\sigma)=\rhon(\inv)f(\partial\invK\sigma)
=\rhon(\inv)f(\invK\partial\sigma)=(d(\inv^{\ast}_{\mathrm{geo}}f))(\sigma)$.
\emph{Involutivity.} $\rhon(\inv)^2=I$ by \eqref{eq:rhoniota} and $\invK^2=\id$ on
chains. \emph{Transport.} By \ref{I1},
$\inv_{\Xbulk}(\lift\sigma)=\widetilde{\invK\sigma}$, so
$\Psi(\inv^{\#}_{\mathrm{geo}}\phi)(\sigma)=\rhon(\inv)\phi(\inv_{\Xbulk}\lift\sigma)
=\rhon(\inv)(\Psi\phi)(\invK\sigma)=\inv^{\ast}_{\mathrm{geo}}(\Psi\phi)(\sigma)$; that
$\inv^{\#}_{\mathrm{geo}}\phi$ is again $\Ga$-equivariant is the displayed computation
with arbitrary $\gamma\in\Ga$ in place of a relation element.

\emph{(3a).} $\conj_{\inv}(\Ga)=\Ga$ and $\conj_{\inv}(g)=g$ give
$\conj_{\inv}(\Ga'')=\conj_{\inv}(\Ga\cap g\Ga g^{-1})=\Ga''$, so
$\Ga=\bigsqcup_j\conj_{\inv}(\Ga''s_j)=\bigsqcup_j\Ga''s^{\inv}_j$. Properties (T2) and
(T3) of $\tau^{\inv}$ are immediate from
$\inv_{\Xbulk}\partial=\partial\inv_{\Xbulk}$, $\aug\inv_{\Xbulk}=\aug$ and
$\inv_{\Xbulk}^2=\id$. For (T1), let $h\in\Ga''$; then $\conj_{\inv}(h)\in\Ga''$ and
\[
  \tau^{\inv}(h\cdot x)
  =\inv_{\Xbulk}\tau\bigl(\conj_{\inv}(h)\cdot\inv_{\Xbulk}x\bigr)
  =\inv_{\Xbulk}\bigl((g^{-1}\conj_{\inv}(h)g)\cdot\tau(\inv_{\Xbulk}x)\bigr)
  =\conj_{\inv}\bigl(g^{-1}\conj_{\inv}(h)g\bigr)\cdot\tau^{\inv}(x)
  =(g^{-1}hg)\cdot\tau^{\inv}(x),
\]
the last step by $\conj_{\inv}(g)=g$ and $\conj_{\inv}^2=\id$. For the intertwining it
suffices to treat $\inv^{\#}_{\mathrm{geo}}$: $\inv^{\#}=-\inv^{\#}_{\mathrm{geo}}$
occurs once on each side of the asserted identity, so the signs cancel. Write
$\psi:=\inv^{\#}_{\mathrm{geo}}\phi$, so that $\phi(\inv_{\Xbulk}y)=\rhon(\inv)\psi(y)$, and note
$s_j\cdot\inv_{\Xbulk}x=\inv_{\Xbulk}(s^{\inv}_j\cdot x)$ and
$\tau\circ\inv_{\Xbulk}=\inv_{\Xbulk}\circ\tau^{\inv}$. Then by \eqref{eq:heckechain}
\[
  \bigl(\inv^{\#}(\widetilde T\phi)\bigr)(x)
  =\rhon(\inv)\sum_{j}\rhon(\delta_j)\,
     \phi\bigl(\inv_{\Xbulk}\tau^{\inv}(s^{\inv}_j\cdot x)\bigr)
  =\sum_{j}\rhon\bigl(\inv\delta_j\inv\bigr)\,
     \psi\bigl(\tau^{\inv}(s^{\inv}_j\cdot x)\bigr),
\]
and $\conj_{\inv}(\delta_j)=\conj_{\inv}(s_j)^{-1}\conj_{\inv}(g)=(s^{\inv}_j)^{-1}g$ is
exactly the coset element attached to $s^{\inv}_j$, so the right-hand side is
$\bigl(\widetilde T[g,\{s^{\inv}_j\},\tau^{\inv}]\,\psi\bigr)(x)$.

\emph{(3b).} If $\tau^{\inv}=\tau$ the two data of (3a) differ only in the
representatives, so $s^{\inv}_j=h_js_{j'}$ for some $h_j\in\Ga''$ and some permutation
$j\mapsto j'$ of $\{1,\dots,m\}$. With $\tau$ held fixed, the computation in the proof
of Theorem~\ref{thm:welldef}(a) shows that each summand of \eqref{eq:heckechain} is
unchanged under $s\mapsto hs$, $h\in\Ga''$; permuting the summands does not change the
sum; hence $\widetilde T[g,\{s^{\inv}_j\},\tau]=\widetilde T[g,\{s_j\},\tau]$ as cochain
maps, and (3a) becomes the asserted commutation.

\emph{(3c).} By (3a) the two sides of $[\inv^{\#},\widetilde T]$ are the operators of
two admissible data for the \emph{same} $g$, so by Theorem~\ref{thm:welldef}(a),(b) they
induce the same map on $\Heq{q}$ --- (b) being where $6\in\cring^{\times}$ enters, and
the only place. The extension to the operator family is
Definition~\ref{def:operators}: every double coset of determinant $p^{\nu}$ has a
diagonal representative $\mathrm{diag}(p^{a},p^{\nu-a})$, which $\conj_{\inv}$ fixes, so
each summand of \eqref{eq:fullop} is covered, $\diam=T_{\Ga pI\Ga}$ included;
$\conj_{\inv}$ moreover preserves both the determinant and the content of an integer
matrix, hence the primitive double coset $\Deprim_{p^{\nu}}$ of \S\ref{sec:primfull};
and $T_N$ is a $\Z$-combination of the $T_{p^{\nu}}$ by \eqref{eq:fullassembly} and
inherits the commutation.
\end{proof}

\begin{remark}[(3b) is the only route at $\ell\in\{2,3\}$; its hypothesis was \emph{false} for the canonical $\tau$ of v0.13.0--v1.0.0 and is \emph{true} for the canonical $\tau$ since v1.1.0]\label{rmk:iotatausym}
Theorem~\ref{thm:welldef}(b) is unavailable for $6\notin\cring^{\times}$, and so
therefore is Proposition~\ref{prop:iotacommutes}(3c): at $\ell\in\{2,3\}$ the
commutation $[\inv^{\ast},\Tdc{g}]=0$ \textbf{does not follow} from the general
argument. What is available there is (3b), whose hypothesis $\tau^{\inv}=\tau$ is a
finite comparison of two integer tables (Obligation~\ref{obl:iota}(iv)).

\emph{The finding (v0.13.0; run of 2026-09-16, \code{HeckeIotaTauSymmetryValidation}).}
\COMP{} On the canonical $\tau$ of Algorithm~\ref{alg:tau} \emph{as it stood before
v1.1.0} --- kept in the tree as \code{TauVertexPolicy::GeometricNearestLegacy} and re-run
by the same validation, so the number below is reproduced by every run --- at
$\ell\in\{2,3\}$, over $(p,\nu)\in\{(2,1),(3,1),(5,1),(7,1),(2,2),(3,2)\}$ and
$n\in\{0,1,2,3,4,10\}$, $3600$ of the $10764$ stored table slots satisfy
$\inv_{\Xbulk}\,\tau\,\inv_{\Xbulk}\neq\tau$ --- by degree, $456/3276$ at $q=0$,
$2340/5148$ at $q=1$, $804/2340$ at $q=2$ --- every comparison exact over $\Z$ (no table
used the mod-$\ell$ fallback of Remark~\ref{rmk:tauintegral}). A minimal witness,
sign-free because every vertex carries orientation sign $+1$ in
Lemma~\ref{lem:iotamodel}'s table: at $p=2$, $\nu=1$, $\ell=2$, $n=0$, in the slot
$(j,q,\sigma)=(1,0,\{v_0\})$ with $s_1=I$,
\[
  \tau(\lift{v_0})=\gamma\cdot\lift{v_1},
  \qquad
  (\inv_{\Xbulk}\,\tau\,\inv_{\Xbulk})(\lift{v_0})=\gamma\cdot\lift{v_2},
  \qquad \gamma=\smat{-1}{1}{0}{-1},
\]
the two right-hand sides being distinct as \emph{chains} (the validation canonicalises
every term through the public key the $\tau$-table uses). The mechanism is step~1 of
Algorithm~\ref{alg:tau}: with $g=\mathrm{diag}(1,2)$ and $s_1=I$ the reduced point of
$\adjg(g)\cdot\mathrm{pos}(v_0)=-1+4i$ is $4i$, \textbf{on the mirror}
$\mathrm{Re}\,z=0$, equidistant from $v_1$ and $v_2$, which $\invK$ exchanges; the
then tie-break --- ascending vertex index --- picked $v_1$. That negative was
\emph{certified}: $\tau^{\inv}$ satisfies (T1)--(T3) of Definition~\ref{def:tau} on
every stored slot, so it exhibited two admissible data for one $g$, exactly as
Proposition~\ref{prop:iotacommutes}(3a) predicts, and not a defect.

\emph{Why no tie-break of the tied set can repair it.} On the mirror the tied set
$S=\{v_1,v_2\}$ is $\invK$-stable without a fixed point, so $c(S)=\invK c(S)$ is
impossible for any choice $c(S)\in S$: the repair this remark proposed at v0.13.0
(``among tied vertices prefer the one of smaller index in the orbit of $\invK$'') cannot
work as stated, and an equivariant rule must leave the tied set. The only
$\invK$-fixed vertex of $K$ is $v_6=i$, and $i$ is the only vertex of $\Xbulk$ on the
mirror at all ($\gamma\cdot i\in i\Rl_{>0}$ forces $\gamma\in\{\pm I,\pm S\}$, and no
horocycle vertex lies on the axis), so a mirror point \emph{must} be sent to $v_6$,
off-nearest. The tied sets that actually occur on the validation's grid are three:
$\{v_1,v_2\}$ on the mirror ($288$ slots), and the $\inv$-paired $\{v_0,v_1\}$ /
$\{v_2,v_3\}$ off it ($72$ each); no other tie arises.

\emph{The repair, taken at v1.1.0 (owner decision of 2026-09-16).} Two changes to
Algorithm~\ref{alg:tau}, neither touching Definition~\ref{def:tau} or
Obligation~\ref{obl:tau}: the $\invK$-equivariant tie-break of step~(1b), and the
tree-routed templates of step~2. \COMP{} \textbf{Step~1 alone was not enough}: with the
tie-break changed and the templates as before, degree $0$ was clean ($0/3276$) while
degrees $1$ and $2$ still had $1848/5148$ and $1068/2340$ differing slots. The
breadth-first templates of the old step~2, in ascending neighbour order, are not
$\invK$-equivariant, and an edge with $\invK$-sign $-1$ was built in the \emph{reverse}
direction; a $1$-chain with prescribed boundary in $\Xbulk$ is unique only up to the
boundary of a $2$-chain, so this was a second, genuine obstruction, invisible in the
degree-$0$ witness. Tree routing removes it: in the tree $Y\subset\Xbulk$ of step~2 a
$1$-chain with prescribed boundary is \emph{unique}, so the edge chain depends only on
its two endpoints and is $\Ga$-equivariant, $\inv_{\Xbulk}$-equivariant and antisymmetric;
degree $2$ then follows, $\partial_2$ being injective on $C_2(\Xbulk)$.

\COMP{} \textbf{The verdict since v1.1.0.} $\inv_{\Xbulk}\,\tau\,\inv_{\Xbulk}=\tau$ on
every stored slot of every table: $0$ of $21528$ slots over the $144$ tables of the grid
($\ell\in\{2,3,5,691\}$), $0$ of $10764$ on the binding subgrid, every comparison exact
over $\Z$; Obligation~\ref{obl:iota}(iv) is \emph{gating} there. The legacy $\tau$
reproduces its $3600$. Hence Proposition~\ref{prop:iotacommutes}(3b) applies to the
canonical $\tau$: $\inv^{\#}\circ\widetilde T=\widetilde T\circ\inv^{\#}$ on the nose
and $\inv^{\ast}\Th_g=\Th_g\inv^{\ast}$ as cochain maps for every $\ell$ --- so at
$\ell\in\{2,3\}$, where (3c) is silent, $[\inv^{\ast},\Tdc{g}]=0$ on $\Heq{q}$ for every
operator whose $\tau$-table lies in the grid ($p\in\{2,3,5,7\}$, $\nu\le2$). Two limits
are part of the statement: the identity is verified table by table by a run, not proved
for every $(p,\nu)$ --- the structural reason above is an argument, not a theorem of this
document, and the obligation is the predicate a run discharges; and $\inv^{\ast}$ as a
cochain map is not implemented (clauses (ii), (iii), (v)), so the commutation is a
consequence of the proposition, not yet an observed matrix identity.

\emph{Two consequences of the tree routing.} (a) \textbf{The canonical $\tau$ vanishes
in degree $2$}: $\tau(\partial\lift t)$ is a closed loop of $Y$-paths, hence $0$, and the
unique fill is $0$; so $\Th^{(2)}_{\mathrm{flat}}=0$. Admissible --- (T1)--(T3) hold ---
and by Theorem~\ref{thm:welldef}(b) without effect on $\Heq{2}$ where
$6\in\cring^{\times}$; on the golden grid $\dim\Heq{2}=0$ at every point, $\ell\in\{2,3\}$
included. (b) \textbf{The re-baseline.} $2097$ of the $14049$ frozen reference values
changed --- every $\Th^{(q)}_{\mathrm{flat}}$, its rank, every $\Th^{(q)}_{\mathrm{eq}}$
and the three JSON records containing them --- and \emph{no} dimension, no $\Amat{q}$
characteristic polynomial, rank or trace, no decomposition invariant, no basis and no
projector, at any point, $\ell\in\{2,3\}$ included (Remark~\ref{rmk:indeplevel}; the
object-by-object record is
\file{docs/matrix-engine/rebaseline-2026-09-16-iota-symmetric-tau.md}). That nothing
moved on $\Heq{q}$ at $\ell\in\{2,3\}$ is an observation, not a theorem:
Theorem~\ref{thm:welldef}(b) promises it only for $6\in\cring^{\times}$.
\end{remark}

\subsection*{The $\pm$ decomposition, and what happens at $\ell=2$}

\noindent
Write $H^{q}_{\mathrm{eq},\pm}$ for the $\pm1$ eigenspaces of $\inv^{\ast}$ on
$\Heq{q}(K,\Rel;L)$, i.e.\ $H^{q}_{\mathrm{eq},\pm}:=\ker(\inv^{\ast}\mp\id)$, and
likewise $C^{q}_{\mathrm{eq},\pm}$, $Z^{q}_{\mathrm{eq},\pm}$,
$B^{q}_{\mathrm{eq},\pm}$; this is one symbol with one more index, not a new symbol
(\ref{N2}). Since $\inv^{\ast}=-\inv^{\ast}_{\mathrm{geo}}$
(Definition~\ref{def:iota}(iv), v1.2.0), these are the eigenspaces of the geometric
involution with the labels exchanged, $H^{q}_{\mathrm{eq},\pm}=\ker(\inv^{\ast}_{\mathrm{geo}}\pm\id)$;
every $\pm$ statement of this document, the live specification and the validation
output is in terms of $\inv^{\ast}$. A $\pm$ \emph{computed} before v1.2.0 --- the
v1.1.1 run record of Remark~\ref{rmk:iotasign} and
\file{docs/implementation/oq3\_iota\_sign\_verification.md} --- is the geometric one and
must be negated; a $\pm$ \emph{predicted} from owner decision OQ3 before v1.2.0 (the
split of Obligation~\ref{obl:iota}(v), family (C) of \S\ref{sec:oraclesnew}, the
$H^{1}_{\pm}$ statements of \file{hecke\_extensions.tex}) was written in the convention
v1.2.0 adopted and holds as printed.

\begin{proposition}[Eigenspace decomposition --- and it needs $2\in\cring^{\times}$]\label{prop:iotasplit}
\PROVED{} Assume \ref{H4} and \ref{I1}, let $\cring$ be a commutative ring and fix a
degree $q$.
\begin{enumerate}
\item \textbf{If $2\in\cring^{\times}$}, the maps
  $\mathrm{pr}_{\pm}:=\tfrac12(\id\pm\inv^{\ast})$ are complementary idempotents
  commuting with $d$, and
  \[
    \Heq{q}(K,\Rel;L)=H^{q}_{\mathrm{eq},+}\oplus H^{q}_{\mathrm{eq},-},
    \qquad H^{q}_{\mathrm{eq},\pm}=\im\,\mathrm{pr}_{\pm},
  \]
  with the same decomposition of $\Cq{q}$, $\Zq{q}$ and $\Bq{q}$. If in addition
  $6\in\cring^{\times}$, both summands are stable under every operator of
  \S\ref{sec:heckedatum} by Proposition~\ref{prop:iotacommutes}(3c), so $\Amat{q}$ is
  block diagonal in any basis adapted to the splitting --- which is what halves the
  matrices.
\item \textbf{If $2\notin\cring^{\times}$ the decomposition is false in general}, and
  what holds instead is a filtration. Put $\mathfrak{n}:=\inv^{\ast}-\id$; then
  $\mathfrak{n}^2=-2\mathfrak{n}$, so in characteristic $2$, $\mathfrak{n}^2=0$ and
  \[
    0\ \subseteq\ \im\mathfrak{n}\ \subseteq\ \ker\mathfrak{n}\ \subseteq\ \Heq{q},
  \]
  a filtration by submodules stable under every operator for which
  Proposition~\ref{prop:iotacommutes}(3b) or (3c) applies. Over $\cring=\F_2$ the module
  $\Heq{q}$ over $\F_2[\inv^{\ast}]\cong\F_2[x]/(x^2)$ ($x\mapsto\mathfrak{n}$) is a
  direct sum of $a$ copies of the trivial module and $b$ copies of the free module,
  with $b=\rank\mathfrak{n}$ and $a=\dim\ker\mathfrak{n}-b$; the two eigenspaces
  coincide, $H^{q}_{\mathrm{eq},+}=H^{q}_{\mathrm{eq},-}=\ker\mathfrak{n}$, and they
  exhaust $\Heq{q}$ if and only if $b=0$, i.e.\ $\inv^{\ast}=\id$. \textbf{At $\ell=2$
  no $\pm$ splitting of $\Htot$ and no halving of any matrix may be claimed.}
\end{enumerate}
\end{proposition}

\begin{proof}
(1) $(\inv^{\ast})^2=\id$ (Proposition~\ref{prop:iotacommutes}(2)) gives
$\mathrm{pr}_+\mathrm{pr}_-=\tfrac14\bigl(\id-(\inv^{\ast})^2\bigr)=0$,
$\mathrm{pr}_{\pm}^2=\mathrm{pr}_{\pm}$ and $\mathrm{pr}_++\mathrm{pr}_-=\id$; each
$\mathrm{pr}_{\pm}$ commutes with $d$ because $\inv^{\ast}$ does, and
$\im\,\mathrm{pr}_{\pm}=\ker(\inv^{\ast}\mp\id)$ since $\inv^{\ast}\mathrm{pr}_{\pm}
=\pm\mathrm{pr}_{\pm}$ and $\mathrm{pr}_{\pm}$ is the identity on that kernel. Stability
under the operators is immediate from commutation.

(2) $\mathfrak{n}^2=(\inv^{\ast})^2-2\inv^{\ast}+\id=2\id-2\inv^{\ast}
=-2\mathfrak{n}$, which vanishes in characteristic $2$; hence
$\im\mathfrak{n}\subseteq\ker\mathfrak{n}$, and both are submodules stable under
anything commuting with $\inv^{\ast}$. Over a field a square-zero endomorphism is, in a
suitable basis, a direct sum of $b=\rank\mathfrak{n}$ Jordan blocks $\smat0100$ and
$a=\dim\ker\mathfrak{n}-b$ zero blocks; the block $\smat0100$ is the free
$\F_2[x]/(x^2)$-module of rank $1$ and the zero block is the trivial module. Finally
$-1=+1$ in characteristic $2$, so $\ker(\inv^{\ast}-\id)=\ker(\inv^{\ast}+\id)$, and
this is all of $\Heq{q}$ exactly when $\mathfrak{n}=0$.
\end{proof}

\begin{remark}[The three degeneracies at $\ell=2$, and the odd-$n$ case]\label{rmk:iotaelltwo}
Three distinct things happen at $\ell=2$ and should not be conflated. First
$\rhon(\inv)=I$ by \eqref{eq:rhoniota}, and $-1=+1$, so
$\inv^{\ast}=\inv^{\ast}_{\mathrm{geo}}$ is the bare simplicial pullback and carries
neither coefficient information nor the sign of Definition~\ref{def:iota}(iv) --- the
v1.2.0 normalisation is invisible at $\ell=2$. Second the splitting fails
(Proposition~\ref{prop:iotasplit}(2)). Third, Remark~\ref{rmk:negI}'s odd-$n$ vanishing
argument fails as well --- it used $2\in\cring^{\times}$ --- so $\ell=2$ is precisely
the regime in which $\Heq{\bullet}$ can be non-zero for odd $n$ \emph{and} has no
$\inv$-splitting. For $\ell\neq2$ and $n$ odd the question does not arise:
$\Heq{\bullet}=0$ there and both eigenspaces are $0$. The additional degree bound and
truncation policy that $\ell=2$ forces on the \emph{boundary} are construction
item~\ref{open:ell2}, whose mathematics is \S\ref{sec:ell2} and whose policy was decided
by the owner on 2026-09-16 (\ref{oc1a}); nothing in this chapter depends on either, $\inv^{\ast}$ being
defined degreewise.
\end{remark}

\begin{remark}[$\inv$, the contragredient, and the pairing]\label{rmk:iotapairing}
\COMP\textsuperscript{r} for the mirror split of $\Hcpt{1}$ (v1.6.0, the closing
paragraph); the rest is commentary on tagged statements it cites.
At coefficient level everything is proved and explicit.
Proposition~\ref{prop:pairing}(4) with $\det\inv=-1$, and $\adjg\inv
=\mathrm{diag}(1,-1)=-\inv$ with $\rhon(\inv)$ diagonal, give
\[
  \bigl\langle\rhon(\inv)v,\ \rhodual(\inv)\varphi\bigr\rangle
  =(-1)^{\,n}\,\langle v,\varphi\rangle,
  \qquad
  \rhodual(\inv)=\rhon(\adjg\inv)^{\top}=(-1)^{\,n}\rhon(\inv).
\]
So $\inv$ acts on the contragredient module $\Dn$ of
Definition~\ref{def:contragredient} by the same diagonal matrix up to $(-1)^{n}$, the
comparison map $\Phi_n$ of Lemma~\ref{lem:symdual} intertwines the two actions (that
lemma is stated for every $g\in\Mat_2(\Z)$, $\inv$ included), and for $n$ even $\inv$ is
an isometry of the evaluation pairing on the nose. The sign of
Definition~\ref{def:iota}(iv) does not disturb this: it applies to \emph{both}
coefficient modules, $\Vcoef$ and $\Dn$, so at cochain level the two involutions carry
one factor $-1$ each and $\langle\inv^{\ast}x,\inv^{\ast}y\rangle
=(-1)(-1)(-1)^{n}\langle x,y\rangle=(-1)^{n}\langle x,y\rangle$ --- the same statement
as for the geometric pair, and every $\pm$ consequence below is invariant under negating
both involutions at once.

\textbf{What this does not establish.} Remark~\ref{rmk:selfadjoint} is explicit that
Proposition~\ref{prop:pairing} proves the coefficient-level half and says nothing about
cohomology; the cohomological pairings are Claim~\ref{cl:pd} of \S\ref{sec:six}, not
available here. Neither does this chapter construct the action of $\inv$ on
$H^{\bullet}(\partial)$ or on $\Hcpt{\bullet}$. That $\inv$ normalises
$\Gbdry=\langle T\rangle\times\langle-I\rangle$ is immediate from \eqref{eq:iotaconj},
so a boundary involution does exist; but its datum belongs to
\S\ref{sec:boundarydatum} and the cone to \S\ref{sec:cone}. \emph{Contingent} on both,
the consequence to expect is this. The coefficient factor $(-1)^{n}$ is not the whole
sign: $z\mapsto-\bar z$ reverses the orientation of $\Hyp$, so $\inv_{\Xbulk}$ reverses
the orientation of the surface and of the pair, and a cohomological pairing through a
fundamental class of the pair --- which is what the perfect pairings of
Claim~\ref{cl:pd}(1) are --- picks up a further factor $-1$. The equivariance to expect is
therefore $-(-1)^{n}$: if the pairing of Claim~\ref{cl:pd} satisfies
$\langle\inv^{\ast}x,\inv^{\ast}y\rangle=-(-1)^{n}\langle x,y\rangle$ and
$2\in\cring^{\times}$, then for $n$ even \emph{equal} signs annihilate each other ---
$\langle x,y\rangle=\varepsilon^{2}\langle x,y\rangle=\langle\inv^{\ast}x,\inv^{\ast}y\rangle
=-\langle x,y\rangle$ for $x\in H^{t}_{c,\varepsilon}$, $y\in H^{2-t}_{\mathrm{eq},\varepsilon}$ ---
so duality would restrict to a perfect pairing of each sign with the \emph{opposite}
one, as for modular symbols, where complex conjugation is orientation-reversing and
$H^{1}_{+}$ is paired with $H^{1}_{-}$ \STD{\cite[Ch.~2]{cremona}}. The dimensions this
document predicts agree with that and not with self-pairing: Obligation~\ref{obl:iota}(v)
puts $(\dim H^{1}_{\mathrm{eq},+},\dim H^{1}_{\mathrm{eq},-})=(d_{\,n+2},d_{\,n+2}-1)$,
and since $H^{0}(\partial;\Mn)$ is one-dimensional and odd under the boundary involution
(v1.2.0), $H^{1}_{c}$ carries the mirror split $(d_{\,n+2}-1,d_{\,n+2})$, so a perfect
pairing must match $+$ with $-$. (Until v1.3.2 this remark predicted the self-pairing,
having omitted the orientation factor.) That sentence is a prediction about a pairing
this document does not have --- Claim~\ref{cl:pd} is a stub of \S\ref{sec:six} --- and is
asserted here as nothing more.

\textbf{Its dimension half is no longer contingent} (v1.6.0; hecke-extensions slot-4
WP-3, \file{docs/implementation/slot4/WP-3\_report.md} findings 4.1 and 5). The two
objects the paragraph above called unavailable now exist: $\inv^{\ast}_{\partial}$ on the
boundary model and $\inv_{c}=\inv^{\ast}\oplus\inv^{\ast}_{\partial}$ on $\Cone{\rres}$,
with $\rres^{q}\inv^{\ast}=\inv^{\ast}_{\partial}\rres^{q}$ and
$\Dcone\inv_{c}=\inv_{c}\Dcone$ \emph{asserted} (\code{HeckeInvolutionValidation}, tier
\code{standard}, run 2026-09-21: the $77$-cell $\inv$ grid and a $15$-cell cone grid
$\ell\in\{2,3,5,7,691\}\times n\in\{0,2,10\}$). \textbf{Recorded outcome}: on the
\textbf{eight} cells of that cone grid with $\ell\in\{3,5,7,691\}$ and $n\in\{2,10\}$ ---
among them $(n,\ell)=(10,691)$, where $\dim\Hcpt{1}=3$ splits $(1,2)$ against
$\dim\Heq{1}=3$'s $(2,1)$; at $n=2$ the split is $(0,1)$ against $\Heq{1}$'s $(1,0)$ and at
$n=10$ it is $(1,2)$ against $(2,1)$ --- $\Hcpt{1}$ carries \emph{exactly} the mirror
of the split Obligation~\ref{obl:iota}(v) puts on $\Heq{1}$, which is what the paragraph
above predicts and the only prediction this document makes about $\inv$ on the cone.
(v1.6.1: v1.6.0 wrote ``thirteen cells --- $(10,691)$ and the twelve cells of the cone grid
with $\ell\in\{3,5,7,691\}$'', which counts $(10,691)$ twice, that cell being one of the
twelve, and folds in the four at $n=0$, where $\Hcpt{1}=0$ and
Obligation~\ref{obl:iota}(v) makes no prediction at all, its $d_k$ being defined for even
$k\ge4$ and $n=0$ being $k=2$. At $\ell=3$ that obligation's dimension statement is also
outside its own hypothesis $\ell\ge5$, so the four cells there compare a measured split
with a measured split; they are kept in the eight because the mirror relation is what is
being recorded, not the obligation's prediction.) At
$\ell=2$ the recorded pairs are $(k,k)$, Proposition~\ref{prop:iotasplit}(2) removing the
split and not the involution, and they are correctly not asserted to sum to
$\dim\Hcpt{q}$. Rule~\ref{A6}: this is a record, it licenses nothing, no code path may
branch on it, and \COMP{} never upgrades to \PROVED{} (\S\ref{sec:status}) --- the
\emph{pairing} statement above stays contingent on Claim~\ref{cl:pd}, and what would
make $\inv_{c}$ commute with $\Acone{q}$ is stated, and left open, at the end of
\S\ref{sec:coneaction}.
\end{remark}

\begin{remark}[The record field \code{iota\_sign}: what is chosen and what is not]\label{rmk:iotasign}
Owner decision OQ3 --- proposed 2026-09-15 as an assertion, \textbf{resolved 2026-09-16
as a definition} --- fixes the convention ``the Eisenstein class lies in
$H^{1}_{\mathrm{eq},+}$'', recorded as
\code{"iota\_sign": "eisenstein\_in\_H1\_plus"} in the extension record. Its status has
to be stated exactly, because Definition~\ref{def:iota} leaves \emph{no} sign free.
\begin{itemize}
\item The only ambiguity in the \emph{geometric} datum is the lift of $z\mapsto-\bar z$
  to $\GL_2(\Z)$ --- $\inv$ versus $-\inv=\mathrm{diag}(1,-1)$, both inducing the same
  map on $\Hyp$ because $-I$ acts trivially there. The two give
  $\rhon(-\inv)=(-1)^{n}\rhon(\inv)$, hence the same $\inv^{\ast}_{\mathrm{geo}}$ for $n$
  even and opposite $\inv^{\ast}_{\mathrm{geo}}$ for $n$ odd. For $\ell\neq2$ and $n$ odd
  $\Heq{\bullet}=0$ (Remark~\ref{rmk:negI}), so on every group this document computes
  with $6\in\cring^{\times}$ the two lifts induce the same involution, and \emph{no}
  choice of lift can move the Eisenstein class from one eigenspace to the other.
\item \textbf{The sign was therefore an assertion about the geometric involution, and
  it was checked.} At v1.1.1 \code{HeckeIotaSignValidation} built
  $\inv^{\ast}_{\mathrm{geo}}$ from Definition~\ref{def:iota} alone and found the
  Eisenstein class in $\ker(\inv^{\ast}_{\mathrm{geo}}+\id)$ on all $74$ signed cells of
  a $77$-cell grid, by two independent routes that agree wherever both apply, with
  $\inv^{\ast}_{\mathrm{geo}}=+\id$ on $H^{0}(\Ga;\Fl)$ pinning the global sign
  (\file{docs/implementation/oq3\_iota\_sign\_verification.md}; that report's $\pm$ are
  geometric). The mechanism is elementary: $\inv$ reverses the boundary circle, and the
  boundary class $[y^{n}]$ is $\rhon(\inv)$-fixed, so $\inv^{\ast}_{\mathrm{geo}}=-1$ on
  $H^{1}(\partial)$ for $\ell>n$; restriction intertwines the two involutions and the
  Eisenstein line is boundary-visible.
\item \textbf{Resolution (v1.2.0): the definition was negated, not the record.}
  Definition~\ref{def:iota}(iv) now reads $\inv^{\ast}:=-\inv^{\ast}_{\mathrm{geo}}$, so
  \code{iota\_sign = eisenstein\_in\_H1\_plus} is \textbf{true by definition}, and the
  field's meaning changes accordingly: it \emph{names the normalisation}
  $\inv^{\ast}=-\inv^{\ast}_{\mathrm{geo}}$ relative to the geometric involution. A third
  party reproduces any $\pm$ statement of a record by that one formula. The choice was
  an owner decision on a normative statement with an $X$ bump (\ref{V2}) and a
  changelog line, exactly as this remark had required; the alternatives and their costs
  --- keeping the geometric involution and renaming its eigenspaces, or a
  degree-parity twist $(-1)^{q}$, which anticommutes with $d$ --- are weighed in
  \file{docs/implementation/oq3\_resolution.md}. \textbf{It was not flipped at a call
  site}, which is the failure mode that a recorded convention field invites, and no
  pipeline value moved: $\inv^{\ast}$ occurs in no golden object.
\item What is \emph{not} chosen: that the geometric sign is uniform across cells.
  ``$+$ is the side of the Eisenstein classes'' is well-posed only because every cell
  gives the same geometric sign, and that is a fact about the model, not a convention.
  Obligation~\ref{obl:iota}(v) is now the gate on it: a cell whose Eisenstein class lies
  in $H^{1}_{\mathrm{eq},-}$ would mean the definition is ill-posed, and the repair
  would again be a normative one, never a per-cell sign.
\end{itemize}
\end{remark}

\begin{obligation}[Sign convention and runtime assertions for $\inv$]\label{obl:iota}
Assume \ref{H4} and \ref{I1}. An implementation of $\inv$ emits the signed block
permutation of Definition~\ref{def:iota}(iv) in every degree, in the flat coordinates
\eqref{eq:flatindex}, and must discharge the following. Each clause names the stored
data it reads, the predicate, the ring it is evaluated over, and the action on failure.
\begin{enumerate}
\item[(i)] \textbf{Simplicial automorphism.} The stored vertex permutation is an
  involution carrying $K$ to $K$, and the recomputed orientation sign of every simplex
  agrees with Lemma~\ref{lem:iotamodel}'s table. Over $\Z$; a failure \textbf{aborts}.
\item[(ii)] \textbf{Local-system automorphism.} For every row
  $(\sigma_{\mathrm{lhs}},\sigma_{\mathrm{rhs}},\gamma)$ of $\Rel$, the triple
  $(\invK\sigma_{\mathrm{lhs}},\invK\sigma_{\mathrm{rhs}},\inv\gamma\inv)$ is the
  inverse of a row of $\Rel$ --- compared as exact $\Mat_2(\Z)$ data together with
  stored simplex indices and their signs --- and
  $\rhon(\inv)\rhon(\gamma)=\rhon(\inv\gamma\inv)\rhon(\inv)$ entrywise over $\cring$. A
  failure \textbf{aborts} (Proposition~\ref{prop:iotacommutes}(1)).

  \smallskip
  \emph{The two halves, and the complex that has only one of them} (v1.7.0;
  \file{docs/implementation/slot4/WP-3\_report.md} finding~3, the rider it asked for).
  The clause has an \emph{exact} half --- the triple compared as $\Mat_2(\Z)$ data --- and
  a \emph{fibre} half --- the identity in $\rhon$. The fibre half names data every local
  system carries and is asserted \textbf{unconditionally, always}: it is what makes
  $\inv_{C}$ preserve $\Ceq$, and no implementation is excused from it. The exact half
  names a row's $\gamma$, which a local system presented by its monodromy matrices alone
  need not store --- the shipped $\Rel$ does not, the exact rows living in the
  model-specific datum --- and it is asserted on \textbf{every row for which the complex
  stores that $\gamma$}. This is a statement about the data the clause reads, not a
  licence: an implementation whose complex carries the exact rows \textbf{may not} skip
  the exact half, and every implementation \textbf{records} which rows it ran it on, so
  that a record says whether the exact comparison happened rather than leaving a reader
  to assume it. A complex carrying no exact group data discharges the fibre half and
  records the exact half as not evaluated; it does not report clause (ii) as passed in
  the exact sense. Rule~\ref{A6} governs the record as usual.
\item[(iii)] \textbf{Cochain involution.} $\inv^{\ast}\inv^{\ast}=\id$ and
  $\inv^{\ast}d=d\inv^{\ast}$ as matrices over $\cring$ in every degree, and
  $\inv^{\ast}$ carries the stored bases of $\Cq{q}$, $\Zq{q}$, $\Bq{q}$ into their own
  spans. A failure \textbf{aborts} (Proposition~\ref{prop:iotacommutes}(2)).
\item[(iv)] \textbf{$\tau$-symmetry --- a diagnostic, never branched on.} Record whether
  $\inv_{\Xbulk}\,\tau\,\inv_{\Xbulk}=\tau$ holds on every stored table entry, compared
  over $\Z$. This is the hypothesis of Proposition~\ref{prop:iotacommutes}(3b) and the
  \emph{only} route to $[\inv^{\ast},\Tdc{g}]=0$ when $6\notin\cring^{\times}$. It is
  \textbf{recorded}, and \textbf{a true value licenses nothing beyond (3b)} --- in
  particular it may not be used to skip clause (v) (rule~\ref{A6};
  Remark~\ref{rmk:iotatausym}). On an arbitrary admissible $\tau$ a false value is not a
  failure. \textbf{On the canonical $\tau$ of Algorithm~\ref{alg:tau} it is, since
  v1.1.0}: steps (1b) and 2 of that algorithm are specified so that the identity holds,
  so a differing slot is a defect of the construction and the validation fails. This is
  not a branch on the diagnostic --- nothing downstream reads the verdict --- but a
  regression gate on the construction.
  \COMP{} \textbf{Discharged --- negatively at v0.13.0 on the then-canonical $\tau$,
  positively since v1.1.0 on the canonical $\tau$}, by
  \code{HeckeIotaTauSymmetryValidation} (runs of 2026-09-16): on the canonical $\tau$
  the recorded value is \emph{true} on every stored slot of every table of the grid
  ($0$ of $21528$ slots over $144$ tables, $\ell\in\{2,3,5,691\}$), every comparison
  exact over $\Z$; the pre-v1.1.0 $\tau$, kept as the legacy policy and re-run alongside,
  differs on $3600$ of $10764$ binding slots (Remark~\ref{rmk:iotatausym} for the grid,
  the witness and the mechanism). The clause remains in force as the predicate a
  \emph{re-run} on any other $\tau$, on any other vertex policy or on a refined
  triangulation must evaluate and record. The validation fails on a differing slot of
  the canonical $\tau$, on clause (i), or if $\tau^{\inv}$ is not an admissible datum
  (Proposition~\ref{prop:iotacommutes}(3a)).
\item[(v)] \textbf{Commutation, the sign, and the $\pm$ dimensions.} For
  $6\in\cring^{\times}$: $\inv^{\ast}\Amat{q}=\Amat{q}\inv^{\ast}$ over $\cring$ in every
  degree and for every computed operator; a failure \textbf{aborts}
  (Proposition~\ref{prop:iotacommutes}(3c)). \textbf{The sign --- a
  definition-consistency oracle, aborting since v1.2.0.} Definition~\ref{def:iota}(iv)
  defines $\inv^{\ast}=-\inv^{\ast}_{\mathrm{geo}}$ so that ``$+$'' is the eigenspace of
  the Eisenstein classes; the implementation must exhibit that this is well posed: on
  every cell $(n,\ell)$ with $\ell\ge5$, $n$ even, on which the Eisenstein class of
  Obligation~\ref{obl:oracle} can be located --- as the joint $T_p$-eigenline of
  eigenvalue $1+p^{\,n+1}$ where that eigenspace is a line, or as the class visible
  under restriction to the cusp where $\rank\rres^{1}=1$ --- it lies in
  $H^{1}_{\mathrm{eq},+}$, the two routes agree wherever both apply, and
  $\inv^{\ast}=-\inv^{\ast}_{\mathrm{geo}}$ entrywise in every degree. The degree-$0$
  instance is exact and cheap and is what pins the global sign: $\inv^{\ast}$ on
  $\Heq{0}=\VR^{\Ga}$ is $-\rhon(\inv)$ restricted to the invariants, which for $n=0$ is
  $-\id$ (the geometric $\inv^{\ast}_{\mathrm{geo}}$, which passes every other clause,
  gives $+\id$ there). A failure \textbf{aborts}, and what it would mean is not a wrong
  sign at a call site but a \emph{non-uniform} geometric sign across cells, i.e.\
  Definition~\ref{def:iota}(iv)'s ``$+$'' ill-posed --- a defect of the definition, to be
  repaired normatively (Remark~\ref{rmk:iotasign}). \textbf{The dimensions --- recorded
  rather than aborting:} for $\ell\ge5$ off the Dickson locus of Remark~\ref{rmk:dickson}
  and $n$ even,
  \[
    \dim H^{1}_{\mathrm{eq},+}=d_{\,n+2},\qquad
    \dim H^{1}_{\mathrm{eq},-}=d_{\,n+2}-1,
  \]
  where for even $k\ge4$ the integer $d_k=\lfloor k/12\rfloor+1$ if
  $k\not\equiv2\pmod{12}$ and $d_k=\lfloor k/12\rfloor$ otherwise is the dimension of
  the space of level-$1$ modular forms of weight $k$, its subspace of cusp forms having
  dimension $d_k-1$ \STD{\cite[Ch.~VII, \S3]{serre_cours}}. Given the gated Eisenstein
  sign this is the statement that the cuspidal part $\ker\rres^{1}$ splits evenly, the
  mod-$\ell$ Eichler--Shimura prediction discussed after this obligation, which is why it
  records.

  \COMP{} \textbf{Checked 2026-09-16 (v1.2.0): every predicate of this clause holds.}
  \code{HeckeIotaSignValidation} (tier \code{standard}, $1.4$\,s, $4\,104\,267$ exact
  assertions, $0$ gating failures) builds both $\inv^{\ast}_{\mathrm{geo}}$ and
  $\inv^{\ast}$ from Definition~\ref{def:iota} alone --- the live $\rhon$ of
  $\mathrm{diag}(-1,1)$, $\invK$ and Lemma~\ref{lem:iotamodel}'s signs re-derived from
  the shipped exact coordinates, the global sign of \eqref{eq:iotacochain} --- and
  evaluates this clause over $n\in\{2,4,\dots,30\}$ even ($11$ values) and
  $\ell\in\{5,7,11,13,101,134217757,2^{31}-1\}$: $77$ cells. Commutation
  $\inv^{\ast}\Amat{q}=\Amat{q}\inv^{\ast}$ for $p\in\{2,3\}$ in every degree,
  $(\inv^{\ast})^{2}=\id$ on $\Heq{q}$, $\inv^{\ast}=-\inv^{\ast}_{\mathrm{geo}}$ on
  cochains and on $\Heq{q}$, and the degree-$0$ anchor $\inv^{\ast}=-1$ on
  $\Heq{0}=H^{0}(\Ga;\Fl)$ at $n=0$ for $\ell\in\{5,13,134217757\}$; the two lifts
  $\inv$ and $-\inv$ of Remark~\ref{rmk:iotasign} induce the same involution on all $77$
  cells, as $n$ even predicts. On every one of the $74$ cells on which a sign is defined
  ($74$ by the $T_p$-eigenline, $53$ by the cusp, $53$ by both, agreeing) the Eisenstein
  class lies in $H^{1}_{\mathrm{eq},+}$; the three cells without a sign are the Dickson
  cells $(20,5)$, $(24,5)$, $(30,5)$, where the joint eigenspace has dimension $2$, $2$,
  $3$ (Fermat: $1+p^{\,n+1}\equiv1+p^{\,n+1-(\ell-1)}$) and $\rank\rres^{1}\neq1$. The
  recorded split is $(d_{\,n+2},d_{\,n+2}-1)$ on all $66$ off-Dickson cells, with
  $\ker\rres^{1}$ splitting evenly wherever $\ell>n$. \emph{Under the geometric
  involution every sign is the opposite}: Eisenstein in
  $\ker(\inv^{\ast}_{\mathrm{geo}}+\id)$ on $74/74$, the mirror split $(d_{\,n+2}-1,d_{\,n+2})$
  on $66/66$, $\inv^{\ast}_{\mathrm{geo}}=+1$ on $H^{0}(\Ga;\Fl)$ --- the v1.1.1 finding
  that motivated the decision, recorded by the same run
  (\code{IOTA\_SIGN eisenstein=plus (geometric=minus) ... h0\_anchor=-1 (geometric=+1)}).
  Evidence: \file{docs/implementation/oq3\_resolution.md} (the decision and its cost)
  and \file{docs/implementation/oq3\_iota\_sign\_verification.md} (the $77$-cell table,
  whose $\pm$ are geometric).
\end{enumerate}
The classical statement behind the $\Heq{1}$ oracle is the Eichler--Shimura
decomposition of $H^1(\Ga;\Sym^n\Cx^2)$ into a space of weight-$(n+2)$ modular forms and
a conjugate space of cusp forms
\STD{\cite[Ch.~8]{shimura}; \cite[Ch.~2]{cremona} for the involution and the $\pm$ parts
in the modular-symbol realisation}: each cuspidal eigenform contributes a
$2$-dimensional $\inv^{\ast}$-stable subspace on which $\inv^{\ast}$ has eigenvalues
$+1$ and $-1$, and the Eisenstein line contributes the one further dimension, which
Definition~\ref{def:iota}(iv) places in $H^{1}_{\mathrm{eq},+}$ by construction. That the
mod-$\ell$ reduction inherits the two dimensions is a \emph{prediction}, not a theorem of
this document (\CONJ{} as a citation: the Dickson locus of Remark~\ref{rmk:dickson} is
exactly where reduction is known to misbehave) --- which is why (v)'s dimension
statement records rather than aborts. \COMP{} The total it must
reproduce is already observable without $\inv$: at $\ell=134217757$ the shipped pipeline
gives $\dim\Heq{1}=1,1,3,5,7$ for $n=2,4,10,22,34$ and
$\dim\Heq{1}=0$ for $n=0$, matching $2d_{n+2}-1$ (respectively $0$) in every case
(runs of 2026-09-15, \code{--bench cohom <n> 134217757}).
\end{obligation}

\paragraph{Code map.}
\textbf{All five clauses are implemented, by two validations.}
\code{Hecke/HeckeIotaTauSymmetryValidation.\{h,cpp\}} (added 2026-09-16, registered as
\code{HeckeIotaTauSymmetryValidation}, tier \code{fast}, $0.5$\,s, $64699$
exact-integer assertions over $432$ $\tau$-tables --- the canonical, the legacy and the
\code{FixedV6} policy at each of $144$ grid points) builds $\invK$ and $\inv_{\Xbulk}$
from Definition~\ref{def:iota}(i),(iii) alone --- the vertex permutation and the sign of
every simplex are \emph{recomputed}, not read from a stored table --- and discharges
\textbf{Obligation~\ref{obl:iota}(i)}, which passes, and \textbf{Obligation~\ref{obl:iota}(iv)},
whose recorded verdict is positive on the canonical $\tau$ --- gating there since v1.1.0
--- and negative on the legacy $\tau$ (Remark~\ref{rmk:iotatausym}). It also verifies
Lemma~\ref{lem:iotamodel}'s unstated hypothesis $\mathrm{Re}\,v_1=-\mathrm{Re}\,v_2$
exactly in the shipped coordinates ($v_1=-\tfrac16+2i$, $v_2=\tfrac16+2i$), and
certifies, as a gate on its own construction, that $\tau^{\inv}$ satisfies (T1)--(T3) of
Definition~\ref{def:tau} --- i.e.\ Proposition~\ref{prop:iotacommutes}(3a) on the
shipped tables. \code{Hecke/HeckeIotaSignValidation.\{h,cpp\}} (added 2026-09-16,
registered as \code{HeckeIotaSignValidation}, tier \code{standard}, $1.4$\,s,
$4\,104\,267$ assertions, $0$ gating failures) covers the remaining three:
it builds both $\inv^{\ast}_{\mathrm{geo}}$ and $\inv^{\ast}=-\inv^{\ast}_{\mathrm{geo}}$
from Definition~\ref{def:iota} alone as explicit $(N_q d_V)\times(N_q d_V)$ matrices
over $\cring$ --- the block $\pm\varepsilon_{\sigma}\rhon(\inv)$ at block position
$(\sigma,\abs{\invK\sigma})$ in the flat coordinates \eqref{eq:flatindex}, the two
asserted to be negatives of each other --- and discharges
\textbf{Obligation~\ref{obl:iota}(ii)} (all $31$ relation rows of the model, the $23$
central $-I$ rows included, entrywise over $\cring$),
\textbf{Obligation~\ref{obl:iota}(iii)} ($\inv^{\ast}\inv^{\ast}=\id$ and
$\inv^{\ast}d=d\inv^{\ast}$ as full matrices in every degree, and the stored $\Cq{q}$,
$\Zq{q}$, $\Bq{q}$ carried into their own spans) and
\textbf{Obligation~\ref{obl:iota}(v)} --- commutation, the gating Eisenstein sign
(``$+$'' on every signed cell, by both routes) and the degree-$0$ anchor
$\inv^{\ast}=-1$, all passing, with the geometric sign recorded alongside; see that
clause's \COMP{} record and Remark~\ref{rmk:iotasign}. It reaches $\Heq{\bullet}$
through the stored $P_q$ and
$H_{\mathrm{eq}}$ of Definition~\ref{def:storedbases}, and the cusp through the
shipped \code{EquivariantCochainComplexMorphismOnK} restriction, on which the same
construction is re-run and \eqref{eq:iotacochain}'s naturality
$\rres\,\inv^{\ast}=\inv^{\ast}_{\partial}\,\rres$ is asserted. \textbf{Code map (v1.6.2, the slot-4 close; this replaces the 2026-09-16 census, which
found no \code{HeckeInvolution} translation unit and is now spent).} Three of the $108$
registered validations (v1.7.4 count; $106$ at v1.7.3, $100$ at v1.7.2, $94$ at v1.6.2, $89$ at v1.5.2, $80$ at v1.4.2,
$77$ at v1.3.3, $68$ when that census was taken) touch this chapter:
\code{HeckeIotaTauSymmetryValidation} and \code{HeckeIotaSignValidation} as before, and
\code{HeckeInvolutionValidation} since hecke-extensions slot-4 WP-3. The
\emph{production} objects are \file{Hecke/HeckeInvolution} --- the morphism specified
below, whose \code{compute()} discharges Obligation~\ref{obl:iota}(i),(ii),(iii)
\emph{blockwise and exactly over $\Z$} and aborts on failure --- and
\file{HeckeExt/HeckeSignSplit}, which carries clause (v) and the $\pm$ decomposition.
The acceptance is \emph{entrywise against an independent construction}: WP-3 exports
\code{HeckeIotaSignValidation}'s builders as \code{hecke::iota\_oracle}
value-preservingly, and the production object reproduces that validation's reference
line field for field. Remark~\ref{rmk:iotapairing}'s closing paragraph and the closing
paragraph of \S\ref{sec:coneaction} record the same.
The consequence for the rest of the
document is unchanged: every \PROVED{} tag above other than
Proposition~\ref{prop:iotacommutes}(1),(2),(3a) and Lemma~\ref{lem:iotamodel} rests on
its printed proof alone. The \emph{specification} recorded here for a \emph{production}
$\inv^{\ast}$ --- as opposed to the validation's own matrix, which exists to be checked
and not to be depended on --- stands unchanged as a \emph{specification} (E3 of the
extension programme, gated by amendment C-A1). Since v1.6.2 it has been \emph{built} to
this specification, so what follows is what the shipped \file{Hecke/HeckeInvolution}
satisfies and is no longer a description of something absent:
$\inv^{\ast}$ is an \code{EquivariantCochainComplexMorphismOnK} built through the
general constructor, its \code{FiberBlockTable} holding exactly one
\code{PerSimplexBlock} per simplex --- source simplex $\abs{\invK\sigma}$, block
$\mp\rhon(\inv)$, the sign of Lemma~\ref{lem:iotamodel} \emph{negated} per
Definition~\ref{def:iota}(iv) and folded into the matrix ---
which is the degenerate single-entry case of the per-simplex sparse block rows that the
repair of Remark~\ref{rmk:nogo} introduced, so no new datum shape is required. That
class's documented contravariant naming applies: its \code{domain} is the \emph{target}
of the induced map on $H^{\bullet}$. Clause (iv) of Obligation~\ref{obl:iota} reads
\code{hecke::HeckeTauTable}'s stored chains and compares them with their
$\inv_{\Xbulk}$-images, so it needs $\inv$ on the simplices of $\Xbulk$ carried by
\code{hecke::SL2ZTessellationTables}, not only on those of $K$.

\paragraph{Formalization note.}
Definition~\ref{def:iota}(i)--(ii) is a $2\times2$ integer matrix and a closed-form
diagonal matrix, and (iii)--(iv) are a vertex permutation with a computable sign; so
Lemma~\ref{lem:iotamodel}, Proposition~\ref{prop:iotacommutes}(1),(2) and clauses
(i)--(iii) of Obligation~\ref{obl:iota} are \code{decidable} predicates on finite data
and are the natural first targets of this chapter.
Proposition~\ref{prop:iotacommutes}(3a) is a finite calculation with the $\tau$-table
and needs no hypothesis on $\ell$; (3c) is the only statement here that invokes an
abstract input --- Theorem~\ref{thm:welldef}(b), and through it the fundamental lemma of
relative homological algebra --- and it inherits that theorem's hypothesis
$6\in\cring^{\times}$ verbatim. Proposition~\ref{prop:iotasplit}(2)'s classification is
the Jordan form of a square-zero endomorphism over a field. The statement to be careful
with is (1): its $\tfrac12$ must be carried as an
$\mathrm{Invertible}\,(2:\cring)$ instance argument and not as a field hypothesis,
since $\cring=\Z/\ell^{e}$ is in scope.

%======================================================================
\part*{Part III --- The integral Hecke datum}
\addcontentsline{toc}{section}{Part III --- The integral Hecke datum}
%======================================================================

%----------------------------------------------------------------------
\section{The integral Hecke datum for $\GL(2)/\SL_2(\Z)$}\label{sec:heckedatum}
%----------------------------------------------------------------------
% SOURCE:  hecke_action.tex:384-569 (RH)
% NOTATION: rows nu, T_N, "t,j,s (running index)", sigma, Mat_2(Z), cring, Rel
%           of sec:notation-crosswalk applied.  In particular the coset index is
%           j, never i (C15); the prime-power exponent is nu, never r (C11).
% ERRATA:  E1 (definitional half), E6, E16 (App. A)

\paragraph{What this chapter fixes.}
The group-theoretic datum out of which every Hecke operator of this document is
built: the element $g=\smat{1}{0}{0}{N}$ with $N=p^{\nu}$, the two congruence
subgroups it cuts out of $\Ga$, the coset spaces of the correspondence with the
\emph{explicit} representatives, class map and cocycle the implementation uses, and
the operators themselves --- $T_p$, the \textbf{full} $T_{p^{\nu}}$, the
\textbf{primitive} $T_{\Deprim_{p^{\nu}}}$, the composite $T_N$, the central $\diam$
and the bad-prime $U_\ell$ --- with the relations among them. Nothing here says what
an operator \emph{is} as a map: that is \S\ref{sec:chainlevel}. This chapter says
which double cosets are fed to that construction, with which auxiliary data, and what
is asserted of the output.

\paragraph{Standing hypotheses.}
\ref{H1}--\ref{H4} of \S\ref{sec:conventions}. Per rule~\ref{A1} each is repeated in
every statement that uses it; in particular \ref{H3} ($6\in\cring^{\times}$, i.e.\
$\ell\ge5$ for $\cring$ of characteristic $\ell$) is \emph{not} in force by default and
is written out wherever it is needed --- which is every statement of
\S\ref{sec:relations} and no statement of \S\S\ref{sec:correspondencedatum}--\ref{sec:operatordefs}.

\paragraph{Notation used throughout.} The coset index is $j\in\{1,\dots,m\}$ (never
$i$: \S\ref{sec:notation-collisions}, C15) and the prime-power exponent is $\nu$ (never
$r$: C11). The composite-index operator is $T_N$ with $N\in\Z_{>0}$, never $T_n$
(Appendix~\ref{app:errata}, entry~\ref{err:e16}).

\subsection{Primitive and full: the two operators at a prime power}
\label{sec:primfull}

The single most consequential naming decision of this chapter is made first, because
the superseded contract stated two numbered prescriptions that conflated the two
objects and were therefore both wrong by $p^{\,n}I$
(Appendix~\ref{app:errata}, entry~\ref{err:e1}).

\begin{contractbox}
At a prime power $p^{\nu}$ with $\nu\ge2$ there are \textbf{two different operators},
and this document names them differently and never abbreviates either away:
\begin{itemize}
\item the \textbf{primitive double-coset operator}
  $T_{\Deprim_{p^{\nu}}}:=\Tdc{\mathrm{diag}(1,p^{\nu})}$, attached to the \emph{single}
  double coset $\Deprim_{p^{\nu}}=\Ga\,\mathrm{diag}(1,p^{\nu})\,\Ga$ of the
  \emph{primitive} determinant-$p^{\nu}$ matrices, with
  $m=p^{\nu}+p^{\nu-1}$ cosets;
\item the \textbf{full operator} $T_{p^{\nu}}$ of \eqref{eq:fullop}, the sum over
  \emph{all} double cosets of determinant $p^{\nu}$, with $\sigma_1(p^{\nu})=1+p+\dots+p^{\nu}$
  right cosets. This is Shimura's $T(p^{\nu})$ (his $T(1,p^{\nu})$ is the primitive $T_{\Deprim_{p^{\nu}}}$), it is the operator that satisfies the
  recursion of Theorem~\ref{thm:relations}, and \textbf{``$T_{p^{\nu}}$'' with no
  qualifier always means this one.}
\end{itemize}
They agree at $\nu\in\{0,1\}$ and differ for $\nu\ge2$, by
\eqref{eq:fullassembly}; at $\nu=2$ the difference is exactly the central term
$\diam=p^{\,n}I$:
\[
  T_{p^{2}}\;=\;T_{\Deprim_{p^{2}}}\;+\;p^{\,n}I .
\]
Where both appear in one line the emphatic forms $T^{\mathrm{full}}_{p^{\nu}}:=T_{p^{\nu}}$
and $T^{\mathrm{prim}}_{p^{\nu}}:=T_{\Deprim_{p^{\nu}}}$ may be used. The implementation's
comments write $T'_{p^{r}}$ for the primitive operator; \emph{this document does not},
because a prime is a decoration and decorations do not make new symbols
(\S\ref{sec:notation-rules}, \ref{N2}).
\end{contractbox}

\noindent
Which entry point of the implementation emits which object is tabulated in
Appendix~\ref{app:errata}, entry~\ref{err:e1}, together with the evidence; the
code map at the end of this chapter repeats the table in short form. The arithmetic relations that distinguish them are
Theorem~\ref{thm:relations} and its caveat.

\subsection{The correspondence datum}\label{sec:correspondencedatum}

Fix a prime $p$, an exponent $\nu\ge1$, and put
\[
  N:=p^{\nu},\qquad g:=\smat{1}{0}{0}{N}\in\De ,
\]
$\De$ the monoid of Definition~\ref{def:semigroup}. Matrices act on column vectors from
the left, so $\Ga h$ is a \emph{right} coset and $h\Ga$ a \emph{left} coset.

\begin{proposition}[Coset decompositions]\label{prop:cosets}
\PROVED{} Let $p$ be a prime, $\nu\ge1$, $N=p^{\nu}$, $g=\smat{1}{0}{0}{N}$, and
$m:=N+p^{\nu-1}=p^{\nu}+p^{\nu-1}$. Then:
\begin{enumerate}
\item \textbf{Right cosets.} $\Ga g\Ga$ is the set of all $\delta\in\Mat_2(\Z)$ with
  $\det\delta=N$ and $\gcd$ of the entries of $\delta$ equal to $1$
  (\emph{primitive}), and
  \[
    \Ga g\Ga\;=\;\bigsqcup_{d\mid N}\ \bigsqcup_{\substack{0\le b<d\\ \gcd(N/d,\,b,\,d)=1}}
    \Ga\,\smat{N/d}{b}{0}{d}\,,
  \]
  a disjoint union of exactly $m$ right cosets. For $\nu=1$ this is
  $\bigsqcup_{0\le b<p}\Ga\beta_b\sqcup\Ga\beta_{\infty}$ with
  $\beta_b:=\smat{1}{b}{0}{p}$ and $\beta_{\infty}:=\smat{p}{0}{0}{1}$.
  \LEAN{prop:cosets} (the first sentence, the count $m$ and the list at $\nu=1$; not the list for general $\nu$)
\item \textbf{Left cosets.} $\Ga g\Ga=\bigsqcup\,\delta\Ga$ over the transposes of the
  matrices in (1); there are again $m$ of them, and for $\nu=1$ they are
  $\smat{1}{0}{b}{p}\Ga$ $(0\le b<p)$ and $\smat{p}{0}{0}{1}\Ga$. \LEAN{prop:cosets}
\item \textbf{The two congruence subgroups.}
  \[
    \begin{aligned}
      \Ga'&:=\Ga\cap g^{-1}\Ga g=\Ga^{0}(N)=\Bigl\{\smat{a}{b}{c}{d}\in\Ga:\ b\equiv0\ (N)\Bigr\},\\
      \Ga''&:=\Ga\cap g\Ga g^{-1}=\Ga_0(N)=\Bigl\{\smat{a}{b}{c}{d}\in\Ga:\ c\equiv0\ (N)\Bigr\},
    \end{aligned}
  \]
  both of index $m$ in $\Ga$, and the conjugation
  $\conj_g\colon\Ga'\to\Ga''$, $h\mapsto ghg^{-1}$, is a group isomorphism with integral
  values in both directions:
  $g\smat{a}{Nb}{c}{d}g^{-1}=\smat{a}{b}{Nc}{d}$.
  \LEAN{prop:cosets-3-iso} (since v1.18.0 including the words ``group isomorphism'')
\item \textbf{The coset spaces.} The bottom-row map
  $\smat{a}{b}{c}{d}\mapsto(c:d)$ induces a bijection
  \[
    \Ga''\backslash\Ga\;\xrightarrow{\ \sim\ }\;\mathbb{P}^1(\Z/N)
    :=\bigl\{(c:d)\in(\Z/N)^2:\gcd(c,d,N)=1\bigr\}\big/(\Z/N)^{\times},
  \]
  \[
    \abs{\mathbb{P}^1(\Z/N)}=m .
  \]
  The right $\Ga$-set $\Ga\backslash\Ga g\Ga$ (cosets $\Ga\delta$, action
  $\Ga\delta\cdot\gamma=\Ga\delta\gamma$) is isomorphic to $\Ga'\backslash\Ga$ via
  $\Ga'\gamma\mapsto\Ga g\gamma$, and the left $\Ga$-set $\Ga g\Ga/\Ga$ is isomorphic to
  $\Ga/\Ga''$ via $\gamma\Ga''\mapsto\gamma g\Ga$.
  \LEAN{prop:cosets} (the bijection and $\abs{\mathbb{P}^1(\Z/N)}=m$; not the two isomorphisms of $\Ga$-sets)
\item \textbf{Transporters and coset elements.} If $s_1,\dots,s_m\in\Ga$ represent
  $\Ga''\backslash\Ga$, then
  \[
    \delta_j:=s_j^{-1}g,\qquad j=1,\dots,m,
  \]
  represent the left cosets: $\Ga g\Ga=\bigsqcup_{j=1}^{m}\delta_j\Ga$. Every
  $\delta_j$ is integral of determinant $N$, and no $\delta_j$ is ever inverted anywhere
  in this document. \LEAN{prop:cosets-5-any} (since v1.18.0 for any transversal, not only the canonical $s_j$ of Definition~\ref{def:cosetreps})
  \LEAN{lem:m4-cosetdata} (since v1.20.0: $\Ga''=\Ga\cap g\Ga g^{-1}$ has finite index, so such $s_j$ exist, for every $g\in\De$)
\end{enumerate}
\end{proposition}

\begin{proof}
(1) \emph{Hermite reduction.} Every $\delta\in\Mat_2(\Z)$ with $\det\delta=N>0$ lies in
$\Ga\smat{a}{b}{0}{d}$ for exactly one triple $(a,b,d)$ with $a,d>0$, $ad=N$,
$0\le b<d$. Existence: let $(\alpha,\gamma_{21})^{\top}$ be the first column of $\delta$
--- not $(0,0)^{\top}$, since $\det\delta\neq0$ --- and $e:=\gcd(\alpha,\gamma_{21})>0$;
the row $(-\gamma_{21}/e,\ \alpha/e)$ has coprime entries, so it is the second row of some
$\gamma_0\in\Ga$, and $\gamma_0\delta$ has lower-left entry $0$; replacing $\gamma_0$ by
$-\gamma_0$ if necessary (note $-I\in\Ga$) makes both diagonal entries positive, and
multiplying on the left by $T^{-\lfloor b/d\rfloor}$ moves $b$ into $[0,d)$ without
touching $a$ or $d$. Uniqueness: if two such normal forms are $\Ga$-equivalent, the
connecting element is upper triangular with positive diagonal and determinant $1$, hence
equals $T^{k}$, which forces $a=a'$, $d=d'$ and $b\equiv b'\ (d)$, hence $b=b'$. So the
listed matrices are pairwise inequivalent and the $\Ga\smat{a}{b}{0}{d}$ over all
admissible triples exhaust $\{\det=N\}$.

\emph{Which of them are primitive.} The content of a matrix is $\Ga$-biinvariant, so
$\delta$ is primitive iff its Hermite form is, i.e.\ iff $\gcd(a,b,d)=\gcd(N/d,b,d)=1$;
and the primitive matrices of determinant $N$ form the single double coset
$\Ga\,\mathrm{diag}(1,N)\,\Ga$, because the elementary divisors $d_1\mid d_2$ of $\delta$
--- which classify $\Ga\delta\Ga$, and which satisfy $d_1=\mathrm{content}(\delta)$,
$d_1d_2=\det\delta$ --- are then $1$ and $N$
\STD{\cite[\S3.2, pp.~55--57, Lemmas~3.11--3.13]{shimura}: the Hecke ring of $(\Ga,\De)$, where double cosets are
parametrised by elementary divisors; locator verified (R20-L; v1.21.0-pre.1, in place of ``Ch.~3, \S3.1'')}. The count, with $N=p^{\nu}$ and $d=p^{\beta}$: for
$\beta=0$ the only $b$ is $0$ and the condition reads $\gcd(p^{\nu},0,1)=1$, giving $1$;
for $\beta=\nu$ the condition is vacuous, giving $p^{\nu}$; for $0<\beta<\nu$ it reads
$p\nmid b$, giving $p^{\beta}-p^{\beta-1}$. The total is
$1+\sum_{\beta=1}^{\nu-1}(p^{\beta}-p^{\beta-1})+p^{\nu}=p^{\nu}+p^{\nu-1}=m$.

(2) Transposition is an anti-automorphism of $\Mat_2(\Z)$ with $\Ga^{\top}=\Ga$ and
$g^{\top}=g$, so it carries $\Ga g\Ga$ onto itself and exchanges left with right cosets.

(3) $g\smat{a}{b}{c}{d}g^{-1}=\smat{a}{b/N}{Nc}{d}$, which is integral iff $N\mid b$ and
then has lower-left entry divisible by $N$; this gives $\Ga\cap g\Ga g^{-1}=\Ga_0(N)$ and
simultaneously the stated formula for $\conj_g$. Likewise
$g^{-1}\smat{a}{b}{c}{d}g=\smat{a}{Nb}{c/N}{d}$ gives $\Ga\cap g^{-1}\Ga g=\Ga^{0}(N)$.
Both are groups, $\conj_g$ is a bijection with inverse $\conj_{g^{-1}}$, and the index
statement follows from (4).

(4) The map is well defined into $\mathbb{P}^1(\Z/N)$: $ad-bc=1$ forces
$\gcd(c,d,N)=1$. It is constant on $\Ga''$-orbits: for
$h=\smat{\alpha}{\beta}{N\gamma_0}{\epsilon}\in\Ga_0(N)$ the bottom row of $h\gamma$ is
$(N\gamma_0a+\epsilon c,\ N\gamma_0b+\epsilon d)\equiv\epsilon\,(c,d)\ (N)$, and
$\alpha\epsilon\equiv1\ (N)$ makes $\epsilon$ a unit. It is injective on orbits: if
$\gamma_1,\gamma_2\in\Ga$ have bottom rows $(c_1,d_1)\equiv u\,(c_2,d_2)$ with $u$ a
unit, then the bottom row of $\gamma_1\gamma_2^{-1}$ is
$(c_1d_2-d_1c_2,\ -c_1b_2+d_1a_2)$, whose first entry is
$\equiv u(c_2d_2-d_2c_2)=0\ (N)$, so $\gamma_1\gamma_2^{-1}\in\Ga_0(N)$. It is
surjective, with the count $m$: a class $(c:d)$ with $\gcd(c,d,p)=1$ has either $d$ a
unit --- normalise to $(k:1)$, $k\in\Z/N$, giving $N$ classes --- or $p\mid d$, in which
case $c$ is a unit and one normalises to $(1:pc_0)$, $c_0\in\Z/p^{\nu-1}$, giving
$p^{\nu-1}$ classes; the two families are disjoint, because $(k:1)=(1:pc_0)$ would force
the unit $u$ with $u\cdot1=pc_0$, and $pc_0$ is not a unit. Each class is hit by the
closed-form lift of Definition~\ref{def:cosetreps}, which is in $\Ga$ by inspection. The
two $\Ga$-set isomorphisms are the orbit--stabiliser bijections for the transitive
actions on $\Ga\backslash\Ga g\Ga$ and $\Ga g\Ga/\Ga$, whose stabilisers of $\Ga g$ and
of $g\Ga$ are $\Ga'$ and $\Ga''$ by (3).

(5) The map $\Ga''s_j\mapsto s_j^{-1}g\Ga$ is the composite of the bijection
$\Ga''\backslash\Ga\xrightarrow{\,\sim\,}\Ga/\Ga''$, $\Ga''\gamma\mapsto\gamma^{-1}\Ga''$,
with the bijection $\Ga/\Ga''\xrightarrow{\,\sim\,}\Ga g\Ga/\Ga$ of (4); it is therefore
a bijection, so the $\delta_j=s_j^{-1}g$ represent the $m$ left cosets. (Directly: if
$s'=h\,s_j$ with $h\in\Ga''$ then $s'^{-1}g=s_j^{-1}g\,(g^{-1}h^{-1}g)$ with
$g^{-1}h^{-1}g\in\Ga'\subseteq\Ga$, so the left coset does not depend on the
representative.) Finally $\det\delta_j=\det g=N$, and $\delta_j$ is integral because
$s_j^{-1}\in\Ga$.
\end{proof}

\begin{definition}[The canonical representatives, the class map and the cocycle]
\label{def:cosetreps}
\LEAN{def:cosetreps} (in [GT] with a $0$-based index, the class map an index)
Let $p,\nu,N,m$ be as in Proposition~\ref{prop:cosets}. The \emph{canonical
transporters} are, in the $1$-based index $j$ that this document and the implementation
share,
\[
  s_j:=\smat{1}{0}{k}{1}\ \ \text{with }k=j-1,\quad 1\le j\le N
  \qquad\text{(class $(k:1)$)},
\]
\[
  s_{N+1+c}:=\smat{0}{-1}{1}{pc},\quad 0\le c<p^{\nu-1}
  \qquad\text{(class $(1:pc)$)},
\]
each of determinant $1$ and with bottom row its assigned class; they exhaust
$\Ga''\backslash\Ga$ by Proposition~\ref{prop:cosets}(4). At $\nu=1$ this is
$s_1=I$, $s_{1+k}=\smat{1}{0}{k}{1}$ $(1\le k\le p-1)$, $s_{p+1}=S$. The
\emph{coset elements} are $\delta_j:=s_j^{-1}g$. The \emph{class map} is
\[
  \mathrm{cls}\colon\Ga\to\mathbb{P}^1(\Z/N),\qquad
  \mathrm{cls}\smat{a}{b}{c}{d}=(c:d),
\]
and $j\cdot\gamma$ denotes the index with
$\mathrm{cls}(s_j\gamma)=\mathrm{cls}(s_{j\cdot\gamma})$, i.e.\ the right action of $\Ga$
on $\{1,\dots,m\}\cong\Ga''\backslash\Ga$; concretely $(c,d)\cdot\gamma$ is the
\emph{row} vector product $(c,d)\gamma$ reduced mod $N$. The \emph{coset cocycle} is
\begin{equation}\label{eq:cocycle}
  h_j(\gamma)\;:=\;s_j\,\gamma\,s_{j\cdot\gamma}^{-1}\ \in\ \Ga''=\Ga_0(N),
  \qquad\text{equivalently}\qquad
  s_j\gamma=h_j(\gamma)\,s_{j\cdot\gamma}.
\end{equation}
Finally, the \emph{upper-triangular transporters} $t_j$ representing
$\Ga^{0}(N)\backslash\Ga$ (class map $=$ the \emph{top} row $(a:b)$) are
$t_j:=\smat{1}{k}{0}{1}=T^{k}$ with $k=j-1$ for $1\le j\le N$, and
$t_{N+1+c}:=\smat{-pc}{-1}{1}{0}$ for $0\le c<p^{\nu-1}$; at $\nu=1$, $t_{p+1}=S$. They
are used only in \S\ref{sec:induced} (Proposition~\ref{prop:alpha}).
\end{definition}

\begin{contractbox}
Definition~\ref{def:cosetreps} is \textbf{binding as written}, not up to an equally valid
choice. The general recipe ``lift $(c:d)$ to $\Ga$ by an extended-Euclid B\'ezout
solve'' does \emph{not} determine the matrix: at $(c,d)=(1,1)$ the textbook iterative
extended gcd returns $\smat{0}{-1}{1}{1}$, which is a correct lift of the class and is
\emph{not} $s_2=\smat{1}{0}{1}{1}$. The superseded contract pinned only the $\nu=1$
table and left $\nu\ge2$ to the implementation, which therefore chose and documented the
closed forms above (\file{Hecke/HeckeCosetTables.h}); they are recorded here so that the
$\nu\ge2$ tables, the $\tau$-table built on them and every byte-level cross-check have a
contractual source. Changing them is a change to this document, not a refactor: the
cochain-level matrices depend on the choice (Remark~\ref{rmk:indeplevel}), only
$\Amat{q}$ does not.
\end{contractbox}

\begin{lemma}[The cocycle is a cocycle]\label{lem:cocycle}
\PROVED{} \LEAN{lem:cocycle} With the notation of Definition~\ref{def:cosetreps}: $j\cdot\gamma$ and
$h_j(\gamma)$ are well defined, $(j,\gamma)\mapsto j\cdot\gamma$ is a right action of
$\Ga$ on $\{1,\dots,m\}$, $h_j(\gamma)\in\Ga_0(N)$, and
\[
  h_j(\gamma_1\gamma_2)=h_j(\gamma_1)\,h_{j\cdot\gamma_1}(\gamma_2),
  \qquad h_j(1)=I,
  \qquad \gamma\in\Ga,\ 1\le j\le m .
\]
Moreover $\abs{s_j}\le\max(p^{\nu},1)$ entrywise, and the same bound holds for
$s_j^{-1}$ and for $\delta_j=s_j^{-1}g$.
\end{lemma}

\begin{proof}
$j\cdot\gamma$ is well defined because $\mathrm{cls}$ induces the bijection of
Proposition~\ref{prop:cosets}(4) and the $\Ga$-action on $\Ga''\backslash\Ga$ is by right
translation; it is an action because $\mathrm{cls}$ is $\Ga$-equivariant.
$h_j(\gamma)=s_j\gamma s_{j\cdot\gamma}^{-1}$ lies in $\Ga$ and has
$\mathrm{cls}(h_j(\gamma)s_{j\cdot\gamma})=\mathrm{cls}(s_j\gamma)=\mathrm{cls}(s_{j\cdot\gamma})$,
so $h_j(\gamma)\in\Ga''$ by injectivity in Proposition~\ref{prop:cosets}(4). For the
cocycle identity, expand
$s_j\gamma_1\gamma_2=h_j(\gamma_1)s_{j\cdot\gamma_1}\gamma_2
=h_j(\gamma_1)h_{j\cdot\gamma_1}(\gamma_2)s_{(j\cdot\gamma_1)\cdot\gamma_2}$ and compare
with $s_j(\gamma_1\gamma_2)=h_j(\gamma_1\gamma_2)s_{j\cdot(\gamma_1\gamma_2)}$, using
that the $s_j$ are a fixed set of representatives. The bounds are read off
Definition~\ref{def:cosetreps}: the entries of $s_j$ lie in
$\{0,\pm1,k,pc\}\subseteq[-p^{\nu},p^{\nu}]$, $s_j^{-1}$ negates the off-diagonal and
swaps the diagonal, and right multiplication by $g=\mathrm{diag}(1,N)$ scales the second
column by $N=p^{\nu}$ --- whence the sharper statement recorded as
Lemma~\ref{lem:bounds}.
\end{proof}

\noindent
The matrices $\rhon(h_j(\gamma))$ are the \emph{block transition matrices} of the
induced local system; they are assembled in Proposition~\ref{prop:blockmatrices} and
\eqref{eq:MH}, and the permutation $j\mapsto j\cdot\gamma$ is what makes the induced
system a matrix local system on $(K,\Rel)$ (Proposition~\ref{prop:lhcompliant}). The
worked $p=2$ table of $(j\cdot\gamma,h_j(\gamma))$ on the generator multiset
$\mathcal{G}$ is \S\ref{sec:example}.

\subsection{The operators}\label{sec:operatordefs}

\begin{definition}[Operators]\label{def:operators}
For a double coset $\Ga\delta\Ga\subseteq\De$ let $\Tdc{\delta}$ be the operator on
$\Heq{q}(K,\Rel;L)$ constructed in \S\ref{sec:chainlevel}
(Definition~\ref{def:hecke}). Write
\[
  \Deprim_{p^{\nu}}:=\Ga\,\mathrm{diag}(1,p^{\nu})\,\Ga
  \qquad(\nu\ge0,\ \Deprim_{1}=\Ga)
\]
for the \emph{primitive} determinant-$p^{\nu}$ double coset, of $m=p^{\nu}+p^{\nu-1}$
cosets for $\nu\ge1$ and $1$ coset for $\nu=0$
(Proposition~\ref{prop:cosets}(1)). Then:
\begin{enumerate}
\item $T_p:=T_{\Deprim_{p}}=\Tdc{\mathrm{diag}(1,p)}$; \LEAN{def:m4-hecke}
\item $\diam:=\Tdc{pI}$, the \emph{central} operator; \LEAN{def:m4-hecke}
\item the \textbf{full} operator of degree $p^{\nu}$ is
  \begin{equation}\label{eq:fullop}
    T_{p^{\nu}}\;:=\!\!\sum_{\substack{\Ga\delta\Ga\ \subseteq\ \De\\ \det\delta=p^{\nu}}}\!\!
    T_{\Ga\delta\Ga}\;,
    \qquad T_{p^{0}}:=I,
  \end{equation}
  the sum over \emph{all} double cosets of determinant $p^{\nu}$ --- equivalently, the
  operator attached to the characteristic function of $\{\delta\in\De:\det\delta=p^{\nu}\}$,
  whose right-coset count is $\sigma_1(p^{\nu})=1+p+\dots+p^{\nu}$. This is Shimura's
  $T(p^{\nu})$ \STD{\cite[\S3.2, Thm.~3.24(1), (2), (7), p.~63]{shimura}}; Shimura's $T(1,p^{\nu})$ is the primitive $\Deprim_{p^{\nu}}$ (v1.19.0-pre.5: Shimura's notation corrected --- his $T(1,p^{\nu})$ is the single double coset $\Deprim_{p^{\nu}}$, his $T(p^{\nu})$ the full sum; v1.21.0-pre.1: \textbf{verified} on the owner's scan, R20-L --- Thm.~3.24(1) $T(m)=\sum_{ad=m,\,a\mid d}T(a,d)$, (2) $T(1,p^{k})=T(p^{k})-T(p,p)T(p^{k-2})$, (7) $\deg T(m)=\sigma_1(m)$);
  \LEAN{def:m4-hecke} (only in the assembled form \eqref{eq:fullassembly}, by which [GT] defines the full operator on the model;
  \eqref{eq:fullop} itself carries no marker)
\item for $N=\prod_{t}p_t^{\nu_t}\in\Z_{>0}$ the composite operator is
  $T_N:=\prod_{t}T_{p_t^{\nu_t}}$, well defined by Theorem~\ref{thm:relations}(1);
  $T_1=I$;
\item $U_\ell:=T_\ell$, the operator of (1) in the bad case $p=\ell$ of \ref{H2}.
\end{enumerate}
\end{definition}

\begin{proposition}[Central operator]\label{prop:diamond}
\PROVED{} \LEAN{prop:diamond-model} (on the model, for any scalar $NI$ with $\rho(NI)$ the scalar $c$: with $\tau=\id$ and no hypothesis
on $\ell$, and for any $\tau$ under $6\in\cring^{\times}$) Assume \ref{H4}. For every prime $p$, every $\nu\ge0$, every commutative
coefficient ring $\cring$ and every degree $q$ --- \emph{with no hypothesis on
$\ell$} --- the double coset $\Ga\,p^{\nu}I\,\Ga=p^{\nu}I\cdot\Ga$ is a single left and a
single right coset, and
\[
  T_{\Ga\,p^{\nu}I\,\Ga}\;=\;\rhon(p^{\nu}I)\;=\;p^{\,\nu n}\cdot I
  \qquad\text{on }\Heq{q}(K,\Rel;L).
\]
In particular $\diam=p^{\,n}I$. Since $\Ga=\SL_2(\Z)$ has level $1$ and hence no
non-trivial nebentypus, the classical diamond $\langle p\rangle$ is trivial and
$\diam=p^{\,n}\langle p\rangle=p^{\,n}I$ is the classical normalisation with weight
$k=n+2$: $p\cdot p^{\,n}=p^{\,k-1}$.
\end{proposition}

\begin{proof}
$p^{\nu}I$ is central in $\GL_2(\Q)$, so $\Ga p^{\nu}I\Ga=p^{\nu}I\Ga=\Ga p^{\nu}I$ is one
coset, $m=1$, $\Ga'=\Ga''=\Ga$, $s_1=I$, $\delta_1=p^{\nu}I$, and
$\conj_{p^{\nu}I}=\id$. Hence $\tau=\id$ satisfies (T1)--(T3) of
Definition~\ref{def:tau} over every $\cring$ --- (T1) reads $\id(hx)=h\,\id(x)$ --- and
it is this explicit datum, not Lemma~\ref{lem:tauexists}, that the construction uses, which
is why no hypothesis on $\ell$ is needed. Equation~\eqref{eq:heckechain} then reads
$(\widetilde T\phi)(x)=\rhon(p^{\nu}I)\phi(x)=p^{\,\nu n}\phi(x)$ by
Definition~\ref{def:rhon} applied to a scalar matrix, and a scalar cochain map induces
the same scalar on cohomology.
\end{proof}

\begin{proposition}[Primitive-to-full assembly]\label{prop:fullassembly}
\PROVED{} Assume \ref{H4}. For every prime $p$, every $\nu\ge0$, every commutative
coefficient ring $\cring$ and every degree $q$ --- again with \emph{no} hypothesis on
$\ell$ ---
\begin{equation}\label{eq:fullassembly}
  T_{p^{\nu}}\;=\;\sum_{0\le 2j\le\nu} p^{\,jn}\;T_{\Deprim_{p^{\nu-2j}}}\,,
  \qquad T_{\Deprim_{p^{0}}}:=I .
\end{equation}
For $\nu=0,1$ the sum has one term and $T_{p^{\nu}}=T_{\Deprim_{p^{\nu}}}$. For $\nu=2$
it reads $T_{p^{2}}=T_{\Deprim_{p^{2}}}+p^{\,n}I$, the extra term being the central
coset $pI$ acting by $\diam=p^{\,n}I$ (Proposition~\ref{prop:diamond}). \LEAN{prop:fullassembly-model} (the summand only:
$T_{\Ga p^{j}\mathrm{diag}(1,p^{\mu})\Ga}=p^{\,jn}\,T_{\Deprim_{p^{\mu}}}$, at cochain level with no hypothesis on $\ell$, for an arbitrary $\tau$
under $6\in\cring^{\times}$, on the model; not \eqref{eq:fullassembly} as an identity with \eqref{eq:fullop}: the enumeration of the double
cosets of determinant $p^{\nu}$ is not formalised, and [GT]'s full operator is defined by the assembled form)
\end{proposition}

\begin{proof}
\emph{The index set.} Every $\delta\in\Mat_2(\Z)$ with $\det\delta=p^{\nu}$ has content
$\gcd(\delta_{11},\delta_{12},\delta_{21},\delta_{22})=p^{j}$ for a unique $j\ge0$, and
$\delta=p^{j}\delta_0$ with $\delta_0$ primitive of determinant $p^{\nu-2j}\ge1$, whence
$2j\le\nu$. Content is $\Ga$-biinvariant, so
\[
  \{\delta\in\De:\det\delta=p^{\nu}\}=\bigsqcup_{0\le2j\le\nu}p^{j}I\cdot\Deprim_{p^{\nu-2j}},
\]
and the double cosets of determinant $p^{\nu}$ are exactly the $p^{j}I\cdot\Deprim_{p^{\nu-2j}}$,
one for each admissible $j$. (Consistency check on right-coset counts:
$\sum_{0\le2j\le\nu}\bigl(p^{\nu-2j}+p^{\nu-2j-1}\bigr)$, with the $j=\nu/2$ term read as
$1$, telescopes to $1+p+\dots+p^{\nu}=\sigma_1(p^{\nu})$.)

\emph{The scalar.} \emph{Within this proof only}, the coset index is written $t$: the
letter $j$ is taken by the content exponent of \eqref{eq:fullassembly}, whose spelling is
fixed by Appendix~\ref{app:errata}, entry~\ref{err:e1}. Fix $j$ and put
$g_0:=\mathrm{diag}(1,p^{\nu-2j})$, $g:=p^{j}g_0$. Since
$p^{j}I$ is central, $\Ga\cap g\Ga g^{-1}=\Ga\cap g_0\Ga g_0^{-1}$ and
$\conj_g=\conj_{g_0}$, so the \emph{same} transporters $s_t$ and the \emph{same} twisted
chain datum $\tau$ serve both (Definition~\ref{def:tau} depends on $g$ only through
$\conj_g$), while $\delta_t=s_t^{-1}g=p^{j}\,(s_t^{-1}g_0)$ and therefore
$\rhon(\delta_t)=p^{\,jn}\rhon(s_t^{-1}g_0)$ by Proposition~\ref{prop:symconv}(2) and
$\rhon(p^{j}I)=p^{\,jn}I$. Substituting into \eqref{eq:heckechain} gives
$T_{p^{j}I\cdot\Deprim_{p^{\nu-2j}}}=p^{\,jn}\,T_{\Deprim_{p^{\nu-2j}}}$ already at
cochain level. Summing over $j$ as in \eqref{eq:fullop} gives \eqref{eq:fullassembly}.
\end{proof}

\begin{remark}[The raw pipeline computes the primitive operator]\label{rmk:rawprimitive}
The coset tables of Definition~\ref{def:cosetreps} --- the $\mathbb{P}^1(\Z/p^{\nu})$
representatives --- enumerate the \emph{primitive} double coset $\Deprim_{p^{\nu}}$ and
nothing else. The pipeline
\code{HeckeCosetTables($p,\nu$)}$\to$\code{HeckeTauTable}$\to$\code{HeckeOperatorOnK}
therefore emits $T_{\Deprim_{p^{\nu}}}$, exposed as
\code{T\_p\_pow\_primitive\_direct}. The \emph{standard} operator $T_{p^{\nu}}$ is
assembled from it by \eqref{eq:fullassembly} (\code{T\_p\_pow\_direct}) or, equivalently,
obtained by the recursion of Theorem~\ref{thm:relations} (\code{T\_p\_pow}); the two
agree by construction and are asserted equal at runtime. This is the whole content of
Appendix~\ref{app:errata}, entry~\ref{err:e1}: a reader who takes the raw pipeline's
output for $T_{p^{\nu}}$ at $\nu\ge2$ is off by the terms of \eqref{eq:fullassembly}, and
two numbered prescriptions of the superseded contract did exactly that.
\end{remark}

\subsection{The Hecke algebra relations}\label{sec:relations}

\begin{theorem}[Hecke relations, integral normalisation]\label{thm:relations}
Assume \ref{H4} and $6\in\cring^{\times}$ --- for $\cring$ of characteristic $\ell$,
$\ell\ge5$ (hypothesis \ref{H3}). Then on $\Heq{q}(K,\Rel;\Sym^n(\cring^2))$, for every
degree $q$, all primes $p,p'$ and all $\nu\ge1$:
\begin{enumerate}
\item $T_pT_{p'}=T_{p'}T_p$, \LEANlit{thm:relations-comm-model}{the commutativity of the Hecke ring, \cite[Prop.~3.8, p.~54, applied in \S3.2, p.~56]{shimura} (there for $\det>0$; for $\det<0$ the same argument with a diagonal representative), transported by \cite[\S1.1, p.~194]{ash_stevens}, as the hypothesis \texttt{hcomm}; locator verified (R20-L)} (as an instance of the commutativity of
  all $\Tdc{g}$, Obligation~\ref{obl:algebra}) and the \emph{full} operators of
  Definition~\ref{def:operators}(3) satisfy
  \[
    T_{p^{\nu+1}}\;=\;T_p\,T_{p^{\nu}}\;-\;p\,\diam\,T_{p^{\nu-1}}
    \;=\;T_p\,T_{p^{\nu}}\;-\;p^{\,n+1}\,T_{p^{\nu-1}} ;
  \]
  equivalently, for $\nu=1$: $(T_p)^2=T_{p^2}+p^{\,n+1}I$ --- note that $(T_p)^2$ and
  $T_{p^2}$ are different objects and are \emph{not} written alike in this document.
  \LEANlit{thm:relations-1-model}{the primitive products of \cite[Thm.~3.24(5), p.~63, with Prop.~3.17, p.~58]{shimura}, transported by \cite[\S1.1, p.~194]{ash_stevens}, as the hypothesis \texttt{hTp}; locator verified (R20-L)} (on the model, for $\tau$ over $\cring$, for the full operators in the assembled
  form \eqref{eq:fullassembly}, every $\nu\ge1$; not the ring homomorphism of the \textsf{[standard]} note below)
  \STD{\cite[Thm.~3.24(4), p.~63]{shimura} for the identity in $\mathcal{H}(\De,\Ga)$ (verified, R20-L);
  \cite[\S1.1, p.~194]{ash_stevens} for the action of $\mathcal{H}$ on cohomology (stated there as well known, with a reference to Kuga--Parry--Sah, not
  proved; read, R20-L) and \cite[Ch.~III.9--10]{brown} for the transfer (\cite[Prop.~III.9.5, p.~82]{brown}), which together
  make $\Ga\delta\Ga\mapsto\Tdc{\delta}$ a ring homomorphism}
\item $T_{M}T_{N}=T_{MN}$ for $\gcd(M,N)=1$; \LEANlit{thm:relations-2-model}{the coprime product of double cosets, \cite[Prop.~3.16, p.~57]{shimura} (Thm.~3.24(3) is its full-operator form), transported by \cite[\S1.1, p.~194]{ash_stevens}, as the hypothesis \texttt{hcop}; locator verified (R20-L)} (in the primitive form
  $T_{\mathrm{diag}(1,M)}T_{\mathrm{diag}(1,N)}=T_{\mathrm{diag}(1,MN)}$, on the model; not $T_N$ of Definition~\ref{def:operators}(4)) with (1) this makes
  Definition~\ref{def:operators}(4) independent of the order of the factors.
\end{enumerate}
\emph{Caveat, and it is necessary.} $T_{p^{\nu}}$ here is the \textbf{full} operator
$T(p^{\nu})$ of \eqref{eq:fullop} (Shimura's notation), \emph{not} the primitive double-coset operator
$T_{\Deprim_{p^{\nu}}}$; for $\nu\ge2$ the two differ by the terms of
\eqref{eq:fullassembly}, e.g.\ $T_{p^{2}}=T_{\Deprim_{p^{2}}}+p^{\,n}I$. They coincide
only at $\nu\le1$, so $T_p$ is unambiguous. Substituting the primitive operator into the
recursion falsifies it by $p^{\,n}I$ (Appendix~\ref{app:errata},
entry~\ref{err:e1}). \LEANlit{lem:m4-hTp-hrel}{the primitive products of \cite[Thm.~3.24(5), p.~63, with Prop.~3.17, p.~58]{shimura}, transported by \cite[\S1.1, p.~194]{ash_stevens}, as the hypothesis \texttt{hTp}; locator verified (R20-L)} (the $p^{\,n}I$ difference, derived on the model; its half
$T_{p^{2}}=T_{\Deprim_{p^{2}}}+p^{\,n}I$ in the assembled form needs no hypothesis) The exponent is pinned by the weight: with $k=n+2$,
$p\cdot p^{\,n}=p^{\,k-1}$ is the classical Eichler--Shimura factor
(Obligation~\ref{obl:oracle}).
\end{theorem}

\begin{proof}[Proof sketch]
Under $6\in\cring^{\times}$ the assignment $\Ga\delta\Ga\mapsto\Tdc{\delta}$ extends
linearly to a ring homomorphism from the abstract Hecke algebra $\mathcal{H}(\De,\Ga)$
(the free $\Z$-module on double cosets with the convolution product) to
$\End\bigl(\Heq{q}(K,\Rel;L)\bigr)$; this is Obligation~\ref{obl:algebra}, and the
hypothesis $6\in\cring^{\times}$ enters twice --- once through
Theorem~\ref{thm:classical}, which identifies $\Tdc{\delta}$ with the classical operator
\eqref{eq:corconjres} on $H^{\bullet}(\Ga;\VR)$, and once through
Remark~\ref{rmk:orbifold}, which identifies the two cohomologies at all. In
$\mathcal{H}(\De,\Ga)$ one has the classical identities
$T(p)\,T(p^{\nu})=T(p^{\nu+1})+p\,T(p,p)\,T(p^{\nu-1})$ and commutativity
(locally at each prime $\mathcal{H}$ is the polynomial ring on $T(p)=T(1,p)$ and
$T(p,p)$, not inverted --- v1.21.0-pre.1, R20-L: ``$T(p,p)^{\pm1}$'' corrected --- and $\mathcal{H}$ is commutative)
\STD{\cite[Thm.~3.20, p.~59; Thm.~3.24(4); Prop.~3.8 with p.~56]{shimura}, verified, R20-L}. Shimura's $T(p^{\nu})$ \emph{is} the full
operator $T_{p^{\nu}}$ of \eqref{eq:fullop} --- both are the characteristic function of
$\{\det=p^{\nu}\}$ --- and this identification is exactly what makes the cited relation a
statement about $T_{p^{\nu}}$ and not about $T_{\Deprim_{p^{\nu}}}$. Applying the
homomorphism and Proposition~\ref{prop:diamond} gives (1); (2) is the coprime
factorisation $T(M)T(N)=T(MN)$ \STD{\cite[Thm.~3.24(3), p.~63]{shimura}; in the primitive form, the product of the double cosets of coprime determinant,
\cite[Prop.~3.16, p.~57]{shimura}; verified, R20-L, in place of ``Ch.~3, \S3.3''} transported the
same way.
\end{proof}

\begin{remark}[The coefficient of $T_{p^{\nu-1}}$ is $p^{\,n+1}$, not $(p+1)p^{\,n}$]
\label{rmk:relcoeff}
A working note in this project carried $(T_p)^2=T_{p^2}+(p+1)p^{\,n}$, attributed to a
corollary of the superseded contract that says no such thing
(Appendix~\ref{app:errata}, entry~\ref{err:e6}). The degree count settles it without any
convention: $T_p$ has $p+1$ right cosets, $(T_p)^2$ therefore $(p+1)^2$, the full
$T_{p^2}$ has $1+p+p^2$, and $(p+1)^2-(1+p+p^2)=p$ --- one factor $p$, multiplying the
central $\diam=p^{\,n}I$, i.e.\ $p^{\,n+1}$.
\end{remark}

\begin{corollary}[$p=\ell$: the bad prime]\label{cor:ulrecursion}
Let $\cring=\Fl$, $n\ge0$, and let $p=\ell$ (hypothesis \ref{H2}), so that
$U_\ell=T_\ell$.
\begin{enumerate}
\item \emph{(Assuming in addition $\ell\ge5$, i.e.\ \ref{H3}.)} Since
  $p^{\,n+1}=\ell^{\,n+1}\equiv0\pmod\ell$, the recursion of
  Theorem~\ref{thm:relations}(1) collapses to $T_{\ell^{\nu+1}}=U_\ell\,T_{\ell^{\nu}}$,
  whence the \emph{full} operators satisfy
  \[
    T_{\ell^{\nu}}\;=\;U_\ell^{\,\nu}\qquad\text{for all }\nu\ge1 .
  \]
\item \emph{(No hypothesis on $\ell$.)} For $n\ge1$, \eqref{eq:fullassembly} has every
  correction term multiplied by $\ell^{\,jn}\equiv0$ for $j\ge1$, so
  \[
    T_{\ell^{\nu}}\;=\;T_{\Deprim_{\ell^{\nu}}}\quad\text{on }\Fl\text{-coefficients},
    \qquad n\ge1,\ \nu\ge0 :
  \]
  at the bad prime, and only there, the full and the primitive prime-power operators
  \emph{coincide}. For $n=0$ they do not: the scalars $\ell^{\,jn}$ are then all $1$.
\end{enumerate}
\end{corollary}

\begin{proof}
(1) is Theorem~\ref{thm:relations}(1) with $p=\ell$ and an induction on $\nu$ from
$T_{\ell^{0}}=I$, $T_{\ell^{1}}=U_\ell$. (2) is Proposition~\ref{prop:fullassembly}, which
carries no hypothesis on $\ell$, read modulo $\ell$.
\end{proof}

\begin{remark}[Singularity at $p=\ell$ is not a defect of the construction]
\label{rmk:badprimesingular}
For $p=\ell$ and $n\ge1$ every $\rhon(\delta_j)\bmod\ell$ is singular
(Lemma~\ref{lem:detsym}); e.g.\
$\rhon(\mathrm{diag}(1,\ell))=\mathrm{diag}(1,0,\dots,0)\bmod\ell$ kills all of
$\Sym^n$ except the top monomial $u_0=x^{n}$, and $U_\ell$ may be nilpotent on
$\Heq{\bullet}$. Nothing in \S\S\ref{sec:chainlevel}--\ref{sec:induced} requires
invertibility of $\rhon(\delta)$ for $\delta\notin\Ga$: the only matrices ever inverted
are $\rhon(\gamma)$ with $\gamma\in\Ga$, which are unimodular
(Obligation~\ref{obl:integral}). $U_\ell$ is therefore defined integrally by the same
pipeline, with no branch and no special case.
\end{remark}

\paragraph{Code map.}
Proposition~\ref{prop:cosets}, Definition~\ref{def:cosetreps},
Lemma~\ref{lem:cocycle}: \file{Hecke/HeckeCosetTables.h/.cpp}
(\code{s(j)}, \code{s\_inv(j)}, \code{delta(j)}, \code{class\_index},
\code{perm(gen,j)} for $j\cdot\gamma$, \code{cocycle(gen,j)} for $h_j(\gamma)$ and
\code{rho\_cocycle(gen,j)} for $\rhon(h_j(\gamma))$; the index is $1$-based, matching
this chapter's subscripts verbatim and deliberately so), gated by
\code{HeckeCosetTablesValidation} and \code{HeckeCosetTablesSelfCheckValidation}.
Definition~\ref{def:operators}, Proposition~\ref{prop:diamond},
Proposition~\ref{prop:fullassembly}, Theorem~\ref{thm:relations},
Corollary~\ref{cor:ulrecursion}: \file{Hecke/HeckeOperatorFamily.h/.cpp}, with
\begin{center}
\begin{tabular}{@{}lp{0.50\textwidth}@{}}
\toprule
entry point & emits \\
\midrule
\code{T\_p($p$)} & $T_p$ (unambiguous) \\
\code{T\_p\_pow($p,\nu$)} & the \textbf{full} $T_{p^{\nu}}$, by the recursion of Thm.~\ref{thm:relations} \\
\code{T\_p\_pow\_direct($p,\nu$)} & the \textbf{full} $T_{p^{\nu}}$, by the assembly \eqref{eq:fullassembly} \\
\code{T\_p\_pow\_primitive\_direct($p,\nu$)} & the \textbf{primitive} $T_{\Deprim_{p^{\nu}}}$ (the raw pipeline) \\
\code{T\_n($N$)} & $T_N$ of Def.~\ref{def:operators}(4); the code identifier keeps the retired spelling \\
\code{T\_pp($p$)} & $\diam=p^{\,n}I$, by the closed formula of Prop.~\ref{prop:diamond}, with no pipeline run \\
\bottomrule
\end{tabular}
\end{center}
There is \emph{no} \code{U\_ell} entry point, and that is the contract's own invariant
made visible: $U_\ell$ is \code{T\_p(ell)}, the same pipeline and the same code path,
with no branch on $p=\ell$ anywhere (Remark~\ref{rmk:badprimesingular},
Obligation~\ref{obl:integral}). \COMP{} A recorded reach limit of that code path, not a
defect of the contract (v1.3.3; \file{docs/implementation/slot1/WP-2\_report.md} \S4):
at $n=16$, \code{--bench hecke 16 <ell> <ell>} ($U_\ell$, $\nu=1$) completes at
$\ell\in\{691,1009,2003,3617\}$ and throws \code{HeckeOverflow} (``exact multiply
overflow'', \file{Hecke/ExactInt.h}) in the coset construction at
$\ell\in\{5003,7001,10007,43867\}$, independently of $n$ (checked at $n\in\{2,4,8,16\}$),
while $\ell=43867$ with a good prime runs in milliseconds --- so the \code{int64} tier of
Theorem~\ref{thm:backend} is exhausted somewhere in
$3617<\ell<5003$, and a wider tier is what an $\ell$ beyond that needs
(\file{docs/KNOWN\_ISSUES.md} \S3).

\smallskip
\textbf{Repaired and re-measured at v1.5.0} (slot-3 package WP-1,
\file{docs/implementation/slot3/WP-1\_report.md}; the repair was planned as slot-2 WP-3
and did not run in slot 2). Two corrections to the record above and one new number.
\emph{(a) The site.} The overflow was never in the coset construction: it is in
Algorithm~\ref{alg:tau} step~1's fundamental-domain reduction, on the quadratic point
Lemma~\ref{lem:bounds}(6) bounds, and the quantity that overflows is a transient local,
not a stored transporter (Remark~\ref{rmk:btabsmall}). \emph{(b)
Lemma~\ref{lem:bounds} did not fail}, and neither did
Theorem~\ref{thm:backend} on the statements it made: the certified $B$ was inside the
tier throughout, and the gap was that no clause covered step~1 --- which is clause~(6)'s
reason for existing and Theorem~\ref{thm:backend}(iv)'s. \emph{(c) The new reach}, after
the per-operation escalation to the \code{\_\_int128} tier at that site: \COMP{} the
binding constraint is again Lemma~\ref{lem:bounds}(5), through clause~(e) of the
$\tau$-certification, which on this model reduces to $4N^{2}\le2^{31}-1$, i.e.\
$\boldsymbol{\ell\le23170}$ (largest prime $23167$); past it the table is \textbf{refused
by Obligation~\ref{obl:tau}'s certification}, by the named clause that owns the bound,
rather than by an anonymous overflow --- the same input now fails where this document
says it should. Measured: $(n,\ell)=(16,5003)$ completes in $73$\,s / $0.95$\,GB,
$(16,10007)$ in $612$\,s / $1.9$\,GB, and $\ell=43867$ builds its $\tau$-table with no
overflow and is then refused by clause~(e). Raising $\ell$ beyond $23170$ is a change to
the \emph{stored} type and is deliberately not made: no cell the programme needs is
behind it, and the memory admission of $U_\ell$ refuses those cells first
(\code{HeckeUEllReachValidation}, \code{HeckeUEllBudgetValidation}).
\textbf{Lifted at v1.23.0-pre.9} (slot-24 ELL-W) \emph{without} changing the stored type: by
Theorem~\ref{thm:backend}(4) a table with $B>2^{31}-1$ is certified in the stored wide tier,
so clause~(e) now refuses only when $2NB>2^{63}-1$, which with $B=4N^{2}$ is $\ell>1048575$;
\COMP{} \code{--bench hecke 0 43867 43867} completes ($28$\,s, $2.1$\,GB) where it was refused,
and $U_{30011}$ at $n=10$ matches $\tau(30011)\bmod30011$ (\code{HeckeWideTauValidation}). The
memory admission of \code{u\_ell = auto} (the frozen model \code{u\_ell\_rss/v2-2026-09-17}) is a
mathematical setting and is \emph{not} recalibrated here; the module route of the saturation
cells does not consult it. The gates are
\code{HeckeOperatorFamilySmokeValidation} --- which asserts
\code{T\_p\_pow} $=$ \code{T\_p\_pow\_direct} and the decomposition witness
\code{T\_p\_pow\_direct} $=$ \code{T\_p\_pow\_primitive\_direct} $+\,p^{\,n}I$ at
$\nu=2$ --- and \code{HeckeRelationsValidation} (oracle O9: commutativity, the $\nu=2$
recursion, and the primitive-versus-full witness). The hypothesis $6\in\cring^{\times}$
of Theorem~\ref{thm:relations} is respected by those gates and not merely by this text:
at $\ell\in\{2,3\}$ the recursion and the $U_{\ell^{\nu}}$ identities are \emph{recorded
and not asserted}, the hard assertions there being only the hypothesis-free ones
(pipeline completes, $\tau$ certifies, well-definedness), and the $p=\ell$ recursion is
asserted only at $(p,\ell)=(5,5)$
(\file{docs/implementation/hecke\_action\_implementation\_plan.md}, amendment A-F1).
Remark~\ref{rmk:relcoeff} and Remark~\ref{rmk:rawprimitive} are the prose form of
\file{Hecke/HeckeOperatorFamily.h}:136--173.

\paragraph{Formalization note.}
Proposition~\ref{prop:cosets} is elementary and constructive throughout: (1) is Hermite
reduction of a $2\times2$ integral matrix, (4) is a finite computation once $N$ is fixed
and a decidable predicate in general, and Definition~\ref{def:cosetreps} gives the
inverse of the bijection in closed form, so \code{P1(ZMod N)} together with the two
closed-form lifts is a complete formalisation of the coset datum with no choice
principle. Lemma~\ref{lem:cocycle} is then a one-line consequence of uniqueness of coset
representatives and is the natural first target. Proposition~\ref{prop:diamond} and
Proposition~\ref{prop:fullassembly} are proved \emph{at cochain level} from the content
decomposition and $\rhon(\lambda I)=\lambda^{n}I$, so neither needs cohomology, neither
needs a hypothesis on $\ell$, and both are formalisable before
\S\ref{sec:chainlevel} is. Theorem~\ref{thm:relations} is the one statement of this
chapter that is \emph{not} self-contained: it rests on the ring-homomorphism property
(Obligation~\ref{obl:algebra}), which rests on the classical comparison
(Theorem~\ref{thm:classical}); a formalisation would either import
$\mathcal{H}(\De,\Ga)$ and its structure theorem or --- the cheaper route for this model
--- verify the $\nu=2$ instance directly as a finite identity of cochain maps and induct,
which is what \code{HeckeRelationsValidation} does numerically.

%----------------------------------------------------------------------
\section{The primary definition: $\Tdc{g}$ in all degrees via a twisted chain map}\label{sec:chainlevel}
%----------------------------------------------------------------------
% SOURCE:  hecke_action.tex:570-786 (RH)
% NOTATION: rows nu, "t,j,s (running index)", sigma, d (coboundary), cring, Rel
%           of sec:notation-crosswalk applied.  Coset index j, chain-term index t,
%           never i (C15); tau is the chain datum and nothing else (C4).
% ERRATA:  E1 (which object the pipeline emits), E16 (App. A)

\paragraph{What this chapter fixes.}
The \emph{normative} definition of the Hecke operator: a $g$-twisted chain datum $\tau$
on the universal cover, the cochain map $\widetilde T$ built from $\tau$ and the coset
datum of \S\ref{sec:heckedatum}, and the operator $\Tdc{g}$ it induces in \emph{every}
cohomological degree at once. Also: what is and is not independent of the choices, at
which level; the exact condition under which the shipped $\tau$ is integral; and the
comparison with the classical operator of the theory of arithmetic groups. Degree-specific
recipes --- modular symbols, bar cochains --- appear in this document only as validation
oracles (\S\ref{sec:equivalence}), never as definitions.

\paragraph{Standing hypotheses.}
\ref{H1}--\ref{H4} of \S\ref{sec:conventions}, each repeated where it is used
(rule~\ref{A1}). The reader should note in advance the shape of the dependence, because
the superseded document stated it only in passing: \emph{the construction} of
$\widetilde T$ (Definition~\ref{def:hecke}) needs no hypothesis on $\ell$ at all;
\emph{the existence} of an abstract $\tau$ needs none either, over $\Z$ and hence over
every $\cring$ (Lemma~\ref{lem:tauexists}(1)); \emph{homotopy uniqueness} of $\tau$
needs $6\in\cring^{\times}$ (Lemma~\ref{lem:tauexists}(2)); and \emph{the identification
with the classical operator} (Theorem~\ref{thm:classical}) needs $6\in\cring^{\times}$
twice over. Remark~\ref{rmk:tauexistsscope} says which downstream statements that
apportionment does and does not move.

\subsection{The correspondence in cohomology}\label{sec:correspondence}

Topologically the datum of \S\ref{sec:heckedatum} is a correspondence. Write
$\mathrm{pr}'\colon\Ga'\backslash\Hyp\to\Ga\backslash\Hyp$ and
$\mathrm{pr}''\colon\Ga''\backslash\Hyp\to\Ga\backslash\Hyp$ for the two covering
projections (both of degree $m$) and $\theta_g\colon\Ga'\backslash\Hyp\to
\Ga''\backslash\Hyp$, $\Ga'z\mapsto\Ga''gz$, for the translation, well defined because
$g\Ga'g^{-1}=\Ga''$ (Proposition~\ref{prop:cosets}(3)); then
$\Tdc{g}=(\mathrm{pr}''\circ\theta_g)_{\ast}\circ(\mathrm{pr}')^{\ast}$. (The projections
are named operators rather than letters, per \ref{N3}: $\pi_1$ is the fundamental group
and $t_j$ is an upper-triangular transporter, \S\ref{sec:notation-framework} and
Definition~\ref{def:cosetreps}.) Algebraically --- and this is the form every later
statement uses ---
\begin{equation}\label{eq:corconjres}
  \Tdc{g}\;=\;\cor^{\Ga}_{\Ga''}\ \circ\ (\conj_g)_{\ast}\ \circ\ \res^{\Ga}_{\Ga'}
  \;\colon\; H^{\bullet}(\Ga;\VR)\longrightarrow H^{\bullet}(\Ga;\VR),
\end{equation}
where $(\conj_g)_{\ast}\colon H^{\bullet}(\Ga';\VR)\to H^{\bullet}(\Ga'';\VR)$ is induced by
the compatible pair $\bigl(\conj_g^{-1}\colon\Ga''\to\Ga',\ \rhon(g)\colon\VR\to\VR\bigr)$.
That this pair \emph{is} compatible is exactly the semigroup extension of
Definition~\ref{def:semigroup}: for $h''\in\Ga''$ one has $g^{-1}h''g\in\Ga'$ and, by
Proposition~\ref{prop:symconv}(2) applied to the integral identity
$g\,(g^{-1}h''g)=h''g$ in $\Mat_2(\Z)$,
\[
  \rhon(g)\,\rhon(g^{-1}h''g)\;=\;\rhon(h'')\,\rhon(g),
\]
with every factor integral and \emph{no} inverse of a determinant-$N$ matrix taken
(Obligation~\ref{obl:integral}). Equation~\eqref{eq:corconjres} is the standard
description of a double-coset operator
\STD{\cite[Ch.~III.9--10]{brown} for $\res$, $\cor$ and the conjugation action;
\cite[\S1.1, p.~194]{ash_stevens} for the double-coset operator (there in its right-module form on homogeneous cochains; v1.17.0-pre.5,
slot-17 R17-4);
\cite[Ch.~2]{harder3}}; it is recorded here as the object
Theorem~\ref{thm:classical} compares with, not as the definition, which is
Definition~\ref{def:hecke}.
\LEAN{thm:m6-cor-res} \LEAN{thm:m6-corH-indep} \LEAN{thm:m6-cor-res-index} \LEAN{thm:m6-Tcls-res} ($\cor^{\Ga}_{\Ga''}$ as Brown's transfer
on any projective resolution, independent of the transversal, $\cor\circ\res=[\Ga:\Ga'']$, and $\cor\circ(\conj_g)_{\ast}\circ\res$ for the pair
$(\conj_g^{-1},\rhon(g))$, in every degree and over every commutative ring; the topological correspondence form is not formalised; v1.24.0-rc.1)

\subsection{The twisted chain datum $\tau$}\label{sec:taudatum}

Throughout this subsection $C_{\bullet}(\Ku)$ is the simplicial chain complex of the
$\Ga$-tessellation $\Ku\cong\Xbulk$ of \ref{H4}, an augmented complex of
$\Z[\Ga]$-modules, acyclic because $\Xbulk$ is contractible (\ref{H4}(ii),
Proposition~\ref{prop:model7h4}); and
$\Ceq(K,\Rel;L)\cong\Hom_{\Z[\Ga]}\bigl(C_{\bullet}(\Ku),\VR\bigr)$ via the isomorphism
$\Psi$ of Theorem~\ref{thm:explicit_complex_equiv}.

\begin{definition}[$g$-twisted chain datum]\label{def:tau}
Let $\cring$ be a commutative ring and $g\in\De$, $\Ga''=\Ga\cap g\Ga g^{-1}$. A
\emph{$g$-twisted chain map over $\cring$} is a degree-preserving $\cring$-linear map
\[
  \tau\colon C_{\bullet}(\Ku)\otimes_\Z\cring\longrightarrow C_{\bullet}(\Ku)\otimes_\Z\cring
\]
such that
\begin{enumerate}
\item[(T1)] $\tau(h\cdot x)=(g^{-1}hg)\cdot\tau(x)$ for all $h\in\Ga''$ and all $x$ (the
  right-hand side is defined because $g^{-1}hg\in\Ga'\subseteq\Ga$ by
  Proposition~\ref{prop:cosets}(3));
\item[(T2)] $\partial\circ\tau=\tau\circ\partial$;
\item[(T3)] $\aug(\tau(x))=\aug(x)$ in degree $0$.
\end{enumerate}
A \emph{$g$-twisted chain datum} is such a $\tau$ together with the finite table of its
values on a set of $\Ga''$-orbit representatives of the simplices of $\Ku$, from which
(T1) determines it everywhere; that table is what Algorithm~\ref{alg:tau} emits and what
Obligation~\ref{obl:tau} certifies.
\end{definition}

\begin{remark}[The ring in Definition~\ref{def:tau}, and why it is not fixed to $\Fl$]
\label{rmk:taurings}
The superseded document stated Definition~\ref{def:tau} with $\cring=\Fl$ fixed. Nothing
in (T1)--(T3) uses that, and two consumers need more: the implementation's table is a
table of \emph{integer} chains ($\cring=\Z$), and the boundary and cone data of
\S\S\ref{sec:boundarydatum}--\ref{sec:cone} are required to be one integral datum whose
$\cring$-versions are its base changes (``tables once, matrices per $\cring$''). The
definition is therefore stated over a general $\cring$; the label
\code{def:tau} names the same notion and the $\Fl$ case is recovered verbatim by
$\cring=\Fl$, so this is a widening and not a retargeting (rule~\ref{L2}). Every
statement below says which ring it is over.
\end{remark}

\noindent
Geometrically $\tau$ is a simplicial approximation of $z\mapsto g^{-1}z$ on $\Hyp$, which
satisfies (T1) on the nose because $g^{-1}(hz)=(g^{-1}hg)(g^{-1}z)$. That is how
Algorithm~\ref{alg:tau} \emph{finds} a candidate; it is not how any statement about
$\tau$ is proved. Definition, existence, uniqueness and certification are purely
algebraic, and the geometry is discharged entirely by
Obligation~\ref{obl:tau} --- the point of that obligation being precisely that no
geometric assertion survives into the correctness of $T_p$.

\begin{lemma}[Existence, and homotopy uniqueness]\label{lem:tauexists}
Assume \ref{H4} and let $g\in\De$. The two halves carry \emph{different} hypotheses and
are therefore stated separately (rule~\ref{A1}); the apportionment is
Remark~\ref{rmk:tauexistsscope}.
\begin{enumerate}
\item \textbf{(Existence over $\Z$; \emph{no} hypothesis on $\ell$.)} \PROVED{} A
  $g$-twisted chain map over $\Z$ exists, and is produced by the terminating procedure
  (V), (E), (F) of the proof below (rule~\ref{A3}). Consequently a $g$-twisted chain map
  over $\cring$ exists for \emph{every} commutative ring $\cring$, namely the base change
  of one over $\Z$.
\item \textbf{(Homotopy uniqueness; needs $6\in\cring^{\times}$.)} \PROVED{} \LEAN{lem:tauexists-2-model} (on the model, for $\cring$-linear chain maps
  of the restricted resolutions $C_\bullet(\Xbulk)\otimes\cring$ over $\mathrm{id}_\cring$, including the last sentence: existence over such $\cring$, re-proved abstractly;
  not (1)) Let $\cring$
  be a commutative ring with $6\in\cring^{\times}$ --- for $\cring$ of characteristic
  $\ell$ this is $\ell\ge5$ (hypothesis \ref{H3}). Then any two $g$-twisted chain maps
  over $\cring$ are homotopic \emph{through $g$-twisted maps}: if $\tau,\tau'$ both
  satisfy (T1)--(T3) there is an additive degree-$(+1)$ map $H$ satisfying (T1) with
  $\tau-\tau'=\partial H+H\partial$. The argument re-proves existence over such an
  $\cring$ as well, abstractly; it is (1) that gives existence over $\Z$ and for every
  $\ell$. \LEAN{lem:m4-tauR-exists} (since v1.20.0: the existence over such $\cring$ for every $g\in\De$, in the form the operator
  on the model reads, a $\tau$ over $\cring$; (1), existence over $\Z$, still not)
\end{enumerate}
\end{lemma}

\begin{proof}[Proof of (1)]
Choose a set of $\Ga''$-orbit representatives of the simplices of $\Ku$: finitely many,
because $\Ga$ has finitely many orbits (\ref{H4}(iii)) and $[\Ga:\Ga'']=m<\infty$
(Proposition~\ref{prop:cosets}(3)). Define $\tau$ on them in ascending dimension and
extend by (T1) at each step, $\tau(h\cdot\sigma):=(g^{-1}hg)\cdot\tau(\sigma)$ for
$h\in\Ga''$ --- the conjugate lies in $\Ga$ by Proposition~\ref{prop:cosets}(3) --- which
is well defined exactly when $\tau(\sigma)$ is fixed by $g^{-1}\Ga''_{\sigma}g$,
$\Ga''_{\sigma}$ being the $\Ga''$-stabiliser of $\sigma$. Since $\Ga$ acts without
inversions (\ref{H4}(i)) no stabiliser reverses an orientation, so this determines an
additive map of the free $\Z$-modules $C_{q}(\Ku)$. \emph{Nothing below inverts a
stabiliser order or solves a system over anything but $\Z$.}
\begin{enumerate}
\item[(V)] \textbf{Vertices.} Let $v$ be a representative vertex. By \ref{H4}(i),(iv)
  $\Ga_{v}$ is cyclic of order $2$, $4$ or $6$ and contains $-I$; so is
  $\Ga''_{v}=\Ga''\cap\Ga_{v}$ (note $-I\in\Ga''$, since $-I\in\Ga$ and
  $g(-I)g^{-1}=-I$), and so is its conjugate $G:=g^{-1}\Ga''_{v}g\subseteq\Ga$. If
  $G=\{\pm I\}$, put $\tau(v):=w$ for an arbitrary vertex $w$ of $\Ku$; it is fixed,
  because $-I$ acts trivially (\ref{H4}(iv)). If $\abs{G}\in\{4,6\}$, a generator
  $\gamma$ of $G$ is a torsion element of $\Ga$ other than $\pm I$, hence elliptic, hence
  fixes a unique point $z_0\in\Hyp$, which $G=\langle\gamma\rangle$ therefore fixes. The
  stabiliser of $z_0$ modulo $\pm I$ is non-trivial, so $z_0$ is $\Ga$-equivalent to $i$,
  to $\zeta$ or to $\zeta^{2}$
  \STD{\cite[Ch.~VII, \S1.2, Thm.~1(1),(3)]{serre_cours}: every point of $\Hyp$ is
  $\Ga$-equivalent to a point of the standard fundamental domain, and the points there
  with non-trivial stabiliser modulo $\pm I$ are exactly $i$, $\zeta^{2}=e^{2\pi i/3}$
  and $\zeta=e^{\pi i/3}$}, i.e.\ to the vertex $v_6$, $v_5$ or $v_4$ of
  Definition~\ref{def:model7}; every $\Ga$-translate of those is a vertex of $\Ku$
  (\ref{H4}(i),(iv), their imaginary parts being $\le1<Y$). Put $\tau(v):=$ that vertex,
  which $G$ fixes. In both cases $\tau(v)$ is a \emph{single} vertex, so (T3) holds.
\item[(E)] \textbf{Edges.} Let $e$ be a representative edge, $\partial e=v''-v'$. Then
  $\tau(\partial e)=\tau(v'')-\tau(v')$ has augmentation $0$, hence is a boundary,
  because $\aug\colon H_0(\Ku)\to\Z$ is an isomorphism ($\Ku\cong\Xbulk$ is contractible,
  \ref{H4}(ii); the first of the three consequences collected at
  \eqref{eq:contractibleuse}). Let $\tau(e)$ be \emph{any} $1$-chain with
  $\partial\tau(e)=\tau(\partial e)$. The (T1)-extension is well defined because
  $\Ga''_{e}\subseteq\Ga_{e}=\{\pm I\}$ (\ref{H4}(iv)) and $g^{-1}(-I)g=-I$ acts
  trivially; and $\partial\tau=\tau\partial$ holds on the whole orbit because both sides
  satisfy (T1) and agree on $e$.
\item[(F)] \textbf{Triangles.} Let $\sigma$ be a representative triangle. By (T2) in
  degree $1$, $\partial\tau(\partial\sigma)=\tau(\partial\partial\sigma)=0$, so
  $\tau(\partial\sigma)$ is a $1$-cycle, hence a boundary, and the $2$-chain with that
  boundary is \emph{unique} ($Z_1(\Ku)=B_1(\Ku)$ and $Z_2(\Ku)=0$,
  \eqref{eq:contractibleuse}); take it. The (T1)-extension and (T2) are as in (E), the
  stabiliser of a triangle again being $\{\pm I\}$. There is nothing to do in degrees
  $\ge3$, $\Ku$ being two-dimensional (\ref{H4}(ii)).
\end{enumerate}
The result satisfies (T1)--(T3) over $\Z$, and those are identities between
$\Z$-linear maps of free $\Z$-modules, so they survive $\otimes_\Z\cring$.

This is the construction of Proposition~\ref{prop:strictpair} with its step (B) ---
which pre-assigns the boundary slots from $\taub$, and is what makes the result a map of
\emph{pairs} --- omitted; that proposition therefore contains (1), at the cost of the
further hypotheses \ref{B1}--\ref{B3}. Neither statement uses (2), so there is no
circularity.
\end{proof}

\begin{proof}[Proof of (2)]
Write $C_{\bullet}:=C_{\bullet}(\Ku)\otimes_\Z\cring$ and let $C^{(c)}_{\bullet}$ be the same
complex with the $\Ga''$-action $h\ast x:=(g^{-1}hg)\cdot x$, which is defined because
$g^{-1}hg\in\Ga$ (Proposition~\ref{prop:cosets}(3)) and is an action because $\conj_g$ is
a group isomorphism. Conditions (T1)--(T3) say precisely that $\tau\colon C_{\bullet}\to
C^{(c)}_{\bullet}$ is an augmentation-preserving map of complexes of
$\cring[\Ga'']$-modules.

\emph{The source is a complex of relatively $\cring[\Ga'']$-projective modules.} By
\ref{H4}(i),(iv) each $C_q(\Ku)$ is a direct sum of permutation modules
$\Z[\Ga/\Ga_\sigma]$ with $\abs{\Ga_\sigma}$ dividing $12$: the stabilisers occurring are
$\{\pm I\}$ of order $2$, $\langle S\rangle$ of order $4$ at $v_6$, and
$\langle ST\rangle,\langle TS\rangle$ of order $6$ at $v_4,v_5$. This --- the
\emph{orders $4$ and $6$ of the elliptic stabilisers} --- is where the hypothesis enters,
and it is the only place: $6\in\cring^{\times}$ makes $2,3$ and hence $4,6,12$ units, so
each $\abs{\Ga_\sigma}$ is invertible in $\cring$ and the averaging idempotent
$\abs{\Ga_\sigma}^{-1}\sum_{u\in\Ga_\sigma}u$ exhibits $\cring[\Ga/\Ga_\sigma]$ as a direct
summand of $\cring[\Ga]$. Restriction to the finite-index subgroup $\Ga''$ preserves
being a summand of a free module.

\emph{The target is acyclic and augmented.} $\Xbulk$ is contractible (\ref{H4}(ii),
Proposition~\ref{prop:model7h4}), so the augmented complex $C_{\bullet}(\Ku)\to\Z$ is
contractible as a complex of abelian groups, hence stays exact after $\otimes_\Z\cring$;
the $\Ga''$-action is changed by $\conj_g$ but the underlying complex is the same one, so
$C^{(c)}_{\bullet}$ is acyclic over $\cring$ and augmented over $\cring$ as well.

Existence and uniqueness up to $\cring[\Ga'']$-homotopy now follow from the fundamental
lemma of relative homological algebra
\STD{\cite[Ch.~I.7.4]{brown}: a map of augmented complexes of relatively projective
modules over an acyclic augmented target exists and is unique up to homotopy}. A homotopy
over $\cring[\Ga'']$ is exactly an additive $H$ with property (T1) --- the translation of
$\Ga''$-equivariance for the $\ast$-action into (T1) is the same one made for $\tau$
above.
\end{proof}

\begin{contractbox}
For $6\notin\cring^{\times}$ --- in particular $\ell\in\{2,3\}$ --- what is available is
\textbf{Lemma~\ref{lem:tauexists}(1) and not (2)}: a $g$-twisted chain datum \emph{exists},
over $\Z$ and hence over every $\cring$, but \textbf{homotopy uniqueness is not
asserted}, and this document does not repair it by an abstract argument. Also holding
there: the \emph{canonical} $\tau$ built by Algorithm~\ref{alg:tau} still exists (that
construction never divides by $2$ or $3$; its only divisions are exact-divisibility
checks in the integer fill) and is still certified by Obligation~\ref{obl:tau}.
Consequently, for $6\notin\cring^{\times}$ the operator is \textbf{defined} as the output
of the pipeline run with the canonical $\tau$; its independence of $\tau$
(Theorem~\ref{thm:welldef}(b)) and its identification with \eqref{eq:corconjres}
(Theorem~\ref{thm:classical}) are asserted only for $6\in\cring^{\times}$. No
cross-$\tau$ comparison is made anywhere at $\ell\in\{2,3\}$, in this document or in the
implementation.
\end{contractbox}

\begin{remark}[The two halves apportioned: what (1) buys, and what it does not]
\label{rmk:tauexistsscope}
\PROVED{} (Rule~\ref{A6}.) Up to and including v0.15.0 Lemma~\ref{lem:tauexists} carried
$6\in\cring^{\times}$ on \emph{both} halves, and the contract box above declared the
whole lemma unavailable at $\ell\in\{2,3\}$. Clause (1) removes the hypothesis from the
existence half --- by a construction, not by the comparison theorem. Three things must be
kept apart from it.
\begin{enumerate}
\item \textbf{Uniqueness does not follow, and the argument does not transfer.} The proof
  of (1) may choose a value \emph{and then} impose (T1), because a \emph{vertex} value
  can be chosen inside the fixed-point set of the twisted stabiliser (Serre's elliptic
  points) and because edge and triangle stabilisers are central. A twisted homotopy $H$
  has no such freedom: $H(v)$ must be a $1$-chain with a \emph{prescribed} boundary
  \emph{and} be fixed by $g^{-1}\Ga''_{v}g$, and for $\abs{\Ga''_{v}}\in\{4,6\}$ nothing
  in the construction produces an invariant fill. Supplying one is exactly what the
  averaging idempotent $\abs{\Ga_\sigma}^{-1}\sum_{u}u$ does in the proof of (2), and it
  is why (2) keeps $6\in\cring^{\times}$. Whether homotopy uniqueness \emph{holds} at
  $\ell\in\{2,3\}$ is \CONJ{} either way and is not decided here; what would settle it is
  a twisted-equivariant fill at the elliptic vertices, or a counterexample --- two
  certified $\tau$ at $\ell=2$ inducing different $\Amat{q}$, which
  \code{HeckeIotaTauSymmetryValidation} already has the two data for
  (Remark~\ref{rmk:iotatausym}) --- and on the golden grid those two, the legacy and
  the canonical $\tau$, induce the \emph{same} $\Amat{q}$ at every $\ell\in\{2,3\}$
  point (the v1.1.0 re-baseline moved no cohomology-level value), so no counterexample
  is known there.
\item \textbf{No downstream hypothesis is weakened by (1).} Every statement of this
  document that carries $6\in\cring^{\times}$ ``via Lemma~\ref{lem:tauexists}'' uses the
  \emph{uniqueness} half, or relative projectivity of $C_{\bullet}(\Ku)\otimes\cring$, or
  both; each was read against its own proof and \textbf{none} of them moves:
  Theorem~\ref{thm:welldef}(b), and through it (c) for arbitrary admissible choices;
  Proposition~\ref{prop:iotacommutes}(3c) (which is (3a) plus
  Theorem~\ref{thm:welldef}(a),(b)); Theorem~\ref{thm:classical}, which uses the
  hypothesis twice and whose \emph{first} use --- relative projectivity, hence also
  $\Heq{\bullet}\cong H^{\bullet}(\Ga;\VR)$ (Remark~\ref{rmk:orbifold}, \ref{H3}) --- is
  untouched by (1); Proposition~\ref{prop:barformula}'s agreement clause, which is a
  comparison of two resolutions; Proposition~\ref{prop:coneindep}(2), and hence
  Theorem~\ref{thm:conewelldef}(c) and (e); and
  Remark~\ref{rmk:conesteinberg}(i),(ii), and hence
  Theorem~\ref{thm:conewelldef}(c$'$). \textbf{What (1) does buy} is that
  Definition~\ref{def:hecke} is non-vacuous over $\Z$ for every $\ell$ --- before it, the
  only datum known to exist at $\ell\in\{2,3\}$ was the output of
  Algorithm~\ref{alg:tau} --- and the corresponding correction inside
  Remark~\ref{rmk:tauintegral}.
\item \textbf{Existence over $\Z$ is not integrality of the canonical table.} (1)
  produces \emph{some} datum over $\Z$; it says nothing about whether
  Algorithm~\ref{alg:tau}'s triangle fills solve over $\Z$, which is a property of that
  algorithm's choices, is verified per table, and is governed by owner decision OP2
  (Remark~\ref{rmk:tauintegral}). The two must not be conflated --- that conflation is
  the class of over-claim Appendix~\ref{app:errata} records at entry~\ref{err:e3}.
\end{enumerate}
\end{remark}

\begin{remark}[Integrality of the emitted $\tau$-table: the claim and its condition]
\label{rmk:tauintegral}
Algorithm~\ref{alg:tau} emits, for each orbit representative, a chain
$\tau(s_j\cdot\lift{\sigma})=\sum_{t}n_t\,\gamma_t\cdot\lift{\sigma_t}$ with
$n_t\in\Z$, $\gamma_t\in\Ga$. Steps 1 and 2 are integral \emph{unconditionally}, because
neither solves a linear system: a vertex value is the single term
$1\cdot\gamma_{\mathrm{red}}^{-1}\lift{w}$, and an edge value is a concatenation of
stored edge-path templates, hence a $\Z$-combination of oriented edges obtained by
addition alone. Step 3, the triangle fill, solves
$\partial_2|_{\mathcal N}x=\tau(\partial\lift{\sigma})$ over $\Z$ on a finite subcomplex
$\mathcal N\subseteq\Ku$, growing $\mathcal N$ on failure; that solve is where
integrality can be lost, because the algorithm permits, as a fallback, solving the
\emph{same} system mod $\ell$.
\begin{itemize}
\item \textbf{What is claimed.} If every fill of a table solved over $\Z$, the table is a
  $g$-twisted chain datum over $\cring=\Z$ in the sense of Definition~\ref{def:tau}, and
  the datum over any $\cring$ --- in particular every $\Fl$, every $\Z/\ell^{e}$, every
  $\Z/m$ --- is its base change; the certification of Obligation~\ref{obl:tau} is then
  exact over $\Z$. This is verifiable per table, and it is verified: the regime is
  recorded per chain and the table-level disjunction is exposed
  (\code{used\_mod\_ell\_solve()}).
\item \textbf{What is \emph{not} claimed.} That the fills always solve over $\Z$. There
  is no theorem here. Lemma~\ref{lem:tauexists}(1) does give existence of \emph{some}
  $g$-twisted chain datum over $\Z$, for every $\ell$ --- but by a different
  construction, which chooses its vertex values at the elliptic points and its triangle
  fills anywhere in $\Ku$, whereas Algorithm~\ref{alg:tau} chooses by nearest vertex and
  fills inside a growing finite neighbourhood $\mathcal N$. \textbf{Existence of an
  integral datum is not integrality of \emph{this} datum}: the $\Z$-fill's feasibility on
  a given $\mathcal N$ is a property of the table, not of the model
  (Remark~\ref{rmk:tauexistsscope}(3)). A chain obtained by the
  mod-$\ell$ fallback defines $\tau$ over $\Fl$ only, and Obligation~\ref{obl:tau}'s
  certification for every identity touching it is then exact mod $\ell$ rather than over
  $\Z$ --- which is sound, and is a strictly weaker datum.
\item \textbf{Policy (owner decision OP2, 2026-09-15).} The mod-$\ell$ fill fallback is
  \textbf{forbidden by default}. It exists only because of a fixed-width integer tier and
  never for a mathematical reason, so the remedy is to widen that tier; the fallback stays
  reachable behind an explicit setting and a datum that uses it is marked
  \code{integral = false}. \textbf{\code{integral = false} is a recorded failure to reach
  $\Z$, not the default escape from a hard fill}, and a datum claiming integrality must
  have solve regime $\Z$ on every one of its tables.
\item \COMP{} On the shipped grid --- $(p,\nu)\in\{(2,1),(3,1),(5,1),(7,1),(2,2)\}$,
  $\ell\in\{5,7,691,134217757\}$, $n\in\{0,2,10\}$, plus the bad-prime points
  $p=\ell\in\{2,5\}$ at the same $n$, i.e.\ $66$ points --- \textbf{$0$ of $66$ used the
  mod-$\ell$ triangle solve}: every table is integral, with certified entry bound
  $B\le196$ (Lemma~\ref{lem:bounds}). Validation
  \code{HeckeTauTableSelfCheckValidation}, $69/69$ points passed, run of 2026-09-15 on
  the \code{RelWithDebInfo} build.
\end{itemize}
\end{remark}

\subsection{Definition of the operator}\label{sec:operatordef}

For $\sigma\in\orient{q}(K)$ write $\lift{\sigma}\in\orient{q}(\Ku)$ for its
distinguished lift under the strict fundamental domain $K\subseteq\Ku$ of \ref{H4}(iii)
(macro \code{\textbackslash lift}, \S\ref{sec:notation-framework}); the implementation
writes \code{i(sigma)} for $\lift{\sigma}$, a spelling this document cannot use because
$i$ is $\sqrt{-1}$ (\S\ref{sec:notation-collisions}, C15, C16).

\begin{definition}[Normative definition of $\Tdc{g}$]\label{def:hecke}
\LEAN{def:heckeP} (the definition \eqref{eq:heckechain}, and $\delta_k=s_k^{-1}g$ for the canonical $s_k$; since v1.19.0: $\tau$ over $\cring$ is formalised
on the model (\texttt{TauR}, \texttt{heckeTildeR}), [GT]'s operator on a general descriptor takes $\tau$ over $\Z$)
Let $g\in\De$, $\Ga''=\Ga\cap g\Ga g^{-1}$ of index $m$, let $s_1,\dots,s_m$ represent
$\Ga''\backslash\Ga$ and $\delta_j:=s_j^{-1}g$ (Proposition~\ref{prop:cosets}(5);
canonically Definition~\ref{def:cosetreps}) \LEAN{lem:m4-cosetdata} (since v1.20.0: such data exist for every $g\in\De$ ---
$\Ga''$ of finite index, a transversal, $\delta_j\in\De$; a non-constructive datum, not the canonical table of
Definition~\ref{def:cosetreps}), and let $\tau$ be a $g$-twisted chain datum
over $\cring$ (Definition~\ref{def:tau}). For
$\phi\in\Hom_{\Z[\Ga]}\bigl(C_{\bullet}(\Ku),\VR\bigr)$ define
\begin{equation}\label{eq:heckechain}
  (\widetilde T\phi)(x)\;:=\;\sum_{j=1}^{m}\ \rhon(\delta_j)\,\phi\bigl(\tau(s_j\cdot x)\bigr),
  \qquad x\in C_q(\Ku),\ q\ge0 .
\end{equation}
Then $\Tdc{g}$ on $\Cq{q}(K,\Rel;L)$ is $\Psi\circ\widetilde T\circ\Psi^{-1}$
(Theorem~\ref{thm:explicit_complex_equiv}), and on $\Heq{q}(K,\Rel;L)$ the induced map.
The sum is finite, all its factors are integral, and no determinant-$N$ matrix is
inverted (Obligation~\ref{obl:integral}); the definition is uniform in $q$ --- one
formula, one pipeline, all degrees --- and needs \emph{no} hypothesis on $\ell$.
\end{definition}

\begin{theorem}[Well-definedness]\label{thm:welldef}
Assume \ref{H4} and let $g,\{s_j\},\tau,\cring$ be as in Definition~\ref{def:hecke}.
\begin{itemize}
\item[] \PROVED{} \textbf{(No hypothesis on $\ell$.)} \LEAN{thm:welldef-lean} (equivariance and the chain-map property; the descent to $\Heq{q}$ is the generic one) $\widetilde T$ carries
  $\Ga$-equivariant cochains to $\Ga$-equivariant cochains and commutes with the
  coboundary; hence it descends to $\Heq{q}(K,\Rel;L)$ for every $q$.
\end{itemize}
Moreover the induced map on $\Heq{q}$ is independent of
\begin{enumerate}
\item[(a)] the choice of representatives $s_j$ --- \PROVED{} \textbf{with no hypothesis on
  $\ell$}, and in fact at cochain level, \emph{provided $\tau$ is held fixed}; \LEAN{thm:welldef-lean} (at cochain level, $\tau$ fixed; since v1.19.0 for any two choices of representatives)
  \LEAN{prop:fullassembly-model} (since v1.20.0 also for $\tau$ over $\cring$, on the model, at cochain level)
\item[(b)] the choice of $\tau$ --- \PROVED{} under $6\in\cring^{\times}$ \LEAN{thm:welldef-b} (on the model, the representatives held fixed;
  $\tau,\tau'$ over $\cring$, as $\cring$-linear maps of $C_q(\Xbulk)\otimes\cring$ satisfying (T1)--(T3); on $\Heq{n}(K,\Rel;L)$ itself, every $n\ge0$)
  (hypothesis \ref{H3}; via the \emph{uniqueness} half Lemma~\ref{lem:tauexists}(2),
  which is where the hypothesis is used and the only place --- existence,
  Lemma~\ref{lem:tauexists}(1), needs none);
\item[(c)] the choice of $g$ inside $\Ga g\Ga$ --- \PROVED{} \LEAN{thm:welldef-c} (since v1.19.0 the whole of (c): the transported half with no hypothesis on $\ell$;
  the general case on the model, under $6\in\cring^{\times}$, with $\tau_1$ over $\Z$) at cochain level for the
  \emph{transported} choices constructed in the proof, with no hypothesis on $\ell$; and
  for arbitrary admissible choices under $6\in\cring^{\times}$, via (a) and (b).
\end{enumerate}
None of (a)--(c) is claimed to fix the cochain map $\widetilde T$ itself: each is a
statement about $\Heq{q}$, and Remark~\ref{rmk:indeplevel} says exactly what fails at
cochain level and what a validation may therefore compare.
\end{theorem}

\begin{proof}
\emph{Equivariance.} Let $\gamma\in\Ga$. Write $s_j\gamma=h_j(\gamma)\,s_{j\cdot\gamma}$
with $h_j(\gamma)\in\Ga''$ as in \eqref{eq:cocycle} (Lemma~\ref{lem:cocycle}). Using
(T1), $\Ga$-equivariance of $\phi$, and $g^{-1}h_j(\gamma)g\in\Ga'\subseteq\Ga$:
\[
  \rhon(\delta_j)\,\phi\bigl(\tau(s_j\gamma x)\bigr)
  =\rhon(\delta_j)\,\rhon\bigl(g^{-1}h_j(\gamma)g\bigr)\,\phi\bigl(\tau(s_{j\cdot\gamma}x)\bigr)
  =\rhon\bigl(s_j^{-1}h_j(\gamma)g\bigr)\,\phi\bigl(\tau(s_{j\cdot\gamma}x)\bigr),
\]
by Proposition~\ref{prop:symconv}(2) and $\delta_j\,g^{-1}h_j(\gamma)g=s_j^{-1}h_j(\gamma)g$
as an identity in $\Mat_2(\Z)$. Since $s_j^{-1}h_j(\gamma)=\gamma s_{j\cdot\gamma}^{-1}$
by \eqref{eq:cocycle}, the right-hand side is
$\rhon(\gamma)\,\rhon(\delta_{j\cdot\gamma})\,\phi\bigl(\tau(s_{j\cdot\gamma}x)\bigr)$.
Summing over $j$ and using that $j\mapsto j\cdot\gamma$ is a bijection
(Lemma~\ref{lem:cocycle}) gives
$(\widetilde T\phi)(\gamma x)=\rhon(\gamma)(\widetilde T\phi)(x)$.

\emph{Chain map.} With $(d\phi)(x)=\phi(\partial x)$, property (T2) and additivity give
$\widetilde T(d\phi)=d(\widetilde T\phi)$ term by term.

\emph{(a).} The general change of admissible representatives is $s_j\mapsto h_js_j$ with
$h_j\in\Ga''$. Holding $\tau$ fixed, the $j$-th summand of \eqref{eq:heckechain} becomes
\[
  \rhon\bigl(s_j^{-1}h_j^{-1}g\bigr)\,\phi\bigl(\tau(h_js_jx)\bigr)
  =\rhon\bigl(s_j^{-1}h_j^{-1}g\bigr)\,\rhon\bigl(g^{-1}h_jg\bigr)\,\phi\bigl(\tau(s_jx)\bigr)
  =\rhon(\delta_j)\,\phi\bigl(\tau(s_jx)\bigr),
\]
unchanged summand by summand; so for fixed $\tau$ the cochain map is literally the same
map. \textbf{The proviso is not cosmetic.} In an admissible construction $\tau$ is built
\emph{relative to} the chosen $s_j$ (Algorithm~\ref{alg:tau} stores one chain per pair
$(j,\sigma)$ and derives the rest through (T1)), so a change of representatives rebuilds
$\tau$; the summand-by-summand invariance then does not apply and only the induced map on
$\Heq{q}$ is invariant, via (b). This is why (a) is stated at the $\Heq{q}$ level.

\emph{(b).} Let $\tau'=\tau+\partial H+H\partial$ with $H$ as in
Lemma~\ref{lem:tauexists}(2), and set
$(\widetilde H\phi)(x):=\sum_{j}\rhon(\delta_j)\phi\bigl(H(s_jx)\bigr)$. Then
$\widetilde T'-\widetilde T=d\circ\widetilde H+\widetilde H\circ d$, and $\widetilde
H\phi$ is $\Ga$-equivariant by the computation above, which used only property (T1) of
its argument --- which $H$ has. So the two cochain maps are cochain homotopic and induce
the same map on $\Heq{\bullet}$.

\emph{(c).} Let $g_1=\gamma_1g\gamma_2$ with $\gamma_1,\gamma_2\in\Ga$. Then
$\Ga''_1:=\Ga\cap g_1\Ga g_1^{-1}=\gamma_1\Ga''\gamma_1^{-1}$, the elements
$\gamma_1s_j$ represent $\Ga''_1\backslash\Ga$ ($\bigsqcup_j\Ga''_1\gamma_1s_j=\gamma_1\bigsqcup_j\Ga''s_j=\Ga$), with
$(\gamma_1s_j)^{-1}g_1=\delta_j\gamma_2$, and
$\tau_1:=\gamma_2^{-1}\cdot\tau(\gamma_1^{-1}\cdot(-))$ is $g_1$-twisted: for
$h_1\in\Ga''_1$, using $\gamma_1^{-1}h_1\gamma_1\in\Ga''$ and
$g_1^{-1}h_1g_1=\gamma_2^{-1}\bigl(g^{-1}\gamma_1^{-1}h_1\gamma_1g\bigr)\gamma_2$,
\[
  \tau_1(h_1x)=\gamma_2^{-1}\cdot\tau\bigl((\gamma_1^{-1}h_1\gamma_1)\,\gamma_1^{-1}x\bigr)
  =\gamma_2^{-1}\bigl(g^{-1}\gamma_1^{-1}h_1\gamma_1g\bigr)\cdot\tau(\gamma_1^{-1}x)
  =(g_1^{-1}h_1g_1)\cdot\tau_1(x),
\]
and (T2), (T3) are inherited. Substituting $(g_1,\{\gamma_1s_j\},\tau_1)$ into
\eqref{eq:heckechain} reproduces the operator for $(g,\{s_j\},\tau)$ verbatim:
$\rhon(\delta_j\gamma_2)\,\phi\bigl(\tau_1(\gamma_1s_jx)\bigr)=\rhon(\delta_j)\rhon(\gamma_2)\,\phi\bigl(\gamma_2^{-1}\cdot\tau(s_jx)\bigr)=\rhon(\delta_j)\,\phi\bigl(\tau(s_jx)\bigr)$,
$\gamma_1$ cancelling inside $\tau_1$'s argument and $\rhon(\gamma_2)$ against $\Ga$-equivariance of $\phi$. That
identity is unconditional. (\textbf{Erratum, v1.18.0}, confirmed by the owner on 2026-10-03; found by slot-18 L18-2, the Lean proof
\texttt{GL2.Model.thetaG\_transported}: through v1.18.0-pre.1 this paragraph took $\{s_j\gamma_1^{-1}\}$ as the representatives of $\Ga''_1\backslash\Ga$ --- not a
transversal in general (its cosets $\Ga''_1s_j\gamma_1^{-1}=\gamma_1\bigl(\Ga''\gamma_1^{-1}s_j\bigr)\gamma_1^{-1}$ are distinct iff $\{\gamma_1^{-1}s_j\}$ represents $\Ga''\backslash\Ga$, which fails in general), and with $\tau_1$ its summands read $\tau(\gamma_1^{-1}s_j\gamma_1^{-1}x)$,
not $\tau(s_jx)$ --- and said that $\rhon(\gamma_1)$ and $\rhon(\gamma_2)$ both cancel against equivariance. The statement~(c) and its conclusion are unchanged. The slot-18 Lean audit F3 confirmed the erratum with a counterexample: $\Ga''=\Ga_0(p)$, $s_j=\smat{1}{0}{j}{1}$ ($j<p$), $s_p=S$, $\gamma_1=S$; the class of $\Ga_0(p)x$ is the bottom row of $x$ in $\mathbb P^1(\F_p)$, and $(-1,0)\smat{1}{0}{j}{1}=(-1,0)$ for every $j$, so the $p$ elements $\gamma_1^{-1}s_j$, $j<p$, lie in one coset.) To pass from it to independence for an \emph{arbitrary}
admissible datum at $g_1$ one must compare that datum with the transported one, which is
(a) and (b) --- hence the hypothesis $6\in\cring^{\times}$ in the general case.
\end{proof}

\begin{remark}[What (T3) is for]\label{rmk:T3}
Property (T3) is used in degree $0$ and nowhere else, and it is what pins the classical
normalisation there. A $0$-cocycle corresponds to $\phi$ with $\phi(\tilde v)=v$ for every
vertex $\tilde v$ of $\Ku$ and a fixed $v\in\VR^{\Ga}$; (T3) forces
$\phi(\tau(s_j\tilde v))=v$, so \eqref{eq:heckechain} gives
\[
  \Tdc{g}\,v=\sum_{j=1}^{m}\rhon(\delta_j)\,v ,
\]
the standard double-coset action on invariants --- for $n=0$, multiplication by
$m$, e.g.\ $(p+1)v$ for $g=\mathrm{diag}(1,p)$. Without (T3) the degree-$0$ operator would
still be well defined but would not be normalised, and the $H^0$ eigenvalue oracles of
\S\ref{sec:example} would be off by the augmentation defect.
\end{remark}

\begin{remark}[The level at which each independence holds]\label{rmk:indeplevel}
All three independence claims of Theorem~\ref{thm:welldef} --- coset representatives (a),
twisted datum (b), double-coset representative (c) --- are claims about
\emph{cohomology}: the induced matrix $\Amat{q}$ on $\Heq{q}$ is invariant, bit for bit in
every tested setting whenever $6\in\cring^{\times}$. The cochain map $\Th_g$ and its flat
matrix $\Th^{(q)}_{\mathrm{flat}}$ are \emph{not} invariant under any of them:
$\Th^{(q)}_{\mathrm{flat}}$ depends on the chosen $\tau$, and in an admissible
construction $\tau$ is coupled to the $s_j$ (Algorithm~\ref{alg:tau} builds $\tau$ from
the representatives), so even a relabelling $s_j\mapsto h_js_j$ with a non-trivial
$h_j\in\Ga_0(N)$ rebuilds $\tau$ and in general changes
$\Th^{(q)}_{\mathrm{flat}}$. The summand-by-summand invariance in the proof of (a) applies
only when $\tau$ is held fixed, which the coupled construction does not do. Consequently:
\begin{itemize}
\item \textbf{Variants in $s_j$, $\tau$ or $g$ are compared at the $\Amat{q}$ level,
  never bit for bit at cochain level.} A validation that compares
  $\Th^{(q)}_{\mathrm{flat}}$ across such a variant is testing a statement this document
  does not make.
\item \textbf{Cochain-level bit-identity is asserted only between genuinely identical
  constructions}: the trivial relabelling $h_j=I$ (same $s_j$, same $\tau$); the two
  evaluation routes $\omega_{\bullet}\circ\Xig$ and direct \eqref{eq:heckechain} on the
  \emph{same} data (Obligation~\ref{obl:agreement}); $\omega_{\bullet}\circ
  \mu_{\bullet}\circ\alpha_{\bullet}$ against $\Th_g$ (Proposition~\ref{prop:mu}(3)),
  because $\mu_{\bullet}$ and $\alpha_{\bullet}$ are likewise built from the \emph{same}
  $g$, the \emph{same} $s_j$ and the \emph{same} $\tau$-table
  (Obligation~\ref{obl:agreement}); and the \code{HECKE\_PURE\_INTEGER} on-versus-off
  determinism check (Obligation~\ref{obl:fp}).
\end{itemize}
This remark is a statement about what may be asserted, not a diagnostic to branch on
(rule~\ref{A6}).
\end{remark}

\subsection{Comparison with the classical operator}\label{sec:classicalcomparison}

\begin{theorem}[Comparison with the classical operator]\label{thm:classical}
Assume \ref{H4} and $6\in\cring^{\times}$ --- for $\cring$ of characteristic $\ell$,
$\ell\ge5$ (hypothesis \ref{H3}). Then under the isomorphisms
$\Heq{q}(K,\Rel;L)\cong H^q(\Ga;\VR)$ of Theorem~\ref{thm:explicit_complex_equiv} and
Remark~\ref{rmk:orbifold}, the operator of Definition~\ref{def:hecke} equals
\eqref{eq:corconjres}: it is the standard double-coset Hecke operator of the theory of
arithmetic groups
\STD{\cite[Ch.~III.9--10]{brown} for $\res$, $\cor$, the conjugation action and the
transfer formula; \cite[\S1]{ash_stevens} and \cite[Ch.~2--3]{harder3} for the
double-coset operator on the cohomology of arithmetic groups; \cite[Ch.~8]{shimura} and
\cite{deng} for the classical modular-forms normalisation}.
\LEAN{thm:m6-classical-theta} \LEAN{thm:m6-Tcls-res} \LEAN{thm:m6-hecke-res} (on the $7$-vertex model, for every commutative ring with $6\in\cring^{\times}$,
every representation $\rho$ of $\De$ (in particular $\rhon$), every $g\in\De$, coset datum, $\tau$ and degree: $\Theta\circ\Tdc{g}=
\cor\circ(\conj_g)_{\ast}\circ\res\circ\Theta$ (the right side is M5's \texttt{Tcls}, identified with $\cor\circ(\conj_g)_{\ast}\circ\res$ by [GT]~\texttt{thm:m6-Tcls-res}) with
$\Theta$ the comparison of Remark~\ref{rmk:orbifold} (mathlib's $\mathrm{Ext}$-built one, equal to the comparison induced by any lift of the identity,
[GT]~\texttt{thm:m6-theta-ext}); the operator does not depend on the projective resolution or on $\tau$; the classical modular-forms normalisation of
Shimura / Deng is not formalised; v1.24.0-rc.1)
\end{theorem}

\begin{proof}[Proof sketch]
The hypothesis $6\in\cring^{\times}$ is used twice, and it is worth separating the uses.
First, it makes $C_{\bullet}(\Ku)\otimes\cring$ a complex of relatively projective
$\cring[\Ga]$-modules (the computation in the proof of Lemma~\ref{lem:tauexists}(2)), so that
this one complex computes $H^{\bullet}$ of $\Ga$ \emph{and} of each of its finite-index
subgroups $\Ga',\Ga''$, and that $\Heq{\bullet}(K,\Rel;L)\cong H^{\bullet}(\Ga;\VR)$ at
all (Remark~\ref{rmk:orbifold}). Second, it supplies the homotopy uniqueness of
Lemma~\ref{lem:tauexists}(2), which is what makes the middle map below independent of the
choice of $\tau$. \emph{Neither} use is met by Lemma~\ref{lem:tauexists}(1), so this
theorem does not weaken (Remark~\ref{rmk:tauexistsscope}(2)). On that resolution the three constituents of \eqref{eq:corconjres} are:
\begin{itemize}
\item $\res^{\Ga}_{\Ga'}$ is the inclusion
  $\Hom_{\Ga}(C_{\bullet},\VR)\hookrightarrow\Hom_{\Ga'}(C_{\bullet},\VR)$;
\item $(\conj_g)_{\ast}\colon\Hom_{\Ga'}(C_{\bullet},\VR)\to\Hom_{\Ga''}(C_{\bullet},\VR)$
  is $\psi\mapsto\rhon(g)\circ\psi\circ\tau$: by (T1), $\tau$ is exactly a chain map over
  $\conj_g^{-1}\colon\Ga''\to\Ga'$, so $\rhon(g)\psi\tau$ is $\Ga''$-equivariant by the
  computation in the proof of Theorem~\ref{thm:welldef}, and by
  Lemma~\ref{lem:tauexists}(2) any two choices of $\tau$ give homotopic maps, i.e.\ the
  map induced in cohomology is the canonical one;
\item $\cor^{\Ga}_{\Ga''}$ is the transfer
  $\psi\mapsto\bigl(x\mapsto\sum_{j}\rhon(s_j)^{-1}\psi(s_jx)\bigr)$
  \STD{\cite[Ch.~III.9]{brown}}, well defined and independent of the $s_j$ by the
  computation in the proof of Theorem~\ref{thm:welldef}(a) with $g$ replaced by $I$.
\end{itemize}
Composing, for $\Ga$-equivariant $\phi$,
\[
  x\ \longmapsto\ \sum_{j}\rhon(s_j)^{-1}\rhon(g)\,\phi\bigl(\tau(s_jx)\bigr)
  \;=\;\sum_{j}\rhon(s_j^{-1}g)\,\phi\bigl(\tau(s_jx)\bigr)\;=\;(\widetilde T\phi)(x),
\]
which is \eqref{eq:heckechain}; the middle equality is
Proposition~\ref{prop:symconv}(2) together with
$\rhon(s_j)^{-1}=\rhon(s_j^{-1})$, legitimate because $s_j\in\Ga$ is unimodular.
\end{proof}

\begin{remark}[Which object of \S\ref{sec:heckedatum} this chapter defines]
\label{rmk:whichoperator}
Definition~\ref{def:hecke} defines $\Tdc{g}$ for \emph{one} double coset. Fed
$g=\mathrm{diag}(1,p^{\nu})$ --- which is what the coset tables of
Definition~\ref{def:cosetreps} enumerate --- it therefore produces the \textbf{primitive}
operator $T_{\Deprim_{p^{\nu}}}$, and for $\nu\ge2$ that is \emph{not} $T_{p^{\nu}}$: the
full operator is the sum \eqref{eq:fullop}, assembled by \eqref{eq:fullassembly}
(Remark~\ref{rmk:rawprimitive}; Appendix~\ref{app:errata}, entry~\ref{err:e1}). At
$\nu=1$, and hence for $T_p$ and $U_\ell$, the distinction is empty.
\end{remark}

\paragraph{Code map.}
Definition~\ref{def:tau}, Lemma~\ref{lem:tauexists}, Remark~\ref{rmk:tauintegral}:
\file{Hecke/HeckeTauTable.h/.cpp} --- the stored \code{ChainTerm} triples
$(n_t,\gamma_t,\sigma_t)$, one chain per pair $(j,\sigma)$, with every
$\Ga''$-related slot derived through the single shared (T1) helper
\code{eval\_tau\_on\_translate} used by \emph{both} construction and certification;
the solve regime is recorded per chain and reported by \code{used\_mod\_ell\_solve()}.
The constructor certifies (Obligation~\ref{obl:tau}) and throws on failure, so an
uncertified table cannot exist; this is the implementation's invariant I8,
``certification over geometry'', and it is the code-level form of the paragraph after
Definition~\ref{def:tau}. Gated by \code{HeckeTauTableValidation},
\code{HeckeTauTableSelfCheckValidation}. Definition~\ref{def:hecke},
\eqref{eq:heckechain}: \file{Hecke/HeckeOperatorOnK.h/.cpp}, whose fused route
$\Th_g=\omega_{\bullet}\circ\Xig$ and whose second, independently bracketed route
evaluate \eqref{eq:heckechain} and are asserted bit-identical
(\S\ref{sec:induced}, Obligation~\ref{obl:agreement}); the induced
$\Amat{q}$ is \code{get\_H\_matrices()}. Theorem~\ref{thm:welldef} and
Remark~\ref{rmk:indeplevel}: \code{HeckeIndependenceValidation}, which varies the $s_j$
(through \code{HeckeCosetTables}' custom-representative constructor), $\tau$ and $g$ and
compares at the $\Amat{q}$ level only --- exactly as Remark~\ref{rmk:indeplevel}
requires. Theorem~\ref{thm:classical}: cross-checked, not assumed, by
\code{HeckeCrossChecks} and the bar-cocycle route of \S\ref{sec:equivalence}
(Obligation~\ref{obl:agreement}); the $\ell\in\{2,3\}$ arms make no cross-$\tau$ and no
classical-comparison assertion, per the contract box after
Lemma~\ref{lem:tauexists}.

\paragraph{Formalization note.}
Definition~\ref{def:tau} is three equational conditions on a map of complexes and is
directly expressible. The two halves of Lemma~\ref{lem:tauexists} want different
treatments. For (2) the useful reformulation is the one used in its proof --- $\tau$ is a
morphism of augmented complexes of $\cring[\Ga'']$-modules $C_{\bullet}\to
C^{(c)}_{\bullet}$ --- since \code{mathlib4}'s homological algebra then supplies the
comparison theorem; there the one hypothesis a formalisation must not lose is the
source's relative projectivity, whose entire content here is that $4$ and $6$ --- the
orders of the elliptic stabilisers, \ref{H4}(iv) --- are units, and stating (2) for a
general $\cring$ with $6\in\cring^{\times}$, rather than for $\Fl$ with $\ell\ge5$, makes
that visible in the statement instead of in the proof. Clause (1) needs \emph{none} of
that machinery: it is a terminating construction over $\Z$ on a finite set of orbit
representatives whose only non-combinatorial inputs are \eqref{eq:contractibleuse} and
the elliptic-point classification of \cite{serre_cours}, and it is therefore the cheaper
of the two to formalise and the one to do first. Definition~\ref{def:hecke} is a finite sum over a
finite index set and is constructive given the tables of
Definition~\ref{def:cosetreps} and Algorithm~\ref{alg:tau}; the equivariance half of
Theorem~\ref{thm:welldef} is a finite manipulation using only
Lemma~\ref{lem:cocycle} and Proposition~\ref{prop:symconv}(2), needs no hypothesis on
$\ell$, and is the natural first target of Part~III. Parts (b) and (c) of
Theorem~\ref{thm:welldef} and all of Theorem~\ref{thm:classical} depend on the comparison
theorem --- through Lemma~\ref{lem:tauexists}(2), never (1) --- and should be attempted
only after it. Obligation~\ref{obl:tau} is already a
decidable predicate on emitted finite data, so it needs no idealisation to formalise ---
which is the shape rule~\ref{A4} asks every obligation to have.

%----------------------------------------------------------------------
\section{Induced local systems and the maps $\alpha,\mu,\omega$}\label{sec:induced}
%----------------------------------------------------------------------
% SOURCE:  hecke_action.tex:787-1063 (RH)
% NOTATION: rows sigma, "t,j,s (running index)", "s_j, delta_j, n^{(j,sigma)}_t",
%           lift(sigma), d (coboundary), mu, cring, Rel, Mat_2(Z)
%           of sec:notation-crosswalk applied.  The coset/block index is j and
%           the chain-term index t, never i (C15); the lift is \lift{\sigma},
%           never i(sigma) (C16); tau is the twisted chain datum and nothing
%           else -- in particular the oriented simplices of the proof of
%           thm:xi are sigma (C4); the middle map is mu, since beta_j is a
%           right-coset representative (C8).
% ERRATA:  E9 (the corestriction convention), E19 (the no-go) (App. A)

\paragraph{What this chapter fixes.}
The same operator as \S\ref{sec:chainlevel}, re-organised as a computation on the
\emph{fixed} complex $(K,\Rel)$ --- which is what makes it executable, since
\code{Equi\-variant\-Cohomology\-OnK} runs on $(K,\Rel)$ and on nothing else. Three
things are fixed here: the \emph{induced} matrix local system $L_H$, which is the
Shapiro-theoretic home of the coset bookkeeping; the two honest morphisms of local
systems $\alpha$ (restriction unit) and $\omega$ (corestriction counit); and the
\textbf{binding} fused cochain map
\[
  \Th_g\;=\;\omega_{\bullet}\circ\Xig ,
\]
with its closed form \eqref{eq:xiexplicit} and the assembly of the cohomology-level
matrix $\Amat{q}$ (Algorithm~\ref{alg:assembly}). Between the second and the third sits
the one place where the superseded framework document was not merely imprecise but
wrong: there is \emph{no} middle morphism of local systems, by a Schur argument, and
Remark~\ref{rmk:nogo} states the no-go rather than repairing the formula. The
recovered three-map factorisation --- which exists, at cochain level, and is a
cross-check rather than a definition --- is \S\ref{sec:equivalence}.

\paragraph{Standing hypotheses.}
\ref{H1}--\ref{H4} of \S\ref{sec:conventions}, each repeated where it is used
(rule~\ref{A1}). \emph{No construction} of
\S\S\ref{sec:inducedmodule}--\ref{sec:xitheta} needs a hypothesis on $\ell$: the induced
module, its block matrices, $\alpha$, $\omega$, $\Xig$ and $\Th_g$ are built over an
arbitrary commutative coefficient ring $\cring$ and are integral; so is every
well-definedness statement about them. The hypothesis $6\in\cring^{\times}$ (\ref{H3}) enters at
exactly two points, both flagged in their own statements: the Shapiro identification of
Proposition~\ref{prop:lhcompliant}(4), and the identification of the induced map with
$\Tdc{g}$ in Theorem~\ref{thm:xi}(3) --- which inherits it from
Theorem~\ref{thm:classical} and not from anything in this chapter.

\paragraph{Notation used throughout.}
The coset index and the induced-module block index are both $j\in\{1,\dots,m\}$ and the
chain-term index is $t$ (never $i$: \S\ref{sec:notation-collisions}, C15); the
distinguished lift of $\sigma\in\orient{q}(K)$ into $\Ku$ is $\lift{\sigma}$, which the
implementation spells \code{i(sigma)} (C16). \textbf{A relation of $\Rel$ is written
$(\sigma_{\mathrm{lhs}},\sigma_{\mathrm{rhs}},\gamma)$ with $\gamma\in\Ga$}, matching
Definition~\ref{def:model7rel}: in this chapter $g=\smat{1}{0}{0}{N}$ is the Hecke
element of \S\ref{sec:correspondencedatum}, and the two symbols appear in one formula
(Theorem~\ref{thm:xi}(1)) --- \S\ref{sec:notation-collisions}, C17. Every relation
element of the model does lie in $\Ga$ (Definition~\ref{def:model7rel}: $T$, $S$,
$ST$, $TS$, $-I$).

\paragraph{Parametric in the triangulation.}
Every statement of this chapter is stated for an arbitrary finite $(K,\Rel)$ satisfying
\ref{H4}, never for the seven vertices of Definition~\ref{def:model7}: the input is a
simplex set, a relation set whose elements are group elements, and the coset datum of
\S\ref{sec:heckedatum}, which is a statement about $\mathbb{P}^1(\Z/N)$ and is
independent of the triangulation altogether. A refinement of the model
(Remark~\ref{rmk:orbifold}, owner direction OD-T1) is therefore admissible input here
with no change to any definition, proof or algorithm; what a refinement changes is the
numbers $N_q$ and the $\tau$-table, not this chapter.

\subsection{The induced module and its block matrices}\label{sec:inducedmodule}

\begin{definition}[Induced module, ordered basis]\label{def:inducedmodule}
Let $\cring$ be a commutative ring, $H\le\Ga$ a subgroup of finite index $m$, and fix
representatives $w_1,\dots,w_m\in\Ga$ of $H\backslash\Ga$ with $w_1=I$. The
\emph{induced module} is
\[
  \Ind_H^{\Ga}\VR
  :=\bigl\{f\colon\Ga\to\VR\ :\ f(h\gamma)=\rhon(h)f(\gamma)\ \ \forall h\in H,\
  \gamma\in\Ga\bigr\},
  \qquad (\gamma\cdot f)(x):=f(x\gamma),
\]
a left $\cring[\Ga]$-module. Restriction to the representatives,
\[
  f\;\longleftrightarrow\;\bigl(f(w_1),\dots,f(w_m)\bigr)\in\VR^{\oplus m},
\]
is an isomorphism of $\cring$-modules, and the \emph{ordered basis} of
$\Ind_H^{\Ga}\VR$ is the concatenation, in the order $w_1,\dots,w_m$, of $m$ copies of
the ordered basis $(u_0,\dots,u_n)$ of $\VR$ --- so the coordinate of block $j$ and
in-block index $k$ is $(j-1)(n+1)+k$, and
$\dim\Ind_H^{\Ga}\VR=m(n+1)=[\Ga:H]\cdot d_V$. \LEAN{def:levelellmodel-model} (the module in the coordinates $(f(w_j))_j$, for any
table of representatives)
\end{definition}

\noindent
The restriction map is bijective because every $\gamma\in\Ga$ is uniquely $hw_j$ with
$h\in H$, so prescribing $f(w_j)$ prescribes $f(hw_j)=\rhon(h)f(w_j)$ and no
consistency condition remains; $\rhon(h)$ is invertible for $h\in\Ga$ by
Proposition~\ref{prop:symconv}(2).

\begin{proposition}[Block transition matrices]\label{prop:blockmatrices}
\PROVED{} Let $\cring$, $H$, $m$, $(w_j)$ be as in
Definition~\ref{def:inducedmodule}. For $\gamma\in\Ga$ define the index $j\cdot\gamma$
and the element $h_j(\gamma)\in H$ by
\[
  w_j\,\gamma\;=\;h_j(\gamma)\,w_{j\cdot\gamma},\qquad j=1,\dots,m .
\]
Both are well defined; $(j,\gamma)\mapsto j\cdot\gamma$ is a right action of $\Ga$ on
$\{1,\dots,m\}$, and
\[
  h_j(I)=I,\qquad h_j(\gamma_1\gamma_2)=h_j(\gamma_1)\,h_{j\cdot\gamma_1}(\gamma_2).
\]
In the coordinates of Definition~\ref{def:inducedmodule} the action of $\gamma$ on
$\Ind_H^{\Ga}\VR$ is the block matrix
\begin{equation}\label{eq:MH}
  M_H(\gamma)\;=\;\sum_{j=1}^{m}\ E_{j,\;j\cdot\gamma}\otimes\rhon\bigl(h_j(\gamma)\bigr)
  \qquad\text{(block row $j$, block column $j\cdot\gamma$)},
\end{equation}
i.e.\ $(\gamma\cdot f)(w_j)=\rhon(h_j(\gamma))\,f(w_{j\cdot\gamma})$; it has exactly one
non-zero block per block row. The map $\gamma\mapsto M_H(\gamma)$ is a group
homomorphism $\Ga\to\GL_{m(n+1)}(\cring)$, and $M_H(\gamma^{-1})=M_H(\gamma)^{-1}$. \LEAN{def:levelellmodel-model} (\eqref{eq:MH}
with $h_j(\gamma)=w_j\gamma w_{j\cdot\gamma}^{-1}$, the two cocycle identities and the homomorphism, for any table; ``one non-zero block per
block row'' in the form of a product of maps)
\end{proposition}

\begin{proof}
Well-definedness and the cocycle identity are the argument of
Lemma~\ref{lem:cocycle} with $(\Ga'',s_j)$ replaced by $(H,w_j)$ and the class map
replaced by the coset map $\gamma\mapsto H\gamma$: $Hw_j\gamma$ is one of the $m$ cosets,
which names $j\cdot\gamma$ uniquely, and $h_j(\gamma)=w_j\gamma w_{j\cdot\gamma}^{-1}\in
H$ by construction; expanding $w_j\gamma_1\gamma_2$ in two ways and using uniqueness of
the decomposition $\gamma=hw_j$ gives the action property and the cocycle identity.
For \eqref{eq:MH},
\[
  (\gamma\cdot f)(w_j)=f(w_j\gamma)=f\bigl(h_j(\gamma)w_{j\cdot\gamma}\bigr)
  =\rhon\bigl(h_j(\gamma)\bigr)f(w_{j\cdot\gamma}),
\]
which is block row $j$ of the displayed matrix. Homomorphy is then a computation in
which the only input is multiplicativity of $\rhon$
(Proposition~\ref{prop:symconv}(2)): with $E_{j,a}E_{j',b}=\delta_{a j'}E_{j,b}$,
\[
  M_H(\gamma_1)M_H(\gamma_2)
  =\sum_{j}E_{j,\,(j\cdot\gamma_1)\cdot\gamma_2}\otimes
    \rhon\bigl(h_j(\gamma_1)\bigr)\rhon\bigl(h_{j\cdot\gamma_1}(\gamma_2)\bigr)
  =\sum_{j}E_{j,\,j\cdot(\gamma_1\gamma_2)}\otimes
    \rhon\bigl(h_j(\gamma_1\gamma_2)\bigr),
\]
which is $M_H(\gamma_1\gamma_2)$; and $M_H(I)=I$. Invertibility follows with
$\gamma_2=\gamma_1^{-1}$.
\end{proof}

\begin{contractbox}
For the Hecke operator at $p^{\nu}$ the implementation instantiates
Proposition~\ref{prop:blockmatrices} twice, and \textbf{both instances are the coset
datum of \S\ref{sec:heckedatum}, not a second construction}:
\begin{itemize}
\item $L'':=L_{\Ga''}$, $\Ga''=\Ga_0(N)$, with $w_j:=s_j$ the canonical transporters of
  Definition~\ref{def:cosetreps}. Then $j\cdot\gamma$ is the bottom-row
  $\mathbb{P}^1(\Z/N)$ index action and $h_j(\gamma)=s_j\gamma s_{j\cdot\gamma}^{-1}$ is
  the cocycle \eqref{eq:cocycle}, both by Lemma~\ref{lem:cocycle}. Only the four
  generator values $\gamma\in\{T,S,ST,TS\}$ and $\gamma=-I$ are ever needed
  (Definition~\ref{def:model7rel}).
\item $L':=L_{\Ga'}$, $\Ga'=\Ga^{0}(N)$, with $w_j:=t_j$ the upper-triangular
  transporters of Definition~\ref{def:cosetreps} and the \emph{top}-row class map
  $\smat{a}{b}{c}{d}\mapsto(a:b)$. This one is used only by $\alpha$
  (Proposition~\ref{prop:alpha}) and by the cross-check of \S\ref{sec:equivalence};
  the binding pipeline never constructs it.
\end{itemize}
\end{contractbox}

\begin{definition}[Induced matrix local system]\label{def:inducedls}
Let $L=(\VR,\{\rhon(\gamma)\}_{(\cdot,\cdot,\gamma)\in\Rel})$ be the matrix local system
of the model (Definition~\ref{def:matrix_ls}) and let $H\le\Ga$ be of finite index $m$
with representatives as in Definition~\ref{def:inducedmodule}. Then
\[
  L_H\;:=\;\Bigl(\Ind_H^{\Ga}\VR,\ \bigl\{M_H(\gamma)\bigr\}_{(\cdot,\cdot,\gamma)\in\Rel}\Bigr)
\]
is the matrix local system on $(K,\Rel)$ obtained from $L$ by replacing the matrix
$\rhon(\gamma)$ of each relation by $M_H(\gamma)$ of \eqref{eq:MH}, \textbf{keeping the
same relation list, the same simplex pairs
$(\sigma_{\mathrm{lhs}},\sigma_{\mathrm{rhs}})$ and the same order}. Its representation
dimension is $m(n+1)$.
\end{definition}

\begin{proposition}[$L_H$ is a matrix local system; Shapiro]\label{prop:lhcompliant}
\PROVED{} Assume \ref{H4} and let $H\le\Ga$ be of finite index $m$, with
$L_H$ as in Definition~\ref{def:inducedls} over a commutative ring $\cring$. Then:
\begin{enumerate}
\item every $M_H(\gamma)$ is invertible over $\cring$ and the assignment
  $\gamma\mapsto M_H(\gamma)$ is a group homomorphism, so $L_H$ satisfies the cocycle
  condition of Definition~\ref{def:matrix_ls} and is a matrix local system on
  $(K,\Rel)$ --- with \emph{no} hypothesis on $\ell$;
\item consequently the data of Definition~\ref{def:inducedls} is a valid
  \code{LocalSystemSimplicialComplex} input (each relation is again of the single-matrix
  form $f(\sigma_{\mathrm{lhs}})=M\,f(\sigma_{\mathrm{rhs}})$ over $\cring$), and
  \code{EquivariantCohomologyOnK} computes $\Heq{\bullet}(K,\Rel;L_H)$ \emph{unchanged};
\item if $-I\in H$ --- which holds for $H\in\{\Ga_0(N),\Ga^{0}(N)\}$ --- then the
  central relations $(\sigma,\sigma,-I)$ of Definition~\ref{def:model7rel} are carried by
  \[
    M_H(-I)\;=\;I_m\otimes\rhon(-I)\;=\;(-1)^{n}\,I_{m(n+1)} ,
  \]
  so they impose $f\equiv0$ for odd $n$ on every block exactly as they do for $L$
  (Remark~\ref{rmk:negI});
\item if moreover $6\in\cring^{\times}$ --- for $\cring$ of characteristic $\ell$,
  $\ell\ge5$ (hypothesis \ref{H3}) --- then
  \[
    \Heq{q}(K,\Rel;L_H)\;\cong\;H^{q}\bigl(\Ga;\Ind_H^{\Ga}\VR\bigr)
    \;\cong\;H^{q}(H;\VR)\qquad\text{for every }q ,
  \]
  the first isomorphism by Theorem~\ref{thm:explicit_complex_equiv} and
  Remark~\ref{rmk:orbifold} --- comparison of the projective $\cring[H]$-resolution $C_{\bullet}(\Xbulk)\otimes\cring$ with any
  projective resolution (v1.16.0-pre.3) ---, the second by Shapiro's lemma
  \STD{\cite[Ch.~III.6.2 and Exercise~III.8.2]{brown}}. For $6\notin\cring^{\times}$
  \textbf{neither isomorphism is asserted} and the left-hand side is the orbifold model
  cohomology of \ref{H3}. \LEAN{def:levelellmodel-model} (since v1.20.0: both isomorphisms for $6\in\cring^{\times}$, on the model, for any
  $H\le\Ga$ of finite index with a table of representatives, $\Ga_0(N)$ included; not (1)--(3))
\end{enumerate}
\end{proposition}

\begin{proof}
(1) is Proposition~\ref{prop:blockmatrices}. The cocycle condition of
Definition~\ref{def:matrix_ls} asks that every relation among the relation elements that
holds in the group presented by $(K,\Rel)$ --- which is $\Ga$, by
Theorem~\ref{thm:fund_group_generators} and \ref{H4}(iii) --- be satisfied by the
assigned matrices; a group homomorphism satisfies every such relation, so there is
nothing to check relation by relation. This is the induced-system half of
Proposition~\ref{prop:matrix_rep_equiv}.
(2) is (1) plus the observation that \eqref{eq:MH} is a single square matrix of size
$m(n+1)$, which is the format \code{SimplexRelation} accepts; the relation list is
unchanged by construction.
(3) $-I$ is central in $\Mat_2(\Z)$ and lies in $H$, so $w_j(-I)=(-I)w_j$ gives
$j\cdot(-I)=j$ and $h_j(-I)=w_j(-I)w_j^{-1}=-I$; \eqref{eq:MH} then reads
$M_H(-I)=I_m\otimes\rhon(-I)$, and $\rhon(-I)u_k=(-x)^{n-k}(-y)^{k}=(-1)^{n}u_k$ by
Definition~\ref{def:rhon}.
(4) The first isomorphism is Theorem~\ref{thm:explicit_complex_equiv} applied to the
local system $L_H$ --- whose hypotheses are \ref{H4}, not a property of the coefficient
module --- followed by Remark~\ref{rmk:orbifold}, whose \emph{only} use of
$6\in\cring^{\times}$ is the invertibility of the simplex-stabiliser orders, again
independent of the coefficient module. The second is Shapiro's lemma for the finite-index
subgroup $H$.
\end{proof}

\begin{remark}[What Proposition~\ref{prop:lhcompliant}(4) is for, and what it is not]
\label{rmk:shapirouse}
Part (4) is an \emph{oracle}, not a step of the pipeline: it predicts
$\dim\Heq{q}(K,\Rel;L_H)$ from classical data about the modular curve $X_0(N)$ and is
checked against the computed dimension (Obligation~\ref{obl:shapiro}). The binding
construction of $\Th_g$ (Definition~\ref{def:xi}) never builds $L'$ or $L''$ and never
computes the cohomology of either. This is why the whole of \S\ref{sec:inducedmodule}
is, for the pipeline, a \emph{coordinate system} rather than a computation.
\end{remark}

\subsection{\texorpdfstring{$\alpha$}{alpha} and \texorpdfstring{$\omega$}{omega}: two honest morphisms of local systems}
\label{sec:alphaomega}

\begin{proposition}[Restriction unit $\alpha$]\label{prop:alpha}
\PROVED{} Assume \ref{H4}, let $\cring$ be a commutative ring, $N=p^{\nu}$,
$\Ga'=\Ga^{0}(N)$ of index $m$ with the upper-triangular representatives
$t_1=I,\dots,t_m$ of Definition~\ref{def:cosetreps}. Then
\[
  \alpha\colon \VR\longrightarrow\Ind_{\Ga'}^{\Ga}\VR,\qquad
  \alpha(v)(\gamma):=\rhon(\gamma)\,v,
\]
is a morphism of $\cring[\Ga]$-modules; in the coordinates of
Definition~\ref{def:inducedmodule} it is the \emph{block column}
$\bigl(\rhon(t_j)\bigr)_{j=1,\dots,m}$. Equivalently
(Theorem~\ref{thm:rep_ls_equiv}) it is a morphism of matrix local systems
$L\to L'$, i.e.\ $\alpha\,\rhon(\gamma)=M_{\Ga'}(\gamma)\,\alpha$ for every relation
element $\gamma$ of $\Rel$ --- with \emph{no} hypothesis on $\ell$. If moreover
$6\in\cring^{\times}$, then under the identifications of
Proposition~\ref{prop:lhcompliant}(4) the induced map on cohomology is
$\res^{\Ga}_{\Ga'}$ \STD{\cite[Ch.~III.6.2 and Exercise~III.8.2]{brown}}.
\end{proposition}

\begin{proof}
$\alpha(v)$ lies in the induced module: $\alpha(v)(h\gamma)=\rhon(h\gamma)v=
\rhon(h)\alpha(v)(\gamma)$ by Proposition~\ref{prop:symconv}(2). Equivariance:
$\alpha(\rhon(\gamma)v)(x)=\rhon(x\gamma)v=\alpha(v)(x\gamma)=(\gamma\cdot\alpha(v))(x)$,
again by Proposition~\ref{prop:symconv}(2). Coordinates: $\alpha(v)(t_j)=\rhon(t_j)v$,
which is the stated block column. The local-system reformulation is the equivalence
$\Rep\simeq\LS$ of Theorem~\ref{thm:rep_ls_equiv}, applicable because every relation
element of $\Rel$ lies in $\Ga$. The Shapiro statement identifies the unit of the
restriction--induction adjunction with $\res$.
\end{proof}

\begin{proposition}[Corestriction counit $\omega$]\label{prop:omega}
\PROVED{} Assume \ref{H4}, let $\cring$ be a commutative ring, $N=p^{\nu}$,
$\Ga''=\Ga_0(N)$ of index $m$, and let $s_1,\dots,s_m\in\Ga$ represent
$\Ga''\backslash\Ga$ (canonically Definition~\ref{def:cosetreps}). Then the trace map
\[
  \omega\colon \Ind_{\Ga''}^{\Ga}\VR\longrightarrow\VR,\qquad
  \omega(f):=\sum_{j=1}^{m}\rhon(s_j)^{-1}f(s_j)=\sum_{j=1}^{m}\rhon(s_j^{-1})f(s_j),
\]
is independent of the choice of the $s_j$ and is a morphism of $\cring[\Ga]$-modules;
in the coordinates of Definition~\ref{def:inducedmodule} it is the \emph{block row}
$\bigl(\rhon(s_j^{-1})\bigr)_{j=1,\dots,m}$. Equivalently
(Theorem~\ref{thm:rep_ls_equiv}) it is a morphism of matrix local systems $L''\to L$ ---
with \emph{no} hypothesis on $\ell$. Every $s_j$ lies in $\Ga$, so $\rhon(s_j^{-1})$ is
integral and invertible over every $\cring$ and \textbf{no determinant-$N$ matrix is
inverted} (Obligation~\ref{obl:integral}). If moreover $6\in\cring^{\times}$, the
induced map on cohomology is the transfer $\cor^{\Ga}_{\Ga''}$
\STD{\cite[Ch.~III.9--10]{brown}}.
\end{proposition}

\begin{proof}
Independence: replacing $s_j$ by $h s_j$ with $h\in\Ga''$ replaces $f(s_j)$ by
$\rhon(h)f(s_j)$ and $\rhon(s_j^{-1})$ by $\rhon(s_j^{-1})\rhon(h)^{-1}$, by
Proposition~\ref{prop:symconv}(2), so the summand is unchanged. Equivariance: with
$s_j\gamma=h_j(\gamma)s_{j\cdot\gamma}$ and hence
$s_j^{-1}h_j(\gamma)=\gamma s_{j\cdot\gamma}^{-1}$ (Lemma~\ref{lem:cocycle}),
\[
  \omega(\gamma\cdot f)=\sum_{j}\rhon(s_j^{-1})f(s_j\gamma)
  =\sum_{j}\rhon(s_j^{-1})\rhon\bigl(h_j(\gamma)\bigr)f(s_{j\cdot\gamma})
  =\rhon(\gamma)\sum_{j}\rhon(s_{j\cdot\gamma}^{-1})f(s_{j\cdot\gamma})
  =\rhon(\gamma)\,\omega(f),
\]
using Proposition~\ref{prop:symconv}(2) twice and that $j\mapsto j\cdot\gamma$ is a
bijection. The coordinate description and the local-system reformulation are as in
Proposition~\ref{prop:alpha}. The identification with $\cor$ is the transfer formula on
the resolution $C_{\bullet}(\Ku)\otimes\cring$, which is the computation already made in
the proof of Theorem~\ref{thm:classical}.
\end{proof}

\begin{remark}[The corestriction convention, stated once]\label{rmk:omegaconvention}
The superseded framework document writes the corestriction as
$\omega_{\ast}(\psi)=\sum_{\gamma}\rho_L(\gamma)\,\psi(\gamma^{-1})$, the sum over
representatives $\gamma$ of $\Ga''\backslash\Ga$
(\file{monodromy\_v5.tex}:1536). That formula and the one in
Proposition~\ref{prop:omega} agree, but \emph{only} after the substitution
$\gamma=s_j^{-1}$ --- which also replaces the index set $\{s_j\}$ of
$\Ga''\backslash\Ga$ by the set $\{s_j^{-1}\}$ of representatives of $\Ga/\Ga''$. The
substitution is recorded here because it is unmarked at the source, where a transcriber
meets it first (Appendix~\ref{app:errata}, entry~\ref{err:e9}). This document uses the
form of Proposition~\ref{prop:omega} and no other.
\end{remark}

\subsection{The middle map: a no-go, not a repair}\label{sec:nogo}

\begin{remark}[The middle map $\mu$ is not a morphism of local systems]\label{rmk:nogo}
\PROVED{} The superseded framework document \file{monodromy\_v5.tex} \S12.5 proposes a
``conjugation'' map $\mu_{\ast}(\phi)(\gamma):=\rho_L(g)\,\phi(g^{-1}\gamma g)$ --- it
calls it $\beta$; here $\beta_j$ is a right-coset representative
(\S\ref{sec:notation-collisions}, C8) --- as a morphism
$\Ind_{\Ga'}^{\Ga}\VR\to\Ind_{\Ga''}^{\Ga}\VR$, and builds the Hecke operator as
$\omega\circ\mu\circ\alpha$ (its \S\S12.6--12.7). This is \textbf{not correct as
stated}, for two independent reasons, the second of which forecloses every repair.

\emph{(i) Ill-typed.} An element of $\Ind_{\Ga'}^{\Ga}\VR$ is a function on all of
$\Ga$, but $g^{-1}\gamma g\in\Ga$ only for $\gamma\in\Ga''$
(Proposition~\ref{prop:cosets}(3)); the formula does not define a function on $\Ga$, and
no $\Ga$-equivariant extension of it is specified.

\emph{(ii) No morphism of local systems can do the job.} Let $\cring=\Q$ and
$\Vcoef=\Sym^{10}\Q^2$. \emph{Any} composite of three morphisms of
$\Q[\Ga]$-modules
\[
  \Vcoef\;\xrightarrow{\ \alpha\ }\;\Ind_{\Ga'}^{\Ga}\Vcoef
  \;\xrightarrow{\ \mu\ }\;\Ind_{\Ga''}^{\Ga}\Vcoef
  \;\xrightarrow{\ \omega\ }\;\Vcoef
\]
is an element $u\in\End_{\Ga}(\Vcoef)$, and the map it induces on $H^{\bullet}(\Ga;\Vcoef)$
is $u_{\ast}$. Now $\Sym^{10}\Q^2$ is an irreducible $\Ga$-module, because $\Ga$ is
Zariski dense in $\SL_2$ \STD{\cite[Ch.~VII, \S1]{serre_cours} for the density of
$\SL_2(\Z)$, and the irreducibility of $\Sym^{n}$ for $\SL_2$ in characteristic $0$},
so $\End_{\Ga}(\Vcoef)=\Q$ and $u_{\ast}$ is a \emph{scalar} in every degree. But $T_2$
on $H^1(\SL_2(\Z);\Sym^{10}\Q^2)$ has eigenvalues $-24,-24,2^{11}+1=2049$
\STD{Eichler--Shimura in weight $k=n+2=12$, \cite[Ch.~8]{shimura};
Obligation~\ref{obl:oracle}}, which is not scalar. Hence \textbf{no} morphism of the two
induced local systems whatsoever --- in particular no repair of the displayed formula ---
realises $T_p$ through $\omega\circ\mu\circ\alpha$.

\emph{Resolution adopted by this document.} The pipeline does not factor the operator
through a third local system at all: it is the fused cochain map
$\Th_g=\omega_{\bullet}\circ\Xig$ of Definition~\ref{def:xi}, which needs neither $\Ga'$
nor $L'$ nor any $\mu$. A middle map \emph{does} exist one level down, as a cochain map
given by per-simplex block rows rather than by one fibre matrix; it is constructed in
Proposition~\ref{prop:mu} and used only as a cross-check. The distinction is exactly the
one (ii) makes: a morphism of local systems is a single fibre matrix applied at every
simplex, and that is what cannot exist. Recorded as
Appendix~\ref{app:errata}, entry~\ref{err:e19}.
\end{remark}

\subsection{The binding cochain map \texorpdfstring{$\Xig$}{Xi\_g} and \texorpdfstring{$\Th_g$}{Theta\_g}}\label{sec:xitheta}

Throughout this subsection $g=\smat{1}{0}{0}{N}$, $N=p^{\nu}$, $\Ga''=\Ga_0(N)$ of index
$m$, $s_1,\dots,s_m$ and $\delta_j=s_j^{-1}g$ are as in
Definition~\ref{def:cosetreps}, and $\tau$ is a $g$-twisted chain datum over $\cring$
(Definition~\ref{def:tau}).

Assume \ref{H4}. For $c\in\Cq{q}(K,\Rel;L)$ write $\phi_c$ for the corresponding
$\Ga$-equivariant cochain on $\Ku$ under the isomorphism $\Psi$ of
Theorem~\ref{thm:explicit_complex_equiv}:
\begin{equation}\label{eq:psiinv}
  \phi_c\bigl(\gamma\cdot\lift{\sigma}\bigr)\;=\;\rhon(\gamma)\,c(\sigma),
  \qquad \gamma\in\Ga,\ \sigma\in\orient{q}(K),
\end{equation}
extended $\cring$-linearly to $C_q(\Ku)\otimes\cring$. \textbf{Two hypotheses make
\eqref{eq:psiinv} a definition rather than a wish, and both are used below.} First,
$K\subseteq\Ku$ is a \emph{strict} fundamental domain (\ref{H4}(iii)), so every simplex
of $\Ku$ is $\gamma\cdot\lift{\sigma}$ for a unique $\sigma$; second, if
$\gamma_1\lift{\sigma}=\gamma_2\lift{\sigma}$ then $\gamma_2^{-1}\gamma_1$ lies in the
stabiliser $\Ga_{\lift{\sigma}}$, and the right-hand side is unambiguous \emph{because}
$c$ satisfies the stabiliser relations and the central relation of
Definition~\ref{def:model7rel} --- which is precisely what ``equivariant'' means for
$(K,\Rel)$ and is the content of Theorem~\ref{thm:explicit_complex_equiv}. Consistency
with the coboundary, $\phi_{dc}=d\,\phi_c$ where $(d\phi)(x)=\phi(\partial x)$, is part
of the same theorem.

\begin{definition}[The map $\Xig$ and the fused map $\Th_g$]\label{def:xi}
Assume \ref{H4}. For $c\in\Cq{q}(K,\Rel;L)$ define
$\Xig(c)\in C^{q}\bigl(K;\Ind_{\Ga''}^{\Ga}\VR\bigr)$ blockwise by
\begin{equation}\label{eq:xidef}
  \bigl(\Xig c\bigr)(\sigma)_j\;:=\;\rhon(g)\ \phi_c\Bigl(\tau\bigl(s_j\cdot\lift{\sigma}\bigr)\Bigr)
  \ \in\ \VR,
  \qquad j=1,\dots,m,\ \ \sigma\in\orient{q}(K).
\end{equation}
Explicitly, and this is the form the implementation evaluates: the table of
Algorithm~\ref{alg:tau} stores, for each pair $(j,\sigma)$, the finite chain
\[
  \tau\bigl(s_j\cdot\lift{\sigma}\bigr)
  \;=\;\sum_{t} n^{(j,\sigma)}_{t}\ \gamma^{(j,\sigma)}_{t}\cdot\lift{\sigma^{(j,\sigma)}_{t}},
  \qquad n_t\in\Z,\ \gamma_t\in\Ga,\ \sigma_t\in\orient{q}(K),
\]
whence, by \eqref{eq:psiinv} and Proposition~\ref{prop:symconv}(2),
\begin{equation}\label{eq:xiexplicit}
  \bigl(\Xig c\bigr)(\sigma)_j
  \;=\;\sum_{t} n^{(j,\sigma)}_{t}\ \rhon\bigl(g\,\gamma^{(j,\sigma)}_{t}\bigr)\
  c\bigl(\sigma^{(j,\sigma)}_{t}\bigr).
\end{equation}
Every matrix occurring in \eqref{eq:xiexplicit} is $\rhon$ of an \emph{integral} matrix
of determinant $N$ (Definition~\ref{def:semigroup}), precomputed once over $\cring$; no
determinant-$N$ matrix is inverted (Obligation~\ref{obl:integral}). Finally
\[
  \Th_g\;:=\;\omega_{\bullet}\circ\Xig,\qquad
  \bigl(\Th_g c\bigr)(\sigma)\;=\;\sum_{j=1}^{m}\rhon(s_j^{-1})\,\bigl(\Xig c\bigr)(\sigma)_j ,
\]
with $\omega_{\bullet}$ the simplexwise application of the block row of
Proposition~\ref{prop:omega}.
\end{definition}

\noindent
Definition~\ref{def:xi} is constructive in the sense of rule~\ref{A3}:
\eqref{eq:xiexplicit} is a finite sum over a finite table, and it is literally what
\code{HeckeTauTable::xi\_terms(j,\,q,\,sigma)} stores --- one triple
$\bigl(n_t,\ \sigma_t,\ \rhon(g\gamma_t)\bigr)$ per term, with the $\rhon(g\gamma_t)$
shared through a content-addressed pool.

\begin{theorem}[$\Xig$ lands in the induced system; $\Th_g$ realises $\Tdc{g}$]
\label{thm:xi}
Assume \ref{H4}, let $\cring$ be a commutative ring, and let
$g,\{s_j\},\tau$ be as above. Then:
\begin{enumerate}
\item \PROVED{} \textbf{(No hypothesis on $\ell$.)}
  $\Xig(c)\in\Cq{q}(K,\Rel;L'')$ for every $c\in\Cq{q}(K,\Rel;L)$: for every relation
  $(\sigma_{\mathrm{lhs}},\sigma_{\mathrm{rhs}},\gamma)\in\Rel$ --- the stabiliser and
  central relations included ---
  \[
    (\Xig c)(\sigma_{\mathrm{lhs}})\;=\;M_{\Ga''}(\gamma)\,(\Xig c)(\sigma_{\mathrm{rhs}}).
  \]
\item \PROVED{} \textbf{(No hypothesis on $\ell$.)} $\Xig\circ d=d\circ\Xig$ and
  $\omega_{\bullet}\circ d=d\circ\omega_{\bullet}$; hence
  $\Th_g\colon\Cq{\bullet}(K,\Rel;L)\to\Cq{\bullet}(K,\Rel;L)$ is a map of cochain
  complexes.
\item \PROVED{} $\Th_g=\Psi\circ\widetilde T\circ\Psi^{-1}$ with $\widetilde T$ the
  operator \eqref{eq:heckechain} of Definition~\ref{def:hecke}, as an identity of
  cochain maps and with \emph{no} hypothesis on $\ell$. Consequently the map induced on
  $\Heq{q}(K,\Rel;L)$ is the operator of Definition~\ref{def:hecke} in every degree $q$;
  and if moreover $6\in\cring^{\times}$ (hypothesis \ref{H3}) it is the classical
  double-coset operator $\Tdc{g}$ of \eqref{eq:corconjres}, by
  Theorem~\ref{thm:classical}.
\end{enumerate}
\end{theorem}

\begin{proof}
(1) Let $(\sigma_{\mathrm{lhs}},\sigma_{\mathrm{rhs}},\gamma)\in\Rel$, so that
$\lift{\sigma_{\mathrm{lhs}}}=\gamma\cdot\lift{\sigma_{\mathrm{rhs}}}$ in $\Ku$
(\ref{H4}(iii) and Theorem~\ref{thm:explicit_complex_equiv}; for a stabiliser or central
relation $\sigma_{\mathrm{lhs}}=\sigma_{\mathrm{rhs}}$ and $\gamma$ fixes the lift).
Write $s_j\gamma=h_j(\gamma)s_{j\cdot\gamma}$ with $h_j(\gamma)\in\Ga''$
(Lemma~\ref{lem:cocycle}). Then, using (T1) of Definition~\ref{def:tau} --- legitimate
because $h_j(\gamma)\in\Ga''$ --- and $\Ga$-equivariance of $\phi_c$,
\begin{align*}
  (\Xig c)(\sigma_{\mathrm{lhs}})_j
  &=\rhon(g)\,\phi_c\bigl(\tau(s_j\gamma\cdot\lift{\sigma_{\mathrm{rhs}}})\bigr)
   =\rhon(g)\,\phi_c\Bigl(\tau\bigl(h_j(\gamma)\,s_{j\cdot\gamma}\cdot
     \lift{\sigma_{\mathrm{rhs}}}\bigr)\Bigr)\\
  &=\rhon(g)\,\phi_c\Bigl(\bigl(g^{-1}h_j(\gamma)g\bigr)\cdot
     \tau\bigl(s_{j\cdot\gamma}\cdot\lift{\sigma_{\mathrm{rhs}}}\bigr)\Bigr)
   =\rhon(g)\,\rhon\bigl(g^{-1}h_j(\gamma)g\bigr)\,
     \phi_c\Bigl(\tau\bigl(s_{j\cdot\gamma}\cdot\lift{\sigma_{\mathrm{rhs}}}\bigr)\Bigr)\\
  &=\rhon\bigl(h_j(\gamma)\bigr)\,(\Xig c)(\sigma_{\mathrm{rhs}})_{j\cdot\gamma},
\end{align*}
the last step by Proposition~\ref{prop:symconv}(2) applied to the identity
$g\,(g^{-1}h_j(\gamma)g)=h_j(\gamma)\,g$ in $\Mat_2(\Z)$ --- an identity between
integral matrices, with no inverse of a determinant-$N$ matrix taken. By \eqref{eq:MH}
the right-hand side is block $j$ of
$M_{\Ga''}(\gamma)\,(\Xig c)(\sigma_{\mathrm{rhs}})$.

(2) For $\Xig$: the coboundary of the equivariant complex is
$(dc)(\sigma)=c(\partial\sigma)$, and $K\subseteq\Ku$ is a simplicial subcomplex
(\ref{H4}(iii)), so $\lift{\partial\sigma}=\partial\lift{\sigma}$. Hence, for
$\sigma\in\orient{q+1}(K)$,
\[
  \bigl(d\,\Xig c\bigr)(\sigma)_j
  =(\Xig c)(\partial\sigma)_j
  =\rhon(g)\,\phi_c\bigl(\tau(s_j\cdot\partial\lift{\sigma})\bigr)
  =\rhon(g)\,\phi_c\bigl(\partial\,\tau(s_j\cdot\lift{\sigma})\bigr)
  =\rhon(g)\,(d\phi_c)\bigl(\tau(s_j\cdot\lift{\sigma})\bigr),
\]
using that $s_j$ acts simplicially and (T2) of Definition~\ref{def:tau}; since
$\phi_{dc}=d\phi_c$ (Theorem~\ref{thm:explicit_complex_equiv}), the right-hand side is
$(\Xig\,dc)(\sigma)_j$. For $\omega_{\bullet}$: it is one fixed fibre matrix applied at
every simplex, and the coboundary acts as $d_1\otimes\id$ on
$C^{\bullet}(K)\otimes\VR$, so the two commute. Composing gives the statement for
$\Th_g$, and (1) gives that $\Th_g$ maps $\Cq{q}(K,\Rel;L)$ into itself, since
$\omega_{\bullet}$ carries $L''$-equivariant cochains to $L$-equivariant ones
(Proposition~\ref{prop:omega}).

(3) For $\sigma\in\orient{q}(K)$,
\[
  (\Th_gc)(\sigma)
  =\sum_{j=1}^{m}\rhon(s_j^{-1})\,\rhon(g)\,\phi_c\bigl(\tau(s_j\cdot\lift{\sigma})\bigr)
  =\sum_{j=1}^{m}\rhon(\delta_j)\,\phi_c\bigl(\tau(s_j\cdot\lift{\sigma})\bigr)
  =(\widetilde T\phi_c)\bigl(\lift{\sigma}\bigr),
\]
the middle equality by Proposition~\ref{prop:symconv}(2) applied to
$s_j^{-1}g=\delta_j$ in $\Mat_2(\Z)$ (again integral, and $\rhon(s_j)^{-1}=\rhon(s_j^{-1})$
because $s_j\in\Ga$ is unimodular), and the last by \eqref{eq:heckechain}. Since
$\Psi$ is evaluation at the lifts \eqref{eq:psiinv}, this says
$\Th_gc=\Psi(\widetilde T\phi_c)=\Psi\widetilde T\Psi^{-1}c$. The remaining assertions
are Definition~\ref{def:hecke} and Theorem~\ref{thm:classical}, the latter carrying the
hypothesis $6\in\cring^{\times}$ and being the \emph{only} place it is needed.
\end{proof}

\begin{remark}[Which choices $\Th_g$ depends on]\label{rmk:thetachoices}
$\Th_g$ depends on $g$, on the representatives $s_j$ and on the chosen $\tau$, and it
is not invariant under changing any of them --- only the induced $\Amat{q}$ is, and only
under the hypotheses of Theorem~\ref{thm:welldef}. Remark~\ref{rmk:indeplevel} says
exactly what may therefore be compared bit for bit and what may not; nothing in this
chapter weakens it.
\end{remark}

\subsection{Assembly of the Hecke matrix on \texorpdfstring{$\Heq{q}$}{H\^{}q\_eq}}\label{sec:assembly}

\begin{algorithm}[Hecke matrix in degree $q$]\label{alg:assembly}
Assume \ref{H4}. \textbf{Input.} The computed equivariant cohomology of $L$
(\S\ref{sec:systematiccomputation}): the stored basis
$\bigl[z^{(q)}_1\bigr],\dots,\bigl[z^{(q)}_{d_q}\bigr]$ of $\Heq{q}(K,\Rel;L)$ with
representative cocycles $z^{(q)}_j\in\Zq{q}\subseteq\cring^{N_qd_V}$ in the flat
layout, and the projector
\[
  P_q\colon\cring^{N_qd_V}\longrightarrow\cring^{d_q},
  \qquad P_q\,z^{(q)}_j=e_j,\quad P_q\,\Bq{q}=0 ,
\]
i.e.\ $P_q$ reads off coordinates in the stored basis modulo $\Bq{q}$ (this is the
defining property of the stored projector, and it is certified as
$P_qH_{\mathrm{eq}}=I$, $P_qB_{\mathrm{eq}}=0$); plus the $\tau$-table of
\S\ref{sec:reduction} for $g=\smat{1}{0}{0}{N}$ and the coset tables of
Definition~\ref{def:cosetreps}.
\textbf{Steps.}
\begin{enumerate}
\item For $j=1,\dots,d_q$ compute $y_j:=\Th_g\,z^{(q)}_j\in\cring^{N_qd_V}$ by
  \eqref{eq:xiexplicit} followed by the $\omega$-row of
  Proposition~\ref{prop:omega}.
\item Set column $j$ of $\Amat{q}\in\Mat_{d_q}(\cring)$ to $P_q\,y_j$.
\end{enumerate}
\textbf{Output.} $\Amat{q}$, the matrix of the operator of
Definition~\ref{def:hecke} on $\Heq{q}(K,\Rel;L)$ in the stored basis, in the
\emph{column} convention
\[
  T\bigl[z^{(q)}_j\bigr]\;=\;\sum_{k}\Amat{q}_{kj}\,\bigl[z^{(q)}_k\bigr].
\]
\textbf{Correctness.} $y_j\in\Cq{q}(K,\Rel;L)$ by Theorem~\ref{thm:xi}(1) and
Proposition~\ref{prop:omega}; $d\,y_j=0$ by Theorem~\ref{thm:xi}(2) and
$d\,z^{(q)}_j=0$, so $y_j\in\Zq{q}$ and $P_qy_j$ is the coordinate vector of
$[y_j]=T[z^{(q)}_j]$ (Theorem~\ref{thm:xi}(3)) modulo $\Bq{q}$.
\textbf{Degenerate degrees.} $d_q=0$ yields the $0\times0$ matrix and
$\dim\Cq{q}=0$ the $0\times0$ $\Th^{(q)}_{\mathrm{eq}}$; nothing is special-cased. This
is the normal situation at $q=2$ for $6\in\cring^{\times}$ and at every $q$ for odd $n$
(Remark~\ref{rmk:negI}).
\textbf{No normalisation.} $\Amat{q}$ is neither symmetrised nor rescaled: the integral
normalisation of the contract box after Proposition~\ref{prop:symconv} is what fixes the
eigenvalue normalisation of the oracles (\S\ref{sec:example}).
\end{algorithm}

\begin{contractbox}
The three runtime assertions of Algorithm~\ref{alg:assembly} are stated as predicates
on emitted data in Obligation~\ref{obl:welldef} (for every $j$: $y_j$ satisfies every
relation of $\Rel$, and $d\,y_j=0$, over $\cring$, a failure aborting the run) and
Obligation~\ref{obl:agreement} (the two evaluation routes of
\S\ref{sec:equivalence}). They are \emph{not} restated here, per rule~\ref{A4}: an
obligation has one home.
\end{contractbox}

\paragraph{Code map.}
Definitions~\ref{def:inducedmodule}, \ref{def:inducedls},
Proposition~\ref{prop:blockmatrices} and \eqref{eq:MH}:
\file{Hecke/InducedLocalSystem.h/.cpp} --- \code{MakeInducedLocalSystemGamma0} builds
$L''$ from \code{HeckeCosetTables::perm}/\code{rho\_cocycle} (bottom-row class map) and
\code{MakeInducedLocalSystemGamma0Upper} builds $L'$ from
\code{perm\_upper}/\code{rho\_cocycle\_upper} (top-row class map), both reusing the base
relation list of \code{GenerateSL2ZRelationsWithGroupData} unchanged and replacing only
the matrix, which is Definition~\ref{def:inducedls} verbatim. \eqref{eq:MH} is assembled
as an \code{engine::BlockMat} with one block write per block row --- the one-non-zero-block-per-row
structure of Proposition~\ref{prop:blockmatrices} is the storage
model, not an optimisation. Proposition~\ref{prop:lhcompliant}(3) is the code path for
the \code{negI\_*} relations, which computes $M_H(-I)=I_m\otimes\rhon(-I)$ directly
rather than through the generator dictionary; (4) is
\code{HeckeShapiroValidation} (Obligation~\ref{obl:shapiro}), which hard-asserts
$\dim\Heq{0}=\dim\Heq{1}=1$ for $H=\Ga_0(2),\Ga_0(3)$ at $n=0$, $\ell\in\{5,101\}$
against classical genus and cusp data of $X_0(2)$, $X_0(3)$.
Propositions~\ref{prop:alpha} and \ref{prop:omega}, and
Proposition~\ref{prop:mu}: \file{Hecke/HeckeCrossChecks.h/.cpp}.
Definition~\ref{def:xi}, \eqref{eq:xiexplicit}, Theorem~\ref{thm:xi} and
Algorithm~\ref{alg:assembly}: \file{Hecke/HeckeOperatorOnK.h/.cpp} --- $\Th_g$ is stored
as \code{Theta\_block(q)}, an $N_q\times N_q$ block matrix whose pattern is exactly the
set of pairs $(\sigma,\sigma_t)$ named by the $\tau$-table
(\code{HeckeThetaBlockValidation} checks that the structural non-zeros are confined to
it), the flat form being \code{get\_Theta\_flat\_matrices()[q]}; the cocycle and
relation assertions of Algorithm~\ref{alg:assembly} are the compute-time checks (i) and
(ii), and $\Amat{q}$ is \code{get\_H\_matrices()[q]}. The flat layout used by $P_q$ and by the $z^{(q)}_j$ is
simplex-major, fibre-minor --- flat index $s\cdot d_V+k$ for the $s$-th simplex in the
stored order and fibre coordinate $k$ --- and is the $\code{nbc}=1$ block-matrix layout
of \code{BlockMat::to\_flat}, which is its single implementation.

\paragraph{Formalization note.}
Definition~\ref{def:inducedmodule} and Proposition~\ref{prop:blockmatrices} are pure
finite group theory --- a transversal, an index action and a $1$-cocycle --- and are the
easiest part of Part~III to formalise; \code{mathlib4} supplies induced modules and
Shapiro's lemma, so Proposition~\ref{prop:lhcompliant}(4) is a matter of matching
conventions rather than of proof. The load-bearing content of this chapter for a proof
assistant is \eqref{eq:psiinv}: it is the one place where a \emph{choice} (the lift
$\lift{\sigma}$) is turned into a well-defined map, and its well-definedness is a
statement about stabilisers, not about coefficients. A formalisation should therefore
take $\Psi$ of Theorem~\ref{thm:explicit_complex_equiv} as the primitive and define
$\Xig$ by \eqref{eq:xidef}, deriving \eqref{eq:xiexplicit} as a lemma --- the reverse of
the implementation's order, which computes \eqref{eq:xiexplicit} and never forms
$\phi_c$. Theorem~\ref{thm:xi}(1) and (2) are finite manipulations using only
Lemma~\ref{lem:cocycle}, Proposition~\ref{prop:symconv}(2) and (T1)--(T2), need no
hypothesis on $\ell$, and are the natural companion target to the equivariance half of
Theorem~\ref{thm:welldef}. Remark~\ref{rmk:nogo}(ii) is a genuinely non-constructive
argument (irreducibility plus a numerical non-scalarity) and should be formalised, if at
all, as a statement about $\End_{\Ga}(\Sym^{10}\Q^2)$ together with an imported
eigenvalue fact, not re-derived.
%----------------------------------------------------------------------
\section{Equivalent formulations and cross-checks}\label{sec:equivalence}
%----------------------------------------------------------------------
% SOURCE:  hecke_action.tex:1064-1161 (RH), plus the "Recovering the three-map
%          factorization" remark of hecke_action.tex:1040-1063, re-homed here
%          because it is a cross-check and not part of the binding pipeline.
% NOTATION: rows sigma, "t,j,s (running index)", "s_j, delta_j", mu, cring,
%           Rel, lift(sigma) of sec:notation-crosswalk applied.  The coset
%           index is j and the chain-term index t (C15); the middle map is mu
%           (C8); the bar-cochain arguments are gamma, gamma_1, gamma_2.
% ERRATA:  E20 (the missing coset permutation) (App. A)

\paragraph{What this chapter fixes.}
Three formulations that compute the \emph{same} operator as
Definition~\ref{def:hecke}/Definition~\ref{def:xi}, what each is good for, and at which
level each is asserted to agree. None of them is a definition; the normative operator is
Definition~\ref{def:hecke}, realised by Definition~\ref{def:xi} and
Algorithm~\ref{alg:assembly}. The three are: the equivariant-chain-map picture and its
two bracketings (\S\ref{sec:chainpicture}), which agree \emph{bit for bit} on the same
data; the classical three-map factorisation
$\Tdc{g}=\omega_{\ast}\circ\mu_{\ast}\circ\alpha_{\ast}$
(\S\ref{sec:threemap}), whose middle map is a cochain map with per-simplex blocks and
\emph{not} a morphism of local systems (Remark~\ref{rmk:nogo}); and the corrected
group-cochain formula in degree $1$ (\S\ref{sec:barformula}), which is an independent
route through the bar resolution and agrees only up to similarity.

\paragraph{Standing hypotheses.}
\ref{H1}--\ref{H4}, repeated at each point of use (rule~\ref{A1}). The hypothesis
$6\in\cring^{\times}$ (\ref{H3}) is needed by
Proposition~\ref{prop:barformula} --- which computes in $H^1(\Ga;\VR)$ and uses the
amalgam torsion --- and by nothing in \S\S\ref{sec:chainpicture}--\ref{sec:threemap},
whose statements are cochain-level identities over an arbitrary commutative $\cring$.

\subsection{The equivariant-chain-map picture, and the two bracketings}
\label{sec:chainpicture}

Definition~\ref{def:hecke} together with Theorem~\ref{thm:xi} \emph{is} the classical
``Hecke operator as an equivariant chain map'' picture, written out on an explicit
resolution. The datum $(\{s_j\},\tau)$ is a simplicial stand-in for the correspondence
of \S\ref{sec:correspondence}: $\sum_{j}s_j(-)$ plays the role of
$(\mathrm{pr}')^{\ast}$ --- unfolding over the $m$ sheets of
$\Ga''\backslash\Hyp\to\Ga\backslash\Hyp$ --- the map $\tau$ plays the role of
$\theta_g^{\ast}$, and the $\rhon(\delta_j)$-weighted sum plays the role of the trace
$(\mathrm{pr}'')_{\ast}$. Nothing in this paragraph is used in a proof; it is stated so
that a reader arriving from the classical literature can see which factor is which.

What \emph{is} used is that the single formula \eqref{eq:heckechain} admits two
different bracketings, and that they are required to agree exactly.

\begin{proposition}[The two evaluation routes agree on the nose]\label{prop:tworoutes}
\PROVED{} \LEAN{prop:tworoutes} Assume \ref{H4}, let $\cring$ be a commutative ring, and fix
$g=\smat{1}{0}{0}{N}$, the transporters $s_j$ and the $\tau$-table of
Definition~\ref{def:xi}. For $c\in\Cq{q}(K,\Rel;L)$ the two expressions
\[
  \underbrace{\sum_{j=1}^{m}\rhon(s_j^{-1})\sum_{t}n^{(j,\sigma)}_{t}\,
    \rhon\bigl(g\,\gamma^{(j,\sigma)}_{t}\bigr)\,c\bigl(\sigma^{(j,\sigma)}_{t}\bigr)}
    _{\text{fused: }\ \omega_{\bullet}\circ\Xig,\ \text{Definition~\ref{def:xi}}}
  \qquad\text{and}\qquad
  \underbrace{\sum_{j=1}^{m}\sum_{t}n^{(j,\sigma)}_{t}\,
    \rhon\bigl(\delta_j\,\gamma^{(j,\sigma)}_{t}\bigr)\,c\bigl(\sigma^{(j,\sigma)}_{t}\bigr)}
    _{\text{route 2: direct }\ \eqref{eq:heckechain}}
\]
are equal term by term, for every $\sigma\in\orient{q}(K)$, every $j$ and every $t$ ---
with \emph{no} hypothesis on $\ell$, and \emph{as an identity of matrices, not only of
their values}.
\end{proposition}

\begin{proof}
$\delta_j\gamma_t=s_j^{-1}(g\gamma_t)$ is an identity in $\Mat_2(\Z)$, between integral
matrices and with no inverse of a determinant-$N$ matrix taken
(Proposition~\ref{prop:cosets}(5)). Apply
$\rhon$ and Proposition~\ref{prop:symconv}(2), which is multiplicativity of $\rhon$ as a
monoid homomorphism on all of $\Mat_2(\Z)$ and in particular on the non-invertible
elements: $\rhon(\delta_j\gamma_t)=\rhon(s_j^{-1})\rhon(g\gamma_t)$.
\end{proof}

\begin{contractbox}
Proposition~\ref{prop:tworoutes} is the implementer's cross-check, and it is one of
\textbf{two} places in this document where \textbf{bit-identity of cochain-level
matrices} is asserted rather than equality in cohomology --- the other being
Proposition~\ref{prop:mu}(3), $\omega_{\bullet}\circ\mu_{\bullet}\circ\alpha_{\bullet}
=\Th_g$ --- legitimately, because the two routes consume the \emph{same} $g$, the
\emph{same} $s_j$ and the \emph{same} $\tau$-table, so Remark~\ref{rmk:indeplevel}'s
prohibition does not apply. The two sides are formed
differently on purpose: the fused route multiplies two \emph{already reduced} stored
matrices $\rhon(s_j^{-1})$ and $\rhon(g\gamma_t)$ over $\cring$, whereas route 2 forms
the integer product $\delta_j\gamma_t$ \emph{exactly} in $\Mat_2(\Z)$ and reduces once,
at the end. Agreement therefore tests the reduction discipline and the monoid identity
of Proposition~\ref{prop:symconv}(2) simultaneously, and a discrepancy is a bug in one
of them and in nothing else. Recorded as Obligation~\ref{obl:agreement}.
\end{contractbox}

\subsection{The three-map factorisation, at cochain level}\label{sec:threemap}

Remark~\ref{rmk:nogo} showed that the middle arrow cannot be a morphism of local
systems. It \emph{can} be a cochain map, and the following makes it explicit. The price
is that it is given by one block matrix per pair of simplices --- a
\code{FiberBlockTable} --- rather than by a single fibre matrix, which is exactly the
generalisation Remark~\ref{rmk:nogo} says is needed.

\begin{proposition}[The middle cochain map $\mu_{\bullet}$]\label{prop:mu}
\PROVED{} Assume \ref{H4}, let $\cring$ be a commutative ring, and let
$g$, $\{s_j\}$, $\tau$, $L'=L_{\Ga^{0}(N)}$ with representatives $t_1=I,\dots,t_m$, and
$L''=L_{\Ga_0(N)}$ with representatives $s_1,\dots,s_m$ be as in
\S\ref{sec:inducedmodule}. For $c'\in\Cq{q}(K,\Rel;L')$ define
\[
  \bigl(\mu_{\bullet}c'\bigr)(\sigma)_j
  \;:=\;\rhon(g)\cdot\Bigl[\,\phi^{L'}_{c'}\bigl(\tau(s_j\cdot\lift{\sigma})\bigr)\Bigr]_{1},
  \qquad j=1,\dots,m,
\]
where $\phi^{L'}_{c'}$ is the $\Ind_{\Ga'}^{\Ga}\VR$-valued equivariant cochain attached
to $c'$ by \eqref{eq:psiinv} applied to the local system $L'$, i.e.\
$\phi^{L'}_{c'}(\gamma\cdot\lift{\sigma})=M_{\Ga'}(\gamma)\,c'(\sigma)$, and
$[\,\cdot\,]_{1}$ is evaluation at $1\in\Ga$, i.e.\ block coordinate $1$ (legitimate
because $t_1=I$). \textbf{The order of the two operations is part of the definition}:
$[\,\cdot\,]_{1}$ lands in $\VR$ and $\rhon(g)$ is applied there, so no action of
$\rhon(g)$ on $\Ind_{\Ga'}^{\Ga}\VR$ is presupposed. Evaluation at $1$ is a
$\Ga'$-morphism $\Ind_{\Ga'}^{\Ga}\VR\to\VR$:
\begin{equation}\label{eq:evalone}
  \bigl[u\cdot f\bigr]_{1}=\rhon(u)\,\bigl[f\bigr]_{1}\qquad\text{for all }u\in\Ga' ,
\end{equation}
because $1\cdot u=1$ and $h_1(u)=t_1ut_1^{-1}=u$ in \eqref{eq:MH}. Then:
\begin{enumerate}
\item \textbf{Closed form.} \LEAN{prop:mu-assembled} (the three factor clauses, assembled; not (2), (3)) With $\tau(s_j\lift{\sigma})=\sum_t n_t\,\gamma_t\cdot
  \lift{\sigma_t}$ as in Definition~\ref{def:xi}, and with $a_t:=1\cdot\gamma_t$ the
  index of the $\Ga^{0}(N)$-coset of $\gamma_t$,
  \[
    \bigl(\mu_{\bullet}c'\bigr)(\sigma)_j
    \;=\;\sum_{t} n_t\ \rhon\bigl(g\,\gamma_t\,t_{a_t}^{-1}\bigr)\
    \bigl[c'(\sigma_t)\bigr]_{a_t},
  \]
  all factors integral of determinant $N$, with $t_{a_t}^{-1}\in\Ga$ and no
  determinant-$N$ matrix inverted.
\item \textbf{It is a cochain map into $L''$.} $\mu_{\bullet}$ carries
  $\Cq{\bullet}(K,\Rel;L')$ into $\Cq{\bullet}(K,\Rel;L'')$ and commutes with $d$ ---
  with no hypothesis on $\ell$.
\item \textbf{It splits $\Xig$.} $\mu_{\bullet}\circ\alpha_{\bullet}=\Xig$ as maps
  $C^{q}(K;\VR)\to C^{q}(K;\Ind_{\Ga''}^{\Ga}\VR)$, \emph{entry for entry}, where
  $\alpha_{\bullet}$ is the simplexwise application of the block column of
  Proposition~\ref{prop:alpha}. Consequently
  $\omega_{\bullet}\circ\mu_{\bullet}\circ\alpha_{\bullet}=\Th_g$, and if moreover
  $6\in\cring^{\times}$ then on cohomology
  $\Tdc{g}=\omega_{\ast}\circ\mu_{\ast}\circ\alpha_{\ast}$ under the identifications of
  Proposition~\ref{prop:lhcompliant}(4).
\end{enumerate}
\end{proposition}

\begin{proof}
(1) By \eqref{eq:psiinv} for $L'$,
$\phi^{L'}_{c'}(\tau(s_j\lift{\sigma}))=\sum_t n_t\,M_{\Ga'}(\gamma_t)\,c'(\sigma_t)$.
By \eqref{eq:MH} the block row $1$ of $M_{\Ga'}(\gamma_t)$ has its single non-zero block
$\rhon(h_1(\gamma_t))$ in block column $a_t=1\cdot\gamma_t$, and
$h_1(\gamma_t)=t_1\gamma_tt_{a_t}^{-1}=\gamma_tt_{a_t}^{-1}\in\Ga^{0}(N)$
(Proposition~\ref{prop:blockmatrices}). Hence
$\bigl[\phi^{L'}_{c'}(\tau(s_j\lift{\sigma}))\bigr]_{1}
=\sum_tn_t\,\rhon(\gamma_tt_{a_t}^{-1})\bigl[c'(\sigma_t)\bigr]_{a_t}$; multiplying by
$\rhon(g)$ and using Proposition~\ref{prop:symconv}(2) on the integral identity
$g\cdot(\gamma_tt_{a_t}^{-1})=g\gamma_tt_{a_t}^{-1}$ gives the display.

(2) Let $(\sigma_{\mathrm{lhs}},\sigma_{\mathrm{rhs}},\gamma)\in\Rel$ and write
$s_j\gamma=h_j(\gamma)s_{j\cdot\gamma}$ with $h_j(\gamma)\in\Ga''$
(Lemma~\ref{lem:cocycle}), so that $u:=g^{-1}h_j(\gamma)g\in\Ga'$
(Proposition~\ref{prop:cosets}(3)). By (T1) of Definition~\ref{def:tau} and
$\Ga$-equivariance of $\phi^{L'}_{c'}$,
\[
  \bigl(\mu_{\bullet}c'\bigr)(\sigma_{\mathrm{lhs}})_j
  =\rhon(g)\Bigl[\,u\cdot\phi^{L'}_{c'}\bigl(\tau(s_{j\cdot\gamma}\cdot
    \lift{\sigma_{\mathrm{rhs}}})\bigr)\Bigr]_{1}
  \overset{\eqref{eq:evalone}}{=}\rhon(g)\,\rhon(u)\,
  \Bigl[\phi^{L'}_{c'}\bigl(\tau(s_{j\cdot\gamma}\cdot\lift{\sigma_{\mathrm{rhs}}})\bigr)\Bigr]_{1}
  =\rhon\bigl(h_j(\gamma)\bigr)\bigl(\mu_{\bullet}c'\bigr)(\sigma_{\mathrm{rhs}})_{j\cdot\gamma},
\]
the last step by Proposition~\ref{prop:symconv}(2) applied to $g\,u=h_j(\gamma)\,g$ in
$\Mat_2(\Z)$; by \eqref{eq:MH} this is block $j$ of
$M_{\Ga''}(\gamma)(\mu_{\bullet}c')(\sigma_{\mathrm{rhs}})$. Commutation with $d$ is the
computation of Theorem~\ref{thm:xi}(2) verbatim, which used only (T2),
$\lift{\partial\sigma}=\partial\lift{\sigma}$ and linearity of $[\,\cdot\,]_{1}$.

(3) By Proposition~\ref{prop:alpha}, $\bigl[\alpha_{\bullet}c(\sigma)\bigr]_{a}
=\rhon(t_{a})c(\sigma)$. Substituting into (1) and using
Proposition~\ref{prop:symconv}(2) on the integral identity
$(g\gamma_tt_{a_t}^{-1})\,t_{a_t}=g\gamma_t$ in $\Mat_2(\Z)$,
\[
  \bigl(\mu_{\bullet}\alpha_{\bullet}c\bigr)(\sigma)_j
  =\sum_t n_t\,\rhon\bigl(g\gamma_tt_{a_t}^{-1}\bigr)\rhon\bigl(t_{a_t}\bigr)c(\sigma_t)
  =\sum_t n_t\,\rhon\bigl(g\gamma_t\bigr)c(\sigma_t)
  =\bigl(\Xig c\bigr)(\sigma)_j
\]
by \eqref{eq:xiexplicit}. Composing with $\omega_{\bullet}$ gives $\Th_g$
(Definition~\ref{def:xi}); the cohomological statement is
Theorem~\ref{thm:xi}(3) together with Propositions~\ref{prop:alpha} and
\ref{prop:omega}, whose identifications with $\res$ and $\cor$ carry
$6\in\cring^{\times}$ and are the only place it is used.
\end{proof}

\begin{contractbox}
Proposition~\ref{prop:mu} is a \textbf{cross-check route only}. It is what makes each of
the three classical factors separately testable --- $\alpha_{\bullet}$ and
$\omega_{\bullet}$ through the single-fibre-matrix machinery of
\S\ref{sec:systematiccomputation}, $\mu_{\bullet}$ through per-simplex block rows ---
and it is asserted against the binding pipeline at three levels of increasing strength:
$\mu_{\bullet}$ alone maps $L'$-equivariant cochains to $L''$-equivariant ones;
$\mu_{\bullet}\alpha_{\bullet}=\Xig$ on the full standard basis of $C^q(K;\VR)$, not
merely on $\im\alpha_{\bullet}$; and
$\omega_{\bullet}\mu_{\bullet}\alpha_{\bullet}=\Th_g$ bit-identically as matrices for
every $q$. The binding operator is and remains $\Th_g$ of
Definition~\ref{def:xi}, which constructs neither $L'$ nor $\mu_{\bullet}$ and needs
neither. Recorded as Obligation~\ref{obl:agreement}.
\end{contractbox}

\subsection{The corrected group-cochain formula}\label{sec:barformula}

\begin{remark}[Deviation/clarification vs.\ \file{monodromy\_v5.tex} \S12.3]
\label{rmk:barformula}
\PROVED{} The superseded framework document states, as its ``Algebraic Hecke
Operator'',
\[
  (T_g\phi)(\gamma_1,\dots,\gamma_s)
  =\sum_{j}\rho_L(g_j)\,\phi\bigl(g_j^{-1}\gamma_1g_j,\dots,g_j^{-1}\gamma_sg_j\bigr).
\]
As written this is not well defined on inhomogeneous cochains --- $g_j^{-1}\gamma g_j$
need not lie in $\Ga$ --- and, even read inside the monoid $\De$, it \textbf{does not
commute with the bar differential}, because right multiplication by $\gamma$
\emph{permutes} the cosets $\delta_j\Ga$ and the formula holds every $j$ fixed. It is
the special case of the correct formula in which every $\gamma_u$ stabilises every
coset. Proposition~\ref{prop:barformula} carries the permutation explicitly. Recorded as
Appendix~\ref{app:errata}, entry~\ref{err:e20}.
\end{remark}

\begin{proposition}[Corrected bar-cochain formula, degree $1$]\label{prop:barformula}
\PROVED{} Assume \ref{H4} and let $\cring$ be a commutative ring. Write
$\Ga g\Ga=\bigsqcup_{j=1}^{m}\delta_j\Ga$ for the left-coset decomposition of
Proposition~\ref{prop:cosets}(2),(5), with $\delta_j=s_j^{-1}g$. For $\gamma\in\Ga$ let
$\gamma\ast j$ be defined by $\gamma\delta_j\Ga=\delta_{\gamma\ast j}\Ga$ and put
\[
  \kappa_j(\gamma)\;:=\;\delta_{\gamma\ast j}^{-1}\,\gamma\,\delta_j\ \in\ \Ga .
\]
Then on inhomogeneous $1$-cochains $\phi\colon\Ga\to\VR$ the assignment
\begin{equation}\label{eq:barcorrected}
  (T_g\phi)(\gamma)\;:=\;\sum_{j=1}^{m}\rhon\bigl(\delta_{\gamma\ast j}\bigr)\,
  \phi\bigl(\kappa_j(\gamma)\bigr)
\end{equation}
carries cocycles to cocycles and coboundaries to coboundaries --- with \emph{no}
hypothesis on $\ell$ --- and if moreover $6\in\cring^{\times}$ (hypothesis \ref{H3}) it
induces on $H^1(\Ga;\VR)$ the same operator as Definition~\ref{def:hecke}.
\end{proposition}

\begin{proof}
\emph{Chain rule.} $\gamma\ast j$ and $\kappa_j(\gamma)$ are well defined because the
$\delta_j\Ga$ are $m$ disjoint cosets permuted by left multiplication
(Proposition~\ref{prop:cosets}(2)), and $\kappa_j(\gamma)\in\Ga$ because
$\gamma\delta_j$ and $\delta_{\gamma\ast j}$ lie in the same left coset. From
$\gamma_1\gamma_2\delta_j=\gamma_1\delta_{\gamma_2\ast j}\kappa_j(\gamma_2)
=\delta_{\gamma_1\ast(\gamma_2\ast j)}\kappa_{\gamma_2\ast j}(\gamma_1)\kappa_j(\gamma_2)$
and uniqueness of the coset:
\[
  (\gamma_1\gamma_2)\ast j=\gamma_1\ast(\gamma_2\ast j),\qquad
  \kappa_j(\gamma_1\gamma_2)=\kappa_{\gamma_2\ast j}(\gamma_1)\,\kappa_j(\gamma_2).
\]
\emph{Cocycles.} Let $\phi(ab)=\phi(a)+\rhon(a)\phi(b)$. Then
\begin{align*}
  (T_g\phi)(\gamma_1\gamma_2)
  &=\sum_j\rhon\bigl(\delta_{\gamma_1\ast(\gamma_2\ast j)}\bigr)
    \Bigl[\phi\bigl(\kappa_{\gamma_2\ast j}(\gamma_1)\bigr)
      +\rhon\bigl(\kappa_{\gamma_2\ast j}(\gamma_1)\bigr)\phi\bigl(\kappa_j(\gamma_2)\bigr)\Bigr]\\
  &=(T_g\phi)(\gamma_1)
    +\sum_j\rhon\bigl(\delta_{\gamma_1\ast j'}\kappa_{j'}(\gamma_1)\bigr)\Big|_{j'=\gamma_2\ast j}
      \phi\bigl(\kappa_j(\gamma_2)\bigr)
  =(T_g\phi)(\gamma_1)+\rhon(\gamma_1)\,(T_g\phi)(\gamma_2),
\end{align*}
using $\delta_{\gamma_1\ast j'}\kappa_{j'}(\gamma_1)=\gamma_1\delta_{j'}$, then
Proposition~\ref{prop:symconv}(2), and that $j\mapsto\gamma_2\ast j$ is a bijection --- the
reindexing that Remark~\ref{rmk:barformula} says the superseded formula omits, and
without which the first sum is not $(T_g\phi)(\gamma_1)$.
\emph{Coboundaries.} For $\phi=dv$, i.e.\ $\phi(\gamma)=\rhon(\gamma)v-v$,
\[
  (T_g\,dv)(\gamma)
  =\sum_j\Bigl[\rhon\bigl(\delta_{\gamma\ast j}\kappa_j(\gamma)\bigr)v
    -\rhon\bigl(\delta_{\gamma\ast j}\bigr)v\Bigr]
  =\sum_j\Bigl[\rhon(\gamma\delta_j)v-\rhon(\delta_{\gamma\ast j})v\Bigr]
  =\rhon(\gamma)w-w,\quad w:=\sum_j\rhon(\delta_j)v,
\]
using Proposition~\ref{prop:symconv}(2) and, in the last step, that
$j\mapsto\gamma\ast j$ is a bijection; so $T_g(dv)=dw$.
\emph{Agreement.} Both operators compute $\cor\circ(\conj_g)_{\ast}\circ\res$ of
\eqref{eq:corconjres}: \eqref{eq:barcorrected} is the standard double-coset action on
the bar resolution \STD{\cite[Ch.~III.9--10]{brown} for $\res$, $\cor$ and the bar
resolution; \cite[\S1.1, p.~194]{ash_stevens} for the double-coset operator (the right-module form on homogeneous cochains; the
left-module form by the $S^{-1}$ device on the same page)}, with $\{\delta_j\}=\{s_j^{-1}g\}$ the representatives of
Proposition~\ref{prop:cosets}(5), and Definition~\ref{def:hecke} computes the same thing
on the resolution $C_{\bullet}(\Ku)\otimes\cring$ by Theorem~\ref{thm:classical}. Both
resolutions are complexes of relatively projective $\cring[\Ga]$-modules when
$6\in\cring^{\times}$ --- for $C_{\bullet}(\Ku)\otimes\cring$ this is the computation in
the proof of Lemma~\ref{lem:tauexists}(2), for the bar resolution it is freeness --- so the
canonical comparison equivalence between them is unique up to homotopy and transports
one operator to the other \STD{\cite[Ch.~I.7.4]{brown}}. This is the \emph{only} use of
$6\in\cring^{\times}$ in this proposition.
\end{proof}

\begin{remark}[$\kappa_j(\gamma)$ is computed by an exact division, not by an inversion]
\label{rmk:kappaexact}
The expression $\kappa_j(\gamma)=\delta_{\gamma\ast j}^{-1}\gamma\delta_j$ displays an
inverse of a determinant-$N$ matrix, which Obligation~\ref{obl:integral} forbids in the
pipeline. It is not needed: writing $\delta_k=s_k^{-1}g$ and
$N\,g^{-1}=\adjg(g)$, the identity
\[
  N\cdot\kappa_j(\gamma)\;=\;\adjg(g)\,s_{k}\,\gamma\,s_j^{-1}\,g,
  \qquad k=\gamma\ast j,
\]
holds in $\Mat_2(\Z)$, so $\kappa_j(\gamma)$ is obtained by forming an integer product
and dividing by $N$ \emph{exactly}; the index $k$ is the unique one for which the
division is exact, and the result is verified by the reconstruction identity
$\delta_k\,\kappa_j(\gamma)=\gamma\,\delta_j$ over $\Z$. This is what the implementation
does (rule~\ref{A3}), and it is why an oracle route inherits the integrality discipline
of the pipeline rather than relaxing it.
\end{remark}

\begin{contractbox}
Proposition~\ref{prop:barformula} is a \textbf{validation oracle only}: an independent
implementation path for $H^1$ through explicit group cocycles on a presentation of
$\Ga$, touching neither the $\tau$-table nor the induced systems nor the simplicial
cochain spaces. The normative operator in all degrees is Definition~\ref{def:hecke} /
Algorithm~\ref{alg:assembly}. Because the two routes produce matrices in
\emph{different, unrelated bases} of $H^1$, they are compared by \textbf{similarity
invariants --- characteristic polynomial and minimal polynomial --- and not entrywise};
asserting an entrywise identity here would be asserting something this document does not
claim (Remark~\ref{rmk:indeplevel}). Any discrepancy of those invariants is a bug.
Recorded as Obligation~\ref{obl:agreement}.
\end{contractbox}

\subsection{Degree-wise specialisation}\label{sec:degreewise}

The pipeline is uniform in $q$ --- one formula, one implementation, every degree --- and
this subsection records what it specialises to, for orientation only. No statement below
is used elsewhere.

\begin{itemize}
\item \textbf{$q=0$.} $\Heq{0}(K,\Rel;L)=\{v\in\VR:\rhon(\gamma)v=v$ for every relation
  element $\gamma$ of $\Rel\}$ and, by Remark~\ref{rmk:T3},
  $T\,v=\sum_{j=1}^{m}\rhon(\delta_j)v$. For $n=0$ this is multiplication by $m$:
  $T_p=p+1$, and at the bad prime $U_\ell=\ell+1\equiv1\pmod\ell$ (hypothesis \ref{H2}).
\item \textbf{$q=1$.} The operator matches the classical modular-symbol and
  group-cocycle computations: Proposition~\ref{prop:barformula} is the second route, and
  the Eichler--Shimura eigenvalues are the oracle (Obligation~\ref{obl:oracle}).
\item \textbf{$q=2$.} No classical shortcut exists and the pipeline applies unchanged.
  For $6\in\cring^{\times}$ the model has $\Heq{2}(K,\Rel;L)=0$ --- $\Xbulk$ is
  $2$-dimensional and $\Ga$ has virtual cohomological dimension $1$
  \STD{\cite[Ch.~VIII, \S\S2--3]{brown}; \cite{deng} for this model} --- so
  $\Amat{2}$ is the $0\times0$ matrix, which is itself a useful consistency assertion.
  For $\ell\in\{2,3\}$, $\Heq{2}$ may be non-zero and Algorithm~\ref{alg:assembly}
  produces its matrix by the same code path, with the caveat of \ref{H3} on what that
  group is.
\end{itemize}

\paragraph{Code map.}
Proposition~\ref{prop:tworoutes}: \file{Hecke/HeckeOperatorOnK.cpp} --- the two
predefined programs over the $\tau$-pattern (\code{program\_fused} and
\code{program\_route2}, the latter in its own block pool so that content-addressed
deduplication cannot make the comparison a tautology), compared as mathematical matrices
by \code{blk::equal} over all degrees, with \code{apply\_Theta} and
\code{apply\_route2} the vector forms; any mismatch throws, and
\code{HeckeThetaBlockValidation} additionally checks
$\code{apply\_Theta}(q,e_b)=\Th^{(q)}_{\mathrm{flat}}[:,b]=\code{apply\_route2}(q,e_b)$
on the whole standard basis. Proposition~\ref{prop:mu}:
\file{Hecke/HeckeCrossChecks.h/.cpp} --- $\mu_{\bullet}$ is the
\code{FiberBlockTable} whose $(\sigma,\sigma')$ entry accumulates
$\sum_{t:\sigma_t=\sigma',\,a_t=a}n_t\,\rhon(g\gamma_tt_{a}^{-1})$ into block $(j,a)$,
the product $g\gamma_tt_{a}^{-1}$ formed exactly in $\Mat_2(\Z)$ and reduced once; the
three assertions of the contract box are the class's three checks, gated by
\code{HeckeCrossCheckValidation}. Proposition~\ref{prop:barformula} and
Remark~\ref{rmk:kappaexact}: \file{Hecke/BarCocycleH1.h/.cpp}, which presents
$\Ga=\langle S\rangle\ast_{\langle S^2\rangle}\langle ST\rangle$, represents a
$1$-cocycle by the pair $(\phi(S),\phi(ST))$ subject to the norm and amalgam-matching
conditions, evaluates $\phi$ on an arbitrary $\kappa\in\Ga$ by reducing $\kappa$ to a
word in $T,S$ and expanding with the cocycle rule (asserting both that the word
reconstructs $\kappa$ exactly and that the accumulated representation matrix equals
$\rhon(\kappa)$), enforces $\ell\ge5$ at construction, and compares
characteristic and minimal polynomials against $\Amat{1}$; the exact division of
Remark~\ref{rmk:kappaexact} and its reconstruction check are in \code{hecke\_matrix()}.

\paragraph{Formalization note.}
Proposition~\ref{prop:tworoutes} is a one-line consequence of
Proposition~\ref{prop:symconv}(2) and is the cheapest statement in Part~III to
formalise; it is also among the most valuable to state separately, because it is one of
only \textbf{two} claims in the document where a \emph{bit}-level agreement is
mathematically forced rather than empirically observed --- the other being
Proposition~\ref{prop:mu}(3), proved by the same route: an integral identity in
$\Mat_2(\Z)$ pushed through the monoid homomorphism $\rhon$
(Proposition~\ref{prop:symconv}(2)). Proposition~\ref{prop:mu} is finite linear algebra
over the same data and needs nothing beyond \eqref{eq:MH}. Proposition~\ref{prop:barformula}
splits cleanly for a proof assistant: the cocycle and coboundary computations are
elementary and unconditional --- the chain rule for $(\ast,\kappa)$ is the same
$1$-cocycle structure as Lemma~\ref{lem:cocycle} and should share its formalisation ---
whereas the agreement clause is a comparison of two resolutions and should be stated as
an instance of the same comparison theorem used by Lemma~\ref{lem:tauexists}(2), not
re-proved. The degree-wise specialisations of \S\ref{sec:degreewise} are orientation and
carry no formalisation obligation, with one exception worth recording: the vanishing
$\Heq{2}=0$ at $6\in\cring^{\times}$ is carried in \S\ref{sec:degreewise} by a citation
(rule~\ref{A2}) and is \emph{not} derived here, so a formalisation must import it rather
than take the emitted $0\times0$ matrix as a proof of it.

%======================================================================
\part*{Part IV --- The boundary, the cone, and compact supports}
\addcontentsline{toc}{section}{Part IV --- The boundary, the cone, and compact supports}
%======================================================================

%----------------------------------------------------------------------
\section{The boundary datum}\label{sec:boundarydatum}
%----------------------------------------------------------------------
% SOURCE:  RC+NEW.  Implementation plan SS0.4 C-A2 (hecke_extensions_implementation_plan.md:
%          312-382) and C-A3(ii.2) (:412-429); OQ2 verdict [V] SS2.6, 5.1-5.5
%          (hecke_extensions_OQ2_verdict.md); specification (A1)
%          (hecke_extensions.tex:761-788); owner decisions OP1, OP2, OP5 (plan S11).
%          No superseded document states the Gamma-level tau_bdry, its existence proof,
%          or the ell=2 unboundedness in this form.
% NOTATION: C1 (q is the cohomological degree, p the Hecke prime), C7 (delta is a coset
%          element; d is the coboundary), C8 (beta_t is the RIGHT-coset representative),
%          C11 (nu is the exponent, r is the restriction cochain map), C12 (the
%          coefficient module is V_R; boundary (co)invariants are V_R^T and (V_R)_T),
%          C14/C15 (s_j the transporters, t a running index, never i), C16 (\lift).
% ERRATA:  E4 (tau_bdry uniqueness argued circularly), E12 (H(bdry) defined as one
%          group and constructed as another), E14's boundary half (the characters in
%          the beta_j convention) -- App. A.

\paragraph{What this chapter fixes.}
The boundary of the model, as a Hecke-equivariant object, and nothing beyond it. The
boundary group $\Gbdry$ and the \emph{definition} $H^{\bullet}(\partial;\VR):=
H^{\bullet}(\Gbdry;\VR)$; the circle model, and the exact hypothesis --- $2\in
\cring^{\times}$ --- under which it computes that group; the upper-triangular set $B_N$,
its left cosets, and why the right cosets of Proposition~\ref{prop:cosets}(1) are the
wrong ones; the boundary $g$-twisted chain datum $\taub$ on $\Xbdry$, which is the
\textbf{normative} boundary datum, its existence and homotopy uniqueness over $\Z$ for
\emph{every} $\ell$, and its unconditional integrality; the boundary operator
$\Thb_g$ and the boundary characters $\chi_b$; and what $\ell=2$ costs, where
$H^{\bullet}(\partial)$ is unbounded above in degree and the six-term sequence of
\S\ref{sec:six} is a \emph{window} of a sequence that does not terminate --- and
where the bulk, unlike the boundary, is \emph{periodic} above degree $2$ (Lemma~\ref{lem:bulkperiodic},
Hecke-equivariantly). The $\ell=2$ truncation policy is construction
item~\ref{open:ell2} (\textbf{OC-1}), decided by the owner on 2026-09-16 as
\ref{oc1a} and stated in \S\ref{sec:ell2} with its alternatives kept reachable.
\textbf{Since v1.3.3 this chapter is implemented}: the certified $\taub$-table, the
boundary coset table and the operator $\Thb_g$ exist in the tree, and
Obligation~\ref{obl:taubdry} is discharged in full by two registered validations ---
the code map at the end of the chapter says what was built, which choices the
implementation fixed where the statements below leave them open, and what it recorded.

\textbf{Three errata land here} (Appendix~\ref{app:errata}): \ref{err:e12}, a definition
and its own normative construction disagreeing inside one document, repaired by
Definition~\ref{def:bdrygroup} together with Proposition~\ref{prop:circlemodel};
\ref{err:e4}, a uniqueness statement derived from a theorem whose hypothesis the
construction exists to avoid, repaired by Lemma~\ref{lem:taubdryexists}, which never
cites that theorem; and the boundary half of \ref{err:e14}, the characters presented in
the right-coset convention as though it were the answer, repaired by
Proposition~\ref{prop:bdrychar}.

\textbf{What this chapter does not fix.} The pair $(\Xbulk,\Xbdry)$, the restriction map
$\rres=j^{\ast}$, the pair homotopy $\etapair$, the cone $\Cone{\rres}$ and the cone
action $\Phig$ are \S\ref{sec:cone}; the six groups and the long exact sequence are
\S\ref{sec:six}; the free $\Ga$-model that \ref{open:ell2} will be constructed on is
\S\ref{sec:orbifoldmodel}. $\taub$ is the boundary \emph{datum}; every statement below
that mentions the cone is a forward reference and is marked as one.

\paragraph{Standing hypotheses.}
\ref{H1}--\ref{H4} of \S\ref{sec:conventions}, each repeated where it is used
(rule~\ref{A1}), together with three hypotheses proper to this chapter. Per the design
direction recorded in Remark~\ref{rmk:orbifold} they are stated for \emph{the boundary
triangulation in force} --- an arbitrary $\Ga$-stable boundary subcomplex satisfying
\ref{B1}--\ref{B3} --- and not for the four horocycle vertices $v_0v_1v_2v_3$ of
Definition~\ref{def:model7}, so that a refinement of that complex is admissible input.
Nothing in this chapter is a statement about seven vertices.
\LEAN{def:model7bd} (on the model, for $\Xbdry$ given by the boundary list $v_0,v_1,v_2,v_3,\{v_0,v_1\},\{v_1,v_2\},\{v_2,v_3\}$: \ref{B1} as a
$\Ga$-saturated, face-closed subcomplex; of \ref{B2} only $C_q(\Xbdry)=0$ for $q\ge2$ --- the component map and the contractibility of the
components are not formalised; \ref{B3} as stabilisers of order $2$ containing $\pm I$; [GT] numbers these conditions differently)

\begin{enumerate}[label=(B\arabic*)]
\item\label{B1} \textbf{The boundary subcomplex.} $\Xbdry\subseteq\Xbulk$ is a
  $\Ga$-stable full subcomplex of the tessellation of \ref{H4}, namely the union of the
  $\Ga$-translates of the horocycle at $\infty$ at the truncation height $Y$. Writing
  $\Xbdry_{\infty}$ for the connected component lying over the cusp $\infty$, the
  setwise stabiliser $\{\gamma\in\Ga:\gamma\Xbdry_{\infty}=\Xbdry_{\infty}\}$ is the
  group $\Gbdry$ of Definition~\ref{def:bdrygroup}, and
  $\Xbdry=\Ga\cdot\Xbdry_{\infty}$. Moreover $\Ga$ has \emph{finitely many} orbits of
  simplices on $\Xbdry$ --- equivalently $\Gbdry$ has finitely many on
  $\Xbdry_{\infty}$ --- which is what makes the tables below finite.
\item\label{B2} \textbf{Components.} The map $\pi_0(\Xbdry)\to\mathbb{P}^1(\Q)$ sending a
  component to the cusp it lies over is a $\Ga$-equivariant bijection, and \textbf{every
  component is a $1$-dimensional contractible simplicial complex} --- a triangulated
  line. In particular $Z_1(\Xbdry_c)=H_1(\Xbdry_c)=0$ and $C_q(\Xbdry)=0$ for $q\ge2$.
\item\label{B3} \textbf{Stabilisers are central.} The $\Ga$-stabiliser of every simplex
  of $\Xbdry$ is exactly $\{\pm I\}$, which acts trivially on $\Xbulk$ (\ref{H4}(iv)).
  \textbf{No elliptic point lies on $\Xbdry$}, and no simplex of $\Xbdry$ has a
  stabiliser of order $4$ or $6$.
\end{enumerate}

\begin{remark}[The shipped triangulation satisfies \ref{B1}--\ref{B3}, and so does every equivariant refinement]\label{rmk:bdrymodel}
\PROVED{} For the model of Definition~\ref{def:model7}, $\Kbdry$ is the path
$v_0v_1v_2v_3$ on the horocycle $\{\mathrm{Im}\,z=Y\}$ with $v_3=T\cdot v_0$, and
$\Xbdry_{\infty}=\Gbdry\cdot\Kbdry$ is the whole horocycle, a line with three
$\Gbdry$-orbits of edges and three of vertices ($v_0$ and $v_3$ being in one). \ref{B1}
holds because the horocycle at $\infty$ is a line and its setwise stabiliser in $\Ga$ is
$\mathrm{Stab}_{\Ga}(\infty)=\Gbdry$; \ref{B2} holds because $\Ga$ acts transitively on
$\mathbb{P}^1(\Q)$ with stabiliser $\Gbdry$ at $\infty$ and the horoballs removed in
\ref{H4} are indexed by the cusps, so $\pi_0(\Xbdry)\cong\Ga/\Gbdry\cong\mathbb{P}^1(\Q)$
equivariantly; \ref{B3} holds because \ref{H4} fixes $Y>1$ while
$\mathrm{Im}\,\zeta=\mathrm{Im}\,\zeta^{2}=\sqrt3/2$ and $\mathrm{Im}\,i=1$, so the three
elliptic points of \S\ref{sec:model} lie strictly below the horocycle and off $\Xbdry$,
and \ref{H4}(iv)'s ``every other simplex stabiliser is exactly $\{\pm I\}$'' applies to
every simplex of $\Xbdry$.

If $\Xbulk'$ is any $\Ga$-equivariant refinement of $\Xbulk$ and $\Xbdry'$ the induced
subdivision of $\Xbdry$, then \ref{B1}--\ref{B3} again hold: subdividing a line gives a
line, refinement does not change $\pi_0$, and the stabiliser of a simplex of a
subdivision is contained in the stabiliser of the simplex containing it, hence still
$\{\pm I\}$ (it contains $-I$, which acts trivially). \textbf{Every statement of this
chapter is therefore inherited by any refinement}, which is what
Remark~\ref{rmk:orbifold}'s design direction asks of it. What is not inherited is
triangulation-indexed data: the number of $\Gbdry$-orbits of boundary simplices, the
size of the $\taub$-table, and the bound of Lemma~\ref{lem:bounds}.
\end{remark}

\subsection{The boundary group, and what the circle model computes}\label{sec:bdrygroup}

\begin{definition}[The boundary group and the boundary cohomology]\label{def:bdrygroup}
$\Gbdry:=\mathrm{Stab}_{\Ga}(\infty)=\bigl\{\smat{a}{b}{0}{d}\in\Ga\bigr\}$, the
subgroup of upper-triangular elements of $\Ga=\SL_2(\Z)$. Since $ad=1$ with
$a,d\in\Z$,
\begin{equation}\label{eq:bdrysplit}
  \Gbdry=\{\pm T^{k}:k\in\Z\}=\langle T\rangle\times\langle -I\rangle\;\cong\;\Z\times\Z/2 ,
\end{equation}
a direct product because $\langle T\rangle\cap\langle-I\rangle=\{I\}$ and $-I$ is
central. For a commutative ring $\cring$ the \textbf{boundary cohomology} is
\[
  H^{\bullet}(\partial;\VR)\;:=\;H^{\bullet}(\Gbdry;\VR),
\]
the group cohomology of $\Gbdry$ with coefficients in $\VR=\Sym^n(\Z^2)\otimes\cring$ and
$\Gbdry$ acting through $\rhon$ (Definition~\ref{def:rhon}). Owner decision OP5,
2026-09-15; it agrees with \S\ref{sec:six}'s own definition of $X_2$ and $X_5$, from
which it is transported unchanged. It is the half of erratum~\ref{err:e12} that was
already right; the other half is Proposition~\ref{prop:circlemodel}.
\end{definition}

\begin{decisionbox}
The boundary group is $\Gbdry=\langle T\rangle\times\langle-I\rangle$, \emph{not}
$\langle T\rangle$ (OP5, amended and decided 2026-09-15; the verdict gave no
recommendation and the owner decided with its argument section, \mbox{[V]~\S5.5}). The
one-dimensional ``circle model'' on $\Kbdry$ is the \emph{specialisation} of this
datum at $2\in\cring^{\times}$ (Proposition~\ref{prop:circlemodel}), exactly as the
$7$-vertex complex is the specialisation of the orbifold model at $6\in\cring^{\times}$
(\S\ref{sec:orbifoldmodel}). Under \ref{H3} with $\ell\ge5$ the two coincide and nothing
in Parts~III--V is affected; at $\ell=2$ they part company, and substantively ---
\S\ref{sec:ell2}.
\end{decisionbox}

\begin{proposition}[The cohomology of $\Gbdry$, in every degree]\label{prop:bdrycohom}
\PROVED{} Let $\cring$ be a commutative ring and $n\ge0$. Write $\VR^{T}:=\ker(\rhon(T)-\id)$
and $(\VR)_{T}:=\VR/\im(\rhon(T)-\id)$ (notation C12). Then for every $q\ge0$ there is a
short exact sequence of $\cring$-modules
\begin{equation}\label{eq:bdrylhs}
  0\longrightarrow\bigl(H^{q-1}(\langle-I\rangle;\VR)\bigr)_{T}
   \longrightarrow H^{q}(\Gbdry;\VR)
   \longrightarrow\bigl(H^{q}(\langle-I\rangle;\VR)\bigr)^{T}\longrightarrow0
\end{equation}
(with the left term read as $0$ for $q=0$), natural in $\VR$. If moreover
$\rhon(-I)=\id$ on $\VR$ --- equivalently $n$ is even, or $2=0$ in $\cring$ --- then
$H^{q}(\langle-I\rangle;\VR)=\VR$ for $q=0$, $=\VR[2]$ for $q$ odd and $=\VR/2\VR$ for
$q\ge2$ even.
\end{proposition}

\begin{proof}
\emph{Two elementary resolutions, used throughout this chapter.} (i) For
$\langle T\rangle\cong\Z$,
\[
  0\longrightarrow\Z[\langle T\rangle]\xrightarrow{\ T-1\ }\Z[\langle T\rangle]
   \xrightarrow{\ \aug\ }\Z\longrightarrow0
\]
is a free resolution of length $1$; hence for every $\Z[\langle T\rangle]$-module $M$,
$H^{0}(\langle T\rangle;M)=M^{T}$, $H^{1}(\langle T\rangle;M)=M_{T}$ and
$H^{s}(\langle T\rangle;M)=0$ for $s\ge2$, so $\langle T\rangle$ has cohomological
dimension $\le1$. \PROVED{} (ii) For $\langle-I\rangle\cong\Z/2$ with generator $t$, the
periodic free resolution
$\cdots\to\Z[\Z/2]\xrightarrow{\,1+t\,}\Z[\Z/2]\xrightarrow{\,t-1\,}\Z[\Z/2]
\xrightarrow{\ \aug\ }\Z\to0$ gives, for $M$ with \emph{trivial} $t$-action, the complex
$M\xrightarrow{0}M\xrightarrow{2}M\xrightarrow{0}M\xrightarrow{2}\cdots$ and hence
$H^{0}=M$, $H^{q}=M[2]$ for $q$ odd, $H^{q}=M/2M$ for $q\ge2$ even. \PROVED{}

Now $\Gbdry=\langle T\rangle\times\langle-I\rangle$ by \eqref{eq:bdrysplit}, so the
Lyndon--Hochschild--Serre spectral sequence of $1\to\langle-I\rangle\to\Gbdry\to
\langle T\rangle\to1$ has
$E_2^{s,q}=H^{s}\bigl(\langle T\rangle;H^{q}(\langle-I\rangle;\VR)\bigr)$
\STD{\cite[Ch.~VII.6]{brown}}, and the conjugation action of $\langle T\rangle$ on
$\langle-I\rangle$ is trivial because $-I$ is central and the extension splits, so the
coefficient module is $H^{q}(\langle-I\rangle;\VR)$ with the $T$-action induced from
$\VR$. By (i), $E_2^{s,q}=0$ for $s\ge2$; the spectral sequence therefore degenerates at
$E_2$ and its two-column filtration in total degree $q$ is \eqref{eq:bdrylhs}, the outer
terms being $E_2^{1,q-1}$ and $E_2^{0,q}$ read off by (i). The last sentence of the
statement is (ii) applied to $M=\VR$, whose $t$-action is trivial exactly when
$\rhon(-I)=\id$.
\end{proof}

\begin{proposition}[What the circle model computes, and when it is the boundary cohomology]\label{prop:circlemodel}
\PROVED{} (Appendix~\ref{app:errata}, entry~\ref{err:e12}: the specialisation hypothesis
$2\in\cring^{\times}$ is here part of the statement, not of a preamble.) Assume \ref{H4} and \ref{B1}--\ref{B3}, let $\cring$ be a commutative ring and
$n\ge0$, and let $C^{\bullet}_{\partial}:=\Hom_{\Z[\Ga]}\bigl(C_{\bullet}(\Xbdry),\VR\bigr)$ be
the equivariant cochain complex of the boundary. For the model of
Definition~\ref{def:model7} let $\Kbdry=v_0v_1v_2v_3$ be the boundary path and let
$\Rel_{\partial}$ be the relation set on $\Kbdry$ consisting of the \emph{cusp relation}
$(v_3,v_0,T)$, i.e.\ $f(v_3)=\rhon(T)f(v_0)$, \textbf{together with the $-I$
self-relations} $(\sigma,\sigma,-I)$, i.e.\ $f(\sigma)=\rhon(-I)f(\sigma)$, on every
simplex $\sigma$ of $\Kbdry$ --- the boundary half of \eqref{eq:negI}; the
\emph{circle model} is $\Ceq(\Kbdry,\Rel_{\partial};L)$. (Until v1.4.0 this proposition
called it ``the circle model with its single relation $f(v_3)=\rhon(T)f(v_0)$'', which
describes $\Rel_{\partial}$ only where the self-relations are vacuous --- clause~(4);
open point (3) of the v1.3.3 changelog.) Then:
\begin{enumerate}
\item \textbf{(Shapiro, unconditionally.)}
  $C^{\bullet}_{\partial}\cong\Hom_{\Z[\Gbdry]}\bigl(C_{\bullet}(\Xbdry_{\infty}),\VR\bigr)
   =\Hom_{\Z[\langle T\rangle]}\bigl(C_{\bullet}(\Xbdry_{\infty}),\VR^{-I}\bigr)$,
  where $\VR^{-I}:=\ker(\rhon(-I)-\id)$, by restriction of a cochain to the
  $\infty$-component.
\item \textbf{(What it computes.)} If $n$ is even, or if $2=0$ in $\cring$, then
  $\VR^{-I}=\VR$ and
  $H^{q}(C^{\bullet}_{\partial})=H^{q}(\langle T\rangle;\VR)$ for every $q$; explicitly
  $H^0=\VR^{T}$, $H^1=(\VR)_{T}$ and $H^{q}=0$ for $q\ge2$. If $n$ is odd and
  $2\in\cring^{\times}$ then $\VR^{-I}=0$ and $C^{\bullet}_{\partial}=0$.
\item \textbf{(When it is $H^{\bullet}(\partial;\VR)$.)} Let $n$ be even. If
  $2\in\cring^{\times}$ the inflation map
  $H^{q}(C^{\bullet}_{\partial})=H^{q}(\langle T\rangle;\VR)\to H^{q}(\Gbdry;\VR)$,
  along $\Gbdry\twoheadrightarrow\langle T\rangle$ and legitimate because $-I$ acts
  trivially on $\VR$, is an isomorphism for every $q$. If $\cring=\Z/\ell^{e}$ with $\ell=2$ and $\VR\neq0$ it is
  \emph{not}: $H^{2}(C^{\bullet}_{\partial})=0$ while $H^{2}(\Gbdry;\VR)\neq0$.
  \textbf{So the circle model computes $H^{\bullet}(\partial;\VR)$ exactly when
  $2\in\cring^{\times}$}, and at $\ell=2$ it computes a different, bounded object.
\item \textbf{(The model computes $C^{\bullet}_{\partial}$ --- for every $n$ and every
  $\cring$.)} For the model of Definition~\ref{def:model7},
  $\Ceq(\Kbdry,\Rel_{\partial};L)=C^{\bullet}_{\partial}$ as cochain complexes, by
  restriction of a cochain to the simplices of $\Kbdry$, with inverse
  $f\mapsto\bigl(\gamma\cdot\lift\sigma\mapsto\rhon(\gamma)f(\sigma)\bigr)$. Consequently
  the circle model is the single-relation model --- the cusp relation alone --- with
  coefficients $\VR^{-I}$ in place of $\VR$; the self-relations are vacuous, and the
  identification with the single-relation model on $\VR$ holds literally, exactly when
  $\rhon(-I)=\id$, i.e.\ for $n$ even or $2=0$ in $\cring$; for $n$ odd and
  $2\in\cring^{\times}$ the circle model is the zero complex
  (Remark~\ref{rmk:bdryoddn}). Without the self-relations, the single-relation model on
  $\VR$ is $\Hom_{\Z[\langle T\rangle]}(C_{\bullet}(\Xbdry_{\infty}),\VR)$, which is
  \emph{not} $C^{\bullet}_{\partial}$ for $n$ odd, $2\in\cring^{\times}$.
\end{enumerate}
\end{proposition}

\begin{proof}
(1) By \ref{B1} and \ref{B2}, $C_{\bullet}(\Xbdry)=\bigoplus_{c}C_{\bullet}(\Xbdry_c)$ with
$\Ga$ permuting the summands transitively with stabiliser $\Gbdry$, i.e.\
$C_{\bullet}(\Xbdry)\cong\Z[\Ga]\otimes_{\Z[\Gbdry]}C_{\bullet}(\Xbdry_{\infty})$ as
complexes of $\Z[\Ga]$-modules; the tensor--hom adjunction
$\Hom_{\Z[\Ga]}(\Z[\Ga]\otimes_{\Z[\Gbdry]}P,\VR)\cong\Hom_{\Z[\Gbdry]}(P,\VR)$,
$\Phi\mapsto\Phi|_{P}$, is the first isomorphism
\STD{\cite[Ch.~III.6]{brown}, where this adjunction is the content of Shapiro's lemma}. For the second: $-I$ acts
trivially on $\Xbulk$ (\ref{H4}(iv)), so a $\Gbdry$-equivariant $\varphi$ satisfies
$\varphi(x)=\varphi((-I)x)=\rhon(-I)\varphi(x)$, i.e.\ it takes values in $\VR^{-I}$, and
conversely.

(2) $\rhon(-I)=(-1)^{n}\id$ by Proposition~\ref{prop:symconv}(2) applied to $-I$, so
$\VR^{-I}=\VR$ when $n$ is even and when $2=0$ in $\cring$, and $\VR^{-I}=\VR[2]=0$ when
$n$ is odd and $2\in\cring^{\times}$. By \ref{B2}, $\Xbdry_{\infty}$ is a contractible
$1$-dimensional complex on which $\langle T\rangle$ acts freely: by \ref{B3} the
$\Ga$-stabiliser of every simplex of $\Xbdry$ is $\{\pm I\}$, and $T^{k}\notin\{\pm I\}$
for $k\neq0$. Hence the augmented complex $C_{\bullet}(\Xbdry_{\infty})\to\Z$ is a free
resolution of $\Z$ over $\Z[\langle T\rangle]$ --- free by the free action, exact by
contractibility --- and $\Hom_{\Z[\langle T\rangle]}(C_{\bullet}(\Xbdry_{\infty}),\VR)$
computes $H^{\bullet}(\langle T\rangle;\VR)$ by definition; the displayed values are
step (i) of the proof of Proposition~\ref{prop:bdrycohom}, and they vanish above degree
$1$ because $\Xbdry_{\infty}$ is $1$-dimensional.

(3) With $n$ even, $-I$ acts trivially on $\VR$. If $2\in\cring^{\times}$ then
$H^{q}(\langle-I\rangle;\VR)=0$ for $q\ge1$ (Proposition~\ref{prop:bdrycohom}, last
sentence: $\VR[2]=0$ and $\VR/2\VR=0$), so the $E_2$-page of the spectral sequence used
there is concentrated in the row $q=0$, where it is
$H^{s}(\langle T\rangle;\VR)$; a spectral sequence concentrated in one row degenerates
with edge map the inflation, so inflation
$H^{s}(\langle T\rangle;\VR)\to H^{s}(\Gbdry;\VR)$ is an isomorphism for every $s$, and
$H^{s}(\langle T\rangle;\VR)=H^{s}(C^{\bullet}_{\partial})$ by (2). If $\cring=\Z/\ell^{e}$, $\ell=2$, then $\VR$ is free of rank $n+1>0$ over
$\cring$, so $\VR/2\VR\neq0$; $\rhon(T)$ is unipotent (Definition~\ref{def:rhon}: $T\cdot
x=x$, $T\cdot y=x+y$, so $\rhon(T)u_j=u_j+\sum_{t<j}\binom{j}{t}u_t$ is unitriangular),
hence so is its reduction on $\VR/2\VR\neq0$, hence
$\bigl(\VR/2\VR\bigr)^{T}\neq0$; by \eqref{eq:bdrylhs} at $q=2$,
$H^{2}(\Gbdry;\VR)$ surjects onto it and is non-zero, while
$H^{2}(C^{\bullet}_{\partial})=0$ by (2).

(4) $\Xbdry=\Ga\cdot\Kbdry$: by Remark~\ref{rmk:bdrymodel},
$\Xbdry_{\infty}=\Gbdry\cdot\Kbdry$ and $\Xbdry=\Ga\cdot\Xbdry_{\infty}$ (\ref{B1}).
This is \ref{FD2} for the pair $(\Kbdry,\Xbdry)$. The coincidence groupoid
$\mathcal{C}(\Kbdry,\Xbdry,\Ga)$ of Definition~\ref{def:relationdatum} --- the triples
$(\sigma,\sigma',\gamma)$ with $\gamma\cdot\lift{\sigma'}=\lift\sigma$ and
$\sigma,\sigma'$ simplices of $\Kbdry$ --- is generated by $\Rel_{\partial}$: both
simplices lie in $\Xbdry_{\infty}$, whose setwise stabiliser is $\Gbdry$ (\ref{B1}), so
$\gamma\in\Gbdry=\{\pm T^{k}\}$; by Remark~\ref{rmk:bdrymodel} the $\Gbdry$-orbits of
simplices of $\Xbdry_{\infty}$ meet $\Kbdry$ in $\{v_0,v_3\}$, $\{v_1\}$, $\{v_2\}$ and
in each of the three edges separately, and by \ref{B3} every stabiliser is $\{\pm I\}$;
hence a triple is $(\sigma,\sigma,\pm I)$, $(v_3,v_0,\pm T)$ or $(v_0,v_3,\pm T^{-1})$,
each a composite of the cusp relation, its inverse and the self-relations. This is
\ref{FD4} for $(\Kbdry,\Rel_{\partial})$. The proof of
Theorem~\ref{thm:explicit_complex_equiv}(1) uses exactly \ref{FD2} and \ref{FD4} --- it
says so at its well-definedness step --- and therefore applies verbatim to
$(\Kbdry,\Rel_{\partial})$ and $\Xbdry$: $\Psi\colon\phi\mapsto\phi|_{\Kbdry}$ is an
isomorphism of cochain complexes $C^{\bullet}_{\partial}\to\Ceq(\Kbdry,\Rel_{\partial};L)$
with the stated inverse. (Assumption~\ref{ass:fund_domain} as a whole is not invoked:
its standing hypothesis that $\abs{\Ku}$ be connected fails for $\Xbdry$, and neither it
nor \ref{FD1}, \ref{FD3} is needed; $\Rel_{\partial}$ is closed under faces,
Assumption~\ref{ass:consistent_face_id}, the faces of a self-relation being
self-relations and the cusp relation a vertex relation.) The self-relations read
$(\id-\rhon(-I))f(\sigma)=0$, i.e.\ $f(\sigma)\in\VR^{-I}$; keeping the cusp relation and
restricting the values to $\VR^{-I}$ is the single-relation model with coefficients
$\VR^{-I}$, which is clause~(1)'s second isomorphism read on $\Kbdry$. The case split is
$\rhon(-I)=(-1)^{n}\id$ (proof of (2)): $\VR^{-I}=\VR$ for $n$ even or $2=0$ in
$\cring$, and $\VR^{-I}=\VR[2]=0$ for $n$ odd and $2\in\cring^{\times}$. For the last
sentence: without the self-relations the well-definedness step of
Theorem~\ref{thm:explicit_complex_equiv}(1) fails at the stabiliser $\{\pm I\}$ exactly
as Remark~\ref{rmk:negI} describes, and for $n$ odd, $2\in\cring^{\times}$, the
single-relation model on $\VR$ has non-zero cochains (any value on $v_1$, say) while
$C^{\bullet}_{\partial}=0$.
\end{proof}

\begin{remark}[Odd $n$, on the boundary]\label{rmk:bdryoddn}
\PROVED{} For $n$ odd and $2\in\cring^{\times}$ the boundary complex is the zero complex
and $H^{\bullet}(\partial;\VR)=0$, for the same reason that $\Heq{\bullet}(K,\Rel;L)=0$
there: $-I$ is in every stabiliser and acts by $(-1)^{n}=-1$
(Remark~\ref{rmk:negI}). The two vanishings are one phenomenon, and both fail at
$\ell=2$, where $(-1)^n=1$. Every statement of this chapter that assumes $n$ even says so.
Since v1.4.0 the model side is literal: Proposition~\ref{prop:circlemodel}(4) identifies
the circle model $\Ceq(\Kbdry,\Rel_{\partial};L)$, with the $-I$ self-relations in
$\Rel_{\partial}$, with $C^{\bullet}_{\partial}$ for every $n$, so the zero complex at odd
$n$, $2\in\cring^{\times}$, is what the implementation computes and not a statement about
a different object (code map of this chapter).
\end{remark}

\subsection{The upper-triangular set $B_N$ and its cosets}\label{sec:bdrycosets}

\begin{proposition}[The cosets of $B_N$, and which decomposition is the right one]\label{prop:bdrycosets}
\PROVED{} Let $p$ be a prime, $\nu\ge1$, $N=p^{\nu}$, $g=\smat{1}{0}{0}{N}$ and
$m=N+p^{\nu-1}$ as in Proposition~\ref{prop:cosets}, and put
\[
  B_N:=\{\delta\in\Ga g\Ga:\ \delta\ \text{upper triangular}\},
\]
the elements of the double coset fixing the cusp $\infty$. Then:
\begin{enumerate}
\item \textbf{(Left cosets.)} $B_N=\bigsqcup\,\delta\Gbdry$ over exactly $m$ cosets. For
  $\nu=1$ they are represented by
  \[
    \delta^{\partial}_{\infty}:=\mathrm{diag}(1,p)=g,
    \qquad
    \delta^{\partial}_{t}:=\smat{p}{t}{0}{1}\ \ (0\le t<p),
  \]
  and $T\delta^{\partial}_{\infty}=\delta^{\partial}_{\infty}T^{p}$,
  $T\delta^{\partial}_{t}=\delta^{\partial}_{t+1}$ for $t<p-1$, and
  $T\delta^{\partial}_{p-1}=\delta^{\partial}_{0}T$. \LEAN{prop:bdrycosets}
\item \textbf{(Bijection with the bulk cosets.)} $\delta\Gbdry\mapsto\delta\Ga$ is a
  bijection $B_N/\Gbdry\to\Ga g\Ga/\Ga$. Hence for each bulk coset element
  $\delta_j=s_j^{-1}g$ of Definition~\ref{def:cosetreps} the set
  $\{\gamma\in\Ga:\delta_j\gamma\in B_N\}$ is a single left coset $\gamma_j\Gbdry$,
  characterised by $\gamma_j\cdot\infty=\delta_j^{-1}\cdot\infty$, and
  $\{\delta_j\gamma_j\}_{j=1}^{m}$ is a set of representatives of $B_N/\Gbdry$.
  \LEAN{prop:bdrycosets-cusp} (since v1.18.0 including the characterisation $\gamma_j\cdot\infty=\delta_j^{-1}\cdot\infty$, as a ratio pair)
\item \textbf{(Two double cosets, not one.)} For $\nu=1$,
  $B_p=\Gbdry\,\delta^{\partial}_{\infty}\Gbdry\ \sqcup\ \Gbdry\,\delta^{\partial}_{0}\Gbdry$,
  with $\Gbdry\cap\delta^{\partial}_{\infty}\Gbdry(\delta^{\partial}_{\infty})^{-1}=\Gbdry$
  (one left coset) and
  $\Gbdry\cap\delta^{\partial}_{0}\Gbdry(\delta^{\partial}_{0})^{-1}=\langle T^{p}\rangle\times\langle-I\rangle$
  of index $p$ ($p$ left cosets). \LEAN{prop:bdrycosets}
\item \textbf{(The right cosets are the wrong ones.)} The upper-triangular right-coset
  representatives $\beta_b=\smat{1}{b}{0}{p}$ $(0\le b<p)$ and
  $\beta_{\infty}=\mathrm{diag}(p,1)$ of Proposition~\ref{prop:cosets}(1) represent
  $\Gbdry\backslash B_p$, not $B_p/\Gbdry$. They are \textbf{not} to be used: they give
  the $\det$-twisted operator, $\chi_{b+1}$ in place of $\chi_b$
  (Proposition~\ref{prop:bdrychar}). \LEAN{prop:bdrycosets}
\end{enumerate}
\end{proposition}

\begin{proof}
An upper-triangular $\delta=\smat{a}{b}{0}{d}$ lies in $\Ga g\Ga$ iff $ad=N$ and
$\gcd(a,b,d)=1$ (Proposition~\ref{prop:cosets}(1)), so $a=\pm p^{u}$, $d=\pm p^{\nu-u}$
with equal signs, $0\le u\le\nu$. Since
$\Gbdry=\{\pm T^{k}\}$ and $\smat{a}{b}{0}{d}T^{k}=\smat{a}{b+ka}{0}{d}$, the left coset
$\delta\Gbdry$ is determined by $(u,\,b\bmod a)$ up to the common sign, which
$-I\in\Gbdry$ absorbs. Counting admissible $b$: for $u=0$ one coset; for $0<u<\nu$ both
$a$ and $d$ are divisible by $p$ so primitivity forces $p\nmid b$, giving
$p^{u}-p^{u-1}$ cosets; for $u=\nu$, $d=\pm1$ and every $b\bmod N$ occurs, giving $N$.
The total is $1+\sum_{u=1}^{\nu-1}(p^{u}-p^{u-1})+N=p^{\nu-1}+N=m$. At $\nu=1$ the
cosets are $u=0$, represented by $\smat{1}{0}{0}{p}=\delta^{\partial}_{\infty}$, and
$u=1$, represented by $\smat{p}{t}{0}{1}=\delta^{\partial}_{t}$, $0\le t<p$; the three
displayed identities are the matrix products
$\smat{1}{1}{0}{1}\smat{1}{0}{0}{p}=\smat{1}{p}{0}{p}=\smat{1}{0}{0}{p}\smat{1}{p}{0}{1}$,
$\smat{1}{1}{0}{1}\smat{p}{t}{0}{1}=\smat{p}{t+1}{0}{1}$ and
$\smat{1}{1}{0}{1}\smat{p}{p-1}{0}{1}=\smat{p}{p}{0}{1}=\smat{p}{0}{0}{1}\smat{1}{1}{0}{1}$.

(2) The map is well defined since $\Gbdry\subseteq\Ga$. It is injective: if
$\delta'=\delta\gamma$ with $\delta,\delta'\in B_N$ and $\gamma\in\Ga$, then
$\gamma=\delta^{-1}\delta'$ is upper triangular
($\delta^{-1}=\frac1N\smat{d}{-b}{0}{a}$, a product of upper-triangular matrices), hence
lies in $\Gbdry$. Both sides have $m$ elements --- the left by (1), the right by
Proposition~\ref{prop:cosets}(2) --- so it is a bijection. For the last sentence,
$\delta_j\gamma$ is upper triangular iff it fixes $\infty$ iff
$\gamma\cdot\infty=\delta_j^{-1}\cdot\infty$; the set of such $\gamma\in\Ga$ is a left
coset of $\mathrm{Stab}_\Ga(\delta_j^{-1}\infty)$, which for any choice $\gamma_j$ is
$\gamma_j\Gbdry$ (conjugate stabilisers), and $\delta_j\gamma_j\Gbdry$ is then one
well-defined element of $B_N/\Gbdry$. These $m$ elements are distinct because the
$\delta_j\Ga$ are, and the bijection just proved carries one family to the other.

(3) $\smat{1}{1}{0}{1}$ conjugated by $\delta^{\partial}_{\infty}$ is
$\delta^{\partial}_{\infty}T(\delta^{\partial}_{\infty})^{-1}=\smat{1}{1/p}{0}{1}$, so
$\Gbdry\cap\delta^{\partial}_{\infty}\Gbdry(\delta^{\partial}_{\infty})^{-1}=\Gbdry$ and
the double coset $\Gbdry\delta^{\partial}_{\infty}\Gbdry$ is a single left coset; it is
the $u=0$ coset. Likewise $\delta^{\partial}_{0}T(\delta^{\partial}_{0})^{-1}=T^{p}$, so
the intersection is $\langle T^{p}\rangle\times\langle-I\rangle$ of index $p$, and by
(1)'s identities the $p$ cosets $\delta^{\partial}_{t}\Gbdry$ are permuted transitively
by $\Gbdry$; they are the $u=1$ cosets. The two families exhaust $B_p$.

(4) $\Gbdry\beta_b=\{\pm\smat{1}{b+kp}{0}{p}\}$ exhausts the $u=0$ family as $b$ runs
over $0,\dots,p-1$, and $\Gbdry\beta_{\infty}=\{\pm\smat{p}{k}{0}{1}\}$ is the whole
$u=1$ family; so these $p+1$ sets are the right cosets $\Gbdry\backslash B_p$. They
split the two families in the opposite way to (1), which is exactly why the two sums
differ; Proposition~\ref{prop:bdrychar}(2) computes by how much.
\end{proof}

\begin{remark}[General $\nu$: what this chapter does not state]\label{rmk:bdrynugeneral}
Clauses (1) and (2) of Proposition~\ref{prop:bdrycosets} hold for every $\nu$; clauses
(3) and (4), Proposition~\ref{prop:bdrychar}(3) with \eqref{eq:bdrygroupform}, and the
$\Gbdry$-level datum of the decision box after Definition~\ref{def:bdryop} are $\nu=1$
statements, and \textbf{no double-coset decomposition $\Gbdry\backslash B_N/\Gbdry$ for
$\nu\ge2$ is stated or proved in this document}. That is why
Obligation~\ref{obl:taubdry}(v) is degree $0$ alone at $\nu\ge2$ (owner confirmation 2
of the slot-2 plan, 2026-09-16: narrow now, add when a phase needs $T^{\partial}_{p^{\nu}}$
on $H^{1}$; hecke-extensions slot 1, WP-4b, finding 11). The shape a later addition would
have to state and prove is recorded so that it is not re-derived from scratch, and is
\emph{not asserted here}: writing
$\delta_{u,b}:=\begin{psmallmatrix}p^{u}&b\\0&p^{\nu-u}\end{psmallmatrix}$ for the
representatives of (1), the computation in the proof of (1) gives
$T^{k}\delta_{u,b}=\delta_{u,\,b+kp^{\nu-u}}$ and $\delta_{u,b}T^{k}=\delta_{u,\,b+kp^{u}}$,
which suggests that the double cosets are indexed by $u$ and $b\bmod p^{\min(u,\nu-u)}$
(with $p\nmid b$ for $0<u<\nu$), that the double coset of $\delta_{u,b}$ is the union of
$p^{\max(0,\,2u-\nu)}$ left cosets, and that the transporters are the corresponding powers
of $T$. The general-$\nu$ analogue of \eqref{eq:bdrygroupform} on $H^{1}$ and a
$\Gbdry$-level twisted datum for route (v.c) at $\nu\ge2$ would follow from it and are
equally not stated. At $\nu=1$ the shape specialises to clause (3). (v1.12.0-pre.3: an $H^{1}$
formula at every $\nu$ \emph{is} now stated and proved, Proposition~\ref{prop:bdrytranscription}(2),
\eqref{eq:bdrygroupformnu}; it is a sum over the left cosets of clause (1), which hold at every
$\nu$, and uses none of the double-coset shape above, which stays unasserted, as does the
$\Gbdry$-level datum at $\nu\ge2$. Obligation~\ref{obl:taubdry}(v) is not extended by it.)
\end{remark}

\subsection{The boundary twisted chain datum $\taub$}\label{sec:taubdry}

\begin{definition}[$g$-twisted chain datum on $\Xbdry$]\label{def:taubdry}
Let $\cring$ be a commutative ring, $g\in\De$ and $\Ga''=\Ga\cap g\Ga g^{-1}$. A
\emph{$g$-twisted chain map on the boundary over $\cring$} is a degree-preserving
$\cring$-linear map
$\taub\colon C_{\bullet}(\Xbdry)\otimes_\Z\cring\to C_{\bullet}(\Xbdry)\otimes_\Z\cring$
such that
\begin{enumerate}
\item[(T1$\partial$)] $\taub(h\cdot x)=(g^{-1}hg)\cdot\taub(x)$ for all $h\in\Ga''$ and
  all $x$;
\item[(T2$\partial$)] $\partial\circ\taub=\taub\circ\partial$;
\item[(T3$\partial$)] the induced map on
  $H_0(\Xbdry)\otimes\cring=\cring[\pi_0(\Xbdry)]$ is $\pi_{\ast}$ for a map of sets
  $\pi\colon\pi_0(\Xbdry)\to\pi_0(\Xbdry)$: each component is carried into a single
  component, with total coefficient $1$.
\end{enumerate}
A \emph{boundary $g$-twisted chain datum} $\taub$ is such a map together with the finite
table of its values on a set of $\Ga''$-orbit representatives of the simplices of
$\Xbdry$, from which (T1$\partial$) determines it everywhere.
\end{definition}

\noindent
(T3$\partial$) is the correct analogue of (T3) of Definition~\ref{def:tau}: there
$H_0(\Ku)=\Z$ and ``$\aug\tau=\aug$'' says the induced map is the identity; here
$H_0(\Xbdry)$ is free on the cusps and the condition says the induced map is a
\emph{permutation-like} map of basis elements, not merely degree-preserving. Both say
that $\tau$ does not spread a point over several components. On each component
separately, (T3$\partial$) is the old (T3).

\begin{lemma}[(T1$\partial$) and (T3$\partial$) force the component map: it is $c\mapsto g^{-1}c$]\label{lem:bdrycomponent}
\PROVED{} Assume \ref{H4} and \ref{B1}--\ref{B3}, and let $g=\smat{1}{0}{0}{N}$,
$\Ga''=\Ga_0(N)$ as in Proposition~\ref{prop:cosets}. Identify $\pi_0(\Xbdry)$ with
$\mathbb{P}^1(\Q)$ by \ref{B2}, and let $g^{-1}$ act on $\mathbb{P}^1(\Q)$ by M\"obius
transformations ($\det g=N>0$, so this is defined and
$g^{-1}\mathbb{P}^1(\Q)=\mathbb{P}^1(\Q)$).
Then every $\taub$ satisfying (T1$\partial$) and (T3$\partial$) has
\[
  \pi(c)\;=\;g^{-1}\cdot c\qquad\text{for every cusp }c .
\]
In particular $\pi(\infty)=\infty$, $\pi$ is a bijection, and
$\pi$ is determined by $g$ alone: \textbf{no choice is made in it}, and it is the same
for every admissible $\taub$.
\end{lemma}

\begin{proof}
Write $\Ga_c:=\mathrm{Stab}_{\Ga}(c)$. $\Ga$ acts transitively on
$\mathbb{P}^1(\Q)$ (a primitive column $(a,c)^{\mathsf T}$ completes to a matrix of
$\Ga$), so $\Ga_c=\gamma\Gbdry\gamma^{-1}$ for any $\gamma$ with $\gamma\infty=c$; hence
$\Ga_c=\langle P_c\rangle\times\langle-I\rangle$ with $P_c=\gamma T\gamma^{-1}$
parabolic, by \eqref{eq:bdrysplit}, and $\Ga_c$ is abelian. $\Ga''=\Ga_0(N)$ has finite
index $m$ in $\Ga$ (Proposition~\ref{prop:cosets}(3)), so $\Ga''\cap\Ga_c$ has finite index in
$\Ga_c$ and contains $P_c^{\,k}$ for some $k\ge1$. By Proposition~\ref{prop:cosets}(3),
$\conj_g^{-1}\colon\Ga''\to\Ga'$ is an isomorphism with integral values, so
$P:=g^{-1}P_c^{\,k}g$ lies in $\Ga$; it has the same trace $\pm2$ as $P_c^{\,k}$ and is
not $\pm I$, hence is a non-trivial parabolic element of $\Ga$ and fixes exactly one
point of $\mathbb{P}^1(\Q)$. That point is $g^{-1}c$, because $P_c^{\,k}$ fixes $c$ and
$P(g^{-1}c)=g^{-1}P_c^{\,k}c=g^{-1}c$.

Now let $\taub$ satisfy (T1$\partial$) and (T3$\partial$), with component map $\pi$.
For $h\in\Ga''\cap\Ga_c$ and $x$ supported in $\Xbdry_c$, the chain $h\cdot x$ is again
supported in $\Xbdry_c$, so comparing the component of each side of (T1$\partial$) ---
which (T3$\partial$) says is well defined --- gives
$\pi(c)=(g^{-1}hg)\cdot\pi(c)$. Thus $\pi(c)$ is fixed by $g^{-1}(\Ga''\cap\Ga_c)g$, in
particular by $P$, whose unique fixed cusp is $g^{-1}c$. Hence $\pi(c)=g^{-1}c$, and
$\pi$ is a bijection because $g^{-1}$ acts invertibly on $\mathbb{P}^1(\Q)$. For
$g=\smat{1}{0}{0}{N}$, $g^{-1}\infty=\infty$.
\end{proof}

\begin{lemma}[Existence, homotopy uniqueness and integrality of $\taub$ --- over $\Z$, for every $\ell$]\label{lem:taubdryexists}
\PROVED{} (Appendix~\ref{app:errata}, entry~\ref{err:e4}: part (3) is proved by the
explicit argument below and \textbf{never} by the comparison theorem, whose
relative-projectivity hypothesis is exactly what this construction exists to avoid.)
Assume \ref{H4} and \ref{B1}--\ref{B3}, and let $g=\smat{1}{0}{0}{N}$,
$\Ga''=\Ga_0(N)$. Then:
\begin{enumerate}
\item \textbf{(Existence over $\Z$, no hypothesis on $\ell$.)} A boundary $g$-twisted
  chain datum over $\cring=\Z$ exists, and may be built by the following terminating
  procedure (rule~\ref{A3}). For each $\Ga''$-orbit representative vertex $v$ of
  $\Xbdry$, lying in the component $\Xbdry_c$, choose any vertex $w$ of the forced
  target component $\Xbdry_{g^{-1}c}$ (Lemma~\ref{lem:bdrycomponent}; the canonical
  choice, the one the implementation makes, is step~2 of Algorithm~\ref{alg:taubdry})
  and set $\taub(v):=w$; for each $\Ga''$-orbit representative edge $e$
  of $\Xbdry$, let $\taub(e)$ be the \emph{unique} $1$-chain of $\Xbdry_{g^{-1}c}$ with
  $\partial\taub(e)=\taub(\partial e)$; extend to all of $C_{\bullet}(\Xbdry)$ by
  (T1$\partial$). There is nothing to do in degrees $\ge2$.
\item \textbf{(Integrality, no hypothesis on $\ell$ and no linear solve.)} Every
  coefficient of the emitted table lies in $\{0,\pm1\}$, and no system of equations is
  solved anywhere in (1). Consequently the table \emph{is} a datum over $\Z$, and the
  datum over any $\cring$ is its base change; \textbf{the mod-$\ell$ fallback of
  Algorithm~\ref{alg:tau} step~3 cannot arise on the boundary, for any $\ell$, because
  there is no step~3: $\Xbdry$ has no $2$-simplices (\ref{B2}).}
\item \textbf{(Homotopy uniqueness over $\Z$, no hypothesis on $\ell$.)} If
  $\taub,\taub'$ both satisfy (T1$\partial$)--(T3$\partial$) over $\Z$ then there is an
  additive degree-$(+1)$ map $N_{\partial}\colon C_{\bullet}(\Xbdry)\to
  C_{\bullet+1}(\Xbdry)$ satisfying (T1$\partial$) with
  $\taub-\taub'=\partial N_{\partial}+N_{\partial}\partial$; moreover
  $N_{\partial}$ vanishes in degree $1$ and is unique in degree $0$.
\end{enumerate}
\end{lemma}

\begin{proof}
\emph{(1).} By Lemma~\ref{lem:bdrycomponent} the target component $g^{-1}c$ is forced, and
$\Xbdry_{g^{-1}c}\neq\emptyset$, so a $w$ exists. The (T1$\partial$)-extension
$\taub(h\cdot v):=(g^{-1}hg)\cdot\taub(v)$, $h\in\Ga''$, is well defined precisely when
$(g^{-1}hg)\cdot w=w$ for every $h$ in the $\Ga''$-stabiliser of $v$. \textbf{By
\ref{B3} that stabiliser is $\Ga''\cap\{\pm I\}=\{\pm I\}$}, $g^{-1}(-I)g=-I$, and $-I$
acts trivially on $\Xbulk$ (\ref{H4}(iv)); so the condition is vacuous. \emph{This is the
whole of the difference from the abstract route of Lemma~\ref{lem:tauexists}(2)}: there
the stabilisers have orders $4$ and $6$ and the averaging idempotent
$\abs{\Ga_\sigma}^{-1}\sum_{u}u$ is needed, which is why $6\in\cring^{\times}$ is required
there and is not required here. (Lemma~\ref{lem:tauexists}(1) avoids the idempotent in
the bulk as well, but by a different device --- choosing the vertex value inside the
fixed-point set of the twisted stabiliser --- and it buys existence only, never
uniqueness.)
For an edge $e$ of $\Xbdry_c$ with endpoints $\partial e=w''-w'$:
$\taub(\partial e)=\taub(w'')-\taub(w')$ is a $0$-chain of total coefficient $0$ in
$\Xbdry_{g^{-1}c}$, hence a boundary, because $\Xbdry_{g^{-1}c}$ is connected
(\ref{B2}); the $1$-chain with that boundary is unique because
$Z_1(\Xbdry_{g^{-1}c})=0$ (\ref{B2}), and it is the oriented path from $\taub(w')$ to
$\taub(w'')$. The same stabiliser argument makes the
(T1$\partial$)-extension of $\taub$ on edges well defined. (T2$\partial$) holds by
construction in degree $1$ and is vacuous elsewhere; (T3$\partial$) holds with
$\pi(c)=g^{-1}c$ by construction. The procedure terminates because $\Ga''$ has finitely
many orbits on the simplices of $\Xbdry$: $\Ga$ has finitely many by \ref{B1}, and
$[\Ga:\Ga'']=m<\infty$ by Proposition~\ref{prop:cosets}(3).

\emph{(2).} A vertex value is a single vertex with coefficient $1$; an edge value is an
oriented path, i.e.\ a sum of edges with coefficients $\pm1$. No other operation occurs.

\emph{(3).} Both data have the same component map $\pi$ (Lemma~\ref{lem:bdrycomponent}),
so for a $\Ga''$-orbit representative vertex $v$ of $\Xbdry_c$ the $0$-chain
$(\taub-\taub')(v)$ has total coefficient $0$ and is supported in $\Xbdry_{g^{-1}c}$; let
$N_{\partial}(v)$ be the unique $1$-chain there with
$\partial N_{\partial}(v)=(\taub-\taub')(v)$ (existence by connectedness, uniqueness by
$Z_1=0$; \ref{B2}), extended by (T1$\partial$), which is well defined by the stabiliser
argument of (1). Put $N_{\partial}:=0$ in degree $1$, which is forced since
$C_2(\Xbdry)=0$. It remains to check
$(\taub-\taub')(e)=N_{\partial}(\partial e)$ for every edge $e$. The difference
$z:=(\taub-\taub')(e)-N_{\partial}(\partial e)$ is a $1$-chain of $\Xbdry_{g^{-1}c}$ with
\[
  \partial z=(\taub-\taub')(\partial e)-\partial N_{\partial}(\partial e)
           =(\taub-\taub')(\partial e)-(\taub-\taub')(\partial e)=0,
\]
using (T2$\partial$) for both data; so $z\in Z_1(\Xbdry_{g^{-1}c})=0$. Uniqueness in degree
$0$ is the uniqueness of the fill just used.
\end{proof}

\begin{remark}[Integrality of $\taub$: the claim, its reason, and what it is not]\label{rmk:taubdryintegral}
Two distinct facts are in play and have been conflated before, so they are separated
here: conflating them is the class of over-claim Appendix~\ref{app:errata} records at
entry~\ref{err:e3}, a hypothesis that was dropped when a statement was carried into a
regime it was not proved in.
\begin{itemize}
\item \textbf{What is claimed.} $\taub$ exists over $\Z$, with coefficients in
  $\{0,\pm1\}$, for \emph{every} $\ell$ --- including $\ell\in\{2,3\}$ ---
  and is homotopy-unique over $\Z$ for every $\ell$
  (Lemma~\ref{lem:taubdryexists}). The two reasons are different and both are
  necessary: \ref{B3} (the stabilisers of boundary simplices are central and act
  trivially) is what makes the (T1$\partial$)-extension well defined \emph{without
  inverting a stabiliser order}, and \ref{B2} (each component is a contractible
  $1$-complex) is what makes the fills exist, be unique, and be paths.
\item \textbf{What is \emph{not} claimed.} That the \emph{uniqueness} half
  Lemma~\ref{lem:tauexists}(2) holds over $\Z$, or without $6\in\cring^{\times}$, for the
  \emph{bulk} $\tau$. It is not claimed, and nothing here repairs it: the bulk homotopy
  argument needs relative projectivity and therefore needs the elliptic stabiliser orders
  $4$ and $6$ to be units. (The \emph{existence} half is a different matter: it does hold
  over $\Z$ for every $\ell$, by a construction, Lemma~\ref{lem:tauexists}(1) --- but the
  device there is fixed points of the twisted stabilisers, not \ref{B3}, and it says
  nothing about uniqueness. See Remark~\ref{rmk:tauexistsscope}.) Nor is the abstract
  comparison theorem used above: Lemma~\ref{lem:taubdryexists} is a construction and an
  explicit uniqueness argument, not an application of \cite[Ch.~I.7.4]{brown}.
\item \textbf{Policy (owner decision OP2, 2026-09-15).} The mod-$\ell$ fill fallback is
  forbidden by default and \code{integral = false} is not the escape from a hard fill
  (Remark~\ref{rmk:tauintegral}). \textbf{On the boundary the policy is vacuous}: by
  Lemma~\ref{lem:taubdryexists}(2) there is no fill to fall back from, so a
  $\taub$-table always has solve regime $\Z$, and a record in which a $\taub$-table
  carries \code{integral = false} is reporting an implementation defect, not a
  mathematical regime. That is a checkable assertion and it is
  Obligation~\ref{obl:taubdry}(iv).
\end{itemize}
\end{remark}

\begin{algorithm}[TAU-BDRY-TABLE: the canonical $\taub$]\label{alg:taubdry}
Assume \ref{H4} and \ref{B1}--\ref{B3}, and let $g=\smat{1}{0}{0}{N}$, $\Ga''=\Ga_0(N)$.
\textbf{Input:} $p$, $\nu$, $N=p^{\nu}$; the boundary subcomplex $\Kbdry$ of the model
with the exact positions of its vertices (Definition~\ref{def:model7}) and its
$\infty$-line coordinates --- the vertices of $\Xbdry_{\infty}=\Gbdry\cdot\Kbdry$ are
$T^{k}v_i$ with $i$ in the first period, numbered by the \emph{position}
$P=k\cdot\mathrm{per}+i$ ($\mathrm{per}=3$ for the model, $v_3=Tv_0$ at position $3$),
with real parts $\mathrm{Re}(T^{k}v_i)=\mathrm{Re}\,v_i+k$; the representatives
$\{s_j\}_{j=1}^{m}$ of Definition~\ref{def:cosetreps}; and a cusp reduction
$\mathrm{red}\colon\mathbb{P}^{1}(\Q)\to\Ga$ with $\mathrm{red}(c)\cdot c=\infty$ (the
implementation's is the nearest-integer Euclid on the primitive column of (P2); by
property~(1) below the output does not depend on this choice). \textbf{Output:} for every
$j\in\{1,\dots,m\}$ and every simplex $\sigma$ of $\Kbdry$ a chain
$\taub(s_j\cdot\lift\sigma)$ in the stored-chain format \eqref{eq:taurecord}, over $\Z$,
with coefficients in $\{0,\pm1\}$. This is the procedure of
Lemma~\ref{lem:taubdryexists}(1) with its one free choice --- ``any vertex $w$ of the
target component'' --- fixed; the output is \textbf{the canonical $\taub$}, and every
statement of this document about ``the'' $\taub$-table refers to it (rule~\ref{A3}).
Hecke-extensions slot 1, WP-4a, finding 2; \file{Hecke/HeckeBoundaryDatum}.
\begin{enumerate}
\item[\textbf{0.}] \textbf{Slots, and the single (T1$\partial$) helper.} As in
  Algorithm~\ref{alg:tau} step~0. In ascending degree, and within a degree in the
  model's order of boundary simplices and $j=1,\dots,m$: a slot $(j,\sigma)$ is
  \emph{derived} whenever a stored slot $(j',\sigma')$ with
  $s_j\cdot\lift\sigma=\gamma\cdot\lift{\sigma'}$ for a transporter $\gamma=s_jt$,
  $\gamma=h\,s_{j'}$ with $h\in\Ga''$, exists --- then
  $\taub(s_j\sigma):=(g^{-1}hg)\cdot\taub(s_{j'}\sigma')$ through the one shared
  (T1$\partial$) helper, with the orientation sign of the transporter tracked --- and is
  \emph{constructed} by steps~1--3 otherwise. Step~4 uses the same helper. By
  property~(1) the order in which the slots are visited does not reach the output.
\item[\textbf{1.}] \textbf{Vertices: the transported real part.} For a constructed slot
  $(j,v)$: $z:=\bigl(\adjg(g)\,s_j\bigr)\cdot\mathrm{pos}(v)\in\Hyp$ with
  $\adjg(g)=\mathrm{diag}(N,1)$, projectively, as in Algorithm~\ref{alg:tau} step~1 (no
  rational number is formed); $c':=g^{-1}\cdot(s_j\infty)$, the forced target component
  (Lemma~\ref{lem:bdrycomponent}); $\gamma_{\mathrm{red}}:=\mathrm{red}(c')$; and
  $x:=\mathrm{Re}(\gamma_{\mathrm{red}}z)\in\Q$, the real part of the image transported
  to the $\infty$-component --- its vertical projection onto the horocycle.
\item[\textbf{2.}] \textbf{The vertex rule.} Let $P$ be the position of the line vertex
  nearest to $x$, i.e.\ with $\abs{x-\mathrm{Re}(T^{k}v_i)}$ minimal, decided exactly;
  \textbf{if two vertices are equidistant from $x$, take the one with the larger real
  part.} Set $\taub(s_j\cdot\lift v):=\gamma_{\mathrm{red}}^{-1}\cdot T^{k}\lift{v_i}$ for
  $P=k\cdot\mathrm{per}+i$, one term with coefficient $1$, so (T3$\partial$) holds by
  inspection.
\item[\textbf{3.}] \textbf{Edges.} For a constructed slot $(j,e)$ with
  $\partial e=v_b-v_a$: the two vertex values $\gamma_a\cdot\lift{w_a}$ and
  $\gamma_b\cdot\lift{w_b}$ are already stored (a derived face's value being the one the
  step-0 helper produced) and lie in one component $c'$ (step~1,
  Lemma~\ref{lem:bdrycomponent}); with
  $\gamma_{\mathrm{red}}:=\mathrm{red}(c')$ let $P_a,P_b$ be the positions of
  $\gamma_{\mathrm{red}}\gamma_a\lift{w_a}$ and $\gamma_{\mathrm{red}}\gamma_b\lift{w_b}$
  on the line. Set $\taub(s_j\cdot\lift e):=\gamma_{\mathrm{red}}^{-1}\cdot(\text{the
  oriented walk from }P_a\text{ to }P_b)$, the sum of the line edges between the two
  positions with coefficient $\pm1$ according to their orientation --- the unique
  $1$-chain of the line with boundary $\gamma_b\lift{w_b}-\gamma_a\lift{w_a}$
  (Lemma~\ref{lem:taubdryexists}(1)); empty if $P_a=P_b$. No linear solve occurs.
\item[\textbf{4.}] \textbf{Certification, and the bulk bijection.}
  Obligation~\ref{obl:taubdry}(i)--(iv) is evaluated on the finished table in the
  constructor; a failure aborts. The element $\gamma_j$ of
  Proposition~\ref{prop:bdrycosets}(2), fixed there only up to $\gamma_j\Gbdry$, is
  normalised as $\gamma_j:=\mathrm{red}(g^{-1}s_j\infty)^{-1}$ (WP-4a, finding 5), and
  the bijection $j\mapsto\delta_j\gamma_j\Gbdry$ is certified.
\end{enumerate}
\textbf{Properties.} \PROVED{}
\begin{enumerate}
\item \textbf{(Well defined, and independent of the derivation order.)} Steps~1--2
  define a function $w$ on the vertices $x$ of $\Xbdry$,
  $w(x)=\gamma_{\mathrm{red}}^{-1}\cdot\mathrm{nearest}\bigl(\mathrm{Re}(\gamma_{\mathrm{red}}\,g^{-1}\mathrm{pos}(x))\bigr)$
  with $\gamma_{\mathrm{red}}$ \emph{any} reduction of the cusp of $g^{-1}\mathrm{pos}(x)$,
  and $w(h\cdot x)=(g^{-1}hg)\cdot w(x)$ for every $h\in\Ga''$. Hence a slot derived in
  step~0 carries the chain its direct construction would give, in every degree, and the
  table is independent of the order of step~0 and of the cusp reduction.
\item \textbf{(It is a boundary datum over $\Z$.)} The output satisfies
  (T1$\partial$)--(T3$\partial$), with coefficients in $\{0,\pm1\}$ and solve regime
  $\Z$, for every $\ell$: it is an instance of Lemma~\ref{lem:taubdryexists}(1)--(2).
\item \textbf{(Any other admissible $\taub$ is homotopic to it)} over $\Z$, for every
  $\ell$, by Lemma~\ref{lem:taubdryexists}(3); the vertex rule affects the chains and
  the size of \S\ref{sec:cone}'s $\etapair$, and nothing on cohomology
  (Proposition~\ref{prop:bdryop}(1)).
\item \textbf{(No $\ell$ and no $n$ in the table.)} Steps~0--4 read neither $\ell$ nor
  $n$: the table is a datum over $\Z$, the same for every coefficient ring and element
  width, and the datum over $\cring$ is its base change.
\end{enumerate}
\end{algorithm}

\begin{proof}[Proof of the properties]
(1) For a vertex $x$ of $\Xbdry$ in the component over the cusp $c$ put
$y:=g^{-1}\mathrm{pos}(x)\in\Hyp$, which lies over $c'=g^{-1}c$; a reduction
$\gamma_{\mathrm{red}}$ is any element of $\Ga$ with $\gamma_{\mathrm{red}}c'=\infty$,
unique up to left multiplication by $\mathrm{Stab}_{\Ga}(\infty)=\Gbdry=\{\pm T^{k}\}$
\eqref{eq:bdrysplit}. Replacing $\gamma_{\mathrm{red}}$ by $\pm T^{k}\gamma_{\mathrm{red}}$
replaces $x=\mathrm{Re}(\gamma_{\mathrm{red}}y)$ by $x+k$ ($-I$ acts trivially on
$\Hyp$), and the nearest position by $P+k\cdot\mathrm{per}$: the distances
$\abs{x-\mathrm{Re}(T^{k'}v_i)}$ and the tie rule ``larger real part'' are both invariant
under a common translation. So the vertex becomes $T^{k}\cdot(\text{vertex})$, which
$(T^{k}\gamma_{\mathrm{red}})^{-1}=\gamma_{\mathrm{red}}^{-1}T^{-k}$ carries to the same
vertex of $\Xbdry_{c'}$ as before: $w(x)$ is well defined. For $h\in\Ga''$ put
$\gamma:=g^{-1}hg\in\Ga$ (Proposition~\ref{prop:cosets}(3)); then
$g^{-1}\mathrm{pos}(hx)=\gamma y$ lies over $\gamma c'$, $\gamma_{\mathrm{red}}\gamma^{-1}$
is a reduction of $\gamma c'$, and
$w(hx)=(\gamma_{\mathrm{red}}\gamma^{-1})^{-1}\cdot\mathrm{nearest}\bigl(\mathrm{Re}(\gamma_{\mathrm{red}}\gamma^{-1}\gamma y)\bigr)=\gamma\cdot w(x)$.
A derived vertex slot is therefore
$(g^{-1}hg)\cdot\taub(s_{j'}v')=\gamma\,w(s_{j'}v')=w(hs_{j'}v')$, the directly
constructed value; a derived edge slot is the (T1$\partial$)-translate of a walk, and
the translate of the unique $1$-chain with a given boundary is the unique $1$-chain with
the translated boundary ($\Ga$ acts by chain maps, and $Z_1(\Xbdry_{c})=0$ by
\ref{B2}), again the directly constructed value. (2) is
Lemma~\ref{lem:taubdryexists}(1)--(2) with the choice of $w$ made by step~2:
(T1$\partial$) holds on every slot by step~0 and (1), (T2$\partial$) by step~3,
(T3$\partial$) by step~2 and Lemma~\ref{lem:bdrycomponent}, and the coefficients are
those of a single vertex and of an oriented walk. (3) is Lemma~\ref{lem:taubdryexists}(3)
and Proposition~\ref{prop:bdryop}(1). (4) is read off the steps: the inputs are $N$, the
cosets and the exact geometry of the boundary path.
\end{proof}

\begin{remark}[$\taub$ is not made $\inv$-symmetric, and no vertex rule can make it so]\label{rmk:taubdrynoiota}
\PROVED{} Let $\inv_{\Xbulk}$ be the involution of \S\ref{sec:involution}
(Definition~\ref{def:iota}: $z\mapsto-\bar z$, realised by
$\smat{-1}{0}{0}{1}\in\GL_2(\Z)$, with $\inv(\gamma z)=\conj_{\inv}(\gamma)\inv(z)$ and
$\conj_{\inv}\smat{a}{b}{c}{d}=\smat{a}{-b}{-c}{d}$). It preserves $\Xbdry$, each
horocycle and the cusp $\infty$, and acts on $\Xbdry_{\infty}$ as $\invK$ on $\Kbdry$
extended by $\inv(T^{k}x)=T^{-k}\inv(x)$. Hecke-extensions slot 1, WP-4a, finding 3;
open point (7) of the v1.3.3 changelog.
\begin{enumerate}
\item \textbf{(The conjugate datum is admissible, and homotopic.)}
  $\taub':=\inv_{\Xbulk}\,\taub\,\inv_{\Xbulk}$ is again a boundary $g$-twisted chain
  datum over $\Z$ for the same $g=\mathrm{diag}(1,N)$, and $\taub'\simeq\taub$ by
  Lemma~\ref{lem:taubdryexists}(3): the two induce the same $\Thb_g$ on cohomology.
\item \textbf{(No vertex rule is $\inv$-equivariant.)} Let $w$ be \emph{any} rule that
  assigns to a rational $x$ --- the transported real part of Algorithm~\ref{alg:taubdry}
  step~1 --- a vertex $w(x)$ of $\Xbdry_{\infty}$, $T$-equivariantly
  ($w(x+1)=T\,w(x)$), with any tie rule whatever or with no reference to distances at
  all. Then $w(-x)=\inv_{\Xbulk}\,w(x)$ fails at $x=0$ --- and, by $T$-equivariance, at
  every integer $x$ --- because no vertex of $\Xbdry_{\infty}$ is $\inv_{\Xbulk}$-fixed.
  An integer transported real part occurs at every $(2,\nu)$ (the vertex $v_0$ of
  $\Kbdry$ has $x=-N/2$; (3) below), so at $p=2$ the $\inv$-symmetry of $\taub$ cannot
  be \emph{arranged by the choice of tie rule} in step~2, nor by any other
  $T$-equivariant vertex rule, in contrast with the bulk, where Algorithm~\ref{alg:tau}
  step~(1b) breaks ties $\invK$-equivariantly through the $\invK$-fixed vertex $v_6$
  (Remark~\ref{rmk:iotatausym}). For odd $p$ no transported real part is an integer
  (proof of (3)), and this clause obstructs nothing there.
\item \textbf{(The canonical table is not $\inv$-symmetric at $(p,\nu)=(2,1)$.)} With
  the shipped coordinates $\mathrm{Re}\,v_1=-\mathrm{Re}\,v_2=-\tfrac16$ (the ones
  Lemma~\ref{lem:iotamodel}'s hypothesis is verified in, \S\ref{sec:involution}), the
  canonical $\taub$ of Algorithm~\ref{alg:taubdry} at $N=2$ has
  $\inv_{\Xbulk}\taub\inv_{\Xbulk}\neq\taub$, at two places on the $\infty$-line itself:
  on the vertex $v_1$ of $\Kbdry$ (the slot $s_1=I$, Definition~\ref{def:cosetreps}) the
  transported real part is $N\,\mathrm{Re}\,v_1=-\tfrac13$, a tie between $v_0$ and
  $v_1$ resolved to $v_1$, while on its mirror $v_2$ it is $+\tfrac13$, a tie between
  $v_2$ and $v_3$ resolved to $v_3$, and $\invK v_3=v_0\neq v_1$; and on $v_0$ it is
  $N\,\mathrm{Re}\,v_0=-1$, an integer, where (2) obstructs every rule: the tie
  $\{T^{-1}v_1,T^{-1}v_2\}$ is resolved to $T^{-1}v_2$, on the mirror $v_3=Tv_0$ the
  tie $\{Tv_1,Tv_2\}$ to $Tv_2$, and $\inv_{\Xbulk}(T^{-1}v_2)=Tv_1\neq Tv_2$. (v1.4.0
  placed the first exhibit on the slots $S\cdot v_1$, $S\cdot v_2$ over the cusp $0$,
  where the transported real parts are $\mp\tfrac1{12}$ and no tie occurs; corrected in
  v1.4.1.)
\item \textbf{(What is and is not asserted.)} Whether the emitted table satisfies
  $\inv_{\Xbulk}\taub\inv_{\Xbulk}=\taub$ at a given $(p,\nu)$ is a property of that
  table: it fails at $(2,1)$ by (3) --- and at every $(2,\nu)$, under every
  $T$-equivariant vertex rule, since $x(v_0)=-2^{\nu-1}$ is an integer --- it is
  asserted nowhere, Obligation~\ref{obl:iota} says nothing about $\taub$, and no
  validation records it. For odd $p$ the mirror point is never a transported real part
  (proof of (3)), so (2) obstructs nothing there; whether a \emph{mirror-consistent}
  $T$-equivariant tie rule (``ties away from $\mathrm{Re}\,z=0$'', read modulo $1$)
  would then give an $\inv$-symmetric table at every odd $(p,\nu)$ is not settled here
  and not claimed. If
  \S\ref{sec:cone} or \S\ref{sec:involution} ever need an $\inv$-symmetric boundary
  datum, that is a contract addition and is not made here (owner close-out decision (a)
  of slot 1, 2026-09-16).
\end{enumerate}
\end{remark}

\begin{proof}
(1) $\conj_{\inv}$ preserves $\Ga$, preserves $\Ga''=\Ga_0(N)$ ($c\mapsto-c$ keeps
$N\mid c$) and fixes $g=\mathrm{diag}(1,N)$. By \ref{I1} $\inv_{\Xbulk}$ is a simplicial
automorphism of $\Xbulk$; it preserves $\mathrm{Im}\,z$, hence the horocycles at height
$Y$ and their $\Ga$-translates (a translate $\gamma\cdot\{\mathrm{Im}\,z=Y\}$ goes to
$\conj_{\inv}(\gamma)\cdot\{\mathrm{Im}\,z=Y\}$), so it preserves $\Xbdry$ and acts on
$C_{\bullet}(\Xbdry)$ as a chain map with
$\inv_{\Xbulk}(\gamma x)=\conj_{\inv}(\gamma)\,\inv_{\Xbulk}(x)$. For $h\in\Ga''$ and
$\inv:=\inv_{\Xbulk}$:
$\taub'(hx)=\inv\taub(\conj_{\inv}(h)\,\inv x)=\inv\bigl((g^{-1}\conj_{\inv}(h)g)\cdot\taub(\inv x)\bigr)
=\conj_{\inv}(g^{-1}\conj_{\inv}(h)g)\cdot\inv\taub\inv(x)=(g^{-1}hg)\cdot\taub'(x)$,
using $\conj_{\inv}(g)=g$ and $\conj_{\inv}^{2}=\id$: (T1$\partial$). (T2$\partial$)
holds because $\inv$ is a chain map. (T3$\partial$): $\inv$ permutes the components
($c\mapsto\conj_{\inv}(\gamma)\infty$ for $c=\gamma\infty$, i.e.\ $c\mapsto-c$ on
$\mathbb{P}^{1}(\Q)$) and carries a single vertex with coefficient $1$ to a single vertex
with coefficient $1$ (vertices carry no orientation sign), so $\taub'$ carries each
component into one component with total coefficient $1$, with component map
$c\mapsto\inv\,g^{-1}\,\inv\,c=g^{-1}c$, as Lemma~\ref{lem:bdrycomponent} forces in any
case. So $\taub'$ is admissible over $\Z$; Lemma~\ref{lem:taubdryexists}(3) and
Proposition~\ref{prop:bdryop}(1) give the rest.

(2) On the horocycle $\{\mathrm{Im}\,z=Y\}$ the involution is $x\mapsto-x$ on real
parts, so $\inv$-equivariance of $w$ at $x=0$ reads $w(0)=\inv_{\Xbulk}w(0)$. The
vertices of $\Xbdry_{\infty}$ are $T^{k}v_i$, $i\in\{0,1,2\}$, and
$\inv_{\Xbulk}(T^{k}v_i)=T^{-k}\invK(v_i)$ with $\invK$ the permutation
$0\leftrightarrow3$, $1\leftrightarrow2$ of Lemma~\ref{lem:iotamodel}: $T^{k}v_0\mapsto
T^{-k}v_3=T^{1-k}v_0$, fixed only if $2k=1$; $T^{k}v_1\mapsto T^{-k}v_2$ and
$T^{k}v_2\mapsto T^{-k}v_1$, never fixed because $v_1$ and $v_2$ lie in different
$\Gbdry$-orbits (Remark~\ref{rmk:bdrymodel}). So no vertex is fixed, and no $w$ satisfies
the identity at $0$. At an integer $k$ the identity reads $w(-k)=\inv_{\Xbulk}w(k)$,
i.e.\ $T^{-k}w(0)=T^{-k}\inv_{\Xbulk}w(0)$ by $T$-equivariance and
$\inv_{\Xbulk}T^{k}=T^{-k}\inv_{\Xbulk}$, which is the identity at $0$. (For the
nearest-vertex rule the tied set at $x=0$ is $\{v_1,v_2\}$, one $\invK$-orbit; the
argument does not need it.)

(3) The vertices of $\Kbdry\subset\Xbdry_{\infty}$ lie over the cusp $\infty$ and are the
slots $(1,v)$, $s_1=I$; by property~(1) of Algorithm~\ref{alg:taubdry} the value at any
vertex is $w$ of it, whichever slot or (T1$\partial$)-translate the implementation
constructs. For $v\in\Kbdry$ step~1 gives $z=g^{-1}\mathrm{pos}(v)=N\,\mathrm{pos}(v)$
($\adjg(g)=\mathrm{diag}(N,1)$ acts as $z\mapsto Nz$), $c'=\infty$,
$\gamma_{\mathrm{red}}=I$ and $x=N\,\mathrm{Re}\,v$. At $N=2$: for $v=v_1$, $x=-\tfrac13$,
equidistant ($\tfrac16$) from $\mathrm{Re}\,v_0=-\tfrac12$ and $\mathrm{Re}\,v_1=-\tfrac16$
and farther from every other line vertex; the tie rule takes $v_1$, so $\taub(v_1)=v_1$.
For $v=v_2$, $x=+\tfrac13$, equidistant from $v_2$ and $v_3$; the tie rule takes $v_3$,
so $\taub(v_2)=v_3$. Now $\inv_{\Xbulk}v_1=\invK v_1=v_2$, so $\inv$-symmetry at $v_1$
would require $\taub(v_2)=\inv_{\Xbulk}\taub(v_1)=\inv_{\Xbulk}v_1=v_2$; but
$\taub(v_2)=v_3\neq v_2$. For $v=v_0$, $x=-1$, equidistant ($\tfrac16$) from
$T^{-1}v_1$ ($\mathrm{Re}=-\tfrac76$) and $T^{-1}v_2$ ($\mathrm{Re}=-\tfrac56$); the tie
rule takes $T^{-1}v_2$, so $\taub(v_0)=T^{-1}v_2$, and
$\taub(v_3)=\taub(Tv_0)=T^{N}\taub(v_0)=Tv_2$ by (T1$\partial$) with $h=T$,
$g^{-1}Tg=T^{N}$ (directly: $x(v_3)=+1$, tie $\{Tv_1,Tv_2\}$, resolved to $Tv_2$).
$\inv$-symmetry at $v_0$ would require
$\taub(v_3)=\inv_{\Xbulk}\taub(v_0)=\inv_{\Xbulk}(T^{-1}v_2)=T\,\invK v_2=Tv_1\neq Tv_2$.
This second failure is clause (2)'s obstruction, attained: $x(v_0)=-N/2$ is an integer
at every $(2,\nu)$. (On the slots $S\cdot v_1$, $S\cdot v_2$ over the cusp $0$, where
v1.4.0 had placed the exhibit, $z=g^{-1}S\,\mathrm{pos}(v)=-N/\mathrm{pos}(v)$,
$\gamma_{\mathrm{red}}=S$, $\gamma_{\mathrm{red}}z=\mathrm{pos}(v)/N$ and
$x=\mathrm{Re}(v)/N=\mp\tfrac1{12}$: no tie, $\taub(Sv_i)=S^{-1}v_i$ for $i=1,2$, which
\emph{is} mirror-consistent; v1.4.0 had applied $g$ in place of $g^{-1}$.)

\emph{Which slots have an integer transported real part.} For a slot $(j,v)$ the matrix
$N\gamma_{\mathrm{red}}g^{-1}s_j=\gamma_{\mathrm{red}}\,\adjg(g)\,s_j$ is integral of
determinant $N$ and fixes $\infty$, hence is $\smat{A}{B}{0}{D}$ with $AD=N$, $A=p^{a}$,
$D=p^{\nu-a}$ after normalising the common sign, and $x=(A\,\mathrm{Re}\,v+B)/D$. An
integer $x$ needs $A\,\mathrm{Re}\,v\in\Z$, i.e.\ $p^{a}/2\in\Z$ ($v\in\{v_0,v_3\}$) or
$p^{a}/6\in\Z$ ($v\in\{v_1,v_2\}$); the second never holds and the first holds exactly
for $p=2$, $a\ge1$. So for odd $p$ no transported real part is an integer, while at every
$(2,\nu)$ the vertex $v_0$ of $\Kbdry$ realises one ($A=N$, $D=1$, $x=-2^{\nu-1}$), and
at $(2,2)$ the slot $\bigl(s_j=\smat{1}{0}{2}{1},v_0\bigr)$ even has $A=D=2$, $B=1$ and
$x=0$. The tie at $x=\pm\tfrac13$ is where $2\,\mathrm{Re}\,v_1$ is the midpoint of
$\mathrm{Re}\,v_0$ and $\mathrm{Re}\,v_1$, i.e.\ exactly at $\mathrm{Re}\,v_1=-\tfrac16$;
with another admissible placement of $v_1$ the first exhibit may disappear, the second
cannot, which is why (3) is stated for the shipped coordinates and (4) claims nothing
beyond what (2) proves at odd $p$.
\end{proof}

\subsection{The boundary operator}\label{sec:bdryop}

\begin{definition}[$\Thb_g$ and $T^{\partial}_p$]\label{def:bdryop}
Let $g\in\De$, let $s_1,\dots,s_m$ and $\delta_j=s_j^{-1}g$ be \emph{the same}
transporters and coset elements as in Definition~\ref{def:hecke}
(Definition~\ref{def:cosetreps}), and let $\taub$ be a boundary $g$-twisted chain datum
over $\cring$. For
$\varphi\in\Hom_{\Z[\Ga]}\bigl(C_{\bullet}(\Xbdry),\VR\bigr)=C^{\bullet}_{\partial}$ define
\begin{equation}\label{eq:bdrychain}
  (\Thb_g\varphi)(x)\;:=\;\sum_{j=1}^{m}\ \rhon(\delta_j)\,\varphi\bigl(\taub(s_j\cdot x)\bigr),
  \qquad x\in C_q(\Xbdry),\ q\ge0 .
\end{equation}
This is \eqref{eq:heckechain} with $\taub$ in place of $\tau$ and $C_{\bullet}(\Xbdry)$
in place of $C_{\bullet}(\Ku)$, and nothing else changed. $T^{\partial}_p$ denotes the
induced map on $H^{\bullet}(C^{\bullet}_{\partial})$ for $g=\mathrm{diag}(1,p)$; the
operator family of Definition~\ref{def:operators} is assembled from these exactly as in
the bulk.
\end{definition}

\begin{decisionbox}
\textbf{The $\Ga$-level $\taub$ on $\Xbdry$ is the normative boundary datum} (owner
decision OP1, 2026-09-15). The one-dimensional datum on $\Xbdry_{\infty}$ obtained by
transporting \eqref{eq:bdrychain} through Proposition~\ref{prop:circlemodel}(1), together
with the left-coset representatives $\delta^{\partial}_{\infty},\delta^{\partial}_{t}$ of
Proposition~\ref{prop:bdrycosets}(1), is its \textbf{Shapiro transcription and a
validation route only} (Obligation~\ref{obl:taubdry}(v)). The reason is economy of
certified objects: with the $\Ga$-level datum, $\taub$ is one more table of the same
shape as $\tau$, certified by the same predicates, and the pair homotopy of
\S\ref{sec:cone} compares two data on \emph{one} model; with a separately constructed
boundary object one would need a fourth certified table, a homotopy between two boundary
cochain maps. (v1.12.0-pre.3: the agreement of \eqref{eq:bdrychain} with the group-level formulas of
Obligation~\ref{obl:taubdry}(v.a)/(v.b) on cohomology is now a theorem ---
Proposition~\ref{prop:bdryop}(2) in degree $0$, Proposition~\ref{prop:bdrytranscription} in
degree $1$; (v.c)'s $\Gbdry$-level chain datum remains a checked route. None of this changes the
decision: \eqref{eq:bdrychain} stays the normative object and the transcription a route.)
\end{decisionbox}

\begin{proposition}[$\Thb_g$ is well defined; and what it is in degree $0$]\label{prop:bdryop}
\PROVED{} Assume \ref{H4}, \ref{B1}--\ref{B3}, and let $g,\{s_j\},\taub,\cring$ be as in
Definition~\ref{def:bdryop}.
\begin{enumerate}
\item \textbf{(No hypothesis on $\ell$.)} $\Thb_g$ carries $C^{\bullet}_{\partial}$ into
  itself and commutes with $d$; hence it descends to $H^{q}(C^{\bullet}_{\partial})$ for
  every $q$. Its independence of $\{s_j\}$, of $\taub$ and of $g\in\Ga g\Ga$ on
  $H^{\bullet}(C^{\bullet}_{\partial})$ holds as in Theorem~\ref{thm:welldef}(a),(c),
  \emph{and (b) holds here over $\Z$ for every $\ell$}, by
  Lemma~\ref{lem:taubdryexists}(3) in place of Lemma~\ref{lem:tauexists}(2).
\item \textbf{(Degree $0$.)} $H^{0}(C^{\bullet}_{\partial})=\VR^{\Gbdry}=\VR^{T}\cap\VR^{-I}$,
  and on it
  \[
    \Thb_g\,v\;=\;\sum_{\delta\in B_N/\Gbdry}\rhon(\delta)\,v ,
  \]
  the classical double-coset action, the sum being over the representatives of
  Proposition~\ref{prop:bdrycosets}(1) and independent of their choice.
\end{enumerate}
\end{proposition}

\begin{proof}
(1) The proof of Theorem~\ref{thm:welldef} uses only (T1), (T2), Lemma~\ref{lem:cocycle}
and Proposition~\ref{prop:symconv}(2); (T1$\partial$) and (T2$\partial$) are the same two
conditions on a different $\Z[\Ga]$-complex, and every step goes through verbatim. Part
(b) there invokes Lemma~\ref{lem:tauexists}(2) \emph{only} to produce the twisted homotopy
$H$; Lemma~\ref{lem:taubdryexists}(3) produces it over $\Z$ with no hypothesis on
$\ell$, and the rest of the argument --- that
$(\widetilde H\varphi)(x)=\sum_j\rhon(\delta_j)\varphi(N_{\partial}(s_jx))$ is
$\Ga$-equivariant, and that the two boundary operators built from $\taub$ and
$\taub'=\taub+\partial N_{\partial}+N_{\partial}\partial$ differ by
$d\widetilde H+\widetilde H d$ --- uses only (T1$\partial$).

(2) A $0$-cochain $\varphi$ is a cocycle iff $\varphi(\partial e)=0$ for every edge,
i.e.\ iff $\varphi$ is constant on each component; $\Ga$-equivariance then makes it
determined by $m_0:=\varphi(v)$ for one vertex $v\in\Xbdry_{\infty}$, subject to
$\rhon(h)m_0=m_0$ for every $h\in\Gbdry$ (the setwise stabiliser of $\Xbdry_{\infty}$,
whose elements move $v$ inside the component on which $\varphi$ is constant); so
$H^0=\VR^{\Gbdry}$, which is $\VR^{T}\cap\VR^{-I}$ by \eqref{eq:bdrysplit}.

For the formula fix $j$ and put $c:=s_j\cdot\infty$, so that $s_j\cdot v$ is a vertex of
$\Xbdry_c$. By (T3$\partial$), $\taub(s_j\cdot v)$ is a single vertex $w_j$ with
coefficient $1$, and by Lemma~\ref{lem:bdrycomponent} it lies in $\Xbdry_{c'}$ with
$c'=g^{-1}c=g^{-1}s_j\infty=\delta_j^{-1}\cdot\infty$. Let $\gamma_j\in\Ga$ be as in
Proposition~\ref{prop:bdrycosets}(2), i.e.\ $\gamma_j\cdot\infty=c'$ and
$\delta_j\gamma_j\in B_N$. Then $\gamma_j\cdot\Xbdry_{\infty}=\Xbdry_{c'}$, and since
$\varphi$ is constant on $\Xbdry_{c'}$,
$\varphi(w_j)=\varphi(\gamma_j\cdot v)=\rhon(\gamma_j)m_0$; by
Proposition~\ref{prop:symconv}(2),
$\rhon(\delta_j)\varphi(w_j)=\rhon(\delta_j\gamma_j)m_0$. Summing over $j$ and using
that $\{\delta_j\gamma_j\}$ represents $B_N/\Gbdry$
(Proposition~\ref{prop:bdrycosets}(2)) gives the displayed sum. It is independent of the
representatives because $\rhon(\delta h)m_0=\rhon(\delta)\rhon(h)m_0=\rhon(\delta)m_0$
for $h\in\Gbdry$ and $m_0\in\VR^{\Gbdry}$.
\end{proof}

\begin{proposition}[The boundary characters, and which coset system gives them]\label{prop:bdrychar}
(Appendix~\ref{app:errata}, entry~\ref{err:e14}, boundary half: the characters are
stated in the normative left-coset convention, and the right-coset convention's
$\det$-twist is stated as part of the proposition rather than hedged.)
Let $\cring$ be a commutative ring of characteristic $\ell$, let $n\ge0$, let $p$ be a
prime with $p\neq\ell$ (\ref{H1}) --- \emph{except in the first sentence of} (1), the
identification of $\VR^{T}$, which uses no hypothesis on $p$ at all --- put
$g=\mathrm{diag}(1,p)$, and write
\[
  \mathrm{Nm}\;:=\;\prod_{c\in\F_\ell}(y+cx)\;=\;y^{\ell}-x^{\ell-1}y\ \in\ \Sym^{\ell}(\Z^2)\otimes\cring,
  \qquad \chi_b(p):=p^{b}+p^{\,n+1-b} .
\]
\emph{Notation.} Claim~\ref{cl:boundary} writes this invariant $N$; in this document
$N=p^{\nu}$ is the Hecke level (\S\ref{sec:notation-primes}), so it is written
$\mathrm{Nm}$ here --- it is the norm, i.e.\ the $\langle T\rangle$-orbit product, of
$y$. This is a genuine collision of the superseded corpus; it was resolved locally in
v0.14.0 and is \textbf{C18} of \S\ref{sec:notation-collisions} since v0.16.0.
\begin{enumerate}
\item \PROVED{} $\VR^{T}=\cring[x,\mathrm{Nm}]_n=\bigoplus_{0\le b\le\lfloor n/\ell\rfloor}
  \cring\cdot x^{\,n-b\ell}\mathrm{Nm}^{\,b}$, of rank $\lfloor n/\ell\rfloor+1$ \LEAN{prop:bdrychar} (this first sentence; not the eigenvector sentence, which needs Proposition~\ref{prop:bdryop}(2)); and for
  $n$ even each $x^{\,n-b\ell}\mathrm{Nm}^{\,b}$ is an eigenvector of $\Thb_g$ on
  $H^{0}(C^{\bullet}_{\partial})=\VR^{T}$ with eigenvalue $\chi_b(p)$.
\item \PROVED{} With the \emph{right}-coset representatives $\beta_b,\beta_{\infty}$ of
  Proposition~\ref{prop:bdrycosets}(4) in place of the left ones, the same monomial is an
  eigenvector with eigenvalue $\chi_{b+1}(p)$. \textbf{The two operators are different},
  and the left-coset one is the one \eqref{eq:bdrychain} computes. \emph{What the
  right-coset sum is} (v1.4.0; open point (6) of the v1.3.3 changelog):
  $M_{\beta}:=\sum_{0\le b<p}\rhon(\beta_b)+\rhon(\beta_{\infty})$ is a well-defined
  endomorphism of $\VR^{T}$ --- each summand maps $\VR^{T}$ into itself, because
  $T\beta_b=\beta_bT^{p}$, $T^{p}\beta_{\infty}=\beta_{\infty}T$ and
  $\VR^{T^{p}}=\VR^{T}$ for $p\neq\ell$ --- and the monomials
  $x^{\,n-b\ell}\mathrm{Nm}^{\,b}$ are an eigenbasis for it. It is asserted as an
  operator on $\VR^{T}$ \textbf{only}: it is not asserted to be induced by any cochain
  map of $C^{\bullet}_{\partial}$, and none is exhibited --- a sum over
  $\Gbdry\backslash B_p$ is not $\Gbdry$-equivariant as a formula on cochains, unlike the
  sum over $B_p/\Gbdry$ of Proposition~\ref{prop:bdryop}(2) --- so ``the right-coset
  $\Thb_g$'' at cochain level is not an object of this document, and a reader should not
  look for one. The implementation builds $M_{\beta}$ as a fibre-level \emph{negative
  control} only (\code{beta\_control\_fiber\_operator}, never used by the pipeline; code
  map). \LEAN{prop:bdrychar} (clause (2) only: the eigenvalue $\chi_{b+1}(p)$ and $M_{\beta}(\VR^{T})\subseteq\VR^{T}$; that the monomials are an eigenbasis rests on (1), which is not formalised)
\item \PROVED{} Assume in addition $n<\ell$, $n$ even, and that $\cring$ is a field.
  Then $H^{0}(C^{\bullet}_{\partial})$ and $H^{1}(C^{\bullet}_{\partial})$ are both
  $1$-dimensional, spanned by $[x^{n}]\in\VR^{T}$ and $[y^{n}]\in(\VR)_{T}$
  respectively, and the classical double-coset operator attached to $B_p/\Gbdry$ acts
  on both by $\chi_0(p)=1+p^{\,n+1}$ --- the Eisenstein eigenvalue of
  Obligation~\ref{obl:oracle}. In degree $0$ that operator is $\Thb_g$
  (Proposition~\ref{prop:bdryop}(2)); in degree $1$ its agreement with $\Thb_g$ is the
  Shapiro transcription, \textbf{proved since v1.12.0-pre.3} as
  Proposition~\ref{prop:bdrytranscription}(3) --- for every $n$, every $\cring$ and every $p$, the
  coordinate $\Phi$ below being an isomorphism whenever $\rhon(-I)=\id$
  (Proposition~\ref{prop:bdrytranscription}(1)) --- and re-checked on every run by
  Obligation~\ref{obl:taubdry}(v.b) (owner decision OP1). (Until v1.12.0-pre.2 this sentence read
  ``which this document does not prove and which Obligation~\ref{obl:taubdry}(v.b) checks''.) \emph{The coordinate}
  (v1.4.0; open point (4) of the v1.3.3 changelog): the identification used here and in
  \eqref{eq:bdrygroupform} is
  \[
    \Phi\colon H^{1}(C^{\bullet}_{\partial})\xrightarrow{\ \sim\ }(\VR)_{T},\qquad
    [\varphi]\longmapsto\bigl[\varphi(\langle v_0,Tv_0\rangle)\bigr],
  \]
  where $\langle v_0,Tv_0\rangle$ is the oriented path from $v_0$ to $v_3=Tv_0$ in
  $\Kbdry$, so that $\varphi(\langle v_0,Tv_0\rangle)$ is the signed sum of $\varphi$
  over its three edges; equivalently $[f(T)]$ for the group $1$-cocycle
  $f(\gamma):=\varphi(\langle v_0,\gamma v_0\rangle)$ of $\langle T\rangle$, which satisfies
  $f(\gamma\gamma')=f(\gamma)+\rhon(\gamma)f(\gamma')$. Its inverse sends $[m]$ to the
  class of the equivariant cochain equal to $m$ on the edge $\langle v_0,v_1\rangle$ and
  to $0$ on the other two edges of $\Kbdry$. Reversing the orientation of the path
  replaces $\Phi$ by $-\Phi$, which leaves the transported operator unchanged; what a
  representative map and a lift must \emph{share} is that sign --- built on opposite
  orientations, the transport of \eqref{eq:bdrygroupform} into the stored basis is off by
  $-1$ and does not agree with \eqref{eq:bdrychain} wherever $2\Abdry{1}\neq0$. On $H^1$
  the operator is
  \begin{equation}\label{eq:bdrygroupform}
    [m]\ \longmapsto\ \Bigl[\ \rhon(\delta^{\partial}_{\infty})
       \Bigl(\sum_{a=0}^{p-1}\rhon(T)^{a}\Bigr)m\ +\ \rhon(\delta^{\partial}_{0})\,m\ \Bigr].
  \end{equation}
\item The statement of (1) and (3) for $H^{1}$ and \emph{arbitrary} $n$ --- that
  $(\VR)_{T}$ carries the same multiset $\{\chi_b\}_{0\le b\le\lfloor n/\ell\rfloor}$ ---
  together with the analysis of when two characters coincide, is
  Claim~\ref{cl:chi} of \S\ref{sec:six} and is not re-proved here. (v1.12.0-pre.3: that the
  operator on $H^{1}$ is, for $n$ even and $p\neq\ell$, diagonalisable over $\Fl$, with eigenvalue $\chi_{n-k}(p)$ on a torus
  weight vector of weight $a^{n-k}d^{k}$, is Lemma~\ref{lem:bdrytorus}; which weights occur --- the
  multiset --- is not.)
\end{enumerate}
\end{proposition}

\begin{proof}
\emph{(1), the invariants.} $T$ acts on $\cring[x,y]$ by $x\mapsto x$, $y\mapsto x+y$
(Definition~\ref{def:rhon}), an automorphism of order $\ell$ generating an additive
$\F_\ell$-action. The identity $\prod_{c\in\F_\ell}(Y+c)=Y^{\ell}-Y$ in $\F_\ell[Y]$,
homogenised by $Y=y/x$ and multiplied by $x^{\ell}$, gives
$\prod_{c\in\F_\ell}(y+cx)=y^{\ell}-x^{\ell-1}y=\mathrm{Nm}$; $T$ permutes the factors,
so $\mathrm{Nm}$ is $T$-invariant, as is $x$. Since $\mathrm{Nm}$ is monic of degree
$\ell$ in $y$ over $\cring[x]$, the elements $x,\mathrm{Nm}$ are algebraically
independent and $\cring[x,y]$ is free over $\cring[x,\mathrm{Nm}]$ with basis
$1,y,\dots,y^{\ell-1}$ (division with remainder by a monic polynomial). Write an
invariant $f$ as
$f=\sum_{t<\ell}a_t(x,\mathrm{Nm})\,y^{t}$. Then
$0=Tf-f=\sum_t a_t\bigl[(x+y)^{t}-y^{t}\bigr]$, and $(x+y)^t-y^t$ has $y$-degree
$t-1$ with leading coefficient $t\,x$. If some $a_t\neq0$ with $t\ge1$, take the largest
such $t$; the coefficient of $y^{t-1}$ in the sum is $t\,x\,a_t$ (all $t'<t$ contribute
$y$-degree $\le t'-1\le t-2$), so $t\,x\,a_t=0$. Here $t$ is a unit in $\cring$ for
$1\le t<\ell$, because $\cring$ has characteristic $\ell$ and therefore contains
$\Fl$; and $x$ is a non-zero-divisor in $\cring[x,y]$ for \emph{every} commutative
$\cring$, the monomials being a basis. So $a_t=0$ --- a contradiction. Hence
$f=a_0(x,\mathrm{Nm})$, and the degree-$n$ part of $\cring[x,\mathrm{Nm}]$ is as
displayed.

\emph{(1), the eigenvalue.} By Proposition~\ref{prop:bdryop}(2) the operator on $H^0$ is
$\sum_{\delta\in B_p/\Gbdry}\rhon(\delta)$, i.e.\
$\rhon(\delta^{\partial}_{\infty})+\sum_{t=0}^{p-1}\rhon(\delta^{\partial}_{t})$. Now
$\delta^{\partial}_{\infty}=\smat{1}{0}{0}{p}$ acts by $x\mapsto x$, $y\mapsto py$, so
$\delta^{\partial}_{\infty}\cdot\mathrm{Nm}=p^{\ell}y^{\ell}-p\,x^{\ell-1}y=p\,\mathrm{Nm}$
in characteristic $\ell$, since $p^{\ell}\equiv p$ (Fermat); hence
$\delta^{\partial}_{\infty}\cdot x^{n-b\ell}\mathrm{Nm}^{b}=p^{b}x^{n-b\ell}\mathrm{Nm}^{b}$.
And $\delta^{\partial}_{t}=\smat{p}{t}{0}{1}$ acts by $x\mapsto px$, $y\mapsto tx+y$, so
using $u^{\ell}\equiv u$ and $p^{\ell-1}\equiv1$ (here $p\neq\ell$, \ref{H1}),
\[
  \delta^{\partial}_{t}\cdot\mathrm{Nm}=(tx+y)^{\ell}-(px)^{\ell-1}(tx+y)
   =t\,x^{\ell}+y^{\ell}-t\,x^{\ell}-x^{\ell-1}y=\mathrm{Nm},
\]
while $\delta^{\partial}_{t}\cdot x^{n-b\ell}=p^{\,n-b\ell}x^{n-b\ell}$ and
$p^{\,n-b\ell}\equiv p^{\,n-b}$. Each of the $p$ elements $\delta^{\partial}_{t}$ therefore
contributes $p^{\,n-b}$, and the total is $p^{b}+p\cdot p^{\,n-b}=\chi_b(p)$.

\emph{(2).} $\beta_{\infty}=\smat{p}{0}{0}{1}$ acts by $x\mapsto px$, $y\mapsto y$, so it
fixes $\mathrm{Nm}$ (again $p^{\ell-1}\equiv1$) and scales $x^{n-b\ell}$ by
$p^{\,n-b}$; $\beta_b=\smat{1}{b}{0}{p}$ acts by $x\mapsto x$, $y\mapsto bx+py$, so
$\beta_b\cdot\mathrm{Nm}=(bx+py)^{\ell}-x^{\ell-1}(bx+py)=b\,x^{\ell}+p\,y^{\ell}-b\,x^{\ell}
-p\,x^{\ell-1}y=p\,\mathrm{Nm}$ and $\beta_b$ fixes $x$. The $p$ elements $\beta_b$
contribute $p\cdot p^{b}$ and $\beta_{\infty}$ contributes $p^{\,n-b}$, total
$p^{\,b+1}+p^{\,n+1-(b+1)}=\chi_{b+1}(p)$. \emph{That $M_{\beta}$ preserves $\VR^{T}$:}
$\smat{1}{1}{0}{1}\smat{1}{b}{0}{p}=\smat{1}{b+p}{0}{p}=\smat{1}{b}{0}{p}\smat{1}{p}{0}{1}$
and $\smat{1}{p}{0}{1}\smat{p}{0}{0}{1}=\smat{p}{p}{0}{1}=\smat{p}{0}{0}{1}\smat{1}{1}{0}{1}$
are the two identities; for $v\in\VR^{T}$ they give
$\rhon(T)\rhon(\beta_b)v=\rhon(\beta_b)\rhon(T)^{p}v=\rhon(\beta_b)v$ and
$\rhon(T)^{p}\rhon(\beta_{\infty})v=\rhon(\beta_{\infty})\rhon(T)v=\rhon(\beta_{\infty})v$,
so $\rhon(\beta_b)v\in\VR^{T}$ and $\rhon(\beta_{\infty})v\in\VR^{T^{p}}$. Finally
$\VR^{T^{p}}=\VR^{T}$: $\rhon(T)$ is unipotent (Definition~\ref{def:rhon}), so in
characteristic $\ell$, $(\rhon(T)-\id)^{\ell^{k}}=\rhon(T)^{\ell^{k}}-\id=0$ for
$\ell^{k}\ge n+1$, i.e.\ $\rhon(T)$ has order a power of $\ell$, and $p$ being prime to
$\ell$, $\rhon(T)^{p}$ generates the same cyclic group as $\rhon(T)$.

\emph{(3).} For $n<\ell$ one has $\lfloor n/\ell\rfloor=0$, so $\VR^{T}=\cring x^{n}$ is
$1$-dimensional and, $\cring$ being a field, so is $(\VR)_{T}$ (both are the kernel and
the cokernel of the same endomorphism $\rhon(T)-\id$ of a finite-dimensional space). The
$H^0$ eigenvalue is (1) at $b=0$. For $H^1$: $H^{1}(C^{\bullet}_{\partial})=
H^{1}(\langle T\rangle;\VR)=(\VR)_{T}$ (Proposition~\ref{prop:circlemodel}(2)), and the
identification is the $\Phi$ of the statement: every equivariant $1$-cochain $\varphi$ on
$\Xbdry_{\infty}$ is a cocycle ($C_2=0$), $f(\gamma):=\varphi(\langle v_0,\gamma v_0\rangle)$
satisfies the cocycle rule because
$\langle v_0,\gamma\gamma'v_0\rangle=\langle v_0,\gamma v_0\rangle+\gamma\langle v_0,\gamma'v_0\rangle$
as chains and $\varphi(\gamma\,c)=\rhon(\gamma)\varphi(c)$; a coboundary $d\psi$ has
$f(T)=\psi(Tv_0)-\psi(v_0)=(\rhon(T)-\id)\psi(v_0)\in\im(\rhon(T)-\id)$, so $\Phi$ is
well defined; it is onto because the cochain $\varphi_m$ of the inverse formula has
$\Phi[\varphi_m]=m$; and it is injective because a $\varphi$ with
$\varphi(\langle v_0,Tv_0\rangle)\in\im(\rhon(T)-\id)$, say $=(\rhon(T)-\id)u$, equals
$d\psi$ for the equivariant $\psi$ with $\psi(v_0):=u$, $\psi(v_1):=u+\varphi(\langle v_0,v_1\rangle)$,
$\psi(v_2):=\psi(v_1)+\varphi(\langle v_1,v_2\rangle)$, which is consistent at
$v_3=Tv_0$ since $\psi(v_2)+\varphi(\langle v_2,v_3\rangle)=u+\varphi(\langle v_0,Tv_0\rangle)=\rhon(T)u=\rhon(T)\psi(v_0)$.
Reversing the orientation replaces $f$ by $-f$ and $\Phi$ by $-\Phi$, and
$(-\Phi)\,A\,(-\Phi)^{-1}=\Phi A\Phi^{-1}$ for any operator $A$. On $(\VR)_{T}$ the
double-coset operator attached to $B_p/\Gbdry$ is computed one double coset at a
time (Proposition~\ref{prop:bdrycosets}(3)). For
$\Gbdry\delta^{\partial}_{\infty}\Gbdry$ the index is $1$, the operator on a $1$-cocycle
$f$ of $\langle T\rangle$ is $f\mapsto\rhon(\delta^{\partial}_{\infty})\circ f\circ
\conj_{\delta^{\partial}_{\infty}}^{-1}$, and
$(\delta^{\partial}_{\infty})^{-1}T\delta^{\partial}_{\infty}=T^{p}$, so
$m=f(T)\mapsto\rhon(\delta^{\partial}_{\infty})f(T^{p})=
\rhon(\delta^{\partial}_{\infty})\bigl(\sum_{a<p}\rhon(T)^{a}\bigr)m$ by the cocycle rule
$f(\gamma\gamma')=f(\gamma)+\gamma\cdot f(\gamma')$. For
$\Gbdry\delta^{\partial}_{0}\Gbdry$ take the transporters $T^{t}$, $0\le t<p$, of
$\langle T^p\rangle\times\langle-I\rangle$ in $\Gbdry$; then $T^{t}\cdot T=T^{t+1}$ has
trivial cocycle for $t<p-1$ and cocycle $T^{p}$ for $t=p-1$, so only the last summand
survives and gives
$m\mapsto\rhon(T^{-(p-1)}\delta^{\partial}_{0})m$, which is
$\rhon(\delta^{\partial}_{0})m$ modulo $\im(\rhon(T)-\id)$. This is
\eqref{eq:bdrygroupform}. Evaluating at $m=y^{n}$: $\delta^{\partial}_{0}=\smat{p}{0}{0}{1}$
sends $y\mapsto y$, so the second term is $y^{n}$; and
$T^{a}y^{n}=(ax+y)^{n}$, $\delta^{\partial}_{\infty}(ax+y)^{n}=(ax+py)^{n}$, so the first
term is $\sum_{a<p}\sum_{k}\binom{n}{k}a^{\,n-k}p^{k}x^{\,n-k}y^{k}$. For $n<\ell$,
$\rhon(T)-\id$ is a single nilpotent Jordan block of rank $n$ whose image is
$\mathrm{span}\{x^{\,n-k}y^{k}:k<n\}$, so every term with $k<n$ dies in $(\VR)_{T}$ and
what is left is $\sum_{a<p}p^{n}y^{n}=p^{\,n+1}y^{n}$. Total $1+p^{\,n+1}=\chi_0(p)$.
\end{proof}

\noindent
\emph{Corroboration.} All three parts were additionally checked by direct polynomial
arithmetic in $\Fl[x,y]$ --- the $T$-invariance of every $x^{\,n-b\ell}\mathrm{Nm}^{\,b}$,
$\dim\VR^{T}=\lfloor n/\ell\rfloor+1$ as $\dim\ker(\rhon(T)-\id)$, the $\chi_b$ of (1),
the $\chi_{b+1}$ of (2) and the $H^{1}$ statement of (3) --- over
$\ell\in\{2,3,5,7,11,13\}$, $n\le24$ and $p\in\{2,3,5,7,11,13\}$ with $p\neq\ell$: $0$
disagreements. This carries \textbf{no} \COMP{} tag, because it was an \emph{ad hoc}
computation and not a registered validation; there is nothing in the tree to re-run.
Obligation~\ref{obl:taubdry}(vi) is where it becomes one.

\begin{proposition}[The Shapiro transcription in degree $1$]\label{prop:bdrytranscription}
\PROVED{} (v1.12.0-pre.3, hecke-extensions slot-12 WP-16: the degree-$1$ identity that
Proposition~\ref{prop:bdrychar}(3) stated and left unproved --- the degree-$1$ twin of
Proposition~\ref{prop:bdryop}(2).) Assume \ref{H4} and \ref{B1}--\ref{B3}. Let $\cring$ be a
commutative ring, $n\ge0$, $p$ a prime, $\nu\ge1$, $N=p^{\nu}$, $g=\smat{1}{0}{0}{N}$, and let
$\{s_j\}$, $\delta_j=s_j^{-1}g$ and $\taub$ be as in Definition~\ref{def:bdryop} --- \emph{any}
admissible choices. Fix a vertex $v_0$ of $\Xbdry_{\infty}$ (for the model of
Definition~\ref{def:model7}, the vertex $v_0$ of $\Kbdry$). For vertices $u,w$ of one component of
$\Xbdry$ write $\langle u,w\rangle$ for the unique $1$-chain of that component with
$\partial\langle u,w\rangle=w-u$ (step~0 of the proof), write $\equiv$ for equality in
$(\VR)_{T}=\VR/\im(\rhon(T)-\id)$, and for $k\in\Z$ put
$[k]_{X}:=\sum_{t=0}^{k-1}X^{t}$ if $k\ge0$ and $[k]_{X}:=-\sum_{t=k}^{-1}X^{t}$ if $k<0$, so that
$X^{k}-1=(X-1)[k]_{X}$ in $\Z[X^{\pm1}]$.
\begin{enumerate}
\item \textbf{(The coordinate, over every $\cring$.)} The map
  $\Phi\colon H^{1}(C^{\bullet}_{\partial})\to(\VR)_{T}$, $[\varphi]\mapsto[\varphi(\langle v_0,Tv_0\rangle)]$,
  is well defined and $\cring$-linear, does not depend on $v_0$, and changes sign with the orientation
  of the path. If $\rhon(-I)=\id$ on $\VR$ --- $n$ even, or $2=0$ in $\cring$ --- it is an isomorphism.
  For the model it is the $\Phi$ of Proposition~\ref{prop:bdrychar}(3), which is therefore an
  isomorphism without that clause's hypotheses $n<\ell$ and ``$\cring$ a field''.
\item \textbf{(Every $\nu$, every $p$.)} Let $\{\beta_C\}$ be any set of representatives of the left
  cosets $C\in B_N/\Gbdry$ of Proposition~\ref{prop:bdrycosets}(1), and for each $C$ let $k_C\in\Z$ be
  defined by $\beta_C^{-1}\,T\,\beta_{T^{-1}C}=\pm T^{k_C}$ (the left-hand side lies in $\Gbdry$,
  because $T\beta_{T^{-1}C}\in C=\beta_C\Gbdry$). Then for every $\varphi\in C^{1}_{\partial}$, with
  $m:=\varphi(\langle v_0,Tv_0\rangle)$,
  \begin{equation}\label{eq:bdrygroupformnu}
    \Phi\bigl(\Thb_g[\varphi]\bigr)\;=\;\Bigl[\ \sum_{C\in B_N/\Gbdry}\rhon(\beta_C)\,[k_C]_{\rhon(T)}\,m\ \Bigr].
  \end{equation}
  Consequently, if $\rhon(-I)=\id$ on $\VR$, the right-hand side is a well-defined endomorphism of
  $(\VR)_{T}$, independent of the representatives, and is the operator induced by
  \eqref{eq:bdrychain} on $H^{1}(C^{\bullet}_{\partial})$ in the coordinate $\Phi$. (Without
  $\rhon(-I)=\id$, over every $\cring$: for representatives with positive diagonal, e.g.\
  $\smat{p^{u}}{b}{0}{p^{\nu-u}}$, every $\beta_C^{-1}T\beta_{T^{-1}C}$ is $+T^{k_C}$, and then
  $\sum_C\rhon(\beta_C)[k_C]_{\rhon(T)}(\rhon(T)-\id)=\sum_C\bigl(\rhon(T\beta_{T^{-1}C})-\rhon(\beta_C)\bigr)
  =(\rhon(T)-\id)\sum_C\rhon(\beta_C)$, so the right-hand side is a well-defined endomorphism of
  $(\VR)_{T}$ for such representatives; v1.12.0-pre.3, referee P4.)
\item \textbf{($\nu=1$: the transcription.)} Let $\nu=1$. Then \eqref{eq:bdrygroupform} is a
  well-defined endomorphism of $(\VR)_{T}$, and
  \[
    \Phi\circ T^{\partial}_{p}\;=\;\eqref{eq:bdrygroupform}\circ\Phi\qquad\text{on }H^{1}(C^{\bullet}_{\partial}).
  \]
  If moreover $\rhon(-I)=\id$ on $\VR$ ($n$ even, or $2=0$ in $\cring$), the operator induced by
  \eqref{eq:bdrychain} on $H^{1}(C^{\bullet}_{\partial})$ \textbf{is} \eqref{eq:bdrygroupform} under
  the isomorphism $\Phi$.
\end{enumerate}
No hypothesis is made on $\ell$, on $n$ against $\ell$, on $p$ against $\ell$ (the boundary operator at $p=\ell$ is
included), or that $\cring$ is a field.
\end{proposition}

\begin{proof}
\emph{Conventions, read off the definitions.} $C^{q}_{\partial}=\Hom_{\Z[\Ga]}(C_{q}(\Xbdry),\VR)$
(Proposition~\ref{prop:circlemodel}), so $\varphi(\gamma x)=\rhon(\gamma)\varphi(x)$ for
$\gamma\in\Ga$ --- the convention of the proof of Proposition~\ref{prop:bdryop}(2); $(d\psi)(x)=\psi(\partial x)$
(proof of Theorem~\ref{thm:welldef}); a cochain is extended $\cring$-linearly to
$C_{\bullet}(\Xbdry)\otimes\cring$, where $\taub$ takes its values (Definition~\ref{def:taubdry}). By
Proposition~\ref{prop:bdryop}(1), $\Thb_g\varphi\in C^{1}_{\partial}$ and $\Thb_g[\varphi]=[\Thb_g\varphi]$.
Since $C_{2}(\Xbdry)=0$ (\ref{B2}), every $1$-cochain is a cocycle. $-I$ fixes every chain
(\ref{H4}(iv)).

\emph{Step 0: chains on $\Xbdry$.} By \ref{B2} each component $\Xbdry_c$ is a contractible
$1$-dimensional complex, so $0\to C_{1}(\Xbdry_c)\xrightarrow{\partial}C_{0}(\Xbdry_c)\xrightarrow{\aug}\Z\to0$
is an exact sequence of free abelian groups, hence split, and it stays exact after
$\otimes\cring$. Summing over the components: (a) $\partial\colon C_{1}(\Xbdry)\otimes\cring\to
C_{0}(\Xbdry)\otimes\cring$ is injective; (b) a $0$-chain is a boundary iff its total coefficient on
every component is $0$, i.e.\ iff its class in $H_{0}(\Xbdry)\otimes\cring=\cring[\pi_{0}(\Xbdry)]$
vanishes. Hence $\langle u,w\rangle$ exists and is unique; $\langle u,w\rangle+\langle w,x\rangle=\langle
u,x\rangle$; $\gamma\langle u,w\rangle=\langle\gamma u,\gamma w\rangle$ for $\gamma\in\Ga$, since $\Ga$ acts
by chain automorphisms permuting the components (\ref{B2}); and
$\langle v_0,T^{k}v_0\rangle=[k]_{T}\cdot\langle v_0,Tv_0\rangle$ for every $k\in\Z$, both sides having
boundary $T^{k}v_0-v_0$ because $(T-1)[k]_{T}=T^{k}-1$.

\emph{(1).} Write $e:=\langle v_0,Tv_0\rangle$. A coboundary has
$(d\psi)(e)=\psi(Tv_0)-\psi(v_0)=(\rhon(T)-\id)\psi(v_0)$, so $\Phi$ is well defined; it is clearly
$\cring$-linear. For another vertex $v_0'$ of $\Xbdry_{\infty}$, step~0 gives
$\langle v_0',Tv_0'\rangle=\langle v_0',v_0\rangle+e+T\langle v_0,v_0'\rangle$, so the two values
differ by $(\rhon(T)-\id)\varphi(\langle v_0,v_0'\rangle)$. Reversing the orientation replaces $e$ by
$-e$. Now let $\rhon(-I)=\id$. A simplicial automorphism of a triangulated line is a shift or a
reflection, and a reflection fixes a vertex or an edge setwise; by \ref{B3} (and \ref{H4}(i), no
inversions) $T$ fixes no simplex of $\Xbdry_{\infty}$, so it is a shift by some $d\neq0$ positions, and
the $\abs{d}$ edges of the segment from $v_0$ to $Tv_0$ --- the support of $e$, each with coefficient
$\pm1$ --- sit at $\abs{d}$ consecutive positions and so lie in pairwise distinct $\langle T\rangle$-orbits.
\emph{Onto:} given $m\in\VR$, pick an edge $\epsilon$ of that segment, with coefficient $c=\pm1$ in $e$;
the $\Gbdry$-equivariant cochain on $\Xbdry_{\infty}$ equal to $\rhon(T)^{k}cm$ on $T^{k}\epsilon$ and
to $0$ on the other orbits exists ($\langle T\rangle$ acts freely and $-I$ trivially, on the edges and on
$\VR$), extends to a unique $\varphi\in C^{1}_{\partial}$ by Proposition~\ref{prop:circlemodel}(1), and has
$\varphi(e)=m$. \emph{One-to-one:} if $\varphi(e)=(\rhon(T)-\id)u$, put $\psi(x):=u+\varphi(\langle v_0,x\rangle)$
for the vertices $x$ of $\Xbdry_{\infty}$. Then $\psi(Tx)=u+\varphi(e)+\rhon(T)\varphi(\langle
v_0,x\rangle)=\rhon(T)\psi(x)$ and $\psi((-I)x)=\psi(x)=\rhon(-I)\psi(x)$, so $\psi$ is
$\Gbdry$-equivariant and extends to a $0$-cochain of $C^{\bullet}_{\partial}$
(Proposition~\ref{prop:circlemodel}(1)); on an edge $\langle x,x'\rangle$ of $\Xbdry_{\infty}$,
$\psi(x')-\psi(x)=\varphi(\langle x,x'\rangle)$, so $d\psi$ and $\varphi$ agree on $\Xbdry_{\infty}$, hence
everywhere by equivariance, and $[\varphi]=0$. For the model, the path of $\Kbdry$ from $v_0$ to
$v_3=Tv_0$ is a $1$-chain with boundary $v_3-v_0$, hence equals $e$, and $\Phi$ is the map of
Proposition~\ref{prop:bdrychar}(3).

\emph{(2), step 1: the coset attached to each transporter.} For each $j$ choose $\gamma_j\in\Ga$ with
$\delta_j\gamma_j\in B_N$ and put $\beta_j:=\delta_j\gamma_j$, $C_j:=\beta_j\Gbdry$. By
Proposition~\ref{prop:bdrycosets}(2) (whose proof uses of $\delta_j$ only that $\delta_j\Ga\in\Ga g\Ga/\Ga$,
so it applies to any admissible $\{s_j\}$), $j\mapsto C_j$ is a bijection onto $B_N/\Gbdry$, the admissible
$\gamma_j$ form the coset $\gamma_j\Gbdry$, and $\gamma_j\cdot\infty=\delta_j^{-1}\cdot\infty$; replacing
$\gamma_j$ within that coset replaces $\beta_j$ by any element of $C_j$, so a choice of the $\gamma_j$
realises the given family $\{\beta_C\}$, and we fix it: $\beta_j=\beta_{C_j}$. For each $j$ let $j\cdot T$ be
the index with $\Ga''s_jT=\Ga''s_{j\cdot T}$ (right translation on $\Ga''\backslash\Ga$, so
$j\mapsto j\cdot T$ is a permutation), $h_j:=s_jTs_{j\cdot T}^{-1}\in\Ga''$ and
$h_j':=g^{-1}h_jg\in\Ga'\subseteq\Ga$ (Proposition~\ref{prop:cosets}(3)). In $\Mat_2(\Z)$,
\begin{equation*}
  \delta_j\,h_j'\;=\;s_j^{-1}h_jg\;=\;T\,s_{j\cdot T}^{-1}g\;=\;T\,\delta_{j\cdot T}. \tag{i}
\end{equation*}
So $\delta_j(h_j'\gamma_{j\cdot T})=T\beta_{j\cdot T}$ is upper triangular and lies in $\Ga g\Ga$, i.e.\ in
$B_N$, and by Proposition~\ref{prop:bdrycosets}(2) $h_j'\gamma_{j\cdot T}=\gamma_j\varepsilon_j$ with
$\varepsilon_j\in\Gbdry$; then $\varepsilon_j=\beta_j^{-1}T\beta_{j\cdot T}$. Hence $T\beta_{j\cdot T}\in
C_j$, i.e.\ $C_{j\cdot T}=T^{-1}C_j$, and $\varepsilon_j=\pm T^{k_{C_j}}$; write $k_j:=k_{C_j}$. \hfill(ii)

\emph{Step 2: the chains.} Put $w_i:=\taub(s_iv_0)$ and $y_j:=\taub(s_je)$. As
$s_je=\langle s_jv_0,s_jTv_0\rangle$ and $s_jTv_0=h_js_{j\cdot T}v_0$, (T2$\partial$) and
(T1$\partial$) give
\[
  \partial y_j\;=\;\taub(h_js_{j\cdot T}v_0)-\taub(s_jv_0)\;=\;h_j'w_{j\cdot T}-w_j .
\]
By (T3$\partial$) and Lemma~\ref{lem:bdrycomponent}, the class of $w_i$ in
$\cring[\pi_{0}(\Xbdry)]$ is that of the cusp $g^{-1}s_i\infty=\delta_i^{-1}\infty=\gamma_i\infty$ (as
$s_iv_0$ lies in the component over $s_i\infty$: $v_0\in\Xbdry_{\infty}$ and the bijection
$\pi_{0}(\Xbdry)\cong\mathbb{P}^1(\Q)$ of \ref{B2} is $\Ga$-equivariant), which is
the class of the vertex $\gamma_iv_0$ (\ref{B2}); \emph{only this class is used}, not that $w_i$ is a single
vertex. By step~0(b), $w_i=\gamma_iv_0+\partial z_i$ for some $z_i\in C_{1}(\Xbdry)\otimes\cring$. The
chain $y_j':=\gamma_j\langle v_0,T^{k_j}v_0\rangle+h_j'z_{j\cdot T}-z_j$ has boundary
$\gamma_jT^{k_j}v_0-\gamma_jv_0+h_j'(w_{j\cdot T}-\gamma_{j\cdot T}v_0)-(w_j-\gamma_jv_0)=h_j'w_{j\cdot T}-w_j$,
because $\gamma_jT^{k_j}v_0=\gamma_j\varepsilon_jv_0=h_j'\gamma_{j\cdot T}v_0$ by (ii); so $y_j=y_j'$ by
step~0(a).

\emph{Step 3: evaluation.} With $m=\varphi(e)$, step~0 and equivariance give $\varphi(\gamma_j\langle
v_0,T^{k_j}v_0\rangle)=\rhon(\gamma_j)[k_j]_{\rhon(T)}m$. Using (i) and
Proposition~\ref{prop:symconv}(2),
\[
  \rhon(\delta_j)\varphi(y_j)\;=\;\rhon(\beta_j)[k_j]_{\rhon(T)}m\;+\;\rhon(T)\rhon(\delta_{j\cdot T})\varphi(z_{j\cdot T})\;-\;\rhon(\delta_j)\varphi(z_j).
\]
\emph{Step 4: the sum.} By \eqref{eq:bdrychain}, $(\Thb_g\varphi)(e)=\sum_j\rhon(\delta_j)\varphi(y_j)$.
Since $j\mapsto j\cdot T$ is a permutation, the last two terms sum over $j$ to
$(\rhon(T)-\id)\sum_j\rhon(\delta_j)\varphi(z_j)\equiv0$: the correction terms telescope. Hence
$\Phi(\Thb_g[\varphi])=[(\Thb_g\varphi)(e)]=\bigl[\sum_j\rhon(\beta_j)[k_j]_{\rhon(T)}m\bigr]$, which is
\eqref{eq:bdrygroupformnu} since $j\mapsto C_j$ is a bijection and $\beta_j=\beta_{C_j}$. The family
$\{\beta_C\}$ was arbitrary. If $\rhon(-I)=\id$, $\Phi$ is an isomorphism by (1), so the right-hand side
equals $\Phi\,\Thb_g\,\Phi^{-1}[m]$: it depends only on $[m]$ and not on the representatives.

\emph{(3).} Take the representatives $\delta^{\partial}_{\infty}$, $\delta^{\partial}_{t}$ of
Proposition~\ref{prop:bdrycosets}(1) and read $T^{-1}C$ and $k_C$ off its identities
$T\delta^{\partial}_{\infty}=\delta^{\partial}_{\infty}T^{p}$, $T\delta^{\partial}_{t}=\delta^{\partial}_{t+1}$
($t<p-1$), $T\delta^{\partial}_{p-1}=\delta^{\partial}_{0}T$: $T^{-1}C_{\infty}=C_{\infty}$ with $k=p$;
$T^{-1}C_{t}=C_{t-1}$ with $k=0$ for $1\le t<p$; $T^{-1}C_{0}=C_{p-1}$ with $k=1$. So the sum in
\eqref{eq:bdrygroupformnu} is
$\rhon(\delta^{\partial}_{\infty})\sum_{a=0}^{p-1}\rhon(T)^{a}m+\rhon(\delta^{\partial}_{0})m$ ---
\eqref{eq:bdrygroupform} on the nose, and (2) gives the displayed identity. That
\eqref{eq:bdrygroupform} is well defined on $(\VR)_{T}$, for every $\cring$ and every $p$: by the
identities just quoted and $\delta^{\partial}_{0}T=T^{p}\delta^{\partial}_{0}$,
$\rhon(\delta^{\partial}_{\infty})\bigl(\sum_{a<p}\rhon(T)^{a}\bigr)(\rhon(T)-\id)=
\rhon(\delta^{\partial}_{\infty})(\rhon(T)^{p}-\id)=(\rhon(T)-\id)\rhon(\delta^{\partial}_{\infty})$ and
$\rhon(\delta^{\partial}_{0})(\rhon(T)-\id)=(\rhon(T)^{p}-\id)\rhon(\delta^{\partial}_{0})$, both with image
in $\im(\rhon(T)-\id)$. The last sentence of (3) is (1).
\end{proof}

\noindent
\emph{What the proof uses, and what it does not} (v1.12.0-pre.3). Only (T1$\partial$), (T2$\partial$),
(T3$\partial$) --- the last only through the $H_{0}$-class of $\taub(s_iv_0)$, so no single-vertex
normalisation of the table and no base change from a $\Z$-datum is needed ---, \ref{B2}, \ref{B3}
(for (1)), \ref{H4}(iv), Lemma~\ref{lem:bdrycomponent}, Proposition~\ref{prop:bdrycosets}(1)--(2),
Proposition~\ref{prop:cosets}(3) and Proposition~\ref{prop:symconv}(2). The choices of $\{s_j\}$, of
$\taub$, of $v_0$ and of the $\gamma_j$ drop out, as in Proposition~\ref{prop:bdryop}(2); the independence
of $\taub$ of Proposition~\ref{prop:bdryop}(1) is not invoked. No division occurs: the correction terms of
step~4 telescope under $j\mapsto j\cdot T$, which is why $p=\ell$ is covered. \emph{At $\nu\ge2$}, (2) is an
$H^{1}$ formula over the \emph{left} cosets of Proposition~\ref{prop:bdrycosets}(1), which hold at every
$\nu$; it needs none of the double-coset structure that Remark~\ref{rmk:bdrynugeneral} records as not
stated, and it does not extend Obligation~\ref{obl:taubdry}(v), which stays (v.a) alone at $\nu\ge2$.
\emph{For $n$ odd and $2\in\cring^{\times}$} both sides of (2) and (3) vanish ($C^{\bullet}_{\partial}=0$,
Proposition~\ref{prop:circlemodel}(2)) while $(\VR)_{T}$ need not; that is why the isomorphism clauses
carry $\rhon(-I)=\id$ and nothing else does.

\noindent
\emph{Corroboration} (v1.12.0-pre.3). The chain-level operator \eqref{eq:bdrychain} was evaluated literally
on $\Xbdry$ --- with random transporters $\{h_js_j\}$, $h_j\in\Ga_0(N)$, a random boundary datum whose
vertex values are $0$-chains of total coefficient $1$ on the forced component plus total-zero chains on it
and on a second component, lines with one and with three vertex orbits per period, and random cochains ---
and compared with \eqref{eq:bdrygroupform} and \eqref{eq:bdrygroupformnu} in $(\VR)_{T}$ over $\Fl$, for
$\ell\in\{2,3,5,7,11\}$, $p\in\{2,3,5,7\}$ ($p=\ell$ included), even $n\le22$ at $\nu=1$, $n\le12$ at
$\nu=2$, and odd $n\le7$ at $\ell=2$: $1488$ trials, $0$ disagreements ($288$ of them at $p=\ell$). The negative control that tests part~(3)'s
coset bookkeeping --- the sum $\sum_{a<p}\rhon(T)^{a}$ attached to $\delta^{\partial}_{0}$ instead of
$\delta^{\partial}_{\infty}$ --- was rejected in $572$ of the $1248$ trials at $\nu=1$ (it coincides with
the true operator on the rest); a second one, $\Phi$ reversed on one side only, is a pure sign control
(rejected exactly when $2\,\Thb_g[\varphi]\neq0$) and tests no step of the proof. The WP-16 referee's
re-run extended the grid ($p=\ell$: $405$ trials, $370$ with $[m]\neq0$) and added a control that does
target the datum axioms: a table with (T3$\partial$)'s total coefficient broken to $c=2$ is rejected
against the formula wherever $(c-1)\Thb_g[\varphi]\neq0$ and matches $c$ times it in every trial
(v1.12.0-pre.3, referee P5). Like the corroboration of
Proposition~\ref{prop:bdrychar} this is an \emph{ad hoc} computation
(\file{prop\_t\_check.py}, outside the tree), not a
registered validation, and carries no \COMP{} tag; the proposition is \PROVED{}, and
Obligation~\ref{obl:taubdry}(v.b) re-checks it on every run at $\nu=1$.

\begin{lemma}[Boundary operators away from $\ell$ factor through the torus]\label{lem:bdrytorus}
\PROVED{} (v1.12.0-pre.3, hecke-extensions slot-12 WP-16; ``Lemma~B$'$'' of the slot-12 finality
proposal.) Let $\cring=\Fl$, $n$ even, $p\neq\ell$ a prime (\ref{H1}), $g=\mathrm{diag}(1,p)$,
$\bar p:=p\bmod\ell$, and assume \ref{H4}, \ref{B1}--\ref{B3}. For $A\in\Mat_2(\Fl)$ write $\rhon(A)$ for
$\rhon(\tilde A)\bmod\ell$, $\tilde A$ any integral lift (well defined by \eqref{eq:rhondef}, and
multiplicative by Proposition~\ref{prop:symconv}(2)). Let $U:=\{\rhon(T)^{k}\}$ and let $\mathbf
T(\Fl)$ be the diagonal torus of $\mathrm{GL}_2(\Fl)$; it normalises the unipotent upper-triangular
group, so $\rhon(\mathbf T(\Fl))$ preserves $\VR^{U}=\VR^{T}=H^{0}(C^{\bullet}_{\partial})$ and
$\im(\rhon(T)-\id)$, hence acts on $(\VR)_{T}\cong H^{1}(C^{\bullet}_{\partial})$ (the latter through
$\Phi$ of Proposition~\ref{prop:bdrytranscription}(1)).
\begin{enumerate}[start=0]
\item On $H^{0}(C^{\bullet}_{\partial})=\VR^{T}$,
  $T^{\partial}_{p}=\rhon(\mathrm{diag}(1,\bar p))+\bar p\,\rhon(\mathrm{diag}(\bar p,1))$.
\item On $H^{1}(C^{\bullet}_{\partial})$, in the coordinate $\Phi$,
  $T^{\partial}_{p}=\bar p\,\rhon(\mathrm{diag}(1,\bar p))+\rhon(\mathrm{diag}(\bar p,1))$.
\item Consequently, in both degrees, $T^{\partial}_{p}$ lies in the image of $\Fl[\mathbf T(\Fl)]$, is
  diagonalisable over $\Fl$; the $T^{\partial}_{p}$ for the primes $p\neq\ell$ commute pairwise; and
  $T^{\partial}_{p}=T^{\partial}_{p'}$ whenever $p\equiv p'\pmod\ell$.
\item On a vector of $\mathbf T(\Fl)$-weight $\mathrm{diag}(a,d)\mapsto a^{\,n-k}d^{\,k}$ the eigenvalue
  of $T^{\partial}_{p}$ is $\chi_{k}(p)$ in degree $0$ and $\chi_{n-k}(p)$ in degree $1$
  ($\chi_b(p)=p^{b}+p^{\,n+1-b}$ as in Proposition~\ref{prop:bdrychar}).
\end{enumerate}
\end{lemma}

\begin{proof}
The preliminaries: $\rhon(T)^{\ell}=\rhon(T^{\ell})=\rhon(I)$ as $T^{\ell}\equiv I\pmod\ell$, so $U$ is
the image of the unipotent upper-triangular group of $\mathrm{SL}_2(\Fl)$;
$\mathrm{diag}(a,d)\smat{1}{u}{0}{1}\mathrm{diag}(a,d)^{-1}=\smat{1}{au/d}{0}{1}$, so for
$t\in\mathbf T(\Fl)$ one has $\rhon(t)\rhon(T)=\rhon(T)^{k}\rhon(t)$ for some $k$, and
$\rhon(T)^{k}-\id=(\rhon(T)-\id)[k]_{\rhon(T)}$; this gives both preservation statements, and
$\VR^{U}=\VR^{T}$. $H^{0}(C^{\bullet}_{\partial})=\VR^{T}$ because $\rhon(-I)=\id$
(Proposition~\ref{prop:bdryop}(2)).

(0) By Proposition~\ref{prop:bdryop}(2) with the representatives of
Proposition~\ref{prop:bdrycosets}(1), $T^{\partial}_{p}=\rhon(\delta^{\partial}_{\infty})
+\sum_{t<p}\rhon(T)^{t}\rhon(\delta^{\partial}_{0})$ on $\VR^{T}$. From
$\delta^{\partial}_{0}T=T^{p}\delta^{\partial}_{0}$, $\rhon(\delta^{\partial}_{0})v\in\VR^{T^{p}}$ for
$v\in\VR^{T}$, and $\VR^{T^{p}}=\VR^{T}$ for $p\neq\ell$ (proof of Proposition~\ref{prop:bdrychar}(2)); so
the sum is $p\,\rhon(\delta^{\partial}_{0})$, and reducing the matrices modulo $\ell$ gives (0).

(1) By Proposition~\ref{prop:bdrytranscription}(3), $T^{\partial}_{p}$ is \eqref{eq:bdrygroupform} under
$\Phi$. On $(\VR)_{T}$, $\rhon(T)$ acts as the identity and $\rhon(\delta^{\partial}_{\infty})=
\rhon(\mathrm{diag}(1,\bar p))$ preserves $\im(\rhon(T)-\id)$ (preliminaries), so
$\rhon(\delta^{\partial}_{\infty})\sum_{a<p}\rhon(T)^{a}m\equiv\bar p\,\rhon(\mathrm{diag}(1,\bar p))m$,
and $\rhon(\delta^{\partial}_{0})=\rhon(\mathrm{diag}(\bar p,1))$.

(2) $\mathbf T(\Fl)$ is abelian and $t^{\ell-1}=1$ for each of its elements, so each $\rhon(t)$ is
annihilated by $X^{\ell-1}-1=\prod_{c\in\Fl^{\times}}(X-c)$, which is split with distinct roots: each
$\rhon(t)$ is diagonalisable over $\Fl$, a commuting family of such is simultaneously diagonalisable, and
so is every $\Fl$-linear combination --- in each degree, on the induced actions. The formulas of (0) and
(1) depend on $p$ only through $\bar p$.

(3) $\mathrm{diag}(1,\bar p)$ and $\mathrm{diag}(\bar p,1)$ act on that weight by $\bar p^{\,k}$ and
$\bar p^{\,n-k}$: in degree $0$, $p^{k}+p\cdot p^{\,n-k}=\chi_{k}(p)$; in degree $1$, $p\cdot
p^{k}+p^{\,n-k}=p^{\,n-k}+p^{\,n+1-(n-k)}=\chi_{n-k}(p)$ (in $\Fl$).
\end{proof}

\noindent
Lemma~\ref{lem:bdrytorus}(2) is the semisimplicity half of Claim~\ref{cl:chi}(2), now proved for
$X_5=H^{1}(C^{\bullet}_{\partial})$ at $n$ even, $p\neq\ell$; the \emph{multiset} half --- which torus
weights $(\VR)_{T}$ carries --- does \textbf{not} follow and stays \CONJ{} there. (3) reads the eigenvalues
off a weight; it does not compute the weights. The degree-$1$ weight rule $\chi_{n-k}$, not $\chi_{k}$,
is consistent with Proposition~\ref{prop:bdrychar}(3): for $n<\ell$ the class of $y^{n}$ ($k=n$) carries
$\chi_{0}(p)$.


\begin{obligation}[Certified $\taub$-table and boundary operator]\label{obl:taubdry}
\LEAN{chk:taubd} \LEAN{chk:taubd-p2} \LEAN{chk:pairgrid} (clauses (i)--(iv) as a predicate, verified at $(p,\nu)=(2,1)$ on the golden table and at the five grid points of [GT]; not (v))
Assume \ref{H4} and \ref{B1}--\ref{B3}. An implementation of the boundary datum emits,
for each $g=\mathrm{diag}(1,N)$ and each $\Ga''$-orbit representative simplex $\sigma$ of
$\Xbdry$, a chain $\taub(s_j\cdot\sigma)=\sum_t n^{(j,\sigma)}_t\,\gamma_t\cdot\sigma_t$
with $n^{(j,\sigma)}_t\in\Z$, $\gamma_t\in\Ga$, in the stored-chain format of
Algorithm~\ref{alg:tau}, and must discharge the following. Each clause names the stored
data it reads, the predicate, the ring it is evaluated over, and the action on failure.
\begin{enumerate}
\item[(i)] \textbf{Structural.} Every emitted chain is supported in a single component,
  the component map agrees with $c\mapsto g^{-1}c$ of Lemma~\ref{lem:bdrycomponent}
  recomputed from $g$ alone, and in degree $0$ each chain is a single simplex with
  coefficient $1$ --- i.e.\ (T3$\partial$). Over $\Z$; a failure \textbf{aborts}.
\item[(ii)] \textbf{(T2$\partial$).} $\partial\taub(s_j\sigma)=\taub(\partial(s_j\sigma))$
  for every stored $(j,\sigma)$ in degree $1$, both sides evaluated through the
  \emph{same} (T1$\partial$) helper that the construction uses. Over $\Z$; a failure
  \textbf{aborts}.
\item[(iii)] \textbf{(T1$\partial$) bookkeeping.} For every stored transporter
  $h\in\Ga''$ occurring in the table, $\taub(h\cdot x)=(g^{-1}hg)\cdot\taub(x)$ as
  \emph{chains}, compared after reduction to canonical $(\,n_t,\gamma_t,\sigma_t)$ form.
  Over $\Z$; a failure \textbf{aborts}. This is Obligation~\ref{obl:tau}'s clause on the
  boundary complex, with $\Xbdry$ in place of $\Ku$.
\item[(iv)] \textbf{Integrality is unconditional here.} Every coefficient
  $n^{(j,\sigma)}_t$ lies in $\{0,\pm1\}$ and the recorded solve regime of a
  $\taub$-table is $\Z$, for \emph{every} $\ell$. Over $\Z$; a failure \textbf{aborts}
  --- unlike Remark~\ref{rmk:tauintegral}'s bulk clause, where a mod-$\ell$ regime is a
  recorded weaker datum, a mod-$\ell$ $\taub$-table is an implementation defect
  (Lemma~\ref{lem:taubdryexists}(2), Remark~\ref{rmk:taubdryintegral}).
\item[(v)] \textbf{The Shapiro transcription agrees --- a cross-check, and a route, not
  a definition.} The normative object is \eqref{eq:bdrychain}; per owner decision OP1
  each comparison below is a \emph{validation route}, and a disagreement is a defect in
  the transcription, never a redefinition of $\Thb_g$ (since v1.12.0-pre.3, with (v.a) and
  (v.b) resting on proved identities --- Propositions~\ref{prop:bdryop}(2) and
  \ref{prop:bdrytranscription}(3) --- a disagreement there is a defect of the implementation). Three comparisons, each between
  the operator induced by \eqref{eq:bdrychain} on $H^{\bullet}(C^{\bullet}_{\partial})$ and
  a second object, as matrices over $\cring$ in the stored basis, entrywise; a failure of
  any \textbf{aborts}. (Until v1.4.0 this clause named one comparison where the chapter
  offers three; open point (1) of the v1.3.3 changelog.)
  \begin{itemize}\setlength{\itemsep}{0pt}
  \item[(v.a)] \textbf{Degree $0$, every $\nu$.} The second object is
    $\sum_{\delta\in B_N/\Gbdry}\rhon(\delta)$ on $H^{0}(C^{\bullet}_{\partial})=\VR^{\Gbdry}$
    over the representatives of Proposition~\ref{prop:bdrycosets}(1)--(2)
    (Proposition~\ref{prop:bdryop}(2)), transported into the stored basis through
    Proposition~\ref{prop:circlemodel}(1): a $0$-cocycle is $(m,m,m,\rhon(T)m)$ on
    $v_0,v_1,v_2,v_3$ with $m\in\VR^{\Gbdry}$.
  \item[(v.b)] \textbf{Degree $1$, $\nu=1$.} The second object is \eqref{eq:bdrygroupform}
    on $H^{1}(C^{\bullet}_{\partial})\cong(\VR)_{T}$ in the coordinate $\Phi$ of
    Proposition~\ref{prop:bdrychar}(3), representative map and lift on the same
    orientation. \emph{Status note} (v1.12.0-pre.3): the identity this comparison asserts is a
    theorem, Proposition~\ref{prop:bdrytranscription}(3), for every $n$, $\cring$ and $p$ (the
    coordinate being an isomorphism at $\rhon(-I)=\id$). (v.b) is therefore
    \textbf{corroboration of a theorem}: what it now tests is the implementation --- the
    $\taub$-table, the stored basis, the transport into it, and that representative map and lift
    share the orientation --- and a failure is an implementation defect. Its predicate, ring and
    action on failure (abort) are unchanged. Until v1.12.0-pre.2 it was the only support of the
    degree-$1$ identity, checked per run on the primes a run builds.
  \item[(v.c)] \textbf{Every computed degree, $\nu=1$: the $\Gbdry$-level datum.} The
    second object is the operator of the decision box after
    Definition~\ref{def:bdryop}: a $\Gbdry$-level twisted chain datum on $\Xbdry_{\infty}$
    built from the line coordinates alone --- one $\delta^{\partial}_{\infty}$-twisted and
    one $\delta^{\partial}_{0}$-twisted chain map, with the transporters $T^{t}$ of
    Proposition~\ref{prop:bdrycosets}(3) and the left cosets of (1) --- each certified
    before use (twisted equivariance on a window of positions, chain map, preservation of
    $C^{\bullet}_{\partial}$). At cochain level the two operators are not asserted equal,
    and no chain homotopy between them is claimed --- the $\Gbdry$-level datum is not a
    boundary $g$-twisted chain datum in the sense of Definition~\ref{def:taubdry}, so
    Lemma~\ref{lem:taubdryexists}(3) does not speak about it; only the cohomology-level
    identity is asserted.
  \end{itemize}
  \textbf{At $\nu\ge2$ clause (v) is (v.a) alone.} (v.b) and (v.c) rest on
  Proposition~\ref{prop:bdrycosets}(3) and \eqref{eq:bdrygroupform}, which are $\nu=1$
  statements (Remark~\ref{rmk:bdrynugeneral}); $H^{1}$ at $\nu\ge2$ is covered by (vii)
  and by the operator's own certificates, Proposition~\ref{prop:bdryop}(1) evaluated in
  \code{compute()}. Owner confirmation 2 of the slot-2 plan (2026-09-16): narrowed to what
  is stated, not extended; open point (2) of the v1.3.3 changelog. (v1.12.0-pre.3: the
  general-$\nu$ formula \eqref{eq:bdrygroupformnu} of Proposition~\ref{prop:bdrytranscription}(2)
  would give a (v.b) at $\nu\ge2$; it is \emph{not} added here --- no implementation computes it,
  and adding a clause is the owner's call under the same confirmation.)
\item[(vi)] \textbf{Oracle --- recorded, and gating only where it is proved.} For
  $n$ even and $p\neq\ell$: the eigenvalues of $\Thb_g$ on
  $H^{0}(C^{\bullet}_{\partial})$ in the basis $x^{\,n-b\ell}\mathrm{Nm}^{\,b}$ are
  $\chi_b(p)$ (Proposition~\ref{prop:bdrychar}(1)); a failure \textbf{aborts}. For
  $n<\ell$ the class of $y^{n}$ in $H^{1}(C^{\bullet}_{\partial})$ has eigenvalue
  $1+p^{\,n+1}$ (Proposition~\ref{prop:bdrychar}(3)); a failure \textbf{aborts}. The
  general-$n$ $H^{1}$ multiset of Claim~\ref{cl:chi} is \textbf{recorded}, not gating,
  until that claim is proved.
\item[(vii)] \textbf{Equivariance of the restriction (forward).} $\rres^{q}\circ
  \Tdc{g}=T^{\partial}_{p}\circ\rres^{q}$ on cohomology in every computed degree, where
  $\rres=j^{\ast}$ is the restriction of \S\ref{sec:cone}. This clause is the
  \emph{boundary datum}'s form of a statement expounded in full at
  Obligation~\ref{obl:coneexact}, whose clause (iii) is the \emph{pair}'s form of it
  (v1.7.1: through v1.6.2 the two forms had the same content and this sentence read
  ``belongs to the pair and is stated in full at''; v1.7.0 strengthened (iii) alone, so
  the two are now distinct clauses binding distinct implementations --- the paragraph
  below, and the open point of the v1.7.1 changelog entry); it is named here because
  it is the only assertion that ties \eqref{eq:bdrychain}'s choice of coset system to
  the bulk, and because it is what Proposition~\ref{prop:bdrychar}(2) settles: with the
  left-coset operator it holds, both sides acting on $\rres^{1}(H^{1})$ by
  $\chi_0(p)=1+p^{\,n+1}$ for $n<\ell$; with the right-coset eigenvalue $\chi_1(p)$ of
  Proposition~\ref{prop:bdrychar}(2) in place of $\Abdry{1}=\chi_0(p)$ on the
  one-dimensional $H^{1}$ ($n<\ell$, Proposition~\ref{prop:bdrychar}(3); the sum
  $M_{\beta}$ itself is an operator on $\VR^{T}=H^{0}$) it fails \textbf{whenever
  $\chi_0(p)\neq\chi_1(p)$ in $\cring$ and $\rres^{1}\neq0$} --- i.e.\ $p\not\equiv1$ and
  $p^{n}\not\equiv1\pmod\ell$; Claim~\ref{cl:chi}(3) says when the two coincide as
  functions of $p$ --- and cannot distinguish the two where $\chi_0(p)=\chi_1(p)$. On
  $H^{0}$ the control distinguishes them only on a class whose restriction to $\VR^{T}$
  has a component $x^{\,n-b\ell}\mathrm{Nm}^{\,b}$ with $\chi_b(p)\neq\chi_{b+1}(p)$,
  which no power of the first Dickson class supplies: the $b$-th power has
  $n=b(\ell+1)$ and restricts to the single monomial $x^{b}\mathrm{Nm}^{b}$, and
  $b\equiv n-b\pmod{\ell-1}$, so $\chi_b\equiv\chi_{b+1}$ identically by
  Claim~\ref{cl:chi}(3) (Remark~\ref{rmk:chicoincide} is the case $b=1$). At
  $6\in\cring^{\times}$ every bulk invariant of degree $n<\ell(\ell-1)$ is such a power
  (Claim~\ref{cl:dimH1}(2)'s Dickson locus with its second generator absent), so on the
  model complexes the $H^{0}$ control cannot bite below that degree --- and did not, on
  the $n\le16$ grid. (Until v1.4.0 the clause said ``fails with the right-coset one'' without
  the qualifier; open point (5) of the v1.3.3 changelog.) At $\ell\in\{2,3\}$ the clause
  is recorded, not asserted; the record is Remark~\ref{rmk:conerestrictsmallell}.
  \textbf{The carve-out stays here, and stays here alone} (v1.7.0). Since v1.7.0
  Obligation~\ref{obl:coneexact}(iii) asserts the \emph{same} identity at \emph{every}
  $\ell$, because an implementation of the \emph{pair} carries the derived homotopy
  $\hder$ and that obligation's clause (i) then proves it with no hypothesis on $\ell$.
  This clause is \textbf{not} strengthened with it: what \emph{this} obligation requires
  an implementation to build is $\taub$ and $\Thb_g$, and this clause reads in addition
  the bulk $\Tdc{g}$ and the restriction $\rres$ --- which is why it is the only clause
  here that ties this chapter to the bulk --- but it requires no $\etapair$ and no
  $\hder$ (\S\ref{sec:cone} builds both), so the proof is unavailable to an
  implementation of the boundary datum, which is what this obligation binds, and a
  clause may not rest on data its own obligation does not require. (v1.7.1: through
  v1.7.0 this sentence read ``the data \emph{this} chapter stores are $\taub$ and
  $\Thb_g$ alone'', which this clause's own statement falsifies.) An implementation that also discharges
  Obligation~\ref{obl:coneexact}(i) asserts the identity at every $\ell$ under that
  clause, and this carve-out is vacuous for it; the asymmetry, and why the two clauses do
  not contradict each other, are stated in full at Obligation~\ref{obl:coneexact}(iii).
\end{enumerate}
\end{obligation}

\subsection{What $\ell=2$ costs, and the truncation policy}\label{sec:ell2}

\noindent
Everything above is uniform in $\ell$. What is not uniform is the \emph{size} of the
answer. This subsection states the degree behaviour at $\ell=2$ --- the boundary is
non-zero in every degree (Lemma~\ref{lem:ell2unbounded}), the bulk is $2$-periodic
above degree $2$, Hecke-equivariantly (Lemma~\ref{lem:bulkperiodic}) --- and then fixes
the policy an implementation follows there: \ref{open:ell2}, \textbf{decided by the
owner on 2026-09-16}.

\begin{lemma}[At $\ell=2$ the boundary cohomology is non-zero in every degree]\label{lem:ell2unbounded}
\PROVED{} Let $\cring$ be a commutative ring with $2=0$ in $\cring$ (in particular
$\cring$ of characteristic $\ell=2$), let $n\ge0$ and $\VR\neq0$. Then
\[
  H^{0}(\partial;\VR)=\VR^{T},\qquad
  0\to(\VR)_{T}\to H^{q}(\partial;\VR)\to\VR^{T}\to0\ \ \text{for every }q\ge1,
\]
and $\VR^{T}\neq0\neq(\VR)_{T}$. Hence $H^{q}(\partial;\VR)\neq0$ for \emph{every}
$q\ge0$: \textbf{$H^{\bullet}(\partial;\VR)$ is unbounded above in degree.} If $\cring$
is a field then $\dim H^{0}(\partial)=\dim\VR^{T}$ and
$\dim H^{q}(\partial)=2\dim\VR^{T}$ for every $q\ge1$; for $\cring=\F_2$ that is
$\lfloor n/2\rfloor+1$ and $2(\lfloor n/2\rfloor+1)$.
\end{lemma}

\begin{proof}
With $2=0$ in $\cring$, $\rhon(-I)=(-1)^{n}\id=\id$, so
Proposition~\ref{prop:bdrycohom}'s last sentence applies and gives
$H^{q}(\langle-I\rangle;\VR)=\VR$ for every $q\ge0$: for $q$ odd it is $\VR[2]=\VR$ and
for $q\ge2$ even it is $\VR/2\VR=\VR$, both because $2=0$. Substituting into
\eqref{eq:bdrylhs} gives the displayed sequences, with the $T$-action on
$H^{q}(\langle-I\rangle;\VR)$ the natural one on $\VR$ because $\Gbdry$ is a direct
product (proof of Proposition~\ref{prop:bdrycohom}). Non-vanishing: $\rhon(T)-\id$ is
nilpotent (Definition~\ref{def:rhon}; the proof of
Proposition~\ref{prop:circlemodel}(3)), say $(\rhon(T)-\id)^{k}=0$; if $\VR^{T}=\ker
(\rhon(T)-\id)$ were $0$ then, taking the largest $i$ with $(\rhon(T)-\id)^{i}\VR\neq0$
--- which exists since $\VR\neq0$ --- that module would be a non-zero submodule of the
kernel, a contradiction; and if $(\rhon(T)-\id)\VR=\VR$ then
$\VR=(\rhon(T)-\id)^{k}\VR=0$, also a contradiction. The dimension count over a field is
$\dim\ker=\dim\coker$ for an endomorphism of a finite-dimensional space; for
$\cring=\F_2$, $\dim\VR^{T}=\lfloor n/\ell\rfloor+1=\lfloor n/2\rfloor+1$ by
Proposition~\ref{prop:bdrychar}(1).
\end{proof}

\begin{remark}[Three complexes, and which one the six-term sequence is a window of]\label{rmk:ell2window}
At $\ell=2$ three objects that coincide for $6\in\cring^{\times}$ come apart, and the
statement ``the sequence is unbounded'' is true of exactly one of them. They must not be
conflated.
\begin{enumerate}
\item \textbf{The model complexes.} $\Ceq(K,\Rel;L)$ vanishes above degree $2$ and
  $C^{\bullet}_{\partial}$ above degree $1$, for every $\ell$, because $K$ is
  $2$-dimensional and $\Kbdry$ is $1$-dimensional. Their cohomologies, and the cone of
  $\rres$ between them, are finite and terminate. \textbf{At $\ell\in\{2,3\}$ they do
  not compute $H^{\bullet}(\Ga;\VR)$, $H^{\bullet}(\partial;\VR)$ or $\Hcpt{\bullet}$}
  (\ref{H3}, Remark~\ref{rmk:orbifold}, Proposition~\ref{prop:circlemodel}(3)); they
  compute the orbifold model cohomology of \ref{H3} and its boundary analogue, which are
  well-defined invariants of the triangulation and are what the shipped pipeline emits.
  Nothing in this document claims more for them, and \ref{open:ell2} does not apply to
  them.
\item \textbf{The group-cohomological objects.} $H^{\bullet}(\partial;\VR)=
  H^{\bullet}(\Gbdry;\VR)$ of Definition~\ref{def:bdrygroup} is non-zero in every degree
  (Lemma~\ref{lem:ell2unbounded}). Consequently the long exact sequence of
  Claim~\ref{cl:les}, read with these objects, \textbf{cannot terminate}: for every
  $q\ge0$, exactness at $H^{q}(\partial)$ gives
  $\coker\bigl(\rres^{q}\colon H^{q}(\Ga;\VR)\to H^{q}(\partial;\VR)\bigr)\hookrightarrow
  \Hcpt{q+1}$, so \emph{at least one} of $H^{q}(\Ga;\VR)$ and $\Hcpt{q+1}$ is non-zero.
  The six terms $X_1,\dots,X_6$ of \S\ref{sec:six} are then a \textbf{window} of a
  sequence with a non-zero term in every degree, and the hypotheses that close the
  window for $6\in\cring^{\times}$ and $n$ even --- $H^{2}(\Ga;\VR)=0$ and
  $\Hcpt{0}=0$ --- are exactly what fails. \textbf{Whether $H^{q}(\Ga;\VR)$ is itself
  non-zero in every degree is not asserted here}, and nothing below needs it: the
  display already puts a non-zero term in every degree of the sequence.
  Claim~\ref{cl:amalgam}'s twelve-torsion vanishing, which does bound the bulk for
  $\ell\ge5$, is unavailable at $\ell=2$, and no argument in this document bounds
  $H^{\bullet}(\Ga;\VR)$ above there. What \emph{is} proved is that it need not be
  bounded to be finite information: Lemma~\ref{lem:bulkperiodic} gives
  $H^{q}(\Ga;\VR)\cong H^{q+2}(\Ga;\VR)$ for every $q\ge2$, every $\cring$ and every
  $n$, by cup product with a canonical class $e\in H^{2}(\Ga;\Z)\cong\Z/12$ and
  compatibly with every Hecke operator $\Tdc{g}$, so the bulk is determined in degrees
  $\ge2$ by $H^{2}$ and $H^{3}$ and is bounded above exactly when both vanish. (Farrell
  cohomology, the packaging in which such periodicity is usually stated, agrees with
  $H^{q}$ only \emph{above} the virtual cohomological dimension --- $1$ for
  $\SL_2(\Z)$ --- and is \textbf{not used}: the lemma is about ordinary cohomology and is
  proved from the tree sequence \eqref{eq:treeses} of Claim~\ref{cl:amalgam}(1) ---
  the Mayer--Vietoris sequence of Claim~\ref{cl:amalgam}(2), re-derived in the proof for
  an arbitrary coefficient module --- and the cup product, with no Farrell theory and
  no vcd bookkeeping.)
\item \textbf{The free $\Ga$-model.} The complexes of \S\ref{sec:orbifoldmodel} compute
  (2) and are infinite-dimensional in the degree direction for every $\ell$; they are
  the only object of the three to which a truncation policy applies. \ref{open:ell2} is
  a statement about them.
\end{enumerate}
\end{remark}

\begin{lemma}[Bulk periodicity above the virtual cohomological dimension, and its Hecke equivariance]\label{lem:bulkperiodic}
Let $\Ga=\SL_2(\Z)=\langle S\rangle\ast_{\langle-I\rangle}\langle U\rangle$ as in
Claim~\ref{cl:amalgam}, with the Bass--Serre tree $\mathcal T$ of its clause (1): vertex
set $\Ga/\langle S\rangle\sqcup\Ga/\langle U\rangle$, edge set $\Ga/\langle-I\rangle$,
the edge $\gamma\langle-I\rangle$ joining $\gamma\langle S\rangle$ to
$\gamma\langle U\rangle$, on which $\Ga$ acts by left translation. All cohomology is
\emph{ordinary} group cohomology. Write $\Ga_{\mathrm{fin}}$ for the set of elements of
finite order of $\Ga$ and, for $h\in\Ga_{\mathrm{fin}}$,
\[
  \zeta_+(h)\in\Cx^{\times}\ :=\ \text{the eigenvalue of $h$ on the vector
  $(z_0,1)^{\mathsf T}$, where $z_0\in\Hyp$ is a fixed point of $h$ on $\Hyp$}
\]
($\zeta_+(\pm I)=\pm1$). For a finite cyclic subgroup $C\le\Ga$ let
$\kappa_C\colon H^{2}(C;\Z)\to\Hom(C,\Q/\Z)$ be the inverse of the Bockstein of
$0\to\Z\to\Q\to\Q/\Z\to0$, and let $\chi_C\colon C\to\Q/\Z$ be $\zeta_+|_C$ read through
$\exp(2\pi i\,\cdot)\colon\Q/\Z\xrightarrow{\ \sim\ }\mu_{\infty}$.
\begin{enumerate}
\item \PROVED{} \textbf{(The class $e$.)} $\zeta_+$ is well defined on
  $\Ga_{\mathrm{fin}}$, takes values in the roots of unity, and satisfies
  $\zeta_+(ghg^{-1})=\zeta_+(h)$ whenever $g\in\GL_2(\Q)^{+}$ and
  $h,\,ghg^{-1}\in\Ga_{\mathrm{fin}}$. On every finite cyclic $C\le\Ga$, $\chi_C$ is a
  well-defined \emph{faithful} character, and $\kappa_C$ is an isomorphism, natural in
  $C$. There is a \textbf{unique} class $e\in H^{2}(\Ga;\Z)$ with
  \[
    \kappa_{\langle S\rangle}\bigl(\res^{\Ga}_{\langle S\rangle}e\bigr)=\chi_{\langle S\rangle},
    \qquad
    \kappa_{\langle U\rangle}\bigl(\res^{\Ga}_{\langle U\rangle}e\bigr)=\chi_{\langle U\rangle};
  \]
  $H^{2}(\Ga;\Z)\cong\Z/12$ and $e$ generates it; and
  $\kappa_C(\res^{\Ga}_{C}e)=\chi_C$ for \emph{every} finite cyclic $C\le\Ga$, so that
  $\res^{\Ga}_{C}e$ generates $H^{2}(C;\Z)\cong\Z/\abs{C}$.
\item \PROVED{} \LEAN{thm:mp-tree-periodic} (only the existence of an isomorphism $H^q(\Ga;V)\cong H^{q+2}(\Ga;V)$ for every
  $q\ge2$, every commutative ring and every $V$, from the literal $2$-periodicity of $P_\bullet$
  (Remark~\ref{rmk:treeconeres}), natural in $V$ on the cohomology of the explicit complex $V^2\to V^3\to\cdots$)
  \LEAN{thm:mp-periodic-natural} (a periodicity isomorphism $H^q(\Ga;-)\cong H^{q+2}(\Ga;-)$, $q\ge2$, that is a natural
  isomorphism of functors on \texttt{groupCohomology}, built through a comparison with the bar resolution that is itself natural
  in $V$; not shown to coincide with the isomorphism of the preceding marker; \textbf{not} that it is $e\cup-$, not its
  compatibility with restriction, not the $q=1$ clause)
  \textbf{(Periodicity.)} Let $V$ be \emph{any} $\Z[\Ga]$-module. The cup
  product $e\cup-\colon H^{q}(\Ga;V)\to H^{q+2}(\Ga;V)$ is an isomorphism for every
  $q\ge2$; it is natural in $V$ --- hence $\cring$-linear for $V=\VR$ --- and compatible
  with restriction to every subgroup. For $q=1$ it is surjective, with kernel
  $\partial\bigl(N_{-I}V\bigr)$, where $N_{-I}V:=(1+(-I))\,V\subseteq V^{-I}$ and
  $\partial\colon V^{-I}=H^{0}(\langle-I\rangle;V)\to H^{1}(\Ga;V)$ is the connecting
  map of the Mayer--Vietoris sequence (\eqref{eq:mvamalgam}, in the coefficient form
  $(\ast)$ of the proof); it is not injective in general: for $V=\VR$ with
  $6\in\cring^{\times}$ it is the zero map onto $H^{3}(\Ga;\VR)=0$
  (Claim~\ref{cl:amalgam}(3)) on a source that is non-zero for $\cring=\Fl$, $\ell\ge5$,
  $n=2$ (the Euler-characteristic identity of that claim gives
  $\dim H^{1}=\dim H^{0}+3-1-1\ge1$), while for $n$ odd, $N_{-I}\VR=0$ and it is
  injective. Consequently
  \[
    H^{q}(\Ga;\VR)\;\cong\;H^{q+2}(\Ga;\VR)\qquad\text{for every }q\ge2,\ \text{every
    commutative ring }\cring\text{ and every }n\ge0 :
  \]
  $H^{\bullet}(\Ga;\VR)$ is determined in degrees $\ge2$ by $H^{2}(\Ga;\VR)$ and
  $H^{3}(\Ga;\VR)$, and vanishes in every degree $\ge2$ if and only if both vanish.
\item \PROVED{} \textbf{(Hecke equivariance.)} Let $g\in\De$, $\Ga'=\Ga\cap g^{-1}\Ga g$,
  $\Ga''=\Ga\cap g\Ga g^{-1}$, and let
  $\Tdc{g}=\cor^{\Ga}_{\Ga''}\circ(\conj_g)_{\ast}\circ\res^{\Ga}_{\Ga'}$ be the operator
  \eqref{eq:corconjres} on the ordinary group cohomology $H^{\bullet}(\Ga;\VR)$, for any
  commutative ring $\cring$ and any $n\ge0$ --- it is defined with no hypothesis on
  $\ell$ and with no model. (It is the operator the shipped pipeline computes when
  $6\in\cring^{\times}$, Theorem~\ref{thm:classical}, and the one a free $\Ga$-model of
  \S\ref{sec:orbifoldmodel} computes for every $\ell$, Remark~\ref{rmk:borelfree};
  nothing in this clause depends on either identification.) Then
  \begin{enumerate}[label=(\alph*)]
  \item $e$ is $\De$-invariant: $\iota_g^{\ast}e=\res^{\Ga}_{\Ga''}e$ in
    $H^{2}(\Ga'';\Z)$, where $\iota_g\colon\Ga''\to\Ga$, $h\mapsto g^{-1}hg$;
  \item $\Tdc{g}(e\cup x)=e\cup\Tdc{g}(x)$ for every $x\in H^{q}(\Ga;\VR)$, $q\ge0$.
  \end{enumerate}
  Consequently, for every $q\ge2$, $e\cup-\colon H^{q}(\Ga;\VR)\to H^{q+2}(\Ga;\VR)$ is
  an isomorphism of modules over the $\cring$-algebra generated by
  $\{\Tdc{g}:g\in\De\}$: \textbf{the Hecke module $H^{q}(\Ga;\VR)$ is $2$-periodic in
  $q\ge2$}, and $\Amat{q+2}(\Tdc{g})$ is $\Amat{q}(\Tdc{g})$ written in the transported
  basis $e\cup(\text{basis of }H^{q})$ --- $\Amat{q}(\Tdc{g})$ denoting here the matrix
  of the operator of this clause on $H^{q}(\Ga;\VR)$, which at $\ell\in\{2,3\}$ is not
  the shipped model's matrix (Remark~\ref{rmk:ell2window}(1)).
\end{enumerate}
\end{lemma}

\noindent\textbf{External inputs.} Four standard facts are used in the form
\ref{Ecup}--\ref{Ecor} below, together with \cite[Ch.~III.10.2]{brown} ($\abs{G}$
annihilates $H^{q}(G;-)$ for finite $G$ and $q\ge1$, a numbered statement already used in
Claim~\ref{cl:amalgam}(3)) in Steps~0 and~5; the fixed-vertex statement of Step~3 is
cited to Serre and proved in place. Each is cited at chapter-and-section level; where the
pointer is to a section rather than to a numbered statement it has \emph{not} been
verified against the text (rule~\ref{A2}). What is then relied on is exactly the
displayed form: \ref{Ecup} and the first two items of \ref{Epair} are direct checks on
the cochain formulae \eqref{eq:cupcochain} and
$\psi\mapsto f\circ\psi\circ(\alpha\times\dots\times\alpha)$, which the reader may take as
the definitions; the last item of \ref{Epair} (an inner automorphism acts trivially),
\ref{Eshapiro} and \ref{Ecor} are structural and are used as stated, \ref{Ecor} being the
description of the corestriction $\cor^{\Ga}_{\Ga''}$ that \eqref{eq:corconjres} means
(the transfer formula on a resolution is the form the proof of
Theorem~\ref{thm:classical} computes with). (Referee pass, v1.3.1.)
\begin{enumerate}[label=(E\arabic*)]
\item\label{Ecup} \STD{\cite[Ch.~V, \S3]{brown}: the cup product} For a group $G$ and
  $G$-modules $A$, $B$ there is a product
  $\cup\colon H^{p}(G;A)\otimes H^{q}(G;B)\to H^{p+q}(G;A\otimes_{\Z}B)$, induced on
  inhomogeneous cochains by
  \begin{equation}\label{eq:cupcochain}
    (u\cup\varphi)(g_1,\dots,g_{p+q})
      =u(g_1,\dots,g_p)\otimes(g_1\cdots g_p)\,\varphi(g_{p+1},\dots,g_{p+q}),
  \end{equation}
  satisfying the Leibniz rule $d(u\cup\varphi)=du\cup\varphi+(-1)^{p}u\cup d\varphi$ and
  natural in $A$ and in $B$. With $A=\Z$ (trivial action, $\Z\otimes B=B$) it makes
  $H^{\bullet}(G;B)$ a module over $H^{\bullet}(G;\Z)$, with $1\in H^{0}(G;\Z)=\Z$ acting
  as the identity.
\item\label{Epair} \STD{\cite[Ch.~III, \S8]{brown}: compatible pairs} A pair
  $(\alpha\colon G'\to G,\ f\colon A\to A')$ with $f(\alpha(g')a)=g'f(a)$ induces
  $(\alpha,f)^{\ast}\colon H^{\bullet}(G;A)\to H^{\bullet}(G';A')$,
  $\psi\mapsto f\circ\psi\circ(\alpha\times\dots\times\alpha)$, functorially in the pair;
  $\res^{G}_{H}=(\mathrm{incl},\id)^{\ast}$; and for $\gamma\in G$ and trivial $A$ the pair
  $(c_\gamma,\id_A)$, $c_\gamma(h)=\gamma^{-1}h\gamma$, induces the \emph{identity} of
  $H^{\bullet}(G;A)$.
\item\label{Eshapiro} \STD{\cite[Ch.~III.6.2 and Exercise~III.8.2]{brown}: Shapiro's
  lemma in the form already used at Propositions~\ref{prop:lhcompliant} and
  \ref{prop:alpha}} For $H\le G$ and an $H$-module $B$ put
  $\mathrm{Coind}^{G}_{H}B:=\Hom_{\Z[H]}(\Z[G],B)$ with $(\gamma\varphi)(x)=\varphi(x\gamma)$.
  Then $\mathrm{sh}\colon H^{q}(G;\mathrm{Coind}^{G}_{H}B)\xrightarrow{\ \sim\ }H^{q}(H;B)$ is
  $\res^{G}_{H}$ followed by the map induced by the $H$-map
  $\mathrm{ev}\colon\varphi\mapsto\varphi(1)$. In particular $H^{q}(G;\mathrm{Coind}^{G}_{1}B)=0$
  for $q\ge1$.
\item\label{Ecor} \STD{\cite[Ch.~III, \S\S5, 9]{brown}: induction, coinduction and the
  transfer} For $[G:H]<\infty$, $\Ind^{G}_{H}B=\Z[G]\otimes_{\Z[H]}B\cong\mathrm{Coind}^{G}_{H}B$
  as $G$-modules, and $\cor^{G}_{H}=\mathrm{tr}_{\ast}\circ\mathrm{sh}^{-1}$ for the $G$-map
  \begin{equation}\label{eq:trmap}
    \mathrm{tr}\colon\mathrm{Coind}^{G}_{H}B\to B,\qquad
    \mathrm{tr}(\varphi)=\sum_{s\in G/H}s\,\varphi(s^{-1})
  \end{equation}
  (independent of the representatives $s$, $\varphi$ being $H$-equivariant); on a
  projective resolution this is the transfer formula used in the proof of
  Theorem~\ref{thm:classical}.
\end{enumerate}

\begin{proof}
\emph{Step 0: $\kappa_C$.} For finite $C$ and $q\ge1$, $\abs{C}$ annihilates
$H^{q}(C;\Q)$ \STD{\cite[Ch.~III.10.2]{brown}, as in the proof of
Claim~\ref{cl:amalgam}(3)} and also acts invertibly on it (it is induced by the
automorphism $\abs{C}\cdot$ of $\Q$), so $H^{q}(C;\Q)=0$. The long exact sequence of
$0\to\Z\to\Q\to\Q/\Z\to0$ therefore gives an isomorphism
$\delta\colon H^{1}(C;\Q/\Z)=\Hom(C,\Q/\Z)\to H^{2}(C;\Z)$; $\kappa_C:=\delta^{-1}$.
(v1.10.0: $\delta$ is the snake-lemma connecting map, $\delta[c]=[d\tilde c]$ for a $\Q$-valued
cochain $\tilde c$ lifting the $\Q/\Z$-valued cocycle $c$, with the cochain differential
$f\mapsto f\circ\partial$ on a projective resolution carrying no sign --- the same convention for
$C$ and for $\Ga$. This fixes the sign of the class $e$ of~(1): under the opposite convention the
class so characterised is $-e$. Nothing in~(2) or~(3) depends on the sign, since $-e$ also
generates; hecke-extensions slot-10 WP-14 review, finding m-6, where the built class is checked
against this convention.)
Naturality: for a homomorphism $\alpha\colon C'\to C$ the connecting maps commute with
$(\alpha,\id)^{\ast}$ (\ref{Epair}), so $\kappa_{C'}(\alpha^{\ast}w)=\kappa_C(w)\circ\alpha$.
For $C$ cyclic of order $m$, $\Hom(C,\Q/\Z)\cong\Z/m$, generated by any faithful
character.

\emph{Step 1: $\zeta_+$.} Let $h\in\Ga_{\mathrm{fin}}$ have order $m$. Then $h$ is
diagonalisable over $\Cx$ (its minimal polynomial divides the separable $X^{m}-1$) with
eigenvalues $\zeta,\zeta^{-1}$, roots of unity. If $h=\pm I$ every $z_0\in\Hyp$ is fixed
and every vector is an eigenvector with eigenvalue $\pm1$, so $\zeta_+(h)=\pm1$ is well
defined. Otherwise $\zeta\neq\zeta^{-1}$ (a diagonalisable matrix with a single eigenvalue
$\pm1$ is $\pm I$), so $\zeta\notin\Rl$, $t:=\operatorname{tr}h=\zeta+\bar\zeta$ satisfies
$\abs t<2$, and writing $h=\smat{a}{b}{c}{d}$: $4bc=4ad-4<4ad-(a+d)^{2}=-(a-d)^{2}\le0$,
so $c\neq0$. The fixed points of $z\mapsto(az+b)/(cz+d)$ are the two roots of
$cz^{2}+(d-a)z-b=0$, whose discriminant is $(a+d)^{2}-4<0$: one root $z_0$ in $\Hyp$, the
other $\bar z_0$. From $h(z_0)=z_0$,
\[
  h\begin{pmatrix}z_0\\1\end{pmatrix}
  =\begin{pmatrix}az_0+b\\ cz_0+d\end{pmatrix}
  =(cz_0+d)\begin{pmatrix}h(z_0)\\1\end{pmatrix}
  =(cz_0+d)\begin{pmatrix}z_0\\1\end{pmatrix},
\]
so $(z_0,1)^{\mathsf T}$ is an eigenvector, $\zeta_+(h)=cz_0+d$ is well defined (the fixed
point in $\Hyp$ is unique), and $(\bar z_0,1)^{\mathsf T}$ is the eigenvector for the
other eigenvalue $\overline{\zeta_+(h)}=\zeta_+(h)^{-1}$; the two form a basis of $\Cx^{2}$
since $z_0\neq\bar z_0$. (Explicitly $\zeta_+(h)=\tfrac t2+i\,\mathrm{sgn}(c)\tfrac{\sqrt{4-t^{2}}}2$;
for instance $\zeta_+(S)=i$ and $\zeta_+(U)=\zeta_+(ST)=e^{\pi i/3}$.)

\emph{Invariance.} Let $g=\smat{\alpha}{\beta}{\gamma}{\delta}\in\GL_2(\Q)^{+}$ with
$h':=ghg^{-1}\in\Ga_{\mathrm{fin}}$. Since $\det g>0$, $g$ acts on $\Hyp$ by M\"obius
transformations, $\operatorname{Im}(g\cdot z)=\det(g)\operatorname{Im}(z)/\abs{\gamma z+\delta}^{2}>0$,
and $g\,(z_0,1)^{\mathsf T}=(\gamma z_0+\delta)\,(g\cdot z_0,1)^{\mathsf T}$ with
$\gamma z_0+\delta\neq0$ ($z_0\notin\Rl$). Now $h'$ fixes $g\cdot z_0\in\Hyp$, so
$\zeta_+(h')$ is the eigenvalue of $h'$ on $(g\cdot z_0,1)^{\mathsf T}$, a non-zero
multiple of $g(z_0,1)^{\mathsf T}$; and
$h'\,g(z_0,1)^{\mathsf T}=g\,h\,(z_0,1)^{\mathsf T}=\zeta_+(h)\,g(z_0,1)^{\mathsf T}$.
Hence $\zeta_+(h')=\zeta_+(h)$ (trivially so if $h=\pm I$).

\emph{Characters.} Let $C=\langle h_0\rangle\le\Ga$ be finite cyclic. If $h_0=\pm I$ then
$C\subseteq\{\pm I\}$ and $\zeta_+|_C$, being $\pm I\mapsto\pm1$, is a faithful character.
Otherwise let
$z_0\in\Hyp$ be the fixed point of $h_0$; every $h\in C$ is a power of $h_0$, fixes $z_0$,
and its eigenvalue on $(z_0,1)^{\mathsf T}$ is $\zeta_+(h)$ (for $h\neq\pm I$ because $z_0$
is then \emph{its} fixed point in $\Hyp$; for $h=\pm I$ trivially). Eigenvalues on a
common eigenvector multiply, so $\zeta_+|_C$ is a homomorphism $C\to\mu_{\infty}$; it is
faithful because $\zeta_+(h)=1$ forces $h$ to act as $1$ on both basis vectors
$(z_0,1)^{\mathsf T}$, $(\bar z_0,1)^{\mathsf T}$ (the second with eigenvalue
$\zeta_+(h)^{-1}$), i.e.\ $h=I$. Thus $\chi_C$ is a faithful character.

\emph{Step 2: $e$, and $H^{2}(\Ga;\Z)\cong\Z/12$.} Claim~\ref{cl:amalgam}(2) with
$\cring=\Z$, $n=0$ --- $\VR=\Z$ with trivial action --- reads, in degrees $1$ to $2$,
\[
  H^{1}(\langle-I\rangle;\Z)\to H^{2}(\Ga;\Z)\xrightarrow{(\res,\res)}
  H^{2}(\langle S\rangle;\Z)\oplus H^{2}(\langle U\rangle;\Z)\xrightarrow{\res-\res}
  H^{2}(\langle-I\rangle;\Z),
\]
and $H^{1}(\langle-I\rangle;\Z)=\Hom(\Z/2,\Z)=0$. So $(\res,\res)$ is injective with image
$\ker(\res-\res)$. Through $\kappa$ (Step 0, naturality) $\res-\res$ is
$(\chi,\chi')\mapsto\chi|_{\langle-I\rangle}-\chi'|_{\langle-I\rangle}$ on
$\Hom(\langle S\rangle,\Q/\Z)\oplus\Hom(\langle U\rangle,\Q/\Z)\cong\Z/4\oplus\Z/6$, each
restriction being onto $\Hom(\langle-I\rangle,\Q/\Z)\cong\Z/2$ ($\Q/\Z$ is divisible), so
the kernel has order $24/2=12$; it contains $(\chi_{\langle S\rangle},\chi_{\langle U\rangle})$,
because both restrict to $\zeta_+|_{\langle-I\rangle}$, and that element has order
$\operatorname{lcm}(4,6)=12$ (the two characters are faithful, Step 1). Hence
$H^{2}(\Ga;\Z)\cong\ker(\res-\res)\cong\Z/12$, and $e$ --- the unique preimage of
$(\chi_{\langle S\rangle},\chi_{\langle U\rangle})$ --- generates it.

\emph{Step 3: $\kappa_C(\res^{\Ga}_Ce)=\chi_C$ for every finite cyclic $C\le\Ga$.} First,
\emph{$C$ fixes a vertex of $\mathcal T$} \STD{\cite[Ch.~I, \S4.3]{serre_trees}, a
finite group acting without inversion on a tree fixes a vertex; the pointer is not
verified against the text (rule~\ref{A2}) and the argument is given}: the action is
without inversion (the two ends of an edge have different types); pick a vertex $v$,
let $\mathcal T_0$ be the finite subtree spanned by the finite orbit $Cv$, which is
$C$-invariant; the set of leaves of a finite tree is invariant under its automorphisms,
and a finite tree with at least three vertices remains a non-empty tree when all its
leaves are deleted (the interior vertices of a path have degree $\ge2$); iterating, $C$
preserves a subtree with one vertex, or with two vertices joined by an edge, whose ends
it then fixes since it does not invert. So $C\le\Ga_v=\gamma D\gamma^{-1}$ for some
$\gamma\in\Ga$ and $D\in\{\langle S\rangle,\langle U\rangle\}$. Put
$D':=\gamma D\gamma^{-1}$ and let $c_\gamma\colon\Ga\to\Ga$, $h\mapsto\gamma^{-1}h\gamma$,
which maps $D'$ isomorphically onto $D$. By \ref{Epair}, $c_\gamma^{\ast}e=e$ and
$\res^{\Ga}_{D'}e=\res^{\Ga}_{D'}(c_\gamma^{\ast}e)=(c_\gamma|_{D'})^{\ast}(\res^{\Ga}_{D}e)$,
the two composites $D'\to\Ga$ being the same homomorphism. Applying $\kappa_{D'}$ and
Step 0, $\kappa_{D'}(\res^{\Ga}_{D'}e)=\chi_D\circ c_\gamma|_{D'}=\zeta_+(\gamma^{-1}(\cdot)\gamma)|_{D'}=\zeta_+|_{D'}=\chi_{D'}$,
by Step 2 and the invariance of Step 1 (with $\gamma^{-1}\in\Ga\subseteq\GL_2(\Q)^{+}$).
Restricting further to $C\le D'$ (Step 0 again) gives $\kappa_C(\res^{\Ga}_Ce)=\chi_{D'}|_C=\chi_C$.
Since $\chi_C$ is faithful, $\res^{\Ga}_Ce$ generates $H^{2}(C;\Z)\cong\Z/\abs C$. This
proves (1).

\emph{Step 4: five formal consequences of \ref{Ecup}--\ref{Ecor}.} Let $G$ be a group,
$u\in H^{p}(G;\Z)$.
\begin{enumerate}[label=(P\arabic*)]
\item\label{Pnat} For a compatible pair $(\alpha\colon G'\to G,f\colon B\to B')$ and
  $x\in H^{q}(G;B)$: $(\alpha,f)^{\ast}(u\cup x)=\alpha^{\ast}u\cup(\alpha,f)^{\ast}x$,
  where $\alpha^{\ast}=(\alpha,\id_\Z)^{\ast}$. On cochains, by \eqref{eq:cupcochain} and
  $f(\alpha(g')b)=g'f(b)$:
  $f\bigl(u(\alpha g'_1,\dots,\alpha g'_p)\,\alpha(g'_1\cdots g'_p)\,\varphi(\alpha g'_{p+1},\dots)\bigr)
  =u(\alpha g'_1,\dots)\,(g'_1\cdots g'_p)\,f\varphi(\alpha g'_{p+1},\dots)$. In particular
  $\res^{G}_{H}(u\cup x)=\res^{G}_{H}u\cup\res^{G}_{H}x$.
\item\label{Pconn} \emph{Connecting homomorphisms.} For a short exact sequence
  $0\to B'\to B\to B''\to0$ of $G$-modules and $x\in H^{q}(G;B'')$:
  $\delta(u\cup x)=(-1)^{p}\,u\cup\delta x$. For a short exact sequence
  $0\to A'\to A\to A''\to0$ of $G$-modules that stays exact after $\otimes_{\Z}B$ (e.g.\
  one of free abelian groups), $w\in H^{p}(G;A'')$ and $x\in H^{q}(G;B)$:
  $\delta(w\cup x)=\delta w\cup x$ in $H^{p+q+1}(G;A'\otimes B)$. Both by the Leibniz
  rule: lift the cocycle of the class being shifted to a cochain of the middle term, cup
  it with a cocycle of the other class, and take $d$; the term with $d$ on the cocycle
  vanishes, leaving $(-1)^{p}u\cup d\varphi$, resp.\ $dw\cup\varphi$.
\item\label{Pshapiro} \emph{Shapiro.} For $H\le G$, $X\in H^{q}(G;\mathrm{Coind}^{G}_{H}B)$:
  $\mathrm{sh}(u\cup X)=\res^{G}_{H}u\cup\mathrm{sh}(X)$. Indeed
  $\mathrm{sh}(u\cup X)=\mathrm{ev}_{\ast}\res(u\cup X)=\mathrm{ev}_{\ast}(\res u\cup\res X)
  =\res u\cup\mathrm{ev}_{\ast}\res X$ by \ref{Pnat} and naturality in $B$ (\ref{Ecup}).
\item\label{Pproj} \emph{Projection formula.} For $[G:H]<\infty$ and
  $y\in H^{q}(H;B)$: $\cor^{G}_{H}(\res^{G}_{H}u\cup y)=u\cup\cor^{G}_{H}y$. Indeed
  \ref{Pshapiro} with $X=\mathrm{sh}^{-1}y$ gives $\mathrm{sh}^{-1}(\res u\cup y)=u\cup\mathrm{sh}^{-1}y$,
  and $\cor(\res u\cup y)=\mathrm{tr}_{\ast}(u\cup\mathrm{sh}^{-1}y)=u\cup\mathrm{tr}_{\ast}\mathrm{sh}^{-1}y=u\cup\cor y$
  by \ref{Ecor} and naturality in $B$.
\item\label{Pfree} For $H$ finite and any $H$-module $M$, $H^{q}(H;\Z[H]\otimes_{\Z}M)=0$
  for $q\ge1$, the tensor product carrying the diagonal action: $h\otimes m\mapsto
  h\otimes h^{-1}m$ is an $H$-isomorphism onto $\Z[H]\otimes M$ with $H$ acting on the left
  factor only, which is $\Ind^{H}_{1}M\cong\mathrm{Coind}^{H}_{1}M$ (\ref{Ecor}), whose
  cohomology vanishes in positive degrees (\ref{Eshapiro}).
\end{enumerate}

\emph{Step 5: cyclic periodicity.} Let $C=\langle t\rangle$ be cyclic of order $m$,
$N:=1+t+\dots+t^{m-1}\in\Z[C]$, $I:=\ker(\aug\colon\Z[C]\to\Z)$. Then
\begin{equation}\label{eq:cyc2ext}
  0\longrightarrow\Z\xrightarrow{\ 1\mapsto N\ }\Z[C]\xrightarrow{\ t-1\ }\Z[C]
   \xrightarrow{\ \aug\ }\Z\longrightarrow0
\end{equation}
is exact: $\ker\aug=I=(t-1)\Z[C]$ since $\Z[C]=\Z[t]/(t^{m}-1)$; $\ker(t-1)=\Z[C]^{C}=\Z N$
(an element $\sum a_jt^{j}$ fixed by $t$ has all $a_j$ equal); and $1\mapsto N$ is
injective with image $\Z N$. This is the first segment of the periodic resolution of
\cite[Ch.~III.1]{brown} used in the proof of Proposition~\ref{prop:bdrycohom}. Split it
as $0\to I\to\Z[C]\to\Z\to0$ and $0\to\Z\to\Z[C]\to I\to0$, with connecting maps
$\delta_1\colon H^{0}(C;\Z)\to H^{1}(C;I)$ and $\delta_2\colon H^{1}(C;I)\to H^{2}(C;\Z)$,
and put $u_0:=\delta_2\delta_1(1)\in H^{2}(C;\Z)$. Let $M$ be any $C$-module. Both short
exact sequences consist of free abelian groups, so they stay exact after $\otimes_{\Z}M$;
write $\delta_1'\colon H^{q}(C;M)\to H^{q+1}(C;I\otimes M)$ and
$\delta_2'\colon H^{q+1}(C;I\otimes M)\to H^{q+2}(C;M)$ for the connecting maps of the
tensored sequences (diagonal actions). By \ref{Pconn} (first-variable form, twice) and
$1\cup x=x$,
\[
  u_0\cup x=\delta_2\bigl(\delta_1(1)\bigr)\cup x=\delta_2'\bigl(\delta_1(1)\cup x\bigr)
  =\delta_2'\delta_1'(1\cup x)=\delta_2'\delta_1'(x)\qquad(x\in H^{q}(C;M)).
\]
By \ref{Pfree} the terms $H^{j}(C;\Z[C]\otimes M)$ flanking $\delta_1'$ and $\delta_2'$
vanish for $j\ge1$, so $\delta_2'$ is an isomorphism for every $q\ge0$, and $\delta_1'$
is an isomorphism for $q\ge1$ and, for $q=0$, surjective with kernel the image of
$H^{0}(C;\Z[C]\otimes M)\to H^{0}(C;M)=M^{C}$ under $\aug\otimes\id$. That image is
$N\cdot M$: an element $\sum_{c}c\otimes m_c$ is $C$-invariant iff $m_c=c\,m_1$ for all
$c$, and $\aug\otimes\id$ sends $\sum_c c\otimes cm_1$ to $Nm_1$. Hence
\emph{$u_0\cup-\colon H^{q}(C;M)\to H^{q+2}(C;M)$ is an isomorphism for $q\ge1$, and for
$q=0$ a surjection $M^{C}\to H^{2}(C;M)$ with kernel $NM$.}
Taking $M=\Z$, $q=0$: $u_0\cup k=k\,u_0$ and the map is onto, so $u_0$ generates
$H^{2}(C;\Z)$, which is $\cong\Z/m$ by Step 0; thus $u_0$ has order $m$. If $u$ is any
generator, $u=a\,u_0$ with $\gcd(a,m)=1$, and since $m$ kills $H^{j}(C;M)$ for $j\ge1$
\STD{\cite[Ch.~III.10.2]{brown}} multiplication by $a$ is invertible there; so
\emph{cup product with any generator of $H^{2}(C;\Z)$ is an isomorphism
$H^{q}(C;M)\to H^{q+2}(C;M)$ for every $q\ge1$ and every $M$, and for $q=0$ a surjection
with kernel $NM$}.

\emph{Step 6: proof of (2).} For $D\in\{\langle S\rangle,\langle U\rangle,\langle-I\rangle\}$
put $P_D:=\Z[\Ga/D]$. Every term of \eqref{eq:treeses} is a free abelian group, so the
sequence is $\Z$-split and $\Hom_{\Z}(-,V)$ --- with the diagonal action
$(\gamma f)(x)=\gamma f(\gamma^{-1}x)$ --- turns it into a short exact sequence of
$\Z[\Ga]$-modules
\begin{equation}\label{eq:treesesdual}
  0\to V\to\Hom_{\Z}(P_{\langle S\rangle},V)\oplus\Hom_{\Z}(P_{\langle U\rangle},V)
   \to\Hom_{\Z}(P_{\langle-I\rangle},V)\to0 ,
\end{equation}
whose long exact cohomology sequence we call $(\ast)$. For each $D$,
$\Phi_D\colon\Hom_{\Z}(P_D,V)\to\mathrm{Coind}^{\Ga}_{D}V$, $\Phi_D(f)(\gamma):=\gamma f(\gamma^{-1}D)$,
is a $\Ga$-isomorphism (it is $D$-equivariant in the argument, $\Ga$-equivariant, and
inverted by $\varphi\mapsto(\gamma D\mapsto\gamma\varphi(\gamma^{-1}))$, all by direct
substitution), so by \ref{Eshapiro} the terms of $(\ast)$ are $H^{q}(\Ga;V)$,
$H^{q}(\langle S\rangle;V)\oplus H^{q}(\langle U\rangle;V)$ and $H^{q}(\langle-I\rangle;V)$:
those of \eqref{eq:mvamalgam}, of which $(\ast)$ is the coefficient-sequence form (under
$\Ext^{q}_{\Z[\Ga]}(P,V)\cong H^{q}(\Ga;\Hom_{\Z}(P,V))$ for $\Z$-free $P$ it \emph{is}
\eqref{eq:mvamalgam}; nothing below needs that identification). Cup product with $e$
acts on every term of $(\ast)$ and commutes with every map of it: with the maps induced
by \eqref{eq:treesesdual} by naturality in the coefficients (\ref{Ecup}), with the
connecting homomorphism by \ref{Pconn} with $p=2$, sign $(-1)^{2}=+1$. On
$H^{q}(\Ga;\Hom_{\Z}(P_D,V))$ it is transported by $\Phi_{D\ast}$ and $\mathrm{sh}$ to
$\res^{\Ga}_{D}e\cup-$ on $H^{q}(D;V)$ (\ref{Pshapiro}), where $\res^{\Ga}_De$ generates
$H^{2}(D;\Z)$ by (1); so by Step 5 the vertical map is an isomorphism on these terms in
every degree $q\ge1$, and in degree $0$ a surjection $V^{D}\to H^{2}(D;V)$ with kernel
$N_DV$. The five lemma on
\[
  H^{q-1}(\oplus)\to H^{q-1}(\langle-I\rangle;V)\to H^{q}(\Ga;V)\to H^{q}(\oplus)\to H^{q}(\langle-I\rangle;V)
\]
(with $H^{j}(\oplus):=H^{j}(\langle S\rangle;V)\oplus H^{j}(\langle U\rangle;V)$), mapped
by $e\cup-$ to the same window two degrees up, gives that $e\cup-\colon H^{q}(\Ga;V)\to H^{q+2}(\Ga;V)$
is an isomorphism as soon as the four outer maps are, i.e.\ for $q-1\ge1$: \textbf{this
is where $q\ge2$ enters}. For $q=1$ the window's left-hand terms are in degree $0$, where
the vertical maps are only surjective; the four lemma still gives surjectivity of
$e\cup-\colon H^{1}(\Ga;V)\to H^{3}(\Ga;V)$, and a diagram chase gives the kernel: if
$e\cup x=0$ then the image of $x$ in $H^{1}(\oplus)$ dies under an injective map, so
$x=\partial y$ with $y\in V^{-I}$; then $\partial(e\cup y)=0$, so $e\cup y=(\res-\res)(e\cup z)$
for some $z\in V^{S}\oplus V^{U}$ (surjectivity in degree $0$), so
$y-(\res-\res)z\in\ker(e\cup-)=N_{-I}V$ and $x=\partial\bigl(y-(\res-\res)z\bigr)\in\partial(N_{-I}V)$;
conversely $e\cup\partial(N_{-I}V)=\partial(e\cup N_{-I}V)=0$. Naturality in $V$ and
compatibility with restriction are \ref{Ecup} and \ref{Pnat}; $\cring$-linearity for
$V=\VR$ is naturality applied to the endomorphisms $v\mapsto rv$. The statements about
$V=\VR$ follow: $N_{-I}\VR=(1+\rhon(-I))\VR=(1+(-1)^{n})\VR$
(Proposition~\ref{prop:symconv}(2)), and for $6\in\cring^{\times}$, $H^{3}(\Ga;\VR)=0$ by
Claim~\ref{cl:amalgam}(3) while
$\dim H^{1}(\Ga;\VR)=\dim H^{0}(\Ga;\VR)+\dim\VR-\dim\VR^{S}-\dim\VR^{U}=\dim H^{0}+3-1-1\ge1$
for $\cring=\Fl$, $\ell\ge5$, $n=2$ (the eigenvalues of $\rhon(S)$ are $-1,1,-1$, those
of $\rhon(U)$ are $1$ and the two primitive cube roots of unity, both matrices semisimple
since $4,6\in\Fl^{\times}$).

\emph{Step 7: proof of (3a).} $[\Ga:\Ga'']<\infty$ (Proposition~\ref{prop:cosets}).
Restricted to $\Z[\Ga'']$, each $P_D$ of \eqref{eq:treeses} is the finite direct sum, over
the $\Ga''$-orbits $\Ga''x$ in $\Ga/D$, of the permutation modules
$\Z[\Ga''x]\cong\Z[\Ga''/\Ga''_x]$ with $\Ga''_x=\Ga''\cap\gamma D\gamma^{-1}$ for
$x=\gamma D$ --- the stabiliser of the vertex or edge $x$ of $\mathcal T$ in $\Ga''$, a
subgroup of a cyclic group, hence finite cyclic. Applying $\Ext^{\bullet}_{\Z[\Ga'']}(-,\Z)$
to \eqref{eq:treeses} exactly as in the proof of Claim~\ref{cl:amalgam}(2), with
Shapiro's lemma (\ref{Eshapiro}, \ref{Ecor}) on each summand, gives an exact segment
\[
  \bigoplus_{\text{edges }x}H^{1}(\Ga''_x;\Z)\longrightarrow H^{2}(\Ga'';\Z)
  \xrightarrow{\ (\res^{\Ga''}_{\Ga''_x})_x\ }\bigoplus_{\text{vertices }x}H^{2}(\Ga''_x;\Z),
\]
the middle map being restriction for the same reason the maps of \eqref{eq:mvamalgam}
are: Shapiro's isomorphism is restriction followed by evaluation at the base coset
(\ref{Eshapiro}), and the augmentation $\Z[\Ga''/\Ga''_x]\to\Z$ becomes, under $\Phi$,
the unit $V\to\mathrm{Coind}\,V$, whose evaluation at the base coset is the identity.
Since $H^{1}(\text{finite};\Z)=0$, \textbf{$H^{2}(\Ga'';\Z)$ is detected on the finite
cyclic subgroups $\Ga''_x$}. It therefore suffices to show
$\res^{\Ga''}_{C}(\iota_g^{\ast}e)=\res^{\Ga''}_{C}(\res^{\Ga}_{\Ga''}e)$ for every finite
cyclic $C\le\Ga''$. Put $C':=g^{-1}Cg=\iota_g(C)\le\Ga$, finite cyclic. By functoriality
(\ref{Epair}), $\res^{\Ga''}_{C}\iota_g^{\ast}e=(\iota_g|_C)^{\ast}e=(\iota_g|_C\colon C\to C')^{\ast}\res^{\Ga}_{C'}e$,
and by Step 0 and (1) applied to $C'$,
$\kappa_C\bigl((\iota_g|_C)^{\ast}\res^{\Ga}_{C'}e\bigr)=\chi_{C'}\circ\iota_g|_C=\zeta_+(g^{-1}(\cdot)g)|_C=\zeta_+|_C=\chi_C=\kappa_C(\res^{\Ga}_Ce)$,
the middle equality being the invariance of Step 1 for $g^{-1}\in\GL_2(\Q)^{+}$. As
$\kappa_C$ is injective and $\res^{\Ga}_C=\res^{\Ga''}_C\res^{\Ga}_{\Ga''}$, this is the
required equality.

\emph{Step 8: proof of (3b) and the consequences.} Let $x\in H^{q}(\Ga;\VR)$. By
\ref{Pnat} for $\res^{\Ga}_{\Ga'}$, $\res^{\Ga}_{\Ga'}(e\cup x)=\res^{\Ga}_{\Ga'}e\cup\res^{\Ga}_{\Ga'}x$.
The map $(\conj_g)_{\ast}$ is $(\alpha,f)^{\ast}$ for the compatible pair
$\alpha=\conj_{g^{-1}}\colon\Ga''\to\Ga'$, $f=\rhon(g)$ of \eqref{eq:corconjres}, so by
\ref{Pnat} again
$(\conj_g)_{\ast}(\res^{\Ga}_{\Ga'}e\cup y)=\alpha^{\ast}(\res^{\Ga}_{\Ga'}e)\cup(\conj_g)_{\ast}y
=\iota_g^{\ast}e\cup(\conj_g)_{\ast}y=\res^{\Ga}_{\Ga''}e\cup(\conj_g)_{\ast}y$, using
$\mathrm{incl}_{\Ga'}\circ\alpha=\iota_g$ and (3a). Finally \ref{Pproj} for
$\Ga''\le\Ga$ gives $\cor^{\Ga}_{\Ga''}(\res^{\Ga}_{\Ga''}e\cup z)=e\cup\cor^{\Ga}_{\Ga''}z$.
Composing the three, $\Tdc{g}(e\cup x)=e\cup\Tdc{g}x$. For $q\ge2$, $e\cup-$ is an
isomorphism by (2) and commutes with every $\Tdc{g}$ by (3b), hence is an isomorphism of
modules over the algebra they generate; the statement about $\Amat{q}$ is the same
sentence in a basis.
\end{proof}

\begin{remark}[What Lemma~\ref{lem:bulkperiodic} does and does not settle for \ref{oc1c}]\label{rmk:bulkperiodicscope}
\begin{enumerate}
\item \textbf{Ordinary cohomology, no Farrell theory.} The lemma is a statement about
  $H^{\bullet}(\Ga;-)$. Farrell cohomology $\widehat H^{\bullet}(\Ga;-)$
  (\cite[Ch.~X]{brown}) agrees with it only in degrees \emph{above}
  $\operatorname{vcd}(\SL_2(\Z))=1$ and is neither used nor asserted about; the
  periodicity is proved from Claim~\ref{cl:amalgam}(2), the cup product and the five
  lemma, and is exact about its lower edge: it is an isomorphism from $q=2$, a surjection
  at $q=1$, and says nothing at $q=0$. The reader may recognise $e$ as the Euler class of
  the central extension $1\to\Z\to B_3\to\SL_2(\Z)\to1$ (the universal covering of
  $\SL_2(\Rl)$ restricted to $\Ga$), whose restriction to $\langle S\rangle$ is
  $0\to\Z\xrightarrow{4}\Z\to\Z/4\to0$, matching $\zeta_+(S)=i$; this document neither
  verifies nor uses that identification --- $e$ is \emph{defined} in clause (1), and its
  $\De$-invariance is proved there, not inherited from $\GL_2(\Rl)^{+}$.
\item \textbf{What is now a theorem in \ref{oc1c}'s hypothesis.} Its bulk half, in a
  form stronger than it asked for: $H^{q}(\Ga;\VR)\cong H^{q+2}(\Ga;\VR)$ for $q\ge2$,
  every $\cring$ and every $n$, \emph{as Hecke modules} for every $\Tdc{g}$, $g\in\De$,
  by one canonical isomorphism. On the free $\Ga$-model of \S\ref{sec:orbifoldmodel}
  the bulk column of \eqref{eq:coneles} in degrees $\ge4$ is therefore a transport of
  degrees $2$ and $3$, Hecke matrices included. On the shipped $7$-vertex model at
  $6\in\cring^{\times}$ the lemma is vacuous ($H^{q}=0$ for $q\ge2$,
  Claim~\ref{cl:amalgam}(3)).
\item \textbf{What remains open in \ref{oc1c}, precisely.} Three things, none of them
  asserted here.
  \begin{enumerate}[label=(\alph*)]
  \item \emph{The boundary.} Lemma~\ref{lem:ell2unbounded} gives the same
    $\cring$-module $H^{q}(\partial;\VR)$ in every degree $q\ge1$ at $\ell=2$, and
    Proposition~\ref{prop:bdrycohom} exhibits $H^{q}(\Gbdry;\VR)$ for every $\cring$ as an
    extension of two pieces that are $2$-periodic for $q\ge2$ (Step 5 of the proof); but
    no \emph{map} $H^{q}(\partial;\VR)\to H^{q+2}(\partial;\VR)$ is proved to be an
    isomorphism. The candidate is $\res^{\Ga}_{\Gbdry}e\cup-$; its compatibility with
    $\rres^{q}=\res^{\Ga}_{\Gbdry}$ is automatic ((P1) of the proof:
    $\rres^{q}(e\cup x)=\res^{\Ga}_{\Gbdry}e\cup\rres^{q}x$), and its
    $T^{\partial}_{p}$-equivariance would follow by the argument of clause (3) run on
    $\Gbdry$ with the boundary coset system of Proposition~\ref{prop:bdrycosets}, given
    the boundary analogue of (3a). Nothing of this is proved here.
  \item \emph{The cone.} Given (a) with an isomorphism from some degree $q_0^{\partial}$
    on, the five lemma applied to \eqref{eq:coneles} would give
    $\Hcpt{q}\cong\Hcpt{q+2}$ for $q\ge\max(3,q_0^{\partial}+1)$, \emph{provided} cup
    product with $e$ acts on $\Hcpt{\bullet}=H^{\bullet}_{\Ga}(Y,\partial Y;\VR)$
    (Definition~\ref{def:borelmodel}; Theorem~\ref{thm:specialisation}(2) where
    $6\in\cring^{\times}$) compatibly with the maps of that sequence. That
    is the standard module structure of the long exact sequence of a pair over
    $H^{\bullet}(\Ga;\Z)$, and it is not written out in this document.
  \item \emph{Hecke equivariance on the boundary and on the cone}, by the same
    projection-formula route as clause (3); not proved here.
  \end{enumerate}
  If (a) holds from $q_0^{\partial}=1$ --- the value Lemma~\ref{lem:ell2unbounded}
  suggests at $\ell=2$ --- then every group of \eqref{eq:coneles} in degree $\ge5$ is a
  transport of one in degree $\le4$, which is why the policy of \ref{open:ell2} recommends
  $q_{\max}\ge4$ at $\ell\in\{2,3\}$ where the stabilisation is to be \emph{observed}: a
  run with $q_{\max}<4$ cannot exhibit even the first bulk isomorphism
  $H^{2}(\Ga;\VR)\cong H^{4}(\Ga;\VR)$, and the cone's first, $\Hcpt{3}\cong\Hcpt{5}$,
  needs $q_{\max}\ge5$.
\end{enumerate}
\end{remark}

\begin{openitem}[OC-1: the $\ell=2$ truncation policy --- DECIDED 2026-09-16]\label{open:ell2}
\textbf{Decided by the owner on 2026-09-16: the policy is \ref{oc1a} below, with the three
riders stated after the box} (owner decision ``OC-1: adoption per brief''; the brief is
\file{docs/implementation/oc1\_decision\_brief.md}, the record is
\file{docs/implementation/hecke\_extensions\_owner\_decisions.md}, entry OC-1). What
follows is therefore a \emph{policy} --- a definition-like statement, which carries no
status tag of its own (\S\ref{sec:status}, Appendix~\ref{app:status}); the truncation
lemma inside it keeps its own tag. From v0.14.0 to v1.2.0 this item carried the tag
\code{[open]} and the header ``reserved by the owner; a tagged working assumption, not a
decision''; the label and the content of the boxed statement are unchanged
(rules~\ref{L1}, \ref{L2}), its status is. At $\ell\in\{2,3\}$ the complexes of \S\ref{sec:orbifoldmodel} are
unbounded above (Remark~\ref{rmk:ell2window}(3)), so an implementation computes a
truncation, and the contract must fix four things: the bound $q_{\max}$, its place in
the record schema, the certificate obligations at the top of the truncation, and what
``exactness'' asserts there. The enabling fact is elementary and is what makes any
truncation legitimate at all:

\smallskip
\noindent\emph{Truncation lemma.} \PROVED{} For a cochain complex $C^{\bullet}$ and
$q\le q_{\max}$, $H^{q}(C^{\bullet})$ depends only on $C^{\le q_{\max}+1}$: it is
$\ker d^{q}/\im d^{q-1}$, and $d^{q}$ is defined on $C^{q}$ with values in $C^{q+1}$,
$q+1\le q_{\max}+1$. \emph{The same is false for $q=q_{\max}+1$}, where the truncated
complex has $\ker d^{q_{\max}+1}=C^{q_{\max}+1}$ and computes a quotient that is in
general strictly larger than $H^{q_{\max}+1}$.

\smallskip
\noindent\textbf{The policy taken here (decided 2026-09-16; the box is verbatim the v0.14.0
working assumption, except that its record sentence says ``per complex'' since v1.4.0 ---
\S\ref{sec:recordschema} item~5, owner confirmation 1 of the slot-2 plan).}

\smallskip
\noindent\fbox{\parbox{0.95\linewidth}{\textbf{[OC-1: $q_{\max}$ is an explicit,
recorded parameter of a run, defaulting to $2$.]} The free-$\Ga$-model complexes of
\S\ref{sec:orbifoldmodel} --- bulk, boundary and cone --- are built in degrees
$\le q_{\max}+1$ and the twisted chain data $\tau$, $\taub$, $\etapair$ are certified in
every degree in which they are built. Cohomology is \textbf{reported only in degrees
$\le q_{\max}$}; degree $q_{\max}+1$ is computed as scaffolding and is never reported as
a cohomology group. The long exact sequence is asserted \textbf{exact at every spot all
three of whose terms are reported}, and asserted to be a \emph{complex} (consecutive
composites zero) at the remaining spots inside the computed range. The record carries,
\emph{per complex it is about} (v1.4.0), $q_{\max}$, the flag \code{truncated}, the
derived/set marker and that complex's list of degrees; a record with
\code{truncated = true} may not report a vanishing above $q_{\max}$. At
$6\in\cring^{\times}$ the model complexes are already finite, \code{truncated = false},
and the policy is vacuous.}}

\smallskip
\noindent Spelled out, ``exact at every spot all three of whose terms are reported''
means, for the sequence
$\cdots\to\Hcpt{q}\to H^{q}(\Ga)\xrightarrow{\rres^{q}}H^{q}(\partial)\to\Hcpt{q+1}\to\cdots$:
exactness is asserted at $\Hcpt{q}$ and at $H^{q}(\Ga)$ for $q\le q_{\max}$, and at
$H^{q}(\partial)$ for $q\le q_{\max}-1$; at $H^{q_{\max}}(\partial)$ only
$\rres^{q_{\max}}\circ(\text{previous})=0$ is asserted, because the next term
$\Hcpt{q_{\max}+1}$ is not reported. Likewise $d\circ d=0$ is checkable in degrees
$\le q_{\max}-1$ and the cochain-map identity $d\Phig=\Phig d$ in degrees $\le q_{\max}$.

\smallskip
\noindent\textbf{Three riders (owner decision, 2026-09-16).}
\begin{enumerate}[label=(\roman*)]
\item \textbf{$q_{\max}$ is derived on a finite model and settable only where
  \code{truncated = true}.} On a model whose complexes are finite --- the $7$-vertex
  model and every refinement of it at $6\in\cring^{\times}$ --- $q_{\max}$ is the top
  degree of the complex, \emph{not} a choice, and the record marks it as derived; a user
  may set $q_{\max}$ only on a truncated free model, where it is the bound of the box.
  $q_{\max}$ is distinct from the \emph{reported-degree list}, which is an independent
  filter on the degrees a record prints (already present in the campaign specification)
  and may be a proper subset of $\{0,\dots,q_{\max}\}$; a record with
  \code{truncated = true} may still not report a vanishing above $q_{\max}$.
\item \textbf{The bulk periodicity above the virtual cohomological dimension is a
  theorem, not a citation}: Lemma~\ref{lem:bulkperiodic}, with its Hecke equivariance,
  replaces the ``standard route, not verified'' that Remark~\ref{rmk:ell2window}(2)
  carried until v1.2.0; what it leaves of \ref{oc1c} is Remark~\ref{rmk:bulkperiodicscope}(3).
\item \textbf{Obligation~\ref{obl:coneexact}(iv) asserts composites zero only where both
  terms are reported}, i.e.\ it carries the box's own qualification (``at
  $H^{q_{\max}}(\partial)$ only $\rres^{q_{\max}}\circ(\text{previous})=0$'') instead of
  its former blanket ``in every computed degree''.
\end{enumerate}

\smallskip
\noindent\emph{Non-normative recommendation.} At $\ell\in\{2,3\}$ a run with
$q_{\max}<4$ cannot exhibit the stabilisation of Lemma~\ref{lem:bulkperiodic} at all ---
its first instance, $H^{2}(\Ga;\VR)\cong H^{4}(\Ga;\VR)$, needs degree $4$ reported and
degree $5$ built, i.e.\ $q_{\max}\ge4$; the second, $H^{3}\cong H^{5}$, and the cone's
first (Remark~\ref{rmk:bulkperiodicscope}(3)(b), from $\Hcpt{3}$) need $q_{\max}\ge5$
--- so $q_{\max}\ge4$ is the recommended minimum there. The default stays $2$: it is not a knob on the shipped finite models (rider (i)),
and at $\ell\ge5$ the extra degrees are known to vanish (Claim~\ref{cl:amalgam}(3)).

\smallskip
\noindent\textbf{The alternatives: \ref{oc1a} taken, \ref{oc1c} kept reachable.}
\begin{enumerate}[label=(OC-1\alph*)]
\item\label{oc1a} \emph{Taken (owner decision 2026-09-16):} $q_{\max}$ a recorded run
  parameter, default $2$, with the riders above.
\item\label{oc1b} $q_{\max}=2$ fixed by the contract. Rejected: it writes the current
  six-term window into a normative statement, which is precisely what OD-T1 and
  \S\ref{sec:orbifoldmodel} are being kept open against, and it would have to be undone
  by a normative amendment the first time anyone asks for $H^{3}$.
\item\label{oc1c} $q_{\max}$ determined by a \emph{stabilisation theorem}: by
  Lemma~\ref{lem:ell2unbounded} the boundary cohomology is, as an $\cring$-module,
  \emph{the same in every degree $q\ge1$} at $\ell=2$; the bulk \emph{is} periodic
  above the virtual cohomological dimension, Hecke-equivariantly ---
  Lemma~\ref{lem:bulkperiodic}: $e\cup-\colon H^{q}(\Ga;\VR)\xrightarrow{\sim}H^{q+2}(\Ga;\VR)$
  for $q\ge2$, commuting with every $\Tdc{g}$, so the bulk half of this hypothesis,
  Hecke action included, is a theorem since v1.3.0; and if a boundary periodicity
  isomorphism compatible with $\rres$ and with $T^{\partial}_{p}$, and the cone
  periodicity it implies through \eqref{eq:coneles}, are proved as well, then a finite
  computation determines every degree and $q_{\max}$ is a \emph{theorem}, not a
  parameter. \textbf{What remains unproved is exactly Remark~\ref{rmk:bulkperiodicscope}(3)}:
  the boundary isomorphism (the candidate $\res^{\Ga}_{\Gbdry}e\cup-$ is named there),
  the module structure of \eqref{eq:coneles} over $H^{\bullet}(\Ga;\Z)$ that the five
  lemma needs on the cone, and the Hecke equivariance on boundary and cone. This is a
  strengthening of \ref{oc1a} and not an alternative to it: adopting \ref{oc1a} changes
  nothing that \ref{oc1c} would later add.
\item\label{oc1d} No truncation: decline to compute the cone at $\ell\in\{2,3\}$ and
  report only the model objects of Remark~\ref{rmk:ell2window}(1). This is the status quo
  and remains available; it is a refusal, not a policy, and it forecloses
  \S\ref{sec:orbifoldmodel}'s whole purpose.
\end{enumerate}

\smallskip
\noindent\textbf{Decision, and why.} \ref{oc1a}, as recommended from v0.14.0 and decided
by the owner on 2026-09-16. It is the only one of the four that is correct for every
$\ell$, degree-generic as \S\ref{sec:orbifoldmodel} and the implementation's data model
require, vacuous exactly where the shipped $6\in\cring^{\times}$ pipeline lives, and ---
the decisive point --- \textbf{the one that keeps \ref{oc1c} free}. A truncation bound
that is a recorded parameter can be replaced by a proved bound without amending a single
statement; a truncation bound written into the contract cannot. The cost of \ref{oc1a}
is that a record must carry three extra fields and that ``$H^{q}=0$ for $q>2$'' stops
being sayable at $\ell\in\{2,3\}$, which is not a cost but a correction. The decision
leaves \ref{oc1c} exactly as reachable as before, and nearer: its bulk half is
Lemma~\ref{lem:bulkperiodic}, so \ref{oc1b} would have written a normative bound that a
theorem already half supersedes.
\end{openitem}

\subsection*{Where the \ref{open:ell2} tag is used}\label{sec:oc1sites}

\noindent
So that the policy can be changed mechanically, these are every site in this document
that depends on it --- items (1)--(6) on its \emph{content}, item (7) on its
\emph{status}. Changing \ref{open:ell2} means changing these and nothing else; the
adoption of 2026-09-16 (v1.3.0) changed exactly (1), (3), (6) and (7).
\begin{enumerate}
\item \ref{open:ell2} itself --- the boxed policy, the spelled-out exactness clause,
  the three riders, and the decision.
\item Remark~\ref{rmk:ell2window}(3), the sentence naming the free
  $\Ga$-model as the object a truncation applies to.
\item Obligation~\ref{obl:taubdry}(v.c), (vi)--(vii) and, forward,
  Obligation~\ref{obl:coneexact}: ``in every computed degree'' is the phrase that reads
  $q_{\max}$ (since v1.4.0 clause (v) is split into (v.a)--(v.c), and only (v.c) carries
  the phrase; (v.a) and (v.b) name their degree), and the exactness assertion of the
  latter is the one \ref{open:ell2} narrows. In \S\ref{sec:cone} also Definition~\ref{def:coneaction}, stated for all
  $q\in\Z$ and read in degrees $\le q_{\max}$ on a truncated model, and
  Remark~\ref{rmk:conetorsion}, which names these two sites.
\item \S\ref{sec:orbifoldmodel}: \ref{def:borelmodel} and
  \ref{thm:specialisation} are stated in a degree range, not in degrees
  $\{0,1,2\}$ (written in v0.17.0; \ref{def:borelmodel}'s $q_{\max}$ is this
  assumption's parameter), and Remark~\ref{rmk:orbifoldrefine}(4) --- the sentence
  saying that a truncation policy is needed whatever the triangulation --- is that
  chapter's third dependent site.
\item \S\ref{sec:six}: Claim~\ref{cl:les} and \eqref{eq:les} --- the six-term
  presentation is the $q_{\max}=2$, $6\in\cring^{\times}$ instance, and
  Remark~\ref{rmk:ell2window} says so.
\item The record schema of \S\ref{sec:implementation}: the four per-complex fields
  $q_{\max}$, \code{truncated}, the derived/set marker and the complex's degree list
  (per complex since v1.4.0, \S\ref{sec:recordschema} item~5), and the job-level
  reported-degree selection.
\item \textbf{The sites that assert the item's \emph{status}} rather than read its
  content --- every one of which changed on adoption, and which the list above had
  omitted until v1.3.0: the abstract (``one open construction item''); the closing
  paragraph of \S\ref{sec:scope}'s scope statement (``the single open construction
  item''); rule~\ref{V5} (``\ref{open:ell2} is open''); the paragraph of
  \S\ref{sec:model} on what happens at $\ell=2$ on the boundary (``the policy there is
  a tagged working assumption''); the ``What this chapter fixes'' paragraph of
  \S\ref{sec:boundarydatum}; the opening of \S\ref{sec:ell2};
  Remark~\ref{rmk:iotaelltwo}'s closing sentence and the remark before
  Obligation~\ref{obl:coneexact} that names this chapter's two reading sites (both
  missed at v1.3.0 and repaired at v1.3.2); the ``What this appendix
  does not supply'' list of Appendix~\ref{app:formalization} (``A decision on
  \ref{open:ell2}''); Appendix~\ref{app:status}'s row for \ref{open:ell2} and its counts;
  and \S\ref{app:changelog-gaps} (``one open construction item is open''). The
  historical version lines of the Release record (the first subsection of Appendix~\ref{app:changelog}; the date line until v1.25.0-pre.7) and the v0.14.0 changelog entry are
  \emph{not} status sites and do not change.
\end{enumerate}

\paragraph{Code map.}
\textbf{Implemented since v1.3.3} (hecke-extensions slot 1, packages WP-4a and WP-4b,
2026-09-16; \file{docs/implementation/slot1/WP-4a\_report.md}, \file{WP-4b\_report.md}).
Until v1.3.2 this paragraph recorded a census of 2026-09-16 that found no $\taub$, no
$\Thb$, no $B_N$, no boundary coset table and no \code{HeckeBoundary} translation unit in
the tree, and only the boundary model without Hecke structure:
\code{TestSL2Z\_To\_S1Cusp\_RestrictionMorphism} (\file{LocalSystemSL2Z.cpp}), which
builds $\Kbdry$ as the four-vertex path, its \code{EquivariantCohomologyOnK}, the inclusion
$\Kbdry\hookrightarrow K$ and the induced $\rres^{q}$ with identity fibre map, pinned by
\code{HeckeMorphismRegressionValidation}. That demo is unchanged. The specification the
census closed with --- a unit beside \code{hecke::HeckeCosetTables}, a $\taub$-table in the
\code{ChainTerm} format of \code{hecke::HeckeTauTable} with every $\Ga''$-related slot
derived through the single shared (T1$\partial$) helper, certified in the constructor,
over the boundary subcomplex of \code{hecke::SL2ZTessellationTables} and degree-generic ---
is realised, unit by unit:
\begin{itemize}\setlength{\itemsep}{0pt}
\item \file{Hecke/SL2ZTessellationTables} (additive): the boundary subcomplex of
  \ref{B1}--\ref{B3} --- the horocycle path $v_0v_1v_2v_3$ at height $Y=2$ (vertices
  $\{0,1,2,3\}$ and edges $e_{01},e_{12},e_{23}$ of $K$), the $\infty$-line coordinates
  (the position of $T^{k}v_i$ is $3k+i$), the $\mathbb{P}^{1}(\mathbb{Q})$ cusp key with
  the reduction to $\infty$ by nearest-integer Euclid on the primitive column, all exact
  and overflow-checked, and a constructor-time \code{validate\_boundary()} asserting
  \ref{B1}--\ref{B3} on the shipped tables.
\item \file{Hecke/HeckeCosetTables} (additive): \code{twisted\_conjugate}, the one (T1)
  split-and-conjugate step, which \code{HeckeTauTable} now also calls --- the bulk and the
  boundary datum literally share the helper, as this paragraph asked.
\item \file{Hecke/HeckeBoundaryDatum} (\code{hecke::HeckeBoundaryTauTableT<Elt>}): the
  $\Ga$-level $\taub$ of Definition~\ref{def:taubdry}, one chain per $(j,\sigma)$ in
  \code{ChainTerm} format, degree-generic (no $q\in\{0,1,2\}$ literal; a boundary with
  cells above dimension~$1$ is refused, not truncated), every $\Ga''$-related slot derived
  through the shared (T1$\partial$) helper \code{tau\_on\_translate}; and the
  $B_N/\Gbdry$ table of Proposition~\ref{prop:bdrycosets} --- the canonical
  representatives $\begin{psmallmatrix}p^{u}&b\\0&p^{\nu-u}\end{psmallmatrix}$,
  $0\le b<p^{u}$, and the bijection of (2) --- certified in the constructor (invariant
  I8: codes 11--16 structural, 17 $=$ (iv), 20--22 $=$ (i), 30 $=$ (ii), 40 $=$ (iii),
  60--66 the coset table). The table contains no $\ell$: it is a datum over $\Z$, built
  for every $\ell$ the coset tables are built for, $\ell\in\{2,3\}$ and $p=\ell$
  included, and \code{solve\_regime() == "Z"} is a constant of the type, not the outcome
  of a fill (Remark~\ref{rmk:taubdryintegral}).
\item \file{Hecke/HeckeBoundaryOperator} (\code{hecke::HeckeBoundaryOperatorT<Elt>}):
  \eqref{eq:bdrychain} on the circle model $C^{\bullet}_{\mathrm{eq}}(\Kbdry,\Rel_{\partial};L)$
  --- now a library call, \code{LocalSystemSL2Z::CreateSL2ZBoundaryPathTriangulation} /
  \code{GenerateSL2ZBoundaryRelationsElt}, lifted out of the demo --- assembled from the
  bulk's own $s_j$, $\delta_j$ and $\taub$ in two bracketings asserted byte-identical;
  runtime assertions that $\Thb$ preserves the relations and the cocycles, the relation
  closure $\Thb C^{q}_{\mathrm{eq}}\subseteq C^{q}_{\mathrm{eq}}$ and $d\Thb=\Thb d$
  (Proposition~\ref{prop:bdryop}(1)); $\Abdry{q}=P_q\Thb z_i$ per
  Algorithm~\ref{alg:assembly}; a JSON record \code{hecke\_boundary\_operator/v1} carrying
  the certificate bundle (not a golden object).
\end{itemize}

\textbf{Obligation~\ref{obl:taubdry} is discharged in full, (i)--(vii)}, by
\code{HeckeBoundaryTauValidation} (\code{standard}; $138$ tables, $52\,686$ checks;
(i)--(iv) re-derived independently of the table's own helpers over
$(p,\nu)\in\{(2,1),(3,1),(5,1),(7,1),(11,1),(2,2),(3,2)\}$,
$\ell\in\{2,3,5,7,691,134217757\}$ at \code{uint32\_t} and $\ell\in\{2^{61}-1,2^{32}+15\}$
at \code{uint64\_t}, $n\in\{0,2,10\}$; the $\ell$-, $n$- and width-independence of the chains
asserted; seven tampered chains rejected) and \code{HeckeBoundaryOperatorValidation}
(\code{standard}; $546$ cells, $19\,669$ checks;
$(p,\nu)\in\{(2,1),(3,1),(5,1),(7,1),(2,2),(3,2)\}$,
$\ell\in\{2,3,5,7,11,13,101,691,134217757,2^{31}-1\}$, even $n\le16$, plus $2^{61}-1$ at
\code{uint64\_t}). Clause by clause: \COMP{}
\begin{itemize}\setlength{\itemsep}{0pt}
\item (v) --- two routes, both asserted equal to \eqref{eq:bdrychain} on cohomology,
  entrywise, inside \code{compute()} itself (a mismatch aborts): \emph{route A}, the
  Shapiro transcription of the decision box after Definition~\ref{def:bdryop} --- a
  $\Gbdry$-level twisted chain datum on the $\infty$-line built from the line coordinates
  alone, with the canonical left cosets $\delta_\infty,\delta_t$ of
  Proposition~\ref{prop:bdrycosets}(1) and the double-coset structure of (3), hence
  available at $\nu=1$; and \emph{route B}, the closed form of
  Proposition~\ref{prop:bdryop}(2) on $H^{0}=\VR^{T}$ (every $\nu$) and
  \eqref{eq:bdrygroupform} on $H^{1}=(\VR)_{T}$ ($\nu=1$), transported into the stored
  $H_{\mathrm{eq}}$ basis through the coordinates of Proposition~\ref{prop:circlemodel}(1)
  --- a $0$-cocycle is $(m,m,m,\rhon(T)m)$ with $m\in\VR^{T}$, and
  $H^{1}\to(\VR)_{T}$ is $\varphi\mapsto\varphi([v_0,Tv_0])$, the signed sum over the
  fundamental path, i.e.\ the group cocycle $f(T)$ of the proof of
  Proposition~\ref{prop:bdrychar}(3); with the opposite orientation the form would not
  agree. The validation re-derives route B on its own and re-asserts both on every cell.
  Since v1.4.0 the clause names the three comparisons the implementation runs: route B
  on $H^{0}$ is (v.a), asserted at every $\nu$ of the grid; route B on $H^{1}$ is (v.b)
  and route A is (v.c), both asserted at $\nu=1$; at $\nu=2$ the clause is (v.a) alone,
  and the $H^{1}$ operator there is covered by (vii) and by \code{compute()}'s own
  certificates, exactly as the clause now says.
\item (vi) --- on the $330$ cells with $n$ even and $p\neq\ell$: every
  $x^{\,n-b\ell}\mathrm{Nm}^{\,b}$ (own polynomial arithmetic; $575$ vectors) has eigenvalue
  $\chi_b(p)$ on $H^{0}$ and $\operatorname{charpoly}\Abdry{0}=\prod_b(X-\chi_b)$ ---
  asserted; for $n<\ell$ ($232$ cells) $[y^{n}]$ has eigenvalue $1+p^{\,n+1}$ on $H^{1}$ ---
  asserted; the general-$n$ $H^{1}$ multiset of Claim~\ref{cl:chi}(2) \textbf{recorded},
  $330/330$ cells agree, the claim still \CONJ{}. The $\beta_j$ right-coset substitution of
  Proposition~\ref{prop:bdrycosets}(4) is built as the \emph{negative control} (never used
  by the pipeline): it acts by $\chi_{b+1}(p)$ on every monomial, as
  Proposition~\ref{prop:bdrychar}(2) proves, differs from $\Abdry{0}$ on $212$ cells, and
  breaks clause (vii) on $H^{1}$ on the $170$ cells with $\chi_0(p)\neq\chi_1(p)$ and
  $\rank\rres^{1}>0$; on $H^{0}$ it cannot bite on this model, because the bulk
  $H^{0}(\Ga)\neq0$ only at $n=0$, where $\chi_0=\chi_1$, and on the Dickson cells
  $n=\ell+1$, where the invariant $x\,\mathrm{Nm}$ is the $b=1$ monomial with
  $\chi_1\equiv\chi_2$ (Claim~\ref{cl:chi}(3)).
\item (vii) --- $\rres^{q}\Amat{q}=\Abdry{q}\rres^{q}$ against \code{HeckeOperatorOnK}
  through the public restriction morphism, \textbf{asserted} on $864$ $(\text{cell},q)$
  pairs at $6\in\cring^{\times}$, $q\in\{0,1\}$ ($q=2$ vacuous, $H^{2}(\partial)=0$
  asserted), $\nu\in\{1,2\}$ and $p=\ell$ included; \textbf{recorded} at $\ell\in\{2,3\}$,
  where the chapter does not assert it: $216/216$ pairs hold on the model complexes ---
  offered to \S\ref{sec:cone} as a recorded fact, not as a claim, and recorded there as
  Remark~\ref{rmk:conerestrictsmallell} since v1.4.0. $\rres^{1}\neq0$ on $480$
  cells. The $\beta_j$ control's failure to bite on $H^{0}$ is what the qualified clause
  (vii) predicts for $n<\ell(\ell-1)$.
\end{itemize}
Family (A) of \S\ref{sec:oraclesnew} has been run with these numbers; its first instance
$(n,\ell,p)=(2,5,2)$ prints $\Abdry{0}=[4]$, $\Abdry{1}=[4]$, bulk $\Amat{1}(T_2)=[4]$,
$\rank\rres^{1}=1$, as that paragraph predicts.

\textbf{Choices the implementation fixed where the statements leave them open} ---
recorded as the implementation's choices, consistent with the statements, not as
changes to them:
\begin{itemize}\setlength{\itemsep}{0pt}
\item Lemma~\ref{lem:taubdryexists}(1) allows ``any vertex $w$ of the target component''.
  The canonical $\taub$ takes the line vertex nearest to $\mathrm{Re}(\gamma_{\mathrm{red}}z)$
  (the vertical projection onto the horocycle), ties to the larger real part --- a
  $T$-equivariant rule, hence independent of which $\gamma_{\mathrm{red}}$ represents
  the cusp reduction, and $-I$ acts trivially. \textbf{Since v1.4.0 this rule is
  Algorithm~\ref{alg:taubdry}}, the contract's definition of the canonical $\taub$
  (\code{HeckeBoundaryTauTableT::construct\_vertex}, \code{construct\_edge},
  \code{SL2ZTessellationTables::boundary\_line\_nearest}), with the well-definedness of
  the derivation order as its property~(1). Any other admissible choice is homotopic
  over $\Z$ for every $\ell$ (Lemma~\ref{lem:taubdryexists}(3)); the choice matters only
  for the size of \S\ref{sec:cone}'s $\etapair$. The rule is \emph{not} $\inv$-symmetric,
  and no vertex rule can be: Remark~\ref{rmk:taubdrynoiota}, which proves it at the
  level of rules (no line vertex is $\inv$-fixed) and exhibits the failure of the shipped
  table at $(p,\nu)=(2,1)$; Obligation~\ref{obl:iota} says nothing about $\taub$.
\item Proposition~\ref{prop:bdrycosets}(2) fixes $\gamma_j$ only up to $\gamma_j\Gbdry$;
  the implementation takes $\gamma_j:=\mathrm{reduce\_to\_infinity}(g^{-1}s_j\infty)^{-1}$,
  deterministic through the (P2) tie rule --- Algorithm~\ref{alg:taubdry} step~4 since
  v1.4.0.
\item Clause (iii)'s ``every stored transporter $h\in\Ga''$ occurring in the table'' is
  read as Obligation~\ref{obl:tau} reads it --- the transporter sets $S(\sigma,\sigma')$
  of the model between boundary simplices, through $s_jt$ --- and the validation covers
  arbitrary $\gamma\in\Ga$ by $40$ fixed-seed random words per table.
\item $\Rel_{\partial}$ of the circle model is the cusp relation
  $f(v_3)=\rhon(T)f(v_0)$ \emph{and} the $-I$ self-relations $f(\sigma)=\rhon(-I)f(\sigma)$
  on every simplex (the boundary half of Remark~\ref{rmk:negI}'s \code{FIX-negI}), so that
  $C^{\bullet}_{\mathrm{eq}}(\Kbdry,\Rel_{\partial};L)=\Hom_{\Z[\Gbdry]}(C_{\bullet}(\Xbdry),\VR)$
  literally for every $n$: vacuous for $n$ even or $\ell=2$, the zero complex for $n$ odd
  at $2\in\cring^{\times}$ (Remark~\ref{rmk:bdryoddn}, asserted by the odd-$n$ arm).
  \textbf{Since v1.4.0 this is Proposition~\ref{prop:circlemodel}(4)}: $\Rel_{\partial}$
  is defined in the proposition with the self-relations, and the identification is
  proved for every $n$; ``the circle model with its single relation'' survives only as
  the description of the case $\rhon(-I)=\id$.
\end{itemize}

\textbf{What it records} (owner confirmation 11 of the slot-1 plan: compute at every
$\ell$, never refuse): \code{boundary\_model = "circle"},
\code{boundary\_is\_group\_cohomology} $=(2\in\cring^{\times})=(\ell\neq2)$ --- the flag
Proposition~\ref{prop:circlemodel}(3) asks for, so the code now records which of the two
objects it computes --- \code{q\_max = 1} with \code{q\_max\_source = "derived"} (the top
degree of the finite circle model), \code{truncated = false}, reported degrees $\{0,1\}$;
\ref{open:ell2} has nothing to truncate on this complex (its rider (i) applies to the free
model). At $\ell=2$ the dimensions are Proposition~\ref{prop:circlemodel}(2)'s,
$\dim H^{0}=\dim H^{1}=\dim\VR^{T}=\lfloor n/2\rfloor+1$, and the validation asserts that
they are \emph{not} Lemma~\ref{lem:ell2unbounded}'s. The bulk object derives
$q_{\max}=2$; since v1.4.0, \S\ref{sec:recordschema} item~5 rules that a record carries
these fields \emph{per complex} and never as a maximum over a record (owner confirmation
1 of the slot-2 plan): bulk $2$ / \code{false} / \code{"derived"} / $[0,1,2]$, boundary
$1$ / \code{false} / \code{"derived"} / $[0,1]$.

The two shape requirements the specification stated hold as built, with one wording
corrected: the table is over $\Xbdry$ --- whose simplices \emph{are} simplices of $\Ku$
(the $\Ga$-translates of the boundary simplices of $K$), addressed by the same canonical
keys; what differs is the \emph{slot set}, supplied by
\code{SL2ZTessellationTables::boundary\_simplices} --- and it is degree-generic, not
$q\in\{0,1,2\}$, because \ref{open:ell2} indexes the free model by $q_{\max}$. What the
pair still lacks --- $\etapair$, the cone and $\Phig$ --- is \S\ref{sec:cone}'s code map.

\paragraph{Formalization note.}
Definition~\ref{def:bdrygroup} and \eqref{eq:bdrysplit} are a two-line computation with
$2\times2$ integer matrices. Proposition~\ref{prop:bdrycosets} is a finite
matrix-arithmetic statement with an explicit counting argument and is decidable for each
$N$; it is the natural first target here. Lemma~\ref{lem:bdrycomponent} and
Lemma~\ref{lem:taubdryexists} are the two that carry the content, and both are
\emph{constructive}: the second is a terminating procedure plus two uniqueness
statements, each of which reduces to $Z_1=0$ on a contractible $1$-complex, so a
formalisation needs simplicial chains and $H_1$ of a tree, not homological algebra ---
in particular it needs \textbf{neither} the fundamental lemma of relative homological
algebra \textbf{nor} any invertibility hypothesis, which is the whole point of stating it
separately from Lemma~\ref{lem:tauexists}(2). The hypothesis a formalisation must not lose
is \ref{B3}: it is what makes the (T1$\partial$)-extension well defined, and it is used
three times (existence on vertices, on edges, and in the uniqueness argument).
Proposition~\ref{prop:bdrycohom} needs the Lyndon--Hochschild--Serre spectral sequence of
a direct product with a cyclic factor, which \code{mathlib4} does not currently carry in
this form; the two-column degeneration can be replaced by the explicit
$\langle T\rangle$-equivariant periodic resolution of $\Z$ over $\Z[\Z/2]$ tensored with
the two-term resolution over $\Z[\Z]$, which is a finite construction and avoids the
spectral sequence entirely --- the recommended route. Proposition~\ref{prop:bdrychar} is
polynomial arithmetic in $\cring[x,y]$ modulo the Frobenius identities $u^{\ell}=u$ and
$p^{\ell-1}=1$, both of which are \code{ZMod} lemmas. Obligation~\ref{obl:taubdry} is a
decidable predicate on emitted finite data in every clause except (vii), which quantifies
over cohomology classes and is decidable only through the stored bases.
%----------------------------------------------------------------------
\section{The pair homotopy and the Hecke action on the mapping cone}\label{sec:cone}
%----------------------------------------------------------------------
% SOURCE:  RC+NEW.  OQ2 verdict [V] SS0, 2.1-2.10, 4, 6 (hecke_extensions_OQ2_verdict.md);
%          implementation plan SS0.4 C-A3(ii)-(iii) and S11, S11.1
%          (hecke_extensions_implementation_plan.md:412-566, 1994-2098); OQ2 revision
%          audit F5, B.1-B.4 (hecke_extensions_OQ2_revision_audit.md); specification (A2)
%          (hecke_extensions.tex:789-880); monodromy_v5.tex:1148-1265 (the cone
%          convention, the only RH input); owner decisions OP2, OP3, OP6 (plan S11).
%          No superseded document states the pair homotopy, the derived h, the cone
%          action, or any of the well-definedness statements.
% NOTATION: C2 (the cone degree is q), C3 (\coneinc/\coneproj, never \inv), C6 (\etapair is
%          the pair homotopy), C11 (\rres), C12 (\Lam, not M), C13 (\cring), the cone
%          element is (c, c_bdry) not (a, b); the coset transport of a chain-level map F is
%          written \widetilde F -- a tilde on a MAP, as in thm:welldef -- and is distinct
%          from the lift \lift{\sigma}, a tilde on a SIMPLEX (see Definition def:cosettransport).
% ERRATA:  E2, E3, E5, E13, E15 (App. A).  E14's cone half (the T_p^2 coefficient) is
%          applied by citing thm:relations and never restating the relation.

\paragraph{What this chapter fixes.}
The Hecke action on compactly supported cohomology, as a definition and not as a
slogan. The pair $(\Xbulk,\Xbdry)$, the restriction $\rres=j^{\ast}$ and the mapping cone
$\Cone{\rres}$ in the convention of \file{monodromy\_v5.tex} \S8; the \emph{pair homotopy}
$\etapair$, a certified chain-level datum one level below the cochains, where $\tau$ and
$\taub$ live, and the only new table this chapter asks for; the cochain homotopy $\hder$,
\textbf{derived} from $\etapair$ by coset transport and never chosen; the cone cochain
map $\Phig$ and the operator $T^{c}_{\Ga g\Ga}$ it induces on $\Hcpt{\bullet}$; and the
theorem that makes $\Hcpt{1}=X_3$ a Hecke module at all --- that the operator is
independent of $\etapair$ over $\Z$ for every $\ell$, and, for $6\in\cring^{\times}$, of
every other choice as well, and coincides with the correspondence on the relative
complex. It fixes, too, why the construction has this shape and not a simpler one: a
freely chosen cochain homotopy leaves the operator on $X_3$ undetermined by a
$\Hom(\ker\rres^{1},\coker\rres^{0})$ that is realised in full over a field and is
non-zero in a real case; and the entrywise identity $\rres\,\Th_g=\Thb_g\,\rres$, however
tempting, decides nothing (rule~\ref{A6}). Finally: integrality --- the three tables
are over $\Z$ once and every $\cring$-computation is their reduction, which is what
makes reductions and Bocksteins equivariant on the nose --- and the two obligations
that make all of this a checkable property of emitted data.

\textbf{Five errata land here} (Appendix~\ref{app:errata}): \ref{err:e2} and \ref{err:e3}
in Proposition~\ref{prop:coneambiguity}, which states the realisation argument with
$Z^{0}$ and never $B^{0}$, and \emph{over a field}, with $\Z/\ell^{e}$ explicitly not
claimed; \ref{err:e5} in Lemma~\ref{lem:etaexists}, which states the augmentation
argument and not the ``both are vertices'' recipe; \ref{err:e13} in
Lemma~\ref{lem:conesign} and Lemma~\ref{lem:derivedhomotopy}, where the cochain-level
sign is \emph{derived} from the chain-level convention and is therefore asserted in one
place only; and \ref{err:e15}, which this whole chapter answers --- ``cone by
functoriality'' is not a definition, and Definition~\ref{def:coneaction} together with
Proposition~\ref{prop:coneindep} is.

\textbf{What this chapter does not fix.} The boundary datum $\taub$ and the boundary
operator $\Thb_g$ are \S\ref{sec:boundarydatum} and are consumed here, not redefined.
The six groups $X_1,\dots,X_6$ and the six-term sequence are \S\ref{sec:six}; this
chapter proves the general long exact sequence and says under which hypotheses it
closes to that window. The free $\Ga$-model on which all of this is to be repeated at
$6\notin\cring^{\times}$ is \S\ref{sec:orbifoldmodel}; here $6\notin\cring^{\times}$ is
treated as \S\ref{sec:chainlevel} treats it (Remark~\ref{rmk:conetorsion}).

\paragraph{Standing hypotheses.}
\ref{H4} of \S\ref{sec:conventions} --- in particular \ref{H4}(ii), \emph{$\Xbulk$ is
contractible and two-dimensional}, and \ref{H4}(iv), \emph{$-I$ lies in every stabiliser
and acts trivially} --- and \ref{B1}--\ref{B3} of \S\ref{sec:boundarydatum}, each repeated
where it is used (rule~\ref{A1}). \LEAN{thm:model7-pairhyp} (on the seven-vertex model with its horocycle boundary: [GT]'s pair hypothesis
\texttt{PairHyp}, the form in which this chapter's lemmas use \ref{H4} and \ref{B1}--\ref{B3}, holds with no further hypothesis, for every $g\in\GL_2(\Q)$) Per the design direction of Remark~\ref{rmk:orbifold}
every statement is made for \emph{the tessellation and boundary triangulation in force}
and not for seven vertices; Remark~\ref{rmk:bdrymodel} is what makes the shipped
triangulation, and every equivariant refinement of it, admissible input. The
consequences of \ref{H4}(ii) that are used, and used by name, are
\begin{equation}\label{eq:contractibleuse}
  H_0(\Xbulk)\xrightarrow{\ \aug\ }\Z\ \text{is an isomorphism},\qquad
  Z_1(\Xbulk)=B_1(\Xbulk),\qquad
  Z_2(\Xbulk)=0 ,
\end{equation}
all over $\Z$ \LEAN{prop:model7h4-ii} (on the model) (and hence, $C_{\bullet}(\Xbulk)$ being free, after $\otimes\cring$ for every
$\cring$ \LEAN{thm:xbulk-repaug} \LEAN{lem:xbulk-contractz} (on the model, as the split exactness of the augmented $\cring$-complex)): the first two because $\Xbulk$ is connected with $H_1(\Xbulk)=0$, the third
because $H_2(\Xbulk)=0$ and $C_3(\Xbulk)=0$. The OQ2 revision audit (finding F5) found
these three facts load-bearing and \emph{unwritten} in every source of this chapter;
they are written here once and cited by \eqref{eq:contractibleuse} at each use. Throughout,
$g=\smat{1}{0}{0}{N}$, $\Ga''=\Ga_0(N)$, $s_1,\dots,s_m$ and $\delta_j=s_j^{-1}g$ are the
canonical data of Definition~\ref{def:cosetreps}, $\tau$ is a $g$-twisted chain datum
over $\Z$ on $\Xbulk$ (Definition~\ref{def:tau}) and $\taub$ a boundary $g$-twisted chain
datum over $\Z$ on $\Xbdry$ (Definition~\ref{def:taubdry}); ``(T1)'' means the twisted
equivariance $F(h\cdot x)=(g^{-1}hg)\cdot F(x)$ for $h\in\Ga''$, for whichever additive
map $F$ is under discussion.

\subsection{The pair, the restriction map and the cone}\label{sec:pair}

\begin{definition}[The pair $(\Xbulk,\Xbdry)$, and $\rres:=j^{\ast}$]\label{def:pair}
\LEAN{def:xbdry} \LEAN{thm:xbdry-pair} ($\Xbdry$ as a $\Ga$-saturated, face-closed sub-descriptor; $j$ equivariant, a chain map over the augmentations)
Assume \ref{H4} and \ref{B1}--\ref{B3}, and let $\cring$ be a commutative ring. The
\emph{pair} is $(\Xbulk,\Xbdry)$ with $j\colon\Xbdry\hookrightarrow\Xbulk$ the inclusion,
a $\Ga$-equivariant simplicial map, written also for the induced chain map
$j\colon C_{\bullet}(\Xbdry)\to C_{\bullet}(\Xbulk)$. As the inclusion of a subcomplex, $j$ commutes with the augmentations:
$\aug\circ j=\aug_\partial$, where $\aug_\partial\colon C_0(\Xbdry)\to\Z$ sends every vertex to~$1$ (stated explicitly at v1.18.0-pre.2, slot-18 L18-4: the Lean
statement \texttt{GL2.Cone.eta\_exists} needs it as the hypothesis \texttt{hjaug}, an abstract $j$ not being an inclusion; Lemma~\ref{lem:etaexists}(V) uses it; no
statement changes). Its two equivariant cochain complexes
are
\begin{align*}
  \Ceq(K,\Rel;L)&\;\cong\;\Hom_{\Z[\Ga]}\bigl(C_{\bullet}(\Xbulk),\VR\bigr)
  &&(\text{Theorem~\ref{thm:explicit_complex_equiv}, }\Ku\cong\Xbulk),\\
  C^{\bullet}_{\partial}&\;=\;\Hom_{\Z[\Ga]}\bigl(C_{\bullet}(\Xbdry),\VR\bigr)
  &&(\text{Proposition~\ref{prop:circlemodel}}),
\end{align*}
and the \emph{restriction} is the cochain map
$\rres:=j^{\ast}\colon\Ceq(K,\Rel;L)\to C^{\bullet}_{\partial}$, $\phi\mapsto\phi\circ j$.
The \emph{relative complex} is $\ker\rres$, with the differential restricted from
$\Ceq$. The \emph{mapping cone} of $\rres$ is the cochain complex
\begin{equation}\label{eq:conediff}
  \Cone{\rres}^{q}:=\Cq{q+1}(K,\Rel;L)\oplus C^{q}_{\partial},
  \qquad
  \Dcone(c,c_\partial):=\bigl(-dc,\ dc_\partial+\rres c\bigr),
\end{equation}
with $\coneinc\colon C^{\bullet}_{\partial}\to\Cone{\rres}$, $c_\partial\mapsto(0,c_\partial)$,
and $\coneproj\colon\Cone{\rres}\to\Ceq(K,\Rel;L)[1]$, $(c,c_\partial)\mapsto c$, where
$\Ceq[1]^{q}:=\Cq{q+1}$ carries the differential $-d$. \textbf{Compactly supported
cohomology} is
\[
  \Hcpt{q}\;:=\;H^{q-1}\bigl(\Cone{\rres}\bigr),\qquad q\in\Z .
\]
This is the convention of \file{monodromy\_v5.tex} \S8 (\code{:1152--1163}) with its
$(a,b')$ written $(c,c_\partial)$, its cone degree $n$ written $q$, and its
$\iota,\pi$ written $\coneinc,\coneproj$ (\S\ref{sec:notation-cone}); for the model of
Definition~\ref{def:model7} it is exactly the table entry
$\Cone{\rres}^{q}=\Cq{q+1}(K)\oplus\Cq{q}(\Kbdry)$ of \S\ref{sec:notation-cone}, since
there $C^{\bullet}_{\partial}=\Ceq(\Kbdry,\Rel_{\partial};L)$
(Proposition~\ref{prop:circlemodel}).
\end{definition}

\begin{lemma}[$\rres$ is onto; the cone is a complex; $\Hcpt{\bullet}$ is relative cohomology]\label{lem:conerelative}
\PROVED{} Assume \ref{H4} and \ref{B1}--\ref{B3} and let $\cring$ be a commutative ring;
no hypothesis on $\ell$ or $n$.
\begin{enumerate}
\item \textbf{(Surjectivity.)} $\rres$ is surjective in every degree: the
  \emph{extension by zero} $\mathrm{ext}_0\colon C^{q}_{\partial}\to\Cq{q}(K,\Rel;L)$,
  which gives a cochain the value $\varphi(\sigma)$ on a simplex $\sigma$ of $\Xbdry$ and
  $0$ on every other simplex of $\Xbulk$, is $\cring$-linear, $\Ga$-equivariant, and a
  section of $\rres$. It is \emph{not} a cochain map, and is never used as one.
\item \textbf{(The cone.)} \LEAN{lem:conerelative} $\Dcone^{2}=0$; $\coneinc$ and $\coneproj$ are cochain maps,
  and $0\to C^{\bullet}_{\partial}\xrightarrow{\coneinc}\Cone{\rres}\xrightarrow{\coneproj}
  \Ceq(K,\Rel;L)[1]\to0$ is a short exact sequence of complexes of $\cring$-modules
  whose connecting homomorphism $H^{q}(\Ceq[1])=\Heq{q+1}\to H^{q+1}(C^{\bullet}_{\partial})$
  is $+\rres^{q+1}$. Hence the \textbf{long exact sequence}
  \begin{equation}\label{eq:coneles}
    \cdots\longrightarrow\Hcpt{q}\xrightarrow{\ \coneproj{}_{\ast}\ }\Heq{q}(K,\Rel;L)
    \xrightarrow{\ \rres^{q}\ }H^{q}(C^{\bullet}_{\partial})
    \xrightarrow{\ \coneinc{}_{\ast}\ }\Hcpt{q+1}\longrightarrow\cdots ,
  \end{equation}
  in every degree $q\in\Z$.
\item \textbf{(Relative cohomology.)} \LEAN{lem:conerelative} $c\mapsto(c,0)$ is a quasi-isomorphism
  $(\ker\rres)[1]\to\Cone{\rres}$; consequently
  \[
    \Hcpt{q}\;\cong\;H^{q}(\ker\rres)\;=\;H^{q}\Bigl(\Hom_{\Z[\Ga]}\bigl(C_{\bullet}(\Xbulk,\Xbdry),\VR\bigr)\Bigr)
    \qquad\text{for every }q,
  \]
  where $C_{\bullet}(\Xbulk,\Xbdry):=C_{\bullet}(\Xbulk)/j\,C_{\bullet}(\Xbdry)$ is the
  relative chain complex, a complex of $\Z[\Ga]$-permutation modules on the simplices of
  $\Xbulk$ not in $\Xbdry$.
\item \textbf{(Vanishing.)} $\Hcpt{q}=0$ for $q\le0$ and for $q\ge3$. In particular
  $\Hcpt{0}=0$ \emph{with no hypothesis on $\ell$ or $n$}: it is the kernel of
  $\rres^{0}\colon\Heq{0}\to H^{0}(C^{\bullet}_{\partial})$, and that map is injective.
\item \textbf{(The window, and the two dimension counts.)} \LEAN{lem:conerelative} (given (4), which is not formalised: a statement about the model) Let $\cring=\F$ be a field.
  \begin{equation}\label{eq:conedims}
    \dim\Hcpt{1}=\dim H^{0}(C^{\bullet}_{\partial})-\dim\Heq{0}+\dim\ker\rres^{1},\qquad
    \dim\Hcpt{2}=\dim H^{1}(C^{\bullet}_{\partial})-\rank\rres^{1} .
  \end{equation}
  \textbf{The first identity holds with no further hypothesis} (v1.6.0): it uses only the
  exactness of \eqref{eq:coneles} and $\Hcpt{0}=0$ from (4), both of which are
  unconditional. \textbf{The second identity is asserted under $\Heq{2}(K,\Rel;L)=0$},
  which is also exactly the hypothesis under which \eqref{eq:coneles} collapses to the
  six-term exact sequence
  $0\to\Heq{0}\to H^{0}(\partial)\to\Hcpt{1}\to\Heq{1}\to H^{1}(\partial)\to\Hcpt{2}\to0$
  of Claim~\ref{cl:les}. That hypothesis is Claim~\ref{cl:amalgam}'s business (it holds
  for $6\in\cring^{\times}$) and is \emph{not} proved here. (Through v1.5.2 both
  identities were stated under it; \file{docs/implementation/slot3/WP-5\_report.md}
  finding~3.)
\end{enumerate}
\end{lemma}

\begin{proof}
(1) $\Ga$ preserves $\Xbdry$ (\ref{B1}), so $\gamma\sigma\in\Xbdry$ iff $\sigma\in\Xbdry$,
and $(\mathrm{ext}_0\varphi)(\gamma\sigma)=\rhon(\gamma)(\mathrm{ext}_0\varphi)(\sigma)$
holds in both cases; $\rres\circ\mathrm{ext}_0=\id$ by construction. That
$\mathrm{ext}_0$ is not a cochain map: a boundary edge is a face of an interior triangle,
so $d(\mathrm{ext}_0\varphi)$ can be non-zero on a simplex where
$\mathrm{ext}_0(d\varphi)$ vanishes.

(2) $\Dcone^{2}(c,c_\partial)=\Dcone(-dc,dc_\partial+\rres c)
=\bigl(d^{2}c,\ d(dc_\partial+\rres c)-\rres\,dc\bigr)=(0,\ d\rres c-\rres dc)=0$, using
$\rres d=d\rres$. $\Dcone\coneinc(c_\partial)=(0,dc_\partial)=\coneinc(dc_\partial)$ and
$\coneproj\Dcone(c,c_\partial)=-dc$, the differential of the shift. Levelwise exactness
is the definition of a direct sum. The long exact sequence is
\STD{\cite[Thm.~1.3.1]{weibel}: the long exact sequence of a short exact sequence of
complexes, with the connecting map of its proof}. For the connecting map: a class in
$H^{q}(\Ceq[1])$ is $[c]$ with $c\in\Zq{q+1}$; lift to $(c,0)\in\Cone{\rres}^{q}$, apply
$\Dcone$ to get $(0,\rres c)=\coneinc(\rres c)$; so the connecting map is
$[c]\mapsto[\rres c]=\rres^{q+1}[c]$, with the sign $+$ forced by \eqref{eq:conediff}.
Re-indexing $H^{q-1}(\Cone{\rres})=\Hcpt{q}$ gives \eqref{eq:coneles}.

(3) The map is a cochain map because $\Dcone(c,0)=(-dc,\rres c)=(-dc,0)$ for
$c\in\ker\rres$, and $-d$ is the differential of the shift. Its cokernel is identified
by $(c,c_\partial)\mapsto(\rres c,c_\partial)$ with the complex
$\Cone{\id}$ of $\id\colon C^{\bullet}_{\partial}\to C^{\bullet}_{\partial}$,
$\Cone{\id}^{q}=C^{q+1}_{\partial}\oplus C^{q}_{\partial}$,
$\Dcone_{\id}(b,b')=(-db,db'+b)$: the map is a cochain map
($\Dcone_{\id}(\rres c,c_\partial)=(-d\rres c,dc_\partial+\rres c)$ is the image of
$\Dcone(c,c_\partial)$), it is surjective by (1), and its kernel is $\{(c,0):\rres c=0\}
=(\ker\rres)[1]$. $\Cone{\id}$ is contractible: with $s(b,b'):=(b',0)$ one computes
$(\Dcone_{\id}s+s\Dcone_{\id})(b,b')=(-db',b')+(db'+b,0)=(b,b')$. So in the long exact
sequence of $0\to(\ker\rres)[1]\to\Cone{\rres}\to\Cone{\id}\to0$ every third term
vanishes and the first map is an isomorphism on cohomology
\STD{\cite[\S1.5]{weibel} for the mapping cone; \cite[Thm.~1.3.1]{weibel} for the
sequence}. Finally $\Hom_{\Z[\Ga]}(-,\VR)$ is left exact, so applying it to
$0\to C_{\bullet}(\Xbdry)\xrightarrow{j}C_{\bullet}(\Xbulk)\to C_{\bullet}(\Xbulk,\Xbdry)\to0$
(exact because $j$ is the inclusion of a subcomplex, \ref{B1}) identifies
$\Hom_{\Z[\Ga]}(C_{\bullet}(\Xbulk,\Xbdry),\VR)$ with $\ker\rres$.

(4) $\ker\rres$ vanishes outside degrees $0,1,2$ because $\Xbulk$ is two-dimensional
(\ref{H4}(ii)), so $H^{q}(\ker\rres)=0$ for $q<0$ and $q>2$; with (3) this is
$\Hcpt{q}=0$ for $q\le-1$ and $q\ge3$. For $q=0$: $\Hcpt{0}=H^{0}(\ker\rres)$ consists
of the equivariant $0$-cocycles $c$ with $\rres c=0$. A $0$-cocycle satisfies
$c(\partial e)=0$ for every edge, hence is constant on the vertices of the connected
complex $\Xbulk$ \eqref{eq:contractibleuse}; a constant that vanishes on the non-empty
set of boundary vertices (\ref{B1}) is $0$. Equivalently: $\rres^{0}$ is the inclusion
$\VR^{\Ga}\hookrightarrow\VR^{\Gbdry}$ (Remark~\ref{rmk:T3}, Proposition~\ref{prop:bdryop}(2)),
which is injective.

(5) \emph{The first identity, with no hypothesis.} $\Hcpt{0}=0$ by (4), so exactness of
\eqref{eq:coneles} at $\Heq{0}$ makes $\rres^{0}$ injective, and exactness at
$H^{0}(\partial)$, at $\Hcpt{1}$ and at $\Heq{1}$ gives the exact sequence
\[
  0\longrightarrow\Heq{0}\xrightarrow{\ \rres^{0}\ }H^{0}(C^{\bullet}_{\partial})
  \xrightarrow{\ \coneinc{}_{\ast}\ }\Hcpt{1}
  \xrightarrow{\ \coneproj{}_{\ast}\ }\ker\rres^{1}\longrightarrow0 ,
\]
whose last term is $\im(\coneproj{}_{\ast}\colon\Hcpt{1}\to\Heq{1})=\ker\rres^{1}$, again
by exactness at $\Heq{1}$. Over a field the alternating sum of dimensions of an exact
sequence vanishes, which is the first identity. \textbf{Nothing here uses $\Heq{2}$},
$H^{q}(C^{\bullet}_{\partial})$ for $q\ge1$, or any hypothesis on $\ell$ or $n$ beyond
\ref{H4} and \ref{B1}--\ref{B3}.

\emph{The six-term window and the second identity, under $\Heq{2}=0$.} With that
hypothesis the terms $\Hcpt{0}$ (by (4)), $\Heq{2}$ (by hypothesis),
$H^{2}(C^{\bullet}_{\partial})$ ($\Xbdry$ is one-dimensional, \ref{B2}) and $\Hcpt{3}$
(by (4)) of \eqref{eq:coneles} all vanish, leaving the displayed six-term sequence; its
second half is $0\to\coker\rres^{1}\to\Hcpt{2}\to0$, i.e.\ the second identity. That half
is where the hypothesis is used: without $\Heq{2}=0$ the map $\Hcpt{2}\to\Heq{2}$ need
not be zero and $\coker\rres^{1}\to\Hcpt{2}$ need not be onto.
\end{proof}

\noindent
Every map built in this chapter is a \emph{coset transport}: the device of
Theorem~\ref{thm:welldef}(b), where it produced $\widetilde H$ from a twisted homotopy
$H$, applied to maps between the chain complexes of the pair. It is given a name and
four identities once, so that every later proof is a two-line computation.

\begin{definition}[Coset transport]\label{def:cosettransport}
Let $Y_{\bullet},Y'_{\bullet}$ each be one of $C_{\bullet}(\Xbulk)$, $C_{\bullet}(\Xbdry)$,
and let $F\colon Y_{\bullet}\to Y'_{\bullet+k}$ be an additive map of degree $k\in\Z$
satisfying (T1). Its \emph{coset transport} is the $\cring$-linear map
\begin{equation}\label{eq:cosettransport}
  \widetilde F\colon\Hom_{\Z[\Ga]}\bigl(Y'_{q+k},\VR\bigr)\longrightarrow
  \Hom_{\Z[\Ga]}\bigl(Y_{q},\VR\bigr),\qquad
  \bigl(\widetilde F\phi\bigr)(y)\;:=\;\sum_{j=1}^{m}\rhon(\delta_j)\,\phi\bigl(F(s_j\cdot y)\bigr),
\end{equation}
of \emph{cochain} degree $-k$. Thus $\widetilde\tau$ is the $\widetilde T$ of
\eqref{eq:heckechain}, i.e.\ $\Th_g$ under $\Psi$ (Theorem~\ref{thm:xi}(3));
$\widetilde{\taub}$ is $\Thb_g$ of \eqref{eq:bdrychain}; and for the twisted homotopy
$H$ of Lemma~\ref{lem:tauexists}(2), $\widetilde H$ is the map of that name in the proof
of Theorem~\ref{thm:welldef}(b). \emph{Notation.} The tilde on a \emph{map} denotes coset
transport; the tilde on a \emph{simplex}, $\lift{\sigma}$, denotes the lift of
\S\ref{sec:notation-framework} (C16). The two never meet in one formula, but a reader
should know they are unrelated. The decoration has two further readings --- the
cochain-level operator $\widetilde T$ of Definition~\ref{def:hecke}, which is
$\widetilde\tau$ and so agrees with the first, and Part~I's saturation $\widetilde K$ and
universal cover $\widetilde X$, which occur nowhere near either. All four are
\textbf{C19} of \S\ref{sec:notation-collisions} since v0.16.0: the decoration is read
from what it decorates, nothing is renamed, and no new tilde may be introduced without a
row there.
\end{definition}

\begin{lemma}[Coset transport: equivariance and four identities]\label{lem:cosettransport}
\PROVED{} Assume \ref{H4} and \ref{B1}, let $\cring$ be a commutative ring, and let
$F,G$ be additive maps with (T1) as in Definition~\ref{def:cosettransport}. Then, with no
hypothesis on $\ell$:
\begin{enumerate}
\item $\widetilde F\phi$ is $\Ga$-equivariant for every $\Ga$-equivariant $\phi$, so
  \eqref{eq:cosettransport} is well typed; and $F\mapsto\widetilde F$ is additive.
\item $\widetilde{\partial\circ F}\,\phi=\widetilde F(d\phi)$ and
  $\widetilde{F\circ\partial}\,\phi=d\bigl(\widetilde F\phi\bigr)$.
\item If $G\colon C_{\bullet}(\Xbdry)\to C_{\bullet+k}(\Xbdry)$, then
  $\widetilde{j\circ G}\,\phi=\widetilde G(\rres\phi)$ for
  $\phi\in\Hom_{\Z[\Ga]}(C_{\bullet}(\Xbulk),\VR)$.
\item If $G\colon C_{\bullet}(\Xbulk)\to C_{\bullet+k}(\Xbulk)$, then
  $\widetilde{G\circ j}\,\phi=\rres\bigl(\widetilde G\phi\bigr)$.
\end{enumerate}
\end{lemma}

\begin{proof}
(1) Let $\gamma\in\Ga$ and write $s_j\gamma=h_j(\gamma)s_{j\cdot\gamma}$ with
$h_j(\gamma)\in\Ga''$ (Lemma~\ref{lem:cocycle}). By (T1) of $F$, $\Ga$-equivariance of
$\phi$, and $g^{-1}h_j(\gamma)g\in\Ga$,
\[
  \rhon(\delta_j)\,\phi\bigl(F(s_j\gamma y)\bigr)
  =\rhon(\delta_j)\,\rhon\bigl(g^{-1}h_j(\gamma)g\bigr)\,\phi\bigl(F(s_{j\cdot\gamma}y)\bigr)
  =\rhon\bigl(\gamma\,s_{j\cdot\gamma}^{-1}g\bigr)\,\phi\bigl(F(s_{j\cdot\gamma}y)\bigr),
\]
by Proposition~\ref{prop:symconv}(2) on the integral identity
$\delta_j\,g^{-1}h_j(\gamma)g=s_j^{-1}h_j(\gamma)g=\gamma s_{j\cdot\gamma}^{-1}g$ in
$\Mat_2(\Z)$ \eqref{eq:cocycle}. Summing over $j$, and using that $j\mapsto j\cdot\gamma$
is a bijection, gives $(\widetilde F\phi)(\gamma y)=\rhon(\gamma)(\widetilde F\phi)(y)$.
This is the equivariance computation of Theorem~\ref{thm:welldef} with $F$ in place of
$\tau$; it used (T1) of $F$ and nothing else about it. Additivity in $F$ is clear.

(2) $(\widetilde{\partial F}\phi)(y)=\sum_j\rhon(\delta_j)\phi(\partial F(s_jy))
=\sum_j\rhon(\delta_j)(d\phi)(F(s_jy))=(\widetilde F d\phi)(y)$, since
$(d\phi)(x)=\phi(\partial x)$. And, $s_j$ acting simplicially so that
$\partial(s_jy)=s_j\partial y$,
$(\widetilde{F\partial}\phi)(y)=\sum_j\rhon(\delta_j)\phi(F(\partial(s_jy)))
=\sum_j\rhon(\delta_j)\phi(F(s_j\,\partial y))=(\widetilde F\phi)(\partial y)
=(d\widetilde F\phi)(y)$.

(3) $(\widetilde{jG}\phi)(y)=\sum_j\rhon(\delta_j)\phi(jG(s_jy))
=\sum_j\rhon(\delta_j)(\phi\circ j)(G(s_jy))=(\widetilde G(\rres\phi))(y)$.

(4) $(\widetilde{Gj}\phi)(y)=\sum_j\rhon(\delta_j)\phi(G(js_jy))
=\sum_j\rhon(\delta_j)\phi(G(s_j\,jy))=(\widetilde G\phi)(jy)=(\rres\widetilde G\phi)(y)$,
using that $j$ is $\Ga$-equivariant.
\end{proof}

\subsection{The pair homotopy}\label{sec:pairhomotopy}

\begin{definition}[Pair homotopy $\etapair$: $\partial\etapair+\etapair\partial=j\taub-\tau j$]\label{def:pairhomotopy}
\LEAN{def:pairhomotopy} (the augmentation compatibility is part of Definition~\ref{def:pair})
Assume \ref{H4} and \ref{B1}--\ref{B3}, let $\cring$ be a commutative ring, and let
$\tau$, $\taub$ be $g$-twisted chain data over $\cring$ on $\Xbulk$ and $\Xbdry$
(Definitions~\ref{def:tau}, \ref{def:taubdry}), for the same $g$, the same $\Ga''$ and
the same twist $h\mapsto g^{-1}hg$. A \emph{pair homotopy} for $(\tau,\taub)$ is an
additive map of degree $+1$
\[
  \etapair\colon C_{\bullet}(\Xbdry)\otimes\cring\longrightarrow C_{\bullet+1}(\Xbulk)\otimes\cring
\]
satisfying (T1) and
\begin{equation}\label{eq:pairhomotopy}
  \partial\circ\etapair+\etapair\circ\partial\;=\;j\circ\taub-\tau\circ j
  \qquad\text{as maps }C_{\bullet}(\Xbdry)\otimes\cring\to C_{\bullet}(\Xbulk)\otimes\cring .
\end{equation}
Both sides of \eqref{eq:pairhomotopy} are twisted chain maps $C_{\bullet}(\Xbdry)\to
C_{\bullet}(\Xbulk)$ lifting the same map of components (Lemma~\ref{lem:bdrycomponent}),
and $\etapair$ is a twisted chain homotopy between them. An \emph{$\etapair$-table} is
such a map together with the finite table of its values on a set of $\Ga''$-orbit
representatives of the simplices of $\Xbdry$ --- a $1$-chain of $\Xbulk$ for each
representative vertex, a $2$-chain for each representative edge --- from which (T1)
determines it everywhere. The \emph{Hecke datum of the pair} is the triple
$(\tau,\taub,\etapair)$.
\end{definition}

\begin{decisionbox}
\textbf{The sign convention is \eqref{eq:pairhomotopy}, $\partial\etapair+\etapair\partial
=j\taub-\tau j$, and the derived homotopy of Definition~\ref{def:coneaction} is
$\hder:=\widetilde{\etapair}$ with no sign} (owner decision OP6, 2026-09-15). The
opposite convention would give $\hder=-\widetilde{\etapair}$; it is used nowhere in
this document or in the implementation plan. What this decision does \emph{not} move is
the cochain identity $d\hder+\hder d=\Thb_g\rres-\rres\Th_g$ and the cone map $\Phig$,
which are forced by $\Dcone$ (Lemma~\ref{lem:conesign}) and are the same under either
convention.
\end{decisionbox}

\begin{lemma}[$\etapair$ exists over $\Z$ for every $\ell$]\label{lem:etaexists}
\PROVED{} \LEAN{lem:etaexists} \LEAN{cor:model7-etaexists} (on the seven-vertex model with its horocycle boundary, every $g\in\GL_2(\Q)$: \ref{H4} and \ref{B1}--\ref{B3} discharged in [GT]) (Appendix~\ref{app:errata}, entry~\ref{err:e5}: the vertex step is an
augmentation argument, not ``a path between two vertices''.) Assume \ref{H4} and
\ref{B1}--\ref{B3}, and let $\tau$, $\taub$ be $g$-twisted chain data \emph{over $\Z$}.
Then a pair homotopy for $(\tau,\taub)$ over $\Z$ exists, and one is built by the
following terminating procedure (rule~\ref{A3}). No hypothesis on $\ell$ is used, no
stabiliser order is inverted, and no linear system is solved over anything but $\Z$.
\begin{enumerate}
\item[(V)] \textbf{Vertices.} For each $\Ga''$-orbit representative vertex $x$ of
  $\Xbdry$: the $0$-chain $(j\taub-\tau j)(x)$ of $\Xbulk$ has augmentation
  $\aug(j\taub x)-\aug(\tau jx)=1-1=0$ --- by $\aug\circ j=\aug_\partial$ (Definition~\ref{def:pair}) with (T3$\partial$) for $\taub$, and (T3) for
  $\tau$ --- hence is a boundary \eqref{eq:contractibleuse}. Choose \emph{any}
  $1$-chain $\etapair(x)\in C_1(\Xbulk)$ with $\partial\etapair(x)=(j\taub-\tau j)(x)$,
  and extend to all boundary vertices by (T1).
\item[(E)] \textbf{Edges.} For each $\Ga''$-orbit representative edge $e$ of $\Xbdry$:
  the $1$-chain $z_e:=(j\taub-\tau j)(e)-\etapair(\partial e)$ of $\Xbulk$ is a cycle,
  hence a boundary, and there is a \emph{unique} $2$-chain $\etapair(e)$ with
  $\partial\etapair(e)=z_e$ \eqref{eq:contractibleuse}. Extend by (T1).
\item[($\ge2$)] There is nothing to do: $C_{q}(\Xbdry)=0$ for $q\ge2$ (\ref{B2}).
\end{enumerate}
The map so defined satisfies (T1) and \eqref{eq:pairhomotopy}; the only choice made
is the vertex chain in (V), and Lemma~\ref{lem:etaunique} says exactly what that choice
costs. \emph{The choice made} (v1.5.0, naming what the implementation does; the
statement is unchanged, (V) still admitting \emph{any} $1$-chain with the stated
boundary): the augmentation-$0$ $0$-chain $(j\taub-\tau j)(x)$ is written as a sum of
pairs (one vertex with coefficient $+1$, one with $-1$) and each pair is joined by the
\textbf{pairwise path routing} of Algorithm~\ref{alg:tau} step~2 --- the stored path
template between the two vertices --- which is the augmentation argument made
constructive, needs no linear solve, and reduces to \emph{the one path joining the two
vertices} when $j\taub x$ and $\tau jx$ are single vertices, as they are on the shipped
tables (v1.6.0, corrected from ``the single edge path'', which could be read as ``one
edge''; the stored path template between two vertices may be several edges long).
\S\ref{sec:cone}'s code map records it as the implemented choice, and
Lemma~\ref{lem:etaunique} is what makes the naming harmless.
The procedure terminates because $\Ga''$ has finitely many orbits of simplices
on $\Xbdry$ (\ref{B1} and $[\Ga:\Ga'']=m<\infty$, Proposition~\ref{prop:cosets}(3)).
\end{lemma}

\begin{proof}
\emph{(T1)-extension is well defined.} Both in (V) and in (E) the value on an orbit
$\Ga''\cdot\sigma$ is set by $\etapair(h\sigma):=(g^{-1}hg)\cdot\etapair(\sigma)$,
which is well defined iff $(g^{-1}hg)\cdot\etapair(\sigma)=\etapair(\sigma)$ for every
$h$ in the $\Ga''$-stabiliser of $\sigma$. By \ref{B3} that stabiliser is
$\Ga''\cap\{\pm I\}$, $g^{-1}(-I)g=-I$, and $-I$ acts trivially on $\Xbulk$
(\ref{H4}(iv)); so the condition is vacuous. \emph{This is where the difference from
the abstract route of Lemma~\ref{lem:tauexists}(2) lies}: no elliptic stabiliser ever
occurs on the source, so no averaging idempotent and no invertibility of $4$ or $6$ is
needed --- the same remark as in Lemma~\ref{lem:taubdryexists}(1).

\emph{(V) and the identity on vertices.} $(j\taub-\tau j)(x)$ is an integer $0$-chain of
augmentation $0$, and $\aug\colon H_0(\Xbulk)\to\Z$ is an isomorphism
\eqref{eq:contractibleuse}, so it is $\partial$ of an integer $1$-chain; this is the
statement entry~\ref{err:e5} corrects --- $\tau(jx)$ is a $0$-chain of augmentation $1$
and not, in general, a single vertex, and the argument never needs it to be one. Once
$\partial\etapair(x)=(j\taub-\tau j)(x)$ holds on representatives it holds on every
boundary vertex, because both sides satisfy (T1) (the right-hand side by (T1) of
$\tau$, of $\taub$, and equivariance of $j$) and agree on representatives. This is
\eqref{eq:pairhomotopy} in degree $0$, where $\etapair\partial=0$.

\emph{(E).} Using (T2) for $\tau$ and $\taub$, that $j$ is a chain map, and the degree-$0$
identity just proved,
\[
  \partial z_e=(j\taub-\tau j)(\partial e)-\partial\etapair(\partial e)
  =(j\taub-\tau j)(\partial e)-(j\taub-\tau j)(\partial e)=0 ,
\]
so $z_e\in Z_1(\Xbulk)=B_1(\Xbulk)$ and a $2$-chain with boundary $z_e$ exists; it is
unique because two such differ by an element of $Z_2(\Xbulk)=0$
\eqref{eq:contractibleuse}. Both facts hold over $\Z$, so the fill is an integer
$2$-chain. Then $\partial\etapair(e)+\etapair(\partial e)=z_e+\etapair(\partial e)
=(j\taub-\tau j)(e)$ on representatives, hence on all boundary edges by (T1) of both
sides: this is \eqref{eq:pairhomotopy} in degree $1$. Degrees $\ge2$ are vacuous by
\ref{B2}, and additivity is by construction.
\end{proof}

\begin{lemma}[$\etapair$ is unique up to $\partial\Lam-\Lam\partial$, $\Lam$ unique]\label{lem:etaunique}
\PROVED{} \LEAN{lem:etaunique} \LEAN{cor:model7-etaexists} (on the seven-vertex model with its horocycle boundary, every $g\in\GL_2(\Q)$: \ref{H4} and \ref{B1}--\ref{B3} discharged in [GT]) Assume \ref{H4} and \ref{B1}--\ref{B3}, let $\cring$ be a commutative ring, and
let $\etapair,\etapair'$ be two pair homotopies for the same $(\tau,\taub)$ over
$\cring$. Then there is a \textbf{unique} additive map of degree $+2$
\[
  \Lam\colon C_{\bullet}(\Xbdry)\otimes\cring\longrightarrow C_{\bullet+2}(\Xbulk)\otimes\cring
  \qquad\text{with}\qquad
  \etapair'-\etapair\;=\;\partial\Lam-\Lam\partial ,
\]
and it satisfies (T1) automatically. Explicitly $\Lam$ is non-zero only on $C_0(\Xbdry)$,
where $\Lam(x)$ is the unique $2$-chain with $\partial\Lam(x)=(\etapair'-\etapair)(x)$;
in degree $1$ the identity reads $(\etapair'-\etapair)(e)=-\Lam(\partial e)$ and is
\emph{forced}. Over $\cring=\Z$ this holds with integer chains. The hypotheses used are
\ref{H4}(ii) through \eqref{eq:contractibleuse} --- for the existence and the uniqueness of
the fills --- and \ref{B2} for the vanishing of $C_2(\Xbdry)$; \ref{B3} is \emph{not}
needed here, because uniqueness makes (T1) automatic.
\end{lemma}

\begin{proof}
Put $\xi:=\etapair'-\etapair$. Subtracting the two instances of
\eqref{eq:pairhomotopy} gives $\partial\xi+\xi\partial=0$, and $\xi$ has (T1) as a
difference of two maps with (T1).

\emph{Degree $0$.} For a boundary vertex $x$, $\partial\xi(x)=-\xi(\partial x)=0$, so
$\xi(x)\in Z_1(\Xbulk)=B_1(\Xbulk)$ \eqref{eq:contractibleuse}: there is a $2$-chain
$\Lam(x)$ with $\partial\Lam(x)=\xi(x)$, and it is unique because $Z_2(\Xbulk)=0$.
Define $\Lam$ on \emph{every} boundary vertex this way (not only on representatives)
and extend additively to $C_0(\Xbdry)$. (T1): for $h\in\Ga''$,
$\partial\bigl((g^{-1}hg)\Lam(x)\bigr)=(g^{-1}hg)\xi(x)=\xi(hx)=\partial\Lam(hx)$, so
$(g^{-1}hg)\Lam(x)=\Lam(hx)$ by uniqueness. No stabiliser argument is needed.

\emph{Degree $1$.} For a boundary edge $e$,
$\partial\bigl(\xi(e)+\Lam(\partial e)\bigr)=\partial\xi(e)+\xi(\partial e)=0$, so
$\xi(e)+\Lam(\partial e)\in Z_2(\Xbulk)=0$, i.e.\ $\xi(e)=-\Lam(\partial e)$. Since
$C_3(\Xbulk)=0$ the only possible value of $\Lam$ on $C_1(\Xbdry)$ is $0$, and with it
$(\partial\Lam-\Lam\partial)(e)=-\Lam(\partial e)=\xi(e)$.

\emph{Degrees $\ge2$.} $C_q(\Xbdry)=0$ (\ref{B2}). So $\xi=\partial\Lam-\Lam\partial$.
\emph{Uniqueness of $\Lam$.} Any $\Lam'$ with $\xi=\partial\Lam'-\Lam'\partial$ has
$\partial\Lam'(x)=\xi(x)$ on vertices, hence $\Lam'(x)=\Lam(x)$ by $Z_2(\Xbulk)=0$, and
$\Lam'=0=\Lam$ in degree $1$.
\end{proof}

\subsection{The derived homotopy and the cone action}\label{sec:coneaction}

The next lemma is pure algebra about the cone and is stated for an \emph{arbitrary}
degree-$(-1)$ map $\hder$; it is where the sign of the cochain identity comes from, and
it is the only place it is asserted.

\begin{lemma}[The sign of $\Dcone\Phig-\Phig\Dcone$]\label{lem:conesign}
\PROVED{} (Appendix~\ref{app:errata}, entry~\ref{err:e13}.) Assume \ref{H4},
\ref{B1}--\ref{B3}, let $\cring$ be a commutative ring, let $\Th\colon\Ceq(K,\Rel;L)\to
\Ceq(K,\Rel;L)$ and $\Thb\colon C^{\bullet}_{\partial}\to C^{\bullet}_{\partial}$ be
cochain maps, and let $\hder\colon\Cq{q}(K,\Rel;L)\to C^{q-1}_{\partial}$ ($q\in\Z$) be
any $\cring$-linear maps. Define
$\Phi\colon\Cone{\rres}^{q}\to\Cone{\rres}^{q}$ by
$\Phi(c,c_\partial):=(\Th c,\ \Thb c_\partial+\hder c)$. Then
\begin{equation}\label{eq:conesign}
  (\Dcone\Phi-\Phi\Dcone)(c,c_\partial)\;=\;\bigl(0,\ (d\hder+\hder d+\rres\Th-\Thb\rres)\,c\bigr),
\end{equation}
so $\Phi$ is a cochain map \textbf{iff} $d\hder+\hder d=\Thb\rres-\rres\Th$. Moreover
$\Phi\circ\coneinc=\coneinc\circ\Thb$ and $\coneproj\circ\Phi=\Th\circ\coneproj$ hold
for every $\hder$.
\end{lemma}

\begin{proof}
$\Dcone\Phi(c,c_\partial)=\Dcone(\Th c,\Thb c_\partial+\hder c)
=\bigl(-d\Th c,\ d\Thb c_\partial+d\hder c+\rres\Th c\bigr)$ and
$\Phi\Dcone(c,c_\partial)=\Phi(-dc,dc_\partial+\rres c)
=\bigl(-\Th dc,\ \Thb dc_\partial+\Thb\rres c-\hder dc\bigr)$. Subtracting, and using
$d\Th=\Th d$, $d\Thb=\Thb d$, gives \eqref{eq:conesign}. The two intertwinings are read
off the definitions: $\Phi(0,c_\partial)=(0,\Thb c_\partial)$ and
$\coneproj\Phi(c,c_\partial)=\Th c$.
\end{proof}

\noindent
This is the sign three independent derivations agree on --- the implementation plan
v1.2, the OQ2 verdict \S2.1 and the OQ2 revision audit B.1 --- and it was re-derived
for this chapter from \eqref{eq:conediff} alone; the plan audit's
$d\hder+\hder d=\rres\Th-\Thb\rres$ is entry~\ref{err:e13}, and nothing downstream of it
carried the wrong orientation.

\begin{definition}[The derived cochain homotopy $\hder:=\widetilde{\etapair}$, and $\Phig$]\label{def:coneaction}
Assume \ref{H4} and \ref{B1}--\ref{B3}, let $\cring$ be a commutative ring, and let
$(\tau,\taub,\etapair)$ be a Hecke datum of the pair (Definition~\ref{def:pairhomotopy}),
with $\Th_g=\widetilde\tau$ and $\Thb_g=\widetilde{\taub}$ the cochain maps of
Definitions~\ref{def:hecke}/\ref{def:xi} and \ref{def:bdryop} built from the same $g$,
$s_j$, $\delta_j$. The \textbf{derived cochain homotopy} is the coset transport of
$\etapair$,
\begin{equation}\label{eq:derivedh}
  \hder:=\widetilde{\etapair}\colon\Cq{q}(K,\Rel;L)\longrightarrow C^{q-1}_{\partial},
  \qquad
  (\hder\phi)(y)=\sum_{j=1}^{m}\rhon(\delta_j)\,\phi\bigl(\etapair(s_j\cdot y)\bigr),
  \quad y\in C_{q-1}(\Xbdry),
\end{equation}
and the \textbf{cone cochain map} is
\begin{equation}\label{eq:conemap}
  \Phig\colon\Cone{\rres}^{q}\to\Cone{\rres}^{q},\qquad
  \Phig(c,c_\partial):=\bigl(\Th_g\,c,\ \Thb_g\,c_\partial+\hder\,c\bigr).
\end{equation}
The operator $T^{c}_{\Ga g\Ga}$ on $\Hcpt{q}=H^{q-1}(\Cone{\rres})$ is the map induced by
$\Phig$ (Theorem~\ref{thm:conewelldef}(a) says it is one), and its matrix $\Acone{q}$ in
the stored basis of $\Hcpt{q}$ is assembled by Algorithm~\ref{alg:assembly} run on the
complex $\Cone{\rres}$ in place of $\Ceq$ --- the stored cocycle basis of
$Z^{q-1}(\Cone{\rres})$, the projector modulo $B^{q-1}(\Cone{\rres})$, and \eqref{eq:conemap}
applied columnwise --- with the same column convention. For the shipped model
$\hder$ is non-zero in cochain degrees $1$ and $2$ only: $\hder^{1}\colon\Cq{1}\to
C^{0}_{\partial}$ is the transport of $\etapair$ on vertices and $\hder^{2}\colon\Cq{2}\to
C^{1}_{\partial}$ of $\etapair$ on edges, and $\hder^{0}=\hder^{3}=0$ because $\Xbdry$ is
one-dimensional (\ref{B2}). \textbf{$\hder$ is never chosen at cochain level}; the only
choices are the chain-level tables, and Proposition~\ref{prop:coneindep} says what they
cost.
\end{definition}

\begin{lemma}[The derived homotopy satisfies the cone identity]\label{lem:derivedhomotopy}
\PROVED{} Under the hypotheses of Definition~\ref{def:coneaction}, and with no hypothesis
on $\ell$:
\[
  d\hder+\hder d\;=\;\Thb_g\,\rres-\rres\,\Th_g
  \qquad\text{on }\Cq{q}(K,\Rel;L),\ \text{for every }q .
\]
\end{lemma}

\begin{proof}
Apply Lemma~\ref{lem:cosettransport} to \eqref{eq:pairhomotopy}. By (2),
$\widetilde{\partial\etapair}\phi=\widetilde{\etapair}(d\phi)=\hder d\phi$ and
$\widetilde{\etapair\partial}\phi=d(\widetilde{\etapair}\phi)=d\hder\phi$; by (3),
$\widetilde{j\taub}\phi=\widetilde{\taub}(\rres\phi)=\Thb_g\rres\phi$; by (4),
$\widetilde{\tau j}\phi=\rres(\widetilde\tau\phi)=\rres\Th_g\phi$. Coset transport is
additive in the map (Lemma~\ref{lem:cosettransport}(1)), so transporting both sides of
$\partial\etapair+\etapair\partial=j\taub-\tau j$ gives
$\hder d+d\hder=\Thb_g\rres-\rres\Th_g$.
\end{proof}

\begin{theorem}[Well-definedness of the cone action]\label{thm:conewelldef}
Assume \ref{H4} and \ref{B1}--\ref{B3}, let $\cring$ be a commutative ring, and let
$(\tau,\taub,\etapair)$ be a Hecke datum of the pair over $\Z$, with $\Phig$ as in
\eqref{eq:conemap} over $\cring$.
\begin{enumerate}
\item[(a)] \PROVED{} \textbf{(No hypothesis on $\ell$ or $n$.)} \LEAN{thm:conewelldef} $\Phig$ is a cochain map
  of complexes of $\cring$-modules, $\Phig\coneinc=\coneinc\Thb_g$ and
  $\coneproj\Phig=\Th_g\coneproj$. Hence $\Phig$ induces $T^{c}_{\Ga g\Ga}$ on
  $\Hcpt{q}$ for every $q$, and \textbf{every map of \eqref{eq:coneles} intertwines}
  $(T^{\partial}_{\Ga g\Ga},T^{c}_{\Ga g\Ga},\Tdc{g})$: $\coneproj{}_{\ast}$,
  $\coneinc{}_{\ast}$ and $\rres^{q}$ commute with the three operators. This is not a
  separate obligation; it is a consequence of (a).
\item[(b)] \PROVED{} \LEAN{thm:conewelldef} \LEAN{cor:model7-etaexists} (on the seven-vertex model with its horocycle boundary, every $g\in\GL_2(\Q)$: \ref{H4} and \ref{B1}--\ref{B3} discharged in [GT]) \textbf{(Independence of $\etapair$ --- over $\Z$, for every
  $\ell$.)} Two pair homotopies $\etapair,\etapair'$ for the same $(\tau,\taub)$ give
  cone-homotopic $\Phig$: $\Phig^{\etapair'}-\Phig^{\etapair}=\Dcone\Upsilon+\Upsilon\Dcone$
  for the integral, filtration-preserving homotopy
  $\Upsilon(c,c_\partial)=(0,-\widetilde{\Lam}\,c)$ of
  Proposition~\ref{prop:coneindep}(1). Hence $T^{c}_{\Ga g\Ga}$ on $\Hcpt{\bullet}$, and
  every map it induces on \eqref{eq:coneles}, is independent of $\etapair$.
\item[(c)] \textbf{(For $6\in\cring^{\times}$ --- hypothesis \ref{H3}.)} \PROVED{}
  $T^{c}_{\Ga g\Ga}$ is independent of $\tau$ and $\taub$, of the transporters $s_j$
  (in the sense of Theorem~\ref{thm:welldef}(a): at cochain level with the tables held
  fixed, on cohomology in general) and of $g$ within $\Ga g\Ga$
  (Proposition~\ref{prop:coneindep}(2)--(4)); and under the identification of
  Lemma~\ref{lem:conerelative}(3) it is the operator $\cor\circ(\conj_g)_{\ast}\circ\res$
  on $H^{\bullet}\bigl(\Hom_{\Z[\Ga]}(C_{\bullet}(\Xbulk,\Xbdry),\VR)\bigr)$ induced by any
  \emph{strict} pair (Proposition~\ref{prop:strictpair}), which exists; the
  identification of $\res$, $\cor$, $\conj$ on the relative complex is that of
  Theorem~\ref{thm:classical}, on the same external inputs (\cite[\S1]{ash_stevens},
  \cite[Ch.~III.9--10]{brown}).
\item[(c$'$)] \textbf{(For $6\in\cring^{\times}$; the two consequences that reach outside
  this document, with their status fixed in Remark~\ref{rmk:conesteinberg}.)} The
  \textbf{Hecke relations of Theorem~\ref{thm:relations} hold on $\Hcpt{\bullet}$ on the
  nose}, in the same normalisation and with the same full-versus-primitive caveat ---
  \PROVED{} through the Steinberg-module form of Remark~\ref{rmk:conesteinberg}(ii), on
  the same external inputs as Theorem~\ref{thm:relations}. The operator of (c) is the
  classical Hecke correspondence on the relative cohomology of the Borel--Serre pair,
  i.e.\ on the compactly supported cohomology of the open orbifold ---
  \STD{\cite[\S\S1--2]{harder_eisenstein} and \cite{borel_serre}, neither of whose
  section numbers is verified here (Remark~\ref{rmk:conesteinberg}(iii))}. $T_p$ on $\Hcpt{1}$ is
  adjoint to $T_p$ on $\Heq{1}(K,\Rel;\Dn)$ under the pairing of
  Proposition~\ref{prop:pairing} --- \OPENSTAT{} here, being Claim~\ref{cl:pd}'s
  (Remark~\ref{rmk:conesteinberg}(iv), Proposition~\ref{prop:coneroutes}(3)).
\item[(d)] \PROVED{} \textbf{(Strictness.)} \LEAN{thm:conewelldef} $\etapair\equiv0$ is admissible iff
  $\tau j=j\taub$ (``$\tau$ is a map of pairs''); then $\hder=0$ and
  $\Phig=(\Th_g,\Thb_g)$. \textbf{The cochain identity $\rres\Th_g=\Thb_g\rres$ alone
  does not license $\hder=0$} (Remark~\ref{rmk:conestrict}).
\item[(e)] \textbf{(For $6\notin\cring^{\times}$, on the complexes of this model.)} (a),
  (b) and (d) hold as stated; $T^{c}_{\Ga g\Ga}$ is \emph{defined} as the output of the
  pipeline run with the canonical $(\tau,\taub,\etapair)$, and (c) is not asserted ---
  exactly as for $\Tdc{g}$ itself in the contract box after
  Lemma~\ref{lem:tauexists} (Remark~\ref{rmk:conetorsion}).
\end{enumerate}
\end{theorem}

\begin{proof}
(a) Lemma~\ref{lem:conesign} with $\Th=\Th_g$, $\Thb=\Thb_g$ (cochain maps by
Theorem~\ref{thm:xi}(2) and Proposition~\ref{prop:bdryop}(1)) and $\hder$ of
\eqref{eq:derivedh}: the bracket in \eqref{eq:conesign} vanishes by
Lemma~\ref{lem:derivedhomotopy}, so $\Phig$ is a cochain map; it is $\cring$-linear
because $\Th_g$, $\Thb_g$, $\hder$ are. The two intertwinings are the last sentence of
Lemma~\ref{lem:conesign}, so $(\Thb_g,\Phig,\Th_g)$ is a map of the short exact sequence
of Lemma~\ref{lem:conerelative}(2) to itself, and a map of short exact sequences of
complexes induces a map of long exact sequences
\STD{\cite[Thm.~1.3.1]{weibel}, naturality of the connecting homomorphism}.

(b), (c), (c$'$) are Proposition~\ref{prop:coneindep}, Proposition~\ref{prop:strictpair}
and Remark~\ref{rmk:conesteinberg}. (d): if $\tau j=j\taub$ then $\etapair=0$ satisfies
\eqref{eq:pairhomotopy} and (T1); conversely if $\etapair=0$ is a pair homotopy then
\eqref{eq:pairhomotopy} reads $0=j\taub-\tau j$. With $\etapair=0$,
\eqref{eq:derivedh} gives $\hder=0$. (e): (a), (b), (d) used no hypothesis on $\ell$;
(c) used Lemma~\ref{lem:tauexists}(2), which is unavailable there.
\end{proof}

\begin{proposition}[Independence of $\etapair$ over $\Z$; of $\tau,\taub,s_j,g$ for $6\in\cring^{\times}$]\label{prop:coneindep}
Assume \ref{H4} and \ref{B1}--\ref{B3}, and let $\cring$ be a commutative ring.
\begin{enumerate}
\item \PROVED{} \textbf{($\etapair$; over $\Z$, every $\ell$.)} \LEAN{prop:coneindep} Let $\etapair,\etapair'$
  be pair homotopies for $(\tau,\taub)$ and $\Lam$ the unique map of
  Lemma~\ref{lem:etaunique}, so $\etapair'-\etapair=\partial\Lam-\Lam\partial$. Then
  $\hder'-\hder=\widetilde{\Lam}\,d-d\,\widetilde{\Lam}$, where
  $\widetilde\Lam\colon\Cq{q}\to C^{q-2}_{\partial}$ is the coset transport of $\Lam$, and
  with $\Upsilon\colon\Cone{\rres}^{q}\to\Cone{\rres}^{q-1}$,
  $\Upsilon(c,c_\partial):=(0,-\widetilde\Lam c)$,
  \begin{equation}\label{eq:conehomotopyeta}
    \Phig^{\etapair'}-\Phig^{\etapair}\;=\;\Dcone\Upsilon+\Upsilon\Dcone .
  \end{equation}
  $\Upsilon$ is $\cring$-linear, is defined over $\Z$ when the tables are, and is
  non-zero only on $\Cone{\rres}^{1}$.
\item \PROVED{} \textbf{($\tau$ and $\taub$; needs $6\in\cring^{\times}$ for the bulk
  half only.)} Let $\tau'=\tau+\partial H+H\partial$ with $H$ a degree-$(+1)$ map with
  (T1) (Lemma~\ref{lem:tauexists}(2): any two $\tau$ are so related when
  $6\in\cring^{\times}$), and $\taub'=\taub+\partial N_{\partial}+N_{\partial}\partial$
  with $N_{\partial}$ as in Lemma~\ref{lem:taubdryexists}(3) (any two $\taub$ are so
  related over $\Z$, for every $\ell$). Then $\etapair':=\etapair+jN_{\partial}-Hj$ is a
  pair homotopy for $(\tau',\taub')$, and with
  $\Upsilon(c,c_\partial):=(-\widetilde Hc,\ \widetilde{N_{\partial}}c_\partial)$,
  \begin{equation}\label{eq:conehomotopytau}
    \Phig^{(\tau',\taub',\etapair')}-\Phig^{(\tau,\taub,\etapair)}\;=\;\Dcone\Upsilon+\Upsilon\Dcone .
  \end{equation}
  Any other pair homotopy for $(\tau',\taub')$ is then covered by (1). In particular
  the induced operator is independent of $\taub$ over $\Z$ for every $\ell$, and of
  $\tau$ for $6\in\cring^{\times}$.
\item \PROVED{} \textbf{($s_j$; no hypothesis on $\ell$ when the tables are held
  fixed.)} Replacing $s_j$ by $h_js_j$ with $h_j\in\Ga''$ and keeping
  $(\tau,\taub,\etapair)$ fixed leaves $\Th_g$, $\Thb_g$, $\hder$ and hence $\Phig$
  unchanged, summand by summand. In an admissible construction the tables are built
  relative to the $s_j$ and are rebuilt with them (Theorem~\ref{thm:welldef}(a), the
  proviso); the induced operator is then invariant by (2), hence for
  $6\in\cring^{\times}$.
\item \PROVED{} \textbf{($g$ within $\Ga g\Ga$.)} For $g_1=\gamma_1g\gamma_2$,
  $\gamma_1,\gamma_2\in\Ga$, the transported triple
  $\bigl(\gamma_2^{-1}\tau(\gamma_1^{-1}\,\cdot\,),\ \gamma_2^{-1}\taub(\gamma_1^{-1}\,\cdot\,),\
  \gamma_2^{-1}\etapair(\gamma_1^{-1}\,\cdot\,)\bigr)$ with transporters
  $s_j\gamma_1^{-1}$ is a Hecke datum of the pair for $g_1$ and reproduces $\Phig$
  verbatim, with no hypothesis on $\ell$; an arbitrary admissible datum at $g_1$ is
  compared with the transported one by (2) and (3), hence under $6\in\cring^{\times}$.
\end{enumerate}
\end{proposition}

\begin{proof}
(1) By Lemma~\ref{lem:cosettransport}(1),(2) applied to $\partial\Lam-\Lam\partial$:
$\hder'-\hder=\widetilde{\etapair'-\etapair}=\widetilde{\partial\Lam}-\widetilde{\Lam\partial}
=\widetilde\Lam d-d\widetilde\Lam$. Now, for $(c,c_\partial)\in\Cone{\rres}^{q}$, using
$\Th_g$, $\Thb_g$ unchanged,
\[
  \Dcone\Upsilon(c,c_\partial)=\Dcone(0,-\widetilde\Lam c)=(0,-d\widetilde\Lam c),\qquad
  \Upsilon\Dcone(c,c_\partial)=\Upsilon(-dc,dc_\partial+\rres c)=(0,\widetilde\Lam dc),
\]
whose sum is $(0,(\widetilde\Lam d-d\widetilde\Lam)c)=(0,(\hder'-\hder)c)
=(\Phig^{\etapair'}-\Phig^{\etapair})(c,c_\partial)$. Degrees: $c\in\Cq{q+1}$,
$\widetilde\Lam c\in C^{q-1}_{\partial}$, so $\Upsilon$ lands in $\Cone{\rres}^{q-1}$;
$\Lam$ is non-zero only on $C_0(\Xbdry)$, so $\widetilde\Lam$ is non-zero only on
$\Cq{2}$ and $\Upsilon$ only on $\Cone{\rres}^{1}$. Integrality: $\Lam$ is an integer
table when $\etapair,\etapair'$ are (Lemma~\ref{lem:etaunique}), and
\eqref{eq:cosettransport} has integral coefficients.

(2) $\etapair'$ has (T1) as a sum of maps with (T1) ($j$ is equivariant). And
$\partial\etapair'+\etapair'\partial
=(j\taub-\tau j)+j(\partial N_{\partial}+N_{\partial}\partial)-(\partial H+H\partial)j
=j\taub'-\tau'j$, using that $j$ is a chain map. By Lemma~\ref{lem:cosettransport},
$\Th_g'=\Th_g+d\widetilde H+\widetilde Hd$ (this is Theorem~\ref{thm:welldef}(b)),
$\Thb_g{}'=\Thb_g+d\widetilde{N_{\partial}}+\widetilde{N_{\partial}}d$
(Proposition~\ref{prop:bdryop}(1)), and by (3),(4) of that lemma
$\hder'=\hder+\widetilde{N_{\partial}}\rres-\rres\widetilde H$. Then
\[
  \Dcone\Upsilon(c,c_\partial)=\bigl(d\widetilde Hc,\ d\widetilde{N_{\partial}}c_\partial-\rres\widetilde Hc\bigr),
  \qquad
  \Upsilon\Dcone(c,c_\partial)=\bigl(\widetilde Hdc,\ \widetilde{N_{\partial}}dc_\partial+\widetilde{N_{\partial}}\rres c\bigr),
\]
and the sum is
$\bigl((d\widetilde H+\widetilde Hd)c,\ (d\widetilde{N_{\partial}}+\widetilde{N_{\partial}}d)c_\partial
+(\widetilde{N_{\partial}}\rres-\rres\widetilde H)c\bigr)$, which is the difference of the
two cone maps. Two pair homotopies for $(\tau',\taub')$ are covered by (1), and cone
homotopies compose. The hypothesis $6\in\cring^{\times}$ enters only through the
existence of $H$ (Lemma~\ref{lem:tauexists}(2)); $N_{\partial}$ exists over $\Z$ for every
$\ell$ (Lemma~\ref{lem:taubdryexists}(3)).

(3) The computation in the proof of Theorem~\ref{thm:welldef}(a) shows that the $j$-th
summand of \eqref{eq:cosettransport} is unchanged under $s_j\mapsto h_js_j$ for any $F$
with (T1); apply it to $F=\tau,\taub,\etapair$. The proviso is that of
Theorem~\ref{thm:welldef}(a) and Remark~\ref{rmk:indeplevel}, unchanged.

(4) As in the proof of Theorem~\ref{thm:welldef}(c): the transported $\tau_1,\taub{}_1$
are $g_1$-twisted (that computation, verbatim, on each complex), and for the transported
$\etapair_1$ one has $\partial\etapair_1+\etapair_1\partial
=\gamma_2^{-1}(j\taub-\tau j)(\gamma_1^{-1}\,\cdot\,)=j\taub{}_1-\tau_1j$ because $j$ is
equivariant. Substituting into \eqref{eq:cosettransport} with transporters
$s_j\gamma_1^{-1}$ reproduces each of $\Th_g$, $\Thb_g$, $\hder$, the factors
$\rhon(\gamma_1)$, $\rhon(\gamma_2)$ cancelling against equivariance of $\phi$.
\end{proof}

\begin{proposition}[Strict pairs exist over $\Z$, for every $\ell$]\label{prop:strictpair}
\PROVED{} Assume \ref{H4} and \ref{B1}--\ref{B3}, and let $\taub$ be a boundary
$g$-twisted chain datum over $\Z$ (Lemma~\ref{lem:taubdryexists}(1)). Then there is a
$g$-twisted chain datum $\tau^{s}$ over $\Z$ on $\Xbulk$ with
\[
  \tau^{s}\circ j\;=\;j\circ\taub ,
\]
so that $\etapair\equiv0$ is a pair homotopy for $(\tau^{s},\taub)$ and
$\Phig=(\Th^{s}_g,\Thb_g)$ preserves $(\ker\rres)[1]\subseteq\Cone{\rres}$. It is built by
the following terminating procedure, which never inverts a stabiliser order and never
solves a system over anything but $\Z$; the only external fact it uses is that every
finite subgroup of $\Ga$ of order $4$ or $6$ fixes a vertex of $\Xbulk$
\STD{\cite[Ch.~VII, \S1.2, Thm.~1(1),(3)]{serre_cours}: a point of $\Hyp$ with
non-trivial stabiliser modulo $\pm I$ is $\Ga$-equivalent to $i$, to $\zeta^{2}$ or to
$\zeta$; these are the vertices $v_6$, $v_4$, $v_5$ of Definition~\ref{def:model7}, whose
$\Ga$-translates all lie in $\Xbulk$ since $\mathrm{Im}\le1<Y$. The pointer read
``Thm.~2'' until v0.16.0; Serre's Theorem~2 is the generation by $S$ and $T$ and is not
what is used here. That such a subgroup has a fixed point at all is the step spelled out
in the proof of Lemma~\ref{lem:tauexists}(1): it is cyclic, and a generator is a torsion
element of $\Ga$ other than $\pm I$, hence elliptic}.
\begin{enumerate}
\item[(B)] \textbf{Boundary simplices.} $\tau^{s}(j\sigma):=j\taub(\sigma)$ for every
  simplex $\sigma$ of $\Xbdry$. This has (T1) and (T2) because $\taub$ does and $j$ is an
  equivariant chain map.
\item[(V)] \textbf{Interior vertices.} For each $\Ga''$-orbit representative vertex $v$
  of $\Xbulk$ not in $\Xbdry$: the twisted stabiliser $g^{-1}\Ga''_{v}g$ is a finite
  subgroup of $\Ga$ (Proposition~\ref{prop:cosets}(3)); choose $\tau^{s}(v):=w$, a vertex
  of $\Xbulk$ fixed by it --- any vertex if the group is $\{\pm I\}$, the fixed vertex
  supplied by the cited fact otherwise. Extend by (T1).
\item[(E)] \textbf{Interior edges.} For each $\Ga''$-orbit representative edge $e$ of
  $\Xbulk$ not in $\Xbdry$: $\tau^{s}(\partial e)$ is a $0$-chain of augmentation $0$,
  hence a boundary \eqref{eq:contractibleuse}; choose any $1$-chain with that boundary.
  Extend by (T1); this is well defined because the stabiliser of an edge is $\{\pm I\}$
  (\ref{H4}(iv)).
\item[(F)] \textbf{Triangles.} For each $\Ga''$-orbit representative triangle $\sigma$:
  $\tau^{s}(\partial\sigma)$ is a $1$-cycle, hence the boundary of a \emph{unique}
  $2$-chain \eqref{eq:contractibleuse}; take it. Extend by (T1), well defined since the
  stabiliser is $\{\pm I\}$.
\end{enumerate}
\end{proposition}

\begin{proof}
(B) is well defined on the $\Ga''$-stable subcomplex $\Xbdry$ and gives (T1), (T2), (T3)
there. (V): the (T1)-extension $\tau^{s}(hv):=(g^{-1}hg)w$ is well defined iff $w$ is
fixed by $g^{-1}\Ga''_{v}g$, which is how $w$ was chosen; if $\Ga''_{v}$ has order $4$
or $6$ then so does its conjugate $g^{-1}\Ga''_{v}g\subseteq\Ga$, which by the cited
fact fixes a point of $\Ga\cdot i\cup\Ga\cdot\zeta$, a vertex of $\Xbulk$ by
\ref{H4}(i),(iv). (T3) holds since $\tau^{s}(v)$ is one vertex. (E), (F): (T2) is
imposed on representatives and propagates to orbits because both sides have (T1);
existence of the fills is \eqref{eq:contractibleuse} over $\Z$, exactly as in
Lemma~\ref{lem:etaexists}, and uniqueness in (F) is $Z_2(\Xbulk)=0$. (T1) for the
extension needs the stabiliser condition, vacuous for $\{\pm I\}$ by \ref{H4}(iv).
Termination: finitely many $\Ga''$-orbits (\ref{H4}(iii), $[\Ga:\Ga'']<\infty$). The
consequences: $\etapair=0$ satisfies \eqref{eq:pairhomotopy} for $(\tau^{s},\taub)$;
$\hder=0$; and for $c\in\ker\rres$,
$\rres\Th^{s}_gc=\widetilde{\tau^{s}j}\,c=\widetilde{j\taub}\,c=\Thb_g\rres c=0$ by
Lemma~\ref{lem:cosettransport}(3),(4), so $\Th^{s}_g$ preserves $\ker\rres$ and
$\Phig(c,0)=(\Th^{s}_gc,0)$.
\end{proof}

\noindent
Proposition~\ref{prop:strictpair} contains, as its bulk half, exactly the statement
Lemma~\ref{lem:tauexists}(1) makes: existence of a $g$-twisted chain datum over $\Z$ with
no hypothesis on $\ell$, by a construction. (Historically it came first; that steps (V),
(E), (F) prove existence on their own, with step~(B) dropped and hence without
\ref{B1}--\ref{B3}, is what Lemma~\ref{lem:tauexists}(1) records. Neither proof uses the
other, so there is no circularity.) What this proposition adds is \emph{strictness} ---
the boundary slots pre-assigned from $\taub$, so that $\etapair\equiv0$ is admissible.
What neither does is repair the \emph{uniqueness} half Lemma~\ref{lem:tauexists}(2),
which needs $6\in\cring^{\times}$ and is used in Proposition~\ref{prop:coneindep}(2); see
Remark~\ref{rmk:tauexistsscope}. And the shipped $\tau$ of Algorithm~\ref{alg:tau} is
\emph{not} built this way (Remark~\ref{rmk:conestrict}).

\begin{remark}[The relative complex, the Steinberg module, and where the Hecke relations come from]\label{rmk:conesteinberg}
This remark fixes the status of the assertions of Theorem~\ref{thm:conewelldef}(c$'$)
that reach outside this document, and gives the algebraic form under which one of them
becomes provable on inputs this document already uses. Assume \ref{H4},
\ref{B1}--\ref{B3} and $6\in\cring^{\times}$ throughout.

\emph{(i) The operator on the relative complex is canonical.} \PROVED{} For a strict pair
$(\tau^{s},\taub)$ the cone operator is, by Lemma~\ref{lem:conerelative}(3) and
Proposition~\ref{prop:strictpair}, the map induced by $\Th^{s}_g|_{\ker\rres}$ on
$H^{\bullet}(\ker\rres)=H^{\bullet}(\Hom_{\Z[\Ga]}(C_{\bullet}(\Xbulk,\Xbdry),\VR))$, and
$\Th^{s}_g|_{\ker\rres}$ is $\cor\circ(\conj_g)_{\ast}\circ\res$ on that complex with
$(\conj_g)_{\ast}$ realised by $\psi\mapsto\rhon(g)\circ\psi\circ\tau^{s}_{\mathrm{rel}}$,
$\tau^{s}_{\mathrm{rel}}$ the map induced on $C_{\bullet}(\Xbulk,\Xbdry)$ --- the proof
sketch of Theorem~\ref{thm:classical} verbatim, on the relative complex, which for
$6\in\cring^{\times}$ is again a complex of relatively projective $\cring[\Ga]$-modules
(its terms are the permutation modules on the interior simplices, with the same
stabilisers). Two strict pairs induce \emph{homotopic} $\tau_{\mathrm{rel}}$ over
$\Ga''$: with $H$, $N_{\partial}$ as in Proposition~\ref{prop:coneindep}(2), the map
$Hj-jN_{\partial}$ satisfies $\partial(Hj-jN_\partial)+(Hj-jN_\partial)\partial=0$, so by
Lemma~\ref{lem:etaunique} (applied to the two pair homotopies $0$ and
$jN_{\partial}-Hj$ for $(\tau^{s\prime},\taub')$) it equals $\partial\Lam-\Lam\partial$
for a unique $\Lam$; extending $\Lam$ by zero on interior vertices to
$\widehat\Lam\colon C_0(\Xbulk)\to C_2(\Xbulk)$ (a (T1)-map, the boundary vertex set
being $\Ga''$-stable) and replacing $H$ by $H':=H-(\partial\widehat\Lam-\widehat\Lam\partial)$
gives a twisted homotopy $\tau^{s\prime}-\tau^{s}=\partial H'+H'\partial$ with
$H'j=jN_{\partial}$, i.e.\ a homotopy \emph{of pairs}, which descends to the relative
complex. So the homotopy class of $\tau_{\mathrm{rel}}$, and with it
$(\conj_g)_{\ast}$, does not depend on the strict pair; and by
Proposition~\ref{prop:coneindep} every admissible triple gives the same operator on
$\Hcpt{\bullet}$ as every strict pair.

\emph{(ii) The Steinberg form.} \PROVED{} (sketch; the derived-category bookkeeping is
routine and is the part a formalisation would have to supply). $\Xbulk$ is contractible
and $\Xbdry$ is a disjoint union of contractible components indexed by
$\mathbb{P}^1(\Q)$ (\ref{H4}(ii), \ref{B2}), so the long exact homology sequence of the
pair gives $H_q(\Xbulk,\Xbdry)=0$ for $q\neq1$ and
$H_1(\Xbulk,\Xbdry)\cong\ker\bigl(\aug\colon\Z[\mathbb{P}^1(\Q)]\to\Z\bigr)=:\mathrm{St}$,
the Steinberg module of $\GL_2(\Q)$, as $\Z[\Ga]$-modules; hence
$C_{\bullet}(\Xbulk,\Xbdry)\simeq\mathrm{St}[1]$ in the derived category of
$\Z[\Ga]$-modules. Since $C_{\bullet}(\Xbulk,\Xbdry)\otimes\cring$ is a bounded complex of
relatively projective modules and $\mathrm{St}$ is $\Z$-free,
\[
  \Hcpt{q}\;\cong\;H^{q}\bigl(\Hom_{\Z[\Ga]}(C_{\bullet}(\Xbulk,\Xbdry),\VR)\bigr)
  \;\cong\;\Ext^{q}_{\Z[\Ga]}\bigl(\mathrm{St}[1],\VR\bigr)
  \;\cong\;H^{q-1}\bigl(\Ga;\Hom_{\Z}(\mathrm{St},\VR)\bigr)
  \qquad(6\in\cring^{\times}),
\]
\STD{\cite[Ch.~VII.7]{brown} for equivariant cohomology of a $\Ga$-complex with
invertible stabiliser orders being computed by $\Hom_{\Ga}(C_{\bullet},-)$, as in
Remark~\ref{rmk:orbifold}; \cite[\S\S10.7, 5.7]{weibel} for hypercohomology and
$\Ext$ of a complex}. The monoid $\De$ acts on $\mathbb{P}^1(\Q)$ by M\"obius
transformations and hence on $\mathrm{St}$, and $\Hom_{\Z}(\mathrm{St},\VR)$ is a
$\De$-module via $(\delta\cdot f)(s):=\rhon(\delta)f(\delta^{-1}s)$, $\delta^{-1}$ acting
through $\GL_2(\Q)$. Under the displayed isomorphisms, a twisted chain map of pairs
induces on $H_1=\mathrm{St}$ the map $s\mapsto g^{-1}s$ (Lemma~\ref{lem:bdrycomponent}:
the component map of $\taub$ is $c\mapsto g^{-1}c$), so
$\cor\circ(\conj_g)_{\ast}\circ\res$ on the relative complex corresponds to the
double-coset operator of $\Ga g\Ga$ on $H^{\bullet}(\Ga;\Hom_{\Z}(\mathrm{St},\VR))$ in the
sense of \cite[\S1]{ash_stevens}. \textbf{Consequently the Hecke relations of
Theorem~\ref{thm:relations} hold for $T^{c}_{\Ga g\Ga}$ on $\Hcpt{\bullet}$}, by the same
argument and on the same external inputs as that theorem
\STD{\cite[Thm.~3.24]{shimura}; \cite[\S1]{ash_stevens}}, since those inputs concern
$H^{\bullet}(\Ga;M)$ for an \emph{arbitrary} $\De$-module $M$ and $\Hom_{\Z}(\mathrm{St},\VR)$
is one. This is the sense in which ``the Hecke relations hold on $\Hcpt{\bullet}$ on
the nose'' is a theorem here and not a citation of geometry.

\emph{(iii) The geometric identification.} That
$H^{\bullet}(\Hom_{\Z[\Ga]}(C_{\bullet}(\Xbulk,\Xbdry),\VR))$ is the cohomology of the
Borel--Serre pair $(\Ga\backslash\overline{\Hyp},\partial)$ with coefficients in the
local system of $\VR$, that this is the compactly supported cohomology of the open
orbifold $\Ga\backslash\Hyp$, and that the operator of (i) is the action of the Hecke
correspondence on that pair, is \STD{\cite[\S\S1--2]{harder_eisenstein}, where the
long exact sequence $H^{\bullet}_c\to H^{\bullet}\to H^{\bullet}(\partial)$ and the Hecke
action on all three terms are set up for $\GL_2$ over a number field, the case here being
$\Q$; \cite{borel_serre} for the bordification} --- \emph{with the caveat}, per rule~\ref{A2}, that this document has
not verified the section numbers of \cite{borel_serre} --- nor those of
\cite{harder_eisenstein} --- against the text, and cites both works as a whole; where a
precise pointer is required this half is to be read as \CONJ{} \emph{as a citation}.
(The \emph{bibliographic} data of \cite{harder_eisenstein} \emph{have} been checked, in
v0.16.0: the entry had carried a wrong title until then.) Nothing in this document depends on (iii): every
statement made about $\Hcpt{\bullet}$ is proved on the complexes of
Definition~\ref{def:pair}, and (iii) only names what they compute.

\emph{(iv) Poincar\'e duality.} The adjointness of $T_p$ on $\Hcpt{1}(\VR)$ and on
$\Heq{1}(K,\Rel;\Dn)$ under a perfect pairing $\Hcpt{1}\times\Heq{1}(\Dn)\to\Hcpt{2}(\cring)
\cong\cring$ is Claim~\ref{cl:pd}, a stub of \S\ref{sec:six} at this version; its
consequence for this chapter, Proposition~\ref{prop:coneroutes}(3), is stated
conditionally on it and carries \OPENSTAT{} until it is written.
\end{remark}

\begin{remark}[$\etapair\equiv0$ iff $\tau j=j\taub$; the entrywise identity decides nothing]\label{rmk:conestrict}
\PROVED{} (Rule~\ref{A6}.) Theorem~\ref{thm:conewelldef}(d) says when $\hder=0$ is
admissible: exactly when $\tau$ is a map of pairs. Two things must be kept apart from
it.
\begin{enumerate}
\item \textbf{The shipped $\tau$ is not a map of pairs.} \COMP{} Verified for this
  chapter on the emitted table \file{docs/implementation/tau\_table\_p2.txt} ($p=2$,
  $g=\mathrm{diag}(1,2)$, $m=3$, every fill over $\Z$), reading its stored slots
  against the vertex and edge lists of Definition~\ref{def:model7}. \emph{On the
  canonical $\tau$ since v1.1.0} (policy \code{GeometricNearest}, the $\inv$-symmetric
  re-baseline; the counts re-read at v1.4.2 by slot-2 WP-7's cusp-support walk,
  \code{hecke::cusp\_support\_walk}, and confirmed by hand on the table): of the $12$
  stored slots at boundary vertices $v_0,\dots,v_3$, \textbf{$10$} map to an interior
  vertex --- every $j$ at $v_0$ and $v_3$ (all to $v_6=i$; e.g.\
  $\tau(s_1\lift{v_0})=\smat{-1}{-1}{1}{0}\cdot\lift{v_6}$, \code{:13}) and $j\in\{2,3\}$
  at $v_1$, $v_2$ (to $v_4=\zeta^{2}$ and $v_6$), the two $j=1$ slots at $v_1$, $v_2$
  staying on $\Xbdry$ (both to $v_0$); of the $9$ stored slots at boundary edges,
  \textbf{$5$} carry interior edges (\code{:34}, \code{:35}, \code{:40}, \code{:49},
  \code{:50} --- the $j=1$ slots at all three boundary edges carry $\{v_0,v_4\}$ and
  $\{v_4,v_6\}$, the $j=2$ slots at $\{v_0,v_1\}$ and $\{v_2,v_3\}$ carry $\{v_4,v_6\}$),
  the other four being the zero chain. (v1.3.3--v1.4.1 gave $8$ of $12$ and $2$ of $9$
  with ``the $j=1$ slots all stay on $\Xbdry$'': the pre-v1.1.0 table.) A different
  statistic on the same table is how many slots carry a \emph{non-zero} $\etapair$ once
  the $\etapair$-table exists (v1.4.2, slot-2 WP-8): $11$ of the $12$ vertex slots and
  $6$ of the $9$ edge slots --- more than leave the boundary, because
  $\etapair(s_j\sigma)\neq0$ also where $\tau(s_j\sigma)$ stays on $\Xbdry$ but differs
  from $j\taub(s_j\sigma)$. Both counts are $\ell$- and $n$-independent. So
  $\tau j\neq j\taub$ for \emph{every} $\taub$,
  $\etapair\equiv0$ is unavailable, and the $\etapair$-table is genuinely needed. This
  is a property of Algorithm~\ref{alg:tau}'s nearest-vertex policy (step~1), not of
  the model: Proposition~\ref{prop:strictpair} builds a strict $\tau^{s}$ on the same
  model.
\item \textbf{$\rres\Th_g=\Thb_g\rres$ on cochains does not license $\hder=0$.} Suppose
  the identity holds entrywise for some $(\tau,\taub)$ with $\tau j\neq j\taub$. Then
  $\hder=\widetilde{\etapair}$ satisfies $d\hder+\hder d=0$
  (Lemma~\ref{lem:derivedhomotopy}), i.e.\ $\hder$ is a degree-$(-1)$ cochain map, and
  by the computation of Proposition~\ref{prop:coneambiguity}(1),
  $\Phig^{\ast}-(\Th_g,\Thb_g)_{\ast}=\coneinc{}_{\ast}\circ\hder_{\ast}\circ\coneproj{}_{\ast}$
  on $\Hcpt{\bullet}$, an element of $\Hom(\ker\rres^{q+1},\coker\rres^{q})$ that
  \emph{need not vanish}. Whether it vanishes for some actual $(\tau,\taub)$ with the
  identity is \CONJ{} either way and is not decided here; what is \PROVED{} is that
  the identity does not imply it. The correct rule is (d): $\hder=0$ is admissible iff
  the $\etapair$-table is identically zero, i.e.\ iff $\tau j=j\taub$ \emph{exactly}.
  The quantities \code{eta\_zero} $:=(\etapair\equiv0)$ and \code{cochain\_commutes}
  $:=(\rres\Th_g=\Thb_g\rres$ entrywise$)$ are therefore \textbf{recorded diagnostics
  and gate nothing}; \code{eta\_zero} implies \code{cochain\_commutes}, not conversely.
\end{enumerate}
\end{remark}

\begin{proposition}[Why a free cochain-level $\hder$ will not do: the $\Hom(\ker\rres^{1},\coker\rres^{0})$ ambiguity]\label{prop:coneambiguity}
(Appendix~\ref{app:errata}, entries~\ref{err:e2} and \ref{err:e3}: the realisation is
stated with $Z^{0}$, never $B^{0}$, and \emph{over a field}.) Assume \ref{H4},
\ref{B1}--\ref{B3}, let $\Th_g$, $\Thb_g$ be as in Definition~\ref{def:coneaction}, and
let $\hder,\hder'\colon\Cq{\bullet}\to C^{\bullet-1}_{\partial}$ be \emph{any} two
$\cring$-linear maps with $d\hder+\hder d=d\hder'+\hder'd=\Thb_g\rres-\rres\Th_g$, so
that both $\Phi:=(\Th_g,\Thb_g+\hder)$ and $\Phi':=(\Th_g,\Thb_g+\hder')$ are cochain
maps of $\Cone{\rres}$ (Lemma~\ref{lem:conesign}).
\begin{enumerate}
\item \PROVED{} \textbf{(The ambiguity, over any commutative $\cring$.)}
  $\xi:=\hder'-\hder$ is a cochain map of degree $-1$ ($d\xi+\xi d=0$), and on
  $\Hcpt{q+1}=H^{q}(\Cone{\rres})$
  \[
    \Phi'_{\ast}-\Phi_{\ast}\;=\;\coneinc{}_{\ast}\circ\xi_{\ast}\circ\coneproj{}_{\ast}
    \;\colon\;\Hcpt{q+1}\twoheadrightarrow\ker\rres^{q+1}\xrightarrow{\ \xi_{\ast}\ }
    H^{q}(C^{\bullet}_{\partial})\twoheadrightarrow\coker\rres^{q}\hookrightarrow\Hcpt{q+1},
  \]
  an element of $\Hom_{\cring}(\ker\rres^{q+1},\coker\rres^{q})$; for $\Hcpt{1}=X_3$
  this is $\Hom(\ker\rres^{1},\coker\rres^{0})$.
\item \PROVED{} \textbf{(Realisation --- over a field only, and inside $\Ceq$.)} \LEAN{prop:coneambiguity} (over a field; not over $\Z/\ell^{e}$) Let
  $\cring=\F$ be a \emph{field}, and read every cochain space of the construction below
  as the \textbf{equivariant} one: $C^{1}:=\Cq{1}(K,\Rel;L)$, $B^{1}:=d\,\Cq{0}$,
  $Z^{1}:=\Zq{1}$ (v1.6.0, stated explicitly;
  \file{docs/implementation/slot3/WP-6\_report.md} finding~3 --- on the contractible
  $\Xbulk$ the flat reading is \emph{vacuous}, the flat complex $C^{\bullet}(K;\VR)$
  being acyclic in degree $1$ \eqref{eq:contractibleuse}, so that the flat $B^{1}$ would
  already contain every flat $1$-cocycle and a $\xi$ annihilating it would be zero on
  every cocycle).
  For every $\F$-linear $\lambda\colon\Heq{1}(K,\Rel;L)\to
  Z^{0}(C^{\bullet}_{\partial})=H^{0}(C^{\bullet}_{\partial})$ there is a degree-$(-1)$
  cochain map $\xi$ of $\Ceq$ with $\xi_{\ast}=\lambda$; hence for every
  $\mu\in\Hom_{\F}(\ker\rres^{1},\coker\rres^{0})$ there is a $\hder'$ satisfying the
  cone identity with $\Phi'_{\ast}-\Phi_{\ast}=\mu$ on $\Hcpt{1}$. \textbf{The
  statement is not claimed over $\Z/\ell^{e}$ or $\Z/m$}, where the splittings used
  below need not exist (entry~\ref{err:e3}).
\item \textbf{(Consequence.)} For a freely chosen cochain-level $\hder$ the operator on
  $X_3$ is undetermined to the full extent of $\Hom(\ker\rres^{1},\coker\rres^{0})$,
  which is non-zero in real cases (Remark~\ref{rmk:coneambiguity691}); $X_3$ would not
  be a Hecke module. For the \emph{derived} $\hder=\widetilde{\etapair}$ the ambiguity
  vanishes identically, over $\Z$, for every $\ell$ (Proposition~\ref{prop:coneindep}(1)):
  the difference of two derived homotopies is not merely a chain map but a cone
  \emph{homotopy}. That is the whole reason the datum is a chain-level $\etapair$ and
  not a cochain-level $\hder$.
\end{enumerate}
\end{proposition}

\begin{proof}
(1) $d\xi+\xi d=0$ by subtraction. Let $(c,c_\partial)\in Z^{q}(\Cone{\rres})$, i.e.\
$dc=0$ and $dc_\partial=-\rres c$; then $[c]\in\ker\rres^{q+1}$ (since $\rres c$ is a
coboundary) and $\coneproj{}_{\ast}[(c,c_\partial)]=[c]$; the image of
$\coneproj{}_{\ast}$ is $\ker\rres^{q+1}$ and the kernel of $\coneinc{}_{\ast}$ is
$\im\rres^{q}$ by exactness of \eqref{eq:coneles}. Now
$(\Phi'-\Phi)(c,c_\partial)=(0,\xi c)$, and $d(\xi c)=-\xi(dc)=0$, so
$[(0,\xi c)]=\coneinc{}_{\ast}[\xi c]=\coneinc{}_{\ast}\xi_{\ast}[c]$.

(2) Over a field choose complements \emph{inside the equivariant complex}:
$\Cq{1}=\Bq{1}\oplus\widetilde H^{1}\oplus W^{1}$
with $\widetilde H^{1}\subseteq\Zq{1}$ mapping isomorphically onto $\Heq{1}$ and
$d\colon W^{1}\to\Bq{2}$ an isomorphism. Define $\xi$ on $\Cq{1}$ by $\xi:=\lambda$ on
$\widetilde H^{1}$ (through $\widetilde H^{1}\cong\Heq{1}$) and $\xi:=0$ on
$\Bq{1}\oplus W^{1}$; define $\xi:=0$ on $\Cq{q}$ for $q\neq1$. Then $d\xi=0$ because
$\xi$ takes values in $Z^{0}(C^{\bullet}_{\partial})$ --- note $Z^{0}\subseteq
C^{0}_{\partial}$ is the kernel of $d$ and is in general a proper subspace
(entry~\ref{err:e2}), and it equals $H^{0}(C^{\bullet}_{\partial})$ because there are no
$0$-coboundaries; and $\xi d=0$ because $d\,\Cq{0}=\Bq{1}\subseteq\ker\xi$ and $\xi$
vanishes on $\Cq{2}$. So $d\xi+\xi d=0$ and $\xi_{\ast}=\lambda$ on $\Heq{1}$. Given
$\mu$, extend it to $\lambda\colon\Heq{1}\to H^{0}(C^{\bullet}_{\partial})$ by choosing a
complement of $\ker\rres^{1}$ in $\Heq{1}$ and a lift $\coker\rres^{0}\to
H^{0}(C^{\bullet}_{\partial})$ (a field again); then put $\hder':=\hder+\xi$ and apply (1).

(3) is (1), (2) and Proposition~\ref{prop:coneindep}(1).
\end{proof}

\begin{remark}[The ambiguity is non-zero at $(n,\ell)=(10,691)$]\label{rmk:coneambiguity691}
\COMP\textsuperscript{r} for the dimensions, \PROVED{} for $X_2$, and \OPENSTAT{}
\emph{in this document} for one input --- the reason $\rres^{1}$ has the rank it has,
which is a claim of the companion [E] and is not re-derived here. The arithmetic is
recorded so that Proposition~\ref{prop:coneambiguity}(3) has a witness. \textbf{Rewritten
at v1.5.0 from the measured values} (Appendix~\ref{app:changelog}, v1.4.2, deferred item~1;
slot-2 WP-7 and slot-3 WP-5): through v1.4.2 it was computed under the assumption
$\rres^{1}=0$, which is false here, and the old numbers are quoted at the end.

At $(n,\ell)=(10,691)$, $\cring=\F_{691}$, $p$ any good prime:
\begin{itemize}\setlength{\itemsep}{0pt}
\item $X_1=\Heq{0}=0$: $10$ is not in the Dickson semigroup
  $\langle\ell+1,\ell^{2}-\ell\rangle$ (Claim~\ref{cl:dimH1}(2),
  Remark~\ref{rmk:dickson}), so there is no invariant at all;
\item $X_2=H^{0}(C^{\bullet}_{\partial})=\VR^{T}\cong\F_{691}$ --- \PROVED{},
  Proposition~\ref{prop:bdrychar}(1) with $n<\ell$;
\item $\dim X_4=\dim\Heq{1}=3=2d_{12}-1$ (Claim~\ref{cl:dimH1}(1));
\item $\rank\rres^{0}=0$, since its source $X_1$ is $0$; hence
  $\coker\rres^{0}=X_2$ is $1$-dimensional;
\item $\rank\rres^{1}=1$, hence $\dim\ker\rres^{1}=\dim H^{1}_{!}=2$. This is the
  \emph{measured} value (\code{HeckeBoundaryParabolicValidation}, slot-2 WP-7;
  \code{HeckeConeValidation}, slot-3 WP-5, which finds it on all ten Kummer cells of
  [E, Claim 1.8], with $\dim H^{1}_{!}=2\dim S_{n+2}$), and it is where v1.4.2 and
  earlier were wrong.
\end{itemize}
\noindent
So by \eqref{eq:conedims}, $\dim X_3=\dim X_2-\dim X_1+\dim\ker\rres^{1}=1-0+2=3$ ---
which \code{HeckeConeValidation} computes directly as $\dim\Hcpt{1}=3$, the two routes
agreeing --- and $\dim\Hcpt{2}=\dim H^{1}(C^{\bullet}_{\partial})-\rank\rres^{1}=1-1=0$.
The ambiguity is
$\Hom_{\F}(\ker\rres^{1},\coker\rres^{0})\cong\F_{691}^{2\times1}=\F_{691}^{2}\neq0$:
\textbf{a free cochain-level $\hder$ leaves the $T_p$-action on this $3$-dimensional
$X_3$ undetermined by a $2$-dimensional space, and the derived one does not} ---
Proposition~\ref{prop:coneambiguity}(3)'s witness, which the repair does not touch. Only
the two dimensions move: the conclusion is the same, and is the reason this remark
exists.

\emph{Why $\rres^{1}\neq0$, and what ``the Eisenstein class'' means mod $\ell$.} $691$
divides the numerator of $B_{12}/12$, which is what makes $(10,691)$ a Kummer cell; the
inference drawn from it through v1.4.2 --- that the Eisenstein class restricts to zero,
so $\rres^{1}=0$ --- does not hold mod $\ell$, because there is no single Eisenstein
class there. The joint $1+p^{\,n+1}$-eigenspace $E$ of $\{T_2,T_3\}$ on $\Heq{1}$ is
\emph{two}-dimensional at every Kummer cell (the integral Eisenstein and one cuspidal
eigenline reduce into the same socle of a non-split block; [E] v1.4.0
\code{cl:harder}), $\rres^{1}$ vanishes on all of $E$ and is non-zero on the top of the
block, so $\rank\rres^{1}=1$. Wherever this document speaks of ``the Eisenstein class''
modulo $\ell$ --- Obligation~\ref{obl:oracle}'s last factor above all --- read the joint
eigenspace $E$. Obligation~\ref{obl:iota}(v) already carries the qualification this
needs, locating the Eisenstein class only \emph{where that eigenspace is a line}, or
through restriction to the cusp where $\rank\rres^{1}=1$; what a Kummer cell removes is
the first of those two routes, $E$ being two-dimensional there. This is the input that
keeps the \OPENSTAT{} tag: [E]'s claim is not proved in this document.

\emph{What v1.4.2 and earlier said}, kept so that the repair is visible and a reader of
the old PDF is not left guessing: ``$\rres^{1}=0$ \dots{} so $\ker\rres^{1}=X_4$,
$\Hom(\ker\rres^{1},\coker\rres^{0})\cong\F_{691}^{3}$, and $\dim X_3=1-0+3=4$''. The
$4$ propagated to [X] \S6.2 E5 acceptance (c), which slot~3 corrected ([X] v1.7). The
measured $\dim\Hcpt{1}=3$ at $(10,691)$ is \code{HeckeConeValidation}'s (v1.5.2).
\end{remark}

\paragraph{What this chapter does \emph{not} say about $\inv$ and the cone action
(v1.6.0; an open point, recorded).}
\S\ref{sec:involution} builds $\inv^{\ast}$ on $\Ceq(K,\Rel;L)$ and, on the boundary
model, $\inv^{\ast}_{\partial}$; $\rres^{q}\inv^{\ast}=\inv^{\ast}_{\partial}\rres^{q}$ is
an identity of the stored matrices, and with it
$\inv_{c}:=\inv^{\ast}\oplus\inv^{\ast}_{\partial}$ satisfies
$\Dcone\inv_{c}=\inv_{c}\Dcone$ and $\inv_{c}^{2}=\id$, so $\inv_{c}$ is an involution of
$\Cone{\rres}$ and acts on $\Hcpt{\bullet}$. \textbf{That $\inv_{c}$ commutes with
$\Acone{q}$ does not follow from any statement of this document.} $\Phig$ is
$(\Th_g,\Thb_g+\hder)$ by Definition~\ref{def:coneaction}, and
$\inv_{c}\Phig=\Phig\inv_{c}$ needs, beyond the two bulk and boundary commutations of
Proposition~\ref{prop:iotacommutes} and Obligation~\ref{obl:taubdry}, the third identity
$\inv^{\ast}_{\partial}\hder=\hder\,\inv^{\ast}$ on the stored $\Ceq$ basis --- i.e.\ that
the pair homotopy $\etapair$ intertwines $\inv_{\Xbulk}$ and $\inv_{\Xbdry}$, up to the
homotopy Lemma~\ref{lem:etaunique} allows. Neither Lemma~\ref{lem:etaexists}'s
construction nor Algorithm~\ref{alg:taubdry}'s vertex policy is built to respect $\inv$
--- the $\inv$-symmetry that \emph{is} imposed is on $\tau$ alone
(Obligation~\ref{obl:iota}(iv), the policy of Algorithm~\ref{alg:tau} step~1 since
v1.1.0) --- so the identity is neither proved nor expected on a general choice of
$\etapair$. \emph{Measured} (\code{HeckeInvolutionValidation},
hecke-extensions slot-4 WP-3, run 2026-09-21): at $(n,\ell)=(10,691)$ the commutator
$[\inv_{c},\Acone{1}(T_2)]$ vanishes, on the one degree with $\Hcpt{q}\neq0$; and on the
$15$-cell cone grid $\ell\in\{2,3,5,7,691\}\times n\in\{0,2,10\}$ both
$\rres^{q}\inv^{\ast}=\inv^{\ast}_{\partial}\rres^{q}$ and
$\Dcone\inv_{c}=\inv_{c}\Dcone$ are \emph{asserted} and hold, while
$\dim H^{q}_{c,\pm}$ is \textbf{recorded} (v1.6.1: that grid's arm builds $\inv_{c}$,
$\Dcone$ and the $\pm$ dimensions and no cone Hecke action at all, so the commutator is
recorded at $(10,691)$ \emph{alone}, as the next sentence says and as v1.6.0's
``$\dim H^{q}_{c,\pm}$ and the commutator are recorded'' did not). One cell is not evidence,
and rule~\ref{A6} applies: the recorded commutator licenses nothing, and no code path
may branch on it; the tagged record of this measurement is
Remark~\ref{rmk:iotapairing}. \textbf{For the next $X$}: either a clause of
Obligation~\ref{obl:eta} requiring $\inv^{\ast}_{\partial}\hder=\hder\,\inv^{\ast}$ where
an $\inv$-equivariant $\etapair$ has been constructed, or a counterexample cell. Until
then $[\inv_{c},\Acone{q}]=0$ is asserted nowhere in this document and must be asserted
nowhere in an implementation.

\subsection{Integrality: one integral cochain map, and the Bocksteins}\label{sec:coneintegral}

\begin{proposition}[Integrality: one integral cochain map $\Phi_{g,\Z}$; Bockstein equivariance]\label{prop:coneintegral}
Assume \ref{H4} and \ref{B1}--\ref{B3}, and let $(\tau,\taub,\etapair)$ be a Hecke datum
of the pair \emph{over $\Z$} --- $\tau$ with every fill over $\Z$
(Remark~\ref{rmk:tauintegral}), $\taub$ and $\etapair$ unconditionally
(Lemmas~\ref{lem:taubdryexists}, \ref{lem:etaexists}).
\begin{enumerate}
\item \PROVED{} \textbf{(Tables once.)} $\tau$, $\taub$, $\etapair$ are tables of
  integer chains, independent of $\cring$ and of $n$. Writing
  $\Cone{\rres}_{\Z}$ for the cone over $\cring=\Z$, its terms
  $\Hom_{\Z[\Ga]}(C_q(\Xbulk),\Vcoef)\cong\bigoplus_{\sigma}\Vcoef^{\Ga_\sigma}$ and
  $\Hom_{\Z[\Ga]}(C_q(\Xbdry),\Vcoef)\cong\bigoplus_{\sigma}\Vcoef^{\{\pm I\}}$ (sums over
  $\Ga$-orbit representatives, finitely many by \ref{H4}(iii), \ref{B1}) are finitely
  generated \emph{free} $\Z$-modules, each summand being a saturated sublattice of
  $\Vcoef=\Z^{n+1}$; and $\Phi_{g,\Z}:=\Phig$ over $\Z$ is a single cochain map of
  $\Cone{\rres}_{\Z}$ with integer matrix entries.
\item \PROVED{} \textbf{(Matrices per $\cring$.)} For every commutative $\cring$ the
  reduction $\Cone{\rres}_{\Z}\otimes\cring\to\Cone{\rres}_{\cring}$ is a map of complexes
  carrying $\Phi_{g,\Z}\otimes\id$ to $\Phig$ over $\cring$; it is an isomorphism when
  $6\in\cring^{\times}$. The matrices $\Th_g$, $\Thb_g$, $\hder$ over $\Z$ have entries of
  size $\sim N^{n}$ and are never stored; over $\cring=\Z/m$ they are formed from the
  tables and the reduced $\rhon(g\gamma_t)\bmod m$, exactly as \eqref{eq:xiexplicit}
  does for $\Th_g$.
\item \PROVED{} \textbf{(Reductions and Bocksteins.)} For $a,b\ge1$ the reduction maps
  $\Cone{\rres}_{\Z}\otimes\Z/\ell^{a+b}\to\Cone{\rres}_{\Z}\otimes\Z/\ell^{a}$ and the short
  exact sequences of complexes
  $0\to\Cone{\rres}_{\Z}\otimes\Z/\ell^{a}\to\Cone{\rres}_{\Z}\otimes\Z/\ell^{a+b}\to
  \Cone{\rres}_{\Z}\otimes\Z/\ell^{b}\to0$ are $\Phi_{g,\Z}$-equivariant; hence so are
  the induced maps on cohomology and the connecting (Bockstein) homomorphisms
  $H^{q}(\Cone{\rres}_{\Z}\otimes\Z/\ell^{b})\to H^{q+1}(\Cone{\rres}_{\Z}\otimes\Z/\ell^{a})$.
  \emph{On the nose}, not up to homotopy.
\item \PROVED{} \textbf{(Independence, on every page.)} Two Hecke data of the pair over
  $\Z$ give $\Phi'_{g,\Z}=\Phi_{g,\Z}+\Dcone\Upsilon+\Upsilon\Dcone$ with $\Upsilon$
  \emph{integral} and preserving the filtration $\ell^{k}\Cone{\rres}_{\Z}$, provided
  $6\in\cring^{\times}$ for the $\tau$-half (for two $\etapair$ and two $\taub$ with the
  same $\tau$, unconditionally). Consequently the two induce the same maps on every page
  $E_r$, $r\ge1$, of the spectral sequence of the $\ell$-adically filtered complex
  $\Cone{\rres}_{\Z}$ \STD{\cite[\S5.4]{weibel} for the spectral sequence of a filtered
  complex}, so the pages carry a well-defined $T^{c}_{\Ga g\Ga}$-action.
\item \textbf{(What a mod-$m$ fill would cost --- and why it cannot occur here.)} A
  table solved only modulo $m$ defines $\Phig$ over $\Z/m$ only, and a record built on
  it must say so. For $\taub$ there is no fill (Lemma~\ref{lem:taubdryexists}(2)); for
  $\etapair$ the edge fills are \emph{unique} over $\Z$ and over every $\Z/m$
  \eqref{eq:contractibleuse}, so a mod-$m$ solve returns the reduction of the integer
  fill and the fallback is never mathematically necessary; the only way a
  non-integral $\etapair$-table can arise is a fixed-width integer tier that fails to
  find the integer fill. \textbf{Owner decision OP2 (2026-09-15) applies}: the fallback
  is forbidden by default, the tier is widened instead, and a record claiming
  integrality must have solve regime $\Z$ on all three tables
  (Remark~\ref{rmk:tauintegral}, Obligation~\ref{obl:eta}(iv)).
\end{enumerate}
\end{proposition}

\begin{proof}
(1) The tables are integer chains by hypothesis and involve neither $\cring$ nor $n$; the
coefficients $\rhon(\delta_j)$, $\rhon(g\gamma_t)$ are integral (Definition~\ref{def:semigroup}).
An equivariant cochain is determined by its values on orbit representatives, each value
constrained to $\Vcoef^{\Ga_\sigma}=\bigcap_{u\in\Ga_\sigma}\ker(\rhon(u)-\id)$, a
kernel of a map between free $\Z$-modules of finite rank, hence a saturated sublattice
and free. (2) The map $\Vcoef^{\Ga_\sigma}\otimes\cring\to(\VR)^{\Ga_\sigma}$ is injective
because a saturated sublattice is a direct summand, and surjective when
$\abs{\Ga_\sigma}\in\cring^{\times}$: for $f\in(\VR)^{\Ga_\sigma}$ with any lift
$\hat f\in\Vcoef$, $\sum_{u\in\Ga_\sigma}\rhon(u)\hat f\in\Vcoef^{\Ga_\sigma}$ reduces to
$\abs{\Ga_\sigma}f$. With $\abs{\Ga_\sigma}\mid12$ (\ref{H4}(iv), \ref{B3}) this is
$6\in\cring^{\times}$. Compatibility with $\Phig$ is that \eqref{eq:cosettransport} is
a formula with integer coefficients. (3) Freeness of the terms (1) makes
$-\otimes_{\Z}\Cone{\rres}_{\Z}$ preserve the exactness of $0\to\Z/\ell^{a}\to\Z/\ell^{a+b}
\to\Z/\ell^{b}\to0$; $\Phi_{g,\Z}\otimes\id$ commutes with maps of the form
$\id\otimes(\text{ring map})$; the connecting homomorphism is natural
\STD{\cite[Thm.~1.3.1]{weibel}}. (4) Proposition~\ref{prop:coneindep}(1),(2) with
integer tables give integer $\Upsilon$; an integral map preserves
$\ell^{k}\Cone{\rres}_{\Z}$. A filtration-preserving homotopy induces, on each graded
piece $\ell^{k}\Cone{\rres}_{\Z}/\ell^{k+1}\Cone{\rres}_{\Z}$, a homotopy between the two
induced maps, so they agree on $E_1=H(\mathrm{gr})$; the maps on $E_{r+1}=H(E_r,d_r)$ are
induced from those on $E_r$, so they agree on every page. (5) is
Lemma~\ref{lem:taubdryexists}(2), \eqref{eq:contractibleuse} over $\Z/m$ (the complex
$C_{\bullet}(\Xbulk)\otimes\Z/m$ is still acyclic, being the base change of a contractible
complex of free modules), and the cited policy.
\end{proof}

\begin{remark}[$6\notin\cring^{\times}$ on this model, and the truncation policy]\label{rmk:conetorsion}
At $\ell\in\{2,3\}$ the complexes of Definition~\ref{def:pair} are the \emph{model}
complexes of Remark~\ref{rmk:ell2window}(1): finite, terminating, and computing the
orbifold model cohomology of \ref{H3} and its boundary and cone analogues --- \emph{not}
$H^{\bullet}(\Ga;\VR)$, $H^{\bullet}(\Gbdry;\VR)$ or the compactly supported cohomology
of the orbifold. On them, Theorem~\ref{thm:conewelldef}(a),(b),(d),
Lemma~\ref{lem:conerelative}, Proposition~\ref{prop:coneintegral}(1),(3) and the
$\etapair$/$\taub$-halves of Proposition~\ref{prop:coneindep} hold verbatim, none having
used a hypothesis on $\ell$; $T^{c}_{\Ga g\Ga}$ is then \emph{defined} as the output of
the pipeline with the canonical triple, and its independence of $\tau$, its
identification with any correspondence, and the Hecke relations are not asserted ---
the contract box after Lemma~\ref{lem:tauexists}, applied to the cone. The systematic
treatment is \S\ref{sec:orbifoldmodel}, where the same definitions are to be repeated on
a free $\Ga$-model of the pair; there the complexes are unbounded above at $\ell=2$
(Remark~\ref{rmk:ell2window}(2),(3)) and the truncation policy of construction
item~\ref{open:ell2}, decided by the owner on 2026-09-16 (\ref{oc1a}), governs. \textbf{This chapter reads that tag at two sites}, listed
in \S\ref{sec:oc1sites}: Obligation~\ref{obl:coneexact}, whose ``every computed degree''
and whose exactness clause are narrowed by \ref{open:ell2}; and
Definition~\ref{def:coneaction}, which is stated for all $q\in\Z$ and, on a truncated
model, is to be read in degrees $\le q_{\max}$ with $\Cone{\rres}$ built to
$q_{\max}+1$. On the model complexes of this chapter the tag is vacuous
(\code{truncated = false}).
\end{remark}

\subsection{The two cross-routes and the Poincar\'e-duality gate}\label{sec:coneroutes}

\begin{proposition}[The two cross-routes and the Poincar\'e-duality gate]\label{prop:coneroutes}
Assume \ref{H4} and \ref{B1}--\ref{B3}, let $\cring$ be a commutative ring and
$(\tau,\taub,\etapair)$ a Hecke datum of the pair over $\Z$. The following are
\emph{cross-checks} of Definition~\ref{def:coneaction}, never definitions; the cone of
\eqref{eq:conediff} remains normative. Both routes (1) and (2) are committed validation
arms, and (3) is the acceptance gate for $X_3$ at $\ell\ge5$ --- owner decision OP3,
2026-09-15, which \emph{amended} the verdict's ``one of the two'' to both.
\begin{enumerate}
\item \PROVED{} \textbf{(Strict-pair route; needs $6\in\cring^{\times}$.)} Let $\tau^{s}$
  be a strict datum for $\taub$ (Proposition~\ref{prop:strictpair}) and
  $\Phig^{s}:=(\Th^{s}_g,\Thb_g)$ with $\etapair=0$. Then $\Phig^{s}$ and $\Phig$ induce
  the same $T^{c}_{\Ga g\Ga}$, hence the same $\Acone{q}$ in the same stored basis, for
  every $q$ --- by Proposition~\ref{prop:coneindep}(2), which is where
  $6\in\cring^{\times}$ enters; at $\ell\in\{2,3\}$ agreement is \emph{recorded}, not
  asserted. On $(\ker\rres)[1]$ the operator is $\Th^{s}_g|_{\ker\rres}$, so this route
  computes $\Hcpt{\bullet}$ as relative cohomology with the bulk pipeline run on
  $\ker\rres$.
\item \PROVED{} \textbf{(Relative-complex route, ``route 3''; no hypothesis on $\ell$.)}
  With $\mathrm{ext}_0$ of Lemma~\ref{lem:conerelative}(1) put
  $E:=\mathrm{ext}_0\circ\hder\colon\Cq{q}\to\Cq{q-1}$ and
  \begin{equation}\label{eq:route3}
    \Th'_g:=\Th_g+dE+Ed .
  \end{equation}
  Then $\Th'_g$ is a cochain map, cochain-homotopic to $\Th_g$ (so it induces the same
  $\Amat{q}$), satisfies $\rres\Th'_g=\Thb_g\rres$ \emph{on the nose}, hence preserves
  $\ker\rres$; and $\Phig':=(\Th'_g,\Thb_g)$ with $\hder=0$ is cone-homotopic to $\Phig$
  via $\Upsilon(c,c_\partial):=(Ec,0)$:
  $\Phig'-\Phig=-(\Dcone\Upsilon+\Upsilon\Dcone)$. Consequently $T^{c}_{\Ga g\Ga}$ is the
  map induced by $\Th'_g|_{\ker\rres}$ on $H^{\bullet}(\ker\rres)\cong\Hcpt{\bullet}$,
  for the \emph{same} triple, over $\Z$ and for every $\ell$; and the $\Acone{q}$ so
  computed equals the normative one.
\item \textbf{(Poincar\'e-duality gate; $6\in\cring^{\times}$, conditional on
  Claim~\ref{cl:pd}.)} \OPENSTAT{} If the pairing of Proposition~\ref{prop:pairing}
  induces a perfect pairing $\Hcpt{1}(K,\Rel;\Vcoef_{\cring})\times\Heq{1}(K,\Rel;\Dn)\to
  \cring$ under which $T_p$ is self-adjoint (Claim~\ref{cl:pd}, with
  $\adjg(g)=\mathrm{diag}(N,1)\in\Ga g\Ga$), then
  $\langle\Acone{1}x,y\rangle=\langle x,\Amat{1}(\Dn)\,y\rangle$ for all $x,y$, and
  $\Acone{1}$ is determined by $\Amat{1}(\Dn)$ --- an operator the bulk pipeline
  already computes. This is a \emph{complete independent pin} of $\Acone{1}$ and the
  acceptance gate once the pairing exists at cochain level (owner decision OP3). It is
  an oracle: a disagreement is a defect, never a redefinition. Its status here is that
  of Claim~\ref{cl:pd}, a stub at this version.
\end{enumerate}
\end{proposition}

\begin{proof}
(1) Proposition~\ref{prop:strictpair} gives $\tau^{s}$; Proposition~\ref{prop:coneindep}(2)
with $\tau'=\tau^{s}$ (and $\taub'=\taub$, $N_{\partial}=0$) gives the cone homotopy,
whose existence needs $H$ with $\tau^{s}-\tau=\partial H+H\partial$, i.e.\
Lemma~\ref{lem:tauexists}(2) --- the uniqueness half, which is why this route keeps
$6\in\cring^{\times}$ even though both $\tau^{s}$ and $\tau$ exist over $\Z$; the last sentence is Lemma~\ref{lem:conerelative}(3) and the
end of the proof of Proposition~\ref{prop:strictpair}.

(2) $E$ is $\cring$-linear and equivariant (Lemma~\ref{lem:conerelative}(1)), and
$\rres E=\rres\,\mathrm{ext}_0\,\hder=\hder$. Then
$\rres\Th'_g=\rres\Th_g+d(\rres E)+(\rres E)d=\rres\Th_g+d\hder+\hder d=\Thb_g\rres$ by
Lemma~\ref{lem:derivedhomotopy}, using $\rres d=d\rres$. $\Th'_g$ is a cochain map
because $dE+Ed$ is one ($d(dE+Ed)=dEd=(dE+Ed)d$), and $\Th'_g-\Th_g=dE+Ed$ is a cochain
homotopy, so $\Amat{q}$ is unchanged. For $c\in\ker\rres$, $\rres\Th'_gc=\Thb_g\rres c=0$.
Finally $\Dcone\Upsilon(c,c_\partial)=\Dcone(Ec,0)=(-dEc,\rres Ec)=(-dEc,\hder c)$ and
$\Upsilon\Dcone(c,c_\partial)=\Upsilon(-dc,\cdot)=(-Edc,0)$, so
$\Dcone\Upsilon+\Upsilon\Dcone=\bigl(-(dE+Ed)c,\ \hder c\bigr)=-(\Phig'-\Phig)(c,c_\partial)$.
Hence $\Phig'$ and $\Phig$ induce the same map on $H^{\bullet}(\Cone{\rres})$; $\Phig'$
preserves $(\ker\rres)[1]$ and acts there by $\Th'_g$, and the quasi-isomorphism of
Lemma~\ref{lem:conerelative}(3) transports the induced maps. No hypothesis on $\ell$ was
used. (3) is the stated conditional, the adjointness identity being the definition of
``self-adjoint'' read in the stored bases.
\end{proof}

\subsection{Obligations}\label{sec:coneobligations}

\noindent
The four certificates of the extension programme's data model, and what each one is
\emph{for}. The first two are the content; the last two are consequences re-checked.
\begin{center}
\begin{tabular}{@{}p{0.20\textwidth}p{0.36\textwidth}p{0.07\textwidth}p{0.29\textwidth}@{}}
\toprule
certificate & stores / \code{verify()} re-derives & ring & what it pins \\
\midrule
\code{TauBoundary\-Certificate} & the $\taub$-table; Obligation~\ref{obl:taubdry}(i)--(iv) & $\Z$ & $\Thb_g$ up to homotopy (Lemma~\ref{lem:taubdryexists}(3)) \\
\code{PairHomotopy\-Certificate} & the $\etapair$-table; Obligation~\ref{obl:eta}(i)--(iv) & $\Z$ & \textbf{the operator on $\Hcpt{\bullet}$} (Prop.~\ref{prop:coneindep}(1)) \\
\code{ConeAction\-Certificate} & $\Thb_g$, $\hder$, $\Phig$, $\Acone{q}$ over $\cring$; Obligation~\ref{obl:eta}(v) and \ref{obl:coneexact}(i)--(iii) & $\cring$ & that the stored $\hder$ \emph{is} the transport of the stored $\etapair$ \\
\code{ConeCertificate} & $\Dcone$, the maps of \eqref{eq:coneles}, $\Amat{q},\Abdry{q},\Acone{q}$; Obligation~\ref{obl:coneexact}(iv)--(vi) & $\cring$ & the \emph{sequence}, not the module: any $\Phi$ of Prop.~\ref{prop:coneambiguity} passes it \\
\bottomrule
\end{tabular}
\end{center}

\begin{obligation}[Certified $\etapair$-table]\label{obl:eta}
\LEAN{chk:eta} \LEAN{chk:eta-p2} \LEAN{chk:pairgrid} (clauses (i)--(iv) as a predicate, verified at the same $(p,\nu)$; not (v), (vi); the link to Definition~\ref{def:pairhomotopy}: \LEAN{thm:etaspec-pairhomotopy} \LEAN{cor:etap2-pairhomotopy} (the latter at $(p,\nu)=(2,1)$) --- an accepted
table is a pair homotopy for every $\tau$, $\taub$, $\etapair$ realising the three tables on the generators $s_j\cdot\sigma$; that such realisations exist is not formalised)
Assume \ref{H4} and \ref{B1}--\ref{B3}. An implementation of the cone action emits, for
each $g=\mathrm{diag}(1,N)$ and each $\Ga''$-orbit representative simplex $\sigma$ of
$\Xbdry$, the chain $\etapair(s_j\cdot\sigma)=\sum_t n^{(j,\sigma)}_t\,\gamma_t\cdot\sigma_t$
with $n^{(j,\sigma)}_t\in\Z$, $\gamma_t\in\Ga$, $\sigma_t$ a simplex of $\Xbulk$ of
dimension $\dim\sigma+1$, in the stored-chain format of Algorithm~\ref{alg:tau} and in
canonical form, together with the \textbf{table hashes} of the stored $\tau$- and
$\taub$-tables it was built against and the solve regime of every fill. It must
discharge the following; each
clause names the data read, the predicate, the ring, and the action on failure.

\smallskip
\noindent\emph{The three hashes} (v1.5.0; \S\ref{sec:recordschema}): what identifies
``the tables it was built against'' in a record is a hash of each stored table ---
\code{table\_fnv64()} on the $\etapair$-table itself and the same function on each base
table, so that the three are one function and not three
(\file{Hecke/HeckePairHomotopy.h}: \code{table\_fnv64()}, \code{tau\_table\_fnv64()},
\code{tau\_boundary\_table\_fnv64()}; the byte-wise FNV-1a-64 of
\file{ModArithmetic/Fnv.h}). They are \emph{identifying and never certifying}: two
admissible $\etapair$ for the same $(\tau,\taub)$ differ in the first hash while
$\Acone{q}$ does not (Lemma~\ref{lem:etaunique},
Proposition~\ref{prop:coneindep}(1)), so a hash may be compared and recorded but not
branched on (rule~\ref{A6}), and clause~(i) --- not a hash --- is what ties the table to
its certified inputs.
\begin{enumerate}
\item[(i)] \textbf{Structural.} $\det\gamma_t=1$; every $\sigma_t$ is a simplex of
  $\Xbulk$ of the right dimension; and the tables it refers to are the certified
  $\tau$ (Obligation~\ref{obl:tau}) and $\taub$ (Obligation~\ref{obl:taubdry}) of the
  same $g$, $\Ga''$ and $s_j$. Over $\Z$; a failure \textbf{aborts}.
\item[(ii)] \textbf{The homotopy identity \eqref{eq:pairhomotopy}.}
  $\partial\etapair(s_j\sigma)=(j\taub-\tau j)(s_j\sigma)$ for every stored vertex slot,
  and $\partial\etapair(s_j\sigma)+\etapair(\partial(s_j\sigma))=(j\taub-\tau j)(s_j\sigma)$
  for every stored edge slot, \textbf{both sides evaluated only through each table's
  own (T1) helper, all built on the one (P1) split} from the stored $\tau$-, $\taub$-
  and $\etapair$-tables --- the same helpers the construction uses, so that
  construction and certification cannot
  disagree on how a non-representative slot is derived --- and compared after
  reduction to canonical $(n_t,\gamma_t,\sigma_t)$ form. Over $\Z$; a failure
  \textbf{aborts}. (v1.5.0, corrected from ``the shared (T1) helper --- the same
  helper the construction uses'': there is no one helper shared by the three tables.
  Each has its own --- \code{eval\_tau\_on\_translate}, \code{tau\_on\_translate},
  \code{eta\_on\_translate} --- and all three route through the one conjugation
  primitive of (P1), \code{HeckeCosetTables::twisted\_conjugate}, which is what the
  clause is about: that construction and certification share the derivation, not that
  one function serves three tables. On the shipped tables a $\tau$-value on a
  translate is never asked for, (T1) of the defect supplying it.)
\item[(iii)] \textbf{(T1) bookkeeping.} For every stored transporter $h\in\Ga''$
  occurring in the table, $\etapair(h\cdot x)=(g^{-1}hg)\cdot\etapair(x)$ as chains in
  canonical form --- Obligation~\ref{obl:tau}'s clause with $\etapair$ in place of
  $\tau$; and the entry bound of Lemma~\ref{lem:bounds} recounted on the
  $\etapair$-chains, so that the arithmetic tier of Theorem~\ref{thm:backend} covers
  the derived $\hder$. Over $\Z$; a failure \textbf{aborts}.
\item[(iv)] \textbf{Solve regime (owner decision OP2).} The recorded regime of every
  edge fill is $\Z$; a fill recorded as solved modulo $m$ marks the table
  \code{integral = false}, is \textbf{forbidden by default} and \textbf{aborts} under the
  default policy; where an implementation offers a mod-$m$ path at all it is behind an
  explicit setting, and then the record carries the
  flag, may not claim integrality, and states that $\Phig$ is defined over $\Z/m$ only
  (Proposition~\ref{prop:coneintegral}(5)). Since the fills are unique
  \eqref{eq:contractibleuse}, a mod-$m$ regime is a tier failure and never a
  mathematical one --- and therefore \textbf{an implementation MAY omit the mod-$m$
  regime altogether} (v1.5.0): it then has no such path and registers no such setting,
  a tier failure aborts unconditionally, and the recorded regime is $\Z$ by
  construction. The shipped implementation is of that kind
  (\file{Hecke/HeckePairHomotopy}, \file{Hecke/HeckeIntegerSolve}; \code{HeckeSettings}
  registers no key for it), and the absence of the setting is conformance, not a gap.
\item[(v)] \textbf{The derived $\hder$ is the transport of the stored $\etapair$.} The
  stored matrix $\hder$ over $\cring$ equals \eqref{eq:derivedh} recomputed from the
  $\etapair$-table and $\rhon(\delta_j\gamma_t)\bmod m$ --- the analogue of the
  \code{xi\_terms} recomputation for $\Th_g$ --- entrywise, in every computed degree.
  Over $\cring$; a failure \textbf{aborts}. \textbf{This clause, and not any identity
  among matrices, is what pins $\Acone{q}$}: a tampered $\hder'=\hder+\xi$ with $\xi$
  a degree-$(-1)$ cochain map passes every clause of Obligation~\ref{obl:coneexact}
  and fails this one (Proposition~\ref{prop:coneambiguity}).
\item[(vi)] \textbf{Diagnostics --- recorded, never branched on (rule~\ref{A6}).}
  \code{eta\_zero} $:=(\text{every stored }\etapair\text{-chain is }0)$;
  \code{cochain\_commutes} $:=(\rres\Th_g=\Thb_g\rres$ entrywise over $\cring)$; the
  three solve regimes. The implication \code{eta\_zero} $\Rightarrow$
  \code{cochain\_commutes} is asserted (its failure \textbf{aborts}, being a
  contradiction in the stored data); nothing else about them is asserted, and no
  code path may select $\hder=0$ on the strength of \code{cochain\_commutes}
  (Remark~\ref{rmk:conestrict}(2)).
\end{enumerate}
\end{obligation}

\begin{obligation}[Cone exactness and intertwining, re-derived over $\cring$]\label{obl:coneexact}
Assume \ref{H4} and \ref{B1}--\ref{B3}, and let the stored data be: the bases of
$Z^{q}(\Cone{\rres})$, $B^{q}(\Cone{\rres})$, $H^{q}(\Cone{\rres})$ and the projector, in the
flat layout of Algorithm~\ref{alg:assembly} on $\Cone{\rres}$; the matrices $\Dcone$,
$\rres$, $\Th_g$, $\Thb_g$, $\hder$, $\Phig$ over $\cring$; $\Amat{q}$, $\Abdry{q}$,
$\Acone{q}$; and the \textbf{three map families} $\coneproj{}_{\ast}^{\,q}$,
$\rres^{q}$, $\coneinc{}_{\ast}^{\,q}$ of \eqref{eq:coneles}, stored per degree in the
stored bases (v1.6.0, corrected from ``the six maps of \eqref{eq:coneles}'': the
sequence is degree-indexed and the code stores the three families per degree, while
``six'' counted the six-term window of Lemma~\ref{lem:conerelative}(5) --- on the model
complexes there are eight non-trivial map instances,
$\coneproj{}_{\ast}^{0},\rres^{0},\coneinc{}_{\ast}^{0},\coneproj{}_{\ast}^{1},\rres^{1},
\coneinc{}_{\ast}^{1},\coneproj{}_{\ast}^{2},\rres^{2}$, of which two are
zero-dimensional and $\coneinc{}_{\ast}^{2}$ lands in $\Hcpt{3}=0$;
\file{docs/implementation/slot3/WP-5\_report.md} finding~11). ``Every
computed degree'' means every $q$ in the reported range, which is all of
$\{0,1,2\}$ on the model complexes and is read from $q_{\max}$ under
\ref{open:ell2} on a truncated free model (\S\ref{sec:oc1sites}).
\begin{enumerate}
\item[(i)] \textbf{The cone identity.} $d\hder+\hder d=\Thb_g\rres-\rres\Th_g$ and
  $\Dcone\Phig=\Phig\Dcone$ \textbf{as matrices on the stored $\Ceq$ basis} ---
  equivalently, the bracket of Lemma~\ref{lem:conesign} annihilates the stored
  equivariant cochains --- in every computed degree; and $\Dcone^{2}=0$, which involves
  only $d$ and $\rres$ and holds on the flat spaces as well. Over $\cring$; a failure
  \textbf{aborts}.

  \smallskip
  \emph{Why the basis is named, and what is recorded instead} (v1.6.0;
  \file{docs/implementation/slot3/WP-6\_report.md} finding~1, the reading implemented).
  The coset transport of Definition~\ref{def:cosettransport} is defined on
  $\Hom_{\Z[\Ga]}(-,\VR)$, and the step that turns $\varphi(\gamma_t\sigma_t)$ into
  $\rhon(\gamma_t)\varphi(\sigma_t)$ --- the step that makes $\Th_g$, $\Thb_g$ and
  $\hder$ finite matrices at all --- \emph{is} the $\Ga$-equivariance of $\varphi$. The
  stored matrices are honest linear maps of the flat cochain spaces and agree with the
  matrix-free routes on every flat vector, but the two identities above are consequences
  of the transport only on the equivariant subspace, and \textbf{the same two identities
  on the full flat cochain space are recorded, never asserted}: on the shipped grid the
  bracket vanishes on $\Ceq$ in $424$ of $424$ computed degrees and on the full flat
  space in $29$ of them, and $\Dcone\Phig=\Phig\Dcone$ holds flat on $1$ of $212$
  operators. A record that reports the cone carries the two as the diagnostics
  \code{phi\_flat\_identity} and \code{residual\_flat\_zero}
  (\S\ref{sec:recordschema} item~6(e)), under rule~\ref{A6}: a \code{false} there is
  conformance and not a defect, and no code path may branch on either. The same qualifier governs the tampered $\hder'$ of clause (viii)(b), which is
  required to satisfy this clause on $\Ceq$ and not flat. (Through v1.5.2 the clause read
  ``as matrices, in every computed degree'', which an implementer can --- and the first
  one did --- read as a statement about the stored flat matrices, and which then aborts
  on correct data.)
\item[(ii)] \textbf{Assembly.} $\Phig$ maps every stored cocycle of $Z^{q-1}(\Cone{\rres})$
  into $Z^{q-1}(\Cone{\rres})$ (checked as $\Dcone\Phig z=0$), and $\Acone{q}$ is the
  projection of those images (Algorithm~\ref{alg:assembly} step~2, recomputed). Over
  $\cring$; a failure \textbf{aborts}.
\item[(iii)] \textbf{Intertwining.} Each map of \eqref{eq:coneles} in the stored bases
  commutes with the corresponding pair among $(\Abdry{q},\Acone{q},\Amat{q})$
  (Theorem~\ref{thm:conewelldef}(a)); in particular
  $\rres^{q}\Amat{q}=\Abdry{q}\rres^{q}$, which is the pair's form of
  Obligation~\ref{obl:taubdry}(vii). Over $\cring$ \textbf{at every $\ell$, the
  $\rres$-half included}; a failure \textbf{aborts} (v1.7.0; through v1.6.2 the
  $\rres$-half carried an $\ell\in\{2,3\}$ carve-out, asserted for
  $6\in\cring^{\times}$ and \emph{recorded} otherwise, which clause (i) already proved
  away).

  \smallskip
  \emph{Why the $\rres$-half needs no hypothesis on $\ell$ \emph{here}} (v1.7.0;
  \file{docs/implementation/slot5/WP-4\_report.md} \S4, which re-checked the hypothesis
  this proof uses against the shipped code rather than against prose; owner ruling of
  2026-09-22, \code{[OD]} slot-4 close-out 3). An implementation that reaches this clause
  has already discharged clause (i), which carries no hypothesis on $\ell$ and whose
  failure aborts: $d\hder+\hder d=\Thb_g\rres-\rres\Th_g$ on the stored $\Ceq$ basis, in
  every computed degree. Let $z\in\Zq{q}$ be a stored \emph{equivariant} cocycle --- an
  element of the span of the stored $\Ceq$ basis, which is where clause (i) holds and
  which is what Algorithm~\ref{alg:assembly} step~2 reads --- so that $dz=0$. Then
  \[
    \Thb_g\rres z-\rres\Th_g z=(d\hder+\hder d)z=d(\hder z),
  \]
  a coboundary of the boundary complex --- so $\rres\Th_g z$ and $\Thb_g\rres z$ have the
  same class, and since $\rres$ is a cochain map it carries $\Bq{q}$ into the boundary
  coboundaries, whence $\rres^{q}\Amat{q}=\Abdry{q}\rres^{q}$ on cohomology in the stored
  bases, both sides assembled by Algorithm~\ref{alg:assembly} step~2 as clause (ii)
  requires. Nothing in the argument inverts $2$, $3$ or $6$: $\hder$ is the coset
  transport \eqref{eq:derivedh} of an $\etapair$-table built over $\Z$ for every $\ell$
  (Obligation~\ref{obl:eta}(iv)), clause (i) is asserted at every $\ell$, and the
  projection is linear over $\cring$. The hypothesis is the \emph{existence of the stored
  $\hder$}, not the invertibility of $6$.

  \smallskip
  \emph{What this clause does not say.} It is an assertion about \textbf{the objects this
  document's implementation computes} --- $\Heq{\bullet}(K,\Rel;L)$, its boundary
  analogue $H^{\bullet}(C^{\bullet}_{\partial})$, and the stored $\rres^{\bullet}$,
  $\Amat{q}$, $\Abdry{q}$ --- and \textbf{nothing more}. At $\ell\in\{2,3\}$ those groups
  are the orbifold model cohomology and its boundary analogue and are \emph{not} the
  group cohomology of $\Ga$ and of $\Gbdry$ (Remark~\ref{rmk:conetorsion},
  \ref{open:ell2}); no statement of this document identifies the two there, and asserting
  the identity above at $\ell\in\{2,3\}$ does not assert, imply or license the
  corresponding statement for $H^{\bullet}(\Ga;\Mn)$ and its boundary restriction. A
  reader who transports it there is reading a clause about the computed complexes as a
  clause about group cohomology, which is exactly the step Remark~\ref{rmk:conetorsion}
  forbids. Recorded outcome, unchanged by the strengthening: $424/424$
  $(\text{cell},q)$ pairs, the $74$ operators at $\ell\in\{2,3\}$ included
  (Remark~\ref{rmk:conerestrictsmallell}).

  \smallskip
  \emph{The asymmetry with Obligation~\ref{obl:taubdry}(vii), which keeps its carve-out}
  (v1.7.0). (vii) states the same identity in the boundary chapter, whose stored data are
  $\taub$ and $\Thb_g$ alone: an implementation of the boundary datum has no
  $\etapair$ (Definition~\ref{def:pairhomotopy}) and no $\hder$ \eqref{eq:derivedh}, so
  the proof above is \emph{not available to it}, and (vii) therefore continues to assert
  the identity for $6\in\cring^{\times}$ and to \emph{record} it at $\ell\in\{2,3\}$. The
  two clauses do not contradict each other because they bind different implementations:
  (vii) binds an implementation of the boundary datum, (iii) binds an implementation of
  the \emph{pair}, which by clause (i) carries $\hder$. For an implementation that
  discharges both --- the shipped one --- the identity is asserted at every $\ell$ under
  (iii), and (vii)'s carve-out is vacuous for it. This is the only place in this document
  where a strictly stronger statement is available to the richer datum, and it is stated
  here rather than moved into (vii) so that the boundary chapter stays readable, and
  implementable, on its own.
  The other maps of \eqref{eq:coneles} were, and are, asserted at every $\ell$.
\item[(iv)] \textbf{Exactness --- where it is asserted.} Each composite of consecutive
  maps of \eqref{eq:coneles} \emph{both of whose end terms are reported} is zero: on
  the model complexes that is every composite, and on a truncated free model it is
  every composite in degrees $\le q_{\max}$ except the outgoing one at
  $H^{q_{\max}}(\partial)$, whose target $\Hcpt{q_{\max}+1}$ is scaffolding with no
  stored basis --- there only the incoming composite
  $\rres^{q_{\max}}\circ(\Hcpt{q_{\max}}\to\Heq{q_{\max}})=0$ is asserted
  (\ref{open:ell2}, the spelled-out clause; rider (iii)). \emph{What ``reported''
  means, per complex} (v1.5.0; \file{docs/implementation/slot3/WP-5\_report.md}
  finding~1, the reading implemented): a term $H^{q}(F)$ of one of the three complexes
  $F$ --- bulk, boundary, cone --- is \textbf{reported} iff $q<0$ (vacuously: the term
  is $0$), or $q\le q_{\max}(F)$, or $F$ is not truncated; an untruncated complex of
  finite length \emph{does} report its vanishing above its top degree, the box of
  \ref{open:ell2} forbidding that only where \code{truncated = true}. The three
  $q_{\max}$ are the per-complex fields of \S\ref{sec:recordschema} item~5, and the
  sentence before this one is exactly this rule in the case
  $q_{\max}(\text{bulk})=q_{\max}(\text{boundary})=q_{\max}$ that \ref{open:ell2}'s
  spelled-out clause states. On the shipped models nothing changes: all three complexes
  are finite and untruncated, so every term is reported and every composite and every
  spot is asserted. At every spot all three of
  whose terms are reported, $\ker=\im$: over a field by the rank identity, over
  $\Z/\ell^{e}$ or $\Z/m$ by a solve exhibiting each kernel generator as an image (a
  non-field elimination, [OD~OQ18]; for a ring over which no such solve is built the
  non-field case is \textbf{recorded} as unchecked, never reported as passed). A failure
  at an asserted spot \textbf{aborts}. \emph{The non-field case, as built} (v1.9.0; through
  v1.8.1 the parenthesis read ``until it exists the non-field case is recorded as
  unchecked'', stale since the elimination exists): over $\cring=\Z/\ell^{e}$, $\ell\ge5$,
  $\ell^{e}<2^{64}$, the non-field elimination is \file{src/ModArithmetic/ENGINE\_API.md}
  \S13.6 (hecke-extensions slot~8), and the check is \textbf{made}: at every spot, every
  composite of consecutive maps is shown zero by a solve into the target's relations, and
  every Howell generator of the kernel of the outgoing map is exhibited as the image of the
  incoming map plus a relation, the witnesses stored and re-derived
  (\code{ConeModuleCertificate} item~14, hecke-extensions slot-9 WP-4; measured by
  \code{HeckeReductionValidation}, arm (4b-C), at five of the seven points of
  \Espec{sec:ring}~(e) and at the two word-regime cells, the point $(60,5)$ at both levels
  $k\in\{2,3\}$ running in \code{HeckeReductionCoverageSlowValidation} since hecke-extensions
  slot-9 WP-4t (v1.9.1), which also covers the extension stage's case list at $\ell\ge5$). Beside it the length identity is asserted at every term --- the length of
  the term is the sum of the lengths of the two images --- and with the composite-zero check
  it implies exactness on its own, a submodule of a finite module of the same length being
  the whole; so the witness list is a second, independent discharge, not the only one. At
  $\ell\in\{2,3\}$ and for $\Z/m$ no cone over the ring is built (\Espec{sec:ring}~(b1)). $\im(\coneproj{}_{\ast}^{\,q})=\ker\rres^{q}$ is checked
  \textbf{explicitly} --- both inclusions, by solves, and not only through the rank
  identity --- at \textbf{every} bulk spot of \eqref{eq:coneles} all three of whose terms
  are reported, and not at $q=1$ alone (v1.6.0, read degree-generically as implemented;
  \file{docs/implementation/slot3/WP-5\_report.md} finding~5: through v1.5.2 the clause
  named the single instance $\im(\Hcpt{1}\to\Heq{1})=\ker\rres^{1}$, a degree-$1$ literal
  in a degree-generic obligation. On the model complexes that is $q\in\{0,1,2\}$; on a
  truncated free model every $q\le q_{\max}$ whose spot is asserted. A ring that is not a
  field, and over which no such solve is built, records the check unchecked, as the rank
  identity above does, and never reports it as passed; over $\Z/\ell^{e}$, $\ell\ge5$, both
  inclusions are made by solves since v1.9.0, as above).
  At $H^{q_{\max}}(\partial)$ on a truncated model the rank of $\rres^{q_{\max}}$ is
  \textbf{recorded}, and with it the consequence of exactness that needs no basis of
  the unreported term, $\coker\rres^{q_{\max}}\hookrightarrow\Hcpt{q_{\max}+1}$
  (Remark~\ref{rmk:ell2window}(2)); at $\ell=2$, where $H^{q}(\partial;\VR)$ is known
  in closed form in every degree (Lemma~\ref{lem:ell2unbounded}), that cokernel's
  dimension is a lower bound for $\dim\Hcpt{q_{\max}+1}$ and is recorded as such, never
  as a reported group. \ref{open:ell2} narrows the set of asserted spots on a truncated
  model and nothing else.
\item[(v)] \textbf{Bookkeeping, where the window closes.} $\Hcpt{0}=0$ is asserted
  unconditionally (Lemma~\ref{lem:conerelative}(4)); $\Heq{2}=0$ itself is asserted for
  $6\in\cring^{\times}$ (Claim~\ref{cl:amalgam}) and \textbf{recorded} otherwise. Of the
  two identities of \eqref{eq:conedims}, over a field and wherever the terms they name
  are reported: \textbf{the first is asserted without consulting $\Heq{2}$ at all}
  (v1.6.0) --- Lemma~\ref{lem:conerelative}(5) proves it from the exactness of
  \eqref{eq:coneles} and $\Hcpt{0}=0$ alone, so an implementation must assert it on
  \emph{every} computed cell and may not make the assertion conditional on a stored
  $\Heq{2}=0$; \textbf{the second is asserted where $\Heq{2}=0$ holds in the stored data}
  and \textbf{recorded} otherwise. A failure of an asserted item \textbf{aborts} --- the
  stored object's own certificate for \textbf{both} identities, each under its own
  minimal guard, \emph{and} the validation arm that carries the first, in the sense of
  clause (viii) (v1.7.0, restated: through v1.6.2 this sentence assigned the second
  identity to the certificate and the unconditional first one to ``the validation arm
  that carries it'' alone, which was the division of labour the shipped certificate then
  had; the certificate now carries the first identity too, so the clause names both and
  neither is the other's substitute).

  \smallskip
  \emph{What the strengthening does and does not say} (v1.6.0;
  \file{docs/implementation/slot3/WP-5\_report.md} finding~3; through v1.5.2 both
  identities were gated on $\Heq{2}=0$). It is a statement about the \emph{dependence},
  not about a known counterexample: on the shipped model complexes $\Heq{2}=0$ holds on
  every computed cell --- including all $18$ at $\ell\in\{2,3\}$, where
  Claim~\ref{cl:amalgam} does not apply and it is therefore \emph{recorded} rather than
  asserted (\code{HeckeConeValidation}: $74$ cells, $\ell\in\{2,3\}$ cells with
  $\Heq{2}\neq0$: none) --- so no cell currently separates the two readings. What the
  strengthening buys is that a cell which did would still be covered, and that a conforming
  implementation may not silently skip the count where its own $\Heq{2}$ is non-zero or
  unreported.

  \smallskip
  \emph{Carried by the certificate since v1.7.0} (hecke-extensions slot-5 WP-4, phase
  OI-2, owner ruling \code{[OD]} slot-4 close-out 2;
  \file{docs/implementation/slot5/WP-4\_report.md} \S1, \S3). \code{ConeCertificate}'s
  item C11 is split into two guards under the one item code: the first identity is
  evaluated without consulting the stored $\Heq{2}=0$, and without the two further
  reported-term conjuncts that belong to the second identity alone; the second keeps
  them. The split is not vacuous --- a \emph{constructed} cell in which $\Heq{2}$ is
  unreported (a bulk complex declared truncated at $q_{\max}=1$, so that the flag is
  structurally false with no stored dimension corrupted) and $\rres^{1}$ is perturbed
  entrywise makes the first identity fire and leaves the second silent, where the shared
  guard would have skipped both. Value-preserving on the shipped grid, as the paragraph
  above predicts: $\Acone{q}$ is bit-identical and both recorded collector digests are
  unchanged.
\item[(vi)] \textbf{Consistency assertions --- theorems re-checked.} For
  $6\in\cring^{\times}$: $[\Acone{q}(T_p),\Acone{q}(T_{p'})]=0$ and the $\nu=2$
  recursion of Theorem~\ref{thm:relations}(1) on $\Hcpt{q}$, with the \emph{full}
  $T_{p^{2}}$ and the coefficient $p^{\,n+1}$ (entries~\ref{err:e6}, \ref{err:e14}),
  in every computed degree; a failure \textbf{aborts}. For $6\notin\cring^{\times}$
  the same quantities are \textbf{recorded}, not asserted
  (Remark~\ref{rmk:conetorsion}). \emph{Where the full operator comes from} (v1.6.0;
  \file{docs/implementation/slot3/WP-6\_report.md} finding~4): the cone pipeline run at
  $g=\mathrm{diag}(1,p^{\nu})$ emits the \emph{primitive} operator
  $T_{\Deprim_{p^{\nu}}}$, exactly as the bulk pipeline does; the full $T_{p^{\nu}}$ this
  clause requires is assembled from the primitive ones by \eqref{eq:fullassembly}
  (Proposition~\ref{prop:fullassembly}), at $\nu=2$ by adding the central $p^{\,n}I$
  (\S\ref{sec:primfull}, entry~\ref{err:e1}). An implementation that checks the recursion
  against the primitive operator checks a false statement. \emph{Over $\Z/\ell^{e}$, as
  built} (v1.9.0, status; hecke-extensions slot-9 WP-4 and its review repair): the commutator
  $[\Acone{q}(T_p),\Acone{q}(T_{p'})]=0$, and likewise $[\Abdry{q}(T_p),\Abdry{q}(T_{p'})]=0$,
  is asserted at every $\ell\ge5$ cell over the ring for every pair of computed primes
  (\code{HeckeReductionValidation}, arm (4b-H)); \textbf{the $\nu=2$ recursion is not computed
  over the ring in this build}, so this clause is \textbf{partial} there
  (\S\ref{sec:notimplemented}).
\item[(vii)] \textbf{Cross-routes (owner decision OP3).} The strict-pair route and
  route 3 of Proposition~\ref{prop:coneroutes}(1),(2) each produce $\Acone{q}$ in the
  stored basis equal to the normative one, entrywise, in every computed degree;
  route 3's agreement \textbf{aborts} on failure for every $\ell$; the strict-pair
  route's agreement \textbf{aborts} for $6\in\cring^{\times}$ and is \textbf{recorded}
  otherwise. Once Claim~\ref{cl:pd} and the cochain-level pairing exist, the
  adjointness of Proposition~\ref{prop:coneroutes}(3) is asserted for $\ell\ge5$ and is
  the acceptance gate for $X_3$; until then it is \textbf{recorded} as unavailable.
  \emph{Over $\Z/\ell^{e}$, as built} (v1.9.0--v1.9.1, status; owner ruling 3~(iii)~(A) of the
  slot-9 decisions, 2026-09-24): only the normative route, route~1 of
  Proposition~\ref{prop:coneroutes}, was built over the ring. The strict-pair route, route~3,
  and the relative sub-complex constructor that route~3 needs refused a composite ring by name,
  each with a validation arm; so the agreement this clause asserts was an $\Fl$ statement in
  that build, and its composite half was open (stale from hecke-extensions slot-10 WP-3's merge,
  \texttt{c9684e9}, until v1.10.0).
  \emph{Over $\Z/\ell^{e}$, since v1.10.0} (status; owner ruling~3~(A) of the slot-10 decisions,
  2026-09-25; hecke-extensions slot-10 WP-3): the strict-pair route, route~3 and the relative
  sub-complex constructor are built over the ring at $\ell\ge5$ (the constructor on the module
  route of the bulk, refusing at $\ell\in\{2,3\}$ by name), and \textbf{the clause's first two
  sentences are discharged over the ring at $\ell^{e}<2^{32}$ --- at every point of
  \Espec{sec:ring}~(e)}: both cross-routes give $\Acone{q}$ entrywise equal to route~1's in its
  cyclic-summand basis, in every computed degree, and both agreements abort on failure (only
  $6\in\cring^{\times}$ reaches them). Route~3's transport uses no inverse matrix: the
  correspondence $M\colon[z]\mapsto[(z,0)]$ and the explicit
  $N\colon[(c,c_{\partial})]\mapsto[c+d\,\mathrm{ext}_0c_{\partial}]$, $\mathrm{ext}_0$ an
  equivariant section of $\rres$, are asserted mutually inverse. The cross-routes follow the
  companion's \texttt{uint32\_t} regime rule: above $2^{32}$ they refuse by the ring context's named
  refusal, and route~1 alone runs there. \textbf{The third sentence over the ring is recorded as
  unavailable}, by its own rule: the cochain-level pairing and the contragredient refuse a
  composite ring (\file{Hecke/HeckePairing}, \file{Hecke/HeckeContragredient}), so the adjointness
  of Proposition~\ref{prop:coneroutes}(3) is not evaluated there (the slot's independent review,
  \file{docs/implementation/slot10/WP-14\_review.md}, finding M-1). Measured by
  \code{HeckeConeCrossRouteRingValidation} ($386$ checks) and the changed rows of
  \code{HeckeReductionValidation}'s arms (X) and (X4b) (\S\ref{sec:notimplemented}).
\item[(viii)] \textbf{Two arms the validation suite must contain.} (a) \emph{Independence}:
  two $\etapair$-tables built with different vertex-chain policies in
  Lemma~\ref{lem:etaexists}(V), and two $\taub$-tables, must give \emph{bit-identical}
  $\Acone{q}$, bit-identical cone bases and $\Dcone$ (which involve only $d$ and
  $\rres$), while $\hder$ and the table hashes may move; a difference in $\Acone{q}$
  \textbf{aborts} (over $\Z$ for the $\etapair$/$\taub$ arm, by
  Proposition~\ref{prop:coneindep}(1),(2)). (b) \emph{Negative control}: a tampered
  $\hder'=\hder+\xi$ with $\xi$ built as in the proof of
  Proposition~\ref{prop:coneambiguity}(2) must \emph{pass} clauses (i)--(vi) and
  \emph{fail} Obligation~\ref{obl:eta}(v), and must change $\Acone{1}$ at a cell where
  $\Hom(\ker\rres^{1},\coker\rres^{0})\neq0$ (Remark~\ref{rmk:coneambiguity691}); the
  arm's failure to behave so \textbf{aborts} the validation. Arm (b) is the executable
  form of the statement that clauses (i)--(vi) certify the sequence and not the
  module.
\end{enumerate}
\end{obligation}

\begin{remark}[Restriction equivariance at $\ell\in\{2,3\}$ on the model complexes: the measurement that preceded the proof]\label{rmk:conerestrictsmallell}
\COMP{} \textbf{Since v1.7.0 this remark is load-bearing for no clause of \S\ref{sec:cone}}: clause (iii)
above asserts $\rres^{q}\Amat{q}=\Abdry{q}\rres^{q}$ at \emph{every} $\ell$, on the
strength of clause (i), and what follows is the record of the outcome \emph{measured
before} that implication was written into the document --- kept because \COMP{} never
upgrades and never disappears (\S\ref{sec:status}), and because it is still the record
that discharges Obligation~\ref{obl:taubdry}(vii) at $\ell\in\{2,3\}$, where the boundary
chapter has no $\hder$ and the proof of clause (iii) is unavailable (clause (iii)'s
asymmetry paragraph). Read as written, for that clause:
Obligation~\ref{obl:taubdry}(vii) --- $\rres^{q}\Tdc{g}=T^{\partial}_{p}\rres^{q}$ --- is
asserted at $6\in\cring^{\times}$ and is
\emph{recorded} at $\ell\in\{2,3\}$, where the model complexes compute the orbifold model
cohomology and its boundary analogue rather than group cohomology
(Remark~\ref{rmk:conetorsion}) and no statement of this document asserts it. Recorded
outcome (\code{HeckeBoundaryOperatorValidation}, tier \code{standard}, run 2026-09-16,
hecke-extensions slot 1 WP-4b, finding 15): on every $(\text{cell},q)$ pair of its grid at
$\ell\in\{2,3\}$ --- $(p,\nu)\in\{(2,1),(3,1),(5,1),(7,1),(2,2),(3,2)\}$, even $n\le16$,
$q\in\{0,1\}$ ($q=2$ vacuous), the bad cases $p=\ell$ and $\nu=2$ included --- the identity
holds entrywise over $\Fl$ in the stored bases: $216/216$ pairs. The observation is
consistent with what a pair homotopy $\etapair$ over $\Z$ (Definition~\ref{def:pairhomotopy},
Lemma~\ref{lem:etaexists}) gives once $\Phig$ exists --- the cohomology-level identity is a
consequence of Theorem~\ref{thm:conewelldef}(a) under its hypotheses and is not proved
without them --- and it is offered to the phase that builds this chapter (E5) as a recorded
fact and as an outcome its $\ell\in\{2,3\}$ arm may expect, \textbf{not} as a claim: a
failure at another $(n,p,\nu)$ or on another triangulation would contradict nothing in
this document. Rule~\ref{A6}: the record licenses nothing, and no code path may branch on
it. \COMP{} never upgrades to \PROVED{} (\S\ref{sec:status}). \emph{Re-recorded through
the cone} (v1.5.2; \code{HeckeConeActionValidation}, hecke-extensions slot-3 WP-6, run
2026-09-17): clause (iii)'s $\rres$-half holds on $424/424$ $(\text{cell},q)$ pairs of that
grid ($\ell\in\{2,3,5,7,691,134217757,2^{31}-1\}$, even $n\le16$, $74$ of its $212$
operators at $\ell\in\{2,3\}$). \emph{Superseded as the clause's authority} (v1.7.0):
the same $424/424$ are now \textbf{asserted} by clause (iii) at every $\ell$ --- not on
the strength of this record, which licenses nothing, but on the strength of clause (i),
which proves the identity with no hypothesis on $\ell$ for any implementation carrying
$\hder$. The sentence above about Theorem~\ref{thm:conewelldef}(a)'s hypotheses
remains correct for the boundary datum alone, which is the case
Obligation~\ref{obl:taubdry}(vii) still covers, and is no longer the last word for the
pair. (Open point (8) of the
v1.3.3 changelog; carried non-normatively in v1.4.0.)
\end{remark}

\paragraph{Code map.}
\textbf{Implemented, except clause (vii), since v1.5.2} (hecke-extensions slot 3, packages
WP-5 and WP-6, 2026-09-17; \file{docs/implementation/slot3/WP-5\_report.md},
\file{WP-6\_report.md}). \emph{The cone complex} (WP-5, phase E5-2):
\file{EquivariantMappingCone} (\code{EquivariantMappingCone}, generic, no Hecke include) ---
$\Cone{\rres}$ of Definition~\ref{def:pair} / \eqref{eq:conediff} on the flat layout of
Algorithm~\ref{alg:assembly}, its $Z/B/H$ bases and projector through the block model with
the pipeline's basis rules, the three map families of \eqref{eq:coneles} in stored bases per
degree, the per-complex degree fields of \S\ref{sec:recordschema} item~5 derived, and
\code{ConeCertificate} (Obligation~\ref{obl:coneexact}(iv),(v) and $\Dcone^{2}=0$, with a
re-deriving \code{verify()}) --- built on \file{Hecke/SL2ZBoundaryPair}, the restriction
$\rres^{\bullet}$ lifted pipeline-side and value-preservingly out of
\file{HeckeExt/HeckeBoundary}; \code{HeckeConeValidation} (tier \code{standard}, $74$ cells:
every spot exact, $\Hcpt{0}=0$ at every $\ell$ and $n$, $\dim\Hcpt{1}=3$ at $(10,691)$ and
$1+2\dim S_{n+2}$ with $\Hcpt{2}=0$ at the ten Kummer cells). \emph{The Hecke action}
(WP-6, phase E5-3): \file{Hecke/HeckeConeAction} (\code{HeckeConeAction}) --- $\hder$ by
the coset transport \eqref{eq:derivedh} of the stored $\etapair$-table in two bracketings
asserted bit-identical, never chosen; $\Phig$ of \eqref{eq:conemap}; the residual of
Lemma~\ref{lem:conesign}; $\Acone{q}$ by Algorithm~\ref{alg:assembly} on $\Cone{\rres}$;
\code{ConeActionCertificate} (Obligation~\ref{obl:eta}(v),(vi) and
Obligation~\ref{obl:coneexact}(i)--(iii)) and \code{ConeRelationsCertificate}
(Obligation~\ref{obl:coneexact}(vi), the \emph{full} $T_{p^{2}}$ of
\S\ref{sec:primfull}); the value-preserving vertex policies \code{EtaVertexPolicy}
(\file{Hecke/HeckePairHomotopy}) and \code{TauBoundaryVertexPolicy}
(\file{Hecke/HeckeBoundaryDatum}, default the canonical Algorithm~\ref{alg:taubdry}, the
second an admissible non-tie choice of Lemma~\ref{lem:taubdryexists}(1)) with
\code{table\_fnv64()} on all three tables, which is what Obligation~\ref{obl:coneexact}(viii)(a)
runs on; a standalone record \code{hecke\_cone\_action/v1};
\code{HeckeConeActionValidation} (tier \code{standard}, $63$ cells, $212$ operators).
Measured there: the identities of clause (i) hold on the stored $\Ceq$ basis in every
degree and on the full flat cochain space only in $29$ of $424$ degrees, which is what the
coset transport's use of $\Ga$-equivariance predicts and is recorded, not asserted (a
wording point for the next $X$). The golden collectors ($\dim\Hcpt{q}$, the tier-A
invariants of $\Acone{q}$, $\hder$, $\Thb_g$, the cone bases, the table hashes) are
\emph{defined and never emitted} --- called from the two validations only; their emission
is phase E2's. \textbf{Since v1.7.3 clause (vii) is implemented} ---
\code{HeckeConeStrictPairRoute} (slot-6 WP-4), \code{HeckeConeRelativeRoute} (WP-5), both in
\file{Hecke/HeckeConeAction}, and the PD gate of \file{Hecke/HeckePairing} (WP-3); the rest of
this sentence is the v1.7.2 status. \textbf{Not implemented} (v1.7.2): clause (vii) of
Obligation~\ref{obl:coneexact} --- the strict-pair route and route~3, phase E5's cross-route
arms (slot~6 of the programme) --- and PD adjointness (recorded unavailable until
Claim~\ref{cl:pd}); the cone's integration into \code{HeckeReport} and the campaign record
(E2). The paragraph below is the v1.4.2 status, kept as the record of the slot-2 half.

\textbf{Implemented in part since v1.4.2} (hecke-extensions slot 2, package WP-8,
2026-09-16; \file{docs/implementation/slot2/WP-8\_report.md}): the $\etapair$-table exists
--- \file{Hecke/HeckePairHomotopy} (\code{hecke::HeckePairHomotopyT<Elt>}), built by
Lemma~\ref{lem:etaexists} for every $\ell$ ((V) by the pairwise path routing of
Algorithm~\ref{alg:tau} step~2, (E) by the unique fill through the one integer solver
\file{Hecke/HeckeIntegerSolve}, moved verbatim out of \file{Hecke/HeckeTauTable.cpp}),
certified in its constructor against Obligation~\ref{obl:eta}(i)--(iv) with no mod-$m$
path (a tier failure aborts; owner decision OP2), \code{eta\_zero} recorded ((vi)),
\code{table\_fnv64()} and a JSON record \code{hecke\_pair\_homotopy/v1} carrying the
per-complex fields of \S\ref{sec:recordschema} item~5 in the boundary degrees $[0,1]$;
\code{HeckePairHomotopyValidation} (tier \code{standard}) re-derives (i)--(iv)
independently on $152$ tables, rejects $27$ tampered chains, exhibits
Lemma~\ref{lem:etaunique} on two vertex-chain policies, and finds \code{eta\_zero} false
on every table (Remark~\ref{rmk:conestrict}(1)). \textbf{Not implemented} at v1.4.2: the cone
$\Cone{\rres}$, $\Phig$, the derived $\hder$, $\Acone{q}$, clause (v) and
Obligation~\ref{obl:coneexact} --- slot 3 of the extension programme, which built them
(above). Until v1.4.1 this
paragraph recorded that nothing in the chapter was implemented; the census it rested on
is kept below. A census of the tree on 2026-09-16 ---
\code{grep} over \file{src/} for
\code{eta}, \code{PairHomotopy}, \code{ConeAction}, \code{HeckeCone},
\code{EquivariantMappingCone}, \code{HeckePairDatum}, \code{TauBoundary} and
\code{Phi\_g} --- found no $\etapair$-table, no equivariant cone
complex, no $\Phig$, no $\Acone{q}$ and no \code{HeckeConeValidation}; none of the $80$
registered validations (v1.4.2 count) computed $\Hcpt{\bullet}$, and the validation of the
$\etapair$-table was \code{HeckePairHomotopyValidation}, the cone's to be a second name
(since v1.5.2 two: \code{HeckeConeValidation} and \code{HeckeConeActionValidation}). (The $\taub$-table the
same census also missed exists since v1.3.3 --- \file{Hecke/HeckeBoundaryDatum},
\S\ref{sec:boundarydatum}'s code map --- and is this chapter's second input; the pair
homotopy $\etapair$ that joins it to the bulk $\tau$ is built since v1.4.2.) What
exists is (a) the
\emph{non-equivariant} algebraic cone of a \code{CochainComplexMorphism},
\code{CochainComplexMorphism::ConeComplex} (\file{SimplicialCohomology.h}:227,
\file{SimplicialCohomology.cpp}:832), the implementation of \file{monodromy\_v5.tex}
\S8 on \code{SimplicialComplex} without a local system --- the model for the
equivariant cone of \eqref{eq:conediff}, not an instance of it; (b) the restriction
morphism $\rres^{q}$ of the circle model as an \code{EquivariantCochainComplexMorphismOnK}
(\S\ref{sec:boundarydatum}, code map); and (c) a comment in
\file{EquivariantCohomologyOnK.h}:115 that names ``the Cone class'' as the intended
user of a protected constructor --- a class that does not exist. (Since v1.5.2 a cone class exists,
\code{EquivariantMappingCone}, and it does not use that constructor; the comment of (c)
still stands in the header.) The
\emph{specification} for the phase that builds this (E5 of the extension programme,
gated by amendment C-A3): an $\etapair$-table builder reusing
\code{hecke::HeckeTauTable}'s step-2/3 machinery, (T1) helper and constructor-time
certification (invariant I8) with $\Xbdry$ as source and $\Xbulk$ as target; a
\code{HeckePairDatum} $\{\tau,\taub,\etapair;\ \text{solve regimes};\ \code{integral}\}$
as tables of \code{ChainTerm}s over $\Z$, degree-generic per \ref{open:ell2}; an
\code{EquivariantMappingCone} realising \eqref{eq:conediff} on
\code{EquivariantCohomologyOnK} objects; the derived $\hder$ and $\Phig$ of
\eqref{eq:derivedh}--\eqref{eq:conemap} assembled through the block model of
\S\ref{sec:systematiccomputation}; the four certificates of the table above with
\code{verify()} re-deriving and never re-running a search; and, per owner decision
OP7, golden tiers --- $\dim\Hcpt{q}$ and the char-poly/rank/trace of $\Acone{q}$ for
$\ell\ge5$ in tier~A, and $\Thb_g$, $\hder$, the cone bases and the three table hashes in
tier~B, because they depend on the choices Proposition~\ref{prop:coneindep} shows to be
immaterial.

\paragraph{Formalization note.}
This is the chapter a formalisation should attack first, because it is the one whose
statements have no external corroboration: every one of
Lemmas~\ref{lem:conerelative}--\ref{lem:derivedhomotopy}, Theorem~\ref{thm:conewelldef}(a),(b),(d),
Propositions~\ref{prop:coneindep}(1),(2)-boundary-half, \ref{prop:strictpair},
\ref{prop:coneambiguity} and \ref{prop:coneintegral}(1)--(4) is proved here from
\eqref{eq:contractibleuse}, \ref{B1}--\ref{B3}, \ref{H4}(iv), Lemma~\ref{lem:cocycle} and
Proposition~\ref{prop:symconv}(2) by finite manipulations, and none of them cites
anything a proof assistant would have to import beyond the long exact sequence of a
short exact sequence of complexes. The recommended order: Lemma~\ref{lem:cosettransport}
first --- it is four one-line identities and its equivariance clause is the
equivariance half of Theorem~\ref{thm:welldef} with a parameter --- then
Lemma~\ref{lem:conesign} (pure algebra on a direct sum), Lemma~\ref{lem:derivedhomotopy}
(a substitution), Lemma~\ref{lem:etaunique} (two uses of $Z_2=0$, one of $Z_1=B_1$), and
Proposition~\ref{prop:coneindep}(1), which is the theorem that matters and is one
displayed computation. The constructive content is Lemmas~\ref{lem:etaexists} and
Proposition~\ref{prop:strictpair}: both are terminating procedures whose correctness
reduces to \eqref{eq:contractibleuse} on a $2$-dimensional complex and a stabiliser
condition, so they need simplicial chains and the homology of a contractible
$2$-complex, not homological algebra; the hypothesis a formalisation must not lose is
\ref{B3} in Lemma~\ref{lem:etaexists} (three uses: the (T1)-extension on vertices, on
edges, and in the strict-pair construction), and \ref{H4}(ii) everywhere. Two things
are genuinely harder. Proposition~\ref{prop:coneambiguity}(2) chooses complements and is
therefore classical linear algebra over a field, correctly so restricted
(entry~\ref{err:e3}). Theorem~\ref{thm:conewelldef}(c) and Remark~\ref{rmk:conesteinberg}
depend on the comparison theorem through Lemma~\ref{lem:tauexists}(2) and, for the
Steinberg form, on derived-category bookkeeping ($C_{\bullet}(\Xbulk,\Xbdry)\simeq
\mathrm{St}[1]$, $\Ext$ of a bounded complex of relative projectives) that
\code{mathlib4} has in pieces; they should be attempted after Part~III's comparison
theorem, and a formalisation may take (iii) of that remark as an axiom about names
rather than a theorem, since nothing in the document depends on it.
Obligations~\ref{obl:eta} and \ref{obl:coneexact} are decidable predicates on emitted
finite data in every clause except \ref{obl:coneexact}(iv) over a non-field ring, which
is decidable through a Howell or Smith form once one is available.

%----------------------------------------------------------------------
\section{The six groups and the long exact sequence}\label{sec:six}
%----------------------------------------------------------------------
% SOURCE:  hecke_extensions.tex:186-338 (RH; moved out of [E] by the scope cut,
%          Appendix~A entry err:e18).
% NOTATION: crosswalk rows q (degree, never the Hecke prime), p (Hecke prime),
%          k = n+2, r -> \rres, R -> \cring, N -> \mathrm{Nm} (C18), Delta -> \cf.
% ERRATA:  E14 (cl:pd: socle, not Loewy), E18 (the move itself),
%          E21 (cl:chi's Dickson-locus sentence), and the two hypothesis
%          repairs recorded in rmk:sixscope.

\paragraph{What this chapter fixes.}
The six cohomology groups this project computes or predicts, the exact sequence that
relates them, the characters by which the Hecke algebra acts on the two boundary terms,
and the duality that pins the two compactly supported ones. Everything here is a
\emph{fact about the objects} constructed in Parts~I--IV; nothing here is a programme.
That is why these nine labels live in this document and not in the companion
specification (\S\ref{sec:scope}; Appendix~\ref{app:errata}, entry~\ref{err:e18}), and
it is why Obligation~\ref{obl:oracle} may read them.

\paragraph{Standing hypotheses.}
\ref{H1}--\ref{H4} and \ref{B1}--\ref{B3} are in force, and per rule~\ref{A1} every
statement repeats the ones it uses. \textbf{Three hypotheses recur and are never
suppressed}: $6\in\cring^{\times}$ (equivalently $\ell\ge5$ for $\cring$ of
characteristic $\ell$), which is \ref{H3} and is what makes the model complexes compute
group cohomology (Remark~\ref{rmk:orbifold}); $2\in\cring^{\times}$, which is what makes
the circle model compute the boundary cohomology (Proposition~\ref{prop:circlemodel}(3));
and $n$ even, without which every group below vanishes for $2\in\cring^{\times}$
(Remark~\ref{rmk:negI}, Remark~\ref{rmk:bdryoddn}).

\subsection{The six groups}

\begin{contractbox}
Let $\cring$ be a commutative ring, $n\ge0$ and $\VR=\Sym^{n}(\Z^{2})\otimes\cring$
(at $\cring=\Fl$ also written $\Mn$). The \textbf{six groups} are
\begin{equation}\label{eq:sixgroups}
\begin{aligned}
  X_1&=H^{0}(\Ga;\VR), & X_2&=H^{0}(\partial;\VR), & X_3&=\Hcpt{1},\\
  X_4&=H^{1}(\Ga;\VR), & X_5&=H^{1}(\partial;\VR), & X_6&=\Hcpt{2},
\end{aligned}
\end{equation}
with $H^{\bullet}(\partial;\VR):=H^{\bullet}(\Gbdry;\VR)$
(Definition~\ref{def:bdrygroup}) and $\Hcpt{q}:=H^{q-1}(\Cone{\rres})$
(Definition~\ref{def:pair}, Lemma~\ref{lem:conerelative}(3)). Write
$\Htot:=\bigoplus_{t=1}^{6}X_t$ and, over a field, $d:=\dim\Htot$.
\end{contractbox}

\noindent
$X_1$ and $X_4$ are the groups the shipped pipeline emits, as $\Heq{0}$ and $\Heq{1}$ of
$(K,\Rel;L)$; the identification with group cohomology is Remark~\ref{rmk:orbifold} and
holds for $6\in\cring^{\times}$. $X_2$ and $X_5$ are computed today as the cohomology of
the circle model $\Kbdry$, which is the $2\in\cring^{\times}$ specialisation of
\eqref{eq:sixgroups} (Proposition~\ref{prop:circlemodel}); \textbf{no Hecke action on
them is implemented} --- the operator is Definition~\ref{def:bdryop} and the code map of
\S\ref{sec:ell2} says what is missing. $X_3$ and $X_6$ are not computable today at all:
the cone of Definition~\ref{def:pair} exists for the non-equivariant
\code{CochainComplex} only. Their Hecke action is \emph{defined} ---
Definition~\ref{def:coneaction}, Theorem~\ref{thm:conewelldef} --- and that definition is
what \S\ref{sec:cone} exists to supply.

\begin{remark}[What this chapter claims, and two hypotheses it does not inherit]\label{rmk:sixscope}
The nine statements below were written in \file{hecke\_extensions.tex}:186--338 under a
blanket ``$\ell\ge5$ throughout, $n$ even'' declared once at :157--160. Re-homed here
under rule~\ref{A1}, each carries its own hypotheses, and two of them are strictly weaker
than the blanket:
\begin{enumerate}
\item \textbf{$\Hcpt{0}=0$ is unconditional.} Lemma~\ref{lem:conerelative}(4) proves it
  with no hypothesis on $\ell$ and none on $n$: $\rres^{0}$ is the injection
  $\VR^{\Ga}\hookrightarrow\VR^{\Gbdry}$. Claim~\ref{cl:les} therefore does \emph{not}
  list it as a hypothesis, where [E] did.
\item \textbf{$n$ even is not needed for exactness.} It is needed for the six groups to
  be non-zero, not for the sequence; Claim~\ref{cl:les} asks only for the two vanishings
  it actually uses.
\end{enumerate}
\textbf{The third hypothesis is not removable.} $H^{2}(\Ga;\VR)=0$ genuinely requires an
invertibility assumption: Lemma~\ref{lem:h2torsion}(2) exhibits $\ell=2$, $n=0$ with
$H^{q}(\Ga;\F_2)\neq0$ for \emph{every} $q\ge0$. Claim~\ref{cl:les} keeps it.
\end{remark}

\subsection{The amalgam, and what it gives in each degree}

\begin{claim}[Amalgam presentation]\label{cl:amalgam}
Let $\cring$ be a commutative ring and $n\ge0$; put $U:=ST$, so that $\langle S\rangle$,
$\langle U\rangle$, $\langle -I\rangle$ have orders $4$, $6$, $2$.
\begin{enumerate}
\item \STD{\cite[Ch.~I, \S4]{serre_trees}: the amalgam decomposition of $\SL_2(\Z)$ and
  the correspondence between amalgams and group actions on trees. The pointer is given at
  chapter-and-section level and is \textbf{not} verified against the text
  (rule~\ref{A2}): no copy is in \file{docs/literature/}. \eqref{eq:treeses} is proved
  independently of this pointer (the markers below), so nothing in this document rests
  on it; the amalgam decomposition itself is used nowhere else (v1.22.0-pre.12, LR-Y)}
  $\Ga=\langle S\rangle\ast_{\langle-I\rangle}\langle U\rangle$, and the augmented
  simplicial chain complex of the associated Bass--Serre tree is a short exact sequence
  of $\Z[\Ga]$-modules
  \begin{equation}\label{eq:treeses}
    0\to\Z[\Ga/\langle-I\rangle]\to\Z[\Ga/\langle S\rangle]\oplus\Z[\Ga/\langle U\rangle]
      \to\Z\to0 .
  \end{equation}
  The vertices of $Y$ (the $\Ga$-triangulation of $\Hyp$ behind $\mathsf M_7$,
  Definition~\ref{def:model7}) lying in the $\Ga$-orbits of $\zeta^{2}$ and $i$,
  with the edges of $Y$ between them, form the tree ($\mathcal T$ in
  Lemma~\ref{lem:bulkperiodic}, $\mathcal Y$ in \S\ref{sec:levelN}): vertex set
  $\mathcal T_0=\Ga/\langle S\rangle\sqcup\Ga/\langle U\rangle$, edge set
  $\mathcal T_1=\Ga/\langle-I\rangle$, the edge $\gamma\langle-I\rangle$ joining
  $\gamma\langle S\rangle$ to $\gamma\langle U\rangle$.
  \PROVED{} \LEAN{thm:mp-treeses-exact} (exactness of the augmented chain complex of the
  tree, over every commutative ring, without Serre: the full subcomplex of $Y$ on the
  $\Ga$-orbits of $\zeta^{2}$ and $i$ is a retract of the acyclic $C_\bullet(Y)$ of
  $\mathsf M_7$) \LEAN{thm:mp-model7-tree-orbits} (the terms:
  $\mathcal T_1\cong\Ga/\langle-I\rangle$,
  $\mathcal T_0\cong\Ga/\langle U\rangle\sqcup\Ga/\langle S\rangle$ as $\Ga$-sets and as
  permutation modules over every ring,
  $k[\Ga/\langle U\rangle\sqcup\Ga/\langle S\rangle]=k[\Ga/\langle S\rangle]\oplus k[\Ga/\langle U\rangle]$
  up to \texttt{Finsupp.sumFinsuppLEquivProdFinsupp} and the order of the summands; the
  differential is the simplicial boundary, $\Ga$-equivariant)
  \LEAN{thm:mp-treeses-rep} (the sequence assembled as one short exact sequence in $\mathrm{Rep}\,k\,\Ga$ over every
  commutative ring $k$, $\gamma\langle-I\rangle\mapsto\gamma\langle S\rangle-\gamma\langle U\rangle$ and every coset $\mapsto1$,
  with middle term the permutation module $k[\Ga/\langle S\rangle\sqcup\Ga/\langle U\rangle]$, which is
  $k[\Ga/\langle S\rangle]\oplus k[\Ga/\langle U\rangle]$ up to \texttt{Finsupp.sumFinsuppLEquivProdFinsupp}; exactness transported
  from [GT] \texttt{thm:mp-treeses-exact} along the orbit--stabiliser bijections).
\item \PROVED{} \LEAN{thm:mp-tree-mv} (the long exact sequence of \eqref{eq:treeses} for every $\cring$, every
  $\cring[\Ga]$-module $V$ (in particular $\VR$) and every $q$, built on the periodic cochain complexes
  $V^2\to V^2\to\cdots$ ($(1-S,1-U)$, $(N_S,N_U)$ alternating) and $V\to V\to\cdots$ ($1-(-I)$, $1+(-I)$ alternating, shifted
  up by one); the middle map is induced by $\varphi^*$ of Remark~\ref{rmk:treeconeres} on cocycles; the maps into and out of
  $H^q(\Ga;V)$ go through the comparison $H^q(\Ga;V)\cong H^q(\Hom_\Ga(P_\bullet,V))$ built on $\mathrm{Ext}$ (mathlib's
  \texttt{groupCohomologyIso}; [GT] \texttt{thm:mp-tree-cochains}))
  \LEAN{thm:mp-mv-terms} (the terms: these periodic complexes compute $H^q(\langle S\rangle;V)\oplus H^q(\langle U\rangle;V)$
  and $H^q(\langle-I\rangle;V)$ as mathlib's \texttt{groupCohomology} of the subgroups, every $q$, through evaluation at the
  generators and a fixed homotopy equivalence with the bar resolution; the identification of the edge term is the identity on
  cocycles)
  \LEAN{thm:mp-mv-resres} (under those identifications the middle map is $\res_S-\res_U$, mathlib's restriction along
  $\langle-I\rangle\subset\langle S\rangle$ and $\langle-I\rangle\subset\langle U\rangle$, with the sign as displayed, every $q$;
  on the periodic cochains restriction is $x\mapsto x$ in even and $x\mapsto(1+S)x$, resp.\ $(1+U+U^2)x$, in odd degrees)
  \LEAN{thm:mp-mv-low} (the beginning $0\to V^\Ga\to V^S\oplus V^U\to V^{-I}$, $a\mapsto(a,a)$, $(x,y)\mapsto x-y$, is exact
  and $V^\Ga=V^S\cap V^U$, proved directly, not as the degree-$0$ case of the sequence; restriction on $H^0$ is the inclusion
  of invariants under mathlib's $H^0(G;V)\cong V^G$)
  \LEAN{thm:mp-mv-res} (through the comparison $H^q(\Ga;V)\cong H^q(\Hom_\Ga(P_\bullet,V))$ built on a fixed homotopy
  equivalence with the bar resolution (not the $\mathrm{Ext}$-built one of the first marker) and the identifications of
  [GT] \texttt{thm:mp-mv-terms}, the map
  $H^q(\Ga;V)\to H^q(\langle S\rangle;V)\oplus H^q(\langle U\rangle;V)$ of the sequence is $(\res_S,\res_U)$, mathlib's
  restriction, every $q$, $q=0$ included; on cochains the projections $((x,y),z)\mapsto x$, $y$)
  \LEAN{thm:mp-mv-les} (the whole sequence \eqref{eq:mvamalgam} in \texttt{groupCohomology} terms, under the same comparison
  and the identifications of [GT] \texttt{thm:mp-mv-terms}, exact at every term, $q=0$
  included ($H^0(\Ga;V)\to H^0(\langle S\rangle;V)\oplus H^0(\langle U\rangle;V)$ injective), with maps $(\res_S,\res_U)$,
  $\res_S-\res_U$ and a connecting map $H^q(\langle-I\rangle;V)\to H^{q+1}(\Ga;V)$ defined by transport of the inclusion of
  the edge part of $\Hom_\Ga(P_\bullet,V)$ and characterised on cocycles ($[z]\mapsto[((0,0),z)]$); the connecting map is
  not identified with that of mathlib's \texttt{groupCohomology} sequence of a coefficient sequence, and the agreement of
  the bar-built and the $\mathrm{Ext}$-built comparisons is not proved)
  \LEAN{thm:mp-mv-natural} (naturality in $V$: each of the three maps commutes with the maps induced by a morphism
  $V\to V'$ of $\cring[\Ga]$-modules and its restrictions to $\langle S\rangle$, $\langle U\rangle$, $\langle-I\rangle$, every
  $q$)
  \textbf{(Mayer--Vietoris, every $\cring$, every $n$, every degree.)}
  \eqref{eq:treeses} yields a long exact sequence, natural in $\VR$,
  \begin{equation}\label{eq:mvamalgam}
    \cdots\to H^{q}(\Ga;\VR)\to H^{q}(\langle S\rangle;\VR)\oplus
    H^{q}(\langle U\rangle;\VR)\xrightarrow{\ \res-\res\ }
    H^{q}(\langle-I\rangle;\VR)\to H^{q+1}(\Ga;\VR)\to\cdots
  \end{equation}
  beginning $0\to H^{0}(\Ga;\VR)\to\VR^{S}\oplus\VR^{U}\to\VR^{-I}$.
\item \PROVED{} \textbf{(Consequences for $6\in\cring^{\times}$.)} If
  $6\in\cring^{\times}$ then $H^{q}(\Ga;\VR)=0$ for every $q\ge2$;
  $H^{0}(\Ga;\VR)=\VR^{S}\cap\VR^{U}$; and, for $n$ even,
  \[
    H^{1}(\Ga;\VR)\;=\;\VR\big/\bigl((1+\rhon(S))\VR+(1+\rhon(U)+\rhon(U)^{2})\VR\bigr).
  \]
  If moreover $\cring$ is a field, then for $n$ even
  $\dim H^{0}(\Ga;\VR)-\dim H^{1}(\Ga;\VR)=\dim\VR^{S}+\dim\VR^{U}-\dim\VR$. (The letter
  $\chi$ is \emph{not} used for this Euler characteristic: it is the boundary character
  of Claim~\ref{cl:chi}, \S\ref{sec:notation-hecke}.)
\item \PROVED{} \textbf{(Torsion bound, every $\cring$, every $n$.)}
  $24\cdot H^{q}(\Ga;\VR)=0$ for every $q\ge2$; the sharp classical constant is $12$,
  $\Ga$ having a free subgroup of index $12$
  \STD{\cite[Ch.~I, \S4]{serre_trees} for free subgroups of a free product of finite
  groups, with the same rule~\ref{A2} caveat; the commutator subgroup of $\SL_2(\Z)$ is
  the classical witness, $\SL_2(\Z)^{\mathrm{ab}}=\Z/12$}. Either
  constant gives (3)'s vanishing, since both are invertible when $6$ is.
\end{enumerate}
\end{claim}

\begin{proof}
(2) For each of the three subgroups $\Ga_0\in\{\langle S\rangle,\langle U\rangle,
\langle-I\rangle\}$, $\Z[\Ga/\Ga_0]=\Z[\Ga]\otimes_{\Z[\Ga_0]}\Z$ is an induced module, so
$\Ext^{q}_{\Z[\Ga]}(\Z[\Ga/\Ga_0],\VR)\cong H^{q}(\Ga_0;\VR)$
\STD{\cite[Ch.~III.6]{brown}, Shapiro's lemma, as in the proof of
Proposition~\ref{prop:circlemodel}(1)}; applying
$\Ext^{\bullet}_{\Z[\Ga]}(-,\VR)$ to \eqref{eq:treeses} and using
$\Ext^{q}_{\Z[\Ga]}(\Z,\VR)=H^{q}(\Ga;\VR)$ gives \eqref{eq:mvamalgam}, with the
connecting maps of the $\Ext$ sequence, and naturality in $\VR$ is naturality of that
sequence. The initial segment is $\Hom_{\Z[\Ga]}(\Z[\Ga/\Ga_0],\VR)=\VR^{\Ga_0}$ together
with $H^{0}(\Ga;\VR)=\VR^{\Ga}$.

(3) For a finite subgroup $\Ga_0\le\Ga$ and $q\ge1$, $\abs{\Ga_0}\cdot
H^{q}(\Ga_0;\VR)=0$
\STD{\cite[Ch.~III.10.2]{brown} (Corollary~10.2, p.~84), from $\cor\circ\res=[\Ga_0:1]$ (\cite[Prop.~III.9.5(ii), p.~82]{brown}) on $H^{q}(\Ga_0;-)$ for
$q\ge1$, where $H^{q}(1;-)=0$}; with $4,6,2\in\cring^{\times}$ every term of
\eqref{eq:mvamalgam} in a positive degree vanishes, so $H^{q}(\Ga;\VR)=0$ for $q\ge2$
and $H^{1}(\Ga;\VR)=\coker(\VR^{S}\oplus\VR^{U}\to\VR^{-I})$. For $n$ even
$\rhon(-I)=\id$ (Proposition~\ref{prop:symconv}(2)), so $\VR^{-I}=\VR$; and
$\rhon(S)^{2}=\rhon(U)^{3}=\rhon(-I)=\id$, so the averaging idempotents of
$\langle S\rangle$ and $\langle U\rangle$ are $\tfrac12(1+\rhon(S))$ and
$\tfrac13(1+\rhon(U)+\rhon(U)^{2})$, whence $\VR^{S}=(1+\rhon(S))\VR$ and
$\VR^{U}=(1+\rhon(U)+\rhon(U)^{2})\VR$; the image of the difference map is their sum.
$H^{0}(\Ga;\VR)=\VR^{S}\cap\VR^{U}$ is exactness at the second term, $S$ and $U$
generating $\Ga$. The Euler characteristic is the alternating count on
$0\to H^{0}\to\VR^{S}\oplus\VR^{U}\to\VR\to H^{1}\to0$.

(4) In \eqref{eq:mvamalgam} at $q\ge2$, exactness gives
$0\to C\to H^{q}(\Ga;\VR)\to A\to0$ with $A$ a submodule of
$H^{q}(\langle S\rangle;\VR)\oplus H^{q}(\langle U\rangle;\VR)$ and $C$ a quotient of
$H^{q-1}(\langle-I\rangle;\VR)$. By the annihilator cited in (3), $\mathrm{lcm}(4,6)=12$
kills $A$ and $2$ kills $C$ (note $q-1\ge1$), so $24$ kills $H^{q}(\Ga;\VR)$. For the
sharp constant: if $F\le\Ga$ is free of index $12$ then $H^{q}(F;\VR)=0$ for $q\ge2$
($F$ has cohomological dimension $\le1$, its Cayley graph being a contractible free
$F$-complex of dimension $1$), and $\cor\circ\res=12$ on $H^{q}(\Ga;\VR)$
\STD{\cite[Ch.~III.9--10]{brown}}.
\end{proof}

\noindent
\emph{Code map.} The presentation of Claim~\ref{cl:amalgam}(3) is the independent
$H^{1}$ route of \file{Hecke/BarCocycleH1.cpp}: it presents
$\Ga=\langle S\rangle\ast_{\langle S^{2}\rangle}\langle ST\rangle$, represents a
$1$-cocycle by the pair $(\phi(S),\phi(ST))$ subject to the norm and amalgam-matching
conditions, enforces $\ell\ge5$ at construction --- which is exactly hypothesis (3) ---
and compares characteristic and minimal polynomials against $\Amat{1}$
(Proposition~\ref{prop:barformula}, \S\ref{sec:degreewise}).

\begin{remark}[The tree resolution over every ring]\label{rmk:treeconeres}
\PROVED{} \LEAN{thm:mp-model7-cone-res} (the free resolution $P_\bullet$, acyclic in every
degree) \LEAN{thm:mp-tree-mp-comparison} ($P_\bullet$ is a projective resolution;
$H^q(\Ga;V)\cong H^q(\Hom_\Ga(P_\bullet,V))\cong
H^q(C^\bullet_{\mathrm{un}}(\mathsf S(\mathsf M_7);V))$; the homotopy equivalence with the
MP resolution is the abstract one; the normalised $C^\bullet$ through
Theorem~\ref{thm:mpcomparison}(4)) \LEAN{thm:mp-tree-cochains}
($\Hom_\Ga(P_\bullet,V)\cong(V^2\to V^3\to V^3\to\cdots)$ by evaluation at the generators,
hence $H^q(\Ga;V)\cong H^q(V^2\to V^3\to\cdots)$, every commutative ring, every $V$, every
$q$, $-I$ included)
Let $P^{\langle S\rangle}$, $P^{\langle U\rangle}$, $P^{\langle-I\rangle}$ be the
$2$-periodic free resolutions of $\Z[\Ga/\langle S\rangle]$, $\Z[\Ga/\langle U\rangle]$,
$\Z[\Ga/\langle-I\rangle]$ induced from the cyclic groups
($\cdots\xrightarrow{N}\Z[\Ga]\xrightarrow{1-t}\Z[\Ga]\to\Z[\Ga/\langle t\rangle]$, right
multiplication). The boundary of \eqref{eq:treeses} is oriented as
$\mathfrak e\mapsto\mathfrak v_S-\mathfrak v_U$ (the edge $[\zeta^{2},i]$ runs from
$\zeta^{2}$ to $i$). This is the \textbf{negative} of the convention
$\partial(\gamma\mathfrak e)=\gamma\mathfrak v_U-\gamma\mathfrak v_S$ of the tree
$\mathcal Y$ of \S\ref{sec:levelN} (its notation paragraph), which moves no statement. The
boundary lifts to $\varphi\colon P^{\langle-I\rangle}\to P^{\langle S\rangle}\oplus
P^{\langle U\rangle}$, with $\varphi=(1,-1)$ in even and $(1+S,-(1+U+U^2))$ in odd degrees,
in the order $(\langle S\rangle,\langle U\rangle)$.
The mapping cone of $\varphi$, augmented to $\Z$, is a free resolution
$P_\bullet\colon\Z[\Ga]^2\leftarrow\Z[\Ga]^3\leftarrow\Z[\Ga]^3\leftarrow\cdots$ of $\Z$
over $\Z[\Ga]$. Its terms are $2$-periodic from degree $1$ and its differentials from
degree $2$: $P_{n+2}=P_n$ for $n\ge1$ and $\partial_{n+2}=\partial_n$ for $n\ge2$; dually
$\Hom_\Ga(P_\bullet,V)$ has $C^{n+2}=C^n$ and $d^{n+2}=d^n$ for every $n\ge1$. For every
commutative ring $\cring$ the same holds over $\cring[\Ga]$. So for every
$\cring[\Ga]$-module $V$ and every $q$,
$H^q(\Ga;V)\cong H^q(\Hom_\Ga(P_\bullet,V))\cong H^q(C^\bullet(\mathsf S(\mathsf M_7);V))$
(Theorem~\ref{thm:mpcomparison}). With $N_S=\sum_{j<4}S^j$, $N_U=\sum_{j<6}U^j$, and $S$,
$U$, $-I$ acting on $V$ through its structure map, $\Hom_\Ga(P_\bullet,V)$ is the cochain
complex $V^2\to V^3\to V^3\to\cdots$ with $d^0(x,y)=((1-S)x,(1-U)y,x-y)$ and
$d^{m+1}((x,y),z)=(v_{m+1}(x,y),\varphi^*_{m+1}(x,y)-e_mz)$, where $v_j=(1-S,1-U)$ and
$\varphi^*_j(x,y)=x-y$ for even $j$, $v_j=(N_S,N_U)$ and
$\varphi^*_j(x,y)=(1+S)x-(1+U+U^2)y$ for odd $j$, and $e_m=1\mp(-I)$ for $m$ even/odd.
$P_\bullet$ and $\cring[\mathfrak F'_\bullet]$ are homotopy equivalent, as any two
projective resolutions of $\cring$ are; the equivalence is not computed. When
$6\in\cring^\times$, Lemma~\ref{lem:treeres}'s length-$1$ resolution
$C_\bullet(\mathcal Y)\otimes\cring$ is a second projective resolution. $P_\bullet$ needs
no unit: its terms are free.
\end{remark}

\begin{lemma}[$H^{2}(\Ga;\VR)=0$ needs its hypothesis]\label{lem:h2torsion}
\begin{enumerate}
\item \PROVED{} If $6\in\cring^{\times}$ then $H^{q}(\Ga;\VR)=0$ for every $q\ge2$ ---
  Claim~\ref{cl:amalgam}(3). This is the hypothesis under which
  \S\ref{sec:degreewise} reports $\Amat{2}$ as the $0\times0$ matrix, and under which
  $\Heq{2}(K,\Rel;L)=0$.
\item \PROVED{} \textbf{It is not removable.} Let $\cring=\F_2$ and $n=0$, so
  $\VR=\F_2$ with trivial $\Ga$-action. Then $H^{q}(\Ga;\F_2)\neq0$ for \emph{every}
  $q\ge0$; in particular $H^{2}(\Ga;\F_2)\neq0$, and the hypothesis $H^{2}(\Ga;M)=0$ of
  Claim~\ref{cl:les} fails at $(\ell,n)=(2,0)$, a pair with $n$ even.
\end{enumerate}
\end{lemma}

\begin{proof}
(2) With trivial coefficients $\F_2$: $H^{q}(\langle-I\rangle;\F_2)=\F_2$ for every
$q\ge0$ (the periodic resolution of step (ii) in the proof of
Proposition~\ref{prop:bdrycohom}, with $2=0$ in $\F_2$, gives
$\F_2\xrightarrow{0}\F_2\xrightarrow{0}\cdots$); $H^{q}(\langle S\rangle;\F_2)=\F_2$ for
every $q\ge0$ by the analogous periodic resolution of a cyclic group of order $4$, whose
two maps are $g-1=0$ and the norm $4=0$; and
$H^{q}(\langle U\rangle;\F_2)=H^{q}(\langle-I\rangle;\F_2)=\F_2$, because
$\langle U\rangle\cong\Z/2\times\Z/3$ with $3\in\F_2^{\times}$, so the
Lyndon--Hochschild--Serre sequence of that product collapses to
$H^{q}(\Z/2;\F_2^{\Z/3})=H^{q}(\Z/2;\F_2)$
\STD{\cite[Ch.~VII.6]{brown}, as in the proof of Proposition~\ref{prop:bdrycohom}}.
Now read \eqref{eq:mvamalgam}: for every $q\ge0$ the map
$H^{q}(\langle S\rangle;\F_2)\oplus H^{q}(\langle U\rangle;\F_2)\to
H^{q}(\langle-I\rangle;\F_2)$ has a $2$-dimensional source and a $1$-dimensional target,
so its kernel is non-zero; by exactness $H^{q}(\Ga;\F_2)$ surjects onto that kernel, and
is therefore non-zero.
\end{proof}

\noindent
\emph{This settles, for $n=0$, the question Remark~\ref{rmk:ell2window}(2) left open.}
That remark records that no argument in this document bounds $H^{\bullet}(\Ga;\VR)$ above
at $\ell=2$ and does not assert non-vanishing in every degree; Lemma~\ref{lem:h2torsion}(2)
asserts it for the trivial module, by the amalgam alone and with no appeal to
Farrell--Tate cohomology. The general-$n$ statement at $\ell=2$ is \emph{not} claimed
here: the same count needs
$\dim\VR^{S}+\dim\VR^{U}>\dim\VR$, which holds at $n=0$ and is not automatic.

\subsection{The dimension of $H^{1}$, and the Dickson correction}

\begin{claim}[Dimension of $H^{1}$]\label{cl:dimH1}
Let $\ell\ge5$, $\cring=\Fl$, $n>0$ even, $k:=n+2$, and let
\[
  d_k:=\begin{cases}\lfloor k/12\rfloor & k\equiv2\pmod{12},\\
                    \lfloor k/12\rfloor+1 & \text{otherwise,}\end{cases}
  \qquad k\ \text{even},\ k\ge4,
\]
be the dimension of the space of level-$1$ modular forms of weight $k$, its subspace of
cusp forms having dimension $d_k-1$
\STD{\cite[Ch.~VII, \S3]{serre_cours}} --- the same $d_k$ as in
Obligation~\ref{obl:iota}(v). Then
\begin{enumerate}
\item \STD{\cite[Ch.~8]{shimura} for Eichler--Shimura over $\Cx$; \cite[Ch.~2]{cremona}
  for its modular-symbol form; \cite{deng} for the mod-$\ell$ statement with the
  correction term}
  \[
    \dim_{\Fl}H^{1}(\Ga;\Mn)
      \;=\;\underbrace{d_k}_{\text{modular forms}}
        \;+\;\underbrace{(d_k-1)}_{\text{cusp forms}}
        \;+\;\dim_{\Fl}H^{0}(\Ga;\Mn)
      \;=\;2d_k-1+\dim_{\Fl}H^{0}(\Ga;\Mn).
  \]
  The first two terms are the Eichler--Shimura decomposition of
  $H^{1}(\Ga;\Sym^{n}\Cx^{2})$ into the weight-$k$ level-$1$ modular forms and a
  conjugate copy of its cusp subspace, of dimensions $d_k$ and $d_k-1$
  (Obligation~\ref{obl:iota}(v) uses the same two numbers); the third is the Dickson
  correction of (2), invisible over $\Cx$.
\item \PROVED{} (v1.23.0-pre.15, proved in place, and since v1.23.0-pre.16 as Lemma~\ref{lem:dicksonSL2}, which holds at every prime $\ell$; until v1.23.0-pre.14 this clause was \STD{} on \cite{wilkerson}, a
  pointer never verified against the text, rule~\ref{A2}; Dickson's theorem, \cite{dickson1911}, \cite{wilkerson}, is the classical source, and
  nothing here rests on either pointer) \quad
  $\Fl[x,y]^{\Ga}=\Fl[x,y]^{\SL_2(\Fl)}$ --- the reduction
  $\Ga\to\SL_2(\Fl)$ being onto, since $\SL_2(\Fl)$ is generated by the elementary
  matrices $\smat{1}{c}{0}{1}$, $\smat{1}{0}{c}{1}$, which lift --- is a polynomial ring
  on two homogeneous invariants of degrees $\ell+1$ and $\ell^{2}-\ell$, so
  \[
    \dim_{\Fl}H^{0}(\Ga;\Mn)=\#\bigl\{(a,b)\in\Z_{\ge0}^{2}:
      a(\ell+1)+b(\ell^{2}-\ell)=n\bigr\},
  \]
  which, for $\ell\ge3$ (the standing hypothesis $\ell\ge5$ is more than needed), is $0$ for $0<n<\ell+1$ and first non-zero at $n=\ell+1$, where the invariant is
  $x^{\ell}y-xy^{\ell}$ (at $\ell=2$ it is non-zero already at $n=2$, the invariant $Q=d_2$: Lemma~\ref{lem:dicksonSL2}). The locus of $(n,\ell)$ on which it is non-zero is the
  \textbf{Dickson locus}; Remark~\ref{rmk:dickson} is where its small-$\ell$
  exceptional behaviour is recorded.
\end{enumerate}
\end{claim}

\begin{proof}[Proof of (2)]
Lemma~\ref{lem:dicksonSL2} below, stated and proved for every prime $\ell$ (v1.23.0-pre.16, referee D-1; in v1.23.0-pre.15 this proof stood here, and
the three sites that used it at $\ell=2$ or at every $\ell$ cited it). $\Fl[x,y]^{\Ga}=S^{G}$ because $\Ga$ acts through its image $\SL_2(\Fl)$, onto as stated, and the
count is that of the basis $\{L^{a}Q^{b}\}$ of the degree-$n$ part.
\end{proof}

\begin{lemma}[The Dickson invariants of $\SL_2(\Fl)$, every prime $\ell$]\label{lem:dicksonSL2}
\PROVED{} For \emph{every} prime $\ell$ (no hypothesis $\ell\ge5$, none on $n$), with $L:=x^{\ell}y-xy^{\ell}$ and $Q:=(x^{\ell^{2}}y-xy^{\ell^{2}})/L$,
\[
  \Fl[x,y]^{\SL_2(\Fl)}=\Fl[L,Q]\quad(\text{these are }\DkL,\DkD\text{ of Definition~\ref{def:dicksoninv}}),
\]
a polynomial ring on homogeneous generators of degrees $\ell+1$ and $\ell^{2}-\ell$, so $\dim_{\Fl}\bigl(\Fl[x,y]^{\SL_2(\Fl)}\bigr)_n=\#\{(a,b)\in\Z_{\ge0}^{2}:
a(\ell+1)+b(\ell^{2}-\ell)=n\}$ for every $n\ge0$; and $\Fl[x,y]^{\GL_2(\Fl)}=\Fl[L^{\ell-1},Q]$. (The count vanishes for $0<n<\ell+1$ only when $\ell\ge3$; at $\ell=2$
the degrees are $3$ and $2$ and the count is $1$ at $n=2$.)
\end{lemma}

\begin{proof}
(v1.23.0-pre.15, in place of the pointer \cite{wilkerson}; finding F1 of \file{docs/implementation/slot24/C-D1\_referee.md}; made a lemma at v1.23.0-pre.16, referee D-1.) Nothing below uses
$\ell\ge5$: it holds for every prime $\ell$. By Definition~\ref{def:rhon}, $\Mn$ is the degree-$n$ part of $S:=\Fl[x,y]$, $\gamma$ acting by the
$\Fl$-algebra automorphism with $(\gamma x,\gamma y)=(x,y)\gamma$, i.e.\ $\gamma x=ax+cy$, $\gamma y=bx+dy$. $\Ga$ acts through its image $\SL_2(\Fl)$ (onto, as
stated), so $\Fl[x,y]^{\Ga}=S^{G}$, where $G$ is the group of the automorphisms $g_A$, $(g_Ax,g_Ay)=(x,y)A$, $A\in\SL_2(\Fl)$; $A$ is read off from $g_Ax$, $g_Ay$, so
$\abs{G}=\abs{\SL_2(\Fl)}=\ell(\ell^{2}-1)$.
\emph{Step 1 (the invariant $L$).} Put $L:=x^{\ell}y-xy^{\ell}=-\det\smat{x}{y}{x^{\ell}}{y^{\ell}}$. As $a^{\ell}=a$ on $\Fl$,
$(g_Ax^{\ell},g_Ay^{\ell})=(x^{\ell},y^{\ell})A$, so the matrix $\smat{g_Ax}{g_Ay}{g_Ax^{\ell}}{g_Ay^{\ell}}$ is $\smat{x}{y}{x^{\ell}}{y^{\ell}}A$, and $g_AL=\det(A)\,L=L$.
\emph{Step 2 (the invariant $Q$; integrality).} Let $V:=\Fl x+\Fl y\subset S$ ($\ell^{2}$ elements) and $F(T):=\prod_{v\in V}(T-v)\in S[T]$. Each $g_A$ restricts to an
$\Fl$-linear bijection of $V$, so it permutes the factors, and the coefficients of $F$ lie in $S^{G}$. Let
$\Delta(T):=\det\begin{pmatrix}x&y&T\\x^{\ell}&y^{\ell}&T^{\ell}\\x^{\ell^{2}}&y^{\ell^{2}}&T^{\ell^{2}}\end{pmatrix}\in S[T]$. For $v=\alpha x+\beta y\in V$,
$v^{\ell^{i}}=\alpha x^{\ell^{i}}+\beta y^{\ell^{i}}$, so the last column at $T=v$ is a combination of the first two and $\Delta(v)=0$. The $v\in V$ are pairwise
distinct and $S$ is a domain, so dividing by the monic $T-v$ one root at a time gives $F\mid\Delta$ in $S[T]$; $\Delta$ has degree $\ell^{2}$ in $T$ with leading
coefficient $\det\smat{x}{y}{x^{\ell}}{y^{\ell}}=-L\ne0$, so $\Delta=-L\,F$. Expanding $\Delta$ along its last column,
$\Delta=(-L)^{\ell}\,T+L_2\,T^{\ell}-L\,T^{\ell^{2}}$ with $L_2:=x^{\ell^{2}}y-xy^{\ell^{2}}$; dividing by $-L$ (signs are immaterial at $\ell=2$),
\[
  F(T)\;=\;T^{\ell^{2}}-Q\,T^{\ell}+L^{\ell-1}\,T,\qquad Q:=L_2/L\in S,
\]
$Q$ homogeneous of degree $\ell^{2}-\ell$ and $G$-invariant, as a coefficient of $F$. ($L$, $Q$ are $\DkL$, $\DkD$ of Definition~\ref{def:dicksoninv},
Lemma~\ref{lem:dicksontwist}(2).) Put $\mathcal D:=\Fl[L,Q]\subseteq S^{G}$. $x$ and $y$ are roots of the monic $F\in\mathcal D[T]$, so $S=\mathcal D[x,y]$ is integral
over $\mathcal D$.
\emph{Step 3 (the fraction fields).} Let $\mathcal K:=\Fl(x,y)$, on which $G$ acts by field automorphisms, and $k:=\mathrm{Frac}\,\mathcal D\subseteq\mathcal K^{G}$.
\emph{Upper bound.} $x\ne0$ is a root of $F(T)/T=T^{\ell^{2}-1}-Q\,T^{\ell-1}+L^{\ell-1}\in\mathcal D[T]$, so $[k(x):k]\le\ell^{2}-1$; and $L=-x\,(y^{\ell}-x^{\ell-1}y)$, so
$y$ is a root of $T^{\ell}-x^{\ell-1}T+L/x\in k(x)[T]$ and $[\mathcal K:k(x)]\le\ell$, as $\mathcal K=k(x,y)$. Hence $[\mathcal K:k]\le\ell(\ell^{2}-1)=\abs{G}$.
\emph{Lower bound.} $\theta:=x+y^{2}\in\mathcal K$ has trivial stabiliser in $G$: $g_A\theta=\theta$ for $A=\smat abcd$ reads $(ax+cy)+(bx+dy)^{2}=x+y^{2}$ in $S$; comparing the coefficients of $x$, $y$ (degree $1$) and of
$x^{2}$, $xy$, $y^{2}$ (degree $2$) gives $a=1$, $c=0$, $b^{2}=0$, $2bd=0$, $d^{2}=1$, so $b=0$ (at $\ell=2$ the $xy$-condition is empty, but $b^{2}=0$ already forces it), and $\det A=d=1$. So the $\abs{G}$ elements $g\theta$ are pairwise
distinct. By the upper bound $\mathcal K$ is algebraic over $k\subseteq\mathcal K^{G}$, so $\theta$ has a minimal polynomial over $\mathcal K^{G}$, whose coefficients every $g$ fixes, so each $g\theta$ is one of its roots; hence $[\mathcal K:\mathcal K^{G}]\ge\abs{G}$. From $k\subseteq\mathcal K^{G}\subseteq\mathcal K$:
$\abs{G}\ge[\mathcal K:k]=[\mathcal K:\mathcal K^{G}]\,[\mathcal K^{G}:k]\ge\abs{G}\,[\mathcal K^{G}:k]$, so $\mathcal K^{G}=k$. Moreover $\mathcal K$, of transcendence
degree $2$ over $\Fl$, is algebraic over $k=\Fl(L,Q)$, so $L$ and $Q$ are algebraically independent: $\mathcal D$ is a polynomial ring in $L$, $Q$.
\emph{Step 4 (conclusion).} $S^{G}\subseteq\mathcal K^{G}=k$, and every $f\in S^{G}$ is integral over $\mathcal D$ (Step~2). The polynomial ring $\mathcal D$ is factorial,
hence integrally closed: if $f=a/b$ with $a,b\in\mathcal D$ coprime and $f^{m}+c_{m-1}f^{m-1}+\dots+c_0=0$ ($c_i\in\mathcal D$), then
$a^{m}=-b\,(c_{m-1}a^{m-1}+\dots+c_0b^{m-1})$, so every prime factor of $b$ divides $a$, and $b$ is a unit. Hence $S^{G}=\mathcal D=\Fl[L,Q]$, and since $L$, $Q$ are
algebraically independent and homogeneous of degrees $\ell+1$, $\ell^{2}-\ell$, the monomials $L^{a}Q^{b}$ with $a(\ell+1)+b(\ell^{2}-\ell)=n$ are a basis of
$(S^{G})_n=H^{0}(\Ga;\Mn)$: the displayed count. At $n=\ell+1$ (for $\ell\ge3$) the basis is $\{L\}$.
\emph{The $\GL_2(\Fl)$ statement} (v1.23.0-pre.16, referee D-3). Steps~1--2 hold verbatim for every $A\in\GL_2(\Fl)$, as $a^{\ell}=a^{\ell^{2}}=a$: $g_AL=\det(A)\,L$ and
$g_AL_2=\det(A)\,L_2$, so $g_AQ=Q$ and $g_A(L^{a}Q^{b})=\det(A)^{a}L^{a}Q^{b}$. The $L^{a}Q^{b}$ are a basis of $S^{\SL_2(\Fl)}$, so they are eigenvectors of $\GL_2(\Fl)$ and
the invariants are spanned by those with $\det(A)^{a}=1$ for all $A$, i.e.\ $(\ell-1)\mid a$ ($\det$ is onto the cyclic group $\Fl^{\times}$ of order $\ell-1$): $S^{\GL_2(\Fl)}=\Fl[L^{\ell-1},Q]$.
(An exact check of Steps~1--3 --- the identity for $F$, the invariance of $L$ and
$Q$, the trivial stabiliser of $\theta$, at $\ell\in\{2,3,5,7\}$ --- is \file{docs/implementation/slot24/cd1fix\_dickson\_check.py}; the referee's $800$-cell check of the
count is \file{C-D1\_referee.md} \S2. Neither is a tag.)
\end{proof}

\noindent
\COMP{} \textbf{Corroboration.} Both clauses were checked against the shipped
campaign table \file{docs/matrix-engine/results/results\_table.tsv} (2026-09-16): over
its $2118$ rows with $\ell\ge5$ and $n>0$ --- $710$ distinct pairs $(n,\ell)$ with
$n\le1000$ and $\ell\in\{5,7,11,13,691,10007,134217757,2^{61}-1,
9223372036854775783\}$, of which $203$ lie on the Dickson locus --- the emitted
$\dim\Heq{0}$ equals the count of (2) and the emitted $\dim\Heq{1}$ equals
$2d_k-1+\dim\Heq{0}$ in \emph{every} row, with $0$ disagreements, and $\dim\Heq{2}=0$
throughout. This is a re-reading of an existing artefact, not a new registered
validation; $\dim\Heq{\bullet}=\dim H^{\bullet}(\Ga;\Mn)$ uses
Remark~\ref{rmk:orbifold} and therefore $\ell\ge5$.

\begin{claim}[Torsion origin of the Dickson classes]\label{cl:bockstein}
\PROVED{} Let $n>0$, $\ell\ge5$, $\Vcoef=\Sym^{n}(\Z^{2})$ and $\Mn=\Vcoef\otimes\Fl$.
\begin{enumerate}
\item $H^{q}(\Ga;\Vcoef)$ is annihilated by $24$ for every $q\ge2$
  (Claim~\ref{cl:amalgam}(4)), hence has no $\ell$-torsion; and
  $H^{0}(\Ga;\Vcoef)=0$, an integral $\SL_2(\Z)$-invariant of positive degree being $0$.
\item Consequently the long exact cohomology sequence of
  $0\to\Vcoef\xrightarrow{\ell}\Vcoef\to\Mn\to0$ gives
  $H^{1}(\Ga;\Mn)=H^{1}(\Ga;\Vcoef)\otimes\Fl$ and
  $H^{1}(\Ga;\Vcoef)[\ell]\cong H^{0}(\Ga;\Mn)$, the second isomorphism being the
  Bockstein. \textbf{The Dickson classes are therefore reductions of $\ell$-torsion}:
  invisible to characteristic-zero Eichler--Shimura, which is why they appear in
  Claim~\ref{cl:dimH1}(1) as a correction term and not as forms.
\item Both isomorphisms are equivariant for every double-coset operator
  $\cor\circ(\conj_g)_{\ast}\circ\res$, $g\in\De$, hence --- under
  $6\in\cring^{\times}$, by Theorem~\ref{thm:classical} --- for $\Tdc{g}$ and every
  operator of \S\ref{sec:heckedatum} built from it.
\end{enumerate}
Reference for the classical form of (2): \cite{deng}.
\end{claim}

\begin{proof}
(1) The first half is Claim~\ref{cl:amalgam}(4) with $\cring=\Z$; $24$ is prime to
$\ell\ge5$. For the second, let $f\in\Vcoef^{\Ga}$, viewed as a binary form of degree $n$
over $\Q$. Invariance under $T$ reads $f(x,y)=f(x,x+y)$ (Definition~\ref{def:rhon}), hence
$f(x,y)=f(x,y+tx)$ for every $t\in\Z$; for fixed $x\neq0$ the polynomial $f(x,-)$ of
degree $\le n$ takes one value at infinitely many points, so it is constant and
$f\in\Q[x]$, i.e.\ $f=c\,x^{n}$. Invariance under $S$ then gives
$c\,x^{n}=\rhon(S)(c\,x^{n})=c\,y^{n}$, so $c=0$ for $n>0$.

(2) No cochain complex is needed: $0\to\Vcoef\xrightarrow{\ \ell\ }\Vcoef\to\Mn\to0$ is
a short exact sequence of $\Z[\De]$-modules ($\Vcoef$ is $\Z$-free, and multiplication by
$\ell$ is $\De$-equivariant), and $H^{\bullet}(\Ga;-)$ is a $\delta$-functor, so there is
a long exact sequence
\[
  \cdots\to H^{q}(\Ga;\Vcoef)\xrightarrow{\ \ell\ }H^{q}(\Ga;\Vcoef)\to
   H^{q}(\Ga;\Mn)\xrightarrow{\ \beta\ }H^{q+1}(\Ga;\Vcoef)\xrightarrow{\ \ell\ }\cdots ,
\]
$\beta$ its connecting map, i.e.\ the Bockstein \STD{\cite[Ch.~III, \S\S0--1]{brown}: the
long exact cohomology sequence of a short exact sequence of coefficient modules}. It
gives $0\to H^{q}(\Ga;\Vcoef)/\ell\to H^{q}(\Ga;\Mn)\to H^{q+1}(\Ga;\Vcoef)[\ell]\to0$ in
every degree. At $q=1$ the right-hand term is $H^{2}(\Ga;\Vcoef)[\ell]=0$ by (1), giving
$H^{1}(\Ga;\Mn)=H^{1}(\Ga;\Vcoef)/\ell=H^{1}(\Ga;\Vcoef)\otimes\Fl$; at $q=0$ the
left-hand term is $H^{0}(\Ga;\Vcoef)/\ell=0$ by (1), giving
$H^{0}(\Ga;\Mn)\cong H^{1}(\Ga;\Vcoef)[\ell]$ through $\beta$.

(3) Each of $\res$, $(\conj_g)_{\ast}$ and $\cor$ is a morphism of $\delta$-functors in
the coefficient module, so the composite commutes with the connecting maps of the
sequence in (2); the sequence is one of $\De$-modules, which is what makes
$(\conj_g)_{\ast}$ act on all three terms at once.
\end{proof}

\begin{remark}[The exponents of the Dickson torsion: $\ell^{2}$-torsion at $n=\ell m$, unbounded along $n=\ell^{j}m$]\label{rmk:bocksteinexp}
Claim~\ref{cl:bockstein} fixes the \emph{number} of $\ell$-primary cyclic summands of
$H^{1}(\Ga;\Vcoef)$ --- $\dim H^{0}(\Ga;\Mn)$, the Dickson count --- and says nothing about
their \emph{exponents}. This remark records what is known about them (finding F-A of
\file{docs/implementation/e8\_ladic\_investigation.md}, [L]; owner ruling 10 of
2026-09-24); \textbf{it changes no claim}. Let $n>0$ and $\ell$ be a prime. Clauses (1)
and (3) assume $\ell\ge5$, where Claim~\ref{cl:bockstein} is; clauses~(2) and (4) hold at
\textbf{every} prime $\ell$, $2$ and $3$ included (clause~(2) since v1.9.0 --- through v1.8.1 the remark assumed
$\ell\ge5$ throughout, clause~(2) included, although its proof never used it, slot-8 WP-9
finding~6, owner ruling~10~(A) of the slot-9 decisions, 2026-09-24; clause~(4) since v1.24.0-rc.1). For $\ell\ge5$ write the
$\ell$-primary torsion of $H^{1}(\Ga;\Vcoef)$ as $\bigoplus_{t=1}^{s}\Z/\ell^{a_t}$,
$s=\dim H^{0}(\Ga;\Mn)$.
\begin{enumerate}
\item \PROVED{} \textbf{(The Bockstein counts the exponent-one summands.)} The mod-$\ell$
  Bockstein $\beta^{0}\colon H^{0}(\Ga;\Mn)\to H^{1}(\Ga;\Mn)$, the connecting map of
  $0\to\Mn\xrightarrow{\ell}\Vcoef\otimes\Z/\ell^{2}\to\Mn\to0$, has rank
  $\#\{t:a_t=1\}$. So $\dim H^{0}(\Ga;\Mn)-\rank\beta^{0}=\#\{t:a_t\ge2\}$, and $\beta^{0}$
  is injective iff $H^{1}(\Ga;\Vcoef)$ has no $\ell^{2}$-torsion.
\item \PROVED{} \textbf{(The Frobenius lower bound, at every prime.)} Let $\ell$ be any
  prime. If $m>0$ and $H^{0}(\Ga;M_m)\neq0$, $M_m=\Sym^{m}\Z^{2}\otimes\Fl$,
  then for every $j\ge0$ the group $H^{1}(\Ga;\Sym^{\ell^{j}m}\Z^{2})$ contains an element of
  order $\ell^{j+1}$. In particular $H^{1}$ has $\ell^{2}$-torsion at $n=\ell m$, and the
  largest $\ell$-primary exponent is unbounded along $n=\ell^{j}m$. \emph{The instances at
  $\ell\in\{2,3\}$} (v1.9.0; the integral-line design brief's finding F-4): at $\ell=2$,
  $m=2$ and the invariant $x^{2}+xy+y^{2}$ of $M_2$ (Proposition~\ref{prop:ringinforce}(2))
  give an element of order $2^{j+1}$ at $n=2^{j+1}$; at $\ell=3$, $m=4$ and $m=6$ and the
  $\SL_2(\F_3)$-invariants $x^{3}y-xy^{3}$ and $x^{6}+x^{4}y^{2}+x^{2}y^{4}+y^{6}$ (the
  reduction $\Ga\to\SL_2(\F_3)$ being onto) give an element of order $9$ at $n=12$ and at
  $n=18$. (Through v1.8.1 this clause read ``the argument uses no hypothesis on $\ell$; it is
  stated here where Claim~\ref{cl:bockstein} is'', and was stated for $\ell\ge5$ only.)
\item \COMP\textsuperscript{a} \textbf{(The probe.)} [L] \S2.5 and \S5.4 record, from the
  exact out-of-tree probe \file{docs/implementation/e8\_ladic\_probe.py} (not a registered
  validation; certified at $\Z/\ell^{8}$), on its sweep $\ell=5$, $n\le130$; $\ell=7$,
  $n\le114$; $\ell\in\{11,13\}$, $n\le98$; and at $(150,5)$, $(180,5)$, $(200,5)$:
  $\Z/25$ at $(30,5)$ with $\rank\beta^{0}=0$ of one Dickson class; $\Z/25\oplus\Z/5$ at
  $(60,5)$, $\rank\beta^{0}=1$ of two; $\Z/125\oplus(\Z/5)^{2}$ at $(150,5)$,
  $\rank\beta^{0}=2$ of three; $\Z/49$ at $(56,7)$ and $(112,7)$; one $\Z/25$ among the
  classes at $(90,5)$, $(100,5)$, $(120,5)$, $(130,5)$, $(180,5)$, $(200,5)$; and exponent
  $1$ on every other torsion cell. On all of them the largest exponent \emph{equals} the
  lower bound of (2) --- $j+1$ at the largest $j$ with $n=\ell^{j}m$, $m$ a Dickson degree,
  and $1$ where $n$ has no such form --- and (1)'s count holds. (v1.9.0, beside this clause and
  not part of its scope $\ell\ge5$: at $\ell\in\{2,3\}$ the integral-line probe's table,
  Appendix~A of \file{docs/implementation/integral\_line\_design\_brief.md}, $n\le20$, shows
  $\Z/2^{j}$ at $n=2^{j}$, $1\le j\le4$, and $\Z/9$ at $n=12$ and $n=18$ --- the instances of
  (2) named there, each at its bound; the table is re-derived in the suite by
  \code{HeckeFreeModelValidation} (hecke-extensions slot-9 WP-11), which asserts the groups and
  not the equality with the bound.) \emph{The $\Z/\ell^{2}$ depth window} (added at v1.14.0-pre.1, taken out at pre.2, restored at
  v1.14.0; the first cluster run, published 2026-09-28 on the owner's sign-off, [S14] Q2, \file{docs/matrix-engine/results/noctua2\_run1/depth\_window\_l2.tsv};
  [EX] \texttt{nc:depth}): over $\Z/\ell^{2}$, at even
  $66\le n\le300$ for $\ell\in\{7,11,13\}$ and $66\le n\le200$ for $\ell\in\{5,691\}$, $\ell^{2}$-torsion occurs at exactly $17$
  cells, which are exactly the loci of (2) in that window --- $n=\ell m$ with $H^{0}(\Ga;M_m)\neq0$ ---, each with exactly one
  summand of exponent $\ge2$ (so (1)'s $\dim H^{0}(\Ga;\Mn)-\rank\beta^{0}=1$), and exactly where the $K=2$ mode is uncertified.
  (2) with $j=2$, $m=6$ forces an element of order $5^{3}$ at $(150,5)=(5^{2}\cdot6,5)$, which $\Z/25$ does not read. That there is no
  $\ell^{2}$-torsion off the loci is observed on this window only; this document does not prove it (item~(4) cites it from the literature).
  \emph{Since v1.24.0-rc.1} every cell of this clause, and of the $\ell\in\{2,3\}$ table beside it, is an instance of the theorem cited in~(4); the clause is a
  computational confirmation of that theorem, not an observation beyond it.
\item \STD{\cite[Thm.~4.5 (= Thm.~3.7(2))]{ccs13}, numbering of arXiv:1204.5390, read verbatim by the phenomenology literature search
  (\file{docs/phenomenology/literature\_search\_B.md} \S1), not by a referee of this document; \VERIFYloc} \textbf{(The exponents, from the literature;
  v1.24.0-rc.1, replacing ``not asserted'' for the first two items of the old list.)} For every prime $\ell$ and every $n>0$ the $\ell$-primary torsion of
  $H^{1}(\Ga;\Vcoef)$ is $\bigoplus\Z/\ell^{\,v_\ell(\gcd(k,h))+1}$, one summand for each Dickson monomial $P_\ell^{k}Q_\ell^{h}$ of degree $n$,
  i.e.\ each $(k,h)$ with $k(\ell+1)+h\,\ell(\ell-1)=n$, where $P_\ell=xy(x^{\ell-1}-y^{\ell-1})$ and $Q_\ell=P_{\ell^{2}}/P_\ell$ are the Dickson
  invariants. Consequences: the number of summands is $\dim H^{0}(\Ga;\Mn)$ (the Dickson monomials of degree $n$ are a basis of the $\SL_2(\Fl)$-invariants, onto which
  $\Ga$ maps), at every prime; and the largest exponent \emph{equals} the lower bound of~(2) --- some $(k,h)$ has $\ell^{j}\mid\gcd(k,h)$ iff $n=\ell^{j}m$ with
  $m=(k/\ell^{j})(\ell+1)+(h/\ell^{j})\ell(\ell-1)$ a Dickson degree --- so it is $j+1$ at the largest such $j$, and $1$ where $n$ has no such form with $j\ge1$ (and $\ell$-torsion at all).
  \emph{Convention.} The module of \cite{ccs13} is $\Z[x,y]$ graded with $\deg x=\deg y=2$ (v1.24.0-rc.2), so their degree-$2n$ part is our $\Vcoef$
  (polynomial degree $n$) and their condition $2k(\ell+1)+2h\,\ell(\ell-1)=2n$ is the one above; it carries the substitutions $\sigma_1\colon x\mapsto x-y$, $\sigma_2\colon y\mapsto x+y$
  (the other variable fixed), acting through $B_3\to\SL_2(\Z)$. In the convention of Proposition~\ref{prop:symconv}(1) these are $\rhon$ of $\smat{1}{0}{-1}{1}$ and
  of $\smat{1}{1}{0}{1}$, two matrices that generate $\Ga$; so their module is $(\Vcoef,\rhon)$ up to an automorphism of $\Ga$, which does not change $H^{1}$. The
  five integral Smith forms of the phenomenology catalogue (PH-07, $n=4,10,12,16,22$) agree with their multiset exactly, in polynomial degree (read literally, without the factor $2$, they would not: at $n=12$, $\ell=2$ the
  exponents are $\{1,2,3\}$, not $\{1,2\}$).
\item \textbf{Not asserted}: [L]'s estimate $1+\log_\ell\bigl(n/(\ell+1)\bigr)$ for the growth of the largest exponent (through v1.24.0-pre.1 this item also
  listed ``that the largest exponent is exactly $j+1$ in general'' and ``that it is $1$ off the loci of (2)'', now (4)). Claim~\ref{cl:bockstein} is
  unchanged. (1) is the \emph{meaning} the companion specification gives its
  $\Z/\ell^{e}$ bookkeeping identity, whose count $\#\{t:a_t\ge2\}$ is asserted where the
  exponent is certified (\Espec{sec:ring}); what (2) adds for those computations is that no
  fixed $e$ reads every exponent correctly, which is Proposition~\ref{prop:certifiede}'s
  subject.
\end{enumerate}
\end{remark}

\begin{proof}
(1) Reduction mod $\ell$, mod $\ell^{2}$ and the identity form a map of short exact
sequences from $0\to\Vcoef\xrightarrow{\ell}\Vcoef\to\Mn\to0$ to
$0\to\Mn\xrightarrow{\ell}\Vcoef\otimes\Z/\ell^{2}\to\Mn\to0$, so by naturality of connecting
maps \STD{\cite[Thm.~1.3.1]{weibel}} $\beta^{0}=\rho\circ\beta$, with $\beta$ the integral
Bockstein, an isomorphism $H^{0}(\Ga;\Mn)\to H^{1}(\Ga;\Vcoef)[\ell]$
(Claim~\ref{cl:bockstein}(2)), and $\rho\colon H^{1}(\Ga;\Vcoef)\to H^{1}(\Ga;\Mn)$ the
reduction, which factors as $H^{1}(\Ga;\Vcoef)\to H^{1}(\Ga;\Vcoef)/\ell\hookrightarrow
H^{1}(\Ga;\Mn)$ (the same claim's sequence at $q=1$). So $\rank\beta^{0}$ is the rank of
$G[\ell]\to G/\ell G$ for $G=H^{1}(\Ga;\Vcoef)$, whose kernel is $G[\ell]\cap\ell G$: the
free and prime-to-$\ell$ parts contribute nothing to $G[\ell]$, and the generator
$\ell^{a_t-1}$ of $(\Z/\ell^{a_t})[\ell]$ lies in $\ell G$ iff $a_t\ge2$.

(2) Lift a non-zero $\Ga$-invariant of $M_m$ to $f\in\Sym^{m}\Z^{2}\subset\Z[x,y]$, so
$\rho_m(\gamma)f-f\in\ell\,\Sym^{m}\Z^{2}$ for every $\gamma\in\Ga$. Each $\rho_\bullet(\gamma)=\Sym^{\bullet}(\gamma)$ is a ring endomorphism of
$\Z[x,y]$ (Definition~\ref{def:rhon}: the multiplicative extension), so for $h_i:=f^{\ell^{i}}$
induction gives $\rho(\gamma)h_i-h_i\in\ell^{i+1}\Z[x,y]$: if $\rho(\gamma)h_i=h_i+\ell^{i+1}u$
then $\rho(\gamma)h_{i+1}=(h_i+\ell^{i+1}u)^{\ell}$, and every term but $h_i^{\ell}$ is divisible
by $\binom{\ell}{k}\ell^{k(i+1)}$ with $1\le k\le\ell-1$, hence by $\ell^{i+2}$, or is
$\ell^{\ell(i+1)}u^{\ell}$, again divisible by $\ell^{i+2}$. So $h_j$ is a $\Ga$-invariant of
$\Vcoef\otimes\Z/\ell^{j+1}$ for $n=\ell^{j}m$, of additive order exactly $\ell^{j+1}$, since
$h_j\bmod\ell=(f\bmod\ell)^{\ell^{j}}\neq0$ in the domain $\Fl[x,y]$. The long exact
sequence of $0\to\Vcoef\xrightarrow{\ell^{j+1}}\Vcoef\to\Vcoef\otimes\Z/\ell^{j+1}\to0$ and
$H^{0}(\Ga;\Vcoef)=0$ (Claim~\ref{cl:bockstein}(1), $n>0$; that half of its proof is about
$\Z$-coefficients and uses no hypothesis on $\ell$) give an injection
$H^{0}(\Ga;\Vcoef\otimes\Z/\ell^{j+1})\hookrightarrow H^{1}(\Ga;\Vcoef)$ --- the argument of
Claim~\ref{cl:bockstein}(2) with $\ell^{j+1}$ in place of $\ell$ --- and the image of $h_j$ has
order $\ell^{j+1}$.
\end{proof}

\subsection{Odd $n$ at $\ell=2$, and the prime-to-$\ell$ Hecke action at $\ell\in\{2,3\}$}\label{sec:oddnltwo}

% v1.22.0-pre.10 (hecke-extensions slot 22, L23-0 extended, 2026-10-10): new subsection.  Source: the L23 master plan v0.2
% (docs/implementation/slot22/L23_master_plan.md, Theorem O) as repaired by its Fable review (L23_master_plan_review.md, dd3254f3,
% sec. 1-2, 5.1), and the census classification review [R] (census_l23_classification_review.md) Theorems A, B(1),(2), D(1),(2) and the
% census plan review [Rv] (census_l23_plan_v02_review.md) sec. 1.2.  Only what that review marks PROVED is stated as a theorem; the one
% CONJ clause is tagged.  Not refereed at pre.10 (Fable referee F-L23 pending).  No \LEAN marker.

\paragraph{What this subsection fixes.} The classes of $H^{\bullet}(\Ga;\Sym^{n})$ at the primes $2$ and $3$ that the rest of this section leaves open: the
complete integral cohomology for $n$ odd (it is an $\F_2$-vector space in every degree, of closed-form dimension, with $H^{1}\neq0$ for $n\ge3$), its
$\Z/2^{k}$ reductions, and the prime-to-$\ell$ Hecke action on $H^{q}(\Ga;\Sym^{n}\F_\ell^{2})$, $\ell\in\{2,3\}$, in every degree, which is the scalar
$1+p$. Throughout $\Ga=\langle S\rangle\ast_{\langle-I\rangle}\langle U\rangle$, $U:=ST$, as in Claim~\ref{cl:amalgam}; $A:=\langle S\rangle$, $B:=\langle U\rangle$,
$C:=\langle-I\rangle$ (orders $4$, $6$, $2$); $\Vcoef=\Sym^{n}\Z^{2}$ with $\rhon$ of Definition~\ref{def:rhon}; for a prime $p$,
$g_p:=\smat{1}{0}{0}{p}$ and $T_p:=\Tdc{g_p}=\cor^{\Ga}_{\Ga''}\circ(\conj_{g_p})_{\ast}\circ\res^{\Ga}_{\Ga'}$, the operator of
Lemma~\ref{lem:bulkperiodic}(3) on ordinary group cohomology, defined for every $\Z[\De]$-module (here $\Ga'=\Ga^{0}(p)$, $\Ga''=\Ga_0(p)$). All
cohomology is ordinary group cohomology. The census and plan documents that motivated this subsection are not cited as authority: every statement
below carries its own proof.

\begin{lemma}[A central element acting by $-1$ forces exponent $2$]\label{lem:centralminusone}
\PROVED{} Let $X$ be a $\Z[\Ga]$-module on which $-I$ acts as $-1$. Then $2\cdot H^{q}(\Ga;X)=0$ for every $q\ge0$. This applies to
$X=\Vcoef$, $\Vcoef/2^{k}\Vcoef$ ($k\ge1$) and $\Vcoef\otimes\Z_{(2)}$ for $n$ odd. Consequently, for $n$ odd and $k\ge1$, every
$H^{q}(\Ga;\Vcoef)$ and every $H^{q}(\Ga;\Vcoef/2^{k})$ is an $\F_2$-vector space, the long exact sequence of
$0\to\Vcoef\xrightarrow{2^{k}}\Vcoef\to\Vcoef/2^{k}\to0$ splits into short exact sequences
\[
  0\to H^{q}(\Ga;\Vcoef)\xrightarrow{\ \mathrm{red}_k\ }H^{q}(\Ga;\Vcoef/2^{k})\xrightarrow{\ \delta_k\ }H^{q+1}(\Ga;\Vcoef)\to0\qquad(q\ge0),
\]
and every $H^{q}(\Ga;\VR)$ vanishes when $2\in\cring^{\times}$ and $n$ is odd. \LEAN{thm:lv3-tree-negone} (only this last clause, and only \textbf{for a field}: if $K$ is a field with $2\ne0$ and $X$ a
finite-dimensional $K[\Ga]$-module on which $-I$ acts as $-1$, then $H^{q}(\Ga;X)=0$ for every $q$, mathlib's \texttt{groupCohomology}, proved on the
complex of Remark~\ref{rmk:treeconeres}; instantiated at $X=\Sym^{n}\F_3^{2}$, $n$ odd, by \texttt{V3\_dims\_odd} ([GT] \texttt{thm:lv3-v3-dims}, marked in Remark~\ref{rmk:lvdimsltwothree}(2)); \textbf{not} for a ring $\cring$ with
$2\in\cring^{\times}$ as stated here, and not $2\cdot H^{q}(\Ga;X)=0$, the $\F_2$-structure or the split sequences). (This generalises Proposition~\ref{prop:badprimes}(2), which is the case
$X=\Vcoef$.)
\end{lemma}

\begin{proof}
$z:=-I$ is central, so $(c_z,\ x\mapsto zx)$, $c_z(h)=z^{-1}hz=h$, is a compatible pair, and it induces the identity of $H^{\bullet}(\Ga;X)$
\STD{\cite[Ch.~III, \S8]{brown}: an inner automorphism, paired with the action of the conjugating element on the module, acts trivially on cohomology
(the form used in the proof of Proposition~\ref{prop:homvanish}(1); the pointer is at section level and is not verified against the text, rule~\ref{A2})}.
As $c_z=\id_{\Ga}$, the induced map is the one induced by the module automorphism $x\mapsto zx=-x$, i.e.\ $-1$. So $-1=1$ on $H^{q}(\Ga;X)$.
For $X=\Vcoef$ etc.\ with $n$ odd, $\rhon(-I)=(-1)^{n}=-1$ (Definition~\ref{def:rhon}). The short exact sequences follow because multiplication by
$2^{k}$ is $0$ on $H^{\bullet}(\Ga;\Vcoef)$; and if $2$ is a unit of $\cring$, multiplication by $2$ is invertible on $H^{q}(\Ga;\VR)$ and zero, so the group vanishes.
\end{proof}

\begin{lemma}[The degree-$0$ Bockstein, at every prime]\label{lem:bocksteinallell}
\PROVED{} Let $\ell$ be any prime (including $2$ and $3$) and $n>0$. The connecting map $\beta$ of $0\to\Vcoef\xrightarrow{\ell}\Vcoef\to\Mn\to0$ is an
isomorphism $H^{0}(\Ga;\Mn)\xrightarrow{\sim}H^{1}(\Ga;\Vcoef)[\ell]$, equivariant for every operator $\cor\circ(\conj_g)_{\ast}\circ\res$, $g\in\De$. The
reduction $H^{1}(\Ga;\Vcoef)\otimes\Fl\to H^{1}(\Ga;\Mn)$ is an isomorphism if and only if $H^{2}(\Ga;\Vcoef)[\ell]=0$ --- which holds for $\ell\ge5$
(Claim~\ref{cl:bockstein}(2)) and \textbf{fails at $\ell=2$ for every odd $n$} (Proposition~\ref{prop:oddnltwo}(3): $H^{2}(\Ga;\Vcoef)\neq0$ is $2$-torsion).
So Claim~\ref{cl:bockstein}(2)'s second isomorphism holds at every prime, its first one only where $H^{2}(\Ga;\Vcoef)$ has no $\ell$-torsion.
\end{lemma}

\begin{proof}
The long exact sequence \STD{\cite[Ch.~III, \S\S0--1]{brown}, as in the proof of Claim~\ref{cl:bockstein}(2)} gives, in every degree,
$0\to H^{q}(\Ga;\Vcoef)/\ell\to H^{q}(\Ga;\Mn)\to H^{q+1}(\Ga;\Vcoef)[\ell]\to0$. At $q=0$ the left term is $0$ because $H^{0}(\Ga;\Vcoef)=0$ for $n>0$ (the argument
of Claim~\ref{cl:bockstein}(1), which is about $\Z$-coefficients and uses no hypothesis on $\ell$); at $q=1$ the right term is $H^{2}(\Ga;\Vcoef)[\ell]$.
Equivariance is the argument of Claim~\ref{cl:bockstein}(3), which uses no hypothesis on $\ell$: $\res$, $(\conj_g)_{\ast}$ and $\cor$ are morphisms of
$\delta$-functors and the sequence is one of $\Z[\De]$-modules.
\end{proof}

\begin{lemma}[Cyclic groups: cohomology and restriction]\label{lem:cyclicres}
\PROVED{} Let $G=\langle t\rangle$ be cyclic of order $m$, $X$ a $\Z[G]$-module, $N_t:=1+t+\dots+t^{m-1}$, and let $d\mid m$, $h:=t^{d}$. Then
$H^{0}(G;X)=X^{G}$, $H^{2j+1}(G;X)=\ker N_t/(t-1)X$ and $H^{2j+2}(G;X)=X^{G}/N_tX$ for $j\ge0$, and
$\res^{G}_{\langle h\rangle}$ is induced in odd degrees by $x\mapsto(1+t+\dots+t^{d-1})x$ and in even degrees $\ge2$ by $x\mapsto x$.
\end{lemma}

\begin{proof}
The cohomology is that of $\Hom_G(-,X)$ on the periodic resolution $\cdots\xrightarrow{N_t}\Z[G]\xrightarrow{t-1}\Z[G]\to\Z$
\STD{\cite[Ch.~I.6 and Ch.~III.1]{brown}: the periodic resolution of a finite cyclic group; section-level pointer, not verified against the text,
rule~\ref{A2}}. A chain map from the periodic resolution of $\langle h\rangle$ (differentials $h-1$, $N_h$) to the restriction of that of $G$ lifting the
identity of $\Z$ is $f_{2j}=1$, $f_{2j+1}=1+t+\dots+t^{d-1}$: indeed $(t-1)(1+\dots+t^{d-1})=h-1$ and $N_t=(1+\dots+t^{d-1})N_h$. Applying
$\Hom(-,X)$ gives the stated maps.
\end{proof}

\begin{proposition}[Odd $n$: the integral cohomology is $\F_2$-elementary, in closed form]\label{prop:oddnltwo}
\PROVED{} Let $n$ be odd, $\Mn=\Sym^{n}\F_2^{2}$, $h^{0}:=\dim_{\F_2}\Mn^{\Ga}$, $c_U:=\dim_{\F_2}\Mn^{U}$.
\begin{enumerate}
\item \textbf{(Mayer--Vietoris collapses.)} $H^{0}(\Ga;\Vcoef)=0$ and, for \emph{every} $j\ge0$, \eqref{eq:mvamalgam} reduces to the exact sequence
  \[
    0\to H^{2j+1}(\Ga;\Vcoef)\to\Vcoef/(S-1)\Vcoef\;\oplus\;\Vcoef/(U-1)\Vcoef\xrightarrow{\ r\ }\Mn\to H^{2j+2}(\Ga;\Vcoef)\to0,
  \]
  $r(a,b)=(1+S)a+(1+U+U^{2})b\bmod 2$, with $\Vcoef/(S-1)\Vcoef\cong\F_2^{(n+1)/2}$, $\Vcoef/(U-1)\Vcoef\cong\Mn^{U}$, and both summands of $r$ injective, with images $\Mn^{S}$ and $\Mn^{U}$.
\item \textbf{(The invariants.)} $\Mn^{S}=(1+S)\Mn$ has dimension $(n+1)/2$; $c_U=2h^{0}$; and
  \[
    h^{0}=\bigl\lfloor\bigl(\lfloor n/3\rfloor+1\bigr)/2\bigr\rfloor ,
  \]
  which is $0$ for $n=1$ and $\ge1$ for $n\ge3$. \LEAN{chk:lv-cells} (only the \emph{value} of $h^{0}$, as $\dim_{\F_2}H^{0}(\Ga;\Mn)$ (mathlib's
  \texttt{groupCohomology}, which is $\Mn^{\Ga}$), at the odd $n\le63$ and at $n=101,129,255$, each by a kernel decision of two ranks; not the formula for general
  $n$, not the Dickson description; the same cells give $\dim_{\F_2}H^{q}(\Ga;\Mn)=(n+1)/2$ for every $q\ge1$, the value that (3) and the split sequences of
  Lemma~\ref{lem:centralminusone} predict; the integral groups $H^{q}(\Ga;\Vcoef)$, and (1), (3), (4), are \textbf{not} formalised; the other odd
  $n\le699$: Remark~\ref{rmk:lvdimsltwothree}) $\Mn^{\Ga}=\Mn^{\GL_2(\F_2)}$ is the degree-$n$ part of the Dickson algebra $\F_2[d_2,d_3]$,
  $d_2=x^{2}+xy+y^{2}$, $d_3=x^{2}y+xy^{2}$ (Lemma~\ref{lem:dicksonSL2} at $\ell=2$, with $L=d_3$, $Q=d_2$ and $\GL_2(\F_2)=\SL_2(\F_2)$, proved in place since v1.23.0-pre.15; no \STD{} tag: \cite{wilkerson} is provenance only, and the dimension formula above does not depend on it), spanned by the
  $d_3^{a}d_2^{b}$ with $3a+2b=n$ ($a$ odd).
\item \textbf{(The groups.)} For every $j\ge0$,
  \[
    H^{2j+1}(\Ga;\Sym^{n}\Z^{2})\cong\F_2^{\,h^{0}}\ \ (\text{zero iff }n=1),\qquad
    H^{2j+2}(\Ga;\Sym^{n}\Z^{2})\cong\F_2^{\,(n+1)/2-h^{0}}\ \ (\text{never zero}),
  \]
  and $H^{0}=0$. So $H^{\bullet}(\Ga;\Vcoef)$ is $2$-periodic from degree $1$ on, not only from degree $2$ as Lemma~\ref{lem:bulkperiodic}(2) gives in
  general: $e\cup-\colon H^{q}(\Ga;\Vcoef)\to H^{q+2}(\Ga;\Vcoef)$ is a Hecke-equivariant isomorphism for every $q\ge1$.
\item \textbf{(The odd-degree classes are Dickson classes.)} For every $j\ge0$ the composite
  $e^{j}\cup\beta\colon H^{0}(\Ga;\Mn)\to H^{2j+1}(\Ga;\Vcoef)$ ($\beta$ of Lemma~\ref{lem:bocksteinallell}) is an isomorphism, equivariant for every
  $\Tdc{g}$, $g\in\De$.
\end{enumerate}
\end{proposition}

\begin{proof}
(1) $V=\Vcoef$ is torsion-free and $-I=S^{2}=U^{3}$ acts on it by $-1$, so $V^{S}=V^{U}=V^{C}\subseteq V^{-I}=0$ and $N_S=(1+S)(1+S^{2})=0$,
$N_U=(1+U+U^{2})(1+U^{3})=0$, $N_{-I}=1+(-I)=0$ on $V$. By Lemma~\ref{lem:cyclicres} every even-degree group of $A$, $B$, $C$ (degree $0$ included)
vanishes, and $H^{2j+1}(A;V)=V/(S-1)V$, $H^{2j+1}(B;V)=V/(U-1)V$, $H^{2j+1}(C;V)=V/2V=\Mn$, with restriction maps to $C=\langle S^{2}\rangle=\langle U^{3}\rangle$
induced by $1+S$ and $1+U+U^{2}$ (the sign of the difference map in \eqref{eq:mvamalgam} is immaterial mod $2$). Exactness of \eqref{eq:mvamalgam} at the
vanishing even-degree terms gives the displayed sequences and $H^{0}(\Ga;V)\subseteq V^{S}=0$; at $j=0$ the left end uses $H^{0}(C;V)=0$.
Through $S\mapsto i$ (as $S^{2}=-1$ on $V$), $V$ is a torsion-free, hence free, $\Z[i]$-module of rank $(n+1)/2$, so $V/(S-1)V\cong(\Z[i]/(i-1))^{(n+1)/2}=\F_2^{(n+1)/2}$.
As $(U-1)(1+U+U^{2})=U^{3}-1=-2$ on $V$, $2V\subseteq(U-1)V$ and $V/(U-1)V=\Mn/(U-1)\Mn$. On $\Mn$, $U$ has order $3$ ($U^{3}=-I$ acts trivially mod $2$), so
$\pi_U:=1+U+U^{2}$ is idempotent ($\pi_U^{2}=3\pi_U=\pi_U$), $\pi_U\Mn=\Mn^{U}$, and $\Mn=\pi_U\Mn\oplus(1-\pi_U)\Mn$ with $U-1$ invertible on $(1-\pi_U)\Mn$ (no fixed vector there)
and zero on $\pi_U\Mn$; hence $(U-1)\Mn=(1-\pi_U)\Mn$ and $\Mn/(U-1)\Mn\cong\Mn^{U}$. The images of the two summands of $r$ are $(1+S)\Mn$ and $\pi_U\Mn=\Mn^{U}$;
$(1+S)\Mn=\Mn^{S}$ has dimension $(n+1)/2$ by (2), so both summands are surjections onto their images between spaces of equal dimension, i.e.\ injective.

(2) $\rhon(S)\bmod2$ substitutes $x\mapsto y$, $y\mapsto x$ (Definition~\ref{def:rhon}), so it swaps $x^{a}y^{n-a}$ and $x^{n-a}y^{a}$, which differ for $n$ odd:
$\Mn$ is a free $\F_2[\langle\bar S\rangle]$-module of rank $(n+1)/2$, and $\Mn^{S}=(1+S)\Mn$. The image of $U$ in $\SL_2(\F_2)\cong S_3$ generates the normal
subgroup $A_3$, so $SUS^{-1}$ acts on $\Mn$ as $U$ or $U^{2}$ and $S$ commutes with $\pi_U$: $\Mn^{U}=\pi_U\Mn$ is an $\F_2[\langle\bar S\rangle]$-direct summand of the free
module $\Mn$, hence projective, hence free ($\F_2[\Z/2]$ is local). Its $S$-invariants are therefore half of it, and they are $\Mn^{\langle S,U\rangle}=\Mn^{\Ga}$:
$c_U=2h^{0}$. Over $\F_4$, $\bar U=\smat{0}{1}{1}{1}$ has the eigenvalues $\omega,\omega^{2}$ ($\omega^{2}+\omega+1=0$), so on $\Sym^{n}$ its eigenvalues are
$\omega^{a}\omega^{2(n-a)}=\omega^{2n-a}$, $0\le a\le n$, and $c_U=\#\{0\le a\le n:\ a\equiv2n\pmod3\}$. With $n=3m+s$, $s\in\{0,1,2\}$: $c_U=m+1$ for $s\in\{0,2\}$ and
$c_U=m$ for $s=1$; $n$ odd forces $m$ odd in the first case and even in the second, so $h^{0}=c_U/2=\lfloor(m+1)/2\rfloor$ in both. For the Dickson description:
$\Mn^{\Ga}=\Mn^{\SL_2(\F_2)}$ since $\Ga\to\SL_2(\F_2)$ is onto, and the count $\#\{(a,b):3a+2b=n\}$ of the cited theorem agrees with $h^{0}$ (for $n$ odd, $a$
runs over the odd $a\le\lfloor n/3\rfloor$).

(3) By (1), $H^{2j+1}\cong\ker r$ and $H^{2j+2}\cong\coker r$; $r$ restricted to each summand is injective, so $\ker r\cong\Mn^{S}\cap\Mn^{U}=\Mn^{\Ga}$, of
dimension $h^{0}$, and $\dim\coker r=(n+1)-\dim(\Mn^{S}+\Mn^{U})=(n+1)-\tfrac{n+1}2-c_U+h^{0}=\tfrac{n+1}2-h^{0}$, which is $\ge1$ since
$h^{0}\le(m+1)/2<(n+1)/2$. The groups are $\F_2$-vector spaces by Lemma~\ref{lem:centralminusone}. For $e\cup-$: Lemma~\ref{lem:bulkperiodic}(2) gives an isomorphism
for $q\ge2$, and for $q=1$ a surjection with kernel $\partial(N_{-I}V)=\partial(0)=0$; equivariance is Lemma~\ref{lem:bulkperiodic}(3)(b).

(4) $H^{1}(\Ga;V)[2]=H^{1}(\Ga;V)$ by Lemma~\ref{lem:centralminusone}, so $\beta$ is an equivariant isomorphism onto $H^{1}(\Ga;V)$
(Lemma~\ref{lem:bocksteinallell}); then apply (3).
\end{proof}

\begin{lemma}[The $\Z/2^{k}$ reductions split, Hecke-equivariantly]\label{lem:oddnsplit}
\PROVED{} Let $n$ be odd, $k\ge2$ and $q\ge0$. The short exact sequence $0\to H^{q}(\Ga;\Vcoef)\to H^{q}(\Ga;\Vcoef/2^{k})\to H^{q+1}(\Ga;\Vcoef)\to0$ of
Lemma~\ref{lem:centralminusone} splits as a sequence of modules over the algebra generated by all $\Tdc{g}$, $g\in\De$ ($T_2$ included):
\[
  H^{q}(\Ga;\Vcoef/2^{k})\;\cong\;H^{q}(\Ga;\Vcoef)\oplus H^{q+1}(\Ga;\Vcoef).
\]
A retraction is $\varsigma:=\mathrm{red}_1^{-1}\circ\mathrm{red}_{k,1}$, $\mathrm{red}_{k,1}\colon H^{q}(\Ga;\Vcoef/2^{k})\to H^{q}(\Ga;\Vcoef/2)$ the reduction. Consequently the
$\Z/2^{k}$-cohomology, $k\ge2$, carries no information beyond $H^{q}(\Ga;\Vcoef)$ and $H^{q+1}(\Ga;\Vcoef)$ with their Hecke actions. For $k=1$ no
equivariant splitting is asserted.
\end{lemma}

\begin{proof}
Multiplication by $2^{k-1}$ on the left term, the identity in the middle and reduction on the right form a morphism of short exact sequences from
$0\to\Vcoef\xrightarrow{2^{k}}\Vcoef\to\Vcoef/2^{k}\to0$ to $0\to\Vcoef\xrightarrow{2}\Vcoef\to\Vcoef/2\to0$, so by naturality of connecting maps
\STD{\cite[Thm.~1.3.1]{weibel}, as in the proof of Remark~\ref{rmk:bocksteinexp}(1)} $\delta_1\circ\mathrm{red}_{k,1}=(2^{k-1})_{\ast}\circ\delta_k=0$ for $k\ge2$
(Lemma~\ref{lem:centralminusone}). Hence $\mathrm{red}_{k,1}$ lands in $\ker\delta_1=\mathrm{red}_1H^{q}(\Ga;\Vcoef)$, and $\mathrm{red}_{k,1}\circ\mathrm{red}_k=\mathrm{red}_1$; so
$\varsigma$ is defined and $\varsigma\circ\mathrm{red}_k=\id$. All of $\mathrm{red}_{k,1}$, $\mathrm{red}_k$, $\mathrm{red}_1$ commute with every $\Tdc{g}$ (maps of
$\Z[\De]$-modules; the argument of Claim~\ref{cl:bockstein}(3)), so $\ker\varsigma$ is a Hecke-stable complement, carried isomorphically onto $H^{q+1}(\Ga;\Vcoef)$ by
$\delta_k$.
\end{proof}

\begin{lemma}[$T_\ell$ has rank at most $2$ at $\ell\in\{2,3\}$]\label{lem:tellrank}
\PROVED{} Let $\ell\in\{2,3\}$, $n\ge0$, and $T_\ell=\Tdc{g_\ell}$ on $H^{q}(\Ga;\Mn)$. If $\ell=2$, or $\ell=3$ and $n$ is even, $T_\ell$ factors through
$H^{q}(\Ga_0(\ell);\Fl)$, whose dimension is
\[
  \dim_{\F_2}H^{q}(\Ga_0(2);\F_2)=1,2,2,2,\dots,\qquad \dim_{\F_3}H^{q}(\Ga_0(3);\F_3)=1,2,1,1,\dots\qquad(q=0,1,2,\dots).
\]
So $\rank T_2\le2$ on every $H^{q}$, $q\ge1$; $\rank T_3\le2$ on $H^{1}$ and $\le1$ on $H^{q}$, $q\ne1$; and $\rank T_\ell\le1$ on $H^{0}$. (For $\ell=3$, $n$ odd,
$H^{\bullet}(\Ga;\Mn)=0$ by Lemma~\ref{lem:centralminusone}.) Whether $T_\ell$ is a rank-one idempotent on $H^{1}$ is not asserted.
\end{lemma}

\begin{proof}
$\Ga'=\Ga\cap g_\ell^{-1}\Ga g_\ell=\Ga^{0}(\ell)$, $\Ga''=\Ga\cap g_\ell\Ga g_\ell^{-1}=\Ga_0(\ell)$, and $(\conj_{g_\ell})_{\ast}$ is induced by the pair
$(\alpha\colon\Ga''\to\Ga',\ h\mapsto g_\ell^{-1}hg_\ell;\ m\mapsto\rhon(g_\ell)m)$. Mod $\ell$, $\rhon(g_\ell)$ substitutes $x\mapsto x$, $y\mapsto0$, so
$\rhon(g_\ell)m=\varpi(m)x^{n}$ with $\varpi(m)$ the coefficient of $x^{n}$. For $h=\smat{a}{b}{\ell c}{d}\in\Ga_0(\ell)$, $\alpha(h)=\smat{a}{\ell b}{c}{d}$
substitutes $x\mapsto ax+cy$, $y\mapsto dy$ mod $\ell$, so $\varpi(\rhon(\alpha(h))m)=a^{n}\varpi(m)$; and $\rhon(h)x^{n}=a^{n}x^{n}$. Under the hypothesis
$a^{n}=1$ ($a\equiv1$ mod $2$; $a\equiv\pm1$ mod $3$ and $n$ even), the pair factors as
$(\alpha,\varpi)\colon(\Ga',\Mn)\to(\Ga'',\Fl)$ followed by $(\id,\ 1\mapsto x^{n})\colon(\Ga'',\Fl)\to(\Ga'',\Mn)$, both compatible, so
$(\conj_{g_\ell})_{\ast}$, and with it $T_\ell$, factors through $H^{q}(\Ga_0(\ell);\Fl)$ (\ref{Epair}, functoriality in the pair).

For the dimensions, restrict \eqref{eq:treeses} to $\Ga_0(\ell)$: each $\Z[\Ga/D]$, $D\in\{A,B,C\}$, becomes $\bigoplus\Z[\Ga_0(\ell)/\mathrm{Stab}]$ over the
$\Ga_0(\ell)$-orbits on $\Ga/D$, i.e.\ over the $D$-orbits on $\Ga_0(\ell)\backslash\Ga\cong\mathbb P^{1}(\F_\ell)$ (rows $(c:d)$, $\Ga$ acting on the right;
Definition~\ref{def:levelellmodel}), with stabilisers conjugate to the stabilisers in $D$ of the points; Shapiro's lemma \STD{\cite[Ch.~III.6]{brown}, as in
Claim~\ref{cl:amalgam}(2)} gives a long exact sequence
$\cdots\to H^{q}(\Ga_0(\ell))\to\bigoplus_{\text{$A$- and $B$-orbits}}H^{q}(\mathrm{Stab})\xrightarrow{r_q}\bigoplus_{\text{points}}H^{q}(C)\to H^{q+1}(\Ga_0(\ell))\to\cdots$
($C$ acts trivially on $\mathbb P^{1}$; $r_q$ is restriction, up to conjugations that act trivially on $H^{\bullet}(C)$). $\bar S\colon(c:d)\mapsto(d:-c)$ and
$\bar U\colon(c:d)\mapsto(d:d-c)$. \emph{$\ell=2$}, points $(0{:}1),(1{:}0),(1{:}1)$: $A$ fixes $(1{:}1)$ (stabiliser $A$) and swaps the other two (stabiliser $C$);
$B$ is transitive (stabiliser $C$). With $\F_2$-coefficients every term is $\F_2$; by Lemma~\ref{lem:cyclicres} $\res^{A}_{C}$ is $0$ in odd and the identity in
even degrees, $\res^{B}_{C}$ the identity ($1+U+U^{2}=3=1$) in odd and even degrees. In the basis (vertices $A{\cdot}(1{:}1)$, $A{\cdot}(0{:}1)$, $B{\cdot}(0{:}1)$;
edges $(1{:}1),(0{:}1),(1{:}0)$) $r_q$ has columns $(\epsilon,0,0)$, $(0,1,1)$, $(1,1,1)$ with $\epsilon=0$ for $q$ odd and $1$ for $q$ even: rank $2$ in every
degree, so $\ker r_q$ and $\coker r_q$ are lines and $\dim H^{q}(\Ga_0(2);\F_2)=\dim\coker r_{q-1}+\dim\ker r_q=2$ for $q\ge1$, $1$ for $q=0$.
\emph{$\ell=3$}, points $(0{:}1),(1{:}0),(1{:}1),(1{:}2)$: $A$ has the two orbits $\{(0{:}1),(1{:}0)\}$, $\{(1{:}1),(1{:}2)\}$ (stabilisers $C$); $B$ fixes $(1{:}2)$
(stabiliser $B$) and permutes the other three cyclically (stabiliser $C$). With $\F_3$-coefficients $H^{q}(C)=H^{q}(A)=0$ and $H^{q}(B)=\F_3$ for $q\ge1$
(Lemma~\ref{lem:cyclicres}); so $r_q=0$ on a line for $q\ge1$, and $r_0$ is the incidence map of a connected graph with four vertices and four edges (rank $3$).
Hence $\dim H^{q}(\Ga_0(3);\F_3)=1,\ 1+1,\ 0+1,\dots$.
\end{proof}

\begin{proposition}[$T_p=1+p$ as a matrix, at $\ell\in\{2,3\}$]\label{prop:ltwothreescalar}
Let $\ell\in\{2,3\}$, $n\ge0$ and $p\ne\ell$ a prime. Then
\[
  T_p\;=\;(1+p)\cdot\id\qquad\text{on }H^{q}(\Ga;\Sym^{n}\Fl^{2})\ \text{for every }q\ge0 ,
\]
\PROVED{} given the inputs \textup{(F1)}, \textup{(I1)} below; in particular $T_p=0$ for every odd $p$ at $\ell=2$; every $T_p$, $p\ne\ell$, is a scalar on
$H^{q}(\Ga;\Mn)$, so the algebra generated by all $T_p$ ($p$ prime) acts on $H^{q}(\Ga;\Mn)$ through $\Fl[T_\ell]$ (v1.22.0-pre.11, F-L23 R2). The statement is the matrix identity, not an eigenvalue statement.
\begin{itemize}
\item[(F1)] \STD{\cite[Ch.~I, \S4]{serre_trees} for the free subgroups of $\SL_2(\Z)$ and \cite[\S1.6]{shimura} for the cusps of $\Ga(N)$; section-level
  pointers, \textbf{not} verified against the text (rule~\ref{A2})} With $\Ga(\ell):=\ker(\Ga\to\SL_2(\Fl))$ and $T':=\smat{1}{0}{1}{1}$: $\Ga(2)=\{\pm I\}\times F$
  with $F$ free on $T^{2}$, $T'^{2}$; $\Ga(3)$ is free of rank $3$; and $\Ga(3)\backslash\Hyp^{\ast}$ has genus $0$ with four cusps. \emph{Derived from these}
  (v1.22.0-pre.11, F-L23 R4): $\Ga(3)^{\mathrm{ab}}=\Z[\text{cusps}]/(\textstyle\sum)$ on the parabolic classes, which $\Ga$ permutes by conjugation as $\SL_2(\F_3)$
  permutes $\mathbb P^{1}(\F_3)$; hence any three of them form a basis --- in particular the class of $T'^{3}$ (the parabolic at the cusp $0$) is non-zero in
  $\Ga(3)^{\mathrm{ab}}\otimes\F_3$ --- and $H^{1}(\Ga(3);\F_3)=\Hom(\Ga(3),\F_3)$ is the augmentation-zero submodule of $\F_3[\mathbb P^{1}(\F_3)]=\F_3\oplus\mathsf{St}_3$,
  $\mathsf{St}_3=\Sym^{2}\F_3^{2}$ (the Steinberg module of $\SL_2(\F_3)$ is projective, so the permutation module splits); so
  $H^{1}(\Ga(3);\F_3)\cong\Sym^{2}\F_3^{2}$ as $\F_3[\SL_2(\F_3)]$-modules ($\Ga$ acting by conjugation).
  \emph{Origin} (v1.24.0-rc.1): the mechanism --- $\Ga(\ell)$ acts trivially on $\Sym^{n}\Fl^{2}$, $X(\ell)$ has genus $0$, and $T_p$ acts on the cusps by $1+p$
  for $p\equiv\pm1\bmod\ell$ --- is \cite[Lemma~3.1]{calegari_emerton_largeindex} ($\ell=2,3,5$, on $W_k/\ell\otimes H^{1}(Y(\ell),\Z/\ell)$), which its authors
  credit to Hatada (their ref.~[11]; \cite{rustom_primepowers} quotes Hatada, Math.\ Ann.\ 239 (1979); that the two are the same paper is not checked; not seen); at $\ell\in\{2,3\}$ every prime $p\ne\ell$ is $\equiv\pm1\bmod\ell$. The
  level-one, every-degree form with (I1) is this document's own.
\item[(I1)] \STD{\cite[Ch.~III, \S9]{brown} (the transfer on inhomogeneous cochains, \ref{Ecor}); section-level pointer, not verified (rule~\ref{A2})} For every
  $\Z[\De]$-module $X$, $\Tdc{g}$ on $H^{1}(\Ga;X)$ is induced by the coset formula \eqref{eq:barcorrected} of Proposition~\ref{prop:barformula}, for any system of
  left-coset representatives $\delta_j$ of $\Ga g\Ga/\Ga$.
  That the class of $T_g\phi$ does not depend on the representatives is a cochain identity (v1.22.0-pre.11, F-L23 R3): replacing $\delta_j$ by $\delta_jb_j$,
  $b_j\in\Ga$, replaces $\kappa_j(\gamma)$ by $b_{\gamma\ast j}^{-1}\kappa_j(\gamma)b_j$, and the cocycle rule with $\phi(b^{-1})=-\rhon(b^{-1})\phi(b)$ gives
  $(T'_g\phi)(\gamma)=(T_g\phi)(\gamma)+\rhon(\gamma)w-w$, $w:=\sum_j\rhon(\delta_j)\phi(b_j)$ (the terms $-\rhon(\delta_{\gamma\ast j})\phi(b_{\gamma\ast j})$ re-index by
  the bijection $j\mapsto\gamma\ast j$). So (I1) is needed only for the representatives $\delta_j=s_jg$, $s_j$ running over $\Ga/\Ga''$ (those of
  Proposition~\ref{prop:cosets}(5) with $s_j^{-1}$ in place of $s_j$), where it is the transfer formula on inhomogeneous $1$-cocycles,
  $(\cor\psi)(\gamma)=\sum_{i}s_{\gamma\cdot i}\,\psi(h_i(\gamma))$ for $\gamma s_i=s_{\gamma\cdot i}h_i(\gamma)$, composed with $(\conj_g)_{\ast}$ and $\res$
  \STD{\cite[Ch.~III.9--10]{brown} and \cite[\S1.1, p.~194]{ash_stevens}, as in the proof of Proposition~\ref{prop:barformula}}. The clause is used: Step~0(c)'s
  $\delta_\infty=\pm\smat{p}{0}{0}{1}$ is not of the shape $s_jg_p$.
\end{itemize}
\end{proposition}

\begin{proof}
\emph{Step 0: two general facts.} (a) $H^{0}$: $T_p$ on $H^{0}(\Ga;X)=X^{\Ga}$ is $m\mapsto\sum_j\rhon(\delta_j)m$ over left-coset representatives of $\Ga g_p\Ga/\Ga$
(by \ref{Ecor}: $\cor$ on $H^{0}$ is $\sum_{s\in\Ga/\Ga''}s$, and $\{sg_p\}$ represents $\Ga g_p\Ga/\Ga$). (b) \emph{Inflated classes.} Let $X$ be an
$\Fl[\GL_2(\Fl)]$-module, $\De$ acting through reduction (on $\Mn$ this is $\rhon$, Definition~\ref{def:rhon}), $d_p:=g_p\bmod\ell$. For $\bar x\in H^{q}(\SL_2(\Fl);X)$,
\[
  T_p\,\mathrm{inf}(\bar x)=(p+1)\,\mathrm{inf}\bigl((c_{d_p},\,d_p)_{\ast}\bar x\bigr),
\]
$(c_{d_p},d_p)_{\ast}$ the automorphism of $H^{q}(\SL_2(\Fl);X)$ induced by the pair $(u\mapsto d_p^{-1}ud_p,\ x\mapsto d_px)$: $\Ga'$ and $\Ga''$ map onto $\SL_2(\Fl)$ (they
contain $\Ga(p)$, and $\SL_2(\Z)\to\SL_2(\Z/p\ell)=\SL_2(\F_p)\times\SL_2(\Fl)$ is onto), $\res$ and $(\conj_{g_p})_{\ast}$ commute with inflation by
functoriality in the pair (\ref{Epair}; $g_p^{-1}hg_p\equiv d_p^{-1}\bar hd_p$), and $\cor^{\Ga}_{\Ga''}\circ\mathrm{inf}=\cor^{\Ga}_{\Ga''}\circ\res^{\Ga}_{\Ga''}\circ\mathrm{inf}
=[\Ga:\Ga'']\,\mathrm{inf}$ \STD{\cite[Prop.~III.9.5(ii), p.~82]{brown}}. For $\ell=2$, and for $\ell=3$ with $p\equiv1$, $d_p=1$ and the pair is the
identity; for $\ell=3$, $p\equiv2$ the factor is $p+1\equiv0$. So
$T_p=1+p$ on $\mathrm{inf}\,H^{q}(\SL_2(\Fl);X)$ in all cases. (c) \emph{The level-$\Ga(\ell)$ operator.} Choose the left-coset representatives
$\delta_c=\smat{1}{0}{\ell c}{p}$ ($0\le c<p$) and $\delta_\infty=\pm\smat{p}{0}{0}{1}$, the sign making $\delta_\infty\equiv d_p$ mod $\ell$ (they represent the $p+1$
column lattices of index $p$; $c\mapsto\ell c$ is a bijection mod $p$). All $\delta_j\equiv d_p$, so $\kappa_j(\gamma)=\delta_{\gamma\ast j}^{-1}\gamma\delta_j\in\Ga(\ell)$
for $\gamma\in\Ga(\ell)$: restricted to $\Ga(\ell)$, \eqref{eq:barcorrected} only reads $\phi$ on $\Ga(\ell)$, and defines an operator $T^{(\ell)}$ on
$H^{1}(\Ga(\ell);X)$ (the proof of Proposition~\ref{prop:barformula} applies verbatim) with $\res^{\Ga}_{\Ga(\ell)}\circ T_p=T^{(\ell)}\circ\res^{\Ga}_{\Ga(\ell)}$ on
$H^{1}$, by (I1). Since $\Ga(\ell)$ acts trivially on $X$, $H^{1}(\Ga(\ell);X)=\Hom(\Ga(\ell),\Fl)\otimes X$ and $T^{(\ell)}=T^{(\ell)}_{0}\otimes d_p$, $T^{(\ell)}_0$
the operator for trivial coefficients. A change of representatives within the same classes mod $\ell$ ($\delta_j\mapsto\delta_ju_j$, $u_j\in\Ga(\ell)$) changes
$T^{(\ell)}_0\phi$ by $\sum_j(\phi(u_j)-\phi(u_{\gamma\ast j}))=0$; using the representatives $\gamma_1^{-1}\delta_j\gamma_0$ ($\gamma_0,\gamma_1\in\Ga$,
$\gamma_1\equiv d_p\gamma_0d_p^{-1}$) one finds $T^{(\ell)}_0\circ\gamma_0=\gamma_1\circ T^{(\ell)}_0$ for the conjugation action $(\gamma\phi)(h)=\phi(\gamma^{-1}h\gamma)$.
Finally, for $\gamma=T'^{\ell}$: $T'^{\ell}\delta_c=\delta_{c+1}$ for $c<p-1$, $T'^{\ell}\delta_{p-1}=\delta_0T'^{\ell}$, $T'^{\ell}\delta_\infty=\delta_\infty T'^{\ell p}$, so
\[
  (T^{(\ell)}_0\phi)(T'^{\ell})=\phi(T'^{\ell})+\phi(T'^{\ell p})=(1+p)\,\phi(T'^{\ell})\qquad(\phi\in\Hom(\Ga(\ell),\Fl)).
\]

\emph{$\ell=3$.} $T^{(3)}_0=(1+p)\id$ on $H^{1}(\Ga(3);\F_3)$: if $p\equiv1$, $d_p=1$ and $T^{(3)}_0$ is $\SL_2(\F_3)$-equivariant, hence (by (F1) and
$\End_{\SL_2(\F_3)}(\Sym^{2}\F_3^{2})=\F_3$: an endomorphism commutes with $\bar T$, whose only Jordan block on $\Sym^{2}$ has size $3$, so it is
$a+bN+cN^{2}$, $N=\bar T-1$, and it preserves the $\bar T'$-fixed line $\F_3y^{2}$, which forces $b=c=0$) a scalar, read off at $T'^{3}$ as $1+p$; if $p\equiv2$,
$\sigma_0:=\mathrm{diag}(1,-1)\in\GL_2(\Z)$ normalises $\Ga(3)$, has $\bar\sigma_0=d_p$, and $T^{(3)}_0\circ\sigma_0^{-1}$ is equivariant, so $T^{(3)}_0=\lambda\,\sigma_0$, and
$(1+p)\phi(T'^{3})=\lambda\phi(\sigma_0^{-1}T'^{3}\sigma_0)=-\lambda\phi(T'^{3})$ forces $\lambda=-(1+p)=0$. So $T^{(3)}=(1+p)\id$ on $H^{1}(\Ga(3);X)$. Now
$\Ga(3)$ is free, so the Lyndon--Hochschild--Serre spectral sequence of $\Ga(3)\to\Ga\to\SL_2(\F_3)$ \STD{\cite[Ch.~VII.6]{brown}} has two rows, $E_2^{a,1}=H^{a}(\SL_2(\F_3);
H^{1}(\Ga(3);\F_3)\otimes X)$, and these vanish for $a\ge1$: restriction to the Sylow subgroup $\langle\bar T\rangle$ is injective (index $8$, \cite[Prop.~III.9.5(ii)]{brown}),
and $\Sym^{2}\F_3^{2}\otimes X$ is free over $\F_3[\langle\bar T\rangle]$ (a free module tensored with any module is free). Hence inflation
$H^{q}(\SL_2(\F_3);X)\to H^{q}(\Ga;X)$ is onto for $q\ge2$ and bijective for $q=0$, and for $q=1$ there is an exact
$0\to\mathrm{inf}\,H^{1}(\SL_2(\F_3);X)\to H^{1}(\Ga;X)\xrightarrow{\res}H^{1}(\Ga(3);X)$. By Step~0(b) $T_p=1+p$ in degrees $q\ne1$ and on
$\mathcal I:=\mathrm{inf}\,H^{1}(\SL_2(\F_3);X)$. Put $D_p:=T_p-(1+p)$ on $H^{1}$: $\res\circ D_p=(T^{(3)}-(1+p))\circ\res=0$, so
$D_p(H^{1})\subseteq\mathcal I$; and $e\cup D_p(x)=D_p(e\cup x)=0$ (Lemma~\ref{lem:bulkperiodic}(3)(b); $T_p=1+p$ on $H^{3}$). $e\cup-$ is
injective on $\mathcal I$: $\res^{\Ga}_{B}$ is injective on $\mathcal I$ ($B\to\SL_2(\F_3)$ is injective with image of index $4$), $\res^{\Ga}_{B}e$ generates $H^{2}(B;\Z)$
(Lemma~\ref{lem:bulkperiodic}(1)), and cup product with a generator is injective on $H^{1}(B;X)$ (Step~5 of the proof of Lemma~\ref{lem:bulkperiodic}). So
$D_p=0$.

\emph{$\ell=2$.} Here $d_p=1$ and $\De_{\mathrm{odd}}$ acts on $\Mn$ through $\SL_2(\F_2)\cong S_3$. Put $\pi_U:=1+U+U^{2}$ (central idempotent of $\F_2[S_3]$, $A_3$ normal).
Then $\Mn=\pi_U\Mn\oplus(1+\pi_U)\Mn$; $\pi_U\Mn$ is an $\F_2[S_3/A_3]$-module, so $\pi_U\Mn\cong\F_2^{t}\oplus\mathsf P^{i}$ with $\mathsf P:=\F_2[S_3/A_3]$; and $(1+\pi_U)\Mn$ is a module over
$\F_2[S_3](1+\pi_U)$, which the action on $\mathsf{St}_2:=\Sym^{1}\F_2^{2}$ identifies with $\Mat_2(\F_2)$ (it is onto, as $\GL_2(\F_2)$ spans $\Mat_2(\F_2)$, and both have
dimension $4$), so $(1+\pi_U)\Mn\cong\mathsf{St}_2^{\,s}$. These are isomorphisms of $\Z[\De_{\mathrm{odd}}]$-modules, and $T_p$ is additive, so it suffices that
$T_p=0$ ($p$ odd) on $H^{q}(\Ga;X)$ for $X\in\{\F_2,\mathsf P,\mathsf{St}_2\}$. Two Mayer--Vietoris counts (\eqref{eq:mvamalgam} and Lemma~\ref{lem:cyclicres}; $\res^{B}_{C}$ is
bijective on $H^{\bullet}(-;\F_2)$ for trivial coefficients, $1+U+U^{2}$ acting as $3=1$): $\dim H^{q}(\Ga;\F_2)=1$ and $\dim H^{q}(\Ga;\mathsf P)=1$ for all $q$
($\mathsf P|_A=\Ind_C^{A}\F_2$, $\mathsf P|_B$, $\mathsf P|_C$ trivial).
\emph{$X=\F_2$:} $T_pe=T_p(e\cup1)=e\cup T_p1=(p+1)e$ in $H^{2}(\Ga;\Z)=\Z/12$ (Lemma~\ref{lem:bulkperiodic}(1),(3)), so $T_p\bar e=0$ for the reduction $\bar e\ne0$
($e\notin2H^{2}$), which spans $H^{2}(\Ga;\F_2)$; $H^{1}(\Ga;\F_2)\cong H^{2}(\Ga;\Z)[2]=\langle6e\rangle$ equivariantly (Bockstein, $H^{1}(\Ga;\Z)=0$), and
$T_p(6e)=6(p+1)e=0$; $H^{0}$: $T_p=p+1=0$; and $e\cup-$ is an equivariant isomorphism $H^{q}\to H^{q+2}$ for $q\ge1$ ($N_{-I}\F_2=0$).
\emph{$X=\mathsf P$:} in $0\to\F_2\to\mathsf P\to\F_2\to0$ (maps of $\Z[\De_{\mathrm{odd}}]$-modules), $H^{q}(\Ga;\mathsf P)$ is a line, so either it injects into $H^{q}(\Ga;\F_2)$ or it is the image of
$H^{q}(\Ga;\F_2)$; in both cases $T_p=0$ on it.
\emph{$X=\mathsf{St}_2$:} $H^{0}=0$. $q=1$: inflation--restriction \STD{\cite[Ch.~VII.6]{brown}} gives
$0\to H^{1}(S_3;\mathsf{St}_2)\to H^{1}(\Ga;\mathsf{St}_2)\xrightarrow{\res}H^{1}(\Ga(2);\mathsf{St}_2)^{S_3}$, and $H^{1}(S_3;\mathsf{St}_2)=0$ (restriction to the
Sylow $\langle\bar S\rangle$ is injective, and $\mathsf{St}_2$ is free over it). By (F1) $\mathcal W:=\Hom(\Ga(2),\F_2)$ is $3$-dimensional with $\phi\mapsto(\phi(-I),\phi(T^{2}),\phi(T'^{2}))$
bijective, and $T^{(2)}_0$ is $S_3$-equivariant (Step~0(c), $d_p=1$), so it preserves $\mathcal W_{\mathsf{St}}:=(1+\pi_U)\mathcal W$. $(\mathcal W\otimes\mathsf{St}_2)^{S_3}\subseteq \mathcal W_{\mathsf{St}}\otimes\mathsf{St}_2$
(on $\pi_U\mathcal W\otimes\mathsf{St}_2$ already $A_3$ has no invariant), and $\mathcal W_{\mathsf{St}}\ne0$ because $H^{1}(\Ga;\mathsf{St}_2)\supseteq\coker(\mathsf{St}_2^{S}\oplus\mathsf{St}_2^{U}\to\mathsf{St}_2)\ne0$
(\eqref{eq:mvamalgam}; $\mathsf{St}_2^{U}=0$, $\dim\mathsf{St}_2^{S}=1$) injects there. So $\mathcal W_{\mathsf{St}}\cong\mathsf{St}_2$ (a sum of copies, of dimension $\le3$), with
$\End_{S_3}(\mathsf{St}_2)=\F_2$ ($\bar U$ makes $\mathsf{St}_2$ a line over $\F_2[\bar U]=\F_4$, on which $\bar S$ is Frobenius-semilinear); $T^{(2)}_0=\lambda$ on it. If every
$\phi\in \mathcal W_{\mathsf{St}}$ vanished at $T'^{2}$, then also at $ST'^{2}S^{-1}=T^{-2}$, so $\mathcal W_{\mathsf{St}}\subseteq\{\phi:\phi(T'^{2})=\phi(T^{2})=0\}$, a line --- impossible;
so Step~0(c) gives $\lambda=1+p=0$, $T^{(2)}=0$ on $\mathcal W_{\mathsf{St}}\otimes\mathsf{St}_2$, and $T_p=0$ on $H^{1}(\Ga;\mathsf{St}_2)$ since $\res$ is injective. $q=2$:
$\mathsf{St}_2=\Sym^{1}\F_2^{2}$, and for $n=1$ Proposition~\ref{prop:oddnltwo}(3) gives $H^{1}(\Ga;\Sym^{1}\Z^{2})=H^{3}=0$, $H^{2}=\F_2$, so
Lemma~\ref{lem:centralminusone}'s sequences give equivariant isomorphisms $H^{1}(\Ga;\mathsf{St}_2)\xrightarrow{\delta_1}H^{2}(\Ga;\Sym^{1}\Z^{2})\xrightarrow{\mathrm{red}_1}H^{2}(\Ga;\mathsf{St}_2)$.
$q\ge3$: $e\cup-$ is an equivariant isomorphism for $q\ge1$ ($N_{-I}\mathsf{St}_2=0$).
\end{proof}

\begin{proposition}[Odd $n$ at $\ell=2$: the Hecke action]\label{prop:oddnhecke}
Let $n$ be odd and $k\ge1$.
\begin{enumerate}
\item \PROVED{} (given Proposition~\ref{prop:ltwothreescalar}'s inputs) For every odd prime $p$, $T_p=0$ on $H^{q}(\Ga;\Vcoef)$ and on $H^{q}(\Ga;\Vcoef/2^{k})$, every $q\ge0$.
\item \PROVED{} $T_2=0$ on $H^{0}(\Ga;\Mn)$ and on $H^{2j+1}(\Ga;\Vcoef)$ for every $j\ge0$.
\item \PROVED{} $\rank T_2\le2$ on every $H^{q}(\Ga;\Vcoef)$; and for $k\ge2$, $T_2$ on $H^{q}(\Ga;\Vcoef/2^{k})$ is $T_2|_{H^{q}}\oplus T_2|_{H^{q+1}}$ under
  Lemma~\ref{lem:oddnsplit}.
\item \CONJ{} $T_2$ is a rank-one idempotent on $H^{2j+2}(\Ga;\Vcoef)$ for every $j\ge0$ --- equivalently, by (2), the sequence of Lemma~\ref{lem:centralminusone}
  at $q=1$ and the periodicity of Proposition~\ref{prop:oddnltwo}(3): $T_2$ on $H^{1}(\Ga;\Mn)$, which kills $\mathrm{red}_1H^{1}(\Ga;\Vcoef)$, induces a rank-one
  idempotent on the quotient $H^{2}(\Ga;\Vcoef)$. The stronger form observed in the census classification review (\file{docs/implementation/slot22/census\_l23\_classification\_review.md}, Theorem~B(2)) and by \file{oddn\_l2.py} --- $T_2$ itself a rank-one idempotent on $H^{1}(\Ga;\Mn)$ whose image is not
  contained in $\mathrm{red}_1H^{1}(\Ga;\Vcoef)$ --- implies it (the image line meets the sub trivially) and is not asserted to be equivalent to it:
  Lemma~\ref{lem:tellrank} allows $\rank T_2=2$ on $H^{1}(\Ga;\Mn)$ with a rank-one induced map (v1.22.0-pre.11, F-L23 R1). Evidence: out-of-tree computations for every odd $n\le101$ (the L23 master plan's
  \file{oddn\_l2.py}); not a registered validation. If (4) holds, the Hecke algebra is $\F_2$ on odd degrees and $\F_2\times\F_2$ on even degrees with
  $(n+1)/2-h^{0}\ge2$; that it is $\F_2$ on odd degrees is (1)--(2).
\end{enumerate}
\end{proposition}

\begin{proof}
(1) $\mathrm{red}_1\colon H^{q}(\Ga;\Vcoef)\hookrightarrow H^{q}(\Ga;\Mn)$ is equivariant (Lemma~\ref{lem:centralminusone}), and $T_p=1+p=0$ on $H^{q}(\Ga;\Mn)$
(Proposition~\ref{prop:ltwothreescalar}); for $k\ge2$ use Lemma~\ref{lem:oddnsplit}, for $k=1$ it is the cited proposition.
(2) By Step~0(a) of the proof of Proposition~\ref{prop:ltwothreescalar}, $T_2m=\sum_j\rhon(\delta_j)m$ over left-coset representatives of determinant $2$, each of rank
$1$ mod $2$, $\delta_j\equiv u_jw_j^{\mathsf T}$; by Definition~\ref{def:rhon} $\rhon(\delta_j)$ substitutes $(x,y)\mapsto\lambda_j(x,y)\,w_j^{\mathsf T}$ with
$\lambda_j$ a linear form, so $\rhon(\delta_j)m=m(w_j)\,\lambda_j^{n}$. For $m\in\Mn^{\Ga}$ the value $m(w)$ is the same for the three non-zero $w\in\F_2^{2}$
($\Ga$ permutes them transitively and $(\rhon(\gamma)m)(w)=m(w\bar\gamma)$), say $c_m$; and $m\in\Mn^{S}=(1+S)\Mn$ (Proposition~\ref{prop:oddnltwo}(2)), $m=f+\rhon(S)f$,
so $3c_m=\sum_{w\ne0}m(w)=\sum_{w\ne0}f(w)+\sum_{w\ne0}f(w\bar S)=0$. Hence $T_2=0$ on $H^{0}(\Ga;\Mn)$, and on $H^{2j+1}(\Ga;\Vcoef)$ by the equivariant
isomorphism of Proposition~\ref{prop:oddnltwo}(4).
(3) $\mathrm{red}_1$ is an equivariant injection into $H^{q}(\Ga;\Mn)$, where Lemma~\ref{lem:tellrank} applies; the second half is Lemma~\ref{lem:oddnsplit}.
\end{proof}

\begin{lemma}[The minimal divisor of $T_p-(1+p)$ is not an independent datum]\label{lem:cprule}
\PROVED{} Let $\ell$ be a prime, $\mathcal N$ a free $\Z_\ell$-module of finite rank and, for $Y\in\End(\mathcal N)$, $\nu(Y):=\sup\{m:Y\mathcal N\subseteq\ell^{m}\mathcal N\}\in\Z_{\ge0}\cup\{\infty\}$,
the least exponent among the elementary divisors of $Y$ ($\infty$ iff $Y=0$).
\begin{enumerate}
\item If $Y\in\End(\mathcal N)$ kills a primitive vector $e_0$ and $c\in\Z_\ell$, then $\nu(Y+c)=\min\bigl(\nu(Y),v_\ell(c)\bigr)$.
\item Let $n>0$ be even, $p\ne\ell$ a prime, $\mathcal N:=\bigl(H^{1}(\Ga;\Vcoef)/\mathrm{tors}\bigr)\otimes\Z_\ell$ with its $T_p$, and assume $\mathcal N\otimes\Q_\ell$ contains a
  $T_p$-eigenvector of eigenvalue $1+p^{\,n+1}$ (the Eisenstein class \STD{\cite[Ch.~2--3]{harder3}, the pointer of Obligation~\ref{obl:iota}}). Then
  \[
    \nu\bigl(T_p-(1+p)\bigr)\;=\;\min\Bigl(\nu\bigl(T_p-(1+p^{\,n+1})\bigr),\ v_\ell(p^{n}-1)\Bigr).
  \]
\end{enumerate}
\end{lemma}

\begin{proof}
(1) If $Y\equiv0$ and $c\equiv0$ mod $\ell^{m}$ then $Y+c\equiv0$. Conversely, if $(Y+c)\mathcal N\subseteq\ell^{m}\mathcal N$ then $ce_0=(Y+c)e_0\in\ell^{m}\mathcal N$, so
$v_\ell(c)\ge m$ ($e_0$ primitive), and then $Y\mathcal N\subseteq\ell^{m}\mathcal N$. (2) The torsion of $H^{1}(\Ga;\Vcoef)$ is stable under every endomorphism, so $\mathcal N$ carries $T_p$;
the saturation of the eigenline in $\mathcal N$ is generated by a primitive $e_0$ with $(T_p-(1+p^{n+1}))e_0=0$. Apply (1) with $Y=T_p-(1+p^{n+1})$ and
$c=p^{n+1}-p=p(p^{n}-1)$, $v_\ell(c)=v_\ell(p^{n}-1)$.
\end{proof}

\begin{remark}[The standard twist: odd $n$ inside even $n$]\label{rmk:standardtwist}
\PROVED{} Let $n$ be even and $\cring=\Z_{(2)}$. The Koszul sequence $0\to\Sym^{n-1}\otimes{\det}\to\Sym^{n}\otimes\Sym^{1}\to\Sym^{n+1}\to0$ (multiplication on the
right, $f\otimes(x\wedge y)\mapsto fx\otimes y-fy\otimes x$ on the left) is exact and natural for every $A\in\Mat_2(\Z)$, and the comultiplication
$\Sym^{n+1}\to\Sym^{n}\otimes\Sym^{1}$, $f\mapsto\partial_xf\otimes x+\partial_yf\otimes y$, composed with multiplication is $(n+1)\id$; as $n+1$ is a unit,
\[
  \Sym^{n}\cring^{2}\otimes\Sym^{1}\cring^{2}\;\cong\;\Sym^{n+1}\cring^{2}\;\oplus\;\Sym^{n-1}\cring^{2}\otimes{\det}
\]
as $\cring[\Mat_2(\Z)]$-modules, so $H^{q}(\Ga;\Sym^{n}\otimes\Sym^{1})$ is the sum of $H^{q}(\Ga;\Sym^{n+1})$ and $H^{q}(\Ga;\Sym^{n-1})$ over $\cring$. On the second summand
$\Tdc{g}$ acts as $\det(g)\cdot\Tdc{g}$ (the pair $(\conj_g)_{\ast}$ carries the scalar $\det g$). For odd $p$ this factor is invisible mod $2$; \textbf{for $p=2$ it is
not}: the operator induced by $T_2$ on the $\det$-summand is $2\,T_2^{(n-1)}$, which is $0$, since $H^{q}(\Ga;\Sym^{n-1}\cring^{2})$ has exponent $2$
(Lemma~\ref{lem:centralminusone}).
\end{remark}

\begin{remark}[An erratum to a project review, and the scope of this subsection]\label{rmk:oddnltwoerr}
The census classification review \file{docs/implementation/slot22/census\_l23\_classification\_review.md}, Theorem~B(6), stated $H^{1}(\Ga;\Sym^{n}\Z^{2})=0$ and
$(\Z/2)^{(n+1)/2}$ in every even degree $\ge2$, for $n$ odd. By Proposition~\ref{prop:oddnltwo}(3) that is true only for $n=1$: for every odd $n\ge3$,
$H^{1}\cong\F_2^{h^{0}}\neq0$ ($H^{1}(\Ga;\Sym^{3}\Z^{2})=\Z/2$, generated by the Bockstein of $d_3=xy(x+y)$) and the even degrees carry $(n+1)/2-h^{0}$.
Its universal-coefficient recursion was started at $\dim H^{0}(\Ga;\Mn)=(n+1)/2$ instead of $h^{0}$; that review carries a dated erratum. No statement of this
document depended on it. Not stated here, because they are not proved: $T_\ell$ a rank-one idempotent on $H^{1}(\Ga;\Mn)$ for even $n$ (the analogue of
Proposition~\ref{prop:oddnhecke}(4)); a growth law for the torsion exponents at $2$ and $3$ (cf.\ Remark~\ref{rmk:bocksteinexp}(4)); the Gouv\^ea bound on slopes; the
$\ell=3$ comparison of the $\pm$ halves of the lattice $H^{1}(\Ga;\Vcoef)/\mathrm{tors}\otimes\Z_3$; and the Hodge--Newton comparison of elementary divisors and slopes, which will be written with its
Eichler--Shimura input in a later version.
\end{remark}

\begin{remark}[The dimensions at $\ell\in\{2,3\}$, $n\le700$, formalised]\label{rmk:lvdimsltwothree}
% v1.22.0-pre.14 (hecke-extensions slot 22, 2026-10-11): new remark, worded within LV2_audit.md sec. 8 (Fable, 3d7b868d); the LV-0..LV-2
% Lean packages (docs/implementation/slot22/LV1_report.md, LV2_report.md), provenance repairs LVprov_audit.md (a2ea84b5).  The L23Cells
% cells n > 64 are cited as data (docs/implementation/slot22/lv/), never by a \LEAN marker: L23Cells is outside lean_census.sh.
\PROVED{} Let $\ell\in\{2,3\}$, $\Mn=\Sym^{n}\F_\ell^{2}$ with $g\mapsto\rhon(g)\bmod\ell$ (\eqref{eq:rhondef}), $d=n+1$, and
$h^{q}:=\dim_{\F_\ell}H^{q}(\Ga;\Mn)$. On the complex $V^{2}\to V^{3}\to V^{3}\to\cdots$ of Remark~\ref{rmk:treeconeres} the $h^{q}$ are ranks of explicit
matrices over $\F_\ell$, by rank--nullity. Write $s=\rank(1-S,1-U)\colon\Mn\to\Mn^{2}$.
\begin{enumerate}
\item \textbf{($\ell=2$.)} $-I$ acts trivially and $N_S=N_U=e_m=0$; with $a=\rank\psi$, $\psi(x,y)=(1+S)x-(1+U+U^{2})y$,
  \[
    h^{0}=d-s,\qquad h^{q}=2d-a-s\quad\text{for every }q\ge1 .
  \]
  \LEAN{thm:lv-tree-f2-dims} (for every finite-dimensional $\F_2[\Ga]$-module on which $-I$ acts trivially, with $\dim$ in place of $d$; the degree-$0$ identity for every finite-dimensional $\F_2[\Ga]$-module, the others when $-I$ acts trivially)
  \LEAN{thm:lv-f1} (the matrices of $1+S$, $1+U+U^{2}$, $1-S$, $1-U$ on $\Mn$, bit-packed, are those of \eqref{eq:rhondef} reduced mod $2$, every $n$)
  \LEAN{thm:lv-v2-dims} (the formula for $\Sym^{n}\F_2^{2}$, every $n$, $a$ and $s$ computed by \texttt{rankF2}, proved equal to the rank over $\F_2$)
\item \textbf{($\ell=3$.)} For $n$ odd, $h^{q}=0$ for every $q$ (Lemma~\ref{lem:centralminusone}'s last clause and Proposition~\ref{prop:ellthreeclosed}(4), at
  $\F_3$). For $n$ even $-I$ acts trivially; with $r_S^{+}=\rank(1+S)$, $r_3=\rank(1+U+U^{2})$, $r_U^{-}=\rank(1-U)$,
  \[
    h^{0}=d-s,\qquad h^{1}=2d-s-r_S^{+}-r_3,\qquad h^{q}=d-r_3-r_U^{-}\quad\text{for every }q\ge2 .
  \]
  \LEAN{thm:lv3-tree-triv} (for every field $K$ with $2\ne0$ and every finite-dimensional $K[\Ga]$-module on which $-I$ acts trivially; proved on the complex of Remark~\ref{rmk:treeconeres}, the resolution of the amalgam of Claim~\ref{cl:amalgam}; over $\F_3$ these are the dimensions of the terms of Proposition~\ref{prop:ellthreeclosed}(2) --- a count, not a
  formalisation of its isomorphisms or of its sequence)
  \LEAN{thm:lv3-rankf3} (\texttt{rankF3} is the rank over $\F_3$)
  \LEAN{thm:lv3-f1} (the matrices of $1\pm S$, $1-U$, $1+U+U^{2}$ on $\Mn$, as two bit-planes, are those of \eqref{eq:rhondef} reduced mod $3$, every $n$)
  \LEAN{thm:lv3-v3-dims} (the vanishing for every odd $n$ and the formula for every even $n$, for $\Sym^{n}\F_3^{2}$)
\item \textbf{(The cells.)} \LEAN{chk:lv-cells} ($\ell=2$: $h^{0}$ and $h^{q}$, $q\ge1$, as numbers, for $n\le64$ and $n=100,101,128,129,255,256$)
  \LEAN{chk:lv3-cells} ($\ell=3$: the same for every even $n\le64$)
  For every other $n\le700$ at $\ell=2$, and every other even $n\le700$ at $\ell=3$, the same kernel decision (\texttt{decide +kernel}, no
  \texttt{native\_decide}) is made in the Lean library \texttt{L23Cells}, which is outside [GT]'s census; those cells are \textbf{data, not \code{\textbackslash LEAN}
  evidence}: values, modules, ledger commits and replay logs are \file{docs/implementation/slot22/lv/lean\_dims\_l2.tsv}, \file{lean\_dims\_l3.tsv} and the ledgers
  \file{lv1\_ledger.tsv}, \file{lv2\_ledger.tsv}, \file{lv3\_ledger.tsv} there. So $h^{0},\dots,h^{3}$ are formalised for $\ell=2,3$ and every $n\le700$ ($701$ values
  of $n$ at each $\ell$; at $\ell=3$, $n$ odd, by the theorem), as Lean theorems on mathlib's \texttt{groupCohomology} with the coefficient module of
  \eqref{eq:rhondef}, standard axioms only.
\end{enumerate}
This is a \textbf{different model} from the census's: the Lean values are computed through the tree complex, the C++ values on the isotropy-resolution model
$\CI$ of Proposition~\ref{prop:isomodel}. They agree wherever both exist --- read at the census snapshot of 2026-10-11, 22:47, recorded in \file{lean\_dims\_l2.tsv} / \file{lean\_dims\_l3.tsv} at commit \texttt{9e07e1a0}: at $\ell=2$, $264$ computed records, all at even $n\le528$; at $\ell=3$, $61$ ($31$ even $n\le60$, $30$ odd); no mismatch --- so the agreement is a cross-model check, not a re-run; later census passes add records, and a \texttt{census\_mismatch} on any of them would be a contradiction between the kernel and a C++ record, to be reported. At $\ell=2$, $n$ odd, the census rows are closed forms, not computations: there the Lean cell is the first
computation of the number, and a per-cell check of Proposition~\ref{prop:oddnltwo}'s values ($h^{0}$, and $h^{q}=(n+1)/2$, $q\ge1$; all $350$ odd $n\le699$ agree with the closed forms, same snapshot).
Nothing here concerns the integral groups $H^{q}(\Ga;\Vcoef)$, the Hecke operators (Proposition~\ref{prop:ltwothreescalar} carries no marker), or any census
column other than $h^{0},\dots,h^{3}$; and no clause of this document is superseded.
\end{remark}

\subsection{The boundary groups, and the long exact sequence}

\begin{claim}[Boundary groups]\label{cl:boundary}
\PROVED{} Let $\cring$ be a field of characteristic $\ell$, $n$ even, and
$2\in\cring^{\times}$ (i.e.\ $\ell\neq2$). Then
\[
  H^{0}(\partial;\VR)=\VR^{T},\qquad H^{1}(\partial;\VR)=(\VR)_{T},
  \qquad H^{q}(\partial;\VR)=0\ \ (q\ge2),
\]
both non-zero groups of dimension $\lfloor n/\ell\rfloor+1$, and
\[
  \VR^{T}=\cring[x,\mathrm{Nm}]_{n}
   =\bigoplus_{0\le b\le\lfloor n/\ell\rfloor}\cring\cdot x^{\,n-b\ell}\mathrm{Nm}^{\,b},
  \qquad \mathrm{Nm}=y^{\ell}-x^{\ell-1}y .
\]
\end{claim}

\begin{proof}
$\Gbdry=\langle T\rangle\times\langle-I\rangle$ and $\rhon(-I)=\id$ for $n$ even, so with
$2\in\cring^{\times}$ Proposition~\ref{prop:circlemodel}(3) identifies
$H^{q}(\Gbdry;\VR)$ with $H^{q}(\langle T\rangle;\VR)$ for every $q$, and
Proposition~\ref{prop:circlemodel}(2) evaluates the latter as $\VR^{T}$, $(\VR)_{T}$,
$0,0,\dots$. The description of $\VR^{T}$ and its rank are
Proposition~\ref{prop:bdrychar}(1); $\dim(\VR)_{T}=\dim\VR^{T}$ because both are the
kernel and the cokernel of one endomorphism of a finite-dimensional space;
non-vanishing is the nilpotence argument in the proof of
Lemma~\ref{lem:ell2unbounded}, which used no hypothesis on $\ell$.
\end{proof}

\begin{remark}[Two repairs in Claim~\ref{cl:boundary}]\label{rmk:bdryrepairs}
\textbf{(i)} [E] states ``$H^{i}(\partial;M)=0$ for $i\ge2$'' under its blanket
$\ell\ge5$. As a statement about $H^{\bullet}(\Gbdry;\VR)$ that is \emph{false at
$\ell=2$}, where every degree is non-zero (Lemma~\ref{lem:ell2unbounded}); the
hypothesis that carries it is $2\in\cring^{\times}$, and it is now in the statement.
\textbf{(ii)} The invariant [E] writes $N$ is written $\mathrm{Nm}$ here: $N=p^{\nu}$ is
the Hecke level (\S\ref{sec:notation-primes}), and the two meet in
Proposition~\ref{prop:bdrychar}. This is collision \textbf{C18} of
\S\ref{sec:notation-collisions}.
\end{remark}

\begin{claim}[The six-term sequence]\label{cl:les}
Assume \ref{H4} and \ref{B1}--\ref{B3} and let $\cring$ be a commutative ring.
\begin{enumerate}
\item \PROVED{} \textbf{(The sequence, in all degrees, with no hypothesis on $\ell$ or
  $n$.)} The short exact sequence of complexes
  $0\to C^{\bullet}_{\partial}\xrightarrow{\coneinc}\Cone{\rres}
   \xrightarrow{\coneproj}\Ceq(K,\Rel;L)[1]\to0$ of
  Lemma~\ref{lem:conerelative}(2) gives the long exact sequence \eqref{eq:coneles},
  \[
    \cdots\to\Hcpt{q}\to\Heq{q}(K,\Rel;L)\xrightarrow{\ \rres^{q}\ }
      H^{q}(C^{\bullet}_{\partial})\to\Hcpt{q+1}\to\cdots ,
  \]
  and \textbf{every map in it intertwines
  $(T^{\partial}_{\Ga g\Ga},T^{c}_{\Ga g\Ga},\Tdc{g})$}, by
  Theorem~\ref{thm:conewelldef}(a) --- which needs only that $\Phig$ is a cochain map,
  and that is a theorem here, not a hypothesis.
\item \PROVED{} \textbf{(Collapse to six terms.)} \LEAN{cl:les} (the injectivity of $\rres^{0}_{\ast}$ and the Euler identity, under $\Heq{2}=0$; exactness is Lemma~\ref{lem:conerelative}(2)) $\Hcpt{q}=0$ for $q\le0$ and $q\ge3$
  \emph{unconditionally}, and $H^{q}(C^{\bullet}_{\partial})=0$ for $q\ge2$ because
  $\Xbdry$ is one-dimensional (\ref{B2}). If in addition $\Heq{2}(K,\Rel;L)=0$, the
  sequence of (1) collapses to
  \begin{equation}\label{eq:les}
    0\to X_1\xrightarrow{\ \rres^{0}\ }X_2\to X_3\to X_4
      \xrightarrow{\ \rres^{1}\ }X_5\to X_6\to0 ,
  \end{equation}
  with $\rres^{0}$ the inclusion $\VR^{\Ga}\hookrightarrow\VR^{\Gbdry}$, and over a field
  \begin{equation}\label{eq:dimbook}
    \dim X_3=\dim X_2-\dim X_1+\dim\ker\rres^{1},\qquad
    \dim X_6=\dim X_5-\rank\rres^{1} .
  \end{equation}
  Write $H^{1}_{!}:=\ker\rres^{1}=\im(X_3\to X_4)$.
\item \textbf{(When \eqref{eq:les} is a statement about the six groups
  \eqref{eq:sixgroups}.)} \PROVED{} Assume $6\in\cring^{\times}$ \emph{and $n$ even}.
  Then $\Heq{q}\cong H^{q}(\Ga;\VR)$ (Remark~\ref{rmk:orbifold}),
  $H^{q}(C^{\bullet}_{\partial})\cong H^{q}(\partial;\VR)$
  (Proposition~\ref{prop:circlemodel}(3), which uses $n$ even and
  $2\in\cring^{\times}$), and $\Heq{2}=0$ (Lemma~\ref{lem:h2torsion}(1)); so
  \eqref{eq:les} is exact with the entries \eqref{eq:sixgroups}. \textbf{The hypothesis
  $H^{2}(\Ga;\VR)=0$ is not removable} (Lemma~\ref{lem:h2torsion}(2)); the hypothesis
  $\Hcpt{0}=0$ that [E] listed beside it \textbf{is}, being unconditional by (2). At
  $\ell\in\{2,3\}$ the three identifications fail and \eqref{eq:les} is a
  \emph{window} of a sequence with a non-zero term in every degree:
  Remark~\ref{rmk:ell2window}, and \ref{open:ell2} for the policy. Read on the model
  complexes, (1) and (2) hold at every $\ell$.
\item \CONJ{} \emph{as a citation.} In the presentation of Claim~\ref{cl:amalgam}(3),
  $\rres^{1}$ is induced by
  \begin{equation}\label{eq:r1formula}
    \rres^{1}[m]\;=\;\bigl[\rhon(S)^{-1}\bigl(\rhon(U)-\id\bigr)m\bigr]
    \ \in\ (\VR)_{T},\qquad U=ST .
  \end{equation}
  The map is \PROVED{} well defined --- it kills both generators of the denominator of
  Claim~\ref{cl:amalgam}(3): $(\rhon(U)-\id)(1+\rhon(U)+\rhon(U)^{2})=\rhon(U)^{3}-\id
  =\rhon(-I)-\id=0$ for $n$ even, and
  $\rhon(S)^{-1}(\rhon(U)-\id)(1+\rhon(S))=(\rhon(T)-\id)+(\rhon(T)-\id)\rhon(S)$ lies in
  $\im(\rhon(T)-\id)$, using $\rhon(S)^{-1}\rhon(ST)=\rhon(T)$,
  $\rhon(S)^{-1}=\rhon(S)$ for $n$ even, and $\rhon(TS)-\rhon(S)=(\rhon(T)-\id)\rhon(S)$.
  \textbf{That it is the restriction map} is \emph{not} proved here: its only source is
  the working note \file{hecke\_extensions\_mathematics.md}, which
  Appendix~\ref{app:errata}, entry~\ref{err:e14}, declares not a source of record.
  Nothing below \eqref{eq:r1formula} depends on it: \eqref{eq:dimbook} is proved from
  (2) alone.
\end{enumerate}
\end{claim}

\begin{proof}
(1) is Lemma~\ref{lem:conerelative}(2) and Theorem~\ref{thm:conewelldef}(a) verbatim.
(2) The two vanishings are Lemma~\ref{lem:conerelative}(4) and (5); the collapse and
\eqref{eq:dimbook} are Lemma~\ref{lem:conerelative}(5), whose hypothesis is exactly
$\Heq{2}=0$; $\rres^{0}$ is identified in the proof of
Lemma~\ref{lem:conerelative}(4). (3) collects the three cited identifications. (4) the
displayed computations, using $\rhon$'s multiplicativity
(Proposition~\ref{prop:symconv}(2)) and $\rhon(-I)=\id$.
\end{proof}

\subsection{The boundary characters}

\begin{claim}[Boundary eigenvalues]\label{cl:chi}
Let $\cring$ be a field of characteristic $\ell$, $n$ even, $2\in\cring^{\times}$ --- so
that $X_2$ and $X_5$ are the two groups of Claim~\ref{cl:boundary}; clauses (1) and (2)
hold verbatim for $H^{\bullet}(C^{\bullet}_{\partial})$ at every $\ell$ --- let
$p\neq\ell$ be a prime (\ref{H1}), $g=\mathrm{diag}(1,p)$, and
$\chi_b(p):=p^{b}+p^{\,n+1-b}$ for $0\le b\le\lfloor n/\ell\rfloor$.
\begin{enumerate}
\item \PROVED{} On $X_2=H^{0}(\partial;\VR)=\VR^{T}$ the monomials
  $x^{\,n-b\ell}\mathrm{Nm}^{\,b}$ are eigenvectors of $\Thb_g$ with eigenvalues
  $\chi_b(p)$. This is Proposition~\ref{prop:bdrychar}(1), and the right-coset
  convention gives $\chi_{b+1}$ instead --- Proposition~\ref{prop:bdrychar}(2), a
  \emph{different operator}. For $n<\ell$ both $X_2$ and $X_5$ are $1$-dimensional and
  carry $\chi_0(p)=1+p^{\,n+1}$, the Eisenstein eigenvalue of
  Obligation~\ref{obl:oracle} (Proposition~\ref{prop:bdrychar}(3)).
\item \CONJ{} $X_5=H^{1}(\partial;\VR)=(\VR)_{T}$ carries the \emph{same} multiset
  $\{\chi_b\}_{0\le b\le\lfloor n/\ell\rfloor}$, with the class of $y^{n}$ carrying
  $\chi_0$. Proposition~\ref{prop:bdrychar}(3) proves the case $n<\ell$; the general
  case is not proved here and is what Obligation~\ref{obl:taubdry}(vi) \emph{records}
  rather than gates. \COMP{} It was checked for every even $n\le24$,
  $\ell\in\{2,3,5,7,11,13\}$ and $p\in\{2,3,5,7,11,13,17\}\setminus\{\ell\}$ --- $468$
  cases --- by forming $\rhon$ over $\Fl$, the operator \eqref{eq:bdrygroupform} on
  $(\VR)_{T}$ and $\sum_{\delta}\rhon(\delta)$ on $\VR^{T}$, and factoring both
  characteristic polynomials: in every case both multisets are exactly
  $\{\chi_b\}$ and both polynomials split. \emph{Ad hoc}, not a registered validation.
  \emph{Status note} (v1.12.0-pre.3): the clause has two halves, and they now differ. The
  \emph{semisimplicity} half --- that $\Thb_g$ on $X_5$ is diagonalisable over $\Fl$, the operators for
  $p\neq\ell$ commuting and depending on $p\bmod\ell$ only, with eigenvalue $\chi_{n-k}(p)$ on a torus
  weight vector of weight $a^{n-k}d^{k}$ --- is \textbf{proved}: Lemma~\ref{lem:bdrytorus}(1)--(3) at
  $\cring=\Fl$, and over any field of characteristic $\ell$ by extension of scalars, which
  rests on Proposition~\ref{prop:bdrytranscription} (the operator of \eqref{eq:bdrychain} on $X_5$
  \emph{is} \eqref{eq:bdrygroupform}). The \emph{multiset} half --- which torus weights $(\VR)_{T}$ carries,
  i.e.\ that the eigenvalues are exactly $\{\chi_b\}_{0\le b\le\lfloor n/\ell\rfloor}$ --- does not follow
  and \textbf{stays \CONJ}; the tag above is that half's. The clause's last part, ``with the class of
  $y^{n}$ carrying $\chi_0$'', \emph{is} proved for every even $n$ and $p\neq\ell$: $y^{n}\notin
  \im(\rhon(T)-\id)$ (every element of that image has $y$-degree $<n$, as $(\rhon(T)-\id)u_j=
  \sum_{t<j}\binom{j}{t}u_t$), and $y^{n}$ has torus weight $d^{n}$, so Lemma~\ref{lem:bdrytorus}(3) with $k=n$
  gives $\chi_0(p)$ (v1.12.0-pre.3, referee P3). Later evidence for it, all ad hoc or recorded,
  none gating: $2080$ group-form cases ($\ell=5$, even $n\le200$; $\ell\in\{7,11,13\}$, even $n\le120$;
  slot-12 WP-13) and $814/814$ shipped operators on $X_5$ (slot-12 WP-14,
  \code{HeckeBoundaryFinalityValidation} arm (L), recorded).
\item \PROVED{} \textbf{(Coincidences.)} As functions of $p$ on primes $p\neq\ell$,
  $\chi_a\equiv\chi_{a'}$ if and only if $a\equiv a'$ or $a\equiv n+1-a'$
  $\pmod{\ell-1}$. In particular $\chi_0\equiv\chi_1$ iff $(\ell-1)\mid n$, and
  $\chi_1\equiv\chi_2$ iff $(\ell-1)\mid(n-2)$.
\end{enumerate}
\end{claim}

\begin{proof}
(3) $\chi_a(p)$ depends on $p$ only through $p\bmod\ell$, and $p\mapsto p\bmod\ell$ is
onto $\Fl^{\times}$ as $p$ runs over the primes $\neq\ell$
\STD{\cite[Ch.~VI, \S4]{serre_cours}: Dirichlet's theorem on primes in arithmetic
progressions; the section pointer is not verified against the text (rule~\ref{A2}), and
the statement is used only through this surjectivity, which is why (3) is really a
statement about functions on $\Fl^{\times}$}. So $\chi_a\equiv\chi_{a'}$ on primes iff
$t^{a}+t^{\,n+1-a}=t^{a'}+t^{\,n+1-a'}$ for every $t\in\Fl^{\times}$.

Put $m:=\ell-1$ and fix a generator $g$ of the cyclic group $\Fl^{\times}$. A function
$\Fl^{\times}\to\Fl$ is the vector of its values at $g^{0},\dots,g^{m-1}$, and
$t\mapsto t^{c}$ has value vector $(g^{cs})_{s<m}$, which depends only on $c\bmod m$. The
$m\times m$ matrix $(g^{cs})_{c,s\in\Z/m}$ is a Vandermonde matrix in the $m$
\emph{distinct} elements $g^{0},\dots,g^{m-1}$ of $\Fl$, hence invertible; so the $m$
functions $t\mapsto t^{c}$, $c\in\Z/m$, are linearly independent over $\Fl$. The
displayed equality is a vanishing $\Fl$-linear combination of them with integer
coefficients in $\{0,\pm1,\pm2\}$, all of which are non-zero in $\Fl$ when they are
non-zero in $\Z$, provided $\ell>2$. Hence for $\ell>2$ the multisets $\{a,n+1-a\}$ and
$\{a',n+1-a'\}$ agree modulo $m$, which is the stated disjunction: if $a\equiv a'$ then
$n+1-a\equiv n+1-a'$, and if $a\equiv n+1-a'$ then $n+1-a\equiv a'$. For $\ell=2$,
$m=1$ and both sides of the equivalence hold for all $a,a'$. Specialising $(a,a')=(0,1)$
gives $0\equiv1$ or $0\equiv n$, and $(1,2)$ gives $1\equiv2$ or $1\equiv n-1$; for
$\ell>2$ the first alternative fails in both.
\end{proof}

\begin{remark}[Which Dickson classes fail to separate the two indexings]\label{rmk:chicoincide}
(Appendix~\ref{app:errata}, entry~\ref{err:e21}.) [E] concludes from
Claim~\ref{cl:chi}(3) that ``$\chi_1\equiv\chi_2$ at every Dickson locus, so Dickson
classes cannot discriminate the two indexings''. \textbf{The first half is false.} On the
Dickson locus $n=a(\ell+1)+b(\ell^{2}-\ell)$ of Claim~\ref{cl:dimH1}(2) one has
$n\equiv2a\pmod{\ell-1}$, since $\ell+1\equiv2$ and $\ell^{2}-\ell\equiv0$; so
$\chi_1\equiv\chi_2$ iff $(\ell-1)\mid 2(a-1)$, \textbf{not} always. It does hold at the
\emph{first} Dickson class $n=\ell+1$ ($a=1$, $b=0$), which is [E]'s example
$(n,\ell)=(6,5)$: there $n-2=\ell-1$ and $\chi_1(2)=2+2^{6}\equiv1\equiv\chi_2(2)$.
It fails already at the next one, $(n,\ell)=(12,5)$ ($a=2$): $\chi_1(2)=2+2^{12}\equiv3$
and $\chi_2(2)=2^{2}+2^{11}\equiv2$ in $\F_5$. \COMP{} Checked for
$\ell\in\{5,7,11,13\}$ and every Dickson locus with $n<60$: the criterion
$(\ell-1)\mid(n-2)$ agrees with direct evaluation in every case, and the criterion of
Claim~\ref{cl:chi}(3) was checked exhaustively for $\ell\le19$, $n<40$ and all
$0\le a,a'\le n+1$ --- $190\,560$ cases, $0$ disagreements. \textbf{What survives} is the
conclusion [E] drew it for: the first Dickson class never separates the left- and
right-coset indexings, so the separation must come from the Eisenstein eigenvalue
$1+p^{\,n+1}$ instead --- which is how Proposition~\ref{prop:bdrychar}(3) and
Obligation~\ref{obl:taubdry}(vii) settle the convention.
\end{remark}

\subsection{Poincar\'e duality}

\begin{claim}[Poincar\'e duality]\label{cl:pd}
\OPENSTAT{} \textbf{This is owner Open Question~4 and is not proved in this document};
its coefficient-level input is proved (Proposition~\ref{prop:pairing}) and its
cohomological content is not. Let $\ell\ge5$, $\cring=\Fl$, $n$ even,
$\Mn=\Sym^{n}\Fl^{2}$ and $\Dn$ its contragredient
(Definition~\ref{def:contragredient}).
\begin{enumerate}
\item \textbf{(Perfect pairings.)} The evaluation pairing of
  Proposition~\ref{prop:pairing} induces perfect pairings
  \[
    \Hcpt{t}(K,\Rel;\Mn)\times\Heq{2-t}(K,\Rel;\Dn)\to\Fl,\qquad
    H^{t}(\partial;\Mn)\times H^{1-t}(\partial;\Dn)\to\Fl ,
  \]
  equivalently, under $6\in\cring^{\times}$ (Remark~\ref{rmk:orbifold},
  Lemma~\ref{lem:conerelative}(3)), $\Hcpt{t}(\Ga;\Mn)\times H^{2-t}(\Ga;\Dn)\to\Fl$ ---
  which is how [E] and Claim~\ref{cl:les} name the same groups.
\item \textbf{(Self-adjointness --- in the adjointness form, and only there.)} Under
  (1), for every $g\in\De$ the operator $\Tdc{g}$ on the $\Mn$-side is adjoint to
  $T_{\Ga(\adjg g)\Ga}$ on the $\Dn$-side:
  \[
    \bigl\langle \Tdc{g}\,x,\ y\bigr\rangle
    =\bigl\langle x,\ T_{\Ga(\adjg g)\Ga}\,y\bigr\rangle ,
  \]
  \textbf{with no determinant factor}; and for $g=\mathrm{diag}(1,N)$ one has
  $\Ga(\adjg g)\Ga=\Ga g\Ga$ (Proposition~\ref{prop:pairing}(5)), so $T_p$, $T_{p^{\nu}}$
  and $T_N$ are \emph{self}-adjoint. This is the form in which
  Proposition~\ref{prop:coneroutes}(3) reads this claim.
\item \textbf{(Consequences.)} For a maximal ideal $\midl$ of the Hecke algebra
  $\Talg$, $\Hcpt{t}(\Mn)_{\midl}$ and $H^{2-t}(\Dn)_{\midl}$ have the same
  characteristic polynomials, and the \emph{Loewy} layers of one are the \emph{socle}
  layers of the other in reverse order --- $\rad^{i}(V^{\vee})=(\soc^{i}V)^{\perp}$ for a
  module $V$ over the local ring $\Talg_{\midl}$. The two series coincide only for rigid
  modules. The classification consequence drawn from this is \Espec{pr:dual} and stays
  there.
\end{enumerate}
\end{claim}

\begin{remark}[What Claim~\ref{cl:pd} rests on, and the $p^{\,n}$ twist it must not inherit]\label{rmk:pdscope}
(Appendix~\ref{app:errata}, entries~\ref{err:e14} and \ref{err:e18}.)
\begin{enumerate}
\item \textbf{The adjointness form is twist-free; the diagonal one is not.}
  Proposition~\ref{prop:pairing} proves two different identities. (3),
  $\langle\rhon(g)v,\varphi\rangle=\langle v,\rhodual(\adjg g)\varphi\rangle$, is exact,
  integral and carries \emph{no} determinant factor. (4),
  $\langle\rhon(g)v,\rhodual(g)\varphi\rangle=(\det g)^{\,n}\langle v,\varphi\rangle$,
  carries the twist $(\det g)^{n}=p^{\,n}$ for $g=\mathrm{diag}(1,p)$. [E]'s wording
  ``$T_q$ self-adjoint once $\Dn$ carries the adjugate $\De$-structure
  $g\mapsto\det(g)g^{-1}$'' is true of the \emph{adjointness} form and false of the
  diagonal one, and Claim~\ref{cl:pd}(2) is stated in the former.
  Note $\det(g)g^{-1}=\adjg g$, so the structure [E] names is
  Definition~\ref{def:contragredient}'s.
\item \textbf{What it depends on.} (i) Proposition~\ref{prop:pairing}(2),(3),(5) --- the
  perfectness of the coefficient pairing over every $\cring$, the adjointness identity,
  and the self-adjugacy of the double coset; (ii) the construction of $\Hcpt{\bullet}$ as
  relative cohomology, Lemma~\ref{lem:conerelative}(3), and of its Hecke action,
  Definition~\ref{def:coneaction}; (iii) the boundary cohomology
  Definition~\ref{def:bdrygroup} and its operator Definition~\ref{def:bdryop}; (iv) the
  hypothesis $6\in\cring^{\times}$, through Remark~\ref{rmk:orbifold} and
  Lemma~\ref{lem:h2torsion}(1). \textbf{Not} yet available: a cochain-level pairing on
  the stored bases. Until it exists, (1) and (2) are assertions with no certificate, and
  Proposition~\ref{prop:coneroutes}(3) --- the acceptance gate for $X_3$ --- inherits
  their \OPENSTAT{} status.
\item \textbf{Socle, not Loewy.} Entry~\ref{err:e14} of Appendix~\ref{app:errata} records
  the working note's statement of (3) for the \emph{Loewy} series. It is the socle series
  that reverses; Claim~\ref{cl:pd}(3) says so.
\item \textbf{$\Mn$ is not self-dual for $n\ge\ell$}, so (1) genuinely relates two
  different local systems; the integral comparison map $\Phi_n$ of
  Lemma~\ref{lem:symdual} is an isomorphism exactly when $n!\in\cring^{\times}$, i.e.\
  $n\le\ell-1$ (Remark~\ref{rmk:divided}). This is why the pipeline, which computes
  $\Mn$ only, cannot check (1) without building $\Dn$ --- phase E7 of the extension
  programme, made unconditional by owner decision OQ4.
\end{enumerate}
\end{remark}

\paragraph{Code map.}
\textbf{Two of the six groups are computed with their Hecke action, two without, two not
at all.} $X_1=\Heq{0}$ and $X_4=\Heq{1}$, their bases, projectors and every matrix
$\Amat{q}(T_p)$, $q\in\{0,1,2\}$: \code{EquivariantCohomologyOnK} and
\code{HeckeOperatorOnK}, with $\Amat{2}$ emitted as the $0\times0$ matrix
(\S\ref{sec:degreewise}, Lemma~\ref{lem:h2torsion}(1)). $X_2$, $X_5$ and the restriction
maps $\rres^{0},\rres^{1}$: the circle model of
\code{LocalSystemPresets::CreateS1Generic} and the morphism
\code{TestSL2Z\_To\_S1Cusp\_RestrictionMorphism}
(\file{LocalSystemSL2Z.cpp}:843--922), pinned by
\code{HeckeMorphismRegressionValidation}; \textbf{no boundary Hecke operator exists}
(census of 2026-09-16, \S\ref{sec:ell2} code map). $X_3$, $X_6$: \textbf{not computable};
the mapping cone exists only for the non-equivariant \code{CochainComplex}
(\code{ConeComplex()}, \file{SimplicialCohomology.h}), not for
\code{EquivariantCochainComplexMorphismOnK}. Claim~\ref{cl:dimH1} is corroborated by
\file{docs/matrix-engine/results/results\_table.tsv} as recorded above;
Claim~\ref{cl:amalgam}(3) is the presentation \code{BarCocycleH1} implements;
Claim~\ref{cl:pd} has \textbf{no implementation}, $\rhodual$ and the pairing not being
built (\S\ref{sec:contragredient} code map).

\paragraph{Formalization note.}
Claim~\ref{cl:amalgam}(2) is the cheapest target in this chapter and the most reusable:
it is $\Ext^{\bullet}_{\Z[\Ga]}(-,\VR)$ applied to one short exact sequence, and
\code{mathlib4} has both the long exact sequence of $\Ext$ and Shapiro's lemma in the
form $\Ext^{\bullet}_{R}(\mathrm{Ind}\,M,N)$. Its input \eqref{eq:treeses} is the one
genuinely external object --- the Bass--Serre tree --- and a formalisation may take
\eqref{eq:treeses} itself as the definition of the amalgam datum rather than derive it.
Claim~\ref{cl:amalgam}(3),(4), Lemma~\ref{lem:h2torsion} and Claim~\ref{cl:boundary} are
then finite consequences, the first needing only $\abs{\Ga_0}H^{q}(\Ga_0;-)=0$ for a
finite subgroup $\Ga_0$. Claim~\ref{cl:chi}(3)
is character theory of a cyclic group and is fully elementary; only its appeal to
Dirichlet is external, and that can be avoided by restating the claim over
$\Fl^{\times}$, which is what the proof actually uses. Claim~\ref{cl:dimH1} and
Claim~\ref{cl:pd} are \emph{not} formalisation targets at this stage: the first imports
Eichler--Shimura and Dickson, the second is open.

%----------------------------------------------------------------------
\section{The general orbifold model, and the torsion primes}\label{sec:orbifoldmodel}
%----------------------------------------------------------------------
% SOURCE:  NEW.  Implementation plan C-A3(ii.7); OQ2 verdict Sec. 5; owner
%          decisions OP4 ("definitions now, construction deferred") and OD-T1
%          ("a refinement of the current triangulation").
% NOTATION: crosswalk rows cring, q (degree), Rel; (H4) reused verbatim.
% ERRATA:  none.  This chapter asserts nothing about a superseded document.

\paragraph{What this chapter fixes.}
\textbf{Definitions only; the construction is deferred} --- owner decision OP4,
2026-09-15. What is fixed here is \emph{which object} the words ``the orbifold cohomology
of the pair'' name when $6\notin\cring^{\times}$, and in what shape its Hecke datum is to
be written, so that the deferred construction cannot later be taken to define something
else. Three things are stated: the Borel construction of the pair and its free
$\Ga$-model (Definition~\ref{def:borelmodel}); the sense in which the $7$-vertex complex
of Definition~\ref{def:model7} is \emph{a} specialisation of it at $6\in\cring^{\times}$
(Theorem~\ref{thm:specialisation}); and what Moerdijk--Pronk's simplicial complex is and
is not cited for (Remark~\ref{rmk:mpcaveat}). What is \emph{not} here: any construction
of a free model, of its $\tau$-tables, or of the $\ell=2$ truncation --- that is one
piece of work with \ref{open:ell2}, and it is deferred with it. (Since v1.8.0 the
construction, on the unrefined $(\Xbulk,\Xbdry)$ and as existence, is
\S\ref{sec:integral}.)

\paragraph{Two standing instructions on how these definitions are phrased.}
\begin{enumerate}
\item \textbf{Degree-generic.} Every definition below is stated in a degree range
  $\{0,\dots,q_{\max}\}$ with $q_{\max}$ a parameter, never in $\{0,1,2\}$. This is
  \ref{open:ell2}'s site~(4) (\S\ref{sec:oc1sites}), and it is what makes the truncation
  policy expressible at all: at $\ell\in\{2,3\}$ the complexes below are unbounded above
  (Remark~\ref{rmk:ell2window}(3)).
\item \textbf{Refinement-parametric} (owner direction OD-T1, 2026-09-15, recorded at
  Remark~\ref{rmk:orbifold}). The input is an \emph{arbitrary} $\Ga$-pair satisfying
  \ref{G1}--\ref{G4} of Definition~\ref{def:borelmodel} --- which are \ref{H4}(i),(ii)
  and \ref{B1},\ref{B2} over an arbitrary input, \emph{without} \ref{H4}(iv)'s
  stabiliser inventory --- not the $7$-vertex one; a $\Ga$-equivariant refinement of the
  $7$-vertex complex is then admissible input rather than an exception. Nothing below
  names a vertex, an edge or a triangle of Definition~\ref{def:model7}, and nothing below
  may be read as ``the $7$-vertex complex plus corrections''.
  Remark~\ref{rmk:orbifoldrefine} says what a refinement would and would not change.
\end{enumerate}

\begin{definition}[The Borel construction of the pair, and its twisted chain data]\label{def:borelmodel}
Let $\cring$ be a commutative ring and $\VR$ a $\Z[\De]$-module, free and finitely
generated over $\Z$ (Definition~\ref{def:rhon}: $\VR=\Sym^{n}(\Z^{2})\otimes\cring$).
Let $(Y,\partial Y)$ be a pair of simplicial $\Ga$-complexes satisfying:
\begin{enumerate}[label=(G\arabic*)]
\item\label{G1} $\Ga$ acts simplicially on $Y$ without inversions, with finite
  stabilisers, and $\partial Y\subseteq Y$ is a $\Ga$-subcomplex;
\item\label{G2} $Y$ is contractible and finite-dimensional;
\item\label{G3} $Y$ has finitely many $\Ga$-orbits of simplices in each degree, and so
  does $\partial Y$;
\item\label{G4} $\partial Y$ is a disjoint union of contractible components permuted
  transitively by $\Ga$, the stabiliser of a distinguished component
  $\partial Y_{\infty}$ being $\Gbdry$.
\end{enumerate}
(\ref{G1}--\ref{G4} are \emph{rule tags}, in the style of \ref{H1}--\ref{H4} and
\ref{B1}--\ref{B3}; they do not reintroduce the letter $G$ of \S\ref{sec:framework} as a
symbol, and no statement of this chapter uses that letter.)
\ref{G1} and \ref{G2} are \ref{H4}(i),(ii), \ref{G4} is \ref{B1},\ref{B2}, and \ref{G3}
is what the finiteness of $K$ supplies in \ref{H4}(iii) --- all four with the $7$-vertex
complex replaced by an arbitrary input. \ref{H4}(iv)'s \emph{inventory} of stabilisers is
deliberately \textbf{not} assumed, only their finiteness; \ref{H4}(iii) itself, the
existence of a strict fundamental domain, is assumed only where
Theorem~\ref{thm:specialisation} needs it. Fix a bound $q_{\max}\ge0$
(\ref{open:ell2}) and a free resolution $F_{\bullet}\to\Z$ of the trivial module over
$\Z[\Ga]$.
\begin{enumerate}
\item \textbf{(The free $\Ga$-model.)} Put
  \begin{equation}\label{eq:borelfree}
    \widehat{C}_{\bullet}(Y):=F_{\bullet}\otimes_{\Z}C_{\bullet}(Y),\qquad
    \widehat{C}_{\bullet}(Y,\partial Y):=F_{\bullet}\otimes_{\Z}
      \bigl(C_{\bullet}(Y)/C_{\bullet}(\partial Y)\bigr),
  \end{equation}
  each with the diagonal $\Ga$-action and the total differential. Both are complexes of
  \textbf{free} $\Z[\Ga]$-modules, because $F_{\bullet}$ is free and the simplicial
  chain modules are free over $\Z$; and $\widehat{C}_{\bullet}(Y)\to\Z$ is a free
  resolution, because $Y$ is contractible (\ref{G2}). This is the simplicial-chain
  incarnation of the Borel construction $E\Ga\times_{\Ga}Y$.
\item \textbf{(The orbifold cochain complexes and their cohomology.)} For
  $0\le q\le q_{\max}+1$ set
  \[
    \widehat{C}^{\,q}(Y;\VR):=\Hom_{\Z[\Ga]}\bigl(\widehat{C}_{q}(Y),\VR\bigr),
  \]
  and likewise for the pair and for $\partial Y$, with the cone
  $\Cone{\rres}$ of the restriction $\rres:=j^{\ast}$ formed exactly as in
  Definition~\ref{def:pair}. The \textbf{orbifold cohomology of the pair} is the
  cohomology of these complexes, reported in degrees $\le q_{\max}$:
  \[
    H^{q}_{\Ga}(Y;\VR),\qquad H^{q}_{\Ga}(\partial Y;\VR),\qquad
    H^{q}_{\Ga}(Y,\partial Y;\VR) .
  \]
  \PROVED{} By (1), $\widehat{C}_{\bullet}(Y)\to\Z$ is a free resolution over $\Z[\Ga]$,
  so $H^{q}_{\Ga}(Y;\VR)=H^{q}(\Ga;\VR)$ by the definition of group cohomology. On the
  boundary, \ref{G4} gives
  $C_{\bullet}(\partial Y)\cong\Z[\Ga]\otimes_{\Z[\Gbdry]}C_{\bullet}(\partial Y_{\infty})$,
  hence $\widehat{C}_{\bullet}(\partial Y)\cong
  \Z[\Ga]\otimes_{\Z[\Gbdry]}\bigl(F_{\bullet}\otimes_{\Z}C_{\bullet}(\partial
  Y_{\infty})\bigr)$ by the projection formula, and the right-hand factor is a free
  $\Z[\Gbdry]$-resolution of $\Z$ ($F_{\bullet}$ restricted to $\Gbdry$ is still free,
  $\partial Y_{\infty}$ is contractible); so by the adjunction of
  Proposition~\ref{prop:circlemodel}(1),
  $H^{q}_{\Ga}(\partial Y;\VR)=H^{q}(\Gbdry;\VR)=H^{q}(\partial;\VR)$ of
  Definition~\ref{def:bdrygroup}. Both hold for every $\cring$, every $n$ and every
  reported degree $q\le q_{\max}$ --- with \emph{no} invertibility hypothesis, which is
  the point of the construction. (Degree $q_{\max}+1$ is scaffolding and is not reported:
  \ref{open:ell2}'s truncation lemma.)
\item \textbf{(Twisted chain data on the free model.)} A \textbf{Hecke datum of the pair
  on $(Y,\partial Y)$} is a triple $(\tau,\taub,\etapair)$ of additive maps
  \[
    \tau\colon\widehat{C}_{\bullet}(Y)\to\widehat{C}_{\bullet}(Y),\quad
    \taub\colon\widehat{C}_{\bullet}(\partial Y)\to\widehat{C}_{\bullet}(\partial Y),\quad
    \etapair\colon\widehat{C}_{\bullet}(\partial Y)\to\widehat{C}_{\bullet+1}(Y),
  \]
  defined in every degree $\le q_{\max}+1$, satisfying (T1)--(T3) of
  Definition~\ref{def:tau} and (T1$\partial$)--(T3$\partial$) of
  Definition~\ref{def:taubdry} verbatim --- with $\widehat{C}_{\bullet}$ in place of
  $C_{\bullet}(\Ku)$ and $C_{\bullet}(\Xbdry)$ --- together with the pair identity
  $\partial\etapair+\etapair\partial=j\taub-\tau j$ of
  Definition~\ref{def:pairhomotopy}. The cochain maps $\Th_g$, $\Thb_g$, the derived
  homotopy $\hder$ and the cone map $\Phig$ are their coset transports
  (Definition~\ref{def:cosettransport}, Definition~\ref{def:coneaction}), unchanged.
\end{enumerate}
\end{definition}

\begin{remark}[Why the free model buys what it buys]\label{rmk:borelfree}
\PROVED{} On $\widehat{C}_{\bullet}$ the source of every comparison argument is
\emph{free} over $\Z[\Ga]$ and the target is acyclic, so the fundamental lemma of
homological algebra applies with no relative projectivity and no invertibility
hypothesis \STD{\cite[Thm.~2.2.6]{weibel}; \cite[Ch.~I.7.4]{brown}}. Consequently, on a
model satisfying \ref{G1}--\ref{G4}, \emph{both} clauses of
Lemma~\ref{lem:tauexists} hold over $\Z$ for every $\ell$ --- existence, which
Lemma~\ref{lem:tauexists}(1) already gives on the $7$-vertex model, \emph{and} homotopy
uniqueness, which on that model needs $6\in\cring^{\times}$
(Lemma~\ref{lem:tauexists}(2), Remark~\ref{rmk:tauexistsscope}). The same applies to
$\taub$ (Lemma~\ref{lem:taubdryexists}, which is already unconditional) and to
$\etapair$ (Lemma~\ref{lem:etaexists}, Lemma~\ref{lem:etaunique}). Hence
Theorem~\ref{thm:conewelldef}(c) --- independence of $\tau$, $\taub$, $s_j$ and $g$, and
the identification with $\cor\circ(\conj_g)_{\ast}\circ\res$ --- holds on this model
\textbf{for every $\ell$}, where on the $7$-vertex model it is asserted only for
$6\in\cring^{\times}$ (Theorem~\ref{thm:conewelldef}(e)). \textbf{That is the entire
reason the free model is wanted}, and it is a statement about the model, not about a
construction of one: no such model is built in this document.
\end{remark}

\begin{theorem}[The $7$-vertex complex is \emph{a} specialisation at $6\in\cring^{\times}$]\label{thm:specialisation}
Let $(Y,\partial Y)$ satisfy \ref{G1}--\ref{G4}, let $\cring$ be a commutative ring with
$6\in\cring^{\times}$ --- more precisely, with $\abs{\Ga_\sigma}\in\cring^{\times}$ for
every simplex $\sigma$ of $Y$ and of $\partial Y$ --- let $q_{\max}\ge0$, and let
$(K_Y,\Rel_Y)$ present the quotient of $Y$ in the sense of \S\ref{sec:framework}
(Assumption~\ref{ass:fund_domain}). Then for every $q$ with $0\le q\le q_{\max}$:
\begin{enumerate}
\item \PROVED{} $H^{q}_{\Ga}(Y;\VR)\cong\Heq{q}(K_Y,\Rel_Y;L)$ and
  $H^{q}_{\Ga}(\partial Y;\VR)\cong
   H^{q}\bigl(\Hom_{\Z[\Ga]}(C_{\bullet}(\partial Y),\VR)\bigr)$ --- the latter being
  $H^{q}(C^{\bullet}_{\partial})$ when $\partial Y=\Xbdry$ --- naturally; i.e.\ the
  equivariant cochain complex of \S\ref{sec:explicitcochaincomplex} is the $t=0$ row of
  the isotropy spectral sequence of the free model, and the rows $t>0$ vanish.
\item \PROVED{} $H^{q}_{\Ga}(Y,\partial Y;\VR)\cong\Hcpt{q}$ of
  Lemma~\ref{lem:conerelative}(3), compatibly with (1) and with the long exact sequence
  \eqref{eq:coneles}.
\item \PROVED{} The two Hecke actions agree: the operator induced on
  $H^{\bullet}_{\Ga}(Y,\partial Y;\VR)$ by a Hecke datum of the pair on $(Y,\partial Y)$
  (Definition~\ref{def:borelmodel}(3)) corresponds under (2) to
  $T^{c}_{\Ga g\Ga}$ of Definition~\ref{def:coneaction}, and likewise in the bulk and on
  the boundary. \emph{Both} are $\cor\circ(\conj_g)_{\ast}\circ\res$
  (Theorem~\ref{thm:classical}, Theorem~\ref{thm:conewelldef}(c),
  Remark~\ref{rmk:borelfree}), which determines them.
\item Taking $(Y,\partial Y)=(\Xbulk,\Xbdry)$ and $(K_Y,\Rel_Y)=(K,\Rel)$ of
  Definitions~\ref{def:model7} and \ref{def:model7rel}, which satisfies \ref{G1}--\ref{G4}
  by Proposition~\ref{prop:model7h4} and \ref{B1}--\ref{B2}, recovers the whole shipped
  pipeline. \textbf{It is one admissible input among others}: any $\Ga$-equivariant
  refinement of $(K,\Rel)$ is another, with the \emph{same} left-hand sides in (1)--(3)
  and different tables (Remark~\ref{rmk:orbifoldrefine}).
\end{enumerate}
\textbf{Not asserted, at any $\ell$}: that a model satisfying \ref{G1}--\ref{G4} with
$\partial Y$ as in \ref{G4} exists for a \emph{refinement} of $(K,\Rel)$, or that one can
be built effectively. That is the deferred construction. \emph{Since v1.8.0} a
finite-rank free model on $(\Xbulk,\Xbdry)$ itself --- which (4) already admits with
$F\otimes C_{\bullet}(X)$ --- is constructed, with its comparison map and a Hecke datum of
the pair on it, as existence (Definition~\ref{def:isomodel},
Proposition~\ref{prop:isomodel}, Proposition~\ref{prop:strictlift}); for a refinement
nothing is asserted and nothing is needed (Remark~\ref{rmk:cafive}). (Worded so at v1.8.1:
v1.8.0's sentence read as if the unrefined case had been deferred.)
\end{theorem}

\begin{proof}
(1) $\widehat{C}_{\bullet}(Y)=F_{\bullet}\otimes C_{\bullet}(Y)$ is a double complex of
$\Z[\Ga]$-modules; its second spectral sequence is the isotropy spectral sequence, with
$E_1^{s,t}=\bigoplus_{\sigma\in\Sigma_s}H^{t}(\Ga_{\sigma};\VR)$, $\Sigma_s$ a set of
representatives of the $\Ga$-orbits of $s$-simplices (finite by \ref{G3}), converging to
$H^{s+t}_{\Ga}(Y;\VR)$ \STD{\cite[Ch.~VII, \S7, (7.10)]{brown}, with the same
rule~\ref{A2} caveat on the pointer as in Remark~\ref{rmk:orbifold}, which is the same
citation}. Each $H^{t}(\Ga_{\sigma};\VR)$ with $t\ge1$ is annihilated by
$\abs{\Ga_{\sigma}}$ \STD{\cite[Ch.~III.10.2]{brown}}, which is a unit, so those rows
vanish; the surviving row $t=0$ is
$s\mapsto\bigoplus_{\sigma\in\Sigma_s}\VR^{\Ga_{\sigma}}$, which is
$\Hom_{\Z[\Ga]}(C_{\bullet}(Y),\VR)$, and that is $\Ceq(K_Y,\Rel_Y;L)$ by
Theorem~\ref{thm:explicit_complex_equiv}. The boundary statement is the same argument on
$\partial Y$. This is Remark~\ref{rmk:orbifold} with $(K,\Rel)$ replaced by
$(K_Y,\Rel_Y)$; no step of it used the $7$-vertex complex.

(2) The relative version of (1) gives
$H^{q}_{\Ga}(Y,\partial Y;\VR)\cong H^{q}(\Hom_{\Z[\Ga]}(C_{\bullet}(Y,\partial Y),\VR))$,
and Lemma~\ref{lem:conerelative}(3) identifies the right-hand side with $\Hcpt{q}$; the
compatibility with \eqref{eq:coneles} is the naturality of both identifications in the
short exact sequence of Lemma~\ref{lem:conerelative}(2), whose three terms are the images
of the three complexes of (1) and (2).

(3) Both operators are $\cor\circ(\conj_g)_{\ast}\circ\res$ on the same group cohomology:
for the free model by Remark~\ref{rmk:borelfree}, for the $7$-vertex model by
Theorem~\ref{thm:conewelldef}(c), whose hypothesis $6\in\cring^{\times}$ is in force. A
map determined by a formula on both sides agrees on both sides.

(4) Proposition~\ref{prop:model7h4} gives \ref{H4}(i)--(iv), hence \ref{G1} and
\ref{G2}, and \ref{G3} because $K$ is finite; \ref{G4} is \ref{B1}--\ref{B2}.
\end{proof}

\begin{remark}[Moerdijk--Pronk as a computational model, and the effectivity caveat]\label{rmk:mpcaveat}
Moerdijk and Pronk \cite{moerdijk_pronk} give a simplicial set $S(\mathcal{T})$ attached
to a triangulated orbifold and prove $H^{\bullet}(\mathcal{M},\mathcal{A})\cong
H^{\bullet}(S,A)$. Line numbers below are of the \LaTeX{} source
\file{simplicial.tex} in \file{docs/literature/orbifolds/}, and every quotation in this
remark \textbf{was checked against that source on 2026-09-16}. \emph{Inside the
quotation marks the symbols are theirs, not this document's} --- in particular their
$G$ is a chart's isotropy group, their $X$ is the space acted on and their $S_k$ is the
degree-$k$ level of a simplicial set, where \S\ref{sec:framework} retired $G$ after
Part~I, $\Xbulk$ is the tessellation and $S=\smat{0}{-1}{1}{0}$.
\begin{enumerate}
\item \textbf{What they are cited for.} (i) The identification of the orbifold cohomology
  of a global quotient with the cohomology of the Borel construction, in their own words:
  ``If the orbifold is given as the orbit space $X/G$ of a group action as above, this
  cohomology is the equivariant sheaf cohomology of the group action. It agrees with the
  (ordinary) cohomology of the Borel construction $EG\times_{G}X$'' (:72--75). That
  sentence is the warrant for Definition~\ref{def:borelmodel} being \emph{the} definition
  and not one option among several. (ii) $S(\mathcal{T})$ as \emph{one computational
  model} of it, for an \emph{adapted} triangulation --- one in which every top simplex
  lies in a chart and every simplex has a vertex of maximal isotropy (:520--528) --- with
  $S_0$ the top simplices and, for $k\ge1$,
  $S_k=\sum_{\sigma_0\cap\cdots\cap\sigma_k\neq\emptyset}
   G^{k}_{v(\sigma_0\cap\cdots\cap\sigma_k)}$ (:589--592), and the comparison theorem at
  :678--684 (Theorem~2.1 of the published version, a \emph{number} this document has not
  checked against the journal text).
\item \textbf{What they are \emph{not} cited for: the normative definition.} Their
  charts are effective by construction --- a chart is ``a finite group $G$ of
  $C^\infty$-automorphisms of $\tilde U$'' (:125), i.e.\ a subgroup of
  $\mathrm{Aut}(\tilde U)$ --- and their construction of the face maps $\mu$ rests on
  ``there exists a unique $h\in H$ such that $\mu=h\circ\lambda$'' (:148), which is
  exactly what fails when a subgroup acts trivially. \textbf{$[\SL_2(\Z)\backslash\Hyp]$
  is not effective}: $-I$ acts trivially on $\Hyp$ (\ref{H4}(iv),
  Remark~\ref{rmk:negI}), and the effective quotient $[\mathrm{PSL}_2(\Z)\backslash\Hyp]$ has
  isotropy $\Z/2,\Z/3$ where this document needs $\Z/4,\Z/6$ --- the distinction that
  Remark~\ref{rmk:negI} shows is not cosmetic. So \emph{as written} their construction
  computes the $\mathrm{PSL}_2(\Z)$-orbifold. That it nevertheless ``does the job in general'' is
  plausible --- their proof (:1063--1440) runs through the translation groupoid and never
  uses the uniqueness of $h$ --- but it is \CONJ{} \emph{as a citation}, and nothing in
  this document rests on it.
\item \textbf{A second reason not to make it normative.} They state the independence of
  $\mu$ from its choices twice without proof: ``It can be shown that this definition is
  independent of the choices made above'' (:772) and ``It can be shown that this
  definition of $\mu$ is again independent of the various choices. \ldots\ We omit the
  details'' (:811--815). A normative definition of a \emph{contract} may not rest on an
  omitted proof.
\item \textbf{The verdict.} Definition~\ref{def:borelmodel} is normative;
  $S(\mathcal{T})$, and equally the isotropy-resolution model (one periodic resolution
  per elliptic orbit, assembled along faces), are computational models of it, to be
  compared with it when either is built. Both are infinite-dimensional in the degree
  direction, which is \ref{open:ell2}'s subject. (Since v1.8.0 the isotropy-resolution
  model is Definition~\ref{def:isomodel}, compared with Definition~\ref{def:borelmodel} by
  Proposition~\ref{prop:isomodel}(5); it is specified, not built. Since v1.9.0 it is built,
  on the unrefined tessellation, by hecke-extensions slot-9 WP-10, \S\ref{sec:codemap}; since
  v1.10.0 with its lifted Hecke datum for $T_2$, $T_3$, $T_5$, Remark~\ref{rmk:strictliftbuilt}.)
\end{enumerate}
\emph{Addendum} (v1.22.0-pre.5, \S\ref{sec:mpmodel}). Items (2) and (3) are \emph{resolved}, not withdrawn. The normative object stays
Definition~\ref{def:borelmodel}. The Moerdijk--Pronk-type set is now defined choice-free (Definition~\ref{def:mpset}: no charts, no $\mu$-maps, no effectivity, no
adaptedness), and its comparison theorem (Theorem~\ref{thm:mpcomparison}) is proved in this document, independently of their argument, for the non-effective
$[\SL_2(\Z)\backslash\Xbulk]$, in every degree and over every ring. It is their own $S(\mathcal T)$ exactly when the action is effective, the triangulation adapted and the
quotient simplicial (Remark~\ref{rmk:mpwhat}(3)); for $\SL_2(\Z)$ it is a non-effective generalisation of it. The \CONJ{} of item (2) concerns their own construction and
stays as stated; nothing rests on it.
\end{remark}

\begin{remark}[What a refinement of the $7$-vertex complex would and would not change]\label{rmk:orbifoldrefine}
Owner direction OD-T1 (Remark~\ref{rmk:orbifold}) expects the correct treatment at
$\ell\in\{2,3\}$ and over $\Z$ to need a refinement --- a subdivision --- of
Definition~\ref{def:model7}'s complex. This remark records the consequence for what is
already written, so that the deferred construction is not mistaken for a revision.
\begin{enumerate}
\item \textbf{Survives verbatim in content.} A $\Ga$-equivariant subdivision of a
  contractible $2$-dimensional $\Ga$-complex is again contractible and $2$-dimensional,
  and it preserves the stabiliser of every point, hence the stabiliser inventory. So
  Proposition~\ref{prop:model7h4}(i)--(iv) holds of the refinement, as do
  \ref{G1}--\ref{G4}, and Theorem~\ref{thm:specialisation} applies to it unchanged.
  Definition~\ref{def:cosetreps} is untouched: it is a statement about $\mathbb{P}^{1}(\Z/N)$
  with no reference to any triangulation.
\item \textbf{Must be recomputed, per refinement.} Everything indexed by the simplices:
  the counts $N_0,N_1,N_2$ and the Euler characteristic of Definition~\ref{def:model7};
  the relation table of Definition~\ref{def:model7rel} and its size; the stored simplex
  order, in which every matrix of this document is written; Algorithm~\ref{alg:tau}'s
  orbit representatives and the bounds of Lemma~\ref{lem:bounds}; and hence the golden
  references of the implementation.
\item \textbf{Is not a revision of anything.} OD-T1 is the owner's stated expectation, a
  design direction and not a theorem (Remark~\ref{rmk:orbifold}, last paragraph); no
  statement of this document depends on it, and none is weakened by it. Its whole force
  is that \ref{G1}--\ref{G4} and Theorem~\ref{thm:specialisation} are phrased over an
  arbitrary input, which is the instruction discharged above.
\item \textbf{What a refinement does \emph{not} by itself buy.} It does not make the
  complexes of Definition~\ref{def:borelmodel} finite: those are unbounded above in
  degree at $\ell\in\{2,3\}$ for \emph{every} triangulation, the elliptic and central
  stabilisers contributing in every degree (Remark~\ref{rmk:ell2window}(2),(3),
  Lemma~\ref{lem:ell2unbounded}, Lemma~\ref{lem:h2torsion}(2)). A truncation policy is
  therefore needed whatever the triangulation, which is why \ref{open:ell2} is stated
  over ``the boundary triangulation in force'' and not over $v_0v_1v_2v_3$.
\end{enumerate}
(v1.8.0: on these grounds OD-T1 is answered ``not needed'' and the free model is built on
the unrefined complex; Remark~\ref{rmk:cafive}.)
(v1.22.0-pre.5, owner ruling Q2 of 2026-10-09: for Moerdijk--Pronk's rule~(ii) on the actual isotropy the $7$-vertex complex is \emph{not} adapted, at every $\ell$ and
over $\Z$, and the refinement $\mathsf R_{\Z}$ --- one subdivision point on $v_4v_6$ and its $S$-image --- is the smallest listed adapted one, the same at $\ell=2$, $\ell=3$ and
over $\Z$ (Definition~\ref{def:mpRZ}, Table~\ref{tab:mprules}, Remark~\ref{rmk:mprulesread}); so OD-T1's expectation is borne out \emph{for their construction}, while every
model proved in \S\ref{sec:mpmodel} already holds on the unrefined complex, which stays the economical model. What distinguishes $\ell=2$, $\ell=3$ and $\Z$ is the
isotropy datum and the coefficient system (Definition~\ref{def:mpell}), not the triangulation.)
\end{remark}

\paragraph{Code map.}
\textbf{Nothing in this chapter is implemented, and that is the point of OP4: the
definitions exist so that what is built later cannot silently be a different object.}
Two shape requirements reach the code and are recorded at
\S\ref{sec:ell2}'s code map as well: the Hecke datum must be typed as tables of twisted
chains on an \emph{abstract} $\Ga$-pair --- orbit representatives, stabilisers,
transporters and the (T1) bookkeeping --- rather than on the seven vertices, and it must
be \textbf{degree-generic}, not the $q\in\{0,1,2\}$ that \code{hecke::HeckeTauTable}
hard-codes today, because \ref{open:ell2} indexes it by $q_{\max}$. Both are properties
of a data model, checkable by inspection, and neither requires the deferred construction.

\paragraph{Formalization note.}
Definition~\ref{def:borelmodel}(1) is the one statement here a proof assistant would find
easy and useful: $F\otimes_{\Z}M$ with the diagonal action is free over $\Z[\Ga]$ when
$F$ is free and $M$ is $\Z$-free, and \code{mathlib4} has the ingredients. It is worth
doing early precisely because Remark~\ref{rmk:borelfree} --- the removal of every
invertibility hypothesis --- is a one-line consequence of it plus the fundamental lemma,
which \code{mathlib4} also has. Theorem~\ref{thm:specialisation}(1) needs the isotropy
spectral sequence, which it does not, and should be deferred; its statement is however
the right interface, since (2) and (3) are formal consequences of (1) together with
Lemma~\ref{lem:conerelative}(3) and Theorem~\ref{thm:conewelldef}(c).
Remark~\ref{rmk:mpcaveat} is bibliography, not mathematics, and carries no obligation.

%======================================================================
\part*{Part V --- Realisation, bounds, and certification}
\addcontentsline{toc}{section}{Part V --- Realisation, bounds, and certification}
%======================================================================

%----------------------------------------------------------------------
\section{The combinatorial reduction primitive and integer bounds}\label{sec:reduction}
%----------------------------------------------------------------------
% SOURCE:  hecke_action.tex:1162-1338 (RH)
% NOTATION: rows nu, T_N, "t,j,s (running index)", lift{sigma}, Rel, cring,
%           Mat_2(Z), gamma (relation element), zeta of sec:notation-crosswalk
%           applied.  In particular: the coset index is j and the chain-term
%           index t, never i (C15, C16); i is sqrt(-1) and the elliptic vertex
%           v_6; the distinguished lift is \lift{sigma}, never i(sigma); the
%           level is N = p^nu and the source's bounds in "p" are bounds in N.
% ERRATA:  none new.  E16 (T_N) applies to the source text re-homed here.

\paragraph{What this chapter fixes.}
Everything in \S\S\ref{sec:chainlevel}--\ref{sec:equivalence} consumes exactly one
combinatorial primitive: \emph{exact bookkeeping of $\Ga$- and $\Ga''$-translates of
simplices and cosets}. This chapter specifies that primitive as exact arithmetic on
bounded integers (Definition~\ref{def:primitives}), gives the terminating construction
of the $\tau$-table whose output Definition~\ref{def:tau} consumes and
Obligation~\ref{obl:tau} certifies (Algorithm~\ref{alg:tau}), proves the magnitude
bounds the construction obeys (Lemma~\ref{lem:bounds}), and fixes the arithmetic backend
that those bounds select (Theorem~\ref{thm:backend}). Nothing here is a statement about
$\Heq{q}$: the chapter's entire output is a certified finite table of integer chains and
a certified integer bound $B$.

\paragraph{Standing hypotheses.}
\ref{H4} throughout; $p$ a prime, $\nu\ge1$, $N=p^{\nu}$, $g=\smat{1}{0}{0}{N}$,
$\Ga''=\Ga_0(N)$, $m=N+p^{\nu-1}$, and the transporters $s_j$, coset elements
$\delta_j=s_j^{-1}g$, class map $\mathrm{cls}$, index action $j\cdot\gamma$ and cocycle
$h_j(\gamma)$ of Definition~\ref{def:cosetreps} and Lemma~\ref{lem:cocycle}. Per
rule~\ref{A1} each statement repeats what it uses. \textbf{No statement of this chapter
carries a hypothesis on $\ell$}: the chapter is integral, and the only place $\ell$ enters
is the recorded mod-$\ell$ fallback of Algorithm~\ref{alg:tau}(3) and the mod-$\ell$
backend box of Theorem~\ref{thm:backend}.

\paragraph{Notation local to this chapter.} Four symbols are used here and nowhere
else in this document, and are declared once rather than added to
chapter~\ref{sec:notation} (rule~\ref{A5} governs \emph{document-wide} symbols; a
chapter-local declaration is the device already used inside the proof of
Proposition~\ref{prop:fullassembly}): $\lVert A\rVert:=\max_{j,t}\abs{A_{jt}}$, the
entrywise sup-norm on $\Mat_2(\Z)$; $\mathcal{F}:=\{z\in\Hyp:\abs{\mathrm{Re}\,z}\le
\tfrac12,\ \abs{z}\ge1\}$, the standard closed fundamental domain of $\Ga$;
$\mathrm{pos}(v)$, the exact position in $\Hyp$ of a vertex $v$ of $K$
(Definition~\ref{def:model7}); and $\mathcal{N}$, a finite subcomplex of $\Ku$. The
generator multiset of the model is
$\mathcal{G}=\{T,S,ST,TS\}$ of \S\ref{sec:model}; its elements and $-I$ all have entries
in $\{0,\pm1\}$, so $\lVert \gamma\rVert\le1$ for $\gamma\in\mathcal{G}\cup\{-I\}$.

\subsection{The sub-primitives}\label{sec:subprimitives}

\begin{definition}[Reduction primitives]\label{def:primitives}
The four primitives below are the \emph{only} places where exact unbounded-in-principle
integer arithmetic occurs. Each is a terminating procedure on explicit data
(rule~\ref{A3}); each is stated with its output specification, and the properties it
guarantees are Lemma~\ref{lem:bounds} and the certification of
Obligation~\ref{obl:tau}, not assumptions.

\begin{enumerate}
\item[(P1)] \textbf{Coset class map, index action and cocycle.} Given $\gamma\in\Ga$ and
  a coset index $j$: compute $\mathrm{cls}(s_j\gamma)=(c,d)\gamma\bmod N$ as a
  \emph{row}-vector product of the bottom row $(c,d)$ of $s_j$ with $\gamma$ (never
  transposed), look the class up in the table of Definition~\ref{def:cosetreps} to get
  $j\cdot\gamma$, and form the cocycle
  $h_j(\gamma)=s_j\gamma s_{j\cdot\gamma}^{-1}\in\Ga''$ by two exact $2\times2$ products
  and one inversion of a unimodular matrix. The class lookup is arithmetic modulo $N$;
  only the cocycle needs exact integers.
\item[(P2)] \textbf{Word reduction in $\Ga$.} Given $\gamma\in\SL_2(\Z)$, produce a word
  \[
    \gamma\;=\;\pm\,T^{a_1}S\,T^{a_2}S\cdots T^{a_k}
  \]
  by the \emph{nearest-integer} Euclidean algorithm on the first column: while
  $c\neq0$, choose the integer $q$ with $\abs{a-qc}\le\abs{c}/2$ (ties resolved by
  rounding the residue up, so that $q$ is a function of $(a,c)$ alone and the word is
  deterministic), replace $M$ by $S^{-1}T^{-q}M$, and emit the letters $T^{q}$
  (omitted when $q=0$) and $S$. The loop terminates because $\abs{c}$ strictly
  decreases; on exit $c=0$ and $\det M=1$ force $M=\pm T^{\pm b}$, which supplies the
  optional leading $-I$ and the trailing $T$-power. Intermediate magnitudes obey
  Lemma~\ref{lem:bounds}(4), and $k=O(\log\lVert \gamma\rVert)$. Used to express the relative
  position of two translates in step~2 of Algorithm~\ref{alg:tau}; \emph{not} used in
  the hot path, where $\rhon(\gamma)$ is evaluated directly from $\gamma\bmod\ell$.
\item[(P3)] \textbf{Hermite reduction at determinant $N$.} Given $\delta\in\Mat_2(\Z)$ with
  $\det\delta=N$, produce $\gamma\in\Ga$ and the normal form
  $\gamma\delta=\smat{N/d}{b}{0}{d}$ of Proposition~\ref{prop:cosets}(1): Euclid on the
  first column, then column normalisation. Intermediate entries are bounded by
  $2\max(\lVert \delta\rVert,N)$ by the same argument as (P2).
  \textbf{This primitive is specified and has no live call site} --- the closed forms of
  Definition~\ref{def:cosetreps} produce $s_j$, and hence
  $\delta_j=s_j^{-1}g$, directly, so the pipeline never needs to reduce an arbitrary
  determinant-$N$ matrix. It is retained because the \emph{proof} of
  Proposition~\ref{prop:cosets}(1) is this procedure, and because any future entry point
  taking an arbitrary $\delta$ needs it.
\item[(P4)] \textbf{Translate bookkeeping, and the canonical key.} A $q$-simplex of
  $\Ku$ is represented as a pair $(\gamma,\sigma)\in\Ga\times\orient{q}(K)$ standing for
  $\gamma\cdot\lift{\sigma}$. Two pairs represent the same oriented simplex iff they
  differ by a stored transporter, so a \emph{canonical key} is computed by ranging over
  the finite transporter set $\mathrm{Tr}(\sigma,\widehat\sigma)$ of the model, taking the
  lexicographically least $(a,b,c,d)$ among the candidates $\gamma\cdot(\text{transporter})$,
  and recording the parity of the induced permutation of the vertex slots as a sign
  $\pm1$. The key quotients by $-I$ (the transporter set is closed under negation and
  $-I$ acts trivially on $\Ku$ by \ref{H4}(iv)); it is idempotent, and it is the
  \emph{only} equality test on simplices of $\Ku$ used anywhere in this document. Every
  group element arising is a product of at most two stored elements.
\end{enumerate}
\end{definition}

\begin{remark}[Why (P1)--(P4) and nothing else]\label{rmk:primitiveclosure}
The Hecke pipeline of \S\S\ref{sec:chainlevel}--\ref{sec:induced} touches $\Z$ in exactly
four places: building the coset tables (P1), building the $\tau$-table
(Algorithm~\ref{alg:tau}, which uses (P2), (P4) and the predicates of
\S\ref{sec:fpfilter}), comparing simplices of $\Ku$ (P4), and reducing a group element
modulo $\ell$ or modulo $N$ for a consumer (Theorem~\ref{thm:backend}, early reduction).
Everything downstream --- $\rhon$, $\Th_g$, $\Th^{(q)}_{\mathrm{eq}}$, $\Amat{q}$, the
eliminations --- is linear algebra over $\cring$. That is the content of the contract box
at the end of this chapter, and it is what makes Theorem~\ref{thm:backend} a statement
about a \emph{finite} arithmetic workload rather than about the whole computation.
\end{remark}

\subsection{Construction of the \texorpdfstring{$\tau$}{tau}-table}\label{sec:tautable}

\begin{algorithm}[TAU-TABLE]\label{alg:tau}
Assume \ref{H4}. \textbf{Input:} $p$, $\nu$, $N=p^{\nu}$; the model $(K,\Rel)$ with its
relation and transporter tables and the exact vertex positions of
Definition~\ref{def:model7}; the representatives $\{s_j\}_{j=1}^{m}$ of
Definition~\ref{def:cosetreps}. \textbf{Output:} for every $j\in\{1,\dots,m\}$ and every
$\sigma\in\orient{q}(K)$, $q=0,1,2$, an integer chain
\begin{equation}\label{eq:taurecord}
  \tau\bigl(s_j\cdot\lift{\sigma}\bigr)\;=\;\sum_{t}n^{(j,\sigma)}_{t}\;
    \gamma_t\cdot\lift{\sigma_t},
  \qquad n^{(j,\sigma)}_{t}\in\Z,\ \gamma_t\in\Ga,\ \sigma_t\in\orient{q}(K),
\end{equation}
stored as triples $(n^{(j,\sigma)}_t,\gamma_t,\sigma_t)$ in the canonical form of
(P4), satisfying (T1)--(T3) of Definition~\ref{def:tau} by construction and by step~4.
By (T1) it suffices to store $\tau$ on the finitely many slots $(j,\sigma)$; that is one
slot per simplex of $\Ga''\backslash\Ku$ \emph{up to the derivation of
step~0}. At $p=2$, $\nu=1$ ($m=3$) the table is $3\cdot(7+11+5)=69$ chains.

\begin{enumerate}
\item[\textbf{0.}] \textbf{Slots, and the single (T1) helper.} The pairs $(j,\sigma)$ are
  \emph{not} pairwise $\Ga''$-inequivalent: $\orient{q}(K)$ is not a section of the
  $\Ga$-groupoid orbits (already $\lift{v_3}=T\cdot\lift{v_0}$, so
  $s_j\cdot\lift{v_3}=h_j(T)\,s_{j\cdot T}\cdot\lift{v_0}$ with $h_j(T)\in\Ga''$).
  Exactly one slot per $\Ga''$-orbit is therefore \emph{constructed} by steps~1--3, and
  every other slot is \emph{derived} from it through one shared helper implementing
  (T1), $\tau(h\cdot x)=(g^{-1}hg)\cdot\tau(x)$, with the sign of (P4) tracked. The
  \emph{same} helper is used by step~4, so construction and certification cannot disagree
  about how a derived slot is obtained.
\item[\textbf{1.}] \textbf{Vertices.} For each constructed slot $(j,v)$: form the
  integral matrix $\adjg(g)\,s_j$ with $\adjg(g)=\mathrm{diag}(N,1)=N\,g^{-1}$ and let
  $z:=\bigl(\adjg(g)s_j\bigr)\cdot\mathrm{pos}(v)\in\Hyp$, where $\mathrm{pos}(v)$ is the
  exact quadratic point of Definition~\ref{def:model7}. (The M\"obius action is
  projective and $\det(\adjg(g)s_j)=N>0$, so $\adjg(g)s_j$ acts exactly as $g^{-1}s_j$;
  \textbf{no rational number is ever formed}.) Reduce $z$ to
  $\mathcal{F}$ by the standard alternation --- shift by $T^{-q}$ with $q$ the
  nearest-integer quotient of the cleared real part, then apply $S$ while $\abs{z}<1$,
  repeat --- obtaining $\gamma_{\mathrm{red}}\in\Ga$ with
  $\gamma_{\mathrm{red}}z\in\mathcal{F}$. Choose a vertex $w$ of $K$:
  \begin{enumerate}
  \item[(1a)] if $\gamma_{\mathrm{red}}z$ is \emph{exactly} the position of a model
    vertex, take that vertex;
  \item[(1b)] otherwise let $S$ be the set of vertices at minimal euclidean distance
    from $\gamma_{\mathrm{red}}z$, decided exactly. If $\abs S=1$ take its element.
    Otherwise the tie is broken $\invK$-\emph{equivariantly} (the vertex involution of
    Lemma~\ref{lem:iotamodel}, read off the exact coordinates): if $S$ contains a full
    $\invK$-orbit --- a swapped pair $\{v,\invK v\}$, or $v_6$ --- take $v_6$, the unique
    $\invK$-fixed vertex of $K$; otherwise $S$ meets every orbit in at most one vertex,
    and take the one in the orbit of smallest index
    ($\{v_0,v_3\}<\{v_1,v_2\}<\{v_4,v_5\}<\{v_6\}$). The first branch is forced: a
    reduced point on the mirror $\mathrm{Re}\,z=0$ has an $\invK$-stable tied set, and
    no choice \emph{within} a swapped pair commutes with $\invK$
    (Remark~\ref{rmk:iotatausym}). Before v1.1.0 ties were broken by ascending vertex
    index, which is not $\invK$-equivariant.
  \end{enumerate}
  Set $\tau(s_j\cdot\lift{v}):=\gamma_{\mathrm{red}}^{-1}\cdot\lift{w}$, a single term
  with coefficient $1$, so (T3) holds by inspection. All branch decisions of the
  reduction and of (1b) are the certified predicates of \S\ref{sec:fpfilter}.
\item[\textbf{2.}] \textbf{Edges.} For a constructed slot $(j,e)$ with $e=(v_0,v_1)$: the
  two faces of $s_j\cdot\lift{e}$ are themselves slots, so their $\tau$-values
  $\gamma_a\cdot\lift{w_a}$ and $\gamma_b\cdot\lift{w_b}$ are already stored (fetched
  through the step-0 helper where the face is derived). If the two canonical keys
  coincide, set $\tau(s_j\cdot\lift{e}):=0$: the empty $1$-chain has exactly the required
  boundary. Otherwise reduce $\gamma_a^{-1}\gamma_b$ to a word in $T^{\pm1},S$ by (P2),
  drop the $-I$ letter (it acts trivially on $\Ku$), and concatenate, for each remaining
  letter $L$, the stored \emph{path template} from $\lift{w}$ to $L\cdot\lift{w'}$
  translated by the running prefix. A path template is a $1$-chain of $\Ku$ supported in
  $\lift{K}\cup L\cdot\lift{K}$, computed once from the model and \emph{certified at
  construction}: its boundary equals $L\cdot\lift{w'}-\lift{w}$ exactly over $\Z$.
  \textbf{Since v1.1.0 the template is tree-routed.} Let $Y\subset\Ku$ be the
  $\Ga$-orbit of the two edges $\{v_4v_6\}$, $\{v_5v_6\}$ --- a tree, the Bass--Serre
  tree of $\mathrm{PSL}_2(\Z)$ with the elliptic points as its vertices --- and let $t$
  be the spoke attachment $t(v_0)=v_4$, $t(v_1)=t(v_2)=v_6$, $t(v_3)=v_5$, $t(v)=v$ on
  $\{v_4,v_5,v_6\}$, a table that commutes with $\invK$ and with $\lift{v_3}=T\lift{v_0}$
  ($t(v_3)=T\,t(v_0)$). The template is
  $[\lift w,\lift{t(w)}]+Y\text{-path}\bigl(\lift{t(w)},L\lift{t(w')}\bigr)-L[\lift{w'},\lift{t(w')}]$,
  the $Y$-path being the unique $1$-chain of $Y$ with that boundary. Concatenated along
  \emph{any} word the templates telescope, because $Y$-paths compose to the unique
  $Y$-path, to $[A,t(A)]+Y\text{-path}(t(A),t(B))-[B,t(B)]$ for the two endpoints
  $A=\gamma_a\lift{w_a}$, $B=\gamma_b\lift{w_b}$: a chain of the endpoints alone,
  $\Ga$-equivariant, $\inv_{\Xbulk}$-equivariant and antisymmetric in $A\leftrightarrow B$,
  so neither the word's tie rule nor any traversal order reaches the output. (Before
  v1.1.0 the templates were breadth-first vertex paths in $K$ in ascending neighbour
  order, which are neither $\invK$-equivariant nor reversal-symmetric;
  Remark~\ref{rmk:iotatausym}.) If
  the word is empty while the two keys differ, use the net-identity detour $T,T^{-1}$,
  whose boundary still telescopes. Set $\tau(s_j\cdot\lift{e}):=$ the accumulated
  $1$-chain. (This is the Manin/continued-fraction device \cite{manin,cremona} in
  general position.)
\item[\textbf{3.}] \textbf{Triangles.} For a constructed slot $(j,t)$:
  $c:=\tau(\partial(s_j\cdot\lift{t}))$ is a $1$-cycle of $\Ku$, by step~2 and (T2) in
  lower degree. Solve $\partial x=c$ for a $2$-chain $x$ supported on a finite subcomplex
  $\mathcal{N}\subseteq\Ku$: initialise $\mathcal{N}$ as the union of the translates
  occurring in $\supp c$, take as columns the canonical keys of the $2$-simplices of
  $\mathcal{N}$ and as rows the canonical keys of $\supp c$ together with all column
  faces, and solve the resulting sparse integer system by column-echelon elimination
  with exact divisibility checks. On failure, grow $\mathcal{N}$ by one ring of adjacent
  translates ($T$, $T^{-1}$, $S$) and retry. \textbf{Termination:} $\varinjlim
  \mathcal{N}=\Ku$ is acyclic (\ref{H4}(ii)), so every $1$-cycle bounds at a finite
  stage. Set $\tau(s_j\cdot\lift{t}):=x$.
  \emph{Solve regime (recorded).} The integer solve is attempted first and is the
  default; solving the same system modulo $\ell$ instead is permitted only behind the
  explicit setting of owner decision OP2, marks the affected chains and the table
  \code{integral = false}, and is recorded --- see Remark~\ref{rmk:tauintegral}, which
  also records that $0$ of $66$ shipped points used it. Both regimes produce a valid
  $\tau$ over their ring, and they agree on cohomology by
  Theorem~\ref{thm:welldef}(b) where that theorem's hypothesis holds.
  No Smith/Hermite normal-form library and no arbitrary-precision arithmetic is used;
  the entry growth of the integer solve is bounded by the fixed-width tier of
  Theorem~\ref{thm:backend} and overflow-checked, and an overflow is a tier failure that
  aborts (\cite{havas} is the classical reference for controlling that growth in
  general; the systems here are small enough that the checked tier suffices).
\item[\textbf{4.}] \textbf{Certification (mandatory).} Verify \emph{exactly} --- over
  $\Z$, or over $\Fl$ for identities involving a chain that step~3 solved modulo $\ell$
  --- the full predicate list of Obligation~\ref{obl:tau}: structure and canonicity,
  (T3) in degree $0$, (T2) in degrees $1$ and $2$ with the right-hand side evaluated
  \emph{only} through the step-0 helper, the (T1) bookkeeping over every stored
  transporter, the magnitude bound of Lemma~\ref{lem:bounds}(5), and the consumer table
  of \eqref{eq:xiexplicit}. \textbf{A table failing any clause is rejected}: the datum
  never exists in an uncertified state.
\end{enumerate}
\end{algorithm}

\begin{remark}[Certification over geometry]\label{rmk:certoverGeometry}
Step~4 converts every heuristic and geometric choice of steps~1--3 --- the nearest
vertex, the tie-break, the path templates, the neighbourhood growth, the
FP-filtered reduction --- into a \emph{proof} that the emitted table is a $g$-twisted
chain datum in the sense of Definition~\ref{def:tau}. Correctness of the resulting Hecke
operator then depends only on Theorems~\ref{thm:welldef} and \ref{thm:xi}, and on no
geometric assertion whatever. Two consequences, both used elsewhere in this document and
neither optional:
\begin{enumerate}
\item \textbf{The choices in steps~1--3 are free, and changing one changes the table.}
  The tie-break of (1b) and the templates of step~2 are, since v1.1.0,
  $\invK$-equivariant, and Remark~\ref{rmk:iotatausym} records the computation showing
  that the canonical $\tau$ is consequently $\inv$-symmetric --- and that the pre-v1.1.0
  choices (ascending-index tie-break, breadth-first templates) were not, each on its own.
  A different admissible choice is a different admissible $\tau$; by
  Theorem~\ref{thm:welldef}(b) it moves nothing on $\Heq{q}$ where $6\in\cring^{\times}$,
  and by Remark~\ref{rmk:indeplevel} it moves $\Th^{(q)}_{\mathrm{flat}}$ --- which is
  why the v1.1.0 change re-baselined every frozen cochain-level reference value and no
  cohomology-level one.
\item \textbf{Elliptic slots are forced, not chosen.} If the slot $s_j\cdot\lift{v}$ has
  a non-trivial stabiliser in $\Ga''$ modulo $\pm I$, then (T1) forces the assigned
  vertex, the reduced point $\gamma_{\mathrm{red}}z$ is exactly an elliptic fixed point,
  and branch~(1a) returns it. This is an explanation of why the construction succeeds,
  not a step of any proof: step~4 is what accepts the table.
\end{enumerate}
\end{remark}

\subsection{Integer bounds}\label{sec:bounds}

\begin{lemma}[Entry bounds]\label{lem:bounds}
\PROVED{} Assume \ref{H4}, and let $\lVert A\rVert=\max_{j,t}\abs{A_{jt}}$ for
$A\in\Mat_2(\Z)$. Then, with $N=p^{\nu}$ and no hypothesis on $\ell$:
\begin{enumerate}
\item $\lVert AB\rVert\le2\lVert A\rVert\lVert B\rVert$ for $A,B\in\Mat_2(\Z)$; hence a product of $k$
  matrices is bounded by $2^{\,k-1}\prod_t\lVert A_t\rVert$. \LEAN{lem:bounds}
\item \textbf{Fixed input data.} $\lVert s_j\rVert$, $\lVert s_j^{-1}\rVert$, $\lVert \delta_j\rVert$,
  $\lVert g\rVert\le\max(N,1)$ for every $j$ (Lemma~\ref{lem:cocycle}) \LEAN{lem:bounds} (this clause), and
  $\lVert \gamma\rVert\le1$ for $\gamma\in\mathcal{G}\cup\{-I\}$.
\item \textbf{Cocycles.} $\lVert h_j(\gamma)\rVert=\lVert s_j\gamma s_{j\cdot\gamma}^{-1}\rVert
  \le4N^{2}\lVert \gamma\rVert\le4N^{2}$ for $\gamma\in\mathcal{G}\cup\{-I\}$.
  \LEAN{lem:bounds} (the first inequality)
\item \textbf{Euclidean reduction.} Every intermediate matrix of (P2) on input $\gamma$
  satisfies $\lVert M\rVert\le\max(\lVert \gamma\rVert,2)$, and in particular
  $\lVert M\rVert\le2\max(\lVert \gamma\rVert,1)$ \LEAN{lem:bounds} (for (P2), on an abstract step relation); the same bound holds for (P3) with
  $\max(\lVert \delta\rVert,N)$ in place of $\lVert \gamma\rVert$.
\item \textbf{$\tau$-table elements.} Every $\gamma_t$ stored by
  Algorithm~\ref{alg:tau} satisfies $\lVert \gamma_t\rVert\le C_K\,N^{2}$ for a constant $C_K$
  of the fixed model, independent of $p$ and $\nu$; and the implementation must
  \emph{certify} the actual attained bound
  $B_{\mathrm{tab}}:=\max_{j,\sigma,t}\lVert \gamma_t\rVert$ at build time, by a full recount
  over the stored table in overflow-checked arithmetic
  (Obligation~\ref{obl:tau}, Obligation~\ref{obl:backend}).
\item \textbf{Reduction intermediates --- which (1)--(5) do \emph{not} bound} (v1.5.0;
  \file{docs/implementation/slot3/WP-1\_report.md} finding~1). Clauses (1)--(5) bound
  the matrices the pipeline \emph{stores}. The fundamental-domain reduction of
  Algorithm~\ref{alg:tau} step~1 also carries a quantity that is not one of them: the
  quadratic point $z=(u+v\sqrt{D}\,\mathrm{i})/w$ on cleared denominators
  ($u,v,w\in\Z$, $w>0$, $v>0$, $D\in\{1,3\}$, $\gcd(u,v,w)=1$ --- the exact
  representation on which (P2) and the predicates of \S\ref{sec:fpfilter} act), with
  $\lVert z\rVert:=\max(\abs u,\abs v,w)$. For it:
  \begin{enumerate}[label=(\alph*)]
  \item \emph{One M\"obius step.} For $\gamma\in\Mat_2(\Z)$ with $\det\gamma>0$, every
    integer formed in evaluating $\gamma\cdot z$ --- the three coordinates and the
    products that build them --- has magnitude at most
    $7\lVert\gamma\rVert^{2}\lVert z\rVert^{2}$, and hence
    $\lVert\gamma\cdot z\rVert\le7\lVert\gamma\rVert^{2}\lVert z\rVert^{2}$ before the
    $\gcd$ normalisation.
  \item \emph{The point step~1 starts from.} With
    $C:=\max_{v}\lVert\mathrm{pos}(v)\rVert$ the largest cleared coordinate of a vertex
    position of the fixed model ($C=12$ on the $7$-vertex model of
    \S\ref{sec:model}), the initial point
    $z_{0}=\bigl(\adjg(g)s_j\bigr)\cdot\mathrm{pos}(v)$ of step~1 satisfies
    $\lVert z_{0}\rVert\le7C^{2}N^{2}$, and the order $C^{2}N^{2}$ is attained.
  \item \emph{Consequence for the backend.} $\lVert z_{0}\rVert$ is of order $N^{2}$
    with a \emph{model} constant in front of it, not of order $B$; so the arithmetic of
    step~1 is bounded by $7\lVert z\rVert^{2}$ and \emph{not} by the $2B^{2}$ of
    Theorem~\ref{thm:backend}(i), and a tier chosen against $B$ alone does not cover
    it. The tier criterion for step~1 is therefore
    $7\max_z\lVert z\rVert^{2}$, the maximum taken over the points the reduction
    actually passes through; (6b) bounds the first of them, and that the others stay at
    that order is \emph{measured} and not asserted here
    (Remark~\ref{rmk:btabsmall}; v1.5.1, which struck v1.5.0's unproved ``with
    $\max_z\lVert z\rVert$ of order $C^{2}N^{2}$'' from the clause and left it where the
    measurement is recorded). In particular the width must be chosen against at least
    $\max\bigl(B,\,7C^{2}N^{2}\bigr)$ and its square, and an implementation that cannot
    bound the reduction a priori must be overflow-checked per operation, which
    Theorem~\ref{thm:backend}(3) requires of it anyway.
  \end{enumerate}
  The gap this clause closes was a real one and not a hypothetical: the measurement
  that found it, and the reach it fixes, are recorded at
  Remark~\ref{rmk:btabsmall} and in \S\ref{sec:heckedatum}'s code map.
\item \textbf{The wide tier --- two stored-class operands in $128$ bits} (v1.23.0-pre.9; slot-24
  package ELL-W, \file{docs/implementation/slot24/ELLW\_report.md}). Let $N\ge2$, let
  $B=\max(B_{\mathrm{tab}},4N^{2})$ be the certified bound of Theorem~\ref{thm:backend}, and assume
  \[
    2NB\;\le\;2^{63}-1 .
  \]
  Then $B<2^{61}$, and:
  \begin{enumerate}[label=(\alph*)]
  \item \emph{Storage.} Every stored $\gamma_t$, the inverse of every stored $\gamma_t$, and every
    matrix of norm $\le2B$ --- in particular every intermediate of (P2) on an input of norm $\le B$,
    by (4) --- has entries of magnitude $<2^{62}$, so signed $64$-bit integers hold them.
  \item \emph{Two stored-class operands.} For $A,A'\in\Mat_2(\Z)$ with
    $\lVert A\rVert,\lVert A'\rVert\le B$, each of the two products in an entry of $AA'$, and each of
    the products $ad$, $bc$ of $\det A$, has magnitude $\le B^{2}<2^{122}$, and each sum of two of them
    has magnitude $<2^{123}$; so $AA'$ and $\det A$ are evaluated exactly in signed $128$-bit
    integers. Clause~(b) bounds the \emph{intermediates}; the \emph{result} is bounded only by
    $2B^{2}$ (by (1)), which need not be below $2^{63}$, and is returned to the $64$-bit tier by a
    checked narrowing, which aborts rather than wraps (Obligation~\ref{obl:backend}(iii)).
  \item \emph{A fixed datum times a stored matrix.} For
    $F\in\{g,\adjg(g),s_j,s_j^{-1},\delta_j\}$ ($\lVert F\rVert\le\max(N,1)=N$ by (2)) and
    $\lVert\gamma\rVert\le B$: $\lVert F\gamma\rVert,\lVert\gamma F\rVert\le2NB\le2^{63}-1$, and every
    intermediate is one of the two products ($\le NB$) or their sum; exact in signed $64$-bit
    integers. These are the products the level-one consumers of $\tau$ form before the early
    reduction (Theorem~\ref{thm:backend}, \emph{Early-reduction scheme}); other consumers (the level
    structures, e.g.\ the three-fold product $v\,\gamma\,w^{-1}$ of the level-$\ell$ route) form
    checked products not covered by (c), and abort beyond $64$ bits.
  \item \emph{Reach.} When $B=4N^{2}$ (i.e.\ $B_{\mathrm{tab}}\le4N^{2}$) the hypothesis reads
    $8N^{3}\le2^{63}-1$, i.e.\ $N\le2^{20}-1=1048575$ (largest prime $1048573$).
  \end{enumerate}
  That the result of each product of (b) the pipeline actually forms is below $2^{63}$ is
  \COMP{} \emph{measured} and not asserted here (Remark~\ref{rmk:btabsmall}): the clause makes no
  value depend on it, because an exceeding result aborts. \textbf{Not bounded a priori by
  (a)--(c)}, and formed in overflow-checked arithmetic whose completion is likewise \COMP{}
  measured (v1.23.0-pre.9, audit F1): the input $\mathrm{rel}=\gamma_a^{-1}\gamma_b$ of (P2) in
  Algorithm~\ref{alg:tau} step~2, a product of two stored matrices ($\lVert\mathrm{rel}\rVert\le2B^{2}$,
  so its (P2) intermediates are $\le4B^{2}$ by (4)), and the running edge prefix
  $P=\gamma_a L_1\cdots L_j$ along its word (a priori of order $B^{3}$; the word is expanded letter
  by letter); the (T1) conjugate $g^{-1}hg$ of step~0 ($\lVert\cdot\rVert$ up to $2N^{2}B$, by
  (1), (2) --- the gap Remark~\ref{rmk:primitiveclosure} leaves, and the same in the narrow tier),
  whose product with a stored $\gamma_t$ goes through the $128$-bit form, itself checked, so it can
  only abort; and the consumers named at the end of (c).
\end{enumerate}
\end{lemma}

\begin{proof}
(1) Each entry of $AB$ is a sum of two products of entries.

(2) is Lemma~\ref{lem:cocycle}: the entries of $s_j$ lie in $\{0,\pm1,k,pc\}$ with
$0\le k<N$ and $0\le pc<N$; $s_j^{-1}$ swaps the diagonal and negates the off-diagonal;
$\delta_j=s_j^{-1}g$ scales the second column by $N$, which does not increase the bound
$\max(N,1)$ because the second column of $s_j^{-1}$ has entries in $\{0,\pm1\}$ for
both branches of Definition~\ref{def:cosetreps}. For $\mathcal{G}\cup\{-I\}$, read
off \S\ref{sec:model}.

(3) Apply (1) twice to (2): $\lVert s_j\gamma s_{j\cdot\gamma}^{-1}\rVert
\le2\lVert s_j\rVert\cdot2\lVert \gamma\rVert\lVert s_{j\cdot\gamma}^{-1}\rVert\le4N^{2}\lVert \gamma\rVert$.

(4) Write $M=\smat{a}{b}{c}{d}$ with $ad-bc=1$ and $c\neq0$, and let $q$ be the
nearest-integer quotient, so $a':=a-qc$ satisfies $\abs{a'}\le\abs{c}/2$. Put
$b':=b-qd$. From $a'd-b'c=\det M=1$ we get $b'=(a'd-1)/c$, whence
\[
  \abs{b'}\;\le\;\frac{\abs{a'}\abs{d}+1}{\abs{c}}
  \;\le\;\frac{\abs{c}\abs{d}/2+1}{\abs{c}}
  \;=\;\frac{\abs{d}}{2}+\frac{1}{\abs{c}}\;\le\;\frac{\abs{d}}{2}+1 .
\]
The intermediate matrix $T^{-q}M=\smat{a'}{b'}{c}{d}$ therefore has
$\lVert T^{-q}M\rVert\le\max\bigl(\abs{c},\abs{d},\abs{c}/2,\abs{d}/2+1\bigr)
\le\max(\lVert M\rVert,2)$, and $S^{-1}T^{-q}M$ has the same entries up to sign and order.
Induction on the loop from $\lVert M_0\rVert=\lVert \gamma\rVert$ gives
$\lVert M_k\rVert\le\max(\lVert \gamma\rVert,2)$ for every $k$. The exit matrix $\pm T^{\pm b}$ is
one of the $M_k$. For (P3) the same computation applies to the first column of
$\delta$ with $\det\delta=N$ in place of $1$, giving the extra additive $N/\abs{c}\le N$.

(5) The reduction word of step~1 for $\bigl(\adjg(g)s_j\bigr)\cdot\mathrm{pos}(v)$ has
convergents bounded by the entries of the integralised input $\adjg(g)s_j$, whose norm is
at most $2N\cdot N=2N^{2}$ by (1) and (2); by (4) the intermediates never exceed twice
that. The path templates of step~2 and the neighbourhood of step~3 contribute a factor
depending on the fixed model only, and products of at most three such elements occur;
both are absorbed into $C_K$. \textbf{The a-priori constant is not what makes the
pipeline sound}: the build-time recount of $B_{\mathrm{tab}}$ is, and it is what
Theorem~\ref{thm:backend} consumes.

(6a) On cleared denominators the M\"obius action of $\gamma=\smat{a}{b}{c}{d}$ with
$\Delta:=\det\gamma>0$ on $z=(u+v\sqrt D\,\mathrm{i})/w$ is
\[
  u'=(au+bw)(cu+dw)+D\,ac\,v^{2},\qquad v'=vw\Delta,\qquad
  w'=(cu+dw)^{2}+D\,c^{2}v^{2},
\]
which is the only formula evaluated. Put $G:=\lVert\gamma\rVert$, $Z:=\lVert z\rVert$.
Then $\abs{au+bw}\le2GZ$ and $\abs{cu+dw}\le2GZ$, so the first summand of $u'$ and of
$w'$ is $\le4G^{2}Z^{2}$ and the second is $\le DG^{2}Z^{2}\le3G^{2}Z^{2}$; and
$\abs{\Delta}\le2G^{2}$ gives $\abs{v'}\le2G^{2}Z^{2}$. Every intermediate the formula
forms is one of these products or their sum, hence $\le7G^{2}Z^{2}$, and so is each of
$\abs{u'},\abs{v'},w'$. The $\gcd$ normalisation only divides.
\emph{The two shapes step~1 applies} (v1.5.1; the bound
Theorem~\ref{thm:backend}(iv) quotes is $7\lVert z\rVert^{2}$, with no $G$, and $G$ is
not $1$ for a translation by a large quotient, so the specialisation is part of the
proof and not a corollary of the displayed inequality). The reduction applies only
$S=\smat{0}{-1}{1}{0}$, where $G=1$ and (6a) gives $7Z^{2}$ outright; and the shift
$T^{-q}=\smat{1}{-q}{0}{1}$ by the \emph{nearest}-integer quotient $q$ of $u$ by $w$,
where $c=0$ and $d=\Delta=1$ collapse the formulas to $u'=(u-qw)w$, $v'=vw$, $w'=w^{2}$
with $\abs{u-qw}\le w/2$, so that $qw$, $u-qw$ and each of $\abs{u'},\abs{v'},w'$ is
$\le2Z^{2}$. Both are therefore $\le7Z^{2}$ as (iv) states, although (6a) applied with
$G=\abs q$ would give only $7q^{2}Z^{2}$. (6b) Apply (6a) with
$z=\mathrm{pos}(v)$, so $Z\le C$, and $\gamma=\adjg(g)s_j$, whose norm is at most $N$:
by Definition~\ref{def:cosetreps} and $\adjg(g)=\mathrm{diag}(N,1)$ the product is
$\smat{N}{0}{k}{1}$ or $\smat{0}{-N}{1}{pc}$ with $0\le k,pc<N$, so $G\le N$ and
$\lVert z_{0}\rVert\le7C^{2}N^{2}$. Attainment (v1.5.1, which names the coordinate that
attains it; v1.5.0 cited $v'=vw\Delta$, which for these matrices is of order $C^{2}N$
and not $C^{2}N^{2}$): it is the \emph{third} coordinate,
$w'=(cu+dw)^{2}+D\,c^{2}v^{2}$, which for $\gamma=\smat{N}{0}{k}{1}$ with $k$ of order
$N$ is of order $k^{2}\lVert\mathrm{pos}(v)\rVert^{2}\asymp C^{2}N^{2}$ --- its second
summand $D\,k^{2}v^{2}$ already is --- so $\lVert z_{0}\rVert\asymp C^{2}N^{2}$,
measured as $\approx144.8\,N^{2}$ against $C^{2}=144$ with $C=12$ on the shipped model
(Remark~\ref{rmk:btabsmall}). (6c) is (6a), (6b) and the definition of
$B=\max(B_{\mathrm{tab}},4N^{2})$: $B$ is a bound on stored matrices, $\lVert
z_{0}\rVert$ is not bounded by it, and $7\lVert z\rVert^{2}$ is what an
overflow-checked step of (P2) must fit.

(7) From $N\ge2$: $B\le(2^{63}-1)/(2N)\le(2^{63}-1)/4<2^{61}$. (7a) The stored $\gamma_t$ have norm
$\le B_{\mathrm{tab}}\le B$; the inverse of a unimodular matrix has the same entries up to sign and
order; and (4) bounds the (P2) intermediates on an input of norm $\le B$ by $\max(B,2)\le2B<2^{62}$.
(7b) Each product is of two entries of magnitude $\le B$, hence $\le B^{2}<2^{122}$, and a sum of two
such is $<2^{123}<2^{127}$; the result bound $2B^{2}$ is (1). (7c) Each entry of $F\gamma$ (or
$\gamma F$) is a sum of two products of an entry of $F$ ($\le N$) and one of $\gamma$ ($\le B$), so each
product is $\le NB$ and the entry $\le2NB$, which is $\le2^{63}-1$ by hypothesis. (7d) $8N^{3}\le2^{63}-1$
iff $N^{3}\le(2^{63}-1)/8$, i.e.\ $N^{3}<2^{60}$, i.e.\ $N<2^{20}$; and $N=2^{20}-1$ gives
$8N^{3}<2^{63}$.
\end{proof}

\begin{remark}[$C_K$ is never evaluated, and $B_{\mathrm{tab}}$ is small]
\label{rmk:btabsmall}
\COMP{} Lemma~\ref{lem:bounds}(5) is stated with an unevaluated model constant on
purpose: no statement of this document depends on its value, because every consumer reads
the certified $B_{\mathrm{tab}}$ instead. Over the shipped grid
$(p,\nu)\in\{(2,1),(3,1),(5,1),(7,1),(2,2)\}$ at
$\ell\in\{5,7,691,134217757\}$ and $n\in\{0,2,10\}$, plus the bad-prime points
$p=\ell\in\{2,5\}$ --- $66$ points --- the certified whole-pipeline bound is
$B\le196$, far inside the first tier of Theorem~\ref{thm:backend}(1)
(\code{HeckeTauTableSelfCheckValidation}, run of 2026-09-15; the same figure is recorded
at Remark~\ref{rmk:tauintegral}).

\smallskip
\textbf{And $B$ is not the whole story} (v1.5.0; slot-3 package WP-1,
\file{docs/implementation/slot3/WP-1\_report.md}, \code{HeckeUEllReachValidation}).
The one place where the certified $B$ does \emph{not} bound the arithmetic is
Lemma~\ref{lem:bounds}(6)'s reduction point, and it was found the hard way. At
$N=\ell=5003$ --- the bad prime, $U_\ell$ --- $B=\max(B_{\mathrm{tab}},4N^{2})=4N^{2}
\approx1.0\cdot10^{8}$, so $2B^{2}\approx2\cdot10^{16}$ and clause~(e) of
Obligation~\ref{obl:tau}'s certification passes with four orders of magnitude to spare,
while Algorithm~\ref{alg:tau} step~1 overflows the $64$-bit tier: the reduction point
was measured at $\lVert z_{0}\rVert\approx144.8\,N^{2}$, i.e.\ $\asymp C^{2}N^{2}$ with
$C^{2}=144$, $C=12$ (v1.5.1: v1.5.0 wrote the two as an equality), attained at the
\emph{first} M\"obius step, with $\max\lVert z\rVert$ constant along the reduction (the
$\gcd$ normalisation returns each step to the entry size, so the growth is one squaring
and not a cascade) --- measured at
$\ell\in\{11,101,691,1009,2003,3617\}$, with $\lVert z\rVert/N^{2}$ rising to $144.8$
and stable there, and independent of $n$, step~1 never seeing $n$. The int64 tier
therefore dies at $(144.8\,N^{2})^{2}\ge2^{63}$, i.e.\ $N\ge4580$, which is exactly the
bracket $3617<N<5003$ the slot-1 record had measured and attributed to the coset
construction: the attribution was wrong, the site is step~1, and the quantity is the one
clause~(6) now bounds. The repair escalates that site to the \code{\_\_int128} tier
\emph{per operation} (Theorem~\ref{thm:backend}(2),(iv)), value-preservingly, and moves
the binding constraint back onto clause~(5) --- see \S\ref{sec:heckedatum}'s code map for
the resulting reach.

\smallskip
\textbf{Past $N=23170$} (v1.23.0-pre.9; slot-24 ELL-W, \file{docs/implementation/slot24/ELLW\_report.md}).
For $U_\ell$ ($N=\ell$, $\nu=1$, $n=0$) the certified $B_{\mathrm{tab}}$ is far below $4N^{2}$ ---
$B_{\mathrm{tab}}=4415$ at $\ell=5003$, $20442$ at $23167$, $20447$ at $23173$, $26480$ at $30011$, i.e.\
$B_{\mathrm{tab}}\approx0.88\,N$ --- so $B=4N^{2}$, the cocycle bound of clause~(3), and clause~(7)'s
hypothesis is its (7d). On every table built past $23170$ (the primes $23173$, $30011$, $43867$ at $n=0$
and $30011$ at $n=10$, $U_\ell$ and $T_{30011}$) the $64$-bit form of every product of clause~(7)(b) the
pipeline formed \emph{completed}: the wide tier escalated $0$ times, so every result was below $2^{63}$
with a margin of many orders of magnitude, and the tier is a certified guarantee there rather than a
path taken. Measured, not asserted (clause~(7), last sentence).
\end{remark}

\subsection{The arithmetic backend}\label{sec:backend}

\begin{theorem}[Arithmetic backend requirements]\label{thm:backend}
\PROVED{} Assume \ref{H4}, and let
\[
  B\;:=\;\max\bigl(B_{\mathrm{tab}},\,4N^{2}\bigr)
\]
be the certified bound of Lemma~\ref{lem:bounds}(5) and (3). All exact integer
arithmetic required by the entire Hecke pipeline consists of
\begin{enumerate}
\item[(i)] $2\times2$ matrix products of at most two stored matrices at a time, of
  magnitude $\le2B^{2}$, and unimodular inversions, of magnitude $\le B$;
\item[(ii)] Euclidean steps (P2), (P3), of magnitude $\le2B$ by
  Lemma~\ref{lem:bounds}(4);
\item[(iii)] the sign predicates of \S\ref{sec:fpfilter}, of magnitude $\le2B^{2}$ on
  their own inputs;
\item[(iv)] \textbf{(v1.5.0)} the M\"obius steps of the fundamental-domain reduction of
  Algorithm~\ref{alg:tau} step~1, of magnitude $\le7B_{\mathrm{red}}^{2}$
  (Lemma~\ref{lem:bounds}(6a) and its proof, which bounds the two shapes step~1 applies),
  where
  $B_{\mathrm{red}}:=\max_z\lVert z\rVert$ over the quadratic points $z$ that reduction passes
  through. Lemma~\ref{lem:bounds}(6b) bounds the \emph{entry} point,
  $\lVert z_{0}\rVert\le7C^{2}N^{2}$, and that $B_{\mathrm{red}}$ is of the same order ---
  i.e.\ that the reduction's own steps do not exceed the norm they start from --- is
  \COMP{} \emph{measured} and not proved here (Remark~\ref{rmk:btabsmall}); nothing below
  needs more, the criteria being stated in $B_{\mathrm{red}}$ itself and the
  implementation being overflow-checked per operation (v1.5.1, replacing v1.5.0's
  ``$B_{\mathrm{red}}\le7C^{2}N^{2}$ by Lemma~\ref{lem:bounds}(6)'', which that lemma
  proves for $z_{0}$ only). \textbf{$B_{\mathrm{red}}$ is not
  bounded by $B$}: it is of order $C^{2}N^{2}$ with the model constant $C$ in front,
  where $B$ is of order $N^{2}$ with the constant $4$ --- and it is squared, so the two
  criteria below are genuinely different.
\end{enumerate}
Consequently:
\begin{enumerate}
\item \textbf{Signed $64$-bit suffices} for (i)--(iii) --- the \emph{stored}-matrix
  arithmetic --- iff $2B^{2}<2^{63}$, i.e.\ $B\le2^{31}$, which holds for every $N$
  whose certified $B_{\mathrm{tab}}$ stays under $2^{31}$ and, since
  $B=\max(B_{\mathrm{tab}},4N^{2})$, for every $N$ with $4N^{2}\le2^{31}-1$, i.e.\
  $N\le23170$, with such a table (v1.5.1, correcting v1.5.0's ``$N\le2^{15}$'': at
  $N=2^{15}$ already $4N^{2}=2^{32}>2^{31}$, and $N\le23170$ is exactly the reach
  \S\ref{sec:heckedatum} records, clause (e) of Obligation~\ref{obl:tau}'s
  certification being this criterion). It
  suffices for (iv) iff $7B_{\mathrm{red}}^{2}<2^{63}$, i.e.\ $B_{\mathrm{red}}\le2^{30}$, which on the
  shipped model ($B_{\mathrm{red}}\approx C^{2}N^{2}\approx144N^{2}$) means $N\lesssim2700$
  \emph{a priori}; the \emph{measured} first overflow is later, at $N\approx4580$,
  because the largest intermediate actually formed is $\approx B_{\mathrm{red}}^{2}$ and
  not $7B_{\mathrm{red}}^{2}$ (Remark~\ref{rmk:btabsmall}, where that threshold is
  computed; v1.5.1, which separated the a-priori criterion from the measured threshold
  --- v1.5.0 quoted the measured $4580$ as though it followed from the criterion, while
  clause~(2) below derives its figure from the criterion).
  \textbf{The second criterion does not follow from the first and is the binding one
  first}: through v1.4.2 this clause read ``with the early-reduction scheme below this
  covers every $N<2^{15}$ outright by Lemma~\ref{lem:bounds}(5)'', which is false ---
  at $N=5003<2^{15}$ the certified $B=4N^{2}$ leaves $2B^{2}$ four orders of magnitude
  inside the tier while step~1's reduction leaves it (Lemma~\ref{lem:bounds}(6),
  Remark~\ref{rmk:btabsmall}).
\item \textbf{$128$-bit integers suffice} for (i)--(iii) iff $B\le2^{63}$, and for (iv)
  iff $7B_{\mathrm{red}}^{2}<2^{127}$, i.e.\ on the shipped model $N\lesssim1.8\cdot10^{8}$.
\item \textbf{Arbitrary precision is never required.} For the target regimes $N<2^{32}$
  and $N<2^{64}$, fixed-width $192$-bit, respectively $320$-bit, integers bound all
  magnitudes a priori ($2B^{2}\le2^{1+2(2\log_2N+\log_2C_K)}$); and the implementation
  must be overflow-checked, so that a violation of the certified bound \textbf{aborts
  rather than corrupts}. Arbitrary-precision arithmetic must not be linked into the
  pipeline: by this clause it is unreachable, and its absence is itself a guard against
  a silent performance cliff.
\item \textbf{The stored wide tier} (v1.23.0-pre.9; slot-24 ELL-W). If $B>2^{31}-1$, clause~(1)
  no longer covers (i) a priori. The pipeline then keeps the \emph{stored} matrices in signed
  $64$-bit integers and forms each product of (i) on two stored-class operands, and each
  determinant of a stored matrix, in $128$-bit integers --- attempting the $64$-bit evaluation
  first and escalating only after its overflow check fires, so that every value the $64$-bit tier
  computes is computed by the same operations --- and narrows the result back with a check.
  Under $2NB\le2^{63}-1$ (Lemma~\ref{lem:bounds}(7)) the intermediates of (i) on two stored-class
  operands are exact in $128$ bits (7b), (ii) on inputs of norm $\le B$ and the unimodular inversions
  are $64$-bit (7a), the level-one consumers' products of a fixed datum with a stored matrix are
  $64$-bit (7c), (iii) is unchanged --- its inputs are step-1 quadratic points, bounded through (iv)
  --- and (iv) is clause~(2) unchanged. This is a sufficient condition, not a characterisation:
  the results of (7b), the (P2) input $\gamma_a^{-1}\gamma_b$ and the edge prefix of step~2, the (T1)
  conjugate, and the level-structure consumers are formed in overflow-checked arithmetic whose
  completion is \COMP{} measured and not proved here (Lemma~\ref{lem:bounds}(7), last paragraph;
  Remark~\ref{rmk:btabsmall}) --- as for (iv), an exceeding operation aborts, it never wraps. On this model,
  with $B=4N^{2}$, that is $N\le2^{20}-1$, below (iv)'s $N\lesssim1.8\cdot10^{8}$, so it is the
  binding constraint of the $\tau$-certification. A \emph{stored} $128$-bit type would not extend
  it: a $128$-bit product of two stored matrices needs $B<2^{63}$, which $64$-bit storage already
  holds; the next step past $N=2^{20}$ is $128$-bit consumer products, not a wider storage.
\end{enumerate}

\noindent\textbf{One predicate needs the wide tier, and says so.} The nearest-vertex
comparison of Algorithm~\ref{alg:tau}(1b) compares two squared \emph{Euclidean} distances
in the coordinates of Lemma~\ref{lem:predshapes}(Q1), on cleared denominators --- a
heuristic choice, which is why the metric it uses is immaterial and only
Algorithm~\ref{alg:tau}(4) disposes; the comparison is formed in the $128$-bit tier, reduced by the
$\gcd$ of its two coefficients, and only then handed to the $64$-bit predicate of
\S\ref{sec:fpfilter}. If it still exceeds the $64$-bit tier after that reduction the
implementation \textbf{aborts}. This is the one place where (iii)'s bound
$2B^{2}$ is on the \emph{predicate's own} inputs rather than on the pipeline's stored
matrices, and it is the reason clause (iii) is phrased that way.

\medskip
\noindent\textbf{Early-reduction scheme (mandatory).} The pipeline needs exact integers
only inside (P2), (P3) and the predicates. All \emph{consumers} of a group element need
it only modulo $\ell$ (for $\rhon$) or modulo $N$ (for the class map (P1)). The
implementation must therefore carry each group element as the triple
\[
  \bigl(\text{exact }\Mat_2(\Z)\ \text{(only while needed)},\ \ \gamma\bmod\ell,\ \
  \gamma\bmod N\bigr)
\]
and reduce eagerly, dropping the exact component at its last use; all recurring
arithmetic is then single-word.

\medskip
\noindent\textbf{Mod-$\ell$ backend.} A modular product $(ab)\bmod\ell$ computed through
a $64$-bit product is exact iff $(\ell-1)^{2}\le2^{64}-1$, i.e.\ for every
$\ell<2^{32}$; no floating-point and no probabilistic arithmetic is involved. For
$2^{32}\le\ell<2^{64}$ this document requires either $128$-bit products or a
Barrett/Montgomery reduction with proved parameters. Fermat exponent reduction
$\exp\bmod(\ell-1)$ remains valid for prime $\ell$ with bases reduced mod $\ell$; note
that it sends $\exp\equiv0$ with base $\not\equiv0$ to $1$, which is correct.
\end{theorem}

\begin{proof}
(i)--(iii) are the closure statement of Remark~\ref{rmk:primitiveclosure} together with
Lemma~\ref{lem:bounds}: every stored matrix has norm $\le B$ by
Lemma~\ref{lem:bounds}(2),(3),(5) and the definition of $B$; a product of two of them has
magnitude $\le2B^{2}$ by Lemma~\ref{lem:bounds}(1); and the predicates of
\S\ref{sec:fpfilter} evaluate $a_1b_1\pm a_2b_2$ on inputs of magnitude $\le B$, hence
of magnitude $\le2B^{2}$. Clauses (1)--(3) are then the representable ranges of the
respective fixed-width signed types. The early-reduction scheme is well defined because
$\rhon$ factors through $\gamma\bmod\ell$ (Definition~\ref{def:rhon} applied to a matrix
with entries reduced mod $\ell$) and the class map through $\gamma\bmod N$
(Definition~\ref{def:cosetreps}); no other consumer of a group element exists, by
Remark~\ref{rmk:primitiveclosure}. The mod-$\ell$ clause is the exactness condition for a
product of two residues in a $64$-bit word. Clause~(4) is Lemma~\ref{lem:bounds}(7) together with
the order of evaluation: the $128$-bit form is reached only after the $64$-bit form's overflow check
fired, and both are exact, so wherever the $64$-bit form completes the result is that form's.
\end{proof}

\begin{contractbox}
\textbf{Separation of concerns (binding).} Floating point may be used \emph{only} inside
the geometric predicates of \S\ref{sec:fpfilter}, and only through the certified filter.
All linear algebra over $\cring$ --- cochains, elimination, projectors, Hecke matrices
--- and all coset, word and translate bookkeeping over $\Z$ and $\Z/N$ is exact integer
arithmetic. A build setting must exist that disables the floating-point fast path
entirely (\emph{pure-integer mode}) for differential testing, and both settings must
produce bit-identical results (Obligation~\ref{obl:fp}).
\end{contractbox}

\paragraph{Code map.}
Definition~\ref{def:primitives}: (P1) \file{Hecke/HeckeCosetTables.h/.cpp}
(\code{class\_index}/\code{class\_index\_of}, the bottom-row class map; \code{perm},
the index action $j\cdot\gamma$; \code{cocycle}, $h_j(\gamma)$, with
\code{rho\_cocycle} its eagerly reduced consumer table); (P2)
\file{Hecke/SL2ZTessellationTables.cpp} \code{word\_in\_T\_S} /
\code{word\_in\_T\_S\_with\_bound}, with \code{evaluate\_word} as the inverse and the
Lemma~\ref{lem:bounds}(4) bound checked at every intermediate step; (P3) --- no call
site, by the clause above; (P4) \code{SL2ZTessellationTables::canonical\_key}.
Algorithm~\ref{alg:tau}: \file{Hecke/HeckeTauTable.h/.cpp} --- step~0
\code{eval\_tau\_on\_translate} / \code{try\_derive}, step~1 \code{construct\_vertex},
\code{reduce\_to\_fundamental\_domain}, \code{choose\_vertex} (the (1b) tie-break) and
\code{compare\_vertex\_distances}, step~2 \code{construct\_edge} with
\code{SL2ZTessellationTables::path\_template} (the tree-routed \code{Canonical} family
and the pre-v1.1.0 \code{LegacyBfs} family, both certified in its own constructor),
step~3 \code{solve\_triangle\_system} with \code{used\_mod\_ell\_solve()} as the recorded
regime, step~4 \code{certify()}, called by the constructor, which throws on a non-zero
result. The exact quadratic points of step~1 are \file{Hecke/QuadPoint.h/.cpp}, and the exact
vertex positions $\mathrm{pos}(v)$ --- with the truncation height of \ref{H4} fixed to
$Y=2$ --- are \code{SL2ZTessellationTables::vertex\_position}.
Lemma~\ref{lem:bounds}(5) and Theorem~\ref{thm:backend}: \file{Hecke/ExactInt.h/.cpp}
(\code{Mat22}, \code{Mat22W}, \code{checked\_mul}/\code{checked\_add}/\code{checked\_sub},
\code{HeckeOverflow}, and \code{GroupElt}, which \emph{is} the early-reduction triple),
with \code{HeckeTauTable::certified\_bound\_B\_tab()} and
\code{certified\_bound\_B()} the certified $B_{\mathrm{tab}}$ and $B$, and clause~(e) of
\code{certify()} the recount and the tier assertion $2B^{2}<2^{63}$. The mod-$\ell$
backend is \code{ModArithmetic::mod\_mul} and, since the engine work,
\code{engine::ModCtx}. Validations: \code{HeckeTessellationValidation} and
\code{HeckeTessellationSmokeValidation} ((P2), (P4), the path templates),
\code{HeckeCosetTablesValidation} and \code{HeckeCosetTablesSelfCheckValidation} ((P1)),
\code{HeckeTauTableValidation} and \code{HeckeTauTableSelfCheckValidation}
(Algorithm~\ref{alg:tau} and the bounds), \code{HeckeExactIntSmokeValidation} and
\code{HeckeExtendedRegimeValidation} (Theorem~\ref{thm:backend}). Since v1.5.2 (slot-3
WP-1): \file{Hecke/QuadPoint.h/.cpp} is \code{QuadPointT<I>} with
$I\in\{\code{int64\_t},\code{\_\_int128}\}$ (\code{QuadPoint}, \code{QuadPointW}) and
\code{HeckeTauTable::construct\_vertex} runs Algorithm~\ref{alg:tau} step~1 through
\code{hecke::escalate\_on\_overflow} --- the \code{int64} tier first, widened per
operation only after a \code{HeckeOverflow} (Lemma~\ref{lem:bounds}(6),
Theorem~\ref{thm:backend}(2),(iv)), \code{hecke::ScopedExactTier} forcing either tier;
\code{HeckeUEllReachValidation} (the forced-wide arm bit-identical to the \code{int64} tier
wherever that tier completes, and the forced-narrow negative control) and
\code{HeckeUEllReachSlow} (the coset tables at $\ell=43867$, and the refusal there by
clause~(e), code~$51$). Since v1.23.0-pre.9 (slot-24 ELL-W): Theorem~\ref{thm:backend}(4) is
\code{hecke::mul\_tiered}, \code{det\_tiered} and \code{inv\_unimodular\_tiered}
(\file{Hecke/ExactInt.h/.cpp}; the \code{int64} \code{Mat22} operation first, the \code{Mat22W} form
after a \code{HeckeOverflow}, the result narrowed with a check; \code{ScopedExactTier} forcing either
tier), used at the stored-class sites of \file{Hecke/HeckeTauTable.cpp} (the (T1) translate of
step~0, the vertex relation of step~2, clause~(a)'s determinant), \code{HeckeCosetTables::twisted\_conjugate}
and \code{SL2ZTessellationTables::word\_in\_T\_S}; clause~(e) of \code{certify()} records the tier
(\code{wide\_tier()}) and asserts $2NB\le2^{63}-1$ there (code~$52$); \textbf{one rollback lever,
value-neutral}: \code{HECKE\_TAU\_WIDE=0} (line \code{EFFECTIVE HeckeTauWide}) restores the
refusal with code~$51$ above $B=2^{31}-1$ and never changes a value at or below it. Gates:
\code{HeckeWideTauValidation} (the helpers with fault injection, forced-wide $=$ narrow on
$\ell\le5003$, the first wide prime $23173$, $U_{30011}$ at $n=10$ against $\tau(30011)$ from the
$q$-expansion of $\Delta$) and \code{HeckeWideTauSlowValidation} (forced-wide $=$ narrow at
$\ell=23167$, $T_{30011}$ modulo $13$); \code{HeckeUEllReachSlow}'s arms (9), (10) follow the lever.
The boundary datum and $\eta$ (\file{Hecke/HeckeBoundaryDatum.cpp}, \file{Hecke/HeckePairHomotopy.cpp})
keep the narrow assertion (code~$51$ above $23170$): the cone route is not widened.

\paragraph{Formalization note.}
Definition~\ref{def:primitives} and Algorithm~\ref{alg:tau} are constructive as written
(rule~\ref{A3}): (P2) is structural recursion on $\abs{c}$ with the decrease exhibited in
the proof of Lemma~\ref{lem:bounds}(4); step~3's termination is the only
non-structural argument, and it rests on acyclicity of $\Ku$ (\ref{H4}(ii)) plus the
finiteness of the support of a chain. Lemma~\ref{lem:bounds}(1)--(4) are arithmetic
inequalities over $\Z$ and need no library beyond \code{Int.lt} and \code{abs}; (5) is a
statement about an emitted table and is best formalised as a predicate on it, not as a
theorem about $C_K$ --- which is how Obligation~\ref{obl:backend} already states it.
Theorem~\ref{thm:backend}(1),(2) are representability facts about fixed-width integer
types; a formalization would carry them as hypotheses on the arithmetic model, and (3)
would become the statement that the abstract algorithm never leaves a stated interval.

%----------------------------------------------------------------------
\section{The certified floating-point filter}\label{sec:fpfilter}
%----------------------------------------------------------------------
% SOURCE:  hecke_action.tex:1339-1407 (RH)
% NOTATION: rows "t,j,s (running index)", cring of sec:notation-crosswalk
%           applied.  The unit roundoff is written as an explicit power of two
%           throughout and is given NO letter: epsilon is the augmentation
%           (\aug, notation ch. sec:notation-framework), and rule A5 forbids a
%           second reading.  This is a deliberate avoidance, not a new symbol.
% ERRATA:  E22 (App. A) -- cor:fp32's totality threshold B <= 2^10 does not
%          follow from cor:fptotal's criterion E(B) < 1/2, which it meets with
%          equality; the threshold is correct and is proved here by exact
%          representability.

\paragraph{What this chapter fixes.}
The reduction steps of Algorithm~\ref{alg:tau} make branch decisions on integer data.
This chapter fixes the \emph{shape} of every one of them, proves a threshold below which
a floating-point evaluation of that shape certifies its own sign, proves the two regimes
in which the exact fallback is provably unreachable, and fixes the priority order and the
admissible constants as a binding contract. Nothing in this chapter is a statement about
$\Ku$, $\tau$ or $\Heq{q}$: a predicate returns the sign of an integer, and by
Remark~\ref{rmk:certoverGeometry} even a \emph{wrong} sign could not survive
Algorithm~\ref{alg:tau}(4). The filter is therefore about speed and about the
\emph{absence} of unverified floating point, not about correctness of the Hecke operator.

\paragraph{Standing hypotheses.}
None on $\ell$, none on $p$, none on the model. The two theorems below are statements
about IEEE-754 arithmetic on integers, and they carry their floating-point hypotheses in
their own statements (rule~\ref{A1}), as (FP1)--(FP3).

\subsection{The predicate inventory}\label{sec:predshapes}

\begin{lemma}[Every branch decision is a two-term integer dot product]\label{lem:predshapes}
\PROVED{} Assume \ref{H4}. Every branch decision taken by Algorithm~\ref{alg:tau} and by
the primitives of Definition~\ref{def:primitives} is the sign of an integer of the form
\begin{equation}\label{eq:dot2}
  x\;=\;a_1b_1\;+\;s\,a_2b_2,\qquad a_1,b_1,a_2,b_2\in\Z,\quad s\in\{+1,-1\},
\end{equation}
and there are exactly four such decisions:
\begin{enumerate}
\item[\textbf{(D)}] \textbf{Orientation / determinant.} $\sgn(ad-bc)$ for
  $\smat{a}{b}{c}{d}\in\Mat_2(\Z)$, i.e.\ \eqref{eq:dot2} with $s=-1$.
\item[\textbf{(Q1)}] \textbf{$\abs{\mathrm{Re}\,z}$ against $\tfrac12$.} For a point
  $z=(u+v\sqrt{-D})/w$ of $\Hyp$ with $u,v,w\in\Z$, $v,w>0$ and $D\in\{1,3\}$ ---
  the exact quadratic points in which the model's vertices and all their translates
  live --- the test $\abs{\mathrm{Re}\,z}\lessgtr\tfrac12$ is $\sgn(2\abs{u}-w)$, i.e.\
  \eqref{eq:dot2} with $(a_1,b_1,a_2,b_2,s)=(2,\abs{u},1,w,-1)$.
\item[\textbf{(Q2)}] \textbf{$\abs{z}$ against $1$.} The test $\abs{z}\lessgtr1$ is
  $\sgn(u^2+Dv^2-w^2)=\sgn\bigl((u-w)(u+w)+(Dv)\,v\bigr)$, i.e.\ \eqref{eq:dot2} with
  $s=+1$.
\item[\textbf{(Q3)}] \textbf{The nearest-vertex comparison.} The difference of two
  squared Euclidean distances on cleared denominators --- $\abs{z-P}^{2}(ww')^{2}$ for
  $z=(u+v\sqrt{-D})/w$ and $P=(u'+v'\sqrt{-D'})/w'$ --- is rational when the two points
  share their $D$ and has the form $A+B\sqrt3$ with $A,B\in\Z$ when they do not
  ($\{D,D'\}=\{1,3\}$). Its sign is decided by the signs of $A$ and $B$ when they agree
  or one vanishes, and otherwise by $\sgn(A^2-3B^2)$, i.e.\ \eqref{eq:dot2} with
  $(a_1,b_1,a_2,b_2,s)=(A,A,3B,B,-1)$; the case $A^2=3B^2$ with $(A,B)\neq(0,0)$ does
  not occur, $\sqrt3$ being irrational.
\end{enumerate}
No other branch decision on integer data occurs: the class map (P1) is arithmetic modulo
$N$, the quotient choice in (P2) and in the reduction is an integer division, and the
linear solves of Algorithm~\ref{alg:tau}(3) are exact eliminations with divisibility
tests. In particular, \textbf{no comparison anywhere in this document is made between
floating-point numbers that are not of the form \eqref{eq:dot2}}.
\end{lemma}

\begin{proof}
(D), (Q1), (Q2) are the displayed identities, valid over $\Z$; for (Q2) expand
$(u-w)(u+w)+Dv^2=u^2-w^2+Dv^2$. (Q1) and (Q2) are the two containment tests defining
$\mathcal{F}$, cleared of denominators; the class of points is closed under the M\"obius
action of a matrix of positive determinant, since for $M=\smat{a}{b}{c}{d}$ with
$\det M>0$ and $r:=au+bw$, $t:=cu+dw$ one has
\[
  M\cdot z=\frac{\bigl(rt+D\,ac\,v^2\bigr)+\bigl(vw\det M\bigr)\sqrt{-D}}
                {t^2+D\,c^2v^2},
\]
with integer numerator and positive denominator (the denominator vanishes only if
$t=cv=0$, which forces $\det M=0$). For (Q3), the squared Euclidean distance from $z$ to a point
$P=(u'+v'\sqrt{-D'})/w'$, cleared by $(ww')^2$, is
$(uw'-u'w)^2+Dv^2w'^2+D'v'^2w^2-2vv'ww'\sqrt{DD'}$; subtracting the corresponding
quantity for a second vertex, after scaling both by the other's $w^2$, leaves
$A+B\sqrt{DD'}$ with $A,B\in\Z$, and $DD'\in\{1,3,9\}$, so $\sqrt{DD'}\in\{1,3,\sqrt3\}$
and the only irrational case is $\{D,D'\}=\{1,3\}$. If $A,B$ have the same strict sign
the sign is that sign; if one vanishes it is the other's; otherwise
$\sgn(A+B\sqrt3)=\sgn(A)\cdot\sgn(A^2-3B^2)$. The enumeration is exhaustive by inspection
of Algorithm~\ref{alg:tau} and Definition~\ref{def:primitives}.
\end{proof}

\subsection{The filter}\label{sec:filterthm}

\begin{theorem}[FP64 filter]\label{thm:fp64}
\PROVED{} Let $a_1,b_1,a_2,b_2\in\Z$ with $\abs{a_1},\abs{b_1},\abs{a_2},\abs{b_2}\le B$
for an integer $B\le2^{53}$, let $s\in\{+1,-1\}$, and let $x=a_1b_1+s\,a_2b_2$, so
$x\in\Z$ and $\abs{x}\le2B^2$. Let $\hat x$ be computed in IEEE-754 binary64 under the
hypotheses
\begin{enumerate}
\item[(FP1)] the four inputs are converted to binary64 exactly --- which holds because
  $B\le2^{53}$;
\item[(FP2)] the rounding attribute is round-to-nearest, so that each elementary
  operation satisfies the standard model $\fl(\alpha\odot\beta)=
  (\alpha\odot\beta)(1+\theta)$ with $\abs{\theta}\le2^{-53}$;
\item[(FP3)] the evaluation order is exactly
  $\hat x=\fl\bigl(\fl(a_1b_1)+s\,\fl(a_2b_2)\bigr)$, with one rounding per step --- no
  fused multiply-add, no contraction of the two operations into one, and no wider
  intermediate precision.
\end{enumerate}
Then
\begin{equation}\label{eq:fperr}
  \abs{\hat x-x}\;\le\;4B^2\cdot2^{-53}\bigl(1+2^{-53}\bigr)\;<\;E(B)
  \qquad\text{where}\qquad E(B):=2^{-50}B^2 ,
\end{equation}
and in fact $\abs{\hat x-x}<\tfrac{1}{2}E(B)\bigl(1+2^{-53}\bigr)$. Consequently:
\begin{enumerate}
\item if $\abs{\hat x}>E(B)$ then $\sgn(\hat x)=\sgn(x)$, so the floating-point verdict
  is \emph{certified} and no exact evaluation is needed;
\item otherwise nothing is claimed, and the sign must be obtained by the exact
  bounded-integer evaluation of Theorem~\ref{thm:backend}.
\end{enumerate}
\end{theorem}

\begin{proof}
By (FP1) the products $a_1b_1$ and $a_2b_2$ are products of exactly represented
operands, so by (FP2) $\fl(a_tb_t)=a_tb_t(1+\theta_t)$ with $\abs{\theta_t}\le2^{-53}$
for $t=1,2$, and by (FP3) the sum rounds once, $\hat
x=\bigl(a_1b_1(1+\theta_1)+s\,a_2b_2(1+\theta_2)\bigr)(1+\theta_3)$ with
$\abs{\theta_3}\le2^{-53}$. Write $\eta:=a_1b_1\theta_1+s\,a_2b_2\theta_2$, so that
$\hat x=(x+\eta)(1+\theta_3)$ and $\abs{\eta}\le2B^2\cdot2^{-53}$. Then
\[
  \abs{\hat x-x}\;=\;\abs{\eta+ (x+\eta)\theta_3}
  \;\le\;\abs{\eta}+\bigl(\abs{x}+\abs{\eta}\bigr)2^{-53}
  \;\le\;2B^2 2^{-53}+2B^2 2^{-53}+2B^2 2^{-106}
  \;=\;4B^2 2^{-53}\bigl(1+2^{-53}\bigr),
\]
which is \eqref{eq:fperr}; and $4B^2 2^{-53}=\tfrac12\cdot2^{-50}B^2=\tfrac12 E(B)$,
giving the sharper form. For (1): if $\abs{\hat x}>E(B)$ then
$\abs{\hat x}>E(B)>\abs{\hat x-x}$, so $x$ and $\hat x$ cannot have opposite signs and
$x\neq0$; hence $\sgn(x)=\sgn(\hat x)$.
(The design --- a cheap filter with a proved threshold and an exact fallback --- is
Shewchuk's \cite{shewchuk}; the bound above is elementary and is proved here rather than
cited.)
\end{proof}

\begin{remark}[The threshold is itself computed in floating point, and that is sound]
\label{rmk:thresholdrounding}
\PROVED{} An implementation cannot form $E(B)=2^{-50}B^2$ exactly for every $B\le2^{53}$;
it computes $\widehat E:=\fl\bigl(\fl(B\cdot B)\cdot2^{-50}\bigr)$. Conversion of $B$ is
exact, the scaling by a power of two is exact (and cannot underflow for $B\ge1$), so
$\widehat E=E(B)(1+\theta)$ with $\abs{\theta}\le2^{-53}$. The certified test
$\abs{\hat x}>\widehat E$ is therefore still sound: $\widehat E\ge
E(B)(1-2^{-53})>\tfrac12E(B)(1+2^{-53})\ge\abs{\hat x-x}$ by \eqref{eq:fperr}, which is
all the proof of Theorem~\ref{thm:fp64}(1) used. \textbf{The factor $2$ of slack between
the proved error bound and the normative constant $E(B)$ exists exactly for this
reason}, and is why $E(B)=2^{-50}B^2$ --- and not the sharper $2^{-51}B^2$ --- is the
binding constant.
\end{remark}

\begin{corollary}[Totality threshold]\label{cor:fptotal}
\PROVED{} In the situation of Theorem~\ref{thm:fp64}, suppose in addition that
\begin{equation}\label{eq:fptotalcrit}
  E(B)+\abs{\hat x-x}\;<\;1 ,
\end{equation}
for which each of the following is sufficient: (i) $E(B)\le\tfrac12$ ---
\textbf{not} $E(B)<\tfrac12$, see Appendix~\ref{app:errata}, entry~\ref{err:e22}; (ii)
$2B^2\le2^{53}$ and $E(B)<1$. Then the filter is \emph{total}: $\abs{\hat x}>E(B)$
certifies $\sgn(x)=\sgn(\hat x)\ne0$ as before, and $\abs{\hat x}\le E(B)$ certifies
$x=0$. The exact fallback is then provably unreachable. Both sufficient conditions hold
for
\[
  B\;\le\;2^{24}\qquad(\text{then }E(B)\le2^{-2}\text{ and }2B^2\le2^{49}),
\]
which is the threshold this document fixes and an implementation must use. Under (ii)
one has moreover $\hat x=x$ exactly.
\end{corollary}

\begin{proof}
$x\in\Z$, so $x\neq0$ implies $\abs{x}\ge1$ and hence
$\abs{\hat x}\ge\abs{x}-\abs{\hat x-x}\ge1-\abs{\hat x-x}>E(B)$ by
\eqref{eq:fptotalcrit}; contrapositively $\abs{\hat x}\le E(B)$ forces $x=0$.
(i) implies \eqref{eq:fptotalcrit} because $\abs{\hat x-x}\le\tfrac12E(B)(1+2^{-53})$
by \eqref{eq:fperr}, so the left side is at most $E(B)\bigl(\tfrac32+2^{-54}\bigr)$,
which for $E(B)\le\tfrac12$ is at most $\tfrac34+2^{-55}<1$. For (ii): $2B^2\le2^{53}$ makes $\abs{a_tb_t}\le B^2\le2^{52}$
and $\abs{x}\le2B^2\le2^{53}$ integers exactly representable in binary64, so under
(FP1)--(FP3) every rounding is exact and $\hat x=x$, whence $\abs{\hat x-x}=0$ and
\eqref{eq:fptotalcrit} reduces to $E(B)<1$. At $B=2^{24}$, $E(B)=2^{-50}2^{48}=2^{-2}$
and $2B^2=2^{49}$.
\end{proof}

\begin{corollary}[FP32 analogue]\label{cor:fp32}
\PROVED{} The statements of Theorem~\ref{thm:fp64} and Corollary~\ref{cor:fptotal} hold
verbatim in IEEE-754 binary32, with $2^{-53}$ replaced by $2^{-24}$ in (FP2), with
(FP1) requiring $B\le2^{24}$, with the constant
\[
  E_{32}(B):=2^{-21}B^2,\qquad
  \abs{\hat x-x}\le4B^2\cdot2^{-24}(1+2^{-24})=\tfrac12E_{32}(B)(1+2^{-24}),
\]
and with the totality threshold
\[
  B\;\le\;2^{10}\qquad(\text{then }E_{32}(B)=\tfrac12<1\text{ and }2B^2=2^{21}\le2^{24}),
\]
which is certified by Corollary~\ref{cor:fptotal}(i) in its corrected, non-strict
form $E_{32}(B)\le\tfrac12$, and independently by
Corollary~\ref{cor:fptotal}(ii) --- the evaluation is \emph{exact} there. It is
\textbf{not} certified by the \emph{strict} criterion ``$E(B)<\tfrac12$'' of the
superseded contract, which $B=2^{10}$ fails with equality; that document nevertheless
derived this threshold from it, and so did not prove it
(Appendix~\ref{app:errata}, entry~\ref{err:e22}). The binary32 path is provided for
future device inner loops on the fixed model constants and is otherwise not used:
$B\le2^{10}$ covers only very small data, and Theorem~\ref{thm:fp64} is preferred
everywhere else.
\end{corollary}

\begin{proof}
Identical to the proofs above with $2^{-24}$ for $2^{-53}$; $4B^2 2^{-24}=\tfrac12\cdot
2^{-21}B^2$. For the threshold, $B=2^{10}$ gives $\abs{a_tb_t}\le2^{20}$ and
$\abs{x}\le2^{21}$, all integers of magnitude $\le2^{24}$ and therefore exactly
representable in binary32, so $\hat x=x$ and Corollary~\ref{cor:fptotal}(ii) applies
with $E_{32}(2^{10})=\tfrac12<1$; criterion (i) applies as well, with equality.
\end{proof}

\begin{contractbox}
\textbf{Priority order (binding).} (1) the FP64 filtered predicate of
Theorem~\ref{thm:fp64}; (2) exact bounded-integer evaluation, $64$-bit then $128$-bit per
Theorem~\ref{thm:backend}; (3) fixed-width wide integers.
\textbf{Arbitrary-precision arithmetic must not be linked into the pipeline}; by
Theorem~\ref{thm:backend}(3) it is unreachable in the target regime, and its absence is
itself a guard against a silent performance cliff.
\textbf{The only admissible thresholds} are $E(B)=2^{-50}B^2$ (binary64) and
$E_{32}(B)=2^{-21}B^2$ (binary32), evaluated with the \emph{certified} $B$ of
Lemma~\ref{lem:bounds}(5), with the \emph{strict} inequality $\abs{\hat x}>E(B)$, and
with the evaluation order (FP3). Each of those four --- constant, certified $B$, strict
inequality, evaluation order --- is load-bearing in the proof above, and none may be
changed as an optimisation.
\textbf{A preconditions violation aborts.} $s\notin\{\pm1\}$, $B$ above the
exact-conversion limit, or an input of magnitude exceeding $B$ leaves the verdict
uncertified; the implementation must throw, never return
(Theorem~\ref{thm:backend}(3), Obligation~\ref{obl:backend}).
\textbf{A pure-integer build must exist} in which the floating-point path is compiled out
entirely, and the two builds must produce bit-identical results
(Obligation~\ref{obl:fp}).
\end{contractbox}

\paragraph{Code map.}
Theorem~\ref{thm:fp64}: \file{Hecke/Predicates.h/.cpp} \code{sign\_dot2}, whose
signature $(a_1,b_1,a_2,b_2,s,B)$ is \eqref{eq:dot2} verbatim; (FP3) is enforced by
\code{volatile} intermediates, (FP1) and the magnitude preconditions by
\code{check\_dot2\_preconditions}, which throws \code{HeckeOverflow}; the threshold is
\code{std::ldexp(dB*dB, -50)}, which is Remark~\ref{rmk:thresholdrounding}'s
$\widehat E$; the exact fallback is \code{sign\_dot2\_exact}, the
\code{int64}$\to$\code{\_\_int128} ladder of Theorem~\ref{thm:backend}.
Corollary~\ref{cor:fptotal}: the branch \code{B <= (1 << 24)}.
Corollary~\ref{cor:fp32}: \code{sign\_dot2\_f32} with constant $2^{-21}$ and branch
\code{B <= (1 << 10)}.
Lemma~\ref{lem:predshapes}: (D) is \code{sign\_det2}; (Q1) is
\code{QuadPoint::cmp\_abs\_re\_half}; (Q2) is \code{QuadPoint::cmp\_norm\_one};
(Q3) is \code{QuadPoint::sign\_lin\_sqrt3}, fed by
\code{HeckeTauTable::compare\_vertex\_distances}, which forms $A,B$ in the $128$-bit
tier, divides by $\gcd(A,B)$ and aborts if the result still exceeds the $64$-bit tier
(Theorem~\ref{thm:backend}, the wide-tier clause). Counters \code{fp\_hits} /
\code{exact\_fallbacks} (\code{PredicateStats}) are the \emph{diagnostic} that
Corollary~\ref{cor:fptotal} is observed, and are never branched on (rule~\ref{A6}).
The build setting is the CMake option \code{HECKE\_PURE\_INTEGER}, which compiles out
both floating-point paths by \code{\#ifndef}, so that every call becomes an
\code{exact\_fallback}.
Validation: \code{HeckeFpFilterValidation} (Obligation~\ref{obl:fp}).

\paragraph{Formalization note.}
Lemma~\ref{lem:predshapes} is a finite case analysis over $\Z$ and needs no
floating-point model at all; it is the statement worth formalising first, because it is
what makes the rest of the chapter a self-contained lemma about IEEE-754 rather than a
claim about the pipeline. Theorem~\ref{thm:fp64} and its corollaries are statements in
the standard model (FP2), which a formalization carries as a hypothesis on an abstract
rounding function rather than as a fact about hardware; mathlib4's
\code{FloatingPoint} development is not needed, since (FP2) plus exact representability
is all the proofs use. Corollary~\ref{cor:fptotal}(ii) is the clause a formalization
should prefer: it replaces an error estimate by the statement that every intermediate is
an integer of bounded magnitude, which is decidable.

%----------------------------------------------------------------------
\section{Correctness obligations}\label{sec:obligations}
%----------------------------------------------------------------------
% SOURCE:  hecke_action.tex:1408-1511 (RH+NEW)
% NOTATION: rows nu, T_N, "t,j,s (running index)", Rel, cring, Mat_2(Z) of
%           sec:notation-crosswalk applied.
% ERRATA:  E1 (obl:algebra, the primitive-vs-full prescription -- the second of
%          its two landing sites), E6 (the coefficient p^{n+1}), E16 (T_N).
%          App. A.

\paragraph{What this chapter fixes.}
The implementation is entitled to rely on the numbered statements of this document; in
exchange it must \emph{discharge} the obligations below. This chapter is the complete
roster: the eleven obligations of the superseded contract, re-homed and repaired, the
remark on the exceptional loci that two of them are gated by, and --- integrated here,
not merely listed --- the four obligations that Chapters~\ref{sec:involution},
\ref{sec:boundarydatum} and \ref{sec:cone} added. It closes with the table
(\S\ref{sec:obligationroster}) that says, for each one, what ring it is evaluated over,
whether a failure \textbf{aborts} or is \textbf{recorded}, and whether a
\emph{registered} validation stands behind it.

\paragraph{How to read an obligation (rule~\ref{A4}).}
An obligation is a checkable predicate on emitted data, not a wish. Each clause names
\emph{(i)} the stored data it reads, \emph{(ii)} the predicate, \emph{(iii)} the ring the
predicate is evaluated over, and \emph{(iv)} what happens on failure. Two failure actions
exist and they are not interchangeable:
\begin{itemize}
\item \textbf{aborts} --- the computation stops and emits nothing. An obligation clause
  that aborts is an assertion this document makes and stakes the run on.
\item \textbf{recorded} --- the value is written into the emitted record and the run
  continues. A recorded quantity is \emph{never} branched on (rule~\ref{A6}): it may not
  select a code path, skip another clause, or be reported as a pass.
\end{itemize}
A clause is \textbf{gated} when it aborts under a stated hypothesis and is recorded
otherwise; every gate in this chapter is $6\in\cring^{\times}$ (\ref{H3}), $p\neq\ell$
(\ref{H1}) or ``off the exceptional locus of Remark~\ref{rmk:dickson}'', and each is
written into its own clause.

\paragraph{Standing hypotheses.}
\ref{H4} throughout, and the datum $g=\smat{1}{0}{0}{N}$, $N=p^{\nu}$, $\Ga''=\Ga_0(N)$,
$\{s_j\}$, $\tau$ of \S\S\ref{sec:heckedatum}--\ref{sec:chainlevel}. Per rule~\ref{A1}
each obligation repeats the hypotheses it needs; in particular \textbf{no obligation is
under a blanket $\ell\ge5$}, and the three that need it say so in the clause that needs
it.

\subsection{The core obligations}\label{sec:coreobligations}

\begin{obligation}[Well-definedness in all degrees]\label{obl:welldef}
Assume \ref{H4}; \textbf{no hypothesis on $\ell$}. Reading the stored $\tau$-table, the
coset tables, and the stored bases and projector of \S\ref{sec:systematiccomputation},
an implementation of $\Th_g$ (Definition~\ref{def:xi}) must discharge:
\begin{enumerate}
\item[(i)] \textbf{Equivariance on emitted data.} For every stored cohomology-basis
  cocycle $z^{(q)}_j$, the image $y_j:=\Th_g z^{(q)}_j$ satisfies \emph{every} relation
  of $\Rel$: $y_j(\sigma_{\mathrm{lhs}})=\rhon(\gamma)\,y_j(\sigma_{\mathrm{rhs}})$ for
  every row $(\sigma_{\mathrm{lhs}},\sigma_{\mathrm{rhs}},\gamma)\in\Rel$, i.e.\
  $y_j\in\Cq{q}(K,\Rel;L)$. Over $\cring$; a failure \textbf{aborts}.
\item[(ii)] \textbf{Cocycle.} $d\,y_j=0$ with $d$ the coboundary of
  \S\ref{sec:explicitcochaincomplex}. Over $\cring$; a failure \textbf{aborts}.
\item[(iii)] \textbf{Restriction to $\Cq{q}$ is exactly solvable.} Each
  $\Th_g$-image of a stored $\Cq{q}$-basis vector lies in the span of that basis, the
  solve that defines $\Th^{(q)}_{\mathrm{eq}}$ succeeds, and its residual, recomputed
  independently, is identically zero. Over $\cring$; a failure \textbf{aborts}.
\end{enumerate}
Clauses (i)--(iii) are the runtime form of Theorems~\ref{thm:welldef} and \ref{thm:xi}:
$\Th_g$ maps $\Cq{q}(K,\Rel;L)$ to itself, commutes with $d$, and induces $\Tdc{g}$ on
$\Heq{q}(K,\Rel;L)$ for $q=0,1,2$ \emph{simultaneously}. They are the three assertions
that Algorithm~\ref{alg:assembly}'s contract box delegates here.
\end{obligation}

\begin{obligation}[Independence of all auxiliary choices]\label{obl:indep}
Assume \ref{H4}. The induced map on $\Heq{q}$ is independent of the choices listed in
Theorem~\ref{thm:welldef}(a)--(c) and Proposition~\ref{prop:omega}. \textbf{Every
comparison is made on $\Amat{q}$, never on $\Th^{(q)}_{\mathrm{flat}}$}
(Remark~\ref{rmk:indeplevel}: a validation comparing cochain matrices across a varied
choice is testing a statement this document does not make).
\begin{enumerate}
\item[(i)] \textbf{Coset representatives.} Rebuild the whole chain --- coset tables,
  $\tau$-table, operator --- against $s_j'=h_js_j$ with $h_j\in\Ga''$, and compare
  $\Amat{q}$ entrywise in every degree. Over $\cring$; a failure \textbf{aborts}. With
  $h_j=I$ for every $j$ the rebuild must additionally reproduce
  $\Th^{(q)}_{\mathrm{flat}}$ bit for bit --- that is the same construction, not a
  varied choice, and a difference is a defect in the rebuild machinery. With
  non-trivial $h_j$ the cochain-level agreement is \textbf{recorded}, never asserted,
  because rebuilding the chain rebuilds $\tau$ and this clause then exercises
  Theorem~\ref{thm:welldef}(a) and (b) together.
\item[(ii)] \textbf{The twisted chain datum.} Rebuild $\tau$ under a different admissible
  vertex policy of Algorithm~\ref{alg:tau}(1) against the \emph{same} cosets and compare
  $\Amat{q}$ entrywise. Over $\cring$; a failure \textbf{aborts} \emph{when
  $6\in\cring^{\times}$}, and is \textbf{recorded} otherwise --- the uniqueness half
  Lemma~\ref{lem:tauexists}(2) is exactly where that hypothesis is used, and no
  cross-$\tau$ assertion is made at $\ell\in\{2,3\}$.
\item[(iii)] \textbf{The double-coset representative $g$.} Theorem~\ref{thm:welldef}(c).
  \textbf{This clause is not reachable through the present interface} and is discharged
  by its proof alone: the coset tables fix $g=\mathrm{diag}(1,N)$ and $\Ga''=\Ga_0(N)$,
  and admit only representatives in the same $\Ga''$-classes, so a genuinely different
  $g$, or a conjugate subgroup $\gamma_1\Ga_0(N)\gamma_1^{-1}$ with
  $\gamma_1\notin\Ga_0(N)$, cannot be built. What \emph{is} run --- replacing every
  $s_j$ by $\gamma s_j$ for one fixed $\gamma\in\Ga''$ --- is the constant-$h_j$ case of
  clause (i) and must be reported as such, not as an instance of (c). Making (iii)
  checkable requires an entry point parameterised by an arbitrary $g\in\De$ and an
  arbitrary finite-index $\Ga''\le\Ga$.
\item[(iv)] \textbf{Reduction words.} The words produced by (P2)/(P3) are functions of
  their input alone (Definition~\ref{def:primitives}), so this clause is discharged by
  determinism of the primitive and needs no run; a non-deterministic quotient rule would
  be a defect in (P2), not an independence failure.
\end{enumerate}
\end{obligation}

\begin{obligation}[Integrality and the two primes]\label{obl:integral}
Assume \ref{H4}; \textbf{no hypothesis on $\ell$}.
\begin{enumerate}
\item[(i)] \textbf{Every matrix is integral.} Every matrix entering the pipeline is
  $\rhon(\delta)\bmod\ell$ for an \emph{integral} $\delta\in\De$ with
  $\det\delta\in\{1,N\}$, read from the stored $\tau$- and coset tables. Over $\Z$ for
  the integrality, over $\cring$ for the reduction; a failure \textbf{aborts}.
\item[(ii)] \textbf{No determinant-$N$ matrix is ever inverted.} Only unimodular
  inversion occurs, and only on $s_j$, $s_{j\cdot\gamma}$ and $\gamma_{\mathrm{red}}$.
  Over $\Z$; a failure \textbf{aborts}. (Lemma~\ref{lem:detsym}: for $p=\ell$ and
  $n\ge1$ every $\rhon(\delta_j)\bmod\ell$ is singular, so an inversion would not merely
  be unjustified, it would fail.)
\item[(iii)] \textbf{The normalisation is the integral one.} $\Amat{0}(T_p)=1+p$ on the
  trivial coefficient system, and the recursion coefficient of
  Theorem~\ref{thm:relations}(1) is $p^{\,n+1}=p^{\,k-1}$ with $k=n+2$ --- neither
  symmetrised nor rescaled (Algorithm~\ref{alg:assembly}, ``no normalisation''). Over
  $\cring$; a failure \textbf{aborts}.
\item[(iv)] \textbf{The bad prime uses the same pipeline.} For $p=\ell$ (\ref{H2}) the
  operator $U_\ell$ is produced by the identical construction with no special case, and
  Corollary~\ref{cor:ulrecursion}(2)'s coincidence of the full and primitive prime-power
  operators at $n\ge1$ holds in the emitted matrices. Over $\cring$; a failure
  \textbf{aborts}. The recursion $T_{\ell^{\nu}}=U_\ell^{\nu}$ of
  Corollary~\ref{cor:ulrecursion}(1) carries $6\in\cring^{\times}$ and is
  \textbf{recorded}, not asserted, at $\ell\in\{2,3\}$.
\end{enumerate}
\end{obligation}

\begin{obligation}[Hecke algebra structure]\label{obl:algebra}
Assume \ref{H4} and $6\in\cring^{\times}$ --- for $\cring$ of characteristic $\ell$,
$\ell\ge5$ (hypothesis \ref{H3}). Then $\Ga\delta\Ga\mapsto\Tdc{\delta}$ extends to a
ring homomorphism $\mathcal{H}(\De,\Ga)\to\End\bigl(\Heq{q}(K,\Rel;L)\bigr)$ for each
$q$ (Theorem~\ref{thm:relations}; \STD{\cite[Thm.~3.24, \S3.2--3.3]{shimura} for the
identities in $\mathcal{H}(\De,\Ga)$ (verified, R20-L), \cite[\S1.1, p.~194]{ash_stevens} (the action, stated as well known) and
\cite[Ch.~III.9--10]{brown} (the transfer) for the computation that makes the assignment a
ring homomorphism}), whence the $\{T_N\}_N$ commute \LEANlit{obl:algebra-model}{the commutativity of the Hecke ring, \cite[Prop.~3.8, p.~54, applied in \S3.2, p.~56]{shimura} (there for $\det>0$; for $\det<0$ the same argument with a diagonal representative), transported by \cite[\S1.1, p.~194]{ash_stevens}, as the hypothesis \texttt{hcomm}; locator verified (R20-L)} (all $\Tdc{g}$, $g\in\De$, commute
and the $\cring$-algebra they generate is commutative, on the model; not the ring homomorphism) and satisfy
$T_{p^{\nu+1}}=T_pT_{p^{\nu}}-p^{\,n+1}T_{p^{\nu-1}}$ and $T_MT_{M'}=T_{MM'}$ for
$\gcd(M,M')=1$.

\medskip\noindent
\textbf{The runtime tests, in the normalisation in which each is true.} Everything below
is over $\cring$, entrywise on $\Amat{q}$, in every computed degree; a failure
\textbf{aborts} under $6\in\cring^{\times}$ and is \textbf{recorded} otherwise.
\begin{enumerate}
\item[(i)] \textbf{Commutation.} $[\Amat{q}(T_p),\Amat{q}(T_{p'})]=0$ for distinct
  primes $p,p'$ (the shipped instance is $p,p'=2,3$).
\item[(ii)] \textbf{The recursion at $\nu=2$, for the \emph{full} operator.}
  \[
    \Amat{q}\bigl(T_{p^{2}}\bigr)\;=\;\Amat{q}(T_p)^{2}\;-\;p^{\,n+1}I ,
  \]
  with $T_{p^{2}}$ the \textbf{full} operator of \eqref{eq:fullop} obtained
  \emph{independently of the recursion} --- by the double-coset assembly
  \eqref{eq:fullassembly}, not by the identity being tested, which would make the check
  vacuous. \LEANlit{lem:m4-hTp-hrel}{the primitive products of \cite[Thm.~3.24(5), p.~63, with Prop.~3.17, p.~58]{shimura}, transported by \cite[\S1.1, p.~194]{ash_stevens}, as the hypothesis \texttt{hTp}; locator verified (R20-L)} (the identity as algebra on the model, the full operator in the assembled form;
  the runtime test itself is the C++)
\item[(iii)] \textbf{The primitive operator, and the exact size of the difference.}
  The raw pipeline --- coset tables, $\tau$-table, operator --- run at $\nu=2$ emits the
  \textbf{primitive} double-coset operator $T_{\Deprim_{p^{2}}}$, whose $p^{2}+p$ cosets
  are what Definition~\ref{def:cosetreps} enumerates, and it satisfies
  \[
    \Amat{q}\bigl(T_{\Deprim_{p^{2}}}\bigr)\;=\;\Amat{q}(T_p)^{2}\;-\;p^{\,n+1}I\;-\;p^{\,n}I ,
    \qquad\text{equivalently}\qquad
    \Amat{q}\bigl(T_{p^{2}}\bigr)=\Amat{q}\bigl(T_{\Deprim_{p^{2}}}\bigr)+p^{\,n}I ,
  \]
  the difference being the diamond scalar $\diam=p^{\,n}I$ of
  Proposition~\ref{prop:diamond}. \LEANlit{lem:m4-hTp-hrel}{the primitive products of \cite[Thm.~3.24(5), p.~63, with Prop.~3.17, p.~58]{shimura}, transported by \cite[\S1.1, p.~194]{ash_stevens}, as the hypothesis \texttt{hTp}; locator verified (R20-L)} (both identities as algebra on the model, the second
  with no hypothesis; the runtime test itself is the C++) \textbf{Both identities are asserted, and the emitting
  entry point of each is named in the record}; the two are \emph{not} interchangeable
  and substituting either operator into the other's identity falsifies it by $p^{\,n}I$.
\item[(iv)] \textbf{Diamond.} $\Amat{q}(\diam)=p^{\,n}I$ (Proposition~\ref{prop:diamond},
  which carries no hypothesis on $\ell$, so this clause aborts for every $\ell$).
\item[(v)] \textbf{Coprime multiplicativity.} $\Amat{q}(T_{MM'})=\Amat{q}(T_M)\Amat{q}(T_{M'})$
  for $\gcd(M,M')=1$ (the shipped instance is $T_6=T_2T_3$).
\end{enumerate}

\begin{erratumbox}
Clauses (ii) and (iii) are \textbf{erratum E1} (Appendix~\ref{app:errata},
entry~\ref{err:e1}). The superseded contract stated this obligation's runtime test as
``compute $T_{\Ga\,\mathrm{diag}(1,p^{2})\,\Ga}$ by the same pipeline with its $p^{2}+p$
cosets and compare with $T_p^{2}-p^{\,n+1}$'': a $p^{2}+p$-coset operator is the
\emph{primitive} one and $T_p^{2}-p^{\,n+1}I$ is the \emph{full} one, so the prescription
as written fails against a correct implementation, by exactly $p^{\,n}I$. \S\ref{sec:primfull}
names the two operators and \eqref{eq:fullassembly} relates them; entry~\ref{err:e1}
carries the evidence, the entry-point table and the runs, and is not restated here.
\end{erratumbox}
\end{obligation}

\begin{obligation}[Agreement of the three formulations]\label{obl:agreement}
Assume \ref{H4}; \textbf{no hypothesis on $\ell$}. The three routes of
\S\S\ref{sec:induced}--\ref{sec:equivalence} must agree, and the level at which each
agreement is asserted is part of the obligation (Remark~\ref{rmk:indeplevel}).
\begin{enumerate}
\item[(i)] \textbf{Bit-identically on cochains, on the same datum.} The fused route
  $\Th_g=\omega_{\bullet}\circ\Xig$ and the direct evaluation of \eqref{eq:heckechain}
  (``route 2''), built from the \emph{same} $g$, the \emph{same} $s_j$ and the
  \emph{same} $\tau$-table, agree on every entry of $\Th^{(q)}_{\mathrm{flat}}$ in every
  degree. Over $\cring$; a failure \textbf{aborts}. Likewise
  $\omega_{\bullet}\circ\mu_{\bullet}\circ\alpha_{\bullet}$ against $\Th_g$
  (Proposition~\ref{prop:mu}(3)), same data, same level.
\item[(ii)] \textbf{Through the public evaluation interface.} Applying the operator to
  each standard basis vector $e_b$ of $C^q(K;\Vcoef)$ reproduces column $b$ of
  $\Th^{(q)}_{\mathrm{flat}}$ exactly, by both routes. Over $\cring$; a failure
  \textbf{aborts}. This clause exists because (i) can be discharged against private
  state; (ii) cannot.
\item[(iii)] \textbf{On $H^{1}$, the bar-cochain route.} The corrected degree-$1$
  group-cochain formula of Proposition~\ref{prop:barformula} --- \emph{with} the coset
  permutation $j\mapsto j\cdot\gamma$, whose omission is entry~\ref{err:e20} --- induces
  the same map on $\Heq{1}$. Over $\cring$, compared at the $\Amat{1}$ level; a failure
  \textbf{aborts}. It is an independent route and not a reformulation, so the comparison
  is on cohomology, not on cochains.
\item[(iv)] \textbf{Across the two arithmetic builds.} The pure-integer build and the
  filtered build produce bit-identical stored objects; see
  Obligation~\ref{obl:fp}, which owns that predicate.
\end{enumerate}
\end{obligation}

\begin{obligation}[Shapiro consistency of $L_H$]\label{obl:shapiro}
Assume \ref{H4}. $\dim_{\cring}\Heq{q}(K,\Rel;L_{\Ga''})$, computed by the equivariant
cohomology of the \emph{induced} local system of Definition~\ref{def:inducedls}, equals
$\dim_{\cring}\Heq{q}$ of the $\Ga_0(N)$-model with coefficients $\VR$, in every computed
degree (Proposition~\ref{prop:lhcompliant}). Over $\cring$; a failure \textbf{aborts}.
For $p=2$, $n=0$, $6\in\cring^{\times}$ the values are $\dim\Heq{0}=1$ and
$\dim\Heq{1}=1$, the latter from $\Ga_0(2)^{\mathrm{ab}}$, the model being the
Borel--Serre-type truncation with both cusps of $X_0(2)$ on the boundary circle covering
the cusp circle of the base. The identification of either side with
$H^{\bullet}(\Ga_0(N);\VR)$ needs $6\in\cring^{\times}$
(Remark~\ref{rmk:orbifold}); the \emph{equality of the two computed dimensions} does
not, and that is what this obligation asserts.
\end{obligation}

\begin{obligation}[Classical eigenvalue oracle]\label{obl:oracle}
Assume \ref{H4}, $6\in\cring^{\times}$ (so $\ell\ge5$ for $\cring=\Fl$), $p\neq\ell$
(\ref{H1}), $n$ even, and $(n,\ell)$ \textbf{off} the exceptional locus of
Remark~\ref{rmk:dickson}. Then, under Eichler--Shimura and reduction modulo $\ell$
\STD{\cite[Ch.~8]{shimura} for the Eichler--Shimura isomorphism over $\Cx$;
\cite[Ch.~2--3]{harder3} for the Eisenstein class and its eigenvalue; \cite{deng} for
the mod-$\ell$ statement}, the characteristic polynomial of $\Amat{1}(T_p)$ on
$\Heq{1}(K,\Rel;\Sym^{n}\Fl^{2})$ is
\begin{equation}\label{eq:esoracle}
  \Bigl(\prod_{f}\bigl(X-a_p(f)\bigr)^{2}\Bigr)\cdot\bigl(X-(1+p^{\,n+1})\bigr)
  \pmod\ell ,
\end{equation}
the product running over the Galois orbit of normalised cuspidal eigenforms $f$ of
level $1$ and weight $k=n+2$ --- each with multiplicity \emph{two}, from the holomorphic
and antiholomorphic summands of the Eichler--Shimura decomposition
(Claim~\ref{cl:dimH1}(1)) --- and the last factor the Eisenstein class, of boundary
eigenvalue $1+p^{\,n+1}=1+p^{\,k-1}$. In particular $\deg=2(d_k-1)+1=2d_k-1$, which is
$\dim\Heq{1}$ off the exceptional locus.

\medskip\noindent
\textbf{``The Eisenstein class'' modulo $\ell$ is a space, not a line, at a Kummer
cell} (v1.5.0; Appendix~\ref{app:changelog}, v1.4.2 deferred item~1). The \emph{factor}
$X-(1+p^{\,n+1})$ of \eqref{eq:esoracle} is well defined wherever this obligation
applies, and nothing above changes. The \emph{class} named by the second standard
pointer of this obligation --- \cite[Ch.~2--3]{harder3}, the Eisenstein class and its eigenvalue,
a characteristic-$0$ statement --- need not survive reduction as a line: where $\ell$
divides the numerator of the relevant Bernoulli quotient, the integral Eisenstein
eigenline and a cuspidal one reduce into the same socle of a non-split block of the
local Hecke algebra, and the joint $1+p^{\,n+1}$-eigenspace $E$ on $\Heq{1}$ is
two-dimensional ([E] v1.4.0 \code{cl:harder}; measured at all ten Kummer cells,
Remark~\ref{rmk:coneambiguity691}). \textbf{Wherever this document, or a record it
governs, speaks of \emph{the} Eisenstein class modulo $\ell$, read the joint eigenspace
$E$}; \eqref{eq:esoracle} singles out no line inside it, and
Obligation~\ref{obl:iota}(v) already states the two conditions --- the eigenspace a
line, or restriction to the cusp with $\rank\rres^{1}=1$ --- under which a class, rather
than a space, is what is being located. Off the Kummer loci $E$ is a line --- measured,
not proved here: on every off-Dickson non-Kummer cell with $k\le26$ ($43$ cells, slot-2
WP-7; the statement in general is [E]'s \code{cl:harder}) --- and the two readings then
coincide. \emph{No assertion of this obligation moves in either case}: what (i)--(iii)
below assert are dimensions, the characteristic polynomial \eqref{eq:esoracle} and the
full/primitive distinction, none of which names a class inside $E$
(v1.5.1, which states that reason rather than resting on the line case).

\medskip\noindent
\textbf{What is asserted and what is recorded.} Over $\cring$, in the stored basis:
\begin{enumerate}
\item[(i)] \textbf{Asserted, aborting:} the dimension $\dim\Heq{1}=2d_k-1$ and the
  characteristic polynomial \eqref{eq:esoracle} at those $(n,\ell,p)$ for which the
  constants $a_p(f)$ have been independently verified --- the shipped set being weights
  $k\in\{12,16,24\}$, the Eisenstein factor at every admissible $(n,\ell,p)$, and the
  Ramanujan congruence at $\ell=691$; $\Amat{0}(T_p)=1+p$ at $n=0$; and
  $\Heq{\bullet}=0$ for odd $n$ at every $\ell\neq2$ (Remark~\ref{rmk:negI}).
\item[(ii)] \textbf{Recorded, never asserted:} the same quantities at weights whose
  $a_p(f)$ this document has not independently verified (the shipped set being
  $k\in\{18,20,22,26\}$); every $(n,\ell)$ on the exceptional locus of
  Remark~\ref{rmk:dickson}; and every $\ell\in\{2,3\}$, where $6\notin\cring^{\times}$
  and the identification of $\Heq{\bullet}$ with $H^{\bullet}(\Ga;\VR)$ is not claimed in
  either direction (\ref{H3}).
\item[(iii)] \textbf{The full/primitive distinction applies here too.} An oracle stated
  for a prime-power index must name which operator it is about; the identity
  $\Amat{q}(T_{p^{2}})=\Amat{q}(T_p)^{2}-p^{\,n+1}I$ is about the \emph{full} operator,
  and the corresponding primitive statement carries the extra $-p^{\,n}I$
  (Obligation~\ref{obl:algebra}(ii),(iii); Appendix~\ref{app:errata},
  entry~\ref{err:e1}). \textbf{The worked oracles of \S\ref{sec:example} are bound by
  this clause}, oracle~6 in particular.
\end{enumerate}
The specialisations used as unit tests are listed in \S\ref{sec:example}; the boundary
analogue is Obligation~\ref{obl:taubdry}(vi), and the $\pm$-eigenspace refinement is
Obligation~\ref{obl:iota}(v). \emph{Provenance of the pointers} (rule~\ref{A2}): all
three are carried unchanged from the places this document already cites them ---
\cite[Ch.~8]{shimura} and \cite{deng} from Claim~\ref{cl:dimH1}(1),
\cite[Ch.~2--3]{harder3} from Theorem~\ref{thm:classical} --- and were not
independently re-verified against the texts here; \cite{deng} carries no section pointer
in either place, which is the weakest of the three.
\end{obligation}

\begin{remark}[Small-$\ell$ exceptional loci: where the classical oracle does not apply]
\label{rmk:dickson}
\PROVED{} Let $\ell\ge5$, $\cring=\Fl$, $n>0$. The dimension and characteristic-polynomial
oracles of Obligation~\ref{obl:oracle} equal the characteristic-$0$ values \emph{except}
where $\Sym^{n}$ acquires genuine mod-$\ell$ $\SL_2(\Fl)$-invariants, i.e.\ where
$\Heq{0}(K,\Rel;\Mn)=(\Mn)^{\SL_2(\Fl)}\neq0$. By Claim~\ref{cl:dimH1}(2) --- which is
where the invariant-theoretic input and its citation live, and which is not restated
here --- that happens exactly when $n$ lies in the numerical semigroup
\[
  \bigl\langle\,\ell+1,\ \ell^{2}-\ell\,\bigr\rangle
  \;=\;\bigl\{a(\ell+1)+b(\ell^{2}-\ell)\ :\ a,b\in\Z_{\ge0}\bigr\},
\]
generated by the degrees of the two Dickson generators; the first instance is
$n=\ell+1$, where the invariant is $x^{\ell}y-xy^{\ell}$. Call this the
\textbf{exceptional locus}. On it, $\Heq{0}\neq0$ for $n>0$ and $\dim\Heq{\bullet}$
legitimately \emph{exceeds} the characteristic-$0$ value, by Claim~\ref{cl:dimH1}(1);
the classical oracle does not apply and must be replaced by a directly computed and
recorded baseline (Obligation~\ref{obl:oracle}(ii)).

\medskip\noindent
\textbf{Both generators matter, and a test on the first alone is unsound.} The
predicate is membership in the semigroup, \emph{not} divisibility by $\ell+1$: at
$(n,\ell)=(20,5)$ one has $20=1\cdot(\ell^{2}-\ell)$ with $\ell+1=6\nmid20$, so a
$(\ell+1)\mid n$ test would assert a characteristic-$0$ dimension at a genuine
exceptional point. \COMP{} The implemented predicate is the semigroup membership
(\code{HeckeEndToEndValidation.cpp}, \code{IsDicksonExceptional}), and $(20,5)$ is in the
shipped grid.

\medskip\noindent
\textbf{The oracle settings of \S\ref{sec:example} avoid the locus}, and it must be
checked rather than assumed, since $n\ge\ell$ there: for $n=10$,
$10\notin\langle6,20\rangle$ ($\ell=5$) and $10\notin\langle8,42\rangle$ ($\ell=7$), so
$\dim\Heq{1}=3=2d_{12}-1$ holds at both.
\end{remark}

\begin{obligation}[FP-filter soundness]\label{obl:fp}
Assume \ref{H4}; \textbf{no hypothesis on $\ell$}.
\begin{enumerate}
\item[(i)] \textbf{No uncertified sign.} Every verdict returned by the filtered path
  satisfies the hypothesis of Theorem~\ref{thm:fp64}(1) or
  Corollary~\ref{cor:fptotal}/\ref{cor:fp32}; equivalently, for every call the returned
  sign equals the sign computed by an \emph{independent} exact evaluation of
  \eqref{eq:dot2}. Over $\Z$; a failure \textbf{aborts}. The exact reference must not be
  the production fallback: this clause is a differential test, and sharing the routine
  would make it vacuous.
\item[(ii)] \textbf{Dead code below the totality threshold.} For $B\le2^{24}$ (binary64)
  and $B\le2^{10}$ (binary32) the exact fallback is never entered. Over $\Z$; a failure
  \textbf{aborts}. The counters that witness it are a diagnostic and are never branched
  on (rule~\ref{A6}).
\item[(iii)] \textbf{Preconditions abort.} $s\notin\{\pm1\}$, an input of magnitude
  exceeding $B$, or $B$ above the exact-conversion limit ($2^{53}$, resp.\ $2^{24}$)
  throws; it never returns a value. Over $\Z$; by construction.
\item[(iv)] \textbf{The two builds are bit-identical.} The build with the
  floating-point path compiled out and the build with it enabled produce
  \emph{bit-identical} stored objects --- $\tau$-tables, $\Th^{(q)}_{\mathrm{flat}}$,
  $\Th^{(q)}_{\mathrm{eq}}$, $\Amat{q}$ and every hashed golden object. Over $\cring$
  and $\Z$; a difference \textbf{aborts}. \textbf{This clause is a two-build protocol,
  not a predicate a single run can evaluate}, and the roster of
  \S\ref{sec:obligationroster} records it as such: a single-process validation can
  establish (i)--(iii) and can establish that \emph{its own} output is
  build-independent, but (iv) as stated is discharged only by building twice and
  comparing.
\end{enumerate}
\end{obligation}

\begin{obligation}[Bounded-integer soundness]\label{obl:backend}
Assume \ref{H4}; \textbf{no hypothesis on $\ell$}.
\begin{enumerate}
\item[(i)] \textbf{The bound is certified, not assumed.} At table-build time
  $B_{\mathrm{tab}}=\max\lVert\gamma_t\rVert$ is recomputed by a full recount over every
  stored chain term and compared with the stored value, and
  $B=\max(B_{\mathrm{tab}},4N^{2})$ is recomputed from it
  (Lemma~\ref{lem:bounds}(3),(5)). Over $\Z$; a failure \textbf{aborts}.
\item[(ii)] \textbf{The selected tier covers $B$.} The assertion for the $64$-bit tier is
  $2B^{2}<2^{63}$, i.e.\ $B\le2^{31}-1$ (Theorem~\ref{thm:backend}(1)); an
  implementation using a wider tier asserts that tier's inequality instead. Over $\Z$;
  a failure \textbf{aborts}. \emph{Since v1.23.0-pre.9} the wider tier is the stored wide tier
  of Theorem~\ref{thm:backend}(4): a $\tau$-table with $B>2^{31}-1$ asserts $2NB\le2^{63}-1$
  (Lemma~\ref{lem:bounds}(7); code~$52$ if it fails), and is refused with code~$51$, as before,
  when the wide tier may not run (the rollback \code{HECKE\_TAU\_WIDE=0}, or a build forced
  narrow).
\item[(iii)] \textbf{Overflow aborts, never wraps.} Every exact product, sum, difference
  and negation is overflow-checked, including the cases whose naive form is undefined
  behaviour (negating the most negative value, taking its magnitude). Over $\Z$; a
  failure \textbf{aborts}.
\item[(iv)] \textbf{Arbitrary precision is absent.} No arbitrary-precision arithmetic is
  linked into the pipeline (Theorem~\ref{thm:backend}(3)); the escalation ladder is
  $64$-bit $\to$ $128$-bit $\to$ fixed-width wide, and the wide tier is unreachable in
  the target regime. This is a property of the build, \textbf{recorded} in the
  implementation map of \S\ref{sec:implementation} rather than asserted at runtime.
  \emph{Scope of the exclusion} (v1.11.0-pre.2; worded at v1.11.0 to the route as built,
  \file{ENGINE\_API.md} \S13.11.2 and \S13.11.8; owner ruling (ii) of 2026-09-25, hecke-extensions
  slot~11): a value of $\Z/\ell^{K}$ carried as $L=\lceil K/m\rceil$ digits in base $\ell^{m}$ of a
  \emph{fixed} word width $w\in\{32,64,128\}$, $m$ the largest exponent with $\ell^{m}<2^{w-1}$ ---
  the $\ell$-adic digit induction of \file{ENGINE\_API.md} \S13.11, in which $L$ is fixed per cell at
  construction from the working precision $K$ (on the pipeline, $K=k+e_{\max}$, fixed before the
  digit run from a certified bound, and capped by the grammar of \Espec{sec:ring}: $k\le4096$,
  $K\lceil\log_2\ell\rceil\le8192$ bits), is never grown during a computation, and in which every
  primitive is a single-word operation or one product into a \emph{fixed} double word of $2w$ bits
  (for $w=128$ the fixed two-limb type of \file{ENGINE\_API.md} \S13.12.2), with every intermediate
  bounded by a constant stated at its definition --- is \textbf{not} the arbitrary-precision
  arithmetic this clause excludes: no digit count is chosen by the data, no allocation follows a
  carry, and no arbitrary-precision library is linked. The escalation ladder above governs the
  exact-integer tables of Theorem~\ref{thm:backend} and is unchanged; the digit route is not a rung
  of it. A route whose digit count grows at run time, or whose digits are not native words, remains
  excluded. (v1.11.0-pre.2 read ``every operation on the digits is a native $w$-bit operation of
  the mod-$\ell^{m}$ backend of (v)''; as built, a digit product goes through the double word, and
  (v) is the mod-$\ell$ backend.)
\item[(v)] \textbf{The mod-$\ell$ backend is exact.} A modular product is computed
  exactly for the $\ell$ in force --- through a $64$-bit product for $\ell<2^{32}$, and
  through a $128$-bit product or a Barrett/Montgomery reduction with proved parameters
  for $2^{32}\le\ell<2^{64}$ (Theorem~\ref{thm:backend}, mod-$\ell$ box). Over $\Fl$;
  a failure \textbf{aborts}.
\end{enumerate}
\end{obligation}

\begin{obligation}[Certified $\tau$-table]\label{obl:tau}
Assume \ref{H4}; \textbf{no hypothesis on $\ell$}. An implementation of
Algorithm~\ref{alg:tau} emits, for each $g$ and each slot $(j,\sigma)$, the chain
\eqref{eq:taurecord} in canonical form, together with the recorded solve regime of every
triangle fill, and must discharge the following. Everything is evaluated \emph{exactly},
over $\Z$, or over $\Fl$ for an identity involving a chain that Algorithm~\ref{alg:tau}(3)
solved modulo $\ell$; \textbf{every clause aborts}, and the certification runs in the
constructor, so that an uncertified table cannot exist
(Remark~\ref{rmk:certoverGeometry}).
\begin{enumerate}
\item[(i)] \textbf{Structure.} Every stored $\gamma_t$ has $\det\gamma_t=1$; every term
  is in the canonical form of Definition~\ref{def:primitives}(P4), with sign $+1$ and
  non-zero coefficient; the canonicalisation is idempotent.
\item[(ii)] \textbf{(T3), augmentation.} $\aug\bigl(\tau(s_j\cdot\lift{v})\bigr)=1$ for
  every vertex slot.
\item[(iii)] \textbf{(T2), chain map.} $\partial\,\tau(x)=\tau(\partial x)$ for every
  stored representative $x$ of degree $1$ and $2$, the right-hand side evaluated
  \emph{only} through the shared (T1) helper of Algorithm~\ref{alg:tau}(0) and compared
  after reduction to canonical form --- so that construction and certification cannot
  disagree about how a derived slot is obtained.
\item[(iv)] \textbf{(T1), twisted equivariance bookkeeping.} For every stored transporter
  $h$ occurring in the table and every $j$, the slot $(j,\sigma)$ equals the sign-tracked
  (T1)-translate of the slot it is related to; this includes the self-constraints imposed
  by the elliptic stabilisers, and it is what makes the non-inequivalent slot set of
  Algorithm~\ref{alg:tau}(0) consistent.
\item[(v)] \textbf{Bounds.} Obligation~\ref{obl:backend}(i),(ii) on this table.
\item[(vi)] \textbf{The consumer table.} The eagerly reduced table of
  \eqref{eq:xiexplicit} matches a recomputation of $\rhon\bigl((g\gamma_t)\bmod\ell\bigr)$
  \emph{from each term's own stored exact $\gamma_t$}, entrywise --- so that a term
  pointing at the wrong cached matrix is caught. Over $\cring$.
\item[(vii)] \textbf{The solve regime is recorded.} Whether any fill used the mod-$\ell$
  fallback is written into the record; a table using it is marked
  \code{integral = false}, may not claim integrality, and is forbidden by default
  (owner decision OP2, Remark~\ref{rmk:tauintegral}). \textbf{Recorded}, and the
  default-policy violation \textbf{aborts}.
\end{enumerate}
\textbf{Given (i)--(vii), no geometric statement enters the correctness of $\Tdc{g}$}:
that is the entire purpose of this obligation, and it is why Algorithm~\ref{alg:tau}'s
steps~1--3 are free to choose.
\end{obligation}

\subsection{The obligations the later chapters added}\label{sec:lateroblig}

Four further obligations are stated in full where their objects are defined, and are
\emph{not} restated here (rule~\ref{A4}: an obligation has one home). What this section
adds is the integration the roster needs --- what each is over, and what its status is.

\begin{itemize}
\item \textbf{Obligation~\ref{obl:iota}} (\S\ref{sec:involution}), the involution.
  Clauses (i)--(iii) and (v) are over $\Z$ and $\cring$ and \textbf{abort}; clause (iv),
  the $\tau$-symmetry $\inv_{\Xbulk}\tau\inv_{\Xbulk}=\tau$, is over $\Z$ and is
  \textbf{recorded}, because it is the hypothesis of
  Proposition~\ref{prop:iotacommutes}(3b) and not a decision procedure (rule~\ref{A6}).
  \textbf{Clause (iv) was discharged negatively at v0.13.0 and is discharged positively
  since v1.1.0}: on the then-canonical $\tau$ the recorded value was \emph{false}, $3600$
  of $10764$ slots differing over the shipped $\ell\in\{2,3\}$ grid, with an exhibited
  witness and mechanism --- the ascending-index tie-break of
  Algorithm~\ref{alg:tau}(1b) on a mirror point, structurally incompatible with $\invK$,
  and, behind it, the breadth-first templates of step~2; on the canonical $\tau$ since
  v1.1.0 the value is \emph{true} on every slot of every table of the grid, and the
  validation gates it (Remark~\ref{rmk:iotatausym};
  Remark~\ref{rmk:certoverGeometry}(1)). The consequence is that at $\ell\in\{2,3\}$
  this document now has a route to $[\inv^{\ast},\Tdc{g}]=0$ for the canonical $\tau$:
  Proposition~\ref{prop:iotacommutes}(3b). The clause stays in force as the predicate any
  other $\tau$, vertex policy or refined triangulation must re-evaluate.
  \textbf{Clauses (ii), (iii) and (v) were implemented and run at v1.1.1}
  (\code{HeckeIotaSignValidation}), so the obligation is evaluated in full. At v1.1.1
  their aborting content passed while the then-\emph{recorded} sign half of (v) came out
  false for the geometric involution (Eisenstein class in $\ker(\inv^{\ast}_{\mathrm{geo}}+\id)$,
  mirror split, on every off-Dickson cell of a $77$-cell grid). \textbf{At v1.2.0 the
  owner resolved it by negating Definition~\ref{def:iota}(iv)},
  $\inv^{\ast}:=-\inv^{\ast}_{\mathrm{geo}}$, so that ``$+$'' is by definition the side of
  the Eisenstein classes; clause (v)'s sign statement is now a gating
  definition-consistency oracle (a failure would mean a non-uniform geometric sign, i.e.\
  an ill-posed definition) and holds on every signed cell, the degree-$0$ anchor is
  $\inv^{\ast}=-1$ on $H^{0}(\Ga;\Fl)$, and the $\pm$ split $(d_{\,n+2},d_{\,n+2}-1)$ is
  recorded as predicted on $66/66$ cells (Remark~\ref{rmk:iotasign};
  \file{docs/implementation/oq3\_resolution.md}). Nothing downstream changed: no
  pipeline value contains $\inv^{\ast}$.
\item \textbf{Obligation~\ref{obl:taubdry}} (\S\ref{sec:boundarydatum}), the boundary
  datum. Clauses (i)--(v) and (vii) are over $\Z$ or $\cring$ and \textbf{abort};
  clause (iv) additionally makes integrality unconditional on the boundary, unlike the
  bulk; clause (vi)'s general-$n$ $H^{1}$ multiset is \textbf{recorded} until
  Claim~\ref{cl:chi} is proved. \textbf{Unimplemented} (\S\ref{sec:boundarydatum}, code
  map).
\item \textbf{Obligation~\ref{obl:eta}} (\S\ref{sec:cone}), the pair homotopy. All
  clauses are over $\Z$ or $\cring$ and \textbf{abort}, except clause (vi)'s three
  diagnostics --- \code{eta\_zero}, \code{cochain\_commutes} and the solve regimes ---
  which are \textbf{recorded} and may never select $\hder=0$. Clause (v), that the
  stored $\hder$ is the transport of the stored $\etapair$, is what pins $\Acone{q}$;
  no identity among matrices does. \textbf{Unimplemented} (\S\ref{sec:cone}, code map).
\item \textbf{Obligation~\ref{obl:coneexact}} (\S\ref{sec:cone}), cone exactness and
  intertwining. Over $\cring$; clauses (i)--(iii) and (viii) \textbf{abort}
  unconditionally, (iv)'s non-field spots and (v)'s $\Heq{2}=0$ at
  $6\notin\cring^{\times}$ are \textbf{recorded}, (vi) aborts for $6\in\cring^{\times}$
  and records otherwise, and (vii) is split by route. \textbf{Unimplemented}
  (\S\ref{sec:cone}, code map). Note that (vi) re-checks
  Theorem~\ref{thm:relations}(1) on $\Hcpt{q}$ \emph{with the full $T_{p^{2}}$ and the
  coefficient $p^{\,n+1}$}: erratum E1 binds there exactly as it binds
  Obligation~\ref{obl:algebra}(ii).
\end{itemize}

\subsection{The roster}\label{sec:obligationroster}

\begin{contractbox}
\textbf{A printed proof and a run are different epistemic positions, and this table keeps
them apart.} An obligation with a registered validation behind it has been evaluated, on
a stated grid, by something that can be re-run; an obligation without one rests on the
proof of the statement it discharges, plus whatever \emph{ad hoc} computation is recorded
at that statement. Neither is worthless and neither substitutes for the other, and
\COMP{} never upgrades to \PROVED{} (\S\ref{sec:status}). ``Registered'' below means:
named in the validation registry of the shipped tree, hence runnable by name and included
in a full suite run.
\end{contractbox}

\noindent
Of the $94$ registered validations (v1.6.2; $89$ at v1.5.2, $80$ at v1.4.2, $77$ at v1.3.3, $68$ at v1.3.2), the column below names
those that evaluate the obligation's own predicate. Status is one of \textbf{discharged} (every clause evaluated,
or proved where the clause is not a run), \textbf{partial} (some clause is not evaluable
through the present interface, and the obligation says which), \textbf{negative} (the
predicate was evaluated and is false, and the document records the consequence), and
\textbf{unimplemented} (the objects the obligation reads do not exist in the tree).

\begin{center}
\footnotesize
\begin{tabular}{@{}p{1.25cm}p{1.0cm}p{3.3cm}p{3.5cm}p{4.5cm}@{}}
\toprule
obligation & over & failure & status & registered validations \\
\midrule
\ref{obl:welldef}  & $\cring$       & aborts & discharged & \code{HeckeOperatorSmoke\-Validation}, \code{HeckeIndependence\-Validation}, \code{HeckeEndToEnd\-Validation} \\
\addlinespace
\ref{obl:indep}    & $\cring$       & aborts; (i) partly records & \textbf{partial}: (iii) not reachable & \code{HeckeIndependence\-Validation} \\
\addlinespace
\ref{obl:integral} & $\Z$, $\cring$ & aborts & discharged & \code{HeckeEndToEnd\-Validation}, \code{HeckeCosetTables\-Validation}, \code{HeckeBadPrime\-Validation} \\
\addlinespace
\ref{obl:algebra}  & $\cring$       & aborts; records at $\ell\in\{2,3\}$ & discharged & \code{HeckeRelations\-Validation}, \code{HeckeOperatorFamily\-SmokeValidation}, \code{HeckeEndToEnd\-Validation} \\
\addlinespace
\ref{obl:agreement}& $\cring$       & aborts & discharged & \code{HeckeCrossCheck\-Validation}, \code{HeckeIndependence\-Validation}, \code{HeckeEndToEnd\-Validation}; also in \code{compute()} \\
\addlinespace
\ref{obl:shapiro}  & $\cring$       & aborts & discharged & \code{HeckeShapiro\-Validation}, \code{HeckeInducedLocal\-SystemSmokeValidation} \\
\addlinespace
\ref{obl:oracle}   & $\cring$       & aborts / records per (i),(ii) & discharged on the shipped grid & \code{HeckeOracle\-Validation}, \code{HeckeEndToEnd\-Validation}, \code{HeckeBlockForm\-Validation} \\
\addlinespace
\ref{rmk:dickson}  & $\Fl$          & gates \ref{obl:oracle} & discharged & \code{HeckeEndToEnd\-Validation} (\code{hecke::IsDicksonExceptionalCount}, \file{Hecke/DicksonCount.h} since v1.3.3 --- one count shared by the three former file-local copies) \\
\addlinespace
\ref{obl:fp}       & $\Z$, $\cring$ & aborts & discharged; (iv) by build protocol & \code{HeckeFpFilter\-Validation}; (iv) needs two builds \\
\addlinespace
\ref{obl:backend}  & $\Z$, $\Fl$    & aborts; (iv) records & discharged & \code{HeckeExactIntSmoke\-Validation}, \code{HeckeTauTable\-Validation}, \code{HeckeExtendedRegime\-Validation} \\
\addlinespace
\ref{obl:tau}      & $\Z$ / $\Fl$   & aborts; (vii) records & discharged; (vi) vacuous on a \code{ChainsOnly} table (no consumer table; code map) & \code{HeckeTauTable\-Validation}, \code{HeckeTauTableSelf\-CheckValidation}, \code{HeckeModuleStream\-Validation} (\code{ChainsOnly}); also in the constructor \\
\midrule
\ref{obl:iota}     & $\Z$, $\cring$ & (i)--(iii),(v) abort; (iv) records; (v)'s $\pm$ dimensions record & (i),(ii),(iii) discharged; (iv) discharged (\textbf{negative} at v0.13.0, positive and gating since v1.1.0); (v) discharged --- the sign gating and \textbf{true by definition} since v1.2.0 (owner decision OQ3, Definition~\ref{def:iota}(iv) negated), the recorded split as predicted & \code{HeckeIotaTauSymmetry\-Validation}, \code{HeckeIotaSign\-Validation} \\
\addlinespace
\ref{obl:taubdry}  & $\Z$, $\cring$ & aborts; (vi) partly records; (vii) records at $\ell\in\{2,3\}$ & discharged since v1.3.3, and re-read against the v1.4.0 clauses with no arm changed: (i)--(iv) over $\Z$ on every table, $\ell\in\{2,3\}$ included; (v.a) route B on $H^{0}$ at every $\nu$ of the grid, (v.b) route B on $H^{1}$ and (v.c) route A at $\nu=1$, all in \code{compute()}; (vi) asserted where proved, the general-$n$ $H^{1}$ multiset recorded ($330/330$); (vii) asserted at $6\in\cring^{\times}$ with the $\beta_j$ control failing on the $170$ cells the qualified clause predicts, recorded at $\ell\in\{2,3\}$ ($216/216$; Remark~\ref{rmk:conerestrictsmallell}) & \code{HeckeBoundaryTau\-Validation} ((i)--(iv)), \code{HeckeBoundaryOperator\-Validation} ((v.a)--(v.c), (vi), (vii)); (i)--(iv) also in the constructor, (v.a)--(v.c) also in \code{compute()} \\
\addlinespace
\ref{obl:eta}      & $\Z$, $\cring$ & aborts; (vi) records & \textbf{discharged} since v1.5.2 (\textbf{partial} at v1.4.2): (i)--(iv) over $\Z$ on every table, $\ell\in\{2,3\}$ included, in the constructor and re-derived by the validation ($152$ tables, $27$ tampered chains rejected, no mod-$m$ path); (v) since slot-3 WP-6: the stored $\hder$ equals \eqref{eq:derivedh} recomputed by an independent coset transport, in every computed degree, on $63$ cells / $212$ operators, $\ell\in\{2,3\}$ included, and the tampered $\hder'=\hder+\xi$ of Proposition~\ref{prop:coneambiguity}(2) fails exactly this clause at $(10,691)$; (vi) recorded (\code{eta\_zero} false on all $152$ tables, and with \code{cochain\_commutes} false on all $212$ operators of the cone-action grid, the implication vacuous) & \code{HeckePairHomotopy\-Validation} ((i)--(iv), (vi)), \code{HeckeConeAction\-Validation} ((v), (vi)); (i)--(iv) also in the constructor, (v) also in \code{HeckeConeAction::compute()} \\
\addlinespace
\ref{obl:coneexact}& $\cring$       & aborts; (iv)--(vii) partly record & \textbf{implemented} since v1.7.3 (slot-6 WP-3, WP-4, WP-5); \textbf{partial} v1.5.2--v1.7.2 (slot-3 WP-5, WP-6; \textbf{unimplemented} at v1.4.2): (iv), (v) and $\Dcone^{2}=0$ on $74$ cells, every spot of \eqref{eq:coneles} exact ($592/592$), $\dim\Hcpt{1}=3$ at $(10,691)$; (i) (the cone identity asserted on the stored $\Ceq$ basis in $424/424$ degrees, the full flat identity \textbf{recorded} --- $29$ of $424$ --- the coset transport being defined on $\Hom_{\Z[\Ga]}(-,V)$; the reading fixed at v1.6.0, the clause naming the basis), (ii), (iii) (the $\rres$-half asserted at \emph{every} $\ell$ since v1.7.0, the $74$ operators at $\ell\in\{2,3\}$ included: $424/424$, and asserted rather than merely measured there by \code{HeckeConeActionValidation}'s own independent re-derivation since slot-5 WP-4), (vi) ($335/335$, the \emph{full} $T_{p^{2}}$), (viii)(a) (two $\etapair$ $\times$ two $\taub$ policies, $\Acone{q}$ bit-identical, both table hashes moving) and (viii)(b) (fails exactly Obligation~\ref{obl:eta}(v), changes $\Acone{1}$ at $(10,691)$); \textbf{(vii) implemented since v1.7.3} (slot-6 WP-4, WP-5, WP-3): the strict-pair route gives $\Acone{q}$ entrywise equal to the normative one on $212/212$ operators ($636$ degrees), asserted at $6\in\cring^{\times}$ and recorded at $\ell\in\{2,3\}$ with $0$ recorded disagreements; route~3 on $212/212$, asserted at every $\ell$; the PD adjointness gate $A_c^{T}G=G\,A(\Dn)$ for $p\in\{2,3,5,7\}$ asserted at $\ell\ge5$ on $160$ (operator, $t$) pairs over $20$ cells, failing $0$, the $p^{n}$-twisted controls rejected --- a test of Claim~\ref{cl:pd}, which stays open (at v1.7.2: unimplemented, the cross-routes phase E5's slot-6 arms, PD adjointness recorded unavailable). \textbf{Over $\Z/\ell^{e}$, $\ell\ge5$, since v1.9.0} (hecke-extensions slot-9 WP-4, phase E8 item~4b, route~1 only): (i)--(iii) as over $\Fl$; (iv) \textbf{including its non-field half}, checked by solves at every spot with the witnesses re-derived (\code{ConeModule\-Certificate} item~14); (v) in lengths; (vi) \textbf{partial} --- the commutator asserted, the $\nu=2$ recursion not computed over the ring; (vii) \textbf{open} over the ring through v1.9.1 --- the two cross-routes and the relative constructor refused, each with an arm --- and \textbf{since v1.10.0} (slot-10 WP-3) its first two sentences \textbf{discharged} at $\ell^{e}<2^{32}$, both cross-routes entrywise equal to route~1 at the seven points, its third (the PD gate) recorded unavailable over the ring; (viii) $\Fl$-only; at the seven points of \Espec{sec:ring}~(e) and on the case list at $\ell\ge5$ & \code{HeckeCone\-Validation} ((iv), (v), $\Dcone^{2}=0$), \code{HeckeConeAction\-Validation} ((i)--(iii), (vi), (viii)); (iv), (v) also in \code{EquivariantMappingCone}'s \code{ConeCertificate}, (i)--(iii) and (vi) in \code{ConeActionCertificate} / \code{ConeRelationsCertificate}; over $\Z/\ell^{e}$ (v1.9.0) \code{HeckeReduction\-Validation} arms (4b-B), (4b-C), (4b-H), (4b-R), (4b-K), (X4b) and \code{HeckeReduction\-CoverageSlow\-Validation}, with \code{ConeModule\-Certificate}; (vii) over $\Z/\ell^{e}$ (v1.10.0) \code{HeckeConeCrossRoute\-Ring\-Validation} \\
\bottomrule
\end{tabular}
\end{center}

\noindent
\textbf{Reading the table.} Twelve of the fifteen rows are discharged by a registered
validation --- ten at v1.3.2, Obligation~\ref{obl:taubdry} since v1.3.3, whose two
validations and the clause-by-clause record are in \S\ref{sec:boundarydatum}'s code map,
and Obligation~\ref{obl:eta} since v1.5.2 (slot-3 WP-6 discharged its clause (v)). One --- Obligation~\ref{obl:indep} --- is partial for a stated structural
reason, and the missing clause is the one that needs an entry point this tree does not
have. One --- Obligation~\ref{obl:iota} --- is evaluated in full and carries \emph{two}
negative records in its history rather than a hope: clause (iv) was discharged
negatively at v0.13.0 before the v1.1.0 repair made it hold and gate, and clause (v)'s
sign prediction was false for the geometric involution at v1.1.1 before the owner
resolved it at v1.2.0 by negating Definition~\ref{def:iota}(iv), after which the clause
gates and holds. One --- Obligation~\ref{obl:coneexact} --- was partial from v1.5.2 to v1.7.2: every clause but
(vii) was evaluated (slot-3 WP-5 and WP-6), and (vii)'s cross-routes were a scheduled phase,
not a missing entry point; \textbf{since v1.7.3 (slot-6 WP-3, WP-4, WP-5) clause (vii) is
evaluated too}, all three of its sentences; until v1.5.1 it and Obligation~\ref{obl:eta} were
\emph{unimplemented}, their clauses a specification written as predicates on emitted data
precisely so that they could be handed to the phase that builds them unchanged (which is
what happened to Obligation~\ref{obl:taubdry}: written at v1.0.0 as such a specification,
handed to the phase unchanged, and discharged at v1.3.3 with no clause moved --- and to
Obligations~\ref{obl:eta} and \ref{obl:coneexact} in slots 2--3, likewise with no clause
moved). One further clause, Obligation~\ref{obl:fp}(iv), is a \emph{protocol} rather
than a predicate and is recorded as such. \COMP{} That protocol has been executed: on
2026-09-16 the tree was configured and built twice out of source, with the floating-point
filter enabled and with \code{HECKE\_PURE\_INTEGER=ON}, and in both builds
\code{GOLDEN\_STRICT=1 ./cohomology\_app GoldenRegressionValidation} reported tier~A
$9446$ checked, $0$ failing and tier~B $4603$ checked, $0$ mismatching against the same
frozen reference --- so all $14049$ hashed objects, $\Th^{(q)}_{\mathrm{flat}}$,
$\Amat{q}$ and the $\tau$-derived data included, agree bit for bit between the two builds
--- and \code{HeckeFpFilterValidation}, \code{HeckeExactIntSmokeValidation},
\code{HeckeTauTableValidation}, \code{HeckeEndToEndValidation} and
\code{HeckeOracleValidation} each exited $0$ in both
(\file{docs/matrix-engine/tools/cross\_config.sh}). This is a recorded run of an existing
script, not a registered validation, and the table says so.

\paragraph{Code map.}
Obligations~\ref{obl:welldef} and \ref{obl:agreement}(i) are asserted inside
\code{HeckeOperatorOnK::compute()}, which throws \code{std::runtime\_error} and leaves
the object in the not-computed state; the executed counts and the route bit-identity flag
are carried in the \code{certificate} object of \code{to\_json()} and in
\code{certificates()}. Obligation~\ref{obl:tau} is \code{hecke::HeckeTauTable::certify()},
called by the constructor, clause by clause: (i)--(iv) its items (a)--(d), (v) its item
(e), (vi) its item (f), (vii) \code{used\_mod\_ell\_solve()}.
\textbf{At v1.23.0-pre.5} (slot-23 G2-F, value-preserving): a table built \code{hecke::TauTableContents::ChainsOnly} --- the module route's table,
built by \code{hecke\_ext::RunModuleT} on the stream route --- keeps the chain records only and carries \emph{no consumer table} (no table of
\eqref{eq:xiexplicit}, no $\rhon$ pool), so on it (vi) has no object and its certification is vacuous (item (f) and the length match of (a) are skipped);
(i)--(v) and (vii) run unchanged in the constructor. The module route reads no consumer table on any path: it evaluates $\rhon$ mod $\ell^{k}$ from each
term's exact $\gamma_t$ (the base change of the integral datum, Remark~\ref{rmk:tauintegral}). Every other consumer builds the table \code{Full}, with (vi) as written
(\file{docs/implementation/slot23/G2F\_report.md}, \file{G2F\_audit.md}).
Obligation~\ref{obl:backend} is \file{Hecke/ExactInt.h} (\code{checked\_*},
\code{HeckeOverflow}, \code{uabs64}) together with that item (e).
Obligation~\ref{obl:fp} is \file{Hecke/Predicates.cpp} and the \code{HECKE\_PURE\_INTEGER}
option. Obligation~\ref{obl:algebra} is \code{HeckeOperatorFamily::T\_p\_pow} (recursion),
\code{T\_p\_pow\_direct} (the full operator by \eqref{eq:fullassembly}) and
\code{T\_p\_pow\_primitive\_direct} (the raw primitive operator) --- the three entry
points of entry~\ref{err:e1}'s table, and the reason clause (iii) can name which
identity belongs to which. Obligation~\ref{obl:oracle}'s gating predicate is
\code{IsDicksonExceptional}. The aggregated gate is
\code{HeckeEndToEndValidation}, which prints an obligation-to-test-identifier coverage
table for the nine it covers; Obligation~\ref{obl:indep} is covered by
\code{HeckeIndependenceValidation} alone, and Obligation~\ref{obl:iota} by the two
validations of Chapter~\ref{sec:involution}'s code map --- (i),(iv) by
\code{HeckeIotaTauSymmetryValidation} and (ii),(iii),(v) by
\code{HeckeIotaSignValidation} --- neither of which is part of that aggregate gate.

\noindent\textbf{v1.23.0-pre.4 (slot-23 G2: the module route over $\Z/\ell^{k}$, a value-preserving re-implementation; no clause relaxed).}
On the module route (\code{HeckeOperatorOnKT::compute\_module()}, \file{Hecke/HeckeOperatorOnK.\{h,cpp\}}, and
\code{EquivariantCohomologyOnKT::compute\_cohomology\_module()}, \file{EquivariantCohomologyOnK.\{h,cpp\}}) Obligations~\ref{obl:welldef} and
\ref{obl:agreement}(i) are discharged as written, in every degree $q=0,1,2$, by a re-implementation (\file{docs/implementation/slot23/G2\_plan.md},
\file{G2AB\_report.md}, \file{G2CDE\_report.md}, \file{G2G\_report.md} and their adversarial audits; merged at \texttt{c5af90e9}, \texttt{cb967760},
\texttt{b33e6fb7}):
(G2-A) the solve that defines $\Th^{(q)}_{\mathrm{eq}}$ (\ref{obl:welldef}(iii)), the solve for $X_q$ in the cohomology and the first solve of
\code{H\_class\_coordinates()} are read off the unit Howell pivots of the $C_{\mathrm{eq}}$ each solves against ($C^{q}_{\mathrm{eq}}$ for $\Th^{(q)}_{\mathrm{eq}}$, $C^{q+1}_{\mathrm{eq}}$ for $X_q$; \code{ModuleDegree::piv}, $C[\mathrm{piv},:]=I$ checked, so $C$ is
injective and the read-off is the unique solution) by \code{engine::elim::solve\_unit\_pivot\_readoff} (\file{src/ModArithmetic/DenseElim.\{h,cpp\}},
\file{ENGINE\_API.md}~\S13.6), with the residual recomputed exactly over every row; a non-unit pivot (Hida's lattice), a failed identity check or a non-zero
residual falls back to the Howell solve \code{solve\_multi\_rhs\_module}, with its error texts; (G2-B) the $d\circ d$ certificate $X_qX_{q-1}=0$ is
decided by one exact GEMM (the scalar loop \code{engine::elim::product\_is\_zero} is the rollback's), and \code{elim::cyclic\_decomposition}'s column clear runs with the loops interchanged;
(G2-C) \code{compute\_module\_stream()} assembles $\Th^{(q)}$ coset by coset ($\Omega_i=\rhon(s_i^{-1})$ evaluated once per degree and coset, memos scoped
to one degree and coset), and on $T_\ell$'s low-rank path sums one coset's terms per block by one GEMM (\code{hecke::lowrank::sum\_of\_low\_rank},
\file{Hecke/HeckeLowRank.h}); (G2-D) route~2 is built in $\sigma$-row chunks and compared with $\Th^{(q)}$ entrywise on every entry of every degree before it is dropped
(\ref{obl:agreement}(i) unchanged, the mismatch text unchanged), the scan of \code{factor\_low\_rank} taking its valuation by \code{FastValuation} (equal to \code{ModCtx::valuation} on every residue); (G2-E) $\Th^{(q)}_{\mathrm{eq}}$ is read off the product $\Th^{(q)}C$
(\code{engine::elim::solve\_unit\_pivot\_readoff\_product}, only the non-pivot rows formed, the residual over every row), and the sat / Loewy runner
\code{RunModuleT} (\file{HeckeExt/HeckeRingHeckeAlgebra.cpp}), alone and on the stream route, sets \code{ModuleRetention::Lean} on its operators and calls
\code{release\_module\_intermediates()} on its cohomology, so that $\Th^{(q)}$, $\Th^{(q)}_{\mathrm{eq}}$ and the cohomology's intermediates $X_q$, \code{Zc} are freed after their last
certificate and their guarded accessors (\code{Theta\_module}, \code{Theta\_eq\_module}, \code{module\_X}, \code{module\_Zc}) throw a named error;
(G2-G) \code{elim::cyclic\_decomposition} and \code{elim::invariant\_factors} keep lean work buffers (one Howell form of $[\Gamma^{\mathsf T}\mid I]$ for the
kernel and the coordinates, \code{kernel\_from\_augmented}, \code{solve\_from\_augmented}). No certificate item is sampled, relaxed or moved out of
\code{compute\_module()}. \textbf{One rollback lever, value-neutral}: \code{HECKE\_MODULE\_STREAM=0} (read once per process by
\code{engine::elim::module\_stream\_enabled()}; any value other than \code{0} or \code{1} warns and keeps the default; scoped as \code{ScopedModuleStream})
runs \code{compute\_module\_legacy()}, the Howell solves, the scalar $d\circ d$ loop, the original loop order and buffers and full retention, i.e.\ the
pre-G2 route (the lever also selects the legacy elimination buffers in every caller, cohomology-only runs and Integral subquotients included); when the variable is set, one \code{EFFECTIVE HECKE\_MODULE\_STREAM} line is printed at process exit with the counters of \code{module\_stream\_totals()}. It
selects a route, never a value: every entry computed is the residue of an exact sum or the unique solution of an injective system, so order and grouping
cannot move a byte; G2-B's loop interchange is neutral because its column operations are independent (none writes the pivot column, and the pivot row's quotients are taken first), and G2-G's shared buffer because the Howell form is canonical; and the gate is \code{HeckeModuleStreamValidation} (\code{standard}) and \code{HeckeModuleStreamSlowValidation} (\code{slow}),
registered in \file{Registry/Slot23/G2.cpp} (source \file{HeckeExt/HeckeModuleStreamValidation.cpp}): stream against legacy entrywise on every object,
byte-equal \code{SAT\_DEFECT} lines, negative controls on the residual, on route~2 and on a released object, and non-vacuity counters. The golden reference
and the digests are unchanged across these packages.
\textbf{At v1.23.0-pre.7} (slot-23 G2-T, value-preserving; \file{docs/implementation/slot23/G2T\_report.md}, \file{G2T\_audit.md}): on the stream route only,
\code{compute\_module\_stream()} evaluates each $\rhon(g)$ mod $\ell^{k}$ straight into row-major buffers (\code{ComputeSymmetricPowerRowsRing},
\file{src/ModArithmetic/FlatMatrix.\{h,cpp\}}: one Pascal row per operator, the column steps in a wavefront, the inverse of the divisor folded into
one reduction per entry where the accumulator cannot overflow) and factors $T_\ell$'s terms by \code{factor\_low\_rank(\dots, fast\_scan)} (a divisibility
scan by $\ell^{v}$, a valuation floor, a rank-one GEMM update; the same pivot sequence); \textbf{one rollback lever, value-neutral}: \code{HECKE\_MODULE\_EVAL=0}
(\code{engine::elim::module\_eval\_enabled()}, scoped as \code{ScopedModuleEval}; inert under \code{HECKE\_MODULE\_STREAM=0}; when set, one
\code{EFFECTIVE HECKE\_MODULE\_EVAL} line at exit) restores the G2-F evaluation; both compute the same canonical residues, so the lever selects a route,
never a value (gate: \code{HeckeModuleStreamValidation}~(V)).
\textbf{At v1.23.0-pre.8} (slot-23 G2-T2 and G2-P, value-preserving; \file{docs/implementation/slot23/G2T2\_report.md}, \file{G2T2\_audit.md},
\file{G2P\_report.md}, \file{G2P\_audit.md}). G2-T2: \code{ComputeSymmetricPowerRowsRing}'s column steps run on AVX2 / AVX-512F lanes where the modulus
allows Shoup products ($\ell^{k}\le2^{31}$; otherwise scalar), each lane computing the scalar step's canonical residue; and route~2 of a degree is built in
one chunk with one $\rhon$ memo per coset shared across the $\sigma$-rows (the memo is keyed by the entries mod $\ell^{k}$, of which $\rhon$ is a function); two
\textbf{rollback levers, value-neutral}: \code{HECKE\_RHO\_SIMD=0} (the scalar fill; \code{RhoFillPath}, scoped as \code{ScopedRhoFillPath}) and
\code{HECKE\_ROUTE2\_SHARED\_MEMO=0} (G2-D's $\sigma$-row chunks; \code{ScopedRoute2SharedMemo}), each with an \code{EFFECTIVE} line at exit when set (gate:
\code{HeckeModuleStreamValidation}~(V)(d)--(g)). G2-P: when the consumer pool (\code{ENGINE\_THREADS}) has more than one worker, \code{compute\_module\_stream()}
assembles $\Theta^{(q)}$ and route~2 in coset batches (the distinct keys listed serially, evaluated by the pool, each $\sigma$-row's block accumulated by one
task in the serial term order and grouping), and \code{howell\_rows} and the product read-off split their disjoint row updates over the elimination hook;
\textbf{the values do not depend on the thread count or the schedule} --- every entry is the canonical residue of an exact sum formed from the same operands,
the Howell form and the read-off are unique, and \code{ENGINE\_THREADS=1} is the serial route itself; \textbf{one rollback lever, value-neutral}:
\code{HECKE\_MODULE\_PARALLEL=0} (\code{engine::elim::module\_parallel\_enabled()}, scoped as \code{ScopedModuleParallel}; the \code{EFFECTIVE
HECKE\_MODULE\_PARALLEL} line at exit when set, also naming \code{ENGINE\_HOST\_ELIM\_HOOK=0}, which keeps the Howell and read-off splits serial) restores
the serial route at any thread count (gate: \code{HeckeModuleStreamValidation}~(P): every object at consumer pools of $1$, $2$, $7$, $16$ workers against
pool~$1$, the rollback, the negative controls of the parallel read-off on a pool worker; \code{HeckeModuleStreamSlowValidation}~(P)).
\textbf{At v1.23.0-pre.10} (slot-24 MT-1, value-preserving; \file{docs/implementation/slot24/MT1\_report.md}, \file{MT1\_audit.md}, MERGE after fixes F1--F3).
(0) At every thread count, \code{elim::cyclic\_decomposition}'s pivot search on the lean module route visits the trailing entries in the same column-major order
with the same strict comparison and stops at the first unit: the serial search keeps the first entry of least valuation, valuation~$0$ is the least possible,
so the held candidate is the one the full scan returns, and everything after the search is unchanged. (a)--(e) When \code{HECKE\_MODULE\_PARALLEL} is on and the
consumer pool has more than one worker, the module route also splits over the pool (all but the assembly's split over the elimination hook, so under
\code{ENGINE\_HOST\_ELIM\_HOOK=0} only the assembly splits, as G2-P's Howell and read-off splits stay serial there): the remaining pivot search (tiles over the columns, folded in tile order
with the same strict comparison, i.e.\ the first minimum of the concatenation) and \code{cyclic\_decomposition}'s disjoint row operations; the products
$Y=\Th H$, $DC$ and $C\cdot\mathrm{basis}$ in row or column blocks (each entry the same exact dot product modulo $\ell^{k}$); the runtime assertions of
Obligation~\ref{obl:welldef}(i), (ii) checked per column by pure reads, the serial loops run verbatim, with their first message, if any column fails (so the
outcome, the message and the executed counts are the serial ones); the assembly as one task per $(\sigma,\sigma_k)$ block, in the serial term order; and the
residual of the read-off behind Obligation~\ref{obl:welldef}(iii) and the right-hand sides of \code{solve\_from\_augmented} in blocks, the smallest failing
index's message thrown. So Obligations~\ref{obl:welldef} and \ref{obl:agreement}(i) are discharged as written and \textbf{the values do not depend on the
thread count or the schedule}; \textbf{one rollback lever, value-neutral}: \code{HECKE\_MODULE\_MT=0} (\code{engine::elim::module\_mt\_enabled()}, scoped as
\code{ScopedModuleMt}; any other value warns and keeps the default; the \code{EFFECTIVE HECKE\_MODULE\_MT} line at exit when set, with the counter
\code{pivot\_exit}; inert under \code{HECKE\_MODULE\_STREAM=0}) restores G2-P's route verbatim, the full scan included (gate:
\code{HeckeModuleMtValidation}, \code{standard}, \file{Registry/Slot24/MT1.cpp}: \code{cyclic\_decomposition} at consumer pools
of $1$, $2$, $3$, $7$, $13$, $16$ workers, the read-off and the solves at pools of $1$, $2$, $7$, $16$, every object of a set of cells and whole \code{SAT\_DEFECT} cells at pools of $1$, $2$, $7$, $16$, each against the lever
off at pool~$1$).
\textbf{At v1.23.0-pre.11} (slot-24 G2-H, value-preserving; \file{docs/implementation/slot24/G2H\_report.md}, \file{G2H\_audit.md}, MERGE after fixes F1, F2).
On the stream route, under the lean retention of the sat / Loewy runner \code{RunModuleT} only, Obligation~\ref{obl:welldef}(iii) is \textbf{certified without
forming $\Th^{(q)}_{\mathrm{eq}}$}: \code{engine::elim::unit\_pivot\_residual\_product} makes the read-off's checks that $C^{q}_{\mathrm{eq}}[\mathrm{piv},:]=I$, in the same order,
and then checks $(\Th^{(q)}[\mathrm{np},:]-C^{q}_{\mathrm{eq}}[\mathrm{np},:]\,\Th^{(q)}[\mathrm{piv},:])\,C^{q}_{\mathrm{eq}}=0$ entrywise, which by associativity and distributivity over
$\Z/\ell^{k}$ (no unit used) is exactly the vanishing of the read-off's residual $\Th^{(q)}C-C\,(\Th^{(q)}C)[\mathrm{piv},:]$; with $C[\mathrm{piv},:]=I$ the read-off
$(\Th^{(q)}C)[\mathrm{piv},:]$ is the unique solution of the solve that defines $\Th^{(q)}_{\mathrm{eq}}$, so the certificate is the read-off's, item for item, and
Obligation~\ref{obl:welldef}(iii) is discharged as written; any other outcome takes the Howell solve with its texts. There, route~2 is subtracted from
$\Th^{(q)}$ in place after $Y=\Th^{(q)}H$ and (i)--(iii), and Obligation~\ref{obl:agreement}(i) is checked as $\Th^{(q)}-(\text{route 2})=0$ entrywise (the same exact
sums, the same mismatch text); the cohomology's $X_q$ and \code{Zc} are freed inside \code{compute\_cohomology\_module()} after their last reader
(\code{set\_module\_early\_release}). On the stream route in every caller, \code{howell\_rows} adopts its store as $H$ (the same rows in the same order) when no
closure row was appended and every row of the store is a pivot row, and the lean \code{cyclic\_decomposition} / \code{invariant\_factors} take the length from
its pivot loop. No certificate item is sampled, relaxed or moved out of \code{compute\_module()}. \textbf{One rollback lever, value-neutral}:
\code{HECKE\_MODULE\_LEAN\_PEAK=0} (\code{engine::elim::module\_lean\_peak\_enabled()}, scoped as \code{ScopedModuleLeanPeak}; any other value warns and keeps the
default; inert under \code{HECKE\_MODULE\_STREAM=0}; when set, one \code{EFFECTIVE HECKE\_MODULE\_LEAN\_PEAK} line at exit) restores the route of before (gate:
\code{HeckeModuleStreamValidation}~(H)).

\paragraph{Formalization note.}
Every obligation here is already a decidable predicate on a finite emitted table
(rule~\ref{A4}), which is the shape a formalization wants: the statement to prove is
``if the predicate holds on the emitted data then the stored $\Amat{q}$ is the matrix of
$\Tdc{g}$'', and the predicate itself needs no proof, only evaluation. Three shapes recur
and are worth factoring once: \emph{entrywise equality of two matrices over $\cring$}
(Obligations~\ref{obl:welldef}(iii), \ref{obl:agreement}, \ref{obl:algebra}),
\emph{equality of two integer chains after canonicalisation}
(Obligation~\ref{obl:tau}(iii),(iv), and identically in
Obligations~\ref{obl:taubdry} and \ref{obl:eta}), and \emph{a magnitude recount against a
stored bound} (Obligation~\ref{obl:backend}(i)). The one obligation that is not of this
shape is \ref{obl:fp}(iv), which quantifies over \emph{builds}; a formalization would
express it as the statement that the pipeline's output is a function of the abstract
sign predicate alone --- which is exactly what Lemma~\ref{lem:predshapes} plus
Theorem~\ref{thm:fp64} give, and is a better target than the two-build diff.

%----------------------------------------------------------------------
\section{Worked micro-example ($p=2$) and the numerical oracles}\label{sec:example}
%----------------------------------------------------------------------
% SOURCE:  hecke_action.tex:1512-1612 (RH+NEW); the shipped tree read and run
%          on 2026-09-16 for every recorded number.
% NOTATION: rows s_j/i (C14, C15), nu, T_N, sigma, Rel, cring, Delta->f_12 (O5)
%          of sec:notation-crosswalk applied.
% ERRATA:  E1 (oracle 6 -- the FIRST of its two landing sites in the running
%          text, and the one the superseded document put the defect in),
%          E23 (oracle 7's bad-prime citation).  App. A.

\paragraph{What this chapter fixes.}
Two things, and they are different in kind. First, the \emph{worked datum}: the coset
tables, the cocycle table and the induced block matrices for $g=\mathrm{diag}(1,2)$,
printed in full, so that \S\S\ref{sec:heckedatum}--\ref{sec:induced} have an instance
small enough to check by hand and so that
\S\ref{sec:chainlevel}'s forward reference has a target. Second, the \emph{oracles}:
the nine numerical unit tests of the superseded contract, re-homed with erratum
\ref{err:e1} applied to oracle~6 and erratum \ref{err:e23} to oracle~7, and the three
new families the boundary, the cone and the involution require. An oracle is not an
obligation: Obligation~\ref{obl:oracle} is the normative statement and this chapter is
the list of instances it reads, together with the registered validation that evaluates
each. \textbf{Nothing here may be cited as a licence to weaken an obligation}, and
Obligation~\ref{obl:oracle}(iii) binds oracle~6 specifically.

\paragraph{Standing hypotheses.}
\ref{H4} throughout, and the datum $g=\smat{1}{0}{0}{N}$, $N=p^{\nu}$,
$\Ga''=\Ga_0(N)$, $\{s_j\}$ of Definition~\ref{def:cosetreps}, $\tau$ of
Algorithm~\ref{alg:tau}. Per rule~\ref{A1} every oracle below repeats the hypotheses it
needs; in particular the eigenvalue oracles carry $6\in\cring^{\times}$, $p\neq\ell$
(\ref{H1}), $n$ even, and ``off the exceptional locus of
Remark~\ref{rmk:dickson}'', and the two oracles that carry none say so.

\subsection{The datum for \texorpdfstring{$T_2$}{T\_2}}\label{sec:t2datum}

Take $p=2$, $\nu=1$, $N=2$, $g=\smat{1}{0}{0}{2}$, $m=p+1=3$
(Proposition~\ref{prop:cosets}(1)). The canonical transporters of
Definition~\ref{def:cosetreps}, their classes in $\mathbb{P}^1(\Z/2)$, the coset
elements $\delta_j=s_j^{-1}g$ and the upper-triangular transporters $t_j$ are
\begin{equation}\label{eq:t2datum}
\begin{array}{c|cccc}
 j & s_j & \mathrm{cls}(s_j) & \delta_j=s_j^{-1}g & t_j\\\hline
 1 & I & (0{:}1) & \smat{1}{0}{0}{2} & I\\[0.4ex]
 2 & \smat{1}{0}{1}{1} & (1{:}1) & \smat{1}{0}{-1}{2} & T\\[0.4ex]
 3 & S=\smat{0}{-1}{1}{0} & (1{:}0) & \smat{0}{2}{-1}{0} & S
\end{array}
\end{equation}
Each $\delta_j$ has determinant $2$ and is integral; the right coset representatives of
Proposition~\ref{prop:cosets}(1) are $\smat{1}{0}{0}{2}$, $\smat{1}{1}{0}{2}$,
$\smat{2}{0}{0}{1}$. The $t_j$ are $T^{k}$ for $k=j-1$ up to $j=N=2$ and
$t_{N+1}=S$, with top-row classes $(1{:}0),(1{:}1),(0{:}1)$.

\begin{proposition}[The $p=2$ cocycle table]\label{prop:t2cocycle}
\PROVED{} With $\mathcal{G}=\{T,S,ST,TS\}$ the generator multiset of
\S\ref{sec:model} and the $s_j$ of \eqref{eq:t2datum}, the permutation
$j\cdot\gamma$ and the cocycle $h_j(\gamma)=s_j\gamma s_{j\cdot\gamma}^{-1}$ of
Definition~\ref{def:cosetreps} are
\begin{equation}\label{eq:t2cocycle}
\begin{array}{c|ccc|ccc}
\gamma & 1\cdot\gamma & 2\cdot\gamma & 3\cdot\gamma
& h_1(\gamma) & h_2(\gamma) & h_3(\gamma)\\\hline
T  & 1 & 3 & 2 & \smat{1}{1}{0}{1} & \smat{-1}{1}{-2}{1} & \smat{1}{-1}{0}{1}\\[1ex]
S  & 3 & 2 & 1 & \smat{1}{0}{0}{1} & \smat{1}{-1}{2}{-1} & \smat{-1}{0}{0}{-1}\\[1ex]
ST & 2 & 3 & 1 & \smat{1}{-1}{0}{1} & \smat{1}{0}{0}{1} & \smat{-1}{-1}{0}{-1}\\[1ex]
TS & 3 & 1 & 2 & \smat{1}{1}{0}{1} & \smat{1}{-1}{2}{-1} & \smat{-1}{0}{2}{-1}
\end{array}
\end{equation}
Every entry has determinant $1$ and even lower-left entry, i.e.\ lies in $\Ga_0(2)$;
the three permutations are the transpositions/$3$-cycles shown, and
Lemma~\ref{lem:cocycle}'s two identities hold on the table:
$j\cdot(TS)=(j\cdot T)\cdot S$ and $h_j(TS)=h_j(T)\,h_{j\cdot T}(S)$ for $j=1,2,3$.
\end{proposition}

\begin{proof}
Twelve products of $2\times2$ integer matrices, each displayed above; the class
$\mathrm{cls}(s_j\gamma)$ determines $j\cdot\gamma$ by
Proposition~\ref{prop:cosets}(4), and $h_j(\gamma):=s_j\gamma s_{j\cdot\gamma}^{-1}$ is
then read off. The two identities are four more products each. \COMP{} The same table
is recomputed from $p$ alone and compared entry by entry --- representatives, classes,
permutations and all twelve cocycles --- by \code{HeckeCosetTablesValidation} (oracle
\code{O16}), which reads it from this chapter and not from
\file{Hecke/HeckeCosetTables.cpp}.
\end{proof}

\noindent
\emph{The induced system.} By \eqref{eq:MH} the induced local system
$L''=L_{\Ga_0(2)}$ of Definition~\ref{def:inducedls} is the relation table of
\S\ref{sec:model} with each $M_\gamma=\rhon(\gamma)$ replaced by the
$3(n{+}1)\times3(n{+}1)$ block matrix
$M_{\Ga''}(\gamma)=\sum_{j}E_{j,\,j\cdot\gamma}\otimes\rhon(h_j(\gamma))$ read off
\eqref{eq:t2cocycle} --- one non-zero block per block row. Because
$\gamma\mapsto M_{\Ga''}(\gamma)$ is a homomorphism
(Proposition~\ref{prop:blockmatrices}) and $S^{2}=(ST)^{3}=-I$ in $\Ga$, the assembled
table must satisfy
\begin{equation}\label{eq:t2relcheck}
  M_{\Ga''}(S)^{4}=M_{\Ga''}(ST)^{6}=I,\qquad
  M_{\Ga''}(S)^{2}=M_{\Ga''}(ST)^{3}=M_{\Ga''}(-I),
\end{equation}
and these are assertions about the \emph{assembly}, not about $\Ga$: they fail if a
permutation or a block is transcribed wrongly. They are checked at
$n\in\{0,2,10\}$, $\ell\in\{5,691\}$ by \code{HeckeCosetTablesValidation}.

\noindent
\emph{$\alpha$ and $\omega$.} With the $t_j$ and $s_j$ of \eqref{eq:t2datum},
Propositions~\ref{prop:alpha} and \ref{prop:omega} give the block column and the block
row
\[
  \alpha=\begin{pmatrix}\rhon(I)\\ \rhon(T)\\ \rhon(S)\end{pmatrix}
  \colon\VR\to\VR^{\oplus3},
  \qquad
  \omega=\bigl(\rhon(I),\ \rhon\smat{1}{0}{-1}{1},\ \rhon\smat{0}{1}{-1}{0}\bigr)
  \colon\VR^{\oplus3}\to\VR ,
\]
the entries of $\omega$ being $\rhon(s_j^{-1})$, all integral and invertible over every
$\cring$ (no determinant-$N$ matrix is inverted; Obligation~\ref{obl:integral}). The
middle map is \emph{not} of this shape (Remark~\ref{rmk:nogo}): it is the per-simplex
block map $\Xig$ of Definition~\ref{def:xi}, determined by the certified $\tau$-table.

\begin{example}[The $\tau$-table at $p=2$, as emitted]\label{ex:t2tau}
\COMP{} Algorithm~\ref{alg:tau} at $p=2$, $\nu=1$ stores one chain per pair
$(j,\sigma)$ with $1\le j\le 3$ and $\sigma\in\orient{q}(K)$, i.e.\
$3\cdot(N_0+N_1+N_2)=3\cdot(7+11+5)=69$ chains --- small enough to read in a code
review, which is why the superseded contract printed the count. Run of 2026-09-16
(\code{HeckeTauTableSelfCheckValidation}, which writes the table to
\file{docs/implementation/tau\_table\_p2.txt}): $69$ stored chains, $15$ triangle fills,
\textbf{all $15$ solved over $\Z$ and none modulo $\ell$}, $B_{\mathrm{tab}}=2$ and
hence the certified bound $B=\max(B_{\mathrm{tab}},4N^{2})=16$ of
Lemma~\ref{lem:bounds}; every chain coefficient in $\{0,\pm1\}$; certification
(Obligation~\ref{obl:tau}) passed at all $69$ grid points, and $0$ of the $66$ grid
points of the wider self-check used the mod-$\ell$ solve.
\textbf{The $69$ slots are not $69$ independent constructions.} Step~0 of
Algorithm~\ref{alg:tau} builds one slot per $\Ga''$-orbit and derives the rest through
the shared (T1) helper; the emitted table marks each derived slot, and $\lift{v_3}=T\cdot\lift{v_0}$
is the first place where the superseded step list's
``one representative per orbit, since $\{(j,\sigma)\}\leftrightarrow$ simplices of
$\Ga''\backslash\Ku$'' is false (Appendix~\ref{app:changelog}, v0.18.0, under
\code{alg:tau}).
\end{example}

\subsection{A fully worked instance: \texorpdfstring{$(n,\ell,p)=(2,5,2)$}{(n,ell,p)=(2,5,2)}}
\label{sec:workedinstance}

At $n=2$ the coefficient module is $\VR=\Sym^2(\Fl^2)$, $d_V=3$, basis
$(u_0,u_1,u_2)=(x^2,xy,y^2)$. By Definition~\ref{def:rhon} the three blocks of $\alpha$
and the three of $\omega$ are, over $\Z$,
\[
  \rhon(I)=I_3,\quad
  \rhon(T)=\begin{pmatrix}1&1&1\\0&1&2\\0&0&1\end{pmatrix},\quad
  \rhon(S)=\begin{pmatrix}0&0&1\\0&-1&0\\1&0&0\end{pmatrix},
\]
\[
  \rhon(s_1^{-1})=I_3,\quad
  \rhon(s_2^{-1})=\begin{pmatrix}1&0&0\\-2&1&0\\1&-1&1\end{pmatrix},\quad
  \rhon(s_3^{-1})=\begin{pmatrix}0&0&1\\0&-1&0\\1&0&0\end{pmatrix},
\]
and the three coset elements act by
$\rhon(\delta_1)=\mathrm{diag}(1,2,4)$,
$\rhon(\delta_2)=\bigl(\begin{smallmatrix}1&0&0\\-2&2&0\\1&-2&4\end{smallmatrix}\bigr)$,
$\rhon(\delta_3)=\bigl(\begin{smallmatrix}0&0&4\\0&-2&0\\1&0&0\end{smallmatrix}\bigr)$.
All of these are reduced modulo $\ell=5$ before use, never after
(\S\ref{sec:reduction}). \COMP{} The pipeline on this instance emits
$\dim\Heq{0}=\dim\Heq{2}=0$, $\dim\Heq{1}=1$ and
\[
  \Amat{1}(T_2)=(9\bmod 5)=(4)=\bigl(1+2^{\,n+1}\bigr)\bmod 5 ,
\]
the Eisenstein eigenvalue of Obligation~\ref{obl:oracle}; and at $n=0$, $\ell=5$,
$\Amat{0}(T_2)=(3)=(p+1)$. Both are asserted, with the two evaluation routes
bit-identical and the runtime assertions of Obligation~\ref{obl:welldef} executed, by
\code{HeckeOperatorSmokeValidation} (run of 2026-09-16, exit code $0$).
\textbf{The cochain-level matrices are not printed here and must not be}: at $n=2$ they
are $21\times21$, $33\times33$ and $15\times15$, they depend on the choice of $s_j$ and
of $\tau$ (Remark~\ref{rmk:indeplevel}), and only $\Amat{q}$ is an invariant of the
datum.

\subsection{The nine oracles}\label{sec:oraclelist}

\noindent
Let $\Amat{q}(T_p)$ be the matrix of Algorithm~\ref{alg:assembly}, and write
$\mathrm{Eis}(p,n):=1+p^{\,n+1}$. Each oracle names its hypotheses, its assertion, and
the registered validation that evaluates it; the numbering $1$--$9$ is the superseded
contract's and is preserved, because Obligation~\ref{obl:oracle}(iii) cites
\emph{oracle~6} by number. Whether a given instance \textbf{aborts} or is
\textbf{recorded} is decided by Obligation~\ref{obl:oracle}(i),(ii) and not here.

\begin{enumerate}
\item \textbf{Trivial coefficients.} $n=0$, $\cring=\Fl$ with $\ell\ge5$, any $p$:
  $\dim\Heq{0}=1$, $\dim\Heq{1}=\dim\Heq{2}=0$, and $\Amat{0}(T_p)=(p+1)$
  --- the normalisation (T3) pins this (Remark~\ref{rmk:T3}). \PROVED{} for
  $\Amat{0}$: Remark~\ref{rmk:T3}'s computation
  $\Tdc{g}v=\sum_j\rhon(\delta_j)v=mv$ on invariants, with $m=p+1$.
  \code{HeckeOracleValidation} (\code{O1}).
\item \textbf{Odd $n$.} $n$ odd, $2\in\cring^{\times}$ (so any $\ell\neq2$), any $p$:
  $\Heq{q}=0$ for every $q$ and the pipeline returns $0\times0$ matrices in every
  degree. \PROVED{}: Remark~\ref{rmk:negI} --- $-I$ is central, lies in the stabiliser
  relation at $v_6$, and $\rhon(-I)=(-1)^{n}I=-I$. \code{HeckeOracleValidation}
  (\code{O2}). This oracle carries \textbf{no} hypothesis $\ell\ge5$ and is the one
  small-$\ell$ statement of the list.
\item \textbf{Weight $4$.} $n=2$, $\ell=5$ (so $6\in\cring^{\times}$), $p=2$:
  $\dim\Heq{1}=1$, $\Amat{1}(T_2)=(\mathrm{Eis}(2,2)\bmod5)=(4)$, $\dim\Heq{2}=0$; and
  $\Amat{1}(T_3)=(28\bmod 5)=(3)$. There are no cusp forms of weight $4$, so the whole
  of $\Heq{1}$ is Eisenstein. \code{HeckeOracleValidation} (\code{O3});
  \S\ref{sec:workedinstance} is this oracle worked out.
\item \textbf{Weight $10$.} $n=8$, $\ell\in\{5,7,101\}$, $p=2$: $\dim\Heq{1}=1$ and
  $\Amat{1}(T_2)=(513\bmod\ell)=(\mathrm{Eis}(2,8)\bmod\ell)$; $\dim S_{10}=0$.
  At $(n,\ell)=(8,7)$ the pair is \emph{on} the exceptional locus of
  Remark~\ref{rmk:dickson} ($8=\ell+1$) and is therefore \textbf{recorded}, not
  asserted. \code{HeckeOracleValidation} (\code{O4}).
\item \textbf{Weight $12$ --- the cusp form.} $n=10$, $\ell\in\{5,7,13,101,10007\}$
  (all off the locus: $10\notin\langle6,20\rangle$, $10\notin\langle8,42\rangle$;
  Remark~\ref{rmk:dickson}), $\ell\neq691$, $n$ even, $p\neq\ell$: $\dim\Heq{1}=3$ and
  \begin{equation}\label{eq:oracle5}
    \det\bigl(X-\Amat{1}(T_p)\bigr)\equiv
      \bigl(X-a_p(\cf)\bigr)^{2}\,\bigl(X-\mathrm{Eis}(p,10)\bigr)\pmod\ell ,
  \end{equation}
  with $\cf$ the normalised weight-$12$ cusp form (LMFDB \code{1.12.a.a}), the
  multiplicity $2$ coming from the two Eichler--Shimura summands
  (Obligation~\ref{obl:oracle}, Claim~\ref{cl:dimH1}(1)). The embedded constants are
  $a_2(\cf)=-24$, $a_3(\cf)=252$, $a_5(\cf)=4830$, $a_7(\cf)=-16744$, and
  $\mathrm{Eis}(p,10)=1+p^{11}$, so at $p=2$ the polynomial is
  $(X+24)^{2}(X-2049)$. \code{HeckeOracleValidation} (\code{O5}),
  \code{HeckeEndToEndValidation}.
\item \textbf{Hecke relations --- with erratum \ref{err:e1}.}
  $6\in\cring^{\times}$, $n\in\{0,2,10\}$, $\ell\in\{5,101\}$, in every degree $q$:
  \begin{enumerate}
  \item[(a)] $[\Amat{q}(T_2),\Amat{q}(T_3)]=0$;
  \item[(b)] $\Amat{q}\bigl(T_4^{\mathrm{full}}\bigr)=\Amat{q}(T_2)^{2}-2^{\,n+1}I$,
    with $T_4^{\mathrm{full}}$ obtained independently of the recursion by the
    double-coset assembly \eqref{eq:fullassembly}
    (\code{HeckeOperatorFamily::T\_p\_pow\_direct});
  \item[(c)] $\Amat{q}\bigl(T_{\Deprim_{4}}\bigr)=\Amat{q}(T_2)^{2}-2^{\,n+1}I-2^{\,n}I$,
    with $T_{\Deprim_{4}}$ the raw $6$-coset pipeline
    (\code{HeckeOperatorFamily::T\_p\_pow\_primitive\_direct});
  \item[(d)] $\Amat{q}(\diam)=p^{\,n}I$ and $\Amat{q}(T_6)=\Amat{q}(T_2)\Amat{q}(T_3)$.
  \end{enumerate}
  \code{HeckeRelationsValidation} (\code{O9}),
  \code{HeckeOperatorFamilySmokeValidation}.
  \begin{erratumbox}
  \textbf{This is the site erratum \ref{err:e1} was found in.} The superseded oracle~6
  read ``$A^{(q)}(T_4)=A^{(q)}(T_2)^2-2^{\,n+1}I$ with $T_4$ computed independently from
  $\Ga\,\mathrm{diag}(1,4)\,\Ga$ ($6$ cosets)'': $6=p^{2}+p$ cosets is the
  \emph{primitive} operator and $T_2^{2}-2^{\,n+1}I$ is the \emph{full} one, so the
  prescription as written fails against a correct implementation by exactly $2^{\,n}I$.
  Clauses (b) and (c) are the two correct statements, each with its emitting entry
  point named. \textbf{Obligation~\ref{obl:oracle}(iii) binds this oracle}: an oracle
  stated for a prime-power index must name which operator it is about, and this one
  does, in both directions. Entry~\ref{err:e1} carries the evidence and the entry-point
  table and is not restated here.
  \end{erratumbox}
\item \textbf{Bad prime $p=\ell$ --- with erratum \ref{err:e23}.} $p=\ell$
  (hypothesis \ref{H2}), $n\in\{0,2,10\}$. \emph{Asserted for every $\ell$:} the
  pipeline runs to completion and certifies, i.e.\ Obligations~\ref{obl:welldef},
  \ref{obl:integral} and \ref{obl:tau} are discharged, and $U_\ell:=T_\ell$ is well
  defined. \emph{Asserted for $\ell\ge5$ only} (the shipped instance is $p=\ell=5$,
  $n=10$): $T_{\ell^{2}}=U_\ell^{2}$ for the \emph{full} operator, which is
  Corollary~\ref{cor:ulrecursion}(1) and needs $\ell\ge5$; and the ordinarity
  stabilisation $\rank\Amat{q}(U_\ell^{r})$ non-increasing and stationary for
  $r\ge\dim\Heq{q}$. \emph{Recorded, never asserted, at $\ell\in\{2,3\}$:} the same two
  quantities, and every eigenvalue statement, because $6\notin\cring^{\times}$
  (\ref{H3}). \code{HeckeRelationsValidation} and \code{HeckeBadPrimeValidation}
  (\code{O10}).
  \begin{erratumbox}
  The superseded oracle~7 asserted $U_4=U_2^{2}$ at $p=\ell=2$ ``(Corollary
  \code{cor:ulrecursion})''. Corollary~\ref{cor:ulrecursion}(1) carries $\ell\ge5$,
  inherited from Theorem~\ref{thm:relations}, and clause~(2) says something else
  entirely (full $=$ primitive for $n\ge1$); neither licenses the assertion at
  $\ell=2$. Entry~\ref{err:e23}; the implementation already records rather than asserts
  it there (amendment A-F1, \code{HeckeRelationsValidation}:273--330), so the defect is
  in the contract text alone.
  \end{erratumbox}
\item \textbf{Ramanujan congruence and non-semisimplicity.} $n=10$, $\ell=691$, $p=2$.
  Since $2049+24=3\cdot691$, \eqref{eq:oracle5} degenerates to $(X+24)^{3}\bmod691$;
  the assertion with content is that $\Amat{1}(T_2)\neq-24\,I$, i.e.\ the minimal
  polynomial is not $X+24$ --- the Eisenstein/cuspidal extension is non-split. This is
  the torsion-sensitive part of the pipeline; cf.\ \cite{deng}. The Jordan partition
  (expected $(2,1)$) is \textbf{recorded}, not asserted. \code{HeckeOracleValidation}
  (\code{O6}), \code{HeckeBlockFormValidation}.
\item \textbf{Internal agreement, and restriction compatibility.}
  (a) Cross-route agreement and pure-integer-versus-filtered agreement are
  Obligations~\ref{obl:agreement} and \ref{obl:fp}; they are \emph{bit}-identity
  statements, not eigenvalue statements, and are listed here only so that the oracle
  roster is complete (\code{HeckeIndependenceValidation}, \code{HeckeFpFilterValidation},
  \code{HeckeCrossCheckValidation}). (b) \emph{Restriction compatibility:} for
  $\ell\ge5$, $n\in\{2,10\}$, $p=2$, and only when the Eisenstein and cuspidal $T_2$
  eigenvalues are distinct modulo $\ell$, the composite
  $\Heq{1}(K,\Rel;L)\to H^{1}(\Kbdry;L|_{\Kbdry})$ has rank $1$: the eigenvector of
  eigenvalue $\mathrm{Eis}(p,n)$ maps to a non-zero class and the cuspidal generalised
  eigenspace maps to $0$ \STD{\cite[Ch.~2--3]{harder3} for the Eisenstein class and its
  boundary eigenvalue}. \code{HeckeTessellationValidation},
  \code{HeckeMorphismRegressionValidation} (\code{O15}).
  \textbf{What this clause is about.} The object on the right is the \emph{circle
  model} $\Kbdry$ of Proposition~\ref{prop:circlemodel}, i.e.\ the
  $2\in\cring^{\times}$ specialisation of Definition~\ref{def:bdrygroup}, and the
  composite is $\rres^{1}$ of \S\ref{sec:cone} \emph{without} any Hecke structure on the
  target. It is therefore \emph{not} the boundary oracle of \ref{obl:taubdry}(vi), and
  it is not evidence for Obligation~\ref{obl:coneexact}(iii).
\end{enumerate}

\begin{remark}[Two implementations of the Dickson predicate]\label{rmk:dicksontwo}
\COMP{} The gating predicate of Remark~\ref{rmk:dickson} occurs twice in the tree under
the same name. \code{HeckeEndToEndValidation.cpp}'s \code{IsDicksonExceptional} is the
semigroup membership $n\in\langle\ell+1,\ell^{2}-\ell\rangle$, which is the predicate
Remark~\ref{rmk:dickson} states and pins; \code{HeckeOracleValidation.cpp}'s copy tests
only the first generator, $(\ell+1)\mid n$. On the grid that second copy is used on ---
$n\in\{0,2,4,6,8,10,12,14,22\}$ against $\ell\in\{5,7,13,101,691,10007,134217757\}$ ---
the two predicates \textbf{agree at every point}, checked exhaustively on 2026-09-16;
the exceptional points there are $(6,5)$, $(12,5)$ and $(8,7)$, all caught by either
test. The weaker copy is therefore sound where it is used and unsound in general
(Remark~\ref{rmk:dickson}'s counterexample $(20,5)$ is outside its grid and inside the
other's). Recorded so that extending \code{HeckeOracleValidation}'s grid is known to
require replacing it.
\end{remark}

\subsection{The new oracle families: \texorpdfstring{$\partial$}{bdry}, the cone,
and \texorpdfstring{$\inv$}{iota}}\label{sec:oraclesnew}

\noindent
\textbf{One of these three families is a specification, not a run; the other two have
been run.} Chapter \ref{sec:cone} is implemented only in its $\etapair$-half --- the
$\etapair$-table exists since v1.4.2, no $\Phig$ and no $\Acone{q}$ does (its code map says
so) --- so family (B) below is what the phase that builds it
must evaluate, written as numbers so that they can be evaluated rather than argued about,
and it may not be read as evidence that its obligation is discharged. Family (A), the
boundary oracles, \textbf{was run at v1.3.3} (2026-09-16, \code{HeckeBoundaryOperatorValidation};
until v1.3.2 it, too, was a specification, chapter \ref{sec:boundarydatum} being
unimplemented) and its own paragraph records the outcome. Family (C), the
involution oracles, \textbf{was run at v1.1.1 and again at v1.2.0} and its own
paragraph records both outcomes --- the second under the normalisation
$\inv^{\ast}=-\inv^{\ast}_{\mathrm{geo}}$ of Definition~\ref{def:iota}(iv) that the first
motivated; Obligation~\ref{obl:iota} is evaluated in full and holds. Each family names
the obligation it serves.

\begin{enumerate}[label=(\Alph*)]
\item \textbf{Boundary oracles} (Obligation~\ref{obl:taubdry}(vi)). For $n$ even,
  $p\neq\ell$: on $H^{0}(C^{\bullet}_{\partial})$ in the basis
  $x^{\,n-b\ell}\mathrm{Nm}^{\,b}$ the eigenvalues of $\Thb_g$ are the characters
  $\chi_b(p)$ of Proposition~\ref{prop:bdrychar}(1) --- \emph{asserted}; for $n<\ell$
  the class of $y^{n}$ in $H^{1}(C^{\bullet}_{\partial})$ has eigenvalue
  $\mathrm{Eis}(p,n)=1+p^{\,n+1}$ (Proposition~\ref{prop:bdrychar}(3)) ---
  \emph{asserted}; the general-$n$ $H^{1}$ multiset of Claim~\ref{cl:chi}(2) ---
  \emph{recorded} until that claim is proved. \textbf{The rank is
  $\lfloor n/\ell\rfloor+1$, not $1$}, so an oracle written for $n<\ell$ does not extend
  by inspection: at $(n,\ell)=(10,5)$ the space $\VR^{T}$ is \emph{three}-dimensional,
  with the three eigenvalues $\chi_0,\chi_1,\chi_2$ of
  Proposition~\ref{prop:bdrychar}(1). The first instance to run is
  $(n,\ell,p)=(2,5,2)$, where $n<\ell$ and $n$ is even, so
  Proposition~\ref{prop:bdrychar}(3) applies: $H^{0}$ and $H^{1}$ are both
  one-dimensional, spanned by $[x^{2}]$ and $[y^{2}]$, and both carry the eigenvalue
  $\chi_0(2)=1+2^{3}=9\equiv4\pmod 5$. That number is the point: it is the same $4$ as
  $\Amat{1}(T_2)$ of \S\ref{sec:workedinstance}, and
  Obligation~\ref{obl:coneexact}(iii) --- $\rres^{1}\Amat{1}=\Abdry{1}\rres^{1}$ ---
  is what forces the agreement. \COMP{} \textbf{Run at v1.3.3}
  (\code{HeckeBoundaryOperatorValidation}, $546$ cells): every asserted number above
  holds --- $\chi_b(p)$ on each $x^{\,n-b\ell}\mathrm{Nm}^{\,b}$ ($575$ vectors on $330$
  cells, the rank $\lfloor n/\ell\rfloor+1$ realised, three eigenvalues at $(10,5)$),
  $1+p^{\,n+1}$ on $[y^{n}]$ for $n<\ell$ ($232$ cells) --- the first instance
  $(2,5,2)$ prints $\Abdry{0}=[4]$, $\Abdry{1}=[4]$, bulk $\Amat{1}(T_2)=[4]$,
  $\rank\rres^{1}=1$, and the general-$n$ $H^{1}$ multiset of Claim~\ref{cl:chi}(2)
  agrees on $330/330$ cells --- recorded, the claim's tag unchanged. The right-coset
  substitution of Proposition~\ref{prop:bdrycosets}(4), run as the negative control,
  returns $\chi_{b+1}$ and breaks Obligation~\ref{obl:coneexact}(iii) on $170$ cells, as
  Proposition~\ref{prop:bdrychar}(2) says it must (\S\ref{sec:boundarydatum}, code map).
\item \textbf{Cone oracles} (Obligation~\ref{obl:coneexact}(iv),(v)). At
  $6\in\cring^{\times}$, $n$ even, off the Dickson locus, the six-term sequence
  \eqref{eq:les} with $\Heq{2}=0$ gives, by \eqref{eq:dimbook}, a dimension book the
  implementation must reproduce exactly. The first instance to run is again
  $(n,\ell,p)=(2,5,2)$, where every term is known independently:
  $\dim X_1=\dim\Heq{0}=0$ (Remark~\ref{rmk:dickson}: $2\notin\langle6,20\rangle$),
  $\dim X_2=\dim H^{0}(C^{\bullet}_{\partial})=1$ and
  $\dim X_5=\dim H^{1}(C^{\bullet}_{\partial})=1$ by
  Proposition~\ref{prop:bdrychar}(3), $\dim X_4=\dim\Heq{1}=1$
  (\S\ref{sec:workedinstance}), and $\rank\rres^{1}=1$ with $\ker\rres^{1}=0$
  (oracle~9(b), and $S_4=0$, so all of $\Heq{1}$ is Eisenstein); whence
  \eqref{eq:dimbook} forces $\dim X_3=\dim\Hcpt{1}=1-0+0=1$ and
  $\dim X_6=\dim\Hcpt{2}=1-1=0$. At $n=10$ the same book is to be run with
  $\dim X_2=\dim X_5=\lfloor10/\ell\rfloor+1$, which is $3$ at $\ell=5$ and $1$ at
  $\ell\in\{13,101\}$: the $\ell\le n$ rows are where a boundary implementation is
  actually exercised, and they are not $n<\ell$ rows in disguise. In addition
  $\Hcpt{0}=0$ unconditionally (Lemma~\ref{lem:conerelative}(4)), every composite of
  consecutive maps vanishes, and $\Acone{q}$ satisfies the relations of
  Obligation~\ref{obl:coneexact}(vi) --- including clause (b) of oracle~6 above with
  $\Amat{q}$ replaced by $\Acone{q}$. The negative control of
  Obligation~\ref{obl:coneexact}(viii)(b) is part of this family and is the one arm
  that must \emph{fail} a clause to pass.
\item \textbf{Involution oracles} (Obligation~\ref{obl:iota}(v)). For $\ell\ge5$ off
  the Dickson locus and $n$ even, $\dim H^{1}_{\mathrm{eq},+}=d_{\,n+2}$ and
  $\dim H^{1}_{\mathrm{eq},-}=d_{\,n+2}-1$ with $d_k$ the dimension of level-$1$
  weight-$k$ modular forms; at $n=10$, $k=12$, $d_{12}=2$, so the split of
  $\dim\Heq{1}=3$ is $2+1$, and the $T_p$-eigenvector of eigenvalue
  $\mathrm{Eis}(p,10)$ lies in the $+$ part. At $n=2$, $k=4$, $d_4=1$, so the split of
  $\dim\Heq{1}=1$ is $1+0$. The dimensions are \textbf{recorded}, not asserted
  (Obligation~\ref{obl:iota}(v)); \emph{asserted} there are
  $\inv^{\ast}\Amat{q}=\Amat{q}\inv^{\ast}$ for $6\in\cring^{\times}$ and, since v1.2.0,
  the sign itself --- the $T_p$-eigenvector of eigenvalue $\mathrm{Eis}(p,n)$ lies in the
  $+$ part on every cell on which it can be located (the Dickson cells carry no sign),
  which under $\inv^{\ast}=-\inv^{\ast}_{\mathrm{geo}}$
  (Definition~\ref{def:iota}(iv)) is a consistency statement about the definition. The
  degree-$0$ instance is exact and needs no cohomology: $\inv^{\ast}$ on
  $\Heq{0}=\VR^{\Ga}$ is $-\rhon(\inv)$ restricted to the invariants, which at $n=0$ is
  $-\id$.
  \COMP{} \textbf{Run at v1.1.1 and at v1.2.0 by \code{HeckeIotaSignValidation}.} At
  v1.1.1, for the then-current involution $\inv^{\ast}_{\mathrm{geo}}$, the asserted part
  held and the recorded prediction came out as the exact mirror --- the Eisenstein
  eigenvector in the $-$ part, the split $(d_{\,n+2}-1,d_{\,n+2})$: $1+2$ at $n=10$,
  $0+1$ at $n=2$ --- on every cell of a $77$-cell grid where a sign is defined, with
  $\inv^{\ast}_{\mathrm{geo}}=+1$ at $n=0$ pinning the global sign. That finding is what
  owner decision OQ3 of 2026-09-16 acted on (Remark~\ref{rmk:iotasign}). \textbf{At
  v1.2.0, under the normalised $\inv^{\ast}$, every number above holds as printed}: the
  Eisenstein eigenvector is in the $+$ part on all $74$ signed cells (gating), the split
  is $(d_{\,n+2},d_{\,n+2}-1)$ on all $66$ off-Dickson cells --- $2+1$ at $n=10$, $1+0$ at
  $n=2$ --- and $\inv^{\ast}=-1$ at $n=0$ (gating). The two runs are one computation read
  through $\inv^{\ast}=-\inv^{\ast}_{\mathrm{geo}}$; the validation prints both.
  \emph{The diagnostic of Obligation~\ref{obl:iota}(iv) had been run before this family
  and was false on the then-canonical $\tau$; on the canonical $\tau$ it holds and
  gates since v1.1.0} (Remark~\ref{rmk:iotatausym}). It is not part of this family and
  may not be used to skip clause (v) --- and was not: (v) was evaluated independently.
\end{enumerate}

\paragraph{Code map.}
The worked datum is \code{hecke::HeckeCosetTables} (representatives, class map,
permutation and cocycle, \file{Hecke/HeckeCosetTables.h}:21 citing this section) and
\code{hecke::HeckeTauTable}; the printed table of Example~\ref{ex:t2tau} is written by
\code{HeckeTauTableSelfCheckValidation}. Of the oracles, $1$--$5$ and $8$ are
\code{HeckeOracleValidation} (its families \code{O1}--\code{O8}, \code{O14}), $6$ and
$7$ are \code{HeckeRelationsValidation} together with
\code{HeckeOperatorFamilySmokeValidation} and \code{HeckeBadPrimeValidation}, $9$(a) is
\code{HeckeIndependenceValidation} / \code{HeckeCrossCheckValidation} /
\code{HeckeFpFilterValidation} and $9$(b) is \code{HeckeTessellationValidation} /
\code{HeckeMorphismRegressionValidation}; \code{HeckeEndToEndValidation} re-runs the
whole matrix and prints the obligation-to-test coverage table. The three families of
\S\ref{sec:oraclesnew} were, until v1.3.3, two without an implementation and one with;
since v1.3.3 family (A) is run (\code{HeckeBoundaryOperatorValidation}), since v1.4.2
the $\etapair$-table of family (B) exists while no cone assembly does, and family (C) is
\code{HeckeIotaSignValidation}, which since v1.1.1 discharges
Obligation~\ref{obl:iota}(ii),(iii),(v).
\code{sec:example} is cited from $16$ sites in $8$ source files
(\S\ref{sec:machinery}, census of 2026-09-16), so this section's numbering is
load-bearing: renaming oracle~6 would silently retarget them.

\paragraph{Formalization note.}
Proposition~\ref{prop:t2cocycle} and \eqref{eq:t2relcheck} are finite computations in
$\Mat_2(\Z)$ and $\GL_{3(n+1)}$, decidable as stated, and are the cheapest thing in
this chapter to formalise --- but also the least valuable, since they instantiate
Lemma~\ref{lem:cocycle} and Proposition~\ref{prop:blockmatrices}, which are the
statements worth having. Oracles $1$ and $2$ are the two worth formalising: each is a
\emph{theorem} of this document (Remark~\ref{rmk:T3}, Remark~\ref{rmk:negI}) with a
two-line proof, not an imported constant, and together they pin the normalisation and
the odd-$n$ vanishing that every other oracle is read against. Oracles $3$--$5$, $8$ and
$9$(b) depend on Eichler--Shimura and on embedded $a_p(\cf)$; a formalisation should
carry the constants as hypotheses and prove only the \emph{shape} of
\eqref{eq:oracle5}, which is Obligation~\ref{obl:oracle}. Oracle~6 is finite linear
algebra over $\cring$ once $\Amat{q}$ exists, and clauses (b),(c) are the pair whose
\emph{difference} $p^{\,n}I$ is the content (Proposition~\ref{prop:diamond}); that
difference, not either identity alone, is what a formalisation should state, because it
is what entry~\ref{err:e1} exists to stop anyone getting wrong again.

%----------------------------------------------------------------------
\section{Implementation notes: mapping to the code base}\label{sec:implementation}
%----------------------------------------------------------------------
% SOURCE:  hecke_action.tex:1613-1690 (RH), extended with the new chapters'
%          code maps; the shipped tree (CMakeLists.txt, Hecke/README.md,
%          Campaign/, ValidationRegistry.cpp) read on 2026-09-16 (NEW).
% NOTATION: rows Rel, cring, nu, T_N of sec:notation-crosswalk applied.
% ERRATA:  none of its own; it records where erratum E1's three entry points
%          live (App. A, err:e1).

\paragraph{What this chapter fixes.}
Three things the per-chapter code maps do not: the \emph{aggregate} map from this
document to the tree and back, in one table; the \emph{configuration surface} --- which
switches exist, and the standing rule that none of them may change a computed value;
and the \emph{record schema}, i.e.\ what an emitted run must contain for the
obligations of \S\ref{sec:obligations} to be checkable after the fact. It closes with
\S\ref{sec:machinery}, the guards on this document itself.

\begin{contractbox}
\textbf{No switch in \S\ref{sec:buildsurface} may change a computed value.} Every build
option, environment variable and command-line flag listed there selects a \emph{route},
a \emph{backend}, a \emph{schedule} or an \emph{amount of certification} --- never an
answer. The two that select a route are exactly the two rollbacks
(\code{HECKE\_PURE\_INTEGER}, Obligation~\ref{obl:fp}(iv); \code{ENGINE\_BACKEND}), and
for both the contract is bit-identity of every stored object, not agreement to within
anything. The one that changes what is \emph{emitted} rather than what is computed is
the certification level (\code{--lean}, \code{certify}), and \S\ref{sec:recordschema}
requires the record to say which level ran.
\end{contractbox}

\subsection{The code map}\label{sec:codemap}

\noindent
Chapter by chapter, the translation units that realise it. ``---'' means the chapter has
no implementation; \S\ref{sec:notimplemented} says what that costs.

\begin{center}
\footnotesize
\begin{tabular}{@{}p{3.1cm}p{11.1cm}@{}}
\toprule
chapter & translation units (paths relative to the repository root) \\
\midrule
\ref{sec:framework}--\ref{sec:systematiccomputation} (Part~I)
 & \file{src/hacs/}: \code{SimplicialCohomology}, \code{LocalSystemSimplicialComplex} (the $(K,\Rel)$ datum, \code{SimplexRelation}), \code{EquivariantCohomologyOnK} (the $C/Z/B/H$ bases and $P_q$ of \ref{def:storedbases}, \ref{prop:ACq}, \ref{prop:AZq}, \ref{prop:projector}, \ref{alg:cohomology}), \code{EquivariantCochainComplexMorphismOnK} (\code{FiberBlockTable}; the \ref{err:e19} fix) \\
\addlinespace
\ref{sec:conventions} (model) & \code{LocalSystemSL2Z} (\code{SL2ZRelationDatum}, \code{GenerateSL2ZRelationsWithGroupData}); \file{Hecke/SL2ZTessellationTables} (the closure and the canonical key of \ref{def:primitives}(P4)); \file{Hecke/HeckeModelDescriptor} (the $7/11/5$ literal is now a descriptor) \\
\addlinespace
\ref{sec:symconv} ($\Sym^n$) & \code{ComputeSymmetricPowerMatrixFM} (\file{src/ModArithmetic/FlatMatrix.cpp}) --- \eqref{eq:rhondef} literally; \code{LocalSystemUtils}. \textbf{Since v1.7.3} (slot-6 WP-1, phase E7a) the contragredient: \code{Mat22::adj} (\file{Hecke/ExactInt}), \code{ComputeContragredientPowerMatrixFM} and \code{ComputeFibreMatrixFM} (\code{LocalSystemUtils}), \code{GenerateSL2ZContragredientRelationsWithGroupData} (\code{LocalSystemSL2Z}), \file{Hecke/HeckeContragredient} ($\Phi_n$ of \ref{lem:symdual}, the pairing of \ref{prop:pairing}), the \code{FibreRepresentation} argument of \file{Hecke/HeckeOperatorOnK} ($A(\Dn)(T_p)^{q}$); since v1.7.3 (slot-6 WP-3, phase E7b) \file{Hecke/HeckePairing} --- the relative fundamental $2$-chain, the cone-level and circle pairings, \code{PairingCertificate}, \code{PairingAdjoint}, i.e.\ \ref{cl:pd}(1)--(3) tested and the PD gate of \ref{obl:coneexact}(vii); \code{HeckeDualityValidation}. \textbf{Since v1.7.4} (slot-7 WP-1) phase E7 is \textbf{accepted}: \ref{cl:pd}(1)--(2) tested on all $54$ cells of the extension stage's case list at $\ell\ge5$ (\code{HeckeDualityValidation} arm A16 on $34$ of them, the \code{slow} \code{HeckeDualityCoverageSlowValidation} on three, two more under an environment switch), failing on none, the twisted controls rejected; \ref{cl:pd} stays open --- a test, not a proof; the E7 objects are golden since the stage's add-only emission of slot~7 ($651$ entries). (At v1.7.2: not built, scheduled as phase E7.) \textbf{Since v1.8.1} (slot-8 WP-3, phase E8, status) $\rho_n$ over $\Z/\ell^{k}$: \code{ComputeSymmetricPowerMatrixRing} (\file{src/ModArithmetic/FlatMatrix.cpp}) --- \eqref{eq:rhondef}'s integer entries reduced modulo $\ell^{k}$, by the unit recurrences where the divisor is a unit, a division-free Pascal route otherwise, and the split $\mathrm{diag}(\ell,1)\cdot g'$ for $U_\ell$; the $\Fl$ builder is unchanged \\
\addlinespace
\ref{sec:involution} ($\inv$) & \file{Hecke/HeckeIotaTauSymmetryValidation} --- $\invK$, its signs and $\tau^{\inv}$, for Obligation~\ref{obl:iota}(i),(iv), (iv) gating on the canonical $\tau$ since v1.1.0 --- and, since v1.1.1, \file{Hecke/HeckeIotaSignValidation} --- $\inv^{\ast}_{\mathrm{geo}}$ and $\inv^{\ast}=-\inv^{\ast}_{\mathrm{geo}}$ on cochains and on $\Heq{\bullet}$, the $\pm$ split and the Eisenstein sign (gating and true by definition since v1.2.0, the geometric sign recorded), for Obligation~\ref{obl:iota}(ii),(iii),(v). Both are validations. \textbf{Since v1.6.2 (slot-4 WP-3, phase E3) the chapter has production code}: \file{Hecke/HeckeInvolution} --- the \code{EquivariantCochainComplexMorphismOnK} this chapter's code map specifies, one $-\varepsilon_\sigma\rhon(\inv)$ block per simplex, \code{compute()} discharging Obligation~\ref{obl:iota}(i),(ii),(iii) blockwise and exactly over $\Z$ and aborting on failure --- and \file{HeckeExt/HeckeSignSplit} --- clause (v), $A^{(q)}_{\pm}$, the per-sign decomposition and the Eisenstein line --- gated entrywise against the two validations' own construction by \code{HeckeInvolutionValidation}. \textbf{Since v1.7.2 (slot-5 WP-1 and WP-12a, phase E2) the chapter's objects are in the record}: \code{hecke\_ext/v2} carries the per-sign objects with \code{iota\_sign}, and, per the owner's ruling of 2026-09-22, the record-level \code{iota\_exact\_comparison} --- Obligation~\ref{obl:iota}(ii)'s rider as an emitted object (\code{ran} / \code{relation\_rows} / \code{rows\_with\_exact} / \code{rows\_compared}), refused by \code{HeckeExtRecord::verify()} code~23 if it misreports. \code{HeckeExtRecordValidation}. \textbf{Planned at v1.20.0-pre.1} (slot-20 WP-LV3(d), \Espec{rmk:lviotak}, \Espec{chk:lvtriangle}): \file{Hecke/HeckeInvolution} and the console honour $\inv^{\ast}$ over $\cring=\Z/\ell^{k}$ (the console's exit~$5$ removed); Proposition~\ref{prop:iotasplit} ($2\in\cring^{\times}$) and Proposition~\ref{prop:iotacommutes} already cover it, so no statement of this chapter changes. \textbf{Built at v1.20.0-pre.2} (WP-LV3, merged \code{c1fda051}): \code{hecke::BuildRingInvolution} (record \code{hecke\_involution\_ring/v1}; $\inv^{2}=\mathrm{id}$ and the length split checked; $\ell\in\{2,3\}$ refused), the console's \code{iota} on a $\Z/\ell^{k}$ handle (\Espec{sec:ring}~(h)), script \code{049\_lv3\_iota\_ring.oc}; gated in \code{HeckeLValueValidation} (arm \code{LV3-I}). \\
\addlinespace
\ref{sec:heckedatum}, \ref{sec:chainlevel} & \file{Hecke/HeckeCosetTables} (\ref{def:cosetreps}, \ref{lem:cocycle}), \file{Hecke/HeckeTauTable} (\ref{def:tau}, \ref{alg:tau}, \ref{obl:tau}, certified in the constructor), \file{Hecke/HeckeOperatorFamily} (\ref{def:operators}, \ref{prop:diamond}, \ref{thm:relations}, \ref{cor:ulrecursion}; the three entry points of \ref{err:e1}); since v1.4.2 \file{Hecke/HeckeIntegerSolve}, \ref{alg:tau} step~3's integer solver, moved verbatim out of \file{HeckeTauTable.cpp} by slot-2 WP-8 so that the $\etapair$ fills use the one solver (value-preserving: golden and \file{tau\_table\_p2.txt} unchanged). Since v1.5.2 (slot-3 WP-1, the $U_\ell$ reach repair recorded in this chapter's code map): \ref{alg:tau} step~1 runs through \code{hecke::escalate\_on\_overflow} on \code{QuadPointT<I>} (\file{Hecke/QuadPoint}, \file{Hecke/ExactInt}; \ref{lem:bounds}(6), \ref{thm:backend}(iv)); \file{HeckeExt/HeckeUEllBudget}, the frozen memory admission of $U_\ell$ (\code{u\_ell\_rss/v2-2026-09-17}), which realises no statement of this document; \code{HeckeUEllReachValidation}, \code{HeckeUEllReachSlow}, \code{HeckeUEllBudgetValidation} \\
\addlinespace
\ref{sec:induced}, \ref{sec:equivalence} & \file{Hecke/InducedLocalSystem} (\ref{def:inducedmodule}, \eqref{eq:MH}, \ref{prop:lhcompliant}), \file{Hecke/HeckeOperatorOnK} (\ref{def:xi}, \eqref{eq:xiexplicit}, \eqref{eq:heckechain}, \ref{alg:assembly}; both routes), \file{Hecke/HeckeCrossChecks} and \file{Hecke/BarCocycleH1} (\ref{prop:barformula}, \ref{rmk:nogo}) \\
\addlinespace
\ref{sec:boundarydatum} & since v1.3.3: \file{Hecke/HeckeBoundaryDatum} (\ref{def:bdrygroup}, \ref{prop:bdrycosets}, \ref{def:taubdry}, \ref{lem:taubdryexists}, \ref{obl:taubdry}(i)--(iv), certified in the constructor), \file{Hecke/HeckeBoundaryOperator} (\ref{def:bdryop}, \ref{prop:bdryop}, \ref{prop:bdrychar}, \ref{prop:circlemodel}(2),(3) recorded, \ref{obl:taubdry}(v)--(vii); both routes), the boundary subcomplex of \file{Hecke/SL2ZTessellationTables} (\ref{B1}--\ref{B3}) and \code{HeckeCosetTables::twisted\_conjugate} (the shared (T1) step); the circle model as \code{LocalSystemSL2Z::CreateSL2ZBoundaryPathTriangulation} / \code{GenerateSL2ZBoundaryRelationsElt}. Since v1.4.2 (slot-2 WP-7, phase E4b): \file{HeckeExt/HeckeBoundary} --- $H^{1}_{!}=\ker\rres^{1}$ of \ref{cl:les}(2) with its Hecke action (an exact solve, $T_p$-stable by \ref{obl:taubdry}(vii)), the joint $1+p^{n+1}$ eigenspace, Harder's criterion and the Harder split, per cell, recorded and never emitted (the record is E2's) --- and \file{Hecke/HeckeBoundaryDiagnostics}, the two recorded diagnostics \code{cochain\_commutes} / \code{eta\_zero} of \ref{rmk:conestrict}(2), both false on every $(\text{cell},p,\nu)$ of the grid; \code{HeckeBoundaryParabolicValidation}. \textbf{Recorded at v1.12.1} (slot-12 S12-FIN): \file{HeckeExt/HeckeFinalityRoutes} --- \code{BoundaryRoute} and \code{BoundaryFinalityCertificate}, the finality certificate (d1) ``boundary'' of \Espec{cert:boundary} ($\rres^{1}$ injective on the block, then each class representative $T_p$ in the local algebra), which rests on Proposition~\ref{prop:bdrytranscription} and Lemma~\ref{lem:bdrytorus}; called by \file{HeckeExt/HeckeExtDriver} (record \code{hecke\_ext/v4}); \code{HeckeBoundaryFinalityValidation} \\
\addlinespace
\ref{sec:cone} & \textbf{implemented in full since v1.7.3} (except \ref{obl:coneexact}(vii) from v1.5.2 to v1.7.2): \file{Hecke/HeckePairHomotopy} (slot-2 WP-8, phase E5-1: \ref{def:pairhomotopy}, \ref{lem:etaexists}, \ref{lem:etaunique}, \ref{obl:eta}(i)--(iv) certified in the constructor, (vi)'s \code{eta\_zero} recorded; since slot-3 WP-6 with \code{EtaVertexPolicy} and \code{table\_fnv64()}) --- the $\etapair$-table over $\Z$ for every $\ell$, no mod-$m$ path; \file{EquivariantMappingCone} and \file{Hecke/SL2ZBoundaryPair} (slot-3 WP-5, phase E5-2: \ref{def:pair}, \eqref{eq:conediff}, \eqref{eq:coneles}, \ref{lem:conerelative}(4), \eqref{eq:conedims}, \ref{obl:coneexact}(iv),(v)); \file{Hecke/HeckeConeAction} (slot-3 WP-6, phase E5-3: \ref{def:coneaction}, \eqref{eq:derivedh}, \eqref{eq:conemap}, \ref{lem:conesign}, \ref{obl:eta}(v),(vi), \ref{obl:coneexact}(i)--(iii),(vi),(viii)); \code{TauBoundaryVertexPolicy} in \file{Hecke/HeckeBoundaryDatum}; \code{HeckeConeValidation}, \code{HeckeConeActionValidation}. v1.6.2 (slot-4 WP-4, E5 cost): three \emph{value-preserving} re-implementations inside those units --- \code{ConeCertificate}'s span inclusions C2 / C3 / C5 computed as the rank identity $\rank[B\mid Y]=\rank B$ off one elimination instead of a materialised solve, the per-degree buffers of \code{compute()} and \code{verify()} released at their last reader, and the action's standard-basis sweep batched on the identity in column blocks with one $\rhon$ memo per degree. \textbf{No clause was sampled, relaxed or moved out of \code{compute()}}: every item of both certificates still runs unconditionally, $A_c^{(q)}$ is bit-identical and both validations' whole printed output is byte-identical to the pre-slot-4 binary's. v1.7.2 (slot-5 WP-4, WP-5, WP-6, WP-12a; all four value-preserving, on the same evidence): the $\ell\in\{2,3\}$ carve-out on \ref{obl:coneexact}(iii)'s $\rres$-half is \emph{dropped in the code}, which now asserts it at every $\ell$ (v1.7.0's ruling, realised); \code{ConeCertificate} item~11 evaluates \eqref{eq:conedims}'s \emph{first} identity without consulting a stored $H^2_{\mathrm{eq}}=0$, each identity under its own minimal guard, the item code kept; the cone's $Z/B/H/P$ loop and \code{ConeCertificate}'s span inclusions run on \code{engine::tape}; \file{Hecke/RhoMemo.h} supplies the three shared $\rhon$ memo tables of \code{HeckeConeAction}'s standard-basis sweep; and \code{ConeActionCertificate} gains \textbf{item~12, record honesty} (rule~\ref{A6}): the emitted \code{restriction\_equivariance\_asserted} of \code{hecke\_cone\_action/v1} must report the assertion item~8 performs, so it reads \code{true} at every $\ell$ and a certificate that says otherwise fails item~12 alone. \textbf{Since v1.7.2 the cone is in the record}: \code{hecke\_ext/v2}'s \code{pair} object carries the complex with its degree fields, \code{les\_maps} per reported spot, and the two flags \code{phi\_flat\_identity} / \code{residual\_flat\_zero}, which v1.6.2 reserved and did not emit (slot-5 WP-1). \textbf{Since v1.7.3 \ref{obl:coneexact}(vii) is implemented}: \code{HeckeConeStrictPairRoute} (slot-6 WP-4: \code{TauVertexPolicy::BoundaryStrict} and the \code{HeckeTauTable} constructor taking the $\taub$ table, \ref{prop:strictpair}, fills over $\Z$ only; \code{solve\_integer\_system\_wide} in \file{Hecke/HeckeIntegerSolve}), \code{HeckeConeRelativeRoute} (WP-5: \eqref{eq:route3}, the relative sub-complex constructor of \code{EquivariantCohomologyOnK}), both in \file{Hecke/HeckeConeAction} and gated by \code{HeckeConeActionValidation} arms A14 / A15, and the PD gate in \file{Hecke/HeckePairing} (WP-3, \code{HeckeDualityValidation}); the relabel of $\Acone{q}$ to tier~B in \code{golden\_objects()} (WP-5) moves no golden value. (At v1.7.2 not implemented: \ref{obl:coneexact}(vii), the cross-routes of slot 6.) \textbf{Since v1.9.0 the boundary and the cone over $\Z/\ell^{k}$} (hecke-extensions slot-9 WP-4, phase E8 item~4b, \Espec{sec:ring}~(f); at $\ell\ge5$, and refused at $\ell\in\{2,3\}$): a \emph{dense module route} beside the tape route, selected by the ring alone --- \code{ModuleConeT} (\file{EquivariantMappingConeModule.cpp}; the cone's cochains as the two Howell bases concatenated, bulk first, asserted to be the block sum's Howell form, $\Dcone$ by an exact module solve, the eight maps of \eqref{eq:coneles} in cyclic-summand coordinates) with \code{ConeModuleCertificate} (the E5 item numbering kept where a ring reading exists, and \textbf{item~14}, \ref{obl:coneexact}(iv)'s non-field half); \code{ModuleConeActionT} in \file{Hecke/HeckeConeAction} (\textbf{route~1 only}: $\hder$ in both bracketings, $\Phig$, $\Acone{q}$ on the cyclic summands, \ref{obl:coneexact}(i)--(iii) and (vi)'s commutator); \code{HeckeBoundaryOperatorT} in \file{Hecke/HeckeBoundaryOperator} ($\Thb_g$ over the ring in both bracketings, $T_p^{\partial}$ on the cyclic summands); \file{Hecke/SL2ZBoundaryPair} and \code{EquivariantCochainComplexMorphismOnK}'s constant-fibre constructor over the ring in force; \file{Hecke/HeckeReduction}'s reductions and Bocksteins on the boundary and the cone. The route is templated on the ring context's element type and runs in every word regime (the \code{U64\_Montgomery} and \code{U64\_Wide} cells exercise it end to end); the four \emph{classes} \code{EquivariantMappingCone}, \code{SL2ZBoundaryPair}, \code{HeckeConeAction} and the morphism class stay \code{uint32\_t}-instantiated, so they honour $\ell^{k}<2^{32}$ and refuse above it by the ring context's named refusal. The strict-pair route, route~3 and the relative sub-complex constructor refused the ring through v1.9.1 (\ref{obl:coneexact}(vii)'s composite half open). The $\Fl$ cone is byte-identical (\file{docs/matrix-engine/tools/cone\_output\_diff.sh} against the slot-8 close binary). \textbf{Since v1.10.0 the two cross-routes over $\Z/\ell^{k}$} (hecke-extensions slot-10 WP-3, merged at \texttt{c9684e9}; \Espec{sec:ring}~(b1), (f)): \code{HeckeConeStrictPairRoute} and \code{HeckeConeRelativeRoute} honour a composite normative action at $\ell\ge5$, $\ell^{k}<2^{32}$ (one guard, \code{CrossRouteRingGuard}; above $2^{32}$ the ring context's named \code{uint32\_t} refusal) --- the strict route's $\Th^{s}_g$ from \code{HeckeOperatorOnK}'s module route on the \code{BoundaryStrict} table and its $\Phig$ from \code{ModuleConeActionT} with $\etapair^{s}=0$; route~3 on \code{EquivariantCohomologyOnK}'s relative sub-complex constructor over the ring (the bulk module route with the vanishing relations of $K_{\partial}$; unit pivots asserted; refused at $\ell\in\{2,3\}$), transported by the explicit mutually inverse $M$, $N$ of \ref{obl:coneexact}(vii)'s as-built sentence --- each $\Acone{q}$ compared entrywise with route~1's; \code{HeckeConeCrossRouteRingValidation} (\code{standard}, $386$ checks: the seven points, a negative control on the strict table, the regime arm at $(10,134217757)$, the flipped rows both ways); the cone-base constructor, \code{HeckeExt/HeckeBoundary}, the pairing and the contragredient keep refusing; the $\Fl$ cone byte-identical (\file{cone\_output\_diff.sh}, $0$ hunks). \code{HeckeReductionValidation} (arms (4b-B), (4b-C), (4b-H), (4b-R), (4b-K), (X4b)), \code{HeckeReductionCoverageSlowValidation}; the $18$ objects per point of \Espec{sec:ring}~(f) are golden (the add-only event R1, $126$ entries). \textbf{At v1.23.0-pre.11} (slot-24 G6, value-preserving; \file{docs/implementation/slot24/G6\_report.md}, \file{G6\_audit.md}, MERGE after fixes F1--F3): \code{ConeCertificate}'s span inclusions C2 / C3 / C5 decide $Y\subseteq\operatorname{span}B$, on operands of at least \code{kSpanTapeMinArea} ($2^{20}$) entries of $[B\mid Y]$ (the tape route's threshold; below it the dense route of before), by the \emph{unit-row route} when every column of $B$ has a unit row $u$ ($B[u,:]=I$, so $y=Bx$ forces $x=y[u]$): $Y[r,:]=B[r,:]\,Y[u,:]$ on the other rows, by one exact product and a comparison, otherwise by the elimination route of before; a degree whose differential has no rows shares $Z$ with $\Ceq$; and in the census campaign the cone's cochain objects are released once its record objects are attached. With the field route's other memory levers (\code{HeckeOperatorOnK}'s program in parts, the release of \code{HeckeOperatorFamily}'s level~(a) after Stage~1, the morphism's $H$-basis read from the block form, the return of freed heap to the system by glibc's \code{malloc\_trim}), which realise no statement of this document, these have \textbf{one rollback lever, value-neutral}: \code{HECKE\_FIELD\_DIET=0} (\file{Util/FieldDiet.h}; the \code{EFFECTIVE HECKE\_FIELD\_DIET} line once when set); \code{tape::factor\_basis}'s two memory-only re-blockings (its stage-1 $R$ released early, stage~2's product in $1024$-row slabs, the same exact entries) have no lever. Every certificate item still runs, $\Acone{q}$ and the records are byte-identical (gate: \code{HeckeThetaRssDietValidation}'s \code{G6\_FIELD\_DIET} arm, \code{HeckeConeValidation}) \\
\addlinespace
\ref{sec:six}, \ref{sec:orbifoldmodel} & \textbf{---}. The related objects in the tree are the boundary datum, operator and $H^{1}_{!}$ of the \ref{sec:boundarydatum} row, the $\etapair$-table and the cone of the row above and the restriction $\rres^{q}$ of \code{TestSL2Z\_To\_S1Cusp\_RestrictionMorphism} (\file{LocalSystemSL2Z.cpp}), lifted pipeline-side since v1.5.2 as \file{Hecke/SL2ZBoundaryPair}; no free model. The free model is phase E9 of \file{hecke\_extensions\_implementation\_plan.md} v1.6, behind a contract chapter (C-A5) fixing the concrete refinement. \textbf{Since v1.8.0} C-A5 is re-scoped into \S\ref{sec:integral}, on the unrefined tessellation (owner ruling 6c of 2026-09-24; Remark~\ref{rmk:cafive}), and phase E9 is absorbed into the integral line's phases (the plan's next revision records it); still no free model in the tree. \textbf{Since v1.8.1} (status) the coefficient sequences of Claim~\ref{cl:bockstein} are computed over the ring in force: \file{Hecke/HeckeReduction} (slot-8 WP-4, phase E8 item~4a) gives the reductions $\Z/\ell^{k+1}\to\Z/\ell^{k}$ on cohomology and the Bockstein of $0\to\Fl\to\Z/\ell^{2}\to\Fl\to0$ as map equalities, with the identity $\mathrm{rank}\,\beta^{0}+h=\dim H^{0}(\Fl)$ asserted on the stated cells, $h$ the recorded \code{higher\_torsion\_count}, equal to the number of summands of exponent $\ge2$ wherever the certified exponent rule applies --- Remark~\ref{rmk:bocksteinexp}(1)'s count, evaluated (\code{HeckeReductionValidation}, \code{HeckeReductionCoverageSlowValidation}); still on the model complexes, not on a free model \\
\addlinespace
\ref{sec:reduction}, \ref{sec:fpfilter} & \file{Hecke/ExactInt} (\code{Mat22}, \code{checked\_*}, \code{HeckeOverflow}; \ref{lem:bounds}, \ref{thm:backend}, \ref{obl:backend}), \file{Hecke/QuadPoint} (\ref{alg:tau} step~1), \file{Hecke/Predicates} (\ref{thm:fp64}, \ref{cor:fptotal}, \ref{cor:fp32}, \ref{obl:fp}) \\
\addlinespace
\ref{sec:obligations} & the constructors and \code{compute()} bodies named in that chapter's code map, plus the $71$ \code{Hecke*} entries of the registry, \file{ValidationRegistry.cpp} and the slot-9 fragments \file{Registry/Slot9/} ($125$ registered, v1.9.1; $65$ at v1.8.1, $61$ at v1.7.4 and v1.8.0, $60$ at v1.7.3, $55$ at v1.7.2, $49$ at v1.6.2, $45$ at v1.5.2, $39$ at v1.4.2, $36$ at v1.3.3, $29$ at v1.3.2) \\
\addlinespace
\ref{sec:example} & \file{Hecke/HeckeCosetTablesValidation} (\code{O16}), \file{Hecke/HeckeOracleValidation} (\code{O1}--\code{O8}, \code{O14}), \file{Hecke/HeckeRelationsValidation} (\code{O9}, \code{O10}), \file{Hecke/HeckeEndToEndValidation} (the aggregate gate and \code{IsDicksonExceptional}) \\
\addlinespace
\ref{sec:integral} & \textbf{Partly built since v1.9.0} (partly tested, not built, v1.8.0--v1.8.1). \file{Integral/HeckeBredonTreeValidation} (slot-8 WP-8) tests \ref{prop:bredon}(2) at $\cring=\Fl$, $\ell\in\{2,3\}$; \file{Integral/IntegralCharpolyCrt} with \code{HeckeIntegralCharpolyValidation} (slot-8 WP-7) uses \ref{prop:crtcharpoly} (Remark~\ref{rmk:integralevidence}). $\Iso_{\bullet}$, $\Isob_{\bullet}$, $\cmpI$, $\tau_{\Iso}$ and the descent assembly did not exist at v1.8.1 (the first three exist since v1.9.0, below); the model complex over $\Z/\ell^{e}$, $\ell\ge5$, is phase E8's (\Espec{sec:ring}), and \textbf{is built since v1.8.1} (status): the third complex of Definition~\ref{def:ringinforce} at $\cring=\Z/\ell^{k}$ by \code{EquivariantCohomologyOnK} and \code{HeckeOperatorOnK} over a composite context (slot-8 WP-1, WP-3; \file{src/ModArithmetic/ENGINE\_API.md} \S13), refused at $\ell\in\{2,3\}$ as Proposition~\ref{prop:certifiede}'s scope clause says it must be. \file{Integral/} holds the two validations above and \file{Integral/IntegralCharpolyCrt}. \textbf{Since v1.9.0: partly built} (hecke-extensions slot~9). (a)~\textbf{At $\ell\in\{2,3\}$ the third complex over $\Z/\ell^{k}$} (slot-9 WP-3): \code{EquivariantCohomologyOnK} honours the ring there under the scoped policy \code{SmallEllPolicy::Bredon} of \file{Hecke/CoefficientRing}, which it alone passes, by a \emph{flat-coordinate module route} (\code{compute\_cohomology\_flat\_module}: $Z^{q}$ the kernel module of the stacked equivariance constraints and $d^{q}$, $B^{q}$ the image of the Howell generators of the equivariant cochains, $H^{q}$ by invariant factors; three re-derived certificates), since the cochain groups are not free there (the invariants computation of \ref{prop:ringinforce}(2)'s proof read over $\Z/2^{e}$, \Espec{sec:ring}~(c1$'$); v1.9.1); it reports lengths and invariant factors only --- the Bredon object of \ref{def:ringsemantics}(2) --- and its $K$-mode, \code{HeckeOperatorOnK} and \file{Hecke/HeckeReduction} refuse there by name (\ref{prop:certifiede}(5)). \code{HeckeBredonTreeRingValidation} compares it with the tree complex \eqref{eq:treemodel} over the same ring (\ref{prop:bredon}(2)), frozen from the design brief's probe (\file{docs/implementation/integral\_line\_probe.py} \code{--model-R}) on $96$ cells. (b)~\textbf{$\Iso_{\bullet}$, $\Isob_{\bullet}$ and $\cmpI$} (slot-9 WP-10, line~2's phase L2-1a): \file{Integral/Model/IsotropyModel} builds \ref{def:isomodel} on the unrefined tessellation to degree $q_{\max}+1$ by \ref{prop:isomodel}(1)'s procedure (\ref{IM1}, \ref{IM2} and $\cmpI$ a chain map certified exactly over $\Z[\Ga]$), with ranks $5,14,19,\ldots$ and $3,6,6,\ldots$ as \ref{prop:isomodel}(2) states, and $\CI(\cring)$ with $\cmpI^{\ast}$ over any word-sized ring; \code{HeckeIsotropyModelValidation}. (c)~\textbf{The integral groups at small $n$} (slot-9 WP-11, line~2's milestone M0): \file{Integral/Model/IntegralGroups} assembles $H^{q}(\Ga;\Sym^{n}\Z^{2})$, $n\le20$, $q\le5$, from ranks at a good prime, the $2$- and $3$-parts of $\sigma^{\le6}\CI(\Z)\otimes\Z/\ell^{K}$ with $K$ certified per differential (\ref{prop:certifiede}, Remark~\ref{rmk:certifiedeperdiff}) and the Dickson-prime parts from the shipped $K$-mode; \code{HeckeFreeModelValidation} finds them equal to the design brief's Appendix~A. $\tau_{\Iso}$ (\ref{prop:strictlift}) and the descent assembly (\ref{prop:descent}) were not built at v1.9.1. \textbf{Since v1.10.0} (hecke-extensions slot~10, line~2's phases L2-1b and M1(a)): (d)~\textbf{$\tau_{\Iso}$ and $\Th_{\Iso,g}$} (slot-10 WP-10, merged at \texttt{cb1c3c4}): \file{Integral/Model/StrictLift} builds \ref{prop:strictlift}(1) for $g=\mathrm{diag}(1,p)$, $p\in\{2,3,5\}$, in degrees $\le q_{\max}+1$ and certifies it exactly over $\Z[\Ga'']$; its coset transport over any word-sized ring; \code{HeckeStrictLiftValidation} (Remark~\ref{rmk:strictliftbuilt}). (e)~\textbf{The class $e$ and $E_q$} (slot-10 WP-13, merged at \texttt{3e371a2}): \file{Integral/Model/PeriodicityClass} builds \ref{lem:bulkperiodic}(1)'s $e$ as a $2$-cocycle on $\Iso_{2}$ (the Bockstein of $\zeta_+$), the Yoneda chain map $E\colon\Iso_{\bullet+2}\to\Iso_{\bullet}$ with $\aug\circ E_0=e$ certified exactly over $\Z[\Ga]$, its matrix on $\CI(\cring)$, and the Mayer--Vietoris connecting cochains through the Bass--Serre tree; \code{HeckePeriodicityClassValidation}. (f)~\textbf{The integral Hecke action, basis-free} (slot-10 WP-11, merged at \texttt{8f82016}): \file{Integral/Model/IntegralHecke} reads cochain operators on $H^{q}(\CI(\cring))$ over any word-sized ring --- operator identities as inclusions $P(Z^{q})\subseteq B^{q}$, induced operators on cyclic decompositions, the universal-coefficient torsion readings, a division-free characteristic polynomial --- with no integral basis; \code{HeckeIntegralActionValidation}. None of (d)--(f) is a member of \Espec{sec:ring}~(b)'s closure: each takes its ring as an explicit argument and reads no setting. The descent assembly (\ref{prop:descent}) and the value $\Z$ of the ring are still not built. \textbf{Since v1.11.0} (slot-11 WP-11): (g)~\file{Integral/Model/StrictLift} also builds the lifts for $\mathrm{diag}(1,p^{2})$, $\mathrm{diag}(1,p^{3})$ and $pI$, and \code{HeckeStrictLiftRelationsValidation} checks \ref{prop:relationsmodel}(2) with them (Remark~\ref{rmk:integralevidence}) \\
\addlinespace
\ref{sec:dicksontheta}, \ref{sec:ordinary} (Part~VII) & \textbf{since v1.11.0}: \code{FibreRepresentation::QuotientL} (\file{LocalSystemUtils}, \file{LocalSystemSL2Z}, \file{Hecke/HeckeOperatorOnK}), \file{HeckeExt/HeckeDicksonTheta} and \file{HeckeExt/HeckeDicksonThetaHecke} (slot-11 WP-5, WP-6), with the validations of Remark~\ref{rmk:thetacheckable}'s last paragraph; the ordinary idempotent of \ref{def:ordinary} at $p=\ell$ in \file{Hecke/LadicCell} (slot-11 WP-18c, the digit route of \Espec{sec:ring}; \code{HeckeLadicCellValidation}). \textbf{Recorded at v1.12.1} (slot-12 S12-FIN): \code{DicksonRoute} in \file{HeckeExt/HeckeFinalityRoutes}, the finality certificate (d0) ``dickson'' of \Espec{cert:dickson}, reads Proposition~\ref{prop:heckeH0}(1)'s eigenvalues $(1+p)p^{a}$ on an $H^{0}$ block. Part~VII's code map has the detail \\
\addlinespace
\ref{sec:hida} (Part~VII) & \textbf{built since v1.12.1's record} (slot-12 WP-4, WP-5; planned at v1.12.0-pre.1, contract first): \file{Hecke/HeckeOrdinaryPart} (Proposition~\ref{prop:torsionprepass}'s pre-pass, Lemma~\ref{lem:ordexponent}; \code{HeckeOrdinaryPrecisionValidation}) and \file{HeckeExt/HeckeHidaTower} (Definition~\ref{def:hidatower}, \ref{Hi2}--\ref{Hi6}; \code{HeckeHidaTowerValidation}, \code{slow} \code{HeckeHidaTowerSlowValidation}); the chapter's code map has the detail \\
\addlinespace
\ref{sec:hidafull} (Part~VII) & \textbf{built at v1.13.0} (contract first, slot-13 WP-2; refereed and repaired, WP-3; not built through v1.13.0-pre.2): \file{Hecke/HeckeLevelEll} (slot-13 WP-6: Proposition~\ref{prop:levelcorrespondence}'s block formula and Proposition~\ref{prop:levelellU}'s on the induced local system at $N=\ell$, the degeneracy and fibre maps, Proposition~\ref{prop:oldspace}, Proposition~\ref{prop:stabilisation}, \ref{Hi10}; \code{HeckeHidaLevelEllValidation}, against an independent group-cocycle route); \file{HeckeExt/HeckeHidaTower} part~2 and the $\rho^{\mathrm H}_n$ fibre (WP-5: Definition~\ref{def:hidalattice} and its variant (H), \ref{Hi7}--\ref{Hi9}; \code{HeckeHidaControlValidation}); the heavy cells (WP-10: \ref{Hi11} held at $(700,691)$, \code{slow} interlocked \code{HeckeHidaHeavyCellProbe}; \ref{Hi12} not run, deferred --- run once on the first cluster after the release and held, v1.14.0's record). The monomial and orbit-wise routes of \S\ref{sec:levelcost} and the Heilbronn-matrix action of \ref{levelO2} are an interface for general $N$, not scheduled (v1.14.0-pre.1: specified in \S\ref{sec:levelN}); built at v1.14.2 except the Heilbronn action (built at v1.15.0-pre.2, slot-15 WP-7). \file{HeckeExt/HeckeHidaComponents} (slot-15 WP-6, \code{HeckeHidaComponentsValidation}): $B_k/k\bmod\ell^{K}$ bignum-free and the readings of [U]'s (U1)--(U3) at tame level~$1$ on the $\Z/\ell^{K}$ module route --- built, a consumer of Theorems~\ref{thm:hidacontrol}, \ref{thm:hidaweight} and of the tower datum, not a new statement of this document. \textbf{Planned at v1.19.0-pre.1} (slot-19 WP-HL, upper core, gated on [U] part~3): the truncated $\Lambda$-adic arm of \file{HeckeExt/HeckeHidaTower}, \code{HeckeHidaLambdaValidation}. \textbf{At v1.19.0-pre.3}: not built; its research prototype (slot-19 R19-2, on $q$-expansions) ran, and [U] 0.3.0-pre.3 carries its readings and its recommendation of a $q$-expansion input route. \textbf{At v1.19.0-pre.4}: built (slot-19 WP-HL) --- \file{HeckeExt/HeckeHidaLambdaForms} (the $q$-expansion input arm, primary by the owner's ruling of 2026-10-03), \file{HeckeExt/HeckeHidaLambda} (the $\Lambda$ stage, record \code{hecke\_lambda/v1}), \code{HeckeLambdaValidation} (\code{standard}; the planned name \code{HeckeHidaLambdaValidation} is superseded); this Part's tower route (\file{HeckeExt/HeckeHidaTower}) enters as a cross-check at $\ell=37$, $103$ only. It realises no statement of this document.\\
\addlinespace
\ref{sec:levelN} (Part~VIII) & \textbf{built at v1.14.2} (slot-14 WP-12; contract first, WP-4; refereed and repaired, WP-5; released unbuilt at v1.14.0): \file{Hecke/HeckeLevelN} --- Definition~\ref{def:cosetrepsN} for every $N$ (\code{Gamma0LevelN}, \code{Gamma0SubTransversalN}), Definition~\ref{def:levelNmodel} over any word-sized ring (\code{InducedRelationsLevel}), Proposition~\ref{prop:levelorbit}(1),(2) orbit-wise; \code{HeckeLevelNValidation}: \ref{LN1}, \ref{LN2}, the $\tau$-route half of \ref{LN3}, \ref{LN4}--\ref{LN6} at the named cells. $H_{N,p}$ (Proposition~\ref{prop:heilbronn}) built at v1.15.0-pre.2 (slot-15 WP-7, \file{Hecke/HeckeLevelNHeilbronn}, \code{HeckeLevelNHeilbronnValidation}). The sparse cokernel elimination and the composite-$N$ degeneracy maps on the model built at
v1.15.0-pre.3 (slot-15 WP-8, \file{Hecke/HeckeLevelNSparse}), \ref{LN6} at scale (WP-9, \file{Hecke/HeckeLevelNR3}). $\alpha_d$,
$d>1$, on the cokernel presentation --- specified at v1.16.0-pre.1 (\S\ref{sec:levelNdegsec}: Propositions~\ref{prop:kappamodel}, \ref{prop:degenpres}, \ref{LN7}--\ref{LN9}) ---
\textbf{built at v1.16.0} (slot-16 WP-13, \file{Hecke/HeckeLevelNDegen}; \ref{LN7}, \ref{LN8} gated, \ref{LN9} reported, in \code{HeckeLevelNSparseValidation}). \textbf{At v1.17.0-pre.3} (slot-17 WP-5, merged and audited): the coset products behind \eqref{eq:MH} in \code{int64} first and \code{\_\_int128} on an overflow, bit-identical below the old reach; composite $N$ reached $90090$ at $k=2$, $60060$ at $k=4$. \textbf{Planned at v1.18.0-pre.1} (slot-18 WP-5; contract first,
rule~21; \textbf{not built}): the structure-aware fill model of Remark~\ref{rmk:levelNscale}(4) (a cost statement), a value-preserving elimination order gated
by bit-identity, the measurements at $N=99330$ ($k=2$) and $90090$ ($k=4$), and the wide-integer route for \code{Gamma1LevelN::contains\_wide} and
\code{HeilbronnChi}. \textbf{At v1.18.0-pre.2} (slot-18 WP-5, merged): the value-free generator order of the sparse elimination (finding W5-F1's root cause),
\code{Gamma1LevelN::contains\_wide}, \code{HeilbronnChi} on the wide route, built; the structure-aware fill model \textbf{failed} its factor-$1.5$ rule (a
finding; \code{hook\_rss/v1} stays); the two window measurements pending \\
\addlinespace
\ref{sec:nebentypus} (Part~VIII) & \textbf{built at v1.16.0} (slot-16 WP-10, WP-11; not built at v1.15.0, owner ruling Q3 of 2026-09-30; contract first, slot-15 WP-3, refereed by WP-4, released at v1.15.0): \file{Hecke/HeckeNebentypus} --- the character model and the $\Ga_1(N)$ model on \file{Hecke/HeckeLevelN}'s plug points, \eqref{eq:nebcorr} on the fixed model, over $\F_q$ and $W_K(\F_q)$ ring contexts; no new coset or $\tau$-table. \code{HeckeNebentypusValidation}: (NB1)--(NB6), \ref{NB4}'s sharp form at $n=0$ gated. \textbf{At v1.17.0-pre.3} (slot-17 WP-6, merged and audited): \file{Hecke/HeckeNebentypusSparse}, the sparse cokernel with a character over $\F_q$ (equal to the dense route entrywise where both run; the cap $D\le4096$ lifted), and $H_{N,p^{\nu},\chi}$ from $\mathrm{Mer}_{p^{\nu}}$, $\nu\ge2$ --- tested against Theorem~\ref{thm:nebrelations}(2),(3), not proved to be $T_{p^{\nu}}$; the model-route double coset of Proposition~\ref{prop:nebcorr} at $\mathrm{diag}(1,p^{\nu})$ and the sparse route over $W_K(\F_q)$ not built.
\textbf{Planned at v1.18.0-pre.1} (slot-18 WP-6; contract first, \textbf{not built}): cohomology over $W_K(\F_q)$ at $d(\chi)\ge2$, $K\le3$, dense then sparse
(Remark~\ref{rmk:nebcost}(4)); the model-route $[\Ga_0(N)\,\mathrm{diag}(1,p^{2})\,\Ga_0(N)]_\chi$ through the shipped table, a second route to
Theorem~\ref{thm:nebrelations}(2). \textbf{At v1.18.0-pre.2} (slot-18 WP-6, merged): both built (\file{Hecke/HeckeNebentypusGalois}; the model-route
$T_{p^{2}}$ equal to the recursion entrywise), and the $D=193\,536$ ladder cell on a second route. \textbf{Planned at v1.19.0-pre.1} (slot-19 WP-LV): the twisted Eisenstein congruence module at $N=\ell$, $\chi=\omega^{j}$, on this model (\Espec{cl:lvtwist}; \code{HeckeLValueValidation}). \textbf{At v1.19.0-pre.3}: not built --- the Eisenstein route refuses this model (a negative control of \code{HeckeLValueValidation}); the untwisted instance is built (the \S\ref{sec:saturation} row). \textbf{At v1.19.0-pre.4}: built (slot-19 WP-LV2, \file{HeckeExt/LValue/}) --- the character model $H^{1}(\Ga_0(\ell);V_{\Z/\ell^{K}}(\omega^{j}))$ on this section's model route, with $\iota^{\ast}$ on the induced model and $T_p$\,/\,$U_\ell$ by Proposition~\ref{prop:nebcorr}; \Espec{cl:lvtwist} computed at four cells ($\ell=37,59,67$) and two controls (\Espec{rmk:lvgates}); the model is accepted at level~$\ell$ only. It realises no new statement of this document.\\
\addlinespace
\ref{sec:saturation} (Part~IX) & \textbf{built as readings} (slot-14 WP-6, \file{HeckeExt/HeckeRingHeckeAlgebra}: $\mathsf A_K$, $\varepsilon^{\mathrm{sat}}_K$, the free part of a torsion block, the \code{stable} verdicts); the verdicts become \emph{certified by Corollary~\ref{cor:stablecert}} with slot-15 WP-5, on the WP-4 referee's ACCEPT of Propositions~\ref{prop:saturation}, \ref{prop:satdefect} (owner Q4); \file{HeckeExt/HeckeFormsHeckeAlgebra} (slot-12 WP-7) is the forms side of Proposition~\ref{prop:cohforms}; \file{HeckeExt/HeckeSaturationDefect} (slot-16 WP-5,
recorded at v1.17.0-pre.1) computes Remark~\ref{rmk:sattor}'s $I$ by two routes. \textbf{Planned at v1.19.0-pre.1} (slot-19 WP-LC): the lattice $H^{1}(\Ga;\Mn\otimes\Z/\ell^{M})_{\mathfrak m}$ as a $\Talg_{\mathfrak m}$-module, \code{HeckeLatticeModuleValidation} (\Espec{chk:latcmp}); (slot-19 WP-LV) the untwisted Eisenstein congruence module on a saturated block (\code{stable} at $K=1$; \Espec{cl:lveis}(3) at [E] 1.19.0-pre.2). \textbf{At v1.19.0-pre.3} (merged): both built --- \file{HeckeExt/HeckeLatticeModule}, \code{HeckeLatticeModuleValidation} (registered \code{standard}, \code{slow} by the owner's ruling of 2026-10-03; \Espec{chk:latcmp} held at $M\le4$, the free summand on the plus sign; the $\ell=7$ controls, \Espec{rmk:latcmpseven}), and \file{HeckeExt/LValue/}, \code{HeckeLValueValidation} (\code{standard}; \Espec{cl:lveis}(3) at four saturated rank-$1$ cells, \Espec{rmk:lvgates}). Neither realises a new statement of this document. \textbf{At v1.19.0-pre.4}: WP-LV2 extends \code{HeckeLValueValidation} by the free part of $H^{1}$ at finite precision (torsion-aware) and the deep rungs $(282,37)$, $(226,103)$ (an env arm; torsion-free there, saturated rank~$1$), and the two-sided bound on planted algebras; the twisted half is the \S\ref{sec:nebentypus} row. \textbf{At v1.23.0-pre.4} (slot 23, value-preserving): \S\ref{sec:satprecisionsec} built (G1, \code{ringalg::LevelCache}, rollback \code{HECKE\_PRECISION\_REUSE=0}); the module route re-implemented (G2-A--E, G2-G, rollback \code{HECKE\_MODULE\_STREAM=0}; the code map of \S\ref{sec:obligations}); the forms cross-route's q-series product by NTT and CRT (G2-X, rollback \code{HECKE\_FORMS\_FASTMUL=0}; the code map of \S\ref{sec:saturation}); each rollback lever selects a route, never a value.\\
\addlinespace
\ref{sec:isomodelell} (Part~IX) & \textbf{built at v1.16.0} at $\ell=3$ (slot-16 WP-16, stretch; not built at v1.16.0-pre.1, contract first, slot-16 WP-12): \file{Integral/Model/IsotropyModelEll} on \file{Integral/Model/IsotropyModel}'s plug point (a parameter $\ell$, columns $F'^{\sigma}$, cochains $M^{Q_\sigma}$), validated against $\Iso$ by \ref{IL1}--\ref{IL4} in \code{HeckeIsotropyModelEllValidation}. \textbf{At v1.17.0-pre.5} (slot-17 WP-7, merged and audited ACCEPT): the model descriptor as data (\code{ModelData}, \code{SubdivideEdgeOrbit}); \ref{IL3}'s level-$N$ clause run against the pipeline over $\F_3$, \ref{IL4} over $\Z/3^{e}$, $e\le3$, the pair and cone half of Proposition~\ref{prop:isomodelell}(5), and Remark~\ref{rmk:isomodellmp}(3)'s refinement half on a synthetic refined descriptor (cohomology only); \ref{IL3}, \ref{IL4} stay \COMP\textsuperscript{a} partial pending the next referee; not built: a strict lift on a refined descriptor. \textbf{At v1.23.0-pre.12} (slot-24 D0, contract first): Proposition~\ref{prop:certifiedeell} states the certified-exponent rule on $\Iso^{(\ell)}$ (one complex over $\Z[1/\mathsf q_\ell]$, the base changes the code builds, the Smith-valuation read-out); what is built and what is not (the \code{Ring2k} / \code{Ring3k} read-out; at $\ell=2$, $\Z/2^{W}$ with $W\le61$ reachable now through \code{ModCtx}, the reduction to $\Z/2^{W}$ for $W\in[62,64]$ not built; v1.23.0-pre.13) is that proposition's code map. \textbf{At v1.23.0-pre.14} (slot-24 D1, contract first): Propositions~\ref{prop:descentN}, \ref{prop:badprimesN} (\S\ref{sec:descentN}) state the descent at level~$N$ and its bad primes; built: the $\Zsix$-corner on $\Iso^{(2)}$, $\Iso^{(3)}$ for any coefficient module and the level-$N$ probe; not built: the direct Smith route on $\Ceq(M_{\Z/\ell^{K}})$, the sparse cokernel over $\Z/\ell^{K}$ (\code{SparseCokernel} is over $\F_p$), the level-$N$ assembly and $M_N(\chi)$ over $\mathsf O_\chi$ (that subsection's code map).
\textbf{At v1.18.0-pre.2} (slot-18 WP-7, merged; planned at pre.1): the strict lift on a synthetic refined descriptor (\file{Integral/Model/RefinedStrictLift}),
\ref{IL5} run (\COMP\textsuperscript{a}); \ref{IL3} over $\F_9$ and $\Z/9$; \ref{IL3}, \ref{IL4} moved to \COMP\textsuperscript{a} (the slot-18 referee F1's C-4) \\
\addlinespace
\ref{sec:descentN} (Part~IX) & \textbf{Contract first} at v1.23.0-pre.14 (Propositions~\ref{prop:descentN}, \ref{prop:badprimesN}; what is built and what not: the code-map paragraph after Proposition~\ref{prop:badprimesN}). D1's level-$N$ assembly is on branch \texttt{s24-d1impl} (\texttt{726c3a67}, audit fixes \texttt{87e8cf32}; audited PASS-with-fixes, \file{docs/implementation/slot24/D1\_code\_audit.md}), \textbf{not merged} at v1.23.0-pre.17. \textbf{Proposition~\ref{prop:descentNthree}} (v1.23.0-pre.17, stated \emph{after} the code, the audit's M1): the $3$-primary part at $\nu_3(N)=0$ is \code{RunSparseCokernelSmith}$(N,n,3)$ of \file{Integral/Descent/DescentLevelN} (gated by \code{ENGINE\_CHECK} on $\nu_3=0$), called by \code{AssembleDescentN} in place of $\Iso^{(3)}$, which keeps $\nu_3(N)>0$ (item~(5)); \code{IntegralDescentLevelNValidation} arm~(S) compares it with the dense $\Iso^{(3)}$ route on five cells; D1 merged at \texttt{1cc9a6d3} after v1.23.0-pre.17. \textbf{Proposition~\ref{prop:descentmaps}} (v1.23.0-pre.18, contract first, milestone D2): maps, the level cone and its long sequence over $\Z$; \textbf{not built} (its code-map paragraph); a check outside the build, \file{docs/implementation/slot24/cd3\_descentmaps\_check.py}. \textbf{At v1.23.0-pre.22} D2 is built, \file{Integral/Descent/DescentMaps} (branch \texttt{s24-d2impl}, audited ACCEPT-with-fixes, \file{docs/implementation/slot24/D2\_code\_audit.md}; item (2)(g) is the audit's F1), merged into \texttt{hecke-extensions} with it; what is built: its code-map paragraph \\
\addlinespace
\ref{sec:isomodeltree} (Part~IX) & \textbf{Contract first} at v1.23.0-pre.20 (Definition~\ref{def:isomodeltree}, Proposition~\ref{prop:isomodeltree}; plan milestone D5, brought forward at gate G-T): \textbf{built and merged} (\texttt{5eb7e9de}, v1.23.0-pre.30; its code-map paragraph: \file{Integral/Model/TreeCochains}, \file{Integral/Model/IsotropyModelTree}, \file{Integral/Descent/DescentTree}, \file{Integral/Descent/TreeComparisonArms}, \code{HeckeIsotropyModelTreeValidation} and its slow companion; the tree is the default route for the dense $\ell$-local bulk pieces, rollback \code{HECKE\_DESCENT\_TREE=0}; Hecke by transport: hook only at v1.23.0-pre.30, its real arm built with D3 and merged at \texttt{dc8261a3}, v1.23.0-pre.33; the tree cusp pair D5c built and merged at \texttt{4064e164}, v1.23.0-pre.36, \S\ref{sec:treepair}); the free tree resolution exists as D0's and D1's comparator; a check outside the build, \file{docs/implementation/slot24/cd6\_tree\_check.py}; (4) effective at v1.23.0-pre.21 (the referee's M1), checked with $\Iso$ by \file{cd6\_iso\_check.py} \\
\addlinespace
\ref{sec:descenthecke} (Part~IX) & \textbf{Contract first} at v1.23.0-pre.24 (Propositions~\ref{prop:strictliftN}, \ref{prop:crtcharpolyN}, \ref{prop:torsionhecke}, Theorem~\ref{thm:descentcert}; plan milestone D3, clauses C-D4, C-D5): \textbf{not built} (its code-map paragraph: \file{Integral/Model/StrictLift} parametrised in $p$, the level-$N$ operator on $\Iso$ and its transports, the output rule on \file{Integral/IntegralCharpolyCrt}, the torsion read, \code{IntegralDescentHeckeValidation}); a check outside the build, \file{docs/implementation/slot24/cd4\_hecke\_check.py}; refereed ACCEPT-with-fixes, fixes at v1.23.0-pre.25, with four code obligations stated in the code-map paragraph (the $3$-local route for every $N$, a certified $\tau_X$, $\tau_{\Iso}$ independent of $r_t$ and $N$, word conditions of the recovery); \textbf{at v1.23.0-pre.31} D3 is \textbf{built on branch \texttt{s24-d3impl}, pending merge and audit} (\file{Integral/Descent/DescentHecke*}, \code{IntegralDescentHeckeValidation} and its slow companion; the transport at $2$ and $3$; Theorem~\ref{thm:descentcert}(4) at $\ell\in\{2,3\}$ not built), and three clauses close D3's contract notes: the cokernel presentation at $\ell\ge5$ (Proposition~\ref{prop:torsionhecke}(3$'$)), obligation (b) through the lift, any admissible lift (Proposition~\ref{prop:strictliftN}(2)); \textbf{at v1.23.0-pre.33} D3 is \textbf{audited and merged} (\texttt{dc8261a3}; \file{docs/implementation/slot24/D3\_code\_audit.md}); \textbf{at v1.23.0-pre.37} the record of Theorem~\ref{thm:descentcert}(1)--(3), (6) with the pair columns is \code{integral\_descent\_record/v1}, item~7 of \S\ref{sec:recordschema} (D6: built on branch \texttt{s24-d6impl} at v1.23.0-pre.37; \textbf{at v1.23.0-pre.38} recorded as \textbf{audited, fixed and merged} --- \file{docs/implementation/slot24/D6\_code\_audit.md}, its fixes \texttt{19a969de}, merged at \texttt{2aa61397} --- with its golden event R1 \textbf{landed} at \texttt{6958a5d2}; the $\ker/\coker$ refusal's code name is \code{[KerCokerPrecision]} since \texttt{19a969de}) \\
\addlinespace
\ref{sec:descentpair} (Part~IX) & \textbf{Contract first} at v1.23.0-pre.26 (Definition~\ref{def:descentpair}, Lemma~\ref{lem:pairlifting}, Propositions~\ref{prop:descentpair}, \ref{prop:descentpairhecke}; plan milestone D4, clause C-D7): \textbf{not built} (its code-map paragraph: the pair on the code's $\Iso$ and $\Iso^{(\ell)}$, the level-$N$ Bredon pair, the certification of the cone, the closed-form boundary oracle, the Hecke operator on the cone from the pair lift, \code{IntegralDescentPairValidation}); a check outside the build, \file{docs/implementation/slot24/cd7\_pair\_check.py}; refereed ACCEPT-with-fixes, fixes at v1.23.0-pre.27, which amend obligations (c) ($\Dcone^{0}$ also by a prime $\ell'\ge5$, $\ell'>n+1$) and (f) (the level-$N$ cone operator at $p\notin\{2,3,5\}$ is refused, not a D4 obligation). \textbf{Built at v1.23.0-pre.28} on branch \texttt{s24-d4impl} (\texttt{d82c7761}), with the D4 code audit's and the pre.28 referee's code fixes \textbf{merged} at \texttt{3991a2f1} (v1.23.0-pre.29): \file{Integral/Descent/DescentPair.\{h,cpp\}}, \file{Integral/Descent/DescentPairHecke.cpp}, \code{IntegralDescentPairValidation}, \code{IntegralDescentPairSlowValidation} (\file{docs/implementation/slot24/DESCENT4\_report.md}); obligation (f) restated at pre.28 for the owner's scope addition of 2026-10-25 ($T_p$ on boundary and cone for every prime $p\nmid N$, the pair lift built on demand and certified at run time); pre.28 refereed ACCEPT-with-fixes (\file{docs/implementation/slot24/C-D7p\_referee.md}), fixes at v1.23.0-pre.29 (the refusal list of Proposition~\ref{prop:descentpairhecke}(5) and of (f) as the code's eleven \code{PairRefusal} codes, the proof of (1), $q_{\max}\le5$) \\
\addlinespace
\ref{sec:treepair} (Part~IX) & \textbf{Contract first} at v1.23.0-pre.33 (Definition~\ref{def:treepair}, Lemma~\ref{lem:treepaircyl}, Proposition~\ref{prop:treepair}, Remark~\ref{rmk:treepairscope}; plan milestone D5c, clause C-D8): \textbf{built, audited and merged} at \texttt{4064e164} (v1.23.0-pre.36; built at v1.23.0-pre.35 on branch \texttt{s24-d5cimpl}; the audit \file{docs/implementation/slot24/D5c\_code\_audit.md}; \file{Integral/Model/TreePair}, \file{Integral/Descent/DescentTreePair}, the default $2$- and $3$-adic pair route, rollbacks \code{HECKE\_DESCENT\_TREE\_PAIR=0}, \code{HECKE\_DESCENT\_TREE=0}; its code-map paragraph: the cusp model and the axis map in closed form on \file{Integral/Model/TreeCochains}, the tree cone as the route of the $2$- and $3$-adic pair pieces with a read-only rollback, the effective comparisons $\Phi^{\Z}_\partial$, $\Psi^{\Z}_\partial$, $\mathsf y_\ell$, $\mathsf y'_\ell$ on the code's $\Iso$, $T_p$ by transport with its off-diagonal $\eta_{\mathcal T}$, a validation against D4's pair and \file{d4\_pair\_oracle.tsv}); a check outside the build, \file{docs/implementation/slot24/cd8\_check.py}; refereed ACCEPT after fixes at v1.23.0-pre.33 (C-D8) and at v1.23.0-pre.35 (C-D8b), the fixes at v1.23.0-pre.34 and v1.23.0-pre.36, not re-refereed \\
\addlinespace
\ref{sec:mpmodel} (Part~IX) & \textbf{built at v1.22.0-pre.5} (slot-22 MP-1a, MP-1b, MP-T, MP-2; audited, F-MP2 ACCEPT WITH REPAIRS; contract text after the code, MP-CT): \file{Orbifold/MP/} --- descriptor and rule predicates, the MP set and its twisted normalised cochains, the comparison lift $\Phi^{\mathsf{MP}}$, $\varepsilon^{\mathsf{MP}}$, the refinement map $\mathsf R_{\Z}\to\mathsf M_7$ with the label factor, the all-degrees certificate; validations \ref{MPC1}--\ref{MPC9} (\code{OrbifoldMPRulesValidation}, \code{OrbifoldMPNerveValidation}, \code{OrbifoldMPTeardropValidation}, \code{OrbifoldMPComparisonValidation}, \code{OrbifoldMPRefinementValidation}, \code{OrbifoldMPAllDegreesValidation}, \code{OrbifoldMPNerveH3Slow}); since v1.22.0-pre.7 also the boundary, $\res$/$\cor$ and the pair of groups (slot-22 MP-3, audited F-MP3), \ref{MPC10}, \ref{MPC11} (\code{OrbifoldMPBoundaryValidation}, \code{OrbifoldMPLevelConeValidation}; widened by MP-3b, v1.22.0-pre.8); since v1.22.0-pre.8 also Hecke $T_p$ on $H^{\le2}$ (slot-22 MP-4, audited F-MP4), \ref{MPC12} (\code{OrbifoldMPHeckeValidation}); no golden object, setting or console verb. \textbf{Not built}: $H^{3}$ and $T_{p^{\nu}}$ on the MP model, level $N$ \\
\addlinespace
level (b), (c) output &\file{Hecke/FpPoly}, \file{Hecke/HeckeModuleDecomposition}, \file{Hecke/HeckeReport} --- the char-poly / primary-decomposition layer and the JSON writer; not the subject of any chapter above, and specified only by \S\ref{sec:recordschema}. Since v1.4.2 (slot-2 WP-1, phase E0b) \file{Hecke/HeckeModuleDecomposition} runs the locality certificate of [E, Cert.~5.2] inside its refinement loop (a candidate with two irreducible factors splits the block, an exhausted budget emits it \code{local = false} with a resumable state, never silently), and \file{Hecke/HeckeSettings} is the mathematical-settings layer --- the key \code{locality} (\code{budget=64} default, or \code{strict}), its source and fingerprint, emitted as \code{level\_c.locality} of \code{hecke\_report/v2}. \code{strict} is \textbf{implemented since v1.5.2} (slot-3 WP-3; until then it was accepted and threw a named not-implemented error): after the same bounded linear witness family budget mode runs, an undecided block is decided by the Frobenius radical and the Berlekamp subalgebra and certified by a residue witness over $\bar A$, terminating, never the default, byte-identical to budget mode wherever budget mode certifies (the whole golden grid re-run under it, \code{HeckeLocalityValidation}); \code{hecke::NotImplementedInThisBuild} survives as the named build-capability error that no settings value now raises. Since v1.5.2 \file{Hecke/HeckeSettings} also registers \code{u\_ell}, \code{sturm} and \code{sturm\_unavailable} (defaults unchanged, fingerprint $0$). The extension layer's \file{HeckeExt/HeckeLocalAlgebra} and \file{HeckeExt/HeckeGenerationCertificates} (slot-2 WP-4 / WP-5, phase E1a: the local algebra $A_S$ of a certified-local block, generation certificates (a) and (b), the socle bound, the resumable commutant) realise [E]'s local-algebra chapter and no statement of this document; \code{HeckeExtAlgebraValidation}; likewise, since v1.5.2 (slot-3 WP-2, phase E1b), \file{HeckeExt/HeckeSturmSweep} --- generation certificate (c) of [E] on real blocks, the sweep to $B'(k,\ell)$ with both \code{sturm\_unavailable} modes and the $U_\ell$ admission of Stage~1; \code{HeckeExtSturmSweepValidation}. Since v1.6.2 (slot-4 WP-1 / WP-2, phase E2) the extension record of \S\ref{sec:recordschema} item~6 is written by \file{HeckeExt/HeckeExtRecord.h} + \file{HeckeExt/HeckeExtReport} --- a \textbf{separate writer}, schema \code{hecke\_ext/v1}, never an extension of \file{Hecke/HeckeReport}, which slot 4 did not edit at all --- with \file{HeckeExt/HeckeExtFinal.h} the one tri-state \code{Final} and \file{HeckeExt/HeckeExtDriver} the Stage~1--5 orchestration (one cohomology, one operator family, a streamed operator for a Sturm prime outside $S_0$, the $U_\ell$ resume); \file{Hecke/HeckeSettings} gains \code{commutant\_budget} and \code{final\_mode}, six keys in all, defaults and fingerprint $0$ unchanged; \code{HeckeExtRecordValidation}, \code{HeckeExtRamanujanValidation}. Also since v1.6.2, and realising no statement of this document, the first E6 units \file{HeckeExt/HeckePredictions} (Bernoulli numbers mod $\ell$ and the Kummer / Harder divisibility predicate) and \file{HeckeExt/GloverFactors} (the composition factors of $\Sym^n(\Fl^2)$ over $\Fl[\GL_2(\Fl)]$); \code{HeckePredictionsValidation}. \textbf{Since v1.7.2 (slot 5) the extension record is \code{hecke\_ext/v2}}: the $u_\ell$ verdict is the \emph{cell's} with a single \code{resumable}, an ideal carries its \code{sign} and a \code{groups} list, and the \code{boundary}, \code{pair} and cone complexes are computed objects with their own degree fields and \code{les\_maps} (WP-1; \file{HeckeExt/HeckeExtComplexRecords} is the one translation layer), written by the same separate writer --- \file{Hecke/HeckeReport} is byte-unchanged by the whole extension stage. \code{verify()} gains code~22 (\S\ref{sec:recordschema} item~5's refusal, WP-11a) and code~23 (WP-12a). The record reaches the campaign per $(n_{\mathrm{sym}},\ell)$ cell (\file{Campaign/CampaignExt}, WP-2) and the golden reference, in the stage's one \textbf{add-only} emission (WP-3: $2617$ entries added, \emph{none changed, none removed}; $14\,049\to16\,666$). \file{Hecke/HeckeSettings} gains \code{field\_polynomial} and \code{split\_field\_report} (WP-7), \textbf{eight} keys in all, defaults and fingerprint $0$ unchanged. Three further units realise no statement of this document: \file{HeckeExt/HeckeCanonicalField} (the project-canonical split field $\Fl[Y]/(C_{\ell,r})$, pipeline-free, with no report consumer in this build), \file{HeckeExt/HeckeDicksonMultiplication} (multiplication by the Dickson invariant $L=x^{\ell}y-xy^{\ell}$, with the certificate that re-derives its entries by an independent convolution and checks the det-twisted intertwining $\rhon(A)L=\det(A)\,L\,\rho_m(A)$ entrywise) and \file{HeckeExt/HeckeM1Table}; \code{HeckeCanonicalFieldValidation}, \code{HeckeExtCampaignValidation}, \code{HeckeExtDeterminismValidation}, \code{HeckeExtKummerValidation}, \code{HeckeExtSmallEllValidation}, \code{HeckeExtSeparationValidation}. \textbf{Since v1.7.3 (slot 6)}, realising no statement of this document and moving no golden value (the record text is byte-unchanged, no settings key was registered): \file{HeckeExt/HeckePresentation} (WP-2, the on-request presentation of a certified-local block --- generators and a minimal relation set with a re-deriving \code{verify()}, reached only through a driver option, no record key); \file{HeckeExt/HeckeAtlas} (WP-8, the residue atlas $\Sigma(\ell)$ for $\ell\le13$, its separating set $S_{\mathrm{sep}}(\ell)$ and the byte-deterministic, write-gated store \code{hecke\_atlas/v1} --- the canonical split field's first consumer); \file{HeckeExt/HeckeIdentityCertificate} items~6 (WP-9, the residue-atlas identity under [E]'s standard-module convention, with the contragredient reading as its negative control) and~7 (WP-3, socle $\leftrightarrow$ Loewy, Claim~\ref{cl:pd}(3) tested); \code{ExtDriverOptions::s\_sep} (WP-9, opt-in); \code{HeckePresentationValidation}, \code{HeckeAtlasSkeletonValidation}, \code{HeckeAtlasValidation}, \code{HeckeExtEndToEndValidation}, and on the console side \code{ConsoleSettingsValidation} (C2a). \textbf{Since v1.7.4 (slot 7)}, realising no statement of this document: \file{Hecke/HeckeSettings} registers \code{presentation} and \code{predictions}, \textbf{ten} keys, fingerprint $0$ at the defaults, and a setter's \code{source} is an argument (the golden harness stamps \code{spec}); the extension record is \code{hecke\_ext/v3} --- the cell-local \code{proof\_range} over the anemic pool, \code{s\_sep}, the atlas \code{label}, the canonical \code{split}, \code{predictions\_mode} and the prediction objects, and the \code{presentation} object under \code{presentation = on}, all [E]'s field inventory (\S\ref{sec:recordschema} item~6(f)) --- written by the same separate writer; the golden texts carry the \code{settings} block reduced to its fingerprint, so a key registered at its default moves no golden line; \file{HeckeExt/HeckeCanonicalField}'s root finder evaluates the splitting polynomial through a precomputed Frobenius, value-preservingly, and routes the defining polynomial through the setting; \file{HeckeExt/HeckeAtlas} reaches $\ell=691$, with \file{HeckeExt/HeckeAtlasCli} (the build outside the suite) and \file{HeckeExt/HeckeTangentPair}; \file{Hecke/HeckeReport.\{h,cpp\}} is still byte-unchanged; the console's harness verbs (C2b), \code{ConsoleDeterminismValidation} \\
\bottomrule
\end{tabular}
\end{center}

\noindent
The arithmetic substrate below all of this --- \file{ModArithmetic/} and the engine
(\code{engine::kernels}, \code{engine::blk}, \code{engine::tape}, \code{engine::cuda})
--- realises no statement of this document. It is bound only by
Theorem~\ref{thm:backend} and by the invariant that $\Fl$-values are integers: the
engine's one floating-point step is a quotient \emph{estimate} corrected by an integer
compare-and-subtract, so it computes an exact residue, and the device kernels are
exact-integer twins of the CPU ones. Consequently no value in this document depends on
the thread count, the batching, the layout, the arena or the device
(\code{EngineDeterminismValidation}, \code{EngineCrossDeviceValidation},
\code{CudaBitIdentityValidation}).

\subsection{What is not implemented}\label{sec:notimplemented}

\begin{contractbox}
\textbf{v1.20.0-pre.1 (hecke-extensions slot-20 WP-2).} Slot~20's code is written against text outside this document, contract first (rule~21), and merges only after that
text is at pre.1 (rule~22). \textbf{WP-LV3} (the $L$-value route at depth~$2$, \Espec{sec:lvdepth2}): a generic $\mathfrak m_{\mathrm{Eis}}$-block reader over $\mathcal O=\Z_\ell[x]/(P)$, the
cross-kind gate (\Espec{cl:lvtriangle}, \Espec{chk:lvtriangle}), the change-of-lattice constant (\Espec{def:lvlattice}), $\inv^{\ast}$ over $\Z/\ell^{k}$ (the \S\ref{sec:involution} row; this
document's part is Definition~\ref{def:iota} and Propositions~\ref{prop:iotacommutes}, \ref{prop:iotasplit}), the twisted rung $(37;62,6)$ on \S\ref{sec:nebentypus}'s character model.
\textbf{WP-MEAS} (the measure-valued \textsf{linear} evaluation, \Espec{sec:lvmeasure}): the Mazur--Tate--Teitelbaum measure on this document's model through the cone of
\S\ref{sec:cone} (its $H^{1}_c$) and the modular symbols of \Espec{rmk:lvshapes}~(B); \code{HeckeLValueMeasureValidation} (\code{standard}). \textbf{NB-4} (the atlas sidecar and
\code{--serve-atlas}): a transport of the console and the atlas store, specified by [P]; it realises no statement of this document, and the console's values, transcripts and digests are
invariants of it. \textbf{WP-JET} (the build on an \code{aarch64} Jetson Orin, \code{sm\_87}): every patch is value-neutral and golden-clean, gated by \code{--all} and
\code{cross\_config.sh}. None of them realises a statement of this document; nothing of this box is built at this version. The milestone-M4 \code{\textbackslash LEAN} markers are not in
this pre-release: they follow slot-20 L20-6 (the Lean audit F3) and its repairs.
\textbf{v1.20.0-pre.2: built} (each merged after its contract text, rule~22; zero golden movement and the pinned digests unchanged at each merge). WP-MEAS (\code{f2897f74}): \code{ModularSymbolMeasure},
\code{LValueRoute::measure()}, \code{HeckeLValueMeasureValidation}; extended by slot-20 S1\,/\,S2 (\code{af3bf133}: the record \code{lfunction/v1} and the \code{read\_only} verb \code{lfunction},
\Espec{sec:lvlfunction}). NB-4 (\code{6218ebc6}): the atlas sidecar \code{oc-atlas-index/1}, \code{--serve-atlas}, \code{ConsoleServeAtlasValidation}. WP-LV3 (\code{c1fda051}): the generic
$\mathfrak m_{\mathrm{Eis}}$-block reader (tame, totally ramified orders; \Espec{chk:lvtriangle}), the cross-kind gate, $c^{\varepsilon}_f$ in the normalisation record (\Espec{prop:lvlattice}), $\inv^{\ast}$ over
$\Z/\ell^{k}$ (the \S\ref{sec:involution} row), the twisted rung $(37;62,6)$ on \S\ref{sec:nebentypus}'s character model. WP-JET: value-neutral, as stated. Still none realises a statement of this document.
\end{contractbox}

\begin{contractbox}
\textbf{v1.19.0-pre.5 (hecke-extensions slot-19 WP-2 continued).} \textbf{Built and merged} since the v1.19.0-pre.4 box: the \textsf{bilinear} and \textsf{linear} instances of
\Espec{def:lvroute} (slot-19 WP-BIL, \code{HidaPairingRoute}, \file{HeckeExt/LValue/HidaPairing}; slot-19 WP-LIN, \code{BirchModularSymbolRoute}, \file{HeckeExt/LValue/BirchModularSymbol};
arms of \code{HeckeLValueValidation}; \Espec{rmk:lvshapes}, \Espec{rmk:lvlegc}). They read this document's model, the cone of \S\ref{sec:cone} and the pairing of
Proposition~\ref{prop:pairing} (\file{Hecke/HeckePairing.h}); their gates observe, consistently with Definition~\ref{def:iota}, that this document's $T_p$ on $H^{1}_c$ is the
modular-symbol $T_p$ (no adjoint, no scalar) and that a geometric-even symbol lies in the minus part, and that the cone-level cup form is $-1$ times Hida's pairing on the
stored bases. None realises a statement of this document; no statement, tag or count of this document changes.
\end{contractbox}

\begin{contractbox}
\textbf{v1.19.0-pre.4 (hecke-extensions slot-19 WP-2 continued).} \textbf{Built and merged} since the v1.19.0-pre.3 box: WP-LV's twisted instance (slot-19 WP-LV2,
\file{HeckeExt/LValue/}, \code{HeckeLValueValidation}; \Espec{cl:lvtwist} on \S\ref{sec:nebentypus}'s character model) and WP-HL (the truncated $\Lambda$-adic arm,
\file{HeckeExt/HeckeHidaLambda}, \file{HeckeExt/HeckeHidaLambdaForms}, \code{HeckeLambdaValidation}; [U] \texttt{sec:u-lambda}). So every slot-19 package of the v1.19.0-pre.1 box is
built. None realises a statement of this document; no statement, tag or count of this document changes.
\end{contractbox}

\begin{contractbox}
\textbf{v1.19.0-pre.3 (hecke-extensions slot-19 WP-2 continued).} Of the v1.19.0-pre.1 box below, \textbf{built and merged}: WP-LC (\file{HeckeExt/HeckeLatticeModule},
\code{HeckeLatticeModuleValidation}, \Espec{chk:latcmp} held); WP-LV's untwisted instance (\file{HeckeExt/LValue/}, \code{HeckeLValueValidation}); WP-NF (the seventh
static target \file{NumberField/}, \code{NumberFieldValidation}); NB-2 (\code{--serve}, \code{ConsoleReadOnlyTagValidation}, \code{ConsoleServeValidation}; compiled only
with \code{-DOC\_BUILD\_NOTEBOOK=ON}). \textbf{Not built}: WP-HL (the truncated $\Lambda$-adic arm; its research prototype, slot-19 R19-2, ran and is read in [U]
0.3.0-pre.3) and WP-LV's twisted instance (\Espec{cl:lvtwist}, refused by the route). As the box below says, none of them realises a statement of this document;
no statement, tag or count of this document changes.
\end{contractbox}

\begin{contractbox}
\textbf{v1.19.0-pre.1 (hecke-extensions slot-19 WP-2).} Slot~19's code is written against text outside this document, contract first (rule~21), and merges only
after that text is at pre.1 (rule~22). \textbf{WP-LC} (the lattice comparison): $H^{1}(\Ga;\Mn\otimes\Z/\ell^{M})_{\mathfrak m}$ as a module over $\Talg_{\mathfrak m}$ on the
\S\ref{sec:saturation} route, gated by \Espec{chk:latcmp} at $(258,37)$, $M\le4$; it uses Proposition~\ref{prop:cohforms} and Corollary~\ref{cor:stablecert} and
proves nothing new here. \textbf{WP-LV} (the $L$-value route, \Espec{sec:lvalue}): the Eisenstein congruence module on the level-one model (its ordinary block, \S\ref{sec:hida}) and,
twisted, of \S\ref{sec:nebentypus}'s character model at $N=\ell$; this document's part is the model and Corollary~\ref{cor:stablecert}. \textbf{WP-HL} (the truncated
$\Lambda$-adic arm, [U] \texttt{sec:u-lambda}): an arm of the Hida tower code of \S\ref{sec:hida}--\S\ref{sec:hidafull}, upper core, built only on the referee F1's ACCEPT of
[U]'s part~3. \textbf{WP-NF} (the number-field core): a seventh static target \file{NumberField/}, STL only, guarded like \code{HeckeCore}; like the engine it
realises \emph{no} statement of this document. \textbf{NB-2} (\code{--serve}, \code{PrintJson}, the \code{read\_only} verb tags): a transport of the console, specified by
[P] v3.10 (slot-19 NB-1); it realises no statement of this document, and the console's values, transcripts and digests are invariants of it. Nothing of this
box is built at this version.
\end{contractbox}

\begin{contractbox}
\textbf{v1.14.0-pre.1 (hecke-extensions slot-14 WP-4).} \S\ref{sec:levelN} is written \emph{before} its implementation
(rule~21), its predictions \ref{LN1}--\ref{LN6} before the code that tests them (rule~20): the composite-$N$ table of
Definition~\ref{def:cosetrepsN}, the orbit-wise fixed spaces and the Heilbronn--Merel operator of Proposition~\ref{prop:heilbronn} are
slot-14 WP-12's, gated on the referee of the chapter (WP-5) and on the owner's [S14] Q4; the sparse elimination at large $N$ is
slot~15's. Nothing of \S\ref{sec:levelN} is built at this version.
\end{contractbox}

\begin{contractbox}
\textbf{v1.15.0-pre.1 (hecke-extensions slot-15 WP-3).} \S\ref{sec:nebentypus} (nebentypus: $\Ga_1(N)$ and $(\Ga_0(N),\chi)$ over $\F_q$ and
$W_K(\F_q)$) is written before its implementation (rule~21), its predictions (NB1)--(NB6) before the code (rule~20); the code is
slot~16's (owner ruling Q3). \S\ref{sec:saturation} turns the readings of slot-14's \file{HeckeExt/HeckeRingHeckeAlgebra} into
statements; its code change (the certified verdict) is slot-15 WP-5's, gated on the referee. \textbf{At v1.16.0} \S\ref{sec:nebentypus} is built
(slot-16 WP-10, WP-11; its code map).
\end{contractbox}

\begin{contractbox}
\textbf{v1.15.0-pre.3 (hecke-extensions slot-15 WP-17).} One sentence in \S\ref{sec:levelN}'s scope paragraph, proposed by slot-15
WP-8: a $\Ga$-module map between induced modules induces the map on the cokernel presentations of
Proposition~\ref{prop:levelorbit}(2), so the presentation-level fibre maps $\alpha^{0}$, $\omega$ of \file{Hecke/HeckeLevelNSparse} rest on
a statement of this document; the degeneracy maps $\alpha_d$, $d>1$, are asserted on the model only
(Proposition~\ref{prop:levelcorrespondence}(4)(c)). The code map of \S\ref{sec:levelN} records the sparse elimination (WP-8) and R3 at
scale (WP-9) as built. No statement added or removed, no label renamed, no tag moved.
\end{contractbox}

\begin{contractbox}
\textbf{v1.16.0-pre.1 (hecke-extensions slot-16 WP-12).} \S\ref{sec:levelNdegsec} (the model-to-presentation map, $\alpha_d$ on the presentation,
higher weight at composite $N$; \ref{LN7}--\ref{LN9}) and \S\ref{sec:isomodelell} (the $\ell$-local isotropy model and the closed form at $\ell=3$;
\ref{IL1}--\ref{IL4}) are written before their implementation (rule~21), their predictions before the code (rule~20). The first is slot-16
WP-13's, gated on the referee (WP-14); the second is not scheduled. v1.15.0-pre.3's sentence below is superseded by
Proposition~\ref{prop:kappamodel}(3) and Proposition~\ref{prop:degenpres}. \textbf{At v1.16.0} (the release) both are built: the first by slot-16 WP-13,
the second at $\ell=3$ by the stretch package WP-16, with \ref{IL3}'s level-$N$ clause and the pair and cone half of Proposition~\ref{prop:isomodelell}(5)
not run (the chapters' code maps).
\end{contractbox}

\begin{contractbox}
\textbf{v1.17.0-pre.1 (hecke-extensions slot-17 WP-2).} Contract first for slot~17 (rule~21), no statement of the three chapters changed: in
\S\ref{sec:levelN} the table's reach beyond composite $N\approx5.5\cdot10^{4}$ (WP-5); in \S\ref{sec:nebentypus} the sparse cokernel with a character and
$T^{(N,\chi)}_{p^{\nu}}$, $\nu\ge2$ (WP-6; the Galois ring at $d(\chi)\ge2$ a stretch); in \S\ref{sec:isomodelell} \ref{IL3} at level $N$, \ref{IL4} over $\Z/3^{e}$,
the pair and cone half of Proposition~\ref{prop:isomodelell}(5) and a synthetic refined descriptor (WP-7). None is built or run at this version; each
chapter's code map says what the package must keep. \textbf{At v1.17.0-pre.3} (slot-17 WP-3 referee F1's m-W3): WP-5 (\S\ref{sec:levelN}: the coset products in \code{int64} then \code{\_\_int128})
and WP-6 (\S\ref{sec:nebentypus}: the sparse cokernel with a character over $\F_q$; $T^{(N,\chi)}_{p^{\nu}}$ from $\mathrm{Mer}_{p^{\nu}}$, tested, not proved) are built,
merged and audited, as their chapters' code maps record; still \textbf{not built}: the sparse route over $W_K(\F_q)$ and the Galois ring at $d(\chi)\ge2$, the
model-route double coset $[\Ga_0(N)\,\mathrm{diag}(1,p^{\nu})\,\Ga_0(N)]_\chi$; WP-7's items as the code map of \S\ref{sec:isomodelell} says.
\textbf{At v1.17.0-pre.5}: WP-7 (\S\ref{sec:isomodelell}) is run, merged and audited (ACCEPT): \ref{IL3} at level $N$ over $\F_3$, \ref{IL4} over $\Z/3^{e}$, the
pair and cone, and the synthetic refined descriptor on cohomology; still \textbf{not built}: a strict lift on a refined descriptor and the Moerdijk--Pronk simplicial
set itself.
\end{contractbox}

\begin{contractbox}
\textbf{v1.18.0-pre.1 (hecke-extensions slot-18 WP-2).} Contract first for slot~18 (rule~21), no statement of the three chapters changed, one checkable item
added (\ref{IL5}, before the code, rule~20): in \S\ref{sec:levelN} the structure-aware fill model as a cost statement (Remark~\ref{rmk:levelNscale}(4)), a
value-preserving elimination order gated by bit-identity, the measurements at $N=99330$ ($k=2$) and $90090$ ($k=4$), and the wide-integer route for the
$\Ga_1(N)$ membership and the Heilbronn action with a character (WP-5); in \S\ref{sec:nebentypus} cohomology over $W_K(\F_q)$ at $d(\chi)\ge2$ and the model-route
$[\Ga_0(N)\,\mathrm{diag}(1,p^{2})\,\Ga_0(N)]_\chi$ (WP-6); in \S\ref{sec:isomodelell} a strict lift on a synthetic refined descriptor (\ref{IL5}) and \ref{IL3} over
$\F_9$, $\Z/9$ (WP-7). None is built or run at this version; each chapter's code map says what the package must keep.
\textbf{At v1.18.0-pre.2} (the WP-3 repairs): WP-5, WP-6 and WP-7 are built and merged, as their chapters' code maps record --- WP-5's generator order,
\code{contains\_wide} and \code{HeilbronnChi}'s wide route (its structure-aware fill model failed its rule, a finding; the two window measurements are
\textbf{not run}); WP-6's cohomology over $W_K(\F_q)$ at $d(\chi)\ge2$ and the model-route $T_{p^{2}}$; WP-7's strict lift on a synthetic refined descriptor
(\ref{IL5}) and \ref{IL3} over $\F_9$, $\Z/9$. Still \textbf{not built}: a Moerdijk--Pronk $S(\mathcal T)$ descriptor (Remark~\ref{rmk:mpcaveat}).
\end{contractbox}

\begin{contractbox}
\textbf{v1.13.0-pre.1 (hecke-extensions slot-13 WP-2).} \S\ref{sec:hidafull} is written \emph{before} its
implementation (rule~21), its predictions (Hi7)--(Hi12) before the code that tests them (rule~20): the
level-$\ell$ operator, the degeneracy maps and \ref{Hi10} are slot-13 WP-6's, the lattice datum and
\ref{Hi7}--\ref{Hi9} WP-5's, \ref{Hi11}, \ref{Hi12} WP-10's (owner-gated). Nothing of \S\ref{sec:hidafull} is
built at this version. v1.13.0-pre.2 (WP-3's repairs) adds, also unbuilt: the variant (H) of
Definition~\ref{def:hidalattice} on Hida's lattice (WP-5), and, not scheduled, the Hecke action on the cokernel presentation
(Remark~\ref{rmk:levelscaling}, \ref{levelO2}). \textbf{At v1.13.0} (the release) what slot~13 built is recorded in
\S\ref{sec:hidafull}'s code map; still not built: the monomial and orbit-wise routes, the Heilbronn-matrix action and the
composite-$N$ tables (\ref{levelO2}, \ref{levelO3}); \ref{Hi12} not run. \textbf{At v1.14.0} \ref{Hi12} has been run once on
the first cluster and held (its \COMP\textsuperscript{a} sentence); \textbf{at v1.14.1} a second cluster run
added \ref{Hi11} at $n=4150$ and \ref{Hi12} at $(4830,2411)$, both held; the level-$N$ interface is specified in \S\ref{sec:levelN}.
\end{contractbox}

\begin{contractbox}
\textbf{v1.12.0-pre.1 (hecke-extensions slot-12 WP-2).} \S\ref{sec:hida} is written \emph{before} its
implementation (the slot plan's rule~21, ``contract first''), its predictions (Hi1)--(Hi6) before the code
that tests them (rule~20): the torsion pre-pass and the projector exponent are slot-12 WP-4's, the towers
and $(10,2411)$ WP-5's. Nothing of \S\ref{sec:hida} is built at this version; what is built will be
recorded at v1.12.0.
\end{contractbox}

\begin{contractbox}
\textbf{v1.11.1 (the slot-11 close, status).} Unchanged in substance from v1.11.0: the WP-15 referee's
repairs are a provenance sentence, a new \PROVED{} remark (Remark~\ref{rmk:ordtorsionfree}, which narrows
\ref{open:ordprecision} to its per-cell bound; the item stays open) and wording. One code change is
recorded: \code{DicksonThetaCell} refuses $\ell\in\{2,3\}$ as the precondition of
Remark~\ref{rmk:thetasmallell} and Lemma~\ref{lem:h2torsion} (a \code{std::logic\_error}, exit~3), as
\Espec{sec:thetaobj}~(h) specifies, where through v1.11.0 it threw \code{NotImplementedInThisBuild}.
The counts are re-derived from \code{--list} at the close: $141$ registered validations ($55$
\code{fast}, $74$ \code{standard}, $12$ \code{slow}), $82$ \code{Hecke*}.
\end{contractbox}

\begin{contractbox}
\textbf{v1.11.0 (hecke-extensions slot-11 WP-15).} What slot~11 built, against the v1.11.0-pre.1 box
below. \textbf{\S\ref{sec:dicksontheta} is built and tested}: $\Qn$ as a third fibre representation, the
six terms and five maps of \eqref{eq:thetales} with $\delta$ by the snake construction, the Hecke
operators on all three terms (slot-11 WP-5, WP-6; the code map of Part~VII); every checkable item of
Remark~\ref{rmk:thetacheckable} held (its last paragraph). \textbf{\ref{prop:relationsmodel}'s lifts are
built} for $\mathrm{diag}(1,p^{2})$, $\mathrm{diag}(1,p^{3})$ and $pI$ and its recursion checked on the
free model (slot-11 WP-11; the note after its proof). \textbf{The ordinary projector} of
\ref{def:ordinary} is built for $U_\ell$ on the free block of $H^{1}$ over $\Z/\ell^{k}$, as a certified
Fitting idempotent, on the $\ell$-adic digit route (slot-11 WP-18c, \Espec{sec:ring}); on $\Fl$ and on
the word-sized composite route it is not built, and \ref{open:ordprecision} stays open (its status
note). \textbf{The digit route} (\file{ENGINE\_API.md} \S13.11, slot-11 WP-18a--c) is what
Obligation~\ref{obl:backend}(iv)'s scope sentence reads; nothing else of this document moves with it. Not
built: $\thL$ over $\Z/\ell^{e}$, $\thL$ at the console, the Hida chapter (slot~12).
\end{contractbox}

\begin{contractbox}
\textbf{v1.11.0-pre.1 (hecke-extensions slot-11 WP-2).} Part~VII is written \emph{before} its
implementation (the slot plan's rule~21, ``contract first''): \S\ref{sec:dicksontheta}'s $\Qn$, the six
terms and five maps of \eqref{eq:thetales} and the Hecke operators on them are slot-11 WP-5 and WP-6's;
\S\ref{sec:ordinary}'s projector has no implementation in slot~11. \ref{prop:relationsmodel} needs lifts
for $\mathrm{diag}(1,p^{2})$ and $pI$, which are slot-11 WP-11's (M-3). What was built is recorded at
v1.11.0.
\end{contractbox}

\begin{contractbox}
\textbf{v1.10.0 (hecke-extensions slot-10 WP-15).} What slot~10 built, against the v1.9.1 text
below (\S\ref{sec:codemap}). \textbf{Obligation~\ref{obl:coneexact}(vii) over $\Z/\ell^{e}$,
$\ell\ge5$}: its first two sentences --- the strict-pair route and route~3, each entrywise equal to
the normative $\Acone{q}$ --- are \textbf{discharged at $\ell^{e}<2^{32}$, every E8 point}
(slot-10 WP-3), the cross-routes refusing above $2^{32}$ by the \texttt{uint32\_t} regime rule;
its third sentence, the Poincar\'e-duality adjointness gate, is \textbf{recorded as unavailable over
the ring}, the pairing and the contragredient refusing a composite ring. So over the ring
Obligation~\ref{obl:coneexact} is partial in (vi) (the $\nu=2$ recursion not computed) and in
(vii)'s third sentence only. Phase E8 is complete: items 1--4b as specified, and (vii) over the
ring as just stated (\Espec{sec:ring}). At $\ell\in\{2,3\}$ the Hecke operator, the reductions, the
boundary and the cone over $\Z/\ell^{e}$ keep refusing, now for the reason of
Remark~\ref{rmk:bredonhecke}. \S\ref{sec:integral}: \textbf{$\tau_{\Iso}$ and $\Th_{\Iso,g}$ are
built} for $T_2$, $T_3$, $T_5$ (line~2's L2-1b, slot-10 WP-10; Remark~\ref{rmk:strictliftbuilt}),
\textbf{the class $e$ and $E_q$} (L2-1b, WP-13), and \textbf{the integral Hecke action read
basis-free} at $n\in\{10,12,16\}$ (M1(a), slot-10 WP-11) --- with $T_pT_{p'}=T_{p'}T_p$ tested on
$H^{\le5}$ and the recursion of Theorem~\ref{thm:relations}(1) \textbf{not} tested there, since it
needs a lifted datum for a double coset of determinant $p^{2}$, which is not built. Not built:
the descent assembly, the value $\Z$ of the ring, integer Hecke matrices in a basis (line~2's
M1(b)), and \S\ref{sec:orbifoldmodel}'s general free model. The other partial item is still
Obligation~\ref{obl:indep}(iii); Claim~\ref{cl:pd} is still tested and not proved. The
registered-validation counts of \S\ref{sec:codemap} are the slot close's.
\textbf{v1.9.0 (hecke-extensions slot-9 WP-15).} What slot~9 built, against the v1.8.1 text
below (\S\ref{sec:codemap}). \textbf{Phase E8 item~4b is built}: the boundary and the cone over
$\Z/\ell^{k}$, $\ell\ge5$ (\Espec{sec:ring}~(f), slot-9 WP-4), so \textbf{the non-field half of
Obligation~\ref{obl:coneexact}(iv) is checked} --- by solves at every spot of
\eqref{eq:coneles}, with the witnesses re-derived --- and its $18$ objects per point are
golden. Over the ring, Obligation~\ref{obl:coneexact} is \textbf{partial} in two clauses:
\textbf{(vi)}, whose commutator is asserted and whose $\nu=2$ recursion is not computed over the
ring; and \textbf{(vii)}, by owner ruling 3~(iii)~(A) of 2026-09-24 --- \textbf{route~1 only}: the
strict-pair route, route~3 and the relative sub-complex constructor refuse the ring, each with an
arm, and \textbf{(vii)'s composite half is open}. Phase E8 is therefore complete as specified:
items 1--4b as specified, 4b over $\Z/\ell^{k}$ by route~1, and (vii)'s composite half open. At
$\ell\in\{2,3\}$ the third complex over $\Z/\ell^{k}$ is built as the Bredon object it is
(Definition~\ref{def:ringsemantics}(2), lengths and invariant factors only; slot-9 WP-3), and
its $K$-mode, the Hecke operator, the reductions, the boundary and the cone keep refusing
there. \S\ref{sec:integral} is \textbf{partly built}: $\Iso_{\bullet}$, $\Isob_{\bullet}$ and
$\cmpI$ (line~2's phase L2-1a, slot-9 WP-10) and the integral groups for $n\le20$, $q\le5$
(line~2's milestone M0, slot-9 WP-11); $\tau_{\Iso}$ and the descent assembly are not built,
and \S\ref{sec:orbifoldmodel}'s general free model is still unbuilt. The other partial item is
still Obligation~\ref{obl:indep}(iii); Claim~\ref{cl:pd} is still tested and not proved.
The registered-validation counts of \S\ref{sec:codemap} are the slot close's.
\textbf{v1.9.1 (the slot-9 close, status).} Unchanged in substance from v1.9.0: the referee's
repairs are wording, and the counts are re-derived from \code{--list} ($125$ registered, $71$
\code{Hecke*}). Since hecke-extensions slot-9 WP-4t the point $(60,5)$ of item~4b's ring arms runs
in the \code{slow} tier (\code{HeckeReductionCoverageSlowValidation}); its golden objects stay in
every run of the suite.
\textbf{v1.8.1 (the slot-8 close).} Unchanged in substance from v1.8.0. Two chapters have no
implementation of their free models: \S\ref{sec:orbifoldmodel} and \S\ref{sec:integral}
($\Iso_{\bullet}$, $\cmpI$, $\tau_{\Iso}$, the descent). What slot~8 built is the complex
over the ring in force at $\Z/\ell^{k}$, $\ell\ge5$ (Definition~\ref{def:ringinforce}'s third
complex; phase E8 items~2, 3 and 4a), with its reductions and Bockstein, on the model
complexes; it is refused at $\ell\in\{2,3\}$. \textbf{The boundary and the cone over
$\Z/\ell^{k}$ are not built} (phase E8 item~4b, the next slot's), so the non-field half of
Obligation~\ref{obl:coneexact}(iv) stays unchecked. The only partial item is still
Obligation~\ref{obl:indep}(iii); Claim~\ref{cl:pd} is still tested and not proved.
\textbf{v1.8.0 (slot-8 WP-9).} The chapter added in this version,
\S\ref{sec:integral}, has no implementation beyond two validations that test statements of
it (\S\ref{sec:codemap}): the free model $\Iso_{\bullet}$, the comparison map $\cmpI$, the lifted
datum $\tau_{\Iso}$ and the descent assembly are specified and not built. The
\emph{prerequisite} the v1.7.4 sentence below calls deferred --- the concrete refinement ---
is re-scoped away: C-A5 is that chapter, on the unrefined tessellation
(Remark~\ref{rmk:cafive}). Everything else in the v1.7.4 text stands, and follows.
\textbf{v1.7.4 (the slot-7 close).} Unchanged from v1.7.3 in substance: one chapter has no
implementation (Chapter~\ref{sec:orbifoldmodel}; its contract prerequisite, the concrete
refinement, is deferred to slot~8), Chapter~\ref{sec:cone} is implemented in full, the only
partial item is Obligation~\ref{obl:indep}(iii), and Claim~\ref{cl:pd} is tested --- since
v1.7.4 on the whole case list at $\ell\ge5$ --- and not proved. The v1.7.3 text follows.
\textbf{v1.7.3 (the slot-6 close).} One chapter has no implementation: Chapter
\ref{sec:orbifoldmodel}. \textbf{Chapter~\ref{sec:cone} is implemented in full}: Obligation~\ref{obl:coneexact}(vii)'s
strict-pair route, route~3 and PD adjointness gate are evaluated since v1.7.3 (slot-6 WP-4,
WP-5, WP-3; \S\ref{sec:codemap}), the last being a \emph{test} of Claim~\ref{cl:pd}, which this
document still does not prove; and \S\ref{sec:contragredient} is implemented since v1.7.3
(slot-6 WP-1). The only partial item left is Obligation~\ref{obl:indep}(iii), below. The
rest of this box is the v1.7.2 text, kept as history. One chapter has no implementation and two have a partial one. Chapter
\ref{sec:orbifoldmodel} (the general orbifold model) is \textbf{unimplemented}: there is
no free $\Ga$-model anywhere in the tree. Chapter~\ref{sec:cone} (the pair homotopy and
the cone) is \textbf{implemented except one clause since v1.5.2} (partial since v1.4.2,
when only the certified $\etapair$-table existed): the equivariant cone
(\file{EquivariantMappingCone}), $\hder$, $\Phig$ and $\Acone{q}$
(\file{Hecke/HeckeConeAction}) exist, Obligation~\ref{obl:eta} is discharged in full and
every clause of Obligation~\ref{obl:coneexact} but (vii) is evaluated
(\code{HeckeConeValidation}, \code{HeckeConeActionValidation}; two of the $108$ registered
validations compute $\Hcpt{\bullet}$). What it still lacks is clause (vii)'s two
cross-routes (the strict pair and route~3, phase E5's slot-6 arms), PD adjointness (which
waits on Claim~\ref{cl:pd}). \textbf{The cone's integration into the record landed in
v1.7.2} (slot-5 WP-1, phase E2): \code{hecke\_ext/v2}'s \code{pair} object carries the
complex, \code{les\_maps} per reported spot and the two flags v1.6.2 reserved.
\textbf{Chapter~\ref{sec:boundarydatum} is no longer among them}: since v1.3.3 the
certified $\Gamma$-level $\taub$, the boundary coset table $B_N/\Gbdry$ and the operator
$\Thb_g$ exist (\file{Hecke/HeckeBoundaryDatum}, \file{Hecke/HeckeBoundaryOperator}), and
every clause of Obligation~\ref{obl:taubdry} is implemented and run, (i)--(iv) by
\code{HeckeBoundaryTauValidation} and (v)--(vii) by
\code{HeckeBoundaryOperatorValidation}; since v1.4.2 the parabolic cohomology
$H^{1}_{!}=\ker\rres^{1}$, its Hecke action and Harder's criterion are computed as well
(\file{HeckeExt/HeckeBoundary}, slot-2 WP-7). What the chapter's objects still lack is
their integration into the report and the campaign record (phase E2), and what the
\emph{pair} lacks is \S\ref{sec:cone}'s cone.
\textbf{Chapter~\ref{sec:involution} is not among them either}: since v1.1.1 every
clause of Obligation~\ref{obl:iota} is implemented and run, (i),(iv) by
\code{HeckeIotaTauSymmetryValidation} and (ii),(iii),(v) by
\code{HeckeIotaSignValidation}. What that chapter lacked up to v1.5.2 was a \emph{production}
$\inv^{\ast}$ --- the \code{EquivariantCochainComplexMorphismOnK} its code map
specifies; the two validations built their own. \textbf{It exists} (v1.6.2, the slot-4
close; built by hecke-extensions slot-4 WP-3 to exactly the specification of
\S\ref{sec:involution}): \file{Hecke/HeckeInvolution} is that morphism and
\file{HeckeExt/HeckeSignSplit} the sign split of clause (v), gated by
\code{HeckeInvolutionValidation} (\textbf{standard}, $2954$ checks), which compares it
\emph{entrywise} against \code{HeckeIotaSignValidation}'s independently built one on that
validation's $77$-cell grid and reproduces its reference line field for field. See
Remark~\ref{rmk:iotapairing}, \S\ref{sec:coneaction} and \S\ref{sec:codemap}'s row. It
has no open item: clause (v)'s sign
question was resolved at v1.2.0 by owner decision OQ3 (Definition~\ref{def:iota}(iv)
negated, the sign gating and holding). Obligation~\ref{obl:indep} is
\textbf{partial}: clause (iii) is not reachable through the present interface, and that
chapter says which. Their obligations are written as predicates on emitted data
precisely so that they can be handed to the phase that builds them unchanged; they are
a specification, and nothing in this section may be read as a claim that they hold.
\end{contractbox}

\subsection{The build and configuration surface}\label{sec:buildsurface}

\noindent
Six library/executable targets (\code{ModArithmetic}, \code{HeckeCore},
\code{MatrixEngine}, \code{EngineCuda}, \code{ConsoleCore} --- the console's language
core, STL-only, realising no statement of this document --- and
\code{cohomology\_app}); the reference toolchain
is recorded in \file{CLAUDE.md} and is not part of this contract. What \emph{is} part of
it:

\begin{itemize}
\item \textbf{The two rollbacks.} \code{cmake -DHECKE\_PURE\_INTEGER=ON} disables the
  certified floating-point filter of \S\ref{sec:fpfilter} and forces exact-integer
  predicates; \code{ENGINE\_BACKEND=legacy} (runtime, no rebuild) sends every retrofitted
  path back to the pre-engine loop. \textbf{Both must produce bit-identical stored
  objects}, which is Obligation~\ref{obl:fp}(iv) for the first and an engine invariant
  for the second; the first is a two-build \emph{protocol}, not a predicate, and
  \S\ref{sec:obligationroster} records its executed run.
\item \textbf{The two shadow builds.} \code{-DENGINE\_SHADOW\_LEGACY=ON} runs both arms
  of every retrofitted call and throws on the first differing entry;
  \code{-DENGINE\_VERIFY\_DEVICE=ON} re-executes every GPU node on the CPU inside a real
  run and compares every written window. Both are diagnostics
  (rule~\ref{A6}): a green shadow build is evidence, not a discharge.
\item \textbf{The execution switches} --- \code{ENGINE\_THREADS}, \code{ENGINE\_BATCH},
  \code{ENGINE\_DEVICE} (\code{cpu} is the default and the golden path),
  \code{ENGINE\_ARENA}, \code{ENGINE\_LAYOUT}, \code{ENGINE\_CSE}, \code{ENGINE\_FUSE},
  \code{ENGINE\_KSPLIT} and the CUDA and tracing families. None can change a value; that
  is asserted, not assumed (\code{EngineDeterminismValidation}: identical digests across
  thread counts, batching, arena, layout, CSE/fusion and device).
\item \textbf{The four entry points.} \code{./cohomology\_app <Name>} runs one
  registered validation (exit code $=$ its own); \code{--all} runs the fast and standard
  tiers and \code{--all-full} adds the slow tier --- \textbf{the Hecke acceptance gate
  needs \code{--all-full}}, because \code{HeckeBenchmarkSmoke} is slow-tier;
  \code{--bench <mode>} drives the shipped pipeline through the public API only; and
  \code{--campaign <spec>} enumerates an $(n,\ell,p,\nu,q)$ grid and writes one record
  per job. All four use the same code path, which is what makes a validation's verdict
  and a campaign's record statements about the same object.
\item \textbf{The certification level.} \code{--lean} and a spec's
  \code{certify = full\,|\,sampled(\ldots)\,|\,none} reduce the \emph{external}
  cross-checks only. The in-process certificates --- $\partial\tau=\tau\partial$ in the
  $\tau$-table constructor, the two-route entrywise comparison and the exact $\Th_{\rm
  eq}$ residual inside \code{compute()}, and the level-(c) certification --- run on every
  job in every mode, \code{certify = none} included
  (\code{CampaignCertificateValidation}). They are Obligations~\ref{obl:tau},
  \ref{obl:agreement}(i) and \ref{obl:welldef}(iii), and they are not sampleable.
\end{itemize}

\subsection{The record schema}\label{sec:recordschema}

\noindent
An emitted record is what makes the obligations of \S\ref{sec:obligations} checkable
\emph{after} the run, by someone who did not make it. The shipped writer is
\code{hecke::BuildHeckeReportJson} (\file{Hecke/HeckeReport.cpp}, schema string
\code{hecke\_report/v2} since v1.4.2 --- hecke-extensions phase E0b, slot-2 package WP-2,
the stage's one golden re-baseline; \code{v1} until then) wrapped, per job, by the
campaign driver. Binding content:

\begin{enumerate}
\item \textbf{The parameters}: $\ell$, $n$, $p$, $\nu$, the model identifier, and the
  reduction regime actually used --- in particular the $\tau$-table's recorded solve
  regime ($\Z$ or mod $\ell$; Obligation~\ref{obl:tau}(vii),
  Remark~\ref{rmk:tauintegral}) and the certified bound $B$ of Lemma~\ref{lem:bounds}.
\item \textbf{The conventions, in band}: polynomial coefficients low-to-high degree,
  matrices as arrays of rows, entries canonical residues in $[0,\ell)$. A record that
  does not say this is not self-describing, and the two orders are exactly the pair a
  consumer gets wrong.
\item \textbf{The three levels}: (a) $\Th^{(q)}_{\mathrm{flat}}$, optional and
  size-gated --- it is $(N_qd_V)^2$; (b) $\Amat{q}(T_p)$ for every computed $p$ and
  $q$; (c) the primary decomposition, with, per block, \emph{either} a proved
  \code{residue\_degree} \emph{or} a \code{residue\_degree\_lower\_bound} --- never
  both --- a machine-readable \code{residue\_degree\_certified} flag, and the certifying
  linear combination. This is rule~\ref{A6} in the schema: a lower bound may not be
  read as a degree.
\item \textbf{Which certificates ran}, and the level they ran at, with counts: the
  $\tau$-certification verdict, the two-route entrywise comparison, the $\Th_{\rm eq}$
  residual (and \code{"theta\_eq\_certified": false} when \code{--lean} turned it off),
  the decomposition certification, and the external cross-checks together with the
  sampling mode and stride. A record in which a certificate is absent and one in which
  it ran and passed must be distinguishable.
\item \textbf{The reported degrees} --- and here the schema carries
  \ref{open:ell2}, decided 2026-09-16. \textbf{Per complex} (v1.4.0; owner confirmation
  1 of the slot-2 plan, 2026-09-16; open point (9) of the v1.3.3 changelog): every
  complex a record is about --- the bulk $\Ceq(K,\Rel;L)$, the boundary
  $C^{\bullet}_{\partial}$, the cone --- carries its own four fields, adopted with the
  policy: \code{q\_max} (an integer, the top degree of \emph{that} complex),
  \code{truncated} (a boolean), \code{q\_max\_source}
  $\in\{\code{"derived"},\code{"set"}\}$ --- rider (i) of \ref{open:ell2}: derived on a
  finite model, and \code{q\_max} may be set only where \code{truncated = true} --- and
  \code{degrees}, the complex's own degree list $[0,\dots,q_{\max}]$, derived.
  \textbf{Never a maximum over the complexes of a record.} A record about one complex
  (the per-degree \code{hecke\_report} record, which is about the bulk) carries the four
  at record level; a record about several (the extension record of phase E2, about
  bulk, boundary and pair) carries them inside each complex's object. The
  \emph{reported-degree selection} of a campaign job --- rider (i)'s independent filter,
  the campaign's \code{point.degrees}, which may be a proper subset --- is a job-level
  field distinct from any complex's \code{degrees} and keeps its meaning. A record's
  identity --- its output path and its resume key --- includes each complex's tuple and
  the reported selection, so that two jobs differing only in them are distinct jobs (the
  campaign job identity carries the reported degree list since 2026-09-16,
  \code{Job::id()}, pinned by \code{CampaignSmokeValidation}). \textbf{On the model
  complexes every field is a constant}: bulk $2$ / \code{false} / \code{"derived"} /
  $[0,1,2]$, boundary $1$ / \code{false} / \code{"derived"} / $[0,1]$ (the boundary
  object derives its own $q_{\max}=1$; Proposition~\ref{prop:circlemodel}, code map of
  \S\ref{sec:boundarydatum}); on a truncated free model they are the bound of the box,
  and ``$H^{q}=0$ for $q>2$'' is not sayable. \textbf{Implemented at
  \code{hecke\_report/v2}} (v1.4.2; hecke-extensions phase E0b, slot-2 package WP-2,
  2026-09-16, commits \texttt{6ced10c} / \texttt{70c37de}; until v1.4.0 this instruction
  was addressed to E5, and v1.4.0--v1.4.1 recorded it as scheduled): the per-degree writer
  \file{Hecke/HeckeReport.\{h,cpp\}} moved from \code{hecke\_report/v1} to \code{v2} and
  carries the four fields at record level beside its per-degree \code{degree}, \emph{per
  complex} and derived --- from \code{HeckeReportInput::num\_degrees}, the degree count of
  the complex the record is about, which the writer now \emph{requires} (a zero count
  throws), so that no caller can obtain the fields as a maximum over anything; the
  \code{"set"} value of \code{q\_max\_source} is reserved for the truncated free model of
  phase E9 and is produced by no shipped caller; the campaign writer
  \file{Campaign/CampaignDriver.cpp}, which embeds one such report per degree and emits
  the job-level \code{"degrees":[\ldots]} selection, emits \code{q\_max},
  \code{truncated} and \code{q\_max\_source} beside it, un-hashed. The schema move was
  the stage's one golden re-baseline --- exactly the $384$ \code{report\_canonical}$^{q}$
  / \code{report\_full}$^{q}$ lines of \file{Golden/GoldenReference.inc} predicted, $0$
  elsewhere, every changed key path allow-listed
  (\file{docs/matrix-engine/rebaseline-2026-09-16-hecke-report-v2.md}); the extension
  record of E2, about several complexes, is item~6 below. This is site~(6) of
  \S\ref{sec:oc1sites}, and it is the only place in this chapter that depends on the
  policy: changing \ref{open:ell2} changes these fields and nothing else here.

  \smallskip
  \textbf{Maps, and complexes built from complexes} (v1.5.0; deferred item~5 of the
  v1.4.2 changelog entry --- \file{docs/implementation/slot2/WP-8\_report.md} finding~4
  and \file{docs/implementation/slot3/WP-5\_report.md} finding~2). The four fields are
  defined for a \emph{complex}; two objects a record is about are not complexes, or not
  primitive ones, and the rule for them is stated here once rather than per object.
  \begin{itemize}\setlength{\itemsep}{0pt}
  \item \textbf{A map} --- the restriction $\rres$ (degree $0$), the pair homotopy
    $\etapair$ (a chain map of degree $+1$ from the boundary chains to the bulk chains,
    Definition~\ref{def:pairhomotopy}), the derived $\hder$ (a cochain map of degree
    $-1$, \eqref{eq:derivedh}) --- carries in its record the four fields of its
    \textbf{source} complex, and no others. \textbf{The degree shift is a constant of
    the schema}, fixed once per map kind by the definition that introduces the map; it
    is not a field, and a record may not vary it. So the $\etapair$-record's fields are
    the boundary's --- on the shipped model $q_{\max}=1$, \code{truncated = false},
    \code{"derived"}, \code{degrees} $[0,1]$ --- with $\etapair(q)$ landing in bulk
    degree $q+1$ (implemented that way at \code{hecke\_pair\_homotopy/v1};
    \S\ref{sec:cone}'s code map). On a truncated free model (phase E9) a map is built,
    and certified, in every source degree in which the box of \ref{open:ell2} builds its
    source complex, i.e.\ $q\le q_{\max}(\text{source})+1$, and \emph{reported} in the
    source degrees $q\le q_{\max}(\text{source})$; its target degrees are the shift of
    those, and the box's bound on the target complex must cover them --- a condition on
    the pair of $q_{\max}$, not a further field.
  \item \textbf{A complex built from others} --- the cone $\Cone{\rres}$ --- carries
    \code{q\_max\_source = "derived"} \emph{unconditionally}, together with the names of
    the complexes it is derived from, and its $q_{\max}$ and \code{truncated} are
    functions of theirs:
    $q_{\max}=\max\bigl(q_{\max}(\text{bulk}),\,q_{\max}(\text{boundary})+1\bigr)$ with
    \code{truncated = false} when both inputs are finite and untruncated, and
    $q_{\max}=\min\bigl(q_{\max}(\text{bulk}),\,q_{\max}(\text{boundary})+1\bigr)$ with
    \code{truncated = true} otherwise. \code{"set"} is legal on a \emph{source} complex
    only: a derived complex's $q_{\max}$ is never a user's to choose, which is
    rider~(i) of \ref{open:ell2} applied one level up. On the shipped model this gives
    the cone $2$ / \code{false} / \code{"derived"} / $[0,1,2]$, from bulk $2$ and
    boundary $1$.

    \smallskip
    \textbf{And the names are the record's own derived flag, so that the rule is
    checkable from the record} (v1.7.0; \file{docs/implementation/slot5/WP-1\_report.md}
    finding~3). The clause above has been a rule about \emph{conforming records} since
    v1.5.0, but a verifier could not apply it, because nothing in a record said which of
    its complexes were derived: a hand-assembled record could declare
    \code{q\_max\_source = "set"} on the cone and be accepted. It is therefore binding
    that \textbf{the list of names is present and non-empty on a derived complex and
    absent or empty on a source complex} --- it is not decoration beside
    \code{"derived"}, it is the predicate --- and consequently that a record declaring
    \code{q\_max\_source = "set"} on a complex whose list of names is non-empty is
    \textbf{non-conforming}, and that a verifier which checks a record's internal
    consistency \textbf{refuses it}. This costs no field: v1.5.0 already required the
    names. What it buys is the property this section exists for, that an obligation be
    checkable ``\emph{after} the run, by someone who did not make it'': the alternative
    reading --- that the rule binds only the builder --- leaves a conforming-looking
    record that claims a user chose a degree bound no user may choose, and leaves it
    undetectable by construction. The builder's own guard (a derived object asserting its
    fields as it emits them) stays; the two are complementary, and neither is the other's
    substitute. \emph{Implementation note}: the shipped extension record already carries
    the names (\code{derived\_from} on its \code{pair} object, empty elsewhere) and the
    cone builder already guards its own fields; what this rule adds is the record-level
    refusal, which \code{HeckeExtRecord::verify()} makes (code~22) and which
    \code{HeckeExtRecordValidation} \emph{asserts}, with a positive control on a source
    complex, since hecke-extensions slot-5 WP-11a (v1.7.1; at v1.7.0 the refusal had not
    yet landed and that validation pinned the gap in its \emph{negative} form).
  \end{itemize}

\item \textbf{The extension record}, \code{hecke\_ext/v1} (v1.6.0),
  \code{hecke\_ext/v2} \textbf{since v1.7.0}. A second writer,
  \emph{not} an extension of the per-degree one: one JSON per $(n,\ell)$ cell, about
  several complexes, whose field inventory is [E, \S7] and whose rules binding on this
  document are the nine below --- (a)--(f) since v1.6.0, (g)--(i) with \code{v2}. It is a \emph{separate} writer by construction, because the
  per-degree writer's text is a golden object and a new key in it is a re-baseline
  (\S\ref{app:changelog}, \ref{CL5}); the two share the conventions of item~2 and the
  degree rule of item~5 and nothing else.
  \begin{enumerate}[label=(\alph*)]
  \item \textbf{Settings are part of the record, and part of its hash input.} The record
    carries every mathematical setting in force --- its \emph{value} and its
    \code{source} $\in\{\code{default},\code{spec},\code{console}\}$ --- together with
    the \textbf{fingerprint} of the \emph{non-default} entries, which is $0$ at the
    defaults. The fingerprint enters every hash the record reports and the identity of
    the job that wrote it, so that only the default configuration is golden and a golden
    mismatch can never be explained by a setting. This is the schema half of owner
    decision~R1: the mathematical settings are the family of switches each of which may
    change a published value, as no \code{ENGINE\_*} variable may. \emph{How} a value was
    set is therefore itself a value-bearing field: a record whose \code{source} says
    \code{console} and one whose \code{source} says \code{default} are different records
    even at the same value.

    \smallskip
    \textbf{The block names what it carries, and the exclusion is by family}
    (v1.7.0; the first half of the item v1.6.0's changelog deferred: whether the settings
    block should name the engine settings it excludes. \emph{It should not}, and the
    reason is a rule rather than a taste). A setting belongs in the block iff it may
    change a published value --- that is the definition of the family, owner
    decision~R1 --- and the \code{ENGINE\_*} variables are excluded \emph{as a family},
    by the rule and not by an enumeration. Enumerating them would make the settings block
    a function of the build: a record would acquire a key whenever the engine gained a
    knob, two builds at identical mathematical settings would emit different blocks, and
    a reader could not tell a new engine knob from a new mathematical one. One sentence
    naming the criterion and the excluded family is the whole content; the list of
    mathematical keys in force is [E, \S7]'s, and grows with its consumers.

    \smallskip
    \textbf{And a resumed run writes the settings block a \emph{fresh} run at the same
    effective values would write} (v1.6.0; slot-4 WP-2 finding~4). A resume is the second
    place a value enters the block, and the rule is not implied by the first: an
    implementation that installs the budget it was resumed \emph{to} makes a resume to the
    default budget produce a different settings text, and therefore a different
    fingerprint and a different hash, from a fresh run at the default --- although the two
    runs computed the same thing. The binding rule is the equality: \emph{a resume to a
    value equal to the default installs the default's canonical text}, any other value
    installs its own, and ``the resumed record equals a fresh default run byte for byte,
    fingerprint included'' is then a statement an implementation can and must assert.
    Without it two conforming implementations differ in a published hash.
  \item \textbf{A budgeted item carries a \code{resumable} object at its own site},
    with its kind, whether it is \code{complete}, the budget, the point reached and the
    budget history; the record carries the count of incomplete items. Inline while the
    state is small; otherwise a \textbf{sidecar} file named by the record layer --- the
    only layer that knows the record's stem --- and referenced by \textbf{a path relative
    to the directory holding the record} (in practice its basename, so that record and
    sidecar move together) together with its FNV-1a-64 from the one byte-wise FNV of
    \file{ModArithmetic/Fnv.h}. \textbf{The name must be injective}: a record may carry
    several oversize states of the \emph{same} kind --- one per ideal --- so the first of
    a kind may take the bare \code{<stem>.state.<kind>.json} and every further one is
    disambiguated by an index. A writer with no stem, which cannot name a file, inlines
    the state and says so with a flag rather than emitting a path that names nothing. \textbf{The record's count of
    incomplete items counts record \emph{entries}}, not distinct decisions (v1.6.0;
    slot-4 WP-2 finding~5): where a schema attaches a verdict that is a property of the
    whole run --- a memory-budget refusal, say --- to each of several sub-objects, one
    refusal is reported once per sub-object and the count rises with them. That is a legal
    schema, and it is what the shipped extension record does; what is \emph{not} legal is
    leaving it unsaid, because a reader who takes the count for a number of decisions will
    be wrong. A schema that carries a run-level verdict per sub-object states that it does.

    \smallskip
    \textbf{And in the extension family, since \code{hecke\_ext/v2}, it may not do it}
    (v1.7.0; owner ruling \code{[OD]} slot-4 close-out 5 of 2026-09-22). A verdict that
    is a property of the \emph{whole run} --- the memory-budget refusal of $U_\ell$ is
    the instance --- is carried \textbf{at record level, once, with a single
    \code{resumable}}, and the record's count of incomplete items is then a count of
    \textbf{decisions}, entries and decisions coinciding. The paragraph above stands as
    the general rule for a schema that does replicate such a verdict, and it is why this
    repair is a \emph{version bump} and not an amendment in place: the value of an
    existing key changes, so a reader must be able to tell the two shapes apart by the
    schema string alone (\ref{CL5}'s discipline, as \code{hecke\_report/v1}
    $\to$ \code{v2} set it in slot~2). The bump is made \textbf{before the first
    emission} of an extension record as a golden object, which is what makes it free.
  \item \textbf{An unavailable complex is a placeholder, never a silent omission.} A
    complex a record is about but has not computed is emitted as
    \code{\{"status": "unavailable"\}} and nothing else --- distinguishable, as item~4
    requires of a certificate, from a complex that was computed. The four degree fields
    of item~5 are emitted \emph{only} for an available complex, because they would
    otherwise be a claim about a complex that does not exist in this record.

    \smallskip
    \textbf{One level down, the same rule for a scalar and for a flag} (v1.7.0;
    \file{docs/implementation/slot5/WP-1\_report.md} findings~2 and~4). A key a record's
    schema requires it to carry but which \emph{this} record has not computed is emitted
    as \code{null} --- never as a default value, and never dropped. Dropping it makes a
    closed key list a function of the cell, which destroys the property that made the
    list a rule; emitting a default conflates two different facts, and item~4's ``absent
    versus ran and passed must be distinguishable'' is the same requirement one level up.
    So a dimension not defined until a complex this record does not carry is \code{null}
    and not $0$; and a \emph{diagnostic} this record never evaluated is \code{null} and
    not \code{false} --- under rule~\ref{A6} a \code{false} on a diagnostic means
    ``evaluated, and it did not hold'', which is conformance, so a record writing
    \code{false} for ``nobody asked'' has destroyed the distinction the diagnostic exists
    to carry.
  \item \textbf{Reasons and phases are named, not numbered ad hoc.} A verdict that is not
    final carries a reason from the fixed list of [E, \S7]; the entry that names a failed
    certificate is written as \code{certificate\_failed:<code or certificate name>},
    where a \emph{code} is a non-zero clause index of the \code{verify()} that failed and
    a \emph{name} is one of the closed list of certificate names of that section --- a
    certificate route may report ``the sweep ran and its answer was no'' without any
    \code{verify()} having failed, and the vocabulary has to be able to say so.
    \emph{Where the list lives} (v1.7.0; the second half of the item v1.6.0's changelog
    deferred). The closed list of certificate names is [E, \S7]'s, with the rest of the
    field inventory, and is not restated here: this document binds that the list is
    \textbf{closed}, that the two payload forms --- a \code{verify()} clause index and a
    certificate name --- are distinguishable by a reader who has only the record, and
    that a payload outside the list makes the record non-conforming. A name added to the
    list is an [E] revision and not a [C] one. Where a
    routine reports the phase it reached, the record carries the \textbf{per-item} value
    and the run's \textbf{high-water mark} under separate keys: the two are different
    facts, and a phase belonging to a block that a split has since removed is the run's,
    not that block's.
  \item \textbf{The two recorded cone flags.} \code{phi\_flat\_identity} and
    \code{residual\_flat\_zero} are the names under which a record reports the flat
    readings of Obligation~\ref{obl:coneexact}(i), inside the object of the complex they
    are about. They are diagnostics under rule~\ref{A6}: recorded, never branched on, and
    a \code{false} is conformance. The names are reserved here so that the value they
    will carry has one spelling from the start.
  \item \textbf{What this document binds and what it does not.} The field inventory,
    the example and the per-block content of the extension record are [E, \S7]'s; this
    item binds only (a)--(e) and (g)--(i), which are schema rules, and item~2's
    conventions and item~5's degree rule, which apply to every record of the programme.
  \item \textbf{A record about several complexes reports the maps between them}
    (v1.7.0, with \code{hecke\_ext/v2}; the objects are [E, \S7]'s). Where a record
    carries two complexes joined by a long exact sequence --- the pair of
    \S\ref{sec:cone} is the instance --- the object of the \emph{derived} complex carries
    one entry per spot of that sequence, and each entry says, per spot, what
    Obligation~\ref{obl:coneexact}(iv) makes of it: the ranks, whether the composite was
    found zero, and whether exactness was \emph{asserted} there. Item~5's reported-degree
    rule governs which spots exist, and rule (c)'s \code{null} governs the rest: a spot
    whose composite was not built, and a spot at which the sequence's exactness was not
    asserted, carry \code{null} and not \code{false}, because clause (iv) asserts
    exactness only where all three terms are reported and a record may not report a
    verdict the obligation did not reach. A Hecke-equivariance flag on such a sequence is
    the \emph{action}'s claim and not the complex's: a record carrying the cone but no
    cone action has not evaluated it, so it is \code{null}.
  \item \textbf{A $\pm$-split record is two further records of the same cell}
    (v1.7.0, with \code{hecke\_ext/v2}). Where an implementation carries the involution's
    splitting of Proposition~\ref{prop:iotasplit}(1), each record --- the unsplit one and
    the two per-sign ones --- names its sign, every per-block object inside a record
    carries the \emph{same} sign as that record, and the two per-sign records' dimensions
    sum, degree by degree, to the unsplit one's. The sum is \textbf{asserted} and not
    recorded: it is Proposition~\ref{prop:iotasplit}(1) itself, and a record family that
    fails it is not a report about one cell. At $2\notin\cring^{\times}$ there is no
    split and the sign field says so, which is item~4's ``absent and ran'' distinction
    again (Obligation~\ref{obl:iota}(v)).
  \item \textbf{A matching a record performed is a matching it reports}
    (v1.7.0, with \code{hecke\_ext/v2}). Where a record's objects are indexed by one
    thing but identified across another --- the extension record's maximal ideals, which
    a run matches across the ambient groups it computed --- the record carries, per
    object, the list of the ambient objects it was matched \emph{against}, the verdict
    per entry, and the \textbf{witness} the test used, including the identity entry for
    the object's own ambient group, which is what makes two such lists comparable. A
    \emph{count} of matches is not a report of the matching: it cannot be re-checked by a
    reader who has only the record, which is what item~4 requires of a certificate and
    this section of everything. A non-match carries the reason it failed, by the same
    rule as (d).
  \end{enumerate}
  \textbf{Implemented at \code{hecke\_ext/v1}} (hecke-extensions slot 4, package WP-1,
  phase E2a, 2026-09-21; \file{docs/implementation/slot4/WP-1\_report.md}):
  \file{HeckeExt/HeckeExtRecord.h} and \file{HeckeExt/HeckeExtReport}, with
  \code{HeckeExtRecordValidation}; \file{Hecke/HeckeReport.\{h,cpp\}} is untouched.
  \textbf{At \code{hecke\_ext/v2} since v1.7.0} (slot 5, package WP-1, phase E2
  \emph{cont.}; \file{docs/implementation/slot5/WP-1\_report.md}), in the same two files
  and with \file{HeckeExt/HeckeExtComplexRecords} the one translation layer between the
  computed objects and the record: (b)'s record-level verdict with its single
  \code{resumable}, (g)'s \code{boundary} and \code{pair} objects --- computed, not (c)'s
  placeholder, so (e)'s two flags are emitted --- (h)'s per-sign records and (i)'s
  matching list. \file{Hecke/HeckeReport.\{h,cpp\}} is still untouched, which is what
  ``a separate writer'' means operationally.
  \emph{Status at v1.6.2 (the slot-4 close)}: the record \textbf{is written}, for the
  \textbf{bulk} complex, end to end from a cell --- slot-4 WP-2 added
  \file{HeckeExt/HeckeExtDriver}, the Stage~1--5 orchestration behind it, gated by
  \code{HeckeExtRamanujanValidation} on [E, Claim~8.1] at $(10,691)$. The boundary, cone
  and per-sign objects, and the campaign's own \code{ext} columns, are slot 5's.
  \emph{Status at v1.7.0}: the boundary, cone and per-sign objects are written (slot-5
  WP-1, above); the campaign's \code{ext} columns and the \emph{emission} of an
  extension record as a golden object are the slot's other packages, and the schema was
  bumped before that emission for exactly the reason (b) gives. The fuller status
  sentences are the slot close's.
  \emph{Status at v1.7.4 (the slot-7 close)}: \textbf{the record is
  \code{hecke\_ext/v3}} (slot-7 WP-3, stage~3), bumped because the default record gained
  nine keys, all of them [E]'s field inventory under (f) --- none of the nine rules above
  moves. Rule (a) holds as written: the record carries every setting with its value and
  \code{source}, and the fingerprint enters every hash the record reports; the golden
  reference hashes the record text with the \code{settings} block reduced to that
  fingerprint (slot-7 stage~1), so that registering a key at its default moves no golden
  line. The first ext campaign publishes $505$ \code{v3} records.
\item \textbf{The integral descent record}, \code{integral\_descent\_record/v1} (v1.23.0-pre.37; plan milestone D6 of
  \file{docs/implementation/slot24/DESCENT\_plan.md} v0.4, the carried L2-4; [IB] OD-9~(a), OD-10~(a) ([X] S8-6d; v1.23.0-pre.38, referee T2);
  \file{docs/implementation/slot24/DESCENT6\_report.md}). The record of Theorem~\ref{thm:descentcert} for one cell $(N,n)$ and a finite
  set $\mathcal P$ of primes $p\nmid N$ (default: the first three of $2,3,5,7,11,13$ not dividing $N$), written under the setting
  \code{coefficient\_ring = Z} and under no other. Binding content:
  \begin{enumerate}[label=(\alph*)]
  \item \textbf{The ring and the object.} \code{coefficient\_ring = "Z"} and \code{object = "group"}: the object is
    $H^{q}(\Ga_0(N);\Vcoef)$, group cohomology, Definition~\ref{def:ringsemantics}(3) --- never the Bredon object, never $\Ceq(\Z)$. The setting
    selects the record and is never implied by it: a builder asked for the record under $\Fl$ or $\Z/\ell^{e}$ refuses (the console verb
    with a runtime error), and the descent's own tools, which build $\F_{\ell'}$- and $\Z/\ell^{K}$-objects of their own, run with the
    ring reset to $\Fl$ for the computation only. \textbf{$\Z$ is un-reserved for the descent only}: every other consumer of the ring
    --- \code{EquivariantCohomologyOnK}, \code{HeckeOperatorOnK}, the console verbs \code{symmetric\_power}, \code{cohomology},
    \code{hecke}, \code{iota}, the field job loop of the campaign --- \textbf{refuses $\Z$ by name} (\code{hecke::NotImplementedInThisBuild}
    naming the site; console exit~$5$), by two mechanisms (v1.23.0-pre.38, referee m1): (i) every construction site reads the exponent through the one
    reader \code{CoefficientRingK(who)}, which refuses $\Z$ naming \code{who}; (ii) the console verbs that compute over $\Fl$ without reading the ring are
    wrapped at registration by \code{console::bind::GateVerbsOnField}, an \emph{opt-in list} of verb names (\file{Console/Bindings/RegisterPipelineVerbs.cpp},
    \file{HeckeExt/RegisterHeckeExtensionVerbs.cpp}), each wrapper calling the same reader; the list is kept by \code{IntegralDescentRecordValidation} arm~(F),
    which probes a fixed list of calls. \textbf{Residual risk}: a future $\Fl$-only verb that neither reads the ring nor is added to the list would run under $\Z$
    as over $\Fl$; a registration-time default (refuse unless a verb is tagged ring-aware) would close it and is a code follow-up, not a clause of this version. The
    general \code{Z/<m>} stays reserved and refused when set (\Espec{sec:ring}~(a)).
  \item \textbf{The degrees.} \code{q\_max = 3}, \code{truncated = true}, \code{q\_max\_source = "derived"}, \code{degrees}
    $[0,1,2,3]$, \code{periodic\_from = 2} --- item~5's four fields at record level, which report the group cohomology of (1)--(3), the bulk, an
    infinite complex. The bulk occurs a second time, as the \code{bulk} column of the pair complex of~(f), with that object's own four fields
    ($q_{\max}=4$, no \code{periodic\_from}): the two are the same complex reported to two ranges for two purposes, and they are compared in degrees
    $\le3$ (the \code{ENGINE\_CHECK} of~(e)); a verifier applying item~5 reads each object's fields for that object (v1.23.0-pre.38, referee m2). $q_{\max}=3$ is \emph{derived}, not set: by Theorem~\ref{thm:descentcert}(6) the data of (1)--(4) in degree $q+2$ are those of
    degree $q$ for every $q\ge2$ (cup product with $\res^{\Ga}_{A}e$, Proposition~\ref{prop:descentN}(5)), so degrees $0,1$ and one period
    $2,3$ determine every degree, and $3$ is the least $q_{\max}$ that does --- degrees $0$ and $1$ are not periodic ($H^{0}$ and the free
    part of $H^{1}$ have no counterpart in degrees $2$, $3$), and below $3$ the period is incomplete. \code{truncated = true} because the
    complex is not bounded; \code{periodic\_from} is what licenses a reader to continue the list, and nothing else does: in particular no
    vanishing above $q_{\max}$ is reported or implied.
  \item \textbf{The groups and the bad primes.} \code{groups}: $H^{0},\dots,H^{3}$ as abelian groups, Theorem~\ref{thm:descentcert}(1);
    \code{bad\_primes}: $\mathcal B(M_N)\cap[5,\infty)$ with its counts, Proposition~\ref{prop:badprimesN}(1); \code{torsion\_primes}: the primes whose
    torsion is read, $\{2\}\cup\{3:n\text{ even}\}\cup\mathcal B(M_N)$. \textbf{$K$ per prime}: for each torsion prime the route of the
    groups' run (\code{tree} or \code{dense} at $2$, $3$; \code{shortcut} at $3$ with $\nu_3(N)=0$; \code{sparse} or \code{bredon} at
    $\ell\ge5$), the certified precision $K$, the widest ring width $W$ read ($W>K$ when a $2$-adic differential escalated) and the
    certificate's verdict. $K$ and the route are \emph{provenance}: recorded, not hashed (item~(h)).
  \item \textbf{The characteristic polynomials.} Per $p\in\mathcal P$, $\chi^{(N)}_{p,1}=\det(x-T_p\mid\Hfr{1}(N))$ by Proposition~\ref{prop:crtcharpolyN}'s
    output rule~(5), Theorem~\ref{thm:descentcert}(2): its degree $r=b_1$, its coefficients (an integer; beyond $2^{63}$ a sign and a digit count, \textbf{and since v1.24.0-pre.1 the exact value} in the entry's list
    \code{big}, item~(l); \code{?} when not certified), \code{determined}, the number undetermined, and \textbf{the digest} $\mathrm{char}_{\Z}$ (the JSON key \code{digest}; the
    golden object \code{charZ\_T<p>}; v1.23.0-pre.38, referee T1). \textbf{Coefficient order} (v1.23.0-pre.38, referee M1): the list \code{coefficients} is
    $c_0,c_1,\dots,c_r$ of Proposition~\ref{prop:crtcharpolyN}(1), $c_i$ the coefficient of $x^{r-i}$ --- \textbf{from the leading coefficient down}, $c_0=1$ first,
    the constant term last (\code{CrtCoefficient::i}, \file{Integral/Descent/DescentHecke.h}; \code{CharpolyOverZ}, \file{DescentHeckeCrt.cpp}); at $(11,0)$,
    $T_2$ reads \code{1,1,-8,-12}, i.e.\ $x^{3}+x^{2}-8x-12=(x+2)^{2}(x-3)$. This is an \textbf{explicit exception to item~2}'s low-to-high rule, which
    governs \code{hecke\_report} and the console's \code{charpoly} objects (constant term first) and not this record. \code{v1} carries no in-band key for the
    order; one outside the hashed body (e.g.\ \code{"coefficient\_order":"high\_to\_low"}) is an open code item, and until it exists this clause is the
    convention's only statement. The digest is recomputable from this text (v1.23.0-pre.38, referee m7; \code{CharpolyZ::digest()},
    \file{DescentHeckeCrt.cpp}, is the implementation): FNV-1a-64 (\file{ModArithmetic/Fnv.h}, offset basis \texttt{0xcbf29ce484222325}) fed a sequence of
    \textbf{unsigned $64$-bit words}, each as its $8$ bytes least significant first: $r$; then, per $i=0,\dots,r$ in this order, $i$, followed --- for a certified
    coefficient --- by the word $1+[c_i<0]$ (so $1$ or $2$), $c_i\bmod2^{64}$ and $c_i\bmod3^{40}$ (canonical residues in $[0,2^{64})$, $[0,3^{40})$), and --- for
    one not certified --- by the word $0$, the number of residue primes, and the residues in the order of the prime list. On a determined polynomial the digest is a function of the integers
    $c_i$ alone, hence of $(N,n,p)$, and \textbf{independent of every choice} (Theorem~\ref{thm:descentcert}); it is a fingerprint, not a
    certificate: $c_i\bmod2^{64}3^{40}$ with the sign does not recover a coefficient beyond about $2^{126}$, and a reader who needs the
    integers recomputes them. On an undetermined coefficient the digest identifies the residue data of the build's fixed rule for the
    residue primes (descending from $2^{30}$), not an integer; two such digests are comparable under that rule only. That case is \textbf{not reachable
    from the record} (v1.23.0-pre.38, referee N1): the record's options carry no prime limit, D3's \code{CrtOptions::max\_primes = 0} takes as many residue
    primes as $b_{p,r}+1$ needs, so every coefficient of an \code{ok} entry is certified and \code{determined = true}; it is reachable only through D3's
    \code{CrtOptions} directly. $\chi^{(N)}_{p,q}$ for $q\ne1$ is not recorded: for $q\ge2$ because $\Hfr{q}=0$; for $q=0$ \emph{by choice}
    ($\Hfr{0}$ is $\Z$ at $n=0$ with $T_p=p+1$, and $0$ otherwise --- nothing a record needs to carry; v1.23.0-pre.38, referee N3).
  \item \textbf{The torsion Hecke modules.} Per $(q,\ell)$ with $\ell$ in (c)'s \code{torsion\_primes} (v1.23.0-pre.38, referee m3: for $n$ odd that is $\{2\}$, so the
    record has the three entries $(q,2)$ and no entry at $3$ or at $\ell\ge5$, $-I$ acting by $-1$ and the cohomology being $2$-primary; an entry absent for that reason
    is not a zero group in the sense of (g)(iii)) --- $1\le q\le3$ at $\ell\in\{2,3\}$, $q=1$ at $\ell\in\mathcal B(M_N)$ (no other
    $(q,\ell)$ carries torsion, Theorem~\ref{thm:descentcert}(1); $H^{0}$ is torsion-free) --- the
    $\Z[x_p:p\in\mathcal P]$-module $(\mathrm{tors}\,H^{q}(A;\Vcoef)_{(\ell)},(T_p))$ of Theorem~\ref{thm:descentcert}(3) by its choice-free invariants
    (Proposition~\ref{prop:torsionhecke}(1)): the exponents, \code{module\_certified} ($2\sigma_{\max}\le K$), and per $p$ one canonical text ---
    the type, the layer polynomials $P_j$, the layer blocks' polynomials modulo $\ell^{e_j}$, the Fitting types, and, when the module is
    certified, the types of $\ker(T_p-c)$, $\coker(T_p-c)$ --- in one system of generators, with \code{commute}. The $\ker/\coker$ types are read, and
    the canonical text carries its \code{ker=} / \code{cok=} parts, exactly when (v1.23.0-pre.38, referee m4) the module is certified, \textbf{and} the number $r$ of
    cyclic summands is at most $400$ (\code{max\_full\_summands}), \textbf{and} $K_{\max}(\ell)\ge a_{\max}+2$, where $K_{\max}(\ell)$ is the largest $K$ with
    $\ell^{K}<2^{31}$ ($30$ at $\ell=2$) and $a_{\max}$ the largest exponent (\code{Invariants}, \file{Integral/Descent/DescentHeckeTorsion.cpp}); otherwise the text
    omits them, the entry stays \code{module\_certified = true} and unrefused, and no flag says so. The read's working width is a ladder (v1.23.0-pre.39, item~(j)):
    $k_0=\min(K_{\max}(\ell),2a_{\max}+3)$, then $K_{\max}(\ell)$ when the read falls short at $k_0$ by a precision shortfall; the width is not in the text and is not one of these caps. These caps enter the hashed canonical text, so they are
    record content: changing one is a golden event (h)(1), not a rollback. A hashed flag for the omission would be a schema revision (\code{v2}). Read on the tree model at $2$,
    $3$ (every $N$) and on the cokernel presentation at $\ell\ge5$, as Theorem~\ref{thm:descentcert}(3) prescribes. \code{module\_certified = false}
    is \emph{not} a refusal: it is the narrower certified content of Theorem~\ref{thm:descentcert}(3) (the layer polynomials and the Fitting types at
    every certified $K$), and the canonical text then omits what needs $2\sigma_{\max}\le K$. A certified read whose exponents differ from the
    groups' $\ell$-primary part, or a pair whose bulk differs from the groups in a degree $\le3$, is a \textbf{defect} (\code{ENGINE\_CHECK}),
    never a refusal and never a record.
  \item \textbf{The boundary pair.} The pair columns that Theorem~\ref{thm:descentcert}(7)(c) assigns to D6: $H^{q}$, $H^{q}(\partial)$,
    $\Hcpt{q}$ of \S\ref{sec:descentpair} (the tree route of \S\ref{sec:treepair} by default) for $0\le q\le4$, \textbf{each complex with its
    own four fields} (item~5: never a maximum over the complexes of a record) --- \code{q\_max = 4}, \code{truncated = true},
    \code{"derived"}, $[0,\dots,4]$ --- and the primes $S$ and the $2$-adic route as provenance. $4$ is the least top degree the pair builds
    admit (\S\ref{sec:descentpair}'s code map; Proposition~\ref{prop:treepair}(6) refuses $q_{\max}\notin\{4,5\}$), and it carries the segment
    $\Hcpt{3}\to H^{3}\to H^{3}(\partial)\to\Hcpt{4}$ of \eqref{eq:descentpairles}, which closes the bulk's reported range $q\le3$. No
    periodicity of the pair columns is asserted here --- Theorem~\ref{thm:descentcert}(6) is a statement about the bulk --- so nothing above $4$ is
    reported or implied. The compact-support complex is the cone, a \emph{derived} complex: by item~5 its object must carry the non-empty list of
    the complexes it is derived from (bulk, boundary), and its fields are those of item~5's rule,
    $q_{\max}=\min(4,4+1)=4$, \code{truncated = true}. \textbf{Met since \texttt{19a969de}} (v1.23.0-pre.38, referee M2): the \code{compact} object carries
    \code{"derived\_from":["bulk","boundary"]} (\code{kCompactDerivedFrom}, \file{Integral/Descent/DescentRecord.h}), asserted by arm~(J); the list is outside the
    hashed body, so it moved no hash and no golden line. (v1.23.0-pre.37: ``not met by \code{v1}'', the record as built on \texttt{538cf881}.)
  \item \textbf{Refusals.} \code{status} $\in\{\code{"ok"},\code{"refused"}\}$ with \code{refusal}, the first refused part's line; each $\chi$ entry, each
    torsion entry and the pair carry their own \code{refused} and \code{refusal} (v1.23.0-pre.38, referee m5, in place of ``every part''): the groups carry none of their
    own --- a refused groups assembly shows as an empty \code{groups} list together with the record-level \code{status} and \code{refusal} --- and the per-prime $K$
    entries carry \code{certified} only. A part's \code{refusal} may be \textbf{non-empty on a part that is not refused}: a \code{[ModulePrecision]} torsion entry is
    \code{"refused":false} (by (e)) with \code{refusal} naming the narrowed certificate; a reader keys on \code{refused}, never on a non-empty \code{refusal}. A refusal is a legitimate limit of the build (a cap, a precision, a degree range,
    rule~19's RSS ceiling), named by its code. \textbf{A refused record is never a valid empty answer}, by three rules: (i) a refused groups
    assembly reports \textbf{no} groups --- the list is empty, never partial --- whereas an \code{ok} record reports exactly $q_{\max}+1$ groups,
    the zero group included (``$0$'', not an absent entry), so an empty list occurs only together with \code{status = "refused"}; (ii) likewise
    for the pair's three lists ($q_{\max}+1=5$ entries each, or none and \code{refused}); (iii) a torsion entry with no exponents is the zero
    group only when it is not \code{refused}. A reader reads \code{status} and the part's flag first; \code{status = "refused"} with parts
    present means that the parts present are certified and the refused part is absent, as at $(1,200)$ below. A defect is never a refusal: it
    throws (\code{hecke::HeckeInternalError}, \code{ENGINE\_CHECK}) and no record is written. A record of this schema is built with every part
    (v1.23.0-pre.38, referee N2): the builder's internal options \code{RecordOptions\{hecke = false\}} and \code{\{pair = false\}} give an \code{ok} record without
    $\chi$ and torsion, or with \code{"pair":\{"ran":false\}}, which is \textbf{not} a conforming \code{v1} record; no caller uses them (dropping them, or recording
    their state, is a code follow-up).
  \item \textbf{The hash: what is hashed, what is not.} \code{record\_fnv1a64} is the one FNV-1a-64 (\file{ModArithmetic/Fnv.h}) of a
    fixed-order text, \code{DescentRecordZ::body()} (\file{Integral/Descent/DescentRecord.cpp}), \textbf{normative for the order}; the partition is exhaustive
    (v1.23.0-pre.38, referee m6). \textbf{Hashed}: the schema, $(N,n,\mathcal P)$ with $\mathcal P$ ascending, ring and object, (b)'s fields, the groups, the bad
    primes with their counts and the torsion primes, per $p$ the degree $r$, \code{determined}, the digest of (d) and the \code{refused} marker, (since v1.24.0-pre.1) the exact value of every coefficient beyond $2^{63}$ (item~(l)), per $(q,\ell)$ the
    exponents, \code{module\_certified}, the canonical texts of (e) and the \code{refused} marker, the pair's $q_{\max}$ with its \code{refused} marker and its
    groups, \code{status}, and the settings' canonical text with their fingerprint. \textbf{Not hashed}: $K$, $W$, the per-prime certificate verdicts and every
    route (of the groups, of the torsion reads, of the pair), the pair's prime set $S$, the coefficient strings of (d) and \code{undetermined}, \code{commute},
    the pair objects' four fields and \code{derived\_from}, the settings' JSON (values with their sources), the refusal texts, timings and RSS. \code{commute} is
    record content (e) without being hashed: since \texttt{19a969de} the golden \code{DESCENT} points \code{ENGINE\_CHECK} it (with \code{module\_certified}), and
    the hash does not see it. The reason is that
    every hashed field is an \emph{invariant}: by Theorem~\ref{thm:descentcert} the data of (1)--(3) are independent of every choice made in
    computing them, models, comparisons, Smith transforms and primes included, and by Remark~\ref{rmk:descentcertscope}(i), (ii) the torsion
    data are certified exactly, so a record's invariants have one value per cell and setting. The routes and the precisions are such
    choices, and the descent's rollbacks --- \code{HECKE\_DESCENT\_TREE=0} (\S\ref{sec:isomodeltree}'s code map) and
    \code{HECKE\_DESCENT\_TREE\_PAIR=0} (\S\ref{sec:treepair}'s code map, obligation (c)) --- are \textbf{read-only and value-preserving},
    the preserved values being exactly ``the record's invariants (the group types and the choice-free Hecke invariants), not Smith
    generators or coordinates''. A field that a rollback may change is therefore a field the hash must not contain, and conversely: with
    $K$ and the routes outside the body, \textbf{a rollback cannot change the hash of a record that is \code{ok} on both routes}. Two
    qualifications, both by design. (1) The certification flags (\code{module\_certified}, \code{determined}) are hashed; they are route-independent
    because the torsion read is on the tree at $\ell\in\{2,3\}$ and on the cokernel at $\ell\ge5$ under every rollback (Theorem~\ref{thm:descentcert}(3)),
    at the build's fixed read width $\ell^{K}<2^{31}$, and $\chi$ is read on the $\mathsf O$-piece alone (Remark~\ref{rmk:descentcertscope}(iv)); changing
    a fixed width, a cap or a canonical text is a change of the record's content, gated as a golden event, not a rollback. (2) A refusal is
    not a value: a cap reached on one route and not on the other gives different \code{status}, hence different hashes, and that is correct.
    A non-default mathematical setting changes the settings text, hence the hash and the campaign job id (\code{locality = strict} changes
    \code{record\_fnv1a64} at $(1,12)$ from \texttt{0x6c841fb033fbd3d1} to \texttt{0x0be335f8f71d98cc}, the invariants unchanged; these two values are
    \textbf{illustrative}, as printed by arm~(J) at \texttt{724b9aaa} (\file{DESCENT6\_report.md} \S4), which asserts only that they differ; no validation pins them,
    and pinning them in~(J) is a code follow-up; v1.23.0-pre.38, referee m8).
  \item \textbf{What is not in the record.} Theorem~\ref{thm:descentcert}(4) (the invariant factors of $T_p-c$ on the free quotient
    $\Hfr{1}(N)\otimes\Z_{(\ell)}$) and~(5) (level pairs $\Ga_0(MN)\subset\Ga_0(N)$ with $T_p$) are \textbf{not recorded and not built} in the record ---
    (4) exists at $\ell\ge5$ as D3's \code{FreeQuotientFactors} and is not called, (5) is read by no built code; (7)'s non-assertions stay
    non-assertions, and integer Hecke matrices in a basis are not named by the ring (Definition~\ref{def:ringsemantics}(3)).
  \item \textbf{The refusal at $(1,200)$, $\ell=2$, and its name.} \textbf{Up to v1.23.0-pre.38} (this paragraph is that version's text, kept as history;
    the paragraph ``The width ladder'' below is the current statement) at $(N,n)=(1,200)$ the record was \code{refused} at
    $(q,\ell)=(1,2)$ with \code{REFUSED [KerCokerPrecision] ker / coker of phi - 0: kernel: I in L not certified (30 < 34)} (since \texttt{19a969de},
    \file{Integral/Descent/DescentHeckeTorsion.cpp}; v1.23.0-pre.37 quoted the code name of \texttt{538cf881}, \code{[NotEndomorphism]}), the groups, $\chi$
    and the other parts present. What fails is the read of the types of $\ker(T_p-c)$, $\coker(T_p-c)$ (Proposition~\ref{prop:torsionhecke}(1)(c))
    on the certified module: the code computes them on the presentation $\bigoplus_i\Z/2^{a_i}$ over $\Z/2^{k}$, $k=\min(30,2a_{\max}+3)$ (pre.37 and pre.38
    called this ``the ring width capped at $2^{30}$''; at $(1,200)$ it is $k=11$, see below), and at $(1,200)$ the Smith form of the image lattice certifies $30$ of the $34$
    pivots its rational rank requires --- $30<34$ counts pivots, not bits. It is a \textbf{precision shortfall of the build}, not a statement
    about $T_p$: by Theorem~\ref{thm:descentcert}(3) $T_p$ is an endomorphism of the certified module, which is what \code{NotEndomorphism} would
    deny. \textbf{The taxonomy of \texttt{538cf881} misnamed it}, and v1.23.0-pre.37 fixed the name: a rank shortfall or an exhausted precision inside the
    $\ker/\coker$ read is \code{[KerCokerPrecision]}, a precision refusal of the family of \code{[ModulePrecision]} (the module and the
    other data certified, the $\ker/\coker$ types not read at this width); \code{[NotEndomorphism]} stays reserved for a failed
    endomorphism safeguard (a relation not carried into the relations, a non-zero free coordinate, an image off the lattice), which by
    Proposition~\ref{prop:heilbronn}(1) and Theorem~\ref{thm:descentcert}(3) signals a defect of the code rather than a limit. The code emits
    \code{[KerCokerPrecision]} since \texttt{19a969de} (\code{D3Refusal::KerCokerPrecision}, \file{Integral/Descent/DescentHecke.h}). (v1.23.0-pre.37's reading rule
    for a record of \texttt{538cf881} --- its \code{NotEndomorphism} with ``not certified'' or ``precision exhausted'' read as \code{KerCokerPrecision} --- is
    vestigial: no merged build writes such a record.) The renaming moves no hash (refusal texts are unhashed) and no golden line (no golden point refuses). Keeping the module's certified parts at
    such an entry --- as \code{[ModulePrecision]} does --- would change the hashed body and is a schema revision (\code{v2}), not part of the fix.
    (Pre.38 ended: ``a wider ring at $2$ for this read ($W=61$, as the groups' escalation) would lift the cell''; the ladder below lifts it at $k=30$, without $W=61$.)
    \textbf{The width ladder} (v1.23.0-pre.39, ICENSUS-P1; \file{docs/implementation/slot24/ICP1\_acceptance.md}). The width of the paragraph above is $k_0=\min(30,2a_{\max}+3)$,
    at $(1,200)$ $k_0=11$ ($a_{\max}=4$), not the word bound. Since v1.23.0-pre.39 a read that falls short at $k_0$ by a \emph{precision shortfall} (a rank over $\Z/\ell^{k}$ below the rational rank, or a
    precision exhausted inside the read; \code{MapRes::precision}) is redone at $k=K_{\max}(\ell)$ ($30$ at $2$, $19$ at $3$, and so on: the word bound $\ell^{K}<2^{31}$ of
    this read), at every $\ell$, not only at $2$; only a shortfall there refuses \code{[KerCokerPrecision]}. $W=61$ is not needed. A failed safeguard inside the read (a relation
    not carried into the relations, $I$ not contained in $L$, the rank input) is not escalated and refuses \code{[NotEndomorphism]}, a defect by the paragraph above. The read stays certified: each width certifies its own read (rank over $\Z/\ell^{k}$ equal to the rational rank, the precision checks), and a certified
    type of $\ker(T_p-c)$, $\coker(T_p-c)$ is choice-free (Proposition~\ref{prop:torsionhecke}(1)), so it does not depend on $k$. An entry that reads at $k_0$ is
    computed as before, so its text and hash are unchanged; an entry that used to refuse now reads, and its \code{status} and hash change, as (h)(2) says they should. No
    golden point refuses, so no golden line moves. The six level-$1$ cells $n\in\{48,64,72,80,120,200\}$, $\mathcal P=\{2,3,5\}$, where the ICENSUS probes refused at $(q,\ell)=(1,2)$,
    now read (\code{refused=0}). Arm~(L) of \code{IntegralDescentHeckeValidation} checks $n\in\{48,64,80,200\}$: the module against the groups' $2$-primary part; each
    escalated read refuses when forced back to $k_0$, and each other read gives the same text when forced to $k=30$; every type of $\ker(T_p-c)$, $\coker(T_p-c)$ against an
    independent $2$-adic Smith form (of $[F\mid\mathrm{diag}\,2^{a_i}]$ for the cokernel, of the Pontryagin dual for the kernel); and that each cell escalates at least one read.
    The canonical text at $(1,48)$ also equals that of the Python oracle \file{docs/implementation/slot24/d3\_hecke\_oracle.py}.
  \item \textbf{Parts not run} (v1.23.0-pre.40). A campaign may compute only a prefix of the record's parts (the spec key \code{parts} $\in\{\code{all},\code{groups,hecke},\code{groups}\}$,
    default \code{all}): a part not run is marked \code{"ran":false} --- the pair already by item~(f)'s \code{pair}, the $\chi$ and torsion entries by \code{"hecke":\{"ran":false\}}, written only then ---
    and its arrays are empty \emph{by absence of the part, not by result}, never a zero group. The hashed body carries the line \code{hecke=not\_run} when the Hecke part was not run, and omits
    the pair lines when the pair was not run (item~(f)); either way a partial record never has the hash of a full one. A full record is unchanged byte for byte, and a non-default \code{parts}
    enters the cell id (\code{\_partsGH}, \code{\_partsG}); the default adds nothing. A \code{groups} record also skips the cross-check of the torsion read against the groups' $\ell$-primary part
    (item~(e)), which lives in the Hecke part: its groups are certified by their own assembly (Theorem~\ref{thm:descentcert}(1)), but not cross-checked as in a full record.
  \item \textbf{The faithful coefficients} (v1.24.0-pre.1; plan item FMTMP, owner request of 2026-11-03,
    \file{docs/implementation/slot24/FMTMP\_acceptance.md}). A $\chi$ entry of~(d) carries \textbf{every certified coefficient exactly}: one with $|c_i|<2^{63}$
    as before by its decimal integer in \code{coefficients}, one with $|c_i|\ge2^{63}$ --- whose \code{coefficients} token stays \code{[}$\pm$\code{,}$D$\code{ digits]}
    ($D$ the number of residue primes used, unchanged) --- \textbf{also by its value}, in the entry's list
    \code{"big":[\{"i":}$i$\code{,"hex":"}$h$\code{"\}},\dots\code{]} ($i$ the index into \code{coefficients}, ascending), written only when the entry has such a
    coefficient. \textbf{The encoding} $h$ is the canonical signed hexadecimal text of the integer $c_i$: \code{-} iff $c_i<0$, then \code{0x}, then $|c_i|$ in
    lowercase hexadecimal digits without a leading zero (the regular expression \verb+^-?0x(0|[1-9a-f][0-9a-f]*)$+, \code{-0x0} excluded); it is a function of the
    integer alone --- not of the residue primes, their number or the mixed-radix digits of Proposition~\ref{prop:crtcharpolyN}(4) from which the builder computes it ---
    hence \textbf{choice-free} in the sense of~(h), and unique. It is produced with word arithmetic only (Horner on a base-$2^{32}$ limb vector; no
    multi-precision library: the owner's ruling [R-Q2], no bignum, stands), and \textbf{checked at build time}: an independent decode of $h$ (Horner on its
    hexadecimal digits with the word \code{mulmod} of Proposition~\ref{prop:crtcharpolyN}(4)) must give D3's $c_i\bmod2^{64}$, $c_i\bmod3^{40}$ and sign, else the
    builder stops (a defect, never a refusal). \textbf{Hashed}: the body of~(h) gains, directly after the line of the $\chi$ entry, one line
    \code{charpoly p=}$p$\code{ i=}$i$\code{ exact=}$h$ per such coefficient, so \code{record\_fnv1a64} covers the values; a record \emph{without} a coefficient beyond
    $2^{63}$ has no such line and no \code{big} key, and keeps its bytes and its hash (\code{0x6c841fb033fbd3d1} at $(1,12)$ stays). For a coefficient below $2^{63}$ the
    digest of~(d) already determines it (its residue modulo $2^{64}$ with its sign), so with the new lines the body determines every coefficient of a determined entry.
    The digest of~(d), its text and the golden objects \code{charZ\_T<p>} do not change. A reader decodes $h$ by one call (Python \code{int(h, 16)}; a decoder without
    sign handling strips the \code{-} first) and re-verifies the entry with the digest of~(d), recomputed from the integers; the reference decoder is
    \file{docs/integral\_census/icensus\_exact.py}. \textbf{Old cells.} A cell written before v1.24.0-pre.1 has tokens without \code{big}: its large coefficients are not in
    the file, and readers report them as unavailable, never guessed. The campaign's cell ids and resume key do not change (the mathematical settings are unchanged; the
    precedent of~(f)'s \code{derived\_from} and~(k)): a resumed campaign directory accepts a complete old-format file, so a re-run that is to carry the values writes
    to a fresh directory. \textbf{Size}: the text grows with $\sum_i\log_{16}|c_i|$; the bound of Proposition~\ref{prop:crtcharpolyN}(3) caps it at about
    $\tfrac14\sum_i b_{p,i}$ characters per $p$, and a cell file can reach several MB --- at $(N,n)=(90,20)$, $r=756$, the cell file measures $7.95$\,MB
    ($2259$ values, the largest of $29\,415$ bits; $72$\,kB before), against the bound's $\le17.5$\,MB (\file{FMTMP\_acceptance.md} \S4); no compression is applied.
  \end{enumerate}
  \textbf{Code map} (v1.23.0-pre.37, \textbf{restated at v1.23.0-pre.38}, referee M2: built on branch \texttt{s24-d6impl}, tip \texttt{538cf881};
  \textbf{audited} --- \file{docs/implementation/slot24/D6\_code\_audit.md}, $2$~M / $3$~m / $8$~N, the fixes at \texttt{19a969de}, N5 included --- and
  \textbf{merged} into \texttt{hecke-extensions} at \texttt{2aa61397}; \file{docs/implementation/slot24/DESCENT6\_report.md} \S\S2--8, \S8 the audit's fixes). The setting in
  \file{Hecke/HeckeSettings} (\code{ParseCoefficientRing("Z")} stores $k=0$; \code{CoefficientRingK(who)} and \code{RequireFieldCoefficientRing(who)}
  refuse $\Z$ by name; \code{CoefficientRingIsZ()}); the record in \file{Integral/Descent/DescentRecord} (\code{BuildDescentRecord}: D1's
  \code{AssembleDescentN}, D3's \code{RunDescentHecke}, D4/D5c's \code{AssembleDescentPair}); the console verb \code{descent(N, n, primes)}
  returning \code{DescentRecord} (\file{Console/Bindings/Descent}), \code{read\_only} with its \code{--serve} twin, bounded by $N\le300$,
  $n\le60$, $\psi(N)(n+1)\le300$ and (the audit's m2) $1\le|\mathcal P|\le6$, every $p\le97$, $\sum_{p\in\mathcal P}p\le100$ (measured worst cases at the
  bound $\approx40$~s and $\approx1.4$~GiB; the default primes sum to at most $31$), $\mathcal P$ sorted ascending; the campaign's descent cells (\file{Campaign/CampaignDescent}, \file{docs/matrix-engine/campaign.md} \S17; one
  cell per $(\code{level},n)$, \code{integral\_descent\_cell/v1}); \code{IntegralDescentRecordValidation} (\code{standard}; arms (S), (F), (R) against
  the four oracle tables outside the build --- \file{descentN\_oracle.tsv}, \file{descent\_oracle\_n128.tsv}, \file{d3\_hecke\_oracle.tsv}, \file{d4\_pair\_oracle.tsv}
  --- on $13$ cells (v1.23.0-pre.38, referee N4, in place of ``three oracles''), (J), (K), (B), (G)). \textbf{The golden event R1 has landed} at \texttt{6958a5d2}: nine points
  \code{DESCENT N=<N> n=<n>}, $156$ tier-A additions, $0$ changed, $0$ removed ($18117\to18273$), predicted before any build and dry-run into scratch
  (\file{docs/implementation/slot24/d6golden/}, ledger \file{d6golden/ledger.md}). \textbf{The four code items of v1.23.0-pre.37 are closed} at \texttt{19a969de}
  (and \texttt{8fb4ab2c} for (4)): (1) the $\ker/\coker$ refusal is \code{[KerCokerPrecision]}, (j); (2) the \code{compact} object carries \code{derived\_from}, (f);
  (3) arm~(B) compares \code{record\_fnv1a64} and the hashed body under \code{HECKE\_DESCENT\_TREE=0} and \code{HECKE\_DESCENT\_TREE\_PAIR=0} with the default,
  in process and through the real environment variables in a child, so (h)'s ``a rollback cannot change the hash'' is pinned on the arm's cells; (4) arm~(G)
  checks the points present in the reference with equal hashes and pins the $156$ values (\file{d6golden/values.tsv}). \textbf{Open code items} (v1.23.0-pre.38;
  none is a clause, none moves a hash): an in-band coefficient-order key outside the hash, (d); pinning the two $(1,12)$ hashes of (h) in~(J); the
  \code{RecordOptions} \code{hecke}/\code{pair} switches, (g); a registration-time ring default for console verbs, (a).
\item \textbf{The integral Hecke structure record}, \code{integral\_hecke\_structure/v1} (v1.25.0-pre.0 on the branch \texttt{s24-wip-heckez-z0}, refereed A7 / A7b; in the main line from v1.25.0-pre.1, 2026-11-06, a pre-release, not owner-released; the HeckeZ plan
  \file{docs/implementation/slot24/heckez/04-record-census.md} v0.4 \S1, owner rulings Q1(a), Q2(a), Q3(a), Q12, Q17 of 2026-11-04). The record of \S\ref{sec:heckeZ} for one
  cell $(N,n)$ and one algebra. \textbf{Not built} (the plan's package HZ-REC, slot~25); this item fixes what it must contain. (v1.25.0-pre.4, 2026-11-06: built by HZ-REC on \texttt{s25-wip-hzrec}, not merged; its schema audit \file{docs/implementation/slot25/HZREC\_schema\_audit.md} led to the clauses marked
  ``pre.4'' below, binding changes \textbf{not yet seen by the owner}. v1.25.0-pre.5: its re-audit \file{docs/implementation/slot25/HZREC\_schema\_reaudit.md} led to the clauses marked ``pre.5'', likewise
  \textbf{not yet seen by the owner}.) Binding content:
  \begin{enumerate}[label=(\alph*)]
  \item \textbf{The ring, the object and the key.} \code{coefficient\_ring = "Z"}, \code{object = "group"}, as item~7~(a): the object is $H^{q}(\Ga_0(N);\Vcoef)$ and the algebras of
    Definition~\ref{def:heckeZ}; a builder asked for the record under any other ring refuses. The key is $(N,n,\code{algebra})$, $\code{algebra}\in\{\code{full},\code{anemic}\}$,
    only \code{full} at $N=1$. The record is \textbf{a new record beside} \code{integral\_descent\_record/v1} (Q2(a)): it references its descent cell by
    \code{descent\_ref} $=\{$\code{record\_fnv1a64}, \code{cell\_sha256}, \code{primes}$\}$ and copies nothing item~7 already holds, so integral census edition~1 and the
    golden entries of item~7 stay byte-identical. It is selected by the campaign key \code{structure = hecke} (default \code{none}: the descent campaign, byte-identical), which
    chooses the record schema and is \emph{not} a twelfth mathematical setting (Q17(v)).
  \item \textbf{The two algebras and the rank.} \code{d} is the a-priori rank of Proposition~\ref{prop:heckeZrat}: $\dim M_k(\Ga_0(N))$ for \code{full}
    (\code{d\_source = "hES\_N"}), $d^{\mathrm{an}}$ of (3$'$) for \code{anemic} (\code{d\_source = "hES\_N\_an"}), with \code{b1} and \code{e\_k}; never the computed
    $\deg f_\vartheta$. The \code{operators} actually used ($T_p$, $p\nmid N$; $U_p$, $p\mid N$, \code{full} only), ascending.
  \item \textbf{Statuses: one closed list, each final or pending; no exact value where only a bound is proved.} Every structural field carries one of: \code{certified},
    \code{upper\_bound\_certified}, \code{non\_reduced\_certified}, \code{bounds} (with both values), \code{zero}, \code{distinct\_certified}, \code{common\_certified},
    \code{agrees\_on\_P} (not a certificate), \code{partial(unfactored)}, \code{open(sandwich)}, \code{open(cap)}, \code{open(precision)}, \code{refused(<reason>)} with
    $\code{<reason>}\in\{\code{by\_theory},\code{non\_reduced},\code{word\_modulus},\code{no\_lift},\code{precision}\}$, \code{not\_run} --- all \textbf{final} (they change only when the mathematics or
    the settings change; two readings apart (v1.26.0-pre.8, confirm-referee D1): the derived \code{maximal\_ideals = not\_run} of item~(d$'$)(ii), on any record (v1.26.0-pre.6, referee-light R1; v1.26.0-pre.7, confirm-referee C3), and the aggregate \code{G3.C\_T.status} while a \code{C\_T} read or \code{G3.C\_T.lower} it depends on is pending, which is rewritten when those move (no byte of a final record moves: the record is \code{pending} by the pending field beside it)) --- and \code{blocked(export)}, \code{blocked(U\_p)}, \code{refused(export\_word)}, \code{refused(k1\_route)}, \code{refused(word\_prime\_cap)} [v1.25.0-pre.7, below], the only \textbf{pending} statuses (a named package moves them:
    EXPORT resp.\ UP-Z the first two (\code{HZ-EISUP} for \code{G1.components} and the \code{maximal\_ideals} of item~(d$'$)(ii); v1.26.0-pre.7, confirm-referee W-A), EXPORT-Z step~2 \code{export\_word}, EXPORT-Z or a slot-26 $K=1$ route \code{k1\_route}; v1.25.0-pre.2, Z0 referee F3, lead decision on O1;
    a later version with a larger count cap \code{word\_prime\_cap}), or a value
    of the closed enumerations of the plan's \S1.2. The overall \code{status} is \code{pending} iff some field is pending (a pending refusal reason included), \code{refused} iff the cell is refused, else
    \code{final}. \code{by\_theory} is (C3) at $2$, $3$, level $1$, $d\ge4$ (Proposition~\ref{prop:heckeZlevelone23}(6)); \code{non\_reduced} (v1.26.0-pre.3, re-referee W1) is the full (C1) refuted by Theorem~\ref{thm:heckeZcert}(1) at a \code{full} record, $N>1$, whose \code{G1.reduced\_Q} is \code{non\_reduced\_certified}, written in \code{G1.C1.status} only (item~(d$'$)(ii)); \code{word\_modulus} a read whose modulus $\ell^{K}$ is
    $\ge2^{63}$ (tested first, before any route is asked; $K=1$ for (C3), $K\ge2$ for the at-root read, so the latter from $\ell\ge2^{31.5}$; Proposition~\ref{prop:heckeZcong}); \code{no\_lift} the named refusal of a
    cochain lift by EXPORT or UP-Z once landed; \code{precision} as item~7~(i), (j). (v1.25.0-pre.1, 2026-11-06, as built by HZ-LOC; restated in v1.25.0-pre.2, Z0 referee F3, F4:) \code{export\_word} a read served by EXPORT (level $N>1$; at
    level~$1$ only a torsion prime of the cell, $\ell\le n+1$ in practice) with $\ell^{K}<2^{63}$ but $K$ above what EXPORT serves: $\ell^{K}\ge2^{31}$, the bound of EXPORT's word ring, or $K>K_{\max}=K_{\mathrm{run}}-\sigma_{\max}-2\sigma^{(1)}_{\max}$,
    the precision that ring loses to the $\ell$-torsion (v1.25.0-pre.3, Z0 re-check; for $\ell<2^{31}$, $K_{\mathrm{run}}$ is the largest $K$ with $\ell^{K}<2^{31}$, so there the first case is contained in the second) --- that is the binding rule.
    As built (\file{Integral/HeckeZLocal.cpp}, \code{LatticeFromExport}) the code writes it at $\ell\ge2^{31}$ for \emph{every} $K$, $K=1$ included (so at level $N>1$, $\ell\in[2^{31},2^{63})$,
    (C3) is refused as well; for $\ell\in[46341,2^{31})$ EXPORT serves $K=1$ only, (C3) runs and the at-root read is refused), and at $\ell\ge46341$ whenever $K$ exceeds EXPORT's
    $K_{\max}$; at $\ell<46341$ with $K$ above EXPORT's $K_{\max}$ it writes \code{open(precision)}, although that is the same limit of EXPORT's ring (the word bound when $\ell^{K}\ge2^{31}$, e.g.\ $K\ge3$ at $\ell\ge1291$,
    else its torsion precision), and HZ-REC writes that case as \code{export\_word}. \code{export\_word} is \textbf{pending}: EXPORT-Z step~2 (owner ruling 2026-11-06; the free block $F_p$ over $\Z$, of which $L/\ell^{K}$
    is a reduction for every $\ell^{K}<2^{63}$) widens the ring and moves these reads (\file{docs/implementation/slot25/EXPORT\_report.md}, ``For HZ-LOC / HZ-TOR'', O1).
    \code{k1\_route} (v1.25.0-pre.2, lead decision on O1): (C3) at $K=1$ at level~$1$ for $\ell\in[2^{31.5},2^{63})$, which the level-one module route does not serve (it runs at $K\ge2$);
    the code writes \code{refused(word\_modulus)} there although $\ell<2^{63}$, and HZ-REC writes \code{refused(k1\_route)}; pending, an EXPORT-Z / slot-26 item; no weight of the
    first golden set is affected. \code{open(precision)} (v1.25.0-pre.2, F4): the at-root read has not closed at the step cap of the settings ((j), count-based; not EXPORT's $K_{\max}$, which is \code{export\_word}), so it is
    final; it is written into \code{root} and \code{C\_T} of the read.
    (v1.25.0-pre.4, HZ-REC schema audit F7, F8; not yet seen by the owner:) No other pending word exists [v1.25.0-pre.7: besides \code{refused(word\_prime\_cap)}, below] (the transient \code{pending(HZ-TOR)} of the HZ-REC brief is retired:
    a Tor read the export refuses is \code{blocked(export)}). A refusal of the export is mapped \emph{by its name} (\code{ExportRefusalName}): \code{[Lift]} with a certificate code [v1.25.0-pre.5, R6, below] $\to$
    \code{refused(no\_lift)} (final), \code{[ModulePrecision]}, and a Tor read whose $K_{\max}$ lies below $e=\sigma_{\max}$, $\to$ \code{refused(export\_word)} (pending), every other
    name $\to$ \code{blocked(export)}; wherever (e)--(h) list \code{blocked(export)}, these two words may stand instead.
    (v1.25.0-pre.5, HZ-REC re-audit R6; not yet seen by the owner:) \code{[Lift]} is \code{refused(no\_lift)} \emph{only when a certificate refuses the lift}: the pair-lift code it carries
    (``the strict lift is refused [\code{<Code>}]'') is one of \code{LiftCertificate}, \code{IsobClosure}, \code{StrictIdentity}, \code{InputTableCertificate} --- then no lift exists on
    the route, which is final. A \code{[Lift]} that carries a cap of the lift (\code{RssCeiling}, \code{IntegerCap}, \code{DomainCap}) or no pair code (a refusal of the cochain-level
    operator that the export files under \code{[Lift]}) is a limit of the route and is written \code{blocked(export)} (pending).
    (v1.25.0-pre.7, HZ-RAT-PERF and its audit F2, F6, \file{docs/implementation/slot25/HZRATPERF\_audit.md}; not yet seen by the owner:) \code{refused(word\_prime\_cap)}: the
    word-prime list of the rational step (Proposition~\ref{prop:heckeZrat}) or of a $U_p$ tag (\code{full}, $p\mid N$) would exceed its count cap \code{word\_prime\_cap} $=2^{21}$
    primes. The list is the \code{CrtPrimes} order (the odd primes downward from $2^{30}-1$, all above $2^{29}$) without the divisors of $N$, and it is prefix-stable: the rational
    step sizes it from its largest per-$\vartheta$ certified target (the bound of Proposition~\ref{prop:heckeZrat} with the routes' slack) and grows it on demand, the $U_p$ tag starts
    at $\min(2048,\code{word\_prime\_cap})$ primes and doubles; both take primes in that order until their certified cover is reached, so neither the initial length nor the cap
    changes which primes are used or any value. The cap counts primes, so only a larger cap (a later version) lifts the word: it is \textbf{pending}. A refused record
    contains G0 (\code{schema}, $N$, $n$, \code{algebra}, \code{descent\_ref.record\_fnv1a64}, \code{status = refused}), \code{G1.C2\_status = refused(word\_prime\_cap)}
    and the provenance line \code{G10.rat.settings}, and no other field (the record's JSON envelope also carries the non-field \code{refusal} text, \code{"refused(word\_prime\_cap): \ldots"}): everything computed before the refusal is discarded, no partial value is written; its overall \code{status}
    is \code{refused} although the word is of the pending class (as for a refused descent). The word is not in the schema frozen at \texttt{heckez-schema-v1-frozen}; it
    enters as an amendment after the freeze (\file{docs/implementation/slot25/HZREC\_report.md} \S11, amendment A1): hash-neutral, since the cap is outside the hashed settings
    text ((i)) and no record of the default campaign reaches it (the largest predicted need, $354\,276$ primes at $(1,702)$ (HZ-KIT's prediction, not a measurement), is below $2^{21}$); a scanner pinned to the frozen list
    rejects a record that carries it.
  \item \textbf{G1, the rational structure.} \code{reduced\_Q}, the result of (C1) (\code{deg\_f} $=d$), the primitive $\vartheta$ and the operator set (both functions of
    cell and settings, hence hashed), the components (type, degree, multiplicity on $H^{1}$, the minimal polynomial of $\lambda_i(X_{p_0})$ for the smallest $p_0$ whose
    eigenvalue generates the coefficient field), the candidate set $\mathsf S_{\mathrm{cand}}$ with its reasons, and the discriminant by its digest, bit count, the factorisation of
    its square part $\{q,e,\code{prime\_proof}\}$, the cofactor (exact; as built, v1.25.0-pre.2, Z0 referee F5: of the gcd over the primitive $\vartheta_j$ that (C2) uses, see \S\ref{sec:heckeZ}, \emph{The cofactor test}) with \code{cofactor\_test} $\in\{\code{one},\code{squarefree\_by\_size},\code{unfactored},\code{upper\_bound}\}$ (v1.25.0-pre.1, 2026-11-06, HZ-RAT audit F4:
    \code{upper\_bound} when a cofactor above $2^{64}$ cannot be certified to divide every $D_j$ by the residue route, or a certified prime divides it; it is then only a
    bound and never \code{complete}) and
    \code{C2\_status} $\in\{\code{complete},\code{partial(unfactored)}\}$ (Theorem~\ref{thm:heckeZcert}(C2): $\mathsf S$ is the set of \emph{all} certified $q_i$)
    or, as built (v1.25.0-pre.7, HZ-RAT-PERF audit F6), one of the pending words \code{blocked(U\_p)} (written, until UP-Z's (d$'$)(ii), for every \code{full} record at $N>1$: (C2) for the full algebra needs $U_p$ at the word primes; v1.26.0-pre.3, re-referee W2) and
    \code{refused(word\_prime\_cap)} ((c)); and (v1.26.0-pre.3, re-referee W1, amendment A2) the final \code{not\_run}, written only on the non-reduced branch of (d$'$)(ii), where no $\vartheta$ satisfies (C1) and (C2) is not asked; the closed list of \code{C2} is these five.
    \emph{Reading} (v1.26.0-pre.3, re-referee W1): \code{partial(unfactored)} beside a \code{G1.C1.status} other than \code{certified} certifies nothing --- Theorem~\ref{thm:heckeZcert}(2) presupposes (C1) --- and says only that no factorisation was attempted; this is the reading of the words an \code{anemic} record with (C1) \code{open(cap)} carries (the code path \file{Integral/HeckeZRecord.cpp}, l.~1076, \code{else} branch; v1.26.0-pre.4, confirm-referee W5: no record of edition~1 reaches it --- every \code{anemic} record of edition~1 has (C1) \code{certified}, and its \code{partial(unfactored)} is the genuine (C2) outcome of Theorem~\ref{thm:heckeZcert}(2)), so no byte of any record moves.
    \code{prime\_proof} $\in\{\code{small},\code{pratt},\code{probable}\}$, \code{pratt} meaning that the C++ build has verified the Pratt certificate itself; a factor whose
    certificate was not found or not verified is \code{probable} and counts as part of the cofactor.
    (v1.25.0-pre.4, F5, F10; not yet seen by the owner:) \emph{The components as v1 writes them.} \code{components[0]} is the rational Eisenstein component (\code{type = eis},
    degree $1$, multiplicity $e_k$ [v1.25.0-pre.5, R1; pre.4 said $1$], $p_0$ the smallest prime of $\mathcal P$, minimal polynomial $x-(1+p_0^{k-1})$; the eigenvalue of $\vartheta$ beside it as \code{lambda\_theta}).
    \code{components[1]} (\code{type = cusp}) is the \emph{whole} cuspidal part as one block, of degree $d-1$, with \code{minpoly} $=c_\vartheta$ and \code{irreducible}
    $\in\{\code{true},\code{open(cap)}\}$: \code{true} when the degree is $1$ or $c_\vartheta$ is irreducible of full degree modulo some prime $q\ge3$ (a deterministic
    certificate; the first $64$ [v1.25.0-pre.5: $256$] odd primes are tried in order, a count cap of the settings text, (j); the prime in G9); \code{open(cap)} is final and says that the block
    may be a union of components (v1 has no factoriser over $\Q$). $p_0\in\{\code{T\_p},\code{theta}\}$: \code{T\_p} when the block is irreducible and $\vartheta=T_p$ (the search
    tries the $T_p$ in ascending order and takes the first with $\deg f=d$, so $p$ is the smallest prime whose eigenvalue generates the coefficient field), else \code{theta}
    (the minimal polynomial is that of $\lambda(\vartheta)$). At $d=1$ [v1.25.0-pre.5, R3; pre.4: $d=e_k$] (at level~$1$: $d=e_k$, no cusp forms) \code{components[1]} is not written. \code{type} $\in\{\code{eis},\code{cusp}\}$.
    (v1.25.0-pre.5, HZ-REC re-audit R1, R3, R4, R8; not yet seen by the owner:) \emph{Corrections.} (1) The multiplicity of \code{components[0]} on $H^{1}$ is $e_k$, not $1$: a rational
    Eisenstein system means one anemic Eisenstein system, which all $e_k$ Eisenstein classes carry ($b_1=2\dim S_k+e_k$, and no cuspidal system equals it); $e_k=1$ at level~$1$,
    $e_k$ the number of cusps of $\Ga_0(N)$ ($-1$ at $k=2$) otherwise; e.g.\ $(2,6)$ \code{anemic}, $\chi_L(T_3)=(x-2188)^{2}(x-12)^{2}$, multiplicity $2$. (2) The criterion is
    $d=1$: \code{components[1]} is written iff $d\ge2$ (at level~$1$ this is $d=e_k$, no cusp forms; at $N>1$ \code{anemic}, $d=e_k=2$ can hold with a cusp form, and $d=1$ with
    $e_k=3$). (3) At $N>1$ v1 does not compute the multiplicity of the cuspidal block and writes \code{components[1].mult = not\_run} (final). (4) The irreducibility cap is
    $256$ odd primes (the settings line \code{heckez.irreducible}); all $256$ primes cost $0.19$\,s at degree $58$ ($n\approx700$), and at most $14$\,ms on every cell measured (a count cap, never time-based).
  \item[(d$'$)] \textbf{The full-$N$ reading of the caps text; frozen with v1} (v1.26.0-pre.1, 2026-11-08, slot~26, UP-Z, route~(i);
    \file{docs/implementation/slot26/UPZ\_design.md} \S1.4 and its critical audit \file{docs/implementation/slot26/UPZ\_design\_audit.md}, ACCEPT route~(i) with conditions C1--C7;
    refereed ACCEPT-with-conditions, \file{docs/implementation/slot26/C126pre1\_referee.md}, its conditions and fixes at v1.26.0-pre.2;
    re-refereed ACCEPT-with-conditions, \file{docs/implementation/slot26/C126pre2\_rereferee.md}, its conditions R1, R2 and should-fix W1, W2 at v1.26.0-pre.3;
confirm-refereed ACCEPT, \file{docs/implementation/slot26/C126pre3\_confirm.md}, its should-fix W3, W4, W5 and note N1 at v1.26.0-pre.4;
    aligned with the audited Z1 code (\file{docs/implementation/slot26/Z1\_code\_audit.md}, should-fix S1 (and S2's key)) at v1.26.0-pre.5;
    referee-light-checked ACCEPT-with-conditions, \file{docs/implementation/slot26/C126pre5\_referee\_light.md}, its conditions R1--R3, should-fix W1, W2 and nits N1--N3 at v1.26.0-pre.6;
    confirm-refereed ACCEPT-with-conditions, \file{docs/implementation/slot26/C126pre6\_confirm\_referee.md}, its conditions C1--C3, should-fix W-A, W-B, W-D and nits N-a, N-b at v1.26.0-pre.7;
    confirm-refereed ACCEPT-with-conditions, \file{docs/implementation/slot26/C126pre7\_confirm\_referee.md}, its conditions D1, D2 at v1.26.0-pre.8;
    extended by Z2's G3 words at v1.26.0-pre.9, \file{docs/implementation/slot26/C126pre9\_disposition.md}, refereed with v1.26.0-pre.8 ACCEPT-with-conditions, \file{docs/implementation/slot26/C126pre9\_referee.md}, its condition E1, should-fix W-1 and nits N-1--N-3 at v1.26.0-pre.10;
    confirm-refereed with v1.26.0-pre.10 ACCEPT-with-conditions, \file{docs/implementation/slot26/C126pre10\_confirm\_referee.md}, its condition C1 and should-fix W-A, W-B at v1.26.0-pre.11;
    \textbf{a binding change not yet seen by the owner}). At level $N>1$ with \code{algebra = full} the phrases of the settings text that speak of $\mathcal P$
    --- \code{heckez.P\_rule} (``enlarged \dots\ until (C1) is certified'', the read set ``\code{P u \{T\_ell\}}''), \code{heckez.maximal\_ideals} (``the first of theta\_search's
    order over P'') --- and HZ-RAT's \code{theta\_search} read as follows. At $N=1$ (no prime divides $N$) and for \code{anemic} records nothing changes (but for the reading of the derived word of~(ii), v1.26.0-pre.7, confirm-referee C3, which moves no byte): no hashed byte of a
    level-$1$ or \code{anemic} record moves, and the settings text, its fingerprint \code{0x8ffc29730c4f9672}, every record id and the schema string stay.
    \begin{enumerate}[label=(\roman*)]
    \item \emph{The operator list.} $\vartheta$'s operator list is \code{G1.P}'s $T_p$ in ascending order, then $U_p$ for the primes $p\mid N$ in ascending order --- \textbf{not}
      \code{G0.operators}, which also holds every $T_\ell$ with \code{t\_ell = added} ((j)). Every ``over P'' of the settings text reads ``over the operator list'', and
      ``\code{P u \{T\_ell\}}'' reads ``the operator list $\cup\{T_\ell\}$'' ($T_\ell$ last). (v1.26.0-pre.2, referee E6.) These are exactly three phrases: the read set
      ``\code{P u \{T\_ell\}}'' of \code{heckez.P\_rule}; ``the first of theta\_search's order over P'' of \code{heckez.maximal\_ideals}; and that line's
      ``(theta\_search = P u \{T\_ell\}) \dots\ over the read set P u \{T\_ell\}''. Nothing else is re-read: $g_{\mathcal P}$ of \code{heckez.S\_rule} keeps its meaning (the bound
      over \code{G1.P}; the $U_p$ bounds are \code{S\_rule}'s own parenthesis), and \code{heckez.K\_read} names no list.
    \item \emph{$\mathcal P$ and (C1).} $\mathcal P$ is selected by \code{P\_rule} exactly as for the \code{anemic} record of the same cell: \code{P\_rule}'s ``(C1)'' is the
      \textbf{anemic} (C1) (Theorem~\ref{thm:heckeZcert}(C1) with $d=d^{\mathrm{an}}$), and $U_p$ never enlarges $\mathcal P$. The full (C1) is the degree condition
      $\deg f_\vartheta=d$ over the operator list, $d=\code{G1.d}=\dim M_k(\Ga_0(N))$ ((b)). (v1.26.0-pre.2, referee E5.) $\mathcal P$ is meant in the record's sense
      throughout: \code{G1.P} is the $\mathcal P$ of Theorem~\ref{thm:heckeZcert} for the \emph{anemic} algebra, while that theorem's $\mathcal P$ for the full algebra is the operator
      list of~(i) (its $X_p$, $p\mid N$, are $U_p$ by Definition~\ref{def:heckeZ}).
      \emph{The non-reduced cells} (v1.26.0-pre.2, referee S4). If the record's $U_p$ tag certifies \code{G1.reduced\_Q = non\_reduced\_certified} (the counting criterion of
      \S\ref{sec:heckeZ}, ``Two forms, and what is reachable today'', asserted for $N$ a power of $p$ only; e.g.\ $(8,10)$, where $U_2$ has a nilpotent block), then $\Talg_\Q$ is not
      reduced, so $\deg f_\vartheta<d$ for \emph{every} $\vartheta$ (Theorem~\ref{thm:heckeZcert}(1): (C1) implies that $\Talg_\Q$ is reduced) and the full (C1) is false. The
      record then runs \textbf{no} $\vartheta$ search (no \code{theta\_cap} is spent) and writes (v1.26.0-pre.3, re-referee W1, R1(a): the accurate words, fixed before any record
      carries them): \code{G1.C1.status = refused(non\_reduced)}, \textbf{final} ((C1) is refuted, not undecided: it can never become \code{certified}, for any $\vartheta$, cap,
      presentation, $\ell$ or $U_p$ convention; the hashed \code{G1.reduced\_Q} beside it is the reason), with no \code{G1.C1.deg\_f}; \code{G1.theta = not\_run} (final; no
      search ran); \code{G1.components = not\_run}; no \code{G1.disc.*} and no \code{G1.discriminants\_used}; \code{G1.C2\_status = not\_run} (final; (C2) presupposes (C1), so
      nothing is certified and nothing is asked); \code{G1.theta\_ops} as in~(vi); \code{G3.partition = U\_p\_tags} with the $U_p$-tag fields of \code{G3.C\_T}; no
      $\vartheta$-derived reason in $\mathsf S_{\mathrm{cand}}$; and the G2 words explicitly: for every $\ell\in\mathsf S$ with $\ell\ge5$,
      \code{G2.ell=<ell>.generation.status = open(sandwich)} (final, no read runs, as for a record whose (C1) is not certified), and for $\ell\in\mathsf S\cap\{2,3\}$
      \code{G2.ell=<ell>.generation.status = blocked(U\_p)} (pending, by~(v)) until Z2 lands, Z2 then writing \code{open(sandwich)} there (or, if Z2-D's negative is recorded and
      audited, \code{refused(no\_lift)} by~(v)). The words \code{refused(non\_reduced)} in \code{C1} and \code{not\_run} in \code{C2} enter the frozen schema as an amendment after the
      freeze (A2, of the kind of \file{docs/implementation/slot25/HZREC\_report.md} \S11, A1): hash-neutral, since no record of edition~1 carries them, the settings text and its
      fingerprint do not move, and every \code{full} record at $N>1$ that will carry them is pending in edition~1 and moves anyway ((ix)); \code{not\_run} in \code{G1.theta} is
      already a word of that field. \emph{The verifier} (\file{docs/integral\_census/hz\_dataset\_verify.py}, its \code{FINAL} set, l.~54--56, copied from \code{StatusTable()})
      gains \code{refused(non\_reduced)} as \textbf{final}, in the same commit as Z1's \code{StatusTable()} row and the closed lists \code{C1} (\code{+refused(non\_reduced)}) and
      \code{C2} (\code{+not\_run}); no word changes class, and \code{PENDING} is unchanged.
      \emph{Otherwise} the search runs as for the \code{anemic} record: \code{G1.C1.status} $\in\{\code{certified},\code{open(cap)}\}$, \code{open(cap)} meaning
      \code{theta\_cap} exhausted, final; \code{G1.reduced\_Q} is \code{true} when (C1) is certified, else \code{true} when the $U_p$ tag certifies it (split degree), else
      \code{undecided} (final). \textbf{These final words only where the search ran} (v1.26.0-pre.3, re-referee R2): \code{G1.C1.status} $\in\{\code{certified},\code{open(cap)}\}$,
      and the final values of \code{G1.reduced\_Q}, \code{G1.theta}, \code{G1.components} and \code{G1.C2\_status} beside it, are written only when the $\vartheta$ search over the operator list ran with
      $U_p$ read at the word primes by~(v). Where Z1's word-prime $U_p$ route did not run --- a predicted RSS above the ceiling, a refused Heilbronn build or another refused route,
      a cell outside Z1's coverage --- these five $\ell$-independent fields (v1.26.0-pre.4, confirm-referee W4: \code{G1.components} is the fifth, $\vartheta$-dependent as the others and \code{blocked(U\_p)} on every \code{full} record at $N>1$ of edition~1, never the final \code{not\_run} of the non-reduced branch here) keep \code{blocked(U\_p)}, or \code{blocked(export)} for the RSS ceiling (pending, (c)); never
      \code{open(cap)}, which would be false-final (a refusal by \code{word\_prime\_cap} refuses the whole record, (c)). Subject to that, \code{blocked(U\_p)} is no longer written in
      \code{G1.C1.status}, \code{G1.reduced\_Q}, \code{G1.theta} or \code{G1.C2\_status} (the other four $\ell$-independent fields).
      \emph{The components and the maximal ideals where the search ran} (v1.26.0-pre.5, Z1 code audit S1(a), \file{docs/implementation/slot26/Z1\_code\_audit.md}; the text
      aligned with the audited code, replacing pre.4's inclusion of \code{G1.components} in the list above): where the search ran and (C1) is \code{certified},
      \code{G1.components} $=$ \code{blocked(U\_p)}, and at every $\ell$ of G2 with \code{generation.status} $=$ \code{certified}, \code{G2.ell=<}$\ell$\code{>.maximal\_ideals}
      $=$ \code{blocked(U\_p)} --- both \textbf{pending} ((c)): the components and the decomposition of~(d) need the Eisenstein $U_p$ systems, which Z1 does not build; they
      stay \code{blocked(U\_p)} until the package that builds them writes their values, a move pending $\to$ value, never final $\to$ final. (v1.26.0-pre.6, referee-light W1; restated at v1.26.0-pre.7, confirm-referee W-D:) that
      package is \textbf{not} Z2, whose card in \file{docs/implementation/slot26/00-slot26-plan.md} (v0.6) is $U_2$, $U_3$ at $\ell\in\{2,3\}$; it is \code{HZ-EISUP}, the card of that
      plan (\S2.2, sequenced right after Z2 by the owner's ruling of 2026-11-08) that builds the Eisenstein $U_p$ systems. Z2 is therefore necessary and not sufficient: a record carrying one of these
      fields stays \code{pending} after Z2 until \code{HZ-EISUP} writes it.
      Where the search ran and (C1) is \code{open(cap)}, \code{G1.components} $=$ \code{not\_run} (final; as for the \code{anemic} record, no component is computed without (C1);
      v1.26.0-pre.6, referee-light W2), and where \code{generation.status} is not \code{certified},
      \code{maximal\_ideals} $=$ \code{not\_run} by the rule of~(d) (\emph{The maximal ideals as v1 writes them}); v1.26.0-pre.5 introduces no final word.
      \emph{The derived word} (v1.26.0-pre.6, referee-light R1; v1.26.0-pre.7, confirm-referee C3: the scope ``\code{full} at $N>1$'' dropped). On any record, \code{maximal\_ideals} $=$ \code{not\_run} beside a \code{generation.status} of the
      pending class is a \emph{derived} word: it records that the decomposition was not attempted because generation is pending, it asserts nothing final, and it is rewritten when
      \code{generation.status} moves. The rule ``stays \code{blocked(U\_p)}'' of~(v) does not apply to this field. \code{not\_run} beside a final \code{generation.status}
      (\code{open(sandwich)}, \code{open(cap)}, \dots) remains final. The word stays of the final class in~(c) and in \code{StatusTable()}, so no byte and no record status moves: the
      pending \code{generation.status} beside it makes the record \code{pending} by~(c).
      This paragraph alone of~(d$'$) reaches beyond \code{full} records at $N>1$ (the word of~(d) stands also beside \code{refused(export\_word)} on an \code{anemic} record and beside
      \code{refused(k1\_route)} on a \code{full} record at $N=1$): only the reading of the word changes, and no byte of a level-$1$ or \code{anemic} record moves.
      Every \code{full} record at $N>1$ of this paragraph is pending in any case, by the pending \code{generation.status} beside the word, and before Z2 also through \code{G3.C\_T.lower} (the rule below).
      \emph{No full record at $N>1$ is final before Z2} (v1.26.0-pre.3, re-referee R1(b); the forecast F10 of \file{docs/implementation/slot26/UPZ\_design\_audit.md} made a rule):
      on \textbf{every} \code{full} record at $N>1$, whatever its $\mathsf S_{\mathrm{cand}}$, until Z2 lands (in the sense of \emph{The G3 words from Z2 on}, below; v1.26.0-pre.10, referee N-1) or Z2-D's negative is recorded in
      \file{docs/implementation/slot26/UPZ\_design.md} and upheld by its audit, the aggregate fields of G3 are written as follows (v1.26.0-pre.4, confirm-referee W3: the list
      replaces pre.3's open-ended ``every other aggregate field whose value may need $U_p$ at $\ell\in\{2,3\}$''; it is exhaustive for this rule): \code{G3.C\_T.lower} $=$
      \code{blocked(U\_p)} (pending); \code{G3.C\_T.status} $=$ \code{upper\_bound\_certified} where every $U_p$ tag of the record certifies a bound, else \code{blocked(U\_p)}
      --- \code{upper\_bound\_certified} is a bound's status, of the final class, and \textbf{is} the planned move of~(ix): it may be sharpened to \code{certified} only together with
      \code{G3.C\_T.lower} (lower $=$ upper at every $\ell$, item~(d), G3), never alone (no package of this version does so: \emph{The G3 words from Z2 on}, below); and \code{G3.C\_T.exact} is absent. So every such record stays \code{pending} ((c))
      through \code{G3.C\_T.lower} and none is golden-eligible before then (nor, after Z2 alone, any record carrying one of the \code{blocked(U\_p)} fields of \emph{The components and the maximal ideals where the
      search ran} above; v1.26.0-pre.6, referee-light W1; v1.26.0-pre.7, confirm-referee C1). Those are the records with \code{G1.C1.status} $=$ \code{certified} ($67$ of the $85$
      records of Z1's recompute). The other $18$ ($9$ with \code{open(cap)}, $9$ with \code{refused(non\_reduced)}) carry no field waiting on \code{HZ-EISUP}: their pending fields are
      \code{generation.status} at $\ell\in\{2,3\}$ and G3's \code{C\_T} fields, which the words of \emph{The G3 words from Z2 on} below make final. The per-$\ell$ words at $\ell\in\mathsf S\cap\{2,3\}$ (G2's \code{generation.status}) are governed
      by~(v) and by the non-reduced branch above, not by this rule. No level-$1$ or \code{anemic} byte moves.
      \emph{G3 on the non-reduced branch} (v1.26.0-pre.4, confirm-referee N1). There \code{G3.C\_T.status} is \code{blocked(U\_p)}, never \code{upper\_bound\_certified}: as built
      (\file{Integral/HeckeZUTag.cpp}) a $U_p$ tag certifies a bound only when $\chi_\Q(U_p)$ splits with $m_0=0$ and \code{distinct} $=d$, i.e.\ (C1) for $\vartheta=U_p$, which
      Theorem~\ref{thm:heckeZcert}(1) excludes when $\Talg_\Q$ is not reduced (e.g.\ the pin $(8,10,\code{full})$ reads \code{C\_T=blocked(U\_p)}); and \code{G3.C\_T.lower} is
      \code{blocked(U\_p)} by the rule above. These are the words before Z2; the final words of these two fields on this branch are those of \emph{The G3 words from Z2 on} (v1.26.0-pre.9).
      \emph{The G3 words from Z2 on} (v1.26.0-pre.9, UP-Z~(b), contract first; \file{docs/implementation/slot26/C126pre9\_disposition.md}; the one place of these words). Z2 has landed for a record iff the
      binary that wrote it carries the route of~(v) at $\ell\in\{2,3\}$, which the record states by \code{G10.contract} ((vii)); a binary carrying that route writes the words of this paragraph, so by~(vii) its records state \code{1.26.0-pre.9} or a later version that moves a word it writes (v1.26.0-pre.10, referee N-3). On every \code{full} record at $N>1$ of such a binary on which the $U_p$ tag step ran at every $p\mid N$ (v1.26.0-pre.10, referee W-1; as built this is every record such a binary writes), on either branch of~(ii), the
      words of \code{G3} are: \code{G3.C\_T[U\_p].status} $=$ \code{upper\_bound\_certified} where that $U_p$ tag certifies its bound, else \code{not\_run}; the rows
      \code{G3.C\_T[U\_p:a,b].status} and \code{.upper} exist only where the tag certifies its bound, with status \code{upper\_bound\_certified}; \code{G3.C\_T.status} $=$
      \code{upper\_bound\_certified} where every $U_p$ tag of the record certifies its bound, else \code{not\_run}; \code{G3.C\_T.lower} $=$ \code{not\_run}; \code{G3.C\_T.exact} absent. All are
      \textbf{final} ((c)). Where the tag step did not run at some $p\mid N$ (no record of this version: a refusal by \code{word\_prime\_cap} refuses the whole record, (c)), \code{G3.C\_T[U\_p].status} of that $p$, \code{G3.C\_T.status} and \code{G3.C\_T.lower} keep \code{blocked(U\_p)} (pending, (c)), as the five fields of~(ii) do where the search did not run; never \code{not\_run}, which would be false-final. (\code{blocked(U\_p)} also at an RSS ceiling: it is the only pending word of the closed list \code{utag}.) \code{not\_run} in the two status fields says that this route ran and certifies no bound, nothing about $C_{\Talg}$ itself. \emph{Why \code{lower} is \code{not\_run}} (v1.26.0-pre.10, referee E1, N-2): the rows of \code{G3.partition}
      $=$ \code{U\_p\_tags} are pairs of the $U_p$-eigen-systems. For $d\ge3$ a pair is not all components, and $\eta_L$, with it the \code{lower} $=\exp\eta_L$ of~G3, is defined by
      Definition~\ref{def:heckeZcong} only when $\mathcal S\sqcup\mathcal S'$ is all components, the hypothesis of every clause of Proposition~\ref{prop:heckeZcong}. For $d=2$ the pair of a certifying tag is all components and $\eta_L$ is defined (e.g.\ $(N,n)=(2,4)$, the pair \code{eis(1),eis(32)}, bound $31$); its read is Proposition~\ref{prop:heckeZcong}(5) with $X_p=U_p$, which no route of this version performs on a \code{U\_p\_tags} row and for which the partition \code{U\_p\_tags} has no field (v1.26.0-pre.11, confirm-referee W-B). On the non-reduced branch that
      Definition does not apply at all (Remark~\ref{rmk:heckeZcaveat}: it assumes $\Talg_\Q$ reduced), so the absence of a $\vartheta$ is not the reason and the question left open at v1.26.0-pre.4 is answered by that scope. So no
      statement of this document prescribes a read of a lower bound on these rows, \code{G3.C\_T.lower} is not a derived word and does not wait on \code{generation.status}, and on such a record no field is sharpened to
      \code{certified}: the rule ``only together with \code{G3.C\_T.lower}'' above is not exercised, and the reading of the aggregate in~(c) does not apply, since \code{G3.C\_T.lower} is not pending. Where no tag certifies there is no row ($77$ of Z1's $85$ records, a count at this version; v1.26.0-pre.11, confirm-referee W-A), and the partition \code{U\_p\_tags} carries only the fields of this paragraph (item~(d), G3), so \code{G3.C\_T.lower} is \code{not\_run} there too: on the non-reduced branch and on the \code{open(cap)} records, which are outside the scope of Definition~\ref{def:heckeZcong} ((C1) does not pass), and also where $\eta_L$ is defined and Proposition~\ref{prop:heckeZcong}(5) could read it with a $T_p$: e.g.\ $(N,n)=(11,0)$, $d=2$, one Eisenstein system, no certifying tag, whose \code{anemic} record reads \code{G3.C\_T.exact} $=$ \code{0x5}. A \code{full} record of this version does not perform that read. A later value of these fields is a
      settings step or a schema step ((viii)), never a package filling a pending field: the read of Proposition~\ref{prop:heckeZcong}(5) for one rational Eisenstein system against the rest (with $X_p=T_p$ where a $T_p$ separates, $U_p$ where only a $U_p$ does; v1.26.0-pre.11, confirm-referee W-A) would be a new
      partition row with its own statement, and \code{HZ-EISUP}, which writes \code{G1.components} and \code{maximal\_ideals}, does not move these fields. At $\ell\in\mathsf S\cap\{2,3\}$ on a record whose \code{G1.C1.status} is
      \code{open(cap)}, \code{G2.ell=<ell>.generation.status} $=$ \code{open(sandwich)} (final; no read without (C1), as at $\ell\ge5$ and on the non-reduced branch), \code{maximal\_ideals} then being the final \code{not\_run} of~(d).
      \emph{Amendment A3}: \code{not\_run} joins the closed list \code{utag} of the $U_p$ tag fields, after the freeze, of the kind of A1 and A2: hash-neutral, since that list is not hashed, no record of edition~1 carries
      \code{not\_run} in such a field, every \code{full} record at $N>1$ is pending in edition~1 and moves anyway, and \code{not\_run} is already in the verifier's \code{FINAL} set (no class moves). The rules of this item for a binary without that route stay
      in force for what such a binary writes, so no record it wrote is invalidated; such a record stays \code{pending} until regenerated. On Z1's recompute of the $85$ records (a count at this version, not a rule) these words leave
      the $18$ records of~(ii) (\code{open(cap)} and \code{refused(non\_reduced)}) with no pending field, hence \code{final}, and the $67$ with (C1) \code{certified} pending through \code{HZ-EISUP}'s fields, $3$ of them also through
      \code{refused(export\_word)} at $\ell\ge5$.
    \item \emph{The $\vartheta$ search.} HZ-RAT's order on the operator list: the units first, in list order, then the combinations, within \code{theta\_cap} and
      \code{theta\_coeff\_max} $=\binom d2$ with $d=\code{G1.d}$ of the record --- the full $d$, not $d^{\mathrm{an}}$.
    \item \emph{Which $U_p$.} $U_p=[A\,\smat100p\,A]$, $A=\Ga_0(N)$: the $U^{(N)}_p$ of Proposition~\ref{prop:levelcorrespondence}(4)(d) and of Proposition~\ref{prop:strictliftN}(9), on
      $p$-old spaces the \textbf{adjoint} $U_p^{\ast}$ of the classical operator (Proposition~\ref{prop:levelconehecke}(3)(d)). The convention is named here and in the header of every
      $U_p$ certificate dump; the $\Z_{(\ell)}$-algebras generated with $U_p$ and with $U_p^{\ast}$ are not asserted to be isomorphic at $\ell\mid N$, so a comparison with another
      convention (the literature, the Lean side) names it, citing Remark~\ref{rmk:stabconvention}(1), (2) ($U^{\ast}$ and $U^{(\ell)}_\ell$ are $W_\ell$-conjugate and not
      interchangeable over $\Z_\ell$; Hida's $T(p)$ is $U^{\ast}$) and [U]'s \emph{Remark on the orientation} (\file{docs/latex/hida\_hecke\_algebra.tex}, after
      \texttt{prop:u-towerorient}). (v1.26.0-pre.2, referee S3.) [U]'s ``[C]'s $U_\ell$ is the classical $U$'' is a statement about Hecke \emph{systems} on the $\chi$-parts of the
      $\Ga_1(\ell^{r})$ model at tame level~$1$: it fixes the dictionary $U\mapsto a_\ell$, $\langle d\rangle\mapsto\chi(d)$, an orientation that remark decides by computation through
      the coupled swap $U\leftrightarrow U^{\ast}$, $\langle d\rangle\leftrightarrow\langle d\rangle^{-1}$. Proposition~\ref{prop:levelconehecke}(3)(d) is a statement about the
      operator against the inclusions $\mathrm{dg}_1$, $\mathrm{dg}_p$ of old classes at $\Ga_0(N)$, $\chi=1$. The two concern different objects and do not contradict each other;
      neither is read by this record, which carries no character and no $\Ga_1(\ell^{r})$ level and names its operator by the double coset $[A\smat100pA]$. That they are one
      operator-level convention on a common model is not asserted.
    \item \emph{Where $U_p$ is read.} At $\ell\ge5$ and at the word primes $q\nmid6N$: the operator $H_{N,p}$ on the cokernel presentation (Proposition~\ref{prop:heilbronn}(1), (2),
      stated for $p\mid N$). For the free block over $\Z$ and every field at $\ell\in\{2,3\}$: $\widetilde U^{(N)}_{\Iso,p}$ of Proposition~\ref{prop:strictliftN}(9) and its
      transports. Until that route is built, every $U_p$-dependent field at $\ell\in\{2,3\}$ stays \code{blocked(U\_p)} (pending, (c)), except \code{maximal\_ideals}, which beside a pending \code{generation.status} carries the derived \code{not\_run} of~(ii) (v1.26.0-pre.6, referee-light R1), and after it is built the fields of~G3 follow \emph{The G3 words from Z2 on} of~(ii) (v1.26.0-pre.9); \code{refused(no\_lift)} is written only for a
      per-cell certificate refusal ((c)), never on the ground that a lift at $p\mid N$ is missing (Proposition~\ref{prop:strictliftN}(9) constructs one); a prediction above the RSS
      ceiling is \code{blocked(export)}. The digests of the $U_p$ matrices are certificate keys \code{G9.Up.p=<p>.mat\_sha256} (v1.26.0-pre.5, Z1 code audit S1(b); v1.26.0-pre.6, referee-light R2): unhashed
      (G9, outside \code{record\_fnv1a64}; inside \code{cert\_sha256} as every G9 key), one key per $p\mid N$, present only on \code{full} records at $N>1$ whose \emph{record build}
      serves the $U_p$ matrix over $\Z$ (its $\tau$ and $M$ included), and written by the package that serves it there; \textbf{no record of Z1 carries it}. Running the $U_p$ dump
      (\code{HECKE\_DUMP\_HZ\_CERT}, \code{WriteUpMatrixFiles}) does not create the key, so \code{cert\_sha256} does not depend on whether a dump ran. That package also adds the key's
      schema row (the table has none at this version) and fixes which digest the key means. The key is absent elsewhere and carries no status word; no level-$1$ or
      \code{anemic} record carries one, so neither of their digests moves.
    \item \emph{Self-description (required).} Every \code{full} record at $N>1$ built under this reading carries the hashed G1 field \code{G1.theta\_ops} (free text, with its own
      schema row): the operator list of~(i), in order, comma-separated (e.g.\ \code{T\_2,T\_3,T\_5,U\_2}), by which the coefficient vector \code{G1.theta} is indexed. Level-$1$ and
      \code{anemic} records do not carry it (their index list is \code{G1.P}), so no byte of theirs moves.
    \item \emph{Provenance (required).} Every \code{full} record at $N>1$ names in G10 (unhashed) the version of this document whose (d$'$) it implements, as
      \code{G10.contract} $=$ \code{gl2\_contract <version>, item 8 (d')}. (v1.26.0-pre.5, Z1 code audit S2:) the key is exactly \code{G10.contract} (not \code{G10.full.contract}); G10
      is outside \code{record\_fnv1a64} and \code{cert\_sha256}, so its text moves no digest. (v1.26.0-pre.6, referee-light R3:) \code{<version>} is the last version of this document,
      among those the binary that wrote the record implements, whose change to (d$'$) moved a word that binary writes, that is (v1.26.0-pre.7, confirm-referee C2), changed, for some record
      and some written key or field other than \code{G10.contract} itself, the word or value that (d$'$) prescribes, or the key's presence; supplying a reason, a class or a pointer, or
      resolving a conflict between two sentences of (d$'$) in favour of a word one of them already prescribed, moves no written word; naming for the first time the word of a field that (d$'$) had left open moves a written word for a binary that writes it (v1.26.0-pre.8, confirm-referee D2). A \emph{written} key or field is one that binary
      writes, so a change to the prescription of a key no binary has written (v1.26.0-pre.6's for \code{G9.Up.*}) moves no written word. It makes no difference whether the code was
      conformed to the text or, as at v1.26.0-pre.5, the text to the code; a later version that moves no written word, and a signature as such, does not change it. At this version
      \code{<version>} is \code{1.26.0-pre.5} for a binary that does not write the words of \emph{The G3 words from Z2 on} in~(ii) (none of v1.26.0-pre.6, v1.26.0-pre.7 and v1.26.0-pre.8 moves a written word by this test, v1.26.0-pre.9 prescribes for that binary what it wrote before, and neither v1.26.0-pre.10 nor v1.26.0-pre.11 moves one), so its value is
      \code{gl2\_contract 1.26.0-pre.5, item 8 (d')}; it is \code{1.26.0-pre.9} for a binary that writes those words (v1.26.0-pre.9 names words (d$'$) had left open and changes others, so by this test it moves a written word for that binary; v1.26.0-pre.10 supplies a reason and pointers and adds a proviso for records no binary writes, and v1.26.0-pre.11 supplies status pointers and a reason, so neither moves one), so its value is \code{gl2\_contract 1.26.0-pre.9, item 8 (d')}. A written value of this key is a historical fact and is never rewritten, at a signature or otherwise. The manifest of the integral Hecke dataset edition~2 states both the version of this document under which it is issued (1.26.0 once
      signed) and the \code{G10.contract} value(s) its \code{full} records carry, and asserts, by the test of this item, that no version after each such value up to the issuing version
      moved a written word of a record carrying that value (v1.26.0-pre.8, confirm-referee D2); if one did, the affected records are regenerated, never edited.
    \item \emph{Frozen with v1.} This reading is part of \code{integral\_hecke\_structure/v1}. Any later change of it --- another order of the list, $U_p^{\ast}$ for $U_p$, another
      list, another (C1) for $\mathcal P$ --- is a settings change (a new caps line, hence a new fingerprint and new record ids) or a schema step to
      \code{integral\_hecke\_structure/v2}; \textbf{never an edit of this item in place}.
    \item \emph{Why this is not a re-baseline.} The \code{blocked(U\_p)} fields are the move the schema planned ((c): ``a named package moves them''); every \code{full} record at
      $N>1$ of edition~1 carries \code{blocked(U\_p)} in exactly the status fields this reading fills (v1.26.0-pre.7, confirm-referee W-B) --- apart from \code{G2.ell=<}$\ell$\code{>.maximal\_ideals}, which edition~1 writes
      \code{not\_run} beside \code{generation.status} $=$ \code{blocked(U\_p)} (v1.26.0-pre.6, referee-light R1): that is the derived word of~(ii), which asserts nothing final, and its
      move \code{not\_run} $\to$ \code{blocked(U\_p)} where \code{generation.status} becomes \code{certified} is the rewriting (ii) prescribes, not a move from a final word --- so no
      final value under another reading exists; and the index list it refers to is
      hashed (\code{G0.operators}, \code{G1.P} and, by (vi), \code{G1.theta\_ops}). Beside the status fields, non-status hashed values of these records also move between
      edition~1 and Z1: \code{G0.operators} (the enlargement of~(i), on $64$ of the $85$ records) and \code{G1.S\_cand.ell=<}$\ell$\code{>.reason} and \code{.v\_disc} ($123$ fields each, the
      discriminant reasons now available). None is a status word and every record concerned was \code{pending}, so the conclusion stands.
    \end{enumerate}
  \item \textbf{G2, per prime $\ell\in\mathsf S$ and per maximal ideal} (sorted by residue degree and the monic residue factor of $\chi(X_{p_0})\bmod\ell$): \code{generation}
    with its \code{status} and \code{route} $\in\{\code{saturation},\code{eis\_single},\code{rational},\code{cyclic},\code{commutant},\code{AK\_unsaturated},\code{mazur}\}$
    (Propositions~\ref{prop:heckeZgen}, \ref{prop:heckeZeiscert}, Theorem~\ref{thm:heckeZcert}, Propositions~\ref{prop:heckeZrat}(5)(c), \ref{prop:heckeZmono}; \code{mazur} is local at the
    Eisenstein ideal of prime level and carries the good prime $q$ it used together with $v_\ell(\mathsf F_q(1+q))$ and $v_\ell(\nu)$; it is \code{certified} \textbf{only when
    the two valuations agree}, the check of Proposition~\ref{prop:heckeZeiscert}(3) --- goodness alone does not exclude a non-Eisenstein block in the Hensel factor
    (v1.25.0-pre.0b, A7-M1)); the maximal ideals with their residue data and \code{support} $\in\{\code{free},\code{torsion},\code{both}\}$
    (at $\ell\in\{2,3\}$ the degree attribution of Proposition~\ref{prop:heckeZuct}(4)); \code{components\_through} (Proposition~\ref{prop:heckeZthrough}). \textbf{Local types
    only under certified generation}: \code{local\_type} (\code{kind} $\in\{\code{etale},\code{monogenic},\code{non\_monogenic}\}$, rank, embedding dimension,
    \code{gorenstein} $\in\{\code{true},\code{false}\}$ from the socle of $\Talg_\midl/\ell$, length, Hilbert function, Loewy length) and \code{pair\_kinds} (depth and kind,
    \code{certified} for a generating $\vartheta$, else \code{upper\_bound\_certified}) are written only when \code{generation.status = certified}; otherwise the mod-$\ell$
    invariants of $A'_1$ appear labelled as a statement about the computed subalgebra (Remark~\ref{rmk:heckeZcaveat}(4)). \textbf{Generation is never \code{certified} from a read
    of $H^{1}(\Ga;\Vcoef/\ell^{K})$ at a torsion prime of the cell} (\S\ref{sec:heckeZ}, ``Two forms, and what is reachable today'').
    (v1.25.0-pre.4, F6, F12; not yet seen by the owner:) \emph{The maximal ideals as v1 writes them}: decomposed (\code{maximal\_ideals = certified}, by the factorisation of
    $f_\vartheta\bmod\ell$) iff \code{generation} is \code{certified} by \code{saturation} or \code{rational} and some $\vartheta$ of \code{theta\_search}'s order over $\mathcal P$
    (at most \code{theta\_cap} candidates) has $\bar\vartheta$ generating $A'_1$ ($1,\bar\vartheta,\dots,\bar\vartheta^{d-1}$ independent in $\End(L/\ell)$, $\dim A'_1=d$; then
    $\Talg\otimes\Z_\ell=\Z_\ell[\vartheta]$ by (C3) and Nakayama, and $f_\vartheta\bmod\ell$ is the Krylov relation of $\bar\vartheta^{d}$); otherwise, and under \code{eis\_single},
    \code{maximal\_ideals = not\_run} (final, except for the derived word of item~(d$'$)(ii) beside a pending \code{generation.status}, v1.26.0-pre.6, on any record since v1.26.0-pre.7, confirm-referee C3; the rule is in the settings text, so a wider rule is a new fingerprint). \code{support} of a maximal ideal is \code{both} when a
    common torsion point of (g) has the same $\vartheta$-residue, \code{not\_run} when that residue is not determined (a non-unit $\vartheta$ with a residue of degree $>1$).
    v1 does not compute the degree attribution at $\ell\in\{2,3\}$ and writes \code{degree\_attr = not\_run} (final).
    (v1.25.0-pre.5, HZ-REC re-audit R4, R9; not yet seen by the owner:) \emph{The $\vartheta$ search} is widened to the read set: the read's own generating $\vartheta$, else the first of
    \code{theta\_search}'s order over $\mathcal P$, else --- when \code{t\_ell = added}, $\ell\le23170$, (j) --- the first of \code{theta\_search}'s order over the read set
    $\mathcal P\cup\{T_\ell\}$ ($T_\ell$ last; its Eisenstein eigenvalue is $1+\ell^{k-1}\equiv1$); where generation is certified by saturation only with $T_\ell$ (no $\vartheta$ over
    $\mathcal P$ generates $A'_1$, e.g.\ $(1,26)$ at $7$) the maximal ideals are then decomposed (there by $\vartheta=T_5+T_7$), and \code{decomposed\_by} names the operators
    (\code{theta=(0 0 1 1) on T\_2,T\_3,T\_5,T\_7}). \emph{\code{components\_through.<i>}}: \code{certified} when the component passes through $\midl$, and \code{not\_run} when it
    is \emph{certified not to pass} ($\pi\nmid f_i\bmod\ell$ at the generating $\vartheta$, Proposition~\ref{prop:heckeZthrough}); under \code{maximal\_ideals = certified} the
    word \code{not\_run} here therefore means ``excluded, certified'', not ``unknown'' (v1 has no separate word for it). \emph{\code{pair\_kinds[eis,cusp].kind}} is
    \code{meeting} when the cuspidal block has degree~$1$ and \code{not\_run} (final) when it has degree $\ge2$: v1 does not decide the kind of a pair with a block that may be
    a union of components.
  \item \textbf{G3, congruence modules: exact, upper, lower.} Per partition treated (all Eisenstein components against the rest; each pair of components meeting somewhere):
    \code{C\_T} with \code{upper} $=g_{\mathcal P}$ (\code{upper\_bound\_certified}), written as the word-size factorisation of $g_{\mathcal P}$ --- each $v_\ell(g_{\mathcal P})$ by residues
    modulo $\ell^{K}$ --- plus a cofactor (as built, v1.25.0-pre.2, F5: that of the gcd over the primitive unit $\vartheta=T_p$, the smallest base cofactor certified by residues to divide every
    $c_{\mathcal S',p}(\lambda(X_p))$, else \code{cofactor\_test = upper\_bound} as in (d)), with the $\ell$-wise form of Proposition~\ref{prop:heckeZcong}(4) (no big gcd, Q1(a); A7-m4), \code{lower} $=\exp\eta_L$ per $\ell\mid g_{\mathcal P}$, and \code{exact}
    \textbf{only} when lower $=$ upper at every $\ell\mid g_{\mathcal P}$ (\code{certified}); \code{partial(unfactored)} when $g_{\mathcal P}$ keeps a cofactor; an
    $\ell$-part whose at-root read is \code{refused(word\_modulus)} stays \code{upper\_bound\_certified}; \code{eta\_L} (the Smith exponents of the at-root read, $r$, the
    $\inv$-sign); \code{fibre\_product} $\in\{\code{certified},\code{not\_certified}\}$ (Propositions~\ref{prop:heckeZcong}(4), \ref{prop:heckeZeiscert}(1));
    \code{second\_route} --- the Bernoulli valuation $v_\ell(B_k/k)$ and whether it agrees (level $1$), Mazur's $v_\ell(\mathrm{num}((N-1)/12))$ (prime level, weight $2$) ---
    which never lifts a status (\code{agrees = false} is a defect, not a status); \code{general\_case} $\in\{\code{bounds},\code{from\_certified\_T}\}$ for non-rational Eisenstein
    components and cusp--cusp partitions.
    (v1.25.0-pre.4, F5, F10, O5, O6; not yet seen by the owner:) \code{partition} $\in\{\code{eisenstein}$ (\code{S} the Eisenstein component, \code{S\_prime} the cuspidal
    block)$,\code{U\_p\_tags}$ (full algebra at $N>1$: the bounds of the $U_p$ tags; its fields are those of (d$'$)(ii))$,\code{general}$ (no rational Eisenstein system)$,\code{none(C1\_open)},\code{none(single\_component)}\}$;
    at $d=1$ [v1.25.0-pre.5; pre.4: $d=e_k$] the partition is \code{none(single\_component)} and G3 carries no congruence module. \code{second\_route.agrees} $\in\{\code{true},\code{false},\code{not\_read}\}$,
    \code{not\_read} when the read it would be compared with has not closed. \code{eta\_L.sign} is \code{+} at level~$1$ and \code{both} at $N>1$ (the read is on $L$, not split by $\inv$).
    (v1.25.0-pre.5, HZ-REC re-audit R2, R3; not yet seen by the owner:) The partition is \code{none(single\_component)} iff $d=1$ (at level~$1$: $d=e_k$), (d). \emph{The second route has word caps},
    in the settings text (\code{heckez.second\_route}, (j)): $v_\ell(B_k/k)$ is read modulo $\ell^{e+2}$, $e=v_\ell(g_{\mathcal P})$, at level~$1$, $\ell\ge5$, $(\ell-1)\nmid k$, by
    the tangent-number route of \file{HeckeExt/HeckeHidaLambda}, and is served iff $\ell^{e+5}<2^{128}$, $\ell^{e+2}+\ell\le2^{84}$ and $\ell^{e+2+V}+\ell\le2^{110}$
    ($V=v_\ell(4^{k/2}-1)$) and the read gives $v<e+2$; otherwise \code{v\_Bk = not\_run} (final). \code{agrees = not\_read} stands where the at-root read has not closed
    \emph{or} \code{v\_Bk = not\_run}; e.g.\ $(1,28)$ at $\ell=1001259881$ and $(1,30)$ at $305065927$ ($e=1$, $\ell^{6}\ge2^{128}$), whose at-root reads are \code{certified}.
  \item \textbf{G4, the glue $J$: \code{zero} or \code{bounds}, never exact.} Per $(\ell,\midl)$ with \code{support = both}: \code{zero} (separated supports, or upper bound of
    order $1$), \code{bounds} with \code{lower} $=\log_\ell|J\cap\Talg'|$ and \code{upper} $=\log_\ell|\Hom_{\Talg'}(L/\ell^{e_\ell},\mathrm{tors}_\ell)|$, or \code{blocked(export)}
    (Proposition~\ref{prop:heckeZglue}(3)); an exact value would need Remark~\ref{rmk:heckeZcbthree}, which no package proves.
    (v1.25.0-pre.4, F3, F12; not yet seen by the owner:) \emph{The keys as v1 writes them}: the \emph{torsion points} at $\ell$ are the joint $\Talg'$-blocks of $H^{1}/\ell^{e}$
    ($\Talg'=\Z[T_p:p$ in the read set$]$, $e=\sigma_{\max}$) with non-zero torsion, keyed \code{pt[i]} in the order of their residual systems; a point with free rank $>0$ is
    \emph{common} (\code{support = both}). G4 has one row group per common point, no aggregate per $\ell$; \code{lower} is $0$, a valid lower bound (equal to $\log_\ell|J\cap\Talg'|$ where
    the read's $j'=1$ certifies $J\cap\Talg'=0$), \code{upper} the Hom bound of that point. Where $\mathcal P$ does not generate on $H^{1}/\ell^{K}$ a torsion point is a point
    of $\Talg'$, not necessarily one maximal ideal of $\Talg$ (Remark~\ref{rmk:heckeZcbthree}).
  \item \textbf{G5--G7, torsion points, Bockstein data, compact support, each with its \code{theory} tag.} G5: per torsion maximal ideal, its system, \code{relation\_to\_free}
    $\in\{\code{distinct\_certified}$ (with the witness $T_p$)$,\code{common\_certified}$ (route \code{T\_ell\_level1} or \code{nonsplit})$,\code{agrees\_on\_P}\}$
    (Theorem~\ref{thm:heckeZdescent}(5)) and \code{generation\_tors}. G6: per common point \code{split} $\in\{\code{split\_by\_support},\code{split},\code{non\_split},
    \code{blocked(export)}\}$ (Propositions~\ref{prop:heckeZsplit}, \ref{prop:heckeZglue}(2)) --- the solver runs at every block size, a Lean cap governs only the certificate
    layer --- and the bound \code{hom\_order}; the universal-coefficient types are derived from item~7's groups, not stored (Proposition~\ref{prop:heckeZuct}). G7: the primes
    $\ell\ge5$ with $\Hcpt{2}[\ell]\ne0$, \code{theory = H\_c}, and \code{hecke = not\_run}: no route computes a Hecke action on $\Hcpt{2}$-torsion, and the record asserts none.
    Every mod-$\ell^{K}$ field of G2--G6 has \code{theory = H}.
    (v1.25.0-pre.4, F1, F3, F12, F13; not yet seen by the owner:) G5 and G6 are written per torsion point of (g), \code{pt[i]}, G5 carrying its system (the residue $\pi_p$ per
    operator); a point that is not common is \code{distinct\_certified} with the witness \code{T\_p} (the first operator whose residue differs from that of every point with a free
    part) or \code{joint}, and its G6 \code{split} is \code{split\_by\_support}. \textbf{\code{split\_on\_P}} (final, a statement about $\Talg'$ like \code{agrees\_on\_P}) is
    added to G6's list: the extension splits $\Talg'$-equivariantly at the point (Proposition~\ref{prop:heckeZsplit}). It is \emph{not} \code{split} for $\Talg$:
    Proposition~\ref{prop:heckeZglue}(2)(b) needs $\mathcal P$ to generate the image of $\Talg$ on $H^{1}_\ell/\ell^{K}$, torsion included, the mixed generation of
    Remark~\ref{rmk:heckeZcbthree} (\CONJ), and G2's generation (of $\Talg^{\mathrm{fr}}$ on $L$) does not give it. v1 writes \code{split\_on\_P} for every $\Talg'$-split point
    and never \code{split}, which is kept for a later theorem; \code{non\_split} and \code{split\_by\_support} are statements about $\Talg$. v1 computes neither
    \code{generation\_tors} nor the Tor terms of $H^{2}$: it writes \code{generation\_tors = not\_run} and, at $q=2$, G5 / G6 \code{not\_run} per $\ell$ (final). The body has
    G4, then G5, then G6.
    (v1.25.0-pre.5, HZ-REC re-audit R4, R11; not yet seen by the owner:) \code{G7.pair = not\_run} (final): the descent cell of \code{descent\_ref} runs \code{parts = groups,hecke}, without the pair, so v1 writes
    no pair datum. \emph{Where the \code{theory} tag is written}: once per $\ell\in\mathsf S$ (\code{G2.ell=<l>.theory = H}) and once on G7 (\code{G7.theory = H\_c}); the rows of G4--G6
    carry no tag of their own and are \code{H}.
    \textbf{G8, the Hida block per maximal ideal} (v1.25.0-pre.6, the slot-25 H1 merge after its audit \file{docs/implementation/slot25/H1\_audit.md}; \textbf{not yet seen by the owner}).
    Written \emph{only} under the non-default settings value \code{heckez.G8 = hida} ((i), (j)); a default record has no G8 field, and any value other than \code{off}, \code{hida} is
    refused. The paths: \code{G8.scope} $\in\{\code{level1},\code{not\_run(level\_N)}\}$ (tame level $1$ only); per prime $\ell\ge5$ of $\mathsf S$, \code{G8.ell=<l>.blocks}
    $\in\{\code{read},\code{not\_run(no\_decomposition)}$ (G2 did not decompose the maximal ideals at $\ell$)$,\code{not\_run(class\_zero)}$ ($(\ell-1)\mid k$)$,
    \code{not\_run(forms\_cap)}\}$; and per maximal ideal \code{m[i]} of G2 at $\ell$ (the same index $i$) \code{G8.ell=<l>.m[i].ordinary} $\in\{\code{true},\code{false}\}$ (the Eisenstein
    ideal is ordinary by theorem and the forms side must agree; another ideal is ordinary iff the residue of $\vartheta$ at $\mathfrak m$ divides $\chi(\vartheta\mid e\,M_k)\bmod\ell$, with
    $T_\ell=U$ on $e\,M_k$; a multiplicity other than $0$ or the full one is a defect, never a value). At an ordinary Eisenstein ideal through the cuspidal component, in addition:
    \code{chain\_key} $=(1,\ell,k\bmod(\ell-1))$ (three decimals); \code{g}, the rank $g_a$ of the ordinary Eisenstein block of $e\,S_k$ ($e$ the ordinary projector) (Fitting split by the full residual system,
    $U$ included); \code{L\_k}, G3's certified $C_T$ $\ell$-length, else \code{not\_run}; \code{v\_Bk} $=v_\ell(B_k/k)$ read modulo $\ell^{W}$, else \code{not\_run};
    \code{structure} $\in\{\code{fibre\_LxL},\code{regular\_deg\_g},\code{open}\}$: \code{fibre\_LxL} iff $g=1$ (the fibre product $\Lambda\times_{\Lambda/(G_a)}\Lambda$,
    [U,~\texttt{prop:u-H11}]), \code{regular\_deg\_g} iff $g\ge2$, some $T_q-(1+q^{k-1})$ generates the block modulo $\ell$ with $v_\ell\det(T_q-(1+q^{k-1})\mid\text{block})=$ \code{v\_Bk}, and
    $v\bigl(G(x_k)-G(x_{k+\ell-1})\bigr)=1$ ($\lambda=1$; then $\mathbf h(a)_{\mathfrak m}\cong\Z_\ell[[Y]]$, [U,~\texttt{prop:u-H12}](1), (2) and [U,~\texttt{rmk:u-H12general}]), else \code{open}
    (the conservative word); \code{I0\_generator} $\in\{$\code{T\_<q>}$,\code{none\_found}\}$, the first such $T_q$ ($q$ ascending, $q\ne\ell$). \emph{Caps}: the forms side is read
    modulo $\ell^{W}$, $W=\max(2,\text{upper}_v+2)$, and served iff $\ell^{W}<2^{84}$, $k-1\ge W$ and the $q$-expansion length
    $\max(\ell,k/12+1,7,\text{the operators of }\vartheta)\cdot(\dim S_k+2)+1\le2^{18}$, else \code{blocks = not\_run(forms\_cap)}; count-based record content, in the hashed
    \code{heckez.G8 = hida} line ((j)). Outside the body: G9 carries K9's integers \code{G9.K9.ell=<l>} (the certificate layer; \code{absent(not\_applicable)} under \code{hida} where no
    ordinary Eisenstein ideal meets a cusp form, \code{absent(H1)} under the default), G10 the read's $W$, time and refusal as \code{G10.hida.ell=<l>} (provenance).
  \item \textbf{The hash partition.} \code{record\_fnv1a64}, the one FNV-1a-64 (\file{ModArithmetic/Fnv.h}) of a fixed-order body --- schema, key, operators,
    \code{descent\_ref.record\_fnv1a64}, G1 to G8 in order, every list sorted by its stated invariant key, every big integer as its canonical hexadecimal text (item~7~(l)), every
    status, then the settings' canonical text with its fingerprint. \textbf{A hashed field is determined by (cell, algebra, settings) and independent of the lattice basis and of
    algorithm internals.} The \emph{certificate layer} G9 (B\'ezout pairs, witness positions, dump hashes, the $K$ used, Pratt trees, the Lean caps; the fields of
    \file{07-lean-plan.md} \S11.1) has its own \code{cert\_sha256}, outside \code{record\_fnv1a64}; \emph{provenance} ($K$ and route of each read, word primes, effort, timings,
    RSS, binary hash) is in neither; the \emph{Lean status} of each certificate lives in a side manifest keyed by (cell, checker), outside both (Q17(i)), so a Lean run never
    moves a record.
    (v1.25.0-pre.4, F2, F11; not yet seen by the owner:) \textbf{Default v1 records carry no G8 field}: G8 (the Hida block per maximal ideal) is written only under a non-default
    value of the settings line \code{heckez.G8} (default \code{off}), by the package H1, so no default record moves when H1 lands (v1.25.0-pre.6: that value is \code{hida}; its
    paths, closed lists and caps are the G8 paragraph of (h)). The word primes that carry the residues
    (\code{CrtPrimes}, below $2^{30}$) are provenance and are not in the settings text.
    (v1.25.0-pre.7, HZ-RAT-PERF; not yet seen by the owner:) Their count cap is printed in the same provenance key,
    \code{word\_primes=CrtPrimes(<2\^{}30,sized\_from\_bound,cap=2\^{}21)}, the last key of HZ-RAT's settings text, which the record erases from its hashed caps text (so
    \code{G10.rat.settings} and \code{G10.rat.word\_primes} show it); it moves no record id, hash, fingerprint or pin.
    (v1.25.0-pre.5, HZ-REC re-audit R10; not yet seen by the owner:) One use of that sequence is not provenance: (C1)'s primitivity probe runs at its first \code{c1\_probe\_primes} primes (the count is hashed, $3$ by default),
    so \code{primitive} could be a false negative if all of them divided a discriminant (never a false positive; the chance is far below $10^{-15}$). The sequence is fixed in v1;
    a build that changed it would change the settings text (a new fingerprint).
  \item \textbf{Caps are record content.} \code{P\_rule}, \code{theta\_search}, \code{factor\_budget} (trial division to $10^{6}$; rho, ECM and the Pratt search by fixed
    counts), \code{bezout\_record\_cap} and $K_{\max}$ per read are in the settings text and count-based, \textbf{never time-based}, so a record is the same on every machine
    (Q17(ii)); a non-default value is a new fingerprint, cell id and hash input (item~7~(h)).
    \emph{The read set at $\ell$} (v1.25.0-pre.1, 2026-11-06; the slot-25 amendment of the plan's \file{04-record-census.md} \S1.5, lead decision of 2026-11-06 on
    the HZ-LOC audit's F1). At each read prime $\ell$ the operator set of the (C3) and at-root reads is $\mathcal P$ plus $T_\ell$ when $\ell\nmid N$ and $\ell\le23170$
    ($T_\ell$ appended, so the primitive $\vartheta$ stays on $\mathcal P$; the cap is count-based, $T_\ell$ having $\ell+1$ cosets, and is the narrow tier of
    Lemma~\ref{lem:bounds}), and $v_\ell$ of the upper bound is taken over $\mathcal P\cup\{\ell\}$: $v_\ell(g_{\mathcal P\cup\{\ell\}})=\min(v_\ell g_{\mathcal P},v_\ell c_\ell(\lambda_\ell))$
    (Definition~\ref{def:heckeZcong} and Proposition~\ref{prop:heckeZcong} hold for every operator set). Otherwise $\mathcal P$ is unchanged. The per-read field
    \code{t\_ell} takes its value in a closed list (v1.25.0-pre.2, Z0 referee F6): \code{in\_P} ($\ell\in\mathcal P$ already), \code{added}, \code{omitted(cap:ell>23170)},
    \code{omitted(ell|N:U\_ell\_not\_exported)} (EXPORT exports no $U_p$, $p\mid N$), \code{omitted(export:<Name>)} with \code{<Name>} the export's refusal \emph{name}
    (\code{ExportRefusalName}) only, never its free-text message. The cap belongs to the settings text: \code{P\_rule}'s default text carries
    ``$\mathcal P\cup\{T_\ell\}$, $\ell\nmid N$, $\ell\le23170$''.
    \code{t\_ell} is a function of (cell, settings, $\ell$) and is hashed; the word regime of a read (\code{u32}, \code{u64}, \dots) is provenance, as (i). Effect on the
    plan's cells: $(1,42)$ at $11$ becomes (C3) \code{certified} with the $11$-part of $C_{\Talg}$ exactly $1$ (on $\mathcal P=\{2,3,5,7\}$ it was \code{open(sandwich)} and
    \code{upper\_bound\_certified}), $(1,40)$ at $3$ closes with $3$-part $1$ (\file{docs/implementation/slot25/HZLOC\_report.md}, ``Audit fixes'').
    \emph{The factor budget as built} (v1.25.0-pre.1, HZ-RAT, \file{docs/implementation/slot25/HZRAT\_report.md} \S4): trial division to $10^{6}$, the word factoriser on
    cofactors below $2^{64}$, the Pratt search and verification, and the cell's pool of certified primes; no multi-limb rho or ECM yet, so the C++ frontier lies below the
    measured one of \S\ref{sec:heckeZ} (level $1$: (C2) complete for $k\le34$ and $k=38$; $g_{\mathcal P}$ \code{partial(unfactored)} from $k=36$ on, $k=44$ excepted). The
    first golden set of (k) lies inside it. Adding rho / ECM later changes \code{factor\_budget}, hence the fingerprint (as above). The Lean caps (\code{k1\_cap}: $f\le8$, $b_1\le40$; \code{k7\_cap}: block rank $\le6$,
    $\dim\mathrm{tors}_\midl\le6$; Q12) sit in G9's header and decide only the Lean status.
    (v1.25.0-pre.4, F2, F11, F16 and the Z0 referee's F10; not yet seen by the owner:) The default caps text has the lines \code{heckez.P\_rule},
    \code{heckez.theta\_search+factor\_budget} (HZ-RAT's text without its \code{word\_primes} key, then \code{pool=fixed\_point;multilimb\_rho=none;ecm=none}: one \code{rho=} key,
    the word factoriser), \code{heckez.S\_rule} (every certified prime of $\gcd_j D_j$, $v_{\mathrm{disc}}=1$ included --- no LC-3b drop rule in v1, so the reason list carries no
    word for it), \code{heckez.K\_read} ((C3) at $K=1$ when the $K_0$ read is not served), \code{heckez.maximal\_ideals} ((e)), \code{heckez.irreducible} ((d)),
    \code{heckez.G8 = off} ((i)) and \code{heckez.bezout\_record\_cap}. (v1.25.0-pre.5, HZ-REC re-audit R2, R8, R9; not yet seen by the owner:) \code{heckez.maximal\_ideals} names the widened search over $\mathcal P\cup\{T_\ell\}$ ((e)),
    \code{heckez.irreducible} the cap of $256$ primes ((d)), and the new line \code{heckez.second\_route} the three word caps of the second route ((f)); the second route is one of the
    widenings below. \textbf{Defaults never change}: every later widening (rho / ECM, a wider decomposition rule, a
    factoriser over $\Q$, a drop rule for $\mathsf S$, H1's G8) enters as a non-default value with a new fingerprint, so the default records and the golden objects of (k) never move.
  \item \textbf{Golden: records with no pending field only, add-only} (Q3(a)). A record becomes a golden object iff its overall status is \code{final} \emph{and it is built under the
    default settings only} (v1.25.0-pre.6, H1 audit F1, not yet seen by the owner: a final record under a non-default value, e.g.\ \code{heckez.G8 = hida}, is never golden; the
    record's \code{golden\_eligible} flag, outside the hash, says so); final statuses that are
    bounds or partial (\code{bounds} for $J$, \code{upper\_bound\_certified}, \code{partial(unfactored)}, \code{open(\dots)}, \code{agrees\_on\_P}, \code{split\_on\_P} (v1.25.0-pre.4), \code{not\_run},
    \code{refused(\dots)} with a final reason) do not prevent it. One tier-A object \code{HZ\_N<N>\_n<n>\_<alg>} $=$ \code{record\_fnv1a64} per cell; the event is add-only, landed at the end of a slot
    (the first planned at the slot-25 close: $(1,10)$, $(1,22)$ and $(1,n)$, $n\in\{2,4,6,8,12,14,16,18,20,24,26,28,30\}$, $15$ objects under Q1(a)).
  \item \textbf{Big integers and the literature routes} (Q1(a), \emph{residue-only}). Every integer of any size in the record is produced from residues modulo word primes
    with a certified bound and written by item~7~(l)'s limb route, the cofactor included (its residues are $D\prod q_i^{-e_i}$); the C++ build certifies factor
    multiplicities of word-size primes (by residues, or by the cofactor route of \S\ref{sec:heckeZ}), the cofactor, Pratt certificates below $2^{64}$ and the cofactor
    test; a factor above $2^{64}$ is \code{probable}; no multi-precision division, gcd, modular power or square root enters the certified build, and the factor finder is
    untrusted. The routes \code{mazur} and \code{second\_route} rest on the literature hypotheses named in Proposition~\ref{prop:heckeZeiscert}(2), (3) ([U,~\texttt{pr:u-U1}];
    \cite{mazur}), and the Lean certificates name every C++ output they trust (the dumped operator matrices, the descent's $\chi_\Z$, $B_k/k\bmod\ell^{2}$; Q10).
  \end{enumerate}
  \textbf{Code map} (v1.25.0-pre.1, 2026-11-06): the record writer HZ-REC is not built (\file{06-work-packages.md} \S2.3). Its inputs are, in the main line: EXPORT,
  \file{Integral/Descent/DescentModuleExport.\{h,cpp\}} (the module $H^{1}\otimes\Z/\ell^{K}$ with every $T_p$, $p\nmid N$; \code{HeckeZExportValidation}); HZ-RAT and HZ-UTAG,
  \file{Integral/HeckeZRational.\{h,cpp\}}, \file{Integral/HeckeZFactorFinder.\{h,cpp\}} (untrusted), \file{Integral/HeckeZUTag.\{h,cpp\}} (G1, $g_{\mathcal P}$, the two
  full-algebra tags; \code{HeckeZRationalValidation}, \code{HeckeZUTagValidation}); HZ-LOC, \file{Integral/HeckeZLocal.\{h,cpp\}} and \file{HeckeExt/HeckeZLocalSource.\{h,cpp\}}
  ((C3), the local types, the at-root read; \code{HeckeZLocalValidation}, \code{HeckeZLocalCensusXcheck}). Not yet in the main line: EXPORT-Z step~1 (the integral blocks
  $\tau_p$, $M_p$ of $T_p$ on $H^{1}$ over $\Z$, the free block pending its step~2), which this item does not use.
\end{enumerate}

\noindent
\textbf{What a record is not.} It is not a proof, and no field of it may be branched on
by the pipeline that wrote it (rule~\ref{A6}). \code{cochain\_commutes},
\code{eta\_zero}, the $\inv$-$\tau$ symmetry flag of Obligation~\ref{obl:iota}(iv) and
the recorded solve regime are diagnostics: they are written so that a later reader can
tell what happened, and a true value licenses nothing beyond the statement that names
it.

\subsection{Document machinery and its guards}\label{sec:machinery}

\noindent
This document has four guards, of which three exist and one does not.

\begin{enumerate}
\item \textbf{The build gate} (\file{docs/latex/build\_gl2\_contract.sh}). Runs
  \code{pdflatex}, \code{bibtex}, \code{pdflatex}, \code{pdflatex} and \textbf{fails}
  on any undefined reference, undefined citation or multiply-defined label. The label
  namespace is part of the contract (\ref{L6}), so a dangling \code{\textbackslash ref}
  is an error and not a warning. Three appendices are written so as to be caught by it:
  Appendix~\ref{app:labelindex} dereferences every externally cited label,
  Appendix~\ref{app:changelog} dereferences every label it names into a discarded box,
  and Appendix~\ref{app:status} prints \code{\textbackslash ref} for every statement it
  indexes --- so an index entry naming a label this document does not define is a build
  failure, exactly as \ref{L6} requires.
\item \textbf{The citation census.} The converse check --- that every label the
  \emph{code} cites as \code{[C, \ldots]} is defined here --- is not gated by the LaTeX
  build, because the citing sites are C++ comments. It is specified as
  \code{label\_census.sh} (architecture \S8.4). Until v1.5.1 it \textbf{was not
  written}; had it existed it would have caught entry~\ref{err:e7} years earlier, and
  Remark~\ref{rmk:labelcensus} records the runs of it done by hand. \textbf{Since v1.5.2
  it exists and is registered} (slot-3 WP-10): \file{docs/matrix-engine/tools/label\_census.sh}
  with its parser \file{label\_census.py} and an allow-list with a reason per entry, run
  inside the suite by \code{ContractCensusValidation}, regenerating
  \file{docs/latex/label\_census.tsv} on every run and failing on any unresolved
  \code{[C, <label>]} or \code{\textbackslash HA\{<label>\}} citation; the site count is
  never gated.
\item \textbf{The two migration guards} (\file{docs/matrix-engine/tools/}).
  \code{raw\_mod\_calls.sh} and \code{check\_matrix\_includes.sh} each own an allow-list
  with a reason per entry and are run inside the validation suite by
  \code{ConsumerMigrationValidation} and \code{MatrixRetirementValidation}, so they
  cannot rot. They guard the substrate, not this document; they are named here because
  they are the pattern item~2 should follow --- a script with an allow-list, registered
  as a validation.
\item \textbf{The version and changelog convention}, \S\ref{sec:versioning}
  (\ref{V1}--\ref{V5}), whose record is Appendix~\ref{app:changelog} and whose rules
  \ref{CL1}--\ref{CL6} are stated there. \ref{V4}'s two supersession steps --- freezing
  \file{hecke\_action.tex} and \file{monodromy\_v5.tex} with a banner --- have not been
  taken.
\end{enumerate}

\begin{remark}[The citation census, run by hand]\label{rmk:labelcensus}
\COMP{} On 2026-09-16 every occurrence of the citation form \code{[C, <label>]} was
extracted from the whole repository outside \file{docs/latex/} and intersected with the
\code{\textbackslash label} set of \file{gl2\_contract.tex} and
\file{gl2\_notation.tex}. Result: \textbf{$64$ distinct labels cited, $761$ citation
sites in $78$ files, and $0$ cited labels that this document does not define} --- of
which $53$ distinct labels and $447$ sites in $68$ files are in C++ translation units
(\code{.h}, \code{.cpp}, \code{.cu}, \code{.inc}), the rest in project documents. The
most-cited are \code{alg:tau} ($37$ C++ sites), \code{thm:backend} ($30$),
\code{obl:tau} ($18$), \code{eq:xiexplicit} ($17$), \code{cor:fp32} ($17$) and
\code{sec:example} ($16$, in $8$ files). Two apparent misses were rejected on
inspection: both are prose in \file{docs/implementation/math\_corpus\_survey.md} that
writes the citation \emph{form} rather than a citation. \textbf{This is a recorded run
of an ad-hoc script, not a registered validation}, and the figures differ in detail from
the architecture's own census of 2026-09-15 ($519$ sites in $71$ files) because the two
used different extraction rules; the load-bearing number --- zero dangling citations ---
is the one that does not depend on the rule. Until item~2 above exists, this check is
re-run by hand or not at all.

\COMP{} \textbf{Re-run at v1.0.0}, after \file{hecke\_extensions.tex} v1.3, by the same
rule: \textbf{$63$ distinct labels cited, $778$ citation sites in $79$ files, and $0$
cited labels that this document does not define}; the sole non-resolving match is again
the prose \code{[C, <label>]} of
\file{docs/implementation/math\_corpus\_survey.md}, which writes the citation form rather
than a citation. Two further checks were run for \ref{V5}\ref{V5ii}, which the census
alone does not settle, because it tests only that a \emph{cited} label resolves:
\textbf{(i)} the \code{\textbackslash label} set of this document together with
\file{gl2\_notation.tex} was intersected with that of \file{hecke\_extensions.tex} --- the
only companion still live --- and the intersection is \textbf{empty}; \textbf{(ii)} in
particular none of the nine labels of \S\ref{sec:six} is defined there any longer. Check~(i)
found, and the companion's v1.3 removed, one collision the nine did not include:
\file{hecke\_extensions.tex} used \code{A1}, \code{A2}, \code{A3} for its three required
additions while this document uses them for authoring rules \ref{A1}--\ref{A3}. No citation
anywhere used either reading, and the companion renamed its three; but a census of citations
would never have found it, which is the argument for making item~2's script compare label
\emph{sets} and not only citations.

\COMP{} \textbf{Registered since v1.5.2} (hecke-extensions slot 3, package WP-10,
\file{docs/implementation/slot3/WP-10\_report.md}): item~2's script exists, so this check
is no longer re-run by hand. \code{ContractCensusValidation}, run in the slot-3 close's
\code{--all} on 2026-09-17: \textbf{$429$ labels defined, $1408$ citation sites in $121$
of $430$ scanned files, $141$ distinct labels cited, $4$ allow-listed (prose writing the
citation form with the placeholder word), $0$ unresolved}; entry~\ref{err:e7}'s historical
dangling citation (\code{thm:explicit\_complex\_equiv}) resolves. The registered script
compares citations against the label set; the label-\emph{set} comparison with the
companion argued for above is not part of it and stays a check run by hand.

\COMP{} \textbf{Content-keyed since v1.6.2} (hecke-extensions slot 4, package WP-6,
\file{docs/implementation/slot4/WP-6\_report.md}). The allow-list is no longer keyed by
\code{relpath:line} but by \code{relpath|<verbatim fragment of the site's source line>},
so that inserting or deleting lines above a cited passage cannot silently move an entry
off its site --- which is what makes the format safe for a revision of this document. An
entry matching \emph{zero} sites is flagged \textsc{stale} and \emph{fails} the census,
by the same ``removal is the unit of progress'' rule the other allow-lists use, and the
retired \code{relpath:line} form is rejected with a named error rather than silently
reinterpreted. \code{ContractCensusValidation} gained four arms for exactly these
properties, all against scratch roots and never the real allow-list. Run in the slot-4
close's \code{--all} on 2026-09-22: \textbf{$429$ labels defined, $1575$ citation sites
in $128$ of $471$ scanned files, $152$ distinct labels cited, $4$ allow-listed, $0$
unresolved, $0$ stale}. The site count is never gated --- it grows every slot; only the
unresolved and stale counts are.
\end{remark}

\paragraph{Formalization note.}
This chapter states no mathematics and is not a formalization target. One thing in it is
worth carrying across, though: the record schema of \S\ref{sec:recordschema} is the
\emph{interface} at which every obligation of \S\ref{sec:obligations} becomes decidable,
and a formalization that models the pipeline as a function from parameters to a record,
with the obligations as predicates on that record, gets the whole of \ref{A4} for free.
The one field that resists this treatment is the certification level, because it changes
which predicates are evaluated rather than what they say --- which is why
\S\ref{sec:buildsurface} requires the record to carry it.

%======================================================================
\part*{Part VI --- The integral line}
\addcontentsline{toc}{section}{Part VI --- The integral line}
%======================================================================

%----------------------------------------------------------------------
\section{The integral line: coefficient rings, the free model on the unrefined tessellation, and descent over \texorpdfstring{$\Z$}{Z}}\label{sec:integral}
%----------------------------------------------------------------------
% SOURCE:  NEW (v1.8.0, hecke-extensions slot-8 WP-9, "L2-C").  Project notes, not
%          sources of record: docs/implementation/integral_line_design_brief.md
%          revision 2 ([IB]) SS2.1-2.6, 3.2, 5.1, 5.2, 6.1-6.5 and its audit
%          docs/implementation/integral_line_design_brief_audit.md ([IBA]);
%          docs/implementation/e8_ladic_investigation.md ([L]) SS2.1-2.6;
%          docs/implementation/hecke_extensions_slot8_plan.md S3 "WP-9", S10
%          (owner rulings 6b, 6c, 6d, 10 of 2026-09-24).
% NOTATION: sec:notation-integral (new at v1.8.0); crosswalk rows cring, e (the
%          ring exponent, never k), q (degree); [IB]'s symbols are read through
%          the table of sec:notation-integral.
% ERRATA:  none.  This chapter asserts nothing about a superseded document.

\paragraph{What this chapter fixes.}
The mathematics of \emph{integral coefficients} on the model, which until v1.7.4 this
document fixed only in pieces (Proposition~\ref{prop:coneintegral},
Claims~\ref{cl:amalgam} and \ref{cl:bockstein}, Lemma~\ref{lem:bulkperiodic}):
\begin{enumerate}
\item which object the shipped complex is \emph{over the ring in force} --- three objects
  live on the model at $\ell\in\{2,3\}$, and the pipeline computes the third
  (Definition~\ref{def:ringinforce}, Proposition~\ref{prop:ringinforce}); what that object
  is, independently of the triangulation (Proposition~\ref{prop:bredon});
\item a free $\Ga$-model on the \emph{unrefined} tessellation $(\Xbulk,\Xbdry)$ --- the
  isotropy-resolution model $\Iso_{\bullet}$ named in Remark~\ref{rmk:mpcaveat}(4), with
  its boundary subcomplex $\Isob_{\bullet}$ and one integral comparison map $\cmpI$
  (Definition~\ref{def:isomodel}, Proposition~\ref{prop:isomodel}); this is the contract
  item C-A5 of the extension plan, \textbf{re-scoped} by owner ruling 6c of 2026-09-24
  (Remark~\ref{rmk:cafive});
\item the descent: the arithmetic fracture square, cohomology-level glueing, the
  vanishing of the one coherence term the bulk could carry, and the realisation of the
  owner's descent by that one integral map
  (Propositions~\ref{prop:fracture}--\ref{prop:homvanish});
\item the Hecke datum on the free model: a strict lift of the strict pair datum of
  Proposition~\ref{prop:strictpair} --- \emph{existence} only
  (Proposition~\ref{prop:strictlift});
\item the bad primes (Proposition~\ref{prop:badprimes}), the certified-exponent rule, stated
  so that it holds at $\ell\in\{2,3\}$ (Proposition~\ref{prop:certifiede}), and the
  residue identity behind the characteristic polynomial over $\Z$
  (Proposition~\ref{prop:crtcharpoly});
\item what each coefficient ring \emph{names} (Definition~\ref{def:ringsemantics}) ---
  in particular what $H^{q}$ over $\Z$ means on the model and what the integral line will
  compute. The setting that selects the ring, its spelling, default, the reservation of
  $\Z$ and the refusal rule are the companion specification's, \Espec{sec:ring}
  (single-editor rule); this chapter never restates them.
\end{enumerate}
\textbf{What is not here}: any implementation of $\Iso_{\bullet}$, $\cmpI$ or the
lifted datum, and in particular their \emph{effective, certified} construction in degrees
$\le q_{\max}+1$, which is the key risk of the integral line and a code matter
(Remark~\ref{rmk:integralscope}) (v1.9.0: $\Iso_{\bullet}$, $\Isob_{\bullet}$ and $\cmpI$ are
built since hecke-extensions slot~9, with the integral groups for $n\le20$; the lifted datum is
not --- \S\ref{sec:codemap}, Remark~\ref{rmk:integralevidence}; v1.10.0: the lifted datum is built
for $g=\mathrm{diag}(1,p)$, $p\in\{2,3,5\}$, in degrees $\le q_{\max}+1$ by slot~10,
Remark~\ref{rmk:strictliftbuilt}, and the proposition itself still claims existence only);
arbitrary-precision arithmetic, which
Theorem~\ref{thm:backend}(3), the contract box of \S\ref{sec:fpfilter} and
Obligation~\ref{obl:backend}(iv) still forbid in the build, unchanged; integer Hecke
matrices in a basis; the divided-power lattice $\Dn$; families. \textbf{Placement.} The
chapter is a Part of its own, after Part~V, so that no chapter number of v1.7.4 moves;
mathematically it continues \S\ref{sec:orbifoldmodel}.

\paragraph{Standing hypotheses.}
\ref{H4} throughout, and \ref{B1}--\ref{B3} wherever the boundary occurs; per
rule~\ref{A1} every statement repeats the ones it uses. $\Vcoef=\Sym^{n}(\Z^{2})$ is the
monomial lattice of Definition~\ref{def:rhon} and nothing else (owner decision OD-4~(a),
confirmed 2026-09-24); $U:=ST$, as in Claim~\ref{cl:amalgam}. For a commutative ring
$\cring$, $\Ceq(\cring)$ is short for $\Ceq(K,\Rel;L)$ with fibre $\VR$ \emph{and the
relations imposed over $\cring$} --- the complex the implementation builds for that ring;
by Theorem~\ref{thm:explicit_complex_equiv} and Remark~\ref{rmk:orbifold},
$\Ceq(\cring)\cong\Hom_{\Z[\Ga]}(C_{\bullet}(\Xbulk),\VR)$, and this identification is used
without further mention.

\subsection{Coefficient rings, and the three objects on the model}\label{sec:integralobjects}

\begin{definition}[The three complexes on the model, and the notation rule]\label{def:ringinforce}
Assume \ref{H4}, let $n\ge0$ and let $\cring$ be a commutative ring.
\begin{enumerate}
\item The \textbf{integral model complex} is $\Ceq(\Z)=\Hom_{\Z[\Ga]}(C_{\bullet}(\Xbulk),\Vcoef)
  \cong\bigoplus_{\sigma}\Vcoef^{\Ga_{\sigma}}$ (sums over $\Ga$-orbit representatives in
  each degree), a complex of finitely generated free $\Z$-modules
  (Proposition~\ref{prop:coneintegral}(1)).
\item Its \textbf{base change} is $\Ceq(\Z)\otimes_{\Z}\cring$.
\item The \textbf{model complex over the ring in force} is $\Ceq(\cring)=
  \Hom_{\Z[\Ga]}(C_{\bullet}(\Xbulk),\VR)\cong\bigoplus_{\sigma}(\VR)^{\Ga_{\sigma}}$.
\end{enumerate}
The reduction map $\Ceq(\Z)\otimes\cring\to\Ceq(\cring)$ is the map of
Proposition~\ref{prop:coneintegral}(2). \textbf{Notation rule} (binding in this document
from v1.8.0; [IB]'s risk K12): $\Ceq(\cring)$ \emph{always} denotes (3), the object the
implementation computes; the base change is always written $\Ceq(\Z)\otimes\cring$; and
no statement may silently identify them --- the identification needs
$6\in\cring^{\times}$ (Proposition~\ref{prop:ringinforce}(1)).
\end{definition}

\begin{proposition}[The reduction map, and where the three objects differ]\label{prop:ringinforce}
\PROVED{} Assume \ref{H4}, let $n\ge0$ and let $\cring$ be a commutative ring.
\begin{enumerate}
\item The reduction map $\Ceq(\Z)\otimes\cring\to\Ceq(\cring)$ is an injective map of
  complexes in every degree, and an isomorphism of complexes when $6\in\cring^{\times}$.
  Consequently, for $\ell\ge5$ and every $e\ge1$, $\Ceq(\Z/\ell^{e})=\Ceq(\Z)\otimes\Z/\ell^{e}$.
\item \textbf{At $\ell=2$ it is not a quasi-isomorphism in general.} For $n=2$ and
  $\cring=\F_2$ the map $H^{0}(\Ceq(\Z)\otimes\F_2)\to H^{0}(\Ceq(\F_2))$ is not onto:
  the class of the Dickson form $x^{2}+xy+y^{2}$ lies in $H^{0}(\Ceq(\F_2))$ and not in the
  image. Already in degree $0$ of the complexes, at $\sigma=v_6$, the reduction
  $\Vcoef^{\langle S\rangle}\otimes\F_2\to(\VR)^{\langle S\rangle}$ has image of rank $1$ in
  a target of rank $2$.
\item \textbf{Odd $n$.} For $n$ odd, $\Ceq(\Z)=0$; for every $\cring$ with
  $2\in\cring^{\times}$, $H^{\bullet}(\Ceq(\cring))=0$; but $\dim_{\F_2}H^{1}(\Ceq(\F_2))=1$
  at $n=1$, and $H^{2}(\Ga;\Sym^{1}\Z^{2})\neq0$. So at odd $n$ the integral model
  complex, its reductions, the model complex over $\F_2$ and group cohomology over $\Z$
  are four different things.
\end{enumerate}
\end{proposition}

\begin{proof}
(1) is Proposition~\ref{prop:coneintegral}(2) and its proof (a saturated sublattice is a
direct summand; the averaging argument needs $\abs{\Ga_\sigma}\in\cring^{\times}$ and
$\abs{\Ga_\sigma}\mid12$), applied to the bulk complex alone; $\Z/\ell^{e}$ has
$6\in(\Z/\ell^{e})^{\times}$ for $\ell\ge5$.

(2) $S$ acts on $\Sym^{2}$ by a linear substitution sending $x\mapsto\pm y$,
$y\mapsto\mp x$ (Definition~\ref{def:rhon}), so $x^{2}\leftrightarrow y^{2}$ and
$xy\mapsto-xy$. Hence $\Vcoef^{\langle S\rangle}=\Z(x^{2}+y^{2})$, while over $\F_2$,
where $-xy=xy$, $(\VR)^{\langle S\rangle}=\F_2(x^{2}+y^{2})\oplus\F_2\,xy$: the image of the
reduction is the first summand. Now $H^{0}$ of either complex is its space of
$0$-cocycles, and the cochain map is injective by (1), so
$H^{0}(\Ceq(\Z)\otimes\F_2)\to H^{0}(\Ceq(\F_2))$ is injective and its image consists of
cocycles whose value at $v_6$ lies in $\F_2(x^{2}+y^{2})$. A $0$-cocycle of $\Ceq(\F_2)$ is
a $\Ga$-invariant vector $f\in\VR$ placed at every vertex ($\Xbulk$ is connected); the
Dickson form $f=x^{2}+xy+y^{2}$ is invariant under $\GL_2(\F_2)$ (it is invariant under
$x\leftrightarrow y$ and under $(x,y)\mapsto(y,x+y)$, which generate), hence under $\Ga$,
whose action on $\VR$ factors through the reduction $\Ga\to\SL_2(\F_2)=\GL_2(\F_2)$. Its
value at $v_6$ is $f\notin\F_2(x^{2}+y^{2})$.

(3) For $n$ odd the central relation $(\sigma,\sigma,-I)$ reads $f(\sigma)=\rhon(-I)f(\sigma)
=-f(\sigma)$ (Remark~\ref{rmk:negI}), so $2f(\sigma)=0$ and $f(\sigma)=0$ in the torsion-free
$\Vcoef$; over $\cring$ with $2\in\cring^{\times}$ the same relation kills every cochain.
At $n=1$, $\cring=\F_2$: by Proposition~\ref{prop:bredon}(2) below,
$\dim H^{1}=\dim\VR^{-I}-\dim(\VR^{S}+\VR^{U})$; $-I$ acts trivially, $S$ acts on $\F_2^{2}$
as the swap, with fixed line $\F_2(x+y)$, and $U-\id$ has determinant $1$ over $\F_2$ (for
any $A\in\SL_2$, $\det(A-\id)=2-\tr A$, and $\tr U=1$), so $\VR^{U}=0$; hence
$\dim H^{1}=2-1=1$. For the last clause read \eqref{eq:mvamalgam} over $\cring=\Z$ with
$\Vcoef=\Z^{2}$, on which $-I$ acts by $-1$. For a cyclic group $\langle g\rangle$ of order
$m$ the periodic free resolution $\cdots\xrightarrow{N}\Z[\langle g\rangle]\xrightarrow{g-1}
\Z[\langle g\rangle]\to\Z\to0$, $N=1+g+\dots+g^{m-1}$, is exact ($\ker(g-1)=\Z N=\im N$ and
$\ker N=\im(g-1)$ in $\Z[\langle g\rangle]$; for $m=2$ it is the resolution (ii) of the proof of
Proposition~\ref{prop:bdrycohom}), so $H^{1}(\langle g\rangle;M)=\ker N/\im(g-1)$ and
$H^{2}(\langle g\rangle;M)=M^{g}/NM$. On $\Z^{2}$: for $\langle-I\rangle$, the norm is $0$ and
$\im(-I-\id)=2\Z^{2}$, so $H^{1}=(\Z/2)^{2}$; for $\langle S\rangle$, the norm
$(1+S)(1+S^{2})$ is $0$ and $\abs{\det(S-\id)}=2$, so $H^{1}=\Z/2$, and
$(\Z^{2})^{S}=0$ gives $H^{2}=0$; for $\langle U\rangle$, $\det(U-\id)=1$, so $H^{1}=0$ and
$H^{2}=0$. The segment $H^{1}(\langle S\rangle)\oplus H^{1}(\langle U\rangle)\to
H^{1}(\langle-I\rangle)\to H^{2}(\Ga)\to H^{2}(\langle S\rangle)\oplus H^{2}(\langle U\rangle)=0$ of
\eqref{eq:mvamalgam} then gives $H^{2}(\Ga;\Z^{2})\cong\coker\bigl(\Z/2\oplus0\to(\Z/2)^{2}\bigr)$,
which is non-zero because $(\Z/2)^{2}$ has four elements and the image at most two.
\end{proof}

\begin{proposition}[The shipped complex is Bredon cohomology over the ring in force]\label{prop:bredon}
Assume \ref{H4}. Let $\mathcal T$ be the Bass--Serre tree of
Lemma~\ref{lem:bulkperiodic} (vertex set $\Ga/\langle S\rangle\sqcup\Ga/\langle U\rangle$,
edge set $\Ga/\langle-I\rangle$, the edge $\gamma\langle-I\rangle$ oriented from
$\gamma\langle S\rangle$ to $\gamma\langle U\rangle$), whose augmented chain complex is
\eqref{eq:treeses} (Claim~\ref{cl:amalgam}(1)). Let $\mathcal D$ be the family of subgroups
of $\Ga$ contained in a conjugate of $\langle S\rangle$ or of $\langle U\rangle$.
($\mathcal D$ is in fact the family of all finite subgroups of $\Ga$, the family for which
Bredon cohomology of proper actions is usually stated; this document does not use that.)
\begin{enumerate}
\item \PROVED{} \textbf{(Fixed sets.)} For every $H\in\mathcal D$: if $H\subseteq\{\pm I\}$
  then $\Xbulk^{H}=\Xbulk$ and $\mathcal T^{H}=\mathcal T$; otherwise $\Xbulk^{H}$ and
  $\mathcal T^{H}$ are each a single vertex. In particular every fixed set is acyclic,
  and every simplex stabiliser of $\Xbulk$ and of $\mathcal T$ lies in $\mathcal D$.
\item \PROVED{} \textbf{(The comparison.)} $C_{\bullet}(\Xbulk)$ and $C_{\bullet}(\mathcal T)$
  are chain homotopy equivalent as complexes of $\Z[\Ga]$-modules, by
  augmentation-preserving maps. Hence, for \emph{every} $\Z[\Ga]$-module $M$ --- in
  particular $M=\VR$ for every commutative $\cring$ and every $n$ --- the complex
  $\Hom_{\Z[\Ga]}(C_{\bullet}(\Xbulk),M)$ is chain homotopy equivalent to the tree complex
  \begin{equation}\label{eq:treemodel}
    M^{\langle S\rangle}\oplus M^{\langle U\rangle}\longrightarrow M^{\langle-I\rangle},
    \qquad (a,b)\longmapsto b-a ,
  \end{equation}
  in degrees $0,1$. So, for $M=\VR$: $H^{0}(\Ceq(\cring))=\VR^{S}\cap\VR^{U}=\VR^{\Ga}$,
  $H^{1}(\Ceq(\cring))=\VR^{-I}/(\VR^{S}+\VR^{U})$ and $H^{q}(\Ceq(\cring))=0$ for $q\ge2$;
  over a field, $\dim H^{1}=\dim\VR^{-I}-\dim\VR^{S}-\dim\VR^{U}+\dim H^{0}$.
\item \PROVED{} \textbf{(What it is.)} $H^{\bullet}(\Ceq(\cring))$ is
  $\Ext^{\bullet}$, over the category of contravariant functors from the orbit category of
  $\mathcal D$ to abelian groups, from the constant functor $\Z$ to the fixed-point functor
  $\Ga/H\mapsto\VR^{H}$: the \textbf{Bredon cohomology} of $\Ga$ for the family $\mathcal D$
  with the fixed-point coefficient system. This is the invariant \ref{H3} calls the
  orbifold model cohomology; it is \textbf{independent of the triangulation} --- any
  simplicial $\Ga$-complex $Y$ acting without inversions, with stabilisers in $\mathcal D$
  and every $Y^{H}$ ($H\in\mathcal D$) acyclic, gives a complex chain homotopy equivalent
  to it by the same argument, a $\Ga$-equivariant subdivision of $\Xbulk$ among them ---
  and it agrees with $H^{\bullet}(\Ga;\VR)$ when $6\in\cring^{\times}$
  (Remark~\ref{rmk:orbifold}) but \emph{not in general} otherwise: at $\ell=2$, $n=0$ group
  cohomology is non-zero in every degree (Lemma~\ref{lem:h2torsion}(2)), while (2)
  vanishes above degree $1$.
\end{enumerate}
\end{proposition}

\begin{proof}
(1) $-I$ is central and lies in every simplex stabiliser of $\Xbulk$ (\ref{H4}(iv)) and
of $\mathcal T$ ($-I=S^{2}=U^{3}$, and the edge stabilisers are $\gamma\langle-I\rangle
\gamma^{-1}=\langle-I\rangle$), so it acts trivially on both, which is the case
$H\subseteq\{\pm I\}$. Let $H\in\mathcal D$ with $H\not\subseteq\{\pm I\}$, say
$H\subseteq\gamma\langle S\rangle\gamma^{-1}=\mathrm{Stab}_{\Ga}(\gamma v_6)$ or
$H\subseteq\gamma\langle U\rangle\gamma^{-1}=\mathrm{Stab}_{\Ga}(\gamma v_4)$
(Lemma~\ref{lem:model7coincidences}(1): $\langle S\rangle$ and $\langle ST\rangle$ are the
stabilisers of $v_6=i$ and $v_4=\zeta^{2}$). \emph{On $\Xbulk$}: $H$ fixes the vertex
$\gamma v_6$ (resp.\ $\gamma v_4$); a non-central $h\in H$ fixes at most one point of
$\Hyp$ (the last step of the proof of Lemma~\ref{lem:model7coincidences}), so $H$ fixes no
other point of $\Xbulk$, and no simplex of positive dimension --- a fixed simplex is fixed
pointwise, the action being without inversions (\ref{H4}(i)). \emph{On $\mathcal T$}: no
edge is fixed, the edge stabilisers being $\{\pm I\}$. The map
$\varphi(\gamma\langle S\rangle):=\gamma v_6$, $\varphi(\gamma\langle U\rangle):=\gamma v_4$
from the vertices of $\mathcal T$ to $\Hyp$ is well defined and $\Ga$-equivariant, with
$\mathrm{Stab}_{\Ga}(x)=\mathrm{Stab}_{\Ga}(\varphi(x))$ for every vertex $x$; it is
injective, because $\gamma v_6=\gamma' v_6$ forces $\gamma^{-1}\gamma'\in\langle S\rangle$,
likewise for $v_4$, and $v_6$, $v_4$ lie in different $\Ga$-orbits
(Lemma~\ref{lem:model7coincidences}(2)). So $H$ fixes the vertex $x$ iff it fixes the
point $\varphi(x)$, which happens for exactly one $x$ by the argument on $\Hyp$ just
given. The stabilisers of $\Xbulk$ are $\{\pm I\}$, $\langle S\rangle$,
$\langle ST\rangle=\langle U\rangle$ and $\langle TS\rangle=T\langle U\rangle T^{-1}$ and
their conjugates (\ref{H4}(iv)); those of $\mathcal T$ are the conjugates of
$\langle S\rangle$, $\langle U\rangle$, $\langle-I\rangle$. All lie in $\mathcal D$.
Acyclicity: $\Xbulk$ is contractible (\ref{H4}(ii)), $\mathcal T$ is acyclic by
\eqref{eq:treeses}, and a point is acyclic.

(2) Let $\mathcal O$ be the category whose objects are the $\Ga$-sets $\Ga/H$,
$H\in\mathcal D$, and whose morphisms are the $\Ga$-maps, and let $\mathcal A$ be the
abelian category of contravariant functors $\mathcal O\to\mathrm{Ab}$. For a simplicial
$\Ga$-complex $Y$ acting without inversions with stabilisers in $\mathcal D$, orient each
simplex $\Ga$-invariantly (possible because a stabiliser fixes its simplex pointwise) and
let $\underline{C}_{q}(Y)\in\mathcal A$ be $\Ga/H\mapsto C_{q}(Y^{H})$, the chains on the
fixed subcomplex, with the maps induced by $\Ga$-maps $\Ga/H\to\Ga/H'$ in the evident
way. Since $\Sigma^{H}\cong\mathrm{Map}_{\Ga}(\Ga/H,\Sigma)$ naturally for a $\Ga$-set
$\Sigma$, $\underline{C}_{q}(Y)\cong\bigoplus_{\sigma}\Z[\mathrm{Mor}_{\mathcal O}(-,\Ga/\Ga_\sigma)]$,
a sum over orbit representatives of $q$-simplices of representable functors, each of which
is projective in $\mathcal A$ by the Yoneda lemma
($\mathrm{Nat}(\Z[\mathrm{Mor}(-,c)],N)\cong N(c)$ is exact in $N$). The augmented complex
$\underline{C}_{\bullet}(Y)\to\underline{\Z}$ is exact at the object $\Ga/H$ iff $Y^{H}$ is
acyclic. By (1), $\underline{C}_{\bullet}(\Xbulk)$ and $\underline{C}_{\bullet}(\mathcal T)$
are therefore two projective resolutions of $\underline{\Z}$ in $\mathcal A$, and the
comparison theorem \STD{\cite[Thm.~2.2.6]{weibel}: a map of a projective complex over $M$
into a resolution of $N$ lifting $M\to N$ exists and is unique up to chain homotopy}
gives chain maps over $\id_{\underline{\Z}}$ in both directions whose composites are
chain homotopic to the identities. Evaluating at the object $\Ga/\{1\}$ --- whose
automorphism group in $\mathcal O$ is $\Ga$, so that naturality there is
$\Ga$-equivariance --- gives the $\Z[\Ga]$-chain homotopy equivalence
$C_{\bullet}(\Xbulk)\simeq C_{\bullet}(\mathcal T)$ with $\Z[\Ga]$-linear homotopies, and
applying the additive functor $\Hom_{\Z[\Ga]}(-,M)$ preserves it. Finally
$\Hom_{\Z[\Ga]}(\Z[\Ga/\Ga_0],M)=M^{\Ga_0}$, and the cochain differential of $\mathcal T$
is evaluation on $\partial(\gamma\langle-I\rangle)=\gamma\langle U\rangle-\gamma\langle S\rangle$,
which is \eqref{eq:treemodel} (note $M^{S}, M^{U}\subseteq M^{-I}$, since $-I=S^{2}=U^{3}$).
$H^{0}=M^{S}\cap M^{U}=M^{\Ga}$ because $S$ and $U$ generate $\Ga$; the dimension formula is
the rank count on the two-term complex.

(3) By Yoneda, $\mathrm{Nat}(\underline{C}_{q}(\Xbulk),\underline{M})\cong\bigoplus_\sigma
M^{\Ga_\sigma}\cong\Hom_{\Z[\Ga]}(C_{q}(\Xbulk),M)$ for the fixed-point functor
$\underline M(\Ga/H)=M^{H}$, compatibly with the differentials; since
$\underline{C}_{\bullet}(\Xbulk)$ is a projective resolution of $\underline{\Z}$, this
computes $\Ext^{\bullet}_{\mathcal A}(\underline{\Z},\underline{M})$. The independence
statement is (2) with $Y$ in place of $\mathcal T$; a $\Ga$-equivariant subdivision of
$\Xbulk$ has the same fixed \emph{point sets} as $\Xbulk$ (subdivided), hence the same
acyclicity, and the same stabilisers. The last clause: for $6\in\cring^{\times}$ it is
Remark~\ref{rmk:orbifold}; for $\ell=2$, $n=0$, (2) gives $H^{q}=0$ for $q\ge2$ while
$H^{q}(\Ga;\F_2)\neq0$ for every $q$ (Lemma~\ref{lem:h2torsion}(2)).
\end{proof}

\noindent
\emph{Evidence.} \COMP\textsuperscript{r} is recorded in Remark~\ref{rmk:integralevidence}:
the tree complex \eqref{eq:treemodel} over $\F_\ell$ reproduces the shipped pipeline's
$\dim\Heq{q}$ at $\ell\in\{2,3\}$ on every $n\le40$, odd $n$ included, and the sixteen
golden cells there. Proposition~\ref{prop:bredon}(2) is also the \emph{oracle} for the
model complex over any ring in force, $\cring=\Z/\ell^{e}$ at $\ell\in\{2,3\}$ included
(then with lengths and invariant factors of \eqref{eq:treemodel} in place of dimensions);
\textbf{it is not} the universal-coefficient reduction of the integral model complex,
which is a different object there by Proposition~\ref{prop:ringinforce}(2).

\subsection{The free model on the unrefined tessellation}\label{sec:isomodel}

\begin{definition}[The isotropy-resolution model $\Iso_{\bullet}$, its boundary subcomplex, and the comparison map]\label{def:isomodel}
Assume \ref{H4} and \ref{B1}--\ref{B3}. For $s\in\{0,1,2\}$ let $\Sigma_s$ be the set of
$s$-simplices of $K$ that represent the $\Ga$-orbits of $s$-simplices of $\Xbulk$ ---
by Lemma~\ref{lem:model7coincidences}(2), $\Sigma_0=\{v_0,v_1,v_2,v_4,v_6\}$, $\Sigma_1$
the nine edges of $K$ other than $v_3v_5$ and $v_5v_6$, and $\Sigma_2$ the five triangles,
lifted into $\Xbulk$ and stored-oriented --- and let $\Sigma^{\partial}_0=\{v_0,v_1,v_2\}$,
$\Sigma^{\partial}_1=\{v_0v_1,v_1v_2,v_2v_3\}$ be the representatives lying in $\Xbdry$.
Each $\Ga_\sigma$ is cyclic (\ref{H4}(iv)); fix the generator $g_{\sigma}=ST$ at $v_4$,
$S$ at $v_6$, $-I$ elsewhere, of order $m_\sigma\in\{6,4,2\}$, and let $F^{\sigma}_{\bullet}$
be the periodic free resolution of $\Z$ over $\Z[\Ga_\sigma]$ ($F^{\sigma}_{t}=\Z[\Ga_\sigma]$,
differential $g_\sigma-1$ from odd $t$ and the norm $\sum_{u\in\Ga_\sigma}u$ from even
$t\ge2$). Put $G_{s,t}:=\bigoplus_{\sigma\in\Sigma_s}\Z[\Ga]\otimes_{\Z[\Ga_\sigma]}F^{\sigma}_{t}$.
An \textbf{isotropy-resolution model} of $(\Xbulk,\Xbdry)$ is a complex of
$\Z[\Ga]$-modules $\Iso_{\bullet}$ with
\[
  \Iso_{q}=\bigoplus_{s+t=q,\ 0\le s\le2}G_{s,t},\qquad
  d=d^{(0)}+d^{(1)}+d^{(2)},
\]
where $d^{(0)}=(-1)^{s}(\id\otimes\partial^{F})$ on $G_{s,\bullet}$ and $d^{(r)}$ maps
$G_{s,t}$ to $G_{s-r,t+r-1}$, together with the map
$\cmpI\colon\Iso_{\bullet}\to C_{\bullet}(\Xbulk)$ that is the augmentation
$\Z[\Ga]\otimes_{\Z[\Ga_\sigma]}\Z[\Ga_\sigma]\to\Z[\Ga/\Ga_\sigma]\cong\Z[\Ga\cdot\sigma]$ on
$G_{s,0}$ and $0$ on $G_{s,t}$, $t\ge1$, such that:
\begin{enumerate}[label=(IM\arabic*)]
\item\label{IM1} $d\circ d=0$ and $\cmpI$ is a chain map;
\item\label{IM2} the sub-sum $\Isob_{\bullet}$ over the representatives in
  $\Sigma^{\partial}_{0}\sqcup\Sigma^{\partial}_{1}$ is a subcomplex, and
  $\cmpI(\Isob_{\bullet})\subseteq C_{\bullet}(\Xbdry)$.
\end{enumerate}
Write $\CI(\cring):=\Hom_{\Z[\Ga]}(\Iso_{\bullet},\VR)$ and likewise
$C^{\bullet}_{\Isob}(\cring)$; and $\cmpI^{\ast}:=\Hom_{\Z[\Ga]}(\cmpI,\VR)\colon
\Ceq(\cring)\to\CI(\cring)$. ($d^{(1)}$ lifts the simplicial boundary through the
transporters, and $d^{(2)}$ is the perturbation correction; $d^{(r)}$ for $r\ge3$ would
leave the columns $0\le s\le2$, $\Xbulk$ being $2$-dimensional.)
\end{definition}

\begin{proposition}[The isotropy-resolution model: existence, what it computes, and the one integral map]\label{prop:isomodel}
\PROVED{} Assume \ref{H4} and \ref{B1}--\ref{B3}, and let $n\ge0$ and $\cring$ be a
commutative ring.
\begin{enumerate}
\item \textbf{(Existence.)} An isotropy-resolution model exists; $d^{(1)}$ and $d^{(2)}$ are
  produced degree by degree by a terminating procedure (lifts along $\Z$-linear
  contracting homotopies of the augmented $F^{\sigma}$), boundary summands first.
\item \textbf{(A free resolution.)} $\Iso_{\bullet}$ is a complex of free $\Z[\Ga]$-modules
  of finite rank, $\Iso_{q}\cong\Z[\Ga]^{w_q}$ with $w_0=5$, $w_1=14$ and $w_q=19$ for
  $q\ge2$; $\cmpI$ is surjective in every degree and a quasi-isomorphism; hence
  $\aug\circ\cmpI\colon\Iso_{\bullet}\to\Z$ is a free resolution of $\Z$ over $\Z[\Ga]$.
  Likewise $\Isob_{q}\cong\Z[\Ga]^{w^{\partial}_q}$ with $w^{\partial}_0=3$ and
  $w^{\partial}_q=6$ for $q\ge1$, and $\cmpI\colon\Isob_{\bullet}\to C_{\bullet}(\Xbdry)$
  is a surjective quasi-isomorphism.
\item \textbf{(What it computes, over every ring, in every degree.)}
  $\CI(\cring)\cong\VR^{\,w_q}$ in degree $q$, $\CI(\Z)\otimes\cring=\CI(\cring)$, and
  $H^{q}(\CI(\cring))=H^{q}(\Ga;\VR)$ for every $q\ge0$; likewise
  $H^{q}(C^{\bullet}_{\Isob}(\cring))=H^{q}(\Gbdry;\VR)=H^{q}(\partial;\VR)$ of
  Definition~\ref{def:bdrygroup}. No invertibility hypothesis is used.
\item \textbf{(The comparison map.)} $\cmpI^{\ast}\colon\Ceq(\cring)\to\CI(\cring)$ is an
  injective cochain map in every degree, defined for every $\cring$ and compatible with
  base change from $\Z$: $\cmpI^{\ast}$ over $\cring$, composed with the reduction
  $\Ceq(\Z)\otimes\cring\to\Ceq(\cring)$ of Definition~\ref{def:ringinforce}, is
  $\cmpI^{\ast}$ over $\Z$ tensored with $\cring$ (as maps into $\CI(\Z)\otimes\cring=\CI(\cring)$). The
  cohomology of its cone is annihilated by $48$; so $\cmpI^{\ast}$ is a quasi-isomorphism
  whenever $6\in\cring^{\times}$, and its cone is $6$-power torsion over $\Z$. The same
  holds on the boundary with $4$ in place of $48$, so there $2\in\cring^{\times}$
  suffices. \emph{The integral line uses $\cmpI^{\ast}$ over $\Z[1/6]$ only, where the three
  objects of Definition~\ref{def:ringinforce} coincide.}
\item \textbf{(Relation to Definition~\ref{def:borelmodel}.)} $\Iso_{\bullet}$ and
  $F_{\bullet}\otimes_{\Z}C_{\bullet}(\Xbulk)$, for any free resolution $F_{\bullet}$ of $\Z$
  over $\Z[\Ga]$, are $\Z[\Ga]$-chain homotopy equivalent by maps over $C_{\bullet}(\Xbulk)$
  (compatibly with $\cmpI$ and with $\aug\otimes\id$). So $\Iso_{\bullet}$ is the
  isotropy-resolution model that Remark~\ref{rmk:mpcaveat}(4) names as a computational
  model of Definition~\ref{def:borelmodel}, compared with it here; and
  Remark~\ref{rmk:borelfree}, whose argument uses only that the source is free and the
  target acyclic, applies to it verbatim.
\end{enumerate}
\end{proposition}

\begin{proof}
\emph{The columns.} For a finite cyclic $\langle g\rangle$ of order $m$ the periodic
complex $F^{\sigma}\to\Z$ is exact: $\ker(g-1)=\Z\cdot N=\im N$ and $\ker N=\im(g-1)$ (the
augmentation ideal) in $\Z[\langle g\rangle]$, $N$ the norm. Being a bounded-below exact
complex of free abelian groups it is split exact, so it has a $\Z$-linear contracting
homotopy; closed formulas for one are standard. Since $\Z[\Ga]$ is free as a right
$\Z[\Ga_\sigma]$-module, $G_{s,\bullet}\to C_{s}(\Xbulk)=\bigoplus_{\sigma\in\Sigma_s}
\Z[\Ga/\Ga_\sigma]$ (the augmentation $\cmpI$) is a free resolution of the $\Z[\Ga]$-module
$C_s(\Xbulk)$, and so is its sub-sum over $\Sigma^{\partial}_{s}$ of $C_s(\Xbdry)$.

(1) \emph{$d^{(1)}$.} By the comparison theorem
\STD{\cite[Thm.~2.2.6]{weibel}}, the boundary $\partial\colon C_s(\Xbulk)\to
C_{s-1}(\Xbulk)$ lifts to a chain map $d^{(1)}\colon G_{s,\bullet}\to G_{s-1,\bullet}$ over
$\partial$; lifting first on the summands of $\Sigma^{\partial}_{1}$ into the resolution of
$C_{0}(\Xbdry)$ --- possible since $\partial C_1(\Xbdry)\subseteq C_0(\Xbdry)$ --- makes
$d^{(1)}$ carry boundary summands to boundary summands. The lift is computed on the free
generators by the contracting homotopies, so it is a terminating procedure in each
degree. \emph{$d^{(2)}$.} $d^{(1)}d^{(1)}\colon G_{2,\bullet}\to G_{0,\bullet}$ is a chain map
over $\partial\partial=0$, hence null-homotopic by the uniqueness half of the same
theorem: there is $d^{(2)}\colon G_{2,t}\to G_{0,t+1}$ with
$\partial^{F}d^{(2)}+d^{(2)}\partial^{F}=-d^{(1)}d^{(1)}$, again computed by the
contracting homotopies. \emph{$d^{2}=0$.} Sorting $d\circ d$ by the drop in $s$: drop $0$ is
$(d^{(0)})^{2}=0$; drop $1$ is $d^{(0)}d^{(1)}+d^{(1)}d^{(0)}=0$ because $d^{(1)}$ commutes
with $\partial^{F}$ and $d^{(0)}$ carries the sign $(-1)^{s}$; drop $2$ occurs only from
$s=2$, where $d^{(0)}=+\partial^{F}$ on both columns $2$ and $0$, and is
$d^{(1)}d^{(1)}+\partial^{F}d^{(2)}+d^{(2)}\partial^{F}=0$; drops $3$ and $4$ leave the
columns. \emph{$\cmpI$ is a chain map}: on $G_{s,0}$, $\cmpI d^{(1)}=\partial\,\cmpI$ is
the lifting property, $d^{(0)}$ vanishes and $d^{(2)}$ lands in $t=1$ where $\cmpI=0$; on
$G_{s,1}$, $\cmpI d^{(0)}=\pm\aug\,\partial^{F}=0$; elsewhere both sides vanish. \emph{The
boundary}: $d^{(2)}$ is defined on column $2$ only, which has no boundary summand, and
$d^{(0)}$, $d^{(1)}$ preserve boundary summands by construction; this is \ref{IM2}.

(2) Each $\Z[\Ga]\otimes_{\Z[\Ga_\sigma]}F^{\sigma}_{t}\cong\Z[\Ga]$, so $w_q$ counts the pairs
$(s,\sigma)$ with $\sigma\in\Sigma_s$, $s\le\min(q,2)$: $\abs{\Sigma_0}=5$,
$\abs{\Sigma_1}=9$, $\abs{\Sigma_2}=5$ by Lemma~\ref{lem:model7coincidences}(2), and
$3,3$ on the boundary (Remark~\ref{rmk:bdrymodel}). Surjectivity of $\cmpI$ is that of the
augmentations. For the quasi-isomorphism, filter both sides by $s$ ---
$F_{s}\Iso:=\bigoplus_{s'\le s}G_{s',\bullet}$ is a subcomplex since no component of $d$
raises $s$, and $F_{s}C:=C_{\le s}(\Xbulk)$ --- and note that $\cmpI$ respects the
filtrations and that on the graded pieces it is the augmentation
$G_{s,\bullet}[s]\to C_{s}(\Xbulk)[s]$, a quasi-isomorphism. The filtrations have length
three, so induction on $s$ with the long exact sequences of $0\to F_{s-1}\to F_{s}\to
\mathrm{gr}_s\to0$ and the five lemma gives the claim; the boundary is the same argument
on $s\le1$. $\Xbulk$ is contractible (\ref{H4}(ii)), so $\aug\circ\cmpI$ is a resolution.

(3) $\Hom_{\Z[\Ga]}(\Z[\Ga]\otimes_{\Z[\Ga_\sigma]}\Z[\Ga_\sigma],\VR)\cong\VR$, and
$\Iso_{q}$ being finite free, $\Hom_{\Z[\Ga]}(\Iso_{q},\Vcoef)\otimes\cring
=\Hom_{\Z[\Ga]}(\Iso_{q},\VR)$. By (2) the bulk complex computes
$\Ext^{\bullet}_{\Z[\Ga]}(\Z,\VR)=H^{\bullet}(\Ga;\VR)$. On the boundary, $\aug$ on each
component of $\Xbdry$ gives a quasi-isomorphism $C_{\bullet}(\Xbdry)\to\Z[\Ga/\Gbdry]$
(\ref{B1}, \ref{B2}: the components are contractible and permuted transitively, with
stabiliser $\Gbdry$), so $\Isob_{\bullet}$ is a free resolution of $\Z[\Ga/\Gbdry]
=\Z[\Ga]\otimes_{\Z[\Gbdry]}\Z$ and computes $H^{\bullet}(\Gbdry;\VR)$
\STD{\cite[Ch.~III.6]{brown}, Shapiro's lemma, as in the proof of
Claim~\ref{cl:amalgam}(2)}.

(4) Injectivity in each degree: $\cmpI$ is surjective and $\Hom_{\Z[\Ga]}(-,\VR)$ is left
exact. Base change: $\cmpI$ is a map of $\Z[\Ga]$-complexes and the reduction of
Definition~\ref{def:ringinforce} is induced by $\Vcoef\to\VR$. The cone: $\Hom_{\Z[\Ga]}(-,\VR)$
applied to the filtrations of (2) --- which are split in each degree as graded modules ---
gives filtrations of $\Ceq(\cring)$ and $\CI(\cring)$ of length three, respected by
$\cmpI^{\ast}$, whose graded map is, summand by summand,
$\VR^{\Ga_\sigma}\to\Hom_{\Z[\Ga_\sigma]}(F^{\sigma}_{\bullet},\VR)$, the inclusion of
$H^{0}(\Ga_\sigma;\VR)$ into a complex computing $H^{\bullet}(\Ga_\sigma;\VR)$
\STD{\cite[Ch.~III.6]{brown}}. Its cone therefore has cohomology $H^{t}(\Ga_\sigma;\VR)$,
$t\ge1$, which $\abs{\Ga_\sigma}$ annihilates \STD{\cite[Ch.~III.10.2]{brown}}: $12$ on
column $0$, $2$ on columns $1$ and $2$. A three-step filtration of a complex whose graded
pieces have cohomology annihilated by $a_0,a_1,a_2$ has cohomology annihilated by
$a_0a_1a_2=48$ (long exact sequences). The boundary has two columns, both with stabilisers
$\{\pm I\}$ (\ref{B3}): $2\cdot2=4$.

(5) $F_{\bullet}\otimes_{\Z}C_{\bullet}(\Xbulk)\to C_{\bullet}(\Xbulk)$, $\aug\otimes\id$, is
surjective, and a quasi-isomorphism because $F_{\bullet}\to\Z$ is a chain homotopy
equivalence of complexes of free abelian groups; its source is free
(Definition~\ref{def:borelmodel}(1)). Both it and $\cmpI$ are surjective
quasi-isomorphisms from bounded-below free complexes onto $C_{\bullet}(\Xbulk)$, so by
Lemma~\ref{lem:strictlifting} below each lifts through the other strictly, over
$C_{\bullet}(\Xbulk)$. Each composite minus the identity is a chain map into the kernel of
the corresponding surjection, which is acyclic; by the comparison theorem
\STD{\cite[Thm.~2.2.6]{weibel}} it is null-homotopic through a homotopy with values in
that kernel, so the composites are homotopic to the identities over $C_{\bullet}(\Xbulk)$.
\end{proof}

\begin{remark}[C-A5, re-scoped; OD-T1 answered ``not needed'']\label{rmk:cafive}
The extension plan's contract item C-A5 was ``the concrete \emph{refinement} of the
$7$-vertex triangulation (owner direction OD-T1) and a free $\Ga$-model satisfying
\ref{G1}--\ref{G4}, its boundary, its stabiliser inventory, its $\tau$/$\taub$/$\etapair$
tables' existence in every degree $\le q_{\max}+1$, and the truncation instance of
\ref{open:ell2} on it''. By owner ruling 6c of 2026-09-24 (the recommendation of the
integral-line brief's OD-8~(a)) it is re-scoped to \textbf{the free model on the unrefined
$\Xbulk$}, and is discharged by this chapter item by item: the model and its boundary are
Definition~\ref{def:isomodel} and Proposition~\ref{prop:isomodel}; \ref{G1}--\ref{G4} hold for
$(\Xbulk,\Xbdry)$ by Theorem~\ref{thm:specialisation}(4); the stabiliser inventory is
\ref{H4}(iv); the Hecke datum of the pair, in every degree, is
Proposition~\ref{prop:strictlift}, with $\etapair\equiv0$; the truncation instance is
Remark~\ref{rmk:truncdevice}. \textbf{OD-T1 is answered ``not needed''}, for the reasons
already written into this document: a $\Ga$-equivariant subdivision stays
$2$-dimensional, preserves every stabiliser and buys nothing at $\ell\in\{2,3\}$
(Remark~\ref{rmk:orbifoldrefine}(1),(4)), the degrees above $2$ come from the stabiliser
resolutions $F^{\sigma}$ and not from any triangulation, and $(\Xbulk,\Xbdry)$ already
satisfies \ref{G1}--\ref{G4} (Theorem~\ref{thm:specialisation}(4)). \PROVED{} Should a
refinement $Y'$ ever be wanted, it compares integrally: the subdivision is covered by
Proposition~\ref{prop:bredon}(3), so $\Hom_{\Z[\Ga]}(C_{\bullet}(\Xbulk),M)\simeq
\Hom_{\Z[\Ga]}(C_{\bullet}(Y'),M)$ for every $\Z[\Ga]$-module $M$, and an
isotropy-resolution model on $Y'$ is related to $\Iso_{\bullet}$ by
Proposition~\ref{prop:isomodel}(5)'s argument. The owner's direction is thereby not
overruled but answered: it was a design direction (Remark~\ref{rmk:orbifold}), and the
definitions of \S\ref{sec:orbifoldmodel} were phrased to admit a refinement; none is
required.
\end{remark}

\subsection{The descent: the fracture square and glueing over \texorpdfstring{$\Z$}{Z}}\label{sec:descent}

Write $\Zsix:=\Z_{2}\times\Z_{3}$ and $\Qsix:=\Zsix[1/6]=\Q_{2}\times\Q_{3}$. The object
the integral line computes is the group cohomology $H^{\bullet}(\Ga;\Vcoef)$, i.e.\ the
cohomology of $\CI(\Z)$, a complex of finitely generated free $\Z$-modules, bounded below
and \emph{not} bounded above (Remark~\ref{rmk:truncdevice}).

\begin{proposition}[The arithmetic fracture square]\label{prop:fracture}
\PROVED{} The sequence of abelian groups
\[
  0\longrightarrow\Z\xrightarrow{\ a\mapsto(a,a)\ }\Z[1/6]\oplus\Zsix
  \xrightarrow{\ (a,b)\mapsto a-b\ }\Qsix\longrightarrow0
\]
is exact. Consequently, for every cochain complex $M$ of abelian groups,
$0\to M\to M[1/6]\oplus(M\otimes\Zsix)\to M\otimes\Qsix\to0$ is an exact sequence of
complexes, and $M$ is quasi-isomorphic to the homotopy fibre product (the mapping cocone)
of $M[1/6]\to M\otimes\Qsix\leftarrow M\otimes\Zsix$ by the canonical map: the square is
a homotopy pullback.
\end{proposition}

\begin{proof}
Injectivity is clear. The middle: $a=b$ in $\Q_2$ and in $\Q_3$ with $a\in\Z[1/6]$,
$b\in\Z_2\times\Z_3$ forces $a\in\Z[1/6]\cap\Z_{(2)}\cap\Z_{(3)}=\Z$. Surjectivity: given
$(x,y)\in\Q_2\times\Q_3$, let $\{x\}_2\in\Z[1/2]$ and $\{y\}_3\in\Z[1/3]$ be their principal
parts (so $x-\{x\}_2\in\Z_2$, $y-\{y\}_3\in\Z_3$) and put $a:=\{x\}_2+\{y\}_3\in\Z[1/6]$; then
$a-x=\{y\}_3-(x-\{x\}_2)\in\Z_2$, since $\Z[1/3]\subset\Z_2$, and symmetrically
$a-y\in\Z_3$, so $(a,(a-x,a-y))\mapsto(x,y)$. $\Qsix$ is flat over $\Z$ (torsion-free
over a principal ideal domain), so $\mathrm{Tor}_1(M^{q},\Qsix)=0$ and the sequence stays
exact after $M^{q}\otimes-$. For a short exact sequence of complexes
$0\to M\to A\to C\to0$ the canonical map from $M$ to the cocone of $A\to C$ is a
quasi-isomorphism \STD{\cite[\S1.5]{weibel}: the cone of a map and the long exact
sequence}.
\end{proof}

\begin{proposition}[Glueing at the cohomology level is exact over $\Z$]\label{prop:cohglue}
\PROVED{} Let $M$ be a cochain complex of abelian groups. For every $q$ the sequence
\[
  0\to H^{q}(M)\to H^{q}(M)[1/6]\oplus\bigl(H^{q}(M)\otimes\Zsix\bigr)\to
  H^{q}(M)\otimes\Qsix\to0
\]
is exact, the middle terms being $H^{q}(M[1/6])$ and $H^{q}(M\otimes\Zsix)$; so $H^{q}(M)$
is the fibre product of $H^{q}(M[1/6])$ and $H^{q}(M\otimes\Zsix)$ over $H^{q}(M\otimes\Qsix)$.
The connecting maps of the long exact sequence of Proposition~\ref{prop:fracture} vanish.
For a chain map $f\colon M\to N$, $H^{q}(f)$ is determined by $H^{q}(f[1/6])$ and
$H^{q}(f\otimes\Zsix)$.
\end{proposition}

\begin{proof}
$\Z[1/6]$, $\Zsix$ and $\Qsix$ are flat, so $H^{q}(M\otimes A)=H^{q}(M)\otimes A$ for each.
Tensoring the exact sequence of Proposition~\ref{prop:fracture} with any abelian group
$G$ stays exact on the left, because $\mathrm{Tor}_1(G,\Qsix)=0$; with $G=H^{q}(M)$ this
is the displayed sequence, and it shows that the map $H^{q}(M)\to H^{q}(M[1/6])\oplus
H^{q}(M\otimes\Zsix)$ in the long exact sequence is injective, so every connecting map
vanishes. A short exact sequence $0\to G\to A\oplus B\xrightarrow{(a,b)\mapsto a-b}C\to0$
exhibits $G$ as $A\times_{C}B$. For $f$: $H^{q}(N)\to H^{q}(N[1/6])\oplus H^{q}(N\otimes\Zsix)$
is injective and $H^{q}(f)$ commutes with it.
\end{proof}

\noindent
Proposition~\ref{prop:cohglue} is [IB]'s Prop.~2.4 with a hypothesis removed: the proof
uses flatness only. Finite generation of $H^{q}(M)$ is what identifies $M\otimes\Zsix$ with
the $6$-adic completion and $H^{q}(M)\otimes\Z_\ell$ with $\varprojlim_{e}H^{q}(M\otimes\Z/\ell^{e})$,
which is how the local runs reach the $\Zsix$-corner (Remark~\ref{rmk:truncdevice}); the
exactness does not need it.

\begin{proposition}[The owner's descent, realised by one integral map]\label{prop:descent}
\PROVED{} Assume \ref{H4} and \ref{B1}--\ref{B3}, let $n\ge0$, and put
\[
  A:=\Ceq(\Z)[1/6]=\Ceq(\Z[1/6]),\qquad B:=\CI(\Z)\otimes\Zsix,\qquad C:=\CI(\Z)\otimes\Qsix,
\]
with $\alpha\colon A\to C$ the composite of $\cmpI^{\ast}[1/6]$ and the canonical
$\CI(\Z)[1/6]\to C$, and $\beta\colon B\to C$ the canonical map.
\begin{enumerate}
\item The mapping cocone of $\alpha-\beta\colon A\oplus B\to C$ is quasi-isomorphic to
  $\CI(\Z)$, hence computes $H^{\bullet}(\Ga;\Vcoef)$.
\item For every $q$,
  $H^{q}(\Ga;\Vcoef)\cong H^{q}(A)\times_{H^{q}(C)}H^{q}(B)$, the fibre product being taken
  along $H^{q}(\alpha)=H^{q}(\cmpI^{\ast})\otimes\Qsix$ and $H^{q}(\beta)$.
\end{enumerate}
\textbf{The only glueing datum is $\cmpI^{\ast}$}: the comparison isomorphisms over $\Q_2$ and
$\Q_3$ are its base changes, and nothing further is chosen or certified beyond
$\cmpI^{\ast}$ being a cochain map (Proposition~\ref{prop:isomodel}(4)).
\end{proposition}

\begin{proof}
$\Ceq(\Z)[1/6]=\Ceq(\Z[1/6])$ is Proposition~\ref{prop:ringinforce}(1). By
Proposition~\ref{prop:isomodel}(4) over $\Z[1/6]$, $\cmpI^{\ast}[1/6]\colon A\to\CI(\Z)[1/6]$ is
a quasi-isomorphism, so $(\cmpI^{\ast}[1/6],\id,\id)$ is a map from the diagram
$(A\to C\leftarrow B)$ to $(\CI(\Z)[1/6]\to C\leftarrow B)$ that is a quasi-isomorphism on
each corner; the induced map of cocones is a quasi-isomorphism by the five lemma on their
long exact sequences. The second cocone is quasi-isomorphic to $\CI(\Z)$ by
Proposition~\ref{prop:fracture}, and $H^{\bullet}(\CI(\Z))=H^{\bullet}(\Ga;\Vcoef)$ by
Proposition~\ref{prop:isomodel}(3). (2) is (1) together with Proposition~\ref{prop:cohglue}
applied to $\CI(\Z)$, transported along the quasi-isomorphism $\cmpI^{\ast}[1/6]$.
\end{proof}

\begin{proposition}[The bulk carries no coherence datum: the term \texorpdfstring{$\Hom^{-1}$}{Hom\^{}-1} vanishes]\label{prop:homvanish}
\PROVED{} Assume \ref{H4} and \ref{B1}--\ref{B3}, let $n\ge0$ and put
$N:=\CI(\Z)\otimes\Qsix$, a complex over the product of fields $\Qsix$.
\begin{enumerate}
\item $H^{\bullet}(N)=H^{\bullet}(\Ga;\Vcoef)\otimes\Qsix$ is concentrated in one degree: in
  degree $1$ for $n>0$ even, in degree $0$ for $n=0$, and it is zero for $n$ odd.
\item Hence $\Hom_{D(\Qsix)}(N,N[-i])=0$ for every $i\ge1$.
\end{enumerate}
\end{proposition}

\begin{proof}
(1) $\Qsix$ is flat, so $H^{q}(N)=H^{q}(\Ga;\Vcoef)\otimes\Qsix$, which vanishes when
$H^{q}(\Ga;\Vcoef)$ is torsion: for $q\ge2$ it is annihilated by $24$
(Claim~\ref{cl:amalgam}(4) with $\cring=\Z$); for $n$ odd it is annihilated by $2$ in every
degree ($-I$ is central, so conjugation by it acts trivially on $H^{\bullet}(\Ga;\Vcoef)$
\STD{\cite[Ch.~III, \S8]{brown}: an inner automorphism acts trivially on cohomology},
while it acts by $-1$ on $\Vcoef$ --- the argument of Remark~\ref{rmk:negI}); for $n>0$, $H^{0}(\Ga;\Vcoef)=0$ (the
argument of Claim~\ref{cl:bockstein}(1), which uses no hypothesis on $\ell$); for $n=0$,
$H^{1}(\Ga;\Z)=\Hom(\Ga,\Z)=0$, $\Ga$ being generated by the elements $S$, $U$ of finite
order. (2) A complex whose cohomology sits in a single degree $d$ is isomorphic in the
derived category to that cohomology placed in degree $d$ (by the two truncations at $d$),
so $\Hom_{D(\Qsix)}(N,N[-i])=\Ext^{-i}_{\Qsix}(H^{d},H^{d})=0$ for $i\ge1$
\STD{\cite[\S10.7]{weibel}: $\Hom_{D(A)}(X,Y[i])=\Ext^{i}_{A}(X,Y)$ for modules $X,Y$, which is
$0$ for $i<0$}; for $n$ odd, $N$ is acyclic.
\end{proof}

\noindent
\emph{Reading.} (2) is the vanishing of the one term, $\Hom^{-1}$ over $\Qsix$, in which
[IB] \S2.3(c) locates the coherence datum a descent could carry: with it zero, a
morphism of glued objects is determined by its two local components and relations among
local Hecke operators hold globally. At the cohomology level, which is all the bulk
needs, no derived category is required: that is Proposition~\ref{prop:cohglue}'s
morphism clause, since $H^{q}(\Ga;\Vcoef)$ injects into the product of the two local
cohomologies. \textbf{The datum is real for the pair}: there the corresponding term
contains $\Hom(\ker\rres^{1},\coker\rres^{0})\otimes\Q\neq0$, which is
Proposition~\ref{prop:coneambiguity}'s ambiguity, present in every model and removed on
the free model by the strict datum (Proposition~\ref{prop:strictlift}).

\begin{remark}[The truncation is a computation device, and the object is not perfect]\label{rmk:truncdevice}
\PROVED{} \textbf{What is glued is the untruncated $\CI(\Z)$}, in
Propositions~\ref{prop:fracture}--\ref{prop:homvanish}; none of them needs a bounded
complex. $\CI(\Z)$ is not perfect over $\Z$ in general: at $n=0$,
$H^{2j}(\Ga;\Z)\cong\Z/12$ for every $j\ge1$ (Lemma~\ref{lem:bulkperiodic}(1),(2)), so its
cohomology is unbounded. \textbf{What a local run builds is the brutal truncation}
$\sigma^{\le q_{\max}+1}\CI$ --- the instance, on $\Iso_{\bullet}$, of \ref{open:ell2}'s
box, which is the ``truncation instance'' C-A5 asked for: built in degrees
$\le q_{\max}+1$, reported in degrees $\le q_{\max}$ by the truncation lemma there,
\code{truncated = true}, $q_{\max}$ a set parameter (rider (i)). Brutal truncation
commutes strictly with base change, so the runs over $\Z/\ell^{e}$ are reductions of one
integral truncated complex, and by Lemma~\ref{lem:bulkperiodic}(2),(3) nothing is lost for
$q_{\max}\ge3$: every degree $\ge2$ is $H^{2}$ or $H^{3}$, Hecke-equivariantly
(\ref{open:ell2}'s non-normative recommendation $q_{\max}\ge4$, and $\ge5$ to see both
periodicity instances, stands). For the $\Zsix$-corner: $H^{q}(\CI(\Z))$ is finitely
generated, so $H^{q}(\CI(\Z)\otimes\Z_\ell)=H^{q}(\Ga;\Vcoef)\otimes\Z_\ell$ is determined by
the $\Z_{(\ell)}$-type of $H^{q}(\Ga;\Vcoef)$, which one certified run over $\Z/\ell^{e}$
yields (Proposition~\ref{prop:certifiede}); no inverse limit is computed.
\end{remark}

\subsection{The Hecke datum on the free model: a strict lift}\label{sec:strictlift}

\begin{lemma}[Strict lifting along a surjective quasi-isomorphism]\label{lem:strictlifting}
\PROVED{} Let $G$ be a group, $P_{\bullet}$ a complex of free $\Z[G]$-modules with
$P_q=0$ for $q<0$, $p\colon E_{\bullet}\to C_{\bullet}$ a surjective quasi-isomorphism of
complexes of $\Z[G]$-modules, $P'\subseteq P$ a subcomplex spanned by part of a basis in
each degree, $E'\subseteq E$ a subcomplex, and $f\colon P\to C$ a chain map. Given a chain map
$f'\colon P'\to E'$ with $pf'=f|_{P'}$, there is a chain map $\hat f\colon P\to E$ with
$p\hat f=f$ and $\hat f|_{P'}=f'$ --- a \emph{strict} lift, not one up to homotopy.
\end{lemma}

\begin{proof}
$\ker p$ is acyclic, by the long exact sequence of $0\to\ker p\to E\to C\to0$. Define
$\hat f$ on the basis elements outside $P'$ by induction on the degree. For such an $x$ in
degree $q$, choose $y_0\in E_q$ with $p(y_0)=f(x)$ ($p$ is onto). Then
$z:=\hat f(\partial x)-\partial y_0$ has $p(z)=f(\partial x)-\partial f(x)=0$ and
$\partial z=0$, so $z$ is a cycle of the acyclic $\ker p$ and $z=\partial k$ with
$k\in(\ker p)_q$ (for $q=0$ the same step applies; in the applications below
$E_{-1}=0$ and $z=0$). Put $\hat f(x):=y_0+k$; then $p\hat f(x)=f(x)$ and
$\partial\hat f(x)=\hat f(\partial x)$. Extend $G$-linearly.
\end{proof}

\begin{proposition}[The strict lift of the strict pair datum]\label{prop:strictlift}
\PROVED{} Assume \ref{H4} and \ref{B1}--\ref{B3}; let $g\in\De$,
$\Ga''=\Ga\cap g\Ga g^{-1}$, $\taub$ a boundary $g$-twisted chain datum over $\Z$
(Lemma~\ref{lem:taubdryexists}(1)), $\tau^{s}$ the strict datum of
Proposition~\ref{prop:strictpair} for it, and $(\Iso_{\bullet},\cmpI)$ an
isotropy-resolution model.
\begin{enumerate}
\item There is an additive map $\tau_{\Iso}\colon\Iso_{\bullet}\to\Iso_{\bullet}$, defined in
  every degree, satisfying (T1)--(T3) of Definition~\ref{def:tau} with $\Iso_{\bullet}$ in
  place of $C_{\bullet}(\Ku)$ and $\aug\circ\cmpI$ as augmentation, with
  \[
    \cmpI\circ\tau_{\Iso}=\tau^{s}\circ\cmpI\quad\text{strictly},\qquad
    \tau_{\Iso}(\Isob_{\bullet})\subseteq\Isob_{\bullet}.
  \]
  Its restriction $\tau^{\partial}_{\Iso}$ to $\Isob_{\bullet}$ satisfies
  $\cmpI\circ\tau^{\partial}_{\Iso}=\taub\circ\cmpI$, and $\etapair\equiv0$ is a pair homotopy
  for $(\tau_{\Iso},\tau^{\partial}_{\Iso})$: this is a Hecke datum of the pair on the free
  model $\Iso_{\bullet}$ --- Definition~\ref{def:borelmodel}(3) with $\Iso_{\bullet}$,
  $\Isob_{\bullet}$ in place of $\widehat{C}_{\bullet}(Y)$, $\widehat{C}_{\bullet}(\partial Y)$ ---
  in every degree.
\item Two such lifts are homotopic through $g$-twisted maps over $\Z$, for every $\ell$
  (Remark~\ref{rmk:borelfree}, whose argument applies by Proposition~\ref{prop:isomodel}(5)).
\item Writing $\Th^{s}_g$ for the coset transport of $\tau^{s}$ on $\Ceq(\cring)$ and
  $\Th_{\Iso,g}$ for that of $\tau_{\Iso}$ on $\CI(\cring)$ --- formula
  \eqref{eq:cosettransport} verbatim --- one has, over every commutative $\cring$,
  \[
    \cmpI^{\ast}\circ\Th^{s}_g=\Th_{\Iso,g}\circ\cmpI^{\ast}\quad\text{on the nose},
  \]
  and $\Th_{\Iso,g}$ induces $\cor\circ(\conj_g)_{\ast}\circ\res$ on
  $H^{\bullet}(\Ga;\VR)$ for every $\cring$, in every degree.
\item On cohomology over any $\cring$ with $6\in\cring^{\times}$, $\Th^{s}_g$ and the shipped
  $\Th_g$ (the coset transport of the canonical $\tau$ of Algorithm~\ref{alg:tau}) induce the
  same operator. \textbf{At $\ell\in\{2,3\}$ they are not compared}: on the Bredon object
  $\Ceq(\Z/\ell^{e})$ they may differ, and no statement of this chapter reads them there.
\end{enumerate}
\textbf{Existence only.} The effective, certified construction of $\tau_{\Iso}$ in degrees
$\le q_{\max}+1$ --- with a (T1)-bookkeeping over $\Ga''$ and the strict identity as
certificates --- is the integral line's key risk (K1 of [IB]) and a code matter; nothing
here bounds its integer sizes or asserts it is built. (v1.10.0: the statement is unchanged; what
hecke-extensions slot~10 built and certified, for $g=\mathrm{diag}(1,p)$, $p\in\{2,3,5\}$, in
degrees $\le q_{\max}+1=6$, and the integer sizes it measured, are
Remark~\ref{rmk:strictliftbuilt}, which is evidence and bounds nothing either.)
\end{proposition}

\begin{proof}
Let $\Iso^{(c)}$, $C^{(c)}_{\bullet}(\Xbulk)$ be the same complexes with the twisted
$\Ga''$-action $h\ast x:=(g^{-1}hg)\cdot x$ ($g^{-1}hg\in\Ga$ by
Proposition~\ref{prop:cosets}(3)); (T1) for a map $\Iso\to\Iso$ is $\Z[\Ga'']$-linearity
$\Iso\to\Iso^{(c)}$. As $[\Ga:\Ga'']<\infty$, $\Iso_{q}$ is free over $\Z[\Ga'']$ on a basis
that restricts to one of $\Isob_{q}$, and $\cmpI\colon\Iso^{(c)}\to C^{(c)}(\Xbulk)$ is still a
$\Ga''$-equivariant surjective quasi-isomorphism, as is its boundary restriction. (1) First
apply Lemma~\ref{lem:strictlifting} to $P=\Isob$, $E=(\Isob)^{(c)}$, $C=C^{(c)}(\Xbdry)$,
$f=\taub\circ\cmpI$ (with $P'=0$), obtaining $\tau^{\partial}_{\Iso}$; then to $P=\Iso$,
$E=\Iso^{(c)}$, $C=C^{(c)}(\Xbulk)$, $f=\tau^{s}\circ\cmpI$, $P'=\Isob$, $E'=(\Isob)^{(c)}$ and
$f'=\tau^{\partial}_{\Iso}$, which is admissible because
$\cmpI\tau^{\partial}_{\Iso}=\taub\cmpI=\tau^{s}\cmpI$ on $\Isob$ ($\tau^{s}\circ j=j\circ\taub$).
(T2) holds because the lift is a chain map; (T3):
$\aug\cmpI\tau_{\Iso}=\aug\tau^{s}\cmpI=\aug\cmpI$ by (T3) for $\tau^{s}$. The pair
identity with $\etapair=0$ is $\tau_{\Iso}\circ j=j\circ\tau^{\partial}_{\Iso}$, which holds by
construction. (2) The difference of two lifts is a (T1)-chain map from the free
$\Z[\Ga'']$-complex $\Iso$ into $\ker\cmpI$, which is acyclic; the fundamental lemma
\STD{\cite[Thm.~2.2.6]{weibel}} makes it null-homotopic through (T1)-maps. (3) For
$\phi\in\Hom_{\Z[\Ga]}(C_{\bullet}(\Xbulk),\VR)$ and $x\in\Iso$,
$(\Th_{\Iso,g}\cmpI^{\ast}\phi)(x)=\sum_{j}\rhon(\delta_j)\,\phi(\cmpI\,\tau_{\Iso}(s_jx))
=\sum_j\rhon(\delta_j)\,\phi(\tau^{s}(s_j\,\cmpI x))=(\cmpI^{\ast}\Th^{s}_g\phi)(x)$,
using $\cmpI\tau_{\Iso}=\tau^{s}\cmpI$ and the $\Ga$-equivariance of $\cmpI$. The coset
transport of Lemma~\ref{lem:cosettransport}(1),(2) applies verbatim on $\Iso$ (its proof used
(T1) and $\partial(s_jy)=s_j\partial y$ only), and on the free resolution $\Iso$ the
operator is $\cor\circ(\conj_g)_{\ast}\circ\res$ by Remark~\ref{rmk:borelfree} (with
Proposition~\ref{prop:isomodel}(5)). (4) $\tau^{s}$ and the canonical $\tau$ are
$g$-twisted chain maps over $\Z$, homotopic through twisted maps when
$6\in\cring^{\times}$ (Lemma~\ref{lem:tauexists}(2)), so they induce the same operator on
$\Heq{\bullet}$ (Theorem~\ref{thm:welldef}(b)).
\end{proof}

\noindent
Lifting the \emph{strict} datum, boundary summands first, is what makes the lifted pair
homotopy vanish; lifting the shipped $\tau$ instead would work the same way for the bulk
and leave a boundary datum to certify ([IB] \S3.2 item~2, the fallback).

\begin{remark}[The strict lift as built (hecke-extensions slot 10)]\label{rmk:strictliftbuilt}
\COMP\textsuperscript{r} (v1.10.0.) For $g=\mathrm{diag}(1,p)$, $p\in\{2,3,5\}$ ---
$\Ga''=\Ga_0(p)$, $m=p+1$ right cosets $\Ga''s_j$ with $s_1=I$ --- the map $\tau_{\Iso}$ of
Proposition~\ref{prop:strictlift}(1) is built in degrees $0,\ldots,q_{\max}+1=6$ on the model
$\Iso_{\bullet}$ of Definition~\ref{def:isomodel}, and certified exactly over $\Z[\Ga'']$
(\file{Integral/Model/StrictLift}, hecke-extensions slot-10 WP-10; registered
\code{HeckeStrictLiftValidation}, \code{standard}, $1824$ checks, $0$ failing;
\file{docs/implementation/slot10/WP-10\_report.md}). The fallback of the paragraph above was not
needed.
\begin{enumerate}
\item \emph{Construction} --- Lemma~\ref{lem:strictlifting}'s induction, degree by degree,
  boundary summands first, on the free $\Z[\Ga'']$-basis $\{s_je_c\}$, extended by (T1). The
  preimage $y_0$ is \emph{read off}, not solved for: at a cell with $t=0$, $\cmpI$ is the
  augmentation of free generators, so the stored chain $\tau^{s}(s_j\,\cmpI x)$ is rewritten term
  by term as an element of $\Iso_{q}$; at $t\ge1$, $y_0=0$. The cycle $z$ of $\ker\cmpI$ is
  contracted \emph{by descent through the columns}: $\ker\cmpI$ is the sum of the augmented
  periodic columns $F^{\sigma}_{\bullet}$, whose closed-form $\Z$-linear contracting homotopy $h$
  (the one Definition~\ref{def:isomodel}'s differentials are built from) gives $k_s=(-1)^{s}h(z_s)$
  on the top column, and the rest is pushed into lower columns. No linear system is solved; the
  result is deterministic.
\item \emph{Certified}, re-derived from the stored values and trusting nothing from the
  construction: (T1) on the Schreier generators of $\Ga''$ (on a free $\Z[\Ga'']$-basis (T1) is
  the \emph{defining} property of the extension, so this certifies its bookkeeping); (T2) and (T3);
  the strict identity $\cmpI\circ\tau_{\Iso}=\tau^{s}\circ\cmpI$, compared as canonical chains of
  $\Xbulk$ through a routine independent of the construction's conversion;
  $\tau_{\Iso}(\Isob_{\bullet})\subseteq\Isob_{\bullet}$; and
  $\cmpI\circ\tau^{\partial}_{\Iso}=\taub\circ\cmpI$ against the boundary table's own chains --- so
  $\etapair\equiv0$ is a pair homotopy, as (1) states. A perturbed entry is refused.
\item \emph{(3) in use.} $\Th_{\Iso,g}$ is a cochain map over $\Fl$, $\Z/\ell^{k}$ and at the
  $64$-bit word ($\Z/2^{40}$, $\Z/3^{30}$), on $165$ and $24$ cells
  (Lemma~\ref{lem:cosettransport}(2) on the model). The identity
  $\cmpI^{\ast}\circ\Th^{s}_g=\Th_{\Iso,g}\circ\cmpI^{\ast}$ holds entrywise on $\Ceq(\cring)$ in
  degrees $q\le2$, $n\le20$, on eight rings among those with a shipped $\Th^{s}_g$ (v1.10.1): $\Fl$,
  $\ell\in\{2,3,5,7\}$, and $\Z/\ell^{k}$, $\ell\in\{5,7\}$, $k\in\{2,3\}$ ($359$ cells; v1.10.1, in
  place of v1.10.0's arithmetic slip ``$349$'') --- a consistency check of the
  transport code, since the identity is a theorem. Over $\Z/\ell^{e}$ at $\ell\in\{2,3\}$ no
  $\Th^{s}_g$ exists in the build (Remark~\ref{rmk:bredonhecke}) and the identity is not separately
  checked; it holds there by the proof of (3), a formal computation valid over every commutative
  $\cring$, whose three inputs are certified over $\Z$ in item~2 or hold by construction: the strict
  identity, (T1) of $\tau_{\Iso}$ (which makes the coset transport well typed,
  Lemma~\ref{lem:cosettransport}(1)) and the $\Ga$-equivariance of $\cmpI$.
\item \emph{(4) in use.} At $\ell\ge5$ the operator $\Th_{\Iso,g}$ induces on $H^{1}$, transported
  by $\cmpI^{\ast}$, equals the shipped $\Amat{1}(T_p)$ --- built on the \emph{canonical} $\tau$
  of Algorithm~\ref{alg:tau}, not on $\tau^{s}$ --- entrywise on $26$ cells, so (4) is exercised
  through Lemma~\ref{lem:tauexists}(2); and $\mathrm{char}(\Th_{\Iso,g}\,|\,H^{1})$, with $H^{1}$
  taken on the model alone, equals the Eichler--Shimura values of the design brief (its oracle O5)
  modulo $\ell$ at $n\in\{10,\ldots,20\}$ ($21$ cells).
\item \emph{Integer sizes} --- the first measurement of [IB]'s key risk K1: largest coefficient
  $3$, $2$, $2$ and largest matrix entry $p$ for $p=2,3,5$; at most $26$ terms in one value; $20$,
  $29$ and $64$ distinct group elements; the term counts stationary with period~$2$ from degree~$2$
  on, as the periodic columns suggest. Nothing approaches the word; no \code{\_\_int128} and no
  big integer occurs (Remark~\ref{rmk:integralscope}(1)).
\end{enumerate}
\textbf{What this remark does not say.} It bounds nothing: Proposition~\ref{prop:strictlift}'s
``existence only'' stands as a statement, and no effectivity is claimed for any other $g$ --- in
particular not for the double cosets of determinant $p^{2}$ that $T_{p^{2}}$ and $\diam$ need
(Theorem~\ref{thm:relations}) --- nor in degrees above $6$
(Remark~\ref{rmk:integralscope}(2)).
\emph{Later builds} (v1.23.0-pre.28; this remark stays the record of what slot~10 built): hecke-extensions slot~24, milestone D4, builds the same lift for $g=\mathrm{diag}(1,p)$ at
\textbf{every prime} $p$ requested, on demand, in degrees $\le6$, by item~1's procedure, and certifies it by item~2's checks, the boundary items included, at run time before first
use (Proposition~\ref{prop:descentpairhecke}(1), its code map (f)); measured to $p=127$ (largest coefficient $\le3$, largest matrix entry $p$; \file{docs/implementation/slot24/DESCENT4\_report.md}
\S4). That too is evidence and bounds nothing; effectivity at a given $p$ is the run-time certificate, refused on failure or at the cap (the full list of refusals: Proposition~\ref{prop:descentpairhecke}(5), v1.23.0-pre.29).
\end{remark}

\begin{remark}[The Hecke operator on the Bredon object at $\ell\in\{2,3\}$ is not compared]\label{rmk:bredonhecke}
(v1.10.0; owner ruling~4~(A) of the hecke-extensions slot-10 decisions, 2026-09-25.) Assume
\ref{H4}, $\ell\in\{2,3\}$ and $\cring=\Z/\ell^{e}$, $e\ge1$. On the Bredon object $\Ceq(\cring)$
of Definition~\ref{def:ringinforce}(3) two Hecke objects could be read: \emph{(a)} the coset
transport \eqref{eq:cosettransport} of an admissible integral $\tau$ --- the canonical one of
Algorithm~\ref{alg:tau}, or the strict $\tau^{s}$ --- a cochain map on $\Ceq(\cring)$ (v1.10.1: the
canonical $\tau$ is integral in practice, not by Algorithm~\ref{alg:tau}'s construction, which records a
mod-$\ell$ triangle-fill fallback in its step~(3); since v1.1.0 it vanishes in degree~$2$,
Remark~\ref{rmk:iotatausym}, consequence~(a), so no fill is solved); and
\emph{(b)} the operator it induces on $\Heq{\bullet}$. This document says nothing about (b) beyond
Proposition~\ref{prop:strictlift}(4): homotopy uniqueness of admissible $\tau$
(Lemma~\ref{lem:tauexists}(2)) needs $6\in\cring^{\times}$, which fails here, so the operator in (b)
is \textbf{not shown to be independent of the choice of $\tau$}, and this document \textbf{compares
none of these operators there} --- neither $\Th^{s}_g$ with the canonical $\Th_g$, nor either with
$\Th_{\Iso,g}$, whose cohomology is $H^{\bullet}(\Ga;\VR)$ (Proposition~\ref{prop:isomodel}(3)), a
different object at $\ell\in\{2,3\}$ (Proposition~\ref{prop:ringinforce}). (v1.23.0-pre.25, the C-D4 referee's M2: at level $N$ with $\nu_3(N)=0$ and over
$\Z_{(3)}$-algebras the Bredon object does compute $H^{\bullet}(\Ga_0(N);\VR)$, and Proposition~\ref{prop:strictliftN}(6)(b) identifies the operator of a datum over $\Z$
certified by (T1)--(T3); the shipped canonical $\tau$ is not so certified, and \S\ref{sec:descenthecke} reads neither --- its $3$-local route is
Proposition~\ref{prop:strictliftN}(5). The level-$1$ statement of this remark is untouched: $\SL_2(\Z)$ has elements of order $3$.) (v1.10.1: the one
comparison it does record there is of another pair, the legacy and the canonical $\tau$, which induce
the same $\Amat{q}$ at every $\ell\in\{2,3\}$ golden point --- an observation, not a theorem,
Remark~\ref{rmk:tauexistsscope}(1).) This is the \textbf{absence of a proof of
independence, not a proof of dependence}: no two admissible $\tau$ with different operators on the
Bredon object are known, and none is claimed. What the build does accordingly: at $e=1$ the
shipped $\Fl$ pipeline computes (a) and (b) for the \emph{canonical} $\tau$ --- the canonical $\tau$
of Algorithm~\ref{alg:tau} since v1.1.0 (Remark~\ref{rmk:iotatausym}; v1.10.1, in place of v1.10.0's
``unchanged since v1.0.0'', which is false of $\tau$ and of (a)); its operator (b) did not move in the
v1.1.0 re-baseline and equals the legacy $\tau$'s at every $\ell\in\{2,3\}$ golden point,
Remark~\ref{rmk:tauexistsscope}(1) --- an observation, not a theorem. It is a definite object, the canonical datum's operator, which no statement of this document
identifies with anything $\tau$-independent at $\ell\in\{2,3\}$; at $e\ge2$ neither is built, by the
companion specification's rule (\Espec{sec:ring}~(b3)) --- (a) is not needed, since the integral line
reads $T_p$ through $\Th_{\Iso,g}$ on $\CI(\cring)$ (Remark~\ref{rmk:strictliftbuilt}), and (b) is
refused by name. Should two admissible $\tau$ with different operators in (b) ever be exhibited, a
later version records that here in place of this remark's reading.
\end{remark}

\subsection{Bad primes, the certified exponent, and the characteristic polynomial over \texorpdfstring{$\Z$}{Z}}\label{sec:badprimes}

\begin{proposition}[The bad primes form a finite, closed-form list]\label{prop:badprimes}
\PROVED{} Let $n\ge0$ and $\Vcoef=\Sym^{n}(\Z^{2})$.
\begin{enumerate}
\item For $n>0$ even, every prime $\ell$ at which $H^{\bullet}(\Ga;\Vcoef)$ has
  $\ell$-torsion lies in
  \[
    \{2,3\}\cup\bigl\{\,5\le\ell\le n-1:\ \#\{(a,b)\in\Z_{\ge0}^{2}:a(\ell+1)+b(\ell^{2}-\ell)=n\}>0\,\bigr\},
  \]
  and at each $\ell\ge5$ of the list the $\ell$-primary torsion lies in $H^{1}$ and has
  exactly that many cyclic summands.
\item For $n$ odd, $2\cdot H^{q}(\Ga;\Vcoef)=0$ for every $q$; for $n=0$ the torsion primes
  are $2$ and $3$.
\end{enumerate}
\end{proposition}

\begin{proof}
(1) For $q\ge2$, $24\cdot H^{q}(\Ga;\Vcoef)=0$ (Claim~\ref{cl:amalgam}(4)); $H^{0}=0$
(Claim~\ref{cl:bockstein}(1)). For $\ell\ge5$, $H^{1}(\Ga;\Vcoef)[\ell]\cong H^{0}(\Ga;\Mn)$
(Claim~\ref{cl:bockstein}(2)), whose dimension is the displayed count
(Claim~\ref{cl:dimH1}(2), proved in place since v1.23.0-pre.15); the number of
cyclic $\ell$-primary summands of a finitely generated group is the dimension of its
$\ell$-torsion over $\Fl$, and the count vanishes for $n<\ell+1$. (2) The centrality
argument in the proof of Proposition~\ref{prop:homvanish}(1); at $n=0$, $H^{1}(\Ga;\Z)=0$ and
Claim~\ref{cl:amalgam}(4) confine the torsion to $2$ and $3$, and $H^{2}(\Ga;\Z)\cong\Z/12$
(Lemma~\ref{lem:bulkperiodic}(1)) exhibits both.
\end{proof}

\noindent
Exponents are \emph{not} bounded by this list. At $\ell\ge5$ they grow along
$n=\ell^{j}m$ (Remark~\ref{rmk:bocksteinexp}, whose lower-bound argument uses no hypothesis
on $\ell$); at $2$ and $3$ the integral-line probe (a computation outside the tree,
\COMP\textsuperscript{a}, [IB] Appendix~A) finds $\Z/2^{j}$ in $H^{1}$ at $n=2^{j}$,
$j\le4$, and $\Z/9$ at $n=12,18$. No fixed exponent is safe at any bad prime, which is
what the next proposition is for.

\begin{proposition}[The certified-exponent rule, at every prime, $2$ and $3$ included]\label{prop:certifiede}
\PROVED{} Let $C^{\bullet}$ be a complex of finitely generated free $\Z$-modules, $q_{\max}\ge0$,
$\ell$ a prime and $e\ge1$. For an integer matrix $D$ let $\ell^{a_1},\dots,\ell^{a_r}$
($r=\rank_{\Q}D$) be its non-zero elementary divisors over $\Z_{(\ell)}$, and let $r_e(D)$ be
the number of diagonal entries $\ell^{a}$ with $a<e$ in the Smith form of $D\bmod\ell^{e}$
over $\Z/\ell^{e}$.
\begin{enumerate}
\item $r_e(D)=\#\{t:a_t<e\}\le\rank_{\Q}D$, with equality iff every $a_t<e$.
\item For each $q$, $\mathrm{tors}\,H^{q}(C)_{(\ell)}\cong\bigoplus_{t}\Z/\ell^{a_t}$ over the
  non-unit elementary divisors $\ell^{a_t}$ of $d^{q-1}$, and
  $\rank H^{q}(C)=\rank C^{q}-\rank_{\Q}d^{q}-\rank_{\Q}d^{q-1}$.
\item \textbf{The rule.} If $r_e(d^{q})=\rank_{\Q}d^{q}$ for every $q\le q_{\max}$ (with
  $d^{-1}=0$), then for every $q\le q_{\max}$ the $\Z_{(\ell)}$-type of $H^{q}(C)$ is read off
  exactly from the Smith forms over $\Z/\ell^{e}$, every exponent being $<e$, and
  $H^{q}(C\otimes\Z/\ell^{f})\cong H^{q}(C)\otimes\Z/\ell^{f}\oplus
  \mathrm{Tor}(H^{q+1}(C),\Z/\ell^{f})$ is thereby known for \emph{every} $f\ge1$. If the
  equality fails for some $q\le q_{\max}$, $e$ is too small: a summand $\Z/\ell^{a}$ with
  $a\ge e$ is indistinguishable from a free one mod $\ell^{e}$.
\item \textbf{Where $\rank_{\Q}$ comes from.} $\rank_{\Q}d^{q}=\rank_{\F_{\ell'}}d^{q}$ at any
  prime $\ell'$ not dividing an elementary divisor of $d^{q}$, i.e.\ at which
  $H^{q+1}(C)$ has no $\ell'$-torsion. For $C=\CI(\Z)$ and for $C=\Ceq(\Z)$, every prime
  $\ell'\ge5$ with $\ell'>n+1$ qualifies (Proposition~\ref{prop:badprimes}, and
  Proposition~\ref{prop:isomodel}(4) for $\Ceq(\Z)$, whose cohomology agrees with group
  cohomology after inverting $6$).
\item \textbf{Scope.} The rule applies to $\sigma^{\le q_{\max}+1}\CI(\Z)$ at every $\ell$ (its
  differentials $d^{q}$, $q\le q_{\max}$, are those of $\CI(\Z)$), and to $\Ceq(\Z)$ at every
  $\ell$ --- at $\ell\ge5$, where $\Ceq(\Z/\ell^{e})=\Ceq(\Z)\otimes\Z/\ell^{e}$, it is [L]'s
  stopping rule. It does \textbf{not} apply to the model complex $\Ceq(\Z/\ell^{e})$ at
  $\ell\in\{2,3\}$: that complex is not the reduction of an integral complex
  (Proposition~\ref{prop:ringinforce}(2)), carries no integral lattice to certify, and is
  checked against the tree complex \eqref{eq:treemodel} over the same ring instead.
\end{enumerate}
\end{proposition}

\begin{proof}
(1) Over the local principal ideal domain $\Z_{(\ell)}$, $D=U\,\mathrm{diag}(\ell^{a_t},0)\,V$
with $U,V$ invertible; reducing mod $\ell^{e}$ keeps $U,V$ invertible, so the Smith form of
$D\bmod\ell^{e}$ has diagonal $\ell^{\min(a_t,e)}$, where $\ell^{e}=0$. (2) If
$x\in C^{q}$ and $mx\in\im d^{q-1}$ for some $m\neq0$, then $m\,d^{q}x=0$, so $d^{q}x=0$
because $C^{q+1}$ is torsion-free; hence $\mathrm{tors}\,H^{q}(C)=\mathrm{tors}\,\coker d^{q-1}$,
whose $\ell$-part is read from the Smith form. The rank count is additivity over $\Q$.
(3) By (1), the hypothesis says every elementary divisor of $d^{q-1}$ and of $d^{q}$ is
$<\ell^{e}$ for $q\le q_{\max}$, so (2) gives the types of $H^{q}$ and $H^{q+1}$'s torsion; the
displayed decomposition is the universal coefficient theorem for a complex of free
modules over a principal ideal domain \STD{\cite[\S3.6]{weibel}}. If $a\ge e$ for some
$\ell^{a}$, (1) counts it as $0$. (4) Reduction mod $\ell'$ preserves the rank exactly when
$\ell'$ divides no elementary divisor, and those are the torsion orders of
$\coker d^{q}$, i.e.\ of $H^{q+1}(C)$, by (2). For $\CI(\Z)$ the torsion of
$H^{\bullet}(\Ga;\Vcoef)$ lies over the primes of Proposition~\ref{prop:badprimes}, all
$\le\max(3,n-1)$; for $\Ceq(\Z)$ the cone of $\cmpI^{\ast}$ is $6$-power torsion, so its
cohomology has the same torsion primes $\ge5$. (5) is (1)--(3) read against
Proposition~\ref{prop:ringinforce}.
\end{proof}

\begin{remark}[The rule, differential by differential]\label{rmk:certifiedeperdiff}
\PROVED{} (v1.9.0.) In the setting of Proposition~\ref{prop:certifiede}, the hypothesis of~(3)
may be checked one differential at a time, each at its own level. If $r_{e_q}(d^{q})=\rank_{\Q}d^{q}$
for some $e_q\ge1$, then every elementary divisor $\ell^{a_t}$ of $d^{q}$ has $a_t<e_q$, and
$\mathrm{tors}\,H^{q+1}(C)_{(\ell)}$ is read off exactly from the Smith form of $d^{q}\bmod\ell^{e_q}$.
The $\Z_{(\ell)}$-type of $H^{q}(C)$ needs $e_{q-1}$ (its torsion comes from $d^{q-1}$) and the
two ranks $\rank_{\Q}d^{q-1}$, $\rank_{\Q}d^{q}$; the displayed decomposition of
$H^{q}(C\otimes\Z/\ell^{f})$ in~(3) needs $e_{q-1}$ and $e_q$. Clause~(3) with one $e$ is the
uniform case $e=\max_{q\le q_{\max}}e_q$. \emph{Proof.} (1) applied to $D=d^{q}$ at $e=e_q$ gives
$a_t<e_q$ for every $t$; (2) reads $\mathrm{tors}\,H^{q+1}(C)_{(\ell)}$ from the non-unit
elementary divisors of $d^{q}$ alone, and $H^{q}$'s rank from the two ranks over $\Q$; the
universal coefficient decomposition involves $H^{q}(C)$ and $\mathrm{Tor}(H^{q+1}(C),\Z/\ell^{f})$
and nothing else. No step couples $d^{q}$ to $d^{q'}$ for $q'\neq q$. \emph{In use}: the
integral groups of line~2's milestone M0 are assembled with $K$ certified per differential
(\code{HeckeFreeModelValidation}, hecke-extensions slot-9 WP-11, which also checks that every
level below each certified one is refused, and the displayed decomposition at one level above the
largest); Definition~\ref{def:ringsemantics}(3)(b)'s ``$e$ certified per degree'' is this reading.
A clarification of~(3), not a correction: no statement or proof of the proposition changes
(\file{docs/implementation/slot9/WP-11\_report.md} finding~2, confirmed by the slot's
independent review, WP-14, \file{docs/implementation/slot9/WP-14\_review.md}).
\end{remark}

\begin{proposition}[The characteristic polynomial over \texorpdfstring{$\Z$}{Z}, by residues]\label{prop:crtcharpoly}
\PROVED{} Let $n\ge0$, $p$ a prime, $g=\mathrm{diag}(1,p)$, and write
$\Hfr{1}:=H^{1}(\Ga;\Vcoef)/\mathrm{tors}$, on which the double-coset operator
$\cor\circ(\conj_g)_{\ast}\circ\res$ acts (it preserves torsion); let $\chi_{p}\in\Z[x]$ be
its characteristic polynomial. Let $\ell'\ge5$ be a prime,
$\ell'\ne p$, at which $H^{1}(\Ga;\Vcoef)$ has no $\ell'$-torsion --- every $\ell'>n+1$
qualifies (Proposition~\ref{prop:badprimes}).
\begin{enumerate}
\item Reduction $\Vcoef\to\Mn$ induces an isomorphism, equivariant for that operator,
  $\Hfr{1}\otimes\F_{\ell'}\cong H^{1}(\Ga;\Mn)$ over $\F_{\ell'}$, and
  $\chi_{p}\bmod\ell'$ is the characteristic polynomial of the shipped matrix
  $\Amat{1}(T_p)$ over $\F_{\ell'}$.
\item Hence, if $P\in\Z[x]$ is monic of the same degree, $\mathcal S$ a finite set of such
  primes, and $\prod_{\ell'\in\mathcal S}\ell'>2B$ where $B$ bounds the absolute values of the
  coefficients of both $P$ and $\chi_{p}$, then $P=\chi_{p}$ iff
  $P\equiv\mathrm{charpoly}\,\Amat{1}(T_p)\pmod{\ell'}$ for every $\ell'\in\mathcal S$. No
  integer beyond a word is formed: the residues of $P$ are reduced digit-wise.
\end{enumerate}
The bound $B$ for $\chi_{p}$ is an input, not a statement of this document; the one the
integral line uses is Deligne's bound on the eigenvalues through Eichler--Shimura
(Remark~\ref{rmk:integralevidence}).
\end{proposition}

\begin{proof}
(1) The long exact sequence of $0\to\Vcoef\xrightarrow{\ell'}\Vcoef\to\Mn\to0$ gives
$0\to H^{1}(\Ga;\Vcoef)\otimes\F_{\ell'}\to H^{1}(\Ga;\Mn)\to H^{2}(\Ga;\Vcoef)[\ell']\to0$
(the argument of Claim~\ref{cl:bockstein}(2)), and the right-hand term is $0$ because $24$
kills $H^{2}$ (Claim~\ref{cl:amalgam}(4)) and $\ell'\ge5$; without $\ell'$-torsion,
$H^{1}(\Ga;\Vcoef)\otimes\F_{\ell'}=\Hfr{1}\otimes\F_{\ell'}$. Every map is induced by a map of
$\Z[\De]$-modules, so it commutes with $\cor\circ(\conj_g)_{\ast}\circ\res$
(Claim~\ref{cl:bockstein}(3)'s argument). At $6\in\F_{\ell'}^{\times}$ the shipped matrix is
that operator on $\Heq{1}\cong H^{1}(\Ga;\Mn)$ (Remark~\ref{rmk:orbifold},
Theorem~\ref{thm:classical}, \ref{H1} as $p\ne\ell'$). (2) $P-\chi_{p}$ has coefficients of
absolute value $\le2B<\prod\ell'$ and is divisible by every $\ell'\in\mathcal S$.
\end{proof}

\subsection{What each coefficient ring names}\label{sec:ringsemantics}

\begin{definition}[The mathematical meaning of the coefficient ring]\label{def:ringsemantics}
Assume \ref{H4}. The setting that selects the coefficient ring --- its spelling, default,
reserved values and refusal rule --- is \Espec{sec:ring}'s; the exponent it writes $k$ is
this document's $e$ (\S\ref{sec:notation-primes}, \ref{levelO1}). For a cell $(n,\ell)$, the ring
in force names the following object.
\begin{enumerate}
\item \textbf{$\Fl$} (the default). $\Ceq(\Fl)$ and its cohomology --- $H^{\bullet}(\Ga;\Mn)$
  when $\ell\ge5$ (Remark~\ref{rmk:orbifold}), the Bredon cohomology of
  Proposition~\ref{prop:bredon} over $\Fl$ when $\ell\in\{2,3\}$. Unchanged since v1.0.0.
\item \textbf{$\Z/\ell^{e}$}, $e\ge2$. $\Ceq(\Z/\ell^{e})$, whose cohomology is reported as
  a $\Z/\ell^{e}$-module (lengths and invariant factors), never as a dimension. At
  $\ell\ge5$ it is $\Ceq(\Z)\otimes\Z/\ell^{e}$ and computes $H^{\bullet}(\Ga;\Vcoef\otimes\Z/\ell^{e})$
  (Proposition~\ref{prop:ringinforce}(1), Remark~\ref{rmk:orbifold} with
  $6\in\cring^{\times}$), with the Hecke action on the nose from the integral tables
  (Proposition~\ref{prop:coneintegral}(2),(3)). At $\ell\in\{2,3\}$ it names the Bredon
  object over $\Z/\ell^{e}$ --- which is \emph{not} a reduction of $\Ceq(\Z)$ and not group
  cohomology. \emph{What this build computes of it there} (v1.9.0; \Espec{sec:ring}~(b1),
  (b3), (c1$'$)): its \textbf{lengths and invariant factors} in degrees $q=0,1,2$, through a
  flat-coordinate module route --- the cochain groups of $\Ceq(\Z/\ell^{e})$ are not free there
  (the invariants computation of Proposition~\ref{prop:ringinforce}(2)'s proof, read over
  $\Z/2^{e}$; \Espec{sec:ring}~(c1$'$) gives it; v1.9.1, the citation through v1.9.0 named
  Proposition~\ref{prop:ringinforce}(2),(3), which do not state it), so no basis of them is used --- measured by
  \code{HeckeBredonTreeRingValidation} against the tree complex \eqref{eq:treemodel} over the
  same ring (Proposition~\ref{prop:bredon}(2)); the certified $K$-mode, which is undefined there
  (Proposition~\ref{prop:certifiede}(5)), the Hecke operator and the reductions are
  \textbf{refused} by name, as build capabilities (\code{hecke::NotImplementedInThisBuild}).
  (Through v1.8.1 this clause ended ``and the implementation of this build refuses it there'',
  stale from hecke-extensions slot-9 WP-3's merge, \texttt{b53836f}, until this version.)
\item \textbf{$\Z$} (reserved through v1.23.0-pre.36; \textbf{un-reserved for the integral descent only} at v1.23.0-pre.37, item~7 of \S\ref{sec:recordschema}:
  the value is honoured by the descent record \code{integral\_descent\_record/v1}, its console verb and its campaign cells, and \textbf{refused by name} by every
  other reader of the ring --- a build capability, \code{hecke::NotImplementedInThisBuild}, never a silent fall-back to $\Fl$; the general $\Z/m$ stays reserved). The object is \textbf{group cohomology}
  $H^{q}(\Ga;\Vcoef)$ in degrees $q\le q_{\max}$, i.e.\ $H^{q}(\CI(\Z))$ --- never the
  integral model complex $\Ceq(\Z)$, which at $\ell\in\{2,3\}$ is a different object
  (Proposition~\ref{prop:ringinforce}) and at odd $n$ is zero while group cohomology is not
  (Proposition~\ref{prop:ringinforce}(3)). It decomposes as
  $H^{q}(\Ga;\Vcoef)\cong\Z^{r_q}\oplus\bigoplus_{\ell}T^{q}_{\ell}$ over the primes of
  Proposition~\ref{prop:badprimes} (at level $N$: $H^{q}(\Ga_0(N);\Vcoef)$ over the primes of Proposition~\ref{prop:badprimesN}, the object of the descent
  record, item~7 of \S\ref{sec:recordschema}; v1.23.0-pre.38, referee m10). What the integral line computes, and from what:
  \begin{enumerate}
  \item $\Z[1/6]$-part and ranks: from $\Ceq$ runs at good primes and at the Dickson
    primes $\ell\ge5$ over $\Z/\ell^{e}$, with $e$ certified
    (Propositions~\ref{prop:ringinforce}(1), \ref{prop:certifiede}); the ranks equally from
    $\CI(\Z)$ at a good prime, which Proposition~\ref{prop:certifiede}(4) covers as well (v1.9.0:
    line~2's milestone M0 takes them there and cross-checks the $\Ceq$ runs,
    \code{HeckeFreeModelValidation});
  \item $2$- and $3$-parts: from $\sigma^{\le q_{\max}+1}\CI(\Z)\otimes\Z/\ell^{e}$ with $e$
    certified per degree (Proposition~\ref{prop:certifiede}(5), Remark~\ref{rmk:truncdevice};
    ``per degree'' read differential by differential, Remark~\ref{rmk:certifiedeperdiff});
  \item the group, glued by Proposition~\ref{prop:descent}, the only glueing datum being
    $\cmpI^{\ast}$;
  \item the Hecke action: by $\Th_{\Iso,g}$ (Proposition~\ref{prop:strictlift}); on each
    $T^{q}_{\ell}$ from the local run (the $\ell$-primary part is canonical and stable
    under $\cor\circ(\conj_g)_{\ast}\circ\res$); on $\Hfr{1}$ as the characteristic
    polynomial over $\Z$ by residues
    (Proposition~\ref{prop:crtcharpoly}), \textbf{basis-free} (v1.10.0: $\Th_{\Iso,g}$ is built
    for $T_2$, $T_3$, $T_5$ and read on $H^{q}(\CI(\cring))$ basis-free at small $n$,
    Remarks~\ref{rmk:strictliftbuilt} and \ref{rmk:integralevidence}; the glued action over $\Z$
    waits for the descent assembly; v1.23.0-pre.38: built at level $N$, \S\ref{sec:descenthecke}, merged at \texttt{dc8261a3}, and carried by the
    record of item~7 of \S\ref{sec:recordschema}).
  \end{enumerate}
  \textbf{Not named by this ring}: integer Hecke matrices on $\Hfr{1}$ in any basis, and the
  global lattice class of $\Hfr{1}$ as a Hecke module --- owner decision OD-11~(b), confirmed
  2026-09-24, puts the integral line basis-free first; a coordinate rule is a later
  revision's (OD-11~(a), small $n$ only). Nor the $\pm$ split at $\ell=2$ (it needs
  $2\in\cring^{\times}$, Obligation~\ref{obl:iota}), nor Poincar\'e duality, which over $\Z$
  holds only after inverting $6$.
\end{enumerate}
Which record carries which of these, and under which setting, is [E]'s: under OD-10~(a)
(confirmed 2026-09-24) the record keys that say ``$\Z$'', the truncation and which object is
reported are emitted only under the value naming $\Z$, so that no golden record moves. (v1.23.0-pre.37: so it is built --- the record's keys
\code{coefficient\_ring = "Z"} and \code{object = "group"} exist only in \code{integral\_descent\_record/v1}, item~7 of \S\ref{sec:recordschema}.)
\end{definition}

\begin{remark}[What is tested today (slots 8 to 11)]\label{rmk:integralevidence}
\COMP\textsuperscript{r} Nine registered validations test statements of this chapter on the
shipped tree --- two since slot~8, three since slot~9, three since slot~10, one since slot~11
(v1.11.0; the title read ``(slots 8, 9 and 10)'' and the count ``eight'' at v1.10.0--v1.11.0-pre.2,
``(slots 8 and 9)'' and ``five'' at v1.9.0--v1.9.1, ``(slot 8)'' and ``two'' through
v1.8.1); none upgrades any of them.
\begin{itemize}
\item \code{HeckeBredonTreeValidation} (\file{Integral/HeckeBredonTreeValidation.cpp},
  slot-8 WP-8, \code{fast}): Proposition~\ref{prop:bredon}(2) at $\cring=\Fl$ ---
  $\dim\Heq{q}$ of the shipped pipeline equals the tree complex \eqref{eq:treemodel},
  computed by an independent $\rhon$ and rank routine, on $82$ cells
  ($\ell\in\{2,3\}$, $0\le n\le40$, odd $n$ included), after first checking the tree
  complex against the sixteen golden cells at $\ell\in\{2,3\}$; $0$ mismatches
  (\file{docs/implementation/slot8/WP-8\_report.md}).
\item \code{HeckeIntegralCharpolyValidation} (\file{Integral/}, slot-8 WP-7,
  \code{standard}): Proposition~\ref{prop:crtcharpoly} in use --- $\chi_{2}$ for
  $n\in\{10,12,14,16,18,20\}$ and $\chi_{3}$ at $n=10$ are certified equal over $\Z$ to the
  Eichler--Shimura values of the integral-line brief (its oracle O5, e.g.\
  $(x-2049)(x+24)^{2}$ at $n=10$, $p=2$) with no integer beyond a word, each cell with one
  held-out prime; the bound $B$ is the validation's input (Deligne's, as a bit length);
  at $n\in\{64,128\}$ the residue vectors and the bound are recorded, never compared with
  an oracle ($n=128$ under the read-only switch \code{HECKE\_INTEGRAL\_CAMPAIGN=1}); $322$
  checks, $0$ failing (\file{docs/implementation/slot8/WP-7\_report.md}).
\item \code{HeckeBredonTreeRingValidation} (\file{Integral/}, slot-9 WP-3, \code{standard}):
  Proposition~\ref{prop:bredon}(2) at $\cring=\Z/\ell^{k}$, $\ell\in\{2,3\}$, $k\in\{2,3,4\}$ ---
  the lengths and invariant factors of $\Heq{q}$ computed by the shipped pipeline over the ring
  (Definition~\ref{def:ringsemantics}(2)) equal the tree complex \eqref{eq:treemodel} over the
  same ring, frozen from the design brief's probe (lattices over $\Z$ and the probe's own Smith
  form, no modular elimination), on $96$ cells (even $n\le20$, odd $n\le9$, $H^{2}=0$
  throughout), with controls at $k=1$ against the shipped $\Fl$ route and at $\ell=5$ against the
  $\ell\ge5$ module route; $1394$ checks, $0$ failing
  (\file{docs/implementation/slot9/WP-3\_report.md}).
\item \code{HeckeIsotropyModelValidation} (\file{Integral/}, slot-9 WP-10, \code{standard}):
  Definition~\ref{def:isomodel} and Proposition~\ref{prop:isomodel}(1)--(3) as built --- the
  ranks $5,14,19,\ldots$ and $3,6,6,\ldots$, \ref{IM1} and \ref{IM2} exactly over $\Z[\Ga]$,
  $\cmpI$ a chain map; $\cmpI^{\ast}$ a cochain map, injective, and an isomorphism on
  $H^{0}$ and $H^{1}$ against the shipped pipeline at $\ell\ge5$ (Proposition~\ref{prop:isomodel}(4),
  where $6\in\cring^{\times}$); $\dim H^{q}(\CI(\F_\ell))$ and the lengths over $\Z/\ell^{k}$ at
  $\ell\in\{2,3\}$ against the design brief's Appendix~A through universal coefficients; the
  boundary model against $H^{\bullet}(\Gbdry)$; $1032$ checks, $0$ failing
  (\file{docs/implementation/slot9/WP-10\_report.md}).
\item \code{HeckeFreeModelValidation} (\file{Integral/}, slot-9 WP-11, \code{standard}): line~2's
  milestone M0 --- $H^{q}(\Ga;\Sym^{n}\Z^{2})$ for $n\le20$, $q\le5$, assembled as
  Definition~\ref{def:ringsemantics}(3)(a)--(b) says (ranks at the good prime $10007$ by
  Proposition~\ref{prop:certifiede}(4); the $2$- and $3$-parts on $\CI$ by
  Proposition~\ref{prop:certifiede} per differential, Remark~\ref{rmk:certifiedeperdiff}; the
  Dickson-prime parts from the shipped $K$-mode, cross-checked on $\CI$), equals Appendix~A
  exactly on $126$ of $126$ $(n,q)$ cells; the cone of $\cmpI^{\ast}\otimes\Z/\ell^{K}$ is
  acyclic through degree $5$ at the ten Dickson cells; $H^{2}\cong H^{4}$, $H^{3}\cong H^{5}$
  (\ref{lem:bulkperiodic}); at odd $n$ the integral model complex vanishes while the group does
  not (Proposition~\ref{prop:ringinforce}(3)); $803$ checks, $0$ failing
  (\file{docs/implementation/slot9/WP-11\_report.md}; v1.10.1: $801$ since slot-10 WP-8, whose arm
  (R) cites the two registered names it used to re-run in process, and whose test~(6) refuses on
  the $7$-vertex complex's own number of degrees). (v1.10.0: the design brief's tests (4)
  and (5) for it, which it declared no slot for, are hosted in \code{HeckeIntegralActionValidation}
  below, and its test (3) in \code{HeckePeriodicityClassValidation}.)
\item \code{HeckeStrictLiftValidation} (\file{Integral/}, slot-10 WP-10, \code{standard}):
  Proposition~\ref{prop:strictlift}(1), (3) and (4) as built, for $T_2$, $T_3$, $T_5$ in degrees
  $\le6$ --- the certificates, the strict compatibility on the rings with a shipped $\Th^{s}_g$,
  the operator on $H^{1}$ against the shipped canonical-$\tau$ operator at $\ell\ge5$, O5 on the
  model, and a negative control; Remark~\ref{rmk:strictliftbuilt}; $1824$ checks, $0$ failing
  (\file{docs/implementation/slot10/WP-10\_report.md}).
\item \code{HeckePeriodicityClassValidation} (\file{Integral/}, slot-10 WP-13, \code{standard}):
  Lemma~\ref{lem:bulkperiodic}(1)--(2) as built. (1): $e$ is a $2$-cocycle in
  $\Hom_{\Z[\Ga]}(\Iso_{2},\Z)$, the Bockstein of the character $\zeta_+$ in the convention of the
  lemma's Step~0; $\kappa_C(\res_C e)=\chi_C$ on the stabilisers, with $\chi=\zeta_+$ re-checked on $224$
  elements of finite order (v1.10.1: v1.10.0 attributed the $224$-element check to the first
  identity); $[e]$ has order $12$ and generates $H^{2}(\Ga;\Z)\cong\Z/12$ ($12e$ a coboundary, $6e$
  and $4e$ not). $E=e\cup-$ is a chain map $\Iso_{\bullet+2}\to\Iso_{\bullet}$ with $\aug\circ
  E_0=e$, certified exactly over $\Z[\Ga]$. (2), clause by clause, on $126$ ring cells ($n\le20$
  over $\Fl$ and $\Z/\ell^{K}$ at $\ell\in\{2,3\}$, and over $\F_5$, $\F_7$): an isomorphism for
  $q\ge2$; for $q=1$ surjective, with kernel equal to $\partial(N_{-I}V)$ as submodules (inclusion
  plus equal lengths), injective at odd $n$; no $q=0$ clause. The connecting map is realised
  through the Bass--Serre tree as the first-variable connecting map of \eqref{eq:treeses}, which
  agrees with the coefficient form $(\ast)$ of the lemma's proof up to one global sign under the
  identification that proof names; the kernel clause reads only its image, which the sign does
  not move, and an independent Mayer--Vietoris length computed in $V$ alone,
  $\mathrm{length}\bigl((N_{-I}V+V^{S}+V^{U})/(V^{S}+V^{U})\bigr)$, agrees on all $126$ cells (the
  slot's review, finding m-7); $2482$ checks, $0$ failing (v1.10.1: $2483$ since slot-10 WP-8's
  oriented-disc arm, in which the fill of $x$ reaches a second round)
  (\file{docs/implementation/slot10/WP-13\_report.md}).
\item \code{HeckeIntegralActionValidation} (\file{Integral/}, slot-10 WP-11, \code{standard}):
  line~2's milestone M1(a) --- $T_p$ as $\Th_{\Iso,g}$, $p\in\{2,3,5\}$, read on
  $H^{q}(\CI(\cring))$ with no integral basis, at $n\in\{10,12,16\}$ (arms on torsion and on
  universal coefficients also for every $n\le20$). The Eichler--Shimura identities of the design
  brief (O6) on $H^{1}(\Fl)=\Hfr{1}\otimes\Fl$ at eight good primes, with O5's at $n=12,16$; the
  $691$-congruence module of $H^{1}_{+}\otimes\Z_{691}$ at $n=10$ is $\Z/691$ on \emph{both} lines
  (the E8 route over $\Z/691^{k}$ and $\Th_{\Iso,g}$ on $\CI\otimes\Z/691^{k}$, $k\in\{2,3\}$), by two
  readings that assume $T_2$ diagonalisable over $\Q_{691}$ on $\Hfr{1}$ --- Eichler--Shimura, of
  which the identity O6 checked modulo $691^{k}$ is the computed evidence (the slot's review,
  finding m-15); on the $2$- and $3$-primary torsion in degrees $\le5$,
  Lemma~\ref{lem:bulkperiodic}(3)(b) for $q=0,\ldots,3$, $T_pT_{p'}=T_{p'}T_p$ on cohomology
  (non-vacuous: the cochain maps do not commute), $\mathrm{char}\,\Amat{3}(T_p)\mid
  \mathrm{char}\,\Amat{1}(T_p)$, and the torsion modules read twice (by reduction and by the
  Bockstein) with exponents equal to M0's; \textbf{not} tested: the recursion of
  Theorem~\ref{thm:relations}(1), which needs a lifted datum for determinant $p^{2}$
  (Remark~\ref{rmk:strictliftbuilt}, last paragraph; the review's finding M-3); the local type of
  $\Hfr{1}\otimes\Z_{\ell}$ at $\ell\in\{2,3\}$ recorded (split at $n=10$ and $16$; at $(16,3)$ with
  no $T_3$ oracle, finding m-17), with O5 and O6 holding on $\Hfr{1}\otimes\Z/\ell^{4}$ there; the
  universal-coefficient lengths at every precision below the run's (a consistency check);
  $4601$ checks, $0$ failing (\file{docs/implementation/slot10/WP-11\_report.md}).
\item \code{HeckeStrictLiftRelationsValidation} (\file{Integral/}, slot-11 WP-11, \code{standard};
  v1.11.0): Proposition~\ref{prop:strictlift}(1) built and certified exactly over $\Z[\Ga'']$ for
  $g=\mathrm{diag}(1,p^{2})$, $\mathrm{diag}(1,p^{3})$ and $pI$ ($p\in\{2,3\}$; $p=5$ for
  $\mathrm{diag}(1,25)$ and $5I$), with $\tau=\id$ of Proposition~\ref{prop:diamond}'s proof as the
  datum at $pI$; the transports are cochain maps over fields, $\Z/\ell^{k}$ and $64$-bit $\Z/\ell^{K}$;
  $\Th_{\Iso,pI}$ induces $p^{\,n}$ on every $H^{q}(\CI(\cring))$, $q\le5$, $\ell\in\{2,3\}$ and $p=\ell$
  included; and \textbf{Proposition~\ref{prop:relationsmodel}(2) for $\nu\in\{1,2\}$} holds on $H^{q}$,
  $q\le5$, basis-free, over $\Z/2^{8}$, $\Z/3^{5}$, $\F_2$, $\F_3$ ($p=\ell$ included), $\F_5$, $\F_7$,
  $\F_{691}$, with the full operators assembled as in the proof of \eqref{eq:fullassembly}
  ($T_{p^{2}}=\Th_{\mathrm{diag}(1,p^{2})}+\Th_{pI}$, $T_{p^{3}}=\Th_{\mathrm{diag}(1,p^{3})}+p^{\,n}\Th_{p}$);
  non-vacuous (the identities fail on cochains in $684$ of $684$ cases), with \ref{err:e1}'s primitive
  operator and \ref{rmk:relcoeff}'s coefficient as negative controls, refused exactly where $p^{\,n}$
  acts non-trivially; the $\ell\ge5$ consistency with the shipped primitive $T_{p^{\nu}}$ through
  $\cmpI^{\ast}$ ($\nu\in\{2,3\}$, \ref{prop:strictlift}(4)); $3619$ checks, $0$ failing
  (\file{docs/implementation/slot11/WP-11\_report.md}). The item ``\textbf{not} tested: the recursion
  of Theorem~\ref{thm:relations}(1)'' of the bullet above is thereby answered on the free model; on
  $\Ceq$ at $\ell\in\{2,3\}$ nothing is compared (Remark~\ref{rmk:bredonhecke}).
\end{itemize}
\end{remark}

\begin{remark}[What this chapter does not say]\label{rmk:integralscope}
Five things a reader might expect here and will not find.
\begin{enumerate}
\item \textbf{No arbitrary precision.} Owner decision OD-3~(c) (confirmed 2026-09-24):
  Theorem~\ref{thm:backend}(3), the box of \S\ref{sec:fpfilter} and
  Obligation~\ref{obl:backend}(iv) are unchanged. Integer outputs of the integral line are
  residue vectors with a certified bound (Proposition~\ref{prop:crtcharpoly}); a decimal
  or rational reconstruction happens outside the build. Re-scoping the ban is a separate
  amendment that this version does not make.
\item \textbf{No effectivity claim for $\tau_{\Iso}$}, and no bound on its integer sizes
  (Proposition~\ref{prop:strictlift}, last paragraph). $\Iso_{\bullet}$ itself is produced by
  a terminating procedure (Proposition~\ref{prop:isomodel}(1)). (v1.10.0: for
  $g=\mathrm{diag}(1,p)$, $p\in\{2,3,5\}$, in degrees $\le6$, $\tau_{\Iso}$ is built and its sizes
  are measured, Remark~\ref{rmk:strictliftbuilt}; that is evidence, not a bound, and for every
  other $g$ and degree this item stands as written. v1.23.0-pre.24: the bulk lift for $g=\mathrm{diag}(1,p)$, every $p$, is \emph{produced} by a terminating
  procedure, Proposition~\ref{prop:strictliftN}(2); its sizes stay unbounded. v1.23.0-pre.28: the pair lift for $g=\mathrm{diag}(1,p)$ is built on demand at every prime $p$ requested,
  degrees $\le6$, and certified at run time, refused on a failed certificate or at the cap, Proposition~\ref{prop:descentpairhecke}(1) (every refusal: its (5), v1.23.0-pre.29); measured to $p=127$, still no bound.)
\item \textbf{No record schema.} The keys are [E]'s (Definition~\ref{def:ringsemantics},
  last paragraph); \S\ref{sec:recordschema} is not amended here.
\item \textbf{No higher level and no families}, and no divided powers: $\Vcoef$ is
  $\Sym^{n}(\Z^{2})$ (OD-4~(a)); $\Dn$ and levels via Shapiro need their own chapter. (v1.23.0-pre.14, the level half of the plan's C-D7:
  levels via Shapiro, for the groups and the descent, are Propositions~\ref{prop:descentN} and~\ref{prop:badprimesN}; families, divided powers and
  the level-$N$ Hecke action over $\Z$ stand as written. v1.23.0-pre.18: restriction, trace, conjugation, the level cone and its long exact sequence over $\Z$,
  and the Hecke squares on cohomology, are Proposition~\ref{prop:descentmaps}. v1.23.0-pre.24: the level-$N$ operator $T_p$, $p\nmid N$, over $\Z$ at cochain level, and
  its invariants, are \S\ref{sec:descenthecke}. v1.23.0-pre.26: the pair and its cone at level $N$ over $\Z$, the boundary by cusps, and $T_p$ on them, are
  \S\ref{sec:descentpair}.)
\item \textbf{Nothing about the pair and the cone over $\Z$} beyond
  Proposition~\ref{prop:strictlift}(1)'s datum; the cone over $\Zsix$ under \ref{open:ell2}
  is a later revision's. (v1.23.0-pre.26, the pair half of the plan's C-D7: the pair models over $\Z$, $\mathsf O_\ell$ and $\Z[1/6]$, the long exact sequence over $\Z$
  glued from them, the boundary in closed form, compact support, the certified exponent under \ref{open:ell2}'s truncation, and the cone with the strict pair datum and
  $T_p$, $p\nmid N$, are \S\ref{sec:descentpair}; the tree cusp pair stands as written, the plan's D5c. v1.23.0-pre.33: the tree cusp pair --- the cusp model, the axis map,
  the tree cone, its comparisons with $(\Iso,\Isob)$ and $T_p$ by transport --- is \S\ref{sec:treepair}, contract first; v1.23.0-pre.36: built, audited and merged at \texttt{4064e164}.)
\end{enumerate}
\end{remark}

\paragraph{Code map.}
Implemented: \file{Integral/HeckeBredonTreeValidation} (Proposition~\ref{prop:bredon}(2) at
$\cring=\Fl$) and \file{Integral/IntegralCharpolyCrt} with
\code{HeckeIntegralCharpolyValidation} (Proposition~\ref{prop:crtcharpoly}), both slot~8.
\textbf{Not implemented}: $\Iso_{\bullet}$, $\Isob_{\bullet}$, $\cmpI$, $\tau_{\Iso}$, the
descent assembly and the value $\Z$ of the ring (\S\ref{sec:notimplemented}). The model
complex over $\Z/\ell^{e}$ at $\ell\ge5$ is E8's (\Espec{sec:ring}).
(v1.10.0: this paragraph was the v1.8.0 text and stale from v1.9.0 on; \S\ref{sec:codemap}'s
row is the current one. Implemented since v1.9.0: \file{Integral/Model/IsotropyModel}
($\Iso_{\bullet}$, $\Isob_{\bullet}$, $\cmpI$) and \file{Integral/Model/IntegralGroups} (the
groups at $n\le20$); since v1.10.0: \file{Integral/Model/StrictLift} ($\tau_{\Iso}$,
$\Th_{\Iso,g}$; Remark~\ref{rmk:strictliftbuilt}), \file{Integral/Model/PeriodicityClass} ($e$,
$E_q$; Lemma~\ref{lem:bulkperiodic}) and \file{Integral/Model/IntegralHecke} (the Hecke action on
$H^{\bullet}(\CI(\cring))$, basis-free). Not implemented: the descent assembly and the value $\Z$ of
the ring.)

\paragraph{Formalization note.}
Propositions~\ref{prop:fracture} and \ref{prop:cohglue} are elementary facts about flat
abelian groups and are cheap and worth formalizing first: they are what makes the descent
exact. Lemma~\ref{lem:strictlifting} is a finite induction. Proposition~\ref{prop:bredon}
needs a functor category and the comparison theorem there, and
Proposition~\ref{prop:isomodel}(1) the assembly of resolutions over a complex; both are
standard homological algebra with no arithmetic.


\subsection{The Hecke relations on the free model}\label{sec:relationsmodel}

\noindent\emph{New at v1.11.0-pre.1} (hecke-extensions slot-11 WP-2; the slot-10 review's finding~6,
and the precondition of line~2's milestone M-3). Theorem~\ref{thm:relations} is a statement about
$\Heq{\bullet}(K,\Rel;L)$ and carries $6\in\cring^{\times}$, because on the shipped complex the
operator is identified with the classical one only then (Theorem~\ref{thm:classical},
Remark~\ref{rmk:orbifold}). On the free model $\CI(\cring)$ that identification holds over
\emph{every} ring (Proposition~\ref{prop:strictlift}(3)), and the relations follow there with no
hypothesis on $\ell$. This subsection states that. It changes nothing about
Theorem~\ref{thm:relations}, whose hypothesis stays, and nothing about the Bredon object
$\Ceq(\cring)$ at $\ell\in\{2,3\}$ (Remark~\ref{rmk:bredonhecke}).

\begin{proposition}[The Hecke relations on $H^{\bullet}(\CI(\cring))$, at every prime]\label{prop:relationsmodel}
Assume \ref{H4} and \ref{B1}--\ref{B3}; let $\cring$ be \emph{any} commutative ring --- no hypothesis
on $6$, so $\cring=\Fl$ and $\cring=\Z/\ell^{e}$ with $\ell\in\{2,3\}$ are included --- let $n\ge0$,
$\VR=\Sym^{n}(\Z^{2})\otimes\cring$, and let $(\Iso_{\bullet},\cmpI)$ be an isotropy-resolution model
(Definition~\ref{def:isomodel}). For $\delta\in\De$ choose a strict lift $\tau_{\Iso}$ for $g=\delta$
(Proposition~\ref{prop:strictlift}(1)) and write $T^{\Iso}_{\Ga\delta\Ga}$ for the map that its coset
transport $\Th_{\Iso,\delta}$ induces on $H^{q}(\CI(\cring))$. Put, for a prime $p$ and $\nu\ge0$,
\[
  T^{\Iso}_{p^{\nu}}:=\sum_{\substack{\Ga\delta\Ga\subseteq\De\\ \det\delta=p^{\nu}}}T^{\Iso}_{\Ga\delta\Ga},
  \qquad T^{\Iso}_{p^{0}}:=\id,
\]
the full operator of \eqref{eq:fullop} read on the free model, and $T^{\Iso}_p:=T^{\Iso}_{p^{1}}$.
\begin{enumerate}
\item \PROVED{} $T^{\Iso}_{\Ga\delta\Ga}$ does not depend on the lift $\tau_{\Iso}$, on the coset
  representatives, or on $\delta$ inside $\Ga\delta\Ga$, and it equals the classical double-coset
  operator $\cor\circ(\conj_\delta)_{\ast}\circ\res$ of \eqref{eq:corconjres} on
  $H^{q}(\Ga;\VR)$, in every degree $q$. In particular $T^{\Iso}_{\Ga\,pI\,\Ga}=p^{\,n}\id$.
\item \PROVED{} (given the citations below) For every degree $q$, all primes $p,p'$ and all
  $\nu\ge1$, on $H^{q}(\CI(\cring))$:
  \[
    T^{\Iso}_pT^{\Iso}_{p'}=T^{\Iso}_{p'}T^{\Iso}_p,\qquad
    T^{\Iso}_{p^{\nu+1}}=T^{\Iso}_p\,T^{\Iso}_{p^{\nu}}-p^{\,n+1}\,T^{\Iso}_{p^{\nu-1}},
  \]
  and $T^{\Iso}_{p^{2}}=T^{\Iso}_{\Ga\,\mathrm{diag}(1,p^{2})\,\Ga}+p^{\,n}\id$. With
  \STD{\cite[Thm.~3.24 and \S3.2]{shimura} for the identities in the abstract Hecke ring
  $\mathcal{H}(\De,\Ga)$; \cite[Thm.~1]{rhie_whaples} and \cite[\S1.1, p.~194]{ash_stevens} for the fact that
  $\Ga\delta\Ga\mapsto\cor\circ(\conj_\delta)_{\ast}\circ\res$ is a ring homomorphism
  $\mathcal{H}(\De,\Ga)\to\End H^{q}(\Ga;W)$ for \emph{every} $\Z[\De]$-module $W$ and every $q$,
  with no invertibility hypothesis. The locator into \cite{rhie_whaples} is not verified against the text (rule~\ref{A2}): no copy is in
  \file{docs/literature/}; \cite[\S1.1, p.~194]{ash_stevens} is verified (slot-17 R17-4, v1.17.0-pre.5) --- it \emph{states} the action as
  ``well known'' for right modules, left modules by the $S^{-1}$ device, without proof and without an invertibility hypothesis}.
\item \textbf{What this does not say.} Nothing about $\Heq{\bullet}$ at $\ell\in\{2,3\}$, where
  Theorem~\ref{thm:relations} keeps its hypothesis and the operator of the Bredon object is not
  compared with anything (Remark~\ref{rmk:bredonhecke}); and nothing at cochain level, where the
  recursion does not hold (the cochain maps of different double cosets are built from different
  lifts). It says nothing about effectivity either: a lift for $\mathrm{diag}(1,p^{2})$ and $pI$ is
  asserted to exist (Proposition~\ref{prop:strictlift}(1)), not to be built.
\end{enumerate}
\end{proposition}

\begin{proof}
(1) By Proposition~\ref{prop:strictlift}(3), $\Th_{\Iso,\delta}$ induces
$\cor\circ(\conj_\delta)_{\ast}\circ\res$ on $H^{\bullet}(\Ga;\VR)=H^{\bullet}(\CI(\cring))$ for every
$\cring$, in every degree --- the identification of the cohomology of the free resolution with group
cohomology needs no hypothesis on $\cring$ (Proposition~\ref{prop:isomodel}(3)). The classical operator
depends on none of the three choices \STD{\cite[Ch.~III.9--10]{brown}, as cited in
\eqref{eq:corconjres}}; so neither does $T^{\Iso}_{\Ga\delta\Ga}$. For $\delta=pI$: $\Ga''=\Ga$,
$\conj_{pI}=\id$, $\cor=\res=\id$, and the coefficient map is $\rhon(pI)=p^{\,n}$ (proof of
Proposition~\ref{prop:diamond}).

(2) By (1) the assignment $\Ga\delta\Ga\mapsto T^{\Iso}_{\Ga\delta\Ga}$ is the classical one, which
extends to a ring homomorphism from $\mathcal{H}(\De,\Ga)$ for the $\Z[\De]$-module $\VR$ (the
monoid $\De$ acts through $\rhon$, Definition~\ref{def:semigroup}), by the cited theorem. In
$\mathcal{H}(\De,\Ga)$, Shimura's $T(p^{\nu})$ is the sum of all double cosets of determinant
$p^{\nu}$ (v1.19.0-pre.5; earlier written $T(1,p^{\nu})$, which in his notation is the primitive coset), the local algebras at distinct primes commute, and
$T(p)T(p^{\nu})=T(p^{\nu+1})+p\,T(p,p)T(p^{\nu-1})$. Apply the homomorphism: $T(p^{\nu})$
goes to $T^{\Iso}_{p^{\nu}}$ by definition and $T(p,p)=\Ga\,pI\,\Ga$ to $p^{\,n}\id$ by (1). The last
identity is the $\nu=2$ case of the content decomposition in the proof of
Proposition~\ref{prop:fullassembly}: the double cosets of determinant $p^{2}$ are
$\Ga\,\mathrm{diag}(1,p^{2})\,\Ga$ and $\Ga\,pI\,\Ga$.
\end{proof}

\noindent
\emph{For line~2 (slot-11 WP-11).} Clause~(2) is the statement the X7 recursion checks: at
$\ell\in\{2,3\}$, $\nu\le2$, on $H^{q}(\CI(\cring))$, $q\le5$, entrywise on cohomology (a
basis-free comparison such as the inclusion $P(Z^{q})\subseteq B^{q}$ of
Remark~\ref{rmk:integralevidence}), with $T^{\Iso}_{p^{2}}$ assembled from the two lifts for
$\mathrm{diag}(1,p^{2})$ and $pI$. It is a theorem there, so a failure is a defect of a lift or of its
transport, not a counterexample. At $\ell\ge5$ the same identity on $\Ceq(\cring)$ is
Theorem~\ref{thm:relations}, and the comparison of the two routes through $\cmpI^{\ast}$ is
Proposition~\ref{prop:strictlift}(4).
\emph{Built at v1.11.0} (slot-11 WP-11, \file{docs/implementation/slot11/WP-11\_report.md}): the lifts
that (3) calls ``asserted to exist, not to be built'' are built for $\mathrm{diag}(1,p^{2})$,
$\mathrm{diag}(1,p^{3})$ and $pI$ at $p\in\{2,3\}$ (and at $p=5$ for $\mathrm{diag}(1,25)$ and $5I$),
certified exactly over $\Z[\Ga'']$; at $pI$ the datum is $\tau=\id$ of Proposition~\ref{prop:diamond}'s
proof, and the induction returns a lift homotopic to $\id$ (either is a strict lift). Clause~(2) holds
for $\nu\in\{1,2\}$ on $H^{q}(\CI(\cring))$, $q\le5$, on every cell tried, at $\ell\in\{2,3\}$ with $p=\ell$
included; the evidence is Remark~\ref{rmk:integralevidence}'s last bullet. It is evidence for the
statement, not a replacement for its proof, and nothing about effectivity for other $\delta$ follows.


%======================================================================
\part*{Part VII --- The Dickson invariant, $\theta$, and the ordinary projector}
\addcontentsline{toc}{section}{Part VII --- The Dickson invariant, $\theta$, and the ordinary projector}
%======================================================================
% NEW at v1.11.0-pre.1 (hecke-extensions slot 11, WP-2; plan
% docs/implementation/hecke_extensions_slot11_plan.md S3 "WP-2", with the audit's M-1 and
% m-1, m-2, m-11, m-12 applied, and M-1 itself corrected: U_ell is NOT zero on all of H^0).
% Placed last, before the appendices, so that no chapter number moves (the v1.8.0
% precedent). Notation: sec:notation-dickson.

\section{The Dickson invariant and the operator $\thL$}\label{sec:dicksontheta}
%----------------------------------------------------------------------

\paragraph{What this chapter fixes.}
Multiplication by the first Dickson invariant $\DkL=x^{\ell}y-xy^{\ell}$ as a map of local systems
$\Sym^{n}\to\Sym^{\nplus}$, $\nplus=n+\ell+1$, twisted by the determinant; the operator $\thL$ it
induces on the equivariant complex and on cohomology, with its Hecke equivariance \emph{in the
code's normalisation} (Definition~\ref{def:rhon}, \ref{H2}): $T_p\thL=p\,\thL T_p$, and
$U_\ell\thL=0$ over $\Fl$; the quotient $\Qn$ with a \emph{fixed} complement basis; the long exact
sequence of $0\to\Sym^{n}\to\Sym^{\nplus}\to\Qn\to0$ and its termination; $H^{0}=\Fl[\DkL,\DkD]$ with
its Hecke action; and which statements of the literature on $\theta$ are checkable on the model
at level~$1$. It is the contract of hecke-extensions slot~11's Dickson/$\theta$ line (WP-5, WP-6):
every claim a validation asserts about $\thL$ is a statement of this section.

\paragraph{Standing hypotheses.}
\ref{H4} throughout; $\ell$ is a prime and the ring is $\Fl$ unless a statement says otherwise.
Where $6\in\Fl^{\times}$ ($\ell\ge5$) is needed it is written into the statement (rule~\ref{A1});
Remark~\ref{rmk:thetasmallell} says what survives at $\ell\in\{2,3\}$. Throughout,
$u_t=x^{m-t}y^{t}$ is the basis of $M_m=\Sym^{m}(\Fl^{2})$ of Definition~\ref{def:rhon}, and
$\rho_m(A)f=f(ax+cy,\,bx+dy)$ for $A=\smat abcd$ is that definition's action, reduced mod~$\ell$;
it is defined for every $A\in\Mat_2(\Z)$, singular modulo $\ell$ or not, and $A\mapsto\rho_\bullet(A)$
is a monoid homomorphism into the \emph{ring} endomorphisms of $\Fl[x,y]$
(Proposition~\ref{prop:symconv}(1),(2) and its proof).

\subsection{The two Dickson invariants}

\begin{definition}[The Dickson invariants]\label{def:dicksoninv}
\[
  \DkL:=x^{\ell}y-xy^{\ell}\in M_{\ell+1},\qquad
  \DkD:=\sum_{i=0}^{\ell}x^{\,i(\ell-1)}\,y^{\,(\ell-i)(\ell-1)}\in M_{\ell^{2}-\ell}.
\]
Write $M_{\DkL}\colon M_m\to M_{m+\ell+1}$ for multiplication by $\DkL$; in the bases $(u_t)$ it
sends $u_i\mapsto u_{i+1}-u_{i+\ell}$, which is the matrix \file{HeckeExt/HeckeDicksonMultiplication}
builds.
\end{definition}

\begin{lemma}[The determinant twist, and what singular matrices do]\label{lem:dicksontwist}
\PROVED{} Let $\ell$ be a prime and $A\in\Mat_2(\Z)$, read modulo $\ell$.
\begin{enumerate}
\item $\rho_{\ell+1}(A)\DkL=\det(A)\,\DkL$. Hence, for every $m\ge0$,
  \begin{equation}\label{eq:dicksontwist}
    \rho_{m+\ell+1}(A)\,M_{\DkL}\;=\;\det(A)\,M_{\DkL}\,\rho_m(A)
    \qquad\text{for every }A\in\Mat_2(\Z)\text{, singular included}.
  \end{equation}
\item $\DkL=xy\prod_{a\in\Fl^{\times}}(x-ay)$ is, up to a unit, the product of one linear form
  for each of the $\ell+1$ points of $\mathbb{P}^{1}(\Fl)$; and $\DkL\cdot\DkD=x^{\ell^{2}}y-xy^{\ell^{2}}$.
\item If $A$ is invertible modulo $\ell$ then $\rho(A)\DkD=\DkD$. If $A$ is singular modulo $\ell$
  then $\rho(A)\DkL=0$, and $\DkD$ need not be fixed: $\rho(\mathrm{diag}(1,0))\DkD=x^{\ell^{2}-\ell}$.
\item Evaluation: $\DkL(v)=0$ and $\DkD(v)=1$ for every $v\in\Fl^{2}\smallsetminus\{0\}$.
\item If $f\in M_m$, $m\ge1$, vanishes at every point of $\Fl^{2}$, then $\DkL\mid f$.
\end{enumerate}
\end{lemma}

\begin{proof}
(1) Modulo $\ell$, $(ax+cy)^{\ell}=ax^{\ell}+cy^{\ell}$ (Frobenius, and $a^{\ell}=a$ in $\Fl$). So
$\rho(A)\DkL=(ax^{\ell}+cy^{\ell})(bx+dy)-(ax+cy)(bx^{\ell}+dy^{\ell})=(ad-bc)(x^{\ell}y-xy^{\ell})$,
the $x^{\ell+1}$ and $y^{\ell+1}$ terms cancelling. Since $\rho_\bullet(A)$ is a ring
endomorphism, $\rho_{m+\ell+1}(A)(\DkL f)=\rho_{\ell+1}(A)\DkL\cdot\rho_m(A)f=\det(A)\,\DkL\,\rho_m(A)f$,
which is \eqref{eq:dicksontwist}. No inverse and no invertibility is used.
(2) $x^{\ell-1}-y^{\ell-1}=\prod_{a\in\Fl^{\times}}(x-ay)$; the forms $x$, $y$ and $x-ay$ vanish on
the $\ell+1$ lines of $\Fl^{2}$. With $X=x^{\ell-1}$, $Y=y^{\ell-1}$: $\DkL=xy(X-Y)$,
$x^{\ell^{2}}y-xy^{\ell^{2}}=xy(X^{\ell+1}-Y^{\ell+1})$ and $\DkD=\sum_{i=0}^{\ell}X^{i}Y^{\ell-i}$.
(3) The computation of (1) with $\ell^{2}$ in place of $\ell$ gives
$\rho(A)(x^{\ell^{2}}y-xy^{\ell^{2}})=\det(A)(x^{\ell^{2}}y-xy^{\ell^{2}})$. For $\det A\neq0$ in
$\Fl$, apply $\rho(A)$ to (2): $\det(A)\DkL\cdot\rho(A)\DkD=\det(A)\DkL\DkD$, and cancel in the
domain $\Fl[x,y]$. For singular $A$, (1) gives $\rho(A)\DkL=0$; $\mathrm{diag}(1,0)$ substitutes
$y\mapsto0$, which leaves the $i=\ell$ term of $\DkD$.
(4) $\DkL(v)=0$ by (2). For $v=(v_1,v_2)\neq0$: $X(v),Y(v)\in\{0,1\}$ and not both $0$; if both are
$1$, $\DkD(v)=\ell+1=1$; if exactly one is, exactly one term of $\DkD$ survives, with value $1$.
(5) A homogeneous $f$ with $f(a,b)=0$, $(a,b)\neq0$, is divisible by $bx-ay$ (dehomogenise). The
$\ell+1$ forms of (2) are pairwise non-associate primes of the factorial ring $\Fl[x,y]$, so their
product divides $f$.
\end{proof}

\begin{proposition}[$H^{0}$ is the Dickson ring]\label{prop:dicksonH0}
Let $\ell$ be a prime and $m\ge0$.
\begin{enumerate}
\item \STD{\cite[\S\S2--3]{dickson1911} for the invariants of $\SL_2(\F_q)$ and $\GL_2(\F_q)$ acting on
  $\F_q[x,y]$; \cite[\S1]{wilkerson} for the modern statement. Neither locator is verified against
  the text (rule~\ref{A2})} $\Fl[x,y]^{\SL_2(\Fl)}=\Fl[\DkL,\DkD]$, a polynomial ring on the two
  generators of Definition~\ref{def:dicksoninv}, and $\Fl[x,y]^{\GL_2(\Fl)}=\Fl[\DkL^{\ell-1},\DkD]$.
  (v1.23.0-pre.16: both statements are proved in place, for every prime $\ell$, by Lemma~\ref{lem:dicksonSL2}; the citations are provenance only.)
  (The action here is by substitution, $f\mapsto f\circ A^{\top}$ on points; a group closed under
  transposition has the same invariants either way.)
\item \PROVED{} (by the $\SL_2$ half of (1), Lemma~\ref{lem:dicksonSL2}) $H^{0}(\Ga;M_m)=H^{0}_{\mathrm{eq}}(K,\Rel;M_m)=M_m^{\SL_2(\Fl)}$ has basis
  \[
    \bigl\{\DkL^{a}\DkD^{b}\ :\ a,b\ge0,\ a(\ell+1)+b(\ell^{2}-\ell)=m\bigr\},
  \]
  so $h^{0}(m):=\dim H^{0}(\Ga;M_m)$ is the count of Lemma~\ref{lem:dicksonSL2} (Claim~\ref{cl:dimH1}(2)), whose two
  generators these are. \emph{The $-I$ parity:} for $\ell$ odd both degrees $\ell+1$ and
  $\ell^{2}-\ell$ are even, so $H^{0}=0$ for $m$ odd, as $\rho_m(-I)=(-1)^{m}$ requires
  (Remark~\ref{rmk:negI}); at $\ell=2$, $\deg\DkL=3$ is odd and $-I\equiv I$.
\item \PROVED{} $M_{\DkL}\colon H^{0}(\Ga;M_m)\to H^{0}(\Ga;M_{m+\ell+1})$ is injective, and its
  cokernel has dimension $[\,\ell(\ell-1)\mid m+\ell+1\,]$, spanned by the class of
  $\DkD^{(m+\ell+1)/(\ell^{2}-\ell)}$ when that exponent is an integer. Equivalently
  $h^{0}(m+\ell+1)-h^{0}(m)=[\,\ell(\ell-1)\mid m+\ell+1\,]$.
\end{enumerate}
\end{proposition}

\begin{proof}
(2) $H^{0}_{\mathrm{eq}}=M_m^{\Ga}$ by Proposition~\ref{prop:bredon}(2) (every $\Z[\Ga]$-module, every
$\ell$), and $\Ga$ acts through the surjection $\Ga\to\SL_2(\Fl)$ of Claim~\ref{cl:dimH1}(2); by Lemma~\ref{lem:dicksonSL2} the invariants are $\Fl[\DkL,\DkD]$. A
polynomial ring on homogeneous generators of degrees $\ell+1$ and $\ell^{2}-\ell$ has the monomials
$\DkL^{a}\DkD^{b}$ of degree $m$ as a basis of its degree-$m$ part. (3) Multiplication by
$\DkL\neq0$ is injective in a domain. The monomials of degree $m+\ell+1$ with $a\ge1$ are $\DkL$
times those of degree $m$; the only other one is $\DkD^{b}$ with $b(\ell^{2}-\ell)=m+\ell+1$.
\end{proof}

\subsection{The Hecke action on $H^{0}$}

The operator of Definition~\ref{def:hecke} on $H^{0}$ is, by Remark~\ref{rmk:T3} (which uses (T3)
only, satisfied by every datum, and no hypothesis on $\ell$),
$\Tdc{g}v=\sum_{j=1}^{m}\rho(\delta_j)v$ for $v\in M^{\Ga}$, with $\delta_j=s_j^{-1}g$ the coset
elements of Definition~\ref{def:cosetreps}. For $g=\mathrm{diag}(1,p)$ they are
$\delta_{1+k}=\smat{1}{0}{-k}{p}$, $0\le k<p$, and $\delta_{p+1}=\smat{0}{p}{-1}{0}$.

\begin{proposition}[Hecke operators on $H^{0}(\Ga;M_m)$]\label{prop:heckeH0}
\PROVED{} Assume \ref{H4}. Let $\ell$ be a prime, $m\ge1$, and $\cring=\Fl$; (1) and (2) need no hypothesis
on $\ell$ beyond primality, and none on the datum; (3) is stated for $\ell$ odd, and (4) holds at every
prime ($\ell=2$ by the direct computation at the end of the proof).
\begin{enumerate}
\item \textbf{(Good primes.)} If $g\in\De$ with $\ell\nmid\det g$ and $\Ga g\Ga$ has $m_g$ left cosets,
  then $\Tdc{g}(\DkL^{a}\DkD^{b})=m_g\,(\det g)^{a}\,\DkL^{a}\DkD^{b}$. In particular, for $p\neq\ell$,
  \[
    T_p(\DkL^{a}\DkD^{b})=(1+p)\,p^{a}\,\DkL^{a}\DkD^{b},\qquad
    \diam(\DkL^{a}\DkD^{b})=p^{2a}\,\DkL^{a}\DkD^{b}.
  \]
  (The second agrees with Proposition~\ref{prop:diamond}: $p^{\,m}=p^{2a}$ in $\Fl$, since
  $m=a(\ell+1)+b(\ell^{2}-\ell)\equiv2a\pmod{\ell-1}$.)
\item \textbf{(The bad prime: rank one, square zero.)} For every $v\in M_m^{\Ga}$,
  \[
    U_\ell\,v\;=\;v(1,0)\cdot G_m,\qquad
    G_m:=\sum_{k\in\Fl}(x-ky)^{m}+(-y)^{m}\in M_m ,
  \]
  where $v(1,0)$ is the coefficient of $u_0=x^{m}$ in $v$. $G_m(1,0)=\ell=0$, so
  $U_\ell^{2}=0$ and $\rank U_\ell\le1$ on $H^{0}(\Ga;M_m)$.
\item \textbf{(Which monomials $U_\ell$ kills.)} $U_\ell(\DkL^{a}\DkD^{b})=0$ for $a\ge1$. For $a=0$,
  $\ell$ odd and $M:=b(\ell-1)$,
  \[
    U_\ell(\DkD^{b})=\bigl(G'_{M}\bigr)^{\ell},\qquad
    G'_{M}:=\sum_{k\in\Fl}(x-ky)^{M}+y^{M}
      =-\sum_{s=1}^{b-1}\binom{M}{s(\ell-1)}x^{\,M-s(\ell-1)}y^{\,s(\ell-1)} ,
  \]
  so $U_\ell(\DkD^{b})=0$ if and only if $\ell\mid\binom{b(\ell-1)}{s(\ell-1)}$ for every
  $0<s<b$. This holds for every $b\le\ell$ and \textbf{fails for $b=\ell+1$}.
\item \textbf{(Consequence.)} $U_\ell=0$ on $H^{0}(\Ga;M_m)$ whenever $m<\ell^{3}-\ell$, and
  $U_\ell\neq0$ on $H^{0}(\Ga;M_{\ell^{3}-\ell})$. The first instance is $(m,\ell)=(120,5)$, where
  $H^{0}$ has basis $\DkL^{20},\DkL^{10}\DkD^{3},\DkD^{6}$ and
  \[
    U_5\,\DkD^{6}=-\DkL^{20},\qquad U_5\,\DkL^{20}=U_5\,\DkL^{10}\DkD^{3}=0 .
  \]
  Within $\ell\in\{5,7,11,13\}$, $m\le200$, it is the only one.
\end{enumerate}
\end{proposition}

\begin{proof}
(1) Each $\delta_j$ is invertible modulo $\ell$, so by Lemma~\ref{lem:dicksontwist}(1),(3) and
multiplicativity $\rho(\delta_j)(\DkL^{a}\DkD^{b})=\det(\delta_j)^{a}\DkL^{a}\DkD^{b}$, and
$\det\delta_j=\det g$; sum over the $m_g$ cosets ($p+1$ for $T_p$, one for $\Ga\,pI\,\Ga$).
(2) Modulo $\ell$, $\delta_{1+k}\equiv\smat{1}{0}{-k}{0}$ substitutes $x\mapsto x-ky$, $y\mapsto0$,
and $\delta_{\ell+1}\equiv\smat{0}{0}{-1}{0}$ substitutes $x\mapsto-y$, $y\mapsto0$. For $v$
homogeneous of degree $m$, $v(X,0)=v(1,0)X^{m}$, so $\sum_j\rho(\delta_j)v=v(1,0)G_m$. At $(1,0)$
every $(x-ky)^{m}$ is $1$ and $(-y)^{m}$ is $0$ ($m\ge1$), so $G_m(1,0)=\ell\cdot1=0$; hence
$U_\ell(U_\ell v)=v(1,0)\,G_m(1,0)\,G_m=0$. (That $G_m\in M_m^{\Ga}$ whenever some invariant has
$v(1,0)\neq0$ follows from $U_\ell$ preserving $M_m^{\Ga}$, Theorem~\ref{thm:welldef}.)
(3) $\DkL(1,0)=0$ and $\DkD(1,0)=1$, so $v(1,0)=[a=0]$. For $a=0$: $\DkD(X,0)=X^{\ell(\ell-1)}$, so
$U_\ell\DkD^{b}=\sum_k(x-ky)^{\ell M}+y^{\ell M}$ ($\ell M$ is even), which is $(G'_M)^{\ell}$ by
Frobenius. Expanding, $\sum_{k\in\Fl}k^{t}=-1$ if $t\ge1$ and $(\ell-1)\mid t$,
and $0$ otherwise ($t=0$ gives $\ell=0$); with $(-1)^{s(\ell-1)}=1$ this leaves
$-\sum_{s=1}^{b}\binom{M}{s(\ell-1)}x^{M-s(\ell-1)}y^{s(\ell-1)}$, whose $s=b$ term $-y^{M}$ cancels the
added $y^{M}$. The surviving monomials are distinct, and $\Fl[x,y]$ is a domain, so $(G'_M)^{\ell}=0$
iff every coefficient vanishes. For $1\le s<b\le\ell$, the base-$\ell$ digits of $b(\ell-1)$ are
$(b-1,\ell-b)$ and those of $s(\ell-1)$ are $(s-1,\ell-s)$, so by Lucas's theorem
\STD{\cite[\S XXI]{lucas1878}: $\binom{A}{B}\equiv\prod_i\binom{A_i}{B_i}\pmod\ell$ on base-$\ell$
digits --- the rule \code{mod\_nCr} implements; locator not verified against the text} the binomial is $\binom{b-1}{s-1}\binom{\ell-b}{\ell-s}\equiv0$, since $\ell-s>\ell-b$. For
$b=\ell+1$, $M=\ell^{2}-1$ has digits $(\ell-1,\ell-1)$ and $s=1$ gives
$\binom{\ell-1}{0}\binom{\ell-1}{\ell-1}=1$.
(4) By (3) a non-zero $U_\ell$ needs a basis monomial $\DkD^{b}$ with $b\ge\ell+1$, i.e.\
$m=b(\ell^{2}-\ell)\ge\ell^{3}-\ell$; at $m=\ell^{3}-\ell$, $b=\ell+1$ gives $U_\ell\DkD^{\ell+1}\neq0$.
At $\ell=5$, $b=6$: every $\binom{24}{4s}$, $1\le s\le5$, is $\equiv1\pmod5$ (digits $(4,4)$ against
$(s-1,5-s)$), so $G'_{24}=-\sum_{s=1}^{5}x^{24-4s}y^{4s}=-(xy)^{4}(x^{4}-y^{4})^{4}=-\DkL^{4}$, using
$(x^{4}-y^{4})^{5}=x^{20}-y^{20}$; hence $U_5\DkD^{6}=(-\DkL^{4})^{5}=-\DkL^{20}$. The next values of
$m$ at which some $\DkD^{b}$ with $b\ge\ell+1$ occurs are $m=140,\ldots,200$ at $\ell=5$ ($b=7,\ldots,10$),
where every $\binom{4b}{4s}$, $0<s<b$, vanishes modulo $5$ by Lucas's theorem (for $b=7$: $28$ has
base-$5$ digits $(1,0,3)$, and each $4s\le24$ has a digit exceeding the corresponding one); at $\ell\ge7$
the first such $m$ is $(\ell+1)(\ell^{2}-\ell)\ge336$. \COMP\textsuperscript{a} An independent script
(slot-11 WP-2 report) lists the $b<3\ell$ with $U_\ell\DkD^{b}\neq0$ as $\{6,11,12\}$, $\{8,15,16\}$,
$\{12,23,24\}$, $\{14,27,28\}$ at $\ell=5,7,11,13$, and confirms $U_5\DkD^{6}=-\DkL^{20}$.
\emph{At $\ell=2$} (v1.11.0-pre.2; the WP-3 referee's m-1), where (3)'s expansion is not used: by (2),
$U_2v=v(1,0)\bigl(x^{m}+(x+y)^{m}+y^{m}\bigr)$; the basis monomials with $v(1,0)\ne0$ are the $\DkD^{b}$,
$m=2b$, and $x^{m}+(x+y)^{m}+y^{m}$ vanishes for $m\in\{2,4\}$ and equals $x^{4}y^{2}+x^{2}y^{4}=\DkL^{2}$ at
$m=6=\ell^{3}-\ell$. So $U_2\DkD=U_2\DkD^{2}=0$ and $U_2\DkD^{3}=\DkL^{2}\neq0$: (4) holds at $\ell=2$,
with $+\DkL^{2}$ in place of the odd-$\ell$ sign.
\end{proof}

\begin{remark}[The plan's statement, and what replaces it]\label{rmk:heckeH0plan}
The slot-11 plan (and its audit's finding M-1, which corrected an earlier ``$(1+\ell)\ell^{a}$'')
stated ``$U_\ell=0$ on all of $H^{0}(\Ga;\Sym^{n}\Fl)$''. That is true for $n<\ell^{3}-\ell$ and false
at $n=\ell^{3}-\ell$ (Proposition~\ref{prop:heckeH0}(4)); the audit's check covered $\DkD$ and $\DkD^{2}$
only. What holds at every $n\ge1$ is $U_\ell^{2}=0$ on $H^{0}$, which is what ``every Dickson class is
non-ordinary'' needs (Corollary~\ref{cor:thetanonordinary}). At $n=0$, $U_\ell=\ell+1=1$ on
$H^{0}(\Ga;\Fl)$: the constant class is ordinary.
\end{remark}

\subsection{The quotient $\Qn$ and the operator $\thL$}

\begin{lemma}[Definition~\ref{def:hecke} with any monoid fibre]\label{lem:heckefibre}
\PROVED{} Assume \ref{H4}. Let $W$ be an $\cring$-module with a monoid homomorphism
$\rho_W\colon\Mat_2(\Z)\to\End_{\cring}(W)$ (for instance $M_m$, $\Qn$, via reduction mod~$\ell$).
Then \eqref{eq:heckechain} with $\rho_W$ in place of $\rhon$ defines $\Tdc{g}$ on
$\Hom_{\Z[\Ga]}(C_{\bullet}(\Xbulk),W)$ for every $g\in\De$ and every datum $(\{s_j\},\tau)$; the
hypothesis-free half of Theorem~\ref{thm:welldef} (equivariance and the chain-map property) and
Remark~\ref{rmk:T3} hold verbatim; and a $\Z[\De]$-linear map $\phi\colon W\to W'$ induces a
cochain map that commutes with $\Tdc{g}$. \emph{The datum $(\{s_j\},\tau)$ does not involve $W$}
(Definition~\ref{def:tau}): one datum serves every fibre at once.
\end{lemma}

\begin{proof}
The proof of Theorem~\ref{thm:welldef}'s first half uses of $\rhon$ only
$\rhon(A_1A_2)=\rhon(A_1)\rhon(A_2)$ on the integral identity
$\delta_j\,g^{-1}h_j(\gamma)g=s_j^{-1}h_j(\gamma)g$; that of Remark~\ref{rmk:T3} uses (T3) only.
For $\phi$: $\phi\rho_W(\delta_j)=\rho_{W'}(\delta_j)\phi$, term by term in \eqref{eq:heckechain}.
\end{proof}

\begin{definition}[The quotient $\Qn$, its fixed basis, and its monoid action]\label{def:Qn}
Let $\ell$ be a prime, $n\ge0$ and $\nplus:=n+\ell+1$. Put $\Qn:=M_{\nplus}/\DkL\cdot M_n$ over $\Fl$,
with $\pi\colon M_{\nplus}\to\Qn$ the projection, and for $A\in\Mat_2(\Z)$ let $\rho_{\Qn}(A)$ be the map
induced by $\rho_{\nplus}(A)$ --- well defined for \emph{every} $A$, singular included, because
$\rho_{\nplus}(A)(\DkL\cdot M_n)\subseteq\DkL\cdot M_n$ by \eqref{eq:dicksontwist}. Then:
\begin{enumerate}
\item $\dim\Qn=\ell+1$ for every $n\ge0$, and
  \[
    \mathcal B_n:=(\bar u_0,\bar u_1,\ldots,\bar u_{\ell-1},\bar u_{\nplus})
  \]
  is a basis --- \textbf{the} basis in which $\rho_{\Qn}$ is written, fixed here. With the class map
  $c(0)=0$, $c(\nplus)=\nplus$ and $c(j)=1+\bigl((j-1)\bmod(\ell-1)\bigr)$ for $1\le j\le\nplus-1$,
  one has $\pi(u_j)=\bar u_{c(j)}$ for every $j$, and hence
  \[
    \rho_{\Qn}(A)_{s,t}\;=\;\sum_{j:\ c(j)=s}\rho_{\nplus}(A)_{j,t}\qquad(s,t\in\{0,1,\ldots,\ell-1,\nplus\}).
  \]
  Equivalently: reduce $u_r\mapsto u_{r-\ell+1}$ for $r=\nplus-1$ down to $\ell$, then read the
  coordinates at $0,\ldots,\ell-1,\nplus$.
\item $A\mapsto\rho_{\Qn}(A)$ is a monoid homomorphism $\Mat_2(\Z)\to\End_{\Fl}(\Qn)$, and
  $\rho_{\Qn}(A)\pi=\pi\rho_{\nplus}(A)$ --- \emph{untwisted}.
\item The local system $L_{\Qn}$ on $(K,\Rel)$ is that of Definition~\ref{def:model7rel} with
  $\rho_{\Qn}$ in place of $\rhon$ on every relation, the central relation included
  ($\rho_{\Qn}(-I)=(-1)^{\nplus}=(-1)^{n}$ for $\ell$ odd); its complex is
  $\Ceq(K,\Rel;\Qn)\cong\Hom_{\Z[\Ga]}(C_{\bullet}(\Xbulk),\Qn)$ by
  Theorem~\ref{thm:explicit_complex_equiv}.
\end{enumerate}
\end{definition}

\noindent
(1) is a statement, not only a definition, and its proof is one line: $\DkL u_i=u_{i+1}-u_{i+\ell}$,
$0\le i\le n$, so $\DkL\cdot M_n$ is spanned by the $n+1$ relations $u_j\equiv u_{j+\ell-1}$,
$1\le j\le n+1$, which are independent (their largest indices $j+\ell-1$ are distinct); every
$u_r$ with $\ell\le r\le\nplus-1$ reduces to $u_{r-(\ell-1)}$, $u_0$ and $u_{\nplus}$ occur in no
relation, and $\dim\Qn=(\nplus+1)-(n+1)=\ell+1=\abs{\mathcal B_n}$. (2) holds because
$\rho_{\nplus}$ is multiplicative and preserves the subspace. This is the basis
\code{FibreRepresentation::QuotientL} implements (slot-11 WP-5); the alternative normal form
$(\bar u_0,\bar u_{n+2},\ldots,\bar u_{\nplus})$ is \emph{not} this document's.

\begin{lemma}[$\Qn$ as homogeneous functions, and its invariants]\label{lem:Qnfunctions}
\PROVED{} Let $\ell$ be a prime and $n\ge0$.
\begin{enumerate}
\item Evaluation $f\mapsto(v\mapsto f(v))$ induces an isomorphism of $\Qn$ onto the space
  $\mathcal F_{\nplus}$ of functions $\phi\colon\Fl^{2}\to\Fl$ with $\phi(0)=0$ and
  $\phi(\lambda v)=\lambda^{\nplus}\phi(v)$ for $\lambda\in\Fl^{\times}$, which intertwines
  $\rho_{\Qn}(A)$ with $(A\star\phi)(v):=\phi(A^{\top}v)$ for every $A\in\Mat_2(\Fl)$. So $\Qn$
  depends, as a $\Mat_2(\Fl)$-module, only on $n\bmod(\ell-1)$. (As a $\GL_2(\Fl)$-module it is the
  principal series of \Espec{cl:glover}(2).)
\item $H^{0}(\Ga;\Qn)=H^{0}_{\mathrm{eq}}(K,\Rel;\Qn)=\Qn^{\SL_2(\Fl)}$ has dimension
  $[\,(\ell-1)\mid n+2\,]$. When it is non-zero it is spanned by
  \[
    \mathbf 1_n:=\bar u_0-\bar u_{\ell-1}+\bar u_{\nplus},
  \]
  the class whose function is $1$ on $\Fl^{2}\smallsetminus\{0\}$. The image of
  $H^{0}(\Ga;M_{\nplus})$ in it is spanned by $\pi(\DkD^{b})=\mathbf 1_n$ when
  $b=\nplus/(\ell^{2}-\ell)$ is an integer, and is $0$ otherwise.
\item On $\mathbf 1_n$: $\Tdc{g}\mathbf 1_n=m_g\,\mathbf 1_n$ if $\ell\nmid\det g$ (so
  $T_p\mathbf 1_n=(p+1)\mathbf 1_n$), and $U_\ell\,\mathbf 1_n=0$.
\end{enumerate}
\end{lemma}

\begin{proof}
(1) $(\rho(A)f)(v)=f(a v_1+cv_2,\,bv_1+dv_2)=f(A^{\top}v)$. The kernel of evaluation on
$M_{\nplus}$ ($\nplus\ge1$) is $\DkL\cdot M_n$ by Lemma~\ref{lem:dicksontwist}(5) and (4); the image lies in
$\mathcal F_{\nplus}$, which has dimension $\ell+1$ (one value per line), and so does $\Qn$. The
space $\mathcal F_{\nplus}$ and the $\star$-action depend on $\nplus$ only through
$\lambda\mapsto\lambda^{\nplus}$, i.e.\ on $\nplus\bmod(\ell-1)$, and $\nplus\equiv n+2$.
(2) $H^{0}_{\mathrm{eq}}=\Qn^{\Ga}$ by Proposition~\ref{prop:bredon}(2), and $\Ga$ acts through
$\SL_2(\Fl)$, which --- like its transpose --- acts transitively on $\Fl^{2}\smallsetminus\{0\}$; so an
invariant function is a constant $c$ there, and a non-zero constant lies in $\mathcal F_{\nplus}$ iff
$\lambda^{\nplus}=1$ for all $\lambda\in\Fl^{\times}$, iff $(\ell-1)\mid\nplus$, iff $(\ell-1)\mid n+2$.
The representative $x^{\nplus}-x^{\nplus-\ell+1}y^{\ell-1}+y^{\nplus}$ of $\mathbf 1_n$ takes the value
$1-1+1$, $1$, $1$ at points with both, only the first, only the second coordinate non-zero (all
exponents are positive multiples of $\ell-1$). $\pi(\DkD^{b})$ has function $\DkD(v)^{b}=1$ by
Lemma~\ref{lem:dicksontwist}(4); the other basis monomials of $H^{0}(M_{\nplus})$ are multiples of
$\DkL$ and map to $0$.
(3) By Remark~\ref{rmk:T3} (with $\rho_{\Qn}$, Lemma~\ref{lem:heckefibre}),
$T\mathbf 1_n=\sum_j\delta_j\star\mathbf 1$ and $(A\star\mathbf 1)(v)=[A^{\top}v\neq0]$. If
$\ell\nmid\det\delta_j$ every term is $\mathbf 1$. For $U_\ell$, with the $\delta_j$ of
Proposition~\ref{prop:heckeH0}(2): $\delta_{1+k}^{\top}v=(v_1-kv_2,0)$ and
$\delta_{\ell+1}^{\top}v=(-v_2,0)$ modulo $\ell$; at $v_2\neq0$ the count of non-zero images is
$(\ell-1)+1=\ell$, at $v_2=0$ it is $\ell+0=\ell$, so the sum is the function $\ell=0$.
\end{proof}



\begin{definition}[The operator $\thL$]\label{def:theta}
For $n\ge0$ let $\thL\colon\Ceq(K,\Rel;M_n)\to\Ceq(K,\Rel;M_{\nplus})$ be $M_{\DkL}$ applied to every
cochain value, and $\pi_{\ast}\colon\Ceq(K,\Rel;M_{\nplus})\to\Ceq(K,\Rel;\Qn)$ likewise. Both are cochain
maps: $M_{\DkL}\rho_n(\gamma)=\rho_{\nplus}(\gamma)M_{\DkL}$ for $\gamma\in\Ga$ by
\eqref{eq:dicksontwist} with $\det\gamma=1$, so $\thL$ carries the relations of $M_n$ to those of
$M_{\nplus}$, and it commutes with $d$ because it acts on values. The same symbol denotes the induced
maps $H^{q}(\Ga;M_n)\to H^{q}(\Ga;M_{\nplus})$. This is the classical $\theta$ in the cohomological
model (\S\ref{sec:thetaliterature}); \Espec{cl:Tlnil}'s ``$L\cdot$''.
\end{definition}

\begin{proposition}[Twisted Hecke equivariance of $\thL$; $U_\ell\thL=0$]\label{prop:thetatwist}
\PROVED{} Assume \ref{H4}; $\ell$ any prime, $\cring=\Fl$. Let $g\in\De$ and let $(\{s_j\},\tau)$ be one
$g$-twisted datum; write $\widetilde T^{(W)}$ for \eqref{eq:heckechain} with fibre $W$
(Lemma~\ref{lem:heckefibre}). Then, \textbf{at cochain level},
\[
  \widetilde T^{(M_{\nplus})}\circ\thL\;=\;\det(g)\cdot\thL\circ\widetilde T^{(M_n)},\qquad
  \pi_{\ast}\circ\widetilde T^{(M_{\nplus})}\;=\;\widetilde T^{(\Qn)}\circ\pi_{\ast}.
\]
Consequently, on $H^{q}$ for every $q$ and every prime $p$:
\[
  T_p\,\thL=p\,\thL\,T_p,\qquad
  T_{p^{\nu}}\,\thL=p^{\nu}\,\thL\,T_{p^{\nu}},\qquad
  \diam\,\thL=p^{2}\,\thL\,\diam,
\]
the left operator on $M_{\nplus}$, the right one on $M_n$; and at $p=\ell$,
$U_\ell\circ\thL=0$ \emph{already on cochains}.
\textbf{The twist is $p$, not $p^{-1}$}: the $\delta_j=s_j^{-1}g$ have determinant $\det g$, are
never inverted (\ref{H2}, Obligation~\ref{obl:integral}), and $\rho_\bullet$ is covariant
(Proposition~\ref{prop:symconv}(2)). The $\diam$ line is consistent with
Proposition~\ref{prop:diamond}: $p^{\nplus}=p^{\,n}\,p^{\ell+1}=p^{2}p^{\,n}$ in $\Fl$.
\end{proposition}

\begin{proof}
$(\widetilde T^{(M_{\nplus})}\thL\phi)(x)=\sum_j\rho_{\nplus}(\delta_j)M_{\DkL}\phi(\tau(s_jx))
=\sum_j\det(\delta_j)\,M_{\DkL}\rho_n(\delta_j)\phi(\tau(s_jx))=\det(g)\,(\thL\widetilde T^{(M_n)}\phi)(x)$
by \eqref{eq:dicksontwist}, with the same $\tau$ and $s_j$ on both sides (Lemma~\ref{lem:heckefibre}).
The second identity is Lemma~\ref{lem:heckefibre} for $\pi$, which is $\Z[\De]$-linear by
Definition~\ref{def:Qn}(2). Summing over the double cosets of \eqref{eq:fullop}, each of determinant
$p^{\nu}$, gives the $T_{p^{\nu}}$ line; $\diam$ is $\Ga\,pI\,\Ga$, of determinant $p^{2}$. At $p=\ell$,
$\det(g)=\ell=0$ in $\Fl$.
\end{proof}

\noindent
\textbf{On cohomology the datum does not matter} when $\ell\ge5$ (Theorem~\ref{thm:welldef}(b)); at
$\ell\in\{2,3\}$ the identities hold for the operators built from one datum, which is what the shipped
pipeline does (the canonical $\tau$ of Algorithm~\ref{alg:tau} is a datum on $C_{\bullet}(\Ku)$, not on
a fibre). This discharges the claim \Espec{cl:Tlnil} marks \STD{}: ``$U_\ell$ kills $\theta$-images''
is Proposition~\ref{prop:thetatwist}, in the stronger form $U_\ell\circ\thL=0$ on cochains.

\subsection{The long exact sequence}

\begin{theorem}[The $\thL$-sequence]\label{thm:thetales}
Assume \ref{H4} and $\ell\ge5$ (so $6\in\Fl^{\times}$); let $n\ge0$ and write $h^{q}(m)=\dim H^{q}(\Ga;M_m)$.
\begin{enumerate}
\item \PROVED{} $0\to\Ceq(K,\Rel;M_n)\xrightarrow{\thL}\Ceq(K,\Rel;M_{\nplus})\xrightarrow{\pi_{\ast}}
  \Ceq(K,\Rel;\Qn)\to0$ is an exact sequence of cochain complexes, and its long exact sequence is
  \begin{equation}\label{eq:thetales}
    0\to H^{0}(M_n)\xrightarrow{\thL}H^{0}(M_{\nplus})\xrightarrow{\pi_{\ast}}H^{0}(\Qn)
    \xrightarrow{\ \delta\ }H^{1}(M_n)\xrightarrow{\thL}H^{1}(M_{\nplus})\xrightarrow{\pi_{\ast}}
    H^{1}(\Qn)\to0,
  \end{equation}
  every group being $\Heq{\bullet}(K,\Rel;-)$, which is $H^{\bullet}(\Ga;-)$ for the two $\Sym$ terms
  (Remark~\ref{rmk:orbifold}) and is written so for all three; $H^{q}=0$ for $q\ge2$ on all three terms, which is
  why it stops. $\delta$ is the connecting map of the snake construction on cochains.
\item \PROVED{} \textbf{(Hecke compatibility.)} For $g\in\De$, the triple
  $(\det(g)\,\Tdc{g}^{(M_n)},\ \Tdc{g}^{(M_{\nplus})},\ \Tdc{g}^{(\Qn)})$ is an endomorphism of
  \eqref{eq:thetales}: besides Proposition~\ref{prop:thetatwist},
  $\delta\circ\Tdc{g}^{(\Qn)}=\det(g)\,\Tdc{g}^{(M_n)}\circ\delta$. At $p=\ell$ this reads
  $\delta\circ U_\ell^{(\Qn)}=0$ and says nothing about $U_\ell$ on $\im\delta$.
\item \PROVED{} \textbf{($\ker\thL=\im\delta$, and what it is.)} On $H^{1}$,
  \[
    \dim\ker\bigl(\thL|_{H^{1}(M_n)}\bigr)=\rank\delta=[\,(\ell-1)\mid n+2\,]-[\,\ell(\ell-1)\mid n+\ell+1\,]
    \in\{0,1\}.
  \]
  When it is $1$, $\ker\thL=\Fl\,\delta(\mathbf 1_n)$ and $T_p\,\delta(\mathbf 1_n)=(1+p^{-1})\,\delta(\mathbf 1_n)$
  for every $p\neq\ell$; since $(\ell-1)\mid k:=n+2$, $1+p^{-1}=1+p^{k-1}$ in $\Fl$ --- \textbf{the
  kernel of $\thL$ on $H^{1}$ is at most the Eisenstein line, and it is exactly that line when
  $(\ell-1)\mid k$ and $\ell(\ell-1)\nmid k+\ell-1$}.
\item \PROVED{} (given Claim~\ref{cl:dimH1}(1), which is \STD{\cite[Ch.~8]{shimura}, \cite{deng}, as cited there}) \textbf{(Closed forms.)} For $n$
  even, with $d_k$ as in Claim~\ref{cl:dimH1} (so $d_2=0$) and $h^{1}(0)=0$:
  \[
  \begin{aligned}
    \dim H^{0}(\Qn)&=[\,(\ell-1)\mid n+2\,],\\
    \rank\thL|_{H^{0}}&=h^{0}(n),\qquad \dim\coker\thL|_{H^{0}}=[\,\ell(\ell-1)\mid n+\ell+1\,],\\
    \rank\thL|_{H^{1}}&=h^{1}(n)-\rank\delta,\\
    \dim H^{1}(\Qn)&=2\bigl(d_{\,n+\ell+3}-d_{\,n+2}\bigr)+[\,(\ell-1)\mid n+2\,],
  \end{aligned}
  \]
  so $\dim H^{1}(\Qn)-\dim H^{0}(\Qn)=2(d_{n+\ell+3}-d_{n+2})$. For $n$ odd all six groups are $0$.
\item \PROVED{} \textbf{(Multiplicities.)} Let $a=(a_p)_{p\neq\ell}$ be a system of eigenvalues in
  $\bar\F_\ell$ and $E_a(H)$ its generalised eigenspace in $H\otimes\bar\F_\ell$ for the commuting
  family $\{T_p\}_{p\neq\ell}$; write $pa:=(p\,a_p)_p$. Then $\thL$ maps $E_a(H^{1}(M_n))$ into
  $E_{pa}(H^{1}(M_{\nplus}))$, injectively unless $a_p=1+p^{-1}$ for all $p$; so
  \[
    \dim E_{pa}\bigl(H^{1}(M_{\nplus})\bigr)\;\ge\;\dim E_a\bigl(H^{1}(M_n)\bigr)-[\,a=(1+p^{-1})_p\,]\cdot\rank\delta .
  \]
\end{enumerate}
\end{theorem}

\begin{proof}
(1) $\Ceq(K,\Rel;W)\cong\bigoplus_{\sigma}W^{\Ga_\sigma}$ over orbit representatives
(Definition~\ref{def:ringinforce}(3) with $W$ for $\VR$; \ref{H4}(i),(iv)), and each $\abs{\Ga_\sigma}\in\{2,4,6\}$
is a unit of $\Fl$ for $\ell\ge5$, so $W\mapsto W^{\Ga_\sigma}$ is exact (averaging). The sequence of
fibres $0\to M_n\to M_{\nplus}\to\Qn\to0$ is exact and $\Ga$-equivariant (Definition~\ref{def:Qn},
\eqref{eq:dicksontwist} with $\det\gamma=1$), so the sequence of complexes is exact in each degree.
Its long exact sequence has $H^{q}=0$ for $q\ge2$ on every term by Proposition~\ref{prop:bredon}(2),
which holds for every $\Z[\Ga]$-module and every $\ell$.
(2) $\Tdc{g}$ on the three terms, twisted by $\det g$ on the first, commute with $\thL$ and $\pi_\ast$
by Proposition~\ref{prop:thetatwist}, so they form an endomorphism of the short exact sequence of
complexes, and connecting maps are natural \STD{\cite[Thm.~1.3.1]{weibel}}.
(3) By exactness $\ker\thL|_{H^{1}}=\im\delta\cong H^{0}(\Qn)/\pi_\ast H^{0}(M_{\nplus})$; by
Lemma~\ref{lem:Qnfunctions}(2) and Proposition~\ref{prop:dicksonH0}(3) the two dimensions are
$[(\ell-1)\mid n+2]$ and $[\ell(\ell-1)\mid\nplus]$, the second implying the first. When $\rank\delta=1$,
$\delta(\mathbf 1_n)\neq0$ spans the image and, by (2) with $g=\mathrm{diag}(1,p)$ and
Lemma~\ref{lem:Qnfunctions}(3), $p\,T_p\,\delta(\mathbf 1_n)=\delta(T_p\mathbf 1_n)=(p+1)\,\delta(\mathbf 1_n)$.
Finally $p^{k}=1$ for $(\ell-1)\mid k$ gives $p^{-1}=p^{k-1}$.
(4) The first two lines are Lemma~\ref{lem:Qnfunctions}(2) and Proposition~\ref{prop:dicksonH0}(3);
the third is exactness at $H^{1}(M_n)$; for the fourth, exactness at the last two spots gives
$\dim H^{1}(\Qn)=h^{1}(\nplus)-h^{1}(n)+\rank\delta$, and Claim~\ref{cl:dimH1}(1),
$h^{1}(m)=2d_{m+2}-1+h^{0}(m)$ for even $m>0$ --- which also holds at $m=0$, where
$H^{1}(\Ga;\Fl)=\Hom(\Ga,\Fl)=0$ ($\Ga^{\mathrm{ab}}=\Z/12$, Claim~\ref{cl:amalgam}(4)'s citation) and
$d_2=0$ by the formula --- together with Proposition~\ref{prop:dicksonH0}(3), gives
$2(d_{\nplus+2}-d_{n+2})+[\ell(\ell-1)\mid\nplus]+\rank\delta$, which is the displayed value. For $n$
odd, $\nplus$ is odd too, $-I$ acts by $-1$ on all three fibres, and every $W^{\Ga_\sigma}=0$
($-I\in\Ga_\sigma$, $2\in\Fl^{\times}$).
(5) $T_p\thL=p\,\thL T_p$ (Proposition~\ref{prop:thetatwist}) maps generalised eigenspaces as stated,
and $\ker\thL=\im\delta$ lies in the eigenspace of $a_p=1+p^{-1}$ by (3).
\end{proof}

\noindent
\emph{Re-derivation note} (slot-11 WP-2, not a registered validation, and not a status tag of
Theorem~\ref{thm:thetales}; v1.11.0-pre.2, the WP-3 referee's m-3). (3)'s rank formula, the three
dimension lines of (4) and $\dim H^{0}(\Qn)$ were re-derived independently of the pipeline, from the
presentation of Claim~\ref{cl:amalgam}(3) applied to $M_n$, $M_{\nplus}$ and $\Qn$ with the basis of
Definition~\ref{def:Qn}, at $\ell\in\{5,7,11,13\}$ and every even $n\le60$ ($0$ disagreements);
$\rho_{\Qn}$'s well-definedness and multiplicativity on random, singular included, $A$ at $n<40$;
Proposition~\ref{prop:heckeH0}(1)--(3) at $\ell\in\{5,7\}$, $n\le220$. The registered check is
slot-11 WP-5 / WP-6's \code{HeckeDicksonThetaValidation}. \emph{v1.11.0}: and
\code{HeckeDicksonThetaHeckeValidation}, each with a \code{slow} companion; their results are the last
paragraph of Remark~\ref{rmk:thetacheckable}.

\begin{remark}[At $\ell\in\{2,3\}$ the sequence is not available on the model]\label{rmk:thetasmallell}
Definition~\ref{def:theta}, Proposition~\ref{prop:thetatwist}, Lemma~\ref{lem:Qnfunctions} and
Proposition~\ref{prop:heckeH0} hold at every prime. Theorem~\ref{thm:thetales} does not: at
$\ell\in\{2,3\}$ the invariants functor $W\mapsto W^{\Ga_\sigma}$ is not exact for a stabiliser whose order
$\ell$ divides --- $\langle U\rangle$ (order $6$) at $\ell=3$, where the order-$4$ $\langle S\rangle$ is harmless,
and every stabiliser at $\ell=2$ ($-I$ lies in each) --- so the sequence of \emph{equivariant cochain complexes} in (1) need not be exact,
and the model computes the Bredon object of Proposition~\ref{prop:bredon}, not $H^{\bullet}(\Ga;-)$
(\ref{H3}; the plan's audit, finding m-12). Group cohomology has its own long exact sequence for
$0\to M_n\to M_{\nplus}\to\Qn\to0$, but it does not terminate --- $H^{2}(\Ga;\F_2)\neq0$
(Lemma~\ref{lem:h2torsion}(2)) --- and the model does not compute it. Nothing in this section is
asserted about either at $\ell\in\{2,3\}$.
\end{remark}

\subsection{$\theta$ in the literature, and what is checkable on the model}\label{sec:thetaliterature}

The operator $\thL$ is the cohomological counterpart of Serre's and Swinnerton-Dyer's $\theta=q\,d/dq$ on
modular forms modulo $\ell$. \textbf{None of the pointers below is verified against the text}: no copy
of the cited works is in \file{docs/literature/} (rule~\ref{A2}); each is at section or theorem level
from the standard literature and is flagged \emph{VERIFY} in the slot-11 WP-2 report for the referee.

\begin{remark}[The classical statements]\label{rmk:thetaclassical}
Let $\ell\ge5$, and let $\widetilde M_k$ be the space of reductions modulo $\ell$ of level-$1$ modular forms
of weight $k$ with $\ell$-integral $q$-expansion; for $f\neq0$ let $w(f)$, the \emph{filtration}, be the
least $k$ with $f\in\widetilde M_k$.
\begin{enumerate}
\item \STD{\cite[\S\S1--2]{serre}, \cite[\S3]{swinnerton_dyer}, and for general level
  \cite[its main theorem]{katz1977}; locators not verified; v1.21.0-pre.2: \cite[Theorem~(2), p.~55]{katz1977} read (slot-21 literature pass,
  \file{docs/implementation/slot21/missing\_literature.md}~\S7): $\ell\nmid w(f)$ implies exact filtration $w(f)+\ell+1$, for $N\ge3$ (level one is Serre and
  Swinnerton-Dyer's case); the ``iff'' is \cite[Lemma~4.1(b), p.~257]{jochnowitz_local}, read} $\theta\colon\widetilde M_k\to\widetilde M_{k+\ell+1}$,
  $a_n(\theta f)=n\,a_n(f)$; on eigenforms $a_p(\theta f)=p\,a_p(f)$ --- Proposition~\ref{prop:thetatwist}'s
  $T_p\thL=p\thL T_p$ is this, in the coefficients of cohomology; and $w(\theta f)\le w(f)+\ell+1$ with
  \textbf{equality iff $w(f)\not\equiv0\pmod\ell$} (the \emph{filtration law}); $E_{\ell-1}\equiv1$, so
  $f\in\widetilde M_k$ implies $f\in\widetilde M_{k+\ell-1}$.
\item \STD{\cite[Thm.~4.1, p.~276]{jochnowitz}, finiteness \cite[Thm.~2.2, p.~272]{jochnowitz}; verified, slot-17 R17-4} every level-$1$ system of Hecke
  eigenvalues modulo $\ell$ away from $\ell$ is a $\theta$-twist $(p^{i}a_p)_p$ of one occurring in weight
  $\le\ell+1$; there are finitely many (\Espec{cl:sigma}).
\item \STD{\cite[Lemma~2.1, p.~204]{ash_stevens} for the inclusion $\subseteq$ (verified, slot-17 R17-4); the equality and the
  $\det^{j}$ / $p^{j}$ bookkeeping are not in \cite{ash_stevens} and rest on \Espec{cl:glover} and Proposition~\ref{prop:thetatwist}} the systems of eigenvalues in
  $H^{1}(\Ga;M_n)$ are those of the Jordan--H\"older factors $\Sym^{r}\otimes\det^{j}$ ($0\le r\le\ell-1$),
  twisted by $p^{j}$, and they are the systems of mod-$\ell$ eigenforms of weight $k=n+2$ up to that
  bookkeeping --- the cohomological form of (2). The factors are those of \Espec{cl:glover}; the
  twist $\det^{j}$ is realised by $\thL^{j}$ (Proposition~\ref{prop:thetatwist}).
\item \STD{\cite[\S3]{edixhoven} (Tate's $\theta$-cycles) and \cite{jochnowitz_local}; locators not
  verified; v1.21.0-pre.2: \cite[\S7, pp.~261--264]{jochnowitz_local} ($\theta$-cycles) read and in agreement (slot-21 literature pass); Edixhoven's not} \textbf{$\theta$-cycles.} For an eigenform $f$ with $\theta f\neq0$ the sequence
  $w(\theta f),w(\theta^{2}f),\ldots$ is periodic of period dividing $\ell-1$ ($\theta^{\ell}=\theta$ on
  $q$-expansions); by the filtration law each step raises $w$ by $\ell+1$ except at a \emph{drop},
  which occurs exactly after a filtration $\equiv0\pmod\ell$ and lowers it to
  $w+\ell+1-s(\ell-1)$ for some $s\ge1$.
\item \STD{\cite[the main theorem, as stated in its introduction]{gross}, re-proved in
  \cite{coleman_voloch}; locators not verified, and whether the latter's range includes $k=\ell-1$ is not
  verified either} \textbf{Companion forms} (Gross):
  for an ordinary eigenform $f$ of weight $2\le k\le\ell-1$ with $\bar\rho_f$ irreducible, there is an
  eigenform $g$ of weight $k'=\ell+1-k$ with $\bar\rho_g\cong\bar\rho_f\otimes\omega^{1-k}$ ($\omega$ the
  mod-$\ell$ cyclotomic character) if and only if $\bar\rho_f|_{G_\ell}$ is split (tamely ramified),
  under Gross's further hypotheses, which are not restated here. In $q$-expansions $\theta g=\theta^{\ell-k+1}f$:
  the companion is the $\theta^{\ell-k}$-twist point of the $\theta$-cycle, and at level~$1$ its
  existence is the statement that the cycle's filtration drops there. \textbf{$\Delta$ modulo $11$ is
  outside this range}: $k=12=\ell+1$ (it is the level-$1$ avatar of the weight-$2$ newform of level $11$).
  The endpoint $k=\ell-1$, where the companion has weight $2$, is the case (C8) of
  Remark~\ref{rmk:thetacheckable} uses; its hypotheses there are not restated (v1.11.0-pre.2).
\end{enumerate}
\end{remark}

\begin{remark}[What is checkable on the model at level $1$, $n\le200$: the prediction list]\label{rmk:thetacheckable}
The list slot-11 WP-6 predicts from, in writing, before it computes (rule~20 of the slot plan). Each item
says what it rests on. $\tau(p)$ is Ramanujan's function ($a_p(\cf)$, \S\ref{sec:notation-hecke}).
\begin{enumerate}
\item[(C1)] \PROVED{} \textbf{On every cell} with $\ell\ge5$: $T_p\thL=p\thL T_p$ entrywise on $H^{0}$
  and $H^{1}$ for $p\in\{2,3,5,7\}\smallsetminus\{\ell\}$; $U_\ell\thL=0$; exactness of
  \eqref{eq:thetales} at all six spots; the closed forms of Theorem~\ref{thm:thetales}(4); $\ker\thL=\im\delta$
  with $T_p=1+p^{-1}$ on it when non-zero; the multiplicity inequality (5).
\item[(C2)] \PROVED{} \textbf{$H^{0}$}: $T_p$ by $(1+p)p^{a}$ on $\DkL^{a}\DkD^{b}$; $U_\ell=0$ on
  $H^{0}(M_n)$ at every grid cell \textbf{except $(n,\ell)=(120,5)$}, where $U_5$ has rank $1$ and
  $U_5\DkD^{6}=-\DkL^{20}$; $U_\ell^{2}=0$ everywhere ($n\ge1$); $T_p=p+1$ and $U_\ell=0$ on $H^{0}(\Qn)$.
\item[(C3)] \PROVED{} (given Claim~\ref{cl:dimH1}(1)) \textbf{Ramanujan mod $5$.} At $\ell=5$,
  $H^{1}(M_4)$ is one-dimensional (weight $6$, no Dickson class), with $T_p=1+p^{5}=1+p$ (the Eisenstein
  class of weight $6$, by the Eichler--Shimura identification that claim cites), and $\rank\delta=0$
  ($4\nmid6$), so $\thL$ embeds it in $H^{1}(M_{10})$ (dimension $3$) as a line with $T_p=p+p^{2}$.
  \STD{the congruence $\tau(n)\equiv n\sigma_1(n)\pmod5$, \cite[\S1]{swinnerton_dyer}} identifies
  $p+p^{2}$ with $\tau(p)$ mod $5$: the line lies in the generalised eigenspace of $\Delta$'s system, which
  is $2$-dimensional (the $\pm$ pair).
\item[(C4)] \PROVED{} (same basis) \textbf{Ramanujan mod $7$.} At $\ell=7$, $\thL$ embeds the Eisenstein
  line of $H^{1}(M_2)$ ($T_p=1+p^{3}$; $\rank\delta=0$, $6\nmid4$) in $H^{1}(M_{10})$ with $T_p=p+p^{4}$,
  which is $\tau(p)$ mod $7$ by $\tau(n)\equiv n\sigma_3(n)\pmod7$ (same citation).
\item[(C5)] \PROVED{} (same basis) \textbf{$\Delta$ mod $11$.} $\rank\delta=0$ at $(n,\ell)=(10,11)$
  ($10\nmid12$), so $\thL$ embeds $H^{1}(M_{10})$ (dimension $3$: Eisenstein $1+p^{11}=1+p$ and the $\Delta$
  pair) in $H^{1}(M_{22})$ (dimension $5$), with systems $p+p^{2}$ (once) and $p\,\tau(p)$ (twice);
  $\dim H^{1}(Q_{10})=2$. \emph{The whole of $H^{1}(M_{22})$} (v1.11.0-pre.2; the WP-3 referee's m-2,
  read on the shipped pipeline): $\thL H^{1}(M_{10})$ (dimension $3$) $\oplus$ the weight-$24$ Eisenstein
  class ($1+p^{23}=1+p^{3}$) $\oplus$ a second class with system $p+p^{2}$ (the $\pm$ twin of
  $\theta E_{12}$ mod $11$, a cusp form of filtration $24$). So the generalised eigenspace of $p+p^{2}$ in
  $H^{1}(M_{22})$ is \textbf{two}-dimensional and contains the $\thL$-image as a line; ``once'' above is
  the image, not the target.
\item[(C6)] \PROVED{} (same basis) \textbf{$\Delta$ mod $13$.} At $(10,13)$, $(\ell-1)=12\mid n+2$ and
  $156\nmid24$, so $\rank\delta=1$: \textbf{$\thL$ kills exactly the Eisenstein line} of $H^{1}(M_{10})$
  ($T_p=1+p^{11}=1+p^{-1}$; consistently $E_{12}\equiv1\pmod{13}$) and embeds the $\Delta$ pair in
  $H^{1}(M_{24})$ (dimension $3$) with system $p\,\tau(p)$; the third class there, $T_p=1+p$, maps onto
  $H^{1}(Q_{10})$ (dimension $1$).
\item[(C7)] \CONJ{} (it rests on Remark~\ref{rmk:thetaclassical}(1),(3) and on reading ``the least $n$
  at which a system occurs in $H^{1}(\Ga;M_n)$'' as $w-2$, a dictionary this document does not prove;
  the Dickson torsion of Claim~\ref{cl:bockstein} is where it could fail) \textbf{Filtration.} A system
  with $\det\bar\rho=\omega^{k-1}$ occurs only in weights $\equiv k\pmod{\ell-1}$. At $\ell=11$:
  $w(\Delta)=12\not\equiv0$, so $w(\theta\Delta)=24$ and $p\,\tau(p)$ is absent from $H^{1}(M_n)$ for
  $n\in\{2,12\}$ and present at $n=22$. At $\ell=13$: $w(\theta\Delta)=26$, so $p\,\tau(p)$ is absent at
  $n=12$ and present at $n=24$; $26\equiv0\pmod{13}$, so the next step drops: $p^{2}\tau(p)$ first occurs at
  $n=14$ or $n=26$ (filtration $16$ or $28$; $4$ is excluded, the system being irreducible), and the model
  decides which. \emph{Recorded at v1.11.0-pre.2} (the WP-3 referee's m-9, so that the answer is on record
  before WP-6 reports): the shipped pipeline gives $T_2|H^{1}(M_{14})$ at $\ell=13$ the eigenvalues
  $9^{1},8^{2},6^{1}$ and $T_3$ the eigenvalues $12^{1},6^{2},2^{1}$, with $8=2^{2}\tau(2)$ and
  $6=3^{2}\tau(3)$ modulo $13$: $p^{2}\tau(p)$ first occurs at $n=14$ (filtration $16$, the $s=1$ drop), with
  multiplicity $2$. The statement above stays the disjunction; this is the model's answer to it.
\item[(C8)] \CONJ{} (it rests on Remark~\ref{rmk:thetaclassical}(5), whose hypotheses are not all
  checked here) \textbf{Gross at $13$.} $\Delta$ mod $13$ is ordinary ($\tau(13)=-577738\equiv8$) of weight
  $12=\ell-1$; its companion would have weight $2$, and $\widetilde M_2=0$ at level~$1$. So Gross's
  criterion, if it applies, says $\bar\rho_\Delta|_{G_{13}}$ is \emph{not} split. On the model this is
  (C7)'s statement that $p\,\tau(p)$ first occurs at $n=24$, not earlier; nothing further is checkable.
  At $\ell=11$ Gross's theorem says nothing ($k=\ell+1$).
\item[(C9)] \PROVED{} \textbf{The Glover multiplicities at $\ell=5$, $n=34,\ldots,44$}: whether the
  multiplicity-$4$ ideals of the census at the cell $n$ (the \emph{target}) are $\thL$-images of
  multiplicity-$\ge4$ systems at $n'=n-6$ (the \emph{source}) is decided by (5) for $\theta_{n'}$ and by
  $\rank\delta$ of that map, i.e.\ at the source: $\rank\delta(n')=1$ exactly at $n'\in\{30,38\}$ (the
  targets $n=36$, $44$), where $4\mid n'+2$ and $20\nmid n'+6$; the explanation by \Espec{cl:glover}'s
  factor count is then confirmed or refuted cell by cell. \emph{Indexing} (v1.11.0; slot-11 WP-6's
  finding~1): through v1.11.0-pre.2 the parenthesis read ``$\rank\delta=1$ exactly at $n=38$ and $n=42$,
  where $4\mid n+2$ and $20\nmid n+6$'' --- $\rank\delta$ of $\theta_n$ \emph{out of} the cell, in
  Definition~\ref{def:theta}'s source indexing. That is true, but it is not the map the question is
  about; both readings are now named.
\end{enumerate}
The ordinary-part predictions at $\ell=691$ and $\ell=2411$ belong to the Hida chapter, which slot~11
does not write (owner ruling on decision~2, 2026-09-25).

\COMP\textsuperscript{r} \emph{Checked at v1.11.0} (slot-11 WP-5, WP-6. \emph{Provenance, v1.11.1} (the
slot-11 WP-15 referee's M-1): WP-6's predictions (its report's \S\S1--2 skeleton) were committed before
their runs; WP-5 records one deviation --- the two closed forms $\dim H^{0}(\Qn)$ and $\rank\delta$ were
checked on their first run, the predictions in hand but not committed before it --- and WP-6's \S1a
re-aligned (C5) and (C7) to v1.11.0-pre.2 after the first runs, (C5)'s assertion tightened from
``$\ge$'' to ``$=$''; every item held as finally stated. Through v1.11.0 this parenthesis read ``each
prediction was committed in writing before its run, and none was edited after it'', which the reports do
not support.) \code{HeckeDicksonThetaValidation} (\code{standard},
$108$ cells, $979$ checks) and \code{HeckeDicksonThetaCoverageSlowValidation} (\code{slow}, $404$ cells,
$\ell\in\{5,7,11,13\}$, even $n\le200$, $3590$ checks): (C1)'s exactness at all six spots as subspaces,
$H^{2}=0$ on the three terms, the closed forms of Theorem~\ref{thm:thetales}(3),(4), and $H^{0}$ spanned by
the $\DkL^{a}\DkD^{b}$. \code{HeckeDicksonThetaHeckeValidation} (\code{standard}, $5482$ checks; $34\,337$
under the read-only interlock \code{HECKE\_THETA\_IMAGE\_ALL=1}, on the $384$ census cells) and its
\code{slow} companion ($3836$ checks): Proposition~\ref{prop:thetatwist} at cochain level and on cohomology
(the $\tau$ datum identical at $n$ and $\nplus$, Lemma~\ref{lem:heckefibre}), the rest of (C1), and (C2) with
$(120,5)$'s exception entrywise. \textbf{Every item held}: (C3)--(C6) as stated, each matched to one system
of $\Sigma(\ell)$ and its twist (\Espec{sec:atlas}); (C7) at $\ell=11$ and $\ell=13$ as stated, the open
alternative decided by the model ($p^{2}\tau(p)$ first at $n=14$, as recorded at v1.11.0-pre.2); (C8)
consistent; (C9) the inequality on every block, and cell by cell the dimension-$4$ ideals at $n=34,38,42$
are partial images ($d_\theta=3$) and those at $n=40,44$ full ($d_\theta=4$). The composite $\thL^{\ell-1}$
on the $\Delta$ block at $n_0=10$ has rank $2$ after every step at every $\ell$ (at $\ell=5$ the possible
loss at $n=22$ does not occur). This is evidence; it changes no \PROVED{} or \CONJ{} tag above.
\end{remark}

%----------------------------------------------------------------------
\section{The ordinary projector}\label{sec:ordinary}
%----------------------------------------------------------------------

\paragraph{What this chapter fixes.}
The ordinary projector of an endomorphism of a finite-length module over $\Fl$ or $\Z/\ell^{e}$, as a
basis-free object, with the properties the $\theta$ line needs (``$\theta$-images are non-ordinary'')
and the ones the Hida line of slot~12 will need (functoriality, exactness, compatibility with
reduction). The precision question --- which $e$ certifies the ordinary part over $\Z_\ell$ --- is
stated as open. No hypothesis on $\ell$.

\begin{definition}[The ordinary projector]\label{def:ordinary}
Let $\ell$ be a prime, $e\ge1$, $M$ a finitely generated $\Z/\ell^{e}$-module (so of finite length
$\lambda(M)$ and finite as a set), and $U\in\End(M)$. Put
\[
  \eord=\eord(U):=\lim_{j\to\infty}U^{\,j!},
\]
the limit in the discrete topology --- the sequence is eventually constant
(Proposition~\ref{prop:ordinary}(1)). $\eord M$ is the \textbf{ordinary part} and $(1-\eord)M$ the
\textbf{non-ordinary part} of $(M,U)$; an element or submodule is \emph{non-ordinary} if $\eord$ kills
it. For $H^{q}$ of the model with $U=U_\ell$ this is the ordinary part in Hida's sense.
\end{definition}

\begin{proposition}[Existence, Fitting, and naturality]\label{prop:ordinary}
\PROVED{} With the notation of Definition~\ref{def:ordinary} and $\lambda=\lambda(M)$:
\begin{enumerate}
\item \textbf{(Fitting.)} $M=\im U^{\lambda}\oplus\ker U^{\lambda}$, both summands $U$-stable, $U$
  bijective on the first and $U^{\lambda}=0$ on the second. For every $j\ge j_0:=\max(\lambda,r)$, where
  $r$ is the order of $U|_{\im U^{\lambda}}$ in the finite group $\mathrm{Aut}(\im U^{\lambda})$,
  $U^{j!}$ is the projection onto $\im U^{\lambda}$ along $\ker U^{\lambda}$. So $\eord$ exists, is
  idempotent, and $\eord M=\im U^{\lambda}$, $(1-\eord)M=\ker U^{\lambda}=\{x:U^{i}x=0\text{ for some }i\}$.
\item \textbf{(Canonical.)} $\eord$ is a polynomial in $U$ (namely $U^{j_0!}$); it depends on no basis
  and commutes with every endomorphism of $M$ commuting with $U$. $\eord M$ is the largest $U$-stable
  submodule on which $U$ is bijective, and $(1-\eord)M$ the largest on which $U$ is nilpotent.
\item \textbf{(Functorial and exact.)} If $f\colon M\to M'$ is linear with $fU=U'f$, then
  $f\,\eord(U)=\eord(U')\,f$. A short exact sequence of such pairs stays exact after applying $\eord$
  (and after applying $1-\eord$).
\item \textbf{(Reduction.)} If $M_{e+1}$ is a $\Z/\ell^{e+1}$-module with $U_{e+1}$, and
  $\mathrm{red}\colon M_{e+1}\to M_{e}$ is a $U$-equivariant map to a $\Z/\ell^{e}$-module with $U_{e}$
  (for instance $M_e=M_{e+1}/\ell^{e}M_{e+1}$, or the reduction $H^{q}(\Z/\ell^{e+1})\to H^{q}(\Z/\ell^{e})$
  of \Espec{sec:ring} with $U=U_\ell$), then $\mathrm{red}\circ\eord(U_{e+1})=\eord(U_{e})\circ\mathrm{red}$.
\end{enumerate}
\end{proposition}

\begin{proof}
(1) The chains $\im U^{i}$ and $\ker U^{i}$ are monotone and stabilise by $i=\lambda$ (each strict
step changes a length by at least one). For $x\in M$, $U^{\lambda}x=U^{2\lambda}y$ for some $y$, so
$x-U^{\lambda}y\in\ker U^{\lambda}$; and $\im U^{\lambda}\cap\ker U^{\lambda}=0$ since $U^{\lambda}$ is
injective on $\im U^{\lambda}=\im U^{2\lambda}$. $U$ is surjective, hence bijective, on the finite set
$\im U^{\lambda}$. For $j\ge j_0$, $r\mid j!$ and $j!\ge\lambda$, so $U^{j!}$ is the identity on the first
summand and $0$ on the second. (2) is read off (1): a $U$-stable submodule on which $U$ is bijective lies
in $\bigcap_i\im U^{i}=\im U^{\lambda}$, one on which $U$ is nilpotent in $\ker U^{\lambda}$; a commuting
endomorphism commutes with $U^{j_0!}$. (3) $fU^{j!}=U'^{j!}f$ for every $j$; take $j$ large for both. For a
short exact sequence $0\to M'\to M\to M''\to0$ of pairs, $\eord$ is a natural idempotent on it, and the
image of a natural idempotent on an exact sequence is a direct summand of it as a complex, hence exact.
(4) is (3) for $f=\mathrm{red}$, a map of $\Z/\ell^{e+1}$-modules with $M_e$ viewed as one.
\end{proof}

\begin{corollary}[$\theta$-images and Dickson classes are non-ordinary]\label{cor:thetanonordinary}
\PROVED{} Assume \ref{H4}; $\cring=\Fl$, $U=U_\ell$.
\begin{enumerate}
\item $\eord(\im\thL)=0$ in $H^{q}(\Ga;M_{\nplus})$ for every $q$ and $n\ge0$ (Proposition~\ref{prop:thetatwist};
  $\im\thL\subseteq\ker U_\ell$ and $\eord$ kills $\ker U_\ell$).
\item $\eord=0$ on $H^{0}(\Ga;M_m)$ for every $m\ge1$ (Proposition~\ref{prop:heckeH0}(2): $U_\ell^{2}=0$), and
  on $H^{0}(\Ga;\Qn)$ (Lemma~\ref{lem:Qnfunctions}(3)). At $m=0$, $\eord=\id$ on $H^{0}(\Ga;\Fl)$.
\item If $\ell\ge5$ and $m\ge1$: $U_\ell$ is nilpotent on the $\ell$-torsion $H^{1}(\Ga;\Sym^{m}\Z^{2})[\ell]$,
  so this subgroup is non-ordinary --- by the $U_\ell$-equivariant Bockstein
  $H^{0}(\Ga;M_m)\cong H^{1}(\Ga;\Sym^{m}\Z^{2})[\ell]$ of Claim~\ref{cl:bockstein}(2),(3) and (2) here.
  $U_\ell$ is zero on it for $m<\ell^{3}-\ell$ and not zero at $m=\ell^{3}-\ell$.
\end{enumerate}
\end{corollary}

\begin{remark}[The ordinary part is torsion-free at level~$1$]\label{rmk:ordtorsionfree}
\PROVED{} (v1.11.1; the slot-11 WP-15 referee's ruling on the author's finding~F-4.) Let $\ell\ge5$,
$n\ge1$, $H:=H^{1}(\Ga;\Sym^{n}\Z^{2})\otimes\Z_\ell$ with $U=U_\ell$, and $T:=H_{\mathrm{tors}}$, a finite
$U$-stable $\ell$-group. On $H$ put $\eord:=\varprojlim_{e}\eord(U\bmod\ell^{e})$ on the $H/\ell^{e}H$
(compatible by Proposition~\ref{prop:ordinary}(4)); on $T$, a module of finite length,
Definition~\ref{def:ordinary} applies directly, and the two agree on $T$ (Proposition~\ref{prop:ordinary}(3)
for $T\hookrightarrow H/\ell^{e}H$, $\ell^{e}T=0$, which is injective since $T\cap\ell^{e}H=0$ for such $e$).
Then:
\begin{enumerate}
\item $\eord(T)=0$: the torsion of $H$ is non-ordinary.
\item $\eord H$ is torsion-free, and the quotient map $H\to H/T$ restricts to an isomorphism
  $\eord H\cong\eord(H/T)$ commuting with every $T_p$: \textbf{the ordinary part of $H$ is the ordinary
  part of its free quotient}, so its rank and the characteristic polynomials of the $T_p$ on it are read
  on the free quotient.
\end{enumerate}
\emph{Proof.} (1) Put $S=\eord(T)$. By Proposition~\ref{prop:ordinary}(1),(2) $U$ is bijective on $S$, hence
on the $U$-stable subgroup $S[\ell]$, which is finite. By Corollary~\ref{cor:thetanonordinary}(3) $U$ is
nilpotent on $T[\ell]=H^{1}(\Ga;\Sym^{n}\Z^{2})[\ell]\otimes\Z_\ell\supseteq S[\ell]$. An endomorphism that is
both bijective and nilpotent exists only on the zero group, so $S[\ell]=0$, and $S=0$, since a non-zero finite
$\ell$-group has non-zero $\ell$-torsion. (2) $\eord$ is $\Z_\ell$-linear, so $\eord(T)\subseteq T$ and
$(\eord H)_{\mathrm{tors}}=\eord H\cap T=\eord(T)=0$ (an $x=\eord x\in T$ lies in $\eord(T)$). For $e$ with
$\ell^{e}T=0$ the sequence $0\to T\to H/\ell^{e}H\to(H/T)/\ell^{e}\to0$ is exact and $U$-equivariant, so
$\eord$ keeps it exact (Proposition~\ref{prop:ordinary}(3)): $\eord(H/\ell^{e}H)\to\eord((H/T)/\ell^{e})$ is
an isomorphism, its kernel being $\eord(T)=0$ by (1). Passing to the limit over $e$ gives
$\eord H\cong\eord(H/T)$; the $T_p$ commute with $U$, hence with $\eord$ (Proposition~\ref{prop:ordinary}(2)).
\qed

\emph{What this does and does not settle.} Question~\ref{open:ordprecision} asks two things: whether the
ordinary part is torsion-free, and which precision, computable from the cell before the run, determines
it. This remark answers the first at $\ell\ge5$, $n\ge1$, from statements of this document alone; the
second --- a per-cell bound known before the run --- stays open (v1.12.0-pre.1: answered at $\ell\ge5$,
$n\ge1$ by Corollary~\ref{cor:ordprecision}, the item's status pending the release). Slot-11 WP-18c's idempotent on the free
block (Question~\ref{open:ordprecision}'s v1.11.0 status note) is therefore the whole ordinary part of
$H^{1}\otimes\Z/\ell^{k}$'s free reading, not a part of it. At $\ell\in\{2,3\}$ nothing is asserted
(Corollary~\ref{cor:thetanonordinary}(3) is stated at $\ell\ge5$).
\end{remark}

\begin{remark}[The Eisenstein line one level up is a theorem]\label{rmk:ordinaryeis}
\PROVED{} Let $\ell\ge5$, $e\ge1$ and let $p\neq\ell$ be a prime whose class generates
$(\Z/\ell^{e+1})^{\times}$ --- for instance a primitive root modulo $\ell^{2}$, which then generates modulo
every $\ell^{e+1}$. Then for weights $k,k'$,
$1+p^{k-1}\equiv1+p^{k'-1}\pmod{\ell^{e+1}}$ if and only if $k\equiv k'\pmod{(\ell-1)\ell^{e}}$: the
Eisenstein eigenvalues of two weights congruent modulo $(\ell-1)\ell^{e-1}$ but not modulo
$(\ell-1)\ell^{e}$ \emph{do} differ modulo $\ell^{e+1}$ at such a $p$. (The order of $p$ is
$(\ell-1)\ell^{e}$.) On a cuspidal ordinary part the corresponding failure is an observation, not a
consequence of this remark (the plan's audit, finding m-13); nothing here predicts it.
\end{remark}

\begin{openitem}[Certified precision for the ordinary part --- ANSWERED 2026-09-26]\label{open:ordprecision}
\OPENSTAT{} For a cell $(n,\ell)$, which exponent $e$ --- computable from the cell, before the run --- suffices
for $\eord(U_\ell)$ on $H^{1}(\Ga;\Sym^{n}\Z^{2}\otimes\Z/\ell^{e})$ to determine the ordinary part of
$H^{1}(\Ga;\Sym^{n}\Z^{2}\otimes\Z_\ell)$ (its rank, its torsion, and the characteristic polynomials of the
$T_p$ on it modulo a target $\ell^{e'}$)? What is known here: Proposition~\ref{prop:ordinary}(4) makes the
projectors compatible along reductions; Corollary~\ref{cor:thetanonordinary}(3) puts the socle of the
Dickson torsion in the non-ordinary part; Remark~\ref{rmk:bocksteinexp}(2) shows the torsion exponents are
unbounded along $n=\ell^{j}m$, so no $e$ independent of $n$ reads every exponent
(Proposition~\ref{prop:certifiede}). Whether the ordinary part is torsion-free at level~$1$, and a
certified per-cell bound, are \textbf{open}; they are the Hida chapter's (slot~12), together with the
engine's per-cell precision rule (\file{src/ModArithmetic/ENGINE\_API.md} \S13.11.11; the
pointer read \S13.12, the $128$-bit draft, through v1.11.0-pre.2).
\emph{Status at v1.11.0} (slot-11 WP-18c, \Espec{sec:ring}, the digit route): a per-cell certified
precision bound for the \emph{free} part is built --- $K=k+e_{\max}$, with $e_{\max}$ read from the
certified $K$-mode at word size (Proposition~\ref{prop:certifiede}(3)) and the stopping rule re-evaluated
at $K$ --- so that the characteristic polynomials of the $T_p$ on $\Hfr{1}\otimes\Z/\ell^{k}$ are
certified to precision $k$; and at $p=\ell$ the idempotent of Definition~\ref{def:ordinary} is built on
that free block as a Fitting idempotent with four certificates. $e_{\max}$ is found by a run at the cell,
not ``computable from the cell, before the run'', and nothing is said about the torsion's ordinary part;
the question above is therefore \textbf{not} answered, and the item stays \OPENSTAT.
\emph{Status at v1.11.1} (the slot-11 WP-15 referee's ruling on F-4): the first half --- whether the
ordinary part is torsion-free at level~$1$ --- is answered at $\ell\ge5$, $n\ge1$ by
Remark~\ref{rmk:ordtorsionfree}: it is, and it equals the ordinary part of the free quotient. The sentence
``What is known here'' above stopped one step short of that corollary. \textbf{What stays open} is the
second half alone: an exponent computable from the cell \emph{before} the run ($e_{\max}$ is read from a
run). The item stays \OPENSTAT, narrowed to that bound; its question text is kept as written (one version,
one text).
\emph{Status at v1.12.0-pre.1} (slot-12 WP-2, \S\ref{sec:hida}): \textbf{answered at $\ell\ge5$, $n\ge1$,
pending release.} By Corollary~\ref{cor:ordprecision}(1) the exponent asked for is the target itself,
$e=e'$ (the rank already at $e=1$; v1.12.0-pre.2, R3): the torsion being non-ordinary (Remark~\ref{rmk:ordtorsionfree}), $\eord(U_\ell)$ on
$H^{1}(\Ga;\Sym^{n}\Z^{2}\otimes\Z/\ell^{e'})$ is the ordinary part of $H^{1}\otimes\Z_\ell$ reduced
modulo $\ell^{e'}$ --- its rank (read already over $\Fl$), its torsion ($0$) and the characteristic
polynomials of the $T_p$ modulo $\ell^{e'}$. For the whole free quotient (slot-11 WP-18c's free block) the
working precision is $K=e'+e_{\max}$ (Corollary~\ref{cor:ordprecision}(2)--(4)), with $e_{\max}$ read before
the run from $\rhon(S)$, $\rhon(T)$ by Proposition~\ref{prop:torsionprepass}. \textbf{The status macro stays
\OPENSTAT{} until the owner-signed release v1.12.0} (slot-12 plan \S5 ruling L3, \S6 item~1): changing a
binding item's status is a release decision. At $\ell\in\{2,3\}$ nothing is asserted.
\emph{Status at v1.12.0} (owner-signed release, 2026-09-26; slot-12 plan \S5 ruling L3, \S6 item~1):
\textbf{the item is answered, and the tag moves \OPENSTAT{} $\to$ \PROVED{}.} Corollary~\ref{cor:ordprecision}
(with Proposition~\ref{prop:torsionprepass} and Remark~\ref{rmk:ordtorsionfree}) answers the question in full
at $\ell\ge5$, $n\ge1$: $e=e'$ for the ordinary part, $K=e'+e_{\max}$ for the free quotient, with $e_{\max}$
computable before the run from $\rhon(S)$, $\rhon(T)$ alone. At $\ell\in\{2,3\}$ nothing is asserted. The
question text and the boxed statement above are unchanged (rules~\ref{L1}, \ref{L2}); only the status tag
moves, as \ref{open:ell2}'s did at v1.3.0.
\end{openitem}

\paragraph{Code map.}
Not built at v1.11.0-pre.1. $\DkL\cdot$ and \eqref{eq:dicksontwist}'s certificate are
\file{HeckeExt/HeckeDicksonMultiplication} (slot~5). Definition~\ref{def:Qn}, Theorem~\ref{thm:thetales} and
Proposition~\ref{prop:heckeH0}(1) are slot-11 WP-5's (\code{FibreRepresentation::QuotientL}, the six terms and
five maps, \code{HeckeDicksonThetaValidation}); Proposition~\ref{prop:thetatwist}, Proposition~\ref{prop:heckeH0}(2)--(4)
and Remark~\ref{rmk:thetacheckable} WP-6's. Definition~\ref{def:ordinary} has no implementation in slot~11
(decision~2 not taken) (through v1.11.0-pre.2; v1.11.1, the WP-15 referee's m-3: superseded by the
``Built at v1.11.0'' paragraph below).
\emph{Built at v1.11.0.} \file{LocalSystemUtils} (\code{FibreRepresentation::QuotientL}: $\rho_{\Qn}$ on
Definition~\ref{def:Qn}'s basis for every $2\times2$ matrix, singular included), \file{LocalSystemSL2Z}
(the relation list of $\Qn$), \file{Hecke/HeckeOperatorOnK} (its \code{QuotientL} arm, on the coset and
$\tau$ tables of $\Sym^{n}$), \file{HeckeExt/HeckeDicksonTheta} (\code{DicksonThetaCell}: the six terms,
the five maps in the stored bases, $\delta$ by the snake construction, the Dickson span; slot-11 WP-5) and
\file{HeckeExt/HeckeDicksonThetaHecke} (the Hecke operators on the three terms and the checks of
Proposition~\ref{prop:thetatwist}, Theorem~\ref{thm:thetales}(2),(3),(5), Proposition~\ref{prop:heckeH0}(2)--(4)
and Remark~\ref{rmk:thetacheckable}; slot-11 WP-6); validations as in Remark~\ref{rmk:thetacheckable}'s last
paragraph. $\ell\in\{2,3\}$ is refused by name (Remark~\ref{rmk:thetasmallell}), and so is a non-default
coefficient ring. (v1.11.1: since slot-11 WP-17 the $\ell\in\{2,3\}$ refusal is the mathematical
precondition \Espec{sec:thetaobj}~(h) specifies --- \code{ENGINE\_CHECK\_MSG} naming
Lemma~\ref{lem:h2torsion}, a \code{std::logic\_error} (console exit~3) --- and no longer
\code{hecke::NotImplementedInThisBuild} (exit~5); the WP-15 referee's ruling on F-2.) Definition~\ref{def:ordinary}: built at $p=\ell$ on the free block of $H^{1}$ over
$\Z/\ell^{k}$, $\ell^{k}\ge2^{64}$, by \file{Hecke/LadicCell} (slot-11 WP-18c; \Espec{sec:ring}, the digit
route, whose record key is \code{hida}); \code{HeckeLadicCellValidation} certifies it at $(10,691)$, where
$U_{691}$ is ordinary on the whole free block, and at $(30,5)$, a genuine Fitting split ($U_5$ with a
nilpotent part of dimension~$4$, $\rank_{\F_5}e=1$).

\paragraph{Formalization note.}
Lemma~\ref{lem:dicksontwist}, Proposition~\ref{prop:heckeH0} and Lemma~\ref{lem:Qnfunctions} are polynomial identities
over $\Fl$ and finite counts; Proposition~\ref{prop:ordinary} is Fitting's lemma. Theorem~\ref{thm:thetales} needs
the long exact sequence of a short exact sequence of complexes and nothing about $\Ga$ beyond
Proposition~\ref{prop:bredon}(2). All are cheap; Dickson's theorem (Proposition~\ref{prop:dicksonH0}(1)) is the one
import.

%----------------------------------------------------------------------
\section{Hida theory at level one: foundation}\label{sec:hida}
%----------------------------------------------------------------------
% NEW at v1.12.0-pre.1 (hecke-extensions slot 12, WP-2; plan
% docs/implementation/slot12/00-slot12-plan.md S2 "WP-2", the "contract first" package, rule 21;
% the predictions of rmk:hidacheckable committed before WP-4 / WP-5's code, rule 20).
% Placed after sec:ordinary, before the appendices, so that no chapter number moves.

\paragraph{What this chapter fixes.}
The three statements every later Hida step needs, and the predictions its first computations test.
(i) The $\ell$-primary torsion of $H^{1}(\Ga;\Vcoef)$ is read \emph{before the run} from the two matrices
$\rhon(S)$, $\rhon(T)$ alone (Proposition~\ref{prop:torsionprepass}); with Remark~\ref{rmk:ordtorsionfree}
this answers Question~\ref{open:ordprecision} at $\ell\ge5$, $n\ge1$ (Corollary~\ref{cor:ordprecision}).
(ii) The ordinary projector of Definition~\ref{def:ordinary} is one explicit power of $U$
(Lemma~\ref{lem:ordexponent}). (iii) The weight towers (Definition~\ref{def:hidatower}) and what is
predicted on them (Remark~\ref{rmk:hidacheckable}). Throughout $\Ga=\SL_2(\Z)$,
$\Vcoef=\Sym^{n}\Z^{2}$ with the action of Definition~\ref{def:rhon}, $n\ge1$, and $H^{\bullet}(\Ga;-)$
is \emph{group} cohomology; the model enters only through Remark~\ref{rmk:orbifold} at $\ell\ge5$. The
exponent of a coefficient ring is $e$ (\S\ref{sec:notation-primes}); a \emph{target} precision is written
$e'$, as in Question~\ref{open:ordprecision} --- it is \Espec{sec:ring}'s $k$, which this document may
not write, $k=n+2$ being the weight. The plan and [X] write the predictions (H1)--(H6) and the
generator count of Lemma~\ref{lem:ordexponent} $r$; here they are (Hi1)--(Hi6) and $\mathfrak g$, since
\ref{H1}--\ref{H4} are the standing hypotheses and $r$ is the restriction map (rule~\ref{A5}).

\begin{proposition}[The torsion pre-pass]\label{prop:torsionprepass}
\PROVED{} Let $\ell$ be a prime and $n\ge1$. Let $J_n\in\Mat_{2(n+1)\times(n+1)}(\Z)$ be the matrix
$\rhon(S)-1$ stacked over $\rhon(T)-1$ (in the basis $(u_j)$ of Definition~\ref{def:rhon}). Write the
$\ell$-primary torsion of $H^{1}(\Ga;\Vcoef)$ as $\bigoplus_{t=1}^{s}\Z/\ell^{a_t}$, $a_t\ge1$, put
$e_{\max}:=\max_t a_t$ ($e_{\max}:=0$ if $s=0$), and for $e\ge0$
\[
  \lambda_e:=\lambda\bigl(H^{0}(\Ga;\Vcoef\otimes\Z/\ell^{e})\bigr)\qquad(\lambda_0=0),
\]
$\lambda$ the length. Then:
\begin{enumerate}
\item \textbf{(The $e$-fold Bockstein.)} For every $e\ge1$ the connecting map of
  $0\to\Vcoef\xrightarrow{\ell^{e}}\Vcoef\to\Vcoef\otimes\Z/\ell^{e}\to0$ is an isomorphism
  \[
    H^{0}(\Ga;\Vcoef\otimes\Z/\ell^{e})\;\xrightarrow{\ \sim\ }\;H^{1}(\Ga;\Vcoef)[\ell^{e}],
    \qquad\text{hence}\qquad \lambda_e=\sum_{t=1}^{s}\min(a_t,e).
  \]
  At $e=1$ and $\ell\ge5$ this is the second isomorphism of Claim~\ref{cl:bockstein}(2), which it
  extends: the argument is the same, and it needs only $H^{0}(\Ga;\Vcoef)=0$.
\item \textbf{(Two matrices suffice.)} $H^{0}(\Ga;\Vcoef\otimes\Z/\ell^{e})=\ker(J_n\bmod\ell^{e})$, and
  $\rank_{\Q}J_n=n+1$. If $\ell^{b_1},\dots,\ell^{b_{n+1}}$ are the elementary divisors of $J_n$ over
  $\Z_{(\ell)}$, the multiset of the positive $b_i$ \emph{is} the multiset of the $a_t$. In particular
  $s=\#\{i:b_i>0\}$ and $e_{\max}=\max_i b_i$; at $\ell\ge5$, $s=\dim_{\Fl}H^{0}(\Ga;\Mn)$ is the Dickson
  count of Claim~\ref{cl:dimH1}(2), so $e_{\max}=0$ whenever $n<\ell+1$ (the hypothesis $\ell\ge5$ is
  needed here: at $\ell=2$, $n=2$ has $\Z/2$, the invariant $x^{2}+xy+y^{2}$ of
  Remark~\ref{rmk:bocksteinexp}(2), as $\ell^{2}-\ell<\ell+1$; v1.12.0-pre.2, R2).
\item \textbf{(The stopping rule.)} $\lambda_{e+1}-\lambda_e=\#\{t:a_t>e\}=(n+1)-r_{e+1}(J_n)$, with
  $r_{e}$ as in Proposition~\ref{prop:certifiede}; it is non-increasing in $e$ and vanishes exactly for
  $e\ge e_{\max}$. Hence
  \[
    e_{\max}=\min\{e\ge0:\lambda_{e+1}=\lambda_e\}=\min\{e\ge0:r_{e+1}(J_n)=n+1\},
  \]
  the minimum exists, and the number of summands $\Z/\ell^{j}$ is
  $(\lambda_j-\lambda_{j-1})-(\lambda_{j+1}-\lambda_j)$ for $j\ge1$. The rule reads $J_n$ over
  $\Z/\ell^{e}$ for $e=1,\dots,e_{\max}+1$ and needs no cochain complex.
\item \textbf{(Hecke.)} The isomorphism of (1) is equivariant for every double-coset operator
  $\cor\circ(\conj_g)_{\ast}\circ\res$, $g\in\De$ --- $U_\ell$'s $g=\mathrm{diag}(1,\ell)$ included ---
  and, under $6\in\cring^{\times}$, for $\Tdc{g}$ (Theorem~\ref{thm:classical}) and every operator of
  \S\ref{sec:heckedatum} built from it.
\item \textbf{(Integrally.)} The torsion subgroup of $H^{1}(\Ga;\Vcoef)$ is isomorphic to
  $\bigoplus_{i}\Z/d_i$, $d_1\mid\cdots\mid d_{n+1}$ the elementary divisors of $J_n$ over $\Z$.
\item \textbf{(Scope.)} (1)--(5) are statements about group cohomology and hold at \emph{every} prime,
  $2$ and $3$ included. What needs $\ell\ge5$ is their \emph{use} as a precision for the model:
  there $H^{1}$ of $\Ceq(\Z/\ell^{e})$ is $H^{1}(\Ga;\Vcoef\otimes\Z/\ell^{e})$ (Remark~\ref{rmk:orbifold},
  Definition~\ref{def:ringsemantics}(2)), and $H^{1}(\Ceq(\Z))$ agrees with $H^{1}(\Ga;\Vcoef)$ after
  inverting $6$ (Proposition~\ref{prop:isomodel}(4)), so its $\Z_{(\ell)}$-type at $\ell\ge5$ is that of
  $H^{1}(\Ga;\Vcoef)$ (v1.12.0-pre.2, the WP-3 referee's R1: through pre.1 this read ``$H^{1}$ of
  $\Ceq(\Z)$ and of $\Ceq(\Z/\ell^{e})$ is group cohomology''); hence $e_{\max}$ is the largest exponent of
  the $\Z_{(\ell)}$-type of $H^{1}$ that Proposition~\ref{prop:certifiede}(3)'s $K$-mode reads from a run; and $H^{2}(\Ga;\Vcoef)$ has no
  $\ell$-torsion (Claim~\ref{cl:bockstein}(1)), so $H^{1}(\Ga;\Vcoef\otimes\Z/\ell^{e})=H^{1}(\Ga;\Vcoef)\otimes\Z/\ell^{e}$.
  At $\ell\in\{2,3\}$ neither holds --- the model computes the Bredon object
  (Proposition~\ref{prop:bredon}; Proposition~\ref{prop:certifiede}(5)), and $H^{2}(\Ga;\Vcoef)$ is
  killed by $24$ (Claim~\ref{cl:amalgam}(4)) but not asserted to vanish --- so (1)--(5) give the $2$- and
  $3$-primary torsion of $H^{1}(\Ga;\Vcoef)$ and \textbf{no} precision for the model, and
  Corollary~\ref{cor:ordprecision} is not asserted there.
\end{enumerate}
\end{proposition}

\begin{proof}
(1) $H^{0}(\Ga;\Vcoef)=\Vcoef^{\Ga}=0$ for $n\ge1$ by the argument of Claim~\ref{cl:bockstein}(1), which
takes place over $\Q$ and uses no hypothesis on $\ell$. The sequence
$0\to\Vcoef\xrightarrow{\ell^{e}}\Vcoef\to\Vcoef\otimes\Z/\ell^{e}\to0$ is a short exact sequence of
$\Z[\De]$-modules ($\Vcoef$ is $\Z$-free and $\ell^{e}$ is central), so it has a long exact cohomology
sequence \STD{\cite[Ch.~III, \S\S0--1]{brown}}, whose first terms
\[
  0=H^{0}(\Ga;\Vcoef)\to H^{0}(\Ga;\Vcoef\otimes\Z/\ell^{e})\xrightarrow{\ \delta\ }
  H^{1}(\Ga;\Vcoef)\xrightarrow{\ \ell^{e}\ }H^{1}(\Ga;\Vcoef)
\]
say that $\delta$ is injective with image $\ker\ell^{e}=H^{1}(\Ga;\Vcoef)[\ell^{e}]$. $H^{1}(\Ga;\Vcoef)$
is finitely generated: $\Ga$ is generated by $S$ and $T$ (Claim~\ref{cl:amalgam}(3) presents it on $S$
and $ST$), so a $1$-cocycle is determined by its values at $S$ and $T$ and $Z^{1}$ embeds in
$\Vcoef\oplus\Vcoef$. Write $H^{1}(\Ga;\Vcoef)\cong\Z^{h}\oplus F\oplus\bigoplus_t\Z/\ell^{a_t}$ with $F$
finite of order prime to $\ell$; the $\ell^{e}$-torsion of $\Z/\ell^{a}$ is $\ell^{a-\min(a,e)}\Z/\ell^{a}$,
of length $\min(a,e)$, and $\Z^{h}$, $F$ have none. Length is additive.

(2) Since $S$ and $T$ generate $\Ga$, an element of $\Vcoef\otimes\Z/\ell^{e}$ is $\Ga$-invariant iff
$\rhon(S)$ and $\rhon(T)$ fix it, i.e.\ iff it lies in $\ker(J_n\bmod\ell^{e})$. Over $\Q$ the kernel of
$J_n$ is $\Vcoef_{\Q}^{\Ga}=0$ (the argument of Claim~\ref{cl:bockstein}(1)), so $\rank_{\Q}J_n=n+1$. Over the
local principal ideal domain $\Z_{(\ell)}$, $J_n=P\left(\begin{smallmatrix}\Lambda\\0\end{smallmatrix}\right)Q$
with $P,Q$ invertible and $\Lambda=\mathrm{diag}(\ell^{b_1},\dots,\ell^{b_{n+1}})$, every $b_i$ finite as the
rank is full. Reduction modulo $\ell^{e}$ keeps $P,Q$ invertible, so
$\ker(J_n\bmod\ell^{e})\cong\bigoplus_i\ker(\ell^{b_i}\colon\Z/\ell^{e}\to\Z/\ell^{e})$, of length
$\sum_i\min(b_i,e)$. With (1): $\sum_t\min(a_t,e)=\sum_i\min(b_i,e)$ for every $e\ge0$; the difference
between $e$ and $e-1$ gives $\#\{t:a_t\ge e\}=\#\{i:b_i\ge e\}$ for every $e\ge1$, and these counts
determine both multisets of positive entries. At $\ell\ge5$, $s=\dim_{\Fl}H^{1}(\Ga;\Vcoef)[\ell]=\lambda_1=
\dim_{\Fl}H^{0}(\Ga;\Mn)$, which is Claim~\ref{cl:dimH1}(2)'s count, zero for $0<n<\ell+1$.

(3) By (1), $\lambda_{e+1}-\lambda_e=\sum_t\bigl(\min(a_t,e+1)-\min(a_t,e)\bigr)=\#\{t:a_t\ge e+1\}$. By (2) and
Proposition~\ref{prop:certifiede}(1) for $D=J_n$ (whose proof uses only the Smith form over $\Z_{(\ell)}$),
$r_{e+1}(J_n)=\#\{i:b_i<e+1\}$, so $(n+1)-r_{e+1}(J_n)=\#\{i:b_i\ge e+1\}$, the same count. It is
non-increasing and zero iff $e\ge\max_t a_t$; the multiplicity of $j$ among the $a_t$ is
$\#\{a_t\ge j\}-\#\{a_t\ge j+1\}$.

(4) As in Claim~\ref{cl:bockstein}(3): the sequence is one of $\Z[\De]$-modules, and $\res$,
$(\conj_g)_{\ast}$ and $\cor$ are morphisms of $\delta$-functors in the coefficient module, so their
composite commutes with $\delta$. Theorem~\ref{thm:classical} identifies the composite with $\Tdc{g}$ under
$6\in\cring^{\times}$.

(5) The torsion subgroup of the finitely generated group $H^{1}(\Ga;\Vcoef)$ is finite and is the product of
its $\ell$-primary parts; the Smith form of $J_n$ over $\Z$ localises to the one over $\Z_{(\ell)}$, with
elementary divisors $\ell^{v_\ell(d_i)}$, so by (2) the $\ell$-primary part is
$\bigoplus_i\Z/\ell^{v_\ell(d_i)}$, the $\ell$-primary part of $\bigoplus_i\Z/d_i$, for every $\ell$.

(6) is read off the cited statements; the identification of the model's $H^{1}(\Z/\ell^{e})$ with
$H^{1}(\Ga;\Vcoef)\otimes\Z/\ell^{e}$ at $\ell\ge5$ is the degree-$1$ half of the sequence in (1), whose next
term is $H^{2}(\Ga;\Vcoef)[\ell^{e}]=0$.
\end{proof}

\noindent
\emph{Re-derivation note} (not a tag; an out-of-tree script, recorded in slot-12 WP-2's report). The elementary divisors of $J_n$ over $\Z_{(\ell)}$,
computed exactly modulo $\ell^{8}$, reproduce Remark~\ref{rmk:bocksteinexp}(3)'s table --- $\Z/25$ at
$(30,5)$, $\Z/25\oplus\Z/5$ at $(60,5)$, $\Z/125\oplus(\Z/5)^{2}$ at $(150,5)$, $\Z/49$ at $(56,7)$, one
$\Z/25$ among the classes at $(90,5)$, $(100,5)$, $(120,5)$ --- and, by clause~(5) at $\ell\in\{2,3\}$, the
$2$- and $3$-primary torsion of $H^{1}$ in [IB] Appendix~A at $n\in\{2,4,6,8,12,16,18,20\}$. This is
evidence for the statement, not part of its proof; the registered check is slot-12 WP-4's.

\begin{corollary}[The certified precision, before the run]\label{cor:ordprecision}
\PROVED{} Let $\ell\ge5$ and $n\ge1$. Let $H:=H^{1}(\Ga;\Vcoef)\otimes\Z_\ell$ with $U=U_\ell$ and $\eord$ as
in Remark~\ref{rmk:ordtorsionfree}, $h:=\rank_{\Z_\ell}H$ --- which is $2d_k-1$ for even $n$
(Claim~\ref{cl:dimH1}(1)) and $0$ for odd $n$, so it too is known before the run; (3) below reads it off a
cyclic decomposition as a check, not as an input (v1.12.0-pre.2, R4) --- and let $e_{\max}$ be
Proposition~\ref{prop:torsionprepass}'s --- computed from $\rhon(S)$, $\rhon(T)$ alone, before any cochain
complex is built. For $e\ge1$ put $H_e:=H^{1}(\Ga;\Vcoef\otimes\Z/\ell^{e})$; then $H_e=H/\ell^{e}H$,
compatibly with every double-coset operator (Proposition~\ref{prop:torsionprepass}(4),(6)), and $H_e$ is
what the model computes over $\Z/\ell^{e}$ (Definition~\ref{def:ringsemantics}(2)). Fix a target
precision $e'\ge1$.
\begin{enumerate}
\item \textbf{(The ordinary part needs no precision beyond the target.)} The image of
  $\eord(U|_{H_{e'}})$ is the image of $\eord H$ in $H_{e'}$, and
  $\eord(U|_{H_{e'}})H_{e'}\cong\eord H/\ell^{e'}\eord H$, compatibly with every $T_p$: a free
  $\Z/\ell^{e'}$-module of rank $h_{\mathrm{ord}}:=\rank_{\Z_\ell}\eord H$ on which $T_p$ acts by the
  reduction of its action on $\eord H$. So the rank of the ordinary part is
  $\dim_{\Fl}\eord(U_\ell)H^{1}(\Ga;\Mn)$ (at $e'=1$), its torsion is $0$
  (Remark~\ref{rmk:ordtorsionfree}), and the characteristic polynomial of $T_p$ on it modulo $\ell^{e'}$
  is that of $T_p$ on $\eord(U|_{H_{e'}})H_{e'}$. \textbf{The exponent Question~\ref{open:ordprecision}
  asks for is $e=e'$} (the rank and the torsion already at $e=1$); it depends on the cell through nothing.
  (v1.12.0-pre.2, R3: through pre.1 this read ``$e=\max(1,e')$'', vacuous under $e'\ge1$.)
\item \textbf{(The free quotient needs $K=e'+e_{\max}$.)} Put $K:=e'+e_{\max}$. The quotient
  $H\to H/H_{\mathrm{tors}}$ induces a surjection $H_K\to(H/H_{\mathrm{tors}})/\ell^{e'}(H/H_{\mathrm{tors}})$, compatible with every $T_p$, whose kernel
  is $H_K[\ell^{e_{\max}}]$; and multiplication by $\ell^{e_{\max}}$ gives
  $H_K/H_K[\ell^{e_{\max}}]\cong\ell^{e_{\max}}H_K$. So
  \[
    \ell^{e_{\max}}H_K\;\cong\;H_K/H_K[\ell^{e_{\max}}]\;\cong\;\Hfr{1}\otimes\Z/\ell^{e'},
  \]
  free of rank $h$ over $\Z/\ell^{e'}$, compatibly with every $T_p$ ($U_\ell$ included): the
  characteristic polynomials of the $T_p$ on the free quotient are read modulo $\ell^{e'}$ on $H_K$, with
  $K$ fixed before the run.
\item \textbf{(Any cyclic decomposition.)} Every decomposition of $H_K$ into cyclic summands has exactly
  $h$ summands of order $\ell^{K}$ and the others of order $\le\ell^{e_{\max}}$; let $F'$ be the sum of the
  former, $T'$ of the latter, $\pi_{F'}$ the projection along $T'$. Then $F'/\ell^{e'}F'\to
  H_K/H_K[\ell^{e_{\max}}]$ is an isomorphism, and for Hecke operators $A,B$ with $F'$-blocks
  $A_{F'}:=\pi_{F'}A|_{F'}$: $A_{F'}\bmod\ell^{e'}$ is the operator (2) induces, and
  $(AB)_{F'}\equiv A_{F'}B_{F'}\pmod{\ell^{e'}}$. This is the free-block reading of slot-11 WP-18c
  (Question~\ref{open:ordprecision}'s v1.11.0 note), whose $e_{\max}$ comes from a run; (2) makes the same
  $K$ available before it.
\item \textbf{(The two readings agree.)} $\eord$ of the operator $U$ induces on the module of (2) is the
  reduction of $\eord(U|_{H/H_{\mathrm{tors}}})$, and its image is isomorphic to $\eord H/\ell^{e'}\eord H$, compatibly with
  every $T_p$ --- the ordinary part of (1).
\end{enumerate}
\end{corollary}

\begin{proof}
$H_e=H/\ell^{e}H$: the degree-$1$ part of the long exact sequence of
Proposition~\ref{prop:torsionprepass}(1) is $0\to H^{1}(\Ga;\Vcoef)/\ell^{e}\to H_e\to
H^{2}(\Ga;\Vcoef)[\ell^{e}]\to0$ and the last term is $0$ (Claim~\ref{cl:bockstein}(1)); tensoring the
finitely generated $H^{1}(\Ga;\Vcoef)$ with $\Z_\ell$ does not change it modulo $\ell^{e}$. Equivariance
is Proposition~\ref{prop:torsionprepass}(4)'s argument.

(1) Write $\epsilon:=\eord$ on $H$. It is Hida's projector: $\epsilon=\lim_{j}U^{j!}$ in the $\ell$-adic
topology of $\End_{\Z_\ell}(H)$, since modulo $\ell^{i}$ the sequence $U^{j!}$ is eventually
$\eord(U\bmod\ell^{i})$ (Proposition~\ref{prop:ordinary}(1)) and these are compatible
(Proposition~\ref{prop:ordinary}(4)), $H$ being finitely generated, hence $\ell$-adically complete; so
``the ordinary part of $H^{1}\otimes\Z_\ell$'' of Question~\ref{open:ordprecision} is the object $\epsilon H$
(v1.12.0-pre.2, R3). By its definition in Remark~\ref{rmk:ordtorsionfree} as the limit of
the $\eord(U\bmod\ell^{j})$, $\epsilon$ is a $\Z_\ell$-linear idempotent commuting with $U$ and every $T_p$
(Proposition~\ref{prop:ordinary}(2)), whose reduction modulo $\ell^{e'}$ is $\eord(U|_{H_{e'}})$. So
$H=\epsilon H\oplus(1-\epsilon)H$ and $H_{e'}=\epsilon H/\ell^{e'}\epsilon H\oplus(1-\epsilon)H/\ell^{e'}(1-\epsilon)H$,
with $\eord(U|_{H_{e'}})$ the projection onto the first summand. $\epsilon H$ is finitely generated and
torsion-free (Remark~\ref{rmk:ordtorsionfree}(2)), hence free of rank $h_{\mathrm{ord}}$, so the first
summand is free of rank $h_{\mathrm{ord}}$ over $\Z/\ell^{e'}$; each $T_p$ preserves $\epsilon H$ and the
characteristic polynomial of the reduction of an endomorphism of a free module is the reduction of its
characteristic polynomial. At $e'=1$, $H_1=H^{1}(\Ga;\Mn)$.

(2) By Proposition~\ref{prop:torsionprepass}(1), $H_{\mathrm{tors}}\cong\bigoplus_t\Z/\ell^{a_t}$ with every $a_t\le
e_{\max}<K$. $H_{\mathrm{tors}}$ is stable under every $T_p$ (a torsion element maps to one), so $H\to H/H_{\mathrm{tors}}$ and hence
$H_K=H/\ell^{K}H\to(H/H_{\mathrm{tors}})/\ell^{K}\to(H/H_{\mathrm{tors}})/\ell^{e'}$ are compatible with every $T_p$; they are onto. To
compute the kernel choose a $\Z_\ell$-module splitting $H=F_0\oplus H_{\mathrm{tors}}$ (not Hecke-stable; used only to
count): $H_K=F_0/\ell^{K}F_0\oplus H_{\mathrm{tors}}$, the map is the reduction $F_0/\ell^{K}\to F_0/\ell^{e'}$ on the first
summand and $0$ on $H_{\mathrm{tors}}$, so its kernel is $\ell^{e'}F_0/\ell^{K}F_0\oplus H_{\mathrm{tors}}$; and
$H_K[\ell^{e_{\max}}]=(F_0/\ell^{K})[\ell^{e_{\max}}]\oplus H_{\mathrm{tors}}[\ell^{e_{\max}}]=\ell^{K-e_{\max}}F_0/\ell^{K}F_0\oplus H_{\mathrm{tors}}$,
the same module since $K-e_{\max}=e'$. Multiplication by $\ell^{e_{\max}}$ is central, with kernel
$H_K[\ell^{e_{\max}}]$ and image $\ell^{e_{\max}}H_K$.

(3) $H_K\cong(\Z/\ell^{K})^{h}\oplus\bigoplus_t\Z/\ell^{a_t}$ with $a_t\le e_{\max}<K$ (as $e'\ge1$), and the
orders of the summands of a cyclic decomposition of a finite abelian $\ell$-group are unique; so $F'\cong(\Z/\ell^{K})^{h}$,
$\ell^{e_{\max}}T'=0$, and $H_K[\ell^{e_{\max}}]=F'[\ell^{e_{\max}}]\oplus T'=\ell^{e'}F'\oplus T'$. Hence
$F'\cap H_K[\ell^{e_{\max}}]=\ell^{e'}F'$ and $F'+H_K[\ell^{e_{\max}}]=H_K$, which is the isomorphism. For
$x\in F'$, $Ax=A_{F'}x+\pi_{T'}Ax$ with $\pi_{T'}Ax\in T'\subseteq H_K[\ell^{e_{\max}}]$, so $A_{F'}$ induces the
operator of (2). $(AB)_{F'}-A_{F'}B_{F'}=\pi_{F'}A\,\pi_{T'}B|_{F'}$, and $A(T')\subseteq A(H_K[\ell^{e_{\max}}])
\subseteq H_K[\ell^{e_{\max}}]$, whose image under $\pi_{F'}$ is $\ell^{e'}F'$.

(4) By Remark~\ref{rmk:ordtorsionfree}(2) the quotient map restricts to $\eord H\cong\eord(H/H_{\mathrm{tors}})$, compatibly
with every $T_p$; apply the argument of (1) to the free module $H/H_{\mathrm{tors}}$ and the projector $\eord(U|_{H/H_{\mathrm{tors}}})$
(the limit of the $\eord(U\bmod\ell^{j})$ on $(H/H_{\mathrm{tors}})/\ell^{j}$, Proposition~\ref{prop:ordinary}(4)).
\end{proof}

\begin{lemma}[The ordinary projector as one power of $U$]\label{lem:ordexponent}
\PROVED{} Let $\ell$ be a prime, $K\ge1$, $\mathfrak g\ge1$, $M$ a $\Z/\ell^{K}$-module generated by at
most $\mathfrak g$ elements, and $U\in\End(M)$. Let $c:=\lceil\log_\ell\mathfrak g\rceil$, the least
$c\ge0$ with $\ell^{c}\ge\mathfrak g$, and put
\[
  N_{\mathrm{ord}}(\mathfrak g,K):=\ell^{\,K-1+c}\prod_{f=1}^{\mathfrak g}(\ell^{f}-1),\qquad
  N^{\mathrm{lcm}}_{\mathrm{ord}}(\mathfrak g,K):=\ell^{\,K-1+c}\operatorname*{lcm}_{1\le f\le\mathfrak g}(\ell^{f}-1).
\]
\begin{enumerate}
\item $N^{\mathrm{lcm}}_{\mathrm{ord}}(\mathfrak g,K)$ divides $N_{\mathrm{ord}}(\mathfrak g,K)$, and
  $N^{\mathrm{lcm}}_{\mathrm{ord}}(\mathfrak g,K)\ge\mathfrak gK\ge\lambda(M)$.
\item Every element of $\mathrm{GL}_{\mathfrak g}(\Z/\ell^{K})$ has order dividing
  $N^{\mathrm{lcm}}_{\mathrm{ord}}(\mathfrak g,K)$.
\item $U^{N}=\eord(U)$ for every positive multiple $N$ of $N^{\mathrm{lcm}}_{\mathrm{ord}}(\mathfrak g,K)$ ---
  in particular for $N=N_{\mathrm{ord}}(\mathfrak g,K)$, and for $N=N_{\mathrm{ord}}(\mathfrak g',K')$
  whenever $\mathfrak g'\ge\mathfrak g$ and $K'\ge K$.
\end{enumerate}
No hypothesis on $\ell$.
\end{lemma}

\begin{proof}
(1) The divisibility is that of the product by the lcm. $\ell^{K-1}\ge K$ for $\ell\ge2$, $K\ge1$ (induction:
$\ell^{K}\ge2\ell^{K-1}\ge2K\ge K+1$), and $\ell^{c}\ge\mathfrak g$, so
$N^{\mathrm{lcm}}_{\mathrm{ord}}\ge\ell^{K-1+c}\ge K\mathfrak g$; and $M$ is a quotient of $(\Z/\ell^{K})^{\mathfrak g}$,
of length $\mathfrak gK$.

(2) Let $g\in\mathrm{GL}_{\mathfrak g}(\Z/\ell^{K})$, $\bar g$ its reduction modulo $\ell$, and
$\chi=\prod_i p_i^{c_i}\in\Fl[x]$ the characteristic polynomial of $\bar g$, the $p_i$ distinct monic
irreducibles of degrees $f_i\le\mathfrak g$, and $c_i\le\mathfrak g$. No $p_i$ is $x$, since
$\chi(0)=\pm\det\bar g\ne0$; so the class of $x$ in the field $\Fl[x]/(p_i)$, of order $\ell^{f_i}$, is a unit,
and $p_i\mid x^{\ell^{f_i}-1}-1\mid x^{m}-1$, $m:=\operatorname{lcm}_{f\le\mathfrak g}(\ell^{f}-1)$. The $p_i$
being distinct irreducibles, $\prod_ip_i\mid x^{m}-1$, and since $c_i\le\mathfrak g\le\ell^{c}$,
\[
  \chi\;\Big|\;\Bigl(\prod_ip_i\Bigr)^{\ell^{c}}\;\Big|\;(x^{m}-1)^{\ell^{c}}=x^{m\ell^{c}}-1
\]
(Frobenius in characteristic $\ell$). By Cayley--Hamilton $\bar g^{\,m\ell^{c}}=1$, i.e.\
$g^{m\ell^{c}}\in1+\ell\Mat_{\mathfrak g}(\Z/\ell^{K})$. For $j\ge1$ and $Y\in\Mat_{\mathfrak g}(\Z/\ell^{K})$,
$(1+\ell^{j}Y)^{\ell}=1+\sum_{i=1}^{\ell-1}\binom{\ell}{i}\ell^{ij}Y^{i}+\ell^{\ell j}Y^{\ell}\in1+\ell^{j+1}\Mat_{\mathfrak g}$,
since $\ell\mid\binom{\ell}{i}$ for $0<i<\ell$ and $\ell j\ge2j\ge j+1$ (this holds at $\ell=2$ too). By
induction $g^{m\ell^{c}\ell^{K-1}}\in1+\ell^{K}\Mat_{\mathfrak g}=\{1\}$.

(3) By Proposition~\ref{prop:ordinary}(1), $M=M_o\oplus M_n$ with $M_o=\im U^{\lambda}$, $M_n=\ker U^{\lambda}$,
$\lambda=\lambda(M)$, $U$ bijective on $M_o$, $U^{\lambda}=0$ on $M_n$, and $\eord(U)$ the projection onto
$M_o$ along $M_n$. On $M_n$, $U^{N}=0$ as $N\ge\lambda$ by (1). $M_o$ is a direct summand, hence a quotient,
of $M$; put $\mathfrak g_o:=\dim_{\Fl}M_o/\ell M_o\le\mathfrak g$ and choose $x_1,\dots,x_{\mathfrak g_o}\in M_o$
lifting a basis of $M_o/\ell M_o$; they generate $M_o$ (Nakayama over the local ring $\Z/\ell^{K}$). Let
$\pi\colon P:=(\Z/\ell^{K})^{\mathfrak g_o}\to M_o$ send the standard basis to the $x_i$; it is onto, and
an isomorphism modulo $\ell$. As $P$ is free there is $\tilde U\in\End(P)$ with $\pi\tilde U=U\pi$ (lift each
$Ux_i$). Modulo $\ell$, $\tilde U$ is the endomorphism $U$ induces on $M_o/\ell M_o$, which is onto since $U$
is onto $M_o$, hence bijective; so $\det\tilde U$ is a unit and $\tilde U\in\mathrm{GL}_{\mathfrak g_o}(\Z/\ell^{K})$.
Since $N^{\mathrm{lcm}}_{\mathrm{ord}}(\mathfrak g_o,K)\mid N^{\mathrm{lcm}}_{\mathrm{ord}}(\mathfrak g,K)\mid N$, (2)
gives $\tilde U^{N}=1$, so $U^{N}\pi=\pi\tilde U^{N}=\pi$ and $U^{N}=1$ on $M_o$ ($\pi$ is onto). Thus
$U^{N}=\eord(U)$. The last sentence: $N^{\mathrm{lcm}}_{\mathrm{ord}}(\mathfrak g,K)$ divides
$N^{\mathrm{lcm}}_{\mathrm{ord}}(\mathfrak g',K')$, hence $N_{\mathrm{ord}}(\mathfrak g',K')$.
\end{proof}

\begin{remark}[The terms of the exponent are needed, and where the lemma is applied]\label{rmk:ordexponentuse}
\PROVED{} (1) \emph{No term can be dropped.} At $K=1$ a unipotent Jordan block $1+E$ of size $\mathfrak g$
over $\Fl$ has $(1+E)^{\ell^{j}}=1+E^{\ell^{j}}$, which is $1$ iff $\ell^{j}\ge\mathfrak g$: its order is exactly
$\ell^{c}$. At $\mathfrak g=1$ and $\ell$ odd, $1+\ell\in(\Z/\ell^{K})^{\times}$ has order $\ell^{K-1}$ and a
generator of $(\Z/\ell^{K})^{\times}$ has order $\ell^{K-1}(\ell-1)=N_{\mathrm{ord}}(1,K)$. A companion matrix
of an irreducible polynomial of degree $f$ whose root generates $\F_{\ell^{f}}^{\times}$ has order
$\ell^{f}-1$. \COMP\textsuperscript{a} (v1.12.0-pre.2, R5; the slot-12 WP-3 referee's out-of-tree
\file{exponent\_check.py}, not a registered validation.) In every case tested --- $\mathrm{GL}_{\mathfrak g}(\Z/\ell^{K})$
exhaustively for $(\ell,\mathfrak g,K)\in\{(2,2,1),(2,2,2),(2,2,3),(2,3,1),(2,3,2),(3,2,1),(3,2,2),(3,3,1),(5,2,1),(7,2,1)\}$
and by sampling for $\mathrm{GL}_2(\Z/25)$ --- the exponent of the group \emph{equals}
$N^{\mathrm{lcm}}_{\mathrm{ord}}(\mathfrak g,K)$ (e.g.\ $168$ for $\mathrm{GL}_3(\Z/4)$, $600$ for
$\mathrm{GL}_2(\Z/25)$); equality in general is not claimed, the examples above showing only that each factor
is attained. (The corrected exponent of [X] v1.15 \S8 --- slot-11 WP-14 pass~2's ruling --- is
$N_{\mathrm{ord}}(r,K)$; an exponent without the $\lceil\log_\ell r\rceil$ term fails on the first example as
soon as $\mathfrak g>\ell^{0}=1$ and $K=1$.)
(2) \emph{Where it is applied} (with Corollary~\ref{cor:ordprecision}; $h^{1}(n):=\dim_{\Fl}H^{1}(\Ga;\Mn)$,
which is $2d_k-1+\lambda_1$ for even $n$ by Claim~\ref{cl:dimH1}(1), and $h=\rank H$). On $H_{e'}$ ---
generated by $h^{1}(n)$ elements, since $H_{e'}/\ell H_{e'}=H/\ell H=H^{1}(\Ga;\Mn)$ --- $\eord=U^{N}$ with
$N=N_{\mathrm{ord}}(h^{1}(n),e')$; on $H_K$ likewise with $N_{\mathrm{ord}}(h^{1}(n),K)$; on the free module of
Corollary~\ref{cor:ordprecision}(2), with $N_{\mathrm{ord}}(h,e')$. On an $F'$-block of
Corollary~\ref{cor:ordprecision}(3) over $\Z/\ell^{K}$, the pair $(F',U_{F'})$ is a $\Z/\ell^{K}$-module generated
by $h$ elements with an endomorphism, so $(U_{F'})^{N}=\eord(U_{F'})$ for every multiple $N$ of
$N^{\mathrm{lcm}}_{\mathrm{ord}}(h,K)$; but the pair depends on the chosen decomposition and $U_{F'}$ is not the
operator $U$ induces on any Hecke module over $\Z/\ell^{K}$, so $\eord(U_{F'})$ means nothing beyond its
reduction modulo $\ell^{e'}$, which is canonical: $U_{F'}\bmod\ell^{e'}$ \emph{is} the induced operator of
Corollary~\ref{cor:ordprecision}(2),(3), and $(U_{F'})^{N}\bmod\ell^{e'}$ is its $\eord$ for every multiple $N$ of
$N^{\mathrm{lcm}}_{\mathrm{ord}}(h,e')$ --- [X]'s $N(r,K)=N_{\mathrm{ord}}(h,K)$ among them. Slot-11 WP-18c's
idempotent is $\eord(U_{F'})$ of such a pair. (v1.12.0-pre.2, R6: through pre.1 this sentence said the block
power ``is \emph{not} an idempotent of anything over $\Z/\ell^{K}$'', which was too strong.)
(3) \emph{Size.} $\log_2N_{\mathrm{ord}}(\mathfrak g,K)<(K+c+\mathfrak g(\mathfrak g+1)/2)\log_2\ell$, so
$U^{N}$ costs $O\bigl((K+c+\mathfrak g^{2})\log\ell\bigr)$ products by repeated squaring; $N$ is never needed as
an integer, only its binary digits (or the factors, one power at a time).
\end{remark}

\begin{definition}[Weight towers at level~$1$]\label{def:hidatower}
Let $\ell\ge5$, $m\ge1$ and $n_0\ge2$ even. The \textbf{tower of level $m$ through $n_0$} is the set of
cells
\[
  \mathrm{Tw}_m(\ell;n_0):=\bigl\{(n,\ell):\ n\ge2\text{ even},\ n\equiv n_0\pmod{(\ell-1)\ell^{m-1}}\bigr\}.
\]
($\ell-1$ is even, so a class is parity-constant; odd $n$ are left out because
$H^{1}(\Ga;\Vcoef)\otimes\Z_\ell=0$ for odd $n$ at $\ell\ge5$, $2$ killing $H^{\bullet}(\Ga;\Vcoef)$
(Proposition~\ref{prop:badprimes}(2)).
$n=0$ is left out as outside Corollary~\ref{cor:ordprecision}.) $\mathrm{Tw}_{m+1}(\ell;n_0)\subseteq
\mathrm{Tw}_m(\ell;n_0)$, and the towers of level~$1$ are the classes of the weight $k=n+2$ modulo $\ell-1$.
To a cell attach, with $H(n):=H^{1}(\Ga;\Sym^{n}\Z^{2})\otimes\Z_\ell$ and $U=U_\ell$: its \textbf{ordinary
rank} $h^{1}_{\mathrm{ord}}(n):=\rank_{\Z_\ell}\eord H(n)$, and for each prime $p\ne\ell$ its
\textbf{ordinary characteristic polynomial} $\chi^{\mathrm{ord}}_{n,p}\in\Z_\ell[X]$, that of $T_p$ on the
free module $\eord H(n)$ (Remark~\ref{rmk:ordtorsionfree}). The \textbf{datum of the tower at level $m$} is
$\bigl(h^{1}_{\mathrm{ord}}(n),\ \chi^{\mathrm{ord}}_{n,p}\bmod\ell^{m}\bigr)_{n}$. By
Corollary~\ref{cor:ordprecision}(1) it is computed on $H^{1}(\Ga;\Vcoef\otimes\Z/\ell^{m})$ --- over the ring
$\Z/\ell^{m}$, with no extra precision --- and the rank already over $\Fl$; equivalently, by
Corollary~\ref{cor:ordprecision}(2)--(4), on $H_K$ at $K=m+e_{\max}(n,\ell)$.
\end{definition}

\begin{remark}[What is checkable on the towers, and at $(10,2411)$: the prediction list]\label{rmk:hidacheckable}
Committed at v1.12.0-pre.1, before the code that tests them (slot-12 WP-4 and WP-5; the plan's rule~20).
Each item carries its own tag. A failure of a \PROVED{} item is a defect of the implementation or of this
document and stops the package; a failure of a \CONJ{} item is a finding, recorded and not repaired.
\begin{enumerate}[label=(Hi\arabic*)]
\item\label{Hi1} \PROVED{} \textbf{Pre-pass versus run.} At every cell $(n,\ell)$, $\ell\ge5$, $n\ge1$,
  where the word-size $K$-mode of Proposition~\ref{prop:certifiede}(3) certifies, the pre-pass $e_{\max}$ of
  Proposition~\ref{prop:torsionprepass}(3) equals the largest torsion exponent of the certified
  $\Z_{(\ell)}$-type of $H^{1}(\Ceq(\Z))$ (slot-11 WP-18c's $e_{\max}$), and $\lambda_e=\sum_t\min(a_t,e)$ for
  its invariant factors $\ell^{a_t}$ at every $e$; the isomorphism of Proposition~\ref{prop:torsionprepass}(1)
  intertwines the $T_p$ on $H^{0}$ with those on $H^{1}[\ell^{e}]$. (Proposition~\ref{prop:torsionprepass}(1),
  (2),(4),(6).)
\item\label{Hi2} \PROVED{} \textbf{The power is the idempotent}, on named modules (v1.12.0-pre.2, R7).
  (i) On $F'/\ell^{e'}F'$ (a free block of Corollary~\ref{cor:ordprecision}(3)): $(U_{F'})^{N}\bmod\ell^{e'}$,
  $N$ any multiple of $N^{\mathrm{lcm}}_{\mathrm{ord}}(h,e')$, equals slot-11 WP-18c's certified idempotent
  $\eord(U_{F'})$ reduced modulo $\ell^{e'}$, and both are $\eord$ of the induced operator
  (Remark~\ref{rmk:ordexponentuse}(2)). (ii) On $H_{e'}$: $U^{N}=\eord(U|_{H_{e'}})$ for
  $N=N_{\mathrm{ord}}(h^{1}(n),e')$, idempotent and commuting with $T_2$, $T_3$. And the ordinary part read on $H_{e'}$
  (Corollary~\ref{cor:ordprecision}(1)) has the same rank and the same characteristic polynomials modulo
  $\ell^{e'}$ as the one read on the free block at $K=e'+e_{\max}$ (Corollary~\ref{cor:ordprecision}(4)).
  (Lemma~\ref{lem:ordexponent}; uniqueness of the Fitting projection.)
\item\label{Hi3} \CONJ{} \textbf{Ordinary ranks are constant on level-$1$ towers.}
  $h^{1}_{\mathrm{ord}}(n)=h^{1}_{\mathrm{ord}}(n')$ for $n,n'\ge2$ even with $n\equiv n'\pmod{\ell-1}$, and it
  is read over $\Fl$ (Corollary~\ref{cor:ordprecision}(1)). \emph{Source}: Hida's control theorem --- the
  rank of the ordinary part of weight-$k$ cusp forms on $\Ga_0(\ell)$, $\ell\ge5$, depends only on $k$ modulo
  $\ell-1$ for $k\ge2$ \cite{hida_ens,hida_invent}, \cite[Ch.~7]{hida_book} (locators \textbf{not verified},
  rule~\ref{A2}; none of the three is in \file{docs/literature/}). The constancy is invoked \emph{from $k=2$}:
  the class $k\equiv2\pmod{\ell-1}$ takes its value from $S_2(\Ga_0(\ell))$, all of which is $\ell$-new with
  $a_\ell=\pm1$ and therefore ordinary --- this is why the genus of $X_0(\ell)$ enters --- while the other
  classes take theirs at a weight $k\ge3$ through the identification that follows (v1.12.0-pre.2, R8).
  (v1.13.0-pre.1 note, no tag moved: the locators of \cite{hida_ens,hida_invent} are now verified, and the
  statements used are Theorems~\ref{thm:hidacontrol}, \ref{thm:hidaweight}; the level-$\ell$ identification that
  follows is Proposition~\ref{prop:stabilisation}. The sentence ``not verified'' above is superseded for those two
  papers, not for \cite{hida_book}.) For $k\ge3$ the $U_\ell$-ordinary part at
  level $\ell$ has the dimension of the level-$1$ one (a level-$1$ eigenform contributes one ordinary
  $\ell$-stabilisation iff $a_\ell$ is a unit; an $\ell$-new eigenform none, as $a_\ell^{2}=\ell^{k-2}$), and
  $\eord H(n)\otimes\Q_\ell$ is the Eisenstein line plus two copies of the ordinary cusp forms
  (Obligation~\ref{obl:oracle}'s Eichler--Shimura decomposition; the Eisenstein class has $U_\ell$-eigenvalue
  $1+\ell^{n+1}$, a unit). Hence the predicted \emph{values}, from $S_2(\Ga_0(\ell))$ (genus of $X_0(\ell)$:
  $0$ at $5,7,13$; $1$ at $11$, the curve $11a$ with $a_{11}=1$) and the small weights ($d_k$ of
  Claim~\ref{cl:dimH1}): $h^{1}_{\mathrm{ord}}(n)=1$ for every even $n\ge2$ at $\ell\in\{5,7\}$ (no ordinary cusp
  form of level~$1$ --- consistent with $\tau(5)\equiv0\bmod5$, $\tau(7)\equiv0\bmod7$); at $\ell=11$,
  $h^{1}_{\mathrm{ord}}(n)=3$ for $n\equiv0\pmod{10}$ and $1$ otherwise; at $\ell=13$, $3$ for $n\equiv10\pmod{12}$
  and $1$ otherwise ($\tau(11)\equiv1\bmod11$ --- the classical congruence $\Delta\equiv f_{11a}\pmod{11}$ ---
  and $\tau(13)\equiv8\bmod13$). \COMP\textsuperscript{a} evidence on the classical side, not a model
  computation (the slot-12 WP-3 referee's out-of-tree \file{hida\_check.py}: $T_\ell$ on $S_k(\SL_2(\Z))$
  from $q$-expansions, exactly): for every even $12\le k\le72$ and $\ell\in\{5,7,11,13\}$ the number of
  $\ell$-ordinary cusp eigenforms is exactly the table's, on all $120$ pairs.
\item\label{Hi4} \CONJ{} \textbf{Ordinary characteristic polynomials are congruent along towers.} For
  $m\in\{1,2\}$, $p\in\{2,3\}$ and two cells $n,n'$ of one $\mathrm{Tw}_m(\ell;n_0)$ (so of equal ordinary
  rank, \ref{Hi3}), $\chi^{\mathrm{ord}}_{n,p}\equiv\chi^{\mathrm{ord}}_{n',p}\pmod{\ell^{m}}$. \emph{Source}: the
  $\Lambda$-adic form of Hida's theorem (the ordinary Hecke module is free over $\Lambda=\Z_\ell[[X]]$ with
  specialisations the weight-$k$ ordinary parts; weights $k\equiv k'\pmod{(\ell-1)\ell^{m-1}}$ give points
  $u^{k}-1\equiv u^{k'}-1\pmod{\ell^{m}}$, $u=1+\ell$) \cite{hida_invent}, \cite[Ch.~7]{hida_book} (locators
  not verified; v1.13.0-pre.1 note, no tag moved: for \cite{hida_invent} and \cite{hida_ens} they are now,
  Theorem~\ref{thm:hidacontrol}). The Eisenstein factor $X-(1+p^{n+1})$ obeys it unconditionally:
  $p^{n+1}\equiv p^{n'+1}\pmod{\ell^{m}}$ as $(\ell-1)\ell^{m-1}$ is the order of $(\Z/\ell^{m})^{\times}$.
\item\label{Hi5} \PROVED{} \STD{\cite[Ch.~2--3]{harder3}, the Eisenstein class of weight $k$ and its
  eigenvalues $1+p^{k-1}$ for every prime $p$, $p=\ell$ included, as in Obligation~\ref{obl:oracle}}
  \textbf{The Eisenstein root, and where the congruence stops.} At every cell, $1+p^{n+1}$ is a root of
  $\chi^{\mathrm{ord}}_{n,p}$ in $\Z_\ell$ (the Eisenstein eigenline of $H(n)\otimes\Q_\ell$ is a \emph{simultaneous} eigenline of $U_\ell$ and of every
  $T_p$ (v1.12.0-pre.2, R9); it lies in $\eord H(n)\otimes\Q_\ell$: $U_\ell$ is
  invertible on $\eord H(n)$ and nilpotent modulo $\ell$ on $(1-\eord)H(n)$
  (Proposition~\ref{prop:ordinary}(1) on $H(n)/\ell H(n)$), so the eigenvalues on the second summand are
  non-units, while the Eisenstein eigenvalue $1+\ell^{n+1}$ is a unit; and $X-(1+p^{n+1})$ divides the monic
  $\chi^{\mathrm{ord}}_{n,p}$ over $\Q_\ell$, hence over $\Z_\ell$); so $\chi^{\mathrm{ord}}_{n,p}(1+p^{n+1})\equiv0
  \pmod{\ell^{e'}}$ at every precision. If $p$ is a primitive root modulo $\ell^{2}$ and $n,n'$ lie in one
  $\mathrm{Tw}_m(\ell;n_0)$ but $n\not\equiv n'\pmod{(\ell-1)\ell^{m}}$, then $1+p^{n+1}\not\equiv1+p^{n'+1}
  \pmod{\ell^{m+1}}$ (Remark~\ref{rmk:ordinaryeis}, $k-1=n+1$): the Eisenstein roots of the two cells agree
  modulo $\ell^{m}$ and not modulo $\ell^{m+1}$. \textbf{Not predicted}: that
  $\chi^{\mathrm{ord}}_{n,p}\not\equiv\chi^{\mathrm{ord}}_{n',p}\pmod{\ell^{m+1}}$ as polynomials --- a cuspidal root
  may be congruent to the other cell's Eisenstein root (Remark~\ref{rmk:ordinaryeis}'s last sentence).
  \emph{Addendum} (v1.21.0-pre.3; no statement or tag of this item changes): for $n<\ell-1$, [U]'s (U17)(a) of \texttt{pr:u-msym} --- the ordinary
  $\mathfrak m_{\mathrm{Eis}}$-block of $H^{+}$ is free of rank one over its Hecke algebra --- is equivalent to Harder's denominator \Espec{cl:lveis}(2),
  $[H^{+}:H^{+}_{!}\oplus\Z_\ell E]=\ell^{v(B_k/k)}$, on that block, as read in [E] after the proof of \Espec{cl:lveis} (the slot-21 referee F4 of [U], M3);
  for $n\ge\ell-1$ both stay conjectural.
\item\label{Hi6} \PROVED{} \STD{\cite[Ch.~8]{shimura}, the Eichler--Shimura isomorphism, for the
  characteristic polynomial of $T_{2411}$ on $H^{1}(\Ga;\Sym^{10}\Q^{2})$} \COMP\textsuperscript{a}
  \textbf{The non-ordinary example $(10,2411)$.} $e_{\max}(10,2411)=0$ ($n<\ell+1$,
  Proposition~\ref{prop:torsionprepass}(2)), $h=h^{1}(10)=3$, and $T_{2411}=U_{2411}$ has characteristic
  polynomial $(X-\tau(2411))^{2}\bigl(X-(1+2411^{11})\bigr)$ on $H(10)$. Since
  $\tau(2411)\equiv0\pmod{2411}$ --- indeed $\tau(2411)\equiv150\cdot2411\pmod{2411^{2}}$, an ad hoc
  computation from the $q$-expansion of $\Delta$ (\COMP\textsuperscript{a}; slot-12 WP-2's report, re-derived
  independently by the WP-3 referee: $\tau(2411)=4\,542\,041\,100\,095\,889\,012$), which the model is to confirm --- the monic factors $f_1=(X-\tau(2411))^{2}$ and $f_2=X-(1+2411^{11})$ are coprime modulo $2411$, so their
  resultant is a unit and $af_1+bf_2=1$ in $\Z_{2411}[X]$ (Bezout), whence
  $H(10)=\ker f_1(U)\oplus\ker f_2(U)$ as $\Z_{2411}$-modules (v1.12.0-pre.2, R10), $U$-stable of ranks $2$ and $1$, and $\Delta$'s rank-$2$ summand is non-ordinary
  ($U^{2}=2\tau(2411)U-\tau(2411)^{2}\in2411\End$ on it by Cayley--Hamilton, so $U^{2m}\equiv0$ modulo $2411^{m}$), and over $\Z/2411^{m}$ for every $m\ge1$ the ordinary
  part is the Eisenstein line: $h^{1}_{\mathrm{ord}}(10)=1$, the Eisenstein count, and
  $\chi^{\mathrm{ord}}_{10,p}=X-(1+p^{11})$.
\end{enumerate}
\end{remark}

\paragraph{Scope: what this chapter does not do.}
It proves no part of Hida theory: \ref{Hi3} and \ref{Hi4} are predictions from the literature, tagged
\CONJ, and the model can only confirm them on a finite grid. Slot~13's (the slot-12 plan's ruling L1):
control and weight independence as statements of this document; the $\Ga_0(\ell)$-stabilisation and the
level-$\ell$ model (at level~$1$, $U_\ell$ is the operator of \ref{H2}, $\Tdc{g}$ for $g=\mathrm{diag}(1,\ell)$);
the $\Lambda$-adic lattice and Hida's Hecke algebra; towers beyond $\ell^{2}$; the cells $(700,691)$ (beyond
the $40$~GB rule) and $(2420,2411)$. At $\ell\in\{2,3\}$ nothing of Corollary~\ref{cor:ordprecision},
Definition~\ref{def:hidatower} or Remark~\ref{rmk:hidacheckable} is asserted
(Proposition~\ref{prop:torsionprepass}(6)).

\paragraph{Code map.}
\textbf{Built} (slot-12 WP-4 and WP-5; recorded at v1.12.1 --- until then this paragraph read ``Not built at
v1.12.0-pre.1 (contract first)'', the plan below being its text).
\begin{itemize}\setlength{\itemsep}{0pt}
\item \file{Hecke/HeckeOrdinaryPart} (\code{hecke::ordinary}, WP-4): Proposition~\ref{prop:torsionprepass}'s
  pre-pass --- $J_n=[\rhon(S)-1;\rhon(T)-1]$ over $\Z/\ell^{e}$, the lengths $\lambda_e$ of its kernel, stopped
  at the first $e$ with $\lambda_{e+1}=\lambda_e$, giving $e_{\max}$ and the torsion multiset before any run ---
  and Lemma~\ref{lem:ordexponent} (the ordinary projector as one power of $U$, the exponent
  $N_{\mathrm{ord}}$); the precision use at $\ell\in\{2,3\}$ is refused by name
  (Proposition~\ref{prop:torsionprepass}(6)). \file{Hecke/LadicCell} takes the pre-pass as an opt-in
  precision source (Corollary~\ref{cor:ordprecision}). \code{HeckeOrdinaryPrecisionValidation}
  (\code{standard}): \ref{Hi1} (pre-pass $=$ the certified word-size $K$-mode on $20$ cells) and \ref{Hi2}
  (the projector against the digit route's certified idempotent, and on the torsion-carrying modules).
\item \file{HeckeExt/HeckeHidaTower} (\code{hecke\_ext::hida}, WP-5): Definition~\ref{def:hidatower} ---
  per cell $K=m_{\max}+e_{\max}$ from the pre-pass, $U_\ell$, $T_2$, $T_3$ on $H_K$ by the shipped module
  route, the ordinary part read on $H_{e'}$ --- and \ref{Hi2}--\ref{Hi6} across the weight towers
  $\mathrm{Tw}_1$, $\mathrm{Tw}_2$, with the non-ordinary cell $(10,2411)$. \code{HeckeHidaTowerValidation}
  (\code{standard}, $163$ cells) and \code{HeckeHidaTowerSlowValidation} (\code{slow}, $280$ cells).
\end{itemize}
Nothing of the chapter's ``not asserted'' list (the $\Lambda$-adic lattice, towers beyond $\ell^{2}$, the
cells $(700,691)$, $(2420,2411)$) is built.

\paragraph{Formalization note.}
Proposition~\ref{prop:torsionprepass} needs the long exact sequence of a coefficient sequence, finite
generation of $\Ga$, and the Smith form over a discrete valuation ring; Lemma~\ref{lem:ordexponent} needs
Fitting's lemma, Cayley--Hamilton and the binomial theorem. Both are cheap. Corollary~\ref{cor:ordprecision}
adds Remark~\ref{rmk:ordtorsionfree}. Hida's theorems (\ref{Hi3}, \ref{Hi4}) are not formalization targets
of this document.

%----------------------------------------------------------------------
\section{Hida theory at level one: control, level \texorpdfstring{$\ell$}{l}, the lattice}\label{sec:hidafull}
%----------------------------------------------------------------------
% NEW at v1.13.0-pre.1 (hecke-extensions slot 13, WP-2; plan docs/implementation/slot13/00-slot13-plan.md
% S2 "WP-2" and S9 "Owner decisions" (the Q2 amendment: Hida's papers read-only at
% the local library (papers/h/Hida/), locators verified against them; the level-Gamma_0(N)
% amendment: the Shapiro route on the fixed model, monomial / orbit structure, fibre maps, a cost model).
% The contract first (rule 21); the predictions (Hi7)--(Hi12) before the code of WP-5 / WP-6 / WP-10
% (rule 20).  Placed after sec:hida, before the appendices, so that no chapter number moves.

\paragraph{What this chapter fixes.}
Everything slot~13 implements on the Hida line, and the predictions it tests. (i) \textbf{Level structure by
Shapiro} (\S\ref{sec:levelshapiro}): the cohomology of $\Ga_0(N)$ is computed on the \emph{fixed} model
$(K,\Rel)$ with the induced local system of Definition~\ref{def:inducedls}, never on a triangulation of
$\Ga_0(N)\backslash\Hyp$ (the owner's direction of 2026-09-26); every level-$N$ Hecke operator, degeneracy
map and level-change map is one block formula on the induced fibre that reuses the \emph{level-$1$}
$\tau$-table (Proposition~\ref{prop:levelcorrespondence}); at $N=\ell$ the operator
$U^{(\ell)}_\ell$ is written out (Proposition~\ref{prop:levelellU}) together with the relations of the old
space (Proposition~\ref{prop:oldspace}); the fibre's monomial structure and the orbit decomposition over
$\mathbb{P}^1(\Z/N)$ give $H^{1}$ as a cokernel of fixed-point spaces (Proposition~\ref{prop:levelorbit}),
and \S\ref{sec:levelcost} is the cost model, with the two regimes it serves and the obstacles to scaling
(Remark~\ref{rmk:levelscaling}). (ii) \textbf{Stabilisation and control}
(\S\ref{sec:stabilisation}): the ordinary part at level $\ell$ is the ordinary part at level $1$, integrally
(Lemma~\ref{lem:levelelltorsion}, Proposition~\ref{prop:stabilisation}); Hida's control and
weight-independence theorems are stated as literature with verified locators
(Theorems~\ref{thm:hidacontrol}, \ref{thm:hidaweight}). (iii) \textbf{The finite-level lattice datum}
(Definition~\ref{def:hidalattice}) and (iv) \textbf{the predictions} (Hi7)--(Hi12)
(Remark~\ref{rmk:hidacheckable2}).

\paragraph{Notation and standing conventions.}
$\Ga=\SL_2(\Z)$, $\Vcoef=\Sym^{n}\Z^{2}$ with $\rhon$ of Definition~\ref{def:rhon}, $\De$ the monoid of
Definition~\ref{def:semigroup}, on which $\rhon$ is multiplicative (Proposition~\ref{prop:symconv}(2)).
$g:=\smat{1}{0}{0}{\ell}$ throughout \S\S\ref{sec:levelshapiro}--\ref{sec:stabilisation} except in
Proposition~\ref{prop:levelcorrespondence}, where $g\in\De$ is arbitrary. For $N\ge1$,
$m_N:=[\Ga:\Ga_0(N)]=\abs{\mathbb{P}^1(\Z/N)}$ (the index $m$ of \S\ref{sec:induced} for $H=\Ga_0(N)$; the
subscript keeps it apart from the tower level $m$ of Definition~\ref{def:hidatower}). For finite-index
$A,B\le\Ga$ and $g\in\De$ write
\[
  [A\,g\,B]\;:=\;\cor^{B}_{B\cap gAg^{-1}}\circ(\conj_g)_{\ast}\circ\res^{A}_{A\cap g^{-1}Bg}
  \colon H^{\bullet}(A;\VR)\to H^{\bullet}(B;\VR),
\]
the form of \eqref{eq:corconjres} between two groups; $[\Ga g\Ga]=\Tdc{g}$ (Theorem~\ref{thm:classical}).
The \textbf{level-$\ell$ operator} is $U^{(\ell)}_\ell:=[\Ga_0(\ell)\,g\,\Ga_0(\ell)]$ (the plan's
$U_\ell^{(\ell)}$); the level-$1$ operator $U_\ell=T_\ell=\Tdc{g}$ keeps its name. The two
\textbf{degeneracy maps} are $\mathrm{dg}_1:=[\Ga\,1\,\Ga_0(\ell)]=\res^{\Ga}_{\Ga_0(\ell)}$ and
$\mathrm{dg}_\ell:=[\Ga\,g\,\Ga_0(\ell)]$ (the plan's $\alpha_1$, $\alpha_\ell$; the letter $\alpha$ is
Proposition~\ref{prop:alpha}'s restriction unit into $\Ind_{\Ga^{0}(N)}$, rule~\ref{A5}). $s_1,\dots,s_{\ell+1}$
are the canonical transporters of Definition~\ref{def:cosetreps} at $N=\ell$ --- $s_{1+b}=\smat{1}{0}{b}{1}$
($0\le b<\ell$), $s_{\ell+1}=S$ --- and $s^{(2)}_b:=\smat{1}{0}{\ell b}{1}$ ($0\le b<\ell$). $\tau$ is the
level-$1$ $g$-twisted chain datum of Definition~\ref{def:tau} for $g=\smat{1}{0}{0}{\ell}$ --- the table the
pipeline already builds for $U_\ell$ --- and $\Ga''_g:=\Ga\cap g\Ga g^{-1}=\Ga_0(\ell)$
(Proposition~\ref{prop:cosets}(3)). Group cohomology is written $H^{\bullet}(A;-)$; the model enters through
Proposition~\ref{prop:lhcompliant}(4), which needs $6\in\cring^{\times}$.

\subsection{Level structure by Shapiro on the fixed model}\label{sec:levelshapiro}

\begin{definition}[The level-$N$ induced model; built at $N=\ell$]\label{def:levelellmodel}
Let $N\ge1$ and let $w_1=I,w_2,\dots,w_{m_N}\in\Ga$ be a fixed table of representatives of
$\Ga_0(N)\backslash\Ga$, one per class of $\mathbb{P}^1(\Z/N)$ under the bottom-row map of
Proposition~\ref{prop:cosets}(4). (Its proof is written for $N=p^{\nu}$; for every $N$, well-definedness and
injectivity hold verbatim, and surjectivity with the count $\abs{\mathbb{P}^1(\Z/N)}=N\prod_{p\mid N}(1+p^{-1})$
follow from the Chinese remainder theorem and the surjectivity of $\SL_2(\Z)\to\SL_2(\Z/N)$; v1.13.0-pre.2,
WP-3 referee m2.)
At $N=p^{\nu}$ the table is Definition~\ref{def:cosetreps}'s; for composite $N$ it is part of the interface and
fixed once, like Definition~\ref{def:cosetreps}, since the cochain-level matrices depend on it
(Remark~\ref{rmk:indeplevel}) (fixed for every $N$ by Definition~\ref{def:cosetrepsN}, v1.14.0-pre.1). The \textbf{level-$N$ model} is the induced matrix local system
$L_{\Ga_0(N)}$ of Definition~\ref{def:inducedls} on $(K,\Rel)$: fibre $\Ind_{\Ga_0(N)}^{\Ga}\VR$ of rank
$m_N(n+1)$, relation matrices $M_{\Ga_0(N)}(\gamma)$ of \eqref{eq:MH}. Each of these is \textbf{monomial}:
the permutation $j\mapsto j\cdot\gamma$ of $\mathbb{P}^1(\Z/N)$ (the row-vector action $(c:d)\mapsto(c:d)\gamma$)
with the $m_N$ blocks $\rhon(h_j(\gamma))$, $h_j(\gamma)=w_j\gamma w_{j\cdot\gamma}^{-1}\in\Ga_0(N)$ ---
Proposition~\ref{prop:blockmatrices}. Its cohomology $\Heq{\bullet}(K,\Rel;L_{\Ga_0(N)})$ is \emph{the level-$N$
cohomology on the model}; for $6\in\cring^{\times}$ it is $H^{\bullet}(\Ga_0(N);\VR)$
(Proposition~\ref{prop:lhcompliant}(4)) \LEAN{def:levelellmodel-model} (this sentence, for any table of $\Ga_0(N)\backslash\Ga$; not the
fixed table nor the monomial matrices as C++ objects), and for $6\notin\cring^{\times}$ nothing is asserted. Slot~13 builds
$N=\ell$ (fibre rank $(\ell+1)(n+1)$) on the shipped \file{Hecke/InducedLocalSystem}; general $N$ is future
work, and every statement below is written for the pair of groups it needs so that nothing precludes it.
\end{definition}

\begin{proposition}[Level-$N$ correspondences on the induced fibre]\label{prop:levelcorrespondence}
\PROVED{} Assume \ref{H4}. Let $A,B\le\Ga$ be of finite index, $g\in\De$, and fix representatives $w_a$ of
$A\backslash\Ga$ and $v_i$ of $B\backslash\Ga$ ($w_1=v_1=I$), $r_t$ of $B''\backslash B$ where
$B'':=B\cap gAg^{-1}$, and $s_j$ of $\Ga''_g\backslash\Ga$ with a $g$-twisted chain datum $\tau$ over $\cring$
(Definition~\ref{def:tau}: the level-$1$ datum of $\Tdc{g}$). For each pair $(i,t)$ write
\[
  r_tv_i=h\,s_j\quad(h\in\Ga''_g),\qquad k:=g^{-1}hg,\qquad a:=\text{the index with }Ak=Aw_a,\qquad
  \delta_{i,t}:=v_i\,s_j^{-1}\,g\,w_a^{-1},
\]
so $j=j(i,t)$, $a=a(i,t)$. Define, for $\phi\in\Hom_{\Z[\Ga]}\bigl(C_{\bullet}(\Ku),\Ind_A^{\Ga}\VR\bigr)$ and
$x\in C_q(\Ku)$, block by block in the coordinates of Definition~\ref{def:inducedmodule},
\begin{equation}\label{eq:levelcorr}
  \bigl(\widetilde{[A\,g\,B]}\,\phi\bigr)(x)_i\;:=\;\sum_{t}\ \rhon(\delta_{i,t})\,
  \phi\bigl(\tau(s_{j(i,t)}\,x)\bigr)_{a(i,t)} .
\end{equation}
Then:
\begin{enumerate}
\item \textbf{(Integral, no new $\tau$-table.)} $B''\subseteq\Ga''_g$; $k\in\Ga$; and
  $\delta_{i,t}=r_t^{-1}\,g\,(kw_a^{-1})$ with $kw_a^{-1}\in A$, so $\delta_{i,t}\in\De$ is integral of
  determinant $\det g$. Formula~\eqref{eq:levelcorr} reads only the level-$1$ $\tau$-table of $\Tdc{g}$ (on the
  $s_j\lift{\sigma}$), the class maps of $A\backslash\Ga$, $B\backslash\Ga$, $\Ga''_g\backslash\Ga$, and
  $\rhon$ of integral matrices; no determinant-$\det g$ matrix is inverted (Obligation~\ref{obl:integral}).
\item \textbf{(A cochain map, every $\cring$.)} $\widetilde{[A\,g\,B]}$ maps $\Ga$-equivariant cochains to
  $\Ga$-equivariant cochains and commutes with the coboundary; under $\Psi$
  (Theorem~\ref{thm:explicit_complex_equiv}) it is a cochain map $\Ceq(K,\Rel;L_A)\to\Ceq(K,\Rel;L_B)$, with
  \emph{no} hypothesis on $\ell$. \LEAN{prop:levelcorrespondence-model} (in the reciprocity form of (3), on the model)
\item \textbf{(What it computes.)} Under the cochain-level Frobenius reciprocity
  $\Hom_{\Ga}(C_{\bullet},\Ind_A^{\Ga}\VR)\cong\Hom_{A}(C_{\bullet},\VR)$, $\phi\mapsto\phi(\cdot)_1$, it is
  the map $\psi\mapsto\bigl(y\mapsto\sum_t\rhon(r_t^{-1}g)\,\psi(\tau(r_ty))\bigr)$ \LEAN{prop:levelcorrespondence-model} (this formula, and
  its factorisation $\cor\circ(\conj_g)_{\ast}\circ\res$ at cochain level and on $H^q$, on the model; not the rest of (3): that it \emph{is}
  $[A\,g\,B]$ on $H^{\bullet}(A;\VR)$, independent of $\tau$ and of every representative); if $6\in\cring^{\times}$ the
  map it induces on $\Heq{\bullet}(K,\Rel;L_A)\cong H^{\bullet}(A;\VR)$ is $[A\,g\,B]$, independent of $\tau$
  and of every choice of representatives \STD{\cite[Ch.~III.9--10]{brown}, as in
  Theorem~\ref{thm:classical}}.
\item \textbf{(The instances.)} (a) $A=B=\Ga$: \eqref{eq:levelcorr} is \eqref{eq:heckechain}. (b)
  \textbf{Level change as fibre maps} ($g=I$, so $\Ga''_I=\Ga$, $s_1=I$, $\tau=\id$): for $A\subseteq B$,
  $[A\,1\,B]=\cor^{B}_{A}$ is the \emph{trace} $\omega_{A/B}\colon f\mapsto\bigl(\sum_{a:\,Aw_a\subseteq Bv_i}
  \rhon(v_iw_a^{-1})f_a\bigr)_i$; for $B\subseteq A$, $[A\,1\,B]=\res^{A}_{B}$ is the \emph{projection}
  (pull-back along $B\backslash\Ga\to A\backslash\Ga$) $\alpha^{0}_{B/A}\colon f\mapsto
  \bigl(\rhon(v_iw_{a(i)}^{-1})f_{a(i)}\bigr)_i$, $Av_i=Aw_{a(i)}$. \LEAN{prop:levelcorrespondence-model} ($\cor$ and $\res$ as the
  factors of (3), at cochain level) Both are morphisms of matrix local
  systems --- one monomial fibre matrix applied at every simplex --- so the level change between
  $\Ga_0(N')\subseteq\Ga_0(N)$ ($N\mid N'$) is a fibre map in both directions; $\omega_{\Ga_0(N)/\Ga}$ is the
  $\omega$ of Proposition~\ref{prop:omega}. (c) \textbf{Degeneracy maps as Hecke-type correspondences}:
  $A=\Ga_0(N)$, $B=\Ga_0(N')$, $g=\smat{1}{0}{0}{d}$, $dN\mid N'$: then $B''=B$ (one $t$) and
  $[A\,g\,B]=(\conj_g)_{\ast}\circ\res$ \LEAN{prop:levelcorrespondence-model} (as the factor $(\conj_g)_{\ast}\circ\res$ of (3), at cochain
  level) --- \eqref{eq:levelcorr} with one term per block, reading, for $d=p^{\nu}$
  a prime power, the level-$1$ $\tau$-table of the primitive double coset $\Ga\smat{1}{0}{0}{d}\Ga$
  (Definition~\ref{def:cosetreps}). For composite $d=d_1d_2$ the degeneracy map is the composite of prime-power
  ones through $C:=\Ga_0(Nd_1)$: $[A\,g\,B]=[C\,g_2\,B]\circ[A\,g_1\,C]$, $g_i=\smat{1}{0}{0}{d_i}$, on cohomology
  for $6\in\cring^{\times}$ (each factor has one term, and $(\conj_{g_2})_\ast\res\circ(\conj_{g_1})_\ast\res=
  (\conj_{g_1g_2})_\ast\res$ by functoriality); at cochain level the composite of the two block formulas is the
  formula used, so no table for composite $d$ is needed (v1.13.0-pre.2, WP-3 referee m2). (d) \textbf{Hecke operators at
  level $N$}: $A=B=\Ga_0(N)$, $g=\smat{1}{0}{0}{p}$ --- $T^{(N)}_p$ for $p\nmid N$, $U^{(N)}_p$ for $p\mid N$
  --- reading the level-$1$ table of $T_p$ (of $U_\ell$ at $p=\ell$).
\end{enumerate}
\end{proposition}

\begin{proof}
(1) $B''=B\cap gAg^{-1}\subseteq\Ga\cap g\Ga g^{-1}=\Ga''_g$, as $A,B\subseteq\Ga$. $r_tv_i\in\Ga$ lies in exactly
one coset $\Ga''_gs_j$, which defines $j$ and $h$; $g^{-1}\Ga''_gg\subseteq\Ga$ by definition of $\Ga''_g$, so
$k\in\Ga$ and $a$ is defined. From $r_tv_i=hs_j$: $v_is_j^{-1}=r_t^{-1}h$, so
$\delta_{i,t}=r_t^{-1}hgw_a^{-1}=r_t^{-1}g\,(kw_a^{-1})$, and $kw_a^{-1}\in A$ because $Ak=Aw_a$. The rest is read
off \eqref{eq:levelcorr}.

(3) The reciprocity $\phi\mapsto\phi_1:=\phi(\cdot)_1$ is a bijection with inverse
$\psi\mapsto\bigl(x\mapsto(\psi(w_ax))_a\bigr)$, by $\phi(w_ax)_1=(w_a\cdot\phi(x))(w_1)=\phi(x)(w_a)$
(Definition~\ref{def:inducedmodule}); it commutes with precomposition by any chain map, hence with the
coboundary. Put $\psi=\phi_1$ and $(T\psi)(y):=\sum_t\rhon(r_t^{-1}g)\psi(\tau(r_ty))$. Then
$(T\psi)(v_ix)=\sum_t\rhon(r_t^{-1}g)\,\phi(\tau(hs_jx))_1$, and by (T1) for $h\in\Ga''_g$,
$\tau(hs_jx)=k\,\tau(s_jx)$; by \eqref{eq:MH}, $\phi(kz)_1=(M_A(k)\phi(z))_1=\rhon(kw_a^{-1})\phi(z)_a$ (block row
$1$ of $M_A(k)$ is the single block $\rhon(h_1(k))$, $h_1(k)=w_1kw_a^{-1}$, in column $a=1\cdot k$). With
Proposition~\ref{prop:symconv}(2), $\rhon(r_t^{-1}g)\rhon(kw_a^{-1})=\rhon(\delta_{i,t})$, which is
\eqref{eq:levelcorr}: so $\widetilde{[A\,g\,B]}\phi$ is the transport of $T\psi$. \emph{$T\psi$ is
$B$-equivariant}: for $b\in B$ write $r_tb=b_tr_{t\cdot b}$ with $b_t\in B''$; then
$\tau(r_tby)=(g^{-1}b_tg)\tau(r_{t\cdot b}y)$ with $g^{-1}b_tg\in A$, so
$\rhon(r_t^{-1}g)\psi(\tau(r_tby))=\rhon(r_t^{-1}b_tg)\psi(\tau(r_{t\cdot b}y))=\rhon(b)\rhon(r_{t\cdot b}^{-1}g)
\psi(\tau(r_{t\cdot b}y))$, and $t\mapsto t\cdot b$ is a bijection --- the computation of
Theorem~\ref{thm:welldef}. $T$ commutes with the coboundary because $\tau$ and each $r_t$ are chain maps. This
proves (2) by transport. If $6\in\cring^{\times}$, $C_{\bullet}(\Ku)\otimes\cring$ is a resolution of $\cring$ by
modules that are projective over $\cring[A]$ and over $\cring[B]$ (the stabiliser orders divide $12$; the proof
of Lemma~\ref{lem:tauexists}(2), restriction to a finite-index subgroup keeping summands of free modules), and
the three factors of $[A\,g\,B]$ are, on it, the inclusion, $\psi\mapsto\rhon(g)\psi\circ\tau$ and the transfer
$\psi\mapsto\sum_t\rhon(r_t^{-1})\psi(r_t\,\cdot)$ --- the proof of Theorem~\ref{thm:classical} with
$(\Ga,\Ga',\Ga'')$ replaced by $(A,A\cap g^{-1}Bg,B'')$; $\tau$ is a chain map over
$\conj_g^{-1}\colon B''\to A\cap g^{-1}Bg$ lifting the identity of $\cring$, and any two such are homotopic by the
argument of Lemma~\ref{lem:tauexists}(2). Composing gives $T$.

(4) (a) $v_1=w_1=I$, $B''=\Ga''_g$, $r_t=s_t$, so $j=t$, $h=I$, $k=I$, $a=1$ and $\delta=s_t^{-1}g$. (b) $g=I$:
$\Ga''_I=\Ga$, one $j$ with $s_1=I$, $h=r_tv_i$, $k=h$; for $A\subseteq B$, $B''=A$ and the $r_tv_i$ run through
representatives of $A\backslash\Ga$ inside $Bv_i$, giving $\rhon(v_iw_a^{-1})$ as $a$ runs over the $A$-cosets in
$Bv_i$; for $B\subseteq A$, $B''=B$, one $t$ ($r_1=I$), $a=a(i)$. A formula with $\tau=\id$ and one fixed matrix
per block pair is a morphism of local systems (Theorem~\ref{thm:rep_ls_equiv}). (c) $gAg^{-1}\cap\Ga=\Ga_0(dN)$
by the computation of Proposition~\ref{prop:cosets}(3), which contains $B=\Ga_0(N')$; so $B''=B$. (d) is the
definition.
\end{proof}

\begin{proposition}[The level-$\ell$ operator as a block formula]\label{prop:levelellU}
\PROVED{} Assume \ref{H4}; $\cring$ any commutative ring; $\ell$ any prime; $H:=\Ga_0(\ell)$ with
$w_i=v_i=s_i$. Put
\[
  \varpi\;:=\;\bigl(\rhon(gs_1^{-1}),\ \rhon(gs_2^{-1}),\ \dots,\ \rhon(gs_\ell^{-1}),\ 0\bigr)
  \colon\Ind_H^{\Ga}\VR\to\VR,\qquad gs_{1+b}^{-1}=\smat{1}{0}{-\ell b}{\ell},
\]
the block row whose last block (the class $(1:0)$, $s_{\ell+1}=S$) is zero. Then the instances of
Proposition~\ref{prop:levelcorrespondence} at $N=\ell$ are:
\begin{equation}\label{eq:Ulevell}
  \bigl(\widetilde U^{(\ell)}_\ell\phi\bigr)(x)_i=\varpi\,\phi\bigl(\tau(s_ix)\bigr),\qquad
  (\widetilde{\mathrm{dg}}_\ell\phi)(x)_i=\rhon(g)\,\phi\bigl(\tau(s_ix)\bigr),\qquad
  (\widetilde{\mathrm{dg}}_1\phi)(x)_i=\rhon(s_i)\,\phi(x),
\end{equation}
$i=1,\dots,\ell+1$, and the trace $\omega\colon f\mapsto\sum_i\rhon(s_i^{-1})f_i$ of
Proposition~\ref{prop:omega}. On the model, with $\tau(s_i\lift{\sigma})=\sum c\,\gamma\,\lift{\sigma'}$ the stored
chain,
\begin{equation}\label{eq:Ulevellmodel}
  \bigl(\widetilde U^{(\ell)}_\ell\phi\bigr)(\lift\sigma)_i\;=\;\sum_{c\,\gamma\,\lift{\sigma'}}c
  \sum_{j\ne(\ell+1)\cdot\gamma}\rhon\bigl(g\,\gamma\,s_j^{-1}\bigr)\,\phi(\lift{\sigma'})_j ,
\end{equation}
every $g\gamma s_j^{-1}$ ($j\ne(\ell+1)\cdot\gamma$) integral of determinant $\ell$. In particular
$U^{(\ell)}_\ell$, $\mathrm{dg}_\ell$ and $\mathrm{dg}_1$ need \textbf{no coset table and no $\tau$-table beyond
level~$1$'s $U_\ell$}: the $\ell$ cosets $s^{(2)}_b$ of $\Ga_0(\ell^{2})\backslash\Ga_0(\ell)$ enter only through the
closed form $s^{(2)-1}_bg=gs_{1+b}^{-1}$. By Proposition~\ref{prop:levelcorrespondence}(2),(3) each is a cochain map
over every $\cring$, and at $6\in\cring^{\times}$ induces $U^{(\ell)}_\ell$, $\mathrm{dg}_\ell$, $\mathrm{dg}_1$ on
$H^{\bullet}(\Ga_0(\ell);\VR)$.
\end{proposition}

\begin{proof}
$U^{(\ell)}_\ell$: $B''=H\cap gHg^{-1}=\Ga_0(\ell^{2})$ (the computation of Proposition~\ref{prop:cosets}(3)), with
representatives $r_b=s^{(2)}_b$ of $\Ga_0(\ell^{2})\backslash\Ga_0(\ell)$: their bottom rows $(\ell b,1)$ are the
$\ell$ classes $(c:d)\in\mathbb{P}^1(\Z/\ell^{2})$ with $\ell\mid c$, which is $\Ga_0(\ell^{2})\backslash\Ga_0(\ell)$
by Proposition~\ref{prop:cosets}(4) at $N=\ell^{2}$. Since $s^{(2)}_b\in\Ga_0(\ell)=\Ga''_g$, $s^{(2)}_bs_i\in\Ga''_gs_i$:
$j=i$, $h=s^{(2)}_b$, $k=g^{-1}s^{(2)}_bg=\smat{1}{0}{b}{1}=s_{1+b}$, $a=1+b$, and
$\delta_{i,b}=s_is_i^{-1}gs_{1+b}^{-1}=gs_{1+b}^{-1}$, whose value is displayed; it equals $s^{(2)-1}_bg$ by
multiplying out. $\mathrm{dg}_\ell$: $A=\Ga$ (so $\Ind_\Ga^\Ga\VR=\VR$, $a=1$), $B=H$,
$B''=H\cap g\Ga g^{-1}=H$, one $t$, $j=i$, $\delta=g$. $\mathrm{dg}_1$ is
Proposition~\ref{prop:levelcorrespondence}(4)(b) with $A=\Ga\supseteq B=H$. For \eqref{eq:Ulevellmodel}: $\phi$ is
$\Ga$-equivariant, so $\phi(\gamma\lift{\sigma'})=M_H(\gamma)\phi(\lift{\sigma'})$, and $\varpi M_H(\gamma)f=
\sum_b\rhon(gs_{1+b}^{-1})\rhon(h_{1+b}(\gamma))f_{(1+b)\cdot\gamma}$ with
$s_{1+b}^{-1}h_{1+b}(\gamma)=\gamma s_{(1+b)\cdot\gamma}^{-1}$ (Lemma~\ref{lem:cocycle}); as $b$ runs over
$0,\dots,\ell-1$ the index $j=(1+b)\cdot\gamma$ runs over all indices but $(\ell+1)\cdot\gamma$, $\gamma$ acting
bijectively. Integrality: $g\gamma s_j^{-1}=gs_{1+b}^{-1}\,h_{1+b}(\gamma)$, a product of integral matrices.
\end{proof}

\begin{proposition}[The old space at level $\ell$]\label{prop:oldspace}
\PROVED{} \COMP\textsuperscript{a} Assume \ref{H4}; $\cring$ any commutative ring; notation of
Proposition~\ref{prop:levelellU}, with $T_\ell=\Tdc{g}$ of level $1$.
\begin{enumerate}
\item $\widetilde U^{(\ell)}_\ell\circ\widetilde{\mathrm{dg}}_1=\ell\cdot\widetilde{\mathrm{dg}}_\ell$ \emph{as cochain
  maps}.
\item $\omega\circ\widetilde{\mathrm{dg}}_1=(\ell+1)\cdot\id$ and $\omega\circ\widetilde{\mathrm{dg}}_\ell=
  \widetilde T_\ell$ (the cochain map of \eqref{eq:heckechain}) \emph{as cochain maps}.
\item If $6\in\cring^{\times}$: $U^{(\ell)}_\ell\circ\mathrm{dg}_\ell=\mathrm{dg}_\ell\circ T_\ell-\ell^{\,n}\,
  \mathrm{dg}_1$ on cohomology.
\item Hence, for $6\in\cring^{\times}$, the map $\Phi\colon H^{\bullet}(\Ga;\VR)^{\oplus2}\to
  H^{\bullet}(\Ga_0(\ell);\VR)$, $(x,y)\mapsto\mathrm{dg}_1x+\mathrm{dg}_\ell y$, satisfies
  $U^{(\ell)}_\ell\Phi=\Phi\,\mathsf M$ with
  \[
    \mathsf M:=\begin{pmatrix}0&-\ell^{\,n}\\ \ell&T_\ell\end{pmatrix},\qquad
    \det(X-\mathsf M)=X^{2}-T_\ell X+\ell^{\,n+1}=X^{2}-T_\ell X+\ell^{\,k-1}.
  \]
\end{enumerate}
\COMP\textsuperscript{a} evidence (slot-13 WP-2's out-of-tree \file{levelell\_check.py}, not a registered
validation; group cohomology by $1$-cocycles on $S,T$, the operators in the form $[A\,g\,B]$ with the
representatives above): (1)--(3) hold on cohomology classes over $\F_{1000003}$ at
$(n,\ell)\in\{(2,5),(4,5),(10,5),(2,7),(6,7),(10,7),(0,5),(0,11)\}$. The script has no $\tau$ and works on
cohomology classes only: the cochain-level forms of (1), (2), and the block formula \eqref{eq:Ulevell} itself, have
no numerical evidence until slot-13 WP-6 evaluates them on the model (v1.13.0-pre.2, WP-3 referee m6, m10).
\end{proposition}

\begin{proof}
(1) By \eqref{eq:Ulevell}, block $i$ of $\widetilde U^{(\ell)}_\ell\widetilde{\mathrm{dg}}_1\phi$ at $x$ is
$\varpi\bigl((\rhon(s_{j})\phi(\tau(s_ix)))_j\bigr)=\sum_{b<\ell}\rhon(gs_{1+b}^{-1})\rhon(s_{1+b})\phi(\tau(s_ix))
=\ell\,\rhon(g)\phi(\tau(s_ix))$. (2) $\sum_i\rhon(s_i^{-1})\rhon(s_i)f=(\ell+1)f$, and
$\sum_i\rhon(s_i^{-1})\rhon(g)\phi(\tau(s_ix))=\sum_i\rhon(s_i^{-1}g)\phi(\tau(s_ix))$ is \eqref{eq:heckechain}.
(3) Block $i$ of $\widetilde{\mathrm{dg}}_\ell\widetilde T_\ell\phi$ is
$\sum_{j=1}^{\ell+1}\rhon(gs_j^{-1}g)\phi(\tau(s_j\tau(s_ix)))$, and block $i$ of
$\widetilde U^{(\ell)}_\ell\widetilde{\mathrm{dg}}_\ell\phi$ is $\sum_{b<\ell}\rhon(gs_{1+b}^{-1})\rhon(g)
\phi(\tau(s_{1+b}\tau(s_ix)))$, the terms $j\le\ell$ of the former. The difference is the term $j=\ell+1$,
$s_{\ell+1}=S$: since $gS^{-1}g=\ell S^{-1}$ and $\rhon(\ell I)=\ell^{\,n}$ (Proposition~\ref{prop:diamond}), it is
$\psi(x)_i=\ell^{\,n}\rhon(S^{-1})\phi(\hat\tau(s_ix))$ with $\hat\tau:=\tau\circ S\circ\tau$. Through the
reciprocity of Proposition~\ref{prop:levelcorrespondence}(3) $\psi$ is the $H$-cochain map
$\phi\mapsto\ell^{\,n}\rhon(S^{-1})\phi\circ\hat\tau$. Now $\hat\tau$ is an augmentation-preserving chain map with
$\hat\tau(hy)=(ShS^{-1})\hat\tau(y)$ for $h\in H$: $\tau(hy)=(g^{-1}hg)\tau(y)$, $S(g^{-1}hg)S^{-1}\in\Ga_0(\ell)$ as
$g^{-1}hg\in\Ga^{0}(\ell)$, so $\tau(S(g^{-1}hg)\tau y)=g^{-1}S g^{-1}hgS^{-1}g\,\hat\tau(y)$, and
$g^{-1}Sg^{-1}=\ell^{-1}S$ gives $ShS^{-1}$. The map $y\mapsto Sy$ has the same three properties. For
$6\in\cring^{\times}$ the two are therefore homotopic through such maps (the argument of
Lemma~\ref{lem:tauexists}(2) for the subgroup $H$ and the homomorphism $\conj_S\colon H\to\Ga$), so $\psi$ is
cochain-homotopic to $\phi\mapsto\ell^{\,n}\rhon(S^{-1})\phi(S\,\cdot)=\ell^{\,n}\phi$, i.e.\ to
$\ell^{\,n}\widetilde{\mathrm{dg}}_1$ in the reciprocity. (4) $U^{(\ell)}_\ell(\mathrm{dg}_1x+\mathrm{dg}_\ell y)=
\ell\,\mathrm{dg}_\ell x+\mathrm{dg}_\ell T_\ell y-\ell^{\,n}\mathrm{dg}_1y$ by (1), (3); $\det(X-\mathsf M)=X(X-T_\ell)+
\ell^{\,n+1}$, and $k-1=n+1$.
\end{proof}

\begin{proposition}[Orbit-wise fixed points; $H^{1}$ as a cokernel]\label{prop:levelorbit}
\PROVED{} Let $N\ge1$, $\cring$ a commutative ring, $M:=\Ind_{\Ga_0(N)}^{\Ga}\VR$ with the monomial action of
Definition~\ref{def:levelellmodel}, and $F\le\Ga$ a finite subgroup.
\begin{enumerate}
\item \textbf{(Mackey, orbit-wise.)} $M^{F}=\bigoplus_{O}M_O^{F}$ over the $F$-orbits $O$ on
  $\mathbb{P}^1(\Z/N)$, $M_O$ the sum of the blocks $j\in O$. For $j_O\in O$, with $F_{j_O}$ its stabiliser acting
  on $\VR$ by $\gamma\mapsto\rhon(h_{j_O}(\gamma))$ ($h_{j_O}(\gamma)\in\Ga_0(N)$), restriction to block $j_O$ is an
  isomorphism $M_O^{F}\to\VR^{F_{j_O}}$, inverse $f_{j_O\cdot\gamma}:=\rhon(h_{j_O}(\gamma))^{-1}f_{j_O}$. A free
  orbit modulo $F\cap\{\pm I\}$ contributes a copy of $\VR^{F\cap\{\pm I\}}$ and \textbf{no solve}.
\item \textbf{(The cokernel.)} If $6\in\cring^{\times}$ and $n$ is even, then
  \[
    H^{1}(\Ga_0(N);\VR)\;\cong\;H^{1}(\Ga;M)\;=\;M\big/\bigl(M^{\langle S\rangle}+M^{\langle ST\rangle}\bigr),\qquad
    M^{\langle S\rangle}=\tfrac12(1+M(S))M,\quad M^{\langle ST\rangle}=\tfrac13\bigl(1+M(ST)+M(ST)^{2}\bigr)M,
  \]
  and $H^{0}=M^{\langle S\rangle}\cap M^{\langle ST\rangle}$; both invariant spaces are computed orbit-wise by (1).
  If $n$ is odd and $2\in\cring^{\times}$, $H^{\bullet}(\Ga_0(N);\VR)=0$.
\item \textbf{(Dimensions.)} If $\cring$ is a field in which $6$ is invertible and $n$ is even,
  \[
    \dim H^{1}-\dim H^{0}\;=\;m_N(n+1)-\Bigl(\tfrac{m_N-\nu_2}{2}(n+1)+\nu_2\,d_S\Bigr)
    -\Bigl(\tfrac{m_N-\nu_3}{3}(n+1)+\nu_3\,d_{ST}\Bigr),
  \]
  where $\nu_2=\nu_2(N)$, $\nu_3=\nu_3(N)$ are the numbers of fixed points of $S$, $ST$ on $\mathbb{P}^1(\Z/N)$ and
  $d_S:=\dim\VR^{\langle S\rangle}$, $d_{ST}:=\dim\VR^{\langle ST\rangle}$ (independent of $N$). At $N=\ell\ge5$:
  $m_\ell=\ell+1$, $\nu_2(\ell)=1+\bigl(\tfrac{-1}{\ell}\bigr)$, $\nu_3(\ell)=1+\bigl(\tfrac{-3}{\ell}\bigr)$.
\end{enumerate}
\COMP\textsuperscript{a} evidence (the same script): (3) equals the cocycle dimension of $H^{1}(\Ga;M)$ over
$\F_{1000003}$ at the eight cells of Proposition~\ref{prop:oldspace}.
\end{proposition}

\begin{proof}
(1) $(\gamma\cdot f)_j=\rhon(h_j(\gamma))f_{j\cdot\gamma}$ (\eqref{eq:MH}), so $F$ preserves each $M_O$, and
$f\in M_O$ is fixed iff $f_j=\rhon(h_j(\gamma))f_{j\cdot\gamma}$ for all $j\in O$, $\gamma\in F$. Taking $j=j_O$:
$f$ is determined by $f_{j_O}$, which must be fixed by $\rhon(h_{j_O}(\gamma))$ for $\gamma\in F_{j_O}$; conversely,
for such $f_{j_O}$ the displayed inverse is well defined (if $j_O\cdot\gamma=j_O\cdot\gamma'$ then
$\gamma'\gamma^{-1}\in F_{j_O}$ and $h_{j_O}(\gamma')=h_{j_O}(\gamma'\gamma^{-1})h_{j_O}(\gamma)$ by
Lemma~\ref{lem:cocycle}'s cocycle identity) and fixed, by the same identity. On a free orbit
$F_{j_O}=F\cap\{\pm I\}$ ($-I$ fixes every point of $\mathbb{P}^1$). (2) Claim~\ref{cl:amalgam}(2),(3) are proved from
\eqref{eq:treeses} alone, which is a sequence of $\Z[\Ga]$-modules; applying $\Hom_{\Z[\Ga]}(-,M)$ instead of
$\Hom_{\Z[\Ga]}(-,\VR)$ gives the same statements for $M$: $H^{q}(\langle S\rangle;M)=H^{q}(\langle ST\rangle;M)=0$
for $q\ge1$ as the orders $4,6$ are units, and $M^{-I}=M$ for $n$ even ($M(-I)=(-1)^{n}$,
Proposition~\ref{prop:lhcompliant}(3)). The averaging idempotents are the displayed ones because $M(S)^{2}=M(-I)=1$
and $M(ST)^{3}=M(-I)=1$. The first isomorphism is Shapiro's lemma \STD{\cite[Ch.~III.6.2]{brown}}. For $n$ odd,
$-I\in\Ga_0(N)$ is central and acts by $-1$, so $2$ kills $H^{\bullet}$ (the argument of
Proposition~\ref{prop:badprimes}(2)). (3) From (2), over a field, $\dim H^{1}-\dim H^{0}=\dim M-\dim M^{\langle
S\rangle}-\dim M^{\langle ST\rangle}$ (the exact sequence $0\to H^{0}\to M^{S}\oplus M^{ST}\to M\to H^{1}\to0$). $S$
and $ST$ act on $\mathbb{P}^1(\Z/N)$ through $\Ga/\{\pm I\}$ with orders $2$ and $3$, so every orbit is free or a
fixed point; a fixed point $j$ has $F_j=F$ acting through $h_j(\gamma)=w_j\gamma w_j^{-1}$, conjugate to $F$, whence
$\dim\VR^{F_j}=d_F$ by (1). At $N=\ell$: $(c:d)\cdot S=(d:-c)$ is fixed iff $d=\lambda c$, $-c=\lambda d$ for a unit
$\lambda$, i.e.\ $\lambda^{2}=-1$ ($c=0$ gives $(0:1)\mapsto(1:0)$), two points or none; $(c:d)\cdot ST=(d:d-c)$ is fixed
iff $\lambda^{2}-\lambda+1=0$, of discriminant $-3$, two roots or none ($\ell\ne3$).
\end{proof}

\subsection{The cost model}\label{sec:levelcost}

\begin{remark}[Dense versus monomial versus orbit-wise; in $m_N$, $k$ and $\ell^{l}$]\label{rmk:levelcost}
A model of operation counts, not an assertion about running times; $d:=n+1=k-1$, $D_N:=m_N\,d$ the fibre rank,
$m_N=N\prod_{p\mid N}(1+p^{-1})$ (Proposition~\ref{prop:cosets}(4) for prime powers, the Chinese remainder theorem
for the rest). All counts are in operations of $\cring$.
\begin{enumerate}
\item \textbf{Dense} (the shipped \code{InducedLocalSystem} + \code{EquivariantCohomologyOnK}, which is what slot~13
  builds at $N=\ell$, $D_\ell=(\ell+1)d\le574$ on WP-6's grid): each relation matrix costs $D_N^{2}$ to store and
  $O(D_N^{3})$ to take a fixed space; the equivariant cochain spaces have dimension $\Theta(N_qD_N)$ and their
  eliminations cost $O((N_qD_N)^{3})$ --- \textbf{cubic in $m_N(k-1)$}. A Hecke \emph{matrix} on $H^{1}$ on this
  route costs, beyond the elimination, the dense projector to cohomology, of size $\Theta((N_qD_N)^{2})$, applied to
  each of the $\dim H^{1}$ image columns of the cochain map of (5) --- so the first level-$N$ Hecke matrix this
  programme computes is cubic in $m_N(k-1)$, whatever (5) says about the operator (v1.13.0-pre.2, WP-3 referee M2).
\item \textbf{Monomial}: $M_{\Ga_0(N)}(\gamma)$ is a permutation of $\mathbb{P}^1(\Z/N)$ and $m_N$ blocks: $m_Nd^{2}$
  to store, $O(m_Nd^{2})$ to apply (against $D_N^{2}=m_N^{2}d^{2}$).
\item \textbf{Orbit-wise} (Proposition~\ref{prop:levelorbit}(1)): the fixed spaces of the stabilisers cost nothing on
  free orbits and one $d\times d$ solve per fixed point, $O((\nu_2+\nu_3)d^{3})$ in all --- and cheaper still, since
  $\VR^{F_j}=\rhon(w_j)\VR^{\langle S\rangle}$ resp.\ $\rhon(w_j)\VR^{\langle ST\rangle}$ at a fixed point $j$: one
  $d\times d$ solve for each of $S$, $ST$ once, then $O(d^{2})$ per fixed point. The counts are exact for every
  $N$: a fixed point $(c:d)$ of $S$ or $ST$ has $c$ a unit (at a prime $p\mid N$ with $p\mid c$, $d$ is a unit modulo
  $p$ and the fixed-point equation of the proof of Proposition~\ref{prop:levelorbit}(3) forces $p\nmid c$), so
  $\nu_2(N)=\#\{\lambda\in\Z/N:\lambda^{2}+1\equiv0\}$ and $\nu_3(N)=\#\{\lambda\in\Z/N:\lambda^{2}-\lambda+1\equiv0\}$;
  by the Chinese remainder theorem both are $\le2^{\omega(N)}$ for every $N$, $\nu_2(N)=0$ if $4\mid N$ or a prime
  $\equiv3\pmod4$ divides $N$, and $\nu_3(N)=0$ if $9\mid N$ or a prime $\equiv2\pmod3$ divides $N$ (v1.13.0-pre.2,
  WP-3 referee m3; the counts checked numerically for $N\le30$ and six composites by the referee's script; checked against
  Definition~\ref{def:cosetrepsN}'s table for every $N\le60$ and eleven $N\le2310$, slot-14 WP-12).
\item \textbf{The cokernel} (Proposition~\ref{prop:levelorbit}(2),(3)): after the orbit-local quotient by
  $M^{\langle S\rangle}$ (on a free $S$-orbit $\{j,j\cdot S\}$, $f\mapsto f_j-\rhon(h_j(S))f_{j\cdot S}$ up to the
  factor $\tfrac12$), what remains is one sparse elimination against $M^{\langle ST\rangle}$: about $D_N/2$ rows,
  $D_N/3$ columns, at most three $d\times d$ blocks per block column, with answer of dimension about $D_N/6$
  ($\dim H^{1}\approx2\dim S_k(\Ga_0(N))$). This elimination is the only super-linear step; its worst case is cubic
  in $D_N$, and this document asserts no better bound (the block sparsity is what modular-symbol implementations
  exploit).
\item \textbf{Hecke operators and level change} (Proposition~\ref{prop:levelcorrespondence}): per simplex and per
  pair $(i,t)$, one level-$1$ $\tau$-chain and one $d\times d$ block per term --- $O([\Ga:B'']\,\bar c\,d^{2})$,
  $\bar c$ the mean number of terms of a level-$1$ chain; for $U^{(\ell)}_\ell$, $\ell+1$ output blocks each costing
  $\ell$ blocks per term by \eqref{eq:Ulevellmodel}. No $D_N\times D_N$ matrix is formed, and the level-$1$
  $\tau$-table --- the dominant memory of $U_\ell$ (slot-13 WP-4's lever) --- is shared, not replicated per level.
  The fibre maps of Proposition~\ref{prop:levelcorrespondence}(4)(b) cost $O(m_{N'}d^{2})$. So \emph{applying} a
  level-$N$ Hecke operator is linear in $m_N$ and quadratic in $d$, and needs only the level-$1$ $\tau$-tables of the
  primes in use (composite degeneracies compose, Proposition~\ref{prop:levelcorrespondence}(4)(c)).
  \textbf{The count is for the cochain map.} Reading it on $H^{1}$ costs, on the dense route (1), the
  $\Theta((N_qD_N)^{2})$ projector per image column --- cubic in $m_N(k-1)$ in all --- and on the cokernel
  presentation (4) it needs a Hecke action \emph{on the presentation}, which the formula of
  Proposition~\ref{prop:levelcorrespondence} is not: that action is the Manin--Merel Heilbronn-matrix formula,
  specified by pointer and not built (Remark~\ref{rmk:levelscaling}, \ref{levelO2}).
\item \textbf{The ring $\Z/\ell^{l}$.} The monomial and orbit structures are ring-independent (a permutation and
  integral $2\times2$ matrices); the orbit-wise averaging needs only $6\in\cring^{\times}$, so $\ell\ge5$. Each
  count above is multiplied by the cost of one ring operation: a word operation for $\ell^{l}<2^{64}$, and beyond
  it the digit route of \file{ENGINE\_API.md} \S13.11, whose per-operation cost grows with the number of digit
  planes. None of (2)--(5) is built at this version.
\end{enumerate}
The regimes this model serves, and the obstacles to scaling, are Remark~\ref{rmk:levelscaling}.
\end{remark}

\begin{remark}[The two regimes, and the obstacles to scaling in $N$, $k$ and $\ell^{l}$]\label{rmk:levelscaling}
\STD{\cite{manin} (Manin symbols and their relations); \cite[\S2]{merel} (the Hecke operators on Manin symbols by
sets of integral matrices of determinant $p$, the Heilbronn--Merel matrices); \cite[Ch.~3 and Ch.~8]{stein} and
\cite[Ch.~2]{cremona} (the algorithmic form); pointers at section and chapter level, \textbf{not} verified against the
texts (rule~\ref{A2}): no copy is available to this document} Added at v1.13.0-pre.2 (WP-3 referee M2 and its
scalability section~\S7), so that the owner's requirement --- the level-$N$ route must scale in large $k$, large $N$
and large $\ell^{l}$ --- is tracked in the contract and not only in a report. Notation of
Remark~\ref{rmk:levelcost}.

\emph{The two regimes.} (R-N) \textbf{Large $N$ at small $k$} --- the modular-symbol regime: the cokernel
presentation of Proposition~\ref{prop:levelorbit}(2) is, for $6\in\cring^{\times}$ and $n$ even, the presentation of
the modular symbols of weight $k$ for $\Ga_0(N)$ by Manin symbols $\mathbb{P}^1(\Z/N)\times\VR$ modulo the relations
$x+xS=0$, $x+xU+xU^{2}=0$ ($U=ST$; the tree of \eqref{eq:treeses} is the Farey tree), and the cost is the block-sparse
elimination of Remark~\ref{rmk:levelcost}(4) plus a Hecke action on the presentation (O2). (R-$k$) \textbf{Large $k$ at
level $1$ or $\ell$} --- the Hida regime of \S\ref{sec:hida} and this chapter: $m_1=1$, $m_\ell=\ell+1$, the dense route
of Remark~\ref{rmk:levelcost}(1) on the fixed model, dominated by the level-$1$ $U_\ell$ $\tau$-table (O5). \textbf{Large
$N$ and large $k$ together are outside every route of Remark~\ref{rmk:levelcost}}: every route has dense $d\times d$
blocks (O4).

\emph{The obstacles}, in order of weight:
\begin{enumerate}[label=(O\arabic*)]
\item\label{levelO1} \textbf{The only built route is dense}: $O((N_qm_Nd)^{3})$ per cohomology computation and a
  $\Theta((N_qm_Nd)^{2})$ projector per Hecke column (Remark~\ref{rmk:levelcost}(1), (5)). Adequate on slot-13 WP-6's
  grid ($D_\ell\le574$, $N_qD_\ell\approx10^{4}$), not beyond.
\item\label{levelO2} \textbf{The cokernel presentation has no Hecke action in this document.} The action is classical:
  $T_p$ on Manin symbols is a fixed finite sum over a set of integral matrices of determinant $p$ (Merel's
  Heilbronn-type matrices), acting on $\mathbb{P}^1(\Z/N)\times\VR$ without any $\tau$-search --- a closed-form Hecke
  action on the tree of \eqref{eq:treeses}. It is \textbf{specified here by pointer only}: transcribing it into this
  document's conventions (the column convention of Definition~\ref{def:rhon}, the right action on
  $\mathbb{P}^1(\Z/N)$, the form $[A\,g\,B]$ of \eqref{eq:corconjres}) and proving it equal to
  Proposition~\ref{prop:levelcorrespondence}(4)(d) on cohomology is a contract task of the size of
  Definition~\ref{def:cosetreps}, not a search like Algorithm~\ref{alg:tau}. Not scheduled; not built. \textbf{Status at
  v1.14.0-pre.1}: transcribed and proved in \S\ref{sec:levelN} (Proposition~\ref{prop:heilbronn}, for $T_p$, $p\nmid N$, and $U_p$,
  $p\mid N$, with its cost in Remark~\ref{rmk:heilbronncost}); not built. \textbf{At v1.14.2}: still not built (slot~15); the $\tau$-route of Proposition~\ref{prop:levelcorrespondence}(4)(d) at
  composite $N$ is built (WP-12).
\item\label{levelO3} \textbf{Composite-$N$ coset tables are not pinned}: the representatives of $\Ga_0(N)\backslash\Ga$ and
  $B''\backslash B$ with their class maps are a binding datum (Remark~\ref{rmk:indeplevel}) fixed by
  Definition~\ref{def:cosetreps} for prime powers only (Definition~\ref{def:levelellmodel}). A table, not a search;
  no new $\tau$-table is needed for any $N$ (Proposition~\ref{prop:levelcorrespondence}(1), (4)(c)). \textbf{Status at
  v1.14.0-pre.1}: pinned for every $N$ by Definition~\ref{def:cosetrepsN}; not built. \textbf{At v1.14.2}: built (slot-14 WP-12, \code{Gamma0LevelN}); \ref{levelO3} is closed.
\item\label{levelO4} \textbf{Large $k$ and large $N$ together}: the cokernel elimination is at best $(m_N/6)$-block-sparse
  with $d^{3}$ per block operation; the Hida regime ($n\approx(\ell-1)\ell^{m-1}$) at a level $N\ge\ell\ge691$ is out of
  reach of every route above.
\item\label{levelO5} \textbf{The level-$1$ $U_\ell$ $\tau$-table at large $\ell$} is shared by every level, not replicated
  (Remark~\ref{rmk:levelcost}(5)) --- but it must exist, and it is the bottleneck already at level $1$. Slot-13 WP-4's
  opt-in rank-$\le k$ lever (\file{docs/implementation/slot13/WP-4\_report.md}) makes it about five times smaller: its
  fitted model puts $(700,691)$ over $\F_{691}$ at $\approx14.7$~GiB (upper band $16.9$~GiB) and $(2420,2411)$ at
  $\approx540$~GiB (band $647$~GiB), the latter beyond a single box; the remaining floor is the coset tables' full-rank
  $\rhon$ of unimodular matrices (a streaming lever, unscheduled). A model of memory, not an assertion; over
  $\Z/\ell^{2}$ at these $n$ no prediction exists.
\item\label{levelO6} \textbf{$\ell\in\{2,3\}$ and $\Z$}: the orbit-wise averaging and the cokernel need $6\in\cring^{\times}$, so
  only the dense route survives there (it does survive: the cochain-level statements of the chapter hold over every
  ring, Scope). \textbf{$\Z/\ell^{l}$, $\ell\ge5$}: the cokernel route survives, by Howell elimination over
  $\Z/\ell^{l}$ and the digit route beyond $2^{64}$ (Remark~\ref{rmk:levelcost}(6)).
\end{enumerate}
\emph{Verdict recorded} (the referee's, adopted): Shapiro on the fixed model is the right route --- the induced model
and a triangulation of $\Ga_0(N)\backslash\Hyp$ give the same cochain complex, the triangulation saves no dimension and
would need a new certified $\tau$-table per $(N,g)$, and the scalable presentation for large $N$ (Manin symbols on the
Farey tree) is itself the Shapiro route on the tree; \ref{levelO1} and \ref{levelO2} decide whether large $N$ is ever reached.
\end{remark}

\subsection{Stabilisation and control}\label{sec:stabilisation}

\begin{lemma}[Level $\ell$: the torsion is non-ordinary]\label{lem:levelelltorsion}
\PROVED{} Let $\ell$ be a prime, $n\ge1$, $H:=\Ga_0(\ell)$, $U:=U^{(\ell)}_\ell$.
\begin{enumerate}
\item For every $e\ge1$, $U$ acts on $H^{0}(H;\Vcoef\otimes\Z/\ell^{e})=(\Vcoef/\ell^{e})^{H}$ by
  $v\mapsto\sum_{b<\ell}\rhon(gs_{1+b}^{-1})v$, over every ring and on the model as well; $U=0$ at $e=1$ and
  $U^{e}=0$ in general.
\item $H^{0}(H;\Vcoef\otimes\Q)=0$, so the connecting map is a $U$-equivariant isomorphism
  $H^{0}(H;\Vcoef/\ell^{e})\cong H^{1}(H;\Vcoef)[\ell^{e}]$; $U$ is nilpotent on the $\ell$-primary torsion
  $T_\ell$ of $H(\ell):=H^{1}(H;\Vcoef)\otimes\Z_\ell$.
\item Consequently $\eord(T_\ell)=0$, $\eord H(\ell)$ is free of finite rank, and $\eord H(\ell)\cong\eord(H(\ell)/T_\ell)$
  (the argument of Remark~\ref{rmk:ordtorsionfree}, with (2) in place of Corollary~\ref{cor:thetanonordinary}(3)).
\item If $\ell\ge5$: $H^{1}(H;\Vcoef\otimes\Z/\ell^{e})=H(\ell)/\ell^{e}H(\ell)$, compatibly with $U$ and every
  operator of Proposition~\ref{prop:levelcorrespondence}(4)(d); so Corollary~\ref{cor:ordprecision}(1) holds at level
  $\ell$ verbatim: the ordinary part is read over $\Z/\ell^{e'}$ at the target precision itself, its rank over $\Fl$.
\end{enumerate}
\end{lemma}

\begin{proof}
(1) A $0$-cocycle of $\Hom_H(C_{\bullet}(\Ku),M)$ is constant, $\Ku$ being connected, with value $v\in M^{H}$; by
(T3) $\psi(\tau(y))=v$ for every vertex $y$, so the formula of Proposition~\ref{prop:levelcorrespondence}(3) gives
$\sum_b\rhon(s^{(2)-1}_bg)v=\sum_b\rhon(gs_{1+b}^{-1})v$ with no hypothesis on $\cring$. Every
$gs_{1+b}^{-1}=\smat{1}{0}{-\ell b}{\ell}$ is $\equiv\smat{1}{0}{0}{0}\pmod\ell$, so modulo $\ell$ the $\ell$ terms are
equal and their sum is $\ell\,\rhon(\smat{1}{0}{0}{0})v\equiv0$: $U=0$ at $e=1$, and in general $Uv\in\ell(\Vcoef/\ell^{e})$.
Filter $F_i:=(\Vcoef/\ell^{e})^{H}\cap\ell^{i}(\Vcoef/\ell^{e})$. For $v=\ell^{i}w\in F_i$, $w$ is $H$-invariant modulo
$\ell^{e-i}$ (multiplication by $\ell^{i}$ is injective $\Vcoef/\ell^{e-i}\to\Vcoef/\ell^{e}$), so
$Uv=\ell^{i}Uw\in\ell^{i+1}(\Vcoef/\ell^{e})$, and $Uv$ is $H$-invariant: $UF_i\subseteq F_{i+1}$, $F_e=0$.
(2) $T=\smat1101$ and $\smat10\ell1$ lie in $H$; over $\Q$ the fixed space of $\rhon(T)$ is $\Q u_0$, $u_0=x^{n}$
(Definition~\ref{def:rhon}: $T\cdot x=x$, $T\cdot y=x+y$), and $\smat10\ell1$ sends $x^{n}$ to $(x+\ell y)^{n}\ne x^{n}$
for $n\ge1$. The Bockstein is the argument of Proposition~\ref{prop:torsionprepass}(1) for the group $H$ (which is
finitely generated, of finite index in $\Ga$), and its $U$-equivariance that of
Proposition~\ref{prop:torsionprepass}(4). With (1), $U^{e}$ kills $H^{1}(H;\Vcoef)[\ell^{e}]$. (3) is
Remark~\ref{rmk:ordtorsionfree}'s proof verbatim. (4) The next term of the sequence is
$H^{2}(H;\Vcoef)[\ell^{e}]=H^{2}(\Ga;\Ind_H^{\Ga}\Vcoef)[\ell^{e}]$ (Shapiro, \STD{\cite[Ch.~III.6.2]{brown}}), which
$24$ kills by the argument of Claim~\ref{cl:amalgam}(4) (it uses \eqref{eq:treeses} alone); so it vanishes for
$\ell\ge5$. Then the proof of Corollary~\ref{cor:ordprecision}(1) applies with (3) in place of
Remark~\ref{rmk:ordtorsionfree}.
\end{proof}

\begin{proposition}[Stabilisation from level $1$ to level $\ell$]\label{prop:stabilisation}
Tags per item (v1.13.0-pre.2, WP-3 referee m1): (1), (2) \PROVED{}; (3) \PROVED{} except the $T_p$-intertwining,
\textsf{[standard]}; (4) the rank \textsf{[standard]}, the isomorphism then \PROVED{} from (3); (5) \textsf{[standard]}. Let $\ell\ge5$, $n\ge1$; $H(1):=H^{1}(\Ga;\Vcoef)\otimes\Z_\ell$ with $e_1:=\eord(T_\ell)$ and
$H(\ell):=H^{1}(\Ga_0(\ell);\Vcoef)\otimes\Z_\ell$ with $e_\ell:=\eord(U^{(\ell)}_\ell)$, both the $\ell$-adic projectors
of Remark~\ref{rmk:ordtorsionfree} (resp.\ Lemma~\ref{lem:levelelltorsion}(3)); $h^{1}_{\mathrm{ord}}(n)=\rank e_1H(1)$
(Definition~\ref{def:hidatower}).
\begin{enumerate}
\item \PROVED{} $e_\ell\,\mathrm{dg}_1=\ell\,(U^{(\ell)}_\ell)^{-1}e_\ell\,\mathrm{dg}_\ell$ ($U^{(\ell)}_\ell$ is invertible on
  $e_\ell H(\ell)$). In particular $e_\ell\circ\mathrm{dg}_1\equiv0\pmod\ell$: \textbf{in this document's convention the
  restriction is not the integral stabilisation} (Remark~\ref{rmk:stabconvention}).
\item \PROVED{} $c_{\mathrm{st}}:=\omega\circ e_\ell\circ\mathrm{dg}_\ell$ restricted to $e_1H(1)$ is an automorphism of the free
  $\Z_\ell$-module $e_1H(1)$; it lies in the closure of $\Z_\ell[T_\ell]$ and is $\equiv T_\ell$ modulo $\ell$ there.
\item \PROVED{} (the last clause \textsf{[standard]}) So $e_\ell\,\mathrm{dg}_\ell\colon e_1H(1)\to e_\ell H(\ell)$ is a split injection of free $\Z_\ell$-modules
  (left inverse $c_{\mathrm{st}}^{-1}\omega$) whose image is $U^{(\ell)}_\ell$-stable; $U^{(\ell)}_\ell$ acts on it as
  $u_{\mathrm{st}}$, the root of $u^{2}-T_\ell u+\ell^{\,n+1}=0$ that is invertible (the unit root), and it
  intertwines $T_p$ with $T^{(\ell)}_p$ for every prime $p\ne\ell$ \STD{\cite[\S3.1 and Ch.~8]{shimura}: the Hecke
  rings act on cohomology as rings, and the degeneracy correspondences commute with $T_p$ for $p$ prime to the
  level; pointers at section level, \textbf{not} verified (rule~\ref{A2}) --- slot-13 WP-6 tests it}.
\item \textbf{(Ranks.)} $\rank_{\Z_\ell}e_\ell H(\ell)=h^{1}_{\mathrm{ord}}(n)$ (\textsf{[standard]}, the citation at the end of the item;
  ranks do not depend on the lattice, Remark~\ref{rmk:stabconvention}(3)). \PROVED{} from this and (3):
  $e_\ell\mathrm{dg}_\ell$ is an
  isomorphism of $\Z_\ell[T_p:p\ne\ell]$-modules $e_1H(1)\cong e_\ell H(\ell)$, and the tower datum of
  Definition~\ref{def:hidatower} (rank, $\chi^{\mathrm{ord}}_{n,p}\bmod\ell^{m}$, $p\ne\ell$) is the same at level~$1$ and at
  level~$\ell$ \STD{\cite[Prop.~4.7, p.~572, with $N=1$, $r=1$; and Lemma~4.6, p.~571]{hida_invent}, locators
  \textbf{verified} against the PDF (slot-13 WP-2), transported by Remark~\ref{rmk:stabconvention}; equivalently
  Eichler--Shimura for $\Ga_0(\ell)$ \cite[Thm.~8.4, p.~234]{shimura} (verified, R20-L: $S_{n+2}(\Ga,\Psi)\cong H^{1}_{P}(\Ga,X_n^{\Psi})$, the
  parabolic part, $\Rl$-linearly) with the newform theory of \cite{atkin_lehner}
  ($a_\ell(f)^{2}=\ell^{\,k-2}$ for an $\ell$-new eigenform of weight $k$ and trivial character), the latter pointer
  not verified}.
\item \textbf{($k=2$.)} \textsf{[standard]} (the citations in the item) At $n=0$: $H(1)=0$, while $H(\ell)=H^{1}(\Ga_0(\ell);\Z_\ell)$ is free of rank
  $2g(X_0(\ell))+1$ and $e_\ell=\id$ on it (the weight-$2$ Eisenstein class has $U^{(\ell)}_\ell$-eigenvalue $1$, an
  $\ell$-new form $a_\ell=\pm1$) \STD{\cite[Thm.~8.4, p.~234]{shimura} for the parabolic part of rank $2g(X_0(\ell))$ (verified, R20-L; the
 Eisenstein rank $1$ is not part of Thm.~8.4); \cite{atkin_lehner}, not verified}; and
  $H^{1}(\Ga_0(\ell);\Fl)=H(\ell)/\ell$. The genus is $0$ at $\ell\in\{5,7,13\}$ and $1$ at $\ell=11$.
\end{enumerate}
\COMP\textsuperscript{a} evidence (the same script, over $\Fl$): $\dim\eord H^{1}(\Ga_0(\ell);\Mn)=
\dim\eord H^{1}(\Ga;\Mn)$ at every even $2\le n\le24$ for $\ell\in\{5,7\}$ and every even $2\le n\le12$ for
$\ell=11$ ($3$ at $n=10$, else $1$); at $n=0$ the level-$\ell$ values are $1,1,3$ at $\ell=5,7,11$; and
$e_\ell\,\mathrm{dg}_\ell\,e_1$ has rank $h^{1}_{\mathrm{ord}}(n)$ while $e_\ell\,\mathrm{dg}_1\,e_1=0$ over $\Fl$ at
$(n,\ell)\in\{(10,5),(12,5),(16,7)\}$ --- (1) and (3) modulo $\ell$. At $\ell=13$ (the WP-3 referee's re-run of the
same script): the level-$13$ value is $1$ at $n=0$, and $1=1$ at $n=2$.
\end{proposition}

\begin{proof}
(1) Proposition~\ref{prop:oldspace}(1) on cohomology, then $e_\ell$ (which commutes with $U:=U^{(\ell)}_\ell$) and
$U$ invertible on $e_\ell H(\ell)$ (Proposition~\ref{prop:ordinary}(1), in the limit).

(2) Let $\mathsf M$ act on $H(1)^{\oplus2}$ as in Proposition~\ref{prop:oldspace}(4), a finitely generated
$\Z_\ell$-module; $U\Phi=\Phi\mathsf M$, so $e_\ell\Phi=\Phi\,\eord(\mathsf M)$ (Proposition~\ref{prop:ordinary}(3)
modulo each $\ell^{j}$, then the limit), and $\omega\Phi(x,y)=(\ell+1)x+T_\ell y$ by
Proposition~\ref{prop:oldspace}(2). Hence $c_{\mathrm{st}}=(\ell+1,\ T_\ell)\,\eord(\mathsf M)\,\binom{0}{1}$, built
from $T_\ell$ by ring operations and limits: it commutes with $T_\ell$, hence preserves $e_1H(1)$, and lies in the
closure of $\Z_\ell[T_\ell]$. Modulo $\ell$ (Proposition~\ref{prop:ordinary}(4)), $n\ge1$ gives $\bar{\mathsf M}=
\mathrm{diag}(0,\bar T_\ell)$ on $(e_1H(1)/\ell)^{\oplus2}$, where $\bar T_\ell$ is invertible; so
$\eord(\bar{\mathsf M})=\mathrm{diag}(0,1)$ and $\bar c_{\mathrm{st}}=\bar T_\ell$, invertible. By Nakayama
$c_{\mathrm{st}}$ is onto the finitely generated $e_1H(1)$, hence bijective; $e_1H(1)$ is free by
Remark~\ref{rmk:ordtorsionfree}(2).

(3) Split injectivity is (2). $\mathrm{dg}_\ell$ and $\omega$ are $\Z_\ell$-linear and $e_\ell H(\ell)$ is free
(Lemma~\ref{lem:levelelltorsion}(3)), so the image $Y$ is a direct summand. By (1) and
Proposition~\ref{prop:oldspace}(3), on $e_1H(1)$: $Ue_\ell\mathrm{dg}_\ell=e_\ell\mathrm{dg}_\ell T_\ell-\ell^{\,n+1}
U^{-1}e_\ell\mathrm{dg}_\ell$, i.e.\ $(U^{2}-UT'+\ell^{\,n+1})e_\ell\mathrm{dg}_\ell=0$ with $T'$ the transport of
$T_\ell$. Over $\Q_\ell$, $Y\otimes\Q_\ell$ is the image of $\eord(\mathsf M)$ on the pairs, which is
$U$-stable; a direct summand whose rationalisation is $U$-stable is $U$-stable. The transported $U$ satisfies
$u^{2}-T_\ell u+\ell^{\,n+1}=0$ and is invertible; it is the unit root because the other root is
$\ell^{\,n+1}u^{-1}\in\ell\End$. The $T_p$-equivariance is the cited standard fact.

(4) $e_\ell H(\ell)$ and $e(U^{\ast})H(\ell)$ have the same rank (Remark~\ref{rmk:stabconvention}(1)), and Hida's
Proposition~4.7 with $N=1$, $r=1$ identifies the latter with the level-$1$ ordinary part (whose rank is
$h^{1}_{\mathrm{ord}}(n)$, the level-$1$ $T(p)$ being $T_\ell$ up to the same convention). A split injection of free
$\Z_\ell$-modules of equal rank is an isomorphism. Through the standard route instead: $H(\ell)\otimes\Cx$ is two
copies of $S_k(\Ga_0(\ell))$ plus two Eisenstein classes; $S_k(\Ga_0(\ell))$ is the old space, image of level~$1$
under the two degeneracies, plus the $\ell$-new space, on which $U^{2}=\ell^{\,k-2}=\ell^{\,n}$, non-ordinary for
$n\ge1$; so the ordinary part is the ordinary part of the old space, of rank $h^{1}_{\mathrm{ord}}(n)$ by (3).

(5) is read off the cited statements: at $n=0$ the ordinary part is everything.
\end{proof}

\begin{remark}[Two conventions for the level-$\ell$ operator]\label{rmk:stabconvention}
\PROVED{} (1) Put $\beta_\ell:=\smat{\ell}{0}{0}{1}$ and $U^{\ast}:=[\Ga_0(\ell)\,\beta_\ell\,\Ga_0(\ell)]$. The Atkin--Lehner
element $W_\ell:=\smat{0}{-1}{\ell}{0}\in\De$ normalises $\Ga_0(\ell)$ and $W_\ell gW_\ell^{-1}=\beta_\ell$; so the
compatible pair $(\conj_{W_\ell},\rhon(W_\ell))$ induces an endomorphism $w$ of $H^{\bullet}(\Ga_0(\ell);\VR)$ with
$wU^{(\ell)}_\ell=U^{\ast}w$, and $w^{2}=\rhon(-\ell I)=\ell^{\,n}$ (Proposition~\ref{prop:diamond}; $n$ even). Over
$\Q_\ell$, $w$ is invertible: $U^{(\ell)}_\ell$ and $U^{\ast}$ are conjugate, and their ordinary parts have equal rank.
Over $\Z_\ell$ they are not interchangeable, as Proposition~\ref{prop:stabilisation}(1) shows. (2) Hida's
$T(p)$ at level $p^{r}$ is the double coset $[\Phi\smat100p\Phi]$ acting through the main involution
$\alpha\mapsto\alpha^{\iota}=\adjg\alpha$ \cite[(4.6a) and p.~563]{hida_invent}; in degree~$0$ it is
$x\mapsto\sum_{i<p}\smat{1}{-i}{0}{p}^{\iota}x$ \cite[proof of Lemma~4.6, p.~571]{hida_invent}, i.e.\ a sum over
matrices $\smat{p}{i}{0}{1}\in\Ga_0(p)\beta_p\Ga_0(p)$ --- in this document's convention $U^{\ast}$, not
$U^{(\ell)}_\ell$ (compare Lemma~\ref{lem:levelelltorsion}(1): here $U=0$ on $H^{0}$ modulo $\ell$, there $T(p)\ne0$ but
$T(p)^{2}=0$, as Hida states). Hida's isomorphism ``$e\circ\res$'' is therefore $e(U^{\ast})\circ\mathrm{dg}_1$ here, and
this document's integral stabilisation is $e_\ell\circ\mathrm{dg}_\ell$ (Proposition~\ref{prop:stabilisation}(3)).
(3) \textbf{Hida's lattice is not $\Vcoef$} (v1.13.0-pre.2, WP-3 referee M1; the pre.1 sentence ``up to the order of
the basis'' was false for $n\ge\ell$). Hida's $L_n(\Z)=\Z^{n+1}$ is the column vector ${}^{t}(x^{n},x^{n-1}y,\dots,y^{n})$
of monomials in a variable vector ${}^{t}(x,y)$, on which $A=\smat abcd\in\Mat_2(\Z)$ acts by
${}^{t}(x,y)\mapsto{}^{t}(ax+by,cx+dy)$ \cite[p.~566]{hida_invent}. Its matrix in that basis is
\[
  \rho^{\mathrm H}_n(A)\;:=\;\rhon(A^{T})^{T}\;=\;D^{-1}\rhon(A)\,D,\qquad D:=\mathrm{diag}\Bigl(\tbinom n0,\tbinom n1,\dots,\tbinom nn\Bigr).
\]
The first equality is \eqref{eq:rhondef} read for $A^{T}$: row $j$ of $\rho^{\mathrm H}_n(A)$ holds the coefficients of
$(ax+by)^{n-j}(cx+dy)^{j}$ in the $u_t$, which is column $j$ of $\rhon(A^{T})$. For the second, the rational symmetric form
$\langle u_j,u_t\rangle:=\delta_{jt}\binom nj^{-1}$ satisfies $\langle v^{n},w^{n}\rangle=\langle v,w\rangle^{n}$ for $v,w\in W\otimes\Q$
(the standard form on $W$; check on $v=x+sy$, $w=x+ty$), the $v^{n}$ span $\Vcoef\otimes\Q$, and $\langle Av,w\rangle=\langle v,A^{T}w\rangle$;
so $\rhon(A^{T})$ is the adjoint of $\rhon(A)$ for a form of Gram matrix $D^{-1}$, $\rhon(A^{T})=D\rhon(A)^{T}D^{-1}$; transpose.
(Checked symbolically at $n\le5$ by the referee and by the author.) So $\rho^{\mathrm H}_n$ is multiplicative on $\Mat_2(\Z)$, and
$f\mapsto Df$ is an injective $\Mat_2(\Z)$-map $L_n(\Z)\to\Vcoef$ with image $\bigoplus_j\binom nj\Z\,u_j$, of index
$\prod_j\binom nj$ --- the image of the divided powers $\Gamma^{n}\Z^{2}$ under the canonical map to $\Sym^{n}\Z^{2}$. Hence
$L_n(\Z_\ell)=\Vcoef\otimes\Z_\ell$ (through $D$) iff $\ell\nmid\binom nj$ for every $j$, iff (Lucas) $n+1=c\,\ell^{s}$ with
$1\le c\le\ell-1$ --- true for $n<\ell$ (so at $(10,691)$, $(10,2411)$), false at $(10,5)$, $(700,691)$ and every other tower cell
with $n\ge\ell$ not of that form. \emph{Lattice-independent} (unaffected): every $\Z_\ell$-rank of an ordinary part and every
characteristic polynomial of a Hecke operator on it (both live on $e H\otimes\Q_\ell$, and $D_\ast\colon H^{1}(-;L_n)\to H^{1}(-;\Vcoef)$ is a
Hecke-equivariant isomorphism after $\otimes\Q$) --- so Proposition~\ref{prop:stabilisation}(4), \ref{Hi3}, \ref{Hi4}, \ref{Hi7},
\ref{Hi9}--\ref{Hi12} and Definition~\ref{def:hidatower}'s datum do not see the difference, and
Proposition~\ref{prop:stabilisation}(1)--(3) are proved for $\Vcoef$ and need nothing from Hida. \emph{Lattice-dependent}: the
isomorphism class of a finite-level module, i.e.\ Definition~\ref{def:hidalattice}'s datum and \ref{Hi8}.
\emph{Addendum} (v1.21.0-pre.2; not an erratum: ``lattice-dependent'' here means ``not shown independent'', and no sentence of this remark is false):
since [U] v0.5.0-pre.2 (\file{docs/latex/hida\_hecke\_algebra.tex}, \texttt{lem:u-latticeord}, \PROVED{} there, elementary; slot-21 referee F4,
\file{docs/implementation/slot21/P8\_F4\_referee\_report.md}, ACCEPT WITH REPAIRS) the finite-level datum is lattice-\emph{independent} on the ordinary part:
$D_\ast\colon e_1H^{q}(\Ga;L_n(\Z_\ell))\cong e_1H^{q}(\Ga;\Sym^{n}\otimes\Z_\ell)$ for every $q$, also on $\partial$ and at level~$\ell$; so
Definition~\ref{def:hidalattice}'s $M_m(n)$ and its variant (H) are isomorphic Hecke modules, \ref{Hi8}'s arms (a) and (b) coincide, and its report (c)
($45/45$ isomorphisms, slot-13 WP-5) is a theorem.
(4) \textbf{In Hida's convention the trace is not an integral left inverse} (v1.13.0-pre.2, WP-3 referee m8). From
$w\,\mathrm{dg}_1=\mathrm{dg}_\ell$, $w\,\mathrm{dg}_\ell=\ell^{\,n}\mathrm{dg}_1$ (from $W_\ell=gS$, $\conj_S$ inner on $\Ga$, and $w^{2}=\ell^{\,n}$),
$U^{\ast}=wUw^{-1}$ and Proposition~\ref{prop:oldspace}(1),(3): $U^{\ast}\Phi=\Phi\,\mathsf M^{\ast}$ with
$\mathsf M^{\ast}=\smat{T_\ell}{\ell^{\,n+1}}{-1}{0}$, the textbook old-space matrix, of the same characteristic polynomial as
$\mathsf M$. On $e_1H(1)$, $n\ge1$, with $u$ the unit root and $u'=\ell^{\,n+1}u^{-1}$ the other: the unit-root eigenspace of
$\mathsf M^{\ast}$ is $\{(-uy,y)\}$, and $\omega\circ e(U^{\ast})\circ\mathrm{dg}_1=\bigl((\ell+1)u-T_\ell\bigr)(u-u')^{-1}
=\ell\,(u-\ell^{\,n}u^{-1})(u-u')^{-1}$, which is $\ell$ times a unit. So the plain trace $\omega$ is \emph{not} a left inverse
of $e(U^{\ast})\circ\mathrm{dg}_1$ over $\Z_\ell$ (how Hida's ``inverse induced by $e_0\circ\mathrm{Tr}$'' is normalised is not
transcribed here; his proof goes through \cite[Cor.~4.5]{hida_invent}, per the WP-3 referee's reading); slot-13 WP-6 tests the pair
$(e_\ell\circ\mathrm{dg}_\ell,\ c_{\mathrm{st}}^{-1}\omega)$ of Proposition~\ref{prop:stabilisation}(3), never a trace inverse in the
$U^{\ast}$ convention. \emph{Proof of (1)}: $W_\ell\smat{a}{b}{\ell c}{d}W_\ell^{-1}=
\smat{d}{-c}{-\ell b}{a}$; $W_\ell gW_\ell^{-1}$ by multiplying out; conjugating each factor of $[A\,g\,B]$ by
$W_\ell$ gives the intertwining, and $W_\ell^{2}=-\ell I$.
\end{remark}

\begin{theorem}[Hida: $\Lambda$-freeness and control]\label{thm:hidacontrol}
\STD{\cite[Thm.~3.1, p.~562; Thm.~6.6, p.~590]{hida_invent}; \cite[Thm.~3.1 and Cor.~3.2, p.~250]{hida_ens};
locators \textbf{verified} against the PDFs in the owner's local library (papers/h/Hida/)
(slot-13 WP-2)} Let $p=\ell\ge5$, Hida's $N$ equal to $1$ (so his $\Phi_r=\Ga_1(N)\cap\Ga_0(p^{r})$ is
$\Ga_0(\ell^{r})$ and his $N_0=1$), $\Lambda=\Z_\ell[[1+\ell\Z_\ell]]$, and $e$ Hida's ordinary idempotent
(Remark~\ref{rmk:stabconvention}(2)). \emph{Hypotheses}: Hida assumes throughout that $p\ge5$ \cite[p.~545]{hida_invent}
(read from the page image; the extracted text misreads it as ``$p>5$''), and admits every $N$ prime to $p$, so $N=1$,
$\Phi_1=\Ga_0(\ell)$, $\ell=5$ are inside them; his (4.4) (no finite subgroups) is used only for the sheaf comparison
(4.5), and at $N=1$ the $\delta\equiv S\pmod p$ of Lemma~4.6 is $S$ itself (v1.13.0-pre.2, WP-3 referee m5).
\emph{Lattice}: $L_n(A)$ is Hida's lattice, $L_n(\Z_\ell)\cong\bigoplus_j\binom nj\Z_\ell\,u_j\subseteq\Vcoef\otimes\Z_\ell$
(Remark~\ref{rmk:stabconvention}(3)), equal to $\Vcoef\otimes\Z_\ell$ for $n<\ell$ and in general \emph{not} for $n\ge\ell$; the
statements below are about it, and are transcribed faithfully (v1.13.0-pre.2, WP-3 referee M1).
\begin{enumerate}
\item \emph{(Invent.\ Thm.~3.1.)} The Pontryagin duals $V^{0}$, $W^{0}$ of the ordinary parts of
  $\varinjlim_rH^{1}(X_r,\Q_\ell/\Z_\ell)$ and $\varinjlim_rH^{1}(Y_r,\Q_\ell/\Z_\ell)$, $Y_r=\Ga_1(\ell^{r})\backslash\Hyp$ and
  $X_r$ its compactification, are free of finite rank over $\Lambda$; and for every $r>0$ restriction identifies the
  level-$r$ ordinary part with the invariants of $1+\ell^{r}\Z_\ell$.
\item \emph{(Invent.\ Thm.~6.6.)} For every $m>0$, $eH^{1}_{P}(\Phi_1,L_n(\Z_\ell))\otimes\Z/\ell^{m}\cong
  eH^{1}_{P}(\Phi_1,L_n(\Z/\ell^{m}))$ ($H^{1}_{P}$ the parabolic cohomology); and for $n\equiv a\pmod{\ell-1}$,
  $eH^{1}_{P}(\Phi_1,L_n(\Z_\ell))\otimes\Q_\ell/\Z_\ell\cong eH^{1}_{P}(\Phi_1,L_n(\Q_\ell/\Z_\ell))\cong
  \mathcal V^{0}(a)[P_n]$ as modules over Hida's universal Hecke algebra, $\mathcal V^{0}$ the (discrete) ordinary part
  of $\varinjlim_rH^{1}(X_r,\Q_\ell/\Z_\ell)$ --- so $V^{0}$ is its Pontryagin dual --- $\mathcal V^{0}(a)$ its
  $a$-component and $P_n$ Hida's weight-$n$ prime of $\Lambda$ (defined in \cite[\S6]{hida_invent}; its normalisation is
  not restated here --- v1.21.0-pre.2: it is \cite[Introduction, p.~545]{hida_invent}, read on the page by the slot-21 referee F4 (M2): $P_{k-2}$ is the kernel of
  the evaluation $X\mapsto u^{k}-1$, $u=1+\ell$ (trivial character), so Hida's $P_n$ is $(1+X)-(1+\ell)^{n+2}$, the $P_k$ of (3) with $k=n+2$; no twist of the
  $\Lambda$-structure intervenes).
\item \emph{(ENS Thm.~3.1, Cor.~3.2; level $N=p$, $N_0=1$.)} The ordinary parts of the universal Hecke algebras
  (of modular forms, of cusp forms) are free of finite rank over $\Lambda$; with $P_k=(1+X)-(1+\ell)^{k}$, the
  cuspidal one satisfies $h^{0}(N,a)/P_kh^{0}(N,a)\cong h^{0}_k(\Ga_1(N_0\ell),a)$ for $k\ge2$, and the rank of
  $h^{0}(N,a)$ over $\Lambda$ equals the $\Z_\ell$-rank of the ordinary weight-$k$ cusp forms for $k\ge3$,
  $k\equiv a\pmod{\ell-1}$.
\end{enumerate}
\end{theorem}

\noindent\emph{To Theorem~\ref{thm:hidacontrol}(2), after ``no twist of the $\Lambda$-structure intervenes} (placed after the theorem, so that its
statement is unchanged). \emph{Addendum} (v1.22.0-pre.1, slot-22 P10-C / F-U5): ``no twist intervenes'' is relative to Hida's own $\Lambda$-structure on $\mathcal V^{0}$, the one
induced by $f|z=z_p^{k}\,(f|_k\sigma_z)$ \cite[p.~549, p.~552]{hida_invent}, i.e.\ $\iota(l)=l^{2}T(l,l)$; on the weight-$2$ tower $\mathcal V^{0}$ this is $1+X\mapsto u^{2}\langle u\rangle$
with $\langle u\rangle$ the natural diamond (``the twist by $z\mapsto z_p^{2}$'', \cite[p.~593]{hida_invent}, read on the page for [U] part~5), and the weight-$k$ classes are those with
$\langle u\rangle=u^{k-2}$ (\cite[p.~589]{hida_invent}: $H^{1}(\Ga_1(Np^{r}),n;\cdot)=\{v:v|z=z^{n}v\}$, $n=k-2$), which is why they sit at $X=u^{k}-1$ --- relative to the natural action alone
they would sit at $u^{k-2}-1$. The index $a$ of $\mathcal V^{0}(a)$ above is the \emph{natural} one (\cite[p.~589]{hida_invent}: $\mathcal V^{0}(a)=\{v:v|\zeta=\zeta^{a}v\}$, $a\equiv n=k-2$): the
weight-$k$ classes of the class $k\bmod(\ell-1)$ land in the component of nebentypus $\omega^{k-2}$ --- Hida's ``$h^{0}_2(a)$ acts on $\mathcal V^{0}(a-2)$'' (p.~593) with $a$ there the Hecke class.
Verified on the page images by the slot-22 referee F-U5 (\file{docs/implementation/slot22/P10C\_F4\_referee\_report.md}, item~7).


\begin{theorem}[Hida: weight independence of the ordinary rank, and the level-$p$ comparison]\label{thm:hidaweight}
\STD{\cite[Thm.~3.1, p.~250]{hida_ens}; \cite[Lemma~4.6, p.~571; Prop.~4.7, p.~572]{hida_invent}; locators
\textbf{verified} (slot-13 WP-2)} Let $\ell\ge5$. (1) For weights $k\ge3$ in one class $a$ modulo $\ell-1$, the
$\Z_\ell$-rank of the ordinary cusp forms of weight $k$ on $\Ga_0(\ell)$ (trivial character, as $k\equiv a$) is the
$\Lambda$-rank of $h^{0}(\ell,a)$, independent of $k$; by the control of Theorem~\ref{thm:hidacontrol}(3) at $k=2$
the same holds from $k=2$ for the cusp forms. (2) For every $n\ge0$ the ordinary parts
$e_0H^{1}(\SL_2(\Z),L_n(\Z_\ell))$ and $eH^{1}(\Ga_0(\ell^{r}),L_n(\Z_\ell))$, $r\ge1$, are $\Z_\ell$-free. (3) For $n>0$,
$r>0$: $e\circ\res\colon e_0H^{1}(\SL_2(\Z),L_n(\Z_\ell))\to eH^{1}(\Ga_0(\ell^{r}),L_n(\Z_\ell))$ is an isomorphism, with
inverse induced by $e_0\circ\mathrm{Tr}$, and likewise for the parabolic cohomology --- in Hida's convention
(Remark~\ref{rmk:stabconvention}). $L_n$ in (2), (3) is Hida's lattice, not $\Vcoef$ for $n\ge\ell$
(Remark~\ref{rmk:stabconvention}(3)); the consequences below use ranks only, which do not depend on the lattice.
\emph{Consequences in this document}: Proposition~\ref{prop:stabilisation}(4); and
with Eichler--Shimura (Obligation~\ref{obl:oracle}), $h^{1}_{\mathrm{ord}}(n)=2\,\rank S^{\mathrm{ord}}_{n+2}+1$ is
constant along $\mathrm{Tw}_1$ --- the source of \ref{Hi3}, whose locators this theorem supplies (its tag is not
moved at this version). Algebra-level consequences (residual ordinary systems and $h^{\mathrm{ord}}_k/\ell$ constant along a class)
do \emph{not} follow from (1); they are Proposition~\ref{prop:hidaalgconst}, from Theorem~\ref{thm:hidacontrol}(3) and
Proposition~\ref{prop:stabilisation} (pointer added at v1.15.0-pre.1, P18; no tag moved).
\end{theorem}

\subsection{The finite-level lattice datum}\label{sec:hidalatticesec}

\begin{definition}[The lattice datum of a tower cell]\label{def:hidalattice}
Let $\ell\ge5$, $n\ge2$ even, $m\ge1$, and $\mathcal P$ a finite set of primes $p\ne\ell$ (default $\{2,3\}$). The lattice is
\textbf{the contract's} $\Vcoef=\Sym^{n}\Z^{2}$ of Definition~\ref{def:rhon}; Hida's $L_n$ is a different $\Z_\ell$-lattice for
$n\ge\ell$ (Remark~\ref{rmk:stabconvention}(3)), and its datum is the variant (H) below. With
$H(n):=H^{1}(\Ga;\Sym^{n}\Z^{2})\otimes\Z_\ell$ and $U=U_\ell$ as in Definition~\ref{def:hidatower}, put
\[
  M_m(n):=\eord H(n)/\ell^{m}\eord H(n),
\]
a free $\Z/\ell^{m}$-module of rank $h^{1}_{\mathrm{ord}}(n)$ (Corollary~\ref{cor:ordprecision}(1): it \emph{is}
$\eord(U|_{H_m})H_m$, $H_m=H^{1}(\Ga;\Sym^{n}\Z^{2}\otimes\Z/\ell^{m})$, computed over $\Z/\ell^{m}$), with the
operators $T_p$, $p\in\mathcal P$, and its Hecke algebra $\Talg^{\mathrm{ord}}_{m,\mathcal P}(n)$, the
$\Z/\ell^{m}$-subalgebra of $\End(M_m(n))$ generated by them. The \textbf{lattice datum of the cell at level $m$} is
the isomorphism class of the $\Z/\ell^{m}[X_p:p\in\mathcal P]$-module $(M_m(n),(T_p)_{p\in\mathcal P})$ ---
equivalently the orbit of the tuple of matrices $(A_p)_{p\in\mathcal P}$ in a basis under simultaneous conjugation by
$\GL_{h}(\Z/\ell^{m})$, $h=h^{1}_{\mathrm{ord}}(n)$. By Proposition~\ref{prop:stabilisation}(4) the same datum is read at
level $\ell$ on $e_\ell H^{1}(\Ga_0(\ell);\Sym^{n}\Z^{2}\otimes\Z/\ell^{m})$ through $e_\ell\mathrm{dg}_\ell$.
$U_\ell$ is \emph{not} part of the datum: the level-$1$ $U_\ell=T_\ell$ transports to $u_{\mathrm{st}}+\ell^{\,n+1}
u_{\mathrm{st}}^{-1}$, not to Hida's $U$ (Proposition~\ref{prop:stabilisation}(3)). What is \emph{not} computed:
Hida's $\Lambda$-adic module and algebra themselves; what is computable is this finite-level datum along
$\mathrm{Tw}_m$.

\emph{Variant (H): the datum on Hida's lattice} (v1.13.0-pre.2, WP-3 referee M1(v)). Replace $\rhon$ by
$\rho^{\mathrm H}_n(A)=\rhon(A^{T})^{T}$ of Remark~\ref{rmk:stabconvention}(3) \emph{everywhere} $\rhon$ enters the model: the relation
matrices of $(K,\Rel)$, the coefficients of the $\tau$-chains and the blocks $\rhon(\delta)$ of \eqref{eq:heckechain}.
$\rho^{\mathrm H}_n$ is integral and multiplicative on $\Mat_2(\Z)$ with $\rho^{\mathrm H}_n(\lambda I)=\lambda^{n}$, and $D$-conjugate to $\rhon$
over $\Q$, so the constructions of \S\ref{sec:chainlevel} apply to it unchanged, and $T_p$ is
\eqref{eq:heckechain} with $\rho^{\mathrm H}_n$ (at level $1$ Hida's $T(p)=[\Phi\smat100p\Phi]$ acting through $\iota$ is the double coset
$\Ga\smat p001\Ga=\Ga\smat100p\Ga$, so no convention question arises there). Put $H^{\mathrm H}(n):=H^{1}(\Ga;L_n(\Z))\otimes\Z_\ell$
and $M^{\mathrm H}_m(n):=\eord H^{\mathrm H}(n)/\ell^{m}$: free of rank $h^{1}_{\mathrm{ord}}(n)$ (Theorem~\ref{thm:hidaweight}(2) for the
freeness, the rank being lattice-independent), and computed over $\Z/\ell^{m}$ as $\eord(U|_{H^{\mathrm H}_m})H^{\mathrm H}_m$,
$H^{\mathrm H}_m=H^{1}(\Ga;L_n(\Z/\ell^{m}))$, by the argument of Corollary~\ref{cor:ordprecision}(1) (the $H^{2}$ term is killed by
$24$, Claim~\ref{cl:amalgam}(4), for any coefficient module). Its datum is the class of $(M^{\mathrm H}_m(n),(T_p)_{p\in\mathcal P})$ as
above. \emph{Implementation}: one \code{FibreRepresentation} variant, like the contragredient and \code{QuotientL} --- the shipped
\code{ComputeSymmetricPowerMatrixFM} with the input matrix transposed and the output transposed. (Through v1.13.0-pre.2
this sentence ended ``nothing else changes''; amended at v1.13.0, as the slot-13 WP-13 code audit's A3 required.)
\emph{One further step changes, on the module route over $\Z/\ell^{m}$} (slot-13 WP-5): the equivariant cochain module
$C^{q}_{\mathrm{eq}}=\ker A_{C^q}$ is free on $L_n$ as on $\Vcoef$ (a direct summand, the stabilisers having order prime to
$\ell$), but its Howell form may carry a non-unit pivot, a property of the coordinate order and not of freeness (on $\Vcoef$
every pivot is a unit, Proposition~\ref{prop:coneintegral}(2), and the step never runs). The route then takes a
\emph{Nakayama basis} --- the Howell rows independent modulo $\ell$, chosen greedily --- and checks it free by
$\mathrm{length}(\ker A_{C^q})=m\cdot\mathrm{rank}$, refusing by name otherwise. Every object computable before is computed
on the unchanged unit-pivot branch, so the step is value-preserving there. The map
$D_\ast\colon M^{\mathrm H}_m(n)\to M_m(n)$ induced by $f\mapsto Df$ is Hecke-equivariant; whether it is an isomorphism (iff it is
surjective modulo $\ell$) is a lattice comparison, recorded per cell.
\end{definition}

\begin{remark}[What is checkable at level $\ell$, along $\mathrm{Tw}_3$, and at the heavy cells: (Hi7)--(Hi12)]\label{rmk:hidacheckable2}
\PROVED{} \CONJ{} \COMP\textsuperscript{a} Committed at
v1.13.0-pre.1, before the code that tests them (slot-13 WP-5, WP-6, WP-10; rule~20); each item carries its own tag,
with the convention of Remark~\ref{rmk:hidacheckable}: a failing \PROVED{} item stops the package, a failing
\CONJ{} item is a finding. The \CONJ{} items' sources are Theorems~\ref{thm:hidacontrol}, \ref{thm:hidaweight}
(locators verified); they stay \CONJ{} because the chain from those statements to this document's objects --- Hida's
duality between Hecke algebras and forms, his parabolic cohomology against the full $H^{1}$, his convention against
this one (Remark~\ref{rmk:stabconvention}(2)), his lattice $L_n$ against $\Vcoef$ (Remark~\ref{rmk:stabconvention}(3); it
matters for \ref{Hi8} only) --- is not carried out here.
\begin{enumerate}[label=(Hi\arabic*),start=7]
\item\label{Hi7} \CONJ{} \textbf{Congruences modulo $\ell^{3}$.} For $p\in\{2,3\}$ and $n,n'$ in one
  $\mathrm{Tw}_3(\ell;n_0)$, $\chi^{\mathrm{ord}}_{n,p}\equiv\chi^{\mathrm{ord}}_{n',p}\pmod{\ell^{3}}$, read over $\Z/\ell^{3}$
  (Corollary~\ref{cor:ordprecision}(1)). The period is $(\ell-1)\ell^{2}$: $100$ at $\ell=5$ ($n=10,110,210$), $294$ at
  $\ell=7$ ($n=10,304$, the slow arm, sized by WP-5 before it runs). The Eisenstein factor obeys it unconditionally
  (\ref{Hi4}'s argument with $m=3$).
\item\label{Hi8} \CONJ{} \textbf{The lattice datum is constant on towers.} For $m\in\{1,2\}$, $\mathcal P=\{2,3\}$ and
  $n,n'$ in one $\mathrm{Tw}_m(\ell;n_0)$: $(M_m(n),T_2,T_3)\cong(M_m(n'),T_2,T_3)$ (Definition~\ref{def:hidalattice}) ---
  strictly stronger than \ref{Hi4}. Two arms: \textbf{(a)} on the contract's lattice $\Vcoef$ (Definition~\ref{def:hidalattice});
  \textbf{(b)} on Hida's lattice, $(M^{\mathrm H}_m(n),T_2,T_3)\cong(M^{\mathrm H}_m(n'),T_2,T_3)$ (Definition~\ref{def:hidalattice}, variant
  (H)); and, reported per cell and \emph{not} a prediction, (c) whether $D_\ast\colon M^{\mathrm H}_m(n)\to M_m(n)$ is an isomorphism.
  \emph{Source, stated honestly} (v1.13.0-pre.2, WP-3 referee M1(iv)): Theorem~\ref{thm:hidacontrol}(1),(2) --- the weight-$n$
  \emph{parabolic} module on \emph{Hida's lattice} is the quotient of a $\Lambda$-free module by $(P_n,\ell^{m})$, and
  $(P_n,\ell^{m})=(P_{n'},\ell^{m})$ in the same $a$-component when $n\equiv n'\pmod{(\ell-1)\ell^{m-1}}$. Arm (b) extends that
  from the parabolic part to the full $H^{1}$ (the Eisenstein line added); arm (a) is in addition \textbf{this document's
  extrapolation} to a different lattice, on which no cited statement bears --- the isomorphism class of a finite-level module
  is lattice-dependent (Remark~\ref{rmk:stabconvention}(3)); for $n,n'<\ell$ the two arms coincide. \emph{The test is a decision
  procedure} (v1.13.0-pre.2, WP-3 referee m4): the solution module $\mathcal H:=\{P\in\Mat_h(\Z/\ell^{m}):PA_p=A'_pP,\ p\in\mathcal P\}$
  ($=\Hom_{\Talg}(M,M')$) is computed by a module solve over $\Z/\ell^{m}$; an isomorphism exists iff some $P\in\mathcal H$ is invertible,
  iff (Nakayama) some element of the reduction $\bar{\mathcal H}\subseteq\Mat_h(\Fl)$ is invertible, iff $\det$ restricted to the
  $\Fl$-space $\bar{\mathcal H}$ (of dimension $\le h^{2}$) is not the zero polynomial. It has degree $h$; for $h<\ell$ it is non-zero iff it
  is non-zero at some point of the grid $\{0,\dots,h\}^{\dim\bar{\mathcal H}}$ (a non-zero polynomial of degree $\le h$ in each variable does
  not vanish on a grid of side $h+1$), $\le4^{9}$ evaluations at the grid's ranks $h\le3$; for $h\ge\ell$ all of $\bar{\mathcal H}$ is
  finite. Outcomes: \emph{isomorphic}, with the certificate $P$ printed, or \emph{not isomorphic}, by the vanishing; no
  ``inconclusive''. The invariants of pre.1 (the lengths of $\Hom(M,M)$, $\Hom(M,M')$, $\Hom(M',M')$, the Fitting ideals of $A_p-c$)
  stay as cheap necessary conditions run first.
\item\label{Hi9} \CONJ{} \textbf{\ref{Hi4} at $p\in\{5,7\}$}, $p\ne\ell$, $m\in\{1,2\}$; same source.
\item\label{Hi10} \PROVED{} (Proposition~\ref{prop:stabilisation} and Lemma~\ref{lem:levelelltorsion}, with their
  \textsf{[standard]} pointers) \textbf{Level $\ell$ against level $1$.} At $\ell\in\{5,7,11,13\}$ and even $n\ge2$ (WP-6's grid: $n\le40$):
  $\dim_{\Fl}\eord H^{1}(\Ga_0(\ell);\Mn)=\rank e_\ell H(\ell)=h^{1}_{\mathrm{ord}}(n)$, the value of \ref{Hi3}; at odd $n$
  both are $0$; at $n=0$ the level-$\ell$ value is $2g(X_0(\ell))+1=1,1,3,1$ against level~$1$'s $0$ --- so the cells
  $(0,\ell)$ enter at level $\ell$ only. Also checkable entrywise on the model: Proposition~\ref{prop:oldspace}(1),(2) at
  cochain level, (3) on cohomology, and $U^{(\ell)}_\ell$ commuting with $T^{(\ell)}_2$, $T^{(\ell)}_3$ --- the commutativity of the
  Hecke ring of $\Ga_0(\ell)$ at primes prime to each other, \STD{\cite[Ch.~3]{shimura}; chapter level, \textbf{not} verified
  (rule~\ref{A2}); v1.13.0-pre.2, WP-3 referee m9}, not proved in this document.
  \COMP\textsuperscript{a}: WP-2's script, recorded in Proposition~\ref{prop:stabilisation}.
\item\label{Hi11} \CONJ{} (the Eisenstein half \PROVED{} by \ref{Hi5}) \textbf{$(700,691)$ lies in
  $\mathrm{Tw}_1(691;10)$} ($700-10=690=\ell-1$): $h^{1}_{\mathrm{ord}}(700)=h^{1}_{\mathrm{ord}}(10)=3$ and
  $\chi^{\mathrm{ord}}_{700,p}\equiv\chi^{\mathrm{ord}}_{10,p}\equiv\bigl(X-(1+p^{11})\bigr)^{3}\pmod{691}$ for $p\in\{2,3\}$ ---
  at $n=10$, $\Delta$ is ordinary at $691$ and $\tau(p)\equiv1+p^{11}$ by the Ramanujan congruence, which the model
  shows at $(10,691)$ (Obligation~\ref{obl:oracle}(i)). Proved: $1+p^{701}\equiv1+p^{11}\pmod{691}$ is a root of
  $\chi^{\mathrm{ord}}_{700,p}$ modulo $691$, and $h^{1}_{\mathrm{ord}}(700)\ge1$. Run only behind the $U_\ell$ memory lever and
  the owner's rule~19 (slot-13 WP-10). \COMP\textsuperscript{a} (added at v1.14.0-pre.1, taken out at pre.2, restored at v1.14.0;
  the first cluster run, jobs \code{34550469}, \code{34550468}, \code{34550470} at commit \texttt{07f413f},
  \file{docs/matrix-engine/results/noctua2\_run1/}, published 2026-09-28 on the owner's sign-off, [S14] Q2; [EX] \texttt{nc:hi11}):
  \textbf{also held at $n=1390$, $2080$ and $2770$}, the next three members of $\mathrm{Tw}_1(691;10)$ ($k$ up to $2772$,
  $\dim H^{1}=233$, $347$, $463$): $h^{1}_{\mathrm{ord}}=3$ and the same three polynomials, $\bigl(X-(1+p^{11})\bigr)^{3}$ for
  $p\in\{2,3\}$ and $(X-1)^{3}$ for $U_{691}$ --- a comparison of printed values, the probe gating (Hi11) at $n=700$ only; and at
  $m=2$ the $(700,691)$ cell reproduces on a second machine, $\chi^{\mathrm{ord}}$ modulo $691^{2}$ equal to slot~13's, with the
  probe's gated \code{HI11 \dots{} verdict=HELD}.
  \COMP\textsuperscript{a} (added at v1.14.1; \emph{after the second cluster run}, second-cluster job \code{3467349} at commit \texttt{e85e25a}, with
  the slot-14 levers on; data signed off by the owner 2026-09-29, [S14] Q-B; [EX] \texttt{nc:hi11}, \texttt{s14:sept};
  the WP-C10 report of the September run, \S1): \textbf{also held at $n=4150$}, the seventh member of
  $\mathrm{Tw}_1(691;10)$ ($k=4152$, $\dim H^{1}=693=2d_k-1$, $d_k=347$): $h^{1}_{\mathrm{ord}}=3$ and the same three polynomials, and
  the probe's own gated verdict line \code{HI11 n=4150 ell=691 h\_ord=3 predicted\_h\_ord=3 verdict=HELD} --- from this build the
  probe prints the (Hi11) verdict at the cell it is given, so ``at $n=700$ only'' above is the first run's build, a dated record. The
  cell took $711.8$~s (the case $715$~s) and $49.9$~GiB (the probe's own peak; the campaign's sampler read the same). The sixth
  member, $n=3460=5(\ell+1)$, is not served: $e_{\max}=1$ there and the probe refuses it by name (second-cluster job \code{3464192}), so the
  computed members are $n=10$, $700$, $1390$, $2080$, $2770$, $4150$.
  \emph{A route to a proof} (v1.24.0-rc.1; the phenomenology catalogue, PH-32): the ordinary rank $3$ and the three polynomials along
  $n\equiv10\bmod690$ would follow from Hida's control theorem modulo $\ell$ at level one ($\ell\ge5$, weight $\ge3$) together with $n=10$; the
  module-level form of that control has not been checked against Hida's text, so the tag stays \CONJ{} and the \COMP\textsuperscript{a} records above stay the evidence.
\item\label{Hi12} \CONJ{} (the Eisenstein half \PROVED{}) \textbf{$(2420,2411)$ lies in $\mathrm{Tw}_1(2411;10)$}
  ($2420-10=2410=\ell-1$): $h^{1}_{\mathrm{ord}}(2420)=h^{1}_{\mathrm{ord}}(10)=1$ --- the Eisenstein line alone, as
  \ref{Hi6} proves at $n=10$ --- and $\chi^{\mathrm{ord}}_{2420,p}=X-(1+p^{2421})\equiv X-(1+p^{11})\pmod{2411}$. Proved:
  $h^{1}_{\mathrm{ord}}(2420)\ge1$ and the Eisenstein root. \COMP\textsuperscript{a} (added at v1.14.0-pre.1, taken out at pre.2,
  restored at v1.14.0; first-cluster job \code{34550471} at commit \texttt{07f413f}, published 2026-09-28 on the owner's sign-off, [S14] Q2;
  [EX] \texttt{nc:hi12}):
  \textbf{run once at $(2420,2411)$, and the prediction held}: $h^{1}_{\mathrm{ord}}=1$, $\chi^{\mathrm{ord}}(T_2)=X+362=X-(1+2^{11})$,
  $\chi^{\mathrm{ord}}(T_3)=X+1266=X-(1+3^{11})$ and $\chi^{\mathrm{ord}}(U)=X-1$ modulo $2411$, equal to $n=10$'s. The Eisenstein half is
  gated by \code{HeckeHidaHeavyCellProbe}; the equality with $n=10$ is a comparison of printed values (the first-cluster build printed no
  (Hi12) verdict; the build \texttt{e85e25a} prints \code{HI12 \dots{} verdict=HELD}, comparing $h^{1}_{\mathrm{ord}}$ with the
  prediction only; v1.14.1).
  \COMP\textsuperscript{a} (added at v1.14.1; \emph{after the second cluster run}, second-cluster job \code{3467349} at commit \texttt{e85e25a},
  levers on; data signed off by the owner 2026-09-29, [S14] Q-B; [EX] \texttt{nc:hi12}, \texttt{s14:sept};
  the WP-C10 report of the September run, \S1): \textbf{also held at $(4830,2411)$}, the third computed
  member of $\mathrm{Tw}_1(2411;10)$ ($4830=10+2\cdot2410$, $k=4832$, $\dim H^{1}=805=2d_k-1$, $d_k=403$): $h^{1}_{\mathrm{ord}}=1$ and
  $\chi^{\mathrm{ord}}(T_2)=X+362$, $\chi^{\mathrm{ord}}(T_3)=X+1266$, $\chi^{\mathrm{ord}}(U)=X-1$ modulo $2411$, equal to $n=10$'s, with
  \code{HI12 n=4830 ell=2411 h\_ord=1 predicted\_h\_ord=1 verdict=HELD}; the cell took $1222.4$~s and $68.4$~GiB (the probe's own peak;
  the campaign's sampler read $65.3$~GiB). $(2420,2411)$ was recomputed with the levers and gave the same values in $189.3$~s and
  $19.2$~GiB (probe; sampler $17.5$~GiB), against $5137$~s and $493.3$~GiB on the first cluster, and the same run's $(1210,2411)$ reproduces the
  degree-$203$ polynomials $\chi^{\mathrm{ord}}(T_2)$, $\chi^{\mathrm{ord}}(T_3)$, $\chi^{\mathrm{ord}}(U)$ of the first-cluster run coefficient
  for coefficient.
\end{enumerate}
\end{remark}

\paragraph{Scope: what this chapter does not do.}
It proves no part of Hida theory: Theorems~\ref{thm:hidacontrol} and \ref{thm:hidaweight} are literature, with
verified locators; \ref{Hi7}--\ref{Hi9}, \ref{Hi11}, \ref{Hi12} are predictions. What is asserted: the level-$N$
correspondences on the induced model for every $N$ (Proposition~\ref{prop:levelcorrespondence}), the level-$\ell$
operator and its old space, the orbit and cokernel structure, and the integral stabilisation at $\ell\ge5$, $n\ge1$.
\emph{At $\ell\in\{2,3\}$} (v1.13.0-pre.2, WP-3 referee m7): the cochain-level statements hold over every ring --- the
reciprocity of Proposition~\ref{prop:levelcorrespondence}(3) is stated over every $\cring$, so the induced model computes the
$\Ga_0(N)$-equivariant (Bredon-type) cohomology of the same tessellation over $\Z/2^{k}$, $\Z/3^{k}$, $\Z$, and
Proposition~\ref{prop:levelcorrespondence}(1),(2), Proposition~\ref{prop:levelellU}, Proposition~\ref{prop:oldspace}(1),(2) and
Lemma~\ref{lem:levelelltorsion}(1)--(3) hold there; what is not asserted there is the cohomology-level independence of $\tau$, the
comparison with $H^{\bullet}(\Ga_0(N);\VR)$ (Proposition~\ref{prop:lhcompliant}(4) needs $6\in\cring^{\times}$), the orbit-wise averaging
and the cokernel of Proposition~\ref{prop:levelorbit}(2),(3), and everything of \S\ref{sec:stabilisation} beyond
Lemma~\ref{lem:levelelltorsion}(1)--(3) --- as at level $1$. Not asserted either: Hida's $\Lambda$-adic module and Hecke algebra as
computed objects; level $\ell^{r}$, $r\ge2$, and general $N$ beyond the interface of Proposition~\ref{prop:levelcorrespondence}
(composite-$N$ coset tables, \ref{levelO3}); a Hecke action on the cokernel presentation of Proposition~\ref{prop:levelorbit}(2) ---
which is the Manin--Merel Heilbronn-matrix formula, specified by pointer in Remark~\ref{rmk:levelscaling}, \ref{levelO2}, and not
built. (v1.14.0-pre.1: both are now stated and proved in \S\ref{sec:levelN}, Definition~\ref{def:cosetrepsN} and
Proposition~\ref{prop:heilbronn}; neither is built.) (v1.14.2: Definition~\ref{def:cosetrepsN} built, Proposition~\ref{prop:heilbronn} not.)

\paragraph{Code map.}
\textbf{Built at v1.13.0} (slot 13; not built at v1.13.0-pre.1 and pre.2, contract first).
\file{Hecke/HeckeLevelEll} (\code{namespace hecke::level}, slot-13 WP-6): one routine for \eqref{eq:levelcorr} ---
$U^{(\ell)}_\ell$, $\widetilde{\mathrm{dg}}_\ell$, $T^{(\ell)}_p$, $T_\ell$ at level~$1$ and, with the identity datum, the fibre maps
$\widetilde{\mathrm{dg}}_1$, $\omega$ --- reading only the level-$1$ $\tau$-table; \eqref{eq:Ulevellmodel} as a second model route; the
operators as block terms applied without a $D_N\times D_N$ matrix, the determinant-$\ell$ blocks stored factored; the induced
relations over $\Z/\ell^{k}$; the reading on $H^{q}$ between two local systems. \code{HeckeHidaLevelEllValidation}: \eqref{eq:Ulevellmodel}
$=$ \eqref{eq:levelcorr} entrywise as cochain maps, both against an independent group-cocycle route (Shapiro on $\langle S,T\rangle$, no
$\tau$) by characteristic polynomials; Proposition~\ref{prop:oldspace}(1), (2) as cochain maps; Proposition~\ref{prop:levelorbit}(3);
Lemma~\ref{lem:levelelltorsion}(1); Proposition~\ref{prop:stabilisation} as the pair $(e_\ell\circ\widetilde{\mathrm{dg}}_\ell, c_{\mathrm{st}}^{-1}\omega)$; \ref{Hi10}
on $92$ cells. \file{HeckeExt/HeckeHidaTower} part~2 (WP-5): towers modulo $\ell^{3}$, $T_5$, $T_7$, Definition~\ref{def:hidalattice}
and its variant (H) (\code{FibreRepresentation::HidaLatticeH}, with the module route's Nakayama basis of the amended
sentence), \ref{Hi8}'s decision procedure and arms (a), (b), (c), \ref{Hi7}, \ref{Hi9}; \code{HeckeHidaControlValidation}. The heavy
cells (WP-10): \code{HeckeHidaHeavyCellProbe} (\code{slow}, interlocked), behind the $U_\ell$ memory lever
(\code{HeckeOperatorOptions::ScopedUEllLowRank}, WP-4, value-preserving; its rank claim is Lemma~\ref{lem:rhorank}, v1.14.1): \ref{Hi11} held at $(700,691)$ over $\F_{691}$ and
$\Z/691^{2}$; \ref{Hi12} not run (predicted beyond the box; deferred to the cluster roadmap). Every \PROVED{} and \CONJ{} item
tested held. Not built: the monomial and orbit-wise routes of \S\ref{sec:levelcost}(2)--(5), the Heilbronn-matrix action of
\ref{levelO2} and the composite-$N$ tables of \ref{levelO3} --- the interface for general $N$, not scheduled (v1.14.0-pre.1: both
specified in \S\ref{sec:levelN}, not built). \textbf{At v1.14.2}: the monomial and orbit-wise routes and the composite-$N$ tables
are built (slot-14 WP-12, \S\ref{sec:levelN}'s code map); the Heilbronn-matrix action is not. \emph{After the release} (v1.14.0's record; the first cluster run,
\file{docs/matrix-engine/results/noctua2\_run1/}, published 2026-09-28 on the owner's sign-off, [S14] Q2): the same probe ran \ref{Hi11} at
$n=1390$, $2080$, $2770$ and $(700,691)$ at $m=2$, and \ref{Hi12} at $(2420,2411)$ (the items' \COMP\textsuperscript{a} sentences).
The latter needed $493.3$~GiB on a $2$~TiB node, $0.912$ of the $U_\ell$-lever fit ($541.2$~GiB) of \ref{levelO5} --- the regime
the cost model placed it in. (v1.14.0-pre.2 kept only this provenance line, under [S14] Q2's default; v1.14.0 restores the rest.)
A second run (v1.14.1's record; second-cluster job \code{3467349}, commit \texttt{e85e25a}, the slot-14 levers on by default --- coset
streaming, the serial phase on the pool and the rank-one closed form of Lemma~\ref{lem:rhorank} beside WP-4's factored blocks;
data signed off by the owner 2026-09-29, [S14] Q-B) added \ref{Hi11} at $n=4150$ and \ref{Hi12} at $(4830,2411)$, both held, and
recomputed $(2420,2411)$ in $19.2$~GiB (the probe's own peaks $49.9$, $68.4$, $19.2$~GiB; the campaign's sampler $49.9$, $65.3$,
$17.5$~GiB). The campaign's admission predictor over-predicted the three cases ($84$, $182$, $46$~GiB) and was refitted
(\texttt{549200e}); the fit of \ref{levelO5} predates the slot-14 levers and is not restated --- it is not a model of a levered run.

\begin{lemma}[The rank of $\rhon$ at a matrix singular modulo $\ell$: the $U_\ell$ lever's rank claim]\label{lem:rhorank}
\PROVED{} (v1.14.1.) Let $\ell$ be a prime, $k\ge1$, $n\ge0$ and $g=\smat{a}{b}{c}{d}\in\Mat_2(\Z)$ with $\ell\mid\det g$.
\begin{enumerate}
\item Over $\Z/\ell^{k}$ the matrix $\rhon(g)$ of Definition~\ref{def:rhon} has at most $\min(k,n+1)$ non-zero invariant factors;
  hence it is a sum of at most $\min(k,n+1)$ outer products $u\,v^{\mathsf T}$, $u,v\in(\Z/\ell^{k})^{n+1}$.
\item Over $\Fl$: if $(a,c)\not\equiv(0,0)$, then $(b,d)\equiv\lambda(a,c)\pmod\ell$ for a unique $\lambda\in\Fl$ and
  $\rhon(g)\equiv w\,v^{\mathsf T}$ with $w_t=\binom{n}{t}a^{\,n-t}c^{\,t}$ (the coefficients of $(a+cy)^{n}$) and $v_j=\lambda^{j}$ ($0^{0}=1$);
  if $(a,c)\equiv(0,0)$, then $\rhon(g)\equiv w\,e_n^{\mathsf T}$ with $w$ the coefficients of $(b+dy)^{n}$. In particular
  $\rank_{\Fl}\rhon(g)\le1$.
\end{enumerate}
\end{lemma}

\begin{proof}
(1) By the Smith normal form over $\Z_\ell$, $g=P\,\mathrm{diag}(\ell^{\alpha},\ell^{\beta})\,Q$ with $P,Q\in\GL_2(\Z_\ell)$ and
$0\le\alpha\le\beta\le\infty$ ($\ell^{\infty}:=0$, $\infty\cdot0:=0$), $\alpha+\beta=v_\ell(\det g)\ge1$. The formula of Definition~\ref{def:rhon} is
polynomial in the entries, so $\rhon$ extends to $\Mat_2(\Z_\ell)$ and stays multiplicative ($\rhon=\Sym^{n}$); $\rhon(P)$, $\rhon(Q)$ are
invertible (with inverses $\rhon(P^{-1})$, $\rhon(Q^{-1})$), and $\rhon(\mathrm{diag}(s,t))$ is diagonal with entries $s^{\,n-j}t^{\,j}$
by \eqref{eq:rhondef}. So the invariant factors of $\rhon(g)$ over $\Z_\ell$, and after reduction over $\Z/\ell^{k}$, are
$\ell^{e_j}$ with $e_j=\alpha(n-j)+\beta j$, $0\le j\le n$, and the non-zero ones modulo $\ell^{k}$ are those with $e_j<k$. If
$\alpha=0$, then $\beta\ge1$, $e_j\ge j$, and at most $\min(k,n+1)$ indices have $e_j<k$. If $\alpha\ge1$, then
$e_j=\alpha n+(\beta-\alpha)j$: for $\beta=\alpha$ all $e_j=\alpha n$, and $e_j<k$ forces $n+1\le\alpha n+1\le k$; for $\beta>\alpha$,
$e_j<k$ forces $j<k-\alpha n\le k$. A matrix $P'DQ'$ with $D$ diagonal with $r$ non-zero entries is a sum of $r$ outer products.
(2) Column $j$ of $\rhon(g)$ holds the coefficients of $(a+cy)^{n-j}(b+dy)^{j}$ \eqref{eq:rhondef}. If $(a,c)\not\equiv0$,
$\det g\equiv0$ makes $(b,d)$ proportional to it, so the column is $\lambda^{j}(a+cy)^{n}$; if $(a,c)\equiv0$, every column $j<n$
vanishes and column $n$ is $(b+dy)^{n}$.
\end{proof}

\noindent\emph{Code} (not a status tag; slot-13 WP-4, slot-14 G0 \texttt{a249a6d}, \texttt{714e553}): \file{Hecke/HeckeLowRank.h}.
\code{hecke::lowrank::factor\_low\_rank} stores a block $\rhon(g)$, $g$ singular modulo $\ell$, as $UV$ by minimum-valuation
pivoting, with exactly as many outer products as non-zero invariant factors, and \code{ENGINE\_CHECK}s their number against
$\min(k,n+1)$ at every block; \code{rank\_one\_rho\_factors} is (2) over $\Fl$ and \code{ENGINE\_CHECK}s $ad\equiv bc$ per key. By the
lemma neither check can fire: they guard the implementation, and the lever cannot change a value (\code{HeckeUEllLowRankValidation},
\code{HeckeUEllStreamValidation}).

\paragraph{Formalization note.}
Propositions~\ref{prop:levelcorrespondence}--\ref{prop:levelorbit} and Lemma~\ref{lem:levelelltorsion} need
Frobenius reciprocity, the comparison of projective resolutions, and the Mayer--Vietoris sequence of an amalgam;
Proposition~\ref{prop:stabilisation}(1)--(3) adds Fitting's lemma and Nakayama. Its part~(4) and
Theorems~\ref{thm:hidacontrol}, \ref{thm:hidaweight} rest on the theory of modular forms and are not
formalization targets.

%======================================================================
\part*{Part VIII --- Level and character}
\addcontentsline{toc}{section}{Part VIII --- Level and character}
%======================================================================

\section{Level \texorpdfstring{$N$}{N}: the composite coset table, the cokernel presentation and its Hecke action}\label{sec:levelN}
%----------------------------------------------------------------------
% NEW at v1.14.0-pre.1 (hecke-extensions slot 14, WP-4; plan docs/implementation/slot14/00-slot14-plan.md
% S2 "WP-4").  The contract first (rule 21): nothing in this chapter is built at this version; the
% predictions (LN1)--(LN6) precede slot-14 WP-12's code (rule 20).  It discharges the obstacles O2 and O3
% of rmk:levelscaling on paper.  Placed after sec:hidafull, before the appendices, so that no chapter
% number moves.  The owner's standing direction (2026-09-26): level N by Shapiro on the fixed model, the
% permutation (x) rho_n block structure exploited orbit by orbit, never a second coset or tau construction.

\paragraph{What this chapter fixes.}
(i) \textbf{The binding coset table for every $N$} (Definition~\ref{def:cosetrepsN}, obstacle \ref{levelO3}): the
representatives of $\Ga_0(N)\backslash\Ga$, their index order and class map, by the Chinese remainder theorem over the
prime-power factors; Definition~\ref{def:cosetreps} verbatim at $N=p^{\nu}$; and the representatives of
$\Ga_0(N')\backslash\Ga_0(N)$ read off the level-$N'$ table. A table, not a search; no new $\tau$-table.
(ii) \textbf{The level-$N$ model} for every $N$ and every coefficient ring (Definition~\ref{def:levelNmodel}), with
$\ell\mid N$ (the Hida setting), $\ell\nmid N$, and the setting $\F_p$, $p\nmid N$, of Mazur's Eisenstein ideal.
(iii) \textbf{The cokernel presentation pinned}: the tree resolution (Lemma~\ref{lem:treeres}), Manin's exact sequence
(Lemma~\ref{lem:manin}), the isomorphism $\varkappa_N$ of Proposition~\ref{prop:levelorbit}(2) made explicit
(Definition~\ref{def:kappaN}), and a duality with cusp-divisor coinvariants that commutes with restriction,
corestriction and conjugation (Lemma~\ref{lem:levelduality}). (iv) \textbf{The Hecke action on the presentation}
(obstacle \ref{levelO2}): Merel's set of integral matrices of determinant $p$, its combinatorial lemma
(Lemma~\ref{lem:merelpath}), and the transcription into this document's conventions, \emph{proved} to induce the
operator of Proposition~\ref{prop:levelcorrespondence}(4)(d) on $H^{1}$ --- for $T^{(N)}_p$, $p\nmid N$, and $U^{(N)}_p$,
$p\mid N$, by one formula (Proposition~\ref{prop:heilbronn}). (v) \textbf{The cost of that route}
(Remark~\ref{rmk:heilbronncost}) and \textbf{the predictions} (LN1)--(LN6) (Remark~\ref{rmk:levelNcheckable}), among
them route R3 of the slot-14 Loewy note: mod-$p$ Loewy length $e+1$ at Mazur's Eisenstein ideal.

\paragraph{Notation and standing conventions.}
Those of \S\ref{sec:hidafull}; $\Ga=\SL_2(\Z)$, $S=\smat0{-1}10$, $T=\smat1101$, $U:=ST=\smat0{-1}11$ (the $U$ of
Claim~\ref{cl:amalgam}), $\rhon$ of Definition~\ref{def:rhon}, $\adjg$ the adjugate of \eqref{eq:adjrules}. For a
prime $p$, $\De_p:=\{\delta\in\Mat_2(\Z):\det\delta=p\}$. \textbf{Cusps}: $\mathbb{P}^1(\Q)$ with the \emph{column}
action $\smat abcd\cdot(x:y):=(ax+by:cx+dy)$ of every $\delta\in\De$; scalars act trivially, so $\delta$ and $\adjg\delta=
(\det\delta)\delta^{-1}$ act as mutually inverse bijections. $[\infty]:=(1:0)$, $[0]:=(0:1)$, $[-1]:=(-1:1)$;
$\mathrm{Div}:=\Z[\mathbb{P}^1(\Q)]$, $\mathrm{Div}^0$ the kernel of the degree, $D_0:=[\infty]-[0]\in\mathrm{Div}^0$;
$\mathrm{Div}_{\cring}:=\mathrm{Div}\otimes\cring$, likewise $\mathrm{Div}^0_{\cring}$. \textbf{Level}: $\mathbb{P}^1(\Z/N)$
with the \emph{row} action $(c:d)\cdot\delta:=\bigl((c,d)\delta\bmod N\bigr)$ of any integral $\delta$ for which the
result is primitive ($\gcd(c',d',N)=1$) --- always when $\gcd(\det\delta,N)=1$. \textbf{The tree}: $\mathcal Y$ is the
Bass--Serre tree of Claim~\ref{cl:amalgam}(1); its chain complex is $C_1(\mathcal Y)=\Z[\Ga/\langle-I\rangle]$ with the
generator $\mathfrak e:=\langle-I\rangle$, $C_0(\mathcal Y)=\Z[\Ga/\langle S\rangle]\oplus\Z[\Ga/\langle U\rangle]$ with the
generators $\mathfrak v_S$, $\mathfrak v_U$, $\partial(\gamma\mathfrak e)=\gamma\mathfrak v_U-\gamma\mathfrak v_S$, and
augmentation $1$ on every vertex --- the sequence \eqref{eq:treeses}, whose maps that claim leaves unnamed (a sign
convention, which moves no statement). The \emph{edges} are $\mathcal Y_1:=\Ga\mathfrak e$, the \emph{vertices}
$\mathcal Y_0:=\Ga\mathfrak v_S\sqcup\Ga\mathfrak v_U$; $-I$ fixes every edge and vertex. For $a,b$ vertices, $\mathrm{path}(a\to b)$ is
the unique $1$-chain with boundary $b-a$ (unique: $\partial$ is injective on $C_1(\mathcal Y)$, \eqref{eq:treeses}).
$\mathfrak e,\mathfrak v_S,\mathfrak v_U,\mathcal Y,\varkappa_N,\Pi_A,\mathsf b,\mathrm{Mer}_p,\mathrm{Mer}_n,W_A,W_{N,p},\mu_N,\zeta_A,H_{N,p},
\mathrm{Div},D_0$ are new symbols of this chapter and mean nothing else (rule~\ref{A5}); $\mathrm{Div}$ is not $\De$, $\mathcal Y$ is not the
tessellation $\Xbulk$, $W_A$ and $W_{N,p}$ are not Definition~\ref{def:rhon}'s $W$, and $\mu_N$ is not Proposition~\ref{prop:mu}'s
$\mu_{\bullet}$ (v1.14.0-pre.2, WP-5 referee F6).

\subsection{The coset table for every \texorpdfstring{$N$}{N}}\label{sec:cosettableN}

\begin{definition}[The canonical table of $\Ga_0(N)\backslash\Ga$ for every $N$]\label{def:cosetrepsN}
Let $N\ge1$ with prime factorisation $N=\prod_{i=1}^{r}q_i^{\nu_i}$, $q_1<\dots<q_r$, $Q_i:=q_i^{\nu_i}$,
$m_i:=Q_i+Q_i/q_i$, $m_N:=\prod_im_i$ ($r=0$, $m_1=1$ at $N=1$).
\begin{enumerate}
\item \textbf{Component index.} For $(c,d)\in\Z^{2}$ primitive modulo $q_i$: $j_i(c,d):=1+(cd^{-1}\bmod Q_i)$ if
  $q_i\nmid d$, and $j_i(c,d):=Q_i+1+(dc^{-1}\bmod Q_i)/q_i$ otherwise --- the index of Definition~\ref{def:cosetreps} at
  $Q_i$. The \emph{component pair} of $j_i$ is $(j_i-1,1)$ for $j_i\le Q_i$ and $(1,q_i(j_i-Q_i-1))$ otherwise.
\item \textbf{Index} (mixed radix, primes ascending): for $(c,d)$ primitive modulo $N$,
  $j(c,d):=1+\sum_{i=1}^{r}\bigl(j_i(c,d)-1\bigr)\prod_{i'<i}m_{i'}\in\{1,\dots,m_N\}$.
\item \textbf{Canonical pair} $(c_j,d_j)\in\{0,\dots,N-1\}^{2}$ of an index $j$: the Chinese-remainder combination of the
  component pairs of its digits $j_i$. At $N=1$: $j=1$, $(c_1,d_1):=(0,1)$, $w_1=I$ (v1.14.0-pre.2, WP-5 referee F4).
\item \textbf{Representative} $w_j\in\Ga$, by the first case that applies: (i) $\gcd(d_j,N)=1$: $w_j:=\smat10k1$,
  $k:=c_jd_j^{-1}\bmod N$; (ii) $\gcd(c_j,N)=1$: $w_j:=\smat0{-1}1{k'}$, $k':=d_jc_j^{-1}\bmod N$; (iii) otherwise
  (a \emph{mixed} class, possible only if $r\ge2$): $t$ the least integer $\ge0$ with $\gcd(c_j,d_j+tN)=1$,
  $d':=d_j+tN$, $a:=d'^{-1}\bmod c_j\in[0,c_j)$, $b:=(ad'-1)/c_j$, $w_j:=\smat ab{c_j}{d'}$.
\item \textbf{Class map and action}: $\mathrm{cls}_N(\gamma):=j(c,d)$ for $\gamma$ with bottom row $(c,d)$; for integral
  $\delta$ with $(c_j,d_j)\delta$ primitive modulo $N$, $j\cdot\delta:=j\bigl((c_j,d_j)\delta\bigr)$ and
  $h_j(\delta):=w_j\,\delta\,w_{j\cdot\delta}^{-1}$ (for $\gamma\in\Ga$ the cocycle of \eqref{eq:cocycle}).
\item \textbf{Representatives of $\Ga_0(N')\backslash\Ga_0(N)$}, $N\mid N'$: the $w^{(N')}_j$ of the level-$N'$ table over
  the indices $j$ whose canonical pair has $N\mid c_j$, in increasing $j$. For $A=B=\Ga_0(N)$ and $g=\smat100p$ they are
  the $r_t$ of Proposition~\ref{prop:levelcorrespondence}, $B''=\Ga_0(Np)$.
\end{enumerate}
\end{definition}

\begin{lemma}[The table is well defined, and is Definition~\ref{def:cosetreps} at prime powers]\label{lem:cosetrepsN}
\PROVED{} With the notation of Definition~\ref{def:cosetrepsN}:
\begin{enumerate}
\item $j$ is constant on $(\Z/N)^{\times}$-classes and induces a bijection $\mathbb{P}^1(\Z/N)\to\{1,\dots,m_N\}$;
  $m_N=N\prod_{q\mid N}(1+q^{-1})$.
\item Every $w_j$ lies in $\Ga$ and has bottom row in the class $j$; in case (iii) $t<c_j$. $w_1=I$. The $w_j$ represent
  $\Ga_0(N)\backslash\Ga$.
\item At $N=p^{\nu}$ index, canonical pairs and representatives are those of Definition~\ref{def:cosetreps}.
\item $(j,\gamma)\mapsto j\cdot\gamma$ is a right action of $\Ga$, computed digit by digit (the digit $i$ of $j\cdot\gamma$
  is the level-$Q_i$ action on the digit $j_i$), and $h_j(\gamma)\in\Ga_0(N)$. For integral $\delta$ with $(c_j,d_j)\delta$
  primitive modulo $N$, $h_j(\delta)$ is integral, $\det h_j(\delta)=\det\delta$, and its lower left entry is $\equiv0\pmod N$.
  \LEAN{prop:heilbronn-model} (this last sentence, for any table of matrices of $\Ga$ with pairwise inequivalent bottom rows)
\item The set of item (6) represents $\Ga_0(N')\backslash\Ga_0(N)$, lies in $\Ga_0(N)$, and starts with $I$; it has
  $m_{N'}/m_N$ elements ($p+1$ for $N'=Np$, $p\nmid N$; $p$ for $p\mid N$). At $N=1$ it is Definition~\ref{def:cosetreps} at
  level $N'$.
\end{enumerate}
\end{lemma}

\begin{proof}
(1) Reduction modulo the $Q_i$ is a bijection $\mathbb{P}^1(\Z/N)\to\prod_i\mathbb{P}^1(\Z/Q_i)$: $(\Z/N)^{\times}=\prod_i
(\Z/Q_i)^{\times}$, and $(c,d)$ is primitive modulo $N$ iff it is modulo every $q_i$ (Chinese remainder theorem). Each $j_i$
is the bijection $\mathbb{P}^1(\Z/Q_i)\to\{1,\dots,m_i\}$ of Proposition~\ref{prop:cosets}(4) and
Definition~\ref{def:cosetreps}, and the mixed radix is a bijection $\prod_i\{1,\dots,m_i\}\to\{1,\dots,m_N\}$.

(2) The canonical pair reduces modulo $Q_i$ to the component pair, which is primitive and of class $j_i$; so $(c_j,d_j)$ is
primitive modulo $N$ and of class $j$. Case (i): the bottom row $(k,1)\equiv d_j^{-1}(c_j,d_j)$; case (ii): $(1,k')\equiv
c_j^{-1}(c_j,d_j)$; both matrices have determinant $1$. Case (iii): $c_j\ne0$ (else $d_j$ is a unit, case (i)) and $c_j\ne1$
(case (ii)), so $2\le c_j<N$. For a prime $s\mid c_j$: if $s\mid N$, then $s\nmid d_j$ by primitivity, so $s\nmid d_j+tN$
for every $t$; if $s\nmid N$, exactly one residue of $t$ modulo $s$ has $s\mid d_j+tN$. By the Chinese remainder theorem some
$t$ modulo $\prod_{s\mid c_j,\,s\nmid N}s\le c_j$ avoids all of them, so $t<c_j$ and $\gcd(c_j,d')=1$; then $ad'\equiv1\pmod{c_j}$,
$b\in\Z$, $\det w_j=ad'-bc_j=1$, and the bottom row $(c_j,d')\equiv(c_j,d_j)\pmod N$. $j=1$ has all digits $1$, component
pairs $(0,1)$, canonical pair $(0,1)$, case (i) with $k=0$: $w_1=I$. The last sentence is the bottom-row bijection
$\Ga_0(N)\backslash\Ga\to\mathbb{P}^1(\Z/N)$ (Definition~\ref{def:levelellmodel}'s parenthetical, for every $N$).

(3) $r=1$: $j=j_1$ is Definition~\ref{def:cosetreps}'s index, the canonical pair is the component pair; $(k,1)$ has $d=1$ a
unit and case (i) gives $\smat10k1=s_{1+k}$; $(1,pc)$ has $c=1$ a unit and $d=pc$ not, and case (ii) gives
$\smat0{-1}1{pc}=s_{N+1+c}$.

(4) The class of $(c_j,d_j)\gamma_1$ is $j\cdot\gamma_1$, so $(c_{j\cdot\gamma_1},d_{j\cdot\gamma_1})\equiv u\,(c_j,d_j)\gamma_1$
with $u$ a unit, and $(j\cdot\gamma_1)\cdot\gamma_2$ is the class of $u(c_j,d_j)\gamma_1\gamma_2$, i.e.\ $j\cdot(\gamma_1\gamma_2)$;
$j\cdot I=j$. The class modulo $Q_i$ of $(c,d)\gamma$ depends only on that of $(c,d)$, whence the digitwise computation.
$h_j(\gamma)\in\Ga_0(N)$ by the argument of Lemma~\ref{lem:cocycle}. For $\delta$: $w_j,w_{j\cdot\delta}\in\Ga$, so
$h_j(\delta)$ is integral of determinant $\det\delta$; its bottom row is $\mathrm{row}_2(w_j)\,\delta\,w_{j\cdot\delta}^{-1}$, and
$\mathrm{row}_2(w_j)\delta\equiv u'\,\mathrm{row}_2(w_{j\cdot\delta})\pmod N$ for a unit $u'$ (both primitive, same class), while
$\mathrm{row}_2(w_{j\cdot\delta})\,w_{j\cdot\delta}^{-1}=(0,1)$.

(5) $\Ga_0(N')\gamma\subseteq\Ga_0(N)$ iff $N$ divides the lower left entry of $\gamma$, a condition on the class of the bottom
row modulo $N'$ ($N\mid N'$); so the classes with $N\mid c_j$ are the image of $\Ga_0(N')\backslash\Ga_0(N)$ in
$\Ga_0(N')\backslash\Ga$. Their lifts lie in $\Ga_0(N)$: at a prime $q_i\mid N$, $q_i^{\nu_i}\mid c$ forces $q_i\nmid d$, so the
component pair modulo $Q'_i$ is $(k,1)$ with $k\equiv cd^{-1}\equiv0\pmod{Q_i}$; hence $c_j\equiv0\pmod N$, and the lift keeps
the lower left entry $\equiv c_j$ in case (i) ($k=c_jd_j^{-1}$), exactly in case (iii), and case (ii) needs $c_j$ a unit
modulo $N'$, so $N=1$. The class of $I$ has index $1$. The count is $[\Ga_0(N):\Ga_0(N')]=m_{N'}/m_N$ by (1) and
Proposition~\ref{prop:cosets}(4).
\end{proof}

\begin{contractbox}
Definition~\ref{def:cosetrepsN} is \textbf{binding as written} (Remark~\ref{rmk:indeplevel}: the cochain-level matrices of
\eqref{eq:levelcorr} depend on it, $H^{\bullet}$ does not), and it closes obstacle~\ref{levelO3}. It is \textbf{a table, not
a search}: cases (i), (ii) are closed forms, case (iii) the least $t<c_j$ with $\gcd(c_j,d_j+tN)=1$ (a bounded loop, at most $c_j$
gcd's, in practice a few; table-building time, never per operator) and one inverse modulo $c_j$;
the class map and the action are digitwise, $O(\omega(N))$ operations modulo the $Q_i$. \textbf{No new $\tau$-table} is
needed for any $N$: every level-$N$ operator of Proposition~\ref{prop:levelcorrespondence} reads the level-$1$ table of
$\Tdc{g}$ (its item (1)), and the Hecke action of Proposition~\ref{prop:heilbronn} reads no $\tau$-table at all.
Any other transversal of $B''\backslash B$ gives the same operator on $H^{\bullet}$ (Proposition~\ref{prop:levelcorrespondence}(3),
$6\in\cring^{\times}$) and, in general, different cochain-level matrices; slot~13's \code{Gamma0SubTransversal} ($\smat10{Nt}1$,
$t<p$, and $\smat xyNp$ for $p\nmid N$) is such an alternative, used by \code{HeckeHidaLevelEllValidation} at $N=\ell$ and
compared with item~(6)'s on $H^{1}$ by \code{HeckeLevelNValidation} (arms (H), (E)); it is not the binding datum, and no shipped
record, report or golden text reads either.
\end{contractbox}

\subsection{The level-\texorpdfstring{$N$}{N} model and the cokernel presentation}\label{sec:levelNpres}

\begin{definition}[The level-$N$ model for every $N$; its coefficient settings]\label{def:levelNmodel}
Let $N\ge1$ and $\cring$ a commutative ring. $M_N:=\Ind_{\Ga_0(N)}^{\Ga}\VR$ in the coordinates of
Definition~\ref{def:inducedmodule} with the $w_j$ of Definition~\ref{def:cosetrepsN}, and $L_N:=L_{\Ga_0(N)}$
(Definition~\ref{def:inducedls}): Definition~\ref{def:levelellmodel} with the table now fixed for every $N$. Its relation
matrices are the monomial $M_N(\gamma)=\sum_jE_{j,\,j\cdot\gamma}\otimes\rhon(h_j(\gamma))$ (\eqref{eq:MH}); the orbit
structure is Proposition~\ref{prop:levelorbit}, stated there for every $N$, with $\nu_2(N)$, $\nu_3(N)$ of
Remark~\ref{rmk:levelcost}(3) multiplicative over the $Q_i$ and the $S$-, $U$-orbits on $\mathbb{P}^1(\Z/N)$ computed
digitwise. The \textbf{coefficient settings}:
\begin{enumerate}
\item[(a)] every $\cring$ at cochain level (Proposition~\ref{prop:levelcorrespondence}(1),(2));
\item[(b)] $6\in\cring^{\times}$ for the comparison with $H^{\bullet}(\Ga_0(N);\VR)$ --- $\cring=\Fl$, $\Z/\ell^{k}$
  ($\ell\ge5$), $\Q$, $\Z[1/6]$ --- in either relative position: $\ell\mid N$ (the Hida setting; $N=\ell$ is slot~13's) or
  $\ell\nmid N$;
\item[(c)] the \textbf{R3 setting} (route R3 of \file{docs/implementation/slot14/WP-M1\_loewy\_note.md} \S6): $\cring=\F_p$
  (or $\Z/p^{k}$), $p\ge5$, $n=0$ (weight $2$), $N$ prime with $p\mid\mathrm{num}\bigl(\frac{N-1}{12}\bigr)$, so $p\nmid N$ ---
  Mazur's Eisenstein ideal (Remark~\ref{rmk:levelNcheckable}, \ref{LN4}--\ref{LN6}). \emph{Since v1.21.0-pre.4}, at $n=0$ only, also
  $p\in\{2,3\}$ on the \textbf{curve object}: $\cring=\F_p$ (or $\Z$), every prime $p$, the fixed spaces $M_N^{\langle S\rangle}$,
  $M_N^{\langle U\rangle}$ taken as the spans of the orbit sums; the cokernel is then $H^{1}(Y_0(N);\F_p)$, the cohomology of the open
  coarse curve, which at $p\in\{2,3\}$ is \textbf{not} $H^{1}(\Ga_0(N);\F_p)$ (Proposition~\ref{prop:levelNcurve},
  Remark~\ref{rmk:levelNcurvefamily}).
\end{enumerate}
The level-$N$ operators are those of Proposition~\ref{prop:levelcorrespondence} with $A=B=\Ga_0(N)$, the $r_t$ of
Definition~\ref{def:cosetrepsN}(6), and the $s_j$ and $\tau$ of Definition~\ref{def:cosetreps} for $\Tdc{g}$ at level~$1$.
\end{definition}

\begin{lemma}[The tree is a projective resolution]\label{lem:treeres}
\PROVED{} Let $6\in\cring^{\times}$ and $A_0\le\Ga$ any subgroup. Then $C_{\bullet}(\mathcal Y)\otimes\cring\to\cring$ is a
projective resolution of length $1$ of $\cring$ over $\cring[A_0]$. Hence $H^{q}(A_0;X)=H^{q}\bigl(\Hom_{A_0}(C_{\bullet}(\mathcal
Y),X)\bigr)$ and $H_q(A_0;X)=H_q\bigl((C_{\bullet}(\mathcal Y)\otimes X)_{A_0}\bigr)$ for every $\cring[A_0]$-module $X$, both
zero for $q\ge2$.
\end{lemma}

\begin{proof}
\eqref{eq:treeses} is an exact sequence of free abelian groups, hence split over $\Z$, so it stays exact after
$\otimes\cring$. As an $\cring[A_0]$-module, $\cring[\Ga/F]$ ($F\in\{\langle-I\rangle,\langle S\rangle,\langle U\rangle\}$, of
orders $2,4,6$) is $\bigoplus\cring[A_0/H']$ over the double cosets $A_0\gamma F$, $H'=A_0\cap\gamma F\gamma^{-1}$ of order
dividing $4$ or $6$, a unit of $\cring$; $e:=\abs{H'}^{-1}\sum_{h\in H'}h$ is an idempotent with $\cring[A_0]e\cong\cring[A_0/H']$,
a direct summand of $\cring[A_0]$.
\end{proof}

\begin{lemma}[Manin's exact sequence]\label{lem:manin}
\PROVED{} \STD{\cite{manin} for the statement (Manin symbols and their relations); pointer \textbf{not} verified against the
text (rule~\ref{A2}); the proof below is complete} Define $\Ga$-maps
\[
  \partial^{t}\colon C_0(\mathcal Y)\to C_1(\mathcal Y),\quad
  \gamma\mathfrak v_S\mapsto\gamma\mathfrak e+\gamma S\mathfrak e,\quad
  \gamma\mathfrak v_U\mapsto\gamma\mathfrak e+\gamma U\mathfrak e+\gamma U^{2}\mathfrak e;
  \qquad
  \mathsf b\colon C_1(\mathcal Y)\to\mathrm{Div},\quad\gamma\mathfrak e\mapsto\gamma\cdot D_0=[\gamma\infty]-[\gamma0].
\]
Then $0\to C_0(\mathcal Y)\xrightarrow{\partial^{t}}C_1(\mathcal Y)\xrightarrow{\mathsf b}\mathrm{Div}^0\to0$ is exact.
\end{lemma}

\begin{proof}
The maps are well defined: $S^{2}=U^{3}=-I$ fixes $\mathfrak e$ and every cusp, so $\gamma S\mathfrak v_S=\gamma\mathfrak v_S$
and $\gamma U\mathfrak v_U=\gamma\mathfrak v_U$ give the same sums. $\mathsf b\,\partial^{t}=0$: $S$ exchanges $[\infty]$ and
$[0]$, and $U$ permutes $[\infty]\mapsto[0]\mapsto[-1]\mapsto[\infty]$, so $\mathsf b(\mathfrak e+S\mathfrak e)=0$ and
$\mathsf b(\mathfrak e+U\mathfrak e+U^{2}\mathfrak e)=([\infty]-[0])+([0]-[-1])+([-1]-[\infty])=0$.

\emph{$\partial^{t}$ is injective.} The edge $\gamma\mathfrak e$ has exactly the two end points $\gamma\mathfrak v_S$,
$\gamma\mathfrak v_U$; if $\partial^{t}(\sum_va_vv)=0$ then $a_v+a_w=0$ for every edge $\{v,w\}$, so $\abs{a_v}$ is constant on
the tree, which is connected (exactness of \eqref{eq:treeses} at $\Z$) and infinite; a finite sum forces $a=0$.

\emph{The Farey complex.} Let $\mathsf F$ be the simplicial complex with vertex set $\mathbb{P}^1(\Q)$, edges
$\{\gamma\infty,\gamma0\}$ and triangles $\{\gamma\infty,\gamma0,\gamma(-1)\}$, $\gamma\in\Ga$. (F1) For coprime pairs,
$\{(a:c),(b:d)\}$ is an edge iff $ad-bc=\pm1$ ($\smat abcd$ or $\smat a{-b}c{-d}$ lies in $\Ga$). (F2) The triangles are
exactly the triples of pairwise adjacent cusps: the sides of $\{\gamma\infty,\gamma0,\gamma(-1)\}$ are $\gamma\{\infty,0\}$,
$\gamma U\{\infty,0\}$, $\gamma U^{2}\{\infty,0\}$; conversely, move two of the three to $\infty,0$ by (F1), and the third,
adjacent to both, is $\pm1$, with $\{\infty,0,1\}=S\{\infty,0,-1\}$. \emph{Height}: $\mathrm{ht}(a:c):=\abs c$ for coprime
$(a,c)$; $\mathsf F_h$ is the full subcomplex on the cusps of height $\le h$. (F3) A cusp $x=(a:c)$ with $c\ge2$ has
exactly two neighbours of smaller height, $y=(b_1:d_1)$ with $ad_1-b_1c=1$, $d_1\in[1,c)$ ($d_1\equiv a^{-1}\bmod c$), and
$z=(a-b_1:c-d_1)$; $\{y,z\}$ is an edge ($b_1(c-d_1)-(a-b_1)d_1=-1$); no neighbour has height $c$ ($c\mid ad-bc$ there); and
by (F2) the only triangle of $\mathsf F_c$ containing $x$ is $\{x,y,z\}$. \emph{Elementary move}: if $L'=L\cup\{x,\{x,y\},
\{x,z\},\{x,y,z\}\}$ with $x\notin L$, $\{y,z\}\in L$, the relative chain complex $C_{\bullet}(L',L)$ is $\Z\{x,y,z\}\to
\Z\{x,y\}\oplus\Z\{x,z\}\to\Z x$ with the maps $\{x,y,z\}\mapsto\{x,z\}-\{x,y\}$ up to orientation signs and
$\{x,y\},\{x,z\}\mapsto\pm x$, which is acyclic; so $L\subseteq L'$ is a homology isomorphism. $\mathsf F_1$ (cusps $\infty$
and $\Z$, triangles $\{\infty,k,k+1\}=T^{k+1}\{\infty,0,-1\}$) is the increasing union of the fans on $\{\infty,-k,\dots,k\}$,
each obtained from the edge $\{\infty,0\}$ by elementary moves; $\mathsf F_h$ is obtained from $\mathsf F_{h-1}$ by one
elementary move per cusp of height $h\ge2$ (two cusps of height $h$ are not adjacent), and $\mathsf F=\bigcup_h\mathsf F_h$.
Homology commutes with increasing unions of subcomplexes, so $\tilde H_{\bullet}(\mathsf F;\Z)=0$.

\emph{Conclusion.} Let $q\colon C_1(\mathcal Y)\to C_1(\mathsf F)$ (oriented $1$-chains, $[x,y]=-[y,x]$),
$q(\gamma\mathfrak e):=[\gamma0,\gamma\infty]$: well defined, onto by (F1), with $\mathsf b=\partial^{\mathsf F}_1\circ q$.
$q(\gamma\mathfrak e)=\pm q(\gamma'\mathfrak e)$ iff $\gamma'\in\pm\gamma$ (sign $+$) or $\gamma'\in\pm\gamma S$ (sign $-$), the
stabiliser of an ordered Farey pair in $\Ga$ being $\{\pm I\}$ (it fixes $(1:0)$ and $(0:1)$); so $\ker q=\partial^{t}\Z[\Ga/\langle
S\rangle]$. And $q\,\partial^{t}(\gamma\mathfrak v_U)=[\gamma0,\gamma\infty]+[\gamma(-1),\gamma0]+[\gamma\infty,\gamma(-1)]=
\partial[\gamma(-1),\gamma0,\gamma\infty]$, so $q\,\partial^{t}\Z[\Ga/\langle U\rangle]=\partial C_2(\mathsf F)$. Then
$\mathrm{im}\,\mathsf b=\mathrm{im}\,\partial^{\mathsf F}_1=\mathrm{Div}^0$ ($\tilde H_0=0$), and $\ker\mathsf b=q^{-1}(\partial
C_2(\mathsf F))=\partial^{t}\Z[\Ga/\langle U\rangle]+\ker q=\mathrm{im}\,\partial^{t}$ ($H_1=0$).
\end{proof}

\begin{definition}[The pinned cokernel isomorphism $\varkappa_N$]\label{def:kappaN}
\PROVED{} (the claims; proof below) Let $6\in\cring^{\times}$, $n$ even, $M:=M_N$ and $\mathrm{Rel}_N:=M^{\langle S\rangle}+M^{\langle U\rangle}=(1+M_N(S))M+
(1+M_N(U)+M_N(U)^{2})M$ (Proposition~\ref{prop:levelorbit}(2)). For $x\in M$ let $\psi_x\in\Hom_{\Ga_0(N)}(C_1(\mathcal
Y),\VR)$, $\psi_x(\gamma\mathfrak e):=\bigl(M_N(\gamma)x\bigr)_1=\rhon(h_1(\gamma))\,x_{1\cdot\gamma}$, and let $\varkappa_N(x)\in
H^{1}(\Ga_0(N);\VR)$ be its class (Lemma~\ref{lem:treeres}). Then $\varkappa_N$ induces an isomorphism
$M/\mathrm{Rel}_N\to H^{1}(\Ga_0(N);\VR)$ --- \textbf{the} isomorphism of Proposition~\ref{prop:levelorbit}(2), which that
proposition did not pin. In group cocycles: the inhomogeneous cocycle of $\psi_x$ transported to $\Ga$ is
$C_x(\gamma):=\phi_x(\mathrm{path}(\mathfrak v_U\to\gamma\mathfrak v_U))$, $\phi_x(\gamma\mathfrak e):=M_N(\gamma)x$, so
$C_x(U)=0$ and $C_x(S)=(M_N(S)-1)x$.
\end{definition}

\noindent
\emph{Proof of the claims.} The reciprocity $\phi\mapsto\phi(\cdot)_1$ of Proposition~\ref{prop:levelcorrespondence}(3) is a
cochain isomorphism $\Hom_{\Ga}(C_{\bullet}(\mathcal Y),M)\to\Hom_{\Ga_0(N)}(C_{\bullet}(\mathcal Y),\VR)$ (its proof uses only
that the domain is a $\Z[\Ga]$-module). $\Hom_{\Ga}(C_1,M)=M^{-I}=M$ by $\phi\mapsto\phi(\mathfrak e)$ ($M_N(-I)=1$, $n$ even),
$\Hom_\Ga(C_0,M)=M^{\langle S\rangle}\oplus M^{\langle U\rangle}$, and the coboundary is $(a,b)\mapsto b-a$; its image is
$\mathrm{Rel}_N$, the equality $M^{\langle S\rangle}=(1+M_N(S))M$ holding because $\tfrac12(1+M_N(S))$ is the averaging idempotent
($M_N(S)^{2}=M_N(-I)=1$), likewise for $U$. The cocycle: $f_0(\gamma):=\gamma\mathfrak v_U$, $f_1(\gamma_0,\gamma_1):=
\mathrm{path}(\gamma_0\mathfrak v_U\to\gamma_1\mathfrak v_U)$ is a $\Ga$-chain map from the standard resolution to
$C_{\bullet}(\mathcal Y)$ over the identity (the boundary of a path is its end minus its start), so $f^{\ast}$ is the comparison
isomorphism, and the inhomogeneous cocycle is $\gamma\mapsto\phi(f_1(1,\gamma))$; $\mathrm{path}(\mathfrak v_U\to S\mathfrak
v_U)=-\mathfrak e+S\mathfrak e$ and $U\mathfrak v_U=\mathfrak v_U$. \qed

\begin{lemma}[Cokernel duality, and its naturality]\label{lem:levelduality}
\PROVED{} Let $6\in\cring^{\times}$, $n$ even, and let $A\le\Ga$ be of finite index with $-I\in A$. Put
$W_A:=H_0(A;\mathrm{Div}^0_{\cring}\otimes\VR)=(\mathrm{Div}^0_{\cring}\otimes\VR)_A$, $A$ acting diagonally ($\gamma\cdot
(D\otimes v)=\gamma D\otimes\rhon(\gamma)v$), and for $\psi\in\Hom_A(C_1(\mathcal Y),\VR)$
\[
  \Pi_A(\psi):=\sum_{\xi\in A\backslash\mathcal Y_1}\bigl[\mathsf b(\xi)\otimes\psi(\xi)\bigr]_A ,
\]
the sum over any set of representatives of the $A$-orbits of edges (finitely many; the summand depends only on the orbit,
the stabiliser $\{\pm I\}$ of an edge acting trivially on $\mathrm{Div}$ and on $\VR$).
\begin{enumerate}
\item $\Pi_A$ vanishes on coboundaries and induces $\Pi_A\colon H^{1}(A;\VR)\to W_A$.
\item \textbf{(Coordinates.)} For $A=\Ga_0(N)$: $\Pi_A(\varkappa_N(x))=\mu_N(x):=\sum_{j}[w_jD_0\otimes x_j]_A$; the
  \emph{Manin map} $\mu_N\colon M_N\to W_A$ is onto with kernel $\mathrm{Rel}_N$; hence $\Pi_A$ is an isomorphism.
\item \textbf{(Restriction.)} For $A'\le A$ of finite index with $-I\in A'$: $\Pi_{A'}\circ\res^{A}_{A'}=\mathrm{tr}\circ\Pi_A$,
  $\mathrm{tr}[y]_A:=\sum_{r\in A'\backslash A}[r\cdot y]_{A'}$.
\item \textbf{(Corestriction.)} Likewise $\Pi_A\circ\cor^{A}_{A'}=\mathrm{pr}\circ\Pi_{A'}$, $\mathrm{pr}\colon W_{A'}\to W_A$ the
  projection.
\item \textbf{(Conjugation.)} For $g\in\De$ and $A'\le\Ga$ with $A'':=gA'g^{-1}\le\Ga$ (both of finite index, containing
  $-I$): $\Pi_{A''}\circ(\conj_g)_{\ast}=\conj_g\circ\Pi_{A'}$, where $(\conj_g)_{\ast}\colon H^{1}(A';\VR)\to H^{1}(A'';\VR)$ is
  induced by the compatible pair $(\conj_g^{-1},\rhon(g))$ of \eqref{eq:corconjres} and $\conj_g[y]_{A'}:=[g\cdot y]_{A''}$,
  $g$ acting on $\mathrm{Div}^0$ by the column action and on $\VR$ by $\rhon(g)$.
\item \textbf{(Hecke operators.)} Consequently, for $g\in\De$ with $A'':=A\cap gAg^{-1}$ of finite index,
  $\Pi_A\circ[A\,g\,A]=[A\,g\,A]_{\ast}\circ\Pi_A$, where $[A\,g\,A]$ is the operator of \S\ref{sec:hidafull}'s notation and
  $[A\,g\,A]_{\ast}[y]_A:=\sum_{A\beta\in A\backslash AgA}[\beta\cdot y]_A$.
\end{enumerate}
\end{lemma}

\begin{proof}
\emph{Star sums.} For an $A$-invariant function $F$ on edges with values in an abelian group in which $6$ is invertible,
$\sum_{\xi\in A\backslash\mathcal Y_1}F(\xi)=\sum_{v}k_v^{-1}\sum_{\xi:\,\mathrm{head}(\xi)=v}F(\xi)$, the outer sum over
representatives $v$ of the $A$-orbits of $U$-vertices, $\mathrm{head}(\gamma\mathfrak e):=\gamma\mathfrak v_U$, and
$k_v:=\abs{\mathrm{Stab}_A(v)/\{\pm I\}}\in\{1,3\}$: $\mathrm{Stab}_A(v)/\{\pm I\}$ acts freely on the three edges at $v$, so
each $A$-orbit of edges meets the star of exactly one representative $v$, in $k_v$ edges. The same holds with tails and
$S$-vertices ($k_v\in\{1,2\}$).

(1) For $\psi=\delta f$, $f\in\Hom_A(C_0(\mathcal Y),\VR)$, $\psi(\xi)=f(\mathrm{head}\,\xi)-f(\mathrm{tail}\,\xi)$;
the star sum turns $\sum_\xi[\mathsf b(\xi)\otimes f(\mathrm{head}\,\xi)]$ into
$\sum_vk_v^{-1}[\mathsf b(\partial^{t}v)\otimes f(v)]=0$ by Lemma~\ref{lem:manin}, and likewise for tails.

(2) $A\backslash\mathcal Y_1\cong A\backslash\Ga$ ($\pm I\subseteq A$ central), with representatives $w_j\mathfrak e$, and
$\psi_x(w_j\mathfrak e)=\rhon(h_1(w_j))x_{1\cdot w_j}=x_j$ since $1\cdot w_j=j$ and $h_1(w_j)=I$; so $\Pi_A(\varkappa_N(x))=
\mu_N(x)$. Tensor the sequence of Lemma~\ref{lem:manin} (split over $\Z$) with $\VR$ and take $A$-coinvariants (right
exact): $(C_0\otimes\VR)_A\to(C_1\otimes\VR)_A\to W_A\to0$ is exact. $x\mapsto\sum_j[w_j\mathfrak e\otimes x_j]$ identifies
$M$ with $(C_1\otimes\VR)_A$ (each $A$-orbit of edges once), with inverse $[hw_i\mathfrak e\otimes v]\mapsto$ the vector with
block $i$ equal to $\rhon(h)^{-1}v$ ($h\in A$). Under it, $[w_i\mathfrak v_S\otimes v]$ maps to $e_i\otimes v+e_{i\cdot S}\otimes
\rhon(h_i(S))^{-1}v=(1+M_N(S)^{-1})(e_i\otimes v)$, and $M_N(S)^{-1}=M_N(S)M_N(-I)=M_N(S)$; the $U$-vertices give
$(1+M_N(U)+M_N(U)^{2})(e_i\otimes v)$ likewise ($U^{-1}=U^{2}\cdot(-I)$). So the image of $(C_0\otimes\VR)_A$ is
$\mathrm{Rel}_N$, and the composite $M\to W_A$ is $\mu_N$. With Definition~\ref{def:kappaN}, $\Pi_A$ is an isomorphism.

(3) $\{r\xi\}$, $r\in A'\backslash A$, $\xi\in A\backslash\mathcal Y_1$, represents $A'\backslash\mathcal Y_1$ (edge stabilisers are
$\{\pm I\}\subseteq A'$), and $[r\cdot(\mathsf b(\xi)\otimes\psi(\xi))]_{A'}=[\mathsf b(r\xi)\otimes\psi(r\xi)]_{A'}$.

(4) The corestriction on the resolution of Lemma~\ref{lem:treeres} is the transfer $\psi'\mapsto\sum_r\rhon(r^{-1})\psi'(r\,\cdot)$
(as in the proof of Proposition~\ref{prop:levelcorrespondence}(3)); then
$[\mathsf b(\xi)\otimes\rhon(r^{-1})\psi'(r\xi)]_A=[r^{-1}\cdot(\mathsf b(r\xi)\otimes\psi'(r\xi))]_A$, and
summing over $\xi$, $r$ gives $\mathrm{pr}\,\Pi_{A'}(\psi')$ by (3)'s bijection.

(5) \emph{A pairing.} For $A$ as in the statement and any $\cring[A]$-module $X$, put $\psi\frown z:=\sum_i
[\mathsf x_i\otimes\psi(y_i)]_A\in(X\otimes\VR)_A$ for $\psi\in\Hom_A(P_1,\VR)$ and a $1$-cycle $z=[\sum_iy_i\otimes\mathsf x_i]$ of
$(P_1\otimes X)_A$, $P$ any projective resolution of length $1$ (such as Lemma~\ref{lem:treeres}'s). It is well defined on
$H^{1}\times H_1$: on $\delta f$ it is the image of $(\partial\otimes1)z=0$ under $[y\otimes\mathsf x]\mapsto[\mathsf x\otimes f(y)]$,
and there are no $1$-boundaries. For a chain map $F\colon P\to P'$ of such resolutions over $\id_{\cring}$,
$(F^{\ast}\psi)\frown z=\psi\frown(F_{\ast}z)$ at chain level, $F^{\ast}$, $F_{\ast}$ being the comparison isomorphisms: so $\frown$ is a
pairing $H^{1}(A;\VR)\times H_1(A;X)\to(X\otimes\VR)_A$ independent of $P$. By definition $\Pi_A(\psi)=\psi\frown\zeta_A$ with
$\zeta_A:=\sum_{\xi\in A\backslash\mathcal Y_1}[\xi\otimes\mathsf b(\xi)]$, a cycle of $(C_1(\mathcal Y)\otimes\mathrm{Div}^0_{\cring})_A$ by the star
sums and Lemma~\ref{lem:manin}.
\emph{Transport.} Let $P^{g}$ be $C_{\bullet}(\mathcal Y)\otimes\cring$ with $h''\in A''$ acting as $g^{-1}h''g\in A'$: a length-$1$
projective resolution over $\cring[A'']$. The maps of the compatible pairs $(\conj_g^{-1},\rhon(g))$ on $H^{1}$ and
$(\conj_g,g)$ on $H_1(-;\mathrm{Div}^0)$ and $H_0(-;\mathrm{Div}^0\otimes\VR)$ are computed on $P$ and $P^{g}$ by
$\psi\mapsto\rhon(g)\circ\psi$ and $y\otimes\mathsf x\mapsto y\otimes g\mathsf x$ \STD{\cite[Ch.~III.8]{brown}, the map of a compatible
pair by a chain map over it --- here the identity of the underlying modules; \textbf{not} verified (rule~\ref{A2}); the dependence is
definitional (the meaning of $(\conj_g)_{\ast}$ in \eqref{eq:corconjres}), no theorem of the source is used; v1.14.0-pre.2, WP-5 referee F7}, and at chain
level $(\rhon(g)\psi)\frown(y\otimes g\mathsf x)=\conj_g\bigl(\psi\frown(y\otimes\mathsf x)\bigr)$. So (5) follows from
$(\conj_g)_{\ast}[\zeta_{A'}]=[\zeta_{A''}]$ in $H_1(A'';\mathrm{Div}^0_{\cring})$.
\emph{The class $[\zeta_A]$.} $H_2(A;\cring)=0$ (Lemma~\ref{lem:treeres}), so $0\to\mathrm{Div}^0_{\cring}\to\mathrm{Div}_{\cring}\to
\cring\to0$ makes $H_1(A;\mathrm{Div}^0_{\cring})\to H_1(A;\mathrm{Div}_{\cring})$ injective, and $\mathrm{Div}_{\cring}=\bigoplus_{c\in
A\backslash\mathbb{P}^1(\Q)}\cring[A/A_c]$, $A_c$ the stabiliser, gives $H_1(A;\mathrm{Div}_{\cring})\cong\bigoplus_cH_1(A_c;\cring)$ by
$[x\otimes aA_c]\mapsto[a^{-1}x]$ on $(C_1(\mathcal Y)\otimes\cring[A/A_c])_A\cong(C_1(\mathcal Y)\otimes\cring)_{A_c}$ (and $P$ is
a resolution over $\cring[A_c]$, Lemma~\ref{lem:treeres}). For $c=\sigma\infty$, $\sigma\in\Ga$: $A_c=\pm\langle\gamma^{+}_c\rangle$
with $\gamma^{+}_c:=\sigma T^{h_c}\sigma^{-1}$, $h_c\ge1$ (independent of $\sigma$). \emph{Translation classes}: for a tree
$\mathcal X$ with an action of $A_c$ making $C_{\bullet}(\mathcal X)\otimes\cring$ a resolution as above and a vertex $x$,
$[\mathrm{path}(x\to\gamma^{+}_cx)]\in H_1(A_c;\cring)$ is a cycle independent of $x$ ($\mathrm{path}(x\to\gamma x)-
\mathrm{path}(y\to\gamma y)=\mathrm{path}(x\to y)-\gamma\,\mathrm{path}(x\to y)$) and of $\mathcal X$: for a chain map $F$ over
$\id$ to another such tree, $F_1\mathrm{path}(x\to\gamma x)$ and $\sum_ia_i\mathrm{path}(y_i\to\gamma y_i)$ ($F_0x=\sum a_iy_i$,
$\sum a_i=1$) have the same boundary, hence are equal ($\partial$ injective). Call it $[\gamma^{+}_c]$. The $c$-component of
$\zeta_A$: at $c=\infty$ the edges with $\mathsf b$-coefficient at $[\infty]$ are $\pm T^{k}\mathfrak e$ (head) and $\pm
T^{k}S\mathfrak e$ (tail), and the component is $\sum_{k<h_\infty}(T^{k}\mathfrak e-T^{k}S\mathfrak e)$, whose boundary is
$\sum_{k<h_\infty}(T^{k}\mathfrak v_U-T^{k+1}\mathfrak v_U)=\mathfrak v_U-T^{h_\infty}\mathfrak v_U$ because $S\mathfrak v_U=T\mathfrak v_U$
($T^{-1}S=U^{2}$); so it is $-\mathrm{path}(\mathfrak v_U\to\gamma^{+}_\infty\mathfrak v_U)$, of class $-[\gamma^{+}_\infty]$;
at $c=\sigma\infty$ translate by $\sigma$. So $[\zeta_A]$ is \textbf{the unique class whose component at every cusp $c$ is
$-[\gamma^{+}_c]$}. Finally $\conj_g$ maps $A'_c$ onto $A''_{gc}$ and $\gamma^{+}_c$ to $\gamma^{+}_{gc}$: writing
$g\sigma=\tau\smat ab0d$ ($\tau\in\Ga$, $a,d>0$, Hermite), $g\gamma^{+}_cg^{-1}=\tau\smat1{ah_c/d}01\tau^{-1}\in A''$ with
$ah_c/d>0$ a positive translation conjugate, which generates $A''_{gc}$ modulo $\pm I$; transported to $P^{g}$ the translation
class of $\gamma^{+}_c$ is that of $\gamma^{+}_{gc}$. Hence both classes have the same image, and are equal.

(6) $[A\,g\,A]=\cor^{A}_{A''}\circ(\conj_g)_{\ast}\circ\res^{A}_{A'}$, $A':=A\cap g^{-1}Ag$, $A''=gA'g^{-1}$ (both contain $-I$), so
by (3)--(5) $\Pi_A[A\,g\,A]\psi=\mathrm{pr}\,\conj_g\,\mathrm{tr}\,\Pi_A\psi$, i.e.\ $[y]_A\mapsto\sum_{r\in A'\backslash A}[gr\cdot y]_A$, and the
$gr$ represent $A\backslash AgA$ ($Agr=Agr'$ iff $rr'^{-1}\in A\cap g^{-1}Ag$).
\end{proof}

\subsection{The Hecke action on the presentation: Merel's set}\label{sec:heilbronnsec}

\begin{lemma}[Merel's set and its path condition]\label{lem:merelpath}
\PROVED{} \STD{\cite[\S2]{merel} for the set and the condition; pointer \textbf{not} verified (rule~\ref{A2}); the proof below
is complete} Let $n\ge1$ and
\[
  \mathrm{Mer}_n:=\Bigl\{\smat abcd\in\Mat_2(\Z):\ ad-bc=n,\ a>b\ge0,\ d>c\ge0\Bigr\},
\]
finite ($n=ad-bc\ge a+d-1$). For every left coset $\eta\Ga$ of determinant-$n$ matrices,
\[
  \sum_{\epsilon\in\mathrm{Mer}_n\cap\eta\Ga}\epsilon\cdot D_0\;=\;D_0\qquad\text{in }\mathrm{Div}^0 ,
\]
\LEAN{lem:merelpath-model} (the set, its finiteness and the path condition for every $\eta\Ga$; not the bound that follows) and $\abs{\mathrm{Mer}_n\cap\eta\Ga}\le x_0(\eta\Z^{2})$, $x_0(L):=\min\{x>0:(x,0)\in L\}$.
\end{lemma}

\begin{proof}
$\epsilon\Ga=\eta\Ga$ iff $\epsilon\Z^{2}=\eta\Z^{2}=:L$ (column lattices), so $B_L:=\mathrm{Mer}_n\cap\eta\Ga$ is the set of positively
oriented bases $(v_1\,|\,v_2)$ of $L$, $v_1=(a,c)^{\top}$, $v_2=(b,d)^{\top}$, with $a>b\ge0$, $d>c\ge0$; and
$\epsilon\cdot D_0=[v_1]-[v_2]$, $[v]$ the cusp of $v$ ("tail" $[v_1]$, "head" $[v_2]$). Both vectors lie in the closed first
quadrant, and $\det(v_1|v_2)>0$ makes the argument strictly increase from $v_1$ to $v_2$.
(i) \emph{The tail determines $\epsilon$}: $v_1$ is the primitive vector of $L$ on its ray; $v'_2-v_2\in L\cap\Q v_1=\Z v_1$, and
$0\le b,b'<a$ forces $v'_2=v_2$. Likewise the head, by $0\le c<d$.
(ii) \emph{Successor}: if $b>0$, then $(v_2\,|\,mv_2-v_1)\in B_L$ for $m:=\lceil a/b\rceil\ge2$: the determinant is $n$, the new
entries are $a'=b>b'=mb-a\in[0,b)$ and $c'=d<d'=md-c$ ($m\ge2$, $d>c$).
(iii) \emph{Predecessor}: if $c>0$, then $(mv_1-v_2\,|\,v_1)\in B_L$ for $m:=\lceil d/c\rceil\ge2$: $a''=ma-b\ge2a-b>a$,
$c''=mc-d\in[0,c)$, $d''=c$.
(iv) \emph{Start}: exactly one $\epsilon\in B_L$ has tail $[\infty]$, namely $v_1=(x_0,0)$, $v_2=(b,n/x_0)$ with $b\in[0,x_0)$
the unique such that $v_2\in L$ ($L/(L\cap\Z\times0)\cong y_1\Z$, $x_0y_1=[\Z^{2}:L]=n$). The head is $[0]$ iff $b=0$.
By (i) the successor of (ii) is the only element whose tail is the given head, and the predecessor of (iii) the only element
whose head is the given tail; arguments increase, $B_L$ is finite, so iterating the predecessor from any $\epsilon$ ends at the
start of (iv), and iterating the successor from the start ends at an element with head $[0]$. So $B_L$ is one chain
$\epsilon_0,\dots,\epsilon_r$ with $\mathrm{head}(\epsilon_k)=\mathrm{tail}(\epsilon_{k+1})$, from $[\infty]$ to $[0]$, and the sum
telescopes to $D_0$. The first coordinates of the tails strictly decrease from $x_0$ ($a>b$), so $r+1\le x_0$.
\end{proof}

\begin{proposition}[The Hecke action on the cokernel presentation: Merel's matrices in this document's conventions]\label{prop:heilbronn}
\PROVED{} Let $\cring$ be a commutative ring with $6\in\cring^{\times}$, $n\ge0$ even, $N\ge1$, $p$ a prime, $A:=\Ga_0(N)$,
$g:=\smat100p$. Define $H_{N,p}\colon M_N\to M_N$ by
\begin{equation}\label{eq:heilbronn}
  H_{N,p}\;:=\;\sum_{\epsilon\in\mathrm{Mer}_p}\ \sum_{j\,:\,(c_j,d_j)\epsilon\ \text{primitive mod }N}
  E_{j\cdot\epsilon,\;j}\otimes\rhon\bigl(\adjg h_j(\epsilon)\bigr),\qquad
  \adjg h_j(\epsilon)=w_{j\cdot\epsilon}\,(\adjg\epsilon)\,w_j^{-1},
\end{equation}
i.e.\ $(H_{N,p}x)_i=\sum\rhon(\adjg h_j(\epsilon))\,x_j$ over the pairs $(\epsilon,j)$ with $j\cdot\epsilon=i$. For $p\nmid N$ every $j$
is admissible; for $p\mid N$ exactly those with $(c_j,d_j)\epsilon\not\equiv0\pmod p$. Every $\adjg h_j(\epsilon)$ is integral of
determinant $p$. \LEAN{prop:heilbronn-model} (the definition, as an executable specification over $\Z/q$, and the integrality and
determinant of every $\adjg h_j(\epsilon)$, with lower left entry $\equiv0\pmod N$) \LEAN{chk:merel-n5} (oracle note: the C++
$H_{5,p}$ at $n=2$ over $\F_{1000003}$, $p\in\{2,5\}$, is this operator on its dumped table, kernel-checked)
\begin{enumerate}
\item $H_{N,p}(\mathrm{Rel}_N)\subseteq\mathrm{Rel}_N$.
\item $\varkappa_N\circ[H_{N,p}]=[A\,g\,A]\circ\varkappa_N$ on $M_N/\mathrm{Rel}_N$; $[A\,g\,A]$ is $T^{(N)}_p$ for $p\nmid N$ and
  $U^{(N)}_p$ for $p\mid N$ --- \textbf{the operator that Proposition~\ref{prop:levelcorrespondence}(4)(d) induces on $H^{1}$}
  (its item (3), with Proposition~\ref{prop:lhcompliant}(4) and Shapiro's lemma for the model), independent of $\tau$ and of
  every table.
\item At $N=1$: $T_p$ on $\VR/(\VR^{S}+\VR^{U})=H^{1}(\Ga;\VR)$ is $x\mapsto\sum_{\epsilon\in\mathrm{Mer}_p}\rhon(\adjg\epsilon)x$.
\item $H_{N,p}$ reads Definition~\ref{def:cosetrepsN}'s table, $\mathrm{Mer}_p$ and $\rhon$ of integral matrices; no $\tau$-table,
  no coset search, and no matrix outside $\Ga$ is inverted (Obligation~\ref{obl:integral}). \LEAN{prop:heilbronn-model} (not (1)--(3))
\end{enumerate}
\end{proposition}

\begin{proof}
By Lemma~\ref{lem:levelduality}(2),(6), $\Pi_A\varkappa_N=\mu_N$ and $\Pi_A[A\,g\,A]=[A\,g\,A]_{\ast}\Pi_A$, with $\Pi_A$ an isomorphism and
$\ker\mu_N=\mathrm{Rel}_N$. It therefore suffices to prove
\[
  (\ast)\qquad[A\,g\,A]_{\ast}\,\mu_N(x)=\mu_N(H_{N,p}x)\qquad(x\in M_N):
\]
then (1) follows ($\mu_N(H_{N,p}r)=[A\,g\,A]_{\ast}\mu_N(r)=0$ for $r\in\mathrm{Rel}_N$), and (2) from
$\Pi_A\varkappa_N(H_{N,p}x)=[A\,g\,A]_\ast\Pi_A\varkappa_N(x)=\Pi_A[A\,g\,A]\varkappa_N(x)$. (3) is $N=1$, (4) is read off
\eqref{eq:heilbronn}; integrality and the determinant by Lemma~\ref{lem:cosetrepsN}(4) and \eqref{eq:adjrules}.

\emph{Step 1: the right cosets.} $A\cap g^{-1}Ag=\{a\in A:p\mid a_{12}\}$, and $a\mapsto(a_{11}:a_{12})\bmod p$ identifies its right cosets
in $A$ with a subset of $\mathbb{P}^1(\F_p)$ ($h$ in it multiplies the top row by the unit $h_{11}$; proportional top rows give $a'a^{-1}$
with $(a'a^{-1})_{12}\equiv0$). If $p\nmid N$ the image is all of $\mathbb{P}^1(\F_p)$ ($T^{y}\in A$ gives $(1:y)$; $\smat pu Nv\in A$,
$pv-uN=1$, gives $(0:1)$), so $\abs{A\backslash AgA}=p+1$; if $p\mid N$ it is $\{(1:y)\}$ ($a_{11}a_{22}\equiv1\bmod p$), so
$A\backslash AgA$ is represented by $\beta_b:=gT^{b}=\smat1b0p$, $0\le b<p$. \emph{Claim}: for each $j$ the right cosets $\Ga\beta w_j$,
$A\beta\in A\backslash AgA$, are pairwise distinct and are all the $p+1$ right cosets of $\Ga$ in $\De_p$ (Proposition~\ref{prop:cosets}(1))
if $p\nmid N$, all but $\Ga\smat p001w_j$ if $p\mid N$. For $p\mid N$: the $\Ga\beta_b$ are distinct Hermite forms and the missing
form is $\smat p001$. For $p\nmid N$: every $\beta\in AgA$ is upper triangular modulo $N$ with unit diagonal ($A$ is, and
$g\equiv\mathrm{diag}(1,p)$), so $\Ga\beta=\Ga\beta'$ gives $\gamma=\beta'\,\adjg\beta/p\in\Ga$ with $\gamma_{21}\equiv0\pmod N$, i.e.\
$A\beta'=A\beta$; $p+1$ distinct cosets exhaust $\Ga\backslash\De_p$.

\emph{Step 2: Merel's decomposition.} For $\eta\in\De_p$, apply Lemma~\ref{lem:merelpath} at the left coset $\adjg(\eta)\Ga$ and then
$\eta$: every $\epsilon\in\mathrm{Mer}_p\cap\adjg(\eta)\Ga$ has $\eta\epsilon=p\gamma_\epsilon$ with $\gamma_\epsilon\in\Ga$, and
\[
  \eta\cdot D_0\;=\;\sum_{\epsilon\in\mathrm{Mer}_p\cap\adjg(\eta)\Ga}\gamma_\epsilon\cdot D_0 .
\]

\emph{Step 3: the computation.} $[A\,g\,A]_{\ast}\mu_N(x)=\sum_j\sum_{A\beta}[\beta w_jD_0\otimes\rhon(\beta)x_j]_A$. Apply Step~2 to
$\eta=\beta w_j$: $\gamma_\epsilon=p^{-1}\beta w_j\epsilon$. Write $\gamma_\epsilon=hw_i$, $h\in A$, $i=\mathrm{cls}_N(\gamma_\epsilon)$:
$[\gamma_\epsilon D_0\otimes\rhon(\beta)x_j]_A=[w_iD_0\otimes\rhon(w_i\gamma_\epsilon^{-1}\beta)x_j]_A$ and $w_i\gamma_\epsilon^{-1}\beta=
w_i(\adjg\epsilon)w_j^{-1}$. The index: for $p\mid N$ and $\beta=\beta_b$, $\mathrm{row}_2(\gamma_\epsilon)=\mathrm{row}_2(w_j)\epsilon$ exactly;
for $p\nmid N$, $\mathrm{row}_2(\beta)\equiv(0,u)\pmod N$ with $u$ a unit and $\mathrm{row}_2(\gamma_\epsilon)\equiv(u/p)\,\mathrm{row}_2(w_j)\epsilon$;
either way $i=j\cdot\epsilon$, and the term is $[w_{j\cdot\epsilon}D_0\otimes\rhon(\adjg h_j(\epsilon))x_j]_A$. The result does not depend
on the representative $\beta$ of $A\beta$ ($\gamma_\epsilon$ moves by $A$ on the left). Finally, for fixed $j$, the left cosets
$\adjg(\beta w_j)\Ga=\adjg(\Ga\beta w_j)$ are by Step~1 pairwise distinct and are all left cosets of $\De_p$ ($p\nmid N$), respectively
all but $\adjg(\smat p001w_j)\Ga=w_j^{-1}g\Ga$; and $\epsilon\in w_j^{-1}g\Ga$ iff $w_j\epsilon\in g\Ga=\{\delta\in\De_p:\mathrm{row}_2(\delta)
\equiv0\bmod p\}$ iff $(c_j,d_j)\epsilon\equiv0\pmod p$ --- exactly the inadmissible pairs, as $\det\epsilon=p$ is a unit modulo the other
primes of $N$. Every $\epsilon\in\mathrm{Mer}_p$ lies in one left coset, so the double sum is over the admissible pairs, which is
$\mu_N(H_{N,p}x)$.
\end{proof}

\begin{remark}[The transcription, and the trap in it]\label{rmk:heilbronnconv}
\PROVED{} \COMP\textsuperscript{a} (1) \textbf{Against Merel's and Stein's conventions} (pointers of Remark~\ref{rmk:levelscaling}, not verified): a
Manin symbol $[P,(c:d)]$ with the right action $P|\delta$ is, here, block $j$ with the class $(c:d)$ and a value in $\VR$; the
index moves by the \emph{same} row action $j\mapsto j\cdot\epsilon$, and the value by $\rhon(\adjg h_j(\epsilon))$ --- the column
convention's form of a right action, $\rhon(\adjg\,\cdot)$ being an anti-homomorphism. Terms with $j\cdot\epsilon$ not
primitive are dropped, which is Merel's rule at $p\mid N$. (2) \textbf{The trap}: what Step~2 of the proof uses is the
\emph{right-coset} condition for $\adjg(\mathrm{Mer}_p)$, which Lemma~\ref{lem:merelpath} gives through $\adjg$; $\mathrm{Mer}_p$ itself
does \emph{not} satisfy it --- at $p=2$ the right coset $\Ga\smat1102$ contains no element of $\mathrm{Mer}_2=\{\smat1002,\smat1012,
\smat2001,\smat2101\}$. Two plausible mis-transcriptions --- (a) $\rhon(h_j(\epsilon))$ in place of $\rhon(\adjg h_j(\epsilon))$; (b)
$\mathrm{Mer}_p$ itself as the right-coset set, i.e.\ index $i=j\cdot\adjg\epsilon$ and value $\rhon(w_i\epsilon w_j^{-1})$ --- do not even
preserve $\mathrm{Rel}_N$ at $(N,n,p)=(1,10,2)$, $(5,2,2)$, $(5,2,5)$, $(6,2,5)$ (\COMP\textsuperscript{a}, slot-14 WP-4's
\file{docs/implementation/slot14/wp4/variants.py}); at weight $2$, where $\rhon$ is trivial, (a) coincides with \eqref{eq:heilbronn} and only (b) is caught, which is
why \ref{LN3} tests weights above $2$. (3)
\textbf{Two routes, one operator.} Proposition~\ref{prop:levelcorrespondence}(4)(d) is a cochain map on the fixed model; $H_{N,p}$ is
a map on $M_N$ that preserves $\mathrm{Rel}_N$ but is not the restriction of a cochain map of the model. They agree on $H^{1}$,
under the canonical identifications, and are compared there --- by characteristic polynomials, or through any $\Ga$-chain map
between $C_{\bullet}(\Ku)$ and $C_{\bullet}(\mathcal Y)$ over the identity. (4) \textbf{Composite $n$}: Lemma~\ref{lem:merelpath} holds for
every $n$, but the double-coset bookkeeping of Step~1 is done for primes only; $T_{p^{r}}$ and $T_m$ follow from $T_p$, $U_p$ by
Theorem~\ref{thm:relations}(1) and its recursion.
\end{remark}

\subsection{The cost of the route}\label{sec:heilbronncostsec}

\begin{remark}[The Heilbronn route's cost; scalability in $N$, $k$ and $\ell^{l}$]\label{rmk:heilbronncost}
\COMP\textsuperscript{a} A model of operation counts, in the notation of Remark~\ref{rmk:levelcost} ($d=n+1$,
$D_N=m_Nd$); not an assertion about running times.
\begin{enumerate}
\item \textbf{Size of $\mathrm{Mer}_p$.} Proved: $\abs{\mathrm{Mer}_p}\le\sum_Lx_0(L)=p^{2}+1$ (Lemma~\ref{lem:merelpath}, over the
  $p+1$ lattices of index $p$: $x_0=1$ for $\{y\equiv0\ (p)\}$ and $x_0=p$ for the $p$ lattices $\{x\equiv by\ (p)\}$). Counted (slot-14 WP-4's \file{docs/implementation/slot14/wp4/merel\_count.py}): $4,7,15,25,49,63$ at
  $p=2,3,5,7,11,13$; $1097$, $13045$, $48653$, $259245$ at $p=101$, $691$, $2003$, $8009$ --- $0.40$--$0.51\,p\log^{2}p$ on
  $101\le p\le8009$, the ratio decreasing, while $\abs{\mathrm{Mer}_p}/(p\log p)$ \emph{increases} ($2.35$ at $101$ to $3.60$ at $8009$):
  Merel's set is not $O(p\log p)$ on this range. The smaller set of Heilbronn matrices built from continued fractions
  \STD{\cite[Ch.~2]{cremona}; \textbf{not} verified, rule~\ref{A2}} ($O(p\log p)$ elements by the cited source) is not specified
  here; any finite set $Y$ with $\sum_{\delta\in Y\cap\Ga\eta}\adjg(\delta)\cdot D_0=D_0$ for every right coset $\Ga\eta\subseteq\De_p$ can
  replace $\adjg(\mathrm{Mer}_p)$ in the proof without change.
\item \textbf{Applying $H_{N,p}$} to $x\in M_N$: per pair $(\epsilon,j)$ the index $j\cdot\epsilon$ ($O(\omega(N))$ operations modulo the
  $Q_i$), the $2\times2$ product $\adjg h_j(\epsilon)$, the $d\times d$ matrix $\rhon(\adjg h_j(\epsilon))$ ($O(d^{2})$ by
  \code{ComputeSymmetricPowerMatrixFM}) and one block times a vector ($O(d^{2})$): $O(\abs{\mathrm{Mer}_p}\,m_N\,d^{2})$ ring operations,
  memory $O(D_N+d^{2})$. \textbf{No $D_N\times D_N$ object}, no $\tau$-table, no coset search.
\item \textbf{A Hecke matrix on $H^{1}$} ($\dim H^{1}\approx D_N/6$): on a Manin-symbol basis (one block, one basis vector of $\VR$) $H_{N,p}$
  costs $O(\abs{\mathrm{Mer}_p}\,d^{2})$ and touches $\le\abs{\mathrm{Mer}_p}$ blocks; expressing the image in the basis of
  $M_N/\mathrm{Rel}_N$ costs $O(\abs{\mathrm{Mer}_p}\,d\cdot\dim H^{1})$ through the normal-form table of $\mathrm{Rel}_N$ ($D_N\times\dim H^{1}$,
  the output of (4); memory $\Theta(D_N\cdot\dim H^{1})$), or one sparse solve per column. In all
  $O(\abs{\mathrm{Mer}_p}\cdot d\cdot(\dim H^{1})^{2})$ for the matrix, \textbf{quadratic in $m_N$}, against the dense route's
  $\Theta((N_qD_N)^{2})$ projector per column (Remark~\ref{rmk:levelcost}(1),(5)). What obstacle~\ref{levelO1} loses for Hecke operators is
  the $N_q^{2}$ factor and the dense model projector, not the quadratic order (v1.14.0-pre.2, WP-5 referee F1).
\item \textbf{The elimination} (Remark~\ref{rmk:levelcost}(4)): the orbit-local quotient by $M^{\langle S\rangle}$, then one
  block-sparse elimination against $M^{\langle U\rangle}$ --- about $D_N/2$ rows, $D_N/3$ columns, $\le3$ blocks of size $d\times d$ per
  block column; worst case cubic in $D_N$, and \textbf{no better bound is asserted}. At $d=1$ (weight $2$) it is the classical Manin-symbol
  elimination, with $\le3$ entries per relation. Nothing in Proposition~\ref{prop:heilbronn} forces a dense $D_N\times D_N$ object,
  so the sparse route at $N\sim10^{4}$, $k=2$ is admitted: $m_N\sim10^{4}$, $\abs{\mathrm{Mer}_5}=15$, about $1.5\cdot10^{5}$ block
  operations per application of $H_{N,5}$. The normal form of $\mathrm{Rel}_N$ is the one $D_N\times\dim H^{1}$ object of the route;
  (2)'s ``no $D_N\times D_N$ object'' is about one application (v1.14.0-pre.2, WP-5 referee F1).
\item \textbf{The three axes.} Large $N$ at small $k$ (the regime (R-N) of Remark~\ref{rmk:levelscaling}): linear in $m_N$ per operator
  application, quadratic per Hecke matrix on $H^{1}$ (3), the elimination as in (4). Large $k$: $d^{2}$ per term, $d^{3}$ once for the fixed spaces of $\rhon(S)$, $\rhon(U)$
  (Remark~\ref{rmk:levelcost}(3)); large $k$ and large $N$ together stay \ref{levelO4} (the fill of (4) with dense $d\times d$
  blocks). $\cring=\Z/\ell^{l}$: every count times one ring operation (a word, or the digit route beyond $2^{64}$), $\mathrm{Mer}_p$
  independent of $\cring$, $N$, $k$; the orbit-wise averaging and the cokernel need $6\in\cring^{\times}$ (\ref{levelO6}).
\end{enumerate}
\COMP\textsuperscript{a} here means the counts of (1), from an out-of-tree script; every other item is a model.
\end{remark}

\subsection{The predictions}\label{sec:levelNpredsec}

\begin{remark}[What is checkable at level $N$: (LN1)--(LN6)]\label{rmk:levelNcheckable}
\PROVED{} \STD{\cite{shimura,mazur,merel96,wwe}} \CONJ{} \COMP\textsuperscript{a} Committed at v1.14.0-pre.1, before slot-14 WP-12's
code (rule~20), with the convention of Remark~\ref{rmk:hidacheckable}: a failing \PROVED{} item stops the package, a failing
\CONJ{} item is a finding. \COMP\textsuperscript{a} evidence: slot-14 WP-4's scripts in \file{docs/implementation/slot14/wp4/} ---
\file{levelN.py}, \file{one\_cell.py}, \file{n1\_dims.py}, \file{eisenstein\_scan.py} (tracked, read-only, not registered validations); their
independent route is \emph{not} the $\tau$-route and \emph{not} the Heilbronn route --- $\cor\circ(\conj_g)_{\ast}\circ\res$ on
inhomogeneous $\Ga_0(N)$-cocycles, carried to $\Ga$ by the cochain-level Shapiro map and normalised as in
Definition~\ref{def:kappaN}.
\begin{enumerate}[label=(LN\arabic*)]
\item\label{LN1} \STD{\cite[Ch.~2, Ch.~8]{shimura}, the dimension formulas and Eichler--Shimura; chapter level, \textbf{not}
  verified} \textbf{Dimensions.} For $n$ even and a field $\cring$ with $6\in\cring^{\times}$, the right-hand side of
  Proposition~\ref{prop:levelorbit}(3) (which is $\dim H^{1}-\dim H^{0}$, independent of the characteristic) equals
  $2\dim S_{n+2}(\Ga_0(N))+c(N)-2\delta_{n,0}$, $c(N)=\sum_{t\mid N}\abs{(\Z/\gcd(t,N/t))^{\times}}$ the number of cusps; over $\Q$, where
  $\dim H^{0}=\delta_{n,0}$, this is $\dim H^{1}=2\dim S_{n+2}+c(N)-\delta_{n,0}$. \COMP\textsuperscript{a}: equal for all $N<200$ and
  even $n\le20$ (\file{n1\_dims.py}, $0$ mismatches). The code's check: $\dim H^{1}-\dim H^{0}$ computed by either route against the
  formula at composite $N$ (over $\Fl$ both dimensions can exceed the characteristic-zero ones by the same amount).
\item\label{LN2} \PROVED{} (Proposition~\ref{prop:levelorbit}(1)) \textbf{Orbit-wise $=$ dense}: the fixed spaces
  $M_N^{\langle S\rangle}$, $M_N^{\langle U\rangle}$ computed orbit by orbit equal the dense ones as subspaces, so both routes give the same
  $\mathrm{Rel}_N$ and the same cokernel, entrywise in the coordinates of Definition~\ref{def:levelNmodel}.
\item\label{LN3} \PROVED{} (Proposition~\ref{prop:heilbronn}) \textbf{Heilbronn $=$ $\tau$-route on $H^{1}$}: at every tested $(N,n,p,\cring)$
  --- $p\nmid N$ and $p\mid N$, prime and composite $N$ including mixed classes, $\cring=\Fl$ with $\ell\mid N$ and $\ell\nmid N$ ---
  $H_{N,p}(\mathrm{Rel}_N)\subseteq\mathrm{Rel}_N$ entrywise, and the characteristic polynomial of $[H_{N,p}]$ on $M_N/\mathrm{Rel}_N$ equals that of
  Proposition~\ref{prop:levelcorrespondence}(4)(d) on $\Heq{1}(K,\Rel;L_N)$. \COMP\textsuperscript{a}: equal to the independent
  cocycle route on every basis vector, and $H_{N,p}(\mathrm{Rel}_N)\subseteq\mathrm{Rel}_N$, at the cells listed in
  \file{docs/implementation/slot14/WP-4\_report.md} (among them $(N,n,p)=(1,10,2)$, $(11,0,2)$, $(11,0,3)$, $(11,0,11)$, $(5,2,2)$,
  $(5,2,5)$, $(6,2,2)$, $(6,2,3)$, $(6,2,5)$, $(35,0,2)$, $(35,0,5)$, $(35,0,7)$ over $\F_{1000003}$), with the classical
  eigenvalues where they are known: $(X+24)^{2}(X-2049)$ for $T_2$ at weight $12$, level~$1$; $(X+2)^{2}(X-3)$, $(X+1)^{2}(X-4)$,
  $(X-1)^{3}$ for $T_2$, $T_3$, $U_{11}$ at weight $2$, level $11$; $(X+4)^{2}(X-9)^{2}$ and $(X+5)^{2}(X-1)(X-125)$ for $T_2$, $U_5$
  at weight $4$, level $5$; $(X+2)^{2}(X-1)^{2}(X-8)^{2}$ for $U_2$ at weight $4$, level $6$; $X^{2}(X-3)^{3}(X^{2}+X-4)^{2}$ for
  $T_2$ at weight $2$, level $35$; and, re-run by the slot-14 WP-5 referee (\file{docs/implementation/slot14/WP-5\_referee.md} \S3.2),
  $(13,2,2)$, $(15,2,5)$, $(4,4,2)$, $(20,2,5)$, $(22,2,11)$, $(8,2,2)$ over $\F_7$, and $(30,0,2)$. That independent route tests
  Proposition~\ref{prop:heilbronn}(2), $H_{N,p}=[A\,g\,A]$; the equality with the $\tau$-route of
  Proposition~\ref{prop:levelcorrespondence}(4)(d) on the model --- this item's other half --- is WP-12's (v1.14.0-pre.2, WP-5
  referee F8).
  \COMP\textsuperscript{a} (v1.14.2, WP-12): the $\tau$-route half --- the characteristic polynomial of
  Proposition~\ref{prop:levelcorrespondence}(4)(d) on $\Heq{1}(K,\Rel;L_N)$ equals the independent cocycle route's on $20$ cells
  (among them the mixed-class levels $6,10,12,15,21,30,35$ and $n=2,4$) and the classical values listed above; the Heilbronn half
  awaits $H_{N,p}$. \COMP\textsuperscript{a} (v1.15.0-pre.2, slot-15 WP-7): the Heilbronn half --- $H_{N,p}$ of \eqref{eq:heilbronn} built
  (\file{Hecke/HeckeLevelNHeilbronn}), $\mathrm{Rel}_N$ preserved and its characteristic polynomials equal to the $\tau$-route's, to the
  independent cocycle route's and to the classical values above, the two mis-transcriptions of Remark~\ref{rmk:heilbronnconv}(2) caught;
  \code{HeckeLevelNHeilbronnValidation} (\code{standard}, 636 checks; \file{docs/implementation/slot15/WP-7\_report.md}). \ref{LN3} is tested
  on both halves.
\item\label{LN4} \STD{\cite[Ch.~II]{mazur}; chapter level, \textbf{not} verified} \CONJ{} \textbf{Mazur's Eisenstein ideal, in
  cohomology.} Let $N$ be prime, $p\ge5$ prime, $p\mid\mathrm{num}\bigl(\frac{N-1}{12}\bigr)$, and $e:=\rank_{\Z_p}\Talg^{0}_{\midl}$ the rank of
  Mazur's cuspidal Hecke algebra completed at the Eisenstein maximal ideal $\midl$ (it is $\ge1$ by the cited source). Prediction: the
  joint generalised eigenspace $W_{N,p}\subseteq H^{1}(\Ga_0(N);\F_p)$ of $(T_q)_{q\ne N}$ at the Eisenstein eigenvalues $1+q$ has
  $\dim W_{N,p}=2e+1\ge3$. (The evidence of \ref{LN6} computes $W_{N,p}$ jointly only for $q\le7$, $q\ne N$; that space contains this
  $W_{N,p}$, and their equality is part of the prediction; v1.14.0-pre.2, WP-5 referee F2.) \CONJ{} because the chain from Mazur's statements about $J_0(N)$ and its Hecke algebra to this document's
  $H^{1}(\Ga_0(N);\F_p)$ is not carried out here.
\item\label{LN5} \STD{\cite[Th\'eor\`eme~2, p.~72]{merel96} (verified on the GDZ page images, slot-17 R17-4, v1.17.0-pre.5; in Merel's notation the prime
  level is $p$ and the residue characteristic $q$), restated as \cite[Thm.~1.5.1, p.~5]{wwe} (verified on the arXiv text, slot-16 WP-8)} \CONJ{} \textbf{Merel's criterion}:
  under Merel's hypotheses --- $p>3$ (here $p\ge5$) and $p\mid(N-1)/\gcd(N-1,12)$, which for $p\ge5$ is \ref{LN4}'s $p\mid\mathrm{num}\bigl(\frac{N-1}{12}\bigr)$ ---
  $e\ge2$ iff $\prod_{i=1}^{(N-1)/2}i^{\,i}$ is a $p$-th power in $(\Z/N)^{\times}$; so, with \ref{LN4}, $\dim W_{N,p}\ge5$ exactly at those cells.
  The criterion is evaluated as $\sum_{i\le(N-1)/2}i\log i\equiv0\pmod p$, $\log$ the discrete logarithm to any primitive root modulo $N$.
\item\label{LN6} \STD{\cite[p.~8 and Cor.~7.2.2, p.~30]{wwe} (verified on the arXiv text, $p>3$, slot-16 WP-8)} \CONJ{} \textbf{Mod-$p$ cohomological Loewy length $e+1$.} The Hecke algebra
  acting on $W_{N,p}$ has Loewy length $e+1$ --- the cuspidal part contributing $e$ (by the cited source $\Talg^{0}_{\midl}/p\cong
  \F_p[x]/(x^{e})$) and the Eisenstein class one more step. This is route R3: level $N$ gives mod-$p$ Loewy length $\ge3$ in
  cohomology, which level~$1$ shows only on classes that are not congruence-free, the classes [E] \Espec{pr:latticesplit}~(R) excludes: on cohomology at Kummer
  classes (\Espec{pr:kfonset}; $(608,547)$), and on the forms side also at irreducible congruence classes (slot-17 R17-2~\S7) (v1.17.0-pre.3, slot-17
  referee F1's m-C3: through v1.17.0-pre.2 this read ``which the census at level~$1$ never shows'', false since slot~15's $(258,37)$; final wording
  v1.17.0-pre.4, after the owner's ruling on the referee's B-1 restated (R) on congruence-free classes).
  \COMP\textsuperscript{a} (\file{eisenstein\_scan.py}; Loewy length by the radical series): $W_{N,p}$ is computed as the joint
  generalised eigenspace of $T_2$, $T_3$, $T_5$, $T_7$ ($q\ne N$) only, which contains \ref{LN4}'s $W_{N,p}$ --- equality is part of the
  prediction --- and $e$ is not computed independently but read off $\dim W_{N,p}$ (v1.14.0-pre.2, WP-5 referee F2). On
  all $168$ cells with $N\le2000$, $p\in\{5,7,11,13\}$, $\dim W_{N,p}\in\{3,5,7\}$, $\dim W_{N,p}\ge5$ exactly at the Merel cells, and
  the Loewy length is $(\dim W_{N,p}+1)/2$ throughout: $\dim W_{N,p}=3$ at $137$ cells, $5$ at $28$, $7$ at $3$ --- under \ref{LN4}
  $e=1$, $2$, $3$; the three $\dim W_{N,p}=7$ cells are \textbf{$(N,p)=(181,5)$, $(1321,11)$, $(1571,5)$, Loewy length $4$}.
  \COMP\textsuperscript{a} (v1.14.2, slot-14 WP-12, \code{HeckeLevelNValidation} arm (R3)): a second, independent route --- the
  $\tau$-route on the fixed model, $H^{1}(\Ga_0(N);\F_p)$ by \eqref{eq:levelcorr} with Definition~\ref{def:cosetrepsN}'s table,
  against the Manin-symbol cokernel of \file{levelN.py} --- reproduces $\dim H^{1}$, $\dim W_{N,p}$, the Loewy length and Merel's
  criterion (recomputed in C++) on all $26$ cells with $N\le211$, and gives $\dim W_{N,p}=7$, radical layers $7,5,3,1$, Loewy
  length $4$ at $(181,5)$, $(1321,11)$, $(1571,5)$; on every cell the layers are $2e+1,2e-1,\ldots,1$. Both routes compute
  $W_{N,p}$ jointly for $q\le7$, $q\ne N$, only, and read $e$ off $\dim W_{N,p}$ (the WP-5 referee F2 caveat stands); the Loewy
  length is that of the subalgebra generated by those $T_q$ --- a single $q$ can under-generate (at $(151,5)$ the nilpotency
  index of $T_2-3$ is $1$ against Loewy length $2$), so the joint $q\le7$ is the least admissible generating set.
  The scan to $N\le10^{4}$ (Merel's criterion only) has $733$ cells, $118$ of them Merel
  cells ($78$, $22$, $11$, $7$ at $p=5,7,11,13$). \textbf{Named for WP-12}: $(181,5)$ ($m_N=182$, weight $2$: $D_N=182$, within the
  dense route's reach), then $(1321,11)$ and $(1571,5)$.
  \emph{Addendum} (v1.21.0-pre.4; the extension of \ref{LN4} and \ref{LN6} to $p\in\{2,3\}$, $n=0$, on the curve object of
  Proposition~\ref{prop:levelNcurve}, Remark~\ref{rmk:levelNcurvefamily}; \CONJ{} as here): the same F2 caveat applies there --- the code
  reads the joint generalised eigenspace $W_{\mathcal F}$ of a finite family $\mathcal F$ of $T_q$ ($q\le13$, $q\ne N$, by default), which
  contains $W_{N,p}$, and a radical series at most $e+1$ long; equality of both is part of the prediction, and a family can fail it by
  generation (no Mazur good prime in $\mathcal F$) or by separation (Remark~\ref{rmk:levelNcurvefamily}(1),(2)).
\end{enumerate}
\end{remark}

\subsection{Degeneracy maps on the presentation, the model-to-presentation map, and higher weight at scale}\label{sec:levelNdegsec}
% NEW at v1.16.0-pre.1 (hecke-extensions slot 16, WP-12; plan docs/implementation/slot16/00-slot16-plan.md S2 "WP-12"
% item (a) and (c)).  The contract first (rule 21) for slot-16 WP-13.  Replaces pre.3's unrefereed sentence in the
% scope paragraph below (fibre maps on the presentation; alpha_d on the model only) by numbered statements:
% prop:kappamodel (the pinned model -> presentation map, kappa_N read on the fixed model, and the naturality that
% carries every fibre map), prop:degenpres (alpha_d, d > 1, on the presentation, proved equal to the model's),
% rmk:degenpresconv, rmk:levelNscale, rmk:levelNcheckabletwo ((LN7)--(LN9)).  Placed at the end of sec:levelN so that
% no chapter or subsection number moves.

This subsection does three things. It reads $\varkappa_N$ of Definition~\ref{def:kappaN} on the fixed model, so that the
model's cocycles and the presentation can be compared \emph{entrywise} (Proposition~\ref{prop:kappamodel}). It puts the
degeneracy maps $\alpha_d=[\Ga_0(M)\,\smat100d\,\Ga_0(N)]$, $d>1$, on the cokernel presentation by Merel's set of
determinant $d$, proved equal to Proposition~\ref{prop:levelcorrespondence}(4)(c) on $H^{1}$
(Proposition~\ref{prop:degenpres}). And it records the reach of the chapter in $N$ and $k$ together
(Remark~\ref{rmk:levelNscale}) with the checks (LN7)--(LN9) (Remark~\ref{rmk:levelNcheckabletwo}).
$\kappa^{\mathrm{mod}}_N$, $e_{46}$, $D_{M,N,d}$, $\mathrm{R3}_k$ are new symbols of this subsection and mean nothing else
(rule~\ref{A5}); $D_{M,N,d}$ is neither the matrix $D$ of Remark~\ref{rmk:stabconvention}(3) nor $D_N$ of
Remark~\ref{rmk:levelcost}.

\begin{proposition}[The model-to-presentation map: $\varkappa_N$ read on the fixed model]\label{prop:kappamodel}
\PROVED{} Assume \ref{H4}. Let $\cring$ be a commutative ring with $6\in\cring^{\times}$, $N\ge1$, and let $M$ be the fibre of
an induced matrix local system $L$ on $(K,\Rel)$ in the sense of Lemma~\ref{lem:inducedchar} --- $M=M_N$ with $n$ even
(Definition~\ref{def:levelNmodel}), or $M=M_{N,\chi}$ with $\chi(-1)=(-1)^{n}$ (Definition~\ref{def:nebmodel}) --- so that
$-I$ acts trivially on $M$, and write $\mathrm{Rel}:=M^{\langle S\rangle}+M^{\langle U\rangle}$ ($\mathrm{Rel}_N$ of
Definition~\ref{def:kappaN} for $M=M_N$). Let $e_{46}:=\{v_4,v_6\}$ be the stored edge of Definition~\ref{def:model7}, so
$\partial\lift{e_{46}}=\lift{v_6}-\lift{v_4}$, and put
\[
  \kappa^{\mathrm{mod}}_N\colon\Cq{1}(K,\Rel;L)\to M,\quad f\mapsto-f(e_{46});\qquad
  \kappa^{\mathrm{mod},0}_N\colon\Cq{0}(K,\Rel;L)\to M,\quad g\mapsto g(v_4).
\]
On a $\Ga$-equivariant refinement of $\Xbulk$ (Remark~\ref{rmk:orbifoldrefine}(1)), $\lift{e_{46}}$ is replaced by the $1$-chain $\mathrm{path}(\lift{v_4}\to\lift{v_6})$ of the
refined edges and $\kappa^{\mathrm{mod}}_N(f):=-f(\mathrm{path}(\lift{v_4}\to\lift{v_6}))$; every item holds verbatim, $F_1(\mathfrak e):=-\mathrm{path}(\lift{v_4}\to\lift{v_6})$
being a $\Ga$-chain map by the same computation (v1.16.0-pre.3, WP-14 referee C-2).
\begin{enumerate}
\item \textbf{(It computes the cokernel.)} $\kappa^{\mathrm{mod}}_N$ maps $\Bq{1}$ into $\mathrm{Rel}$ and induces an isomorphism
  $\Heq{1}(K,\Rel;L)\to M/\mathrm{Rel}$; $\kappa^{\mathrm{mod},0}_N$ induces an isomorphism $\Heq{0}(K,\Rel;L)\to M^{\langle
  S\rangle}\cap M^{\langle U\rangle}$.
\item \textbf{(It is the pinned isomorphism.)} For $M=M_N$, $\varkappa_N\circ[\kappa^{\mathrm{mod}}_N]$ is the identification
  $\Heq{1}(K,\Rel;L_N)\cong H^{1}(\Ga_0(N);\VR)$ of Proposition~\ref{prop:lhcompliant}(4); likewise with $\chi$
  (Proposition~\ref{prop:neborbit}(3)). So $[\kappa^{\mathrm{mod}}_N]$ is the inverse of Definition~\ref{def:kappaN}'s isomorphism
  read on the model, and \textbf{no choice} enters beyond the binding table of Definition~\ref{def:cosetrepsN} and the stored
  orientation of $e_{46}$.
\item \textbf{(Naturality: fibre maps.)} Let $\phi\colon M\to M'$ be a map of $\cring[\Ga]$-modules between two such fibres. Then
  $\phi(\mathrm{Rel})\subseteq\mathrm{Rel}'$, the induced map $M/\mathrm{Rel}\to M'/\mathrm{Rel}'$ is $H^{1}(\Ga;\phi)$ under the
  isomorphisms of (1), and $\kappa^{\mathrm{mod}}_N(\phi\circ f)=\phi\bigl(\kappa^{\mathrm{mod}}_N(f)\bigr)$ \emph{exactly}, with no
  relation used. In particular the fibre maps $\alpha^{0}$ and $\omega$ of Proposition~\ref{prop:levelcorrespondence}(4)(b) and the
  diamonds $\langle d\rangle$ of Proposition~\ref{prop:nebdecomp}(1) act on the presentation by their own block matrices. (This is
  the sentence v1.15.0-pre.3 added to this chapter's scope paragraph, now a numbered statement.)
\item \textbf{(Entrywise comparison of operators.)} For every cocycle $f\in\Zq{1}(K,\Rel;L_N)$, every prime $p$ and $g=\smat100p$,
  \[
    \bigl(\widetilde{[A\,g\,A]}\,f\bigr)(e_{46})-H_{N,p}\,f(e_{46})\in\mathrm{Rel}_N\qquad(A=\Ga_0(N)),
  \]
  with \eqref{eq:levelcorr} on the left and \eqref{eq:heilbronn} on the right; for $dM\mid N$ and $f\in\Zq{1}(K,\Rel;L_M)$,
  $\bigl(\widetilde{[\Ga_0(M)\,\smat100d\,\Ga_0(N)]}\,f\bigr)(e_{46})-D_{M,N,d}\,f(e_{46})\in\mathrm{Rel}_N$ with
  \eqref{eq:degenpres}; and with a character, the same with $M_{N,\chi}$, \eqref{eq:nebcorr} and $H_{N,p,\chi}$ of
  Remark~\ref{rmk:nebheilbronn}.
\end{enumerate}
\end{proposition}

\begin{proof}
\emph{A $\Ga$-chain map from the tree to the model.} Define $F_0\colon C_0(\mathcal Y)\to C_0(\Xbulk)$ by $F_0(\gamma\mathfrak v_S):=
\gamma\lift{v_6}$, $F_0(\gamma\mathfrak v_U):=\gamma\lift{v_4}$, and $F_1\colon C_1(\mathcal Y)\to C_1(\Xbulk)$ by $F_1(\gamma\mathfrak e):=
-\gamma\lift{e_{46}}$. They are well defined because $\langle S\rangle=\mathrm{Stab}_\Ga(\lift{v_6})$, $\langle U\rangle=\langle ST\rangle=
\mathrm{Stab}_\Ga(\lift{v_4})$ (Lemma~\ref{lem:model7coincidences}(1)) and $-I$ fixes $\lift{e_{46}}$ with its orientation (it acts trivially on
$\Xbulk$, \ref{H4}(iv)); they are $\Ga$-linear by construction, and $\partial F_1(\mathfrak e)=-(\lift{v_6}-\lift{v_4})=F_0(\mathfrak v_U-\mathfrak v_S)=
F_0(\partial\mathfrak e)$, so $F$ is a chain map, and it preserves the augmentation (vertices to vertices). For $6\in\cring^{\times}$ both
$C_{\bullet}(\mathcal Y)\otimes\cring$ (Lemma~\ref{lem:treeres}) and $C_{\bullet}(\Xbulk)\otimes\cring$ (the proof of
Proposition~\ref{prop:levelcorrespondence}(3): the stabiliser orders divide $12$) are projective resolutions of $\cring$ over $\cring[\Ga]$,
and over $\cring[A]$ for every finite-index $A\le\Ga$; so $F$ is a comparison map and $F^{\ast}=\Hom_\Ga(F,M)$ induces \textbf{the} comparison
isomorphism on cohomology, for every $\cring[\Ga]$-module $M$ \STD{\cite[Thm.~2.2.6]{weibel}}.
\emph{(1)} By Theorem~\ref{thm:explicit_complex_equiv}(1) a cochain $f$ is $\phi_f$ with $\phi_f(\lift\sigma)=f(\sigma)$, so
$F^{\ast}\phi_f(\mathfrak e)=-f(e_{46})$ and, in degree $0$, $F^{\ast}\phi_g=(g(v_6),g(v_4))\in M^{\langle S\rangle}\oplus M^{\langle U\rangle}$.
In the tree complex $\Hom_\Ga(C_1(\mathcal Y),M)=M^{\langle-I\rangle}=M$ ($\phi\mapsto\phi(\mathfrak e)$) and the coboundary is $(a,b)\mapsto b-a$
(the proof of Definition~\ref{def:kappaN}'s claims; Proposition~\ref{prop:neborbit}(3) with a character), so $H^{1}=M/\mathrm{Rel}$ and $H^{0}=
\{(a,b):a=b\}\cong M^{\langle S\rangle}\cap M^{\langle U\rangle}$ via $b$; $F^{\ast}$ is a quasi-isomorphism. Directly: $(dg)(e_{46})=g(v_6)-g(v_4)$, so
$\kappa^{\mathrm{mod}}_N(dg)=g(v_4)-g(v_6)\in M^{\langle U\rangle}+M^{\langle S\rangle}$ ($g(v_6)$ is fixed by $S$, $g(v_4)$ by $ST$: the stabiliser
rows of Definition~\ref{def:model7rel}).
\emph{(2)} The reciprocity $\phi\mapsto\phi(\cdot)_1$ commutes with precomposition by any $\Ga$-chain map (the proof of
Proposition~\ref{prop:levelcorrespondence}(3)). So $\varkappa_N\circ[\kappa^{\mathrm{mod}}_N]$ is the map $\Hom_A(F,\VR)$ on
$A$-cochains, $A=\Ga_0(N)$, under the identifications of both sides with $H^{1}(A;\VR)$ by comparison with a fixed projective
resolution; $F$ restricted to $A$ is a chain map over $\id_{\cring}$ between projective $\cring[A]$-resolutions, so it induces the identity of
$H^{1}(A;\VR)$ under those identifications --- which are the ones Proposition~\ref{prop:lhcompliant}(4) and Definition~\ref{def:kappaN} use
(Theorem~\ref{thm:explicit_complex_equiv}, Remark~\ref{rmk:orbifold}, Shapiro's lemma, Lemma~\ref{lem:treeres}): the identification of
Proposition~\ref{prop:lhcompliant}(4) is comparison of the projective $\cring[A]$-resolution $C_{\bullet}(\Xbulk)\otimes\cring$ with any projective resolution,
then Shapiro, so ``no choice enters'' is checkable from that item alone (v1.16.0-pre.3, WP-14 referee C-3). With $\chi$ the same, with
$\rho_{n,\chi}$ (Lemma~\ref{lem:inducedchar}).
\emph{(3)} $\phi$ commutes with $M(S)$ and $M(U)$, so it maps $M^{\langle S\rangle}$, $M^{\langle U\rangle}$ into $M'^{\langle S\rangle}$,
$M'^{\langle U\rangle}$; the tree complex is natural in the module, and $F^{\ast}$ commutes with $\phi_{\ast}$. Exactness:
$(\phi\circ f)(e_{46})=\phi(f(e_{46}))$. A morphism of matrix local systems is one fibre matrix at every simplex
(Theorem~\ref{thm:rep_ls_equiv}), which is what (4)(b) of Proposition~\ref{prop:levelcorrespondence} and (1) of Proposition~\ref{prop:nebdecomp}
provide.
\emph{(4)} By Proposition~\ref{prop:levelcorrespondence}(3) the model operator induces $[A\,g\,B]$ under the identification of
Proposition~\ref{prop:lhcompliant}(4); by Proposition~\ref{prop:heilbronn}(2), Proposition~\ref{prop:degenpres}(2) and
Remark~\ref{rmk:nebheilbronn} the presentation operators induce the same map under $\varkappa$; by (2) the two identifications differ by
$[\kappa^{\mathrm{mod}}_N]$. So both sides of each display have the same class in $M/\mathrm{Rel}$ for every cocycle $f$.
\end{proof}

\begin{proposition}[The degeneracy maps $\alpha_d$ on the cokernel presentation]\label{prop:degenpres}
\PROVED{} Let $\cring$ be a commutative ring with $6\in\cring^{\times}$, $n$ even, $M,N,d\ge1$ with $dM\mid N$, $A:=\Ga_0(M)$, $B:=\Ga_0(N)$,
$g:=\smat100d$; let $w'_j$ ($1\le j\le m_M$) and $w_i$ ($1\le i\le m_N$) be the tables of Definition~\ref{def:cosetrepsN} at levels $M$ and $N$,
$\mathrm{Mer}_1:=\{I\}$, and for $\epsilon\in\mathrm{Mer}_d$ (Lemma~\ref{lem:merelpath}) call $(i,\epsilon)$ \emph{admissible} if
$\mathrm{row}_2(w_i)\,\adjg\epsilon\equiv0\pmod d$, and then put $j(i,\epsilon):=\mathrm{cls}_M\bigl(\mathrm{row}_2(w_i)\,\adjg\epsilon/d\bigr)$.
Define $D_{M,N,d}\colon M_M\to M_N$ by
\begin{equation}\label{eq:degenpres}
  D_{M,N,d}\;:=\;\sum_{\epsilon\in\mathrm{Mer}_d}\ \sum_{i\,:\,(i,\epsilon)\ \text{admissible}}
  E_{i,\,j(i,\epsilon)}\otimes\rhon\bigl(w_i\,(\adjg\epsilon)\,w'^{-1}_{j(i,\epsilon)}\bigr),
\end{equation}
i.e.\ $(D_{M,N,d}\,x)_i=\sum_\epsilon\rhon(w_i(\adjg\epsilon)w'^{-1}_{j(i,\epsilon)})\,x_{j(i,\epsilon)}$ over the $\epsilon$ admissible with $i$.
Every $w_i(\adjg\epsilon)w'^{-1}_{j(i,\epsilon)}$ equals $g\,a$ with $a\in A$, so it is integral of determinant $d$; $\mathrm{row}_2(w_i)\adjg\epsilon/d$ is
primitive, and admissibility and $j(i,\epsilon)$ depend only on the class $i$ (the canonical pair $(c_i,d_i)$ may replace $\mathrm{row}_2(w_i)$).
\begin{enumerate}
\item $D_{M,N,d}(\mathrm{Rel}_M)\subseteq\mathrm{Rel}_N$.
\item $\varkappa_N\circ[D_{M,N,d}]=\alpha_d\circ\varkappa_M$ on $M_M/\mathrm{Rel}_M$, where $\alpha_d:=[A\,g\,B]=(\conj_g)_{\ast}\circ\res^{A}_{g^{-1}Bg}$ is
  \textbf{the operator Proposition~\ref{prop:levelcorrespondence}(4)(c) induces on $H^{1}$}, independent of $\tau$ and of every table.
\item \textbf{($d=1$.)} $D_{M,N,1}$ is, entry for entry, the fibre map $\alpha^{0}$ of Proposition~\ref{prop:levelcorrespondence}(4)(b)
  for $B\subseteq A$ ($\epsilon=I$, every $i$ admissible, $j(i,I)=\mathrm{cls}_M(w_i)$).
\item \textbf{(Composite $d$.)} For $d=d_1d_2$, $[D_{M,N,d}]=[D_{Md_1,N,d_2}]\circ[D_{M,Md_1,d_1}]$ on $M_M/\mathrm{Rel}_M$ (by (2) and the
  factorisation of Proposition~\ref{prop:levelcorrespondence}(4)(c)); either side may be used, and they differ at cochain level by an
  element of $\mathrm{Rel}_N$.
\item \textbf{(The trace.)} If $N=dM$, then $[\omega_{B/A}]\circ[D_{M,N,d}]=[A\,g\,A]$ on $M_M/\mathrm{Rel}_M$, $\omega_{B/A}$ the trace of
  Proposition~\ref{prop:levelcorrespondence}(4)(b) acting by its block matrix (Proposition~\ref{prop:kappamodel}(3)): $T_p$ for $d=p\nmid M$, $U_p$
  for $d=p\mid M$, and at $M=1$, $d=p^{2}$, the primitive $T(1,p^{2})=T_{p^{2}}-p^{\,n}=T_p^{2}-p^{\,n}(p+1)$ ($T_{p^{2}}$ the full operator of \S\ref{err:e1}).
\item \textbf{(With a character.)} If $\chi\colon G_M\to\cring^{\times}$, $\chi(-1)=(-1)^{n}$, and $M_{N,\chi}$ is formed with $\chi$ composed
  with $G_N\to G_M$, then (1), (2) hold for $D_{M,N,d,\chi}$, \eqref{eq:degenpres} with each block multiplied by
  $\chi\bigl(a_{w_i(\adjg\epsilon)w'^{-1}_{j}}\bigr)$, on $M_{M,\chi}/\mathrm{Rel}\to M_{N,\chi}/\mathrm{Rel}$.
\item $D_{M,N,d}$ reads only the two tables, $\mathrm{Mer}_d$ and $\rhon$ of integral matrices: no $\tau$-table, no coset search, and no matrix
  outside $\Ga$ is inverted (Obligation~\ref{obl:integral}); $\abs{\mathrm{Mer}_d}\le\sigma_2(d)=\sum_{a\mid d}a^{2}$
  (Lemma~\ref{lem:merelpath}, summed over the lattices of index $d$, $x_0=a$ for the $a$ lattices with Hermite diagonal $(a,d/a)$); and at most
  $2m_N/m_d$ admissible $i$ per $\epsilon$ (the admissible bottom rows modulo $d$ lie in the kernel of $v\mapsto v\,\adjg\epsilon$ on $(\Z/d)^{2}$, of order
  $d$, with at most two primitive projective points), so $D_{M,N,d}$ costs $O(\abs{\mathrm{Mer}_d}\,m_M\,d_V^{2})$ when the admissible classes are enumerated
  per $\epsilon$ rather than scanned (v1.16.0-pre.3, WP-14 referee C-5).
\end{enumerate}
\end{proposition}

\begin{proof}
\emph{The groups.} $g^{-1}Bg=\{\smat ab{c}{e}\in\Ga: d\mid b,\ (N/d)\mid c\}\subseteq A$ (as $M\mid N/d$), and $B\subseteq gAg^{-1}$, so in
Proposition~\ref{prop:levelcorrespondence}(4)(c)'s notation $B''=B$, $A':=A\cap g^{-1}Bg=g^{-1}Bg$, and $[A\,g\,B]=(\conj_g)_{\ast}\circ\res^{A}_{A'}$;
$-I\in A'$, $gA'g^{-1}=B$. By Lemma~\ref{lem:levelduality}(3),(5),
\[
  \Pi_B\circ[A\,g\,B]=[A\,g\,B]_{\ast}\circ\Pi_A,\qquad [A\,g\,B]_{\ast}[y]_A:=\sum_{r\in A'\backslash A}[g r\cdot y]_B .
\]
As in the proof of Proposition~\ref{prop:heilbronn}, (1) and (2) follow from $\Pi_A\varkappa_M=\mu_M$, $\Pi_B\varkappa_N=\mu_N$ (Lemma~\ref{lem:levelduality}(2),
$\ker\mu=\mathrm{Rel}$, $\Pi$ an isomorphism) and
\[
  (\ast)\qquad[A\,g\,B]_{\ast}\,\mu_M(x)=\mu_N(D_{M,N,d}\,x)\qquad(x\in M_M).
\]
\emph{Step 1.} $[A\,g\,B]_{\ast}\mu_M(x)=\sum_j\sum_{r\in A'\backslash A}[g r w'_j D_0\otimes\rhon(g r)x_j]_B$. As $r$ runs over $A'\backslash A$ and $j$ over the
table, $u:=rw'_j$ runs once over $A'\backslash\Ga$, with $j=\mathrm{cls}_M(u)$ and $\rhon(gr)=\rhon(g u w'^{-1}_j)$; the summand
$[guD_0\otimes\rhon(guw'^{-1}_j)x_j]_B$ depends only on $A'u$ ($g a' g^{-1}\in B$ for $a'\in A'$). Put $\eta:=gu\in g\Ga$; $A'u\mapsto B\eta$ is a bijection
$A'\backslash\Ga\to B\backslash g\Ga$ ($Bg=gA'$).
\emph{Step 2 (Merel at determinant $d$).} As in Step~2 of the proof of Proposition~\ref{prop:heilbronn}, Lemma~\ref{lem:merelpath} at the left coset
$\adjg(\eta)\Ga$ gives $\eta\cdot D_0=\sum_{\epsilon\in\mathrm{Mer}_d\cap\adjg(\eta)\Ga}\gamma_\epsilon\cdot D_0$ with $\gamma_\epsilon:=d^{-1}\eta\epsilon\in\Ga$.
Write $\gamma_\epsilon=hw_i$, $h\in B$, $i=\mathrm{cls}_N(\gamma_\epsilon)$: the term is $[w_iD_0\otimes\rhon(w_i\gamma_\epsilon^{-1}\eta\,w'^{-1}_j)x_j]_B$, and
$\gamma_\epsilon^{-1}\eta=d\,\epsilon^{-1}=\adjg\epsilon$.
\emph{Step 3 (re-indexing).} The pairs $(B\eta,\epsilon)$ with $\eta\in g\Ga$, $\epsilon\in\mathrm{Mer}_d\cap\adjg(\eta)\Ga$ correspond bijectively to the pairs
$(B\gamma,\epsilon)$ with $\gamma\in\Ga$, $\epsilon\in\mathrm{Mer}_d$, $\gamma\,\adjg\epsilon\in g\Ga$, by $\gamma=\gamma_\epsilon$, $\eta=\gamma\,\adjg\epsilon$
($\epsilon\in\adjg(\gamma\adjg\epsilon)\Ga=\epsilon\gamma^{-1}\Ga$ always; $b\in B$ multiplies $\eta$ and $\gamma$ on the left together). Taking $\gamma=w_i$:
$w_i\adjg\epsilon\in g\Ga=\{\delta\in\Mat_2(\Z):\det\delta=d,\ \mathrm{row}_2(\delta)\equiv0\ (d)\}$ is admissibility, and then $u=g^{-1}w_i\adjg\epsilon\in\Ga$ has
bottom row $\mathrm{row}_2(w_i)\adjg\epsilon/d$, so $j=j(i,\epsilon)$ and $w_i(\adjg\epsilon)w'^{-1}_j=g\,(uw'^{-1}_j)$ with $uw'^{-1}_j\in A$. The double sum is therefore
$\sum_{(i,\epsilon)\ \text{admissible}}[w_iD_0\otimes\rhon(w_i(\adjg\epsilon)w'^{-1}_{j(i,\epsilon)})x_{j(i,\epsilon)}]_B=\mu_N(D_{M,N,d}x)$, which is $(\ast)$.
\emph{Dependence on the class $i$.} $\mathrm{row}_2(w_i)\equiv\lambda(c_i,d_i)\pmod N$ for a unit $\lambda$ (Lemma~\ref{lem:cosetrepsN}(2)); $d\mid N$ and $M\mid N/d$.
\emph{(3)} $\mathrm{Mer}_1=\{I\}$ (the only matrix with $ad-bc=1$, $a>b\ge0$, $d>c\ge0$), and the block is $\rhon(w_iw'^{-1}_{\mathrm{cls}_M(w_i)})$, which is
Proposition~\ref{prop:levelcorrespondence}(4)(b)'s $\alpha^{0}_{B/A}$ with $v_i=w_i$, $w_{a(i)}=w'_{\mathrm{cls}_M(w_i)}$.
\emph{(4)} is (2) for the three pairs and Proposition~\ref{prop:levelcorrespondence}(4)(c)'s factorisation on cohomology.
\emph{(5)} $A\cap gAg^{-1}=\Ga_0(dM)=B$ (the computation of Proposition~\ref{prop:cosets}(3)) and $A\cap g^{-1}Ag=g^{-1}Bg=A'$, so
$[A\,g\,A]=\cor^{A}_{B}\circ(\conj_g)_{\ast}\circ\res^{A}_{A'}=[B\,1\,A]\circ[A\,g\,B]$; $[B\,1\,A]=\omega_{B/A}$ (Proposition~\ref{prop:levelcorrespondence}(4)(b)), and
Proposition~\ref{prop:kappamodel}(3) carries it to the presentation. The double coset $A\smat100pA$ is $T_p$ resp.\ $U_p$ (Proposition~\ref{prop:levelcorrespondence}(4)(d));
at $M=1$, $\Ga\smat100{p^{2}}\Ga$ is the primitive double coset, $T_{p^{2}}=T(1,p^{2})+T(p,p)$ with $T(p,p)=p^{\,n}$ (Proposition~\ref{prop:diamond}) and
$T_{p^{2}}=T_p^{2}-p^{\,n+1}$ (Theorem~\ref{thm:relations}).
\emph{(6)} As in Remark~\ref{rmk:nebheilbronn}: the proofs above use of $\rhon$ only multiplicativity on the matrices they evaluate --- elements of $A$ and
$B$, $gr$ with $r\in A$, and $w_i\gamma_\epsilon^{-1}\eta w'^{-1}_j=h^{-1}gr$ --- all in $\De_0(M)$ ($a_g=1$, $c_g=0$), where $\rho_{n,\chi}$ is multiplicative
(Lemma~\ref{lem:nebdata}(2) at level $M$); $\VR(\chi)$ restricted to $B$ is the level-$N$ module for $\chi$ composed with $G_N\to G_M$; and
Lemma~\ref{lem:levelduality}(5) is used for $g\in\De_0(M)$ only.
\emph{(7)} is read off \eqref{eq:degenpres}; the count is Lemma~\ref{lem:merelpath}'s bound summed over the column lattices of index $d$, which in column
Hermite form $\smat a b0{d/a}$, $0\le b<a$, are $a$ in number for each $a\mid d$, each with $x_0=a$.
\end{proof}

\begin{remark}[The transcription of $\alpha_d$, its traps, and its evidence]\label{rmk:degenpresconv}
\PROVED{} \COMP\textsuperscript{a} (1) \textbf{The shape.} $D_{M,N,d}$ has the shape of \eqref{eq:heilbronn} read from the target: the source block $j$
is computed from the target class $i$ by the row action of $\adjg\epsilon/d$, and the value is $\rhon$ of $w_i(\adjg\epsilon)w'^{-1}_j$, of the form $g\cdot A$.
At $d=p$ prime and $M=1$, $N=p$ it is $\mathrm{dg}_p$ of Proposition~\ref{prop:levelellU} read on the presentation. (2) \textbf{The traps}, as in
Remark~\ref{rmk:heilbronnconv}(2): (a) $\epsilon$ in place of $\adjg\epsilon$ (Merel's set itself as the right-coset set) and (b) the adjugate of the block in
place of the block do \emph{not} preserve $\mathrm{Rel}$ at $(M,N,d,n)=(1,2,2,10)$, $(1,6,3,4)$, $(2,12,3,2)$, $(1,4,4,6)$, and disagree with the independent
route there (\file{docs/implementation/slot16/wp12/negctl.py}). (3) \textbf{Evidence} (\file{docs/implementation/slot16/wp12/degen.py}, over $\F_{1000003}$,
importing slot-14 WP-4's \file{levelN.py}; the independent route is $(\conj_g)_{\ast}\circ\res$ on inhomogeneous $\Ga_0(M)$-cocycles, carried to $\Ga$ by the
cochain-level Shapiro map at level $N$ and normalised as in Definition~\ref{def:kappaN} --- no Merel set, no presentation formula): (1), (2) and Hecke
equivariance $H_{N,p}D=DH_{M,p}$ modulo $\mathrm{Rel}_N$ ($p\in\{5,7\}$, $p\nmid N$) on $28$ cells $(M,N,d,n)$ --- among them $d=4,6,9$, mixed-class levels
$N=6,12,30,35$, $M>1$ with $\gcd(d,M)>1$ ($(2,4,2)$, $(3,9,3)$, $(2,8,4)$), and $d=1$ entrywise against $\alpha^{0}$; (5) on $8$ cells (four $T_p$, two at $M>1$, two
$T(1,p^{2})$); (4) on $6$ composites; and the old-space relations $U_pD_{M,Mp,1}=p\,D_{M,Mp,p}$ and $U_pD_{M,Mp,p}=D_{M,Mp,p}T_p-p^{\,n}D_{M,Mp,1}$ modulo
$\mathrm{Rel}$ ($p\nmid M$; the analogue of Proposition~\ref{prop:oldspace}(1),(3) with $\Ga_0(M)\to\Ga_0(Mp)$) on $7$ cells. All held; the run of record is in
\file{docs/implementation/slot16/WP-12\_report.md} \S3.
\end{remark}

\begin{remark}[Higher weight and composite $N$ at scale: the reach of this chapter]\label{rmk:levelNscale}
A scope statement and a model of operation counts and memory, in the notation of Remarks~\ref{rmk:levelcost} and \ref{rmk:heilbronncost}
($d_V:=n+1=k-1$, $D_N=m_Nd_V$, $h:=\dim H^{1}$); not an assertion about running times.
\begin{enumerate}
\item \textbf{Scope.} Every statement of this chapter --- Definitions~\ref{def:cosetrepsN}, \ref{def:levelNmodel}, \ref{def:kappaN},
  Lemmas~\ref{lem:treeres}--\ref{lem:merelpath}, Propositions~\ref{prop:heilbronn}, \ref{prop:kappamodel}, \ref{prop:degenpres} --- holds for \emph{every}
  $N\ge1$ and every even $n\ge0$ with $6\in\cring^{\times}$. ``$N\le10^{4}$ at weight~$2$'' is the reach of slot~15's runs (\ref{LN6}'s scan), not a
  hypothesis; higher weight $k>2$ at composite $N$ needs no new statement, only the costs below.
\item \textbf{Fixed spaces.} One $d_V\times d_V$ solve for each of $S$, $U$, then $O(d_V^{2})$ per fixed point (Remark~\ref{rmk:levelcost}(3)); at most
  $2^{\omega(N)+1}$ fixed points.
\item \textbf{The cokernel at $k>2$.} After the orbit-local quotient by $M^{\langle S\rangle}$ the elimination against $M^{\langle U\rangle}$ has $\le3$ blocks
  of size $d_V\times d_V$ per block column (Remark~\ref{rmk:levelcost}(4)); a block operation costs $O(d_V^{3})$. Worst case cubic in $D_N$, no better
  bound asserted; \ref{levelO4} stands: at large $k$ \emph{and} large $N$ the fill is dense $d_V\times d_V$ blocks.
\item \textbf{Memory of the normal form.} The normal-form table of $\mathrm{Rel}_N$ (Remark~\ref{rmk:heilbronncost}(3)) has $D_N\times h$ entries with
  $h\approx D_N/6$ (\ref{LN1}): about $m_N^{2}(k-1)^{2}/6$ words, i.e.\ $\frac23m_N^{2}(k-1)^{2}$ bytes in $32$-bit words --- $\approx67$~MB at
  $(m_N,k)=(10^{4},2)$, $\approx0.6$~GB at $(10^{4},4)$, $\approx6.7$~GB at $(10^{4},11)$ --- an upper bound for a dense table; a sparse normal form stores
  less (the shipped \code{SparseCokernel} of \file{Hecke/HeckeLevelNSparse}: \code{final\_nnz} entries plus the pivots, \code{HeckeLevelNSparseValidation}'s
  cost line --- the figure WP-13's RSS prediction should use). With the rule-19 prediction this is the RSS model for slot-16 WP-13's ladder.
  \emph{Caveat at composite $N$, $k>2$} (v1.17.0-pre.1; slot-16 WP-13's finding F-3, \file{docs/implementation/slot16/WP-13\_report.md}~\S4; a statement
  about the cost model only, no value depends on it): \COMP\textsuperscript{a} the fitted fill law of \code{CokernelCostModel} (fitted on prime $N$)
  \emph{underestimates} the measured \code{final\_nnz} at composite $N$ and $k>2$, by a factor of $1.47$ at $(N,k)=(30030,4)$, $5.70$ at $(4620,6)$ and
  $6.80$ at $(2310,12)$ (WP-13's ladder rungs \code{30030:2}, \code{4620:4}, \code{2310:10} in its $N{:}n$ notation); the dense bound above stays an upper
  bound there ($3.9$~GB against $2.4$~GB measured at $(2310,12)$). A prediction at composite $N$, $k>2$ therefore scales the measured \code{final\_nnz} of
  the last rung run by $m_N^{1.335}d_V^{2}$ rather than reading the fitted law. \emph{Sharpened at v1.17.0-pre.3} (slot-17 WP-5's finding W5-F1,
  \file{docs/implementation/slot17/WP-5\_report.md}~\S3): that scaling, too, fails at composite $N$ with a large prime factor --- at $N=99330=2\cdot3\cdot5\cdot7\cdot11\cdot43$,
  $k=2$, it predicted $42$~M non-zeros and $1.4$~GB, and the sparse presentation grew past $8$~GB, where the run was stopped by its watchdog (rule~19) before it
  finished; this is an open measurement, not a result (no \code{final\_nnz} was recorded), and it is recorded as a failed prediction. Along the primorial rungs
  $30030$, $60060$, $90090$ the measured-to-fitted fill ratio rose ($3.08$, $3.79$, $4.01$). No value depends on it.
  \emph{The structure-aware fill model} (v1.18.0-pre.1, contract first for slot-18 WP-5; \textbf{not built, not run} at this version; a statement about the
  cost model only). The fill of the sparse elimination is to be modelled through the structure the Shapiro presentation already has --- the factorisation of
  $N$ (the number $\omega(N)$ of prime factors and the largest one, $P(N)$), the orbit and block sizes of the permutation $\otimes\,\rhon$ form
  \eqref{eq:MH} (blocks of size $d_V$ permuted by the coset action), and $d_V=k-1$ --- rather than through $m_N$ alone; fitted on the logged \code{SPARSE\_CELL}
  and \code{DEGEN\_CELL} rungs of slot-17 WP-5 and WP-6 \emph{before} any new run (rule~20), and the same model replaces \code{CokernelCostModel}'s fitted law
  (slot-17 WP-6's finding F-1: final fill $4$--$5\times$ the law at $m_N\approx10^{4}$, $d(\chi)\ge5$). It replaces the rule-19 RSS hook (\code{hook\_rss/v1})
  only if it predicts every logged rung within a factor $1.5$ and from above; otherwise it is recorded as a finding and \code{hook\_rss/v1} stays. The two
  open measurements, $N=99330$ at $k=2$ ($m_N=D_N=304\,128$, $10\,991$~MiB predicted by \code{hook\_rss/v1}) and $N=90090$ at $k=4$ ($m_N=290\,304$,
  $D_N=870\,912$, $9622$~MiB predicted by \code{hook\_rss/v1}), both admitted by \code{HECKE\_LEVEL\_N\_HOOK\_RSS\_CAP\_GIB=12} (v1.18.0-pre.2, the slot-18 referee
  F1's m-C1: pre.1 gave the figures of $90090$ at $k=4$ only, attached to both cells), are each predicted by the model first and then run once in an exclusive memory window; their \code{final\_nnz} and peak RSS enter
  this item as \COMP\textsuperscript{a} at the pre-release that records them, a missed prediction as a failed prediction. A change of the elimination order
  (column order by block, as the permutation structure suggests) is admissible only as a choice that changes no value: the presentation and every Hecke
  matrix of this chapter must stay bit-identical on every logged rung (the canonical form is the natural reduced row echelon form in the shipped column order),
  checked as a gate; if the blow-up is not ordering-driven, nothing is changed and the finding is recorded.
  \emph{Outcome at v1.18.0-pre.2} (slot-18 WP-5, \file{docs/implementation/slot18/WP-5\_report.md}; predictions committed first,
  \file{docs/implementation/slot18/WP-5\_predictions.md}). \COMP\textsuperscript{a} \textbf{The structure-aware model failed its rule}: fitted on the $57$ logged
  cokernels of slots~15--17 ($\log\lambda$ against $\log m_N$, $\omega(N)$, $\log P(N)$ and $[d_V>1]$, $\lambda$ the fill per free column), it leaves an rms
  log-residual of $0.43$ and a maximum of $0.81$, beyond the factor~$1.5$; the fill is not even monotone in $P(N)$ (at $N=2310\,P$, $k=2$, a $3.8$-fold spread over
  $P=13,\dots,43$). So \code{hook\_rss/v1} and \code{CokernelCostModel}'s fitted law \textbf{stay}; this is a finding, recorded as the rule above prescribes, and
  no cost statement of this item is replaced. \emph{W5-F1's root cause is the order of the generators, not the fill}: \code{final\_nnz} at $N=99330$, $k=2$, is
  $54\,450\,319$ (computed by WP-5's scratch prototype, which reproduces nine logged \code{final\_nnz} exactly; close to $90090$'s $50.7$~M), and the slot-15--17
  consumption order (all $S$-generators, then all $ST$-generators) made the back-substitutions explode at composite $N$ with a large prime factor ($>16$~min and
  \code{bad\_alloc} at $5.8$~GB in the prototype), while the generators sorted by their largest column, descending, eliminate in $5.7$~s. The shipped
  \code{SparseCokernel} now consumes them in that order (rollback \code{HECKE\_LEVEL\_N\_GEN\_ORDER=as\_built}); the choice changes no value, the stored form being
  the unique natural reduced row echelon form, and the gate held: presentation hash and \code{final\_nnz} equal under both orders on all $21$ logged rungs (the
  column order is \emph{not} changed). At $d_V>1$ with $P(N)=13$ the sorted order costs up to $+16\,\%$ peak RSS ($60060$ at $k=4$); \code{hook\_rss/v1} stays an
  upper bound on every rung. The two window measurements above are \textbf{pending} (deferred in slot~18 behind the LAT37 runs and the disk emergency of
  2026-10-03); WP-5 predicted $54.45$~M non-zeros and $\approx1.9$~GB at $(99330,k=2)$, and $216$~M ($216$--$370$~M) and $\approx5.5$~GB ($5.5$--$9.6$~GB) at
  $(90090,k=4)$.
\item \textbf{Operators.} $H_{N,p}$ costs $O(\abs{\mathrm{Mer}_p}\,m_Nd_V^{2})$ per application, $D_{M,N,d}$ $O(\abs{\mathrm{Mer}_d}\,m_Md_V^{2})$ with the admissible classes
  enumerated per $\epsilon$ (at most $2m_N/m_d$ per $\epsilon$; scanning every $i$ costs $O(\abs{\mathrm{Mer}_d}\,m_N)$ checks, $\approx4\cdot10^{9}$ at $M=1$,
  $d=N\approx10^{4}$) (Proposition~\ref{prop:degenpres}(7)), and a Hecke matrix on $H^{1}$ $O(\abs{\mathrm{Mer}_p}\,d_V\,h^{2})$ (Remark~\ref{rmk:heilbronncost}(3)); every
  fibre map $O(m_Nd_V^{2})$. None forms a $D_N\times D_N$ object.
\item \textbf{The ring.} Over $\Z/\ell^{l}$, $\ell\ge5$, every count times one ring operation, and the elimination by Howell forms
  (Remark~\ref{rmk:levelscaling}, \ref{levelO6}); with a character, over $\F_q$ or $W_K(\F_q)$ (Remark~\ref{rmk:nebcost}).
\end{enumerate}
\end{remark}

\begin{remark}[What is checkable at level $N$, continued: (LN7)--(LN9)]\label{rmk:levelNcheckabletwo}
\PROVED{} \STD{\cite{atkin_lehner}, \cite{billerey_menares}} \CONJ{} \COMP\textsuperscript{a} Committed at v1.16.0-pre.1, before slot-16 WP-13's code (rule~20), with
the convention of Remark~\ref{rmk:hidacheckable}: a failing \PROVED{} item stops the package, a failing \CONJ{} item is a finding.
\begin{enumerate}[label=(LN\arabic*),start=7]
\item\label{LN7} \PROVED{} (Proposition~\ref{prop:kappamodel}) \textbf{Model $=$ presentation, entrywise.} At every tested $(N,n,p,\cring)$, for a basis of
  $\Zq{1}(K,\Rel;L_N)$ modulo $\Bq{1}$: the vectors $-f(e_{46})$ span $M_N/\mathrm{Rel}_N$ and are independent there; $(\widetilde{[A\,g\,A]}f)(e_{46})-
  H_{N,p}f(e_{46})\in\mathrm{Rel}_N$ for $p\nmid N$ and $p\mid N$; $\alpha^{0}$, $\omega$ (and $\langle d\rangle$ with a character) commute with
  $f\mapsto f(e_{46})$ \emph{exactly}. This replaces \ref{LN3}'s comparison of characteristic polynomials by a comparison of classes.
\item\label{LN8} \PROVED{} (Proposition~\ref{prop:degenpres}) \textbf{$\alpha_d$ on the presentation.} $D_{M,N,d}(\mathrm{Rel}_M)\subseteq\mathrm{Rel}_N$;
  $(\widetilde{[A\,g\,B]}f)(e_{46})-D_{M,N,d}f(e_{46})\in\mathrm{Rel}_N$ for $f\in\Zq{1}(K,\Rel;L_M)$ (the model's $\alpha_d$, Proposition~\ref{prop:levelcorrespondence}(4)(c),
  against the presentation's); $d=1$ is $\alpha^{0}$ entrywise; the trace (5) and the composite (4); among WP-13's cells at least one with
  $d=p\mid M$, where (5) gives $U_p$ (proved, so far untested; WP-14 referee). \STD{\cite{atkin_lehner}; not verified (rule~\ref{A2})}
  The old-space relations $U_p\alpha_1=p\,\alpha_p$, $U_p\alpha_p=\alpha_pT_p-p^{\,n}\alpha_1$ at $M\to Mp$, $p\nmid M$ (the analogue of Proposition~\ref{prop:oldspace}(1),(3) with $\Ga_0(M)\to\Ga_0(Mp)$; \STD{\cite{atkin_lehner}, not verified};
  \COMP{} in Remark~\ref{rmk:degenpresconv}(3)), and Hecke equivariance $T_p\alpha_d=\alpha_dT_p$ for $p\nmid N$. \COMP\textsuperscript{a}: Remark~\ref{rmk:degenpresconv}(3).
\item\label{LN9} \STD{\cite{billerey_menares}; \textbf{not} verified (rule~\ref{A2})} \CONJ{} \textbf{$\mathrm{R3}_k$: the Eisenstein reading at weight $k\ge4$.}
  For $N$ prime, $p\ge5$ prime, $p\nmid N$, $p>k+1$, $k=n+2\ge4$ even, let $W^{(k)}_{N,p}\subseteq H^{1}(\Ga_0(N);\Sym^{n}\F_p^{2})$ be the joint generalised
  eigenspace of $(T_q)_{q\le7,\ q\nmid Np}$ at $1+q^{k-1}$ (the code's reading: $\dim W^{(k)}_{N,p}$, the Loewy length and radical layers of the algebra
  generated by those $T_q$ on it, reported beside $v_p(N^{k}-1)$, $v_p(N^{k-2}-1)$ and $v_p(B_k/k)$). Prediction: $\dim W^{(k)}_{N,p}\ge2$ (the two cusps), and
  $\dim W^{(k)}_{N,p}\ge3$ \emph{only if} $p$ divides $(N^{k}-1)(N^{k-2}-1)\,\mathrm{num}(B_k/k)$. Reason (not a proof here; it uses that the mod-$p$ joint generalised eigenspace counts characteristic-$0$ eigensystems with multiplicity, i.e.\ that
  $H^{1}(\Ga_0(N);\Sym^{k-2}\Z_p^{2})$ is torsion-free, which holds since $H^{0}(\Ga_0(N);\Sym^{k-2}\F_p^{2})=0$ for $0<k-2\le p-3$ --- exactly the
  hypothesis $p>k+1$, which is why it is the right one; WP-14 referee C-4): a cuspidal eigensystem congruent to
  the Eisenstein one is old from level~$1$ --- possible only if $p\mid\mathrm{num}(B_k/k)$ (Herbrand--Ribet, $p>k+1$) --- or new of level $N$, where the
  Steinberg shape of $\rho_f|_{G_{\Q_N}}$ with $a_N=\pm N^{k/2-1}$ forces the Frobenius eigenvalues $\{1,N^{k-1}\}$ of $1\oplus\bar\chi_p^{k-1}$ to be
  $\{\pm N^{k/2},\pm N^{k/2-1}\}$ modulo $p$, i.e.\ $N^{k}\equiv1$ or $N^{k-2}\equiv1$. The cited source treats the converse; it is not stated here. At $k=2$,
  \ref{LN4}--\ref{LN6} apply instead ($N^{k-2}-1=0$). \textbf{Reported, not gated} (slot-16 plan, WP-13).
\end{enumerate}
\end{remark}

\paragraph{Scope: what this chapter does not do.}
It asserts, for every $N$, the coset table, the tree resolution, Manin's sequence, the cokernel isomorphism and its duality, and
the Hecke action of $T_p$ ($p\nmid N$) and $U_p$ ($p\mid N$) on the cokernel presentation, and --- since v1.16.0-pre.1,
\S\ref{sec:levelNdegsec} --- the model-to-presentation map with the fibre maps carried by naturality (Proposition~\ref{prop:kappamodel}) and the
degeneracy maps $\alpha_d$, $d\ge1$, on the presentation (Proposition~\ref{prop:degenpres}), all \PROVED{} for $6\in\cring^{\times}$, $n$
even. Not asserted: anything at $\ell\in\{2,3\}$ beyond the cochain level of Definition~\ref{def:levelNmodel}(a) (the tree
resolution and the cokernel need $6\in\cring^{\times}$), except, since v1.21.0-pre.4, the weight-two curve object of
Proposition~\ref{prop:levelNcurve} (which is not group cohomology there); nebentypus ($\Ga_1(N)$; since v1.15.0-pre.1
\S\ref{sec:nebentypus}); composite $n$ in
Proposition~\ref{prop:heilbronn}'s bookkeeping (Remark~\ref{rmk:heilbronnconv}(4)); any bound on the elimination better than cubic;
\ref{LN4}--\ref{LN6}, \ref{LN9}, which are literature and predictions. (v1.16.0-pre.1: the sentence v1.15.0-pre.3 placed here --- a
$\Ga$-module map between induced modules carries $\mathrm{Rel}$ to $\mathrm{Rel}$ and induces $H^{1}(\Ga;\phi)$, so the fibre maps $\alpha^{0}$, $\omega$ act on the
presentation, while $\alpha_d$, $d>1$, was asserted on the model only --- was released \emph{unrefereed} at v1.15.0; it is superseded by
Proposition~\ref{prop:kappamodel}(3), which proves it, and Proposition~\ref{prop:degenpres}, which puts $\alpha_d$ on the presentation.)

\paragraph{Code map.}
\textbf{Built at v1.14.2} (slot-14 WP-12, \texttt{7a5158d}; not built at v1.14.0, the release, and at v1.14.1 --- contract first,
rule~21; refereed by slot-14 WP-5, ACCEPT WITH REPAIRS, repaired at pre.2). \file{Hecke/HeckeLevelN} (\code{namespace
hecke::level}): Definition~\ref{def:cosetrepsN}(1)--(5) as \code{Gamma0LevelN}, a \code{CosetLevel} for every $N$ (at
$N=p^{\nu}$ the shipped table entrywise, Lemma~\ref{lem:cosetrepsN}(3)), item~(6) as \code{Gamma0SubTransversalN};
Definition~\ref{def:levelNmodel} as \code{InducedRelationsLevel} on any \code{CosetLevel} over any word-sized ring (the
monomial $M_N(g_r)$ of \eqref{eq:MH}; \file{Hecke/InducedLocalSystem} is unchanged and stays at prime powers);
Proposition~\ref{prop:levelorbit}(1),(2) orbit-wise (\code{OrbitWiseFixedSpace}, \code{CokernelFromFixedSpaces}); the
operators by \eqref{eq:levelcorr} on the level-$1$ $\tau$-table (\file{Hecke/HeckeLevelEll}, unchanged), the $r_t$ of item~(6).
\code{HeckeLevelNValidation}: \ref{LN1}, \ref{LN2}, the $\tau$-route half of \ref{LN3} (against an independent group-cocycle
route at composite $N$ and the classical characteristic polynomials), \ref{LN4}--\ref{LN6} on the $26$ scan cells $N\le211$
and, under \code{HECKE\_LEVEL\_N\_R3\_CELLS}, the named cells $(1321,11)$, $(1571,5)$. \textbf{Built at v1.15.0-pre.2}
(slot-15 WP-7): $H_{N,p}$ of \eqref{eq:heilbronn}, \file{Hecke/HeckeLevelNHeilbronn} (\code{HeilbronnOperator}), tested by
\code{HeckeLevelNHeilbronnValidation} (\ref{LN3} on both halves). \textbf{Built at v1.15.0-pre.3} (slot-15 WP-8, WP-9): the sparse
cokernel elimination of Remark~\ref{rmk:levelcost}(4) (\file{Hecke/HeckeLevelNSparse}: \code{SparseCokernel}, entrywise equal to the dense
route, the cost model and its route choice, the fibre maps $\alpha^{0}$, $\omega$ on the presentation), the degeneracy maps $\alpha_d$ at
composite $N$ on the model (Proposition~\ref{prop:levelcorrespondence}(4)(c)), \code{HeckeLevelNSparseValidation}; \ref{LN6} at scale over the
Merel candidates $N\le10^{4}$ (\file{Hecke/HeckeLevelNR3}, \code{HeckeLevelNValidation} arm (R3s)). \textbf{Built at v1.16.0} (slot-16 WP-13, \file{docs/implementation/slot16/WP-13\_report.md}; cross-model audit
\file{docs/implementation/slot16/WP-13\_audit.md}, ACCEPT WITH REPAIRS, repaired by WP-17; not built at v1.16.0-pre.1 --- contract first, rule~21):
\S\ref{sec:levelNdegsec} --- \file{Hecke/HeckeLevelNDegen}: $\kappa^{\mathrm{mod}}_N$ of Proposition~\ref{prop:kappamodel} (\code{KappaModelPoints}, which finds $e_{46}$ by its
vertex set and checks its orientation against the model's coboundary, \code{KappaModel}, \code{KappaModel0}) and $D_{M,N,d}$ of \eqref{eq:degenpres} on the sparse
presentation (\code{DegeneracyOperator}, the admissible classes enumerated per $\epsilon$ as in Proposition~\ref{prop:degenpres}(7) and, as a check, by a full scan);
the $\mathrm{R3}_k$ reading \ref{LN9} (\code{ReadEisensteinR3k} in \file{Hecke/HeckeLevelNSparse}); arms of \code{HeckeLevelNSparseValidation} (no new name):
\ref{LN7} on $17$ cells, \ref{LN8} against the model on $40$ cells (the $28$ cells of \file{docs/implementation/slot16/wp12/degen.py} and $12$ new, among them
$d=p\mid M$) with $4$ negative-control cells, the trace (5) on $20$ cells including $U_p$ at $d=p\mid M$ on $12$, $\alpha_d$ at scale on $4$ cells; \ref{LN9}
reported and not gated ($145$ readings, no finding); weight $k>2$ at composite $N$ on the sparse route up to $N=30030$ at $k=4$
(Remark~\ref{rmk:levelNscale}; the ladder of the report).
\textbf{Built at v1.17.0-pre.3} (slot-17 WP-5, \file{docs/implementation/slot17/WP-5\_report.md}, merged; cross-model audit
\file{docs/implementation/slot17/WP-5\_audit.md}, ACCEPT WITH REPAIRS, text only; contract first at v1.17.0-pre.1, rule~21; slot-16 WP-13's finding F-1):
\textbf{the reach of the table at composite $N$.} Definition~\ref{def:cosetrepsN} and its stored table are unchanged (entries $O(N^{2})$, in \code{int64}). Every
coset product $h=XYZ$ behind a block $\rhon(h)$ of \eqref{eq:MH} (\code{MonomialOf}, the orbit-wise fixed spaces and generators of
Proposition~\ref{prop:levelorbit}(1), the level-change maps of Proposition~\ref{prop:levelcorrespondence}(4)(b)) is formed in \code{int64} first and, only on an
overflow, in \code{\_\_int128}, with $h\in\Ga_0(N)$ and $\det h=1$ checked exactly over $\Z$ and the block keyed on $h\bmod\ell^{k}$
(\code{SubgroupProductForRho}; attempt-then-escalate as in Theorem~\ref{thm:backend}). Hence the table and every Hecke matrix of the chapter are
bit-identical wherever the old width completed (for $N\le30030$ the escalation never fires: \code{HeckeLevelNSparseValidation} arm (W5b), a gate), and the
reach is $N<2^{31}$ for the products. \COMP\textsuperscript{a} on the ladder: $N=60060$ and $90090$ at $k=2$, $N=60060$ at $k=4$ ($D_N=580\,608$), \ref{LN3},
\ref{LN8} (for $\alpha_2\colon30030\to60060$, $\alpha_3\colon30030\to90090$) and the trace (5) held; the old width's refusal at $N=60060$ is the negative control
(W5n). The expected \emph{verdicts} of WP-5's committed predictions (\file{docs/implementation/slot17/WP-5\_predictions.md}) held; its \emph{numeric} predictions
did not all hold: \code{final\_nnz} was under-predicted by up to $30\,\%$ ($90090$, $k=2$: $50.7$~M against $39$~M) and over-predicted at $60060$, $k=4$
($118.8$~M against $165$~M), and the $N=99330$ prediction failed (Remark~\ref{rmk:levelNscale}(4)).
\textbf{Built at v1.18.0-pre.2} (slot-18 WP-5, \file{docs/implementation/slot18/WP-5\_report.md}, merged): the sparse elimination consumes the generators of
$\mathrm{Rel}_N$ sorted by their largest column, descending (\code{GeneratorOrder::kLastColumnDescending}, stable; rollback
\code{HECKE\_LEVEL\_N\_GEN\_ORDER=as\_built}) --- a value-free choice, because the stored form is the unique natural reduced row echelon form (all $21$ logged
rungs bit-identical under both orders, a gate; arm (W18o) repeats the comparison entry for entry on $12$ cells in \code{--all}, and the Hecke matrices on $4$).
Finding W5-F1 is that order's cost: the fill at $N=99330$, $k=2$, is $54.45$~M, not a fill-law excess. The structure-aware fill model of
Remark~\ref{rmk:levelNscale}(4) \textbf{fails} its factor-$1.5$ rule (the fill is not monotone in $P(N)$), so \code{hook\_rss/v1} stays; this is a finding (the
remark's outcome paragraph). \code{Gamma1LevelN::contains\_wide} (exact in \code{\_\_int128}) and \code{HeilbronnChi} (its product formed by
\code{Gamma0LevelN::cocycle\_wide}) are on the \code{int64}-then-\code{\_\_int128} route, arm (W18g). Item~(c), the two window measurements, is \textbf{pending}.
The plan, as committed at v1.18.0-pre.1 (contract first, rule~21; ``not built'' at that version; kept for the record;
no statement of this chapter changes): (a)~the structure-aware fill model of Remark~\ref{rmk:levelNscale}(4), with its replacement rule for the RSS hook and
for \code{CokernelCostModel}'s law; (b)~at most one value-preserving change of the elimination order, merged only if every logged slot-16 / slot-17 rung stays
bit-identical (a gate, not a report); (c)~the measurements $N=99330$ ($k=2$) and $N=90090$ ($k=4$) in an exclusive window, each predicted by~(a) first, under a
VmHWM watchdog; (d)~\code{Gamma1LevelN::contains\_wide} and \code{HeilbronnChi}'s products on the \code{int64}-then-\code{\_\_int128} route of the v1.17.0-pre.3
paragraph above (the slot-17 WP-6 \code{\_\_int128} fallback becomes the one route), bit-identical wherever \code{int64} completes, checked by forcing the wide route
on cells where both run. Arms of \code{HeckeLevelNSparseValidation} only (its \code{SPARSE\_CELL} and \code{DEGEN\_CELL} arms; no new name; the slot-18 plan's
``\code{HeckeLevelNDegenValidation}'' is the translation unit of those arms, \file{Hecke/HeckeLevelNDegenValidation}, not a registered name); its \code{EFFECTIVE} lines name the RSS model in use; golden unchanged. The Moerdijk--Pronk line: none of (a)--(d) uses the $7$-vertex model beyond \S\ref{sec:induced}'s parametric statements.

\paragraph{Formalization note.}
Lemma~\ref{lem:manin} needs the Farey complex and a filtration argument; Lemma~\ref{lem:merelpath} is elementary lattice
combinatorics; Lemma~\ref{lem:levelduality} needs projective resolutions of length $1$, Shapiro's lemma for homology and the map of a
compatible pair; Proposition~\ref{prop:heilbronn} adds only coset counting. None depends on the theory of modular forms;
\ref{LN1}, \ref{LN4}--\ref{LN6} do.

% v1.21.0-pre.4 (slot 22, LN23; owner decision O6): prop:levelNcurve and rmk:levelNcurvefamily as delivered by the F4 referee
% (docs/implementation/slot22/LN23_C_proposition.tex.txt, LN23_F4_referee_report.md Sec. 11), placed as the last subsection of sec:levelN
% so that no statement number of v1.21.0-pre.3 moves.
\subsection{The weight-two presentation at every prime: the curve object}\label{sec:levelNcurve}

\begin{proposition}[The $n=0$ cokernel presentation at every prime is the cohomology of the open modular curve]\label{prop:levelNcurve}
\PROVED{} Let $N\ge1$, $A:=\Ga_0(N)$, $E:=A\backslash\Ga\cong\mathbb{P}^1(\Z/N)$ with the right action of $\Ga$
(Definition~\ref{def:cosetrepsN}), $D:=\Z[E]=M_N$ at $n=0$ (Definition~\ref{def:levelNmodel}), and, for a finite
subgroup $F\le\Ga$, $D^{F}$ the fixed submodule, which is the $\Z$-span of the $F$-orbit sums
$\sum_{x\in O}e_x$. Let $\mathcal T$ be the Bass--Serre tree of Claim~\ref{cl:amalgam}(1) (vertices
$\Ga/\langle S\rangle\sqcup\Ga/\langle U\rangle$, edges $\Ga/\langle-I\rangle$, oriented from
$\gamma\langle S\rangle$ to $\gamma\langle U\rangle$), $U=ST$, and $G:=A\backslash\mathcal T$: a finite graph with edge set
$E$ ($-I\in A$ is the edge stabiliser), vertex set the $\langle S\rangle$-orbits $V_2$ and the
$\langle U\rangle$-orbits $V_3$ on $E$, the edge $x$ joining the two orbits containing it. $A$ acts without
inversions (the orientation is preserved), so $G$ is a graph and
$\Hom_{A}(C_{\bullet}(\mathcal T),X)=C^{\bullet}_{\mathrm{cell}}(G;X)$ for every abelian group $X$ with trivial
action. Let $Y_0(N):=A\backslash\mathfrak H$ and $X_0(N)$ its compactification, both as coarse spaces.
\begin{enumerate}
\item \textbf{(Free, and what it is.)} $D/(D^{\langle S\rangle}+D^{\langle U\rangle})=H^{1}(G;\Z)\cong
  H^{1}(Y_0(N);\Z)=\Hom(\Ga_0(N),\Z)$, a free $\Z$-module of rank
  \[
    b_1(G)\;=\;\abs{E}-\abs{V_2}-\abs{V_3}+1\;=\;1+\tfrac{m_N}{6}-\tfrac{\nu_2(N)}{2}-\tfrac{2\nu_3(N)}{3}\;=\;2g(X_0(N))+c(N)-1
  \]
  ($\nu_2,\nu_3$ as in Remark~\ref{rmk:levelcost}(3), $c(N)$ the number of cusps, $g$ the genus); and
  $H^{1}(Y_0(N);\Z)\cong H_1(X_0(N),\mathrm{cusps};\Z)$ by Poincar\'e--Lefschetz duality. The quotient by the
  smaller lattice spanned by $e_x+e_{x\cdot S}$ and $e_x+e_{x\cdot U}+e_{x\cdot U^{2}}$ alone has torsion whenever
  $\nu_2(N)+\nu_3(N)>0$ (\COMP\textsuperscript{a}: $\Z/2$ at $N=5,17,41$, $\Z/3$ at $N=43$, $\Z/6$ at
  $N=13,37,73$); the fixed points contribute $e_x$ itself, not $2e_x$ or $3e_x$.
\item \textbf{(Every characteristic.)} For every prime $p$ and $M:=\F_p[E]$, with $M^{\langle S\rangle}$,
  $M^{\langle U\rangle}$ the fixed subspaces (spanned by the orbit sums --- no $\tfrac12$, no $\tfrac13$),
  \[
    M/(M^{\langle S\rangle}+M^{\langle U\rangle})\;=\;H^{1}(G;\F_p)\;=\;H^{1}(Y_0(N);\F_p)\;\cong\;
    \bigl(D/(D^{\langle S\rangle}+D^{\langle U\rangle})\bigr)\otimes\F_p ,
  \]
  of dimension $2g+c-1$ for every $p$, and $M^{\langle S\rangle}\cap M^{\langle U\rangle}=H^{0}(G;\F_p)=\F_p$.
  For $p\ge5$ this is $H^{1}(\Ga_0(N);\F_p)$ of Proposition~\ref{prop:levelorbit}(2) (the same cochain complex);
  for $p\in\{2,3\}$ it is \textbf{not} group cohomology ($\dim H^{1}(\Ga_0(N);\F_3)=2g+c-1+\nu_3$,
  $\dim H^{1}(\Ga_0(N);\F_2)=2g+c-1+\nu_2+\delta_{\nu_2,0}$), but the Bredon object of hypothesis~\ref{H3} and
  Proposition~\ref{prop:bredon}(3) with the permutation coefficients $\Z[E]$: the cohomology of the coarse
  quotient $Y_0(N)$.
\item \textbf{(Hecke.)} For a prime $q\nmid N$ the operator $H_{N,q}$ of \eqref{eq:heilbronn} at $n=0$ (every block
  $\rhon=1$: $(H_{N,q}x)_i=\sum_{(\epsilon,j):\,j\cdot\epsilon=i}x_j$) is defined over $\Z$, preserves
  $D^{\langle S\rangle}+D^{\langle U\rangle}$, and induces on $H^{1}(Y_0(N);\Z)$ the Hecke operator $T_q$, i.e.\ the
  endomorphism of the lattice $H^{1}(Y_0(N);\Z)\subset H^{1}(Y_0(N);\Q)$ that Proposition~\ref{prop:heilbronn}(2)
  identifies over $\Q$; its reduction is $T_q$ on (2) for every $p$, $p\in\{2,3\}$ included. Under the duality of (1)
  it is $T_q$ on $H_1(X_0(N),\mathrm{cusps};\Z)$: for $q\nmid N$ the Hecke correspondence equals its transpose, and
  $w_N$ commutes with $T_q$.
\end{enumerate}
\end{proposition}

\begin{proof}
(1) \emph{The quotient graph.} The edge $\gamma\langle-I\rangle$ has end points $\gamma\langle S\rangle$ and
$\gamma\langle U\rangle$; its $A$-orbit is the coset $A\gamma\in E$, and the $A$-orbits of vertices are the
$\langle S\rangle$- and $\langle U\rangle$-orbits on $E$ ($A\backslash\Ga/\langle S\rangle$, $A\backslash\Ga/\langle U\rangle$).
A vertex of $G$ coming from a fixed point $x$ (i.e.\ $w_xSw_x^{-1}\in A$, an elliptic point of $A$) has the single edge
$x$ at it; a free orbit gives a vertex of valence $2$ (resp.\ $3$). The cellular coboundary $\delta\colon\Z^{V}\to\Z^{E}$
sends $1_v$ to $\pm$ the indicator of the orbit $v$ (sign $-$ on $V_2$, $+$ on $V_3$): \emph{once} per edge, since
each edge has one tail and one head in $G$. So $\mathrm{im}\,\delta=D^{\langle S\rangle}+D^{\langle U\rangle}$ and
$H^{1}(G;\Z)=\operatorname{coker}\delta=D/(D^{\langle S\rangle}+D^{\langle U\rangle})$. (By contrast the chain map
$\partial^{t}$ of Lemma~\ref{lem:manin} on \emph{coinvariants} counts a fixed edge $2$ or $3$ times: that is the
torsion of the last sentence of (1).) $H^{1}$ of a finite connected graph is free of rank $\abs{E}-\abs{V}+1$;
$\abs{V_2}=(m_N+\nu_2)/2$, $\abs{V_3}=(m_N+2\nu_3)/3$, and $12g=12+m_N-3\nu_2-4\nu_3-6c$ (the Hurwitz formula for
$X_0(N)\to X_0(1)$) give the displayed rank. Alternatively, the incidence matrix of a graph is totally
unimodular, so every invariant factor of $\delta$ is $1$ and $\operatorname{coker}\delta$ is free: the rank over
$\F_p$ equals the rank over $\Q$ for every $p$ (\COMP\textsuperscript{a}: an integer elimination with $\pm1$ pivots
only on every prime $N\le3000$ tested, \file{docs/implementation/slot22/ln23ref/}).

\emph{$G$ and $Y_0(N)$.} $\mathcal T$ embeds in $\mathfrak H$ as the union of the geodesic arcs $\gamma[i,\rho]$
($\gamma\in\Ga$, $\rho=e^{2\pi i/3}$), the dual tree of the Farey tessellation; $\mathfrak H\setminus\mathcal T$ is the
disjoint union, over the cusps $c\in\mathbb{P}^1(\Q)$, of open discs $R_c$ bounded by $\mathcal T$, and flowing each
$R_c$ along the geodesics ending at $c$ is a $\Ga$-equivariant strong deformation retraction of $\mathfrak H$ onto
$\mathcal T$ \STD{\cite[Ch.~I, \S4.2]{serre_trees} for the tree and its embedding; the retraction is the classical
\emph{spine} statement; locator \textbf{not} verified (rule~\ref{A2})}. It descends to a deformation retraction of
$Y_0(N)=A\backslash\mathfrak H$ onto $G=A\backslash\mathcal T$, so $H^{\bullet}(Y_0(N);X)=H^{\bullet}(G;X)$ for every
coefficient group $X$. (A reader who prefers not to rely on the retraction may use Proposition~\ref{prop:bredon}(2)
instead: $C_{\bullet}(\Xbulk)$ and $C_{\bullet}(\mathcal T)$ are chain homotopy equivalent over $\Z[\Ga]$, hence so are
$\Hom_{A}(C_{\bullet}(\Xbulk),X)$ and $\Hom_{A}(C_{\bullet}(\mathcal T),X)=C^{\bullet}(G;X)$, and
$\Hom_{A}(C_{\bullet}(\Xbulk),X)$ is the cellular cochain complex of $A\backslash\Xbulk$, which is $Y_0(N)$ with a
small open disc removed at every cusp, a homotopy equivalence.) $H^{1}(Y_0(N);\Z)=\Hom(\pi_1(Y_0(N)),\Z)=
\Hom(A,\Z)$ because $\pi_1$ of the coarse quotient is $A$ modulo the normal closure of its elliptic elements
(torsion, invisible to $\Hom(-,\Z)$). Finally $X_0(N)$ is a compact orientable surface and $Y_0(N)=X_0(N)\setminus
\mathrm{cusps}$, so Poincar\'e--Lefschetz duality gives $H^{1}(Y_0(N);\Z)\cong H_1(X_0(N),\mathrm{cusps};\Z)$, free of
rank $2g+c-1$ by the long exact sequence of the pair.

(2) The cokernel of a map of free $\Z$-modules commutes with $\otimes\F_p$ (right exactness), and the fixed space of a
permutation module over any ring is the span of the orbit sums (a fixed vector is constant on orbits); so the
$\F_p$-presentation is $\delta\otimes\F_p$ and its cokernel is $H^{1}(G;\F_p)=H^{1}(G;\Z)\otimes\F_p$ ($H^{2}(G)=0$).
$H^{0}(G;\F_p)=\F_p$ as $G$ is connected ($\mathcal T$ is). For $p\ge5$, $C^{\bullet}(G;\F_p)=\Hom_A(C_\bullet(\mathcal
T),\F_p)$ computes $H^{\bullet}(A;\F_p)$ by Lemma~\ref{lem:treeres}, and this is the identification of
Proposition~\ref{prop:levelorbit}(2) and Definition~\ref{def:kappaN} (the same complex). The dimensions of
$H^{1}(A;\F_p)=\Hom(A,\F_p)$ at $p\in\{2,3\}$ follow from $A^{\mathrm{ab}}$: the elliptic elements of order $3$
(resp.\ $4$) modulo $\pm I$ contribute the extra classes (\COMP\textsuperscript{a}: arm (IL3N) of
\file{Integral/HeckeIsotropyModelEllOpenHalves}, Proposition~\ref{prop:ellthreeclosed}).

(3) Proposition~\ref{prop:heilbronn}(1),(2) hold for the ring $\cring=\Z[1/6]$: $H_{N,q}$ preserves
$\mathrm{Rel}_{\Z[1/6]}=(D^{\langle S\rangle}+D^{\langle U\rangle})[1/6]$ and induces $[A\,g\,A]=T_q$ on
$H^{1}(A;\Z[1/6])=H^{1}(G;\Z[1/6])$. Since $D/(D^{\langle S\rangle}+D^{\langle U\rangle})$ is torsion-free by (1),
$(D^{\langle S\rangle}+D^{\langle U\rangle})[1/6]\cap D=D^{\langle S\rangle}+D^{\langle U\rangle}$; hence
$H_{N,q}(D^{\langle S\rangle}+D^{\langle U\rangle})\subseteq D^{\langle S\rangle}+D^{\langle U\rangle}$ and the induced
endomorphism of the lattice $H^{1}(G;\Z)$ is the restriction of $T_q$ on $H^{1}(G;\Q)$. Reduction mod $p$ commutes with
the induced map by right exactness, which gives (2)'s $T_q$ for every $p$. For the duality: the transpose of the
correspondence $T_q=\pi_{2\ast}\pi_1^{\ast}$ on $X_0(N)$ ($q\nmid N$) is $\pi_{1\ast}\pi_2^{\ast}=T_q$ itself, because
$\pi_2=\pi_1\circ w_q$ with $w_q$ an involution of $X_0(Nq)$; and $w_N$ commutes with $T_q$ for $q\nmid N$.
\end{proof}

\begin{remark}[Reading Mazur's Eisenstein ideal on the curve object; the Hecke family]\label{rmk:levelNcurvefamily}
\STD{\cite[II, \S9 (9.7), \S16 Def.\ p.~124, Prop.~(16.6), Cor.~(16.2), \S17 Prop.~(17.1)]{mazur}, read on the journal
scan (slot-22 F4); \cite[\S1 (Thm.~1.1), \S6]{calegari_emerton} (LIT-30, held)} \CONJ{} Let $N$ be prime and
$p\mid\mathrm{num}\bigl(\tfrac{N-1}{12}\bigr)$ --- at $p=2$: $N\equiv1\pmod8$, exactly iff $N\equiv9\pmod{16}$; at $p=3$:
$N\equiv1\pmod9$, exactly iff $N\equiv10,19\pmod{27}$ --- so that Mazur's Eisenstein maximal ideal
$\mathfrak m=(p,\,T_q-1-q\ (q\ne N),\,T_N-1)$ of the Hecke algebra $\mathbb T$ of $M_2(\Ga_0(N))$ has
$\mathbb T^{0}_{\mathfrak m}\ne0$. On the curve object $H:=H_1(X_0(N),\mathrm{cusps};\Z)$ of
Proposition~\ref{prop:levelNcurve}, $H\otimes\Q\cong\mathbb T_\Q\oplus\mathbb T^{0}_\Q$ as a Hecke module, so the
$\mathfrak m$-generalised eigenspace $W_{\mathfrak m}\subseteq H\otimes\F_p$ has $\dim W_{\mathfrak m}=2e_p+1$ when
$p\,\|\,\mathrm{num}$ ($e_p+1=\rank_{\Z_p}\mathbb T_{\mathfrak m}$, Calegari--Emerton's $e_p$) and $2g_p+1$ in general;
with $\mathbb T_{\mathfrak m}/p\cong\F_p[x]/x^{e_p+1}$ the Loewy length of $W_{\mathfrak m}$ is $e_p+1$ iff
$\mathbb T_{\mathfrak m}/p$ acts faithfully (\ref{LN4}, \ref{LN6} extended to $p\in\{2,3\}$). What the code reads
is the joint generalised eigenspace $W_{\mathcal F}$ and the radical series of the ideal generated by
$\eta_q:=T_q-1-q$, $q\in\mathcal F$, for a finite family $\mathcal F$ of primes. Two distinct things can fail, and
only the second is a matter of a bound:
\begin{enumerate}
\item \textbf{Generation.} The radical series of $(\eta_q:q\in\mathcal F)$ on $W_{\mathfrak m}$ has length $e_p+1$ iff
  the ideal is $\mathfrak m\mathbb T_{\mathfrak m}$, iff some $\eta_q\notin\mathfrak m^{2}$ (membership in $\mathfrak m^{2}$ is read modulo $p$ throughout, i.e.\ in $\mathbb T_{\mathfrak m}/p$). By Mazur this holds \emph{if}
  $\mathcal F$ contains a \emph{good} prime: $\ell\ne N$ with $\ell$ not a $p$-th power modulo $N$ and
  $\tfrac{\ell-1}{2}\not\equiv0\pmod p$ (for $p=2$: $\ell\equiv3\pmod4$ a non-residue; $\ell=2$ is \emph{excluded} by
  Prop.~(17.1), although Mazur's definition admits it when $2$ is not a quartic residue). For $p$ odd and $\ell=2$,
  $\tfrac{\ell-1}{2}=\tfrac12$ is a $p$-adic unit, so $2$ is good iff it is not a $p$-th power modulo $N$ (at $(1747,3)$,
  $\{2\}$ alone reads Loewy $4=e_3+1$). The condition is \emph{sufficient, not necessary}, at $p=2$: at $(2473,2)$
  ($e_2=1$, $2$ a quartic residue modulo $2473$, no good prime $\le13$) $\eta_2$ generates $\mathfrak m\mathbb T_{\mathfrak m}/2$,
  while $\eta_3,\eta_5,\eta_7,\eta_{11}$ do not (\COMP); at $(41,2)$ ($2$ not a quartic residue) $\eta_2\in\mathfrak m^{2}$, which is why Prop.~(17.1) excludes $\ell=2$ although Mazur's definition admits it.
  A family without a good prime may therefore read the Loewy length too short (and $\dim W$ right): $(2137,2)$ with $q\le13$ reads $15$ / Loewy $2$; the smallest
  good prime is $23$, and $\{23\}$ alone reads $15$ / Loewy $8$ $=e_2+1$. At $(41,2)$ the good primes $3$ or $7$ alone
  give $7,5,3,1$; $\eta_2$ and $\eta_5$ (not good) give $7,3$.
\item \textbf{Separation.} $W_{\mathcal F}\supseteq W_{\mathfrak m}$, with equality iff no other maximal ideal
  $\mathfrak m'\ni p$ of $\mathbb T$ contains every $\eta_q$, $q\in\mathcal F$. A family that does not separate reads
  $\dim W$ too large and a non-curvilinear series ($(2953,2)$ with $q\le13$: $39$ / $39,33,23,13,8,4$; with $19$ added:
  $7$ / $7,5,3,1$). Every maximal ideal of $\mathbb T^{0}/p$ is determined by the eigenvalues of $T_n$, $n\le
  B:=\lfloor(N+1)/6\rfloor$, because $\mathbb T^{0}=\Hom_\Z(S_2(\Ga_0(N);\Z),\Z)$ under $(t,f)\mapsto a_1(tf)$ and a
  cusp form with $a_n\equiv0\pmod p$ for $n\le B$ vanishes mod $p$ \STD{Sturm, \cite[Thm.~1, p.~276]{sturm}, held
  (LIT-13), every residue characteristic}; so $\mathcal F=\{q\le B\}$ separates, and also contains a good prime when
  $e_p\ge2$ (the $T_n$, $n\le B$, span $\mathbb T^{0}_{\mathfrak m}/p$, and $T_n-\sigma(n)\in\mathfrak m^{2}$ for all
  $n\le B$ would put that span inside $\F_p+\mathfrak m^{2}$). In practice a few primes beyond $13$ suffice; the
  Sturm bound is a proof device, not a cost.
\end{enumerate}
A reading with a family containing a good prime and reading a curvilinear series $2e+1,2e-1,\ldots,1$ of length
$e+1$ is consistent with $e_p=e$; it is a statement only with a separating family. \COMP\textsuperscript{a}: on every
cell of \code{HeckeLevelNCurveValidation}'s default grids ($N\le2000$ at $p=2$, $N\le2500$ at $p=3$) the family
$q\le13$ agrees with Calegari--Emerton; on the Thm~1.1 grid to $3\cdot10^{4}$ it fails at $28$ of $397$ cells in exactly
the two ways above, all repaired by $q\le100$ (\file{docs/implementation/slot22/LN23\_report.md} \S5, the F4 report).
The heavy oracle of that validation, $(21929,2)$ with $e_2=127$ ($\dim W=255$, Loewy length $128$, $h(-4\cdot21929)=256$), is Calegari--Emerton's own example
\cite[Thm.~1.1, \S6]{calegari_emerton} (v1.24.0-rc.1; the phenomenology catalogue, PH-39). Reported, not asserted here (v1.24.0-rc.2): the validation's
heavy arm reproduces it on the curve object (\file{docs/implementation/slot22/LN23\_report.md} Q6), printing as \code{layers=} the radical-series dimensions
$\dim\mathfrak{m}^{i}W=255-2i$, $i=0,\dots,127$ (so the Loewy layers have dimension $2$, the last dimension $1$). Their \S6 sentence that $2143$ is the smallest prime $N\equiv1\bmod9$ with an element of order $9$ in the class group of their sextic field is
reported false by an exact class-number computation out of tree ($1297$ is smaller; PH-40, \file{docs/implementation/slot22/NFQ1\_report.md} \S0, \S1);
this document does not assert it, and nothing here uses either sentence.
\end{remark}

%----------------------------------------------------------------------
\section{Nebentypus: \texorpdfstring{$\Ga_1(N)$}{Gamma1(N)} and \texorpdfstring{$(\Ga_0(N),\chi)$}{(Gamma0(N),chi)} on the fixed model by Shapiro}\label{sec:nebentypus}
%----------------------------------------------------------------------
% NEW at v1.15.0-pre.1 (hecke-extensions slot 15, WP-3; plan docs/implementation/slot15/00-slot15-plan.md
% S2 "WP-3" Part A, owner ruling Q2 of 2026-09-30: GENERAL N, characters chi:(Z/N)^x -> F_q^x, q = ell^d a
% general finite coefficient field, ord chi | q-1, NO layering).  The contract first (rule 21): nothing in
% this chapter is built at this version (the code is slot 16's, owner ruling Q3); the predictions
% (NB1)--(NB6) precede that code (rule 20).  The owner's standing direction on level N applies unchanged:
% Shapiro on the fixed model, the permutation (x) rho_n block structure exploited orbit by orbit, no second
% coset or tau construction.  Placed after sec:levelN, before the appendices, so that no chapter number moves.

\paragraph{What this chapter fixes.}
(i) \textbf{The groups, the characters and the coefficient rings} (Definition~\ref{def:nebdata}): $\Ga_1(N)\le\Ga_0(N)$, a
Dirichlet character $\chi\colon(\Z/N)^{\times}\to\F_q^{\times}$ for a finite field $\F_q$ of characteristic $\ell$, $q=\ell^{d}$
arbitrary, and its Teichm\"uller lift to the Galois ring $W_K(\F_q)=W(\F_q)/\ell^{K}$ (Lemma~\ref{lem:teichmuller}); the
coefficient module $\VR(\chi)$ of $\Ga_0(N)$ and its extension to the semigroup $\De_0(N)$. (ii) \textbf{The two models by
Shapiro} (Definition~\ref{def:nebmodel}): the \emph{character model} $L_{N,\chi}$ on the table of Definition~\ref{def:cosetrepsN},
whose relation matrices are the monomial matrices of Definition~\ref{def:levelNmodel} with every block scaled by a value of
$\chi$ (Lemma~\ref{lem:inducedchar}), and the \emph{$\Ga_1(N)$ model} on a table fibred over it (Lemma~\ref{lem:gammaonetable}).
(iii) \textbf{The diamond operators and the $\chi$-decomposition} (Proposition~\ref{prop:nebdecomp}): $\langle d\rangle$ as a fibre
map, the character model as the $\chi$-eigenspace of the diamonds inside the $\Ga_1(N)$ model, the decomposition over a ring
holding the values of every character, and $\chi=1$ as Definition~\ref{def:levelNmodel} bit for bit. (iv) \textbf{The Hecke
operators with character} (Proposition~\ref{prop:nebcorr}, Theorem~\ref{thm:nebrelations}): $T_p$ for $p\nmid N$, $U_p$ for
$p\mid N$, $\langle p\rangle$ and $T(p,p)$, all read off the shipped block formula \eqref{eq:levelcorr} with a scalar per block, on
the level-$1$ $\tau$-table; the relations with the $\chi(p)$ factor. (v) \textbf{The orbit structure and the cokernel} with a
character (Proposition~\ref{prop:neborbit}), the Heilbronn action with a character (Remark~\ref{rmk:nebheilbronn}), the cost
(Remark~\ref{rmk:nebcost}), the refusal rules (the box after Remark~\ref{rmk:nebcost}) and the predictions (NB1)--(NB6)
(Remark~\ref{rmk:nebcheckable}).

\paragraph{Notation and standing conventions.}
Those of \S\ref{sec:levelN}. For $\gamma=\smat abcd\in\Mat_2(\Z)$ write $a_\gamma:=a$, $d_\gamma:=d$. $N\ge1$;
$(\Z/N)^{\times}$ is written $G_N$, of order $\varphi(N)$ and exponent $\lambda(N)$ (Carmichael). $\ell$ is a prime, $q=\ell^{d}$ with
$d\ge1$, and $\F_q$ is \emph{the} finite field of Remark~\ref{rmk:nebfield} (a fixed model, so that matrices over it are
data). \textbf{Rule~\ref{A5}}: in this chapter an undecorated $\chi$, and $\chi_{-3}$, $\chi_{-4}$, $\chi_{-7}$, denote Dirichlet
characters; the decorated symbols $\chi_b$, $\chi^{\mathrm{ord}}_{n,p}$, $\chi_C$ keep their meanings elsewhere and do not occur
here. $\varsigma_N$, $\De_0(N)$, $\rho_{n,\chi}$, $\VR(\chi)$, $L_{N,\chi}$, $L_{\Ga_1(N)}$, $\sigma_u$, $R_\sigma$, $\langle d\rangle$,
$W_K(\F_q)$, $\mathfrak{t}$ (the Teichm\"uller lift) are new symbols of this chapter and mean nothing else; $\langle d\rangle$ is not the cyclic group $\langle
S\rangle$ (the argument of a diamond is always a residue class, never a matrix).

\subsection{The groups, the characters, the coefficient rings}\label{sec:nebdatasec}

\begin{definition}[Groups, semigroup, characters, coefficients]\label{def:nebdata}
\begin{enumerate}
\item $\varsigma_N\colon\Ga_0(N)\to G_N$, $\gamma\mapsto d_\gamma\bmod N$, and $\Ga_1(N):=\ker\varsigma_N=\{\gamma\in\Ga_0(N):
  d_\gamma\equiv1\ (N)\}$.
\item $\De_0(N):=\{\delta\in\Mat_2(\Z):\det\delta>0,\ N\mid c_\delta,\ \gcd(a_\delta,N)=1\}$, a monoid containing $\Ga_0(N)$
  \cite[\S3.3, (3.3.3), p.~68]{shimura} (his $\De'$ with $t=1$ and $\mathfrak h=(\Z/N\Z)^{\times}$; his $\Ga'$ of (3.3.2), p.~67, is then $\Ga_0(N)$; verified, R20-L). \LEAN{def:nebdata-model} (without the condition
  $\det\delta>0$)
\item A \emph{character} is a homomorphism $\chi\colon G_N\to\cring^{\times}$, extended to $\Z$ by $\chi(a):=\chi(a\bmod N)$ for
  $\gcd(a,N)=1$. Its \emph{parity} is $\chi(-1)\in\{\pm1\}$. \LEAN{def:nebdata-model} (over any commutative $\cring$)
\item The \textbf{coefficient rings} of this chapter are $\cring=\F_q$ and $\cring=W_K(\F_q)$ (Lemma~\ref{lem:teichmuller}),
  $K\ge1$; the characters are $\chi\colon G_N\to\F_q^{\times}$ --- so $\operatorname{ord}\chi\mid q-1$ and $\ell\nmid\operatorname{ord}\chi$
  automatically --- and, over $W_K(\F_q)$, their Teichm\"uller lifts $\mathfrak{t}\circ\chi$. The least $q$ that carries $\chi$ is
  $\ell^{d(\chi)}$, $d(\chi)$ the order of $\ell$ in $(\Z/\operatorname{ord}\chi)^{\times}$; any $q=\ell^{d}$ with $d(\chi)\mid d$ may be
  used.
\item $\VR(\chi):=\VR$ with $\Ga_0(N)$ acting by $\rho_{n,\chi}(\gamma):=\chi(a_\gamma)\,\rhon(\gamma)$, and $\De_0(N)$ acting by the
  same formula. \LEAN{def:nebdata-model}
\end{enumerate}
\end{definition}

\begin{lemma}[The definition is consistent]\label{lem:nebdata}
\PROVED{} (1) $\varsigma_N$ is a surjective homomorphism, so $\Ga_0(N)/\Ga_1(N)\cong G_N$, and $a_\gamma\equiv d_\gamma^{-1}\pmod N$ on
$\Ga_0(N)$ \LEAN{def:nebdata-model} (the congruence $a_\gamma d_\gamma\equiv1$ only; not the surjectivity of $\varsigma_N$). (2) $\delta\mapsto a_\delta\bmod N$ is a homomorphism of monoids $\De_0(N)\to G_N$; hence $\rho_{n,\chi}$ is
multiplicative on $\De_0(N)$ and $\VR(\chi)$ is a $\cring[\De_0(N)]$-module; on $\Ga_0(N)$, $\chi(a_\gamma)=\chi(d_\gamma)^{-1}=
\chi^{-1}(\varsigma_N(\gamma))$ \LEAN{def:nebdata-model}. (3) $-I\in\Ga_0(N)$ acts on $\VR(\chi)$ by the scalar $\chi(-1)(-1)^{n}$. (4) $\Ga_1(N)$ acts on
$\VR(\chi)$ through $\rhon$: $\VR(\chi)|_{\Ga_1(N)}=\VR|_{\Ga_1(N)}$. \LEAN{def:nebdata-model} (over any commutative $\cring$)
\end{lemma}

\begin{proof}
(1),(2) For $\gamma,\gamma'$ with $N\mid c,c'$: $(\gamma\gamma')_{22}=cb'+dd'\equiv dd'$ and $(\gamma\gamma')_{11}=aa'+bc'\equiv aa'
\pmod N$; $\det\gamma=ad-bc\equiv ad$, so $ad\equiv1$ on $\Ga_0(N)$. A unit $u$ modulo $N$ is $d_\gamma$ for a lift
$\gamma\in\Ga$ of $\mathrm{diag}(u^{-1},u)\in\SL_2(\Z/N)$ (surjectivity of $\SL_2(\Z)\to\SL_2(\Z/N)$, as in
Definition~\ref{def:levelellmodel}), and such a lift lies in $\Ga_0(N)$. $\De_0(N)$ is closed under products by the same
congruences, and $\gcd(aa',N)=1$. Multiplicativity of $\rho_{n,\chi}$ is that of $\rhon$ (Proposition~\ref{prop:symconv}(2)) times that
of $\chi\circ a$. (3) $a_{-I}=-1$ and $\rhon(-I)=(-1)^{n}$. (4) $a_\gamma\equiv d_\gamma^{-1}\equiv1$ on $\Ga_1(N)$.
\end{proof}

\begin{lemma}[The Galois ring and the Teichm\"uller lift]\label{lem:teichmuller}
\PROVED{} Let $C\in\F_\ell[Y]$ be monic irreducible of degree $d$ (the project-canonical $C_{\ell,d}$ of \Espec{def:canon}, which
fixes $\F_q=\F_\ell[Y]/(C)$), and $\widetilde C\in\Z[Y]$ its coefficientwise lift with coefficients in $[0,\ell)$. Put
$W_K(\F_q):=(\Z/\ell^{K})[Y]/(\widetilde C)$ (any monic lift of $C$ gives an isomorphic ring, the Galois ring $\mathrm{GR}(\ell^{K},d)$; the coefficientwise lift is fixed so that matrices over it are data, as in Remark~\ref{rmk:nebfield}(1)).
\begin{enumerate}
\item $W_K(\F_q)$ is a local ring, free of rank $d$ over $\Z/\ell^{K}$, with maximal ideal $(\ell)$ and residue field $\F_q$;
  every element is $\ell^{j}u$ with $u$ a unit, and it is a quotient of $W_{K+1}(\F_q)$ by $\ell^{K}$.
\item For each $x\in\F_q^{\times}$ there is exactly one $\mathfrak{t}(x)\in W_K(\F_q)$ with $\mathfrak{t}(x)^{q-1}=1$ and $\mathfrak{t}(x)\equiv x\pmod\ell$;
  $\mathfrak{t}\colon\F_q^{\times}\to W_K(\F_q)^{\times}$ is an injective homomorphism, and for \emph{every} lift $\tilde x$ of $x$,
  $\mathfrak{t}(x)=\tilde x^{\,q^{K-1}}$.
\item $W_K(\F_q)^{\times}=\mathfrak{t}(\F_q^{\times})\times(1+\ell W_K(\F_q))$, the second factor an $\ell$-group. Consequently a homomorphism
  $G_N\to W_K(\F_q)^{\times}$ whose order is prime to $\ell$ is $\mathfrak{t}\circ\chi$ for exactly one $\chi\colon G_N\to\F_q^{\times}$.
\end{enumerate}
\end{lemma}

\begin{proof}
(1) $\widetilde C$ is monic, so the quotient is free of rank $d$; modulo $\ell$ it is the field $\F_q$, so $(\ell)$ is the unique
maximal ideal ($\ell$ is nilpotent), and an element not in $(\ell)$ is a unit. Writing an element as $\ell^{j}u$: take $j$ maximal
with the element in $(\ell^{j})$. (2) The polynomial $Y^{q-1}-1$ has the simple root $x$ in $\F_q$ (its derivative
$(q-1)x^{q-2}$ is a unit), so by Hensel's lemma over the complete (here: $\ell$-adically nilpotent) local ring it has exactly one root
lifting $x$. For a lift $\tilde x$ write $\tilde x=\mathfrak{t}(x)(1+\ell y)$; then $\tilde x^{\,q^{K-1}}=\mathfrak{t}(x)^{q^{K-1}}(1+\ell
y)^{q^{K-1}}$, $\mathfrak{t}(x)^{q}=\mathfrak{t}(x)$ gives $\mathfrak{t}(x)^{q^{K-1}}=\mathfrak{t}(x)$, and $(1+\ell y)^{\ell^{m}}\equiv1\pmod{\ell^{m+1}}$ by
induction on $m$ (binomial theorem, $\ell\ge2$), with $m=d(K-1)\ge K-1$. Uniqueness makes $\mathfrak{t}$ multiplicative; it is injective
because it is a section of reduction. (3) Reduction $W_K(\F_q)^{\times}\to\F_q^{\times}$ is onto with kernel $1+\ell W_K$, of order
$q^{K-1}$, and $\mathfrak{t}$ splits it. A homomorphism of order prime to $\ell$ has trivial component in the $\ell$-group $1+\ell W_K$.
\end{proof}

\begin{remark}[Which field, and which characters]\label{rmk:nebfield}
(1) \textbf{The field is data.} $\F_q$ is $\F_\ell[Y]/(C_{\ell,d})$ of \Espec{def:canon} --- the field the extension layer already
uses for residue fields of degree $d$ --- so a character is a table of elements of that field, and the cochain-level matrices of
this chapter are well defined. Another choice of $C$ changes those matrices by a field isomorphism and changes no dimension, rank
or characteristic polynomial up to that isomorphism. (2) \textbf{Characters of $\ell$-power order} (possible when $\ell\mid\varphi(N)$,
e.g.\ $\ell^{2}\mid N$ or $\ell\mid p-1$ for a prime $p\mid N$) take values in $\mu_{\ell^{j}}$, which lives in no $W_K(\F_q)$ outside
the trivial group and reduces to $1$ in every $\F_q$; they need a \emph{ramified} coefficient ring $\Z_\ell[\zeta_{\ell^{j}}]/\ell^{K}$
and are \textbf{refused by name} (\code{character\_wild}; the box after Remark~\ref{rmk:nebcost}). Over $\F_q$ such a character is
trivial, so it is not a restriction of the owner's scope (Q2: $\operatorname{ord}\chi\mid q-1$). Since v1.22.0-pre.1 such characters are admitted over the ramified chain rings $\mathcal O_{j,e}$ of
Definition~\ref{def:nebwild} (\S\ref{sec:nebwildsec}); over $\F_q$ and $W_K(\F_q)$ the refusal stands. (3) The Galois group of
$\F_q/\F_\ell$ acts on characters by $\chi\mapsto\chi^{\ell}$ and on everything below by conjugating matrix entries; the models of
$\chi$ and $\chi^{\ell}$ are Frobenius conjugates, so one character per Frobenius orbit suffices for dimensions and characteristic
polynomials (up to conjugating their coefficients).
\end{remark}

\subsection{The two models}\label{sec:nebmodelsec}

\begin{lemma}[Induced modules of an arbitrary representation; the character model is monomial]\label{lem:inducedchar}
\PROVED{} Let $\cring$ be a commutative ring, $H\le\Ga$ of finite index $m$ with representatives $w_1=I,\dots,w_m$ of
$H\backslash\Ga$, and $\eta\colon H\to\GL_r(\cring)$ a homomorphism, $X=\cring^{r}$. Definition~\ref{def:inducedmodule} and
Proposition~\ref{prop:blockmatrices} hold verbatim for $\Ind_H^{\Ga}X$ with $\eta$ in place of $\rhon|_H$:
\[
  M^{\eta}_H(\gamma)=\sum_{j}E_{j,\,j\cdot\gamma}\otimes\eta\bigl(h_j(\gamma)\bigr),
\]
a homomorphism $\Ga\to\GL_{mr}(\cring)$ with one non-zero block per block row; Definition~\ref{def:inducedls} and
Proposition~\ref{prop:lhcompliant}(1),(2),(4) hold for the induced matrix local system $L^{\eta}_H$ (in (4), $H^{q}(H;X)$ with the
action $\eta$). In particular, for $H=\Ga_0(N)$ with the table of Definition~\ref{def:cosetrepsN} and $\eta=\rho_{n,\chi}$,
\begin{equation}\label{eq:MHchi}
  M_{N,\chi}(\gamma)\;=\;\sum_{j=1}^{m_N}E_{j,\,j\cdot\gamma}\otimes\chi\bigl(a_{h_j(\gamma)}\bigr)\,\rhon\bigl(h_j(\gamma)\bigr) ,
\end{equation}
the matrix $M_N(\gamma)$ of Definition~\ref{def:levelNmodel} with block $j$ multiplied by the scalar $\chi(a_{h_j(\gamma)})\in
\mu_{\operatorname{ord}\chi}(\cring)$: \textbf{same permutation, same blocks up to a root of unity}.
\end{lemma}

\begin{proof}
The proofs cited use of $\rhon$ only that it is a homomorphism on $H$ (the bijectivity of restriction to the representatives, the
cocycle identity, and homomorphy of $M_H$), and of Proposition~\ref{prop:lhcompliant}(4) only Theorem~\ref{thm:explicit_complex_equiv},
Remark~\ref{rmk:orbifold} (hypotheses on the model and on $6\in\cring^{\times}$, not on the module) and Shapiro's lemma, which holds
for every $\cring[H]$-module. \eqref{eq:MHchi} is the displayed formula with $\eta=\rho_{n,\chi}$ and
Definition~\ref{def:nebdata}(5).
\end{proof}

\begin{lemma}[The table of $\Ga_1(N)\backslash\Ga$]\label{lem:gammaonetable}
\PROVED{} For $u\in G_N$ let $\sigma_1:=I$ and, for $u\ne1$, $\sigma_u:=\smat{\alpha}{\beta}{N}{u'}$ with $u'\in[1,N)$ the
representative of $u$, $\alpha:=u'^{-1}\bmod N\in[1,N)$ and $\beta:=(\alpha u'-1)/N$. Order $G_N$ by the representatives $u'$.
\begin{enumerate}
\item $\sigma_u\in\Ga_0(N)$ and $\varsigma_N(\sigma_u)=u$; the $\sigma_u$ represent $\Ga_1(N)\backslash\Ga_0(N)$.
\item The elements $w_{(j,u)}:=\sigma_u\,w_j$, $1\le j\le m_N$, $u\in G_N$ ($w_j$ of Definition~\ref{def:cosetrepsN}), represent
  $\Ga_1(N)\backslash\Ga$, with $w_{(1,1)}=I$; ordered $j$-major, $u$-minor, index $J(j,u):=(j-1)\varphi(N)+\mathrm{rk}(u)$. For $N\ge3$
  there are $m_N\varphi(N)=N^{2}\prod_{p\mid N}(1-p^{-2})$ of them ($[\Ga:\Ga_1(N)]$, as $-I\notin\Ga_1(N)$); for $N\le2$,
  $\Ga_1(N)=\Ga_0(N)$ and the table is Definition~\ref{def:cosetrepsN}'s.
\item For $\gamma\in\Ga$: $(j,u)\cdot\gamma=\bigl(j\cdot\gamma,\ u\,\varsigma_N(h_j(\gamma))\bigr)$ and
  $h_{(j,u)}(\gamma)=\sigma_u\,h_j(\gamma)\,\sigma_{u\varsigma_N(h_j(\gamma))}^{-1}\in\Ga_1(N)$. The model $L_{\Ga_1(N)}:=L^{\rhon}_{\Ga_1(N)}$ (Lemma~\ref{lem:inducedchar} with $H=\Ga_1(N)$, $\eta=\rhon|_{\Ga_1(N)}$) of
  Definition~\ref{def:inducedls} on this table has fibre rank $m_N\varphi(N)(n+1)$ and monomial relation matrices whose block
  permutation is fibred over that of $M_N$.
\end{enumerate}
\end{lemma}

\begin{proof}
(1) $\alpha u'-\beta N=1$ and $N\mid\alpha u'-1$, so $\beta\in\Z$ and $\det\sigma_u=1$; $c=N$, $d=u'$. Cosets of $\ker\varsigma_N$ are
the fibres of $\varsigma_N$. (2) $\Ga=\bigsqcup_j\Ga_0(N)w_j$ and $\Ga_0(N)=\bigsqcup_u\Ga_1(N)\sigma_u$. The count is
$m_N\cdot\abs{G_N}$, and $m_N\varphi(N)=N\prod(1+p^{-1})\cdot N\prod(1-p^{-1})$. (3) $w_{(j,u)}\gamma=\sigma_uw_j\gamma=\sigma_u
h_j(\gamma)w_{j\cdot\gamma}=\bigl[\sigma_uh_j(\gamma)\sigma_{u'}^{-1}\bigr]\sigma_{u'}w_{j\cdot\gamma}$ with $u'=u\,\varsigma_N(h_j(\gamma))$; the
bracket lies in $\Ga_0(N)$ with $\varsigma_N=u\cdot\varsigma_N(h_j(\gamma))\cdot u'^{-1}=1$. Uniqueness of the decomposition gives the
action and the cocycle.
\end{proof}

\begin{definition}[The level-$N$ models with nebentypus]\label{def:nebmodel}
Let $N\ge1$, $\cring\in\{\F_q,W_K(\F_q)\}$ (Definition~\ref{def:nebdata}(4)) --- or any commutative ring at cochain level --- and
$\chi\colon G_N\to\cring^{\times}$.
\begin{enumerate}
\item The \textbf{character model} $L_{N,\chi}:=L^{\rho_{n,\chi}}_{\Ga_0(N)}$: fibre $\Ind_{\Ga_0(N)}^{\Ga}\VR(\chi)$ of rank $m_N(n+1)$,
  coordinates of Definition~\ref{def:inducedmodule} with the $w_j$ of Definition~\ref{def:cosetrepsN}, relation matrices
  \eqref{eq:MHchi}. For $6\in\cring^{\times}$ its cohomology is $H^{\bullet}(\Ga_0(N);\VR(\chi))$ (Lemma~\ref{lem:inducedchar}).
\item The \textbf{$\Ga_1(N)$ model} $L_{\Ga_1(N)}$ of Lemma~\ref{lem:gammaonetable}(3): for $6\in\cring^{\times}$ its cohomology is
  $H^{\bullet}(\Ga_1(N);\VR)$.
\end{enumerate}
The character model is \textbf{the} model of this chapter: its size is that of the $\Ga_0(N)$ model, whatever $\chi$. The
$\Ga_1(N)$ model is $\varphi(N)$ times larger and is specified as the oracle of (NB3) and for $\ell\mid\varphi(N)$, where no
character decomposition exists (Proposition~\ref{prop:nebdecomp}(3)).
\end{definition}

\begin{proposition}[Diamond operators, and the character model inside the $\Ga_1(N)$ model]\label{prop:nebdecomp}
\PROVED{} Let $\cring$ be a commutative ring and view $\Ind_{\Ga_1(N)}^{\Ga}\VR$ as functions $f\colon\Ga\to\VR$ with
$f(h\gamma)=\rhon(h)f(\gamma)$, $h\in\Ga_1(N)$ (Definition~\ref{def:inducedmodule}).
\begin{enumerate}
\item \textbf{(Diamonds as fibre maps.)} For $\sigma\in\Ga_0(N)$, $(R_\sigma f)(x):=\rhon(\sigma)f(\sigma^{-1}x)$ is an automorphism of
  $\cring[\Ga]$-modules of $\Ind_{\Ga_1(N)}^{\Ga}\VR$ depending only on $\varsigma_N(\sigma)$, and $\sigma\mapsto R_\sigma$ is a
  homomorphism. $\langle d\rangle:=R_{\sigma_d}$ ($d\in G_N$) is therefore a \emph{fibre map} of $L_{\Ga_1(N)}$ --- one matrix at every
  simplex, a morphism of matrix local systems --- and for $6\in\cring^{\times}$ it induces on $H^{\bullet}(\Ga_1(N);\VR)$ the operator
  $[\Ga_1(N)\,\sigma_d\,\Ga_1(N)]=(\conj_{\sigma_d})_{\ast}$ of \S\ref{sec:hidafull}'s notation. $\langle-1\rangle$ is $(-1)^{n}\cdot$
  the action of $-I$, and $\langle d\rangle\langle d'\rangle=\langle dd'\rangle$.
\item \textbf{(The character model is the $\chi$-eigenspace.)} The inclusion of functions identifies $\Ind_{\Ga_0(N)}^{\Ga}\VR(\chi)$
  with the $\cring[\Ga]$-submodule $\{f:\langle d\rangle f=\chi(d)f\ \forall d\}$ of $\Ind_{\Ga_1(N)}^{\Ga}\VR$, and
  $P_\chi f:=\sum_{u\in G_N}\chi(u)^{-1}\langle u\rangle f$ is a $\cring[\Ga]$-map onto a submodule of it with $P_\chi|=\varphi(N)\cdot\id$
  there. No division is used.
\item \textbf{(Decomposition, (NB3).)} If $\varphi(N)\in\cring^{\times}$ and $\cring$ is a domain or a local ring containing a
  primitive $\lambda(N)$-th root of unity whose powers have unit differences (e.g.\ $\F_q$ or $W_K(\F_q)$ with $\lambda(N)\mid q-1$), then
  $e_\chi:=\varphi(N)^{-1}P_\chi$ are orthogonal idempotents with $\sum_\chi e_\chi=1$ over the $\varphi(N)$ characters
  $G_N\to\cring^{\times}$, and
  \[
    \Ind_{\Ga_1(N)}^{\Ga}\VR=\bigoplus_{\chi}\Ind_{\Ga_0(N)}^{\Ga}\VR(\chi),\qquad
    H^{\bullet}(\Ga_1(N);\VR)=\bigoplus_\chi H^{\bullet}(\Ga_0(N);\VR(\chi))\quad(6\in\cring^{\times}),
  \]
  $\langle d\rangle$ acting on the $\chi$-summand by $\chi(d)$, and compatibly with every operator of
  Proposition~\ref{prop:nebcorr}.
\item \textbf{($\chi=1$, (NB2).)} For $\chi=1$, \eqref{eq:MHchi} is $M_N(\gamma)$ entry for entry and $L_{N,1}$ is
  Definition~\ref{def:levelNmodel}'s $L_N$; every operator of Proposition~\ref{prop:nebcorr} at $\chi=1$ is the operator of
  Definition~\ref{def:levelNmodel} on the same cochains.
\end{enumerate}
\end{proposition}

\begin{proof}
(1) For $h\in\Ga_1(N)$: $(R_\sigma f)(hx)=\rhon(\sigma)f(\sigma^{-1}h\sigma\cdot\sigma^{-1}x)=\rhon(\sigma)\rhon(\sigma^{-1}h\sigma)f(\sigma^{-1}x)=
\rhon(h)(R_\sigma f)(x)$, using that $\Ga_1(N)=\ker\varsigma_N$ is normal in $\Ga_0(N)$. $\Ga$-equivariance: $(\gamma\cdot R_\sigma f)(x)=
\rhon(\sigma)f(\sigma^{-1}x\gamma)=R_\sigma(\gamma\cdot f)(x)$. For $h\in\Ga_1(N)$: $R_{\sigma h}f(x)=\rhon(\sigma)\rhon(h)f(h^{-1}\sigma^{-1}x)=
R_\sigma f(x)$. Homomorphy: $R_\sigma R_{\sigma'}=R_{\sigma\sigma'}$ by expanding. A $\cring[\Ga]$-endomorphism of the fibre acts at every
simplex by the same matrix (Theorem~\ref{thm:rep_ls_equiv}). The induced map: under the reciprocity $\phi\mapsto\psi:=\phi(\cdot)(I)$ of
Proposition~\ref{prop:levelcorrespondence}(3), $(R_\sigma\circ\phi)(y)(I)=\rhon(\sigma)\phi(y)(\sigma^{-1})=\rhon(\sigma)\phi(\sigma^{-1}y)(I)=
\rhon(\sigma)\psi(\sigma^{-1}y)$ ($\Ga$-equivariance of $\phi$), which is Proposition~\ref{prop:levelcorrespondence}(3)'s formula for
$A=B=\Ga_1(N)$, $g=\sigma$ (one coset, $\tau=\sigma^{-1}\cdot$, a valid $\sigma$-twisted datum since $\sigma\in\Ga$). $-I$: $R_{-I}f(x)=
(-1)^{n}f(-x)=(-1)^{n}((-I)\cdot f)(x)$ as $-I$ is central.
(2) $F\in\Ind_{\Ga_0(N)}^{\Ga}\VR(\chi)$ means $F(\sigma x)=\chi(a_\sigma)\rhon(\sigma)F(x)$ for $\sigma\in\Ga_0(N)$; restricted to
$\Ga_1(N)$ this is the defining condition of the larger module (Lemma~\ref{lem:nebdata}(4)), and then $R_\sigma F(x)=\rhon(\sigma)
\chi(a_{\sigma^{-1}})\rhon(\sigma^{-1})F(x)=\chi(\varsigma_N(\sigma))F(x)$ (Lemma~\ref{lem:nebdata}(2)); conversely this identity for all
$\sigma$ is the defining condition. $P_\chi$ commutes with $\Ga$ by (1); $\langle d\rangle P_\chi f=\sum_u\chi(u)^{-1}\langle du\rangle f=
\chi(d)P_\chi f$; on the eigenspace $P_\chi$ is $\sum_u\chi(u)^{-1}\chi(u)=\varphi(N)$.
(3) Character orthogonality $\sum_\chi\chi(u)=\varphi(N)[u=1]$ and $\sum_u\chi(u)\chi'(u)^{-1}=\varphi(N)[\chi=\chi']$ hold over such a
ring (the characters of the finite abelian group $G_N$ with values in $\mu_{\lambda(N)}(\cring)$ are $\varphi(N)$ in number and the
standard proof needs only that $1-\zeta$ is a non-zero-divisor for $\zeta\ne1$ in $\mu_{\lambda(N)}$). So the $e_\chi$ are orthogonal
idempotents summing to $1$, each a $\cring[\Ga]$-map with image the $\chi$-eigenspace, i.e.\ the character model by (2).
Cohomology commutes with finite direct sums; the Hecke compatibility is Proposition~\ref{prop:nebcorr}(4).
(4) $\chi(\cdot)=1$ in \eqref{eq:MHchi}; the formulas of Proposition~\ref{prop:nebcorr} at $\chi=1$ are \eqref{eq:levelcorr}.
\end{proof}

\subsection{The Hecke operators with a character}\label{sec:nebheckesec}

\begin{proposition}[Level-$N$ correspondences with a character]\label{prop:nebcorr}
\PROVED{} Let $\cring$ be a commutative ring and $\chi\colon G_N\to\cring^{\times}$. Let $A,B\le\Ga_0(N)$ be of finite index,
$g\in\De_0(N)\cap\De$, and fix $w_a$, $v_i$, $r_t$, $s_j$, $\tau$ exactly as in Proposition~\ref{prop:levelcorrespondence}, with
$\tau$ the level-$1$ $g$-twisted chain datum of $\Tdc{g}$ (Definition~\ref{def:tau}) --- \textbf{no new $\tau$-table}. Define
\begin{equation}\label{eq:nebcorr}
  \bigl(\widetilde{[A\,g\,B]}_\chi\,\phi\bigr)(x)_i\;:=\;\sum_{t}\ \chi\bigl(a_{\delta_{i,t}}\bigr)\,\rhon(\delta_{i,t})\,
  \phi\bigl(\tau(s_{j(i,t)}\,x)\bigr)_{a(i,t)},
\end{equation}
\eqref{eq:levelcorr} with one scalar per term, for cochains with values in $\Ind_A^{\Ga}\VR(\chi)$ (in the coordinates of
Lemma~\ref{lem:inducedchar}). \LEAN{prop:nebcorr-model} (in the reciprocity form of Proposition~\ref{prop:levelcorrespondence}(3), on the
model, for a fixed table and $\tau$)
\begin{enumerate}
\item Every $\delta_{i,t}$ lies in $\De_0(N)$, so the scalar $\chi(a_{\delta_{i,t}})$ is defined; \LEAN{prop:nebcorr-model} (this
  sentence; not the congruence that follows) $a_{\delta_{i,t}}\equiv a_{r_t}^{-1}\,a_g\,
  a_{kw_a^{-1}}\pmod N$.
\item \eqref{eq:nebcorr} is a cochain map $\Ceq(K,\Rel;L^{\chi}_A)\to\Ceq(K,\Rel;L^{\chi}_B)$ over every $\cring$, and for
  $6\in\cring^{\times}$ it induces $[A\,g\,B]\colon H^{\bullet}(A;\VR(\chi))\to H^{\bullet}(B;\VR(\chi))$ --- the map
  $\cor\circ(\conj_g)_{\ast}\circ\res$ of the compatible pair $(\conj_g^{-1},\rho_{n,\chi}(g))$ --- independent of $\tau$ and of every
  choice of representatives. \LEAN{prop:nebcorr-model} (the cochain map over every $\cring$ and the factorisation
  $\cor\circ(\conj_g)_{\ast}\circ\res$, on the model; not the independence of $\tau$ and of the representatives)
\item \textbf{(The operators.)} With $A=B=\Ga_0(N)$ and the $r_t$ of Definition~\ref{def:cosetrepsN}(6):
  $T^{(N,\chi)}_p:=[A\,\smat100p\,A]$ for $p\nmid N$; $U^{(N,\chi)}_p:=[A\,\smat100p\,A]$ for $p\mid N$; for $p\nmid N$,
  $T^{(N,\chi)}(p,p):=[A\,pI\,A]=\rho_{n,\chi}(pI)=\chi(p)\,p^{\,n}$ (one coset, $\tau=\id$, every $\cring$, as in
  Proposition~\ref{prop:diamond}). \LEAN{prop:nebcorr-model} ($T^{(N,\chi)}(p,p)=\chi(p)\,c$, $c$ the scalar of $\rho(pI)$; not the
  $\Ga_1(N)$ model) On the $\Ga_1(N)$ model ($A=B=\Ga_1(N)$, $\chi$ trivial, $g\in\De_0(N)$):
  $T^{(\Ga_1)}_p$, $U^{(\Ga_1)}_p$ likewise, and $T^{(\Ga_1)}(p,p):=[\Ga_1(N)\,p\sigma_p\,\Ga_1(N)]=p^{\,n}\langle p\rangle$.
\item \textbf{(Compatibility.)} Under the inclusion of Proposition~\ref{prop:nebdecomp}(2) the cochain maps of (3) on the
  character model are the restrictions of those on the $\Ga_1(N)$ model, for $g\in\{\smat100p,pI\}$ ($T(p,p)$ with $\sigma_p$), and
  they commute with every $\langle d\rangle$. \LEAN{prop:nebcorr-model} (term-wise only: for $a\equiv1\pmod N$ the scalar
  $\chi(a)$ of a term is $1$; not the restriction form on the $\Ga_1(N)$ model, nor the commutation with $\langle d\rangle$)
\end{enumerate}
\end{proposition}

\begin{proof}
(1) $\delta_{i,t}=r_t^{-1}\,g\,(kw_a^{-1})$ by Proposition~\ref{prop:levelcorrespondence}(1), with $r_t\in B\subseteq\Ga_0(N)$,
$kw_a^{-1}\in A\subseteq\Ga_0(N)$ and $g\in\De_0(N)$; $\De_0(N)$ is a monoid (Lemma~\ref{lem:nebdata}(2)), and $r_t^{-1}\in\Ga_0(N)$.
(2) The proof of Proposition~\ref{prop:levelcorrespondence}(2),(3) goes through with $\rhon$ replaced by $\rho_{n,\chi}$ on
$\De_0(N)$, because every matrix at which it evaluates $\rhon$ in the transport and in the $B$-equivariance step --- $r_t^{-1}g$,
$kw_a^{-1}=h_1(k)$, $r_t^{-1}b_tg=b\,r^{-1}_{t\cdot b}g$ with $b\in B$, $b_t\in B''$ --- lies in $\De_0(N)$ by (1), and it uses of $\rhon$
only multiplicativity on the matrices it multiplies (Lemma~\ref{lem:nebdata}(2)) and \eqref{eq:MH}, here \eqref{eq:MHchi}
(Lemma~\ref{lem:inducedchar}). In particular $\rho_{n,\chi}(r_t^{-1}g)\rho_{n,\chi}(kw_a^{-1})=\rho_{n,\chi}(\delta_{i,t})=\chi(a_{\delta_{i,t}})
\rhon(\delta_{i,t})$, which is \eqref{eq:nebcorr}. The classical identification is the proof of Theorem~\ref{thm:classical} for the pair
$(\conj_g^{-1},\rho_{n,\chi}(g))$, compatible because $\rho_{n,\chi}(ghg^{-1})\rho_{n,\chi}(g)=\rho_{n,\chi}(g)\rho_{n,\chi}(h)$ for
$h\in A\cap g^{-1}Bg$ (all in $\De_0(N)$), with \STD{\cite[Ch.~III.9--10]{brown}} as there.
(3) $pI$ is central, $\Ga_0(N)\,pI\,\Ga_0(N)=pI\,\Ga_0(N)$, and $a_{pI}=p$ is a unit modulo $N$ iff $p\nmid N$. On $\Ga_1(N)$: $p\sigma_p$
has $a\equiv p\cdot p^{-1}=1$, lies in $\De_0(N)$, normalises $\Ga_1(N)$ and acts as $\rhon(p\sigma_p)=p^{\,n}\rhon(\sigma_p)$, which is
$p^{\,n}R_{\sigma_p}$ by the computation in the proof of Proposition~\ref{prop:nebdecomp}(1).
(4) Let $g\in\{\smat100p,\ p\sigma_p\}$, so $a_g\equiv1\pmod N$, and put $A_0'':=\Ga_0(N)\cap g\Ga_0(N)g^{-1}$,
$A_1'':=\Ga_1(N)\cap g\Ga_1(N)g^{-1}=A_0''\cap\Ga_1(N)$ (for $g=\smat100p$: $g^{-1}\gamma'g\in\Ga_0(N)$ iff $Np\mid c_{\gamma'}$).
\emph{The map $A_1''\backslash\Ga_1(N)\to A_0''\backslash\Ga_0(N)$ is a bijection.} Injective: $r,r'\in\Ga_1(N)$ with $rr'^{-1}\in A_0''$
have $rr'^{-1}\in A_0''\cap\Ga_1(N)=A_1''$. Surjective: $\varsigma_N(A_0'')=G_N$, since for $u\in G_N$ a lift $\gamma\in\Ga$ of
$\mathrm{diag}(u^{-1},u)\in\SL_2(\Z/Np)$ (which exists: $u$ is a unit modulo $Np$ when $p\mid N$, and by the Chinese remainder theorem
with $\gamma\equiv I\pmod p$ when $p\nmid N$) has $Np\mid c_\gamma$ and $d_\gamma\equiv u$; hence $A_0''\Ga_1(N)=\Ga_0(N)$ and every coset
$A_0''r$ meets $\Ga_1(N)$. Choose the $r_t$ in $\Ga_1(N)$, for both models. In the function picture of
Proposition~\ref{prop:levelcorrespondence}(3), both operators read $\psi\mapsto\sum_t\rho(r_t^{-1}g)\,\psi(\tau(r_t\,\cdot))$ with the same
$r_t$ and the same level-$1$ $\tau$; on $\psi\in\Hom_{\Ga_0(N)}(C_\bullet,\VR(\chi))\subseteq\Hom_{\Ga_1(N)}(C_\bullet,\VR)$ (the inclusion
of Proposition~\ref{prop:nebdecomp}(2) under reciprocity, $\VR(\chi)|_{\Ga_1(N)}=\VR$) one has
$\rho_{n,\chi}(r_t^{-1}g)=\chi(a_{r_t^{-1}})\chi(a_g)\rhon(r_t^{-1}g)=\rhon(r_t^{-1}g)$, as $a_{r_t^{-1}}=d_{r_t}\equiv1$ and $a_g\equiv1$. So
the two cochain maps agree on the character model, and Proposition~\ref{prop:levelcorrespondence}(3)'s independence of the
representatives makes the choice $r_t\in\Ga_1(N)$ harmless. Commutation with $\langle d\rangle$: on a sum of $\chi$-eigenspaces it
follows, since each operator preserves each character model and $\langle d\rangle$ is the scalar $\chi(d)$ there
(Proposition~\ref{prop:nebdecomp}(3)'s situation); in general it is the statement $\sigma_dg\sigma_d^{-1}\in\Ga_1(N)g\Ga_1(N)$, i.e.\ that
the $\Ga_1(N)$-double cosets in $\Ga_0(N)g\Ga_0(N)$ are indexed by $a_\delta\bmod N$ \STD{\cite[Lemma~3.29(1), (2), p.~66]{shimura} (printed as a Lemma; for $\det\delta$ prime to $N$) and \cite[Prop.~3.32(1), Prop.~3.33, pp.~68--70]{shimura}
(with $\mathfrak h=\{1\}$, $t=1$, i.e.\ $\Ga'=\Ga_1(N)$, for general $\det\delta$); verified, R20-L, v1.21.0-pre.1, in place of ``Prop.~3.29 and \S3.3''} (v1.15.0-pre.2, WP-4 referee A1).
\end{proof}

\begin{theorem}[Hecke relations with a character]\label{thm:nebrelations}
Assume \ref{H4}, $6\in\cring^{\times}$, and $\chi\colon G_N\to\cring^{\times}$. On $H^{\bullet}(\Ga_0(N);\VR(\chi))$, computed by the
character model:
\begin{enumerate}
\item the operators $T^{(N,\chi)}_p$ ($p\nmid N$), $U^{(N,\chi)}_p$ ($p\mid N$) and $T^{(N,\chi)}(p,p)$ commute;
\item for $p\nmid N$ and $\nu\ge1$ the full operators satisfy $T_{p^{\nu+1}}=T_pT_{p^{\nu}}-p\,T(p,p)\,T_{p^{\nu-1}}=T_pT_{p^{\nu}}-
  \chi(p)\,p^{\,n+1}\,T_{p^{\nu-1}}$; \LEANlit{thm:nebrelations-model}{the primitive products of \cite[Thm.~3.24(5), Prop.~3.17]{shimura} carried to the Hecke ring of $(\Ga_0(N),\De_0(N))$ by \cite[Prop.~3.31 and p.~71, before Thm.~3.35]{shimura}
  ($p\nmid N$; Shimura's ring carries no character, $\chi$ enters only through $T(p,p)\mapsto\chi(p)c$), as the hypotheses \texttt{hTpchi0}, \texttt{hTpchi}; locator verified (R20-L)} (on the model, for a fixed table and $\tau$, the full operators in the assembled form $\sum_j(\chi(p)c)^{j}\,T_{\Deprim}$)
\item for $p\mid N$, $U_{p^{\nu}}=(U_p)^{\nu}$. \LEANlit{thm:nebrelations-model}{the products of the $U$-double cosets, \cite[Thm.~3.34(3) with Prop.~3.33, pp.~70--71]{shimura} ($p\mid N$; $T'(p^{\mu})$ is the single double coset of $\mathrm{diag}(1,p^{\mu})$), as the hypothesis \texttt{hUchi}; locator verified (R20-L)}
  (on the model, for a fixed table and $\tau$)
\end{enumerate}
\STD{\cite[Thm.~3.34, Thm.~3.35 and \S3.3]{shimura} for the identities in the Hecke ring $\mathcal{H}(\Ga_0(N),\De_0(N))$ --- for $p\nmid N$
Thm.~3.24 carried over by Prop.~3.31 (p.~71, before Thm.~3.35), for $p\mid N$ Thm.~3.34(3) with Prop.~3.33; Shimura's ring has no character,
$\chi$ enters through $T(p,p)$ (Proposition~\ref{prop:nebcorr}(3)) --- and
\cite[Ch.~III.9--10]{brown} with Obligation~\ref{obl:algebra} for the ring homomorphism (the character twist: \cite[\S1.1, Lemma~1.1.5, p.~196]{ash_stevens});
the Shimura locators are \textbf{verified} on the owner's scan (R20-L, v1.21.0-pre.1)}
\end{theorem}

\begin{proof}[Proof sketch]
As for Theorem~\ref{thm:relations}: the double-coset action of $\mathcal H(\Ga_0(N),\De_0(N))$ on $H^{\bullet}(\Ga_0(N);M)$ for a
$\cring[\De_0(N)]$-module $M$ (here $\VR(\chi)$, Lemma~\ref{lem:nebdata}(2)) is a ring homomorphism (Obligation~\ref{obl:algebra}'s
argument, which uses only that the coefficients form a $\De_0(N)$-module), the cited identities hold in that ring, and
Proposition~\ref{prop:nebcorr}(3) evaluates $T(p,p)$ as the scalar $\chi(p)p^{\,n}$. With $k=n+2$, $\chi(p)p^{\,n+1}=\chi(p)p^{\,k-1}$,
the classical factor.
\end{proof}

\subsection{Orbit structure, the cokernel, and the Heilbronn action}\label{sec:neborbitsec}

\begin{proposition}[Orbit-wise fixed points and the cokernel with a character]\label{prop:neborbit}
\PROVED{} Let $\cring$ be a commutative ring, $\chi\colon G_N\to\cring^{\times}$, $M:=\Ind_{\Ga_0(N)}^{\Ga}\VR(\chi)$ with
\eqref{eq:MHchi}, and $F\le\Ga$ finite.
\begin{enumerate}
\item \textbf{(Mackey, orbit-wise.)} Proposition~\ref{prop:levelorbit}(1) holds with $\rhon(h_{j_O}(\gamma))$ replaced by
  $\rho_{n,\chi}(h_{j_O}(\gamma))$: restriction to block $j_O$ is an isomorphism $M_O^{F}\to\VR^{F_{j_O}}$ for the \emph{twisted}
  action, with inverse $f_{j_O\cdot\gamma}:=\rho_{n,\chi}(h_{j_O}(\gamma))^{-1}f_{j_O}$; a free orbit modulo $F\cap\{\pm I\}$ contributes
  $\VR^{\pm I}$, which is $\VR$ if $\chi(-1)=(-1)^{n}$ and $0$ otherwise, with \textbf{no solve}.
\item \textbf{(Parity.)} If $2\in\cring^{\times}$ and $\chi(-1)\ne(-1)^{n}$ then $H^{\bullet}(\Ga_0(N);\VR(\chi))=0$.
\item \textbf{(The cokernel.)} If $6\in\cring^{\times}$ and $\chi(-1)=(-1)^{n}$ then
  \[
    H^{1}(\Ga_0(N);\VR(\chi))\cong M\big/\bigl(M^{\langle S\rangle}+M^{\langle U\rangle}\bigr),\qquad H^{0}=M^{\langle S\rangle}\cap
    M^{\langle U\rangle},
  \]
  with the averaging idempotents of Proposition~\ref{prop:levelorbit}(2) ($M_{N,\chi}(S)^{2}=M_{N,\chi}(-I)=1$); the isomorphism is
  Definition~\ref{def:kappaN}'s $\varkappa_N$ with $M_{N,\chi}$ in place of $M_N$. More generally, for any finite-index $H\le\Ga$ and
  $\cring[H]$-module $X$ free of finite rank, $H^{1}(H;X)\cong M^{\langle-I\rangle}/(M^{\langle S\rangle}+M^{\langle U\rangle})$, $M=\Ind_H^{\Ga}X$
  --- the form needed for the $\Ga_1(N)$ model, $-I\notin\Ga_1(N)$ for $N\ge3$.
\item \textbf{(Dimensions, (NB1)'s proved half.)} If $\cring$ is a field with $6\in\cring^{\times}$ containing a primitive $12$-th root of
  unity $\zeta_{12}$ (else pass to an extension, which changes no dimension), and $\chi(-1)=(-1)^{n}$:
  \[
    \dim H^{1}-\dim H^{0}=m_N(n+1)-\Bigl(\tfrac{m_N-\nu_2}{2}(n+1)+\sum_{\lambda}d_S(\chi(-\lambda))\Bigr)
    -\Bigl(\tfrac{m_N-\nu_3}{3}(n+1)+\sum_{\mu}d_U(\chi(\mu^{-1}))\Bigr),
  \]
  the first sum over the $\nu_2=\nu_2(N)$ roots $\lambda$ of $\lambda^{2}+1\equiv0\pmod N$, the second over the $\nu_3$ roots $\mu$ of
  $\mu^{2}-\mu+1\equiv0\pmod N$ (the fixed points of $S$, $U$ on $\mathbb{P}^1(\Z/N)$, $(c:d)S=\lambda(c:d)$, $(c:d)U=\mu(c:d)$), and
  $d_S(z):=\#\{0\le m\le n: z\,i^{\,n-2m}=1\}$, $d_U(z):=\#\{0\le m\le n:z\,\zeta_6^{\,n-2m}=1\}$ with $i:=\zeta_{12}^{3}$,
  $\zeta_6:=\zeta_{12}^{2}$, the eigenvalue labels fixed by $\rhon(S)\sim\mathrm{diag}(i^{\,n-2m})_m$, $\rhon(U)\sim\mathrm{diag}(\zeta_6^{\,n-2m})_m$.
  At $\chi=1$ this is Proposition~\ref{prop:levelorbit}(3). The right-hand side depends on $\chi$ only through the values at the
  $\nu_2+\nu_3$ fixed points.
\end{enumerate}
\end{proposition}

\begin{proof}
(1) The proof of Proposition~\ref{prop:levelorbit}(1) uses only the cocycle identity and the block formula, both valid for
$\rho_{n,\chi}$ (Lemma~\ref{lem:inducedchar}); on a free orbit the stabiliser is $F\cap\{\pm I\}$, acting through
$\rho_{n,\chi}(-I)=\chi(-1)(-1)^{n}$ (Lemma~\ref{lem:nebdata}(3)). (2) $-I$ is central in $\Ga_0(N)$ and acts on $\VR(\chi)$ by the
scalar $\varepsilon:=\chi(-1)(-1)^{n}$; conjugation by an element of the group acts trivially on its cohomology, and it acts by
$\varepsilon$ through the coefficients, so $(\varepsilon-1)$ kills $H^{\bullet}$; $\varepsilon-1=-2$ is a unit. (3) Lemma~\ref{lem:treeres}
for $H$ and Shapiro: $H^{\bullet}(H;X)=H^{\bullet}(\Hom_\Ga(C_\bullet(\mathcal Y),M))$ with $\Hom_\Ga(C_1,M)=M^{\langle-I\rangle}$ (the
stabiliser of $\mathfrak e$), $\Hom_\Ga(C_0,M)=M^{\langle S\rangle}\oplus M^{\langle U\rangle}$ and coboundary $(a,b)\mapsto b-a$, exactly as in
the proof of Definition~\ref{def:kappaN}'s claims; for $H=\Ga_0(N)$ and $\chi(-1)=(-1)^{n}$, $M^{\langle-I\rangle}=M$. (4) The exact
sequence $0\to H^{0}\to M^{\langle S\rangle}\oplus M^{\langle U\rangle}\to M\to H^{1}\to0$ and (1). $S$ and $U$ have orders $4$, $6$ with
$S^{2}=U^{3}=-I$ acting trivially on $\mathbb{P}^1(\Z/N)$, so orbits are free modulo $\pm I$ (sizes $2$, $3$) or fixed points. At a fixed
point $j$ of $S$, $h:=h_j(S)=w_jSw_j^{-1}\in\Ga_0(N)$ and its bottom row is $\mathrm{row}_2(w_j)Sw_j^{-1}\equiv\lambda\,\mathrm{row}_2(w_j)w_j^{-1}=
\lambda(0,1)$, so $d_h\equiv\lambda$ and $a_h\equiv\lambda^{-1}=-\lambda$; $h$ is $\Ga$-conjugate to $S$, so $\rhon(h)$ has the
eigenvalues of $\rhon(S)$ (semisimple, its order $4$ being a unit), and the fixed space of $\chi(-\lambda)\rhon(h)$ has dimension
$d_S(\chi(-\lambda))$. Likewise for $U$ with $a_h\equiv\mu^{-1}$. The fixed points of $S$ are the classes $(1:\lambda)$ with
$\lambda^{2}\equiv-1$, of $U$ the $(1:\mu)$ with $\mu^{2}-\mu+1\equiv0$, one each per root (the computation in the proof of
Proposition~\ref{prop:levelorbit}(3), done modulo $N$ by the Chinese remainder theorem; a class $(c:d)$ with $c$ a non-unit is not fixed,
as then $\lambda$ would be a non-unit). $\rhon(S)$ acts on $x^{n-m}y^{m}$ with eigenvalues $i^{\,n-m}(-i)^{m}=i^{\,n-2m}$ after
diagonalising $S\sim\mathrm{diag}(i,-i)$, and $\rhon(U)$ with $\zeta_6^{\,n-2m}$ ($U\sim\mathrm{diag}(\zeta_6,\zeta_6^{-1})$).
\end{proof}

\begin{remark}[The Heilbronn action with a character]\label{rmk:nebheilbronn}
\PROVED{} Proposition~\ref{prop:heilbronn} holds for the character model with
\[
  H_{N,p,\chi}\;:=\;\sum_{\epsilon\in\mathrm{Mer}_p}\ \sum_{j\ \text{admissible}}E_{j\cdot\epsilon,\,j}\otimes\chi\bigl(d_{h_j(\epsilon)}\bigr)\,
  \rhon\bigl(\adjg h_j(\epsilon)\bigr)
\]
in place of \eqref{eq:heilbronn} ($a_{\adjg h}=d_h$), on $M_{N,\chi}/\mathrm{Rel}_{N,\chi}$, $\chi(-1)=(-1)^{n}$, inducing $T^{(N,\chi)}_p$ for $p\nmid
N$ and $U^{(N,\chi)}_p$ for $p\mid N$. \emph{Proof.} Lemma~\ref{lem:levelduality} and the proof of Proposition~\ref{prop:heilbronn} use of
$\rhon$ only its multiplicativity on the matrices they evaluate, and those are: elements of $A=\Ga_0(N)$; the coset representatives
$\beta\in AgA\subseteq\De_0(N)$; and, in Step~3, $w_i\gamma_\epsilon^{-1}\beta=\adjg h_j(\epsilon)$, where $\gamma_\epsilon=hw_i$ with $h\in A$
gives $w_i\gamma_\epsilon^{-1}\beta=h^{-1}\beta\in\De_0(N)$ --- so every one lies in $\De_0(N)$, where $\rho_{n,\chi}$ is multiplicative, and the
conjugation of Lemma~\ref{lem:levelduality}(5) is used only for $g\in\De_0(N)$. With $\rho_{n,\chi}$ for $\rhon$ throughout, the proofs
are unchanged (the hypotheses ``$n$ even, $-I\in A$'' of Lemma~\ref{lem:levelduality} are used there only through $M_N(-I)=1$, which is $\chi(-1)=(-1)^{n}$ here; v1.15.0-pre.2, WP-4 referee A2). In particular $d_{h_j(\epsilon)}$ is a unit modulo $N$ at every admissible pair. \qed
\end{remark}

\subsection{The cost of the route}\label{sec:nebcostsec}

\begin{remark}[Cost; scalability in $N$, $k$, $\operatorname{ord}\chi$ and $q$]\label{rmk:nebcost}
A model of operation counts, in the notation of Remarks~\ref{rmk:levelcost} and \ref{rmk:heilbronncost}; not an assertion about running
times. (1) \textbf{The character costs nothing structurally}: the character model has the fibre rank $D_N=m_N(n+1)$ of the $\Ga_0(N)$
model, the same block permutation and the same number of non-zero blocks; each block is $\rhon(h)$ times a table-lookup scalar
$\chi(a_h)$ ($\chi$ is a table of $\varphi(N)$ entries). Orbit-wise, the free orbits need no solve (Proposition~\ref{prop:neborbit}(1)), the
$\nu_2+\nu_3\le2^{\omega(N)+1}$ fixed points each one $(n+1)$-dimensional fixed space (the $\chi(a_h)^{-1}$-eigenspace of $\rhon(h)$, one $d\times d$ solve per distinct scalar, then $O(d^{2})$ per fixed point as in Remark~\ref{rmk:levelcost}(3)). Every count of Remarks~\ref{rmk:levelcost} and
\ref{rmk:heilbronncost} carries over with $\ell$-operations replaced by $\F_q$-operations. (2) \textbf{The field}: an $\F_q$ operation is
$O(d^{2})$ $\F_\ell$-operations naively ($O(d\log d)$ with fast multiplication); take $q=\ell^{d(\chi)}$, the least field of
Definition~\ref{def:nebdata}(4); $d(\chi)=1$ iff $\operatorname{ord}\chi\mid\ell-1$ (every $\omega^{i}$ at $N=\ell$, and every quadratic
character). One character per Frobenius orbit (Remark~\ref{rmk:nebfield}(3)). (3) \textbf{All characters at once} costs $\varphi(N)$
character-model runs, which is exactly the size of one $\Ga_1(N)$-model run ($D_{\Ga_1}=\varphi(N)D_N$) but with the elimination cubic
in $D_N$ instead of $\varphi(N)D_N$: a gain of $\varphi(N)^{2}$ on the dense route. The $\Ga_1(N)$ model is therefore specified as an
oracle ((NB3)), not as a production route. (4) Over $W_K(\F_q)$ the module route of $\S$\ref{sec:levelN}'s code map applies unchanged
with $\Z/\ell^{K}$ replaced by the Galois ring, whose Howell forms exist because $W_K(\F_q)$ is a chain ring with the same ideal lattice
$(\ell^{j})$ (Lemma~\ref{lem:teichmuller}(1)); \textbf{built at v1.18.0-pre.2} (slot-18 WP-6, \file{Hecke/HeckeNebentypusGalois}, the code map's item~(iv),
dense and sparse, $K\le3$, $d\le4$; v1.18.0-pre.1 read ``specified, not built'').
\end{remark}

\begin{contractbox}
\textbf{Refusal rules of \S\ref{sec:nebentypus}}, each by name. (i) $\ell\in\{2,3\}$: cochain level only (Proposition~\ref{prop:nebcorr}(2)'s
first half, every $\cring$); no comparison with group cohomology, no cokernel, no dimension formula is asserted, and nothing
defers to the Bredon object (\code{nebentypus\_small\_ell}). (ii) A character of $\ell$-power order into $\F_q$ or $W_K(\F_q)$ (\code{character\_wild},
Remark~\ref{rmk:nebfield}(2)): still refused.  Over the wild rings $\mathcal O_{j,e}$ of Definition~\ref{def:nebwild}
a character $\chi_{i,t}$ is admitted; refused there by name are a character whose wild order does not divide $\ell^{j}$
(\code{wild\_ring\_order}), a requested $\pi$-adic precision above the ring's $e$ (\code{wild\_precision}), and a tame
level $N\ne\ell^{r}$ (\code{wild\_tame\_level}: tame level $N>1$ with wild characters is slot-23 P11, not this text). (iii) A character value outside the declared $\F_q$ (the field must contain
$\mu_{\operatorname{ord}\chi}$, \code{character\_field}). (iv) \textbf{Not a refusal}: $\chi(-1)\ne(-1)^{n}$ gives $0$ by
Proposition~\ref{prop:neborbit}(2) --- a theorem, reported as such; and $n$ odd with $\chi$ odd is \emph{allowed} (odd weight), unlike
level~$N$ with trivial character. (v) \textbf{The MP constraint} (owner, 2026-09-28): nothing in this chapter uses the $7$-vertex model,
a trivial central action or Bredon semantics --- every statement is for an arbitrary $(K,\Rel)$ satisfying \ref{H4}
(\S\ref{sec:induced}'s ``parametric in the triangulation''), the central element $-I$ enters only through its scalar
$\chi(-1)(-1)^{n}$ on the coefficients, and the refined Moerdijk--Pronk model with a non-trivial central action is admissible input.
\end{contractbox}

\subsection{The predictions}\label{sec:nebpredsec}

\begin{remark}[What is checkable with nebentypus: (NB1)--(NB6)]\label{rmk:nebcheckable}
Committed before slot~16's code (rule~20), with the convention of Remark~\ref{rmk:hidacheckable}: a failing \PROVED{} item stops the package, a failing \CONJ{} item is a finding. Cells are $(N,\chi,n)$ with $k=n+2$, over $\F_q$, $\ell\ge5$, $\ell\nmid N$ unless said.
\begin{enumerate}[label=(NB\arabic*)]
\item\label{NB1} \PROVED{} (Proposition~\ref{prop:neborbit}(4)) the Euler characteristic of the character model; and for $n>0$
  \STD{\cite[Ch.~8]{shimura} for Eichler--Shimura with a character; \cite{cohen_oesterle} for $\dim S_k(\Ga_0(N),\chi)$; both
  \textbf{not} verified; v1.21.0-pre.2: \cite[Th\'eor\`eme~1, pp.~71--72]{cohen_oesterle} read ($\dim S_k-\dim M_{2-k}$; the Eisenstein dimension is not in
  that paper), and \cite[Thm.~8.4]{shimura} covers the cuspidal part only (R20-L)} the prediction $\dim H^{1}-\dim H^{0}=2\dim S_k(\Ga_0(N),\chi^{\pm1})+\dim\mathcal E_k(\Ga_0(N),\chi^{\pm1})$ ($\mathcal E_k$
  the Eisenstein space), the sign of the exponent immaterial for dimensions (complex conjugation). Reason: over a field of
  characteristic $\ell\ge5$ the right-hand side of Proposition~\ref{prop:neborbit}(4) equals the characteristic-$0$ one for $\mathfrak{t}\circ\chi$
  (the fixed-space dimensions of the order-$4$ and order-$6$ elements are Brauer-character multiplicities, invariant under reduction),
  and over $\Cx$, $H^{0}=0$ for $n>0$.
\item\label{NB2} \PROVED{} (Proposition~\ref{prop:nebdecomp}(4)) $\chi=1$ reproduces Definition~\ref{def:levelNmodel}'s level-$N$
  model and operators \emph{bit for bit}: relation matrices, cochain maps, $h^{0}$, $h^{1}$ and every Hecke matrix.
\item\label{NB3} \PROVED{} (Proposition~\ref{prop:nebdecomp}(3)) with $\lambda(N)\mid q-1$ and $\ell\nmid\varphi(N)$:
  $h^{\bullet}(\Ga_1(N))=\sum_\chi h^{\bullet}(\Ga_0(N),\chi)$, $\langle d\rangle=\chi(d)$ on the $\chi$-part, and the characteristic
  polynomial of $T_p$ on the $\Ga_1(N)$ model is the product over $\chi$ of those on the character models. Cells: $N\in\{5,7,13\}$,
  $n\le10$, $\ell\in\{5,7,11,13,29\}\setminus\{N\}$ with $\lambda(N)\mid q-1$.
\item\label{NB4} \CONJ{} \textbf{The weight-shift lemma, on cohomology.} For $N=\ell\ge5$, $\omega\colon G_\ell=\F_\ell^{\times}\to
  \F_\ell^{\times}$ the identity (so $q=\ell$), $0\le i\le\ell-2$, $n\ge0$ with $(-1)^{i}=(-1)^{n}$: every system of $T_p$-eigenvalues
  ($p\ne\ell$) in $H^{1}(\Ga_0(\ell);V_{\F_\ell}(\omega^{i}))\otimes\bar\F_\ell$ is of the form $(p^{e}b_p)_p$ with $(b_p)_p$ a system in
  $H^{1}(\Ga;\Sym^{m}\bar\F_\ell^{2})$ at level~$1$ for some $m$ and $e$; the prediction to be tested is the sharp form \emph{at $n=0$: the systems of weight $i+2$ at level $1$ ($m=i$, $e=0$) and $p^{\,i}\cdot$ the systems of weight $\ell+1-i$ ($m=\ell-1-i$, $e=i$), the two constituents $\Sym^{i}$ and $\Sym^{\ell-1-i}\otimes\det^{i}$ of $\Ind_{B}^{\SL_2(\F_\ell)}\omega^{i}$ as $\De$-modules; for $n>0$, the same with $\Sym^{n}\otimes$ each constituent in place of it, whose own constituents are left to the first run} (v1.15.0-pre.2, WP-4 referee A4) (the weight ``$k+i$'' of
  \file{docs/implementation/slot14/WP-C16\_dihedral\_over\_extension\_fields\_plan.md} M8, in its normalisation, and the other
  constituent of $\Ind_{B(\F_\ell)}^{\SL_2(\F_\ell)}\omega^{i}$ with the $\theta$-twist $e>0$). Mechanism: $\Ga_0(\ell)\to B(\F_\ell)$ and
  the Jordan--H\"older filtration of $\Sym^{n}\otimes\Ind_{B}^{\SL_2(\F_\ell)}\omega^{i}$ (the long exact sequence of
  Lemma~\ref{lem:treeres}'s length-$1$ resolution makes every system of $H^{1}$ of the middle term a system of $H^{1}$ of a
  constituent, or of an $H^{0}$). \emph{Not proved here}: the constituents with their $\det$-twists as $\De$-modules for $n>0$ (at $n=0$ they are
  derived in (c) below, v1.17.0-pre.3, and the sharp form at $n=0$ is \PROVED{} there given one \STD{} input, v1.17.0-pre.4, whose locator
\cite[Lemmas~1.1.1, 1.1.4]{ash_stevens} is verified since v1.17.0-pre.5). At $n=0$ the prediction reads: weight~$2$, level $\ell$, character $\omega^{i}$
  $\leftrightarrow$ level~$1$, weight $i'+2$.

  \emph{(v1.17.0-pre.1; slot-16 WP-11, \file{docs/implementation/slot16/WP-11\_report.md}~\S1, \S4; proposed for the referee, slot-17 WP-3.)}
  \textbf{(a) The sharp form at $n=0$ as data.} \COMP\textsuperscript{r} (\code{HeckeNebentypusValidation}, arm (X4), gated, every run; the cokernel
  route at level~$\ell$ against the level-$1$ cokernel, $T_p$ for $p\in\{2,3,5,7\}\setminus\{\ell\}$) For every $\ell\in\{5,7,\dots,43\}$ and every even
  $i\le\ell-3$ ($132$ cells), with $\omega^{i}$ acting through $\chi(a_\gamma)$ (Definition~\ref{def:nebdata}(5)),
  \[
    \mathrm{charpoly}\bigl(T_p\mid H^{1}(\Ga_0(\ell);\F_\ell(\omega^{i}))\bigr)=\mathrm{charpoly}\bigl(T_p\mid H^{1}(\Ga;\Sym^{i}\F_\ell^{2})\bigr)\cdot
    \mathrm{charpoly}\bigl(p^{\,i}\,T_p\mid H^{1}(\Ga;\Sym^{\ell-1-i}\F_\ell^{2})\bigr),
  \]
  held on $132$ of $132$. \textbf{The orientation it pins}: under the $\chi(a_\gamma)$ convention $\Sym^{i}$ is the \emph{untwisted} constituent and
  $\Sym^{\ell-1-i}\otimes{\det}^{i}$ the twisted one; the swapped assignment ($\Sym^{\ell-1-i}$ untwisted, $\Sym^{i}\otimes{\det}^{\ell-1-i}$) holds only on
  the $18$ cells where the two coincide.
  \textbf{(b) The reading at $n>0$.} \COMP\textsuperscript{r} (the same arm; \emph{reported, not gated}) At $n\in\{1,2\}$ ($240$ cells, $\ell\le43$) the
  constituents are those of Clebsch--Gordan, $\Sym^{n}\otimes\Sym^{i}=\bigoplus_j\Sym^{n+i-2j}\otimes{\det}^{j}$ and
  $\Sym^{n}\otimes\Sym^{\ell-1-i}\otimes{\det}^{i}=\bigoplus_j\Sym^{n+\ell-1-i-2j}\otimes{\det}^{i+j}$, where it holds over $\F_\ell$ ($n\le i\le\ell-1-n$); the
  product of their $H^{1}$ characteristic polynomials (each twisted constituent's $T_p$ multiplied by $p^{\text{twist}}$) matched on $216$ cells as it
  stands and on the $24$ cells $i=n$ after removing one factor $X-p^{\,n}(1+p)$: the constituent $\Sym^{0}\otimes{\det}^{n}$ carries $H^{0}=\F_\ell$ with
  $T_p=p^{\,n}(1+p)$, and the connecting map of the long exact sequence removes that class ($240$ of $240$). This is the ``or of an $H^{0}$'' clause of the
  mechanism above; its form was found by the run.
  \textbf{(c) A proof of (a), given one standard input.} \PROVED{} (given the \textsf{[standard]} input named; owner ruling of 2026-10-02 on the slot-17
  referee F1's M-1, v1.17.0-pre.4; the input's locator \textbf{verified} at v1.17.0-pre.5, slot-17 R17-4) \emph{(Proposed at v1.17.0-pre.1; completed at
  v1.17.0-pre.3 by the referee F1's repair M-1, \file{docs/implementation/slot17/WP-3\_referee\_F1.md}~\S4: the orientation is now derived, not pinned by the
  data.)} Given \STD{\cite[\S1.1]{ash_stevens}: the long exact
  sequence of $\Ga$-cohomology of an exact sequence of modules over the Hecke monoid, and Shapiro's isomorphism
  $H^{\bullet}(\Ga_0(\ell);\F_\ell(\omega^{i}))\cong H^{\bullet}(\Ga;\Ind_{\Ga_0(\ell)}^{\Ga}\F_\ell(\omega^{i}))$ (Lemma~\ref{lem:inducedchar}, Proposition~\ref{prop:nebcorr}), commute
  with the double-coset operators of a compatible pair --- \cite[Lemma~1.1.1, p.~194]{ash_stevens} (the long exact sequence, for right or left modules over the
  monoid) and \cite[Def.~1.1.2, p.~194; Lemma~1.1.4, p.~195]{ash_stevens} (Shapiro, for a compatible pair with $[\Ga:\Ga_0]<\infty$; its proof uses their
  Lemma~1.1.3); \textbf{verified} on the GDZ page images, slot-17 R17-4 (\file{docs/implementation/slot17/R17-4\_report.md}~\S1a, rows~1--2; v1.17.0-pre.5)},
  used for the pair $(\Ga_0(\ell),\De_0(\ell))\subset(\Ga,\{\delta\in\Mat_2(\Z):\det\delta>0,\ \ell\nmid\det\delta\})$, whose
  induced module is the module $I$ below. The pair is compatible in the sense of \cite[Def.~1.1.2]{ash_stevens}: (b) in their orientation, $\Ga\De_0(\ell)$ is the
  whole of the larger monoid, since $\mathrm{GL}_2(\F_\ell)=\SL_2(\F_\ell)\cdot B(\F_\ell)$ (the other orientation, $\De_0(\ell)\Ga$, holds as well); (c)
  $\Ga\cap\De_0(\ell)\De_0(\ell)^{-1}=\Ga_0(\ell)$, since $B(\F_\ell)B(\F_\ell)^{-1}=B(\F_\ell)$.

  \emph{The constituents, in this document's conventions.} By Proposition~\ref{prop:symconv}, $\rho_m(b)$ for $b=\smat{a}{b'}{0}{d}$ substitutes $x\mapsto ax$,
  $y\mapsto b'x+dy$: so $\rho_m(b)x^{m}=a^{m}x^{m}$, the span of the monomials divisible by $x$ is $B$-stable, and $\varphi_m$ := ``the coefficient of $y^{m}$''
  satisfies $\varphi_m(\rho_m(b)v)=d^{m}\varphi_m(v)$. Let $\psi(b):=a^{i}$, $I:=\{f\colon\mathrm{GL}_2(\F_\ell)\to\F_\ell:\ f(bg)=\psi(b)f(g)\ \forall b\in B(\F_\ell)\}$ with
  $(hf)(g):=f(gh)$ (Definition~\ref{def:inducedmodule}'s convention, with the $\chi(a_\delta)$ of Definition~\ref{def:nebdata}(5)), $0<i<\ell-1$, and $m:=\ell-1-i$;
  on $\Sym^{m}\otimes{\det}^{i}$ an element $g$ acts by $\det(g)^{i}\rho_m(g)$.
  ($\alpha$) $\Sym^{m}\otimes{\det}^{i}\to I$, $v\mapsto(g\mapsto\varphi_m(g\cdot v))$, is well defined ($d^{m}(ad)^{i}=a^{i}d^{\ell-1}=a^{i}$ on $B$), equivariant, and
  injective: $\varphi_m(g\cdot v)=\det(g)^{i}\,v(c_g,d_g)$, so if it vanishes for all $g$ then $v(c,1)=0$ for every $c\in\F_\ell$ and the coefficient of $x^{m}$
  vanishes; then $v(x,1)$ has degree $<m\le\ell-1$ and $\ell$ roots, so $v=0$.
  ($\beta$) $I\to\Sym^{i}$, $f\mapsto\sum_{Bg\in B\backslash\mathrm{GL}_2(\F_\ell)}f(g)\,\rho_i(g)^{-1}x^{i}$, is well defined ($\rho_i(b)^{-1}x^{i}=a^{-i}x^{i}$), equivariant, and onto:
  its image contains $x^{i}$, $y^{i}$ and every $(x+cy)^{i}$, which span $\Sym^{i}$ as $i<\ell$.
  ($\beta\circ\alpha=0$) An equivariant map $\Sym^{m}\otimes{\det}^{i}\to\Sym^{i}$ sends $x^{m}$ to a vector fixed by the upper unipotent group, hence to a
  multiple of $x^{i}$; the torus acts on $x^{m}\otimes{\det}^{i}$ by $a^{m+i}d^{i}=d^{i}$ and on $x^{i}$ by $a^{i}$, which differ for $0<i<\ell-1$; so $x^{m}\mapsto0$,
  and $x^{m}$ generates $\Sym^{m}$.
  Dimensions: $(\ell-i)+(i+1)=\ell+1=\dim I$. Hence
  \[
    0\to\Sym^{\ell-1-i}\otimes{\det}^{i}\to I\to\Sym^{i}\to0
  \]
  is exact: $\Sym^{i}$ is the \emph{untwisted} constituent, as (a) found. For $i=0$ the constant functions split off ($\ell+1\equiv1$ is a unit) and ($\alpha$)
  embeds $\Sym^{\ell-1}$ as a complement (its image contains no non-zero constant, as $(\Sym^{\ell-1})^{\SL_2(\F_\ell)}=0$). Had the coefficient character
  been $\chi^{-1}(d_\delta)$ (equal to $\chi(a_\delta)$ on $\Ga_0(\ell)$, not on $\De_0(\ell)$), the same computation would give $\Sym^{\ell-1-i}$ untwisted --- the
  swapped reading that (a) rejects on $114$ distinguishing cells.

  \emph{The long exact sequence.} $\Ind_{\Ga_0(\ell)}^{\Ga}\F_\ell(\omega^{i})$ is $I$ restricted along $\Ga\to\SL_2(\F_\ell)$, with the Hecke monoid acting through
  $\mathrm{GL}_2(\F_\ell)$; let $0\to A\to\Ind\to Q\to0$ be the sequence just proved. $H^{2}(\Ga;-)=0$ on $\F_\ell[\Ga]$-modules
  (Claim~\ref{cl:amalgam}(3), whose argument uses nothing about the module), so the long exact sequence ends $\to H^{1}(\Ga;Q)\to0$, and in the Grothendieck
  group of $\F_\ell[T_p:p\ne\ell]$-modules $[H^{1}(\Ind)]=[H^{1}(A)]+[H^{1}(Q)]-[H^{0}(A)]+[H^{0}(\Ind)]-[H^{0}(Q)]$. For $0<i\le\ell-3$ every $H^{0}$ vanishes:
  $H^{0}(\Ga;\Sym^{a}\F_\ell^{2})=(\Sym^{a})^{\SL_2(\F_\ell)}=0$ for $0<a\le\ell-1$, and $H^{0}(\Ga_0(\ell);\F_\ell(\omega^{i}))=0$ since $a_\gamma$ takes every value of
  $\F_\ell^{\times}$ on $\Ga_0(\ell)$. For $i=0$, $[H^{0}(\Ind)]=[H^{0}(\F_\ell)]$ cancels the constituent $\F_\ell$'s $H^{0}$ and $H^{1}(\Ga;\F_\ell)=0$. Since
  $H^{1}(\Ga;\Sym^{a}\otimes{\det}^{e})=H^{1}(\Ga;\Sym^{a})$ with $T_p$ multiplied by $p^{e}$, the characteristic polynomials multiply as displayed in (a).

  \emph{Tag at this version.} The sharp form at $n=0$ is \PROVED{} by (c), given the one \STD{} input above, and \COMP\textsuperscript{r} by (a)
  (v1.17.0-pre.4, owner ruling of 2026-10-02 on the referee F1's M-1: the tag the referee proposed, with the locator then flagged \textbf{not verified};
  at v1.17.0-pre.3 the referee's alternative, \COMP\textsuperscript{r} only until the locator is read, was applied instead). \emph{v1.17.0-pre.5:} the locator
  is read and \textbf{verified} on the GDZ page images (slot-17 R17-4, \file{docs/implementation/slot17/R17-4\_report.md}~\S1a): the long exact sequence is
  \cite[Lemma~1.1.1, p.~194]{ash_stevens} (not Lemma~1.1.3, which is restriction / corestriction), and Shapiro is \cite[Lemma~1.1.4, p.~195]{ash_stevens} for
  the compatible pairs of their Def.~1.1.2 (p.~194), checked for this pair above. The ``not verified'' flag is dropped; the \COMP\textsuperscript{r} fallback
  is not triggered. The general
  statement of (NB4) (every $n$, every system) stays \CONJ{}, its mechanism being the same long exact sequence over the constituents of $\Sym^{n}\otimes\Ind$
  plus the identification of the $H^{0}$ systems $p^{e}(1+p)$ with a level-$1$ $H^{1}$ system, which (b) shows is needed and which this version does not write
  out.
\item\label{NB5} \STD{LMFDB newform data; \textbf{not} verified, nothing transcribed at this version} \CONJ{} Classical
  characteristic polynomials: for $(N,\chi,k)\in\{(5,\chi_5,4),(13,\chi_{13},4),(7,\chi_{-7},3),(4,\chi_{-4},5),(3,\chi_{-3},7)\}$
  ($\chi_5$, $\chi_{13}$ the quadratic characters, even as $5\equiv13\equiv1\pmod4$), the characteristic polynomials of $T_2$, $T_3$ on the character model's $H^{1}$ equal, modulo $\ell$, the product over
  the newforms $f\in S_k(\Ga_0(N),\chi)$ of $(X-a_p(f))^{2}$ times the Eisenstein factors --- each Eichler--Shimura half contributes once;
  for real $\chi$ this is the form $(X-a_p(f))(X-\overline{a_p(f)})$ written through v1.16.0-pre.1, which for non-real $\chi$ is
  \textbf{wrong} (it is the characteristic polynomial on the holomorphic halves of the $\chi$- \emph{and} $\bar\chi$-parts, since
  $\overline{a_p(f)}=a_p(f^{\rho})$ with $f^{\rho}\in S_k(\Ga_0(N),\bar\chi)$; v1.16.0-pre.2, slot-16 WP-9 audit
  \file{docs/implementation/slot16/WP-9\_audit.md} \S3: the written form fails in $20$ of $24$ non-real cases, the squared form holds in
  $24$ of $24$). \PROVED{} \emph{(the squaring, in characteristic~$0$ through the Teichm\"uller lift; the identification of $S_k(\Ga_0(N),\chi)$ rather
  than $S_k(\Ga_0(N),\bar\chi)$ as the space whose eigenvalues occur is the eigenvalue dictionary, \PROVED{} modulo the named literature hypotheses
  (L1)--(L5) in \textsf{[GT,~\texttt{bridge:nebentypus-dictionary}]} (\file{lean/blueprint/src/gl2\_groundtruth.tex}; Lean, \code{proved\_mod\_lit}, audited by
  \file{docs/implementation/slot16/L16-2\_fable\_audit.md}), and \CONJ{} only for those literature inputs (Eichler--Shimura with a character,
  \cite[Ch.~8]{shimura}, VERIFY); the squared form holds whichever way it falls; v1.16.0-pre.3, WP-14 referee C-9 at the owner's ruling.)} Let $\varepsilon:=\smat100{-1}$. Then $\varepsilon\Ga_0(N)\varepsilon=\Ga_0(N)$ and $a_{\varepsilon\gamma\varepsilon}=a_\gamma$,
  so $\rho_{n,\chi}(\varepsilon\gamma\varepsilon)=\rhon(\varepsilon)\rho_{n,\chi}(\gamma)\rhon(\varepsilon)$ and the compatible pair
  $(\conj_\varepsilon,\rhon(\varepsilon))$ defines an involution $\iota_\varepsilon$ of $H^{1}(\Ga_0(N);\VR(\chi))$, of the $\chi$-part itself; it
  commutes with $T_p$ and $U_p$ because $\varepsilon\smat100p\varepsilon=\smat100p$ and the double coset is $\varepsilon$-stable. By Eichler--Shimura
  with a character \STD{\cite[Ch.~8]{shimura}}, $H^{1}_{\mathrm{par}}\otimes\Cx=H\oplus H^{\mathrm{anti}}$ with $H$ the Hecke-stable image of the holomorphic
  forms and $H^{\mathrm{anti}}$ the antiholomorphic half --- for non-real $\chi$ \emph{not} the complex conjugate of $H$, which lives in $\VR(\bar\chi)$ (WP-14 referee
  C-8) --- and $\iota_\varepsilon$ maps $H$ onto $H^{\mathrm{anti}}$ (the holomorphic class of $f$ to the antiholomorphic class of $f^{\rho}(z)=\overline{f(-\bar z)}$),
  $H\cap\iota_\varepsilon H=0$, $H\oplus\iota_\varepsilon H=H^{1}_{\mathrm{par}}$. So $h\mapsto h\pm\iota_\varepsilon h$ are Hecke-equivariant
  isomorphisms of $H$ onto the $\pm1$-eigenspaces of $\iota_\varepsilon$ on $H^{1}_{\mathrm{par}}$, and $\mathrm{charpoly}(T\mid H^{1}_{\mathrm{par}})=
  \mathrm{charpoly}(T\mid H)^{2}$, with $H$ carrying the eigenvalues of a single space of forms of character $\chi$ or $\bar\chi$. The passage from characteristic~$0$ (through the
  Teichm\"uller lift) to $\F_q$ uses $H^{0}(\Ga_0(N);V_{\F_q}(\chi))=0$, true on the (NB5) cells ($n\le5$; WP-9 audit \S1); the \PROVED{} covers the squaring
  only.
  \textbf{The reduction rule}: a non-real character given by a label (Conrey $N.m$) is reduced modulo $\ell$ by its value at the least
  primitive root $g$ modulo $N$ ($g=3,2,2$ at $N=7,9,13$): $\psi(g):=s$, the \emph{least} primitive $\operatorname{ord}\chi$-th root of unity in $\F_\ell$
  (in $[1,\ell)$), $\bar\psi$ its inverse --- a rule that presumes $d(\chi)=1$ ($\operatorname{ord}\chi\mid\ell-1$, Definition~\ref{def:nebdata}(4)), which holds at
  the (NB5) cells (C-8). Without this rule the mod-$\ell$ reading is ambiguous, since choosing $s^{-1}$ is the $\chi\leftrightarrow\bar\chi$
  swap itself. The values were transcribed with LMFDB labels by slot-16 WP-9 / WP-9b (\file{docs/implementation/slot16/WP-9\_audit.md}) before
  slot~16's code package ran.
\item\label{NB6} \STD{\cite{ribet77} for CM forms with nebentypus; LMFDB for $\dim=1$; \textbf{not} verified} \CONJ{}
  \textbf{CM cells.} At $(4,\chi_{-4},5)$ (CM by $\Q(i)$), $(3,\chi_{-3},7)$ (by $\Q(\sqrt{-3})$), $(7,\chi_{-7},3)$ (by $\Q(\sqrt{-7})$),
  where the cusp space is predicted to be one CM newform: for every prime $p\nmid N\ell$ \emph{inert} in the CM field, $T^{(N,\chi)}_p$
  acts on the cuspidal part of $H^{1}$ with characteristic polynomial $X^{2}$ ($a_p=0$), and for $p$ split the two eigenvalues are
  $\pi^{k-1}+\bar\pi^{k-1}$, $p=\pi\bar\pi$, reduced modulo $\ell$. These are the sanity cells of slot~16's [U] part~2 (CM Hida
  families need exactly such level and character).
\end{enumerate}
\end{remark}

\paragraph{Scope: what this chapter does not do.}
It asserts, for every $N$, every finite field $\F_q$ of characteristic $\ell$ and every $\chi\colon G_N\to\F_q^{\times}$ (and its
Teichm\"uller lift over $W_K(\F_q)$): the two models, the diamonds and the $\chi$-decomposition, the Hecke operators $T_p$, $U_p$,
$T(p,p)$ on the level-$1$ $\tau$-table, the orbit-wise fixed spaces, the cokernel and the Euler characteristic, all \PROVED{} for
$6\in\cring^{\times}$ (cochain level for every $\cring$); the relations are standard (Theorem~\ref{thm:nebrelations}). Not asserted: anything at $\ell\in\{2,3\}$ beyond the
cochain level; ramified coefficient rings; the degeneracy maps between levels with characters (Proposition~\ref{prop:levelcorrespondence}(4)(b),(c)
with $\rho_{n,\chi}$ applies verbatim when $N'\mid N$ and $\chi$ factors through $G_{N'}$, not written out); Atkin--Lehner operators
(they move $\chi$ to $\bar\chi$ and need $\De$-matrices outside $\De_0(N)$); the weight-shift lemma ((NB4) is a prediction; since v1.17.0-pre.1 its
sharp form at $n=0$ is data, \COMP\textsuperscript{r}, and since v1.17.0-pre.4 \PROVED{} given one \STD{} input, whose locator is verified since v1.17.0-pre.5,
\ref{NB4}(a), (c)); the
comparison with classical forms beyond (NB1), (NB5), (NB6).

\paragraph{Code map.}
\textbf{Built at v1.16.0} (slot-16 WP-10 and WP-11, \file{docs/implementation/slot16/WP-10\_report.md}, \file{docs/implementation/slot16/WP-11\_report.md};
cross-model audit \file{docs/implementation/slot16/WP-10\_WP-11\_audit.md}, repaired; not built at v1.15.0-pre.1 or at v1.15.0 --- contract first, rule~21, owner ruling Q3
of 2026-09-30). \file{Hecke/HeckeNebentypus} (\code{hecke::neb}): an $\F_q$ context with its linear algebra, a Galois-ring context with the Teichm\"uller lift,
Dirichlet characters and the reduction rule of \ref{NB5}, the character model as a cokernel with $H_{N,p,\chi}$ (Remark~\ref{rmk:nebheilbronn}), \eqref{eq:MHchi} and
\eqref{eq:nebcorr} on the fixed model, the $\Ga_1(N)$ table and the diamonds (Lemma~\ref{lem:gammaonetable}, Proposition~\ref{prop:nebdecomp}(1)); no coset or
$\tau$-table is added. \code{HeckeNebentypusValidation} (\code{standard}) carries (NB1)--(NB6): Proposition~\ref{prop:neborbit}(4) on every $\chi$ modulo $N\le30$,
$n\le5$; \ref{NB2} bit for bit; \ref{NB3} on the $\Ga_1(N)$ model; \ref{NB5} against the committed table \file{docs/implementation/slot16/WP-9\_nb\_predictions.tsv},
every non-real distinguishing case reading $\chi$, not $\bar\chi$; \ref{NB6} on its three cells. \ref{NB4} (WP-11): the sharp form at $n=0$ is \textbf{gated} and
held on $132$ cells, $\ell\le43$ (under the $\chi(a_\gamma)$ convention $\Sym^{i}$ is the untwisted constituent); for $n=1,2$ the reading is recorded, not gated. Also
WP-11, on level~$1$ blocks (readings of \S\ref{sec:saturation}, not statements of this chapter): $m=2$ and $H^{1}_{\mathfrak m}\cong A_W^{2}$ on the non-ordinary
dihedral blocks of slot-14 WP-C16 at $(72,47)$, $(118,47)$, $(156,103)$, with $A_W\cong\F_q[\varepsilon]/(\varepsilon^{2})$ the block's local algebra (certified over $\Z_\ell$ as
$\Lambda_{\mathfrak m}\cong\Talg'^{\,2}_{\mathfrak m}$ by Lemma~\ref{lem:heckerankone} and Corollary~\ref{cor:stablecert}); at the F1 hidden blocks $(220,17)$, $(246,19)$
$m=2$ but $H^{1}_{\mathfrak m}\not\cong\Talg_{\mathfrak m}^{2}$ over the full Hecke algebra (a generation certificate by the $T_p$, $p\le\dim S_k$). The tags of (NB1)--(NB6) are unchanged at v1.16.0.
\textbf{Built at v1.17.0-pre.3} (slot-17 WP-6, merged and audited --- \file{docs/implementation/slot17/WP-6\_report.md}, cross-model audit
\file{docs/implementation/slot17/WP-6\_audit.md}, ACCEPT; contract first at v1.17.0-pre.1, rule~21; the text below is WP-6's proposal with the slot-17 referee F1's
repairs M-2, m-W1--m-W3, \file{docs/implementation/slot17/WP-3\_referee\_F1.md}~\S4):
(i) \textbf{the sparse cokernel with a character} --- \file{Hecke/HeckeNebentypusSparse} (\code{NebSparseCokernel}): the elimination of
Remark~\ref{rmk:levelcost}(4) on the generators of Proposition~\ref{prop:neborbit}(1) over $\F_q$, the pivot of every reduced generator its
smallest column, so that the result is the natural reduced row echelon form of $\mathrm{Rel}_{N,\chi}$, hence equal to the dense route entrywise
wherever both run (checked: pivots and free columns, every row, every $H_{N,p,\chi}$ for $p\le7$, $p\ne\ell$, and $T_4$ / $U_4$, on $13$ cells and one
parity-zero cell with $d(\chi)\in\{1,2\}$; the full \ref{NB3} grid of $132$ cells reads the same verdicts on the sparse route). It lifts the cap $D\le4096$
(reached: $N=30030$, $k=4$, $D_N=290\,304$, over $\F_{17}$; $d(\chi)=2$ at $N=2003$, $k=12$). Over $W_K(\F_q)$ it is \textbf{not built} (with (iii)). The route
is a choice, never a value: \code{route=auto} is the dense route at $D\le4096$.
(ii) \textbf{$T^{(N,\chi)}_{p^{\nu}}$, $\nu\ge2$}: the operator $H_{N,p^{\nu},\chi}$ defined by Remark~\ref{rmk:nebheilbronn}'s formula with $\mathrm{Mer}_{p^{\nu}}$
(Lemma~\ref{lem:merelpath} holds for every determinant) in place of $\mathrm{Mer}_p$, every element of determinant $p^{\nu}$ included. That it induces the full
$T_{p^{\nu}}$ is Merel's theorem at composite determinant and is \textbf{not} a statement of this document (Proposition~\ref{prop:heilbronn} is for a prime); it is
\textbf{tested, not proved}: it agrees with the recursion of Theorem~\ref{thm:nebrelations}(2) for $T_{p^{2}}$, $T_{p^{3}}$ at $p\in\{2,3\}$ and $T_{25}$, and
$U_{p^{\nu}}=(U_p)^{\nu}$ (Theorem~\ref{thm:nebrelations}(3)), on six cells; the recursion without $\chi(p)$ fails wherever $\chi(p)\ne1$. The model-route double
coset $[\Ga_0(N)\,\mathrm{diag}(1,p^{\nu})\,\Ga_0(N)]_\chi$ of Proposition~\ref{prop:nebcorr} is \textbf{not built}.
(iii) Galois-ring cohomology at $d(\chi)\ge2$ (Remark~\ref{rmk:nebcost}(4)): \textbf{not built}. No coset or $\tau$-table is added; no statement, tag or label
changes (Theorem~\ref{thm:nebrelations}(2)'s test moves from ``untested'' to ``tested on six cells'').
\textbf{Built at v1.18.0-pre.2} (slot-18 WP-6, \file{docs/implementation/slot18/WP-6\_report.md}, merged; contract first at v1.18.0-pre.1, rule~21):
(iv) \textbf{built}: \file{Hecke/HeckeNebentypusGalois} --- cohomology over $W_K(\F_q)$, $K\le3$, $d\le4$, on the dense and the sparse Howell route (equal
entrywise on $8$ cells, $d\in\{2,3,4\}$); $K=1$ equals the $\F_q$ route entrywise; at $\chi=1$, $d=1$ it equals the fixed model over $\Z/\ell^{K}$ ($5$ cells, $3$
with torsion); the $\Ga_1(N)$ model over $\Z/\ell^{K}$ equals the union over the characters (Proposition~\ref{prop:nebdecomp}(3)) at $N=5,7$. (v) \textbf{built}:
the model-route $T_{p^{2}}$ through the shipped \code{Gamma0LevelN} table on $6$ cells / $11$ pairs, equal to the recursion $T_p^{2}-\chi(p)p^{\,n+1}$ entrywise
and to (ii)'s $\mathrm{Mer}_{p^{2}}$ operator by characteristic polynomial. (vi) \textbf{done}: the $D=193\,536$ ladder cell on a second route (entrywise), and a
natural \code{\_\_int128} escalation in the default run. No statement of this chapter changes; (ii)'s operator stays \emph{tested, not proved}.
The plan, as committed at v1.18.0-pre.1 (\file{docs/implementation/slot18/00-slot18-plan.md}; contract first, rule~21; ``not built'' at that version; kept for
the record; no statement of this chapter changes): (iv) \textbf{cohomology over $W_K(\F_q)$ at $d(\chi)\ge2$}, $K\le3$ (Remark~\ref{rmk:nebcost}(4); scope~(i)'s
$W_K(\F_q)$ half), first on the dense route, then on the sparse route of~(i). What the package must keep: every statement used is already one of this chapter
over a commutative ring ($\cring=W_K(\F_q)$, $\chi$ valued in its Teichm\"uller lift, Lemma~\ref{lem:teichmuller}); the elimination is by Howell forms over the
chain ring (Remark~\ref{rmk:nebcost}(4)), and the result is read as lengths and invariant factors, never as a dimension; two checks are gated where both routes
run --- the reduction $K=1$ equals the $\F_q$ route entrywise, and at $\chi=1$ (so $d=1$, $W_K(\F_\ell)=\Z/\ell^{K}$) the $\Ga_0(N)$ character model equals
\S\ref{sec:levelN}'s model over $\Z/\ell^{K}$ (\ref{NB2} over that ring; and, at $N=1$, \code{EquivariantCohomologyOnK} over $\Z/\ell^{K}$) bit for bit; and the sparse route over
$W_K(\F_q)$ equals the dense one entrywise on every cell where both run. A disagreement is a finding, never a re-pin. (v) \textbf{The model route for
$T_{p^{2}}$}: the double coset $[\Ga_0(N)\,\mathrm{diag}(1,p^{2})\,\Ga_0(N)]_\chi$ of Proposition~\ref{prop:nebcorr} on the fixed model, built from the shipped
\code{Gamma0LevelN} table (\code{BuildLevelCorrespondenceChi} on a $\mathrm{diag}(1,p^{2})$ datum) --- \textbf{never a second coset or $\tau$ construction}. Since
$\mathcal H(\Ga_0(N),\De_0(N))\ni T_{p^{2}}=[\mathrm{diag}(1,p^{2})]+[\mathrm{diag}(p,p)]$ for $p\nmid N$ and $[\mathrm{diag}(p,p)]=T(p,p)$ acts as $\chi(p)p^{\,n}$
(Proposition~\ref{prop:nebcorr}(3)), the model route gives the full operator as $[\mathrm{diag}(1,p^{2})]_\chi+\chi(p)p^{\,n}$; it must equal the recursion
$T_p^{2}-\chi(p)p^{\,n+1}$ of Theorem~\ref{thm:nebrelations}(2) (a \PROVED{} identity given that theorem's \STD{} input: a failure stops the package) and is
compared with~(ii)'s $H_{N,p^{2},\chi}$ from $\mathrm{Mer}_{p^{2}}$ on at most six cells, which turns~(ii)'s test from ``against the recursion'' into ``against a
second route through this chapter's own double coset''; (ii)'s operator stays \emph{tested, not proved}. (vi) The slot-17 WP-6 audit's two minors: a second
route at one ladder cell with $D\ge10^{5}$ (the $\F_q$ reduction of a $W_K$ run, or \ref{NB6}'s cocycle route at a reduced prime), and one synthetic
\code{\_\_int128}-overflow case outside the ladder. Arms of \code{HeckeNebentypusValidation} only (no new name); its default run stays within its tier
target; the slot-18 golden event's $d=2$ cell (an $\F_q$-route point) must read the same after~(iv).

\paragraph{Formalization note.}
Lemmas~\ref{lem:nebdata}, \ref{lem:teichmuller}, \ref{lem:inducedchar}, \ref{lem:gammaonetable} and Propositions~\ref{prop:nebdecomp},
\ref{prop:neborbit} are elementary group theory over a commutative ring; Proposition~\ref{prop:nebcorr}(2) inherits
Proposition~\ref{prop:levelcorrespondence}'s dependence on Brown's double-coset formalism, and (4) the $\Ga_1$/$\Ga_0$ comparison of
Hecke rings; Theorem~\ref{thm:nebrelations} needs the Hecke ring of $(\Ga_0(N),\De_0(N))$. (NB1), (NB4)--(NB6) need modular forms.

% v1.22.0-pre.1 (slot 22, P10-C; refereed by F-U5, docs/implementation/slot22/P10C_F4_referee_report.md, ACCEPT WITH REPAIRS): the text of
% docs/implementation/slot22/P10C_contract_text.tex.txt v0.2, part (A.2), verbatim ("pre.N" numbered 1), placed as the last subsection of
% sec:nebentypus so that no statement number of v1.21.0 moves.
\subsection{Wild characters, the groups \texorpdfstring{$\Ga_{\mathsf H}(N)$}{Gamma\_H(N)}, and the tower at \texorpdfstring{$\ell$}{l}}\label{sec:nebwildsec}
% NEW at v1.22.0-pre.1 (hecke-extensions slot 22, P10-C; refereed by F-U5 before any code of P10, rules 20--22).

\paragraph{What this subsection fixes.} (i) The \textbf{ramified coefficient rings} $\mathcal O_{j,e}=\Zl[\zeta_{\ell^{j}}]/(\pi^{e})$ and the
\textbf{wild characters} $\chi_{i,t}$ of $(\Z/\ell^{r})^{\times}$ with values in them (Definition~\ref{def:nebwild}, Lemma~\ref{lem:nebwild}) --- lifting
refusal (ii) for these rings only. (ii) The \textbf{intermediate groups} $\Ga_0(N)\supseteq\Ga_{\mathsf H}(N)\supseteq\Ga_1(N)$, $\mathsf H\le G_N$, with
a character on $\mathsf H$, their models, diamonds and operators (Definition~\ref{def:nebgammaH}, Propositions~\ref{prop:nebgammaH},
\ref{prop:nebgammaHheilbronn}) --- the level structure of slot-22 WAKE-1 (\file{Hecke/HeckeGammaH}) and of the $\ell$-adic tower of [U] part~5, where
$\mathsf H$ is the Teichm\"uller subgroup of $(\Z/\ell^{r})^{\times}$. (iii) \textbf{Specialisation} from $\Ga_{\mathsf H}(N)$ to $\Ga_0(N)$ with a character, and
\textbf{coinvariants} along the tower, without dividing by $\abs{G_N/\mathsf H}$ (Lemma~\ref{lem:nebshapiro}) --- the replacement of
Proposition~\ref{prop:nebdecomp}(3), which needs $\varphi(N)\in\cring^{\times}$ and fails at $\ell\mid\varphi(N)$. (iv) The \textbf{torsion} of the wild
character models (Lemma~\ref{lem:nebwildtors}) and the \textbf{elimination over a chain ring} (Lemma~\ref{lem:nebchainrref}). (v) The predictions
(NB7)--(NB10) (Remark~\ref{rmk:nebwildcheckable}) and the shared \code{CosetLevel} interface (Remark~\ref{rmk:nebcosetlevel}).
\emph{Notation}: as in \S\ref{sec:nebentypus}; $\ell\ge5$ throughout this subsection; $u:=1+\ell$; for $d\in\Z$ prime to $\ell$, $\omega(d)\in\mu_{\ell-1}(\Zl)$
is the Teichm\"uller representative of $d\bmod\ell$ and $\langle d\rangle_{\ell}:=d\,\omega(d)^{-1}\in1+\ell\Zl$ (the subscript keeps it apart from the diamond
$\langle d\rangle$); the sans-serif $\mathsf H$ is a subgroup of $G_N$, \emph{not} the $H\le\Ga$ of Lemma~\ref{lem:inducedchar}.

\begin{definition}[Wild coefficient rings and wild characters]\label{def:nebwild}
Let $j\ge1$, $\varphi_j:=\varphi(\ell^{j})=\ell^{j-1}(\ell-1)$, $\mathcal O_j:=\Zl[z]/(\Phi_{\ell^{j}}(z))$, $\zeta$ the image of $z$, $\pi:=1-\zeta$, and for
$e\ge1$, $\mathcal O_{j,e}:=\mathcal O_j/\pi^{e}\mathcal O_j$. \textbf{As data}: for $e=K\varphi_j$, $\mathcal O_{j,e}=(\Z/\ell^{K})[z]/(\Phi_{\ell^{j}}(z))$ on the power
basis $1,z,\dots,z^{\varphi_j-1}$ with coefficients in $[0,\ell^{K})$ (the representation of \Espec{rmk:lvwild} and of \code{nf::CycloModRing}); for other $e$, the
quotient of $\mathcal O_{j,\lceil e/\varphi_j\rceil\varphi_j}$ by $\pi^{e}$, an element stored as its reduction against the Howell basis of $\pi^{e}\mathcal O$ (Lemma~\ref{lem:nebwild}(2)).
For $r\ge1$, $N=\ell^{r}$, $i\in\Z/(\ell-1)$ and $t\in\Z$ with $\ell^{j}\mid t\,\ell^{r-1}$, put $L(d)\in\Z/\ell^{r-1}$ by $\langle d\rangle_\ell\equiv u^{L(d)}\pmod{\ell^{r}}$ and
\[
  \chi_{i,t}(d):=\omega(d)^{i}\,\zeta^{\,t\,L(d)}\in\mathcal O_j^{\times}\qquad(d\in G_{\ell^{r}}),\qquad \psi_t:=\chi_{0,t}.
\]
$\chi_{i,t}$ is the \textbf{wild character} of tame part $\omega^{i}$ and wild part $\psi_t$; a character of $G_{\ell^{r}}$ into $\mathcal O_{j,e}^{\times}$ is used only in this
form (P4's $\chi_{i,t}$ of \Espec{rmk:lvwild}, with its $m=r$).
\end{definition}

\begin{lemma}[The wild rings]\label{lem:nebwild}
\PROVED{} (1) $\mathcal O_j$ is a complete discrete valuation ring, totally ramified of degree $\varphi_j$ over $\Zl$, with uniformiser $\pi$, $\ell\mathcal O_j=\pi^{\varphi_j}\mathcal O_j$
and residue field $\F_\ell$ \STD{\cite[Lemma~1.4, Prop.~2.1]{washington}, cyclotomic fields of prime-power conductor; \VERIFYloc}. Hence $\mathcal O_{j,e}$ is a
\emph{chain ring}: its ideals are $\pi^{s}\mathcal O_{j,e}$, $0\le s\le e$; it is local with residue field $\F_\ell$, and $\mathcal O_{j,K\varphi_j}=\mathcal O_j/\ell^{K}$ is free of rank
$\varphi_j$ over $\Z/\ell^{K}$; $\mathcal O_{j,e}\cong\F_\ell[\pi]/(\pi^{e})$ for $e\le\varphi_j$.
(2) Every matrix over $\mathcal O_{j,e}$ is equivalent to $\mathrm{diag}(\pi^{s_1},\dots,\pi^{s_k},0,\dots)$ (Smith), and every submodule of $\mathcal O_{j,e}^{D}$ has a unique
Howell form (reduced row echelon, pivots $\pi^{s}$, closed under multiplication by $\pi^{e-s}$); a finitely generated module is $\bigoplus\mathcal O/\pi^{e_i}$.
(3) $\chi_{i,t}$ is a homomorphism $G_{\ell^{r}}\to\mathcal O_j^{\times}$, $\chi_{i,t}\equiv\omega^{i}\pmod\pi$ and $\psi_t\equiv1\pmod\pi$. For $m\in\Z$ with $\ell^{j}\nmid m$,
$v_\pi(\zeta^{m}-1)=\ell^{v_\ell(m)}$; so $c(\psi_t):=\min_dv_\pi(\psi_t(d)-1)=\ell^{v_\ell(t)}$ if $\psi_t\ne1$. $\psi_t$ has order $\ell^{\,j-v_\ell(t)}$ and, if non-trivial,
conductor $\ell^{\,j-v_\ell(t)+1}$.
(4) $\sigma_s\colon\zeta\mapsto\zeta^{s}$ ($\ell\nmid s$) is a ring automorphism of every $\mathcal O_{j,e}$ carrying the $\chi_{i,t}$-model (Definition~\ref{def:nebmodel}) entrywise
onto the $\chi_{i,st}$-model; one $t$ per orbit suffices for lengths and characteristic polynomials, as in Remark~\ref{rmk:nebfield}(3).
\end{lemma}

\begin{proof}
(1) The ring statements are the cited ones; a quotient of a discrete valuation ring by a power of the maximal ideal is a chain ring. (2) A chain ring is a local
principal ideal ring: elimination with a pivot of least valuation in the remaining block (row and column operations) gives Smith, rows only (and the closure
condition) Howell; uniqueness of the reduced form as for $\Z/\ell^{K}$ (\HA{rmk:nebcost}(4) for $W_K(\F_q)$, the same argument). (3) $d\mapsto\omega(d)$ and $d\mapsto
L(d)$ are homomorphisms ($\langle\cdot\rangle_\ell$ is), and $\zeta^{tL(d)}$ is well defined because $\zeta^{t\ell^{r-1}}=1$. $\zeta^{m}$ is a primitive $\ell^{\,j-v(m)}$-th root
of unity, and $1-\zeta'$ for a primitive $\ell^{s}$-th root $\zeta'$ has $v_\ell=1/\varphi(\ell^{s})$, i.e.\ $v_\pi=\varphi_j/\varphi(\ell^{s})=\ell^{j-s}$. The order of $\psi_t$
is the order of $\zeta^{t}$ ($L$ is onto $\Z/\ell^{r-1}$); a character of $(\Z/\ell^{r})^{\times}$ of exact order $\ell^{s}\ge\ell$, trivial on $\mu_{\ell-1}$, has conductor
$\ell^{s+1}$. (4) $\sigma_s$ preserves $\pi\mathcal O$, and the model's matrices are polynomials in the character values.
\end{proof}

\begin{definition}[The groups $\Ga_{\mathsf H}(N)$ and their models]\label{def:nebgammaH}
Let $N\ge1$, $\mathsf H\le G_N$ a subgroup and $\chi\colon G_N\to\cring^{\times}$ a character (only $\chi|_{\mathsf H}$ enters the module; the extension to $G_N$ fixes the
diamonds, item (4) below). Put
\[
  \Ga_{\mathsf H}(N):=\varsigma_N^{-1}(\mathsf H)=\{\gamma\in\Ga_0(N):d_\gamma\bmod N\in\mathsf H\},
\]
so $\Ga_{\{1\}}(N)=\Ga_1(N)$, $\Ga_{G_N}(N)=\Ga_0(N)$. Let $\mathsf U\subseteq G_N$ be the set of the least representatives $u'\in[1,N)$ of the cosets $G_N/\mathsf H$
(so $1\in\mathsf U$), ordered by $u'$, and $w_{(j,u)}:=\sigma_u\,w_j$ ($w_j$ of Definition~\ref{def:cosetrepsN}, $\sigma_u$ of Lemma~\ref{lem:gammaonetable}), ordered
$j$-major, $u$-minor. The \textbf{$\Ga_{\mathsf H}$-model with character} is $L^{\chi}_{\Ga_{\mathsf H}(N)}:=L^{\rho_{n,\chi}}_{\Ga_{\mathsf H}(N)}$ of Lemma~\ref{lem:inducedchar} on this table,
$\rho_{n,\chi}(\gamma)=\chi(a_\gamma)\rhon(\gamma)$ (Definition~\ref{def:nebdata}(5)), and for $\sigma\in\Ga_0(N)$, $R_\sigma f(x):=\rho_{n,\chi}(\sigma)f(\sigma^{-1}x)$ on
$\Ind_{\Ga_{\mathsf H}(N)}^{\Ga}\VR(\chi)$; the \textbf{diamond} of $d\in G_N$ is $\langle d\rangle:=R_{\sigma_d}$.
\end{definition}

\begin{proposition}[$\Ga_{\mathsf H}(N)$: table, model, diamonds, operators]\label{prop:nebgammaH}
\PROVED{} In the situation of Definition~\ref{def:nebgammaH}:
\begin{enumerate}
\item $\Ga_{\mathsf H}(N)\trianglelefteq\Ga_0(N)$, $\Ga_0(N)/\Ga_{\mathsf H}(N)\cong G_N/\mathsf H$ through $\varsigma_N$; $-I\in\Ga_{\mathsf H}(N)$ iff $-1\in\mathsf H$.
\item The $w_{(j,u)}$ represent $\Ga_{\mathsf H}(N)\backslash\Ga$; there are $m_N\,[G_N:\mathsf H]$ of them; $w_{(1,1)}=I$. The bottom row of $w_{(j,u)}$ is $\equiv u'(c_j,d_j)\pmod N$, and
  the class of $\gamma\in\Ga$ is the $\mathsf H$-orbit of its bottom row modulo $N$. For $\gamma\in\Ga$: $(j,u)\cdot\gamma=(j\cdot\gamma,[u\,\varsigma_N(h_j(\gamma))])$, $[x]$ the
  element of $\mathsf U$ in $x\mathsf H$, and $h_{(j,u)}(\gamma)=\sigma_u\,h_j(\gamma)\,\sigma_{[u\varsigma_N(h_j(\gamma))]}^{-1}\in\Ga_{\mathsf H}(N)$. At $\mathsf H=\{1\}$ this is
  Lemma~\ref{lem:gammaonetable}; at $\mathsf H=G_N$ it is Definition~\ref{def:cosetrepsN}'s table \emph{entry for entry}.
\item The relation matrices are monomial, $M^{\chi}_{\mathsf H}(\gamma)=\sum_{J}E_{J,J\cdot\gamma}\otimes\chi(a_{h_J(\gamma)})\rhon(h_J(\gamma))$, and for $6\in\cring^{\times}$ the model computes
  $H^{\bullet}(\Ga_{\mathsf H}(N);\VR(\chi))$; Proposition~\ref{prop:neborbit}(1),(3) hold for it ((3) in its general form, with $M^{\langle-I\rangle}$ when $-1\notin\mathsf H$).
\item $R_\sigma$ is an automorphism of $\cring[\Ga]$-modules depending only on $\sigma\Ga_{\mathsf H}(N)$, $R_\sigma R_{\sigma'}=R_{\sigma\sigma'}$; so $G_N/\mathsf H$ acts by fibre maps
  $\langle d\rangle$, and $\langle h\rangle=1$ for $h\in\mathsf H$. For a character $\psi\colon G_N/\mathsf H\to\cring^{\times}$ the inclusion of functions identifies
  $\Ind_{\Ga_0(N)}^{\Ga}\VR(\chi\psi)$ with $\{f:\langle d\rangle f=\psi(d)f\ \forall d\}$ (no division).
\item The operators of Proposition~\ref{prop:nebcorr}(3) with $A=B=\Ga_{\mathsf H}(N)$ and the transversals $r_t$ chosen in $\Ga_{\mathsf H}(N)$ ---
  $T_p$ ($p\nmid N$), $U_p$ ($p\mid N$), and $[\Ga_{\mathsf H}\,p\sigma_p\,\Ga_{\mathsf H}]=p^{\,n}\chi(p)\langle p\rangle$ --- commute with every $\langle d\rangle$ and restrict, under (4), to the
  operators of the character model of $\chi\psi$ (Proposition~\ref{prop:nebcorr}(4)'s statement with $\Ga_{\mathsf H}(N)$ for $\Ga_1(N)$).
\item (\ref{NB10}'s oracle) If $\abs{\mathsf H}\in\cring^{\times}$ and $\cring$ holds the values of the characters of $\mathsf H$ with unit differences, then
  $e_{\chi}:=\abs{\mathsf H}^{-1}\sum_{h\in\mathsf H}\chi(h)^{-1}\langle h\rangle_{\Ga_1}$ on the $\Ga_1(N)$ model is an idempotent with image the $\Ga_{\mathsf H}$-model of $\chi$, compatibly with
  every operator of (5).
\end{enumerate}
\end{proposition}

\begin{proof}
(1) $\Ga_{\mathsf H}(N)$ is the preimage of a subgroup under the homomorphism $\varsigma_N$ (Lemma~\ref{lem:nebdata}(1)); $\varsigma_N(-I)=-1$. (2) $\Ga=\bigsqcup_j\Ga_0(N)w_j$ and
$\Ga_0(N)=\bigsqcup_{u\in\mathsf U}\Ga_{\mathsf H}(N)\sigma_u$ ($\varsigma_N(\sigma_u)=u$); the bottom row of $\sigma_u$ is $(N,u')$, so that of $\sigma_uw_j$ is $\equiv u'(c_j,d_j)$; two
elements of $\Ga_0(N)w_j$ lie in one $\Ga_{\mathsf H}$-coset iff their $\varsigma_N$-values agree modulo $\mathsf H$, i.e.\ iff their bottom rows agree up to $\mathsf H$. The action and
the cocycle are the computation of Lemma~\ref{lem:gammaonetable}(3) with $\varsigma_N(\cdot)\in\mathsf H$ in place of $=1$. At $\mathsf H=G_N$, $\mathsf U=\{1\}$, $\sigma_1=I$.
(3) Lemma~\ref{lem:inducedchar} with $H=\Ga_{\mathsf H}(N)$, $\eta=\rho_{n,\chi}|_{\Ga_{\mathsf H}}$; Proposition~\ref{prop:neborbit}(1),(3) use only the cocycle identity and the block
formula. (4) The proof of Proposition~\ref{prop:nebdecomp}(1),(2), with $\rho_{n,\chi}$ (multiplicative on $\Ga_0(N)$) for $\rhon$ and $\Ga_{\mathsf H}(N)$ (normal) for $\Ga_1(N)$:
for $h\in\Ga_{\mathsf H}$, $R_hf(x)=\rho_{n,\chi}(h)f(h^{-1}x)=f(x)$; and for $F\in\Ind_{\Ga_0}\VR(\chi\psi)$, $R_\sigma F=\chi(a_\sigma)\chi\psi(a_\sigma)^{-1}F=\psi(d_\sigma)F$.
(5) The proof of Proposition~\ref{prop:nebcorr}(3),(4) with $\Ga_{\mathsf H}(N)$ for $\Ga_1(N)$: the transversal exists because $\varsigma_N(A''_0)=G_N$ there, and
$a_g\equiv1\pmod N$ for $g\in\{\smat100p,p\sigma_p\}$; $\rho_{n,\chi}(p\sigma_p)=\chi(1)p^{\,n}\rhon(\sigma_p)=p^{\,n}\chi(p)\rho_{n,\chi}(\sigma_p)$ as $a_{\sigma_p}\equiv p^{-1}$.
(6) Proposition~\ref{prop:nebdecomp}(3)'s argument for the finite abelian group $\mathsf H$.
\end{proof}

\begin{proposition}[The Heilbronn action on $\Ga_{\mathsf H}(N)$]\label{prop:nebgammaHheilbronn}
\PROVED{} Proposition~\ref{prop:heilbronn} and Remark~\ref{rmk:nebheilbronn} hold for $A=\Ga_{\mathsf H}(N)$, $-1\in\mathsf H$, $n$ even with $\chi(-1)=1$, on the table of
Proposition~\ref{prop:nebgammaH}(2): with $J\cdot\epsilon$ the class of the $\mathsf H$-orbit of $\mathrm{row}_2(w_J)\epsilon$ and $h_J(\epsilon):=w_J\,\epsilon\,w_{J\cdot\epsilon}^{-1}$,
\[
  H^{\mathsf H}_{N,p,\chi}:=\sum_{\epsilon\in\mathrm{Mer}_p}\ \sum_{J\ \text{admissible}}E_{J\cdot\epsilon,\,J}\otimes\chi\bigl(d_{h_J(\epsilon)}\bigr)\,\rhon\bigl(\adjg h_J(\epsilon)\bigr)
\]
preserves $\mathrm{Rel}$ and induces $T_p$ ($p\nmid N$), $U_p$ ($p\mid N$) of Proposition~\ref{prop:nebgammaH}(5) on $H^{1}(\Ga_{\mathsf H}(N);\VR(\chi))$, $6\in\cring^{\times}$.
\end{proposition}

\begin{proof}
The proof of Proposition~\ref{prop:heilbronn} with two sentences changed; the rest (Lemma~\ref{lem:levelduality} for $A$ of finite index with $-I\in A$, Merel's decomposition,
the left-coset bookkeeping) does not use $A=\Ga_0(N)$. \emph{Step~1}: for $p\nmid N$, every $\beta\in AgA$ is upper triangular modulo $N$ with $a_\beta d_\beta\equiv p$ and
$d_\beta\in p\mathsf H$ ($d$ is multiplicative modulo $N$ on $\Ga_0(N)$-type matrices, $d_g=p$); if $\Ga\beta=\Ga\beta'$, then $\gamma=\beta'\adjg\beta/p\in\Ga$ has $\gamma_{21}\equiv0$ and
$d_\gamma\equiv d_{\beta'}a_\beta/p\equiv d_{\beta'}d_\beta^{-1}\in\mathsf H$, so $A\beta=A\beta'$; the count $\abs{A\backslash AgA}=p+1$ (needed for the exhaustion) holds because
$\Ga_1(N)\subseteq A$ and $\Ga_1(N)\to\SL_2(\F_p)$ is onto for $p\nmid N$ (Chinese remainder theorem), so $(0{:}1)\in\mathbb P^{1}(\F_p)$ is attained inside $A$ (Proposition~\ref{prop:heilbronn}'s element
$\smat pu Nv$ with $v$ arbitrary is \emph{not} available in $A$; v0.6.0-pre.2, F-U5 m6); for $p\mid N$ the representatives $gT^{b}$ are unchanged and the image is $\{(1{:}y)\}$ as before. \emph{Step~3}: for
$p\nmid N$, $\mathrm{row}_2(\gamma_\epsilon)\equiv(d_\beta/p)\,\mathrm{row}_2(w_J)\epsilon$ with $d_\beta/p\in\mathsf H$, so the class is $J\cdot\epsilon$; for $p\mid N$ it is exact. The character
scalar is Remark~\ref{rmk:nebheilbronn}'s, $a_{\adjg h}=d_h$ a unit modulo $N$ at every admissible pair; moreover $\mathrm{row}_2(w_J)\epsilon\equiv\eta\,\mathrm{row}_2(w_{J\cdot\epsilon})\pmod N$ with
$\eta\in\mathsf H$ by the definition of $J\cdot\epsilon$, so $d_{h_J(\epsilon)}\equiv\eta\in\mathsf H$: \emph{only $\chi|_{\mathsf H}$ enters the operator}, as it must (the extension of $\chi$ to $G_N$
normalises the diamonds, not the Hecke operators; F-U5 m7).
\end{proof}

\begin{lemma}[Specialisation and coinvariants without division]\label{lem:nebshapiro}
\PROVED{} Let $\Ga''\trianglelefteq\Ga'\le\Ga$ be of finite index, $C:=\Ga'/\Ga''$ abelian, $\cring$ commutative with $6\in\cring^{\times}$, and $X$ an $\cring[\Ga']$-module free of
finite rank over $\cring$. Let $C$ act on $H^{1}(\Ga'';X)$ by conjugation (the diamonds of Proposition~\ref{prop:nebgammaH}(4) when $\Ga''=\Ga_{\mathsf H}(N)$, $\Ga'=\Ga_0(N)$). Then for
every finitely presented $\cring[C]$-module $P$, $\Ga'$ acting on $P$ through $\Ga'\to C$, $\gamma\mapsto(\gamma\Ga'')^{-1}$ --- \textbf{the inverse of the quotient map} (v0.6.0-pre.2,
F-U5 M1: under Shapiro's isomorphism $\Ind_{\Ga''}^{\Ga'}X\cong\cring[C]\otimes_{\cring}X$, $f\mapsto\sum_{u\in C}u\otimes u^{-1}f(u)$, the conjugation $R_c$ becomes multiplication by $c$ on the
group-ring factor while $\Ga'$ acts there by $\bar\gamma^{-1}\otimes\rho(\gamma)$; for $\Ga''=\Ga_{\mathsf H}(N)$ the map is $\gamma\mapsto a_\gamma\bmod\mathsf H$),
\[
  H^{1}(\Ga'';X)\otimes_{\cring[C]}P\;\cong\;H^{1}(\Ga';X\otimes_{\cring}P),
\]
naturally in $P$; in particular (a) $P=\cring_\psi$ for a character $\psi$ of $C$: $H^{1}(\Ga'';X)\otimes_{\cring[C]}\cring_\psi\cong H^{1}(\Ga';X(\psi))$, where $X(\psi)$ is $X$ with $\gamma$ acting
by $\psi(a_\gamma)\,\rho(\gamma)$ --- for $X=\VR(\chi)$ this is $\VR(\chi\psi)$ of Definition~\ref{def:nebgammaH}, i.e.\ \emph{the $\psi$-coinvariants of the diamonds ($\langle d\rangle\sim\psi(d)$) are the
$\chi\psi$-model}, in agreement with the eigenspace statement of Proposition~\ref{prop:nebgammaH}(4) (checked on the model: $q(f):=\sum_{u}\psi(u)\rho_{n,\chi}(\sigma_u)^{-1}f(\sigma_u)$ has
$q(R_\sigma f)=\psi(d_\sigma)q(f)$ and $q(\gamma f)=\psi(d_\gamma)^{-1}\chi(a_\gamma)\rhon(\gamma)q(f)=\rho_{n,\chi\psi}(\gamma)q(f)$); (b) $P=\cring[C/C']$ for a subgroup
$C'=\Ga'''/\Ga''$: the $C'$-coinvariants of $H^{1}(\Ga'';X)$ are $H^{1}(\Ga''';X)$. The isomorphisms commute with every operator $[\Ga'g\Ga']$ / $[\Ga''g\Ga'']$ of
Proposition~\ref{prop:nebgammaH}(5) whose transversal is chosen in $\Ga''$, and with the ordinary projector $e$ of $U_p$ when both sides carry it.
\end{lemma}

\begin{proof}
Shapiro: $H^{1}(\Ga'';X)\cong H^{1}(\Ga';\cring[C]\otimes_{\cring}X)$, $C$ acting on $H^{1}(\Ga'';X)$ by conjugation corresponding to the right regular action on $\cring[C]$
(Lemma~\ref{lem:inducedchar}, Proposition~\ref{prop:nebdecomp}(1)'s computation). $H^{2}(\Ga';-)=0$ on $\cring[\Ga']$-modules ($6\in\cring^{\times}$, Lemma~\ref{lem:treeres}: a
projective resolution of length one), so $H^{1}(\Ga';-)$ is right exact and additive; choosing a presentation $\cring[C]^{a}\to\cring[C]^{b}\to P\to0$ and tensoring
$\cring[C]\otimes X$ over $\cring[C]$ with it gives the isomorphism (the five lemma on the two right-exact rows). The operators: on both sides they are given by the
same $r_t\in\Ga''$ and the same level-one $\tau$ (Proposition~\ref{prop:nebcorr}(4)'s choice), so Shapiro's isomorphism intertwines them term by term; $e$ is a limit of powers
of $U_p$ in the finite (or profinite) commutative algebra they generate, and $e(M\otimes P)=eM\otimes P$.
\end{proof}

\begin{lemma}[The wild character models: torsion and lengths]\label{lem:nebwildtors}
\PROVED{} Let $N=\ell^{r}$, $n=0$, $\chi=\chi_{i,t}$ with values in $\mathcal O_j$, $\chi$ even ($i$ even), and $EC(\chi)$ the right-hand side of
Proposition~\ref{prop:neborbit}(4) for $\chi$ (over $\mathrm{Frac}\,\mathcal O_j$ with $\zeta_{12}$ adjoined; it depends on $\chi$ only through $\omega^{i}$ on $\mu_4,\mu_6\subset\mu_{\ell-1}$). Put
$c:=0$ if $i\not\equiv0\pmod{\ell-1}$, $c:=c(\psi_t)$ of Lemma~\ref{lem:nebwild}(3) if $i\equiv0$ and $\psi_t\ne1$.
\begin{enumerate}
\item If $\chi\ne1$: $H^{1}(\Ga_0(\ell^{r});\mathcal O_j(\chi))\cong\mathcal O_j^{EC(\chi)}\oplus\mathcal O_j/\pi^{c}$ and, for every $e$,
  $H^{1}(\Ga_0(\ell^{r});\mathcal O_{j,e}(\chi))\cong\mathcal O_{j,e}^{EC(\chi)}\oplus\mathcal O/\pi^{\min(c,e)}$, $H^{0}\cong\mathcal O/\pi^{\min(c,e)}$.
\item The ordinary projector $e$ of $U_\ell$ kills the torsion summand: $eH^{1}(\Ga_0(\ell^{r});\mathcal O_{j,e}(\chi))$ is free over $\mathcal O_{j,e}$ and equals
  $eH^{1}(\Ga_0(\ell^{r});\mathcal O_j(\chi))/\pi^{e}$.
\end{enumerate}
\end{lemma}

\begin{proof}
(1) $H^{2}=0$ ($6$ invertible), so the sequences of $0\to\mathcal O\xrightarrow{\pi^{e}}\mathcal O\to\mathcal O/\pi^{e}\to0$ and $0\to\mathcal O\to\mathcal K\to\mathcal K/\mathcal O\to0$ ($\mathcal K$ the
fraction field) give $H^{1}(\mathcal O/\pi^{e})=H^{1}(\mathcal O)/\pi^{e}$ and $H^{1}(\mathcal O)_{\mathrm{tors}}\cong H^{0}(\mathcal K/\mathcal O(\chi))/\mathrm{im}\,H^{0}(\mathcal K(\chi))$. For $\chi\ne1$,
$H^{0}(\mathcal K(\chi))=0$ (some $a_\gamma$ has $\chi(a_\gamma)\ne1$), and $H^{0}(\mathcal K/\mathcal O(\chi))=\{x:(\chi(a)-1)x=0\ \forall a\in G_{\ell^{r}}\}=\pi^{-c}\mathcal O/\mathcal O$ with
$c=\min_av_\pi(\chi(a)-1)$, which is $0$ when $\omega^{i}\ne1$ (a unit difference) and $c(\psi_t)$ otherwise. The free rank is the dimension over $\mathcal K$, which is $EC(\chi)$ (Proposition~\ref{prop:neborbit}(4),
$H^{0}=0$). $H^{0}(\mathcal O/\pi^{e}(\chi))$ is the $\pi^{\min(c,e)}$-torsion of $\mathcal O/\pi^{e}$. (2) The torsion is the image of the connecting map from
$H^{0}(\Ga_0(\ell^{r});\mathcal K/\mathcal O(\chi))$, which is Hecke-equivariant (\STD{\cite[Lemma~1.1.1, p.~194]{ash_stevens}}, as in \HA{NB4}(c)); on $H^{0}$, $U_\ell$ is the sum over
$b<\ell$ of $\chi(a_{\delta_b})=\chi(1)$, i.e.\ multiplication by $\ell$, nilpotent on a finite module, so $e$ kills it; then $e$ commutes with $/\pi^{e}$ as in (1).
\end{proof}

\begin{lemma}[Reduced echelon form over a chain ring with free quotient]\label{lem:nebchainrref}
\PROVED{} Let $\cring$ be a chain ring with maximal ideal $(\pi)$ and residue field $k$ ($\Z/\ell^{K}$, $W_K(\F_q)$, $\mathcal O_{j,e}$), and $\mathrm{Rel}\subseteq\cring^{D}$ a submodule
generated by $g_1,\dots,g_s$ with $\cring^{D}/\mathrm{Rel}$ free. Then $\mathrm{Rel}$ is a free direct summand; its image in $k^{D}$ is spanned by the $\bar g_i$ and has the same rank; and
$\mathrm{Rel}$ has a unique reduced row echelon basis \emph{with unit pivots}, whose pivot columns are those of the reduced row echelon form of $\langle\bar g_i\rangle\subseteq k^{D}$. The
free columns index a basis of the cokernel.
\end{lemma}

\begin{proof}
A free quotient splits, so $\mathrm{Rel}$ is projective, hence free over the local ring, and $\mathrm{Tor}_1(\cring^{D}/\mathrm{Rel},k)=0$ makes $\mathrm{Rel}\otimes k\to k^{D}$ injective with
image $\langle\bar g_i\rangle$. Let $P$ be the pivot columns over $k$. The projection $\mathrm{Rel}\to\cring^{P}$ is an isomorphism modulo $\pi$, between free modules of rank $\abs P$,
hence an isomorphism (Nakayama); the preimages of the unit vectors are the unique reduced basis with unit pivots on $P$.
\end{proof}

\noindent\emph{Applicability.} For the cokernel of Proposition~\ref{prop:neborbit}(3) the quotient is $H^{1}$, free over $\Z/\ell^{K}$ for the $\Ga_{\mathsf H}(\ell^{r})$-models of [U] part~5 (their
$H^{1}(\Zl)$ is torsion-free: Lemma~\ref{lem:nebwildtors}'s argument with $\chi$ tame, where $c=0$, or $\chi=1$, where $H^{0}(\Ql)\to H^{0}(\Ql/\Zl)$ is onto). So the existing sparse
elimination over $\F_\ell$ (\file{Hecke/HeckeLevelNSparse}, \file{Hecke/HeckeNebentypusSparse}) fixes the pivots and the free columns, and the lift to $\Z/\ell^{K}$ is a back-substitution with
unit pivots. Over $\mathcal O_{j,e}$ at $i\equiv0$ the cokernel has the torsion of Lemma~\ref{lem:nebwildtors}(1) and the lemma does \emph{not} apply: there the elimination is a genuine Howell
form (Lemma~\ref{lem:nebwild}(2)).

\begin{remark}[What is checkable: (NB7)--(NB10)]\label{rmk:nebwildcheckable}
Committed before slot-22 P10's code (rule~20), convention of Remark~\ref{rmk:hidacheckable}. Cells: $\ell\in\{5,7,11,13\}$, $r\le3$, $n=0$; numbers in
\file{docs/implementation/slot22/P10\_spec.md} \S7.
\begin{enumerate}[label=(NB\arabic*),start=7]
\item\label{NB7} \PROVED{} (Lemma~\ref{lem:nebwild}(3), Proposition~\ref{prop:neborbit}(3)) \textbf{Reduction modulo $\pi$.} The relation matrices and every operator of the $\chi_{i,t}$-model over
  $\mathcal O_{j,1}=\F_\ell$ are entrywise those of the $\omega^{i}$-model over $\F_\ell$ (\HA{NB2}'s bit-for-bit standard), so the cokernel and the operators on it agree.
\item\label{NB8} \PROVED{} (Lemma~\ref{lem:nebwildtors}) \textbf{Lengths.} $\mathrm{length}\,H^{1}(\Ga_0(\ell^{r});\mathcal O_{j,e}(\chi_{i,t}))=e\cdot EC+\min(c,e)$, $\mathrm{length}\,H^{0}=\min(c,e)$, and the
  ordinary part is free over $\mathcal O_{j,e}$; its rank is $2d(a)+1$ ($a\equiv i+2$, [U] Theorem~\texttt{thm:u-towercontrol}(5)) --- the rank clause \PROVED{} \emph{given} [U]
  Proposition~\texttt{prop:u-towerorient} (\COMP{} at $r=1$, gated by (U18)(o); F-U5 M3).
\item\label{NB9} \PROVED{} (Lemma~\ref{lem:nebshapiro}(a)) \textbf{Wild specialisation.} $eH^{1}(\Ga_{\mathsf H_r}(\ell^{r});\Z/\ell^{K}(\omega^{i}))\otimes_{(\Z/\ell^{K})[G_r]}\mathcal O_{j,e}$, $\langle u\rangle\mapsto\zeta^{t}$,
  is $eH^{1}(\Ga_0(\ell^{r});\mathcal O_{j,e}(\chi_{i,t}))$ ($e\le K\varphi_j$, $\mathsf H_r$ the Teichm\"uller subgroup, $G_r=G_{\ell^{r}}/\mathsf H_r$), with equal $T_p$ and $U_\ell$; computed by the two
  independent routes (the $\Ga_0$ character model over $\mathcal O_{j,e}$ with a Howell elimination; the $\Ga_{\mathsf H}$ model over $\Z/\ell^{K}$ tensored down).
\item\label{NB10} \PROVED{} (Proposition~\ref{prop:nebgammaH}(2),(6)) \textbf{The $\Ga_{\mathsf H}$ model against its two neighbours.} At $\mathsf H=G_N$ it is the character model entry for entry; at
  $\mathsf H=\mathsf H_r$, $r\le2$, its $H^{1}$, $T_p$, $U_\ell$ and diamonds equal those of $e_{\omega^{i}}$ applied to the $\Ga_1(\ell^{r})$ model.
\end{enumerate}
\end{remark}

\begin{remark}[One level interface: $\Ga_{\mathsf H}(N)$ for WAKE-1 and for the tower]\label{rmk:nebcosetlevel}
(Not a statement; the code contract.) \code{hecke::level::CosetLevel} (\file{Hecke/HeckeLevelEll.h}: \code{N}, \code{m}, \code{rep}, \code{class\_of}, \code{contains},
\code{contains\_wide}, \code{name}) is the one interface. Slot-22 WAKE-1 adds \code{GammaHLevelN} (\file{Hecke/HeckeGammaH}; $\mathsf H$ by generators) with the table of
Proposition~\ref{prop:nebgammaH}(2), its sub-transversal (the analogue of \code{Gamma1SubTransversal}) and its diamonds; P10 \emph{reuses} it with $N=\ell^{r}$ and $\mathsf H=\mathsf H_r$
generated by the Teichm\"uller lift of a primitive root modulo $\ell$, and adds only what is level-independent: the tame scalar $\chi(a_\delta)$ (\code{neb::BlockScalar}, already generic in
$\delta$), the chain-ring elimination (Lemma~\ref{lem:nebchainrref}) and Proposition~\ref{prop:nebgammaHheilbronn}'s operator. \code{Gamma1LevelN} is the case $\mathsf H=\{1\}$ and stays the
oracle (\ref{NB10}); \code{Gamma0LevelN} the case $\mathsf H=G_N$. A second $\Ga_{\mathsf H}$ class is not to be written (the slot-22 runbook: ``reuse WAKE-1's \code{CosetLevel} additions,
do not fork them''); if WAKE-1's class lacks the diamonds of a quotient $G_N/\mathsf H$ of order prime to $6$ but divisible by $\ell$, P10 extends that class.
\end{remark}

\paragraph{Code map of \S\ref{sec:nebwildsec}} \textbf{Planned} (slot-22 P10, after F-U5's ACCEPT): \file{Hecke/HeckeChainRing.h} (header: $\Z/\ell^{K}$ and $\mathcal O_{j,e}$ element
arithmetic, $v_\pi$, unit inverse, Howell / Smith), \file{Hecke/HeckeNebentypusWild.\{h,cpp\}} (the $\chi_{i,t}$ character model over $\mathcal O_{j,e}$: Definition~\ref{def:nebwild},
Lemmas~\ref{lem:nebwild}, \ref{lem:nebwildtors}), WAKE-1's \file{Hecke/HeckeGammaH} (Proposition~\ref{prop:nebgammaH}) extended by P10 with
Proposition~\ref{prop:nebgammaHheilbronn} and Lemma~\ref{lem:nebchainrref}; validation \code{HeckeNebentypusWildValidation} (\code{standard}) gating (NB7)--(NB10).

% v1.22.0-pre.3 (slot 22, LES-C; spec docs/implementation/slot22/LES0_spec.md v0.1, owner questions Q1--Q7 taken at the spec's recommendations by the
% lead pending owner objection): NEW last subsection of sec:nebentypus, after sec:nebwildsec, so that no statement number of v1.22.0-pre.2 moves.
% NOT yet refereed (rule 22: a Fable referee F-LES reads it before any LES code); predictions docs/implementation/slot22/LES_predictions.md.
% v1.23.0-pre.1 (slot 23, LES-2; referee F-C123pre1): prop:levelconehecke(4) replaced (independence of the transversal, by (T1)); obl:levelconeexact(ix) criterion-fails control.
% v1.23.0-pre.2: repairs of F-C123pre1 (MINOR-1, MINOR-2, NIT-1..-5); wording only.
\subsection{The level-change cone and inflation--restriction}\label{sec:nebles}
% NEW at v1.22.0-pre.3 (hecke-extensions slot 22, LES-C).

\paragraph{What this subsection fixes.} (i) For a pair $B\le A$ of finite-index subgroups of $\Ga$ with compatible characters, the
\textbf{restriction} and the \textbf{trace} as fibre maps of the induced models (Definition~\ref{def:levelpair}, Lemmas~\ref{lem:nebcompat}, \ref{lem:levelresmorph}).
(ii) The \textbf{mapping cone} of the restriction and its \textbf{six-term exact sequence} on the fixed model, at every ring and every $\ell$
(Definition~\ref{def:levelcone}, Lemma~\ref{lem:levelconeexact}); what it computes (Proposition~\ref{prop:levelconeshapiro}); when the model computes group
cohomology at $\ell\in\{2,3\}$ (Proposition~\ref{prop:modeltorsion}); the split case (Proposition~\ref{prop:levelsplit}). (iii) The \textbf{Hecke operators} on the
sequence (Proposition~\ref{prop:levelconehecke}). (iv) The \textbf{normal case}: the cone sequence is inflation--restriction plus two pieces
(Definition~\ref{def:normalpair}, Proposition~\ref{prop:levelinfres}), made explicit for a cyclic quotient of order $p$ (Proposition~\ref{prop:cycliccone},
Lemmas~\ref{lem:levelfreeeuler}, \ref{lem:unipotentinv}). (v) \textbf{Which pairs are normal}: $\Ga_0(MN)\trianglelefteq\Ga_0(N)$ iff $M\mid\gcd(N,24)$
(Proposition~\ref{prop:gamma0normal}), the steps of the $p$-tower (Proposition~\ref{prop:hidanormalsteps}), level one (Remark~\ref{rmk:normalN1}). (vi) The
obligations (Obligation~\ref{obl:levelconeexact}), the predictions \ref{LES1}--\ref{LES9} (Remark~\ref{rmk:lescheckable}), the deferred old-space cone
(Remark~\ref{rmk:oldspacecone}) and the code map (planned). \emph{Scope} (slot-22 LES-0, owner questions Q2--Q5 at the recommendation): coefficients
$\F_\ell$ (and $\F_q$) and $\Z/\ell^{K}$; at $\ell\in\{2,3\}$ the sequence is the model's (Bredon) sequence, the group-cohomology sequence in all degrees
is deferred to the free model of \S\ref{sec:integral} (stated over $\Z$ at v1.23.0-pre.18: Proposition~\ref{prop:descentmaps}); the old-space $U_p$ cone is deferred; no console verb.
\emph{Notation}: as in \S\ref{sec:nebentypus}; $B\le A\le\Ga$ (the names of Proposition~\ref{prop:levelcorrespondence}), $m_A:=[\Ga:A]$, $m_B:=[\Ga:B]$,
$U:=ST$; $\VR=\Sym^{n}$ over $\cring$ with $\rhon$. $\mathsf G:=A/B$ in the normal case (sans-serif, \emph{not} the $G_N=(\Z/N)^{\times}$ of
\S\ref{sec:nebentypus}); $\mathrm{Nm}_{\mathsf G}$ its norm element; $\alpha_{B/A}$, $\omega_{B/A}$, $Q_{A/B}$, $\mathsf G$, $\Ga^{2}$, $\Ga^{3}$ are new symbols of this
subsection and mean nothing else.

\begin{definition}[Level pair with characters]\label{def:levelpair}
A \emph{character datum} on a finite-index $A\le\Ga$ is a level $N_A\ge1$ with $A\subseteq\Ga_0(N_A)$ and a character $\chi\colon G_{N_A}\to\cring^{\times}$
(Definition~\ref{def:nebdata}(3); $N_A=1$, $\chi=1$ is allowed for every $A$); $\VR(\chi)$ is $\VR$ with $A$ acting by $\rho_{n,\chi}(\gamma)=\chi(a_\gamma)\rhon(\gamma)$,
a homomorphism on $\Ga_0(N_A)$ (Lemma~\ref{lem:nebdata}(2)). A \emph{level pair} is $(A,\chi)\supseteq(B,\chi')$: $B\le A$ of finite index, with character
data $(N_A,\chi)$ on $A$ and $(N_B,\chi')$ on $B$. The levels in use are $\Ga_{\mathsf H}(N)$ (Definition~\ref{def:nebgammaH}; $\Ga_0(N)$, $\Ga_1(N)$ among
them), their intersections $\Ga_{\mathsf H}(N)\cap\Ga_0(N')$, and preimages of subgroups of $\SL_2(\Z/M)$ ($\Ga(M)$, and $\Ga^{2}$, $\Ga^{3}$ of
Remark~\ref{rmk:normalN1}), the last with $\chi=1$. $\Ind_A^{\Ga}\VR(\chi)$ and its model $L^{\chi}_A$ are those of Lemma~\ref{lem:inducedchar}, on fixed
tables $w_1=I,\dots,w_{m_A}$ of $A\backslash\Ga$ and $v_1=I,\dots,v_{m_B}$ of $B\backslash\Ga$; $a(i)$ is the index with $Av_i=Aw_{a(i)}$ and
$h_i:=v_iw_{a(i)}^{-1}\in A$.
\end{definition}

\begin{lemma}[Compatibility]\label{lem:nebcompat}
\PROVED{} In a level pair, $\VR(\chi)|_B=\VR(\chi')$ as $\cring[B]$-modules iff $\chi(a_\gamma)=\chi'(a_\gamma)$ for every $\gamma\in B$. For
$B=\Ga_{\mathsf H'}(N')\subseteq A=\Ga_{\mathsf H}(N)$ ($N\mid N'$) this is $\chi'(u)=\chi(u\bmod N)$ for every $u\in\mathsf H'$; in particular, for
$\Ga_0(MN)\subseteq\Ga_0(N)$ it is $\chi'=\chi\circ\mathrm{red}_{MN\to N}$.
\end{lemma}

\begin{proof}
Both actions are $\gamma\mapsto(\text{scalar})\cdot\rhon(\gamma)$ and $\rhon(\gamma)$ is invertible, so they agree iff the scalars agree. On $\Ga_{\mathsf H'}(N')$
the values $a_\gamma\bmod N'$ are exactly $\mathsf H'$: $a_\gamma\equiv d_\gamma^{-1}$ (Lemma~\ref{lem:nebdata}(1)), $d_\gamma\bmod N'$ ranges over $\mathsf H'$ by the surjectivity
of $\varsigma_{N'}$, and $\mathsf H'$ is a group.
\end{proof}

\noindent A level pair whose characters are \emph{not} compatible has no restriction map of the kind below; the code refuses it by name
(\code{les\_incompatible\_character}). This is the hypothesis of the owner's request (``compatible Nebentypes'').

\begin{definition}[The restriction and the trace on the fibre]\label{def:levelres}
For a compatible level pair, view $\Ind_A^{\Ga}\VR(\chi)=\{f\colon\Ga\to\VR:f(a x)=\rho_{n,\chi}(a)f(x),\ a\in A\}$ (Definition~\ref{def:inducedmodule}) and let
$r_1=I,\dots,r_{[A:B]}$ represent $B\backslash A$. Put
\[
  \alpha_{B/A}f:=f\ \ (\text{the inclusion of functions}),\qquad (\omega_{B/A}g)(x):=\sum_{t}\rho_{n,\chi}(r_t)^{-1}\,g(r_tx),
\]
$\alpha_{B/A}\colon\Ind_A^{\Ga}\VR(\chi)\to\Ind_B^{\Ga}\VR(\chi')$ and $\omega_{B/A}\colon\Ind_B^{\Ga}\VR(\chi')\to\Ind_A^{\Ga}\VR(\chi)$. In the coordinates of
Definition~\ref{def:inducedmodule},
\begin{equation}\label{eq:levelres}
  (\alpha_{B/A}f)_i=\chi(a_{h_i})\,\rhon(h_i)\,f_{a(i)},\qquad
  (\omega_{B/A}g)_a=\sum_{i:\,a(i)=a}\chi(a_{h_i})^{-1}\rhon(h_i)^{-1}g_i .
\end{equation}
\end{definition}

\begin{lemma}[Restriction and trace are fibre maps]\label{lem:levelresmorph}
\PROVED{} For a compatible level pair over any commutative $\cring$, no hypothesis on $\ell$ or $n$:
\begin{enumerate}
\item $\alpha_{B/A}$ and $\omega_{B/A}$ are well defined, independent of the $r_t$, morphisms of $\cring[\Ga]$-modules given by \eqref{eq:levelres}, hence (one monomial
  fibre matrix at every simplex) morphisms of matrix local systems $L^{\chi}_A\rightleftarrows L^{\chi'}_B$ (Theorem~\ref{thm:rep_ls_equiv}); $\alpha_{B/A}$ is
  injective and split over $\cring$, $\omega_{B/A}$ is surjective.
\item $\omega_{B/A}\circ\alpha_{B/A}=[A:B]\cdot\id$.
\item Under the reciprocity $\phi\mapsto\phi(\cdot)(I)$ of Proposition~\ref{prop:levelcorrespondence}(3), $(\alpha_{B/A})_{\ast}$ is $\psi\mapsto\psi$ (an $A$-equivariant
  cochain read as a $B$-equivariant one) and $(\omega_{B/A})_{\ast}$ is the transfer $\psi\mapsto\sum_t\rho_{n,\chi}(r_t)^{-1}\psi(r_t\,\cdot)$; both are cochain maps
  $\Ceq(K,\Rel;L^{\chi}_A)\rightleftarrows\Ceq(K,\Rel;L^{\chi'}_B)$, and for $6\in\cring^{\times}$ they induce $\res^{A}_{B}$ and $\cor^{A}_{B}$ on
  $H^{\bullet}(A;\VR(\chi))$, $H^{\bullet}(B;\VR(\chi'))$ \STD{\cite[Ch.~III.9--10]{brown}, as in Proposition~\ref{prop:omega}; locator read at an earlier R20-L pass
(\cite[Prop.~III.9.5, p.~82]{brown})}.
\item At $\chi=\chi'=1$, \eqref{eq:levelres} is Proposition~\ref{prop:levelcorrespondence}(4)(b)'s $\alpha^{0}_{B/A}$ and trace, entry for entry on the same tables.
\end{enumerate}
\end{lemma}

\begin{proof}
(1) For $b\in B$, $f(bx)=\rho_{n,\chi}(b)f(x)=\rho_{n,\chi'}(b)f(x)$ (Lemma~\ref{lem:nebcompat}), so $\alpha_{B/A}f\in\Ind_B$; right translation commutes with the
inclusion. For $\omega_{B/A}$: replacing $r_t$ by $br_t$ ($b\in B$) changes the term by $\rho(r_t)^{-1}\rho(b)^{-1}\rho(b)$, i.e.\ not at all ($\rho:=\rho_{n,\chi}$, multiplicative on
$A$); for $a\in A$ write $r_ta=b_tr_{t\cdot a}$ ($b_t\in B$, $t\mapsto t\cdot a$ a permutation), then $(\omega g)(ax)=\sum_t\rho(r_t)^{-1}\rho(b_t)g(r_{t\cdot a}x)=
\sum_t\rho(ar_{t\cdot a}^{-1})g(r_{t\cdot a}x)=\rho(a)(\omega g)(x)$; right translation commutes. Coordinates: $(\alpha f)_i=f(v_i)=f(h_iw_{a(i)})=\rho(h_i)f_{a(i)}$; for
$\omega$, the $r_tw_a$ represent the $B$-cosets in $Aw_a$, i.e.\ $r_tw_a=bv_i$ with $a(i)=a$, and then $r_t=bh_i$ and $\rho(r_t)^{-1}g(r_tw_a)=\rho(h_i)^{-1}g(v_i)$.
Each block row of $\alpha$ has one non-zero block, an invertible one, and every $a$ occurs as some $a(i)$, so $\alpha$ is injective; $f\mapsto(f_{a(i_a)})_a$ for one
$i_a$ with $a(i_a)=a$, corrected by $\rho(h_{i_a})^{-1}$, is a retraction over $\cring$. Given $f\in\Ind_A$, the $g\in\Ind_B$ with $g_{i_a}:=\chi(a_{h_{i_a}})\rhon(h_{i_a})f_a$
and $g_i:=0$ otherwise has $\omega g=f$. (2) $(\omega\alpha f)(x)=\sum_t\rho(r_t)^{-1}\rho(r_t)f(x)$. (3) $(\omega\circ\phi)(y)(I)=\sum_t\rho(r_t)^{-1}\phi(y)(r_t)=
\sum_t\rho(r_t)^{-1}\phi(r_ty)(I)$ by the $\Ga$-equivariance of $\phi$; the statement for $\alpha$ is immediate. Cochain maps: both are post-composition with module maps. For
$6\in\cring^{\times}$, $C_{\bullet}(\Xbulk)\otimes\cring$ is a projective resolution of $\cring$ over $\cring[A]$ and over $\cring[B]$ (Remark~\ref{rmk:orbifold},
Lemma~\ref{lem:treeres}), on which restriction of equivariance and the displayed transfer compute $\res$ and $\cor$ (the cited section). (4) Compare with
Proposition~\ref{prop:levelcorrespondence}(4)(b): $\rhon(v_iw_{a(i)}^{-1})=\rhon(h_i)$, and $\rhon(h_i)^{-1}=\rhon(w_av_i^{-1})$ is its trace block.
\end{proof}

\begin{definition}[The level cone and relative cohomology]\label{def:levelcone}
For a compatible level pair, $\Cone{\alpha_{B/A}}^{q}:=\Cq{q+1}(K,\Rel;L^{\chi}_A)\oplus\Cq{q}(K,\Rel;L^{\chi'}_B)$ with
$\Dcone(c,c'):=(-dc,\ dc'+(\alpha_{B/A})_{\ast}c)$ --- \eqref{eq:conediff} with $\rres$ replaced by $(\alpha_{B/A})_{\ast}$ --- and the \emph{relative cohomology of
the pair} $H^{q}(A,B):=H^{q-1}(\Cone{\alpha_{B/A}})$ (the indexing of $\Hcpt{q}$ in Definition~\ref{def:pair}). $H^{q}(A)$, $H^{q}(B)$ abbreviate
$\Heq{q}(K,\Rel;L^{\chi}_A)$, $\Heq{q}(K,\Rel;L^{\chi'}_B)$, and $\res^{q}:=(\alpha_{B/A})_{\ast}$ on $H^{q}$, $\cor^{q}:=(\omega_{B/A})_{\ast}$.
\end{definition}

\begin{lemma}[The six-term sequence on the model]\label{lem:levelconeexact}
\PROVED{} For a compatible level pair, over every commutative $\cring$, at every $\ell$ and every $n$:
\begin{enumerate}
\item $\Dcone^{2}=0$, and $0\to\Ceq(L^{\chi'}_B)\to\Cone{\alpha_{B/A}}\to\Ceq(L^{\chi}_A)[1]\to0$ is a short exact sequence of complexes whose connecting map is
  $+\res^{q+1}$.
\item $H^{q}(A)=H^{q}(B)=0$ for $q\ge2$, and $\res^{0}$ is injective.
\item Hence $H^{q}(A,B)=0$ for $q\notin\{1,2\}$ and the long exact sequence is the \textbf{six-term exact sequence}
  \begin{equation}\label{eq:levelsixterm}
    0\to H^{0}(A)\xrightarrow{\res^{0}}H^{0}(B)\to H^{1}(A,B)\to H^{1}(A)\xrightarrow{\res^{1}}H^{1}(B)\to H^{2}(A,B)\to0 ;
  \end{equation}
  over a field $\dim H^{1}(A,B)=\dim H^{0}(B)-\dim H^{0}(A)+\dim\ker\res^{1}$ and $\dim H^{2}(A,B)=\dim H^{1}(B)-\rank\res^{1}$.
\end{enumerate}
\end{lemma}

\begin{proof}
(1) is the proof of Lemma~\ref{lem:conerelative}(2) verbatim: it uses of $\rres$ only that it is a cochain map, which $(\alpha_{B/A})_{\ast}$ is
(Lemma~\ref{lem:levelresmorph}(3)). (2) By Theorem~\ref{thm:explicit_complex_equiv} the model complex of $L^{\chi}_A$ is $\Hom_{\Z[\Ga]}(C_{\bullet}(\Xbulk),M)$ with
$M=\Ind_A^{\Ga}\VR(\chi)$, and Proposition~\ref{prop:bredon}(2) --- which holds for \emph{every} $\Z[\Ga]$-module $M$ --- makes it chain homotopy equivalent to the
two-term complex \eqref{eq:treemodel}; so $H^{q}=0$ for $q\ge2$ and $H^{0}=M^{\langle S\rangle}\cap M^{\langle U\rangle}=M^{\Ga}$, likewise for $B$. The equivalence is a
chain map on the chain side, so it is natural in $M$ and $\res^{q}$ is the map induced by $\alpha_{B/A}$ on the tree complexes; on $H^{0}$ this is the restriction of
the injective $\alpha_{B/A}$ to $M^{\Ga}$. (3) Exactness of the long sequence and the two vanishings; $H^{3}(A,B)$ sits between $H^{2}(B)=0$ and $H^{3}(A)=0$.
\end{proof}

\noindent\emph{Contrast with the boundary cone of \S\ref{sec:cone}.} There $\rres$ is onto and $\Heq{2}=0$ is a hypothesis (Lemma~\ref{lem:conerelative}(5)); here
$\alpha_{B/A}$ is injective and the vanishing is unconditional, because both ends are cohomologies of the same model with different coefficient modules. The analogue
of $\ker\rres$ is $\coker\alpha_{B/A}$ (Proposition~\ref{prop:levelconeshapiro}).

\begin{proposition}[What the cone computes]\label{prop:levelconeshapiro}
\PROVED{} For a compatible level pair over a commutative $\cring$:
\begin{enumerate}
\item $Q_{A/B}:=\Ind_B^{A}\VR(\chi')\big/\VR(\chi)$ (the cokernel of the unit $v\mapsto(x\mapsto\rho_{n,\chi}(x)v)$) is an $\cring[A]$-module, free over $\cring$ of rank
  $([A:B]-1)(n+1)$, and $\coker\alpha_{B/A}\cong\Ind_A^{\Ga}Q_{A/B}$.
\item If $6\in\cring^{\times}$, then $\Cone{\alpha_{B/A}}\to\Ceq(K,\Rel;L_{\coker\alpha})$, $(c,c')\mapsto[c']$, is a quasi-isomorphism, and
  $H^{q}(A,B)\cong H^{q-1}(A;Q_{A/B})$ for every $q$ (Shapiro, Lemma~\ref{lem:inducedchar}).
\item At $\ell\in\{2,3\}$ (2) is \emph{not} asserted: the sequence \eqref{eq:levelsixterm} is a sequence of Bredon groups (Proposition~\ref{prop:bredon}(3)), and an end is
  group cohomology exactly when Proposition~\ref{prop:modeltorsion} applies to it.
\end{enumerate}
(For orientation only, not used: with this shift $H^{\bullet}(A;Q_{A/B})$ is the relative cohomology of the group pair $(A,B)$ in the sense of \cite{adamson1954} and of
\cite{bieri_eckmann1978}, $\Hom(\Delta,V)=Q_{A/B}$ for $\Delta=\ker(\Z[A/B]\to\Z)$; neither paper is held, no locator is claimed, \VERIFYloc.)
\end{proposition}

\begin{proof}
(1) Shapiro coordinates $f\mapsto(f(r_t))_t$ identify $\Ind_B^{A}\VR(\chi')$ with $\VR^{[A:B]}$ over $\cring$; the unit is $v\mapsto(\rho(r_t)v)_t$, whose image meets
$\{f:f(r_1)=0\}$ trivially and spans with it ($r_1=I$), so the quotient is free of the stated rank. Induction from $A$ to $\Ga$ is exact (it is a finite direct sum
over $\cring$ of copies of the module, Definition~\ref{def:inducedmodule}) and transitive, $\Ind_A^{\Ga}\Ind_B^{A}=\Ind_B^{\Ga}$, with $\Ind_A^{\Ga}(\text{unit})=\alpha_{B/A}$.
(2) For $6\in\cring^{\times}$ the fixed-point functors $M\mapsto M^{\Ga_\sigma}$ ($\abs{\Ga_\sigma}\mid12$) are exact (averaging), so
$0\to\Ceq(L_A)\to\Ceq(L_B)\to\Ceq(L_{\coker\alpha})\to0$ is exact (Theorem~\ref{thm:explicit_complex_equiv}: the model complex is $\bigoplus M^{\Ga_\sigma}$ in each degree).
The map of the statement is a cochain map ($\Dcone(c,c')\mapsto[dc']$), surjective, with kernel $\{(c,\alpha_{\ast}c''):c\in C^{q+1}_A,c''\in C^{q}_A\}\cong\Cone{\id_{C_A}}$,
which is acyclic; so it is a quasi-isomorphism. Then $H^{q-1}(\Ceq(L_{\Ind_AQ}))\cong H^{q-1}(\Ga;\Ind_A^{\Ga}Q_{A/B})\cong H^{q-1}(A;Q_{A/B})$ by
Lemma~\ref{lem:inducedchar} (Proposition~\ref{prop:lhcompliant}(4) for an arbitrary representation).
(3) is Proposition~\ref{prop:bredon}(3) applied to both ends.
\end{proof}

\begin{proposition}[When the model computes group cohomology at $\ell\in\{2,3\}$]\label{prop:modeltorsion}
\PROVED{} Let $B\le\Ga$ have finite index, $X$ an $\cring[B]$-module, and suppose that \emph{every finite subgroup of $B$ has order in $\cring^{\times}$}. Then
$\Heq{q}(K,\Rel;L^{X}_B)\cong H^{q}(B;X)$ for every $q$ (so $H^{q}(B;X)=0$ for $q\ge2$), whatever $\ell$. The finite subgroups of $\Ga$ have orders $1,2,3,4,6$, and every
subgroup of even order contains $-I$; so at $\ell=2$ the hypothesis is $-I\notin B$, and at $\ell=3$ it is ``$B$ contains no element of order $3$''. It holds at every
$\ell$ for torsion-free $B$: $\Ga_1(N)$ for $N\ge4$, and every subgroup of it (among them $\Ga(N)$, $N\ge4$, and $\Ga_1(p^{r})\cap\Ga_0(p^{r+1})$, $p^{r}\ge4$), and $\Ga(3)$.
$\Ga_0(N)$ contains an element of order $3$ iff $x^{2}+x+1\equiv0\pmod N$ is solvable (not for $2\mid N$ or $9\mid N$; yes for $N=3$; the full characterisation ---
iff $9\nmid N$ and no prime $p\equiv2\pmod3$ divides $N$ --- is Proposition~\ref{prop:descentNthree}, v1.23.0-pre.18, referee m2).
\end{proposition}

\begin{proof}
By Proposition~\ref{prop:bredon}(2) the model complex is chain homotopy equivalent to $\Hom_{\Z[\Ga]}(C_{\bullet}(\mathcal T),\Ind_B^{\Ga}X)\cong
\Hom_{\Z[B]}(C_{\bullet}(\mathcal T),X)$ (Frobenius reciprocity, $\phi\mapsto\phi(\cdot)(I)$, as in Proposition~\ref{prop:levelcorrespondence}(3)). The proof of
Lemma~\ref{lem:treeres} shows that $C_{\bullet}(\mathcal T)\otimes\cring\to\cring$ is a projective resolution over $\cring[B]$ as soon as the orders of the groups
$B\cap\gamma F\gamma^{-1}$ ($F\in\{\langle-I\rangle,\langle S\rangle,\langle U\rangle\}$) are units --- these are finite subgroups of $B$, which is our hypothesis. So the
complex computes $H^{\bullet}(B;X)$. Finite subgroups: they are cyclic, contained in a conjugate of $\langle S\rangle$ ($\cong\Z/4$) or of $\langle U\rangle$ ($\cong\Z/6$)
(Proposition~\ref{prop:bredon}: $\mathcal D$ is the family of finite subgroups), and the unique element of order $2$ in either is $-I$. Torsion: a torsion element
$\gamma\ne\pm I$ of $\Ga$ has eigenvalues primitive $k$-th roots of unity, $k\in\{3,4,6\}$, so $\tr\gamma\in\{-1,0,1\}$; on $\Ga_1(N)$, $\tr\gamma\equiv2\pmod N$, which is
$\not\equiv-1,0,1$ for $N\ge4$, and $-I\notin\Ga_1(N)$ for $N\ge3$. $\Ga(3)$: $-I\notin\Ga(3)$; order $4$ needs trace $0\not\equiv2\pmod3$; an element of order $6$ has a
square of order $3$; and $\gamma=I+3X$ of order $3$ would satisfy $\gamma^{2}+\gamma+I=0$ (Cayley--Hamilton, trace $-1$), but $\gamma^{2}+\gamma+I=3I+9(X+X^{2})\equiv3I\pmod9$. Order $3$ in $\Ga_0(N)$: $\gamma=\smat abcd$ of order $3$ has $a+d=-1$, and $N\mid c$ gives
$ad\equiv1$, i.e.\ $a^{2}+a+1\equiv0\pmod N$; conversely, for such an $a$, $\smat{a}{-(a^{2}+a+1)/N}{N}{-1-a}$ has determinant $1$ and trace $-1$. $a^{2}+a+1$ is odd, and
takes no value $\equiv0\pmod9$ (check $a\bmod9$); $a=1$ works for $N=3$.
\end{proof}

\begin{proposition}[Index prime to $\ell$: the split case]\label{prop:levelsplit}
\PROVED{} If $[A:B]\in\cring^{\times}$, then $e:=[A:B]^{-1}\alpha_{B/A}\omega_{B/A}$ is an idempotent $\cring[\Ga]$-endomorphism of $\Ind_B^{\Ga}\VR(\chi')$ with image
$\alpha_{B/A}(\Ind_A^{\Ga}\VR(\chi))$ and kernel $\ker\omega_{B/A}$, so $\Ind_B^{\Ga}\VR(\chi')=\alpha_{B/A}(\Ind_A)\oplus\ker\omega_{B/A}$ \emph{on the fibre}, at every $\ell$.
Hence every $\res^{q}$ is split injective, $H^{q}(B)=\res^{q}H^{q}(A)\oplus\ker\cor^{q}$, $H^{1}(A,B)\cong\coker\res^{0}\cong\ker\cor^{0}$ and $H^{2}(A,B)\cong\coker\res^{1}\cong
\ker\cor^{1}$. Instances: $\Ga_0(p)\subset\Ga$ at $\ell\nmid p+1$; the first step $\Ga_0(N)\cap\Ga_1(p)\trianglelefteq\Ga_0(pN)$ (index $p-1$) at $\ell=p$.
\end{proposition}

\begin{proof}
$e^{2}=[A:B]^{-2}\alpha(\omega\alpha)\omega=e$ by Lemma~\ref{lem:levelresmorph}(2), $e\alpha=\alpha$ and $\omega e=\omega$, so $\im e=\im\alpha$ and $\ker e=\ker\omega$. The
model complex is additive in the fibre module (Theorem~\ref{thm:explicit_complex_equiv}), so $\Ceq(L_B)=\alpha_{\ast}\Ceq(L_A)\oplus\Ceq(L_{\ker\omega})$ as complexes;
$\cor\circ\res=[A:B]$ is invertible, so $\res$ is split injective and $\ker\cor^{q}=H^{q}(L_{\ker\omega})$. Insert in \eqref{eq:levelsixterm}.
\end{proof}

\begin{proposition}[Hecke operators on the six-term sequence]\label{prop:levelconehecke}
Let $(A,\chi)\supseteq(B,\chi')$ be a compatible level pair with $A=\Ga_{\mathsf H}(N_A)\cap\Ga_0(N_A')$, $B=\Ga_{\mathsf H'}(N_B)\cap\Ga_0(N'_B)$ with $N_A\mid N'_A$, $N_B\mid N'_B$ (any of
the levels of Definition~\ref{def:levelpair} of this form; so $\Ga_1(N'_A)\subseteq A$, $\Ga_1(N'_B)\subseteq B$), $g=\smat100p$, and the operators of Propositions~\ref{prop:levelcorrespondence}, \ref{prop:nebcorr} in the reciprocity form
$\psi\mapsto\sum_t\rho_{n,\chi}(r_t^{-1}g)\,\psi(\tau(r_t\,\cdot))$, $r_t$ a transversal of $A''\backslash A$, $A'':=A\cap gAg^{-1}$, with the level-$1$ $\tau$ of $\Tdc{g}$.
\begin{enumerate}
\item \PROVED{} \textbf{(The criterion.)} If $B\cap A''=B'':=B\cap gBg^{-1}$ and $A=A''B$, then a transversal $r_t$ of $B''\backslash B$ is one of $A''\backslash A$, and with this
  common transversal $T^{B}_g\circ(\alpha_{B/A})_{\ast}=(\alpha_{B/A})_{\ast}\circ T^{A}_g$ \textbf{as cochain maps}, over every $\cring$; so $\Phi_g:=(T^{A}_g,T^{B}_g)$ is a
  cochain map of $\Cone{\alpha_{B/A}}$ commuting with the two maps of Lemma~\ref{lem:levelconeexact}(1), and the whole sequence \eqref{eq:levelsixterm} is equivariant ---
  \textbf{no pair homotopy} (contrast Definition~\ref{def:pairhomotopy}).
\item \PROVED{} \textbf{(Instances.)} The criterion holds for $T_p$, $p\nmid N_B$ (any pair), with $r_t\in B$ representing the $p+1$ classes of
  $\mathbb P^{1}(\F_p)$ --- and likewise for $A=\Ga$ and $B$ the preimage of a subgroup of $\SL_2(\Z/M)$ normalised by $\mathrm{diag}(1,p)$, $p\nmid M$ (the level-one cells $\Ga^{2}$, $\Ga^{3}$, $\Ga(M)$ of \ref{LES4}); for $U_p$, $p\mid N_A$, when $A$ and $B$ have the \emph{same} $\Ga_0$-level ($N'_A=N'_B$, the $\Ga_{\mathsf H}$/diamond direction), with
  $r_t=\smat{1}{0}{tN'_A}{1}\in B$, $t<p$; and the diamonds $R_\sigma$ ($\sigma\in A\cap\Ga_0(N'_B)$, Definition~\ref{def:nebgammaH}) commute with $\alpha_{B/A}$ on the fibre.
  So the shipped sub-transversals of \file{Hecke/HeckeGammaH} (\code{Gamma1SubTransversal}, \code{GammaHSubTransversal}, inside $\Ga_1\subseteq B$) are the ones to use:
  both levels' operators are built from the \textbf{smaller} group's transversal.
\item \PROVED{} \COMP\textsuperscript{a} \textbf{(Level-raising $U_p$: the square does not commute in general.)} If $p\mid N'_A$ and $pN'_A\mid N'_B$, then $B\subseteq A''$, so
  $A''B=A''\ne A$ and the criterion fails. Let $\mathrm{dg}_g:=[A\,g\,B]=(\conj_g)_{\ast}\circ\res$, $\psi\mapsto\rhon(g)\,\psi\circ\tau$, the one-term cochain map of
  Proposition~\ref{prop:levelcorrespondence}(4)(c) ($B\cap gAg^{-1}=B$, $r_1=I$).
  \begin{enumerate}
  \item[(a)] $T^{B}_g\circ(\alpha_{B/A})_{\ast}=[B:B'']\cdot\mathrm{dg}_g=p\cdot\mathrm{dg}_g$ \textbf{as cochain maps}, over every $\cring$. In particular
    $U^{B}_p\circ(\alpha_{B/A})_{\ast}=0$ on cochains when $p=0$ in $\cring$ (e.g.\ $\cring=\F_p$, $\ell=p$).
  \item[(b)] $T^{A}_g=(\omega_{A''/A})_{\ast}\circ[A\,g\,A'']$ as cochain maps (the factorisation $\cor\circ(\conj_g)_{\ast}\circ\res$ of
    Proposition~\ref{prop:levelcorrespondence}(3), read at cochain level), so $(\alpha_{B/A})_{\ast}\circ T^{A}_g=(\alpha_{B/A})_{\ast}(\omega_{A''/A})_{\ast}\circ[A\,g\,A'']$.
  \item[(c)] Let $B=A''$ and $6,p\in\cring^{\times}$. On $H^{q}$, $\res\circ U^{A}_p-U^{B}_p\circ\res=p\,(e-1)\,\mathrm{dg}_g$ with the idempotent $e$ of
    Proposition~\ref{prop:levelsplit}; so the square commutes on $H^{q}$ \textbf{iff $\mathrm{dg}_g(H^{q}(A))\subseteq\res H^{q}(A)$}. This \textbf{fails} on $H^{1}$ at
    $(\Ga_0(25)\subset\Ga_0(5),\ \ell=7,\ n=0,\ \chi=1)$.
  \item[(d)] On $p$-old spaces this document's $U^{(N)}_p$ is the \emph{adjoint} $U_p^{\ast}=W_pU_pW_p^{-1}$ of the classical operator: the relations
    $U\circ\mathrm{dg}_1=p\,\mathrm{dg}_p$, $U\circ\mathrm{dg}_p=\mathrm{dg}_p\circ T_p-p^{\,n}\mathrm{dg}_1$ of Proposition~\ref{prop:oldspace}(1),(3) are those of $U_p^{\ast}$ (the
    classical $U_p$ has $U_p\circ\mathrm{dg}_p=\mathrm{dg}_1$, $U_p\circ\mathrm{dg}_1=\mathrm{dg}_1\circ T_p-p^{\,k-1}\mathrm{dg}_p$), and $\res$ is the inclusion of old forms
    under any Eichler--Shimura normalisation, so the identification is convention-free. The classical $U_p$, which commutes with the inclusion of old forms, is
    $[A\,g'A]$ with $g':=\smat p001=SgS^{-1}$ and the twisted datum $\tau':=S\circ\tau\circ S^{-1}$ (it satisfies (T1)--(T3) of Definition~\ref{def:tau} for $g'$, on the same
    level-$1$ table); at $\chi=1$ the criterion (1) holds for it on \emph{every} level-raising pair, with the common transversal $T^{b}\in B$, $b<p$
    ($A\cap g'Ag'^{-1}=A\cap\Ga^{0}(p)$, $B\cap g'Bg'^{-1}=B\cap(A\cap\Ga^{0}(p))$, $A=(A\cap\Ga^{0}(p))B$), so its square is strict. That construction (and the
    character convention for $g'$, whose $a$-entry is $p$) is not needed in slot~22 and stays deferred (Remark~\ref{rmk:oldspacecone}).
  \end{enumerate}
  Obligation~\ref{obl:levelconeexact}(vi) gates (a) and records the square on $H^{1}$. \COMP\textsuperscript{a} evidence:
  \file{docs/implementation/slot22/les0/levelraise\_up\_check.py} ($H^{1}$ at $n=0$ as $\Hom(K,\F_\ell)$, independent of the tree complex and of $\tau$): (a) on $H^{1}$, the
  failure of (c), and the commuting classical square at $(\Ga_0(25)\subset\Ga_0(5),\ \ell\in\{5,7,11\})$, $(\Ga_0(49)\subset\Ga_0(7),\ \ell=5)$,
  $(\Ga_0(121)\subset\Ga_0(11),\ \ell=7)$, output \file{les0/out\_levelraise\_up.txt}, $0$ mismatches.
\item \PROVED{} \textbf{(Independence of the transversal.)} Let $\psi$ be $A$-equivariant and $b\in A''\subseteq\Ga''_g$. By (T1) of Definition~\ref{def:tau},
  $\tau(b\,r_t\,y)=(g^{-1}bg)\,\tau(r_t\,y)$ with $g^{-1}bg\in A$, hence
  $\rho_{n,\chi}\bigl((br_t)^{-1}g\bigr)\psi(\tau(br_ty))=\rho_{n,\chi}(r_t^{-1}g)\,\psi(\tau(r_ty))$ (Proposition~\ref{prop:symconv}(2); the scalars agree by multiplicativity of $a$ modulo $N_A$ on $\Delta_0(N_A)$, Lemma~\ref{lem:nebdata}(2), equivalently $a_{xg}=a_x$ for $g=\mathrm{diag}(1,p)$ together with multiplicativity on $\Ga_0(N_A)$, and $a_{g^{-1}bg}=a_b$). So each term of the reciprocity form depends only on the coset $A''r_t$, and $T^{A}_g$ on $\Ceq(K,\mathrm{Rel};L^{\chi}_A)$ does \emph{not} depend on the
  transversal. Consequently, under the criterion of (1), $T^{B}_g\circ(\alpha_{B/A})_{\ast}=(\alpha_{B/A})_{\ast}\circ T^{A}_g$ as cochain maps for \textbf{every} pair of
  transversals of $A''\backslash A$ and $B''\backslash B$ --- the common transversal of (1) is a device of the proof, not a requirement --- and no homotopy $\hder$ and no ambiguity of
  Proposition~\ref{prop:coneambiguity}(1) arise. (The operator of \eqref{eq:levelcorr} as implemented is transversal-independent even on the whole flat cochain space
  $C^{q}(K;L^{\chi}_A)$: it uses $r_t$ only through the indices $j(i,t)$, $a(i,t)$ of \eqref{eq:levelcorr} (the class of $r_tv_i$ in $\Ga''_g\backslash\Ga$ and the class in $A\backslash\Ga$ of the conjugate $g^{-1}r_tv_is_j^{-1}g$, read off \code{BuildLevelCorrespondence} and \code{BuildLevelCorrespondenceChi}), both functions of the coset $A''r_t$; so the independence
  is not specific to equivariant cochains in the implementation, and no record on that formula separates it from (T1).) When the criterion of (1) fails the square need not
  commute even on cohomology, (3)(c). \textbf{Rule} (as in (2)): LES code builds both operators from the \textbf{smaller} group's transversal, for definiteness; no certificate depends on that choice.
\end{enumerate}
\end{proposition}

\begin{proof}
(1) $B\cap A''=B''$ makes $B''\backslash B\to A''\backslash A$ injective, and $A=A''B$ makes it surjective, so a transversal of the former is one of the latter. In the
reciprocity form both operators are $\psi\mapsto\sum_t\rho_{n,\chi}(r_t^{-1}g)\psi(\tau(r_t\,\cdot))$ (the scalar $\chi(a_{r_t^{-1}g})=\chi(a_{r_t})^{-1}$ equals $\chi'(a_{r_t})^{-1}$, as
$r_t\in B$ and $a_g=1$; Lemma~\ref{lem:nebcompat}), with the same $\tau$, and $(\alpha_{B/A})_{\ast}$ is $\psi\mapsto\psi$ (Lemma~\ref{lem:levelresmorph}(3)); so the two
composites are the same formula. $\Dcone\Phi_g=\Phi_g\Dcone$ is then the computation of Lemma~\ref{lem:conesign} with $\hder=0$.
(2) $gAg^{-1}\cap\Ga=\{\smat{a}{b'}{c'}{d}:c'=pc,\ \smat{a}{pb'}{c}{d}\in A\}$, so $A''=A\cap\Ga_0(pN'_A)$ for $p\mid N'_A$ and $A''=A\cap\Ga_0(p)$ for $p\nmid N_A N'_A$ (the
$a$-, $d$-conditions are unchanged), and likewise for $B$. \emph{$T_p$, $p\nmid N_B$}: $B\cap A''=B\cap\Ga_0(p)=B''$; $A''\backslash A$ is the set of the $p+1$ points of
$\mathbb P^{1}(\F_p)$ (bottom rows modulo $p$), and $B\to\SL_2(\F_p)$ is onto ($B\supseteq\Ga_1(\mathrm{lcm})$, Chinese remainder theorem, $p$ prime to the levels of $B$), so
$A=A''B$. The same argument applies to $A=\Ga$ ($A''=\Ga_0(p)$) and $B$ the preimage of a subgroup $\bar B\le\SL_2(\Z/M)$ normalised by $D:=\mathrm{diag}(1,p)\bmod M$, $p\nmid M$:
$B\supseteq\Ga(M)$ and $\Ga(M)\to\SL_2(\F_p)$ is onto, and $x\in\Ga$ lies in $gBg^{-1}$ iff $p\mid c_x$ and $D^{-1}xD\in\bar B$ modulo $M$, so $B\cap gBg^{-1}=B\cap\Ga_0(p)$. This covers
$\Ga(M)$ and $\Ga^{2}$, $\Ga^{3}$ (preimages of the characteristic subgroups $A_3$, $Q_8$), i.e.\ the level-one cells of Remark~\ref{rmk:normalN1}. \emph{$U_p$, same level $N':=N'_A=N'_B$, $p\mid N'$}: here $N_A\mid N'_A$ and
$N_B\mid N'_B$ (the levels of $\Ga_{\mathsf H}(N)\cap\Ga_0(N')$ are written with $N\mid N'$), so $\Ga_1(N')\subseteq B$; $B\cap A''=B\cap\Ga_0(pN')=B''$; the right cosets of
$A''=A\cap\Ga_0(pN')$ in $A$ are the classes of the invariant $a_\gamma c_\gamma/N'\bmod p$ (the map of Proposition~\ref{prop:gamma0normal}, here with $M=p$; it is a
coset invariant because $(h\gamma)_{11}(h\gamma)_{21}/N'\equiv a_\gamma c_\gamma/N'\pmod p$ for $h\in A''$), and $\smat10{tN'}1$ ($t<p$) lies in $\Ga_1(N')\subseteq B$, has
invariant $t$, and so meets each class. \emph{Diamonds}: $\alpha_{B/A}$ is the inclusion of functions and $R_\sigma f(x)=\rho_{n,\chi}(\sigma)f(\sigma^{-1}x)$ is the same
formula on both modules.

(3) For $\gamma\in B$, $pN'_A\mid c_\gamma$, so $B\subseteq A\cap\Ga_0(pN'_A)=A''$ by the computation of (2). (a) $B''=B\cap\Ga_0(pN'_B)$ by (2), of index $p$ in $B$ (the
invariant $a_\gamma c_\gamma/N'_B\bmod p$, as in (2)); let $r_t$ ($t<p$) be a transversal. For an $A$-equivariant $\psi$ and $r_t\in B\subseteq A''\subseteq\Ga''_g$, (T1) of
Definition~\ref{def:tau} gives $\tau(r_ty)=(g^{-1}r_tg)\,\tau(y)$ with $g^{-1}r_tg\in A$ (as $r_t\in gAg^{-1}$), so $\psi(\tau(r_ty))=\rho_{n,\chi}(g^{-1}r_tg)\psi(\tau(y))$
and the $t$-th term of $T^{B}_g(\alpha_{B/A})_{\ast}\psi$ is $\rho_{n,\chi'}(r_t^{-1}g)\rho_{n,\chi}(g^{-1}r_tg)\psi(\tau(y))=\rhon(g)\psi(\tau(y))$: the $\rhon$-factors
multiply to $\rhon(g)$ (Proposition~\ref{prop:symconv}), and the scalars $\chi'(a_{r_t})^{-1}\chi(a_{g^{-1}r_tg})$ cancel, since $a_{g^{-1}r_tg}=a_{r_t}$ and
$\chi'(a_{r_t})=\chi(a_{r_t})$ (Lemma~\ref{lem:nebcompat}). All $p$ terms equal $\mathrm{dg}_g\psi$. This is the proof of Proposition~\ref{prop:oldspace}(1) with
$(\Ga,\Ga_0(\ell))$ replaced by $(A,B)$. (b) With $r_t$ a transversal of $A''\backslash A$ and $a_g=1$, $\rho_{n,\chi}(r_t^{-1}g)\psi(\tau(r_ty))=\rho_{n,\chi}(r_t)^{-1}
\bigl(\rhon(g)\psi(\tau(r_ty))\bigr)$; the bracket is $([A\,g\,A'']\psi)(r_ty)$ (Proposition~\ref{prop:levelcorrespondence}(4)(c): one term) and the sum over $t$ is the
transfer $(\omega_{A''/A})_{\ast}$ of Lemma~\ref{lem:levelresmorph}(3). (c) By (a), (b) and $B=A''$, on $H^{q}$: $\res\circ U^{A}_p=\res\circ\cor\circ\mathrm{dg}_g$ and
$U^{B}_p\circ\res=p\,\mathrm{dg}_g$. As $[A:A'']=p\in\cring^{\times}$ (the invariant of (2)), $\res\circ\cor=(\alpha_{B/A}\omega_{B/A})_{\ast}=p\,e$ with $e$ the idempotent of
Proposition~\ref{prop:levelsplit}, whose image on $H^{q}(B)$ is $\res H^{q}(A)$ \STD{\cite[Ch.~III.9--10]{brown}, through Lemma~\ref{lem:levelresmorph}(3); locator read at
an earlier R20-L pass, Prop.~III.9.5, p.~82}. So the defect is $p(e-1)\mathrm{dg}_g$, zero iff $\mathrm{dg}_g(H^{q}(A))\subseteq\ker(e-1)=\res H^{q}(A)$. \emph{The cell.} At
$n=0$, $\chi=1$, $\cring=\F_7$ the coefficients are trivial, so $H^{1}(K;\F_7)=\Hom(K,\F_7)$ for $K\in\{A,B\}$. $A=\Ga_0(5)$ is generated by $-I$, $T$, $e_1:=\smat2{-1}5{-2}$,
$e_2:=\smat{-2}{-1}52$ (both of order $4$), and $P:=\smat1051=e_1Te_2$ (generation: \file{levelraise\_up\_check.py} reduces every Schreier generator of $\Ga_0(5)$ to
$\pm I$ by these). As $4\in\F_7^{\times}$, every $E\in\Hom(A,\F_7)$ kills $-I,e_1,e_2$, so $\dim H^{1}(A;\F_7)\le1$, and $=1$ (\ref{LES3}, arm~(i) $p=5$, $n=0$, read from
$B=\Ga_0(5)$ there); for $E\ne0$ normalise $E(T)=1$; then $E(P)=E(e_1)+E(T)+E(e_2)=1$.
Then $\mathrm{dg}_gE(\gamma)=E(g^{-1}\gamma g)=E\smat{a}{5b}{c/5}{d}$, so $\mathrm{dg}_gE(T)=E(T^{5})=5$ and $\mathrm{dg}_gE\smat10{25}1=E(P)=1$, while $\res E(T)=1$ and
$\res E\smat10{25}1=E(P^{5})=5$. $\mathrm{dg}_gE=\lambda\,\res E$ would need $\lambda=5$ and $5\lambda=1$, i.e.\ $25\equiv1\pmod7$, which is false; so
$\mathrm{dg}_gE\notin\res H^{1}(A)$ and the square fails on $H^{1}$. (The argument needs only $E(P)\ne0$ and $p^{2}\not\equiv1\pmod\ell$.) (d) The two sets of relations
are read off Proposition~\ref{prop:oldspace}(1),(3) and the classical $q$-expansion identities $U_p(f(pz))=f(z)$, $U_pf=T_pf-p^{k-1}f(pz)$, $k=n+2$. For $g'$:
$SgS^{-1}=\smat p001$; $\tau'$ is a chain map preserving the augmentation, and for $h\in\Ga\cap g'\Ga g'^{-1}=S\Ga''_gS^{-1}$,
$\tau'(hx)=S\tau\bigl((S^{-1}hS)S^{-1}x\bigr)=(g'^{-1}hg')\tau'(x)$. $g'\smat abcd g'^{-1}=\smat{a}{pb}{c/p}{d}$, so $g'Ag'^{-1}\cap\Ga\supseteq A\cap\Ga^{0}(p)$ and
$A\cap g'Ag'^{-1}=A\cap\Ga^{0}(p)$ (intersecting with $A$ restores the $c$-condition; the $a$-, $d$-conditions are unchanged), likewise for $B$; for $\gamma\in A\subseteq\Ga_0(p)$
the right coset $(A\cap\Ga^{0}(p))\gamma$ is determined by $b_\gamma a_\gamma^{-1}\bmod p$ (left multiplication by $h\in\Ga^{0}(p)\cap\Ga_0(p)$ scales $(a,b)$ by $a_h$ modulo
$p$, and $\gamma_1\gamma_2^{-1}\in\Ga^{0}(p)$ iff $b_1a_2\equiv a_1b_2$), and $T^{t}$ has invariant $t$, so the $T^{b}\in\Ga_1(N'_B)\subseteq B$ represent
$(A\cap\Ga^{0}(p))\backslash A$ and $(B\cap\Ga^{0}(p))\backslash B$ alike.
(4) Replacing $r_t$ by $br_t$, $b\in A''$, the displayed computation of the statement uses (T1) for $A''\subseteq\Ga''_g$, $g^{-1}bg\in A$ (as $b\in gAg^{-1}$), multiplicativity of
$\rho_n$ on $\Delta$ (both factors $\rho_n((br_t)^{-1}g)$ and $\rho_n(g^{-1}bg)$ being integral with product $\rho_n(r_t^{-1}g)$), and multiplicativity of $a$ modulo $N_A$ on $\Ga_0(N_A)$
(Lemma~\ref{lem:nebdata}(2)); this is the well-definedness of the double-coset sum, performed for $r_t\in B\subseteq A''$ in the proof of Proposition~\ref{prop:levelcorrespondence}(3)(a).
Under the criterion, pick a common transversal $c$; then for any transversals $a$, $b$, $T^{B}(b)\alpha_{\ast}\psi=T^{B}(c)\alpha_{\ast}\psi=\alpha_{\ast}T^{A}(c)\psi=\alpha_{\ast}T^{A}(a)\psi$,
using only that $\alpha_{\ast}$ maps $\Ceq(L_A)$ into $\Ceq(L_B)$ (Lemma~\ref{lem:levelresmorph}(3)). The parenthetical remark is read off \code{BuildLevelCorrespondence}.
\end{proof}

\begin{definition}[Normal pair]\label{def:normalpair}
A compatible level pair is \emph{normal} if $B\trianglelefteq A$; put $\mathsf G:=A/B$, finite, and $\mathrm{Nm}_{\mathsf G}:=\sum_{g\in\mathsf G}R_g$. For $\sigma\in A$ the
\emph{conjugation} $R_\sigma$ on $\Ind_B^{\Ga}\VR(\chi')$ is $(R_\sigma f)(x):=\rho_{n,\chi}(\sigma)f(\sigma^{-1}x)$ --- Proposition~\ref{prop:nebdecomp}(1), Definition~\ref{def:nebgammaH}
for an arbitrary normal $B$; for $B=\Ga_{\mathsf H'}(N)\trianglelefteq A=\Ga_0(N)$ it is the diamond. In coordinates, $(R_\sigma f)_J=\rho_{n,\chi}(\sigma h)f_{J'}$ where
$\sigma^{-1}v_J=hv_{J'}$, $h\in B$.
\end{definition}

\begin{proposition}[The normal case is inflation--restriction]\label{prop:levelinfres}
Let the pair be normal, $\mathsf G=A/B$.
\begin{enumerate}
\item \PROVED{} \textbf{(Fibre identities, every $\cring$, every $\ell$.)} $R_\sigma$ is a $\cring[\Ga]$-automorphism of $\Ind_B^{\Ga}\VR(\chi')$ depending only on
  $\sigma B$, with $R_\sigma R_{\sigma'}=R_{\sigma\sigma'}$; $\alpha_{B/A}(\Ind_A)=(\Ind_B)^{\mathsf G}$; $\omega_{B/A}\alpha_{B/A}=[A:B]$ and
  $\alpha_{B/A}\omega_{B/A}=\mathrm{Nm}_{\mathsf G}$; $\varphi(x):=\rho_{n,\chi}(x)^{-1}f(x)$ ($x\in A$) identifies $\Ind_B^{A}\VR(\chi')$ with $\mathrm{Map}(\mathsf G,\VR)$, $R_\sigma$
  becoming left translation and $\omega_{B/A}$ the sum, so $\ker\omega_{B/A}=\sum_{g}(R_g-1)\Ind_B$. For $6\in\cring^{\times}$, $R_\sigma$ induces the conjugation action of
  $\sigma$ on $H^{\bullet}(B;\VR(\chi'))$, and $\res^{q}$ lands in $H^{q}(B)^{\mathsf G}$.
\item \STD{\cite[Ch.~VII, \S6, Thm.~6.3 and Cor.~6.4]{brown}, the Hochschild--Serre spectral sequence and its five-term sequence; \VERIFYloc} \textbf{(Two rows.)}
  For $6\in\cring^{\times}$: $E_2^{a,b}=H^{a}(\mathsf G;H^{b}(B;\VR(\chi')))\Rightarrow H^{a+b}(A;\VR(\chi))$ has only the rows $b\in\{0,1\}$ and abuts to zero in total degree
  $\ge2$. Hence $d_2\colon H^{a}(\mathsf G;H^{1}(B))\to H^{a+2}(\mathsf G;H^{0}(B))$ is an isomorphism for every $a\ge1$, and
  \begin{equation}\label{eq:levelfiveterm}
    0\to H^{1}(\mathsf G;H^{0}(B))\xrightarrow{\ \inf\ }H^{1}(A)\xrightarrow{\ \res\ }H^{1}(B)^{\mathsf G}\xrightarrow{\ \mathrm{tg}\ }H^{2}(\mathsf G;H^{0}(B))\to0
  \end{equation}
  is exact ($H^{0}(B)=\VR(\chi')^{B}$).
\item \PROVED{} \textbf{(The cone sequence is \eqref{eq:levelfiveterm} plus two pieces.)} For $6\in\cring^{\times}$, \eqref{eq:levelsixterm} gives
  $0\to H^{0}(B)/H^{0}(A)\to H^{1}(A,B)\to\ker\res^{1}\to0$ with $\ker\res^{1}=\inf H^{1}(\mathsf G;H^{0}(B))\cong H^{1}(\mathsf G;H^{0}(B))$, and
  $H^{2}(A,B)=\coker\res^{1}\supseteq H^{1}(B)^{\mathsf G}/\res H^{1}(A)\cong H^{2}(\mathsf G;H^{0}(B))$ with quotient $H^{1}(B)/H^{1}(B)^{\mathsf G}$: inflation--restriction is the
  part of the cone sequence inside the $\mathsf G$-invariants, and $H^{0}(B)/H^{0}(A)$, $H^{1}(B)/H^{1}(B)^{\mathsf G}$ are the two pieces it leaves out.
\item \PROVED{} If $\abs{\mathsf G}\in\cring^{\times}$: $H^{\ge1}(\mathsf G;-)=0$, $H^{q}(A)\cong H^{q}(B)^{\mathsf G}$, and Proposition~\ref{prop:levelsplit} applies; for
  $B=\Ga_{\mathsf H}(N)\trianglelefteq A=\Ga_0(N)$ and $\cring$ holding the values of the characters of $G_N/\mathsf H$, $H^{q}(B)$ is the sum of the character models
  (Propositions~\ref{prop:nebdecomp}(3), \ref{prop:nebgammaH}(6)).
\item \PROVED{} \textbf{(Coinvariants.)} For $6\in\cring^{\times}$, $\cor^{1}$ induces $H^{1}(B)_{\mathsf G}\xrightarrow{\ \sim\ }H^{1}(A)$, and $\res^{1}\circ\cor^{1}$ is the norm map
  $H^{1}(B)_{\mathsf G}\to H^{1}(B)^{\mathsf G}$ of the $\mathsf G$-module $H^{1}(B)$.
\end{enumerate}
\end{proposition}

\begin{proof}
(1) $R_\sigma$: the proof of Proposition~\ref{prop:nebdecomp}(1) with $B$ for $\Ga_1(N)$ and $\rho_{n,\chi}$ for $\rhon$: it uses only $\sigma^{-1}B\sigma=B$, the multiplicativity of
$\rho_{n,\chi}$ on $A$ and $\rho_{n,\chi}|_B=\rho_{n,\chi'}$ (Lemma~\ref{lem:nebcompat}); it also gives the induced action on cohomology. Coordinates: $\sigma^{-1}v_J=hv_{J'}$ gives
$f(\sigma^{-1}v_J)=\rho(h)f_{J'}$. $f$ is $\mathsf G$-fixed iff $f(\sigma x)=\rho(\sigma)f(x)$ for all $\sigma\in A$, i.e.\ iff $f\in\Ind_A$. With $B\backslash A=\mathsf G$ the $r_t$
of Definition~\ref{def:levelres} run over $\mathsf G$, and $(\alpha\omega f)(x)=\sum_t\rho(r_t^{-1})f(r_tx)=\sum_tR_{r_t^{-1}}f(x)$. $\varphi(bx)=\rho(x)^{-1}\rho(b)^{-1}\rho(b)f(x)=
\varphi(x)$, so $\varphi$ is a function on $B\backslash A=\mathsf G$, arbitrary; $(R_\sigma f)$ has $\varphi'(x)=\rho(x)^{-1}\rho(\sigma)f(\sigma^{-1}x)=\varphi(\sigma^{-1}x)$, and
$\omega f(I)=\sum_t\varphi(r_t)$. The kernel of the sum on $\mathrm{Map}(\mathsf G,\VR)$ is spanned by the $(\delta_g-\delta_1)v=(R_g-1)(\delta_1v)$. Since $R_\sigma\alpha=\alpha$
(fixed points), $R_{\sigma\ast}\res^{q}=\res^{q}$. (2) The spectral sequence of the extension $1\to B\to A\to\mathsf G\to1$ (the citation); for $6\in\cring^{\times}$,
$H^{b}(B;-)=0$ for $b\ge2$ and $H^{a+b}(A;-)=0$ for $a+b\ge2$ (Lemma~\ref{lem:treeres}). With two rows, $E_3=E_\infty$, and $E_\infty^{a,1}=\ker d_2$, $E_\infty^{a+2,0}=\coker d_2$
vanish for $a\ge1$; for $a=0$, $E^{2,0}_\infty=0$ makes $\mathrm{tg}=d_2$ onto. (3) Read \eqref{eq:levelsixterm} with (2); $H^{0}(B)/H^{0}(A)=\coker\res^{0}$, and $\res$ lands in
the invariants by (1). (4) $H^{\ge1}(\mathsf G;M)$ is killed by $\abs{\mathsf G}$ \STD{\cite[Ch.~III, Cor.~10.2]{brown}; \VERIFYloc}. (5) $H^{1}(A;-)$ is right exact and additive
($H^{2}(A;-)=0$), so it preserves the cokernel $\Ind_B^{A}\VR\to(\Ind_B^{A}\VR)_{\mathsf G}=\Ind_B^{A}\VR/\ker\omega\cong\VR$ of (1); by Shapiro
$H^{1}(A;\Ind_B^{A}\VR)=H^{1}(B;\VR)$, the map induced by $\omega$ being $\cor$ (Lemma~\ref{lem:levelresmorph}(3)) and $R_g$ the conjugation. $\res\circ\cor=\mathrm{Nm}_{\mathsf G}$ by
(1).
\end{proof}

\begin{proposition}[A cyclic quotient of order $p$]\label{prop:cycliccone}
Let $\mathsf G=\langle\sigma\rangle\cong\Z/p$ and $N:=\sum_{i<p}\sigma^{i}$.
\begin{enumerate}
\item \STD{\cite[Ch.~III, \S1, Example~2]{brown}, the periodic resolution of a cyclic group; \VERIFYloc} For an $\cring[\mathsf G]$-module $M$: $H^{0}=\ker(\sigma-1)$,
  $H^{2k-1}=\ker N/(\sigma-1)M$, $H^{2k}=\ker(\sigma-1)/NM$ ($k\ge1$).
\item \PROVED{} Over a field $\F$ of characteristic $p$: $\F[\mathsf G]=\F[t]/(t^{p})$, $t=\sigma-1$, $N=t^{p-1}$; a finite-dimensional module is $\bigoplus J_k$ (Jordan blocks of $\sigma$,
  $1\le k\le p$), $J_p$ the free module; $\dim H^{a}(\mathsf G;J_k)=1$ for $k<p$, $a\ge1$, and $=0$ for $k=p$. So $\dim H^{a}(\mathsf G;M)$ ($a\ge1$) is the number of
  \emph{non-free} blocks, and the Jordan type is read from $\rank t^{k}$, $k=1,\dots,p$.
\item \PROVED{} \textbf{(The level step.)} Let the pair be normal with $\mathsf G\cong\Z/p$ and $\cring=\F$ of characteristic $p\ge5$. Then
  \begin{enumerate}
  \item the number of non-free Jordan blocks of $\sigma$ on $H^{1}(B)$ equals that on $H^{0}(B)=\VR(\chi')^{B}$;
  \item if $H^{0}(B)=0$, then $H^{1}(B)$ is a \textbf{free} $\F[\mathsf G]$-module, and $H^{1}(A)\cong H^{1}(B)^{\mathsf G}\cong H^{1}(B)_{\mathsf G}$ (one $\Z/p$ step of the
    freeness of [U,~\texttt{prop:u-towerfree}], for the whole of $H^{1}$);
  \item if $A$ is torsion-free and $H^{0}(B)=H^{0}(A)$, of dimension $c$, then with $h:=\dim H^{1}(A)$: $\dim H^{1}(B)=ph-(p-1)c$, the Jordan type of $\sigma$ on $H^{1}(B)$ is
    $J_p^{\,h-c}\oplus J_1^{\,c}$, $\dim\ker\res^{1}=c$, $\rank\res^{1}=h-c$, $\dim H^{1}(A,B)=c$, $\dim H^{2}(A,B)=(p-1)h-(p-2)c$, and
    $\dim H^{1}(B)^{\mathsf G}=\dim H^{1}(B)_{\mathsf G}=h$;
  \item if $H^{0}(B)=H^{0}(A)=\F$ (trivial action) and $A$ is arbitrary: exactly one non-free block, of size $s=\dim H^{1}(B)-p(h-1)$, and $h-1$ free blocks.
  \end{enumerate}
\end{enumerate}
\end{proposition}

\begin{proof}
(2) $\sigma^{p}-1=(\sigma-1)^{p}$ in characteristic $p$, so $N=(\sigma^{p}-1)/(\sigma-1)=t^{p-1}$; $\F[t]/(t^{p})$ is a principal ideal ring with ideals $(t^{k})$, so a finitely
generated module is a sum of cyclic ones $J_k=\F[t]/(t^{k})$. On $J_k$, $k<p$: $N=0$, $\ker N/tJ_k$ and $\ker t/NJ_k=\ker t$ are one-dimensional; on $J_p$: $\ker N=tJ_p$,
$\ker t=t^{p-1}J_p=NJ_p$. The number of blocks of size $\ge k$ is $\rank t^{k-1}-\rank t^{k}$. (3)(a) By Proposition~\ref{prop:levelinfres}(2) ($6\in\F^{\times}$),
$H^{1}(\mathsf G;H^{1}(B))\cong H^{3}(\mathsf G;H^{0}(B))$, and $\dim H^{3}=\dim H^{1}$ for cyclic $\mathsf G$ by (1); apply (2) to both. (b) No non-free block; with $H^{0}(B)=0$
\eqref{eq:levelfiveterm} gives $H^{1}(A)\cong H^{1}(B)^{\mathsf G}$, and Proposition~\ref{prop:levelinfres}(5) the coinvariants. (c) $\mathsf G$ acts trivially on
$H^{0}(B)=H^{0}(A)=H^{0}(B)^{\mathsf G}$, so $H^{0}(B)\cong J_1^{c}$ and $\dim H^{a}(\mathsf G;H^{0}(B))=c$ ($a=1,2$). Lemma~\ref{lem:levelfreeeuler} for the torsion-free $A$ and its
subgroup $B$ of index $p$ gives $\dim H^{1}(B)-c=p(h-c)$. By (a) $H^{1}(B)$ has exactly $c$ non-free blocks; by \eqref{eq:levelfiveterm} the number of all blocks,
$\dim H^{1}(B)^{\mathsf G}$, is $\rank\res^{1}+c=(h-c)+c=h$, since $\ker\res^{1}=\inf H^{1}(\mathsf G;H^{0}(B))$ has dimension $c$. So $h-c$ blocks are free, and the $c$ non-free
blocks have total size $\dim H^{1}(B)-p(h-c)=c$: each has size $1$. The cone dimensions are Lemma~\ref{lem:levelconeexact}(3); the coinvariants
Proposition~\ref{prop:levelinfres}(5). (d) The same count with $c=1$, without Lemma~\ref{lem:levelfreeeuler}: one non-free block by (a), $h$ blocks in all, so $h-1$ free ones
and the rest has dimension $\dim H^{1}(B)-p(h-1)$.
\end{proof}

\begin{lemma}[Torsion-free groups: the Euler characteristic]\label{lem:levelfreeeuler}
\PROVED{} If $A\le\Ga$ has finite index and is torsion-free, then over every field, at every $\ell$ and for every character datum,
$\dim H^{1}(A)-\dim H^{0}(A)=m_A(n+1)/12$ on the model. In particular it is multiplicative along torsion-free pairs: $\dim H^{1}(B)-\dim H^{0}(B)=[A:B](\dim H^{1}(A)-\dim H^{0}(A))$.
\end{lemma}

\begin{proof}
Let $M=\Ind_A^{\Ga}\VR(\chi)$, of dimension $m_A(n+1)$. For $F\in\{\langle-I\rangle,\langle S\rangle,\langle U\rangle\}$, of order $f\in\{2,4,6\}$, Mackey decomposes $M|_F$ over
the double cosets $A\gamma F$ into modules induced from $F\cap\gamma^{-1}A\gamma=\{1\}$; so $M|_F$ is free over $\F[F]$ and $\dim M^{F}=\dim M/f$ (one fixed line per free rank-one
summand, spanned by its norm). Proposition~\ref{prop:bredon}(2): $\dim H^{1}-\dim H^{0}=\dim M^{-I}-\dim M^{S}-\dim M^{U}=\dim M\,(\tfrac12-\tfrac14-\tfrac16)$.
\end{proof}

\begin{lemma}[Invariants of a unipotent image]\label{lem:unipotentinv}
\PROVED{} Let $\ell=p$, $\cring=\F_p$, $\chi'=1$, and let $B\le\Ga$ contain $T$ and lie in $\Ga_1(p)$ (e.g.\ $\Ga_1(p^{r})\cap\Ga_0(N)$, $\Ga_1(p^{r})\cap\Ga_0(p^{r+1}N)$, $r\ge1$).
Then $\dim\VR^{B}=\lfloor n/p\rfloor+1$. If $A\supseteq B$ also has these two properties, then $\VR^{B}=\VR^{A}$, with trivial action of $A/B$. For the character
$\chi=\omega^{a-2}$ ($n=0$) on $B=\Ga_{\mathsf H}(p^{r})$ with $\mathsf H\supseteq\mu_{p-1}$ (the tower of \S\ref{sec:nebwildsec}): $\F_p(\chi)^{B}=\F_p$ if $a\equiv2\pmod{p-1}$ and $0$
otherwise.
\end{lemma}

\begin{proof}
$\VR=\Sym^{n}\F_p^{2}$ is a representation of $\SL_2(\F_p)$, and the image of $B$ there is $\{\smat1b01\}$ ($c\equiv0$, $d\equiv1$ force $a\equiv1$), generated by $T$.
$\rhon(T)$ is the substitution $x\mapsto x$, $y\mapsto y+x$ (Proposition~\ref{prop:symconv}). Put $\mathcal N:=\prod_{u\in\F_p}(y+ux)=y^{p}-x^{p-1}y$, which is invariant;
$\F_p[x,y]$ is free over $\F_p[x,\mathcal N]$ on $1,y,\dots,y^{p-1}$ ($\mathcal N$ is monic of degree $p$ in $y$). For an invariant $P=\sum_{k<p}c_ky^{k}$, $c_k\in\F_p[x,\mathcal N]$,
$0=P(x,y+x)-P(x,y)=\sum_{k}c_k((y+x)^{k}-y^{k})$, and if $k_0\ge1$ is the largest $k$ with $c_k\ne0$, the coefficient of $y^{k_0-1}$ in this basis is $k_0xc_{k_0}\ne0$ ($k_0<p$).
So the invariants are $\F_p[x,\mathcal N]$, whose degree-$n$ part has the basis $x^{n-pj}\mathcal N^{j}$, $0\le j\le\lfloor n/p\rfloor$. The second sentence: both images are the
same group. The third: $\gamma\in B$ acts by $\omega(a_\gamma)^{a-2}$, and $a_\gamma\bmod p$ runs over $\F_p^{\times}$ ($d_\gamma\bmod p$ runs over the reduction of $\mathsf H$,
which contains $\mu_{p-1}$).
\end{proof}

\begin{proposition}[Which $\Ga_0$-steps are normal]\label{prop:gamma0normal}
\PROVED{} \COMP\textsuperscript{a} Let $N\ge1$, $M\ge2$. Then $\Ga_0(MN)\trianglelefteq\Ga_0(N)$ \textbf{iff $M\mid\gcd(N,24)$}, and then $\gamma\mapsto a_\gamma c_\gamma/N\bmod M$ is an
isomorphism $\Ga_0(N)/\Ga_0(MN)\cong\Z/M$. \COMP\textsuperscript{a} evidence: \file{docs/implementation/slot22/les0/normality\_check.py} (all $150$ pairs $N\le24$, $MN\le48$, by
brute force in $\SL_2(\Z/MN)$) and \file{les\_predict.py} (all $1380$ pairs $N\le60$, $2\le M\le24$, by conjugating generators), $0$ mismatches.
\end{proposition}

\begin{proof}
$\Ga_0(MN)$ is normal iff it acts trivially, by right multiplication, on $\Ga_0(MN)\backslash\Ga_0(N)$. By Definition~\ref{def:levelellmodel} (bottom rows), this set is
$\mathcal F:=\{(c:d)\in\mathbb P^{1}(\Z/MN):N\mid c\}$ (every such row lifts to $\Ga_0(N)$, as $\SL_2(\Z)\to\SL_2(\Z/MN)$ is onto), with the right action of $\Ga_0(N)$ on rows;
$\Ga_0(MN)$ acts through its image, the upper triangular group of $\SL_2(\Z/MN)$, generated by $T$ and $D_u:=\mathrm{diag}(u^{-1},u)$, $u\in(\Z/MN)^{\times}$. By the Chinese
remainder theorem $\mathcal F=\prod_q\mathcal F_q$ over the primes $q\mid MN$, the generators act componentwise and the $D_u$ realise every tuple of local units. Fix $q$,
$r:=v_q(N)$, $s:=v_q(M)$. If $s=0$, $\mathcal F_q$ is the single point $(0:1)$. If $s\ge1$, $r=0$: $\mathcal F_q=\mathbb P^{1}(\Z/q^{s})$ and $T$ moves $(1:0)$ to $(1:1)$. If
$s,r\ge1$: $\mathcal F_q=\{(q^{r}x:1):x\in\Z/q^{s}\}$, $T$ sends $x$ to $x(1+q^{r}x)^{-1}$, i.e.\ moves it by $q^{r}x^{2}(1+q^{r}x)^{-1}$ in $\Z/q^{s}$ --- trivial for all $x$ iff
$s\le r$ --- and $D_u$ sends $x$ to $u^{-2}x$, trivial for all $x$ iff $u^{2}\equiv1\pmod{q^{s}}$ for every unit $u$, i.e.\ iff $q^{s}\in\{2,3,4,8\}$. So normality holds iff
$q^{s}\mid q^{r}$ and $q^{s}\mid24$ for every $q\mid M$, i.e.\ iff $M\mid\gcd(N,24)$. The map: for $\gamma,\gamma'\in\Ga_0(N)$, $(\gamma\gamma')_{21}/N=(c/N)a'+d(c'/N)$ and
$(\gamma\gamma')_{11}\equiv aa'\pmod N$, so $\phi(\gamma\gamma')\equiv aa'^{2}(c/N)+ada'(c'/N)\equiv\phi(\gamma)+\phi(\gamma')\pmod M$, using $ad\equiv1$ and $a'^{2}\equiv1$ modulo
$M\mid24$; $\phi\smat10N1=1$, and $\phi(\gamma)=0$ iff $MN\mid c$ ($a$ is a unit modulo $M$).
\end{proof}

\begin{proposition}[The steps of the $p$-tower]\label{prop:hidanormalsteps}
\PROVED{} \COMP\textsuperscript{a} Let $p\ge5$ be prime, $p\nmid N$, $r\ge1$, all groups intersected with $\Ga_0(N)$, $\mu:=\mu_{p-1}\subset(\Z/p^{k})^{\times}$ the Teichm\"uller subgroup,
$\Ga^{t}_r:=\Ga_{\mu}(p^{r})$ ([U] part~5's tame tower) and $\mathsf H_{r+1,r}:=\{u\in(\Z/p^{r+1})^{\times}:u\bmod p^{r}\in\mu\}$.
\begin{center}\small
\begin{tabular}{@{}llll@{}}
\toprule
step & index & normal & quotient \\
\midrule
$\Ga_0(pN)\supset\Ga_0(N)\cap\Ga_1(p)$ & $p-1$ & yes & $(\Z/p)^{\times}$, cyclic of order $p-1$ \\
$\Ga_1(p^{r})\supset\Ga_1(p^{r})\cap\Ga_0(p^{r+1})$ & $p$ & yes & $\Z/p$, via $c/p^{r}\bmod p$ \\
$\Ga_1(p^{r})\cap\Ga_0(p^{r+1})\supset\Ga_1(p^{r+1})$ & $p$ & yes & $\Z/p$, via $d\bmod p^{r+1}$ ($d\equiv1\bmod p^{r}$) \\
$\Ga_1(p^{r})\supset\Ga_1(p^{r+1})$ & $p^{2}$ & \textbf{no} & --- \\
$\Ga_0(p^{r})\supset\Ga^{t}_r$ & $p^{r-1}$ & yes & $(1+p\Z)/(1+p^{r}\Z)\cong\Z/p^{r-1}$, filtered by $\Z/p$ steps \\
$\Ga^{t}_r\supset\Ga^{t}_r\cap\Ga_0(p^{r+1})=\Ga_{\mathsf H_{r+1,r}}(p^{r+1})$ & $p$ & \textbf{no} & --- \\
$\Ga_{\mathsf H_{r+1,r}}(p^{r+1})\supset\Ga^{t}_{r+1}$ & $p$ & yes & $\mathsf H_{r+1,r}/\mu\cong\Z/p$ \\
$\Ga_0(p^{r})\supset\Ga_0(p^{r+1})$ & $p$ & \textbf{no} & --- \\
\bottomrule
\end{tabular}
\end{center}
More generally, for $M\mid N'$: $\Ga_1(N')\cap\Ga_0(MN')\trianglelefteq\Ga_1(N')$ with quotient $\Z/M$ ($c/N'\bmod M$), $\Ga_1(MN')\trianglelefteq\Ga_1(N')\cap\Ga_0(MN')$ with quotient
$\{d\equiv1\bmod N'\}\cong\Z/M$, and $\Ga_1(MN')\not\trianglelefteq\Ga_1(N')$ for $M>1$ prime to $N$. So raising the $d$-condition (the diamonds, the $\Lambda$-direction) is normal with
quotient $\Z/p$, raising the $c$-condition ($\Ga_0$, $\Ga^{t}$) is not, for $p\ge5$; $\Ga^{t}_1=\Ga_0(p)$, and the first step has quotient $(\Z/p)^{\times}$, not $\Z/p$.
\COMP\textsuperscript{a} evidence: \file{les\_predict.py}, $p\in\{5,7\}$, $r\le3$, $N\in\{1,2,3\}$, and the $\Ga_1$ steps for all $M\mid N'\le24$; $0$ mismatches.
\end{proposition}

\begin{proof}
Row 1: $\Ga_0(N)\cap\Ga_1(p)=\varsigma_{pN}^{-1}(\ker((\Z/pN)^{\times}\to(\Z/p)^{\times}))$, normal with that quotient (Proposition~\ref{prop:nebgammaH}(1)). The general $\Ga_1$
statements: for $\gamma,\gamma'\in\Ga_1(N')$, $(\gamma\gamma')_{21}/N'=(c/N')a'+d(c'/N')\equiv c/N'+c'/N'\pmod M$ ($a'\equiv d\equiv1\bmod N'$, $M\mid N'$), a homomorphism onto
$\Z/M$ ($\smat10{N'}1\mapsto1$) with kernel $\Ga_1(N')\cap\Ga_0(MN')$; $\varsigma_{MN'}$ restricted to $\Ga_1(N')\cap\Ga_0(MN')$ is a homomorphism onto
$\{d\equiv1\bmod N'\}=\{1+N'k\}$, $(1+N'k)(1+N'k')\equiv1+N'(k+k')\pmod{MN'}$, with kernel $\Ga_1(MN')$; and with $x:=NN'$ and $\sigma=\smat10{x}1\in\Ga_1(N')\cap\Ga_0(N)$ (inside $\Ga_0(N)$, as the proposition requires),
$\sigma T\sigma^{-1}=\smat{1-x}{1}{-x^{2}}{1+x}\notin\Ga_1(MN')$ ($1+NN'\not\equiv1\bmod MN'$, as $\gcd(N,M)=1$ and $M>1$; in particular for $M=p\nmid N$) while
$T\in\Ga_1(MN')$. Rows 2--4 are the case $N'=p^{r}$, $M=p$. Row 5:
Proposition~\ref{prop:nebgammaH}(1) with $(\Z/p^{r})^{\times}=\mu\times(1+p\Z/p^{r})$, the second factor cyclic of order $p^{r-1}$ ($p$ odd), filtered by
$\Ga_{\mu(1+p^{s})}(p^{r})$, $s=r,\dots,1$. Row 6: $\gamma\in\Ga_0(p^{r+1})$ has $d\bmod p^{r}\in\mu$ iff $d\bmod p^{r+1}\in\mathsf H_{r+1,r}$. The cosets are $x\in\Z/p$ (bottom rows
$(p^{r}x:1)$), and $\gamma\in\Ga^{t}_r$ acts by $x\mapsto a^{2}x+ac/p^{r}$ modulo $p$ (the local computation of Proposition~\ref{prop:gamma0normal}'s proof with $s=1\le r$); the
stabiliser of $x=0$ contains a lift of $\mathrm{diag}(\omega(g)^{-1},\omega(g))$, $g$ a primitive root modulo $p$, which sends $x=1$ to $g^{-2}\ne1$ ($p\ge5$). Row 7: both groups
are $\varsigma_{p^{r+1}}^{-1}$ of subgroups of $(\Z/p^{r+1})^{\times}$, hence normal in $\Ga_0(p^{r+1})$. Row 8: Proposition~\ref{prop:gamma0normal} with $M=p\nmid24$.
\end{proof}

\begin{remark}[Level one]\label{rmk:normalN1}
\PROVED{} (Given Claim~\ref{cl:amalgam}(1).) $\Ga^{\mathrm{ab}}\cong\Z/12$: the abelianisation of $\langle S\rangle\ast_{\langle-I\rangle}\langle U\rangle$ is
$(\Z/4\oplus\Z/6)/\langle(2,-3)\rangle\cong\Z/12$, $S\mapsto3$, $U\mapsto2$. So a normal subgroup of $\Ga$ with quotient $\Z/p$ exists only for $p\in\{2,3\}$, and is unique:
$\Ga^{2}$, the preimage of $A_3\subset\SL_2(\F_2)\cong S_3$ ($S\mapsto1$, $U\mapsto0$: it contains $-I$, $U$ and every element of order $3$, not $S$), and $\Ga^{3}$, the preimage of
the quaternion group $Q_8\trianglelefteq\SL_2(\F_3)$ ($S\mapsto0$, $U\mapsto2$: it contains $-I$ and $S$, no element of order $3$ or $6$). At $\ell=3$, $\Ga^{3}$ satisfies
Proposition~\ref{prop:modeltorsion} while $\Ga$ does not: the cell where one end of \eqref{eq:levelsixterm} is group cohomology and the other Bredon; at $\ell=2$ both are Bredon.
Further normal pairs at level one: $\Ga^{4},\Ga^{6},\Ga^{12}=[\Ga,\Ga]$ (not built in slot~22) and $\Ga(M)\trianglelefteq\Ga$, $\mathsf G=\SL_2(\Z/M)$, torsion-free for $M\ge3$
(Proposition~\ref{prop:modeltorsion}). At $M=\ell\ge5$, $\Ga(\ell)$ acts trivially on $\Sym^{n}\F_\ell$, so $H^{0}(B)=\VR$ and
$\ker\res^{1}\cong H^{1}(\SL_2(\F_\ell);\Sym^{n}\F_\ell)$ (Proposition~\ref{prop:levelinfres}(3)) --- computable from the finite group alone, the oracle O3 of \ref{LES8}.
\end{remark}

\begin{proof}[Proof of the identifications]
The pushout of $\Z/4\leftarrow\Z/2\to\Z/6$ has order $12$ and is generated by the image of $(1,1)$; a surjection $\Ga\to\Z/p$ factors through $\Z/12$. The reductions
$\Ga\to\SL_2(\F_2)\to\{\pm1\}$ (sign; $S\bmod2$ is a transposition) and $\Ga\to\SL_2(\F_3)\to\SL_2(\F_3)/Q_8\cong\Z/3$ ($U\bmod3$ has order $6$, so $U\notin Q_8$, whose
elements have order dividing $4$) are onto, so their kernels are $\Ga^{2}$, $\Ga^{3}$ by uniqueness. Membership is read off $S\mapsto3$, $U\mapsto2$: an element of order $3$ is
conjugate to $U^{\pm2}\mapsto\pm4$, of order $4$ to $S^{\pm1}\mapsto\pm3$, of order $6$ to $U^{\pm1}\mapsto\pm2$.
\end{proof}

\begin{obligation}[Level-cone certificates, re-derived over $\cring$]\label{obl:levelconeexact}
Data: the fibre matrices $M_A(\gamma)$, $M_B(\gamma)$ ($\gamma\in\{S,T,-I\}$, \eqref{eq:MHchi}), $\alpha_{B/A}$, $\omega_{B/A}$ of \eqref{eq:levelres} and, in the normal case,
$R_\sigma$ for a generator set of $\mathsf G$; the stored bases of $\Ceq$, $Z^{q}$, $B^{q}$, $H^{q}$ of both models and of $\Cone{\alpha_{B/A}}$ (Route~A, the model) or the tree
complexes \eqref{eq:treemodel} (Route~B, $6\in\cring^{\times}$); the operators of Proposition~\ref{prop:levelconehecke}(1). Ring: $\cring$ ($\F_\ell$ by ranks; $\Z/\ell^{K}$ by
Howell lengths). Each failure \textbf{aborts} unless marked.
\begin{enumerate}
\item[(i)] \textbf{Fibre.} $\alpha M_A(\gamma)=M_B(\gamma)\alpha$, $M_A(\gamma)\omega=\omega M_B(\gamma)$, $\omega\alpha=[A:B]$; normal case: $\alpha\omega=\mathrm{Nm}_{\mathsf G}$,
  $R_\sigma M_B(\gamma)=M_B(\gamma)R_\sigma$, $R_\sigma^{\operatorname{ord}\sigma}=1$; at $\chi=1$ on $\Ga_0$-pairs, $\alpha$ equals \code{LevelChangeMap} entry for entry.
\item[(ii)] \textbf{Cone.} $\Dcone^{2}=0$ and $d\alpha_{\ast}=\alpha_{\ast}d$ on the stored equivariant bases.
\item[(iii)] \textbf{Exactness} at the six terms of \eqref{eq:levelsixterm} (image $=$ kernel, injectivity of $\res^{0}$, surjectivity onto $H^{2}(A,B)$), $H^{0}(A,B)=H^{3}(A,B)=0$,
  and the connecting map re-derived by the zig-zag equal to $+\res^{1}$.
\item[(iv)] \textbf{Route~A $=$ Route~B} ($\ell\ge5$): the six dimensions and $\rank\res^{1}$ at every overlapping cell.
\item[(v)] \textbf{Hecke.} $\Phi_g$ of Proposition~\ref{prop:levelconehecke}(1) commutes with $\Dcone$ \emph{as matrices} on the stored equivariant basis, for $T_p$
  ($p\nmid N_B$) and for same-level $U_p$, and the three map families intertwine it.
\item[(vi)] \textbf{Level-raising $U_p$} (Proposition~\ref{prop:levelconehecke}(3)): \textbf{gate} the strict identity $U^{B}_p\circ(\alpha_{B/A})_{\ast}=p\cdot\mathrm{dg}_g$
  \emph{as matrices} on the stored equivariant bases ($0$ failing entries; $U^{B}_p\circ(\alpha_{B/A})_{\ast}=0$ at $\ell=p$); the square $U^{B}_p\res^{1}$ versus
  $\res^{1}U^{A}_p$ on $H^{1}$ is computed and \textbf{recorded} (its defect, Proposition~\ref{prop:levelconehecke}(3)(c)), not gated (rule~\ref{A6}: no code path branches
  on it).
\item[(vii)] \textbf{Normal case.} $\im\res^{q}\subseteq H^{q}(B)^{\mathsf G}$; $\dim\ker\res^{1}=\dim H^{1}(\mathsf G;H^{0}(B))$;
  $\dim H^{1}(B)^{\mathsf G}-\rank\res^{1}=\dim H^{2}(\mathsf G;H^{0}(B))$; for $\mathsf G\cong\Z/p$ the Jordan type of $\sigma$ on $H^{1}(B)$ and the equalities
  $\dim H^{a}(\mathsf G;H^{1}(B))=\dim H^{a+2}(\mathsf G;H^{0}(B))$, $a\in\{1,2\}$ (Proposition~\ref{prop:cycliccone}); $\cor^{1}$ an isomorphism on $H^{1}(B)_{\mathsf G}$.
\item[(viii)] \textbf{Model tags.} Every term at $\ell\in\{2,3\}$ carries \code{model=group} or \code{model=bredon} by Proposition~\ref{prop:modeltorsion}; no
  $\ell\in\{2,3\}$ dimension is compared with a group-cohomology formula.
\item[(ix)] \textbf{Negative controls}, each caught: $\chi$ replaced by $\chi^{-1}$ in $\alpha$ (unless $\chi$ is real), $\rhon(h_i)$ replaced by $I$, the sign of $\Dcone$
  flipped, an incompatible $\chi'$ (refused, \code{les\_incompatible\_character}), $R_\sigma$ for $\sigma$ not normalising $B$ (refused, \code{les\_not\_normal}), and the
  \textbf{criterion-fails control}: at a pair where the criterion of Proposition~\ref{prop:levelconehecke}(1) fails (level-raising, (3)), $U^{B}_p(\alpha_{B/A})_{\ast}\ne(\alpha_{B/A})_{\ast}U^{A}_p$ as
  matrices on the stored equivariant bases, $q=0,1,2$, at $(\Ga_0(25)\subset\Ga_0(5),\ \ell=7,\ n=0,\ \chi=1)$ --- caught. (The former control ``operators on independently fixed transversals: the strict square fails
  at some cell'' is \textbf{retired}: by Proposition~\ref{prop:levelconehecke}(4) it can never fire. Operators built on two transversals are \textbf{gated} equal on $\Ceq$ ($0$ differing entries, together with the strict square on $A$'s own transversal) at a cell with
  $m_A>1$, e.g.\ $(\Ga_0(25)\subset\Ga_0(5),\ T_2,\ \ell=7,\ n=2)$; the comparison on the flat cochains is only \emph{recorded} ($0$, by construction of \eqref{eq:levelcorr}); neither is a negative control, since the identity is a theorem.)
\end{enumerate}
\end{obligation}

\begin{remark}[What is checkable: \ref{LES1}--\ref{LES9}]\label{rmk:lescheckable}
Committed before any LES code (rule~20) in \file{docs/implementation/slot22/LES\_predictions.md} with the numbers of
\file{docs/implementation/slot22/les0/out\_predict.txt} (\file{les\_predict.py}, \file{les\_model.py}: the tree complex of Proposition~\ref{prop:bredon}(2) in exact arithmetic,
and the oracle O3, an independent crossed-homomorphism computation on $\SL_2(\F_\ell)$); convention of Remark~\ref{rmk:hidacheckable}.
\begin{enumerate}[label=(LES\arabic*)]
\item\label{LES1} \PROVED{} (Propositions~\ref{prop:gamma0normal}, \ref{prop:hidanormalsteps}, Remark~\ref{rmk:normalN1}) \textbf{Normality}: the table of
  Proposition~\ref{prop:hidanormalsteps}, $\Ga_0(MN)\trianglelefteq\Ga_0(N)$ iff $M\mid\gcd(N,24)$ on $N\le60$, $M\le24$, the content of $\Ga^{2}$, $\Ga^{3}$, and the fibre
  identities of Obligation~\ref{obl:levelconeexact}(i) at every cell below.
\item\label{LES2} \PROVED{} (Lemma~\ref{lem:levelconeexact}) \textbf{Exactness} of \eqref{eq:levelsixterm} at every cell, every $\ell$.
\item\label{LES3} \PROVED{} (Proposition~\ref{prop:levelorbit}(3), Proposition~\ref{prop:neborbit}(4), Proposition~\ref{prop:levelsplit}) \textbf{Dimensions} at $\ell\ge5$:
  the ends by the closed forms, the split cells by $\dim H^{1}(A,B)=\dim H^{0}(B)-\dim H^{0}(A)$, $\dim H^{2}(A,B)=\dim H^{1}(B)-\dim H^{1}(A)$; the non-split cells
  \COMP\textsuperscript{a}.
\item\label{LES4} \PROVED{} (Proposition~\ref{prop:levelconehecke}(1),(2)) \textbf{Strict Hecke squares}: $0$ failing entries.
\item\label{LES5} \PROVED{} (Proposition~\ref{prop:bredon}(2)) \textbf{Route~A $=$ Route~B} at every overlapping cell.
\item\label{LES6} \PROVED{} (Lemma~\ref{lem:levelconeexact}, Proposition~\ref{prop:modeltorsion}) \textbf{$\ell\in\{2,3\}$}: exactness and the \code{model=} tags; the model
  dimensions \COMP\textsuperscript{a} (tree complex), never compared with O1.
\item\label{LES7} \PROVED{} (Proposition~\ref{prop:cycliccone}(3)(c), Lemmas~\ref{lem:levelfreeeuler}, \ref{lem:unipotentinv}) \textbf{The $\Ga_1$-type $\Z/p$ steps} at $\ell=p$:
  $c=\lfloor n/p\rfloor+1$, $\dim H^{1}(B)=p\dim H^{1}(A)-(p-1)c$, Jordan type $J_p^{\,h-c}\oplus J_1^{\,c}$, and the cone dimensions; \COMP\textsuperscript{a} at the
  smaller cells.
\item\label{LES8} \PROVED{} (Proposition~\ref{prop:cycliccone}(3)(b),(d), Lemma~\ref{lem:unipotentinv}, Remark~\ref{rmk:normalN1}) \textbf{The tower cells}
  $\Ga^{t}_2\trianglelefteq\Ga_0(p^{2})$, $n=0$, $\chi=\omega^{a-2}$: $H^{1}(\Ga^{t}_2)$ free over $\F_p[\Z/p]$ for $a\not\equiv2$, exactly one non-free block at $a\equiv2$ (its size
  \COMP\textsuperscript{a}); and $\Ga(\ell)\trianglelefteq\Ga$: $\dim\ker\res^{1}=\dim H^{1}(\SL_2(\F_\ell);\Sym^{n}\F_\ell)$ (O3, \COMP\textsuperscript{a}).
\item\label{LES9} \PROVED{} (Proposition~\ref{prop:levelconehecke}(2),(3)) \textbf{Hecke squares on $H^{1}$} by Route~B: for $T_p$ and same-level $U_p$ they commute;
  for level-raising $U_p$ the strict identity $U^{B}_p\circ\res=p\cdot\mathrm{dg}_g$ holds with $0$ failing entries ($=0$ at $\ell=p$), and the square's defect is
  \textbf{non-zero} at $(\Ga_0(25)\subset\Ga_0(5),\ \ell=7,\ n=0,\ \chi=1)$ (\COMP\textsuperscript{a} \file{les0/out\_levelraise\_up.txt}, also at $\ell\in\{5,11\}$,
  $(\Ga_0(49)\subset\Ga_0(7),\ \ell=5)$, $(\Ga_0(121)\subset\Ga_0(11),\ \ell=7)$); other level-raising cells recorded.
\end{enumerate}
\end{remark}

\begin{remark}[The old-space cone, deferred]\label{rmk:oldspacecone}
(Not a statement.) For $p\mid N_B$, $p\nmid N_A$ the restriction does not intertwine $T^{A}_p$ and $U^{B}_p$ (Proposition~\ref{prop:oldspace}(1),(3)); the $U_p$-equivariant
object is the cone of $\Phi=(\mathrm{dg}_1,\mathrm{dg}_p)\colon\Ceq(L_A)^{\oplus2}\to\Ceq(L_B)$ with $U_p$ acting on the source by $\mathsf M$ of Proposition~\ref{prop:oldspace}(4),
whose $H^{1}$-cokernel is the $p$-new quotient. Since Proposition~\ref{prop:oldspace}(3) holds only up to homotopy, its cone action needs an $\etapair$-type certificate;
deferred (slot-22 LES-0, owner question Q4). The alternative is the classical $U_p=[A\,\mathrm{diag}(p,1)A]$ of Proposition~\ref{prop:levelconehecke}(3)(d), whose
square with $\res$ is strict by Proposition~\ref{prop:levelconehecke}(1); this document's $U^{(N)}_p$ is its adjoint $U_p^{\ast}$ on old spaces.
\end{remark}

\paragraph{Code map of \S\ref{sec:nebles}} \textbf{Planned} (slot-22 LES-1..3, after the referee F-LES): \code{CongruenceLevel(M,K)} (for $\Ga(M)$, $\Ga^{2}$, $\Ga^{3}$, only
for levels not of $\Ga_{\mathsf H}$ type), \code{RestrictionFibreMap}, \code{ConjugationFibreMap} and the trace (Definition~\ref{def:levelres}, Lemma~\ref{lem:levelresmorph}) in
\file{Hecke/}; the cone on the model (an audit of \file{EquivariantMappingCone}, or a sibling \code{LevelPairCone}); the tree-complex route over $\F_\ell$ and $\Z/\ell^{K}$
reusing \code{OnCokernelLevelChange} and Lemma~\ref{lem:nebchainrref}; validations \code{HeckeLevelNormalityValidation} (\code{fast}, \ref{LES1}),
\code{HeckeLevelConeValidation} (\code{standard}, \ref{LES2}--\ref{LES6}, \ref{LES9}) and \code{HeckeLevelInfResValidation} (\code{standard}, \ref{LES7}, \ref{LES8}).

%======================================================================
\part*{Part IX --- Integral and local structure}
\addcontentsline{toc}{section}{Part IX --- Integral and local structure}
%======================================================================

%----------------------------------------------------------------------
\section{The \texorpdfstring{$\Z/\ell^{K}$}{Z/l^K} Hecke algebra: saturation, the defect, and the comparison with modular forms}\label{sec:saturation}
%----------------------------------------------------------------------
% NEW at v1.15.0-pre.1 (hecke-extensions slot 15, WP-3 Part B; plan docs/implementation/slot15/00-slot15-plan.md
% S2 "WP-3" Part B, owner rulings Q-C / Q-E of slot 14 and Q4, Q15 of 2026-09-30).  The proposals P15 (= P-WP6-1,
% P-WP6-2 of docs/implementation/slot14/WP-6_report.md S5; [M1] = docs/implementation/slot14/WP-M1_loewy_note.md,
% Prop. 1), P17 and P18 (docs/implementation/slot14/WP-D1_binding_doc_proposals.md), P20 (ibid., addendum; source
% docs/implementation/slot14/WP-C9_weight_shift_report.md S3.3), and the torsion-aware comparison of plan WP-19 (c)
% as corrected in c1b45cd.  Placed after sec:nebentypus, before the appendices.

\paragraph{What this chapter fixes.}
(i) \textbf{Saturation} (Definition~\ref{def:satdata}, Proposition~\ref{prop:saturation}, P15): for a commutative order $\Talg$
acting faithfully on a free $\Z_\ell$-lattice, the image $\mathsf A_K$ of $\Talg$ on the lattice modulo $\ell^{K}$ is
$\Talg/(\Talg\cap\ell^{K}\Talg^{\mathrm{sat}})$; reducing the \emph{module} first loses exactly $(\Talg\cap\ell\Talg^{\mathrm{sat}})/\ell\Talg$;
and the stopping rule ``$\dim\mathsf A_K/\ell=\rank\Talg$'' is \emph{equivalent} to $\mathsf A_K/\ell\cong\Talg/\ell$ --- which makes the
\code{stable} verdict of \file{HeckeExt/HeckeRingHeckeAlgebra} a theorem, not a reading. (ii) \textbf{The defect}
(Proposition~\ref{prop:satdefect}, P17): $\varepsilon^{\mathrm{sat}}_K=\mathrm{length}\bigl((\Talg^{\mathrm{sat}}/\Talg)[\ell^{K}]\bigr)$. (iii)
\textbf{The level-$1$ inputs} (Lemma~\ref{lem:heckerankone}, Lemma~\ref{lem:freepart}): the rank $\lceil f/2\rceil$ that the stopping
rule compares with, and the free part of a torsion block. (iv) \textbf{The certified reading} (Corollary~\ref{cor:stablecert}),
with its generation caveat. (v) \textbf{The comparison of the cohomological and the forms Hecke algebra}, torsion-aware
(Proposition~\ref{prop:cohforms}). (vi) \textbf{P18}: the algebra-level constancy along a Hida class that
Theorem~\ref{thm:hidaweight}(1), a statement about ranks, does not give (Proposition~\ref{prop:hidaalgconst}). (vii) \textbf{P20}: the
forms-side $\theta$ and its Frobenius kernel, against Theorem~\ref{thm:thetales}(3) (Remark~\ref{rmk:thetaforms}). (viii) (v1.17.0-pre.1) \textbf{The defect as a
$\mathrm{Tor}$ term} (\S\ref{sec:sattorsec}, Remark~\ref{rmk:sattor}). (ix) (v1.23.0-pre.3) \textbf{One module run per cell}
(\S\ref{sec:satprecisionsec}, Lemma~\ref{lem:precisionreduction}, Obligation~\ref{obl:precisionreuse}): the levels $K$ the route reads are reductions of the
module run at the largest of them, which licenses one run per cell in place of one per level.

\paragraph{Notation.}
$\ell$ a prime; $\Z_\ell$-modules are finitely generated. For a finite $\Z_\ell$-module $G$, $G[\ell^{K}]$ is its $\ell^{K}$-torsion and
$\mathrm{length}(G)$ its length ($=\log_\ell\abs G$). $\Talg$ (\S\ref{sec:notation}'s symbol for a Hecke algebra) is decorated here:
$\Talg^{\mathrm{sat}}$, $\Talg^{\mathrm{coh}}$, $\Talg^{\mathbb L}$, $\Talg^{\mathrm f}$. $\mathbb L:=H^{1}(\Ga;\Vcoef\otimes\Z_\ell)$ and $\mathbb L':=\mathbb L/\mathbb L_{\mathrm{tors}}$;
$\mathsf A_K$, $\varepsilon^{\mathrm{sat}}_K$, $\delta^{\mathrm{sat}}$, $c^{\mathrm{sat}}$ are new symbols of this chapter and mean nothing else (rule~\ref{A5};
$\mathbb L$ is not Hida's $\Lambda$ nor $\Lam$, $\mathsf A_K$ not $\Amat{q}$, and $\varepsilon^{\mathrm{sat}}_K$ not the augmentation $\aug$).

\subsection{Saturation and the defect}\label{sec:satsec}

\begin{definition}[Saturation data]\label{def:satdata}
\LEAN{def:satdata}
Let $L$ be a free $\Z_\ell$-module of finite rank and $\Talg\subseteq\End_{\Z_\ell}(L)$ a commutative $\Z_\ell$-subalgebra (so $\Talg$ is
free of finite rank $r:=\rank_{\Z_\ell}\Talg$ and acts faithfully). Put
\[
  \Talg^{\mathrm{sat}}:=(\Talg\otimes\Q)\cap\End_{\Z_\ell}(L)\quad(\text{inside }\End_{\Q_\ell}(L\otimes\Q)),\qquad
  \mathsf A_K:=\mathrm{im}\bigl(\Talg\to\End(L/\ell^{K}L)\bigr)\ (K\ge1),
\]
$\delta^{\mathrm{sat}}:=\mathrm{length}(\Talg^{\mathrm{sat}}/\Talg)$, $c^{\mathrm{sat}}$ the least $c\ge0$ with $\ell^{c}\Talg^{\mathrm{sat}}\subseteq\Talg$, and
$\varepsilon^{\mathrm{sat}}_K:=K\,r-\mathrm{length}(\mathsf A_K)$. $\Talg$ is \emph{saturated} if $\Talg=\Talg^{\mathrm{sat}}$.
\end{definition}

\noindent
$\Talg^{\mathrm{sat}}$ is a commutative $\Z_\ell$-algebra containing $\Talg$, free of rank $r$ (a saturated submodule of $\End(L)$ with
$\Talg^{\mathrm{sat}}\otimes\Q=\Talg\otimes\Q$), so $\Talg^{\mathrm{sat}}/\Talg$ is finite, $\delta^{\mathrm{sat}}$ and $c^{\mathrm{sat}}$ are finite, and $c^{\mathrm{sat}}$ is the
exponent of $\Talg^{\mathrm{sat}}/\Talg$. $L$ is a $\Talg^{\mathrm{sat}}$-module.

\begin{proposition}[Saturation: what reducing the module loses (P15)]\label{prop:saturation}
\PROVED{} In the notation of Definition~\ref{def:satdata}:
\begin{enumerate}
\item $\ker\bigl(\Talg\to\End(L/\ell^{K})\bigr)=\Talg\cap\ell^{K}\End(L)=\Talg\cap\ell^{K}\Talg^{\mathrm{sat}}$, so
  $\mathsf A_K=\Talg/(\Talg\cap\ell^{K}\Talg^{\mathrm{sat}})$; in particular $\ker(\Talg/\ell\Talg\to\mathsf A_1)=(\Talg\cap\ell\Talg^{\mathrm{sat}})/\ell\Talg$. \LEAN{prop:saturation}
\item The following are equivalent: (a) $\Talg/\ell\Talg\to\mathsf A_1$ is an isomorphism; (b) $L/\ell L$ is a faithful $\Talg/\ell\Talg$-module;
  (c) $\Talg$ is saturated. They hold if $L$ contains a $\Talg$-submodule $N\cong\Talg$ (free of rank one) with $L/N$ torsion-free.
  \LEAN{prop:saturation} ((a) as the injectivity of the canonical surjection)
\item \textbf{(The stopping rule.)} For every $K\ge1$, $\dim_{\Fl}\mathsf A_K/\ell\mathsf A_K\le r$, and the following are equivalent:
  (a) $\dim_{\Fl}\mathsf A_K/\ell\mathsf A_K=r$; (b) the canonical surjection $\Talg/\ell\Talg\to\mathsf A_K/\ell\mathsf A_K$ is an isomorphism (of
  $\Fl$-algebras); (c) $\Talg\cap\ell^{K}\Talg^{\mathrm{sat}}\subseteq\ell\Talg$. If they hold at $K$, they hold at every $K'\ge K$; they hold
  for every $K\ge c^{\mathrm{sat}}+1$. \LEAN{prop:saturation} ((b) as the injectivity of the canonical surjection)
\item For $K\ge c^{\mathrm{sat}}$: $\Talg\cap\ell^{K}\Talg^{\mathrm{sat}}=\ell^{K}\Talg^{\mathrm{sat}}$ and $\mathrm{length}(\mathsf A_K)=K\,r-\delta^{\mathrm{sat}}$. \LEAN{prop:saturation}
\item Every maximal ideal $\midl$ of $\Talg$ contains $\ell$ \LEAN{prop:saturation} (this sentence only); in (3)(b) the isomorphism carries the $\midl$-adic filtration to the
  $\midl$-adic filtration, so the Hilbert function and the Loewy length of $\mathsf A_K/\ell\mathsf A_K$ at $\midl$ are those of
  $\Talg/\ell\Talg$. All statements localise: for $\Talg=\prod_\midl\Talg_\midl$, $L=\bigoplus_\midl L_\midl$, they hold for each pair
  $(\Talg_\midl,L_\midl)$, with $(\Talg^{\mathrm{sat}})_\midl=(\Talg_\midl)^{\mathrm{sat}}$.
\end{enumerate}
\end{proposition}

\begin{proof}
(1) $t\in\Talg$ acts as $0$ on $L/\ell^{K}L$ iff $tL\subseteq\ell^{K}L$ iff $t\in\ell^{K}\End(L)$ ($L$ free: the matrix entries). If
$t=\ell^{K}e$ with $e\in\End(L)$, then $e=\ell^{-K}t\in\Talg\otimes\Q$, so $e\in\Talg^{\mathrm{sat}}$; the converse inclusion holds as
$\Talg^{\mathrm{sat}}\subseteq\End(L)$.
(2) (a)$\Leftrightarrow$(b): both say $\Talg\cap\ell\End(L)=\ell\Talg$, by (1). (c)$\Rightarrow$(a): $\Talg\cap\ell\Talg=\ell\Talg$.
(a)$\Rightarrow$(c): let $x\in\Talg^{\mathrm{sat}}\setminus\Talg$ and $j\ge1$ least with $\ell^{j}x\in\Talg$ (exists, $\Talg^{\mathrm{sat}}/\Talg$ being
finite); then $\ell^{j}x\in\Talg\cap\ell\Talg^{\mathrm{sat}}=\ell\Talg$ by (a) and (1), so $\ell^{j-1}x\in\Talg$ ($\Talg$ is torsion-free), against
minimality. The sufficient condition: let $N=\Talg v$ with $t\mapsto tv$ injective, and $x\in\Talg^{\mathrm{sat}}$. Then $xv\in L$ and
$xv\in(\Talg\otimes\Q)v=N\otimes\Q$; as $L/N$ is torsion-free, $L\cap(N\otimes\Q)=N$, so $xv=sv$ with $s\in\Talg$. Then $(x-s)v=0$ with
$x-s\in\Talg\otimes\Q$, and $t\mapsto tv$ is injective on $\Talg\otimes\Q$ ($N\otimes\Q\cong\Talg\otimes\Q$), so $x=s\in\Talg$.
(3) By (1), $\mathsf A_K/\ell\mathsf A_K=\Talg/(\ell\Talg+\Talg\cap\ell^{K}\Talg^{\mathrm{sat}})$, a quotient of the $r$-dimensional $\Talg/\ell\Talg$; it has
dimension $r$ iff $\Talg\cap\ell^{K}\Talg^{\mathrm{sat}}\subseteq\ell\Talg$. The map is induced by a ring map, hence is one of algebras.
Monotonicity: $\Talg\cap\ell^{K'}\Talg^{\mathrm{sat}}\subseteq\Talg\cap\ell^{K}\Talg^{\mathrm{sat}}$ for $K'\ge K$. For $K\ge c^{\mathrm{sat}}+1$:
$\ell^{K}\Talg^{\mathrm{sat}}=\ell^{K-c^{\mathrm{sat}}}\ell^{c^{\mathrm{sat}}}\Talg^{\mathrm{sat}}\subseteq\ell^{K-c^{\mathrm{sat}}}\Talg\subseteq\ell\Talg$.
(4) $\ell^{K}\Talg^{\mathrm{sat}}\subseteq\ell^{K-c^{\mathrm{sat}}}\Talg\subseteq\Talg$, so the intersection is $\ell^{K}\Talg^{\mathrm{sat}}$, and
$\mathrm{length}(\Talg/\ell^{K}\Talg^{\mathrm{sat}})=\mathrm{length}(\Talg^{\mathrm{sat}}/\ell^{K}\Talg^{\mathrm{sat}})-\mathrm{length}(\Talg^{\mathrm{sat}}/\Talg)=Kr-\delta^{\mathrm{sat}}$.
(5) $\Talg$ is finite over the local ring $\Z_\ell$, so $\ell$ lies in every maximal ideal. A surjective algebra map carries $\midl^{i}$
onto the $i$-th power of the image of $\midl$; an isomorphism preserves the filtration. $\Talg$ is semilocal and $\ell$-adically
complete, hence the product of its localisations, and $L$ decomposes by the corresponding idempotents; each idempotent lies in
$\Talg$, so saturation, intersections and images are computed factorwise.
\end{proof}

\begin{proposition}[The defect (P17)]\label{prop:satdefect}
\PROVED{} In the notation of Definition~\ref{def:satdata}, for every $K\ge1$:
\begin{enumerate}
\item $\varepsilon^{\mathrm{sat}}_K=\mathrm{length}\bigl((\Talg^{\mathrm{sat}}/\Talg)[\ell^{K}]\bigr)$; \LEAN{prop:satdefect}
\item hence $0\le\varepsilon^{\mathrm{sat}}_K\le\varepsilon^{\mathrm{sat}}_{K+1}\le\delta^{\mathrm{sat}}$, $\varepsilon^{\mathrm{sat}}_K=\delta^{\mathrm{sat}}$ for $K\ge c^{\mathrm{sat}}$, and
  $\varepsilon^{\mathrm{sat}}_K=\varepsilon^{\mathrm{sat}}_{K+1}$ implies $\varepsilon^{\mathrm{sat}}_K=\delta^{\mathrm{sat}}$ (a plateau is final); \LEAN{prop:satdefect}
\item $\delta^{\mathrm{sat}}=0$ iff $\Talg$ is saturated iff $\varepsilon^{\mathrm{sat}}_1=0$. \LEAN{prop:satdefect}
\end{enumerate}
The quantity $K\dim(\mathsf A_K/\ell)-\mathrm{length}(\mathsf A_K)$ equals $\varepsilon^{\mathrm{sat}}_K$ exactly when $\dim\mathsf A_K/\ell=r$, and not in
general.
\end{proposition}

\begin{proof}
(1) $\ell^{K}\Talg\subseteq\Talg\cap\ell^{K}\Talg^{\mathrm{sat}}$, so by Proposition~\ref{prop:saturation}(1)
$\mathrm{length}(\mathsf A_K)=\mathrm{length}(\Talg/\ell^{K}\Talg)-\mathrm{length}\bigl((\Talg\cap\ell^{K}\Talg^{\mathrm{sat}})/\ell^{K}\Talg\bigr)$, and $\Talg$ free of rank
$r$ gives $\mathrm{length}(\Talg/\ell^{K}\Talg)=Kr$. Multiplication by $\ell^{-K}$ inside $\Talg\otimes\Q$ is an isomorphism
$(\Talg\cap\ell^{K}\Talg^{\mathrm{sat}})/\ell^{K}\Talg\to(\ell^{-K}\Talg\cap\Talg^{\mathrm{sat}})/\Talg=\{x\in\Talg^{\mathrm{sat}}/\Talg:\ell^{K}x=0\}$.
(2) $G[\ell^{K}]\subseteq G[\ell^{K+1}]\subseteq G$ for $G=\Talg^{\mathrm{sat}}/\Talg$, and $G[\ell^{K}]=G$ for $K\ge c^{\mathrm{sat}}$, the exponent. If
$G[\ell^{K}]=G[\ell^{K+1}]$ and $x\in G$ had order $\ell^{j}$, $j>K$, then $\ell^{j-K-1}x$ would lie in $G[\ell^{K+1}]\setminus G[\ell^{K}]$.
(3) $G=0$ iff $G[\ell]=0$. The last sentence: $K\dim(\mathsf A_K/\ell)-\mathrm{length}(\mathsf A_K)=\varepsilon^{\mathrm{sat}}_K-K(r-\dim\mathsf A_K/\ell)$.
\end{proof}

\subsection{The level-\texorpdfstring{$1$}{1} inputs}\label{sec:satlevelone}

\paragraph{Setting.}
$\ell\ge5$, $n>0$ even, $k:=n+2$, $\mathbb L=H^{1}(\Ga;\Vcoef\otimes\Z_\ell)$, $\mathbb L_{\mathrm{tors}}$ its torsion, $\mathbb L'=\mathbb L/\mathbb L_{\mathrm{tors}}$.
$\Talg^{\mathrm{coh}}\subseteq\End(\mathbb L)$ is the $\Z_\ell$-algebra generated by all $T_p$, $p$ prime, with $T_\ell=\Tdc{g}$ of level~$1$
($g=\mathrm{diag}(1,\ell)$; the pipeline's $U_\ell$, Proposition~\ref{prop:oldspace}'s $T_\ell$); it preserves $\mathbb L_{\mathrm{tors}}$, and
$\Talg^{\mathbb L}\subseteq\End(\mathbb L')$ is its image. $\Talg^{\mathrm f}\subseteq\End(M_k(\Z_\ell))$ is the $\Z_\ell$-algebra generated by the
classical $T_p$ on the level-$1$ modular forms with $q$-expansion coefficients in $\Z_\ell$. For a maximal ideal $\midl$ of
$\Talg^{\mathrm{coh}}$ (it lies over one of $\Talg^{\mathbb L}$ when $\mathbb L'_\midl\ne0$), subscripts $\midl$ denote localisation; $f_\midl:=\rank\mathbb L'_\midl$.

\begin{lemma}[The rank of the free Hecke algebra of a block]\label{lem:heckerankone}
\PROVED{} (given the \textsf{[standard]} inputs named in the proof) \LEANlit{lem:heckerankone}{Eichler--Shimura at level one (\cite[Thm.~12.2.2, Prop.~12.4.10]{diamond_im}), as the hypothesis \texttt{hES}} (over a coefficient field $E\supseteq\Q_\ell$) In the Setting, let $e\in\Talg^{\mathbb L}$ be any idempotent (e.g.\ the
idempotent of a maximal ideal of $\Talg^{\mathbb L}$, or of a maximal ideal $\midl'$ of a subalgebra $\Talg'\subseteq\Talg^{\mathbb L}$, whose block
$e\mathbb L'$ may be a sum of several blocks of $\Talg^{\mathbb L}$). Then $\rank_{\Z_\ell}e\Talg^{\mathbb L}=\lceil\rank(e\mathbb L')/2\rceil$; in
particular $\rank_{\Z_\ell}\Talg^{\mathbb L}_\midl=\lceil f_\midl/2\rceil$ (v1.15.0-pre.2, WP-4 referee B1). And
$\Talg^{\mathbb L}\otimes\Q_\ell$ is a product of fields, one factor per $\mathrm{Gal}(\bar\Q_\ell/\Q_\ell)$-orbit of systems of Hecke eigenvalues
of normalised eigenforms in $M_k$.
\end{lemma}

\begin{proof}
Eichler--Shimura \STD{\cite[Thm.~12.2.2, p.~108]{diamond_im} (the isomorphism $\mathcal M_k(\Ga)\oplus\bar{\mathcal S}_k(\Ga)\to H^{1}(\Ga,V_k)$, for every congruence
subgroup~$\Ga$, $k\ge2$) with its Hecke equivariance \cite[Prop.~12.4.10, p.~117]{diamond_im} (compatibility with the action of $R(\Ga,\De)$), both
\textbf{verified} on page images of the library's copy (slot-17 L17-5 repairs, the Lean entry \texttt{DiamondImEichlerShimura}; v1.18.0-pre.1, the slot-17 referee
F1's m-C4); \cite[Ch.~8, Prop.~8.5, pp.~237--238]{shimura}, which Diamond--Im cite for the equivariance, kept as a secondary pointer, \textbf{verified} for the
cuspidal part ($S_k\to H^{1}_{P}$; R20-L, v1.21.0-pre.1);
the pipeline's $T_\ell$ is the classical $T_\ell$ by
Theorem~\ref{thm:classical}; Claim~\ref{cl:dimH1}(1) and Obligation~\ref{obl:oracle} state the same decomposition} gives
$\mathbb L'\otimes\Cx\cong M_k(\Cx)\oplus\overline{S_k(\Cx)}$ as modules over the polynomial ring in the $T_p$ (after fixing
$\bar\Q_\ell\cong\Cx$). $M_k(\Cx)$ has a basis of normalised eigenforms $E_k,f_1,\dots,f_s$ \STD{two inputs, cited separately (v1.17.0-pre.1, slot-17 L17-5's
finding, and v1.17.0-pre.2, the citation repair of the WP-3 brief's addendum: the earlier pointer to \cite[Ch.~VII, \S5]{serre_cours} alone does not give the
existence). \emph{The eigenbasis}: the $T_m$ commute and are self-adjoint for the Petersson inner product on $S_k(\Cx)$, so $S_k(\Cx)$ has a basis of
simultaneous eigenvectors, \cite[\S3.6, p.~55]{diamond_im} (the $T_n$, $(n,N)=1$, self-adjoint for the Petersson product, and $S_k(N,\varepsilon)$ an orthogonal
direct sum of $\mathbb T^{(N)}$-eigenspaces) with \cite[\S3.5, p.~54]{diamond_im} (every eigenspace at most one-dimensional); level~$1$ is the case $N=1$ there;
\textbf{verified} on the page images of the library's copy (slot-17 L17-5 audit and R17-4, v1.17.0-pre.5). Secondary pointers: \cite[Thm.~3.41, Thm.~3.42 and the paragraph after it, p.~77]{shimura} (\textbf{verified}, R20-L, v1.21.0-pre.1: at $N=1$ the
$[\Ga\alpha\Ga]_k$ on $S_k(\Ga)$ commute and are self-adjoint, hence simultaneously diagonalisable with real eigenvalues),
\cite[\S5.5, Thm.~5.5.4]{diamond_shurman} (\textbf{not verified}, rule~\ref{A2}: not in the library);
and $E_k$ spans the Eisenstein complement and is an eigenform. \emph{Multiplicity one}: a non-zero eigenform has $a_1\ne0$ and, normalised,
$a_m=\lambda_m$, \cite[Ch.~VII, \S5, Thm.~7, p.~102]{serre_cours}, with \cite[Ch.~VII, \S5, Cor.~1 of Thm.~7, p.~102; Prop.~11, p.~101]{serre_cours},
verified on page images (slot-17 L17-5); Serre's \S5 gives this distinctness half only, not the existence of the basis}
whose eigenvalue systems are pairwise distinct: for two normalised eigenforms with the same $T_m$-eigenvalues for all $m\ge1$
(the $T_m$ are polynomials in the $T_p$ by Theorem~\ref{thm:relations}), $a_m=\lambda_m$ for $m\ge1$, so their difference is a constant
modular form of weight $k\ge4$, hence $0$, and a cusp form cannot coincide with $E_k$, whose constant term is $\ne0$. So the $T_p$
act semisimply on $\mathbb L'\otimes\bar\Q_\ell$ with each cuspidal system twice and the Eisenstein system once, and
$\Talg^{\mathbb L}\otimes\bar\Q_\ell\cong\bar\Q_\ell^{\,s+1}$ (an algebra generated by commuting semisimple operators is the product over its
distinct systems). Localising at $\midl$ keeps the systems reducing to $\midl$: $f_\midl=2s_\midl+e_\midl$, $e_\midl\in\{0,1\}$ the
number of Eisenstein systems in $\midl$, and $\rank\Talg^{\mathbb L}_\midl=s_\midl+e_\midl=\lceil f_\midl/2\rceil$. For an idempotent $e$ the same
count reads $\rank(e\mathbb L')=2s+\epsilon$ with $\epsilon\in\{0,1\}$ the number of Eisenstein systems in the block --- at most one, $M_k$
having one Eisenstein series --- so $\rank\,e\Talg^{\mathbb L}=s+\epsilon=\lceil\rank(e\mathbb L')/2\rceil$ for every $e$.
\end{proof}

\begin{lemma}[The free part of a torsion block (P-WP6-2)]\label{lem:freepart}
\PROVED{} In the Setting, let $t\ge0$ with $\ell^{t}\mathbb L_{\mathrm{tors},\midl}=0$ and $K>t$. Put $N:=\mathbb L_\midl/\ell^{K}\mathbb L_\midl
=H^{1}(\Ga;\Vcoef\otimes\Z/\ell^{K})_\midl$ and $\mathsf A_K:=\mathrm{im}(\Talg^{\mathrm{coh}}_\midl\to\End N)$.
\begin{enumerate}
\item Multiplication by $\ell^{t}$ induces a Hecke-equivariant isomorphism $\mathbb L'_\midl/\ell^{K-t}\mathbb L'_\midl\cong\ell^{t}N$. \LEAN{lem:freepart}
\item The image of $\Talg^{\mathrm{coh}}_\midl$ in $\End(\ell^{t}N)$ is $\mathsf A_K/\mathsf A_K[\ell^{t}]$, and $a\mapsto\ell^{t}a$ is an isomorphism of
  $\Z_\ell$-modules $\mathsf A_K/\mathsf A_K[\ell^{t}]\to\ell^{t}\mathsf A_K\subseteq\End N$; by (1) this image is the image $\mathsf A'_{K-t}$ of
  $\Talg^{\mathbb L}_\midl$ in $\End(\mathbb L'_\midl/\ell^{K-t})$, to which Proposition~\ref{prop:saturation} applies with $(L,\Talg)=(\mathbb L'_\midl,
  \Talg^{\mathbb L}_\midl)$ at level $K-t$. \LEAN{lem:freepart} (in full since v1.18.0)
\item If $\mathbb L'_\midl=0$ (a pure torsion block) and $K\ge t$, then $N=\mathbb L_\midl$ and $\mathsf A_K=\Talg^{\mathrm{coh}}_\midl$. \LEAN{lem:freepart}
\end{enumerate}
The identification $H^{1}(\Ga;\Vcoef\otimes\Z/\ell^{K})=\mathbb L/\ell^{K}$ is Claim~\ref{cl:bockstein}'s argument with $\ell^{K}$ in place of
$\ell$ ($H^{2}(\Ga;\Vcoef)$ has no $\ell$-torsion).
\end{lemma}

\begin{proof}
(1) $\mathbb L_\midl\cong\mathbb L'_\midl\oplus\mathbb L_{\mathrm{tors},\midl}$ as $\Z_\ell$-modules, so $\mathbb L_\midl[\ell^{t}]=\mathbb L_{\mathrm{tors},\midl}$ and
$\ell^{t}\colon\mathbb L_\midl\to\ell^{t}\mathbb L_\midl$ has kernel $\mathbb L_{\mathrm{tors},\midl}$, inducing $\mathbb L'_\midl\cong\ell^{t}\mathbb L_\midl$. Since
$K\ge t$, $\ell^{t}N=\ell^{t}\mathbb L_\midl/\ell^{K}\mathbb L_\midl$, and the preimage of $\ell^{K}\mathbb L_\midl=\ell^{t}(\ell^{K-t}\mathbb L_\midl)$ in
$\mathbb L'_\midl$ is $\ell^{K-t}\mathbb L'_\midl$. Multiplication by $\ell^{t}$ commutes with every endomorphism. (2) $a\in\mathsf A_K$ kills $\ell^{t}N$
iff $\ell^{t}a=0$ in $\End N$; and $\Talg^{\mathrm{coh}}_\midl\to\Talg^{\mathbb L}_\midl$ is onto. (3) $N=\mathbb L_\midl/\ell^{K}\mathbb L_\midl=\mathbb L_\midl$, and
$\Talg^{\mathrm{coh}}_\midl$ acts faithfully on $\mathbb L_\midl$ by definition.
\end{proof}

\begin{corollary}[The certified reading of the $\Z/\ell^{K}$ Hecke-algebra route]\label{cor:stablecert}
\PROVED{} (given Lemma~\ref{lem:heckerankone}) In the Setting, let $S$ be a finite set of primes, $\Talg':=\Z_\ell[T_p:p\in S]\subseteq\Talg^{\mathbb L}$,
$\midl'$ a maximal ideal of $\Talg'$ (the code's joint system), $e$ its idempotent, $f:=\rank(e\mathbb L')$, and $\mathsf A'_K$ the image of
$\Talg'_{\midl'}$ in $\End(e\mathbb L'/\ell^{K})$ (v1.15.0-pre.2, WP-4 referee B1: the block of the code is one of $\Talg'$, which without the
certificate of (2) may be a sum of several blocks of $\Talg^{\mathbb L}$).
\begin{enumerate}
\item If $\dim_{\Fl}\mathsf A'_K/\ell\mathsf A'_K=\lceil f/2\rceil$ at some $K$, then $\rank\Talg'_{\midl'}=\rank\,e\Talg^{\mathbb L}$ and
  $\mathsf A'_K/\ell\mathsf A'_K\cong\Talg'_{\midl'}/\ell\Talg'_{\midl'}$ as $\Fl$-algebras (Hilbert function and Loewy length included), at this and every
  larger $K$; such a $K$ exists iff $\rank\Talg'_{\midl'}=\rank\,e\Talg^{\mathbb L}$, and then $K\le c^{\mathrm{sat}}(\Talg'_{\midl'})+1$. \LEANlit{cor:stablecert}{Eichler--Shimura at level one (\cite[Thm.~12.2.2, Prop.~12.4.10]{diamond_im}), as the hypothesis \texttt{hES}} (since v1.19.0 at $T_e=e\,T^{\mathbb L}$ on $e\mathbb L'$, for every $\Talg'\hookrightarrow T_e$; the rank hypothesis is discharged from Lemma~\ref{lem:heckerankone})
\item $\Talg'_{\midl'}=e\Talg^{\mathbb L}$ (then $e\mathbb L'$ is one block of $\Talg^{\mathbb L}$) iff the images of the $T_p$, $p\in S$, generate
  $e\Talg^{\mathbb L}/\ell$ as an $\Fl$-algebra (Nakayama). \LEAN{cor:stablecert-uncond} (this sentence)
  This \emph{generation certificate} is a separate check; without it, (1) certifies $\Talg'_{\midl'}/\ell$, not $e\Talg^{\mathbb L}/\ell$, and a
  proper suborder of full rank can have a different, even a larger, Loewy length ($\{(a,b,c)\in\Z_\ell^{3}:a\equiv b\equiv c\ (\ell)\}$ has
  reduction of Loewy length $2$, its suborder $\Z_\ell[(0,\ell,\ell^{2})]$ one of Loewy length $3$).
\item Applied to a torsion-free block ($\mathbb L_\midl=\mathbb L'_\midl$) this is the \code{stable} verdict; applied through
  Lemma~\ref{lem:freepart}(2) to the free part of a torsion block, the \code{torsion\_free\_stable} verdict; Lemma~\ref{lem:freepart}(3) is
  the \code{torsion\_stable} verdict. The full algebra of a torsion block with $f>0$ is \emph{not} covered: $\Talg^{\mathrm{coh}}_\midl$ may
  have $\Z_\ell$-torsion, and no stopping rule is asserted for it.
\end{enumerate}
\end{corollary}

\begin{proof}
(1) $\dim\mathsf A'_K/\ell\le\rank\Talg'_{\midl'}\le\rank\,e\Talg^{\mathbb L}=\lceil f/2\rceil$ (Proposition~\ref{prop:saturation}(3) for
$(e\mathbb L',\Talg'_{\midl'})$; Lemma~\ref{lem:heckerankone} applied to $e$; $\Talg'_{\midl'}=e\Talg'\subseteq e\Talg^{\mathbb L}$); equality forces both
equalities, and Proposition~\ref{prop:saturation}(3),(5) give the rest. (2) $\Talg'_{\midl'}\subseteq e\Talg^{\mathbb L}$ are finite $\Z_\ell$-modules;
$\Talg'_{\midl'}+\ell e\Talg^{\mathbb L}=e\Talg^{\mathbb L}$ implies
equality by Nakayama, and the image of $\Talg'_{\midl'}$ in $e\Talg^{\mathbb L}/\ell$ is the subalgebra generated by the images of the $T_p$. The
example: the reduction of the first order is $\Fl[x,y]/(x,y)^{2}$, and $t=(0,\ell,\ell^{2})$ satisfies $t(t-\ell)(t-\ell^{2})=0$ with
$\Z_\ell[t]/\ell=\Fl[t]/t^{3}$. (3) is the dictionary of \file{HeckeExt/HeckeRingHeckeAlgebra.h} ((ii)--(iv) and its stopping rule):
its block $e_\midl M_K$ is $e\mathbb L/\ell^{K}$ for the idempotent $e$ of its joint system $\midl'$ of $\Talg'$ (Lemma~\ref{lem:freepart} read with
$e$ in place of the localisation at a maximal ideal of $\Talg^{\mathrm{coh}}$, which changes nothing in its proof), its $\mathsf A_K$ is the algebra generated
by $T_p$, $p\in S$, and its $\lceil f/2\rceil$ is Lemma~\ref{lem:heckerankone} for $e$.
\end{proof}

\subsection{Cohomology against modular forms, with torsion}\label{sec:cohformssec}

\begin{proposition}[The cohomological and the forms Hecke algebra]\label{prop:cohforms}
In the Setting:
\begin{enumerate}
\item \PROVED{} $H^{1}(\Ga;\Mn)=\mathbb L/\ell\mathbb L$ Hecke-equivariantly, an extension
  $0\to\mathbb L_{\mathrm{tors}}/\ell\to H^{1}(\Ga;\Mn)\to\mathbb L'/\ell\to0$; $\mathbb L_{\mathrm{tors}}/\ell\cong\mathbb L[\ell]\cong H^{0}(\Ga;\Mn)$ as
  $\Fl$-spaces (Claim~\ref{cl:bockstein}). \LEAN{prop:cohforms} (the exact sequence and the equality of lengths; not the identification $H^{1}(\Ga;\Mn)=\mathbb L/\ell\mathbb L$, an interface of [GT])
\item \PROVED{} Restriction $\Talg^{\mathrm{coh}}\to\Talg^{\mathbb L}$ is onto, with kernel the operators mapping $\mathbb L$ into $\mathbb L_{\mathrm{tors}}$. \LEAN{prop:cohforms}
\item \PROVED{} (given the \textsf{[standard]} inputs of Lemma~\ref{lem:heckerankone}) $T_p\mapsto T_p$ extends to an isomorphism
  $\Talg^{\mathbb L}\cong\Talg^{\mathrm f}$: \textbf{the natural map runs from cohomology to forms}, $\Talg^{\mathrm{coh}}\twoheadrightarrow\Talg^{\mathbb L}
  \cong\Talg^{\mathrm f}$. \LEANlit{prop:cohforms}{Eichler--Shimura at level one and the forms-side Hecke algebra, as the hypotheses \texttt{hES}, \texttt{hF}}
\item \PROVED{} Modulo $\ell$: with $\mathsf A:=\mathrm{im}(\Talg^{\mathrm{coh}}\to\End H^{1}(\Ga;\Mn))$ (the census algebra) and $\mathsf A^{\mathbb L}:=
  \mathrm{im}(\Talg^{\mathrm{coh}}\to\End(\mathbb L'/\ell))$, there are natural surjections $\mathsf A\twoheadrightarrow\mathsf A^{\mathbb L}$ and
  $\Talg^{\mathrm f}/\ell\cong\Talg^{\mathbb L}/\ell\twoheadrightarrow\mathsf A^{\mathbb L}$ \LEAN{prop:cohforms} (the first surjection $\mathsf A\twoheadrightarrow\mathsf A^{\mathbb L}$) \LEANlit{prop:cohforms}{the hypotheses \texttt{hES}, \texttt{hF}} (the second surjection, since v1.18.0). \textbf{No natural map between $\mathsf A$ and
  $\Talg^{\mathrm f}/\ell$ is asserted}: both are quotients of $\Talg^{\mathrm{coh}}/\ell$, with kernels $(\Talg^{\mathrm{coh}}\cap\ell\End\mathbb L)/\ell$
  and $(\ker(\Talg^{\mathrm{coh}}\to\Talg^{\mathbb L})+\ell\Talg^{\mathrm{coh}})/\ell$, and neither contains the other in general --- an operator may vanish on
  $\mathbb L'$ without vanishing on a torsion class, and conversely.
\item \PROVED{} If $\mathbb L_\midl$ is torsion-free: $\mathsf A_\midl=\mathsf A^{\mathbb L}_\midl=\Talg^{\mathbb L}_\midl/(\Talg^{\mathbb L}_\midl\cap\ell\Talg^{\mathbb L,\mathrm{sat}}_\midl)$,
  a quotient of $\Talg^{\mathrm f}_\midl/\ell$ with kernel $(\Talg^{\mathbb L}_\midl\cap\ell\Talg^{\mathbb L,\mathrm{sat}}_\midl)/\ell\Talg^{\mathbb L}_\midl$, and
  $\mathsf A_\midl\cong\Talg^{\mathrm f}_\midl/\ell$ iff $\Talg^{\mathbb L}_\midl$ is saturated in $\End(\mathbb L_\midl)$. This holds if $\mathbb L_\midl$ contains a
  $\Talg^{\mathbb L}_\midl$-submodule $N\cong\Talg^{\mathbb L}_\midl$ with torsion-free quotient --- e.g.\ $\mathbb L_\midl$ free over $\Talg^{\mathbb L}_\midl$,
  or one of the $\pm$-parts of $\inv^{*}$ (\S\ref{sec:involution}) free of rank one over it (multiplicity one).
\item \PROVED{} On every block, $\Talg^{\mathrm f}_\midl/\ell$ is recovered from cohomology as $\mathsf A_K/\ell\mathsf A_K$ of
  Corollary~\ref{cor:stablecert} (torsion-free blocks) or of its free part (Lemma~\ref{lem:freepart}), with the generation caveat of
  Corollary~\ref{cor:stablecert}(2).
\end{enumerate}
\end{proposition}

\begin{proof}
(1) Claim~\ref{cl:bockstein}(2) with $\Z_\ell$-coefficients ($\Z_\ell$ is flat over $\Z$); the extension is $\mathbb L/\ell$ modulo the image of
$\mathbb L_{\mathrm{tors}}$, which is $\mathbb L_{\mathrm{tors}}/(\mathbb L_{\mathrm{tors}}\cap\ell\mathbb L)=\mathbb L_{\mathrm{tors}}/\ell\mathbb L_{\mathrm{tors}}$, the torsion being a
direct summand as a $\Z_\ell$-module; its dimension is the number of cyclic summands, $\dim\mathbb L[\ell]$. Hecke operators preserve
torsion, so the sequence is equivariant. (2) Every $\Z_\ell$-linear endomorphism preserves $\mathbb L_{\mathrm{tors}}$; the image of the generators
generates the image. (3) By Lemma~\ref{lem:heckerankone}'s proof, $\mathbb L'\otimes\bar\Q_\ell$ and $M_k(\bar\Q_\ell)$ carry the same set of
systems, each semisimply; so a polynomial in the $T_p$ with $\Q_\ell$-coefficients vanishes on one iff it vanishes on the other (both
iff it vanishes at every system). $\mathbb L'$ and $M_k(\Z_\ell)$ are torsion-free, so the kernels of $\Z_\ell[X_p]\to\End(\mathbb L')$ and
$\to\End(M_k(\Z_\ell))$ are the $\Q_\ell$-kernels intersected with $\Z_\ell[X_p]$, hence equal. (4) An endomorphism of $\mathbb L/\ell$ induced
by $\Talg^{\mathrm{coh}}$ preserves $\mathbb L_{\mathrm{tors}}/\ell$ and induces one of the quotient $\mathbb L'/\ell$, zero if the first is zero; the
second surjection is $\Talg^{\mathbb L}\to\End(\mathbb L'/\ell)$ reduced modulo $\ell$. The kernels are read off
Proposition~\ref{prop:saturation}(1) and (2). (5) $\mathbb L_\midl=\mathbb L'_\midl$; Proposition~\ref{prop:saturation}(1),(2) for
$(\mathbb L_\midl,\Talg^{\mathbb L}_\midl)$, and (3). (6) Corollary~\ref{cor:stablecert}(1) and (3).
\end{proof}

\noindent
\emph{Where the torsion lives.} $\mathbb L_{\mathrm{tors}}\ne0$ only on the Dickson locus ($\dim H^{0}(\Ga;\Mn)>0$, Claim~\ref{cl:dimH1}(2)); the
blocks it meets are those of the systems of Proposition~\ref{prop:heckeH0}, and it is non-ordinary (Remark~\ref{rmk:ordtorsionfree}). On a
block without such a system, (5) applies. The census rows of slot-13 WP-9 were torsion-free on all $654$ non-reduced rows tested
(\COMP\textsuperscript{a}, \file{docs/implementation/slot13/WP-9\_census\_lattice.tsv}); that is data, not a statement of this document.

\subsection{Hida's control at the level of algebras (P18)}\label{sec:hidaalgsec}

\begin{proposition}[The ordinary cuspidal Hecke algebra modulo $\ell$ is constant along a class]\label{prop:hidaalgconst}
Let $\ell\ge5$, $a\bmod\ell-1$, and $k,k'\ge2$ with $k\equiv k'\equiv a\pmod{\ell-1}$. Write $h^{0}_k:=h^{0}_k(\Ga_1(\ell),a)$ and
$h^{0}:=h^{0}(\ell,a)$ in the notation of Theorem~\ref{thm:hidacontrol}(3).
\begin{enumerate}
\item \PROVED{} (from Theorem~\ref{thm:hidacontrol}(3), \textsf{[standard]}) $(\ell,P_k)=(\ell,X)$ in $\Lambda=\Z_\ell[[X]]$, so
  $h^{0}_k/\ell h^{0}_k\cong h^{0}/(\ell,X)h^{0}\cong h^{0}_{k'}/\ell h^{0}_{k'}$ as $\Fl$-algebras, each Hecke operator of the universal
  algebra going to itself. \LEANlit{prop:hidaalgconst}{Hida's control theorem, Theorem~\ref{thm:hidacontrol}(3)} (as an isomorphism of rings
  $h^{0}_k/\ell h^{0}_k\cong h^{0}_{k'}/\ell h^{0}_{k'}$; not ``each Hecke operator going to itself'')
\item \PROVED{} (from (1)) The maximal ideals of $h^{0}_k$ with their residual systems of Hecke eigenvalues, and the Hilbert functions and
  Loewy lengths of the local factors of $h^{0}_k/\ell$, do not depend on $k$ in the class.
\item \textsf{[standard]} (the level-$1$ reading; the citations of Proposition~\ref{prop:stabilisation}(4)) For $k\ge3$ the ordinary cuspidal
  Hecke algebra of level~$1$ and weight $k$ (on $e_1S_k(\Z_\ell)$, $T_\ell$ the level-one operator) is isomorphic to $h^{0}_k$, with
  $T_p\mapsto T_p$ ($p\ne\ell$) and $T_\ell\mapsto U_\ell+\ell^{\,k-1}U_\ell^{-1}$; so modulo $\ell$, $T_\ell\mapsto U_\ell$, and (2) holds for the
  level-$1$ ordinary cuspidal algebras for $k,k'\ge3$.
\end{enumerate}
Theorem~\ref{thm:hidaweight}(1) is a statement about \emph{ranks}; every algebra-level consequence used in this project --- that the
residual ordinary systems and $h^{\mathrm{ord}}_k/\ell$ are constant along a class --- rests on this proposition, i.e.\ on
Theorem~\ref{thm:hidacontrol}(3) with Proposition~\ref{prop:stabilisation}, not on Theorem~\ref{thm:hidaweight}(1). The statements are on the
forms side (Hecke algebras of cusp forms), where Hida's lattice $L_n\ne\Vcoef$ (Remark~\ref{rmk:stabconvention}(3)) plays no role; they
reach the cohomology of $\Vcoef$ through Proposition~\ref{prop:cohforms}(3), for the free part only.
\end{proposition}

\begin{proof}
(1) $P_k=(1+X)-(1+\ell)^{k}=X-\bigl((1+\ell)^{k}-1\bigr)$ and $(1+\ell)^{k}-1\in\ell\Z_\ell$, so $(\ell,P_k)=(\ell,X)$. Theorem~\ref{thm:hidacontrol}(3)
gives $h^{0}/P_kh^{0}\cong h^{0}_k$; reducing modulo $\ell$, $h^{0}_k/\ell\cong h^{0}/(\ell,P_k)=h^{0}/(\ell,X)$, and the same for $k'$. The
isomorphisms are of $\Lambda$-algebras sending the universal $T(m)$ to $T(m)$. (2) An isomorphism of algebras preserves maximal ideals,
their residue fields, the images of each $T(m)$ and the $\midl$-adic filtrations. (3) For $k\ge3$ the ordinary cusp forms of weight $k$ on
$\Ga_0(\ell)$ are the ordinary $\ell$-stabilisations of the ordinary level-$1$ forms (an $\ell$-new form has $a_\ell^{2}=\ell^{\,k-2}$, not a
unit): the stabilisation $f\mapsto f(z)-\beta f(\ell z)$ with $\beta=\ell^{\,k-1}/u$ has $U_\ell$-eigenvalue the unit root $u$ of
$X^{2}-a_\ell X+\ell^{\,k-1}$ and the same $a_p$, $p\ne\ell$ (the sources cited in Proposition~\ref{prop:stabilisation}(4)). $u$ lies in
$\Z_\ell[T_\ell]$ (Hensel: $X^{2}-T_\ell X+\ell^{\,k-1}\equiv X(X-T_\ell)$ with $T_\ell$ a unit on the ordinary part), and $T_\ell=u+\ell^{\,k-1}u^{-1}$, so
the two algebras are generated by corresponding operators; $\ell^{\,k-1}u^{-1}\equiv0\pmod\ell$.
\end{proof}

\subsection{The forms-side \texorpdfstring{$\theta$}{theta} (P20)}\label{sec:thetaformssec}

\begin{remark}[The $\theta$ of modular forms modulo $\ell$ has a Frobenius kernel; the cohomological $\thL$ does not see it]\label{rmk:thetaforms}
(1) \PROVED{} On $q$-expansions, $\theta:=q\,\frac{d}{dq}$ maps $\sum a_mq^{m}\in\Fl[[q]]$ to $\sum ma_mq^{m}$, so $\ker\theta=\Fl[[q^{\ell}]]$. If
$g\in M_{k_0}(\Fl)$ then $V(g):=g(q^{\ell})=g^{\ell}$ (Frobenius on coefficients in $\Fl$) is a mod-$\ell$ modular form of weight $\ell k_0$,
and $E_{\ell-1}^{\,j}V(g)$, $E_{\ell-1}\equiv1$, one of weight $\ell k_0+j(\ell-1)$ with the same $q$-expansion; it lies in $\ker\theta$. So the
forms-side $\theta$ on $M_k(\Fl)_\midl$ is \emph{not} injective whenever such a $V(g)$ lies in the block --- although
Theorem~\ref{thm:thetales}(3) makes the cohomological $\thL$ injective on $H^{1}(\Ga;\Mn)$ away from the Eisenstein line.
\LEAN{rmk:thetaforms} (the two power-series facts $\ker\theta=\Fl[[q^{\ell}]]$ and $g(q^{\ell})=g^{\ell}$ only, nothing about modular forms)
(2) \COMP\textsuperscript{a} (slot-14 WP-C9 \S3.3, \file{theta\_check.py} on $q$-expansions; the cluster run of that report is not published,
the forms-side check is local and re-runnable) At $k=176$, $\ell=7$ (the cell $(174,7)$), $\theta$ has rank $6$ on the $7$-dimensional
block $W_{176}$ of the system $(a_2,a_3)\equiv(1,0)$, its kernel spanned by $V(g)$ with $g$ the weight-$20$ eigenform in that system;
likewise at $(406,11)$ (rank $6$ of $7$); at $(144,7)$ rank $5$ of $6$ with no change of Hilbert function. Along the cohomologically
bijective $\thL$-steps $(174\to182,7)$ and $(406\to418,11)$ (equal $\dim$, $\mathsf A_1$ preserved) the forms-side $\Talg_\midl/\ell$ changes Hilbert
function $(1,3,2,1)\to(1,2,2,2)$ at equal dimension and Loewy length.
(3) This is consistent with Proposition~\ref{prop:cohforms}(4),(5): $\Talg^{\mathrm f}_\midl/\ell$ is not a function of $H^{1}(\Ga;\Mn)_\midl$ (the
saturation defect intervenes), so a map bijective on $H^{1}(\Ga;\Mn)$ need not preserve it. No statement relating the forms-side $\theta$
and $\thL$ as maps is made here.
\end{remark}

\subsection{The saturation defect as a \texorpdfstring{$\mathrm{Tor}$}{Tor} term}\label{sec:sattorsec}
% v1.17.0-pre.1 (slot-17 WP-2): placed at the end of the chapter so that no statement number of v1.16.0 moves.

\begin{remark}[The saturation defect as a $\mathrm{Tor}$ term]\label{rmk:sattor}
\PROVED{} (v1.17.0-pre.1; proposed by slot-16 WP-5, \file{docs/implementation/slot16/WP-5\_tor\_note.md}~\S5) In the notation of
Definition~\ref{def:satdata} put $G:=\Talg^{\mathrm{sat}}/\Talg$ and $I:=\mathrm{Tor}_1^{\Z_\ell}(G,\Fl)$.
\begin{enumerate}
\item $I=G[\ell]\cong(\Talg\cap\ell\Talg^{\mathrm{sat}})/\ell\Talg=\ker\bigl(\Talg/\ell\Talg\to\End(L/\ell L)\bigr)$, canonically; the last two are ideals of
  $\Talg/\ell\Talg$, and $(\Talg/\ell\Talg)/I\cong\mathsf A_1$ as $\Fl$-algebras.
\item $\dim_{\Fl}I=\varepsilon^{\mathrm{sat}}_1=r-\dim_{\Fl}\mathsf A_1$, the number of cyclic summands of $G$; and $I=0$ iff $\Talg$ is saturated. More
  generally the number of summands $\Z/\ell^{e}$ of $G$ with $e\ge j$ is $\varepsilon^{\mathrm{sat}}_j-\varepsilon^{\mathrm{sat}}_{j-1}$ ($\varepsilon^{\mathrm{sat}}_0:=0$).
\item At every $K$ at which the equivalent conditions of Proposition~\ref{prop:saturation}(3) hold, $I$ is the kernel of the composite
  $\mathsf A_K/\ell\mathsf A_K\to\End\bigl((L/\ell^{K}L)/\ell\bigr)=\End(L/\ell L)$, read through the isomorphism $\Talg/\ell\Talg\cong\mathsf A_K/\ell\mathsf A_K$ of
  Proposition~\ref{prop:saturation}(3)(b).
\item If $L$ contains a $\Talg$-submodule $N\cong\Talg$ with $L/N$ torsion-free, then $I=0$ (Proposition~\ref{prop:saturation}(2)); the converse fails:
  a saturated $\Talg$ need not have such an $N$.
\item $I$ is not $\mathrm{Tor}_1^{\Talg}(L,\Talg/\midl)$: for a maximal ideal $\midl$ of a local $\Talg$ the latter vanishes iff $L$ is free over
  $\Talg$ (the multiplicity-one question), so by (4) $I\ne0$ implies that $L$ is not free over $\Talg$, while $I=0$ says nothing about freeness.
\end{enumerate}
\end{remark}

\begin{proof}
(1) Tensoring $0\to\Z_\ell\xrightarrow{\ell}\Z_\ell\to\Fl\to0$ with $G$ gives $\mathrm{Tor}_1^{\Z_\ell}(G,\Fl)=\ker(\ell\colon G\to G)=G[\ell]$ ($\Z_\ell$ is
free). Inside $\Talg\otimes\Q$, $G[\ell]=(\ell^{-1}\Talg\cap\Talg^{\mathrm{sat}})/\Talg$, and multiplication by $\ell$ is a $\Talg$-linear isomorphism
$(\ell^{-1}\Talg\cap\Talg^{\mathrm{sat}})/\Talg\to(\Talg\cap\ell\Talg^{\mathrm{sat}})/\ell\Talg$ ($\Talg\otimes\Q$ is torsion-free). By
Proposition~\ref{prop:saturation}(1) at $K=1$, $\ker(\Talg\to\End(L/\ell L))=\Talg\cap\ell\Talg^{\mathrm{sat}}\supseteq\ell\Talg$, whose image in $\Talg/\ell\Talg$ is
the displayed kernel; a kernel of a ring map is an ideal, and the quotient is the image $\mathsf A_1$.
(2) $\dim_{\Fl}G[\ell]=\mathrm{length}(G[\ell])=\varepsilon^{\mathrm{sat}}_1$ by Proposition~\ref{prop:satdefect}(1), and $\varepsilon^{\mathrm{sat}}_1=r-\mathrm{length}(\mathsf A_1)$ by
definition; for a finite $\Z_\ell$-module $G\cong\bigoplus_i\Z/\ell^{e_i}$, $\mathrm{length}(G[\ell^{j}])=\sum_i\min(e_i,j)$, so
$\mathrm{length}(G[\ell^{j}])-\mathrm{length}(G[\ell^{j-1}])=\#\{i:e_i\ge j\}$; at $j=1$ this is the number of summands. $I=0$ iff $G=0$
(Proposition~\ref{prop:satdefect}(3)).
(3) The composite $\Talg\to\mathsf A_K\to\End(L/\ell^{K}L)\to\End(L/\ell L)$ is the reduction map of (1), and $\ell\mathsf A_K$ acts as $0$ on $L/\ell L$, so it
factors through $\mathsf A_K/\ell\mathsf A_K$; under (3)(b) its kernel there is the kernel of $\Talg/\ell\Talg\to\End(L/\ell L)$, which is $I$ by (1).
(4) The first sentence is Proposition~\ref{prop:saturation}(2) with (2). For the second, take the order $\Talg:=\{(a,b,c)\in\Z_\ell^{3}:a\equiv b\equiv c
\pmod\ell\}$ of Corollary~\ref{cor:stablecert}(2), with $\Talg/\ell\Talg\cong\Fl[x,y]/(x,y)^{2}$, and $L:=\Talg^{\vee}=\Hom_{\Z_\ell}(\Talg,\Z_\ell)$ with
$(t\lambda)(s):=\lambda(ts)$; $\Talg$ acts faithfully ($t\lambda=0$ for all $\lambda$ gives $\lambda(t)=0$ for all $\lambda$). If $x\in\Talg\otimes\Q$
preserves $L$ (the action extended $\Q$-linearly), then $\lambda(x)=(x\lambda)(1)\in\Z_\ell$ for every $\lambda\in\Talg^{\vee}$, so $x\in\Talg^{\vee\vee}=\Talg$
($\Talg$ is free): $\Talg$ is saturated (the argument uses nothing about this $\Talg$: every order is saturated on its $\Z_\ell$-dual). A submodule $N\cong\Talg$ has rank $3=\rank L$, so $L/N$
torsion-free forces $N=L$, i.e.\ $\Talg^{\vee}\cong\Talg$; then $\Talg^{\vee}/\ell\cong\Talg/\ell$ would be a free $\Talg/\ell$-module of rank one, so
$\Talg/\ell\cong(\Talg/\ell)^{\vee}$ would be Gorenstein, but $\Fl[x,y]/(x,y)^{2}$ has a two-dimensional socle.
(5) For $L$ finitely generated over the local Noetherian ring $\Talg$, $\mathrm{Tor}_1^{\Talg}(L,\Talg/\midl)=0$ iff $L$ is free (Nakayama and a minimal
presentation); a free $L\ne0$ contains $N\cong\Talg$ as a direct summand with free, hence torsion-free, quotient, so (4) applies.
\end{proof}

\paragraph{Scope: what this chapter does not do.}
It proves the saturation identities for any commutative order on a free lattice, the certified reading of the $\Z/\ell^{K}$ route on
free and torsion blocks, and the comparison of cohomological and forms Hecke algebras at level~$1$, $\ell\ge5$, given Eichler--Shimura.
Not asserted: why $\mathsf A$ has Loewy length $\le2$ on the non-reduced blocks where it has ([E] \Espec{pr:latticesplit}, \CONJ{} --- refuted as stated at $\ell=37$
(slot-15 WP-12~(a), \file{docs/implementation/slot15/WP-4\_referee.md} \S8), its v1.16.0 form refuted as stated at $(608,547)$ and contradicted on the forms side
(slot-17 R17-2; v1.17.0-pre.3, m-C3), its v1.17.0-pre.2 Kummer-free form contradicted on the forms side at irreducible congruence classes, the regular
$\ell=139$ among them, and (R) now asserted on congruence-free classes only (v1.17.0-pre.4, owner ruling on the referee F1's B-1); no bound on the Loewy
length of $\mathsf A$ is asserted here); any bound on
$c^{\mathrm{sat}}$ or $\delta^{\mathrm{sat}}$ in terms of $(n,\ell)$; a stopping rule for $\Talg^{\mathrm{coh}}_\midl$ on a torsion block with
$f_\midl>0$; (the faithfulness of $\Talg^{\mathrm f}/\ell$ on $M_k(\Fl)$ at $\ell\ge5$ is \Espec{cl:mkfaithful}, slot-15 WP-19, not restated here; at $\ell\in\{2,3\}$ nothing is asserted); levels other
than~$1$.

\paragraph{Code map.}
\file{HeckeExt/HeckeRingHeckeAlgebra} (slot-14 WP-6) computes $\mathsf A_K$, its length, the Hilbert function of $\mathsf A_K/\ell$,
$\varepsilon^{\mathrm{sat}}_K$ (its \code{eps}), the free part of a torsion block on $\ell^{t}$-scaled generators (Lemma~\ref{lem:freepart}(2)), and
the verdicts \code{stable}, \code{torsion\_free\_stable}, \code{torsion\_stable} of Corollary~\ref{cor:stablecert}(3) with $\Talg'$ the
algebra of its generator set $S$. At v1.15.0-pre.1 these verdicts were \textbf{readings}; since slot-15 WP-5 (built at v1.15.0; WP-4's ACCEPT of
Propositions~\ref{prop:saturation}, \ref{prop:satdefect}, owner ruling Q4) they are \emph{certified by Corollary~\ref{cor:stablecert}} in its
B1 form (for $\Talg'_{\midl'}$; for $\Talg^{\mathbb L}_\midl$ only with a generation certificate, which is not run), printed as
\code{stable\_cert=certified} (Fable audit ACCEPT WITH REPAIRS, repaired). The digit route for $\ell^{K}\ge2^{64}$ is built (WP-5).
\textbf{Built at v1.16.0, recorded at v1.17.0-pre.1} (slot-16 WP-5, \file{docs/implementation/slot16/WP-5\_report.md}; cross-model audit
\file{docs/implementation/slot16/WP-3\_WP-5\_audit.md}, repaired): \file{HeckeExt/HeckeSaturationDefect} computes Remark~\ref{rmk:sattor}'s $I$ on a
torsion-free block at a stable $K$ by two routes compared per block --- the kernel of (3) in the regular representation of $\mathsf A_K/\ell\mathsf A_K$, and
the length formula of (2) from the route's own $\mathsf A_1$ --- with its minimal generators, radical layers, Loewy length, socle and annihilator, and the
filtration $(\varepsilon^{\mathrm{sat}}_j)_j$ with the plateau of Proposition~\ref{prop:satdefect}(2) as \code{final}; torsion blocks are refused by name
(Corollary~\ref{cor:stablecert}(3) asserts no stopping rule there); \code{HeckeSaturationDefectValidation} (\code{standard}).
\textbf{At v1.23.0-pre.4} (slot 23, value-preserving, no statement moves): the route's levels come from one module run per cell in the common case (G1;
\S\ref{sec:satprecisionsec}'s code map), and the module route under it is the G2 re-implementation mapped in the code map of \S\ref{sec:obligations}. The forms
side of Proposition~\ref{prop:cohforms} (the sat route's forms cross-route, \code{BuildForms} / \code{RunForms} in \file{HeckeExt/HeckeSaturationDefect.cpp},
and every q-series product of \file{HeckeExt/HeckeFormsHeckeAlgebra.\{h,cpp\}}) multiplies through \code{forms::MulSeries}, which since slot-23 G2-X
(\file{docs/implementation/slot23/G2X\_report.md}, \file{G2X\_audit.md} with its fixes; merged at \texttt{1ba46525}) is \code{MulSeriesFast}: the exact integer
convolution of the residues by NTT over the primes $5\cdot2^{25}+1$ and $7\cdot2^{26}+1$ and CRT, reduced modulo $\ell$, inside its checked exactness bounds;
short operands take \code{MulSeriesSchoolbook} (delayed reduction), products past the transform length \code{MulSeriesChunked} (blocks of $2^{24}$, each an
NTT product within both bounds), and $\ell\notin[2,2^{16})$ the naive product; \code{EisensteinSeries} takes its divisor sums by a sieve. \textbf{The rollback
lever is value-neutral}: \code{HECKE\_FORMS\_FASTMUL=0} restores the $O(N^{2})$ product \code{MulSeriesNaive} and the trial-division loop (any other value warns
and keeps the default; scoped as \code{ScopedFormsFastMul}; its \code{EFFECTIVE} line \code{FormsFastMulEffectiveLine()}, counters
\code{FormsFastMulCountersNow()}), and since the truncated product modulo $\ell$ and the divisor sums are unique, no route can move a byte; gate
\code{HeckeFormsFastMulValidation} (\code{standard}; \file{Registry/Slot23/G2X.cpp}, source \file{HeckeExt/HeckeFormsFastMulValidation.cpp}): every route against
the naive oracle byte for byte, the identities $E_4^{2}=E_8$ and $E_4^{3}-E_6^{2}=1728\Delta$ beyond the oracle's reach, and whole \code{SAT\_DEFECT} cells.
\textbf{At v1.23.0-pre.10} (slot-24 FNTT, value-preserving; \file{docs/implementation/slot24/FNTT\_report.md}, \file{FNTT\_audit.md} with its fixes): \code{BuildCuspSpace} /
\code{BuildModularSpace} compute the same basis series in the transform domain of the same two primes (the spectra of $E_6^{2}$ and $\Delta$ cached, $S_k$ and
the extra basis series of $M_k$ in one pass, no power tables), and \code{HeckeOnCusp} reduces modulo $\ell$ once per accumulated sum (whenever both levers \code{HECKE\_FORMS\_FASTMUL}, \code{HECKE\_FORMS\_NTT\_CACHE} are on and
$(m+1)\ell^{2}<2^{64}$, $m$ the number of columns, at any $N$); the transform-domain route applies when the fast product is on, $N\ge1024$, $2N-1\le2^{25}$ and $N(\ell-1)^{2}<(5\cdot2^{25}+1)(7\cdot2^{26}+1)$
(\code{FnttApplicable}; every operand is a residue vector of length $\le N$, so every coefficient of the exact integer convolution lies in $[0,N(\ell-1)^{2}]$ and
the CRT recovers it), the earlier route elsewhere. Since the truncated product modulo $\ell$ is unique, no route moves a byte; \textbf{the rollback lever is
value-neutral}: \code{HECKE\_FORMS\_NTT\_CACHE=0} restores the G2-X route verbatim (any other value warns and keeps the default; scoped as
\code{ScopedFormsNttCache}; its \code{EFFECTIVE} line \code{FormsNttEffectiveLine()}). \textbf{The forms $\ell$-range}: the sentence above on
$\ell\notin[2,2^{16})$ stays true of \code{MulSeriesFast}, but no longer bounds the forms cross-route. A forms space is built (\code{forms::FormsSpaceSupported}$(\ell,N)$)
for every $5\le\ell<2^{16}$, and for $2^{16}\le\ell<2^{27}$ only when $N<1024$ or the route above applies and is exact; any other space is refused by
\code{BuildCuspSpace} / \code{BuildModularSpace} (an \code{ENGINE\_CHECK} naming the bounds), and a cell whose forms space would be refused is refused by
\code{ComputeSaturationDefectCell} and the \code{HECKE\_SAT\_DEFECT\_CELLS} parser before the module run, so no forms product at $\ell\ge2^{16}$, $N\ge1024$
runs by the $O(N^{2})$ product. At the sat route's precision $N=2(k/12+1)p_{\max}+2$ (\code{forms::SatFormsPrecision}) the condition is roughly
$2(k/12+1)\ell^{3}<7.9\cdot10^{16}$: $\ell=100003$ is admitted for $n\le464$, the largest admitted integer $\ell$ is $114868$ at $n=300$ and $77702$ at $n=1000$ (the largest admitted primes, the only $\ell$ a cell takes, $114859$ and $77699$; the audit's
derivation, \file{FNTT\_audit.md}~F5). The gate is \code{HeckeFormsFastMulValidation} (arms (N1)--(N3), (N5), (N6): the transform-domain route against the schoolbook, the naive
product and the G2-X route byte for byte, at $\ell$ up to $1000003$, and the bound at its edge).

\paragraph{Formalization note.}
Propositions~\ref{prop:saturation}, \ref{prop:satdefect}, Lemma~\ref{lem:freepart}, Corollary~\ref{cor:stablecert} and Remark~\ref{rmk:sattor}
(v1.17.0-pre.1) are commutative algebra over $\Z_\ell$ (finite modules, Nakayama, $\mathrm{Tor}$ over a DVR) and a good first target; Lemma~\ref{lem:heckerankone} and Proposition~\ref{prop:cohforms}(3) need
Eichler--Shimura, Proposition~\ref{prop:hidaalgconst} Hida theory.
(v1.23.0-pre.3) Lemma~\ref{lem:precisionreduction} is homological algebra over $\Z_\ell$ (one long exact sequence and its naturality) on top of
Proposition~\ref{prop:ringinforce}(1), Remark~\ref{rmk:orbifold} (with its naturality in the ring), Claim~\ref{cl:amalgam}(3) and
Proposition~\ref{prop:coneintegral}(1),(2), Proposition~\ref{prop:coneindep}(2)--(4) with Lemma~\ref{lem:tauexists}(2) over $\Z_\ell$, applied to the bulk
complex alone, and Theorem~\ref{thm:relations}(1) (the $T_p$ commute on $M_K$; used by (2), (3) and the module structure in (1); its proof uses
Theorem~\ref{thm:classical}, and its Lean marker is a \code{\textbackslash LEANlit} with the hypothesis \code{hcomm}); the compatibility of its (1) with the classical $T_p$ on $\mathbb L$ (as opposed to the
$T_p$ transported along Remark~\ref{rmk:orbifold}) needs Theorem~\ref{thm:classical} as well, which is unformalized; it carries no \code{\textbackslash LEAN} marker.

\subsection{One module run per cell: the precision levels as reductions}\label{sec:satprecisionsec}
% v1.23.0-pre.3 (hecke-extensions slot 23, G1; owner decision 2026-10-16 approving lever G1 of docs/implementation/slot23/ELL_scout.md sec. 3;
% note docs/implementation/slot23/G1_clause_note.md): placed at the end of the chapter so that no statement number moves.

\paragraph{Why this section exists.}
The route of Corollary~\ref{cor:stablecert} reads each block at the levels $K=1$ (through $M_K/\ell$), $K=2,3,\dots$ up to the first stable,
final level, and at $K_{\mathrm{cls}}:=\max(2,e_{\max}+1)$, $e_{\max}$ of Proposition~\ref{prop:torsionprepass}, never above $K_{\mathrm{need}}:=\max(\mathrm{KMAX},K_{\mathrm{cls}})$.
Here $\mathrm{KMAX}$ is the route's level cap, a run setting, not a mathematical quantity (the largest precision read for the stopping rule: the code's
\code{LoewyOptions::kmax} and \code{SatOptions::kmax}, default $6$; a block not stable by it reads \code{not\_stable\_at\_KMAX}), and $K_{\mathrm{stab}}$ is the
level at which the stopping rule (read from $K=2$; on a torsion block decided only from $K_{\mathrm{cls}}$) first reads stable in the sense of Corollary~\ref{cor:stablecert}(1), i.e.\ $\dim_{\Fl}\mathsf A'_K/\ell\mathsf A'_K=\lceil f/2\rceil$ (on a
torsion block, for its free part, Corollary~\ref{cor:stablecert}(3)); so a block stable already at $K=1$ reads $2$. It is the code's \code{LoewyBlock::K\_stab}, printed as \code{K\_stab=},
on a stable verdict read at $K\ge2$ (on a not-stable block the field holds the last level read, on \code{reduced\_rank1}, read at $K=1$ only, it holds $0$).
Through v1.23.0-pre.2 the implementation ran the whole module route once per level read: at $(n,\ell)=(300,7)$ five runs, $K=2,\dots,6$,
$377.1$ of the cell's $378.0$\,s (\file{ELL\_scout.md}~\S1.3). Lemma~\ref{lem:precisionreduction} says that one run at the largest level read determines
every lower level up to isomorphism of Hecke modules, which is all the route reads; Obligation~\ref{obl:precisionreuse} says how the implementation
may use that. Nothing in this section changes a value.

\begin{lemma}[The precision levels of the module route are reductions of one run]\label{lem:precisionreduction}
\PROVED{} Assume \ref{H4}, $\ell\ge5$ and $n\ge0$. Let $S$ be a finite set of primes and, for $p\in S$, let $T_p$ act on $\Ceq(\Z/\ell^{K})$ through
$\Th_g$ (Definition~\ref{def:xi}; $T_\ell=\Tdc{g}$ with $g=\mathrm{diag}(1,\ell)$, as in the Setting of \S\ref{sec:satlevelone}) formed, at every level $K$,
from a $\tau$-table with every fill over $\Z$ (Remark~\ref{rmk:tauintegral}; the table may be computed afresh at each level, each one integral) and the
reduced $\rhon(g\gamma)$. For $K\ge1$ put
$M_K:=\Heq{1}(\Ceq(\Z/\ell^{K}))$ with the induced operators $T_p$, and for $1\le K\le K'$ let $r_{K',K}\colon M_{K'}\to M_K$ be the map induced by the
reduction $\Z/\ell^{K'}\to\Z/\ell^{K}$.
\begin{enumerate}
\item \textbf{(Reduction.)} $r_{K',K}$ is surjective, commutes with every $T_p$, and has kernel $\ell^{K}M_{K'}$. Hence it induces an isomorphism
  $M_{K'}/\ell^{K}M_{K'}\cong M_K$ of $\Z[X_p:p\in S]$-modules ($X_p$ acting as $T_p$). With $\mathbb L=H^{1}(\Ga;\Vcoef\otimes\Z_\ell)$ there are
  isomorphisms $M_K\cong\mathbb L/\ell^{K}\mathbb L$, natural in $K$ and carrying $r_{K',K}$ to the canonical projection, compatible with the identification
  $H^{1}(\Ga;\Vcoef\otimes\Z/\ell^{K})=\mathbb L/\ell^{K}$ of Lemma~\ref{lem:freepart} through Remark~\ref{rmk:orbifold} (whose comparison is natural in the ring,
  being induced by one complex $C_{\bullet}(\Ku)$ for every $\cring$). They commute with the $T_p$ when the $T_p$ on $\mathbb L$ are transported along
  Remark~\ref{rmk:orbifold} over $\Z_\ell$ (the pipeline's operators, as in the Setting); that these are the classical double-coset operators on $\mathbb L$ is
  Theorem~\ref{thm:classical}, which this sentence then also rests on. The first sentence of (1), and (2)--(5), use neither directly; but (2) and (3), and
  the module structure over the polynomial ring $\Z[X_p:p\in S]$ in (1), use that the $T_p$ commute on $M_K$, Theorem~\ref{thm:relations}(1), whose proof uses
  Theorem~\ref{thm:classical} (through Remark~\ref{rmk:orbifold}); (4) and (5) use no Hecke relation.
\item \textbf{(Presentations.)} If $M_{K'}=\bigoplus_{i=1}^{D}(\Z/\ell^{a_i})g_i$, $1\le a_i\le K'$, with $T_pg_j=\sum_iA^{(p)}_{ij}g_i$, then
  $M_K=\bigoplus_{i=1}^{D}(\Z/\ell^{\min(a_i,K)})\bar g_i$, $\bar g_i:=r_{K',K}(g_i)$, with $T_p\bar g_j=\sum_i\bigl(A^{(p)}_{ij}\bmod\ell^{\min(a_i,K)}\bigr)\bar g_i$.
  No summand vanishes, so $D=\dim_{\Fl}M_1$ at every level. Write $M_{K'}{\downarrow}K$ for this presentation (the \emph{derived module}). If
  $e'$ is an idempotent of the image $\mathsf B_{K'}$ of $\Z_\ell[X_p:p\in S]$ in $\End(M_{K'})$, its image $e$ in $\mathsf B_K$ is an idempotent, $r_{K',K}$
  restricts to an isomorphism $e'M_{K'}/\ell^{K}\cong eM_K$, and $e$ is the unique idempotent of $\mathsf B_K$ with the same image in $\mathsf B_1$; so the
  block of a joint system at level $K$ is the reduction of its block at level $K'$.
\item \textbf{(Readings.)} Let $\Phi$ be any function of a finite $\Z_\ell$-module with commuting endomorphisms indexed by $S$ that is constant on
  isomorphism classes of $\Z_\ell[X_p:p\in S]$-modules. Then $\Phi(M_{K'}{\downarrow}K)=\Phi(M_K)$, and likewise on blocks; and, by (1), likewise for every
  reading of the tower $(M_K,r_{K',K})_{K\le K'}$ that is constant on isomorphism classes of towers of $\Z_\ell[X_p:p\in S]$-modules (one that uses the maps
  between levels: the derived tower $(M_{K'}{\downarrow}K)_{K}$ with the reductions of (2) is isomorphic to it, compatibly with the $r_{K',K}$). In particular the joint
  $\Fl$-rational systems of $M_1$ and their dimensions, the block $eM_K$ and its cyclic exponents (at $K_{\mathrm{cls}}$: the free rank $f$ and the torsion
  exponents), the algebra $\mathsf A_K$ generated by the $T_p$ in $\End(eM_K)$ with its length, cyclic exponents and the Hilbert function of
  $\mathsf A_K/\ell\mathsf A_K$ at $\midl$, the free-part reading of Lemma~\ref{lem:freepart}(2), $\varepsilon^{\mathrm{sat}}_K$, the ideal $I$ of
  Remark~\ref{rmk:sattor}(3) with its radical layers, socle and annihilator, and the hypotheses and conclusions of Corollary~\ref{cor:stablecert}(1),(3)
  are the same read on $M_{K'}{\downarrow}K$ as on $M_K$; so are the monomial basis of $\mathsf A_K/\ell$ chosen greedily in a fixed order of the monomials in the
  $T_p$ and the generators of $I$ as the reduced echelon form of $I$ in that basis (the reduced echelon form of a subspace in fixed coordinates is unique,
  and the monomial coordinates are fixed by $(\mathsf A_K,(T_p)_p)$), hence the \code{SAT\_DEFECT} fields \code{generators}, \code{monomial\_basis},
  \code{quotient\_hilbert}, \code{dim\_image}, \code{tor\_routes} and \code{sat\_structure}, and the \code{forms=(}\dots\code{)} reading, which sees only the
  $K=1$ block through the dimension of a joint algebra.
\item \textbf{(Downward only.)} Nothing here determines a level above the levels run: the $\Z_\ell$-lattices with one endomorphism $(\Z_\ell,0)$ and
  $(\Z_\ell,\ell^{K})$ have isomorphic reductions modulo $\ell^{K}$ and non-isomorphic ones modulo $\ell^{K+1}$.
\item \textbf{(The hypothesis is the absence of $\ell$-torsion one degree up, and it is needed.)} Let $C$ be a complex of free $\Z_\ell$-modules of
  finite rank with a cochain map $\varphi$, $1\le K\le K'$, $q$ any degree. There is an exact sequence, natural in $K$ and $\varphi$-equivariant,
  \[
    0\to H^{q}(C)/\ell^{K}\to H^{q}(C/\ell^{K})\xrightarrow{\ \delta_K\ }H^{q+1}(C)[\ell^{K}]\to0,
  \]
  under which $r_{K',K}$ induces the reduction on the left and multiplication by $\ell^{K'-K}$ on the right. So $r_{K',K}$ is onto with kernel
  $\ell^{K}H^{q}(C/\ell^{K'})$ for every $K\le K'$ if and only if $H^{q+1}(C)[\ell]=0$; if $H^{q+1}(C)$ has a cyclic summand $\Z/\ell^{a}$, $a\ge1$, then
  $r_{K+1,K}$ is not onto for every $K\ge a$. Example: $C=(\Z_\ell\xrightarrow{\ell}\Z_\ell)$ in degrees $q,q+1$ has $H^{q}(C/\ell^{K})=\ell^{K-1}\Z/\ell^{K}\cong\Z/\ell$
  for every $K$ and $r_{K',K}=0$ for $K'>K$, although the groups are isomorphic. On the model this excludes $\Heq{0}$ on the Dickson locus, where
  $\mathbb L$ has $\ell$-torsion: there $\Heq{0}(\Ceq(\Z/\ell^{K}))\cong\mathbb L[\ell^{K}]$ and $r_{K',K}$ is multiplication by $\ell^{K'-K}$. Nothing is asserted at
  $\ell\in\{2,3\}$ (where $\Ceq(\Z/\ell^{K})$ is not a base change, Proposition~\ref{prop:ringinforce}(2)), for the cone, or at other levels.
\end{enumerate}
\end{lemma}

\begin{proof}
\emph{The integral complex.} Put $C:=\Ceq(\Z_\ell)$. By Proposition~\ref{prop:ringinforce}(1) ($6\in\Z_\ell^{\times}$) $C=\Ceq(\Z)\otimes\Z_\ell$, a complex of
free $\Z_\ell$-modules of finite rank (Definition~\ref{def:ringinforce}(1)), and $C/\ell^{K}C=\Ceq(\Z)\otimes\Z/\ell^{K}=\Ceq(\Z/\ell^{K})$ for every $K\ge1$. Fix one
integral table. It makes $\Th_g$ the base change of one integral cochain map $\Th_{g,\Z}$ of $\Ceq(\Z)$ (Proposition~\ref{prop:coneintegral}(1),(2), applied to
the bulk complex alone: \eqref{eq:cosettransport} has integer coefficients); write $\varphi$ for $\Th_{g,\Z}\otimes\Z_\ell$. The table used at level $K$ may be
another one, but it is integral too, and two integral tables give cochain maps of $C=\Ceq(\Z_\ell)$ (each the base change of its integral one, as above)
differing by $d\Upsilon+\Upsilon d$ with $\Upsilon$ $\Z_\ell$-linear (Proposition~\ref{prop:coneindep}(2) with Lemma~\ref{lem:tauexists}(2) over
$\cring=\Z_\ell$, $6\in\Z_\ell^{\times}$: the two $\tau$ differ by $\partial H+H\partial$ with $H$ additive with (T1); replacing $H$ by the $\Z_\ell$-linear extension of its restriction to
$C_{\bullet}(\Ku)\otimes1$ if necessary, $H$ is $\Z_\ell$-linear, keeps (T1), and still gives $\tau-\tau'=\partial H+H\partial$, both sides being $\Z_\ell$-linear, agreeing on
$C_{\bullet}(\Ku)\otimes1$, which spans, and $\partial$ preserving it; and
$\Th_g'=\Th_g+d\widetilde H+\widetilde Hd$, the bulk component of \eqref{eq:conehomotopytau} (Theorem~\ref{thm:welldef}(b)), so $\Upsilon:=\widetilde H$; Proposition~\ref{prop:coneindep}(3),(4), again under $6\in\Z_\ell^{\times}$, if the transporters $s_j$ or $g\in\Ga g\Ga$
are rechosen; applied to the bulk complex alone); reduced modulo $\ell^{K}$ the homotopy stays a homotopy ($C/\ell^{K}C=\Ceq(\Z/\ell^{K})$), so the operator the level-$K$ table induces on $H^{\bullet}(C/\ell^{K})$ is the one $\varphi\otimes\Z/\ell^{K}$ induces. That map commutes with the
reductions $C/\ell^{K'}\to C/\ell^{K}$; the matrices of $T_p$ on $M_K$ are those it induces (Algorithm~\ref{alg:assembly}). By Remark~\ref{rmk:orbifold} over
$\cring=\Z_\ell$ and Claim~\ref{cl:amalgam}(3), $H^{2}(C)\cong H^{2}(\Ga;\Vcoef\otimes\Z_\ell)=0$ and $H^{1}(C)\cong\mathbb L$; the comparison of
Remark~\ref{rmk:orbifold} is induced by the one complex $C_{\bullet}(\Ku)$ for every $\cring$, hence natural along $\Z_\ell\to\Z/\ell^{K}$.

(5) $C$ is torsion-free, so $0\to C\xrightarrow{\ell^{K}}C\to C/\ell^{K}\to0$ is a short exact sequence of complexes, and its long exact sequence gives the
displayed sequence \STD{\cite[Thm.~1.3.1]{weibel}}. The triple $(\ell^{K'-K},\id,\text{reduction})$ is a map from the sequence at $K'$ to the one at $K$
(left square: $\ell^{K}\circ\ell^{K'-K}=\id\circ\ell^{K'}$), so naturality of the connecting map gives $\delta_K\circ r_{K',K}=\ell^{K'-K}\circ\delta_{K'}$, and on the
left $r_{K',K}$ is induced by $\id_C$, i.e.\ it is the reduction $H^{q}(C)/\ell^{K'}\to H^{q}(C)/\ell^{K}$. $\varphi$ commutes with all vertical maps and with
$\ell^{K}$, hence with every map in sight. If $H^{q+1}(C)[\ell]=0$, then $H^{q+1}(C)[\ell^{K}]=0$ for all $K$, the left maps are isomorphisms, and $r_{K',K}$ is
the reduction $H^{q}(C)/\ell^{K'}\to H^{q}(C)/\ell^{K}$: onto, with kernel $\ell^{K}(H^{q}(C)/\ell^{K'})$, which corresponds to $\ell^{K}H^{q}(C/\ell^{K'})$.
Conversely, write the torsion of $H^{q+1}(C)$ as $\bigoplus_j\Z/\ell^{a_j}$ and let $K\ge a=a_{j_0}$. Multiplication by $\ell$ from
$H^{q+1}(C)[\ell^{K+1}]$ to $H^{q+1}(C)[\ell^{K}]$ respects the decomposition, and on the summand $j_0$, which lies in both, it is
$\ell\colon\Z/\ell^{a}\to\Z/\ell^{a}$, not onto. So $\delta_K\circ r_{K+1,K}=\ell\circ\delta_{K+1}$ is not onto, and neither is $r_{K+1,K}$, $\delta_K$ being onto.
The example: $Z^{q}(C/\ell^{K})=\ker(\ell\colon\Z/\ell^{K}\to\Z/\ell^{K})=\ell^{K-1}\Z/\ell^{K}$, $B^{q}=0$, and the reduction carries $\ell^{K'-1}$ to $0$ modulo
$\ell^{K}$ when $K'>K$. The Dickson sentence is the case $q=0$ on $C$ (the Dickson locus has $n>0$): $H^{0}(C)=H^{0}(\Ga;\Vcoef\otimes\Z_\ell)=0$ (Claim~\ref{cl:bockstein}(1), $\Z_\ell$ flat), so
$\delta_K$ is an isomorphism onto $H^{1}(C)[\ell^{K}]=\mathbb L[\ell^{K}]$.

(1) is (5) at $q=1$, $H^{2}(C)=0$; the isomorphism $M_K\cong H^{1}(C)/\ell^{K}\cong\mathbb L/\ell^{K}$ is the left map of (5), and $r_{K',K}$ corresponds to the
reduction; every $T_p$ is induced by a cochain map of the form $\varphi$, so all of this is $T_p$-equivariant. The identification with $\mathbb L/\ell^{K}$ and
its compatibility with Lemma~\ref{lem:freepart} come from the naturality of Remark~\ref{rmk:orbifold} along $\Z_\ell\to\Z/\ell^{K}$; it is $T_p$-equivariant for the
transported $T_p$ on $\mathbb L$ by construction, and for the classical ones by Theorem~\ref{thm:classical} over $\Z_\ell$ ($6\in\Z_\ell^{\times}$).

(2) $\ell^{K}M_{K'}=\bigoplus_i\ell^{K}(\Z/\ell^{a_i})g_i$, so $M_{K'}/\ell^{K}M_{K'}=\bigoplus_i(\Z/\ell^{\min(a_i,K)})\bar g_i$ with the reduced matrices, and (1)
carries it isomorphically and equivariantly onto $M_K$; $\min(a_i,K)\ge1$. The map $\Z_\ell[X_p]\to\End(M_{K'})\to\End(M_{K'}/\ell^{K})\cong\End(M_K)$ shows that
$\mathsf B_{K'}\to\End(M_K)$ has image $\mathsf B_K$, a ring map, so $e$ is idempotent. $e'$ commutes with $\ell^{K}$, so $e'M_{K'}\cap\ell^{K}M_{K'}=\ell^{K}e'M_{K'}$
and $e'M_{K'}/\ell^{K}$ maps isomorphically onto $e(M_{K'}/\ell^{K})=eM_K$. Uniqueness: $\mathsf B_K$ is commutative (the $T_p$ commute on cohomology,
Theorem~\ref{thm:relations}(1)), and the kernel of $\mathsf B_K\to\mathsf B_1$ consists of operators $x$ with $xM_K\subseteq\ell M_K$, hence $x^{K}=0$; in a commutative
ring two idempotents congruent modulo a nil ideal are equal ($(e_1-e_2)^{3}=e_1-e_2$ for commuting idempotents, and a nilpotent $y$ with $y^{3}=y$ is $0$).

(3) By (2) the derived module and $M_K$ are isomorphic $\Z_\ell[X_p]$-modules, and blocks correspond by (2). Every listed reading is defined from the module
and the operators alone: the systems and dimensions from $M_1=M_K/\ell$; the block from its idempotent; $\mathsf A_K$ as an abstract algebra with its
generators $T_p$ is carried by any equivariant isomorphism $\phi$ to the corresponding one ($a\mapsto\phi a\phi^{-1}$), hence so are its length, exponents,
$\midl$-adic filtration and $I$; the hypotheses and conclusions of Corollary~\ref{cor:stablecert}(1),(3) are statements about these readings. Readings across levels: for
$K\le K''\le K'$ the isomorphisms $M_{K'}/\ell^{K}M_{K'}\cong M_K$ of (1) are induced by the $r_{K',K}$, and $r_{K'',K}\circ r_{K',K''}=r_{K',K}$, so they carry
the canonical projections $M_{K'}/\ell^{K''}\to M_{K'}/\ell^{K}$ (on the derived modules: the reduction of the coordinates) to $r_{K'',K}$; the derived tower and
$(M_K,r_{K'',K})$ are isomorphic towers.

(4) Both reductions modulo $\ell^{K}$ are $(\Z/\ell^{K},0)$; modulo $\ell^{K+1}$ the endomorphisms $0$ and $\ell^{K}\ne0$ are not conjugate.
\end{proof}

\begin{obligation}[One module run per cell; the levels below it derived]\label{obl:precisionreuse}
Data: per cell $(n,\ell)$ of the route of Corollary~\ref{cor:stablecert} (\file{HeckeExt/HeckeRingHeckeAlgebra}, \file{HeckeExt/HeckeSaturationDefect}): the
set $\mathrm{Lv}_{\mathrm{read}}$ of levels at which a module is read (the per-$K$ route's \code{K\_run}), the set $\mathrm{Lv}_{\mathrm{run}}$ of levels at
which the module route is run, the run modules $M_{K^{\ast}}$, the derived modules $M_{K^{\ast}}{\downarrow}K$ of Lemma~\ref{lem:precisionreduction}(2), and the
per-$K$ route (one run per level read, the route through v1.23.0-pre.2) as the oracle. Each failure \textbf{aborts} unless marked.
\begin{enumerate}
\item[(i)] \textbf{Scope.} A level is derived only in the situation of Lemma~\ref{lem:precisionreduction}(1): $\ell\ge5$, level~$1$, degree~$1$, the module route
  at the run level (native; on the digit route the derived module keeps the lowest $\min(a_i,K)$ base-$\ell$ digits of row $i$ and is reshaped to $K$ digits per
  entry), at every run level a $\tau$-table with every fill over $\Z$ (\code{used\_mod\_ell\_solve()} false; the module route refuses a mod-$\ell$ table at
  every run level, \file{Hecke/HeckeOperatorOnK.cpp}, \code{RefuseCoefficientRing}, so this is enforced, not only checked). Anything else --- $\ell\in\{2,3\}$,
  $\Heq{0}$, the cone, level $N>1$, a table with a mod-$\ell$ fill --- takes the per-$K$ route, by name. So does every failure of a run of this route, the
  intended run or a restart of (ii): a refusal, an exception (\code{ENGINE\_CHECK}, a certificate of the digit route), or a resource cap (the memory cap,
  a word or digit-route limit). A run at $K''$ can meet such a failure where the per-$K$ route, never running above $\max\mathrm{Lv}_{\mathrm{read}}$ of the block,
  would not: a crossing from the $32$- to the $64$-bit word, or into the digit route when $\ell^{K''}\ge2^{64}$. The levels still unread are then taken per $K$;
  the readings already taken are kept.
\item[(ii)] \textbf{Downward only; the restart rule.} A level $K$ is read from $M_{K^{\ast}}{\downarrow}K$ with $K\le K^{\ast}:=\max\mathrm{Lv}_{\mathrm{run}}$ so far
  ($K=K^{\ast}$: the run module itself). A level $K>K^{\ast}$ is never derived (Lemma~\ref{lem:precisionreduction}(4)): each such event triggers a new run at
  some $K''$ with $K\le K''\le K_{\mathrm{need}}$, which becomes $K^{\ast}$ (so a later level above $K''$ triggers a further run); the readings taken before the
  restart are kept (by Lemma~\ref{lem:precisionreduction}(1) they agree with those the new run would give). The choice of run levels --- $K^{\ast}=K_{\mathrm{need}}$
  in advance, or $K_{\mathrm{cls}}$ first and a restart at $K_{\mathrm{need}}$ --- is free and moves no printed value; it can change the run's word regime, its
  memory peak and whether it fails, and a failure falls under (i). $K_{\mathrm{cls}}\le K^{\ast}$ always (the free-rank reading needs it).
\item[(iii)] \textbf{The derived module, exactly.} Exponents $\min(a_i,K)$; entry $(i,j)$ reduced modulo $\ell^{\min(a_i,K)}$ (the row-exponent convention of
  \code{OpModule}); checked: the number of summands is $D$, every exponent lies in $[1,K]$, the $T_p$ commute, and at $K=1$ the result equals \code{ReduceModEll}
  entry for entry.
\item[(iv)] \textbf{Invariance of what is printed.} Every field printed or recorded from a derived level (the \code{LOEWY\_RING} and \code{SAT\_DEFECT} lines,
  the verbs' records) is a reading in the sense of Lemma~\ref{lem:precisionreduction}(3). The shipped readings are: the joint systems and their dimensions;
  lengths, cyclic exponents, $\Fl$-ranks and Hilbert functions (from Howell spans); $\varepsilon^{\mathrm{sat}}_K$, the stopping verdicts and
  \code{stable\_cert}; the monomial basis, chosen greedily in a fixed monomial order by $\Fl$-linear independence in $\mathsf A_K/\ell$ (\code{ReadDefect});
  the generators of $I$ as the reduced echelon form of $I$ in that basis, filtered against $\midl I$; the \code{SAT\_DEFECT} fields \code{generators},
  \code{monomial\_basis}, \code{quotient\_hilbert}, \code{dim\_image}, \code{tor\_routes}, \code{sat\_structure}; \code{forms=(}\dots\code{)}, which reads the $K=1$
  block only; and the \code{LOEWY\_RING} fields computed from these (\code{D}, \code{rank\_T}, \code{free\_level}, \code{free\_hilbert}, \code{hidden},
  \code{cf9}, \code{cert}) (the reason each is invariant: Lemma~\ref{lem:precisionreduction}(3)). The presentation-ordered \code{exps\_K2} is stored and printed nowhere, and
  may not be printed from a derived level. A new field read from a module must be shown to be a reading before it is printed from a derived level.
  \textbf{Recorded} (a review item, not a run).
\item[(v)] \textbf{Oracle, rollback and gate.} The per-$K$ route stays in the tree as the oracle and the rollback arm, behind a read-only environment switch
  that prints its \code{EFFECTIVE} line. A registered validation compares the whole printed output of the two routes byte for byte, timings and the fields
  of (vi) excepted, on a grid containing at least: a cell with $K_{\mathrm{stab}}\ge3$ and a plateau ($(300,7)$, \code{eps=[9,16,21,25,27,27]}); a cell with two
  levels ($(300,11)$); a Dickson-locus cell with a torsion block (\code{torsion\_free\_stable}; \code{refused:torsion\_block}); a cell with only
  \code{reduced\_rank1} blocks at $\ell\ge37$; a \code{not\_stable\_at\_KMAX} block (KMAX lowered); a restart; run levels with $\ell^{K^{\ast}}<2^{32}$ and
  $\ge2^{32}$; a cell run on the digit route (\code{ModuleRoute::Digit} with \code{native\_range}, as \code{HeckeRingLoewyValidation}'s (C3) does), derived to
  levels with $\ell^{K}<2^{64}$ and compared with the per-$K$ route's native reading there. Negative controls, each of which must fire the named check (a
  control that passes is a failure of the gate):
  \begin{itemize}\setlength{\itemsep}{0pt}
  \item a derived module keeping $a_i$ in place of $\min(a_i,K)$, at a level $K<\max_ia_i$: caught by (iii)'s exponent check (an exponent above $K$);
  \item one operator entry $(i,j)$ of $T_p$ on a derived module at a level $K\ge2$ changed by $\ell^{K-1}$, chosen in a row with $\min(a_i,K)=K$, inside a block of
    dimension $\ge2$, at an entry where $\ell^{K-1}[E_{ij},T_q]\ne0$ on the derived module for some $q\in S\setminus\{p\}$: caught by (iii)'s commutation check.
  \end{itemize}
  \textbf{Tier.} The cell $(300,7)$ costs $\approx378$\,s single-thread on the per-$K$ oracle alone (\file{ELL\_scout.md}~\S1.3), so an arm carrying it is
  \code{slow}-tier; a \code{standard}-tier arm uses a cheaper cell with $K_{\mathrm{stab}}\ge3$ and a plateau in its place, and keeps every other grid item.
  Each new registered name moves the roster, \file{docs/cluster/noctua2/EXPECT\_NAMES} and the manual's tables in the same commit.
\item[(vi)] \textbf{Record honesty} (rule~\ref{A6}). \code{K\_run} (of the \code{hecke\_ring\_algebra} record, and \code{SatCell::K\_run}) keeps the per-$K$ route's
  list, i.e.\ $\mathrm{Lv}_{\mathrm{read}}$, so the record stays byte-stable; its documented meaning changes from ``the precisions actually run'' to ``the levels read'', and
  the header comments (\file{HeckeExt/HeckeRingHeckeAlgebra.h}, \file{HeckeExt/HeckeSaturationDefect.h}) and the \code{hecke\_ring\_algebra} verb help
  (\file{HeckeExt/RegisterHeckeExtensionVerbs.cpp}) say so in the commit that lands the route. $\mathrm{Lv}_{\mathrm{run}}$ is reported in a new field of
  both cells (the \code{hecke\_ring\_algebra} record and \code{SatCell}), and \code{module\_seconds} counts the runs only. \textbf{Recorded}.
\item[(vii)] \textbf{Certificates.} \code{CertifyStable} runs on the modules the route read, derived or run, unchanged; on a derived module its
  \code{certified} rests on Lemma~\ref{lem:precisionreduction}(1) as well as on Corollary~\ref{cor:stablecert}, and the run level it was derived from is
  recorded beside it (not in \code{stable\_cert}'s value).
\end{enumerate}
\end{obligation}

\paragraph{Code map of \S\ref{sec:satprecisionsec}.} \textbf{Built at v1.23.0-pre.4} (slot-23 lever G1, owner decision 2026-10-16; \file{docs/implementation/slot23/G1\_clause\_note.md};
implemented and audited, \file{docs/implementation/slot23/G1\_report.md}, \file{G1\_impl\_audit.md} with its fixes; merged at \texttt{e4a5b2b3}; planned through
v1.23.0-pre.3). One class, \code{ringalg::LevelCache} (\file{HeckeExt/HeckeRingHeckeAlgebra.\{h,cpp\}}), is the level cache of both cell drivers,
\code{ComputeRingLoewyCell} (\file{HeckeExt/HeckeRingHeckeAlgebra.cpp}) and \code{ComputeSaturationDefectCell} (\file{HeckeExt/HeckeSaturationDefect.cpp}): its
\code{get(K)} returns the module read at level $K$, derived from the latest run $M_{K^{\ast}}$ when $K\le K^{\ast}$ and otherwise after a (re)run through
\code{RunModule} (\code{RunModuleT}, or \code{RunModuleDigit} on the digit route) --- one run of the module route per cell in the common case. Obligation~\ref{obl:precisionreuse} maps clause by clause:
(i) the scope test and the fallback in \code{LevelCache} ($\ell<5$ and \code{ModuleRoute::Hida} take the per-$K$ route by name; a failed run above the request drops the
cell to the per-$K$ route for the levels still unread, with one \code{WARNING HeckeRingHeckeAlgebra: precision reuse fell back to per\_K} line; the
validation-only \code{ScopedInjectRunFailure} exercises it); (ii) \code{RunLevelPolicy} (\code{Auto}, \code{AtRequest}, \code{Lazy}, \code{SameWord},
\code{Max}; \code{Auto} looks ahead to $K_{\mathrm{need}}$ within the word for $\ell\le$ \code{kEagerEllMax}, $=11$ since slot-23 G2-A, $13$ before, a choice that
moves only $\mathrm{Lv}_{\mathrm{run}}$; the word test \code{FitsU32Word}); (iii) \code{ReduceToLevel} (Lemma~\ref{lem:precisionreduction}(2), native and digit
forms; it never derives upwards) with \code{CheckLevelModule}, \code{CheckDerivedModule} and \code{CheckRunModule} (the checks \code{summands},
\code{exponent}, \code{commute}, \code{range}, \code{ReduceModEll}); (v) the rollback arm \code{HECKE\_PRECISION\_REUSE=0} (exactly \code{0} selects the
per-$K$ route, one run per level read; any other value warns and keeps the default), scoped as \code{ScopedPrecisionReuse}, its \code{EFFECTIVE} line
\code{PrecisionReuseEffectiveLine()}, and the gate \code{HeckePrecisionReuseValidation} (\code{standard}) and \code{HeckePrecisionReuseSlowValidation}
(\code{slow}; $(300,7)$, $(300,11)$), registered in \file{Registry/Slot23/G1.cpp} (source \file{HeckeExt/HeckePrecisionReuseValidation.cpp}); (vi) the fields
\code{Lv\_run} and \code{level\_mode} of \code{LoewyCell} and \code{SatCell} and of the \code{hecke\_ring\_algebra} record
(\file{HeckeExt/RegisterHeckeExtensionVerbs.cpp}), \code{K\_run} keeping $\mathrm{Lv}_{\mathrm{read}}$; (vii) \code{LoewyBlock::cert\_from\_cls} /
\code{cert\_from\_K}. \textbf{The lever is value-neutral}: \code{HECKE\_PRECISION\_REUSE} selects a route, never a value; the byte-identity of the two routes'
printed output, timings and the run list excepted, is the gate of (v). G1 does not change \code{compute\_module()} (\file{Hecke/HeckeOperatorOnK.cpp}); its
slot-23 G2 re-implementation is mapped in the code map of \S\ref{sec:obligations}. Through v1.23.0-pre.2 every level read was a run, and the implementation already used the
case $K=1$ of Lemma~\ref{lem:precisionreduction}(2) (\code{ReduceModEll}; ``$M_K/\ell=H^{1}(\Fl)$ at every $K$''). The module route builds a fresh
$\tau$-table per run and per $p$ (\code{RunModuleT}, \code{hida::RunCellOn}) and refuses a table with a mod-$\ell$ fill at every run level $K\ge2$
(\file{Hecke/HeckeOperatorOnK.cpp}, \code{RefuseCoefficientRing}; \code{RunModuleT} requires $K\ge2$), so Obligation~\ref{obl:precisionreuse}(i)'s $\tau$ condition
is enforced at every run level.

%======================================================================
%----------------------------------------------------------------------
\section{The \texorpdfstring{$\ell$}{l}-local isotropy model, and the closed form at \texorpdfstring{$\ell=3$}{l = 3}}\label{sec:isomodelell}
%----------------------------------------------------------------------
% NEW at v1.16.0-pre.1 (hecke-extensions slot 16, WP-12 item (b); plan docs/implementation/slot16/00-slot16-plan.md S2 "WP-12";
% source: docs/implementation/slot14/WP-Q1_small_ell_orbifold_note.md S2(b)-(d), S3.3, S4).  Part VI's subject (the integral line at
% ell in {2,3}), placed after sec:saturation, before the appendices, so that no chapter number moves.  Contract only: nothing here
% is built at this version.  The owner's standing direction of 2026-09-28 (a Moerdijk--Pronk-conformant refined model with a
% non-trivial central action, not now, and no design decision that would block it) is stated as Remark rmk:isomodellmp.

\paragraph{What this chapter fixes.} Over a $\Z_{(\ell)}$-algebra the prime-to-$\ell$ part of every simplex stabiliser is a unit, so its
column in the isotropy-resolution model of Definition~\ref{def:isomodel} can be replaced by its invariants. (i) \textbf{The
$\ell$-local model} $\Iso^{(\ell)}_{\bullet}$ (Definition~\ref{def:isomodelell}): Definition~\ref{def:isomodel} with the column of $\sigma$ the
inflation of the periodic resolution of the $\ell$-Sylow subgroup, projective and not free, and truncated to degree $0$ where that subgroup is
trivial. (ii) \textbf{What it computes and how it compares} (Proposition~\ref{prop:isomodelell}): group cohomology in every degree, cochain groups
$M^{Q_\sigma}$, comparison maps with $\Iso_{\bullet}\otimes\cring$ in both directions, the Hecke action transported, level~$N$ by Shapiro. (iii)
\textbf{At $\ell=3$} (Proposition~\ref{prop:ellthreeclosed}): one non-trivial column (the orbit of $\zeta^{2}$), the closed form
$H^{q}(\Ga;M)\cong H^{q}(\Z/3;M^{-I})$ for $q\ge2$, and the exact comparison with the shipped Bredon route: \emph{that route computes $H^{0}$
and a submodule of $H^{1}$, and misses exactly $H^{\ge1}(\Z/3;M^{-I})$}. (iv) Its place beside the Moerdijk--Pronk refined model
(Remark~\ref{rmk:isomodellmp}) and the checks (IL1)--(IL4) (Remark~\ref{rmk:isomodellcheckable}). (v) (v1.23.0-pre.12) \textbf{The certified-exponent rule on it} (Proposition~\ref{prop:certifiedeell}): one complex over
$\Z[1/\mathsf q_\ell]$ whose base changes are the complexes built over $\Z/\ell^{K}$ and $\F_{\ell'}$, Proposition~\ref{prop:certifiede} over $\Z_{(\ell)}$, and the Smith-valuation read-out.
\textbf{Notation}: those of \S\ref{sec:isomodel}; $\cring$ is a commutative $\Z_{(\ell)}$-algebra ($\Z_{(\ell)}$, $\Z/\ell^{e}$, $\Fl$, \dots), $M$ an
$\cring[\Ga]$-module. $P_\sigma$, $Q_\sigma$, $\bar g_\sigma$, $F'^{\sigma}$, $G'_{s,t}$, $\Iso^{(\ell)}$, $\cmpI'$, $\Phi_{\ell}$, $\Psi_{\ell}$ are new symbols of this
chapter and mean nothing else (rule~\ref{A5}); $\Phi_\ell$ is not the cone map $\Phi_g$, nor $\Psi_\ell$ the map $\Psi$ of
Theorem~\ref{thm:explicit_complex_equiv}. $\mathsf q_\ell$, $\mathsf O_\ell$, $C^{\bullet}_\ell$, $F^{q}$, $P_q$, $\delta^{q}$ are new symbols of Proposition~\ref{prop:certifiedeell} (v1.23.0-pre.12), and so are $b_j$ and $s$ (v1.23.0-pre.13); they mean
nothing else there.

\begin{definition}[The $\ell$-local isotropy model]\label{def:isomodelell}
Assume \ref{H4} and \ref{B1}--\ref{B3}; let $\ell$ be a prime and $\cring$ a $\Z_{(\ell)}$-algebra. For $\sigma\in\Sigma_s$ (Definition~\ref{def:isomodel}) the
stabiliser $\Ga_\sigma=\langle g_\sigma\rangle$ is cyclic of order $m_\sigma$; let $P_\sigma$ be its $\ell$-Sylow subgroup, $Q_\sigma$ its subgroup of order prime
to $\ell$ (so $\Ga_\sigma=P_\sigma\times Q_\sigma$ and $\abs{Q_\sigma}\in\cring^{\times}$), $\bar g_\sigma$ the image of $g_\sigma$ in $\Ga_\sigma/Q_\sigma\cong P_\sigma$,
and $e_\sigma:=\abs{Q_\sigma}^{-1}\sum_{w\in Q_\sigma}w$. The \textbf{column} $F'^{\sigma}_{\bullet}$ is: if $P_\sigma\ne1$, the periodic complex
$F'^{\sigma}_t:=\cring[\Ga_\sigma/Q_\sigma]=\cring[\Ga_\sigma]e_\sigma$ ($t\ge0$) with differential $\bar g_\sigma-1$ from odd $t$ and the norm $\sum_{x\in\Ga_\sigma/Q_\sigma}x$
from even $t\ge2$; if $P_\sigma=1$, $F'^{\sigma}:=\cring$ in degree $0$ alone. Put $G'_{s,t}:=\bigoplus_{\sigma\in\Sigma_s}\cring[\Ga]\otimes_{\cring[\Ga_\sigma]}F'^{\sigma}_t$.
An \textbf{$\ell$-local isotropy model} is a complex of $\cring[\Ga]$-modules $\Iso^{(\ell)}_{\bullet}$ with $\Iso^{(\ell)}_q=\bigoplus_{s+t=q,\ 0\le s\le2}G'_{s,t}$ and
$d=d'^{(0)}+d'^{(1)}+d'^{(2)}$ of the shape of Definition~\ref{def:isomodel} ($d'^{(0)}=(-1)^{s}(\id\otimes\partial^{F'})$, $d'^{(r)}\colon G'_{s,t}\to G'_{s-r,t+r-1}$), with
the augmentation $\cmpI'\colon\Iso^{(\ell)}_{\bullet}\to C_{\bullet}(\Xbulk)\otimes\cring$ on $G'_{s,0}$ ($\cring[\Ga]\otimes_{\cring[\Ga_\sigma]}F'^{\sigma}_0\to\cring[\Ga/\Ga_\sigma]$) and $0$
on $t\ge1$, satisfying \ref{IM1} and \ref{IM2} (the sub-sum over $\Sigma^{\partial}_0\sqcup\Sigma^{\partial}_1$ a subcomplex mapping into $C_{\bullet}(\Xbdry)\otimes\cring$).
Write $C^{\bullet}_{\Iso^{(\ell)}}(M):=\Hom_{\cring[\Ga]}(\Iso^{(\ell)}_{\bullet},M)$.
\end{definition}

\begin{proposition}[The $\ell$-local model: existence, what it computes, the comparison, the Hecke action]\label{prop:isomodelell}
\PROVED{} In the setting of Definition~\ref{def:isomodelell}:
\begin{enumerate}
\item \textbf{(Existence.)} $\Iso^{(\ell)}$ exists; $d'^{(1)}$, $d'^{(2)}$ are produced degree by degree by lifts along $\cring$-linear contracting homotopies of the
  augmented columns, boundary summands first --- the procedure of Proposition~\ref{prop:isomodel}(1).
\item \textbf{(A projective resolution.)} Each $G'_{s,t}\cong\bigoplus_{\sigma}\cring[\Ga/Q_\sigma]$ is a projective, in general not free, $\cring[\Ga]$-module;
  $\cmpI'$ is a surjective quasi-isomorphism, so $\aug\circ\cmpI'\colon\Iso^{(\ell)}_{\bullet}\to\cring$ is a projective resolution of $\cring$ over $\cring[\Ga]$;
  likewise on the boundary.
\item \textbf{(What it computes.)} $C^{q}_{\Iso^{(\ell)}}(M)=\bigoplus_{s+t=q}\bigoplus_{\sigma\in\Sigma_s,\ (t=0\text{ or }P_\sigma\ne1)}M^{Q_\sigma}$, and
  $H^{q}(C^{\bullet}_{\Iso^{(\ell)}}(M))=H^{q}(\Ga;M)$ for \textbf{every} $q\ge0$ and every $\cring[\Ga]$-module $M$; on the boundary $H^{q}(\Gbdry;M)$. In
  particular, for $M=\Ind_{\Ga_0(N)}^{\Ga}\VR$ (Definition~\ref{def:levelNmodel}) or $\Ind_{\Ga_0(N)}^{\Ga}\VR(\chi)$ (Definition~\ref{def:nebmodel}) it computes
  $H^{q}(\Ga_0(N);\VR)$, resp.\ $H^{q}(\Ga_0(N);\VR(\chi))$, in every degree (Shapiro's lemma, \STD{\cite[Ch.~III.6.2]{brown}}).
\item \textbf{(Comparison with the free model.)} There are $\cring[\Ga]$-chain maps $\Phi_\ell\colon\Iso_{\bullet}\otimes\cring\to\Iso^{(\ell)}_{\bullet}$ and
  $\Psi_\ell\colon\Iso^{(\ell)}_{\bullet}\to\Iso_{\bullet}\otimes\cring$ with $\cmpI'\Phi_\ell=\cmpI\otimes\cring$ and $(\cmpI\otimes\cring)\Psi_\ell=\cmpI'$ (strict lifts over
  $C_{\bullet}(\Xbulk)\otimes\cring$, boundary to boundary), unique up to homotopy over $C_{\bullet}(\Xbulk)\otimes\cring$, mutually inverse up to homotopy; so
  $\Phi_\ell^{\ast}\colon C^{\bullet}_{\Iso^{(\ell)}}(M)\to\CI(M)$ and $\Psi_\ell^{\ast}$ are mutually inverse quasi-isomorphisms, $\CI(M):=\Hom_{\Z[\Ga]}(\Iso_{\bullet},M)$. Column by
  column, the map $\varphi^{\sigma}_t:=\abs{Q_\sigma}^{\lfloor t/2\rfloor}\pi_\sigma\colon\cring[\Ga_\sigma]\to\cring[\Ga_\sigma/Q_\sigma]$ ($\pi_\sigma$ the projection; $0$ for
  $t\ge1$ when $P_\sigma=1$) is a chain map $F^{\sigma}\otimes\cring\to F'^{\sigma}$ over $\id_\cring$, defined over $\Z$ already (no inverse of $\abs{Q_\sigma}$).
\item \textbf{(The Hecke action.)} For $g\in\De$ and $T_{\Ga g\Ga}$ any cochain-level Hecke operator on $\CI(M)$ inducing $[\Ga g\Ga]$ --- the strict lift of
  Proposition~\ref{prop:strictlift}, $p\in\{2,3,5\}$, and its extensions in Proposition~\ref{prop:relationsmodel} --- the cochain map
  $T'_{\Ga g\Ga}:=\Psi_\ell^{\ast}\circ T_{\Ga g\Ga}\circ\Phi_\ell^{\ast}$ induces $[\Ga g\Ga]$ on $H^{\bullet}(\Ga;M)$, independently of the choices of
  $\Phi_\ell$, $\Psi_\ell$; likewise for the pair and the cone (the boundary stabilisers are $\{\pm I\}$, \ref{B3}), and $U_\ell$ needs nothing new (the
  $\De$-module structure of $M$ is unchanged, \ref{H2}).
\end{enumerate}
\end{proposition}

\begin{proof}
\emph{The columns.} $e_\sigma$ is an idempotent of $\cring[\Ga_\sigma]$ ($\abs{Q_\sigma}\in\cring^{\times}$) with $\cring[\Ga_\sigma]e_\sigma\cong\cring[\Ga_\sigma/Q_\sigma]$, so each
$F'^{\sigma}_t$ is a projective $\cring[\Ga_\sigma]$-module. $F'^{\sigma}\to\cring$ is exact: for $P_\sigma\ne1$ it is the periodic resolution of $\cring$ over the cyclic group
$\Ga_\sigma/Q_\sigma$ (the argument of Proposition~\ref{prop:isomodel}'s proof, over $\cring$), and for $P_\sigma=1$ it is $\cring\to\cring$. Bounded below and consisting of free
$\cring$-modules, it is split exact over $\cring$, so it has an $\cring$-linear contracting homotopy. $\cring[\Ga]$ is free as a right $\cring[\Ga_\sigma]$-module, so
$\cring[\Ga]\otimes_{\cring[\Ga_\sigma]}-$ is exact and preserves projectives, and $G'_{s,\bullet}\to C_s(\Xbulk)\otimes\cring$ is a projective resolution, with
$\cring[\Ga]\otimes_{\cring[\Ga_\sigma]}\cring[\Ga_\sigma/Q_\sigma]\cong\cring[\Ga/Q_\sigma]$.
(1) The proof of Proposition~\ref{prop:isomodel}(1) uses only that the column resolutions are projective with $\cring$-linear contracting homotopies and the comparison theorem
\STD{\cite[Thm.~2.2.6]{weibel}}, which needs projective sources; it applies with one change: the generator $1\otimes e_\sigma$ of $\cring[\Ga]e_\sigma$ is
$Q_\sigma$-fixed, so every lift produced by the $\cring$-linear contracting homotopies is replaced by its $Q_\sigma$-average, $y\mapsto e_\sigma y$ ($\abs{Q_\sigma}\in\cring^{\times}$).
The average still lifts the same element, because the element lifted is the image of a $Q_\sigma$-fixed element under an $\cring[\Ga]$-map and the
differentials commute with $e_\sigma$; this is the identity $\Hom_{\cring[\Ga]}(\cring[\Ga]e_\sigma,Y)=e_\sigma Y$ (v1.16.0-pre.3, WP-14 referee C-1).
The average is kept in this procedure for uniformity and as the general statement; on $\Xbulk$ (and on any $\Ga$-equivariant refinement of it) it changes
nothing on $d$ itself: a lift of $d'^{(1)}$ or $d'^{(2)}$ starts at a source of dimension $s\ge1$, whose stabiliser fixes a segment of $\Hyp$ pointwise (no
inversions, \ref{H4}(i)) and is therefore $\{\pm I\}$, a non-central element of $\Ga$ fixing only isolated points of $\Hyp$; so $Q_\sigma\subseteq\{\pm I\}$ ($=\{\pm I\}$
at $\ell=3$, trivial at $\ell=2$) is central and contained in every target's $Q_\tau$, and $e_\sigma y=y$ for every $y\in\cring[\Ga]e_\tau$. Where the average
is load-bearing is the comparison of (4): the source generators of $\Psi_\ell$ are $Q_x$-fixed while the free target $\Iso\otimes\cring$ is not
(v1.16.0, slot-16 WP-16 audit F-1). (2) The filtration and five-lemma argument of Proposition~\ref{prop:isomodel}(2), with
$C_{\bullet}(\Xbulk)\otimes\cring$ exact over $\cring$ (split over $\Z$, \ref{H4}(ii)). (3) $\Hom_{\cring[\Ga]}(\cring[\Ga/Q_\sigma],M)=M^{Q_\sigma}$, and
$H^{q}(\Hom_{\cring[\Ga]}(P_{\bullet},M))=\Ext^{q}_{\cring[\Ga]}(\cring,M)=H^{q}(\Ga;M)$ for any projective resolution $P$ ($\cring\otimes_\Z$ of a free $\Z[\Ga]$-resolution is a
free $\cring[\Ga]$-resolution). The boundary as in Proposition~\ref{prop:isomodel}(3).
(4) $\Iso\otimes\cring$ is free and $\Iso^{(\ell)}$ projective over $\cring[\Ga]$, both bounded below, and both augmentations onto $C_{\bullet}(\Xbulk)\otimes\cring$ are surjective
quasi-isomorphisms; Lemma~\ref{lem:strictlifting}'s proof gives the strict lifts --- with a projective source, the step ``choose $y_0$'' and the step ``$z=\partial k$'' are
lifts of a map out of a projective module through a surjection --- constructively, $y_0:=e_\sigma y_0'$ and $k:=e_\sigma k'$ for any $y_0'$, $k'$ produced as in
Lemma~\ref{lem:strictlifting}'s proof (C-1) ---, onto $(\ker p)$'s boundaries in the second case ($\ker p$ is acyclic) --- boundary first, as there. Uniqueness and
inverse up to homotopy: Proposition~\ref{prop:isomodel}(5)'s argument (a chain map into an acyclic kernel from a projective complex is null-homotopic through that kernel).
The column maps: on odd $t$, $\pi_\sigma(g_\sigma-1)=(\bar g_\sigma-1)\pi_\sigma$, so $\varphi_{t-1}=\varphi_t$ is needed and holds ($\lfloor(t-1)/2\rfloor=\lfloor t/2\rfloor$ for odd $t$);
on even $t\ge2$, $\pi_\sigma(\sum_{u\in\Ga_\sigma}u)=\abs{Q_\sigma}\sum_{x\in\Ga_\sigma/Q_\sigma}x$, so $\varphi_{t-1}\circ N=N'\circ\varphi_t$ needs $\varphi_t=\abs{Q_\sigma}\varphi_{t-1}$,
which is the exponent; augmentations agree in degree $0$. For $P_\sigma=1$ the target has $\bar g_\sigma=1$, and the degree-$0$ map composed with $g_\sigma-1$ is $0$.
(5) The Hecke operator on $H^{\bullet}(\Ga;M)$ is $\cor\circ(\conj_g)_{\ast}\circ\res$, defined independently of the projective resolution \STD{\cite[Ch.~III.9--10]{brown}};
$\Phi^{\ast}_\ell$ and $\Psi^{\ast}_\ell$ are the comparison isomorphisms between the resolutions, and conjugating a cochain map inducing it by them induces it again. The pair and the
cone: (4) holds on the boundary sub-sum, so the same conjugation applies to the pair datum of Proposition~\ref{prop:strictlift}.
\end{proof}

\begin{proposition}[At $\ell=3$: one column, the closed form, and the comparison with the shipped Bredon route]\label{prop:ellthreeclosed}
\PROVED{} Let $\ell=3$, $\cring$ a $\Z_{(3)}$-algebra, $M$ an $\cring[\Ga]$-module, and $\Iso^{(3)}$ on the model of Definition~\ref{def:model7}.
\begin{enumerate}
\item \textbf{(One column.)} $P_\sigma\ne1$ for exactly one representative, $\sigma=v_4$ ($\Ga_{v_4}=\langle ST\rangle$, $P_{v_4}=\langle(ST)^{2}\rangle\cong\Z/3$,
  $Q_{v_4}=\{\pm I\}$); every other column is $\cring[\Ga/\Ga_\sigma]$ in degree $0$ ($\Ga_{v_6}=\langle S\rangle$ has order $4$, the rest $\{\pm I\}$). So
  $C^{q}_{\Iso^{(3)}}(M)$ differs from the shipped (Bredon) cochain group $\Cq{q}$ (Proposition~\ref{prop:bredon}) only at $v_4$: degree $0$ is
  $3\cdot M^{-I}\oplus M^{\langle S\rangle}\oplus M^{-I}$ ($M^{-I}$ at $v_4$ in place of $M^{\langle U\rangle}$), degree $1$ is $9\cdot M^{-I}\oplus M^{-I}$, degree $2$ is
  $5\cdot M^{-I}\oplus M^{-I}$, and every degree $q\ge3$ is one $M^{-I}$.
  The boundary sub-model is the Bredon boundary complex ($P_\sigma=1$ on $\Xbdry$, \ref{B3}).
\item \textbf{(The closed form.)} With $U=ST$ and $\langle U\rangle=\{\pm I\}\times\langle U^{2}\rangle$,
  \[
    H^{q}(\Ga;M)\;\cong\;H^{q}(\langle U\rangle;M)\;=\;H^{q}(\langle U^{2}\rangle;M^{-I})\qquad(q\ge2),
  \]
  which is $(M^{-I})^{U^{2}}/N(M^{-I})$ for even $q$ and $\ker N/(U^{2}-1)M^{-I}$ for odd $q$, $N=1+U^{2}+U^{4}$; and there is an exact sequence
  \[
    0\to M^{-I}\big/\bigl(M^{\langle S\rangle}+M^{\langle U\rangle}\bigr)\xrightarrow{\ \iota\ }H^{1}(\Ga;M)\to H^{1}(\langle U^{2}\rangle;M^{-I})\to0 ,
  \]
  with $H^{0}(\Ga;M)=M^{\langle S\rangle}\cap M^{\langle U\rangle}=M^{\Ga}$.
\item \textbf{(The shipped Bredon route.)} The complex the shipped pipeline computes over $\cring$, $\Ceq(K,\Rel;L)$ with fibre $M$ (Proposition~\ref{prop:bredon}(2),(3):
  Bredon cohomology for the finite subgroups with fixed-point coefficients), has $H^{0}=M^{\Ga}$, $H^{1}=M^{-I}/(M^{\langle S\rangle}+M^{\langle U\rangle})$ and
  $H^{q}=0$ for $q\ge2$; the map $\cmpI^{\ast}\colon\Ceq\to\CI(M)$ of Proposition~\ref{prop:isomodel}(4) induces the identity on $H^{0}$ and, on $H^{1}$, an
  injection whose cokernel is $H^{1}(\langle U^{2}\rangle;M^{-I})$ (the same terms as (2)'s sequence; that its image is the image of $\iota$ is not asserted). So at $\ell=3$ \textbf{the Bredon route is exact in degree $0$, injective in degree $1$ with cokernel $H^{1}(\Z/3;M^{-I})$, and misses
  $H^{q}(\Z/3;M^{-I})$ in every degree $q\ge2$} --- nothing else.
\item \textbf{(Parity.)} \LEAN{thm:lv3-tree-negone} (only the vanishing, and only for $\cring=K$ a field with $2\ne0$ and $M$ finite-dimensional: $H^{q}(\Ga;M)=0$
  for every $q$, proved on the complex of Remark~\ref{rmk:treeconeres}, not through the terms of (1); instantiated at $M=\Sym^{n}\F_3^{2}$, $n$ odd)
  If $-I$ acts on $M$ by $-1$ (e.g.\ $\VR$ with $n$ odd, or $\VR(\chi)$ with $\chi(-1)\ne(-1)^{n}$), every term of (1) vanishes and
  $H^{\bullet}(\Ga;M)=0$; this needs only $2\in\cring^{\times}$ (Proposition~\ref{prop:neborbit}(2)'s argument), not $6$.
\end{enumerate}
\end{proposition}

\begin{proof}
(1) The stabilisers are those of \ref{H4}(iv) and Lemma~\ref{lem:model7coincidences}(1) on the representatives $\Sigma_s$ of Definition~\ref{def:isomodel}
($v_5$ is not a representative); $\langle ST\rangle$ has order $6=2\cdot3$, $\langle S\rangle$ order $4$, $\{\pm I\}$ order $2$. The ranks are read off
Proposition~\ref{prop:isomodelell}(3) with $M^{Q_\sigma}=M^{\Ga_\sigma}$ where $P_\sigma=1$, and $M^{Q_{v_4}}=M^{-I}$.
(2) Claim~\ref{cl:amalgam}(2), \eqref{eq:mvamalgam}, with $H^{q}(\langle S\rangle;M)=H^{q}(\langle-I\rangle;M)=0$ for $q\ge1$ ($4$ and $2$ are units;
\STD{\cite[Ch.~III.10.2]{brown}}): for $q\ge2$ the sequence gives $H^{q}(\Ga;M)\cong H^{q}(\langle U\rangle;M)$, and at $q=1$ it reads $M^{S}\oplus M^{U}\to M^{-I}\to
H^{1}(\Ga;M)\to H^{1}(\langle U\rangle;M)\to0$. The Lyndon--Hochschild--Serre sequence of $\{\pm I\}\trianglelefteq\langle U\rangle$ collapses ($2$ a unit,
\STD{\cite[Ch.~VII.6]{brown}}): $H^{q}(\langle U\rangle;M)=H^{q}(\langle U^{2}\rangle;M^{-I})$; the cyclic group $\langle U^{2}\rangle$ of order $3$ has the periodic resolution.
(3) The Bredon complex is Proposition~\ref{prop:bredon}(2). The map: filter $\CI(M)$ and $\Ceq$ by $s$ as in the proof of Proposition~\ref{prop:isomodel}(4); the
$E_1$-page of $\CI(M)$ is $E_1^{s,t}=\bigoplus_{\sigma\in\Sigma_s}H^{t}(\Ga_\sigma;M)$, its row $t=0$ is $\Ceq$, and $\cmpI^{\ast}$ is the inclusion of that row, so on cohomology
it is the edge map $E_2^{q,0}\to H^{q}$. Over a $\Z_{(3)}$-algebra $E_1^{s,t}=0$ for $t\ge1$ except $E_1^{0,t}=H^{t}(\langle U\rangle;M)$ at $\sigma=v_4$. The differentials
out of $E_r^{0,t}$, $t\ge1$, land in $E_r^{r,t-r+1}$, which is $0$ for $t-r+1\ge1$ (no column $s\ge1$ has higher cohomology) and is $E_2^{2,0}=H^{2}(\Ceq)=0$ for $r=t+1=2$
(Proposition~\ref{prop:bredon}(2)), and $E_r^{r,t-r+1}=0$ for $r\ge3$ ($s\le2$). So the sequence degenerates at $E_2$: $0\to E_2^{1,0}\to H^{1}\to E_2^{0,1}\to0$, $H^{0}=E_2^{0,0}$,
$H^{q}=E_2^{0,q}$ for $q\ge2$; $E_2^{1,0}=H^{1}(\Ceq)$, the edge map $E_2^{1,0}\to H^{1}$ is injective, and $E_2^{0,1}=H^{1}(\langle U\rangle;M)$.
(4) $-I$ is central and in every $\Ga_\sigma$; $M^{-I}=0$ and $M^{Q_\sigma}\subseteq M^{-I}=0$ when $-I$ acts by $-1$ and $2$ is a unit.
\end{proof}

\begin{remark}[Its place beside the Bredon route, the free model, and the Moerdijk--Pronk refined model]\label{rmk:isomodellmp}
\PROVED{} (1) \textbf{The Bredon route is not replaced.} What $\F_\ell$ and $\Z/\ell^{e}$ \emph{name} at $\ell\in\{2,3\}$ is Definition~\ref{def:ringsemantics}(1) and \Espec{sec:ring}:
the shipped route computes the Bredon object, and switching the default object to group cohomology would change published values (golden-gated, an owner decision).
$\Iso^{(\ell)}$ is a second, smaller computational model of the \emph{group} cohomology that $\Iso$ already computes (Proposition~\ref{prop:isomodel}(3)); $\Iso$ is its oracle in every
degree. Proposition~\ref{prop:ellthreeclosed}(3) is the exact statement of the difference at $\ell=3$.
(2) \textbf{Size.} At $\ell=3$ the degrees $q\ge3$ cost one copy of $M^{-I}$ instead of $19$ copies of $M$ (Proposition~\ref{prop:isomodel}(2)); at $\ell=2$ only the corner column
shrinks ($P_{v_4}=\{\pm I\}$, $Q_{v_4}=\langle U^{2}\rangle$, $M^{Q}$ of rank about $\frac13$), every simplex keeping a $\Z/2$ column --- the generic central isotropy, which no
$2$-local reduction removes ($\Ga\to\mathrm{PSL}_2(\Z)$ does not split, Lemma~\ref{lem:bulkperiodic}'s class of order $12$). So $\ell=2$ is specified by
Definition~\ref{def:isomodelell} and gains little; $\ell=3$ is where it pays.
(3) \textbf{The MP constraint} (owner, 2026-09-28). Definition~\ref{def:isomodelell} and Proposition~\ref{prop:isomodelell} use of the model only \ref{H4}, \ref{B1}--\ref{B3} and the
stabiliser inventory: no vertex count, no edge list, no special role of $v_4$ beyond ``the representatives whose stabiliser has order divisible by $\ell$''. A
$\Ga$-equivariant subdivision $Y'$ of $\Xbulk$ --- in particular the Moerdijk--Pronk-conformant refinement the owner's direction names (subdividing the edge from $\zeta^{2}$ to $i$
and the edge from $\zeta^{2}$ towards the cusp) --- preserves every point stabiliser (Remark~\ref{rmk:orbifoldrefine}(1)); the new vertices have stabiliser $\{\pm I\}$,
so the simplices with $\abs{P_\sigma}>2$ (i.e.\ $P_\sigma\supsetneq\{\pm I\}$) are again the orbit of $\zeta^{2}$ at $\ell=3$, and of $i$ at $\ell=2$
(at $\ell=2$ every simplex has $P_\sigma\supseteq\{\pm I\}$; WP-14 referee C-7), and the model on $Y'$ is homotopy equivalent to the one on $\Xbulk$ over
$\cring$ (Proposition~\ref{prop:isomodelell}(4)'s argument with $C_{\bullet}(Y')\simeq C_{\bullet}(\Xbulk)$, Remark~\ref{rmk:cafive}). The centre is \textbf{not} trivialised: $-I$
lies in every $\Ga_\sigma$, enters $Q_\sigma$ at $\ell\ne2$ as a unit-order factor whose invariants are taken (\emph{not} deleted), and at $\ell=2$ keeps its column; a
coefficient module on which $-I$ acts by a scalar other than $1$ (odd $n$, a character with $\chi(-1)=-1$) is admissible input, which is the ``non-trivial action of the
centre everywhere'' the refined model is to carry. Nothing here is Bredon-only: the Bredon object appears only as the $t=0$ row and in the comparison (3).
(v1.22.0-pre.5: two refinements are to be kept apart. Slot~17's synthetic descriptor, written $\mathsf R_{\zeta}$ in \S\ref{sec:mpmodel}, subdivides all three edge orbits at
$\zeta^{2}$; the Moerdijk--Pronk-conformant refinement of the owner's ruling Q2 is $\mathsf R_{\Z}$, which subdivides only the edge orbit of $v_4v_6$
(Definition~\ref{def:mpRZ}). Both are adapted; the edge from $\zeta^{2}$ towards the cusp is forced by none of Moerdijk--Pronk's rules, and goodness forces, among others,
the \emph{other} edges from $\zeta^{2}$ and $\zeta$ towards the horocycle, $v_1v_4$ and $v_2v_5$ (Remark~\ref{rmk:mprulesread}, Table~\ref{tab:mprules}).)
\end{remark}

\begin{remark}[What is checkable: (IL1)--(IL4)]\label{rmk:isomodellcheckable}
\PROVED{} \COMP\textsuperscript{a} Committed at v1.16.0-pre.1, before any code (rule~20); a failing \PROVED{} item stops the package. Cells over $\Z/3^{e}$ and $\F_3$, $n\le200$,
$q\le q_{\max}$.
\begin{enumerate}[label=(IL\arabic*)]
\item\label{IL1} \PROVED{} (Proposition~\ref{prop:isomodelell}(1),(3),(4)) the built $\Iso^{(3)}$ --- every lift averaged by $e_\sigma$ (C-1) --- satisfies
  $d\circ d=0$, \ref{IM1} and \ref{IM2}; and $\Iso^{(3)}$ and $\Iso$ give the same lengths and invariant factors of $H^{q}(\Ga;\Sym^{n}(\Z/3^{e})^{2})$ for
  $q\le q_{\max}$, and those of [IB] Appendix~A for $n\le20$, $q\le5$ (\code{HeckeFreeModelValidation}'s oracle).
\item\label{IL2} \PROVED{} (Proposition~\ref{prop:ellthreeclosed}(2)) the closed form: for $q\ge2$, $H^{q}(\Ga;\Sym^{n}\Z_{(3)}^{2})\cong H^{q}(\Z/3;\Sym^{n}\Z_{(3)}^{2})$ ($n$ even), $2$-periodic
  in $q$. \COMP\textsuperscript{a} (slot-14 WP-Q1's \file{z3check.py}, kept in the slot-14 scratch, not tracked; \file{docs/implementation/slot14/WP-Q1\_small\_ell\_orbifold\_note.md}
  \S3.3): the $3$-parts of $H^{2}$ and $H^{3}$ in [IB] Appendix~A agree with $H^{2}$, $H^{3}$ of $\Z/3$ on $22$ / $22$ cells $n=0,2,\dots,20$ ($\Z/3$ in $H^{2}$ exactly at
  $n\equiv0\pmod6$, in $H^{3}$ exactly at $n\equiv4\pmod6$).
\item\label{IL3} \PROVED{} (Proposition~\ref{prop:ellthreeclosed}(3)) over $\F_3$: $\dim H^{1}(\Ga;\Sym^{n}\F_3^{2})=\dim H^{1}_{\mathrm{Bredon}}+\dim H^{1}(\Z/3;\Sym^{n}\F_3^{2})$ for
  $n$ even, the first summand being the shipped pipeline's $\dim\Heq{1}$ (\code{HeckeBredonTreeValidation}'s cells), and $H^{0}$ equal; the same at every level $N$ with
  $M=\Ind_{\Ga_0(N)}^{\Ga}\VR$. \COMP\textsuperscript{a}: [IB] Appendix~A's $H^{1}$ defect (group minus model) has $3$-part $\Z/3$ exactly at $n=4,10,16$, which is
  $H^{1}(\Z/3;M)$ (the same note, \S3.3). \COMP\textsuperscript{a} \emph{partial} (v1.16.0 at level~$1$, slot-16 WP-16; v1.17.0-pre.5 at level~$N$, slot-17 WP-7,
  \code{HeckeIsotropyModelEllValidation} arm (IL3N)): held at level~$1$ over $\F_3$, $n\le40$, with the map-level comparison over $\F_3$ and $\Z/9$; and at level $N$
  with $M=\Ind_{\Ga_0(N)}^{\Ga}\VR$ built from the shipped \code{Gamma0LevelN} table (\S\ref{sec:levelN}; no second coset construction), \emph{against the shipped
  level-$N$ model's $\Heq{1}$ over $\F_3$ only} (Definition~\ref{def:levelNmodel}(a), the cochain level), $N\in\{1,2,3,4,5,7,11,13\}$, $n\in\{0,1,2,4,6\}$ ($N\le16$,
  $n\le6$ under \code{HECKE\_ISOMODEL\_ELL\_FULL=1}); over $\Z/9$ and $\Z/27$ the level-$N$ cells compare $\Iso^{(3)}$ with $\Iso$ (lengths) and the closed form, not the
  pipeline, which is mod~$p$; the $\Z/3$ summand is non-zero at $N=7$. \COMP\textsuperscript{a} (v1.18.0-pre.2, slot-18 WP-7, arm (IL3Q), code map item~(vi)): over
  $\F_9$ and, in lengths, over $\Z/9$ against the shipped level-$N$ model's $\Heq{1}$ over $\Z/9$, $12$ cells. \textbf{Tag moved} from \COMP\textsuperscript{a}
  \emph{partial} to \COMP\textsuperscript{a} at v1.18.0-pre.2 (pre.1 through v1.17.0: ``\emph{partial} until the next referee rules''), on the slot-18 referee F1's
  ruling C-4 (``move both at the close on WP-7's \ref{IL5} + (vi) pass'', \file{docs/implementation/slot18/WP-3\_F1\_referee\_report.md}) and WP-7's pass
  (\file{docs/implementation/slot18/WP-7\_report.md}; \code{--all} $156/156$).
\item\label{IL4} \PROVED{} (Proposition~\ref{prop:isomodelell}(5)) the transported $T'_2$, $T'_5$ (and $T'_3=U_3$) on $\Iso^{(3)}$ have the characteristic polynomials of $T_2$, $T_5$, $T_3$
  on $\Iso$ (Proposition~\ref{prop:strictlift}, \code{HeckeIntegralActionValidation}), degree by degree, $q\le q_{\max}$. \COMP\textsuperscript{a} \emph{partial} (v1.16.0
  over $\F_3$, slot-16 WP-16; v1.17.0-pre.5, slot-17 WP-7, arms (IL4e), (PC)): held over $\F_3$, even $n\le12$, $q\le5$, $p\in\{2,3,5\}$; over $\Z/3^{e}$, $e\in\{2,3\}$,
  even $n\le12$ resp.\ $n\le8$ ($n\le20$ / $16$ under the full grid), in lengths and invariant factors of $H^{q}$ and of $\operatorname{im}(T'_p-a)$, $a\in\{0,1,-1\}$
  (basis-free), and in characteristic polynomials where $H^{q}$ is free; and the pair and cone half of Proposition~\ref{prop:isomodelell}(5): $\Phi^{\ast}_\ell$,
  $\Psi^{\ast}_\ell$ and $T_p$ commute with the restriction to the boundary sub-model, and the transported $T'_p$ on the mapping cone of the restriction and on the
  boundary has the characteristic polynomials (over $\F_3$) and image invariant factors (over $\Z/9$) of $T_p$ on $\Iso$. \textbf{Tag moved} from
  \COMP\textsuperscript{a} \emph{partial} to \COMP\textsuperscript{a} at v1.18.0-pre.2, with \ref{IL3} and on the same ground (the slot-18 referee F1's C-4, WP-7's
  pass of \ref{IL5}, which re-checks the transported $T'_p$ against $T_p$ on the free model over $\Z/\ell^{e}$ and $\F_\ell$, and of~(vi)).
\item\label{IL5} \PROVED{} (Proposition~\ref{prop:strictlift}(3)'s argument with Proposition~\ref{prop:isomodelell}(4),(5) in place of
  Remark~\ref{rmk:borelfree} --- the source is projective, not free --- given that $\Iso^{(\ell)}(Y')$ is an $\ell$-local isotropy model over $Y'$ and that the
  built map satisfies (T1)--(T3); committed at v1.18.0-pre.1, before slot-18 WP-7's code, rule~20; route and hypothesis repaired at v1.18.0-pre.2, the slot-18
  referee F1's M-1: pre.1 read ``Proposition~\ref{prop:strictlift}(3)'s argument and Proposition~\ref{prop:isomodelell}(5)'', whose last step passes through
  Remark~\ref{rmk:borelfree}, a statement about a model \emph{free} over $\Z[\Ga]$) \textbf{The strict action on a refined descriptor.} Let $Y'$ be a synthetic
  refined descriptor as in Remark~\ref{rmk:isomodellmp}(3) (a $\Ga$-equivariant subdivision of $\Xbulk$, given as data), and suppose $\Iso^{(\ell)}(Y')$ is an
  $\ell$-local isotropy model over $Y'$ in the sense of Definition~\ref{def:isomodelell} with $Y'$ in place of $\Xbulk$ (it exists by
  Proposition~\ref{prop:isomodelell}(1), Remark~\ref{rmk:isomodellmp}(3): projective, in general not free, over $\cring[\Ga]$, with its augmentation onto
  $C_{\bullet}(Y')\otimes\cring$ a surjective quasi-isomorphism). Let $g=\mathrm{diag}(1,p)$, $\Ga''=\Ga\cap g\Ga g^{-1}$, and $\tau'$ an additive map on the
  resolution side over $Y'$ satisfying (T1) (over $\Ga''$), (T2) and (T3) of Definition~\ref{def:tau} with that model in place of $C_{\bullet}(\Ku)$. (How
  slot-18 WP-7 obtains $\tau'$ --- a strict lift in the sense of Lemma~\ref{lem:strictlifting}, certified by the strict identity relative to the comparison with
  $\Xbulk$ as in Remark~\ref{rmk:strictliftbuilt} --- is a certificate of the built map, not a hypothesis of this item.) Then its coset transport $T_p^{Y'}$ and the transported $T'_p$ of Proposition~\ref{prop:isomodelell}(5) (through the
  comparison of the models over $Y'$ and $\Xbulk$) induce the same operator $\cor\circ(\conj_g)_{\ast}\circ\res$ on $H^{q}(\Ga;\VR)$ (\emph{proof}: by (T1), $\tau'$ is an
  $\cring[\Ga'']$-chain map between projective resolutions (restriction to $\Ga''$ keeps projectivity, $\cring[\Ga]$ being free over $\cring[\Ga'']$), unique up to
  $\Ga''$-homotopy by the fundamental lemma; (T3) fixes the augmentation; Lemma~\ref{lem:cosettransport}(1),(2) apply verbatim (they use (T1) and
  $\partial(s_jy)=s_j\partial y$ only); so the coset transport induces $\cor\circ(\conj_g)_{\ast}\circ\res$, by Proposition~\ref{prop:isomodelell}(4),(5)
  (Brown III.9--10: independent of the projective resolution), and so does the transported $T'_p$ by~(5)); checked degree by degree,
  $q\le q_{\max}$: equal characteristic polynomials over $\F_3$ (and $\F_2$), and equal lengths and invariant factors of $H^{q}$ and of $\operatorname{im}(T-a)$,
  $a\in\{0,1,-1\}$, over $\Z/3^{e}$, $e\le3$, at $\ell\in\{2,3\}$, $p\in\{2,3,5\}$, on every refined descriptor the package builds, with one control cell with
  $-I\mapsto-1$ (every $H^{q}$ zero on both sides). A failure stops the package. (v1.18.0-pre.2, the referee F1's m-C2: at $\ell=2$, where $\Iso^{(2)}$ ``gains
  little'' (Remark~\ref{rmk:isomodellmp}(2)), the checks over $\F_2$ and $\Z/2^{e}$ are a check of the code path --- the projective repair of the lift, the
  periodic column at $v_4$ --- not new mathematics, as~(vi) is for $\F_9$.)
  \COMP\textsuperscript{a} (v1.18.0-pre.2, slot-18 WP-7, \code{HeckeIsotropyModelEllValidation} arm (IL5)): held on the two synthetic refined descriptors of the
  code map's item~(iv), $\ell\in\{3,2\}$, $p\in\{2,3,5\}$: the hypothesis certified on each model ($d\circ d=0$, well defined, \ref{IM1}, \ref{IM2}); $T_p^{Y'}$ and
  the transported $T'_p$ equal on $H^{q}$, $q\le5$ (their difference maps cocycles to coboundaries), with equal image invariant factors of $T-a$ over
  $\Z/\ell^{e}$, $e\le3$, and equal characteristic polynomials over $\F_\ell$, both also equal to $T_p$ on $\Iso$; the control $n=1$ at $\ell=3$; $684$ ($1080$
  under the full grid) degree comparisons. The \PROVED{} tag is unchanged.
\end{enumerate}
\end{remark}

% NEW at v1.23.0-pre.12 (hecke-extensions slot 24, milestone D0 of docs/implementation/slot24/DESCENT_plan.md v0.2, its clause C-D0, the
% audit finding M3 of DESCENT_plan_audit.md).  Placed at the end of the chapter's statements so that no statement number moves.
% A pre-release draft, NOT refereed; the plan requires a referee before D0's validation lands.
\begin{proposition}[The certified-exponent rule on the $\ell$-local model: one complex over {$\Z[1/\mathsf q_\ell]$}]\label{prop:certifiedeell}
\PROVED{} (v1.23.0-pre.12; refereed ACCEPT-with-fixes, \file{docs/implementation/slot24/C-D0\_referee.md}; its fixes F1--F7 applied at v1.23.0-pre.13,
not re-refereed.) Assume \ref{H4} and \ref{B1}--\ref{B3}; let $\ell$ be a prime, $\Sigma_s$, $\Ga_\sigma$, $P_\sigma$, $Q_\sigma$ and $e_\sigma$ as
in Definition~\ref{def:isomodelell}, and put
\[
  \mathsf q_\ell:=\operatorname{lcm}\{\abs{Q_\sigma}:\sigma\in\Sigma_0\sqcup\Sigma_1\sqcup\Sigma_2\},\qquad \mathsf O_\ell:=\Z[1/\mathsf q_\ell]\subset\Z_{(\ell)} .
\]
On the model of Definition~\ref{def:model7}: at $\ell=2$, $\abs{Q_\sigma}\in\{1,3\}$ ($Q_{v_4}=\langle U^{2}\rangle$, $U=ST$, the only non-trivial one), so $\mathsf q_2=3$
and $\mathsf O_2=\Z[1/3]$; at $\ell=3$, $\abs{Q_\sigma}\in\{2,4\}$ ($Q_{v_6}=\Ga_{v_6}=\langle S\rangle$ of order $4$, every other $Q_\sigma=\{\pm I\}$), so $\mathsf q_3=4$ and
$\mathsf O_3=\Z[1/2]$; at $\ell\ge5$, $Q_\sigma=\Ga_\sigma$, $\mathsf q_\ell=12$ and $\mathsf O_\ell=\Z[1/6]$. Let $M$ be a $\Z[\Ga]$-module that is finitely generated and
free over $\Z$ ($\Vcoef=\Sym^{n}\Z^{2}$, $\Ind_{\Ga_0(N)}^{\Ga}\Vcoef$, \dots), $M_S:=M\otimes_\Z S$ for a ring $S$, and $q_{\max}\ge0$.
\begin{enumerate}
\item \textbf{(One complex, and its base changes.)} Definition~\ref{def:isomodelell} and Proposition~\ref{prop:isomodelell}(1)--(3) hold verbatim with $\cring$ any
  commutative ring in which $\mathsf q_\ell$ is a unit, not necessarily a $\Z_{(\ell)}$-algebra. Fix one such model $\Iso^{(\ell)}_{\mathsf O_\ell}$ over
  $\cring=\mathsf O_\ell$. For every $\mathsf O_\ell$-algebra $S$, the complex $\Iso^{(\ell)}_{\mathsf O_\ell}\otimes_{\mathsf O_\ell}S$ with the base-changed $d$ and
  $\cmpI'$ is a model in this sense over $S$ --- for $S$ a $\Z_{(\ell)}$-algebra an $\ell$-local isotropy model of Definition~\ref{def:isomodelell} --- and so computes
  $H^{q}(\Ga;M_S)$ for every $q$ (Proposition~\ref{prop:isomodelell}(3)). In particular the complexes over $S=\Z/\ell^{K}$ (every $K\ge1$) and over $S=\F_{\ell'}$
  ($\ell'$ a prime, $\ell'\nmid\mathsf q_\ell$; every $\ell'\ge5$) are base changes of \textbf{one} complex: no lift $d'^{(1)}$, $d'^{(2)}$ is chosen per ring.
\item \textbf{(The integral cochains.)} $C^{\bullet}_\ell:=\Hom_{\mathsf O_\ell[\Ga]}(\Iso^{(\ell)}_{\mathsf O_\ell},M_{\mathsf O_\ell})
  =\bigoplus_{s+t=q}\bigoplus_{\sigma}M_{\mathsf O_\ell}^{Q_\sigma}$ (the index set of Proposition~\ref{prop:isomodelell}(3)) is a complex of finitely generated free
  $\mathsf O_\ell$-modules; $C^{\bullet}_\ell\otimes_{\mathsf O_\ell}S=C^{\bullet}_{\Iso^{(\ell)}\otimes S}(M_S)$ for every $\mathsf O_\ell$-algebra $S$;
  $H^{q}(C^{\bullet}_\ell)\cong H^{q}(\Ga;M)\otimes\mathsf O_\ell$ and $H^{q}(C^{\bullet}_\ell\otimes\Z_{(\ell)})\cong H^{q}(\Ga;M)\otimes\Z_{(\ell)}$ for every $q$.
\item \textbf{(The rule over the local principal ideal domain $\Z_{(\ell)}$.)} Proposition~\ref{prop:certifiede}(1)--(3) and Remark~\ref{rmk:certifiedeperdiff} hold for
  $C:=C^{\bullet}_\ell\otimes_{\mathsf O_\ell}\Z_{(\ell)}$, a complex of finitely generated free $\Z_{(\ell)}$-modules, in place of a complex of free $\Z$-modules, at the
  prime $\ell$: the elementary divisors $\ell^{a_t}$ are those of $d^{q}$ over $\Z_{(\ell)}$, $r_e(d^{q})$ is read on $C\otimes\Z/\ell^{e}$, which is the complex of~(1)
  over $\Z/\ell^{e}$, and what is read off is the $\Z_{(\ell)}$-type of $H^{q}(\Ga;M)\otimes\Z_{(\ell)}$, together with
  $H^{q}(\Ga;M\otimes\Z/\ell^{f})\cong H^{q}(\Ga;M)\otimes\Z/\ell^{f}\oplus\mathrm{Tor}(H^{q+1}(\Ga;M),\Z/\ell^{f})$ for every $f\ge1$. As in
  Proposition~\ref{prop:certifiede}(5), it applies to the truncation $\sigma^{\le q_{\max}+1}$ (its $d^{q}$, $q\le q_{\max}$, are those of the full complex).
\item \textbf{(Where $\rank_{\Q}$ comes from.)} $\rank_{\Q}d^{q}=\rank_{\F_{\ell'}}(d^{q}\otimes\F_{\ell'})$ for every prime $\ell'\nmid\mathsf q_\ell$ at which
  $H^{q+1}(\Ga;M)$ has no $\ell'$-torsion (such an $\ell'$ \emph{qualifies for $d^{q}$}), and $\rank C^{q}_\ell=\dim_{\F_{\ell'}}C^{q}_\ell\otimes\F_{\ell'}$ for every
  $\ell'\nmid\mathsf q_\ell$. Put $b_j:=\dim_{\Q}H^{j}(\Ga;M\otimes\Q)$. (v1.23.0-pre.13, the C-D0 referee's F1, F2:) for every $M$ as above,
  \begin{enumerate}
  \item[(a)] \textbf{($q\ge1$.)} $24\cdot H^{j}(\Ga;M)=0$ for every $j\ge2$, so every prime $\ell'\ge5$ qualifies for $d^{q}$, $q\ge1$ (the case $q=1$, where
    $H^{q+1}=H^{2}$, included); $\ell'\ge5$ divides neither $24$ nor $\mathsf q_\ell\mid12$. Moreover $b_j=0$ for $j\ge2$.
  \item[(b)] \textbf{($q=0$.)} A prime $\ell'\nmid\mathsf q_\ell$ qualifies for $d^{0}$ iff $\dim_{\F_{\ell'}}H^{0}(\Ga;M\otimes\F_{\ell'})=b_0$, the left side measured on this
    complex as $\rank P_0-\rank\delta^{0}$ over $\F_{\ell'}$ (with (5)); and $\rank_{\Q}d^{0}=\rank C^{0}_\ell-b_0$ in any case. $b_0$ is an input stated per $M$:
    for $M=\Vcoef$ and for $M=\Ind_{\Ga_0(N)}^{\Ga}\Vcoef$ ($N\ge1$), $b_0=0$ if $n>0$ and $b_0=1$ if $n=0$. For $M=\Vcoef$ every prime $\ell'\ge5$ with $\ell'>n+1$
    qualifies for every $d^{q}$ (Proposition~\ref{prop:badprimes}), the set of Proposition~\ref{prop:certifiede}(4).
  \item[(c)] \textbf{(The admissible sources.)} $\rank_{\Q}d^{q}$ may be taken from exactly two sources: (i) $\rank_{\F_{\ell'}}$ of the reduction mod $\ell'$ of
    \textbf{this} $\mathsf O_\ell$-complex $C^{\bullet}_\ell$ (equivalently of $\delta^{q}$, by (5)) --- the same model $\Iso^{(\ell)}_{\mathsf O_\ell}$ reduced mod $\ell'$ ---
    at an $\ell'$ that qualifies for $d^{q}$ by (a) or (b); (ii) the rational Betti numbers, through
    $\rank_{\Q}d^{q}=\rank C^{q}_\ell-b_q-\rank_{\Q}d^{q-1}$ ($d^{-1}:=0$; $b_j=0$ for $j\ge2$ by (a)), each $b_j$ used taken from a certified source. The rank of a
    differential of any \textbf{other} complex is \textbf{not} admissible --- not a census rank, not the shipped $C_7$ (Bredon) pipeline's rank at $\ell'$, and not
    the rank of the $\ell'$-local model $\Iso^{(\ell')}$ over $\F_{\ell'}$, whose columns and differentials are not those of $\Iso^{(\ell)}$ although its cohomology is
    the same; a computation on another complex enters only through~(ii), as a certified value of some $b_j$.
  \end{enumerate}
\item \textbf{(The full layout.)} Let $F^{q}:=\bigoplus_{x}M_{\mathsf O_\ell}$ over the cells $x$ of degree $q$, $P_q:=\mathrm{diag}_x\,\rho_M(e_x)$, and
  $\delta^{q}\colon F^{q}\to F^{q+1}$ the matrix whose block at $(y,x)$ is $\rho_M(a_{yx})$, $a_{yx}$ the coefficient of $d$ from the cell $y$ of degree $q+1$ to the cell
  $x$ of degree $q$ (so $a_{yx}=e_y\,a_{yx}\,e_x$, its well-definedness as a map $\cring[\Ga]e_y\to\cring[\Ga]e_x$). Then $P_q$ is an idempotent over $\mathsf O_\ell$ with
  image $C^{q}_\ell$, $\delta^{q}=P_{q+1}\delta^{q}P_q$, and there are $U$, $V$ invertible over $\mathsf O_\ell$ with $U\delta^{q}V=\mathrm{diag}(d^{q},0)$. Hence $\delta^{q}$ and
  $d^{q}$ have the same rank over $\Q$ and over every $\F_{\ell'}$, $\ell'\nmid\mathsf q_\ell$, and the same non-zero elementary divisors over $\Z_{(\ell)}$;
  $\rank C^{q}_\ell=\rank P_q$ over $\Q$ and over every $\F_{\ell'}$; and (3), (4) and (6) may be applied to $\delta^{q}$, $P_q$ in place of $d^{q}$, $C^{q}_\ell$.
\item \textbf{(The valuation read-out is the invariant-factor count.)} Let $A$ be $d^{q}$ or $\delta^{q}$ reduced to $\Z/\ell^{K}$, $K\ge1$ (at $\ell=2$ possibly through
  $\Z/2^{W}$, $W\ge K$, and then read mod $2^{K}$). The Smith valuations of $A$ over $\Z/\ell^{K}$ form the \textbf{list, with multiplicity},
  $(\sigma_1,\dots,\sigma_s)$, $\sigma_1\le\dots\le\sigma_s<K$, of the exponents of the non-zero diagonal entries $\ell^{\sigma_i}$ of its Smith form, $s$ the number of
  those entries (counted with multiplicity; repeated values, $\sigma_i=0$ above all, are the rule) --- equivalently the invariant factors $\ell^{K-\sigma_i}$ of its column
  span. They are exactly the exponents $a_t<K$, with multiplicity, of the elementary divisors of $d^{q}$ over $\Z_{(\ell)}$, and $s=r_K(d^{q})$ of
  Proposition~\ref{prop:certifiede}(1). Hence, if $s=\rank_{\Q}d^{q}$, read from a source admissible by~(4)(c), then every $a_t<K$, $(\sigma_1,\dots,\sigma_s)$ is the
  full list of exponents of the non-zero elementary divisors of $d^{q}$ over $\Z_{(\ell)}$, and $\mathrm{tors}\,H^{q+1}(\Ga;M)\otimes\Z_{(\ell)}\cong\bigoplus_{\sigma_i>0}\Z/\ell^{\sigma_i}$
  exactly (Remark~\ref{rmk:certifiedeperdiff} with $e_q=K$); if $s<\rank_{\Q}d^{q}$, then $K$ is too small and nothing about $d^{q}$'s divisors $\ell^{a}$, $a\ge K$, is
  read off. (v1.23.0-pre.13, the referee's F3, F4:) \textbf{A-priori exponents.} For $q\ge1$ every $a_t\le v_\ell(24)$: $a_t\le3$ at $\ell=2$, $a_t\le1$ at $\ell=3$,
  $a_t=0$ at $\ell\ge5$; so for $q\ge1$ every $K\ge v_\ell(24)+1$ certifies, and $\sigma_i\le v_\ell(24)$ is a consistency check of the read-out. For $q=0$ (the
  $\ell$-primary torsion of $H^{1}$) no a-priori bound is asserted here (Proposition~\ref{prop:badprimes}'s closing note, Remark~\ref{rmk:bocksteinexp}). \textbf{The
  stopping rule:} if $s<\rank_{\Q}d^{q}$, raise $K$ and rerun, up to the cap of the ring in use (the code map); at the cap, refuse --- never report a group from
  $s<\rank_{\Q}d^{q}$. By the a-priori bound a cap $\ge v_\ell(24)+1$ binds only at $d^{0}$.
\end{enumerate}
\end{proposition}

\begin{proof}
(1) The proof of Proposition~\ref{prop:isomodelell}(1)--(3) uses of $\cring$ only that $\abs{Q_\sigma}\in\cring^{\times}$ for every $\sigma$: $e_\sigma$ is then an
idempotent of $\cring[\Ga_\sigma]$ and $F'^{\sigma}_t=\cring[\Ga_\sigma]e_\sigma$ is projective; for $P_\sigma=1$, $Q_\sigma=\Ga_\sigma$ and the column $\cring=\cring[\Ga_\sigma]e_\sigma$ is
projective; the averaging $y\mapsto e_\sigma y$ is defined. The exactness of $F'^{\sigma}\to\cring$ holds over every ring (the periodic resolution of a finite cyclic
group, resp.\ $\cring\to\cring$), the columns consist of free $\cring$-modules, so the contracting homotopies are $\cring$-linear; the comparison theorem, the
filtration argument and $\Ext^{q}_{\cring[\Ga]}(\cring,M)=H^{q}(\Ga;M)$ are ring-independent. Nowhere is $\ell$ inverted or localised at, so the hypothesis
``$\Z_{(\ell)}$-algebra'' may be replaced by ``$\mathsf q_\ell\in\cring^{\times}$''; $\mathsf O_\ell$ is such a ring, and Proposition~\ref{prop:isomodelell}(1) produces
$\Iso^{(\ell)}_{\mathsf O_\ell}$. \emph{Base change.} $G'_{s,t}\otimes_{\mathsf O_\ell}S=\bigoplus_\sigma S[\Ga/Q_\sigma]$ and $F'^{\sigma}\otimes S$ is the column over $S$
(the same group-ring coefficients), so the base-changed complex has the terms and the shape $d=d'^{(0)}+d'^{(1)}+d'^{(2)}$ of Definition~\ref{def:isomodelell} over $S$;
$d\circ d=0$, \ref{IM1} and \ref{IM2} are identities between finitely many group-ring coefficients and survive every ring map. The augmentation $\cmpI'$ over
$\mathsf O_\ell$ is a surjective quasi-isomorphism (Proposition~\ref{prop:isomodelell}(2)); its mapping cone is a bounded-below exact complex of free
$\mathsf O_\ell$-modules ($\mathsf O_\ell[\Ga/Q_\sigma]$ and $C_{s}(\Xbulk)\otimes\mathsf O_\ell$ are free over $\mathsf O_\ell$), hence contractible over $\mathsf O_\ell$, and
$-\otimes_{\mathsf O_\ell}S$ keeps it contractible; so $\cmpI'\otimes S$ is a quasi-isomorphism, surjective by right exactness. Hence
$\Iso^{(\ell)}_{\mathsf O_\ell}\otimes S$ satisfies Definition~\ref{def:isomodelell} over $S$ (with the replaced hypothesis), and Proposition~\ref{prop:isomodelell}(3)
applies. A model produced over $S$ directly by the lifting procedure would be another model, homotopy equivalent to this one
(Proposition~\ref{prop:isomodelell}(4)'s uniqueness argument with both sides projective); that gives the same cohomology but \emph{not} a common integral form,
which is what (3) needs and why (1) fixes one complex.

(2) $\Hom_{\mathsf O_\ell[\Ga]}(\mathsf O_\ell[\Ga/Q_\sigma],M_{\mathsf O_\ell})=M_{\mathsf O_\ell}^{Q_\sigma}=\rho_M(e_\sigma)M_{\mathsf O_\ell}$, a direct summand of the
finitely generated free $\mathsf O_\ell$-module $M_{\mathsf O_\ell}$, hence finitely generated projective over the principal ideal domain $\mathsf O_\ell$ (a localisation of
$\Z$), hence free. Taking the image of an idempotent commutes with $-\otimes S$, so $\rho_M(e_\sigma)M_{\mathsf O_\ell}\otimes S=\rho_M(e_\sigma)M_S=M_S^{Q_\sigma}$, and the
cochain differential, given by the coefficients of $d$ acting through $\rho_M$, is base-changed with them. By Proposition~\ref{prop:isomodelell}(3) over
$\mathsf O_\ell$, $H^{q}(C^{\bullet}_\ell)=H^{q}(\Ga;M_{\mathsf O_\ell})$. Since $\Iso_q$ is finitely generated free over $\Z[\Ga]$ (Proposition~\ref{prop:isomodel}(2)),
$\CI(M_{\mathsf O_\ell})=\CI(M)\otimes\mathsf O_\ell$, and $\mathsf O_\ell$ is flat over $\Z$, so $H^{q}(\Ga;M_{\mathsf O_\ell})=H^{q}(\Ga;M)\otimes\mathsf O_\ell$; likewise for the
flat $\mathsf O_\ell$-algebra $\Z_{(\ell)}$, with $C^{\bullet}_\ell\otimes\Z_{(\ell)}=C^{\bullet}_{\Iso^{(\ell)}\otimes\Z_{(\ell)}}(M_{\Z_{(\ell)}})$ by the first part.

(3) The proof of Proposition~\ref{prop:certifiede}(1)--(3) uses of $\Z$ only: the Smith form over the local principal ideal domain $\Z_{(\ell)}$ of a matrix with
entries in $\Z_{(\ell)}$ (an integer matrix is one case of it); that $C^{q+1}$ is torsion-free, and additivity of ranks over the fraction field $\Q$; and the
universal coefficient theorem for a complex of free modules over a principal ideal domain \STD{\cite[\S3.6]{weibel}}. All three hold for $C$, which is a complex of
finitely generated free $\Z_{(\ell)}$-modules by (2). $C\otimes_{\Z_{(\ell)}}\Z/\ell^{e}=C^{\bullet}_\ell\otimes_{\mathsf O_\ell}\Z/\ell^{e}$, which by (1), (2) is the complex
built over $\Z/\ell^{e}$ and computes $H^{q}(\Ga;M\otimes\Z/\ell^{e})$; and $H^{q}(C)=H^{q}(\Ga;M)\otimes\Z_{(\ell)}$ by (2), so its torsion is the $\ell$-primary torsion
of $H^{q}(\Ga;M)$. Remark~\ref{rmk:certifiedeperdiff}'s proof is (1), (2) and the universal coefficient theorem again, one differential at a time. The truncation
statement is Proposition~\ref{prop:certifiede}(5)'s observation, unchanged.

(4) Over the principal ideal domain $\mathsf O_\ell$, $d^{q}=U\,\mathrm{diag}(c_1,\dots,c_r,0)\,V$ with $U$, $V$ invertible and $c_i\neq0$, $r=\rank_{\Q}d^{q}$. The proof of
Proposition~\ref{prop:certifiede}(2), over $\mathsf O_\ell$, gives $\mathrm{tors}\,\coker d^{q}=\mathrm{tors}\,H^{q+1}(C^{\bullet}_\ell)=\mathrm{tors}\,H^{q+1}(\Ga;M)\otimes\mathsf O_\ell$,
whose order ideal is generated by $\prod c_i$. A prime $\ell'\nmid\mathsf q_\ell$ is a prime of $\mathsf O_\ell$, with residue field $\F_{\ell'}$; the reduction keeps the rank
$r$ iff $\ell'$ divides no $c_i$, i.e.\ iff $H^{q+1}(\Ga;M)\otimes\mathsf O_\ell$, equivalently (as $\ell'$ is not inverted) $H^{q+1}(\Ga;M)$, has no $\ell'$-torsion. A free
module of rank $m$ reduces to an $\F_{\ell'}$-space of dimension $m$. For $M=\Vcoef$: by Proposition~\ref{prop:badprimes}, the torsion primes of
$H^{\bullet}(\Ga;\Vcoef)$ are $2$, $3$ and primes $\le n-1$ ($n>0$ even), $2$ ($n$ odd), $2$, $3$ ($n=0$); every $\ell'\ge5$, $\ell'>n+1$, avoids them and does not divide
$\mathsf q_\ell\mid12$.

\emph{(4)(a)} (v1.23.0-pre.13). Apply $\Ext^{\bullet}_{\Z[\Ga]}(-,M)$ to the short exact sequence \eqref{eq:treeses} of Claim~\ref{cl:amalgam}(1). For a subgroup
$H\le\Ga$, $\Z[\Ga/H]=\Z[\Ga]\otimes_{\Z[H]}\Z$ and $\Z[\Ga]$ is free over $\Z[H]$, so $\Ext^{j}_{\Z[\Ga]}(\Z[\Ga/H],M)\cong\Ext^{j}_{\Z[H]}(\Z,M)=H^{j}(H;M)$ (Shapiro). This is the
Mayer--Vietoris sequence \eqref{eq:mvamalgam} for the arbitrary $\Z[\Ga]$-module $M$ (Claim~\ref{cl:amalgam}(2), whose argument uses nothing about the module). For $j\ge2$ it
gives $0\to C\to H^{j}(\Ga;M)\to A\to0$ with $C$ a quotient of $H^{j-1}(\langle-I\rangle;M)$ and $A$ a submodule of $H^{j}(\langle S\rangle;M)\oplus H^{j}(\langle U\rangle;M)$.
Since $j-1\ge1$ and $\abs{H}$ kills $H^{i}(H;-)$ for $i\ge1$ and $H$ finite \STD{\cite[Ch.~III.10]{brown}}, $2$ kills $C$ and $\operatorname{lcm}(4,6)=12$ kills $A$, so $24$
kills $H^{j}(\Ga;M)$ --- the argument of Claim~\ref{cl:amalgam}(4), which uses nothing about $\Vcoef$. For $q\ge1$, $j=q+1\ge2$ (at $q=1$, $j=2$ and $j-1=1\ge1$), so
$H^{q+1}(\Ga;M)$ has no $\ell'$-torsion for $\ell'\ge5$, and $\ell'\nmid\mathsf q_\ell$ as $\mathsf q_\ell\mid12$. As $H^{j}(\Ga;M\otimes\Q)=H^{j}(\Ga;M)\otimes\Q$ (flatness, as in (2)),
$b_j=0$ for $j\ge2$. A level-$N$ module $\Ind_{\Ga_0(N)}^{\Ga}V$ is a $\Z[\Ga]$-module, so the bound holds for it with no further argument.

\emph{(4)(b).} $C^{\bullet}_\ell$ is a complex of finitely generated free modules over the principal ideal domain $\mathsf O_\ell$ (by (2)), so the universal coefficient
theorem of (3) gives $H^{0}(C^{\bullet}_\ell\otimes\F_{\ell'})\cong H^{0}(C^{\bullet}_\ell)\otimes\F_{\ell'}\oplus\mathrm{Tor}_1^{\mathsf O_\ell}(H^{1}(C^{\bullet}_\ell),\F_{\ell'})$. Here
$H^{0}(C^{\bullet}_\ell)=\ker d^{0}$ is a submodule of the free $C^{0}_\ell$, hence free, of rank $\dim_{\Q}H^{0}(C^{\bullet}_\ell)\otimes\Q=b_0$ (by (2) and flatness), and
$\mathrm{Tor}_1(H^{1}(C^{\bullet}_\ell),\F_{\ell'})=H^{1}(C^{\bullet}_\ell)[\ell']$, which vanishes iff $H^{1}(\Ga;M)\otimes\mathsf O_\ell$, equivalently $H^{1}(\Ga;M)$, has no
$\ell'$-torsion. The left side is $H^{0}(\Ga;M\otimes\F_{\ell'})$ by (1), (2), of dimension $\rank C^{0}_\ell-\rank_{\F_{\ell'}}d^{0}=\rank P_0-\rank_{\F_{\ell'}}\delta^{0}$ by (5). Over
$\Q$, rank--nullity gives $\rank_{\Q}d^{0}=\rank C^{0}_\ell-b_0$. \emph{The values of $b_0$.} By Shapiro, $H^{0}(\Ga;\Ind_{\Ga_0(N)}^{\Ga}V_{\Q})=H^{0}(\Ga_0(N);V_{\Q})$ ($N=1$ is
$M=\Vcoef$; induction and coinduction agree at finite index). For $n=0$ this is $\Q$. For $n>0$, $\Ga_0(N)$ contains $T=\smat1101$ and $T_N:=\smat10N1$. In
$\Vcoef\otimes\Q=\Sym^{n}\Q^{2}$, homogeneous polynomials of degree $n$ in $x$, $y$, one of the two acts by a substitution $f(x,y)\mapsto f(x,y+cx)$ and the other by
$f(x,y)\mapsto f(x+c'y,y)$ ($c,c'\in\{\pm1,\pm N\}$; which is which, and the signs, depend on the convention of Definition~\ref{def:rhon}, not on the argument). If
$f(x,y+cx)=f(x,y)$, then $f(x,y+mcx)=f(x,y)$ for every $m\in\Z$, so $t\mapsto f(x,y+tx)-f(x,y)$, a polynomial in $t$ with infinitely many zeros, vanishes; hence
$x\,\partial_yf=0$, and $f\in\Q x^{n}$. Likewise the other fixes only $\Q y^{n}$. As $\Q x^{n}\cap\Q y^{n}=0$ for $n>0$, $b_0=0$.

\emph{(4)(c).} (i) is (4) at a qualifying $\ell'$. (ii) is rank--nullity over $\Q$, $b_q=\dim\ker(d^{q}\otimes\Q)-\rank_{\Q}d^{q-1}=\rank C^{q}_\ell-\rank_{\Q}d^{q}-\rank_{\Q}d^{q-1}$,
with $b_q=H^{q}(C^{\bullet}_\ell)\otimes\Q$'s dimension by (2). \emph{Why nothing else.} (6)'s certificate compares $s=r_K(d^{q})$, which never exceeds $\rank_{\Q}d^{q}$, with
a number that must \emph{equal} $\rank_{\Q}d^{q}$ of this complex; a smaller number (an $\F_{\ell'}$ rank at a non-qualifying $\ell'$) can certify a too-small $K$, and the
rank of another complex's differential is a different number in general --- cohomology is invariant under the choice of model, the ranks of the differentials are
not (they depend on the ranks of the cochain modules, which differ from model to model). Betti numbers are invariants, which is why (ii) admits them from any certified
source.

(5) $e_x$ is an idempotent of $\mathsf O_\ell[Q_x]$ and $\rho_M$ is multiplicative, so $P_q$ is an idempotent with image $\bigoplus_x\rho_M(e_x)M_{\mathsf O_\ell}=C^{q}_\ell$
(by (2), cell by cell). By Proposition~\ref{prop:isomodelell}(3)'s identification $\Hom(\cring[\Ga]e_x,M)=e_xM$, the cochain differential $d^{q}$ of $C^{\bullet}_\ell$ sends the
component at $x$ to the component at $y$ by $\rho_M(a_{yx})$; this is $\delta^{q}$ restricted to $C^{q}_\ell$, and $a_{yx}=e_ya_{yx}e_x$ gives
$\rho_M(a_{yx})=\rho_M(e_y)\rho_M(a_{yx})\rho_M(e_x)$, i.e.\ $\delta^{q}=P_{q+1}\delta^{q}P_q$: $\delta^{q}$ vanishes on $\ker P_q$ and lands in $\operatorname{im}P_{q+1}$.
$F^{q}=\operatorname{im}P_q\oplus\ker P_q$ with both summands free over $\mathsf O_\ell$ (direct summands of a finitely generated free module over a principal ideal domain); in
adapted bases of $F^{q}$ and $F^{q+1}$, $\delta^{q}=\mathrm{diag}(d^{q},0)$, and the base-change matrices $U$, $V$ are invertible over $\mathsf O_\ell$, hence over every
$\mathsf O_\ell$-algebra ($\Z_{(\ell)}$, $\Z/\ell^{K}$, $\F_{\ell'}$, $\Q$). Ranks and elementary divisors are invariant under such $U$, $V$, and the zero block adds none.
$\rank P_q=\rank\operatorname{im}P_q$ over every field $\mathsf O_\ell\to k$, the image of an idempotent commuting with base change.

(6) By (5), $A$ and the reduction of $d^{q}$ have the same Smith form up to zero rows and columns, so it suffices to treat $d^{q}$. Over $\Z_{(\ell)}$,
$d^{q}=U\,\mathrm{diag}(\ell^{a_t},0)\,V$; reducing mod $\ell^{K}$ keeps $U$, $V$ invertible, and the Smith form over the local ring $\Z/\ell^{K}$ is unique, so it is
$\mathrm{diag}(\ell^{\min(a_t,K)},0)$ with $\ell^{K}=0$: its non-zero diagonal entries are the $\ell^{a_t}$ with $a_t<K$ (this is Proposition~\ref{prop:certifiede}(1)'s
proof). At $\ell=2$ the composite $\mathsf O_2\to\Z/2^{W}\to\Z/2^{K}$ is the reduction map $\mathsf O_2\to\Z/2^{K}$, so a matrix formed mod $2^{W}$ and read mod $2^{K}$ is
the reduction to $\Z/2^{K}$. The count, with multiplicity, is $s=\#\{t:a_t<K\}=r_K(d^{q})$; by (4)(c) the number compared with it is $\rank_{\Q}d^{q}$; and
Proposition~\ref{prop:certifiede}(1) with (3) and Remark~\ref{rmk:certifiedeperdiff} give the rest. \emph{A-priori exponents.} By (3) and the proof of (4),
$\bigoplus_t\Z/\ell^{a_t}\cong\mathrm{tors}\,\coker(d^{q}\otimes\Z_{(\ell)})=\mathrm{tors}\,H^{q+1}(\Ga;M)\otimes\Z_{(\ell)}$, which for $q\ge1$ is killed by $\ell^{v_\ell(24)}$ by
(4)(a); so $a_t\le v_\ell(24)$, and with $K\ge v_\ell(24)+1$ every $a_t<K$, i.e.\ $s=\rank_{\Q}d^{q}$. The stopping rule only restates the two cases of the
certificate. That a given routine returns the Smith valuations of its input over $\Z/\ell^{K}$, with multiplicity, is that routine's specification, not part of this
proposition.
\end{proof}


\paragraph{Scope: what this chapter does not do.} It asserts the $\ell$-local model for every prime $\ell$ over every $\Z_{(\ell)}$-algebra, its comparisons and Hecke transport,
and the closed form and the Bredon comparison at $\ell=3$, all \PROVED. Not asserted: anything over $\Z$ (the model exists only $\ell$-locally --- at best over $\Z[1/\mathsf q_\ell]$, Proposition~\ref{prop:certifiedeell} (v1.23.0-pre.12); the glueing is
\S\ref{sec:descent}'s); a strict lift on $\Iso^{(\ell)}$ of $\Xbulk$ itself (the transport (5) is a conjugation; v1.18.0-pre.2 builds one on a synthetic refinement $Y'$, item~(v)
of the code map); the Moerdijk--Pronk simplicial set $S(\mathcal T)$ as a
normative object (Remark~\ref{rmk:mpcaveat}); a change of what the default coefficient ring names at $\ell\in\{2,3\}$.

\paragraph{Code map.} \textbf{Built at v1.16.0} at $\ell=3$ (slot-16 WP-16, the stretch package, audited by a cross-model audit
\file{docs/implementation/slot16/WP-16\_audit.md}, ACCEPT WITH REPAIRS, repaired by WP-17; not built at v1.16.0-pre.1 --- contract first, rule~21).
\file{Integral/Model/IsotropyModelEll} (\code{integral\_model::LocalModel}) on \file{Integral/Model/IsotropyModel}'s plug point, with a parameter $\ell$ that
builds the columns $F'^{\sigma}$ and the cochain groups $M^{Q_\sigma}$ ([Q1n]: \file{docs/implementation/slot14/WP-Q1\_small\_ell\_orbifold\_note.md} \S4 item~2), exact over
$\Q[\Ga]$ with denominators prime to $\ell$; the strict comparisons $\Phi_\ell$, $\Psi_\ell$ of Proposition~\ref{prop:isomodelell}(4). \code{HeckeIsotropyModelEllValidation}
(\code{standard}; \file{docs/implementation/slot16/WP-16\_report.md}) runs \ref{IL1}--\ref{IL4} against $\Iso$: \ref{IL1} on $\Z/3^{e}$, $e\le3$, $n\le20$ against
$\Iso$ and [IB] Appendix~A, and $\Iso^{(3)}=\Iso$ on the $\F_3$ cells $20<n\le200$; \ref{IL2} on $103$ cells up to $n=200$; \ref{IL3} at level~$1$ and
(v1.17.0-pre.5, slot-17 WP-7) at level~$N$ against the pipeline over $\F_3$, \ref{IL4} over $\F_3$ and (v1.17.0-pre.5) over $\Z/3^{e}$ with the pair and
cone (both \COMP\textsuperscript{a} partial, above); the case $\ell=2$ through the same code (the MP arm). The lifting procedure averages every
lift by $e_\sigma$ (Proposition~\ref{prop:isomodelell}(1)); \ref{IL1} tests $d\circ d=0$ and \ref{IM1}, \ref{IM2} on the built complex (v1.16.0-pre.3, WP-14 referee
C-1), where the average is a no-op on $d$ (the reason is in the proof of (1)). \textbf{Where the average is tested} is the comparison map $\Psi_\ell$'s
well-definedness check $e_x\Psi_\ell(x)=\Psi_\ell(x)$ in \ref{IL4}: without the average $\Psi_3$ still passes the chain-map identity and strictness but fails
well-definedness on $22$ of its $25$ source generators, the validation's negative control (v1.16.0, WP-16 audit F-1).
\textbf{Run at v1.17.0-pre.5} (slot-17 WP-7, merged \texttt{ea3a102}, audited ACCEPT, \file{docs/implementation/slot17/WP-7\_report.md},
\file{docs/implementation/slot17/WP-7\_audit.md}): (i) \ref{IL3}'s level-$N$ clause, (ii) \ref{IL4} over $\Z/3^{e}$, $e\le3$, (iii) the pair and cone half of
Proposition~\ref{prop:isomodelell}(5), as stated in \ref{IL3}, \ref{IL4}; (iv) Remark~\ref{rmk:isomodellmp}(3)'s refinement half on a synthetic refined
descriptor. The descriptor is data (\file{Integral/Model/IsotropyModel}'s \code{ModelData}: representatives, stabiliser inventory, transports, horocycle, and the
\file{Hecke/HeckeModelDescriptor} header, whose $N_q$ are checked against the data), read by the one constructor that also reads the shipped descriptor;
\code{SubdivideEdgeOrbit} subdivides every edge orbit at $\zeta^{2}$ (and, in a second descriptor, a horocycle edge), the new vertices have stabiliser
$\{\pm I\}$ (the code gives a midpoint its edge's stabiliser, correct for the pointwise stabilisers of the shipped data), $-I$ lies in every stabiliser
(checked at load), and a module with $-I\mapsto-1$ is admitted (cohomology $0$). The comparison is on cohomology only: lengths and invariant factors of
$H^{q}$, $q\le5$, of $\Iso^{(3)}$ on $Y'$ against $\Iso^{(3)}$ on $\Xbulk$ and against $\Iso$ on both, over $\Z/3^{e}$, $e\le3$, on the boundary, with a
level-$N$ module, and at $\ell=2$ through the same code; the orbifold Euler characteristic $-1/12$ is unchanged. That the data reproduce the pre-slot-17
constructor is guarded indirectly: arm (REF) is a \emph{descriptor round trip} (the shipped model is now built from the same data, so it compares one code
path with itself), and the real guard is that \code{HeckeIsotropyModelValidation}, \code{HeckeStrictLiftValidation}, \code{HeckeFreeModelValidation},
\code{IsotropyModelDeltaSquaredValidation}, the arm (IL0), \ref{IL1}--\ref{IL4} and the golden and determinism gates pass unchanged (WP-7 audit MINOR-1). Not built: a strict
lift on a refined descriptor, and the Moerdijk--Pronk simplicial set itself (Remark~\ref{rmk:mpcaveat}).
\emph{The plan of the runs, as committed at v1.17.0-pre.1 (contract first, rule~21; kept for the record):}
(i) \ref{IL3}'s clause at level $N$: $M=\Ind_{\Ga_0(N)}^{\Ga}\VR$ through the shipped \code{Gamma0LevelN} table (\S\ref{sec:levelN}) --- no second coset
construction --- read as \ref{IL3} states it, its first summand the model's $\Heq{1}$ at level $N$ over $\F_3$ (Definition~\ref{def:levelNmodel}(a), the
cochain level, which is all \S\ref{sec:levelN} asserts at $\ell=3$; its tree resolution and cokernel need $6\in\cring^{\times}$ and are not used); (ii) \ref{IL4} over $\Z/3^{e}$, $e\le3$, with the $T'_p$ of
Proposition~\ref{prop:isomodelell}(5) against $\Iso$ in lengths and invariant factors; (iii) the pair and cone half of Proposition~\ref{prop:isomodelell}(5);
(iv) Remark~\ref{rmk:isomodellmp}(3)'s refinement half, on a \textbf{synthetic refined descriptor}. What the MP constraint requires of (iv): the descriptor
is a $\Ga$-equivariant subdivision $Y'$ of $\Xbulk$ given as data (vertices, simplices, stabiliser inventory, relations), read by the same code path as
the shipped descriptor (\file{Hecke/HeckeModelDescriptor}), so that no vertex count, edge list or special role of $v_4$ is hard-wired; it subdivides at least
one edge at $\zeta^{2}$ (the Moerdijk--Pronk-conformant refinement of Remark~\ref{rmk:isomodellmp}(3)), so the new vertices have stabiliser $\{\pm I\}$; the
centre is kept ($-I$ acts on the coefficients by its scalar, a module with $-I\mapsto-1$ is admissible input); and the comparison is on cohomology
(lengths and invariant factors of $H^{q}$, $q\le q_{\max}$, with $\Iso^{(3)}$ on $Y'$ against $\Iso^{(3)}$ on $\Xbulk$ and against $\Iso$), never on cochains;
and one control cell with $-I\mapsto-1$ (odd $n$), on which every $H^{q}$ is $0$ on $Y'$ and on $\Xbulk$, so that ``the centre is kept'' is run, not only stated
(v1.17.0-pre.3, slot-17 referee F1's m-C5).
A pass moves \ref{IL3} and \ref{IL4} from \COMP\textsuperscript{a} partial to \COMP\textsuperscript{a}; the tag move is recorded as a proposal for this
document's next referee (the runs passed at v1.17.0-pre.5; the tags are not moved here).
\textbf{Run at v1.18.0-pre.2} (slot-18 WP-7, \file{docs/implementation/slot18/WP-7\_report.md}; contract first, rule~21):
(v) \textbf{a strict lift on a synthetic refined descriptor} (\file{Integral/Model/RefinedStrictLift}): the descriptor carries a
subdivision record (\code{ModelData}'s \code{collapse}, \code{refines}, written by \code{SubdivideEdgeOrbit}), from which the collapse
$c\colon C_{\bullet}(Y')\to C_{\bullet}(\Xbulk)$ and the subdivision $\mathrm{sd}$ are built and certified ($dc=cd$, $d\,\mathrm{sd}=\mathrm{sd}\,d$,
$c\,\mathrm{sd}=\id$); $\tau'$ is Lemma~\ref{lem:strictlifting}'s induction on $\Iso^{(\ell)}(Y')$ for the base datum
$\mathrm{sd}\circ\tau^{s}\circ c$ ($\tau^{s}$ the shipped strict pair datum, read through the shipped lift's strict identity; the shipped
coset tables; no new coset or $\tau$-table), made well defined on the projective generators by the $Q_\sigma$-average of the twisted
extension to the free cover (load-bearing at $\ell=2$, where $Q_{v_4}=\langle U^{2}\rangle$ has a periodic column; a no-op at $\ell=3$);
certified exactly: (T1) on the Schreier generators of $\Ga_0(p)$, (T2), (T3) and the strict identity over $Y'$. The comparison of the
models over $Y'$ and $\Xbulk$ is the strict lift of $\mathrm{sd}$, resp.\ $c$ (Lemma~\ref{lem:strictlifting} along the surjective
quasi-isomorphisms onto the chain complexes of $Y'$ and $\Xbulk$). Checked by \ref{IL5} (above). (vi) \textbf{\ref{IL3} beyond $\F_3$}: over $\F_9$ with
$\Ind_{\Ga_0(N)}^{\Ga}\VR$ through the shipped twisted-monomial route (trivial character), $N\in\{1,2,5,7\}$, $n\le4$ ($N\le13$, $n\le6$ under the full grid):
dimensions equal the $\F_3$ ones and $\Iso$'s, and $\dim H^{1}=\dim\Heq{1}+\dim H^{1}(\Z/3;M^{-I})$ with the first summand the
shipped pipeline's over $\F_3$ (it does not run over $\F_9$); over $\Z/9$ in lengths against the shipped level-$N$ model's $\Heq{1}$
\emph{over $\Z/9$} (its flat-coordinate route, Definition~\ref{def:levelNmodel}(a)), $12$ cells ($32$). The MP items of (v) hold
(the report's \S4 names what a Moerdijk--Pronk $S(\mathcal T)$ descriptor would still need). On the pass of \ref{IL5} with~(vi) \ref{IL3} and \ref{IL4} are
moved to \COMP\textsuperscript{a} at v1.18.0-pre.2 (above; the slot-18 referee F1's C-4); \ref{IL5} keeps its \PROVED{} tag and gains the \COMP\textsuperscript{a} line.
\emph{The plan of (v), (vi), as committed at v1.18.0-pre.1 (contract first, rule~21; kept for the record; ``not built'' at that version):}
(v) \textbf{a strict lift on a synthetic refined descriptor}, so that $T_p$ acts on the model over $Y'$ directly and not only by transport: $\tau'$ as in
\ref{IL5}, built from the descriptor's data by the code path that reads the shipped descriptor; it may use the shipped strict pair datum $\tau^{s}$
(Proposition~\ref{prop:strictlift}) through a $\Ga$-equivariant comparison between the models over $Y'$ and over $\Xbulk$ (Lemma~\ref{lem:strictlifting} lifts
along any surjective quasi-isomorphism), but builds no new coset or $\tau$-table. Checked by \ref{IL5}; the integer sizes are recorded as in
Remark~\ref{rmk:strictliftbuilt} (evidence, bounding nothing). (vi) \textbf{\ref{IL3} beyond $\F_3$}: Proposition~\ref{prop:ellthreeclosed}(2),(3) hold over
every $\Z_{(3)}$-algebra, so \ref{IL3}'s identity is read over $\F_9$ (flat over $\F_3$: the dimensions must equal the $\F_3$ ones, a check of the $\F_q$ code
path, not new mathematics) and over $\Z/9$ in lengths ($\operatorname{length}H^{1}=\operatorname{length}H^{1}_{\mathrm{Bredon}}+\operatorname{length}H^{1}(\Z/3;M^{-I})$,
from (3)'s exact sequence), against the shipped level-$N$ model's $\Heq{1}$ over that ring where the shipped route runs over it (Definition~\ref{def:levelNmodel}(a),
the cochain level), and against $\Iso$ otherwise, as at v1.17.0-pre.5. \textbf{What the MP constraint requires of~(v)}, beyond item~(iv) above: the refined
descriptor stays out of \file{Hecke/HeckeModelDescriptor}'s registry and out of the golden harness; nothing in~(v) hard-wires the $7$-vertex model, a trivial
centre or Bredon-only semantics (the centre $-I$ lies in every stabiliser of $Y'$ and acts on the coefficients by its scalar); and the package report names what a
Moerdijk--Pronk $S(\mathcal T)$ descriptor would still need beyond~(v) (Remark~\ref{rmk:mpcaveat}). Arms of \code{HeckeIsotropyModelEllValidation} only (no new
name), its default run within the \code{standard} tier target. A pass of \ref{IL5} on every cell, with~(vi), is proposed to the next referee as the ground for
moving \ref{IL3}, \ref{IL4} to \COMP\textsuperscript{a}; the tags are not moved here. (End of the pre.1 plan.)

\paragraph{Code map of Proposition~\ref{prop:certifiedeell}.} \textbf{Stated at v1.23.0-pre.12, before D0's code} (contract first, rule~21; slot-24 milestone D0 of
\file{docs/implementation/slot24/DESCENT\_plan.md} v0.2, clause C-D0; audit finding M3 of \file{DESCENT\_plan\_audit.md}); refereed ACCEPT-with-fixes
(\file{docs/implementation/slot24/C-D0\_referee.md}), the fixes applied at v1.23.0-pre.13 (F2, F5--F7 in this code map). What the shipped code already does, read against the clauses:
(1)--(2) \textbf{hold of \file{Integral/Model/IsotropyModelEll} as built.} \code{LocalModel} builds $d$ once, in its constructor, exactly over $\Q[\Ga]$
(\code{Rat}, overflow-checked \code{int64}; an overflow throws and refuses) before any coefficient ring is named, and \code{coboundary}, \code{projector} and
\code{Pullback} reduce the stored coefficients entrywise into the \code{ModCtx} they are given (\code{reduce\_rat}); so the complexes over $\Z/\ell^{K}$ and over
$\F_{\ell'}$ are base changes of one complex, as~(1) requires. The only denominators are the coefficients $1/\abs{Q_\sigma}$ of $e_\sigma$ (\code{build\_columns}); every
stored coefficient therefore lies in $\mathsf O_\ell[\Ga]$, which is what the proof of~(1) gives. The code does \emph{not} certify this at build time: a denominator that is
not a unit of the target ring is refused at reduction (\code{inverse\_mod}'s \code{ENGINE\_CHECK}), never computed; a stray denominator prime to both $\ell$ and $\ell'$
would pass silently. \textbf{D0 obligation, the denominator check} (v1.23.0-pre.13, referee F5; not built at this version): at build time, every \code{Rat} denominator
in \code{diff\_} and in every $e_\sigma$ divides a power of $\mathsf q_\ell$ (the per-term denominators; \code{QGroupRingElt::max\_den} exists), else the build
refuses. (5) is the code's layout: cochains on the full
$\bigoplus_xM$ with $\delta^{q}$ and $P_q=\mathrm{diag}\,\rho_M(e_x)$, $\delta^{q}=P_{q+1}\delta^{q}P_q$ resting on \code{certify()}'s \code{well\_defined} check
($e_{\mathrm{col}}\,a\,e_{\mathrm{row}}=a$ in $\Q[\Ga]$, i.e.\ $a_{yx}=e_ya_{yx}e_x$); \code{cohomology\_lengths} returns
$\operatorname{len}P_q-\operatorname{len}\operatorname{im}\delta^{q}-\operatorname{len}\operatorname{im}\delta^{q-1}$.
\textbf{Not built} (D0's, \file{Integral/Descent/DescentAssembly} as planned): (i) the read-out of~(3), (6) on $\Iso^{(\ell)}$ --- \code{LocalModel} returns lengths, not
Smith valuations, and no code compares the count $s$ with $\rank_{\Q}\delta^{q}$ on $\Iso^{(\ell)}$ (on $\CI$ it is \code{HeckeFreeModelValidation}'s). The
$\F_{\ell'}$ rank of~(4)(c)(i) is that of the \textbf{same} \code{LocalModel}$(\mathrm{desc},\ell)$ object reduced into \code{ModCtx}$(\ell')$, $\ell'\nmid\mathsf q_\ell$; nothing
reads it yet. \textbf{D0 requirement} (v1.23.0-pre.13, referee F2; the header's audit note F-5: nothing ties \code{ell} to the \code{ModCtx} later passed in, so a
\code{LocalModel}$(\mathrm{desc},\ell')$ over \code{ModCtx}$(\ell')$ runs without complaint and is $\Iso^{(\ell')}$, inadmissible by~(4)(c)): the read-out guards that the
model's \code{ell()} is the $\ell$ whose $K$ is being certified and that the $\F_{\ell'}$ rank is read on that same object, never on another \code{LocalModel}, the census
or the shipped $C_7$ route; D0's validation prints, per cell, (4)(b)'s comparison $\rank P_0-\rank\delta^{0}$ over $\F_{\ell'}$ against $b_0$, and asserts
$\sigma_i\le v_\ell(24)$ for $q\ge1$ ((6)'s a-priori bound). (i$'$) (v1.23.0-pre.13, referee F7) The hand-off from \code{LocalModel}'s \code{engine::DenseBuf<Elt>} over a
\code{ModCtx} to \code{ring2k::Mat<uint64\_t>} and to \code{Ring3k}'s canonical \code{uint32\_t} input is not written. \emph{Degrees:} \code{LocalModel} inherits the
descriptor's construction parameter $q_{\max}$ (\file{Integral/Model/IsotropyModel.h}, default $5$, built to $q_{\mathrm{top}}=q_{\max}+1=6$; $d^{q}$ exists for
$q\le q_{\max}$); the plan's (P) arm at degree~$6$ needs $q_{\max}\ge5$, i.e.\ at least the default, never the production value $3$. (ii) At $\ell=2$ (v1.23.0-pre.13,
referee F6), the one-run read-out on \code{engine::ring2k} (rung \code{uint64\_t}, \file{src/ModArithmetic/Binary/Ring2k.h}) through $\Z/2^{W}$: \textbf{$W\le61$ is reachable
now} --- \code{LocalModel} is instantiated at \code{Elt = uint64\_t}, \code{ModCtx<uint64\_t>::prime\_power(2, W)} exists, \code{reduce\_rat} accepts moduli $<2^{62}$, and
\code{smith(A, k)} reads its input by masking (\file{Ring2k.h}: ``nothing ever reduces''), so these residues are valid input at $k=W$ and one run certifies every
exponent $<W$, as~(6)'s ``formed mod $2^{W}$, read mod $2^{K}$'' allows; \textbf{$W\in[62,64]$ is not built}: \code{reduce\_rat} refuses a modulus $\ge2^{62}$, and the reduction
$\mathsf O_2\to\Z/2^{W}$ (the odd denominators $3^{j}$ inverted to width $W$) is to be written. (iii) At $\ell=3$, \code{engine::ring3k} carries $K\le19$ only ($3^{19}<2^{31}$, \file{Ring3k.h}); a $3$-adic exponent
$\ge19$ is then a refusal under~(6) (``$K$ too small''), and no wider rung at $3$ exists. Both engines' Smith routines are specified in
\file{src/ModArithmetic/ENGINE\_API.md} \S\S8.2, 8.4 and gated by \code{Ring2kLadderValidation} and \code{OrbifoldMod3kValidation}; their correctness is the
hypothesis of (6)'s last sentence, not a consequence of it. By (6)'s a-priori bound both caps ($K\le19$ at $3$; $W\le61$, later $64$, at $2$) bind only at $d^{0}$;
reaching one is a refusal by the stopping rule. (iv) The plan's ``$\abs{Q}\in\{1,2\}$ at $\ell=3$'' (\file{DESCENT\_plan.md} \S2.7, row C-D0) is
corrected here to $\abs{Q_\sigma}\in\{2,4\}$ ($Q_{v_6}=\langle S\rangle$; \code{build\_columns} computes $Q_{v_6}$ of order $4$); the ring $\mathsf O_3=\Z[1/2]$ is the same.
The plan row is amended to match at v1.23.0-pre.13 (referee F8, \file{DESCENT\_plan.md} v0.3).

\paragraph{Formalization note.} Definition~\ref{def:isomodelell} and Proposition~\ref{prop:isomodelell} are homological algebra over $\cring[\Ga]$ (projective resolutions,
the comparison theorem, Shapiro's lemma); Proposition~\ref{prop:ellthreeclosed} adds the amalgam's Mayer--Vietoris sequence and the cohomology of cyclic groups. None
depends on modular forms. Proposition~\ref{prop:certifiedeell} (v1.23.0-pre.12) is base change and Smith forms over the principal ideal domains $\Z[1/\mathsf q_\ell]$ and
$\Z_{(\ell)}$; it carries \textbf{no} \code{\textbackslash LEAN} marker, since none of it is formalised.

% NEW at v1.23.0-pre.14 (hecke-extensions slot 24, milestone D1 of docs/implementation/slot24/DESCENT_plan.md v0.3, its clauses C-D1 and
% C-D2, the audit findings M8, m3, m4, m7 of DESCENT_plan_audit.md).  Placed after the last statement of sec:isomodelell, so that no statement
% number moves.  A pre-release draft; refereed ACCEPT-with-fixes (C-D1_referee.md), fixes F1-F7 applied at v1.23.0-pre.15, not
% re-refereed.  No \LEAN marker.
\subsection{The descent at level \texorpdfstring{$N$}{N}, and its bad primes}\label{sec:descentN}

(v1.23.0-pre.14; milestone D1 of \file{docs/implementation/slot24/DESCENT\_plan.md} v0.3, clauses C-D1 and C-D2 of its \S2.7, the audit findings M8, m3, m4,
m7 of \file{DESCENT\_plan\_audit.md}.) \S\ref{sec:descent} glues $H^{\bullet}(\Ga;\Vcoef)$ from one integral map; nothing in its proofs uses that the coefficient
module is $\Vcoef$ rather than any $\Z[\Ga]$-module that is finitely generated and free over $\Z$. This subsection states the descent in that generality ---
so for $\Ind_{\Ga_0(N)}^{\Ga}\Vcoef$ and its nebentypus twists, by Shapiro's lemma the level-$N$ groups --- with the two places where the generality changes
a statement: the coherence term of Proposition~\ref{prop:homvanish}, which survives only when $b_0b_1=0$, and the bad primes, which at level $N$ are not
bounded by $n-1$. (It sits in \S\ref{sec:isomodelell}'s section, after that section's last statement, only so that no statement number moves; it uses that
section's Proposition~\ref{prop:certifiedeell}, and nothing in it is specific to the $\ell$-local model. Referee F7, v1.23.0-pre.15.) \emph{Notation.} For a $\Z[\Ga]$-module $M$ and a ring $S$, $M_S:=M\otimes_\Z S$; $\CI(M):=\Hom_{\Z[\Ga]}(\Iso_{\bullet},M)$ and
$\Ceq(M):=\Hom_{\Z[\Ga]}(C_{\bullet}(\Xbulk),M)$, so that $\CI(\cring)$, $\Ceq(\cring)$ of \S\ref{sec:integral} are $\CI(\VR)$, $\Ceq(\VR)$ (the model complex
with fibre $M$, Theorem~\ref{thm:explicit_complex_equiv}(1)); $\mathrm T(M)$ is the tree complex \eqref{eq:treemodel}; $b_j(M):=\dim_{\Q}H^{j}(\Ga;M_{\Q})$;
$\mathsf O:=\Z[1/6]$. $M_N:=\Ind_{\Ga_0(N)}^{\Ga}\Vcoef$ is Definition~\ref{def:levelNmodel}'s fibre over $\Z$. For a character $\chi\colon G_N\to\Cx^{\times}$ of
order $m_\chi$, $\mathsf O_\chi:=\Z[\zeta_{m_\chi}]$ and $\chi$ is read with values in $\mu_{m_\chi}\subset\mathsf O_\chi^{\times}$; $V_{\mathsf O_\chi}(\chi)$ is
Definition~\ref{def:nebdata}(5) over $\cring=\mathsf O_\chi$ (Lemma~\ref{lem:nebdata}'s proof is ring-independent), and
$M_N(\chi):=\Ind_{\Ga_0(N)}^{\Ga}V_{\mathsf O_\chi}(\chi)$ in the coordinates \eqref{eq:MHchi}, free of rank $\varphi(m_\chi)m_N(n+1)$ over $\Z$; $M_N(1)=M_N$.
This is an integral form of \S\ref{sec:nebentypus}'s character model, whose coefficient rings are $\F_q$ and $W_K(\F_q)$ (item~(8) below).

\begin{proposition}[The descent for every finite free coefficient module; level $N$ and nebentypus]\label{prop:descentN}
\PROVED{} (v1.23.0-pre.14; refereed ACCEPT-with-fixes, \file{docs/implementation/slot24/C-D1\_referee.md}, the fixes applied at v1.23.0-pre.15, not
re-refereed.) Assume \ref{H4} and \ref{B1}--\ref{B3}, and let $M$ be a $\Z[\Ga]$-module that is finitely generated and free
over $\Z$ ($M_N$, $M_N(\chi)$, \dots).
\begin{enumerate}
\item \textbf{(The three complexes.)} $\CI(M)$ is a complex of finitely generated free $\Z$-modules with $\CI(M)\otimes S=\CI(M_S)$ and
  $H^{q}(\CI(M_S))=H^{q}(\Ga;M_S)$ for every ring $S$ and every $q$. $\Ceq(M)$ is a complex of finitely generated free $\Z$-modules in degrees $0,1,2$ with
  $H^{2}(\Ceq(M))=0$, and $\cmpI^{\ast}\colon\Ceq(M)\to\CI(M)$ is an injective cochain map whose cone has cohomology annihilated by $48$. Over $\mathsf O$:
  $\Ceq(M)[1/6]=\Ceq(M_{\mathsf O})$ and $\mathrm T(M)[1/6]=\mathrm T(M_{\mathsf O})$ are complexes of finitely generated free $\mathsf O$-modules, and for every
  $\mathsf O$-algebra $S$ --- $\Z/\ell^{K}$ and $\F_{\ell'}$ for primes $\ell,\ell'\ge5$ among them --- $\Ceq(M_{\mathsf O})\otimes_{\mathsf O}S=\Ceq(M_S)$,
  $\mathrm T(M_{\mathsf O})\otimes_{\mathsf O}S=\mathrm T(M_S)$, and both compute $H^{q}(\Ga;M_S)$ for every $q$ ($0$ for $q\ge2$).
\item \textbf{(The descent.)} With $A:=\Ceq(M_{\mathsf O})$, $B:=\CI(M)\otimes\Zsix$, $C:=\CI(M)\otimes\Qsix$ and $\alpha$, $\beta$ as in
  Proposition~\ref{prop:descent}, the mapping cocone of $\alpha-\beta$ is quasi-isomorphic to $\CI(M)$, and
  $H^{q}(\Ga;M)\cong H^{q}(A)\times_{H^{q}(C)}H^{q}(B)$ for every $q$. The only glueing datum is $\cmpI^{\ast}$. By Shapiro's lemma
  $H^{q}(\Ga;M_N)=H^{q}(\Ga_0(N);\Vcoef)$ and $H^{q}(\Ga;M_N(\chi))=H^{q}(\Ga_0(N);V_{\mathsf O_\chi}(\chi))$.
\item \textbf{(Maps and Hecke operators are determined locally.)} For a chain map $f$ between complexes of abelian groups (those of~(1), (2) among them),
  $H^{q}(f)$ is determined by $H^{q}(f[1/6])$ and $H^{q}(f\otimes\Zsix)$ (Proposition~\ref{prop:cohglue}'s morphism clause). An endomorphism $\phi$ of a finitely
  generated abelian group $G$ is determined by $\phi\otimes\Z[1/6]$ and $\phi\otimes\Zsix$; so every Hecke operator on $H^{q}(\Ga_0(N);\Vcoef)$ or
  $H^{q}(\Ga_0(N);V_{\mathsf O_\chi}(\chi))$, on whatever model it is computed, is determined by its two local components. A $\Z[\Ga]$-map $u\colon M\to M'$
  (the fibre maps of restriction, trace and $R_\sigma$; the $\mathsf O_\chi$-action on $M_N(\chi)$) induces $\Hom(\Iso_{\bullet},u)$ and
  $\Hom(C_{\bullet}(\Xbulk),u)$, which commute strictly with $\cmpI^{\ast}$. \textbf{No coherence datum is used}: neither (2) nor (3) needs
  Proposition~\ref{prop:homvanish} or~(4).
\item \textbf{(The term $\Hom^{-1}$, for $M$.)} Put $\mathcal N_M:=\CI(M)\otimes\Qsix$. Then
  $\Hom_{D(\Qsix)}(\mathcal N_M,\mathcal N_M[-1])\cong\Hom_{\Qsix}\bigl(H^{1}(\mathcal N_M),H^{0}(\mathcal N_M)\bigr)\cong\Qsix^{\,b_0(M)b_1(M)}$ and
  $\Hom_{D(\Qsix)}(\mathcal N_M,\mathcal N_M[-i])=0$ for $i\ge2$; so Proposition~\ref{prop:homvanish}(2) holds for $M$ \textbf{iff} $b_0(M)\,b_1(M)=0$. For
  $M=M_N(\chi)$: $b_0=1$ if $n=0$ and $\chi=1$, and $b_0=0$ otherwise; so the vanishing holds whenever $n>0$ or $\chi\ne1$ --- for $n>0$ with
  $\chi(-1)=(-1)^{n}$, and for $\chi(-1)\ne(-1)^{n}$, where $\mathcal N_M$ is acyclic --- and \textbf{fails} at $n=0$, $\chi=1$ as soon as $b_1>0$: at $N=2$,
  Proposition~\ref{prop:levelorbit}(3) over $\Q$ gives $b_1=1$ (the audit's finding M8), and for every $N>1$, at $n=0$, $\chi=1$,
  $6\,b_1=6+m_N-3\nu_2(N)-4\nu_3(N)\ge6$, so $b_1\ge1$ and the vanishing fails at every level $N>1$ (proved below, v1.23.0-pre.15; $b_1=1$ exactly at
  $N\in\{2,3,5,7,13\}$; nothing else in this proposition uses it).
\item \textbf{(Truncation and periodicity.)} Remark~\ref{rmk:truncdevice} holds for $M$ verbatim: what is glued is the untruncated $\CI(M)$, a local run
  builds $\sigma^{\le q_{\max}+1}$, and brutal truncation commutes with base change. For $q\ge2$, $e\cup-\colon H^{q}(\Ga;M)\to H^{q+2}(\Ga;M)$ is an
  isomorphism, natural in $M$ (Lemma~\ref{lem:bulkperiodic}(2)). \textbf{At chain level} (v1.23.0-pre.15): if $E\colon\Iso_{\bullet+2}\to\Iso_{\bullet}$ is a
  $\Z[\Ga]$-chain map with $\aug\circ E_0$ a cocycle representing $e$ (\file{Integral/Model/PeriodicityClass}'s $E$, Remark~\ref{rmk:integralevidence}), then for
  every $M$, $\Hom(E,M)\colon\CI(M)^{\bullet}\to\CI(M)^{\bullet+2}$ induces $e\cup-$ in every degree, and so, through Shapiro's isomorphism below,
  $\res^{\Ga}_{A}e\cup-$ on $H^{\bullet}(A;X)$ for $M=\Ind_A^{\Ga}X$. Let $A\le\Ga$ have finite index, $X$ a $\Z[A]$-module and
  $\mathrm{sh}\colon H^{\bullet}(\Ga;\Ind_A^{\Ga}X)\to H^{\bullet}(A;X)$ Shapiro's isomorphism; then $\mathrm{sh}(e\cup y)=\res^{\Ga}_{A}e\cup\mathrm{sh}(y)$. Hence
  \textbf{$\res^{\Ga}_{A}e\cup-\colon H^{q}(A;X)\to H^{q+2}(A;X)$ is an isomorphism for every $q\ge2$}, and $\res^{\Ga}_{C}e$ generates $H^{2}(C;\Z)$ for every
  finite cyclic $C\le A$ (Lemma~\ref{lem:bulkperiodic}(1)). \textbf{Hecke equivariance}: let $g\in\De$, $A'':=A\cap gAg^{-1}$, $A':=A\cap g^{-1}Ag$, both of
  finite index, $\iota_g\colon A''\to A'$, $h\mapsto g^{-1}hg$, and $f\colon X\to X$ additive with $f(g^{-1}hg\cdot x)=h\cdot f(x)$ for $h\in A''$ (a compatible
  pair; $f=\rhon(g)$ for $X=\Vcoef$, $f=\rho_{n,\chi}(g)$, $g\in\De_0(N)$, for $X=V_{\mathsf O_\chi}(\chi)$). For
  $[AgA]:=\cor^{A}_{A''}\circ(\iota_g,f)^{\ast}\circ\res^{A}_{A'}$,
  \[
    [AgA]\bigl(\res^{\Ga}_{A}e\cup x\bigr)=\res^{\Ga}_{A}e\cup[AgA](x)\qquad(x\in H^{q}(A;X),\ q\ge0).
  \]
  In particular, for $A=\Ga_0(N)$, $T_p$ ($p\nmid N$) on $H^{\bullet}(\Ga_0(N);\Vcoef)$ and on $H^{\bullet}(\Ga_0(N);V_{\mathsf O_\chi}(\chi))$ commutes with
  $\res^{\Ga}_{A}e\cup-$, which for $q\ge2$ is an isomorphism of Hecke modules; $q_{\max}=3$ determines every degree $\ge2$ at level~$N$, Hecke-equivariantly.
  The argument does not use $p\nmid N$: it covers $U_p:=[\Ga_0(N)\smat100p\Ga_0(N)]$, $p\mid N$, with $f=\rhon(\smat100p)$, resp.\ $\rho_{n,\chi}(\smat100p)$ --- the
  double coset of Proposition~\ref{prop:levelcorrespondence}, resp.\ \ref{prop:nebcorr}, at $A=B=\Ga_0(N)$ --- on cohomology as well (its cochain operator is
  deferred, [R-Q4] of the plan; lifted for the free model by Proposition~\ref{prop:strictliftN}(9), v1.26.0-pre.1).
\item \textbf{(The good prime.)} A prime $\ell'\ge5$ \emph{qualifies for $d^{q}$} if $H^{q+1}(\Ga;M)$ has no $\ell'$-torsion; then
  $\rank_{\Q}d^{q}=\rank_{\F_{\ell'}}d^{q}$ on each of $\CI(M)$, $\Ceq(M_{\mathsf O})$, $\mathrm T(M_{\mathsf O})$ and $C^{\bullet}_\ell$ of
  Proposition~\ref{prop:certifiedeell} (every $\ell'\ge5$ qualifies for $d^{q}$, $q\ge1$, by Proposition~\ref{prop:certifiedeell}(4)(a)). \textbf{For $M_N$}
  ($\chi=1$, every $N$, every $n$) every prime $\ell'\ge5$ with $\ell'>n+1$ qualifies for every $d^{q}$ (Proposition~\ref{prop:badprimesN}(2)); the plan's rule
  ``$\ell'>n+1$ and $\ell'\nmid6N$'' is a subset, its condition $\ell'\nmid N$ not needed for $M_N$. \textbf{For $M_N(\chi)$}, $\chi\ne1$, every prime $\ell'\ge5$ with
  $\ell'\nmid m_\chi N$ qualifies for every $d^{q}$ (Proposition~\ref{prop:badprimesN}(3)(a)), and $\ell'\nmid N$ is needed: $\ell'\mid N$, $\ell'>n+1$ may carry
  torsion (the example of Proposition~\ref{prop:badprimesN}(3)). One rule valid for every $\chi$: $\ell'\ge5$, $\ell'>n+1$, $\ell'\nmid m_\chi N$.
\item \textbf{(The certified exponent on the $\mathsf O$-piece at $\ell\ge5$.)} Let $\ell\ge5$ and $K\ge1$.
  \begin{enumerate}
  \item[(a)] \textbf{Direct Smith form on the Bredon complex.} Proposition~\ref{prop:certifiede}(1)--(3) and Remark~\ref{rmk:certifiedeperdiff} hold for
    $\Ceq(M_{\mathsf O})\otimes_{\mathsf O}\Z_{(\ell)}$ in place of a complex of free $\Z$-modules, at the prime $\ell$; its reduction mod $\ell^{K}$ is
    $\Ceq(M_{\Z/\ell^{K}})$, its target $C^{1}$ written in a basis over $\Z/\ell^{K}$ of the free modules $(M_{\Z/\ell^{K}})^{\Ga_\sigma}$, its source in a basis or
    by any generating family, as in~(b) (only the image of $d^{0}$ enters). $d^{1}$ is onto (its non-zero elementary
    divisors are units), so all $\ell$-torsion is read from $d^{0}$: if the number $s$ of non-zero diagonal entries $\ell^{\sigma_i}$ of the Smith form of
    $d^{0}\bmod\ell^{K}$ equals $\rank_{\Q}d^{0}=\rank C^{0}_{\mathrm{eq}}(M_{\mathsf O})-b_0(M)$, then $\mathrm{tors}\,H^{1}(\Ga;M)\otimes\Z_{(\ell)}\cong
    \bigoplus_{\sigma_i>0}\Z/\ell^{\sigma_i}$ exactly; if $s<\rank_{\Q}d^{0}$, $K$ is too small, and the stopping rule of Proposition~\ref{prop:certifiedeell}(6)
    applies (raise $K$; at the cap of the ring, refuse). No a-priori bound on the exponents of $d^{0}$ is asserted (Remark~\ref{rmk:bocksteinexp}).
  \item[(b)] \textbf{The sparse cokernel route at level $N$} (\textbf{[m7]}). Let $\bar R$ be any matrix over $\Z/\ell^{K}$, written in a basis of the free module
    $(M_{\Z/\ell^{K}})^{\langle-I\rangle}$, whose rows \emph{generate} the submodule $(M_{\Z/\ell^{K}})^{\langle S\rangle}+(M_{\Z/\ell^{K}})^{\langle U\rangle}$ ---
    for instance the columns of the averaging idempotents, or the orbit-wise generators of Propositions~\ref{prop:levelorbit}(1), \ref{prop:neborbit}(1)
    provided that at each fixed point they generate $(V_{\Z/\ell^{K}})^{F_j}$ over $\Z/\ell^{K}$ (a kernel basis computed over $\F_\ell$ need not) --- and $s$ the
    number of non-zero diagonal entries $\ell^{\sigma_i}$ of its Smith form. Then $s=r_K(d^{0})$ for $d^{0}$ of $\mathrm T(M_{\mathsf O})$; if
    $s=\rank\mathrm T^{1}-b_1(M)\;\bigl(=\rank\mathrm T^{0}-b_0(M)\bigr)$, then $\mathrm{tors}\,H^{1}(\Ga;M)\otimes\Z_{(\ell)}\cong\bigoplus_{\sigma_i>0}\Z/\ell^{\sigma_i}$
    exactly; otherwise $K$ is too small, with the same stopping rule. For $M_N$, $n$ even: $\rank\mathrm T^{1}=m_N(n+1)$ and $b_1=b_0+{}$the right-hand side of
    Proposition~\ref{prop:levelorbit}(3) over $\Q$; the comparison value is this closed form, from Proposition~\ref{prop:levelorbit}(1)'s orbit counts, never the
    number of rows of $\bar R$. \emph{The certificate that the rows generate} (v1.23.0-pre.15, referee F4): the averaging-idempotent columns need none ($e_Fx$ is
    fixed, and $x=e_Fx$ for fixed $x$); any other family of vectors in block $j$ generates $(V_{\Z/\ell^{K}})^{F_j}$ iff (i) each is fixed mod $\ell^{K}$ by a
    generator of the cyclic group $F_j$, and (ii) their reductions mod $\ell$ span $(V_{\F_\ell})^{F_j}$, of dimension $d_{F_j}$.
  \item[(c)] \textbf{Proposition~\ref{prop:certifiedeell}(4) at level $N$.} (4)(a) holds for every $M$ as stated. (4)(b)'s values extend: $b_0(M_N(\chi))=1$ if
    $n=0$ and $\chi=1$, and $0$ otherwise; a prime $\ell'\ge5$ qualifies for $d^{0}$ iff $\ell'\notin\mathcal B(M)$ of Proposition~\ref{prop:badprimesN}(1). (4)(c)'s
    two admissible sources apply on $\Ceq(M_{\mathsf O})$ and on $\mathrm T(M_{\mathsf O})$ with ``this complex'' read as the complex whose differential is
    being certified; the rank of one of $\CI$, $\Ceq$, $\mathrm T$, $C^{\bullet}_\ell$ is not admissible for another.
  \end{enumerate}
  (At $\ell=3$, (7)(b) holds verbatim for $M_N$ when $\nu_3(N)=0$: Proposition~\ref{prop:descentNthree}(3), v1.23.0-pre.17.)
\item \textbf{(Nebentypus, $\ell$-locally.)} Let $\ell$ be a prime, $\ell\nmid6m_\chi$, and $K\ge1$. Then
  $M_N(\chi)\otimes\Z/\ell^{K}\cong\bigoplus_{\lambda\mid\ell}\Ind_{\Ga_0(N)}^{\Ga}V_{W_\lambda}(\mathfrak t\circ\bar\chi_\lambda)$ as $\Z/\ell^{K}[\Ga]$-modules, the sum over
  the primes $\lambda$ of $\mathsf O_\chi$ above $\ell$, with $W_\lambda\cong W_K(\F_{q_\lambda})$, $q_\lambda:=\abs{\mathsf O_\chi/\lambda}$, and
  $\bar\chi_\lambda:=\chi\bmod\lambda\colon G_N\to\F_{q_\lambda}^{\times}$. So the $\ell$-primary part of $H^{\bullet}(\Ga_0(N);V_{\mathsf O_\chi}(\chi))$ is read on
  \S\ref{sec:nebentypus}'s character models over $W_K(\F_{q_\lambda})$, one per $\lambda$, at the characters $\bar\chi_\lambda$.
\end{enumerate}
\end{proposition}

\begin{proof}
(1) $\Iso_{q}$ is finitely generated free over $\Z[\Ga]$ (Proposition~\ref{prop:isomodel}(2)), so $\CI(M)\cong M^{w_q}$ in degree $q$, $\CI(M)\otimes S=\CI(M_S)$,
and $H^{q}(\CI(M_S))=\Ext^{q}_{\Z[\Ga]}(\Z,M_S)=H^{q}(\Ga;M_S)$ --- the proof of Proposition~\ref{prop:isomodel}(3), which uses nothing about $\VR$.
$\Ceq(M)=\bigoplus_{\sigma}M^{\Ga_\sigma}$, a submodule of a finitely generated free $\Z$-module in each degree, hence free; $H^{2}(\Ceq(M))=0$ by
Proposition~\ref{prop:bredon}(2), which is stated for every $\Z[\Ga]$-module. The injectivity and the bound $48$ are the proof of
Proposition~\ref{prop:isomodel}(4) with $M$ for $\VR$: the graded pieces of the cone have cohomology $H^{t}(\Ga_\sigma;M)$, $t\ge1$, annihilated by $\abs{\Ga_\sigma}$
for every module \STD{\cite[Ch.~III.10.2]{brown}}. Over $\mathsf O$: invariants are a kernel and $\mathsf O$ is flat, so $M^{\Ga_\sigma}[1/6]=M_{\mathsf O}^{\Ga_\sigma}$;
as $\abs{\Ga_\sigma}\mid12$ is a unit, $M_{\mathsf O}^{\Ga_\sigma}=e_\sigma M_{\mathsf O}$ for the averaging idempotent $e_\sigma$, a direct summand of a finitely
generated free $\mathsf O$-module, hence free over the principal ideal domain $\mathsf O$, and $e_\sigma M_{\mathsf O}\otimes S=e_\sigma M_S=M_S^{\Ga_\sigma}$ (the image of an
idempotent commutes with base change); the differentials, given by group-ring coefficients acting on $M$, are base-changed with them. The same holds for
$\mathrm T$, whose terms are $M^{\langle S\rangle}\oplus M^{\langle U\rangle}$ and $M^{\langle-I\rangle}$. $\Ceq(M_S)$ computes $H^{q}(\Ga;M_S)$ because the cone of
$\cmpI^{\ast}$ over $S$ is annihilated by $48\in S^{\times}$; $\mathrm T(M_S)$ does by Lemma~\ref{lem:treeres} with $A_0=\Ga$ and $\cring=S$; and $H^{q}(\Ga;M_S)=0$
for $q\ge2$ by the same lemma.

(2) The proof of Proposition~\ref{prop:descent} with (1) in place of Proposition~\ref{prop:isomodel}(3),(4): $\cmpI^{\ast}[1/6]\colon A\to\CI(M)[1/6]$ is a
quasi-isomorphism, the induced map of cocones is one by the five lemma, and Propositions~\ref{prop:fracture} and~\ref{prop:cohglue} apply to $\CI(M)$, a
complex of abelian groups. Shapiro: \STD{\cite[Ch.~III.6.2]{brown}}, induction and coinduction agreeing at finite index, as in the proof of
Proposition~\ref{prop:levelorbit}(2).

(3) The first sentence is Proposition~\ref{prop:cohglue}. For the second apply Proposition~\ref{prop:cohglue} to $G$ placed in degree $0$: $G\to G[1/6]\oplus
(G\otimes\Zsix)$ is injective, and $\phi$ commutes with it. A Hecke operator is an endomorphism of the finitely generated group $H^{q}(\Ga;M)$
(finitely generated by (1)). $\cmpI^{\ast}=\Hom(\cmpI,M)$, and $\Hom(-,u)\circ\Hom(\cmpI,M)=\Hom(\cmpI,M')\circ\Hom(-,u)$ is functoriality of $\Hom$ in two variables.

(4) $\Qsix=\Q_2\times\Q_3$ is a product of two fields, so every $\Qsix$-module is projective and injective, and a complex $X$ of $\Qsix$-modules is
isomorphic in $D(\Qsix)$ to $\bigoplus_qH^{q}(X)[-q]$ (split $X^{q}=Z^{q}\oplus X^{q}/Z^{q}$ and $Z^{q}=B^{q}\oplus H^{q}$). Hence
$\Hom_{D}(X,X[-i])=\prod_{q,p}\Ext^{q-p-i}_{\Qsix}(H^{q},H^{p})=\prod_q\Hom_{\Qsix}(H^{q},H^{q-i})$, since $\Ext^{j}$ vanishes for $j\ne0$ over a semisimple ring
\STD{\cite[\S10.7]{weibel}}. For $X=\mathcal N_M$, $H^{q}(\mathcal N_M)=H^{q}(\Ga;M)\otimes\Qsix\cong\Qsix^{\,b_q}$ (flatness; $H^{q}$ finitely generated), and $b_q=0$
for $q\ge2$ (Proposition~\ref{prop:certifiedeell}(4)(a)); so only $q=1$, $i=1$ survives, and $\Hom_{\Qsix}(\Qsix^{b_1},\Qsix^{b_0})\cong\Qsix^{b_0b_1}$. \emph{The values of
$b_0$.} $H^{0}(\Ga;M_N(\chi)_{\Q})=H^{0}(\Ga_0(N);V(\chi)_{\Q})$ (Shapiro in degree $0$). For $n>0$, $\Ga_1(N)\subseteq\Ga_0(N)$ acts on $V(\chi)$ through $\rhon$
(Lemma~\ref{lem:nebdata}(4)) and contains $\smat1101$ and $\smat10N1$; the argument of Proposition~\ref{prop:certifiedeell}(4)(b), over the field $\Q(\zeta_{m_\chi})$ of
characteristic $0$, leaves no non-zero fixed vector. For $n=0$, $V(\chi)=\mathsf O_\chi$ with $\gamma$ acting by $\chi(a_\gamma)$, and $\gamma\mapsto a_\gamma\bmod N$ is
onto $G_N$ (Lemma~\ref{lem:nebdata}(1)), so the invariants vanish unless $\chi=1$, where they are $\Q$. If $\chi(-1)\ne(-1)^{n}$, $-I$ acts on the coefficients by
$-1$ and trivially on cohomology, so $2$ kills $H^{\bullet}(\Ga_0(N);V_{\mathsf O_\chi}(\chi))$ (the proof of Proposition~\ref{prop:neborbit}(2), over $\Z$) and every
$b_j=0$. At $N=2$, $n=0$: $m_2=3$, $\nu_2(2)=1$ ($(1:1)$ is the fixed point of $S$), $\nu_3(2)=0$, and Proposition~\ref{prop:levelorbit}(3) gives $\dim H^{1}-\dim
H^{0}=3-2-1=0$, so $b_1=b_0=1$. \emph{Every $N>1$} (v1.23.0-pre.15, referee F2). At $n=0$, $d_S=d_{ST}=1$ and $b_0=1$, so
Proposition~\ref{prop:levelorbit}(3) over $\Q$ gives $b_1=1+m_N/6-\nu_2/2-2\nu_3/3$, and $b_1\ge1$ iff $m_N\ge3\nu_2+4\nu_3$. By the Chinese remainder theorem
$\mathbb P^1(\Z/N)=\prod_{p^{e}\Vert N}\mathbb P^1(\Z/p^{e})$ compatibly with $\Ga$, so $m_N$, $\nu_2$, $\nu_3$ are multiplicative, $m_{p^{e}}=p^{e-1}(p+1)$. Let $g\in\{S,ST\}$,
or its transpose (whichever side \eqref{eq:MH}'s action uses; $g$ and $g^{\mathsf T}$ have the same characteristic polynomial $\chi_g$, $\chi_S=X^{2}+1$,
$\chi_{ST}=X^{2}-X+1$, and each has an off-diagonal entry $\pm1$). A point of $\mathbb P^1(\Z/p^{e})$ fixed by $g$ is a line spanned by a primitive $v$ with $gv=uv$;
then $\chi_g(u)v=\mathrm{adj}(g-u)(g-u)v=0$ and $v$ has a unit coordinate, so $\chi_g(u)=0$. Conversely, for a root $u$, $g-u$ has determinant $0$ and an entry
$\pm1$, so its kernel is free of rank $1$, spanned by a primitive vector (solve the row with the entry $\pm1$; the other row is then automatic); and the line
determines $u$. So $\nu_2(p^{e})$, $\nu_3(p^{e})$ are the numbers of roots of $X^{2}+1$, $X^{2}-X+1$ in $\Z/p^{e}$. At a root the derivative is a unit when $p\ne2$,
resp.\ $p\ne3$ ($(2u)^{2}=-4$, $(2u-1)^{2}=-3$), so Hensel's lemma gives $\nu_2(p^{e})=1+\bigl(\frac{-1}{p}\bigr)$ for $p$ odd and
$\nu_3(p^{e})=1+\bigl(\frac{-3}{p}\bigr)$ for $p\ne2,3$; directly, $\nu_2(2)=1$, $\nu_2(2^{e})=0$ for $e\ge2$ ($u^{2}\not\equiv-1\bmod4$), $\nu_3(2^{e})=0$ ($X^{2}+X+1$ has
no root mod~$2$), $\nu_3(3)=1$ and $\nu_3(3^{e})=0$ for $e\ge2$ (the one root $-1$ mod~$3$, and $\chi_{ST}(-1+3t)\equiv3\bmod9$). Let $k\ge1$ be the number of primes
dividing $N$. If $\nu_2\nu_3\ne0$, every local factor of both is non-zero, which excludes $p=2,3$ and forces $p\equiv1\pmod{12}$: $\nu_2=\nu_3=2^{k}$ and
$m_N\ge14^{k}\ge7\cdot2^{k}=3\nu_2+4\nu_3$. If $\nu_2\ne0=\nu_3$, every $p^{e}\Vert N$ is $2$ or has $p\equiv1\pmod4$, and $m_{p^{e}}\ge3\nu_2(p^{e})$ ($3=3\cdot1$,
resp.\ $p+1\ge6=3\cdot2$), so $m_N\ge3^{k}\nu_2\ge3\nu_2$. If $\nu_3\ne0=\nu_2$, every $p^{e}\Vert N$ is $3$ or has $p\equiv1\pmod3$, and $m_{p^{e}}\ge4\nu_3(p^{e})$
($4=4\cdot1$, resp.\ $p+1\ge8=4\cdot2$), so $m_N\ge4\nu_3$. If $\nu_2=\nu_3=0$ there is nothing to prove. (Equality, $b_1=1$, \emph{exactly} at $N=2,3,5,7,13$ (v1.23.0-pre.16, referee D-5: in each case above equality forces $k=1$ and equality in the one local factor; the referee's direct count to $N=2000$ agrees); at $N=1$,
$6b_1=6+1-3-4=0$. The referee's numerical check, $2\le N\le20000$, minimum $6b_1=6$: \file{C-D1\_referee.md} \S2; the local counts against
direct fixed points, $N\le120$: \file{docs/implementation/slot24/cd1fix\_f2\_check.py}. Not tags.)

(5) The first part is Remark~\ref{rmk:truncdevice}'s argument, which uses of $\CI(\Z)$ only that it is a complex of free modules, and
Lemma~\ref{lem:bulkperiodic}(2), stated for every $\Z[\Ga]$-module. \emph{Chain level.} $E^{\ast}:=\Hom(E,-)$ is a chain map $\CI(M)\to\CI(M)[2]$, natural in $M$; as
each $\Iso_{q}$ is free, $\Hom_{\Z[\Ga]}(\Iso_{q},-)$ is exact, so a short exact sequence of modules gives one of complexes, $E^{\ast}$ maps it to the shifted one, and the
induced maps commute with the connecting maps. $e\cup-$ is natural in $M$ and commutes with the connecting maps as well (\ref{Ecup}: the Leibniz rule, $e$ a cocycle
of even degree, \STD{\cite[Ch.~V, \S3]{brown}, section level, not verified against the text, rule~\ref{A2}}). In degree $0$ they agree: $x\in M^{\Ga}$ is the cocycle
$x\circ\aug$, and $x\circ\aug\circ E_0$ is the image of $e$ under $x_\ast\colon\Z\to M$, $1\mapsto x$, which is $e\cup x$ (\ref{Ecup}'s naturality in the second
variable and its unit). Dimension shifting: for $0\to M\to\mathrm{Coind}\,M\to M''\to0$ with $H^{q}(\Ga;\mathrm{Coind}\,M)=0$, $q\ge1$, the connecting map
$H^{q}(\Ga;M'')\to H^{q+1}(\Ga;M)$ is onto for every $q\ge0$, so two such families that agree in degree $q$ on every module agree in degree $q+1$. Shapiro then
transports $e\cup-$ to $\res^{\Ga}_{A}e\cup-$ by the identity below. The Shapiro identity is~\ref{Pshapiro} of that lemma's proof, at $G=\Ga$, $H=A$,
$u=e$ ($\Ind=\mathrm{Coind}$ at finite index, \ref{Ecor}). So $\res^{\Ga}_Ae\cup-=\mathrm{sh}\circ(e\cup-)\circ\mathrm{sh}^{-1}$, an isomorphism for $q\ge2$. The
generator statement is Lemma~\ref{lem:bulkperiodic}(1) for $C\le A\le\Ga$. \emph{Equivariance.} $\iota_g(A'')\subseteq A'$: for $h\in A''$, $g^{-1}hg\in A$ as
$h\in gAg^{-1}$, and $g^{-1}hg\in g^{-1}Ag$ as $h\in A$. Since $A''\subseteq\Ga'':=\Ga\cap g\Ga g^{-1}$, functoriality of pull-backs and
Lemma~\ref{lem:bulkperiodic}(3)(a) give $\iota_g^{\ast}\res^{\Ga}_{A'}e=\res^{\Ga''}_{A''}\iota_g^{\ast}e=\res^{\Ga''}_{A''}\res^{\Ga}_{\Ga''}e=\res^{\Ga}_{A''}e$. Then, by
\ref{Pnat} (twice: for $\res^{A}_{A'}$ and for the compatible pair $(\iota_g,f)$) and the projection formula \ref{Pproj},
$[AgA](\res^{\Ga}_Ae\cup x)=\cor^{A}_{A''}\bigl(\iota_g^{\ast}\res^{\Ga}_{A'}e\cup(\iota_g,f)^{\ast}\res^{A}_{A'}x\bigr)
=\cor^{A}_{A''}\bigl(\res^{A}_{A''}(\res^{\Ga}_Ae)\cup y\bigr)=\res^{\Ga}_Ae\cup\cor^{A}_{A''}y$ with $y:=(\iota_g,f)^{\ast}\res^{A}_{A'}x$; this is the proof of
Lemma~\ref{lem:bulkperiodic}(3)(b) with $A$ for $\Ga$ and $\res^{\Ga}_Ae$ for $e$. $T_p$ at level $N$ is $[AgA]$ with $g=\smat100p$ and $f=\rhon(g)$, resp.\
$\rho_{n,\chi}(g)$ (Propositions~\ref{prop:levelcorrespondence}, \ref{prop:nebcorr}; $\smat100p\in\De_0(N)$ for every $p$), and nothing used $p\nmid N$. The
consequences for $q_{\max}$ are those of Remark~\ref{rmk:truncdevice}.

(6) Proposition~\ref{prop:certifiedeell}(4) and its proof, over $\Z$ for $\CI(M)$ and over $\mathsf O$ for $\Ceq(M_{\mathsf O})$, $\mathrm T(M_{\mathsf O})$: a free
complex over a principal ideal domain reduces mod a prime $\ell'$ (not inverted) with the same ranks iff $\ell'$ divides no elementary divisor, i.e.\ iff the
torsion of $H^{q+1}$ of the complex has no $\ell'$-part; for $\Ceq(M_{\mathsf O})$ and $\mathrm T(M_{\mathsf O})$ that cohomology is $H^{q+1}(\Ga;M)\otimes\mathsf O$ by
(1). The rest is Proposition~\ref{prop:badprimesN}: for $M_N$, every $\ell'\ge5$ outside $\mathcal B(n,N)$ carries no torsion and every element of
$\mathcal B(n,N)\cap[5,\infty)$ is $\le n+1$ ((2) there, and $n=0$, $n$ odd); for $\chi\ne1$, (3)(a) there gives $s_{\ell'}=0$ for $\ell'\nmid6m_\chi N$, and the
example in (3) has $\ell'=5\mid N$, $\ell'>n+1=3$.

(7)(a) $\Ceq(M_{\mathsf O})\otimes\Z_{(\ell)}$ is a complex of finitely generated free $\Z_{(\ell)}$-modules (by (1), $\Z_{(\ell)}$ an $\mathsf O$-algebra for
$\ell\ge5$) with cohomology $H^{q}(\Ga;M)\otimes\Z_{(\ell)}$; the proof of Proposition~\ref{prop:certifiede}(1)--(3) uses of $\Z$ only Smith forms over $\Z_{(\ell)}$,
torsion-freeness of the terms and the universal coefficient theorem, exactly as in the proof of Proposition~\ref{prop:certifiedeell}(3). Its reduction mod
$\ell^{K}$ is $\Ceq(M_{\Z/\ell^{K}})$ by (1). A basis over $\Z/\ell^{K}$ of the free module $(M_{\Z/\ell^{K}})^{\Ga_\sigma}=e_\sigma M_{\Z_{(\ell)}}\otimes\Z/\ell^{K}$ lifts to a
basis of $e_\sigma M_{\Z_{(\ell)}}$ (Nakayama over the local ring $\Z_{(\ell)}$, the lifts taken in $e_\sigma M_{\Z_{(\ell)}}$, i.e.\ after applying $e_\sigma$), so the matrix of $d^{0}\bmod\ell^{K}$ in such bases is the reduction of the matrix of
$d^{0}$ in bases over $\Z_{(\ell)}$, and its Smith form over $\Z/\ell^{K}$ is $\mathrm{diag}(\ell^{\min(a_t,K)},0)$ (the proof of Proposition~\ref{prop:certifiede}(1)).
$H^{2}(\Ceq(M_{\mathsf O}))=0$ means $d^{1}$ is onto, so $\coker d^{1}$ has no torsion and every non-zero elementary divisor of $d^{1}$ is a unit;
$\mathrm{tors}\,H^{1}$ is $\mathrm{tors}\coker d^{0}$ (Proposition~\ref{prop:certifiede}(2)). For a generating family of the source the image of $d^{0}\bmod\ell^{K}$, so its
cokernel, is unchanged, and (b)'s comparison of multiplicities applies verbatim. $\rank_{\Q}d^{0}=\rank C^{0}_{\mathrm{eq}}-b_0$ is rank--nullity over $\Q$, the kernel
of $d^{0}\otimes\Q$ being $H^{0}(\Ga;M_{\Q})$ by (1); it is source (ii) of Proposition~\ref{prop:certifiedeell}(4)(c). The certificate and the stopping rule are
Proposition~\ref{prop:certifiede}(1),(3) with Remark~\ref{rmk:certifiedeperdiff} at $e_0=K$, and Proposition~\ref{prop:certifiedeell}(6)'s last paragraph; the absence of
an a-priori bound at $d^{0}$ is Proposition~\ref{prop:badprimes}'s closing note, which applies to $M=\Vcoef$, a case of $M$.

(7)(b) By (1), $\mathrm T(M_{\mathsf O})\otimes\Z/\ell^{K}=\mathrm T(M_{\Z/\ell^{K}})$, so the image of $d^{0}\bmod\ell^{K}$ is the submodule
$\mathsf R:=(M_{\Z/\ell^{K}})^{\langle S\rangle}+(M_{\Z/\ell^{K}})^{\langle U\rangle}$, and $\coker\bar R=(M_{\Z/\ell^{K}})^{\langle-I\rangle}/\mathsf R=\coker(d^{0})\otimes
\Z/\ell^{K}$ (right exactness). Over $\Z_{(\ell)}$, $\coker d^{0}\cong\Z_{(\ell)}^{D-r}\oplus\bigoplus_t\Z/\ell^{a_t}$ ($D:=\rank\mathrm T^{1}$, $r:=\rank_{\Q}d^{0}$, the $a_t\ge1$
the non-unit elementary divisors), so $\coker\bar R\cong(\Z/\ell^{K})^{D-r+\#\{a_t\ge K\}}\oplus\bigoplus_{a_t<K}\Z/\ell^{a_t}$. A Smith form of $\bar R$ with $s$ non-zero
diagonal entries $\ell^{\sigma_i}$ gives $\coker\bar R\cong(\Z/\ell^{K})^{D-s}\oplus\bigoplus_i\Z/\ell^{\sigma_i}$. Over the local principal ideal ring $\Z/\ell^{K}$ a finitely
generated module is a direct sum of cyclic modules $\Z/\ell^{c}$, $1\le c\le K$, with uniquely determined multiplicities; comparing the multiplicity of
$c=K$ gives $s=r-\#\{a_t\ge K\}=r_K(d^{0})$, and comparing $c<K$ gives $\{\sigma_i>0\}=\{a_t<K\}$ as multisets. So $s=r$ iff every $a_t<K$, and then the torsion is
read off. $r=\rank\mathrm T^{0}-b_0=\rank\mathrm T^{1}-b_1$ is rank--nullity over $\Q$ for the two-term complex $\mathrm T(M_{\Q})$, whose cohomology is $b_0$, $b_1$ by
(1). For $M_N$ with $n$ even, $M_N(-I)=1$, so $\mathrm T^{1}=M_N$, and Proposition~\ref{prop:levelorbit}(3) over the field $\Q$ gives $b_1-b_0$. The orbit-wise
generators generate $\mathsf R$ iff they generate each orbit's fixed module, which on a free orbit holds by Proposition~\ref{prop:levelorbit}(1) with no solve, and
at a fixed point is the stated condition. \emph{The certificate.} As $\ell\nmid\abs{F_j}$ ($\abs{F_j}\mid12$), $W:=(V_{\Z/\ell^{K}})^{F_j}=e_{F_j}V_{\Z/\ell^{K}}$ is free with
$W/\ell W=e_{F_j}V_{\F_\ell}=(V_{\F_\ell})^{F_j}$ (the image of an idempotent commutes with base change, as in (1)). A vector fixed by a generator of $F_j$ lies in $W$;
if the reductions of a family in $W$ span $W/\ell W$, the family generates $W$ by Nakayama's lemma over the local ring $\Z/\ell^{K}$; conversely a generating family
of $W$ lies in $W$ and its reductions span $W/\ell W$.

(7)(c) (4)(a) of Proposition~\ref{prop:certifiedeell} is stated for every $M$. The values of $b_0$ are proved in (4) above. A prime $\ell'\ge5$ qualifies for
$d^{0}$ iff $H^{1}(\Ga;M)$ has no $\ell'$-torsion, i.e.\ iff $s_{\ell'}(M)=0$ (Proposition~\ref{prop:badprimesN}(1)). The admissibility statement is
Proposition~\ref{prop:certifiedeell}(4)(c)'s reason: ranks of differentials depend on the complex, Betti numbers do not.

(8) $\mathsf O_\chi=\Z[Y]/(\Phi_{m_\chi})$. As $\ell\nmid m_\chi$, $\Phi_{m_\chi}\bmod\ell$ divides the separable $Y^{m_\chi}-1$, so it is a product of pairwise distinct
monic irreducibles $\bar g_\lambda$, one per $\lambda\mid\ell$, of degree $f_\lambda$ with $\ell^{f_\lambda}=q_\lambda$; by Hensel's lemma
$\Phi_{m_\chi}\equiv\prod_\lambda g_\lambda\pmod{\ell^{K}}$ with monic, pairwise comaximal lifts $g_\lambda$, and the Chinese remainder theorem gives
$\mathsf O_\chi/\ell^{K}\cong\prod_\lambda W_\lambda$, $W_\lambda:=(\Z/\ell^{K})[Y]/(g_\lambda)$, a local ring free of rank $f_\lambda$ over $\Z/\ell^{K}$ with residue field
$\F_{q_\lambda}$ (as in Lemma~\ref{lem:teichmuller}(1)). \emph{$W_\lambda\cong W_K(\F_{q_\lambda})$}: $\bar g_\lambda$ has a simple root in $\F_{q_\lambda}$, which lifts by
Hensel's lemma to a root of $g_\lambda$ in $W_K(\F_{q_\lambda})$; the induced ring map $W_\lambda\to W_K(\F_{q_\lambda})$ is an isomorphism modulo $\ell$ (an
injective map of fields of the same size), hence onto by Nakayama's lemma, hence bijective (free modules of the same rank). $\Ga$ acts
$\mathsf O_\chi$-linearly on $M_N(\chi)$, so $M_N(\chi)\otimes\Z/\ell^{K}=M_N(\chi)\otimes_{\mathsf O_\chi}(\mathsf O_\chi/\ell^{K})=\bigoplus_\lambda
M_N(\chi)\otimes_{\mathsf O_\chi}W_\lambda$, and the coordinates \eqref{eq:MHchi} show $M_N(\chi)\otimes_{\mathsf O_\chi}W_\lambda=\Ind V_{W_\lambda}(\chi_\lambda)$, $\chi_\lambda$
the composite $G_N\to\mathsf O_\chi^{\times}\to W_\lambda^{\times}$. Its order divides $m_\chi$, prime to $\ell$, so $\chi_\lambda=\mathfrak t\circ\bar\chi_\lambda$ by
Lemma~\ref{lem:teichmuller}(3). Cohomology is additive in the coefficients.
\end{proof}

\noindent
\emph{Reading.} (2) and (3) are the level-$N$ descent with the Hecke action determined locally, by Proposition~\ref{prop:cohglue} alone; (4) records why the
plan's C-D1 may not extend Proposition~\ref{prop:homvanish} unrestricted: at weight $2$ and trivial character a level-$N$ group has rational cohomology in
two degrees. Nothing in (2) or (3) depends on it. (7)(a) is the route the D0 report recommends for D1 at $\ell\ge5$ in place of the shipped $C_7$ $K$-mode
(\file{docs/implementation/slot24/DESCENT0\_report.md} \S1), and (7)(b) the sparse cokernel route at level~$N$; both certify with the closed-form
$b_0$ and need no second prime.

\begin{proposition}[The bad primes at level $N$: a closed-form list, and the one degree of torsion at $\ell\ge5$]\label{prop:badprimesN}
\PROVED{} (v1.23.0-pre.14; refereed ACCEPT-with-fixes, \file{docs/implementation/slot24/C-D1\_referee.md}, the fixes applied at v1.23.0-pre.15, not
re-refereed.) For a prime $\ell\ge5$, $n\ge0$ and $t\in\Z/(\ell-1)$ put
\[
  \delta_\ell(n):=\#\{(a,b)\in\Z_{\ge0}^{2}:a(\ell+1)+b(\ell^{2}-\ell)=n\},\qquad
  \beta_\ell(n,t):=\#\{\,j\in\Z:0\le j\le n/\ell,\ n-2j+t\equiv0\pmod{\ell-1}\,\},
\]
and $\beta_\ell(n):=\beta_\ell(n,0)$; $\delta_\ell(n)$ is the count of Proposition~\ref{prop:badprimes}(1).
\begin{enumerate}
\item \textbf{(Every finite free $M$.)} Let $M$ be a $\Z[\Ga]$-module, finitely generated and free over $\Z$, and $\ell$ a prime. Then
  $H^{1}(\Ga;M)[\ell]\cong H^{0}(\Ga;M/\ell M)\big/\bigl(H^{0}(\Ga;M)/\ell\bigr)$, of dimension $s_\ell(M):=\dim_{\F_\ell}H^{0}(\Ga;M/\ell M)-b_0(M)$. If $\ell\ge5$,
  $H^{q}(\Ga;M)$ has no $\ell$-torsion for $q\ne1$, and the $\ell$-primary torsion of $H^{1}(\Ga;M)$ has exactly $s_\ell(M)$ cyclic summands. Put
  $\mathcal B(M):=\{\ell\ge5:s_\ell(M)>0\}$: the primes $\ge5$ at which $H^{\bullet}(\Ga;M)$ has torsion.
\item \textbf{(Level $N$, trivial character: the list.)} Let $N\ge1$, $n\ge0$, and $\ell\ge5$. If $n>0$:
  \[
    s_\ell(M_N)=\delta_\ell(n)\ \ (\ell\nmid N),\qquad s_\ell(M_N)=\beta_\ell(n)\ \ (\ell\mid N);
  \]
  if $n=0$, $s_\ell(M_N)=0$. Hence, for $n>0$ even, every prime at which $H^{\bullet}(\Ga_0(N);\Sym^{n}\Z^{2})$ has torsion lies in
  \[
    \mathcal B(n,N):=\{2,3\}\cup\{\,5\le\ell\le n-1:\ \ell\nmid N,\ \delta_\ell(n)>0\,\}\cup\{\,5\le\ell\le n+1:\ \ell\mid N,\ \beta_\ell(n)>0\,\};
  \]
  at each $\ell\ge5$ of $\mathcal B(n,N)$ the $\ell$-primary torsion lies in $H^{1}$ and has exactly $s_\ell(M_N)$ cyclic summands, so
  $\mathcal B(n,N)\cap[5,\infty)=\mathcal B(M_N)$; \textbf{every prime $\ell\ge5$ outside $\mathcal B(n,N)$ --- in particular every $\ell>n+1$ --- carries no torsion in
  any degree}, by this proof and not by omission (for $\ell>n$ by Claims~A and~B of the proof alone; at $5\le\ell\le n$, $\ell\nmid N$, by
  Claim~\ref{cl:dimH1}(2), proved in place since v1.23.0-pre.15). A prime $\ell=n+1\ge5$ lies in $\mathcal B(n,N)$ iff $\ell\mid N$, with $s_\ell=1$; e.g.\ $n=4$, $\ell=N=5$, the
  fixed vector $x^{4}$ (the audit's finding m3). For $n$ odd, $2\cdot H^{q}(\Ga_0(N);\Sym^{n}\Z^{2})=0$ for every $q$; for $n=0$ the torsion primes lie in $\{2,3\}$.
  At $N=1$ this is Proposition~\ref{prop:badprimes}.
\item \textbf{(Nebentypus.)} Let $\chi$, $m_\chi$, $\mathsf O_\chi$ and $M_N(\chi)$ be as in Proposition~\ref{prop:descentN}, $(n,\chi)\ne(0,1)$, and $\ell\ge5$ a
  prime with $\ell\nmid m_\chi$. Write $N=\ell^{k}N'$ with $\ell\nmid N'$, $\chi=\chi_\ell\chi'$ along $G_N=(\Z/\ell^{k})^{\times}\times(\Z/N')^{\times}$, and for a prime
  $\lambda\mid\ell$ of $\mathsf O_\chi$ let $\F_\lambda:=\mathsf O_\chi/\lambda\supseteq\F_\ell$ and $\bar\chi_\lambda:=\chi\bmod\lambda$. Then
  $s_\ell(M_N(\chi))=\sum_{\lambda\mid\ell}[\F_\lambda:\F_\ell]\,s_\lambda$, $s_\lambda:=\dim_{\F_\lambda}H^{0}(\Ga_0(N);V_{\F_\lambda}(\bar\chi_\lambda))$, and
  \begin{enumerate}
  \item[(a)] if $\ell\nmid N$: $s_\lambda=\delta_\ell(n)$ if $\chi=1$, and $s_\lambda=0$ if $\chi\ne1$;
  \item[(b)] if $\ell\mid N$: $s_\lambda=0$ if $\chi'\ne1$; if $\chi'=1$, then $\chi$ factors through $(\Z/\ell)^{\times}$, $\bar\chi_\lambda(a)=(a\bmod\ell)^{t_\lambda}$ for a
    unique $t_\lambda\in\Z/(\ell-1)$, and $s_\lambda=\beta_\ell(n,t_\lambda)$.
  \end{enumerate}
  So for $\chi\ne1$ the torsion at primes $\ell\ge5$, $\ell\nmid m_\chi$, lies in $H^{1}$ and occurs only at $\ell\mid N$ with $\chi'=1$ (the conductor of $\chi$ a power
  of $\ell$); such an $\ell$ need not be $\le n+1$: for $N=5$, $\chi=\bigl(\frac{\cdot}{5}\bigr)$ (even, $t=2$) and $n=2$, $s_5=\beta_5(2,2)=1$. The primes $\ell\mid6m_\chi$
  are not treated by this clause (at $\ell\mid m_\chi$, $\mathsf O_\chi$ is ramified over $\ell$); they are bad primes until a later clause says otherwise.
\end{enumerate}
\end{proposition}

\begin{proof}
(1) $M$ is free over $\Z$, so $0\to M\xrightarrow{\ell}M\to M/\ell M\to0$ is exact, and its long exact sequence
$0\to H^{0}(M)\xrightarrow{\ell}H^{0}(M)\to H^{0}(M/\ell M)\to H^{1}(M)\xrightarrow{\ell}H^{1}(M)$ gives the displayed isomorphism. $H^{0}(\Ga;M)=M^{\Ga}$ is a submodule
of the finitely generated free $\Z$-module $M$, hence free, of rank $b_0(M)$ (flatness of $\Q$), and $H^{0}(M)/\ell\to H^{0}(M/\ell M)$ is injective by the sequence;
so $\dim H^{1}[\ell]=s_\ell(M)$. The number of cyclic $\ell$-primary summands of a finitely generated abelian group $G$ is $\dim_{\F_\ell}G[\ell]$. For $\ell\ge5$:
$H^{0}$ is torsion-free, and $24\cdot H^{q}(\Ga;M)=0$ for $q\ge2$ (Proposition~\ref{prop:certifiedeell}(4)(a), proved there for every $\Z[\Ga]$-module).

(2) $M_N/\ell M_N=\Ind_{\Ga_0(N)}^{\Ga}\Mn$ (the coordinates \eqref{eq:MH} commute with reduction), so $H^{0}(\Ga;M_N/\ell)=H^{0}(\Ga_0(N);\Mn)$ by Frobenius
reciprocity in degree $0$; $b_0(M_N)=0$ for $n>0$ and $1$ for $n=0$ (Proposition~\ref{prop:certifiedeell}(4)(b)). For $n=0$, $H^{0}(\Ga_0(N);\F_\ell)=\F_\ell$ and
$s_\ell=0$. Let $n>0$. $\rhon$ has entries that are integer polynomials in those of $\gamma$ (\eqref{eq:rhondef}), so $\Ga_0(N)$ acts on $\Mn$ through its image
$\bar\Ga$ in $\SL_2(\F_\ell)$, and $H^{0}(\Ga_0(N);\Mn)=\Mn^{\bar\Ga}$. \emph{The image} (with~(3)'s refinement). $\Ga_0(N)$ is the preimage of the upper-triangular
group $B(\Z/N)$ under $\SL_2(\Z)\to\SL_2(\Z/N)$, and $\SL_2(\Z)\to\SL_2(\Z/M')$ is onto for every $M'\ge1$ (as in the proof of Lemma~\ref{lem:nebdata}). If
$\ell\nmid N$, $\SL_2(\Z/N\ell)=\SL_2(\Z/N)\times\SL_2(\F_\ell)$ (Chinese remainder theorem), so $\Ga_0(N)$, the preimage of $B(\Z/N)\times\SL_2(\F_\ell)$, maps onto
$B(\Z/N)\times\SL_2(\F_\ell)$ and $\bar\Ga=\SL_2(\F_\ell)$. If $\ell\mid N$, $\Ga_0(N)\to B(\Z/N)=B(\Z/\ell^{k})\times B(\Z/N')$ is onto, and so is
$B(\Z/\ell^{k})\to B(\F_\ell)$ (lift the diagonal unit $a$ to a unit $\tilde a$ and take $\smat{\tilde a}{\tilde b}0{\tilde a^{-1}}$); so $\bar\Ga=B(\F_\ell)$.
\emph{Case $\ell\nmid N$.} $\Mn^{\SL_2(\F_\ell)}=H^{0}(\Ga;\Mn)$, as $\Ga\to\SL_2(\F_\ell)$ is onto as well, and its dimension is $\delta_\ell(n)$ by the proof of
Proposition~\ref{prop:badprimes}(1) (Claim~\ref{cl:dimH1}(2), proved in place since v1.23.0-pre.15 --- the same input as at level~$1$).
For $\ell>n$ the count needs no Claim~\ref{cl:dimH1}(2): see \emph{the list} below.
\emph{Case $\ell\mid N$.} By Definition~\ref{def:rhon}, $\Mn$ is the degree-$n$ part of $\F_\ell[x,y]$, $\gamma$ acting by the algebra automorphism with
$\gamma\cdot x=ax+cy$, $\gamma\cdot y=bx+dy$. $B(\F_\ell)$ is generated by $u_b:=\smat1b01$ and $h_a:=\mathrm{diag}(a,a^{-1})$; $u_b$ fixes $x$ and sends $y\mapsto y+bx$,
and $h_a$ sends $x\mapsto ax$, $y\mapsto a^{-1}y$. Put $P:=\prod_{c\in\F_\ell}(y+cx)=y^{\ell}-x^{\ell-1}y$. \emph{Claim A: the $u_b$-invariants ($b\in\F_\ell$) of
degree $n$ are spanned by the $x^{i}P^{j}$, $i+\ell j=n$, and these are linearly independent.} $u_b$ permutes the factors of $P$ ($c\mapsto c+b$), so $x$ and $P$
are invariant. Independence: $x^{i}P^{j}$ has $y$-degree $\ell j$. Spanning, by induction on $n$: let $f$ be invariant of degree $n$; divide by $P$, monic of
degree $\ell$ in $y$ over $\F_\ell[x]$: $f=gP+r$ with $\deg_yr<\ell$, uniquely, $g$ and $r$ homogeneous. $u_b$ preserves $y$-degrees and fixes $f$ and $P$, so by
uniqueness $u_bg=g$ and $u_br=r$; $g\in\F_\ell[x,P]$ by induction. Let $e:=\deg_yr$ and suppose $e\ge1$, $r=\sum_{k\le e}r_k(x)y^{k}$ with $r_e\ne0$. The
polynomial $R(t):=r(x,y+tx)-r(x,y)\in\F_\ell[x,y][t]$ has degree $\le e<\ell$ in $t$, its $t^{e}$-coefficient $r_e(x)x^{e}\ne0$, and the $\ell$ distinct roots
$t\in\F_\ell$ --- impossible over the domain $\F_\ell[x,y]$. So $e=0$ and $r\in\F_\ell x^{n}$. \emph{Claim B:} $h_a(x^{i}P^{j})=a^{i-j}x^{i}P^{j}$: $h_aP=
\prod_c(a^{-1}y+cax)=a^{-\ell}\prod_c(y+ca^{2}x)=a^{-\ell}P$ ($c\mapsto ca^{2}$ is a bijection of $\F_\ell$) and $a^{\ell}=a$. So the torus $\{h_a\}$ acts diagonally on the basis of
Claim~A, the $B(\F_\ell)$-invariants are spanned by the $x^{i}P^{j}$ with $a^{i-j}=1$ for all $a\in\F_\ell^{\times}$, i.e.\ $(\ell-1)\mid i-j$, and with $i=n-\ell j$,
$i-j=n-(\ell+1)j\equiv n-2j\pmod{\ell-1}$: the dimension is $\beta_\ell(n)$. \emph{The list.} $\delta_\ell(n)>0$, $n>0$, needs $n\ge\ell+1$ (the least positive value of
$a(\ell+1)+b(\ell^{2}-\ell)$, as $\ell^{2}-\ell\ge\ell+1$); $\beta_\ell(n)>0$ needs a $j$ with $(\ell-1)\mid n-2j$: $j=0$ forces $\ell-1\le n$, $j\ge1$ forces
$\ell\le n$. At $\ell=n+1$, $j\le n/\ell<1$ leaves $j=0$, and $(\ell-1)=n$ divides $n$: $\beta_\ell(n)=1$; $\delta_{n+1}(n)=0$. \emph{The good primes without
Claim~\ref{cl:dimH1}(2)} (v1.23.0-pre.15, referee F1(a)): for $\ell>n>0$ and $\ell\nmid N$, Claim~A leaves only $x^{n}$ among the $u_b$-invariants of degree $n$,
since $\deg_yP=\ell>n$; the lower unipotent $\smat1011\in\SL_2(\Fl)=\bar\Ga$ sends $x\mapsto x+y$, and $(x+y)^{n}\ne x^{n}$ (its coefficient of $x^{n-1}y$ is
$n\not\equiv0$), so $x^{n}$ is not invariant and $s_\ell(M_N)=0$. So the rule ``every $\ell>n+1$ is good'' rests on Claims~A and~B alone. The rest of the list
statement is (1). \emph{$n$ odd:} $M_N(-I)=(-1)^{n}=-1$ (the proof of Proposition~\ref{prop:levelorbit}(2)), $-I$ is central and acts trivially on cohomology, so $2$ kills
$H^{\bullet}$ (Proposition~\ref{prop:badprimes}(2)'s argument). \emph{$n=0$:} $s_\ell=0$ for $\ell\ge5$, so by (1) there is no $\ell$-torsion for $\ell\ge5$. At $N=1$,
$\ell\nmid N$ always and the list is Proposition~\ref{prop:badprimes}(1)'s.

(3) $b_0(M_N(\chi))=0$ by Proposition~\ref{prop:descentN}(4). $M_N(\chi)/\ell=\Ind V_{\mathsf O_\chi/\ell}(\chi\bmod\ell)$ and $\mathsf O_\chi/\ell=\prod_\lambda\F_\lambda$ (the
proof of Proposition~\ref{prop:descentN}(8) at $K=1$), so $H^{0}(\Ga;M_N(\chi)/\ell)=\bigoplus_\lambda H^{0}(\Ga_0(N);V_{\F_\lambda}(\bar\chi_\lambda))$, of
$\F_\ell$-dimension $\sum_\lambda[\F_\lambda:\F_\ell]s_\lambda$; then apply (1). $\gamma\in\Ga_0(N)$ acts on $V_{\F_\lambda}(\bar\chi_\lambda)$ by
$\bar\chi_\lambda(a_\gamma)\,\rhon(\gamma\bmod\ell)$. $\bar\chi_\lambda$ has order prime to $\ell$ (its values lie in $\F_\lambda^{\times}$), so its $(\Z/\ell^{k})^{\times}$-component
is trivial on the $\ell$-group $\ker\bigl((\Z/\ell^{k})^{\times}\to\F_\ell^{\times}\bigr)$, $k\ge1$. Moreover $\mu_{m_\chi}\to\F_\lambda^{\times}$ is injective ($Y^{m_\chi}-1$ is
separable mod $\ell$), so $\bar\chi_\lambda=1$ iff $\chi=1$, and $\bar\chi'_\lambda=1$ iff $\chi'=1$. (a) $\ell\nmid N$: the action factors through
$\phi\colon\Ga_0(N)\to\SL_2(\F_\ell)\times G_N$, $\gamma\mapsto(\gamma\bmod\ell,a_\gamma\bmod N)$, which is onto by the computation in (2) ($\Ga_0(N)\to
B(\Z/N)\times\SL_2(\F_\ell)$ is onto, and so is $B(\Z/N)\to G_N$, $\beta\mapsto a_\beta$). The fixed space of $\SL_2(\F_\ell)\times G_N$ on $\Sym^{n}\F_\lambda^{2}\otimes
\F_\lambda(\bar\chi_\lambda)$ is $(\Sym^{n}\F_\lambda^{2})^{\SL_2(\F_\ell)}$ if $\bar\chi_\lambda=1$ and $0$ otherwise; the former has dimension $\delta_\ell(n)$ (a fixed
space of matrices over $\F_\ell$ commutes with the extension $\F_\ell\subseteq\F_\lambda$). (b) $\ell\mid N$: the action factors through
$\Ga_0(N)\to B(\F_\ell)\times(\Z/N')^{\times}$, $\gamma\mapsto(\gamma\bmod\ell,a_\gamma\bmod N')$, onto by (2) ($B(\Z/N')\to(\Z/N')^{\times}$ onto). The second factor acts
by $\bar\chi'_\lambda$, so the fixed space is $0$ unless $\chi'=1$. If $\chi'=1$, $\chi=\chi_\ell$ has order prime to $\ell$, so it factors through
$(\Z/\ell)^{\times}=\F_\ell^{\times}$; $\bar\chi_\lambda$ maps the cyclic group $\F_\ell^{\times}$ into $\mu_{\ell-1}(\F_\lambda)=\F_\ell^{\times}$, hence is $a\mapsto a^{t_\lambda}$ for a
unique $t_\lambda\bmod\ell-1$. Claim~A holds over $\F_\lambda$ verbatim ($P$ and the $u_b$, $b\in\F_\ell$, are unchanged; $\F_\lambda[x,y]$ is a domain), and
$u_b$ has $a_{u_b}=1$, so the twist does not change the $u_b$-invariants; $h_a$ has $a_{h_a}=a$ and acts on $x^{i}P^{j}$ by $a^{t_\lambda}a^{i-j}$. So the fixed space
has dimension $\#\{j:(\ell-1)\mid n-2j+t_\lambda\}=\beta_\ell(n,t_\lambda)$. \emph{The example.} $\chi=\bigl(\frac{\cdot}{5}\bigr)$ has $\chi(-1)=1$, $m_\chi=2$, $\mathsf O_\chi=\Z$,
$\bar\chi(a)=a^{2}$ in $\F_5$ (Euler's criterion), so $t=2$; $\beta_5(2,2)=\#\{j=0:4\mid4\}=1$. The parity is built in: if $\chi(-1)\ne(-1)^{n}$, $n+t_\lambda$ is odd
($\bar\chi_\lambda(-1)=(-1)^{t_\lambda}$) and $\beta_\ell(n,t_\lambda)=0$, as Proposition~\ref{prop:neborbit}(2) requires.
\end{proof}

\noindent
\emph{Reading.} At $\ell\ge5$ the torsion of any finite free coefficient module sits in $H^{1}$ and is counted by $H^{0}$ modulo $\ell$, so the bad-prime
list is the list of primes at which the reduction has more invariants than the lattice; at level $N$ the reduction of $\Ga_0(N)$ is $\SL_2(\F_\ell)$ for
$\ell\nmid N$ (the level-$1$ Dickson count) and the Borel subgroup for $\ell\mid N$ (the count $\beta_\ell$, which reaches $\ell=n+1$). The descent's
$\Z[1/6]$-piece at level~$N$ therefore needs certified $\Z/\ell^{K}$ runs at the primes of $\mathcal B(n,N)\cap[5,\infty)$ only, in $H^{1}$ only
(Proposition~\ref{prop:descentN}(7)); every other prime $\ge5$ is torsion-free by Proposition~\ref{prop:badprimesN}, not by omission.

\begin{proposition}[The $3$-adic shortcut at level $N$: no element of order $3$]\label{prop:descentNthree}
\PROVED{} (v1.23.0-pre.17; \textbf{not refereed}; the D1 code audit's finding M1, \file{docs/implementation/slot24/D1\_code\_audit.md}.) Assume \ref{H4} and
\ref{B1}--\ref{B3}. Let $N\ge1$ with $\nu_3(N)=0$, $n\ge0$, $M:=M_N=\Ind_{\Ga_0(N)}^{\Ga}\Vcoef$, $U:=ST$, and $\mathrm T(M_S)$ the tree complex
\eqref{eq:treemodel} with fibre $M_S$. By Proposition~\ref{prop:modeltorsion} and the local counts in the proof of Proposition~\ref{prop:descentN}(4),
$\nu_3(N)=0$ \textbf{iff} $\Ga_0(N)$ contains no element of order~$3$ \textbf{iff} $9\mid N$ or $N$ has a prime factor $p\equiv2\pmod3$ ($p=2$ included).
\begin{enumerate}
\item \textbf{(The stabilisers.)} Every finite subgroup of $\Ga_0(N)$ has order $2$ or $4$ (or $1$); for every $\gamma\in\Ga$,
  $\Ga_0(N)\cap\gamma\langle U\rangle\gamma^{-1}=\{\pm I\}$, and $\Ga_0(N)\cap\gamma\langle S\rangle\gamma^{-1}$ is $\{\pm I\}$ or $\gamma\langle S\rangle\gamma^{-1}$. So
  \textbf{$3$ is prime to the order of every stabiliser of $\Ga_0(N)$} on $\mathcal T$ and on $\Xbulk$. Equivalently, $U$ has no fixed point on
  $\mathbb P^1(\Z/N)$ (the $\langle U\rangle$-stabiliser of a point contains $U^{3}=-I$, so it is $\{\pm I\}$ or $\langle U\rangle$; v1.23.0-pre.18, referee m1), every $\langle U\rangle$-orbit is free modulo $\{\pm I\}$ (of length $3$), and $M$ restricted to $\langle U^{2}\rangle\cong\Z/3$ is induced from
  the trivial subgroup, so $H^{t}(\langle U^{2}\rangle;M_S)=0$ for every $t\ge1$ and every ring $S$.
\item \textbf{(What the tree complex computes $3$-locally.)} For every $\Z_{(3)}$-algebra $S$: $\mathrm T(M_{\Z_{(3)}})$ is a complex of finitely generated free
  $\Z_{(3)}$-modules in degrees $0,1$, $\mathrm T(M_{\Z_{(3)}})\otimes_{\Z_{(3)}}S=\mathrm T(M_S)$, and
  \[
    H^{q}(\mathrm T(M_S))\;\cong\;H^{q}(\Ceq(M_S))\;\cong\;H^{q}(\Ga;M_S)\;=\;H^{q}(\Ga_0(N);\Sym^{n}S^{2})\qquad\text{for every }q ;
  \]
  moreover $\cmpI^{\ast}\otimes\Z_{(3)}\colon\Ceq(M_{\Z_{(3)}})\to\CI(M)\otimes\Z_{(3)}$ is a quasi-isomorphism. Hence
  \textbf{$H^{q}(\Ga_0(N);\Vcoef)$ has no $3$-torsion for $q\ge2$}, and $\mathrm{tors}\,H^{1}(\Ga_0(N);\Vcoef)\otimes\Z_{(3)}=\mathrm{tors}\coker(d^{0})$ for
  $d^{0}$ of $\mathrm T(M_{\Z_{(3)}})$. This is Proposition~\ref{prop:descentN}(1)'s statement for $\mathrm T$ (for $\Ceq$: its cohomology) with $\mathsf O$ replaced by $\Z_{(3)}$, for
  this $M$; it is false for $M_N$ when $\nu_3(N)>0$ (item~(5)).
\item \textbf{(The certified exponent at $\ell=3$.)} Let $n$ be even and $K\ge1$. Proposition~\ref{prop:descentN}(7)(b) holds at $\ell=3$ verbatim: let $\bar R$ be any matrix
  over $\Z/3^{K}$, written in a basis of $M_{\Z/3^{K}}=(M_{\Z/3^{K}})^{\langle-I\rangle}$, whose rows generate
  $\mathsf R:=(M_{\Z/3^{K}})^{\langle S\rangle}+(M_{\Z/3^{K}})^{\langle U\rangle}$, and $s$ the number of non-zero diagonal entries $3^{\sigma_i}$ of its Smith form.
  Then $s=r_K(d^{0})\le r:=\rank\mathrm T^{0}-b_0=m_N(n+1)-b_1$, with the closed forms
  $\rank\mathrm T^{0}=\tfrac{m_N-\nu_2}{2}(n+1)+\nu_2\,d_S+\tfrac{m_N}{3}(n+1)$ (Proposition~\ref{prop:levelorbit}(1) at $\nu_3=0$), $b_0=[n=0]$ and
  $b_1=b_0+{}$the right-hand side of Proposition~\ref{prop:levelorbit}(3) over $\Q$ --- never the number of rows of $\bar R$. If $s=r$, then
  $\mathrm{tors}\,H^{1}(\Ga_0(N);\Vcoef)\otimes\Z_{(3)}\cong\bigoplus_{\sigma_i>0}\Z/3^{\sigma_i}$ exactly; if $s<r$, $K$ is too small, and the stopping rule of
  Proposition~\ref{prop:certifiedeell}(6) applies (raise $K$; at the cap of the ring, refuse). No a-priori bound on the exponents of $d^{0}$ is asserted.
  \emph{The generators and their certificate.} (a) On a free orbit --- every $\langle U\rangle$-orbit, and the free $\langle S\rangle$-orbits --- the orbit-wise
  generators $\sum_{i<o}M(g)^{i}e_{(j_0,a)}$, $a=0,\dots,n$ ($g=U$, $o=3$; $g=S$, $o=2$), are a basis of the orbit's fixed module over every ring
  (Proposition~\ref{prop:levelorbit}(1): no solve, no division). (b) At each of the $\nu_2(N)$ fixed points $j$ of $S$, $F_j=\langle S\rangle$ has order $4$, prime
  to $3$: the columns of the averaging idempotent $e_{F_j}=\tfrac12\bigl(1+\rhon(h_j(S))\bigr)$ ($n$ even) need no certificate, and any other family generates
  $(V_{\Z/3^{K}})^{F_j}$ iff (i) each vector is fixed mod $3^{K}$ by a generator of $F_j$ and (ii) the reductions mod $3$ span $(V_{\F_3})^{F_j}$, of dimension
  $d_S$ --- (7)(b)'s certificate (referee F4), whose only use of $\ell\ge5$ is $\ell\nmid\abs{F_j}$. (c) $\langle U\rangle$ has \textbf{no} fixed points, so the case
  $\abs{F_j}=6$, where $3\mid\abs{F_j}$ and (7)(b)'s certificate would fail, does not occur. So the rows built by (a) and (b) are admissible with no certificate.
\item \textbf{($n$ odd.)} $2$ kills $H^{\bullet}(\Ga_0(N);\Vcoef)$ (Proposition~\ref{prop:badprimesN}(2)), so its $3$-primary part is $0$ in every degree, at every
  $N$; nothing is computed.
\item \textbf{(Scope: the hypothesis is needed, and nothing else is.)} $\nu_2(N)$ plays no role (the stabilisers $\{\pm I\}$, $\langle S\rangle$ have orders
  $2$, $4$), and $9\mid N$ is a case of the proposition ($\nu_3(9m)=0$), not an exception. If $\nu_3(N)>0$, $\Ga_0(N)$ contains a conjugate of $\langle U\rangle$
  (Proposition~\ref{prop:modeltorsion}), $H^{q}(\Ga_0(N);\Vcoef)$ may carry $3$-torsion in degrees $q\ge2$, which $\mathrm T$ and $\Ceq$ cannot see
  ($H^{q}=0$ for $q\ge2$ on both; Proposition~\ref{prop:ellthreeclosed}(3) at $N=1$), and the $3$-primary part needs the $\ell$-local model $\Iso^{(3)}$
  (Proposition~\ref{prop:certifiedeell} at $\ell=3$). Instances (the independent oracle \file{docs/implementation/slot24/descentN\_oracle.tsv}):
  $H^{2}(\SL_2(\Z);\Z)\supseteq\Z/3$ ($N=1$, $n=0$), $H^{2}(\Ga_0(7);\Z)\supseteq(\Z/3)^{2}$, $H^{3}(\Ga_0(3);\Sym^{4}\Z^{2})\supseteq\Z/3$.
\end{enumerate}
In the descent (Proposition~\ref{prop:descentN}(2)) the $3$-primary part of $H^{q}(\Ga;M)\otimes\Zsix$ may therefore be read on $\mathrm T(M_{\Z_{(3)}})$ in place of
$\Iso^{(3)}$ when $\nu_3(N)=0$: by~(2) the identification is induced by $\cmpI^{\ast}$, the descent's only glueing datum, and the fibre maps $u$ of
Proposition~\ref{prop:descentN}(3) (restriction, trace, $R_\sigma$) commute strictly with it, so their $3$-local components may be read on $\Ceq(M_{\Z_{(3)}})$
as well. No cochain-level Hecke operator on $\mathrm T$ at level~$N$ is asserted here.
\end{proposition}

\begin{proof}
(1) The finite subgroups of $\Ga$ are cyclic of order $1,2,3,4,6$, each contained in a conjugate of $\langle S\rangle\cong\Z/4$ or of $\langle U\rangle\cong\Z/6$, and
every one of even order contains $-I$ (Proposition~\ref{prop:modeltorsion} and its proof). A subgroup of order $3$ or $6$ contains an element of order $3$, which
$\Ga_0(N)$ lacks: by Proposition~\ref{prop:modeltorsion}, $\Ga_0(N)$ has one iff $a^{2}+a+1\equiv0\pmod N$ is solvable, and $a\mapsto-a$ maps the roots of
$X^{2}+X+1$ bijectively onto those of $\chi_{ST}=X^{2}-X+1$, whose number modulo $N$ is $\nu_3(N)$ (the proof of Proposition~\ref{prop:descentN}(4)); that proof's
local counts give $\nu_3(p^{e})=0$ exactly for $p\equiv2\pmod3$ ($p=2$ included) and for $p^{e}=3^{e}$, $e\ge2$, and $\nu_3$ is multiplicative, whence the
stated equivalence. As $-I\in\Ga_0(N)$ is central, $\Ga_0(N)\cap\gamma\langle U\rangle\gamma^{-1}$ is a subgroup of the cyclic group $\gamma\langle U\rangle\gamma^{-1}$ of
order $6$ containing $-I$ and no element of order $3$, i.e.\ $\{\pm I\}$; likewise $\Ga_0(N)\cap\gamma\langle S\rangle\gamma^{-1}\supseteq\{\pm I\}$ has order $2$ or $4$.
The stabilisers of $\Ga_0(N)$ on cells of $\mathcal T$ (conjugates of $\langle-I\rangle$, $\langle S\rangle$, $\langle U\rangle$ intersected with $\Ga_0(N)$) and of
$\Xbulk$ (finite subgroups of $\Ga$, Proposition~\ref{prop:bredon}(1), intersected with $\Ga_0(N)$) are finite subgroups of $\Ga_0(N)$. The stabiliser in
$\langle U\rangle$ of the coset $\Ga_0(N)w$ is $\langle U\rangle\cap w^{-1}\Ga_0(N)w=w^{-1}\bigl(\Ga_0(N)\cap w\langle U\rangle w^{-1}\bigr)w=\{\pm I\}$, so
$\langle U\rangle/\{\pm I\}\cong\langle U^{2}\rangle$ acts freely on $\Ga_0(N)\backslash\Ga=\mathbb P^1(\Z/N)$ (the block indices $j\mapsto j\cdot\gamma$ of \eqref{eq:MH}, Definition~\ref{def:levelNmodel}), with orbits of
length $3$. $M=\bigoplus_j V_j$ is the sum of its blocks (one per point of $\mathbb P^1(\Z/N)$), permuted by $\Ga$ as the points are; on an orbit $O=\{j,U^{2}j,U^{4}j\}$
the block sum $M_O$ is $\bigoplus_{i<3}U^{2i}V_j$, so $M_O\otimes S\cong S[\langle U^{2}\rangle]\otimes_S(V_j\otimes S)$ is induced from the trivial subgroup. For a cyclic
group $C=\langle c\rangle$ of order $3$ and such an $X=\bigoplus_{i<3}c^{i}X_0$, the periodic resolution gives $H^{2k}(C;X)=X^{C}/\mathcal NX$ and
$H^{2k+1}(C;X)=\ker\mathcal N/(c-1)X$ ($k\ge1$, resp.\ $k\ge0$; $\mathcal N=1+c+c^{2}$); $x\in X^{C}$ is $x=\sum_ic^{i}x_0=\mathcal Nx_0$ with $x_0$ its $X_0$-component, and
$y=\sum_ic^{i}y_i\in\ker\mathcal N$ (each $y_i\in X_0$) has $y_0+y_1+y_2=0$ (the $X_0$-component of $\mathcal Ny$), so $y=(c-1)\bigl(cy_2-y_0\bigr)$ --- indeed
$(c-1)(cy_2-y_0)=c^{2}y_2-cy_2-cy_0+y_0=c^{2}y_2+cy_1+y_0$. So $H^{t}(C;X)=0$ for $t\ge1$, over every ring.

(2) Let $A_0:=\Ga_0(N)$ and $S$ a $\Z_{(3)}$-algebra; $2$ and $4$ are units in $S$. \emph{The tree.} The proof of Lemma~\ref{lem:treeres} uses of
$6\in\cring^{\times}$ only that the orders of the groups $A_0\cap\gamma F\gamma^{-1}$ are units (as noted in the proof of Proposition~\ref{prop:modeltorsion}); by (1)
they are $2$ or $4$, so $C_{\bullet}(\mathcal T)\otimes S\to S$ is a projective resolution over $S[A_0]$ and $H^{q}(A_0;V_S)=H^{q}(\Hom_{A_0}(C_{\bullet}(\mathcal
T),V_S))$, $0$ for $q\ge2$. Frobenius reciprocity (as in the proof of Proposition~\ref{prop:modeltorsion}) identifies $\Hom_{A_0}(C_{\bullet}(\mathcal T),V_S)$ with
$\Hom_{\Z[\Ga]}(C_{\bullet}(\mathcal T),\Ind_{A_0}^{\Ga}V_S)=\mathrm T(M_S)$ (the derivation of \eqref{eq:treemodel}; $\Ind_{A_0}^{\Ga}V_S=M_S$, induction commuting
with $\otimes S$), and Shapiro's lemma \STD{\cite[Ch.~III.6.2]{brown}} gives $H^{q}(A_0;V_S)=H^{q}(\Ga;M_S)$. \emph{The Bredon complex.} Proposition~\ref{prop:modeltorsion}
with $B=A_0$, $X=V_S$, $\cring=S$ (its hypothesis is (1)) gives $H^{q}(\Ceq(M_S))\cong H^{q}(A_0;V_S)$. \emph{The comparison.} $C_s(\Xbulk)\otimes S=\bigoplus_\sigma
S[\Ga/\Ga_\sigma]$ restricts over $S[A_0]$ to a sum of $S[A_0/H']$, $H'=A_0\cap\gamma\Ga_\sigma\gamma^{-1}$ over the double cosets, of order a unit by (1), hence
projective (the idempotent of Lemma~\ref{lem:treeres}'s proof); $C_{\bullet}(\Xbulk)\otimes S\to S$ is exact (\ref{H4}(ii), split over $\Z$). So $C_{\bullet}(\Xbulk)\otimes
S$ and $\Iso_{\bullet}\otimes S$ (free over $S[\Ga]$, hence over $S[A_0]$) are projective resolutions of $S$ over $S[A_0]$, and $\cmpI\otimes S$ is a chain map between
them over $\id_S$, hence a homotopy equivalence over $S[A_0]$ (the comparison theorem, \STD{\cite[Thm.~2.2.6]{weibel}}). $\Hom_{A_0}(-,V_S)$ preserves homotopy
equivalences, and Frobenius reciprocity is natural in the first variable, so $\cmpI^{\ast}\colon\Hom_{\Z[\Ga]}(C_{\bullet}(\Xbulk),M_S)\to\Hom_{\Z[\Ga]}(\Iso_{\bullet},M_S)$
is a quasi-isomorphism; at $S=\Z_{(3)}$ the target is $\CI(M)\otimes\Z_{(3)}$ (Proposition~\ref{prop:descentN}(1)), the source $\Ceq(M_{\Z_{(3)}})$, and the
cohomology is $H^{q}(\Ga;M)\otimes\Z_{(3)}$ ($\Z_{(3)}$ flat). \emph{Freeness and base change of $\mathrm T$.} $\mathrm T^{1}=M_S^{\langle-I\rangle}=e_{-I}M_S$ and
$M_S^{\langle S\rangle}=e_{\langle S\rangle}M_S$ for the averaging idempotents ($2$, $4$ units): direct summands of finitely generated free modules, free over the
principal ideal domain $\Z_{(3)}$, and $e(M_{\Z_{(3)}})\otimes S=eM_S$ (the image of an idempotent commutes with base change). $\langle U\rangle$ has order $6$,
not a unit, but by (1) every $\langle U\rangle$-orbit is free modulo $\{\pm I\}=\langle U\rangle\cap\{\pm I\}$, so by Proposition~\ref{prop:levelorbit}(1), over every
ring, restriction to the blocks $j_O$ is an isomorphism $M_S^{\langle U\rangle}\to\bigoplus_OV_S^{\{\pm I\}}=\bigoplus_Oe_{-I}V_S$, with the inverse given by the
transported formula with coefficients in $\Z$; so $M_{\Z_{(3)}}^{\langle U\rangle}$ is free and $M_{\Z_{(3)}}^{\langle U\rangle}\otimes S\to M_S^{\langle U\rangle}$ is an
isomorphism. The differential $(a,b)\mapsto b-a$ is base-changed with them. \emph{The consequences.} $H^{q}(\Ga_0(N);\Vcoef)\otimes\Z_{(3)}=0$ for $q\ge2$, a
finitely generated group, so it has no $3$-torsion; $H^{1}(\mathrm T(M_{\Z_{(3)}}))=\coker d^{0}$.
\emph{Against Proposition~\ref{prop:ellthreeclosed}.} With $\cring=S$ there, (2) of that proposition gives $H^{q}(\Ga;M_S)\cong H^{q}(\langle U^{2}\rangle;M_S^{-I})$ for
$q\ge2$, and (3) gives that the Bredon route is exact in degree $0$ and injective in degree~$1$ with cokernel $H^{1}(\langle U^{2}\rangle;M_S^{-I})$. $M_S^{-I}=e_{-I}M_S$
is a $\langle U^{2}\rangle$-summand of $M_S$ ($-I$ is central), so its $H^{t}$, $t\ge1$, vanish by (1): the two statements agree. At $N=1$ ($\nu_3(1)=1$) neither term vanishes in general: $H^{2}(\SL_2(\Z);\Z)\supseteq\Z/3$.

(3) The proof of Proposition~\ref{prop:descentN}(7)(b), with (2) in place of~(1) there: $\mathrm T(M_{\Z_{(3)}})\otimes\Z/3^{K}=\mathrm T(M_{\Z/3^{K}})$, so the image of
$d^{0}\bmod3^{K}$ is $\mathsf R$ and $\coker\bar R=\coker(d^{0})\otimes\Z/3^{K}$; over $\Z_{(3)}$, $\coker d^{0}\cong\Z_{(3)}^{D-r}\oplus\bigoplus_t\Z/3^{a_t}$ ($D=m_N(n+1)$,
the $a_t\ge1$), and comparing multiplicities of the cyclic summands $\Z/3^{c}$ over the local principal ideal ring $\Z/3^{K}$ gives $s=r-\#\{a_t\ge K\}=r_K(d^{0})$ and
$\{\sigma_i>0\}=\{a_t<K\}$. By (2), $\mathrm{tors}\coker d^{0}=\mathrm{tors}\,H^{1}(\Ga;M)\otimes\Z_{(3)}$, and $H^{1}(\Ga;M)=H^{1}(\Ga_0(N);\Vcoef)$
(Proposition~\ref{prop:descentN}(2)). $r=\rank\mathrm T^{0}-b_0$ is rank--nullity for $\mathrm T(M_{\Q})$, whose $H^{0}$ is $M_{\Q}^{\langle S\rangle}\cap
M_{\Q}^{\langle U\rangle}=M_{\Q}^{\Ga}$, of dimension $b_0$ (Proposition~\ref{prop:descentN}(4)); $\rank\mathrm T^{0}=\dim M_{\Q}^{\langle S\rangle}+\dim
M_{\Q}^{\langle U\rangle}$ by (2), and Proposition~\ref{prop:levelorbit}(1) over $\Q$ gives the free orbits' $(n+1)$ each and the fixed points' $d_S$; $r=m_N(n+1)-b_1$ is
rank--nullity in degree $1$ ($\mathrm T^{1}=M$ for $n$ even, $M(-I)=1$). \emph{The generators.} The rows generate $\mathsf R$ iff they generate each orbit's fixed module.
(a) On an orbit free modulo $\{\pm I\}$ the restriction to block $j_0$ of $\sum_{i<o}M(g)^{i}e_{(j_0,a)}$ is $e_a$ (the other terms lie in other blocks), so by
Proposition~\ref{prop:levelorbit}(1) they map to a basis of $V^{\{\pm I\}}=V$ ($n$ even) under an isomorphism defined over every ring. (b) The proof of
(7)(b)'s certificate, with $\ell=3\nmid4=\abs{F_j}$: $W:=(V_{\Z/3^{K}})^{F_j}=e_{F_j}V_{\Z/3^{K}}$ is free with $W/3W=(V_{\F_3})^{F_j}$, the idempotent's columns
generate $W$, and Nakayama's lemma over $\Z/3^{K}$ gives the iff. $e_{F_j}=\tfrac14\sum_{h\in F_j}\rhon(h)=\tfrac12(1+\rhon(h_j(S)))$ since $\rhon(-I)=1$ for $n$ even.
(c) is (1). The stopping rule and the absence of an a-priori bound at $d^{0}$ are Proposition~\ref{prop:certifiedeell}(6) (its bound $a_t\le v_\ell(24)$ is for $q\ge1$ only).

(4) Proposition~\ref{prop:badprimesN}(2)'s proof ($M_N(-I)=-1$, $-I$ central acting trivially on cohomology), for every $N$.

(5) The orders are (1). If $\nu_3(N)>0$, $\Ga_0(N)$ contains an element $\gamma$ of order $3$, so $\langle-I,\gamma\rangle$, a cyclic group of order $6$, is a
conjugate of $\langle U\rangle$ (the classification in (1)). $H^{q}=0$ for $q\ge2$ on $\mathrm T$ and $\Ceq$: $\mathrm T$ has length $1$, and
$H^{2}(\Ceq(M))=0$ (Proposition~\ref{prop:descentN}(1)) with $\Ceq$ in degrees $\le2$. The instances are rows of the oracle table: $(N,n,q)=(1,0,2)$:
$\Z/2^{2}+\Z/3$; $(7,0,2)$: $\Z/2+\Z/3+\Z/3$; $(3,4,3)$: $\Z/2+\Z/2+\Z/2+\Z/2+\Z/2+\Z/3$; $\nu_3(1)=\nu_3(3)=1$,
$\nu_3(7)=2$. The closing paragraph: Proposition~\ref{prop:descentN}(2)'s $H^{q}(B)=H^{q}(\Ga;M)\otimes\Zsix$, a finitely generated $\Zsix$-module, is the sum of its
$2$- and $3$-primary torsion and its free part; its $3$-primary part is $\mathrm{tors}\,H^{q}(\Ga;M)\otimes\Z_{(3)}$, which (2) computes through $\cmpI^{\ast}$, the
descent's only glueing datum.
\emph{A check outside the build, not a tag} (v1.23.0-pre.17): \file{docs/implementation/slot24/cd1b3\_three\_shortcut\_check.py} computes, independently of
\code{cohomology\_app}, the $3$-primary torsion of $\coker\bigl(\tfrac12(1+M(S))\mid1+M(U)+M(U)^{2}\bigr)$ over $\Z/3^{6}$, certified by the exact-trace value of $r$,
and compares it with the oracle's $3$-part of $H^{1}$, and the absence of $3$-torsion in the oracle's $H^{q}$, $q=2,\dots,5$, on every even-$n$ cell with
$\nu_3(N)=0$ of \file{descentN\_oracle.tsv} ($80$ cells, $N\in\{2,4,5,6,8,9,10,11,12,14,15,16,25\}$, $9\mid N$ at $N=9$) and on $(18,4)$, $(27,2)$, $(27,4)$,
$(45,2)$ (the oracle run for this check): $84=80+4$ cells in two runs, $0$ mismatches; over the grid's $148$ cells with $\nu_3(N)=0$ ($n$ odd included) the oracle has no
$3$-torsion in $H^{\ge2}$, and $22$ of its $32$ even-$n$ cells with $\nu_3(N)>0$ have some (v1.23.0-pre.18, referee m3: the header of
\file{cd1b3\_three\_shortcut\_check.out} records one command per figure). The script imports the oracle's construction of $M_N$, so a convention error there would
cancel (referee m4); the referee's own script, with its own $M_N$ and its own resolution (\file{docs/implementation/slot24/C-D1\_three\_referee.md} \S4), closes that
dependence: $162$ cells with $\nu_3(N)=0$ and $20$ control cells, $0$ mismatches.
\end{proof}

\begin{proposition}[Maps over $\Z$: fibre maps, the level cone and its long exact sequence, Hecke squares, the normal case]\label{prop:descentmaps}
\PROVED{} \STD{\cite[Ch.~III.6.2, III.8--10, VII.6]{brown}, \cite[Thm.~2.2.6]{weibel}; \VERIFYloc} (v1.23.0-pre.18; milestone D2 of
\file{docs/implementation/slot24/DESCENT\_plan.md} v0.3, clause C-D3 of its \S2.7, the owner's LES-0 request of slot~22; refereed ACCEPT-with-fixes, \file{docs/implementation/slot24/C-D3\_referee.md}; its minors m1--m5 at
v1.23.0-pre.19, among them the new item (3)(b$'$), \textbf{not re-refereed}; item (2)(g) and (2)(b)'s signs at v1.23.0-pre.22, from the D2 code audit,
refereed ACCEPT-with-fixes (\file{docs/implementation/slot24/C-D3g\_referee.md}); its fixes M1, m1--m5 at v1.23.0-pre.23, \textbf{not re-refereed}.) Assume
\ref{H4} and \ref{B1}--\ref{B3}. \emph{Models.} For a commutative ring $S$, an \emph{$S$-model} is a complex $\mathsf P_{\bullet}$ of finitely generated projective
$S[\Ga]$-modules, $\mathsf P_q=0$ for $q<0$, with an augmentation $\mathsf P_0\to S$ making it a resolution of $S$; $\mathsf P(M):=\Hom_{S[\Ga]}(\mathsf P_{\bullet},M_S)$. A
\emph{comparison} is an $S[\Ga]$-chain map $\varphi\colon\mathsf P\to\mathsf P'$ over $\id_S$; one exists between any two $S$-models, unique up to homotopy
\STD{\cite[Thm.~2.2.6]{weibel}}, and $\varphi^{\ast}\colon\mathsf P'(M)\to\mathsf P(M)$ is a quasi-isomorphism compatible with the identifications of both with
$H^{\bullet}(\Ga;M_S)=\Ext^{\bullet}_{S[\Ga]}(S,M_S)$. The models in use: $\Iso_{\bullet}\otimes S$ (every $S$); the free tree resolution of
Remark~\ref{rmk:treeconeres} (every $S$); $C_{\bullet}(\Xbulk)\otimes S$ and $C_{\bullet}(\mathcal T)\otimes S$ for $6\in S^{\times}$ (Remark~\ref{rmk:orbifold},
Lemma~\ref{lem:treeres}); $\Iso^{(\ell)}\otimes S$ for $\mathsf O_\ell$-algebras $S$ (Proposition~\ref{prop:certifiedeell}(1)); and the $\ell$-local tree models of the
plan's later clause C-D6. $\cmpI$ (whenever $C_{\bullet}(\Xbulk)\otimes S$ is a model), $\Phi_\ell$, $\Psi_\ell$ and $\cmpI'$ are comparisons. \emph{Fibres.} Let
$(A,\chi)\supseteq(B,\chi')$ be a compatible level pair (Definition~\ref{def:levelpair}, Lemma~\ref{lem:nebcompat}) whose characters take values in
$\mathsf O:=\mathsf O_{m}=\Z[\zeta_{m}]$ for some $m$ ($\mathsf O=\Z$ if $\chi=\chi'=1$), $V(\chi):=V_{\mathsf O}(\chi)$, $M_A:=\Ind_A^{\Ga}V(\chi)$,
$M_B:=\Ind_B^{\Ga}V(\chi')$, $\alpha:=\alpha_{B/A}$, $\omega:=\omega_{B/A}$ (Definition~\ref{def:levelres}, over $\mathsf O$), and $Q:=\coker\alpha\cong\Ind_A^{\Ga}Q_{A/B}$
(Proposition~\ref{prop:levelconeshapiro}(1)); $M_A$, $M_B$, $Q$ are finitely generated and free over $\Z$. $H^{q}(A):=H^{q}(A;V(\chi)_S)$, $H^{q}(B):=H^{q}(B;V(\chi')_S)$,
the ring $S$ read from the context; $\beta_j:=b_j(M_B)-b_j(M_A)$ for $j\in\{0,1\}$ and $\beta_j:=0$ otherwise (the $b_j$ of \S\ref{sec:descentN}).
\begin{enumerate}
\item \textbf{(Fibre maps commute strictly; they glue to maps over $\Z$.)} Let $u\colon M\to M'$ be a $\Z[\Ga]$-map of finitely generated $\Z$-free modules and
  $u_{\mathsf P}:=\Hom(\mathsf P_{\bullet},u_S)\colon\mathsf P(M)\to\mathsf P(M')$.
  \begin{enumerate}
  \item[(a)] $\varphi^{\ast}\circ u_{\mathsf P'}=u_{\mathsf P}\circ\varphi^{\ast}$ \textbf{as cochain maps}, for every comparison $\varphi\colon\mathsf P\to\mathsf P'$, every $S$
    and every degree --- in particular for $\cmpI^{\ast}$ ($=\varepsilon^{\ast}$ of the plan), $\Phi_\ell^{\ast}$, $\Psi_\ell^{\ast}$, $\cmpI'^{\ast}$, and for the chain
    homotopy equivalence $\Ceq(M_S)\simeq\mathrm T(M_S)$ of Proposition~\ref{prop:bredon}(2).
  \item[(b)] If $\mathsf P_q=\bigoplus_x S[\Ga/Q_x]$ ($\abs{Q_x}\in S^{\times}$; every model in use), then $\mathsf P(M)^{q}=\bigoplus_xM_S^{Q_x}$ and $u_{\mathsf P}$ is $u$
    restricted block by block: no lift, no solve.
  \item[(c)] $H^{q}(u_{\mathsf P})=H^{q}(\Ga;u_S)$ under these identifications, so $H^{q}(\Ga;u)\colon H^{q}(\Ga;M)\to H^{q}(\Ga;M')$ is the map of fibre products
    of Proposition~\ref{prop:descentN}(2) induced by $u$ on the $\mathsf O$-piece $\Ceq(M_{\Z[1/6]})$ and on the $\ell$-local pieces, $\ell\in\{2,3\}$, \emph{each read on
    any $\Z_{(\ell)}$- or $\mathsf O_\ell$-model}: $u$ maps the descent diagram of $M$ to that of $M'$, hence the mapping cocones, compatibly with the
    quasi-isomorphisms of Proposition~\ref{prop:descentN}(2). No coherence datum enters.
  \item[(d)] \textbf{(Instances, in every degree, over every ring.)} Under Shapiro's isomorphism, $H^{q}(\Ga;\alpha)=\res^{A}_{B}$, $H^{q}(\Ga;\omega)=\cor^{A}_{B}$ and, for a
    normal pair (Definition~\ref{def:normalpair}), $H^{q}(\Ga;R_\sigma)=c_\sigma$, the conjugation action of $\sigma\in A$ on $H^{q}(B)$ --- for every $q$ and every $S$,
    with no hypothesis on $6$ (contrast Lemma~\ref{lem:levelresmorph}(3) and Proposition~\ref{prop:levelinfres}(1), stated on the model for $6\in S^{\times}$). Hence over $\Z$:
    $\cor\circ\res=[A:B]$, $c_\sigma\circ\res=\res$, $\res\circ\cor=\mathrm{Nm}_{\mathsf G}$ (normal pair), in every degree; and if $[A:B]\in S^{\times}$,
    $H^{q}(B)=\res H^{q}(A)\oplus\ker\cor^{q}$ for \textbf{every} $q$ (Proposition~\ref{prop:levelsplit} in all degrees). The diamonds $R_\sigma$ of
    Proposition~\ref{prop:levelconehecke}(2) and the $\mathsf O$-action on $M_A$, $M_B$ are fibre maps as well.
  \end{enumerate}
\item \textbf{(The level cone on every model; the long exact sequence over $\Z$.)}
  \begin{enumerate}
  \item[(a)] For every $S$-model $\mathsf P$, $0\to\mathsf P(M_A)\xrightarrow{\alpha_{\mathsf P}}\mathsf P(M_B)\to\mathsf P(Q)\to0$ is exact, and
    $\pi_{\mathsf P}\colon\Cone{\alpha_{\mathsf P}}\to\mathsf P(Q)$, $(c,c')\mapsto[c']$, is a surjective quasi-isomorphism, where
    $\Cone{\alpha_{\mathsf P}}^{q}=\mathsf P(M_A)^{q+1}\oplus\mathsf P(M_B)^{q}$ ($q\ge-1$), $\Dcone(c,c')=(-dc,\,dc'+\alpha_{\mathsf P}c)$, as in Definition~\ref{def:levelcone}. For a
    comparison $\varphi$, $(\varphi^{\ast}[1],\varphi^{\ast})$ is a cochain map of cones commuting with $\pi$ and with the two maps of the cone's short exact sequence, and a
    quasi-isomorphism.
  \item[(b)] $H^{q}(A,B):=H^{q-1}(\Cone{\alpha_{\mathsf P}})\cong H^{q-1}(\Ga;Q_S)\cong H^{q-1}(A;Q_{A/B}\otimes S)$, independently of $\mathsf P$, and the long sequence
    \begin{equation}\label{eq:descentmapsles}
      0\to H^{0}(A)\xrightarrow{\res}H^{0}(B)\to H^{1}(A,B)\to H^{1}(A)\xrightarrow{\res}H^{1}(B)\to H^{2}(A,B)\to H^{2}(A)\xrightarrow{\res}H^{2}(B)\to H^{3}(A,B)\to\cdots
    \end{equation}
    is exact \textbf{in every degree} over every $S$; over $\Z$ ($\mathsf P=\Iso$) it is the long exact sequence of $0\to V(\chi)\to\Ind_B^{A}V(\chi')\to Q_{A/B}\to0$ over $A$,
    and over $\Z_{(\ell)}$ its localisation. $24\cdot H^{q}(A,B)=0$ for $q\ge3$, and $\res^{\Ga}_{A}e\cup-$ is an isomorphism $H^{q}(A,B)\to H^{q+2}(A,B)$ for $q\ge3$ (so
    $q_{\max}=4$ determines every degree $\ge3$). Over a $\Z[1/6]$-algebra \eqref{eq:descentmapsles} is the six-term sequence \eqref{eq:levelsixterm}, read as group
    cohomology (Proposition~\ref{prop:levelconeshapiro}(2)); at $\ell\in\{2,3\}$ it is \textbf{not} six-term in general: e.g.\ $H^{3}(\Ga_0(4),\Ga_0(8))=\Z/2$ at $n=0$, and
    $H^{3}(\Ga,\Ga_0(11))=\Z/3$ at $n=4$ (on every one of the $74$ cells of the check below some $H^{q}(A,B)$, $q\ge3$, is non-zero).
    \textbf{Signs} (v1.23.0-pre.22): the maps $H^{q}(B)\to H^{q+1}(A,B)$ and $H^{q}(A,B)\to H^{q}(A)$ of \eqref{eq:descentmapsles} are induced by the two maps $c'\mapsto(0,c')$ and
    $(c,c')\mapsto c$ of the cone's short exact sequence; under $\pi_{\mathsf P}$ the second is $-\delta$, $\delta\colon H^{q-1}(\Ga;Q_S)\to H^{q}(\Ga;M_{A,S})$ the connecting map of
    $0\to M_A\to M_B\to Q\to0$, with the convention $\delta[x]=[d\tilde x]$ ($\tilde x$ a lift, the cochain differential carrying no sign, as for the snake-lemma
    connecting maps elsewhere in this document); the ``$-$'' depends on that convention. Immaterial for groups, kernels, images, cokernels and exactness, and for
    equivariance ($T\circ\delta=\delta\circ T$ is unchanged under $\delta\mapsto-\delta$, so (3)(c) is unaffected); it matters only when the map is compared entrywise with an
    independently computed $\delta$, or composed with maps defined through $\delta$ (LES-1's predictions, a Bockstein compatibility). (v1.23.0-pre.23, referee m5.)
  \item[(c)] \textbf{(Glueing.)} $\Cone{\alpha_{\Iso}}$ is a complex of finitely generated free $\Z$-modules, and Proposition~\ref{prop:cohglue} applies to it: $H^{q}(A,B)$
    over $\Z$ is the fibre product of the cohomology of $\Cone{\alpha_{C_{\bullet}(\Xbulk)}}$ over $\Z[1/6]$ (the $\mathsf O$-piece; the only glueing datum is
    $(\cmpI^{\ast}[1],\cmpI^{\ast})$) and of the cones on any $\Z_{(2)}$- and $\Z_{(3)}$-models, over $H^{q}(\Cone{\alpha_{\Iso}}\otimes\Qsix)$. The maps of
    \eqref{eq:descentmapsles} (both maps of the cone's short exact sequence and $\res$) commute strictly with every comparison, so the sequence glued from the local
    sequences is \eqref{eq:descentmapsles} over $\Z$, and every local sequence is exact.
  \item[(d)] \textbf{(The certified exponent on the cone.)} On every $\Z_{(\ell)}$-model whose cochains are finitely generated free $\Z_{(\ell)}$-modules
    ($\Iso\otimes\Z_{(\ell)}$, $\Iso^{(\ell)}\otimes\Z_{(\ell)}$, and the models of~(f)), $\Cone{\alpha_{\mathsf P}}$ is a bounded below complex of finitely generated free
    $\Z_{(\ell)}$-modules, and Proposition~\ref{prop:certifiede}(1)--(3), Remark~\ref{rmk:certifiedeperdiff} and Proposition~\ref{prop:certifiedeell}(6) (the Smith-valuation
    read-out and the stopping rule) hold for its differentials $\Dcone^{q}$, $q\ge-1$, verbatim. $\rank_{\Q}\Dcone^{q}$ is admissible from exactly two sources: (i) the rank
    over $\F_{\ell'}$ of the reduction mod $\ell'$ of \textbf{this} cone on the model's integral form --- $\Iso$ over $\Z$, $\Iso^{(\ell)}$ over $\mathsf O_\ell$ with
    $\ell'\nmid\mathsf q_\ell$ (as in Proposition~\ref{prop:certifiedeell}(4)(c); a $\Z_{(\ell)}$-module reduced mod $\ell'\ne\ell$ is $0$), the tree and Bredon complexes of~(f)
    over $\Z[1/6]$ with $\ell'\ge5$ (v1.23.0-pre.19, referee m4) ---, at an $\ell'$ that qualifies for $\Dcone^{q}$, i.e.\ at which $H^{q+1}(\Ga;Q)$ has no $\ell'$-torsion --- every
    prime $\ell'\ge5$ when $q\ne0$, and for $q=0$ iff $s_{\ell'}(Q)=0$; (ii) the Betti recursion
    $\rank_{\Q}\Dcone^{q}=\rank\Cone{\alpha_{\mathsf P}}^{q}-\beta_q-\rank_{\Q}\Dcone^{q-1}$ ($\Dcone^{-2}:=0$), with $b_j(Q)=\beta_j$. The rank of a differential of
    $\mathsf P(M_A)$ or $\mathsf P(M_B)$ alone is not admissible for the cone (Proposition~\ref{prop:certifiedeell}(4)(c)). Equivalently the rule may be run on the
    smaller complex $\mathsf P(Q)$ (Proposition~\ref{prop:descentN} for $M=Q$): the groups are the same by~(a); the maps of \eqref{eq:descentmapsles} need the cone.
  \item[(e)] \textbf{(Bad primes of the cone.)} For $\ell\ge5$, $H^{q}(A,B)$ has no $\ell$-torsion for $q\ne2$, and the $\ell$-primary torsion of
    $H^{2}(A,B)=H^{1}(\Ga;Q)$ has exactly $s_\ell(Q)=\dim_{\F_\ell}H^{0}(\Ga;Q/\ell)-\beta_0$ cyclic summands (Proposition~\ref{prop:badprimesN}(1) for $M=Q$). If moreover
    $\ell\nmid[A:B]$, then $s_\ell(Q)=s_\ell(M_B)-s_\ell(M_A)$; for $\Ga_0(MN)\subseteq\Ga_0(N)$, $\chi=1$, these are the closed forms of
    Proposition~\ref{prop:badprimesN}(2).
  \item[(f)] \textbf{(The cheap routes.)} \emph{At $\ell\ge5$} $C_{\bullet}(\mathcal T)\otimes\Z_{(\ell)}$ and $C_{\bullet}(\Xbulk)\otimes\Z_{(\ell)}$ are $\Z_{(\ell)}$-models, so
    (a)--(d) hold for the tree and the Bredon cones; $\mathrm{tors}\,H^{2}(A,B)\otimes\Z_{(\ell)}=\mathrm{tors}\coker d^{0}$ for $d^{0}$ of $\mathrm T(Q_{\Z_{(\ell)}})$, read by
    the sparse cokernel route of Proposition~\ref{prop:descentN}(7)(b) for $M=Q$: the images under $M_B\to Q$ of generators of
    $(M_B)^{\langle S\rangle}+(M_B)^{\langle U\rangle}$ over $\Z/\ell^{K}$ generate $Q^{\langle S\rangle}+Q^{\langle U\rangle}$, and the comparison value is
    $r=\rank\mathrm T^{0}(Q)-\beta_0$ with $\rank\mathrm T^{j}(Q)=\rank\mathrm T^{j}(M_B)-\rank\mathrm T^{j}(M_A)$. \emph{At $\ell=3$}, for $A=\Ga_0(N)$, $B=\Ga_0(MN)$,
    $\chi=1$ and $\nu_3(N)=0$ (so $\nu_3(MN)=0$): the same holds over $\Z_{(3)}$ --- $\mathrm T$ and $\Ceq$ compute $H^{\bullet}(A;Q_{A/B}\otimes S)$ for every
    $\Z_{(3)}$-algebra $S$, $\cmpI^{\ast}\otimes\Z_{(3)}$ is a quasi-isomorphism of cones, $H^{q}(A,B)\otimes\Z_{(3)}=0$ for $q\ge3$, the $3$-local sequence
    \eqref{eq:descentmapsles} is six-term, and $\mathrm{tors}\,H^{2}(A,B)\otimes\Z_{(3)}$ is read as in Proposition~\ref{prop:descentNthree}(3) for $M=Q$ (the images of
    (3)(a), (3)(b)'s generators of $M_B$; $r=\rank\mathrm T^{0}(Q)-\beta_0$, $\rank\mathrm T^{0}(Q)$ the difference of the closed forms of
    Proposition~\ref{prop:descentNthree}(3) at $MN$ and at $N$). At $\ell=2$, and at $\ell=3$ when $\nu_3(N)>0$, the cone is read on $\Iso^{(\ell)}$ or another
    $\Z_{(\ell)}$-model.
  \item[(g)] \textbf{(Generators, coordinates and maps from a certified Smith form: the precision rule.)} (v1.23.0-pre.22; the D2 code audit's F1; refereed ACCEPT-with-fixes, \file{docs/implementation/slot24/C-D3g\_referee.md}, its fixes at v1.23.0-pre.23.)
    (d) and Proposition~\ref{prop:certifiedeell}(6) certify Smith \emph{valuations}, i.e.\ group types; the maps of~(1) and \eqref{eq:descentmapsles} are read from the
    \emph{transforms}, and this item says how far those are exact. Let $A$ be an $r\times c$ matrix over $\Z_{(\ell)}$, $K\ge1$, and $LAR=\Sigma$ over $\Z/\ell^{K}$ with
    $L\in\GL_r(\Z/\ell^{K})$, $R\in\GL_c(\Z/\ell^{K})$, $\Sigma$ zero but for the diagonal entries $u_1\ell^{\sigma_1},\dots,u_s\ell^{\sigma_s}$ in its first $s$ places, $u_i$ units (the code does not normalise its
    pivots; v1.23.0-pre.23, referee m4), every $\sigma_i<K$; put $D:=\mathrm{diag}(u_i\ell^{\sigma_i})$. Suppose the form is \textbf{certified}: $s=\rank_{\Q}A$. Put $\sigma_{\max}:=\max_i\sigma_i$ ($0$ if $s=0$) and $K':=K-\sigma_{\max}\ge1$. Then:
    \begin{enumerate}
    \item[(i)] There are $\tilde L\in\GL_r(\Z_{(\ell)})$, $\tilde R\in\GL_c(\Z_{(\ell)})$ with $\tilde LA\tilde R=\mathrm{diag}(D,0)$
      \textbf{exactly}, and $\tilde L\equiv L$, $\tilde L^{-1}\equiv L^{-1}$, $\tilde R\equiv R$, $\tilde R^{-1}\equiv R^{-1}\pmod{\ell^{K'}}$. \textbf{Moreover}
      (v1.23.0-pre.23, referee M1) the first $s$ rows of $\tilde L$, the last $r-s$ columns of $\tilde L^{-1}$ and the last $c-s$ rows of $\tilde R^{-1}$ are $\equiv$ those of
      $L$, $L^{-1}$, $R^{-1}$ \textbf{modulo $\ell^{K}$}; only the other blocks (the last $r-s$ rows of $L$, the first $s$ columns of $L^{-1}$, all of $R$, the first $s$ rows of
      $R^{-1}$) are exact merely modulo $\ell^{K'}$.
    \item[(ii)] Hence, modulo $\ell^{K'}$ and for one exact choice: the last $c-s$ columns of $R$ are a basis of $\ker A$ (saturated in $\Z_{(\ell)}^{c}$), the matching rows of
      $R^{-1}$ its coordinates; the columns of $L^{-1}$ are generators of $\coker A\cong\bigoplus_{\sigma_i>0}\Z/\ell^{\sigma_i}\oplus\Z_{(\ell)}^{r-s}$, the $i$-th of order
      $\ell^{\sigma_i}$ ($i\le s$), and the rows of $L$ the coordinate functionals: the class of $y$ is $((Ly)_i\bmod\ell^{\sigma_i})_{i\le s}$, $((Ly)_i)_{i>s}$. By the
      ``moreover'' of (i), the torsion coordinates $(Ly)_i\bmod\ell^{\sigma_i}$, $i\le s$, are exact for \textbf{every} $\sigma_i<K$; the torsion generators (the first $s$
      columns of $L^{-1}$) and the free part (coordinates and generators) are exact modulo $\ell^{K'}$. A read-out that refuses unless $2\sigma_{\max}\le K$ (the code's
      \code{BuildPres}) applies a \textbf{conservative policy}, sufficient and not necessary (v1.23.0-pre.23, referee m1). (This is about reading exact vectors. A read
      that applies a map to the torsion generators, Proposition~\ref{prop:torsionhecke}(3), \textbf{does} need $2\sigma_{\max}\le K$ for the module; v1.23.0-pre.25, the
      C-D4 referee's m1.) \textbf{Lattice coordinates.} For $y$ in the lattice
      $A\Z_{(\ell)}^{c}$ spanned by the columns of $A$, its coordinates in the basis $(\tilde L^{-1}De_i)_{i\le s}$ of that lattice are $D^{-1}(\tilde Ly)_{\le s}$, known modulo
      $\ell^{K-\sigma_i}$ from $(Ly)_{\le s}$, and $(Ly)_{>s}\equiv0\pmod{\ell^{K}}$ (v1.23.0-pre.23, referee M1; the code's \code{LatticeQuotient}).
    \item[(iii)] \textbf{(Composition.)} A matrix formed by sums and products from data so read is known modulo $\ell^{K'}$ and is the reduction of the same formula in
      $\tilde L,\tilde R$ over $\Z_{(\ell)}$; once its own $\rank_{\Q}$ is certified, (i), (ii) apply to it with $K'$ in place of $K$. \textbf{Other operations}
      (v1.23.0-pre.23, referee m2): an exact division by $\ell^{\sigma_i}$ costs $\sigma_i$ digits (a datum known modulo $\ell^{k}$ and divisible by $\ell^{\sigma_i}$ gives a
      quotient known modulo $\ell^{k-\sigma_i}$); reducing a torsion row modulo $\ell^{a_i}$ costs nothing (its entries are defined only modulo $\ell^{a_i}$, and any lift is
      admissible when the relation matrix $\Delta$ is adjoined); data from several forms combine at the \textbf{minimum} of their precisions. \textbf{Certifying ranks}
      (v1.23.0-pre.23, referee m3): since the exact matrix of a derived form depends on the unknown $\tilde L$, $\tilde R$, its certifying $\rank_{\Q}$ must be an invariant of
      the map, not of a matrix; the admissible ones are, besides (d)'s two sources for $\rank_{\Q}D^{q}_{\mathrm{cone}}$: for an induced map $f\colon X\to Y$ with presentation
      relations $\Delta_Y$ ($t_Y$ torsion generators), $\rank_{\Q}[F\mid\Delta_Y]=t_Y+\rank_{\Q}f$; for the kernel lattice $L_f=\{x:Fx\in\operatorname{im}\Delta_Y\}$,
      $\rank_{\Q}L_f=g_X-\rank_{\Q}f$ ($g_X$ generators of $X$); for the coordinate matrix $C$ of a sublattice $I\subseteq L_f$, $\rank_{\Q}I$: $t_X$ for
      $I=\operatorname{im}\Delta_X$ (the kernel of $f$), $t_X+\rank_{\Q}g$ for $I=F_g\Z_{(\ell)}^{g_W}+\operatorname{im}\Delta_X$ (the comparison
      $\operatorname{im}g=\ker f$); and $\rank_{\Q}f$ of the maps themselves: $\res^{q}$, $\cor^{q}$: $b_q(A)$; $j^{q}\colon H^{q}(B)\to H^{q+1}(A,B)$: $\beta_q$;
      $(c,c')\mapsto c$: $0$; $c_\sigma$: $b_q(B)$; for a normal pair the invariants map $x\mapsto(c_\sigma x-x)_\sigma$ on $H^{1}(B)$: $b_1(B)-b_1(A)$ (its rational
      kernel is $H^{1}(B)^{\mathsf G}_{\Q}=\res H^{1}(A)_{\Q}$). These hold because rationally
      $\cor\circ\res=[A:B]$ ((1)(d)), so $\res$ is injective and the rational long sequence splits into short exact sequences
      $0\to H^{q}(A)_{\Q}\to H^{q}(B)_{\Q}\to H^{q+1}(A,B)_{\Q}\to0$. Along a chain of certified Smith forms
      the precisions therefore decrease additively, $K\ge K'\ge K''\ge\cdots$, and kernels, images and cokernels of induced maps, as well as the subgroup comparisons
      ``$\operatorname{im}=\ker$'', computed so are exact \textbf{provided} every exponent read stays below the precision left at its step; otherwise the read-out is
      refused (the stopping rule of Proposition~\ref{prop:certifiedeell}(6)), never reported.
    \end{enumerate}
    Without the certification nothing is asserted: if $s<\rank_{\Q}A$, $(L,\Sigma,R)$ need not be the reduction of any exact Smith form.
  \end{enumerate}
\item \textbf{(Hecke squares.)} Let $A=\Ga_0(N)\supseteq B=\Ga_0(MN)$, $\chi'=\chi\circ\mathrm{red}_{MN\to N}$, $g=\smat100p$, and $[AgA]$, $[BgB]$ the operators of
  Proposition~\ref{prop:descentN}(5) on $H^{\bullet}(A)$, $H^{\bullet}(B)$ ($f=\rho_{n,\chi}(g)$, resp.\ $\rho_{n,\chi'}(g)$).
  \begin{enumerate}
  \item[(a)] If $p\nmid M$, then $\res\circ[AgA]=[BgB]\circ\res$ and $\cor\circ[BgB]=[AgA]\circ\cor$ on $H^{q}$ for \textbf{every} $q$, over every ring. For $p\nmid MN$ these
    are $T_p$ at both levels; for $p\mid N$, $p\nmid M$, they are $U_p$ at both levels --- \textbf{this document's} $U_p=[AgA]$, which on $p$-old spaces is the adjoint
    $U_p^{\ast}$ of the classical operator (Proposition~\ref{prop:levelconehecke}(3)(d)). Over $\Z$ the squares hold on the glued groups and are read on any of their
    local components ((1)(c), Proposition~\ref{prop:descentN}(3)); on the model over every ring they hold as cochain maps by the criterion
    Proposition~\ref{prop:levelconehecke}(1), whose two hypotheses $B\cap A''=B''$ and $A=A''B$ the proof verifies (its instance (2) for $U_p$ asks for equal
    $\Ga_0$-levels, which $\Ga_0(MN)\subset\Ga_0(N)$, $M>1$, does not have; v1.23.0-pre.19, referee m2).
  \item[(b)] If $p\mid M$ (level raising at $p$) the square with $\res$ fails over $\Z$ in general \textbf{for this document's operator} $[AgA]$, $g=\smat100p$: at
    $(\Ga_0(25)\subset\Ga_0(5),n=0,\chi=1)$ on $H^{1}$ (Proposition~\ref{prop:levelconehecke}(3)(c) at $\ell=7$). For $p\mid N$ this is the convention, not level raising:
    $[AgA]$ is the adjoint $U_p^{\ast}$ on $p$-old spaces, and the \textbf{classical} $U_p=[A\,g'A]$, $g'=\smat p001$, satisfies the criterion
    Proposition~\ref{prop:levelconehecke}(1) on every level-raising pair at $\chi=1$, so its square with $\res$ commutes, as cochain maps
    (Proposition~\ref{prop:levelconehecke}(3)(d)). For $p\nmid N$, $p\mid M$, $[AgA]$ is $T_p$ at level $N$ and $U_p$ at level $MN$, and both squares fail in
    general: at $(\Ga_0(2)\subset\Ga,\ n=2,\ p=2,\ \chi=1)$ on $H^{1}$ over $\Z$ (v1.23.0-pre.19, referee m1).
  \item[(b$'$)] \textbf{(Corestriction and $U_p$, every $M$.)} If $p\mid N$, then
    \[
      \cor^{A}_{B}\circ[BgB]=[AgA]\circ\cor^{A}_{B}\quad\text{on }H^{q}\text{ for \textbf{every} }q\text{, every }M\text{ (}p\mid M\text{ included) and over every ring,}
    \]
    and over $\Z$ on the glued groups ((1)(c)); here $[AgA]$, $[BgB]$ are this document's $U_p$ at levels $N$ and $MN$. In particular, on the Hida $p$-tower
    $\Ga_0(p^{r+1})\subset\Ga_0(p^{r})$, $r\ge1$, \textbf{corestriction is $U_p$-equivariant over $\Z$ in every degree}, and restriction is not, for this $U_p$ ((b)). If $p\nmid N$ and
    $p\mid M$, the square with $\cor$ fails in general, at the cell of~(b) (v1.23.0-pre.19, referee m1).
  \item[(c)] \textbf{(The relative term.)} If $p\nmid M$, then over $\Z[1/6]$-algebras the pair $(T^{A}_g,T^{B}_g)$ of Proposition~\ref{prop:levelconehecke}(1) is a cochain endomorphism of the cone
    on the model commuting with its short exact sequence, so \eqref{eq:descentmapsles}$\otimes\Z[1/6]$ is Hecke-equivariant. At $\ell\in\{2,3\}$ a Hecke operator on
    $H^{q}(A,B)\otimes\Z_{(\ell)}$ and the equivariance of the connecting map are \textbf{not asserted}: they need cochain-level operators $T^{A}$, $T^{B}$ on an $\ell$-local
    model with $T^{B}\alpha_{\mathsf P}=\alpha_{\mathsf P}T^{A}$ \emph{strictly} --- the level-$N$ operators of the plan's C-D4. Given such operators on $\Iso$, their transports
    $\Psi^{\ast}T\Phi^{\ast}$ to any other model satisfy the same strict identity by~(1)(a), and $(T^{A}[1],T^{B})$ makes \eqref{eq:descentmapsles}$\otimes\Z_{(\ell)}$ equivariant.
    (v1.23.0-pre.24: such operators are Proposition~\ref{prop:strictliftN}(3), with the strict identity in its (7).)
  \end{enumerate}
\item \textbf{(The normal case: inflation--restriction over $\Z$.)} Let the pair be normal, $\mathsf G:=A/B$.
  \begin{enumerate}
  \item[(a)] Over every ring and in every degree, $\mathsf G$ acts on $H^{q}(B)$ through $c_\sigma=H^{q}(\Ga;R_\sigma)$, glued strictly by~(1); $\res$ lands in
    $H^{q}(B)^{\mathsf G}$, $\res\circ\cor=\mathrm{Nm}_{\mathsf G}$, $\cor\circ\res=\abs{\mathsf G}$; invariants commute with the flat base changes of the descent, so
    $H^{q}(B)^{\mathsf G}$ over $\Z$ is glued from its local components.
  \item[(b)] \STD{\cite[Ch.~VII, \S6]{brown}; \VERIFYloc} The Hochschild--Serre spectral sequence $E_2^{a,b}=H^{a}(\mathsf G;H^{b}(B))\Rightarrow H^{a+b}(A)$ exists over every ring, with
    the five-term exact sequence $0\to H^{1}(\mathsf G;H^{0}(B))\xrightarrow{\inf}H^{1}(A)\xrightarrow{\res}H^{1}(B)^{\mathsf G}\xrightarrow{\mathrm{tg}}H^{2}(\mathsf G;H^{0}(B))
    \xrightarrow{\inf}H^{2}(A)$. Over $\Z[1/6]$-algebras its last map is $0$ and the rest is Proposition~\ref{prop:levelinfres}(2)--(5); over $\Z_{(2)}$, $\Z_{(3)}$ the rows
    $b\ge2$ need not vanish when $B$ has torsion ($\Ga_0(MN)\trianglelefteq\Ga_0(N)$, $M\mid\gcd(N,24)$, contains $-I$), and neither the surjectivity of $\mathrm{tg}$ nor
    Proposition~\ref{prop:levelinfres}(2), (3)'s second half, (5) is asserted there.
  \item[(c)] \textbf{(Torsion-free $B$: two rows over every ring.)} If $B$ is torsion-free ($\Ga_1(N)$, $N\ge4$, and its subgroups, $\Ga(3)$:
    Proposition~\ref{prop:modeltorsion}), then $H^{b}(B;X)=0$ for $b\ge2$, every $X$ and every ring, and over every $S$
    \[
      0\to H^{1}(\mathsf G;H^{0}(B))\to H^{1}(A)\to H^{0}(\mathsf G;H^{1}(B))\xrightarrow{d_2}\cdots\to H^{a}(\mathsf G;H^{0}(B))\xrightarrow{\inf}H^{a}(A)\to
      H^{a-1}(\mathsf G;H^{1}(B))\xrightarrow{d_2}H^{a+1}(\mathsf G;H^{0}(B))\to\cdots
    \]
    is exact. If $n>0$ and $S$ is $\Z$-torsion-free ($\Z$, $\Z_{(\ell)}$, $\Z[1/6]$, \dots), then $H^{0}(B)=0$ and
    \[
      H^{q}(A;V(\chi)_S)\;\cong\;H^{q-1}\bigl(\mathsf G;H^{1}(B;V(\chi')_S)\bigr)\qquad\text{for every }q\ge1
    \]
    ($q=1$: $\res$ is an isomorphism onto $H^{1}(B)^{\mathsf G}$). Hence $\abs{\mathsf G}\cdot H^{q}(A;V(\chi))=0$ for $q\ge2$; for $\ell\nmid\abs{\mathsf G}$,
    $H^{q}(A;V(\chi))\otimes\Z_{(\ell)}=0$ for $q\ge2$ and $H^{1}(A)\otimes\Z_{(\ell)}=(H^{1}(B)\otimes\Z_{(\ell)})^{\mathsf G}$; for $\mathsf G$ cyclic,
    $H^{q}(A)\cong H^{q+2}(A)$ for $q\ge2$ by the periodic resolution of Proposition~\ref{prop:cycliccone}(1). \textbf{The primary instance} is the bottom of the Hida
    $p$-tower: $A=\Ga_0(N)$, $B=\Ga_1(N)$, $N\ge4$, $\mathsf G=(\Z/N)^{\times}$, in particular $N=p^{r}\ge4$ and $n>0$:
    $H^{q}(\Ga_0(p^{r});\Sym^{n}\Z^{2})\cong H^{q-1}\bigl((\Z/p^{r})^{\times};H^{1}(\Ga_1(p^{r});\Sym^{n}\Z^{2})\bigr)$ for $q\ge2$, all of it $2$- and $3$-primary
    ($24$ kills $H^{\ge2}(A)$, Proposition~\ref{prop:certifiedeell}(4)(a)); its $\ell$-primary part vanishes for $\ell\in\{2,3\}$, $\ell\nmid p^{r-1}(p-1)$, and
    $H^{\ge1}\bigl((\Z/p^{r})^{\times};H^{1}(\Ga_1(p^{r});\Sym^{n}\Z^{2})\bigr)$ has no $\ell$-torsion for $\ell\ge5$, even when $\ell\mid p^{r-1}(p-1)$ (v1.23.0-pre.19, referee m5). If $A$ is
    torsion-free as well (the steps $\Ga_1(p^{r+1})\trianglelefteq\Ga_1(p^{r})\cap\Ga_0(p^{r+1})\trianglelefteq\Ga_1(p^{r})$, $p^{r}\ge4$, of
    Proposition~\ref{prop:hidanormalsteps}), the model computes group cohomology over every ring (Proposition~\ref{prop:modeltorsion}), and
    Lemma~\ref{lem:levelconeexact} and Proposition~\ref{prop:levelinfres}(2)--(5) hold over $\Z$ verbatim.
  \end{enumerate}
\item \textbf{(What the tree must supply.)} For an $\ell$-local tree model $\mathsf P_{\mathcal T}$ of the plan's D5 (clause C-D6), $\ell\in\{2,3\}$, to carry items (1), (2) and
  (4) it suffices that (T1) $\mathsf P_{\mathcal T}$ is an $S$-model for $S=\mathsf O'$ a localisation of $\Z$ in which $\ell$ is not inverted, or $S=\Z_{(\ell)}$; (T2)
  $\mathsf P_{\mathcal T,q}=\bigoplus_xS[\Ga/Q_x]$ with $\abs{Q_x}\in S^{\times}$, so that $\mathsf P_{\mathcal T}(M)^{q}=\bigoplus_xM_S^{Q_x}$ is finitely generated and free
  over $\Z_{(\ell)}$ after base change and fibre maps act block by block ((1)(b), (2)(d)); (T3) a comparison $\Phi_{\mathcal T}\colon\Iso\otimes S\to\mathsf P_{\mathcal T}$
  (or $\Psi_{\mathcal T}$ the other way), given explicitly when classes are to be transported or the model is to be checked against $\Iso$ at a test cell. Nothing else:
  no strictness over $C_{\bullet}(\Xbulk)$, no Hecke datum, no boundary. Item (3)(c) at $\ell$ needs in addition the operators of C-D4, transported as there.
\end{enumerate}
\end{proposition}

\begin{proof}
(1)(a) $\varphi^{\ast}(u_{\mathsf P'}c)=u\circ c\circ\varphi=u_{\mathsf P}(\varphi^{\ast}c)$: $\Hom(-,-)$ is a bifunctor, and nothing is assumed of $S$, $\varphi$ or the degree.
$\cmpI\colon\Iso\to C_{\bullet}(\Xbulk)$, $\Phi_\ell$, $\Psi_\ell$, $\cmpI'$ are $S[\Ga]$-chain maps over $\id_S$ (Propositions~\ref{prop:isomodel}(4),
\ref{prop:isomodelell}(4)), and the equivalence of Proposition~\ref{prop:bredon}(2) is induced by chain maps on the chain side (as used in the proof of
Lemma~\ref{lem:levelconeexact}(2)). (b) $\Hom_{S[\Ga]}(S[\Ga/Q],M_S)\to M_S^{Q}$, $c\mapsto c(Q)$, is an isomorphism, and $(u\circ c)(Q)=u(c(Q))$. (c) $H^{q}(\mathsf P(M))\cong
\Ext^{q}_{S[\Ga]}(S,M_S)$ through comparisons, naturally in $M$ \STD{\cite[Thm.~2.2.6]{weibel}}; $\Ext$ is a functor in the second variable, so $H^{q}(u_{\mathsf P})$ is
$H^{q}(\Ga;u_S)$. In Proposition~\ref{prop:descentN}(2) the corners $\Ceq(M_{\mathsf O})$, $\CI(M)\otimes\Zsix$, $\CI(M)\otimes\Qsix$ and the maps $\alpha$ (base change after
$\cmpI^{\ast}[1/6]$) and $\beta$ are natural in $M$ by~(a), so $u$ induces a morphism of diagrams and of their mapping cocones; the quasi-isomorphisms of
Proposition~\ref{prop:fracture} ($M\to$ cocone) and of the proof of Proposition~\ref{prop:descent} (the map of cocones induced by $\cmpI^{\ast}[1/6]$) are natural as well. By
Proposition~\ref{prop:cohglue}'s morphism clause $H^{q}(\Ga;u)$ is determined by its two components, and the $\ell$-primary part of the $\Zsix$-component may be read through
any comparison by~(a). (d) Take $\mathsf P=\Iso\otimes S$: it is free over $S[\Ga]$, hence over $S[A]$ and $S[B]$, a projective resolution of $S$ over both, so the
reciprocity $\phi\mapsto\phi(\cdot)(I)$ of Proposition~\ref{prop:levelcorrespondence}(3) identifies $\Hom_{S[\Ga]}(\mathsf P,\Ind_A^{\Ga}X)$ with
$\Hom_{S[A]}(\mathsf P,X)$ and computes Shapiro's isomorphism \STD{\cite[Ch.~III.6.2]{brown}}. The computation in the proof of Lemma~\ref{lem:levelresmorph}(3) is
algebraic and uses only the $\Ga$-equivariance of $\phi$, so it holds for $\mathsf P$ verbatim: $\alpha_{\mathsf P}$ is ``read an $A$-equivariant cochain as a $B$-equivariant
one'', which induces $\res$, and $\omega_{\mathsf P}$ is the transfer formula, which induces $\cor$ \STD{\cite[Ch.~III.9]{brown}}; the role of $6\in S^{\times}$ there was only
to make $C_{\bullet}(\Xbulk)\otimes S$ a projective resolution, which $\Iso\otimes S$ is over every $S$. For $R_\sigma$: $\phi(\sigma^{-1}y)(I)=(\sigma^{-1}\phi(y))(I)=
\phi(y)(\sigma^{-1})$, so $(R_\sigma\circ\phi)(y)(I)=\rho_{n,\chi}(\sigma)\,\phi(y)(\sigma^{-1})$ and the reciprocity carries $R_{\sigma,\mathsf P}$ to
$\psi\mapsto\rho_{n,\chi}(\sigma)\psi(\sigma^{-1}\,\cdot)$ on $\Hom_{S[B]}(\mathsf P,V(\chi'))$, the conjugation $c_\sigma$ computed on a projective resolution over the
bigger group $A$ \STD{\cite[Ch.~III.8]{brown}}. The identities are the fibre identities of Lemma~\ref{lem:levelresmorph}(2) and Proposition~\ref{prop:levelinfres}(1), which
hold over every ring, pushed through $H^{q}(\Ga;-)$. For $[A:B]\in S^{\times}$ the idempotent $e$ of Proposition~\ref{prop:levelsplit} is a $\Ga$-endomorphism of
$M_{B,S}$, and $\mathsf P(-)$ is additive, so $\mathsf P(M_B)=\alpha_{\mathsf P}\mathsf P(M_A)\oplus\mathsf P(\ker\omega)$ as complexes in every degree, and
$\cor\circ\res=[A:B]$ is invertible. The diamonds commute with $\alpha$ on the fibre (Proposition~\ref{prop:levelconehecke}(2)); the $\mathsf O$-action commutes with
the $\Ga$-action by definition.

(2)(a) $\alpha$ is injective and split over $\mathsf O$ (Lemma~\ref{lem:levelresmorph}(1)), hence over $\Z$ ($\mathsf O$ is free over $\Z$), so $0\to M_{A,S}\to M_{B,S}\to
Q_S\to0$ is exact for every $S$, with $Q_S=Q\otimes S$; each $\mathsf P_q$ is projective, so $\Hom_{S[\Ga]}(\mathsf P_q,-)$ is exact. The rest is the proof of
Proposition~\ref{prop:levelconeshapiro}(2) with this exactness in place of the averaging over $\abs{\Ga_\sigma}\mid12$: $\pi_{\mathsf P}$ is a surjective cochain map whose
kernel is $\{(c,\alpha_{\mathsf P}c'')\}\cong\Cone{\id}$, acyclic. For a comparison, $\varphi^{\ast}$ commutes with $\alpha$ by (1)(a), so $(\varphi^{\ast}[1],\varphi^{\ast})$
commutes with $\Dcone$ and with both maps of $0\to\mathsf P(M_B)\to\Cone{\alpha_{\mathsf P}}\to\mathsf P(M_A)[1]\to0$; it is a quasi-isomorphism by the five lemma on the two
long exact sequences. (b) The first isomorphism is (a); the second is Shapiro's lemma for $Q=\Ind_A^{\Ga}Q_{A/B}$ \STD{\cite[Ch.~III.6.2]{brown}}. Independence of
$\mathsf P$ is (a) for a comparison. The long sequence of the short exact sequence of complexes in (a) is that of $0\to M_A\to M_B\to Q\to0$ in $H^{\bullet}(\Ga;-)$, whose
connecting map is that of the cone (Lemma~\ref{lem:levelconeexact}(1)); Shapiro's isomorphism is natural, and $\Ind_A^{\Ga}$ is exact and carries
$0\to V(\chi)\to\Ind_B^{A}V(\chi')\to Q_{A/B}\to0$ to $0\to M_A\to M_B\to Q\to0$ (Proposition~\ref{prop:levelconeshapiro}(1)). Localisation is flat. $24\cdot H^{j}(\Ga;Q)=0$ for
$j\ge2$ is Proposition~\ref{prop:certifiedeell}(4)(a) for $M=Q$, and the periodicity is Proposition~\ref{prop:descentN}(5) for $M=Q$ (natural in $M$, so it is compatible with
Shapiro's isomorphism). Over a $\Z[1/6]$-algebra $H^{q}(\Ga;-)=0$ for $q\ge2$, which leaves six terms. The examples are rows of the check below.
(c) $\Cone{\alpha_{\Iso}}$ is built from the finitely generated free $\Z$-modules $\CI(M_A)$, $\CI(M_B)$. Proposition~\ref{prop:cohglue} for it, with
$\Cone{\alpha_{\Iso}}[1/6]=\Cone{\alpha_{\Iso\otimes\Z[1/6]}}$ and (a) for the comparison $\cmpI$ over $\Z[1/6]$ (where $C_{\bullet}(\Xbulk)\otimes\Z[1/6]$ is a model) and for
any comparison $\Iso\otimes\Z_{(\ell)}\to\mathsf P$ at $\ell\in\{2,3\}$ ($H^{q}(\Cone{}\otimes\Zsix)=\bigoplus_{\ell\in\{2,3\}}H^{q}(\Cone{}\otimes\Z_{(\ell)})\otimes\Z_\ell$). The
maps: (a) and (1)(a).
(d) The proofs of Proposition~\ref{prop:certifiede}(1)--(3) and Remark~\ref{rmk:certifiedeperdiff} use only that $C$ is a complex of finitely generated free modules over
$\Z$, and Proposition~\ref{prop:certifiedeell}(3), (6) carry them to $\Z_{(\ell)}$; $\Cone{\alpha_{\mathsf P}}^{q}$ is a sum of two such modules. (i): by
Proposition~\ref{prop:certifiede}(2), $\rank_{\F_{\ell'}}\Dcone^{q}=\rank_{\Q}\Dcone^{q}$ iff no elementary divisor of $\Dcone^{q}$ is divisible by $\ell'$, iff
$H^{q+1}(\Cone{\alpha_{\mathsf P}})=H^{q+1}(\Ga;Q)$ has no $\ell'$-torsion; for $q+1\ge2$ the group is killed by $24$, for $q+1=1$ the condition is $s_{\ell'}(Q)=0$ by
Proposition~\ref{prop:badprimesN}(1), and for $q+1=0$ it is a subgroup of the free group $Q$. (ii): rank--nullity over $\Q$ in each degree, with
$\dim_{\Q}H^{j}(\Cone{}\otimes\Q)=b_j(Q)$; over $\Q$ the idempotent $e$ splits $M_{B,\Q}=\alpha M_{A,\Q}\oplus\ker\omega_{\Q}$ $\Ga$-equivariantly and $\ker\omega_{\Q}\cong Q_{\Q}$,
so $b_j(Q)=b_j(M_B)-b_j(M_A)$, which is $0$ for $j\ge2$ (Proposition~\ref{prop:certifiedeell}(4)(a)) and for $j=-1$. The last sentence is (a).
(e) Proposition~\ref{prop:badprimesN}(1) for the finitely generated free module $Q$, with $b_0(Q)=\beta_0$ by (d). If $\ell\nmid[A:B]$, $e$ splits the fibre sequence over
$\Z_{(\ell)}$ and over $\F_\ell$, so $H^{0}(\Ga;M_B/\ell)=H^{0}(\Ga;M_A/\ell)\oplus H^{0}(\Ga;Q/\ell)$; subtract $b_0(M_B)=b_0(M_A)+\beta_0$.
(f) At $\ell\ge5$ the stabilisers have orders dividing $12$, units in $\Z_{(\ell)}$: $C_{\bullet}(\mathcal T)\otimes\Z_{(\ell)}$ is a model by Lemma~\ref{lem:treeres}, and
$C_{\bullet}(\Xbulk)\otimes\Z_{(\ell)}$ by Remark~\ref{rmk:orbifold}; (a) applies. Exactness of $\mathrm T(M_B)\to\mathrm T(Q)\to0$ in degree $0$ is (a), i.e.\
$(M_B)^{F}\to Q^{F}$ is onto over $\Z/\ell^{K}$ for $F\in\{\langle S\rangle,\langle U\rangle\}$, so images of generators generate; the ranks are additive by the same
exactness. Proposition~\ref{prop:descentN}(7)(b) is stated for every finitely generated free $M$, here $M=Q$, with $b_0(Q)=\beta_0$. At $\ell=3$: by
Proposition~\ref{prop:descentNthree}(1) every finite subgroup of $A=\Ga_0(N)$ has order $1$, $2$ or $4$, so for every $\Z_{(3)}$-algebra $S$,
$C_{\bullet}(\mathcal T)\otimes S$ and $C_{\bullet}(\Xbulk)\otimes S$ are projective resolutions of $S$ over $S[A]$ (the proof of Proposition~\ref{prop:descentNthree}(2)),
as is $\Iso\otimes S$, and $\cmpI\otimes S$ is a comparison over $S[A]$. Frobenius reciprocity identifies $\mathrm T(\Ind_A^{\Ga}X_S)$ with
$\Hom_{S[A]}(C_{\bullet}(\mathcal T)\otimes S,X_S)$, naturally in the $A$-module $X$; so (a) and (b) hold verbatim with $S[A]$-models in place of $S[\Ga]$-models and the
$A$-sequence $0\to V\to\Ind_B^{A}V\to Q_{A/B}\to0$ in place of the fibre sequence (the middle term is $\Hom_{S[B]}(C_{\bullet}(\mathcal T)\otimes S,V_S)=\mathrm T(M_B)$ by
transitivity of induction, and $B\subseteq A$ has no element of order~$3$ either). $\mathrm T$ has length $1$, so $H^{j}(A;Q_{A/B}\otimes\Z_{(3)})=0$ for $j\ge2$, i.e.\
$H^{q}(A,B)\otimes\Z_{(3)}=0$ for $q\ge3$. The generators: $C_0(\mathcal T)\otimes S$ restricted to $S[A]$ is projective, so $\mathrm T^{0}(M_B)\to\mathrm T^{0}(Q)$ is onto,
and the images of the generators of Proposition~\ref{prop:descentNthree}(3)(a), (b) for $M_B$ generate $Q^{\langle S\rangle}+Q^{\langle U\rangle}$; the proof of
Proposition~\ref{prop:descentNthree}(3) applies to $M=Q$ with (a) here in place of its~(2), and $\rank\mathrm T^{0}$ is additive by the same surjectivity and the
injectivity of $\alpha$.
(b), signs. A cocycle $(c,c')$ of $\Cone{\alpha_{\mathsf P}}^{q-1}$ has $dc'=-\alpha_{\mathsf P}c$. $\delta$ lifts $\pi_{\mathsf P}(c,c')=[c']$ to $c'$, and $dc'=\alpha_{\mathsf P}(-c)$, so
$\delta[c']=[-c]$.
(g) (v1.23.0-pre.22; the D2 code audit's Appendix~B, made integral over $\Z_{(\ell)}$.) (i) Lift $L$, $R$ to $L_0$, $R_0$ over $\Z_{(\ell)}$; their determinants are units mod
$\ell$, so they are invertible over the local ring $\Z_{(\ell)}$. Then $L_0AR_0=\Sigma_0+\ell^{K}E$ over $\Z_{(\ell)}$, $\Sigma_0$ the lift of $\Sigma$ with entries
$u_i\ell^{\sigma_i}$, $u_i\in\Z_{(\ell)}^{\times}$ lifting the units (all that follows is verbatim with the units; v1.23.0-pre.23, referee m4); in blocks of sizes $s$ and $r-s$, resp.\ $s$ and $c-s$,
\[
  L_0AR_0=\begin{pmatrix}P&\ell^{K}E_{12}\\\ell^{K}E_{21}&\ell^{K}E_{22}\end{pmatrix},\qquad P=D+\ell^{K}E_{11}=DU,\quad D=\mathrm{diag}(u_i\ell^{\sigma_i}),\quad
  U=1+D^{-1}\ell^{K}E_{11}.
\]
$D^{-1}\ell^{K}$ has entries $u_i^{-1}\ell^{K-\sigma_i}$, all divisible by $\ell^{K'}$, $K'\ge1$; so $U\equiv1\pmod{\ell^{K'}}$, $U\in\GL_s(\Z_{(\ell)})$, and $P^{-1}=U^{-1}D^{-1}$. The
matrices $X:=-P^{-1}\ell^{K}E_{12}=-U^{-1}(\ell^{K}D^{-1})E_{12}$ and $Y:=-\ell^{K}E_{21}P^{-1}=-E_{21}U^{-1}(\ell^{K}D^{-1})$ ($\ell^{K}$ is a scalar) are integral over
$\Z_{(\ell)}$ and $\equiv0\pmod{\ell^{K'}}$. Then
\[
  \begin{pmatrix}1&0\\Y&1\end{pmatrix}L_0AR_0\begin{pmatrix}1&X\\0&1\end{pmatrix}=\begin{pmatrix}P&0\\0&\Sigma'\end{pmatrix},\qquad
  \Sigma'=\ell^{K}\bigl(E_{22}-\ell^{K}E_{21}P^{-1}E_{12}\bigr).
\]
The left side has $\rank_{\Q}=\rank_{\Q}A=s=\rank_{\Q}P$, so $\Sigma'=0$ --- \textbf{this is where the certification enters}. Right-multiplying by $\mathrm{diag}(U^{-1},1)$ turns $P$
into $D$. So $\tilde L:=\bigl(\begin{smallmatrix}1&0\\Y&1\end{smallmatrix}\bigr)L_0$ and $\tilde R:=R_0\bigl(\begin{smallmatrix}1&X\\0&1\end{smallmatrix}\bigr)\mathrm{diag}(U^{-1},1)$
satisfy $\tilde LA\tilde R=\mathrm{diag}(D,0)$; each factor besides $L_0$, $R_0$ is $\equiv1\pmod{\ell^{K'}}$, hence so is its inverse, which gives the four congruences. \emph{Moreover} (v1.23.0-pre.23, referee M1): $\tilde L=\bigl(\begin{smallmatrix}1&0\\Y&1\end{smallmatrix}\bigr)L_0$
leaves the first $s$ rows of $L_0$ unchanged; $\tilde L^{-1}=L_0^{-1}\bigl(\begin{smallmatrix}1&0\\-Y&1\end{smallmatrix}\bigr)$ leaves the last $r-s$ columns of $L_0^{-1}$
unchanged; $\tilde R^{-1}=\mathrm{diag}(U,1)\bigl(\begin{smallmatrix}1&-X\\0&1\end{smallmatrix}\bigr)R_0^{-1}$ leaves the last $c-s$ rows of $R_0^{-1}$ unchanged; and
$L_0$, $L_0^{-1}$, $R_0^{-1}$ reduce to $L$, $L^{-1}$, $R^{-1}$ modulo $\ell^{K}$. (ii)
For an exact Smith form $\tilde LA\tilde R=\mathrm{diag}(D,0)$ over the principal ideal domain $\Z_{(\ell)}$, $\ker A=\tilde R(0\oplus\Z_{(\ell)}^{c-s})$, a direct summand, and
$y\mapsto\tilde Ly$ identifies $\coker A$ with $\bigoplus_{i\le s}\Z_{(\ell)}/\ell^{\sigma_i}\oplus\Z_{(\ell)}^{r-s}$, with inverse given by the columns of $\tilde L^{-1}$; reduce by (i).
A torsion coordinate $(\tilde Ly)_i$, $i\le s$, is a row of the first block of $\tilde L$, hence $\equiv(Ly)_i\pmod{\ell^{K}}$ by the ``moreover'' of (i), and
$\sigma_i<K$; the generators (columns of $\tilde L^{-1}$, first block) and the free coordinates (last rows of $\tilde L$) are exact modulo $\ell^{K'}$ only, which is
why $2\sigma_{\max}\le K$ is sufficient (all torsion exponents $\le K'$) and not necessary (v1.23.0-pre.23, referee m1). Lattice coordinates: for $y=Ax$,
$\tilde Ly=\mathrm{diag}(D,0)\tilde R^{-1}x$, so $A\Z_{(\ell)}^{c}=\tilde L^{-1}(D\Z_{(\ell)}^{s}\oplus0)$ with basis $(\tilde L^{-1}De_i)_{i\le s}$, and the coordinates of $y$
are $D^{-1}(\tilde Ly)_{\le s}=D^{-1}(L_0y)_{\le s}$; $(L_0y)_{\le s}$ is known modulo $\ell^{K}$ and its $i$-th entry is divisible by $\ell^{\sigma_i}$, so the $i$-th
coordinate is known modulo $\ell^{K-\sigma_i}$. The lower block of $\tilde Ly$ vanishes: $Y(L_0y)_{\le s}+(L_0y)_{>s}=0$, so
$(L_0y)_{>s}=-Y(L_0y)_{\le s}=\ell^{K}E_{21}U^{-1}\cdot D^{-1}(L_0y)_{\le s}\equiv0\pmod{\ell^{K}}$, and $(L_0y)_{>s}\equiv(Ly)_{>s}$. (iii) Sums and products commute with reduction mod $\ell^{K'}$; a matrix
known modulo $\ell^{K'}$ that is the reduction of a $\Z_{(\ell)}$-matrix $A'$ is all (i) needs, with $K'$ in place of $K$ and $\rank_{\Q}A'$ in place of $\rank_{\Q}A$. The other operations: if $x\equiv x_0\pmod{\ell^{k}}$ and $\ell^{\sigma}\mid x,x_0$, then
$x/\ell^{\sigma}\equiv x_0/\ell^{\sigma}\pmod{\ell^{k-\sigma}}$; a torsion row is an element of $\Z/\ell^{a_i}$, so its reduction loses nothing; a formula in data known
modulo $\ell^{k_1}$ and $\ell^{k_2}$ is known modulo $\ell^{\min(k_1,k_2)}$. The certifying ranks: rationally $\cor\circ\res=[A:B]$ ((1)(d)) makes $\res$ injective, so
the rational long sequence breaks into the short exact sequences stated, which give the ranks of $\res$, $j$, the connecting map ($0$), $\cor$ (surjective, as
$\cor\circ\res$ is invertible) and $c_\sigma$ (an automorphism); for a normal pair $H^{1}(B)^{\mathsf G}_{\Q}=\res H^{1}(A)_{\Q}$ ($\res\circ\cor=\sum_\sigma c_\sigma$).
The columns $\ell^{a_i}e_i$ of $\Delta$ span the torsion coordinates rationally, and on the free coordinates $F$ is $f\otimes\Q$; this gives $t_Y+\rank_{\Q}f$ for
$[F\mid\Delta_Y]$, $g_X-\rank_{\Q}f$ for $L_f$ (rank--nullity), and $t_X$, resp.\ $t_X+\rank_{\Q}g$, for the sublattices $I$. Kernels,
images and cokernels of an induced map, and a subgroup comparison, are lattice computations (kernels and quotients of the matrices so formed), each a certified Smith form
read under (ii). The last sentence: $A=(\ell^{K})$, $L=R=1$, $s=0<1=\rank_{\Q}A$; then $\Sigma=0$ over $\Z/\ell^{K}$, but every exact Smith form of $A$ is $(\ell^{K})$.

(3)(a) Put $A'':=A\cap gAg^{-1}$, $A':=A\cap g^{-1}Ag$, $B''$, $B'$ likewise, $\iota_g\colon A''\to A'$, $h\mapsto g^{-1}hg$, so that
$[AgA]=\cor^{A}_{A''}\circ(\iota_g,f)^{\ast}\circ\res^{A}_{A'}$ (Proposition~\ref{prop:descentN}(5)). From $g\smat abcd g^{-1}=\smat{a}{b/p}{pc}{d}$ and
$g^{-1}\smat abcd g=\smat{a}{pb}{c/p}{d}$: $A''=\Ga_0(pN)$, $A'=\Ga_0(N)\cap\Ga^{0}(p)$, $B''=\Ga_0(pMN)$, $B'=\Ga_0(MN)\cap\Ga^{0}(p)$. If $p\nmid M$, then
$B\cap A''=\Ga_0(\operatorname{lcm}(MN,pN))=B''$ and $B\cap A'=B'$, and $[A:A'']=[A:A']=[B:B'']=[B:B']$ ($=p+1$ if $p\nmid N$, $=p$ if $p\mid N$); so
$[A''B:A'']=[B:B\cap A'']=[A:A'']$, i.e.\ $A=A''B$, and likewise $A=A'B$. The double coset formula \STD{\cite[Ch.~III.9]{brown}, Mackey} with a single double coset
gives $\res^{A}_{B}\cor^{A}_{A''}=\cor^{B}_{B''}\res^{A''}_{B''}$ and $\res^{A}_{A'}\cor^{A}_{B}=\cor^{A'}_{B'}\res^{B}_{B'}$. With the naturality of $\res$ and $\cor$ under the
isomorphism of pairs $(\iota_g,f)$, which maps $B''$ onto $B'$:
$\res^{A}_{B}[AgA]=\cor^{B}_{B''}(\iota_g,f)^{\ast}\res^{A'}_{B'}\res^{A}_{A'}=\cor^{B}_{B''}(\iota_g,f)^{\ast}\res^{B}_{B'}\res^{A}_{B}=[BgB]\res^{A}_{B}$, and
$[AgA]\cor^{A}_{B}=\cor^{A}_{A''}(\iota_g,f)^{\ast}\cor^{A'}_{B'}\res^{B}_{B'}=\cor^{A}_{A''}\cor^{A''}_{B''}(\iota_g,f)^{\ast}\res^{B}_{B'}=\cor^{A}_{B}[BgB]$. The
character is harmless: $a_{g^{-1}hg}=a_h$, so $f=\rho_{n,\chi}(g)$ is compatible on every subgroup, and $\chi'|_B=\chi|_B$. Over $\Z$: (1)(c). The cochain statement is
the criterion Proposition~\ref{prop:levelconehecke}(1), whose hypotheses $B\cap A''=B''$ and $A=A''B$ were just verified. (b) $H^{1}(\cdot;\Z)\otimes\F_7\to H^{1}(\cdot;\F_7)$ is an isomorphism, natural in the group and compatible with the operators, at
both levels: its cokernel $\mathrm{Tor}(H^{2},\F_7)$ vanishes, since $24$ kills $H^{2}$. A square commuting over $\Z$ would commute over $\F_7$. The classical $U_p$: Proposition~\ref{prop:levelconehecke}(3)(d), whose hypotheses $p\mid N'_A=N$ and
$pN'_A\mid N'_B=MN$ are $p\mid N$, $p\mid M$. The cell $(\Ga_0(2)\subset\Ga,n=2,p=2)$ is a row of the referee's check (\file{C-D3\_referee.md} \S2.2).
(b$'$) Let $p\mid N$ and $M$ arbitrary; $A'=A\cap g^{-1}Ag=\Ga_0(N)\cap\Ga^{0}(p)$, $B'=B\cap g^{-1}Bg=\Ga_0(MN)\cap\Ga^{0}(p)$ and $B''=B\cap gBg^{-1}$ as in (a), so
$B\cap A'=\Ga_0(MN)\cap\Ga^{0}(p)=B'$ for every $M$. Let $\gamma=\smat abcd\in A$. As $p\mid N\mid c$, $a$ is a unit mod $p$; for $j\equiv ba^{-1}\pmod p$,
$\gamma T^{-j}=\smat{a}{b-ja}{c}{d-jc}$ ($T=\smat1101$) lies in $\Ga_0(N)\cap\Ga^{0}(p)=A'$, so $\gamma\in A'T^{j}\subseteq A'B$: $A=A'B$, whatever $M$ is. The double coset formula
\STD{\cite[Ch.~III.9]{brown}, Mackey} with the single double coset $A'\cdot1\cdot B$ gives $\res^{A}_{A'}\cor^{A}_{B}=\cor^{A'}_{A'\cap B}\res^{B}_{A'\cap B}=
\cor^{A'}_{B'}\res^{B}_{B'}$. The isomorphism of pairs $(\iota_g,f)$, $\iota_g\colon A''\to A'$, carries $B''$ onto $g^{-1}B''g=g^{-1}Bg\cap B=B'$ (for every $M$),
and $\cor$ is natural under it: $(\iota_g,f)^{\ast}\cor^{A'}_{B'}=\cor^{A''}_{B''}(\iota_g,f)^{\ast}$. Since $B''\subseteq A''\cap B$, transitivity gives
$\cor^{A}_{A''}\cor^{A''}_{B''}=\cor^{A}_{B''}=\cor^{A}_{B}\cor^{B}_{B''}$. Hence
$[AgA]\cor^{A}_{B}=\cor^{A}_{A''}(\iota_g,f)^{\ast}\res^{A}_{A'}\cor^{A}_{B}=\cor^{A}_{A''}(\iota_g,f)^{\ast}\cor^{A'}_{B'}\res^{B}_{B'}
=\cor^{A}_{A''}\cor^{A''}_{B''}(\iota_g,f)^{\ast}\res^{B}_{B'}=\cor^{A}_{B}\cor^{B}_{B''}(\iota_g,f)^{\ast}\res^{B}_{B'}=\cor^{A}_{B}[BgB]$, the last step
Proposition~\ref{prop:descentN}(5) at level $MN$. The character is harmless as in (a). Nothing used $p\nmid M$: the condition $B\cap A''=B''$ of (a), which fails for
$p\mid M$ ($B\subseteq A''$), enters only the square with $\res$. Over $\Z$: (1)(c). If $p\nmid N$, $p\mid M$, then $[A:A']=p+1$ while $[B:B\cap A']=[B:B']=p$, so
$A'B\ne A$, the double coset formula has a second term, and the referee's cell shows it is non-zero. (c) The first sentence is
Proposition~\ref{prop:levelconehecke}(1) with Proposition~\ref{prop:levelconeshapiro}(2); the transport: $\Psi^{\ast}T^{B}\Phi^{\ast}\alpha=\Psi^{\ast}T^{B}\alpha\Phi^{\ast}=
\Psi^{\ast}\alpha T^{A}\Phi^{\ast}=\alpha\Psi^{\ast}T^{A}\Phi^{\ast}$ by (1)(a) twice; a cochain endomorphism of the cone commuting with its short exact sequence makes the long
sequence equivariant.

(4)(a) is (1)(d) and (1)(c); $(-)^{\mathsf G}$ is a kernel, and $\Z[1/6]$, $\Z_{(\ell)}$, $\Zsix$ are flat. (b) The citation; over a $\Z[1/6]$-algebra
$H^{2}(A)=0$. (c) Torsion-free $B$: every finite subgroup has order $1$, so Proposition~\ref{prop:modeltorsion} gives $H^{b}(B;X)=0$ for $b\ge2$ over every ring. With two
rows $d_r=0$ for $r\ge3$, and $H^{m}(A)$ has the filtration $0\to E_\infty^{m,0}\to H^{m}(A)\to E_\infty^{m-1,1}\to0$ with $E_\infty^{m,0}=\coker(d_2\colon
E_2^{m-2,1}\to E_2^{m,0})$ and $E_\infty^{m-1,1}=\ker(d_2\colon E_2^{m-1,1}\to E_2^{m+1,0})$; splicing gives the sequence, the first map being the edge map $\inf$ and, in
degree $1$, the second the restriction to $E_2^{0,1}=H^{1}(B)^{\mathsf G}$. $H^{0}(B)=V(\chi')_S^{B}=0$: $B$ has finite index, so it contains $T^{k}$ and $\smat10k1$ for
some $k\ge1$, both with $a=1$, so they act untwisted; over the $\Q$-algebra $S\otimes\Q$ the $T^{k}$-invariants of $\Sym^{n}$ ($x\mapsto x$, $y\mapsto y+kx$) are
$S_{\Q}x^{n}$, and $\smat10k1$ sends $x^{n}$ to $(x+ky)^{n}$, whose coefficient of $x^{n-1}y$ is $nk\in(S\otimes\Q)^{\times}$; so the invariants vanish over $S\otimes\Q$, and
over $S\subseteq S\otimes\Q$. Then $E_2^{\bullet,0}=0$, and the sequence gives $H^{q}(A)\cong E_2^{q-1,1}$. $\abs{\mathsf G}$ kills $H^{\ge1}(\mathsf G;-)$
\STD{\cite[Ch.~III.10.2]{brown}}, which gives the next two consequences, and the periodic resolution the last. $\Ga_1(N)\trianglelefteq\Ga_0(N)$ is the kernel of
$\gamma\mapsto d_\gamma\bmod N$, onto $(\Z/N)^{\times}$; $\abs{(\Z/p^{r})^{\times}}=p^{r-1}(p-1)$. If $A$ is torsion-free too, Proposition~\ref{prop:modeltorsion} applies to both
ends at every $\ell$, and the proofs of Lemma~\ref{lem:levelconeexact} and Proposition~\ref{prop:levelinfres}(2)--(5) use $6\in\cring^{\times}$ only through that and through
$H^{\ge2}(A)=H^{\ge2}(B)=0$.

(5) Items (1), (2) and (4) use of a model only that it is an $S$-model and that comparisons exist ((T1); they exist by the comparison theorem, and (T3) makes them
explicit), and (2)(d) uses that its cochains are finitely generated free over $\Z_{(\ell)}$: under (T2), $M_S^{Q_x}=e_{Q_x}M_S$ is a direct summand of a finitely generated
free module, free over the local principal ideal domain $\Z_{(\ell)}$, and (1)(b) applies. (3)(c)'s transport is its last sentence.

\emph{A check outside the build, not a tag} (v1.23.0-pre.18): \file{docs/implementation/slot24/cd3\_descentmaps\_check.py}, independent of \code{cohomology\_app} and of
\file{descentN\_oracle.py} (own coset tables, induced modules, the free tree resolution of $\Z$ over $\Z[\Ga]$ in every degree, exact integer Smith forms; the oracle table is
read as data only), output \file{cd3\_descentmaps\_check.out} (its header records the command). On $74$ cells
$\Ga_0(MN)\subseteq\Ga_0(N)$ ($16$ pairs, $N\le6$, $MN\le16$, even $n\le8$, $\psi(MN)(n+1)\le140$; $29$ with $\nu_3(N)=0$, $21$ normal): the fibre identities of~(1)(d) (and of (4)(a) on
the normal pairs), $H^{j}(\Cone{\alpha})=H^{j}(\Ga;Q)$ as groups for $-1\le j\le4$ ((2)(a)), the Betti recursion of (2)(d)(ii) for every $\Dcone^{q}$, the bad-prime count of
(2)(e) at every $\ell\ge5$ that occurs and at $5,7,11,13$, the sparse cokernel of (2)(f) against $H^{1}(\Ga;Q)$ at every $\ell\ge5$ and at $\ell=3$ when $\nu_3(N)=0$, and no
$3$-torsion in $H^{\ge3}(A,B)$ when $\nu_3(N)=0$; the ends against the oracle on $888$ rows. On $12$ cells $\Ga_1(N)\trianglelefteq\Ga_0(N)$ ($N\in\{4,5,7,9\}$, $1\le n\le4$),
$H^{q}(\Ga_0(N);\Sym^{n}\Z^{2})=H^{q-1}((\Z/N)^{\times};H^{1}(\Ga_1(N);\Sym^{n}\Z^{2}))$ for $q=1,\dots,5$ ((4)(c); $H^{1}(\Ga_1(N))$ on the tree complex of the torsion-free
$\Ga_1(N)$, checked against the free resolution), e.g.\ $H^{3}(\Ga_0(7);\Sym^{4}\Z^{2})\supseteq(\Z/3)^{2}$, while $3\nmid\abs{(\Z/5)^{\times}}$ leaves no $3$-torsion in
$H^{\ge2}(\Ga_0(5);\Sym^{n}\Z^{2})$. $0$ mismatches (\code{CD3\_CHECK \dots\ mismatch=0 \dots\ verdict=OK}). Item~(3) is not covered by that script; the
referee's independent script (\file{docs/implementation/slot24/C-D3\_referee.md} \S2.2: $H^{1}$ over $\Z$ as crossed homomorphisms, $[AgA]$ through Shapiro, torsion
included) checks both squares on $32$ cells, and at v1.23.0-pre.19 on $10$ more: the square with $\cor$ commutes on all $15$ cells with $p\mid N$, $p\mid M$ (e.g.\
$\Ga_0(8)\subset\Ga_0(2)$ and $\Ga_0(8)\subset\Ga_0(4)$ at $p=2$, $n\le6$; $\Ga_0(18)\subset\Ga_0(6)$, $\Ga_0(9)\subset\Ga_0(3)$ at $p=3$; $\Ga_0(25)\subset\Ga_0(5)$), where the
square with $\res$ fails on every one, and fails on all $5$ cells with $p\nmid N$, $p\mid M$, $n>0$ ((b), (b$'$)).
\end{proof}

\noindent\emph{What is deferred, and why} (v1.23.0-pre.18). (i) The \textbf{second degeneracy map} $\mathrm{dg}_p=[A\,g\,B]$ and every Hecke-type map: they are not fibre maps
(conjugation by $g\notin\Ga$ is not an automorphism of $\Ga$), so (1) does not reach them; on cohomology they are determined locally (Proposition~\ref{prop:descentN}(3)),
and their cochain lifts on $\Iso$ or $\Iso^{(\ell)}$ at level~$N$ are C-D4's, as is the Hecke action on the relative term at $\ell\in\{2,3\}$ ((3)(c)). (ii) The
\textbf{Hochschild--Serre spectral sequence over $\Z_{(2)}$, $\Z_{(3)}$ when $B$ has torsion}, beyond the five-term sequence (4)(b): its higher differentials are neither
computed nor needed, because $H^{q}(A)$ is computed directly; Proposition~\ref{prop:levelinfres}(5) (coinvariants) is not asserted there. (iii) The \textbf{old-space cone}
(Remark~\ref{rmk:oldspacecone}) stays deferred: its $U_p$-action needs an $\etapair$-type certificate. (iv) The \textbf{boundary pair and its cone over $\Z$} are the plan's D4,
not this clause. (v) Integral nebentypus modules $M_N(\chi)$ over $\mathsf O_\chi$ are not built (Proposition~\ref{prop:descentN}(8)); (1)--(2) and (4) are stated for them,
(2)(f) at $\ell=3$ and (3)(b)'s example for $\chi=1$ only. (vi) The tree models (5) are the plan's D5, gated by G-T.

\paragraph{Code map of Propositions~\ref{prop:descentN}, \ref{prop:badprimesN}, \ref{prop:descentNthree} and~\ref{prop:descentmaps}.} \textbf{Stated at v1.23.0-pre.14, before D1's code} (contract first,
rule~21). \textbf{Built} (and used by D0 at level~$1$; \file{Integral/Descent/DescentAssembly}, \file{docs/implementation/slot24/DESCENT0\_report.md}): the
$\Zsix$-corner on $\Iso^{(2)}$, $\Iso^{(3)}$ for any \code{CoefficientModule}, under Proposition~\ref{prop:certifiedeell}; the Smith-valuation hand-off
\code{smith\_valuations}; the shipped $C_7$ $K$-mode (\code{EquivariantCohomologyOnK} over a composite context, linear $K$ search) at level~$1$; the tree Bredon
complex $\mathrm T(M)$ over $\Z/\ell^{K}$ and $\F_{\ell'}$ as a \emph{measured comparator} (D0's (c$'$), feeding no result), and \code{RunLevelNProbe} (the
level-$N$ probe, \code{MonomialOf}, which checks $b_0$ against~(7)(c)'s closed form at the good prime); \file{Integral/Model/PeriodicityClass}'s chain map $E$ on
$\Iso_{\bullet}$, which is natural in the coefficient module. \textbf{Not built}: (7)(a), the direct Smith form of $d^{0}$ of $\Ceq(M_{\Z/\ell^{K}})$ as a production
route at any level; (7)(b), the sparse cokernel over $\Z/\ell^{K}$ --- \file{Hecke/HeckeLevelNSparse}'s \code{SparseCokernel} eliminates over $\F_p$ by RREF and is gated to a prime field
(\code{ctx.is\_field()}, \code{InvP(order, p)}, the constructor's \code{IsPrimeU64(p)}); its generators (\code{OrbitWiseFixedGenerators}: transported blocks on
free orbits, the columns of the averaging idempotent $\tfrac1{\abs{F_j}}\sum_{h\in F_j}\rhon(h)$ at fixed points, orbit sums at $n=0$) are (7)(b)-admissible over
$\Z/\ell^{K}$ once $1/\abs{F_j}$ is taken mod $\ell^{K}$; what is not built is the elimination over $\Z/\ell^{K}$ (minimal-valuation pivoting, the Smith read-out)
(v1.23.0-pre.15, referee F3: pre.14 described those generators as kernel bases over $\F_p$, which they are not); the level-$N$ descent assembly (D1's
\code{IntegralDescentLevelNValidation}, with its arm~(B), Proposition~\ref{prop:badprimesN}'s list against $\F_\ell$ runs at every prime $\le n+2$ including an
$\ell=N=n+1$ cell); the matrix of $E$ on $\CI(M_N)$; (8)'s integral form $M_N(\chi)$ (the build carries characters over $\F_q$ and $W_K(\F_q)$ only); every level-$N$
Hecke operator over $\Z$ (C-D4). \textbf{A check outside the build, not a tag}: a script in exact modular arithmetic (\file{docs/implementation/slot24/cd2\_badprimes\_check.py})
compares $s_\ell$ with $\rank_{\Q}-\rank_{\F_\ell}$ ($\rank_{\Q}$ taken as the common rank at two primes $>10^{6}$) of the cokernel presentation over $\Z$ on $280$ cells ($N\in\{1,2,3,4,5,6,7,10,11,13\}$, even $n\le12$,
$\ell\in\{5,7,11,13\}$) and, for the Legendre characters modulo $5,7,11,13$, on $112$ cells, and the twisted Borel count of~(3)(b) with direct fixed spaces on
$984$ cells: $0$ mismatches.
\textbf{Proposition~\ref{prop:descentNthree}} (v1.23.0-pre.17, stated \emph{after} D1's code, finding M1 of \file{docs/implementation/slot24/D1\_code\_audit.md};
on branch \texttt{s24-d1impl}, \texttt{726c3a67} and \texttt{87e8cf32}, not merged at pre.17). \file{Integral/Descent/DescentLevelN}: \code{RunSparseCokernelSmith}$(N,n,3)$
is~(3) --- \code{FixedGenerators} builds (3)(a)'s orbit sums ($S$ with $o=2$, $ST$ with $o=3$, no division) and (3)(b)'s idempotent columns
$\tfrac12(1+B)$ at the $S$-fixed points (an \code{ENGINE\_CHECK} that $\gcd(o,3^{K})=1$ at a fixed point, and that the fixed-point counts are $\nu_2$, $\nu_3$), the
elimination is sparse with unit pivots and the residual goes to \code{SmithLevelsEll}; $K=$ \code{ring3k::kMaxK} $=19$ at once (no $K$ search: $s<r$ is a
refusal, the stopping rule at the cap), $r=\code{rank\_T0}-b_0$ from \code{LevelNClosedForm} (the closed forms of~(3)); \code{AssembleDescentN} takes it
when $\nu_3=0$ and puts nothing $3$-primary in $H^{\ge2}$, as~(2) says, and runs $\Iso^{(3)}$ otherwise (item~(5)); \code{force\_dense\_three} is the
comparison arm~(S) of \code{IntegralDescentLevelNValidation}, at $(N,n)=(5,4)$, $(11,6)$, $(2,12)$, $(35,4)$, $(8,6)$. $n$ odd: only the $2$-adic piece runs, item~(4). (D1's code was merged into \texttt{hecke-extensions} after v1.23.0-pre.17, merge \texttt{1cc9a6d3}.)
\textbf{Proposition~\ref{prop:descentmaps}} (v1.23.0-pre.18, \textbf{stated before D2's code}, contract first, rule~21; plan milestone D2). \textbf{Built} and used: the
fibre maps on the $7$-vertex model over $\F_\ell$, $\Z/\ell^{K}$ (\S\ref{sec:nebles}'s code), the $\Zsix$-corner on $\Iso^{(2)}$, $\Iso^{(3)}$ and the $\nu_3=0$ shortcut of
D1 for one module at a time. \textbf{Not built}: (1) $\alpha$, $\omega$, $R_\sigma$ on $\CI(-)$ and on $C^{\bullet}_{\Iso^{(\ell)}}(-)$ (block by block, (1)(b); a
\code{FibreMap} applied to the column layout of Proposition~\ref{prop:certifiedeell}(5)); (2) the level cone $\Cone{\alpha_{\mathsf P}}$ on $\Iso$, $\Iso^{(2)}$, $\Iso^{(3)}$ and on
$\mathrm T$ at $\ell\ge5$ and at $\ell=3$, $\nu_3(N)=0$, its certified-$K$ read-out with the Betti recursion (2)(d)(ii), the sparse cokernel of $Q$ over $\Z/\ell^{K}$ ((2)(f);
the same missing elimination as for (7)(b)), and the glued sequence \eqref{eq:descentmapsles} to $q_{\max}$; (3) the Hecke squares (3)(a) over $\Z$ need the level-$N$
operators of C-D4 (only their local components on the model over $\F_\ell$, $\ell\ge5$, exist, \S\ref{sec:nebles}); (4) the two-row read-out (4)(c) for
$\Ga_1(N)\trianglelefteq\Ga_0(N)$. The plan's validation \code{IntegralDescentLESValidation} (arms (E) exactness, (S) split case, (I) $\cor\circ\res=[A:B]$, (L) against
LES-1's predictions after $\otimes\Z[1/6]$, (2) the $2$-adic pieces at $\Ga_0(4)\trianglelefteq\Ga_0(2)$, $\Ga_0(9)\trianglelefteq\Ga_0(3)$) is \textbf{not built}; the
check \file{docs/implementation/slot24/cd3\_descentmaps\_check.py} (outside the build) is its oracle.
(v1.23.0-pre.22:) \textbf{Built at D2} (\file{Integral/Descent/DescentMaps}, \code{IntegralDescentMapsValidation} and its slow companion, the plan's
\code{IntegralDescentLESValidation} under another name): (1) block by block on $\Iso^{(2)}$, $\Iso^{(3)}$ and on $M$ for the cokernel route; (2) the cone on $\Iso^{(2)}$, on
$\Iso^{(3)}$ for $\nu_3(N)>0$, the cokernel routes (2)(f) at $\ell\ge5$ and at $\ell=3$, $\nu_3(N)=0$, the certified read-out (2)(d)(ii), the glued sequence to $q=4$ (degree~$6$ in
one arm), exactness as subgroups; every generator, coordinate and induced map read under (2)(g), its precisions tracked and refused when exhausted; (1)(d)'s identities and
(4)(a), (4)(b) in degree~$1$ on normal pairs. Still not built: (3), (4)(c), (5). Oracles: \file{d2\_maps\_oracle.py} and, independent of it, the D2 code audit's
\file{d2audit/d2audit\_oracle.py}; report \file{DESCENT2\_report.md}.

\paragraph{Formalization note.} Propositions~\ref{prop:descentN} and~\ref{prop:badprimesN} (v1.23.0-pre.14), \ref{prop:descentNthree} (v1.23.0-pre.17) and~\ref{prop:descentmaps} (v1.23.0-pre.18) are base change, Shapiro's lemma, the comparison theorem for projective resolutions, the semisimple
derived category of $\Qsix$, invariant theory of $B(\F_\ell)$ on $\F_\ell[x,y]$, the double coset formula and a two-row spectral sequence; they carry \textbf{no} \code{\textbackslash LEAN} marker.

% NEW at v1.23.0-pre.20 (hecke-extensions slot 24, milestone D5 of docs/implementation/slot24/DESCENT_plan.md v0.4, its clause C-D6; the owner's
% ruling at gate G-T of 2026-10-25: D5 forward, before D3/D4; the evidence DESCENT1_report.md).  Placed after the last statement of
% sec:descentN, before sec:mpmodel, so that no statement number moves.  Refereed ACCEPT-with-fixes (C-D6_referee.md); its fixes M1, m1-m5
% applied at v1.23.0-pre.21 (M1: prop:isomodeltree(4) made effective), not re-refereed.  No \LEAN marker.  Bulk only:
% the boundary, the pair and the cusp are clause C-D8 (D5c), not stated here.
\subsection{The tree models: the \texorpdfstring{$\ell$}{l}-local resolutions on the Bass--Serre tree, bulk only}\label{sec:isomodeltree}

(v1.23.0-pre.20; milestone D5 of \file{docs/implementation/slot24/DESCENT\_plan.md} v0.4, clause C-D6 of its \S2.7; the owner's ruling at gate G-T,
2026-10-25, on the evidence of \file{DESCENT1\_report.md}.) \S\ref{sec:isomodelell} replaces the columns of the free model $\Iso_{\bullet}$ on $\Xbulk$ by their
$\ell$-local forms; this subsection does the same on the Bass--Serre tree $\mathcal T$ of Claim~\ref{cl:amalgam}(1), whose free resolution is
Remark~\ref{rmk:treeconeres}. It fixes (i) \textbf{the $\ell$-local tree models} $\Iso^{(\ell)}_{\mathcal T}$ (Definition~\ref{def:isomodeltree}): the cone of a lift of the
tree's boundary between $\ell$-local columns, with closed-form lifts at $\ell=2$, $\ell=3$ and $\ell\ge5$; (ii) (Proposition~\ref{prop:isomodeltree}) what they compute,
in every degree and for every coefficient module, in particular $\Ind_{\Ga_0(N)}^{\Ga}$ of \S\ref{sec:descentN}; their integral form over $\mathsf O_\ell$ and the
certified-exponent rule on it; their cochains in closed form; closed-form comparisons with the free tree resolution and \emph{effective} comparisons with $\Iso$ over $\Z$,
i.e.\ over the augmentation, constructed through the inclusion of the tree in $\Xbulk$ and an explicit retraction back (v1.23.0-pre.21); the fibre maps and the level cone of Proposition~\ref{prop:descentmaps} on them; and the Hecke action by
transport. \textbf{Bulk only}: nothing here concerns $\Xbdry$, $\Isob_{\bullet}$, the pair $(\Ga,\Ga_\infty)$, its long exact sequence or its cone --- the tree cusp
pair is the plan's clause C-D8, not stated in this document. \emph{Notation.} $\mathbf c:=-I$ \emph{as an element of $\Ga$} (so $1-\mathbf c\ne0$ in $\Z[\Ga]$; it acts on
a module $M$ by $\rho_M(-I)$); $U=ST$ as in Claim~\ref{cl:amalgam}; $N_S=\sum_{j<4}S^{j}$, $N_U=\sum_{j<6}U^{j}$ as in Remark~\ref{rmk:treeconeres}; for $a\in\cring[\Ga]$,
$r_a$ is right multiplication $y\mapsto ya$, a left $\cring[\Ga]$-module map; for a subgroup $H$ of finite order invertible in $\cring$, $e_H:=\abs H^{-1}\sum_{h\in H}h$ and
$M^{H}=\rho_M(e_H)M$. The symbols $x\in\{S,U,\mathfrak e\}$, $\Ga_x$, $g_x$, $P_x$, $Q_x$, $e_x$, $N'_x$, $F^{x}$, $\varphi'$, $\Iso_{\mathcal T}$, $\Iso^{(\ell)}_{\mathcal T}$,
$C^{\bullet}_{\mathcal T}$, $C^{\bullet}_{\mathcal T,\ell}$, $\Phi^{\mathcal T}_\ell$, $\Psi^{\mathcal T}_\ell$, $\Phi^{\Z}_{\mathcal T}$, $\Psi^{\Z}_{\mathcal T}$, $\Phi_{\mathcal T}$, $\Psi_{\mathcal T}$,
$\mathsf k_j$, $\kappa_t$, $\kappa'_t$, $\mathsf w$, $T'_{\mathcal T}$ and (v1.23.0-pre.21) $\mathsf i_{\mathcal T}$, $\mathsf r_{\mathcal T}$, $\mathsf G_{\mathcal T}$, $\mathsf R_\sigma$, $\mathsf h$, $\mathsf h_p$,
(L), (L$'$), $\mathsf H_{\mathcal T}$, $\mathsf H_X$, $\mathsf H_\ell$, $\mathsf K$, $\mathsf K'$ are new symbols of this subsection and mean nothing else (rule~\ref{A5}); $P_x$, $Q_x$, $e_x$ are
those of Definition~\ref{def:isomodelell} for the three tree stabilisers, $\Phi_{\mathcal T}$, $\Psi_{\mathcal T}$ are the plan's names, and $\mathsf q_\ell$, $\mathsf O_\ell$
are Proposition~\ref{prop:certifiedeell}'s.

\begin{definition}[The free and the $\ell$-local tree models]\label{def:isomodeltree}
Let $x$ run over the two vertices and the edge of the tree, $x\in\{S,U,\mathfrak e\}$ for $\mathfrak v_S$, $\mathfrak v_U$, $\mathfrak e$ of Remark~\ref{rmk:treeconeres},
with stabilisers $\Ga_x=\langle g_x\rangle$, $g_S=S$, $g_U=U$, $g_{\mathfrak e}=\mathbf c$, of orders $4,6,2$. For a prime $\ell$ let $P_x$ be the $\ell$-Sylow subgroup of
$\Ga_x$, $Q_x$ its subgroup of order prime to $\ell$, and $e_x:=e_{Q_x}$; let $\cring$ be a commutative ring in which every $\abs{Q_x}$ is a unit, i.e.\ an
$\mathsf O_\ell$-algebra, $\mathsf O_\ell=\Z[1/\mathsf q_\ell]$, $\mathsf q_\ell:=\operatorname{lcm}_x\abs{Q_x}$:
\begin{center}
\begin{tabular}{@{}lllll@{}}
 & $Q_S$, $P_S$ & $Q_U$, $P_U$ & $Q_{\mathfrak e}$, $P_{\mathfrak e}$ & $\mathsf q_\ell$ \\ \hline
$\ell=2$ & $1$, $\langle S\rangle$ & $\langle U^{2}\rangle$, $\langle U^{3}\rangle=\langle\mathbf c\rangle$ & $1$, $\langle\mathbf c\rangle$ & $3$ \\
$\ell=3$ & $\langle S\rangle$, $1$ & $\langle\mathbf c\rangle$, $\langle U^{2}\rangle$ & $\langle\mathbf c\rangle$, $1$ & $4$ \\
$\ell\ge5$ & $\langle S\rangle$, $1$ & $\langle U\rangle$, $1$ & $\langle\mathbf c\rangle$, $1$ & $12$
\end{tabular}
\end{center}
(the values of $\mathsf q_\ell$ of Proposition~\ref{prop:certifiedeell}).
\begin{enumerate}
\item \textbf{(Columns.)} $F^{x}_t:=\cring[\Ga]e_x$ for every $t\ge0$ if $P_x\ne1$, and for $t=0$ only if $P_x=1$; its differential is
  $\partial^{x}_t:=r_{(1-g_x)e_x}$ for odd $t$ and $r_{N'_x}$ for even $t\ge2$, $N'_x:=\sum_{j<\abs{P_x}}g_x^{j}e_x$; its augmentation is the
  restriction to $F^{x}_0=\cring[\Ga]e_x$ of the projection $\cring[\Ga]\to\cring[\Ga/\Ga_x]$, $u\mapsto u\Ga_x$, so $ye_x\mapsto y\Ga_x$.
\item \textbf{(The model.)} An \textbf{$\ell$-local tree model} is the mapping cone $\Iso^{(\ell)}_{\mathcal T}$ of a $\cring[\Ga]$-chain map $\varphi'\colon F^{\mathfrak e}_{\bullet}\to
  F^{S}_{\bullet}\oplus F^{U}_{\bullet}$ that lifts the boundary of \eqref{eq:treeses} in the orientation of Remark~\ref{rmk:treeconeres}, $\gamma\langle\mathbf c\rangle\mapsto
  \gamma\langle S\rangle-\gamma\langle U\rangle$, through the augmentations:
  \[
    \Iso^{(\ell)}_{\mathcal T,q}=F^{S}_q\oplus F^{U}_q\oplus F^{\mathfrak e}_{q-1},\qquad \partial(a,z)=\bigl(\partial a+\varphi'_{q-1}z,\;-\partial^{\mathfrak e}z\bigr),
  \]
  augmented by $F^{S}_0\oplus F^{U}_0\to\cring$, $(u,u')\mapsto\aug(u)+\aug(u')$, $\aug$ the augmentation of $\cring[\Ga]$ (so
  $ye_x\mapsto\aug(y)$). For a $\cring[\Ga]$-module $M$ put
  $C^{\bullet}_{\mathcal T,\ell}(M):=\Hom_{\cring[\Ga]}(\Iso^{(\ell)}_{\mathcal T},M)$.
\item \textbf{(The canonical lifts.)} Writing $\varphi'_t=(r_{a_t},r_{b_t})$ (components into $F^{S}_t$, $F^{U}_t$):
  \begin{enumerate}
  \item[(a)] $\ell=2$: $e_S=e_{\mathfrak e}=1$, $e_U=\tfrac13(1+U^{2}+U^{4})$; $(a_t,b_t)=(1,-e_U)$ for even $t$ and $(1+S,-e_U)$ for odd $t$.
  \item[(b)] $\ell=3$: $e_S=\tfrac14N_S$, $e_U=e_{\mathfrak e}=\tfrac12(1+\mathbf c)$; $(a_0,b_0)=(e_S,-e_{\mathfrak e})$, and $F^{\mathfrak e}_t=0$ for $t\ge1$.
  \item[(c)] $\ell\ge5$: $e_S=\tfrac14N_S$, $e_U=\tfrac16N_U$, $e_{\mathfrak e}=\tfrac12(1+\mathbf c)$; $(a_0,b_0)=(e_S,-e_U)$, and $F^{x}_t=0$ for $t\ge1$.
  \end{enumerate}
  \textbf{The} $\ell$-local tree model $\Iso^{(\ell)}_{\mathcal T}$ is the cone of the canonical lift.
\item \textbf{(The free tree model.)} $\Iso_{\mathcal T}$ is the same construction with $Q_x=1$ for every $x$ (every column free and periodic, $e_x=1$, $N'_x=\sum_{u\in\Ga_x}u$),
  over $\Z$ or any $\cring$, with the lift $(a_t,b_t)=(1,-1)$ for even and $(1+S,-(1+U+U^{2}))$ for odd $t$: the resolution $P_{\bullet}$ of Remark~\ref{rmk:treeconeres};
  $C^{\bullet}_{\mathcal T}(M):=\Hom_{\Z[\Ga]}(\Iso_{\mathcal T},M)$.
\end{enumerate}
\end{definition}

\begin{proposition}[The tree models: what they compute, closed forms, the comparisons with $\Iso$ over $\Z$, fibre maps, Hecke by transport]\label{prop:isomodeltree}
\PROVED{} \STD{\cite[Thm.~2.2.6]{weibel}, \cite[Ch.~III.6.2]{brown}} (v1.23.0-pre.20; milestone D5 of \file{docs/implementation/slot24/DESCENT\_plan.md} v0.4, clause C-D6;
refereed ACCEPT-with-fixes, \file{docs/implementation/slot24/C-D6\_referee.md}; its fixes M1, m1--m5 applied at v1.23.0-pre.21 --- (4) made effective ---, not re-refereed.) Assume \ref{H4} and \ref{B1}--\ref{B3}; let $\ell$ be a prime, $\cring$ an $\mathsf O_\ell$-algebra and $M$ a $\cring[\Ga]$-module; for the
integral statements let $M$ be a $\Z[\Ga]$-module finitely generated and free over $\Z$ (as in Proposition~\ref{prop:descentN}: $\Vcoef$, $M_N$, $M_N(\chi)$) and
$M_\cring:=M\otimes\cring$.
\begin{enumerate}
\item \textbf{(What the models compute.)}
  \begin{enumerate}
  \item[(a)] The canonical lifts of Definition~\ref{def:isomodeltree}(3) are well defined ($a_t\in e_{\mathfrak e}\cring[\Ga]e_S$, $b_t\in e_{\mathfrak e}\cring[\Ga]e_U$) and are
    chain maps lifting the boundary; so is the lift of (4). Any two lifts are homotopic, and their cones are isomorphic by an isomorphism that is the identity on
    $F^{S}\oplus F^{U}$ and on the quotient $F^{\mathfrak e}[-1]$; so every $\ell$-local tree model is isomorphic to \emph{the} one.
  \item[(b)] $\Iso^{(\ell)}_{\mathcal T}\to\cring$ is a resolution of $\cring$ by finitely generated projective $\cring[\Ga]$-modules, $\cring[\Ga]e_x\cong\cring[\Ga/Q_x]$ --- an
    $\cring$-model in the sense of Proposition~\ref{prop:descentmaps} --- and, restricted, a projective resolution of $\cring$ over $\cring[A]$ for every subgroup $A\le\Ga$.
    $\Iso_{\mathcal T}\to\Z$ is a free resolution over $\Z[\Ga]$ (Remark~\ref{rmk:treeconeres}).
  \item[(c)] $C^{q}_{\mathcal T,\ell}(M)=\bigoplus_xM^{Q_x}$, the sum over the cells $x$ of degree $q$ (the columns $F^{S}_q$, $F^{U}_q$, $F^{\mathfrak e}_{q-1}$ that are non-zero), and
    $H^{q}(C^{\bullet}_{\mathcal T,\ell}(M))=H^{q}(\Ga;M)$ for \textbf{every} $q\ge0$ and every $\cring[\Ga]$-module $M$; for $M=\Ind_{\Ga_0(N)}^{\Ga}V$ this is $H^{q}(\Ga_0(N);V)$
    (Shapiro's lemma, \STD{\cite[Ch.~III.6.2]{brown}}), in particular for $M_{N,\cring}$ and $M_N(\chi)_{\cring}$ of \S\ref{sec:descentN}. If $\cring$ is a principal ideal
    domain of characteristic $0$ and $M$ is finitely generated and free over $\cring$, each $M^{Q_x}$ is free of rank
    $\frac1{\abs{Q_x}}\sum_{w\in Q_x}\operatorname{tr}\rho_M(w)$.
  \item[(d)] \textbf{(One integral form; the certified-exponent rule.)} The canonical models have all coefficients in $\mathsf O_\ell[\Ga]$ (the free one in $\Z[\Ga]$). For
    every $\mathsf O_\ell$-algebra $\cring$, $\Iso^{(\ell)}_{\mathcal T,\mathsf O_\ell}\otimes_{\mathsf O_\ell}\cring$ is the canonical model over $\cring$; no lift is chosen per ring.
    Proposition~\ref{prop:certifiedeell}(2)--(6) hold verbatim with $C^{\bullet}_{\mathcal T,\ell}(M_{\mathsf O_\ell})$ in place of $C^{\bullet}_\ell$, with the same $\mathsf q_\ell$:
    the integral cochains are free over $\mathsf O_\ell$ with $H^{q}=H^{q}(\Ga;M)\otimes\mathsf O_\ell$; the rule over $\Z_{(\ell)}$; $\rank_{\Q}d^{q}$ from exactly the two
    sources of (4)(c) --- the reduction mod a qualifying $\ell'\nmid\mathsf q_\ell$ of \textbf{this} complex, or the Betti recursion --- the rank of a differential of
    $\Iso^{(\ell)}$, of $\Iso$ or of the free tree being \textbf{not} admissible for it, nor conversely; the full layout; the Smith-valuation read-out with the a-priori
    bound $\sigma_i\le v_\ell(24)$ for $q\ge1$ and the stopping rule.
  \end{enumerate}
\item \textbf{(The cochains in closed form.)} For every $\cring[\Ga]$-module $M$, with $x\in M$ or $M^{Q_S}$, $y\in M^{Q_U}$, $z\in M^{Q_{\mathfrak e}}$ the components on
  $F^{S}_q$, $F^{U}_q$, $F^{\mathfrak e}_{q-1}$:
  \begin{enumerate}
  \item[(a)] $\ell=2$: $C^{0}=M\oplus M^{U^{2}}$, $C^{q}=M\oplus M^{U^{2}}\oplus M$ ($q\ge1$), and
    \begin{align*}
      d^{0}(x,y)&=\bigl((1-S)x,\,(1-\mathbf c)y,\;x-y\bigr),\\
      d^{q}(x,y,z)&=\bigl(N_Sx,\,(1+\mathbf c)y,\;(1+S)x-y-(1-\mathbf c)z\bigr)\qquad(q\text{ odd}),\\
      d^{q}(x,y,z)&=\bigl((1-S)x,\,(1-\mathbf c)y,\;x-y-(1+\mathbf c)z\bigr)\qquad(q\ge2\text{ even}).
    \end{align*}
  \item[(b)] $\ell=3$: $C^{0}=M^{\langle S\rangle}\oplus M^{-I}$, $C^{1}=M^{-I}\oplus M^{-I}$ (the $U$-column in degree $1$ and the edge), $C^{q}=M^{-I}$ ($q\ge2$), and
    \begin{align*}
      d^{0}(x,y)&=\bigl((1-U)y,\;x-y\bigr),\qquad d^{1}(y,z)=(1+U+U^{2})y,\\
      d^{q}&=1-U\quad(q\ge2\text{ even}),\qquad d^{q}=1+U+U^{2}\quad(q\ge3\text{ odd}).
    \end{align*}
    In degrees $\ge2$ this is the periodic complex of the cyclic group $\langle U^{2}\rangle$ acting on $M^{-I}$ (on which $U$ acts as $U^{-2}$), whence
    $H^{q}(\Ga;M)\cong H^{q}(\langle U^{2}\rangle;M^{-I})$ for $q\ge2$ --- Proposition~\ref{prop:ellthreeclosed}(2), read off the cochains.
  \item[(c)] $\ell\ge5$: $C^{0}=M^{\langle S\rangle}\oplus M^{\langle U\rangle}$, $C^{1}=M^{-I}$, $d^{0}(x,y)=x-y$, $C^{q}=0$ ($q\ge2$): the tree complex \eqref{eq:treemodel} up to the
    sign of $d^{0}$.
  \item[(d)] \textbf{(Periodicity.)} At $\ell=2$, $C^{q+2}=C^{q}$ and $d^{q+2}=d^{q}$ for $q\ge1$; at $\ell=3$ for $q\ge2$. In both cases $H^{q+2}(\Ga;M)=H^{q}(\Ga;M)$ as
    groups for $q\ge2$, read on the same cochains, so a run built to degree $q_{\max}+1$ with $q_{\max}=3$ yields every degree (that this identification is $e\cup-$ of
    Proposition~\ref{prop:descentN}(5) is \textbf{not} asserted).
  \item[(e)] \textbf{(Ranks.)} With $r_H:=\rank M^{H}$ (for $M$ as in (1)(c); $r_1=\rank M$): $\rank C^{q}$ is $r_1+r_{U^{2}}$, $2r_1+r_{U^{2}}$ ($q\ge1$) at $\ell=2$;
    $r_{\langle S\rangle}+r_{-I}$, $2r_{-I}$, $r_{-I}$ ($q\ge2$) at $\ell=3$; $r_{\langle S\rangle}+r_{\langle U\rangle}$, $r_{-I}$ at $\ell\ge5$; and $2r_1$, $3r_1$ for the free tree.
  \end{enumerate}
\item \textbf{(Comparison with the free tree, in closed form.)} Let $\ell$ be a prime and $\cring$ an $\mathsf O_\ell$-algebra. Define $\Phi^{\mathcal T}_\ell\colon\Iso_{\mathcal T}\otimes\cring\to
  \Iso^{(\ell)}_{\mathcal T}$ and $\Psi^{\mathcal T}_\ell\colon\Iso^{(\ell)}_{\mathcal T}\to\Iso_{\mathcal T}\otimes\cring$ by
  \[
    \Phi^{\mathcal T}_\ell(a,z)=\bigl(f\,a+k\,z,\;f\,z\bigr),\qquad\Psi^{\mathcal T}_\ell(a,z)=\bigl(g\,a+k'z,\;g\,z\bigr),
  \]
  where on each column $f^{x}_t:=r_{\abs{Q_x}^{\lfloor t/2\rfloor}e_x}$ and $g^{x}_t$ is $\abs{Q_x}^{-\lfloor t/2\rfloor}$ times the inclusion $\cring[\Ga]e_x\subseteq\cring[\Ga]$ (both
  $0$ where the $\ell$-local column vanishes), and the corrections $k_t\colon F^{\mathfrak e}_t\to F^{U}_{t+1}$ (free source, $\ell$-local target), $k'_t$ from the
  $\ell$-local $F^{\mathfrak e}_t$ to the free vertex columns of degree $t+1$, are:
  \begin{enumerate}
  \item[(a)] $\ell=2$: $k_t=r_{-\mathsf k_{\lceil t/2\rceil}e_U}$ with $\mathsf k_j:=\frac{3^{j}-1}2$; $k'_t=r_{\kappa'_t}$ into $F^{U}_{t+1}$, with
    $\kappa'_t:=\frac{1-3^{-\lceil t/2\rceil}}2\,e_U+[t\text{ even}]\,\mathsf w$ and $\mathsf w:=\frac13(2+2U+U^{2}+U^{3})-(1+\mathbf c)e_U$;
  \item[(b)] $\ell=3$: $k_t=r_{-2^{(t-1)/2}e_{\mathfrak e}}$ for odd $t$, $k_t=0$ for even $t$; $k'_0=r_{-\frac12e_{\mathfrak e}}$ into $F^{S}_1$ (and $k'_t=0$, $t\ge1$, its source being $0$);
  \item[(c)] $\ell\ge5$: $k=0$; $k'_0=\bigl(r_{-\frac12e_{\mathfrak e}},\,r_{\frac13e_{\mathfrak e}(2+2U+U^{2}+U^{3})}\bigr)$ into $F^{S}_1\oplus F^{U}_1$.
  \end{enumerate}
  Then $\Phi^{\mathcal T}_\ell$ and $\Psi^{\mathcal T}_\ell$ are $\cring[\Ga]$-chain maps over $\id_\cring$, defined over $\mathsf O_\ell$, with
  \[
    \Phi^{\mathcal T}_\ell\circ\Psi^{\mathcal T}_\ell=\id\quad\text{strictly},\qquad\Psi^{\mathcal T}_\ell\circ\Phi^{\mathcal T}_\ell\simeq\id .
  \]
  Hence $\Psi^{\mathcal T\ast}_\ell\Phi^{\mathcal T\ast}_\ell=\id$ on $C^{\bullet}_{\mathcal T,\ell}(M)$ and $\Phi^{\mathcal T\ast}_\ell\Psi^{\mathcal T\ast}_\ell\simeq\id$ on $C^{\bullet}_{\mathcal T}(M)$:
  \textbf{$C^{\bullet}_{\mathcal T,\ell}(M)$ is a retract of $C^{\bullet}_{\mathcal T}(M)$ with $\Psi^{\mathcal T\ast}_\ell\Phi^{\mathcal T\ast}_\ell=\id$ strictly and $\Phi^{\mathcal T\ast}_\ell\Psi^{\mathcal T\ast}_\ell\simeq\id$},
  the homotopy being $\mathsf H_\ell^{\ast}$ for the effective $\mathsf H_\ell$ of (4)(c) (no side conditions such as $\mathsf h\Phi^{\ast}=0$, $\Psi^{\ast}\mathsf h=0$, $\mathsf h^{2}=0$ are asserted), and $\Phi^{\mathcal T\ast}_\ell$, $\Psi^{\mathcal T\ast}_\ell$ are mutually inverse on
  cohomology. On cochains, with $m:=\lfloor q/2\rfloor$ and the components of (2): at $\ell=2$, $\Phi^{\mathcal T\ast}_2(x,y,z)=(x,\,3^{m}y,\,z-\mathsf k_my)$ and
  $\Psi^{\mathcal T\ast}_2(x,y,z)=(x,\,3^{-m}e_Uy,\,z+\kappa'_{q-1}y)$ for $q\ge1$, the identity and $(x,y)\mapsto(x,e_Uy)$ for $q=0$; at $\ell=3$,
  $\Phi^{\mathcal T\ast}_3$ is $(x,y)\mapsto(x,y)$, $(y,z)\mapsto(0,y,z)$ and, for $q\ge2$, $y\mapsto(0,\,2^{m}y,\,-[q\text{ even}]\,2^{m-1}y)$, and $\Psi^{\mathcal T\ast}_3$ is
  $(x,y)\mapsto(e_Sx,e_{\mathfrak e}y)$, $(x,y,z)\mapsto\bigl(e_{\mathfrak e}y,\,e_{\mathfrak e}(z-\tfrac12x)\bigr)$ and, for $q\ge2$, $(x,y,z)\mapsto2^{-m}e_{\mathfrak e}y$.
\item \textbf{(Comparison with $\Iso$, over $\Z$: effective.)} (v1.23.0-pre.21; the referee's M1.) Simplices of $K$ stand for their lifts into $\Xbulk$
  (Lemma~\ref{lem:model7coincidences}), $\Sigma_s$, $\Ga_\sigma$, $g_\sigma$, $m_\sigma$, $G_{s,t}$, $d^{(r)}$ and $\cmpI$ are those of Definition~\ref{def:isomodel}, and
  $\Iso_{\bullet}$ is any isotropy-resolution model (e.g.\ the one of Proposition~\ref{prop:isomodel}(1)'s procedure, which the code builds).
  \begin{enumerate}
  \item[(a)] \textbf{(The tree in $\Xbulk$, a retraction, a homotopy.)} The vertices in the $\Ga$-orbits of $v_6=i$ and $v_4=\zeta^{2}$ and the edges in the orbit of
    $\{v_4,v_6\}$ form a $\Ga$-subcomplex of $\Xbulk$, and
    $\mathsf i_{\mathcal T}\colon C_{\bullet}(\mathcal T)\to C_{\bullet}(\Xbulk)$, $\gamma\mathfrak v_S\mapsto\gamma v_6$, $\gamma\mathfrak v_U\mapsto\gamma v_4$, $\gamma\mathfrak e\mapsto\gamma\{v_4,v_6\}$ (stored
    orientation), is an injective $\Z[\Ga]$-chain map preserving augmentations. Define $\mathsf r_{\mathcal T}\colon C_{\bullet}(\Xbulk)\to C_{\bullet}(\mathcal T)$ and
    $\mathsf G_{\mathcal T}\colon C_s(\Xbulk)\to C_{s+1}(\Xbulk)$ $\Z[\Ga]$-linearly by their values on the representatives:
    \begin{center}
    \begin{tabular}{@{}lll@{}}
    $\sigma\in\Sigma_s$ & $\mathsf r_{\mathcal T}(\sigma)$ & $\mathsf G_{\mathcal T}(\sigma)$ \\ \hline
    $v_0$ & $\mathfrak v_U$ & $-\{v_0,v_4\}$ \\
    $v_1$; $v_2$ & $\mathfrak v_S$; $\mathfrak v_S$ & $-\{v_1,v_6\}$; $-\{v_2,v_6\}$ \\
    $v_4$; $v_6$ & $\mathfrak v_U$; $\mathfrak v_S$ & $0$; $0$ \\
    $\{v_0,v_1\}$ & $\mathfrak e$ & $\{v_0,v_1,v_4\}-\{v_1,v_4,v_6\}$ \\
    $\{v_1,v_2\}$; $\{v_1,v_4\}$ & $0$; $-\mathfrak e$ & $\{v_1,v_2,v_6\}$; $\{v_1,v_4,v_6\}$ \\
    $\{v_2,v_3\}$ & $-S\mathfrak e$ & $\{v_2,v_3,v_5\}+\{v_2,v_5,v_6\}$ \\
    $\{v_2,v_5\}$ & $-S\mathfrak e$ & $\{v_2,v_5,v_6\}$ \\
    $\{v_4,v_6\}$ & $\mathfrak e$ & $0$ \\
    $\{v_0,v_4\}$, $\{v_1,v_6\}$, $\{v_2,v_6\}$ & $0$ & $0$ \\
    the five triangles & $0$ & $0$
    \end{tabular}
    \end{center}
    Then $\mathsf r_{\mathcal T}$ and $\mathsf G_{\mathcal T}$ are well defined, $\mathsf r_{\mathcal T}$ is a $\Z[\Ga]$-chain map preserving augmentations,
    $\mathsf r_{\mathcal T}\circ\mathsf i_{\mathcal T}=\id$, and $\id-\mathsf i_{\mathcal T}\circ\mathsf r_{\mathcal T}=\partial\mathsf G_{\mathcal T}+\mathsf G_{\mathcal T}\partial$ on
    $C_{\bullet}(\Xbulk)$: the equivalence of Proposition~\ref{prop:bredon}(2), in closed form.
  \item[(b)] \textbf{(Contractions and the lifting formula.)} Fix for each $\sigma\in\Sigma_0\sqcup\Sigma_1\sqcup\Sigma_2\sqcup\{S,U,\mathfrak e\}$ a set $\mathsf R_\sigma\subset\Ga$ of
    representatives of $\Ga/\Ga_\sigma$ (any choice; the code's is \code{decompose()} of \file{Integral/Model/PeriodicColumn.h}, shared by \file{IsotropyModel},
    \file{IsotropyModelEll}, \file{StrictLift}, \file{RefinedStrictLift} and \file{PeriodicityClass} (v1.23.0-pre.23, referee t2) and documented at \code{GKey} in \file{IsotropyModel.h}: the least matrix of the coset in the \textbf{signed} lexicographic
    order on $(a,b,c,d)\in\Z^{4}$, and $j$ taken relative to the \textbf{fixed} generator $g_\sigma$ --- v1.23.0-pre.22, re-referee N1); write
    $\gamma=\rho g_\sigma^{j}$ with $\rho\in\mathsf R_\sigma$, $0\le j<m_\sigma$. On a column $\Z[\Ga]$ with differentials $r_{g-1}$ (odd) and $r_N$ (even, $N$ the norm of
    $\langle g\rangle$), as in Definition~\ref{def:isomodel}, put
    \[
      \mathsf h(\rho g^{j}):=\rho\textstyle\sum_{i<j}g^{i}\ \text{(from even degrees)},\qquad \mathsf h(\rho g^{j}):=[j=m-1]\,\rho\ \text{(from odd degrees)},
    \]
    and $\mathsf h_{-1}(\gamma\Ga_\sigma):=\rho$; for the tree's columns ($r_{1-g}$, $r_N$) negate $\mathsf h$ from even degrees; on a column that enters a total complex with a global sign $\epsilon$ ($G_{s,\bullet}$ with
    $\epsilon=(-1)^{s}$; the edge column of a cone, $\epsilon=-1$) use $\epsilon\mathsf h$. Let $p$ be $\cmpI\colon\Iso_{\bullet}\to C_{\bullet}(\Xbulk)$ or the column augmentation
    $\varepsilon_{\mathcal T}\colon\Iso_{\mathcal T}\to C_{\bullet}(\mathcal T)$ (vertex columns $s=0$, edge column $s=1$; a surjective quasi-isomorphism, (1)(b)'s proof). For a cycle $z$ of
    $\ker p$ define $\mathsf h_p(z)$ by descending induction on the column index $s$: if $z_s$ is the component of $z$ on the largest $s$, put $k_s:=\epsilon\mathsf h(z_s)$
    column by column, replace $z$ by $z-\partial k_s$ and continue; after at most three steps $z=0$, and $\mathsf h_p(z):=\sum k_s$ satisfies $\partial\mathsf h_p(z)=z$ with
    $\mathsf h_p(z)$ in column degrees $\ge1$, so $p\,\mathsf h_p(z)=0$. Let $P'$ be a bounded-below complex of free $\Z[\Ga]$-modules with given bases and $f\colon P'\to C$ a chain map into
    the target of $p$. For a basis element $x$, with $y_0$ any preimage of $f(x)$ under $p$ (e.g.\ $\rho\otimes1$ in $G_{s,0}$, resp.\ $F^{x}_0$, for the basis element $\rho\sigma$),
    \[
      \text{(L)}\qquad \hat f(x):=y_0+\mathsf h_p\bigl(\hat f(\partial x)-\partial y_0\bigr),
    \]
    extended $\Z[\Ga]$-linearly, is a chain map with $p\hat f=f$: Lemma~\ref{lem:strictlifting}, made effective. Likewise, for chain maps $A,B\colon P'\to E$ ($E$ the source of $p$) with
    $pA=a\,p'$, $pB=b\,p'$ for a $\Z[\Ga]$-chain map $p'\colon P'\to C'$ (no surjectivity is needed; v1.23.0-pre.22, re-referee N3) and $a-b=\partial g+g\partial$ for a $\Z[\Ga]$-map $g\colon C'_\bullet\to C_{\bullet+1}$, and $y_1$ a preimage of
    $g\,p'(x)$,
    \[
      \text{(L$'$)}\qquad H(x):=y_1+\mathsf h_p\bigl(Ax-Bx-H(\partial x)-\partial y_1\bigr)
    \]
    gives $\partial H+H\partial=A-B$ and $pH=g\,p'$. (L) and (L$'$) are finite in each degree, but the maps they produce are \textbf{not periodic} in the degree: with the
    code's representatives the number of terms per generator is bounded, but the integer coefficients grow linearly in $q$ (so (2)'s periodicity applies to the cochains,
    not to the comparisons); they are built to the degree in use --- to $q+1$ for $\Phi^{\Z}_{\mathcal T}$, $\Psi^{\Z}_{\mathcal T}$ and to $q$ for the homotopies when an operator is
    transported in degree $q$ (v1.23.0-pre.22, re-referee N2).
  \item[(c)] \textbf{(The integral pair and its homotopies.)} $\Phi^{\Z}_{\mathcal T}\colon\Iso_{\bullet}\to\Iso_{\mathcal T}$ is (L) for $f=\mathsf r_{\mathcal T}\circ\cmpI$ along
    $p=\varepsilon_{\mathcal T}$, and $\Psi^{\Z}_{\mathcal T}\colon\Iso_{\mathcal T}\to\Iso_{\bullet}$ is (L) for $f=\mathsf i_{\mathcal T}\circ\varepsilon_{\mathcal T}$ along $p=\cmpI$: $\Z[\Ga]$-chain maps
    over $\mathsf r_{\mathcal T}$, resp.\ $\mathsf i_{\mathcal T}$, hence over $\id_\Z$ (compatible with $\aug\circ\cmpI$ and the augmentation of $\Iso_{\mathcal T}$). By (L$'$):
    $\mathsf H_{\mathcal T}$ ($A=\Phi^{\Z}_{\mathcal T}\Psi^{\Z}_{\mathcal T}$, $B=\id$, along $\varepsilon_{\mathcal T}$, $g=0$ as $\mathsf r_{\mathcal T}\mathsf i_{\mathcal T}=\id$) with
    $\partial\mathsf H_{\mathcal T}+\mathsf H_{\mathcal T}\partial=\Phi^{\Z}_{\mathcal T}\Psi^{\Z}_{\mathcal T}-\id$; $\mathsf H_X$ ($A=\id$, $B=\Psi^{\Z}_{\mathcal T}\Phi^{\Z}_{\mathcal T}$, along $\cmpI$,
    $g=\mathsf G_{\mathcal T}$) with $\partial\mathsf H_X+\mathsf H_X\partial=\id-\Psi^{\Z}_{\mathcal T}\Phi^{\Z}_{\mathcal T}$; and, for each $\ell$ and $\cring$, $\mathsf H_\ell$
    ($A=\Psi^{\mathcal T}_\ell\Phi^{\mathcal T}_\ell$, $B=\id$ on $\Iso_{\mathcal T}\otimes\cring$, along $\varepsilon_{\mathcal T}\otimes\cring$, $g=0$) with
    $\partial\mathsf H_\ell+\mathsf H_\ell\partial=\Psi^{\mathcal T}_\ell\Phi^{\mathcal T}_\ell-\id$. $\Phi^{\Z}_{\mathcal T}$, $\Psi^{\Z}_{\mathcal T}$ are built once over $\Z$, independently of $\ell$;
    another choice of $\Iso_{\bullet}$, of the $\mathsf R_\sigma$ or of the preimages changes them by a homotopy; as comparisons over $\id_\Z$ they are unique up to homotopy.
  \item[(d)] \textbf{(Over $\cring$.)} Put
  \[
    \Phi_{\mathcal T}:=\Phi^{\mathcal T}_\ell\circ(\Phi^{\Z}_{\mathcal T}\otimes\cring)\colon\Iso_{\bullet}\otimes\cring\to\Iso^{(\ell)}_{\mathcal T},\qquad
    \Psi_{\mathcal T}:=(\Psi^{\Z}_{\mathcal T}\otimes\cring)\circ\Psi^{\mathcal T}_\ell\colon\Iso^{(\ell)}_{\mathcal T}\to\Iso_{\bullet}\otimes\cring .
  \]
  Then $\Phi_{\mathcal T}\Psi_{\mathcal T}-\id=\partial\mathsf K+\mathsf K\partial$ with $\mathsf K:=\Phi^{\mathcal T}_\ell\mathsf H_{\mathcal T}\Psi^{\mathcal T}_\ell$, and
  $\id-\Psi_{\mathcal T}\Phi_{\mathcal T}=\partial\mathsf K'+\mathsf K'\partial$ with $\mathsf K':=\mathsf H_X-\Psi^{\Z}_{\mathcal T}\mathsf H_\ell\Phi^{\Z}_{\mathcal T}$ (all $\otimes\cring$).
  So $\Phi^{\ast}_{\mathcal T}\colon C^{\bullet}_{\mathcal T,\ell}(M)\to\CI(M):=\Hom_{\Z[\Ga]}(\Iso_{\bullet},M)$ and $\Psi^{\ast}_{\mathcal T}$ are mutually inverse cochain homotopy equivalences, with the
  explicit homotopies $\mathsf K^{\ast}$, $\mathsf K'^{\ast}$, compatible with the identifications of both sides with $H^{\bullet}(\Ga;M)$; every other comparison between the two (e.g.\ a strict lift
  through $\Iso^{(\ell)}_{\mathcal T}$ directly, or one through $\Iso^{(\ell)}$ and $\Phi_\ell$, $\Psi_\ell$ --- Proposition~\ref{prop:isomodelell}(4), whose proof uses of $\cring$ only that every
  $\abs{Q_\sigma}$ is a unit, so it holds over every $\mathsf O_\ell$-algebra) is homotopic to these. The comparisons are maps \textbf{over the augmentation to $\Z$} (resp.\ $\cring$);
  (4) needs no chain map between $\Iso^{(\ell)}_{\mathcal T}$ and $C_{\bullet}(\Xbulk)$: the effective construction passes through $C_{\bullet}(\mathcal T)\rightleftarrows C_{\bullet}(\Xbulk)$,
  at the integral level only.
  \end{enumerate}
\item \textbf{(Fibre maps and the level cone: (T1)--(T3) hold.)} For $\ell\in\{2,3\}$, $\Iso^{(\ell)}_{\mathcal T}$ satisfies (T1)--(T3) of Proposition~\ref{prop:descentmaps}(5): it is an
  $S$-model for $S=\mathsf O_\ell$, a localisation of $\Z$ in which $\ell$ is not inverted, and for $S=\Z_{(\ell)}$ ((1)(b), (1)(d)); its terms are
  $\bigoplus_xS[\Ga/Q_x]$ with $\abs{Q_x}\in S^{\times}$ ((1)(c)); and (4) gives the comparisons and the homotopies between their composites explicitly, the integral pair $\Phi^{\Z}_{\mathcal T}$, $\Psi^{\Z}_{\mathcal T}$
  included (v1.23.0-pre.21).
  Hence Proposition~\ref{prop:descentmaps}(1), (2) and (4) hold on it verbatim: a $\Z[\Ga]$-map $u$ of fibres (restriction $\alpha$, trace $\omega$, $R_\sigma$, the diamonds,
  the $\mathsf O_\chi$-action) acts on $C^{\bullet}_{\mathcal T,\ell}$ block by block, on each $M^{Q_x}$, and commutes \textbf{strictly} with $\Phi^{\ast}_{\mathcal T}$, $\Psi^{\ast}_{\mathcal T}$,
  $\Phi^{\mathcal T\ast}_\ell$, $\Psi^{\mathcal T\ast}_\ell$; the level cone $\Cone{\alpha_{\mathcal T}}$ on $\Iso^{(\ell)}_{\mathcal T}$ is exact in the sense of (2)(a) there, computes $H^{\bullet-1}(A;Q_{A/B}\otimes
  S)$, carries the long exact sequence \eqref{eq:descentmapsles} over $\Z_{(\ell)}$ in every degree, and every comparison above induces a quasi-isomorphism of cones commuting
  with both maps of the cone's short exact sequence; the certified-exponent rule of (2)(d) there holds for $\Dcone^{q}$ on it, with source (i) read on the reduction mod a
  qualifying $\ell'\nmid\mathsf q_\ell$ of \textbf{this} cone on $\Iso^{(\ell)}_{\mathcal T,\mathsf O_\ell}$. So the $\ell$-local pieces of Proposition~\ref{prop:descentN}(2) and of
  Proposition~\ref{prop:descentmaps}(1)(c), (2)(c) may be read on the tree models at $\ell=2$ and $\ell=3$, with no change to the glueing datum $\cmpI^{\ast}$ of the
  $\Z[1/6]$-piece. This is a statement about the groups and the maps induced by fibre maps; any other map between the glued groups (an explicit fibre-product
  isomorphism, a Hecke operator through Proposition~\ref{prop:descentN}(3)) enters the $\ell$-local piece through $\Phi_{\mathcal T}$, $\Psi_{\mathcal T}$ of (4)(d).
\item \textbf{(Hecke by transport.)} Let $T$ be \emph{any} cochain endomorphism of $\CI(M_\cring)$ inducing an endomorphism $t$ of $H^{\bullet}(\Ga;M_\cring)$ --- e.g.\ $\Th_{\Iso,g}$ of
  Proposition~\ref{prop:strictlift}(3), $g\in\De$, at level~$1$; (a)--(d) are stated for every such $T$, $U_p$-type cochain operators included once they exist, and are
  instantiated at level $N$ by the cochain-level $T_p$, $p\nmid N$, of the plan's clause C-D4 once it is stated (the plan's ruling [R-Q4]: $U_p$ deferred) --- and put
  \[
    T'_{\mathcal T}:=\Psi^{\ast}_{\mathcal T}\circ T\circ\Phi^{\ast}_{\mathcal T}\colon C^{\bullet}_{\mathcal T,\ell}(M_\cring)\to C^{\bullet}_{\mathcal T,\ell}(M_\cring).
  \]
  \begin{enumerate}
  \item[(a)] $T'_{\mathcal T}$ is a cochain map inducing $t$ under the identification of (1)(c), in every degree; up to cochain homotopy it does not depend on the choice of
    $\Phi^{\Z}_{\mathcal T}$, $\Psi^{\Z}_{\mathcal T}$ (nor of any comparison homotopic to $\Phi_{\mathcal T}$, $\Psi_{\mathcal T}$).
  \item[(b)] \textbf{(It equals the operator on $\Iso^{(\ell)}$.)} For $T'':=\Psi_\ell^{\ast}\circ T\circ\Phi_\ell^{\ast}$ on $C^{\bullet}_{\Iso^{(\ell)}}(M_\cring)$
    (Proposition~\ref{prop:isomodelell}(5), which holds over every $\mathsf O_\ell$-algebra $\cring$: its proof, and that of (4) there, use of $\cring$ only that every $\abs{Q_\sigma}$ is a
    unit) and every comparison $\vartheta\colon\Iso^{(\ell)}_{\mathcal T}\to\Iso^{(\ell)}$ over $\id_\cring$, $H(\vartheta^{\ast})\circ H(T'')=H(T'_{\mathcal T})\circ
    H(\vartheta^{\ast})$: the two transports are the same operator on $H^{\bullet}(\Ga;M_\cring)$. An effective such comparison is $\vartheta:=\Phi_\ell\circ\Psi_{\mathcal T}$
    ($\Phi_\ell$ of Proposition~\ref{prop:isomodelell}(4), $\Psi_{\mathcal T}$ of (4)(d); built from constructed pieces only), the one a numerical check uses (v1.23.0-pre.22,
    re-referee N4; wording v1.23.0-pre.23, referee t1).
  \item[(c)] For two such $T_1$, $T_2$, $T'_{1,\mathcal T}T'_{2,\mathcal T}\simeq(T_1T_2)'_{\mathcal T}$; so every relation that holds among the $T_i$ up to cochain homotopy holds among their
    transports up to homotopy, and on cohomology exactly.
  \item[(d)] If $T$ commutes strictly with $\Hom(\Iso_{\bullet},u)$ for a fibre map $u$, then $T'_{\mathcal T}$ commutes strictly with $u$ on $C^{\bullet}_{\mathcal T,\ell}$; given level-$A$,
    level-$B$ operators with $T^{B}\alpha_{\Iso}=\alpha_{\Iso}T^{A}$ strictly (the plan's C-D4), $(T'^{A}_{\mathcal T}[1],T'^{B}_{\mathcal T})$ is a cochain endomorphism of
    $\Cone{\alpha_{\mathcal T}}$ commuting with its short exact sequence, so Proposition~\ref{prop:descentmaps}(3)(c)'s equivariance holds on the tree.
  \end{enumerate}
  No $\tau$-type table on the tree is used: the Hecke datum stays on $\Iso$, and applying $T'_{\mathcal T}$ composes cochain maps, all of them effective by (4).
\end{enumerate}
\textbf{Scope.} Bulk only. Nothing here is stated about the boundary $\Xbdry$, $\Isob_{\bullet}$, the pair $(\Ga,\Ga_\infty)$, its long exact sequence, its cone, or a Hecke
action on any of them (the plan's C-D8, tree cusp pair; v1.23.0-pre.33: \S\ref{sec:treepair}); nothing about the Bredon object $\Ceq$ at $\ell\in\{2,3\}$; and no native Hecke datum on the tree.
\end{proposition}

\begin{proof}
\emph{A cone lemma} ($\ast$). Let $\varphi\colon E\to V$, $\varphi'\colon E'\to V'$ be chain maps, $f^{V}\colon V\to V'$, $f^{E}\colon E\to E'$ chain maps and $k\colon E_t\to V'_{t+1}$ maps
with $\partial k+k\partial=f^{V}\varphi-\varphi'f^{E}$. Then $F(a,z):=(f^{V}a+kz,\,f^{E}z)$ is a chain map $\Cone{\varphi}\to\Cone{\varphi'}$ (cones as in
Definition~\ref{def:isomodeltree}(2)): $\partial F(a,z)=(\partial f^{V}a+\partial kz+\varphi'f^{E}z,\,-\partial f^{E}z)$ and $F\partial(a,z)=(f^{V}\partial a+f^{V}\varphi z-k\partial z,\,-f^{E}\partial z)$
agree exactly when the displayed identity holds. With $V=V'$, $E=E'$, $f=\id$ and $\varphi_2-\varphi_1=\partial h+h\partial$ it gives the isomorphism $(a,z)\mapsto(a+hz,z)$
from $\Cone{\varphi_2}$ to $\Cone{\varphi_1}$, inverse $(a,z)\mapsto(a-hz,z)$. \emph{Conventions.} Every map below between the modules $\cring[\Ga]e_x$ is $r_a$ for an
$a\in e_y\cring[\Ga]e_x$ (well defined on $\cring[\Ga]e_y$ exactly then); $r_b\circ r_a=r_{ab}$; all elements that occur lie in the commutative rings $\cring[\langle S\rangle]$ or
$\cring[\langle U\rangle]$, which contain the central $\mathbf c=S^{2}=U^{3}$, and in them $e_U$, $e_{\mathfrak e}$, $e_S$ commute with everything. Repeatedly used:
$U e_U=\mathbf c e_U$ at $\ell=2$ ($U=\mathbf cU^{-2}$ and $U^{-2}e_U=e_U$), so $(1-U)e_U=(1-\mathbf c)e_U$, $(1+U+U^{2})e_U=(2+\mathbf c)e_U$ and $N_Ue_U=3(1+\mathbf c)e_U$;
$(1+S)(1-S)=1-\mathbf c$, $(1+S)(1+\mathbf c)=N_S$, $(1+U+U^{2})(1-U)=1-\mathbf c$, $(1+U+U^{2})(1+\mathbf c)=N_U$; and $N_xe_x=\abs{Q_x}N'_x$, since $\Ga_x\to\Ga_x/Q_x$ has
fibres of size $\abs{Q_x}$ and $\{g_x^{j}\}_{j<\abs{P_x}}$ maps onto $\Ga_x/Q_x$, which is cyclic of order $\abs{P_x}$ generated by the image of $g_x$.

(1)(a) \emph{Well defined}: at $\ell=2$, $e_{\mathfrak e}=e_S=1$ and $e_U^{2}=e_U$; at $\ell=3$ and $\ell\ge5$, $\mathbf c\in\langle S\rangle$ and $\mathbf c\in\langle U\rangle$, so
$e_{\mathfrak e}e_S=e_S$, $e_{\mathfrak e}e_U=e_U$. \emph{Chain maps}: the condition is $a_t\,\delta^{S}_t=\varepsilon_t\,a_{t-1}$ and $b_t\,\delta^{U}_t=\varepsilon_t\,b_{t-1}$ for
$t\ge1$, where $\delta^{x}_t$, $\varepsilon_t$ are the coefficients of $\partial^{x}_t$, $\partial^{\mathfrak e}_t$. At $\ell=2$: for odd $t$,
$(1+S)(1-S)=1-\mathbf c=\varepsilon_t\cdot1$ and $-e_U(1-U)e_U=-(1-\mathbf c)e_U=\varepsilon_t(-e_U)$; for even $t\ge2$, $1\cdot N_S=(1+\mathbf c)(1+S)$ and $-e_U\cdot N'_U=-(1+U)e_U=-(1+\mathbf c)e_U=
\varepsilon_t(-e_U)$. At $\ell=3$ and $\ell\ge5$ the edge column vanishes in degrees $\ge1$, so the condition is empty. The free lift is Remark~\ref{rmk:treeconeres}'s, checked by
the same identities with $1+U+U^{2}$ in place of $e_U$. \emph{Lifting}: in degree $0$, $(a_0,b_0)$ applied to $e_{\mathfrak e}$ augments to $\aug(a_0)\langle S\rangle-\aug(b_0)\langle U\rangle
=\langle S\rangle-\langle U\rangle$, the image of $\langle\mathbf c\rangle$. \emph{Homotopic}: two lifts of the same map from the projective complex $F^{\mathfrak e}$ to the
resolution $F^{S}\oplus F^{U}$ of $C_0(\mathcal T)\otimes\cring$ are homotopic \STD{\cite[Thm.~2.2.6]{weibel}} (projectivity: (b)); the cone lemma gives the isomorphism.

(1)(b) $e_x$ is an idempotent of $\cring[\Ga_x]$ ($\abs{Q_x}\in\cring^{\times}$), $\cring[\Ga]e_x\oplus\cring[\Ga](1-e_x)=\cring[\Ga]$, and $ye_x\mapsto yQ_x$ identifies $\cring[\Ga]e_x$ with
$\cring[\Ga]\otimes_{\cring[Q_x]}\cring=\cring[\Ga/Q_x]$. Under $\cring[\Ga_x]e_x\cong\cring[\Ga_x/Q_x]$, $g_xe_x\mapsto\bar g_x$ and $N'_x\mapsto\sum_{y\in\Ga_x/Q_x}y$, so the column is
$\cring[\Ga]\otimes_{\cring[\Ga_x]}$ of the periodic resolution of $\cring$ over the cyclic group $\Ga_x/Q_x\cong P_x$ inflated to $\Ga_x$ (for $P_x\ne1$; for $P_x=1$, of
$\cring\to\cring$), with the sign of the odd differentials reversed relative to Definition~\ref{def:isomodelell}'s $\bar g-1$ (Remark~\ref{rmk:treeconeres} already uses $1-t$; an
isomorphic complex: multiply degree $t$ by $(-1)^{\lceil t/2\rceil}$). That complex is exact over every
ring (the proof of Proposition~\ref{prop:isomodelell}, ``the columns'', and Proposition~\ref{prop:certifiedeell}(1)), and $\cring[\Ga]\otimes_{\cring[\Ga_x]}-$ is exact ($\cring[\Ga]$ is
free over $\cring[\Ga_x]$), so $F^{x}\to\cring[\Ga/\Ga_x]$ is a projective resolution. The column augmentations give a map of cones from $\Iso^{(\ell)}_{\mathcal T}$ to the cone
of the boundary $C_1(\mathcal T)\otimes\cring\to C_0(\mathcal T)\otimes\cring$, i.e.\ to $C_{\bullet}(\mathcal T)\otimes\cring$; it is a quasi-isomorphism by the five lemma on the long
sequences of $0\to F^{S}\oplus F^{U}\to\Cone{\varphi'}\to F^{\mathfrak e}[-1]\to0$ and of its target, the outer maps being quasi-isomorphisms. $C_{\bullet}(\mathcal T)\otimes\cring\to\cring$
is a quasi-isomorphism because \eqref{eq:treeses} is split exact over $\Z$ (the proof of Lemma~\ref{lem:treeres}). The terms are finitely generated projective. For $A\le\Ga$,
$\cring[\Ga]$ is free over $\cring[A]$, so projective $\cring[\Ga]$-modules are projective over $\cring[A]$. The free case is Remark~\ref{rmk:treeconeres}.

(1)(c) $\Hom_{\cring[\Ga]}(\cring[\Ga]e_x,M)\to e_xM=M^{Q_x}$, $c\mapsto c(e_x)$, is an isomorphism, and $H^{q}(\Hom(\mathsf P,M))=\Ext^{q}_{\cring[\Ga]}(\cring,M)=H^{q}(\Ga;M)$ for
every projective resolution $\mathsf P$ (Proposition~\ref{prop:isomodelell}(3)'s argument). Shapiro: \STD{\cite[Ch.~III.6.2]{brown}}, induction and coinduction agreeing at finite
index. Ranks: $\rho_M(e_x)$ is an idempotent endomorphism of the finitely generated free module $M$ over the principal ideal domain $\cring$; its image is a direct summand,
hence free, of rank $\dim$ of the image over the fraction field, which for an idempotent in characteristic $0$ is its trace, $\abs{Q_x}^{-1}\sum_w\operatorname{tr}\rho_M(w)$.

(1)(d) The denominators of the canonical lifts, of the $e_x$ and of $N'_x$ are divisors of $\abs{Q_x}$, which divide $\mathsf q_\ell$. Base change replaces each coefficient
by its image and keeps the shape; the augmented model over $\mathsf O_\ell$ is a bounded-below exact complex of free $\mathsf O_\ell$-modules ($\mathsf O_\ell[\Ga]e_x\cong
\mathsf O_\ell[\Ga/Q_x]$), hence split exact, and stays exact after $-\otimes_{\mathsf O_\ell}\cring$ (the argument of Proposition~\ref{prop:certifiedeell}(1)'s proof). The proofs of
Proposition~\ref{prop:certifiedeell}(2)--(6) use of $C^{\bullet}_\ell$ only: (i) $C^{q}=\bigoplus_xM^{Q_x}_{\mathsf O_\ell}$, images of idempotents, free over the principal ideal domain
$\mathsf O_\ell$ and commuting with base change --- (1)(c) here; (ii) $H^{q}(C^{\bullet}_\ell)=H^{q}(\Ga;M)\otimes\mathsf O_\ell$, which there comes from a free model and flatness --- here
(1)(c) over $\mathsf O_\ell$, and $C^{\bullet}_{\mathcal T}(M)\otimes\mathsf O_\ell=C^{\bullet}_{\mathcal T}(M_{\mathsf O_\ell})$ ($\Iso_{\mathcal T}$ finitely generated free) with $\mathsf O_\ell$ flat; (iii) the
statements (4)(a), (4)(b) about $H^{\bullet}(\Ga;M)$, which do not see the model; (iv) for (5), $a_{yx}=e_ya_{yx}e_x$ for every coefficient: for the lifts $a_t$, $b_t$ this is (1)(a), and for the column differentials $(1-g_x)e_x$, $N'_x$ and the edge's
$-\varepsilon_t$ it holds because $e_x$ is central in the commutative ring $\cring[\Ga_x]$ and idempotent; (v) Smith forms over
$\mathsf O_\ell$ and $\Z_{(\ell)}$ and the universal coefficient theorem. All hold for the tree model, and (4)(c)'s exclusion of other complexes is by its own argument (the
rank of a differential is a property of that differential).

(2) Dualise (1)(c): for $a\in e_y\cring[\Ga]e_x$, $\Hom(r_a,M)\colon M^{Q_x}\to M^{Q_y}$ is $m\mapsto am$ ($c(e_ya)=c(ae_x)=a\,c(e_x)$). So
$d^{q}(u,w)$ is $\delta^{x}_{q+1}u_x$ on $F^{x}_{q+1}$ and $a_qu_S+b_qu_U-\varepsilon_qw$ on $F^{\mathfrak e}_q$ (the sign of $-\partial^{\mathfrak e}$). Substituting the coefficients and
$e_Uy=y$ for $y\in M^{U^{2}}$, $e_{\mathfrak e}y=y$ for $y\in M^{-I}$, $e_Sx=x$ for $x\in M^{\langle S\rangle}$ gives (a)--(c); at $\ell=3$ the columns $F^{S}$, $F^{\mathfrak e}$ live in degree
$0$ only, which removes every term but those displayed. On $M^{-I}$, $U$ and $U^{-2}=\mathbf c^{-1}U$ agree, $U^{-2}$ generates $\langle U^{2}\rangle$, and $1+U+U^{2}$ is its
norm; $\operatorname{im}d^{1}=(1+U+U^{2})M^{-I}$, so the cohomology in degrees $\ge2$ is that of the periodic resolution of $\langle U^{2}\rangle$ \STD{\cite[Ch.~III.6.2]{brown}} on
$M^{-I}$ (the case $q=2$ included, its incoming image being the norm's). (d) is inspection of (a), (b): $H^{q}$ uses $d^{q}$ and $d^{q-1}$, both periodic once $q-1\ge1$
at $\ell=2$; at $\ell=3$ once $q-1\ge2$, and $H^{2}$ has $\operatorname{im}d^{1}=\operatorname{im}d^{3}$. (e) is (1)(c) read on (a)--(c).

(3) \emph{Column maps.} $f^{x}$: for odd $t$, $(1-g_x)\cdot e_x=e_x\cdot(1-g_x)e_x$ with equal scalars ($\lfloor t/2\rfloor=\lfloor(t-1)/2\rfloor$); for even $t$,
$N_x\cdot\abs{Q_x}^{t/2-1}e_x=\abs{Q_x}^{t/2}e_x\cdot N'_x$ by $N_xe_x=\abs{Q_x}N'_x$; where the $\ell$-local column vanishes in degree $t\ge1$, $f^{x}_{t-1}\partial_t=r_{(1-g_x)e_x}=0$ as
$g_x\in Q_x$. $g^{x}$: the same identities with the inverse scalars ($\abs{Q_x}^{-1}\in\cring$). Both preserve the augmentations. $f^{x}g^{x}=\id$ column by column ($e_x^{2}=e_x$,
the scalars cancel). \emph{The corrections.} By ($\ast$) $\Phi^{\mathcal T}_\ell$ is a chain map iff $\kappa_t\delta^{U}_{t+1}+\varepsilon_t\kappa_{t-1}=D_t:=b_t\,f^{U}_t-f^{\mathfrak e}_t\,b'_t$
(coefficients in $F^{U}_t$; the $S$-components of $D$ vanish: at $\ell=2$ because $f^{S}=f^{\mathfrak e}=\id$ and the $a_t$ agree, at $\ell\ge3$ because $f^{S}_t=0$ for $t\ge1$ and
$e_{\mathfrak e}e_S=e_S$ at $t=0$), $b_t$ the free and $b'_t$ the local lift, $\kappa_{-1}:=0$. At $\ell=2$: $D_t=(1-3^{t/2})e_U$ ($t$ even) and $(1-3^{(t-1)/2}(2+\mathbf c))e_U$ ($t$ odd);
with $\kappa_{2j-1}=\kappa_{2j}=-\mathsf k_je_U$, $t=2j+1$ gives $-\mathsf k_{j+1}(1+\mathbf c)-\mathsf k_j(1-\mathbf c)=1-3^{j}(2+\mathbf c)$ (as $\mathsf k_{j+1}+\mathsf k_j=2\cdot3^{j}-1$,
$\mathsf k_{j+1}-\mathsf k_j=3^{j}$), and $t=2j+2$ gives $-\mathsf k_{j+1}(1-\mathbf c)-\mathsf k_{j+1}(1+\mathbf c)=1-3^{j+1}$, both times $e_U$; $t=0$: $D_0=0=\kappa_0$. At $\ell=3$:
$D_0=0$, $D_t=-2^{(t-1)/2}(1+U+U^{2})e_{\mathfrak e}$ ($t$ odd), $-2^{t/2}e_{\mathfrak e}$ ($t\ge2$ even), and $\delta^{U}_{t+1}$ is $(1+U+U^{2})e_{\mathfrak e}$, resp.\ $(1-U)e_{\mathfrak e}$; for odd
$t$, $\kappa_{t-1}=0$ and $\kappa_t(1+U+U^{2})e_{\mathfrak e}=D_t$; for even $t\ge2$, $\kappa_t=0$ and $\varepsilon_t\kappa_{t-1}=-(1+\mathbf c)2^{(t-2)/2}e_{\mathfrak e}=-2^{t/2}e_{\mathfrak e}=D_t$. At $\ell\ge5$, $D=0$ and $k=0$. For $\Psi^{\mathcal T}_\ell$ the identity is
$\kappa'_t\delta_{t+1}+\varepsilon_t\kappa'_{t-1}=D'_t:=b'_t\,g_t-g^{\mathfrak e}_t\,b_t$ with free target coefficients $\delta$. At $\ell=2$, $D'_t$ has $S$-component $0$ and $U$-component
$1-3^{-t/2}e_U$ ($t$ even), $(1+U+U^{2})-3^{-(t-1)/2}e_U$ ($t$ odd); split $\cring[\langle U\rangle]=\cring[\langle U\rangle]e_U\times\cring[\langle U\rangle](1-e_U)$. On the second factor
$N_U=3(1+\mathbf c)e_U$ vanishes and $\mathsf w(1-U)=1-e_U$ (as $(1-U)(2+2U+U^{2}+U^{3})=2-U^{2}-U^{4}=3(1-e_U)$ and $(1-U)(1+\mathbf c)e_U=(1-\mathbf c)(1+\mathbf c)e_U=0$), so the
identity reads $\mathsf w(1-U)=1-e_U$ for even $t$ and $(1-\mathbf c)\mathsf w=(1+U+U^{2})(1-U)\mathsf w=(1+U+U^{2})(1-e_U)$ for odd $t$; on the first factor, with $\kappa'_t=\lambda_te_U$,
$\lambda_{2j-1}=\lambda_{2j}=\frac{1-3^{-j}}2$, it reads $3(1+\mathbf c)\lambda_{2j-1}+(1-\mathbf c)\lambda_{2j-2}=2+\mathbf c-3^{-(j-1)}$ and $(1-\mathbf c)\lambda_{2j}+(1+\mathbf c)\lambda_{2j-1}=1-3^{-j}$,
both direct. At $\ell=3$ and $\ell\ge5$ the only condition is in degree $t=0$, $\kappa'_0(1-S)=e_S-e_{\mathfrak e}$ and (at $\ell\ge5$) $\kappa'_{0,U}(1-U)=e_{\mathfrak e}-e_U$, and
$-\frac12e_{\mathfrak e}(1-S)=-\frac14(1-S+S^{2}-S^{3})=e_S-e_{\mathfrak e}$, $\frac13e_{\mathfrak e}(2+2U+U^{2}+U^{3})(1-U)=e_{\mathfrak e}(1-e_{\langle U^{2}\rangle})=e_{\mathfrak e}-e_U$.
\emph{$\Phi\Psi=\id$.} $\Phi^{\mathcal T}_\ell\Psi^{\mathcal T}_\ell(a,z)=(fga+(fk'+kg)z,\,fgz)=(a+(fk'+kg)z,\,z)$. At $\ell=2$, $f^{U}_{t+1}k'_t=r_{3^{\lceil t/2\rceil}\kappa'_te_U}=
r_{\mathsf k_{\lceil t/2\rceil}e_U}$ ($\mathsf we_U=0$) and $k_tg^{\mathfrak e}_t=r_{-\mathsf k_{\lceil t/2\rceil}e_U}$; at $\ell=3$, $k'$ lands in $F^{S}_1$, where $f^{S}_1=0$, and $k_0=0$; at
$\ell\ge5$, $f_1=0$ and $k=0$. \emph{$\Psi\Phi\simeq\id$}: $\Psi^{\mathcal T}_\ell\Phi^{\mathcal T}_\ell$ and $\id$ are chain maps over $\id_\cring$ from the projective (free) complex
$\Iso_{\mathcal T}\otimes\cring$ to the resolution $\Iso_{\mathcal T}\otimes\cring$, hence homotopic \STD{\cite[Thm.~2.2.6]{weibel}}; an explicit homotopy is $\mathsf H_\ell$ of (4)(c),
since $\Psi^{\mathcal T}_\ell\Phi^{\mathcal T}_\ell$ is over the identity of $C_{\bullet}(\mathcal T)\otimes\cring$ (the column maps preserve the augmentations and $k$, $k'$ land in column degrees $\ge1$). The cochain statements are $\Hom(-,M)$ of these; the
cochain formulas are the dual of (2)'s kind, $(F^{\ast}c)(y)=\sum_xb_{yx}c(x)$.

(4) (a) \emph{Well defined}: by Lemma~\ref{lem:model7coincidences}(1) the stabilisers are $\Ga_{v_6}=\langle S\rangle$, fixing $\mathfrak v_S$, $\Ga_{v_4}=\langle ST\rangle=\langle U\rangle$,
fixing $\mathfrak v_U$, and $\{\pm I\}$ for every other representative, acting trivially on $C_{\bullet}(\mathcal T)$, on $C_{\bullet}(\Xbulk)$'s orientations and on the tables; so
$\mathsf i_{\mathcal T}$, $\mathsf r_{\mathcal T}$, $\mathsf G_{\mathcal T}$ extend $\Z[\Ga]$-linearly. The subcomplex is Claim~\ref{cl:amalgam}(1)'s tree, whose vertices and edges lie on the unit arc
of $D_Y\subset\Xbulk$, and $\Ga_{\{v_4,v_6\}}=\{\pm I\}$; so $\mathsf i_{\mathcal T}$ is a bijection of the bases onto that subcomplex's, and a chain map since
$\partial\{v_4,v_6\}=v_6-v_4$. The faces of the representatives that are not representatives are $v_3=Tv_0$, $v_5=Tv_4$, $\{v_3,v_5\}=T\{v_0,v_4\}$ and
$\{v_5,v_6\}=S^{-1}\{v_4,v_6\}$ (Lemma~\ref{lem:model7coincidences}(2)), so $\mathsf r_{\mathcal T}(v_3)=\mathsf r_{\mathcal T}(v_5)=T\mathfrak v_U=S\mathfrak v_U$ ($S^{-1}T=-U\in\langle U\rangle$),
$\mathsf r_{\mathcal T}\{v_3,v_5\}=0$, $\mathsf r_{\mathcal T}\{v_5,v_6\}=S\mathfrak e$, $\mathsf G_{\mathcal T}(v_3)=-\{v_3,v_5\}$, $\mathsf G_{\mathcal T}(v_5)=0$, $\mathsf G_{\mathcal T}\{v_3,v_5\}=\mathsf G_{\mathcal T}\{v_5,v_6\}=0$.
\emph{Chain map}: on edges $\partial\mathsf r_{\mathcal T}(\{a,b\})=\mathsf r_{\mathcal T}(b)-\mathsf r_{\mathcal T}(a)$ is read off the table (e.g.\ $\partial(-S\mathfrak e)=S\mathfrak v_U-\mathfrak v_S=
\mathsf r_{\mathcal T}(v_3)-\mathsf r_{\mathcal T}(v_2)$); on triangles $\mathsf r_{\mathcal T}(\partial\Delta)=0$, e.g.\ $\partial\{v_2,v_3,v_5\}\mapsto0-(-S\mathfrak e)+(-S\mathfrak e)$,
$\partial\{v_0,v_1,v_4\}\mapsto-\mathfrak e-0+\mathfrak e$. $\mathsf r_{\mathcal T}\mathsf i_{\mathcal T}=\id$ on $\mathfrak v_S$, $\mathfrak v_U$, $\mathfrak e$ by the table. \emph{Homotopy}: on vertices
$\partial\mathsf G_{\mathcal T}(v)=v-\mathsf i_{\mathcal T}\mathsf r_{\mathcal T}(v)$ (e.g.\ $\partial(-\{v_0,v_4\})=v_0-v_4$); on edges, e.g.\
$\partial(\{v_0,v_1,v_4\}-\{v_1,v_4,v_6\})+\mathsf G_{\mathcal T}(v_1)-\mathsf G_{\mathcal T}(v_0)=\{v_0,v_1\}-\{v_4,v_6\}$ and
$\partial(\{v_2,v_3,v_5\}+\{v_2,v_5,v_6\})+\mathsf G_{\mathcal T}(v_3)-\mathsf G_{\mathcal T}(v_2)=\{v_2,v_3\}+\{v_5,v_6\}$, i.e.\ $x-\mathsf i_{\mathcal T}\mathsf r_{\mathcal T}(x)$; on triangles
($C_3=0$) $\mathsf G_{\mathcal T}(\partial\Delta)=\Delta$, e.g.\ $\mathsf G_{\mathcal T}(\partial\{v_2,v_3,v_5\})=0-\{v_2,v_5,v_6\}+(\{v_2,v_3,v_5\}+\{v_2,v_5,v_6\})$. The other
cases are the same bookkeeping (all $19$ are checked by \file{cd6\_iso\_check.py}).
(b) $\Z[\Ga]$ is free as a right $\Z[\langle g\rangle]$-module on $\mathsf R_\sigma$, so the column is $\bigoplus_{\rho}\rho\otimes F$ over $\Z$, and $\mathsf h$ is the standard contraction
of each summand: $(g-1)\sum_{i<j}g^{i}=g^{j}-1$ gives $\partial\mathsf h+\mathsf h_{-1}\aug=\id$ in degree $0$; in odd degrees $N[j=m-1]+\mathsf h((g-1)g^{j})=g^{j}$ (for
$j<m-1$ the second term is $g^{j}$, for $j=m-1$ it is $-\sum_{i<m-1}g^{i}=g^{m-1}-N$); in even degrees $\ge2$, $(g-1)\sum_{i<j}g^{i}+\mathsf h(Ng^{j})=g^{j}-1+1$. Negating the even
$\mathsf h$ for $r_{1-g}$, and multiplying by $\epsilon$ for a global sign, preserves these identities. $\mathsf h_p$: no component of $\partial$ raises $s$ (Definition~\ref{def:isomodel};
the cone), so the top component $z_s$ of a cycle is a cycle of its columns, in column degree $0$ in the kernel of their augmentation (since $pz=0$ and $p$ is the direct sum
of the column augmentations), whence $z_s=\partial^{\mathrm{col}}\epsilon\mathsf h(z_s)$ and $z-\partial k_s$ is a cycle of $\ker p$ without component on $s$ or above.
(L) is the proof of Lemma~\ref{lem:strictlifting} with $k:=\mathsf h_p(z)$. (L$'$): by induction $pH=g\,p'$ in lower degrees, so with $w:=Ax-Bx-H(\partial x)$,
$p(w-\partial y_1)=(a-b)p'x-g\,p'\partial x-\partial g\,p'x=0$, and $\partial(w-\partial y_1)=(A-B)\partial x-\partial H\partial x=H\partial\partial x=0$ by induction; then
$\partial H(x)=w$, which is $\partial H+H\partial=A-B$ on $x$, and $pH(x)=py_1=g\,p'x$. (c) $\varepsilon_{\mathcal T}\Phi^{\Z}_{\mathcal T}\Psi^{\Z}_{\mathcal T}=\mathsf r_{\mathcal T}\cmpI\Psi^{\Z}_{\mathcal T}=
\mathsf r_{\mathcal T}\mathsf i_{\mathcal T}\varepsilon_{\mathcal T}=\varepsilon_{\mathcal T}$ and $\cmpI\Psi^{\Z}_{\mathcal T}\Phi^{\Z}_{\mathcal T}=\mathsf i_{\mathcal T}\mathsf r_{\mathcal T}\cmpI$, so (L$'$) applies as
stated, with $g=0$, resp.\ $g=\mathsf G_{\mathcal T}$ by (a); $\Psi^{\mathcal T}_\ell\Phi^{\mathcal T}_\ell$ is over the identity of $C_{\bullet}(\mathcal T)\otimes\cring$ (the proof of (3)). Both
$\Iso_{\bullet}$ and $\Iso_{\mathcal T}$ are bounded-below free resolutions of $\Z$ (Proposition~\ref{prop:isomodel}(2), Remark~\ref{rmk:treeconeres}); uniqueness up to homotopy of comparisons over
$\id_\Z$: the argument of Proposition~\ref{prop:isomodel}(5) (a chain map from a projective bounded-below complex into the acyclic kernel of an augmentation is null-homotopic).
(d) $\Phi_{\mathcal T}\Psi_{\mathcal T}-\id=\Phi^{\mathcal T}_\ell(\Phi^{\Z}_{\mathcal T}\Psi^{\Z}_{\mathcal T}-\id)\Psi^{\mathcal T}_\ell+(\Phi^{\mathcal T}_\ell\Psi^{\mathcal T}_\ell-\id)$, the second term $0$ by (3) and
the first $\partial\mathsf K+\mathsf K\partial$ as $\Phi^{\mathcal T}_\ell$, $\Psi^{\mathcal T}_\ell$ are chain maps; $\id-\Psi_{\mathcal T}\Phi_{\mathcal T}=(\id-\Psi^{\Z}_{\mathcal T}\Phi^{\Z}_{\mathcal T})-
\Psi^{\Z}_{\mathcal T}(\Psi^{\mathcal T}_\ell\Phi^{\mathcal T}_\ell-\id)\Phi^{\Z}_{\mathcal T}=\partial\mathsf K'+\mathsf K'\partial$. The cochain statements are $\Hom(-,M)$ of these; each comparison is
compatible with the identification $H^{q}(\Hom(\mathsf P,M))=\Ext^{q}(\cring,M)$, which is itself computed through comparisons. A comparison through $\Iso^{(\ell)}$
($\Phi_\ell$, $\Psi_\ell$ of Proposition~\ref{prop:isomodelell}(4), with any comparison $\Iso^{(\ell)}\rightleftarrows\Iso^{(\ell)}_{\mathcal T}$ between projective resolutions) is a comparison
$\Iso_{\bullet}\otimes\cring\rightleftarrows\Iso^{(\ell)}_{\mathcal T}$ again, so homotopic to $\Phi_{\mathcal T}$, $\Psi_{\mathcal T}$; Proposition~\ref{prop:isomodelell}(4)'s proof (strict lifts
along $\cring$-linear contractions, $e_\sigma$-averaged) needs of $\cring$ only $\abs{Q_\sigma}\in\cring^{\times}$ and the exactness of $C_{\bullet}(\Xbulk)\otimes\cring$.

(5) (T1): $\mathsf q_2=3$, $\mathsf q_3=4$, so $\ell\nmid\mathsf q_\ell$ and $\mathsf O_\ell\subset\Z_{(\ell)}$; (1)(b), (1)(d). (T2): (1)(b), (1)(c), with
$M^{Q_x}_{\Z_{(\ell)}}$ free over the local principal ideal domain $\Z_{(\ell)}$. (T3): (4). Proposition~\ref{prop:descentmaps}(5) then gives (1), (2), (4) there, whose proofs use
nothing else of a model. The source (i) of (2)(d) there is read on a model's integral form; here that form is $\Iso^{(\ell)}_{\mathcal T,\mathsf O_\ell}$ by (1)(d), and a prime
$\ell'\nmid\mathsf q_\ell$ is a prime of $\mathsf O_\ell$, as in Proposition~\ref{prop:certifiedeell}(4). The last sentence is Proposition~\ref{prop:descentmaps}(1)(c), (2)(c), which allow
any $\Z_{(\ell)}$-model for the $\ell$-local pieces.

(6) (a) $\Phi^{\ast}_{\mathcal T}$, $\Psi^{\ast}_{\mathcal T}$ are cochain maps inducing the comparison isomorphisms, mutually inverse on cohomology (4); conjugating a cochain map that
induces $t$ by them induces $t$ again (the argument of Proposition~\ref{prop:isomodelell}(5)). Homotopic comparisons give homotopic $\Phi^{\ast}$, $\Psi^{\ast}$, and composition
respects homotopy. (b) $H(\vartheta^{\ast})$ is the comparison isomorphism between the two identifications with $H^{\bullet}(\Ga;M_\cring)$, and both $T''$ and $T'_{\mathcal T}$
induce $t$ through them, by (a) and Proposition~\ref{prop:isomodelell}(5) (over every $\mathsf O_\ell$-algebra: its proof uses the resolution-independence of $\cor\circ(\conj_g)_\ast\circ\res$
and the comparisons of (4) there, nothing about $\ell$). (c) $T'_{1}T'_{2}=\Psi^{\ast}T_1(\Phi^{\ast}\Psi^{\ast})T_2\Phi^{\ast}$ and $\Phi^{\ast}_{\mathcal T}\Psi^{\ast}_{\mathcal T}=
(\Psi_{\mathcal T}\Phi_{\mathcal T})^{\ast}\simeq\id$ (4). (d) Proposition~\ref{prop:descentmaps}(1)(a) for $\Phi_{\mathcal T}$ and $\Psi_{\mathcal T}$:
$uT'_{\mathcal T}=\Psi^{\ast}uT\Phi^{\ast}=\Psi^{\ast}Tu\Phi^{\ast}=T'_{\mathcal T}u$; the cone is the computation in the proof of Proposition~\ref{prop:descentmaps}(3)(c).

\emph{A check outside the build, not a tag} (v1.23.0-pre.20): \file{docs/implementation/slot24/cd6\_tree\_check.py}, independent of \code{cohomology\_app}, of
\file{descentN\_oracle.py} and of \file{integral\_line\_probe.py} (its own group ring $\Q[\Ga]$, its own \emph{induced} module $M_N$ on left cosets by first column --- the
oracle's is coinduced on right cosets by bottom row ---, the four models built from Definition~\ref{def:isomodeltree} as chain data, $\ell$-adic Smith forms mod $\ell^{K}$
certified by exact ranks from the Betti recursion; the oracle tables are read as data only), output \file{cd6\_tree\_check.out} (its header records the command). It checks, exactly in
$\Q[\Ga]$ to degree $8$, (1)(a) and (3)'s chain-map identities and $\Phi^{\mathcal T}_\ell\Psi^{\mathcal T}_\ell=\id$ at $\ell=2,3,\ge5$; and on $145$ cells --- level $1$, $0\le n\le30$ and
$n\in\{40,56\}$, and $\Ga_0(N)$ for $2\le N\le16$, $N\ne15$, $n\le8$, $\psi(N)(n+1)\le120$, and $(N,n)\in\{(25,2),(49,2)\}$ --- at $\ell=2$ and $3$: (2)'s closed forms against the chain data entry by entry, (2)(e)'s ranks, the groups
$H^{q}$, $q\le5$, of the $\ell$-local and of the free tree model against the $\ell$-parts of \file{descentN\_oracle.tsv} and \file{descent\_oracle\_n128.tsv}, the certification and
the a-priori bound of (1)(d), $\Psi^{\mathcal T\ast}_\ell\Phi^{\mathcal T\ast}_\ell=\id$ on cochains, and $\Phi^{\mathcal T\ast}_\ell\Psi^{\mathcal T\ast}_\ell$ the identity on $H^{q}(-\otimes\Z/\ell^{K})$ (on every cell by that left inverse and equal finite lengths, and directly on generators of the cocycles, $1\le q\le4$, where $3\psi(N)(n+1)\le90$); at
$\ell\in\{5,7\}$ the groups on the Bredon tree; and on $18$ level pairs $\Ga_0(MN)\subset\Ga_0(N)$ the strict commutation of the fibre maps $\alpha$, $\omega$ with
$\Phi^{\mathcal T\ast}_\ell$, $\Psi^{\mathcal T\ast}_\ell$ and the equality of the level-cone groups (degrees $-1$ to $4$) on the $\ell$-local and the free tree model: $49\,303$ checks,
$0$ mismatches (\code{CD6\_CHECK \dots\ mismatch=0 \dots\ verdict=OK}). It does \textbf{not} build $\Iso$. Its companion (v1.23.0-pre.21)
\file{docs/implementation/slot24/cd6\_iso\_check.py} (output \file{cd6\_iso\_check.out}) does: it builds an $\Iso_{\bullet}$ to degree $9$ by Proposition~\ref{prop:isomodel}(1)'s procedure with the
representatives of (4)(b), and checks exactly in $\Z[\Ga]$, resp.\ $\Q[\Ga]$, to degree $8$: $d\circ d=0$, $\cmpI$ a chain map, $w_q=5,14,19,\dots$; (4)(a) on all $19$ representatives;
$\Phi^{\Z}_{\mathcal T}$, $\Psi^{\Z}_{\mathcal T}$ chain maps over $\mathsf r_{\mathcal T}$, $\mathsf i_{\mathcal T}$; the identities of $\mathsf H_{\mathcal T}$, $\mathsf H_X$, and of $\mathsf H_\ell$, $\mathsf K$,
$\mathsf K'$ for $\ell\in\{2,3,5\}$; and on $10$ cells ($\Ind_{\Ga_0(N)}^{\Ga}\Sym^{n}$, $N\le5$, $n\le6$) mod $p\in\{2,3,5\}$, degrees $\le6$, equal dimensions of $H^{q}(\CI)$ and $H^{q}(C^{\bullet}_{\mathcal T})$ and
$\Phi^{\Z\ast}_{\mathcal T}$, $\Psi^{\Z\ast}_{\mathcal T}$ mutually inverse on them: $12\,870$ checks, $0$ mismatches (\code{CD6\_ISO\_CHECK \dots\ mismatch=0 \dots\ verdict=OK}); four
mutations of the tables (a sign of $\mathsf r_{\mathcal T}$, a value of $\mathsf r_{\mathcal T}$, of $\mathsf G_{\mathcal T}$, of $\mathsf i_{\mathcal T}$) are each caught. (6) is not exercised (no
cochain Hecke operator on $\Iso$ is built there).
\end{proof}

\noindent\emph{Code map} (v1.23.0-pre.20, \textbf{stated before D5's code}, contract first, rule~21; \textbf{updated at v1.23.0-pre.30}: D5 is \textbf{built and merged}, merge \texttt{5eb7e9de};
source \file{docs/implementation/slot24/DESCENT5\_report.md}, including its \S8). \textbf{Built.} The files are \file{Integral/Model/TreeCochains.\{h,cpp\}} (the $\ell$-local tree
cochains $C^{\bullet}_{\mathcal T,\ell}(M)$ of (2) in closed form, block by block, for any module given by its block-monomial action), \file{Integral/Model/IsotropyModelTree.\{h,cpp\}} (the free tree
model and the canonical $\ell$-local models of Definition~\ref{def:isomodeltree}; $\Phi^{\mathcal T}_\ell$, $\Psi^{\mathcal T}_\ell$ of (3); and (4) made effective on the code's $\Iso$: the tables of (4)(a), the
solvers $\mathsf h_p$, $\Phi^{\Z}_{\mathcal T}$, $\Psi^{\Z}_{\mathcal T}$ by (L), $\mathsf H_{\mathcal T}$, $\mathsf H_X$ by (L$'$), $\mathsf H_\ell$, $\mathsf K$, $\mathsf K'$), \file{Integral/Descent/DescentTree.\{h,cpp\}}
(\code{RunTreeLocalPiece}: the $\ell$-primary torsion of $C^{\bullet}_{\mathcal T,\ell}(M)\otimes\Z/\ell^{K}$ by the certified-exponent rule of (1)(d)) and
\file{Integral/Descent/TreeComparisonArms.\{h,cpp\}} (the validations' arms); the validations are \code{HeckeIsotropyModelTreeValidation} (standard) and
\code{HeckeIsotropyModelTreeSlowValidation} (slow), registered in \file{Registry/Slot24/DESCENT5.cpp}, sources \file{cmake/slot24/DESCENT5.cmake}. \textbf{The tree is the default route for the dense
$\ell$-local bulk pieces}: D0's $2$- and $3$-adic pieces, D1's $2$-adic piece and its dense $3$-adic piece at $\nu_3(N)>0$, and D2's dense pair pieces at $2$ and (for $\nu_3>0$) at $3$;
\textbf{rollback} \code{HECKE\_DESCENT\_TREE=0} (read-only, value-preserving; read once per call) returns to the dense route. Unchanged: the $\nu_3(N)=0$ shortcut at $3$
(Proposition~\ref{prop:descentNthree}) and every $\ell\ge5$ route (the sparse route is cheaper than a dense tree there); the Bredon tree at level $1$, $\ell\ge5$, is built and
checked against the direct Smith form of $C_{\mathrm{eq}}$ at every Dickson prime, but is not the production route. $\Phi^{\Z}_{\mathcal T}$, $\Psi^{\Z}_{\mathcal T}$ and the homotopies are (4)(b)--(c) on the code's $\Iso$ (\file{Integral/Model/IsotropyModel},
whose column contractions and signed-lexicographically least coset representatives, \code{decompose()} and \code{column\_homotopy} of \file{PeriodicColumn.h}, are those of (4)(b)). Both routes agree with the oracles
on every cell where both run, and D0's, D1's and D2's validations pass unchanged on both settings of the switch. \textbf{Hecke by transport} ((6)): the hook \code{TransportToTree} ($\mathsf T'_{\mathcal T}=\Psi^{\ast}_{\mathcal T}\,T\,\Phi^{\ast}_{\mathcal T}$ on any
cochain endomorphism of $C_{\mathrm I}$) \textbf{exists} and is smoke-tested; its \textbf{real arm}, ``the transported $T'_p$ equals the one on $\Iso^{(2)}$'' ((6)(b)), is \textbf{pending D3}'s level-$N$ cochain operators (C-D4). (v1.23.0-pre.31: that arm is built with D3 on branch \texttt{s24-d3impl}, pending merge and audit --- the code map of \S\ref{sec:descenthecke}; v1.23.0-pre.33: audited and merged at \texttt{dc8261a3}.) \textbf{Not built here}: the tree cusp pair
(C-D8, D5c; contract first at v1.23.0-pre.33, \S\ref{sec:treepair}; v1.23.0-pre.36: built, audited and merged at \texttt{4064e164}, \S\ref{sec:treepair}'s code map); at v1.23.0-pre.30 D4's cone was the pair $(\Iso,\Isob)$ and the tree had no boundary, so the tree bulk did not move D4's
wall (since v1.23.0-pre.36 the tree cusp pair is the default $2$- and $3$-adic pair route, \file{Integral/Descent/DescentTreePair}; v1.23.0-pre.38, referee m9). The free tree resolution in closed form remains D0's and D1's comparator (c)
(\file{DESCENT0\_report.md}, \file{DESCENT1\_report.md}).

\paragraph{Formalization note.} Definition~\ref{def:isomodeltree} and Proposition~\ref{prop:isomodeltree} are homological algebra over $\cring[\Ga]$ (mapping cones, the comparison
theorem, Shapiro's lemma) and identities in the group rings of $\Z/4$ and $\Z/6$; the free tree resolution they build on carries three \code{\textbackslash LEAN} markers
(Remark~\ref{rmk:treeconeres}). They carry \textbf{no} \code{\textbackslash LEAN} marker.

% NEW at v1.23.0-pre.24 (hecke-extensions slot 24, milestone D3 of docs/implementation/slot24/DESCENT_plan.md v0.4, its clauses C-D4 and C-D5
% (thm:descentcert, as restricted by the audit's M5); the owner's rulings [R-Q2] (no bignum: residues + a certified bound), [R-Q3] (invariants only),
% [R-Q4] (T_p, p not dividing N; U_p deferred), [R-GT] (the tree is the 2-adic route)).  Placed after sec:isomodeltree, before sec:mpmodel, so
% that no statement number moves.  A pre-release draft, NOT refereed.  No \LEAN marker.  C-D4' (a strict lift directly on Iso^(3)) is not drafted.
% v1.23.0-pre.25: refereed ACCEPT-with-fixes (docs/implementation/slot24/C-D4_referee.md); its fixes M1, M2 (resolution (a): the l = 3 piece is the
% transport on Iso^(3) for every N, the shipped Bredon operator is never read at l in {2,3}), m1-m5, n1, t1-t3 applied; NOT re-refereed.
\subsection{Hecke operators \texorpdfstring{$T_p$}{T\_p}, \texorpdfstring{$p\nmid N$}{p not dividing N}, on the descent over \texorpdfstring{$\Z$}{Z}: the operator at level \texorpdfstring{$N$}{N}, the characteristic polynomial, the torsion}\label{sec:descenthecke}

(v1.23.0-pre.24; milestone D3 of \file{docs/implementation/slot24/DESCENT\_plan.md} v0.4, its clauses C-D4 and C-D5 as restricted by the audit's M5; the owner's rulings
of 2026-10-24: residues and a certified bound, no big integer [R-Q2]; Hecke output as invariants only --- characteristic polynomials over $\Z$ and data on certified
groups, integer Hecke matrices later (D8) [R-Q3]; $T_p$ for $p\nmid N$ only, every $U_p$ deferred [R-Q4] (lifted for the free model by Proposition~\ref{prop:strictliftN}(9), v1.26.0-pre.1); and of 2026-10-25: the tree is the $2$-adic route [R-GT].)
\S\ref{sec:descentN} glues the \emph{groups} $H^{q}(\Ga_0(N);\Vcoef)$ and the maps induced by fibre maps; Proposition~\ref{prop:descentN}(3) says that a Hecke operator
on them is determined by its two local components, but no cochain operator at level $N$ exists on a model that computes the $2$- and $3$-adic parts, and
Proposition~\ref{prop:descentmaps}(3)(c) leaves the $2$- and $3$-adic level sequence without a Hecke action for that reason. This subsection fixes: (i) the level-$N$
operator $T_p$, $p\nmid N$, at cochain level on the free model $\Iso_{\bullet}$, built from a bulk strict lift that exists \emph{and is constructed} for every $p$, its
compatibility with the glueing map $\cmpI^{\ast}$ ($=\varepsilon^{\ast}$ of the plan), its transports to $\Iso^{(\ell)}$ and to the tree models --- the route at $2$ ---,
the cheap routes at $\ell\ge5$ and at $3$ (the latter for an integral datum only, and not used by the record; v1.23.0-pre.25), and its strict commutation with the restriction of D2 (Proposition~\ref{prop:strictliftN}); (ii) the characteristic polynomial
over $\Z$ at level $N$ from residues and a certified bound, with the output rule for a coefficient beyond the word (Proposition~\ref{prop:crtcharpolyN}); (iii) which
invariants of an endomorphism of a finite $\ell$-group are choice-free, and how they are read exactly from a local run (Proposition~\ref{prop:torsionhecke}); (iv) the
record those make certified (Theorem~\ref{thm:descentcert}, the plan's C-D5), and what it does not assert.

\emph{What is used.} The Hecke clauses of this document that the subsection rests on, and nothing else of the Hecke chapters: Definition~\ref{def:tau}
(the $g$-twisted chain datum), Proposition~\ref{prop:strictpair} with Lemma~\ref{lem:taubdryexists}(1) (an integral datum on $\Xbulk$ for every $g\in\De$),
Lemma~\ref{lem:strictlifting} and the bulk half of Proposition~\ref{prop:strictlift} and of its proof, Remark~\ref{rmk:strictliftbuilt}(1), (2) (the construction and its
certificates, read for every $p$), Lemma~\ref{lem:cosettransport}, Proposition~\ref{prop:levelcorrespondence} (formula \eqref{eq:levelcorr}, items (1)--(4)), Proposition~\ref{prop:levelconehecke}(1),
(2), (4) (the criterion, its instance for $T_p$, $p\nmid N_B$, and the independence of the transversal; (4) with $\tau_{\Iso}$ in place of $\tau$, its proof using (T1)
only, v1.23.0-pre.25, referee t3), Proposition~\ref{prop:relationsmodel} (the level-$1$ relations on the free
model, and its citations), Propositions~\ref{prop:isomodelell}(4), (5) and \ref{prop:isomodeltree}(3), (4), (6) (comparisons and transport), Propositions~\ref{prop:descentN}(3),
(5) (local determination, periodicity) and \ref{prop:descentmaps}(1), (2)(g), (3) (fibre maps, the precision rule, the Hecke squares), and Proposition~\ref{prop:crtcharpoly}
(level~$1$); at v1.23.0-pre.31 also Propositions~\ref{prop:levelorbit}(1), (2) and \ref{prop:heilbronn}(1), (2) with Definition~\ref{def:kappaN} (the cokernel presentation at $\ell\ge5$
and at the residue primes, with $6$ inverted) and Proposition~\ref{prop:descentpairhecke}(1) (the pair lift, one admissible bulk lift); at v1.23.0-pre.32 (referee N6), for completeness of D3's proofs,
Proposition~\ref{prop:isomodel}(2) ($\cmpI$ onto, the check through the lift), Proposition~\ref{prop:descentN}(1), (7)(b) (flat base change and $b_1$, the certifying rank) and
Lemma~\ref{lem:levelduality} (through Proposition~\ref{prop:heilbronn}); at v1.26.0-pre.1 (item~(9) of Proposition~\ref{prop:strictliftN}; v1.26.0-pre.2, referee S2) also
Proposition~\ref{prop:levelconehecke}(3)(d) (the adjoint convention) and Proposition~\ref{prop:descentmaps}(3)(b$'$) (within its (3) above), and Proposition~\ref{prop:heckeZspec}(0)
is named by (9)(d); Definition~\ref{def:heckeZ} and Proposition~\ref{prop:heckeZspec}, of a later chapter that rests on this one, are cited by item~(9) only to \emph{identify}
the operator and to point at a relation, not used. \emph{Not used}: Theorems~\ref{thm:relations}, \ref{thm:classical} at $\ell\in\{2,3\}$, the shipped operator on the Bredon object at $\ell\in\{2,3\}$
(Remark~\ref{rmk:bredonhecke}: never read here --- Proposition~\ref{prop:strictliftN}(6)(b) is a statement about an integral datum and is read by no item of
Theorem~\ref{thm:descentcert}, whose $3$-local piece is the transport of Proposition~\ref{prop:strictliftN}(5) for every $N$; v1.23.0-pre.25, referee M2), Corollary~\ref{cor:stablecert},
and any $\tau$-type table on the tree.

\emph{Notation.} $A:=\Ga_0(N)$, $g:=\smat100p$, $\Ga''_g:=\Ga\cap g\Ga g^{-1}$, $M_N:=\Ind_{A}^{\Ga}\Vcoef$ and $M_{N,\cring}:=M_N\otimes\cring$ (\S\ref{sec:descentN}),
$\mathsf O:=\Z[1/6]$; $\Hfr{q}(N):=H^{q}(\Ga_0(N);\Vcoef)/\mathrm{tors}$ (so $\Hfr{q}(1)=\Hfr{q}$). The symbols $\tau_X$, $\tau_{\Iso}$ (here: an admissible bulk lift of
Proposition~\ref{prop:strictliftN}(2)), $\widetilde T^{(N)}_{X,p}$, $\widetilde T^{(N)}_{\Iso,p}$, $T'^{(\ell)}_p$, $T'_{\mathcal T,p}$, $\chi^{(N)}_{p,q}$, $r_N$, $\mathsf B_p$, $B_{p,i}$,
$\mathsf b_{p,i}$, $\mathcal S$, $V_j(G)$, $\phi_j$, $P_j(\phi)$, $d_j$, $e_j$, $G_{[c]}$, $G_{\circ}$, $\mathrm{len}$ and $\tilde x_j$, $\tilde\Phi$ are new symbols of this subsection and mean
nothing else (rule~\ref{A5}). $T_p$ is always $[AgA]$, $p\nmid N$; $U_p:=[AgA]$, $p\mid N$, is defined and used only in item~(9) of Proposition~\ref{prop:strictliftN} and its proof
(v1.26.0-pre.1; item~(1) names it, and $U_p^{\ast}$, through Proposition~\ref{prop:levelconehecke}(3)(d); v1.26.0-pre.2, referee E1), where $\widetilde U^{(N)}_{\Iso,p}$,
$\widetilde U^{(N)}_{X,p}$, $U'^{(\ell)}_p$, $U'_{\mathcal T,p}$ and $U^{\mathrm{cl}}_p$ (the classical operator) are new symbols as well.

\begin{proposition}[The level-$N$ operator $T_p$ on the free model, every $p\nmid N$ (and $U_p$, $p\mid N$, item~(9)): the bulk lift, the induced-fibre form, the glueing map, the transports]\label{prop:strictliftN}
\PROVED{} \STD{\cite[Ch.~III.8--10]{brown}, \cite[Thm.~2.2.6]{weibel}} (v1.23.0-pre.24; milestone D3, clause C-D4 of the plan; refereed ACCEPT-with-fixes,
\file{docs/implementation/slot24/C-D4\_referee.md}, its fixes at v1.23.0-pre.25.) Assume \ref{H4} and
\ref{B1}--\ref{B3}. Let $N\ge1$, $n\ge0$, $p$ a prime with $p\nmid N$ (in item~(9): $p\mid N$), $\cring$ a commutative ring, and $(\Iso_{\bullet},\cmpI)$ an isotropy-resolution model.
\begin{enumerate}
\item \textbf{(The double coset; no adjoint.)} With $u,b\in\Z$, $up-bN=1$, $\gamma_1:=\smat{p}{-b}{-N}{u}$ and $\gamma_2:=\smat{pu}{b}{N}{1}$ lie in $A$ and
  $\gamma_1g\gamma_2=\smat p001$. So $AgA=A\smat p001A$, and $T_p:=[AgA]=\cor^{A}_{A''}\circ(\iota_g,\rhon(g))^{\ast}\circ\res^{A}_{A'}$ on $H^{\bullet}(A;\VR)$
  (Proposition~\ref{prop:descentN}(5)) is the double-coset operator of $\smat p001$ as well: for $p\nmid N$ the distinction between this document's operator and the
  classical one, which separates $U_p$ from $U_p^{\ast}$ when $p\mid N$ (Proposition~\ref{prop:levelconehecke}(3)(d)), does not arise.
\item \textbf{(A bulk strict lift for every $p$, constructed.)} Let $\tau_X$ be a $g$-twisted chain datum over $\Z$ on $C_{\bullet}(\Xbulk)$ (Definition~\ref{def:tau} with
  $\Xbulk$ for $\Ku$). The construction below and items (3)--(8) use of $\tau_X$ nothing else (v1.23.0-pre.25, referee m3). One exists for every $p$: the strict datum
  $\tau^{s}$ of Proposition~\ref{prop:strictpair}, with the boundary datum of Lemma~\ref{lem:taubdryexists}(1), both built by terminating procedures; and the shipped
  canonical $\tau$ of Algorithm~\ref{alg:tau} qualifies \emph{once its stored table is certified integral and checked against (T1)--(T3) over $\Z$} on $C_{\bullet}(\Xbulk)$
  (it is integral in practice, not by construction, Remark~\ref{rmk:bredonhecke}), in which case no boundary datum is built at all --- the cheaper input. \textbf{Code
  obligation}: whichever $\tau_X$ is read, its table is checked integral and against (T1)--(T3) over $\Z$ before the lift is used, and the lift is refused otherwise.
  \textbf{The check may be made through the lift} (v1.23.0-pre.31, D3's contract note~2): let the table read be a family of \emph{integral} chains $t_{j,c}\in C_{s}(\Xbulk)$,
  one for each element $s_je_c$ of the basis below with $e_c\in G_{s,0}$, and let $\tau_{\Iso}$ be additive and certified by (T1) on Schreier generators of $\Ga''_g$, (T2),
  (T3) and the strict identity on \textbf{every} basis element --- $\cmpI\tau_{\Iso}(s_je_c)=t_{j,c}$ for $e_c\in G_{s,0}$, and $\cmpI\tau_{\Iso}(s_je_c)=0$ for
  $e_c\in G_{s,t}$, $t\ge1$, where $\cmpI(s_je_c)=0$. Then the table \textbf{is} the table of a $g$-twisted chain datum over $\Z$ on $C_{\bullet}(\Xbulk)$: there is exactly one
  additive $\tau'_X$ with $\tau'_X\circ\cmpI=\cmpI\circ\tau_{\Iso}$, it satisfies (T1)--(T3) over $\Z$, and $\tau'_X(s_j\,\cmpI e_c)=t_{j,c}$ (the proof, \emph{the check through
  the lift}). So this certificate discharges the obligation for $\tau_X:=\tau'_X$, no separate check of the table on $\Xbulk$ is needed, and a table that is not the table of
  any such datum fails the certificate for \emph{every} $\tau_{\Iso}$. Then there is an additive $\tau_{\Iso}\colon\Iso_{\bullet}\to\Iso_{\bullet}$ with (T1)--(T3) of
  Definition~\ref{def:tau} over $\Z$ ($\Iso_{\bullet}$ in place of $C_{\bullet}(\Ku)$, augmentation $\aug\circ\cmpI$) and
  $\cmpI\circ\tau_{\Iso}=\tau_X\circ\cmpI$ strictly; two such are homotopic through (T1)-maps over $\Z$. \textbf{Every} such $\tau_{\Iso}$ is admissible below (v1.23.0-pre.31,
  D3's contract note~3): items (3)--(8) read of it only (T1)--(T3) over $\Z$ and the strict identity, and, for the cochain-level statements of (3)(c) and (7), that it is
  fixed independently of the transversal $r_t$ and of $N$; the bulk component of a pair lift (Proposition~\ref{prop:descentpairhecke}(1), built with the boundary step first)
  is one, and it is the one the code reads. Every read made from the induced operator is independent of the choice (\emph{independence of the lift}, below). One is \textbf{produced}, in every degree $q$, by the procedure of
  Remark~\ref{rmk:strictliftbuilt}(1) with no boundary step: on the free $\Z[\Ga''_g]$-basis $\{s_je_c\}$ of $\Iso_q$, the preimage of $\tau_X(s_j\cmpI e_c)$ is read off
  (at $t=0$; it is $0$ at $t\ge1$), and the cycle of $\ker\cmpI$ left over is contracted by descent through the columns. The procedure is deterministic, solves no
  linear system, terminates, and is uniform in $p$; its input $\tau_X$ and the coset table of $\Ga''_g\backslash\Ga$ ($p+1$ cosets) depend on $p$. Its output is certified by
  Remark~\ref{rmk:strictliftbuilt}(2)'s checks without the boundary items ((T1) on Schreier generators of $\Ga''_g$, (T2), (T3), and the strict identity). \textbf{No bound on
  its integer sizes is asserted} (Remark~\ref{rmk:integralscope}(2)); a coefficient that leaves the word is refused, never truncated (Theorem~\ref{thm:backend}(3)).
  \textbf{Independence of the lift} (v1.23.0-pre.31). Let $\tau_{\Iso}$, $\tau'_{\Iso}$ be admissible (for the same $\tau_X$, or for two data $\tau_X$, $\tau'_X$ as above), and
  $\widetilde T$, $\widetilde T'$ the operators of~(3) built from them with the same tables. Then $\widetilde T-\widetilde T'=d_{\mathrm I}\mathsf H+\mathsf Hd_{\mathrm I}$, $d_{\mathrm I}$ the coboundary of $\CI(M_{N,\cring})$, with
  $\mathsf H=\Theta_{h_\Iso}$ (formula \eqref{eq:levelcorr} with $h_\Iso$ for $\tau$), integral, of degree $-1$, for a (T1)-homotopy $h_\Iso$ from $\tau_{\Iso}$ to $\tau'_{\Iso}$
  (constructible by the comparison procedure; v1.23.0-pre.32, referee m1, N4), over every $\cring$; so $\widetilde T$ and $\widetilde T'$, and their transports~(5) (conjugated by the same comparisons),
  induce the same endomorphism $T_p$ of $H^{q}(A;\VR)$ in every degree. Hence every datum that the record of Theorem~\ref{thm:descentcert} reads from the induced operator ---
  the torsion modules and every invariant of Proposition~\ref{prop:torsionhecke}(1), the invariant factors of its~(4) --- is the same for both lifts; for a fixed certified
  Smith form, even the matrix $\tilde\Phi$ of Proposition~\ref{prop:torsionhecke}(3) is the same when $2\sigma_{\max}\le K$, and agrees modulo $\ell^{\min(\sigma_i,K')}$ in row $i$
  otherwise. (The characteristic polynomial of Proposition~\ref{prop:crtcharpolyN} reads no lift at all.)
\item \textbf{(The induced-fibre form.)} Let $\widetilde T^{(N)}_{\Iso,p}$ be formula \eqref{eq:levelcorr} with $A=B=\Ga_0(N)$ ($w_a=v_i$ the table of
  Definition~\ref{def:levelNmodel}, $r_t$ a transversal of $A''\backslash A$, $A''=\Ga_0(pN)$), with $\tau_{\Iso}$ in place of $\tau$ and $x$ running over $\Iso_{\bullet}$: for
  $\phi\in\Hom_{\Z[\Ga]}(\Iso_q,M_{N,\cring})$,
  \[
    \bigl(\widetilde T^{(N)}_{\Iso,p}\phi\bigr)(x)_i=\sum_t\rhon(\delta_{i,t})\,\phi\bigl(\tau_{\Iso}(s_{j(i,t)}x)\bigr)_{a(i,t)} .
  \]
  \begin{enumerate}
  \item[(a)] It is integral and reads only the level-$1$ table of $\tau_{\Iso}$ (its values on the $s_je_c$), the coset tables, and $\rhon$ of integral matrices of determinant
    $p$ (Proposition~\ref{prop:levelcorrespondence}(1)); no level-$N$ table of any lift is needed.
  \item[(b)] It is a cochain map $\CI(M_{N,\cring})\to\CI(M_{N,\cring})$ over \textbf{every} $\cring$.
  \item[(c)] Under Shapiro's isomorphism (the reciprocity $\phi\mapsto\phi(\cdot)_1$) it induces $T_p$ on $H^{q}(A;\VR)$ for \textbf{every} $q$ and every $\cring$, with no
    hypothesis on $6$; on cohomology it is independent of $\tau_X$, $\tau_{\Iso}$, the $s_j$, the $r_t$ and the table, and at cochain level it is independent of the transversal
    $r_t$ ($\tau_{\Iso}$ fixed). The $\tau_{\Iso}$ produced in~(2) \textbf{is} fixed in this sense: it is built once, on $\{s_je_c\}$, independently of $r_t$ and of $N$ (a level-$1$
    object), which is what (7) needs at both levels; so is the bulk component of a pair lift (Proposition~\ref{prop:descentpairhecke}(1), built once per $p$; v1.23.0-pre.31). A datum built from the transversal (as the natural one of the bar model) changes $\widetilde T^{(N)}_{\Iso,p}$ by a
    coboundary when $r_t$ changes; then only the cohomology statement survives (v1.23.0-pre.25, referee m4).
  \item[(d)] At $N=1$ it is the coset transport $\Th_{\Iso,g}$ of Proposition~\ref{prop:strictlift}(3).
  \end{enumerate}
\item \textbf{(The glueing map.)} Let $\widetilde T^{(N)}_{X,p}$ be \eqref{eq:levelcorr} on the shipped level-$N$ model with the datum $\tau_X$ of~(2). Then, over every $\cring$,
  \[
    \cmpI^{\ast}\circ\widetilde T^{(N)}_{X,p}=\widetilde T^{(N)}_{\Iso,p}\circ\cmpI^{\ast}\qquad\text{on the nose.}
  \]
  If the shipped model reads another admissible datum (the canonical $\tau$ of Algorithm~\ref{alg:tau}), its operator and $\widetilde T^{(N)}_{X,p}$ agree on cohomology whenever
  $6\in\cring^{\times}$ (Proposition~\ref{prop:levelcorrespondence}(3)). Hence $T_p$ on $H^{q}(A;\Vcoef)$ over $\Z$ is the endomorphism glued
  (Proposition~\ref{prop:descentN}(3)) from the shipped operator on the $\mathsf O$-piece $\Ceq(M_{N,\mathsf O})$ and from $\widetilde T^{(N)}_{\Iso,p}\otimes\Z_{(\ell)}$, or any of its
  transports~(5), on the $\ell$-local pieces, $\ell\in\{2,3\}$; no coherence datum enters.
\item \textbf{(The transports; the route at $2$.)} Let $\ell$ be a prime and $\cring$ an $\mathsf O_\ell$-algebra. Then
  \[
    T'^{(\ell)}_p:=\Psi_\ell^{\ast}\circ\widetilde T^{(N)}_{\Iso,p}\circ\Phi_\ell^{\ast}\ \text{on }C^{\bullet}_{\Iso^{(\ell)}}(M_{N,\cring}),\qquad
    T'_{\mathcal T,p}:=\Psi_{\mathcal T}^{\ast}\circ\widetilde T^{(N)}_{\Iso,p}\circ\Phi_{\mathcal T}^{\ast}\ \text{on }C^{\bullet}_{\mathcal T,\ell}(M_{N,\cring})
  \]
  are cochain maps defined over $\mathsf O_\ell$ (Propositions~\ref{prop:isomodelell}(4), \ref{prop:isomodeltree}(3), (4)) that induce $T_p$ on $H^{q}(A;\VR)$ in every degree,
  independently of the comparisons up to cochain homotopy, and that agree with each other on cohomology (Proposition~\ref{prop:isomodeltree}(6)(a), (b) with
  $T:=\widetilde T^{(N)}_{\Iso,p}$). At $\ell=2$ the route is $T'_{\mathcal T,p}$ on $\Iso^{(2)}_{\mathcal T}$ ([R-GT]); every factor is constructed, and applying it to a cochain
  composes three cochain maps --- the Hecke datum stays on $\Iso_{\bullet}$, no $\tau$-type table on the tree is used.
\item \textbf{(The cheap routes.)} (a) If $6\in\cring^{\times}$ --- the $\mathsf O$-piece, and its reductions modulo $\ell^{K}$, $\ell\ge5$ --- the shipped operator induces
  $T_p$ on $H^{q}(\Ceq(M_{N,\cring}))=H^{q}(A;\VR)$, and so does $\mathsf i_{\mathcal T}^{\ast}\circ\widetilde T^{(N)}_{X,p}\circ\mathsf r_{\mathcal T}^{\ast}$ on the tree complex
  $\mathrm T(M_{N,\cring})$ \eqref{eq:treemodel}, with the integral chain maps $\mathsf i_{\mathcal T}$, $\mathsf r_{\mathcal T}$ of Proposition~\ref{prop:isomodeltree}(4)(a).
  (b) If $\nu_3(N)=0$ (Proposition~\ref{prop:descentNthree}: $\Ga_0(N)$ has no element of order $3$) and $\cring$ is a $\Z_{(3)}$-algebra, the same two statements hold
  \textbf{for an integral datum}: \eqref{eq:levelcorr} with a $g$-twisted chain datum over $\Z$ on $\Xbulk$, certified by (T1)--(T3) over $\Z$, induces $T_p$ on
  $H^{q}(A;\VR)$ for every $q$. The shipped canonical $\tau$ is such a datum only once so certified (Remark~\ref{rmk:bredonhecke}: integral in practice, not by construction).
  \textbf{This item is read by no statement of this subsection} (v1.23.0-pre.25, referee M2, resolution (a)): the $3$-local piece of Theorem~\ref{thm:descentcert} is the
  transport $T'^{(3)}_p$ (or $T'_{\mathcal T,p}$) of~(5) for \emph{every} $N$, $\nu_3(N)=0$ included, so the shipped operator on the Bredon object is never read at
  $\ell\in\{2,3\}$.
\item \textbf{(Strict squares with the restriction; the level sequence.)} Let $B:=\Ga_0(MN)\subseteq A$ with $p\nmid MN$, and build $\widetilde T^{(N)}_{\Iso,p}$ and
  $\widetilde T^{(MN)}_{\Iso,p}$ from the same $\tau_{\Iso}$ and a common transversal $r_t\in B$ of $\Ga_0(pMN)\backslash B$ (which is one of $A''\backslash A$,
  Proposition~\ref{prop:levelconehecke}(2)). Then $\widetilde T^{(MN)}_{\Iso,p}\circ\alpha_{\Iso}=\alpha_{\Iso}\circ\widetilde T^{(N)}_{\Iso,p}$ \textbf{as cochain maps}, over every
  $\cring$; so $(\widetilde T^{(N)}_{\Iso,p}[1],\widetilde T^{(MN)}_{\Iso,p})$ is a cochain endomorphism of $\Cone{\alpha_{\Iso}}$ commuting with its short exact sequence, its
  transports to every model of Proposition~\ref{prop:descentmaps} satisfy the same strict identity, and the long exact sequence \eqref{eq:descentmapsles} is
  $T_p$-equivariant \textbf{over $\Z$ and over every $\Z_{(\ell)}$, $\ell\in\{2,3\}$ included}, in every degree --- which supplies what
  Proposition~\ref{prop:descentmaps}(3)(c) left unasserted at $\ell\in\{2,3\}$. The square with $\cor$ commutes on cohomology (Proposition~\ref{prop:descentmaps}(3)(a)); no
  cochain-level statement is made for it.
\item \textbf{(Relations.)} For primes $p,p'\nmid N$, $T_pT_{p'}=T_{p'}T_p$ on $H^{q}(A;\VR)$, every $q$, every $\cring$.
\item \textbf{($p\mid N$: the operator $U_p$ on the free model.)} (v1.26.0-pre.1, slot~26, UP-Z; \file{docs/implementation/slot26/UPZ\_design.md} \S7.1 and its critical audit
  \file{docs/implementation/slot26/UPZ\_design\_audit.md} \S2.3, condition~C5; refereed ACCEPT-with-conditions, \file{docs/implementation/slot26/C126pre1\_referee.md}, its conditions and fixes at
  v1.26.0-pre.2; \textbf{not yet seen by the owner}.) In this item $p$ is a prime with $p\mid N$,
  in place of $p\nmid N$ in the hypothesis above; everything else is as there, $g=\smat100p$. Put
  $U_p:=[AgA]=\cor^{A}_{A''}\circ(\iota_g,\rhon(g))^{\ast}\circ\res^{A}_{A'}$ on $H^{\bullet}(A;\VR)$.
  \begin{enumerate}
  \item[(a)] \textbf{(The coset data.)} $A''=A\cap gAg^{-1}=\Ga_0(pN)$, of index $p$ in $A$, with the transversal $r_t:=\smat10{tN}1$, $0\le t<p$, of $A''\backslash A$; and
    $A''\subseteq\Ga''_g=\Ga_0(p)$. (For $p\nmid N$ the index is $p+1$, the transversal of~(3).)
  \item[(b)] \textbf{(Which $U_p$.)} $U_p$ is this document's $U^{(N)}_p$ of Proposition~\ref{prop:levelcorrespondence}(4)(d), the $[A\,g\,A]$ of Proposition~\ref{prop:heilbronn}(2) at
    $p\mid N$, and the operator to which Definition~\ref{def:heckeZ} sends $X_p$ (through Proposition~\ref{prop:descentmaps}(3)(b$'$)). On $p$-old spaces it is the \emph{adjoint}
    $U_p^{\ast}=W_pU^{\mathrm{cl}}_pW_p^{-1}$ of the classical operator $U^{\mathrm{cl}}_p$ (Proposition~\ref{prop:levelconehecke}(3)(d); v1.26.0-pre.2, referee E2); item~(1) has no
    analogue at $p\mid N$: $A\smat p001A\ne AgA$ as double cosets, since by Step~1 of the proof of Proposition~\ref{prop:heilbronn} the Hermite form $\smat p001$ is the one right
    coset of $\De_p$ absent from $AgA$ at $p\mid N$ (v1.26.0-pre.2, referee E9), and $[A\smat p001A]$ is not asserted to equal $U_p$. Every statement and record of this document that reads $U_p$, $p\mid N$, reads this operator.
  \item[(c)] \textbf{(Items (2)--(6) hold, their references to (7), (8) dropped.)} Write $\widetilde U^{(N)}_{\Iso,p}$ for formula \eqref{eq:levelcorr} as in~(3) with $A''=\Ga_0(pN)$, the transversal of~(a)
    and \textbf{the same} $\tau_{\Iso}$ of~(2) --- the level-$1$ lift, built and certified once per $p$, independently of $N$ (the bulk component of a pair lift,
    Proposition~\ref{prop:descentpairhecke}(1), is one; the coset table of $\Ga''_g\backslash\Ga$ still has $p+1$ cosets, the level-$N$ transversal has $p$) ---
    $\widetilde U^{(N)}_{X,p}$ for the shipped operator of~(4), and $U'^{(\ell)}_p$, $U'_{\mathcal T,p}$ for the transports of~(5). Then items (2)--(6), and the independence of
    the lift stated in~(2), hold (with their references to (7) and (8) dropped: item~(2)'s ``items (3)--(8)'' reads ``items (3)--(6)'', its ``(3)(c) and (7)'' reads
    ``(3)(c)''; v1.26.0-pre.2, referee E7) with $U_p$ for $T_p$ and these operators for $\widetilde T^{(N)}_{\Iso,p}$, $\widetilde T^{(N)}_{X,p}$, $T'^{(\ell)}_p$, $T'_{\mathcal T,p}$; (3)(d) is
    vacuous (no prime divides $1$). In particular $\widetilde U^{(N)}_{\Iso,p}$ is integral and reads only the level-$1$ table of $\tau_{\Iso}$, is a cochain map over every
    $\cring$, induces $U_p$ on $H^{q}(A;\VR)$ for \textbf{every} $q$ and every $\cring$ with no hypothesis on $6$, and is independent of $r_t$ at cochain level; $U_p$ over $\Z$ is
    glued from the $\mathsf O$-piece and the $\ell$-local pieces, $\ell\in\{2,3\}$ ((4)); at $\ell\in\{2,3\}$ the route is the transport~(5), at $\ell=2$ on the tree. Every datum
    the record reads from the induced operator (Proposition~\ref{prop:torsionhecke}) is independent of the lift. \textbf{Code obligation}: the level-$N$ coset count is $[A:A'']=p$
    at $p\mid N$ (not $p+1$), and the lift is the certified level-$1$ one of~(2); a lift refused by a certificate is a per-cell refusal, never evidence that no lift exists.
  \item[(d)] \textbf{(Not asserted at $p\mid N$.)} Items (1), (7) and (8). Nothing in this subsection, or in the record of item~8 of \S\ref{sec:recordschema}, reads them for $U_p$.
    The commutativity of $U_p$ with the $T_p$ that the record's $\vartheta$ needs is Proposition~\ref{prop:heckeZspec}(0), not item~(8) (v1.26.0-pre.2, referee E3).
  \end{enumerate}
\end{enumerate}
\end{proposition}

\begin{proof}
(1) $\det\gamma_1=pu-bN=1$, $\det\gamma_2=pu-bN=1$, the lower-left entries are $\mp N$, so both lie in $A$; and
$\gamma_1g=\smat{p}{-bp}{-N}{up}$, $\gamma_1g\gamma_2=\smat{p^{2}u-bpN}{pb-bp}{-Npu+upN}{-Nb+up}=\smat p001$. The operator $[AgA]$ depends only on the double coset: replacing
$g$ by $\gamma_1g\gamma_2$ composes the three factors with conjugations by $\gamma_1$, $\gamma_2\in A$, which act trivially on $H^{\bullet}(A;-)$ \STD{\cite[Ch.~III, \S8]{brown}:
an inner automorphism acts trivially on cohomology}, together with the compatible coefficient maps $\rhon(\gamma_i)$ (the computation of Theorem~\ref{thm:welldef}(c); the compatibility is the multiplicativity of $\rhon$ on
$\Mat_2(\Z)$, Proposition~\ref{prop:symconv}(2), v1.23.0-pre.25, referee t2).

(2) Proposition~\ref{prop:strictpair} builds $\tau^{s}$ for every $g\in\De$ from a boundary datum, and Lemma~\ref{lem:taubdryexists}(1) builds that for $g=\smat100p$; a
certified canonical table is a $g$-twisted chain datum over $\Z$ by definition of the certificate. Nothing below reads more of $\tau_X$ than (T1)--(T3) over $\Z$.
Then the proof of Proposition~\ref{prop:strictlift}(1) with $P'=0$ (no boundary step): $\Ga''_g$ has finite index in $\Ga$, so $\Iso_q$ is free over $\Z[\Ga''_g]$ on
$\{s_je_c\}$, and $\cmpI\colon\Iso^{(c)}\to C^{(c)}_{\bullet}(\Xbulk)$ (twisted actions) is a $\Ga''_g$-equivariant surjective quasi-isomorphism; Lemma~\ref{lem:strictlifting}
with $f=\tau_X\circ\cmpI$ gives $\tau_{\Iso}$, (T2) because it is a chain map, (T3) from (T3) of $\tau_X$; uniqueness up to (T1)-homotopy is
Proposition~\ref{prop:strictlift}(2)'s argument. \emph{The procedure.} The induction of Lemma~\ref{lem:strictlifting} needs a preimage $y_0$ of $f(x)$ and a contraction of the
cycle $z$ of the acyclic $\ker\cmpI$. At $t=0$, $\cmpI$ is the augmentation of free generators, so a chain $\sum c\,\gamma\,\sigma'$ of $\Xbulk$ is the image of
$\sum c\,\gamma\,(1\otimes1)_{\sigma'}\in G_{\bullet,0}$; at $t\ge1$ or in degrees $\ge3$, $f(x)=0$ and $y_0=0$. For the contraction filter $\ker\cmpI$ by $s$: if $s$ is the
largest column on which $z$ has a component $z_s$, then $d^{(0)}z_s=0$ (it is the top-column part of $dz=0$) and $z_s$ lies in the augmented kernel of
$G_{s,\bullet}\to C_s(\Xbulk)$, whose closed-form $\Z$-linear contracting homotopy $h$ (Proposition~\ref{prop:isomodel}'s proof, the columns) gives $k_s:=(-1)^{s}h(z_s)$ with
$d^{(0)}k_s=z_s$; $k_s$ lies in $t\ge1$ or in the augmentation kernel, so in $\ker\cmpI$, and $z-dk_s$ has no component on columns $\ge s$. Descending on $s\le2$ contracts $z$.
This is the formula (L) of Proposition~\ref{prop:isomodeltree}(4)(b) with $\cmpI$ as its surjection, $f=\tau_X\circ\cmpI$, the basis $\{s_je_c\}$ and the twisted
$\Z[\Ga''_g]$-linear extension (T1) in place of the $\Z[\Ga]$-linear one --- (L) uses of the extension only that $\mathsf h_p$ is $\Z$-linear with $\partial\mathsf h_p(z)=z$ and
$\cmpI\,\mathsf h_p=0$. No step solves a system or depends on $p$ except through the coset table. The certificates are those of Remark~\ref{rmk:strictliftbuilt}(2), which re-derive (T1)--(T3) and the strict
identity from the stored values; their boundary items concern $\tau^{\partial}_{\Iso}$, not built here.

\emph{The check through the lift} (v1.23.0-pre.31). Put $T:=\cmpI\circ\tau_{\Iso}\colon\Iso_{\bullet}\to C_{\bullet}(\Xbulk)$, additive. On $\Iso$, (T1) is the defining property of the
extension from the free $\Z[\Ga''_g]$-basis (the Schreier check certifies its bookkeeping), and $\cmpI$ is $\Ga$-equivariant, so $T(hy)=(g^{-1}hg)\,T(y)$ for $h\in\Ga''_g$.
\emph{Step 1: $T=0$ on $G_{s,t}$, $t\ge1$.} $G_{s,t}=\bigoplus_\sigma\Z[\Ga]\otimes_{\Z[\Ga_\sigma]}F^{\sigma}_t$ with $F^{\sigma}_t=\Z[\Ga_\sigma]$ is free over $\Z[\Ga]$ on its
$e_c$, hence spanned over $\Z[\Ga'']$ by the $s_je_c$ in it, on which $T$ vanishes by hypothesis; twisted linearity gives the rest. \emph{Step 2: $T=0$ on
$\ker\cmpI\cap\bigoplus_sG_{s,0}$.} On $G_{s,0}=\bigoplus_\sigma\Z[\Ga]\otimes_{\Z[\Ga_\sigma]}\Z[\Ga_\sigma]$, $\cmpI$ is $\id\otimes\aug$ (Definition~\ref{def:isomodel}); $\Z[\Ga]$ is free,
hence flat, over $\Z[\Ga_\sigma]$, so its kernel is $\bigoplus_\sigma\Z[\Ga]\otimes_{\Z[\Ga_\sigma]}I_\sigma$, $I_\sigma$ the augmentation ideal of the cyclic group $\Ga_\sigma$, which is
$(g_\sigma-1)\Z[\Ga_\sigma]=\partial^{F}(F^{\sigma}_1)$. So such an $x$ is $x=d^{(0)}y$ for some $y\in G_{s,1}$ (the sign $(-1)^{s}$ of $d^{(0)}$ is absorbed into $y$). Now
$dy=d^{(0)}y+d^{(1)}y+d^{(2)}y$ with $d^{(r)}y\in G_{s-r,r}$, at $t=r\ge1$ for $r\ge1$; by Step~1, (T2) of $\tau_{\Iso}$ and the chain map $\cmpI$,
$T(x)=T(dy)=\cmpI\,\tau_{\Iso}(dy)=\cmpI\,d\,\tau_{\Iso}(y)=\partial\,\cmpI\,\tau_{\Iso}(y)=\partial\,T(y)=0$. \emph{Step 3.} $\cmpI$ vanishes on $t\ge1$, so
$\ker\cmpI=(\ker\cmpI\cap\bigoplus_sG_{s,0})\oplus\bigoplus_{t\ge1}G_{s,t}$, and $T(\ker\cmpI)=0$ by Steps~1--2. As $\cmpI$ is onto in every degree
(Proposition~\ref{prop:isomodel}(2)), $\tau'_X(\cmpI x):=T(x)$ defines an additive map on $C_{\bullet}(\Xbulk)$, the only one with $\tau'_X\cmpI=\cmpI\tau_{\Iso}$; it is integral
($\cmpI$ and $\tau_{\Iso}$ are). For $y=\cmpI x$ and $h\in\Ga''_g$: (T1) $\tau'_X(hy)=T(hx)=(g^{-1}hg)T(x)=(g^{-1}hg)\tau'_X(y)$; (T2)
$\partial\tau'_X(y)=\partial\cmpI\tau_{\Iso}x=\cmpI\tau_{\Iso}dx=\tau'_X(\cmpI dx)=\tau'_X(\partial y)$; (T3) $\aug\tau'_X(y)=(\aug\circ\cmpI)(\tau_{\Iso}x)=(\aug\circ\cmpI)(x)=\aug(y)$,
by (T3) of $\tau_{\Iso}$ with the augmentation $\aug\circ\cmpI$. Finally $\tau'_X(s_j\cmpI e_c)=T(s_je_c)=t_{j,c}$; as $\Ga=\bigsqcup_j\Ga''_gs_j$ and every simplex of $\Xbulk$
is $\gamma\,\cmpI e_c$, $\gamma\in\Ga$ (the stored-oriented representatives $\Sigma_s$), the simplices $s_j\cmpI e_c$ contain a set of $\Ga''_g$-orbit representatives, and the
(T1)-extension of the table from any such set is $\tau'_X$: the table is the table of the datum $\tau'_X$. Conversely, if the table is not the table of any $g$-twisted chain datum
over $\Z$, no additive $\tau_{\Iso}$ passes the certificate, since one that did would produce such a datum.

\emph{Independence of the lift.} Both $\tau_{\Iso}$ and $\tau'_{\Iso}$ are chain maps $\Iso_{\bullet}\to\Iso_{\bullet}$ that are $\Z[\Ga''_g]$-linear for the twisted action $h\mapsto g^{-1}hg$ on
the target (T1) and lift $\id_\Z$ (T3); the source is free over $\Z[\Ga''_g]$ and the target, restricted along $h\mapsto g^{-1}hg$, is an acyclic augmented complex, so the fundamental
lemma \STD{\cite[Thm.~2.2.6]{weibel}} gives a $\Z[\Ga''_g]$-linear $h_\Iso$ of degree $+1$, i.e.\ with (T1), and $\tau_{\Iso}-\tau'_{\Iso}=\partial h_\Iso+h_\Iso\partial$ --- whether or not the
two lifts start from the same $\tau_X$. For any additive $u$ of degree $k$ with (T1) put $(\Theta_u\phi)(x)_i:=\sum_t\rhon(\delta_{i,t})\,\phi(u(s_{j(i,t)}x))_{a(i,t)}$, the formula
\eqref{eq:levelcorr} with $u$ for $\tau$; the proof of Proposition~\ref{prop:levelcorrespondence}(3) shows, with (T1) of $u$ only, that $\Theta_u$ maps $\Ga$-equivariant cochains of degree
$q+k$ to $\Ga$-equivariant cochains of degree $q$. $\Theta_u$ is additive in $u$, $\Theta_{\tau_{\Iso}}=\widetilde T$, and, because $\partial(s_jx)=s_j\partial x$ and the indices
$j(i,t)$, $a(i,t)$, $\delta_{i,t}$ do not depend on $x$, $\Theta_{\partial\circ h_\Iso}\phi=\Theta_{h_\Iso}(\phi\circ\partial)$ and $\Theta_{h_\Iso\circ\partial}\phi=(\Theta_{h_\Iso}\phi)\circ\partial$. With
$(d_{\mathrm I}\phi)(x)=\phi(\partial x)$, unsigned (the convention of Lemma~\ref{lem:cosettransport}(2); v1.23.0-pre.32, referee m1), $\mathsf H:=\Theta_{h_\Iso}$ gives
$\widetilde T-\widetilde T'=\Theta_{\partial h_\Iso}+\Theta_{h_\Iso\partial}=\mathsf Hd_{\mathrm I}+d_{\mathrm I}\mathsf H$ in every degree, with no sign; these are the two identities of
Lemma~\ref{lem:cosettransport}(2) in the level-$N$ form. (Under a signed convention $d_{\mathrm I}=\varepsilon_q(\,\cdot\circ\partial)$ on $C^{q}$ one would take $\mathsf H:=\varepsilon_{q-1}\Theta_{h_\Iso}$ on
$C^{q}$; not the convention here.) It is
integral because $h_\Iso$ and the $\delta_{i,t}$ are. For the transports, $\Psi^{\ast}(\widetilde T-\widetilde T')\Phi^{\ast}=d(\Psi^{\ast}\mathsf H\Phi^{\ast})+(\Psi^{\ast}\mathsf H\Phi^{\ast})d$, $d$ the
coboundary of the target model, the comparisons being cochain maps. So the induced endomorphisms of $H^{q}$ coincide (both are $T_p$ by (3)(c), (5)). The torsion module $(G,\phi)$ and the lattice $(\Lambda,\psi)$ of
Proposition~\ref{prop:torsionhecke}(3), (4) are defined from the induced endomorphism, so their invariants agree; for a fixed certified Smith form, $\tilde\Phi_{ij}$ is the exact
coordinate of $\phi[x_j]$ when $2\sigma_{\max}\le K$ and agrees with it modulo $\ell^{\min(\sigma_i,K')}$ otherwise (ibid.\ (3) and its proof), which gives the last sentence.

(3) (a) is Proposition~\ref{prop:levelcorrespondence}(1), whose proof reads nothing of the complex. (b) is Proposition~\ref{prop:levelcorrespondence}(2)'s proof with
$\Iso_{\bullet}$ for $C_{\bullet}(\Ku)$: it uses (T1), (T2) of $\tau$ and $\partial(s_jx)=s_j\partial x$, both true on $\Iso_{\bullet}$ (Lemma~\ref{lem:cosettransport}). (c) The proof of
Proposition~\ref{prop:levelcorrespondence}(3) assumes $6\in\cring^{\times}$ only to make $C_{\bullet}(\Ku)\otimes\cring$ a projective resolution of $\cring$ over $\cring[A]$, over
$\cring[A']$ and over $\cring[A'']$; $\Iso_{\bullet}\otimes\cring$ is a \emph{free} $\cring[\Ga]$-resolution (Proposition~\ref{prop:isomodel}(2)), hence free over every subgroup of
finite index. On it the inclusion computes $\res^{A}_{A'}$, $\psi\mapsto\rhon(g)\psi\circ\tau_{\Iso}$ computes $(\iota_g,\rhon(g))^{\ast}$ --- $\tau_{\Iso}$ is, by (T1) and (T3),
a chain map along $\iota_g\colon A''\to A'$ lifting $\id_\cring$, unique up to homotopy by the fundamental lemma, the source being free over $\cring[A'']$ and the target
acyclic \STD{\cite[Thm.~2.2.6]{weibel}} --- and the transfer computes $\cor^{A}_{A''}$ \STD{\cite[Ch.~III.9--10]{brown}}; the composite is the reciprocity form, which
equals \eqref{eq:levelcorr} by Proposition~\ref{prop:levelcorrespondence}(3)'s computation. Independence on cohomology: the double-coset operator does not depend on these choices
(Brown, loc.\ cit.). At cochain level, the independence of the transversal is Proposition~\ref{prop:levelconehecke}(4), whose proof uses (T1) of $\tau$ only, for a $\tau$ that does not depend on $r_t$;
the $\tau_{\Iso}$ of~(2) is built on $\{s_je_c\}$ from $\tau_X$ alone. (d) is
Proposition~\ref{prop:levelcorrespondence}(4)(a).

(4) For $\phi$ on $C_q(\Xbulk)$ and $x\in\Iso_q$, using $\cmpI\,\tau_{\Iso}=\tau_X\,\cmpI$ and the $\Ga$-equivariance of $\cmpI$ ($s_j\in\Ga$):
$(\widetilde T^{(N)}_{\Iso,p}\cmpI^{\ast}\phi)(x)_i=\sum_t\rhon(\delta_{i,t})\phi(\cmpI\,\tau_{\Iso}(s_jx))_a=\sum_t\rhon(\delta_{i,t})\phi(\tau_X(s_j\,\cmpI x))_a=
(\cmpI^{\ast}\widetilde T^{(N)}_{X,p}\phi)(x)_i$ --- the proof of Proposition~\ref{prop:strictlift}(3), blockwise. For $6\in\cring^{\times}$ both shipped operators induce
$[AgA]$ by Proposition~\ref{prop:levelcorrespondence}(3), and $H(\cmpI^{\ast})$ is an isomorphism intertwining $\widetilde T^{(N)}_{X,p}$ with $\widetilde T^{(N)}_{\Iso,p}$. So on
$H^{q}(A;\Vcoef)\otimes\mathsf O$ and on $H^{q}(A;\Vcoef)\otimes\Z_{(\ell)}$ the local operators are $T_p\otimes\mathsf O$, $T_p\otimes\Z_{(\ell)}$, and
Proposition~\ref{prop:descentN}(3) glues them to $T_p$.

(5) Proposition~\ref{prop:isomodeltree}(6)(a)--(c) is stated for any cochain endomorphism $T$ of $\CI(M_\cring)$ inducing an endomorphism $t$ of $H^{\bullet}(\Ga;M_\cring)$; take
$T=\widetilde T^{(N)}_{\Iso,p}\otimes\cring$ and $t$ the transport of $T_p$ under Shapiro's isomorphism, which is natural and independent of the resolution. The same argument with
$\Phi_\ell$, $\Psi_\ell$ (Proposition~\ref{prop:isomodelell}(4), (5): conjugating a cochain map that induces $t$ by comparisons induces $t$) gives $T'^{(\ell)}_p$. The coefficients
lie in $\mathsf O_\ell$ because $\Phi_\ell$, $\Psi_\ell$, $\Phi_{\mathcal T}$, $\Psi_{\mathcal T}$ do and $\widetilde T^{(N)}_{\Iso,p}$ is integral.

(6) (a) For $6\in\cring^{\times}$ the first statement is Proposition~\ref{prop:levelcorrespondence}(3); $C_{\bullet}(\Xbulk)\otimes\cring$ and $C_{\bullet}(\mathcal T)\otimes\cring$ are
projective resolutions of $\cring$ (Remark~\ref{rmk:orbifold}, Lemma~\ref{lem:treeres}) and $\mathsf i_{\mathcal T}$, $\mathsf r_{\mathcal T}$ are augmentation-preserving
$\Z[\Ga]$-chain maps between them (Proposition~\ref{prop:isomodeltree}(4)(a)), i.e.\ comparisons; conjugating by them preserves the induced operator. (b) By
Proposition~\ref{prop:descentNthree}(1) every stabiliser in $A$ of a simplex of $\Xbulk$, or of a cell of $\mathcal T$, has order $1$, $2$ or $4$, a unit in $\cring$; so
$C_{\bullet}(\Xbulk)\otimes\cring$ and $C_{\bullet}(\mathcal T)\otimes\cring$ are projective resolutions of $\cring$ over $\cring[A]$ and over its subgroups $A'$, $A''$ (the proof of
Proposition~\ref{prop:descentNthree}(2)), which, with a datum over $\Z$ satisfying (T1)--(T3), is all the proof of Proposition~\ref{prop:levelcorrespondence}(3) uses; a datum known
only modulo $\ell$, or not checked, is outside the hypothesis.

(7) The proof of Proposition~\ref{prop:levelconehecke}(1) reads, for both operators, the reciprocity form $\psi\mapsto\sum_t\rhon(r_t^{-1}g)\psi(\tau(r_t\,\cdot))$ with the same
$\tau$ and the same $r_t$, and $(\alpha_{B/A})_{\ast}$ as $\psi\mapsto\psi$; it uses nothing of the complex, so it holds with $\Iso_{\bullet}$, $\tau_{\Iso}$. Its hypotheses
$B\cap A''=B''$ and $A=A''B$ for $T_p$, $p\nmid MN$, are verified in Proposition~\ref{prop:levelconehecke}(2). The cone: the computation of Lemma~\ref{lem:conesign} with
$\hder=0$, as there. Transport: Proposition~\ref{prop:descentmaps}(1)(a) ($\alpha$ commutes strictly with every comparison), so
$T'^{B}\alpha=\Psi^{\ast}T^{B}\Phi^{\ast}\alpha=\Psi^{\ast}T^{B}\alpha\Phi^{\ast}=\Psi^{\ast}\alpha T^{A}\Phi^{\ast}=\alpha T'^{A}$; on the tree this is
Proposition~\ref{prop:isomodeltree}(6)(d). Over $\Z$ the cone $\Cone{\alpha_{\Iso}}$ itself carries the operator; Proposition~\ref{prop:descentmaps}(2)(c) glues its cohomology from the
local cones, compatibly by the strict identities.

(8) $\Vcoef$ is a module over the semigroup $\De_0(N)$ of integral matrices with lower-left entry $\equiv0\pmod N$ and upper-left entry prime to $N$ ($\rhon$ is multiplicative on
$\Mat_2(\Z)$), so the Hecke ring $\mathcal H(\Ga_0(N),\De_0(N))$ acts on $H^{q}(A;\VR)$ through $AgA\mapsto[AgA]$, a ring homomorphism for every module and every $q$ (the citations of
Proposition~\ref{prop:relationsmodel}(2): \cite[Thm.~1]{rhie_whaples}, \cite[\S1.1, p.~194]{ash_stevens}); and $T(1,p)$, $T(1,p')$ commute in that ring
\STD{\cite[Thm.~3.34, Thm.~3.35 and \S3.3]{shimura}, as cited in Theorem~\ref{thm:nebrelations}; commutativity among the $T(1,p)$, $p\nmid N$, \VERIFYloc}.

(9) (v1.26.0-pre.1.) (a) For $\gamma=\smat abcd\in A$, $g\gamma g^{-1}=\smat{a}{b/p}{pc}{d}$; so the integral elements of $gAg^{-1}$ are the $\smat{a'}{b'}{c'}{d'}\in\Ga$ with
$c'=pc$, $N\mid c$, i.e.\ $pN\mid c'$, and $A''=A\cap gAg^{-1}=\Ga_0(pN)$; with $\Ga$ in place of $A$ the same computation gives $\Ga''_g=\Ga_0(p)\supseteq\Ga_0(pN)$. As
$[\Ga:\Ga_0(M)]=M\prod_{q\mid M}(1+1/q)$, $[A:A'']=p$ when $p\mid N$. The $r_t$ lie in $A$ (lower-left entry $tN$), and $r_sr_t^{-1}=\smat10{(s-t)N}1\in\Ga_0(pN)$ iff
$p\mid s-t$; so the $p$ matrices $r_t$ represent distinct classes of $A''\backslash A$, hence all of them. (b) Proposition~\ref{prop:levelcorrespondence}(4)(d) is
\eqref{eq:levelcorr} with $A=B=\Ga_0(N)$ and this $g$, named $U^{(N)}_p$ for $p\mid N$; Proposition~\ref{prop:heilbronn}(2) and Definition~\ref{def:heckeZ} (through
Proposition~\ref{prop:descentmaps}(3)(b$'$)) use the double coset of the same $g$; the adjointness is Proposition~\ref{prop:levelconehecke}(3)(d). The proof of item~(1) needs
$u,b$ with $up-bN=1$, which do not exist when $p\mid N$. (c) Item by item, the hypothesis $p\nmid N$ is not used. (2): the statement and its proof read $\tau_X$, $\Ga''_g$ and the
coset table of $\Ga''_g\backslash\Ga$ only --- no $N$. The lift the code reads, the bulk component of a pair lift, is admissible at $p\mid N$ as well (v1.26.0-pre.2, referee S1):
Proposition~\ref{prop:descentpairhecke}(1)'s construction, certificate and homotopy uniqueness read $\Ga''_g=\Ga_0(p)$, $g$, $\taub$, $\tau^{s}$ and $\Isob$ only, nothing of $N$
(the datum is a level-$1$ object, built once per $p$); the hypothesis $p\nmid N$ of that proposition is the scope of its other items and is not used by its item~(1). (3)(a), (b): Proposition~\ref{prop:levelcorrespondence}(1), (2) and their proofs hold for any finite-index $A$, $B$ and any
$g\in\De$; Lemma~\ref{lem:cosettransport} reads nothing of $N$. (3)(c): the proof uses that $\Iso_{\bullet}\otimes\cring$ is free over every subgroup of finite index, that
$\tau_{\Iso}$ is a chain map along $\iota_g\colon A''\to A'$ lifting $\id_\cring$ (by (T1) on $\Ga''_g\supseteq A''$, (a), and (T3)), the fundamental lemma and the transfer; and,
for the independence of $r_t$, Proposition~\ref{prop:levelconehecke}(4), whose proof uses (T1) and $A''\subseteq\Ga''_g$ only. (4): the displayed identity is the computation of its
proof, which uses the strict identity $\cmpI\tau_{\Iso}=\tau_X\cmpI$ and the $\Ga$-equivariance of $\cmpI$; the comparison for $6\in\cring^{\times}$ is
Proposition~\ref{prop:levelcorrespondence}(3), for any $g\in\De$; the glueing is Proposition~\ref{prop:descentN}(3), with the same argument as there. (5): Propositions~\ref{prop:isomodeltree}(6)
and \ref{prop:isomodelell}(4), (5) hold for any cochain endomorphism that induces an endomorphism of $H^{\bullet}$. (6): (a) is Proposition~\ref{prop:levelcorrespondence}(3) with
Proposition~\ref{prop:isomodeltree}(4)(a), (b) is Proposition~\ref{prop:descentNthree} (orders of stabilisers); neither uses $p\nmid N$. The independence of the lift uses the
fundamental lemma over $\Z[\Ga''_g]$ and the formal identities of $\Theta_u$, which read the indices $j(i,t)$, $a(i,t)$, $\delta_{i,t}$ only. (d) Item (1) is excluded by (b); item (7)
rests on Proposition~\ref{prop:levelconehecke}(2)'s instance for $T_p$, $p\nmid MN$, and item (8) on the Hecke-ring relations among the $T(1,p)$, $p\nmid N$; neither is claimed for
$U_p$ here.
\end{proof}

\begin{proposition}[The characteristic polynomial over $\Z$ at level $N$, by residues; the certified bound; the output rule]\label{prop:crtcharpolyN}
\PROVED{} \STD{\cite[Thm.~12.2.2, p.~108, Prop.~12.4.10, p.~117]{diamond_im}, \cite[Thm.~8.2]{deligne_weil1}\VERIFYloc, \cite[\S4.5--4.6, \S5.2]{diamond_shurman}\VERIFYloc}
(given the eigenvalue bound of~(3) as the input, as in Proposition~\ref{prop:crtcharpoly}; v1.23.0-pre.24; clause C-D4 of the plan; refereed ACCEPT-with-fixes, its fixes at
v1.23.0-pre.25.) Let $N\ge1$,
$n\ge0$, $p\nmid N$ a prime, $T_p=[AgA]$ (Proposition~\ref{prop:strictliftN}(1)), and $\chi^{(N)}_{p,q}:=\det\bigl(x-T_p\mid\Hfr{q}(N)\bigr)\in\Z[x]$ ($T_p$ preserves torsion).
\begin{enumerate}
\item \textbf{(Only degree $1$ carries anything.)} $\Hfr{q}(N)=0$ for $q\ge2$; $\Hfr{0}(N)=H^{0}(A;\Vcoef)$ is $0$ for $n>0$ and $\Z$ for $n=0$, where $T_p=p+1$. So
  $\chi^{(N)}_{p,q}=1$ for $q\ge2$, $\chi^{(N)}_{p,0}=x-(p+1)$ if $n=0$ and $1$ otherwise, and $\chi^{(N)}_{p,1}$ is monic of degree $r_N:=b_1(M_N)$, the rank of
  Proposition~\ref{prop:descentN}(7)(b).
\item \textbf{(Residues on the $\mathsf O$-piece.)} Let $\ell'\ge5$ be a prime with $s_{\ell'}(M_N)=0$ (Proposition~\ref{prop:badprimesN}(1), (2); every
  $\ell'>n+1$ qualifies; $\ell'=p$ is admissible --- neither the argument nor the model needs $\ell'\ne p$ --- and the code may exclude it, v1.23.0-pre.25, referee n1). Reduction $\Vcoef\to\Vcoef\otimes\F_{\ell'}$ induces a $T_p$-equivariant isomorphism $\Hfr{1}(N)\otimes\F_{\ell'}\cong H^{1}(A;\Vcoef\otimes\F_{\ell'})$, so
  $\chi^{(N)}_{p,1}\bmod\ell'$ is the characteristic polynomial of $T_p$ on $H^{1}(\Ceq(M_{N,\F_{\ell'}}))$ --- the shipped level-$N$ operator, or its transport to the tree complex
  $\mathrm T(M_{N,\F_{\ell'}})$ (Proposition~\ref{prop:strictliftN}(6)(a)), or, for $n$ even, the endomorphism $[H_{N,p}]$ induced on the \textbf{cokernel presentation}
  $M_{N,\F_{\ell'}}/\mathrm{Rel}_N$ (Proposition~\ref{prop:heilbronn}(1), (2), through $\varkappa_N$ of Definition~\ref{def:kappaN}; v1.23.0-pre.31, D3's contract note~1: the route the code
  reads); for $n$ odd, $r_N=0$ and nothing is read. \textbf{No $2$- or $3$-adic input enters $\chi^{(N)}_{p,1}$.}
\item \textbf{(The bound --- the input.)} Every complex eigenvalue $\lambda$ of $T_p$ on $\Hfr{1}(N)\otimes\Cx$ satisfies $\abs\lambda\le\mathsf B_p:=1+p^{\,n+1}$. Hence the
  coefficient $c_i$ of $x^{r_N-i}$ in $\chi^{(N)}_{p,1}$ satisfies $\abs{c_i}\le B_{p,i}:=\binom{r_N}{i}\mathsf B_p^{\,i}\le2^{\mathsf b_{p,i}}$ with the word-size integer
  $\mathsf b_{p,i}:=r_N+i\,(n+1)\,\mathrm{bitlen}(p)$, $\mathrm{bitlen}(p)=\lfloor\log_2p\rfloor+1$. The bound is attained, by the Eisenstein eigenvalue $1+p^{\,n+1}$,
  whenever $r_N>0$, i.e.\ $n$ even and $(N,n)\ne(1,0)$ ($r_N=0$ for $n$ odd; v1.23.0-pre.25, referee t1).
\item \textbf{(Recovery in words.)} Let $\mathcal S$ be a finite set of primes as in (2), $P:=\prod_{\ell'\in\mathcal S}\ell'$, and $\rho_{\ell'}\in\F_{\ell'}[x]$ the residues of~(2).
  If $\sum_{\ell'\in\mathcal S}(\mathrm{bitlen}(\ell')-1)\ge\mathsf b_{p,i}+1$, then $P>2B_{p,i}$ and $c_i$ is the unique integer in $(-P/2,P/2]$ with $c_i\equiv\rho_{\ell',i}$
  for every $\ell'\in\mathcal S$. It is computed without a big integer, \textbf{provided every $\ell'\in\mathcal S$ is $<2^{32}$} (v1.23.0-pre.25, referee m2): Garner's
  mixed-radix digits $d_k<\ell'_k$ of the representative $c^{+}\in[0,P)$, each digit a word operation (its products are $<\ell'^{\,2}_k<2^{64}$); the sign by comparing the
  digits, from the top, with those of $(P-1)/2$, which are $(\ell'_k-1)/2$; $\abs{c_i}$ in mixed radix by digit-wise subtraction from $P$; and $c_i\bmod Q$, for any word-size
  $Q\le2^{64}$ (for instance $2^{64}$, $3^{40}$, $\ell^{K}$), by Horner's rule on the digits modulo $Q$, \textbf{with a word mulmod}: the product of an accumulator $<Q$ and a
  radix $\ell'_k<2^{32}$ exceeds the word for $Q>2^{32}$ (about $2^{91}$ at $Q=3^{40}$), so it is formed by shift-and-add or Montgomery multiplication modulo $Q$, never by a
  double-word integer (Remark~\ref{rmk:integralscope}(1)). \textbf{The negative case}: if $c_i=c^{+}-P<0$, then $c_i\bmod Q=(c^{+}\bmod Q-P\bmod Q)\bmod Q$, with $P\bmod Q$ by
  the same Horner rule on the radices (equivalently: Horner on the digits of $\abs{c_i}$, then negated modulo $Q$).
\item \textbf{(The output rule, [R-Q2].)} Per coefficient: \emph{certified} (the inequality of (4) holds) and $\abs{c_i}<2^{63}$: the integer; \emph{certified} and
  $\abs{c_i}\ge2^{63}$ (decided on the digits): the sign, the mixed-radix digits and their primes, and $\mathsf b_{p,i}$ --- determined, not printed in decimal;
  \emph{not certified}: the residue vector $(\rho_{\ell',i})_{\ell'\in\mathcal S}$ and $\mathsf b_{p,i}$, flagged undetermined, \textbf{never} a lifted integer. $\chi^{(N)}_{p,1}$ is
  reported as determined iff every coefficient is certified. A held-out prime is a consistency check, not a certificate.
\item \textbf{(Other moduli.)} Residues of $\chi^{(N)}_{p,1}$ modulo $\ell^{k}$, $\ell$ any prime, read on an $\ell$-local lattice known modulo $\ell^{k}$
  (Proposition~\ref{prop:torsionhecke}(4)), may join $\mathcal S$ as further pairwise coprime moduli (at $\ell=2$ with the sign test adapted to an even $P$); this is admissible
  and never needed.
\end{enumerate}
At $N=1$ this is Proposition~\ref{prop:crtcharpoly} with its bound made explicit.
\end{proposition}

\begin{proof}
(1) $24\cdot H^{q}(\Ga;M_N)=0$ for $q\ge2$ (Proposition~\ref{prop:certifiedeell}(4)(a)), and Shapiro. $H^{0}(A;\Vcoef)\subseteq\Vcoef$ is torsion-free, of rank
$b_0(M_N)$, which is $1$ for $n=0$ and $0$ for $n>0$ (Proposition~\ref{prop:descentN}(4), (7)(c)). For $n=0$ the coefficients are trivial and $\rho_0(g)=1$, so on $H^{0}$
$\res$ and $(\iota_g,1)^{\ast}$ are the identity of $\Z$ and $\cor^{A}_{A''}$ is multiplication by $[A:A'']=\abs{\mathbb P^{1}(\F_p)}=p+1$ \STD{\cite[Ch.~III.9]{brown}}.

(2) The long exact sequence of $0\to\Vcoef\xrightarrow{\ell'}\Vcoef\to\Vcoef\otimes\F_{\ell'}\to0$ over $A$ gives
$0\to H^{1}(A;\Vcoef)\otimes\F_{\ell'}\to H^{1}(A;\Vcoef\otimes\F_{\ell'})\to H^{2}(A;\Vcoef)[\ell']\to0$; the right-hand term vanishes ($24$ kills it, $\ell'\ge5$), and
$H^{1}(A;\Vcoef)$ has no $\ell'$-torsion when $s_{\ell'}(M_N)=0$ (Proposition~\ref{prop:badprimesN}(1)), so the left-hand term is $\Hfr{1}(N)\otimes\F_{\ell'}$. The map is
induced by the coefficient map, which commutes with $\rhon(g)$, so it commutes with each factor of $[AgA]$ (naturality of $\res$, $\cor$ and of $(\iota_g,f)^{\ast}$ in compatible
pairs, \STD{\cite[Ch.~III.8--9]{brown}}). The model statements are Proposition~\ref{prop:strictliftN}(6)(a) at $\cring=\F_{\ell'}$. The presentation (v1.23.0-pre.31): for $n$ even
and $\cring=\F_{\ell'}$, $6\in\cring^{\times}$ as $\ell'\ge5$, so $\varkappa_N\colon M_{N,\cring}/\mathrm{Rel}_N\to H^{1}(A;\VR)$ is an isomorphism (Definition~\ref{def:kappaN}), $H_{N,p}$ preserves
$\mathrm{Rel}_N$ and $\varkappa_N\circ[H_{N,p}]=[A\,g\,A]\circ\varkappa_N$ (Proposition~\ref{prop:heilbronn}(1), (2)), with $g=\smat100p$ there as here and $[A\,g\,A]=T_p$ for $p\nmid N$; conjugate
endomorphisms have the same characteristic polynomial. For $n$ odd, $-I\in A$ is central and acts on $\Vcoef$ by $-1$, so $2$ kills $H^{\bullet}(A;\Vcoef)$ (the argument of
Proposition~\ref{prop:levelorbit}(2)), $r_N=0$ and $\chi^{(N)}_{p,1}=1$.

(3) By Eichler--Shimura \STD{\cite[Thm.~12.2.2, p.~108]{diamond_im}, verified (Lemma~\ref{lem:heckerankone})} $H^{1}(A;\Vcoef\otimes\Cx)\cong M_{k}(\Ga_0(N))\oplus\bar S_{k}(\Ga_0(N))$,
$k=n+2$, for every $n\ge0$, compatibly with the Hecke ring \STD{\cite[Prop.~12.4.10, p.~117]{diamond_im}, verified}, $T_p$ acting on $\bar S_k$ through complex conjugation; and
$\Hfr{1}(N)\otimes\Cx=H^{1}(A;\Vcoef\otimes\Cx)$. For $p\nmid N$: on $S_k(\Ga_0(N))$ the eigenvalues of $T_p$ are the $a_p$ of normalised newforms of levels dividing $N$ \STD{\cite[\S5.6--5.8]{diamond_shurman}, Atkin--Lehner--Li, \VERIFYloc}, and
$\abs{a_p}\le2p^{(k-1)/2}$ \STD{\cite[Thm.~8.2]{deligne_weil1}, the Ramanujan--Petersson bound, \VERIFYloc}; on the Eisenstein subspace they are $\psi(p)+\varphi(p)p^{\,k-1}$ with Dirichlet
characters $\psi$, $\varphi$ (the eigenforms $E_k^{\psi,\varphi,t}$, and for $k=2$ the differences of $E_2$, eigenvalue $1+p$) \STD{\cite[\S4.5--4.6, \S5.2]{diamond_shurman}, \VERIFYloc}, of
absolute value $\le1+p^{\,k-1}$. As $2p^{(n+1)/2}\le1+p^{\,n+1}$, every $\abs\lambda\le\mathsf B_p$. The coefficient $c_i$ is $\pm$ the $i$-th elementary symmetric function of the
$r_N$ eigenvalues, a sum of $\binom{r_N}{i}$ products of $i$ of them. Bits: $\binom{r_N}{i}\le2^{r_N}$, and $p\le2^{\mathrm{bitlen}(p)}-1$ gives
$1+p^{\,n+1}\le2^{(n+1)\mathrm{bitlen}(p)}$. Sharpness: $E_k$ ($k\ge4$), resp.\ a difference of $E_2$ at $N>1$, contributes $1+p^{\,n+1}$.

(4) $\mathrm{bitlen}(\ell')-1=\lfloor\log_2\ell'\rfloor$, so $P\ge2^{\sum\lfloor\log_2\ell'\rfloor}\ge2^{\mathsf b_{p,i}+1}>2B_{p,i}$; two integers in $(-P/2,P/2]$ congruent modulo every
$\ell'$ are equal (Chinese remainder theorem), and $c_i$ is one of them. Garner: $c^{+}:=\sum_kd_k\prod_{j<k}\ell'_j$ with $d_k:=(\rho_{\ell'_k,i}-\sum_{m<k}d_m\prod_{j<m}\ell'_j)\cdot(\prod_{j<k}\ell'_j)^{-1}
\bmod\ell'_k$, every quantity reduced modulo $\ell'_k$ before it is formed, so no intermediate exceeds $\ell'^{\,2}_k$. Mixed-radix representations with digits $d_k<\ell'_k$ are unique
and ordered lexicographically from the top, and $(P-1)/2=\sum_k\frac{\ell'_k-1}2\prod_{j<k}\ell'_j$ (all $\ell'_k$ odd; the sum telescopes); $c_i=c^{+}$ if $c^{+}\le(P-1)/2$ and $c^{+}-P$
otherwise. Horner modulo $Q$ forms only residues modulo $Q$ and products of two of them, which the word mulmod reduces without leaving the word; $c^{+}-P\equiv c^{+}\bmod Q-P\bmod Q$.

(5) is the reading rule of~(4); it forms no integer beyond a word (Theorem~\ref{thm:backend}(3), Remark~\ref{rmk:integralscope}(1)). (6): $\chi\bmod\ell^{k}=\det(x-X)\bmod\ell^{k}$ for any
matrix $X$ of $T_p$ on $\Hfr{1}(N)\otimes\Z_{(\ell)}$ (base change of a characteristic polynomial), and the Chinese remainder theorem holds for pairwise coprime moduli.
\end{proof}

\begin{proposition}[Hecke on finite $\ell$-groups: what is choice-free, and how it is read exactly from a local run]\label{prop:torsionhecke}
\PROVED{} (v1.23.0-pre.24; clause C-D5 of the plan; refereed ACCEPT-with-fixes, its fixes M1, m1 at v1.23.0-pre.25.) Let $\ell$ be a prime, $G$ a finite abelian $\ell$-group, $\mathrm{len}(G):=\log_\ell\abs G$, and
$\phi\in\End(G)$. Write $G=\bigoplus_i\Z/\ell^{a_i}e_i$ and $\phi e_k=\sum_i\phi_{ik}e_i$, $\phi_{ik}\in\Z/\ell^{a_i}$ (so $\ell^{a_k}\phi_{ik}=0$ in $\Z/\ell^{a_i}$).
\begin{enumerate}
\item \textbf{(Choice-free invariants.)} Each of the following depends only on the isomorphism class of the $\Z[x]$-module $(G,\phi)$:
  \begin{enumerate}
  \item[(a)] \textbf{the layers}: for $j\ge1$, $V_j(G):=(\ell^{j-1}G)[\ell]\big/(\ell^{j}G)[\ell]$, an $\F_\ell$-space of dimension $\#\{i:a_i=j\}$, with the endomorphism $\phi_j$
    induced by $\phi$, and $P_j(\phi):=\det(x-\phi_j)\in\F_\ell[x]$. In the coordinates above, $\phi_j$ is the block $(\phi_{ik}\bmod\ell)_{a_i=a_k=j}$;
  \item[(b)] \textbf{the Fitting components}: for $c\in\Z$, $G_{[c]}:=\ker(\phi-c)^{\mathrm{len}(G)}$, which depends only on $c\bmod\ell$; $G=\bigoplus_{\bar c\in\F_\ell}G_{[c]}\oplus G_{\circ}$,
    $G_\circ:=\bigcap_{c}\operatorname{im}(\phi-c)^{\mathrm{len}(G)}$, all summands $\phi$-stable; and $G_{[c]}$ has exactly $\operatorname{mult}_{x-\bar c}P_j(\phi)$ cyclic summands of
    order $\ell^{j}$, so its type is read off the $P_j$;
  \item[(c)] the types of $\ker(\phi-c)$ and $\coker(\phi-c)$, $c\in\Z$ (they depend on $c$ modulo the exponent of $G$);
  \item[(d)] for commuting $\phi_1,\dots,\phi_m$, the isomorphism class of the $\Z[x_1,\dots,x_m]$-module $G$, and every invariant derived from it;
  \item[(e)] \textbf{the layer blocks modulo $\ell^{e_j}$} (v1.23.0-pre.25, referee M1): with $d_j:=\min\{\abs{a_i-j}:a_i\ne j\}$ ($d_j:=\infty$ if $G$ has one layer) and
    $e_j:=\min(j,d_j)$, the characteristic polynomial modulo $\ell^{e_j}$ of the diagonal block $(\phi_{ik})_{a_i=a_k=j}$; its reduction modulo $\ell$ is $P_j(\phi)$.
  \end{enumerate}
\item \textbf{(Choice-free only modulo $\ell^{e_j}$.)} The characteristic polynomial over $\Z/\ell^{j}$ of the diagonal block $(\phi_{ik})_{a_i=a_k=j}$ --- ``the characteristic
  polynomial on the Smith summands of order $\ell^{j}$'' --- depends on the decomposition: on $G=\Z/\ell\,e_1\oplus\Z/\ell^{2}e_2$, $\phi e_1=\ell e_2$, $\phi e_2=0$, it is $x$,
  and in the basis $e_1$, $e_1+e_2$ it is $x-\ell$. It is an invariant modulo $\ell^{e_j}$ (1)(e), and in general no finer (v1.23.0-pre.25, referee M1, replacing ``only its
  reduction modulo $\ell$''): for every $1\le d<j$, on $G=\Z/\ell^{j-d}e_1\oplus\Z/\ell^{j}e_2$ with $\phi e_1=\ell^{d}e_2$, $\phi e_2=0$ the layer-$j$ block is $(0)$ in the
  basis $(e_1,e_2)$ and $(\ell^{d})$ in the basis $(e_1,e_1+e_2)$, and on $G=\Z/\ell^{j+d}e_1\oplus\Z/\ell^{j}e_2$ with $\phi e_2=\ell^{d}e_1$, $\phi e_1=0$ it is $(0)$ and
  $(-\ell^{d})$ in the basis $(e_1+e_2,e_2)$; in both $d_j=d=e_j$, and the two polynomials differ modulo $\ell^{d+1}$ (the example above is $j=2$, $d=1$). If $d_j\ge j$ the
  invariance modulo $\ell^{j}$ is the whole block's. Modulo $\ell$ it is $P_j(\phi)$; if $G$ has one layer ($G$ free over $\Z/\ell^{a}$), the characteristic polynomial over
  $\Z/\ell^{a}$ is an invariant. (For instance $G=\Z/2\oplus\Z/8$: the layer-$3$ polynomial is an invariant modulo $4$, not modulo $8$.)
\item \textbf{(The exact read from a local run.)} Let $C^{\bullet}$ be a complex of finitely generated free $\Z_{(\ell)}$-modules, $q\ge1$, $F$ a cochain endomorphism of $C^{\bullet}$
  (over $\Z_{(\ell)}$), $G:=\mathrm{tors}\,H^{q}(C)$ and $\phi:=H^{q}(F)|_G$. Let $L\,(d^{q-1}\bmod\ell^{K})\,R=\Sigma$ be a Smith form over $\Z/\ell^{K}$ with pivots
  $u_i\ell^{\sigma_i}$, certified ($s=\rank_{\Q}d^{q-1}$, Proposition~\ref{prop:descentmaps}(2)(g)), $J:=\{i\le s:\sigma_i>0\}$, $\sigma_{\max}:=\max_J\sigma_i$, $\tilde x_j$ the
  $j$-th column of $L^{-1}$ ($j\in J$), and $\tilde\Phi_{ij}:=(L\,F\,\tilde x_j)_i\bmod\ell^{\sigma_i}$ ($i,j\in J$). If $2\sigma_{\max}\le K$ (i.e.\
  $\sigma_i\le K-\sigma_{\max}$ for every $i\in J$), then $G=\bigoplus_{j\in J}\Z/\ell^{\sigma_j}[x_j]$ for cocycles $x_j\equiv\tilde x_j\pmod{\ell^{K-\sigma_{\max}}}$, and
  $\phi[x_j]=\sum_{i\in J}\tilde\Phi_{ij}[x_i]$ exactly: $(G,\phi)\cong(\bigoplus_J\Z/\ell^{\sigma_i},\tilde\Phi)$, and every invariant of (1) computed from $(\sigma,\tilde\Phi)$ is
  certified. Otherwise the read is refused (or $K$ raised; at the cap of the ring, refused). It applies $F$ to the $\abs J$ vectors $\tilde x_j$ only.
  \textbf{The condition is necessary} for the module (v1.23.0-pre.25, referee m1): for every $\sigma\ge1$ and $\sigma<K<2\sigma$ there is a certified form whose read is wrong (the
  proof), and on real data $K=2\sigma_{\max}-1$ gives a wrong module in $33$ (cell, $p$, $\ell$) cases of the referee's check. It is \emph{not} the conservative policy of
  Proposition~\ref{prop:descentmaps}(2)(g)(ii), whose ``exact for every $\sigma_i<K$'' concerns the coordinates of an \emph{exact} vector: here the vector is $F\tilde x_j$, and
  $\tilde x_j$ is exact only modulo $\ell^{K'}$, $K':=K-\sigma_{\max}$; a read may not be relaxed by analogy with \code{BuildPres}. \textbf{What needs only a certified form}
  ($K>\sigma_{\max}$, so $K'\ge1$): $\tilde\Phi_{ij}$ agrees with the true coordinate modulo $\ell^{\min(\sigma_i,K')}$, so the layer polynomials $P_j$, through (1)(b) the types of
  the Fitting components $G_{[c]}$, and the layer blocks' polynomials modulo $\ell^{\min(e_j,K')}$ are certified at \textbf{every} certified $K$; the module, (1)(c), (d) and
  (1)(e) beyond $\ell^{K'}$ need $2\sigma_{\max}\le K$. The read uses of the complex only $d^{q-1}$ and $F$ on $C^{q}$.
  \textbf{(3$'$) A presentation in place of a complex} (v1.23.0-pre.31, D3's contract note~1). Let $\Delta$ be an $r\times g$ matrix over $\Z_{(\ell)}$, $Q:=\coker\Delta=\Z_{(\ell)}^{r}/\operatorname{im}\Delta$,
  and $F\in\End(\Z_{(\ell)}^{r})$ with $F(\operatorname{im}\Delta)\subseteq\operatorname{im}\Delta$; put $G:=\mathrm{tors}\,Q$, $\phi$ the endomorphism of $G$ induced by $F$, $\Lambda:=Q/G$ and
  $\psi$ the endomorphism of $\Lambda$ induced by $F$. Then everything in~(3) holds verbatim with $d^{q-1}:=\Delta$, $C^{q}:=\Z_{(\ell)}^{r}$ and the certifying rank $s=\rank_{\Q}\Delta$:
  the same read $L(\Delta\bmod\ell^{K})R=\Sigma$, $\tilde\Phi_{ij}=(LF\tilde x_j)_i\bmod\ell^{\sigma_i}$, the same precisions, the same condition $2\sigma_{\max}\le K$ for the module.
  \textbf{Moreover the free part is read}: with $\tilde x_j$ the $j$-th column of $L^{-1}$, $s<j\le r$, the classes of \emph{any} exact lifts $x_j\equiv\tilde x_j\pmod{\ell^{K}}$ form a basis of
  $\Lambda$ (one such is $x_j:=\tilde L^{-1}e_j$, the proof), and $X_{ij}:=(LF\tilde x_j)_i\bmod\ell^{K'}$, $s<i,j\le r$, is the matrix of $\psi$ in that basis modulo $\ell^{K'}$, at every certified $K$ --- an input of~(4) with
  $k=K'$ (v1.23.0-pre.32, referee m3). The precision $K'$ is \textbf{sharp} (referee N3): the rows $>s$ of $\tilde L$ are exact only modulo $\ell^{K'}$, and on the referee's certified
  instances the read differs from the exact matrix at exactly $\ell^{K'}$; unlike the torsion coordinates of Proposition~\ref{prop:descentmaps}(2)(g)(ii), it may not be relaxed. \textbf{The read is presentation-independent}: if $(\Delta',F')$ is another such datum and $\theta\colon Q\to Q'$ an isomorphism with $\theta\circ[F]=[F']\circ\theta$ (for instance
  two presentations of one module with one operator, or a presentation and a complex $C^{\bullet}$ as in~(3) with $H^{q}(C)\cong Q$ compatibly with the operators), then $\theta$ restricts to an
  isomorphism $(G,\phi)\cong(G',\phi')$ and induces $(\Lambda,\psi)\cong(\Lambda',\psi')$, so every invariant of~(1) and the invariant factors of~(4) agree, each computed from a read that certifies it --- the module and (1)(c), (d), (e) from a read with
  $2\sigma_{\max}\le K$, the rest as in~(3), the factors of~(4) when its pivot condition holds (v1.23.0-pre.32, referee N1).
\item \textbf{(The free quotient: the invariant factors of $\phi-c$.)} Let $\Lambda$ be a free $\Z_{(\ell)}$-module of rank $r$ with an endomorphism $\psi$, $\chi_\psi\in\Z_{(\ell)}[x]$
  its characteristic polynomial, $X$ its matrix in some basis, known modulo $\ell^{k}$ (e.g.\ $\Lambda=\Hfr{1}(N)\otimes\Z_{(\ell)}$, $\psi=T_p$, $X$ read on a local run with the
  precision left by Proposition~\ref{prop:descentmaps}(2)(g)(iii)), and $c\in\Z$. If the Smith form of $(X-c)\bmod\ell^{k}$ has $r$ pivots of valuation $<k$, these valuations are
  the exponents of the invariant factors of $\psi-c$ on $\Lambda$, i.e.\ the type of the finite $\Z_{(\ell)}[x]$-module $\coker(\psi-c)$ --- choice-free. This happens iff $k$
  exceeds the largest exponent; the exponents sum to $v_\ell(\chi_\psi(c))$, so $k>v_\ell(\chi_\psi(c))$ suffices when $\chi_\psi(c)\ne0$. If $\chi_\psi(c)=0$ it never happens,
  and nothing is asserted.
\end{enumerate}
\end{proposition}

\begin{proof}
(1)(a) $\ell^{m}G$ and $G[\ell]$ are $\phi$-stable, so $\phi$ induces $\phi_j$, and an isomorphism of $\Z[x]$-modules induces isomorphisms of these subquotients intertwining the
$\phi_j$. In coordinates, $(\ell^{j-1}G)[\ell]=\bigoplus_{a_i\ge j}\Z\,\ell^{a_i-1}e_i$, so the classes of $\ell^{j-1}e_k$, $a_k=j$, form a basis of $V_j(G)$, and
$\phi(\ell^{j-1}e_k)=\sum_i\ell^{j-1}\phi_{ik}e_i$: the terms with $a_i<j$ vanish; for $a_i>j$, $\ell^{a_k}\phi_{ik}=0$ in $\Z/\ell^{a_i}$ gives $v_\ell(\phi_{ik})\ge a_i-j$, so
$\ell^{j-1}\phi_{ik}e_i\in(\ell^{a_i-1}G)[\ell]\subseteq(\ell^{j}G)[\ell]$ vanishes in $V_j$; for $a_i=j$ it is $(\phi_{ik}\bmod\ell)\,\ell^{j-1}e_i$.
(b) \emph{Fitting's lemma.} For $\psi\in\End(G)$ the chains $\ker\psi^{m}$, $\operatorname{im}\psi^{m}$ are monotone in the finite group $G$, and each strict step changes the order
by a factor $\ge\ell$, so both are stationary from $m=\mathrm{len}(G)$ on; then $G=\ker\psi^{L}\oplus\operatorname{im}\psi^{L}$ ($L:=\mathrm{len}(G)$): if $y=\psi^{L}x$ and
$\psi^{L}y=0$ then $x\in\ker\psi^{2L}=\ker\psi^{L}$, so $y=0$, and $\abs{\ker\psi^{L}}\abs{\operatorname{im}\psi^{L}}=\abs G$. $\psi$ is nilpotent on the first summand and bijective
on the second. Take $\psi=\phi-c$. If $c'\equiv c\pmod\ell$, $\phi-c'=(\phi-c)-(c'-c)$ with $c'-c$ nilpotent on the $\ell$-group $G$ and central, so $\phi-c'$ is nilpotent on
$G_{[c]}$ and $G_{[c]}\subseteq G_{[c']}$, and conversely. If $c'\not\equiv c$, $c-c'$ is a unit on $G$ and $\phi-c'=(\phi-c)+(c-c')$ is bijective on $G_{[c]}$, so $G_{[c]}\cap G_{[c']}=0$
and $G_{[c]}\subseteq\operatorname{im}(\phi-c')^{L}$. Applying Fitting's lemma to $\phi-c_2$ on $\operatorname{im}(\phi-c_1)^{L}$, and so on over the classes of $\F_\ell$, gives the decomposition,
with $G_\circ$ the last image. The functor $V_j$ is additive on $\phi$-stable direct sums; on $V_j(G_{[c]})$, $\phi_j-\bar c$ is nilpotent, and on $V_j(G_{[c']})$, $c'\not\equiv c$,
and on $V_j(G_\circ)$ it is induced by an automorphism, hence bijective (an automorphism maps $(\ell^{m}G)[\ell]$ onto itself). So
$\operatorname{mult}_{x-\bar c}P_j(\phi)=\dim V_j(G_{[c]})$, the number of summands of order $\ell^{j}$ of $G_{[c]}$ by~(a).
(c), (d) Kernels, cokernels and module structure are carried by isomorphisms of modules. $\phi-c$ depends on $c$ modulo the exponent.
(e) A homomorphism $\Z/\ell^{b}\to\Z/\ell^{b'}$ is multiplication by an integer $m$, defined modulo $\ell^{b'}$, with $\ell^{b}m\equiv0\pmod{\ell^{b'}}$, so
$v_\ell(m)\ge\max(0,b'-b)$; composites multiply. A second decomposition of the same type, ordered alike, is $e'_k=he_k$ for an automorphism $h$, and the matrix of $\phi$ in it is
$h^{-1}\phi h$ (entries of the $i$-th row reduced modulo $\ell^{a_i}$). An entry of the layer-$j$ block of a product of such matrices is a sum over paths
$j=b_0\to b_1\to\dots\to b_m=j$ of layers; a term is divisible by $\ell$ to the sum of the ascents $\max(0,b_{t+1}-b_t)$, and a closed path from $j$ that visits a layer $a\ne j$
ascends by at least $\abs{a-j}\ge d_j$ (for $a>j$ on the way up, for $a<j$ on the way back). Hence, writing $X_j$ for the layer-$j$ block and reading in $\Z/\ell^{j}$,
$(h^{-1}\phi h)_j\equiv(h^{-1})_j\,\phi_j\,h_j$ and $(h^{-1})_jh_j\equiv(h^{-1}h)_j=1\pmod{\ell^{e_j}}$: $h_j$ is invertible modulo $\ell^{e_j}$ and the block changes by a
conjugation modulo $\ell^{e_j}$, which leaves its characteristic polynomial unchanged. Its reduction modulo $\ell$ is $\det(x-\phi_j)=P_j(\phi)$ by~(a).

(2) $\phi(e_1+e_2)=\ell e_2=\ell(e_1+e_2)$, as $\ell e_1=0$; so in the basis $(e_1,e_1+e_2)$ --- orders $\ell$, $\ell^{2}$, a basis of the same type --- the $\Z/\ell^{2}$-block is $(\ell)$,
against $(0)$ before. The general pairs: in $\Z/\ell^{j-d}e_1\oplus\Z/\ell^{j}e_2$ the basis $(e_1,e_1+e_2)$ has the same type (orders $\ell^{j-d}$, $\ell^{j}$; it generates), and
$\phi(e_1+e_2)=\ell^{d}e_2=\ell^{d}(e_1+e_2)-\ell^{d}e_1$, whose coefficient on the layer-$j$ generator is $\ell^{d}$; in $\Z/\ell^{j+d}e_1\oplus\Z/\ell^{j}e_2$ the basis
$(e_1+e_2,e_2)$ has orders $\ell^{j+d}$, $\ell^{j}$, and $\phi e_2=\ell^{d}e_1=\ell^{d}(e_1+e_2)-\ell^{d}e_2$, coefficient $-\ell^{d}$ (both $\phi$ are endomorphisms: $\ell^{j-d}\cdot\ell^{d}=\ell^{j}$,
$\ell^{j}\cdot\ell^{d}=\ell^{j+d}$). As $d<j$, $x$ and $x\mp\ell^{d}$ differ modulo $\ell^{d+1}$. One layer: a change of decomposition is then a change of basis of a free
$\Z/\ell^{a}$-module, under which the characteristic polynomial is invariant.

(3) The certified form gives $\mathrm{tors}\coker d^{q-1}\cong\bigoplus_J\Z/\ell^{\sigma_i}$, with exact transforms $\tilde L$, $\tilde R$ over $\Z_{(\ell)}$, the first $s$ rows of $\tilde L$
congruent to those of $L$ modulo $\ell^{K}$ and the first $s$ columns of $\tilde L^{-1}$ congruent to those of $L^{-1}$ modulo $\ell^{K'}$, $K':=K-\sigma_{\max}$
(Proposition~\ref{prop:descentmaps}(2)(g)(i)); the classes of $x_j:=\tilde L^{-1}e_j$ generate, with coordinates $y\mapsto(\tilde Ly)_i\bmod\ell^{\sigma_i}$ (ibid.\ (ii)). As $C^{q+1}$
is torsion-free, $\mathrm{tors}\,H^{q}(C)=\mathrm{tors}\coker d^{q-1}$ and the $x_j$ are cocycles (the proof of Proposition~\ref{prop:certifiede}(2)). $F$ is $\Z_{(\ell)}$-linear and
preserves torsion classes, and
\[
  (L\,F\,\tilde x_j)_i-(\tilde L\,F\,x_j)_i=\bigl((L-\tilde L)F\tilde x_j\bigr)_i+\bigl(\tilde L\,F(\tilde x_j-x_j)\bigr)_i\equiv0\pmod{\ell^{K'}}
\]
for $i\le s$ (the first term is $\equiv0\bmod\ell^{K}$, the second $\bmod\ell^{K'}$). If $\sigma_i\le K'$, this is $0$ modulo $\ell^{\sigma_i}$, so $\tilde\Phi_{ij}$ is the $i$-th coordinate of
$\phi[x_j]=[Fx_j]$. For every $i$, the same congruence gives $\tilde\Phi_{ij}\equiv$ the true coordinate modulo $\ell^{\min(\sigma_i,K')}$, with $K'\ge1$ for a certified form;
so the layer blocks are read exactly modulo $\ell^{\min(j,K')}$, which gives the $P_j$ (modulo $\ell$), the Fitting types by (1)(b), and the block polynomials modulo
$\ell^{\min(e_j,K')}$ by (1)(e). \emph{The condition is necessary} (v1.23.0-pre.25, referee m1). Proposition~\ref{prop:descentmaps}(2)(g)(ii)'s exactness of torsion coordinates for every
$\sigma_i<K$ is for an exact vector $y$; here $y=F\tilde x_j$ and the error $\ell^{K'}w$ of $\tilde x_j$ passes through $F$. The minimal instance: $C^{q-1}=C^{q}=\Z_{(\ell)}^{2}$,
$C^{q+1}=0$, $d^{q-1}=\mathrm{diag}(\ell^{\sigma},0)$, $F^{q-1}=0$ and $F^{q}e_1=0$, $F^{q}e_2=e_1$ (a cochain map), so $G=\Z/\ell^{\sigma}[e_1]$ and $\phi=0$; for $\sigma<K$ the form
with $L=\smat10{\ell^{K-\sigma}}1$, $R=1$ is certified ($Ld^{q-1}=\smat{\ell^{\sigma}}0{\ell^{K}}0\equiv d^{q-1}\bmod\ell^{K}$), $\tilde x_1=(1,-\ell^{K-\sigma})^{T}$, and
$\tilde\Phi_{11}=(LF\tilde x_1)_1=-\ell^{K-\sigma}$, which is $\ne0$ modulo $\ell^{\sigma}$ exactly when $K<2\sigma$.

(3$'$) (v1.23.0-pre.31.) \emph{Reduction to~(3).} Let $C^{\bullet}$ be the two-term complex $C^{q-1}:=\Z_{(\ell)}^{g}\xrightarrow{\ \Delta\ }C^{q}:=\Z_{(\ell)}^{r}$, all other terms $0$. Then
$H^{q}(C)=Q$, and $C^{q+1}=0$ is torsion-free. Put $F^{q}:=F$ and $F^{m}:=0$ for $m\ne q-1,q$; for each basis vector $e_k$ of $\Z_{(\ell)}^{g}$, $F\Delta e_k\in\operatorname{im}\Delta$, so
choose $F^{q-1}e_k$ with $\Delta F^{q-1}e_k=F\Delta e_k$ and extend linearly ($\Z_{(\ell)}^{g}$ is free). Then $\Delta F^{q-1}=F\Delta$, and $d^{q}=0$, so $(F^{m})$ is a cochain
endomorphism of $C^{\bullet}$ over $\Z_{(\ell)}$, and $H^{q}(F^{\bullet})$ is the endomorphism $[F]$ of $Q$; its restriction to $G=\mathrm{tors}\,H^{q}(C)$ is $\phi$. The read of~(3) uses only
$d^{q-1}=\Delta$ and $F^{q}=F$, never the chosen $F^{q-1}$; so it is the read of~(3) for this complex, and~(3), with its proof, applies verbatim (the $x_j$ are cocycles trivially, $d^{q}=0$).
\emph{The free part.} By Proposition~\ref{prop:descentmaps}(2)(g)(i), there are exact $\tilde L$, $\tilde R$ with $\tilde L\Delta\tilde R=\mathrm{diag}(D,0)$, the last $r-s$ columns of
$\tilde L^{-1}$ congruent to those of $L^{-1}$ modulo $\ell^{K}$ (its ``moreover'') and all of $\tilde L$ congruent to $L$ modulo $\ell^{K'}$. By (2)(g)(ii), $y\mapsto\tilde Ly$ identifies $Q$ with
$\bigoplus_{i\le s}\Z_{(\ell)}/\ell^{\sigma_i}\oplus\Z_{(\ell)}^{r-s}$; the first summand is $G$, so $y\mapsto((\tilde Ly)_i)_{i>s}$ identifies $\Lambda=Q/G$ with $\Z_{(\ell)}^{r-s}$, and
$x_j:=\tilde L^{-1}e_j$, $j>s$, map to the standard basis. $F$ induces $\psi$ because it preserves $\operatorname{im}\Delta$ and every endomorphism of $Q$ preserves $G$; the matrix of $\psi$
in the basis $([x_j])_{j>s}$ is $((\tilde LFx_j)_i)_{i,j>s}$. For $i,j>s$,
\[
  (LF\tilde x_j)_i-(\tilde LFx_j)_i=\bigl((L-\tilde L)F\tilde x_j\bigr)_i+\bigl(\tilde LF(\tilde x_j-x_j)\bigr)_i\equiv0\pmod{\ell^{K'}},
\]
the first term because $L\equiv\tilde L\pmod{\ell^{K'}}$, the second because $\tilde x_j\equiv x_j\pmod{\ell^{K}}$. So $X$ is the reduction modulo $\ell^{K'}$ of the exact matrix of $\psi$, which
is the input of~(4) with $k=K'$. \emph{Any exact lifts} (v1.23.0-pre.32, referee m3). Let $x'_j$, $s<j\le r$, be exact with $x'_j\equiv\tilde x_j\pmod{\ell^{K}}$. As
$x_j\equiv\tilde x_j\pmod{\ell^{K}}$ too, $x'_j=x_j+\ell^{K}w_j$ with $w_j\in\Z_{(\ell)}^{r}$, so the coordinates of $[x'_j]$ in the basis $([x_j])_{j>s}$ of $\Lambda$ form the matrix
$P=I+\ell^{K}M$, $M:=((\tilde Lw_j)_i)_{i,j>s}$, which is invertible over $\Z_{(\ell)}$ ($\det P\equiv1\bmod\ell$); so $([x'_j])$ is a basis, the matrix of $\psi$ in it is $P^{-1}X_{\mathrm{ex}}P$,
$X_{\mathrm{ex}}$ the exact matrix in the basis $([x_j])$, and $P^{-1}X_{\mathrm{ex}}P\equiv X_{\mathrm{ex}}\pmod{\ell^{K}}$ ($P\equiv P^{-1}\equiv I$); in particular $X$ is its reduction modulo
$\ell^{K'}$ as well. \emph{Independence.} $\theta$ maps torsion onto torsion, so it restricts to $G\to G'$ and induces $\Lambda\to\Lambda'$, isomorphisms intertwining the induced
maps because $\theta[F]=[F']\theta$; the invariants of~(1) depend only on the isomorphism class of $(G,\phi)$ (by~(1)), and those of~(4) only on that of $(\Lambda,\psi)$ (the type of
$\coker(\psi-c)$ is not changed by an isomorphism).

(4) Proposition~\ref{prop:certifiede}(1) for the exact matrix $X-c$ over $\Z_{(\ell)}$, whose reduction modulo $\ell^{k}$ is $(X\bmod\ell^{k})-c$: the number of pivots of valuation $<k$ is
$\#\{t:a_t<k\}\le\rank_{\Q}(X-c)\le r$; equality with $r$ forces $\rank_{\Q}(X-c)=r$ and every $a_t<k$, and then the pivots are the elementary divisors up to units. Their product is
$\det(X-c)=\pm\chi_\psi(c)$ up to a unit, so $\sum a_t=v_\ell(\chi_\psi(c))$. If $\chi_\psi(c)=0$, $\rank_{\Q}(X-c)<r$. The type of $\coker(X-c)$ does not depend on the basis.
\end{proof}

\begin{theorem}[The certified descent record: groups, characteristic polynomials over $\Z$, Hecke data on certified groups]\label{thm:descentcert}
\PROVED{} (v1.23.0-pre.24; clause C-D5 of the plan, restricted by the audit's M5 and corrected as in Remark~\ref{rmk:descentcertscope}; refereed ACCEPT-with-fixes, its fixes
M2, m5 at v1.23.0-pre.25.) Assume \ref{H4} and
\ref{B1}--\ref{B3}. Let $N\ge1$, $n\ge0$, $q_{\max}\ge3$, $\mathcal P$ a finite set of primes not dividing $N$, and $A=\Ga_0(N)$. For every $q\le q_{\max}$ the following are
determined by the descent pieces --- the $\mathsf O$-piece ($\Ceq(M_{N,\mathsf O})$ with the shipped operator or its tree transport, or the sparse route: the cokernel presentation
$M_{N,\mathsf O}/\mathrm{Rel}_N$ of Proposition~\ref{prop:levelorbit}(2) with $[H_{N,p}]$ of Proposition~\ref{prop:heilbronn}, v1.23.0-pre.31) and the $\ell$-local
pieces, $\ell\in\{2,3\}$ ($\Iso^{(\ell)}$ or the tree models, with the transported operators) ---, certified by the rules cited, and \textbf{independent of every choice} made
in computing them (models, comparisons, $\tau_X$, $\tau_{\Iso}$, tables, transversals, Smith transforms, primes):
\begin{enumerate}
\item \textbf{(Groups.)} $H^{q}(A;\Vcoef)\cong\Z^{b_q(M_N)}\oplus\bigoplus_{\ell\in\{2,3\}\cup\mathcal B(M_N)}\mathrm{tors}\,H^{q}(A;\Vcoef)_{(\ell)}$, with the certified
  exponents of Propositions~\ref{prop:descentN}(7), \ref{prop:certifiedeell}, \ref{prop:isomodeltree}(1)(d) and~\ref{prop:descentNthree}; no prime $\ell\ge5$ outside
  $\mathcal B(M_N)$ (the closed form $\mathcal B(n,N)\cap[5,\infty)$ for $n>0$ even, empty for $n=0$ and for $n$ odd) contributes, by Proposition~\ref{prop:badprimesN}(1), (2); at
  $\ell\ge5$ only $q=1$ carries torsion.
\item \textbf{(Characteristic polynomials over $\Z$.)} For $p\in\mathcal P$, $\chi^{(N)}_{p,q}$ as in Proposition~\ref{prop:crtcharpolyN}(1), with $\chi^{(N)}_{p,1}$ reported by its
  output rule (5) from the $\mathsf O$-piece alone.
\item \textbf{(Torsion.)} For each $\ell$ with $\mathrm{tors}\,H^{q}(A;\Vcoef)_{(\ell)}\ne0$, the isomorphism class of the $\Z[x_p:p\in\mathcal P]$-module
  $\bigl(\mathrm{tors}\,H^{q}(A;\Vcoef)_{(\ell)},(T_p)_{p\in\mathcal P}\bigr)$ --- in particular the layer polynomials $P_j(T_p)$, the types of the Fitting components and of
  $\ker(T_p-c)$, $\coker(T_p-c)$, and the layer blocks' polynomials modulo $\ell^{e_j}$ (Proposition~\ref{prop:torsionhecke}(1)(e)) --- read by
  Proposition~\ref{prop:torsionhecke}(3) (the layer polynomials and the Fitting types at every certified $K$, the rest at $2\sigma_{\max}\le K$): at $\ell\ge5$ ($q=1$) on the
  $\mathsf O$-piece modulo $\ell^{K}$ --- the shipped operator on $\Ceq$ or $\mathrm T$ (Proposition~\ref{prop:strictliftN}(6)(a)), \textbf{or} $[H_{N,p}]$ on the cokernel presentation
  $M_{N,\Z_{(\ell)}}/\mathrm{Rel}_N$ read by Proposition~\ref{prop:torsionhecke}(3$'$) (Proposition~\ref{prop:heilbronn}(1), (2); v1.23.0-pre.31, D3's contract note~1: the route the code
  reads); at $\ell=3$ on $\Iso^{(3)}$ with $T'^{(3)}_p$ or on the tree model with $T'_{\mathcal T,p}$, for
  \textbf{every} $N$, $\nu_3(N)=0$ included --- never on $\Ceq$ or $\mathrm T$ at $\ell=3$ (v1.23.0-pre.25, referee M2, resolution (a)); at $\ell=2$ on $\Iso^{(2)}_{\mathcal T}$ with
  $T'_{\mathcal T,p}$ (Proposition~\ref{prop:strictliftN}(5)).
\item \textbf{(The free quotient.)} For $p\in\mathcal P$, a prime $\ell$ and $c\in\Z$ with $\chi^{(N)}_{p,1}(c)\ne0$: the invariant factors of $T_p-c$ on $\Hfr{1}(N)\otimes\Z_{(\ell)}$,
  by Proposition~\ref{prop:torsionhecke}(4) (the lattice $\Hfr{1}(N)\otimes\Z_{(\ell)}$ read on the $\ell$-piece at $\ell\in\{2,3\}$, on the $\mathsf O$-piece otherwise --- at $\ell\ge5$
  on the cokernel presentation by Proposition~\ref{prop:torsionhecke}(3$'$), with $k=K'$; v1.23.0-pre.31) --- so $K>\sigma_{\max}+v_\ell(\chi^{(N)}_{p,1}(c))$ suffices there,
  $\sigma_{\max}$ that of the certified form of $\Delta$ ($\sigma_{\max}=0$ and $K'=K$ when $\mathrm{tors}\,H^{1}(A;\Vcoef)_{(\ell)}=0$, e.g.\ $\ell\notin\mathcal B(M_N)$); a read that falls
  short of (4)'s pivot condition is refused, never misread (v1.23.0-pre.32, referee m2).
\item \textbf{(D2's maps.)} For $B=\Ga_0(MN)$ and $p\in\mathcal P$ with $p\nmid M$: $\res$ and $\cor$ commute with $T_p$ on $H^{q}$ over $\Z$, the sequence \eqref{eq:descentmapsles}
  is $T_p$-equivariant over $\Z$ in every degree (Proposition~\ref{prop:strictliftN}(7)), and
  \[
    \chi^{(MN)}_{p,1}=\chi^{(N)}_{p,1}\cdot\det\bigl(x-T_p\mid H^{2}(A,B)/\mathrm{tors}\bigr),\qquad\text{in particular }\chi^{(N)}_{p,1}\mid\chi^{(MN)}_{p,1}\text{ in }\Z[x].
  \]
  For $p\mid M$ nothing is asserted (the squares fail in general, Proposition~\ref{prop:descentmaps}(3)(b)); for $p\mid N$, nothing is asserted here of $U_p$ ([R-Q4]; on cohomology,
  Proposition~\ref{prop:descentmaps}(3)(b$'$); its lift on the free model is Proposition~\ref{prop:strictliftN}(9), v1.26.0-pre.1, whether a binary serves it over $\Z$ is item~8~(d$'$)(v) of \S\ref{sec:recordschema}, v1.26.0-pre.11, confirm-referee C1).
\item \textbf{(Every degree.)} For $q\ge2$ the data of (1)--(4) in degree $q+2$ are those of degree $q$, through $\res^{\Ga}_{A}e\cup-$, an isomorphism of Hecke modules
  (Proposition~\ref{prop:descentN}(5)); $q_{\max}=3$ determines every degree.
\item \textbf{(Not asserted.)} (a) the $\Z_{(\ell)}[T_p:p\in\mathcal P]$-module structure of $\Hfr{1}(N)\otimes\Z_{(\ell)}$ beyond~(4) --- in particular any invariant of $T_p-c$ at a root
  $c$ of $\chi^{(N)}_{p,1}$ --- and the class of the extension $0\to\mathrm{tors}\to H^{1}\to\Hfr{1}(N)\to0$ as Hecke modules (the plan's D-$\delta$7 and Q-R1): no certified read
  is specified here. When $\chi^{(N)}_{p,1}$ is squarefree over $\Q$ the $\Z_{(\ell)}$-similarity class should be determined by $T_p$ modulo a finite power of $\ell$ (finite
  determinacy for a separable characteristic polynomial), but no bound is given and none is asserted; for repeated roots --- which occur at level $N$ whenever there are old
  forms, $(\chi^{(1)}_{p,1})^{2}\mid\chi^{(N)}_{p,1}$ --- no finite read is known to suffice (v1.23.0-pre.25, referee m5, in place of ``open mathematics'').
  \textbf{The read at a root} (v1.25.0-pre.0; wording v1.25.0-pre.0b, A7-m3): the invariant at a root that \S\ref{sec:heckeZ} uses --- the lattice congruence module $\eta_L$ at the root
  $c=\lambda(T_p)$ of a rational Eisenstein component, read as one Smith form of $c_{\mathcal S',p}(T_p)$ modulo $\ell^{K}$, $K>v_\ell c_{\mathcal S',p}(\lambda(T_p))$ --- is specified
  and proved by Proposition~\ref{prop:heckeZcong}(5) (with the minimal polynomial on the complement, never $T_p-c$ alone); the extension class is determined by finite data,
  $\varepsilon_K$ for any $K\ge e_\ell$ (Proposition~\ref{prop:heckeZglue}(2)(a)), and from computed operators a non-split verdict is certified, a split one only under
  Remark~\ref{rmk:heckeZcbthree} (Proposition~\ref{prop:heckeZglue}(2)(b)); both are statements of \S\ref{sec:heckeZ}, recorded by item~8 of \S\ref{sec:recordschema}, \textbf{not} data of this record, which still
  asserts neither;
  (b) integer Hecke matrices in a basis (tier B; the plan's D8); (c) $U_p$, nebentypus $\chi\ne1$, groups other than $\Ga_0(N)$, the boundary pair and its cone (D4, D5c). (v1.23.0-pre.26: the
  pair and its cone are \S\ref{sec:descentpair}, a separate statement; this record is unchanged, and the record's pair columns are D6's. v1.23.0-pre.37: the record is \code{integral\_descent\_record/v1}, item~7 of \S\ref{sec:recordschema}, which fixes its
  degrees, its hashed and unhashed fields and its refusals; (4) and (5) are not in it; v1.23.0-pre.38: built, audited and merged at \texttt{2aa61397}, its golden
  points landed at \texttt{6958a5d2}.)
\end{enumerate}
\end{theorem}

\begin{proof}
(1) is the cited statements. (2) is Proposition~\ref{prop:crtcharpolyN}, whose residues are read on the $\mathsf O$-piece (its (2)). (3) Each local operator named in (3) ---
Proposition~\ref{prop:strictliftN}(5) at $\ell\in\{2,3\}$, (6)(a) at $\ell\ge5$; (6)(b) is not used --- is a cochain endomorphism, over $\mathsf O_\ell\subset\Z_{(\ell)}$, of a complex of finitely generated free $\Z_{(\ell)}$-modules
(Propositions~\ref{prop:certifiedeell}(1), \ref{prop:isomodeltree}(1)(d)), inducing $T_p\otimes\Z_{(\ell)}$ under the comparison with $H^{q}(A;\Vcoef)\otimes\Z_{(\ell)}$
(flat base change, Proposition~\ref{prop:descentN}(1)); Proposition~\ref{prop:torsionhecke}(3) reads $(G,T_p)$ exactly, one $p$ at a time, in one system of generators, so
the joint module is read as well, and the parts that need only a certified form are those of Proposition~\ref{prop:torsionhecke}(3). Two computations differ by an isomorphism of $\Z_{(\ell)}$-modules that intertwines the operators (both induce the same $T_p$), so their
invariants agree (Proposition~\ref{prop:torsionhecke}(1)). \emph{The cokernel presentation at $\ell\ge5$} (v1.23.0-pre.31). For $n$ odd, $2$ kills $H^{\bullet}(A;\Vcoef)$ (the
proof of Proposition~\ref{prop:crtcharpolyN}(2)), so there is no $\ell$-torsion and nothing is read. For $n$ even put $\cring:=\Z_{(\ell)}$; $6\in\cring^{\times}$. $\mathrm{Rel}_N\subseteq M_{N,\cring}$
is spanned by finitely many vectors with entries in $\Z[1/6]\subseteq\cring$, the columns of a matrix $\Delta$ (the generators of $M^{\langle S\rangle}$, $M^{\langle U\rangle}$ ($U=ST$, written $M^{\langle ST\rangle}$ there; v1.23.0-pre.32, referee N2) of
Proposition~\ref{prop:levelorbit}(1), (2), orbit by orbit, or the columns of $1+M_N(S)$ and $1+M_N(U)+M_N(U)^{2}$); $H_{N,p}$ is an endomorphism of $M_{N,\cring}$ with integral entries
($\rhon$ of integral matrices) and $H_{N,p}(\mathrm{Rel}_N)\subseteq\mathrm{Rel}_N$ (Proposition~\ref{prop:heilbronn}(1)). The certifying rank is $\rank_{\Q}\Delta=m_N(n+1)-b_1(M_N)$:
$\coker(\Delta\otimes\Q)=M_{N,\Q}/\mathrm{Rel}_N\cong H^{1}(A;\Vcoef\otimes\Q)$ by $\varkappa_N$ over $\Q$, of dimension $b_1(M_N)$ (Proposition~\ref{prop:descentN}(7)(b) and flat base change,
Proposition~\ref{prop:descentN}(1)). Proposition~\ref{prop:torsionhecke}(3$'$) therefore reads $(\mathrm{tors}\coker\Delta,[H_{N,p}])$, for all $p\in\mathcal P$ in the one system of generators of one
certified Smith form; and $\varkappa_N$ over $\cring$ (Definition~\ref{def:kappaN}) is an isomorphism $\coker\Delta\to H^{1}(A;\Vcoef\otimes\cring)$ with $\varkappa_N[H_{N,p}]=T_p\varkappa_N$
(Proposition~\ref{prop:heilbronn}(2), $[A\,g\,A]=T_p$ as $p\nmid N$), so by the independence clause of~(3$'$) the invariants read are those of
$(\mathrm{tors}\,H^{1}(A;\Vcoef\otimes\Z_{(\ell)}),T_p)=(\mathrm{tors}\,H^{1}(A;\Vcoef)_{(\ell)},T_p)$ (flat base change). (4) Proposition~\ref{prop:torsionhecke}(4) with $\psi=T_p$, whose characteristic polynomial is
$\chi^{(N)}_{p,1}$ by base change; $\chi^{(N)}_{p,1}(c)\bmod\ell^{K}$ is read in words by Proposition~\ref{prop:crtcharpolyN}(4). At $\ell\ge5$ on the presentation: $\varkappa_N$ induces an
isomorphism of $\coker\Delta/\mathrm{tors}$ with $H^{1}(A;\Vcoef\otimes\Z_{(\ell)})/\mathrm{tors}=\Hfr{1}(N)\otimes\Z_{(\ell)}$ intertwining $[H_{N,p}]$ with $T_p$, and
Proposition~\ref{prop:torsionhecke}(3$'$) gives the matrix of $T_p$ on it modulo $\ell^{K'}$ in the basis of its free part. \emph{The precision} (v1.23.0-pre.32, referee m2):
Proposition~\ref{prop:torsionhecke}(4) with $k=K'=K-\sigma_{\max}$ certifies once $K'$ exceeds the largest exponent, and the exponents sum to $v_\ell(\chi^{(N)}_{p,1}(c))$, so
$K>\sigma_{\max}+v_\ell(\chi^{(N)}_{p,1}(c))$ suffices; $\sigma_{\max}$ is over \emph{all} pivots of $\Delta$'s certified form, and the word read of $\chi^{(N)}_{p,1}(c)\bmod\ell^{K}$ above
only locates $v_\ell(\chi^{(N)}_{p,1}(c))$, it is not the precision of the read. If $\mathrm{tors}\,H^{1}(A;\Vcoef)_{(\ell)}=0$ then $\mathrm{tors}\coker\Delta=0$ (through $\varkappa_N$), every
pivot is a unit, $\sigma_{\max}=0$ and $K'=K$. The bound is sufficient, not necessary; a read with fewer than $r-s$ pivots of valuation $<K'$ is refused by~(4), never wrong. (5) The squares: Proposition~\ref{prop:descentmaps}(3)(a)
on cohomology, and Proposition~\ref{prop:strictliftN}(7) for the sequence. Rationally $\res$ is injective ($\cor\circ\res=[A:B]$) and the sequence splits into
$0\to H^{1}(A)_{\Q}\to H^{1}(B)_{\Q}\to H^{2}(A,B)_{\Q}\to0$ (Proposition~\ref{prop:descentmaps}(2)(g)(iii)), $T_p$-equivariantly, so the characteristic polynomials multiply;
the quotient of two monic integer polynomials that is monic in $\Q[x]$ lies in $\Z[x]$ (Gauss). (6) Proposition~\ref{prop:descentN}(5). (7) states scope.
\end{proof}

\begin{remark}[What the plan stated, and what holds]\label{rmk:descentcertscope}
\PROVED{} (v1.23.0-pre.24; (i)--(iv) amended at v1.23.0-pre.25 after the referee's M1, M2, m1, m5.) The plan's \S2.3 and its clause rows C-D4, C-D5 are right in substance; four points are corrected here. (i) \textbf{Torsion is not open.} The plan
lists ``a certified $\Z/\ell^{K}[T]$-module type \dots\ on torsion'' among the things not asserted. On a \emph{certified finite} group the operators are read exactly once
$K\ge\sigma_i+\sigma_{\max}$ (Proposition~\ref{prop:torsionhecke}(3)), so the whole module structure under any finite set of $T_p$ is certified
(Theorem~\ref{thm:descentcert}(3)); what stays unasserted of D-$\delta$7 is the \emph{lattice} $\Hfr{1}(N)\otimes\Z_{(\ell)}$ under $T$ at the roots of $\chi$, and the extension
class --- not asserted because no certified read is specified, and for repeated roots none is known to suffice (Theorem~\ref{thm:descentcert}(7)(a)). (ii) \textbf{The bound.} ``$K$ above the group's exponent plus a certified bound on the congruence exponent'' becomes: for torsion,
$K\ge2\sigma_{\max}$ for the module (necessary, Proposition~\ref{prop:torsionhecke}(3)), while the layer polynomials $P_j$ and the Fitting types need only a certified form
($K>\sigma_{\max}$), and no congruence bound; for the lattice, $k>v_\ell(\chi(c))$ at a non-root $c$, with $\chi$ known exactly --- or simply the self-certifying count of
Proposition~\ref{prop:torsionhecke}(4). (iii) \textbf{``The characteristic polynomial over $\Z/\ell^{k}$ on each Smith summand''} is not well defined
(Proposition~\ref{prop:torsionhecke}(2)); its well-defined replacement is the layer block's characteristic polynomial modulo $\ell^{e_j}$, $e_j=\min(j,d_j)$
(Proposition~\ref{prop:torsionhecke}(1)(e); in general no finer), whose reduction modulo $\ell$ is $P_j$. (iv) \textbf{The lever of D-$\delta$2 is a theorem}: $\chi^{(N)}_{p,1}$ needs no
$2$- or $3$-adic input at all (Proposition~\ref{prop:crtcharpolyN}(2)); the local pieces at $2$ and $3$ enter only (3) and (4). The plan's C-D4 ``strict lift for $\mathrm{diag}(1,p)$,
every $p$'' is existence already (Proposition~\ref{prop:strictlift}(1)); what is new is the bulk lift without a boundary step, its construction for every $p$, and its level-$N$
form (Proposition~\ref{prop:strictliftN}(2), (3)). C-D4$'$ (a strict lift directly on $\Iso^{(3)}$) is not drafted; it is a cost lever. The cheap route of
Proposition~\ref{prop:strictliftN}(6)(b) at $\ell=3$, $\nu_3(N)=0$, holds for an integral datum certified over $\Z$ only, and the shipped canonical $\tau$ is integral in practice,
not by construction; so the record does not use it, and its $3$-local piece is the transport $T'^{(3)}_p$ for every $N$ (v1.23.0-pre.25, referee M2, resolution (a): one route
at $\ell=3$, and the shipped Bredon operator never read at $\ell\in\{2,3\}$).
\end{remark}

\noindent\emph{A check outside the build, not a tag} (v1.23.0-pre.24): \file{docs/implementation/slot24/cd4\_hecke\_check.py} (\code{python3 -I}, pure integers, no computer algebra
system), output \file{cd4\_hecke\_check.out}. It reads no line of \code{cohomology\_app} and imports nothing from the tree: it rebuilds the cosets of $\Ga_0(N)$, the coinduced module
$M_N$, the group $1$-cocycles of $\Ga=\langle S,U\mid S^{4},S^{2}U^{-3}\rangle$ with Shapiro, $T_p$ as $\cor\circ(\iota_g,\rhon(g))^{\ast}\circ\res$ on cocycles of $\Ga_0(N)$,
$\res$ and $\cor$ between levels, exact Smith forms over $\Z$ and over $\Z/\ell^{K}$. On $21$ cells (level $1$, $n\in\{1,2,3,4,6,8,10,12,16\}$; $\Ga_0(N)$, $N\in\{2,3,4,5,7,11\}$, $n=0$;
$(N,n)\in\{(2,2),(2,10),(3,4),(5,4),(7,2),(11,2)\}$) and $p\le7$, $p\nmid N$, it checks, in degree $1$: $\chi^{(N)}_{p,1}$ over $\Z$ against \emph{independent} oracles (level $1$:
$q$-expansions of $E_4^{a}E_6^{b}$, e.g.\ $(x-2049)(x+24)^{2}$ at $n=10$, $p=2$; $\Ga_0(11)$, $n=0$: point counts on $X_0(11)$; genus-$0$ levels at $n=0$:
$(x-1-p)^{c-1}$); (2) on the Bredon quotient $M^{-I}/(M^{S}+M^{U})$ over $\F_{\ell'}$, with equal dimensions; (3)'s bound coefficientwise; (4)'s recovery in $64$-bit words by
Garner's digits, the sign test and $\chi(7)\bmod2^{64}$, $\bmod3^{40}$ from the digits, including coefficients reported as residues and bound; Proposition~\ref{prop:torsionhecke}(1)
under random automorphisms and (3)'s read on the local Smith form at every $K\ge2\sigma_{\max}$ with admissibly perturbed lifts (every torsion prime: $2$, $3$, $5$, $7$, $11$);
(4) at $\ell\in\{691,3\}$, $n=10$, and $\ell=5$, $N=11$; Theorem~\ref{thm:descentcert}(5) on eight level pairs over $\Z$ (torsion coordinates included), the square with $\cor$ and
$U_p$, $p\mid N$ (Proposition~\ref{prop:descentmaps}(3)(b$'$)), the failing square at $p\mid M$ as a control, and $(\chi^{(1)}_{p,1})^{2}\mid\chi^{(N)}_{p,1}$ at prime $N$; at $n=10$ the
relations $(T_2-2049)(T_2+24)=0$, $3(T_3-252)=256(T_2+24)$ on $\Hfr{1}$ and $T_2T_3=T_3T_2$ on $H^{1}$ over $\Z$. Planted faults, each detected: a dropped coset, a single
CRT prime, a bound without the Eisenstein term, a perturbed residue (held-out prime), the torsion read below $2\sigma_{\max}$, and a per-summand characteristic polynomial. $931$ checks, $0$ failing (\code{CD4\_CHECK \dots\ failing=0 verdict=OK}). It
covers $H^{1}$ only (the cocycle model has no higher degrees) and does not build $\Iso_{\bullet}$, $\Iso^{(\ell)}$ or a tree model: it checks the statements, not their models.
\emph{The referee's check} (v1.23.0-pre.25), \file{docs/implementation/slot24/cd4referee/cd4referee\_check.py}, shares no code with it (Reidemeister--Schreier on $\Ga_0(N)$,
$H^{2}$ through the Bockstein, $[AgA]$ for $\smat100p$ and $\smat p001$ separately, adversarial certified Smith transforms, $18$ cells and the heavy cell $(11,10)$, $14$ level
pairs); at pre.25 the author added an arm \code{M}: the sharpness pairs of Proposition~\ref{prop:torsionhecke}(2) for every $1\le d<j\le6$ ($\ell=2$), $\le4$ ($\ell=3$), on both
sides, the referee's $\Z/2\oplus\Z/8$, the minimal example of (3) for $\sigma\le4$, $K\in(\sigma,2\sigma+1]$, and on every real torsion read the block polynomials modulo
$\ell^{\min(e_j,K')}$ (\code{CD4R\_CHECK \dots\ failing=0 verdict=OK}, output \file{cd4referee\_check.out}).

\noindent\emph{Code map} (v1.23.0-pre.24, \textbf{stated before D3's code}, contract first, rule~21). \textbf{Not built.} The plan's D3 builds: \file{Integral/Model/StrictLift}
parametrised in $p$, bulk only, through degree $q_{\max}=3$ (Proposition~\ref{prop:strictliftN}(2)); the level-$N$ operator $\widetilde T^{(N)}_{\Iso,p}$ on the code's $\Iso$ with
\file{Hecke/Gamma0LevelN}'s table, applied to columns, and its transports to $\Iso^{(\ell)}$ and to \file{Integral/Model/IsotropyModelTree} (D5) by the code's comparisons; the
characteristic polynomial with the output rule (Proposition~\ref{prop:crtcharpolyN}(4), (5)), on the existing \file{Integral/IntegralCharpolyCrt}; the torsion read
(Proposition~\ref{prop:torsionhecke}(3)) on \file{Integral/Descent/DescentMaps}'s presentations; and \code{IntegralDescentHeckeValidation} with the plan's arms (R), (C), (E), (O),
(M), (L), (K). The existing \code{HeckeIntegralCharpolyValidation} is level~$1$, $N=1$ of Proposition~\ref{prop:crtcharpolyN}. \textbf{Obligations of the code} (v1.23.0-pre.25,
from the referee's M2, m2, m3, m4): (a) the $3$-local piece --- torsion read and lattice --- is computed with $T'^{(3)}_p=\Psi_3^{\ast}\widetilde T^{(N)}_{\Iso,p}\Phi_3^{\ast}$ on
$\Iso^{(3)}$ (or $T'_{\mathcal T,p}$) for \textbf{every} $N$; the code does not branch on $\nu_3(N)$, and the shipped operator on $\Ceq$ or $\mathrm T$ is read only with
$6$ inverted (the $\mathsf O$-piece, and its reductions modulo $\ell^{K}$, $\ell\ge5$), never at $\ell\in\{2,3\}$; (b) the datum $\tau_X$ fed to the lift, strict or canonical,
is checked integral and against (T1)--(T3) over $\Z$ first, and the lift refused otherwise; (c) $\tau_{\Iso}$ is built once per $p$, on $\{s_je_c\}$, independently of $r_t$
and of $N$; (d) the residue primes are $<2^{32}$, and every Horner step modulo $Q>2^{32}$ uses a word mulmod; a negative $c_i$ is reduced as $c^{+}\bmod Q-P\bmod Q$.

\noindent\emph{Code map, status at v1.23.0-pre.31.} \textbf{Built on branch \texttt{s24-d3impl}, pending merge and audit} (not merged into \texttt{hecke-extensions}; tip
\texttt{a1375159}; source \file{docs/implementation/slot24/DESCENT3\_report.md}, read at its \S\S2, 3, 6). The files, all new, in \file{Integral/Descent/}: \file{DescentHecke.h} (the API
and the refusal taxonomy); \file{DescentHeckeCrt.cpp} (Proposition~\ref{prop:crtcharpolyN}: residues of $[H_{N,p}]$ on the shipped sparse cokernel presentation over $\F_{\ell'}$ ---
the presentation form of its~(2) ---, residue primes descending from $2^{30}$, the bound of~(3), Garner's digits, the sign test, $c_i\bmod Q$ with a shift-and-add word mulmod and
the negative case, the output rule~(5), a held-out prime); \file{DescentHeckeOps.cpp} ($\widetilde T^{(N)}_{\Iso,p}$ of Proposition~\ref{prop:strictliftN}(3) with the bulk component
of D4's pair lift as $\tau_{\Iso}$ --- admissible by~(2) ---, certified by \code{StrictLift::certify()}, built once per $p$ on $\{s_je_c\}$, applied to cochains block by block; and
the \textbf{transport at $\ell=2$ and at $\ell=3$}: $T'_{\mathcal T,p}=\Psi_{\mathcal T}^{\ast}\,\widetilde T^{(N)}_{\Iso,p}\,\Phi_{\mathcal T}^{\ast}$ on the $\ell$-local tree model with D5's
certified comparisons, for every $N$, with no branch on $\nu_3(N)$, the shipped Bredon $\tau$ never read at $2$ or $3$); \file{DescentHeckeTorsion.cpp} (the read of
Proposition~\ref{prop:torsionhecke}(3): \code{TorsionOnTree} at $\ell\in\{2,3\}$, $q\le3$, production; \code{TorsionOnCokernel} at $\ell\ge5$, $q=1$, production, the presentation
form~(3$'$) with $\Delta$ the generators of $\mathrm{Rel}_N$ and the certifying rank $m_N(n+1)-b_1(M_N)$; \code{TorsionOnFreeIso}, \code{TorsionOnIsoEll}, oracles;
\code{FreeQuotientFactors}, Theorem~\ref{thm:descentcert}(4) at $\ell\ge5$ on the presentation with $k=K'$; \code{RunDescentHecke}, the record of Theorem~\ref{thm:descentcert}(2),
(3)); the validations \code{IntegralDescentHeckeValidation} (\code{standard}) and \code{IntegralDescentHeckeSlowValidation} (\code{slow}), registered in
\file{Registry/Slot24/DESCENT3.cpp}, sources \file{cmake/slot24/DESCENT3.cmake}; an oracle outside the build, \file{docs/implementation/slot24/d3\_hecke\_oracle.py}. The obligations:
(a) met (the $3$-local piece is the tree transport for every $N$); (b) met \textbf{through the lift} (Proposition~\ref{prop:strictliftN}(2), the check through the lift): the table
read is the strict datum $\tau^{s}$'s integral table, and \code{StrictLift::certify()} checks (T1), (T2), (T3) and the strict identity on every basis element, those at $t\ge1$
included, which makes that table a $g$-twisted chain datum over $\Z$; $\tau^{s}$ is not re-checked separately on $\Xbulk$, and need not be. (v1.23.0-pre.32, referee N5:
\code{certify()} checks the strict identity for $q\le2$ only; for $q\ge3$ every summand of $\Iso_q$ has $t\ge1$ ($s\le2$) and $C_q(\Xbulk)=0$, so the identity holds trivially there,
and the code checks exactly what the clause requires. $H_{N,p}(\mathrm{Rel}_N)\subseteq\mathrm{Rel}_N$ is a theorem, Proposition~\ref{prop:heilbronn}(1), re-checked at two residue primes only;
\code{TorsionOnCokernel}'s \code{NotEndomorphism} tests are partial safeguards, not a certificate.) (c) met (the pair lift, once per $p$,
independent of $r_t$ and $N$); (d) met. \textbf{Not built}: Theorem~\ref{thm:descentcert}(4) at $\ell\in\{2,3\}$ (the lattice $\Hfr{1}(N)\otimes\Z_{(\ell)}$ on the $\ell$-piece; it
needs the free part of $H^{1}=\ker d^{1}/\operatorname{im}d^{0}$ in the tree's coordinates, which the presentation form~(3$'$) does not supply, $d^{1}\ne0$ there), and
Theorem~\ref{thm:descentcert}(5) (no level pair is read with $T_p$ by D3). Theorem~\ref{thm:descentcert}(6) is exercised by an arm of the standard validation ($H^{4}$ against
$H^{2}$ with $T_p$ on the tree). Whatever the branch builds is recorded here as built only on that branch; the merge and the audit are separate events.
(v1.23.0-pre.33: \textbf{audited and merged} into \texttt{hecke-extensions} at \texttt{dc8261a3} --- the audit \file{docs/implementation/slot24/D3\_code\_audit.md} (\texttt{73a00c9b}), its fixes
M1, m1--m3 at \texttt{51a4b295}, \file{DESCENT3\_report.md} \S9.) The $\ell\ge5$ cokernel torsion read (\code{TorsionOnCokernel}) refuses \code{[RssCeiling]} on its predicted RSS
\textbf{before building} anything (rule~19). \code{HeldOutMismatch} and a residue-prime rank defect (\code{[ResidueRank]}: the presentation's $\dim H^{1}\ne b_1$ at a residue prime $\ell'>n+1$)
throw \code{hecke::HeckeInternalError}: they are \textbf{failures of the code, not refusals}, and propagate. \textbf{Open code item} (v1.23.0-pre.34, the C-D8 referee's N7): the enumerator
\code{HeldOutMismatch} survives in \code{enum class D3Refusal} (\file{DescentHecke.h}) although it is never a refusal; it is to be removed or marked vestigial, so that counts of refusal codes stay
honest. Not a statement of this document. (v1.23.0-pre.35: \textbf{closed} on branch \texttt{s24-d5cimpl} --- the enumerator is removed; \code{[HeldOutMismatch]} is still thrown as a
\code{hecke::HeckeInternalError}; \S\ref{sec:treepair}'s code map; v1.23.0-pre.36: audited and merged at \texttt{4064e164}.) D5c extends this section's code in two places, with no value
changed (v1.23.0-pre.36): the comparison of Proposition~\ref{prop:isomodeltree}(4) on the code's $\Iso$ (now \code{SharedTreeComparison()}) and \code{TreeTransport} are built to degree $5$
(\code{kTreeComparisonTop}; the maps in degrees $\le4$ are those of the degree-$4$ build, the construction being degree by degree), and \code{IsoHeckeN::apply\_boundary} applies
$\widetilde T^{(N)}_{\Iso,p}$ on $\Isob$'s rows, reporting a term of $\tau_{\Iso}(\Isob)$ off $\Isob$ (the closure, \code{[IsobClosure]}). (v1.23.0-pre.37, D6: the record that carries this section's data is item~7 of
\S\ref{sec:recordschema}; its (j) renames the refusal of the $\ker/\coker$ read at the ring width --- \code{TorsionOnTree}'s \code{Invariants}, reported as \code{[NotEndomorphism]} --- to
\code{[KerCokerPrecision]}, an open code item of this section's \code{D3Refusal} at v1.23.0-pre.37, closed at \texttt{19a969de}: \code{D3Refusal::KerCokerPrecision},
merged at \texttt{2aa61397}; v1.23.0-pre.38.)

\paragraph{Formalization note.} Propositions~\ref{prop:torsionhecke} and \ref{prop:crtcharpolyN}(1), (4), (5) are finite algebra (Fitting's lemma, the Chinese remainder theorem,
mixed radix) and cheap to formalise; Proposition~\ref{prop:strictliftN} is the comparison theorem and the double-coset formalism; \ref{prop:crtcharpolyN}(3) rests on
Eichler--Shimura and Deligne. No \code{\textbackslash LEAN} marker.

\subsection{The boundary pair and the cone over \texorpdfstring{$\Z$}{Z}: the Borel--Serre pair on the descent}\label{sec:descentpair}

(v1.23.0-pre.26; milestone D4 of \file{docs/implementation/slot24/DESCENT\_plan.md} v0.4, clause C-D7 of its \S2.7, its \S2.2 row ``Boundary LES'' and its
D-$\delta$6; \file{DESCENT\_plan\_audit.md} on D-$\delta$6.) \S\ref{sec:descentN}--\S\ref{sec:descenthecke} glue the \emph{bulk} $H^{\bullet}(\Ga_0(N);\Vcoef)$, its maps and
its Hecke operators; \S\ref{sec:cone} builds the pair and its cone on the cellular model, where they compute the Borel--Serre long exact sequence only with $6$
inverted (Lemma~\ref{lem:conerelative}, Remark~\ref{rmk:conesteinberg}). This subsection fixes, over $\Z$: (i) the \emph{pair models} --- $(\Iso_{\bullet},\Isob_{\bullet})$ over
$\Z$, $(\Iso^{(\ell)},{\Isob}^{(\ell)})$ over $\mathsf O_\ell$ (a free boundary at $\ell=2$, the Bredon boundary at $\ell\ge3$), the cellular pair over $\Z[1/6]$ --- and the
comparisons of pairs, which exist and are unique up to homotopy of pairs (Definition~\ref{def:descentpair}, Lemma~\ref{lem:pairlifting}); (ii) what they compute, the
long exact sequence $\Hcpt{q}\to H^{q}(\Ga)\to H^{q}(\Gbdry)\to\Hcpt{q+1}$ over $\Z$ glued from the three pieces, the boundary in closed form at every level, compact support
($\Hcpt{1}$ torsion-free, the rational Betti numbers, exponent bounds, periodicity), and the certified exponent on the boundary and the cone, the infinite $2$-adic boundary
included (Proposition~\ref{prop:descentpair}); (iii) the cone with the strict pair datum ($\hder=0$) and $T_p$, $p\nmid N$, on it, with exactly what is certified at which
$p$, its relation to \S\ref{sec:descenthecke}'s bulk operator, and D2's maps on the pair (Proposition~\ref{prop:descentpairhecke}); (iv) where the plan's wording is corrected
(Remark~\ref{rmk:descentpairscope}). \textbf{Not here}: the tree cusp pair (the plan's \S5, D5c, clause C-D8; v1.23.0-pre.33: \S\ref{sec:treepair}), $U_p$ ([R-Q4]), Hecke operators with nebentypus.

\emph{What is used.} Definition~\ref{def:pair}, \eqref{eq:conediff}, Lemma~\ref{lem:conerelative}, Definition~\ref{def:bdrygroup}, Propositions~\ref{prop:bdrycohom},
\ref{prop:circlemodel}(4), Lemmas~\ref{lem:ell2unbounded}, \ref{lem:bdrycomponent}, \ref{lem:taubdryexists}(1), \ref{lem:manin}, Propositions~\ref{prop:strictpair},
\ref{prop:coneindep}, Theorem~\ref{thm:conewelldef}, \S\ref{sec:integral}'s \ref{def:isomodel}--\ref{rmk:strictliftbuilt}, \ref{lem:bulkperiodic}, and
Propositions~\ref{prop:isomodelell}, \ref{prop:certifiedeell}, \ref{prop:descentN}, \ref{prop:descentmaps}, \ref{prop:strictliftN}. \emph{Notation.} $\mathbb{P}^1:=\mathbb{P}^1(\Q)$
with $\Ga$, and $\De$, acting by M\"obius transformations; $\Z[\mathbb{P}^1]=\Z[\Ga/\Gbdry]$ (\ref{B2}); $\mathrm{St}:=\ker(\aug\colon\Z[\mathbb{P}^1]\to\Z)$, a $\Z$-free
$\Z[\Ga]$-module ($=\mathrm{Div}^0$ of Lemma~\ref{lem:manin}), $\mathrm{St}_S:=\mathrm{St}\otimes S$; $c(N)$ is the number of cusps of $\Ga_0(N)$, as in \S\ref{sec:levelN}. $M$ is a $\Z[\Ga]$-module, finitely generated and free over $\Z$, $d:=\rank_\Z M$,
$M_S:=M\otimes S$, $M^{\vee}:=\Hom_\Z(M,\Z)$ with the contragredient action; $b_j(M):=\dim_{\Q}H^{j}(\Ga;M_{\Q})$ and $b_j(\partial;M):=\dim_\Q H^{j}(\Gbdry;M_\Q)$. When $-I$ acts on
$M$ by a sign we write it $\varepsilon_M$ ($\varepsilon_{M_N(\chi)}=\chi(-1)(-1)^{n}$). For a $\Z[\Gbdry]$-module $X$, $X^{T}$, $X_{T}$ are the invariants and coinvariants of
$T$ (notation C12). $\mathsf P$, $\mathsf P^{\partial}$, $\rres_{\mathsf P}$, $\Hcpt{q}(\mathsf P;M)$, $c$ (a cusp), $w_c$, $\Gbdry^{(w)}$, $F_{n,w}$, $\varepsilon_M$,
$\mathrm{St}$, $\Upsilon_k$, $j_{\mathrm{St}}$, $R_{\bullet}$, $\varepsilon_R$, $\varepsilon^{\partial}$ (v1.23.0-pre.27), $\tau^{\mathrm{rel}}$, $\widetilde T^{(N)}_{\Isob,p}$, $\Phi^{(N)}_p$ are new symbols of this subsection and mean nothing else (rule~\ref{A5}); $\mathsf P$ is
\emph{not} the $P$ of Lemma~\ref{lem:strictlifting}.

\begin{definition}[Pair models, their cochains and cones; comparisons and homotopies of pairs]\label{def:descentpair}
Assume \ref{H4} and \ref{B1}--\ref{B3}, and let $S$ be a commutative ring. An \textbf{$S$-pair model} is an $S$-model $\mathsf P_{\bullet}$ in the sense of
Proposition~\ref{prop:descentmaps} (finitely generated projective $S[\Ga]$-modules, $\mathsf P_q=0$ for $q<0$, an augmentation making it a resolution of $S$) together with a
subcomplex $\mathsf P^{\partial}_{\bullet}\subseteq\mathsf P_{\bullet}$ whose term in each degree is a direct summand with projective complement, and an augmentation
$\mathsf P^{\partial}_0\to S[\mathbb{P}^1]$ making $\mathsf P^{\partial}_{\bullet}$ a resolution of $S[\mathbb{P}^1]$, such that $\mathsf P^{\partial}_0\to\mathsf P_0\to S$ equals
$\mathsf P^{\partial}_0\to S[\mathbb{P}^1]\xrightarrow{\aug}S$. Put $\mathsf P(M):=\Hom_{S[\Ga]}(\mathsf P_{\bullet},M_S)$, $\mathsf P^{\partial}(M)$ likewise,
$\rres_{\mathsf P}\colon\mathsf P(M)\to\mathsf P^{\partial}(M)$ the restriction, $\Cone{\rres_{\mathsf P}}$ by \eqref{eq:conediff} (with $\mathsf P(M)$, $\mathsf P^{\partial}(M)$ for
$\Ceq$, $C^{\bullet}_{\partial}$), $\coneinc$, $\coneproj$ as there, and
\[
  \Hcpt{q}(\mathsf P;M):=H^{q-1}\bigl(\Cone{\rres_{\mathsf P}}\bigr),\qquad q\in\Z .
\]
A \textbf{comparison of pairs} $\varphi\colon(\mathsf P,\mathsf P^{\partial})\to(\mathsf P',\mathsf P'^{\partial})$ is an $S[\Ga]$-chain map $\mathsf P\to\mathsf P'$ over $\id_S$ with
$\varphi(\mathsf P^{\partial})\subseteq\mathsf P'^{\partial}$, its restriction over $\id_{S[\mathbb{P}^1]}$. Two comparisons are \textbf{homotopic as pairs} if their difference is
$\partial k+k\partial$ with $k(\mathsf P^{\partial})\subseteq\mathsf P'^{\partial}$. The pair models in use:
\begin{enumerate}
\item[(a)] $(\Iso_{\bullet},\Isob_{\bullet})\otimes S$, every $S$ (Definition~\ref{def:isomodel}, \ref{IM2}: $\Isob$ is the sub-sum over $\Sigma^{\partial}_0\sqcup\Sigma^{\partial}_1$;
  its augmentation $\Isob_0\xrightarrow{\cmpI}C_0(\Xbdry)\to\Z[\pi_0\Xbdry]=\Z[\mathbb{P}^1]$, \ref{B2}; Proposition~\ref{prop:isomodel}(2));
\item[(b)] $(\Iso^{(\ell)},{\Isob}^{(\ell)})\otimes_{\mathsf O_\ell}S$ for $\mathsf O_\ell$-algebras $S$, ${\Isob}^{(\ell)}$ the sub-sum of the summands $G'_{s,t}$ over $\Sigma^{\partial}_0\sqcup\Sigma^{\partial}_1$
  (Definition~\ref{def:isomodelell}, \ref{IM2}), the model fixed once over $\mathsf O_\ell$, boundary summands first (Propositions~\ref{prop:isomodelell}(1),
  \ref{prop:certifiedeell}(1));
\item[(c)] $(C_{\bullet}(\Xbulk),C_{\bullet}(\Xbdry))\otimes S$ for $6\in S^{\times}$ (Remark~\ref{rmk:orbifold}; $C_s(\Xbdry)\otimes S=S[\Ga/\{\pm I\}]^{3}$ is projective once $2\in S^{\times}$,
  \ref{B3}, and resolves $S[\mathbb{P}^1]$ because the components of $\Xbdry$ are lines, \ref{B1}, \ref{B2}); here $\Cone{\rres_{\mathsf P}}$ is Definition~\ref{def:pair}'s cone with fibre $M$.
\end{enumerate}
\end{definition}

\begin{lemma}[Comparisons of pairs exist, boundary first, and are unique up to homotopy of pairs]\label{lem:pairlifting}
\PROVED{} \STD{\cite[Thm.~2.2.6]{weibel}} Let $(\mathsf P,\mathsf P^{\partial})$, $(\mathsf P',\mathsf P'^{\partial})$ be $S$-pair models.
\begin{enumerate}
\item Every $S[\Ga]$-chain map $\varphi^{\partial}\colon\mathsf P^{\partial}\to\mathsf P'^{\partial}$ over $\id_{S[\mathbb{P}^1]}$ (one exists) extends to a comparison of pairs $\varphi$ with
  $\varphi|_{\mathsf P^{\partial}}=\varphi^{\partial}$.
\item Any two comparisons of pairs $\varphi,\varphi'$ are homotopic as pairs; a homotopy $k^{\partial}$ between $\varphi|_{\mathsf P^{\partial}}$ and $\varphi'|_{\mathsf P^{\partial}}$ with
  values in $\mathsf P'^{\partial}$ exists and extends to a homotopy $k$ of pairs.
\item For a comparison $\varphi$: $\varphi^{\ast}$ and $\varphi^{\partial\ast}$ are quasi-isomorphisms; $\rres_{\mathsf P}\varphi^{\ast}=\varphi^{\partial\ast}\rres_{\mathsf P'}$ \textbf{strictly};
  $(\varphi^{\ast}[1],\varphi^{\partial\ast})\colon\Cone{\rres_{\mathsf P'}}\to\Cone{\rres_{\mathsf P}}$ is a cochain map commuting with $\coneinc$ and $\coneproj$, and a quasi-isomorphism. A
  homotopy of pairs $k$ between $\varphi,\varphi'$ gives the cone homotopy $\Upsilon_k(c,c_\partial):=(-k^{\ast}c,\ k^{\partial\ast}c_\partial)$:
  $(\varphi^{\ast}[1],\varphi^{\partial\ast})-(\varphi'^{\ast}[1],\varphi'^{\partial\ast})=\Dcone\Upsilon_k+\Upsilon_k\Dcone$. So the maps induced on the three terms of the long exact sequence, and
  on the sequence, do not depend on the comparison.
\end{enumerate}
\end{lemma}

\begin{proof}
(1) $\mathsf P^{\partial}$ is a bounded below complex of projectives and $\mathsf P'^{\partial}$ a resolution of $S[\mathbb{P}^1]$, so $\varphi^{\partial}$ exists \STD{\cite[Thm.~2.2.6]{weibel}}. Apply
Lemma~\ref{lem:strictlifting} with $P:=\mathsf P$, $P':=\mathsf P^{\partial}$, $E:=\mathsf P'$, $C:=S$ in degree $0$, $p$ the augmentation (a surjective quasi-isomorphism),
$E':=\mathsf P'^{\partial}$, $f:=\aug$ and $f':=\varphi^{\partial}$; $pf'=f|_{\mathsf P^{\partial}}$ is the compatibility of the augmentations. The lemma is stated for a free source with $P'$
spanned by part of a basis; with a projective source and a projective complement its two steps are lifts of maps out of a projective module through surjections, exactly as in
the proof of Proposition~\ref{prop:isomodelell}(4). (2) $u:=\varphi-\varphi'$ is a chain map with $\aug u=0$, and $u|_{\mathsf P^{\partial}}$ is a chain map into $\mathsf P'^{\partial}$
over $0\colon S[\mathbb{P}^1]\to S[\mathbb{P}^1]$, hence $u|_{\mathsf P^{\partial}}=\partial k^{\partial}+k^{\partial}\partial$ with $k^{\partial}$ into $\mathsf P'^{\partial}$ (the uniqueness half of
the comparison theorem). Extend $k$ by induction on the degree over a projective complement $Q_q$ of $\mathsf P^{\partial}_q$: for $x\in Q_q$ the element $u(x)-k(\partial x)$ is a cycle
($\partial u(x)-\partial k\partial x=u\partial x-(u\partial x-k\partial\partial x)=0$ by induction), of augmentation $0$ if $q=0$; $\mathsf P'$ is acyclic in positive degrees and
$\ker\aug=\im\partial$ in degree $0$, so the map $x\mapsto u(x)-k(\partial x)$ out of the projective $Q_q$ lifts through $\partial\colon\mathsf P'_{q+1}\to\im\partial$; call the lift $k$. (3) Both
complexes compute $\Ext^{\bullet}_{S[\Ga]}(S,M_S)$, resp.\ $\Ext^{\bullet}_{S[\Ga]}(S[\mathbb{P}^1],M_S)$, and $\varphi^{\ast}$, $\varphi^{\partial\ast}$ are the comparison maps; both
sides of the strict identity send $c$ to $c\circ\varphi|_{\mathsf P^{\partial}}$. The cone: $\Dcone(\varphi^{\ast}c,\varphi^{\partial\ast}c_\partial)=(-\varphi^{\ast}dc,\
\varphi^{\partial\ast}dc_\partial+\rres\varphi^{\ast}c)=(\varphi^{\ast}[1],\varphi^{\partial\ast})\Dcone(c,c_\partial)$, and the five lemma on the two long exact sequences gives the
quasi-isomorphism. The homotopy: $\varphi^{\ast}-\varphi'^{\ast}=dk^{\ast}+k^{\ast}d$, likewise on the boundary, and $\Dcone\Upsilon_k+\Upsilon_k\Dcone$ has first component
$(dk^{\ast}+k^{\ast}d)c$, second component $(dk^{\partial\ast}+k^{\partial\ast}d)c_\partial+(k^{\partial\ast}\rres-\rres k^{\ast})c$, and the last term vanishes because
$k(\mathsf P^{\partial})\subseteq\mathsf P'^{\partial}$.
\end{proof}

\begin{proposition}[The pair over $\Z$: what it computes, the glued long exact sequence, the boundary in closed form, compact support, the certified exponent]\label{prop:descentpair}
\PROVED{} \STD{\cite[Ch.~III.2, III.5--6, III.8--9, V.1--3, VI.9]{brown}, \cite[Thm.~1.3.1, Thm.~2.2.6]{weibel}; \VERIFYloc} (v1.23.0-pre.26; milestone D4, clause C-D7 of the plan;
refereed ACCEPT-with-fixes, \file{docs/implementation/slot24/C-D7\_referee.md}, its fixes M1, M2, m1--m5 at v1.23.0-pre.27; the pre.26 locator VII.6 is withdrawn, referee m1.)
Assume \ref{H4} and \ref{B1}--\ref{B3}; $M$ as in the notation.
\begin{enumerate}
\item \textbf{(What a pair model computes, over every ring.)} Let $(\mathsf P,\mathsf P^{\partial})$ be an $S$-pair model.
  \begin{enumerate}
  \item[(a)] $H^{q}(\mathsf P(M))=H^{q}(\Ga;M_S)$, $H^{q}(\mathsf P^{\partial}(M))=H^{q}(\Gbdry;M_S)$ (Shapiro's lemma), and $H^{q}(\rres_{\mathsf P})=\res^{\Ga}_{\Gbdry}$, for every $q$.
  \item[(b)] $\rres_{\mathsf P}$ is surjective in every degree (split by extension by zero along the complement); with Lemma~\ref{lem:conerelative}(2)'s argument the sequence
    \begin{equation}\label{eq:descentpairles}
      \cdots\to\Hcpt{q}(\mathsf P;M)\xrightarrow{\ \coneproj{}_{\ast}\ }H^{q}(\Ga;M_S)\xrightarrow{\ \res\ }H^{q}(\Gbdry;M_S)\xrightarrow{\ \coneinc{}_{\ast}\ }\Hcpt{q+1}(\mathsf P;M)\to\cdots
    \end{equation}
    is exact in every degree $q\in\Z$.
  \item[(c)] $\Hcpt{q}(\mathsf P;M)\cong H^{q}\bigl(\Hom_{S[\Ga]}(\mathsf P/\mathsf P^{\partial},M_S)\bigr)\cong\Ext^{q-1}_{S[\Ga]}(\mathrm{St}_S,M_S)\cong H^{q-1}\bigl(\Ga;\Hom_\Z(\mathrm{St},M_S)\bigr)$,
    independently of the model, and under these identifications \eqref{eq:descentpairles} is the long exact $\Ext$-sequence of
    $0\to\mathrm{St}\xrightarrow{\,j_{\mathrm{St}}\,}\Z[\mathbb{P}^1]\xrightarrow{\aug}\Z\to0$ in the first variable, \textbf{with these signs} (v1.23.0-pre.27, referee m4, in the manner of
    Proposition~\ref{prop:descentmaps}(2)(b)): identify through $R_{\bullet}:=(\mathsf P/\mathsf P^{\partial})_{\bullet+1}$ with differential $-\partial$ and augmentation
    $\varepsilon_R[x]:=\varepsilon^{\partial}(\partial x)\in\mathrm{St}$ ($\varepsilon^{\partial}\colon\mathsf P^{\partial}_0\to S[\mathbb{P}^1]$ the augmentation; the connecting map of $0\to\mathsf P^{\partial}\to\mathsf P\to\mathsf P/\mathsf P^{\partial}\to0$
    on $H_1$, no sign), the cochains $\Hom(\mathsf P/\mathsf P^{\partial},M_S)^{q}=\Hom(R,M_S)^{q-1}$ identified identically. Then $\res$ is induced by $\aug$; $\coneinc{}_{\ast}=j_{\mathrm{St}}^{\ast}$; and
    $\coneproj{}_{\ast}=-\delta$, $\delta\colon\Ext^{q-1}(\mathrm{St},M_S)\to\Ext^{q}(S,M_S)$ the connecting map with the convention $\delta[x]=[d\tilde x]$ on a resolution of the short exact
    sequence ($\tilde x$ a lift, the cochain differential carrying no sign). The ``$-$'' depends on the convention: identifying $H_0(R)$ with $\mathrm{St}$ through $-\varepsilon_R$ moves it to
    $\coneinc{}_{\ast}$. Immaterial for groups, kernels, images, exactness and $T_p$-equivariance; it matters only when a map is compared entrywise with an independently computed
    $\Ext$ map. Hence $\Hcpt{q}=0$ for $q\le0$ over every $S$; over $S=\Z$,
    $\Hcpt{1}(\Ga;M)=\Hom_{\Z[\Ga]}(\mathrm{St},M)$ is \textbf{torsion-free}; and for the model (c) this is Lemma~\ref{lem:conerelative}'s $\Hcpt{\bullet}$ with fibre $M$.
  \item[(d)] \textbf{(The models and their integral forms.)} On (a), $\Cone{\rres_{\Iso}}$ is a complex of finitely generated free $\Z$-modules of ranks $5d$, $17d$ and $25d$ in degrees
    $-1$, $0$ and $\ge1$, with $\Cone{\rres_{\Iso}}\otimes S$ the cone over $S$. On (b), the cochains are finitely generated free $\mathsf O_\ell$-modules
    (Proposition~\ref{prop:certifiedeell}(2)) with every base change; \textbf{at $\ell=2$} every boundary stabiliser is its own $2$-Sylow ($P_\sigma=\Ga_\sigma=\{\pm I\}$, \ref{B3}),
    so ${\Isob}^{(2)}_q\cong\mathsf O_2[\Ga]^{w^{\partial}_q}$ ($w^{\partial}_q=3,6,6,\dots$) is free and computes $H^{\bullet}(\Gbdry;M_{\mathsf O_2})$, non-zero in infinitely many degrees by~(4); \textbf{at
    $\ell\ge3$} every boundary column is $\mathsf O_\ell$ in degree $0$ ($P_\sigma=1$, $Q_\sigma=\{\pm I\}$), $\cmpI'$ restricts to an \textbf{isomorphism of complexes}
    ${\Isob}^{(\ell)}\to C_{\bullet}(\Xbdry)\otimes\mathsf O_\ell$, and ${\Isob}^{(\ell)}(M)$ is the circle model with fibre $M_{\mathsf O_\ell}$ (Proposition~\ref{prop:circlemodel}(4)): \textbf{on
    $\Iso^{(3)}$ the boundary is Bredon}. On (c) over $\Z[1/6]$, at level $1$ with $M=\Vcoef$, the cone tensored with $\Z/\ell^{K}$, $\ell\ge5$, is the shipped E8 item-4b pair
    (Definition~\ref{def:pair} over $\cring=\Z/\ell^{K}$; Proposition~\ref{prop:ringinforce}(1)) --- an identity of complexes.
  \end{enumerate}
\item \textbf{(The glued long exact sequence over $\Z$.)} $\cmpI$ is a comparison of $\Z[1/6]$-pair models $(\Iso,\Isob)\otimes\Z[1/6]\to(C_{\bullet}(\Xbulk),C_{\bullet}(\Xbdry))\otimes\Z[1/6]$
  (\ref{IM2}), and $\alpha_c:=(\cmpI^{\ast}[1],\cmpI^{\partial\ast})$ is a quasi-isomorphism of cones over $\Z[1/6]$ commuting strictly with $\rres$, $\coneinc$, $\coneproj$. With
  $A_c:=\Cone{\rres_{C(\Xbulk)}}$ over $\Z[1/6]$, $B_c:=\Cone{\rres_{\Iso}}\otimes\Zsix$, $C_c:=\Cone{\rres_{\Iso}}\otimes\Qsix$:
  \begin{enumerate}
  \item[(a)] the mapping cocone of $A_c\oplus B_c\to C_c$ ($\alpha_c$ followed by base change, minus the canonical map) is quasi-isomorphic to $\Cone{\rres_{\Iso}}$, and
    $\Hcpt{q}(\Ga;M)\cong H^{q-1}(A_c)\times_{H^{q-1}(C_c)}H^{q-1}(B_c)$ for every $q$; \textbf{the only glueing datum is the pair map $\cmpI$};
  \item[(b)] every map of \eqref{eq:descentpairles} commutes strictly with every comparison of pairs (Lemma~\ref{lem:pairlifting}(3)) and with base change, so the sequence over $\Z$ is
    determined by its localisations: the $\ell$-primary parts of its terms, and the kernels, images and cokernels of its maps, at $\ell\ge5$ on the $\Z[1/6]$-piece (model (c), or
    (a)), at $\ell\in\{2,3\}$ on model (b) or (a) over $\Z_{(\ell)}$, the free parts on any piece. It is exact over $\Z$; and since a finitely generated abelian group vanishes iff it vanishes
    after $\otimes\Z[1/6]$, $\otimes\Z_{(2)}$ and $\otimes\Z_{(3)}$, \textbf{exactness at a spot of the glued sequence is equivalent to exactness of the three local sequences there};
  \item[(c)] the maps are read on each piece by Proposition~\ref{prop:descentmaps}(2)(g)'s precision rule, from Smith forms certified by~(6); kernels, images and the comparisons
    ``$\im=\ker$'' so computed are exact provided every exponent read stays below the precision left at its step, and are refused otherwise.
  \end{enumerate}
  No coherence datum enters (a)--(c): every map is induced by a comparison or by the inclusion $\mathsf P^{\partial}\subseteq\mathsf P$. Proposition~\ref{prop:coneambiguity}'s ambiguity
  concerns a cochain-level Hecke homotopy, not these maps; it is removed for the Hecke operator in Proposition~\ref{prop:descentpairhecke}.
\item \textbf{(Level $N$; the cusps.)} For $A\le\Ga$ of finite index and $X$ a $\Z[A]$-module, finitely generated and free over $\Z$, Shapiro's lemma (in the second variable;
  induced and coinduced modules agree in finite index, Definition~\ref{def:inducedmodule}) gives $H^{q}(\Gbdry;\Ind_A^{\Ga}X)\cong\Ext^{q}_{\Z[A]}(\Z[\mathbb{P}^1],X)\cong
  \bigoplus_{c\in A\backslash\mathbb{P}^1}H^{q}(A_c;X)$, $A_c$ the stabiliser of the cusp $c$, and $\Hcpt{q}(\Ga;\Ind_A^{\Ga}X)\cong\Ext^{q-1}_{\Z[A]}(\mathrm{St},X)$: the group-theoretic
  pair of $A$, \textbf{block-diagonal by cusps on the boundary}. For $A=\Ga_0(N)$ and the coordinates of Definition~\ref{def:cosetrepsN}: the cusps are the $\langle T\rangle$-orbits on
  $\Ga_0(N)\backslash\Ga\cong\mathbb{P}^1(\Z/N)$ (bottom rows, right action), $c(N)$ their number; the orbit of the coset of $w$ has length $w_c$ (the \emph{width}),
  $\sum_cw_c=m_N$; its stabiliser in $\Gbdry$ is $\Gbdry^{(w_c)}:=\langle T^{w_c}\rangle\times\langle-I\rangle$; and
  $\operatorname{Res}_{\Gbdry}M_N\cong\bigoplus_c\Ind_{\Gbdry^{(w_c)}}^{\Gbdry}\Vcoef$ with $T^{w_c}$ acting on $\Vcoef$ by $\rhon(T^{w_c})$ and $-I$ by $(-1)^{n}$ (for $M_N(\chi)$ by $\chi(-1)(-1)^{n}$ on the
  $-I$ part; the $T^{w_c}$-action is then twisted by the value of $\chi$ on $wT^{w_c}w^{-1}$). So $H^{q}(\Gbdry;M_N)\cong\bigoplus_cH^{q}(\Gbdry^{(w_c)};\Vcoef)$.
\item \textbf{(The boundary in closed form.)} Let $X$ be a $\Z[\Gbdry^{(w)}]$-module, finitely generated and free over $\Z$, on which $-I$ acts by $\varepsilon=\pm1$, $T^{w}$ by $t$. Then, with
  $X^{t}$, $X_t$ the invariants and coinvariants of $t$,
  \[
    \varepsilon=+1:\ H^{0}=X^{t},\ H^{1}=X_t,\ H^{2j}=(X/2)^{t},\ H^{2j+1}=(X/2)_t\ (j\ge1);\qquad
    \varepsilon=-1:\ H^{0}=0,\ H^{2j-1}=(X/2)^{t},\ H^{2j}=(X/2)_t\ (j\ge1),
  \]
  with $\dim_{\F_2}(X/2)^{t}=\dim_{\F_2}(X/2)_t$. Hence $2\,H^{q}(\Gbdry^{(w)};X)=0$ for $q\ge2$ (for $q\ge1$ if $\varepsilon=-1$), $H^{q}\otimes\Z[1/2]=0$ for $q\ge2$ (\textbf{the boundary is
  bounded away from $2$}), and $H^{q}\otimes\Z_{(2)}\ne0$ for every $q\ge2$ as soon as $X\ne0$ and $t$ is unipotent modulo $2$ (as for $\operatorname{Res}_{\Gbdry}M_N$, (3)): then $(X/2)^{t}\ne0$. For $X=\Vcoef$ and $t=\rhon(T^{w})$, $w\ge1$:
  \[
    \Vcoef^{T^{w}}=\Z\,x^{n},\qquad (\Vcoef)_{T^{w}}\cong\Z\oplus F_{n,w},\quad\abs{F_{n,w}}=n!\,w^{n},\qquad
    \dim_{\F_2}(\Vcoef/2)^{T^{w}}=\begin{cases}n+1,&w\text{ even},\\ \lfloor n/2\rfloor+1,&w\text{ odd}.\end{cases}
  \]
  So for $M_N$, $n$ even: $H^{0}(\Gbdry;M_N)=\Z^{c(N)}$, $H^{1}=\bigoplus_c(\Z\oplus F_{n,w_c})$, and $H^{q}$, $q\ge2$, is elementary abelian of $\F_2$-rank
  $\sum_c\dim(\Vcoef/2)^{T^{w_c}}$; $n$ odd: $H^{0}=0$ and $H^{q}$, $q\ge1$, elementary abelian of that rank. At $N=1$ this is Proposition~\ref{prop:bdrycohom} and
  Lemma~\ref{lem:ell2unbounded}, with $H^{1}$ explicit: for $n$ even, $\mathrm{tors}\,H^{1}(\Gbdry;\Vcoef)$ has order $n!$; for $n$ odd, $H^{1}(\Gbdry;\Vcoef)=(\Vcoef/2)^{T}$, of order
  $2^{\lfloor n/2\rfloor+1}$ (v1.23.0-pre.27, referee m3). \textbf{Periodicity}: $\res_{\Gbdry^{(w)}}e\cup-\colon H^{q}\to H^{q+2}$ is an
  isomorphism for $q\ge2$ ($q\ge1$ if $\varepsilon=-1$).
\item \textbf{(Compact support.)}
  \begin{enumerate}
  \item[(a)] \textbf{(Low degrees.)} $\rres^{0}$ induces the inclusion $M_S^{\Ga}\subseteq M_S^{\Gbdry}$ over every $S$, so $\Hcpt{0}=0$ and $\Dcone^{-1}$ is injective over every field: on
    $\Cone{\rres_{\Iso}}$ over $\Z$ every elementary divisor of $\Dcone^{-1}$ is a unit, $\rank_\Q\Dcone^{-1}=5d$, and $\mathrm{tors}\,\Hcpt{1}=0$ (the integral form of (1)(c)).
  \item[(b)] \textbf{(Rational ranks.)} Let $F$ be a field with $6\in F^{\times}$. Then $\Hcpt{q}(\Ga;M_F)=0$ for $q\ge3$, $\dim_F\Hcpt{2}(\Ga;M_F)=\dim_F(M_F^{\vee})^{\Ga}$, and
    \[
      \dim_F\Hcpt{1}=\dim H^{1}-\dim H^{0}+\dim(M_F^{\vee})^{\Ga},\qquad \rank\bigl(\res^{1}\colon H^{1}(\Ga;M_F)\to H^{1}(\Gbdry;M_F)\bigr)=\dim H^{1}(\Gbdry;M_F)-\dim(M_F^{\vee})^{\Ga},
    \]
    using $\dim H^{0}(\Gbdry;M_F)=\dim H^{1}(\Gbdry;M_F)$. For $M=M_N(\chi)$ over $\Q$: $(M^{\vee}_\Q)^{\Ga}\ne0$ iff $n=0$ and $\chi=1$, so $b_0(M^{\vee})=b_0(M)$, $\dim\Hcpt{1}=b_1(M)$,
    $\dim\Hcpt{2}=b_0(M)$; for $M_N$, $n$ even, $\rank\res^{1}=c(N)-b_0(M_N)$.
  \item[(c)] \textbf{(Exponents and the primes.)} For $q\ge3$, $48\,\Hcpt{q}(\Ga;M)=0$. At $\ell\ge5$, $\Hcpt{q}\otimes\Z_{(\ell)}=0$ for $q\ge3$ and the $\ell$-local sequence \eqref{eq:descentpairles}
    is six-term, $0\to H^{0}\to H^{0}(\Gbdry)\to\Hcpt{1}\to H^{1}\to H^{1}(\Gbdry)\to\Hcpt{2}\to0$. At $\ell=3$, $\coneproj{}_{\ast}\colon\Hcpt{q}\otimes\Z_{(3)}\to H^{q}(\Ga;M)\otimes\Z_{(3)}$ is an
    isomorphism for $q\ge3$ (at level $1$: $H^{q}(\Z/3;M^{-I})$, Proposition~\ref{prop:ellthreeclosed}(2), non-zero e.g.\ for $q$ even at $n=0$). At $\ell=2$ the boundary and, in general, the
    cone are non-zero in infinitely many degrees (D-$\delta$6). (Observed on every cell of the two checks below: the exponent of $\Hcpt{q}$, $q\ge3$, divides $12$; not asserted,
    v1.23.0-pre.27, referee m5.)
  \item[(d)] \textbf{(Periodicity.)} $e\cup-$, $e\in H^{2}(\Ga;\Z)$ of Lemma~\ref{lem:bulkperiodic}, acts on the three terms of \eqref{eq:descentpairles} in their coefficient form $H^{\bullet}(\Ga;X)$,
    $X\in\{M,\Hom_\Z(\Z[\mathbb{P}^1],M),\Hom_\Z(\mathrm{St},M)\}$ (on the boundary term it is $\res_{\Gbdry}e\cup-$, the periodicity of (4)), and commutes with its maps; it is an
    isomorphism on $H^{q}(\Ga;M)$ and on $H^{q}(\Gbdry;M)$ for $q\ge2$, and on $\Hcpt{q}$ for $q\ge3$. Degrees $\le4$ of the pair determine every degree: \textbf{$q_{\max}\ge4$ for
    the pair}, one more than the bulk's $3$.
  \item[(e)] \textbf{($\Hcpt{2}$ at $\ell\ge5$; v1.23.0-pre.27, referee M2.)} For every prime $\ell\ge5$,
    \[
      \Hcpt{2}(\Ga;M)\otimes\Z_{(\ell)}\;\cong\;H_0(\Ga;M)\otimes\Z_{(\ell)}=M_{\Ga}\otimes\Z_{(\ell)} .
    \]
    For $M=M_N$, $M_{\Ga}\cong\Vcoef_{A}$ ($A=\Ga_0(N)$), which is $\Z$ for $n=0$ and is killed by $\otimes\Z_{(\ell)}$ for every $\ell>n+1$ if $n>0$, every $N$: so
    $\mathrm{tors}\,\Hcpt{2}(\Ga;M_N)$ has \textbf{no $\ell$-torsion for $\ell\ge5$, $\ell>n+1$}, and its $\ell$-part for $\ell\ge5$ is that of $\mathrm{tors}\,\Vcoef_A$, a cokernel of a $d\times2d$ integer
    matrix. (5)(b)'s $\dim\Hcpt{2}=\dim(M^{\vee}_F)^{\Ga}$ is the $\otimes\Q$ shadow ($\dim M_{\Ga,\Q}=\dim(M^{\vee}_\Q)^{\Ga}$).
  \end{enumerate}
\item \textbf{(The certified exponent on the pair; D-$\delta$6.)} Let $\mathsf P$ be a $\Z_{(\ell)}$-pair model whose cochains are finitely generated free $\Z_{(\ell)}$-modules ((a) or (b) over
  $\Z_{(\ell)}$). Propositions~\ref{prop:certifiede}(1)--(3), Remark~\ref{rmk:certifiedeperdiff} and Proposition~\ref{prop:certifiedeell}(6) (the Smith-valuation read-out and the stopping rule) hold
  verbatim for the differentials $\Dcone^{\partial,q}$ of $\mathsf P^{\partial}(M)$ and $\Dcone^{q}$ of $\Cone{\rres_{\mathsf P}}$ ($q\ge-1$), read on the brutal truncations $\sigma^{\le q_{\max}+1}$ of
  bulk and boundary; the cone of the truncated restriction equals the cone in degrees $\le q_{\max}$, so $\Hcpt{q}$ is read for $q\le q_{\max}$. \textbf{Admissible sources for
  $\rank_\Q$}, besides the Betti recursion $\rank_\Q D^{q}=\rank C^{q}-b_q-\rank_\Q D^{q-1}$ on the complex in question:
  \begin{enumerate}
  \item[(a)] boundary: the closed form, $b_0(\partial;M)=b_1(\partial;M)=\rank M^{\Gbdry}$ ($=c(N)$ for $M_N$, $n$ even; $0$ for $n$ odd), $b_q(\partial;M)=0$ for $q\ge2$; or the same boundary
    complex modulo a prime $\ell'$ (not dividing $\mathsf q_\ell$ on (b)) at which $H^{q+1}(\Gbdry;M)$ has no $\ell'$-torsion --- every $\ell'\ge3$ for $q\ge1$;
  \item[(b)] cone, $\Dcone^{-1}$: $\rank_\Q=\rank\mathsf P(M)^{0}$, no source needed ((5)(a));
  \item[(c)] cone, $\Dcone^{0}$: the Betti recursion with $b_0(\Cone{\rres})=\dim_\Q\Hcpt{1}$ of (5)(b), its $b_j(M)$ from a source admissible by Propositions~\ref{prop:certifiedeell}(4),
    \ref{prop:descentN}(7); \textbf{or} the same cone modulo a prime $\ell'\ge5$ (not dividing $\mathsf q_\ell$ on (b)) at which $H_0(\Ga;M)$ has no $\ell'$-torsion, which by (5)(e) is
    exactly when $\Hcpt{2}$ has none --- \textbf{for $M_N$, every $\ell'\ge5$ with $\ell'>n+1$}, the rule of Proposition~\ref{prop:descentN}(7) for the bulk; for $M_N(\chi)$, $\chi\ne1$, the
    computed $\mathrm{tors}\,H_0(\Ga;M)$ decides (no closed-form rule is asserted). At $\ell'\in\{2,3\}$ there is no such rule, and a prime dividing $\mathrm{tors}\,\Hcpt{2}$ gives a deficient rank
    (v1.23.0-pre.27, referee M2, replacing ``not known a priori'');
  \item[(d)] cone, $\Dcone^{q}$, $q\ge1$: the same cone modulo any prime $\ell'\ge5$ (not dividing $\mathsf q_\ell$; $\Hcpt{q+2}$ is killed by $48$), or the recursion with
    $b_1(\Cone{\rres})=\dim\Hcpt{2}$, $b_q(\Cone{\rres})=0$ for $q\ge2$.
  \end{enumerate}
  Not admissible: the rank of a differential of $\mathsf P(M)$ or $\mathsf P^{\partial}(M)$ alone for the cone, of another model, or of the cone at a non-qualifying prime.
  \textbf{A-priori exponents} (consistency checks, and caps): $a_t\le v_\ell(2)$ on $\Dcone^{\partial,q}$, $q\ge1$; on $\Dcone^{\partial,0}$, $\sum_ta_t=v_\ell\abs{\mathrm{tors}\,H^{1}(\Gbdry;M)}$, for $M_N$,
  $n$ even, $\sum_cv_\ell(n!\,w_c^{n})$; $a_t=0$ on $\Dcone^{-1}$; $a_t\le v_\ell(48)$ on $\Dcone^{q}$, $q\ge1$; on $\Dcone^{0}$ ($\mathrm{tors}\,\Hcpt{2}$), at $\ell\ge5$, $\sum_ta_t=v_\ell\abs{\mathrm{tors}\,H_0(\Ga;M)}$ ((5)(e); $0$ for $M_N$ when $\ell>n+1$),
  \textbf{no cap at $\ell\in\{2,3\}$} (v1.23.0-pre.27, referee M2). \textbf{The truncation}
  is \ref{open:ell2}'s box: $H^{q}(\Gbdry)$ and $\Hcpt{q}$ are reported for $q\le q_{\max}$, \eqref{eq:descentpairles} is asserted exact at every spot whose three terms are reported; by (4) and (5)(d)
  nothing is lost for $q_{\max}\ge4$ --- \textbf{the infinite $2$-adic boundary is a closed form in every degree, and the truncation reports a finite, periodic part of it}.
\end{enumerate}
\end{proposition}

\begin{proof}
(1)(a) $\mathsf P$ resolves $S$, and $\mathsf P^{\partial}$ resolves $S[\mathbb{P}^1]=S[\Ga]\otimes_{S[\Gbdry]}S$, so $H^{\bullet}(\mathsf P^{\partial}(M))=\Ext^{\bullet}_{S[\Ga]}(S[\Ga/\Gbdry],M_S)=
H^{\bullet}(\Gbdry;M_S)$ \STD{\cite[Ch.~III.6]{brown}}; the inclusion $\mathsf P^{\partial}\subseteq\mathsf P$ lies over $\aug\colon S[\mathbb{P}^1]\to S$, which induces the restriction
\STD{\cite[Ch.~III.8--9]{brown}}. (b) Restriction to a direct summand is split surjective. Lemma~\ref{lem:conerelative}(2)'s proof of the cone sequence and its connecting map uses only that
$\rres$ is a cochain map. (c) $\ker\rres_{\mathsf P}=\Hom_{S[\Ga]}(\mathsf P/\mathsf P^{\partial},M_S)$, and $c\mapsto(c,0)$ is a quasi-isomorphism $(\ker\rres)[1]\to\Cone{\rres}$ by
Lemma~\ref{lem:conerelative}(3)'s argument (surjectivity of $\rres$). The long exact homology sequence of $0\to\mathsf P^{\partial}\to\mathsf P\to\mathsf P/\mathsf P^{\partial}\to0$, with
$H_0(\mathsf P^{\partial})=S[\mathbb{P}^1]\xrightarrow{\aug}S=H_0(\mathsf P)$ onto and all other homology zero, shows that $\mathsf P/\mathsf P^{\partial}$, a bounded below complex of
projectives, has homology $\mathrm{St}_S$ in degree $1$ only ($\mathrm{St}$ is $\Z$-free, so $\mathrm{St}\otimes S=\ker(S[\mathbb{P}^1]\to S)$). Splitting off the contractible part in degrees
$\le1$ ($H_0=0$, $\partial_1$ onto the projective degree-$0$ term) leaves a projective resolution of $\mathrm{St}_S$ shifted by one, so the Hom complex computes
$\Ext^{\bullet-1}_{S[\Ga]}(\mathrm{St}_S,M_S)$; and $\Ext_{S[\Ga]}(\mathrm{St}_S,M_S)=H^{\bullet}(\Ga;\Hom_\Z(\mathrm{St},M_S))$ because $\mathrm{St}$ is $\Z$-free
\STD{\cite[Ch.~III, Prop.~2.2]{brown}}. The identification with the $\Ext$-sequence: $\mathsf P^{\partial}\to\mathsf P\to\mathsf P/\mathsf P^{\partial}$ is a short exact sequence of projective
resolutions over $0\to\mathrm{St}\to\Z[\mathbb{P}^1]\to\Z\to0$ shifted, and $\Hom(-,M_S)$ of it is the cone sequence \STD{\cite[Thm.~1.3.1, \S1.5]{weibel}}. \emph{The signs} (v1.23.0-pre.27,
referee m4). Write $\mathsf P=\mathsf P^{\partial}\oplus C$ as graded modules, $C$ the projective complement, and $\partial=\smat{\partial^{\partial}}{\beta}{0}{\partial^{C}}$; $\partial^{2}=0$ gives
$\partial^{\partial}\beta=-\beta\partial^{C}$, and $C\cong\mathsf P/\mathsf P^{\partial}$, so $R_q=C_{q+1}$ with $\partial^{R}=-\partial^{C}$ (a complex of projectives with homology $\mathrm{St}_S$ in degree $0$ only;
$R_{-1}=C_0$ splits off, $H_0(\mathsf P/\mathsf P^{\partial})=0$). (i) $\beta\colon R\to\mathsf P^{\partial}$ is a chain map ($\partial^{\partial}\beta=\beta\partial^{R}$) with $\varepsilon^{\partial}\beta=\varepsilon_R$, a lift of $j_{\mathrm{St}}$,
so $j_{\mathrm{St}}^{\ast}[x]=[x\circ\beta]$. Under $c\mapsto(c,0)$, $\coneinc{}_{\ast}[c_\partial]=[(0,c_\partial)]=[(d\tilde c,0)]$ ($\tilde c$ the extension by zero: $(0,c_\partial)-\Dcone(\tilde c,0)=(d\tilde c,0)$), and
$d\tilde c=\tilde c\circ\partial$ vanishes on $\mathsf P^{\partial}$ and is $c_\partial\circ\beta$ on $C$: $\coneinc{}_{\ast}=j_{\mathrm{St}}^{\ast}$. (ii) $E:=R\oplus\mathsf P$ with $\partial^{E}(r,y):=(\partial^{R}r+\pi_Cy,\ \partial y)$
($\pi_C$ the projection, a chain map $\mathsf P\to C$) is a complex ($\partial^{R}\pi_C+\pi_C\partial=0$) in a short exact sequence $0\to R\to E\to\mathsf P\to0$ of bounded below complexes of
projectives, each with homology in degree $0$ only; $\kappa(r,a,c):=a-\beta r$ ($a\in\mathsf P^{\partial}$, $c\in C$) is a chain map $E\to\mathsf P^{\partial}$, and $\varepsilon^{\partial}\kappa$ makes $H_0(E)\cong S[\mathbb{P}^1]$ with
$R\to E$ over $-j_{\mathrm{St}}$ and $E\to\mathsf P$ over $\aug$ (on a $0$-cycle, $c=\partial^{C}r$, the two augmentations differ by $\aug(\partial r)=0$). For this resolution of
$0\to\mathrm{St}\xrightarrow{-j_{\mathrm{St}}}S[\mathbb{P}^1]\to S\to0$, $\tilde x(r,y):=x(r)$ lifts a cocycle $x$ of $\Hom(R,M_S)$ and $d\tilde x=x\circ\pi_C$, the extension by zero, which is
$\coneproj{}_{\ast}[x]$. The isomorphism $(-1,\id,\id)$ from the sequence with $j_{\mathrm{St}}$ to the one with $-j_{\mathrm{St}}$ and the naturality of $\delta$ give $\coneproj{}_{\ast}=-\delta$. $\Hcpt{1}=\Hom_{\Z[\Ga]}(\mathrm{St},M)\subseteq\Hom_\Z(\mathrm{St},M)$, torsion-free as $M$ is. For (c) of Definition~\ref{def:descentpair}
the cone is Definition~\ref{def:pair}'s. (d) The ranks: Proposition~\ref{prop:isomodel}(2), $\Cone{\rres_{\Iso}}^{q}=C^{q+1}_{\Iso}(M)\oplus C^{q}_{\Isob}(M)$. At $\ell=2$: $Q_\sigma=1$ on the boundary, so
$F'^{\sigma}=F^{\sigma}\otimes\mathsf O_2$, free. At $\ell\ge3$: $G'_{s,0}=\mathsf O_\ell[\Ga]\otimes_{\mathsf O_\ell[\Ga_\sigma]}\mathsf O_\ell=\mathsf O_\ell[\Ga/\{\pm I\}]$ for $\sigma\in\Sigma^{\partial}$, and
$\cmpI'$ on it is the identification with $C_s(\Xbdry)\otimes\mathsf O_\ell$; there is no $t\ge1$ on the boundary, so $d'|_{{\Isob}^{(\ell)}}=d'^{(1)}$, and $\cmpI'd'^{(1)}=\partial\cmpI'$ (\ref{IM1}) with
$\cmpI'$ bijective forces $d'^{(1)}=\partial$ under it. Then ${\Isob}^{(\ell)}(M)=\Hom_{\Ga}(C_{\bullet}(\Xbdry),M_{\mathsf O_\ell})$, which Proposition~\ref{prop:circlemodel}(4) identifies with the circle
model. The E8 pair: the shipped item-4b complexes are $\Ceq(\Z/\ell^{K})$ and $\Ceq(\Kbdry,\Rel_\partial;L)$ over $\Z/\ell^{K}$, which are Definition~\ref{def:pair}'s with $\cring=\Z/\ell^{K}$
(Proposition~\ref{prop:circlemodel}(4)), and $\Ceq(\Z[1/6])\otimes\Z/\ell^{K}=\Ceq(\Z/\ell^{K})$ (Proposition~\ref{prop:ringinforce}(1)), the same on the boundary.

(2) $\cmpI(\Isob)\subseteq C_{\bullet}(\Xbdry)$ is \ref{IM2}; on $\Isob_0$ the augmentations to $\Z[\mathbb{P}^1]$ agree by definition; so $\cmpI$ is a comparison of pairs once both are
$\Z[1/6]$-pair models, and Lemma~\ref{lem:pairlifting}(3) applies. (a) Proposition~\ref{prop:descent}'s proof with the cones in place of $\Ceq$, $\CI$: Proposition~\ref{prop:fracture} for the
complex of free $\Z$-modules $\Cone{\rres_{\Iso}}$, $\alpha_c$ a quasi-isomorphism over $\Z[1/6]$, the five lemma on the cocones, and Proposition~\ref{prop:cohglue}. (b) The strict identities
are Lemma~\ref{lem:pairlifting}(3) and $\cmpI^{\ast}$'s base-change compatibility (Proposition~\ref{prop:isomodel}(4)); Proposition~\ref{prop:cohglue}'s morphism clause determines each map by its
two local components; the localisation $G_{(\ell)}$ of a finitely generated abelian group at $\ell\ge5$ is $(G[1/6])_{(\ell)}$, and $G=0$ iff $G_{(\ell)}=0$ for every $\ell$; exactness at a spot is the
vanishing of a subquotient, and localisation is exact. (c) is Proposition~\ref{prop:descentmaps}(2)(g) applied to the differentials and maps of the pieces.

(3) $\Ext^{\bullet}_{\Z[\Ga]}(Y,\Ind_A^{\Ga}X)\cong\Ext^{\bullet}_{\Z[A]}(\operatorname{Res}Y,X)$ for every $\Z[\Ga]$-module $Y$ (Shapiro, \STD{\cite[Ch.~III.6]{brown}}: $\Ind=\mathrm{CoInd}$ in finite index),
with $Y=\Z[\mathbb{P}^1]$ and $Y=\mathrm{St}$; as an $A$-set $\mathbb{P}^1=\bigsqcup_{c}A/A_c$, so $\Z[\mathbb{P}^1]=\bigoplus_c\Ind_{A_c}^{A}\Z$ and Shapiro again. For $A=\Ga_0(N)$:
$A\backslash\Ga/\Gbdry\to A\backslash\mathbb{P}^1$, $A\gamma\Gbdry\mapsto A\gamma\infty$, is a bijection, and $A\backslash\Ga$ is $\mathbb{P}^1(\Z/N)$ through bottom rows
(Lemma~\ref{lem:cosetrepsN}(1)) with $\Gbdry$ acting on the right through $T$ ($-I$ acts trivially on $\mathbb{P}^1(\Z/N)$); the stabiliser of $Aw$ in $\Gbdry$ is $\Gbdry\cap w^{-1}Aw\ni-I$, of
index $w_c$, hence $\langle T^{w_c}\rangle\times\langle-I\rangle$. In the coordinates of \eqref{eq:MH} the orbit's blocks span a $\Gbdry$-submodule permuted transitively with block stabiliser
$\Gbdry^{(w_c)}$, acting on the block of $w$ by $\rhon(h_w(\gamma))$, $h_w(T^{w_c})=wT^{w_c}w^{-1}\in\Ga_0(N)$; a module whose summands are permuted transitively is induced from a block stabiliser
\STD{\cite[Ch.~III, Prop.~5.3]{brown}}, and $\rhon(w)$ intertwines $\rhon(wT^{w_c}w^{-1})$ with $\rhon(T^{w_c})$ ($\rhon$ is multiplicative, Proposition~\ref{prop:symconv}(2)). With nebentypus the
character value on $wT^{w_c}w^{-1}$ multiplies the block.

(4) $\Gbdry^{(w)}=\langle T^{w}\rangle\times\langle-I\rangle$ and Proposition~\ref{prop:bdrycohom}'s proof (the Lyndon--Hochschild--Serre sequence with two columns) applies with $T^{w}$ for $T$:
$0\to(H^{q-1}(\langle-I\rangle;X))_t\to H^{q}\to(H^{q}(\langle-I\rangle;X))^{t}\to0$. For $X$ $\Z$-free, the periodic complex $X\xrightarrow{\varepsilon-1}X\xrightarrow{1+\varepsilon}X\xrightarrow{\varepsilon-1}\cdots$
gives $H^{0}(\langle-I\rangle;X)=X$, $H^{\rm odd}=0$, $H^{\rm even\ge2}=X/2$ for $\varepsilon=1$, and $H^{0}=0$, $H^{\rm odd}=X/2$, $H^{\rm even}=0$ for $\varepsilon=-1$; in each degree one outer term
vanishes, which gives the display without an extension problem. $(X/2)^{t}$ and $(X/2)_t$ are the kernel and cokernel of an endomorphism of a finite-dimensional $\F_2$-space. The vanishing
statements follow. For $\Vcoef$: dehomogenise $u_j=x^{n-j}y^{j}\mapsto y^{j}$, $\Vcoef\cong\Z[y]_{\le n}$; by \eqref{eq:rhondef} with $T^{w}=\smat1w01$, $\rhon(T^{w})u_j=x^{n-j}(wx+y)^{j}\mapsto(y+w)^{j}$, so
$\rhon(T^{w})-1$ is $\Delta_w\colon P(y)\mapsto P(y+w)-P(y)$. $\Delta_wP=0$ forces $P$ constant ($P(y+kw)=P(y)$ for all $k$), so $\Vcoef^{T^{w}}=\Z x^{n}$ ($u_0=x^{n}$); $\Delta_w$ maps $\Z[y]_{\le n}$
into $\Z[y]_{\le n-1}$, and in the bases $(y^{1},\dots,y^{n})$, $(y^{0},\dots,y^{n-1})$ its matrix is triangular with diagonal $(jw)_{j=1}^{n}$ ($\Delta_wy^{j}=jwy^{j-1}+\cdots$); so
$\Z[y]_{\le n-1}/\Delta_w\Z[y]_{\le n}$ is finite of order $\prod_jjw=n!\,w^{n}$, and $(\Vcoef)_{T^{w}}=\Z y^{n}\oplus$ it. Modulo $2$: if $w$ is even, every coefficient
$\binom ji w^{j-i}$, $i<j$, of $\Delta_w$ is even and $\Delta_w\equiv0$; if $w$ is odd, $P(y+w)=P(y+1)$ over $\F_2$, and the kernel of $\Delta_1$ on $\F_2[y]_{\le n}$ is $\F_2[y^{2}+y]_{\le n}$ (a
polynomial invariant under $y\mapsto y+1$ is a polynomial in $y^{2}+y$, by division with remainder), of dimension $\lfloor n/2\rfloor+1$. $M_N$ by (3). \emph{Periodicity}: $H^{2}(\Gbdry^{(w)};\Z)\cong
H^{2}(\langle-I\rangle;\Z)=\Z/2$ by restriction (the case $X=\Z$, $\varepsilon=1$, $t=1$ of the display: $(X/2)^{t}=\Z/2$), and $\res e$ restricts to a generator on $\langle-I\rangle$
(Lemma~\ref{lem:bulkperiodic}(1)); so $\res_{\Gbdry^{(w)}}e=\mathrm{pr}^{\ast}u_0$, $u_0$ the generator of $H^{2}(\langle-I\rangle;\Z)$, $\mathrm{pr}$ the projection onto $\langle-I\rangle$. (v1.23.0-pre.27,
referee m1: the tensor resolution, no multiplicativity of a spectral sequence.) Let $F$ be $\Z[\langle T^{w}\rangle]\xrightarrow{T^{w}-1}\Z[\langle T^{w}\rangle]$ and $F'$ the periodic resolution of $\Z$ over
$\langle-I\rangle$ ($F'_b=\Z[\langle-I\rangle]$, differentials $\sigma-1$, $\sigma+1$ alternately, $\sigma=-I$); $F\otimes F'$ is a free resolution of $\Z$ over $\Gbdry^{(w)}$ \STD{\cite[Ch.~V.1]{brown}}. With
$K:=\Hom_{\langle-I\rangle}(F',X)$, $\Hom_{\Gbdry^{(w)}}(F\otimes F',X)$ is the cocone of $t-1\colon K\to K$, whose long exact sequence $H^{q-1}(K)\xrightarrow{t-1}H^{q-1}(K)\to H^{q}\to H^{q}(K)\xrightarrow{t-1}H^{q}(K)$
is the two-column sequence above. The shift $s\colon F'_{b+2}\to F'_b$ (the identity of $\Z[\langle-I\rangle]$) is a chain map; $s^{\ast}$ is cup product with the generator of
$H^{2}(\langle-I\rangle;\Z)=\Z/2$ \STD{\cite[Ch.~VI.9]{brown}} (a unique generator, so no sign), and $(\id\otimes s)^{\ast}$ is cup product with $\mathrm{pr}^{\ast}u_0$ (the cross product with $u_0$,
\STD{\cite[Ch.~V.2--3]{brown}}). It commutes with $t-1$, so maps the sequence of degree $q$ to that of degree $q+2$; on $H^{b}(K)=H^{b}(\langle-I\rangle;X)$ it is an isomorphism for $b\ge1$, and for $b=0$
when $\varepsilon=-1$ (both sides vanish). The five lemma gives the claim for $q\ge2$, and for $q\ge1$ if $\varepsilon=-1$.

(5)(a) Under (1)(a), $\rres^{0}$ on $H^{0}$ is the inclusion of invariants; $H^{-1}(\Gbdry)=0$. So $\Hcpt{0}=0$ over every $S$, i.e.\ $\ker(\Dcone^{-1}\otimes F)=H^{-1}(\Cone{\rres}\otimes F)=0$ for every
field $F$; an integer matrix of full column rank modulo every prime has unit elementary divisors. (b) With $6\in F^{\times}$, Lemma~\ref{lem:manin}'s sequence $0\to F[\Ga/\langle S\rangle]\oplus
F[\Ga/\langle U\rangle]\to F[\Ga/\{\pm I\}]\to\mathrm{St}_F\to0$ is a projective resolution of length $1$ (the stabilisers have orders $4$, $6$, $2$), so by (1)(c) $\Hcpt{q}=\Ext^{q-1}(\mathrm{St}_F,M_F)=0$ for
$q\ge3$ and $\Hcpt{2}=\coker\bigl(\mu\colon M_F^{\{\pm I\}}\to M_F^{\langle S\rangle}\oplus M_F^{\langle U\rangle}\bigr)$, $\mu(v)=((1+S)v,(1+U+U^{2})v)$ (pull back along
$\gamma\mathfrak v_S\mapsto\gamma\mathfrak e+\gamma S\mathfrak e$, $\gamma\mathfrak v_U\mapsto\gamma\mathfrak e+\gamma U\mathfrak e+\gamma U^{2}\mathfrak e$). For a finite subgroup $H$ the invariants $M_F^{H}$ and
$(M_F^{\vee})^{H}$ are in perfect duality (the averaging idempotent is self-adjoint), and under it the adjoint of $v\mapsto(1+S)v$, $M^{\{\pm I\}}\to M^{\langle S\rangle}$, is
$\lambda\mapsto\lambda+S^{-1}\lambda=2\lambda$, $(M^{\vee})^{\langle S\rangle}\to(M^{\vee})^{\{\pm I\}}$, that of $1+U+U^{2}$ is $3\lambda$. So $\dim\coker\mu=\dim\ker\bigl((\lambda,\lambda')\mapsto2\lambda+3\lambda'\bigr)$,
and $2\lambda=-3\lambda'$ lies in $(M^{\vee})^{\langle S\rangle}\cap(M^{\vee})^{\langle U\rangle}=(M^{\vee})^{\Ga}$ ($S$, $U$ generate $\Ga$), the kernel being isomorphic to it. The two other formulas are
the exactness of the rational sequence $0\to H^{0}\to H^{0}(\Gbdry)\to\Hcpt{1}\to H^{1}\to H^{1}(\Gbdry)\to\Hcpt{2}\to H^{2}=0$ (Proposition~\ref{prop:certifiedeell}(4)(a): $24H^{2}(\Ga;M)=24H^{3}(\Ga;M)=0$, $24\in F^{\times}$), and $\dim H^{0}(\Gbdry;M_F)=\dim H^{1}(\Gbdry;M_F)$
(kernel and cokernel of $t-1$ on $M_F^{-I}$, (4) with $2\in F^{\times}$). For $M_N(\chi)$: $(M^{\vee}_\Q)^{\Ga}=((V(\chi)^{\vee})_\Q)^{\Ga_0(N)}$ (Shapiro in degree $0$), and $V(\chi)^{\vee}_\Q\cong\Sym^{n}\otimes\chi^{-1}$
has invariants iff $n=0$, $\chi=1$ (the argument of Proposition~\ref{prop:descentN}(4)); for $M_N$, $n$ even, $\dim H^{1}(\Gbdry;M_{N,\Q})=c(N)$ by (4). (c) For $q\ge3$: $H^{q-1}(\Gbdry)$ is killed by $2$ ((4))
and $H^{q}(\Ga;M)$ by $24$ (Proposition~\ref{prop:certifiedeell}(4)(a)), and $\Hcpt{q}$ is an extension of a subgroup of the second by a quotient of the first. At $\ell\ge5$ both vanish for
$q\ge2$, resp.\ $q-1\ge2$, and $H^{2}(\Ga;M)\otimes\Z_{(\ell)}=0$; at $\ell=3$ the boundary terms vanish for $q-1\ge2$ ((4)), so $\coneproj{}_\ast$ is an isomorphism for $q\ge3$. (d) (v1.23.0-pre.27, referee m1 and n1: no extension of \cite[Ch.~V.4]{brown} to $\Ext(Y,-)$ is used, and the dangling ``Lemma~\ref{lem:bulkperiodic}(4)'' of pre.26 is gone.) By
\STD{\cite[Ch.~III, Prop.~2.2]{brown}} the three terms are $H^{\bullet}(\Ga;X)$ for $X=M$, $X=\Hom_\Z(\Z[\mathbb{P}^1],M)=\mathrm{Coind}_{\Gbdry}^{\Ga}M$ and $X=\Hom_\Z(\mathrm{St},M)$ (the last in
degree $q-1$), and \eqref{eq:descentpairles} is, up to the signs of (1)(c), the long exact sequence of the coefficient sequence $0\to M\to\Hom_\Z(\Z[\mathbb{P}^1],M)\to\Hom_\Z(\mathrm{St},M)\to0$,
which is $\Z$-split ($\mathrm{St}$ is $\Z$-free). $e\cup-$ on each is Lemma~\ref{lem:bulkperiodic}(2)'s (any $\Z[\Ga]$-module); it is induced by a cochain map natural in the coefficients (cup
product with a cocycle $\epsilon$ representing $e$ through a diagonal approximation, \STD{\cite[Ch.~V.1--3]{brown}}), so it commutes with the coefficient maps and, the sequence of cochain
complexes being exact, with the connecting maps ($d(\epsilon\cup\tilde x)=\epsilon\cup d\tilde x$: $\deg\epsilon=2$, $d\epsilon=0$); a sign of (1)(c) does not affect commutation. The Shapiro
isomorphism $H^{\bullet}(\Ga;\mathrm{Coind}\,M)\cong H^{\bullet}(\Gbdry;M)$ is $\res^{\Ga}_{\Gbdry}$ followed by the coefficient map ``evaluation at $1$'' \STD{\cite[Ch.~III.6]{brown}}, so it carries
$e\cup-$ to $\res_{\Gbdry}e\cup-$ (Lemma~\ref{lem:bulkperiodic}(2): compatible with restriction, natural in the coefficients) --- the Shapiro--cup compatibility for $\mathrm{Coind}$, here stated once;
it holds at infinite index. The isomorphisms: Lemma~\ref{lem:bulkperiodic}(2) for $X=M$ and $X=\mathrm{Coind}\,M$ in degrees $q\ge2$ (for the latter, equivalently item (4) above), and for
$X=\Hom_\Z(\mathrm{St},M)$ in degrees $q-1\ge2$, i.e.\ on $\Hcpt{q}$, $q\ge3$, directly.

(5)(e) Over $\Z_{(\ell)}$, $\ell\ge5$, $6$ is a unit, and the cone over $\Z_{(\ell)}$ is the cone over $\Z$ tensored with the flat $\Z_{(\ell)}$; so by (1)(c) and Lemma~\ref{lem:manin}'s
resolution, as in (b), $\Hcpt{2}\otimes\Z_{(\ell)}=\coker\mu$ over $\Z_{(\ell)}$, $\mu(v)=((1+S)v,(1+U+U^{2})v)$, $M^{\{\pm I\}}\to M^{\langle S\rangle}\oplus M^{\langle U\rangle}$. For a finite
$H\le\Ga$, $\abs H\in\{2,4,6\}$ is a unit, and the norm $N_H\colon M_H\to M^{H}$ is an isomorphism (its composites with the canonical $M^{H}\to M_H$ are multiplication by $\abs H$). Since
$S^{2}=U^{3}=-I$, $N_{\langle S\rangle}=(1+S)N_{\{\pm I\}}$ and $N_{\langle U\rangle}=(1+U+U^{2})N_{\{\pm I\}}$, so $\mu\circ N_{\{\pm I\}}=(N_{\langle S\rangle}\oplus N_{\langle U\rangle})\circ\nu$ with
$\nu\colon M_{\{\pm I\}}\to M_{\langle S\rangle}\oplus M_{\langle U\rangle}$, $[v]\mapsto([v],[v])$, and $\coker\mu\cong\coker\nu$. And $\coker\nu\cong M_{\Ga}$ through $([a],[b])\mapsto[a-b]$: it is onto, and if
$a-b=(S-1)x+(U-1)y$ (which is all of $\sum_{g}(g-1)M$, as $S$, $U$ generate $\Ga$: $gh-1=g(h-1)+(g-1)$) then $([a],[b])=\nu[v]$ with $v=a-(S-1)x=b+(U-1)y$. For $M_N$:
$(\Ind_A^{\Ga}\Vcoef)_{\Ga}=\Z\otimes_{\Z[A]}\Vcoef=\Vcoef_A$. Modulo $\ell$, $\Vcoef_A\otimes\F_\ell=(\Vcoef/\ell)_A$ is a quotient of $(\Vcoef/\ell)_B$, $B$ the upper triangular subgroup of
$\SL_2(\F_\ell)$, which lies in the image of $A$ for every $N$ ($\SL_2(\Z)\to\SL_2(\Z/N\ell)$ is onto, and the upper triangular matrices modulo $N\ell$ lift into $\Ga_0(N)$). The unipotent
coinvariants are one line for $n<\ell$ (by (4)'s computation, $\Delta_1$ is triangular with diagonal $1,\dots,n$, units modulo $\ell$), spanned by the class of $u_n$, on which $\mathrm{diag}(a,a^{-1})$
acts by $a^{\pm n}$; so $(\Vcoef/\ell)_B=0$ unless $(\ell-1)\mid n$, in particular for $0<n<\ell-1$. By Nakayama's lemma for the finitely generated $\Vcoef_A$, $\Vcoef_A\otimes\Z_{(\ell)}=0$. The
$\otimes\Q$ shadow: coinvariants and invariants of the dual are dual over $\Q$. (6) The cone is a bounded below complex of finitely generated free $\Z_{(\ell)}$-modules, so
the cited statements apply as in Proposition~\ref{prop:descentmaps}(2)(d). The sources: (a) by (4); (b) by (5)(a); (c) by (5)(b) and the Betti recursion, and its prime by (5)(e) with the criterion that follows; a reduction at $\ell'$ has the rational rank iff
$\ell'$ divides no elementary divisor, i.e.\ iff $\mathrm{tors}\,H^{q+1}(\Cone{\rres})=\mathrm{tors}\,\Hcpt{q+2}$ has no $\ell'$-torsion (Proposition~\ref{prop:certifiede}(4)); (d) by (5)(c). The
exponents: the elementary divisors of $\Dcone^{q}$ are the exponents of $\mathrm{tors}\,\Hcpt{q+2}$, those of $\Dcone^{\partial,q}$ of $\mathrm{tors}\,H^{q+1}(\Gbdry)$; then (4), (5)(a), (5)(c). The
truncation lemma of \ref{open:ell2} on bulk and boundary, and $\Cone{\rres_{\mathsf P}}^{q}=\mathsf P(M)^{q+1}\oplus\mathsf P^{\partial}(M)^{q}$.
\end{proof}

\begin{proposition}[The cone over $\Z$ with the strict pair datum, and $T_p$, $p\nmid N$, on the pair; D2's maps]\label{prop:descentpairhecke}
\PROVED{} \STD{\cite[Ch.~III.2, III.8--10]{brown}, \cite[Thm.~2.2.6]{weibel}; \VERIFYloc} (v1.23.0-pre.26; milestone D4, clause C-D7 of the plan; refereed ACCEPT-with-fixes, its fix M3 at
v1.23.0-pre.27; (1) and (5) amended at v1.23.0-pre.28 for the owner's scope addition of 2026-10-25: the pair lift for every $p\nmid N$ on demand, certified at run time; pre.28 refereed
ACCEPT-with-fixes, \file{docs/implementation/slot24/C-D7p\_referee.md}, its fixes M1, m1--m3, N1 at v1.23.0-pre.29.) Assume \ref{H4} and
\ref{B1}--\ref{B3}. Let $N\ge1$, $n\ge0$, $p\nmid N$ prime, $g=\smat100p$, $\Ga''_g=\Ga\cap g\Ga g^{-1}=\Ga_0(p)$, $A=\Ga_0(N)$, $M_N$ as in \S\ref{sec:descentN}.
\begin{enumerate}
\item \textbf{(The pair lift: for every $p$, produced; built on demand and certified at run time; v1.23.0-pre.28.)} Let $\taub$ be the boundary datum of Lemma~\ref{lem:taubdryexists}(1) for $g$ (integral, coefficients in
  $\{0,\pm1\}$, by a terminating procedure, every $p$), and $\tau^{s}$ the strict datum of Proposition~\ref{prop:strictpair} for it (terminating, every $p$; $\tau^{s}j=j\taub$). A \textbf{pair lift}
  is a $\tau_{\Iso}$ as in Proposition~\ref{prop:strictlift}(1): (T1)--(T3) over $\Z$, $\cmpI\tau_{\Iso}=\tau^{s}\cmpI$, $\tau_{\Iso}(\Isob)\subseteq\Isob$; $\tau^{\partial}_{\Iso}$ its restriction. For
  \textbf{every} $p$ one is \textbf{produced}, in every degree, by Remark~\ref{rmk:strictliftbuilt}(1)'s procedure (Proposition~\ref{prop:strictliftN}(2)'s, with the boundary step first):
  deterministic, no system solved, terminating; it is certified by Remark~\ref{rmk:strictliftbuilt}(2)'s checks, the boundary items included. Two pair lifts --- for any admissible $\taub$,
  $\tau^{s}$ --- are homotopic as pairs through (T1)-maps over $\Z$. \textbf{Built and certified} (v1.23.0-pre.28, the owner's scope addition of 2026-10-25, replacing pre.27's
  ``$p\in\{2,3,5\}$; for every other $p$ \dots\ not built''): for \textbf{every prime} $p$ requested, the lift is built \emph{on demand} by that procedure in degrees $\le6$, so the pair
  is served for $q_{\max}\le5$ ($q_{\max}+1\le6$; obligation (g) of the code map asks $q_{\max}\ge4$); a request for cochains above degree $6$ ($q_{\mathrm{top}}>6$, or $q_{\max}>5$
  for the pair) is \textbf{refused}, and higher degrees come from (7), \textbf{never from extending the lift} (v1.23.0-pre.29, referee m2). It is certified by \textbf{all} of
  Remark~\ref{rmk:strictliftbuilt}(2)'s checks --- (T1) on the Schreier generators of $\Ga_0(p)$, (T2), (T3), the strict identity
  $\cmpI\tau_{\Iso}=\tau^{s}\cmpI$, $\tau_{\Iso}(\Isob)\subseteq\Isob$ and $\cmpI\tau^{\partial}_{\Iso}=\taub\cmpI$ --- before its first use, and kept for later requests. A lift is
  \textbf{refused}, never truncated or used (Theorem~\ref{thm:backend}(3)), when (v1.23.0-pre.29, referee M1; pre.28 named only a failed certificate and an int64 product) its
  \textbf{predicted RSS} exceeds the ceiling ($4$\,GiB; rule~19, decided before anything is built), when its construction \textbf{does not complete within the code's checked bounds}
  --- the \emph{cap}: an exact product outside $64$-bit signed integers or a $\tau$-table chain coefficient outside $32$-bit signed integers (the exact-integer layer's overflow
  exception), $p$ (or $p^{r}$) outside the coset tables' $32$-bit unsigned domain, or a failed mandatory self-certification of an input table (Obligation~\ref{obl:tau}) --- or when its
  certificate fails. \textbf{Any other exception} during the construction (an internal check, an internal error) is a \textbf{failure of the code}, not a refusal: it propagates, and
  nothing is cached. \textbf{The certificate suffices on its own}, whatever $p$: the operator of (2) and the proof of (3) read of the lift only (T1)--(T3) over $\Z$
  and $\tau_{\Iso}(\Isob)\subseteq\Isob$, each re-derived by the certificate from the stored values, and these \textbf{force} the component map of $\tau^{\partial}_{\Iso}$ on
  $H_0(\Isob)=\Z[\mathbb{P}^1]$ to be $c\mapsto g^{-1}c$ (the proof of (1), without the strict identities; v1.23.0-pre.29, referee m1). The strict identities with $\tau^{s}$, $\taub$ tie the lift
  to its inputs, so that an input built wrongly at a new $p$ is refused, and are what the cellular comparisons of (4) read. Since the datum is a level-$1$ object, built once per $p$, independently of $N$ and of the transversal
  (Proposition~\ref{prop:strictliftN}(3)(c)), that is all level $N$ needs at $p\nmid N$. \textbf{No bound on its sizes is asserted} (Remark~\ref{rmk:integralscope}(2)); what was measured
  (evidence, bounding nothing): largest coefficient $\le3$ and largest matrix entry $p$ for $p\le127$, certification time $0.12$\,s at $p=13$, $26$\,s at $p=127$, growing
  $\approx p^{2.4}$ (\file{docs/implementation/slot24/DESCENT4\_report.md} \S4). These sizes are a property of the \emph{shipped model}, not of the procedure (v1.23.0-pre.29, referee N1): on another
  free model of the shape of Definition~\ref{def:isomodel} the same procedure gives largest coefficients $1,1,3,4,9,13,27,49,81,157,243,481$ in degrees $0$--$11$ at $p=7$, growing
  $\approx\times3$ per two degrees (\file{docs/implementation/slot24/C-D7p\_referee.md}, N1), and the shipped model's stationary term counts (Remark~\ref{rmk:strictliftbuilt}(5)) are
  model-specific; hence (7), not a longer lift, is the route above degree $6$. Which $p$ a run has certified is a run-time fact, not a statement of this document.
\item \textbf{(The level-$N$ operator on the pair.)} For a pair lift, $\widetilde T^{(N)}_{\Iso,p}$ is Proposition~\ref{prop:strictliftN}(3)'s operator on $\CI(M_{N,\cring})$, and
  $\widetilde T^{(N)}_{\Isob,p}$ the same formula \eqref{eq:levelcorr} with $\tau^{\partial}_{\Iso}$ and $x$ running over $\Isob$, on $C^{\bullet}_{\Isob}(M_{N,\cring})$. Over every $\cring$ both are integral cochain
  maps, $\rres\circ\widetilde T^{(N)}_{\Iso,p}=\widetilde T^{(N)}_{\Isob,p}\circ\rres$ \textbf{on the nose}, and
  \[
    \Phi^{(N)}_p:=\bigl(\widetilde T^{(N)}_{\Iso,p}[1],\ \widetilde T^{(N)}_{\Isob,p}\bigr)
  \]
  is a cochain endomorphism of $\Cone{\rres_{\Iso}}(M_{N,\cring})$ with \textbf{$\hder=0$}, commuting with $\coneinc$ and $\coneproj$ and preserving $(\ker\rres)[1]$. At $N=1$ it is the cone
  operator of the Hecke datum of the pair $(\tau_{\Iso},\tau^{\partial}_{\Iso},0)$ (Proposition~\ref{prop:strictlift}(1)).
\item \textbf{(What it induces; independence; the long exact sequence.)} On $H^{q}(A;\VR)$, $\coneproj$-compatibly, $\Phi^{(N)}_p$ induces $T_p$ --- \textbf{the operator of \S\ref{sec:descenthecke}}:
  the bulk component of a pair lift is a bulk lift in the sense of Proposition~\ref{prop:strictliftN}(2) (with $\tau_X=\tau^{s}$), and the induced operator does not depend on the lift
  (Proposition~\ref{prop:strictliftN}(3)(c)). On $\Hcpt{q}(A;\VR)=\Ext^{q-1}_{\Z[A]}(\mathrm{St},\VR)=H^{q-1}(A;\Hom_\Z(\mathrm{St},\VR))$ it induces the double-coset operator
  $[AgA]=\cor^{A}_{A''}\circ(\iota_g,f)^{\ast}\circ\res^{A}_{A'}$ of Proposition~\ref{prop:descentN}(5) for the compatible pair $(\Hom_\Z(\mathrm{St},\VR),f)$,
  $f(\phi):=\rhon(g)\circ\phi\circ g^{-1}$ ($g^{-1}$ acting on $\mathrm{St}\subseteq\Z[\mathbb{P}^1]$ by M\"obius transformations); on $H^{q}(\Gbdry;M_N)=\Ext^{q}_{\Z[A]}(\Z[\mathbb{P}^1],\VR)$ the same with
  $\Z[\mathbb{P}^1]$ for $\mathrm{St}$. The induced operators are independent of the pair lift, $\taub$, $\tau^{s}$, the $s_j$ and the transversal, \textbf{over $\Z$, for every $\ell$}; so
  \eqref{eq:descentpairles} for $A$ is $T_p$-equivariant over $\Z$ and over every $\Z_{(\ell)}$, \textbf{$\ell\in\{2,3\}$ included}, in every degree, and $T_pT_{p'}=T_{p'}T_p$ on its three terms.
\item \textbf{(The pieces.)} For an $\mathsf O_\ell$-algebra $S$ and comparisons of pairs $\Phi_\ell\colon(\Iso,\Isob)\otimes S\to(\Iso^{(\ell)},{\Isob}^{(\ell)})\otimes S$, $\Psi_\ell$ back
  (Lemma~\ref{lem:pairlifting}; Proposition~\ref{prop:isomodelell}(4)), the transport $(\Psi_\ell^{\ast}[1],\Psi_\ell^{\partial\ast})\circ\Phi^{(N)}_p\circ(\Phi_\ell^{\ast}[1],\Phi_\ell^{\partial\ast})$ is a cochain
  endomorphism of the $\ell$-local cone inducing the same operator, independently of the comparisons. On the $\Z[1/6]$-piece: at $N=1$ the cone operator of Definition~\ref{def:coneaction} with the shipped canonical
  $(\tau,\taub,\etapair)$ and $\hder=\widetilde{\etapair}$ (where built, \S\ref{sec:codemap}) induces the same operator on $\Hcpt{\bullet}\otimes\Z[1/6]$, for \textbf{every} $p$ (Theorem~\ref{thm:conewelldef}(b), (c),
  Proposition~\ref{prop:coneindep}, Lemma~\ref{lem:pairlifting}(3) for $\cmpI$); at level $N$ the Bredon pair carries \eqref{eq:levelcorr} with the strict cellular pair $(\tau^{s},\taub)$, and
  $\alpha_c$ of Proposition~\ref{prop:descentpair}(2) intertwines it with $\Phi^{(N)}_p$ \textbf{on the nose} (Proposition~\ref{prop:strictliftN}(4)'s computation on bulk and boundary). $T_p$ on the
  glued sequence over $\Z$ is determined by these local components (Proposition~\ref{prop:cohglue}'s morphism clause on the cone).
\item \textbf{(What is certified at which $p$.)} The groups and maps of \eqref{eq:descentpairles}: no $p$ enters. $T_p$ on the bulk: every $p\nmid N$ (\S\ref{sec:descenthecke}). $T_p$ on the boundary and
  the cone (v1.23.0-pre.27, referee M3: weakened from pre.26's ``at $\ell\ge5$, every $p\nmid N$'', so that statement and code map agree; v1.23.0-pre.28, the owner's scope addition:
  the pair lift is built on demand, (1)): \textbf{certified, at every $\ell$ and every $N$, for every prime $p\nmid N$ for which a pair lift has been built and its certificate passed (at
  run time)} ((1), (2) over every $\cring$), with the strict identity $\rres\circ\widetilde T^{(N)}_{\Iso,p}=\widetilde T^{(N)}_{\Isob,p}\circ\rres$ checked entrywise in every degree used (a
  theorem by (2), checked as a consistency check of the code; a failure refuses); \textbf{at $\ell\ge5$ and level $1$, also at every $p$ at which the shipped cone operator of
  Definition~\ref{def:coneaction} is built} (the $\Z[1/6]$-piece of (4), \S\ref{sec:codemap}). \textbf{Refused exactly when} (v1.23.0-pre.29, referee M1; each refusal names its
  cause, by one of the code's eleven refusal codes, listed as the code enumerates them): (a) $p$ is not prime; (b) $p\mid N$ ($U_p$ deferred; outside this proposition); (c) the bulk-only lift is requested for the cone
  (below); (d) cochains above degree $6$ are requested ($q_{\mathrm{top}}>6$; for the pair, $q_{\max}>5$; (1)); (e) the pair lift's predicted RSS exceeds the ceiling, decided before
  building ((1)); (f)--(h) the cap of (1), during the lift's construction \emph{or the operator's assembly} --- (f) an exact product outside $64$-bit or a $\tau$-table coefficient outside
  $32$-bit signed integers (and, sharing that exception type in the code, an internal guard of the exact integer geometry: the code map's (f)), (g) $p$ ($p^{r}$) outside the coset tables' $32$-bit unsigned domain, (h) a failed self-certification of an input table; (i) the pair lift's certificate fails;
  (j) on the operator's layout a term of $\tau_{\Iso}(\Isob)$ lies off $\Isob$; (k) the strict identity $\rres\circ\widetilde T^{(N)}_{\Iso,p}=\widetilde T^{(N)}_{\Isob,p}\circ\rres$ fails
  entrywise, $\widetilde T^{(N)}_{\Isob,p}$ assembled separately from $\tau^{\partial}_{\Iso}$. (i)--(k) cannot occur for a correct implementation ((j) is implied by (i)'s item
  $\tau_{\Iso}(\Isob)\subseteq\Isob$, and (k) is a theorem by (2)); they are consistency checks of the code. \textbf{No other condition refuses}: any other exception is a failure of the
  code, never a refusal, and no refused value is truncated or used. \textbf{In cases (d)--(k)} the operator still \textbf{exists and is determined} ((1)--(4), and (7) above degree $6$; at
  $\ell\ge5$ it is also computable from the strict cellular pair $(\tau^{s},\taub)$ at $p$, both terminating, with \eqref{eq:levelcorr} on the level-$N$ Bredon pair, (4)) but is \textbf{not
  asserted as computed}: that cellular route is not built (the code map's (f)). The bulk-only
  lift of Proposition~\ref{prop:strictliftN}(2), in particular one from the canonical $\tau$, need not preserve $\Isob$ (Remark~\ref{rmk:conestrict}); it is \textbf{not} a pair datum, and is never
  used on the cone (that would need a pair homotopy on $\Iso$, not constructed here).
\item \textbf{(D2's maps on the pair.)} Let $u\colon M\to M'$ be a $\Z[\Ga]$-map of finitely generated $\Z$-free modules ($\alpha_{B/A}$, $\omega_{B/A}$, $R_\sigma$ of
  Proposition~\ref{prop:descentmaps}). On every pair model $(u_{\mathsf P}[1],u_{\mathsf P^{\partial}})$ is a cochain map of cones commuting with $\rres$, $\coneinc$, $\coneproj$ and with every comparison
  of pairs, strictly; so $\res^{A}_{B}$, $\cor^{A}_{B}$ and $c_\sigma$ act on \eqref{eq:descentpairles} over $\Z$ as a commuting ladder of long exact sequences, glued from the pieces, in every degree
  (on the boundary: $\res$, $\cor$ between the cusp blocks of $A$ and of $B$). For $B=\Ga_0(MN)$ with $p\nmid MN$ and a common transversal (Proposition~\ref{prop:strictliftN}(7)):
  $\widetilde T^{(MN)}_{\Isob,p}\alpha=\alpha\,\widetilde T^{(N)}_{\Isob,p}$ as well, so the ladder for $\res$ is $T_p$-equivariant, over $\Z$ and every $\Z_{(\ell)}$, at the $p$ of (5).
  $\cor\circ T_p=T_p\circ\cor$ on the three terms on cohomology (Proposition~\ref{prop:descentmaps}(3)(a)'s argument, which is formal in the coefficient module). For $p\mid M$ nothing is asserted
  (Proposition~\ref{prop:descentmaps}(3)(b)); of $U_p$, $p\mid N$, nothing is asserted here (its lift on the free model: Proposition~\ref{prop:strictliftN}(9), v1.26.0-pre.1; v1.26.0-pre.11, confirm-referee C1).
\item \textbf{(Every degree, Hecke-equivariantly.)} For $q\ge3$, $e\cup-\colon\Hcpt{q}\to\Hcpt{q+2}$ commutes with $T_p$ (Proposition~\ref{prop:descentN}(5) for the compatible pair of (3)); so
  $q_{\max}=4$ determines the Hecke modules $\Hcpt{q}$ in every degree, and $q_{\max}=3$ those of $H^{q}(\Gbdry)$, $q\ge2$, likewise.
\end{enumerate}
\end{proposition}

\begin{proof}
(1) Existence and the procedure: Proposition~\ref{prop:strictlift}(1)'s proof (Lemma~\ref{lem:strictlifting}, first on $\Isob$ with $f=\taub\circ\cmpI$, then on $\Iso$ relative to it) for $g=\smat100p$, whose
inputs $\taub$ (Lemma~\ref{lem:taubdryexists}(1), stated for $\smat100N$, every $N$) and $\tau^{s}$ (Proposition~\ref{prop:strictpair}, every $g\in\De$) exist for every $p$; the procedure is
Remark~\ref{rmk:strictliftbuilt}(1)'s, whose read-off of $y_0$ and contraction by descent through the columns use nothing of $p$ but the coset table of $\Ga''_g\backslash\Ga$, exactly as in
Proposition~\ref{prop:strictliftN}(2)'s proof. The boundary step is that procedure on $\Isob$ (v1.23.0-pre.28, spelled out for general $p$): $\Isob$ is the sub-sum of the columns over
$\Sigma^{\partial}_0\sqcup\Sigma^{\partial}_1$ (Definition~\ref{def:descentpair}(a)), a subcomplex on which $\cmpI$ is a surjective quasi-isomorphism onto $C_{\bullet}(\Xbdry)$; the read-off of $y_0$ at
$t=0$ lands in a boundary column, and the column homotopy $h$ maps each column to itself, so the descent contracts a cycle of $\ker\cmpI\cap\Isob$ inside $\Isob$ --- whence
$\tau_{\Iso}(\Isob)\subseteq\Isob$ by construction. Termination (finitely many $\Ga''_g$-orbits of basis elements, $p+1$ cosets; finitely many columns in each degree) and integrality (the inputs
$\taub$, $\tau^{s}$ are integral by Lemma~\ref{lem:taubdryexists}(2) and Proposition~\ref{prop:strictpair}, $h$ is $\Z$-linear, (T1) multiplies by elements of $\Ga$) hold with no
restriction on $p$: \textbf{no step of the proof was verified only at $p\in\{2,3,5\}$}; what was, before v1.23.0-pre.28, is the \emph{implementation} (Remark~\ref{rmk:strictliftbuilt}), and
at a new $p$ the run-time certificate of (1) stands in for it. \emph{The certificate suffices; the component map is forced} (v1.23.0-pre.29, referee m1: in pre.28 only a parenthesis of the
statement, and the proof read the component map off the strict identity, i.e.\ off (B-$\varepsilon$), which the statement says is not needed). Let $\tau_{\Iso}$ satisfy (T1)--(T3) over $\Z$
and $\tau_{\Iso}(\Isob)\subseteq\Isob$, and nothing else. (i) By (T2) and $\tau_{\Iso}(\Isob)\subseteq\Isob$, $\tau^{\partial}_{\Iso}$ is a chain map on $\Isob$; by (T1) it is twisted-equivariant along
$\iota_g$; so it induces a map $\phi$ on $H_0(\Isob)=\Z[\mathbb{P}^1]$, twisted along $\iota_g$. (ii) Let $c$ be a cusp, $P_c$ a generator of its stabiliser in $\Ga$ modulo $\pm I$, and $k\ge1$ with
$P_c^{\,k}\in\Ga_0(p)$. By (i), $\phi([c])$ is a finite formal sum of cusps fixed by $P:=g^{-1}P_c^{\,k}g\in\Ga$; $P$ is parabolic, fixes the single cusp $g^{-1}c$ and has infinite orbits on every
other point of $\mathbb{P}^1(\Q)$, so a finitely supported $P$-fixed sum is $a\,[g^{-1}c]$, $a\in\Z$. (iii) Definition~\ref{def:descentpair} makes the augmentation of $\Isob$ (onto $\Z[\mathbb{P}^1]$,
followed by $\Z[\mathbb{P}^1]\to\Z$) the restriction of that of $\Iso$; as $\tau_{\Iso}$ extends $\tau^{\partial}_{\Iso}$ and preserves the augmentation by (T3), $a=\varepsilon(\phi[c])=\varepsilon([c])=1$. So
$\phi(c)=g^{-1}c$ for every lift with these properties; no step uses (S) or (B-$\varepsilon$), nor $p\in\{2,3,5\}$. \emph{Second proof}, under the strict identities: $\cmpI\tau^{\partial}_{\Iso}=\taub\cmpI$ and
Lemma~\ref{lem:bdrycomponent} for $\taub$. Homotopy of pairs: both $\tau^{\partial}_{\Iso}$ induce on $H_0(\Isob)=\Z[\mathbb{P}^1]$ the component map $c\mapsto g^{-1}c$ (by (i)--(iii)); $\Iso$ is free over $\Z[\Ga''_g]$ on a basis restricting to one of $\Isob$, and the twisted targets $\Iso^{(c)}$, ${\Isob}^{(c)}$ are $\Ga''_g$-resolutions of $\Z$, $\Z[\mathbb{P}^1]^{(c)}$; the proof of
Lemma~\ref{lem:pairlifting}(2) over $\Z[\Ga''_g]$, with (T1) as $\Z[\Ga''_g]$-linearity into the twisted complexes. (2) Integrality, cochain maps and the $N=1$ case are Proposition~\ref{prop:strictliftN}(3)(a),
(b), (d), whose proofs read of the complex only (T1), (T2) and $\partial(s_jx)=s_j\partial x$, true on $\Isob$. The strict identity: for $x\in\Isob$, $(\widetilde T^{(N)}_{\Iso,p}\phi)(x)_i=\sum_t\rhon(\delta_{i,t})
\phi(\tau_{\Iso}(s_jx))_a$ and $\tau_{\Iso}(s_jx)=\tau^{\partial}_{\Iso}(s_jx)\in\Isob$, so it equals $(\widetilde T^{(N)}_{\Isob,p}\rres\phi)(x)_i$ --- Lemma~\ref{lem:cosettransport}(3), (4) blockwise. The cone:
Lemma~\ref{lem:conesign}'s computation with $\hder=0$. (3) The bulk: Proposition~\ref{prop:strictliftN}(3)(c). The cone: $\Phi^{(N)}_p$ preserves $(\ker\rres)[1]$ and $c\mapsto(c,0)$ is a quasi-isomorphism
(Proposition~\ref{prop:descentpair}(1)(c)), so the operator on $\Hcpt{\bullet}$ is that of the coset transport of $\tau^{\mathrm{rel}}$, the map induced by $\tau_{\Iso}$ on $\Iso/\Isob$, on
$\Hom(\Iso/\Isob,M_N)$. $\Iso/\Isob$ is a complex of free $\Z[\Ga]$-modules, free over every subgroup of finite index, with homology $\mathrm{St}$ in degree $1$, and $\tau^{\mathrm{rel}}$ is a chain map along
$\iota_g$ (by (T1), (T3)) lifting $g^{-1}\colon\mathrm{St}\to\mathrm{St}$ (the component map on $H_0(\Isob)$, restricted to the kernel of the augmentation); unique up to homotopy by the fundamental lemma.
Proposition~\ref{prop:strictliftN}(3)(c)'s proof applies verbatim with $\mathrm{St}$ in place of $\Z$ --- the inclusion computes $\res$, $\psi\mapsto\rhon(g)\psi\circ\tau^{\mathrm{rel}}$ computes $(\iota_g,f)^{\ast}$
on $\Ext(\mathrm{St},-)$, the transfer computes $\cor$ \STD{\cite[Ch.~III.9--10]{brown}} --- and $\Ext_{\Z[A]}(\mathrm{St},\VR)=H^{\bullet}(A;\Hom_\Z(\mathrm{St},\VR))$ \STD{\cite[Ch.~III, Prop.~2.2]{brown}}
carries these maps to those of the compatible pair $(\Hom_\Z(\mathrm{St},\VR),f)$; Shapiro for the level. The boundary is the same argument on $\Isob$ with $\Z[\mathbb{P}^1]$. Independence: two pair lifts differ by a
homotopy of pairs $k$ through (T1)-maps ((1)); its coset transport gives the cone homotopy $\Upsilon_k$ of Lemma~\ref{lem:pairlifting}(3) (with Lemma~\ref{lem:cosettransport}(2)), and the transversal by
Proposition~\ref{prop:levelconehecke}(4) on both summands. Equivariance of \eqref{eq:descentpairles}: $\Phi^{(N)}_p$ commutes with $\coneinc$, $\coneproj$, and the identities hold over every $\cring$.
Commutativity: the three terms are $H^{\bullet}(A;X)$ for $X\in\{\VR,\Hom_\Z(\Z[\mathbb{P}^1],\VR),\Hom_\Z(\mathrm{St},\VR)\}$, Hecke modules over the same Hecke ring, and Proposition~\ref{prop:strictliftN}(8)'s
argument is formal in $X$. (4) Lemma~\ref{lem:pairlifting}(3) and Proposition~\ref{prop:isomodelell}(4), (5). At $N=1$ over $\Z[1/6]$: Theorem~\ref{thm:conewelldef}(b), (c) and Proposition~\ref{prop:coneindep}
say the shipped operator does not depend on $(\tau,\taub,\etapair)$ when $6$ is inverted, and the strict datum $(\tau^{s},\taub,0)$ is one admissible choice; Proposition~\ref{prop:strictlift}(3)'s
identity $\cmpI^{\ast}\Th^{s}_g=\Th_{\Iso,g}\cmpI^{\ast}$, and its boundary analogue from $\cmpI\tau^{\partial}_{\Iso}=\taub\cmpI$, make $\alpha_c$ intertwine on the nose. At level $N$:
Proposition~\ref{prop:strictliftN}(4)'s computation with $\cmpI\tau_{\Iso}=\tau^{s}\cmpI$ and $\cmpI\tau^{\partial}_{\Iso}=\taub\cmpI$. (5) collects (1)--(4) and Proposition~\ref{prop:strictliftN}; the bulk-only lift:
Remark~\ref{rmk:conestrict}. (6) $u_{\mathsf P}$ is $\Hom(\mathsf P,u_S)$, which commutes with restriction to $\mathsf P^{\partial}$ and with precomposition by any comparison; the cone identity is immediate. The
square with $\alpha$: Proposition~\ref{prop:strictliftN}(7)'s proof (Proposition~\ref{prop:levelconehecke}(1)'s reciprocity computation reads nothing of the complex) on $\Isob$ with $\tau^{\partial}_{\Iso}$.
$\cor$: Proposition~\ref{prop:descentmaps}(3)(a). (7) Proposition~\ref{prop:descentN}(5) is stated for any compatible pair $(X,f)$; by (3) and Proposition~\ref{prop:descentpair}(5)(d) the cup product on
$\Hcpt{\bullet}$ is $\res^{\Ga}_{A}e\cup-$ on $H^{\bullet-1}(A;\Hom_\Z(\mathrm{St},\VR))$.
\end{proof}

\begin{remark}[What the plan stated, and what holds]\label{rmk:descentpairscope}
\PROVED{} (v1.23.0-pre.26.) The plan's D4 and its rows are right in substance; six points are corrected or sharpened. (i) \textbf{C-D7 is a clause, not an amendment.} The plan's row
``(amend \ref{rmk:integralscope}(4), (5))'' undersells it: the pair needs Definition~\ref{def:descentpair}, Lemma~\ref{lem:pairlifting} and Propositions~\ref{prop:descentpair},
\ref{prop:descentpairhecke}; the amendment of Remark~\ref{rmk:integralscope}(4), (5) is made as well. (ii) \textbf{D-$\delta$6, sharpened.} The boundary is unbounded \emph{only} at $2$: at $3$ and at
$\ell\ge5$ it vanishes above degree $1$ (Proposition~\ref{prop:descentpair}(4)); the cone is, in general, unbounded at $2$ ((5)(c)), and \textbf{can be} unbounded at $3$ as well, \textbf{exactly when}
$H^{q}(A;\Vcoef)\otimes\Z_{(3)}\ne0$ for some $q\ge2$ ($\Hcpt{q}\otimes\Z_{(3)}\cong H^{q}\otimes\Z_{(3)}$ for $q\ge3$, (5)(c), and the bulk is $2$-periodic for $q\ge2$): e.g.\ level $1$, $n\in\{0,4,6,10,12\}$;
not for $n$ odd, nor for $n\in\{2,8\}$ at level $1$ (v1.23.0-pre.27, referee m2; pre.26 said ``unbounded at $2$ and at $3$''). Both are periodic from degree $3$ with closed-form boundary
terms, so the truncation loses nothing at $q_{\max}\ge4$ ((5)(d)) --- the pair needs one degree more than the bulk. ``With a certified exponent'': every exponent is bounded a priori ((6)),
that of $\mathrm{tors}\,\Hcpt{2}$ at $\ell\ge5$ only (by $\mathrm{tors}\,H_0(\Ga;M)$, (5)(e)), not at $\ell\in\{2,3\}$; its differential $\Dcone^{0}$ is certified by the Betti recursion \textbf{or} by a
prime $\ell'\ge5$ at which $H_0(\Ga;M)$ has no torsion --- for $M_N$ every $\ell'\ge5$, $\ell'>n+1$, as in the bulk (v1.23.0-pre.27, referee M2; pre.26 said ``not by a prime''). (iii) \textbf{``Built for
$p\in\{2,3,5\}$''} is right for \emph{built}; the pair lift \emph{exists and is produced} for every $p$ (Proposition~\ref{prop:descentpairhecke}(1)), and what is \emph{asserted} at other $p$ is
Proposition~\ref{prop:descentpairhecke}(5): existence and determinacy for every $p$; computed only at the $p$ with a built pair lift, and at $\ell\ge5$, level $1$, where the shipped cone
operator is built (v1.23.0-pre.27, referee M3). \emph{Since v1.23.0-pre.28} (the owner's scope addition of 2026-10-25) the pair lift is built on demand for every prime $p\nmid N$ and certified at
run time (Proposition~\ref{prop:descentpairhecke}(1)), so $T_p$ on the boundary and the cone is computed, at every $\ell$, level $1$ and level $N$, at every such $p$ whose lift passes its
certificate within the cap; refused only on the eleven conditions of Proposition~\ref{prop:descentpairhecke}(5) (v1.23.0-pre.29, referee M1; pre.28 listed four). The plan's arm (S),
$\tau_{\Iso}(\Isob)\subseteq\Isob$, is Remark~\ref{rmk:strictliftbuilt}(2)'s existing certificate. (iv) \textbf{Arm (Q)}, ``the $\Z[1/6]$ part equals the shipped E8 4b pair after $\otimes\Z/\ell^{K}$'', holds
at level~$1$ as an identity of complexes (Proposition~\ref{prop:descentpair}(1)(d)); at level~$N$ no shipped pair is the $\Z[1/6]$-piece (D1's $\mathsf O$-piece at level $N$ is the sparse cokernel route, which carries no cone),
and the Bredon pair with fibre $M_N$ is a new object of D4 (the code map). (v) \textbf{Coherence.} ``The cone is where coherence is real'' (the plan's \S2.2) is right for the
\emph{Hecke operator} on the cone and wrong for its groups and maps, which glue with no datum (Proposition~\ref{prop:descentpair}(2)); on the free model the strict datum makes the operator
choice-free over $\Z$ (Proposition~\ref{prop:descentpairhecke}(3)). (vi) \textbf{Level $N$ needs no new geometry}: the boundary is block-diagonal by the cusps of $\Ga_0(N)$, read from the
$T$-orbits on $\mathbb{P}^1(\Z/N)$, and its groups depend only on $n$ and the widths (Proposition~\ref{prop:descentpair}(3), (4)) --- the same statement the plan's \S5.5 makes for the tree.
\end{remark}

\noindent\emph{A check outside the build, not a tag} (v1.23.0-pre.26): \file{docs/implementation/slot24/cd7\_pair\_check.py} (\code{python3 -I}, pure integers, no computer algebra system), output
\file{cd7\_pair\_check.out}. It reads no line of \code{cohomology\_app}, imports nothing from the tree and builds none of $\Iso$, $\Isob$, $\Iso^{(\ell)}$: its pair model is the free tree resolution
of $\Z/4*_{\Z/2}\Z/6$ (ranks $2,3,3,\dots$) with the minimal resolution of $\Gbdry$ (ranks $1,2,2,\dots$) and a comparison $\varphi$ \emph{computed} by solving $\partial y=\varphi(\partial x)$ over $\Z$ on balls of
$\Ga$ (Hermite forms) and certified ($\partial\varphi=\varphi\partial$ through degree $8$). A second complex, Lemma~\ref{lem:manin}'s sequence resolved column by column, computes $\Ext^{\bullet}(\mathrm{St},M)$; a
third, the tree Bredon complex with the axis of $T$, is the $\Z[1/6]$-piece. On $24$ cells (level $1$, $n\in\{0,1,2,3,4,6,8,10,12\}$; $\Ga_0(N)$, $N\in\{2,\dots,9,11,13\}$, $n=0$; $(N,n)\in\{(2,1),(2,2),(2,3),
(2,6),(3,1),(3,2),(4,2),(5,2),(7,2)\}$), degrees $\le6$, it checks: the boundary groups against the closed form read from the cusp widths alone, $\abs{F_{n,w}}=n!\,w^{n}$ and $2H^{q}(\Gbdry)=0$ (arm B);
$\Hcpt{q}=\Ext^{q-1}(\mathrm{St},M)$ on the two independent complexes (X); exactness of \eqref{eq:descentpairles} over $\Z$ at every spot as equalities of lattices (E); $\Hcpt{0}=0$, $\Hcpt{1}$ torsion-free,
$\rank\Hcpt{1}=b_1$, $\rank\Hcpt{2}=b_0$, $\rank\res^{1}=b_1(\partial)-b_0$, $48\Hcpt{q}=0$, the periodicities and the $3$- and $\ell\ge5$-local statements of (5)(c) (C); the certified local reads at $2$ and $3$
with the sources of (6) and the exponent caps (L); the Bredon pair against the free one at every $\ell\ge5$ dividing a term, and the $3$-local boundary as the Bredon one (Q); model independence
under other comparisons $T\varphi$, $(-I)\varphi$ (I). Planted faults, each detected: the circle boundary at $\ell=2$; the Bredon bulk at $\ell=3$ for $\Hcpt{q}$, $q\ge3$; a perturbed entry of $\rres$; the
truncated top degree reported; $K$ too small (refused); a prime dividing $\mathrm{tors}\,\Hcpt{2}$ as rank source for $\Dcone^{0}$; the closed form without the widths; the bulk-only rank for the cone.
Arm R (added at v1.23.0-pre.27, referee M2): $\mathrm{tors}\,\Hcpt{2}\otimes\Z_{(\ell)}=H_0(\Ga;M)\otimes\Z_{(\ell)}$ for every $\ell\ge5$ with $H_0$ the cokernel of $[S-1\mid T-1]$, the
free ranks equal, every such $\ell\le n+1$, and $\Dcone^{0}$ modulo $\ell'\in\{5,\dots,31\}$ of rational rank iff $\ell'\nmid\mathrm{tors}\,H_0$ (Proposition~\ref{prop:descentpair}(5)(e), (6)(c)).
\code{CD7\_CHECK checks=2528 \dots\ R=308 \dots\ failing=0 verdict=OK} at pre.27 ($2220$ at pre.26, arm R new). \emph{Corrected at pre.27} (referee M1): pre.26 said here that on these cells
$\Hcpt{2}$ has no torsion at $\ell\ge5$; the script's own output refutes it --- $\Hcpt{2}=\Z/4\oplus\Z/60$ at $(N,n)=(1,8)$ and $\Z/2\oplus\Z/12\oplus\Z/1260$ at $(1,12)$. By
Proposition~\ref{prop:descentpair}(5)(e) the $\ell\ge5$ torsion of $\Hcpt{2}$ is that of $H_0(A;\Vcoef)$, zero for $\ell>n+1$: on these cells $\Z/5$ at $(1,8)$, $\Z/5\oplus\Z/7$ at $(1,12)$, none
elsewhere. It is not an Eisenstein-congruence phenomenon; $691$ is absent at $n=10$ because $691>n+1$. The script checks no Hecke operator; (S) is the code's certificate. The referee's
independent check, \file{docs/implementation/slot24/cd7referee/cd7referee\_check.py} (a $\Ga$-CW pair, the mapping cylinder of the horocycle lines into the Bass--Serre tree; $32$ cells, Hecke
on $\Hcpt{1}$), rerun at pre.27: \code{CD7R\_CHECK checks=2457 \dots\ failing=0 verdict=OK}.

\noindent\emph{Code map} (v1.23.0-pre.26, amended at v1.23.0-pre.27 in (c) and (f); \textbf{stated before D4's code}, contract first, rule~21; (f) restated and the build recorded at
v1.23.0-pre.28; (f), (g) and the validation's controls updated at v1.23.0-pre.29). \textbf{Built} (v1.23.0-pre.28) on branch \texttt{s24-d4impl} (\texttt{d82c7761}); the D4 code
audit's fixes (\file{docs/implementation/slot24/D4\_code\_audit.md}, \texttt{46fc0c73}) and the code side of the pre.28 referee's report (\file{C-D7p\_referee.md}, \texttt{6daf1e03})
applied there, and \textbf{merged} into \texttt{hecke-extensions} at \texttt{3991a2f1} (v1.23.0-pre.29):
\file{Integral/Descent/DescentPair.\{h,cpp\}} (the pieces --- \code{Iso2}, \code{Iso3}, the shipped E8 pair at level $1$, the Bredon pair at level $N$, the free \code{IsoFree} arm ---, the cone, its
certification, the closed-form boundary \code{BoundaryClosedForm} from \code{CuspWidths}, the long exact sequence by D2's \code{ExactAt}), \file{Integral/Descent/DescentPairHecke.cpp}
(\code{BuildConeHecke}: $T_p$ on the boundary and the cone from the pair lift, built on demand, certified and cached per $p$; the refusal codes \code{PairRefusal}, printed
\code{REFUSED [<code>]}; \code{PredictPairLiftRss}), \file{Integral/Model/StrictLift.\{h,cpp\}} (the validation hook \code{StrictLift::with\_homotopy}), sources in
\file{cmake/slot24/DESCENT4.cmake}, registered in \file{Registry/Slot24/DESCENT4.cpp}: \code{IntegralDescentPairValidation}
(\code{standard}) and \code{IntegralDescentPairSlowValidation} (\code{slow}); report \file{docs/implementation/slot24/DESCENT4\_report.md}. \textbf{Obligations of the code} (plan D4,
\code{IntegralDescentPairValidation}): (a) the pair on the code's $\Iso$: $\Isob(M_N)$ as the block selection of the boundary summands (no lift: $\rres$ is a restriction), the cone of
\eqref{eq:conediff} over $\Z/\ell^{K}$ and over $\F_{\ell'}$, built to degree $q_{\max}+1\ge5$; on $\Iso^{(\ell)}$ the boundary sub-sum of the \emph{same} $\mathsf O_\ell$-model, at $\ell=3$ checked to be the circle model
(Proposition~\ref{prop:descentpair}(1)(d)), at $\ell=2$ free; (b) the $\Z[1/6]$-piece: level $1$, the shipped E8 pair (arm (Q)); level $N$, the Bredon pair with fibre $M_N$ (dense $\Ceq(M_N)$ in degrees $\le2$ and the circle
model), a new object --- the sparse cokernel route carries no cone; (c) certification by Proposition~\ref{prop:descentpair}(6): $\Dcone^{-1}$ by $\rank\mathsf P(M)^{0}$; $\Dcone^{0}$ by the Betti recursion with
$\dim\Hcpt{1}=b_1$ from a certified source, \textbf{or} by the same cone modulo a prime $\ell'\ge5$, $\ell'>n+1$ (for $M_N$; for $M_N(\chi)$, $\chi\ne1$, only at an $\ell'\ge5$ checked not to divide the
computed $\mathrm{tors}\,H_0(\Ga;M)$; never at $\ell'\in\{2,3\}$; v1.23.0-pre.27, referee M2); $\Dcone^{q}$, $q\ge1$, by the same cone modulo $\ell'\ge5$; the boundary by its closed form; the a-priori
exponents as consistency checks, at $\ell\ge5$ on $\Dcone^{0}$ the $\ell$-part of $\mathrm{tors}\,H_0(\Ga;M)$; refusal at the cap;
(d) arm (B): the boundary against the closed form from the cusp widths on $\mathbb{P}^1(\Z/N)$, including $\abs{\mathrm{tors}\,H^{1}(\Gbdry;M_N)}=\prod_cn!\,w_c^{n}$; (e) arm (E): exactness of
\eqref{eq:descentpairles} at every reported spot, by Proposition~\ref{prop:descentmaps}(2)(g)'s precision rule, on each piece; (f) Hecke on the boundary and the cone (v1.23.0-pre.28, the owner's scope addition of 2026-10-25, replacing pre.27's ``\dots\ for
$p\in\{2,3,5\}$ \dots; not a D4 obligation: the level-$N$ boundary and cone operator at $p\notin\{2,3,5\}$ \dots\ is refused''): for \textbf{every prime} $p\nmid N$, level $1$ and level $N$, the pair
lift of \file{Integral/Model/StrictLift} is built on the first request (boundary step first, degrees $\le6$), certified by Remark~\ref{rmk:strictliftbuilt}(2)'s checks with the boundary items
before first use, and cached per $p$; $\widetilde T^{(N)}_{\Iso,p}$ and $\widetilde T^{(N)}_{\Isob,p}$ by \eqref{eq:levelcorr} with $\tau_{\Iso}$, $\tau^{\partial}_{\Iso}$; the strict identity
$\rres\widetilde T_{\Iso}=\widetilde T_{\Isob}\rres$ checked entrywise in every degree used, on every call, $\widetilde T_{\Isob}$ assembled separately on $\Isob$'s own layout from
$\tau^{\partial}_{\Iso}$; \textbf{refused only} on the eleven conditions of Proposition~\ref{prop:descentpairhecke}(5), each with its \code{PairRefusal} code (v1.23.0-pre.29, referee M1,
replacing pre.28's four): \code{NotPrime}, \code{PDividesN}, \code{BulkOnly} (never the bulk-only lift on the cone), \code{DegreeRange} ($q_{\mathrm{top}}>6$ in \code{BuildConeHecke};
$q_{\max}>5$ in \code{AssembleDescentPair}, where $q_{\max}<4$ is a programming error and throws), \code{RssCeiling} (\code{PredictPairLiftRss}$(p)>4$\,GiB, before building),
\code{IntegerCap}, \code{DomainCap}, \code{InputTableCertificate} (the cap: \code{ClassifyPairException} maps \code{hecke::HeckeOverflow} to \code{IntegerCap},
\code{std::overflow\_error} to \code{DomainCap}, a \code{HeckeTauTable} ``certification FAILED'' to \code{InputTableCertificate}, while the lift is built \emph{and} while the operator is
assembled, referee m3), \code{LiftCertificate}, \code{IsobClosure}, \code{StrictIdentity}; every other exception (an \code{ENGINE\_CHECK}, \code{solve\_regime()}${}\ne{}$\code{"Z"}, an
internal error) \textbf{propagates} as a failure and leaves the per-$p$ slot empty. \emph{Known width of \code{IntegerCap}} (v1.23.0-pre.29, \textbf{resolved in the code at v1.23.0-pre.30}, commit \texttt{bc128cb7}, merged at \texttt{5eb7e9de}): at pre.29 \code{hecke::HeckeOverflow} was also
thrown by internal guards of the exact integer geometry on the lift's path, which were therefore refused as \code{IntegerCap}. The new type \code{hecke::HeckeInternalError}
(\file{Hecke/ExactInt.h}, a subclass of \code{HeckeOverflow}, so every existing handler is unchanged) is now thrown at the non-cap sites (\code{CuspKey}: $(0:0)$, $\det\ne1$, singular, \code{reduce\_to\_infinity}'s termination, primitivity and internal checks;
\code{word\_in\_T\_S}'s determinant and internal error; \code{evaluate\_word}; \code{boundary\_line\_nearest}'s denominator; \code{QuadPoint}'s domain checks; the $\tau$-table reduction's termination; the integer solvers' termination and
self-verification), and \code{ClassifyPairException} returns ``no refusal'' for it: a genuine internal error \textbf{propagates} as a failure; \textbf{only a real overflow} (\code{HeckeOverflow} proper) is refused as \code{IntegerCap}. The classification check of \code{IntegralDescentPairValidation} covers the new type. The cellular route of
Proposition~\ref{prop:descentpairhecke}(4), (5) at $\ell\ge5$ is not built (not needed: the pair lift serves every $\ell$); at level $1$, $\ell\ge5$, the shipped cone operator is read only where
built (not read by D4); (g) $4\le q_{\max}\le5$ for the pair (the lift's degree $6$; v1.23.0-pre.29, referee m2); (h) DESCENT4's report prices the pair on $\Isob$, $\Iso^{(2)}$ (the third
input to gate G-T); (i) \emph{the negative controls of (f)} in \code{IntegralDescentPairValidation} (v1.23.0-pre.29, referee m4 and the D4 code audit): a perturbed lift value
(breaks (T2)) refused as \code{LiftCertificate}; \textbf{homotopy faults} $\tau'=\tau_{\Iso}+dK+Kd$, $K(s_1e_v)=e_e$ extended by (T1) (\code{StrictLift::with\_homotopy}; $v$ the first
$t=0$ boundary $0$-cell, $e$ a $t=0$ $1$-cell in $\Isob$ or in the bulk), which keep (T1)--(T3) --- asserted --- and are refused as \code{LiftCertificate} by (S) and the boundary
items alone at $p=2,7$ (flags printed: into $\Isob$, (S) and (B-$\varepsilon$) fail; into the bulk, (S), (B)-preserved and (B-$\varepsilon$) fail --- the referee's Arm~F); an off-$\Isob$
term planted with the closure test ignored, refused by the strict identity alone (\code{StrictIdentity}) at $(1,10)$, $p=2,7$, and $(11,0)$, $p=13$; $q_{\mathrm{top}}=7$ and $q_{\max}=6$
(\code{DegreeRange}); $p=100003$ (\code{RssCeiling}, before building); the bulk-only lift (\code{BulkOnly}); and the exception classification on six synthetic exceptions (three
caps refuse; a logic, an internal and a range error propagate).

\paragraph{Formalization note.} Lemma~\ref{lem:pairlifting} is the comparison theorem relative to a subcomplex; Proposition~\ref{prop:descentpair}(4) and (5)(a), (b) are finite algebra and cheap; (1)(c)
and Proposition~\ref{prop:descentpairhecke}(3) need $\Ext$ with Shapiro and the double-coset formalism. No \code{\textbackslash LEAN} marker.

% NEW at v1.23.0-pre.33 (hecke-extensions slot 24, milestone D5c of docs/implementation/slot24/DESCENT_plan.md v0.4, its clause C-D8 and its \S5;
% the owner's ruling [R-GT]).  Placed after sec:descentpair, before sec:mpmodel, so that no statement number moves.  A pre-release draft, NOT
% refereed (an adversarial referee follows).  No \LEAN marker.  The plan's "native tree Hecke datum" (its \S5.4(c)2) is not drafted.
\subsection{The tree cusp pair: the boundary pair and the cone over \texorpdfstring{$\Z$}{Z} on the tree models}\label{sec:treepair}

(v1.23.0-pre.33; milestone D5c of \file{docs/implementation/slot24/DESCENT\_plan.md} v0.4, clause C-D8 of its \S2.7 and its \S5; the owner's ruling at gate G-T.)
\S\ref{sec:descentpair} computes the Borel--Serre pair of $\Ga_0(N)$ over $\Z$ on $(\Iso,\Isob)$ and its $\ell$-local forms; its $2$-adic piece meets a memory wall at
$D:=\rank_\Z M_N\approx1060$--$1240$ (\file{DESCENT4\_report.md} \S4), the dense bulk $\Iso^{(2)}$ inside the cone. \S\ref{sec:isomodeltree} moved the bulk to the tree but
carries no boundary: the $\Ga$-translates of the axis of $T$ cover $\mathcal T$, so the tree has no subcomplex indexed by the cusps (the plan's \S5.1). This subsection
fixes, over every $\mathsf O_\ell$-algebra: (i) a \emph{cusp model} $\mathsf B^{(\ell)}_{\mathcal T}$, a projective resolution of $\cring[\mathbb{P}^1]$ built from two copies of the
tree's edge column, and an explicit chain map $\mathsf a^{(\ell)}\colon\mathsf B^{(\ell)}_{\mathcal T}\to\Iso^{(\ell)}_{\mathcal T}$ over $\cring[\mathbb{P}^1]\xrightarrow{\aug}\cring$ (the \emph{axis map}),
whose cochain cone $\Cone{\mathsf a^{(\ell)\ast}}$ is the computed object (Definition~\ref{def:treepair}); (ii) the mapping cylinder, which makes it a pair model in the sense of
Definition~\ref{def:descentpair}, and the strict deformation of that pair's cone onto $\Cone{\mathsf a^{(\ell)\ast}}$ (Lemma~\ref{lem:treepaircyl}); (iii) what it computes, in every
degree, with the signs of Proposition~\ref{prop:descentpair}(1)(c), its cochains in closed form, the certified exponent and the precision rule, explicit comparisons with
$(\Iso,\Isob)$ made effective as in Proposition~\ref{prop:isomodeltree}(4), and $T_p$ by transport, with the off-diagonal term it acquires on $\Cone{\mathsf a^{(\ell)\ast}}$ (Proposition~\ref{prop:treepair}; v1.23.0-pre.34, referee m1); (iv)
where the plan's \S5 is corrected (Remark~\ref{rmk:treepairscope}). \emph{What is used.} Definitions~\ref{def:descentpair}, \ref{def:isomodeltree}, Lemma~\ref{lem:pairlifting},
Propositions~\ref{prop:descentpair}, \ref{prop:descentpairhecke}, \ref{prop:isomodeltree}, \ref{prop:descentmaps}(2)(g), \ref{prop:torsionhecke}(3), Lemma~\ref{lem:model7coincidences},
Definition~\ref{def:isomodel}, \ref{B1}--\ref{B3}, \eqref{eq:conediff}, Definition~\ref{def:coneaction} (the shape of a cone datum) and Proposition~\ref{prop:coneambiguity} (why the five lemma does not
suffice). \emph{Notation.} As in \S\ref{sec:isomodeltree} and \S\ref{sec:descentpair}; $\cring$ is an $\mathsf O_\ell$-algebra,
$\ell$ a prime or $\ell=0$ for the free models over every ring ($\mathsf O_0:=\Z$); $T=\smat1101$, $S^{-1}=-S$, $U=ST$, so $S^{-1}U=T$; a map written $r_a$ is right multiplication as in
\S\ref{sec:isomodeltree}, and on cochains $\Hom(r_a,M)$ is $m\mapsto am$. $\mathcal L$, $\mathfrak v_\infty$, $\mathfrak e_\infty$, $\varepsilon_{\mathcal L}$, $\mathsf B_{\mathcal T}$,
$\mathsf B^{(\ell)}_{\mathcal T}$, $\mathsf e_t$, $\mathsf z_t$, $\mathsf v^{x}_t$, $\varepsilon_{\mathsf B}$, $\varepsilon^{(\ell)}_{\mathcal T}$, $\varepsilon^{(\ell)}_{\mathsf B}$, $\mathsf a$, $\mathsf a^{(\ell)}$, $\mathsf f_{\mathsf a}$, $\mathsf i_\partial$, $\mathsf r_\partial$, $\mathsf G_\partial$,
$\mathsf g_\square$, $\mathsf g'$, $\mathrm{Cyl}$, $\pi_{\mathrm{Cyl}}$, $\varpi$, $\varsigma$, $\varsigma'$, $\mathsf h_{\mathrm{Cyl}}$, $\Phi^{\Z}_\partial$, $\Psi^{\Z}_\partial$, $\Phi^{\partial}_\ell$, $\Psi^{\partial}_\ell$,
$\Phi_\partial$, $\Psi_\partial$, $\mathsf H^{\mathsf B}_\partial$, $\mathsf H^{X}_\partial$, $\mathsf y_\ell$, $\mathsf y'_\ell$, $\mathsf Y_\Phi$, $\mathsf Y_\Psi$, $\mathsf j_\Phi$, $\mathsf j_\Psi$, $\Theta_\Phi$, $\Theta_\Psi$,
$\mathrm{pr}_\partial$, $T^{\mathrm c}_{\mathcal T}$, $T'_\partial$, $\eta_{\mathcal T}$ are new symbols of this subsection and mean nothing else (rule~\ref{A5}); $\varpi$, $\varsigma$ are not the symbols of
other chapters.

\begin{definition}[The cusp model, the axis map, the tree cone; the geometric tables]\label{def:treepair}
Assume \ref{H4} and \ref{B1}--\ref{B3}. Let $\ell$ be a prime, $\cring$ an $\mathsf O_\ell$-algebra, and $F^{\mathfrak e}_{\bullet}$, $e_{\mathfrak e}$, $e_S$, $e_U$, $F^{x}$, $(a_t,b_t)$ as in
Definition~\ref{def:isomodeltree} (so $F^{\mathfrak e}_t=\cring[\Ga]e_{\mathfrak e}$ for every $t$ at $\ell=2$ and for the free model, for $t=0$ only at $\ell\ge3$; $\varepsilon_t$ the coefficient of its
differential: $(1-\mathbf c)e_{\mathfrak e}$ for odd $t$, $(1+\mathbf c)e_{\mathfrak e}$ for even $t\ge2$).
\begin{enumerate}
\item \textbf{(The cusp model.)} $\mathsf B^{(\ell)}_{\mathcal T}$ is the mapping cone of $r_{(1-T)e_{\mathfrak e}}\colon F^{\mathfrak e}_{\bullet}\to F^{\mathfrak e}_{\bullet}$:
  \[
    \mathsf B^{(\ell)}_{\mathcal T,q}=F^{\mathfrak e}_q\oplus F^{\mathfrak e}_{q-1},\qquad \partial\mathsf e_t=\varepsilon_t\,\mathsf e_{t-1},\qquad \partial\mathsf z_t=(1-T)\,\mathsf e_t-\varepsilon_t\,\mathsf z_{t-1},
  \]
  $\mathsf e_t$ the generator $e_{\mathfrak e}$ of the first summand in degree $t$, $\mathsf z_t$ that of the second in degree $t+1$ (absent where $F^{\mathfrak e}_t=0$; $\varepsilon_0\mathsf z_{-1}:=0$); augmented by
  $\mathsf B^{(\ell)}_{\mathcal T,0}=\cring[\Ga]e_{\mathfrak e}\to\cring[\mathbb{P}^1]$, $ye_{\mathfrak e}\mapsto y\cdot\infty$. Ranks over $\cring[\Ga]e_{\mathfrak e}$: $1,2,2,\dots$ at $\ell=2$ (free) and for the
  free model $\mathsf B_{\mathcal T}$ ($e_{\mathfrak e}=1$, every $\cring$); $1,1,0,\dots$ at $\ell\ge3$.
\item \textbf{(The axis map.)} $\mathsf a^{(\ell)}\colon\mathsf B^{(\ell)}_{\mathcal T}\to\Iso^{(\ell)}_{\mathcal T}$ is $\cring[\Ga]$-linear with
  \[
    \mathsf a^{(\ell)}(\mathsf e_t)=a_t\,\mathsf v^{S}_t,\qquad
    \mathsf a^{(\ell)}(\mathsf z_t)=-[t\text{ even}]\,S^{-1}e_S\,\mathsf v^{S}_{t+1}+[t\text{ even}]\,S^{-1}e_U\,\mathsf v^{U}_{t+1}+(S^{-1}-T)e_{\mathfrak e}\,\mathsf v^{\mathfrak e}_t,
  \]
  $\mathsf v^{x}_t:=e_x$ the generator of $F^{x}_t$ in $\Iso^{(\ell)}_{\mathcal T}$ ($\mathsf v^{\mathfrak e}_t$ in total degree $t+1$) and $a_t$ the $S$-component of the canonical lift, every term dropped
  whose target column vanishes; $\mathsf a:=\mathsf a^{(0)}$ on the free models, with $e_x=1$. At $\ell\ge3$ \textbf{the choice of $S^{-1}$ rather than $S$ is immaterial} (v1.23.0-pre.34,
  referee N8): only $\mathsf z_0$ occurs, every surviving term of $\mathsf a^{(\ell)}(\mathsf z_0)$ carries $e_{\mathfrak e}=\tfrac12(1+\mathbf c)$ on the right ($e_U=e_{\mathfrak e}$ at $\ell=3$; the
  $\mathsf v^{S}_1$-term is dropped, $F^{S}_1=0$), and $S^{-1}=\mathbf cS$ with $\mathbf c$ central gives $Se_{\mathfrak e}=S^{-1}e_{\mathfrak e}$; writing $S$ there yields the same map. At $\ell\in\{0,2\}$ it does not as a chain map, but
  \textbf{on cochains it is immaterial wherever $\mathbf c$ acts trivially} (v1.23.0-pre.35): $S^{-1}$ enters $\mathsf a^{(\ell)}$ only as the left factor of the coefficients $S^{-1}e_x$
  ($x\in\{S,U,\mathfrak e\}$), and $(S-S^{-1})e_x=(1-\mathbf c)Se_x$, so on $\Hom_{\cring[\Ga]}(-,M)$ (where $\Hom(r_b,M)$ is $m\mapsto bm$) the two maps differ by $m\mapsto Se_x(1-\mathbf c)m$. For
  $M_N=\Ind_{\Ga_0(N)}^{\Ga}\Sym^n$ (any $M$ induced from a subgroup containing $-I$), $\mathbf c(g\otimes v)=g\otimes(-I)v$ gives $\rho_{M_N}(\mathbf c)=(-1)^n$: for $n$ \textbf{even} the two axis maps have
  the \textbf{same cochains} at every $\ell$ (every group, map and operator computed from these cochains agrees; the chain-level constructions of Proposition~\ref{prop:treepair}(3) presuppose a chain map
  and use $\mathsf a^{(\ell)}$; v1.23.0-pre.36, the C-D8b referee's N6); for $n$ odd they differ by $2Se_x$, which vanishes modulo $2$.
\item \textbf{(The tree cone.)} For a $\cring[\Ga]$-module $M$ put $\mathsf B^{\bullet}_{\ell}(M):=\Hom_{\cring[\Ga]}(\mathsf B^{(\ell)}_{\mathcal T},M)$ and
  \[
    \Cone{\mathsf a^{(\ell)\ast}}^{q}:=C^{q+1}_{\mathcal T,\ell}(M)\oplus\mathsf B^{q}_{\ell}(M),\qquad\Dcone(\alpha,c):=\bigl(-d\alpha,\ dc+\mathsf a^{(\ell)\ast}\alpha\bigr),
  \]
  \eqref{eq:conediff} with $\mathsf a^{(\ell)\ast}$ for $\rres$, with $\coneinc(c)=(0,c)$, $\coneproj(\alpha,c)=\alpha$, and $\Hcpt{q}(\mathcal T_\ell;M):=H^{q-1}(\Cone{\mathsf a^{(\ell)\ast}})$.
\item \textbf{(The geometric tables.)} $C_{\bullet}(\mathcal L)$ is the complex $\Z[\Ga/\langle\mathbf c\rangle]\mathfrak e_\infty\to\Z[\Ga/\langle\mathbf c\rangle]\mathfrak v_\infty$, $\partial\mathfrak e_\infty=(1-T)\mathfrak v_\infty$, with
  $\varepsilon_{\mathcal L}(\gamma\mathfrak v_\infty)=\gamma\cdot\infty$ (the cellular chains of $\Ga\times_{\Gbdry}\mathbb R$, one vertex and one edge per period); $\varepsilon_{\mathsf B}\colon\mathsf B_{\mathcal T}\to C_{\bullet}(\mathcal L)$ is
  $\mathsf e_0\mapsto\mathfrak v_\infty$, $\mathsf z_0\mapsto\mathfrak e_\infty$, $0$ on $\mathsf e_t,\mathsf z_t$, $t\ge1$. \textbf{The $\ell$-local column augmentations} (v1.23.0-pre.34, referee m4):
  $\varepsilon^{(\ell)}_{\mathcal T}\colon\Iso^{(\ell)}_{\mathcal T}\to C_{\bullet}(\mathcal T)\otimes\cring$ is $ye_x\mapsto y\Ga_x$ ($=y\mathfrak v_x$) on $F^{x}_0$, $x\in\{S,U\}$, and $ye_{\mathfrak e}\mapsto y\langle\mathbf c\rangle$
  ($=y\mathfrak e$) on $F^{\mathfrak e}_0$ (total degree $1$), $0$ on every column degree $t\ge1$ (the vertex part is Definition~\ref{def:isomodeltree}(1)'s augmentation); $\varepsilon^{(\ell)}_{\mathsf B}\colon
  \mathsf B^{(\ell)}_{\mathcal T}\to C_{\bullet}(\mathcal L)\otimes\cring$ is $ye_{\mathfrak e}\mapsto y\mathfrak v_\infty$ on $\mathsf e_0$, $ye_{\mathfrak e}\mapsto y\mathfrak e_\infty$ on $\mathsf z_0$, $0$ on $\mathsf e_t,\mathsf z_t$, $t\ge1$.
  Both are well defined ($Q_x\subseteq\Ga_x$, $Q_{\mathfrak e}\subseteq\langle\mathbf c\rangle$ fixes $\mathfrak v_\infty$, $\mathfrak e_\infty$); $\varepsilon^{(0)}_{\mathcal T}=\varepsilon_{\mathcal T}$ of
  Proposition~\ref{prop:isomodeltree}(4)(b), $\varepsilon^{(0)}_{\mathsf B}=\varepsilon_{\mathsf B}$, and at $\ell=2$ $\varepsilon^{(2)}_{\mathsf B}=\varepsilon_{\mathsf B}\otimes\cring$. With the boundary simplices $v_0,v_1,v_2$, $v_0v_1$, $v_1v_2$, $v_2v_3$ of $\Sigma^{\partial}$
  (Definition~\ref{def:isomodel}; $v_3=Tv_0$), the edges and triangles of $\Sigma_1$, $\Sigma_2$, and $\mathfrak v_S$, $\mathfrak v_U$, $\mathfrak e$ of the tree, define $\Z[\Ga]$-linearly on representatives:
  \begin{center}
  \begin{tabular}{@{}ll@{}}
  $\mathsf i_\partial\colon C_{\bullet}(\mathcal L)\to C_{\bullet}(\Xbdry)$ & $\mathfrak v_\infty\mapsto v_0$, $\mathfrak e_\infty\mapsto-(v_0v_1+v_1v_2+v_2v_3)$ \\
  $\mathsf r_\partial\colon C_{\bullet}(\Xbdry)\to C_{\bullet}(\mathcal L)$ & $v_0,v_1,v_2\mapsto\mathfrak v_\infty$; $v_0v_1,v_1v_2\mapsto0$, $v_2v_3\mapsto-\mathfrak e_\infty$ \\
  $\mathsf G_\partial\colon C_0(\Xbdry)\to C_1(\Xbdry)$ & $v_0\mapsto0$, $v_1\mapsto v_0v_1$, $v_2\mapsto v_0v_1+v_1v_2$ \\
  $\mathsf f_{\mathsf a}\colon C_{\bullet}(\mathcal L)\to C_{\bullet}(\mathcal T)$ & $\mathfrak v_\infty\mapsto\mathfrak v_S$, $\mathfrak e_\infty\mapsto(S^{-1}-T)\mathfrak e$ \\
  $\mathsf g_\square\colon C_0(\Xbdry)\to C_1(\mathcal T)$ & $v_0\mapsto-\mathfrak e$, $v_1,v_2\mapsto0$ \\
  $\mathsf g'\colon C_s(\mathcal L)\to C_{s+1}(\Xbulk)$ & $\mathsf g':=-\bigl(\mathsf G_{\mathcal T}\,j_X\,\mathsf i_\partial+\mathsf i_{\mathcal T}\,\mathsf g_\square\,\mathsf i_\partial\bigr)$: $\mathfrak v_\infty\mapsto v_0v_4+v_4v_6$,\\
   & $\mathfrak e_\infty\mapsto(v_0v_1v_4-v_1v_4v_6)+v_1v_2v_6+(v_2v_3v_5+v_2v_5v_6)$
  \end{tabular}
  \end{center}
  ($j_X$ the inclusion $C_{\bullet}(\Xbdry)\subseteq C_{\bullet}(\Xbulk)$; $\mathsf i_{\mathcal T}$, $\mathsf r_{\mathcal T}$, $\mathsf G_{\mathcal T}$ those of Proposition~\ref{prop:isomodeltree}(4)(a); $\mathsf g_\square$ is $0$ on edges, $\mathsf G_\partial$
  on edges likewise, $C_2(\mathcal T)=C_2(\Xbdry)=0$).
\end{enumerate}
\end{definition}

\begin{lemma}[The cylinder: a pair model, and its cone deformed onto the tree cone]\label{lem:treepaircyl}
\PROVED{} \STD{\cite[\S1.5]{weibel}} Let $\mathsf a\colon\mathsf B\to\mathsf P$ be a chain map from a complex of \textbf{finitely generated} projective $\cring[\Ga]$-modules, \textbf{zero in negative degrees}, resolving $\cring[\mathbb{P}^1]$ (as Definition~\ref{def:descentpair} requires; $\mathsf B^{(\ell)}_{\mathcal T}$ satisfies this, Proposition~\ref{prop:treepair}(1)(a); v1.23.0-pre.34, referee N1) to an
$\cring$-model $\mathsf P$ (Proposition~\ref{prop:descentmaps}), over $\aug\colon\cring[\mathbb{P}^1]\to\cring$. Put $\mathrm{Cyl}_q:=\mathsf P_q\oplus\mathsf B_q\oplus\mathsf B_{q-1}$,
$\partial(x,b,u):=(\partial x-\mathsf au,\ \partial b+u,\ -\partial u)$, with $\mathsf B$ the middle summand, and $\pi_{\mathrm{Cyl}}(x,b,u):=x+\mathsf ab$, augmented by $\aug\circ\pi_{\mathrm{Cyl}}$.
\begin{enumerate}
\item $(\mathrm{Cyl},\mathsf B)$ is an $\cring$-pair model (Definition~\ref{def:descentpair}); $\pi_{\mathrm{Cyl}}$ and the inclusion $i_{\mathsf P}$ of $\mathsf P$ are mutually inverse comparisons over $\id_\cring$
  ($\pi_{\mathrm{Cyl}}i_{\mathsf P}=\id$); on cochains $\rres_{\mathrm{Cyl}}\circ\pi_{\mathrm{Cyl}}^{\ast}=\mathsf a^{\ast}$. Its cochains are $\mathrm{Cyl}(M)^{q}=\mathsf P(M)^{q}\oplus\mathsf B(M)^{q}\oplus\mathsf B(M)^{q-1}$ with
  $d(\alpha,\beta,\gamma)=(d\alpha,\ d\beta,\ -\mathsf a^{\ast}\alpha+\beta-d\gamma)$, and $\rres_{\mathrm{Cyl}}(\alpha,\beta,\gamma)=\beta$.
\item On $\Cone{\rres_{\mathrm{Cyl}}}^{q}=\mathrm{Cyl}(M)^{q+1}\oplus\mathsf B(M)^{q}$ put $\varpi(\alpha,\beta,\gamma,c):=(\alpha,\ c+\gamma)$, $\varsigma(\alpha,c):=(\alpha,0,c,0)$, $\varsigma'(\alpha,c):=(\alpha,\mathsf a^{\ast}\alpha,0,c)$,
  $\mathsf h_{\mathrm{Cyl}}(\alpha,\beta,\gamma,c):=(0,c,0,0)$. Then $\varpi\colon\Cone{\rres_{\mathrm{Cyl}}}\to\Cone{\mathsf a^{\ast}}$, $\varsigma$, $\varsigma'$ are cochain maps,
  $\varpi\varsigma=\varpi\varsigma'=\id$, $\id-\varsigma\varpi=\Dcone\mathsf h_{\mathrm{Cyl}}+\mathsf h_{\mathrm{Cyl}}\Dcone$, and $\varsigma'-\varsigma=\Dcone(\mathsf h_{\mathrm{Cyl}}\varsigma')+(\mathsf h_{\mathrm{Cyl}}\varsigma')\Dcone$.
  \textbf{Strictly}, $\varpi\circ\coneinc=\coneinc$ and $\coneproj\circ\varpi=i_{\mathsf P}^{\ast}\circ\coneproj$.
\item Hence Proposition~\ref{prop:descentpair}(1)(a)--(c) hold for $\Cone{\mathsf a^{\ast}}$: $H^{q}(\mathsf B(M))=H^{q}(\Gbdry;M)$, $H^{q}(\mathsf a^{\ast})=\res^{\Ga}_{\Gbdry}$,
  $H^{q-1}(\Cone{\mathsf a^{\ast}})=\Hcpt{q}\cong\Ext^{q-1}_{\cring[\Ga]}(\mathrm{St}_\cring,M)$, and its long exact sequence is \eqref{eq:descentpairles}, exact in every degree, with
  $\coneinc{}_\ast=j_{\mathrm{St}}^{\ast}$ and $\coneproj{}_\ast=-\delta$ under the identifications through $H(\varpi)$ \textbf{with the signs of Proposition~\ref{prop:descentpair}(1)(c)}; its
  connecting map $H^{q}(\Ga;M)\to H^{q}(\Gbdry;M)$ is $H(\mathsf a^{\ast})$, with no sign.
\end{enumerate}
\end{lemma}

\begin{proof}
(1) $\partial^{2}=0$: the first component of $\partial^{2}(x,b,u)$ is $\partial(\partial x-\mathsf au)-\mathsf a(-\partial u)=0$, the second $\partial(\partial b+u)-\partial u=0$. The middle summand is a subcomplex
with projective complement $\mathsf P\oplus\mathsf B[-1]$. $\pi_{\mathrm{Cyl}}\partial(x,b,u)=\partial x-\mathsf au+\mathsf a\partial b+\mathsf au=\partial\pi_{\mathrm{Cyl}}(x,b,u)$, and with $s(x,b,u):=(0,0,b)$,
$\partial s+s\partial=\id-i_{\mathsf P}\pi_{\mathrm{Cyl}}$ (both send $(x,b,u)$ to $(-\mathsf ab,b,u)$), so $\mathrm{Cyl}\simeq\mathsf P$ resolves $\cring$; on $\mathsf B_0$ the two augmentations agree because $\mathsf a$
lies over $\aug$. The cochain formulas are $\Hom(-,M)$ of $\partial$, the dual of $(x,b,u)\mapsto\alpha x+\beta b+\gamma u$. (2) With $\Dcone(\alpha,\beta,\gamma,c)=(-d\alpha,-d\beta,\mathsf a^{\ast}\alpha-\beta+d\gamma,
dc+\beta)$: $\varpi\Dcone=(-d\alpha,\ dc+d\gamma+\mathsf a^{\ast}\alpha)=\Dcone\varpi$; $\Dcone\varsigma(\alpha,c)=(-d\alpha,0,\mathsf a^{\ast}\alpha+dc,0)=\varsigma\Dcone(\alpha,c)$;
$\Dcone\varsigma'(\alpha,c)=(-d\alpha,-\mathsf a^{\ast}d\alpha,0,dc+\mathsf a^{\ast}\alpha)=\varsigma'\Dcone(\alpha,c)$; $\Dcone\mathsf h_{\mathrm{Cyl}}+\mathsf h_{\mathrm{Cyl}}\Dcone$ sends $(\alpha,\beta,\gamma,c)$ to
$(0,-dc,-c,c)+(0,dc+\beta,0,0)=(0,\beta,-c,c)=(\id-\varsigma\varpi)(\alpha,\beta,\gamma,c)$; and $\varsigma'-\varsigma=(\id-\varsigma\varpi)\varsigma'$. The strict identities are read off. (3) By (1)
Proposition~\ref{prop:descentpair}(1) applies to $(\mathrm{Cyl},\mathsf B)$; $\varpi$ is a quasi-isomorphism by (2), commuting with $\coneinc$ and, through the comparison $i_{\mathsf P}$ over $\id_\cring$ (which
induces the identity under the identifications with $H^{\bullet}(\Ga;M)$), with $\coneproj$. The connecting map of $0\to\mathsf B(M)\to\Cone{\mathsf a^{\ast}}\to\mathsf P(M)[1]\to0$: a cocycle $\alpha$ lifts to $(\alpha,0)$,
and $\Dcone(\alpha,0)=(0,\mathsf a^{\ast}\alpha)$; $\mathsf a^{\ast}$ is the comparison along $\aug$, which induces $\res$ (the proof of Proposition~\ref{prop:descentpair}(1)(a)).
\end{proof}

\begin{proposition}[The tree cusp pair: what it computes, closed forms, comparisons with $(\Iso,\Isob)$ over $\Z$, Hecke by transport, precision]\label{prop:treepair}
\PROVED{} \STD{\cite[Thm.~2.2.6, \S1.5]{weibel}, \cite[Ch.~III.6, III.8--9]{brown}} (v1.23.0-pre.33; milestone D5c, clause C-D8 of the plan; refereed ACCEPT after fixes, \file{docs/implementation/slot24/C-D8\_referee.md}; the fixes at v1.23.0-pre.34, not re-refereed.) Assume \ref{H4} and
\ref{B1}--\ref{B3}; $\ell$, $\cring$ as in Definition~\ref{def:treepair}, $M$ a $\cring[\Ga]$-module; for the integral statements $M$ as in \S\ref{sec:descentpair} ($\Z$-free of finite rank) and $M_\cring:=M\otimes\cring$.
\begin{enumerate}
\item \textbf{(What it computes, over every ring.)}
  \begin{enumerate}
  \item[(a)] $\mathsf B^{(\ell)}_{\mathcal T}\to\cring[\mathbb{P}^1]$ is a resolution by finitely generated projective $\cring[\Ga]$-modules ($\cring[\Ga]e_{\mathfrak e}\cong\cring[\Ga/Q_{\mathfrak e}]$), free at
    $\ell=2$ and on the free model; $\varepsilon^{(\ell)}_{\mathsf B}$ is a surjective quasi-isomorphism onto $C_{\bullet}(\mathcal L)\otimes\cring$, which resolves $\cring[\mathbb{P}^1]$ through $\varepsilon_{\mathcal L}$.
  \item[(b)] $\mathsf a^{(\ell)}$ is well defined (every coefficient lies in $e_{\mathfrak e}\cring[\Ga]e_y$), a chain map, and lies over $\aug$; on column degrees $0$ it lies over $\mathsf f_{\mathsf a}$
    ($\varepsilon^{(\ell)}_{\mathcal T}\mathsf a^{(\ell)}=\mathsf f_{\mathsf a}\varepsilon^{(\ell)}_{\mathsf B}$, the $\ell$-local column augmentations of Definition~\ref{def:treepair}(4); v1.23.0-pre.34, referee m4).
  \item[(c)] Lemma~\ref{lem:treepaircyl} applies to $\mathsf a^{(\ell)}$: $(\mathrm{Cyl}(\mathsf a^{(\ell)}),\mathsf B^{(\ell)}_{\mathcal T})$ is an $\cring$-pair model, and $\Cone{\mathsf a^{(\ell)\ast}}$ computes, in every degree
    and for every $M$, $H^{\bullet}(\Ga;M)$, $H^{\bullet}(\Gbdry;M)$, $\Hcpt{\bullet}(\Ga;M)$ and \eqref{eq:descentpairles}, with the signs of Proposition~\ref{prop:descentpair}(1)(c). For $M=\Ind_{\Ga_0(N)}^{\Ga}V$
    (Shapiro) these are the groups of $\Ga_0(N)$, and Proposition~\ref{prop:descentpair}(3)--(5) hold verbatim (they concern the groups, not a model).
  \end{enumerate}
\item \textbf{(The cochains in closed form.)} With the components of Proposition~\ref{prop:isomodeltree}(2) on $C^{\bullet}_{\mathcal T,\ell}(M)$ ($x$ on $F^{S}_q$, $y$ on $F^{U}_q$, $z$ on $F^{\mathfrak e}_{q-1}$) and
  $(u,w)$ on $(\mathsf e_q,\mathsf z_{q-1})$:
  \begin{enumerate}
  \item[(a)] $\ell=2$ (and the free model, $y\in M$): $\mathsf B^{0}=M$, $\mathsf B^{q}=M\oplus M$ ($q\ge1$), $d^{0}u=((1-\mathbf c)u,(1-T)u)$,
    $d^{q}(u,w)=(\varepsilon_{q+1}u,\ (1-T)u-\varepsilon_qw)$; and $\mathsf a^{(2)\ast}(x,y)=x$ in degree $0$,
    \[
      \mathsf a^{(2)\ast}(x,y,z)=\bigl(a_q\,x,\ [q\text{ odd}]\,S^{-1}(y-x)+(S^{-1}-T)z\bigr)\qquad(q\ge1;\ a_q=1\text{ for even, }1+S\text{ for odd }q).
    \]
  \item[(b)] $\ell=3$: $\mathsf B^{0}=\mathsf B^{1}=M^{-I}$, $d^{0}=1-T$, $\mathsf B^{q}=0$ ($q\ge2$); $\mathsf a^{(3)\ast}(x,y)=x$ ($x\in M^{\langle S\rangle}\subseteq M^{-I}$), $\mathsf a^{(3)\ast}(y,z)=S^{-1}y+(S^{-1}-T)z$ on
    $C^{1}=M^{-I}\oplus M^{-I}$, and $0$ above.
  \item[(c)] $\ell\ge5$: $\mathsf B$ as at $\ell=3$; $\mathsf a^{\ast}(x,y)=x$, $\mathsf a^{\ast}z=(S^{-1}-T)z$ on $C^{1}=M^{-I}$; the tree cone is $M^{\langle S\rangle}\oplus M^{\langle U\rangle}\to M^{-I}\oplus M^{-I}\to M^{-I}$, and
    $\Hcpt{2}=M^{-I}/\bigl((1-T)M^{-I}+(1-S)M^{-I}\bigr)$, the coinvariants of Proposition~\ref{prop:descentpair}(5)(e).
  \item[(d)] \textbf{(Ranks.)} $\rank\Cone{\mathsf a^{(\ell)\ast}}^{q}=\rank C^{q+1}_{\mathcal T,\ell}+\rank\mathsf B^{q}_\ell$; at $\ell=2$, with $d=\rank M$ and $r:=\rank M^{U^{2}}=\tfrac13(d+2\operatorname{tr}\rho_M(U^{2}))$:
    $d+r$, $3d+r$, $4d+r$, $4d+r,\dots$ in degrees $-1,0,1,2,\dots$ --- with $r\approx d/3$, $\approx\tfrac43d,\tfrac{10}3d,\tfrac{13}3d,\dots$, against $\approx16\tfrac13d,\,24\tfrac13d,\dots$ in degrees $0,1,\dots$ on
    $(\Iso^{(2)},{\Isob}^{(2)})$ (\file{DESCENT\_plan.md} \S5.3).
  \end{enumerate}
\item \textbf{(Comparisons with $(\Iso,\Isob)$ over $\Z$: effective.)}
  \begin{enumerate}
  \item[(a)] \textbf{(The tables.)} $\mathsf i_\partial$, $\mathsf r_\partial$, $\mathsf f_{\mathsf a}$, $\varepsilon_{\mathsf B}$ are chain maps over $\Z[\mathbb{P}^1]$, resp.\ (for $\mathsf f_{\mathsf a}$) over $\aug$; $\mathsf r_\partial\mathsf i_\partial=\id$;
    $\id-\mathsf i_\partial\mathsf r_\partial=\partial\mathsf G_\partial+\mathsf G_\partial\partial$ on $C_{\bullet}(\Xbdry)$; $\mathsf r_{\mathcal T}j_X-\mathsf f_{\mathsf a}\mathsf r_\partial=\partial\mathsf g_\square+\mathsf g_\square\partial$ on $C_{\bullet}(\Xbdry)$;
    $\mathsf i_{\mathcal T}\mathsf f_{\mathsf a}-j_X\mathsf i_\partial=\partial\mathsf g'+\mathsf g'\partial$ on $C_{\bullet}(\mathcal L)$.
  \item[(b)] \textbf{(The boundary comparisons.)} \emph{The solver, generalised} (v1.23.0-pre.34, referee m2). Let $E$ be a complex of $\Z[\Ga]$-modules filtered by finitely many
    columns $E=\bigoplus_{s\le s_{\max}}E_{s,\bullet}$ (total degree $s+t$), each column free and periodic as in Proposition~\ref{prop:isomodeltree}(4)(b) (generators $\Z[\Ga/\Ga_\sigma]$-free, differentials
    $r_{g-1}$, $r_N$ or the tree's $r_{1-g}$, $r_N$), entering with a global sign $\epsilon_s$, the differential lowering $s$ off the columns, and let $p\colon E\to C$ be its \emph{column augmentation} (the augmentations
    of the columns $s$ in column degree $0$, a surjective quasi-isomorphism onto a complex $C$ with $C_s$ the column's $H_0$). Then $\mathsf h_p$ of loc.\ cit.\ --- descending in $s$, $\epsilon_s\mathsf h$
    column by column (sign rule: $\epsilon=-1$ on a cone's shifted column; the tree's columns with even-degree $\mathsf h$ negated), at most $s_{\max}+1$ steps --- satisfies $\partial\mathsf h_p(z)=z$ and
    $p\,\mathsf h_p(z)=0$ on cycles $z$ of $\ker p$, and (L), (L$'$) hold verbatim along $p$; everything is $\Z$-linear, so all of it holds over every ring $\cring$ after $\otimes\cring$, for
    $\cring[\Ga]$-maps $A$, $B$, $f$ with coefficients in $\cring[\Ga]$ (they may carry the denominators of $\mathsf O_\ell$ and the idempotents $e_x$) and $P'$ free over $\cring[\Ga]$ with given bases.
    Loc.\ cit.\ is the case $p\in\{\cmpI,\varepsilon_{\mathcal T}\}$; here it is used for $p=\varepsilon_{\mathsf B}$ (two columns, $\mathsf e$ ($s=0$) and $\mathsf z$ ($s=1$, on top, $\epsilon=-1$), the tree's sign
    convention $r_{1-g}$), for $p=\cmpI|_{\Isob}$ (the columns $s\le1$ over $\Sigma^{\partial}$; each $\epsilon_s\mathsf h$ maps a column into itself, so the descent stays in $\Isob$), and, over $\cring$,
    for $\varepsilon_{\mathcal T}\otimes\cring$ and $\cmpI\otimes\cring$ in (c). Put $\Phi^{\Z}_\partial:=$ (L) for $f=\mathsf r_\partial\circ\cmpI$ along $\varepsilon_{\mathsf B}$, $\Psi^{\Z}_\partial:=$ (L) for
    $f=\mathsf i_\partial\circ\varepsilon_{\mathsf B}$ along $\cmpI|_{\Isob}$: chain maps $\Isob\rightleftarrows\mathsf B_{\mathcal T}$ over $\id_{\Z[\mathbb{P}^1]}$, over $\mathsf r_\partial$, resp.\ $\mathsf i_\partial$; by (L$'$),
    $\mathsf H^{\mathsf B}_\partial\colon\Phi^{\Z}_\partial\Psi^{\Z}_\partial\simeq\id$ ($g=0$) and $\mathsf H^{X}_\partial\colon\id\simeq\Psi^{\Z}_\partial\Phi^{\Z}_\partial$ ($g=\mathsf G_\partial$). $\ell$-locally, $\Phi^{\partial}_\ell\colon\mathsf B_{\mathcal T}\otimes\cring\to
    \mathsf B^{(\ell)}_{\mathcal T}$ is $r_{\abs{Q_{\mathfrak e}}^{\lfloor t/2\rfloor}e_{\mathfrak e}}$ on both columns ($0$ where the $\ell$-local column vanishes) and $\Psi^{\partial}_\ell$ the inclusions with the inverse scalars
    --- the scalars are $1$ in every degree that occurs ($Q_{\mathfrak e}=1$ at $\ell=2$; only $t=0$ at $\ell\ge3$; v1.23.0-pre.34, referee N3) ---; both lie over the identity of $C_{\bullet}(\mathcal L)\otimes\cring$
    ($\varepsilon^{(\ell)}_{\mathsf B}\Phi^{\partial}_\ell=\varepsilon_{\mathsf B}\otimes\cring$, $(\varepsilon_{\mathsf B}\otimes\cring)\Psi^{\partial}_\ell=\varepsilon^{(\ell)}_{\mathsf B}$):
    chain maps over $\id$, $\Phi^{\partial}_\ell\Psi^{\partial}_\ell=\id$ \textbf{strictly}, no correction terms ($r_{(1-T)e}$ commutes with the column maps); at $\ell=2$ both are the identity.
    $\Phi_\partial:=\Phi^{\partial}_\ell\Phi^{\Z}_\partial$, $\Psi_\partial:=\Psi^{\Z}_\partial\Psi^{\partial}_\ell$.
  \item[(c)] \textbf{(The two squares.)} With $\Phi_{\mathcal T}$, $\Psi_{\mathcal T}$, $\Phi^{\mathcal T}_\ell$, $\Psi^{\mathcal T}_\ell$ of Proposition~\ref{prop:isomodeltree}(3), (4) and $j\colon\Isob\subseteq\Iso$: $\mathsf y_\ell:=$
    (L$'$) along $\varepsilon_{\mathcal T}\otimes\cring$ on $\Isob\otimes\cring$ with $A=\Phi^{\Z}_{\mathcal T}j$, $B=\Psi^{\mathcal T}_\ell\mathsf a^{(\ell)}\Phi_\partial$, $p'=\cmpI|_{\Isob}$, $g=\mathsf g_\square$; $\mathsf y'_\ell:=$ (L$'$) along
    $\cmpI\otimes\cring$ on $\mathsf B_{\mathcal T}\otimes\cring$ with $A=\Psi_{\mathcal T}\mathsf a^{(\ell)}\Phi^{\partial}_\ell$, $B=j\Psi^{\Z}_\partial$, $p'=\varepsilon_{\mathsf B}$, $g=\mathsf g'$. Put $\mathsf Y_\Phi:=\Phi^{\mathcal T}_\ell\mathsf y_\ell$,
    $\mathsf Y_\Psi:=\mathsf y'_\ell\Psi^{\partial}_\ell$. ($\Phi^{\mathcal T}_\ell$, $\Psi^{\mathcal T}_\ell$ lie over the identity of $C_{\bullet}(\mathcal T)\otimes\cring$ --- $\varepsilon^{(\ell)}_{\mathcal T}\Phi^{\mathcal T}_\ell=\varepsilon_{\mathcal T}\otimes\cring$,
    $(\varepsilon_{\mathcal T}\otimes\cring)\Psi^{\mathcal T}_\ell=\varepsilon^{(\ell)}_{\mathcal T}$: in column degree $0$ they are $r_{e_x}$, resp.\ the inclusion, $ye_x\Ga_x=y\Ga_x$, and the corrections $k$, $k'$ land in
    column degree $1$ --- and $\Phi^{\partial}_\ell$, $\Psi^{\partial}_\ell$ over that of $C_{\bullet}(\mathcal L)\otimes\cring$, (b); v1.23.0-pre.34, referee m4.) Then
    \[
      \Phi_{\mathcal T}\,j-\mathsf a^{(\ell)}\Phi_\partial=\partial\mathsf Y_\Phi+\mathsf Y_\Phi\partial,\qquad \Psi_{\mathcal T}\,\mathsf a^{(\ell)}-j\,\Psi_\partial=\partial\mathsf Y_\Psi+\mathsf Y_\Psi\partial .
    \]
  \item[(d)] \textbf{(The cone maps.)}
    \[
      \mathsf j_\Phi(\alpha,c):=\bigl(\Phi_{\mathcal T}^{\ast}\alpha,\ \Phi_\partial^{\ast}c-\mathsf Y_\Phi^{\ast}\alpha\bigr),\qquad \mathsf j_\Psi(\alpha',c'):=\bigl(\Psi_{\mathcal T}^{\ast}\alpha',\ \Psi_\partial^{\ast}c'-\mathsf Y_\Psi^{\ast}\alpha'\bigr)
    \]
    are cochain maps $\Cone{\mathsf a^{(\ell)\ast}}\rightleftarrows\Cone{\rres_{\Iso}}(M_\cring)$, and \textbf{strictly} $\mathsf j_\Phi\coneinc=\coneinc\Phi_\partial^{\ast}$, $\coneproj\mathsf j_\Phi=\Phi_{\mathcal T}^{\ast}\coneproj$, likewise for
    $\mathsf j_\Psi$. They are quasi-isomorphisms, \textbf{mutually inverse on cohomology}, $H(\mathsf j_\Phi)$ is the comparison isomorphism of Lemma~\ref{lem:pairlifting}(3) (through $\varpi$ and a comparison of
    pairs, the proof), and $(H(\Phi_{\mathcal T}^{\ast}),H(\Phi_\partial^{\ast}),H(\mathsf j_\Phi))$ is an isomorphism of long exact sequences from the tree's \eqref{eq:descentpairles} to that of $(\Iso,\Isob)$.
    On cohomology none of these depends on the choices (representatives, preimages, $\Iso$, the solvers' order).
  \end{enumerate}
\item \textbf{(Hecke by transport.)} Let $T=(T_{\mathrm b}[1],T_\partial)$ be a cochain endomorphism of $\Cone{\rres_{\Iso}}(M_\cring)$ with $\rres\,T_{\mathrm b}=T_\partial\,\rres$ strictly --- e.g.\ $\Phi^{(N)}_p$ of
  Proposition~\ref{prop:descentpairhecke}(2) for $M=M_N$, $p\nmid N$ --- and put $T^{\mathrm c}_{\mathcal T}:=\mathsf j_\Psi\circ T\circ\mathsf j_\Phi$.
  \begin{enumerate}
  \item[(a)] $T^{\mathrm c}_{\mathcal T}(\alpha,c)=\bigl(T'_{\mathcal T}\alpha,\ T'_\partial c+\eta_{\mathcal T}\alpha\bigr)$ with $T'_{\mathcal T}=\Psi_{\mathcal T}^{\ast}T_{\mathrm b}\Phi_{\mathcal T}^{\ast}$ (the operator of
    Proposition~\ref{prop:isomodeltree}(6)), $T'_\partial:=\Psi_\partial^{\ast}T_\partial\Phi_\partial^{\ast}$ and
    \[
      \eta_{\mathcal T}:=-\Psi_\partial^{\ast}T_\partial\mathsf Y_\Phi^{\ast}-\mathsf Y_\Psi^{\ast}T_{\mathrm b}\Phi_{\mathcal T}^{\ast},\qquad d\,\eta_{\mathcal T}+\eta_{\mathcal T}\,d=T'_\partial\,\mathsf a^{(\ell)\ast}-\mathsf a^{(\ell)\ast}\,T'_{\mathcal T}.
    \]
    $(T'_{\mathcal T},T'_\partial,\eta_{\mathcal T})$ has the shape of Definition~\ref{def:coneaction}'s datum. \textbf{Where $\eta_{\mathcal T}$ comes from} (v1.23.0-pre.34, referee m1): $\Cone{\mathsf a^{(\ell)\ast}}$ is not the
    cone of a restriction to a subcomplex. On the pair model $(\mathrm{Cyl}(\mathsf a^{(\ell)}),\mathsf B^{(\ell)}_{\mathcal T})$ of (1)(c) the transported datum is \textbf{strict}:
    $T^{\mathrm{Cyl}}:=(\Theta_\Psi^{\ast}[1],\Psi_\partial^{\ast})\circ T\circ(\Theta_\Phi^{\ast}[1],\Phi_\partial^{\ast})$ is diagonal on $\Cone{\rres_{\mathrm{Cyl}}}$ ($\Theta_\Phi$, $\Theta_\Psi$ the comparisons of pairs of the proof
    of (3)(d)), and $T^{\mathrm c}_{\mathcal T}\simeq\varpi\circ T^{\mathrm{Cyl}}\circ\varsigma'$; the off-diagonal term enters through $\varpi$ and $\varsigma'$. The diagonal $(T'_{\mathcal T}[1],T'_\partial)$ \textbf{need not} be a
    cochain map and $\eta_{\mathcal T}$ \textbf{need not} vanish --- \emph{observed}, not part of the proved statement: for $T=\theta^{\ast}$ in the check below and for the C-D8 referee's $T_\varepsilon$
    (\file{docs/implementation/slot24/cd8referee/}); for fibre maps $\eta_{\mathcal T}$ may be non-zero while the diagonal is a cochain map (the check below). $\eta_{\mathcal T}$ is \textbf{determined} by $T$ and the
    comparisons; the ambiguity of Proposition~\ref{prop:coneambiguity} does not arise.
  \item[(b)] $T^{\mathrm c}_{\mathcal T}$ is a cochain endomorphism of $\Cone{\mathsf a^{(\ell)\ast}}$, commuting with $\coneinc$ through $T'_\partial$ and with $\coneproj$ through $T'_{\mathcal T}$, and under $H(\mathsf j_\Phi)$ it
    induces $H(T)$ on $\Hcpt{\bullet}$; $T'_\partial$, $T'_{\mathcal T}$ induce $H(T_\partial)$, $H(T_{\mathrm b})$. For $T=\Phi^{(N)}_p$: $T_p$ of Proposition~\ref{prop:descentpairhecke}(3) on the three terms of
    \eqref{eq:descentpairles}, over every $\cring$, in every degree $\le q_{\max}$ for which the pair lift is built; the tree's sequence is $T_p$-equivariant. By (e), the comparisons
    are needed only in chain degrees $\le q_{\max}$ for this.
  \item[(c)] On cohomology $T^{\mathrm c}_{\mathcal T}$ does not depend on the comparisons; $H(T^{\mathrm c}_{1,\mathcal T})H(T^{\mathrm c}_{2,\mathcal T})=H((T_1T_2)^{\mathrm c}_{\mathcal T})$. (No cochain homotopy between
    $T^{\mathrm c}_{1}T^{\mathrm c}_{2}$ and $(T_1T_2)^{\mathrm c}$ is constructed here.)
  \item[(d)] A $\Z[\Ga]$-map $u$ of fibres acts on every term blockwise and commutes \textbf{strictly} with $\mathsf a^{(\ell)\ast}$, $\mathsf j_\Phi$, $\mathsf j_\Psi$; so
    Proposition~\ref{prop:descentpairhecke}(6) (D2's maps on the pair, the $T_p$-equivariant ladder for $\res$) holds on the tree pair verbatim, and the level ladder's signs (Proposition~\ref{prop:descentmaps}(2)(b))
    are unchanged.
  \item[(e)] \textbf{(Truncation: what the read of $\Hcpt{q}$ uses; v1.23.0-pre.35.)} Let $m\ge1$. Suppose given, over $\cring$, only: $\Phi_{\mathcal T}$, $\Psi_{\mathcal T}$, $\Phi_\partial$,
    $\Psi_\partial$ in chain degrees $\le m$, commuting with $\partial$ there and lying over the identities of (3)(b), (c); $\mathsf Y_\Phi$, $\mathsf Y_\Psi$ in chain degrees $\le m-1$ (into chain degrees
    $\le m$) with the two identities of (3)(c) in chain degrees $\le m-1$; and $T$ as in (4), or the truncation of one to cochain degrees $\le m$ --- the maps below use $T_{\mathrm b}$ on
    $C^{\le m}_{\mathrm I}$ and $T_\partial$ on ${\Isob}^{\le m}$ only, so (i)--(iii) hold for it as for the $T$ it truncates; the pair lift of $\Phi^{(N)}_p$, which the code builds to chain degree $6$, is such a truncation,
    Lemma~\ref{lem:pairlifting}'s induction continuing it (v1.23.0-pre.36, the C-D8b referee's N2). Then:
    \begin{enumerate}
    \item[(i)] these comparison data extend to all degrees with every identity of (3)(b), (c); the extension is not unique, and nothing below depends on it;
    \item[(ii)] in cone degree $j$, $\mathsf j_\Phi$, $\mathsf j_\Psi$, $T'_{\mathcal T}$, $T'_\partial$, $\eta_{\mathcal T}$ and $T^{\mathrm c}_{\mathcal T}$ use $\Phi$, $\Psi$ in chain degrees $\le j+1$ and $\mathsf Y$ in chain degree $j$
      only --- more precisely the bulk $\Phi_{\mathcal T}$, $\Psi_{\mathcal T}$ in chain degree $j+1$, and $\Phi_\partial$, $\Psi_\partial$, $\mathsf Y_\Phi$, $\mathsf Y_\Psi$ in chain degree $j$ ---, so in cone degrees $j\le m-1$ they are the
      maps of every extension; likewise $T'_{\mathcal T}$ on $C^{q}_{\mathcal T,\ell}$ and $T'_\partial$ on $\mathsf B^{q}_\ell$ use the comparisons in chain degree $q$ only, so $H^{q}$ and $H^{q}(\partial)$, $q\le m$, with the
      operators of (b) on them, are read from the data as well (v1.23.0-pre.36, the C-D8b referee's N4);
    \item[(iii)] hence, for $j\le m-1$, $T^{\mathrm c}_{\mathcal T}$ maps the cocycles of $\Cone{\mathsf a^{(\ell)\ast}}^{j}$ to cocycles and the coboundaries to coboundaries, and induces on
      $H^{j}(\Cone{\mathsf a^{(\ell)\ast}})=\Hcpt{j+1}$ the operator of (b): \textbf{the read of $\Hcpt{q}$, $q\le m$, needs the comparisons in chain degrees $\le m$ only};
    \item[(iv)] the identity of (a), $d\eta_{\mathcal T}+\eta_{\mathcal T}d=T'_\partial\mathsf a^{(\ell)\ast}-\mathsf a^{(\ell)\ast}T'_{\mathcal T}$ on $C^{j+1}_{\mathcal T,\ell}$, involves $\eta_{\mathcal T}$ in cone degree $j+1$, i.e.\
      $\Phi$, $\Psi$ in chain degree $j+2$: it can be formed from the data iff $j\le m-2$. The read of $\Hcpt{q}$ does not use it: it holds for the maps of (iii) by (a). \emph{As a certificate} of that read
      it bears on cone degree $q-2$, where (given $T'_{\mathcal T}$, $T'_\partial$ cochain maps and the strict block) it is equivalent to $T^{\mathrm c}_{\mathcal T}$ commuting with $\Dcone$ on cone degree $q-2$, which
      gives the coboundary half of (iii) (and is stronger than it; v1.23.0-pre.36, the C-D8b referee's N1); in cone degree $q-1$ the read needs only the cocycle half of (iii), which can be checked directly. At $\ell=3$ it is empty in cone degrees $j\ge1$ ($\mathsf B^{j}_3=0$ for $j\ge2$).
    \end{enumerate}
    With $m=5$ (the comparisons of the code, built to degree $5$): $\Hcpt{q}$ for $q\le5$, so $q_{\max}=5$ is read with no comparison in degree $6$, and the identity of (a) can be formed in cone degrees
    $j\le3$.
  \end{enumerate}
\item \textbf{(The precision rule.)} Over $\Z_{(\ell)}$, $\ell\in\{2,3\}$:
  \begin{enumerate}
  \item[(a)] $\Cone{\mathsf a^{(\ell)\ast}}$ on $\Iso^{(\ell)}_{\mathcal T,\mathsf O_\ell}$, $\mathsf B^{(\ell)}_{\mathcal T,\mathsf O_\ell}$ is a bounded below complex of finitely generated free $\mathsf O_\ell$-modules, one integral form for
    every $\mathsf O_\ell$-algebra; Proposition~\ref{prop:descentpair}(6) holds for it \textbf{verbatim} --- the Smith-valuation read-out, the stopping rule, the a-priori exponents, the truncation ---,
    with the rank sources read on \textbf{this} complex: the Betti recursion with the dimensions of Proposition~\ref{prop:descentpair}(5)(b), or the reduction modulo a qualifying $\ell'\nmid\mathsf q_\ell$ of this tree
    cone ((6)(c), (d) there), of $\mathsf B^{(\ell)}_{\ell}(M)$ for the boundary ((6)(a)); (6)(b)'s ``$\rank\mathsf P(M)^{0}$'' reads $\rank C^{0}_{\mathcal T,\ell}(M)$ here. The rank of a differential of
    $\Cone{\rres_{\Iso}}$, of $\Cone{\rres_{\Iso^{(\ell)}}}$, of the cylinder's $\Cone{\rres_{\mathrm{Cyl}}}$ (its differentials' rational ranks differ from those of $\Cone{\mathsf a^{(\ell)\ast}}$ by the contractible
    summand $\mathsf B\oplus\mathsf B[-1]$; v1.23.0-pre.34, referee N4) or of the bulk tree alone is \textbf{not} admissible for it, nor conversely.
  \item[(b)] The maps of \eqref{eq:descentpairles} are read by Proposition~\ref{prop:descentmaps}(2)(g); the certifying ranks of (g)(iii) for them are invariants of the maps: over $\Q$,
    $\rank\res^{0}=b_0(M)$, $\rank\res^{1}=b_1(\partial;M)-\dim(M^{\vee}_\Q)^{\Ga}$, $\rank\res^{q}=0$ ($q\ge2$), $\rank\coneinc{}^{q}_\ast=b_q(\partial;M)-\rank\res^{q}$,
    $\rank\coneproj{}^{q}_\ast=b_q(M)-\rank\res^{q}$ (Proposition~\ref{prop:descentpair}(5)(b)).
  \item[(c)] Every matrix of $\mathsf a^{(\ell)\ast}$, $\Phi^{\ast}_{\mathcal T}$, $\Psi^{\ast}_{\mathcal T}$, $\Phi^{\ast}_\partial$, $\Psi^{\ast}_\partial$, $\mathsf Y^{\ast}_\Phi$, $\mathsf Y^{\ast}_\Psi$, $\mathsf j_\Phi$, $\mathsf j_\Psi$, $T^{\mathrm c}_{\mathcal T}$ has its
    entries in $\mathsf O_\ell$ and is \textbf{exact} modulo every $\ell^{K}$ (exact data, no Smith form: no digit is lost). A torsion operator is read on the generators of the tree's presentation by
    Proposition~\ref{prop:torsionhecke}(3) ($2\sigma_{\max}\le K$), and its free part by (3$'$) there.
  \item[(d)] \textbf{(Glueing.)} $(\mathrm{Cyl}(\mathsf a^{(\ell)}),\mathsf B^{(\ell)}_{\mathcal T})\otimes\Z_{(\ell)}$ is a $\Z_{(\ell)}$-pair model, so Proposition~\ref{prop:descentpair}(2)(b) admits it for the $\ell$-primary
    parts and the $\ell$-local maps of \eqref{eq:descentpairles}, read on $\Cone{\mathsf a^{(\ell)\ast}}$ through $\varpi$; the glueing datum of the $\Z[1/6]$-piece ($\cmpI$) is unchanged. A map between the
    tree's $\ell$-local piece and anything computed on $\Iso$ (a Hecke operator, an explicit comparison) passes through $\mathsf j_\Phi$, $\mathsf j_\Psi$.
  \end{enumerate}
\item \textbf{(Refusal and scope.)} A run on the tree pair \textbf{refuses}, never truncates: when its predicted RSS exceeds the ceiling (rule~19, before building); when a read exhausts the
  precision (the stopping rule of (5)(a), (c)); for Hecke, on every condition of Proposition~\ref{prop:descentpairhecke}(5) (the pair lift's codes), and when $q_{\max}\notin\{4,5\}$ ($5$ is the top chain degree $m$ of the
  comparisons, (4)(e); v1.23.0-pre.35); when an
  exact product of the (L)-, (L$'$)-built maps leaves $64$-bit signed integers (the cap). \textbf{A failed certificate of a constructed map} --- a chain-map identity of $\mathsf a^{(\ell)}$, $\Phi^{\Z}_\partial$,
  $\Psi^{\Z}_\partial$, the identities of (3)(a), (3)(c), the strict blocks of (3)(d), (4)(a) --- is a \textbf{failure of the code} (\code{hecke::HeckeInternalError}), not a refusal: by (1)--(4) it cannot occur
  for a correct implementation. \textbf{Scope}: the groups and maps for every $\ell$, in every degree; the production pieces are $\ell\in\{2,3\}$ ($\ell\ge5$: the Bredon tree pair (2)(c) is a
  check; D4's pairs stay the route); $T_p$, $p\nmid N$, by transport only; not $U_p$, not nebentypus Hecke operators, not $\Ga_1(N)$; no native $\tau$-type datum on the tree.
\end{enumerate}
\end{proposition}

\begin{proof}
(1)(a) The column $F^{\mathfrak e}\to\cring[\Ga/\langle\mathbf c\rangle]$ is a projective resolution (Proposition~\ref{prop:isomodeltree}(1)(b)); $r_{(1-T)e_{\mathfrak e}}$ is a chain endomorphism ($T$ commutes
with $\mathbf c$ and $e_{\mathfrak e}$), lying over $y\langle\mathbf c\rangle\mapsto y\langle\mathbf c\rangle-yT\langle\mathbf c\rangle$ on $\cring[\Ga/\langle\mathbf c\rangle]$. $\Gbdry/\langle\mathbf c\rangle=\langle\bar T\rangle\cong\Z$ acts freely on the
right of $\Ga/\langle\mathbf c\rangle$ with orbit set $\Ga/\Gbdry=\mathbb{P}^1$ (\ref{B2}), so on each orbit $\cong\Z$ the map $1-T$ is injective on $\cring[\Z]$ with cokernel $\cring$: it is injective with
cokernel $\cring[\mathbb{P}^1]$. The cone of a map of resolutions has homology the cone of the map on $H_0$ \STD{\cite[\S1.5]{weibel}}: $H_1=\ker=0$, $H_0=\coker=\cring[\mathbb{P}^1]$. $\varepsilon^{(\ell)}_{\mathsf B}$ is the
column augmentation, a map of cones to the cone $C_{\bullet}(\mathcal L)\otimes\cring$ of $1-T$ on $\cring[\Ga/\langle\mathbf c\rangle]$, a quasi-isomorphism by the five lemma, surjective.
(1)(b) $e_{\mathfrak e}$ is central in $\cring[\langle S\rangle]$ and $\cring[\langle U\rangle]$ and $e_{\mathfrak e}e_S=e_S$, $e_{\mathfrak e}e_U=e_U$ (Proposition~\ref{prop:isomodeltree}(1)(a)), and $e_{\mathfrak e}$ commutes with
$S^{-1}$, $T$. \emph{Chain map}, on $\mathsf e_t$: $a_t\delta^{S}_t=\varepsilon_ta_{t-1}$, the $S$-half of the chain-map identity of the canonical lift (loc.\ cit.). On $\mathsf z_t$, write
$\gamma:=(S^{-1}-T)e_{\mathfrak e}=S^{-1}(1-U)e_{\mathfrak e}$ ($S^{-1}U=T$) and $\alpha_t$, $\beta_t$ for the $\mathsf v^{S}_{t+1}$-, $\mathsf v^{U}_{t+1}$-coefficients; $\partial\mathsf a(\mathsf z_t)=\mathsf a(\partial\mathsf z_t)$ reads, on
$F^{\mathfrak e}_{t-1}$: $-\gamma\varepsilon_t=-\varepsilon_t\gamma$ (central $\varepsilon_t$); on $F^{S}_t$: $\alpha_t\delta^{S}_{t+1}=(1-S^{-1})a_t-\varepsilon_t\alpha_{t-1}$ ($1-T-\gamma=1-S^{-1}$ on $e_S$); on $F^{U}_t$:
$\beta_t\delta^{U}_{t+1}=-\gamma b_t-\varepsilon_t\beta_{t-1}$. For even $t$ ($\alpha_t=-S^{-1}e_S$, $\beta_t=S^{-1}e_U$, $\alpha_{t-1}=\beta_{t-1}=0$, $\delta^{S}_{t+1}=(1-S)e_S$, $\delta^{U}_{t+1}=(1-U)e_U$): $-S^{-1}(1-S)=1-S^{-1}$
and $a_t=1$ (resp.\ $e_S$), and $S^{-1}(1-U)e_U=-\gamma b_t$ as $b_t=-e_U$ ($-1$ on the free model). For odd $t$ ($\alpha_t=\beta_t=0$, $\alpha_{t-1}=-S^{-1}e_S$, $\beta_{t-1}=S^{-1}e_U$,
$\varepsilon_t=1-\mathbf c$): $(1-S^{-1})(1+S)=S-S^{-1}=-(1-\mathbf c)S^{-1}$ ($\mathbf cS^{-1}=S$), so $(1-S^{-1})a_t-\varepsilon_t\alpha_{t-1}=0$; and $-\gamma b_t=S^{-1}(1-U)(1+U+U^{2})=S^{-1}(1-\mathbf c)=\varepsilon_t\beta_{t-1}$ on the
free model, $S^{-1}(1-U)e_U=S^{-1}(1-\mathbf c)e_U$ at $\ell=2$ ($Ue_U=\mathbf ce_U$), so $-\gamma b_t-\varepsilon_t\beta_{t-1}=0$. Where a target
column vanishes the corresponding equation is empty or reads $(1-S^{-1})e_S=0$ ($S\in Q_S$ at $\ell\ge3$), resp.\ $\gamma e_U=S^{-1}(1-U)e_U=0$ ($U\in Q_U$ at $\ell\ge5$). \emph{Over $\aug$ and $\mathsf f_{\mathsf a}$}: in
degree $0$, $\mathsf a(\mathsf e_0)=a_0s_0$ augments to $\mathfrak v_S$, $\aug(a_0)=1$ ($\varepsilon^{(\ell)}_{\mathcal T}(a_0\mathsf v^{S}_0)=a_0\Ga_S=\mathfrak v_S=\mathsf f_{\mathsf a}\varepsilon^{(\ell)}_{\mathsf B}(\mathsf e_0)$, as $a_0\in\{1,e_S\}$); in degree $1$ only the edge component $\gamma\mathsf v^{\mathfrak e}_0$ survives $\varepsilon^{(\ell)}_{\mathcal T}$, giving $(S^{-1}-T)\mathfrak e=\mathsf f_{\mathsf a}(\mathfrak e_\infty)=\mathsf f_{\mathsf a}\varepsilon^{(\ell)}_{\mathsf B}(\mathsf z_0)$; both sides vanish on column degrees $\ge1$.
(1)(c) Lemma~\ref{lem:treepaircyl}; $\Iso^{(\ell)}_{\mathcal T}$ is an $\cring$-model (Proposition~\ref{prop:isomodeltree}(1)(b)), $\mathsf B^{(\ell)}_{\mathcal T}$ resolves $\cring[\mathbb{P}^1]$ by (a), $\mathsf a^{(\ell)}$ lies over $\aug$ by (b).

(2) $\Hom(r_a,M)$ is $m\mapsto am$, and $\Hom(\cring[\Ga]e_x,M)=M^{Q_x}$ (Proposition~\ref{prop:isomodeltree}(1)(c), (2)'s proof): $d^{q}(u,w)$ is $\varepsilon_{q+1}u$ on $\mathsf e_{q+1}$ and $(1-T)u-\varepsilon_qw$ on $\mathsf z_q$;
$(\mathsf a^{\ast}c)(\mathsf e_q)=a_qx$, $(\mathsf a^{\ast}c)(\mathsf z_{q-1})=-[q\text{ odd}]S^{-1}x+[q\text{ odd}]S^{-1}y+(S^{-1}-T)z$ ($e_Uy=y$, $e_Sx=x$). At $\ell\ge3$ only degrees $\le1$ remain. At $\ell\ge5$,
$(1-S^{-1})M^{-I}=(1-S)M^{-I}$ ($S^{-1}$ is invertible and commutes with $1-S$), and $(1-T)M^{-I}+(1-S^{-1})M^{-I}$ is $\sum_g(g-1)M^{-I}$ ($S$, $T$ generate $\Ga$); with $2$ invertible $M^{-I}\cong M_{\{\pm I\}}$, whose
$\Ga$-coinvariants are $M_\Ga$ ($\sum_g(g-1)M=(S-1)M+(T-1)M$ by $(gh-1)m=(g-1)hm+(h-1)m$). (d): Proposition~\ref{prop:isomodeltree}(2)(e), (1)(c).

(3)(a) Inspection on the representatives, as for Proposition~\ref{prop:isomodeltree}(4)(a): $\partial\mathsf i_\partial(\mathfrak e_\infty)=-(Tv_0-v_0)=\mathsf i_\partial((1-T)\mathfrak v_\infty)$;
$\mathsf r_\partial(\partial v_2v_3)=T\mathfrak v_\infty-\mathfrak v_\infty=\partial(-\mathfrak e_\infty)$; $\partial\mathsf G_\partial(v_2)=v_2-v_0$; on $v_2v_3$, $\mathsf G_\partial(Tv_0-v_2)=-(v_0v_1+v_1v_2)=v_2v_3-\mathsf i_\partial\mathsf r_\partial(v_2v_3)$;
$\mathsf r_{\mathcal T}(v_0)-\mathsf f_{\mathsf a}\mathsf r_\partial(v_0)=\mathfrak v_U-\mathfrak v_S=\partial\mathsf g_\square(v_0)$; on $v_2v_3$, $-S\mathfrak e+(S^{-1}-T)\mathfrak e=-T\mathfrak e=\mathsf g_\square(Tv_0-v_2)$ ($S^{-1}\mathfrak e=S\mathfrak e$, $\mathbf c$
fixing $\mathfrak e$); $\mathsf g'$: the identity $\mathsf i_{\mathcal T}\mathsf r_{\mathcal T}=\id-\partial\mathsf G_{\mathcal T}-\mathsf G_{\mathcal T}\partial$ of Proposition~\ref{prop:isomodeltree}(4)(a) composed with $j_X\mathsf i_\partial$, and the
$\mathsf g_\square$-identity composed with $\mathsf i_\partial$ (using $\mathsf r_\partial\mathsf i_\partial=\id$), give $j_X\mathsf i_\partial-\mathsf i_{\mathcal T}\mathsf f_{\mathsf a}=\partial(-\mathsf g')+(-\mathsf g')\partial$; the displayed values are
$\mathsf G_{\mathcal T}$'s table. (b) The columns of $\mathsf B_{\mathcal T}$ are the tree's edge column, contracted by Proposition~\ref{prop:isomodeltree}(4)(b)'s $\mathsf h$ with the same sign rules, and the descent
through the two columns terminates as there. \emph{The generalised solver} (v1.23.0-pre.34, referee m2; loc.\ cit.\ states it for $p\in\{\cmpI,\varepsilon_{\mathcal T}\}$ only): the proof of
Proposition~\ref{prop:isomodeltree}(4)(b) uses of $p$ only that it is the column augmentation of finitely many free periodic columns. On one column, $\mathsf h$ (with $\mathsf h_{-1}$) is a contracting
homotopy of the augmented column, $\partial\mathsf h+\mathsf h\partial=\id$ in degrees $\ge1$ and $\partial\mathsf h=\id-\mathsf h_{-1}\varepsilon$ in degree $0$, a $\Z$-linear identity on the basis $\rho g^{j}$; the global sign $\epsilon_s$ of
the column in the total complex is compensated by $\epsilon_s\mathsf h$. If $z\in\ker p$ is a cycle with top component $z_s$ in column $s$, then $z_s$ is a cycle of column $s$ in column degree $\ge1$, or
in column degree $0$ with augmentation $0$ (as $pz=0$), so $\partial(\epsilon_s\mathsf h(z_s))$ has column-$s$ component $z_s$ and its other components lie in columns $<s$; $z-\partial k_s$ is a
cycle of $\ker p$ with no column-$s$ component, and induction on $s$ ends after $s_{\max}+1$ steps with $\partial\mathsf h_p(z)=z$; $\mathsf h_p(z)$ lies in column degrees $\ge1$, so $p\,\mathsf h_p(z)=0$.
(L) and (L$'$) then follow from these two properties alone, by the induction on the degree of loc.\ cit. Every identity is $\Z$-linear on bases, so it survives $\otimes\cring$; the
$\cring[\Ga]$-linear extensions of (L), (L$'$) need only that $P'\otimes\cring$ be free over $\cring[\Ga]$ on the given basis, and nothing is asked of the coefficients of $A$, $B$, $f$, $g$
beyond $\cring[\Ga]$-linearity, so denominators in $\mathsf O_\ell$ and the idempotents $e_x$ are allowed. For $p=\varepsilon_{\mathsf B}$ the columns are $\mathsf e$ ($s=0$) and $\mathsf z$ ($s=1$, $\epsilon=-1$, the cone's
shifted column); for $p=\cmpI|_{\Isob}$ the columns $G_{s,\bullet}$ over $\Sigma^{\partial}_s$, $s\le1$, and $\epsilon_s\mathsf h$ maps each into itself, so the descent stays in $\Isob$ (as in the
boundary step of the proof of Proposition~\ref{prop:descentpairhecke}(1)).
$\Phi^{\partial}_\ell$, $\Psi^{\partial}_\ell$: the column computation of Proposition~\ref{prop:isomodeltree}(3) for $x=\mathfrak e$, and $r_{(1-T)e}$ commutes with $r_{ce}$ for central scalars $c$; the corrections $k$, $k'$ of
loc.\ cit.\ arise only between columns of different stabilisers, of which $\mathsf B_{\mathcal T}$ has none. (c) (L$'$)'s hypotheses: $\varepsilon_{\mathcal T}\Phi^{\Z}_{\mathcal T}j=\mathsf r_{\mathcal T}\cmpI j=\mathsf r_{\mathcal T}j_X\cmpI|_{\Isob}$;
$\varepsilon_{\mathcal T}\Psi^{\mathcal T}_\ell\mathsf a^{(\ell)}\Phi_\partial=\varepsilon^{(\ell)}_{\mathcal T}\mathsf a^{(\ell)}\Phi^{\partial}_\ell\Phi^{\Z}_\partial=\mathsf f_{\mathsf a}\varepsilon^{(\ell)}_{\mathsf B}\Phi^{\partial}_\ell\Phi^{\Z}_\partial=\mathsf f_{\mathsf a}\varepsilon_{\mathsf B}\Phi^{\Z}_\partial=\mathsf f_{\mathsf a}\mathsf r_\partial\cmpI|_{\Isob}$
($\varepsilon_{\mathcal T}\Psi^{\mathcal T}_\ell=\varepsilon^{(\ell)}_{\mathcal T}$, $\varepsilon^{(\ell)}_{\mathsf B}\Phi^{\partial}_\ell=\varepsilon_{\mathsf B}$, (1)(b) with the $\ell$-local augmentations; tensored with $\cring$ throughout; v1.23.0-pre.34, referee m4); and $\cmpI\Psi_{\mathcal T}\mathsf a^{(\ell)}\Phi^{\partial}_\ell=\mathsf i_{\mathcal T}\mathsf f_{\mathsf a}\varepsilon_{\mathsf B}$, $\cmpI j\Psi^{\Z}_\partial=j_X\mathsf i_\partial\varepsilon_{\mathsf B}$, with (a)'s homotopies. Then
$\Phi_{\mathcal T}j-\mathsf a\Phi_\partial=\Phi^{\mathcal T}_\ell(\Phi^{\Z}_{\mathcal T}j-\Psi^{\mathcal T}_\ell\mathsf a\Phi_\partial)$ by $\Phi^{\mathcal T}_\ell\Psi^{\mathcal T}_\ell=\id$, and
$\Psi_{\mathcal T}\mathsf a-j\Psi_\partial=(\Psi_{\mathcal T}\mathsf a\Phi^{\partial}_\ell-j\Psi^{\Z}_\partial)\Psi^{\partial}_\ell$ by $\Phi^{\partial}_\ell\Psi^{\partial}_\ell=\id$.
(d) \emph{Cochain maps}: the cone lemma ($\ast$) of Proposition~\ref{prop:isomodeltree}'s proof, dualised: $(\alpha,c)\mapsto(F^{\ast}\alpha,G^{\ast}c+K\alpha)$ commutes with the two $\Dcone$ iff
$dK+Kd=G^{\ast}\mathsf a^{\ast}-j^{\ast}F^{\ast}$, and (c) dualises to $j^{\ast}\Phi^{\ast}_{\mathcal T}-\Phi^{\ast}_\partial\mathsf a^{\ast}=\mathsf Y^{\ast}_\Phi d+d\mathsf Y^{\ast}_\Phi$ ($(\partial Y)^{\ast}=Y^{\ast}d$, $(Y\partial)^{\ast}=dY^{\ast}$), so $K=-\mathsf Y^{\ast}_\Phi$;
likewise for $\mathsf j_\Psi$. The strict identities are read off. \emph{The comparison of pairs.} Put $\mathrm{pr}_\partial\colon\Iso\to\Isob$ the projection along the columns over $\Sigma\setminus\Sigma^{\partial}$,
$\widehat{\mathsf Y}:=\mathsf Y_\Phi\mathrm{pr}_\partial$, and $\Theta_\Phi:=i_{\mathsf P}\Phi_{\mathcal T}-\partial K_\Phi-K_\Phi\partial\colon\Iso\otimes\cring\to\mathrm{Cyl}(\mathsf a^{(\ell)})$ with $K_\Phi(y):=(\mathsf Y_\Phi y,0,-\Phi_\partial y)$ on $\Isob$, $0$
on the complement. Since $i_{\mathsf P}\mathsf a-i_{\mathsf B}=-(\partial s+s\partial)$ on $\mathsf B$ ($s$ of Lemma~\ref{lem:treepaircyl}'s proof), $i_{\mathsf P}\Phi_{\mathcal T}j-i_{\mathsf B}\Phi_\partial=\partial K_\Phi j+K_\Phi\partial j$, so
$\Theta_\Phi|_{\Isob}=i_{\mathsf B}\Phi_\partial$: $\Theta_\Phi$ is a comparison of pairs $(\Iso,\Isob)\to(\mathrm{Cyl},\mathsf B)$ (over $\id_\cring$: $\partial K_\Phi$ has augmentation $0$). As $\pi_{\mathrm{Cyl}}\Theta_\Phi=\Phi_{\mathcal T}-\partial\widehat{\mathsf Y}-\widehat{\mathsf Y}\partial$,
\[
  \mathsf j_\Phi-(\Theta_\Phi^{\ast}[1],\Phi_\partial^{\ast})\circ\varsigma'=\Dcone\Upsilon+\Upsilon\Dcone,\qquad\Upsilon(\alpha,c):=(-\widehat{\mathsf Y}^{\ast}\alpha,0)
\]
(both sides send $(\alpha,c)$ to $(d\widehat{\mathsf Y}^{\ast}\alpha+\widehat{\mathsf Y}^{\ast}d\alpha,\,-\mathsf Y_\Phi^{\ast}\alpha)$, as $\rres\widehat{\mathsf Y}^{\ast}=\mathsf Y^{\ast}_\Phi$). And $\Theta_\Psi(x,b,u):=\Psi_{\mathcal T}x+j\Psi_\partial b-\mathsf Y_\Psi u$ is a
comparison of pairs $(\mathrm{Cyl},\mathsf B)\to(\Iso,\Isob)\otimes\cring$ ($\partial\Theta_\Psi-\Theta_\Psi\partial=(\Psi_{\mathcal T}\mathsf a-j\Psi_\partial-\partial\mathsf Y_\Psi-\mathsf Y_\Psi\partial)u=0$) with
$\mathsf j_\Psi=\varpi\circ(\Theta_\Psi^{\ast}[1],\Psi_\partial^{\ast})$ \textbf{exactly}. Hence, by Lemma~\ref{lem:pairlifting}(2), (3) for the self-comparisons $\Theta_\Phi\Theta_\Psi$ of $(\mathrm{Cyl},\mathsf B)$ and
$\Theta_\Psi\Theta_\Phi$ of $(\Iso,\Isob)$, and Lemma~\ref{lem:treepaircyl}(2): $\mathsf j_\Psi\mathsf j_\Phi\simeq\varpi\bigl((\Theta_\Phi\Theta_\Psi)^{\ast}[1],(\Phi_\partial\Psi_\partial)^{\ast}\bigr)\varsigma'\simeq\varpi\varsigma'=\id$ and
$\mathsf j_\Phi\mathsf j_\Psi\simeq(\Theta^{\ast}_\Phi[1],\Phi^{\ast}_\partial)\varsigma'\varpi(\Theta^{\ast}_\Psi[1],\Psi^{\ast}_\partial)\simeq\bigl((\Theta_\Psi\Theta_\Phi)^{\ast}[1],(\Psi_\partial\Phi_\partial)^{\ast}\bigr)\simeq\id$. (The five lemma alone would not
do: a map of short exact sequences of cones whose outer maps induce the identity need not induce the identity on the cone, Proposition~\ref{prop:coneambiguity}.) Independence: two choices give two
comparisons of pairs, homotopic as pairs (Lemma~\ref{lem:pairlifting}(2)).

(4)(a) Expand $\mathsf j_\Psi T\mathsf j_\Phi(\alpha,c)=\mathsf j_\Psi(T_{\mathrm b}\Phi^{\ast}_{\mathcal T}\alpha,\ T_\partial\Phi^{\ast}_\partial c-T_\partial\mathsf Y^{\ast}_\Phi\alpha)$. A triangular endomorphism $(\alpha,c)\mapsto(A\alpha,Bc+E\alpha)$ of
$\Cone{\mathsf a^{\ast}}$ commutes with $\Dcone$ iff $A$, $B$ are cochain maps and $dE+Ed=B\mathsf a^{\ast}-\mathsf a^{\ast}A$ (compare the second components of $\Dcone(A\alpha,Bc+E\alpha)$ and of the image of
$(-d\alpha,dc+\mathsf a^{\ast}\alpha)$); $T^{\mathrm c}_{\mathcal T}$ is a composite of cochain maps. \emph{The pair model} (v1.23.0-pre.34, referee m1): $\Theta_\Phi$, $\Theta_\Psi$ are comparisons of pairs, so
$\rres_{\mathrm{Cyl}}\Theta_\Psi^{\ast}=\Psi_\partial^{\ast}\rres$ and $\rres\,\Theta_\Phi^{\ast}=\Phi_\partial^{\ast}\rres_{\mathrm{Cyl}}$ strictly, and with $\rres\,T_{\mathrm b}=T_\partial\rres$ every factor of $T^{\mathrm{Cyl}}$ is diagonal; by the
proof of (3)(d), $\mathsf j_\Psi=\varpi(\Theta_\Psi^{\ast}[1],\Psi_\partial^{\ast})$ exactly and $\mathsf j_\Phi\simeq(\Theta_\Phi^{\ast}[1],\Phi_\partial^{\ast})\varsigma'$ (by $\Upsilon$), so
$T^{\mathrm c}_{\mathcal T}-\varpi T^{\mathrm{Cyl}}\varsigma'=\Dcone\Upsilon'+\Upsilon'\Dcone$ with $\Upsilon':=\mathsf j_\Psi T\Upsilon$. That $\eta_{\mathcal T}$ need not vanish and the diagonal operator need not be a cochain map is
an \emph{observation} of the checks (arm H below, for $T=\theta^{\ast}$, $\theta=\id+\partial K+K\partial$, $K(\Isob)\subseteq\Isob$; the referee's arm H for $T_\varepsilon$), not a proved assertion; nothing
else in (4) depends on it. (b) The blocks are read off; $H(\mathsf j_\Psi)=H(\mathsf j_\Phi)^{-1}$ by (3)(d), and Proposition~\ref{prop:descentpairhecke}(3) for
$\Phi^{(N)}_p$. (c) By (3)(d), $H(T^{\mathrm c}_1)H(T^{\mathrm c}_2)=H(\mathsf j_\Psi)H(T_1)H(\mathsf j_\Phi)H(\mathsf j_\Psi)H(T_2)H(\mathsf j_\Phi)$ and $H(\mathsf j_\Phi\mathsf j_\Psi)=\id$. (d) $u$ acts by $\Hom(-,u)$, which
commutes with precomposition by every chain map; Proposition~\ref{prop:descentpairhecke}(6)'s proof then applies.
(e) (v1.23.0-pre.35.) \emph{(i)} The targets $\Iso\otimes\cring$, $\Isob\otimes\cring$, $\Iso^{(\ell)}_{\mathcal T}$, $\mathsf B^{(\ell)}_{\mathcal T}$, $\mathsf B_{\mathcal T}\otimes\cring$ are acyclic in positive degrees
($\Iso$, $\Isob$ are resolutions by free abelian groups, hence $\Z$-split, and stay exact after $\otimes\cring$; the $\ell$-local models by Proposition~\ref{prop:isomodeltree}(1)(b) and (1)(a) here), and
the sources are projective over $\cring[\Ga]$. The comparison theorem's induction \STD{\cite[Thm.~2.2.6]{weibel}} extends a family $f_q$, $q\le m$, with $\partial f_q=f_{q-1}\partial$: for a summand $\cring[\Ga]e$ of
$P_{m+1}$ ($e$ idempotent; $e=1$ where the source is free; v1.23.0-pre.36, the C-D8b referee's N3), $f_m\partial e$ is a cycle ($\partial f_m\partial=f_{m-1}\partial\partial=0$) of positive degree $m$, hence a boundary $\partial y$,
and $f_{m+1}(x):=xy$ for $x\in\cring[\Ga]e$ ($\Ga$-linear, and $\partial f_{m+1}(x)=x\,\partial y=xf_m\partial e=f_m\partial x$ since $x=xe$; equivalently, by projectivity). A homotopy $Y_q$, $q\le m-1$, with
$f_q-g_q=\partial Y_q+Y_{q-1}\partial$ extends likewise: $z:=f_m-g_m-Y_{m-1}\partial$ satisfies $\partial z=(f_{m-1}-g_{m-1}-\partial Y_{m-1})\partial=Y_{m-2}\partial\partial=0$, so $z(e)=\partial y'$ and $Y_m(x):=xy'$ on each summand $\cring[\Ga]e$, as before; then the maps in
degree $m+1$, and so on. The conditions ``over'' the augmentations and the tables live in degrees $0$, $1\le m$, and the proofs of (3)(d) and (4)(a)--(c) use of $\Phi$, $\Psi$, $\mathsf Y$ only the identities
(3)(b), (c) (not that they were built by (L), (L$'$)), so they apply to the extension. \emph{(ii)} By the formulas of (3)(d) and (4)(a): on $(\alpha,c)\in C^{j+1}_{\mathcal T,\ell}\oplus\mathsf B^{j}_\ell$,
$\mathsf j_\Phi$ uses $\Phi_{\mathcal T}^{\ast}$ in cochain degree $j+1$, $\Phi_\partial^{\ast}$ in degree $j$ and $\mathsf Y_\Phi^{\ast}\colon C^{j+1}\to{\Isob}^{j}$, the dual of $\mathsf Y_\Phi$ in chain degree $j$; likewise
$\mathsf j_\Psi$ and $T$ in cone degree $j$. \emph{(iii)} Let $\widetilde T^{\mathrm c}$ be built from an extension; it is a cochain map by (4)(a), (b), and agrees with $T^{\mathrm c}_{\mathcal T}$ in cone degrees $\le m-1$ by (ii).
For a cocycle $g$ of cone degree $j\le m-1$, $\Dcone T^{\mathrm c}_{\mathcal T}g=\Dcone\widetilde T^{\mathrm c}g=\widetilde T^{\mathrm c}\Dcone g=0$ (the extension enters only through $\widetilde T^{\mathrm c}(0)=0$); for $g=\Dcone x$, $x$ of cone degree $j-1$,
$T^{\mathrm c}_{\mathcal T}g=\widetilde T^{\mathrm c}\Dcone x=\Dcone\widetilde T^{\mathrm c}x=\Dcone T^{\mathrm c}_{\mathcal T}x$; and the induced map is $H(\widetilde T^{\mathrm c})$, which is (b)'s and, by (c), does not depend on the extension.
\emph{(iv)} $\eta^{j+1}_{\mathcal T}$ is the off-diagonal block of $T^{\mathrm c}_{\mathcal T}$ in cone degree $j+1$, so by (ii) it uses $\Phi$, $\Psi$ in chain degree $j+2$; formable iff $j+2\le m$. By the triangular
criterion of the proof of (a), $T^{\mathrm c}_{\mathcal T}$ commutes with $\Dcone$ on cone degree $j$ iff $T'_{\mathcal T}$ on $C^{j+1}$, $T'_\partial$ on $\mathsf B^{j}$ commute with $d$, the first component of
$T^{\mathrm c}_{\mathcal T}(0,c)$ vanishes, and the identity of (a) holds on $C^{j+1}_{\mathcal T,\ell}$; with $j=q-2$ this gives the coboundary half of (iii) for $\Hcpt{q}$ (commuting with $\Dcone$ maps coboundaries to coboundaries; the converse need not hold). At $\ell=3$, $\eta^{j}_{\mathcal T}$ has target
$\mathsf B^{j}_3=0$ for $j\ge2$.

(5)(a) The integral forms are Proposition~\ref{prop:isomodeltree}(1)(d) and (1)(a) here (every coefficient of $\mathsf B^{(\ell)}_{\mathcal T}$ and $\mathsf a^{(\ell)}$ in $\mathsf O_\ell[\Ga]$). Proposition~\ref{prop:descentpair}(6)'s proof
uses of the cone only that it is a bounded below complex of finitely generated free $\Z_{(\ell)}$-modules computing $\Hcpt{\bullet}$ and the boundary term (here through $\varpi$, Lemma~\ref{lem:treepaircyl}), since
the elementary divisors of a differential of any such complex are the torsion exponents of the next cohomology group, and the rank sources' criteria are properties of the groups; the
exclusion of other complexes' ranks is (6)'s own argument (the rank of a differential is a property of that differential). (b) The rational long sequence of
Proposition~\ref{prop:descentpair}(5)(b) is exact with $\Hcpt{0}=0$, $\res^{0}$ injective and $H^{q}(\Ga;M_\Q)=H^{q}(\Gbdry;M_\Q)=0$ for $q\ge2$. (c) The comparisons are built over $\Z$ (resp.\
$\mathsf O_\ell$: $\Phi^{\mathcal T}_\ell$, $\Psi^{\mathcal T}_\ell$, $\Phi^{\partial}_\ell$, $\Psi^{\partial}_\ell$ have coefficients in $\mathsf O_\ell[\Ga]$, and (L), (L$'$) preserve integrality); products of exact matrices are exact.
(d) Proposition~\ref{prop:descentpair}(2)(b) and Lemma~\ref{lem:treepaircyl}. (6) collects (1)--(5).
\end{proof}

\begin{remark}[What the plan's \S5 stated, and what holds]\label{rmk:treepairscope}
\PROVED{} (v1.23.0-pre.33.) (i) \textbf{The construction.} The owner's candidate --- the cylinder or cone of $\Ind_{\Gbdry}^{\Ga}C(\text{axis of }T)\to C(\mathcal T)$ --- holds in the algebraic form of the
plan's \S5.1, with the \emph{minimal} cusp resolution ($1,2,2,\dots$; the plan's ``$F^\infty$''), not the cellular axis ($2,4,4,\dots$): the axis map $\mathsf a$ needs no geometric axis, only
the path from $T\mathfrak v_S$ through $S^{-1}\mathfrak v_U$ ($=T\mathfrak v_U$) to $\mathfrak v_S$, i.e.\ $\mathfrak e_\infty\mapsto(S^{-1}-T)\mathfrak e$, in the orientation of $\mathfrak e_\infty$
($\partial\mathfrak e_\infty=(1-T)\mathfrak v_\infty$, $\partial\mathsf f_{\mathsf a}(\mathfrak e_\infty)=\mathfrak v_S-T\mathfrak v_S$; v1.23.0-pre.34, referee N2). (ii) \textbf{The computed object is the cone of $\mathsf a^{\ast}$, not the cylinder pair's cone}: the
cylinder carries the boundary twice and is only the pair model that lets \S\ref{sec:descentpair} apply (Lemma~\ref{lem:treepaircyl}). (iii) \textbf{``No $\eta$ appears on the tree side'' (plan
\S5.4(c)1) holds for $\Cone{\rres_{\mathrm{Cyl}}}$ and fails for $\Cone{\mathsf a^{\ast}}$}, the object the code computes (v1.23.0-pre.34, referee m1): transported along the strict comparisons of pairs
$\Theta_\Phi$, $\Theta_\Psi$, the operator on the pair model's cone is diagonal, $\eta=0$, as the plan says; on $\Cone{\mathsf a^{\ast}}$, which is not the cone of a restriction to a subcomplex,
$T^{\mathrm c}_{\mathcal T}\simeq\varpi T^{\mathrm{Cyl}}\varsigma'$ carries the off-diagonal $\eta_{\mathcal T}$ of Proposition~\ref{prop:treepair}(4)(a), observed non-zero in the checks (not proved in general); it is
determined, so the ambiguity of Proposition~\ref{prop:coneambiguity} does not return. (iv) The plan's ``strict maps of pairs $(\Iso,\Isob)\rightleftarrows(\mathrm{Cyl},\mathsf B)$'' exist ($\Theta_\Phi$,
$\Theta_\Psi$ of the proof of (3)(d)); the code uses the cone maps $\mathsf j_\Phi$, $\mathsf j_\Psi$ of the homotopy-commutative squares instead, which are smaller and need no cylinder; that they are
mutually inverse on cohomology is proved through $\Theta_\Phi$, $\Theta_\Psi$, not by the five lemma. (v) The comparison is over $\Z$ and $\Z[\mathbb{P}^1]$ through $C_{\bullet}(\mathcal T)$, $C_{\bullet}(\mathcal L)$ and
$C_{\bullet}(\Xbdry)$, as the plan said; the geometric statement through Ash's retraction stays \textbf{not needed}; ``every edge lies on exactly two axes'' (plan \S5.1) is not used and not verified
here. (vi) \textbf{Cost.} The groups need the tree cone only (ranks (2)(d)); the comparisons and the transport touch $\Iso$-sized cochains, as a product, never an elimination; the plan's proxy
($\approx\tfrac1{170}$ of the $\Iso^{(2)}$ pair in time) is a prediction, not a statement of this document.
\end{remark}

\noindent\emph{A check outside the build, not a tag} (v1.23.0-pre.33): \file{docs/implementation/slot24/cd8\_check.py} (\code{python3 -I}, pure integers and fractions, no computer algebra system), output
\file{cd8\_check.out}. It imports the group ring, the tree models and $\Phi^{\mathcal T}_\ell$, $\Psi^{\mathcal T}_\ell$ of \file{cd6\_tree\_check.py} and the model $\Iso$, the solvers and (L), (L$'$) of
\file{cd6\_iso\_check.py}, reads \file{d4\_pair\_oracle.tsv} as data, and builds everything of this subsection from its text: $\mathsf B^{(\ell)}_{\mathcal T}$, $\mathsf a^{(\ell)}$ at $\ell\in\{0,2,3,5\}$, the tables of
Definition~\ref{def:treepair}(4), $\Phi^{\Z}_\partial$, $\Psi^{\Z}_\partial$, $\mathsf H^{\mathsf B}_\partial$, $\mathsf H^{X}_\partial$, $\mathsf y_\ell$, $\mathsf y'_\ell$, $\mathsf j_\Phi$, $\mathsf j_\Psi$, the cylinder maps and $T^{\mathrm c}_{\mathcal T}$. Exactly in
$\Q[\Ga]$ to degree $6$: $\partial^{2}=0$ on the four cusp models, $\mathsf a^{(\ell)}$ a chain map over $\mathsf f_{\mathsf a}$, the identities of (3)(a)--(c); on the $32$ cells of the D4 oracle (level $1$,
$n\le12$; $\Ga_0(N)$, $N\le13$), over $\Z/\ell^{K}$: (2)'s closed forms entry by entry, $H^{q}(\mathsf B^{(\ell)})$ against the oracle's boundary groups and $\Hcpt{q}$, $q\le4$, of the tree cone against its
$\Hcpt{q}$ ($\ell$-parts at $2$, $3$, $5$; ranks from the Betti recursion, certified pivots, a-priori exponents), the lengths of $H^{q}$ of the tree cone against $(\Iso,\Isob)$'s cone, $\mathsf j_\Phi$, $\mathsf j_\Psi$
cochain maps with the strict $\coneinc$, $\coneproj$ blocks and mutually inverse on cohomology modulo $\ell$, Lemma~\ref{lem:treepaircyl}(2)'s identities; and (4) for $T=\theta^{\ast}$ ($\theta=\id+\partial K+K\partial$,
$K(\Isob)\subseteq\Isob$, not a fibre map) and for fibre maps $u\in\{1,1+2J\}$ ($J$ the all-ones endomorphism of $\Z[\mathbb{P}^1(\F_N)]$, $N\in\{5,7\}$): $T^{\mathrm c}_{\mathcal T}$ a cochain map with the stated
blocks, $\eta_{\mathcal T}\ne0$, inducing $H(T)$. Planted faults, each detected: $S$ for $S^{-1}$ in $\mathsf a$ (at chain level, exactly in $\Q[\Ga]$, at $\ell\in\{0,2\}$ only: at $\ell\ge3$ it is the same map,
Definition~\ref{def:treepair}(2); the C-D8 referee's check detects instead the dropped term $S^{-1}e_U\mathsf v^{U}_1$ at $\ell=0,2,3$; v1.23.0-pre.34, referee m3); a sign of $\mathsf r_\partial$; $+\mathsf Y^{\ast}_\Phi$ in $\mathsf j_\Phi$; $1-T^{2}$ for $1-T$ in
$\mathsf B$; $\eta_{\mathcal T}$ dropped (for $T=\theta^{\ast}$ the diagonal operator fails to be a cochain map on all $32$ cells at $\ell\in\{0,2\}$ and on $24$ of the $32$ at $\ell=3$ --- at $\ell=3$ the $7$ cells with $n$ odd have
$M^{-I}=0$, so no $3$-local boundary and $\eta_{\mathcal T}=0$, and at $(1,0)$ it happens to commute; for fibre maps $\eta_{\mathcal T}$ is a cocycle and the diagonal operator is a cochain map). \code{CD8\_CHECK degree=6 cells=32 checks=67252 mismatch=0 \dots\ verdict=OK} (about $15$ minutes). Not exercised: $T=\Phi^{(N)}_p$ (the pair lift is in the
code only); the code's own $\Iso$ (the check's $\Iso$ is built by the same procedure, not read from \code{cohomology\_app}). \emph{Extended at v1.23.0-pre.35} by two arms: (M3) on every
cell at $\ell\in\{0,2\}$, $\rho_{M_N}(\mathbf c)=(-1)^n$, and the closed form of $\mathsf a^{\ast}$ with $S$ for $S^{-1}$ equals the true one modulo $2^{20}$ in every degree $q\le4$ iff $n$ is even
(Definition~\ref{def:treepair}(2)); (Q) the read at $q_{\max}=5$ (Proposition~\ref{prop:treepair}(4)(e)): in cone degree $4$, $T^{\mathrm c}_{\mathcal T}$ built from the maps of chain degree $\le5$ maps every
cocycle modulo $p$ to a cocycle and induces $H(T)$ on $\Hcpt{5}$ (non-zero modulo $p$ on cells at $\ell=0$ and $\ell=2$), and with the degree-$6$ extension (the check builds its comparisons to degree $6$)
the identity of (4)(a) holds in cone degree $4$ as well. \code{CD8\_CHECK degree=6 cells=32 checks=68168 mismatch=0 \dots\ verdict=OK} (about $17$ minutes; (M3) $256+64$ checks, (Q) $3\times198+2$, $\Hcpt{5}\ne0$ modulo $p$ on all $32$ cells at $\ell=0$ and at $\ell=2$).

\noindent\emph{Code map} (v1.23.0-pre.33, \textbf{stated before D5c's code}, contract first, rule~21; the status of v1.23.0-pre.35). \textbf{Built, audited and merged}
into \texttt{hecke-extensions} at \texttt{4064e164} (v1.23.0-pre.36; built at v1.23.0-pre.35 on branch \texttt{s24-d5cimpl}, then at \texttt{e720f11c}; the audit
\file{docs/implementation/slot24/D5c\_code\_audit.md} at \texttt{f7af1744}, ACCEPT after fixes, 0~M / 4~m / 7~N, its m1--m4, N1, N2, N4, N6 fixed at \texttt{e75c3937}, N3, N7 this document's, N5
deferred (below); the C-D8b referee's Arm (G) at \texttt{3e4bc339}; \file{docs/implementation/slot24/DESCENT5c\_report.md}, \S9 for the fixes). The chain-level half is \file{Integral/Model/TreePair.\{h,cpp\}}: the cusp
models (\code{CuspModelChain}), the axis map (\code{AxisMapChain}), the tables of Definition~\ref{def:treepair}(4), and \code{TreePairComparison} ($\Phi^{\Z}_\partial$, $\Psi^{\Z}_\partial$,
$\Phi^{\partial}_\ell$, $\Psi^{\partial}_\ell$, $\mathsf y_\ell$, $\mathsf y'_\ell$, $\mathsf Y_\Phi$, $\mathsf Y_\Psi$), whose \code{certify()} checks every identity of (1)(b) and (3)(a)--(c) exactly in $\Q[\Ga]$, once per
$\ell$ and process, in chain degrees $\le5$, on D5's comparison now built to degree $5$ (\code{kTreeComparisonTop}). The production half is \file{Integral/Descent/DescentTreePair.\{h,cpp\}}:
\code{BuildTreePairComplex} (the closed forms of (2), read by D4's \code{RunPairLes} verbatim, with $d_{\mathsf B}d_{\mathsf B}=0$ and $d_{\mathsf B}\mathsf a^{\ast}=\mathsf a^{\ast}d_{\mathcal T}$ checked entrywise on
every call), \code{TreeConeRanksModPrime} (the rank source on \textbf{this} cone, (5)(a)), \code{TreePairCochains} (every map of (3)(d), (4) applied to vectors, never as a dense $\Iso$-sized matrix) and
\code{ReadTreePairHecke}. \textbf{The tree cone is the default route} of the $2$- and $3$-adic pair pieces of \code{AssembleDescentPair}; the read-only rollbacks \code{HECKE\_DESCENT\_TREE\_PAIR=0}
and \code{HECKE\_DESCENT\_TREE=0} select D4's dense pair (\code{EFFECTIVE HeckeDescentTreePair route=\dots}). Validations: \code{IntegralDescentTreePairValidation} (standard; $2934$ checks, $0$ failing at the merge) and
\code{IntegralDescentTreePairSlowValidation} (slow), registered in \file{Registry/Slot24/DESCENT5c.cpp}, sources \file{cmake/slot24/DESCENT5c.cmake}; D0--D5's and D3's standard names pass on the default and on
both rollbacks, and \code{IntegralDescentPairValidation} prints its route (\code{EFFECTIVE IntegralDescentPairValidation}, \code{route=tree} or \code{route=dense}; the audit's N4). \textbf{Reach} (the report's \S4): the $2$-adic pair wall moves from $D\approx1060$--$1240$ (D4: refusal, resp.\ $40$~GiB) to the refusal at
$D\approx6340$ (rule~19, before building), $\approx8000$ at $40$~GiB; the next wall is set by D2's dense Smith form with transforms, not by the model; largest measured cell $D=6000$. (Re-timed nearly idle at the merge, the report's \S9.4: $D=852$ in
$11.3$~s, $D=4500$ in $1325$~s at a $12.8$~GiB peak; the filed time model, $\propto D^{2.97}$, is an upper estimate at $D\ge4500$, the two idle points giving exponent $2.86$; the memory model is unchanged.)
\textbf{Deferred} (the audit's N5): the validation-only fault hooks (\code{ActiveTreePairFault()}, a process-wide atomic as D3's \code{ScopedD3Fault}, and the planted error of
\code{ScopedPlantedCocycleEta}) are read inside production paths (\code{BuildTreePairComplex}, the cusp model of \file{TreePair.cpp}, \code{j\_phi}, \code{apply\_hecke}, \code{y\_psi\_star}); safe while the
validations run serially, as they do; making them thread-local or injected changes production signatures, and is left. The
$\ell$-local column augmentations $\varepsilon^{(\ell)}_{\mathcal T}$, $\varepsilon^{(\ell)}_{\mathsf B}$ of Definition~\ref{def:treepair}(4) are the code's (\code{TreeModelChain::augment}: the cosets of the
column groups in column degree $0$, $0$ above). \code{D3Refusal::HeldOutMismatch} is removed, which closes the open code item of \S\ref{sec:descenthecke}'s code map (the C-D8 referee's N7).
Below, the obligations as stated before the code, with (e) and (f) amended at v1.23.0-pre.35 and (c), (e), (f), (g) at v1.23.0-pre.36; on D5's \file{Integral/Model/TreeCochains} and
\file{IsotropyModelTree} and D4's \file{Integral/Descent/DescentPair}: \textbf{Obligations of the code}: (a) $\mathsf B^{(\ell)}_{\mathcal T}(M)$ and $\mathsf a^{(\ell)\ast}$ in the closed form of (2), block by block for a
module given by its block-monomial action, and the tree cone to cochain degree $q_{\max}+1$, $4\le q_{\max}\le5$ when Hecke is asked; (b) certification by (5)(a) on \textbf{this} cone, never with a
rank of $(\Iso,\Isob)$'s cone; the boundary by its closed form (Proposition~\ref{prop:descentpair}(4), \code{BoundaryClosedForm}); (c) the tree cone as the route of the $2$- and $3$-adic pair pieces, \textbf{every $N$}
(amended at v1.23.0-pre.36, the C-D8b referee's N7, the audit's N3: the code routes every $3$-adic pair piece to the tree, $\nu_3(N)=0$ included, and, since the audit's m3, $3\in$ \code{extra\_primes}
with $n$ odd, where the $3$-local cusp cochains have rank $\rank e_{\mathfrak e}M=D\cdot[n\text{ even}]=0=\rank M^{-I}$, as in (2)(b); D4's $3$-adic piece was dense everywhere, so nothing is lost), with a \textbf{read-only, value-preserving rollback} switch to D4's dense pair, read once per call --- the preserved values are \textbf{the record's invariants}
(the group types and the choice-free Hecke invariants of Proposition~\ref{prop:torsionhecke}), not Smith generators or coordinates, which the two cones do not share; ``invariants only'', as ruled for the
descent (v1.23.0-pre.34, referee N6); (d) $\Phi^{\Z}_\partial$, $\Psi^{\Z}_\partial$, $\mathsf y_\ell$, $\mathsf y'_\ell$ by (L),
(L$'$) on the code's $\Iso$, built once per run to the degree in use, with their certificates (chain maps, over $\mathsf r_\partial$, $\mathsf i_\partial$, the two square identities) checked when built, a failure
thrown as \code{hecke::HeckeInternalError}; (e) $T^{\mathrm c}_{\mathcal T}=\mathsf j_\Psi\Phi^{(N)}_p\mathsf j_\Phi$ for every $p\nmid N$ with a certified pair lift, applied to vectors (the presentation's generators),
never as a dense $\Iso$-sized matrix. \textbf{Certificates} (v1.23.0-pre.35; \textbf{restated at v1.23.0-pre.36} for the merged code, the D5c audit's m1 at \texttt{e75c3937} and the C-D8b referee's Arm (G) at
\texttt{3e4bc339}, the C-D8b referee's m1; Proposition~\ref{prop:treepair}(4)(e) with $m=5$): the read of $\Hcpt{q}$, $q\le q_{\max}\le5$, rests on (1)--(4), the chain-level certificates of (d) (exact,
independent of $M$ and of $D$) and the pair lift's own; the cochain-level checks (\code{certify\_transport}) are tests of the module-level assembly, a failure being a \code{HeckeInternalError}. In
every cone degree $j\le q_{\max}-1$ of the read they run on \textbf{the operator as applied} (\code{apply\_hecke}, whose values on the presentation's generators are the values read): (i) its blocks equal
the separately assembled $T'_{\mathcal T}$, $\eta_{\mathcal T}$, $T'_\partial$ (\code{bulk\_hecke}, \code{eta}, \code{bdry\_hecke}) --- $T^{\mathrm c}_{\mathcal T}(\alpha,0)=(T'_{\mathcal T}\alpha,\eta_{\mathcal T}\alpha)$ and
$T^{\mathrm c}_{\mathcal T}(0,c)=(0,T'_\partial c)$, the strict block included --- exactly modulo $\ell^{K}$; (ii) $\Dcone T^{\mathrm c}_{\mathcal T}g=0$ on the presentation's generators $g$ modulo $\ell^{K-\sigma_{\max}}$;
(iii) the identity of (4)(a) on the same $\alpha$, with the values $\eta_{\mathcal T}\alpha$, $T'_{\mathcal T}\alpha$ just compared in (i), in every cone degree $j\le\min(q_{\max}-1,3)$ --- all of them at $q_{\max}=4$; at
$q_{\max}=5$ the top cone degree $4$, which would need $\Phi_{\mathcal T}$ in degree $6$ and which the read does not use ((4)(e)(iv)), is \textbf{counted as skipped} (\code{eta\_skipped}, printed; (i) and (ii)
still run there), and an arm (Q5) of the standard validation compares $T_p$ on $\Hcpt{q}$, $H^{q}$, $H^{q}(\partial)$, $q\le5$, with D4's \code{BuildConeHecke(\dots,6)} through $H(\mathsf j_\Phi)$, $H(\Phi^{\ast}_{\mathcal T})$,
$H(\Phi^{\ast}_\partial)$. \textbf{The vectors} of (i) and (iii): for $D\le1300$ (\code{kEtaEntrywiseMaxD}) \textbf{every basis vector} of $C^{j+1}_{\mathcal T,\ell}$ and of $\mathsf B^{j}_\ell$ (entrywise); above,
the $\alpha$- and $c$-parts of the generators \textbf{and} $64$ random vectors (\code{kFreivaldsVectors}) of $C^{j+1}_{\mathcal T,\ell}$ and of $\mathsf B^{j}_\ell$, uniform in $\Z/\ell^{K}$ (rejection sampling
from splitmix64), seeded deterministically by $(\ell,N,n,j)$.
\emph{What is deterministic, what is probabilistic.} \textbf{Deterministic}: (1)--(4) and (4)(e); the chain-level certificates; (ii); and (i), (iii) on the basis vectors and on the generators. Entrywise ($D\le1300$),
(iii) is the identity of (4)(a) on all of $C^{j+1}_{\mathcal T,\ell}$, which with (i) (the applied operator \emph{is} the block operator) and $T'_{\mathcal T}$, $T'_\partial$ cochain maps (composites of certified ones)
is equivalent to the applied operator commuting with $\Dcone$ on cone degree $j$: both halves of (4)(e)(iii) are checked at cochain level. \textbf{Probabilistic}: (i), (iii) on the random vectors above
$D=1300$, a Freivalds test. For a non-zero error matrix $E$ over $\Z/\ell^{K}$ write $E=\ell^{s}E'$, $E'\not\equiv0\bmod\ell$, $s<K$; for $v$ uniform, $Ev\equiv0$ needs $E'v\equiv0\bmod\ell$, so
$\Pr[Ev=0]\le1/\ell$, and one check on $64$ independent vectors misses with probability $\le\ell^{-64}\le2^{-64}$. The bound is over the seed, for an error not depending on it (the stream taken as uniform); the
seed is fixed per $(\ell,N,n,j)$, so a re-run repeats the same vectors.
\emph{What ``on the generators'' certifies} (the C-D8b referee's m1, N1, N5). (i) on the generators compares the values read with an independent assembly: it catches an error of the applied operator not shared by
the separately assembled maps, whether or not its values are cocycles --- such an error changes the read only by being non-zero on a generator. It cannot see an error in a primitive shared by both sides. (iii) on
the generators checks the identity, both terms, on their $\alpha$-parts exactly modulo $\ell^{K}$; these are $d$-cocycles modulo $\ell^{K-\sigma_{\max}}$ only, and to that precision, given (i) and
$T'_\partial$ a cochain map, it is the second component of (ii) on $g$: nothing beyond (i) and (ii). So an off-diagonal error shared by the applied operator and $\eta_{\mathcal T}$, with $d_{\mathsf B}$-cocycle values,
passes (i), (ii) and (iii) on the generators and can change the read (the C-D8b referee's Arm (G)); the identity on the random vectors sees it, through its term $\eta^{j+1}_{\mathcal T}d\alpha$. The validation pins
one such plant (\code{ScopedPlantedCocycleEta}, $\ell=2$, $(N,n)=(11,0)$, $T_3$, $q=2$): clean on the generators alone, with the read of $\Hcpt{3}$ changed; caught with the random vectors; caught entrywise.
\textbf{Above $D=1300$ the coboundary half of (4)(e)(iii) is not checked deterministically at cochain level}: the read takes $T^{\mathrm c}_{\mathcal T}$ on the generators only, its proof is (4)(e)(iii) with the
chain-level certificates, and the identity on the random vectors tests it with the bound above.
\emph{Why this is admissible.} No value is read from a check, from the separately assembled maps or from the random vectors: the values read are the images of the generators under the applied operator,
products of exact matrices ((5)(c)), computed deterministically; the precision rule (5) governs them --- the Smith read-out, the rank sources, the exactness modulo $\ell^{K}$ ---; and their correctness for every $D$
is a deterministic statement, (4)(e)(iii) with the chain-level certificates (exact, not weakening with $D$). The checks, deterministic or random, can only add a failure (\code{HeckeInternalError}); they never
change, select or refuse a value. A miss of the random vectors therefore leaves the read where the deterministic checks leave it. What is \emph{tested and not proved} at cochain level above $D=1300$ is the
module-level assembly of the off-diagonal block outside the generators' span; that is the part checked with probability of a miss $\le\ell^{-64}$. (Forming $\eta_{\mathcal T}$ as a matrix
$C^{q+1}_{\mathcal T,\ell}\to\mathsf B^{q}_\ell$ costs about $\rank C^{q+1}_{\mathcal T,\ell}$ ($\approx(2+\tfrac13)D$ at $\ell=2$) applications of $T_{\mathrm b}$ to $\Iso$-sized cochains per degree and per $p$, of the order
$p\cdot D^{3}$, that of the tree cone's own elimination; v1.23.0-pre.34, referee N5; measured at D5c: $\approx D^{2}$, the applications being block-monomial, $\approx15$~s per degree and per $p$ at D4's wall; v1.23.0-pre.36, measured at the merge (the report's \S9.6): (i) about doubles the
entrywise transport, $24$~s for $T_3$ at $D=1152$ against the groups' $\approx29$~s; the random-vector mode grows as $\approx D^{1.1}$, $2.5$~s at $D=1152$, extrapolated --- not measured --- to under $2\,\%$ of the
groups at $D\approx3000$--$6000$.)
Refusals exactly those of (6); (f) a validation with arms: groups against D4's pair on every cell
where both run and against \file{d4\_pair\_oracle.tsv}; $\mathsf j_\Psi\mathsf j_\Phi=\id$, $\mathsf j_\Phi\mathsf j_\Psi=\id$ on cohomology at a few cells; $T_p$ on the three terms through the tree equal to D4's through
$H(\mathsf j_\Phi)$, and the characteristic polynomial of $T_p$ on $\Hcpt{1}$ against the oracle's \code{HECKE} rows; negative controls, each restricted to where it is visible (v1.23.0-pre.34,
referee m3): $+\mathsf Y^{\ast}_\Phi$ in $\mathsf j_\Phi$; $1-T^{2}$ for $1-T$ in $\mathsf B$; $S$ for $S^{-1}$ in $\mathsf a$ \textbf{only where it can bite} (v1.23.0-pre.35): at chain level
over $\Q[\Ga]$ at $\ell\in\{0,2\}$ (the certificate of (1)(b)), or on the cochains at $\ell=2$ of a module on which $\mathbf c$ acts non-trivially --- for $M_N$, $n$ odd, modulo $2^{K}$, $K\ge2$ --- and
\textbf{asserted clean} (the same cochains) for $n$ even at every $\ell$ and at $\ell\ge3$ (Definition~\ref{def:treepair}(2)); the code pins it caught at $(N,n)=(1,5)$, $(3,1)$ and at chain level, clean at
$(1,4)$ and at $\ell=3$, $(2,2)$; and at $\ell=3$ instead the term $S^{-1}e_U\mathsf v^{U}_1$ of $\mathsf a^{(3)}(\mathsf z_0)$ dropped; $\eta_{\mathcal T}$ dropped \textbf{for an operator that is not a fibre map}
(for a fibre map the diagonal is a cochain map), detected on at least one cell per $\ell\in\{2,3\}$ for an operator for which $\eta_{\mathcal T}$ is not a cocycle --- $\theta^{\ast}$ is one; $\Phi^{(N)}_p$ at the
cells where this has been observed --- and never asked at $\ell=3$ with $n$ odd ($M^{-I}=0$, no $3$-local boundary, $\eta_{\mathcal T}=0$); each control prints its detection count, in the
manner of the \code{EFFECTIVE} convention (v1.23.0-pre.36: met by the production certificate (i) itself, the audit's m1 --- each fault line names the certificate that caught it, every operator fault caught
at $\ell=2$ and at $\ell=3$; \code{TREEPAIR\_FAULTS} caught$=19$ clean$=3$, the clean ones the two asserted-clean $S$-for-$S^{-1}$ controls and the generators-only mode of the Arm (G) plant); every refusal code once; and $\mathsf j_\Psi\mathsf j_\Phi=\id$ on cohomology at a cell with non-trivial $2$-torsion in $\Hcpt{\ge3}$, where the five-lemma caveat bites
(the referee's \S5); (g) RSS predicted per cell (rule~19) and the cost lines at D4's wall ($D\approx1060$--$1240$) and beyond, \textbf{including the cost of forming $\eta_{\mathcal T}$} ((e); its dominant
term, part of the RSS prediction), in a report \file{DESCENT5c\_report.md} (v1.23.0-pre.36: met --- \code{TreePairCochains::transport\_bytes} adds the transport's columns on the built layout to the
prediction before any operator is formed, refusing \code{[RssCeiling]} above $0.9\times$ the ceiling, the audit's N2).

\paragraph{Formalization note.} Definition~\ref{def:treepair}, Lemma~\ref{lem:treepaircyl} and Proposition~\ref{prop:treepair} are homological algebra (cones, cylinders, the comparison theorem relative to a
subcomplex) and identities in the group rings of $\Z/4$, $\Z/6$ and $\langle T,\mathbf c\rangle$; no \code{\textbackslash LEAN} marker.

%----------------------------------------------------------------------
\subsection{The Hecke algebra over \texorpdfstring{$\Z$}{Z}: structure, certificates, congruence modules}\label{sec:heckeZ}
%----------------------------------------------------------------------
% NEW at v1.25.0-pre.0 (hecke-extensions slot 24, HeckeZ package Z0; BRANCH PRE-RELEASE on s24-wip-heckez-z0, enters the main line
% as v1.25.0-pre.1 only after the owner has released v1.24.0, owner ruling Q8 of 2026-11-04).  Source: the HeckeZ plan
% docs/implementation/slot24/heckez/ (signed by the owner 2026-11-04, Q1-Q20 at the recommendation): 01-math-core.md v0.6
% (Def. 1.1 ... Prop. 5.6, Lemma 1.6, Remarks 4.5, 5.6'), 03-bockstein-tor.md v0.3 (B1, B7, B8, B8', B9, B11-B15, CB-1/2/3),
% 04-record-census.md v0.4 (sec. 1, record item 8), 06-work-packages.md v0.4 (Z0); audits A1, A3, P2a, P2b, P3 (all findings applied
% in the stage files before this text was written).  Statements appended at the end of their section: no number moves.

(v1.25.0-pre.0 on the branch \texttt{s24-wip-heckez-z0}, refereed A7 / A7b; \textbf{in the main line from v1.25.0-pre.1} (2026-11-06), a pre-release, not owner-released; the HeckeZ plan \file{docs/implementation/slot24/heckez/}, signed by the owner on 2026-11-04 with every
question at its recommendation --- in particular Q1(a), \emph{residue-only} big integers, Q6(a), \emph{the anemic algebra first at level $N>1$}, and
Q8(a), this text on its own branch; stages \file{01-math-core.md} v0.6, \file{03-bockstein-tor.md} v0.3, \file{04-record-census.md} v0.4, package Z0
of \file{06-work-packages.md} v0.4; the adversarial audits \file{A1-audit.md}, \file{A3-audit.md}, \file{P2a-final-math-audit.md},
\file{P3-reaudit.md}.) \S\ref{sec:descentN}--\S\ref{sec:treepair} certify the \emph{groups} $H^{q}(\Ga_0(N);\Vcoef)$, the characteristic polynomials
$\chi^{(N)}_{p,1}$ over $\Z$ and the torsion Hecke modules (Theorem~\ref{thm:descentcert}). This subsection fixes what these data, together with
reads modulo $\ell^{K}$ on the free part, say about the \emph{ring} generated by the Hecke operators over $\Z$: (i) the algebras and their coarse
structure --- the prime spectrum does not see the square-zero glue between the free part and the torsion (Definition~\ref{def:heckeZ},
Propositions~\ref{prop:heckeZspec}, \ref{prop:heckeZrat}, Lemma~\ref{lem:heckeZbridge}); (ii) how the local structure at a prime $\ell$ is read from
mod-$\ell^{K}$ data, and exactly when it can and cannot be (Propositions~\ref{prop:heckeZmax}--\ref{prop:heckeZthrough}); (iii) the universal
coefficient sequences with their Hecke action, their splitting, the level-one picture at $2$ and $3$, and the glue $J$, which computed data bound but
do not determine (Propositions~\ref{prop:heckeZuct}--\ref{prop:heckeZglue}, three named conjectures in Remarks~\ref{rmk:heckeZcbone}--\ref{rmk:heckeZcbthree});
(iv) the descent of the structure to $\Z[1/6]$ and to $2$, $3$ (Theorem~\ref{thm:heckeZdescent}); (v) the global certificate and what is reachable
with the shipped tools (Theorem~\ref{thm:heckeZcert}, Remark~\ref{rmk:heckeZcaveat}); (vi) congruence modules and their certificates, Ramanujan's
congruence modulo $691$ for every $m$ from one finite computation (Definition~\ref{def:heckeZcong}, Propositions~\ref{prop:heckeZcong},
\ref{prop:heckeZeiscert}). The record that carries the structure is item~8 of \S\ref{sec:recordschema}. \textbf{Not here}: $U_p$, $p\mid N$, over $\Z$
(Theorem~\ref{thm:descentcert}(5), (7)(c) assert nothing of it; its lift is Proposition~\ref{prop:strictliftN}(9), and whether a binary serves it is item~8~(d$'$)(v); v1.26.0-pre.11, confirm-referee C1); non-rational Eisenstein components and cusp--cusp congruence modules beyond bounds (the
plan's Prop.~5.7, a sketch); Hida families ([U]); $p$-adic $L$-functions (postponed by the owner).

\emph{What is used.} Theorem~\ref{thm:descentcert} and the statements it cites; Definition~\ref{def:satdata}, Propositions~\ref{prop:saturation},
\ref{prop:satdefect}, Corollary~\ref{cor:stablecert}, Proposition~\ref{prop:cohforms}; Definition~\ref{def:ringsemantics}(3); Proposition~\ref{prop:ltwothreescalar},
Lemmas~\ref{lem:tellrank}, \ref{lem:bocksteinallell}, Proposition~\ref{prop:heckeH0}, Claim~\ref{cl:bockstein}; Proposition~\ref{prop:torsionhecke};
[U,~\texttt{lem:u-fibre}], [U,~\texttt{def:u-congmod}]; \Espec{cl:sturm} (as corrected at [E] 1.24.0-pre.1). \emph{Standing assumptions.} \ref{H4},
\ref{B1}--\ref{B3}; $N\ge1$; $n\ge0$ \textbf{even} (for $n$ odd, $2$ kills $H^{\bullet}(\Ga_0(N);\Vcoef)$, the proof of Proposition~\ref{prop:crtcharpolyN}(2),
and everything below is empty); $k:=n+2$; $\Ga:=\Ga_0(N)$; $H^{q}:=H^{q}(\Ga;\Vcoef)$, group cohomology (Definition~\ref{def:ringsemantics}(3)).
\emph{Notation.} $\Talg$, $\Talg^{(q)}$, $\Talg^{\mathrm{an}}$, $\Talg^{\mathrm{fr}}$, $\Talg^{\mathrm{an,fr}}$, $\Talg^{\mathrm{tors}}$, $\Talg^{\mathrm{cusp}}$, $\Talg'$,
$\Talg_\ell$, $\Talg_{\Z_\ell}$, $\mathcal H_N$, $X_p$, $L$, $L_\ell$, $I_{\mathrm{fr}}$, $I_{\mathrm{tors}}$, $J$, $A_K(M)$, $A'_K(M)$, $d$, $d^{\mathrm{an}}$, $e_k$,
$e^{\mathrm{an}}_k(N)$, $\vartheta$, $f_\vartheta$, $f_i$, $g_i$, $P_i$, $\mathcal O_i$, $\operatorname{disc}_\vartheta$, $\mathsf S$, $\mathsf S_{\mathrm{cand}}$, $\mathsf c$,
$\mathcal S$, $\mathcal S'$, $I_{\mathcal S}$, $\Talg_{\mathcal S}$, $C_{\Talg}$, $\eta_L$, $L_{\mathcal S}$, $c_{\mathcal S',p}$, $g_{\mathcal P}$, $C_{\mathrm{ft}}$, and in
Proposition~\ref{prop:heckeZeiscert}(3) $\nu$, $\eta_q$, $e^{\mathrm{cusp}}_\midl$ (the rank of the cuspidal block at $\midl$, not an exponent), $\mathsf h_q$, $\mathsf F_q$,
$f_{q,\midl}$ (v1.25.0-pre.0b, A7-m7) are new symbols of this subsection and mean nothing else (rule~\ref{A5}). In particular $\vartheta$ is a primitive element, \textbf{not} the operator $\theta$ of
\S\ref{sec:dicksontheta}; $\eta_L$ (always with its subscript) is a congruence module, not the pair homotopy $\etapair$; $d$ is a rank, not the
$d=\rank_\Z M$ of \S\ref{sec:descentpair}; $\mathcal P$ is a finite set of primes, as in Theorem~\ref{thm:descentcert}; $\Talg$ here is the Hecke algebra of
Definition~\ref{def:heckeZ}, and Definition~\ref{def:satdata}'s general pair is written $(M,\Talg_\ell)$ when it is applied. Statement numbers of the
plan (Prop.~2.2, B8, \dots) are those of \file{01-math-core.md} and \file{03-bockstein-tor.md}.

\begin{definition}[The Hecke algebras over $\Z$]\label{def:heckeZ}
Let $\mathcal H_N:=\Z[X_p:p\text{ prime}]$ act on $H^{q}$ by $X_p\mapsto T_p$ for $p\nmid N$ (Propositions~\ref{prop:strictliftN}, \ref{prop:heilbronn}; $p=\ell$
included) and $X_p\mapsto U_p$ for $p\mid N$ (on cohomology, Proposition~\ref{prop:descentmaps}(3)(b$'$); over $\Z$ the lift of Proposition~\ref{prop:strictliftN}(9), v1.26.0-pre.1, served by a binary or not according to item~8~(d$'$)(v) (v1.26.0-pre.11, confirm-referee C1); Theorem~\ref{thm:descentcert}(5),
(7)(c) assert nothing of it); $\mathcal H^{\mathrm{an}}_N:=\Z[X_p:p\nmid N]$.
\begin{enumerate}
\item $\Talg^{(q)}:=\im(\mathcal H_N\to\End_\Z H^{q})$, and $\Talg:=\Talg^{(1)}$, the \textbf{full} Hecke algebra; $\Talg^{\mathrm{an}}$, the \textbf{anemic} algebra, is the
  image of $\mathcal H^{\mathrm{an}}_N$. At $N=1$ they agree.
\item $L:=H^{1}/\mathrm{tors}=\Hfr{1}(N)$; $\Talg^{\mathrm{fr}}:=\im(\Talg\to\End L)$, $\Talg^{\mathrm{tors}}:=\im(\Talg\to\End(\mathrm{tors}\,H^{1}))$, with kernels
  $I_{\mathrm{fr}}$, $I_{\mathrm{tors}}$, and $J:=I_{\mathrm{fr}}\cap I_{\mathrm{tors}}$ (\textbf{the glue}); $\Talg^{\mathrm{an,fr}}:=\im(\Talg^{\mathrm{an}}\to\End L)\subseteq\Talg^{\mathrm{fr}}$.
  Everything stated below for $\Talg^{\mathrm{fr}}$ and $d$ holds verbatim for $\Talg^{\mathrm{an,fr}}$ and $d^{\mathrm{an}}$ (Proposition~\ref{prop:heckeZrat}(3$'$)), except
  where Propositions~\ref{prop:heckeZcong}(5) and~\ref{prop:heckeZeiscert} say otherwise.
\item $\Talg^{\mathrm{cusp}}$: the image of $\Talg^{\mathrm{fr}}$ on the cuspidal part $L\cap(L\otimes\Q)_{\mathrm{cusp}}$ (equivalently on $\Hcpt{1}$'s image
  $H^{1}_{!}$ modulo torsion, \S\ref{sec:descentpair}; the two agree by Proposition~\ref{prop:heckeZrat}(1)).
\item For a prime $\ell$: $L_\ell:=L\otimes\Z_\ell$, $\Talg_\ell:=\Talg^{\mathrm{fr}}\otimes\Z_\ell$ (acting on $L_\ell$ only), $\Talg_{\Z_\ell}:=\Talg\otimes\Z_\ell$ (acting on
  $H^{1}_\ell:=H^{1}\otimes\Z_\ell$, torsion included); the two are never identified. For a finitely generated $\Talg$-module $M$ and $K\ge1$,
  $A_K(M):=\im(\Talg\to\End(M/\ell^{K}M))$; for the subalgebra $\Talg'\subseteq\Talg^{\mathrm{fr}}$ generated by the operators actually computed --- $X_p$,
  $p\in\mathcal P$, and integer combinations of them --- $A'_K(M)$ is the image of $\Talg'$. Without an argument, $A_K:=A_K(L_\ell)$; it is
  Definition~\ref{def:satdata}'s $\mathsf A_K$ for the pair $(L_\ell,\Talg_\ell)$.
\end{enumerate}
\end{definition}

\noindent
Definition~\ref{def:satdata} and Propositions~\ref{prop:saturation}, \ref{prop:satdefect} are stated for a general pair (a free $\Z_\ell$-module of finite rank and a
commutative subalgebra of its endomorphisms), and their proofs use nothing else; the level-one setting of \S\ref{sec:saturation} is that section's
\emph{application}. They are applied below at level $N$, at $\ell\in\{2,3\}$ and to subalgebras $\Talg'_\ell$ (the plan's A1-m10). The census algebra of
Proposition~\ref{prop:cohforms}(4) is $A'_1(H^{1})$, not $A_1(L)$. \emph{Other degrees}: for $n>0$, $H^{0}=0$; for $n=0$, $H^{0}=\Z$ with $T_p=p+1$,
$U_p=p$; for $q\ge2$, $H^{q}$ is finite, supported at $2$ and $3$ and periodic (Theorem~\ref{thm:descentcert}(1), (6)), so $\Talg^{(q)}$, $q\ne1$, is a finite
algebra at $2$ and $3$, and nothing beyond Proposition~\ref{prop:heckeZspec}(1) is asserted about it.

\begin{proposition}[Coarse structure]\label{prop:heckeZspec}
\PROVED{} (the plan's Prop.~1.2.)
\begin{enumerate}
\item[(0)] The operators $X_p$ commute on $H^{q}$: at level $1$ by Proposition~\ref{prop:relationsmodel}; at level $N$ \STD{\cite[Ch.~3]{shimura}} (double cosets of
  $\Ga_0(N)$ in its Hecke semigroup). So $\Talg^{(q)}$ is commutative.
\item $\Talg$ is a finite $\Z$-algebra; $\Talg^{\mathrm{fr}}$ is free of finite rank $d:=\rank_\Z\Talg^{\mathrm{fr}}$; $\Talg^{\mathrm{tors}}$ is finite, supported at the torsion primes
  $\{2,3\}\cup\mathcal B(M_N)$ (Theorem~\ref{thm:descentcert}(1)).
\item $I_{\mathrm{fr}}I_{\mathrm{tors}}=0$; hence $J^{2}=0$, $\Talg/J\hookrightarrow\Talg^{\mathrm{fr}}\times\Talg^{\mathrm{tors}}$, and
  $\operatorname{Spec}\Talg=V(I_{\mathrm{fr}})\cup V(I_{\mathrm{tors}})=\operatorname{Spec}\Talg^{\mathrm{fr}}\cup\operatorname{Spec}\Talg^{\mathrm{tors}}$. \textbf{$\Talg\to\Talg/J$ induces a
  homeomorphism of spectra with equal residue fields}: the glue plays no role for prime ideals and components.
\item Every prime $P$ with $P\cap\Z=0$ contains $I_{\mathrm{fr}}$. The minimal primes of $\Talg$ are (a) the primes of characteristic $0$ --- the minimal primes of
  $\Talg^{\mathrm{fr}}$, i.e.\ the primes of $\Talg\otimes\Q$, the \textbf{components} --- and (b) the \textbf{torsion-only points}: the maximal ideals
  $\midl\supseteq I_{\mathrm{tors}}$ with $\midl\not\supseteq I_{\mathrm{fr}}$. Each of the latter is minimal and maximal, an isolated point, i.e.\ a finite local factor of $\Talg$.
\end{enumerate}
\end{proposition}

\begin{proof}
(0) is cited. (1) $\End_\Z$ of a finitely generated abelian group is a finitely generated $\Z$-module; a submodule of $\End(L)$ is free; $\Talg^{\mathrm{tors}}\subseteq
\End(\mathrm{tors}\,H^{1})$ is finite. (2) For $a\in I_{\mathrm{fr}}$, $aH^{1}\subseteq\mathrm{tors}$, and $b\in I_{\mathrm{tors}}$ kills $\mathrm{tors}$, so $ba=0$. Then
$V(I_{\mathrm{fr}})\cup V(I_{\mathrm{tors}})=V(I_{\mathrm{fr}}I_{\mathrm{tors}})=V(0)$; $J$ is nilpotent, so it lies in every prime, and $\Talg\to\Talg/J$ is a homeomorphism on spectra
with equal residue fields. (3) $m:=|\Talg^{\mathrm{tors}}|$ has $m\cdot1\in I_{\mathrm{tors}}$; if $P\cap\Z=0$ then $m\notin P$, so $P\not\supseteq I_{\mathrm{tors}}$ and by (2)
$P\supseteq I_{\mathrm{fr}}$. Let $\midl$ be as in (b) and $P\subseteq\midl$ prime. $P\supseteq I_{\mathrm{fr}}$ would give $\midl\supseteq I_{\mathrm{fr}}$; so
$P\supseteq I_{\mathrm{tors}}$, and $P$ is a prime of the finite ring $\Talg^{\mathrm{tors}}$, hence maximal: $P=\midl$. The primes of (a) are minimal: a prime below one of
them has characteristic $0$, contains $I_{\mathrm{fr}}$, and the primes of the finite $\Q$-algebra $\Talg^{\mathrm{fr}}\otimes\Q$ are pairwise incomparable. Conversely (A7-m2), let
$P$ be minimal of characteristic $\ell$. If $P\supseteq I_{\mathrm{fr}}$, then $P$ is a prime of $\Talg^{\mathrm{fr}}$ over $\ell$; $\Talg^{\mathrm{fr}}$ is $\Z$-free, hence flat, so
going-down gives a prime of $\Talg^{\mathrm{fr}}$ over $0$ strictly inside $P$, against minimality. So $P\not\supseteq I_{\mathrm{fr}}$, hence $P\supseteq I_{\mathrm{tors}}$ by (2), and
$P$ is a maximal ideal as in (b).
\end{proof}

\begin{proposition}[Rational structure; reducedness; the generation inputs]\label{prop:heckeZrat}
\PROVED{} given the cited standard statements (the plan's Props.~1.3, 1.4, 1.5(2)--(4)); (4)(c) \CONJ{} for $k>2$.
\begin{enumerate}
\item $\Talg^{\mathrm{fr}}$ is the subring $\Z[X_p]$ of $\Talg_\Q:=\Talg\otimes\Q\subseteq\End(L\otimes\Q)$ generated by the operators. It does not depend on the
  lattice: on every $\Talg$-stable lattice $M\subset L\otimes\Q$ the image of $\mathcal H_N$ is the same subring of $\Talg_\Q$. (Saturation,
  Proposition~\ref{prop:heckeZgen}, does depend on the lattice.)
\item \STD{\cite[Thm.~12.2.2, Prop.~12.4.10]{diamond_im}, as in Lemma~\ref{lem:heckerankone}; the compatibility with $U_p$, $p\mid N$, is the verified [GT] row
  \texttt{lit:diamond-im-es}} $H^{1}\otimes\Cx\cong M_k(\Ga)\oplus\overline{S_k(\Ga)}$ as $\mathcal H_N$-modules. Hence $\Talg_\Q\cong\Talg(M_k(\Ga))_\Q$; its systems are
  those of $M_k(\Ga)$, Eisenstein ones included; a cuspidal system occurs in $H^{1}\otimes\overline\Q$ twice, an Eisenstein one once (for the full algebra, when
  $\Talg_\Q$ is reduced).
\item $d=\dim M_k(\Ga)=(b_1+e_k)/2$, $b_1=\rank H^{1}$ (the descent record), $e_k=\dim\mathrm{Eis}_k(\Ga)$: the number $c(N)$ of cusps for $k\ge4$, $c(N)-1$ for $k=2$.
\item[(3$'$)] \STD{\cite[Thm.~6.2.3, Rmks.~6.3.3--6.3.4]{diamond_im} (strong multiplicity one and the old-space decomposition for the cuspidal part),
  \cite[Thms.~5.8--5.10]{stein} (the basis $E_k^{\psi,\bar\psi}(q^{t})$ of $\mathrm{Eis}_k(\Ga)$ for the rest); quoted in the verified [GT] rows \texttt{lit:hz-es-anemic},
  \texttt{lit:hz-eis-dim} (v1.25.0-pre.0b, A7-m5: in place of \cite{diamond_shurman}, not in the library, and \cite{atkin_lehner}, unquoted)} The anemic algebra has rank
  \[
    d^{\mathrm{an}}:=\rank_\Z\Talg^{\mathrm{an,fr}}=\sum_{M\mid N}\dim S_k^{\mathrm{new}}(\Ga_0(M))+e^{\mathrm{an}}_k(N),
  \]
  \[
    e^{\mathrm{an}}_k(N):=\#\{\psi\text{ primitive}:\mathrm{cond}(\psi)^{2}\mid N\}-[k=2,\,N=1].
  \]
  For $(N,k)\ne(1,2)$ (where $e^{\mathrm{an}}_2(1)=0$, A7-N9), $e^{\mathrm{an}}_k(N)=1$ iff $16\nmid N$ and $p^{2}\nmid N$ for every odd $p$; then the anemic Eisenstein part is \textbf{one rational system}
  $T_p\mapsto1+p^{k-1}$, carried by several Eisenstein series at once. $d^{\mathrm{an}}=d$ iff there are no old forms and every anemic Eisenstein system is carried by
  one series (always at $N=1$).
\item \emph{Reducedness} (a predictor only: every cell checks it, Theorem~\ref{thm:heckeZcert}(C1)). (a) \STD{Petersson normality and the explicit Eisenstein basis:
  \cite[\S6]{diamond_im}, \cite[Thms.~5.9, 5.10]{stein}, the verified [GT] row \texttt{lit:hz-semisimple-an}} $\Talg^{\mathrm{an}}_\Q$ is reduced. (b) \STD{the old-space
  decomposition \cite[Rmk.~6.3.4]{diamond_im}; \cite{atkin_lehner}\VERIFYloc} If $p^{3}\mid N$ for some $p$, the full $\Talg_\Q$ is \textbf{not}
  reduced. (c) If $N$ is cube-free and $\alpha_p\ne\beta_p$ (the roots of $x^{2}-a_px+p^{k-1}$) for every $p\mid N$ and every newform of level $M\mid N$, $p\nmid M$,
  then $\Talg_\Q$ is reduced; the hypothesis $\alpha_p\ne\beta_p$ is known for $k=2$ \STD{\cite{coleman_edixhoven}\VERIFYloc} and is \CONJ{} for $k>2$.
\item \emph{Generation inputs.} (a) \STD{\cite{sturm}; \cite[Thm.~9.23]{stein}; the integral duality $\Talg^{\mathrm{cusp}}\times S_k(\Z)\to\Z$,
  $(t,f)\mapsto a_1(tf)$, \cite[Prop.~12.4.13]{diamond_im}; quoted in the verified [GT] row \texttt{lit:hz-sturm-integral} (A7-N2)} $\Talg^{\mathrm{cusp}}$ is spanned \emph{as a $\Z$-module} by the $T_n$, $n\le B_N:=k\mu_N/12$,
  $\mu_N:=[\SL_2(\Z):\Ga_0(N)]$ (the $T_n$ with $p\mid\gcd(n,N)$ are words in $U_p$). (b) For $M_k$ the $\Z$-module statement is false at primes $\ge5$
  (\Espec{cl:sturm}(b)); that $\Talg(M_k)_\Q$ is spanned by $T_m$, $m\le\lfloor k/12\rfloor+1$, is \CONJ{} (\Espec{cl:sturm}(a) as corrected at [E] 1.24.0-pre.1)
  and \textbf{is used nowhere here}. (c) Let $\Talg'\subseteq\Talg''$ be commutative rings and $M$ a faithful $\Talg''$-module. \emph{Cyclic-faithful}: if $M=\Talg'v$, then
  $\Talg'=\Talg''$ as images in $\End M$. \emph{Commutant}: if $\End_{\Talg'}(M)$ is the image of $\Talg'$, then $\Talg''=\Talg'$.
\end{enumerate}
\end{proposition}

\begin{proof}
(1) $\mathcal H_N\to\End(L)\subset\End(L\otimes\Q)$ and $\mathcal H_N\to\End(M)\subset\End(L\otimes\Q)$ have the same image as $\mathcal H_N\to\End(L\otimes\Q)$. (2) is cited.
(3) By (2), $\Talg_\Q$ acts faithfully on $M_k(\Ga)$, and the action on $\overline{S_k}$ factors through it. The pairing $\Talg(M_k)_\Q\times M_k(\Q)\to\Q$,
$(t,f)\mapsto a_1(tf)$, is perfect: if $a_1(tf)=0$ for all $t$ then $a_m(f)=a_1(T_mf)=0$ for $m\ge1$, so $f$ is constant, and $f=0$ as $k>0$; if $a_1(tf)=0$ for all
$f$ then $a_m(tf)=a_1(tT_mf)=0$ for all $f$ and $m$, so $t=0$. So $d=\dim M_k$, and $b_1=\dim M_k+\dim S_k$, $\dim M_k=\dim S_k+e_k$. No reducedness is used. (3$'$)
$\Talg^{\mathrm{an}}_\Q$ is reduced ((4)(a)) and acts faithfully on $M_k(\Ga)$, so its dimension is the number of distinct anemic systems of $M_k(\Ga)$ over $\overline\Q$.
Cuspidal: every eigenform lies in the old space of a unique newform of level $M\mid N$, the old space shares its anemic system, and distinct newforms of levels dividing $N$
have distinct anemic systems (strong multiplicity one). Eisenstein: $\mathrm{Eis}_k(\Ga_0(N))$ has the basis $E_k^{\psi,\bar\psi}(q^{t})$, $\mathrm{cond}(\psi)^{2}t\mid N$ (for
$k=2$, $\psi=1$: the combinations $E_2(q)-tE_2(q^{t})$, $t>1$), with anemic system $T_p\mapsto\psi(p)+\bar\psi(p)p^{k-1}$, independent of $t$; distinct $\psi$ give distinct
systems (the semisimple representations $\psi\oplus\bar\psi\varepsilon^{k-1}$, $\varepsilon$ the cyclotomic character, are distinct, and are determined by their traces
at almost all primes), and no Eisenstein system is cuspidal (\cite{deligne_weil1}: at any prime $p\ge7$, $p\nmid N$,
$|\psi(p)+\bar\psi(p)p^{k-1}|\ge p^{k-1}-1>2p^{(k-1)/2}\ge|a_p|$; A7-m1: the bound $1+p^{k-1}$ is the trivial-$\psi$ eigenvalue only, and at $k=2$, $p\le5$ the inequality fails). (4) $\Talg_\Q$ is reduced iff each generator is
semisimple (commuting semisimple operators generate a reduced algebra; conversely a reduced finite $\Q$-algebra is a product of number fields, whose elements have
separable minimal polynomials). (a) The $T_p$, $p\nmid N$, are normal for the Petersson product on $S_k$ and diagonal on the Eisenstein basis. (b), (c) On the old
space of a newform $f$ of level $M$, $N=Mp^{r}\cdot(\text{prime to }p)$, $U_p$ acts on $v_j:=f(q^{p^{j}})$, $0\le j\le r$, by $U_pv_j=v_{j-1}$ ($j\ge1$),
$U_pv_0=a_pv_0-p^{k-1}v_1$ ($p\nmid M$), $U_pv_0=a_pv_0$ ($p\mid M$). $v_r$ is cyclic, so the minimal polynomial is the characteristic polynomial:
$x^{r-1}(x^{2}-a_px+p^{k-1})$ ($p\nmid M$), $x^{r}(x-a_p)$ with $a_p\ne0$ ($p\,\|\,M$), $x^{r+1}$ ($p^{2}\mid M$, $a_p=0$). So $U_p$ fails to be semisimple iff
$v_p(N)\ge3$ in each case, or $r\ge1$ and $\alpha_p=\beta_p$. The Eisenstein old spaces behave the same way: for $k\ge4$, $\alpha_p=\psi(p)\ne\bar\psi(p)p^{k-1}=\beta_p$
(different absolute values); for $k=2$ the functional $v_j\mapsto p^{-j}$ on $\mathrm{span}\{E_2(q^{p^{j}})\}$ is a $U_p$-eigenvector of eigenvalue $p$ cutting out the modular
subspace, so the generalized $0$-eigenspace, one Jordan block of size $r-1$, lies in that subspace. At $p^{3}\mid N$ the level-$1$ series $E_k$ ($k\ge4$) or the
$E_2$-family ($k=2$) contributes a nilpotent block of size $\ge2$, which is (b). (5)(a) By the duality, $\{T_n:n\le B_N\}$ spans $\Talg^{\mathrm{cusp}}/\ell$ iff every
$f\in S_k(\Z)$ with $a_n(f)\equiv0\pmod\ell$ for $n\le B_N$ is $\equiv0\pmod\ell$, which is Sturm's theorem modulo $\ell$ (the form he proved; the characteristic-$0$
vanishing statement would give only full rank). So $\Sigma+\ell\Talg^{\mathrm{cusp}}=\Talg^{\mathrm{cusp}}$ for $\Sigma:=\sum_{n\le B_N}\Z T_n$ and every $\ell$; $\Sigma$ has
finite index prime to every prime, hence $\Sigma=\Talg^{\mathrm{cusp}}$. (b) is cited. (c) Cyclic-faithful: for $t\in\Talg''$ write $tv=t'v$, $t'\in\Talg'$; then
$t(sv)=s(tv)=s(t'v)=t'(sv)$ for $s\in\Talg'$, so $t-t'$ kills $M$. Commutant: $\Talg''$ commutes with $\Talg'$.
\end{proof}

\noindent
Checks (the plan's \file{s01r/dan.py}): $(N,n)=(1,10)$: $b_1=3$, $c(1)=1$, $d=d^{\mathrm{an}}=2$; $(2,6)$: $d=3$, $d^{\mathrm{an}}=1+1=2$; $(8,10)$: $d=13=\dim M_{12}(\Ga_0(8))$,
$d^{\mathrm{an}}=(1+0+1+3)+1=6$. \textbf{The rank is an a-priori input of every certificate}: $d$ for the full algebra, $d^{\mathrm{an}}$ for the anemic one, never the computed
$\deg f_\vartheta$ --- an operator that fails to separate two systems is otherwise indistinguishable from one that separates all (the plan's A1-M3). The
cyclic-faithful and commutant certificates exist in code over $\Fl$ (\code{HeckeExt/HeckeGenerationCertificates}); on a direct summand of $L$ at odd $\ell$ cyclic-faithful
implies saturation (Proposition~\ref{prop:saturation}(2), sufficient condition), so for $\Talg^{\mathrm{fr}}$ Proposition~\ref{prop:heckeZgen} subsumes it; it matters for
$\Talg^{\mathrm{tors}}$ and at $\ell=2$.

\begin{lemma}[The integral/$\ell$-adic bridge]\label{lem:heckeZbridge}
\PROVED{} (the plan's Lemma~1.6, HZ-0.) Let $L$ be a free $\Z$-module of finite rank, $\Talg^{\mathrm{fr}}\subseteq\End_\Z L$ the subalgebra generated by a family of operators,
$\Talg'\subseteq\Talg^{\mathrm{fr}}$ a subalgebra, $\ell$ a prime, $L_\ell$, $\Talg_\ell$ as in Definition~\ref{def:heckeZ}(4), $\Talg'_\ell:=\Talg'\otimes\Z_\ell$.
\begin{enumerate}
\item $\Talg_\ell\to\End_{\Z_\ell}(L_\ell)$ is injective, with image the $\Z_\ell$-subalgebra generated by the base-changed operators; likewise for $\Talg'$.
\item $\rank_{\Z_\ell}\Talg_\ell=\rank_\Z\Talg^{\mathrm{fr}}$; so the $d$ of Proposition~\ref{prop:heckeZrat}(3) is the rank the local statements compare with.
\item If $\rank\Talg'=\rank\Talg^{\mathrm{fr}}$, then $\ell\nmid[\Talg^{\mathrm{fr}}:\Talg']$ iff $\Talg'_\ell=\Talg_\ell$. Hence $\Talg'=\Talg^{\mathrm{fr}}$ iff $\Talg'_\ell=\Talg_\ell$ for
  every $\ell$.
\end{enumerate}
\end{lemma}

\begin{proof}
$\Z_\ell$ is flat over $\Z$, so $\Talg^{\mathrm{fr}}\hookrightarrow\End_\Z L$ stays injective after $\otimes\Z_\ell$, and $\End_\Z(L)\otimes\Z_\ell=\End_{\Z_\ell}(L_\ell)$ ($L$
finite free); the image is the $\Z_\ell$-span of the words in the operators, which is (1). (2) is base change of a free module. (3) Tensoring
$0\to\Talg'\to\Talg^{\mathrm{fr}}\to\Talg^{\mathrm{fr}}/\Talg'\to0$ with $\Z_\ell$ gives $\Talg_\ell/\Talg'_\ell\cong(\Talg^{\mathrm{fr}}/\Talg')\otimes\Z_\ell$, the $\ell$-primary
part of a finite group.
\end{proof}

\begin{proposition}[Maximal ideals are detected modulo $\ell$]\label{prop:heckeZmax}
\PROVED{} \STD{Cayley--Hamilton for finitely generated modules, \cite[Prop.~2.4]{atiyah_macdonald}} (the plan's Prop.~2.1.) Let $R$ be a commutative ring, $M$ a
finitely generated faithful $R$-module, $\ell\in R$, $K\ge1$ and $I_K:=\ker(R\to\End(M/\ell^{K}M))$. Then $\ell^{K}R\subseteq I_K\subseteq\sqrt{\ell R}$. Hence
$\operatorname{Spec}(R/I_K)\to\operatorname{Spec}R$ is a homeomorphism onto $V(\ell)$ with equal residue fields, for every $K\ge1$.
\end{proposition}

\begin{proof}
For $x\in I_K$, $xM\subseteq\ell^{K}M\subseteq(\ell R)M$, so Cayley--Hamilton gives $x^{m}+a_1x^{m-1}+\dots+a_m\in\Ann(M)=0$ with $a_i\in(\ell R)^{i}$; hence $x^{m}\in\ell R$.
\end{proof}

\noindent
\emph{How it is applied, and its limits.} With $R=\Talg\otimes\Z_{(\ell)}$, $M=H^{1}\otimes\Z_{(\ell)}$ (faithful by flat base change; $H^{1}\otimes\Z_{(\ell)}$ is \emph{not}
faithful over $\Talg$ itself, an operator mapping $H^{1}$ into its prime-to-$\ell$ torsion acting as $0$), and with $R=\Talg^{\mathrm{fr}}\otimes\Z_{(\ell)}$, $M=L\otimes\Z_{(\ell)}$,
which detects $V(\ell)\subset\operatorname{Spec}\Talg^{\mathrm{fr}}$ and misses the torsion-only points. \textbf{A computation produces $A'_K(M)$, not $A_K(M)$}:
$\operatorname{Spec}A'_K$ is the $\ell$-fibre of $\operatorname{Spec}\Talg'$, onto which $V(\ell)\subset\operatorname{Spec}\Talg$ maps surjectively, not necessarily injectively.
The two agree under \emph{generation modulo $\ell$}, $A'_1(M)=A_1(M)$, which Proposition~\ref{prop:heckeZgen} certifies. So ``the census lists exactly the maximal
ideals of $\Talg$ over $\ell$'' holds under the census's generation certificate (Corollary~\ref{cor:stablecert}(2)), not unconditionally; two maximal ideals that
differ on a computed operator are certainly distinct. At $\ell\ge5$, $H^{1}(\Ga;\Vcoef/\ell)=H^{1}/\ell$ (Proposition~\ref{prop:heckeZuct}(4)). At $\ell\in\{2,3\}$,
$H^{2}[\ell]$ is a \emph{quotient} of $H^{1}(\Ga;\Vcoef/\ell)$, not a submodule, and the census algebra there can carry maximal ideals of $\Talg^{(2)}$ that are not
maximal ideals of $\Talg^{(1)}$; the filter is applied on the integral model (Proposition~\ref{prop:heckeZuct}(4)).

\begin{proposition}[Generation by saturation]\label{prop:heckeZgen}
\PROVED{} (the plan's Prop.~2.2; it strengthens Corollary~\ref{cor:stablecert}(2).) Let $M$ be a $\Talg^{\mathrm{fr}}$-stable lattice in a $\Talg$-stable subspace of $L\otimes\Q_\ell$
on which $\Talg_\ell$ acts faithfully (e.g.\ $M=L_\ell$, or one $\inv$-part when every system occurs in it; A7-N6), and $d=\rank\Talg^{\mathrm{fr}}$ from Proposition~\ref{prop:heckeZrat}(3) ---
\textbf{an a-priori input}; for the anemic algebra read $\Talg^{\mathrm{an,fr}}$ and $d^{\mathrm{an}}$ of Proposition~\ref{prop:heckeZrat}(3$'$); a rank read off a
computed $\deg f_\vartheta$ is not admissible here.
\begin{enumerate}
\item $\dim_{\Fl}A'_1(M)\le\rank\Talg'_\ell\le d$. If $\dim_{\Fl}A'_1(M)=d$, then $\Talg'_\ell=\Talg_\ell$, $\Talg_\ell$ is saturated in $\End(M)$, and $\Talg_\ell/\ell=A'_1(M)$.
\item Conversely, if $\Talg_\ell$ is saturated in $\End M$ and $\Talg'+\ell\Talg_\ell=\Talg_\ell$ (true for a large enough finite set of operators), then
  $\dim A'_1(M)=d$. So the test succeeds for some finite operator set \textbf{iff} $\Talg_\ell$ is saturated on $M$.
\item If $\Talg_\ell$ is not saturated on $M$, no operator set passes. If $\rank\Talg'=d$, mod-$\ell^{K}$ data on $M$ then give only
  $\Talg'_\ell\subseteq\Talg_\ell\subseteq\Talg'^{\,\mathrm{sat}}:=(\Talg'_\Q\otimes\Q_\ell)\cap\End(M)$, with $\mathrm{length}(\Talg'^{\,\mathrm{sat}}/\Talg'_\ell)$ read by
  Proposition~\ref{prop:satdefect}.
\end{enumerate}
\end{proposition}

\begin{proof}
(1) The first inequality is Proposition~\ref{prop:saturation}(3) at $K=1$ for the pair $(M,\Talg'_\ell)$; the second holds as $\Talg'\subseteq\Talg^{\mathrm{fr}}$
(Lemma~\ref{lem:heckeZbridge}(2)). Equality forces $\rank\Talg'=d$, so $\Talg'_\Q=\Talg_\Q$, and it forces $\Talg'_\ell$ saturated (Proposition~\ref{prop:saturation}(2),
(a)$\Leftrightarrow$(c)), i.e.\ $\Talg'_\ell=(\Talg_\Q\otimes\Q_\ell)\cap\End M\supseteq\Talg_\ell\supseteq\Talg'_\ell$ ($\Talg_\ell\subseteq\End M$ as $M$ is
$\Talg^{\mathrm{fr}}$-stable). (2) Nakayama over $\Z_\ell$ gives $\Talg'_\ell=\Talg_\ell$, and saturation gives $\dim A_1(M)=d$ (Proposition~\ref{prop:saturation}(2)(a)).
$\Talg^{\mathrm{fr}}$ is a finitely generated $\Z$-module, spanned by finitely many words in finitely many $X_p$. (3) A pass would make $\Talg_\ell=\Talg'_\ell$ saturated,
by (1). With $\rank\Talg'=d$, $\Talg'_\Q\otimes\Q_\ell=\Talg_\Q\otimes\Q_\ell$, so $\Talg_\ell\subseteq\Talg'^{\,\mathrm{sat}}$, the saturation of the pair $(M,\Talg'_\ell)$.
\end{proof}

\noindent
The hypothesis is $\dim A'_1=d$ with $d$ from Proposition~\ref{prop:heckeZrat}, not $\dim A'_1=\rank\Talg'$: the latter makes $\Talg'$ saturated and says nothing about
$\Talg'_\Q=\Talg_\Q$. Saturation holds at every multiplicity-one maximal ideal at odd $\ell$ (Proposition~\ref{prop:cohforms}(5)); census rows with positive
\code{sat\_defect\_dim} show that it fails in practice. \emph{What saturation means} (not used by any certificate). $(\Talg_\Q\otimes\Q_\ell)\cap\End M$ is the
multiplier ring $\{x:xM\subseteq M\}$, so $\Talg_\ell$ is saturated on $M$ iff $M$ is a module over no larger order; $\Talg_\ell$ is always saturated on its dual
$\Hom_{\Z_\ell}(\Talg_\ell,\Z_\ell)$ (if $x\Talg^{\vee}\subseteq\Talg^{\vee}$ then $x\cdot1\in\Talg^{\vee\vee}=\Talg_\ell$), which on the forms side is the lattice of $f\in
M_k(\Q_\ell)$ with $a_m(f)\in\Z_\ell$ for $m\ge1$ --- not computable without the congruence modules. Unsaturation is a property of the \emph{cohomology lattice}, not of $\Talg$.

\begin{proposition}[\'Etale primes, and the candidate bad set]\label{prop:heckeZetale}
\PROVED{} (the plan's Prop.~2.3.) Let $\mathcal O$ be a finite free $\Z_\ell$-algebra with $\mathcal O\otimes\Q_\ell$ reduced, and $\mathcal R\subseteq\mathcal O$ a subalgebra of full
rank (e.g.\ $\mathcal R=\Z_\ell[\vartheta]$ or $\Talg'_\ell$, $\mathcal O=\Talg_\ell$); $\operatorname{disc}$ is the Gram determinant of the trace form of the regular
representation.
\begin{enumerate}
\item $\operatorname{disc}(\mathcal R)=[\mathcal O:\mathcal R]^{2}\operatorname{disc}(\mathcal O)$ up to a unit. In particular, for $\vartheta\in\Talg^{\mathrm{fr}}$ with
  $\Q[\vartheta]=\Talg_\Q$, $\operatorname{disc}_\vartheta:=\operatorname{disc}\Z[\vartheta]=\operatorname{disc}(f_\vartheta)=[\Talg^{\mathrm{fr}}:\Z[\vartheta]]^{2}\operatorname{disc}\Talg^{\mathrm{fr}}$,
  $f_\vartheta$ the minimal polynomial (the trace is that of $\Talg_\Q$, not of $L\otimes\Q$, which counts cuspidal systems twice).
\item If $v_\ell(\operatorname{disc}\mathcal R)\le1$, then $\mathcal R=\mathcal O$. If $v_\ell(\operatorname{disc}\mathcal R)=0$, $\mathcal O$ is \'etale over $\Z_\ell$, every faithful
  lattice is saturated, and each maximal ideal lies on exactly one component. If $v_\ell(\operatorname{disc}\mathcal R)=1$, then $\ell$ is odd, exactly one maximal ideal
  $\midl$ is ramified, it lies on one component, and $\mathcal O_\midl$ is a DVR with ramification index $2$.
\item \textbf{Hence}, for $\ell$ with $\ell^{2}\nmid\operatorname{disc}_\vartheta$ (resp.\ $\ell^{2}\nmid\operatorname{disc}\Talg'$), $\Talg_\ell=\Z_\ell[\vartheta]\cong\Z_\ell[x]/(f_\vartheta)$
  (resp.\ $=\Talg'_\ell$) \textbf{with no generation hypothesis}; its maximal ideals are $(\ell,h(\vartheta))$ for the irreducible factors $\bar h$ of $\bar f_\vartheta$, and
  component $i$ passes through $(\ell,h(\vartheta))$ iff $\bar h\mid\bar f_i$ ($f_\vartheta=\prod_if_i$ over $\Q$, one $f_i$ per component).
\item Components $i\ne j$ meet above $\ell$ only if $\ell\mid\operatorname{Res}(f_i,f_j)$ for \textbf{every} such $\vartheta$. If no two meet above $\ell$, then
  $\Talg_\ell=\prod_i\mathcal O_i\otimes\Z_\ell$, $\mathcal O_i:=\Talg^{\mathrm{fr}}/P_i$ the Hecke order of component $i$ ($P_i$ its minimal prime).
\item If components meet at $\ell$, then $v_\ell(\operatorname{disc}\Talg^{\mathrm{fr}})\ge2\,v_\ell|C_{\Talg}(\mathcal S,\mathcal S')|$ for the congruence module of
  Definition~\ref{def:heckeZcong} ($\mathcal S\sqcup\mathcal S'$ all components).
\end{enumerate}
\end{proposition}

\begin{proof}
(1) A $\Z_\ell$-basis of $\mathcal R$ is the image of one of $\mathcal O$ under a matrix of determinant $[\mathcal O:\mathcal R]$ up to a unit, and the Gram matrices
transform accordingly. $\Z[\vartheta]\cong\Z[x]/(f_\vartheta)$ has Gram determinant $\operatorname{disc}(f_\vartheta)$ in the basis $1,\vartheta,\dots,\vartheta^{d-1}$.
(2) By (1), $2\,\mathrm{length}(\mathcal O/\mathcal R)\le v_\ell\operatorname{disc}\mathcal R$, so $\mathcal R=\mathcal O$ when $v\le1$. For $v=0$ the trace form of $\mathcal O/\ell$ is
non-degenerate, so $\mathcal O/\ell$ is a separable $\Fl$-algebra and $\mathcal O$ is \'etale, a product of unramified DVRs, hence a maximal order; the multiplier ring of
a faithful lattice is an order containing $\mathcal O$, hence equal to it. For $v=1$: the nilradical $\mathfrak n$ of $\mathcal O/\ell$ lies in the kernel of the reduced
trace form ($\operatorname{Tr}(nx)=0$ for nilpotent $nx$), and $v_\ell\det$ of the Gram matrix is at least the dimension of that kernel, so $\dim_{\Fl}\mathfrak n\le1$; and
$\mathfrak n\ne0$, since a reduced finite algebra over the perfect field $\Fl$ is separable, with non-degenerate trace form, which would give $v=0$. So $\mathfrak n$ is
one-dimensional and sits in one local factor $\mathcal O_\midl/\ell$; its radical $\mathfrak n$ is a non-zero vector space over the residue field modulo $\mathfrak n^{2}$, of
$\Fl$-dimension $1$, so the residue field is $\Fl$, $\mathfrak n^{2}=0$, $\mathcal O_\midl/\ell\cong\Fl[\varepsilon]/(\varepsilon^{2})$, and (Nakayama)
$\mathcal O_\midl\cong\Z_\ell[x]/(F)$ with $F$ monic of degree $2$, $\bar F=(x-a)^{2}$. If $\mathcal O_\midl\otimes\Q$ were not a field, $\mathcal O_\midl$ would be of index
$\ge\ell$ in $\Z_\ell\times\Z_\ell$ and $v\ge2$ by (1); so it is a quadratic order with $v_\ell(\operatorname{disc})=1$, hence maximal (by (1)), ramified (an unramified one
has unit discriminant), a DVR with ramification index $2$, and it lies on one component. At $\ell=2$: $F=x^{2}+bx+c$ with $\bar F=(x-a)^{2}$ has $b$ even, so
$v_2(b^{2}-4c)\ge2$; $v=1$ does not occur. (3) is (2) with $\mathcal O/\ell=\Fl[x]/(\bar f_\vartheta)$; the kernel of $\Talg_\ell\to\Z_\ell[x]/(f_i)$ is $P_i\Talg_\ell$, generated by
$f_i(\vartheta)$, and $f_i(\vartheta)\in(\ell,h(\vartheta))$ iff $\bar h\mid\bar f_i$. (4) Write $\operatorname{Res}(f_i,f_j)=af_i+bf_j$, $a,b\in\Z[x]$, and evaluate at $\vartheta$:
$f_i(\vartheta)\in P_i$, $f_j(\vartheta)\in P_j$, so $\operatorname{Res}\in(P_i+P_j)\cap\Z\subseteq\midl\cap\Z=\ell\Z$ if $\midl\supseteq P_i+P_j$. If no two meet, every block
$\Talg_\midl=e_\midl\Talg_\ell$ of the complete semilocal ring $\Talg_\ell$ lies on exactly one component, and the component idempotent $\varepsilon_i\in\Talg_\Q\otimes\Q_\ell$ is
$\sum_{\midl\supseteq P_i}e_\midl\in\Talg_\ell$ ($\varepsilon_ie_\midl\ne0$ iff $(\mathcal O_i)_\midl\ne0$ iff $P_i\subseteq\midl$); then
$\Talg_\ell=\prod_i\varepsilon_i\Talg_\ell$ and $\varepsilon_i\Talg_\ell$ is the image of $\Talg_\ell$ in $\mathcal O_i\otimes\Q_\ell$. (5) [U,~\texttt{lem:u-fibre}] gives
$\Talg^{\mathrm{fr}}=\Talg_{\mathcal S}\times_{C}\Talg_{\mathcal S'}$ with both projections onto and $(\Talg_{\mathcal S}\times\Talg_{\mathcal S'})/\Talg^{\mathrm{fr}}\cong C$; apply (1)
to $\Talg^{\mathrm{fr}}\subseteq\Talg_{\mathcal S}\times\Talg_{\mathcal S'}$.
\end{proof}

\noindent
So the \textbf{computable candidate bad set} is $\mathsf S_{\mathrm{cand}}:=\{\ell:\ell^{2}\mid\operatorname{disc}\Talg'\}\cup\{\text{torsion primes}\}$, contained in
$\{\ell:\ell^{2}\mid\operatorname{disc}_\vartheta\}\cup\{\text{torsion primes}\}$ and shrinking as the operator set grows; the set $\{\ell\mid\operatorname{disc}\Talg^{\mathrm{fr}}\}$ is
not computable before $\Talg^{\mathrm{fr}}$ is known, and primes with $v_\ell=1$ need nothing.

\begin{proposition}[Monogenic blocks: $K=1$ suffices]\label{prop:heckeZmono}
\PROVED{} \STD{a finite flat local $\Z_\ell$-algebra is Gorenstein iff its reduction modulo $\ell$ is, \cite[Prop.~3.1.19(b)]{bruns_herzog}\VERIFYloc} (the plan's
Prop.~2.4, its hypothesis as corrected by A1-M2.) Assume generation at $\ell$ certified ($\Talg'_\ell=\Talg_\ell$, by Proposition~\ref{prop:heckeZgen} or by a theorem
such as Proposition~\ref{prop:heckeZeiscert}), $\vartheta\in\Talg'$ with $\Q[\vartheta]=\Talg_\Q$, and that $\bar\vartheta$ generates \textbf{$\Talg_\ell/\ell\Talg_\ell$} as an
$\Fl$-algebra. Then:
\begin{enumerate}
\item $\Talg_\ell=\Z_\ell[\vartheta]\cong\Z_\ell[x]/(f_\vartheta)$.
\item The maximal ideals are the irreducible factors $\bar h$ of $\bar f_\vartheta$; component $i$ passes through $(\bar h)$ iff $\bar h\mid\bar f_i$.
\item For disjoint sets $\mathcal S,\mathcal S'$ of components, $C_{\Talg}(\mathcal S,\mathcal S')\otimes\Z_\ell\cong\Z_\ell[x]/(f_{\mathcal S},f_{\mathcal S'})$, of length
  $v_\ell\operatorname{Res}(f_{\mathcal S},f_{\mathcal S'})$, $f_{\mathcal S}:=\prod_{i\in\mathcal S}f_i$.
\item Every $\Talg_\midl$ is a complete intersection and Gorenstein, and $\Talg_\midl/\ell$ has embedding dimension $\le1$ over its residue field.
\end{enumerate}
Conversely, if some $\Talg_\midl/\ell$ has embedding dimension $\ge2$ --- in particular if $\Talg_\midl$ is not Gorenstein --- no $\vartheta$ generates.
\emph{Computable form of the hypothesis.} By Proposition~\ref{prop:saturation}(3) for $(L_\ell,\Talg_\ell)$, $\Talg_\ell/\ell\Talg_\ell\to A_K/\ell A_K$ is an isomorphism of
$\Fl$-algebras exactly when $\dim_{\Fl}A_K/\ell A_K=d$, and for every $K\ge c^{\mathrm{sat}}+1$. So raise $K$ until $\dim A_K/\ell A_K=d$ ($K=c^{\mathrm{sat}}+1$ at the
latest; $c^{\mathrm{sat}}$ need not be known in advance), then test whether $\bar\vartheta$ generates $A_K/\ell A_K$; under generation $A_K=A'_K$. On the
Proposition~\ref{prop:heckeZgen} path ($\Talg_\ell$ saturated) $K=1$. On an unsaturated lattice the test at $K=1$ is \textbf{wrong}: $A_1$ is then a proper quotient of
$\Talg_\ell/\ell$ (the plan's A1-M2: $\Talg_\ell=\{(a,b)\in\Z_\ell^{2}:a\equiv b\ (\ell^{2})\}$ on $\Z_\ell^{2}$, $A_1=\Fl$, $\vartheta=(1+\ell^{3},1)$ generates $A_1$ but
$\Z_\ell[\vartheta]=\{a\equiv b\ (\ell^{3})\}\subsetneq\Talg_\ell$; there $c^{\mathrm{sat}}=2$ and $A_3/\ell A_3\cong\Fl[\varepsilon]/(\varepsilon^{2})$, which $\bar\vartheta=1$ does not
generate).
\end{proposition}

\begin{proof}
(1) $\Z_\ell[\vartheta]+\ell\Talg_\ell=\Talg_\ell$, so Nakayama gives equality, and $\Z_\ell[\vartheta]\cong\Z_\ell[x]/(f_\vartheta)$ as $f_\vartheta$ is the minimal polynomial and
is monic. (2) $\Talg_\ell/\ell=\Fl[x]/(\bar f_\vartheta)$, and Proposition~\ref{prop:heckeZetale}(3)'s argument. (3) The kernel of $\Z_\ell[x]/(f_\vartheta)\to\Z_\ell[x]/(f_{\mathcal S})$
is generated by $f_{\mathcal S}$ ($f_{\mathcal S}\mid f_\vartheta$), so $I_{\mathcal S}\otimes\Z_\ell=f_{\mathcal S}(\vartheta)\Talg_\ell$, likewise for $\mathcal S'$, and
$C\otimes\Z_\ell=\Z_\ell[x]/(f_\vartheta,f_{\mathcal S},f_{\mathcal S'})=\Z_\ell[x]/(f_{\mathcal S},f_{\mathcal S'})$. This is the cokernel of multiplication by $f_{\mathcal S'}(x)$ on
the free $\Z_\ell$-module $\Z_\ell[x]/(f_{\mathcal S})$ of rank $\deg f_{\mathcal S}$, whose determinant is the norm $\prod_{f_{\mathcal S}(\alpha)=0}f_{\mathcal S'}(\alpha)=
\operatorname{Res}(f_{\mathcal S},f_{\mathcal S'})$; its length is the valuation. (4) $\Talg_\midl\cong\Z_\ell[x]/(F_\midl)$ with $F_\midl$ the Hensel factor of $f_\vartheta$ at
$\midl$, a quotient of a two-dimensional regular local ring by one element; $\Talg_\midl/\ell\cong\Fl[x]/(\bar h^{a})$, whose maximal ideal is generated by $\bar h(x)$ and
whose socle, $\bar h^{a-1}\Fl[x]/(\bar h^{a})$, is one-dimensional over the residue field; so $\Talg_\midl/\ell$ is Gorenstein, and so is $\Talg_\midl$ (cited). Converse: a local
$\Fl$-algebra generated by one element has a principal maximal ideal.
\end{proof}

\noindent
\emph{The kind of a block.} For $\Talg_\midl=\Z_\ell[x]/(F_\midl)$ the kind --- DVR, non-maximal order of one component, or a meeting of several components --- is decided in
two steps at finite precision: the $\ell$-adic factorisation of $F_\midl$ (Hensel, from $f_\vartheta$ modulo a finite power of $\ell$) gives the components through $\midl$ and
their residue fields, which separates ``meeting'' from ``one component''; for one component, whether $\Z_\ell[x]/(F_\midl)$ is maximal is its index in the maximal order
(Dedekind's criterion, or one Ore--Montes step when that is inconclusive). \emph{Per block} (not asserted as a statement: the bookkeeping when the Hensel factors of distinct
blocks are not coprime modulo $\ell$ is not written out): if for each block some computed $\vartheta_\midl$ generates $\Talg_\midl/\ell$ (tested in $e_\midl A_K/\ell A_K$,
$K\ge c^{\mathrm{sat}}_\midl+1$, or at $K=1$ on a saturated block), then $\Talg_\midl=\Z_\ell[e_\midl\vartheta_\midl]$.
\emph{Example} ($691$). At $(1,10)$, $\Talg^{\mathrm{fr}}=\{(a,b)\in\Z^{2}:a\equiv b\ (691)\}$ (Eisenstein, $\Delta$) has the $\Z$-basis $(1,1)$, $(691,0)$, and
$T_2=(2049,-24)=-24\cdot(1,1)+3\cdot(691,0)$. Modulo $691$, $T_2\equiv-24+3\varepsilon$, $\varepsilon:=(691,0)$, $\varepsilon^{2}\equiv0$, so $\bar T_2$ generates
$\Talg_{691}/691=\F_{691}[\varepsilon]/(\varepsilon^{2})$ (which is $A_1$ there: the lattice is saturated at $691$). Hence $\Talg_{691}=\Z_{691}[T_2]$, and the depth of the
congruence is $v_{691}\operatorname{Res}(x-2049,x+24)=v_{691}(2073)=1$: \textbf{no $K>1$ is needed at $691$}. At $3$, $\Z[T_2]$ has index $3$ in $\Talg^{\mathrm{fr}}$
($2073=3\cdot691$) and $\bar T_2$ does not generate. (Sign: \Espec{cl:ramanujan} uses $(0,691)\equiv-(691,0)$, hence its $-3$ against $+3$ here.)

\begin{proposition}[Reconstruction from $A_K$]\label{prop:heckeZlocal}
\PROVED{} (the plan's Prop.~2.5.) Let $\vartheta\in\Talg'$ with $\Q[\vartheta]=\Talg_\Q$, $\mathcal R:=\Z_\ell[\vartheta]\subseteq\Talg_\ell$, $\ell^{e}$ the exponent of
$\Talg_\ell/\mathcal R$, and assume $A_K$ is the image of $\Talg_\ell$ itself (generation, certified or by a theorem). Let $c^{\mathrm{sat}}$ be that of
Definition~\ref{def:satdata} for $(L_\ell,\Talg_\ell)$ ($0$ if saturated) and $\pi\colon\mathcal R\to A_K$ the restriction of the reduction.
\begin{enumerate}
\item For every $E\ge e$ and $K\ge E+c^{\mathrm{sat}}$: $\ell^{E}\Talg_\ell=\pi^{-1}(\ell^{E}A_K)$. So $\Talg_\ell=\ell^{-E}\pi^{-1}(\ell^{E}A_K)\subset\Q_\ell[x]/(f_\vartheta)$ is
  an explicit order.
\item $e\le\lfloor v_\ell(\operatorname{disc}_\vartheta)/2\rfloor$. \textbf{Stopping rule}, saturated case: $Q_K:=A_K/\pi(\mathcal R)\cong(\Talg_\ell/\mathcal R)/\ell^{K}$ has exponent
  $\ell^{\min(e,K)}$, so the first $K$ with $\exp Q_K<\ell^{K}$ gives $e=\log_\ell\exp Q_K$, without $\operatorname{disc}_\vartheta$. Unsaturated case: the stopping rule does
  \textbf{not} apply ($Q_K$ also carries the defect); use $E:=\lfloor v_\ell(\operatorname{disc}_\vartheta)/2\rfloor$ and $K\ge E+c^{\mathrm{sat}}$, $c^{\mathrm{sat}}$ off the
  plateau of Proposition~\ref{prop:satdefect}(2).
\end{enumerate}
\end{proposition}

\begin{proof}
(1) ``$\subseteq$'': $\ell^{E}t\in\mathcal R$ for $t\in\Talg_\ell$, and $\pi(\ell^{E}t)=\ell^{E}\bar t$. ``$\supseteq$'': if $\pi(r)=\ell^{E}\bar a$ with $a\in\Talg_\ell$, then
$r-\ell^{E}a\in\ker(\Talg_\ell\to A_K)=\Talg_\ell\cap\ell^{K}\Talg^{\mathrm{sat}}_\ell\subseteq\ell^{K-c^{\mathrm{sat}}}\Talg_\ell\subseteq\ell^{E}\Talg_\ell$
(Proposition~\ref{prop:saturation}(1)). (2) $e\le\mathrm{length}(\Talg_\ell/\mathcal R)\le v_\ell(\operatorname{disc}_\vartheta)/2$ by Proposition~\ref{prop:heckeZetale}(1); in the
saturated case $A_K=\Talg_\ell/\ell^{K}\Talg_\ell$, so $Q_K=\Talg_\ell/(\mathcal R+\ell^{K}\Talg_\ell)$, and a finite $\ell$-group $G$ of exponent $\ell^{e}$ has $G/\ell^{K}G$ of
exponent $\ell^{\min(e,K)}$.
\end{proof}

\noindent
Only \emph{generation} is needed for reconstruction; saturation is one way to certify it. With generation certified there is also a \textbf{rational route} that needs no
$K>1$: $\Talg'$ is the $\Z$-algebra generated by the $g_p(\vartheta)\in\Q[x]/(f_\vartheta)$ with $X_p=g_p(\vartheta)$, computed by a Hermite normal form over $\Z$; the two
routes cross-check each other.

\begin{proposition}[Which components pass through a block]\label{prop:heckeZthrough}
\PROVED{} (the plan's Prop.~2.6.) Let $\vartheta$ be as in Proposition~\ref{prop:heckeZlocal}, $f_\vartheta=\prod_if_i$, $g_i:=f_\vartheta/f_i$, $\varepsilon_i\in\Talg_\Q$ the
idempotent of component $i$, and $e\in\Talg_\ell$ an idempotent (e.g.\ of a block of $\Talg'_\ell$). Then $g_i(\vartheta)e\ne0$ iff $\varepsilon_ie\ne0$ iff component $i$ passes
through some maximal ideal in the support of $e$. If $g_i(\vartheta)e\ne0$, then $g_i(\vartheta)e\not\equiv0\pmod{\ell^{K}}$ on $L_\ell$ for every $K>v_\ell\operatorname{Res}(f_i,g_i)$;
more generally for every $K>v_\ell(D)$, $D\in(f_i,g_i)+\ell^{K}\Z[x]$ an integer. No generation and no saturation is needed.
\end{proposition}

\begin{proof}
$g_i(\vartheta)$ kills the components $j\ne i$ ($f_j\mid g_i$), and $g_i(\vartheta)\varepsilon_i$ is a unit of $\varepsilon_i\Talg_\Q$ ($\gcd(f_i,g_i)=1$); so $g_i(\vartheta)e=0$ iff
$\varepsilon_ie=0$. For a block idempotent $e_\midl$, $\varepsilon_ie_\midl\ne0$ iff $(\mathcal O_i)_\midl\ne0$ iff $P_i\subseteq\midl$. Precision: $X:=eL_\ell\cap(L\otimes\Q_\ell)_i$ is
non-zero when $\varepsilon_ie\ne0$ and saturated in $L_\ell$ ($eL_\ell$ is a direct summand). Write $D=uf_i+vg_i+\ell^{K}h$; on $X$, $f_i(\vartheta)=0$, so $Dx=v(\vartheta)g_i(\vartheta)x
+\ell^{K}h(\vartheta)x$. If $g_i(\vartheta)eL_\ell\subseteq\ell^{K}L_\ell$, then $DX\subseteq\ell^{K}L_\ell\cap X=\ell^{K}X$, i.e.\ $v_\ell(D)\ge K$. $D=\operatorname{Res}(f_i,g_i)$,
$h=0$, is a case.
\end{proof}

\noindent
\emph{The idempotent} (the plan's A1-N7, F07-5). In practice $e$ is computed in $A'_K$ and lifted. $\ker(\Talg'_\ell\to A'_K)$ consists of elements $\ell^{K}x$ with $x$ in the
saturation of $\Talg'_\ell$, a finite $\Z_\ell$-algebra, so it lies in the Jacobson radical; idempotents lift uniquely modulo it, and exist because $\Talg'_\ell$ is
$\ell$-adically complete (lift at each finite level $A'_K\to A'_1$, a nilpotent kernel, and take the limit).

\paragraph{The coefficient sequences, their splitting, and the glue.} The next four statements are the plan's stage 03 (\file{03-bockstein-tor.md} v0.3: B1, B7,
B8, B8$'$, B9, B11--B15). Notation: $H^{q}_\ell:=H^{q}\otimes\Z_\ell$, $\mathrm{tors}_\ell:=\mathrm{tors}\,H^{1}_\ell$, $e_\ell$ its exponent ($\ell^{e_\ell}$), $L_q:=H^{q}/\mathrm{tors}$.

\begin{proposition}[Universal coefficients with the Hecke action]\label{prop:heckeZuct}
\PROVED{} \STD{the universal coefficient theorem for a complex of free abelian groups, \cite[\S3.6]{weibel}\VERIFYloc; naturality of connecting maps,
\cite[Thm.~1.3.1]{weibel}} (the plan's B1, with B4(1) and the consequences of 01 \S2.) Let $\ell$ be a prime and $K\ge1$.
\begin{enumerate}
\item For every $q\ge0$ there is a natural exact sequence
  \[
    0\to H^{q}\otimes\Z/\ell^{K}\xrightarrow{\ \mathrm{red}_K\ }H^{q}(\Ga;\Vcoef/\ell^{K})\xrightarrow{\ \beta_K\ }H^{q+1}[\ell^{K}]=\mathrm{Tor}^{\Z}_1(H^{q+1},\Z/\ell^{K})\to0,
  \]
  $\beta_K$ the connecting map of $0\to\Vcoef\xrightarrow{\ell^{K}}\Vcoef\to\Vcoef/\ell^{K}\to0$; it splits as a sequence of abelian groups, not naturally.
\item Every map in it commutes with every double-coset operator $\cor\circ(\conj_g)_\ast\circ\res$, $g$ in the Hecke semigroup of $\Ga_0(N)$: $T_p$ ($p\nmid N$, $p=\ell$
  included), $U_p$ ($p\mid N$) and their words.
\item The same holds for any complex $C$ of finitely generated free $\Z_{(\ell)}$-modules quasi-isomorphic in degrees $\le q_{\max}$ to
  $\Hom_\Ga(P_\bullet,\Vcoef)\otimes\Z_{(\ell)}$, $P_\bullet$ a resolution by finitely generated free modules: $H^{q}(C/\ell^{K})=H^{q}(\Ga;\Vcoef/\ell^{K})$ for $q<q_{\max}$, with
  the sequence of (1), and a cochain endomorphism $F$ of $C$ induces compatible endomorphisms of the three terms. \textbf{That endomorphism is $T_p$ on
  $H^{q}(\Ga;\Vcoef/\ell^{K})$ only if $F$ is a cochain-level lift of $T_p$ over $\Z_{(\ell)}$} in the sense of Proposition~\ref{prop:strictliftN}(2) --- not merely an $F$
  inducing $T_p$ on $H^{\bullet}(C)$: the ambiguity is the component of $\End_{D(\Z_{(\ell)})}(C)$ in $\Ext^{1}(H^{q+1}(C),H^{q}(C))$, which acts as
  $H^{q+1}[\ell^{K}]\to H^{q}/\ell^{K}$, exactly on the $\mathrm{Tor}$ term (example: $C^{0}=\Z a\oplus\Z b$, $C^{1}=\Z c$, $da=0$, $db=\ell c$; $F(a)=a$, $F(b)=b+a$,
  $F(c)=c$ is the identity on $H^{\bullet}(C)$ but not on $H^{0}(C/\ell)$). The shipped operators are such lifts (Propositions~\ref{prop:strictliftN}(2),
  \ref{prop:descentpairhecke}(1)); an export built from integral-cohomology data alone would be wrong on the $\mathrm{Tor}$ part.
\item \emph{Consequences.} $0\subseteq\mathrm{tors}_q/\ell^{K}\subseteq H^{q}/\ell^{K}\subseteq H^{q}(\Ga;\Vcoef/\ell^{K})$ is a filtration by Hecke submodules with graded pieces
  $\mathrm{tors}_q/\ell^{K}$, $L_q/\ell^{K}$, $H^{q+1}[\ell^{K}]$; the middle step is the reduction of $0\to\mathrm{tors}_q\to H^{q}\to L_q\to0$, not a universal-coefficient
  sequence; $H^{q}/\ell^{K}=\ker\beta_K$ (on a free model: $x$ lies in it iff $\ell^{-K}d\tilde x$ is a coboundary for a lift $\tilde x$), and there is no Hecke-equivariant
  complement in general. At $\ell\ge5$, $H^{1}(\Ga;\Vcoef/\ell^{K})=H^{1}/\ell^{K}$ ($H^{2}$ has no $\ell$-torsion, Theorem~\ref{thm:descentcert}(1)), which is
  $L/\ell^{K}$ exactly when $\ell\nmid|\mathrm{tors}\,H^{1}|$. At $\ell\in\{2,3\}$, $H^{1}/\ell\subset H^{1}(\Ga;\Vcoef/\ell)\twoheadrightarrow H^{2}[\ell]$: $H^{2}[\ell]$ is a
  \emph{quotient}, there is no canonical Hecke-equivariant copy of it inside, one exists iff the sequence splits $\Talg$-equivariantly (Proposition~\ref{prop:heckeZsplit}
  guarantees it under disjoint supports and does not characterise it), and a mod-$\ell$ maximal ideal must be attributed to its degree through $\beta_1$ before it is
  reported as a prime of $\Talg$.
\end{enumerate}
\end{proposition}

\begin{proof}
(1) $\Vcoef$ is $\Z$-free, so the coefficient sequence is an exact sequence of $\Z[\De]$-modules, and $H^{\bullet}(\Ga;-)$ is a $\delta$-functor; the long exact sequence
breaks into the displayed short ones because $H^{q}\xrightarrow{\ell^{K}}H^{q}$ has cokernel $H^{q}/\ell^{K}$ and kernel $H^{q}[\ell^{K}]$ (the case $K=1$ is
Lemma~\ref{lem:bocksteinallell} and, at $\ell\ge5$, Claim~\ref{cl:bockstein}(2)). $\Ga_0(N)$ is of type $\mathrm{FP}_\infty$, so $\Hom_\Ga(P_\bullet,\Vcoef)$ is a complex of free
abelian groups and the universal coefficient theorem gives the splitting. (2) $\res$, $(\conj_g)_\ast$ and $\cor$ are morphisms of $\delta$-functors in the coefficient
module (Claim~\ref{cl:bockstein}(3); at level $N$ the operator of Lemma~\ref{lem:bulkperiodic}(3), realised on the models by Proposition~\ref{prop:strictliftN}; $U_p$ on
cohomology, Proposition~\ref{prop:descentmaps}(3)(b$'$)). (3) A quasi-isomorphism of bounded-below complexes of free $\Z_{(\ell)}$-modules stays one after $\otimes\Z/\ell^{K}$
(its cone is acyclic, bounded below and free over a principal ideal domain, hence split exact); for a comparison in degrees $\le q_{\max}$ apply this to the brutal
truncations, or compare the universal coefficient sequences by the five lemma in degrees $<q_{\max}$. $C$ free gives the sequence for $C$, and $F$ commutes with
$\mathrm{red}_K$ and $\beta_K$ by naturality, on the nose. The example: $H^{0}(C)=\Z a$, $H^{1}(C)=\Z/\ell$, $F$ is the identity on both, and on $H^{0}(C/\ell)=\Fl a\oplus\Fl b$
it sends $b\mapsto b+a$. (4) $\mathrm{Tor}_1(L_q,\Z/\ell^{K})=0$ makes the first inclusion injective; the rest is (1), (2) and the definition of $\beta_K$.
\end{proof}

\noindent
\emph{Nothing here depends on $U_p$ over $\Z$ being built}: the statement is about the operators on cohomology. The record of Theorem~\ref{thm:descentcert} predicts every
mod-$\ell^{K}$ type at no cost (the plan's B2: $H^{q}(\Ga;\Vcoef/\ell^{K})\cong(\Z/\ell^{K})^{b_q}\oplus\bigoplus_i\Z/\ell^{\min(a_{q,i},K)}\oplus\bigoplus_i\Z/\ell^{\min(a_{q+1,i},K)}$
for $H^{q}_{(\ell)}\cong\Z_{(\ell)}^{b_q}\oplus\bigoplus_i\Z/\ell^{a_{q,i}}$; checked against the Lean-kernel dimensions of Remark~\ref{rmk:lvdimsltwothree} on $16$ of $16$ cases,
the plan's \S6.1).

\begin{proposition}[Splitting by support]\label{prop:heckeZsplit}
\PROVED{} (the plan's B7.) Let $0\to S\to M\to Q\to0$ be an exact sequence of finite modules over a commutative ring $R$ (here $\Talg$, or the image of a subalgebra
$\Talg'$). If $\supp_R(S)\cap\supp_R(Q)=\emptyset$, it splits uniquely as $R$-modules, $M=M_{\supp S}\oplus M_{\supp Q}$, cut out by the idempotents of the Artinian image of
$R$ in $\End(M)$; and $\Hom_R(Q,S)=\Hom_R(S,Q)=\Ext^{1}_R(Q,S)=0$. For a subring $R'\subseteq R$, disjointness for $R'$ implies disjointness for $R$ (maximal ideals of $R$
contract to those of $R'$), so \textbf{a separation computed from the operators in $\mathcal P$ is a certificate for all of $\Talg$}. The same holds for modules finitely
generated over $\Z_\ell$ (e.g.\ $L_\ell$), with the complete semilocal ring $\Talg_{\Z_\ell}$ (resp.\ $\Talg_\ell$) in place of the Artinian image.
\end{proposition}

\begin{proof}
The image $A$ of $R$ in $\End M$ is finite, hence Artinian, $A=\prod_\midl A_\midl$; its idempotents act on the whole sequence, and on each factor one of the outer terms
vanishes. Homomorphisms and extensions between modules supported on disjoint sets of idempotents vanish. A maximal ideal of $R$ in both supports would contract to one of
$R'$ in both supports. In the $\Z_\ell$-finite case the same argument runs with the block idempotents of the complete semilocal ring.
\end{proof}

\noindent
\emph{The separation certificate} (the plan's K6). The torsion residues of $T_p$ are the roots of $\prod_jP_j(x)$, the layer polynomials of
Proposition~\ref{prop:torsionhecke}(1)(a) --- possibly several, possibly in an extension of $\Fl$ --- so separation at $T_p$ is
$\gcd\bigl(\chi^{(N)}_{p,1}\bmod\ell,\ \prod_jP_j\bigr)=1$ in $\Fl[x]$, certified by a B\'ezout pair or by $\operatorname{Res}\not\equiv0\pmod\ell$; a certificate of $O(b_1)$ small
integers, not ``a few''. The per-operator test is weaker than the joint one (at $(1,18)$, $\ell=5$, every $T_p$ alone agrees with some free eigenvalue, but no free joint
system agrees with the torsion's); the joint test needs the matrices. \emph{When the sequences of Proposition~\ref{prop:heckeZuct}(1) split.} At $\ell\ge5$ the only
non-trivial one is $q=0$, an isomorphism; the question is the torsion/free extension inside $H^{1}/\ell^{K}$ (Proposition~\ref{prop:heckeZglue}). At $\ell\in\{2,3\}$, level $1$:
never by support for the prime-to-$\ell$ algebra (Proposition~\ref{prop:heckeZlevelone23}(1)); $T_\ell$ may separate. At $\ell\in\{2,3\}$, level $N$:
Proposition~\ref{prop:ltwothreescalar} is a level-one statement and fails at level $N$ ($\Ga_0(11)$, $n=0$, $\ell=2$: $T_3$ acts on the newform's line by $-1\not\equiv1+3$);
supports must be computed.

\begin{proposition}[Level one at $\ell\in\{2,3\}$]\label{prop:heckeZlevelone23}
\PROVED{} given Proposition~\ref{prop:ltwothreescalar} and Lemma~\ref{lem:tellrank} (the plan's B8, B8$'$, B9, with A3-M1's correction). Let $N=1$, $n$ even,
$\ell\in\{2,3\}$.
\begin{enumerate}
\item For every prime $p\ne\ell$ and every $q$: $T_p-(1+p)\in\ell\End(L_q\otimes\Z_{(\ell)})$, and $T_p-(1+p)$ is nilpotent on $\mathrm{tors}\,H^{q}_{(\ell)}$, on $L_q/\ell^{K}$,
  on $H^{q+1}[\ell^{K}]$ and on $H^{q}(\Ga;\Vcoef/\ell^{K})$. It is \textbf{not} nilpotent on $H^{q}_{(\ell)}$ in general (at $(1,10)$, $\ell=3$, $T_2-3$ has the non-zero
  eigenvalues $2046$ and $-27$ on $L\otimes\Q$); on $L_{(\ell)}$ it is only topologically nilpotent. Equivalently: $T_p-(1+p)$ lies in every maximal ideal of $\Talg$ above
  $\ell$.
\item The maximal ideals of $\Talg$ above $\ell$ are $\midl_\psi:=(\ell,\ T_p-(1+p)\ (p\ne\ell),\ \psi(T_\ell))$, $\psi$ running over the irreducible factors modulo $\ell$
  of the characteristic polynomial of $T_\ell$ on $H^{1}/\ell$, which is $\chi^{(1)}_{\ell,1}\bmod\ell$ (free part) times $\prod_jP_j$ (the torsion layers of $T_\ell$,
  Proposition~\ref{prop:torsionhecke}(1)(a)) --- both in the descent record. $\midl_\psi\in\operatorname{Spec}\Talg^{\mathrm{fr}}$ iff $\psi\mid\chi^{(1)}_{\ell,1}\bmod\ell$;
  $\midl_\psi\in\operatorname{Spec}\Talg^{\mathrm{tors}}$ iff $\psi\mid\prod_jP_j$. No generation hypothesis is needed.
\item $T_\ell$ is nilpotent on $\mathrm{tors}_\ell$; so $\supp(\mathrm{tors}_\ell)\subseteq\{\midl_x\}$.
\item (a) The torsion point $\midl_x$ \emph{coincides} with a free point iff $\chi^{(1)}_{\ell,1}$ has the root $0$ modulo $\ell$, i.e.\ iff some cusp eigenform is
  non-ordinary at $\ell$ (the Eisenstein series has $a_\ell=1+\ell^{k-1}\equiv1$); a certificate for all $p$. (b) $\Talg$ has a torsion-only point above $\ell$ iff
  $\mathrm{tors}_\ell\ne0$ and every cusp eigenform is $\ell$-ordinary. (c) The torsion never meets the Eisenstein point $\midl_{x-1}$.
\item \textbf{No cusp eigenform of level $1$ is $\ell$-ordinary.} Hence if $S_k\ne0$ the torsion point $\midl_x$ is common with the cuspidal part, and a torsion-only point above
  $\ell$ exists iff $\mathrm{tors}_\ell\ne0$ and $S_k=0$ ($k\in\{4,6,8,10,14\}$).
\item $\dim_{\Fl}A_1(L/\ell)\le1+\rank(T_\ell\mid L/\ell)\le3$. Hence $\Talg_\ell$ is \textbf{not saturated on $L_\ell$ whenever $d=\dim M_k\ge4$} ($k=36$ or $k\ge40$),
  at $\ell=2$ and at $\ell=3$, and Theorem~\ref{thm:heckeZcert}(C3) cannot pass at $2$ or $3$ for any operator set. This is a statement about the cohomology lattice $L$
  (and its direct summands), not about every $\Talg^{\mathrm{fr}}$-stable lattice.
\end{enumerate}
\end{proposition}

\begin{proof}
(1) Proposition~\ref{prop:ltwothreescalar} gives $T_p=1+p$ on $H^{q}(\Ga;\Vcoef/\ell)$, and $H^{q}/\ell$ (a sub), $\mathrm{tors}/\ell$ (a sub of that), $L_q/\ell$ (a
subquotient) and $H^{q+1}[\ell]$ (a quotient) are subquotients of it (Proposition~\ref{prop:heckeZuct}(4)); $0$ on $L_q/\ell$ is the statement
$T_p-(1+p)\in\ell\End(L_q\otimes\Z_{(\ell)})$. On a finite $\ell$-group, or on a lattice modulo $\ell^{K}$, an endomorphism that is $0$ on every $\ell$-adic layer is nilpotent.
(2) Proposition~\ref{prop:heckeZmax} with $M=H^{1}\otimes\Z_{(\ell)}$, faithful over $\Talg\otimes\Z_{(\ell)}$: $V(\ell)\cong\operatorname{Spec}A_1(H^{1}/\ell)$, and $\Talg$ acts on
$H^{1}/\ell$ through $\Fl[T_\ell]$ by (1). In the Grothendieck group of finite $\Talg$-modules $[G/\ell G]=[G[\ell]]$ ($0\to G[\ell]\to G\xrightarrow{\ell}G\to G/\ell G\to0$), and
$G[\ell]$ is filtered by the layers, whence the torsion factor. (3) $\beta\colon H^{0}(\Ga;\Vcoef/\ell)\to H^{1}[\ell]$ is a $T_\ell$-equivariant isomorphism
(Lemma~\ref{lem:bocksteinallell}), and $T_\ell=U_\ell$ has square $0$ on $H^{0}(\Ga;\Mn)$ (Proposition~\ref{prop:heckeH0}(2)); by the Grothendieck-group identity of (2),
every composition factor of $\mathrm{tors}_\ell$ is one of $H^{1}[\ell]$, on which $T_\ell$ acts by $0$, so $T_\ell$ is nilpotent. (4) From (2), (3). (5) On $L/\ell$, a
quotient of the submodule $H^{1}/\ell$ of $H^{1}(\Ga;\Mn)$, $\rank T_\ell\le2$ (Lemma~\ref{lem:tellrank}: $\ell=2$ every $q\ge1$; $\ell=3$, $n$ even, $q=1$; rank does not grow
on subquotients). The invertible Fitting part of $T_\ell$ on $L/\ell$ has dimension at most that rank, and it counts the eigenvalues of $T_\ell$ on $L\otimes\Q$ that are
$\ell$-adic units, with multiplicity. The Eisenstein system contributes $1$; each cuspidal system occurs twice in $L$ (Proposition~\ref{prop:heckeZrat}(2)), so an
$\ell$-ordinary cusp eigenform would add $\ge2$, total $\ge3>2$. $S_k=0$ exactly for $k\in\{4,6,8,10,14\}$. (6) $A_1(L/\ell)=\Fl[\bar T_\ell]$ by (1). If $Y$ has rank $r$ and
$g$ is the characteristic polynomial of $Y$ on $\im Y$, then $g(Y)Y=0$, so $\deg m_Y\le r+1$. Saturation is $\dim A_1=d$ (Proposition~\ref{prop:saturation}(2)); and $\dim
M_k\ge4$ iff $k=36$ or $k\ge40$.
\end{proof}

\noindent
Checks (the plan's \S6.2): at the eight level-one cells $n\in\{8,\dots,30\}$ the torsion's $T_\ell$ residues are $\{0\}$ at $2$ and $3$; the free ones are $\{1\}$ at
$n=8,12$ ($S_k=0$) and $\{0,1\}$ otherwise --- torsion-only points exactly at $n=8,12$, and at $(1,10)$ the torsion points over $2$ and $3$ \textbf{are} $\Delta$'s maximal
ideals. The $2$- and $3$-adic structure of $\Talg^{\mathrm{fr}}$ at level $1$ for $d\ge4$ must therefore come from Proposition~\ref{prop:heckeZeiscert} or from $A_K/\ell A_K$ at
$K\ge c^{\mathrm{sat}}+1$ under a generation theorem (Proposition~\ref{prop:heckeZmono}, computable form).

\begin{remark}[CB-1: $T_p=1+p$ on the $2$-primary torsion at level one]\label{rmk:heckeZcbone}
\CONJ{} (the plan's CB-1; recorded, \textbf{used by nothing}, no hypothesis structure.) For $N=1$, $n$ even and every odd prime $p$, $T_p=1+p$ exactly on
$\mathrm{tors}\,H^{1}(\SL_2(\Z);\Vcoef)_{(2)}$. Evidence: all eight checked cells, $n\le30$, exponents up to $2^{4}$ (e.g.\ $(1,16)$: $\Z/2\oplus\Z/4\oplus\Z/16$,
$\ker(T_3-4)=\ker(T_5-6)=\ker(T_7-8)$ everything); consistent with the Frobenius classes $h_j$ of Remark~\ref{rmk:bocksteinexp}(2), on which $T_p\equiv(1+p)p^{a2^{j}}$ modulo $2^{j+1}$ (the plan's
B10(2), from Lemma~\ref{lem:dicksontwist} and Proposition~\ref{prop:heckeH0}(1)), since $p^{2^{j}}\equiv1\pmod{2^{j+1}}$. The analogue at $\ell=3$ is \textbf{false}: the $\Z/9$ at $n=12$ and $n=30$ carries $(1+p)p^{3}\bmod9$, not $1+p$.
\end{remark}

\begin{proposition}[The glue $J$: what it is, the extension class, what computed data give, coincidence]\label{prop:heckeZglue}
\PROVED{} \STD{the Artin--Rees lemma, \cite[Cor.~10.10]{atiyah_macdonald}} for (5)(b)'s limit; \textbf{no stopping rule for that limit is asserted} (the plan's B11--B15,
with A3-M2's restriction).
\begin{enumerate}
\item \emph{(What $J$ is.)} $I_{\mathrm{fr}}=\Talg_{\mathrm{tors}}$, the additive torsion subgroup of the ring $\Talg$; so $J=\{t\in\Talg_{\mathrm{tors}}:t\cdot\mathrm{tors}\,H^{1}=0\}$ and
  $|\Talg_{\mathrm{tors}}|=|J|\cdot|\im(\Talg_{\mathrm{tors}}\to\Talg^{\mathrm{tors}})|$. Every $t\in J$ induces $\varphi_t\colon L\to\mathrm{tors}\,H^{1}$, and $t\mapsto\varphi_t$ is an injective
  $\Talg$-linear map $J\hookrightarrow\Hom_{\Talg}(L,\mathrm{tors}\,H^{1})=\bigoplus_\ell\Hom_{\Talg}(L/\ell^{e_\ell},\mathrm{tors}_\ell)$. $J_\midl\ne0$ only at maximal ideals in
  $\supp(L/\ell)\cap\supp(\mathrm{tors}_\ell)$.
\item \emph{(The extension class.)} Let $\varepsilon$ be the class of $0\to\mathrm{tors}\to H^{1}\to L\to0$ in $\Ext^{1}_{\Talg}(L,\mathrm{tors})$ and, for $K\ge e_\ell$,
  $\varepsilon_K$ that of $0\to\mathrm{tors}_\ell\to H^{1}_\ell/\ell^{K}\to L_\ell/\ell^{K}\to0$. (a) The $\ell$-part of $\varepsilon$ is the pull-back of $\varepsilon_K$ along
  $L_\ell\to L_\ell/\ell^{K}$, $H^{1}_\ell\cong(H^{1}_\ell/\ell^{K})\times_{L_\ell/\ell^{K}}L_\ell$; so $\varepsilon_\ell$ splits iff $\varepsilon_K$ does, for every $K\ge e_\ell$.
  (b) \emph{Split test}: in a basis adapted to $\mathrm{tors}_\ell\subset H^{1}_\ell/\ell^{K}$ each operator is $\smat{B_p}{0}{\beta_p}{c_p}$ ($B_p$ on $L/\ell^{K}$, $c_p$ on
  $\mathrm{tors}_\ell$, $\beta_p\in\Hom(L/\ell^{e_\ell},\mathrm{tors}_\ell)$ the derivation part); $\varepsilon_K$ splits for $\Talg'$ iff $hB_p-c_ph=\beta_p$ ($p\in\mathcal P$) has a
  solution $h\in\Hom(L/\ell^{e_\ell},\mathrm{tors}_\ell)$. Unsolvable for $\Talg'$ implies non-split for $\Talg$ (a certificate); solvable for $\Talg'$ implies split for $\Talg$
  only if $\mathcal P$ generates the image of $\Talg$ on $H^{1}_\ell/\ell^{K}$. (c) Split implies $J_\ell=0$; \textbf{the converse fails}: on $H=\Z_\ell^{2}\oplus\Fl$ let
  $x(a,b,c)=(0,\ell b,\bar a)$; $\Z_\ell[x]\cong\Z_\ell[x]/(x^{2}-\ell x)$ acts faithfully on $L=\Z_\ell^{2}$ by $\mathrm{diag}(0,\ell)$, so $I_{\mathrm{fr}}=0=J$, but no complement
  is $x$-stable. (d) If $\Talg^{\mathrm{fr}}_\midl$ acts on $L_\midl$ by scalars and $\mathrm{tors}_\midl$ is one-dimensional with the same residual character, then
  $t-\lambda(t)\in J_\midl$ for every $t$, $J_\midl$ is spanned by the $\beta(T_p)$, and \textbf{non-split $\Leftrightarrow J_\midl\ne0$} there.
\item \emph{(Reading $J$ from mod-$\ell^{K}$ data.)} Let $c:=c^{\mathrm{sat}}$ for $(L_\ell,\Talg_\ell)$, $K\ge e_\ell+c$, $A_K:=$ the image of $\Talg_{\Z_\ell}$ in
  $\End(H^{1}_\ell/\ell^{K})$ and $J_K:=\ker\bigl(A_K\to\End(L_\ell/\ell^{K})\times\End(\mathrm{tors}_\ell)\bigr)$. Then $J_\ell:=J\otimes\Z_\ell\to A_K$ is injective with image
  $J_K$. For a subalgebra $\Talg'$: $J\cap\Talg'\subseteq J\hookrightarrow\Hom_{\Talg'}(L/\ell^{e_\ell},\mathrm{tors}_\ell)$. \textbf{What computed data give}: (i) this
  \emph{sandwich} (a non-zero element of $J\cap\Talg'$ certifies $J_\ell\ne0$; record status \code{bounds}); (ii) $J_\ell=0$ \emph{exactly} where the supports separate
  ((1), Proposition~\ref{prop:heckeZsplit}) or where the upper bound has order $1$ (status \code{zero}); (iii) $J_\ell=\ker(A'_K\to\cdots)$ \textbf{only under} the
  mixed generation statement of Remark~\ref{rmk:heckeZcbthree}, which is \CONJ: so $J$ is never recorded as exact.
\item \emph{(Coincidence certificates.)} If $\varepsilon_K$ is non-split for some $\Talg'$, or $J\cap\Talg'\ne0$, or $\Hom_{\Talg}(L,\mathrm{tors}_\ell)\ne0$, then some maximal
  ideal of $\Talg$ contains $I_{\mathrm{fr}}$ and $I_{\mathrm{tors}}$: a torsion system \emph{equal for all $p$} to a free one. If moreover $\mathrm{tors}_{\ell,\midl'}$ is
  one-dimensional over $\Fl$ for the $\Talg'$-block $\midl'$ on which the non-split extension or the witness was found, then $\Talg$ acts on it by one character, a residual
  system of $\Talg^{\mathrm{fr}}$. (Coincidence is a $\operatorname{Spec}$ statement; this does not compute $J$.)
\item \emph{(The fibre product and the local rings.)} (a) With $C_{\mathrm{ft}}:=\Talg/(I_{\mathrm{fr}}+I_{\mathrm{tors}})$, the torsion--free congruence ring, finite and supported
  exactly on the common points: $\Talg/J\cong\Talg^{\mathrm{fr}}\times_{C_{\mathrm{ft}}}\Talg^{\mathrm{tors}}$. So $\Talg$ is determined by $\Talg^{\mathrm{fr}}$, $\Talg^{\mathrm{tors}}$, the
  two maps $\Talg^{\mathrm{fr}}\to C_{\mathrm{ft}}\leftarrow\Talg^{\mathrm{tors}}$ (not the ring $C_{\mathrm{ft}}$ alone), the $\Talg/J$-module $J$ and the class of the square-zero
  extension $0\to J\to\Talg\to\Talg/J\to0$ (the trivial extension $\Talg/J\oplus J$ is in general a different ring). (b) $\Talg_{\Z_\ell}=\varprojlim_KA_K(H^{1}_\ell/\ell^{K})$, and
  there is $c'$ with $\Talg_{\Z_\ell}\cap\ell^{K}\End(H^{1}_\ell)\subseteq\ell^{K-c'}\Talg_{\Z_\ell}$; for $K\ge e_\ell$, $E_K:=\{f:f(H^{1}_\ell)\subseteq\ell^{K}H^{1}_\ell\}\subseteq
  \ell^{K-e_\ell}\End(H^{1}_\ell)$; so $\Talg_{\Z_\ell}/\ell^{K-c'-e_\ell}$ is a quotient of $A_K$: the local rings at torsion-touching ideals, glue included, are a finite
  computation once $c'$ is known \emph{and} $A_K$ itself, the image of all of $\Talg$, is known --- i.e.\ under Remark~\ref{rmk:heckeZcbthree} when only $\Talg'$ is computed.
\end{enumerate}
\end{proposition}

\begin{proof}
(1) $\Talg^{\mathrm{fr}}$ is torsion-free of rank $\rank\Talg$ ($\mathrm{tors}\otimes\Q=0$), so $I_{\mathrm{fr}}=\ker(\Talg\to\Talg^{\mathrm{fr}})$ is finite, hence
$\subseteq\Talg_{\mathrm{tors}}$, and $\Talg_{\mathrm{tors}}\subseteq I_{\mathrm{fr}}$ as $\Talg^{\mathrm{fr}}$ is torsion-free. $t\in I_{\mathrm{fr}}$ maps $H^{1}$ into $\mathrm{tors}$, and
$t\in I_{\mathrm{tors}}$ kills $\mathrm{tors}$, so $t\in J$ factors through $L$; $\varphi_t=0$ means $t=0$. For $s\in\Talg$, $\varphi_{st}=s|_{\mathrm{tors}}\circ\varphi_t$,
$\varphi_{ts}=\varphi_t\circ s|_L$ and $st=ts$, so $\varphi_t$ is $\Talg$-linear and so is $t\mapsto\varphi_t$. The support statement is
Proposition~\ref{prop:heckeZsplit} ($\Hom$ between disjoint supports vanishes). (2)(a) $x\mapsto(\bar x,\pi x)$ is a map of extensions of $L_\ell$ by $\mathrm{tors}_\ell$
(identity on $\mathrm{tors}_\ell=\mathrm{tors}_\ell/\ell^{K}$ for $K\ge e_\ell$, and on $L_\ell$), hence an isomorphism (five lemma); a section of $\varepsilon_K$ pulls back to one of
$\varepsilon_\ell$, and a section of $\varepsilon_\ell$ reduces to one of $\varepsilon_K$. (b) A complement is the graph of an $h$, and its $T_p$-stability is the displayed
equation. (c) For $t\in J$ and a $\Talg$-stable complement $L'$: $t(L')\subseteq L'\cap\mathrm{tors}=0$ and $t(\mathrm{tors})=0$. The example: $x^{2}(a,b,c)=(0,\ell^{2}b,0)$, so
$x(x-\ell)=0$ and there is no relation of degree $1$ (evaluate on $(1,0,0)$ and $(0,1,0)$); $x$ acts on $L$ by $\mathrm{diag}(0,\ell)$, whose minimal polynomial is $x(x-\ell)$, so
$\Z_\ell[x]\to\End L$ is injective; a complement $\{(a,b,h(a,b))\}$ is $x$-stable only if $\bar a=h(0,\ell b)=0$ for all $a$. (d) On the block the inner derivations are
$(c_p-\lambda_p)h\equiv0$; for $t\in\Talg_\midl$, $t-\lambda(t)$ is $0$ on $L_\midl$ and on $\mathrm{tors}_\midl$, so lies in $J_\midl$, with $\varphi_{t-\lambda(t)}=\beta(t)$; every
$j\in J_\midl$ has $\lambda(j)=0$; $\beta$ is a derivation, so its values are spanned by those on the generators $T_p$; and the split equation reads $\beta_p=0$ for all $p$.
(3) Injective: $t\in J_\ell$ zero on $H^{1}_\ell/\ell^{K}$ has $t(H^{1}_\ell)\subseteq\ell^{K}H^{1}_\ell\cap\mathrm{tors}_\ell=0$ ($K\ge e_\ell$). Onto $J_K$: lift $a\in J_K$ to
$t\in\Talg_{\Z_\ell}$; $t|_{\mathrm{tors}}=0$ and $t|_L\in\Talg_\ell\cap\ell^{K}\End(L_\ell)\subseteq\ell^{K-c}\Talg_\ell\subseteq\ell^{e_\ell}\Talg_\ell$
(Proposition~\ref{prop:saturation}(1)); choose $s\in\Talg_{\Z_\ell}$ with $s|_L=\ell^{-e_\ell}t|_L$ ($\Talg_{\Z_\ell}\to\Talg_\ell$ is onto). Then $t-\ell^{e_\ell}s\in J_\ell$, and
$\ell^{e_\ell}s$ is $0$ on $H^{1}_\ell/\ell^{K}$: for a free complement $L'$, $\ell^{e_\ell}H^{1}_\ell=\ell^{e_\ell}L'$, and for $y\in L'$, $s(y)=w+\tau$ ($w\in L'$, $\tau$
torsion) gives $\ell^{e_\ell}s(y)=\ell^{e_\ell}w$, whose image in $L$ is $t|_L(\bar y)\in\ell^{K}L$, so $\ell^{e_\ell}w\in\ell^{K}L'$. So $t-\ell^{e_\ell}s$ has image $a$. The
sandwich is (1) with $\Hom_{\Talg}\subseteq\Hom_{\Talg'}$; (ii) is (1) and the order bound; (iii) is the case $A'_K=A_K$. (4) Non-split for $\Talg'$ implies non-split for
$\Talg$, which by Proposition~\ref{prop:heckeZsplit} (contrapositive) implies a common point; $J\ne0$ or $\Hom\ne0$ implies one by (1). A one-dimensional module has one maximal
ideal. (5)(a) $\Talg/(I\cap I')\to\Talg/I\times_{\Talg/(I+I')}\Talg/I'$ is injective, and onto: given $(x+I,y+I')$ with $x-y=i+i'$, $z:=x-i=y+i'$ maps to the pair; $I_{\mathrm{fr}}\cap
I_{\mathrm{tors}}=J$. (b) The constant $c'$ is Artin--Rees for $\Talg_{\Z_\ell}\subseteq\End(H^{1}_\ell)$ over $\Z_\ell$. With the bound on $E_K$ below,
$\ell^{K}\Talg_{\Z_\ell}\subseteq\ker(\Talg_{\Z_\ell}\to A_K)=\Talg_{\Z_\ell}\cap E_K\subseteq\ell^{K-c'-e_\ell}\Talg_{\Z_\ell}$, so the kernels define the $\ell$-adic topology and the
inverse limit is the completion, $\Talg_{\Z_\ell}$ itself. For $f\in E_K$, $K\ge e_\ell$: $f(H^{1}_\ell)\subseteq
\ell^{K}H^{1}_\ell=\ell^{K}L'$, and $g:=\ell^{-(K-e_\ell)}f$ is a well-defined endomorphism with values in $\ell^{e_\ell}L'$ ($L'$ free), so $f\in\ell^{K-e_\ell}\End(H^{1}_\ell)$; the
inclusion is strict in general (on $\Z/\ell\oplus\Z/\ell^{2}$, $e_1\mapsto\ell e_2$ lies in $E_1$, not in $\ell\End$).
\end{proof}

\noindent
Checks (the plan's \S6.4, level $1$, $\ell=5$): at $(1,18)$ and $(1,20)$ $T_5$ separates the torsion from the free part --- split, $J_5=0$, certified; at $(1,26)$ and $(1,30)$
the torsion meets a non-ordinary cuspidal point, the extension splits for the computed $\Talg'$ (every one of the $25$ candidate torsion lines), $J_5\cap\Talg'=0$ and $|J_5|\le25$
(\code{bounds}); the exact $J_5$ needs Remark~\ref{rmk:heckeZcbthree}. The torsion image $\mathrm{tors}_\ell\subset H^{1}/\ell^{K}$ is canonical but \textbf{not determined by the
Hecke module $H^{1}/\ell^{K}$ alone} when $\Hom_{\Talg}(\mathrm{tors},L/\ell^{e_\ell})\ne0$: either the verdict is checked for every candidate, or the torsion cocycles of
Proposition~\ref{prop:torsionhecke}(3) are exported with the matrices. Theorem~\ref{thm:descentcert}(7)(a) leaves the extension class open; by (2)(a) it is determined by
finite data, $\varepsilon_K$ for any $K\ge e_\ell$, and from computed operators (2)(b) certifies a non-split verdict, a split one only under Remark~\ref{rmk:heckeZcbthree}
(A7-m3).

\begin{remark}[CB-2: the $\inv$-refinement of the bound on $J$]\label{rmk:heckeZcbtwo}
\CONJ{} (the plan's CB-2, a sketch; recorded, used by nothing.) At odd $\ell$ the involution $\inv$ commutes with $\Talg$ (Obligation~\ref{obl:iota}), $\varphi_t$ respects the
$\inv$-signs, and at a cuspidal common point $L_\midl=L^{+}_\midl\oplus L^{-}_\midl$ with one sign matching the torsion's; so $|J_\midl|\le|\Hom_{\Talg,\inv}(L_\midl,\mathrm{tors}_\midl)|$
(e.g.\ $|J_5|\le5$ at $(1,26)$). It needs $\inv$ on the $\Z/\ell^{K}$ handle (\Espec{sec:lvdepth2} has it over $\Z/\ell^{k}$) and the sign of the torsion class; not run.
\end{remark}

\begin{remark}[CB-3: mixed generation]\label{rmk:heckeZcbthree}
\CONJ{} (the plan's CB-3, A3-M2; \textbf{nothing is certified under it}, and no hypothesis structure of the Lean lane may carry it silently.) For the computed set
$\mathcal P$ (with an explicit bound), the image $A'_K$ of $\Talg'=\Z[X_p:p\in\mathcal P]$ on the whole $H^{1}_\ell/\ell^{K}$ --- free part, torsion and glue --- equals the image
$A_K$ of $\Talg_{\Z_\ell}$; equivalently, the cotangent space of $\Talg_\midl/\ell$, glue included, is spanned by the $T_p-\lambda(T_p)$, $p\in\mathcal P$. It is a Sturm-type
statement for the mixed module; by Proposition~\ref{prop:heckeZglue}(2)(d), at a scalar block the derivation $\beta$ vanishes once it vanishes on a set of generators of
that cotangent space. Under it, $J_\ell$ would be the kernel $J_K$ of the computed algebra (Proposition~\ref{prop:heckeZglue}(3)(iii)) and the local rings at
torsion-touching ideals would be computed (Proposition~\ref{prop:heckeZglue}(5)(b)). A non-split verdict does not need it (Proposition~\ref{prop:heckeZglue}(4)).
\end{remark}

\begin{theorem}[Descent of the Hecke structure]\label{thm:heckeZdescent}
(The plan's Thm.~3.1.) Write $\Z_{(6)}$ for the semilocalisation at $2$, $3$.
\begin{enumerate}
\item \emph{(Fracture.)} \PROVED{} For every finitely generated abelian group $X$, $X\to X[1/6]\times X_{(6)}$ is injective with image $X[1/6]\times_{X\otimes\Q}X_{(6)}$
  ($\Z_{(6)}[1/6]=\Q$); for $X=\Talg$ this is a fibre product of rings. Images commute with flat base change, so $\Talg\otimes R=\im(\mathcal H_N\otimes R\to\End_R(H^{1}\otimes R))$
  for $R\in\{\Z[1/6],\Z_{(2)},\Z_{(3)},\Z_\ell\}$: the algebra form of Propositions~\ref{prop:fracture} and~\ref{prop:cohglue}.
\item \emph{(Spectra.)} \PROVED{} As sets, $\operatorname{Spec}\Talg=\operatorname{Spec}\Talg_\Q\sqcup\bigsqcup_\ell V(\ell)$; every prime of characteristic $\ell$ is maximal;
  $\operatorname{Spec}\Talg[1/6]$ is the open complement of $V(2)\sqcup V(3)$; a component $P_i$ specialises to $\midl$ iff $P_i\subseteq\midl$
  (Proposition~\ref{prop:heckeZthrough}).
\item \emph{(Local--global.)} \PROVED{} $\Talg^{\mathrm{fr}}=\{x\in\Talg_\Q:x\in\Talg_\ell\ \forall\ell\}$, and for $\vartheta\in\Talg^{\mathrm{fr}}$ with $\Q[\vartheta]=\Talg_\Q$,
  $\Talg_\ell=\Z_\ell[\vartheta]$ at every $\ell$ with $\ell^{2}\nmid\operatorname{disc}_\vartheta$ (Proposition~\ref{prop:heckeZetale}(3)). So $\Talg^{\mathrm{fr}}$ is determined
  by $f_\vartheta$ and the finitely many $\Talg_\ell$, $\ell^{2}\mid\operatorname{disc}_\vartheta$; and $\Talg'=\Talg^{\mathrm{fr}}$ iff $\Talg'_\ell=\Talg_\ell$ for every $\ell$
  (Lemma~\ref{lem:heckeZbridge}(3)).
\item \emph{(Above $2$ and $3$: what the truncated rings see.)} \PROVED{} given the model statements of Theorem~\ref{thm:descentcert}(3), (4) and
  Propositions~\ref{prop:isomodeltree}, \ref{prop:certifiedeell}. For $\ell\in\{2,3\}$ the inputs --- $L\otimes\Z_{(\ell)}$ and $\mathrm{tors}\,H^{1}_{(\ell)}$ with the operators ---
  are computed on the $\ell$-local models, and:
  \begin{enumerate}[label=(\alph*)]
  \item the maximal ideals over $\ell$ and their residue fields are given by $A'_K$ for \textbf{every} $K\ge1$, once generation modulo $\ell$ holds
    (Propositions~\ref{prop:heckeZmax}, \ref{prop:heckeZgen}); without it a computation sees the primes of $\Talg'$, onto which $\operatorname{Spec}\Talg$ maps;
  \item which components pass through each maximal ideal: at $K>v_\ell\operatorname{Res}(f_i,g_i)$ (Proposition~\ref{prop:heckeZthrough}), with no generation needed;
  \item $\Talg_\ell$, hence the local rings and the $\ell$-parts of the congruence modules, always under generation: in the monogenic case at $K=1$ when $\Talg_\ell$ is
    saturated on $L_\ell$ and at $K=c^{\mathrm{sat}}+1$ when it is not (Proposition~\ref{prop:heckeZmono}); at $K\ge e+c^{\mathrm{sat}}$ otherwise
    (Proposition~\ref{prop:heckeZlocal});
  \item $\Talg^{\mathrm{tors}}\otimes\Z_\ell$ exactly, as the image on the finite module (any $K$ at least its exponent), under a generation certificate for the torsion
    (Proposition~\ref{prop:heckeZrat}(5)(c)) over $\Z/\ell^{e}$, $e\ge2$ in general ($\Z/4$ at $(1,10)$, $(\Z/4)^{2}$ at $(2,6)$); its input is the torsion \emph{module} with its
    operators (Theorem~\ref{thm:descentcert}(3)), not the layer polynomials (at $(1,10)$ the layers of $\mathrm{tors}_2=\Z/2\oplus\Z/4$ read ``$x$'' for every $T_p$, while $T_5$
    acts by $2$ and $\ker T_2\cong\Z/4$);
  \item \textbf{at level $1$, $\ell\in\{2,3\}$}: the maximal ideals above $\ell$ are the $\midl_\psi$ of Proposition~\ref{prop:heckeZlevelone23}(2), \textbf{known without any
    generation hypothesis}, and $\dim A_1(L/\ell)\le3$, so Theorem~\ref{thm:heckeZcert}(C3) fails at $2$, $3$ for every $d\ge4$; the $2$- and $3$-adic structure of
    $\Talg^{\mathrm{fr}}$ comes from Proposition~\ref{prop:heckeZeiscert} or from $A_K/\ell A_K$ at $K\ge c^{\mathrm{sat}}+1$ under a generation theorem.
  \end{enumerate}
  The same holds verbatim at $\ell\ge5$, with inputs on the $\mathsf O$-piece (Propositions~\ref{prop:levelorbit}(2), \ref{prop:heilbronn}; the free part modulo
  $\ell^{K'}$ of Theorem~\ref{thm:descentcert}(4)). \textbf{So ``everything above $p$ survives in the truncated rings'' holds with the precise meaning (a)--(d); its only
  extra hypothesis is generation at $p$, and the same hypothesis is needed at every $\ell\ge5$.}
\item \emph{(Torsion and glue.)} \PROVED{} $\operatorname{Spec}\Talg$, its components and its isolated torsion points do not depend on $J$
  (Proposition~\ref{prop:heckeZspec}(2)). Whether a maximal ideal of $\Talg^{\mathrm{tors}}$ \emph{equals} one of $\Talg^{\mathrm{fr}}$ is decided by agreement on \textbf{all}
  $T_p$: a disagreement on a computed operator certifies that they differ (``differs from system $X$ at $T_p$'' is a certificate, ``agrees on $\mathcal P$'' is not).
  Equality is certified (i) at level $1$, $\ell\in\{2,3\}$, for every cell, by $T_\ell$ (Proposition~\ref{prop:heckeZlevelone23}(4), (5): the torsion point is $\midl_x$, common
  with the cuspidal part exactly when $S_k\ne0$); (ii) everywhere by Proposition~\ref{prop:heckeZglue}(4). The residual case ``split for $\Talg'$, common on $\mathcal P$'' is
  not decided. The local rings at maximal ideals touching both parts are $\varprojlim_KA_K(H^{1}_\midl/\ell^{K})$ (Proposition~\ref{prop:heckeZglue}(5)(b)), computable from
  computed data only under Remark~\ref{rmk:heckeZcbthree}; without it $J$ is \code{bounds} or \code{zero} (Proposition~\ref{prop:heckeZglue}(3)), never exact.
\end{enumerate}
\end{theorem}

\begin{proof}
(1) Mayer--Vietoris for $\Z=\Z[1/6]\times_\Q\Z_{(6)}$, tensored with $X$ (flat): the kernel consists of elements killed by a power of $6$ and by an integer prime to $6$; and
$\Talg\otimes R\to\End(H^{1})\otimes R=\End_R(H^{1}\otimes R)$ stays injective with the base-changed image. (2) $\Talg/P$ is a finite domain when $P\cap\Z=\ell\Z$, hence a
field; the rest is Proposition~\ref{prop:heckeZspec}(3). (3) A lattice in a finite-dimensional $\Q$-vector space is the intersection of its $\ell$-adic closures; then
Proposition~\ref{prop:heckeZetale}(3) and Lemma~\ref{lem:heckeZbridge}. (4) Propositions~\ref{prop:heckeZmax}--\ref{prop:heckeZthrough} with $M=L_\ell$ from the local model;
the model statements say that the operators there induce $T_p\otimes\Z_{(\ell)}$; (e) is Proposition~\ref{prop:heckeZlevelone23}. (5) Proposition~\ref{prop:heckeZspec}(2); a
maximal ideal of $\Talg$ is determined by the residues of all its generators; the certificates are Propositions~\ref{prop:heckeZlevelone23}, \ref{prop:heckeZglue}.
\end{proof}

\noindent
\emph{What is not available at $2$ and $3$ today} (as written at v1.25.0-pre.0, before EXPORT and UP-Z; what a record reads at $2$ and $3$ is item~8, and for a \code{full} record at $N>1$ item~8~(d$'$)(v), by the binary that wrote it; v1.26.0-pre.11, confirm-referee C1). $U_p$, $p\mid N$, on the local models over $\Z_{(\ell)}$; the $\inv$-splitting at $\ell=2$ (so multiplicity one cannot be read
per sign there); and \textbf{any export of the operators on $H^{1}_{(\ell)}/\ell^{K}$}, at every level --- the console's $\Z/\ell^{k}$ mode needs $\ell\ge5$, and the descent's
$2$- and $3$-local reads are internal. So (4)(a)--(d) at $2$ and $3$ are reachable today only through (4)(e) (level $1$) and Proposition~\ref{prop:heckeZeiscert}; the full
algebra at level $N>1$ stays partial at $2$ and $3$ until $U_p$ exists there.

\begin{theorem}[The global certificate]\label{thm:heckeZcert}
\PROVED{} (the plan's Thm.~4.1 and Prop.~4.2.) Let $\mathcal P$ be a finite set of primes, $\Talg':=\Z[X_p:p\in\mathcal P]\subseteq\Talg^{\mathrm{fr}}$, and $\vartheta\in\Talg'$.
Let $d$ be the a-priori rank of Proposition~\ref{prop:heckeZrat}(3) (resp.\ $d^{\mathrm{an}}$ of (3$'$), with $\mathcal P\not\ni p$ for $p\mid N$ and $\Talg^{\mathrm{an,fr}}$ for
$\Talg^{\mathrm{fr}}$ throughout). Suppose:
\begin{itemize}
\item[\textbf{(C1)}] \emph{rational}: $f_\vartheta:=$ the squarefree part of the characteristic polynomial $\chi_\vartheta$ of $\vartheta$ on $L$ (assembled over $\Z$ by CRT)
  has $\deg f_\vartheta=d$;
\item[\textbf{(C2)}] \emph{bad set}: for one or more $\vartheta_j\in\Talg'$, \textbf{each passing (C1) itself}, with $D_j:=\operatorname{disc}(f_{\vartheta_j})$, a factorisation
  $D_{j_0}=\pm\prod_iq_i^{e_i}\cdot\mathsf c$ of \textbf{one} of them whose factors $q_i$ are \textbf{certified primes}, and the integer $\mathsf c\ge1$ the \emph{cofactor};
  $\mathsf S$ is the set of \textbf{all} the $q_i$, except that a $q_i$ with $v_{q_i}(D_j)\le1$ for some $j$ ($v_{q_i}(D_j)$ read by residues) \textbf{may} be left out of $\mathsf S$
  (v1.25.0-pre.0b, A7-m4: no gcd of the big integers $D_j$ is formed; v1.25.0-pre.0c, A7b-m1: this drop rule is the only way the other $\vartheta_j$ enter);
\item[\textbf{(C3)}] \emph{local}: for each $\ell\in\mathsf S$, $\dim_{\Fl}A'_1(L/\ell)=d$.
\end{itemize}
Then:
\begin{enumerate}
\item $f_\vartheta(\vartheta)=0$ on $L$, $f_\vartheta$ is the minimal polynomial of $\vartheta$, $\Talg_\Q=\Q[\vartheta]$, and $\Talg_\Q$ is \textbf{reduced} --- for a full
  $\vartheta$ (with $U_p$) as for an anemic one; no semisimplicity hypothesis is used.
\item $\Talg'_\ell=\Talg_\ell$ for every $\ell\nmid\mathsf c$, and for every $\ell$ when $\mathsf c$ is squarefree (in particular $\mathsf c=1$). Then $\Talg'=\Talg^{\mathrm{fr}}$, and
  every prime ideal, component, local ring and congruence module of $\Talg^{\mathrm{fr}}$ is a finite computation from the structure constants of $\Talg'$ alone.
\item \emph{(Termination of the rational step.)} (a) If the systems of $\Talg_\Q$ are separated by $\mathcal P$, then $\vartheta=\sum_{p\in\mathcal P}c_pX_p$ fails to be
  primitive only for $c$ on at most $\binom d2$ proper hyperplanes, so some $c\in\{0,\dots,\binom d2\}^{\mathcal P}$ works; a finite separating $\mathcal P$ exists (the $d$
  systems are distinct ring maps out of $\mathcal H_N$, resp.\ $\mathcal H^{\mathrm{an}}_N$; one $p$ per pair), and no a-priori bound on $p$ is needed, since the search stops
  on $\deg f_\vartheta=d$, which it can test. (b) \STD{\cite{deligne_weil1}} The roots of $\chi_\vartheta$ are bounded by $R:=\sum_p|c_p|(1+p^{k-1})$ and its coefficients by
  $\binom{b_1}{j}R^{j}$, so CRT over word primes stops after finitely many primes. (c) At a fixed $\ell\in\mathsf S$, enlarging $\mathcal P$ reaches (C3) after finitely many
  steps \textbf{iff} $\Talg_\ell$ is saturated on $L_\ell$ (Proposition~\ref{prop:heckeZgen}(2), (3)); the algorithm cannot know which case holds and uses a cap.
\end{enumerate}
\end{theorem}

\begin{proof}
(1) The minimal polynomial $m_\vartheta$ of $\vartheta$ on $L\otimes\Q$ has the irreducible factors of $\chi_\vartheta$, so $f_\vartheta\mid m_\vartheta$ and $\deg m_\vartheta\ge d$;
and $\deg m_\vartheta=\dim\Q[\vartheta]\le\dim\Talg_\Q=d$ (Proposition~\ref{prop:heckeZrat}(3), the perfect pairing, no reducedness used; resp.\ (3$'$)). Hence
$m_\vartheta=f_\vartheta$, squarefree, and $\Q[\vartheta]=\Talg_\Q$ is reduced. (Taking $f_\vartheta$ to be the squarefree part of $\chi_\vartheta$ matters: ``$f$ squarefree,
$f(\vartheta)=0$, $\deg f=d$'' for an arbitrary $f$ would admit $f=m_\vartheta h$.) (2) By (1) for $\vartheta_j$, $\rank\Talg'=d=\rank\Talg^{\mathrm{fr}}$ and
$\operatorname{disc}\Talg'\mid D_j$ (Proposition~\ref{prop:heckeZetale}(1) for $\Z[\vartheta_j]\subseteq\Talg'$). Let $\ell\notin\mathsf S$. If $\ell=q_i$ was left out of $\mathsf S$,
then $v_\ell(D_j)\le1$ for some $j$ by the drop rule of (C2). Otherwise $\ell\ne q_i$ for every $i$; if $\ell^{2}\mid D_j$ for all $j$, then
$\ell^{2}\mid D_{j_0}=\pm\prod q_i^{e_i}\mathsf c$, so $\ell^{2}\mid\mathsf c$: excluded when $\ell\nmid\mathsf c$ or $\mathsf c$ is squarefree. So
$v_\ell(\operatorname{disc}\Talg'_\ell)\le v_\ell(D_j)\le1$ for some $j$, and $\Talg'_\ell=\Talg_\ell$ by Proposition~\ref{prop:heckeZetale}(2). At $\ell\in\mathsf S$ it is
Proposition~\ref{prop:heckeZgen}(1) through (C3). $\Talg'=\Talg^{\mathrm{fr}}$ is Lemma~\ref{lem:heckeZbridge}(3). (3)(a) Over $\overline\Q$ the reduced $\Talg_\Q$ has $d$
characters; $\vartheta$ is primitive iff it separates every pair, and for a pair separated by some $X_p$ the failing $c$ form a proper hyperplane; a non-zero product of
$\binom d2$ linear forms does not vanish on the whole grid $\{0,\dots,\binom d2\}^{\mathcal P}$. Distinct characters of $\mathcal H_N\otimes\overline\Q$ differ on a
generator. (b) Deligne's bound for the cusp eigenvalues, $|a_p|\le2p^{(k-1)/2}\le1+p^{k-1}$; Eisenstein eigenvalues $\psi(p)+\bar\psi(p)p^{k-1}$; at $p\mid N$,
$|a_p|\le p^{k-1}$; then Vieta. (c) Proposition~\ref{prop:heckeZgen}(2), (3).
\end{proof}

\noindent
\emph{The cofactor test and the big-integer rule} (owner ruling Q1(a) of 2026-11-04, \emph{residue-only}). Only the square part of $\operatorname{disc}\Talg'$ matters. A
cofactor $\mathsf c<B^{3}$ with no prime factor $\le B$ that is neither prime nor a perfect square is a product of two distinct primes, hence squarefree and harmless;
each step is single-word arithmetic for $B=10^{6}$ ($B^{3}<2^{64}$): ``no factor $\le B$'' by the trial division itself, ``not prime'' by a Fermat witness
$a^{\mathsf c-1}\not\equiv1\pmod{\mathsf c}$, ``not a square'' by an exact integer square root. A cofactor that cannot be cleared leaves $\Talg^{\mathrm{fr}}$ certified at every
$\ell\nmid\mathsf c$ (record status \code{partial(unfactored)}); a factor whose primality is not proved is part of the cofactor. Under Q1(a) the C++ build certifies what
residues and single-word arithmetic decide: $\chi_\vartheta$, discriminants and resultants as exact values (residues modulo word primes, a certified bound, the limb output of
\S\ref{sec:recordschema} item~7~(l)); $v_q(D)$ for word-size $q$ --- from $D\bmod q^{e+1}$ while $q^{e+1}<2^{64}$, and in general by the cofactor route: reconstruct $X$ from
the residues $Dq^{-e}\bmod p_i$, check $q^{e}X\equiv D$ modulo word primes whose product exceeds the certified bound, and read $X\bmod q\ne0$ from its mixed-radix digits; the
cofactor exactly (its residues are $D\prod q_i^{-e_i}$); Pratt certificates for $q<2^{64}$; the cofactor test. A factor above $2^{64}$ is \code{probable} in the record and part
of the cofactor; the Lean kernel may check its Pratt certificate, but that result lives beside the record. Exact division, gcds, modular powers and square roots of big
integers are multi-precision arithmetic and stay out of the certified build ([R-Q2] unchanged). Two consequences (the plan's \file{01-math-core.md} \S4 Q1(a), restored at
A7-m4): \textbf{the gcd over several $\vartheta_j$ is used only through its word-size primes} ((C2) factors one $D_{j_0}$; $\ell^{2}\mid\gcd_jD_j$ implies $\ell^{2}\mid D_{j_0}$, so the
proof of (2) needs nothing more), and \textbf{$g_{\mathcal P}$ is never formed as a big gcd}: its $\ell$-parts for word-size $\ell$ are read by residues modulo $\ell^{K}$
($v_\ell(g_{\mathcal P})=\min_{p\in\mathcal P}v_\ell c_{\mathcal S',p}(\lambda(X_p))$), and as an upper bound one $c_{\mathcal S',p_0}(\lambda(X_{p_0}))$ with its cofactor serves. \emph{As built} (v1.25.0-pre.2, Z0 referee F5;
\file{Integral/HeckeZRational.cpp}, \code{GcdOf}): HZ-RAT takes (C2) on the gcd of the $D_j$ over \emph{every} primitive $\vartheta_j$ it found, and $g_{\mathcal P}$ on the gcd
of the $c_{\mathcal S',p}(\lambda(X_p))$ over the primitive unit $\vartheta=T_p$, in both cases through word-size primes only ($v_q$ of the gcd is $\min_jv_q(\cdot)$, read by residues);
its cofactor is the smallest base cofactor, kept only if the residue route above shows that it divides every term and no certified prime divides it, otherwise
\code{cofactor\_test = upper\_bound} (a bound, never \code{complete}). This is the same certificate: $v_\ell(\gcd_jD_j)=\min_jv_\ell(D_j)$, so \code{complete} on the gcd discharges
(C2)'s drop rule for every $j$ at once, and nothing but residue arithmetic enters. The mathematics and the soundness are the same under every option of Q1: a wrong
C++ big integer can only make a Lean check fail.

\paragraph{Two forms, and what is reachable today.} The theorem is used in two forms, both proved, differing only in the a-priori rank: \textbf{full} ($\mathcal P$ contains
the $p\mid N$, $X_p=U_p$; $d=\dim M_k(\Ga_0(N))$; conclusion $\Talg'=\Talg^{\mathrm{fr}}$) and \textbf{anemic} ($p\nmid N$ only; $d^{\mathrm{an}}$; conclusion
$\Talg'=\Talg^{\mathrm{an,fr}}$). ``Certifiable'' means certifiable by a proved statement once its inputs exist; what the shipped tools reach \textbf{today} is narrower,
because every step that reads a module --- (C3), Propositions~\ref{prop:heckeZmono}--\ref{prop:heckeZthrough}, the at-root read of
Proposition~\ref{prop:heckeZcong}(5) --- is stated for $L_\ell$ and needs the matrices of the operators on $L/\ell^{K}$ (for the glue, on the whole $H^{1}_{(\ell)}/\ell^{K}$).
The console's $\Z/\ell^{k}$ mode at level $1$, $\ell\ge5$, delivers $H^{1}(\Ga;\Vcoef\otimes\Z/\ell^{K})=H^{1}/\ell^{K}$, \textbf{not} $L/\ell^{K}$; the two differ exactly at the
cell's torsion primes ($0\to\mathrm{tors}_\ell/\ell\to H^{1}/\ell\to L/\ell\to0$; at level $1$ every $\ell\ge5$ with $\ell\le n+1$ is one,
$H^{1}[\ell]\cong H^{0}(\Ga;\mathrm{Sym}^{n}\Fl^{2})\ne0$), where $\dim A'_1(H^{1}/\ell)\ge\dim A'_1(L/\ell)$, possibly strictly, and $L/\ell^{K}$ is not recoverable from the Hecke
module $H^{1}/\ell^{K}$ alone. So \textbf{a certified status for (C3), Propositions~\ref{prop:heckeZmono}--\ref{prop:heckeZthrough} or the at-root read is never read from
$H^{1}(\Ga;\Vcoef/\ell^{K})$ at a torsion prime of the cell}. With $U_p$ over $\Z$ not built (the state when this list was written, v1.25.0-pre.0, before EXPORT and UP-Z: for the \code{anemic} line (C3) and the lower bound are read since EXPORT, e.g.\ the \code{anemic} $(11,0)$, item~8; for the \code{full} line see that line; v1.26.0-pre.11, confirm-referee C1):
\begin{itemize}\setlength{\itemsep}{0pt}
\item \emph{level $1$}: the full algebra ($=\Talg^{\mathrm{an}}$) is certifiable by this theorem and Propositions~\ref{prop:heckeZcong}, \ref{prop:heckeZeiscert}, and
  \textbf{reachable today} at every $\ell\ge5$ not dividing $|\mathrm{tors}\,H^{1}|$; at a torsion prime $\ell\ge5$ only \textbf{blockwise}, at the blocks outside
  $\supp(\mathrm{tors}_\ell)$ (the torsion support is non-ordinary, so every ordinary block qualifies; the per-block (C3) compares with the per-block rank from the Hensel
  factors attributed by Proposition~\ref{prop:heckeZthrough}); at $\ell\in\{2,3\}$ through Proposition~\ref{prop:heckeZeiscert} only;
\item \emph{level $N>1$, anemic}: certifiable by (C1)--(C3) with $d^{\mathrm{an}}$, and the congruence module between the cuspidal systems and the anemic Eisenstein system by
  Proposition~\ref{prop:heckeZcong}(4) when $e^{\mathrm{an}}_k(N)=1$; Proposition~\ref{prop:heckeZeiscert}(1) is not available (Sturm needs every $T_n$, the $U_p$-words
  included). \textbf{Reachable today: (C1), (C2) and the upper bound $g_{\mathcal P}$ only}; (C3) and the lower bound after the export of the whole module
  $H^{1}_{(\ell)}/\ell^{K}$ with the matrices of every computed $T_p$ by cochain-level lifts (Proposition~\ref{prop:heckeZuct}(3)) --- the free quotient alone is not enough;
\item \emph{level $N>1$, full} (at v1.25.0 this line read ``nothing is certified but two tags''; v1.26.0-pre.11, confirm-referee C1): what a record writes is item~8~(d$'$)(ii), (v), by the binary that wrote it. (C1) for the full algebra needs no $U_p$ over $\Z$: the $\vartheta$ search reads $U_p$ at the word primes by Proposition~\ref{prop:heilbronn}, so \code{G1.C1.status} can be \code{certified} ($67$ of the $85$ records of Z1's recompute, a count at this version), and the words of G2 at $\ell\ge5$ are read on that operator; the fields at $\ell\in\{2,3\}$ follow item~8~(d$'$)(v), according to whether the binary carries the route over $\Z$. Two tags are computable without any lattice from $\chi_\Q(U_p)$ assembled by CRT from the level-$N$ route over word primes:
  \code{upper\_bound\_certified} (Proposition~\ref{prop:heckeZcong}(2) for $\vartheta=U_p\in\Talg^{\mathrm{fr}}$; e.g.\ $(2,6)$: $C(f,E^{(1)})=0$, $C(f,E^{(2^{7})})\mid2^{3}\cdot17$,
  $C(E^{(1)},E^{(2^{7})})\mid127$) and, \textbf{for $N$ a power of $p$}, \code{non\_reduced\_certified} (the counting criterion: e.g.\ $(8,10)$, $\chi_\Q(U_2)=x^{16}(x-1)(x-2048)(x^{2}+24x+2048)^{2}$, $d=13$,
  $d^{\mathrm{an}}=6$; a reduced full $\Talg_\Q$ would need $\ge7$ distinct systems with $U_2\ne0$, but $H^{1}\otimes\Cx\supseteq M_{12}$ carries only $6$ non-zero $U_2$-roots with
  multiplicity; the attribution of $U_p$-roots to systems uses Proposition~\ref{prop:heckeZrat}(2) for $U_p$; for other $N$ the criterion is \textbf{not proved}
  (v1.25.0-pre.2, Z0 referee F2; HZ-UTAG audit F1: with a second prime $q\mid N$ several systems over one anemic class may share $U_p=0$, e.g.\ $N=12$, $k=4$:
  $U_2\in\{1,8,0\}$, $U_3\in\{1,27\}$, so ``reduced $\Rightarrow b_1-m_0\ge d-d^{\mathrm{an}}$'' is unproved, with margin $0$ at $(12,2)$); the code implements exactly this scope
  (at $(24,2)$ the raw inequality holds and the tag is withheld). $d^{\mathrm{an}}=d$ alone does not make the full algebra
  certified: the index $[\Talg^{\mathrm{fr}}:\Talg^{\mathrm{an,fr}}]$ needs the integral $U_p$, unless each $U_p$ is \emph{known} to lie in $\Talg^{\mathrm{an,fr}}$ (at $(11,0)$
  $U_{11}=\mathrm{id}$ on $L$, a literature input).
\end{itemize}
\emph{Lean form} (not a change of the theorem). The Lean certificate uses Proposition~\ref{prop:heckeZetale}(2) only in its $v_\ell=0$ clause (a B\'ezout witness
$uf+vf'=D$); a prime with $v_\ell(D_j)=1$ for every $j$ is put into $\mathsf S$ (record reason \code{v1\_pending\_LC3b}) and closed by (C3) until the $v_\ell=1$ clause is
formalised. \emph{The measured (C2) frontier} (v1.25.0-pre.0c, A7b-m1: re-measured on what (C2) forms, \textbf{one} $D_{j_0}$, with $\vartheta_{j_0}=T_2$; level $1$, the
discriminant of $\chi(T_2\mid M_k)$, trial division to $10^5$ and Pollard rho at $8$\,s per number, \file{docs/implementation/slot24/heckez/a7b\_frontier\_single.py}): the
factorisation is complete for every $k\le64$ except $k=52$, $60$; at $k=36$, $42$, $44$, $48$, $50$, $54$, $56$, $58$, $62$, $64$ it contains a prime of $65$--$204$ bits, so under
Q1(a) these records read \code{partial(unfactored)} while the Lean lane can close them; $k=52$, $60$ keep cusp cofactors of $109$ and $263$ bits, and $k=84$, $98$ the composite
cofactors of $\mathrm{num}(B_k/k)$ of $189$ and $206$ bits. This is the list of the earlier measurement on $\gcd_{p\le7}$, weight by weight: the index factor
$[\Talg^{\mathrm{fr}}:\Z[T_2]]^{2}$ of $D_{j_0}$ splits into primes below $2^{64}$ except for the square of a $65$-, $72$- resp.\ $67$-bit prime at $k=60$, $62$, $64$, weights already
in the list. For $\vartheta_{j_0}=T_3$, $T_5$, $T_7$ the index factor is larger and often unsplit at this budget (\code{partial(unfactored)} for $T_5$, $T_7$ at most weights $k\ge36$, for $T_3$ at
$k\ge52$; a new prime above $2^{64}$ at $k=32$, $38$, $40$ or $46$), so $T_2$ is the measured choice of $\vartheta_{j_0}$. More operators barely help: $\operatorname{disc}\Talg^{\mathrm{fr}}$ divides
every $\operatorname{disc}_\vartheta$ and contains the discriminant of the cuspidal Hecke field. The congruence part (Propositions~\ref{prop:heckeZcong}, \ref{prop:heckeZeiscert})
certifies beyond the frontier. \emph{As built} (v1.25.0-pre.1, 2026-11-06): the C++ factor finder has no multi-limb rho or ECM yet, so the frontier the build reaches
is lower than this measured one; see record item~8~(j) of \S\ref{sec:recordschema}. Every status it writes is one this subsection proves; the difference is coverage.

\begin{remark}[The $691$ caveat, revised]\label{rmk:heckeZcaveat}
\PROVED{} for (5) and (3) (the plan's Rem.~4.5, Props.~4.3, 4.4). \textbf{Scope}: Propositions~\ref{prop:heckeZgen}--\ref{prop:heckeZthrough},
Theorem~\ref{thm:heckeZcert}(2), (3) and Definition~\ref{def:heckeZcong} onwards assume $\Talg_\Q$ reduced --- always for the anemic algebra, for the full one when (C1) passes;
the full algebra at $p^{3}\mid N$ is refused (Proposition~\ref{prop:heckeZrat}(4)(b); $\Talg_\Q$ may then be non-monogenic, e.g.\ $\Q[x,y]/(x^{2},y^{2})$ when two primes have cube level).
\begin{enumerate}
\item \emph{One operator overstates.} $T_2$ at $(1,10)$ gives $2073=3\cdot691$; the $3$ is an index prime of $\Z[T_2]$, not a congruence. A congruence prime divides
  $\operatorname{Res}(f_{\vartheta,\mathcal S},f_{\vartheta,\mathcal S'})$ for \emph{every} $\vartheta$, and the depth is $v_\ell\operatorname{Res}$ for a $\vartheta$ that generates
  modulo $\ell$ (Proposition~\ref{prop:heckeZmono}(3)).
\item \emph{Bad primes are not small.} Level-one Eisenstein congruence primes divide $\mathrm{num}(B_k/k)$: $691$ ($k=12$), $3617$ ($16$), $43867$ ($18$), $283\cdot617$ ($20$),
  $131\cdot593$ ($22$), $103\cdot2294797$ ($24$), $657931$ ($26$), $9349\cdot362903$ ($28$), $1721\cdot1001259881$ ($30$). They come from rational data, not from a mod-$\ell$ census range.
\item \emph{Existence is not depth; when $K>1$ is needed.} $K=1$ suffices for the maximal ideals and residue fields (Proposition~\ref{prop:heckeZmax}), the saturation test
  (Proposition~\ref{prop:heckeZgen}), the \'etale verdict (Proposition~\ref{prop:heckeZetale}: no computation), and the whole local structure at a saturated block where a
  computed $\vartheta$ generates $\Talg_\midl/\ell$ (Proposition~\ref{prop:heckeZmono}; on an unsaturated block under a generation theorem the test runs at
  $K=c^{\mathrm{sat}}+1$, and the rational route needs no $K>1$). $K>1$ is needed when no computed $\vartheta$ generates (embedding dimension $\ge2$, every non-Gorenstein
  $\Talg_\midl$, unless the rational route is used: Proposition~\ref{prop:heckeZlocal}), for Proposition~\ref{prop:heckeZthrough} ($K>v_\ell\operatorname{Res}$), and for the
  at-root read (Proposition~\ref{prop:heckeZcong}(5), $K>v_\ell c(\lambda)$). $K=1$ alone cannot separate the kinds: $\Z_\ell\times_{\Fl}\Z_\ell$ (meeting, depth $1$),
  $\Z_\ell\times_{\Z/\ell^{2}}\Z_\ell$ (depth $2$), $\Z_\ell[\sqrt{\ell u}]$ (ramified DVR) and $\Z_\ell[\ell\sqrt u]$ (for $u$ a non-square unit a non-maximal order of one component; for $u$ a square and $\ell$ odd it is
  $\Z_\ell\times_{\Z/\ell}\Z_\ell$ up to isomorphism, a meeting of two components; A7-N7) all reduce to $\Fl[\varepsilon]/(\varepsilon^{2})$.
\item \emph{The census algebra is not $\Talg/\ell$.} It is $A'_1(H^{1})$, which differs by the saturation defect (Remark~\ref{rmk:sattor}), mixes in torsion
  (Proposition~\ref{prop:cohforms}(1), (4)) and needs a generation certificate (Proposition~\ref{prop:heckeZmax}); a census ``non-Gorenstein'' row is a statement about $A'_1$,
  and about $\Talg_\midl$ only where $\Talg$ is saturated \emph{and} generation is certified --- both hold at $(1,10)$, $\ell=691$ (\code{length\_A}$=2$, \code{hilbert}$=[1,1]$);
  at $(1,22)$, $\ell=103$, the census's second maximal ideal is the non-Eisenstein conjugate.
\item \emph{Other primes say nothing at $\ell$.} There are orders $\mathcal O_1\ne\mathcal O_2$ in the same $\Talg_\Q$, with faithful lattices, whose mod-$\ell'^{K}$ data
  ($\ell'\ne\ell$, all $K$) and rational data all coincide: $\{(a,b)\in\Z^{2}:a\equiv b\ (\ell)\}$ and $\{a\equiv b\ (\ell^{2})\}$ on $\Z^{2}$ (both $\Z_{\ell'}^{2}$ at every $\ell'\ne\ell$,
  characteristic polynomial $(x-a)(x-b)$ for every element of both). \textbf{Adding primes $\ell'$ helps only to find the bad set and to assemble rational data, never to
  certify structure at a bad $\ell$}; there one adds operators or raises $K$.
\item \emph{Anemic versus full.} The anemic algebra merges systems that differ only at $U_p$: at $(2,6)$ the two Eisenstein series ($U_2=1$, $U_2=2^{7}$) are anemically one;
  the Eisenstein--Eisenstein congruence at $127$ is invisible anemically, and the anemic $2$-part $2^{7}$ of the cuspidal--Eisenstein congruence overstates the full one
  ($\le2^{3}$, from $U_2$: $128-(-8)=2^{3}\cdot17$).
\item \emph{Factoring.} Only the square part of $\operatorname{disc}\Talg'$ must be factored (Theorem~\ref{thm:heckeZcert}(C2)).
\item \emph{Saturation fails outside multiplicity one}; then mod-$\ell^{K}$ data on the cohomology lattice cannot certify generation (Proposition~\ref{prop:heckeZgen}(3)), but
  Ramanujan-type congruences are still certified by Propositions~\ref{prop:heckeZcong}(4) and~\ref{prop:heckeZeiscert}, which need no generation.
\item \emph{Torsion systems.} Coincidence of a torsion maximal ideal with a free one is not certified by agreement on finitely many operators; it is certified at level $1$
  at $2$, $3$ by $T_\ell$, and anywhere by a non-split extension (Theorem~\ref{thm:heckeZdescent}(5)).
\item \emph{$\mathrm{Tor}$ terms.} $691$ is not a torsion prime: at $(1,10)$ no group $H^{q}$, $H^{q}(\partial)$, $\Hcpt{q}$ has $691$-torsion (integral census edition~1). It
  is invisible to every Bockstein and is the lattice index $\eta_L$ of Definition~\ref{def:heckeZcong}: the congruence module is the \emph{cokernel} of the lattice sequence
  $0\to L_{\mathcal S}\oplus L_{\mathcal S'}\to L\to\eta_L\to0$, not a term of a coefficient sequence; reduced modulo $\ell^{K}$ it appears as
  $\mathrm{Tor}_1(\eta_L,\Z/\ell^{K})=\eta_L[\ell^{K}]$, the $\ell$-part of $\eta_L$ once $K\ge v_\ell(\exp\eta_L)$ (A7-N7). In degree $\ge1$, $\mathrm{Tor}$ terms of $H^{q}(\Ga;\Vcoef)$ occur only at $2$
  and $3$, but $\Hcpt{2}$ has $\ell\ge5$ torsion ($78$ of the $210$ pair cells of integral census edition~1, always $\ell\le n$ or $\ell\mid N$) and no route computes a Hecke action
  on it.
\end{enumerate}
\end{remark}

\begin{proof}
(5): every element $(a,b)$ of either order acts on $\Z^{2}$ diagonally with characteristic polynomial $(x-a)(x-b)$, and both orders become $\Z_{\ell'}^{2}$ after
$\otimes\Z_{\ell'}$. (3): the four rings are free of rank $2$ over $\Z_\ell$ with one maximal ideal and a square-zero cotangent line modulo $\ell$; their kinds differ by
Proposition~\ref{prop:heckeZetale}(4), (5) and the index in the maximal order. The other items cite the statements named in them.
\end{proof}

\begin{definition}[Congruence modules]\label{def:heckeZcong}
For disjoint sets $\mathcal S,\mathcal S'$ of components, $I_{\mathcal S}:=\ker(\Talg^{\mathrm{fr}}\to\Talg_{\mathcal S})$, $\Talg_{\mathcal S}$ the image in $\prod_{i\in\mathcal S}K_i$
($K_i$ the field of component $i$). The \textbf{algebra congruence module} is $C_{\Talg}(\mathcal S,\mathcal S'):=\Talg^{\mathrm{fr}}/(I_{\mathcal S}+I_{\mathcal S'})$, the finite-level
integral form of [U,~\texttt{def:u-congmod}]. When $\mathcal S\sqcup\mathcal S'$ is all components and $L$ is a faithful lattice, the \textbf{lattice congruence module} is
$\eta_L(\mathcal S,\mathcal S'):=L/(L_{\mathcal S}\oplus L_{\mathcal S'})$, $L_{\mathcal S}:=L\cap(L\otimes\Q)_{\mathcal S}$; equivalently $\eta_L\cong\pi_{\mathcal S}(L)/L_{\mathcal S}$,
$\pi_{\mathcal S}$ the rational projection. For $p\in\mathcal P$, $c_{\mathcal S',p}$ is the \textbf{minimal} polynomial of $X_p$ on $(L\otimes\Q)_{\mathcal S'}$ (in $\Z[x]$, a
monic factor of $\chi$), and when $\Talg_{\mathcal S}=\Z$ with character $\lambda$, $g_{\mathcal P}:=\gcd_{p\in\mathcal P}c_{\mathcal S',p}(\lambda(X_p))$.
\end{definition}

\noindent
The characteristic polynomial on $(L\otimes\Q)_{\mathcal S'}$ counts each cuspidal system twice ($(x+24)^{2}$ at $(1,10)$, giving $2073^{2}$): a valid but not sharp bound, and
the certificates below close only with the minimal polynomial. For non-rational Eisenstein components ($\Talg_{\mathcal S}=\mathcal O$, characters at level $N$) the bounds
become ideals of $\mathcal O$ and the analogue of the at-root read has not been written; for two cuspidal sets neither side is cyclic, and $C_{\Talg}$ is computed from a
certified $\Talg^{\mathrm{fr}}$ (Theorem~\ref{thm:heckeZcert}) or bounded by Proposition~\ref{prop:heckeZcong}(2) --- \textbf{nothing more is asserted} (the plan's Prop.~5.7, a sketch;
prototype \Espec{cl:258alg}).

\begin{proposition}[Bounds, the Eisenstein certificate, the read at a root]\label{prop:heckeZcong}
\PROVED{} \STD{\cite{deligne_weil1}} for (5)'s existence clause (the plan's Props.~5.2--5.5, with A1-m2, m3, m4, m6 and P2a-5). Let $\mathcal S\sqcup\mathcal S'$ be all
components.
\begin{enumerate}
\item $I_{\mathcal S}L\subseteq L_{\mathcal S'}$ and $I_{\mathcal S'}L\subseteq L_{\mathcal S}$; so $\eta_L$ is a $C_{\Talg}$-module and $\exp\eta_L\mid\exp C_{\Talg}$. This holds for
  every $\Talg'\subseteq\Talg^{\mathrm{fr}}$ as well. \textbf{The lower bound uses no generation.}
\item Every element of $(I_{\mathcal S}+I_{\mathcal S'})\cap\Z$ kills $C_{\Talg}$, in particular $\operatorname{Res}(f_{\vartheta,\mathcal S},f_{\vartheta,\mathcal S'})$ for every
  $\vartheta\in\Talg^{\mathrm{fr}}$ with $\Q[\vartheta]=\Talg_\Q$; upper bounds from computed $\vartheta\in\Talg'$ are valid for $\Talg^{\mathrm{fr}}$.
\item $|C_{\Talg}|=[\Talg_{\mathcal S}\times\Talg_{\mathcal S'}:\Talg^{\mathrm{fr}}]$ and $\Talg^{\mathrm{fr}}=\Talg_{\mathcal S}\times_{C_{\Talg}}\Talg_{\mathcal S'}$.
\item \emph{(The Eisenstein certificate.)} Let $\Talg_{\mathcal S}=\Z$ (one rational Eisenstein component, character $\lambda$). Then $C_{\Talg}\cong\Z/g$ is cyclic,
  $g\Z=\{a:(a,0)\in\Talg^{\mathrm{fr}}\}$, and
  \[
    \exp\eta_L\ \big|\ g\ \big|\ g_{\mathcal P}.
  \]
  \textbf{If the outer terms agree, $g$ equals them and $\Talg^{\mathrm{fr}}=\Z\times_{\Z/g}\Talg_{\mathcal S'}$}: $\lambda(t)\equiv\phi(t|_{\mathcal S'})\pmod g$ for every
  $t\in\Talg^{\mathrm{fr}}$, $\phi\colon\Talg_{\mathcal S'}\to\Z/g$ the induced ring map --- every Hecke operator satisfies the congruence, not only the computed ones.
  \emph{$\ell$-wise form}: since $\exp\eta_L\mid g\mid g_{\mathcal P}$, ``$\exp\eta_L=g_{\mathcal P}$'' is equivalent to $v_\ell(\exp\eta_L)\ge v_\ell(g_{\mathcal P})$ for
  every prime $\ell\mid g_{\mathcal P}$; each is one read (5).
\item \emph{(Reading $\eta_L$ at a root.)} In the setting of (4) let $c:=c_{\mathcal S',p}$ with $c(\lambda(X_p))\ne0$. Such a $p$ exists when $\mathcal S'$ is cuspidal and
  $k\ge3$ ($|a_p|\le2p^{(k-1)/2}<1+p^{k-1}$; at $k=2$ choose $p$ with $a_p\ne1+p$ for every cusp form). Let $r:=\rank L_{\mathcal S}$ ($r=1$ for the full algebra; $r\ge1$ for the
  merged anemic Eisenstein system at level $N>1$, e.g.\ $r=2$ at $(2,6)$). Then $c(\lambda)L_{\mathcal S}\subseteq c(X_p)L\subseteq L_{\mathcal S}$, and with $a_1\mid\dots\mid a_r$
  the elementary divisors of $c(X_p)L$ in $L_{\mathcal S}$,
  \[
    \eta_L\cong\textstyle\bigoplus_{i\le r}\Z/\bigl(c(\lambda(X_p))/a_i\bigr),
  \]
  independently of $p$; for $r=1$, $|\eta_L|=c(\lambda(X_p))/a$. Locally: one Smith form of $c(X_p)\bmod\ell^{K}$ on $L$, $K>v_\ell c(\lambda(X_p))$. At level $N>1$ for the
  \emph{full} algebra with several Eisenstein systems the old Eisenstein partner of $\lambda$ lies in $\mathcal S'$ and agrees with $\lambda$ at every $T_p$, $p\nmid N$, so only a
  $U_p$ separates ($(2,6)$: $127=2^{7}-1$ from $U_2$), so the read needs $U_p$ over $\Z$, and no route of this version performs it (item~8~(d$'$)(ii), \emph{Why \code{lower} is \code{not\_run}}; v1.26.0-pre.11, confirm-referee C1); for the anemic algebra with $e^{\mathrm{an}}_k(N)=1$, $\mathcal S'$ is cuspidal and
  the read works.
\item \emph{(Free lattices.)} If $L_\midl\cong\Talg_\midl^{\,a}$, then $\eta_{L,\midl}\cong C_{\Talg,\midl}^{\,a}$. At odd $\ell$, an $\inv$-part $L^{\pm}_\midl$ through which every
  component of $\midl$ passes and which is \emph{cyclic} over $\Talg_\midl$ is free of rank one and gives $\eta_{L^{\pm},\midl}=C_{\Talg,\midl}$. (At level $1$ the Eisenstein class lies
  in $L^{+}$ only, so $L^{-}_\midl$ is not faithful at the Eisenstein ideal: \Espec{cl:ramanujan}, $H^{1}_{+}\cong\Talg_\midl$, $H^{1}_{-}\cong\F_{691}$.)
\end{enumerate}
\end{proposition}

\begin{proof}
(1) $t\in I_{\mathcal S}$ kills $(L\otimes\Q)_{\mathcal S}$, so $tL\subseteq L\cap(L\otimes\Q)_{\mathcal S'}$; an integer $m\in I_{\mathcal S}+I_{\mathcal S'}$ therefore kills $\eta_L$, and
$C_{\Talg}$ acts on $\eta_L=L/(L_{\mathcal S}\oplus L_{\mathcal S'})$ through $(I_{\mathcal S}+I_{\mathcal S'})L\subseteq L_{\mathcal S}\oplus L_{\mathcal S'}$. (2) B\'ezout for resultants,
$\operatorname{Res}=af_{\mathcal S}+bf_{\mathcal S'}$, evaluated at $\vartheta$: $f_{\mathcal S}(\vartheta)\in I_{\mathcal S}$, $f_{\mathcal S'}(\vartheta)\in I_{\mathcal S'}$. (3)
[U,~\texttt{lem:u-fibre}]. (4) [U,~\texttt{lem:u-fibre}] with $A=\Z$ gives $C\cong\Z/\pi_A(I_{\mathcal S'})$ and $\pi_A(I_{\mathcal S'})=\{a:(a,0)\in\Talg^{\mathrm{fr}}\}$.
Cayley--Hamilton on the $\mathcal S'$-part gives $c_{\mathcal S',p}(X_p)\in I_{\mathcal S'}$, with $\pi_A$-image $c_{\mathcal S',p}(\lambda(X_p))$, so $g\mid g_{\mathcal P}$; and $C$ cyclic of
order $g$ has exponent $g$, so (1) gives $\exp\eta_L\mid g$. Equal outer terms give $g=g_{\mathcal P}$, and the fibre product is (3). The $\ell$-wise form is the divisibility chain.
(5) $c(X_p)$ is $0$ on $(L\otimes\Q)_{\mathcal S'}$ and $c(\lambda)$ on $(L\otimes\Q)_{\mathcal S}$, so $c(X_p)L\subseteq L\cap(L\otimes\Q)_{\mathcal S}=L_{\mathcal S}$,
$c(\lambda)L_{\mathcal S}=c(X_p)L_{\mathcal S}\subseteq c(X_p)L$, and $\pi_{\mathcal S}=c(\lambda)^{-1}c(X_p)$; hence $\eta_L\cong\pi_{\mathcal S}(L)/L_{\mathcal S}\cong c(X_p)L/c(\lambda)L_{\mathcal S}$,
read in a basis of $L_{\mathcal S}$ adapted to $c(X_p)L$ (each $a_i\mid c(\lambda)$). $L_{\mathcal S}$ is saturated in $L$, so the non-zero elementary divisors of $c(X_p)$ on $L$ are
the $a_i$, each of valuation $\le v_\ell c(\lambda)$, and $r$ is known; a Smith form modulo $\ell^{K}$ with $K>v_\ell c(\lambda)$ reads them. The existence clause is Deligne's bound.
(6) $L_{\mathcal S}=(\Talg\cap\varepsilon_{\mathcal S}\Talg_\Q)^{a}=I_{\mathcal S'}^{\,a}$ locally, likewise $L_{\mathcal S'}=I_{\mathcal S}^{\,a}$, and $I_{\mathcal S}\cap I_{\mathcal S'}=0$, so
$\eta=(\Talg/(I_{\mathcal S}+I_{\mathcal S'}))^{a}$; a cyclic faithful module is free of rank one.
\end{proof}

\noindent
\emph{At $(1,10)$}: $c_{\Delta,2}(2049)=2073=3\cdot691$, $c_{\Delta,3}(177148)=176896=2^{8}\cdot691$, $g_{\mathcal P}=691$; one read at $\ell=691$, $K=2$, with
$v_{691}(\exp\eta_L)\ge1$ certifies $\Talg^{\mathrm{fr}}=\Z\times_{\Z/691}\Z$ and $\tau(m)\equiv\sigma_{11}(m)\pmod{691}$ for \textbf{every} $m$ (\Espec{cl:ramanujan} cites this
for all $m$; here it is derived from one lattice computation). No $2$- or $3$-adic read is needed, and at level $1$ none ever is once $2,3\in\mathcal P$: every level-one cusp eigenform is non-ordinary at $2$ and $3$
(Proposition~\ref{prop:heckeZlevelone23}(5)), so $c_{\mathcal S',2}(1+2^{k-1})$ and $c_{\mathcal S',3}(1+3^{k-1})$ are norms of elements $\equiv1$ modulo the primes above $2$, resp.\ $3$, and
$g_{\mathcal P}$ is prime to $6$ (A7-N4) --- which is why Proposition~\ref{prop:heckeZeiscert}(1) reaches $\ell\in\{2,3\}$. $|\eta_L|$ from $T_2$ ($a=3$) is a cross-check, not part of the
certificate. The pair complex (Proposition~\ref{prop:descentpair}; the denominator of the Eisenstein class) gives an independent reading, kept as a cross-check.
\emph{Factored $g_{\mathcal P}$} (P2a-5). The $\ell$-wise form needs the prime factorisation of $g_{\mathcal P}$, obtained without a big gcd: the word-size primes of one
$c_{\mathcal S',p_0}(\lambda(X_{p_0}))$, each $v_\ell(g_{\mathcal P})$ read by residues modulo $\ell^{K}$ (Q1(a), A7-m4); the certificate closes \textbf{at every certified
prime $\ell\mid g_{\mathcal P}$}, and when $g_{\mathcal P}$ keeps a cofactor (unfactored, or a factor only \code{probable}) the global fibre-product statement is
\code{partial(unfactored)}, as is Proposition~\ref{prop:heckeZeiscert}. \emph{The word modulus} (the plan's FX-1, R-1). A read modulo $\ell^{K}$ needs $\ell^{K}$ below the
engine's largest word modulus ($<2^{63}$); (C3) has $K=1$ (limit $\ell<2^{63}$), the at-root read $K>v_\ell c(\lambda)\ge1$, i.e.\ $K\ge2$ (limit $\ell<2^{31.5}$; $K=2$ at $691$
and at every Eisenstein prime of a level-one cell with $k\le46$). A certified $\ell\mid g_{\mathcal P}$ whose at-root read needs $\ell^{K}\ge2^{63}$ cannot be read on any route of the
plan; its $\ell$-part stays \code{upper\_bound\_certified} (read reason \code{refused(word\_modulus)}), final. At level $1$, $k\le64$: no read at all at
$k\in\{36,42,48,56,60,62,64\}$ ($\mathrm{num}(B_k/k)$ has a prime factor of $65$--$107$ bits); (C3) yes and the at-root read refused at $k\in\{34,38,40,46,50,52,54,58\}$; the
at-root read at $K=2$ on the $64$-bit regimes for $2294797$ ($k=24$) and the Eisenstein primes of $k=26,28,30,32,44$.
\emph{As built} (v1.25.0-pre.1, 2026-11-06; \file{docs/implementation/slot25/HZLOC\_report.md}). At level $1$ the at-root reads run on the $64$-bit regimes of the console's
module route: the seven named primes ($2294797$; $657931$, $362903$, $1001259881$, $305065927$; $2947939$, $1798482437$) are read at $K=2$ and every one closes. That route needs
$K\ge2$, so a (C3) read at $K=1$ for $\ell\in[2^{31.5},2^{63})$ at level $1$ is not served today; the code writes \code{refused(word\_modulus)} although its modulus is below $2^{63}$,
and the record carries the \textbf{pending} \code{refused(k1\_route)} (v1.25.0-pre.2, lead decision on O1: no new $K=1$ route now, an EXPORT-Z / slot-26 item; none of these
weights is in the first golden set). At level $N>1$ the lattice comes from EXPORT, whose word ring has $\ell^{K}<2^{31}$: for $\ell\in[46341,2^{31})$ it serves $K=1$ only,
so (C3) runs and the at-root read is \code{refused(export\_word)}; for $\ell\ge2^{31}$ it serves no $K$ and (C3) is \code{refused(export\_word)} too; \code{word\_modulus}
($\ell^{K}\ge2^{63}$) is tested first. \code{export\_word} is pending: EXPORT-Z step~2 moves it (record item~8~(c)). The read set
at $\ell$ contains $T_\ell$ for $\ell\nmid N$, $\ell\le23170$ (item~8~(j)).

\begin{proposition}[The single-Eisenstein-component certificate; two further routes]\label{prop:heckeZeiscert}
(The plan's Prop.~5.6 and Rem.~5.6$'$.)
\begin{enumerate}
\item \PROVED{} given Proposition~\ref{prop:heckeZrat}(5)(a). Suppose the Eisenstein part is one rational component (level $1$, $k\ge4$; prime level $N$, $k=2$), let
  $\mathcal P\supseteq\{p\le B_N\}$ (with $U_N$ if $N>1$), $\mathcal S'$ the cuspidal components and $g_{\mathcal P}$ as in Definition~\ref{def:heckeZcong}. If
  $v_\ell(\exp\eta_L)\ge v_\ell(g_{\mathcal P})$ for every prime $\ell\mid g_{\mathcal P}$, then
  \[
    \Talg^{\mathrm{fr}}=\Z\times_{\Z/g_{\mathcal P}}\Talg^{\mathrm{cusp}},\qquad\Talg^{\mathrm{cusp}}=\textstyle\sum_{n\le B_N}\Z\,T_n|_{\mathrm{cusp}},
  \]
  glued by $T_n\mapsto\lambda(T_n)\bmod g_{\mathcal P}$. This certifies $\Talg^{\mathrm{fr}}$ completely, at every prime, $2$ and $3$ included, \textbf{without saturation}.
  \textbf{Available today at level $1$ only}: at prime level the $\Z$-span needs the $T_n$ with $N\mid n$, i.e.\ $U_N$ over $\Z$; for the anemic algebra Sturm gives no spanning
  set. If $\exp\eta_L<g_{\mathcal P}$ the certificate does not close ($L$ not free at the Eisenstein ideal); $C_{\Talg}$ is then sandwiched, and
  Proposition~\ref{prop:heckeZcong}(6) on the Eisenstein $\inv$-part, or Propositions~\ref{prop:heckeZgen}, \ref{prop:heckeZmono} at that $\ell$, may close it.
\item \emph{(Bernoulli route, level $1$.)} \PROVED{} given [U,~\texttt{pr:u-U1}] (proved in [U] from Emerton's theorem) and [U,~\texttt{rmk:u-hidaM}]. At level $1$, $\ell\ge5$,
  $(\ell-1)\nmid k$, the $\ell$-part of $C_{\Talg}(\mathrm{Eis},\mathrm{cusp})$ is $\Z/\ell^{v_\ell(B_k/k)}$: the lower bound is [U,~\texttt{thm:u-sandwich}], the equality
  [U,~\texttt{pr:u-U1}], and $\Talg^{\mathrm{fr}}\otimes\Z_\ell$ localised at the Eisenstein ideal is [U]'s $\Talg_k(a)_\midl$ ([U,~\texttt{rmk:u-hidaM}] with
  Proposition~\ref{prop:heckeZrat}). So where $v_\ell(g_{\mathcal P})=v_\ell(B_k/k)$ the $\ell$-part of (1) closes from $B_k/k\bmod\ell^{K}$ alone. \textbf{The at-root read stays
  the primary certificate}; this is an independent second one, a disagreement is a defect, and where the at-root read is \code{refused(word\_modulus)} this route does
  \textbf{not} lift the $\ell$-part above \code{upper\_bound\_certified} (a future option, owner-gated, with a Lean form still to be planned). At $(1,10)$:
  $v_{691}(B_{12}/12)=1=v_{691}(g_{\mathcal P})$.
\item \emph{(Mazur route, prime level, weight $2$.)} \PROVED{} given \STD{\cite[II (9.7), Cor.~(16.2), (16.6)]{mazur}; the page images are quoted in the
  verified [GT] row \texttt{lit:hz-mazur-eis}} (v1.25.0-pre.0b: the valuation check, A7-M1). Let $N$ be prime, $n=0$, $\ell\ge5$, $\ell\mid\nu:=\mathrm{num}((N-1)/12)$,
  $\midl$ the Eisenstein maximal ideal of $\Talg^{\mathrm{cusp}}$ over $\ell$, $\Talg^{\mathrm{cusp}}_\midl$ Mazur's cuspidal local ring and $e^{\mathrm{cusp}}_\midl$ its rank over
  $\Z_\ell$ (a block rank, \textbf{not} known a priori). For a prime $q\ne N$ let $\mathsf h_q\in\Z[x]$ be the characteristic polynomial of $T_q$ on $S_2(\Ga_0(N))$ --- the monic
  square root of the cuspidal factor $\chi^{(N)}_{q,1}/(x-1-q)$ (each cuspidal system occurs twice in $L$, Proposition~\ref{prop:heckeZrat}(2)) --- and $\mathsf F_q\in\Z_\ell[x]$ its
  Hensel factor at $1+q$: the monic factor whose roots are exactly the roots of $\mathsf h_q$ congruent to $1+q$ modulo the maximal ideal of $\overline\Z_\ell$, with multiplicity ($\mathsf h_q$ need not be squarefree; A7b-N1). Let $f_{q,\midl}$ be the
  characteristic polynomial of $T_q$ on $\Talg^{\mathrm{cusp}}_\midl$. \textbf{The check}:
  \[
    v_\ell\bigl(\mathsf F_q(1+q)\bigr)=v_\ell(\nu).
  \]
  If it holds, then $\mathsf F_q=f_{q,\midl}$, $e^{\mathrm{cusp}}_\midl=\deg\mathsf F_q$, $\Talg^{\mathrm{cusp}}_\midl=\Z_\ell[\eta_q]\cong\Z_\ell[x]/(\mathsf F_q)$, $\eta_q:=T_q-1-q$, the
  Eisenstein ideal of $\Talg^{\mathrm{cusp}}_\midl$ is $(\eta_q)$, and the $\midl$-block $\Talg^{\mathrm{fr}}_\midl$ of $\Talg^{\mathrm{fr}}\otimes\Z_\ell$ (the Eisenstein component
  included) has rank $e^{\mathrm{cusp}}_\midl+1$:
  \[
    \Talg^{\mathrm{fr}}_\midl=\Z_\ell[T_q]_\midl\cong\Z_\ell[x]/\bigl((x-1-q)\mathsf F_q\bigr)\cong\Z_\ell\times_{\Z/\ell^{v_\ell(\nu)}}\Z_\ell[x]/(\mathsf F_q),
  \]
  $\Z_\ell[T_q]_\midl$ the image of $\Z_\ell[T_q]$ in that block; so $U_N\in\Talg^{\mathrm{an,fr}}_\midl=\Talg^{\mathrm{fr}}_\midl$. If it fails, nothing is certified by this $q$; a computed $v_\ell(\mathsf F_q(1+q))<v_\ell(\nu)$ is impossible by the proof and is reported as a defect (A7b-N4). For
  a \textbf{good} $q$ in Mazur's sense --- (i) $q$ is not an $\ell$-th power modulo $N$, i.e.\ $q^{(N-1)/\ell}\not\equiv1\pmod N$, and (ii) $\frac{q-1}{2}\not\equiv0\pmod\ell$, i.e.\
  $q\not\equiv1\pmod\ell$, one modular power and one residue (A7b-N4) --- the check fails \textbf{exactly} when a non-Eisenstein maximal ideal $\midl'\ni\ell$ of $\Talg^{\mathrm{cusp}}$ has a $T_q$-eigenvalue
  $\equiv1+q$, and this happens for good $q$: at $(N,\ell,q)=(61,5,37)$, $q$ is good, $\mathsf h_{37}=(x-8)(x^{3}+6x^{2}-36x-108)$, the rational newform of level $61$ ($a_{37}=8\equiv38$,
  $a_2=-1\not\equiv3\pmod5$) enters $\mathsf F_{37}$, $\deg\mathsf F_{37}=2>1=e^{\mathrm{cusp}}_\midl$ and $v_5(\mathsf F_{37}(38))=2>1=v_5(\nu)$; at the same cell $q=2$ passes. So
  goodness is how $q$ is chosen and the check is the certificate: the record carries $q$, checks that it is good, and certifies only when the check holds, else tries the next
  good $q$ (or splits $\mathsf h_q$ by the block attribution of Proposition~\ref{prop:heckeZthrough}). Everything is residue arithmetic modulo $\ell^{K}$, $K>v_\ell(\nu)$:
  $\chi^{(N)}_{q,1}$ by residues, the exact division by $x-1-q$, the square root by the triangular recursion on the coefficients ($\ell$ odd), the Hensel factor by lifting
  the coprime factorisation modulo $\ell$ --- no squarefree part and no gcd of big integers (owner ruling Q1(a)). For $\ell\,\|\,\nu$, $\mathsf F_q(x+1+q)$ is Eisenstein and
  $\Talg^{\mathrm{cusp}}_\midl$ a totally ramified DVR; for $\ell^{2}\mid\nu$ (e.g.\ $N=101$, $\nu=25$) only $\Talg^{\mathrm{cusp}}_\midl/(\eta_q)=\Z/\ell^{v_\ell(\nu)}$ and
  $\mathsf F_q$ are certified, and the kind is read by Proposition~\ref{prop:heckeZmono}(4). This needs only integral characteristic polynomials: \textbf{no $U_N$ and no
  export}; it is local at $\midl$, and the global $\Z$-span still needs $U_N$.
\end{enumerate}
\end{proposition}

\begin{proof}
(1) Proposition~\ref{prop:heckeZcong}(4) gives $C=\Z/g_{\mathcal P}$ and the fibre product; Proposition~\ref{prop:heckeZrat}(5)(a) gives $\Talg^{\mathrm{cusp}}$ (the $T_n$,
$n\le B_N$, are words in the $X_p$, $p\le B_N$); the gluing map is the ring map $\Talg^{\mathrm{cusp}}\to C$ of [U,~\texttt{lem:u-fibre}], determined on the $\Z$-spanning set.
(2) is the cited statements. (3) Write $\Talg:=\Talg^{\mathrm{cusp}}_\midl$ (Mazur's $\mathbf T_{\mathfrak P}$, a completion and a direct factor of $\Talg^{\mathrm{cusp}}\otimes\Z_\ell$,
containing $w=-U_N$; [GT] \texttt{lit:hz-mazur-eis}), $e:=e^{\mathrm{cusp}}_\midl$ and $I_\midl$ its Eisenstein ideal. Mazur (9.7) gives $\Talg/I_\midl\cong\Z_\ell/\nu=\Z/\ell^{v_\ell(\nu)}$, and
$\eta_q\in I_\midl$. $\Talg$ is free of rank $e$ over $\Z_\ell$, and $\mathsf h_q$ is the product over the maximal ideals of $\Talg^{\mathrm{cusp}}$ over $\ell$ of the characteristic polynomials of
$T_q$ on their blocks of $S_2(\Ga_0(N);\Q_\ell)$, each free of rank one over its block of $\Talg^{\mathrm{cusp}}\otimes\Q_\ell$ (the perfect pairing,
Proposition~\ref{prop:heckeZrat}(5)(a)); so $f_{q,\midl}\mid\mathsf h_q$, and as $\Talg$ is local with $\eta_q\in\midl$, every root of $f_{q,\midl}$ is $\equiv1+q$: $\mathsf F_q=f_{q,\midl}G$,
$G$ the Hensel factor at $1+q$ of $\mathsf h_q/f_{q,\midl}$. $f_{q,\midl}(1+q)=\pm\mathrm{N}(\eta_q)\ne0$ ($|a_q|\le2\sqrt q<1+q$), so
$v_\ell(f_{q,\midl}(1+q))=\mathrm{length}(\Talg/\eta_q\Talg)\ge\mathrm{length}(\Talg/I_\midl)=v_\ell(\nu)$, with equality iff $\eta_q\Talg=I_\midl$; and every root $\alpha$ of $G$ has
$v(1+q-\alpha)>0$, so $v_\ell(G(1+q))>0$ unless $G=1$. Hence the check holds iff $G=1$ and $\eta_q\Talg=I_\midl$. Then $\Talg=\Z_\ell+I_\midl=\Z_\ell+\eta_q\Talg$, hence
$\Talg=\Z_\ell[\eta_q]+\eta_q^{j}\Talg$ for every $j$, and $\Talg=\Z_\ell[\eta_q]$ ($\midl$-adic completeness; $\Z_\ell[\eta_q]$ is closed); its rank is $e$, so the minimal
polynomial of $T_q$ on $\Talg$ is $f_{q,\midl}=\mathsf F_q$ and $\Talg\cong\Z_\ell[x]/(\mathsf F_q)$. For a good $q$, (16.1), (16.2) and (16.6) give $\eta_q\Talg=I_\midl$, so the check
reduces to $G=1$; at $(61,5,37)$ the exact computation of the A7 fix round (\file{docs/implementation/slot24/heckez/Z0-report.md} \S5) gives $G=x-8$. The image of $\Z_\ell[T_q]$ in the
$\midl$-block is $\Z_\ell[x]/((x-1-q)\mathsf F_q)$ (the minimal polynomial there; $\mathsf F_q(1+q)\ne0$), and $\Z_\ell[T_q]_\midl\subseteq\Talg^{\mathrm{fr}}_\midl\subseteq\Z_\ell\times\Talg$,
both inclusions of index $\ell^{v_\ell(\nu)}$ in the outer ring --- the first by Proposition~\ref{prop:heckeZmono}(3) for $\Z_\ell[x]/((x-1-q)\mathsf F_q)$, the second by
Proposition~\ref{prop:heckeZcong}(3) with $C_\midl=\Talg/I_\midl=\Z/\ell^{v_\ell(\nu)}$ ($I_{\mathcal S}$ maps onto the Eisenstein ideal, $U_N-1=-(1+w)$) --- so they are equal, and $U_N$
lies in the anemic algebra at $\midl$. For $\ell\,\|\,\nu$: $\bar{\mathsf F}_q(x+1+q)=x^{e}$ ($\Talg$ local, $\eta_q\in\midl$) and its constant term has valuation $1$.
\end{proof}

\paragraph{Code map and reach} (v1.25.0-pre.0). \textbf{Nothing of this subsection is built yet}; the C++ programme is the plan's slot-25 packages (\file{06-work-packages.md}
\S2.3): EXPORT (the whole module $H^{1}_{(\ell)}/\ell^{K}$ with the cochain-level operators, every level, $\ell\in\{2,3\}$ included; it lifts every
\code{blocked(export)} of item~8), HZ-RAT (combinations $\vartheta=\sum c_pX_p$, $\chi_\vartheta$, discriminants, resultants, the cofactor route and the Pratt verifier below
$2^{64}$, under Q1(a)), HZ-UTAG (the two tags at $N>1$), HZ-LOC (the local reads, the at-root read with the $64$-bit regimes), HZ-TOR (separation, split test, the sandwich
for $J$), HZ-REC (item~8), HZ-EMIT (the Lean certificates); UP-Z ($U_p$ over $\Z$) in slot~26. Built today and used: Theorem~\ref{thm:descentcert}'s record (\code{integral\_descent\_record/v1},
item~7); \code{HeckeExt/HeckeRingHeckeAlgebra} and \code{HeckeSaturationDefect} ($\mathsf A'_K$ at level $1$, $\ell\ge5$); \code{HeckeExt/HeckeLocalAlgebra},
\code{HeckePresentation}, \code{HeckeGenerationCertificates} (over $\Fl$). The first golden objects of item~8 are planned for the slot-25 close (item~8~(k)).

\paragraph{Formalization note.} No \code{\textbackslash LEAN} marker is placed in this subsection. The plan's Lean lane (\file{07-lean-plan.md} v0.4) formalises
Lemma~\ref{lem:heckeZbridge} (HZ-0), Propositions~\ref{prop:heckeZspec}, \ref{prop:heckeZmax}, \ref{prop:heckeZgen}, \ref{prop:heckeZetale} ($v_\ell=0$ clause),
\ref{prop:heckeZmono}, \ref{prop:heckeZlocal}, \ref{prop:heckeZthrough}, Theorems~\ref{thm:heckeZdescent}(1)--(3), \ref{thm:heckeZcert}, Propositions~\ref{prop:heckeZcong},
\ref{prop:heckeZeiscert}(1), \ref{prop:heckeZrat}(5)(c), and of the stage-03 block Propositions~\ref{prop:heckeZuct}, \ref{prop:heckeZsplit}, \ref{prop:heckeZlevelone23},
\ref{prop:heckeZglue}(1)--(4), (5)(a), as [GT] entries, with per-cell certificates (K1--K10) re-checked by a verified checker; a marker enters only after the census line of
\S\ref{sec:status}'s rule, and only after the owner has released [C] 1.24.0 (owner ruling Q8). Remarks~\ref{rmk:heckeZcbone}--\ref{rmk:heckeZcbthree} are never hypotheses of a
Lean theorem.


%----------------------------------------------------------------------
\section{The Moerdijk--Pronk model: the simplicial set in all degrees, refinements, morphisms}\label{sec:mpmodel}
%----------------------------------------------------------------------
% NEW at v1.22.0-pre.5 (hecke-extensions slot 22, MP-CT).  Source: docs/implementation/slot22/MP_math_plan.md v0.2.1 ([MP] below;
% Opus audit MP_math_plan_audit.md, repairs R1-R8 applied there; the F-MP2 repair of v0.2.1, the label factor rho(h_{sigma_0})), the
% implementation plan MP_impl_plan.md v0.2.1 sec. 4 and its review (R1, R2, R5, R6), the owner rulings Q1-Q7 of 2026-10-09 (taken at the
% recommendations), the F-MP2 code audit F-MP2_code_audit.md sec. 1.3, and the built code and validations (MP1a_report.md, MP1b_report.md,
% MPT_report.md, MP2_report.md; docs/validation-catalogue.md, slot 22 MP-1a .. MP-2).  Placed after sec:isomodelell, before the appendices,
% so that no chapter or statement number of v1.22.0-pre.4 moves (the implementation plan's sec. 4 item 2 named the place after
% sec:orbifoldmodel, which would renumber fourteen chapters; the precedent of sec:isomodelell is followed instead).  Refereed by
% F-MP2 (contract part), ACCEPT WITH REPAIRS, repairs R1-R8 applied at v1.22.0-pre.6 (docs/implementation/slot22/F-MP2_contract_referee.md).
% v1.22.0-pre.6: the first \LEAN markers (30 on 11 labels, [GT] Part G12, MPL1_report.md v1.1 sec. 5); def:mprules(3) goodness repaired.
% v1.22.0-pre.7: 25 further \LEAN markers (MPL3_report.md v1.0 sec. 5), unnormalised clauses, MP-3 COMP^r records (MPC10, MPC11); refereed at
% pre.8 (C_pre7_referee.md, Fable, ACCEPT; minor repairs F1-F4 applied at pre.8).
% v1.22.0-pre.9: code-vs-chapter wording repairs (a)-(e) in sec:mpmodel; not refereed.
% v1.22.0-pre.8: rmk:mphecke rewritten for MP-4 (F-MP4 sec. 3) with its COMP^r record and MPC12; MP-3b coverage of MPC10, MPC11; not refereed.
% v1.22.0-pre.11: def:mpqd and prop:mpqreduced (the Q-reduced nerve, F-MPT referee sec. 3, proof checked and written out by the editor), Gap G1
% closed in rmk:mphybrid(3); not refereed.

\paragraph{What this chapter fixes.} Moerdijk and Pronk attach a simplicial set to a triangulated orbifold whose simplicial cohomology is the
orbifold cohomology in \emph{every} degree, from a \emph{finite} triangulation \cite{moerdijk_pronk}; Remark~\ref{rmk:mpcaveat} records why
their construction could not be made normative here (effective charts, $\mu$-maps whose independence of choices is stated without proof).
This chapter states the global-quotient form of their object for $[\SL_2(\Z)\backslash\Xbulk]$ \emph{without} those two obstacles, proves its
comparison theorem independently of their argument, and fixes the maps that the code builds. (i) \textbf{Orbifold triangulations as finite data}
and Moerdijk--Pronk's rules as decidable predicates (Definitions~\ref{def:mptri}, \ref{def:mprules}, \ref{def:mpell}); the triangulations in use,
$\mathsf M_7$ and the refinement $\mathsf R_{\Z}$ (Definition~\ref{def:mpRZ}, Table~\ref{tab:mprules}, Remark~\ref{rmk:mprulesread}). (ii)
\textbf{The choice-free MP complex} (Definition~\ref{def:mpset}) and what it is (Remark~\ref{rmk:mpwhat}). (iii) \textbf{The comparison theorem}
(Theorem~\ref{thm:mpcomparison}): its cohomology is $H^{q}(\Ga;M)$ for every $q$, every ring and every module, $-I$ included, through the free
bisimplicial $\Ga$-set $\mathfrak F$ (Definition~\ref{def:mpfree}) and one assembly lemma used twice (Lemmas~\ref{lem:mprows}, \ref{lem:mpcols},
\ref{lem:mpassembly}). (iv) \textbf{The $\ell$-local hybrid} (Remark~\ref{rmk:mphybrid}). (v) \textbf{All degrees}: the comparison lift
$\Phi^{\mathsf{MP}}$ (Definition~\ref{def:mpphi}) and the exact scope of the all-degrees statement (Proposition~\ref{prop:mpalldegrees}). (vi)
\textbf{Morphisms} with the label factor (Definition~\ref{def:mpmorphism}), naturality and refinement invariance (Theorem~\ref{thm:mpnatural}), the
Bredon comparison (Proposition~\ref{prop:mpbredon}), the boundary (Proposition~\ref{prop:mpboundary}), restriction and corestriction
(Proposition~\ref{prop:mprescor}), Hecke operators (Remark~\ref{rmk:mphecke}). (vii) Moerdijk--Pronk's teardrop (Example~\ref{ex:mpteardrop}) and
what is checked (Remark~\ref{rmk:mpcheckable}).

\begin{contractbox}
\textbf{No shipped value or default changes} (owner ruling Q4 of 2026-10-09). What $\Fl$ and $\Z/\ell^{e}$ \emph{name} at $\ell\in\{2,3\}$ stays
Definition~\ref{def:ringsemantics}: the Bredon object of Proposition~\ref{prop:bredon}, not group cohomology; in particular $\F_2$ and $\F_3$ keep
their Bredon semantics. The MP model is a new object with its own validation names; its output lines carry \code{model=mp}; it adds no golden
object, record, setting, console verb or env hook that changes a value. Whether $\F_2$ / $\Z/2^{e}$ should ever \emph{mean} group cohomology is a
separate, golden-gated owner question (the implementation plan's D1), not decided here.
\end{contractbox}

\noindent\textbf{Normative object.} As in \S\ref{sec:orbifoldmodel}: ``the orbifold cohomology of $[\Ga\backslash\Xbulk]$'' is the cohomology of the Borel
construction, i.e.\ group cohomology $H^{q}(\Ga;M)$ in every degree (Definition~\ref{def:borelmodel}(2)); every complex of this chapter is a
computational model of it. Moerdijk--Pronk's sheaf-side identification of orbifold sheaf cohomology with the Borel construction is cited in
Remark~\ref{rmk:mpcaveat}(1)(i) and is \textbf{not} re-proved; what is proved is the computational half, the analogue of their comparison theorem,
for the global quotient and without their effectivity assumption.
\textbf{Notation.} $\Ga=\SL_2(\Z)$ unless a statement says ``a group $\Ga$''; $Y$ a simplicial $\Ga$-complex as in \ref{G1}; $\cring$ a commutative
ring, $M$ an $\cring[\Ga]$-module ($M=\VR$ with $\rhon$ in the numbers). $\mathfrak T$, $\mathrm{Ar}$, $\mathsf S$, $\underline{M}$, $\mathfrak F$,
$\mathfrak F'$, $O(x)$, $\mathsf h$, $\Phi^{\mathsf{MP}}$, $\varepsilon^{\mathsf{MP}}$, $\mathsf M_7$, $\mathsf R_{\Z}$, $\mathsf R_{\Z}^{+}$,
$\mathsf R_{\zeta}$, $\mathsf R_{\mathrm{good}}$, $\mathrm{Tot}$, $C^{\bullet}_{\mathrm{un}}$, $D$, and (v1.22.0-pre.11) $Q_0$, $\mathcal C$, $\mathsf p(x)$, $\mathfrak D$, $\mathfrak F^{Q}$,
$\mathsf S^{(\ell)}$, and the hypothesis tags (SC), (AC), (IO), (PD), (QD) are new symbols of this chapter and mean nothing else (rule~\ref{A5}): $D$ denotes the degree from which periodicity transports every higher degree, $D=3$ (Proposition~\ref{prop:mpalldegrees}(2);
neither the matrix $D$ of Remark~\ref{rmk:stabconvention}(3) nor $D_{M,N,d}$), $\mathsf S$ is not the matrix $S$, $\mathfrak F'$ is not the column $F'^{\sigma}$ of
Definition~\ref{def:isomodelell}, $\Phi^{\mathsf{MP}}$ is not the cone map $\Phi_g$ nor $\Phi_\ell$, and $\varepsilon^{\mathsf{MP}}$ is not the
augmentation $\aug$ (the implementation's reports write $S_k$, $F'$, $\Phi$, $\varepsilon^{\ast}$ and $h_{\tilde\sigma}$). $P_\sigma$, $Q_\sigma$ are
those of Definition~\ref{def:isomodelell}. The elliptic vertices are written $v_4=\zeta^{2}$, $v_5=\zeta$, $v_6=i$ as in Definition~\ref{def:model7};
[MP] is \file{docs/implementation/slot22/MP\_math\_plan.md} v0.2.1, cited by section for the full derivations.

\subsection{Orbifold triangulations as finite data}\label{sec:mptrisec}

\begin{definition}[Orbifold triangulation; descriptor; the four hypotheses]\label{def:mptri}
Let $\Ga$ be a group and $Y$ a finite-dimensional simplicial $\Ga$-complex satisfying \ref{G1} (simplicial action without inversions, finite
stabilisers) and \ref{G3} (finitely many orbits of simplices in each degree).
\begin{enumerate}
\item \textbf{(Descriptor.)} A \emph{descriptor} of $Y$ is the finite datum: for each degree $q$ a set $\Sigma_q(Y)$ of orbit representatives
  $\iota(\sigma)$, the stabilisers $\Ga_\sigma$, the transporters $\mathrm{Tr}(\sigma,\sigma')=\{\gamma\in\Ga:\gamma\,\iota(\sigma')=\iota(\sigma)\}$ (finite unions of
  cosets of stabilisers), and the faces $\partial_j\iota(\sigma)=h\cdot\iota(\tau)$, $h\in\Ga$. The descriptors of this document are
  \code{integral\_model::ModelData} (\S\ref{sec:isomodel}, \S\ref{sec:isomodelell}'s code map) and, in Lean, [GT]'s \code{XBulkDatum}.
  $\Sigma_{\max}(Y)$ denotes the representatives of the maximal simplices.
\item \textbf{(Isotropy datum; orbifold triangulation.)} The \emph{isotropy datum} is $\sigma\mapsto\Ga_\sigma$ (the pointwise stabiliser, which is the
  setwise one by \ref{G1}) with, for a face $\tau$ of $\sigma$ realised as $h\iota(\tau)\subseteq\iota(\sigma)$, the monomorphism
  $\Ga_\sigma\to\Ga_\tau$, $\gamma\mapsto h^{-1}\gamma h$ (well defined up to conjugation in $\Ga_\tau$). An \textbf{orbifold triangulation}
  $\mathfrak T$ of $[\Ga\backslash Y]$ is a descriptor of $Y$ together with its isotropy datum. \emph{Effectivity is not assumed}: a normal subgroup
  acting trivially on $Y$ --- for $\Ga=\SL_2(\Z)$ on $\Xbulk$ the centre $\{\pm I\}$, \ref{H4}(iv) --- lies in every $\Ga_\sigma$ and is kept.
\item \textbf{(The four hypotheses.)} (SC) \emph{simplicial-complex property}: two closed simplices of $Y$ meet in the empty set or a common face;
  (AC) \emph{acyclicity}: the augmented chain complex $C_{\bullet}(Y)\to\Z$ is exact; (IO) \emph{invariant order}: no simplex of $Y$ has two vertices
  in one $\Ga$-orbit, so ordering the vertices of each simplex by the index of their orbit is a $\Ga$-invariant total order on every simplex;
  (PD) \emph{purity}: every maximal simplex of $Y$ has dimension $\dim Y$, so $\Sigma_{\max}$ is the set of top-dimensional representatives and every simplex
  lies in one. On a descriptor (SC), (IO) and (PD) are finite, decidable statements.
\end{enumerate}
\end{definition}

\begin{definition}[Moerdijk--Pronk's rules as predicates]\label{def:mprules}
Let $\mathfrak T$ be an orbifold triangulation of $[\Ga\backslash Y]$. Following the rules (i), (ii) of an adapted triangulation in
\cite{moerdijk_pronk} (quoted in Remark~\ref{rmk:mpcaveat}(1)(ii)), read on the isotropy datum:
\LEAN{def:mp-datum} ((1)--(4) as decidable predicates on a finite descriptor, [GT]'s \code{MPDatum}, not on \code{XBulkDatum}; (3) in its finite-group form and in the strong form below)
\begin{enumerate}
\item \textbf{Adapted} (their rule (ii), \emph{subgroup form}): every simplex $\sigma$ of $Y$ has a vertex $v(\sigma)$ with $\Ga_\tau\subseteq\Ga_{v(\sigma)}$ for
  every face $\tau$ of $\sigma$ (point stabilisers, as subgroups of $\Ga$ after transport --- \emph{not} as orders).
\item \textbf{Adapted for $\mathfrak T^{(\ell)}$}: the same with $P_\tau\subseteq P_{v(\sigma)}$, the $\ell$-Sylow datum of Definition~\ref{def:mpell}. This is
  rule (ii) applied to the $\ell$-relevant datum, \emph{not} Moerdijk--Pronk-conformance of $[\Ga\backslash Y]$, which concerns the actual isotropy.
\item \textbf{Good} (their chart condition): for every vertex $v$, the set $\mathsf T(v):=\{\gamma\in\Ga:\gamma\,\mathrm{St}(v)\cap\mathrm{St}(v)\ne\emptyset\}$,
  $\mathrm{St}(v)$ the closed star, generates a finite subgroup of $\Ga$. The \emph{strong} form is $\mathsf T(v)\subseteq\Ga_p$ for one vertex $p$ of $\mathrm{St}(v)$
  ($\mathsf T(v)$ in the stabiliser of one point of $\mathrm{St}(v)$; the stabiliser of any point lies, under (IO), in that of a vertex of its carrier). It is \emph{not}
  $\mathsf T(v)\subseteq\Ga_v$, which is strictly stronger (v1.22.0-pre.6; pre.5 wrote this form, which no code or table uses).
  \LEAN{thm:mp-goodvtx} (the strong form as stated here, \code{goodAtB\_iff\_onePt}; and that $\mathsf T(v)\subseteq\Ga_v$ fails at every vertex orbit of $\mathsf M_7$
  and at $v_0,v_1,v_2,m$ on $\mathsf R_{\Z}$)
\item \textbf{Collar}: no simplex meets both the boundary subcomplex and a vertex whose stabiliser is not the kernel of the action.
\end{enumerate}
\end{definition}

\begin{definition}[The $\ell$-relevant structure]\label{def:mpell}
Let $\Ga=\SL_2(\Z)$, so every $\Ga_\sigma$ is cyclic, $\Ga_\sigma=P_\sigma\times Q_\sigma$ with $P_\sigma$ its $\ell$-Sylow subgroup and $\abs{Q_\sigma}$ prime to $\ell$
(Definition~\ref{def:isomodelell}). The \textbf{$\ell$-relevant orbifold structure} of $\mathfrak T$ is the pair $(\mathfrak T^{(\ell)},M^{Q})$: $\mathfrak T^{(\ell)}$
the isotropy datum with $\Ga_\sigma$ replaced by $P_\sigma$, and $M^{Q}$ the \emph{coefficient system} $\sigma\mapsto M^{Q_\sigma}$, contravariant along faces
($\Ga_\sigma\subseteq\Ga_\tau$ for $\tau\le\sigma$ gives $Q_\sigma=Q_\tau\cap\Ga_\sigma\subseteq Q_\tau$, hence $M^{Q_\tau}\subseteq M^{Q_\sigma}$), with $P_\sigma$ acting on
$M^{Q_\sigma}$. \LEAN{def:mp-datum} (the $\ell$-Sylow datum $\sigma\mapsto P_\sigma$ on a finite descriptor, as the input of Definition~\ref{def:mprules}(2); not the
coefficient system $M^{Q}$, not the table below) On $\Xbulk$ (owner ruling Q3 of 2026-10-09):
\begin{center}\small
\begin{tabular}{@{}l>{\raggedright\arraybackslash}p{0.33\textwidth}l>{\raggedright\arraybackslash}p{0.33\textwidth}@{}}
\toprule
& $\ell$-singular simplices, $P_\sigma$ & elsewhere & $Q$-invariants kept \\
\midrule
$\ell=3$ & the orbit of $\zeta^{2}$: $P=\langle(ST)^{2}\rangle\cong\Z/3$ & $P=1$ & $M^{-I}$, except $M^{\langle S\rangle}$ at $i$; at $\zeta^{2}$ $M^{-I}$ with its $\Z/3$-action \\
$\ell=2$ & every simplex has $P\supseteq\{\pm I\}$; $P=\langle S\rangle\cong\Z/4$ at $i$ & $P=\{\pm I\}$ & $M$, except $M^{\langle(ST)^{2}\rangle}$ at $\zeta^{2}$ \\
over $\Z$ & $\langle ST\rangle\cong\Z/6$ at $\zeta^{2}$, $\langle S\rangle$ at $i$ & $\{\pm I\}$ & none (no order is a unit) \\
\bottomrule
\end{tabular}
\end{center}
``Relevant at $\ell$'' is thus a statement about which columns of a model may be truncated to degree $0$, not about a different orbifold.
\end{definition}

\begin{definition}[The triangulations in use]\label{def:mpRZ}
All are $\Ga$-equivariant subdivisions of $\Xbulk$ (\ref{H4}); they differ only in the descriptor.
\LEAN{def:mp-models} ((1) and (2) as MP descriptors: $\mathsf M_7$ computed from the dump that defines [GT]'s \code{model7}, $\mathsf R_{\Z}$ from the C++ dump
\code{ORBIFOLD\_MP\_DUMP}; the identification of the dumps with the geometric complexes is the trust boundary, as for \ref{H4}(i)) \LEAN{def:mp-rz-tau} ((2) as a
$\Ga$-complex descriptor with the $\Xbulk$ axioms: $9$ vertices, $15$ edges, $7$ triangles of the lifted domain, all transporters; not (3))
\begin{enumerate}
\item $\mathsf M_7$: $\Xbulk$ with the triangulation of Definition~\ref{def:model7} (vertex orbits $\{v_0,v_3\},\{v_1\},\{v_2\},\{v_4,v_5\},\{v_6\}$; quotient counts
  $(V,E,F)=(5,9,5)$, the representatives $\Sigma_s$ of Definition~\ref{def:isomodel}).
\item $\mathsf R_{\Z}$: the subdivision of the edge orbit of $v_4v_6$ at one interior point $m$ (its stabiliser is $\{\pm I\}$, which fixes the edge pointwise), so
  that $v_5v_6=S\cdot v_4v_6$ is cut at $m':=S\,m$. Lifted domain: vertices $v_0,\dots,v_6,m,m'$; triangles $v_0v_1v_4$, $v_1v_2v_6$, $v_2v_3v_5$, $v_1v_4m$,
  $v_1mv_6$, $v_2v_5m'$, $v_2m'v_6$; relation rows those of Definition~\ref{def:model7rel} with the row of $\{v_4,v_6\}\leftarrow\{v_5,v_6\}$ replaced by
  $v_5m'\leftarrow v_4m$, $m'v_6\leftarrow mv_6$ and the vertex row $m'\leftarrow m$, all by $S$, and the central rows on every simplex. Stabilisers: $\langle ST\rangle$,
  $\langle TS\rangle$ at $v_4,v_5$, $\langle S\rangle$ at $v_6$, $\{\pm I\}$ elsewhere, $m$, $m'$ included. Quotient counts $(V,E,F,\chi)=(6,12,7,1)$. In code
  \code{SubdivideEdgeOrbit} applied to $\mathsf M_7$ at the edge orbit whose ends have stabilisers of orders $6$ and $4$, found from the inventory, not by name.
\item Predicate rows only (no model is built on them): $\mathsf R_{\zeta}$, slot~17's subdivision of the three edge orbits at $\zeta^{2}$
  (Remark~\ref{rmk:isomodellmp}(3)); $\mathsf R_{\Z}^{+}$, $\mathsf R_{\Z}$ with the edges $v_0v_4$ (and $v_3v_5$), $v_1v_4$, $v_2v_5$ cut once; $\mathsf R_{\mathrm{good}}$,
  $v_4v_6$ cut twice (and its $S$-image), $v_1v_4$, $v_1v_6$, $v_2v_5$, $v_2v_6$ cut once and the two triangles at $i$ coned off at their centres
  ([MP] \S3.1; the table of \file{docs/implementation/slot22/mp/mp\_rules.py}).
\end{enumerate}
\end{definition}

\begin{table}[ht]
\centering\small
\begin{tabular}{@{}llccccl@{}}
\toprule
triangulation & $(V,E,F,\chi)$ & adapted for $\mathfrak T^{(3)}$ & adapted for $\mathfrak T^{(2)}$ & adapted (MP's rule) & good & fails at \\
\midrule
$\mathsf M_7$ (shipped) & $(5,9,5,1)$ & yes & yes & \textbf{no} & no & $v_1,v_2,\zeta^{2},i$ \\
$\mathsf R_{\Z}$ & $(6,12,7,1)$ & yes & yes & \textbf{yes} & no & $v_1,v_2,m$ \\
$\mathsf R_{\zeta}$ & $(8,18,11,1)$ & yes & yes & yes & no & three vertices \\
$\mathsf R_{\Z}^{+}$ & $(9,21,13,1)$ & yes & yes & yes & no & $m$ \\
$\mathsf R_{\mathrm{good}}$ & $(13,33,21,1)$ & yes & yes & yes & \textbf{yes} & --- \\
\bottomrule
\end{tabular}
\caption[Moerdijk--Pronk's rules on the candidate triangulations of $\Xbulk$]{Moerdijk--Pronk's rules on the candidate triangulations of $\Xbulk$ (Definitions~\ref{def:mprules}, \ref{def:mpRZ}); quotient counts. The collar
predicate is false on every row. ``Fails at'' lists the vertices whose closed star is not good (finite-group and strong form of Definition~\ref{def:mprules}(3) alike).
\code{OrbifoldMPRulesValidation}. Rows $\mathsf M_7$ and $\mathsf R_{\Z}$ only, every column: \LEAN{thm:mp-model7-rules} \LEAN{thm:mp-rz-rules} \LEAN{thm:mp-rz-dump}
(the C++ predicate rows of these two rows are the Lean predicates).}\label{tab:mprules}
\end{table}

\begin{remark}[Reading the rules on $\Xbulk$]\label{rmk:mprulesread}
\PROVED{} \COMP\textsuperscript{r}
\begin{enumerate}
\item \textbf{$\mathsf M_7$ is not adapted, at every $\ell$ and over $\Z$, for $\SL_2(\Z)$ and for the effective $\mathrm{PSL}_2(\Z)$ alike.} The edge $v_4v_6$ has stabiliser
  $\{\pm I\}$ (resp.\ $1$) inside, $\langle ST\rangle\cong\Z/6$ and $\langle S\rangle\cong\Z/4$ (resp.\ $\Z/3$, $\Z/2$) at its ends (Lemma~\ref{lem:model7coincidences}(1)),
  neither containing the other, so neither $v_1v_4v_6$ nor $v_2v_5v_6$ has a dominating vertex. Rule (ii) knows no prime.
\item \textbf{$\mathsf R_{\Z}$ is adapted}, and is the smallest listed refinement that is: each of its seven triangles has at most one vertex with stabiliser
  $\ne\{\pm I\}$, which dominates. So for Moerdijk--Pronk's construction \emph{as they define it}, a refinement is forced, and it is the same $\mathsf R_{\Z}$ at $\ell=2$,
  $\ell=3$ and over $\Z$ (owner ruling Q2 of 2026-10-09). The Sylow data $\mathfrak T^{(3)}$, $\mathfrak T^{(2)}$ of $\mathsf M_7$ are adapted ($\Z/3$ at $\zeta^{2}$ only,
  resp.\ $\{\pm I\}$ everywhere and $\Z/4$ at $i$ only); this enters no theorem below and is not conformance.
\item \textbf{Goodness} fails on $\mathsf R_{\Z}$ and $\mathsf R_{\Z}^{+}$, at different vertices (\code{good\_fails\_at}): on $\mathsf R_{\Z}^{+}$ at $m$ only, on $\mathsf R_{\Z}$ at $v_1$, $v_2$ and $m$
  (the labels $[1,2,7]$ of the oracle); there the closed star meets both the orbit of $\zeta^{2}$ and that of $i$, so $\mathsf T(v)$
  contains generators of $\langle ST\rangle$ and $\langle S\rangle$, which generate $\Ga$. Decidability: the finite subgroups of $\SL_2(\Z)$ are cyclic of order
  $\le6$ (a finite subgroup fixes a point of $\Hyp$, so lies in that point's stabiliser, conjugate to a subgroup of $\langle S\rangle$ or $\langle ST\rangle$, \ref{H4}(iv);
  the orders $1,2,3,4,6$ as listed in Proposition~\ref{prop:modeltorsion}'s proof), so the closure of $\mathsf T(v)$ under products may be declared infinite once it
  exceeds $6$ elements; on the four oracle rows the finite-group form and the strong form of Definition~\ref{def:mprules}(3) (some $p\in\mathrm{St}(v)$ with $\mathsf T(v)\subseteq\Ga_p$) agree (\code{good\_forms\_agree}). The stronger reading $\mathsf T(v)\subseteq\Ga_v$ would fail at every vertex orbit of $\mathsf M_7$ and at $v_0,v_1,v_2,m$ on $\mathsf R_{\Z}$
  (Lean, \code{goodVtx\_discrepancy}), which is why it is not the definition. Goodness is Moerdijk--Pronk's hypothesis for \emph{their} sheaf-side proof only; nothing in this chapter
  uses it, nor the collar, nor adaptedness (Theorem~\ref{thm:mpcomparison}, Remark~\ref{rmk:mphybrid}, Theorem~\ref{thm:mpnatural}).
\item \textbf{The subgroup form is necessary}: comparing stabiliser orders (as \file{mp\_rules.py} does) agrees with it on these rows only because every
  non-maximal stabiliser is the central $\{\pm I\}$; with a planted non-central sub-stabiliser the order form accepts what the subgroup form rejects.
\end{enumerate}
\COMP\textsuperscript{r}: \code{OrbifoldMPRulesValidation} (\code{fast}) reproduces Table~\ref{tab:mprules} on its four oracle rows against \file{mp\_rules.py} ($4\times3=12$ gated cells, \code{ORACLE
mp\_rules.py cells=12 agree=12}); the $\mathsf R_{\zeta}$ row is reported only (an ungated \code{OUTPUT P1.12} line). (SC) and (IO) are checked on all five, the descriptor's quotient counts against \code{ModelData}'s representatives, and the negative
controls of (4) and of a refused two-vertices-in-one-orbit order.
\end{remark}

\subsection{The Moerdijk--Pronk set, choice-free}\label{sec:mpsetsec}

\begin{definition}[The MP set $\mathsf S(Y;\mathfrak T)$, its local system and its cochains]\label{def:mpset}
Let $\mathfrak T$ be an orbifold triangulation of $[\Ga\backslash Y]$ with (SC) and (PD). For top representatives $\sigma,\sigma'\in\Sigma_{\max}$ put
$\mathrm{Ar}(\sigma,\sigma'):=\{\gamma\in\Ga:\iota(\sigma)\cap\gamma\,\iota(\sigma')\ne\emptyset\}$ (the arrows $\sigma\xleftarrow{\gamma}\sigma'$). For $k\ge0$,
\[
  \mathsf S_k:=\bigl\{(\sigma_0,\dots,\sigma_k;\gamma_1,\dots,\gamma_k):\ \sigma_i\in\Sigma_{\max},\ \gamma_i\in\Ga,\
    \iota(\sigma_0)\cap\gamma_1\iota(\sigma_1)\cap\gamma_1\gamma_2\iota(\sigma_2)\cap\cdots\cap\gamma_1\cdots\gamma_k\iota(\sigma_k)\ne\emptyset\bigr\}.
\]
Faces: $d_0$ drops $(\sigma_0,\gamma_1)$; $d_i$ ($0<i<k$) drops $\sigma_i$ and replaces $(\gamma_i,\gamma_{i+1})$ by $\gamma_i\gamma_{i+1}$; $d_k$ drops $(\sigma_k,\gamma_k)$.
Degeneracies insert $(\sigma_i,1)$. The \textbf{local system} $\underline M$ has fibre $M$ at every $\sigma\in\Sigma_{\max}$, the arrow $\sigma_0\xleftarrow{\gamma}\sigma_1$
acting by $m\mapsto\gamma\cdot m$. The \textbf{MP complex} is the normalised cochain complex
\[
  C^{k}(\mathsf S;\underline M):=\{f\colon\mathsf S_k^{\mathrm{nd}}\to M\},\qquad
  (df)(s):=\gamma_1\cdot f(d_0s)+\sum_{i=1}^{k+1}(-1)^{i}f(d_is)\quad(s\in\mathsf S_{k+1}^{\mathrm{nd}}),
\]
with $f:=0$ on degenerate strings; $\mathsf S_k^{\mathrm{nd}}$ is the set of non-degenerate strings (no $i$ with $\sigma_i=\sigma_{i+1}$ and $\gamma_{i+1}=1$). The same
formula on all functions $\mathsf S_k\to M$ defines the \textbf{unnormalised MP complex} $C^{\bullet}_{\mathrm{un}}(\mathsf S;\underline M)$, which contains
$C^{\bullet}(\mathsf S;\underline M)$ as the subcomplex of functions vanishing on degenerate strings (v1.22.0-pre.6: named, so that the Lean markers can be clause-scoped). On the
triangulations of Definition~\ref{def:mpRZ}, $\mathsf S(\mathsf M_7)$ and $\mathsf S(\mathsf R_{\Z})$, with $M=\VR$, $\gamma\cdot m=\rhon(\gamma)m$.
\LEAN{def:mp-cover} \LEAN{lem:mp-set-gamma} \LEAN{lem:mp-face-face} ($\mathsf S_k$ as the normal forms $g_0=1$ of $\mathfrak F'_k$ with the coordinates $(\sigma_i;\gamma_i)$ as an explicit bijection, the
intersection read as a common vertex, which is the same under (SC), Remark~\ref{rmk:mpwhat}(1); its faces and degeneracies and the simplicial identities, the middle faces as deletions)
\LEAN{thm:mp-dd} (the twisted differential and $d\circ d=0$ on $C^{\bullet}_{\mathrm{un}}$)
\LEAN{thm:mp-norm-dgn} (the normalised complex: the functions vanishing on the degenerate strings $s_i(s)$ form a subcomplex of $C^{\bullet}_{\mathrm{un}}$ for the twisted differential,
isomorphic to the cochains of the simplicial object of Definition~\ref{def:mpfree}'s marker that vanish on its degeneracies, and its inclusion is a quasi-isomorphism; v1.22.0-pre.7)
\end{definition}

\begin{remark}[What the MP set is, and when it is Moerdijk--Pronk's]\label{rmk:mpwhat}
\PROVED{} \COMP\textsuperscript{r}
\begin{enumerate}
\item \textbf{Finite and decidable.} $\mathrm{Ar}(\sigma,\sigma')$ is the finite union, over vertex pairs $v$ of $\iota(\sigma)$, $w$ of $\iota(\sigma')$ in one orbit, of the cosets
  $\mathrm{Tr}(v,w)$ of $\Ga_w$; under (SC) a string is in $\mathsf S_k$ iff some vertex $p$ of $\iota(\sigma_0)$ has $(\gamma_1\cdots\gamma_j)^{-1}p\in\iota(\sigma_j)$ for all $j$
  (a finite intersection of closed simplices is a common face, hence non-empty iff it has a vertex). The faces of a string are strings; $\mathsf S$ is a simplicial set.
  No chart, lift, $\mu$-map, effectivity or adaptedness enters the definition.
  \LEAN{thm:mp-datum-arrows} \LEAN{thm:mp-models-cover} ($\mathrm{Ar}(\sigma,\sigma')$ the finite union of the cosets $\mathrm{Tr}(v,w)$ on a descriptor; on $\mathsf M_7$ and $\mathsf R_{\Z}$
  the counts of all strings $\abs{\mathsf S_0},\abs{\mathsf S_1}=5,132$ and $7,172$)
\item \textbf{What it is.} Under (SC), $\mathsf S=\pi_0$ of the nerve of the action groupoid $\Ga\ltimes Y$ restricted to the transversal $\bigsqcup_{\sigma\in\Sigma_{\max}}\iota(\sigma)$:
  the component of a $k$-simplex of that nerve is the closed simplex $\iota(\sigma_0)\cap\gamma_1\iota(\sigma_1)\cap\cdots$, which is contractible.
\item \textbf{When it is Moerdijk--Pronk's $S(\mathcal T)$.} If $\Ga$ acts effectively, $\mathfrak T$ is adapted and the quotient is a simplicial complex, the intersection is a face
  $\tau$ with a maximal vertex $v(\tau)$ and the $\gamma_i$ range over a coset of $\Ga_{v(\tau)}$ --- under adaptedness the stabilisers of the vertices shared by
  $\iota(\sigma_{i-1})$ and $\gamma_i\iota(\sigma_i)$ are nested in $\Ga_{v(\tau)}$, so the union of cosets in (1) is a single coset $h\Ga_{v(\tau)}$, and effectivity and the
  simplicial quotient make the shared vertex set the unique lift of $\bar\sigma_{i-1}\cap\bar\sigma_i$: $\mathsf S_k\cong\bigsqcup\Ga_{v(\tau)}^{k}$, their formula (Remark~\ref{rmk:mpcaveat}(1)(ii)).
  For $\Ga=\SL_2(\Z)$, \textbf{not} effective, their own construction gives the $\mathrm{PSL}_2(\Z)$-set (Remark~\ref{rmk:mpcaveat}(2)); $\mathsf S$ keeps $-I$ as an arrow
  $\sigma\xleftarrow{-I}\sigma$ of every object and is a \textbf{non-effective generalisation} of their object, not their object (owner ruling Q1 of 2026-10-09: this
  choice-free form is normative for the MP model). On a non-adapted $\mathfrak T$ $\mathrm{Ar}(\sigma,\sigma)$ need not be a group (on $\mathsf M_7$, $\mathrm{Ar}(\sigma,\sigma)$ at
  $v_1v_4v_6$ is $\langle ST\rangle\cup\langle S\rangle$, eight elements), and their formula no longer describes $\mathsf S$, which is still correct
  (Theorem~\ref{thm:mpcomparison}).
\end{enumerate}
\COMP\textsuperscript{r} (\code{OrbifoldMPRulesValidation}, \code{OrbifoldMPNerveValidation}): $\abs{\mathsf S_0^{\mathrm{nd}}},\dots,\abs{\mathsf S_3^{\mathrm{nd}}}=5,127,2541,53143$ on
$\mathsf M_7$ and $7,165,2927,56349$ on $\mathsf R_{\Z}$; all strings $\abs{\mathsf S_1},\abs{\mathsf S_2}=172,3264$ on $\mathsf R_{\Z}$ by both $\mathsf S$ and the $\mu$-form count
$\sum_\tau\abs{\Ga_{v(\tau)}}^{k}$, against, on the non-adapted $\mathsf M_7$, $132,2800$ by $\mathsf S$ versus $124,2704$ by the $\mu$-form count (\code{ORACLE adapted\_form\_check.py cells=6 agree=6}); faces closed,
simplicial identities exhaustive to $k=4$, every string once.
\end{remark}

\subsection{The comparison theorem}\label{sec:mpcompsec}

\begin{definition}[The free bisimplicial $\Ga$-set $\mathfrak F$ and the labelled Čech nerve $\mathfrak F'$]\label{def:mpfree}
Let $Y$, $\mathfrak T$ satisfy (SC), (IO), (PD). Let $Y_s$ be the set of non-degenerate ordered $s$-simplices of $Y$ (one increasing vertex tuple per geometric simplex,
for the order of (IO)), a $\Delta$-set with $0\le s\le\dim Y$. For $k\ge0$ put
\[
  \mathfrak F_{k,s}:=\bigl\{\bigl((g_0,\sigma_0),\dots,(g_k,\sigma_k);x\bigr):\ g_i\in\Ga,\ \sigma_i\in\Sigma_{\max},\ x\in Y_s,\ x\subseteq g_0\iota(\sigma_0)\cap\cdots\cap g_k\iota(\sigma_k)\bigr\},
\]
$\mathfrak F'_k:=\{((g_0,\sigma_0),\dots,(g_k,\sigma_k)):\bigcap_ig_i\iota(\sigma_i)\ne\emptyset\}$, with $\Ga$ acting by left multiplication on every $g_i$ and on $x$. Faces in $k$ drop a
pair, degeneracies repeat one (simplicial in $k$); faces in $s$ are those of $Y_{\bullet}$ (semi-simplicial in $s$). $\Z[\mathfrak F_{\bullet,\bullet}]$ is a double complex,
bounded in $s$; $\mathrm{Tot}$ is its total complex; $\Z[\mathfrak F'_{\bullet}]$ is the chain complex of the simplicial set $\mathfrak F'$, augmented by
$(g_0,\sigma_0)\mapsto1$, and $N\Z[\mathfrak F'_\bullet]$ its normalised quotient. For $x\in Y_s$, $O(x):=\{(g,\sigma):x\subseteq g\iota(\sigma)\}$.
\LEAN{def:mp-bisimp} ((IO), $Y_s$, $\mathfrak F_{k,s}$ in the bounded form $\dim Y\le2$ with both kinds of faces, and the augmented complexes $\cring[\mathfrak F_{\bullet,s}]\to\cring[Y_s]$,
$\cring[\mathfrak F'_\bullet]\to\cring$, $\cring[Y_\bullet]\to\cring$ over every commutative ring; not the degeneracies in $k$, not $N\Z[\mathfrak F'_\bullet]$)
\LEAN{def:mp-simpK} \LEAN{thm:mp-simpK-alt} (v1.22.0-pre.7: $[k]\mapsto\cring[\mathfrak F'_k]$ as a simplicial object in $\cring[\Ga]$-modules, faces the deletions and degeneracies the
repetitions, whose alternating face map complex is $\cring[\mathfrak F'_\bullet]$, over every commutative ring; not $N\Z[\mathfrak F'_\bullet]$ as a quotient complex)
\end{definition}

\begin{lemma}[Rows]\label{lem:mprows}
\PROVED{} \LEAN{lem:mp-rows} (the augmented row contracted by the extra degeneracy, non-equivariantly, over every commutative ring) For each $s$, the augmentation $\Z[\mathfrak F_{\bullet,s}]\to\Z[Y_s]$, $((g_i,\sigma_i);x)\mapsto x$, is a quasi-isomorphism of complexes of $\Z[\Ga]$-modules.
\end{lemma}

\begin{proof}
It is $\Ga$-equivariant, and a quasi-isomorphism is checked non-equivariantly. $\mathfrak F_{\bullet,s}=\bigsqcup_{x\in Y_s}O(x)^{\bullet+1}$, and $(g_i,\sigma_i)\in O(x)$ for each $i$
is exactly the condition $x\subseteq\bigcap_ig_i\iota(\sigma_i)$. $O(x)\ne\emptyset$ by (PD). $O(x)^{\bullet+1}$ is the nerve of the chaotic groupoid on $O(x)$ (the Čech nerve of
$O(x)\to\mathrm{pt}$), and for a fixed $o\in O(x)$ the extra degeneracy $(o_0,\dots,o_k)\mapsto(o,o_0,\dots,o_k)$ contracts the augmented complex
$\Z[O(x)^{\bullet+1}]\to\Z$. The augmentation kills boundaries: $((a,b);x)\mapsto x-x$.
\end{proof}

\begin{lemma}[Columns]\label{lem:mpcols}
\PROVED{} \LEAN{lem:mp-cols} (the cone contraction on the first vertex $v^{\ast}$, non-equivariant, for $\dim Y\le2$) For each $k$, the augmentation $\Z[\mathfrak F_{k,\bullet}]\to\Z[\mathfrak F'_k]$ (concentrated in $s=0$), $((g_i,\sigma_i);x)\mapsto((g_i,\sigma_i))$, is a quasi-isomorphism
of complexes of $\Z[\Ga]$-modules.
\end{lemma}

\begin{proof}
Again checked non-equivariantly, fibre by fibre. Over a string $t=((g_i,\sigma_i))\in\mathfrak F'_k$ the $s$-fibre is the set of non-degenerate ordered faces of
$\tau_t:=\bigcap_ig_i\iota(\sigma_i)$, which by (SC) (and induction on $k$) is a closed simplex, non-empty by the definition of $\mathfrak F'_k$; with the order of (IO) its
augmented ordered chain complex is that of a standard simplex, contracted by the cone on its first vertex $v^{\ast}$ ($w_0\cdots w_j\mapsto v^{\ast}w_0\cdots w_j$ if
$w_0\ne v^{\ast}$, else $0$). \emph{No equivariant homotopy is claimed}: $v^{\ast}(\tau_t)$ is not compatible with the $k$-faces (a $k$-face enlarges $\tau_t$ and may change its
first vertex). The augmentation kills boundaries: an edge $(w_0w_1)$ maps to $t-t$.
\end{proof}

\begin{lemma}[Bounded assembly]\label{lem:mpassembly}
\PROVED{} Let $f\colon A\to B$ be a map of double complexes of modules over a ring, $A_{k,s}=B_{k,s}=0$ unless $k\ge0$ and $0\le s\le N$. If $f$ is a quasi-isomorphism on
every row (fixed $s$), or on every column (fixed $k$), then $\mathrm{Tot}\,f$ is a quasi-isomorphism.
\LEAN{thm:mp-assembly-grid} (only the acyclic form used in the proof of Theorem~\ref{thm:mpcomparison}(2): in a four-row bounded double complex with exact columns and
acyclic upper rows the bottom row is acyclic; not the quasi-isomorphism form)
\end{lemma}

\begin{proof}
Replace $f$ by its mapping cone formed in the direction in which $f$ is a quasi-isomorphism: a double complex $C$, bounded in the same way up to one step, with exact rows,
resp.\ exact columns, and $\mathrm{Tot}\,C=\mathrm{Cone}(\mathrm{Tot}\,f)$ up to sign; it suffices that $\mathrm{Tot}\,C$ is acyclic. Exact rows: filter $\mathrm{Tot}\,C$ by $s$; the
filtration has finitely many steps, each graded piece is a shifted row, hence acyclic, and induction over the long exact sequences gives the claim. Exact columns: filter by
$k\le p$; the filtration is exhaustive with $F_{-1}=0$, its graded pieces are shifted columns, so every $F_p$ is acyclic by induction, and $F_{q+1}$ contains $\mathrm{Tot}\,C$ in
every degree $\le q+1$, so $H_q(\mathrm{Tot}\,C)=H_q(F_{q+1})=0$. (Equivalently the acyclic assembly lemma, \STD{\cite[Lemma~2.7.3]{weibel}, locator not checked against the text, \VERIFYloc}, the two totalisations
coinciding because $s$ is bounded.)
\end{proof}

\begin{theorem}[Comparison: the MP complex computes group cohomology in every degree]\label{thm:mpcomparison}
\PROVED{} Let $\Ga$ be a group and $Y$ a simplicial $\Ga$-complex with \ref{G1}, \ref{G3}, (SC), (AC), (IO), (PD) (Definition~\ref{def:mptri}), and $\mathfrak T$ its orbifold
triangulation. Then:
\begin{enumerate}
\item $\Ga$ acts freely on $\mathfrak F'_k$ (already on $g_0$); the normal form $g_0=1$ identifies $\Ga\backslash\mathfrak F'_k$ with $\mathsf S_k$, by
  $((1,\sigma_0),(\gamma_1,\sigma_1),(\gamma_1\gamma_2,\sigma_2),\dots)\leftrightarrow(\sigma_i;\gamma_i)$, $\gamma_i=g_{i-1}^{-1}g_i$; and $\mathsf S_k$ is finite.
  \LEAN{lem:mp-free} (the free action and the normal form as a transversal) \LEAN{lem:mp-set-finite} (finiteness, given finitely many maximal representatives and finite arrow sets)
\item $\Z[\mathfrak F'_\bullet]\to\Z$ and $N\Z[\mathfrak F'_\bullet]\to\Z$ are resolutions of $\Z$ by free $\Z[\Ga]$-modules, the latter of rank $\abs{\mathsf S_k^{\mathrm{nd}}}$ in degree
  $k$ --- the labelled Čech nerve of the cover of $Y$ by its closed maximal simplices.
  \LEAN{thm:mp-resolution} \LEAN{thm:mp-projres} (the first: $\cring[\mathfrak F'_\bullet]\to\cring$ a resolution of $\cring$ by projective $\cring[\Ga]$-modules, for every commutative ring $\cring$
  with (AC) over $\cring$, (IO), finite vertex sets and $\dim Y\le2$; not $N\Z[\mathfrak F'_\bullet]$)
  \LEAN{thm:mp-norm-heq} (v1.22.0-pre.7: the normalised clause dually, at the level of cochains --- the cochains vanishing on the degeneracies are a homotopy equivalent
  subcomplex of $\Hom_{\cring[\Ga]}(\cring[\mathfrak F'_\bullet],M)$, every ring, every module; not $N\Z[\mathfrak F'_\bullet]\to\Z$ as a resolution, not its rank $\abs{\mathsf S_k^{\mathrm{nd}}}$)
\item $\Hom_{\Z[\Ga]}(N\Z[\mathfrak F'_\bullet],M)=C^{\bullet}(\mathsf S;\underline M)$ \textbf{on the nose}, $F\mapsto(f(s):=F(x_s))$, $x_s$ the normal form of $s$; the twist $\gamma_1$ of
  the differential is the re-normalisation of $d_0x_s=\gamma_1\cdot x_{d_0s}$. Likewise $\Hom_{\Z[\Ga]}(\Z[\mathfrak F'_\bullet],M)=C^{\bullet}_{\mathrm{un}}(\mathsf S;\underline M)$.
  \LEAN{thm:mp-homequiv} (the unnormalised identity, over every commutative ring)
\item Hence, for every commutative ring $\cring$ and every $\cring[\Ga]$-module $M$, naturally in $M$,
  \[
    H^{k}\bigl(C^{\bullet}(\mathsf S(Y;\mathfrak T);\underline M)\bigr)\;\cong\;H^{k}(\Ga;M)\qquad\text{for every }k\ge0 ,
  \]
  and the same for $C^{\bullet}_{\mathrm{un}}(\mathsf S(Y;\mathfrak T);\underline M)$.
  No hypothesis on $\abs{\Ga_\sigma}$ being a unit, no effectivity, no adaptedness and no choice of charts is used; the central $-I$ of $\SL_2(\Z)$ is included.
  \LEAN{thm:mp-comparison} (for $C^{\bullet}_{\mathrm{un}}$: every group, every commutative ring $\cring$, every $\cring[\Ga]$-module, every $k$, with (AC) over $\cring$ --- implied by (AC) over $\Z$ ---,
  (IO), finite vertex sets and $\dim Y\le2$ as hypotheses) \LEAN{thm:mp-comparison-natural} (naturality in $M$, for the isomorphism built through a homotopy equivalence with the bar resolution)
  \LEAN{thm:mp-comparison-normalised} (v1.22.0-pre.7: the displayed isomorphism for the normalised $C^{\bullet}$, under the same hypotheses; on $\mathsf M_7$ and $\mathsf R_{\Z}$ with no
  hypothesis left)
\item For $\Ga=\SL_2(\Z)$, both triangulations $\mathsf M_7$ and $\mathsf R_{\Z}$ of Definition~\ref{def:mpRZ} satisfy the hypotheses: \ref{G1}, (AC) and (SC) by \ref{H4}
  ($\Xbulk$ is a contractible simplicial $\Ga$-complex, and a $\Ga$-equivariant subdivision is again one, Remark~\ref{rmk:orbifoldrefine}(1)), \ref{G3} by the finite
  quotient, (PD) since $\Xbulk$ is a $2$-dimensional pseudo-manifold with boundary built from triangles, and (IO) by the vertex orbits listed in Definition~\ref{def:mpRZ},
  each triangle meeting each orbit at most once.
  \LEAN{thm:mp-models-cover} ((IO), finite vertex sets and $\dim\le2$ for both, on the descriptors of Definition~\ref{def:mpRZ}'s marker) \LEAN{thm:mp-ac-bridge} \LEAN{thm:mp-model7-ac}
  \LEAN{thm:mp-retr} \LEAN{thm:mp-rz-ac} ((AC) for both over every commutative ring: for $\mathsf M_7$ from \ref{H4}(ii) through the identification of the ordered with the oriented chains; for
  $\mathsf R_{\Z}$ \ref{H4}(ii) of the subdivision, by a deformation retraction onto $\mathsf M_7$ given by finite tables; so, with (4), the comparison for $C^{\bullet}_{\mathrm{un}}$ on $\mathsf M_7$
  and on $\mathsf R_{\Z}$ with no hypothesis left. \ref{G1}, \ref{G3} are not restated)
\end{enumerate}
\end{theorem}

\begin{proof}
(1) $\Ga$ acts freely on the first label. A string is in $\mathfrak F'_k$ iff $\bigcap g_i\iota(\sigma_i)\ne\emptyset$, iff (multiplying by $g_0^{-1}$) its normal form satisfies the
condition of $\mathsf S_k$; distinct normal forms lie in distinct orbits. Finiteness: $\gamma_i\in\mathrm{Ar}(\sigma_{i-1},\sigma_i)$, a finite set (Remark~\ref{rmk:mpwhat}(1),
\ref{G1}, \ref{G3}).
(2) By Lemma~\ref{lem:mprows} the row augmentations form a map of double complexes $\Z[\mathfrak F_{\bullet,\bullet}]\to\Z[Y_\bullet]$ (the target concentrated in $k=0$; the
$s$-faces of both are those of $Y_\bullet$) that is a quasi-isomorphism on every row, so $\mathrm{Tot}\to C^{\mathrm{ord}}_{\bullet}(Y)$ is a quasi-isomorphism by
Lemma~\ref{lem:mpassembly}; under (IO) the ordered chain complex of $Y_\bullet$ is the oriented chain complex $C_\bullet(Y)$ (one generator per simplex, the orientation induced
by the order; a stored orientation differs by signs only), and $C_\bullet(Y)\to\Z$ is a quasi-isomorphism by (AC). By Lemma~\ref{lem:mpcols} the column augmentations form a
map $\Z[\mathfrak F_{\bullet,\bullet}]\to\Z[\mathfrak F'_\bullet]$ (concentrated in $s=0$; the $k$-faces of both drop a pair) that is a quasi-isomorphism on every column, so
$\mathrm{Tot}\to\Z[\mathfrak F'_\bullet]$ is one by Lemma~\ref{lem:mpassembly}. The two composites $\mathrm{Tot}\to\Z$ agree (both send $((g,\sigma);v)\in\mathfrak F_{0,0}$ to $1$ and
vanish elsewhere), so $\Z[\mathfrak F'_\bullet]\to\Z$ is a quasi-isomorphism. It is free over $\Z[\Ga]$ by (1). The projection onto the normalised quotient is a quasi-isomorphism
\STD{\cite[\S8.3]{weibel}, the normalised chain complex of a simplicial abelian group; section-level pointer, \VERIFYloc}, $\Ga$ preserves degeneracy, and the
quotient is free on the non-degenerate normal forms.
(3) $\Hom_{\Z[\Ga]}(\Z[\Ga\cdot x_s],M)=M$ by $F\mapsto F(x_s)$. For $i\ge1$, $d_ix_s$ is already in normal form and is $x_{d_is}$; $d_0x_s=((\gamma_1,\sigma_1),(\gamma_1\gamma_2,\sigma_2),
\dots)=\gamma_1\cdot x_{d_0s}$, so $F(d_0x_s)=\gamma_1\cdot f(d_0s)$; degenerate strings are killed in the normalised quotient. Without the quotient the same computation
gives the unnormalised identity.
(4) By (2), $H^{k}(\Hom_{\Z[\Ga]}(N\Z[\mathfrak F'],M))=\Ext^{k}_{\Z[\Ga]}(\Z,M)=H^{k}(\Ga;M)$, the definition of group cohomology, independent of the free resolution
\STD{\cite[Thm.~2.2.6]{weibel}}, natural in $M$; over $\cring$, $\cring\otimes N\Z[\mathfrak F']\to\cring$ is a free $\cring[\Ga]$-resolution ($N\Z[\mathfrak F']\to\Z$ is a bounded-below
exact complex of free abelian groups, hence split over $\Z$), and $\Hom_{\cring[\Ga]}(\cring\otimes N\Z[\mathfrak F'],M)=\Hom_{\Z[\Ga]}(N\Z[\mathfrak F'],M)$. The same
argument on the free resolution $\Z[\mathfrak F']$ of (2) gives the statement for $C^{\bullet}_{\mathrm{un}}$. Nothing in (1)--(3) used
an order being invertible, effectivity or adaptedness.
(5) \ref{H4}(i),(ii) and Remark~\ref{rmk:orbifoldrefine}(1); the orbits are read off Definition~\ref{def:model7rel} (Lemma~\ref{lem:model7coincidences}(2)) and, for $\mathsf R_{\Z}$, the rows of
Definition~\ref{def:mpRZ}(2): $v_0v_1v_4$, $v_1v_2v_6$, $v_2v_3v_5$, $v_1v_4v_6$, $v_2v_5v_6$, $v_1v_4m$, $v_1mv_6$, $v_2v_5m'$, $v_2m'v_6$ each meet the orbits
$\{v_0,v_3\}$, $\{v_4,v_5\}$, $\{m,m'\}$ at most once.
\end{proof}

\noindent\emph{Reading.} The resolution direction is the nerve direction $k$; the triangulation direction $s$ stays finite ($0\le s\le2$). That is the precise sense in which a
finite triangulation accounts for all degrees: $\mathsf S_k\ne\emptyset$ for every $k$ (strings $\sigma\xleftarrow{w}\sigma\xleftarrow{w}\cdots$, $w\in\Ga_\sigma$), and
$H^{k}(\Ga;\F_2)=\F_2$ for every $k$ is seen by $\mathsf S(\mathsf M_7)$. Read the other way, row-wise the $\Ga$-cochains of $\Z[\mathfrak F]$ compute
$\bigoplus_{x\in\Ga\backslash Y_s}H^{k}(\Ga_x;M)$, the $E_1$ page of the isotropy spectral sequence of Theorem~\ref{thm:specialisation}'s proof; the MP set is that spectral sequence's
total complex made finite in each bidegree by passing to $\pi_0$. Refinement is \emph{not} needed for correctness: the theorem holds on $\mathsf M_7$ at every $\ell$ and over $\Z$;
what adaptedness buys is the description $\mathsf S_k\cong\bigsqcup\Ga_{v(\tau)}^{k}$, i.e.\ being Moerdijk--Pronk's set in their form (Remark~\ref{rmk:mpwhat}(3)).

\begin{remark}[The $\ell$-local hybrid is Proposition~\ref{prop:isomodelell}]\label{rmk:mphybrid}
\PROVED{} \COMP\textsuperscript{r}
\begin{enumerate}
\item \textbf{What it is.} Over a $\Z_{(\ell)}$-algebra $\cring$, the hybrid complex of the $\ell$-relevant structure (Definition~\ref{def:mpell}) --- the column of $\sigma$ the
  inflation to $\Ga_\sigma$ of a resolution of $\cring$ over $\cring[P_\sigma]$, with $\Hom$ into $M$ equal to cochains of $P_\sigma$ with values in $M^{Q_\sigma}$ --- is an
  $\ell$-local isotropy model in the sense of Definition~\ref{def:isomodelell} with the column of $\sigma$ any projective $\cring[\Ga_\sigma]$-resolution of $\cring$ through
  $\cring[\Ga_\sigma/Q_\sigma]$: \textbf{the shipped $\Iso^{(\ell)}$ (\code{integral\_model::LocalModel}) is its periodic-column instance}, and
  Proposition~\ref{prop:isomodelell}(2),(3) is its theorem (a projective resolution; $H^{q}(\Ga;M)$ in every degree). No adaptedness is assumed. The bar-column form ([MP]
  \S2.3: the inflated bar resolutions of $P_\sigma$, projective because $\cring[\Ga_\sigma/Q_\sigma]=\cring[\Ga_\sigma]e_\sigma$) satisfies the hypotheses of the same proof, which
  uses only projective columns with $\cring$-linear contracting homotopies; so it is again a projective resolution of $\cring$, homotopy equivalent to the periodic-column one
  over $\cring[\Ga]$ \STD{\cite[Thm.~2.2.6]{weibel}}, and both compare with $\mathsf S$ through Theorem~\ref{thm:mpcomparison}. They are compared \emph{on cohomology}, never on
  cochains: the cochain-level models are choice-dependent.
\item \textbf{The three regimes} (Proposition~\ref{prop:ellthreeclosed}, Remark~\ref{rmk:isomodellmp}(2)). $\ell=3$: one resolved column, at $\zeta^{2}$ ($P=\Z/3$), with row~$0$
  $\bigoplus_\sigma M^{Q_\sigma}$ --- \emph{not} the Bredon value $M^{\langle ST\rangle}$ at the corner --- and $H^{q}(\Ga;M)=H^{q}(\Z/3;M^{-I})$ for $q\ge2$. $\ell=2$: every simplex keeps
  the central column $C^{\bullet}(\{\pm I\};M^{Q_\sigma})$ and $i$ carries $C^{\bullet}(\Z/4;M)$, a $\Z/2$-gerbe that no $2$-local reduction removes (the class $e$ of
  Lemma~\ref{lem:bulkperiodic}(1) has order $12$). Over $\Z$: every column is the full one; nothing is a unit, and $H^{2}(\Ga;\Z)\cong\Z/12$ (Lemma~\ref{lem:bulkperiodic}(1))
  takes its $3$-part from $\langle ST\rangle$ and its $2$-part from $\{\pm I\}$ and $\langle S\rangle$.
\item \textbf{Gap G1, closed} (v1.22.0-pre.11; recorded as open until pre.10). In nerve form the $\ell$-reduced object is a simplicial set with a \emph{coefficient system}
  $\sigma\mapsto M^{Q_\sigma}$, not a local system. Pre.10 recorded that its description needs adaptedness for $\mathfrak T^{(\ell)}$ and a choice of lifts making the twists
  $\ell$-elements. Neither is needed: Proposition~\ref{prop:mpqreduced} shows that the \emph{choice-free} quotient $\mathfrak D=\mathfrak F'/{\sim}$ of the labelled nerve by the
  prime-to-$\ell$ parts at a dominant point is a projective resolution under the domination hypothesis (QD) of Definition~\ref{def:mpqd}, which bounds the $Q$'s, not the Sylow
  parts, and holds on $\mathsf M_7$ (not adapted) as well as on $\mathsf R_{\Z}$ at $\ell\in\{2,3\}$. The subgroup form with labels in $P_{v(\tau)}$ does not exist for any choice
  of lifts (the cone-point holonomy, [MPT] \file{docs/latex/mp\_triangulations.tex} \S2.5). The hybrid double complex stays the normative definition of the $\ell$-relevant MP
  model; the $Q$-reduced nerve is a second projective resolution, compared with it \emph{on cohomology} (Proposition~\ref{prop:mpqreduced}(4)).
\end{enumerate}
\COMP\textsuperscript{r} (\code{OrbifoldMPComparisonValidation}, \code{OrbifoldMPAllDegreesValidation}): $\Iso^{(\ell)}$ at $\ell\in\{2,3\}$ has the dimensions of $\CI$ and of the oracle
on $H^{0},\dots,H^{7}$ for $n\le12$, the strict comparison of Proposition~\ref{prop:isomodelell}(4) is an isomorphism on $H^{\le2}$, and composed with $\Phi^{\mathsf{MP}}$
(Definition~\ref{def:mpphi}) on to the nerve for $n\le3$.
\end{remark}

\begin{definition}[$Q$-domination]\label{def:mpqd}
Let $\Ga$, $Y$ satisfy the hypotheses of Theorem~\ref{thm:mpcomparison} with every vertex stabiliser $\Ga_v$ finite abelian, and let $\ell$ be a prime. For a simplex $x$ of $Y$ let
$Q_x\subseteq\Ga_x=\bigcap_{w\in x}\Ga_w$ be the subgroup of elements of order prime to $\ell$ (a subgroup, $\Ga_x$ being abelian; for $\Ga=\SL_2(\Z)$ it is the $Q_\sigma$ of
Definition~\ref{def:isomodelell}). Then $Q_x\subseteq Q_w$ for every vertex $w$ of $x$, $\gamma Q_x\gamma^{-1}=Q_{\gamma x}$, and $Q_{x\cup x'}=Q_x\cap Q_{x'}$ whenever
$x\cup x'$ is a simplex. The triangulation is \textbf{$Q$-dominated at $\ell$}, \textbf{(QD)}, if there are a subgroup $Q_0\subseteq\Ga$ acting trivially on $Y$ and a $\Ga$-stable set
$\mathcal C$ of vertices (the \emph{$Q$-cone vertices}) such that
(a) no simplex of $Y$ has two vertices in $\mathcal C$; (b) $Q_w=Q_0$ for every vertex $w\notin\mathcal C$; (c) $Q_0\subseteq Q_v$ for every $v\in\mathcal C$.
The \emph{dominant vertex} $\mathsf p(x)$ of a simplex $x$ is its vertex in $\mathcal C$ if it has one, otherwise its first vertex for the order of (IO). Under (QD):
\begin{enumerate}
\item $Q_w\subseteq Q_{\mathsf p(x)}$ for every vertex $w$ of $x$;
\item $\mathsf p(x')=\mathsf p(x)$ whenever $\mathsf p(x)\in x'\subseteq x$;
\item $\mathsf p(\gamma x)=\gamma\,\mathsf p(x)$ for $\gamma\in\Ga$.
\end{enumerate}
For $\Ga=\SL_2(\Z)$ on $\mathsf M_7$ and on $\mathsf R_{\Z}$: at $\ell=3$, $Q_0=\{\pm I\}$ and $\mathcal C$ the orbit of $i$ ($Q_i=\langle S\rangle$, $Q=\{\pm I\}$ at every other vertex,
$\zeta^{2}$ included); at $\ell=2$, $Q_0=1$ and $\mathcal C$ the orbit of $\zeta^{2}$ ($Q_{\zeta^{2}}=\langle(ST)^{2}\rangle$, $Q=1$ elsewhere, $i$ included). (QD) is a finite,
decidable statement on a descriptor; it is \emph{not} adaptedness (Definition~\ref{def:mprules}(1),(2)): $\mathsf M_7$ is not adapted, yet (QD) holds on it at $\ell\in\{2,3\}$,
since a triangle of $\mathsf M_7$ containing both $\zeta^{2}$ and $i$ contains only one of them in $\mathcal C$. At $\ell\ge5$, (QD) holds on $\mathsf R_{\Z}$ ($Q_0=\{\pm I\}$,
$\mathcal C$ the two elliptic orbits, which no triangle of $\mathsf R_{\Z}$ joins) but not on $\mathsf M_7$.
\end{definition}

\begin{proof}[Proof of (1)--(3) and of the $\SL_2(\Z)$ clause]
$Q_x\subseteq Q_w$ because $\Ga_x\subseteq\Ga_w$; conjugation carries $\Ga_x$ to $\Ga_{\gamma x}$ and preserves orders; $\Ga_{x\cup x'}=\Ga_x\cap\Ga_{x'}$. (1) If $\mathsf p(x)\in\mathcal C$,
every other vertex $w$ of $x$ is outside $\mathcal C$ by (a), so $Q_w=Q_0\subseteq Q_{\mathsf p(x)}$ by (b), (c); otherwise every $Q_w=Q_0$. (2) If $\mathsf p(x)\in\mathcal C$, it is the
unique vertex of $x'$ in $\mathcal C$; otherwise $x$, hence $x'$, has no vertex in $\mathcal C$, and the first vertex of $x$ lies in $x'$, so it is the first vertex of $x'$.
(3) $\mathcal C$ is $\Ga$-stable and the order of (IO) is $\Ga$-invariant. For $\SL_2(\Z)$: the stabilisers are $\langle S\rangle\cong\Z/4$ at the orbit of $i$, $\langle ST\rangle\cong\Z/6$
at the orbit of $\zeta^{2}$ and $\{\pm I\}$ at every other vertex ($m$ included), and $\{\pm I\}$ acts trivially (\ref{H4}(iv)); (a) holds by (IO), each triangle meeting each orbit at
most once (Theorem~\ref{thm:mpcomparison}(5)). At $\ell\ge5$, $Q_x=\Ga_x$, so every vertex with stabiliser $\neq Q_0$ must lie in $\mathcal C$, and $Q_0\subseteq\{\pm I\}$; on $\mathsf M_7$ the
triangle $v_1v_4v_6$ then has two vertices in $\mathcal C$, on $\mathsf R_{\Z}$ no triangle has two non-central vertices (Remark~\ref{rmk:mprulesread}(2)).
\end{proof}

\begin{proposition}[The $Q$-reduced nerve: a choice-free projective resolution]\label{prop:mpqreduced}
\PROVED{} Let $\Ga$, $Y$, $\mathfrak T$ satisfy the hypotheses of Theorem~\ref{thm:mpcomparison} with finite abelian vertex stabilisers, let $\ell$ be a prime with (QD)
(Definition~\ref{def:mpqd}), and let $\cring$ be a commutative ring with $\abs{Q_v}\in\cring^{\times}$ for every vertex $v$. For $t=((g_i,\sigma_i))_{i\le k}\in\mathfrak F'_k$
(Definition~\ref{def:mpfree}) put $\tau_t:=\bigcap_ig_i\iota(\sigma_i)$, a closed simplex by (SC), and let $t\sim t'$ iff $t'=((q_ig_i,\sigma_i))_i$ with $q_i\in Q_{\mathsf p(\tau_t)}$
($k+1$ independent elements).
\begin{enumerate}
\item $\sim$ is an equivalence relation, compatible with the faces, the degeneracies and the $\Ga$-action; $\mathfrak D:=\mathfrak F'/{\sim}$ is a simplicial $\Ga$-set, the stabiliser of
  the class $[t]$ is $Q_{\mathsf p(\tau_t)}$, and $\mathsf p(\tau_{t'})=\mathsf p(\tau_t)$ for $t'\sim t$. In the normal form $g_0=1$ of Theorem~\ref{thm:mpcomparison}(1),
  $\mathsf S^{(\ell)}_k:=\Ga\backslash\mathfrak D_k$ is the quotient of $\mathsf S_k$ by $(\sigma_i;\gamma_i)\mapsto(\sigma_i;c_i\gamma_i)$, $c_i\in Q_{\mathsf p_{i-1}}$ independent,
  $\mathsf p_{i-1}:=(\gamma_1\cdots\gamma_{i-1})^{-1}\mathsf p(\tau_s)$; a class is degenerate iff some representative has $\sigma_i=\sigma_{i+1}$ and $\gamma_{i+1}\in Q_{\mathsf p_i}$.
\item $\cring[\mathfrak D_k]\cong\bigoplus_{[s]\in\mathsf S^{(\ell)}_k}\cring[\Ga/Q_{\mathsf p(\tau_s)}]$ is a projective $\cring[\Ga]$-module, and $\cring[\mathfrak D_\bullet]\to\cring$
  and its normalised quotient $N\cring[\mathfrak D_\bullet]\to\cring$ are resolutions of $\cring$ by projective $\cring[\Ga]$-modules, the latter projective on the non-degenerate classes.
\item $\Hom_{\cring[\Ga]}(N\cring[\mathfrak D_\bullet],M)$ is, on the nose, the complex $C^{\bullet}(\mathsf S^{(\ell)};M^{Q})$ of functions $f$ on the non-degenerate classes with
  $f(s)\in M^{Q_{\mathsf p(\tau_s)}}$ ($\mathsf p(\tau_s)$ a vertex of $\iota(\sigma_0)$) and the differential of Definition~\ref{def:mpset},
  $(df)(s)=\gamma_1\cdot f(d_0s)+\sum_{i\ge1}(-1)^{i}f(d_is)$. Hence, for every $\cring[\Ga]$-module $M$, naturally in $M$,
  $H^{k}(C^{\bullet}(\mathsf S^{(\ell)};M^{Q}))\cong H^{k}(\Ga;M)$ for every $k\ge0$; likewise for the unnormalised complex.
\item The projection $\mathfrak F'\to\mathfrak D$ is a map of projective resolutions over $\id_\cring$, so it induces the identity of $H^{\bullet}(\Ga;M)$ under (3) and
  Theorem~\ref{thm:mpcomparison}(4). The $Q$-reduced nerve, the full MP complex and the hybrid of Remark~\ref{rmk:mphybrid} are compared \emph{on cohomology} by the comparison
  theorem \STD{\cite[Thm.~2.2.6]{weibel}}, never on cochains.
\item For $\Ga=\SL_2(\Z)$ on $\mathsf M_7$ and on $\mathsf R_{\Z}$, at $\ell\in\{2,3\}$, the hypotheses hold for every $\Z_{(\ell)}$-algebra $\cring$
  ($\abs{Q_v}\in\{1,2,4\}$ at $\ell=3$, $\in\{1,3\}$ at $\ell=2$). No adaptedness and no choice of lifts is used.
\end{enumerate}
\end{proposition}

\begin{proof}
(1) Let $t'=((q_ig_i,\sigma_i))$ with $q_i\in Q_{\mathsf p}$, $\mathsf p:=\mathsf p(\tau_t)$. Each $q_i$ fixes $\mathsf p$, so $\mathsf p\in\tau_{t'}$. If $\mathsf p\in\mathcal C$, it is the vertex of
$\tau_{t'}$ in $\mathcal C$, so $\mathsf p(\tau_{t'})=\mathsf p$; otherwise $Q_{\mathsf p}=Q_0$ acts trivially on $Y$ and $\tau_{t'}=\tau_t$. Hence the relation is symmetric and transitive (the
group $Q_{\mathsf p}^{k+1}$ acting). A face $d_j$ enlarges $\tau$: $\mathsf p(\tau_t)$ is a vertex of $\tau_{d_jt}$, so $Q_{\mathsf p(\tau_t)}\subseteq Q_{\mathsf p(\tau_{d_jt})}$ by
Definition~\ref{def:mpqd}(1), and $t\sim t'$ implies $d_jt\sim d_jt'$; degeneracies keep $\tau$. For $\gamma\in\Ga$, $\gamma(q_ig_i)=(\gamma q_i\gamma^{-1})(\gamma g_i)$ with
$\gamma Q_{\mathsf p}\gamma^{-1}=Q_{\gamma\mathsf p}=Q_{\mathsf p(\tau_{\gamma t})}$ (Definition~\ref{def:mpqd}(3)). As $\Ga$ acts freely on the first label of $\mathfrak F'$, $\gamma t\sim t$
iff $\gamma g_i=q_ig_i$ for all $i$, i.e.\ iff $\gamma\in Q_{\mathsf p(\tau_t)}$. In normal form, $t'$ re-normalised by $q_0^{-1}$ has $\gamma'_i=g_{i-1}^{-1}(q_{i-1}^{-1}q_i)g_{i-1}\,\gamma_i$, and
$g_{i-1}^{-1}Q_{\mathsf p}g_{i-1}=Q_{\mathsf p_{i-1}}$ ($Q_{\mathsf p}$ abelian, the $q_{i-1}^{-1}q_i$ independent). A class is degenerate iff it is $[s_iu]$, i.e.\ iff some
representative repeats a pair up to $Q_{\mathsf p}$, which in normal form is the stated condition.

(2) The $\Ga$-orbit of $[t]$ is $\Ga/Q_{\mathsf p(\tau_t)}$ by (1), and $\cring[\Ga/Q]=\cring[\Ga]e_Q$, $e_Q:=\abs{Q}^{-1}\sum_{q\in Q}q$, is a direct summand of $\cring[\Ga]$. The
degenerate classes form a $\Ga$-stable subset, so $N\cring[\mathfrak D_k]$ is the permutation module on the non-degenerate classes, projective likewise, and the projection
$\cring[\mathfrak D_\bullet]\to N\cring[\mathfrak D_\bullet]$ is a quasi-isomorphism \STD{\cite[\S8.3]{weibel}, as in the proof of Theorem~\ref{thm:mpcomparison}(2), \VERIFYloc}.
Exactness of $\cring[\mathfrak D_\bullet]\to\cring$ is the argument of Theorem~\ref{thm:mpcomparison}(2) with one quotient in each lemma. Let
$\mathfrak F^{Q}_{k,s}:=\mathfrak F_{k,s}/{\equiv}$, where $((g_i,\sigma_i);x)\equiv((q_ig_i,\sigma_i);x)$ for $q_i\in Q_x$ --- the relation by the \emph{simplex} $x$, not by the dominant
point. It is well defined ($Q_x\subseteq\Ga_x$ fixes $x$, so $x\subseteq q_ig_i\iota(\sigma_i)$); the $k$-faces and degeneracies keep $x$, the $s$-faces $x\mapsto\partial_jx$ coarsen it
($Q_x\subseteq Q_{\partial_jx}$), and $\Ga$ acts as in (1); so $\mathfrak F^{Q}$ is a bisimplicial $\Ga$-set (semi-simplicial in $s$) with
$\mathfrak F^{Q}_{k,s}=\bigsqcup_{x\in Y_s}(O(x)/Q_x)^{k+1}$, $Q_x$ acting on $O(x)$ by $(g,\sigma)\mapsto(qg,\sigma)$.
\emph{Rows} (fixed $s$): $O(x)/Q_x\ne\emptyset$ by (PD), and $\cring[(O(x)/Q_x)^{\bullet+1}]\to\cring$ is contracted by the extra degeneracy with a fixed class, as in
Lemma~\ref{lem:mprows}; so $\mathrm{Tot}\,\cring[\mathfrak F^{Q}]\to\cring[Y_\bullet]$ is a quasi-isomorphism by Lemma~\ref{lem:mpassembly}, and $\cring[Y_\bullet]\to\cring$ is one by (IO)
and (AC) (over $\Z$, hence over $\cring$: a bounded-below exact complex of free abelian groups is split).
\emph{Columns} (fixed $k$): the augmentation $\mathfrak F^{Q}_{k,s}\to\mathfrak D_k$ (in $s=0$), $((g_i,\sigma_i);x)\mapsto[((g_i,\sigma_i))]$, is well defined, since for a vertex $w$ of
$x\subseteq\tau_t$, $Q_x\subseteq Q_w\subseteq Q_{\mathsf p(\tau_t)}$ (Definition~\ref{def:mpqd}(1)), and it commutes with the $k$-faces. The column splits over $[t]\in\mathfrak D_k$; the
summand over $[t]$ is the augmented ordered chain complex of the $\Delta$-set $Z_{[t]}:=\{(t';x):t'\sim t,\ x\subseteq\tau_{t'}\}/{\equiv}$. Put $\mathsf p:=\mathsf p(\tau_t)$, which is
$\mathsf p(\tau_{t'})$ for every $t'\sim t$ by (1). Define $h([t]):=[(t;\mathsf p)]$ and $h[(t';x)]:=(-1)^{j}[(t';\mathsf p\cup x)]$ if $\mathsf p\notin x$, $j$ the position of $\mathsf p$ in the
ordered simplex $\mathsf p\cup x$, and $0$ if $\mathsf p\in x$. This is well defined: $\mathsf p\cup x\subseteq\tau_{t'}$ is a simplex, its gluing group is
$Q_{\mathsf p\cup x}=Q_{\mathsf p}\cap Q_x=Q_x$ (as $Q_x\subseteq Q_{\mathsf p}$), and the classes $[(t';\mathsf p)]$, $t'\sim t$, coincide because $t'$ and $t$ differ by elements of
$Q_{\mathsf p}$. As $h$ changes only the simplex coordinate, $dh+hd=\id$ is the identity of the cone on a vertex of an ordered simplex, checked term by term (the faces of
$\mathsf p\cup x$ other than $x$ pair off with those of $h(\partial x)$, with signs $(-1)^{j}(-1)^{i}$ and $(-1)^{i}(-1)^{j-1}$ for $i<j$, $(-1)^{j}(-1)^{i+1}$ and $(-1)^{i}(-1)^{j}$
for $i\ge j$), including the augmentation degree; this is Lemma~\ref{lem:mpcols} with the cone point $\mathsf p$ in place of the first vertex. So every column is acyclic in
positive degrees with $H_0=\cring[\mathfrak D_k]$, and $\mathrm{Tot}\,\cring[\mathfrak F^{Q}]\to\cring[\mathfrak D_\bullet]$ is a quasi-isomorphism by Lemma~\ref{lem:mpassembly}. The two
composites $\mathrm{Tot}\,\cring[\mathfrak F^{Q}]\to\cring$ agree (both send $((g,\sigma);v)\in\mathfrak F^{Q}_{0,0}$ to $1$), so $\cring[\mathfrak D_\bullet]\to\cring$ is a
quasi-isomorphism. (Adaptedness was not used; (QD) entered twice: $\mathsf p(\tau_{t'})=\mathsf p(\tau_t)$ on a class, and $Q_x\subseteq Q_{\mathsf p(\tau_t)}$.)

(3) $\Hom_{\cring[\Ga]}(\cring[\Ga/Q_{\mathsf p}]\cdot[x_s],M)=M^{Q_{\mathsf p}}$ by $F\mapsto F([x_s])$, $x_s$ a normal form of $s$; two normal forms of one class differ by an element
of $Q_{\mathsf p}$, which fixes $F([x_s])$, so $f(s):=F([x_s])$ is well defined. As in Theorem~\ref{thm:mpcomparison}(3), $d_ix_s=x_{d_is}$ for $i\ge1$ and
$d_0x_s=\gamma_1\cdot x_{d_0s}$, which gives the differential; $\mathsf p(\tau_s)\in\tau_s\subseteq\iota(\sigma_0)$. By (2) the cohomology is $\Ext^{k}_{\cring[\Ga]}(\cring,M)$ \STD{\cite[Thm.~2.2.6]{weibel}},
which the free resolution $\cring\otimes\Z[\mathfrak F'_\bullet]$ identifies with $H^{k}(\Ga;M)$ (Theorem~\ref{thm:mpcomparison}(4)), naturally in $M$. Without the normalisation the same
computation gives the unnormalised statement.
(4) $\cring\otimes\Z[\mathfrak F'_\bullet]\to\cring[\mathfrak D_\bullet]$ is $\Ga$-equivariant, commutes with faces and degeneracies (by (1)) and with the augmentations; two maps of
projective resolutions over $\id_\cring$ are homotopic, so the induced map on $\Ext$ is the canonical identification.
(5) Definition~\ref{def:mpqd} and Theorem~\ref{thm:mpcomparison}(5); the orders of the $Q_v$ are read off its $\SL_2(\Z)$ clause.
\end{proof}

\noindent\COMP\textsuperscript{a} (not a registered validation): \file{docs/implementation/slot22/fmptri/fmptri\_spotcheck.py} (the F-MPT referee's spot-check, check D) builds
the choice-free $\mathfrak D$ on $\mathsf M_7$ and on $\mathsf R_{\Z}$ and computes $\dim H^{0},H^{1},H^{2}$ of $C^{\bullet}(\mathsf S^{(\ell)};M^{Q})$, $M=\Sym^{n}\F_\ell^{2}$: equal to the tree oracle
of \file{mpcore.py} on every cell, $\ell=3$, $n\le4$ and $\ell=2$, $n\le3$, on both triangulations, with $d\circ d=0$, the $d_0$-twist representative-independent on all of degree $1$ and
every value in its fibre asserted; $\abs{\mathsf S^{(\ell),\mathrm{nd}}_k}$, $k\le3$, is $5,52,510,5386$ ($\ell=3$) and $5,63,621,5847$ ($\ell=2$) on $\mathsf M_7$, $7,70,588,5656$ and
$7,101,1007,9053$ on $\mathsf R_{\Z}$. On $\mathsf R_{\Z}$ these classes are the quotient form of [MPT] (\file{docs/latex/mp\_triangulations.tex} \S2.5, in the $\mu$-coordinates
of its Proposition \texttt{prop:muform}), whose tables agree. Re-run of 2026-10-10: \code{FMPTRI\_SPOTCHECK failing=0 verdict=OK}.

\subsection{All degrees}\label{sec:mpallsec}

\begin{definition}[The comparison lift $\Phi^{\mathsf{MP}}$]\label{def:mpphi}
Let $\Ga=\SL_2(\Z)$, $\mathfrak T\in\{\mathsf M_7,\mathsf R_{\Z}\}$, and $\Iso_\bullet$ the isotropy-resolution model of Definition~\ref{def:isomodel} on $\Xbulk$ with the
representatives of $\mathsf M_7$ (the target for both triangulations). $\Phi^{\mathsf{MP}}\colon N\Z[\mathfrak F'_\bullet]\to\Iso_\bullet$ is the $\Z[\Ga]$-linear map defined on the
free generators $e_s$ ($s\in\mathsf S_k^{\mathrm{nd}}$, the normal form $x_s$) by
\[
  \Phi_0(e_s):=h\cdot e_{0,v,0},\qquad \Phi_k(e_s):=c\bigl(\Phi_{k-1}(d\,e_s)\bigr)\quad(k\ge1),
\]
where $h\cdot v$ is the first vertex of $\iota(\sigma_0)$ lying in the orbit of a representative $v\in\Sigma_0$ of $\mathsf M_7$ ($h$ any element with that property; the
construction depends on the choices of the $h$ and of $c$, and Proposition~\ref{prop:mpalldegrees}(1) shows that the homotopy class does not; the code's choice is recorded by
\code{PhiLift}), $e_{0,v,0}$ the generator $1\otimes1$ of
$\Z[\Ga]\otimes_{\Z[\Ga_v]}F^{v}_0\subseteq G_{0,0}$, $d\,e_s=\gamma_1\cdot e_{d_0s}+\sum_{i\ge1}(-1)^{i}e_{d_is}$, and $c$ a $\Z$-linear contracting homotopy of the augmented
$\Iso_\bullet\to\Z$ (it exists: Proposition~\ref{prop:isomodel}(2), split over $\Z$). The code builds it to degree~$3$ (\code{orbifold\_mp::PhiLift}, $c=$ \code{PeriodicityClass::fill}).
\end{definition}

\begin{proposition}[All degrees from a finite triangulation: what holds, and its exact scope]\label{prop:mpalldegrees}
Let $\Ga=\SL_2(\Z)$, $\mathfrak T\in\{\mathsf M_7,\mathsf R_{\Z}\}$, $\cring$ a commutative ring and $V$ an $\cring[\Ga]$-module.
\begin{enumerate}
\item \PROVED{} $\Phi^{\mathsf{MP}}$ is a $\Z[\Ga]$-chain map over the augmentations, unique up to homotopy, and $(\Phi^{\mathsf{MP}})^{\ast}\colon\CI(V)\to C^{\bullet}(\mathsf S;\underline V)$ is a
  quasi-isomorphism in \textbf{every} degree over every $\cring$; both sides are $H^{\bullet}(\Ga;V)$ (Theorem~\ref{thm:mpcomparison}, Proposition~\ref{prop:isomodel}(3)).
\item \PROVED{} \textbf{($D=3$, the degree bound of the notation paragraph.)} For $q\ge2=\mathrm{vcd}(\Ga)+1$, cup product with the class $e\in H^{2}(\Ga;\Z)$ of Lemma~\ref{lem:bulkperiodic}(1) is an isomorphism
  $H^{q}(\Ga;V)\to H^{q+2}(\Ga;V)$ (Lemma~\ref{lem:bulkperiodic}(2)). Hence $H^{\le3}(\mathsf S;\underline V)$ determines every degree: $H^{2j}\cong H^{2}$, $H^{2j+1}\cong H^{3}$ for
  $j\ge1$, transported through (1).
\item \PROVED{} \textbf{(The torsion bound; why $\Z/2^{k}$, $k\le3$, and $\F_3$ suffice in degrees $\ge2$.)} $24\cdot H^{q}(\Ga;\VR)=0$ for $q\ge2$, every $\cring$ and $n$
  (Claim~\ref{cl:amalgam}(4)), and for $\cring=\Z$ there is no $\ell$-torsion for $\ell\ge5$ in degrees $\ge2$ (Proposition~\ref{prop:badprimes}). So for $q\ge2$ and
  $k\ge v_\ell(24)$ ($k\ge3$ at $\ell=2$, $k\ge1$ at $\ell=3$), $H^{q}(\Ga;\Sym^{n}(\Z/\ell^{k})^{2})\cong H^{q}(\Ga;\Vcoef)_{(\ell)}\oplus H^{q+1}(\Ga;\Vcoef)_{(\ell)}$ (the
  universal coefficient sequence of the finite free complex $\CI(\Z)$, both terms finite of exponent dividing $\ell^{v_\ell(24)}$), independently of $k$: the cells $k\le3$ at
  $\ell=2$ and $k=1$ at $\ell=3$ are complete in degrees $\ge2$, not a choice.
\item \COMP\textsuperscript{r} \textbf{(What is certified.)} $\Phi^{\mathsf{MP}}$ is built to degree $3$ on both triangulations ($55\,816$ generators on $\mathsf M_7$, $59\,448$ on
  $\mathsf R_{\Z}$; $20$ distinct group elements, coefficients $\le2$ in absolute value) and certified exactly: $d\Phi=\Phi d$ over $\Z[\Ga]$ and $\aug\circ\Phi_0=\aug$, a
  \emph{coefficient-free} certificate serving every $n$ and every ring. On $\CI$, the cup product with $e$ is certified (\code{PeriodicityClass} on $q_{\max}=7$) and is an
  isomorphism $H^{q}\to H^{q+2}$, $q=2,\dots,5$, over $\Fl$, $\ell\in\{2,3,5,7\}$, $n=0,\dots,12$ ($52$ rows). $(\Phi^{\mathsf{MP}})^{\ast}$ is measured an isomorphism on $H^{\le1}$ for
  $n\le12$, on $H^{2}$ for $n\le3$, and on $H^{3}$ for $\ell\in\{2,3\}$, $n\in\{0,2\}$. The certificate line
  \code{MP\_ALL\_DEGREES model=<M7|R\_Z> ring=F\_$\ell$ nerve\_to=2 periodic\_on=C\_I from=2 via=thm:mpcomparison certified=4/4 scope=bulk verdict=OK} records these, and
  \code{certified=4/4} counts $n=0,\dots,3$.
\item \textbf{(Scope.)} The statement is about $\Ga=\SL_2(\Z)$ and the \textbf{bulk} only: not the boundary (where $2$-periodicity starts in degree $1$ and is a separate statement,
  Remark~\ref{rmk:bulkperiodicscope}(3)(a)), not the pair, and not a finite-index subgroup $B\le\Ga$ (which would need its own periodicity clause). On the MP complex the degrees
  above $\texttt{nerve\_to}$ are \emph{transported} by (1) and (2), not computed.
\end{enumerate}
\end{proposition}

\begin{proof}
(1) The construction is well posed: $\aug(\Phi_0(e_s))=1$, and by induction $\Phi_{k-1}(d\,e_s)$ is a cycle of $\Iso_{k-1}$ (augmentation $1-1=0$ at $k=1$), so $c$ fills it; extended
$\Ga$-linearly from the free generators it is a chain map over $\id_{\Z}$. Both $N\Z[\mathfrak F']$ (Theorem~\ref{thm:mpcomparison}(2)) and $\Iso$ (Proposition~\ref{prop:isomodel}(2)) are free
resolutions of $\Z$, so a chain map over $\id_\Z$ is a homotopy equivalence, unique up to homotopy \STD{\cite[Thm.~2.2.6]{weibel}}, and $\Hom_{\Z[\Ga]}(-,V)$ of it is a
quasi-isomorphism. (2) is Lemma~\ref{lem:bulkperiodic}(2), which holds for every $\Z[\Ga]$-module $V$, read through (1). (3) $\CI(\Z)$ is a complex of free abelian groups of finite rank
(Proposition~\ref{prop:isomodel}(2),(3)) with $\CI(\Z)\otimes\Z/\ell^{k}=\CI(\Z/\ell^{k})$, so the universal coefficient sequence $0\to H^{q}(\Z)\otimes\Z/\ell^{k}\to H^{q}(\Z/\ell^{k})\to
\mathrm{Tor}(H^{q+1}(\Z),\Z/\ell^{k})\to0$ holds and splits. For $q\ge2$ both $H^{q}(\Z)$ and $H^{q+1}(\Z)$ are finite and killed by $24$; for a finite abelian group $G$ with
$\ell^{a}G_{(\ell)}=0$ and $k\ge a$, $G\otimes\Z/\ell^{k}=G_{(\ell)}=\mathrm{Tor}(G,\Z/\ell^{k})$. (4) is the record of the validations named in Remark~\ref{rmk:mpcheckable}. (5) is a
statement of scope.
\end{proof}

\noindent\emph{Which part is load-bearing} (the F-MP2 audit's finding~6). The all-degrees statement (1)--(2) rests on Theorem~\ref{thm:mpcomparison}, the existence of
$\Phi^{\mathsf{MP}}$ and Lemma~\ref{lem:bulkperiodic}(2); the certificate of (4) shows that the \emph{built} $\Phi^{\mathsf{MP}}$ is such a lift. The measured isomorphisms on
$H^{\le3}$ are consistency checks of the implementation ($\rhon$, layouts, the length machinery) against the frozen oracle; they are not what makes the statement true.

\subsection{Morphisms, naturality, and the maps that are needed}\label{sec:mpmorphsec}

\begin{definition}[Morphisms of MP data; the label factor]\label{def:mpmorphism}
Let $(\Ga,Y,\mathfrak T)$ and $(\Ga',Y',\mathfrak T')$ satisfy the hypotheses of Theorem~\ref{thm:mpcomparison}.
\begin{enumerate}
\item \textbf{(Simplicial morphism.)} A morphism $(\varphi,f)$ is a homomorphism $\varphi\colon\Ga\to\Ga'$ and a $\varphi$-equivariant simplicial map $f\colon Y\to Y'$, recorded on the
  descriptors by a \textbf{label table} $\sigma\mapsto(\mathsf h_\sigma,\sigma')$, $\sigma\in\Sigma_{\max}$, $\mathsf h_\sigma\in\Ga'$, $\sigma'\in\Sigma'_{\max}$, with
  $f(\iota(\sigma))\subseteq\mathsf h_\sigma\,\iota'(\sigma')$, which is all that the induced maps of (3) use.
\item \textbf{(Refinement.)} For a $\Ga$-equivariant subdivision $\tilde Y$ of $Y$ (not a simplicial map $\tilde Y\to Y$): $\varphi=\id$ and a label table
  $\tilde\sigma\mapsto(\mathsf h_{\tilde\sigma},\sigma)$ with $\tilde\iota(\tilde\sigma)\subseteq\mathsf h_{\tilde\sigma}\,\iota(\sigma)$ --- the geometric simplex is unique, the label
  only up to right multiplication by $\Ga_\sigma$ ($=\{\pm I\}$ for a triangle of $\Xbulk$ under (IO), the setwise stabiliser being the pointwise one).
  The table need not come from a simplicial map; it is certified equally by a $\Ga$-equivariant map of vertex sets sending the vertices of $\tilde\iota(\tilde\sigma)$ into
  $\mathsf h_{\tilde\sigma}\,\iota(\sigma)$ (a common vertex then goes to a common vertex), which is how the Lean marker below certifies $\mathsf R_{\Z}\to\mathsf M_7$
  (v1.22.0-pre.7).
  \LEAN{def:mp-morph} \LEAN{def:mp-orbmap} (for (1) and (2): a morphism as a label table carrying non-empty intersections to non-empty ones, over any group homomorphism
  $\varphi$, given in particular by an equivariant vertex map with the containments; on descriptors, from decided orbit and label tables; the simplicial map $f$ and the
  geometric containment in $\abs Y$ are not formalised, only what the induced maps use)
\item \textbf{(The induced maps.)} In both cases
  \[
    u\colon\mathfrak F'_k\to\mathfrak F'{}'_k,\qquad ((g_i,\sigma_i))\mapsto((\varphi(g_i)\,\mathsf h_{\sigma_i},\sigma'_i)),
  \]
  and on the quotients $u\colon(\sigma_i;\gamma_i)\mapsto(\sigma'_i;\ \mathsf h_{\sigma_{i-1}}^{-1}\varphi(\gamma_i)\,\mathsf h_{\sigma_i})$, a map of simplicial sets
  $\mathsf S(Y)\to\mathsf S(Y')$, covered by the morphism of local systems $\mathsf h_{\sigma}\colon\underline{M'}_{u(\sigma)}=M'\to M'=\underline{\varphi^{\ast}M'}_{\sigma}$. The
  \textbf{induced cochain map} is
  \[
    (u^{\ast}f)(s):=\mathsf h_{\sigma_0}\cdot f(u(s))\qquad(s=(\sigma_i;\gamma_i)\in\mathsf S_k^{\mathrm{nd}};\ f:=0\text{ on degenerate images}),
  \]
  $C^{\bullet}(\mathsf S(Y');\underline{M'})\to C^{\bullet}(\mathsf S(Y);\underline{\varphi^{\ast}M'})$. The factor $\mathsf h_{\sigma_0}$ --- the \textbf{label factor} --- is part of the map for
  every morphism: refinement, the forget-coset map of Proposition~\ref{prop:mprescor}, the boundary inclusion of Proposition~\ref{prop:mpboundary} and the $g$-twist alike (for the $g$-twist the label factor is $\rho_n(g)$ on a $\conj_g$-twisted
  chain map of one $\mathfrak F'$, Remark~\ref{rmk:mphecke}, not a morphism of MP data; v1.22.0-pre.8).
  The same formula on all strings $s\in\mathsf S_k$ defines $u^{\ast}\colon C^{\bullet}_{\mathrm{un}}(\mathsf S(Y');\underline{M'})\to C^{\bullet}_{\mathrm{un}}(\mathsf S(Y);\underline{\varphi^{\ast}M'})$,
  of which the map above is the restriction ($u$ commutes with the degeneracies; v1.22.0-pre.7).
  \LEAN{lem:mp-morph-simplicial} \LEAN{def:mp-pull} ($u$ on $\mathfrak F'$, equivariant and commuting with the faces, $u(x_s)=\mathsf h_{\sigma_0}\cdot x_{u(s)}$, the $\gamma$-coordinates of
  $u(s)$; $u^{\ast}$ with the label factor on $C^{\bullet}_{\mathrm{un}}$, equal to $F\mapsto F\circ u$ on the equivariant cochains; not the normalised restriction)
\end{enumerate}
\end{definition}

\begin{theorem}[Naturality and refinement invariance]\label{thm:mpnatural}
\PROVED{} In the setting of Definition~\ref{def:mpmorphism}:
\begin{enumerate}
\item $u$ is a $\varphi$-equivariant map of simplicial sets commuting with the augmentations to $\Z$; $u^{\ast}$ is the map $F\mapsto F\circ\Z[u]$ of
  $\Hom_{\Z[\Ga']}(N\Z[\mathfrak F'{}'],M')\to\Hom_{\Z[\Ga]}(N\Z[\mathfrak F'],\varphi^{\ast}M')$ read on the normal forms, hence a cochain map; the assignment is compatible with
  composition when labels compose by $\mathsf h''_\sigma=\psi(\mathsf h_\sigma)\,\mathsf h'_{\sigma'}$. The same holds for the unnormalised complexes: $u^{\ast}$ of
  Definition~\ref{def:mpmorphism}(3) on $C^{\bullet}_{\mathrm{un}}$ is $F\mapsto F\circ\Z[u]$ on $\Hom_{\Z[\Ga']}(\Z[\mathfrak F'{}'],M')$, a cochain map preserving the normalised
  subcomplexes (v1.22.0-pre.7).
  \LEAN{thm:mp-pull-cochain} \LEAN{lem:mp-pull-comp} \LEAN{def:mp-zu} (the unnormalised clause: $\Z[u]$, over every commutative ring, a $\Ga$-chain map over the augmentations, $u^{\ast}$
  precomposition with it, a cochain map for every representation, and composition; the normalised $u^{\ast}$ is not stated in Lean)
\item The isomorphism of Theorem~\ref{thm:mpcomparison}(4) is natural: under it, $H^{\bullet}(u^{\ast})$ is the restriction $\varphi^{\ast}\colon H^{\bullet}(\Ga';M')\to H^{\bullet}(\Ga;\varphi^{\ast}M')$;
  likewise for $C^{\bullet}_{\mathrm{un}}$, and the inclusions $C^{\bullet}\subseteq C^{\bullet}_{\mathrm{un}}$ are quasi-isomorphisms commuting with $u^{\ast}$ and inducing the identity of
  $H^{\bullet}(\Ga;M)$, resp.\ $H^{\bullet}(\Ga';M')$ (v1.22.0-pre.7).
  \LEAN{thm:mp-pull-natural} \LEAN{def:mp-barmap} \LEAN{thm:mp-homotopy-exact} (for $C^{\bullet}_{\mathrm{un}}$, every group homomorphism $\varphi$, every commutative ring, every
  representation $M'$, every degree, with (AC) over the ring, (IO), finiteness and $\dim\le2$ for both data: under the isomorphisms of Theorem~\ref{thm:mpcomparison}(4)'s
  naturality marker, $H^{\bullet}(u^{\ast})$ is \code{mathlib}'s \texttt{groupCohomology.map}\,$\varphi\,(\mathrm{id})$; the transfer to $C^{\bullet}$ is by the normalised markers of
  Definition~\ref{def:mpset} and Theorem~\ref{thm:mpcomparison}(4) on cohomology, not stated for the normalised $u^{\ast}$)
\item \textbf{(Refinement invariance.)} For a refinement, $u^{\ast}\colon C^{\bullet}(\mathsf S(Y);\underline M)\to C^{\bullet}(\mathsf S(\tilde Y);\underline M)$ is a quasi-isomorphism in every
  degree, inducing the identity of $H^{\bullet}(\Ga;M)$, and independent of the label choices on cohomology (different labels give chain-homotopic maps); likewise on
  $C^{\bullet}_{\mathrm{un}}$.
  \LEAN{thm:mp-refine-invariance} \LEAN{thm:mp-label-indep} (for $C^{\bullet}_{\mathrm{un}}$ and $\varphi=\id$: a quasi-isomorphism in every degree, the identity under the comparison
  isomorphisms; label independence \emph{on cohomology} only --- the chain homotopy between the two label choices is not formalised)
\item \textbf{(The label factor cannot be dropped.)} $f\mapsto f\circ u$ is a cochain map when all labels act on $M'$ by one and the same central automorphism (all by $+1$,
  all by $-1$, or all in $\ker\rho'$, $\rho'$ the action of $\Ga'$ on $M'$), but not in general: with mixed labels $\pm I$ it fails already at odd $n$, and with a constant
  non-central label $h$ it would need $\gamma$ and $h^{-1}\gamma h$ to act alike on $M'$ for every arrow $\gamma$ (v1.22.0-pre.6, referee F-MP2 F1: pre.5's sufficient
  condition ``every label acts centrally, e.g.\ by $\pm1$'' was false). On $\mathsf R_{\Z}\to\mathsf M_7$ every label is $\pm I$,
  the factor is $(\pm1)^{n}$ per simplex, and it is visible for mixed labels at odd $n$.
  \LEAN{thm:mp-untwisted} (the defect formula $d(f\circ u)(s)-(df)(u(s))=(\varphi(\gamma_1)-\mathsf h_{\sigma_0}^{-1}\varphi(\gamma_1)\mathsf h_{\sigma_1})\cdot f(u(d_0s))$, the sufficient condition
  of one central automorphism, and the failure as soon as one string has a vector on which the two act differently) \LEAN{thm:mp-rz-untwisted} (on $\mathsf R_{\Z}\to\mathsf M_7$: labels
  all $I$ give a cochain map for every representation; mixed labels $I$, $-I$ do not, at $n=1$ over $\F_3$)
\end{enumerate}
\end{theorem}

\begin{proof}
(1) $f(g\iota(\sigma))=\varphi(g)f(\iota(\sigma))\subseteq\varphi(g)\mathsf h_\sigma\iota'(\sigma')$ (for a refinement, $g\tilde\iota(\tilde\sigma)\subseteq g\mathsf h_{\tilde\sigma}\iota(\sigma)$ as subsets of
$\abs{\tilde Y}=\abs Y$), so a non-empty intersection maps into a non-empty one --- and this is all that the rest of the proof uses of $f$, resp.\ of the subdivision: any
label table carrying non-empty intersections to non-empty ones gives the same $u$ and $u^{\ast}$ (v1.22.0-pre.7); dropping and repeating pairs commute with $u$; $u(\gamma\cdot x)=\varphi(\gamma)\cdot u(x)$. On the normal
form, $u(x_s)=((\mathsf h_{\sigma_0},\sigma'_0),(\varphi(\gamma_1)\mathsf h_{\sigma_1},\sigma'_1),\dots)=\mathsf h_{\sigma_0}\cdot x_{u(s)}$, since
$\mathsf h_{\sigma_0}^{-1}\varphi(\gamma_1\gamma_2)\mathsf h_{\sigma_2}=(\mathsf h_{\sigma_0}^{-1}\varphi(\gamma_1)\mathsf h_{\sigma_1})(\mathsf h_{\sigma_1}^{-1}\varphi(\gamma_2)\mathsf h_{\sigma_2})$; so
$(F\circ u)(x_s)=\mathsf h_{\sigma_0}\cdot F(x_{u(s)})$, which is the displayed formula (Theorem~\ref{thm:mpcomparison}(3)); it is a cochain map because $\Z[u]$ is a chain map. The local
systems are compatible with the arrows: $\mathsf h_{\sigma_0}\,(\mathsf h_{\sigma_0}^{-1}\varphi(\gamma_1)\mathsf h_{\sigma_1})=\varphi(\gamma_1)\,\mathsf h_{\sigma_1}$. Composition: substitute. Without the normalised quotient the same
computation gives the unnormalised clause; $u$ commutes with the degeneracies, so $\Z[u]$ descends to $N\Z[\mathfrak F']$ and $u^{\ast}$ preserves the functions vanishing on degenerate
strings.
(2) $N\Z[\mathfrak F']\to\Z$ is a free $\Z[\Ga]$-resolution and $\varphi^{\ast}N\Z[\mathfrak F'{}']\to\Z$ is acyclic, so $\Z[u]$ is \emph{the} lift of $\id_\Z$ up to homotopy
\STD{\cite[Thm.~2.2.6]{weibel}}, and $F\mapsto F\circ\Z[u]$ computes the restriction along $\varphi$ on $\Ext$ \STD{\cite[Ch.~III.8]{brown}, the map $(\varphi,\text{id})$ of pairs
(group, module); chapter-level pointer, \VERIFYloc}. The same argument on the free resolutions $\Z[\mathfrak F']$, $\varphi^{\ast}\Z[\mathfrak F'{}']$ gives the unnormalised statement;
the projection $\pi\colon\Z[\mathfrak F']\to N\Z[\mathfrak F']$ is a map of free resolutions over $\id_\Z$, hence a homotopy equivalence (the same theorem), so the inclusion
$C^{\bullet}\subseteq C^{\bullet}_{\mathrm{un}}$, which is $\Hom(\pi,M)$, is one inducing the identity of $H^{\bullet}(\Ga;M)$, and $\pi\circ\Z[u]=N\Z[u]\circ\pi$ makes it commute with $u^{\ast}$.
(3) With $\varphi=\id$, (2) gives $\id^{\ast}=\id$ on $H^{\bullet}(\Ga;M)$, both sides being $H^{\bullet}(\Ga;M)$ by Theorem~\ref{thm:mpcomparison}. Two label tables give two $\Ga$-maps
$N\Z[\mathfrak F'_{\tilde Y}]\to N\Z[\mathfrak F'_Y]$ over $\id_\Z$, chain-homotopic by the same comparison theorem; likewise $\Z[\mathfrak F'_{\tilde Y}]\to\Z[\mathfrak F'_Y]$ for $C^{\bullet}_{\mathrm{un}}$.
(4) Without the factor, $d(f\circ u)(s)-(df)(u(s))=\gamma_1\cdot f(u(d_0s))-(\mathsf h_{\sigma_0}^{-1}\gamma_1\mathsf h_{\sigma_1})\cdot f(d_0u(s))$, which vanishes for all $f$ iff
$\gamma_1$ and $\mathsf h_{\sigma_0}^{-1}\gamma_1\mathsf h_{\sigma_1}$ act alike on $M'$ --- for labels $\pm I$ iff $\rho'(\mathsf h_{\sigma_0})=\rho'(\mathsf h_{\sigma_1})$ on every arrow,
i.e.\ ($Y$ being connected by (AC)) iff the labels are constant or $n$ is even; with the factor, $\mathsf h_{\sigma_0}\mathsf h_{\sigma_0}^{-1}\gamma_1\mathsf h_{\sigma_1}=\gamma_1\mathsf h_{\sigma_1}$, which
is the factor of $d_0s$. For labels $\pm I$, $\rhon(\pm I)=(\pm1)^{n}$.
\end{proof}

\noindent\COMP\textsuperscript{r} for (3), (4) on $\mathsf R_{\Z}\to\mathsf M_7$ (\code{OrbifoldMPRefinementValidation}): the seven labels (each $\mathsf R_{\Z}$ triangle in one $\mathsf M_7$
triangle; the label coset $\{h,h\cdot(-I)\}$ of each simplex, both members passing containment; $h=I$ by the subdivision record); every string of $\mathsf S_{\le3}(\mathsf R_{\Z})$ maps into $\mathsf S(\mathsf M_7)$ ($1044$ non-degenerate strings have a
degenerate image); $u^{\ast}$ a cochain map on $27$ products over $\F_2$, $\F_3$, $\Z/4$; with $h\cdot(-I)$ on one simplex only, at $\F_3$, $n=1$, the untwisted $f\circ u$ is \emph{not}
a cochain map (\code{REFINE\_TWIST}); $u^{\ast}$ an isomorphism on $H^{\le2}$ at $\ell\in\{2,3,5\}$, $n\in\{0,2,4\}$, label-independent; a wrong label $h\cdot T$ fails containment
and sends $7146$ strings outside $\mathsf S(\mathsf M_7)$.
\LEAN{def:mp-rz-refine} (the seven labels, all $\mathsf h=I$, with the orbit table of the collapse of the new vertices onto vertices of $\mathsf M_7$ recorded by \texttt{RZTau.cTab}, and the
coset member $\mathsf h\cdot(-I)$ on one triangle; kernel-decided) \LEAN{thm:mp-rz-refine} ($u^{\ast}$ a quasi-isomorphism on $C^{\bullet}_{\mathrm{un}}$ for every commutative ring and every
representation, in every degree, the identity under the comparisons; both label choices alike on cohomology) \LEAN{thm:mp-rz-untwisted} (the untwisted failure at $\F_3$, $n=1$, on one string)

\begin{proposition}[The comparison with the shipped Bredon complex]\label{prop:mpbredon}
Let $\Ga=\SL_2(\Z)$, $\mathfrak T\in\{\mathsf M_7,\mathsf R_{\Z}\}$ and put $\varepsilon^{\mathsf{MP}}:=(\Phi^{\mathsf{MP}})^{\ast}\circ\cmpI^{\ast}\colon\Ceq(\cring)\to C^{\bullet}(\mathsf S;\underline{\VR})$, with
$\cmpI^{\ast}$ the map of Proposition~\ref{prop:isomodel}(4) and $\Ceq(\cring)=\Hom_{\Z[\Ga]}(C_\bullet(\Xbulk),\VR)$ (Theorem~\ref{thm:explicit_complex_equiv}).
\begin{enumerate}
\item \PROVED{} $\varepsilon^{\mathsf{MP}}$ is a cochain map, compatible with base change, and its homotopy class is canonical: any two $\Z[\Ga]$-maps $N\Z[\mathfrak F']\to C_\bullet(\Xbulk)$ over
  $\id_\Z$ are homotopic; in particular it is, up to homotopy, the map induced by the augmentation $\mathrm{Tot}\to C_\bullet(\Xbulk)$ composed with an inverse of
  $\mathrm{Tot}\to\Z[\mathfrak F']$. Likewise, for \emph{any} $\Ga$-map $\Phi\colon\cring[\mathfrak F'_\bullet]\to C_\bullet(\Xbulk)\otimes\cring$ over $\id_\cring$, the map
  $\Phi^{\ast}\colon\Ceq(\cring)\to C^{\bullet}_{\mathrm{un}}(\mathsf S;\underline{\VR})$ is homotopic to $\varepsilon^{\mathsf{MP}}$ followed by the inclusion $C^{\bullet}\subseteq C^{\bullet}_{\mathrm{un}}$
  (v1.22.0-pre.7).
  \LEAN{def:mp-epsmp} \LEAN{thm:mp-epsmp-canonical} ($\varepsilon^{\mathsf{MP}}$ built as $\Phi^{\ast}$, $\Phi$ the lift of $\id_\cring$ from $\cring[\mathfrak F'_\bullet]$ into $C_\bullet(\Xbulk)\otimes\cring$,
  into the unnormalised MP complex, a cochain map at every ring; any two such lifts homotopic, so the class on cohomology is canonical; $\Phi^{\mathsf{MP}}$ and $\cmpI$ themselves and
  the compatibility with base change are not formalised)
\item \PROVED{} It is a quasi-isomorphism \textbf{if} $6\in\cring^{\times}$ (the cone of $\cmpI^{\ast}$ is killed by $48$, Proposition~\ref{prop:isomodel}(4)): this is
  Theorem~\ref{thm:specialisation} re-proved on the MP model. \textbf{Not only if}: for $\cring=\Z[1/2]$ and $n$ odd both sides are acyclic (Remark~\ref{rmk:negI}, second bullet, as
  corrected at v1.22.0-pre.4; $H^{\bullet}(\Ga;\VR)=0$ since $2\in\cring^{\times}$), while $3\notin\cring^{\times}$.
  \LEAN{thm:mp-epsmp-quasiiso} (the ``if'' only, with the stabiliser orders dividing $12$ of \ref{H4}(iv): $\Phi$ a homotopy equivalence and $\Phi^{\ast}$ a quasi-isomorphism for every
  representation, the Bredon side on $\mathsf M_7$ and the MP side on $\mathsf M_7$ or $\mathsf R_{\Z}$; not the ``not only if'' witness)
\item \PROVED{} At every $\cring$, $\varepsilon^{\mathsf{MP}}$ is injective on $H^{0}$ and $H^{1}$: the Bredon groups are the row $t=0$ of the isotropy spectral sequence of $\CI$, and the
  edge maps $E_2^{p,0}\to H^{p}$ are injective for $p\le1$ (no differential reaches $E^{p,0}$ from a column $p-r<0$), as in the proof of Proposition~\ref{prop:ellthreeclosed}(3).
  So at $\ell\in\{2,3\}$ the Bredon defect lives in the cokernel, never in the kernel, in degrees $\le1$, and in degree $\ge2$ the Bredon side is $0$.
\item \COMP\textsuperscript{r} (\code{OrbifoldMPComparisonValidation}, from the shipped Bredon pipeline \code{EquivariantCohomologyOnK}): quasi-isomorphism at $\ell\in\{5,7\}$,
  $n\le12$ (non-trivially on $H^{0},H^{1}$; at $q=2$ both sides are $0$); at $(\ell,n)=(2,0),(2,2),(3,0),(3,2)$, $q=0,1,2$: rank of $\varepsilon^{\mathsf{MP}}$ equal to the Bredon dimension,
  defects $[0,1,1]$, $[0,1,2]$, $[0,1,1]$, $[0,0,0]$; the ``if, not iff'' witness at $\ell=3$, $n\in\{1,3,5\}$.
\end{enumerate}
\end{proposition}

\begin{proof}
(1) Composite of cochain maps; base change from Proposition~\ref{prop:isomodel}(4) and the definition of $\Phi^{\mathsf{MP}}$ over $\Z$. $\cmpI\circ\Phi^{\mathsf{MP}}$ is a $\Ga$-map from a free
complex to the acyclic augmented $C_\bullet(\Xbulk)\to\Z$ over $\id_\Z$, unique up to homotopy \STD{\cite[Thm.~2.2.6]{weibel}}. The same theorem, applied over $\cring$ to the
projective resolution $\cring[\mathfrak F'_\bullet]$ (Theorem~\ref{thm:mpcomparison}(2)) and the exact $C_\bullet(\Xbulk)\otimes\cring\to\cring$ (\ref{H4}(ii)), makes any $\Phi$ over $\id_\cring$
homotopic to $(\cmpI\circ\Phi^{\mathsf{MP}}\circ\pi)\otimes\cring$, $\pi\colon\Z[\mathfrak F']\to N\Z[\mathfrak F']$ the projection; dualised, this is the second sentence of (1), and with the
inclusion $C^{\bullet}\subseteq C^{\bullet}_{\mathrm{un}}$ a quasi-isomorphism (Theorem~\ref{thm:mpnatural}(2)), clause (2) below holds for $\Phi^{\ast}$ as well. (2) $(\Phi^{\mathsf{MP}})^{\ast}$ is always a
quasi-isomorphism (Proposition~\ref{prop:mpalldegrees}(1)) and $\cmpI^{\ast}$ is one for $6\in\cring^{\times}$. The witness: the Bredon complex vanishes by \eqref{eq:negI} with $2$
invertible, and $2H^{\bullet}(\Ga;\VR)=0$ by centrality. (3) Filter $\CI(\VR)$ by $s$ as in the proof of Proposition~\ref{prop:isomodel}(4): $E_1^{s,t}=\bigoplus_{\sigma\in\Sigma_s}H^{t}(\Ga_\sigma;\VR)$,
row $t=0$ is $\Ceq$, $\cmpI^{\ast}$ is the inclusion of that row, so on cohomology it is the edge map $E_2^{p,0}\to E_\infty^{p,0}=F^{p}H^{p}\subseteq H^{p}$; $E_2^{p,0}\to E_\infty^{p,0}$
is the quotient by the images of $d_r\colon E_r^{p-r,r-1}\to E_r^{p,0}$, which vanish for $p\le1<r$; then compose with the isomorphism $(\Phi^{\mathsf{MP}})^{\ast}$. Bredon $H^{\ge2}=0$:
Proposition~\ref{prop:bredon}(2).
\end{proof}

\begin{proposition}[The boundary and the pair]\label{prop:mpboundary}
\PROVED{} \COMP\textsuperscript{r} Let $\Ga=\SL_2(\Z)$ and $\Xbdry$ the horocycle subcomplex (\ref{B1}--\ref{B3}), $\Gbdry=\langle T\rangle\times\langle-I\rangle$, with the top simplices of the boundary descriptor
the three edges $v_0v_1$, $v_1v_2$, $v_2v_3$.
\begin{enumerate}
\item (AC) fails for $\Xbdry$ ($H_0=\Z[\Ga/\Gbdry]$), and the proof of Theorem~\ref{thm:mpcomparison}, parts (1)--(2), gives a free resolution $N\Z[\mathfrak F'(\Xbdry)]\to\Z[\Ga/\Gbdry]$ of the
  \emph{permutation module}; hence $H^{\bullet}(\mathsf S(\Xbdry);\underline M)\cong\Ext^{\bullet}_{\Z[\Ga]}(\Z[\Ga/\Gbdry],M)\cong H^{\bullet}(\Gbdry;M)=H^{\bullet}(\partial;M)$ of
  Definition~\ref{def:bdrygroup}, in every degree, by Shapiro's lemma.
\item The inclusion is the labelled morphism $j\colon\mathfrak F'(\Xbdry)\to\mathfrak F'(\Xbulk)$, a boundary edge to the unique triangle of $\Xbulk$ containing it, with its label
  (Definition~\ref{def:mpmorphism}); $j^{\ast}$ carries the label factor, its cone computes the relative object of Definition~\ref{def:borelmodel} (the MP-model form of
  Definition~\ref{def:pair}) --- in the relative indexing $\mathrm{Cone}^{q}:=C^{q}(\mathsf S(\Xbulk);\underline M)\oplus C^{q-1}(\mathsf S(\Xbdry);\underline M)$,
  $d(c,x):=(dc,\ j^{\ast}c-dx)$, so that $H^{q}(\mathrm{Cone})=H^{q}(\Xbulk,\Xbdry;M)$; $(c,x)\mapsto(c,-x)$ identifies it with $\Cone{\rres}[-1]$ in the convention of
  Definition~\ref{def:pair} (the shift negating the differential), so $H^{q}(\mathrm{Cone})=H^{q-1}(\Cone{\rres})=\Hcpt{q}$ (v1.22.0-pre.7, the indexing of the code) ---, and the long
  exact sequence of the pair holds
  unconditionally, at every ring, in every degree; under Theorem~\ref{thm:mpcomparison} and (1) $j^{\ast}$ is the restriction $H^{\bullet}(\Ga;M)\to H^{\bullet}(\Gbdry;M)$.
\end{enumerate}
\end{proposition}

\begin{proof}
(1) Lemmas~\ref{lem:mprows}, \ref{lem:mpcols}, \ref{lem:mpassembly} hold for $\Xbdry$ verbatim ((SC), (IO), (PD) are inherited by the subcomplex), so $\Z[\mathfrak F'(\Xbdry)]\simeq C_\bullet(\Xbdry)$,
and $C_\bullet(\Xbdry)=\Z[\Ga]\otimes_{\Z[\Gbdry]}C_\bullet(\partial\Xbulk_{\infty})$ with the distinguished component contractible, so $C_\bullet(\Xbdry)\to H_0=\Z[\Ga/\Gbdry]$ is a
quasi-isomorphism; the source is free. Shapiro: \STD{\cite[Ch.~III.6.2]{brown}}, as in Definition~\ref{def:borelmodel}(2) and Proposition~\ref{prop:circlemodel}(1). (2)
Theorem~\ref{thm:mpnatural}(1) gives the cochain map; $\Z[j]$ is a $\Ga$-map from the free resolution of (1) to the free resolution of $\Z$ over the augmentation
$\varepsilon\colon\Z[\Ga/\Gbdry]\to\Z$, unique up to homotopy \STD{\cite[Thm.~2.2.6]{weibel}}, so $H^{\bullet}(j^{\ast})$ is
$\Ext^{\bullet}_{\Z[\Ga]}(\varepsilon,M)\colon\Ext^{\bullet}_{\Z[\Ga]}(\Z,M)\to\Ext^{\bullet}_{\Z[\Ga]}(\Z[\Ga/\Gbdry],M)$, which under Shapiro is $\res^{\Ga}_{\Gbdry}$
\STD{\cite[Ch.~III.9]{brown}, restriction as the map induced by $\Z[\Ga/H]\to\Z$; \VERIFYloc}; the cone and its sequence as in
Lemma~\ref{lem:conerelative}(1),(2), which use of the restriction only that it is a cochain map.
\end{proof}

\noindent\COMP\textsuperscript{r} for (1), (2) (\code{OrbifoldMPBoundaryValidation}, \code{standard}, $332$ checks; slot-22 MP-3, audited F-MP3, v1.22.0-pre.7; since MP-3b,
v1.22.0-pre.8, also the digest of the oracle file actually read, gated against the frozen one):
$H^{0},\dots,H^{3}(\mathsf S(\Xbdry);\Sym^{n}\cring)=H^{\bullet}(\Gbdry;\Sym^{n}\cring)$ on $134$ cells ($\F_2,\F_3,\F_5,\F_7$, $n\le12$; $\Z/4,\Z/8,\Z/9$ by lengths, $n\le4$; the boundary
descriptors of $\mathsf M_7$ and $\mathsf R_{\Z}$) against an explicit resolution of $\Gbdry$ and the shipped circle model of Proposition~\ref{prop:circlemodel}; $j^{\ast}$ a cochain map, its rank on $H^{\le3}$ that of the restriction
of $\CI$; the cone exact, with $H^{\bullet}(\Xbulk,\Xbdry;\F_2)=[0,0,1,1]$ for $q=0,\dots,3$ at $n=0$ on $\mathsf S_{\le4}(\mathsf M_7)$, in particular $H^{1}(\Xbulk,\Xbdry;\F_2)=0$.

\begin{proposition}[Restriction, corestriction and the sequence of a pair of groups]\label{prop:mprescor}
\PROVED{} \COMP\textsuperscript{r} Let $B\le A\le\Ga=\SL_2(\Z)$ have finite index, acting on the same $\Xbulk$ with the triangulation $\mathfrak T$, and $\mathfrak F'_A$, $\mathfrak F'_B$ the labelled Čech nerves built
from the $A$-, resp.\ $B$-orbit representatives.
\begin{enumerate}
\item \textbf{(Shapiro on the nose, on one $\mathfrak F'$.)} $\mathfrak F'_A$ is a free $B$-set, and $\Hom_{\Z[B]}(N\Z[\mathfrak F'_A],M)=\Hom_{\Z[A]}(N\Z[\mathfrak F'_A],\mathrm{Coind}_B^{A}M)$ is
  an identity of cochain complexes; on it $\res^{A}_{B}$ is the restriction of equivariance and $\cor^{A}_{B}$ the transfer $F\mapsto\sum_t r_t^{-1}\cdot F(r_t\,\cdot)$ over a
  transversal $r_t$ of $B\backslash A$, both cochain maps, with $\cor\circ\res=[A:B]$. The same holds with $\mathfrak F'_A$ replaced by any free $\Ga$-set resolving $\Z$ as in Theorem~\ref{thm:mpcomparison}(2) (free over every subgroup), in
  particular by $\mathfrak F'_{\Ga}$, on which $\Hom_{\Z[B]}(N\Z[\mathfrak F'_{\Ga}],M)=\Hom_{\Z[\Ga]}(N\Z[\mathfrak F'_{\Ga}],\mathrm{Coind}_B^{\Ga}M)$,
  $\res^{A}_{B}$ is induced by the $\Ga$-map $\mathrm{Coind}_A^{\Ga}M\to\mathrm{Coind}_B^{\Ga}M$ (the inclusion of $A$-equivariant into $B$-equivariant functions $\Ga\to M$) and $\cor^{A}_{B}$
  by $\phi\mapsto(\gamma\mapsto\sum_t r_t^{-1}\phi(r_t\gamma))$; either nerve serves, the two realisations agreeing on $\Ext$ by transitivity of coinduction and the independence of
  $\res$, $\cor$ from the projective resolution (v1.22.0-pre.7; the code works on $\mathfrak F'_{\Ga}$).
\item \textbf{(Forget the coset.)} $\mathsf S(\mathfrak F'_B)=B\backslash\mathfrak F'_B$ and $A\backslash\mathfrak F'_A$ are different simplicial sets; the forget-coset map is a labelled morphism
  $\mathfrak F'_B\to\mathfrak F'_A$ of Definition~\ref{def:mpmorphism} (labels the coset representatives, \emph{not} central, so the label factor is load-bearing), a comparison of two
  free $B$-resolutions, unique up to homotopy; under Theorem~\ref{thm:mpcomparison} it induces $\res^{A}_{B}$.
\item \textbf{(The pair of groups.)} The cone of $\res$ gives the long exact sequence of the \emph{pair of groups} $(A,B)$ in every degree and at every $\ell$ --- the group-cohomology
  form of \S\ref{sec:nebles}'s level cone, which is the model's (Bredon) sequence at $\ell\in\{2,3\}$ (Proposition~\ref{prop:levelconeshapiro}(3)). It is \emph{not}
  inflation--restriction (Proposition~\ref{prop:levelinfres}), a different statement with different terms.
\end{enumerate}
\end{proposition}

\begin{proof}
(1) $B$ acts freely on the first label; $\mathrm{Coind}_B^{A}=\Ind_B^{A}$ for finite index, and the identity is the adjunction, termwise and compatible with the differentials;
$\res$, $\cor$ on $\Ext$ are computed on any projective resolution by these formulas \STD{\cite[Ch.~III.9--10]{brown}}. (2) Theorem~\ref{thm:mpnatural}(1),(2) with $\varphi$ the
inclusion. (3) The cone of a cochain map has a long exact sequence (Lemma~\ref{lem:conerelative}(2)'s argument); with (1), (2) its terms are $H^{\bullet}(A;M)$, $H^{\bullet}(B;M)$
and the relative groups of the pair. For $\mathfrak F'_{\Ga}$ in (1): it is free over $A$ and $B$ (already on the first label), so the same adjunction gives
$\Hom_{\Z[B]}(N\Z[\mathfrak F'_{\Ga}],M)=\Hom_{\Z[\Ga]}(N\Z[\mathfrak F'_{\Ga}],\mathrm{Coind}_B^{\Ga}M)$ with $\mathrm{Coind}_A^{\Ga}\mathrm{Coind}_B^{A}=\mathrm{Coind}_B^{\Ga}$, and the
displayed maps are the transported $\res$, $\cor$ \STD{\cite[Ch.~III.9--10]{brown}}.
\end{proof}

\noindent\COMP\textsuperscript{r} for (1)--(3) (\code{OrbifoldMPLevelConeValidation}, \code{standard}, $723$ checks; slot-22 MP-3, audited F-MP3, v1.22.0-pre.7; MP-3b, audited
F-MP3b, v1.22.0-pre.8): the five pairs
$\Ga_0(4)<\Ga_0(2)$, $\Ga_0(9)<\Ga_0(3)$, $\Ga_1(4)<\Ga_0(4)$, $\Ga^{2}<\Ga$, $\Ga^{3}<\Ga$ over $\Fl$, $\ell\in\{2,3\}$, $n\le2$, on $\mathsf M_7$: the $B$-descriptor with
$\abs{\mathsf S_k^{\mathrm{nd}}(B)}=[\Ga:B]\,\abs{\mathsf S_k^{\mathrm{nd}}(\mathsf M_7)}$ when $-I\in B$; $\res$, $\cor$ on $\mathfrak F'_{\Ga}$ with $\cor\circ\res=[A:B]$ on the cochains; the
forget-coset map with its label factor (non-central labels, e.g.\ $10$ of $15$ for $\Ga_0(2)\to\Ga$; for $\Ga_1(4)\to\Ga_0(4)$ all labels are $+I$) equal to Shapiro $\circ\res$ on
$H^{\le2}$, and dropping the factor breaks the cochain map; $H^{q}(A,B;M)$ from the cone equal to the long exact sequence of $0\to\mathrm{Coind}_AM\to\mathrm{Coind}_BM\to Q\to0$ on $21$ cells
($q\le2$; $q=3$ for $\Ga^{2},\Ga^{3}<\Ga$ at $n=0$), e.g.\ $\Ga^{2}<\Ga$, $\F_2$, $n=0$: $\res=0$ on $H^{1}$, $H^{\bullet}(A,B)=[0,1,1,1]$. Since MP-3b (v1.22.0-pre.8): the pair sequence over $\Z/\ell^{k}$ --- $\Z/4$
($\Ga_0(4)<\Ga_0(2)$, $n\le1$ on $\mathsf M_7$, $n=1$ on $\mathsf R_{\Z}$), $\Z/9$ ($\Ga_0(9)<\Ga_0(3)$) and $\Z/25$ ($\Ga_0(25)<\Ga_0(5)$, a step of the Hida $5$-tower), $n=0$ unless
stated: the lengths of $H^{q}(A)$, $H^{q}(B)$, of the image of $\res$ and of $H^{q}(A,B)$ (MP $q\le2$, $\CI$ $q\le3$) against an integral Mayer--Vietoris oracle (Smith normal form,
universal coefficients), with torsion visible (over $\Z/4$ the lengths of $H^{\bullet}(\Ga_0(2))$ are $[2,4,3,3]\ne2\cdot[1,2,2,2]$); the pairs $\Ga^{2}<\Ga$ ($\F_2$) and
$\Ga_0(4)<\Ga_0(2)$ ($\F_3$, $\Z/4$) on $\mathsf R_{\Z}$, with their $B$-descriptors; and $\Ga_1(8)<\Ga_1(4)$, $-I\notin A$ ($\F_2$; $\F_3$, $n=1$; MP $q\le1$, $\CI$ $q\le3$):
labels with $u=-I$ occur only when $-I\notin A$, and do occur, $\Ga_1(8)\to\Ga_1(4)$, $60$ of $120$ (representative-dependent); replacing them by their $u=I$ candidates moves
strings out of $\mathsf S_A$ --- a negative control that is structurally guaranteed ($-h\notin A$ is automatic since $-I\notin A$), which shows that the containment gate detects a
$u=I$-only labelling. \textbf{Not gated at this version}: $q=3$ on the nerve for the $\Ga_0$ pairs, $\Z/\ell^{k}$ at $n\ge2$, $\Z/49$ ($\Ga_0(49)$ has $280$ top orbits on
$\mathsf M_7$, beyond the nerve's bound), inflation--restriction.

\begin{remark}[Hecke operators on the MP model]\label{rmk:mphecke}
(Not a statement.) For $g\in\De$ the operator $\Tdc{g}=\cor\circ(\conj_g)_{\ast}\circ\res$ needs a $g$-twisted chain datum
(Definition~\ref{def:tau}) on a projective resolution. On the MP model this is built \emph{on the MP resolution itself}: $N\Z[\mathfrak F'_\bullet]$
is free over $\Z[\Ga'']$, $\Ga''=\Ga\cap g\Ga g^{-1}$, on $\{s_j e_s\}$ ($s_j$ the representatives of $\Ga''\backslash\Ga$), and acyclic, so a
$\conj_g$-twisted chain map $\tau$ lifting $\id_\Z$ exists and is unique up to $\conj_g$-twisted homotopy \emph{over every commutative ring}
(the comparison theorem; unlike Lemma~\ref{lem:tauexists}(2) no hypothesis $6\in\cring^{\times}$ is needed, because the source is free).
It is built generator by generator with the MP model's own $\Z$-linear filling of $N\Z[\mathfrak F'_\bullet]$ (the double complex of the
proof of Theorem~\ref{thm:mpcomparison} made explicit), $\tau_0(s_je_s)=e_s$, $\tau_k(s_je_s)=\mathrm{fill}(\tau_{k-1}(\partial\,s_je_s))$;
no datum of $C_\bullet(\Xbulk)$ enters. The three maps are, in the Shapiro coordinates of Proposition~\ref{prop:mprescor}(1), $\res$ the
inclusion of invariants, $(\conj_g)_\ast\colon F\mapsto\rho_n(g)\,F\circ\tau$ --- the factor $\rho_n(g)$ is the label factor of the
$g$-twist, it makes $F'$ $\Ga''$-equivariant and cannot be dropped (Theorem~\ref{thm:mpnatural}(4)) --- and $\cor$ the transfer;
fused, $(Tf)(e_s)=\sum_j\sum c\,\rho_n(s_j^{-1}g\gamma)f(e_{s'})$ over the terms $c\gamma e_{s'}$ of $\tau(s_je_s)$, the form of
\eqref{eq:cosettransport}. Since $\Phi\circ\tau^{\mathrm{MP}}$ and $\tau_{\mathrm{Iso}}\circ\Phi$ are two such twisted maps into the acyclic
$\mathrm{Iso}_\bullet$, $T_p\Phi^{\ast}=\Phi^{\ast}\hat\Theta_g$ on $H^{\bullet}$ over every ring. The $g$-twist is \emph{not} a labelled
morphism of MP data (Definition~\ref{def:mpmorphism}): $\Xbulk$ and $g\Xbulk$ are different truncations of $\Hyp$, and a \emph{canonical}
geometric datum through a common truncation is gap G3, off the critical path --- this remark shows it is not needed to compute.
\COMP\textsuperscript{r} \code{OrbifoldMPHeckeValidation} (slot 22, MP-4): $\tau$ certified exactly over $\Z[\Ga]$ on $\mathsf M_7$ ($p\in\{2,3,5,7\}$,
degrees $\le 2$) and $\mathsf R_{\Z}$ ($p\in\{2,3\}$, degrees $\le1$); $T_p$ a cochain map; $T_p\Phi^{\ast}=\Phi^{\ast}\hat\Theta_g$ on
$H^{\le1}$ over $\F_\ell$, $\ell\in\{2,3,5,7\}$, $n\le12$, on $H^{\le2}$ over $\Z/4,\Z/8,\Z/9,\Z/27$ ($n\in\{0,2\}$) and over $\F_2,\F_3$
($n\le3$); the characteristic polynomials on $H^0,H^1(\Ga;\Sym^n\F_\ell)$ equal an independent amalgam-cocycle oracle, with
$(x-\tau(p))^2(x-1-p^{11})$ at $n=10$, $\ell\in\{5,7\}$; the Bredon operator of Section~\ref{sec:chainlevel} is recovered on $\operatorname{im}\varepsilon^{\ast}$.
Not built: $H^3$, $T_{p^\nu}$, level $N$, the $\iota$-sign.
Notation (v1.22.0-pre.8; the text above is the audit F-MP4 \S3's verbatim, except that its \COMP{} is written \COMP\textsuperscript{r} and its undefined ``Section~sec:hecke'' is \S\ref{sec:chainlevel}): $g=\mathrm{diag}(1,p)$, so $\Ga''=\Ga_0(p)$ and $T_p=T=\Tdc{g}$;
$\Phi=\Phi^{\mathsf{MP}}$ of Definition~\ref{def:mpphi}; $\mathrm{Iso}_\bullet=\Iso_\bullet$, $\tau_{\mathrm{Iso}}$ and $\hat\Theta_g=\Th_{\Iso,g}$ as in Proposition~\ref{prop:strictlift}(1),(3);
$\tau(p)$ in the last formula is Ramanujan's $\tau$, not the chain map.
\end{remark}

\noindent\COMP\textsuperscript{r} for Remark~\ref{rmk:mphecke} (\code{OrbifoldMPHeckeValidation}, \code{standard}, $2059$ checks; slot-22 MP-4, audited F-MP4 (ACCEPT WITH
REPAIRS); predictions \file{MP4\_predictions.md} and the independent oracle \file{mp/mp4/mp4\_oracle.py} committed before the first run; v1.22.0-pre.8). The scope of the record
in the remark, with the \emph{non-vacuous} counts (cells on which the operator acts on a non-zero space) printed and gated: $T_p\Phi^{\ast}=\Phi^{\ast}\hat\Theta_g$ on $354$ of
$592$ cells ($240$ of the $448$ cells on $\mathsf M_7$ over $\Fl$; $114$ of the $144$ cells over $\Z/\ell^{k}$ and on $\mathsf R_{\Z}$); the characteristic polynomials against the
oracle on $254$ of $464$ cells ($220$ of $416$ on $\mathsf M_7$, $34$ of $48$ on $\mathsf R_{\Z}$); the shipped Bredon $T_p$ on $66$ of $108$ cells ($15$ of $24$ over $\F_2,\F_3$,
on $\operatorname{im}\varepsilon^{\ast}$, $n\in\{0,2\}$; $51$ of $84$ over $\F_5,\F_7$, $\varepsilon^{\ast}$ an isomorphism, $n$ even $\le12$, $p\ne\ell$). The cochain-map identity
$T\,d^{q}=d^{q}\,T$ is gated for $q\le1$ on $\mathsf M_7$ and $q=0$ on $\mathsf R_{\Z}$. The uniqueness over every ring is what makes the cells over $\Z/4,\Z/8,\Z/9,\Z/27$ and
$\F_2,\F_3$ meaningful, where Lemma~\ref{lem:tauexists}(2) does not apply. In degree $2$ the datum stops at $\tau_2$: $d\tau_2=\tau_1d$ is certified, so $T$ preserves
$2$-coboundaries and the class of $T_p\Phi^{\ast}z$ is defined, but $T\,d^{2}=d^{2}\,T$ needs $\tau_3$ and is checked only on the tested classes ($T_p\Phi^{\ast}z$ a cocycle,
cohomologous to $\Phi^{\ast}\hat\Theta_gz$) --- all of $H^{2}$, $\Phi^{\ast}$ being onto (\ref{MPC4}); at $\ell\in\{2,3\}$, $H^{2}$ has no independent Hecke oracle (only
$\hat\Theta_g$ through $\Phi^{\ast}$). Negative controls: the label factor $\rho_n(g)$ dropped, or mismatched with the datum ($gT$ for $g$), breaks the cochain map and the
identity on $H^{1}$; a perturbed $\tau_1$ fails the certificate. \textbf{Not gated at this version}: $H^{3}$ ($\tau_3$, $d^{3}$); $T_{p^{\nu}}$, $T(p,p)$ and the full operators
\eqref{eq:fullassembly}; characteristic polynomials over $\Z/\ell^{k}$; $\mathsf R_{\Z}$ beyond $p\in\{2,3\}$, degrees $\le1$; a second filling (the independence of the filling
is argued, not tested); level $N$ ($T_p$, $U_p$ on the subgroup MP sets of Proposition~\ref{prop:mprescor}); compatibility with $j^{\ast}$ of Proposition~\ref{prop:mpboundary};
the $\iota$-sign.

\begin{example}[Moerdijk--Pronk's teardrop in the $G^{k}$-form]\label{ex:mpteardrop}
\COMP\textsuperscript{r} \emph{Tear-$n$} is the $2$-sphere with one cone point of order $n$, not a global quotient, so it tests the $G^{k}$-form of their $S(\mathcal T)$ (not
Definition~\ref{def:mpset}): six top simplices around and below the cone point, $\mathsf S_k=\bigsqcup C_n^{k}$ over the strings whose intersection contains the cone point, with
their $\mu$-maps on the edge $ac$ ([MP] \S6.1--6.2, read from the author source quoted in Remark~\ref{rmk:mpcaveat}). \code{OrbifoldMPTeardropValidation} (\code{fast}), against
\file{docs/implementation/slot22/mp/teardrop.py}: for $n\in\{2,3,4\}$, $\ell\in\{2,3,5\}$, $\dim H^{0},\dots,H^{3}=1,0,1,[\ell\mid n]$ --- Moerdijk--Pronk's values $A,0,A,H^{3}(C_n;A)$
--- and $\abs{\mathsf S_0},\dots,\abs{\mathsf S_4}=6,39,177,705,2877$; $6,48,294,1866,13290$; $6,57,465,4323,44925$; $d\circ d=0$ on every non-degenerate string to $k=4$; negative
controls ($\mu_{ac}=\mu_{ca}$, a second non-inverse $\mu$ pair, which does not model a $C_{n+1}$ vertex group, trivial $\mu$) fail (the first two break $d\circ d=0$; trivial $\mu$ keeps it and fails the dimension gate). ``Infinitely many degrees from one cone point'' in its simplest form; the $\ell=3$ structure of
Definition~\ref{def:mpell} is the disk with one cone point of order $3$ and a boundary circle, i.e.\ Tear-$3$ without its lower hemisphere ([MP] \S6.5).
\end{example}

\subsection{What is checked}\label{sec:mpchecksec}

\begin{remark}[What is checkable: \ref{MPC1}--\ref{MPC12}]\label{rmk:mpcheckable}
\PROVED{} \COMP\textsuperscript{r} Predictions committed before each first run (rule~20) in \file{docs/implementation/slot22/MP\_predictions.md} (P1--P4 and the MP-2 addendum).
Every validation below is registered in \file{Registry/Slot22/MP.cpp}, adds $0$ golden objects, tags its output \code{model=mp}, ends with
\code{ORBIFOLD\_MP <name> checks=<k> failing=0}, and where a Python oracle is reproduced prints \code{ORACLE <script> cells=<k> agree=<k>}. The oracles are independent of the code
under test: the frozen \file{audit/mv\_check.out} (an independent Mayer--Vietoris on $\Z/4\ast_{\Z/2}\Z/6$, own $\Sym^{n}$, own ranks), the shipped Bredon pipeline, and the shipped
$\CI$ (\file{Integral/Model/IsotropyModel}).
\begin{enumerate}[label=(MPC\arabic*)]
\item\label{MPC1} \PROVED{} (Definitions~\ref{def:mprules}, \ref{def:mpRZ}, Remark~\ref{rmk:mprulesread}) \textbf{Rules}: Table~\ref{tab:mprules}, (SC), (IO), the quotient counts, the
  subgroup-form and negative controls. \code{OrbifoldMPRulesValidation} (\code{fast}, $77$ checks).
\item\label{MPC2} \PROVED{} (Definition~\ref{def:mpset}, Remark~\ref{rmk:mpwhat}) \textbf{The set}: $\abs{\mathsf S_k^{\mathrm{nd}}}$, $k\le3$, the $\mu$-form counts, faces and simplicial
  identities, $d\circ d=0$ formally in $\Z[\Ga]\otimes\Z[\mathsf S_q^{\mathrm{nd}}]$ through degree $4$. \code{OrbifoldMPRulesValidation}, \code{OrbifoldMPNerveValidation}.
\item\label{MPC3} \PROVED{} (Theorem~\ref{thm:mpcomparison}) \textbf{The nerve against the oracle}: $H^{0},H^{1}$ over $\Fl$, $\ell\in\{2,3,5,7\}$, $n=0,\dots,12$; $H^{2}$, $\ell\in\{2,3,5\}$,
  $n\le3$ ($n\in\{4,6\}$ under the read-only hook \code{ORBIFOLD\_MP\_CELLS}); $\Z/\ell^{k}$ lengths, $\ell\in\{2,3\}$, $k\le3$, against $\CI$'s integral groups through the universal
  coefficient theorem; odd $n$ with $2$ a unit: non-zero acyclic complexes; on $\mathsf M_7$ and $\mathsf R_{\Z}$. \code{OrbifoldMPNerveValidation} (\code{standard}, $584$ checks);
  $H^{3}$ at $\ell\in\{2,3\}$, $n\in\{0,2\}$ in \code{OrbifoldMPNerveH3Slow} (\code{slow}). Negative controls: the untwisted $d_0$, the inverse twist $\rhon(\gamma_1^{-1})$, and a planted
  stabiliser $\Ga_{\zeta^{2}}=\langle(ST)^{2}\rangle$ without $-I$, which changes $H^{2}(\F_2)$ from $1$ to $2$ (the gated effect); that $H^{0}$, $H^{1}$ and the $\Z/\ell^{k}$ cells are unchanged is an ungated observation of the F-MP1 probe.
  \textbf{Hence the default $H^{2}$ cells are load-bearing for the central gerbe and never move behind a hook or into the slow tier; (AC), the one hypothesis of
  Theorem~\ref{thm:mpcomparison} not decided by a predicate, is guarded by this gate} (and, since v1.22.0-pre.6, proved in Lean for $\mathsf M_7$ and $\mathsf R_{\Z}$ over every
  ring: the markers of Theorem~\ref{thm:mpcomparison}(5)).
\item\label{MPC4} \PROVED{} (Definition~\ref{def:mpphi}, Proposition~\ref{prop:mpalldegrees}(1)) \textbf{$\Phi^{\mathsf{MP}}$}: certified to degree $3$; $(\Phi^{\mathsf{MP}})^{\ast}\delta_I=d_{\mathsf{MP}}(\Phi^{\mathsf{MP}})^{\ast}$
  exactly ($54$ products); an isomorphism $H^{q}(\CI)\to H^{q}(\mathsf S)$ as in Proposition~\ref{prop:mpalldegrees}(4), and over $\Z/\ell^{k}$, $\ell\in\{2,3\}$, $k\in\{2,3\}$, $n\in\{0,2\}$;
  a perturbed $\Phi_1$ fails the certificate. \code{OrbifoldMPComparisonValidation} (\code{standard}, $901$ checks).
\item\label{MPC5} \PROVED{} (Proposition~\ref{prop:mpbredon}) \textbf{$\varepsilon^{\mathsf{MP}}$}: Proposition~\ref{prop:mpbredon}(4). \code{OrbifoldMPComparisonValidation}.
\item\label{MPC6} \PROVED{} (Remark~\ref{rmk:mphybrid}) \textbf{The hybrid}: as stated there. \code{OrbifoldMPComparisonValidation}, \code{OrbifoldMPAllDegreesValidation}.
\item\label{MPC7} \PROVED{} (Theorem~\ref{thm:mpnatural}(3),(4)) \textbf{Refinement $\mathsf R_{\Z}\to\mathsf M_7$}: as recorded after the theorem. \code{OrbifoldMPRefinementValidation}
  (\code{standard}, $115$ checks).
\item\label{MPC8} \PROVED{} (Proposition~\ref{prop:mpalldegrees}(2)--(4)) \textbf{All degrees}: $\CI$ and $\Iso^{(\ell)}$ $H^{0},\dots,H^{7}$ over $\Fl$, $\ell\in\{2,3,5,7\}$, $n\le12$, against the oracle
  ($832$ values); the cup product with $e$ an isomorphism $H^{q}\to H^{q+2}$, $q=2,\dots,5$; the eight \code{MP\_ALL\_DEGREES} lines; negative control: $\ell\cdot E^{\ast}$ is not an
  isomorphism wherever $H^{2}\ne0$. \code{OrbifoldMPAllDegreesValidation} (\code{standard}, $1086$ checks).
\item\label{MPC9} \COMP\textsuperscript{r} (Example~\ref{ex:mpteardrop}) \textbf{The teardrop}. \code{OrbifoldMPTeardropValidation} (\code{fast}, $66$ checks).
\item\label{MPC10} \PROVED{} (Proposition~\ref{prop:mpboundary}; v1.22.0-pre.7) \textbf{The boundary and the pair}: as recorded after the proposition.
  \code{OrbifoldMPBoundaryValidation} (\code{standard}, $332$ checks; $331$ at v1.22.0-pre.7, MP-3b adds the oracle-file digest gate).
\item\label{MPC11} \PROVED{} (Proposition~\ref{prop:mprescor}; v1.22.0-pre.7) \textbf{Restriction, corestriction, the pair of groups}: as recorded after the proposition; the LES-0
  arm (iv) on the MP model; since MP-3b (v1.22.0-pre.8) also the pair sequence over $\Z/4$ ($\Ga_0(4)<\Ga_0(2)$, $n\le1$), $\Z/9$ ($\Ga_0(9)<\Ga_0(3)$) and $\Z/25$
  ($\Ga_0(25)<\Ga_0(5)$) against an integral Mayer--Vietoris oracle (Smith normal form, UCT), the pairs $\Ga^{2}<\Ga$, $\Ga_0(4)<\Ga_0(2)$ on $\mathsf R_{\Z}$, and
  $\Ga_1(8)<\Ga_1(4)$ with $-I\notin A$ (forget-coset labels with $u=-I$ reached and load-bearing). \code{OrbifoldMPLevelConeValidation} (\code{standard}, $723$ checks; $537$ at
  v1.22.0-pre.7).
\item\label{MPC12} \COMP\textsuperscript{r} (Remark~\ref{rmk:mphecke}; v1.22.0-pre.8) \textbf{Hecke operators on the MP model}: $T_p$, $p\in\{2,3,5,7\}$, on $H^{\le2}$ from the
  twisted chain map on the MP resolution, as recorded after the remark (non-vacuous counts included). \code{OrbifoldMPHeckeValidation} (\code{standard}, $2059$ checks).
\end{enumerate}
(v1.22.0-pre.7: \ref{MPC10}, \ref{MPC11} added --- slot-22 MP-3, audited by F-MP3 (ACCEPT WITH REPAIRS, the audit file \file{F-MP3code\_audit.md}), with the predictions
\file{MP3\_predictions.md}, both in \file{docs/implementation/slot22/}. v1.22.0-pre.8: \ref{MPC10}, \ref{MPC11} widened by MP-3b, audited F-MP3b, and \ref{MPC12} added ---
slot-22 MP-4, audited F-MP4, predictions \file{MP4\_predictions.md}; both audits ACCEPT WITH REPAIRS.) \textbf{Not checked at this version}: for \ref{MPC11}, $q=3$ on the nerve
for the $\Ga_0$ pairs and $\Z/\ell^{k}$ at $n\ge2$; for \ref{MPC12}, the list after Remark~\ref{rmk:mphecke} ($H^{3}$, $T_{p^{\nu}}$, characteristic polynomials over $\Z/\ell^{k}$,
level $N$, the $\iota$-sign).
\end{remark}

\paragraph{Scope: what this chapter does not do.} It does not re-prove Moerdijk--Pronk's sheaf-side theorem, and does not prove their $\mu$-independence (it avoids it); it does not
make any triangulation of this chapter, or the MP model, the shipped object or a default (the contract box above); it asserts the all-degrees statement for $\SL_2(\Z)$ in the bulk only
(Proposition~\ref{prop:mpalldegrees}(5)); it does not need goodness, a collar, or adaptedness for any theorem; it builds no \emph{canonical geometric} Hecke datum on the MP model (gap G3; the computational datum of Remark~\ref{rmk:mphecke} is a twisted chain map on the MP
resolution, built to degree $2$; v1.22.0-pre.8). The
non-developable case (complexes of groups that are not global quotients, the teardrop) is recorded by Example~\ref{ex:mpteardrop} only.

\paragraph{Code map.} \textbf{Built} (slot-22 MP-1a, MP-1b, MP-T, MP-2; audited: F-MP1, the MP-T audit, F-MP2 ACCEPT WITH REPAIRS, \file{docs/implementation/slot22/F-MP2\_code\_audit.md}),
in \file{src/hacs/Orbifold/MP/}, compiled into \code{cohomology\_app}: \file{Descriptor} (Definition~\ref{def:mptri}, the predicates of Definition~\ref{def:mprules},
read off \code{integral\_model::ModelData}; $\mathsf R_{\Z}$ by \code{SubdivideEdgeOrbit}), \file{MPSet} and \file{Nerve} (Definition~\ref{def:mpset}: the strings, faces, the
twisted normalised coboundaries over $\Z/m$), \file{SparseLength} (lengths over $\Z/\ell^{e}$), \file{AmalgamOracle} (the C++ port of \file{mv\_check.py}, and its frozen output as
\file{MvCheckOracle.inc}), \file{Comparison} (\code{PhiLift}, Definition~\ref{def:mpphi}), \file{Refinement} (Definition~\ref{def:mpmorphism}(2), \code{RefinementPullback} with the
label factor), \file{CochainMaps}, \file{Models}, \file{MP2Common}; the seven validations of Remark~\ref{rmk:mpcheckable}; registration \file{Registry/Slot22/MP.cpp},
\file{cmake/slot22/MP.cmake}. The scratch emitter \code{ORBIFOLD\_MP\_DUMP=<dir>} (no-op unless set, refused inside the tree) writes the $\mathsf R_{\Z}$ descriptor for the
Lean cross-checks. \textbf{Built at slot 22} (MP-3, audited F-MP3; v1.22.0-pre.7): the boundary descriptor, $j^{\ast}$ and the cone (Proposition~\ref{prop:mpboundary}), the
subgroup descriptors, $\res$, $\cor$ on $\mathfrak F'_{\Ga}$, the forget-coset map and the Shapiro comparison with the label factor, the cone of a pair of groups
(Proposition~\ref{prop:mprescor}), with \code{OrbifoldMPBoundaryValidation}, \code{OrbifoldMPLevelConeValidation}; MP-3b (audited F-MP3b; v1.22.0-pre.8): \file{MP3bLevelArms},
\file{MP3bOracle} (the $\Z/\ell^{k}$, $\mathsf R_{\Z}$ and $-I\notin A$ pairs). \textbf{Built at slot 22} (MP-4, audited F-MP4; v1.22.0-pre.8): \file{MPFill} (the MP model's
$\Z$-linear filling), \file{MPHecke} (\code{MPHeckeDatum}: the twisted datum, its certificate, the operator with the label factor, the factorisation $\cor\circ(\conj_g)_\ast\circ\res$),
\file{MP4Oracle.inc}, with \code{OrbifoldMPHeckeValidation} (Remark~\ref{rmk:mphecke}). No console verb,
setting, record or golden object (owner ruling Q4).

\paragraph{Formalization note.} The chapter is homological algebra over $\Z[\Ga]$ on finite data, written to be Lean-friendly: every object is a \code{Finset} given a descriptor,
$d\circ d=0$ is a finite identity per degree, Theorem~\ref{thm:mpcomparison} is uniform in the coefficient ring with no invertibility hypothesis, and the per-fibre contractions
are non-equivariant. The one homological input not in \code{mathlib} at the lane's pin is Lemma~\ref{lem:mpassembly} (row-/column-wise quasi-isomorphism of a bounded double
complex implies one on totals). The Lean rows L0--L10 of [MP] \S8.2 (descriptor predicates, the set, the cochains, $\mathfrak F$, rows, columns, assembly, the resolution,
Theorem~\ref{thm:mpcomparison}, Theorem~\ref{thm:mpnatural}, refinement $\mathsf M_7\leftarrow\mathsf R_{\Z}$, Proposition~\ref{prop:mpbredon}(2)) are the milestone M6 ``orbifold'' item
(owner ruling Q5), the teardrop is not formalised (Q7). \textbf{Markers} (v1.22.0-pre.6; pre.5 carried none): the slot-22 Lean packages MP-L1 / MP-L2 ([GT] Part~G12,
\file{docs/implementation/slot22/MPL1\_report.md} v1.1) prove the rows L0--L7 and L0b, and $30$ clause-scoped \code{\textbackslash LEAN} markers on $11$ labels record them:
Definitions~\ref{def:mprules}, \ref{def:mpell}, \ref{def:mpRZ}, Table~\ref{tab:mprules} (rows $\mathsf M_7$, $\mathsf R_{\Z}$), Definition~\ref{def:mpset}, Remark~\ref{rmk:mpwhat}(1),
Definition~\ref{def:mpfree}, Lemmas~\ref{lem:mprows}, \ref{lem:mpcols}, \ref{lem:mpassembly} (acyclic form) and Theorem~\ref{thm:mpcomparison}(1)--(5), the comparison for the
\emph{unnormalised} complex $C^{\bullet}_{\mathrm{un}}$, natural in $M$, unconditional on $\mathsf M_7$ and $\mathsf R_{\Z}$. \textbf{v1.22.0-pre.7}: the slot-22 Lean package
MP-L3 ([GT] Part~G12, \file{docs/implementation/slot22/MPL3\_report.md} v1.0) proves the rows L8--L10 and the normalised complex, and $25$ further clause-scoped markers record
them: Definition~\ref{def:mpset} (the normalised $C^{\bullet}$ as a subcomplex, its inclusion a quasi-isomorphism), Definition~\ref{def:mpfree} ($\mathfrak F'$ as a simplicial object),
Theorem~\ref{thm:mpcomparison}(2) (normalised $\simeq$ unnormalised on cochains) and (4) (the comparison for $C^{\bullet}$), Definition~\ref{def:mpmorphism}(1)--(3),
Theorem~\ref{thm:mpnatural}(1)--(4) and its \COMP\textsuperscript{r} record on $\mathsf R_{\Z}\to\mathsf M_7$, Proposition~\ref{prop:mpbredon}(1) and the ``if'' of (2). Lean states
Theorem~\ref{thm:mpnatural}(1)--(3) and Proposition~\ref{prop:mpbredon}(1),(2) for $C^{\bullet}_{\mathrm{un}}$ (whence the unnormalised clauses written into them at pre.7), the label
independence of Theorem~\ref{thm:mpnatural}(3) on cohomology only, $\varepsilon^{\mathsf{MP}}$ as $\Phi^{\ast}$ for a lift $\Phi$ of the identity (homotopic to the written one by
Proposition~\ref{prop:mpbredon}(1)). Not formalised, so not marked: the normalised $u^{\ast}$ as a cochain map in its own right, the chain homotopy between label choices, the ranks
$\abs{\mathsf S_k^{\mathrm{nd}}}$, the quasi-isomorphism form of Lemma~\ref{lem:mpassembly}, Proposition~\ref{prop:mpbredon}(2)'s ``not only if'' witness, (3), (4) and the base change,
Proposition~\ref{prop:mpalldegrees}, Remark~\ref{rmk:mphybrid}, Propositions~\ref{prop:mpboundary}, \ref{prop:mprescor} (Lean rows L11--L16, MP-L4).


\appendix
%======================================================================

%----------------------------------------------------------------------
\section{Errata to the superseded documents}\label{app:errata}
%----------------------------------------------------------------------
% SOURCE:  00-unified-contract-architecture.md SS6 (E1-E17, NEW);
%          hecke_extensions_implementation_plan.md SS11.1-11.2;
%          math_corpus_survey.md SS3;  monodromy_v5.tex:1345-1392
% NOTATION: full sec:notation-crosswalk applied
% ERRATA:  this appendix IS the erratum record

\paragraph{What this appendix fixes.}
One entry per defect of the superseded corpus that this document corrects, retires or
inherits the duty to record. Each entry gives, in this order: \textbf{the defective
statement} with its \code{path:line}; \textbf{the correct statement}; \textbf{the
evidence}; and \textbf{where in this document} the correct statement is now made. The
entries are the authority every later chapter cites for its corrections: a chapter that
re-homes text from a superseded document is not finished until every entry touching its
source has been applied.

Three of these are \emph{contradictions} --- two statements that cannot both be true
(\ref{err:e1}, \ref{err:e6}, \ref{err:e12}); the rest are loose arguments, stale
statements, dangling citations, or notation that cannot survive the merge. None of them
is a defect of the shipped implementation: where a document and the code disagreed, the
code was right in every case (\ref{err:e1} most pointedly).

\begin{center}
\begin{longtable}{@{}p{0.05\textwidth}p{0.52\textwidth}p{0.36\textwidth}@{}}
\toprule
\# & defect & corrected in \\
\midrule
\endhead
\ref{err:e1} & primitive vs.\ full $T_{p^{\nu}}$: two numbered prescriptions off by $p^{n}I$ & \S\ref{sec:heckedatum}, \S\ref{sec:obligations} (\ref{obl:algebra}), \S\ref{sec:example} \\
\ref{err:e2} & ``$H^0(B)=Z^0(B)=B^0$'' & \S\ref{sec:cone} \\
\ref{err:e3} & a realisation theorem proved over a field, applied over $\Z/\ell^{e}$ & \S\ref{sec:cone} (hypothesis at point of use) \\
\ref{err:e4} & $\taub$-uniqueness argued from a theorem whose hypothesis it avoids & \S\ref{sec:boundarydatum} (\ref{lem:taubdryexists}) \\
\ref{err:e5} & ``$\tau(jx)$ \ldots both are vertices'' & \S\ref{sec:cone} (\ref{lem:etaexists}) \\
\ref{err:e6} & $T_p^2=T_{p^2}+(p+1)p^{n}$ --- wrong coefficient, false attribution & \S\ref{sec:heckedatum} (\ref{thm:relations}) \\
\ref{err:e7} & \code{[C, thm:explicit\_complex\_equiv]} cited from five source files, label not in [C] & \S\ref{sec:explicitcochaincomplex} \\
\ref{err:e8} & ``issues with torsion, which we ignore'' & \S\ref{sec:conventions} (\ref{rmk:orbifold}, \ref{rmk:negI}), \S\ref{sec:orbifoldmodel} \\
\ref{err:e9} & corestriction agrees only after $\gamma\mapsto\gamma^{-1}$, unmarked & \S\ref{sec:induced} (\ref{prop:omega}) \\
\ref{err:e10} & garbled ``Simplicial interpretation'' display & not re-homed; recorded only \\
\ref{err:e11} & basis vectors named $\eta$, $\delta$, $\beta$ & \S\ref{sec:notation-cochains}, \S\ref{sec:systematiccomputation}, \S\ref{sec:six} \\
\ref{err:e12} & $H^{\bullet}(\partial)$ defined as $\Gbdry$, constructed as $\langle T\rangle$ & \S\ref{sec:boundarydatum} (\ref{def:bdrygroup}, \ref{prop:circlemodel}) \\
\ref{err:e13} & the cone-homotopy sign & \S\ref{sec:cone} (derived, not asserted) \\
\ref{err:e14} & five uncorrected statements in a working note & \S\ref{sec:boundarydatum} (\ref{prop:bdrychar}), \S\ref{sec:six} (\ref{cl:pd}); rest in [E] \\
\ref{err:e15} & ``cone by functoriality'' & \S\ref{sec:cone} is the definition \\
\ref{err:e16} & $T_n$ vs.\ $T_N$ for the composite index & \S\ref{sec:notation-primes}, \S\ref{sec:heckedatum} \\
\ref{err:e17} & no version marker on either superseded document & \S\ref{sec:versioning} \\
\ref{err:e18} & the scope cut: nine labels move from [E] to [C] & \S\ref{sec:six} \\
\ref{err:e19} & [v5] \S12's map $\beta$ is not a morphism of local systems & \S\ref{sec:induced} (\ref{rmk:nogo}) \\
\ref{err:e20} & [v5] \S12.3's bar-cochain formula omits the coset permutation & \S\ref{sec:equivalence} (\ref{rmk:barformula}, \ref{prop:barformula}) \\
\ref{err:e21} & ``$\chi_1\equiv\chi_2$ at every Dickson locus'' --- false as stated & \S\ref{sec:six} (\ref{cl:chi}, \ref{rmk:chicoincide}) \\
\ref{err:e22} & the FP32 totality threshold $B\le2^{10}$ does not follow from the strict criterion it is derived from & \S\ref{sec:fpfilter} (\ref{cor:fptotal}, \ref{cor:fp32}) \\
\ref{err:e23} & the bad-prime oracle asserts $U_4=U_2^{2}$ at $\ell=2$ citing a corollary that carries $\ell\ge5$ & \S\ref{sec:example} (oracle~7) \\
\bottomrule
\end{longtable}
\end{center}

%----------------------------------------------------------------------
\subsection{Primitive versus full \texorpdfstring{$T_{p^{\nu}}$}{T\_\{p\^{}nu\}}: two numbered prescriptions off by \texorpdfstring{$p^{n}I$}{p\^{}n I}}\label{err:e1}
%----------------------------------------------------------------------

\emph{Survey finding D1; plan \S11.2. The most serious defect of the superseded
contract: it contradicts itself, and the two defective sites are a numbered obligation
and a numbered oracle.}

\paragraph{The defective statements.}
\begin{itemize}
\item \file{hecke\_action.tex}:1443--1444, inside Obligation~\ref{obl:algebra}:
  ``compute $T_{\Ga\,\mathrm{diag}(1,p^2)\,\Ga}$ by the same pipeline with its $p^2+p$
  cosets and compare with $T_p^2-p^{\,n+1}$''.
\item \file{hecke\_action.tex}:1595--1596, oracle 6 of the validation plan:
  ``$A^{(q)}(T_4)=A^{(q)}(T_2)^2-2^{\,n+1}I$ with $T_4$ computed independently from
  $\Ga\,\mathrm{diag}(1,4)\,\Ga$ ($6$ cosets)''.
\end{itemize}
In both, $p^2+p$ cosets (at $p=2$: $6$) is the \textbf{primitive} double coset
$\Deprim_{p^2}$, while $T_p^2-p^{\,n+1}I$ is the \textbf{full} operator. The two sides
differ by exactly $p^{\,n}I$, so each prescription, read literally, \emph{fails} against
a correct implementation.

\paragraph{The correct statements.} Both, with the entry point that emits each named:
\begin{contractbox}
For every degree $q$ and every prime $p$, on $H^q(K,\Rel;\Sym^n(\Fl^2))$:
\[
  A^{(q)}\bigl(T_{p^2}^{\mathrm{full}}\bigr) = A^{(q)}(T_p)^2 - p^{\,n+1}I,
  \qquad
  A^{(q)}\bigl(T_{\Deprim_{p^2}}\bigr) = A^{(q)}(T_p)^2 - p^{\,n+1}I - p^{\,n}I,
\]
the two differing by the diamond scalar $\diam=p^{\,n}I$ of
Proposition~\ref{prop:diamond}. At $p=2$: $A^{(q)}(T_4^{\mathrm{full}})=A^{(q)}(T_2)^2-2^{\,n+1}I$
and $A^{(q)}(T_4^{\mathrm{prim}})=A^{(q)}(T_2)^2-2^{\,n+1}I-2^{\,n}I$.
\end{contractbox}
\noindent
\textbf{Which the pipeline emits, from which entry point} (this is the part the
superseded text never said, and the part a reader must have):
\begin{center}
\begin{tabular}{@{}lll@{}}
\toprule
entry point & emits & at $\nu=1$ \\
\midrule
\code{HeckeOperatorFamily::T\_p\_pow($p,\nu$)} & the \textbf{full} $T_{p^{\nu}}$, by the recursion of Thm.~\ref{thm:relations} & $=T_p$ \\
\code{HeckeOperatorFamily::T\_p\_pow\_direct($p,\nu$)} & the \textbf{full} $T_{p^{\nu}}$, by the double-coset assembly \eqref{eq:fullassembly} & $=T_p$ \\
\code{HeckeOperatorFamily::T\_p\_pow\_primitive\_direct($p,\nu$)} & the \textbf{primitive} $T_{\Deprim_{p^{\nu}}}$, the raw pipeline & $=T_p$ \\
\bottomrule
\end{tabular}
\end{center}
\noindent
The raw pipeline \code{HeckeCosetTables($p,\nu$)}$\to$\code{HeckeTauTable}$\to$\code{HeckeOperatorOnK}
computes the \emph{primitive} operator --- which is what made the defective wording
plausible --- but the entry point named \code{T\_p\_pow\_direct} does not return it: it
assembles the full operator from primitive terms by
$T_{p^{\nu}}^{\mathrm{full}}=\sum_{0\le 2j\le\nu}p^{\,jn}\,T_{\Deprim_{p^{\nu-2j}}}$,
with $T_{\Deprim_{p^{0}}}:=I$. The three coincide at $\nu=1$, so $T_p$ is unambiguous.

\paragraph{Evidence.} Four independent confirmations, obtained from the code and from a
run rather than from prose.
\begin{enumerate}
\item \textbf{The classical decomposition} of the determinant-$p^{\nu}$ Hecke operator
  into primitive double cosets, \STD{\cite[Thm.~3.24(1), (2), p.~63]{shimura}, the $T(p^{\nu})$
  decomposition; v1.19.0-pre.5: notation corrected; v1.21.0-pre.1, R20-L: locator corrected --- the earlier ``Thm.~3.34'' is the structure of the
  level-$N$ ring $R(\Ga',\De')$, not this decomposition}, is what makes the full operator --- and not the primitive one ---
  satisfy the recursion of Theorem~\ref{thm:relations}.
\item \textbf{The degree count}, which settles it with no convention at all: $\Deprim_p$
  has $p+1$ cosets, the full $\De_{p^2}$ has $1+p+p^2$, and $\Deprim_{p^2}$ has $p^2+p$.
  The missing coset is the central one, $\diam$, contributing $p^{\,n}$.
\item \textbf{The implementation}, whose header states the distinction normatively and
  whose \code{T\_p\_pow\_direct} implements the sum above
  (\file{Hecke/HeckeOperatorFamily.h}:136--173,
  \file{Hecke/HeckeOperatorFamily.cpp}:234--268).
\item \textbf{Two runs} on the tree of 2026-09-15
  (\code{RelWithDebInfo}, \code{make -j16}), each exit code $0$:
  \begin{itemize}
  \item \code{./cohomology\_app HeckeOperatorFamilySmokeValidation} $\to$ ``PASSED (for
    $p\in\{2,3\}$, \code{ell} $\in\{5,13,101\}$, $n\in\{0,2,10\}$: \ldots
    \code{T\_p\_pow(p,2)==T\_p\_pow\_direct(p,2)} [FULL recursion == FULL double-coset],
    and \code{T\_p\_pow\_direct(p,2)==T'\_\{p\^{}2\}+p\^{}n*I} [primitive-vs-full
    decomposition witness])''.
  \item \code{./cohomology\_app HeckeRelationsValidation} $\to$ ``O9 PASSED at
    \code{(ell,n)} $=(13,0),(13,2),(13,10),(101,0),(101,2),(101,10)$ for all
    $q=0..2$ [commutativity, \code{T\_4=T\_2\^{}2-2\^{}\{n+1\}I} (full, independent),
    \ldots, decomposition witness]''. Oracle O9.2 \emph{is} the assertion
    $A^{(q)}(T_4^{\mathrm{full}})=A^{(q)}(T_2)^2-2^{\,n+1}I$ and O9.5 \emph{is}
    $A^{(q)}(T_4^{\mathrm{full}})=A^{(q)}(T_4^{\mathrm{prim}})+2^{\,n}I$.
  \end{itemize}
  A hand-checkable instance, from \code{HeckeOperatorSmokeValidation} (also exit $0$) at
  $(p,\ell,n)=(2,5,0)$, where $A^{(0)}(T_2)=(3)=(1+p)$: the full operator gives
  $3^2-2^{0+1}=7\equiv 2\pmod 5$, which is $\sigma_1(4)=1+2+4$ (degree $1+p+p^2=7$); the
  primitive gives $7-2^0=6\equiv 1$ (degree $p+p^2=6$).
\end{enumerate}

\paragraph{Where the correct statement is now made.} \S\ref{sec:heckedatum}
(Definition~\ref{def:operators}, \eqref{eq:fullop}, \eqref{eq:fullassembly} and the
caveat to Theorem~\ref{thm:relations}) --- which the superseded document had already
reworded correctly; \S\ref{sec:obligations}, Obligation~\ref{obl:algebra}, with the
corrected prescription; and \S\ref{sec:example}, oracle 6, likewise. The half-applied
state was itself recorded at the time
(\file{docs/implementation/hecke\_action\_implementation\_plan.md}:45 directed the
rewording of \S\code{sec:heckedatum} only; \S\code{sec:obligations} and
\S\code{sec:example} were never in its scope, and \code{:81} carries the same stale
phrasing).

%----------------------------------------------------------------------
\subsection{``\texorpdfstring{$H^0(B)=Z^0(B)=B^0$}{H0(B)=Z0(B)=B0}''}\label{err:e2}
%----------------------------------------------------------------------
\emph{Plan \S11.1 item 1; a loose step in the cone realisation argument.}

\textbf{Defective statement.} \file{hecke\_extensions\_OQ2\_verdict.md} \S2.2:
``$H^0(B)=Z^0(B)=B^0$ (no coboundaries in degree $0$)''.
\textbf{Correct statement.} $Z^0(B)=\ker\bigl(d\colon B^0\to B^1\bigr)\subseteq B^0$,
with equality only if $d^0=0$. The first equality is right (there are indeed no
coboundaries in degree $0$); the second is not.
\textbf{Evidence.} Immediate. The construction is unaffected: it only ever needs values
in $Z^0(B)$.
\textbf{Corrected in} \S\ref{sec:cone}, which states the realisation argument with
$Z^0(B)$ and never with $B^0$.

%----------------------------------------------------------------------
\subsection{A realisation theorem proved over a field, applied over \texorpdfstring{$\Z/\ell^{e}$}{Z/ell\^{}e}}\label{err:e3}
%----------------------------------------------------------------------
\emph{Plan \S11.1 item 2. The rule~\ref{A1} failure mode, in its purest form.}

\textbf{Defective statement.} The realisation theorem of
\file{hecke\_extensions\_OQ2\_verdict.md} \S2.2 --- ``every element of
$\Hom(\ker\rres^1,\coker\rres^0)$ is realised'' --- uses a splitting of the cochain
groups and is stated for $\cring$ \emph{a field}. The hypothesis is not repeated in the
statement, and the surrounding architecture has since put $\cring=\Z/\ell^{e}$ and
$\Z/m$ in scope.
\textbf{Correct statement.} The theorem holds as stated for $\cring=\Fl$. It is
\textbf{not} claimed over $\Z/\ell^{e}$ or $\Z/m$, where the splitting is unavailable.
\textbf{Evidence.} The proof's use of a splitting of $\cring$-modules.
\textbf{Corrected in} \S\ref{sec:cone}, where the hypothesis $\cring=\Fl$ is stated in
the statement that uses it, and the $\Z/\ell^{e}$ case is explicitly not claimed.

%----------------------------------------------------------------------
\subsection{\texorpdfstring{$\taub$}{tau\_bdry}-uniqueness argued from a theorem whose hypothesis it avoids}\label{err:e4}
%----------------------------------------------------------------------
\emph{Plan \S11.1 item 3. Circular as written.}

\textbf{Defective statement.} \file{hecke\_extensions\_OQ2\_verdict.md} \S2.6.1 derives
the homotopy-uniqueness of $\taub$ from ``the comparison theorem''.
\textbf{Why it is defective.} That theorem's relative-projectivity hypothesis is
precisely what forces $\ell\ge5$ in Lemma~\ref{lem:tauexists}(2). Citing it makes the
conclusion conditional on exactly the hypothesis the boundary construction is supposed
to avoid.
\textbf{Correct statement.} The conclusion is right over $\Z$ \emph{for every} $\ell$,
by an explicit argument: $N_0(x)$ is the unique $1$-chain of the contractible,
one-dimensional component with $\partial N_0(x)=(\taub-\taub')(x)$; $N_1=0$ because
$Z_1(\Xbdry)=0$; and the (T1)-extension is well defined because the boundary
stabilisers are central and act trivially.
\textbf{Corrected in} \S\ref{sec:boundarydatum}, Lemma~\ref{lem:taubdryexists}, which
gives that explicit argument and never the abstract citation.

%----------------------------------------------------------------------
\subsection{``\texorpdfstring{$\tau(jx)$}{tau(jx)} \ldots both are vertices''}\label{err:e5}
%----------------------------------------------------------------------
\emph{Plan \S11.1 item 4; loose wording, construction unaffected.}

\textbf{Defective statement.} \file{hecke\_extensions\_OQ2\_verdict.md} \S2.3.1:
``$\etapair_0(x):=$ an edge path in $\Xbulk$ from $\tau(jx)$ to $j\taub(x)$ (both are
vertices)''.
\textbf{Correct statement.} $\tau(jx)$ is a $0$-chain of augmentation $1$, not
necessarily a single vertex. The construction survives because $\Xbulk$ is connected and
$\aug\bigl(j\taub(x)-\tau(jx)\bigr)=0$ (both (T3)s give $\aug=1$), so the difference is
a boundary.
\textbf{Evidence.} The augmentation computation above.
\textbf{Corrected in} \S\ref{sec:cone}, Lemma~\ref{lem:etaexists}, which states the
identity and the augmentation argument rather than the recipe.

%----------------------------------------------------------------------
\subsection{\texorpdfstring{$T_p^2=T_{p^2}+(p+1)p^{n}$}{Tp\^{}2 = Tp\^{}2 + (p+1)p\^{}n} --- wrong coefficient, false attribution}\label{err:e6}
%----------------------------------------------------------------------
\emph{Survey finding D2; plan \S11.1 item 5. A wrong formula, attributed to a document
that does not contain it.}

\textbf{Defective statements.}
\file{hecke\_extensions\_OQ2\_verdict.md}:278 --- ``$T_p^2=T_{p^2}+(p+1)p^n$ with
$T(p,p)=p^n$ on $\Sym^n$, [C] Corollary on $T_{p^2}$''. The same error stood at
\file{hecke\_extensions\_mathematics.md}:63, as
``$T(1,q)^2=T(1,q^2)+(q+1)T(q,q)$'' --- in Shimura's notation (adopted since v1.19.0-pre.5) this is the \emph{true} identity for the primitive $T(1,q^2)=T_{\Deprim_{q^2}}$ ($T(q^2)=T(1,q^2)+T(q,q)$); the defect was its reading with $T(1,q^2)$ as the full $T_{q^2}$, the inversion corrected at v1.19.0-pre.5, so the erratum concerns the naming, not the identity.
\textbf{Correct statement.}
\[
  T_p^{\,2} \;=\; T_{p^2} + p\cdot\diam \;=\; T_{p^2} + p^{\,n+1},
\]
the coefficient being $p$, not $p+1$.
\textbf{Evidence}, three ways.
\begin{enumerate}
\item The superseded contract has it right --- \file{hecke\_action.tex}:507--516,
  Theorem~\ref{thm:relations}(1): $T_{p^{\nu+1}}=T_pT_{p^{\nu}}-p\,\diam\,T_{p^{\nu-1}}
  =T_pT_{p^{\nu}}-p^{\,n+1}T_{p^{\nu-1}}$, which at $\nu=1$ with $T_{p^0}=I$ is the
  display above. So the attribution ``[C] Corollary'' is false: [C] says the opposite.
\item The classical statement, \STD{\cite[Ch.~3]{shimura}}: $T_p^2=T_{p^2}+p^{\,k-1}
  \langle p\rangle$ in weight $k$; here $k=n+2$, so $p^{\,k-1}=p^{\,n+1}$, and
  $\langle p\rangle=\diam$ acts on $\Sym^n$ by $p^{\,n}$.
\item The degree count, which needs no convention: $\deg T_p^2=(p+1)^2=p^2+2p+1$,
  $\deg T_{p^2}=1+p+p^2$ (the \emph{full} operator), $\deg\diam=1$, so the coefficient is
  $(p+1)^2-(1+p+p^2)=p$. The value $p+1$ is the \emph{number of cosets} of $\Deprim_p$
  --- a different quantity, and presumably how it got in.
\end{enumerate}
\textbf{Corrected in} \S\ref{sec:heckedatum}, Theorem~\ref{thm:relations}, which
already carried the correct relation and carries it here unchanged. The defective
instance is recorded because the review document that contains it stands as published
and is not edited; \file{hecke\_extensions\_mathematics.md}:63 was corrected in place on
2026-09-15. The \emph{conclusion} of the passage --- that the Hecke relations hold on
the nose on $H^{\bullet}_c$ --- is unaffected.

%----------------------------------------------------------------------
\subsection{A \texorpdfstring{\code{[C, \ldots]}}{[C, ...]} citation to a label that was not in [C]}\label{err:e7}
%----------------------------------------------------------------------
\emph{Found during the architecture pass of 2026-09-15; not in the corpus survey.}

\textbf{Defective state.} Five C++ source files cite
\code{[C, thm:explicit\_complex\_equiv]} --- the load-bearing isomorphism $\Psi$ ---
namely \file{LocalSystemSL2Z.h}:60, \file{LocalSystemSL2Z.cpp}:190,
\file{InducedLocalSystem.cpp}:150, \file{HeckeTessellationValidation.h}:36 and
\file{.cpp}:396, and \file{HeckeTessellationSmokeValidation.cpp}:64. That label was in
\file{monodromy\_v5.tex}:587, \textbf{not} in [C]. \file{LocalSystemSL2Z.cpp}:190 hedges
by naming both files (``in \file{docs/latex/hecke\_action.tex} /
\file{monodromy\_v5.tex}'').
\textbf{Correct state.} The citations were correct in intent and are now correct in
fact: the framework of [v5] is re-homed here in full and the label comes with it.
\textbf{Evidence.} The label census of the architecture pass; the label is absent from
\file{hecke\_action.tex} and present at \file{monodromy\_v5.tex}:587.
\textbf{Corrected in} \S\ref{sec:explicitcochaincomplex}, which carries
Theorem~\ref{thm:explicit_complex_equiv} verbatim under that label. No source file
changes. A label census run against this document
(\S\ref{sec:machinery}) would have caught this defect at the time it was introduced,
and is registered for exactly that reason.

%----------------------------------------------------------------------
\subsection{``issues with torsion, which we ignore''}\label{err:e8}
%----------------------------------------------------------------------
\emph{Survey finding S1, first bullet. The [v5] \S12 erratum block does not cover it.}

\textbf{Defective statement.} \file{monodromy\_v5.tex}:1490: ``Since
$\piOne(\Xbulk)\cong\Ga$ (up to issues with torsion, which we ignore for this
theoretical overview)''.
\textbf{Why it is defective.} That torsion is exactly what the model makes
load-bearing. The central relation $-I$ is the difference between $H^{\bullet}=0$ and a
spuriously non-zero $H^{\bullet}$ for odd $n$; ``ignore'' is the wrong instruction, and
the sentence appears in a passage a reader reaches while looking for the framework.
\textbf{Correct statement.} $\Ga=\SL_2(\Z)$ has torsion, $\piOne$ of the quotient
\emph{orbifold} is $\Ga$ only in the orbifold sense, and the elliptic and central
stabilisers enter the equivariance constraints explicitly.
\textbf{Corrected in} \S\ref{sec:conventions}: Remark~\ref{rmk:orbifold} and
Remark~\ref{rmk:negI} state it as a hypothesis with its consequences, and
\S\ref{sec:orbifoldmodel} treats the $6\notin\cring^{\times}$ regime properly rather than
by exclusion.

%----------------------------------------------------------------------
\subsection{The corestriction convention, unmarked at the source}\label{err:e9}
%----------------------------------------------------------------------
\emph{Survey finding S1, second bullet.}

\textbf{Defective statement.} \file{monodromy\_v5.tex}:1536 gives the corestriction as
$\omega_{\ast}(\psi)=\sum_{\gamma\in\Ga''\backslash\Ga}\rho_L(\gamma)\psi(\gamma^{-1})$,
which agrees with the convention used here only after the substitution
$\gamma\mapsto\gamma^{-1}$ in the summation --- a fact recorded at
\file{hecke\_action.tex}:906 but \emph{not} marked at [v5], where a transcriber will meet
it first.
\textbf{Correct statement.} True as stated, under its own convention; not
interchangeable with the one used here without the substitution.
\textbf{Corrected in} \S\ref{sec:induced}, Proposition~\ref{prop:omega}: one convention,
fixed once, with the substitution shown explicitly.

%----------------------------------------------------------------------
\subsection{A garbled display in ``Simplicial interpretation''}\label{err:e10}
%----------------------------------------------------------------------
\emph{Survey finding S1, third bullet.}

\textbf{Defective statement.} \file{monodromy\_v5.tex}:1545--1553 contains
\code{C\_*(K\^{}u; Z), W) -> C\_eq\^{}*(K,R;W)}: an unbalanced parenthesis and a map
with no source.
\textbf{Disposition.} \textbf{Not re-homed.} [v5] \S12 is superseded in full
(\ref{err:e19}, \ref{err:e20}); this passage is recorded here so that the supersession is
complete and no reader mistakes the omission for an oversight. The mathematics the
passage was reaching for is Theorem~\ref{thm:explicit_complex_equiv}
(\S\ref{sec:explicitcochaincomplex}), stated correctly.

%----------------------------------------------------------------------
\subsection{Basis vectors named \texorpdfstring{$\eta$, $\delta$, $\beta$}{eta, delta, beta}}\label{err:e11}
%----------------------------------------------------------------------
\emph{Survey collisions C6--C8. Three reserved symbols used for basis vectors, inside the
cone section that also needs all three for something else.}

\textbf{Defective usage.} \file{monodromy\_v5.tex} \S7--\S8 writes $\eta_p^{(q)}$,
$\delta_s^{(q)}$ and $\beta_j^{(q)}$ for the stored basis vectors of $\Bq{q}$, $\Zq{q}$
and $\Cq{q}$ (e.g.\ \code{:1098--1109}, \code{:1182}, \code{:1208--1211},
\code{:1325--1328}) --- colliding with the pair homotopy $\etapair$, the coset element
$\delta_j$ and the right-coset representative $\beta_j$ respectively, and additionally
indexing the first by $p$, which is the Hecke prime.
\textbf{Correct usage.} The basis-name scheme $c^{(q)}_j$ / $z^{(q)}_j$ / $b^{(q)}_j$ for
$\Cq{q}$ / $\Zq{q}$ / $\Bq{q}$, indexed by $j$.
\textbf{Corrected in} \S\ref{sec:notation-cochains} (the scheme),
\S\ref{sec:systematiccomputation} and \S\ref{sec:six} (the text that uses it).
Typographic only: no label and no code identifier moves.

%----------------------------------------------------------------------
\subsection{\texorpdfstring{$H^{\bullet}(\partial)$}{H(bdry)} defined as one group and constructed as another}\label{err:e12}
%----------------------------------------------------------------------
\emph{Survey finding D4; verdict \S5.5. A definition and its own normative construction
disagree inside one document.}

\textbf{Defective state.} \file{hecke\_extensions.tex}:189 \emph{defines}
$H^{\bullet}(\partial;M):=H^{\bullet}(\Gbdry;M)$ with $\Gbdry=\langle T\rangle\times
\langle-I\rangle$, while its normative construction (A1), \code{:755--768}, is the
circle model, which computes $H^{\bullet}(\langle T\rangle;M)$. The two agree only when
$2\in\cring^{\times}$.
\textbf{Correct statement.} The definition is $H^{\bullet}(\Gbdry;\VR)$. The circle
model computes it \emph{under the hypothesis} $2\in\cring^{\times}$, and that hypothesis
is part of the statement, not of a preamble.
\textbf{Corrected in} \S\ref{sec:boundarydatum}: Definition~\ref{def:bdrygroup} is the
definition, Proposition~\ref{prop:circlemodel} is the specialisation statement carrying
$2\in\cring^{\times}$ at the point of use, and \S\ref{sec:ell2} states what is
\emph{not} settled at $\ell=2$ and the policy decided for it (\ref{open:ell2}) rather than assuming it away.

%----------------------------------------------------------------------
\subsection{The cone-homotopy sign}\label{err:e13}
%----------------------------------------------------------------------
\emph{Survey finding D5; historical and annotated, but the wrong line is in the body.}

\textbf{Defective statement.}
\file{hecke\_extensions\_plan\_audit.md}:271--274 asserts
$d\hder+\hder d=\rres\,\Th_g-\Thb_g\,\rres$. A dated correction box three lines later
(\code{:276--288}) gives the right sign, but a reader who stops at the bullet does not
reach it.
\textbf{Correct statement.} $d\hder+\hder d=\Thb_g\,\rres-\rres\,\Th_g$, which this
document does not assert but \emph{derives}, from the chain-level identity
$\partial\etapair+\etapair\partial=j\taub-\tau j$ together with $\hder:=\widetilde{\etapair}$.
\textbf{Evidence.} Re-derived independently twice in the corpus; the audit's own sweep
found $17+14$ consistent occurrences elsewhere and \textbf{zero} silent uses of the
wrong orientation, so nothing downstream was contaminated.
\textbf{Corrected in} \S\ref{sec:cone}, where the sign is fixed once at chain level and
the cochain identity is a consequence --- so there is no second place for it to be
asserted differently.

%----------------------------------------------------------------------
\subsection{Five uncorrected statements in a working note}\label{err:e14}
%----------------------------------------------------------------------
\emph{Survey finding S2. The note is what two other documents cite as their source.}

\textbf{Defective state.} \file{hecke\_extensions\_mathematics.md} carries five
statements corrected only in [E] and in the [E] audit, never in the note itself:
\code{:32} (``the exact `+1' should be double-checked'' --- the $\Z$-module Sturm
statement for $M_k$ is false at $\ell\ge5$, counterexample $(n,\ell)=(32,131)$);
\code{:35} (the duality consequence stated for the \emph{Loewy} series);
\code{:39} (the Bella\"iche--Khare recall, and ``no single $T_p$ ever suffices'');
\code{:63} (the $(p+1)$ coefficient --- \ref{err:e6}, since fixed in place);
\code{:67} (the boundary characters presented in the $\beta_j$, i.e.\ right-coset,
convention as though it were \emph{the} answer, with only a ``Convention warning''
as a hedge).
\textbf{Correct statements.} It is the \emph{socle} series that reverses under duality,
not the Loewy series; the boundary characters must be stated in the normative
(left-coset $\delta_j$) convention, the $\beta_j$ convention giving the $\det$-twisted
operator and \emph{swapping} the two families; the Sturm and Bella\"iche--Khare items are
as [E] carries them.
\textbf{Corrected in} \S\ref{sec:boundarydatum}, Proposition~\ref{prop:bdrychar} (the
characters, in the normative convention, with the twist stated) and \S\ref{sec:six},
Claim~\ref{cl:pd} (socle, not Loewy). The Sturm bound and the Bella\"iche--Khare item are
programme material and stay in [E], which already carries them correctly.
\textbf{Status of the note.} \file{hecke\_extensions\_mathematics.md} is a working note.
It is \textbf{not} a source of record, and nothing in this document or in [E] may cite
it as one.

%----------------------------------------------------------------------
\subsection{``cone by functoriality''}\label{err:e15}
%----------------------------------------------------------------------
\emph{Survey finding S3.}

\textbf{Defective statement.}
\file{hecke\_extensions\_computational\_strategy.md}:15 lists, as a Stage~0 deliverable,
``a routine \code{hecke(q)} returning the matrices of $T_p$ on all $X_i$ (\ldots\ cone by
functoriality)''.
\textbf{Why it is defective.} ``By functoriality'' is not a definition. The restriction
$\rres$ is Hecke-equivariant only \emph{up to homotopy}, and with a freely chosen
cochain homotopy the induced operator on $X_3=H^1_c$ is genuinely undetermined. Pinning
it down is the entire content of \S\ref{sec:cone}.
\textbf{Corrected in} \S\ref{sec:cone}, which \emph{is} the definition: $\hder$ is
derived from $\etapair$ and never freely chosen (rule~\ref{A6}), and
Proposition~\ref{prop:coneindep} is what makes the result independent of the remaining
freedom.

%----------------------------------------------------------------------
\subsection{\texorpdfstring{$T_n$}{T\_n} versus \texorpdfstring{$T_N$}{T\_N} for the composite index}\label{err:e16}
%----------------------------------------------------------------------
\emph{Survey collision C2. The superseded contract contradicted itself on one symbol.}

\textbf{Defective usage.} \file{hecke\_action.tex}:80 (the abstract) and \code{:1439}
write $T_n$ for the composite-index operator while $n$ is, throughout the same document,
the symmetric power; \code{:518} correctly writes $T_N$.
\textbf{Correct usage.} $T_N$ with $N\in\Z_{>0}$, everywhere; $n$ is the symmetric power
and nothing else.
\textbf{Corrected in} \S\ref{sec:notation-primes} (the rule) and
\S\ref{sec:heckedatum} (the text). Typographic; no label moves.

%----------------------------------------------------------------------
\subsection{No version marker on either superseded document}\label{err:e17}
%----------------------------------------------------------------------
\emph{Survey finding S6.}

\textbf{Defective state.} \file{hecke\_action.tex}:72 and
\file{monodromy\_v5.tex}:41 both carry \code{\textbackslash date\{\textbackslash today\}},
so neither has a citable version and every rebuild silently re-dates the PDF. The
committed \file{hecke\_action.pdf} is dated 2026-07-03 and the \code{.tex} has not
changed since, so nothing is out of sync today --- but one rebuild would make the PDF
and every ``as of'' statement about it unverifiable, with no diff to show for it.
\textbf{Correct state.} An explicit version line and a changelog.
\textbf{Corrected in} \S\ref{sec:versioning}, which fixes the convention for this
document, and in the frozen banners both superseded documents received at v1.0.0
(\ref{V4}). \file{hecke\_action.tex} now carries \emph{Version 1.0 --- frozen 2026-07-02}
and \file{monodromy\_v5.tex} \emph{Version 5.0 --- frozen 2026-07-03}, each date being that
of the last commit that changed the file's content; neither file uses
\code{\textbackslash today} any more, so a rebuild no longer re-dates either PDF and ``as
of \file{hecke\_action.tex} v1.0'' is a verifiable statement.

%----------------------------------------------------------------------
\subsection{The scope cut: nine labels move from [E] into [C]}\label{err:e18}
%----------------------------------------------------------------------
\emph{Not a defect: a recorded relocation, entered here because
rule~\ref{L1} makes the move of a label an event that must be minuted.}

\textbf{What moved.} \code{sec:six}, \code{cl:amalgam}, \code{cl:dimH1},
\code{cl:bockstein}, \code{cl:boundary}, \code{cl:les}, \code{eq:les}, \code{cl:chi} and
\code{cl:pd} were defined in \file{hecke\_extensions.tex} and are defined here, in
\S\ref{sec:six}.
\textbf{Why.} They are facts about the objects, not a programme (\S\ref{sec:scope}), and
they are the oracles Obligation~\ref{obl:oracle} reads. A normative obligation may not
depend on a versioned, churning companion.
\textbf{What it costs.} Eleven citation sites in project documents, written
\code{[E, \ldots]}, become \code{[C, \ldots]} --- \code{cl:les} at five,
\code{cl:pd} at two, the remaining six at one each. \textbf{No C++ source file cites any
of the nine}, so no code changes. The labels themselves are unchanged, per
rule~\ref{L1}.
\textbf{Completed at v1.0.0.} Until then the move was one-sided: the nine were defined
here \emph{and} still defined in the live \file{hecke\_extensions.tex}, which is why
\ref{V5}\ref{V5ii} held the leading version field at $0$.
\file{hecke\_extensions.tex} v1.3 removes its nine copies, replaces them with a pointer
section citing this document, and names as superseded the four of the nine whose content
this document \emph{corrected} rather than transported: \code{cl:chi}'s Dickson-locus
sentence (entry~\ref{err:e21}), \code{cl:boundary}'s unhypothesised
$H^{i}(\partial)=0$ for $i\ge2$ (Remark~\ref{rmk:bdryrepairs}), \code{cl:les}'s
hypothesis pair (Remark~\ref{rmk:sixscope}) and \code{cl:pd}'s self-adjointness wording
(Remark~\ref{rmk:pdscope}). One \code{[E, cl:les]} citation survives in the corpus, in
the quoted body of \file{docs/implementation/hecke\_extensions\_OQ2\_verdict.md}; it is a
verbatim quotation inside a dated review record and is left as written.

%----------------------------------------------------------------------
\subsection{[v5] \S12: the map \texorpdfstring{$\beta$}{beta} is not a morphism of local systems}\label{err:e19}
%----------------------------------------------------------------------
\emph{Recorded at \file{monodromy\_v5.tex}:1345--1368 and superseded there; carried here
because [v5] \S12 is dropped and this document is now the only record.}

\textbf{Defective statement.} \file{monodromy\_v5.tex}:1532--1534 defines a
``conjugation'' map
$\beta_{\ast}(\phi)(\gamma)=\rho_L(g)\phi(g^{-1}\gamma g)$ intended as a morphism
$\Ind_{\Ga'}^{\Ga}\Vcoef\to\Ind_{\Ga''}^{\Ga}\Vcoef$, and builds the Hecke operator
as $\omega\circ\beta\circ\alpha$ (\S\S12.5--12.7).
\textbf{Correct statement --- a no-go, not a repair.} Two independent reasons.
\emph{(i)} The formula does not define a function on $\Ga$ at all: $g^{-1}\gamma g\in\Ga$
only for $\gamma\in\Ga''$. \emph{(ii)} More decisively, \emph{no} morphism of
$\Ga$-modules $\Ind_{\Ga'}^{\Ga}\Vcoef\to\Ind_{\Ga''}^{\Ga}\Vcoef$ whatsoever can make
$\omega\circ\mu\circ\alpha$ realise $T_p$: any such composite
$\Vcoef\to\Ind_{\Ga'}^{\Ga}\Vcoef\to\Ind_{\Ga''}^{\Ga}\Vcoef\to\Vcoef$ lies in
$\End_{\Ga}(\Vcoef)$, a scalar ring for irreducible $\Vcoef=\Sym^{10}\Q^2$, whereas
$T_2$ on $H^1(\SL_2(\Z);\Sym^{10}\Q^2)$ has eigenvalues $-24,-24,2049$
(\STD{Eichler--Shimura in weight $12$, \cite[Ch.~8]{shimura}}) and is therefore not
scalar.
\textbf{Corrected in} \S\ref{sec:induced}, Remark~\ref{rmk:nogo}, which keeps that label
and states the no-go; the middle map is written $\mu$ throughout (it is $\beta$ in [v5]
\S12, and $\beta_j$ is a coset representative here). The construction that replaces the
sketch is the fused cochain map $\Th_g=\omega_{\bullet}\circ\Xig$ of
Definition~\ref{def:xi}, assembled by Algorithm~\ref{alg:assembly}, which needs neither
the intermediate group $\Ga'$ nor any $\mu$.
\textbf{Consequence for the code.} The fix is implemented: the generalisation of a
single constant fibre matrix to per-simplex sparse block rows
(\code{EquivariantCochainComplexMorphismOnK}, \code{FiberBlockTable}); the legacy
single-matrix constructor is unchanged and still produces byte-identical results.

%----------------------------------------------------------------------
\subsection{[v5] \S12.3: the bar-cochain formula omits the coset permutation}\label{err:e20}
%----------------------------------------------------------------------
\emph{Recorded at \file{monodromy\_v5.tex}:1369--1379; carried here for the same reason
as \ref{err:e19}.}

\textbf{Defective statement.} The ``Algebraic Hecke Operator'' formula of [v5] \S12.3,
$(T_g\phi)(\gamma_1,\dots,\gamma_s)=\sum_j\rho_L(g_j)\phi(g_j^{-1}\gamma_1g_j,\dots,
g_j^{-1}\gamma_sg_j)$, silently assumes that every $\gamma_i$ stabilises every coset
$\Ga g_j$.
\textbf{Correct statement.} The classical formula must carry the coset permutation
$j\mapsto\gamma\ast j$ explicitly; without it the formula does not commute with the bar
differential.
\textbf{Corrected in} \S\ref{sec:equivalence}: Remark~\ref{rmk:barformula} states the
deviation and Proposition~\ref{prop:barformula} gives the corrected degree-$1$ formula,
both keeping their labels. The corrected formula is what the independent bar-cocycle
route implements, and that route is cross-checked against the fused pipeline
(Obligation~\ref{obl:agreement}).

%----------------------------------------------------------------------
\subsection{``$\chi_1\equiv\chi_2$ at every Dickson locus''}\label{err:e21}
%----------------------------------------------------------------------
\emph{Found while Chapter~\ref{sec:six} was written (2026-09-16), in a statement being
re-homed from [E] by the scope cut of entry~\ref{err:e18}. It is a defect in a
\emph{claim}, not in a working note, and it is corrected at the label that carries it.}

\textbf{Defective statement.} \file{hecke\_extensions.tex}:323--324, inside
\code{cl:chi}: ``$\chi_1\equiv\chi_2$ at every Dickson locus, so Dickson classes cannot
discriminate the two indexings ($T_2=1$ on $L$ at $(6,5)$: $\chi_1(2)=2+2^{6}\equiv1$)''.
\textbf{Why it is defective.} By that same claim's own coincidence criterion ---
$\chi_a\equiv\chi_{a'}$ iff $a\equiv a'$ or $a\equiv n+1-a'\pmod{\ell-1}$, which is
correct and is proved at Claim~\ref{cl:chi}(3) --- one has $\chi_1\equiv\chi_2$ iff
$(\ell-1)\mid(n-2)$. On the Dickson locus $n=a(\ell+1)+b(\ell^{2}-\ell)$ one has
$n\equiv2a\pmod{\ell-1}$, so the condition reads $(\ell-1)\mid2(a-1)$, which is
\emph{not} automatic. First counterexample: $(n,\ell)=(12,5)$, a Dickson locus
($a=2$, $b=0$, $\dim H^{0}=1$), where $\chi_1(2)=2+2^{12}\equiv3$ and
$\chi_2(2)=2^{2}+2^{11}\equiv2$ in $\F_5$ --- different.
\textbf{Correct statement.} $\chi_1\equiv\chi_2$ exactly when $(\ell-1)\mid(n-2)$; in
particular at the \emph{first} Dickson class $n=\ell+1$, which is [E]'s own example
$(n,\ell)=(6,5)$, and there always, since then $n-2=\ell-1$. The conclusion [E] drew the
sentence for survives in that form: the first Dickson class never separates the
left-coset from the right-coset indexing, so the separation must come from the Eisenstein
eigenvalue $1+p^{\,n+1}$ instead.
\textbf{Corrected in} \S\ref{sec:six}: Claim~\ref{cl:chi}(3) states the criterion and
proves it, and Remark~\ref{rmk:chicoincide} states what is and is not true on the Dickson
locus, with the counterexample and the computational check. \code{cl:chi} keeps its
label (rule~\ref{L2}); what changes is one sentence of its commentary, and the change is
recorded here and in Appendix~\ref{app:changelog}.
\textbf{What it does not affect.} Nothing in this document or in [E] branches on the
sentence. Proposition~\ref{prop:bdrychar}(2) and Obligation~\ref{obl:taubdry}(vii)
already settle the coset convention by the Eisenstein eigenvalue and not by a Dickson
class.

%----------------------------------------------------------------------
\subsection{The FP32 totality threshold does not follow from the criterion it is derived from}\label{err:e22}
%----------------------------------------------------------------------
\emph{Found while Chapter~\ref{sec:fpfilter} was written (2026-09-16), in a statement
being re-homed. The threshold is correct and the shipped constant is right; the
\emph{derivation} of it from the neighbouring corollary is off by the strictness of one
inequality, and the two statements therefore do not compose as written.}

\textbf{Defective statements.} \file{hecke\_action.tex}:1385--1391, inside
\code{cor:fptotal}: ``If $E(B)<\tfrac12$ \ldots\ the FP64 filter is \emph{total}'';
and :1394--1398, inside \code{cor:fp32}: ``the filter is total iff $2^{-21}B^2<\tfrac12$,
i.e.\ $B\le2^{10}$''.
\textbf{Why it is defective.} At $B=2^{10}$ the binary32 threshold is
$E_{32}(2^{10})=2^{-21}\cdot2^{20}=\tfrac12$ exactly, which does \emph{not} satisfy the
strict criterion $E<\tfrac12$ that \code{cor:fptotal} supplies; read literally,
\code{cor:fp32}'s own criterion yields $B\le1023$, not $B\le2^{10}$. Separately, ``iff''
overstates: the criterion is sufficient for totality, not necessary.
\textbf{Correct statement.} The criterion is the non-strict $E(B)\le\tfrac12$ --- the
proof needs $E(B)+\abs{\hat x-x}<1$, and $\abs{\hat x-x}\le\tfrac12E(B)(1+2^{-53})$, so
$E(B)\le\tfrac12$ gives $\tfrac34+2^{-55}<1$ with room to spare. Under it, $B\le2^{10}$
is certified in binary32 and $B\le2^{24}$ in binary64, both as stated.
\textbf{Evidence, and a second and stronger route.} Independently of any error estimate,
both totality regimes are \emph{exact}-arithmetic regimes: $2B^2\le2^{\mathrm{prec}}$
makes every product and the sum an exactly representable integer, so $\hat x=x$ and
$\abs{\hat x}\le E(B)$ decides $x=0$ whenever $E(B)<1$. At $B=2^{10}$ in binary32,
$2B^2=2^{21}\le2^{24}$; at $B=2^{24}$ in binary64, $2B^2=2^{49}\le2^{53}$. That is the
justification the implementation's own validation records at its binary32
totality-boundary case (\file{Hecke/HeckeFpFilterValidation.cpp}, Phase E: ``$E_{32}(2^{10})
= 0.5$ exactly --- not $<0.5$, so the totality guarantee here rests on the products/sums
at this $B$ staying exactly representable in binary32''), which is the same observation,
made in the code and not in the contract.
\textbf{Corrected in} \S\ref{sec:fpfilter}: Corollary~\ref{cor:fptotal} states the
criterion as $E(B)+\abs{\hat x-x}<1$ with the two sufficient conditions (i)
$E(B)\le\tfrac12$ and (ii) exact representability, and Corollary~\ref{cor:fp32} derives
$B\le2^{10}$ from both. Both labels are preserved and both thresholds are unchanged
(rule~\ref{L2}: what changes is the criterion, not the constant).
\textbf{What it does not affect.} Nothing computed. The constants $2^{-50}B^2$,
$2^{-21}B^2$, $2^{24}$ and $2^{10}$ are exactly what \file{Hecke/Predicates.cpp} uses and
what the contract box of \S\ref{sec:fpfilter} keeps binding; the $18$ source sites citing
\code{cor:fp32} and the $12$ citing \code{cor:fptotal} continue to assert what they
asserted.

%----------------------------------------------------------------------
\subsection{The bad-prime oracle cites a corollary that excludes the bad prime}\label{err:e23}
%----------------------------------------------------------------------
\emph{Found while Chapter~\ref{sec:example} was written (2026-09-16), in the oracle list
being re-homed. The shipped implementation is right; the contract text asserted, at
$\ell\in\{2,3\}$, an identity whose stated justification does not reach there.}

\textbf{Defective statement.} \file{hecke\_action.tex}:1604--1607, oracle~7 of the
validation plan: ``\textbf{Bad prime} $p=\ell=2$, $n\in\{0,2,10\}$: pipeline must run and
produce well-defined $U_2$ \ldots; $U_2$-eigenvalue assertions are \emph{not} made
(hypothesis (H3): $\ell=2$), only $U_4=U_2^2$ (Corollary \code{cor:ulrecursion}) and
idempotent-stability checks.''

\textbf{Why it is defective.} The oracle is careful to withhold eigenvalue assertions at
$\ell=2$ and then asserts one identity anyway, on the strength of a corollary that does
not cover it. Corollary~\ref{cor:ulrecursion}(1) --- $T_{\ell^{\nu}}=U_\ell^{\,\nu}$ for
the \emph{full} operators --- is derived from the recursion of
Theorem~\ref{thm:relations}(1), which carries $6\in\cring^{\times}$; at $\ell=2$ that
hypothesis fails and the recursion is not available, so the collapse
$p^{\,n+1}=\ell^{\,n+1}\equiv0$ has nothing to collapse. Clause~(2) of the same corollary
carries no hypothesis on $\ell$ but says something different --- that the full and the
primitive prime-power operators coincide on $\Fl$-coefficients for $n\ge1$ --- and it too
fails to give $U_4=U_2^{2}$. The statement may well be true at $\ell=2$; what is wrong is
that a numbered oracle asserts it with a citation that excludes its own setting, which is
rule~\ref{A1}'s failure mode exactly.

\textbf{Correct statement.} Split the oracle by what each clause needs, which is what
oracle~7 of \S\ref{sec:example} now does. For \emph{every} $\ell$ and $n$: the pipeline
runs to completion and certifies, and $U_\ell:=T_\ell$ is well defined
(Obligations~\ref{obl:welldef}, \ref{obl:integral}, \ref{obl:tau}, none of which carries
a hypothesis on $\ell$). For $\ell\ge5$ only: $T_{\ell^{2}}=U_\ell^{2}$ and the
rank-stabilisation check, by Corollary~\ref{cor:ulrecursion}(1). At $\ell\in\{2,3\}$ both
are \textbf{recorded}, not asserted.

\textbf{Evidence.} The implementation already draws this line and records the reason.
\code{HeckeRelationsValidation} hard-asserts $U_{\ell^{2}}=U_\ell^{2}$ and the
$\rank U_\ell^{r}$ stabilisation at $p=\ell=5$, and marks the same two quantities
``\code{[A-F1 RECORD ONLY, NOT asserted}'' at $p=\ell\in\{2,3\}$, where it asserts only
that the pipeline completes and the objects are well typed
(\file{Hecke/HeckeRelationsValidation.cpp}:273--330, 389--408); amendment A-F1 of the
implementation plan is where that gating was decided. So the gap was noticed in the code
and never fixed in the contract --- the same shape as \ref{err:e1}, and the reason
\S\ref{sec:example} is now written with the hypothesis on every clause.

\textbf{Corrected in} \S\ref{sec:example}, oracle~7, which carries the split above and
an erratum box naming this entry.
\textbf{What it does not affect.} Nothing computed, and no label moves.
Corollary~\ref{cor:ulrecursion} itself is unchanged: it was already apportioned in
v0.5.0 (Appendix~\ref{app:changelog}), and this entry is about the oracle that cited it,
not about the corollary.

%----------------------------------------------------------------------
\section{Index of statements with STATUS}\label{app:status}
%----------------------------------------------------------------------
% SOURCE:  this document itself, read mechanically on 2026-09-16: every
%          \begin{<theorem-env>} ... \label{} in gl2_contract.tex and
%          gl2_notation.tex, with the status macros occurring in the statement
%          and in the proof that immediately follows it; the provenance column
%          from Appendix E's tagged changelog lines (rule CL4).
% NOTATION: labels are printed verbatim AND dereferenced (the number in the
%          first column is a real \ref), so a label named here that this
%          document does not define is an undefined reference and therefore a
%          build failure -- the gate rule (L6) requires of this appendix.
% ERRATA:  none of its own.  Two gaps it detected in this document are
%          recorded below and are NOT errata: Appendix~A is the record of the
%          superseded corpus's defects (rule CL2), not of this one's.

\newcommand{\BNEW}{\textsf{[new]}}
\newcommand{\BRH}{\textsf{[re-homed]}}
\newcommand{\BFIX}{\textsf{[repaired]}}
\newcommand{\BSTD}{\textsf{[standard]}}
\newcommand{\Bdef}{\textsf{---}\ \emph{definition}}
\newcommand{\Bhyp}{\textsf{---}\ \emph{hypothesis}}
\newcommand{\Balg}{\textsf{---}\ \emph{procedure}}
\newcommand{\Bobl}{\textsf{---}\ \emph{predicate}}
\newcommand{\Brmk}{\textsf{---}\ \emph{commentary}}
\newcommand{\Buntag}{\textbf{\textsf{[untagged]}}}

\paragraph{What this appendix fixes.}
\S\ref{sec:status} declares that every assertion that is not a definition carries
exactly one status tag and that this appendix indexes them. It does: one line per
numbered statement of this document, giving its number, its kind, its label, the status
tags it actually carries, and --- as rule~\ref{CL4} directs --- its provenance. It is
\textbf{mechanically derived}, not curated: the status column reports the tags that
occur in the statement and in its proof, and nothing is rounded up. Where the
document's own rules are not met, the line says so rather than supplying the tag the
rule would require.

\paragraph{How to read a line.}
\begin{itemize}
\item \textbf{Status.} One or more of \PROVED, \BSTD{}, \CONJ, \COMP, \OPENSTAT, as
  \S\ref{sec:status} defines them. \textbf{A statement with several clauses may carry
  different tags on different clauses}, and many do (Claim~\ref{cl:chi} carries all
  four of \PROVED, \CONJ, \COMP{} and a \BSTD{} pointer, one per clause); this index
  does not resolve clause structure, and a multi-tag line means ``these tags occur in
  this statement'', never ``this statement is all of these''. The tags on a clause are
  at the clause.
\item \textbf{\BSTD{} carries its citation}, in the compact form
  \code{\textbackslash cite[locator]\{key\}}: the same pointer the statement makes,
  with the sentence of justification left at the statement (rule~\ref{A2} binds the
  statement, and the terseness licence of \ref{A4} applies here).
\item \textbf{\COMP{} is split, because the two kinds are not the same evidence.}
  \COMP\textsuperscript{r} means a \emph{registered} validation re-runs the check, so
  the claim is re-executable by name in the shipped tree;
  \COMP\textsuperscript{a} means an \emph{ad hoc} run recorded in the statement that
  nothing re-runs --- a sweep done once while the chapter was written. Fifteen of the
  twenty-one are \COMP\textsuperscript{r} (fourteen of nineteen until v1.8.0, which added
  \ref{rmk:integralevidence} and the ad hoc \ref{rmk:bocksteinexp}) (eight until v1.3.3, which added
  \ref{prop:bdrycosets}, \ref{prop:bdrychar} and \ref{obl:taubdry}; eleven until v1.4.0,
  which added \ref{rmk:conerestrictsmallell}; twelve until v1.5.0, which added
  \ref{rmk:coneambiguity691}, whose dimensions are measured and re-runnable since the
  rewrite of its arithmetic; thirteen until v1.6.0, which added \ref{rmk:iotapairing},
  whose dimension half is measured by \code{HeckeInvolutionValidation}) and six are
  \COMP\textsuperscript{a}; the
  six are Remarks~\ref{rmk:conestrict}, \ref{rmk:chicoincide},
  \ref{rmk:dicksontwo}, \ref{rmk:labelcensus}, \ref{rmk:bocksteinexp}(3) (since v1.8.0)
  and Claim~\ref{cl:chi}(2), and each of them says so in its own text. Neither kind upgrades to \PROVED{}
  (\S\ref{sec:status}).
\item \textbf{A dash} is not a missing tag: it is a statement that makes no assertion
  of its own. \emph{definition} and \emph{hypothesis} assert nothing;
  \emph{procedure} (an \code{alg}) and \emph{predicate} (an \code{obl}) are
  specifications, whose correctness is the statement they discharge and whose
  \emph{status} is in \S\ref{sec:obligationroster}; \emph{commentary} is a remark that
  derives from, restates or points at tagged statements it cites. \Buntag{} marks the
  one place where a dash is a defect.
\item \textbf{Provenance} (\ref{CL4}) is \BNEW, \BRH{} or \BFIX{}, collapsed to the
  statement's \emph{current} content: a statement re-homed in one version and repaired
  in a later one is \BFIX{} here, and one whose re-homed form was strictly generalised
  is \BNEW. \BRH/\BFIX{} marks the one statement whose clauses split
  (Assumption~\ref{ass:fund_domain}: \ref{FD1}--\ref{FD3} re-homed then repaired,
  \ref{FD4} new). A star marks a statement for which
  Appendix~\ref{app:changelog} carries \emph{no} provenance line at all.
\end{itemize}

\begin{contractbox}
\textbf{The column totals of this appendix and of
Appendix~\ref{app:changelog}'s provenance table do not agree, and neither is wrong.}
The changelog counts \emph{tag instances} --- a statement written in one version and
revised in another is counted in both, which is what \ref{CL3} requires of a
per-version record --- and counts container labels (\code{sec:}, \code{eq:}) named on a
tagged line. This appendix counts \emph{statements}, once each, at their current
content. The $226$ statements here split $146$ \BNEW{} (one of them starred) / $45$
\BRH{} / $34$ \BFIX{}, plus the one whose clauses split both ways ($225$ and $145$ until
v1.9.0 added one \BNEW{} line, $205$ and $125$ until v1.8.0 added twenty), against the
changelog's $142/45/38$ tag instances over a larger label set ($141/45/38$ until v1.9.0,
$121/45/38$ until v1.8.0). The two \BFIX{} figures
are the informative case and they differ for a stated reason: $38$ statements were
\BFIX{} in \emph{some} version, and $34$ of those are \BFIX{} at current content ---
three were later generalised and are now \BNEW{}, and one is the split case. Quoting
either total as the other is the misreading \ref{CL4} exists to stop.
\end{contractbox}

\paragraph{The counts.}
$434$ numbered statements across Chapters \ref{sec:framework}--\ref{sec:implementation},
\ref{sec:integral}, \ref{sec:dicksontheta}, \ref{sec:ordinary}, \ref{sec:hida}, \ref{sec:hidafull}, \ref{sec:levelN}, \ref{sec:nebentypus},
\ref{sec:saturation}, \ref{sec:isomodelell} and \ref{sec:mpmodel} (Chapters
\ref{sec:scope} and \ref{sec:notation} carry rules and tables, not numbered statements). By kind: $119$
\code{rmk}, $115$ \code{prop}, $71$ \code{def}, $64$ \code{lem}, $15$ \code{obl}, $22$ \code{thm}, $7$ \code{cl}, $5$
\code{ass}, $7$ \code{cor}, $4$ \code{alg}, $3$ \code{ex}, $2$ \code{open} (the first decided at
v1.3.0, the kind kept, rule~\ref{L1}; the second \ref{open:ordprecision}, answered at v1.12.0, the kind
kept likewise). By status, counting tag
\emph{occurrences} (so a multi-tag statement is counted once per tag): \PROVED{} $261$, \BSTD{} $116$,
\COMP{} $58$ ($29$ registered, $29$ ad hoc), \CONJ{} $19$, \OPENSTAT{} $6$. $299$ statements carry at
least one tag and $135$ carry none --- $68$ definitions, $11$ obligations, $5$
assumptions, $4$ algorithms, $1$ example, $45$ remarks, and \textbf{one corollary}.
(v1.25.0-pre.0, a branch pre-release: twenty-three lines added, \BNEW{} --- \S\ref{sec:heckeZ}'s \ref{def:heckeZ} and \ref{def:heckeZcong} (untagged), \ref{prop:heckeZspec}, \ref{prop:heckeZrat} (with \CONJ), \ref{prop:heckeZmax}, \ref{prop:heckeZmono}, \ref{prop:heckeZuct}, \ref{prop:heckeZglue}, \ref{thm:heckeZcert}, \ref{prop:heckeZcong}, \ref{prop:heckeZeiscert} (\PROVED, \BSTD), \ref{lem:heckeZbridge}, \ref{prop:heckeZgen}, \ref{prop:heckeZetale}, \ref{prop:heckeZlocal}, \ref{prop:heckeZthrough}, \ref{prop:heckeZsplit}, \ref{prop:heckeZlevelone23}, \ref{thm:heckeZdescent}, \ref{rmk:heckeZcaveat} (\PROVED) and \ref{rmk:heckeZcbone}, \ref{rmk:heckeZcbtwo}, \ref{rmk:heckeZcbthree} (\CONJ); so $411\to434$ lines, \code{def} $69\to71$, \code{lem} $63\to64$, \code{prop} $101\to115$, \code{thm} $20\to22$, \code{rmk} $115\to119$; \PROVED{} $243\to261$, \BSTD{} $107\to116$, \CONJ{} $15\to19$; tagged $278\to299$, untagged $133\to135$ (definitions $66\to68$). No tag set of an existing line moves; the statements are the last of their section, so no number moves.)
(v1.23.0-pre.33: four lines added, \BNEW{} --- \S\ref{sec:treepair}'s \ref{def:treepair} (untagged), \ref{lem:treepaircyl} (\PROVED, \BSTD), \ref{prop:treepair} (\PROVED, \BSTD) and \ref{rmk:treepairscope} (\PROVED); so $407\to411$ lines, \code{def} $68\to69$, \code{lem} $62\to63$, \code{prop} $100\to101$, \code{rmk} $114\to115$; \PROVED{} $240\to243$, \BSTD{} $105\to107$; tagged $275\to278$, untagged $132\to133$ (definitions $65\to66$). No tag set of an existing line moves; the statements are the last of their section, so no number moves.)
(v1.23.0-pre.24: five lines added, \BNEW{} --- \S\ref{sec:descenthecke}'s \ref{prop:strictliftN} (\PROVED, \BSTD), \ref{prop:crtcharpolyN} (\PROVED, \BSTD), \ref{prop:torsionhecke} (\PROVED), \ref{thm:descentcert} (\PROVED) and \ref{rmk:descentcertscope} (\PROVED); so $397\to402$ lines, \code{prop} $95\to98$, \code{thm} $19\to20$, \code{rmk} $112\to113$; \PROVED{} $231\to236$, \BSTD{} $100\to102$; tagged $266\to271$. No tag set of an existing line moves; the statements are the last of their section, so no number moves.)
(v1.23.0-pre.26: five lines added, \BNEW{} --- \S\ref{sec:descentpair}'s \ref{def:descentpair} (untagged), \ref{lem:pairlifting} (\PROVED, \BSTD), \ref{prop:descentpair} (\PROVED, \BSTD), \ref{prop:descentpairhecke} (\PROVED, \BSTD) and \ref{rmk:descentpairscope} (\PROVED); so $402\to407$ lines, \code{def} $67\to68$, \code{lem} $61\to62$, \code{prop} $98\to100$, \code{rmk} $113\to114$; \PROVED{} $236\to240$, \BSTD{} $102\to105$; tagged $271\to275$, untagged $131\to132$ (definitions $64\to65$). No tag set of an existing line moves; the statements are the last of their section, so no number moves.)
(v1.23.0-pre.29: no line added, no tag set moves --- the C-D7p referee's fixes restate \ref{prop:descentpairhecke}(1), (5)'s refusals and (1)'s proof inside the existing tags; the counts stand.)
(v1.23.0-pre.28: no line added, no tag set moves --- \ref{prop:descentpairhecke}(1), (5) restated for every prime $p\nmid N$ inside its existing tags; the counts stand.)
(v1.23.0-pre.27: no line added, no tag set moves --- the C-D7 referee's fixes sharpen \S\ref{sec:descentpair}'s statements inside their existing tags (\ref{prop:descentpair}'s locator list
corrected, VII.6 withdrawn); the counts stand.)
(v1.23.0-pre.25: no line added, no tag set moves --- the C-D4 referee's fixes sharpen \S\ref{sec:descenthecke}'s five statements inside their existing tags; the counts stand.)
(v1.23.0-pre.23: no line added, no tag set moves --- the C-D3g referee's fixes sharpen \ref{prop:descentmaps}(2)(g) and (2)(b) inside their existing tags; the counts stand.)
(v1.23.0-pre.22: no line added, no tag set moves --- \ref{prop:descentmaps} gains item (2)(g), proved by elementary row and column operations over $\Z_{(\ell)}$ inside its
existing tags; the counts stand.)
(v1.23.0-pre.21: no line added, no tag set moves --- \ref{prop:isomodeltree}(4) is made effective inside its existing tags; the counts stand.)
(v1.23.0-pre.20: two lines added, \BNEW{} --- \S\ref{sec:isomodeltree}'s \ref{def:isomodeltree} (untagged) and \ref{prop:isomodeltree} (\PROVED, \BSTD); so $395\to397$ lines, \code{def} $66\to67$, \code{prop} $94\to95$; \PROVED{} $230\to231$, \BSTD{} $99\to100$; tagged $265\to266$, untagged $130\to131$ (definitions $63\to64$). No tag set of an existing line moves; the statements are the last of their section, so no number moves.)
(v1.23.0-pre.19: no line added, no tag set moves --- \ref{prop:descentmaps} gains item (3)(b$'$), proved by the double coset formula of \cite[Ch.~III.9]{brown},
already in its tag; the counts stand.)
(v1.23.0-pre.18: one line added, \BNEW{} --- \S\ref{sec:descentN}'s \ref{prop:descentmaps} (\PROVED, \BSTD); so $394\to395$ lines, \code{prop} $93\to94$; \PROVED{} $229\to230$, \BSTD{} $98\to99$; tagged $264\to265$. No tag set of an existing line moves; the statement is the last of its section, so no number moves.)
(v1.23.0-pre.17: one line added, \BNEW{} --- \S\ref{sec:descentN}'s \ref{prop:descentNthree} (\PROVED, \BSTD); so $393\to394$ lines, \code{prop} $92\to93$; \PROVED{} $228\to229$, \BSTD{} $97\to98$; tagged $263\to264$. No tag set of an existing line moves; the statement is the last of its section, so no number moves.)
(v1.23.0-pre.16: one line added, \BNEW{} --- \ref{lem:dicksonSL2} (\PROVED, no \STD); so $392\to393$ lines, \code{lem} $60\to61$; \PROVED{} $227\to228$; tagged $262\to263$. The tag of \ref{prop:oddnltwo}(2) loses its \STD{} on \cite{wilkerson} (the proposition keeps \BSTD{} via \cite{weibel}, $97$ unchanged); \ref{prop:dicksonH0}'s statements keep their citations as provenance. \textbf{Statement numbers after \ref{cl:dimH1} in its section shift by one.})
(v1.23.0-pre.15: no line added; \ref{cl:dimH1}'s tag set gains \PROVED{} (clause (2), proved in place) and drops \cite{wilkerson}; so \PROVED{} $226\to227$, tagged
$262$ unchanged. \ref{prop:descentN}'s citation list gains \cite[Ch.~V.3]{brown}.)
(v1.23.0-pre.14: two lines added, \BNEW{} --- \S\ref{sec:descentN}'s \ref{prop:descentN} (\PROVED, \BSTD) and \ref{prop:badprimesN} (\PROVED); so $390\to392$ lines, \code{prop} $90\to92$; \PROVED{} $224\to226$, \BSTD{} $96\to97$; tagged $260\to262$. No tag set of an existing line moves.)
(v1.23.0-pre.12: one line added, \BNEW{} --- \S\ref{sec:isomodelell}'s \ref{prop:certifiedeell} (\PROVED, \BSTD); so $389\to390$ lines, \code{prop} $89\to90$;
\PROVED{} $223\to224$, \BSTD{} $95\to96$; tagged $259\to260$. No tag set of an existing line moves. \textbf{Not counted, a gap of v1.23.0-pre.3 left to the 1.23.0 cut}:
\ref{lem:precisionreduction} and \ref{obl:precisionreuse} have no line in this appendix yet.)
(v1.22.0-pre.14: one line added, \BNEW{} --- \S\ref{sec:oddnltwo}'s \ref{rmk:lvdimsltwothree} (\PROVED); so $388\to389$ lines, \code{rmk} $111\to112$; \PROVED{}
$222\to223$; tagged $258\to259$. No tag set of an existing line moves; the twelve \code{\textbackslash LEAN} markers are not counted.)
(v1.22.0-pre.12: one line added, \BNEW{} --- \S\ref{sec:six}'s \ref{rmk:treeconeres} (\PROVED); so $387\to388$ lines, \code{rmk} $110\to111$; \PROVED{}
$221\to222$; tagged $257\to258$. \ref{cl:amalgam} gains a \PROVED{} at its clause (1) and \ref{lem:bulkperiodic} none, so no tag set of an existing line moves; the seven
\code{\textbackslash LEAN} markers are not counted.)
(v1.22.0-pre.11: two lines added, \BNEW{} --- \S\ref{sec:mpmodel}'s \ref{def:mpqd} (definition, untagged) and \ref{prop:mpqreduced} (\PROVED, \BSTD, \COMP\textsuperscript{a});
so $385\to387$ lines, \code{prop} $88\to89$, \code{def} $65\to66$; \PROVED{} $220\to221$, \BSTD{} $94\to95$, \COMP{} $57\to58$ (ad hoc $28\to29$); tagged $256\to257$, untagged
$129\to130$ (definitions $62\to63$). The F-L23 repairs change no tag set.)
(v1.22.0-pre.10: eleven lines added, \BNEW{} --- \S\ref{sec:oddnltwo}'s \ref{lem:centralminusone}, \ref{lem:bocksteinallell}, \ref{lem:cyclicres},
\ref{prop:oddnltwo}, \ref{lem:oddnsplit}, \ref{lem:tellrank}, \ref{prop:ltwothreescalar} (\PROVED, \BSTD), \ref{prop:oddnhecke} (\PROVED, \CONJ), \ref{lem:cprule} (\PROVED, \BSTD),
\ref{rmk:standardtwist} (\PROVED) and \ref{rmk:oddnltwoerr} (commentary); so $374\to385$ lines, \code{lem} $54\to60$, \code{prop} $85\to88$, \code{rmk} $108\to110$; \PROVED{}
$210\to220$, \BSTD{} $86\to94$ (eight lines; \ref{prop:ltwothreescalar} counted once), \CONJ{} $14\to15$; tagged $246\to256$, untagged $128\to129$ (remarks $44\to45$). No tag set of an
existing line moves. v1.22.0-pre.9 added no line.)
(v1.22.0-pre.8: no line added; one tag set grows --- \ref{rmk:mphecke}, commentary until pre.7, gains \COMP\textsuperscript{r} (\ref{MPC12}, slot-22 MP-4), so \COMP{} $56\to57$,
registered $28\to29$, tagged $245\to246$, untagged $129\to128$ (remarks $45\to44$); no other count moves.)
(v1.22.0-pre.7: no line added; two tag sets grow --- \ref{prop:mpboundary} and \ref{prop:mprescor} gain \COMP\textsuperscript{r} (\ref{MPC10}, \ref{MPC11}, slot-22 MP-3), so
\COMP{} $54\to56$, registered $26\to28$; no other count moves.)
(v1.22.0-pre.6: no line added and no tag moves; the line of \ref{prop:mpboundary} gains the locator \cite[Ch.~III.9]{brown} of its repaired proof of (2), referee F-MP2 R2.)
(v1.22.0-pre.5: twenty-three lines added, \BNEW{} --- \S\ref{sec:mpmodel}'s \ref{def:mptri}, \ref{def:mprules}, \ref{def:mpell}, \ref{def:mpRZ},
\ref{rmk:mprulesread} (\PROVED, \COMP\textsuperscript{r}), \ref{def:mpset}, \ref{rmk:mpwhat} (\PROVED, \COMP\textsuperscript{r}), \ref{def:mpfree}, \ref{lem:mprows},
\ref{lem:mpcols} (\PROVED), \ref{lem:mpassembly} (\PROVED, \BSTD), \ref{thm:mpcomparison} (\PROVED, \BSTD), \ref{rmk:mphybrid} (\PROVED, \BSTD, \COMP\textsuperscript{r}), \ref{def:mpphi},
\ref{prop:mpalldegrees} (\PROVED, \BSTD, \COMP\textsuperscript{r}), \ref{def:mpmorphism}, \ref{thm:mpnatural} (\PROVED, \BSTD, \COMP\textsuperscript{r}), \ref{prop:mpbredon}
(\PROVED, \BSTD, \COMP\textsuperscript{r}), \ref{prop:mpboundary} and \ref{prop:mprescor} (\PROVED, \BSTD), \ref{rmk:mphecke} (commentary), \ref{ex:mpteardrop}
(\COMP\textsuperscript{r}) and \ref{rmk:mpcheckable} (\PROVED, \COMP\textsuperscript{r}); so $351\to374$ lines, \code{def} $57\to65$, \code{lem} $51\to54$, \code{thm} $17\to19$,
\code{prop} $81\to85$, \code{rmk} $103\to108$, \code{ex} $2\to3$; \PROVED{} $197\to210$, \BSTD{} $78\to86$, \COMP{} $46\to54$ (registered $18\to26$); tagged $231\to245$,
untagged $120\to129$ (definitions $54\to62$, remarks $44\to45$). No tag set of an existing line moves (the addenda to \ref{rmk:mpcaveat}, \ref{rmk:orbifoldrefine} and
\ref{rmk:isomodellmp} add no tag).)
(v1.22.0-pre.3: twenty-one lines added, \BNEW{} --- \S\ref{sec:nebles}'s \ref{def:levelpair}, \ref{lem:nebcompat}, \ref{def:levelres}, \ref{lem:levelresmorph}
(\PROVED, \BSTD), \ref{def:levelcone}, \ref{lem:levelconeexact}, \ref{prop:levelconeshapiro}, \ref{prop:modeltorsion}, \ref{prop:levelsplit},
\ref{prop:levelconehecke} (\PROVED, \BSTD, \COMP\textsuperscript{a}), \ref{def:normalpair}, \ref{prop:levelinfres} (\PROVED, \BSTD), \ref{prop:cycliccone} (\PROVED, \BSTD),
\ref{lem:levelfreeeuler}, \ref{lem:unipotentinv}, \ref{prop:gamma0normal} and \ref{prop:hidanormalsteps} (\PROVED, \COMP\textsuperscript{a}), \ref{rmk:normalN1} (\PROVED),
\ref{obl:levelconeexact} (predicate), \ref{rmk:lescheckable} (\PROVED, \COMP\textsuperscript{a}) and \ref{rmk:oldspacecone} (commentary); so
$330\to351$ lines, \code{def} $53\to57$, \code{lem} $46\to51$, \code{prop} $73\to81$, \code{obl} $14\to15$, \code{rmk} $100\to103$; \PROVED{} $182\to197$, \BSTD{} $74\to78$,
\COMP{} $42\to46$ (ad hoc $24\to28$), \OPENSTAT{} $6$ unchanged (after the F-LES repair R1, which made \ref{prop:levelconehecke}(3) and \ref{LES9} \PROVED{}); tagged $216\to231$, untagged $114\to120$ (definitions $50\to54$, obligations $10\to11$, remarks $43\to44$).
No tag set of an existing line moves.)
(v1.17.0-pre.1: one line added, \BNEW{} --- \S\ref{sec:saturation}'s \ref{rmk:sattor} (\PROVED); two tag sets grow --- \ref{rmk:nebcheckable} gains
\COMP\textsuperscript{r} (\ref{NB4}(a), \code{HeckeNebentypusValidation} arm (X4)), \ref{rmk:levelNscale} moves from commentary to \COMP\textsuperscript{a}
(slot-16 WP-13's measured fill ratios in its item~(4)); so $317\to318$ lines, \code{rmk} $96\to97$; \PROVED{} $173\to174$, \COMP{} $37\to39$ (registered
$17\to18$, ad hoc $20\to21$); tagged $205\to207$, untagged $112\to111$ (remarks $43\to42$). The \code{\textbackslash LEAN} marker of \S\ref{sec:status} is not counted.)
(v1.17.0-pre.4: \ref{NB4}(c)'s sharp form at $n=0$ tagged \PROVED{} (given one \STD{} input, locator \cite[\S1.1]{ash_stevens} not verified; owner ruling of
2026-10-02). \ref{rmk:nebcheckable}'s tag set already holds \PROVED{}, \BSTD{}, \CONJ{} and \COMP\textsuperscript{r} (the general (NB4) and (NB5) stay \CONJ), so
it is unchanged and \textbf{no count moves}: $318$ lines; \PROVED{} $174$, \BSTD{} $70$, \COMP{} $39$, \CONJ{} $13$, \OPENSTAT{} $6$; tagged $207$, untagged $111$.)
(v1.17.0-pre.5: locators only --- \ref{NB4}(c)'s input verified (\cite[Lemmas~1.1.1, 1.1.4]{ash_stevens}), \ref{lem:heckerankone}'s eigenbasis from
\cite{diamond_im}, \ref{LN5} from \cite[Th\'eor\`eme~2]{merel96}; the code-map sentences of \S\ref{sec:isomodelell}; \textbf{no tag or count moves}.)
(v1.17.0: the first \code{\textbackslash LEAN} / \code{\textbackslash LEANlit} markers, $42$ on $20$ labels; markers are not counted, \textbf{no tag or count moves}.)
(v1.21.0-pre.4: two lines added, \BNEW{} --- \S\ref{sec:levelN}'s \ref{prop:levelNcurve} (\PROVED, one \BSTD{} step with its locator
not verified, \COMP\textsuperscript{a}) and \ref{rmk:levelNcurvefamily} (\BSTD, \CONJ, \COMP\textsuperscript{a}); so $318\to320$ lines, \code{prop} $70\to71$,
\code{rmk} $97\to98$; \PROVED{} $174\to175$, \BSTD{} $70\to72$, \COMP{} $39\to41$ (ad hoc $21\to23$), \CONJ{} $13\to14$; tagged $207\to209$, untagged $111$.
No tag set of an existing line moves (\ref{rmk:levelNcheckable}'s addendum adds no tag; \ref{def:levelNmodel} stays a definition).)
(v1.22.0-pre.1: ten lines added, \BNEW{} --- \S\ref{sec:nebwildsec}'s \ref{def:nebwild}, \ref{lem:nebwild} (\PROVED, \BSTD), \ref{def:nebgammaH},
\ref{prop:nebgammaH}, \ref{prop:nebgammaHheilbronn}, \ref{lem:nebshapiro}, \ref{lem:nebwildtors} (\PROVED, \BSTD), \ref{lem:nebchainrref} (each \PROVED),
\ref{rmk:nebwildcheckable} (\PROVED; its \ref{NB8} names [U]'s orientation, \COMP{} at $r=1$, counted here as \COMP\textsuperscript{a}) and \ref{rmk:nebcosetlevel}
(commentary: the code contract); so $320\to330$ lines, \code{def} $51\to53$, \code{lem} $42\to46$, \code{prop} $71\to73$, \code{rmk} $98\to100$; \PROVED{} $175\to182$,
\BSTD{} $72\to74$, \COMP{} $41\to42$ (ad hoc $23\to24$); tagged $209\to216$, untagged $111\to114$ (definitions $48\to50$, remarks $42\to43$). No tag set of an existing
line moves (\ref{rmk:nebfield}'s sentence and the refusal box add no tag; the addendum after \ref{thm:hidacontrol} is unnumbered prose and adds none).)
(v1.16.0-pre.1: ten lines added, \BNEW{} --- \S\ref{sec:levelN}'s \ref{prop:kappamodel}, \ref{prop:degenpres}, \ref{rmk:degenpresconv},
\ref{rmk:levelNscale}, \ref{rmk:levelNcheckabletwo} and \S\ref{sec:isomodelell}'s \ref{def:isomodelell}, \ref{prop:isomodelell},
\ref{prop:ellthreeclosed}, \ref{rmk:isomodellmp}, \ref{rmk:isomodellcheckable}; so $307\to317$ lines, \code{rmk} $91\to96$, \code{prop} $66\to70$,
\code{def} $50\to51$; \PROVED{} $165\to173$, \BSTD{} $68\to70$, \COMP{} $34\to37$ (ad hoc $17\to20$: slot-16 WP-12's \file{degen.py} and
\file{negctl.py} in \ref{rmk:degenpresconv} and \ref{rmk:levelNcheckabletwo}, slot-14 WP-Q1's \file{z3check.py} in \ref{rmk:isomodellcheckable}),
\CONJ{} $12\to13$; tagged $197\to205$, untagged $110\to112$ (definitions $47\to48$, remarks $42\to43$).)
(v1.15.0-pre.1: twenty-three lines added, \BNEW{} --- \S\ref{sec:nebentypus}'s \ref{def:nebdata}, \ref{lem:nebdata},
\ref{lem:teichmuller}, \ref{rmk:nebfield}, \ref{lem:inducedchar}, \ref{lem:gammaonetable}, \ref{def:nebmodel}, \ref{prop:nebdecomp},
\ref{prop:nebcorr}, \ref{thm:nebrelations}, \ref{prop:neborbit}, \ref{rmk:nebheilbronn}, \ref{rmk:nebcost}, \ref{rmk:nebcheckable} and
\S\ref{sec:saturation}'s \ref{def:satdata}, \ref{prop:saturation}, \ref{prop:satdefect}, \ref{lem:heckerankone}, \ref{lem:freepart},
\ref{cor:stablecert}, \ref{prop:cohforms}, \ref{prop:hidaalgconst}, \ref{rmk:thetaforms}; so $284\to307$ lines, \code{rmk} $86\to91$,
\code{prop} $59\to66$, \code{def} $47\to50$, \code{lem} $36\to42$, \code{thm} $16\to17$, \code{cor} $6\to7$; \PROVED{} $148\to165$,
\BSTD{} $62\to68$, \COMP{} $33\to34$ (ad hoc $16\to17$: slot-14 WP-C9's \file{theta\_check.py} in \ref{rmk:thetaforms}), \CONJ{}
$11\to12$; tagged $179\to197$, untagged $105\to110$ (definitions $44\to47$, remarks $40\to42$). \ref{thm:hidaweight} gains a pointer
sentence (P18) and moves no tag.)
(v1.14.1: one line added, \BNEW{} --- \S\ref{sec:hidafull}'s \ref{lem:rhorank} (the $U_\ell$ lever's rank claim, owner ruling [S14] Q-H); so
$283\to284$ lines, \code{lem} $35\to36$; \PROVED{} $147\to148$; tagged $178\to179$, untagged unchanged. The \COMP\textsuperscript{a}
sentences added to \ref{Hi11} and \ref{Hi12} move no tag set.)
(v1.14.0-pre.1: twelve lines added, \BNEW{} --- \S\ref{sec:levelN}'s \ref{def:cosetrepsN}, \ref{lem:cosetrepsN},
\ref{def:levelNmodel}, \ref{lem:treeres}, \ref{lem:manin}, \ref{def:kappaN}, \ref{lem:levelduality}, \ref{lem:merelpath},
\ref{prop:heilbronn}, \ref{rmk:heilbronnconv}, \ref{rmk:heilbronncost} and \ref{rmk:levelNcheckable}; so $271\to283$ lines, \code{rmk}
$83\to86$, \code{prop} $58\to59$, \code{def} $44\to47$, \code{lem} $30\to35$; \PROVED{} $138\to147$, \BSTD{} $57\to62$, \COMP{}
$30\to33$ (ad hoc $13\to16$: WP-4's scripts in \ref{rmk:heilbronnconv}, \ref{rmk:heilbronncost}, \ref{rmk:levelNcheckable}), \CONJ{}
$10\to11$; tagged $168\to178$, untagged $103\to105$ (definitions $42\to44$). The \COMP\textsuperscript{a} sentences added to
\ref{rmk:hidacheckable2} (\ref{Hi11}, \ref{Hi12}) and \ref{rmk:bocksteinexp} move no tag set: both already carry the tag.)
(v1.13.0-pre.2: one line added, \BNEW{} --- \ref{rmk:levelscaling}, \BSTD{} (WP-3 referee M2); \ref{rmk:hidacheckable2} gains
\BSTD{} (\ref{Hi10}'s commutativity, referee m9); \ref{prop:stabilisation} moves to per-item tags (referee m1) with its tag set
unchanged; so $270\to271$ lines, \code{rmk} $82\to83$, \BSTD{} $55\to57$, tagged $167\to168$, untagged unchanged.)
(v1.13.0-pre.1: thirteen lines added, \BNEW{} --- \S\ref{sec:hidafull}'s \ref{def:levelellmodel},
\ref{prop:levelcorrespondence}, \ref{prop:levelellU}, \ref{prop:oldspace}, \ref{prop:levelorbit},
\ref{rmk:levelcost}, \ref{lem:levelelltorsion}, \ref{prop:stabilisation}, \ref{rmk:stabconvention},
\ref{thm:hidacontrol}, \ref{thm:hidaweight}, \ref{def:hidalattice} and \ref{rmk:hidacheckable2}; so $257\to270$
lines, \code{rmk} $79\to82$, \code{prop} $53\to58$, \code{def} $42\to44$, \code{lem} $29\to30$, \code{thm}
$14\to16$; \PROVED{} $130\to138$, \BSTD{} $49\to55$, \COMP{} $26\to30$ (ad hoc $9\to13$: WP-2's
\file{levelell\_check.py} in \ref{prop:oldspace}, \ref{prop:levelorbit}, \ref{prop:stabilisation} and
\ref{rmk:hidacheckable2}), \CONJ{} $9\to10$; tagged $157\to167$, untagged $100\to103$ (definitions $40\to42$,
remarks $39\to40$: \ref{rmk:levelcost} is commentary). \ref{rmk:hidacheckable}'s \ref{Hi3}, \ref{Hi4} gain a
locator note and keep their tags.)
(v1.12.0: no line added; \ref{open:ordprecision} changes tag \OPENSTAT{} $\to$ \PROVED{} (owner-signed
release, 2026-09-26; \ref{cor:ordprecision} discharges it, as \ref{open:ell2}'s tag moved at v1.3.0); so
\PROVED{} $129\to130$, \OPENSTAT{} $7\to6$; tagged/untagged unchanged, kind \code{open} unchanged.)
(v1.12.0-pre.3: two lines added, \BNEW{} --- \ref{prop:bdrytranscription} and \ref{lem:bdrytorus}, both
\PROVED{}; so $255\to257$ lines, \code{prop} $52\to53$, \code{lem} $28\to29$, \PROVED{} $127\to129$, tagged
$155\to157$. \ref{prop:bdrychar}, \ref{obl:taubdry}, \ref{cl:chi} and \ref{rmk:bdrynugeneral} gain status
notes and keep their tags: \ref{cl:chi}(2)'s \CONJ{} is now its multiset half's alone.)
(v1.12.0-pre.1: six lines added, \BNEW{} --- \S\ref{sec:hida}'s \ref{prop:torsionprepass},
\ref{cor:ordprecision}, \ref{lem:ordexponent}, \ref{rmk:ordexponentuse}, \ref{def:hidatower} and
\ref{rmk:hidacheckable}; so $249\to255$ lines, \code{rmk} $77\to79$, \code{prop} $51\to52$, \code{def}
$41\to42$, \code{lem} $27\to28$, \code{cor} $5\to6$; \PROVED{} $122\to127$, \BSTD{} $47\to49$, \COMP{}
$24\to25$ (ad hoc $7\to8$: \ref{rmk:hidacheckable}'s $\tau(2411)$), \CONJ{} $8\to9$; tagged $150\to155$,
untagged $99\to100$ (definitions $39\to40$). The revised \ref{open:ordprecision} and
\ref{rmk:ordtorsionfree} keep their tags.)
(v1.12.0-pre.2: no line added; \ref{rmk:ordexponentuse} gains \COMP\textsuperscript{a} (the WP-3 referee's
R5), so \COMP{} $25\to26$ (ad hoc $8\to9$); \ref{rmk:hidacheckable}'s new classical evidence (R8) is under its
existing \COMP\textsuperscript{a}; tagged and untagged unchanged.)
(v1.11.0-pre.1: twenty lines added, \BNEW{} --- \ref{prop:relationsmodel} and Part~VII's nineteen; so
$228\to248$ lines, \code{rmk} $71\to76$, \code{prop} $46\to51$, \code{def} $37\to41$, \code{lem}
$24\to27$, \code{thm} $13\to14$, \code{cor} $4\to5$, \code{open} $1\to2$; \PROVED{} $109\to121$,
\BSTD{} $41\to47$, \COMP{} $22\to23$ (ad hoc $6\to7$: \ref{prop:heckeH0}'s script, which its own proof
says is ad hoc), \CONJ{} $7\to8$, \OPENSTAT{} $6\to7$; tagged $135\to149$, untagged $93\to99$
(definitions $35\to39$, remarks $37\to39$). No existing line changes.)
(v1.11.1: one line added, \BNEW{} --- \ref{rmk:ordtorsionfree}, \PROVED{} (the slot-11 WP-15 referee's
ruling on F-4); so $248\to249$ lines, \code{rmk} $76\to77$, \PROVED{} $121\to122$, tagged $149\to150$;
the revised \ref{rmk:thetacheckable} and \ref{open:ordprecision} keep their tags.)
(v1.11.0: no line added; \ref{rmk:thetacheckable} gains \COMP\textsuperscript{r} --- slot-11 WP-5 and WP-6's
registered validations, its last paragraph --- so \COMP{} $23\to24$ (registered $16\to17$); tagged and
untagged unchanged. The revised \ref{obl:backend}, \ref{rmk:integralevidence}, \ref{open:ordprecision}
and \ref{prop:relationsmodel} keep their tags.)
(v1.10.0: two lines added, \BNEW{} --- \ref{rmk:strictliftbuilt}, \COMP\textsuperscript{r}, and
\ref{rmk:bredonhecke}, untagged commentary; so $226\to228$ lines, \code{rmk} $69\to71$, \COMP{}
$21\to22$ (registered $15\to16$), tagged $134\to135$, untagged $92\to93$ (remarks $36\to37$); the
revised \ref{lem:bulkperiodic}, \ref{prop:strictlift}, \ref{obl:coneexact},
\ref{def:ringsemantics}, \ref{rmk:integralevidence}, \ref{rmk:integralscope} and
\ref{rmk:mpcaveat} keep their tags.) (v1.9.0: one line added, \BNEW{} --- \ref{rmk:certifiedeperdiff}, \PROVED{}; so $225\to226$
lines, \code{rmk} $68\to69$, \PROVED{} $108\to109$, tagged $133\to134$; the revised
\ref{rmk:bocksteinexp}, \ref{def:ringsemantics}, \ref{rmk:integralevidence} and
\ref{obl:coneexact} keep their tags.) (v1.8.0, re-counted mechanically from the table: $20$ lines added, all \BNEW{} ---
\S\ref{sec:integral}'s nineteen statements and \ref{rmk:bocksteinexp}; \PROVED{}
$93\to108$, \BSTD{} $34\to41$, \COMP{} $19\to21$ with \ref{rmk:integralevidence}
registered and \ref{rmk:bocksteinexp} ad hoc, tagged/untagged $117/88\to133/92$; the
figures before this sentence were $205$ lines and the v1.6.0 counts below.) (v1.3.0:
\ref{lem:bulkperiodic} added with \PROVED{} and a \BSTD{} pointer; \ref{rmk:ell2window}
lost its \CONJ{} and \ref{open:ell2} its \OPENSTAT{}, both by the owner decision and the
lemma that discharge them; \ref{rmk:bulkperiodicscope} is commentary. v1.3.3 added three
\COMP\textsuperscript{r} tags --- \ref{prop:bdrycosets}, \ref{prop:bdrychar},
\ref{obl:taubdry} --- and left this paragraph at ``$13$ ($8$ registered)'', which is
corrected here to the table's own count. v1.4.0: \ref{alg:taubdry} (procedure, its
properties \PROVED) and \ref{rmk:bdrynugeneral} (commentary) untagged;
\ref{rmk:taubdrynoiota} \PROVED; \ref{rmk:conerestrictsmallell}
\COMP\textsuperscript{r}. v1.4.1: the split stood at $115/90$ with $11$ untagged
obligations, which counted \ref{obl:taubdry} --- \COMP\textsuperscript{r} since v1.3.3,
and counted as such in the $17$ --- among the untagged; recounted from the table,
$116/89$ with $10$. v1.5.0: \ref{rmk:coneambiguity691} gains
\COMP\textsuperscript{r} beside the \PROVED{} and \OPENSTAT{} it already carried, its
arithmetic having been rewritten from measured values, so \COMP{} goes $17\to18$ and
its registered half $12\to13$; the tagged/untagged split $116/89$ is unchanged, the
remark having been tagged before. No other status column moved and no line was added or
removed: v1.5.0 adds no numbered statement. v1.6.0: \ref{rmk:iotapairing} gains
\COMP\textsuperscript{r}, its contingent mirror-split prediction having become a measured
record, so \COMP{} goes $18\to19$ and its registered half $13\to14$; the remark was
untagged commentary before, so the tagged/untagged split goes $116/89\to117/88$ and the
untagged remarks $36\to35$. No other status column moved and no line was added or
removed: v1.6.0 adds no numbered statement either.)

\begin{contractbox}
\textbf{Two gaps in this document, found by compiling this index.} Neither is an
erratum (\ref{CL2}: Appendix~\ref{app:errata} records the superseded corpus), and
neither is repaired here, because repairing another chapter's statement from an index
is exactly the silent edit rule~\ref{V3} exists to prevent.
\begin{enumerate}
\item \textbf{Corollary~\ref{cor:ulrecursion} carries no status tag}, although
  \S\ref{sec:status} requires one of every assertion that is not a definition. Its two
  clauses are proved inline --- (1) from Theorem~\ref{thm:relations}(1) under
  $\ell\ge5$, (2) from \eqref{eq:fullassembly} with no hypothesis on $\ell$ --- so the
  tag it should carry is \PROVED; the index records that it does not carry one rather
  than supplying it. It is the only untagged statement in the document that is not a
  definition, a hypothesis, a procedure, a predicate or commentary.
\item \textbf{Remark~\ref{rmk:taurings} has no line in
  Appendix~\ref{app:changelog}}: it was introduced in v0.5.0 alongside
  Definition~\ref{def:tau}'s widening, which \emph{is} recorded and which names it, but
  it never got a line of its own and appears in no ``re-homed unchanged'' roster. Its
  provenance is \BNEW{} on the evidence of the text; the star on its line records that
  the changelog does not say so. \ref{CL1} requires one line per statement, so this is
  a gap in the record and not in the statement.
\end{enumerate}
\end{contractbox}

\paragraph{What this index is not.}
It is not a summary of what the document proves, and a reader looking for ``what should
be attacked first'' should use Appendix~\ref{app:changelog}'s roster of the $67$
statements that no superseded document and no project note states --- those are the
ones an error in which cannot be caught by comparing documents. Appendix~\ref{app:formalization}
is built on that roster for the same reason. Nor is it a coverage report: whether an
obligation is \emph{discharged} is \S\ref{sec:obligationroster}, which is a different
question from whether the statement it discharges is proved, and the two must not be
read off one another.

\begin{center}
\scriptsize
\begin{longtable}{@{}p{0.9cm}p{0.8cm}p{4.0cm}p{6.6cm}p{2.1cm}@{}}
\toprule
no. & kind & label & status & provenance \\
\midrule
\endfirsthead
\toprule
no. & kind & label & status & provenance \\
\midrule
\endhead
\multicolumn{5}{@{}l}{\textbf{\S\ref{sec:framework} Simplicial complexes with relations, and the universal cover}}\\
\midrule
\ref{ass:K_contractible} & ass & \texttt{\detokenize{ass:K_contractible}} & \Bhyp & \BRH \\
\ref{def:quotientcomplex} & def & \texttt{\detokenize{def:quotientcomplex}} & \Bdef & \BNEW \\
\ref{ass:edge_non_ident} & ass & \texttt{\detokenize{ass:edge_non_ident}} & \Bhyp & \BRH \\
\ref{lemma:pR_simplicial} & lem & \texttt{\detokenize{lemma:pR_simplicial}} & \PROVED & \BRH \\
\ref{def:bdryint} & def & \texttt{\detokenize{def:bdryint}} & \Bdef & \BNEW \\
\ref{ass:open_interior} & ass & \texttt{\detokenize{ass:open_interior}} & \Bhyp & \BRH \\
\ref{ex:S1} & ex & \texttt{\detokenize{ex:S1}} & \Buntag & \BRH \\
\ref{prop:injectivity} & prop & \texttt{\detokenize{prop:injectivity}} & \PROVED & \BFIX \\
\ref{def:saturated} & def & \texttt{\detokenize{def:saturated}} & \Bdef & \BNEW \\
\ref{prop:saturated_complexes} & prop & \texttt{\detokenize{prop:saturated_complexes}} & \PROVED & \BFIX \\
\ref{def:universal_cover} & def & \texttt{\detokenize{def:universal_cover}} & \Bdef & \BRH \\
\ref{prop:universalcover} & prop & \texttt{\detokenize{prop:universalcover}} & \BSTD{}~\cite[\S1.3, Props.~1.33--1.34]{hatcher}, \cite[\S1.3]{hatcher}, \cite[Props.~1.33--1.34]{hatcher} & \BRH \\
\ref{rmk:injectivityofi} & rmk & \texttt{\detokenize{rmk:injectivityofi}} & \Brmk & \BRH \\
\ref{def:orderedrelation} & def & \texttt{\detokenize{def:orderedrelation}} & \Bdef & \BNEW \\
\ref{def:relationdatum} & def & \texttt{\detokenize{def:relationdatum}} & \Bdef & \BNEW \\
\ref{def:kernelN} & def & \texttt{\detokenize{def:kernelN}} & \Bdef & \BNEW \\
\ref{ass:fund_domain} & ass & \texttt{\detokenize{ass:fund_domain}} & \Bhyp & \BRH/\BFIX \\
\ref{rmk:fd13centre} & rmk & \texttt{\detokenize{rmk:fd13centre}} & \Brmk & \BNEW \\
\ref{rmk:fdconsequences} & rmk & \texttt{\detokenize{rmk:fdconsequences}} & \PROVED & \BNEW \\
\ref{rmk:twosettings} & rmk & \texttt{\detokenize{rmk:twosettings}} & \Brmk & \BNEW \\
\ref{prop:fund_group_action} & prop & \texttt{\detokenize{prop:fund_group_action}} & \PROVED \BSTD{}~\cite[\S1.3]{hatcher} & \BRH \\
\ref{lem:fd4criterion} & lem & \texttt{\detokenize{lem:fd4criterion}} & \PROVED & \BNEW \\
\ref{rmk:fd4criterion} & rmk & \texttt{\detokenize{rmk:fd4criterion}} & \Brmk & \BNEW \\
\ref{thm:fund_group_generators} & thm & \texttt{\detokenize{thm:fund_group_generators}} & \PROVED & \BFIX \\
\midrule
\multicolumn{5}{@{}l}{\textbf{\S\ref{sec:localsystems} Matrix local systems and the three equivalent categories}}\\
\midrule
\ref{def:matrix_ls} & def & \texttt{\detokenize{def:matrix_ls}} & \Bdef & \BRH \\
\ref{rmk:derivedmatrices} & rmk & \texttt{\detokenize{rmk:derivedmatrices}} & \Brmk & \BNEW \\
\ref{prop:matrix_rep_equiv} & prop & \texttt{\detokenize{prop:matrix_rep_equiv}} & \PROVED & \BRH \\
\ref{thm:lcs_rep_equiv} & thm & \texttt{\detokenize{thm:lcs_rep_equiv}} & \BSTD{}~\cite[\S1.3, Thm.~1.38, and \S3.H]{hatcher} & \BRH \\
\ref{thm:rep_ls_equiv} & thm & \texttt{\detokenize{thm:rep_ls_equiv}} & \PROVED & \BRH \\
\ref{cor:threecategories} & cor & \texttt{\detokenize{cor:threecategories}} & \PROVED & \BNEW \\
\ref{rmk:lcsnotfree} & rmk & \texttt{\detokenize{rmk:lcsnotfree}} & \Brmk & \BNEW \\
\midrule
\multicolumn{5}{@{}l}{\textbf{\S\ref{sec:explicitcochaincomplex} The explicit equivariant cochain complex}}\\
\midrule
\ref{rmk:universalcovercohomology} & rmk & \texttt{\detokenize{rmk:universalcovercohomology}} & \PROVED & \BRH \\
\ref{thm:cohom_equiv} & thm & \texttt{\detokenize{thm:cohom_equiv}} & \PROVED \BSTD{}~\cite[Ch.~I.4]{brown}, \cite[Ch.~I.7.4]{brown}, \cite[Ch.~III.1]{brown}, \cite[\S3.H]{hatcher} & \BRH \\
\ref{def:naivecochains} & def & \texttt{\detokenize{def:naivecochains}} & \Bdef & \BNEW \\
\ref{def:eq_cochains} & def & \texttt{\detokenize{def:eq_cochains}} & \Bdef & \BRH \\
\ref{lem:eqgroupoid} & lem & \texttt{\detokenize{lem:eqgroupoid}} & \PROVED & \BNEW \\
\ref{ass:consistent_face_id} & ass & \texttt{\detokenize{ass:consistent_face_id}} & \Bhyp & \BRH \\
\ref{rmk:facesautomatic} & rmk & \texttt{\detokenize{rmk:facesautomatic}} & \PROVED & \BNEW \\
\ref{lemma:d_maps_C_eq} & lem & \texttt{\detokenize{lemma:d_maps_C_eq}} & \PROVED & \BRH \\
\ref{def:eqcohomology} & def & \texttt{\detokenize{def:eqcohomology}} & \Bdef & \BNEW \\
\ref{thm:explicit_complex_equiv} & thm & \texttt{\detokenize{thm:explicit_complex_equiv}} & \PROVED & \BFIX \\
\ref{rmk:psiuse} & rmk & \texttt{\detokenize{rmk:psiuse}} & \Brmk & \BNEW \\
\ref{rmk:psiinjectivelift} & rmk & \texttt{\detokenize{rmk:psiinjectivelift}} & \Brmk & \BNEW \\
\midrule
\multicolumn{5}{@{}l}{\textbf{\S\ref{sec:systematiccomputation} Systematic computation: bases, projectors, and the stored objects}}\\
\midrule
\ref{prop:ACq} & prop & \texttt{\detokenize{prop:ACq}} & \PROVED & \BNEW \\
\ref{prop:AZq} & prop & \texttt{\detokenize{prop:AZq}} & \PROVED & \BNEW \\
\ref{def:storedbases} & def & \texttt{\detokenize{def:storedbases}} & \Bdef & \BNEW \\
\ref{prop:Bqimage} & prop & \texttt{\detokenize{prop:Bqimage}} & \PROVED & \BNEW \\
\ref{prop:projector} & prop & \texttt{\detokenize{prop:projector}} & \PROVED & \BNEW \\
\ref{alg:cohomology} & alg & \texttt{\detokenize{alg:cohomology}} & \Balg & \BNEW \\
\ref{rmk:cost} & rmk & \texttt{\detokenize{rmk:cost}} & \Brmk & \BNEW \\
\midrule
\multicolumn{5}{@{}l}{\textbf{\S\ref{sec:conventions} Conventions, the two primes, and the 7-vertex model}}\\
\midrule
\ref{def:model7} & def & \texttt{\detokenize{def:model7}} & \Bdef & \BNEW \\
\ref{def:model7rel} & def & \texttt{\detokenize{def:model7rel}} & \Bdef & \BNEW \\
\ref{lem:model7coincidences} & lem & \texttt{\detokenize{lem:model7coincidences}} & \PROVED \BSTD{}~\cite[Ch.~VII, \S1.2, Thms.~1--2]{serre_cours} & \BNEW \\
\ref{prop:model7h4} & prop & \texttt{\detokenize{prop:model7h4}} & \PROVED & \BNEW \\
\ref{rmk:fd4computed} & rmk & \texttt{\detokenize{rmk:fd4computed}} & \COMP\textsuperscript{r} & \BNEW \\
\ref{rmk:orbifold} & rmk & \texttt{\detokenize{rmk:orbifold}} & \BSTD{}~\cite[\S3.6]{benson}, \cite[Ch.~III.6]{brown}, \cite[Ch.~VII, \S7]{brown} & \BFIX \\
\ref{rmk:negI} & rmk & \texttt{\detokenize{rmk:negI}} & \PROVED & \BFIX \\
\midrule
\multicolumn{5}{@{}l}{\textbf{\S\ref{sec:symconv} The $\Sym^n$ convention, its semigroup extension, and the contragredient}}\\
\midrule
\ref{def:rhon} & def & \texttt{\detokenize{def:rhon}} & \Bdef & \BRH \\
\ref{prop:symconv} & prop & \texttt{\detokenize{prop:symconv}} & \PROVED & \BRH \\
\ref{def:semigroup} & def & \texttt{\detokenize{def:semigroup}} & \Bdef & \BRH \\
\ref{lem:detsym} & lem & \texttt{\detokenize{lem:detsym}} & \PROVED & \BRH \\
\ref{def:contragredient} & def & \texttt{\detokenize{def:contragredient}} & \Bdef & \BNEW \\
\ref{prop:pairing} & prop & \texttt{\detokenize{prop:pairing}} & \PROVED & \BNEW \\
\ref{rmk:selfadjoint} & rmk & \texttt{\detokenize{rmk:selfadjoint}} & \Brmk & \BNEW \\
\ref{lem:symdual} & lem & \texttt{\detokenize{lem:symdual}} & \PROVED & \BNEW \\
\ref{rmk:divided} & rmk & \texttt{\detokenize{rmk:divided}} & \Brmk & \BNEW \\
\midrule
\multicolumn{5}{@{}l}{\textbf{\S\ref{sec:involution} The involution $\inv$}}\\
\midrule
\ref{def:iota} & def & \texttt{\detokenize{def:iota}} & \Bdef & \BNEW \\
\ref{lem:iotamodel} & lem & \texttt{\detokenize{lem:iotamodel}} & \PROVED & \BNEW \\
\ref{rmk:iotasigns} & rmk & \texttt{\detokenize{rmk:iotasigns}} & \Brmk & \BNEW \\
\ref{prop:iotacommutes} & prop & \texttt{\detokenize{prop:iotacommutes}} & \PROVED & \BNEW \\
\ref{rmk:iotatausym} & rmk & \texttt{\detokenize{rmk:iotatausym}} & \COMP\textsuperscript{r} & \BNEW \\
\ref{prop:iotasplit} & prop & \texttt{\detokenize{prop:iotasplit}} & \PROVED & \BNEW \\
\ref{rmk:iotaelltwo} & rmk & \texttt{\detokenize{rmk:iotaelltwo}} & \Brmk & \BNEW \\
\ref{rmk:iotapairing} & rmk & \texttt{\detokenize{rmk:iotapairing}} & \COMP\textsuperscript{r} & \BNEW \\
\ref{rmk:iotasign} & rmk & \texttt{\detokenize{rmk:iotasign}} & \Brmk & \BNEW \\
\ref{obl:iota} & obl & \texttt{\detokenize{obl:iota}} & \CONJ \COMP\textsuperscript{r} \BSTD{}~\cite[Ch.~VII, \S3]{serre_cours}, \cite[Ch.~8]{shimura}, \cite[Ch.~2]{cremona} & \BNEW \\
\midrule
\multicolumn{5}{@{}l}{\textbf{\S\ref{sec:heckedatum} The integral Hecke datum for $\GL(2)/\SL_2(\Z)$}}\\
\midrule
\ref{prop:cosets} & prop & \texttt{\detokenize{prop:cosets}} & \PROVED \BSTD{}~\cite[\S3.2, Lemmas~3.11--3.13]{shimura} & \BNEW \\
\ref{def:cosetreps} & def & \texttt{\detokenize{def:cosetreps}} & \Bdef & \BNEW \\
\ref{lem:cocycle} & lem & \texttt{\detokenize{lem:cocycle}} & \PROVED & \BNEW \\
\ref{def:operators} & def & \texttt{\detokenize{def:operators}} & \BSTD{}~\cite[\S3.2, Thm.~3.24]{shimura} & \BRH \\
\ref{prop:diamond} & prop & \texttt{\detokenize{prop:diamond}} & \PROVED & \BRH \\
\ref{prop:fullassembly} & prop & \texttt{\detokenize{prop:fullassembly}} & \PROVED & \BNEW \\
\ref{rmk:rawprimitive} & rmk & \texttt{\detokenize{rmk:rawprimitive}} & \Brmk & \BNEW \\
\ref{thm:relations} & thm & \texttt{\detokenize{thm:relations}} & \BSTD{}~\cite[Thm.~3.24]{shimura}, \cite[\S1]{ash_stevens}, \cite[Ch.~III.9--10]{brown}, \cite[Thm.~3.24 and \S3.2]{shimura}, \cite[Prop.~3.16]{shimura} & \BFIX \\
\ref{rmk:relcoeff} & rmk & \texttt{\detokenize{rmk:relcoeff}} & \Brmk & \BNEW \\
\ref{cor:ulrecursion} & cor & \texttt{\detokenize{cor:ulrecursion}} & \Buntag & \BFIX \\
\ref{rmk:badprimesingular} & rmk & \texttt{\detokenize{rmk:badprimesingular}} & \Brmk & \BNEW \\
\midrule
\multicolumn{5}{@{}l}{\textbf{\S\ref{sec:chainlevel} The primary definition: $\Tdc{g}$ in all degrees via a twisted chain map}}\\
\midrule
\ref{def:tau} & def & \texttt{\detokenize{def:tau}} & \Bdef & \BFIX \\
\ref{rmk:taurings} & rmk & \texttt{\detokenize{rmk:taurings}} & \Brmk & \BNEW$^{\ast}$ \\
\ref{lem:tauexists} & lem & \texttt{\detokenize{lem:tauexists}} & \PROVED \BSTD{}~\cite[Ch.~VII, \S1.2, Thm.~1(1),(3)]{serre_cours} & \BFIX \\
\ref{rmk:tauexistsscope} & rmk & \texttt{\detokenize{rmk:tauexistsscope}} & \PROVED \CONJ & \BNEW \\
\ref{rmk:tauintegral} & rmk & \texttt{\detokenize{rmk:tauintegral}} & \COMP\textsuperscript{r} & \BNEW \\
\ref{def:hecke} & def & \texttt{\detokenize{def:hecke}} & \Bdef & \BRH \\
\ref{thm:welldef} & thm & \texttt{\detokenize{thm:welldef}} & \PROVED & \BFIX \\
\ref{rmk:T3} & rmk & \texttt{\detokenize{rmk:T3}} & \Brmk & \BNEW \\
\ref{rmk:indeplevel} & rmk & \texttt{\detokenize{rmk:indeplevel}} & \Brmk & \BFIX \\
\ref{thm:classical} & thm & \texttt{\detokenize{thm:classical}} & \BSTD{}~\cite[Ch.~III.9--10]{brown}, \cite[\S1]{ash_stevens}, \cite[Ch.~2--3]{harder3}, \cite[Ch.~8]{shimura}, \cite{deng}, \cite[Ch.~III.9]{brown} & \BRH \\
\ref{rmk:whichoperator} & rmk & \texttt{\detokenize{rmk:whichoperator}} & \Brmk & \BNEW \\
\midrule
\multicolumn{5}{@{}l}{\textbf{\S\ref{sec:induced} Induced local systems and the maps $\alpha,\mu,\omega$}}\\
\midrule
\ref{def:inducedmodule} & def & \texttt{\detokenize{def:inducedmodule}} & \Bdef & \BRH \\
\ref{prop:blockmatrices} & prop & \texttt{\detokenize{prop:blockmatrices}} & \PROVED & \BRH \\
\ref{def:inducedls} & def & \texttt{\detokenize{def:inducedls}} & \Bdef & \BRH \\
\ref{prop:lhcompliant} & prop & \texttt{\detokenize{prop:lhcompliant}} & \PROVED \BSTD{}~\cite[Ch.~III.6.2 and Exercise~III.8.2]{brown} & \BFIX \\
\ref{rmk:shapirouse} & rmk & \texttt{\detokenize{rmk:shapirouse}} & \Brmk & \BNEW \\
\ref{prop:alpha} & prop & \texttt{\detokenize{prop:alpha}} & \PROVED \BSTD{}~\cite[Ch.~III.6.2 and Exercise~III.8.2]{brown} & \BRH \\
\ref{prop:omega} & prop & \texttt{\detokenize{prop:omega}} & \PROVED \BSTD{}~\cite[Ch.~III.9--10]{brown} & \BRH \\
\ref{rmk:omegaconvention} & rmk & \texttt{\detokenize{rmk:omegaconvention}} & \Brmk & \BNEW \\
\ref{rmk:nogo} & rmk & \texttt{\detokenize{rmk:nogo}} & \PROVED \BSTD{}~\cite[Ch.~VII, \S1]{serre_cours}, \cite[Ch.~8]{shimura} & \BFIX \\
\ref{def:xi} & def & \texttt{\detokenize{def:xi}} & \Bdef & \BRH \\
\ref{thm:xi} & thm & \texttt{\detokenize{thm:xi}} & \PROVED & \BFIX \\
\ref{rmk:thetachoices} & rmk & \texttt{\detokenize{rmk:thetachoices}} & \Brmk & \BNEW \\
\ref{alg:assembly} & alg & \texttt{\detokenize{alg:assembly}} & \Balg & \BFIX \\
\midrule
\multicolumn{5}{@{}l}{\textbf{\S\ref{sec:equivalence} Equivalent formulations and cross-checks}}\\
\midrule
\ref{prop:tworoutes} & prop & \texttt{\detokenize{prop:tworoutes}} & \PROVED & \BNEW \\
\ref{prop:mu} & prop & \texttt{\detokenize{prop:mu}} & \PROVED & \BNEW \\
\ref{rmk:barformula} & rmk & \texttt{\detokenize{rmk:barformula}} & \PROVED & \BFIX \\
\ref{prop:barformula} & prop & \texttt{\detokenize{prop:barformula}} & \PROVED \BSTD{}~\cite[Ch.~III.9--10]{brown}, \cite[\S1.1, p.~194]{ash_stevens}, \cite[Ch.~I.7.4]{brown} & \BFIX \\
\ref{rmk:kappaexact} & rmk & \texttt{\detokenize{rmk:kappaexact}} & \Brmk & \BNEW \\
\midrule
\multicolumn{5}{@{}l}{\textbf{\S\ref{sec:boundarydatum} The boundary datum}}\\
\midrule
\ref{rmk:bdrymodel} & rmk & \texttt{\detokenize{rmk:bdrymodel}} & \PROVED & \BNEW \\
\ref{def:bdrygroup} & def & \texttt{\detokenize{def:bdrygroup}} & \Bdef & \BRH \\
\ref{prop:bdrycohom} & prop & \texttt{\detokenize{prop:bdrycohom}} & \PROVED \BSTD{}~\cite[Ch.~VII.6]{brown} & \BNEW \\
\ref{prop:circlemodel} & prop & \texttt{\detokenize{prop:circlemodel}} & \PROVED \BSTD{}~\cite[Ch.~III.6]{brown} & \BFIX \\
\ref{rmk:bdryoddn} & rmk & \texttt{\detokenize{rmk:bdryoddn}} & \PROVED & \BNEW \\
\ref{prop:bdrycosets} & prop & \texttt{\detokenize{prop:bdrycosets}} & \PROVED \COMP\textsuperscript{r} & \BNEW \\
\ref{rmk:bdrynugeneral} & rmk & \texttt{\detokenize{rmk:bdrynugeneral}} & \Brmk & \BNEW \\
\ref{def:taubdry} & def & \texttt{\detokenize{def:taubdry}} & \Bdef & \BNEW \\
\ref{lem:bdrycomponent} & lem & \texttt{\detokenize{lem:bdrycomponent}} & \PROVED & \BNEW \\
\ref{lem:taubdryexists} & lem & \texttt{\detokenize{lem:taubdryexists}} & \PROVED & \BFIX \\
\ref{rmk:taubdryintegral} & rmk & \texttt{\detokenize{rmk:taubdryintegral}} & \Brmk & \BNEW \\
\ref{alg:taubdry} & alg & \texttt{\detokenize{alg:taubdry}} & \Balg & \BNEW \\
\ref{rmk:taubdrynoiota} & rmk & \texttt{\detokenize{rmk:taubdrynoiota}} & \PROVED & \BNEW \\
\ref{def:bdryop} & def & \texttt{\detokenize{def:bdryop}} & \Bdef & \BNEW \\
\ref{prop:bdryop} & prop & \texttt{\detokenize{prop:bdryop}} & \PROVED & \BNEW \\
\ref{prop:bdrychar} & prop & \texttt{\detokenize{prop:bdrychar}} & \PROVED \COMP\textsuperscript{r} & \BFIX \\
\ref{prop:bdrytranscription} & prop & \texttt{\detokenize{prop:bdrytranscription}} & \PROVED & \BNEW \\
\ref{lem:bdrytorus} & lem & \texttt{\detokenize{lem:bdrytorus}} & \PROVED & \BNEW \\
\ref{obl:taubdry} & obl & \texttt{\detokenize{obl:taubdry}} & \Bobl \COMP\textsuperscript{r} & \BNEW \\
\ref{lem:ell2unbounded} & lem & \texttt{\detokenize{lem:ell2unbounded}} & \PROVED & \BNEW \\
\ref{rmk:ell2window} & rmk & \texttt{\detokenize{rmk:ell2window}} & \Brmk & \BNEW \\
\ref{lem:bulkperiodic} & lem & \texttt{\detokenize{lem:bulkperiodic}} & \PROVED \BSTD{}~\cite[Ch.~V, \S3]{brown}, \cite[Ch.~III, \S\S5--6, 8--9, 10.2]{brown}, \cite[Ch.~I, \S4.3]{serre_trees} & \BNEW \\
\ref{rmk:bulkperiodicscope} & rmk & \texttt{\detokenize{rmk:bulkperiodicscope}} & \Brmk & \BNEW \\
\ref{open:ell2} & open & \texttt{\detokenize{open:ell2}} & \PROVED & \BNEW \\
\midrule
\multicolumn{5}{@{}l}{\textbf{\S\ref{sec:cone} The pair homotopy and the Hecke action on the mapping cone}}\\
\midrule
\ref{def:pair} & def & \texttt{\detokenize{def:pair}} & \Bdef & \BNEW \\
\ref{lem:conerelative} & lem & \texttt{\detokenize{lem:conerelative}} & \PROVED \BSTD{}~\cite[Thm.~1.3.1]{weibel}, \cite[\S1.5]{weibel} & \BNEW \\
\ref{def:cosettransport} & def & \texttt{\detokenize{def:cosettransport}} & \Bdef & \BNEW \\
\ref{lem:cosettransport} & lem & \texttt{\detokenize{lem:cosettransport}} & \PROVED & \BNEW \\
\ref{def:pairhomotopy} & def & \texttt{\detokenize{def:pairhomotopy}} & \Bdef & \BNEW \\
\ref{lem:etaexists} & lem & \texttt{\detokenize{lem:etaexists}} & \PROVED & \BFIX \\
\ref{lem:etaunique} & lem & \texttt{\detokenize{lem:etaunique}} & \PROVED & \BNEW \\
\ref{lem:conesign} & lem & \texttt{\detokenize{lem:conesign}} & \PROVED & \BFIX \\
\ref{def:coneaction} & def & \texttt{\detokenize{def:coneaction}} & \Bdef & \BNEW \\
\ref{lem:derivedhomotopy} & lem & \texttt{\detokenize{lem:derivedhomotopy}} & \PROVED & \BNEW \\
\ref{thm:conewelldef} & thm & \texttt{\detokenize{thm:conewelldef}} & \PROVED \OPENSTAT \BSTD{}~\cite[\S\S1--2]{harder_eisenstein}, \cite{borel_serre}, \cite[Thm.~1.3.1]{weibel} & \BNEW \\
\ref{prop:coneindep} & prop & \texttt{\detokenize{prop:coneindep}} & \PROVED & \BNEW \\
\ref{prop:strictpair} & prop & \texttt{\detokenize{prop:strictpair}} & \PROVED \BSTD{}~\cite[Ch.~VII, \S1.2, Thm.~1(1),(3)]{serre_cours} & \BNEW \\
\ref{rmk:conesteinberg} & rmk & \texttt{\detokenize{rmk:conesteinberg}} & \PROVED \CONJ \OPENSTAT \BSTD{}~\cite[Ch.~VII.7]{brown}, \cite[\S\S10.7, 5.7]{weibel}, \cite[Thm.~3.24]{shimura}, \cite[\S1]{ash_stevens}, \cite[\S\S1--2]{harder_eisenstein}, \cite{borel_serre} & \BNEW \\
\ref{rmk:conestrict} & rmk & \texttt{\detokenize{rmk:conestrict}} & \PROVED \CONJ \COMP\textsuperscript{a} & \BNEW \\
\ref{prop:coneambiguity} & prop & \texttt{\detokenize{prop:coneambiguity}} & \PROVED & \BFIX \\
\ref{rmk:coneambiguity691} & rmk & \texttt{\detokenize{rmk:coneambiguity691}} & \PROVED \COMP\textsuperscript{r} \OPENSTAT & \BNEW \\
\ref{prop:coneintegral} & prop & \texttt{\detokenize{prop:coneintegral}} & \PROVED \BSTD{}~\cite[\S5.4]{weibel}, \cite[Thm.~1.3.1]{weibel} & \BNEW \\
\ref{rmk:conetorsion} & rmk & \texttt{\detokenize{rmk:conetorsion}} & \Brmk & \BNEW \\
\ref{prop:coneroutes} & prop & \texttt{\detokenize{prop:coneroutes}} & \PROVED \OPENSTAT & \BNEW \\
\ref{obl:eta} & obl & \texttt{\detokenize{obl:eta}} & \Bobl & \BNEW \\
\ref{obl:coneexact} & obl & \texttt{\detokenize{obl:coneexact}} & \Bobl & \BNEW \\
\ref{rmk:conerestrictsmallell} & rmk & \texttt{\detokenize{rmk:conerestrictsmallell}} & \COMP\textsuperscript{r} & \BNEW \\
\midrule
\multicolumn{5}{@{}l}{\textbf{\S\ref{sec:six} The six groups and the long exact sequence}}\\
\midrule
\ref{rmk:sixscope} & rmk & \texttt{\detokenize{rmk:sixscope}} & \Brmk & \BNEW \\
\ref{cl:amalgam} & cl & \texttt{\detokenize{cl:amalgam}} & \PROVED \BSTD{}~\cite[Ch.~I, \S4]{serre_trees}, \cite[Ch.~III.6]{brown}, \cite[Ch.~III.10.2]{brown}, \cite[Ch.~III.9--10]{brown} & \BNEW \\
\ref{rmk:treeconeres} & rmk & \texttt{\detokenize{rmk:treeconeres}} & \PROVED & \BNEW \\
\ref{lem:h2torsion} & lem & \texttt{\detokenize{lem:h2torsion}} & \PROVED \BSTD{}~\cite[Ch.~VII.6]{brown} & \BNEW \\
\ref{cl:dimH1} & cl & \texttt{\detokenize{cl:dimH1}} & \PROVED \BSTD{}~\cite[Ch.~VII, \S3]{serre_cours}, \cite[Ch.~8]{shimura}, \cite[Ch.~2]{cremona}, \cite{deng} & \BRH \\
\ref{lem:dicksonSL2} & lem & \texttt{\detokenize{lem:dicksonSL2}} & \PROVED & \BNEW \\
\ref{cl:bockstein} & cl & \texttt{\detokenize{cl:bockstein}} & \PROVED \BSTD{}~\cite[Ch.~III, \S\S0--1]{brown} & \BRH \\
\ref{rmk:bocksteinexp} & rmk & \texttt{\detokenize{rmk:bocksteinexp}} & \PROVED \COMP\textsuperscript{a} \BSTD{}~\cite[Thm.~1.3.1]{weibel}, \cite[Thm.~4.5]{ccs13} (v1.24.0-rc.1) & \BNEW \\
\ref{lem:centralminusone} & lem & \texttt{\detokenize{lem:centralminusone}} & \PROVED \BSTD{}~\cite[Ch.~III, \S8]{brown} & \BNEW \\
\ref{lem:bocksteinallell} & lem & \texttt{\detokenize{lem:bocksteinallell}} & \PROVED \BSTD{}~\cite[Ch.~III, \S\S0--1]{brown} & \BNEW \\
\ref{lem:cyclicres} & lem & \texttt{\detokenize{lem:cyclicres}} & \PROVED \BSTD{}~\cite[Ch.~I.6, Ch.~III.1]{brown} & \BNEW \\
\ref{prop:oddnltwo} & prop & \texttt{\detokenize{prop:oddnltwo}} & \PROVED \BSTD{}~\cite[Thm.~1.3.1]{weibel} & \BNEW \\
\ref{lem:oddnsplit} & lem & \texttt{\detokenize{lem:oddnsplit}} & \PROVED \BSTD{}~\cite[Thm.~1.3.1]{weibel} & \BNEW \\
\ref{lem:tellrank} & lem & \texttt{\detokenize{lem:tellrank}} & \PROVED \BSTD{}~\cite[Ch.~III.6]{brown} & \BNEW \\
\ref{prop:ltwothreescalar} & prop & \texttt{\detokenize{prop:ltwothreescalar}} & \PROVED \BSTD{}~\cite[Ch.~I, \S4]{serre_trees}, \cite[\S1.6]{shimura}, \cite[Ch.~III, \S9]{brown}, \cite[Prop.~III.9.5(ii)]{brown}, \cite[Ch.~VII.6]{brown} & \BNEW \\
\ref{prop:oddnhecke} & prop & \texttt{\detokenize{prop:oddnhecke}} & \PROVED \CONJ & \BNEW \\
\ref{lem:cprule} & lem & \texttt{\detokenize{lem:cprule}} & \PROVED \BSTD{}~\cite[Ch.~2--3]{harder3} & \BNEW \\
\ref{rmk:standardtwist} & rmk & \texttt{\detokenize{rmk:standardtwist}} & \PROVED & \BNEW \\
\ref{rmk:oddnltwoerr} & rmk & \texttt{\detokenize{rmk:oddnltwoerr}} & \Brmk & \BNEW \\
\ref{rmk:lvdimsltwothree} & rmk & \texttt{\detokenize{rmk:lvdimsltwothree}} & \PROVED & \BNEW \\
\ref{cl:boundary} & cl & \texttt{\detokenize{cl:boundary}} & \PROVED & \BNEW \\
\ref{rmk:bdryrepairs} & rmk & \texttt{\detokenize{rmk:bdryrepairs}} & \Brmk & \BNEW \\
\ref{cl:les} & cl & \texttt{\detokenize{cl:les}} & \PROVED \CONJ & \BNEW \\
\ref{cl:chi} & cl & \texttt{\detokenize{cl:chi}} & \PROVED \CONJ \COMP\textsuperscript{a} \BSTD{}~\cite[Ch.~VI, \S4]{serre_cours} & \BNEW \\
\ref{rmk:chicoincide} & rmk & \texttt{\detokenize{rmk:chicoincide}} & \COMP\textsuperscript{a} & \BNEW \\
\ref{cl:pd} & cl & \texttt{\detokenize{cl:pd}} & \OPENSTAT & \BFIX \\
\ref{rmk:pdscope} & rmk & \texttt{\detokenize{rmk:pdscope}} & \OPENSTAT & \BNEW \\
\midrule
\multicolumn{5}{@{}l}{\textbf{\S\ref{sec:orbifoldmodel} The general orbifold model, and the torsion primes}}\\
\midrule
\ref{def:borelmodel} & def & \texttt{\detokenize{def:borelmodel}} & \PROVED & \BNEW \\
\ref{rmk:borelfree} & rmk & \texttt{\detokenize{rmk:borelfree}} & \PROVED \BSTD{}~\cite[Thm.~2.2.6]{weibel}, \cite[Ch.~I.7.4]{brown} & \BNEW \\
\ref{thm:specialisation} & thm & \texttt{\detokenize{thm:specialisation}} & \PROVED \BSTD{}~\cite[Ch.~VII, \S7, (7.10)]{brown}, \cite[Ch.~III.10.2]{brown} & \BNEW \\
\ref{rmk:mpcaveat} & rmk & \texttt{\detokenize{rmk:mpcaveat}} & \CONJ & \BNEW \\
\ref{rmk:orbifoldrefine} & rmk & \texttt{\detokenize{rmk:orbifoldrefine}} & \Brmk & \BNEW \\
\midrule
\multicolumn{5}{@{}l}{\textbf{\S\ref{sec:reduction} The combinatorial reduction primitive and integer bounds}}\\
\midrule
\ref{def:primitives} & def & \texttt{\detokenize{def:primitives}} & \Bdef & \BRH \\
\ref{rmk:primitiveclosure} & rmk & \texttt{\detokenize{rmk:primitiveclosure}} & \Brmk & \BNEW \\
\ref{alg:tau} & alg & \texttt{\detokenize{alg:tau}} & \Balg & \BFIX \\
\ref{rmk:certoverGeometry} & rmk & \texttt{\detokenize{rmk:certoverGeometry}} & \Brmk & \BNEW \\
\ref{lem:bounds} & lem & \texttt{\detokenize{lem:bounds}} & \PROVED & \BFIX \\
\ref{rmk:btabsmall} & rmk & \texttt{\detokenize{rmk:btabsmall}} & \COMP\textsuperscript{r} & \BNEW \\
\ref{thm:backend} & thm & \texttt{\detokenize{thm:backend}} & \PROVED & \BFIX \\
\midrule
\multicolumn{5}{@{}l}{\textbf{\S\ref{sec:fpfilter} The certified floating-point filter}}\\
\midrule
\ref{lem:predshapes} & lem & \texttt{\detokenize{lem:predshapes}} & \PROVED & \BNEW \\
\ref{thm:fp64} & thm & \texttt{\detokenize{thm:fp64}} & \PROVED & \BFIX \\
\ref{rmk:thresholdrounding} & rmk & \texttt{\detokenize{rmk:thresholdrounding}} & \PROVED & \BNEW \\
\ref{cor:fptotal} & cor & \texttt{\detokenize{cor:fptotal}} & \PROVED & \BFIX \\
\ref{cor:fp32} & cor & \texttt{\detokenize{cor:fp32}} & \PROVED & \BFIX \\
\midrule
\multicolumn{5}{@{}l}{\textbf{\S\ref{sec:obligations} Correctness obligations}}\\
\midrule
\ref{obl:welldef} & obl & \texttt{\detokenize{obl:welldef}} & \Bobl & \BRH \\
\ref{obl:indep} & obl & \texttt{\detokenize{obl:indep}} & \Bobl & \BFIX \\
\ref{obl:integral} & obl & \texttt{\detokenize{obl:integral}} & \Bobl & \BRH \\
\ref{obl:algebra} & obl & \texttt{\detokenize{obl:algebra}} & \BSTD{}~\cite[Thm.~3.24, \S3.2--3.3]{shimura}, \cite[\S1]{ash_stevens}, \cite[Ch.~III.9--10]{brown} & \BFIX \\
\ref{obl:agreement} & obl & \texttt{\detokenize{obl:agreement}} & \Bobl & \BRH \\
\ref{obl:shapiro} & obl & \texttt{\detokenize{obl:shapiro}} & \Bobl & \BRH \\
\ref{obl:oracle} & obl & \texttt{\detokenize{obl:oracle}} & \BSTD{}~\cite[Ch.~8]{shimura}, \cite[Ch.~2--3]{harder3}, \cite{deng} & \BFIX \\
\ref{rmk:dickson} & rmk & \texttt{\detokenize{rmk:dickson}} & \PROVED \COMP\textsuperscript{r} & \BRH \\
\ref{obl:fp} & obl & \texttt{\detokenize{obl:fp}} & \Bobl & \BRH \\
\ref{obl:backend} & obl & \texttt{\detokenize{obl:backend}} & \Bobl & \BRH \\
\ref{obl:tau} & obl & \texttt{\detokenize{obl:tau}} & \Bobl & \BRH \\
\midrule
\multicolumn{5}{@{}l}{\textbf{\S\ref{sec:example} Worked micro-example ($p=2$) and the numerical oracles}}\\
\midrule
\ref{prop:t2cocycle} & prop & \texttt{\detokenize{prop:t2cocycle}} & \PROVED \COMP\textsuperscript{r} & \BRH \\
\ref{ex:t2tau} & ex & \texttt{\detokenize{ex:t2tau}} & \COMP\textsuperscript{r} & \BNEW \\
\ref{rmk:dicksontwo} & rmk & \texttt{\detokenize{rmk:dicksontwo}} & \COMP\textsuperscript{a} & \BNEW \\
\midrule
\multicolumn{5}{@{}l}{\textbf{\S\ref{sec:implementation} Implementation notes: mapping to the code base}}\\
\midrule
\ref{rmk:labelcensus} & rmk & \texttt{\detokenize{rmk:labelcensus}} & \COMP\textsuperscript{a} & \BNEW \\
\midrule
\multicolumn{5}{@{}l}{\textbf{\S\ref{sec:integral} The integral line (v1.8.0)}}\\
\midrule
\ref{def:ringinforce} & def & \texttt{\detokenize{def:ringinforce}} & \Bdef & \BNEW \\
\ref{prop:ringinforce} & prop & \texttt{\detokenize{prop:ringinforce}} & \PROVED & \BNEW \\
\ref{prop:bredon} & prop & \texttt{\detokenize{prop:bredon}} & \PROVED \BSTD{}~\cite[Thm.~2.2.6]{weibel} & \BNEW \\
\ref{def:isomodel} & def & \texttt{\detokenize{def:isomodel}} & \Bdef & \BNEW \\
\ref{prop:isomodel} & prop & \texttt{\detokenize{prop:isomodel}} & \PROVED \BSTD{}~\cite[Thm.~2.2.6]{weibel}, \cite[Ch.~III.6]{brown}, \cite[Ch.~III.10.2]{brown} & \BNEW \\
\ref{rmk:cafive} & rmk & \texttt{\detokenize{rmk:cafive}} & \PROVED & \BNEW \\
\ref{prop:fracture} & prop & \texttt{\detokenize{prop:fracture}} & \PROVED \BSTD{}~\cite[\S1.5]{weibel} & \BNEW \\
\ref{prop:cohglue} & prop & \texttt{\detokenize{prop:cohglue}} & \PROVED & \BNEW \\
\ref{prop:descent} & prop & \texttt{\detokenize{prop:descent}} & \PROVED & \BNEW \\
\ref{prop:homvanish} & prop & \texttt{\detokenize{prop:homvanish}} & \PROVED \BSTD{}~\cite[Ch.~III, \S8]{brown}, \cite[\S10.7]{weibel} & \BNEW \\
\ref{rmk:truncdevice} & rmk & \texttt{\detokenize{rmk:truncdevice}} & \PROVED & \BNEW \\
\ref{lem:strictlifting} & lem & \texttt{\detokenize{lem:strictlifting}} & \PROVED & \BNEW \\
\ref{prop:strictlift} & prop & \texttt{\detokenize{prop:strictlift}} & \PROVED \BSTD{}~\cite[Thm.~2.2.6]{weibel} & \BNEW \\
\ref{rmk:strictliftbuilt} & rmk & \texttt{\detokenize{rmk:strictliftbuilt}} & \COMP\textsuperscript{r} & \BNEW \\
\ref{rmk:bredonhecke} & rmk & \texttt{\detokenize{rmk:bredonhecke}} & \Brmk & \BNEW \\
\ref{prop:badprimes} & prop & \texttt{\detokenize{prop:badprimes}} & \PROVED & \BNEW \\
\ref{prop:certifiede} & prop & \texttt{\detokenize{prop:certifiede}} & \PROVED \BSTD{}~\cite[\S3.6]{weibel} & \BNEW \\
\ref{rmk:certifiedeperdiff} & rmk & \texttt{\detokenize{rmk:certifiedeperdiff}} & \PROVED & \BNEW \\
\ref{prop:crtcharpoly} & prop & \texttt{\detokenize{prop:crtcharpoly}} & \PROVED & \BNEW \\
\ref{def:ringsemantics} & def & \texttt{\detokenize{def:ringsemantics}} & \Bdef & \BNEW \\
\ref{rmk:integralevidence} & rmk & \texttt{\detokenize{rmk:integralevidence}} & \COMP\textsuperscript{r} & \BNEW \\
\ref{rmk:integralscope} & rmk & \texttt{\detokenize{rmk:integralscope}} & \Brmk & \BNEW \\
\ref{prop:relationsmodel} & prop & \texttt{\detokenize{prop:relationsmodel}} & \PROVED \BSTD{}~\cite[Thm.~3.24]{shimura}, \cite[Thm.~1]{rhie_whaples}, \cite[\S1]{ash_stevens}, \cite[Ch.~III.9--10]{brown} & \BNEW \\
\midrule
\multicolumn{5}{@{}l}{\textbf{\S\ref{sec:dicksontheta} The Dickson invariant and $\thL$ (v1.11.0-pre.1)}}\\
\midrule
\ref{def:dicksoninv} & def & \texttt{\detokenize{def:dicksoninv}} & \Bdef & \BNEW \\
\ref{lem:dicksontwist} & lem & \texttt{\detokenize{lem:dicksontwist}} & \PROVED & \BNEW \\
\ref{prop:dicksonH0} & prop & \texttt{\detokenize{prop:dicksonH0}} & \PROVED \BSTD{}~\cite[\S\S2--3]{dickson1911}, \cite[\S1]{wilkerson} & \BNEW \\
\ref{prop:heckeH0} & prop & \texttt{\detokenize{prop:heckeH0}} & \PROVED \BSTD{}~\cite[\S XXI]{lucas1878} \COMP\textsuperscript{a} & \BNEW \\
\ref{rmk:heckeH0plan} & rmk & \texttt{\detokenize{rmk:heckeH0plan}} & \Brmk & \BNEW \\
\ref{lem:heckefibre} & lem & \texttt{\detokenize{lem:heckefibre}} & \PROVED & \BNEW \\
\ref{def:Qn} & def & \texttt{\detokenize{def:Qn}} & \Bdef & \BNEW \\
\ref{lem:Qnfunctions} & lem & \texttt{\detokenize{lem:Qnfunctions}} & \PROVED & \BNEW \\
\ref{def:theta} & def & \texttt{\detokenize{def:theta}} & \Bdef & \BNEW \\
\ref{prop:thetatwist} & prop & \texttt{\detokenize{prop:thetatwist}} & \PROVED & \BNEW \\
\ref{thm:thetales} & thm & \texttt{\detokenize{thm:thetales}} & \PROVED \BSTD{}~\cite[Thm.~1.3.1]{weibel}, \cite[Ch.~8]{shimura}, \cite{deng} & \BNEW \\
\ref{rmk:thetasmallell} & rmk & \texttt{\detokenize{rmk:thetasmallell}} & \Brmk & \BNEW \\
\ref{rmk:thetaclassical} & rmk & \texttt{\detokenize{rmk:thetaclassical}} & \BSTD{}~\cite{serre}, \cite[\S3]{swinnerton_dyer}, \cite{katz1977}, \cite[Thm.~4.1]{jochnowitz}, \cite[Lemma~2.1]{ash_stevens}, \cite[\S3]{edixhoven}, \cite{jochnowitz_local}, \cite{gross}, \cite{coleman_voloch} & \BNEW \\
\ref{rmk:thetacheckable} & rmk & \texttt{\detokenize{rmk:thetacheckable}} & \PROVED \BSTD{}~\cite[\S1]{swinnerton_dyer} \CONJ \COMP\textsuperscript{r} & \BNEW \\
\midrule
\multicolumn{5}{@{}l}{\textbf{\S\ref{sec:ordinary} The ordinary projector (v1.11.0-pre.1)}}\\
\midrule
\ref{def:ordinary} & def & \texttt{\detokenize{def:ordinary}} & \Bdef & \BNEW \\
\ref{prop:ordinary} & prop & \texttt{\detokenize{prop:ordinary}} & \PROVED & \BNEW \\
\ref{cor:thetanonordinary} & cor & \texttt{\detokenize{cor:thetanonordinary}} & \PROVED & \BNEW \\
\ref{rmk:ordtorsionfree} & rmk & \texttt{\detokenize{rmk:ordtorsionfree}} & \PROVED & \BNEW \\
\ref{rmk:ordinaryeis} & rmk & \texttt{\detokenize{rmk:ordinaryeis}} & \PROVED & \BNEW \\
\ref{open:ordprecision} & open & \texttt{\detokenize{open:ordprecision}} & \PROVED & \BNEW \\
\midrule
\multicolumn{5}{@{}l}{\textbf{\S\ref{sec:hida} Hida theory at level one: foundation (v1.12.0-pre.1)}}\\
\midrule
\ref{prop:torsionprepass} & prop & \texttt{\detokenize{prop:torsionprepass}} & \PROVED \BSTD{}~\cite[Ch.~III, \S\S0--1]{brown} & \BNEW \\
\ref{cor:ordprecision} & cor & \texttt{\detokenize{cor:ordprecision}} & \PROVED & \BNEW \\
\ref{lem:ordexponent} & lem & \texttt{\detokenize{lem:ordexponent}} & \PROVED & \BNEW \\
\ref{rmk:ordexponentuse} & rmk & \texttt{\detokenize{rmk:ordexponentuse}} & \PROVED \COMP\textsuperscript{a} & \BNEW \\
\ref{def:hidatower} & def & \texttt{\detokenize{def:hidatower}} & \Bdef & \BNEW \\
\ref{rmk:hidacheckable} & rmk & \texttt{\detokenize{rmk:hidacheckable}} & \PROVED \BSTD{}~\cite[Ch.~2--3]{harder3}, \cite[Ch.~8]{shimura} \CONJ \COMP\textsuperscript{a} & \BNEW \\
\midrule
\multicolumn{5}{@{}l}{\textbf{\S\ref{sec:hidafull} Hida theory at level one: control, level $\ell$, the lattice (v1.13.0-pre.1; \ref{rmk:levelscaling} pre.2; \ref{lem:rhorank} v1.14.1)}}\\
\midrule
\ref{def:levelellmodel} & def & \texttt{\detokenize{def:levelellmodel}} & \Bdef & \BNEW \\
\ref{prop:levelcorrespondence} & prop & \texttt{\detokenize{prop:levelcorrespondence}} & \PROVED \BSTD{}~\cite[Ch.~III.9--10]{brown} & \BNEW \\
\ref{prop:levelellU} & prop & \texttt{\detokenize{prop:levelellU}} & \PROVED & \BNEW \\
\ref{prop:oldspace} & prop & \texttt{\detokenize{prop:oldspace}} & \PROVED \COMP\textsuperscript{a} & \BNEW \\
\ref{prop:levelorbit} & prop & \texttt{\detokenize{prop:levelorbit}} & \PROVED \BSTD{}~\cite[Ch.~III.6.2]{brown} \COMP\textsuperscript{a} & \BNEW \\
\ref{rmk:levelcost} & rmk & \texttt{\detokenize{rmk:levelcost}} & \Brmk & \BNEW \\
\ref{rmk:levelscaling} & rmk & \texttt{\detokenize{rmk:levelscaling}} & \BSTD{}~\cite{manin}, \cite[\S2]{merel}, \cite[Ch.~3, 8]{stein}, \cite[Ch.~2]{cremona} & \BNEW \\
\ref{lem:levelelltorsion} & lem & \texttt{\detokenize{lem:levelelltorsion}} & \PROVED \BSTD{}~\cite[Ch.~III.6.2]{brown} & \BNEW \\
\ref{prop:stabilisation} & prop & \texttt{\detokenize{prop:stabilisation}} & \PROVED \BSTD{}~\cite[Prop.~4.7, Lemma~4.6]{hida_invent}, \cite{shimura,atkin_lehner} \COMP\textsuperscript{a} & \BNEW \\
\ref{rmk:stabconvention} & rmk & \texttt{\detokenize{rmk:stabconvention}} & \PROVED & \BNEW \\
\ref{thm:hidacontrol} & thm & \texttt{\detokenize{thm:hidacontrol}} & \BSTD{}~\cite[Thm.~3.1, Thm.~6.6]{hida_invent}, \cite[Thm.~3.1, Cor.~3.2]{hida_ens} & \BNEW \\
\ref{thm:hidaweight} & thm & \texttt{\detokenize{thm:hidaweight}} & \BSTD{}~\cite[Thm.~3.1]{hida_ens}, \cite[Lemma~4.6, Prop.~4.7]{hida_invent} & \BNEW \\
\ref{def:hidalattice} & def & \texttt{\detokenize{def:hidalattice}} & \Bdef & \BNEW \\
\ref{rmk:hidacheckable2} & rmk & \texttt{\detokenize{rmk:hidacheckable2}} & \PROVED \BSTD{}~\cite[Ch.~3]{shimura} \CONJ \COMP\textsuperscript{a} & \BNEW \\
\ref{lem:rhorank} & lem & \texttt{\detokenize{lem:rhorank}} & \PROVED & \BNEW \\
\midrule
\multicolumn{5}{@{}l}{\textbf{\S\ref{sec:levelN} Level $N$: the composite coset table, the cokernel presentation and its Hecke action (v1.14.0-pre.1)}}\\
\midrule
\ref{def:cosetrepsN} & def & \texttt{\detokenize{def:cosetrepsN}} & \Bdef & \BNEW \\
\ref{lem:cosetrepsN} & lem & \texttt{\detokenize{lem:cosetrepsN}} & \PROVED & \BNEW \\
\ref{def:levelNmodel} & def & \texttt{\detokenize{def:levelNmodel}} & \Bdef & \BNEW \\
\ref{lem:treeres} & lem & \texttt{\detokenize{lem:treeres}} & \PROVED & \BNEW \\
\ref{lem:manin} & lem & \texttt{\detokenize{lem:manin}} & \PROVED \BSTD{}~\cite{manin} & \BNEW \\
\ref{def:kappaN} & def & \texttt{\detokenize{def:kappaN}} & \PROVED & \BNEW \\
\ref{lem:levelduality} & lem & \texttt{\detokenize{lem:levelduality}} & \PROVED \BSTD{}~\cite[Ch.~III.8]{brown} & \BNEW \\
\ref{lem:merelpath} & lem & \texttt{\detokenize{lem:merelpath}} & \PROVED \BSTD{}~\cite[\S2]{merel} & \BNEW \\
\ref{prop:heilbronn} & prop & \texttt{\detokenize{prop:heilbronn}} & \PROVED & \BNEW \\
\ref{rmk:heilbronnconv} & rmk & \texttt{\detokenize{rmk:heilbronnconv}} & \PROVED \COMP\textsuperscript{a} & \BNEW \\
\ref{rmk:heilbronncost} & rmk & \texttt{\detokenize{rmk:heilbronncost}} & \BSTD{}~\cite[Ch.~2]{cremona} \COMP\textsuperscript{a} & \BNEW \\
\ref{rmk:levelNcheckable} & rmk & \texttt{\detokenize{rmk:levelNcheckable}} & \PROVED \BSTD{}~\cite{shimura,mazur,wwe}, \cite[Th\'eor\`eme~2]{merel96} \CONJ \COMP\textsuperscript{a} & \BNEW \\
\ref{prop:kappamodel} & prop & \texttt{\detokenize{prop:kappamodel}} & \PROVED & \BNEW \\
\ref{prop:degenpres} & prop & \texttt{\detokenize{prop:degenpres}} & \PROVED & \BNEW \\
\ref{rmk:degenpresconv} & rmk & \texttt{\detokenize{rmk:degenpresconv}} & \PROVED \COMP\textsuperscript{a} & \BNEW \\
\ref{rmk:levelNscale} & rmk & \texttt{\detokenize{rmk:levelNscale}} & \COMP\textsuperscript{a} & \BNEW \\
\ref{rmk:levelNcheckabletwo} & rmk & \texttt{\detokenize{rmk:levelNcheckabletwo}} & \PROVED \BSTD{}~\cite{atkin_lehner,billerey_menares} \CONJ \COMP\textsuperscript{a} & \BNEW \\
\ref{prop:levelNcurve} & prop & \texttt{\detokenize{prop:levelNcurve}} & \PROVED \BSTD{}~\cite[Ch.~I, \S4.2]{serre_trees} \COMP\textsuperscript{a} & \BNEW \\
\ref{rmk:levelNcurvefamily} & rmk & \texttt{\detokenize{rmk:levelNcurvefamily}} & \BSTD{}~\cite{mazur,calegari_emerton}, \cite[Thm.~1]{sturm} \CONJ \COMP\textsuperscript{a} & \BNEW \\
\midrule
\multicolumn{5}{@{}l}{\textbf{\S\ref{sec:nebentypus} Nebentypus: $\Ga_1(N)$ and $(\Ga_0(N),\chi)$ on the fixed model by Shapiro (v1.15.0-pre.1)}}\\
\midrule
\ref{def:nebdata} & def & \texttt{\detokenize{def:nebdata}} & \Bdef & \BNEW \\
\ref{lem:nebdata} & lem & \texttt{\detokenize{lem:nebdata}} & \PROVED & \BNEW \\
\ref{lem:teichmuller} & lem & \texttt{\detokenize{lem:teichmuller}} & \PROVED & \BNEW \\
\ref{rmk:nebfield} & rmk & \texttt{\detokenize{rmk:nebfield}} & \Brmk & \BNEW \\
\ref{lem:inducedchar} & lem & \texttt{\detokenize{lem:inducedchar}} & \PROVED & \BNEW \\
\ref{lem:gammaonetable} & lem & \texttt{\detokenize{lem:gammaonetable}} & \PROVED & \BNEW \\
\ref{def:nebmodel} & def & \texttt{\detokenize{def:nebmodel}} & \Bdef & \BNEW \\
\ref{prop:nebdecomp} & prop & \texttt{\detokenize{prop:nebdecomp}} & \PROVED & \BNEW \\
\ref{prop:nebcorr} & prop & \texttt{\detokenize{prop:nebcorr}} & \PROVED \BSTD{}~\cite[Ch.~III.9--10]{brown}, \cite[\S3.3]{shimura} & \BNEW \\
\ref{thm:nebrelations} & thm & \texttt{\detokenize{thm:nebrelations}} & \BSTD{}~\cite[\S3.3]{shimura}, \cite[Ch.~III.9--10]{brown} & \BNEW \\
\ref{prop:neborbit} & prop & \texttt{\detokenize{prop:neborbit}} & \PROVED & \BNEW \\
\ref{rmk:nebheilbronn} & rmk & \texttt{\detokenize{rmk:nebheilbronn}} & \PROVED & \BNEW \\
\ref{rmk:nebcost} & rmk & \texttt{\detokenize{rmk:nebcost}} & \Brmk & \BNEW \\
\ref{rmk:nebcheckable} & rmk & \texttt{\detokenize{rmk:nebcheckable}} & \PROVED \BSTD{}~\cite{shimura,cohen_oesterle,ribet77}, \cite[Lemmas~1.1.1, 1.1.4]{ash_stevens} \CONJ \COMP\textsuperscript{r} & \BNEW \\
\ref{def:nebwild} & def & \texttt{\detokenize{def:nebwild}} & \Bdef & \BNEW \\
\ref{lem:nebwild} & lem & \texttt{\detokenize{lem:nebwild}} & \PROVED \BSTD{}~\cite[Lemma~1.4, Prop.~2.1]{washington} & \BNEW \\
\ref{def:nebgammaH} & def & \texttt{\detokenize{def:nebgammaH}} & \Bdef & \BNEW \\
\ref{prop:nebgammaH} & prop & \texttt{\detokenize{prop:nebgammaH}} & \PROVED & \BNEW \\
\ref{prop:nebgammaHheilbronn} & prop & \texttt{\detokenize{prop:nebgammaHheilbronn}} & \PROVED & \BNEW \\
\ref{lem:nebshapiro} & lem & \texttt{\detokenize{lem:nebshapiro}} & \PROVED & \BNEW \\
\ref{lem:nebwildtors} & lem & \texttt{\detokenize{lem:nebwildtors}} & \PROVED \BSTD{}~\cite[Lemma~1.1.1, p.~194]{ash_stevens} & \BNEW \\
\ref{lem:nebchainrref} & lem & \texttt{\detokenize{lem:nebchainrref}} & \PROVED & \BNEW \\
\ref{rmk:nebwildcheckable} & rmk & \texttt{\detokenize{rmk:nebwildcheckable}} & \PROVED \COMP\textsuperscript{a} & \BNEW \\
\ref{rmk:nebcosetlevel} & rmk & \texttt{\detokenize{rmk:nebcosetlevel}} & \Brmk & \BNEW \\
\ref{def:levelpair} & def & \texttt{\detokenize{def:levelpair}} & \Bdef & \BNEW \\
\ref{lem:nebcompat} & lem & \texttt{\detokenize{lem:nebcompat}} & \PROVED & \BNEW \\
\ref{def:levelres} & def & \texttt{\detokenize{def:levelres}} & \Bdef & \BNEW \\
\ref{lem:levelresmorph} & lem & \texttt{\detokenize{lem:levelresmorph}} & \PROVED \BSTD{}~\cite[Ch.~III.9--10]{brown} & \BNEW \\
\ref{def:levelcone} & def & \texttt{\detokenize{def:levelcone}} & \Bdef & \BNEW \\
\ref{lem:levelconeexact} & lem & \texttt{\detokenize{lem:levelconeexact}} & \PROVED & \BNEW \\
\ref{prop:levelconeshapiro} & prop & \texttt{\detokenize{prop:levelconeshapiro}} & \PROVED & \BNEW \\
\ref{prop:modeltorsion} & prop & \texttt{\detokenize{prop:modeltorsion}} & \PROVED & \BNEW \\
\ref{prop:levelsplit} & prop & \texttt{\detokenize{prop:levelsplit}} & \PROVED & \BNEW \\
\ref{prop:levelconehecke} & prop & \texttt{\detokenize{prop:levelconehecke}} & \PROVED \BSTD{}~\cite[Ch.~III.9--10]{brown} \COMP\textsuperscript{a} & \BNEW \\
\ref{def:normalpair} & def & \texttt{\detokenize{def:normalpair}} & \Bdef & \BNEW \\
\ref{prop:levelinfres} & prop & \texttt{\detokenize{prop:levelinfres}} & \PROVED \BSTD{}~\cite[Ch.~VII, \S6; Ch.~III, Cor.~10.2]{brown} & \BNEW \\
\ref{prop:cycliccone} & prop & \texttt{\detokenize{prop:cycliccone}} & \PROVED \BSTD{}~\cite[Ch.~III, \S1, Example~2]{brown} & \BNEW \\
\ref{lem:levelfreeeuler} & lem & \texttt{\detokenize{lem:levelfreeeuler}} & \PROVED & \BNEW \\
\ref{lem:unipotentinv} & lem & \texttt{\detokenize{lem:unipotentinv}} & \PROVED & \BNEW \\
\ref{prop:gamma0normal} & prop & \texttt{\detokenize{prop:gamma0normal}} & \PROVED \COMP\textsuperscript{a} & \BNEW \\
\ref{prop:hidanormalsteps} & prop & \texttt{\detokenize{prop:hidanormalsteps}} & \PROVED \COMP\textsuperscript{a} & \BNEW \\
\ref{rmk:normalN1} & rmk & \texttt{\detokenize{rmk:normalN1}} & \PROVED & \BNEW \\
\ref{obl:levelconeexact} & obl & \texttt{\detokenize{obl:levelconeexact}} & \Bobl & \BNEW \\
\ref{rmk:lescheckable} & rmk & \texttt{\detokenize{rmk:lescheckable}} & \PROVED \COMP\textsuperscript{a} & \BNEW \\
\ref{rmk:oldspacecone} & rmk & \texttt{\detokenize{rmk:oldspacecone}} & \Brmk & \BNEW \\
\midrule
\multicolumn{5}{@{}l}{\textbf{\S\ref{sec:saturation} The $\Z/\ell^{K}$ Hecke algebra: saturation, the defect, the comparison with forms (v1.15.0-pre.1)}}\\
\midrule
\ref{def:satdata} & def & \texttt{\detokenize{def:satdata}} & \Bdef & \BNEW \\
\ref{prop:saturation} & prop & \texttt{\detokenize{prop:saturation}} & \PROVED & \BNEW \\
\ref{prop:satdefect} & prop & \texttt{\detokenize{prop:satdefect}} & \PROVED & \BNEW \\
\ref{lem:heckerankone} & lem & \texttt{\detokenize{lem:heckerankone}} & \PROVED \BSTD{}~\cite[\S\S3.5--3.6, Thm.~12.2.2, Prop.~12.4.10]{diamond_im}, \cite[Thm.~3.41, Ch.~8]{shimura}, \cite[Thm.~5.5.4]{diamond_shurman}, \cite[Ch.~VII]{serre_cours} & \BNEW \\
\ref{lem:freepart} & lem & \texttt{\detokenize{lem:freepart}} & \PROVED & \BNEW \\
\ref{cor:stablecert} & cor & \texttt{\detokenize{cor:stablecert}} & \PROVED & \BNEW \\
\ref{prop:cohforms} & prop & \texttt{\detokenize{prop:cohforms}} & \PROVED \BSTD{}~\cite[Ch.~8]{shimura} & \BNEW \\
\ref{prop:hidaalgconst} & prop & \texttt{\detokenize{prop:hidaalgconst}} & \PROVED \BSTD{}~\cite[Thm.~3.1, Cor.~3.2]{hida_ens} & \BNEW \\
\ref{rmk:thetaforms} & rmk & \texttt{\detokenize{rmk:thetaforms}} & \PROVED \COMP\textsuperscript{a} & \BNEW \\
\ref{rmk:sattor} & rmk & \texttt{\detokenize{rmk:sattor}} & \PROVED & \BNEW \\
\midrule
\multicolumn{5}{@{}l}{\textbf{\S\ref{sec:isomodelell} The $\ell$-local isotropy model, and the closed form at $\ell=3$ (v1.16.0-pre.1)}}\\
\midrule
\ref{def:isomodelell} & def & \texttt{\detokenize{def:isomodelell}} & \Bdef & \BNEW \\
\ref{prop:isomodelell} & prop & \texttt{\detokenize{prop:isomodelell}} & \PROVED \BSTD{}~\cite[Ch.~III.6.2]{brown} & \BNEW \\
\ref{prop:ellthreeclosed} & prop & \texttt{\detokenize{prop:ellthreeclosed}} & \PROVED & \BNEW \\
\ref{rmk:isomodellmp} & rmk & \texttt{\detokenize{rmk:isomodellmp}} & \PROVED & \BNEW \\
\ref{rmk:isomodellcheckable} & rmk & \texttt{\detokenize{rmk:isomodellcheckable}} & \PROVED \COMP\textsuperscript{a} & \BNEW \\
\ref{prop:certifiedeell} & prop & \texttt{\detokenize{prop:certifiedeell}} & \PROVED \BSTD{}~\cite[\S3.6]{weibel}, \cite[Ch.~III.10]{brown} & \BNEW \\
\ref{prop:descentN} & prop & \texttt{\detokenize{prop:descentN}} & \PROVED \BSTD{}~\cite[Ch.~III.6.2, III.10.2, V.3]{brown}, \cite[\S10.7]{weibel} & \BNEW \\
\ref{prop:badprimesN} & prop & \texttt{\detokenize{prop:badprimesN}} & \PROVED & \BNEW \\
\ref{prop:descentNthree} & prop & \texttt{\detokenize{prop:descentNthree}} & \PROVED \BSTD{}~\cite[Ch.~III.6.2]{brown}, \cite[Thm.~2.2.6]{weibel} & \BNEW \\
\ref{prop:descentmaps} & prop & \texttt{\detokenize{prop:descentmaps}} & \PROVED \BSTD{}~\cite[Ch.~III.6.2, III.8--10, VII.6]{brown}, \cite[Thm.~2.2.6]{weibel} & \BNEW \\
\ref{def:isomodeltree} & def & \texttt{\detokenize{def:isomodeltree}} & \Bdef & \BNEW \\
\ref{prop:isomodeltree} & prop & \texttt{\detokenize{prop:isomodeltree}} & \PROVED \BSTD{}~\cite[Thm.~2.2.6]{weibel}, \cite[Ch.~III.6.2]{brown} & \BNEW \\
\ref{prop:strictliftN} & prop & \texttt{\detokenize{prop:strictliftN}} & \PROVED \BSTD{}~\cite[Ch.~III.8--10]{brown}, \cite[Thm.~2.2.6]{weibel} & \BNEW \\
\ref{prop:crtcharpolyN} & prop & \texttt{\detokenize{prop:crtcharpolyN}} & \PROVED \BSTD{}~\cite[Thm.~12.2.2, Prop.~12.4.10]{diamond_im}, \cite{deligne_weil1}, \cite{diamond_shurman} & \BNEW \\
\ref{prop:torsionhecke} & prop & \texttt{\detokenize{prop:torsionhecke}} & \PROVED & \BNEW \\
\ref{thm:descentcert} & thm & \texttt{\detokenize{thm:descentcert}} & \PROVED & \BNEW \\
\ref{rmk:descentcertscope} & rmk & \texttt{\detokenize{rmk:descentcertscope}} & \PROVED & \BNEW \\
\ref{def:descentpair} & def & \texttt{\detokenize{def:descentpair}} & \Bdef & \BNEW \\
\ref{lem:pairlifting} & lem & \texttt{\detokenize{lem:pairlifting}} & \PROVED \BSTD{}~\cite[Thm.~2.2.6]{weibel} & \BNEW \\
\ref{prop:descentpair} & prop & \texttt{\detokenize{prop:descentpair}} & \PROVED \BSTD{}~\cite[Ch.~III.2, III.5--6, III.8--9, V.1--3, VI.9]{brown}, \cite[Thm.~1.3.1, Thm.~2.2.6]{weibel} & \BNEW \\
\ref{prop:descentpairhecke} & prop & \texttt{\detokenize{prop:descentpairhecke}} & \PROVED \BSTD{}~\cite[Ch.~III.2, III.8--10]{brown}, \cite[Thm.~2.2.6]{weibel} & \BNEW \\
\ref{rmk:descentpairscope} & rmk & \texttt{\detokenize{rmk:descentpairscope}} & \PROVED & \BNEW \\
\ref{def:treepair} & def & \texttt{\detokenize{def:treepair}} & \Bdef & \BNEW \\
\ref{lem:treepaircyl} & lem & \texttt{\detokenize{lem:treepaircyl}} & \PROVED \BSTD{}~\cite[\S1.5]{weibel} & \BNEW \\
\ref{prop:treepair} & prop & \texttt{\detokenize{prop:treepair}} & \PROVED \BSTD{}~\cite[Thm.~2.2.6, \S1.5]{weibel}, \cite[Ch.~III.6, III.8--9]{brown} & \BNEW \\
\ref{rmk:treepairscope} & rmk & \texttt{\detokenize{rmk:treepairscope}} & \PROVED & \BNEW \\
\ref{def:heckeZ} & def & \texttt{\detokenize{def:heckeZ}} & \Bdef & \BNEW \\
\ref{prop:heckeZspec} & prop & \texttt{\detokenize{prop:heckeZspec}} & \PROVED \BSTD{}~\cite[Ch.~3]{shimura} & \BNEW \\
\ref{prop:heckeZrat} & prop & \texttt{\detokenize{prop:heckeZrat}} & \PROVED \BSTD{}~\cite[Thm.~12.2.2, Props.~12.4.10, 12.4.13, \S6]{diamond_im}, \cite[Thms.~5.8--5.10, 9.23]{stein}, \cite{atkin_lehner}, \cite{coleman_edixhoven}, \cite{sturm} \CONJ & \BNEW \\
\ref{lem:heckeZbridge} & lem & \texttt{\detokenize{lem:heckeZbridge}} & \PROVED & \BNEW \\
\ref{prop:heckeZmax} & prop & \texttt{\detokenize{prop:heckeZmax}} & \PROVED \BSTD{}~\cite[Prop.~2.4]{atiyah_macdonald} & \BNEW \\
\ref{prop:heckeZgen} & prop & \texttt{\detokenize{prop:heckeZgen}} & \PROVED & \BNEW \\
\ref{prop:heckeZetale} & prop & \texttt{\detokenize{prop:heckeZetale}} & \PROVED & \BNEW \\
\ref{prop:heckeZmono} & prop & \texttt{\detokenize{prop:heckeZmono}} & \PROVED \BSTD{}~\cite[Prop.~3.1.19]{bruns_herzog} & \BNEW \\
\ref{prop:heckeZlocal} & prop & \texttt{\detokenize{prop:heckeZlocal}} & \PROVED & \BNEW \\
\ref{prop:heckeZthrough} & prop & \texttt{\detokenize{prop:heckeZthrough}} & \PROVED & \BNEW \\
\ref{prop:heckeZuct} & prop & \texttt{\detokenize{prop:heckeZuct}} & \PROVED \BSTD{}~\cite[\S3.6, Thm.~1.3.1]{weibel} & \BNEW \\
\ref{prop:heckeZsplit} & prop & \texttt{\detokenize{prop:heckeZsplit}} & \PROVED & \BNEW \\
\ref{prop:heckeZlevelone23} & prop & \texttt{\detokenize{prop:heckeZlevelone23}} & \PROVED & \BNEW \\
\ref{rmk:heckeZcbone} & rmk & \texttt{\detokenize{rmk:heckeZcbone}} & \CONJ & \BNEW \\
\ref{prop:heckeZglue} & prop & \texttt{\detokenize{prop:heckeZglue}} & \PROVED \BSTD{}~\cite[Cor.~10.10]{atiyah_macdonald} & \BNEW \\
\ref{rmk:heckeZcbtwo} & rmk & \texttt{\detokenize{rmk:heckeZcbtwo}} & \CONJ & \BNEW \\
\ref{rmk:heckeZcbthree} & rmk & \texttt{\detokenize{rmk:heckeZcbthree}} & \CONJ & \BNEW \\
\ref{thm:heckeZdescent} & thm & \texttt{\detokenize{thm:heckeZdescent}} & \PROVED & \BNEW \\
\ref{thm:heckeZcert} & thm & \texttt{\detokenize{thm:heckeZcert}} & \PROVED \BSTD{}~\cite{deligne_weil1} & \BNEW \\
\ref{rmk:heckeZcaveat} & rmk & \texttt{\detokenize{rmk:heckeZcaveat}} & \PROVED & \BNEW \\
\ref{def:heckeZcong} & def & \texttt{\detokenize{def:heckeZcong}} & \Bdef & \BNEW \\
\ref{prop:heckeZcong} & prop & \texttt{\detokenize{prop:heckeZcong}} & \PROVED \BSTD{}~\cite{deligne_weil1} & \BNEW \\
\ref{prop:heckeZeiscert} & prop & \texttt{\detokenize{prop:heckeZeiscert}} & \PROVED \BSTD{}~\cite[II (9.7), (16.1), Cor.~(16.2), (16.6)]{mazur} & \BNEW \\
\midrule
\multicolumn{5}{@{}l}{\textbf{\S\ref{sec:mpmodel} The Moerdijk--Pronk model (v1.22.0-pre.5)}}\\
\midrule
\ref{def:mptri} & def & \texttt{\detokenize{def:mptri}} & \Bdef & \BNEW \\
\ref{def:mprules} & def & \texttt{\detokenize{def:mprules}} & \Bdef & \BNEW \\
\ref{def:mpell} & def & \texttt{\detokenize{def:mpell}} & \Bdef & \BNEW \\
\ref{def:mpRZ} & def & \texttt{\detokenize{def:mpRZ}} & \Bdef & \BNEW \\
\ref{rmk:mprulesread} & rmk & \texttt{\detokenize{rmk:mprulesread}} & \PROVED \COMP\textsuperscript{r} & \BNEW \\
\ref{def:mpset} & def & \texttt{\detokenize{def:mpset}} & \Bdef & \BNEW \\
\ref{rmk:mpwhat} & rmk & \texttt{\detokenize{rmk:mpwhat}} & \PROVED \COMP\textsuperscript{r} & \BNEW \\
\ref{def:mpfree} & def & \texttt{\detokenize{def:mpfree}} & \Bdef & \BNEW \\
\ref{lem:mprows} & lem & \texttt{\detokenize{lem:mprows}} & \PROVED & \BNEW \\
\ref{lem:mpcols} & lem & \texttt{\detokenize{lem:mpcols}} & \PROVED & \BNEW \\
\ref{lem:mpassembly} & lem & \texttt{\detokenize{lem:mpassembly}} & \PROVED \BSTD{}~\cite[Lemma~2.7.3]{weibel} & \BNEW \\
\ref{thm:mpcomparison} & thm & \texttt{\detokenize{thm:mpcomparison}} & \PROVED \BSTD{}~\cite[\S8.3, Thm.~2.2.6]{weibel} & \BNEW \\
\ref{rmk:mphybrid} & rmk & \texttt{\detokenize{rmk:mphybrid}} & \PROVED \BSTD{}~\cite[Thm.~2.2.6]{weibel} \COMP\textsuperscript{r} & \BNEW \\
\ref{def:mpqd} & def & \texttt{\detokenize{def:mpqd}} & \Bdef & \BNEW \\
\ref{prop:mpqreduced} & prop & \texttt{\detokenize{prop:mpqreduced}} & \PROVED \BSTD{}~\cite[Thm.~2.2.6, \S8.3]{weibel} \COMP\textsuperscript{a} & \BNEW \\
\ref{def:mpphi} & def & \texttt{\detokenize{def:mpphi}} & \Bdef & \BNEW \\
\ref{prop:mpalldegrees} & prop & \texttt{\detokenize{prop:mpalldegrees}} & \PROVED \BSTD{}~\cite[Thm.~2.2.6]{weibel} \COMP\textsuperscript{r} & \BNEW \\
\ref{def:mpmorphism} & def & \texttt{\detokenize{def:mpmorphism}} & \Bdef & \BNEW \\
\ref{thm:mpnatural} & thm & \texttt{\detokenize{thm:mpnatural}} & \PROVED \BSTD{}~\cite[Thm.~2.2.6]{weibel}, \cite[Ch.~III.8]{brown} \COMP\textsuperscript{r} & \BNEW \\
\ref{prop:mpbredon} & prop & \texttt{\detokenize{prop:mpbredon}} & \PROVED \BSTD{}~\cite[Thm.~2.2.6]{weibel} \COMP\textsuperscript{r} & \BNEW \\
\ref{prop:mpboundary} & prop & \texttt{\detokenize{prop:mpboundary}} & \PROVED \BSTD{}~\cite[Ch.~III.6.2, III.9]{brown} \COMP\textsuperscript{r} & \BNEW \\
\ref{prop:mprescor} & prop & \texttt{\detokenize{prop:mprescor}} & \PROVED \BSTD{}~\cite[Ch.~III.9--10]{brown} \COMP\textsuperscript{r} & \BNEW \\
\ref{rmk:mphecke} & rmk & \texttt{\detokenize{rmk:mphecke}} & \COMP\textsuperscript{r} & \BNEW \\
\ref{ex:mpteardrop} & ex & \texttt{\detokenize{ex:mpteardrop}} & \COMP\textsuperscript{r} & \BNEW \\
\ref{rmk:mpcheckable} & rmk & \texttt{\detokenize{rmk:mpcheckable}} & \PROVED \COMP\textsuperscript{r} & \BNEW \\
\bottomrule
\end{longtable}
\end{center}

%----------------------------------------------------------------------
\section{Formalization notes}\label{app:formalization}
%----------------------------------------------------------------------
% SOURCE:  the nineteen (eighteen until v1.8.0) per-chapter "Formalization note" paragraphs of this
%          document, gathered and reconciled (NEW); Appendix E's roster of
%          the statements nothing outside this document states.
% NOTATION: labels are dereferenced, so this appendix is build-gated like
%          Appendix B.
% ERRATA:  none.

\paragraph{What this appendix is, and what it is not.}
The owner intends to commission an assessment of a Lean~4 / \code{mathlib4}
formalization of this document. \textbf{This appendix is the input to that commission,
not the commission.} It says what the document's chapters have each recorded about
their own formalizability, reconciled into one place: which objects a formalization
would have to build, which statements it should import rather than re-prove, which
hypotheses it must not lose, and which statements are the cheapest and the most valuable
to attack. It deliberately \textbf{does not} estimate effort, propose a schedule, or
rank the work by cost: that is the assessment, and an appendix that pre-empted it with
guesses would be supplying the answer it is meant to supply the data for.

\begin{contractbox}
\textbf{Every claim below about what \code{mathlib4} contains is a claim made by the
chapter that made it, and was not re-verified against \code{mathlib4} in this pass.}
No version of the library was consulted while this document was written; the
per-chapter notes were written by authors reasoning about what a library of that scope
plausibly carries. The commission must check each one. They are recorded here because a
wrong belief about the library is cheap to correct and expensive to discover late ---
but a reader must not treat ``\code{mathlib4} has this'' in this document as a fact
about \code{mathlib4}.
\end{contractbox}

\subsection{What a formalization would have to build}

\noindent
These have no counterpart the chapters expect to find, and are the substance of the
work.
\begin{enumerate}
\item \textbf{The $(K,\Rel)$ framework} (Chapter~\ref{sec:framework}):
  Definition~\ref{def:relationdatum}, the quotient complex
  \ref{def:quotientcomplex}, saturation \ref{def:saturated}, and the fundamental-domain
  hypotheses \ref{FD1}--\ref{FD4}. The chapter's own note is that this is
  formalization-\emph{friendly} rather than hard: for finite $K$ with decidable equality
  on $\orient{\bullet}(\Ku)$ the coincidence groupoid is a decidable relation and
  \ref{FD4} is reachability in a finite graph. It is nonetheless entirely new
  construction.
\item \textbf{The isomorphism $\Psi$} of Theorem~\ref{thm:explicit_complex_equiv}(1),
  with $\Phi$ given by \eqref{eq:psiinverse}. It is finite, constructive, needs no
  topology, no field and no freeness, and is the statement six source files cite. It is
  the single most load-bearing object of Part~I.
\item \textbf{$\Sym^n$ as a monoid representation over an arbitrary $\cring$}
  (Definition~\ref{def:rhon}, Proposition~\ref{prop:symconv}). What is needed is not
  $\Sym^n$ of a module --- the chapters expect that, via \code{MvPolynomial} --- but the
  extension of the action from $\GL_2$ to the monoid $\De$ of
  Definition~\ref{def:semigroup}, in the explicit basis $u_j=x^{n-j}y^j$ with the closed
  form \eqref{eq:rhondef}, over $\cring=\Z/m$ as well as over a field. The contragredient
  $\Dn$ and the pairing (Definition~\ref{def:contragredient},
  Proposition~\ref{prop:pairing}) come with it and are finite linear algebra over a
  commutative ring.
\item \textbf{Twisted chain data}: Definition~\ref{def:tau}'s (T1)--(T3), the boundary
  analogue \ref{def:taubdry}, and the pair homotopy \ref{def:pairhomotopy}. These are
  equational conditions on maps of complexes and are directly expressible; the
  \emph{constructions} (Algorithm~\ref{alg:tau}, Lemmas~\ref{lem:taubdryexists} and
  \ref{lem:etaexists}) are terminating procedures on finite data whose correctness
  reduces to $Z_1=0$ on a tree and $Z_2=0$ on a contractible $2$-complex --- simplicial
  chains, not homological algebra.
\item \textbf{The truncated tessellation and its boundary} (\ref{H4},
  \ref{B1}--\ref{B3}), the pair $(\Xbulk,\Xbdry)$ of \ref{def:pair} and the mapping cone
  with the derived homotopy $\hder=\widetilde{\etapair}$ (\ref{def:coneaction}). The
  cone of a cochain map and the long exact sequence of a short exact sequence of
  complexes are expected from the library; the \emph{pair homotopy} and its coset
  transport are not.
\item \textbf{The Hecke ring $\mathcal{H}(\De,\Ga)$} and double cosets of a group in a
  monoid (\ref{thm:relations}, \ref{obl:algebra}). Chapter~\ref{sec:heckedatum}'s note
  offers the cheaper alternative explicitly: verify the $\nu=2$ instance directly as a
  finite identity of cochain maps and induct, which is what
  \code{HeckeRelationsValidation} does numerically, rather than import a structure
  theorem.
\item \textbf{The free $\Ga$-model} $F_{\bullet}\otimes C_{\bullet}(Y)$ of
  \ref{def:borelmodel}, whose clause (1) --- $F\otimes_{\Z}M$ is free over $\Z[\Ga]$ when
  $F$ is free and $M$ is $\Z$-free --- the chapter calls easy, useful and worth doing
  early, because Remark~\ref{rmk:borelfree}, the removal of every invertibility
  hypothesis, is one line from it plus the fundamental lemma.
\item \textbf{An arithmetic model} for Part~V. Theorem~\ref{thm:backend} is about
  fixed-width integer types and Theorem~\ref{thm:fp64} about IEEE-754; both chapters
  recommend carrying these as \emph{hypotheses on an abstract arithmetic model} rather
  than as facts about hardware, with (FP2) plus exact representability being all the
  proofs use.
\item \textbf{The integral line} (Chapter~\ref{sec:integral}, since v1.8.0). Its note
  ranks the fracture square and cohomology-level glueing
  (Propositions~\ref{prop:fracture}, \ref{prop:cohglue}) first --- elementary facts about
  flat abelian groups, and what makes the descent exact --- then the strict lifting lemma
  \ref{lem:strictlifting}, a finite induction; the Bredon identification
  (Proposition~\ref{prop:bredon}) needs a functor category and the comparison theorem in it,
  and the isotropy-resolution model (Proposition~\ref{prop:isomodel}) the assembly of
  resolutions over a complex.
\end{enumerate}

\subsection{What to import rather than prove}

\noindent
The chapters are unanimous that these should be axioms or library imports, and say why
in each case: Theorem~\ref{thm:lcs_rep_equiv} (the only genuinely topological input of
Chapter~\ref{sec:localsystems}, and used nowhere downstream --- Remark~\ref{rmk:lcsnotfree});
Theorem~\ref{thm:cohom_equiv} and Theorem~\ref{thm:explicit_complex_equiv}(2)
(\cite[Ch.~I.4]{brown}, \cite[\S3.H]{hatcher}); Proposition~\ref{prop:universalcover};
the fundamental lemma of relative homological algebra, which is what
Lemma~\ref{lem:tauexists}(2) and, through it, Theorem~\ref{thm:welldef}(b),(c),
Theorem~\ref{thm:classical}, Theorem~\ref{thm:conewelldef}(c) and
Proposition~\ref{prop:iotacommutes}(3c) all rest on; Eichler--Shimura and the Dickson
invariants, which are what Claim~\ref{cl:dimH1} imports and which make it and
Claim~\ref{cl:pd} \emph{not} targets at this stage; the Bass--Serre sequence
\eqref{eq:treeses}, which a formalization may take as the definition of the amalgam
datum; the isotropy spectral sequence behind Theorem~\ref{thm:specialisation}(1); and
the $\Heq{2}=0$ vanishing that \S\ref{sec:degreewise} carries by citation --- which a
formalization must import rather than read off an emitted $0\times0$ matrix.

\noindent
Two further items are flagged by their chapters as things to \emph{avoid} formalising
in the obvious way. Remark~\ref{rmk:nogo}(ii) is a non-constructive argument
(irreducibility plus a numerical non-scalarity) and should be a statement about
$\End_{\Ga}(\Sym^{10}\Q^2)$ with an imported eigenvalue fact, not re-derived.
Proposition~\ref{prop:bdrycohom} needs the Lyndon--Hochschild--Serre spectral sequence
of a product with a cyclic factor; Chapter~\ref{sec:boundarydatum} recommends replacing
it by the explicit $\langle T\rangle$-equivariant periodic resolution tensored with the
two-term resolution over $\Z[\Z]$, which is finite and avoids the spectral sequence
entirely.

\subsection{Hypotheses a formalization must not lose}

\noindent
Rule~\ref{A1} exists because of exactly these, and the OQ2 audit's finding F5 --- an
unwritten load-bearing hypothesis --- is the cautionary case. Five recur:
\begin{itemize}
\item \ref{H4}(ii), \emph{$\abs{\Ku}$ contractible and two-dimensional}: used in
  Lemma~\ref{lem:etaunique} (via $Z_2=0$), in the fills of Algorithm~\ref{alg:tau}
  step~3 and Lemma~\ref{lem:etaexists}, and everywhere in
  Chapter~\ref{sec:cone}.
\item \ref{B3}, \emph{boundary stabilisers are central}: three uses in
  Lemma~\ref{lem:taubdryexists} and three in Lemma~\ref{lem:etaexists}, and it is what
  makes the (T1$\partial$)-extension well defined.
\item \emph{Relative projectivity} in Lemma~\ref{lem:tauexists}(2), whose entire content
  here is that $4$ and $6$ --- the elliptic stabiliser orders of \ref{H4}(iv) --- are
  units. Stating it as $6\in\cring^{\times}$ for a general commutative $\cring$, rather
  than as $\ell\ge5$ for $\Fl$, is what puts it in the statement instead of the proof.
\item Proposition~\ref{prop:iotasplit}(1)'s $\tfrac12$, which must be carried as an
  $\mathrm{Invertible}\,(2:\cring)$ instance argument and \textbf{not} as a field
  hypothesis, since $\cring=\Z/\ell^{e}$ is in scope.
\item The (T3) normalisation of Definition~\ref{def:tau}: without it the degree-$0$
  operator is still well defined but is not normalised, and every $H^0$ eigenvalue
  oracle of \S\ref{sec:example} is off by the augmentation defect
  (Remark~\ref{rmk:T3}).
\end{itemize}

\subsection{Shapes that recur, and should be factored once}

\noindent
Chapter~\ref{sec:obligations} identifies three predicate shapes that every obligation
reduces to --- entrywise equality of two matrices over $\cring$; equality of two integer
chains after canonicalisation; a magnitude recount against a stored bound --- and the
one obligation that is not of that shape, \ref{obl:fp}(iv), which quantifies over
\emph{builds} and which a formalization should replace by the statement that the
pipeline's output is a function of the abstract sign predicate alone
(Lemma~\ref{lem:predshapes} plus Theorem~\ref{thm:fp64} give exactly that, and it is a
better target than the two-build diff). Two further shapes recur across chapters and
are worth one shared development: the \textbf{$1$-cocycle chain rule}
(Lemma~\ref{lem:cocycle}, reused verbatim by Proposition~\ref{prop:blockmatrices}, by
the $(\ast,\kappa)$ rule of Proposition~\ref{prop:barformula} and by
Lemma~\ref{lem:cosettransport}), and the \textbf{shared (T1)-extension helper} that
Algorithm~\ref{alg:tau}, \ref{lem:taubdryexists} and \ref{lem:etaexists} each use for
both construction and certification --- in the implementation it is deliberately one
function, and a formalization that duplicates it loses the property that makes the
certification independent of the construction only where it should be.

\subsection{The targets: cheapest, and most valuable}

\noindent
\textbf{These are two different lists, and the chapters say so.} The cheapest statements
are cheap because they are finite, constructive and self-contained; the most valuable
are those an error in which could not be caught any other way. The second list is
\emph{not} constructed here: it is Appendix~\ref{app:changelog}'s roster of the $67$
statements that no superseded document and no project note states --- the ones a referee
should attack first, for the reason given there, namely that every other statement of
this contract can be checked against a source and these cannot. A formalization's
priority order should be read off that roster, and this appendix adds only the
per-chapter judgement of which of them are also \emph{tractable}.

\begin{center}
\footnotesize
\begin{tabular}{@{}p{3.9cm}p{4.4cm}p{6.4cm}@{}}
\toprule
part & cheapest first target & why, and whether it is in the roster \\
\midrule
I, the framework &
 \ref{thm:explicit_complex_equiv}(1); then \ref{prop:ACq}, \ref{prop:AZq}, \ref{prop:Bqimage}, \ref{prop:projector} &
 $\Psi$ needs no topology, no field and no freeness, and is cited by six source files; the four propositions are kernels, images and quotients of explicit matrices --- ``the most directly formalizable material in the document''. All five are in the roster. The contract box of \S\ref{sec:systematiccomputation} must come with them: the \emph{basis} is part of the specification, and a proof that fixes only the spans would not pin $\Amat{q}$ \\
\addlinespace
II, model and coefficients &
 \ref{rmk:negI}; \eqref{eq:pairadjoint} &
 the odd-$n$ vanishing is two lines from centrality and $2\in\cring^{\times}$, and it is oracle~2 of \S\ref{sec:example}; \eqref{eq:pairadjoint} is what everything about the duality rests on. \ref{prop:pairing} is in the roster, \ref{rmk:negI} is not (it is re-homed) \\
\addlinespace
III, the Hecke datum &
 \ref{lem:cocycle}; then the equivariance half of \ref{thm:welldef} and \ref{thm:xi}(1),(2) &
 one line from uniqueness of coset representatives; the two halves are finite manipulations using only \ref{lem:cocycle} and \ref{prop:symconv}(2), need \textbf{no} hypothesis on $\ell$, and are the chapters' own ``natural first target of Part~III''. \ref{lem:cocycle} is in the roster \\
\addlinespace
III, the two bit-identities &
 \ref{prop:tworoutes}, \ref{prop:mu}(3) &
 the only two claims in the document where a \emph{bit}-level agreement is mathematically forced rather than empirically observed --- an integral identity in $\Mat_2(\Z)$ pushed through $\rhon$. Both are in the roster, and both are cheap \\
\addlinespace
IV, the boundary &
 \ref{prop:bdrycosets}; then \ref{lem:bdrycomponent}, \ref{lem:taubdryexists} &
 a finite matrix-arithmetic statement, decidable for each $N$; the two that carry the content are constructive and need $H_1$ of a tree, \textbf{neither} the fundamental lemma \textbf{nor} any invertibility hypothesis. All three are in the roster \\
\addlinespace
IV, the cone &
 \ref{lem:cosettransport}, \ref{lem:conesign}, \ref{lem:derivedhomotopy}, \ref{lem:etaunique}, then \ref{prop:coneindep}(1) &
 \textbf{the chapter to attack first.} It is the one whose statements have \emph{no} external corroboration at all: its only sources are project notes, which are not sources of record. In that order the work is four one-line identities, pure algebra on a direct sum, a substitution, two uses of $Z_2=0$, and then one displayed computation --- and \ref{prop:coneindep}(1) is the theorem that matters \\
\addlinespace
IV, six groups, orbifold model &
 \ref{cl:amalgam}(2); \ref{def:borelmodel}(1) &
 the first is $\Ext^{\bullet}_{\Z[\Ga]}(-,\VR)$ on one short exact sequence and is the most reusable thing in its chapter; the second is easy and unlocks \ref{rmk:borelfree}. Both are in the roster. \ref{cl:dimH1} and \ref{cl:pd} are \emph{not} targets: the first imports Eichler--Shimura and Dickson, the second is open \\
\addlinespace
V, realisation &
 \ref{lem:predshapes}; \ref{cor:fptotal}(ii) &
 \ref{lem:predshapes} is a finite case analysis over $\Z$ needing no floating-point model at all, and it is the \emph{only} assertion in Part~V that no source states --- the whole of the rest of the part is re-homable. \ref{cor:fptotal}(ii) replaces an error estimate by a decidable integer-magnitude statement \\
\addlinespace
the involution &
 \ref{def:iota}(i)--(ii), \ref{lem:iotamodel}, \ref{prop:iotacommutes}(1),(2) &
 decidable predicates on finite data. \ref{prop:iotacommutes}(3c) is the only statement of that chapter invoking an abstract input, and inherits $6\in\cring^{\times}$ verbatim \\
\addlinespace
the obligations &
 all of \S\ref{sec:obligations} &
 each is already a decidable predicate on a finite emitted table (\ref{A4}), so the predicate needs no proof, only evaluation; what a formalization proves is the implication ``predicate holds on the emitted data $\Rightarrow$ the stored $\Amat{q}$ is the matrix of $\Tdc{g}$'' \\
\bottomrule
\end{tabular}
\end{center}

\subsection{What this appendix does not supply}

\begin{enumerate}
\item \textbf{Any cost.} No effort estimate, no ordering by expense, no schedule.
\item \textbf{Any verified statement about \code{mathlib4}.} See the box at the head of
  this appendix.
\item \textbf{An executable oracle for one chapter and one clause.} A formalization of
  Chapter~\ref{sec:orbifoldmodel} cannot be cross-checked against a run, and one of
  Chapter~\ref{sec:cone} in everything but Obligation~\ref{obl:coneexact}(vii)'s
  cross-routes and PD adjointness (computed since v1.7.3, slot-6 WP-3 / WP-4 / WP-5), because nothing in the tree computed the rest of their
  objects (\S\ref{sec:notimplemented}; Chapter~\ref{sec:boundarydatum} was a third until
  v1.3.3, Chapter~\ref{sec:cone} wholly uncomputed until v1.4.2 and computed in its
  $\etapair$-half only until v1.5.2). For the rest of the document the implementation is a
  second, independent witness --- $89$ registered validations, and
  Appendix~\ref{app:status}'s eleven \COMP\textsuperscript{r} claims --- and a
  formalization can be developed against it. For those two chapters it cannot, which
  is a fact about the input to the commission and not a judgement about the work.
\item \textbf{A decision on \ref{open:ell2} --- taken at v1.3.0, in the weakest form a
  formalisation could want.} The owner decided \ref{oc1a} on 2026-09-16: the decision
  is that $q_{\max}$ \emph{is a parameter}, so a development takes it as a section
  variable and proves nothing about it, which is what \ref{def:borelmodel} and
  \ref{thm:specialisation} already do. What a formalisation of the truncated regime
  would still be formalising is a policy, not a theorem --- unless \ref{oc1c} is
  proved, whose bulk half now is one: \ref{lem:bulkperiodic} is a natural target (a
  five lemma over \ref{cl:amalgam}(2), the cup product with a class in $H^{2}(\Ga;\Z)$,
  and the transfer's projection formula; \code{mathlib4} has group cohomology in low
  degrees but, as of this writing, neither cup products on it nor the transfer), and
  \ref{rmk:bulkperiodicscope}(3) lists the three statements that would complete it.
  The degree-generic statements of Chapters~\ref{sec:cone} and \ref{sec:orbifoldmodel}
  are written so that none of this blocks the rest.
\end{enumerate}

%----------------------------------------------------------------------
\section{Index of preserved labels}\label{app:labelindex}
%----------------------------------------------------------------------
This appendix exists to make the label namespace machine-checkable: every label that
a C++ source file or a project document cites is dereferenced here, so a missing,
renamed or duplicated label is a build error rather than a silent breakage.

\subsection*{The 67 labels of \file{hecke\_action.tex}}
\S\ref{sec:scope},
\S\ref{sec:conventions},
\ref{rmk:orbifold},
\S\ref{sec:model},
\ref{rmk:negI},
\S\ref{sec:symconv},
\ref{def:rhon},
\ref{prop:symconv},
\ref{def:semigroup},
\ref{lem:detsym},
\S\ref{sec:heckedatum},
\ref{prop:cosets},
\ref{def:operators},
\eqref{eq:fullop},
\eqref{eq:fullassembly},
\ref{prop:diamond},
\ref{thm:relations},
\ref{cor:ulrecursion},
\S\ref{sec:chainlevel},
\eqref{eq:corconjres},
\ref{def:tau},
\ref{lem:tauexists},
\ref{rmk:tauexistsscope},
\ref{def:hecke},
\eqref{eq:heckechain},
\ref{thm:welldef},
\ref{rmk:indeplevel},
\ref{thm:classical},
\S\ref{sec:induced},
\ref{def:inducedmodule},
\ref{prop:blockmatrices},
\eqref{eq:MH},
\ref{def:inducedls},
\ref{prop:lhcompliant},
\ref{prop:alpha},
\ref{prop:omega},
\ref{rmk:nogo},
\ref{def:xi},
\eqref{eq:xidef},
\eqref{eq:xiexplicit},
\ref{thm:xi},
\ref{alg:assembly},
\S\ref{sec:equivalence},
\ref{rmk:barformula},
\ref{prop:barformula},
\S\ref{sec:reduction},
\ref{def:primitives},
\ref{alg:tau},
\ref{lem:bounds},
\ref{thm:backend},
\S\ref{sec:fpfilter},
\ref{thm:fp64},
\ref{cor:fptotal},
\ref{cor:fp32},
\S\ref{sec:obligations},
\ref{obl:welldef},
\ref{obl:indep},
\ref{obl:integral},
\ref{obl:algebra},
\ref{obl:agreement},
\ref{obl:shapiro},
\ref{obl:oracle},
\ref{rmk:dickson},
\ref{obl:fp},
\ref{obl:backend},
\ref{obl:tau},
\S\ref{sec:example},
\S\ref{sec:implementation}.

\subsection*{The label of \file{monodromy\_v5.tex} cited from C++ source}
\ref{thm:explicit_complex_equiv}.

\subsection*{Further \file{monodromy\_v5.tex} labels carried with the re-homed text}
\ref{ass:K_contractible},
\ref{ass:edge_non_ident},
\ref{lemma:pR_simplicial},
\ref{ass:open_interior},
\ref{prop:injectivity},
\ref{prop:saturated_complexes},
\ref{def:universal_cover},
\ref{prop:universalcover},
\ref{rmk:injectivityofi},
\ref{ass:fund_domain},
\ref{prop:fund_group_action},
\ref{thm:fund_group_generators},
\ref{def:matrix_ls},
\ref{prop:matrix_rep_equiv},
\ref{thm:lcs_rep_equiv},
\ref{thm:rep_ls_equiv},
\ref{rmk:universalcovercohomology},
\ref{thm:cohom_equiv},
\S\ref{sec:explicitcochaincomplex},
\ref{def:eq_cochains},
\ref{ass:consistent_face_id},
\ref{lemma:d_maps_C_eq},
\S\ref{sec:systematiccomputation}.

\subsection*{Labels re-homed from \file{hecke\_extensions.tex} (cited from project documents as [E])}
\S\ref{sec:six},
\ref{cl:amalgam},
\ref{cl:dimH1},
\ref{cl:bockstein},
\ref{cl:boundary},
\ref{cl:les},
\eqref{eq:les},
\ref{cl:chi},
\ref{cl:pd}.

\subsection*{Labels new in this document}
\S\ref{sec:notation},
\S\ref{sec:status},
\S\ref{sec:labelpolicy},
\S\ref{sec:authoring},
\S\ref{sec:versioning},
\ref{def:model7},
\ref{def:model7rel},
\ref{prop:model7h4},
\eqref{eq:negI},
\eqref{eq:flatindex},
\eqref{eq:rhondef},
\S\ref{sec:contragredient},
\eqref{eq:adjrules},
\ref{def:contragredient},
\ref{prop:pairing},
\eqref{eq:dualaction},
\eqref{eq:pairadjoint},
\ref{rmk:selfadjoint},
\ref{lem:symdual},
\eqref{eq:BnGram},
\ref{rmk:divided},
\S\ref{sec:involution},
\ref{def:iota},
\eqref{eq:iotaconj},
\eqref{eq:rhoniota},
\eqref{eq:iotacochain},
\ref{lem:iotamodel},
\ref{rmk:iotasigns},
\ref{prop:iotacommutes},
\ref{rmk:iotatausym},
\ref{prop:iotasplit},
\ref{rmk:iotaelltwo},
\ref{rmk:iotapairing},
\ref{rmk:iotasign},
\ref{obl:iota},
\S\ref{sec:primfull},
\S\ref{sec:correspondencedatum},
\ref{def:cosetreps},
\eqref{eq:cocycle},
\ref{lem:cocycle},
\S\ref{sec:operatordefs},
\ref{prop:fullassembly},
\ref{rmk:rawprimitive},
\S\ref{sec:relations},
\ref{rmk:relcoeff},
\ref{rmk:badprimesingular},
\S\ref{sec:correspondence},
\S\ref{sec:taudatum},
\ref{rmk:taurings},
\ref{rmk:tauintegral},
\S\ref{sec:operatordef},
\ref{rmk:T3},
\S\ref{sec:classicalcomparison},
\ref{rmk:whichoperator},
\S\ref{sec:boundarydatum},
\ref{B1}, \ref{B2}, \ref{B3},
\ref{rmk:bdrymodel},
\S\ref{sec:bdrygroup},
\ref{def:bdrygroup},
\eqref{eq:bdrysplit},
\ref{prop:bdrycohom},
\eqref{eq:bdrylhs},
\ref{prop:circlemodel},
\ref{rmk:bdryoddn},
\S\ref{sec:bdrycosets},
\ref{prop:bdrycosets},
\ref{rmk:bdrynugeneral},
\S\ref{sec:taubdry},
\ref{def:taubdry},
\ref{lem:bdrycomponent},
\ref{lem:taubdryexists},
\ref{rmk:taubdryintegral},
\ref{alg:taubdry},
\ref{rmk:taubdrynoiota},
\S\ref{sec:bdryop},
\ref{def:bdryop},
\eqref{eq:bdrychain},
\ref{prop:bdryop},
\ref{prop:bdrychar},
\eqref{eq:bdrygroupform},
\ref{prop:bdrytranscription},
\eqref{eq:bdrygroupformnu},
\ref{lem:bdrytorus},
\ref{obl:taubdry},
\S\ref{sec:ell2},
\ref{lem:ell2unbounded},
\ref{rmk:ell2window},
\ref{open:ell2},
\ref{oc1a}, \ref{oc1b}, \ref{oc1c}, \ref{oc1d},
\S\ref{sec:oc1sites},
\S\ref{sec:cone},
\ref{def:pair},
\ref{def:pairhomotopy},
\ref{lem:etaexists},
\ref{lem:etaunique},
\ref{def:coneaction},
\ref{thm:conewelldef},
\ref{prop:coneindep},
\ref{prop:coneintegral},
\ref{rmk:conestrict},
\ref{prop:coneroutes},
\ref{obl:eta},
\ref{obl:coneexact},
\ref{rmk:conerestrictsmallell},
\S\ref{sec:pair},
\eqref{eq:contractibleuse},
\eqref{eq:conediff},
\ref{lem:conerelative},
\eqref{eq:coneles},
\eqref{eq:conedims},
\ref{def:cosettransport},
\eqref{eq:cosettransport},
\ref{lem:cosettransport},
\S\ref{sec:pairhomotopy},
\eqref{eq:pairhomotopy},
\S\ref{sec:coneaction},
\ref{lem:conesign},
\eqref{eq:conesign},
\eqref{eq:derivedh},
\eqref{eq:conemap},
\ref{lem:derivedhomotopy},
\eqref{eq:conehomotopyeta},
\eqref{eq:conehomotopytau},
\ref{prop:strictpair},
\ref{rmk:conesteinberg},
\ref{prop:coneambiguity},
\ref{rmk:coneambiguity691},
\S\ref{sec:coneintegral},
\ref{rmk:conetorsion},
\S\ref{sec:coneroutes},
\eqref{eq:route3},
\S\ref{sec:coneobligations},
\eqref{eq:sixgroups},
\ref{rmk:sixscope},
\eqref{eq:treeses},
\eqref{eq:mvamalgam},
\ref{lem:h2torsion},
\ref{rmk:bdryrepairs},
\eqref{eq:dimbook},
\eqref{eq:r1formula},
\ref{rmk:chicoincide},
\ref{rmk:pdscope},
\S\ref{sec:orbifoldmodel},
\ref{G1}, \ref{G2}, \ref{G3}, \ref{G4},
\ref{def:borelmodel},
\eqref{eq:borelfree},
\ref{rmk:borelfree},
\ref{thm:specialisation},
\ref{rmk:mpcaveat},
\ref{rmk:orbifoldrefine},
\S\ref{sec:subprimitives},
\ref{rmk:primitiveclosure},
\S\ref{sec:tautable},
\eqref{eq:taurecord},
\ref{rmk:certoverGeometry},
\S\ref{sec:bounds},
\ref{rmk:btabsmall},
\S\ref{sec:backend},
\S\ref{sec:predshapes},
\ref{lem:predshapes},
\eqref{eq:dot2},
\S\ref{sec:filterthm},
\eqref{eq:fperr},
\ref{rmk:thresholdrounding},
\eqref{eq:fptotalcrit},
\S\ref{sec:coreobligations},
\eqref{eq:esoracle},
\S\ref{sec:lateroblig},
\S\ref{sec:obligationroster},
\S\ref{sec:machinery},
\S\ref{sec:notation-integral},
\S\ref{sec:integral},
\S\ref{sec:integralobjects},
\ref{def:ringinforce},
\ref{prop:ringinforce},
\ref{prop:bredon},
\eqref{eq:treemodel},
\S\ref{sec:isomodel},
\ref{def:isomodel},
\ref{IM1}, \ref{IM2},
\ref{prop:isomodel},
\ref{rmk:cafive},
\S\ref{sec:descent},
\ref{prop:fracture},
\ref{prop:cohglue},
\ref{prop:descent},
\ref{prop:homvanish},
\ref{rmk:truncdevice},
\S\ref{sec:strictlift},
\ref{lem:strictlifting},
\ref{prop:strictlift},
\ref{rmk:strictliftbuilt},
\ref{rmk:bredonhecke},
\S\ref{sec:badprimes},
\ref{prop:badprimes},
\ref{prop:certifiede},
\ref{rmk:certifiedeperdiff},
\ref{prop:crtcharpoly},
\S\ref{sec:ringsemantics},
\ref{def:ringsemantics},
\ref{rmk:integralevidence},
\ref{rmk:integralscope},
\ref{rmk:bocksteinexp},
\S\ref{sec:relationsmodel},
\ref{prop:relationsmodel},
\S\ref{sec:notation-dickson},
\S\ref{sec:dicksontheta},
\ref{def:dicksoninv},
\ref{lem:dicksontwist},
\eqref{eq:dicksontwist},
\ref{prop:dicksonH0},
\ref{prop:heckeH0},
\ref{rmk:heckeH0plan},
\ref{lem:heckefibre},
\ref{def:Qn},
\ref{lem:Qnfunctions},
\ref{def:theta},
\ref{prop:thetatwist},
\ref{thm:thetales},
\eqref{eq:thetales},
\ref{rmk:thetasmallell},
\S\ref{sec:thetaliterature},
\ref{rmk:thetaclassical},
\ref{rmk:thetacheckable},
\S\ref{sec:ordinary},
\ref{def:ordinary},
\ref{prop:ordinary},
\ref{cor:thetanonordinary},
\ref{rmk:ordtorsionfree},
\ref{rmk:ordinaryeis},
\ref{open:ordprecision},
\S\ref{sec:hida},
\ref{prop:torsionprepass},
\ref{cor:ordprecision},
\ref{lem:ordexponent},
\ref{rmk:ordexponentuse},
\ref{def:hidatower},
\ref{rmk:hidacheckable},
\S\ref{sec:hidafull},
\S\ref{sec:levelshapiro},
\ref{def:levelellmodel},
\ref{prop:levelcorrespondence},
\eqref{eq:levelcorr},
\ref{prop:levelellU},
\eqref{eq:Ulevell},
\eqref{eq:Ulevellmodel},
\ref{prop:oldspace},
\ref{prop:levelorbit},
\S\ref{sec:levelcost},
\ref{rmk:levelcost},
\ref{rmk:levelscaling},
\S\ref{sec:stabilisation},
\ref{lem:levelelltorsion},
\ref{prop:stabilisation},
\ref{rmk:stabconvention},
\ref{thm:hidacontrol},
\ref{thm:hidaweight},
\S\ref{sec:hidalatticesec},
\ref{def:hidalattice},
\ref{rmk:hidacheckable2},
\ref{lem:rhorank},
\S\ref{sec:levelN},
\S\ref{sec:cosettableN},
\ref{def:cosetrepsN},
\ref{lem:cosetrepsN},
\S\ref{sec:levelNpres},
\ref{def:levelNmodel},
\ref{lem:treeres},
\ref{lem:manin},
\ref{def:kappaN},
\ref{lem:levelduality},
\S\ref{sec:heilbronnsec},
\ref{lem:merelpath},
\ref{prop:heilbronn},
\eqref{eq:heilbronn},
\ref{rmk:heilbronnconv},
\S\ref{sec:heilbronncostsec},
\ref{rmk:heilbronncost},
\S\ref{sec:levelNpredsec},
\ref{rmk:levelNcheckable},
\S\ref{sec:levelNdegsec},
\ref{prop:kappamodel},
\ref{prop:degenpres},
\eqref{eq:degenpres},
\ref{rmk:degenpresconv},
\ref{rmk:levelNscale},
\ref{rmk:levelNcheckabletwo},
\ref{LN7},
\ref{LN8},
\ref{LN9},
\S\ref{sec:levelNcurve},
\ref{prop:levelNcurve},
\ref{rmk:levelNcurvefamily},
\S\ref{sec:nebentypus},
\S\ref{sec:nebdatasec},
\ref{def:nebdata},
\ref{lem:nebdata},
\ref{lem:teichmuller},
\ref{rmk:nebfield},
\S\ref{sec:nebmodelsec},
\ref{lem:inducedchar},
\eqref{eq:MHchi},
\ref{lem:gammaonetable},
\ref{def:nebmodel},
\ref{prop:nebdecomp},
\S\ref{sec:nebheckesec},
\ref{prop:nebcorr},
\eqref{eq:nebcorr},
\ref{thm:nebrelations},
\S\ref{sec:neborbitsec},
\ref{prop:neborbit},
\ref{rmk:nebheilbronn},
\S\ref{sec:nebcostsec},
\ref{rmk:nebcost},
\S\ref{sec:nebpredsec},
\ref{rmk:nebcheckable},
\S\ref{sec:nebwildsec},
\ref{def:nebwild},
\ref{lem:nebwild},
\ref{def:nebgammaH},
\ref{prop:nebgammaH},
\ref{prop:nebgammaHheilbronn},
\ref{lem:nebshapiro},
\ref{lem:nebwildtors},
\ref{lem:nebchainrref},
\ref{rmk:nebwildcheckable},
\ref{rmk:nebcosetlevel},
\S\ref{sec:nebles},
\ref{def:levelpair},
\ref{lem:nebcompat},
\ref{def:levelres},
\eqref{eq:levelres},
\ref{lem:levelresmorph},
\ref{def:levelcone},
\ref{lem:levelconeexact},
\eqref{eq:levelsixterm},
\ref{prop:levelconeshapiro},
\ref{prop:modeltorsion},
\ref{prop:levelsplit},
\ref{prop:levelconehecke},
\ref{def:normalpair},
\ref{prop:levelinfres},
\eqref{eq:levelfiveterm},
\ref{prop:cycliccone},
\ref{lem:levelfreeeuler},
\ref{lem:unipotentinv},
\ref{prop:gamma0normal},
\ref{prop:hidanormalsteps},
\ref{rmk:normalN1},
\ref{obl:levelconeexact},
\ref{rmk:lescheckable},
\ref{rmk:oldspacecone},
\S\ref{sec:saturation},
\S\ref{sec:satsec},
\ref{def:satdata},
\ref{prop:saturation},
\ref{prop:satdefect},
\S\ref{sec:satlevelone},
\ref{lem:heckerankone},
\ref{lem:freepart},
\ref{cor:stablecert},
\S\ref{sec:cohformssec},
\ref{prop:cohforms},
\S\ref{sec:hidaalgsec},
\ref{prop:hidaalgconst},
\S\ref{sec:thetaformssec},
\ref{rmk:thetaforms},
\S\ref{sec:sattorsec},
\ref{rmk:sattor},
\S\ref{sec:isomodelell},
\ref{def:isomodelell},
\ref{prop:isomodelell},
\ref{prop:ellthreeclosed},
\ref{rmk:isomodellmp},
\ref{rmk:isomodellcheckable},
\ref{prop:certifiedeell},
\S\ref{sec:descentN},
\ref{prop:descentN},
\ref{prop:badprimesN},
\ref{prop:descentNthree},
\ref{prop:descentmaps},
\eqref{eq:descentmapsles},
\S\ref{sec:isomodeltree},
\ref{def:isomodeltree},
\ref{prop:isomodeltree},
\S\ref{sec:descenthecke},
\ref{prop:strictliftN},
\ref{prop:crtcharpolyN},
\ref{prop:torsionhecke},
\ref{thm:descentcert},
\ref{rmk:descentcertscope},
\S\ref{sec:descentpair},
\ref{def:descentpair},
\ref{lem:pairlifting},
\ref{prop:descentpair},
\eqref{eq:descentpairles},
\ref{prop:descentpairhecke},
\ref{rmk:descentpairscope},
\S\ref{sec:treepair},
\ref{def:treepair},
\ref{lem:treepaircyl},
\ref{prop:treepair},
\ref{rmk:treepairscope},
\ref{IL1},
\ref{IL2},
\ref{IL3},
\ref{IL4},
\ref{IL5},
\S\ref{sec:mpmodel},
\S\ref{sec:mptrisec},
\ref{def:mptri},
\ref{def:mprules},
\ref{def:mpell},
\ref{def:mpRZ},
\ref{tab:mprules},
\ref{rmk:mprulesread},
\S\ref{sec:mpsetsec},
\ref{def:mpset},
\ref{rmk:mpwhat},
\S\ref{sec:mpcompsec},
\ref{def:mpfree},
\ref{lem:mprows},
\ref{lem:mpcols},
\ref{lem:mpassembly},
\ref{thm:mpcomparison},
\ref{rmk:mphybrid},
\S\ref{sec:mpallsec},
\ref{def:mpphi},
\ref{prop:mpalldegrees},
\S\ref{sec:mpmorphsec},
\ref{def:mpmorphism},
\ref{thm:mpnatural},
\ref{prop:mpbredon},
\ref{prop:mpboundary},
\ref{prop:mprescor},
\ref{rmk:mphecke},
\ref{ex:mpteardrop},
\S\ref{sec:mpchecksec},
\ref{rmk:mpcheckable},
\ref{MPC1},
\ref{MPC2},
\ref{MPC3},
\ref{MPC4},
\ref{MPC5},
\ref{MPC6},
\ref{MPC7},
\ref{MPC8},
\ref{MPC9},
\ref{MPC10},
\ref{MPC11},
\ref{MPC12},
\S\ref{app:errata},
\S\ref{app:status},
\S\ref{app:formalization},
\S\ref{app:labelindex},
\S\ref{app:changelog}.

\subsection*{Errata entries of Appendix~\ref{app:errata}}
\S\ref{err:e1},
\S\ref{err:e2},
\S\ref{err:e3},
\S\ref{err:e4},
\S\ref{err:e5},
\S\ref{err:e6},
\S\ref{err:e7},
\S\ref{err:e8},
\S\ref{err:e9},
\S\ref{err:e10},
\S\ref{err:e11},
\S\ref{err:e12},
\S\ref{err:e13},
\S\ref{err:e14},
\S\ref{err:e15},
\S\ref{err:e16},
\S\ref{err:e17},
\S\ref{err:e18},
\S\ref{err:e19},
\S\ref{err:e20},
\S\ref{err:e21},
\S\ref{err:e22}.

%----------------------------------------------------------------------
\section{Changelog}\label{app:changelog}
%----------------------------------------------------------------------
% SOURCE:  the ten commits that built this document (91cf200, 7f1f5f2, 1c84f7c,
%          6d8abae, ef87af6, 3560525, 592c6e8, 24c235e, d5a2b2d, 3e5ffbc), read
%          with `git show', cross-checked against a label census of every
%          revision.  Nothing is recorded here that a commit message or the
%          document itself does not evidence.
% NOTATION: labels are printed as literal strings; every one is also
%          dereferenced invisibly, so a mistyped label fails the build.
% ERRATA:  none.  This appendix records this document's own history;
%          Appendix~\ref{app:errata} records the defects of the documents it
%          supersedes.  The two are disjoint by construction (CL2).

\newcommand{\PNEW}{\textsf{[new]}}
\newcommand{\PRH}{\textsf{[re-homed]}}
\newcommand{\PFIX}{\textsf{[repaired]}}
% \cl{label}: print the label verbatim AND dereference it into a discarded box,
% so a label named here that this document does not define is an undefined
% reference and therefore a build failure (build_gl2_contract.sh).
\newcommand{\cl}[1]{\texttt{\detokenize{#1}}\begingroup\setbox0\hbox{\ref{#1}}\endgroup}

\paragraph{What this appendix fixes.}
The record that rule~\ref{V3} requires and that this document did not have for its
first ten versions: one entry per version, and within it one line per statement whose
content changed, naming the label. It also fixes what a version \emph{is} for a
document still under construction (\ref{V5}), and marks which statements are new
mathematics rather than transported from a superseded document.

\begin{contractbox}
\textbf{Versions 0.1.0--0.10.0 are assigned retrospectively.} The convention of
\S\ref{sec:versioning} was written in v0.2.0 and then honoured by nobody, this appendix
included, until v0.11.0. The first eleven versions below are therefore a numbering imposed
\emph{now} on the ten commits that built the document, reconstructed from those commits
and checked against the document, not a series of releases that were cut at the time.
Where a change is not evidenced by a commit or visible in the text, it is not recorded
here, and \S\ref{app:changelog-gaps} says what was left out.
\end{contractbox}

\begingroup\raggedright\setlength{\parskip}{0.5ex plus 0.2ex}
\subsection*{Release record}
\addcontentsline{toc}{subsection}{Release record}
This is the supersession record of rule~\ref{V1}: one entry per version, newest first. Up to v1.25.0-pre.7 it stood in the date line of the
title page; the entries from v1.25.0-pre.7 down to v0.1.0, the sentence naming the superseded documents and the companion pointer that
follow them are that date line's text, verbatim. The changelog proper, one entry per version with one line per changed statement
(rule~\ref{V3}), follows the rules below.

Version 1.26.4 --- 2026-11-10 (hecke-extensions slot~26, \textbf{signed by the owner 2026-11-10}: the owner-signed release of v1.26.4-pre.1, title page only; header, date and changelog only --- no statement's mathematics, label, number, tag, status, written word or marker count ($229$ on $88$ labels) changes; see the changelog)

Version 1.26.4-pre.1 --- 2026-11-10 (hecke-extensions slot~26, a \textbf{pre-release, not signed}: title page only, no change of mathematical content: the date line of the title page names only the version, the date and this appendix; the sign-off record of every version stays in this record and in the changelog.)

Version 1.26.3 --- 2026-11-09 (hecke-extensions slot~26, \textbf{signed by the owner 2026-11-09}: the owner-signed release of v1.26.3-pre.1; header, date and changelog only --- no statement's mathematics, label, number, tag, status, written word or marker count ($229$ on $88$ labels) changes; see the changelog)

Version 1.26.3-pre.1 --- 2026-11-09 (hecke-extensions slot~26, a \textbf{pre-release, not signed}: text-only: private path removed, no change of mathematical content.)

Version 1.26.2 --- 2026-11-09 (hecke-extensions slot~26, \textbf{signed by the owner 2026-11-09}: the owner-signed release of v1.26.2-pre.1; header, date and changelog only --- no statement's mathematics, label, number, tag, status, written word or marker count ($229$ on $88$ labels) changes; see the changelog)

Version 1.26.2-pre.1 --- 2026-11-09 (hecke-extensions slot~26, a \textbf{pre-release, not signed}: paths only, no change of content (repository layout, owner ruling 2026-11-09).)

Version 1.26.1 --- 2026-11-09 (hecke-extensions slot~26, \textbf{signed by the owner 2026-11-09}: the owner-signed release of v1.26.1-pre.1; header, date and changelog only --- no statement's mathematics, label, number, tag, status, written word or marker count ($229$ on $88$ labels) changes; see the changelog)

Version 1.26.1-pre.1 --- 2026-11-09 (hecke-extensions slot~26, a \textbf{pre-release, not signed}: no change of mathematical content: machine names and private paths removed for the public release (owner ruling 2026-11-09).)

Version 1.26.0 --- 2026-11-09 (hecke-extensions slot~26, \textbf{signed by the owner 2026-11-09}: the owner-signed release of v1.26.0-pre.11; header, date and changelog only --- no statement's mathematics, label, number, tag, status, written word or marker count ($229$ on $88$ labels) changes; see the changelog)

Version 1.26.0-pre.11 --- 2026-11-09 (hecke-extensions slot~26, UP-Z~(b), the condition C1 and should-fix W-A, W-B of the confirm-referee of v1.26.0-pre.10; \textbf{a pre-release, not owner-released}: status sentences of the descent and HeckeZ chapters and item~8~(d$'$)(ii), (vii) of \S\ref{sec:recordschema}, text only (binding changes \textbf{not yet seen by the owner}); no written word, golden entry, level-$1$ or \code{anemic} byte or label moves; no marker ($229$ on $88$) moves; see the changelog);
Version 1.26.0-pre.10 --- 2026-11-08 (hecke-extensions slot~26, UP-Z~(b), the condition E1 of the referee of v1.26.0-pre.8 / pre.9; \textbf{a pre-release, not owner-released}: item~8~(d$'$)(ii), (vii) of \S\ref{sec:recordschema}, text only
(binding changes \textbf{not yet seen by the owner}); no written word, golden entry, level-$1$ or \code{anemic} byte or label moves; no marker ($229$ on $88$) moves; see the changelog);
Version 1.26.0-pre.9 --- 2026-11-08 (hecke-extensions slot~26, UP-Z~(b), Z2's words, contract first; \textbf{a pre-release, not owner-released, not refereed}: item~8~(d$'$)(ii), (v), (vii) of \S\ref{sec:recordschema}, text only
(binding changes \textbf{not yet seen by the owner}); no golden entry, level-$1$ or \code{anemic} byte or label moves; no marker ($229$ on $88$) moves; see the changelog)

Version 1.26.0-pre.8 --- 2026-11-08 (hecke-extensions slot~26, UP-Z; \textbf{a pre-release, not owner-released}; the conditions of the confirm-referee of v1.26.0-pre.7,
\file{docs/implementation/slot26/C126pre7\_confirm\_referee.md}: item~8~(c), (d$'$)(vii) of \S\ref{sec:recordschema}, text only (binding changes \textbf{not yet seen by the owner}); no record word, hash input,
golden entry or label moves; no marker ($229$ on $88$) moves; see the changelog)

Version 1.26.0-pre.7 --- 2026-11-08 (hecke-extensions slot~26, UP-Z; \textbf{a pre-release, not owner-released}; the conditions of the confirm-referee of v1.26.0-pre.6,
\file{docs/implementation/slot26/C126pre6\_confirm\_referee.md}: item~8~(c), (d), (d$'$)(ii), (vii), (ix) of \S\ref{sec:recordschema}, text only (binding changes \textbf{not yet seen by the owner}); no record word, hash input,
golden entry or label moves; no marker ($229$ on $88$) moves; see the changelog)

Version 1.26.0-pre.6 --- 2026-11-08 (hecke-extensions slot~26, UP-Z; \textbf{a pre-release, not owner-released}; the conditions of the referee-light of v1.26.0-pre.5,
\file{docs/implementation/slot26/C126pre5\_referee\_light.md}: item~8~(c), (d), (d$'$)(ii), (v), (vii), (ix) of \S\ref{sec:recordschema}, text only (binding changes \textbf{not yet seen by the owner}); no record word, hash input,
golden entry or label moves; no marker ($229$ on $88$) moves; see the changelog)

Version 1.26.0-pre.5 --- 2026-11-08 (hecke-extensions slot~26, UP-Z; \textbf{a pre-release, not owner-released}; the text aligned with the audited Z1 code, should-fix S1 (and S2's key)
of \file{docs/implementation/slot26/Z1\_code\_audit.md}: item~8~(d$'$)(ii), (v), (vii) of \S\ref{sec:recordschema} (binding changes \textbf{not yet seen by the owner}); no label
is added, renamed or removed; no marker ($229$ on $88$) moves; see the changelog)

Version 1.26.0-pre.4 --- 2026-11-08 (hecke-extensions slot~26, UP-Z; \textbf{a pre-release, not owner-released}; the should-fix items of the confirm-referee of v1.26.0-pre.3
(ACCEPT): item~8~(d), (d$'$)(ii) of \S\ref{sec:recordschema} (W3, W4, W5, N1; binding changes \textbf{not yet seen by the owner}); no label is added, renamed or removed; no marker
($229$ on $88$) moves; see the changelog)

Version 1.26.0-pre.3 --- 2026-11-08 (hecke-extensions slot~26, UP-Z; \textbf{a pre-release, not owner-released}; the conditions of the adversarial re-referee of v1.26.0-pre.2
(ACCEPT-with-conditions): item~8~(c), (d), (d$'$)(ii) of \S\ref{sec:recordschema} (R1, R2, W1, W2; binding changes \textbf{not yet seen by the owner}); no label is added, renamed or
removed; no marker ($229$ on $88$) moves; see the changelog)

Version 1.26.0-pre.2 --- 2026-11-08 (hecke-extensions slot~26, UP-Z; \textbf{a pre-release, not owner-released}; the fixes of the adversarial referee of v1.26.0-pre.1
(ACCEPT-with-conditions): \cl{prop:strictliftN}(9) (S1, E1--E4, E7--E9) and item~8~(d$'$) of \S\ref{sec:recordschema} (S3, S4, E5, E6; S4 a binding change \textbf{not yet seen by the
owner}); no label is added, renamed or removed; no marker ($229$ on $88$) moves; see the changelog)

Version 1.26.0-pre.1 --- 2026-11-08 (hecke-extensions slot~26, UP-Z, contract first; \textbf{a pre-release, not owner-released, not yet refereed}; binding changes \textbf{not yet seen by the owner}:
\cl{prop:strictliftN}(9) ($U_p$, $p\mid N$, on the free model) and item~8~(d$'$) of \S\ref{sec:recordschema} (the full-$N$ reading of the caps text, frozen with v1); no label is added, renamed or
removed; no marker ($229$ on $88$) moves; see the changelog)

Version 1.25.0 --- 2026-11-07 (hecke-extensions slot~25, \textbf{signed by the owner 2026-11-07}: the owner-signed release of v1.25.0-pre.8; header, date and changelog only --- no statement's mathematics, label, number, tag, status or marker count ($229$ on $88$ labels) changes; see the changelog)

Version 1.25.0-pre.8 --- 2026-11-07 (hecke-extensions slot~25, LaTeX polish stage~5c; \textbf{a pre-release, not owner-released}; one binding change,
\textbf{not yet seen by the owner}, approved by the owner in advance (ruling Q1 of 2026-11-06): rule~\ref{V1} amended, so that this record is the first
subsection of Appendix~\ref{app:changelog} and no longer the date line; otherwise no change of content: a one-line date, the uniform author block, a
generic abstract, the old date line and abstract moved here verbatim, two part headings; no label, number, statement, tag set or \verb|\LEAN| marker
($229$ on $88$) moves; see the changelog)

Version 1.25.0-pre.7 --- 2026-11-06 (hecke-extensions slot~25, the HZ-RAT-PERF fix round after its audit; \textbf{a pre-release, not owner-released}; item~8 only, binding changes \textbf{not yet seen by the owner}: the pending word \code{refused(word\_prime\_cap)} (the word-prime count cap $2^{21}$), the closed list of \code{C2\_status} as built, the cap in the unhashed provenance key; no label, tag set, marker, record or golden line moves; see the changelog);

Version 1.25.0-pre.6 --- 2026-11-06 (hecke-extensions slot~25, the H1 merge after its audit; \textbf{a pre-release, not owner-released}; item~8 only, binding changes \textbf{not yet seen by the owner}: the G8 paragraph (the Hida block under \code{heckez.G8 = hida}: paths, closed lists, caps), golden only under the default settings; no label, tag set or marker moves; see the changelog);

Version 1.25.0-pre.5 --- 2026-11-06 (hecke-extensions slot~25, the HZ-REC schema re-audit fix round; \textbf{a pre-release, not owner-released}; item~8 only, binding changes \textbf{not yet seen by the owner}: the Eisenstein multiplicity $e_k$, the cuspidal block iff $d\ge2$, the $\vartheta$ search over $\mathcal P\cup\{T_\ell\}$, the second route's caps and \code{not\_read}, the irreducibility cap $256$, \code{[Lift]} by its certificate, the \code{not\_run} words of v1; no label, tag set or marker moves; see the changelog);

Version 1.25.0-pre.4 --- 2026-11-06 (hecke-extensions slot~25, the HZ-REC schema-audit fix round; \textbf{a pre-release, not owner-released}; item~8 only, binding changes \textbf{not yet seen by the owner}: the new final value \code{split\_on\_P}, G4--G6 per torsion point, no G8 in default records, the components, maximal ideals and caps text as written; no label, tag set or marker moves; see the changelog);

Version 1.25.0-pre.3 --- 2026-11-06 (hecke-extensions slot~25, the Z0 re-check; \textbf{a pre-release, not owner-released}; text only: \code{export\_word} also for $K$ above EXPORT's $K_{\max}$ below the word bound; no label, tag set, marker or value moves; see the changelog);

Version 1.25.0-pre.2 --- 2026-11-06 (hecke-extensions slot~25, the Z0 merge after its referee (REVISE); \textbf{a pre-release, not owner-released}; text only: the counting criterion of the full-algebra tag narrowed to $N$ a power of $p$, \code{export\_word} pending, the as-built refusal rule, the cofactor route as built, the closed \code{t\_ell} list; no label, tag set, marker or value moves; see the changelog);

Version 1.25.0-pre.1 --- 2026-11-06 (hecke-extensions slot~25, the Z0 merge; \textbf{a pre-release, not owner-released}, on top of the signed v1.24.0: the refereed \S\ref{sec:heckeZ} of v1.25.0-pre.0c (referees A7, A7b: ACCEPT) enters the main line with record item~8 of \S\ref{sec:recordschema}, reconciled with the code of slot~25 by dated notes; no marker ($229$ on $88$ labels) or value moves; see the changelog);

Version 1.24.0 --- 2026-11-05 (hecke-extensions slot~25, \textbf{signed by the owner 2026-11-05}: the owner-signed release of v1.24.0-rc.2; header, date and changelog only --- no statement's mathematics, label, number, tag, status or marker count ($229$ on $88$ labels) changes; see the changelog);

Version 1.24.0-rc.2 --- 2026-11-05 (hecke-extensions slot~25, \textbf{the release candidate of 1.24.0: not owner-released}; rc.1 refereed ACCEPT\_WITH\_MINORS, \texttt{docs/implementation/slot25/RC124\_referee.md}, its findings applied as text: the marker rule of \S\ref{sec:status} amended by one sentence, Remarks~\ref{rmk:bocksteinexp}, \ref{rmk:levelNcurvefamily} and Proposition~\ref{prop:ltwothreescalar} corrected; no label, tag set, marker or value moves; see the changelog);

Version 1.24.0-rc.1 --- 2026-11-05 (hecke-extensions slot~25, \textbf{the release candidate of 1.24.0: not owner-released, not yet refereed}, on top of v1.24.0-pre.1: the Lean markers on Theorem~\ref{thm:classical} and \eqref{eq:corconjres} ($229$ on $88$ labels), the phenomenology citations in Remarks~\ref{rmk:bocksteinexp}, \ref{rmk:levelNcurvefamily}, Proposition~\ref{prop:ltwothreescalar} and \ref{Hi11}; no label or tag set moves; see the changelog);

Version 1.25.0-pre.0c --- 2026-11-04 (hecke-extensions slot~24, HeckeZ Z0 after the A7b re-audit; \textbf{still a branch pre-release} on \texttt{s24-wip-heckez-z0}, \textbf{not owner-released}: the drop rule of Theorem~\ref{thm:heckeZcert}(C2) and the (C2) frontier measured on one discriminant (A7b-m1); no label, number, tag set or marker moves; see the changelog);

Version 1.25.0-pre.0b --- 2026-11-04 (hecke-extensions slot~24, HeckeZ Z0 fix round after the adversarial referee A7; \textbf{still a branch pre-release} on \texttt{s24-wip-heckez-z0}, \textbf{not owner-released}: the Mazur route of Proposition~\ref{prop:heckeZeiscert}(3) is certified by a valuation check, goodness of $q$ alone being insufficient (A7-M1); the referee's seven MINOR findings and the cheap NOTEs; no label, number, tag set or marker moves; see the changelog);

Version 1.25.0-pre.0 --- 2026-11-04 (hecke-extensions slot~24, HeckeZ package Z0; \textbf{a branch pre-release} on \texttt{s24-wip-heckez-z0}, \textbf{not owner-released, not refereed}; it enters the main line as v1.25.0-pre.1 only after the owner has released v1.24.0; on top of v1.24.0-pre.1: the new \S\ref{sec:heckeZ}, the Hecke algebra over $\Z$ (23 statements appended at the end of their section, no number moves), item~8 of \S\ref{sec:recordschema}, and a pointer in Theorem~\ref{thm:descentcert}(7)(a); no marker moves; see the changelog);

Version 1.24.0-pre.1 --- 2026-11-03 (hecke-extensions slot~24, \textbf{a pre-release, not owner-released, not refereed}, on top of v1.23.0: FMTMP, \S\ref{sec:recordschema} item~7 (l), \emph{The faithful coefficients}; text only, no label moves; see the changelog);

Version 1.23.0 --- 2026-11-03 (hecke-extensions slot~24, \textbf{signed by the owner 2026-11-03}: the owner-signed release of v1.23.0-pre.40; header, date and changelog only --- no statement's mathematics, label, number, tag, status or marker count ($222$ on $86$ labels) changes; see the changelog);

Version 1.23.0-pre.40 --- 2026-11-02 (slot 24, \textbf{a pre-release, not owner-released, not refereed}, on top of v1.23.0-pre.39: ICENSUS-P2, \S\ref{sec:recordschema} item~7 (k), \emph{Parts not run}; text only, no label moves; see the changelog);

Version 1.23.0-pre.39 --- 2026-11-02 (hecke-extensions slot 24, \textbf{a pre-release, not owner-released, not refereed}, on top of v1.23.0-pre.38: ICENSUS-P1, \S\ref{sec:recordschema} item~7 (e), (j), the $\ker/\coker$ read's width ladder; no label, number, tag set or marker moves; see the changelog);

Version 1.23.0-pre.38 --- 2026-10-28 (hecke-extensions slot 24, \textbf{a pre-release, not owner-released, not re-refereed}, on top of v1.23.0-pre.37: the pre.37
referee's fixes M1, M2, m1--m10, N1--N4, T1, T2 (\file{docs/implementation/slot24/CON\_pre37\_referee.md}, PASS-with-fixes), text only --- item~7 of \S\ref{sec:recordschema}: the coefficient order
(leading coefficient first, an exception to item~2), the digest's byte layout, the hash partition, and the code map as merged (\texttt{2aa61397}) with the golden event landed (\texttt{6958a5d2}); no label, number,
tag set or marker moves; see the changelog);

Version 1.23.0-pre.37 --- 2026-10-26 (hecke-extensions slot 24, \textbf{a pre-release, not owner-released, not refereed}, on top of v1.23.0-pre.36: the integral descent record
\code{integral\_descent\_record/v1} gets its contract home, item~7 of \S\ref{sec:recordschema} (D6, built on branch \texttt{s24-d6impl}, pending audit and merge; its golden event pending the owner's go);
$\Z$ un-reserved for the descent only (Definition~\ref{def:ringsemantics}(3)); the $\ker/\coker$ precision refusal renamed \code{[KerCokerPrecision]}; no label, number, tag set or marker moves; see the changelog);

Version 1.23.0-pre.36 --- 2026-10-26 (hecke-extensions slot 24, \textbf{a pre-release, not owner-released, not re-refereed}, on top of v1.23.0-pre.35: the pre.35
referee's fixes m1, N1--N7 (\file{docs/implementation/slot24/C-D8b\_referee.md}, ACCEPT after fixes) and D5c \textbf{merged} (\texttt{4064e164}) --- \S\ref{sec:treepair}'s certificates restated for
the merged code (the applied operator's blocks, entrywise or on the generators and $64$ random vectors), Proposition~\ref{prop:treepair}(4)(e) sharpened, the code map's status; no label, number,
tag set or marker moves; see the changelog);

Version 1.23.0-pre.35 --- 2026-10-26 (hecke-extensions slot 24, \textbf{a pre-release, not owner-released, not refereed}, on top of v1.23.0-pre.34: \S\ref{sec:treepair} brought in line with
the D5c code (branch \texttt{s24-d5cimpl}, pending audit) --- $S$ versus $S^{-1}$ on cochains, the same for $n$ even; a truncation statement, Proposition~\ref{prop:treepair}(4)(e), by which $q_{\max}=5$ is read with the
comparisons to degree $5$; the code map's status, certificates and controls; no label, number, tag set or marker moves; see the changelog);

Version 1.23.0-pre.34 --- 2026-10-26 (hecke-extensions slot 24, \textbf{a pre-release, not owner-released, not re-refereed}, on top of v1.23.0-pre.33: the pre.33
referee's fixes m1--m4, N1--N8 (\file{docs/implementation/slot24/C-D8\_referee.md}, ACCEPT after fixes) --- in \S\ref{sec:treepair} the off-diagonal $\eta_{\mathcal T}$ arises on $\Cone{\mathsf a^{\ast}}$, not on the pair model,
and its non-vanishing is observed, not proved; the solver and (L), (L$'$) for any column augmentation; the code's negative controls restricted where they are vacuous; the $\ell$-local
column augmentations defined; no label, number, tag set or marker moves; see the changelog);

Version 1.23.0-pre.33 --- 2026-10-26 (hecke-extensions slot 24, \textbf{a pre-release draft, not owner-released, not refereed}, on top of v1.23.0-pre.32: \S\ref{sec:treepair}, the tree cusp pair
(plan milestone D5c, clause C-D8), contract first --- the cusp model and the axis map in closed form, the tree cone and the cylinder lemma, effective comparisons with $(\Iso,\Isob)$, $T_p$ by transport with
its off-diagonal term; the D3 merge recorded in \S\ref{sec:descenthecke}'s code map; four statements added at the end of their section, no number, tag set or marker moves; see the changelog);

Version 1.23.0-pre.32 --- 2026-10-26 (hecke-extensions slot 24, \textbf{a pre-release, not owner-released, not re-refereed}, on top of v1.23.0-pre.31: the pre.31
referee's fixes m1--m3, N1--N6 (\file{docs/implementation/slot24/C-D3h\_referee.md}, ACCEPT after fixes) --- in \S\ref{sec:descenthecke} the homotopy of two lifts is $\Theta_{h_\Iso}$ with no sign,
the precision of Theorem~\ref{thm:descentcert}(4) at $\ell\ge5$ is $K>\sigma_{\max}+v_\ell(\chi(c))$, any exact lifts of the free basis are a basis; no label, number, tag set or marker moves; see the changelog);

Version 1.23.0-pre.31 --- 2026-10-26 (hecke-extensions slot 24, \textbf{a pre-release, not owner-released, not refereed}, on top of v1.23.0-pre.30: three clauses of \S\ref{sec:descenthecke} close
the contract notes of D3 (Hecke over $\Z$, built on branch \texttt{s24-d3impl}) --- the torsion read on a cokernel presentation (Proposition~\ref{prop:torsionhecke}(3$'$)), admitted at $\ell\ge5$ and at the residue primes; obligation~(b)
met through the lift; any admissible bulk lift, the read independent of it (Proposition~\ref{prop:strictliftN}(2)); the code map records D3 as built on its branch; no number, tag set or marker moves; see the changelog);

Version 1.23.0-pre.30 --- 2026-10-26 (hecke-extensions slot 24, \textbf{a pre-release, not owner-released, not re-refereed}, on top of v1.23.0-pre.29: a non-mathematical code-map update --- D5 (the tree models) is built and merged
(\S\ref{sec:isomodeltree}), and \code{hecke::HeckeInternalError} separates internal errors, which fail, from a real overflow, which refuses as \code{IntegerCap} (\S\ref{sec:descentpair}); no statement, label, number, tag set or marker moves; see the changelog);

Version 1.23.0-pre.29 --- 2026-10-25 (hecke-extensions slot 24, \textbf{a pre-release, not owner-released, not re-refereed}, on top of v1.23.0-pre.28: the pre.28
referee's fixes M1, m1--m4, N1 (\file{docs/implementation/slot24/C-D7p\_referee.md}, ACCEPT-with-fixes) --- in \S\ref{sec:descentpair} the refusals of $T_p$ on the pair are exactly the
code's eleven refusal codes, the forced component map proved in (1)'s proof, $q_{\max}\le5$; the code map records D4 as merged; no label, number, tag set or marker moves; see the changelog);

Version 1.23.0-pre.28 --- 2026-10-25 (hecke-extensions slot 24, \textbf{a pre-release, not owner-released}, refereed ACCEPT-with-fixes, on top of v1.23.0-pre.27 (C-D7 fixes applied): the owner's
scope addition of 2026-10-25 --- in \S\ref{sec:descentpair} $T_p$ on the boundary and the cone for every prime $p\nmid N$, at level $1$ and level $N$, from the pair lift built on demand and
certified at run time; the code map records D4's build; see the changelog);

Version 1.23.0-pre.27 --- 2026-10-25 (hecke-extensions slot 24, \textbf{a pre-release, not owner-released, not re-refereed}, on top of v1.22.0, \textbf{signed by the
owner 2026-10-16}: the C-D7 referee's fixes M1--M3, m1--m5 (\file{docs/implementation/slot24/C-D7\_referee.md}, ACCEPT-with-fixes) --- in \S\ref{sec:descentpair} $\Hcpt{2}$ at
$\ell\ge5$ is $H_0$, a prime rank source for $\Dcone^{0}$, $T_p$ on the cone certified only where a pair lift is built, the signs of the $\Ext$ identification, elementary periodicity
proofs; see the changelog);

Version 1.23.0-pre.26 --- 2026-10-25 (hecke-extensions slot 24, \textbf{a pre-release draft, not owner-released, not refereed}, on top of v1.22.0, \textbf{signed by the
owner 2026-10-16}: milestone D4 of the descent plan, clause C-D7 --- \S\ref{sec:descentpair} (new, after \S\ref{sec:descenthecke}, so that no statement number moves): the
boundary pair and its cone over $\Z$ on $\Isob\subseteq\Iso$ and $\Iso^{(\ell)}$, glued, the boundary in closed form, compact support and the certified exponent, the cone
with the strict pair datum and $T_p$, $p\nmid N$; see the changelog);

Version 1.23.0-pre.25 --- 2026-10-25 (hecke-extensions slot 24, \textbf{a pre-release, not owner-released, not re-refereed}, on top of v1.22.0, \textbf{signed by the
owner 2026-10-16}: the C-D4 referee's fixes M1, M2, m1--m5, n1, t1--t3 (\file{docs/implementation/slot24/C-D4\_referee.md}, ACCEPT-with-fixes) --- in \S\ref{sec:descenthecke}
the layer blocks choice-free modulo $\ell^{\min(j,d_j)}$, one route at $\ell=3$ (the transport on $\Iso^{(3)}$, no Bredon operator), the necessity of $2\sigma_{\max}\le K$, the
word conditions of the recovery; see the changelog);

Version 1.23.0-pre.24 --- 2026-10-25 (hecke-extensions slot 24, \textbf{a pre-release draft, not owner-released, not refereed}, on top of v1.22.0, \textbf{signed by the
owner 2026-10-16}: milestone D3 of the descent plan, clauses C-D4 and C-D5 --- \S\ref{sec:descenthecke} (new, after \S\ref{sec:isomodeltree}, so that no statement number
moves): the level-$N$ operator $T_p$, $p\nmid N$, on the free model and its transports, the characteristic polynomial over $\Z$ by residues with the output rule, the
choice-free Hecke invariants of finite $\ell$-groups, and the certified record; see the changelog);

Version 1.23.0-pre.23 --- 2026-10-25 (hecke-extensions slot 24, \textbf{a pre-release, not owner-released, not re-refereed}, on top of v1.22.0, \textbf{signed by the
owner 2026-10-16}: the C-D3g referee's fixes M1, m1--m5, t1, t2 (\file{docs/implementation/slot24/C-D3g\_referee.md}, ACCEPT-with-fixes) --- in Proposition~\ref{prop:descentmaps}(2)(g)
the blocks of the transforms that are exact modulo $\ell^{K}$, the lattice coordinates, the conservative refusal policy, the further operations and the certifying ranks of
derived forms; (2)(b)'s sign convention; editorial fixes in Proposition~\ref{prop:isomodeltree}; no label, number, tag set or marker moves; see the changelog);

Version 1.23.0-pre.22 --- 2026-10-25 (hecke-extensions slot 24, \textbf{a pre-release, not owner-released}, refereed ACCEPT-with-fixes, on top of v1.22.0, \textbf{signed by the
owner 2026-10-16}: the D2 code audit's F1 and F8 (\file{docs/implementation/slot24/D2\_code\_audit.md}) --- Proposition~\ref{prop:descentmaps} gains item (2)(g), the precision rule
for generators, coordinates and maps read from a certified Smith form, and (2)(b) states the signs of the cone's maps; no label, number, tag set or marker moves; see the
changelog);

Version 1.23.0-pre.21 --- 2026-10-25 (hecke-extensions slot 24, \textbf{a pre-release, not owner-released}, the C-D6 referee's fixes M1, m1--m5
(\file{docs/implementation/slot24/C-D6\_referee.md}, ACCEPT-with-fixes; not re-refereed), on top of v1.22.0, \textbf{signed by the owner 2026-10-16}: in
Proposition~\ref{prop:isomodeltree}, (4)'s comparisons with $\Iso$ over $\Z$ and every homotopy between their composites are constructed (the tree inside $\Xbulk$, an explicit
retraction and homotopy, the lifting formula (L)), and five minor repairs; no label, number, tag set or marker moves; see the changelog);

Version 1.23.0-pre.20 --- 2026-10-25 (hecke-extensions slot 24, \textbf{a pre-release draft, not owner-released, not refereed}, on top of v1.22.0, \textbf{signed by the owner 2026-10-16}: milestone D5 of the descent plan, clause C-D6, brought forward by the owner's ruling at gate G-T --- \S\ref{sec:isomodeltree} (new, after the last statement of \S\ref{sec:descentN}, so that no statement number moves): Definition~\ref{def:isomodeltree}, the free and the $\ell$-local tree models, and Proposition~\ref{prop:isomodeltree}: they compute group cohomology in every degree for every coefficient module (level $N$ by Shapiro), one integral form over $\mathsf O_\ell$ with the certified-exponent rule, cochains and comparisons with the free tree in closed form, comparisons with $\Iso$ over $\Z$, the fibre maps and the level cone, the Hecke action by transport; bulk only; no existing label, number, tag set or marker moves; see the changelog);

Version 1.23.0-pre.19 --- 2026-10-25 (hecke-extensions slot 24, \textbf{a pre-release, not owner-released}, the C-D3 referee's fixes m1--m5
(\file{docs/implementation/slot24/C-D3\_referee.md}, ACCEPT-with-fixes; not re-refereed), on top of v1.22.0, \textbf{signed by the owner 2026-10-16}: in
Proposition~\ref{prop:descentmaps}, (3)(b)'s failure of the square with restriction is restricted to this document's $U_p$ (the adjoint of the classical one), and the new item
(3)(b$'$): corestriction commutes with $U_p$ for $p\mid N$ and every $M$, over $\Z$ in every degree (the Hida tower); four wording and citation fixes; no number, label or
marker moves; see the changelog);

Version 1.23.0-pre.18 --- 2026-10-25 (hecke-extensions slot 24, \textbf{a pre-release, not owner-released}, refereed ACCEPT-with-fixes (\file{docs/implementation/slot24/C-D3\_referee.md}), on top of v1.22.0, \textbf{signed by the owner 2026-10-16}: milestone D2 of the descent plan, clause C-D3 --- Proposition~\ref{prop:descentmaps}: fibre maps (restriction, trace, conjugation) commute strictly with every comparison and glue to maps over $\Z$; the level cone on every model and its long exact sequence over $\Z$ in every degree, with the certified exponent; the Hecke squares for $p\nmid M$; inflation--restriction over $\Z$ and the two-row case of a torsion-free $B$ (the Hida $p$-tower); what a tree model must supply; one new statement, the last of its section, so no number moves; no marker moves; see the changelog);

Version 1.23.0-pre.17 --- 2026-10-25 (hecke-extensions slot 24, \textbf{a pre-release, not owner-released}, refereed ACCEPT (\file{docs/implementation/slot24/C-D1\_three\_referee.md}), on top of v1.22.0, \textbf{signed by the owner 2026-10-16}: the D1 code audit's finding M1 --- Proposition~\ref{prop:descentNthree}, the $3$-adic shortcut at level $N$ when $\Ga_0(N)$ has no element of order~$3$: the tree complex computes the $3$-primary part, $H^{\ge2}$ has no $3$-torsion, and the sparse cokernel route holds at $\ell=3$; one new statement, the last of its section, so no number moves; no marker moves; see the changelog);

Version 1.23.0-pre.16 --- 2026-10-25 (hecke-extensions slot 24, \textbf{a pre-release, not owner-released}, on top of v1.22.0, \textbf{signed by the owner 2026-10-16}: the Dickson-proof referee's fixes D-1--D-5 (ACCEPT-with-fixes, not re-refereed) --- the proof of Claim~\ref{cl:dimH1}(2) is now Lemma~\ref{lem:dicksonSL2}, valid at every prime $\ell$ and for $\GL_2$ as well, cited by its three users; scope, tag and exposition nits; one new statement, no marker moves, statement numbers after Claim~\ref{cl:dimH1} in its section shift by one; see the changelog);

Version 1.23.0-pre.15 --- 2026-10-25 (hecke-extensions slot 24, \textbf{a pre-release, not owner-released}, on top of v1.22.0, \textbf{signed by the owner 2026-10-16}: the C-D1 / C-D2 referee's fixes F1--F7 to \S\ref{sec:descentN} (ACCEPT-with-fixes, not re-refereed) --- Claim~\ref{cl:dimH1}(2), the Dickson invariants of $\SL_2(\Fl)$, proved in place, so no count of Proposition~\ref{prop:badprimesN} rests on an unverified pointer, and the good-prime rule $\ell>n+1$ proved without it; $b_1\ge1$ for every $N>1$; the sparse cokernel's generators, their certificate, and periodicity at chain level; no label, number or marker moves; see the changelog);

Version 1.23.0-pre.14 --- 2026-10-25 (hecke-extensions slot 24, \textbf{a pre-release draft, not owner-released, not refereed}, on top of v1.22.0, \textbf{signed by the owner 2026-10-16}: milestone D1 of the descent plan, clauses C-D1 and C-D2 --- \S\ref{sec:descentN} (new, after the last statement of \S\ref{sec:isomodelell}, so that no statement number moves): Proposition~\ref{prop:descentN}, the descent for every finite free coefficient module (level $N$, nebentypus), Hecke determined locally, the coherence term only when $b_0b_1=0$, periodicity by $\res(e)\cup$, the good prime and the certified exponent on the direct Smith and sparse cokernel routes; Proposition~\ref{prop:badprimesN}, the closed-form bad-prime list $\mathcal B(n,N)$ with the torsion at $\ell\ge5$ in $H^{1}$ only; no existing label, number, tag set or marker moves; see the changelog);

Version 1.23.0-pre.13 --- 2026-10-25 (hecke-extensions slot 24, \textbf{a pre-release, not owner-released}, on top of v1.22.0, \textbf{signed by the owner 2026-10-16}: the C-D0 referee's fixes F1--F8 to Proposition~\ref{prop:certifiedeell} (ACCEPT-with-fixes, not re-refereed) --- the good prime certified for every $M$ ($24\cdot H^{\ge2}(\Ga;M)=0$; at $q=0$ the $H^{0}$ test against $b_0$), the two admissible sources of $\rank_{\Q}$, the valuations counted with multiplicity, the a-priori exponents for $q\ge1$, and the code map's D0 obligations; no label, number, tag set or marker moves; see the changelog);

Version 1.23.0-pre.12 --- 2026-10-25 (hecke-extensions slot 24, \textbf{a pre-release draft, not owner-released, not refereed}, on top of v1.22.0, \textbf{signed by the owner 2026-10-16}: milestone D0 of the descent plan, clause C-D0 --- Proposition~\ref{prop:certifiedeell} (new, \PROVED{}, at the end of \S\ref{sec:isomodelell}, so that no statement number moves): $\Iso^{(\ell)}$ built once over $\Z[1/\mathsf q_\ell]$, its base changes the complexes the code builds, the certified-exponent rule over $\Z_{(\ell)}$ with $\rank_{\Q}$ read at a good prime, and the Smith-valuation read-out; its code map; no existing statement, label, number, tag set or marker moves; see the changelog);

Version 1.23.0-pre.11 --- 2026-10-21 (hecke-extensions slot 24, \textbf{a pre-release, not owner-released, not refereed}, on top of v1.22.0, \textbf{signed by the owner 2026-10-16}: code-map text only, from two audited packages --- G2-H (Obligation~\ref{obl:welldef}(iii) certified without forming $\Th^{(q)}_{\mathrm{eq}}$ on the sat / Loewy runner's lean route, route~2 compared in place, lever \code{HECKE\_MODULE\_LEAN\_PEAK}) in \S\ref{sec:obligations}'s code map, G6 (the cone certificate's span inclusions by the unit-row route and the field route's memory diet, lever \code{HECKE\_FIELD\_DIET}) in \S\ref{sec:codemap}'s row for \S\ref{sec:cone}; every lever value-neutral; no statement, label, number, tag set or marker moves; see the changelog);

Version 1.23.0-pre.10 --- 2026-10-20 (hecke-extensions slot 24, \textbf{a pre-release, not owner-released, not refereed}, on top of v1.22.0, \textbf{signed by the owner 2026-10-16}: code-map text only, from two audited packages --- FNTT (the forms cross-route's transform-domain route, lever \code{HECKE\_FORMS\_NTT\_CACHE}; the forms $\ell$-range past $2^{16}$ restated, \code{forms::FormsSpaceSupported}) in \S\ref{sec:saturation}'s code map, MT-1 (the module route's pivot-search exit and further pool splits, lever \code{HECKE\_MODULE\_MT}, values independent of the thread count) in \S\ref{sec:obligations}'s; every lever value-neutral; no statement, label, number, tag set or marker moves; see the changelog);

Version 1.23.0-pre.9 --- 2026-10-20 (hecke-extensions slot 24, \textbf{a pre-release draft, not owner-released, not refereed}, on top of v1.22.0, \textbf{signed by the owner 2026-10-16}: package ELL-W --- Lemma~\ref{lem:bounds} gains clause~(7), the wide tier (products of two stored-class matrices formed in $128$-bit integers under $2NB\le2^{63}-1$), Theorem~\ref{thm:backend} gains consequence~(4), Obligation~\ref{obl:backend}(ii) names its assertion; the reach of the $\tau$-certification moves from $\ell\le23170$ to $\ell\le1048573$; no label, existing clause number, tag set or marker moves, and no value at $\ell\le23170$; see the changelog);

Version 1.23.0-pre.8 --- 2026-10-18 (hecke-extensions slot 23, \textbf{a pre-release, not owner-released}, on top of v1.22.0, \textbf{signed by the owner 2026-10-16}: two audited code-map sentences (G2-T2, the SIMD $\rhon$ fill and route~2's shared memo, levers \code{HECKE\_RHO\_SIMD}, \code{HECKE\_ROUTE2\_SHARED\_MEMO}; G2-P, the module route on the thread pool, lever \code{HECKE\_MODULE\_PARALLEL}, values independent of the thread count) in \S\ref{sec:obligations}'s code map; code-map text only, every lever value-neutral, no statement, label, number, tag set or marker moves; see the changelog);

Version 1.23.0-pre.7 --- 2026-10-18 (hecke-extensions slot 23, \textbf{a pre-release, not owner-released}, on top of v1.22.0, \textbf{signed by the owner 2026-10-16}: one audited code-map sentence (G2-T, the module route's term evaluation and its value-neutral rollback lever \code{HECKE\_MODULE\_EVAL=0}) in \S\ref{sec:obligations}'s code map; code-map text only, no statement, label, number, tag set or marker moves; see the changelog);

Version 1.23.0-pre.6 --- 2026-10-18 (hecke-extensions slot 23, \textbf{a repair-only pre-release, not owner-released}, on top of v1.22.0, \textbf{signed by the owner 2026-10-16}: the fixes F1--F6 of the referee report \file{C1230\_pre3\_referee.md} (ACCEPT-WITH-FIXES, MINOR 2, NIT 4) to pre.1--pre.3 --- the proof of Lemma~\ref{lem:precisionreduction} repaired (the $\tau$-homotopy over $\Ceq(\Z_\ell)$), its use of Theorem~\ref{thm:relations}(1) named, $K_{\mathrm{cls}}$, $\mathrm{KMAX}$, $K_{\mathrm{stab}}$ attributed or defined, readings across levels added to (3) as a consequence of (1); otherwise no statement's mathematics, and no label, number, tag set or marker, moves; see the changelog);

Version 1.23.0-pre.5 --- 2026-10-18 (hecke-extensions slot 23, \textbf{a pre-release, not owner-released}; pre.4 refereed (ACCEPT-WITH-FIXES, fixes applied), pre.5 adds one audited code-map sentence (G2-F, \code{ChainsOnly} $\tau$ tables): code-map text only --- \S\ref{sec:satprecisionsec}'s code map built (G1), the module route's G2 re-implementation in \S\ref{sec:obligations}'s code map, the G2-X forms product in \S\ref{sec:saturation}'s, each rollback lever value-neutral; no statement, label, number, tag set or marker moves; see the changelog);

Version 1.23.0-pre.3 --- 2026-10-16 (hecke-extensions slot 23, \textbf{a pre-release, not owner-released}, audited (ACCEPT-WITH-FIXES, fixes applied), on top of v1.22.0, signed by the owner 2026-10-16 (its date line below is left as released): new \S\ref{sec:satprecisionsec} --- Lemma~\ref{lem:precisionreduction} (the precision levels of the module route are reductions of one run) and Obligation~\ref{obl:precisionreuse} (one module run per cell); two new statements, no existing statement, number, tag set or marker moves; see the changelog);

Version 1.23.0-pre.2 --- 2026-10-12 (hecke-extensions slot 23, \textbf{a repair-only pre-release, not owner-released}: the repairs of the referee report F-C123pre1 (ACCEPT; MINOR 2, NIT 5) to Version 1.23.0-pre.1; no label, statement number, tag set or marker moves; see the changelog);

Version 1.23.0-pre.1 --- 2026-10-12 (hecke-extensions slot 23, \textbf{a pre-release, not owner-released; refereed by F-C123pre1}: Proposition~\ref{prop:levelconehecke}(4) replaced --- the level-correspondence operator is independent of the transversal --- and the negative control of Obligation~\ref{obl:levelconeexact}(ix) replaced by the criterion-fails control; no label, statement number, tag set or marker moves; see the changelog);

Version 1.22.0 --- 2026-10-12 (hecke-extensions slot 22, the CLOSE-T release candidate from v1.22.0-pre.17; \textbf{owner sign-off pending}, not yet owner-released; header, date and changelog only --- no statement's mathematics, label, number, tag, status or marker count ($222$ on $86$ labels) changes; see the changelog);

Version 1.22.0-pre.17 --- 2026-10-12 (hecke-extensions slot 22, \textbf{a pre-release, not owner-released}: CLOSE-T T1 repairs --- the F-T1 release referee's MINOR-1 (the changelog of pre.9) and seven wording NITs; no marker, label, statement or value moves; see the changelog);

Version 1.22.0-pre.16 --- 2026-10-12 (hecke-extensions slot 22, \textbf{a pre-release, not owner-released, not refereed}: LR-Y phase 5 --- four \code{\textbackslash LEAN} markers on Claim~\ref{cl:amalgam}(1) (the tree sequence as one short exact sequence of $\Ga$-representations) and (2) (the first map is $(\res_S,\res_U)$, the whole Mayer--Vietoris sequence in group-cohomology terms, naturality in the coefficients); two existing marker scopes shortened; no new statement, no existing statement moves; see the changelog);

Version 1.22.0-pre.15 --- 2026-10-11 (hecke-extensions slot 22, \textbf{a pre-release, not owner-released, not refereed}: F-pre14 referee repairs --- wording only; no marker, label or statement change; see the changelog);

Version 1.22.0-pre.14 --- 2026-10-11 (hecke-extensions slot 22, \textbf{a pre-release, not owner-released, not refereed}: the Lean dimension results at $\ell\in\{2,3\}$ --- twelve \code{\textbackslash LEAN} markers on Lemma~\ref{lem:centralminusone} (its last clause, for fields), Proposition~\ref{prop:oddnltwo}(2) (the value of $h^{0}$ at certified cells), Proposition~\ref{prop:ellthreeclosed}(4) (for fields) and the new Remark~\ref{rmk:lvdimsltwothree} ($\dim H^{q}(\Ga;\Sym^{n}\F_\ell^{2})$, $\ell=2,3$, $n\le700$); one new statement, no existing statement moves; see the changelog);

Version 1.22.0-pre.13 --- 2026-10-11 (hecke-extensions slot 22, \textbf{a pre-release, not owner-released, not refereed}: LR-Y phase 4 --- four \code{\textbackslash LEAN} markers on Claim~\ref{cl:amalgam}(2) (the Mayer--Vietoris terms as group cohomology of the cyclic subgroups, the middle map $\res_S-\res_U$, the beginning) and Lemma~\ref{lem:bulkperiodic}(2) (a natural periodicity isomorphism); two existing marker scopes rewritten; no new statement, no existing statement moves; see the changelog);

Version 1.22.0-pre.12 --- 2026-10-11 (hecke-extensions slot 22, \textbf{a pre-release, not owner-released, not refereed}: LR-Y phases 1--3 --- seven \code{\textbackslash LEAN} markers on Claim~\ref{cl:amalgam}(1),(2), Lemma~\ref{lem:bulkperiodic}(2) and the new Remark~\ref{rmk:treeconeres} (the tree resolution over every ring); the Serre pointer of Claim~\ref{cl:amalgam}(1) no longer offered as a hypothesis; one new statement, no existing statement moves; see the changelog);

Version 1.22.0-pre.11 --- 2026-10-10 (hecke-extensions slot 22, \textbf{a pre-release, not owner-released, not refereed}: the F-L23 referee's repairs in \S\ref{sec:oddnltwo}; \S\ref{sec:mpmodel} --- Definition~\ref{def:mpqd} and Proposition~\ref{prop:mpqreduced}, the $Q$-reduced nerve as a projective resolution, closing Gap~G1; two new statements, no existing statement moves; see the changelog);

Version 1.22.0-pre.10 --- 2026-10-10 (hecke-extensions slot 22, L23-0 extended, \textbf{a pre-release, not owner-released, not refereed}: \S\ref{sec:oddnltwo} --- odd $n$ at $\ell=2$ in closed form, the $\Z/2^{k}$ splitting, $T_p=1+p$ at $\ell\in\{2,3\}$, $\rank T_\ell\le2$, the $c_p$ rule; eleven new statements, no existing statement moves; see the changelog);

Version 1.22.0-pre.9 --- 2026-10-10 (hecke-extensions slot 22, \textbf{a pre-release, not owner-released, not refereed}: \S\ref{sec:mpmodel} --- code-vs-chapter wording repairs in Remark~\ref{rmk:mprulesread}, \ref{MPC3} and Example~\ref{ex:mpteardrop}; no statement number or value moves; see the changelog);

Version 1.22.0-pre.8 --- 2026-10-10 (hecke-extensions slot 22, \textbf{a pre-release, not owner-released, not refereed}: \S\ref{sec:mpmodel} --- the pre.7 referee's
four minor repairs; Remark~\ref{rmk:mphecke} rewritten for the built MP-4 route (a twisted chain map on the MP resolution itself, unique over every ring) with its
\COMP\textsuperscript{r} record and \ref{MPC12}; the MP-3b coverage of \ref{MPC10}, \ref{MPC11}; no statement number or value moves; see the changelog);

Version 1.22.0-pre.7 --- 2026-10-10 (hecke-extensions slot 22, \textbf{a pre-release, not owner-released, not refereed at pre.7} (refereed at pre.8, ACCEPT): \S\ref{sec:mpmodel} --- $25$ further
\code{\textbackslash LEAN} markers (Lean MP-L3: morphisms, naturality, refinement, $\varepsilon^{\mathsf{MP}}$, the normalised complex), the unnormalised clauses they need in
Definition~\ref{def:mpmorphism}, Theorem~\ref{thm:mpnatural} and Proposition~\ref{prop:mpbredon}, and the \COMP\textsuperscript{r} records of Propositions~\ref{prop:mpboundary},
\ref{prop:mprescor} (MP-3); no statement number or value moves; see the changelog);

Version 1.22.0-pre.6 --- 2026-10-10 (hecke-extensions slot 22, \textbf{a pre-release, not owner-released}: \S\ref{sec:mpmodel} refereed by F-MP2 (contract part,
ACCEPT WITH REPAIRS), repairs R1--R8 applied --- Theorem~\ref{thm:mpnatural}(4)'s first sentence corrected; the strong form of goodness in Definition~\ref{def:mprules}(3)
is the one the code and Lean check; the first \code{\textbackslash LEAN} markers of \S\ref{sec:mpmodel} ($30$ on $11$ labels); no statement number, label, tag or value moves; see the changelog);

Version 1.22.0-pre.5 --- 2026-10-09 (hecke-extensions slot 22, \textbf{a pre-release, not owner-released, not refereed at pre.5}: new chapter \S\ref{sec:mpmodel},
the Moerdijk--Pronk model --- the choice-free MP set, its comparison theorem in every degree over every ring, all degrees with its scope, morphisms with the label factor,
the Bredon comparison; no shipped value or default changes, no existing number, label, tag or marker moves; see the changelog);

Version 1.22.0-pre.4 --- 2026-10-09 (hecke-extensions slot 22, \textbf{a pre-release, not owner-released}: Remark~\ref{rmk:negI}, second bullet, the superfluous hypothesis $6\in R^{\times}$ deleted --- erratum, wording only; see the changelog);

Version 1.22.0-pre.3 --- 2026-10-09 (hecke-extensions slot 22, \textbf{a pre-release, not owner-released}, refereed by F-LES with repairs R1--R6 applied: new last subsection
\S\ref{sec:nebles} of \S\ref{sec:nebentypus} --- the level-change cone for $B\le A$ with compatible characters, its six-term sequence, inflation--restriction in
the normal case, which level steps are normal, level-raising $U_p$ settled negatively; no existing number, label, tag or marker moves; see the changelog);

Version 1.22.0-pre.2 --- 2026-10-08 (hecke-extensions slot 22, \textbf{a pre-release, not owner-released}: Remark~\ref{rmk:levelNcurvefamily}(1) reworded, Mazur's good-prime condition sufficient, not necessary, at $p=2$; nothing else changes; see the changelog);

Version 1.22.0-pre.1 --- 2026-10-06 (hecke-extensions slot 22, \textbf{a pre-release, not owner-released}: new last subsection \S\ref{sec:nebwildsec} of
\S\ref{sec:nebentypus} --- wild characters over the chain rings $\mathcal O_{j,e}$, the groups $\Ga_{\mathsf H}(N)$, specialisation without division, \ref{NB7}--\ref{NB10} ---
as refereed by F-U5; an addendum after Theorem~\ref{thm:hidacontrol}; no existing number, label, tag or marker moves; see the changelog);

Version 1.21.0 --- 2026-10-06 (hecke-extensions slot 21, the close: \textbf{the release} from v1.21.0-pre.4 under the owner's decision D16
(accepted as recommended); header, date and changelog only --- no statement's mathematics, label, number, tag, status or marker count ($140$ on $65$ labels) changes; see the changelog);

Version 1.21.0-pre.4 --- 2026-10-06 (hecke-extensions slot 21, \textbf{a pre-release, not owner-released}: new
Proposition~\ref{prop:levelNcurve} --- the weight-two presentation at every prime, $p\in\{2,3\}$ included, is the cohomology of the open modular curve ---
with Remark~\ref{rmk:levelNcurvefamily}, in a new last subsection of \S\ref{sec:levelN}; no existing number, label, tag or marker moves; see the changelog);

Version 1.21.0-pre.3 --- 2026-10-06 (hecke-extensions slot 21, \textbf{a pre-release, not owner-released}: one addendum sentence in \ref{Hi5} --- [U]'s (U17)(a)
equals \Espec{cl:lveis}(2) on the ordinary Eisenstein block for $n<\ell-1$; no clause's mathematics, label, tag or marker count changes; see the changelog);

Version 1.21.0-pre.2 --- 2026-10-05 (hecke-extensions slot 21, \textbf{a pre-release, not owner-released}: addenda only --- [U]'s lattice-independence of
ordinary parts recorded in Remark~\ref{rmk:stabconvention}(3), the normalisation of Hida's $P_n$ in Theorem~\ref{thm:hidacontrol}(2), and literature locators read in the
slot-21 literature pass; no clause's mathematics, label, tag or marker count changes; see the changelog);

Version 1.21.0-pre.1 --- 2026-10-05 (hecke-extensions slot 21, step R20-L, \textbf{a pre-release, not owner-released}: the Shimura locators of the
nine Shimura \code{\textbackslash LEANlit} markers and of the Hecke-ring \textsf{[standard]} notes verified on the owner's scan and restated where the
book numbers differ; no clause's mathematics, label, tag or marker count changes; see the changelog);

Version 1.20.0 --- 2026-10-05 (hecke-extensions slot 20, the close: \textbf{the release} from v1.20.0-pre.2 under the owner's ruling O6,
on the slot-20 referees F1 / F1$'$ (ACCEPT) and the Lean audit F3 of milestone M4 (ACCEPT WITH REPAIRS, repaired); the milestone-M4 \code{\textbackslash LEAN} markers ($140$ on $65$ labels);
no statement's mathematics, label or count change; see the changelog);

Version 1.20.0-pre.2 --- 2026-10-04 (hecke-extensions slot 20, the close delta for the Fable referee F1$'$, \textbf{a pre-release, not owner-released}: code-map text only ---
WP-MEAS, NB-4, WP-LV3 and S1\,/\,S2 \emph{built} in \S\ref{sec:notimplemented}, $\inv^{\ast}$ over $\Z/\ell^{k}$ built in \S\ref{sec:codemap}; no statement, tag, label or count change; see the changelog);

Version 1.20.0-pre.1 --- 2026-10-04 (hecke-extensions slot 20, WP-2, \textbf{a pre-release, not owner-released}, for the slot-20 referee F1: code-map text only ---
the slot-20 packages WP-LV3, WP-MEAS and NB-4 planned in \S\ref{sec:codemap} and \S\ref{sec:notimplemented}, their mathematics in [E] 1.20.0-pre.1; the milestone-M4
\code{\textbackslash LEAN} markers follow the Lean audit F3 of slot~20; no statement, tag, label or count change; see the changelog);

Version 1.19.0 --- 2026-10-04 (hecke-extensions slot 19, the close: \textbf{the release} from v1.19.0-pre.6 under the owner's ruling O3,
on the slot-19 referees F1 / F1-delta's ACCEPT WITH REPAIRS; the milestone-M3 \code{\textbackslash LEAN} markers ($98$ on $48$ labels); no statement's
mathematics, label or count change; see the changelog);

Version 1.19.0-pre.6 --- 2026-10-04 (F1-delta m-C1: one sentence in err:e6, textual only; before that: hecke-extensions slot 19, WP-2 continued, \textbf{a pre-release, not owner-released}: Shimura's notation corrected
(the full operator is his $T(p^{\nu})$, his $T(1,p^{\nu})$ the primitive coset); Proposition~\ref{prop:model7h4}'s \ref{H4}(ii) clause records [GT]'s proof of acyclicity (slot-19 L19-3);
a v1.19.0-pre.5 box in \S\ref{sec:notimplemented} (WP-BIL, WP-LIN built); no tag, label or count change; see the changelog);

Version 1.19.0-pre.4 --- 2026-10-04 (hecke-extensions slot 19, WP-2 continued, \textbf{a pre-release, not owner-released}: code-map text only --- WP-LV2 (the
twisted Eisenstein route on \S\ref{sec:nebentypus}'s character model) and WP-HL (the truncated $\Lambda$-adic arm) built, in \S\ref{sec:codemap} and \S\ref{sec:notimplemented};
no statement, tag, label or count change; see the changelog);

Version 1.19.0-pre.3 --- 2026-10-04 (hecke-extensions slot 19, WP-2 continued, \textbf{a pre-release, not owner-released}: code-map text only --- the
slot-19 packages as built (WP-LC, WP-LV untwisted, WP-NF, NB-2) and as not built (WP-HL, WP-LV twisted) in \S\ref{sec:codemap} and \S\ref{sec:notimplemented};
no statement, tag, label or count change; see the changelog);

Version 1.19.0-pre.2 --- 2026-10-03 (hecke-extensions slot 19, WP-3 repairs, \textbf{a pre-release, not owner-released}: the slot-19 referee F1's
ACCEPT, its typographic note m-C1 applied in three \S\ref{sec:codemap} rows, and the WP-LV pointer follows [E] 1.19.0-pre.2's saturation hypothesis; no statement, tag,
label or count change; see the changelog);

Version 1.19.0-pre.1 --- 2026-10-03 (hecke-extensions slot 19, WP-2, \textbf{a pre-release, not owner-released}, for the slot-19 referee F1:
code-map text only --- the slot-19 packages WP-NF, WP-LC, WP-LV, WP-HL and NB-2 planned in \S\ref{sec:codemap} and \S\ref{sec:notimplemented}, their mathematics in
[E] 1.19.0-pre.1 and [U] 0.3.0-pre.1; no statement, tag, label or count change; see the changelog);

Version 1.18.0 --- 2026-10-03 (hecke-extensions slot 18, the close: \textbf{the release} from v1.18.0-pre.2 under the owner's ruling O3 of 2026-10-03,
on the slot-18 referee F1's ACCEPT WITH REPAIRS; the milestone-M2 \code{\textbackslash LEAN} markers ($81$ on $42$ labels); the \textbf{erratum} in the proof of
Theorem~\ref{thm:welldef}(c), confirmed by the Lean audit F3 and by the owner; no statement, tag, label or count change; see the changelog);

Version 1.18.0-pre.2 --- 2026-10-03 (hecke-extensions slot 18, WP-3 repairs, \textbf{a pre-release, not owner-released}: the slot-18 referee F1's
M-1, m-C1, m-C2 applied (\ref{IL5}'s proof route and hypothesis, Remark~\ref{rmk:levelNscale}(4)); the code maps of WP-5, WP-6, WP-7 as built and run;
\ref{IL3}, \ref{IL4} moved to \COMP\textsuperscript{a}; an \textbf{erratum} in the proof of Theorem~\ref{thm:welldef}(c), for the owner's confirmation; see the
changelog);

Version 1.18.0-pre.1 --- 2026-10-03 (hecke-extensions slot 18, WP-2, \textbf{a pre-release, not owner-released}, for the slot-18 referee F1:
contract-first code maps for slot-18 WP-5 (\S\ref{sec:levelN}, Remark~\ref{rmk:levelNscale}(4)), WP-6 (\S\ref{sec:nebentypus}) and WP-7
(\S\ref{sec:isomodelell}, the new checkable item \ref{IL5}); Lemma~\ref{lem:heckerankone}'s Eichler--Shimura input located in \cite{diamond_im}, verified;
no statement's content, tag, label or count changes; see the changelog);

Version 1.17.0 --- 2026-10-02 (hecke-extensions slot 17, the close: \textbf{the release} from v1.17.0-pre.5, under the owner's ruling O3 of
2026-10-01; the referee F1 (slot-17 WP-3) was ACCEPT WITH REPAIRS, $0$ blocking, its repairs applied at pre.3 and pre.4; over pre.5 the first
\code{\textbackslash LEAN} / \code{\textbackslash LEANlit} markers (Lean milestone M1, clause-scoped, $42$ markers on $20$ labels) --- no statement, tag,
label or count changes; see the changelog);

Version 1.17.0-pre.5 --- 2026-10-02 (hecke-extensions slot 17, WP-2, \textbf{a pre-release, not owner-released}: pointers only --- slot-17 R17-4's
verified locators (\ref{NB4}(c)'s input, \cite[Lemmas~1.1.1, 1.1.4]{ash_stevens}, now verified; Lemma~\ref{lem:heckerankone}'s eigenbasis from
\cite{diamond_im}; \ref{LN5} from \cite[Th\'eor\`eme~2]{merel96}), two narrowed attributions to \cite{ash_stevens}, and \S\ref{sec:isomodelell}'s status sentences
and code map after slot-17 WP-7; no statement, tag, label or count changes; see the changelog);

Version 1.17.0-pre.4 --- 2026-10-02 (hecke-extensions slot 17, WP-2, \textbf{a pre-release, not owner-released}: two owner rulings of 2026-10-02 on the
slot-17 referee F1 --- \ref{NB4}(c)'s sharp form at $n=0$ tagged \PROVED{} given one \STD{} input, its locator \cite[\S1.1]{ash_stevens} not verified; and,
following [E]'s restatement of (R) on congruence-free classes, the final wording of \ref{LN6}'s and \S\ref{sec:saturation}'s dependent sentences; no label or
count changes; see the changelog);

Version 1.17.0-pre.3 --- 2026-10-02 (hecke-extensions slot 17, WP-2 resumed, \textbf{a pre-release, not owner-released}: the slot-17
referee F1's repairs --- \ref{NB4}(c)'s orientation proved, its last \STD{} input located but not verified, so the sharp form at $n=0$ is \COMP\textsuperscript{r};
the code maps of \S\ref{sec:levelN} (slot-17 WP-5) and \S\ref{sec:nebentypus} (slot-17 WP-6) built; Remark~\ref{rmk:levelNscale}(4)'s $N=99330$ point; five minors;
no label, tag set or count changes; see the changelog);

Version 1.17.0-pre.2 --- 2026-10-02 (hecke-extensions slot 17, WP-2, second pass, \textbf{a pre-release, not owner-released}: one
citation repair in Lemma~\ref{lem:heckerankone}'s proof --- the eigenbasis from \cite[Thm.~3.41]{shimura} or \cite[\S5.5]{diamond_shurman}, both not verified,
multiplicity one from \cite[Ch.~VII, \S5]{serre_cours}; no statement, tag or label changes; see the changelog);

v1.17.0-pre.1 --- 2026-10-01 (hecke-extensions slot 17, WP-2, drafted during the slot-16 close, \textbf{a pre-release, not
owner-released}: the deferred items of slot~16 for one referee (slot-17 WP-3, Fable) --- Remark~\ref{rmk:sattor} (new), \ref{NB4}'s sharp form at $n=0$
as data with a proposed proof, the composite-$N$ caveat of Remark~\ref{rmk:levelNscale}(4), the code maps of slot~17's packages (contract first), and the
\code{\textbackslash LEAN} tag defined, not yet used; one numbered statement added, none removed or renamed; see the changelog);

v1.16.0 --- 2026-10-01 (hecke-extensions slot 16, the close: \textbf{the release} from v1.16.0-pre.3, under the owner's ruling O2 of
2026-10-01; the referee (slot-16 WP-14) was ACCEPT WITH REPAIRS, $0$ blocking, every repair applied at pre.3 and confirmed; over pre.3 verified locators, the
reason the $e_\sigma$ average is a no-op on $d$, and code-map and status sentences only --- no statement added, no tag moved; see the changelog);

v1.16.0-pre.3 --- 2026-10-01 (hecke-extensions slot 16, WP-12, \textbf{a pre-release, not owner-released}: the repairs answering the
WP-14 Fable referee (ACCEPT WITH REPAIRS, 0 blocking, 1 major C-1, 8 minors C-2--C-9), all applied; no label added, renamed or removed, no count moves;
see the changelog);

v1.16.0-pre.2 --- 2026-10-01 (hecke-extensions slot 16, WP-12 follow-up, \textbf{a pre-release, not owner-released}: \ref{NB5}'s
product form squared, with its proof and the reduction rule (slot-16 WP-9 audit); no label added, renamed or removed; see the changelog);

v1.16.0-pre.1 --- 2026-10-01 (hecke-extensions slot 16, WP-12, \textbf{a pre-release, not owner-released}: new
\S\ref{sec:levelNdegsec} --- the model-to-presentation map, the degeneracy maps $\alpha_d$ on the cokernel presentation (replacing v1.15.0-pre.3's
unrefereed sentence), higher weight at composite $N$, \ref{LN7}--\ref{LN9} --- and the new chapter \S\ref{sec:isomodelell}, the $\ell$-local
isotropy model and the closed form at $\ell=3$; ten numbered statements added, none removed or renamed, no tag moved; referee slot-16 WP-14; see
the changelog);

v1.15.0 --- 2026-10-01 (hecke-extensions slot 15, WP-18: \textbf{the release at the slot-15 close} from v1.15.0-pre.3, under the
owner's ruling Q2 of 2026-09-30; the referee (slot-15 WP-4) was ACCEPT WITH REPAIRS, $0$ blocking, all applied at pre.2; \S\ref{sec:nebentypus}
released unbuilt; pre.3's one scope sentence released unrefereed; no statement added, no tag moved; see the changelog);

v1.15.0-pre.3 --- 2026-09-30 (hecke-extensions slot 15, WP-17, \textbf{a pre-release, not owner-released}: one scope
sentence in \S\ref{sec:levelN} on the degeneracy maps and the cokernel presentation; the sparse elimination built; no statement added, no
tag moved; see the changelog);

v1.15.0-pre.2 --- 2026-09-30 (hecke-extensions slot 15, WP-3 resumed, \textbf{a pre-release, not owner-released}: the
repairs answering the WP-4 Fable referee, 0 blocking, all applied; $H_{N,p}$ and \file{HeckeHidaComponents} built; no statement added, no tag
moved; see the changelog);

v1.15.0-pre.1 --- 2026-09-30 (hecke-extensions slot 15, package WP-3, \textbf{a pre-release, not owner-released}:
the new chapters \S\ref{sec:nebentypus} (nebentypus for every $N$ over $\F_q$, contract first, not built) and \S\ref{sec:saturation}
(the saturation statements behind the $\Z/\ell^{K}$ Hecke-algebra route, and the torsion-aware comparison with modular forms); no existing
statement changed, no tag moved; see the changelog);

v1.14.2 --- 2026-09-29 (hecke-extensions slot 14, owner ruling Q-F: \S\ref{sec:levelN}'s code maps now say \emph{built}
(slot-14 WP-12, \texttt{7a5158d}); no statement changed, no tag moved; see the changelog);

v1.14.1 --- 2026-09-29 (hecke-extensions slot 14, owner rulings of 2026-09-29 on the contract-currency audit:
the second cluster run's \COMP\textsuperscript{a} sentences in \ref{Hi11}, \ref{Hi12} and \S\ref{sec:hidafull}'s code map ([S14] Q-B), and
Lemma~\ref{lem:rhorank}, the $U_\ell$ lever's rank claim (Q-H); no existing statement changed, no tag moved; see the changelog);

v1.14.0 --- 2026-09-28 (hecke-extensions slot 14, package WP-13: \textbf{the owner-signed
release} of \S\ref{sec:levelN} from v1.14.0-pre.2, released unbuilt. The referee was ACCEPT WITH REPAIRS with $0$ blocking.
The data sentences of the first cluster run are restored in \ref{Hi11}, \ref{Hi12} and Remark~\ref{rmk:bocksteinexp}, citing the
extract \file{docs/matrix-engine/results/noctua2\_run1/}, published on the owner's sign-off ([S14] Q2, Q3). Release note:
\file{docs/implementation/slot14/WP-13\_release\_note.md});

v1.14.0-pre.2: 2026-09-28 (hecke-extensions slot 14, package WP-5, \textbf{a pre-release, not
owner-released}: the repairs answering the Fable referee of \S\ref{sec:levelN}, ACCEPT WITH REPAIRS, $0$ blocking, all eleven
findings applied --- a Hecke matrix on $H^{1}$ costs quadratically in $m_N$ through the normal form of $\mathrm{Rel}_N$
(Remark~\ref{rmk:heilbronncost}), what the Eisenstein scan computed is stated (\ref{LN4}, \ref{LN6}); the data remarks of the first
cluster run taken back out of \ref{Hi11}, \ref{Hi12} and Remark~\ref{rmk:bocksteinexp} ([S14] Q2's default), the provenance line
of \S\ref{sec:hidafull}'s code map kept. Referee: \file{docs/implementation/slot14/WP-5\_referee.md});

v1.14.0-pre.1: 2026-09-28 (hecke-extensions slot 14, package WP-4, \textbf{a pre-release, not
owner-released}: a new chapter \S\ref{sec:levelN}, \textbf{the contract first} --- the binding coset table for every $N$
(obstacle O3), the cokernel presentation pinned with Manin's sequence and a duality natural in restriction, corestriction and
conjugation, and the Hecke action of $T_p$ and $U_p$ on it by Merel's matrices, transcribed and \textbf{proved} (obstacle O2), its
cost and the predictions (LN1)--(LN6), among them mod-$p$ Loewy length $e+1$ at Mazur's Eisenstein ideal; data remarks of the
first cluster run added to \ref{Hi11}, \ref{Hi12} and Remark~\ref{rmk:bocksteinexp}, their publication the owner's (removed at pre.2). Report:
\file{docs/implementation/slot14/WP-4\_report.md});

v1.13.0: 2026-09-27 (hecke-extensions slot 13, the close, package WP-15: \textbf{the owner-signed
release} of \S\ref{sec:hidafull} from v1.13.0-pre.2 --- the referee ACCEPT with $0$ blocking, slot-13 WP-6's arms held;
one sentence of Definition~\ref{def:hidalattice} amended (the module route's Nakayama basis on Hida's lattice), the code
maps read ``built''. Report: \file{docs/implementation/slot13/WP-15\_report.md});

v1.13.0-pre.2: 2026-09-26 (hecke-extensions slot 13, package WP-3: the repairs answering the Fable
referee of \S\ref{sec:hidafull}, ACCEPT WITH REPAIRS, all twelve findings applied --- Hida's lattice
$L_n$ is $\rhon$ conjugated by the binomial diagonal, not $\Vcoef$ (Remark~\ref{rmk:stabconvention}(3), the variant (H) of
Definition~\ref{def:hidalattice}, \ref{Hi8}'s arms), and the cost model's regimes and obstacles, with the Manin--Merel
Hecke action on the cokernel named (Remark~\ref{rmk:levelscaling}). Response:
\file{docs/implementation/slot13/WP-3\_response.md});

v1.13.0-pre.1: 2026-09-26 (hecke-extensions slot 13, package WP-2, the pre-release of an
$X$ bump, \textbf{the contract first}: a new chapter \S\ref{sec:hidafull} --- level structure by Shapiro on
the fixed model (the level-$N$ correspondences of Proposition~\ref{prop:levelcorrespondence}, the level-$\ell$
operator of Proposition~\ref{prop:levelellU}, the old space, the orbit-wise cokernel and a cost model), the
integral stabilisation (Proposition~\ref{prop:stabilisation}), Hida's control and weight independence as
literature with verified locators, the lattice datum and the predictions (Hi7)--(Hi12), committed before
their code. Report: \file{docs/implementation/slot13/WP-2\_report.md});

v1.12.1: 2026-09-26 (hecke-extensions slot 12, post-close, package S12-census2: a
\textbf{non-normative} bump, code maps only --- \S\ref{sec:hida}'s code map records what is built,
\file{Hecke/HeckeOrdinaryPart} and \file{HeckeExt/HeckeHidaTower}, and \S\ref{sec:codemap} names
\file{HeckeExt/HeckeFinalityRoutes}; companion pointer [E] v1.13.1. Report:
\file{docs/implementation/slot12/census2\_report.md});

v1.12.0: 2026-09-26 (hecke-extensions slot 12, post-close, \textbf{the owner-signed
release}, from v1.12.0-pre.3: publishes \S\ref{sec:hida} (the foundation of Hida theory at level~$1$) and
\textbf{the Shapiro transcription in degree~$1$, proved} --- Proposition~\ref{prop:bdrytranscription} and
Lemma~\ref{lem:bdrytorus} --- both accepted by their Fable referees after repairs, all applied, none
outstanding. \textbf{The one status change this version makes}: Question~\ref{open:ordprecision} is
\textbf{answered} at $\ell\ge5$, $n\ge1$ by Corollary~\ref{cor:ordprecision} (with
Proposition~\ref{prop:torsionprepass} and Remark~\ref{rmk:ordtorsionfree}): the tag moves \OPENSTAT{} $\to$
\PROVED{}, and the item is retitled ``--- ANSWERED 2026-09-26'', exactly as \ref{open:ell2} lost its
\code{[open]} tag at v1.3.0. No statement added, no label renamed or removed, one tag moves; the question
text and the boxed statement of \ref{open:ordprecision} are unchanged (rules~\ref{L1}, \ref{L2}). Report:
this release note (owner sign-off 2026-09-26; \file{docs/implementation/slot12/WP-2\_report.md},
\file{WP-3\_referee.md}, \file{WP-16\_report.md} and \file{WP-16\_referee.md} record the refereed
content));

v1.12.0-pre.3: 2026-09-26 (hecke-extensions slot 12, package WP-16, post-close, the third
pre-release: \textbf{the Shapiro transcription in degree~$1$, proved} --- Proposition~\ref{prop:bdrytranscription}
(the coordinate $\Phi$ an isomorphism over every ring at $\rhon(-I)=\id$; an $H^{1}$ formula at every $\nu$ and
every $p$; \eqref{eq:bdrygroupform} at $\nu=1$, with no hypothesis $p\in\cring^{\times}$) and
Lemma~\ref{lem:bdrytorus} (boundary operators away from $\ell$ factor through the torus); status notes on
Proposition~\ref{prop:bdrychar}, Obligation~\ref{obl:taubdry}(v.b) and Claim~\ref{cl:chi}(2). Two statements
added, no label renamed or removed, no tag moved. Report: \file{docs/implementation/slot12/WP-16\_report.md});

v1.12.0-pre.2: 2026-09-26 (hecke-extensions slot 12, package WP-3, the second pre-release:
\textbf{the Fable referee's repairs} (ten minor findings, all applied; none blocking, none major) --- the scope
clause of Proposition~\ref{prop:torsionprepass}, the definition of the $\ell$-adic projector and $e=e'$ in
Corollary~\ref{cor:ordprecision}, the free-block sentence of Remark~\ref{rmk:ordexponentuse} with a \COMP{}
record, and (Hi2), (Hi3), (Hi5), (Hi6) made explicit. No label renamed or removed, no statement added. Report:
\file{docs/implementation/slot12/WP-2\_report.md});

v1.12.0-pre.1: 2026-09-26 (hecke-extensions slot 12, package WP-2, the pre-release of an
$X$ bump, \textbf{the contract first}: a new chapter \S\ref{sec:hida}, the foundation of Hida theory at
level~$1$ --- the torsion pre-pass (Proposition~\ref{prop:torsionprepass}: the $\ell$-primary torsion of
$H^{1}$ from $\rhon(S)$, $\rhon(T)$ alone), the certified precision before the run
(Corollary~\ref{cor:ordprecision}, which answers Question~\ref{open:ordprecision} at $\ell\ge5$, $n\ge1$;
the item's status macro waits for the release), the ordinary projector as one power of $U$
(Lemma~\ref{lem:ordexponent}), the weight towers and the predictions (Hi1)--(Hi6), committed before
their code. Six statements added, no label renamed or removed, no existing tag moved. Report:
\file{docs/implementation/slot12/WP-2\_report.md});

v1.11.1: 2026-09-25 (hecke-extensions slot 11, package WP-17, the slot close, a $Y$ bump:
\textbf{the WP-15 referee's repairs} --- Remark~\ref{rmk:thetacheckable}'s provenance sentence restated to the
reports, a new Remark~\ref{rmk:ordtorsionfree} (the ordinary part is torsion-free at level~$1$, at $\ell\ge5$,
$n\ge1$), Question~\ref{open:ordprecision} narrowed to the per-cell bound (still open), and wording. No
normative statement moves, no label renamed or removed; one \PROVED{} remark added. Report:
\file{docs/implementation/slot11/WP-17\_report.md});

v1.11.0: 2026-09-25 (hecke-extensions slot 11, package WP-15: \textbf{the slot's $X$ bump
published} from v1.11.0-pre.2 with what slot~11 built --- Part~VII's objects and every checkable item of
Remark~\ref{rmk:thetacheckable} (all held; (C9)'s indexing said), the lifts Proposition~\ref{prop:relationsmodel}
needs and its recursion on the free model, the $\ell$-adic digit route and the ordinary idempotent on it;
Obligation~\ref{obl:backend}(iv)'s scope sentence worded to the digit route as built; Question~\ref{open:ordprecision}
still open. No statement added, no label renamed or removed; one \COMP{} tag added. Report:
\file{docs/implementation/slot11/WP-15\_report.md});

v1.11.0-pre.2: 2026-09-25 (hecke-extensions slot 11, package WP-2, the second pre-release:
\textbf{the Fable referee's repairs} (one major, nine minor; Proposition~\ref{prop:heckeH0} confirmed
independently) --- Obligation~\ref{obl:backend}(iv)'s owner-ruled scope sentence for the fixed-width
digit-plane route, Proposition~\ref{prop:heckeH0} at $\ell=2$, a multiplicity in (C5), the model's answer to
(C7) recorded, and wording. No label renamed or removed, no tag moved. Report:
\file{docs/implementation/slot11/WP-2\_report.md});

v1.11.0-pre.1: 2026-09-25 (hecke-extensions slot 11, package WP-2, the pre-release of
an $X$ bump, \textbf{the contract first}: a new Part~VII states multiplication by the Dickson invariant
$\DkL$ and the operator $\thL$ it induces (\S\ref{sec:dicksontheta}) --- the determinant twist on all of
$\Mat_2(\Fl)$, $T_p\thL=p\,\thL T_p$ and $U_\ell\thL=0$ at cochain level, the quotient $\Qn$ with a fixed
basis, the long exact sequence with $\ker\thL=\im\delta$ and its closed forms, $H^{0}=\Fl[\DkL,\DkD]$
with $T_p$ and $U_\ell$ on it (Proposition~\ref{prop:heckeH0}: $U_\ell$ has rank $\le1$ and square zero on
$H^{0}$, and is \emph{not} zero at $n=\ell^{3}-\ell$), the literature on $\theta$ and what of it is
checkable on the model --- and the ordinary projector (\S\ref{sec:ordinary}), with the precision
question open; and \S\ref{sec:relationsmodel} states \ref{thm:relations}'s recursion on the free model
for every ring. Twenty statements added, no label renamed or removed, no existing tag moved. It carries
its own identifier so that no version number is published twice under different text (\ref{V1},
\ref{V2}); the referee's repairs will be v1.11.0-pre.2. Report:
\file{docs/implementation/slot11/WP-2\_report.md});

v1.10.1: 2026-09-25 (hecke-extensions slot 10, package WP-15, phase~2, a $Y$ bump:
\textbf{the phase-1 referee's repairs --- no normative predicate, no hypothesis, no displayed
formula, no tag and no label moves}. The Fable referee of v1.10.0 accepted it after two major
findings, each a one-line repair, and four minor ones on this document: Remark~\ref{rmk:bredonhecke}
dates the canonical $\tau$ to v1.1.0 and cites the one comparison this document has at
$\ell\in\{2,3\}$, the legacy against the canonical $\tau$, as an observation; its ``compares nothing
there'' names the operators meant; Remark~\ref{rmk:strictliftbuilt}(3) counts $359$ cells on eight
rings; Remark~\ref{rmk:integralevidence} attributes its $224$-element check correctly. Report:
\file{docs/implementation/slot10/WP-15\_report.md});

v1.10.0: 2026-09-25 (hecke-extensions slot 10, package WP-15, an $X$ bump on
the owner's rulings of 2026-09-25, all at recommendation~(A): \textbf{the strict lift of
Proposition~\ref{prop:strictlift} is built} for $g=\mathrm{diag}(1,p)$, $p\in\{2,3,5\}$, in degrees
$\le q_{\max}+1$, certified exactly over $\Z$, and recorded with its measured integer sizes in a new
Remark~\ref{rmk:strictliftbuilt}; a new Remark~\ref{rmk:bredonhecke} beside
Proposition~\ref{prop:strictlift}(4) states that the operator on cohomology of the Bredon object at
$\ell\in\{2,3\}$ is \textbf{not shown} independent of $\tau$ and that this document compares nothing
there; Lemma~\ref{lem:bulkperiodic}'s proof fixes the sign of its Bockstein, so its class $e$ is
pinned; Obligation~\ref{obl:coneexact}(vii) over $\Z/\ell^{e}$, $\ell\ge5$, is \textbf{discharged in
its first two sentences at $\ell^{e}<2^{32}$}, every E8 point, its third (the Poincar\'e-duality
gate) recorded unavailable over the ring; \S\ref{sec:integral} gains the lifted datum, the class $e$
with its periodicity map, and the integral Hecke action read basis-free. Two statements added, no
label renamed or removed, no existing tag moved. Report:
\file{docs/implementation/slot10/WP-15\_report.md});

v1.9.1: 2026-09-25 (hecke-extensions slot 9, package WP-17, the slot close,
a $Y$ bump: \textbf{the referee's repairs and status --- no normative predicate, no hypothesis,
no displayed formula, no tag and no label moves}. The Fable referee of v1.9.0 accepted it after
repairs (two major, both one-sentence, and four minor findings on this document): the slot's
independent review is now in the tree and cited
(\file{docs/implementation/slot9/WP-14\_review.md}); Obligation~\ref{obl:coneexact}(iv)'s as-built
sentence follows the slot's trimming of the item-4b checks; two commit hashes are back-filled; one
citation, one sentence's grammar and one code-map lead are repaired; the \code{Hecke*} count is
re-derived. Report: \file{docs/implementation/slot9/WP-17\_report.md});

v1.9.0: 2026-09-24 (hecke-extensions slot 9, package WP-15, an $X$ bump on
the owner's rulings of 2026-09-24, all at recommendation~(A): \textbf{Remark~\ref{rmk:bocksteinexp}(2)
is stated at every prime}, with its instances at $\ell\in\{2,3\}$; a new
Remark~\ref{rmk:certifiedeperdiff} reads Proposition~\ref{prop:certifiede}(3) \textbf{differential
by differential}; Definition~\ref{def:ringsemantics}(2)'s stale refusal clause is replaced by what
this build computes at $\ell\in\{2,3\}$ --- the Bredon object's lengths and invariant factors ---
and what it refuses; Obligation~\ref{obl:coneexact}(iv)'s non-field half is \textbf{checked} over
$\Z/\ell^{e}$, $\ell\ge5$ (phase E8 item~4b), while (vi) is partial there and (vii) holds by
route~1 only, its composite half open; \S\ref{sec:integral} is partly built (its free model,
comparison map and integral groups for $n\le20$). No label renamed or removed, no existing tag
moved; one statement added. Report: \file{docs/implementation/slot9/WP-15\_report.md});

v1.8.1: 2026-09-24 (hecke-extensions slot 8, package WP-14, the slot close,
a $Y$ bump: \textbf{the referee's wording repairs and status --- no normative predicate, no
hypothesis, no displayed formula, no tag and no label moves}. The five minor findings of the
referee of v1.8.0 are applied as wording (Lemma~\ref{lem:strictlifting}'s proof, the notation
row of $F^{\sigma}_{\bullet}$, Theorem~\ref{thm:specialisation}'s v1.8.0 sentence,
Proposition~\ref{prop:bredon}'s family, Remark~\ref{rmk:truncdevice}'s citation); the
abstract's stale sentence is repaired; Appendix~\ref{app:changelog} gains the entries of
v1.7.2--v1.7.4 and its resolvable commit hashes. Phase E8 items 2, 3 and 4a are implemented
--- the complex over the ring in force at $\Z/\ell^{k}$, $\ell\ge5$, its reductions and
Bockstein --- and item 4b (the boundary and the cone over $\Z/\ell^{k}$) is the next slot's;
\S\ref{sec:codemap} and \S\ref{sec:notimplemented} record it, and the registered-validation
counts are re-derived ($114$; $65$ \code{Hecke*}). Close report:
\file{docs/implementation/slot8/WP-14\_report.md});

v1.8.0: 2026-09-24 (hecke-extensions slot 8, package WP-9, an $X$ bump:
\textbf{the integral line}, a new Part~VI, \S\ref{sec:integral} --- what the shipped complex
is over the ring in force (Bredon cohomology, Proposition~\ref{prop:bredon}), the free model
on the \emph{unrefined} tessellation with one integral comparison map
(Definition~\ref{def:isomodel}, Proposition~\ref{prop:isomodel}), the extension plan's
contract item C-A5 re-scoped onto it and owner direction OD-T1 answered ``not needed''
(Remark~\ref{rmk:cafive}, owner ruling 6c of 2026-09-24), the descent over $\Z$
(Propositions~\ref{prop:fracture}--\ref{prop:homvanish}), the strict lift of the Hecke
datum (existence only, Proposition~\ref{prop:strictlift}), the bad primes, the
certified-exponent rule and the characteristic polynomial over $\Z$ by residues, and what
each coefficient ring names (Definition~\ref{def:ringsemantics}); and the F-A remark
beside Claim~\ref{cl:bockstein} (Remark~\ref{rmk:bocksteinexp}, owner ruling 10).
Twenty statements added, \ref{H3} and Theorem~\ref{thm:specialisation}'s closing paragraph
amended, no label renamed or removed, no existing tag moved; the arbitrary-precision ban is
unchanged. Report: \file{docs/implementation/slot8/WP-9\_report.md});

v1.7.4: 2026-09-23 (hecke-extensions slot 7, package WP-14, the slot close:
\textbf{status only --- no normative predicate, no hypothesis, no displayed formula, no tag and
no label moves}. Phase E7 is \textbf{accepted}: Claim~\ref{cl:pd}(1)--(2) are \emph{tested} on
all $54$ cells of the extension stage's case list at $\ell\ge5$ and fail on none, and the claim
stays open (slot-7 WP-1); the extension record of \S\ref{sec:recordschema} item~6 is
\code{hecke\_ext/v3} (slot-7 WP-3, stage~3), its new keys being [E]'s field inventory;
\S\ref{sec:codemap} and \S\ref{sec:notimplemented} record it, and the registered-validation
counts are re-derived ($108$; $61$ \code{Hecke*}). Close report:
\file{docs/implementation/slot7/WP-14\_report.md});

v1.7.3: 2026-09-22 (hecke-extensions slot 6, package WP-14, the slot close:
\textbf{status only --- no normative predicate, no hypothesis, no displayed formula, no tag and
no label moves}. \S\ref{sec:contragredient} is \textbf{implemented} (slot-6 WP-1), and
Obligation~\ref{obl:coneexact}(vii) is \textbf{implemented in all three sentences} --- the
strict-pair route (WP-4), route~3 (WP-5) and the Poincar\'e-duality gate for $X_3$ at
$\ell\ge5$ (WP-3), which Claim~\ref{cl:pd} still only \emph{tests}; \S\ref{sec:codemap} and
\S\ref{sec:notimplemented} record it, and the registered-validation counts are re-derived by
\code{--list} ($106$; $60$ \code{Hecke*}). Close report:
\file{docs/implementation/slot6/WP-14\_report.md});

v1.7.2: 2026-09-22 (hecke-extensions slot 5, package WP-13, the slot close:
\textbf{status only --- no normative predicate, no hypothesis, no displayed formula, no tag and
no label moves}. \S\ref{sec:codemap} records what slot~5 shipped, the registered-validation
counts are re-derived by \code{--list} ($100$; $55$ \code{Hecke*}), and the one point v1.7.1
left open for the owner --- whether Obligation~\ref{obl:taubdry}(vii) is a distinct weaker
clause or a pointer to Obligation~\ref{obl:coneexact}(iii) --- is \textbf{decided: Reading~A,
a distinct weaker clause}, so v1.7.1's text stands as drafted. Close report:
\file{docs/implementation/slot5/WP-13\_report.md});

v1.7.1: 2026-09-22 (hecke-extensions slot 5, package WP-12, the independent
referee pass on v1.7.0, a $Y$ bump: \textbf{pointers and self-contradictions only --- no
normative predicate, no hypothesis, no displayed formula, no tag and no label moves}. Five
of the six repairs are of text v1.7.0 itself introduced: \S\ref{sec:recordschema} item~6(f)
bound ``(a)--(e) and (g)--(h)'' where the item's preamble binds ``the nine below'', leaving
rule~(i) bound and unbound at once; Remark~\ref{rmk:conerestrictsmallell}'s opening said
``load-bearing for no clause'' two lines before saying it still discharges
Obligation~\ref{obl:taubdry}(vii); that clause's new asymmetry paragraph rested on ``the
data \emph{this} chapter stores are $\taub$ and $\Thb_g$ alone'', which the clause's own
statement falsifies, the real missing datum being $\etapair$ and $\hder$; its opening
sentence still called it a clause that ``belongs to the pair and is stated in full at''
Obligation~\ref{obl:coneexact}, which v1.7.0 falsified by strengthening (iii) alone;
and Obligation~\ref{obl:coneexact}(iii)'s proof now names $z\in\Zq{q}$, the ``one more
word'' its author's own residual finding asked for. The sixth records that slot-5 WP-11a
landed \S\ref{sec:recordschema} item~5's refusal. \textbf{Open for the owner}: whether
Obligation~\ref{obl:taubdry}(vii) is a distinct weaker clause or a pointer to
Obligation~\ref{obl:coneexact}(iii) --- this version takes the first reading, which is the
one v1.7.0's own ``the pair's form of'' commits to, and flags the choice. Report:
\file{docs/implementation/referee\_report\_2026-09-22\_contract\_v1\_7\_0.md});

v1.7.0: 2026-09-22 (hecke-extensions slot 5, package WP-11, the slot's one
$X$ revision. \textbf{Obligation~\ref{obl:coneexact}(iii)'s $\ell\in\{2,3\}$ carve-out on the
$\rres$-half is dropped and the identity asserted at \emph{every} $\ell$, with the proof
written in}: clause (i) carries no hypothesis on $\ell$ and aborts on failure, so on a stored
cocycle $z$ the difference $\Thb_g\rres z-\rres\Th_g z$ is $d(\hder z)$, a coboundary, and
$\rres^{q}\Amat{q}=\Abdry{q}\rres^{q}$ on cohomology follows --- the implication that was
re-derived in three successive packages and was never contract text (owner ruling
\code{[OD]} slot-4 close-out 3 of 2026-09-22; the hypothesis it uses re-checked against the
shipped code by slot-5 WP-4). The clause states, in as many words, \emph{what it does not
say}: at $\ell\in\{2,3\}$ it is about the complexes the implementation computes and not
about $H^{\bullet}(\Ga;\Mn)$, which Remark~\ref{rmk:conetorsion} keeps separate. And it
states the \textbf{asymmetry}: Obligation~\ref{obl:taubdry}(vii), in the boundary chapter,
has neither $\etapair$ nor $\hder$ and therefore no proof, so it keeps its carve-out and
gains a pointer; Remark~\ref{rmk:conerestrictsmallell} is demoted to the measurement that
preceded the proof, load-bearing for no clause. Obligation~\ref{obl:coneexact}(v)'s closing
sentence names \emph{both} the stored object's certificate and the validation arm, the
shipped \code{ConeCertificate} now carrying the unconditional first identity of
\eqref{eq:conedims} under its own guard. \S\ref{sec:recordschema}: item~5 makes the
derived-complex \emph{names} the record's own derived flag, so that
\code{q\_max\_source = "set"} on a derived complex is refusable by a verifier and not only
by the builder; item~6 becomes \code{hecke\_ext/v2} with nine binding rules, (b), (c), (g),
(h) and (i) new or amended, and (a) and (d) answer the two halves of v1.6.0's deferred
item~3. Obligation~\ref{obl:iota}(ii) gains the rider for a complex without exact group
data. \emph{Deferred again}: the $\inv$-equivariance of $\hder$, which no accepted package
measured. No label added, renamed or removed; no tag moves. Status sentences stay with the
slot close except where this version's own text falsifies them. Report
\file{docs/implementation/slot5/WP-11\_report.md});

v1.6.2: 2026-09-22 (hecke-extensions slot 4, package WP-11, the slot close,
a $Y$ bump: \textbf{status only --- no normative statement, no hypothesis, no displayed
formula, no asserted predicate, no tag and no label moves}. What is recorded: the
\emph{production} involution exists --- \S\ref{sec:involution}'s 2026-09-16 census, its
specification paragraph, \S\ref{sec:codemap}'s row and \S\ref{sec:notimplemented} are
rewritten from the pointers v1.6.1 left, naming \file{Hecke/HeckeInvolution},
\file{HeckeExt/HeckeSignSplit} and \code{HeckeInvolutionValidation} (slot-4 WP-3, phase
E3); \S\ref{sec:recordschema} item~6's status becomes ``written, bulk complex'' with
\file{HeckeExt/HeckeExtDriver} named (WP-1, WP-2, phase E2); \S\ref{sec:codemap} gains
\file{HeckeExt/HeckeExtRecord}, \file{HeckeExtReport}, \file{HeckeExtFinal},
\file{HeckePredictions} and \file{GloverFactors}, the two new settings keys
\code{commutant\_budget} and \code{final\_mode}, and --- in the
Chapter~\ref{sec:cone} row --- WP-4's three \emph{value-preserving} certificate
re-implementations, with the statement that no clause was sampled, relaxed or moved out
of \code{compute()}; Remark~\ref{rmk:labelcensus} records the content-keyed allow-list
(WP-6) and the slot-4 census run; and the validation counts go $89\to94$ and
$45\to49$ \code{Hecke*}, both counted by \code{--list} and not asserted. Three referee
questions of v1.6.1 (\code{ConeCertificate} C11 versus the strengthened clause (v); the
$\ell\in\{2,3\}$ carve-out in Obligation~\ref{obl:coneexact}(iii); the scope of the
v1.6.1 pointer repair) are \emph{not} decided here --- they are the owner's, and the
first two are slot-5 code. Report
\file{docs/implementation/slot4/WP-11\_report.md});

v1.6.1: 2026-09-22 (hecke-extensions slot 4, package WP-9, the independent
referee pass on v1.6.0, a $Z$ bump: \textbf{no hypothesis, displayed formula, asserted
predicate, tag or label moves}. Two recorded measurements are corrected ---
Remark~\ref{rmk:iotapairing}'s ``thirteen cells'' counts $(10,691)$ twice and folds in four
$n=0$ cells at which Obligation~\ref{obl:iota}(v) makes no prediction (it is \emph{eight}),
and \S\ref{sec:coneaction}'s closing paragraph said the commutator $[\inv_{c},\Acone{q}]$
was recorded on the $15$-cell grid where it is recorded at $(10,691)$ alone. Three status
sentences that v1.6.0's own new text falsifies --- \S\ref{sec:involution}'s census
(``no \code{HeckeInvolution} translation unit''), its specification paragraph and
\S\ref{sec:codemap}'s row, plus \S\ref{sec:notimplemented}, all saying that no production
$\inv^{\ast}$ exists --- carry a pointer, minimally, the rewrite itself remaining the slot
close's; and the title page's companion pointer, three revisions stale, moves to [E]
v1.5.1. Companion repairs at [E] v1.5.1: two false statements
([E]'s \code{cl:harder}, its new predicate paragraph at $\ell\in\{2,3\}$, and the Glover
recursion's factor bound) and two proof defects. Report
\file{docs/implementation/referee\_report\_2026-09-22\_contract\_v1\_6\_0.md});

v1.6.0: 2026-09-21 (hecke-extensions slot 4, package WP-8, an $X$ bump:
the eight items Appendix~\ref{app:changelog} deferred at v1.5.2 are resolved ---
Obligation~\ref{obl:coneexact}(i) ``as matrices'' named as the stored $\Ceq$ basis with
the flat identities \emph{recorded}, (iii) given the $\ell\in\{2,3\}$ carve-out on the
$\rres$-half, (iv)'s explicit $\im=\ker$ check read degree-generically, (v) and
\eqref{eq:conedims}'s first identity freed of the hypothesis $\Heq{2}=0$ (with
Lemma~\ref{lem:conerelative}(5) and its proof restructured), (vi) told where the
\emph{full} $T_{p^{\nu}}$ comes from, the preamble's ``six maps'' replaced by the three
map families, Proposition~\ref{prop:coneambiguity}(2) read inside $\Ceq$ and
Lemma~\ref{lem:etaexists}(V)'s ``single edge path'' repaired --- together with the
normative findings of the slot-4 packages accepted when this revision started:
\S\ref{sec:recordschema} gains item~6, the extension record \code{hecke\_ext/v1} and its
six binding rules --- two of them from the E2b driver, on what a \emph{resumed} run's
settings block contains and on what the incomplete-item count counts --- and
Remark~\ref{rmk:iotapairing}'s mirror-split prediction becomes a measured record while
\S\ref{sec:coneaction} gains a paragraph saying that $[\inv_{c},\Acone{q}]=0$ follows from
nothing here. Owner confirmations of 2026-09-21:
(a) yes, (b) yes, (c) per mode. No label added, renamed or removed; one tag moves. Status
sentences are the slot close's. Report
\file{docs/implementation/slot4/WP-8\_report.md});

v1.5.2: 2026-09-17 (hecke-extensions slot 3, package WP-11, the slot close,
a $Y$ bump: status only --- strict locality mode implemented (the one false status
sentence of v1.5.1, at \S\ref{sec:codemap}); Chapter~\ref{sec:cone} implemented except
Obligation~\ref{obl:coneexact}(vii), Obligation~\ref{obl:eta} discharged in full; the
$U_\ell$ repair, the Sturm sweep and the citation census in the code maps; the validation
count $80\to89$. No normative statement, tag or label moves; the normative findings that
arrived after v1.5.0 started are listed for the next $X$. Report
\file{docs/implementation/slot3/WP-11\_report.md});

v1.5.1: 2026-09-17 (hecke-extensions slot 3, package WP-9, a $Y$ bump:
the referee pass on v1.5.0 --- all ten resolutions accepted, \ref{L1} and
Appendix~\ref{app:status}'s counts re-verified mechanically, and seven defects repaired,
six of them in the $U_\ell$ item: Theorem~\ref{thm:backend}(1)'s ``$N\le2^{15}$'' is
$N\le23170$, the reach \S\ref{sec:heckedatum} already recorded; its $N\lesssim4580$ is a
\emph{measurement}, not a consequence of its criterion; item~(iv)'s
``$B_{\mathrm{red}}\le7C^{2}N^{2}$ by Lemma~\ref{lem:bounds}(6)'' is proved for the entry
point only and the rest is measured; Lemma~\ref{lem:bounds}(6)'s proof gains the two
shapes the reduction applies and names the coordinate that attains the order; and
Obligation~\ref{obl:oracle}'s ``off the Kummer loci $E$ is a line'' is measured, not
proved; the seventh is rule~\ref{A4}'s example, which still quoted
Obligation~\ref{obl:eta}(ii) in the form v1.5.0 corrected.
No normative statement, tag or label moves. Report
\file{docs/implementation/slot3/WP-9\_referee\_report.md});

v1.5.0: 2026-09-17 (hecke-extensions slot 3, package WP-8, an $X$ bump:
the seven items v1.4.2 deferred are resolved --- Remark~\ref{rmk:coneambiguity691}'s
arithmetic rewritten from the measured $\rank\rres^{1}=1$ (its conclusion unchanged) and
``the Eisenstein class'' mod $\ell$ read as the joint eigenspace, here and at
Obligation~\ref{obl:oracle}; Obligation~\ref{obl:eta}(ii), (iv) and its preamble
corrected; Lemma~\ref{lem:etaexists}(V) names the constructive choice;
\S\ref{sec:recordschema} item~5 gains the rule for a map and for a derived complex;
Obligation~\ref{obl:coneexact}(iv) the per-complex meaning of ``reported''; and
Lemma~\ref{lem:bounds} gains clause~(6), with Theorem~\ref{thm:backend}(1) corrected,
which is what the $U_\ell$ reach repair found --- together with the normative findings of
the slot-3 packages accepted at the time (WP-1, WP-5); no label added, renamed or
removed; status sentences are the slot close's; report
\file{docs/implementation/slot3/WP-8\_report.md});

v1.4.2: 2026-09-16 (hecke-extensions slot 2, package WP-13, the slot close,
a $Y$ bump: status only --- \S\ref{sec:recordschema} item~5 implemented at
\code{hecke\_report/v2}, the code maps and the obligation roster brought up to the slot-2
units, Obligation~\ref{obl:eta}(i)--(iv) discharged, no normative statement changed; the
findings the slot's packages handed to the next $X$ revision are listed in
Appendix~\ref{app:changelog}; report \file{docs/implementation/slot2/WP-13\_report.md});

v1.4.1: 2026-09-16 (hecke-extensions slot 2, package WP-10, a $Y$ bump:
referee pass on v1.4.0 --- each of the nine resolutions re-derived and accepted, no
normative statement changed; the prose and proof repairs are listed in
Appendix~\ref{app:changelog}, the largest being Remark~\ref{rmk:taubdrynoiota}(3)'s
$(2,1)$ exhibit, which v1.4.0 had placed on slots where no tie occurs; report
\file{docs/implementation/slot2/WP-10\_referee\_report.md});

v1.4.0: 2026-09-16 (hecke-extensions slot 2, package WP-9, an $X$ bump:
the nine open points that v1.3.3 recorded against chapter~\ref{sec:boundarydatum} are
resolved, as the slot-1 reports proposed and the owner accepted ---
Obligation~\ref{obl:taubdry}(v) is split into the three comparisons the chapter offers,
(v.a)--(v.c), and is degree $0$ alone at $\nu\ge2$; its clause (vii) is qualified by
$\chi_0(p)\neq\chi_1(p)$; Proposition~\ref{prop:circlemodel}(4) identifies the circle
model with $\Rel_{\partial}$ \emph{including} the $-I$ self-relations for every $n$;
Proposition~\ref{prop:bdrychar}(2) says what the right-coset sum is and (3) states the
$H^{1}\cong(\VR)_{T}$ coordinate; the canonical $\taub$ is Algorithm~\ref{alg:taubdry},
and Remark~\ref{rmk:taubdrynoiota} proves that no vertex rule makes it $\inv$-symmetric;
the record fields of \S\ref{sec:recordschema} item~5 are per complex; the $\ell\in\{2,3\}$
restriction equivariance is recorded in \S\ref{sec:cone} as Remark~\ref{rmk:conerestrictsmallell};
Remark~\ref{rmk:bdrynugeneral} records the general-$\nu$ shape as not stated);

v1.3.3: 2026-09-16 (hecke-extensions slot 1 close, a $Y$ bump: chapter
\ref{sec:boundarydatum} is now \emph{implemented} --- the certified $\Gamma$-level $\taub$
and the boundary operator $\Thb_g$ exist in the tree and
Obligation~\ref{obl:taubdry} is discharged in full, (i)--(vii), by two registered
validations --- so every status site that said ``three chapters state objects nothing
implements'' now says two; the code maps, the roster row, the oracle-family (A) run record,
the validation counts ($68\to77$) and the build-surface target list are updated; no normative
statement changed);

v1.3.2: 2026-09-16 (coherence pass after OQ3 and OC-1: ten stale or
over-broad sites repaired, among them the direction of Remark~\ref{rmk:iotapairing}'s
contingent prediction --- duality pairs each $\inv^{\ast}$-sign with the \emph{opposite}
one, the orientation-reversal factor having been omitted --- the ``for $n$ even''
qualification of Definition~\ref{def:iota}(iv)'s ``not a pullback along any lift'', and
the provenance counts of Appendix~\ref{app:changelog}; no normative statement changed);

v1.3.1: 2026-09-16 (referee pass on Lemma~\ref{lem:bulkperiodic}, the
audit the OC-1 decision asked for: every clause and step re-derived, tag \PROVED{}
confirmed, no normative statement changed; four prose repairs listed in
Appendix~\ref{app:changelog}; report
\file{docs/implementation/referee\_report\_2026-09-16\_bulkperiodic.md});

v1.3.0: 2026-09-16 (\textbf{OC-1 decided}: owner decision of
2026-09-16 adopting \ref{oc1a} with three riders --- \ref{open:ell2} loses its
\code{[open]} tag and becomes the binding truncation policy, Obligation~\ref{obl:coneexact}(iv)
and the record schema of \S\ref{sec:implementation} are narrowed to it; and the bulk
periodicity above the virtual cohomological dimension, named in
Remark~\ref{rmk:ell2window}(2) as an unverified citation, is \emph{proved} as
Lemma~\ref{lem:bulkperiodic} together with its Hecke equivariance, so that half of
\ref{oc1c}'s hypothesis is a theorem);

v1.2.0: 2026-09-16 (owner decision OQ3 resolved by \textbf{negating
Definition~\ref{def:iota}(iv)}: $\inv^{\ast}:=-\inv^{\ast}_{\mathrm{geo}}$, the geometric
pullback twisted by the sign character, so that ``$+$'' is \emph{by definition} the
$\inv^{\ast}$-eigenspace containing the Eisenstein classes; Obligation~\ref{obl:iota}(v)
becomes a gating definition-consistency oracle and holds on every cell, the split is
$(d_{\,n+2},d_{\,n+2}-1)$ as predicted, $\inv^{\ast}=-\id$ on $H^{0}(\Ga;\Fl)$; every $\pm$
label switches against the geometric reading; the cost is
\file{docs/implementation/oq3\_resolution.md});

v1.1.1: 2026-09-16 (status only, no normative statement changed:
Obligation~\ref{obl:iota}(ii),(iii),(v) implemented and run by
\code{HeckeIotaSignValidation} --- the aborting content passes; the \emph{recorded}
sign prediction of (v) is false for the then-current involution, the Eisenstein class
lying in its $H^{1}_{\mathrm{eq},-}$ with the mirror split
$(d_{\,n+2}-1,d_{\,n+2})$; resolution deferred to the owner, \code{[OQ3-PENDING]});

v1.1.0: 2026-09-16 (the canonical $\tau$ of Algorithm~\ref{alg:tau}
re-baselined by owner decision to be $\inv$-symmetric --- the $\invK$-equivariant tie-break
of step~(1b) and the tree-routed path templates of step~2; Obligation~\ref{obl:iota}(iv)
discharged \emph{positively} on the canonical $\tau$ and gating there;
Proposition~\ref{prop:iotacommutes}(3b) available at $\ell\in\{2,3\}$; every frozen
cochain-level reference value re-baselined, nothing on $\Heq{q}$ moved);

v1.0.0: 2026-09-16 (\textbf{the leading version field
becomes $1$}: rule~\ref{V5}\ref{V5ii} discharged --- \ref{V4}'s two freezes carried out
(\file{hecke\_action.tex} frozen at v1.0, \file{monodromy\_v5.tex} at v5.0, each with a
superseded-by banner) and \file{hecke\_extensions.tex} v1.3 revised to cite the nine labels
of \S\ref{sec:six} rather than define them, so the label namespace is unambiguous; no
normative statement of this document changed);

v0.19.0: 2026-09-16 (chapters 21--22 and Appendices B--C: the worked
$T_2$ datum and the numerical oracles, with erratum E1 applied to oracle~6 and the new
erratum E23 to oracle~7; the code map, the configuration surface and the record schema;
the STATUS index; the formalization notes. \textbf{The last chapter stub is written, but
the leading version field is held at $0$}: rule~\ref{V5} is amended in this version to
require an unambiguous label namespace as well, and nine labels of \S\ref{sec:six} are
still defined both here and in the live \file{hecke\_extensions.tex}; erratum E23);

v0.18.0: 2026-09-16 (chapters 18--20: the reduction primitive, the
$\tau$-table algorithm and the integer bounds; the certified floating-point filter; the
correctness obligations, with erratum E1 applied to \ref{obl:algebra} and the whole
obligation roster integrated; erratum E22);

v0.17.0: 2026-09-16 (chapters 16--17: the six groups, the long exact
sequence and Poincar\'e duality, re-homed from [E] by the scope cut; the general
orbifold model as definitions, with the construction deferred);

v0.16.0: 2026-09-16 (the $\tau$-existence lemma apportioned: existence
over $\Z$ for every $\ell$, homotopy uniqueness still under $6\in\cring^{\times}$;
notation collisions C18, C19; a corrected bibliography entry);

v0.15.0: 2026-09-16 (chapter 15: the pair homotopy $\etapair$, the
derived $\hder$, the cone action $\Phig$ and its well-definedness);

v0.14.0: 2026-09-16 (chapter 14: the boundary datum, and the
$\ell=2$ truncation policy \textbf{OC-1} as a tagged working assumption);

v0.13.0: 2026-09-16 ($\inv_{\Xbulk}\tau\inv_{\Xbulk}\neq\tau$ computed and recorded;
Obligation~\ref{obl:iota}(i),(iv) discharged);

v0.12.0: 2026-09-15 (chapter 9: the involution $\inv$);

v0.11.0: 2026-09-15 (the Changelog, Appendix~E, back-filled; rule (V5));

v0.10.0: 2026-09-15 (FD4 proved for the model; FD1/FD3 repaired for a non-effective
action);

v0.9.0: 2026-09-15 (Part~I re-homed from \file{monodromy\_v5.tex}; erratum E7
closed);

v0.8.0: 2026-09-15 (notation C17; two stale uniqueness claims);

v0.7.0: 2026-09-15 (chapters 12--13: induced systems, $\Xig$, $\Th_g$, the equivalent
formulations);

v0.6.0: 2026-09-15 (notation C16, the lift macro);

v0.5.0: 2026-09-15
(chapters 10--11: the Hecke datum, $\tau$, the definition);

v0.4.0: 2026-09-15 (notation
C15);

v0.3.0: 2026-09-15 (chapters 7--8: conventions, the model, $\Sym^n$ and the
contragredient);

v0.2.0: 2026-09-15 (chapter 1 and Appendix~A, the errata);

v0.1.0: 2026-09-15 (skeleton).

Supersedes \file{hecke\_action.tex} (frozen at v1.0, 2026-07-02) and
\file{monodromy\_v5.tex} (frozen at v5.0, 2026-07-03).

Companion specification: \file{hecke\_extensions.tex} v1.16.3 (its 1.17.0-pre.4, in the tree beside this version's pre.4, restates
\Espec{pr:latticesplit}~(R) on congruence-free classes, owner ruling on B-1 of 2026-10-02; its 1.17.0-pre.3, in the tree beside this version's pre.3, carries the referee F1's
repairs to \Espec{pr:kfonset} and leaves \Espec{pr:latticesplit}~(R) for the owner's ruling on B-1; its 1.17.0-pre.2, beside this version's pre.2, restates
\Espec{pr:latticesplit}~(R) again; neither is released; v1.17.0-pre.1 moves the pointer there: the companion's patch v1.16.3, owner ruling O4
of 2026-10-01, adds a \PROVED{} remark and a \CONJ{} prediction; its 1.17.0-pre.1, drafted with this version for the same referee, is not released;
until v1.17.0-pre.1 the pointer read v1.16.2: v1.16.0 moves the pointer there: the companion's v1.16.0, owner-signed on
2026-10-01 after the slot-15 close (ruling S15-O2; \Espec{pr:latticesplit} restated), its patch v1.16.1 (census edition~3, evidence and pointers
only) and its patch v1.16.2 (one verified locator, released with this document's v1.16.0); until v1.16.0 the pointer read v1.15.0:
v1.14.0 moves the pointer there: the companion's v1.15.0,
owner-released on 2026-09-28 together with this document's v1.14.0, carries evidence and pointer updates only; until v1.14.0 the
pointer read v1.14.0: v1.14.0-pre.1 moves the pointer there: the companion's
v1.14.0, owner-released on 2026-09-27 at the slot-13 close after this document's v1.13.0, touched only the companion; until
v1.14.0-pre.1 the pointer read v1.13.1: v1.13.0 leaves it there, v1.13.0-pre.1 and pre.2 leave it there, slot-13 WP-2
and WP-3 touching only this document; v1.12.1 moves the pointer there: the
companion's v1.13.0 and v1.13.1, slot-12 S12-FIN, the finality certificates, touched only the companion;
until v1.12.1 the pointer read v1.12.1: v1.12.0-pre.1 leaves it there, slot-12
WP-2 touching only this document; v1.12.0, the owner-signed release, leaves it there as well, the
release touching only this document; v1.11.1, published in the same
package, slot-11 WP-17, as the companion's v1.12.1; v1.11.0 had it at v1.12.0, published in the same
package, slot-11 WP-15; v1.11.0-pre.1 and v1.11.0-pre.2 had it at v1.11.0,
the companion being at its pre-releases v1.12.0-pre.1 (slot-11 WP-4) and v1.12.0-pre.2 (slot-11
WP-18c) meanwhile --- and, before either of them, at v1.11.1 (\texttt{4999a21}, slot-10 WP-4's W2
documents), which the pointer did not follow; the companion's v1.11.1 and v1.12.0-pre.1 both preceded
this document's v1.11.0-pre.1 (v1.11.1: the WP-15 referee's m-4); v1.10.1, at whose commit the
companion reads v1.11.0; v1.10.0; the companion's
v1.11.0-pre.1 of slot-10 WP-2 and v1.11.0 of slot-10 WP-15, the same package as this version,
which publishes it in its second phase --- at this version's commit the companion reads
v1.11.0-pre.2; v1.9.1 had it at v1.10.1, the companion's referee-repair $Y$ bump at the slot-9 close; v1.9.0 had it at v1.10.0, the companion's
v1.10.0-pre.1 of slot-9 WP-2 and v1.10.0 of slot-9 WP-15, the same package as this version,
which publishes it in its second phase --- at this version's commit the companion reads
v1.10.0-pre.2; v1.8.1 had it at v1.9.1, the companion's
v1.9.0-pre.1 and v1.9.0 $X$ bumps of slot~8 (\Espec{sec:ring}) and its v1.9.1 referee pass;
v1.8.0 and v1.7.4 had it at v1.8.2, the companion's own
status bump at the slot-7 close, after its v1.8.0-pre.1 and v1.8.0 $X$ bumps and the v1.8.1
referee pass; v1.7.3 had it at v1.7.2, the companion's status bump at the slot-6 close, after
its v1.7.0 $X$ bump and v1.7.1 referee pass; v1.7.2 had
it at v1.6.2, the companion's slot-5 status bump; v1.7.1 had it at v1.6.1, repaired in the same referee pass,
and v1.7.0 had moved it to v1.6.0 from v1.5.1, the pointer
having stood at v1.3.3 since that companion's v1.4.0, three revisions).

\subsection*{The title page of v1.25.0-pre.7}
\addcontentsline{toc}{subsection}{The title page of v1.25.0-pre.7}
The author line and the abstract of v1.0.0 to v1.25.0-pre.7, verbatim (the abstract as of v1.25.0-pre.7).

\begin{quote}
Orbifold Cohomology of $\GL(2)$
\end{quote}

\begin{quote}
\textbf{Complete as a document, incomplete as a programme.} Every chapter
\ref{sec:scope}--\ref{sec:implementation} and every appendix is written, no stub
remains, and at v1.0.0 the label namespace is unambiguous as well: \ref{V4}'s two
supersession freezes are done, and the nine labels of \S\ref{sec:six} that the live
companion \file{hecke\_extensions.tex} also defined are cited there rather than
restated, so a reference \code{[C,~cl:pd]} resolves to exactly one statement. Both
criteria of \ref{V5} are therefore met and the leading version field is $1$.
\textbf{That is a statement about this document and not about the programme.} Two chapters ---
\ref{sec:orbifoldmodel} (the general orbifold model) and, since v1.8.0, \ref{sec:integral}
(the integral line: its free model, comparison map and descent) --- state objects that
\emph{nothing in the implementation builds}, and say so in their own code maps and in
\S\ref{sec:notimplemented} (v1.9.0: \ref{sec:integral}'s free model, its comparison map and
its integral groups for $n\le20$ are built, its lifted Hecke datum and descent are not;
v1.10.0: the lifted Hecke datum is built for $T_2$, $T_3$, $T_5$, with the class $e$ of
Lemma~\ref{lem:bulkperiodic} and the integral Hecke action read basis-free at small $n$ --- the
descent assembly is still not; v1.11.0-pre.1: the two chapters of the new Part~VII,
\ref{sec:dicksontheta} and \ref{sec:ordinary}, are written \emph{before} their implementation, and
their code maps say so; v1.11.0: \ref{sec:dicksontheta} is built and tested, and the ordinary projector
of \ref{sec:ordinary} is built at $p=\ell$ on the $\ell$-adic digit route, both recorded in their code
map; \ref{open:ordprecision} stays open; v1.11.1: its torsion half is answered by
Remark~\ref{rmk:ordtorsionfree}, the per-cell bound stays open; v1.12.0-pre.1: a third chapter of
Part~VII, \ref{sec:hida}, is written \emph{before} its implementation and its code map says so, and
Corollary~\ref{cor:ordprecision} answers the per-cell bound, the item's status pending the release;
v1.12.0: released, and \ref{open:ordprecision}'s tag moves \OPENSTAT{} $\to$ \PROVED{};
v1.13.0-pre.1: a fourth chapter of Part~VII, \ref{sec:hidafull}, the level-$\ell$ model, the stabilisation and
the lattice datum, is written \emph{before} its implementation and its code map says so; v1.14.0-pre.1: a fifth,
\ref{sec:levelN}, level $N$ for every $N$ --- the composite coset table, the cokernel presentation and its Hecke action --- likewise);
\ref{sec:cone} (the pair homotopy and the cone), which did so
in whole or in part until v1.7.2, is implemented in full since v1.7.3, and
\ref{sec:boundarydatum} (the boundary datum), which did so until v1.3.2, since v1.3.3, its
obligation discharged (\S\ref{sec:obligationroster}) (v1.8.1: this sentence named
\ref{sec:cone} as unbuilt until then, stale since v1.7.3); the one
construction item that was open, \ref{open:ell2}, was \textbf{decided by the owner on
2026-09-16} (\ref{oc1a}, v1.3.0) and its dependent sites are enumerated in
\S\ref{sec:oc1sites}. Every numbered statement is indexed with
its status and its provenance in Appendix~\ref{app:status}. What each version of this
document changed, statement by statement, is Appendix~\ref{app:changelog}. It is cited as
\textbf{[C]} --- the same short key the
$53$ citing sites in the C++ source already use --- and it carries verbatim every
\code{\textbackslash label} those sites cite.
\end{quote}
\par\endgroup

\subsection*{The rules this appendix follows}

\begin{enumerate}[label=(CL\arabic*)]
\item\label{CL1} \textbf{One line per statement whose content changed, naming the
  label} (\ref{V3}). A statement written into a stub that this document had already
  declared counts as changed: before, the label named \code{[to be written]}; after, it
  names an assertion. Statements transported verbatim modulo the notation crosswalk
  from a named source, with nothing added and nothing corrected, are not given a line
  each; they are named together in that version's \emph{re-homed unchanged} roster,
  which lists every one of them, so the test ``the contract must say which statements
  moved'' is still met and ``editorial pass'' is still not an entry.
\item\label{CL2} \textbf{This document's history only.} The changelog says what
  changed between two versions of \emph{this} document --- which for this purpose is
  \file{gl2\_contract.tex} together with the files it is built from:
  \file{gl2\_notation.tex} (chapter~\ref{sec:notation}), \file{gl2\_preamble.tex} (the
  macros \ref{A5} enforces) and \file{gl2\_refs.bib} (the citations \ref{A2} demands). The defects of the superseded
  corpus are Appendix~\ref{app:errata}'s subject and are not restated here: an entry
  that exists because an erratum was applied names the erratum number
  (\cl{err:e1}--\cl{err:e21}) and stops. Adding or amending an erratum entry is itself
  a change to this document and does get a line, under that entry's label.
\item\label{CL3} \textbf{A statement written and later revised gets one line per
  version in which it changed}, not one summary line. The changelog is a per-version
  record, and ``as of [C] v$X.Y$'' must be answerable: \cl{prop:tworoutes} was written in
  v0.7.0 asserting a uniqueness that the same version's \cl{prop:mu}(3) had already made
  false, and was corrected in v0.8.0. It therefore takes a line in each, and collapsing
  the two would hide precisely what v0.7.0 said.
\item\label{CL4} \textbf{Provenance is marked, and only on statements.} A line whose
  subject is a numbered statement --- an environment of kind \code{thm}, \code{lem},
  \code{prop}, \code{cor}, \code{def}, \code{rmk}, \code{alg}, \code{obl}, \code{ass},
  \code{ex}, \code{cl} or \code{open} (the last added to this list at v1.3.0, as the
  v0.14.0 entry asked) --- carries exactly one of three tags. \PNEW{}: the statement
  appears in no superseded document and in no project note; it is first asserted and
  first proved here. \PRH{}: transported from a named superseded document, content
  unchanged beyond notation and any named erratum. \PFIX{}: transported, but the source
  statement was false, circular, ill-typed, or missing a hypothesis its own proof uses;
  the line names the erratum, or says that the defect was found in the pass that
  repaired it. The tag attaches to the \emph{content}, not to the label, so a preserved
  label whose statement is strictly generalised here carries \PNEW{} for the
  generalisation and the line says so.

  \textbf{Section, display, rule, erratum and appendix labels carry no tag} ---
  \code{sec:}, \code{eq:}, the \ref{L1}-style item labels of an enumeration,
  \code{err:}, \code{app:}. A subsection is a container and a numbered display is a
  formula; neither is a statement, and a \emph{new container for re-homed mathematics}
  tagged \PNEW{} would read as new mathematics, which is exactly the misreading this
  rule exists to prevent. Where such a label shares a line with a statement, the tag on
  that line is the statement's and says nothing about the container; where it stands
  alone, the line's prose says what happened to it.

  This is the only provenance marking the document currently carries; when
  Appendix~\ref{app:status} is written it should carry the same three values as a
  column, so that the marking reaches the statement itself and not only its changelog
  line.
\item\label{CL5} \textbf{Normative and non-normative entries are separated}, because
  \ref{V2}'s bump rule reads off the split: a version with any entry in its
  \emph{Normative} list increments $X$. ``Change to a normative statement'' includes
  \emph{adding} one --- otherwise writing a chapter would be a $Y$ bump. Chapter
  \ref{sec:notation} is declared normative and exhaustive and \ref{A5} binds every
  later statement to it, so a change to its rules, its collision table or its crosswalk
  is normative. Recording an erratum, a formalization note, a code map, and a design
  direction that nothing in the document depends on are not.
\item\label{CL6} \textbf{A version is one commit.} The commits below each did one
  coherent thing, so no version splits or merges them (the first ten versions were
  assigned to their commits retrospectively, see the box above). The commit hash is
  given with each version where it is known, and is the evidence; \code{[hash]} marks an
  entry written before its commit existed and not yet back-filled.
\end{enumerate}

\noindent
Under \ref{CL5} every version so far increments $X$: ten writing slots and one repair
pass each added or changed a normative statement. The exceptions are v1.0.0, which moved
the leading field under \ref{V5}, and the $Y$ bumps v1.1.1 (evidence only), v1.3.1 (a
referee pass), v1.3.2 (a coherence pass), v1.3.3 (evidence and status only), v1.4.1
(a referee pass), v1.4.2 (status only), v1.5.1 (a referee pass), v1.5.2 (status
only), v1.6.1 (a referee pass), v1.6.2 (status only), v1.7.1 (a referee pass), v1.7.2,
v1.7.3 and v1.7.4 (status only), v1.8.1 (a referee's wording and status; the list
completed at v1.8.1, having stopped at v1.5.2) and v1.9.1 (a referee's repairs and status), which
changed prose and no statement --- which is itself information about how this document
has been built.

\subsection*{Provenance: the counts, and what they do not measure}

Over all versions to date, with the \ref{CL1} rosters folded in, and counting
\emph{statement-kind labels} as \ref{CL4} defines them:

\begin{center}
\begin{tabular}{@{}lrl@{}}
\toprule
& labels & of which \\
\midrule
\PNEW{}   & $236$ & $148$ assertions or definitions, $88$ remarks (v1.4.0: \cl{alg:taubdry}; \cl{rmk:taubdrynoiota}, \cl{rmk:bdrynugeneral}, \cl{rmk:conerestrictsmallell}; v1.8.0: fifteen and five, \S\ref{sec:integral} and \cl{rmk:bocksteinexp}; v1.9.0: \cl{rmk:certifiedeperdiff}; v1.10.0: \cl{rmk:strictliftbuilt}, \cl{rmk:bredonhecke}; v1.11.0-pre.1: fifteen and five, \cl{prop:relationsmodel} and Part~VII; v1.11.1: \cl{rmk:ordtorsionfree}; v1.12.0-pre.1: four and two, \S\ref{sec:hida}; v1.12.0-pre.3: two, \cl{prop:bdrytranscription} and \cl{lem:bdrytorus}; v1.13.0-pre.1: ten and three, \S\ref{sec:hidafull}; v1.13.0-pre.2: \cl{rmk:levelscaling}; v1.14.0-pre.1: nine and three, \S\ref{sec:levelN}; v1.15.0-pre.1: eighteen and five, \S\ref{sec:nebentypus} and \S\ref{sec:saturation}; v1.16.0-pre.1: five and five, \S\ref{sec:levelNdegsec} and \S\ref{sec:isomodelell}) \\
\PRH{}    & $45$ & $4$ \code{thm}, $3$ \code{lem}, $9$ \code{prop}, $4$ \code{ass},
                   $12$ \code{def}, $3$ \code{rmk}, $1$ \code{ex}, $2$ \code{cl},
                   $7$ \code{obl} \\
\PFIX{}   & $38$ & $7$ \code{thm}, $5$ \code{lem}, $9$ \code{prop}, $3$ \code{cor},
                   $1$ \code{ass}, $1$ \code{def}, $2$ \code{alg}, $6$ \code{rmk},
                   $1$ \code{cl}, $3$ \code{obl} \\
\bottomrule
\end{tabular}
\end{center}

\noindent
Of the $236$ \PNEW{} statements ($226$ until v1.16.0-pre.1, $203$ until v1.15.0-pre.1, $191$ until v1.14.0-pre.1, $190$ until v1.13.0-pre.2, $177$ until v1.13.0-pre.1, $175$ until v1.12.0-pre.3, $169$ until v1.12.0-pre.1, $168$ until v1.11.1, $148$ until v1.11.0-pre.1, $146$ until v1.10.0, $145$ until v1.9.0, $125$ until v1.8.0) the $88$ remarks are new \emph{commentary} on re-homed
mathematics --- what a theorem does not say, which choices an object depends on, what a
diagnostic may not be branched on --- and are new writing rather than new mathematics
(v1.4.0's \cl{rmk:taubdrynoiota} carries a proof and \cl{rmk:conerestrictsmallell} a
\COMP{} record, and both are listed with the boundary roster's remarks below; v1.11.1's
\cl{rmk:ordtorsionfree} carries a \PROVED{} corollary and is listed with Part~VII below).
The $130$ assertions and definitions ($125$ until v1.16.0-pre.1, $116$ until v1.14.0-pre.1, $106$ until v1.13.0-pre.1 --- the figure pre.1 left unchanged here, corrected at
pre.2 ---, $104$ until v1.12.0-pre.3, $100$ until v1.12.0-pre.1, $85$ until v1.11.0-pre.1, $70$ until v1.8.0) are the new mathematics, and they are listed below.
Counting every label named on a tagged line, containers included, $194$ sit on \PNEW{}
lines and $80$ on \PRH{} lines or in a roster (no such figure is carried for \PFIX{}
lines: the versions before v0.18.0 did not record one). Those raw figures are \emph{not} a
provenance split --- the excess over $125$ and $45$ is containers named alongside a
statement --- and quoting them as one is the misreading \ref{CL4} exists to stop. Before
that rule was imposed the tags reached $108$ labels, which would have read as $108$ new
theorems.

\begin{contractbox}
\textbf{Pages and labels are different measurements, and neither refutes the other.}
The architecture costed this document at roughly $62\,\%$ re-homing, $11\,\%$
re-casting and $27\,\%$ new \emph{by page}, and writing slots 5--6 measured the
re-homing rate at $1.0\times$ over $14$ pp of \file{monodromy\_v5.tex} source. Both
figures stand. The table above counts \emph{labels}, and the two disagree for a
structural reason: re-homed theorem-proof prose is label-sparse --- one
\code{\textbackslash label} can carry a page of proof --- while the apparatus this
consolidation added is label-dense, since chapter~\ref{sec:notation}'s tables, the
rules \ref{L1}--\ref{L6}, \ref{A1}--\ref{A6} and \ref{V1}--\ref{V5}, the twenty-two errata
entries and the per-statement remarks each earn a label within a few lines. The
conclusion a reader should draw from the two together is that \textbf{re-homing
dominates by page and by theorem, while the new apparatus dominates by label} --- not
that either measurement is wrong.
\end{contractbox}

\paragraph{The $85$ statements that are new here, and why they are the ones to
scrutinise.} Each asserts something that no superseded document and no project note
states anywhere. Every other statement of this contract can be checked against a source
--- that is what \PRH{} and \PFIX{} mean, and it is how the defects of
Appendix~\ref{app:errata} were found. These cannot be: an error in one of them cannot
be caught by comparing documents, only by reading the proof or by running the
implementation. They are therefore the first thing a referee should attack and the
first thing a formalisation should cover, and they are exactly the lines this appendix
tags \PNEW{} whose subject is not a \code{rmk}.

\begin{itemize}\setlength{\itemsep}{0pt}
\item \emph{Part~I, the framework} ($19$): \cl{def:relationdatum},
  \cl{def:orderedrelation}, \cl{def:quotientcomplex}, \cl{def:saturated},
  \cl{def:bdryint}, \cl{def:kernelN}, \cl{ass:fund_domain}~\ref{FD4},
  \cl{lem:eqgroupoid}, \cl{lem:fd4criterion}, \cl{lem:model7coincidences},
  \cl{def:naivecochains}, \cl{def:eqcohomology}, \cl{cor:threecategories},
  \cl{prop:Bqimage}, \cl{def:storedbases}, \cl{prop:ACq}, \cl{prop:AZq},
  \cl{prop:projector}, \cl{alg:cohomology}.
\item \emph{Part~II, the model and its coefficients} ($11$): \cl{def:model7},
  \cl{def:model7rel}, \cl{prop:model7h4}, \cl{def:contragredient}, \cl{prop:pairing},
  \cl{lem:symdual}, \cl{def:iota}, \cl{lem:iotamodel}, \cl{prop:iotacommutes},
  \cl{prop:iotasplit}, \cl{obl:iota}.
\item \emph{Part~III, the Hecke datum and the maps} ($6$): \cl{prop:cosets} --- the
  extension to every $\nu$; the $\nu=1$ case is re-homed --- \cl{def:cosetreps},
  \cl{lem:cocycle}, \cl{prop:fullassembly}, \cl{prop:mu}, \cl{prop:tworoutes}.
\item \emph{Part~IV, the boundary} ($13$; $11$ until v1.12.0-pre.3): \cl{prop:bdrycohom}, \cl{prop:bdrycosets},
  \cl{def:taubdry}, \cl{lem:bdrycomponent}, \cl{def:bdryop}, \cl{prop:bdryop},
  \cl{obl:taubdry}, \cl{lem:ell2unbounded}, \cl{lem:bulkperiodic} (v1.3.0: the bulk
  periodicity above the vcd and its Hecke equivariance, a five-lemma argument on
  \cl{cl:amalgam}(2) with an explicitly constructed class $e\in H^{2}(\Ga;\Z)$ --- the
  statement no source of this corpus makes), \cl{open:ell2}, \cl{alg:taubdry} (v1.4.0:
  the canonical $\taub$ as a procedure, with its four proved properties),
  \cl{prop:bdrytranscription} (v1.12.0-pre.3: the Shapiro transcription in degree $1$, at every $n$,
  $\cring$ and $p$, with an $H^{1}$ formula at every $\nu$ --- the statement \cl{prop:bdrychar}(3) had
  left unproved), \cl{lem:bdrytorus} (v1.12.0-pre.3: the boundary operators away from $\ell$ factor
  through the torus; a referee should read \cl{prop:bdrytranscription}'s steps~1--4 first).
  (\cl{def:bdrygroup} is \PRH{} and \cl{prop:circlemodel}, \cl{lem:taubdryexists},
  \cl{prop:bdrychar} are \PFIX{}; they are checkable against a source and are therefore
  not in this roster. Three v1.4.0 remarks --- \cl{rmk:taubdrynoiota}, which carries a
  proof, \cl{rmk:bdrynugeneral} and, in \S\ref{sec:cone}, \cl{rmk:conerestrictsmallell}
  --- are \PNEW{} commentary.)
\item \emph{Part~IV, the pair homotopy and the cone} ($15$): \cl{def:pair},
  \cl{lem:conerelative}, \cl{def:cosettransport}, \cl{lem:cosettransport},
  \cl{def:pairhomotopy}, \cl{lem:etaunique}, \cl{def:coneaction},
  \cl{lem:derivedhomotopy}, \cl{thm:conewelldef}, \cl{prop:coneindep},
  \cl{prop:strictpair}, \cl{prop:coneintegral}, \cl{prop:coneroutes}, \cl{obl:eta},
  \cl{obl:coneexact}. \textbf{This is the chapter with no external corroboration at
  all}: its only sources are project notes (the OQ2 verdict and the plan), which are
  not sources of record, so every statement here is either first proved in this
  document or repaired from one (\cl{lem:etaexists}, \cl{lem:conesign},
  \cl{prop:coneambiguity} are \PFIX{}, for \cl{err:e5}, \cl{err:e13},
  \cl{err:e2}/\cl{err:e3}). Four remarks --- \cl{rmk:conesteinberg},
  \cl{rmk:conestrict}, \cl{rmk:coneambiguity691}, \cl{rmk:conetorsion} --- are
  \PNEW{} as well and carry mathematics a referee should read
  (\cl{rmk:conesteinberg}(ii) in particular: it is where the Hecke relations on
  $\Hcpt{\bullet}$ become a theorem rather than a citation).
\item \emph{Part~IV, the six groups and the general orbifold model} ($7$):
  \cl{cl:amalgam} --- the Mayer--Vietoris sequence in \emph{every} degree over
  \emph{every} coefficient ring, and the torsion bound, where [E] stated only the
  $\ell\ge5$ consequence --- \cl{lem:h2torsion}, \cl{cl:boundary}, \cl{cl:les} (the
  hypothesis-apportioned form, with $\Hcpt{0}=0$ unconditional), \cl{cl:chi} (clause~(3)
  proved), \cl{def:borelmodel}, \cl{thm:specialisation}. (\cl{cl:dimH1} and
  \cl{cl:bockstein} are \PRH{} and \cl{cl:pd} is \PFIX{}, so they are not in this
  roster. Seven further remarks --- \cl{rmk:sixscope}, \cl{rmk:bdryrepairs},
  \cl{rmk:chicoincide}, \cl{rmk:pdscope}, \cl{rmk:borelfree}, \cl{rmk:mpcaveat},
  \cl{rmk:orbifoldrefine} --- are \PNEW{} commentary.)
\item \emph{Part~V, realisation and certification} ($1$): \cl{lem:predshapes} --- every
  branch decision of \cl{alg:tau} and \cl{def:primitives} is one of exactly four
  two-term integer dot products, the fourth of which the superseded contract did not
  list. Part~V is the most re-homable part of the document and this is the only
  assertion in it that no source states; its four other \PNEW{} labels
  (\cl{rmk:primitiveclosure}, \cl{rmk:certoverGeometry}, \cl{rmk:btabsmall},
  \cl{rmk:thresholdrounding}) are remarks, two of them carrying a
  \PROVED{} or a \COMP{} claim a referee should still read
  --- \cl{rmk:thresholdrounding}, which is why the binding filter constant is
  $2^{-50}B^2$ and not the sharper bound the proof gives, and \cl{rmk:btabsmall}.
\item \emph{Part~VI, the integral line} ($15$, v1.8.0): \cl{def:ringinforce},
  \cl{prop:ringinforce}, \cl{prop:bredon}, \cl{def:isomodel}, \cl{prop:isomodel},
  \cl{prop:fracture}, \cl{prop:cohglue}, \cl{prop:descent}, \cl{prop:homvanish},
  \cl{lem:strictlifting}, \cl{prop:strictlift}, \cl{prop:badprimes}, \cl{prop:certifiede},
  \cl{prop:crtcharpoly}, \cl{def:ringsemantics}. Like the cone chapter, its sources are
  project notes only (the integral-line design brief and its audit, the $\ell$-adic
  investigation), which are not sources of record; two registered validations test
  \cl{prop:bredon}(2) and \cl{prop:crtcharpoly}, nothing tests the rest (v1.8.0; since v1.9.0
  and v1.10.0 more do, \cl{rmk:integralevidence}, which lists them). Five remarks ---
  \cl{rmk:cafive}, \cl{rmk:truncdevice}, \cl{rmk:integralevidence}, \cl{rmk:integralscope}
  and, in \S\ref{sec:six}, \cl{rmk:bocksteinexp} --- are \PNEW{}; three carry a \PROVED{}
  clause a referee should read (\cl{rmk:bocksteinexp}(1),(2) above all). v1.9.0 adds a sixth,
  \cl{rmk:certifiedeperdiff} (\PROVED{}), and states \cl{rmk:bocksteinexp}(2) at every prime.
  v1.10.0 adds a seventh and an eighth: \cl{rmk:strictliftbuilt} (\COMP{}, registered --- the
  lifted datum as built, with its certificates and measured sizes) and \cl{rmk:bredonhecke}
  (commentary --- what this document does \emph{not} compare at $\ell\in\{2,3\}$); and three
  registered validations now test \cl{prop:strictlift}(1),(3),(4), \cl{lem:bulkperiodic}(1)--(3)
  and the Hecke action on the free model (\cl{rmk:integralevidence}). v1.11.0-pre.1 adds
  \cl{prop:relationsmodel}: \cl{thm:relations}'s recursion on the free model over every ring.
\item \emph{Part~VII, the Dickson invariant, $\thL$ and the ordinary projector} ($14$ assertions or
  definitions and $5$ remarks, v1.11.0-pre.1): \cl{def:dicksoninv}, \cl{lem:dicksontwist},
  \cl{prop:dicksonH0}, \cl{prop:heckeH0}, \cl{lem:heckefibre}, \cl{def:Qn}, \cl{lem:Qnfunctions},
  \cl{def:theta}, \cl{prop:thetatwist}, \cl{thm:thetales}, \cl{def:ordinary}, \cl{prop:ordinary},
  \cl{cor:thetanonordinary}, \cl{open:ordprecision}; remarks \cl{rmk:heckeH0plan},
  \cl{rmk:thetasmallell}, \cl{rmk:thetaclassical}, \cl{rmk:thetacheckable}, \cl{rmk:ordinaryeis};
  v1.11.1 adds a sixth remark, \cl{rmk:ordtorsionfree} (\PROVED{}: the torsion of $H^{1}\otimes\Z_\ell$ is
  non-ordinary at $\ell\ge5$, $n\ge1$ --- a corollary of \cl{cor:thetanonordinary}(3) and \cl{prop:ordinary}).
  Its sources are the slot-11 plan and its audit (project notes, not sources of record) and the
  literature of \cl{rmk:thetaclassical}, none of it in \file{docs/literature/}; a referee should read
  \cl{prop:heckeH0}(2)--(4) first, since it corrects the plan.
\item \emph{Part~VII, \S\ref{sec:hida}, the foundation of Hida theory at level~$1$} ($4$ assertions or
  definitions and $2$ remarks, v1.12.0-pre.1): \cl{prop:torsionprepass}, \cl{cor:ordprecision},
  \cl{lem:ordexponent}, \cl{def:hidatower}; remarks \cl{rmk:ordexponentuse} (\PROVED{}: sharpness) and
  \cl{rmk:hidacheckable} (the prediction list). Its sources are the slot-12 plan (a project note, not a
  source of record) and Hida's papers and book, none of them in \file{docs/literature/}; a referee should
  read \cl{prop:torsionprepass}(1),(2),(6) and \cl{cor:ordprecision}(1),(2) first, since they depart from
  the plan's statements.
\item \emph{Part~VII, \S\ref{sec:hidafull}, control, level $\ell$ and the lattice} ($10$ assertions or definitions
  and $3$ remarks, v1.13.0-pre.1; a fourth remark, \cl{rmk:levelscaling} (\BSTD{}: the regimes, the obstacles
  (O1)--(O6), the Manin--Merel action by pointer), v1.13.0-pre.2): \cl{def:levelellmodel}, \cl{prop:levelcorrespondence}, \cl{prop:levelellU},
  \cl{prop:oldspace}, \cl{prop:levelorbit}, \cl{lem:levelelltorsion}, \cl{prop:stabilisation},
  \cl{thm:hidacontrol}, \cl{thm:hidaweight}, \cl{def:hidalattice}; remarks \cl{rmk:levelcost} (commentary: the
  cost model), \cl{rmk:stabconvention} (\PROVED{}) and \cl{rmk:hidacheckable2} (the prediction list). Its sources are
  the slot-13 plan and its owner decisions (project notes) and Hida's two 1986 papers, read from the owner's
  read-only copies with every locator checked against the page; a referee should read
  \cl{prop:levelellU}, \cl{prop:oldspace}(3) and \cl{rmk:stabconvention} first, since the convention they fix makes
  $e\circ\mathrm{dg}_\ell$, not Hida's $e\circ\res$, the integral stabilisation. v1.14.1 adds \cl{lem:rhorank} (the number of
  non-zero invariant factors of $\rhon$ at a matrix singular modulo $\ell$, and its rank-one closed form over $\Fl$ --- the
  $U_\ell$ lever's rank claim, checked at every block by the implementation).
\item \emph{Part~VII, \S\ref{sec:levelN}, level $N$: the composite coset table, the cokernel presentation and its Hecke action}
  ($9$ assertions or definitions and $3$ remarks, v1.14.0-pre.1): \cl{def:cosetrepsN}, \cl{lem:cosetrepsN}, \cl{def:levelNmodel},
  \cl{lem:treeres}, \cl{lem:manin}, \cl{def:kappaN}, \cl{lem:levelduality}, \cl{lem:merelpath}, \cl{prop:heilbronn}; remarks
  \cl{rmk:heilbronnconv} (\PROVED{}: the transcription and its trap), \cl{rmk:heilbronncost} (the cost model) and
  \cl{rmk:levelNcheckable} (the prediction list). Its sources are the slot-14 plan and the owner's standing direction on level $N$
  (project notes), and Manin, Merel, Stein, Cremona, Mazur, Merel 1996 and Wake--Wang-Erickson, none of them verified here
  ([S14] Q6's default, rule~\ref{A2}) --- every proof in the chapter is complete without them; a referee should read
  \cl{lem:levelduality}(5) and Step~3 of \cl{prop:heilbronn}'s proof first, since the transcription (the adjugate, the row action,
  the admissible pairs at $p\mid N$) lives there.
\item \emph{Part~VII, \S\ref{sec:nebentypus}, nebentypus} ($10$ assertions or definitions and $4$ remarks, v1.15.0-pre.1):
  \cl{def:nebdata}, \cl{lem:nebdata}, \cl{lem:teichmuller}, \cl{lem:inducedchar}, \cl{lem:gammaonetable}, \cl{def:nebmodel},
  \cl{prop:nebdecomp}, \cl{prop:nebcorr}, \cl{thm:nebrelations}, \cl{prop:neborbit}; remarks \cl{rmk:nebfield}, \cl{rmk:nebheilbronn}
  (\PROVED{}), \cl{rmk:nebcost} (the cost model) and \cl{rmk:nebcheckable} (the prediction list). Its sources are the slot-15 plan and
  the owner's ruling Q2 of 2026-09-30 (project notes); the literature (Shimura \S3.3, Cohen--Oesterl\'e, Ribet 1977, Ash--Stevens) is
  cited at chapter level and \textbf{not} verified; a referee should read \cl{prop:nebcorr}(2),(4), \cl{prop:neborbit}(4) and
  \cl{rmk:nebheilbronn} first.
\item \emph{Part~VII, \S\ref{sec:saturation}, the $\Z/\ell^{K}$ Hecke algebra} ($8$ assertions or definitions and $1$ remark,
  v1.15.0-pre.1): \cl{def:satdata}, \cl{prop:saturation}, \cl{prop:satdefect}, \cl{lem:heckerankone}, \cl{lem:freepart},
  \cl{cor:stablecert}, \cl{prop:cohforms}, \cl{prop:hidaalgconst}; remark \cl{rmk:thetaforms}. Its sources are the slot-14 Loewy note
  \file{docs/implementation/slot14/WP-M1\_loewy\_note.md} (Prop.~1, \S2, \S5.0), the proposals P15, P17, P18, P20 of slot-14 WP-6 and
  WP-D1, and the slot-15 plan's WP-19~(c) (commit \texttt{c1b45cd}); a referee should read \cl{prop:saturation}(3) and
  \cl{cor:stablecert} first, since the certified verdict of slot-15 WP-5 rests on them.
\item \emph{Part~VII, \S\ref{sec:levelNdegsec}, degeneracy maps on the presentation and level $N$ at scale} ($2$ assertions and $3$ remarks,
  v1.16.0-pre.1): \cl{prop:kappamodel}, \cl{prop:degenpres}; remarks \cl{rmk:degenpresconv} (\PROVED{}, \COMP\textsuperscript{a}), \cl{rmk:levelNscale} (the
  cost and memory model) and \cl{rmk:levelNcheckabletwo} (\ref{LN7}--\ref{LN9}). Its sources are the slot-16 plan (WP-12 items (a), (c)) and
  v1.15.0-pre.3's unrefereed sentence, which it supersedes; a referee should read Step~3 of \cl{prop:degenpres}'s proof (the re-indexing and the
  admissibility condition) and \cl{prop:kappamodel}(2) first.
\item \emph{Part~VII, \S\ref{sec:isomodelell}, the $\ell$-local isotropy model} ($3$ assertions or definitions and $2$ remarks, v1.16.0-pre.1):
  \cl{def:isomodelell}, \cl{prop:isomodelell}, \cl{prop:ellthreeclosed}; remarks \cl{rmk:isomodellmp} (the MP constraint) and
  \cl{rmk:isomodellcheckable} (\ref{IL1}--\ref{IL4}). Its source is \file{docs/implementation/slot14/WP-Q1\_small\_ell\_orbifold\_note.md}
  \S\S2--4; a referee should read \cl{prop:ellthreeclosed}(3)'s spectral-sequence argument and \cl{prop:isomodelell}(4) first.
\item \emph{Part~VII, \S\ref{sec:nebles}, the level-change cone and inflation--restriction} ($18$ assertions or definitions, $1$ obligation and $3$ remarks,
  v1.22.0-pre.3): \cl{def:levelpair}, \cl{lem:nebcompat}, \cl{def:levelres}, \cl{lem:levelresmorph}, \cl{def:levelcone}, \cl{lem:levelconeexact},
  \cl{prop:levelconeshapiro}, \cl{prop:modeltorsion}, \cl{prop:levelsplit}, \cl{prop:levelconehecke}, \cl{def:normalpair}, \cl{prop:levelinfres}, \cl{prop:cycliccone},
  \cl{lem:levelfreeeuler}, \cl{lem:unipotentinv}, \cl{prop:gamma0normal}, \cl{prop:hidanormalsteps}; \cl{obl:levelconeexact}; remarks \cl{rmk:normalN1} (\PROVED{}),
  \cl{rmk:lescheckable} (\ref{LES1}--\ref{LES9}) and \cl{rmk:oldspacecone} (deferred). Its source is \file{docs/implementation/slot22/LES0\_spec.md} v0.1; a referee
  should read \cl{prop:levelconehecke}(1)--(3) (the convention of $A''=A\cap gAg^{-1}$; (3) is the level-raising $U_p$, \PROVED{} with negative content after the F-LES
  repair R1: the strict identity $U^{B}_p\circ\res=p\cdot\mathrm{dg}_g$ and the failing cell), \cl{prop:cycliccone}(3)(c) and \cl{prop:gamma0normal}'s local count first.
\item \emph{Part~VII, \S\ref{sec:mpmodel}, the Moerdijk--Pronk model} ($17$ assertions or definitions, $5$ remarks and $1$ example, v1.22.0-pre.5): \cl{def:mptri},
  \cl{def:mprules}, \cl{def:mpell}, \cl{def:mpRZ} (with \cl{tab:mprules}), \cl{def:mpset}, \cl{def:mpfree}, \cl{lem:mprows}, \cl{lem:mpcols}, \cl{lem:mpassembly},
  \cl{thm:mpcomparison}, \cl{def:mpphi}, \cl{prop:mpalldegrees}, \cl{def:mpmorphism}, \cl{thm:mpnatural}, \cl{prop:mpbredon}, \cl{prop:mpboundary}, \cl{prop:mprescor};
  remarks \cl{rmk:mprulesread}, \cl{rmk:mpwhat}, \cl{rmk:mphybrid}, \cl{rmk:mphecke} (commentary), \cl{rmk:mpcheckable} (\ref{MPC1}--\ref{MPC9}); example
  \cl{ex:mpteardrop}. Its sources are \file{docs/implementation/slot22/MP\_math\_plan.md} v0.2.1 (audited), \file{MP\_impl\_plan.md} v0.2.1, the F-MP2 code audit and the
  built code; a referee should read \cl{thm:mpcomparison}'s Step (2) with \cl{lem:mpcols} (the non-equivariant fibre contraction and the bounded assembly),
  \cl{def:mpmorphism}(3) with \cl{thm:mpnatural}(1),(4) (the label factor), and \cl{prop:mpalldegrees}(3),(5) (the torsion bound and the scope) first.
\item \emph{\S\ref{sec:six}, \S\ref{sec:oddnltwo}, odd $n$ at $\ell=2$ and the prime-to-$\ell$ Hecke action at $\ell\in\{2,3\}$} ($9$ assertions and $2$ remarks,
  v1.22.0-pre.10): \cl{lem:centralminusone}, \cl{lem:bocksteinallell}, \cl{lem:cyclicres}, \cl{prop:oddnltwo}, \cl{lem:oddnsplit}, \cl{lem:tellrank},
  \cl{prop:ltwothreescalar}, \cl{prop:oddnhecke} (its clause (4) \CONJ), \cl{lem:cprule}; remarks \cl{rmk:standardtwist} (\PROVED{}) and \cl{rmk:oddnltwoerr} (commentary). Its
  sources are \file{docs/implementation/slot22/L23\_master\_plan.md} v0.2 (Theorem~O) as repaired by its Fable review \file{L23\_master\_plan\_review.md}, and the census
  reviews \file{census\_l23\_classification\_review.md} (Theorems~A, B(1),(2), D(1),(2)) and \file{census\_l23\_plan\_v02\_review.md} \S1.2; a referee should read
  \cl{prop:ltwothreescalar}'s Step~0(c) (the level-$\Ga(\ell)$ operator, its twisted equivariance and the cusp value) and its $\ell=2$ $\mathsf{St}_2$ case, then
  \cl{prop:oddnltwo}(2) ($c_U=2h^{0}$) and \cl{lem:tellrank}'s orbit count, first.
\item \emph{Part~VII, \S\ref{sec:mpmodel}, the $Q$-reduced nerve} ($1$ assertion and $1$ definition, v1.22.0-pre.11): \cl{def:mpqd}, \cl{prop:mpqreduced}. Its source is
  the F-MPT referee \file{docs/implementation/slot22/mp\_triangulations\_referee.md} \S3, whose proof the editor checked and wrote out; a referee should read
  \cl{prop:mpqreduced}(2)'s column step (the cone on the dominant vertex, well defined on the $Q_x$-classes) and (1)'s invariance of $\mathsf p(\tau)$ on a class first.
\item \emph{\S\ref{sec:six}, the tree resolution over every ring} ($1$ remark, v1.22.0-pre.12): \cl{rmk:treeconeres} (\PROVED{}, three \code{\textbackslash LEAN}
  markers). Its sources are the LR-Y phase reports \file{docs/implementation/slot22/LRY\_phase1\_report.md}, \file{LRY\_phase2\_report.md}, \file{LRY\_phase3\_report.md} \S C
  and the Fable audits \file{LRY\_audit.md}, \file{LRY3\_audit.md}; a referee should read its orientation sentence against \S\ref{sec:levelN}'s tree $\mathcal Y$ and the
  periodicity indexing ($P_{n+2}=P_n$ for $n\ge1$, $\partial_{n+2}=\partial_n$ for $n\ge2$) first.
\item \emph{\S\ref{sec:oddnltwo}, the dimensions at $\ell\in\{2,3\}$ formalised} ($1$ remark, v1.22.0-pre.14): \cl{rmk:lvdimsltwothree} (\PROVED{}, nine
  \code{\textbackslash LEAN} markers). Its sources are \file{docs/implementation/slot22/LV1\_report.md}, \file{LV2\_report.md}, the Fable audits \file{LV1\_audit.md},
  \file{LV2\_audit.md} (\S8 sets what may be said) and \file{LVprov\_audit.md}; a referee should read its $\ell=3$ formula against Proposition~\ref{prop:ellthreeclosed}(2)
  and the sentence on the \texttt{L23Cells} cells (data, not markers) first.
\item \emph{Part~VII, \S\ref{sec:isomodelell}, the certified-exponent rule on $\Iso^{(\ell)}$} ($1$ assertion, v1.23.0-pre.12): \cl{prop:certifiedeell}. Its sources are
  \file{docs/implementation/slot24/DESCENT\_plan.md} v0.2 \S2.7 (row C-D0), \S2.1 and D0, and the audit finding M3 of \file{DESCENT\_plan\_audit.md}; a referee should read
  (1)'s replacement of ``$\Z_{(\ell)}$-algebra'' by ``$\mathsf q_\ell$ a unit'' in Proposition~\ref{prop:isomodelell}'s proof, and (5)'s reduction of the code's full layout to $d^{q}$, first.
  Refereed ACCEPT-with-fixes (\file{docs/implementation/slot24/C-D0\_referee.md}); the fixes of v1.23.0-pre.13 add (4)(a)--(c) and (6)'s a-priori bound, which a
  re-referee should read first.
\item \emph{Part~VII, \S\ref{sec:descentN}, the descent at level $N$ and its bad primes} ($2$ assertions, v1.23.0-pre.14): \cl{prop:descentN},
  \cl{prop:badprimesN}. Their sources are \file{docs/implementation/slot24/DESCENT\_plan.md} v0.3 \S2.7 (rows C-D1, C-D2), \S2.1, \S2.4 D-$\delta$4 and D1, the
  audit findings M8, m3, m4, m7 of \file{DESCENT\_plan\_audit.md}, and \file{DESCENT0\_report.md} \S1; a referee should read \cl{prop:badprimesN}(2)'s image of
  $\Ga_0(N)$ and Claim~A (the Borel invariants), (3)(b)'s twist, \cl{prop:descentN}(4) (the iff, and the values of $b_0$) and (7)(b) (generators, not a basis) first.
  Refereed ACCEPT-with-fixes (\file{docs/implementation/slot24/C-D1\_referee.md}); the fixes of v1.23.0-pre.15 add the proof of \cl{cl:dimH1}(2) and of $b_1\ge1$,
  which a re-referee should read first.
\item \emph{Part~VII, \S\ref{sec:descentN}, maps and the level cone over $\Z$} ($1$ assertion, v1.23.0-pre.18): \cl{prop:descentmaps}. Its sources are
  \file{docs/implementation/slot24/DESCENT\_plan.md} v0.3 \S2.2 (the rows ``Maps'' and ``Level-change LES''), \S2.7 (row C-D3) and D2, \S\ref{sec:nebles}, and the slot-22
  LES-0 request; a referee should read (2)(f) at $\ell=3$ (the $S[A]$-models), (3)(a)'s group computation ($A''$, $A'$, the indices, $p\nmid M$), (4)(c)'s vanishing of
  $V(\chi')^{B}_S$ and (5)'s sufficiency claim first. Refereed ACCEPT-with-fixes (\file{docs/implementation/slot24/C-D3\_referee.md}); the fixes of v1.23.0-pre.19
  add (3)(b$'$) (the square with $\cor$ for $p\mid N$, every $M$) and restrict (3)(b) to this document's $U_p$, which a re-referee should read first. v1.23.0-pre.22 adds
  (2)(g) (the precision rule, from the D2 code audit's F1) and (2)(b)'s signs; refereed ACCEPT-with-fixes (\file{docs/implementation/slot24/C-D3g\_referee.md}); the fixes of
  v1.23.0-pre.23 add (2)(g)(i)'s ``moreover'' (the blocks exact modulo $\ell^{K}$) with (ii)'s lattice coordinates, and (iii)'s list of certifying ranks, which a re-referee
  should read first.
\item \emph{Part~VII, \S\ref{sec:isomodeltree}, the tree models} ($1$ assertion, v1.23.0-pre.20): \cl{prop:isomodeltree} (with \cl{def:isomodeltree}). Its sources are
  \file{docs/implementation/slot24/DESCENT\_plan.md} v0.4 \S2.6, \S2.7 (row C-D6), gate G-T and D5, \cl{rmk:treeconeres}, \cl{prop:descentmaps}(5), and the evidence of
  \file{DESCENT1\_report.md}; a referee should read (1)(a) (well-definedness and the chain-map identities of the canonical lifts at $\ell=2$), (3)'s corrections $k$, $k'$ (above
  all $\kappa'_t$ and $\mathsf w$ at $\ell=2$) and the strict identity $\Phi^{\mathcal T}_\ell\Psi^{\mathcal T}_\ell=\id$, (1)(d)'s transfer of \cl{prop:certifiedeell}(2)--(6), and (4)'s
  ``over $\Z$, not over $C_{\bullet}(\Xbulk)$'' first. Refereed ACCEPT-with-fixes (\file{docs/implementation/slot24/C-D6\_referee.md}); the fixes of v1.23.0-pre.21 make (4)
  effective --- (4)(a)'s tables of $\mathsf r_{\mathcal T}$, $\mathsf G_{\mathcal T}$, (4)(b)'s solvers and the formulas (L), (L$'$), (4)(d)'s homotopies $\mathsf K$, $\mathsf K'$ ---, which a
  re-referee should read first.
\item \emph{Part~VII, \S\ref{sec:descenthecke}, Hecke $T_p$, $p\nmid N$, on the descent} ($5$ assertions, v1.23.0-pre.24): \cl{prop:strictliftN}, \cl{prop:crtcharpolyN},
  \cl{prop:torsionhecke}, \cl{thm:descentcert}, \cl{rmk:descentcertscope}. Their sources are \file{docs/implementation/slot24/DESCENT\_plan.md} v0.4 \S2.2 (the row ``Hecke
  action''), \S2.3, \S2.4 D-$\delta$1--D-$\delta$3, D-$\delta$7, \S2.7 (rows C-D4, C-D5) and D3, \S5 of \file{DESCENT\_plan\_audit.md} and its finding M5, and the owner's rulings
  [R-Q2]--[R-Q4], [R-GT]; a referee should read \cl{prop:strictliftN}(3)(c) (the identification with $[AgA]$ without $6$ inverted: the free source), (4) and (7) (the strict
  identities), \cl{prop:crtcharpolyN}(3) (the bound is an input; its citations are not all verified) and (5) (the output rule), \cl{prop:torsionhecke}(3) (the read needs
  $2\sigma_{\max}\le K$, stronger than the coordinate rule of \cl{prop:descentmaps}(2)(g)) and (4), and \cl{thm:descentcert}(3), (7) against the plan's M5 first. Refereed
  ACCEPT-with-fixes (\file{docs/implementation/slot24/C-D4\_referee.md}); the fixes of v1.23.0-pre.25 --- \cl{prop:torsionhecke}(1)(e), (2) (the layer blocks modulo
  $\ell^{e_j}$, its proof and sharpness), \cl{prop:strictliftN}(6)(b) with \cl{thm:descentcert}(3) (the one route at $\ell=3$) and the code map's obligations --- are what a
  re-referee should read first. v1.23.0-pre.31 (D3's contract notes: \cl{prop:torsionhecke}(3$'$), \cl{prop:strictliftN}(2)) refereed ACCEPT after fixes
  (\file{docs/implementation/slot24/C-D3h\_referee.md}); the fixes of v1.23.0-pre.32 --- the unsigned $\mathsf H=\Theta_{h_\Iso}$ of \cl{prop:strictliftN}(2), the precision of
  \cl{thm:descentcert}(4) at $\ell\ge5$, and the any-lifts argument of \cl{prop:torsionhecke}(3$'$) --- are what a re-referee should read first.
\item \emph{Part~VII, \S\ref{sec:descentpair}, the boundary pair and the cone over $\Z$} ($5$ assertions, v1.23.0-pre.26): \cl{def:descentpair}, \cl{lem:pairlifting},
  \cl{prop:descentpair}, \cl{prop:descentpairhecke}, \cl{rmk:descentpairscope}. Their sources are \file{docs/implementation/slot24/DESCENT\_plan.md} v0.4 \S2.2 (the row
  ``Boundary LES''), \S2.4 D-$\delta$6, \S2.7 (row C-D7), D4 and \S5.5, \file{DESCENT\_plan\_audit.md} on D-$\delta$6, and the owner's rulings [R-Q2]--[R-Q4]; a referee should
  read \cl{prop:descentpair}(1)(c) (the identification with $\Ext(\mathrm{St},-)$ on a pair model, and the integral torsion-freeness of $\Hcpt{1}$), (1)(d) (the boundary of
  $\Iso^{(\ell)}$ at $2$ and at $\ell\ge3$), (2) (the pair map $\cmpI$ as the only glueing datum), (4) (the closed form and its periodicity at $\varepsilon=\pm1$), (5)(b) (the rational
  $\Hcpt{2}$ by the Manin resolution and the adjoint computation) and (6) (the rank sources for $\Dcone^{0}$), and \cl{prop:descentpairhecke}(1) (the pair lift for every $p$, and
  its homotopy uniqueness), (3) (the identification of the cone operator with the double-coset operator on $\Ext(\mathrm{St},-)$ without $6$ inverted) and (5) (what is
  certified at which $p$) first. Refereed ACCEPT-with-fixes (\file{docs/implementation/slot24/C-D7\_referee.md}); the fixes of v1.23.0-pre.27 --- \cl{prop:descentpair}(5)(e) ($\Hcpt{2}$ at
  $\ell\ge5$ is $H_0$) with (6)(c), (1)(c)'s signs and their proof, the periodicity proofs of (4) and (5)(d), and \cl{prop:descentpairhecke}(5) (weakened) --- are what a re-referee
  should read first. v1.23.0-pre.28 (\emph{not refereed}) restates \cl{prop:descentpairhecke}(1), (5) and the code map (f) for every prime $p\nmid N$ (the pair lift built
  on demand, certified at run time); refereed ACCEPT-with-fixes (\file{docs/implementation/slot24/C-D7p\_referee.md}); the fixes of v1.23.0-pre.29 --- the refusal
  list of \cl{prop:descentpairhecke}(1), (5) and the code map's (f) (the code's eleven refusal codes) and the forced-component-map paragraph of the proof of (1) --- are what a re-referee
  should read first.
\item \emph{Part~VII, \S\ref{sec:treepair}, the tree cusp pair} ($4$ assertions, v1.23.0-pre.33, \textbf{not refereed}): \cl{def:treepair}, \cl{lem:treepaircyl}, \cl{prop:treepair},
  \cl{rmk:treepairscope}. Their sources are \file{docs/implementation/slot24/DESCENT\_plan.md} v0.4 \S5 (the owner's candidate, the cusp model, \S5.3's costs, \S5.4(c)), \S2.7 (row C-D8) and D5c,
  and the evidence of \file{DESCENT4\_report.md} \S4 and \file{DESCENT5\_report.md}; a referee should read \cl{def:treepair}(2) with \cl{prop:treepair}(1)(b) (the axis map: well defined and a chain
  map at $\ell=2$ in odd degrees), \cl{lem:treepaircyl}(2), (3) (the deformation of the cylinder's cone onto the tree cone and the signs), \cl{prop:treepair}(3)(c), (d) (the two squares by (L$'$), and
  $H(\mathsf j_\Psi)=H(\mathsf j_\Phi)^{-1}$ through the comparisons of pairs $\Theta_\Phi$, $\Theta_\Psi$, not the five lemma) and (4)(a) (the off-diagonal $\eta_{\mathcal T}$, which corrects the plan) first.
  Refereed ACCEPT after fixes (\file{docs/implementation/slot24/C-D8\_referee.md}, 0 M / 4 m / 8 N; independent check \file{cd8referee/}); the fixes of v1.23.0-pre.34 --- \cl{prop:treepair}(4)(a) and
  \cl{rmk:treepairscope}(iii) (where $\eta_{\mathcal T}$ arises: on $\Cone{\mathsf a^{\ast}}$, not on the pair model; its non-vanishing observed, not proved), the generalised solver of \cl{prop:treepair}(3)(b), the
  $\ell$-local column augmentations of \cl{def:treepair}(4), and the code map's restricted negative controls (f) --- are what a re-referee should read first.
  v1.23.0-pre.35 (\emph{not refereed}): \cl{def:treepair}(2)'s cochain statement ($\rho_{M_N}(\mathbf c)=(-1)^n$), the new \cl{prop:treepair}(4)(e) (truncation) with its proof, and the code map's
  certificates paragraph (what ``on the generators'' certifies) are what a referee should read first. Refereed ACCEPT after fixes (\file{docs/implementation/slot24/C-D8b\_referee.md}, 0 M / 1 m / 7 N; independent
  check \file{pre35referee/}); the fixes of v1.23.0-pre.36 --- the code map's certificates (the applied operator's blocks; deterministic and probabilistic parts; admissibility) and \cl{prop:treepair}(4)(e)'s
  hypothesis on $T$, (ii) and (iv) --- are what a re-referee should read first.
\item \emph{Part~VII, \S\ref{sec:heckeZ}, the Hecke algebra over $\Z$} ($17$ assertions, $2$ definitions and $4$ remarks, three of them \CONJ; v1.25.0-pre.0, a branch
  pre-release, \textbf{not refereed}): \cl{def:heckeZ}, \cl{prop:heckeZspec}, \cl{prop:heckeZrat}, \cl{lem:heckeZbridge}, \cl{prop:heckeZmax}, \cl{prop:heckeZgen},
  \cl{prop:heckeZetale}, \cl{prop:heckeZmono}, \cl{prop:heckeZlocal}, \cl{prop:heckeZthrough}, \cl{prop:heckeZuct}, \cl{prop:heckeZsplit}, \cl{prop:heckeZlevelone23},
  \cl{prop:heckeZglue}, \cl{thm:heckeZdescent}, \cl{thm:heckeZcert}, \cl{def:heckeZcong}, \cl{prop:heckeZcong}, \cl{prop:heckeZeiscert}; remarks \cl{rmk:heckeZcbone},
  \cl{rmk:heckeZcbtwo}, \cl{rmk:heckeZcbthree} (\CONJ) and \cl{rmk:heckeZcaveat} (\PROVED). Their sources are the HeckeZ plan \file{docs/implementation/slot24/heckez/}
  (signed by the owner 2026-11-04): \file{01-math-core.md} v0.6 \S\S1--5, \file{03-bockstein-tor.md} v0.3 \S\S2--4, \file{04-record-census.md} v0.4 \S1 (item~8), and the audits
  \file{A1-audit.md}, \file{A3-audit.md}, \file{P2a-final-math-audit.md}, \file{P3-reaudit.md}, whose findings the stage files had applied; a referee should read
  \cl{thm:heckeZcert}(C1)'s dimension count and (C2)'s set $\mathsf S$ with the cofactor, \cl{prop:heckeZgen}(1) (the a-priori $d$), \cl{prop:heckeZmono}'s computable form
  (the test at $K\ge c^{\mathrm{sat}}+1$ on an unsaturated lattice), \cl{prop:heckeZetale}(2)'s $v_\ell=1$ clause, \cl{prop:heckeZglue}(3) (the sandwich; $J$ never exact) and
  \cl{prop:heckeZeiscert}(3) (Mazur's good prime and the index count) first. Refereed: A7 REVISE $\to$ A7b ACCEPT
  (\file{docs/implementation/slot24/heckez/A7-Z0-referee.md}, \file{A7b-Z0-reaudit.md}); the main-line merge v1.25.0-pre.1 refereed REVISE
  (\file{docs/implementation/slot25/Z0\_referee.md}), its findings applied in v1.25.0-pre.2. The reconciliation notes dated 2026-11-06 (R1--R9 of
  \file{docs/implementation/slot25/Z0\_merge\_report.md}: item~8~(c), (d), (f), (j) and the code map; \S\ref{sec:heckeZ}'s two \emph{As built} paragraphs, the cofactor route as built,
  the counting criterion's scope) are what a re-referee should read first.
\end{itemize}

\subsection*{v1.26.4 --- 2026-11-10 --- hecke-extensions slot 26: the owner-signed release of v1.26.4-pre.1}

\noindent\emph{This version.} \textbf{Signed by the owner 2026-11-10} (``I sign those title pages.''): the owner-signed release of v1.26.4-pre.1; no change of content. v1.26.4-pre.1 (a pre-release on top of v1.26.3) changes the title page only. \textbf{No statement's content changes; no label is added, renamed or removed; no status tag moves; the marker count stays $229$ on $88$ labels.}
\begin{itemize}\setlength{\itemsep}{0pt}
\item Header and \verb|\date| (the version only; the title page names no sign-off); the Release record's first entry; this entry; the companion pointer of the v1.26.4-pre.1 entry.
\end{itemize}
Companion pointers: [E] 1.25.3, [U] 0.7.3, [X] v1.28; [C], [E] and [U] are signed together.

\subsection*{v1.26.4-pre.1 --- 2026-11-10 --- hecke-extensions slot 26: title page only}

\noindent\emph{This version.} A \textbf{pre-release, not signed}: title page only, no change of mathematical content (public release HACS 0.1.1). \textbf{No statement's content changes; no label is added, renamed or removed; no status tag moves; the marker count stays $229$ on $88$ labels.}
\begin{itemize}\setlength{\itemsep}{0pt}
\item Header and \verb|\date|: the date line of the title page names only the version, the date and this appendix; the words ``binding document; signed by the owner 2026-11-09'' leave the title page. The sign-off record of every version stays in the Release record and in this changelog.
\item The Release record's first entry; this entry.
\end{itemize}
Companion pointers: [E] 1.25.3, [U] 0.7.3 (the same title-page change; signed with v1.26.4), [X] v1.28.

\subsection*{v1.26.3 --- 2026-11-09 --- hecke-extensions slot 26: the owner-signed release of v1.26.3-pre.1}

\noindent\emph{This version.} \textbf{Signed by the owner 2026-11-09} (the owner's pre-signature, effective on the referee's ACCEPT, \file{RELFIX\_signed\_referee.md}): the owner-signed release of v1.26.3-pre.1; no change of mathematical content. v1.26.3-pre.1 (a pre-release on top of v1.26.2) is text-only: a private path removed. \textbf{No statement's content changes; no label is added, renamed or removed; no status tag moves; the marker count stays $229$ on $88$ labels.}
\begin{itemize}\setlength{\itemsep}{0pt}
\item Header and \verb|\date|; the Release record's first entry; this entry; the companion pointer of the v1.26.3-pre.1 entry.
\end{itemize}
Companion pointers: [E] 1.25.2, [U] 0.7.2, [X] v1.28; [C] and [U] are signed together.

\subsection*{v1.26.3-pre.1 --- 2026-11-09 --- hecke-extensions slot 26: text only, a private path removed}

\noindent\emph{This version.} A \textbf{pre-release, not signed}: text-only: private path removed, no change of mathematical content (release audit of 2026-11-09, finding S1). \textbf{No statement's content changes; no label is added, renamed or removed; no status tag moves; the marker count stays $229$ on $88$ labels.}
\begin{itemize}\setlength{\itemsep}{0pt}
\item Header and \verb|\date|; the Release record's first entry; this entry.
\item The \emph{Corroboration} paragraph after the proof of Proposition~\ref{prop:bdrytranscription} names its ad hoc script by its file name \file{prop\_t\_check.py} instead of a path in the author's scratch directory.
\item The two \COMP\textsuperscript{a} data remarks of v1.14.1 (the second-cluster run at $n=4150$ and at $(4830,2411)$) cite ``the WP-C10 report of the September run, \S1'' instead of that report's file name, which carried the cluster's name.
\end{itemize}
Companion pointers: [E] 1.25.2, [U] 0.7.2 (the same scrub; signed with v1.26.3), [X] v1.28.

\subsection*{v1.26.2 --- 2026-11-09 --- hecke-extensions slot 26: the owner-signed release of v1.26.2-pre.1}

\noindent\emph{This version.} \textbf{Signed by the owner 2026-11-09}: the owner-signed release of v1.26.2-pre.1; no change of content. v1.26.2-pre.1 (a pre-release on top of v1.26.1) is a paths-only update for the repository layout (referee \file{LAYOUT\_signed\_referee.md}: ACCEPT). \textbf{No statement's content changes; no label is added, renamed or removed; no status tag moves; the marker count stays $229$ on $88$ labels.}
\begin{itemize}\setlength{\itemsep}{0pt}
\item Header and \verb|\date|; the Release record's first entry; this entry.
\end{itemize}
Companion pointers: [E] 1.25.2, [U] 0.7.1, [X] v1.28; [C] and [E] are signed together.

\subsection*{v1.26.2-pre.1 --- 2026-11-09 --- hecke-extensions slot 26: paths only}

\noindent\emph{This version.} A \textbf{pre-release, not signed}: paths only, no change of content (repository layout, owner ruling 2026-11-09). The program's sources moved under \file{src/} (\file{src/hacs/}, \file{src/ModArithmetic/}, \file{src/Engine/}, \file{src/Cuda/}; the old-to-new map of every file is \file{tools/layout/mapping.tsv}). \textbf{No statement's content changes; no label is added, renamed or removed; no status tag moves; the marker count stays $229$ on $88$ labels.}
\begin{itemize}\setlength{\itemsep}{0pt}
\item Header and \verb|\date|; the Release record's first entry; this entry.
\item Twelve lines name a file by its new path: eleven long-form paths (\file{src/ModArithmetic/\dots}, \file{src/hacs/\dots}) in \cl{sec:codemap}'s table and the paragraphs citing \file{ENGINE\_API.md}, \file{DenseElim}, \file{FlatMatrix}, \file{Binary/Ring2k.h}, \file{Orbifold/MP/} and \file{HeckeIotaSignValidation}; and the 2026-09-16 grep census of the cone chapter, which now names \file{src/} (the same files).
\end{itemize}
Companion pointers: [E] 1.25.2 (paths only; signed with v1.26.2), [U] 0.7.1, [X] v1.28 (keeps its paths: history).

\subsection*{v1.26.1 --- 2026-11-09 --- hecke-extensions slot 26: the owner-signed release of v1.26.1-pre.1}

\noindent\emph{This version.} \textbf{Signed by the owner 2026-11-09}: the owner-signed release of v1.26.1-pre.1; no change of content. v1.26.1-pre.1 (a pre-release on top of v1.26.0) changed no mathematical content: machine names and private paths removed for the public release (owner ruling 2026-11-09); machine and cluster names were replaced by neutral words and private directory paths by prose or placeholders, in the text and in source comments. \textbf{No statement's content changes; no label is added, renamed or removed; no status tag moves; the marker count stays $229$ on $88$ labels.}
\begin{itemize}\setlength{\itemsep}{0pt}
\item Header and \verb|\date|; the Release record's first entry; this entry.
\end{itemize}
Companion pointers: [E] 1.25.1, [U] 0.7.1, [X] v1.28; all three binding documents are signed together.

\subsection*{v1.26.0 --- 2026-11-09 --- hecke-extensions slot 26: the owner-signed release of v1.26.0-pre.11}

\noindent\emph{This version.} \textbf{Signed by the owner 2026-11-09}: the owner-signed release of v1.26.0-pre.11; no change of content. It releases the content of v1.26.0-pre.1 to pre.11 as described in their entries below, including the binding changes there marked ``not yet seen by the owner'' (item~8~(c), (d), (d$'$) of \S\ref{sec:recordschema} and \cl{prop:strictliftN}(9)); the owner signs them with this release. \textbf{No statement's content changes; no label is added, renamed or removed; no status tag moves; the marker count stays $229$ on $88$ labels.} By the test of item~8~(d$'$)(vii) a signature as such moves no written word: the values \code{gl2\_contract 1.26.0-pre.5, item 8 (d')} and \code{gl2\_contract 1.26.0-pre.9, item 8 (d')} of \code{G10.contract} stay, and neither becomes \code{1.26.0}.
\begin{itemize}\setlength{\itemsep}{0pt}
\item Header and \verb|\date|; the Release record's first entry; this entry.
\end{itemize}
Companion pointers: [E] 1.25.0, [U] 0.7.0, [X] v1.28; [E] and [U] are not re-signed.

\subsection*{v1.26.0-pre.11 --- 2026-11-09 --- hecke-extensions slot 26, UP-Z (b): the conditions of the final confirm-referee of v1.26.0-pre.10}
\textbf{This version} is a pre-release, \textbf{not owner-released}, on top of v1.26.0-pre.10. It discharges the condition C1 and the should-fix items W-A, W-B of the confirm-referee \file{docs/implementation/slot26/C126pre10\_confirm\_referee.md} (ACCEPT-with-conditions); disposition \file{docs/implementation/slot26/C126pre11\_disposition.md}. Text only. Its binding changes are \textbf{not yet seen by the owner}. It states nothing of Z2 as done: each edited sentence holds before and after Z2's merge, or is scoped to the binary that wrote the record.
\paragraph{Normative.}
\begin{itemize}\setlength{\itemsep}{0pt}
\item Item~8~(d$'$)(ii), \emph{Why \code{lower} is \code{not\_run}} (W-A): where no tag certifies there is no row; \code{G3.C\_T.lower} is \code{not\_run} there too, on the \code{open(cap)} records, outside the scope of \cl{def:heckeZcong}, and also where $\eta_L$ is defined and \cl{prop:heckeZcong}(5) could read it with a $T_p$ ($(11,0)$), which a \code{full} record of this version does not do. A future row: $X_p=T_p$ where a $T_p$ separates, $U_p$ where only a $U_p$ does. (W-B): ``the partition \code{U\_p\_tags} has no field''. The word \code{not\_run} and its final class do not move.
\item Item~8~(d$'$)(vii): v1.26.0-pre.11 moves no written word by the test of that item, so the values of \code{G10.contract} stay \code{gl2\_contract 1.26.0-pre.5, item 8 (d')} and \code{gl2\_contract 1.26.0-pre.9, item 8 (d')}.
\end{itemize}
\paragraph{Status sentences (C1; the edits promised in the entry of v1.26.0-pre.1, audit F7).} \cl{thm:heckeZcert}, ``what is reachable today'': the list dated (v1.25.0-pre.0, before EXPORT and UP-Z); the \code{full} line at $N>1$ no longer says ``nothing is certified but two tags'' ((C1) is certified without $U_p$ over $\Z$ on $67$ of Z1's $85$ records), and points to item~8~(d$'$)(ii), (v). \cl{thm:descentcert}(5), \cl{prop:descentpairhecke} (``$U_p$ is deferred''): nothing is asserted there of $U_p$, with a pointer to \cl{prop:strictliftN}(9). \cl{def:heckeZ} and the ``Not here'' of \S\ref{sec:heckeZ} (``not built''): the lift of \cl{prop:strictliftN}(9), served by a binary or not according to (d$'$)(v). \cl{prop:heckeZcong}(5): the read needs $U_p$ over $\Z$ and no route of this version performs it. ``What is not available at $2$ and $3$ today'': dated. \cl{prop:descentN}(5), ``its cochain operator is deferred'': lifted for the free model by \cl{prop:strictliftN}(9). The ``Not here'' of \S\ref{sec:heckeZ}, ``What is not available'' and \cl{prop:descentN}(5) were found by a search for the same phrases, beyond the referee's list.
\paragraph{Unchanged.} No written word or hash input moves; no label is added, renamed or removed; no \verb|\LEAN| marker ($229$ on $88$) moves; no golden entry, level-$1$ or \code{anemic} byte, settings text or fingerprint moves; no code change and no rerun.

\subsection*{v1.26.0-pre.10 --- 2026-11-08 --- hecke-extensions slot 26, UP-Z (b): the condition E1 of the referee of v1.26.0-pre.8 / pre.9}
\textbf{This version} is a pre-release, \textbf{not owner-released}, on top of v1.26.0-pre.9. It discharges the condition E1 of the referee \file{docs/implementation/slot26/C126pre9\_referee.md} (ACCEPT-with-conditions), its should-fix W-1 and nits N-1--N-3; disposition \file{docs/implementation/slot26/C126pre10\_disposition.md}. Text only. Its binding changes are \textbf{not yet seen by the owner}.
\paragraph{Normative.}
\begin{itemize}\setlength{\itemsep}{0pt}
\item Item~8~(d$'$)(ii), \emph{The G3 words from Z2 on} (E1, N-2): the reason for \code{G3.C\_T.lower} $=$ \code{not\_run} is corrected. For $d\ge3$, $\eta_L$ of a pair is undefined; for $d=2$ it is defined and no route of this version performs the read of Proposition~\ref{prop:heckeZcong}(5) on a \code{U\_p\_tags} row; on the non-reduced branch, the scope of Definition~\ref{def:heckeZcong}. ``No statement defines a lower bound'' becomes ``no statement prescribes a read of a lower bound''; the same in the entry of v1.26.0-pre.9 below (\code{not\_run} where no bound, resp.\ no lower bound, is read on the pairs). The word \code{not\_run} and its final class do not move.
\item Item~8~(d$'$)(ii) (W-1): those words presuppose that the $U_p$ tag step ran at every $p\mid N$; where it did not (no record of this version), \code{G3.C\_T[U\_p].status} of that $p$, \code{G3.C\_T.status} and \code{G3.C\_T.lower} keep \code{blocked(U\_p)}.
\item Item~8~(d$'$)(ii) (N-1, N-3): the retained rule's ``until Z2 lands'' points at the per-record sense; a binary carrying the route writes these words, which ties (ii) to (vii).
\item Item~8~(d$'$)(vii): v1.26.0-pre.10 moves no written word by the test of that item, so the values of \code{G10.contract} stay \code{gl2\_contract 1.26.0-pre.5, item 8 (d')} and \code{gl2\_contract 1.26.0-pre.9, item 8 (d')}.
\end{itemize}
\paragraph{Unchanged.} No written word or hash input moves; no label is added, renamed or removed; no \verb|\LEAN| marker ($229$ on $88$) moves; no golden entry, level-$1$ or \code{anemic} byte, settings text or fingerprint moves; no code change and no rerun.

\subsection*{v1.26.0-pre.9 --- 2026-11-08 --- hecke-extensions slot 26, UP-Z (b): Z2's G3 words, contract first}
\textbf{This version} is a pre-release, \textbf{not owner-released, not refereed}, on top of v1.26.0-pre.8; disposition \file{docs/implementation/slot26/C126pre9\_disposition.md}. Text only. Its binding changes are \textbf{not yet seen by the owner}.
\paragraph{Normative.}
\begin{itemize}\setlength{\itemsep}{0pt}
\item Item~8~(d$'$)(ii), new paragraph \emph{The G3 words from Z2 on}: \code{G3.C\_T[U\_p].status}, \code{G3.C\_T.status}, \code{G3.C\_T.lower} (final; \code{not\_run} where no bound, resp.\ no lower bound, is read on the \code{U\_p\_tags} pairs; this wording corrected at v1.26.0-pre.10), \code{open(sandwich)} at $\ell\in\{2,3\}$ on an \code{open(cap)} record, and amendment A3 (\code{not\_run} in the closed list \code{utag}). It names the words left open at v1.26.0-pre.4. The rules before Z2 stay in force for a binary without Z2's route.
\item Item~8~(d$'$)(v): pointer to that paragraph.
\item Item~8~(d$'$)(vii): the value of \code{G10.contract} is \code{gl2\_contract 1.26.0-pre.9, item 8 (d')} for a binary that writes those words, and stays \code{gl2\_contract 1.26.0-pre.5, item 8 (d')} for one that does not.
\end{itemize}
\paragraph{Unchanged.} No label is added, renamed or removed; no \verb|\LEAN| marker ($229$ on $88$) moves; no golden entry, level-$1$ or \code{anemic} byte, settings text or fingerprint moves; records already written keep their value of \code{G10.contract} and stay valid, without a code change or a rerun.

\subsection*{v1.26.0-pre.8 --- 2026-11-08 --- hecke-extensions slot 26, UP-Z: the conditions of the confirm-referee of v1.26.0-pre.7}
\textbf{This version} is a pre-release, \textbf{not owner-released}, on top of v1.26.0-pre.7. It discharges the conditions D1, D2 of the
confirm-referee \file{docs/implementation/slot26/C126pre7\_confirm\_referee.md} (ACCEPT-with-conditions); disposition \file{docs/implementation/slot26/C126pre8\_disposition.md}. Text only. Its binding changes are \textbf{not yet seen by the owner}.
\paragraph{Normative.}
\begin{itemize}\setlength{\itemsep}{0pt}
\item Item~8~(c) (D1): ``one reading apart'' becomes two: beside the derived \code{maximal\_ideals = not\_run} of (d$'$)(ii), the aggregate \code{G3.C\_T.status} while a \code{C\_T} read or \code{G3.C\_T.lower} it depends on is pending; it is rewritten when those move. No byte of a final record moves.
\item Item~8~(d$'$)(vii) (D2): naming for the first time the word of a field that (d$'$) had left open moves a written word for a binary that writes it; the manifest asserts that no later version moved a written word \emph{of a record carrying that value}, so the assertion stays true after a partial regeneration.
\end{itemize}
\paragraph{Unchanged.} No label is added, renamed or removed; no \verb|\LEAN| marker ($229$ on $88$) moves; no record word, hash input (\code{record\_fnv1a64}, \code{cert\_sha256}), settings text or
fingerprint, golden entry or level-$1$ / \code{anemic} byte moves; the value \code{gl2\_contract 1.26.0-pre.5, item 8 (d')} that the code writes stays correct, without a code change or a rerun.

\subsection*{v1.26.0-pre.7 --- 2026-11-08 --- hecke-extensions slot 26, UP-Z: the confirm-referee's conditions}
\textbf{This version} is a pre-release, \textbf{not owner-released}, on top of v1.26.0-pre.6. It discharges the conditions C1--C3, the should-fix items W-A, W-B, W-D and the nits N-a, N-b of the
confirm-referee \file{docs/implementation/slot26/C126pre6\_confirm\_referee.md} (ACCEPT-with-conditions); disposition \file{docs/implementation/slot26/C126pre7\_disposition.md}. Text only. Its binding changes are \textbf{not yet seen by the owner}.
\paragraph{Normative.}
\begin{itemize}\setlength{\itemsep}{0pt}
\item Item~8~(d$'$)(ii) (C1): ``not golden-eligible after Z2 alone'' is stated only of a record carrying a field that waits on \code{HZ-EISUP} (the records with (C1) \code{certified}); the other $18$ of Z1's $85$ depend on the words Z2 fixes.
\item Item~8~(d$'$)(ii), (c), (d) (C3): \code{maximal\_ideals} $=$ \code{not\_run} beside a pending \code{generation.status} is a derived word on any record, not only on a \code{full} record at $N>1$; no byte moves.
\item Item~8~(d$'$)(vii) (C2): \code{<version>} is a decidable test (a change of the prescribed word or value of a written key or field other than \code{G10.contract}, or of the key's presence; a reason, a class, a pointer or a conflict resolved in favour of an already prescribed word moves nothing), ``the code'' is the binary that wrote the record, and the manifest applies the test.
\item Item~8~(c) (W-A): \code{HZ-EISUP} named as the mover of \code{G1.components} and of the \code{maximal\_ideals} of (d$'$)(ii).
\item Item~8~(d$'$)(ii) (W-D): \code{HZ-EISUP} has its card in plan v0.6 (\S2.2), sequenced right after Z2 by the owner's ruling of 2026-11-08; Z2 stays necessary and not sufficient.
\item Item~8~(d$'$)(ix) (W-B): ``exactly the status fields''; the non-status moves of \code{G0.operators} and \code{G1.S\_cand.ell=<}$\ell$\code{>.reason}, \code{.v\_disc} are named.
\end{itemize}
\paragraph{Editorial.} (d$'$)(vii) ``A written value of this key'' (N-a); (d$'$)(ii) the antecedent of ``pending in any case'' (N-b, also required by C3).
\paragraph{Unchanged.} No label is added, renamed or removed; no \verb|\LEAN| marker ($229$ on $88$) moves; no record word, hash input (\code{record\_fnv1a64}, \code{cert\_sha256}), settings text or
fingerprint, golden entry or level-$1$ / \code{anemic} byte moves; the value \code{gl2\_contract 1.26.0-pre.5, item 8 (d')} that the code writes stays correct, without a code change or a rerun.

\subsection*{v1.26.0-pre.6 --- 2026-11-08 --- hecke-extensions slot 26, UP-Z: the referee-light's conditions}
\textbf{This version} is a pre-release, \textbf{not owner-released}, on top of v1.26.0-pre.5. It discharges the conditions R1--R3, the should-fix items W1, W2 and the nits N1--N3 of the
referee-light \file{docs/implementation/slot26/C126pre5\_referee\_light.md} (ACCEPT-with-conditions); disposition \file{docs/implementation/slot26/C126pre6\_disposition.md}. Text only. Its binding changes are \textbf{not yet seen by the owner}.
\paragraph{Normative.}
\begin{itemize}\setlength{\itemsep}{0pt}
\item Item~8~(d$'$)(ii) (R1), with pointers in (c), (d), (d$'$)(v), (d$'$)(ix): on a \code{full} record at $N>1$, \code{maximal\_ideals} $=$ \code{not\_run} beside a pending \code{generation.status}
  is a derived word, rewritten when \code{generation.status} moves; \code{not\_run} beside a final \code{generation.status} remains final. (ix) covers edition~1's \code{not\_run} $\to$ \code{blocked(U\_p)}.
\item Item~8~(d$'$)(ii) (W1): the Eisenstein $U_p$ systems are not built by Z2, and no card of plan v0.5 builds them; working name \code{HZ-EISUP}. Z2 is necessary and not sufficient.
\item Item~8~(d$'$)(ii) (W2): \code{G1.components} $=$ \code{not\_run} at (C1) \code{open(cap)} stated explicitly (final; the first normative statement of a word the code already wrote).
\item Item~8~(d$'$)(v) (R2): \code{G9.Up.*} is written by the package that serves the $U_p$ matrix over $\Z$ inside the record build; no record of Z1 carries it; a dump run does not create it.
\item Item~8~(d$'$)(vii) (R3): \code{<version>} is the last version that moved a written word (at this version \code{1.26.0-pre.5}); a written value is never rewritten; the manifest states the
  issuing version and the \code{G10.contract} value(s), and asserts that no written word moved against them.
\end{itemize}
\paragraph{Editorial.} (d$'$)(ii) ``the other four'' (N1); one phrase for the source of v1.26.0-pre.5 (N2); the referee-light check recorded where it stood pending (N3).
\paragraph{Unchanged.} No label is added, renamed or removed; no \verb|\LEAN| marker ($229$ on $88$) moves; no record word, hash input (\code{record\_fnv1a64}, \code{cert\_sha256}), settings text or
fingerprint, golden entry or level-$1$ / \code{anemic} byte moves; the value \code{gl2\_contract 1.26.0-pre.5, item 8 (d')} that the code writes stays correct, without a code change or a rerun.

\subsection*{v1.26.0-pre.5 --- 2026-11-08 --- hecke-extensions slot 26, UP-Z: the text aligned with the audited Z1 code}
\textbf{This version} is a pre-release, \textbf{not owner-released}, on top of v1.26.0-pre.4. It takes should-fix S1 (and S2's key) of the Z1 code audit
\file{docs/implementation/slot26/Z1\_code\_audit.md} (PASS-with-fixes) on the contract side: the text moves, not the code (disposition
\file{docs/implementation/slot26/C126pre5\_disposition.md}; a referee-light check was pending at this version, since done: \file{docs/implementation/slot26/C126pre5\_referee\_light.md}, discharged at v1.26.0-pre.6). Its binding changes are \textbf{not yet seen by the owner}.
\paragraph{Normative.}
\begin{itemize}\setlength{\itemsep}{0pt}
\item Item~8~(d$'$)(ii) (S1(a)): where the search ran and (C1) is \code{certified}, \code{G1.components} $=$ \code{blocked(U\_p)}, and \code{G2.ell=<}$\ell$\code{>.maximal\_ideals}
  $=$ \code{blocked(U\_p)} where \code{generation.status} is \code{certified} --- pending, until the Eisenstein $U_p$ systems are built (Z2 or a named successor); replaces
  pre.4's inclusion of \code{G1.components} in ``\code{blocked(U\_p)} is no longer written in \dots''. Elsewhere the existing \code{not\_run} words of (d$'$)(ii) and~(d). No
  final word is introduced.
\item Item~8~(d$'$)(v) (S1(b)): \code{G9.Up.p=<p>.mat\_sha256} unhashed (outside \code{record\_fnv1a64}; inside \code{cert\_sha256} as every G9 key), present only on \code{full}
  records at $N>1$ where the $U_p$ matrix over $\Z$ was computed, absent elsewhere.
\item Item~8~(d$'$)(vii) (S2): the key is exactly \code{G10.contract} with the value \code{gl2\_contract <version>, item 8 (d')}, \code{<version>} the implemented version.
\end{itemize}
\paragraph{Unchanged.} No label is added, renamed or removed; no \verb|\LEAN| marker ($229$ on $88$) moves; only words written under \code{full\_N}; no \code{record\_fnv1a64},
no golden line and no byte of a level-$1$ or \code{anemic} record moves.

\subsection*{v1.26.0-pre.4 --- 2026-11-08 --- hecke-extensions slot 26, UP-Z: the confirm-referee's should-fix items}
\textbf{This version} is a pre-release, \textbf{not owner-released}, on top of v1.26.0-pre.3. It takes the contract-side should-fix items W3, W4, W5 and the note N1 of the
confirm-referee \file{docs/implementation/slot26/C126pre3\_confirm.md} (ACCEPT: Z1 may merge on the contract side; disposition
\file{docs/implementation/slot26/C126pre4\_disposition.md}); W6 and N2--N4 are Z1 / K3 items. Its binding changes are \textbf{not yet seen by the owner}.
\paragraph{Normative.}
\begin{itemize}\setlength{\itemsep}{0pt}
\item Item~8~(d$'$)(ii), the rule ``no full record at $N>1$ is final before Z2'' (W3): the fields enumerated, exhaustively --- \code{G3.C\_T.lower} $=$ \code{blocked(U\_p)};
  \code{G3.C\_T.status} $=$ \code{upper\_bound\_certified} where every $U_p$ tag certifies a bound, else \code{blocked(U\_p)}, sharpened to \code{certified} only together with
  \code{.lower}, by Z2; \code{G3.C\_T.exact} absent; the per-$\ell$ words at $\ell\in\{2,3\}$ by~(v). Replaces ``every other aggregate field whose value may need $U_p$''.
\item Item~8~(d$'$)(ii), the ``otherwise'' branch (W4): \code{G1.components} is the fifth $\ell$-independent field kept \code{blocked(U\_p)} / \code{blocked(export)} where the
  word-prime $U_p$ route did not run.
\item Item~8~(d$'$)(ii), G3 on the non-reduced branch (N1): \code{G3.C\_T.status} $=$ \code{blocked(U\_p)} there (the tag's bound needs (C1) for $\vartheta=U_p$); the final words of
  \code{G3.C\_T.status} and \code{.lower} on that branch are \textbf{left open for Z2} (\cl{prop:heckeZcong}(1) uses no generation, so they do not follow from the present text),
  and Z2's contract text must name them before any such record is final.
\item Item~8~(d), G1 (W5): the reading sentence of \code{partial(unfactored)} beside a non-certified (C1) --- ``the words the \code{anemic} records with (C1) \code{open(cap)}
  already carry'' corrected: no record of edition~1 reaches that code path (every \code{anemic} record of edition~1 has (C1) \code{certified}).
\end{itemize}
\paragraph{Non-normative.} The status note of (d$'$) (confirm-refereed); the header comment, \verb|\date|, the Release record's first entry; this entry. No label is added, renamed
or removed; no status tag or marker moves ($229$ on $88$); no record, golden line, fingerprint (\code{0x8ffc29730c4f9672}) or digest moves; every changed word is written only
under \code{full\_N} ($\code{algebra = full}$, $N>1$), and no record of edition~1 carries a word this version fixes that it did not carry before.

\subsection*{v1.26.0-pre.3 --- 2026-11-08 --- hecke-extensions slot 26, UP-Z: the re-referee's conditions}
\textbf{This version} is a pre-release, \textbf{not owner-released}, on top of v1.26.0-pre.2. It discharges the conditions R1, R2 and the should-fix items W1, W2 of the adversarial
re-referee \file{docs/implementation/slot26/C126pre2\_rereferee.md} (ACCEPT-with-conditions; disposition \file{docs/implementation/slot26/C126pre3\_disposition.md}). Its binding changes
are \textbf{not yet seen by the owner}.
\paragraph{Normative.}
\begin{itemize}\setlength{\itemsep}{0pt}
\item Item~8~(d$'$)(ii), the non-reduced branch (W1, R1(a); lead decision): \code{G1.C1.status = refused(non\_reduced)} (final; replaces pre.2's \code{open(cap)}),
  \code{G1.theta = not\_run} (replaces \code{open(cap)}), \code{G1.C2\_status = not\_run} (replaces \code{partial(unfactored)}); G2 \code{generation.status = open(sandwich)} at
  $\ell\in\mathsf S$, $\ell\ge5$, and \code{blocked(U\_p)} at $\ell\in\mathsf S\cap\{2,3\}$ until Z2 (then \code{open(sandwich)}; \code{refused(no\_lift)} on Z2-D's audited negative).
  Amendment A2 after the freeze (hash-neutral; no record carries the words yet); the verifier's \code{FINAL} set gains \code{refused(non\_reduced)}.
\item Item~8~(d$'$)(ii), the ``otherwise'' branch (R2): the final \code{C1} words only where the $\vartheta$ search ran with $U_p$ at the word primes; else the four
  $\ell$-independent fields keep \code{blocked(U\_p)} / \code{blocked(export)} (pending), never \code{open(cap)}.
\item Item~8~(d$'$)(ii) (R1(b)): \code{G3.C\_T.lower}, and every aggregate field that may need $U_p$ at $\ell\in\{2,3\}$, is \code{blocked(U\_p)} on every \code{full} record at
  $N>1$ until Z2 lands or Z2-D's negative is recorded and audited: no \code{full} record at $N>1$ is final before Z2.
\item Item~8~(c): the refusal reason \code{non\_reduced} added to the closed list and defined (final).
\item Item~8~(d), G1 (W2, W1): \code{C2\_status}'s as-built \code{blocked(U\_p)} ``until UP-Z's (d$'$)(ii)'', needing $U_p$ ``at the word primes'' (was ``over $\Z$''); \code{not\_run}
  added (closed list of \code{C2}: five words); the reading of \code{partial(unfactored)} beside a non-certified (C1) (certifies nothing; no \code{anemic} byte moves).
\end{itemize}
\paragraph{Non-normative.} The status note of (d$'$) (re-refereed); the header comment, \verb|\date|, the Release record's first entry; this entry. No label is added, renamed or
removed; no status tag or marker moves ($229$ on $88$); no record, golden line, fingerprint (\code{0x8ffc29730c4f9672}) or digest moves.

\subsection*{v1.26.0-pre.2 --- 2026-11-08 --- hecke-extensions slot 26, UP-Z: the referee's conditions and fixes}
\textbf{This version} is a pre-release, \textbf{not owner-released}, on top of v1.26.0-pre.1. It discharges every condition and finding of the adversarial referee
\file{docs/implementation/slot26/C126pre1\_referee.md} (ACCEPT-with-conditions; disposition \file{docs/implementation/slot26/C126pre2\_disposition.md}). The one new binding rule (S4) is \textbf{not yet seen by the owner}.
\paragraph{Normative.}
\begin{itemize}\setlength{\itemsep}{0pt}
\item Item~8~(d$'$)(ii) of \S\ref{sec:recordschema} (S4): at a \code{full} record, $N>1$, with \code{G1.reduced\_Q = non\_reduced\_certified} (e.g.\ $(8,10)$) no $\vartheta$ search is run and
  \code{G1.C1.status = open(cap)} (final), \code{G1.theta = open(cap)}, \code{G1.components = not\_run}, \code{G1.C2\_status = partial(unfactored)}, \code{G3.partition = U\_p\_tags};
  otherwise \code{G1.C1.status} $\in\{\code{certified},\code{open(cap)}\}$ and \code{G1.reduced\_Q} $\in\{\code{true},\code{undecided}\}$ as stated; \code{blocked(U\_p)} no longer in
  these $\ell$-independent fields. Also (E5) $\mathcal P$ in the record's sense; (E6) the three re-read phrases enumerated, $g_{\mathcal P}$ of \code{S\_rule} unchanged.
\item Item~8~(d$'$)(iv) (S3): cites \cl{rmk:stabconvention}(1), (2) and [U]'s orientation remark, and states that [U]'s ``classical $U$'' (Hecke systems, $\Ga_1(\ell^{r})$, $\chi$) and
  \cl{prop:levelconehecke}(3)(d) (the operator against $\mathrm{dg}_1$, $\mathrm{dg}_p$ at $\Ga_0(N)$, $\chi=1$) concern different objects; one operator-level convention is not asserted.
\item \cl{prop:strictliftN}(9): (b) the classical operator written $U^{\mathrm{cl}}_p$ (E2), $A\smat p001A\ne AgA$ at $p\mid N$ (E9); (c) ``their references to (7), (8) dropped''
  (E7); (d) the commutativity of $U_p$ with the $T_p$ is \cl{prop:heckeZspec}(0) (E3); the proof (9)(c) discharges \cl{prop:descentpairhecke}(1) at $p\mid N$ (S1: its item~(1)
  reads nothing of $N$).
\end{itemize}
\paragraph{Non-normative.} \S\ref{sec:descenthecke}: the introduction points at (9) (E4); ``What is used'' names \cl{prop:levelconehecke}(3)(d), \cl{prop:descentmaps}(3)(b$'$) and
says \cl{def:heckeZ}, \cl{prop:heckeZspec} are cited for identification only (S2); the notation sentence (E1, $U^{\mathrm{cl}}_p$). Status notes of (9) and (d$'$): refereed
ACCEPT-with-conditions (E8). E10: nothing to do (no \verb|\LEAN| macro added). The header comment, \verb|\date|, the Release record's first entry; this entry. No label is added,
renamed or removed; no status tag or marker moves ($229$ on $88$); no record, golden line, fingerprint (\code{0x8ffc29730c4f9672}) or digest moves.

\subsection*{v1.26.0-pre.1 --- 2026-11-08 --- hecke-extensions slot 26, UP-Z: $U_p$ at $p\mid N$, contract first}
\textbf{This version} is a pre-release, \textbf{not owner-released} and \textbf{not yet refereed}, on top of the signed v1.25.0; its changes are \textbf{binding changes not yet seen by
the owner}. They are condition~C5 of the critical audit \file{docs/implementation/slot26/UPZ\_design\_audit.md} of the UP-Z design \file{docs/implementation/slot26/UPZ\_design.md}
(ACCEPT route~(i), the invariant route, with conditions C1--C7): the text gates the merge of the code package Z1 and is refereed (adversarially) first. Its wording carries C1
(the freeze), C2 (the self-describing operator list), C4 (a)--(d) and C6 (the contract version in the provenance).
\paragraph{Normative.}
\begin{itemize}\setlength{\itemsep}{0pt}
\item \cl{prop:strictliftN}: new item~(9), $p\mid N$: (a) $A''=\Ga_0(pN)$ of index $p$, the transversal $r_t=\smat10{tN}1$, $t<p$; (b) $U_p:=[A\,\smat100p\,A]$ is the $U^{(N)}_p$ of
  \cl{prop:levelcorrespondence}(4)(d), of \cl{prop:heilbronn}(2) and of \cl{def:heckeZ}, the adjoint $U_p^{\ast}$ on $p$-old spaces (\cl{prop:levelconehecke}(3)(d)); (c) items (2)--(6) and
  the independence of the lift hold verbatim with the same certified level-$1$ $\tau_{\Iso}$, so the induced-fibre operator, the glueing map and the transports to $\ell\in\{2,3\}$
  exist at $p\mid N$; (d) items (1), (7), (8) are not asserted. The proof (9) checks, item by item, that $p\nmid N$ enters only the hypothesis line and items (1), (7), (8) --- the
  audit's \S2.3 list, which also corrects the design's ``(6)--(8)'' to ``(7), (8)''. Title, hypothesis line and the subsection's notation sentence (``$U_p$ does not occur'') amended
  accordingly; the label and the status tag are unchanged.
\item Item~8 of \S\ref{sec:recordschema}, new paragraph (d$'$) after (d): at $N>1$, \code{algebra = full}, the operator list is \code{G1.P}'s $T_p$ then $U_p$, $p\mid N$ (not
  \code{G0.operators}); $\mathcal P$ and its (C1) are the \code{anemic} ones; \code{theta\_coeff\_max} $=\binom d2$ with the full $d$; the operator is the adjoint-convention $U_p$,
  named; the routes (Heilbronn at $\ell\ge5$, the free model at $\ell\in\{2,3\}$; \code{blocked(U\_p)} until built, never \code{refused(no\_lift)} for want of a lift); the required
  hashed field \code{G1.theta\_ops} and unhashed \code{G10.contract} on \code{full} records at $N>1$; the freeze: any later change is a caps line or \code{/v2}, never an edit in place.
\end{itemize}
\paragraph{Non-normative.} The header comment, \verb|\date|, the Release record's first entry; this entry. Not changed (status edits after the Z1 merge, audit F7):
\cl{thm:descentcert}(5) ``$U_p$ is deferred'' and (7)(c), \cl{def:heckeZ}'s ``not built over $\Z$''. No label is added, renamed or removed; no status tag set or \verb|\LEAN| marker moves
($229$ on $88$); no record, golden line, fingerprint (\code{0x8ffc29730c4f9672}) or digest moves. Companion pointers: [E] 1.25.0, [U] 0.7.0, [X] v1.28.

\subsection*{v1.25.0 --- 2026-11-07 --- hecke-extensions slot 25: the owner-signed release of v1.25.0-pre.8}

\noindent\emph{This version.} \textbf{Signed by the owner 2026-11-07}: the owner-signed release of v1.25.0-pre.8; no change of content. It releases the content of v1.25.0-pre.1 to pre.8 as described in their entries below, including the binding changes there marked ``not yet seen by the owner'' (the refereed \S\ref{sec:heckeZ} and record item~8 of \S\ref{sec:recordschema}, rule~\ref{V1} as amended, Parts~VIII and IX); the owner signs them with this release. \textbf{No statement's content changes; no label is added, renamed or removed; no status tag moves; the marker count stays $229$ on $88$ labels.}
\begin{itemize}\setlength{\itemsep}{0pt}
\item Header and \verb|\date|; the Release record's first entry; this entry.
\end{itemize}
Companion pointers: [E] 1.25.0, [U] 0.7.0, [X] v1.28.

\subsection*{v1.25.0-pre.8 --- 2026-11-07 --- hecke-extensions slot 25, LaTeX polish stage 5c: the title page}
\textbf{This version} is a pre-release, \textbf{not owner-released}, on the branch \texttt{s25-wip-latexpolish-s5c}. Its one normative change, to rule~\ref{V1}, is a
\textbf{binding change not yet seen by the owner} in this form; the owner approved it in advance (ruling Q1 of 2026-11-06, \file{docs/hacs/OWNER\_RULINGS\_1106.md})
and signs it at the 1.25.0 release. The plan is stage~5c of \file{docs/hacs/LATEX\_polish\_plan.md};
the report is \file{docs/hacs/LATEX\_polish\_s5c\_report.md}.

\paragraph{Normative.}
\begin{itemize}\setlength{\itemsep}{0pt}
\item \cl{V1}: the supersession record leaves the date line. Old: \emph{The date line is \code{Version $X.Y$ --- YYYY-MM-DD}, never \code{\textbackslash today}, with the
  supersession record appended.} New: \emph{The date line is \code{Version $X.Y$ --- YYYY-MM-DD} and a status clause, never \code{\textbackslash today}. The supersession
  record --- one entry per version --- is the first subsection of Appendix~\ref{app:changelog}.}
\end{itemize}

\paragraph{Non-normative.} No statement, label, number or \verb|\LEAN| marker moves.
\begin{itemize}\setlength{\itemsep}{0pt}
\item The date line is one line under the amended \ref{V1}. Its text up to v1.25.0-pre.7 --- $167$ version entries (v1.25.0-pre.7 down to v0.1.0), the sentence
  ``Supersedes \dots'' and the companion pointer --- is the Release record, the first subsection of this appendix, verbatim. Form only: each entry and the two closing
  sentences became a paragraph; the $53$ line breaks \verb|\\| that ended the first $53$ entries were dropped, and the run-on line of v0.10.0 to v0.2.0 was split before
  each ``v0.'' (whitespace only).
\item The author line ``Orbifold Cohomology of $\GL(2)$'' is replaced by the uniform author block (owner ruling Q3 of 2026-11-06).
\item A generic abstract replaces the old one, which named versions. The old author line and abstract are quoted verbatim in the second subsection of this appendix.
  The old abstract's one label-citation site (of \code{cl:pd}) is therefore still in this file, and the label census row of this file is unchanged.
\item Two unnumbered part headings with their table-of-contents lines (owner ruling Q5 of 2026-11-06): ``Part VIII --- Level and character'' before \S\ref{sec:levelN}, and
  ``Part IX --- Integral and local structure'' before \S\ref{sec:saturation}. No section number or label moves; Part~VII now ends after \S\ref{sec:hidafull}.
\item Consequential edits of the part headings (owner ruling Q5), fix round~1 of the stage's audit; prose only, no statement changes.
  (a)~the code map, \S\ref{sec:codemap}: the rows of \S\ref{sec:levelN} and \S\ref{sec:nebentypus}, ``(Part~VII)'' $\to$ ``(Part~VIII)''; the rows of
  \S\ref{sec:saturation}, \S\ref{sec:isomodelell}, \S\ref{sec:descentN}, \S\ref{sec:isomodeltree}, \S\ref{sec:descenthecke}, \S\ref{sec:descentpair},
  \S\ref{sec:treepair} and \S\ref{sec:mpmodel}, ``(Part~VII)'' $\to$ ``(Part~IX)''.
  (b)~\S\ref{sec:scope}, ``How to read this document''. Old: \emph{Part~VII (\dots, \S\ref{sec:hidafull}, since v1.13.0-pre.1, and \S\ref{sec:levelN}, since
  v1.14.0-pre.1) is \dots\ control and the lattice datum; then level $N$ for every $N$ --- the composite coset table, the cokernel presentation and its Hecke action by
  Merel's matrices; it is placed after Part~VI for the same reason. Chapter~\ref{sec:notation} \dots\ binds all five parts;} New: \emph{Part~VII (\dots, and
  \S\ref{sec:hidafull}, since v1.13.0-pre.1) is \dots\ control and the lattice datum; it is placed after Part~VI for the same reason. Part~VIII (\S\ref{sec:levelN},
  since v1.14.0-pre.1, and \S\ref{sec:nebentypus}) is level $N$ for every $N$ --- the composite coset table, the cokernel presentation and its Hecke action by Merel's
  matrices --- and nebentypus; Part~IX (\S\S\ref{sec:saturation}--\ref{sec:mpmodel}) is the $\Z/\ell^{K}$ Hecke algebra, the $\ell$-local isotropy model with the
  descent at level $N$ and the Hecke algebra over $\Z$, and the Moerdijk--Pronk model. Chapter~\ref{sec:notation} \dots\ binds all nine parts;} (``five'' had been
  stale since Part~VI was added; the document has nine parts.)
\item A consequence of the amended \ref{V1}: in \S\ref{sec:oc1sites}'s list, ``the historical version lines of the date line'' $\to$ ``the historical version lines of
  the Release record (the first subsection of Appendix~\ref{app:changelog}; the date line until v1.25.0-pre.7)''.
\item The abstract is the stage audit's version: the Hecke algebra over $\Z$ is listed with the $\ell$-local isotropy model (it is \S\ref{sec:heckeZ}), and the status
  sentence follows Appendix~\ref{app:status}'s wording.
\item The appendix keeps its title ``Changelog'', which rule~\ref{V3} names; the plan's retitling (``Revision history'') would have changed \ref{V3}'s text.
\item The header comment, \verb|\date|; this entry. Companion pointers unchanged.
\end{itemize}
\noindent\emph{Not changed, listed for the owner:} the counts in \S\ref{sec:scope} (``$53$ C++ source files \dots\ at $519$ sites across $71$ files'') are stale; the label
census (\file{docs/latex/label\_census.tsv}) is the live record. Updating them is a prose change for a later version.

\subsection*{v1.25.0-pre.7 --- 2026-11-06 --- hecke-extensions slot 25, the HZ-RAT-PERF fix round}
\textbf{This version} is a pre-release, \textbf{not owner-released}, on the branch \texttt{s25-wip-ratperffix}; its changes to item~8 of \S\ref{sec:recordschema} are
\textbf{binding changes not yet seen by the owner}, who signs them at the 1.25.0 release. HZ-RAT-PERF (value-neutral speed and memory of the rational step) added the named
refusal \code{refused(word\_prime\_cap)}; its audit (\file{docs/implementation/slot25/HZRATPERF\_audit.md}: PASS\_WITH\_MINORS for the code, F2 MAJOR: the word was code ahead of
the binding text; F6: the draft sentence incomplete) asked for the text with or before the merge:
\begin{itemize}\setlength{\itemsep}{0pt}
\item (c): \code{refused(word\_prime\_cap)} in the list of pending statuses and beside ``No other pending word exists''; a paragraph: what the list is (the \code{CrtPrimes}
  order without the divisors of $N$, prefix-stable, the $U_p$ tag's list from $\min(2048,\text{cap})$ primes), why the word is pending (only a larger cap lifts it), what a
  refused record contains, and that the word enters the frozen schema as amendment A1 (hash-neutral).
\item G1: the closed list of \code{C2\_status} as built (\code{complete}, \code{partial(unfactored)}, \code{blocked(U\_p)}, \code{refused(word\_prime\_cap)}).
\item (i): the count cap in the unhashed \code{word\_primes=} provenance key.
\item Header, \verb|\date|; this entry.
\end{itemize}
No label is added, renamed or removed; no status tag set or \verb|\LEAN| marker moves ($229$ on $88$); no record, golden line, fingerprint (\code{0x8ffc29730c4f9672}) or digest moves.

\subsection*{v1.25.0-pre.6 --- 2026-11-06 --- hecke-extensions slot 25, the H1 merge}
\textbf{This version} is a pre-release, \textbf{not owner-released}; its changes to item~8 of \S\ref{sec:recordschema} are \textbf{binding changes not yet seen by the owner}, who signs
them at the 1.25.0 release. The package H1 (the Eisenstein component over $\Lambda$, [U] part~6) writes G8 under \code{heckez.G8 = hida}; its audit
(\file{docs/implementation/slot25/H1\_audit.md}: PASS\_WITH\_MINORS; F1 MEDIUM, F2 MAJOR (report; an upstream HZ-RAT defect), F3--F5 MINOR, F6--F9 NOTE) recommended the text:
\begin{itemize}\setlength{\itemsep}{0pt}
\item (h): a G8 paragraph beside G5--G7 --- only under \code{heckez.G8 = hida}; the paths \code{G8.scope}, \code{G8.ell=<l>.blocks}, \code{G8.ell=<l>.m[i].\{ordinary, chain\_key, g, L\_k, v\_Bk,
  structure, I0\_generator\}}; their closed lists and meanings, citing [U,~\texttt{prop:u-H11}], [U,~\texttt{prop:u-H12}] and the new [U,~\texttt{rmk:u-H12general}] ([U] 0.7.0-pre.3);
  the caps as count-based record content, the $q$-expansion cap naming $\vartheta$'s operators (F4); G9 \code{K9.ell=<l>} (with \code{absent(not\_applicable)} under \code{hida}, F3) and G10
  \code{hida.ell=<l>}.
\item (i): a pointer to it. (k): golden only under the default settings (F1; the code's \code{golden\_eligible} now requires the default \code{heckez.G8 = off} line).
\item Header, \verb|\date|; this entry.
\end{itemize}
No label is added, renamed or removed; no status tag set or \verb|\LEAN| marker moves ($229$ on $88$); no default record, golden line or digest moves (F4 moves only the \code{hida} settings
fingerprint and the three \code{hida} regression pins, which are never golden).

\subsection*{v1.25.0-pre.5 --- 2026-11-06 --- hecke-extensions slot 25, the HZ-REC schema re-audit fix round}
\textbf{This version} is a pre-release, \textbf{not owner-released}, on the branch \texttt{s25-wip-hzrec}; its changes to item~8 of \S\ref{sec:recordschema} are
\textbf{binding changes not yet seen by the owner}, who signs them at the 1.25.0 release. The re-audit of pre.4's record (\file{docs/implementation/slot25/HZREC\_schema\_reaudit.md}:
FAIL, narrow; R1 MAJOR, R2 MEDIUM, R3--R7 MINOR, R8--R12 NOTE) and the lead's decisions on it, each marked in the item:
\begin{itemize}\setlength{\itemsep}{0pt}
\item R1 (MAJOR): \code{components[0].mult} is $e_k$, the multiplicity of the Eisenstein system on $H^{1}$ (pre.4 said $1$; wrong at $N>1$ with $e_k\ge2$); (d).
\item R2: the second route's word caps in the settings text (\code{heckez.second\_route}) and \code{not\_read} also where \code{v\_Bk = not\_run}; (f), (j). Every record hash moves.
\item R3, R4, R10, R11: the cuspidal block iff $d\ge2$; the \code{not\_run} of \code{components\_through}, \code{pair\_kinds.kind}, \code{G7.pair} and the cusp multiplicity at $N>1$;
  the (C1) probe primes; where the \code{theory} tag is written; (d)--(i).
\item R6: \code{[Lift]} is \code{refused(no\_lift)} only for a certificate code, a cap of the lift is \code{blocked(export)}; (c).
\item R8: the irreducibility cap $256$ (was $64$); R9: the $\vartheta$ search over $\mathcal P\cup\{T_\ell\}$ where $\mathcal P$ fails; both in the settings text; (d), (e), (j).
\item Header, \verb|\date|; this entry. R5, R7, R12 are code and documents (\file{docs/implementation/slot25/HZREC\_report.md}, ``Re-audit fixes'').
\end{itemize}
No label is added, renamed or removed; no status tag set or \verb|\LEAN| marker moves ($229$ on $88$); no golden line or digest moves (the record's golden objects are not yet emitted).

\subsection*{v1.25.0-pre.4 --- 2026-11-06 --- hecke-extensions slot 25, the HZ-REC schema-audit fix round}
\textbf{This version} is a pre-release, \textbf{not owner-released}, on the branch \texttt{s25-wip-hzrec}; its changes to item~8 of \S\ref{sec:recordschema} are
\textbf{binding changes not yet seen by the owner}, who signs them at the 1.25.0 release. The schema audit of the built record
(\file{docs/implementation/slot25/HZREC\_schema\_audit.md}: FAIL as a short fix round, 2 MAJOR, 4 MEDIUM, 8 MINOR, 2 NOTE) and the lead's decisions on it:
\begin{itemize}\setlength{\itemsep}{0pt}
\item F1 (MAJOR): the new final G6 value \code{split\_on\_P} for every $\Talg'$-split point; \code{split} (for $\Talg$) needs the mixed generation of
  Remark~\ref{rmk:heckeZcbthree} (\CONJ) and is not written in v1; (h), (k).
\item F2 (MAJOR): no G8 field in default v1 records; H1 writes G8 under the non-default settings line \code{heckez.G8}; (i), (j).
\item F3: G4--G6 per torsion point (a joint $\Talg'$-block with torsion) with its system, no per-$\ell$ aggregate; (g), (h).
\item F5, F10: the components as written (the cuspidal part as one block, \code{irreducible} certified by a prime or \code{open(cap)}, $p_0$, no block at $d=e_k$), the closed
  lists of \code{type}, $p_0$ and \code{partition}; (d), (f).
\item F6, F12: the maximal-ideal rule (a searched $\vartheta$ generating $A'_1$), \code{degree\_attr}, \code{generation\_tors}, $q=2$ and G4's \code{lower} as v1 writes them; (e)--(h).
\item F7, F8: \code{pending(HZ-TOR)} retired, the export's refusals mapped by name; (c). F11, F16: the caps text before it is frozen, defaults never change; (i), (j).
  O5, O6: \code{agrees = not\_read}, \code{eta\_L.sign}; (f).
\item Header, \verb|\date|, the item's status sentence; this entry. F4, F9, F13--F15 are code (\file{docs/implementation/slot25/HZREC\_report.md}, ``Schema-audit fixes'').
\end{itemize}
No label is added, renamed or removed; no status tag set or \verb|\LEAN| marker moves ($229$ on $88$); no golden line or digest moves.

\subsection*{v1.25.0-pre.3 --- 2026-11-06 --- hecke-extensions slot 25, the Z0 re-check}
\textbf{This version} is a pre-release, \textbf{not owner-released}; text only (\file{docs/implementation/slot25/Z0\_recheck.md}). Item~8~(c): pre.2's binding rule
(\code{export\_word} iff $2^{31}\le\ell^{K}<2^{63}$) left an EXPORT-served read at $\ell^{K}<2^{31}$ with $K$ above EXPORT's $K_{\max}=K_{\mathrm{run}}-\sigma_{\max}-2\sigma^{(1)}_{\max}$
(the precision lost to the $\ell$-torsion; level $N>1$ at $\ell=2,3$ among others) without a record word, since \code{open(precision)} was defined as the settings' step cap.
Now \code{export\_word} iff $\ell^{K}<2^{63}$ and $K>K_{\max}$ of EXPORT (the word bound included), pending as before (EXPORT-Z step~2 serves every $K$);
\code{open(precision)} is the step cap only. No label is added, renamed or removed; no status tag set or \verb|\LEAN| marker moves ($229$ on $88$); no computed value,
record, golden line or digest moves.

\subsection*{v1.25.0-pre.2 --- 2026-11-06 --- hecke-extensions slot 25, the Z0 merge after its referee}
\textbf{This version} is a pre-release, \textbf{not owner-released}. v1.25.0-pre.1 was refereed REVISE (\file{docs/implementation/slot25/Z0\_referee.md}: 0 BLOCKER, 3 MAJOR,
4 MINOR, 4 NOTE); pre.2 applies its findings as text and the lead's decision on its open item O1:
\begin{itemize}\setlength{\itemsep}{0pt}
\item F2 (\textbf{a statement narrowed}): \S\ref{sec:heckeZ}, ``Two forms'', level $N>1$, full --- \code{non\_reduced\_certified} by the counting criterion is asserted only for $N$ a power
  of $p$; for other $N$ it is not proved (HZ-UTAG audit F1). This replaces pre.1's ``no statement's mathematics changes'' for that sentence; the code was already restricted.
\item F3, F4: item~8~(c) --- \code{refused(export\_word)} is \textbf{pending} (EXPORT-Z step~2 widens the ring), not final; the binding rule \code{export\_word} iff
  $\ell^{K}<2^{63}$ and $K$ exceeds EXPORT's $K_{\max}$ (the word bound $\ell^{K}\ge2^{31}$ or its torsion precision; pre.3) at an EXPORT-served read, \code{word\_modulus} tested first; the as-built rule ($\ell\ge2^{31}$ refused at every $K$, $K=1$ included; \code{open(precision)}
  at $\ell<46341$ above EXPORT's $K_{\max}$); \code{open(precision)} added to the closed list (final: the step cap). Same in \S\ref{sec:heckeZ}'s \emph{As built} paragraph.
\item O1 (lead decision): no new $K=1$ route now; (C3) at $K=1$, level~$1$, $\ell\in[2^{31.5},2^{63})$ is the pending \code{refused(k1\_route)}, an EXPORT-Z / slot-26 item, not a
  final narrowing of \code{word\_modulus}.
\item F5: the cofactor of (C2) and of $g_{\mathcal P}$ as built (the gcd over the primitive $\vartheta$'s, residue-checked, else \code{upper\_bound}) in \S\ref{sec:heckeZ}'s cofactor
  paragraph and item~8~(d), (f).
\item F6: item~8~(j) --- \code{t\_ell} $\in\{\code{in\_P},\code{added},\code{omitted(cap:ell>23170)},\code{omitted(ell|N:U\_ell\_not\_exported)},\code{omitted(export:<Name>)}\}$, the
  refusal name only; the cap $23170$ in \code{P\_rule}'s default text. (k): a pending refusal reason keeps a record out of the golden set.
\item F7: Appendix~\ref{app:status}, the referee item ``Part~VII, \S\ref{sec:heckeZ}'' records A7 $\to$ A7b ACCEPT and this referee.
\item Header, \verb|\date|; this entry. F1, F8--F11 concern the merge report, the slot history and HZ-REC (\file{docs/implementation/slot25/Z0\_merge\_report.md}).
\end{itemize}
No label is added, renamed or removed; no status tag set or \verb|\LEAN| marker moves ($229$ on $88$); no computed value, record, golden line or digest moves.
Companion pointers: [E] 1.24.0, [U] 0.6.0, [X] v1.28.

\subsection*{v1.25.0-pre.1 --- 2026-11-06 --- hecke-extensions slot 25, the Z0 merge: the Hecke algebra over $\Z$ in the main line}
\textbf{This version} is a pre-release, \textbf{not owner-released}, on top of the signed v1.24.0. By the owner's ruling Q8 (2026-11-04) the HeckeZ text drafted and refereed on the
branch \texttt{s24-wip-heckez-z0} (v1.25.0-pre.0, pre.0b, pre.0c below; referee A7 REVISE, re-audit A7b ACCEPT: \file{docs/implementation/slot24/heckez/A7-Z0-referee.md},
\file{A7b-Z0-reaudit.md}) enters the main line now that v1.24.0 is released. Taken exactly as Z0 defines it: \S\ref{sec:heckeZ} ($23$ statements), item~8 of
\S\ref{sec:recordschema}, the pointer in \cl{thm:descentcert}(7)(a), the notation list, $23$ lines of Appendix~\ref{app:status}, three bib entries and the pre.0--0c entries
below; everything outside these keeps v1.24.0's text (the hunk list: \file{docs/implementation/slot25/Z0\_merge\_report.md}). Reconciled, each by a note dated 2026-11-06, with the
plan's \file{04-record-census.md} (its slot-25 amendment of \S1.5) and the code of slot~25 (EXPORT, HZ-RAT, HZ-UTAG, HZ-LOC):
\begin{itemize}\setlength{\itemsep}{0pt}
\item Item~8~(c): the refusal reason \code{export\_word} (level $N>1$, $\ell\ge46341$, EXPORT's ring $\ell^{K}<2^{31}$). (d): \code{cofactor\_test} gains \code{upper\_bound}.
  (j): the read set at $\ell$ is $\mathcal P\cup\{T_\ell\}$ for $\ell\nmid N$, $\ell\le23170$, with the hashed field \code{t\_ell}; the factor budget as built. The code map.
\item \S\ref{sec:heckeZ}: the word-modulus paragraph as built (the seven $64$-bit reads; (C3) at $K=1$ for $\ell\in[2^{31.5},2^{63})$ at level~$1$ not served, an open item;
  \code{export\_word}); the counting criterion of \cl{thm:heckeZcert}'s full-algebra tag set only for $N$ a power of $p$; the (C2) frontier as built.
\item Header, \verb|\date|, the body's two version pointers; this entry.
\end{itemize}
No label is added, renamed or removed beyond Z0's; no status tag set or \verb|\LEAN| marker moves ($229$ on $88$); no computed value, record, golden line or digest moves.
Companion pointers: [E] 1.24.0 (the \Espec{cl:sturm}(a) erratum Z0 cites is in it), [U] 0.6.0, [X] v1.28.

\subsection*{v1.24.0 --- 2026-11-05 --- hecke-extensions slot 25: the owner-signed release of v1.24.0-rc.2}

\noindent\emph{This version.} \textbf{Signed by the owner 2026-11-05}: the owner-signed release of v1.24.0-rc.2; no change of content. It releases the content of v1.24.0-pre.1, rc.1 and rc.2 as described in their entries below, each with the referee status stated there (rc.1 refereed ACCEPT\_WITH\_MINORS, all findings applied in rc.2). \textbf{No statement's content changes; no label is added, renamed or removed; no status tag moves; the marker count stays $229$ on $88$ labels.}
\begin{itemize}\setlength{\itemsep}{0pt}
\item Header and \verb|\date|; this entry.
\end{itemize}
Companion pointers: [E] 1.24.0, [U] 0.6.0, [X] v1.28.

\subsection*{v1.24.0-rc.2 --- 2026-11-05 --- hecke-extensions slot 25, RC124: the release candidate of 1.24.0, after its referee}
\textbf{This version} is the release candidate of 1.24.0, \textbf{not owner-released}. v1.24.0-rc.1 was refereed ACCEPT\_WITH\_MINORS
(\file{docs/implementation/slot25/RC124\_referee.md}); rc.2 applies its findings as text (the map finding $\to$ fix: \file{docs/implementation/slot25/RC124\_report.md}, ``rc.1 $\to$ rc.2''):
\begin{itemize}\setlength{\itemsep}{0pt}
\item F1: Remark~\ref{rmk:bocksteinexp}(4), the convention --- \cite{ccs13} grade with $\deg x=\deg y=2$; their degree-$2n$ part is $\Vcoef$ and their condition is the one stated.
\item F2, F6: \S\ref{sec:status} --- the marker rule amended by one sentence (beside \PROVED{} or \STD{}, and clause-scoped on definitions and displayed equations,
  as already applied); the rc.1 sentence shortened to a pointer, ``the full form'' $\to$ ``the group-cohomological form''. No marker or tag moves.
\item F3, F4: Remark~\ref{rmk:levelNcurvefamily} --- the printed \code{layers=} are the dimensions $\dim\mathfrak{m}^{i}W=255-2i$; the curve-object confirmation and the $1297$
  sentence are reworded as reported records this document does not assert (nothing uses them; no tag added).
\item F7: Remark~\ref{rmk:bocksteinexp}, the scope sentence's two parentheses merged. F8: Proposition~\ref{prop:ltwothreescalar}, the Hatada attribution; \file{gl2\_refs.bib}:
  \code{calegari\_emerton\_largeindex}'s note gains Lemma~3.1's hypothesis and ``journal version not checked''.
\item Header, \verb|\date|, this entry.
\end{itemize}
No label is added, renamed or removed; no status tag set or marker moves ($229$ on $88$); no computed value, record, golden line, digest, validation or code path moves.
Companion pointers: [E] 1.24.0-rc.2, [U] 0.6.0, [X] v1.28.

\subsection*{v1.24.0-rc.1 --- 2026-11-05 --- hecke-extensions slot 25, RC124: the release candidate of 1.24.0}
\textbf{This version} is the release candidate of 1.24.0: \textbf{not owner-released and not yet refereed} (its adversarial referee is
\file{docs/implementation/slot25/RC124\_referee.md}; the change list with each change's status is \file{docs/implementation/slot25/RC124\_report.md}).
It contains v1.24.0-pre.1 unchanged (FMTMP, \S\ref{sec:recordschema} item~7 (l), owner request of 2026-11-03) and adds:
\begin{itemize}\setlength{\itemsep}{0pt}
\item \textbf{The Lean markers of question Q4} (approved in principle by the owner on 2026-10-18; \file{docs/implementation/slot24/LEAN\_M6\_report.md} \S4 as
  corrected by \file{LEAN\_M6\_audit.md} F1, F3): $3$ on Theorem~\ref{thm:classical} (\texttt{thm:m6-classical-theta}, \texttt{thm:m6-Tcls-res},
  \texttt{thm:m6-hecke-res}) and $4$ on \eqref{eq:corconjres} (\texttt{thm:m6-cor-res}, \texttt{thm:m6-corH-indep}, \texttt{thm:m6-cor-res-index},
  \texttt{thm:m6-Tcls-res}), each clause-scoped by its parenthesis; the marker list of \S\ref{sec:status} brought up to date. $222$ markers on $86$ labels become
  $229$ on $88$ ($215$ \code{\textbackslash LEAN}, $14$ \code{\textbackslash LEANlit}). The proposal's third marker (the witness of Proposition~\ref{prop:mpbredon}(2)) is not applied.
\item \textbf{The phenomenology citations} (\file{docs/phenomenology/FINDINGS\_SUMMARY.md} \S3, \file{phenomena\_catalogue.md} \S11 items 2, 3, 5, 6; \textbf{not yet seen by the owner}):
  Remark~\ref{rmk:bocksteinexp} --- clause~(4) is now the theorem of Callegaro--Cohen--Salvetti \cite[Thm.~4.5]{ccs13} (\STD, locator \VERIFYloc), the full multiset of
  exponents at every prime, so the largest exponent equals the lower bound of~(2) and is $1$ off its loci; with the module convention checked against
  Proposition~\ref{prop:symconv}; clause~(3) is a confirmation of it; ``Not asserted'' keeps only [L]'s growth estimate.
  Proposition~\ref{prop:ltwothreescalar}~(F1) --- its origin, \cite[Lemma~3.1]{calegari_emerton_largeindex} after Hatada.
  \ref{Hi11} --- a route to a proof through Hida's control modulo $\ell$, not checked; the tag stays \CONJ.
  Remark~\ref{rmk:levelNcurvefamily} --- the locator Thm.~1.1; the heavy oracle $(21929,2)$, $e_2=127$, is Calegari--Emerton's own example; their \S6 sentence on $2143$
  is false ($1297$) and not used. \file{gl2\_refs.bib}: \code{ccs13}, \code{calegari\_emerton\_largeindex}, \code{rustom\_primepowers}.
\item Header and \verb|\date|; Appendix~\ref{app:status}'s row of Remark~\ref{rmk:bocksteinexp} gains the citation (its tag set does not change); this entry.
\end{itemize}
No label is added, renamed or removed; no status tag set moves; no computed value, record, golden line, digest, validation or code path moves.
Companion pointers: [E] 1.24.0-rc.1, [U] 0.6.0, [X] v1.28.
\subsection*{v1.25.0-pre.0c --- 2026-11-04 --- hecke-extensions slot 24, HeckeZ Z0: the A7b re-audit's MINOR (branch pre-release)}
\textbf{This version} is still a \textbf{branch pre-release} on \texttt{s24-wip-heckez-z0}, not owner-released, on top of v1.25.0-pre.0b. The adversarial re-audit A7b
(\file{docs/implementation/slot24/heckez/A7b-Z0-reaudit.md}) returned ACCEPT with $0$ BLOCKER, $0$ MAJOR, $1$ MINOR and $5$ NOTE findings. No label, number, tag set or marker
moves ($222$ on $86$ labels); no golden line, no code.
\begin{itemize}\setlength{\itemsep}{0pt}
\item \cl{thm:heckeZcert}(C2) (A7b-m1). The clause on the other $\vartheta_j$ is now a statement: a certified prime $q_i$ of $D_{j_0}$ with $v_{q_i}(D_j)\le1$ for some $j$ may be
  left out of $\mathsf S$; the proof of (2) covers it (Proposition~\ref{prop:heckeZetale}(2)). The measured (C2) frontier is re-measured on one $D_{j_0}$, $\vartheta_{j_0}=T_2$
  (\file{a7b\_frontier\_single.py}): the same weights as on $\gcd_{p\le7}$; Q1(a)'s loss list $k\in\{36,42,44,48,50,54,56,58,62,64\}$ is unchanged.
\item NOTEs: \cl{prop:heckeZeiscert}(3) ``with multiplicity'' (N1); ``one modular power and one residue'', a read below $v_\ell(\nu)$ is a defect (N4). N2, N3, N5 need no change.
\item Header, \verb|\date|; this entry.
\end{itemize}

\subsection*{v1.25.0-pre.0b --- 2026-11-04 --- hecke-extensions slot 24, HeckeZ Z0 fix round: the A7 referee's findings (branch pre-release)}
\textbf{This version} is still a \textbf{branch pre-release} on \texttt{s24-wip-heckez-z0}, not owner-released, on top of v1.25.0-pre.0. The adversarial Opus referee A7
(\file{docs/implementation/slot24/heckez/A7-Z0-referee.md}) returned REVISE with $0$ BLOCKER, $1$ MAJOR, $7$ MINOR and $10$ NOTE findings; all MAJOR and MINOR findings are
applied, and the NOTEs N1, N2, N4, N6--N9 (N3, N5, N10 need no change); the table is in \file{Z0-report.md} \S5. No label, number, tag set or marker moves ($222$ on $86$
labels); no golden line, no code.
\begin{itemize}\setlength{\itemsep}{0pt}
\item \cl{prop:heckeZeiscert}(3) (A7-M1, MAJOR). A good prime $q$ does not keep a non-Eisenstein block out of the Hensel factor at $1+q$: at $(N,\ell,q)=(61,5,37)$, $q$ good, the
  rational newform of level $61$ ($a_{37}=8\equiv38\bmod5$) enters it, its degree is $2>1$ and its value at $38$ has $5$-adic valuation $2>1=v_5(\nu)$ (exact modular-symbol
  computation, \file{docs/implementation/slot24/heckez/a7fix\_mazur\_cx.py}; some good $q<60$ is contaminated at $27$ of the $46$ pairs $(N,\ell)$ with $N\le260$). The statement now
  rests on the \textbf{valuation check} $v_\ell(\mathsf F_q(1+q))=v_\ell(\nu)$, which is necessary and sufficient for $\mathsf F_q=f_{q,\midl}$ and $(\eta_q)=I_\midl$; goodness is how $q$ is
  chosen. $\mathsf F_q$ is the Hensel factor of $\mathsf h_q$, the square root of the cuspidal factor --- not of a squarefree part --- so the route is residue-only (Q1(a)).
  The rank is $e^{\mathrm{cusp}}_\midl$ (A7-m7). Record item~8~(e): \code{mazur} carries both valuations and certifies only on agreement.
\item \cl{prop:heckeZrat}: (3$'$)'s proof, ``no Eisenstein system is cuspidal'', at $p\ge7$ for every $\psi$ (A7-m1); the citations of (3$'$), (4)(a), (4)(b), (5)(a) by the
  verified [GT] rows \texttt{lit:hz-es-anemic}, \texttt{lit:hz-eis-dim}, \texttt{lit:hz-semisimple-an}, \texttt{lit:hz-sturm-integral} (Diamond--Im Thm.~6.2.3, Rmks.~6.3.3--6.3.4,
  Prop.~12.4.13; Stein Thms.~5.8--5.10, 9.23) in place of Diamond--Shurman, and the Appendix~\ref{app:status} row (A7-m5, N2); (3$'$)'s ``iff'' excludes $(N,k)=(1,2)$ (N9).
\item \cl{prop:heckeZspec}(3): the converse (no other minimal primes) by going-down for the $\Z$-free $\Talg^{\mathrm{fr}}$ (A7-m2).
\item \cl{thm:descentcert}(7)(a)'s pointer and the closing sentence after \cl{prop:heckeZglue}: the extension class is determined by finite data; a split verdict needs
  \cl{rmk:heckeZcbthree}; ``the one invariant \dots\ needs'' becomes ``the invariant \dots\ uses'' (A7-m3). The record of item~7 still asserts neither.
\item \cl{thm:heckeZcert}(C2) factors one $D_{j_0}$, no gcd of big integers; the cofactor paragraph restores the plan's two Q1(a) sentences; item~8~(d), (f) and the
  ``factored $g_{\mathcal P}$'' paragraph accordingly (A7-m4).
\item NOTEs: \cl{prop:heckeZeiscert}(3) cites the verified row \texttt{lit:hz-mazur-eis} without \verb|\VERIFYloc| (N1); $g_{\mathcal P}$ prime to $6$ at level $1$ (N4);
  \cl{prop:heckeZgen}'s lattice in a stable subspace (N6); \cl{rmk:heckeZcaveat}(3) the square-$u$ case, (10) cokernel versus $\mathrm{Tor}$ (N7).
\item Header, \verb|\date|, the notation list, Appendix~\ref{app:status}'s row of \cl{prop:heckeZeiscert}; this entry.
\end{itemize}
Companion pointers: [E] 1.24.0-pre.2 (\Espec{cl:sturm}(a): the condition for $k\equiv2\bmod12$, A7-m6; range $4\le k\le70$, N8), [U] 0.6.0.

\subsection*{v1.25.0-pre.0 --- 2026-11-04 --- hecke-extensions slot 24, HeckeZ package Z0: the Hecke algebra over $\Z$ (branch pre-release)}
\textbf{This version} is a \textbf{branch pre-release} on \texttt{s24-wip-heckez-z0}, not owner-released, not refereed (an adversarial Opus referee follows, the plan's step
A7), on top of v1.24.0-pre.1. By the owner's ruling Q8 of 2026-11-04 it enters the main line as \textbf{v1.25.0-pre.1} only after the owner has released v1.24.0; neither
release waits for the other. Source: the HeckeZ plan \file{docs/implementation/slot24/heckez/}, signed 2026-11-04 with every question at its recommendation
(\file{06-work-packages.md} v0.4 \S5; package Z0, \S2.3). No marker moves ($222$ on $86$ labels); no golden line, no code.
\begin{itemize}\setlength{\itemsep}{0pt}
\item \PNEW{} \S\ref{sec:heckeZ}, appended at the end of \S\ref{sec:isomodelell} (after \S\ref{sec:descentpair} and \S\ref{sec:treepair}), so no number moves. From
  \file{01-math-core.md} v0.6: \cl{def:heckeZ} (untagged; Def.~1.1), \cl{prop:heckeZspec} (\PROVED, \BSTD; Prop.~1.2), \cl{prop:heckeZrat} (\PROVED, \BSTD, \CONJ; Props.~1.3,
  1.4, 1.5(2)--(4)), \cl{lem:heckeZbridge} (\PROVED; Lemma~1.6), \cl{prop:heckeZmax} (\PROVED, \BSTD; Prop.~2.1), \cl{prop:heckeZgen} (\PROVED; Prop.~2.2, Rem.~2.2$'$),
  \cl{prop:heckeZetale} (\PROVED; Prop.~2.3), \cl{prop:heckeZmono} (\PROVED, \BSTD; Prop.~2.4), \cl{prop:heckeZlocal} (\PROVED; Prop.~2.5), \cl{prop:heckeZthrough} (\PROVED;
  Prop.~2.6), \cl{thm:heckeZdescent} (\PROVED; Thm.~3.1), \cl{thm:heckeZcert} (\PROVED, \BSTD; Thm.~4.1, Prop.~4.2), \cl{rmk:heckeZcaveat} (\PROVED; Rem.~4.5, Props.~4.3, 4.4),
  \cl{def:heckeZcong} (untagged; Def.~5.1), \cl{prop:heckeZcong} (\PROVED, \BSTD; Props.~5.2--5.5), \cl{prop:heckeZeiscert} (\PROVED, \BSTD; Prop.~5.6, Rem.~5.6$'$). From
  \file{03-bockstein-tor.md} v0.3: \cl{prop:heckeZuct} (\PROVED, \BSTD; B1 with B4(1)), \cl{prop:heckeZsplit} (\PROVED; B7), \cl{prop:heckeZlevelone23} (\PROVED; B8, B8$'$, B9),
  \cl{prop:heckeZglue} (\PROVED, \BSTD; B11--B15), and the conjectures CB-1, CB-2, CB-3 as \cl{rmk:heckeZcbone}, \cl{rmk:heckeZcbtwo}, \cl{rmk:heckeZcbthree} (\CONJ; used by
  nothing). Not written as statements: the plan's Prop.~1.5(1) (\CONJ, used nowhere; \Espec{cl:sturm} as corrected at [E] 1.24.0-pre.1 is cited instead), Prop.~5.7 (a sketch),
  the per-block remark of Prop.~2.4 (a sketch) and the stopping rule of B15(2) (a sketch) --- each is named in the text as not asserted.
\item \S\ref{sec:recordschema} gains item~8, \code{integral\_hecke\_structure/v1}, clauses (a)--(l) (\file{04-record-census.md} v0.4 \S1.6; (l), the big-integer rule, at Q1(a)).
\item \cl{thm:descentcert}(7)(a): a sentence pointing to \cl{prop:heckeZcong}(5) for the read at a root and to \cl{prop:heckeZglue}(2)(a) for the extension class; the record of
  item~7 still asserts neither.
\item Appendix~\ref{app:status}: $23$ lines; \file{gl2\_refs.bib}: \code{atiyah\_macdonald}, \code{bruns\_herzog}, \code{coleman\_edixhoven}; header and \verb|\date|; this entry.
\end{itemize}
Companion pointers: [E] 1.24.0-pre.1 (the \Espec{cl:sturm}(a) erratum, owner ruling Q9), [U] 0.6.0.

\subsection*{v1.24.0-pre.1 --- 2026-11-03 --- hecke-extensions slot 24, FMTMP: the faithful coefficients of $\chi$}
\textbf{This version} is a pre-release, not owner-released, not refereed, on top of the signed v1.23.0. \S\ref{sec:recordschema} item~7 gains the lettered item~(l),
\emph{The faithful coefficients}: every coefficient of $\chi^{(N)}_{p,1}$ beyond $2^{63}$ is recorded also by its exact value, the canonical signed hexadecimal text of the integer,
in the entry's list \code{big}, and hashed by one body line each; items~(d) and~(h) point to it. A record without such a coefficient keeps its bytes and its \code{record\_fnv1a64};
the digest of~(d) and the golden reference do not move; cell ids and resume keys are unchanged. Text only: no label, number, tag set or marker moves.
\begin{itemize}\setlength{\itemsep}{0pt}
\item \S\ref{sec:recordschema} item~7: (l) new; one clause each in (d) and (h).
\item Header and \verb|\date|; this entry.
\end{itemize}

\subsection*{v1.23.0 --- 2026-11-03 --- hecke-extensions slot 24: the owner-signed release of v1.23.0-pre.40}

\noindent\emph{This version.} \textbf{Signed by the owner 2026-11-03} (``I sign off 1.23.0''): the owner-signed release of v1.23.0-pre.40; no change of content. It releases the content of v1.23.0-pre.1 \dots{} pre.40 as described in their entries below, each with the referee status stated there. \textbf{No statement's content changes; no label is added, renamed or removed; no status tag moves; the marker count stays $222$ on $86$ labels.} The Lean markers on Theorem~\ref{thm:classical} and Equation~\ref{eq:corconjres} (Q4) are not part of this release.
\begin{itemize}\setlength{\itemsep}{0pt}
\item Header and \verb|\date|; this entry.
\end{itemize}
Companion pointers: [E] 1.23.0, [U] 0.6.0, [X] v1.28.

\subsection*{v1.23.0-pre.40 --- 2026-11-02 --- hecke-extensions slot 24, ICENSUS-P2: the campaign's \texttt{parts} key}
\textbf{This version} is a pre-release, not owner-released, not refereed, on top of v1.23.0-pre.39. \S\ref{sec:recordschema} item~7 gains the lettered item~(k), \emph{Parts not run}: a descent campaign cell may compute only the groups
(\code{parts = groups}) or the groups and the Hecke data (\code{parts = groups,hecke}); a part not run is marked \code{"ran":false}, a \code{groups} record's hashed body says \code{hecke=not\_run}, a non-default value is part of the cell id,
and the default \code{all} leaves every id, record and hash byte for byte as before. Text only: no label, number, tag set or marker moves, and no golden line.

\subsection*{v1.23.0-pre.39 --- 2026-11-02 --- hecke-extensions slot 24, the $\ker/\coker$ width ladder (ICENSUS-P1)}

\noindent\emph{This version.} \textbf{A pre-release, not owner-released, not refereed.} On top of v1.23.0-pre.38, item~7 of \S\ref{sec:recordschema}:
\begin{itemize}\setlength{\itemsep}{0pt}
\item (e): the $\ker/\coker$ read's working width is a ladder, $k_0=\min(K_{\max}(\ell),2a_{\max}+3)$, then $K_{\max}(\ell)$ on a precision shortfall; the width is not in the canonical text and is not one of (e)'s caps.
\item (j): the paragraph on the refusal at $(1,200)$ kept as pre.38's text, as history, with its two misstatements marked (the width there is $k_0=11$, not $2^{30}$; $W=61$ is not what lifts the cell); the new
  paragraph ``The width ladder'': escalation on a precision shortfall only, at every $\ell$; a failed safeguard refuses \code{[NotEndomorphism]}; why the read stays certified and no hash of an \code{ok} entry
  and no golden line moves; the six level-$1$ cells that now read; arm~(L) and the oracle check at $(1,48)$ (\file{docs/implementation/slot24/ICP1\_acceptance.md}).
\end{itemize}
The code change it describes: \file{Integral/Descent/DescentHeckeTorsion.cpp} (\code{Invariants}), \file{Integral/Descent/DescentMaps} (\code{MapRes::precision}). No label, number, tag set or marker moves (222 on 86).

\subsection*{v1.23.0-pre.38 --- 2026-10-28 --- hecke-extensions slot 24, the pre.37 referee's fixes (CON)}

\noindent\emph{This version.} \textbf{A pre-release, not owner-released, not re-refereed.} It carries v1.23.0-pre.37, refereed PASS-with-fixes (package CON, P7~(a),
\file{docs/implementation/slot24/CON\_pre37\_referee.md}, $2$~M / $10$~m / $5$~N / $2$~T, the pre.30 code map included), and applies the fixes, each checked against the code;
\file{docs/implementation/slot24/CON\_pre38\_fixes.md} maps finding to fix. Text only: no code, hash, record or golden line moves. No statement is weakened.
\begin{itemize}\setlength{\itemsep}{0pt}
\item M1, item~7 (d): \code{coefficients} is $c_0,\dots,c_r$, $c_i$ the coefficient of $x^{r-i}$, the leading $1$ first --- an explicit exception to item~2's low-to-high rule; the in-band key is an open code item.
\item M2: item~7's code map restated as built, audited (\file{D6\_code\_audit.md}), fixed (\texttt{19a969de}), merged (\texttt{2aa61397}), golden R1 landed (\texttt{6958a5d2}); the console's prime bound
  ($\le6$ primes, each $\le97$, $\sum p\le100$); arm~(B); the four code items of v1.23.0-pre.37 closed; (f) met; (j) in the present tense, its interim reading rule vestigial; the status-table row of
  \S\ref{sec:descenthecke}, the pointer in \cl{thm:descentcert}(7)(c) and \S\ref{sec:descenthecke}'s code map updated.
\item m1 (a): the two refusal mechanisms (\code{CoefficientRingK} at construction; \code{GateVerbsOnField}'s opt-in list, kept by arm~(F)) and the residual risk. m2 (b): the record-level fields and the pair's \code{bulk}
  object, one complex reported twice. m3 (e): indexed by \code{torsion\_primes} ($n$ odd: $\ell=2$ only). m4 (e): the three conditions of the $\ker/\coker$ read, record content. m5 (g): the groups carry no refusal of their
  own; a non-refused entry may carry a refusal text. m6 (h): the hashed / unhashed partition exhaustive, \code{DescentRecordZ::body()} normative for the order, \code{commute} unhashed. m7 (d): the digest's
  byte layout (every field one $64$-bit word, least significant byte first; the sign as the word $1+[c_i<0]$). m8 (h): the $(1,12)$ hashes illustrative, not pinned.
\item m9: the status-table row of \S\ref{sec:isomodeltree} brought up to its code map; ``the tree bulk does not move D4's wall'' put at v1.23.0-pre.30. m10: \cl{def:ringsemantics}(3) at level $N$, and (4)'s ``waits for the
  descent assembly'' marked built.
\item N1 (d): the undetermined case is not reachable from the record. N2 (g): a record built without its Hecke or pair part is not a conforming \code{v1} record. N3 (d): $q=0$ omitted by choice. N4: four oracle
  tables. T1: the JSON key is \code{digest}. T2: [IB]'s OD-9~(a), OD-10~(a), recorded by [X] S8-6d. N5 (a source comment) is outside [C].
\item One referee claim not taken: m7's ``fixed widths'' (u32 for $r$, $i$ and the residues, u64 for the count) --- \code{engine::FnvHasher::feed} takes one \code{uint64\_t}, so every field enters as an $8$-byte word.
\item No label is removed, renamed or added; no number, tag set or marker moves ($222$ on $86$ labels).
\end{itemize}

\subsection*{v1.23.0-pre.37 --- 2026-10-26 --- hecke-extensions slot 24, the integral descent record (D6)}

\noindent\emph{This version.} \textbf{A pre-release, not owner-released, not refereed.} It carries v1.23.0-pre.36 and gives the record of D6 (branch \texttt{s24-d6impl}, tip \texttt{538cf881},
\file{docs/implementation/slot24/DESCENT6\_report.md}; pending audit and merge) its contract home. No statement is weakened.
\begin{itemize}\setlength{\itemsep}{0pt}
\item New item~7 of \S\ref{sec:recordschema}, \code{integral\_descent\_record/v1}: (a) ring $\Z$, object group, $\Z$ refused by name elsewhere, \code{Z/<m>} reserved; (b) $q_{\max}=3$ derived from
  \cl{thm:descentcert}(6) (checked: degrees $0$, $1$ and one period $2$, $3$ determine every degree, and $3$ is the least such), \code{truncated = true}, \code{periodic\_from = 2}; (c) groups, bad primes, $K$ per
  prime as provenance; (d) the \code{char\_Z} digest (FNV over $c_i\bmod2^{64}$, $\bmod3^{40}$, or the residues when undetermined): choice-free when determined, a fingerprint not a certificate; (e) the torsion Hecke
  invariants, \code{module\_certified = false} not a refusal; (f) the pair, each complex $q_{\max}=4$, the compact object's list of names required by item~5 and missing in \code{v1} (code item); (g) refusals:
  a refused record is never a valid empty answer; (h) the hashed and unhashed fields, justified by \cl{thm:descentcert}'s choice-independence, \cl{rmk:descentcertscope} and the read-only, value-preserving
  rollbacks \code{HECKE\_DESCENT\_TREE=0}, \code{HECKE\_DESCENT\_TREE\_PAIR=0}; (i) \cl{thm:descentcert}(4), (5) not recorded and not built; (j) the refusal at $(1,200)$, $\ell=2$.
\item The refusal name: \code{[NotEndomorphism] \dots\ (30 < 34)} at $(1,200)$ is a precision shortfall of the $\ker/\coker$ read at the ring width $2^{30}$ ($30<34$ counts certified pivots); the taxonomy
  misnamed it; the name is \code{[KerCokerPrecision]}, \code{[NotEndomorphism]} kept for a failed endomorphism safeguard. A code item; moves no hash and no golden line.
\item \cl{def:ringsemantics}(3): $\Z$ ``(reserved)'' becomes un-reserved for the integral descent only. \cl{thm:descentcert}(7)(c), the status table's row for \S\ref{sec:descenthecke} and \S\ref{sec:descenthecke}'s code map
  point to item~7.
\item No label is removed, renamed or added; no number, tag set or marker moves ($222$ on $86$ labels).
\end{itemize}

\subsection*{v1.23.0-pre.36 --- 2026-10-26 --- hecke-extensions slot 24, the C-D8b referee's fixes; D5c merged}

\noindent\emph{This version.} \textbf{A pre-release, not owner-released, not re-refereed.} It carries v1.23.0-pre.35, refereed ACCEPT after fixes (0 M / 1 m / 7 N,
\file{docs/implementation/slot24/C-D8b\_referee.md}, with the author's response), and records D5c merged into \texttt{hecke-extensions} at \texttt{4064e164}. No statement is weakened.
\begin{itemize}\setlength{\itemsep}{0pt}
\item m1, code map of \S\ref{sec:treepair}: the certificates restated for the merged code --- \code{certify\_transport} compares the blocks of the \textbf{applied} operator ($T'_{\mathcal T}$, $\eta_{\mathcal T}$,
  $T'_\partial$, the strict block) with separately assembled ones and checks the identity of (4)(a) on those values; entrywise for $D\le1300$; above, on the generators and $64$ random vectors uniform in $\Z/\ell^{K}$,
  seeded by $(\ell,N,n,j)$, a miss $\le\ell^{-64}$ per check. What is deterministic and what probabilistic; admissible because no value is read from a check and the read's correctness is the deterministic
  (4)(e)(iii) with the chain-level certificates. ``Not new coverage'' removed.
\item N1: \cl{prop:treepair}(4)(e)(iv) and its proof: commuting with $\Dcone$ \emph{gives} the coboundary half of (iii); the code map says that above $D=1300$ it is not checked deterministically at cochain level.
\item N2: (4)(e)'s hypothesis admits the truncation of $T$ to cochain degrees $\le m$ (the pair lift, built to degree $6$, is one). N3: the proof of (e)(i) for projective, not free, sources.
  N4: (e)(ii) refined (bulk in chain degree $j+1$, boundary and homotopies in $j$) and extended to $H^{q}$, $H^{q}(\partial)$, $q\le m$. N5: the identity on the generators' $\alpha$-parts is the second component of
  (ii) modulo $\ell^{K-\sigma_{\max}}$. N6: \cl{def:treepair}(2), ``computed from these cochains''; the chain-level constructions use $\mathsf a^{(\ell)}$.
\item N7, code map: merged at \texttt{4064e164} (audit \file{D5c\_code\_audit.md}, its N5 deferred and recorded); obligation (c) every $2$- and $3$-adic pair piece, every $N$, including $3\in$ \code{extra\_primes} with
  $n$ odd; the skipped top degree at $q_{\max}=5$ counted (\code{eta\_skipped}) with the (Q5) arm; (f) met by the production certificate; (g) met by \code{transport\_bytes}. The status table and the remarks of
  \S\ref{sec:isomodeltree} and \S\ref{sec:descenthecke} that said ``pending'' updated; D5c's two additions to D3's code (degree-$5$ comparison, \code{IsoHeckeN::apply\_boundary}) recorded there.
\item Checks re-run: \file{cd8\_check.py} and the referee's \file{pre35referee/pre35\_check.py}, both OK (lines in the author's response).
\item No label is removed or renamed; no number, tag set or marker moves ($222$ on $86$ labels).
\end{itemize}

\subsection*{v1.23.0-pre.35 --- 2026-10-26 --- hecke-extensions slot 24, \S\ref{sec:treepair} in line with the D5c code}

\noindent\emph{This version.} \textbf{A pre-release, not owner-released, not refereed.} It carries v1.23.0-pre.34 and changes \S\ref{sec:treepair} (and one code-map note of
\S\ref{sec:descenthecke}) after the D5c implementation (branch \texttt{s24-d5cimpl}, tip \texttt{e720f11c}, \file{docs/implementation/slot24/DESCENT5c\_report.md} \S6; pending audit and merge). No statement is weakened.
\begin{itemize}\setlength{\itemsep}{0pt}
\item \cl{def:treepair}(2): at $\ell\in\{0,2\}$ $S$ and $S^{-1}$ give different chain maps but, the difference being $(1-\mathbf c)Se_x$, \textbf{the same cochains} on every module on which $\mathbf c$ acts trivially;
  $\rho_{M_N}(\mathbf c)=(-1)^n$, so for $n$ even at every $\ell$, with proof; for $n$ odd the difference $2Se_x$ vanishes modulo $2$. The code map's (f) and the check paragraph restate the negative control:
  at chain level, or on cochains at $\ell=2$ with $n$ odd (modulo $2^{K}$, $K\ge2$); asserted clean for $n$ even and at $\ell\ge3$ --- as the code pins it ($(1,5)$, $(3,1)$ caught; $(1,4)$ clean).
\item New \cl{prop:treepair}(4)(e), with proof: comparisons given in chain degrees $\le m$ (homotopies $\le m-1$) extend (the comparison theorem; the targets are acyclic in positive degrees), and the maps in
  cone degrees $\le m-1$ are those of every extension, so the read of $\Hcpt{q}$, $q\le m$, needs nothing in degree $m+1$; the identity of (4)(a) in cone degree $j$ needs chain degree $j+2$, is formable iff
  $j\le m-2$, and is not used by the read. With the code's $m=5$, \textbf{$q_{\max}=5$ stays admissible}; (6)'s refusal ``$q_{\max}\notin\{4,5\}$'' is unchanged and points to (4)(e). The alternative, restricting the
  transport to $q_{\max}=4$, is not needed.
\item Code map: built on \texttt{s24-d5cimpl}, pending audit; files, the default route of the $2$- and $3$-adic pair pieces, the rollbacks \code{HECKE\_DESCENT\_TREE\_PAIR=0}, \code{HECKE\_DESCENT\_TREE=0}, the two
  validations, the reach ($D\approx1060$--$1240$ to the refusal at $D\approx6340$, $\approx8000$ at $40$~GiB; the next wall D2's dense Smith form with transforms). (e)'s certificates restated per cone degree
  ((i) strict block, (ii) cocycles on the generators in every read degree, (iii) the identity of (4)(a) in cone degrees $\le\min(q_{\max}-1,3)$, entrywise for $D\le1300$, on the generators above); ``on the
  generators'' is deterministic, partial and not probabilistic, and admissible under (5) because no value is read from it. The $\ell$-local augmentations of \cl{def:treepair}(4) are the code's.
  \S\ref{sec:descenthecke}'s open code item N7 (\code{D3Refusal::HeldOutMismatch}) closed on the branch.
\item Check: \file{docs/implementation/slot24/cd8\_check.py} gains arms (M3) and (Q); its line is recorded in the check paragraph.
\item No label is removed or renamed; no number, tag set or marker moves ($222$ on $86$ labels).
\end{itemize}

\subsection*{v1.23.0-pre.34 --- 2026-10-26 --- hecke-extensions slot 24, the C-D8 referee's fixes}

\noindent\emph{This version.} \textbf{A pre-release, not owner-released, not re-refereed.} It carries v1.23.0-pre.33, refereed ACCEPT after fixes (0 M / 4 m / 8 N,
\file{docs/implementation/slot24/C-D8\_referee.md}, independent check \file{cd8referee/}), and changes \S\ref{sec:treepair} (and one code-map note of \S\ref{sec:descenthecke}). No statement is weakened except that
``$\eta_{\mathcal T}\ne0$ in general'' is now an observation of the checks, not a proved assertion.
\begin{itemize}\setlength{\itemsep}{0pt}
\item \textbf{m1} \cl{prop:treepair}(4)(a), its proof, \cl{rmk:treepairscope}(iii) and the subsection's introduction: on the pair model $(\mathrm{Cyl},\mathsf B)$ the transported datum $T^{\mathrm{Cyl}}$ is strict ($\eta=0$; the plan's
  sentence holds there); on $\Cone{\mathsf a^{\ast}}$, not the cone of a restriction to a subcomplex, $T^{\mathrm c}_{\mathcal T}\simeq\varpi T^{\mathrm{Cyl}}\varsigma'$ acquires $\eta_{\mathcal T}$ (homotopy proved);
  ``$\eta_{\mathcal T}\ne0$, the diagonal not a cochain map'' is observed ($\theta^{\ast}$; the referee's $T_\varepsilon$), not proved; ``forced'' and ``no strict pair datum exists on the tree'' are gone.
\item \textbf{m2} \cl{prop:treepair}(3)(b): $\mathsf h_p$, (L), (L$'$) for any column augmentation of a complex filtered by finitely many free periodic columns, with the sign rule per column, over
  every $\cring$ after $\otimes\cring$ (coefficients in $\mathsf O_\ell[\Ga]$ allowed), with proof; applied to $\varepsilon_{\mathsf B}$, $\cmpI|_{\Isob}$, $\varepsilon_{\mathcal T}\otimes\cring$, $\cmpI\otimes\cring$. The
  citation ``stated for any such $p$'' is replaced.
\item \textbf{m3} Code map (f) and the check paragraph: $S$ for $S^{-1}$ only at $\ell\in\{0,2\}$; at $\ell=3$ the dropped term $S^{-1}e_U\mathsf v^{U}_1$; $\eta_{\mathcal T}$ dropped only for non-fibre operators,
  detected on a cell per $\ell\in\{2,3\}$, never at $\ell=3$ with $n$ odd; detection counts printed; an arm for $\mathsf j_\Psi\mathsf j_\Phi=\id$ at a cell with $2$-torsion in $\Hcpt{\ge3}$.
\item \textbf{m4} \cl{def:treepair}(4): $\varepsilon^{(\ell)}_{\mathcal T}$, $\varepsilon^{(\ell)}_{\mathsf B}$ defined; \cl{prop:treepair}(1)(b) stated with them; $\Phi^{\mathcal T}_\ell$, $\Psi^{\mathcal T}_\ell$, $\Phi^{\partial}_\ell$,
  $\Psi^{\partial}_\ell$ lie over the identity of $C_{\bullet}(\mathcal T)\otimes\cring$, resp.\ $C_{\bullet}(\mathcal L)\otimes\cring$; the proof of (3)(c) uses them.
\item \textbf{N1--N8} \cl{lem:treepaircyl}: $\mathsf B$ finitely generated projective, zero in negative degrees; \cl{rmk:treepairscope}(i): the path runs from $T\mathfrak v_S$ through $S^{-1}\mathfrak v_U$ to $\mathfrak v_S$;
  (3)(b): the parenthesis on $\Phi^{\partial}_\ell$, scalars $1$; (5)(a): $\Cone{\rres_{\mathrm{Cyl}}}$ not admissible, $\rank C^{0}_{\mathcal T,\ell}(M)$; code map (e), (g): the cost of forming $\eta_{\mathcal T}$;
  (c): the rollback preserves the record's invariants; \S\ref{sec:descenthecke}'s code map: \code{D3Refusal::HeldOutMismatch} an open code item; \cl{def:treepair}(2): $S$ versus $S^{-1}$ immaterial at $\ell\ge3$.
\item No label is removed or renamed; no number, tag set or marker moves ($222$ on $86$ labels).
\end{itemize}

\subsection*{v1.23.0-pre.33 --- 2026-10-26 --- hecke-extensions slot 24, C-D8: the tree cusp pair (D5c), contract first}

\noindent\emph{This version.} \textbf{A pre-release draft, not owner-released, not refereed} (an adversarial referee follows). It carries v1.23.0-pre.32 and adds \S\ref{sec:treepair}, the plan's clause
C-D8 for milestone D5c, \textbf{stated before its code} (rule~21); the D3 code-map status of \S\ref{sec:descenthecke} records the merge. Source: \file{docs/implementation/slot24/DESCENT\_plan.md} v0.4
\S5, \S2.7 (row C-D8), D5c; \file{DESCENT4\_report.md} \S4 (the $2$-adic pair wall), \file{DESCENT5\_report.md}; the referee reports \file{C-D6\_referee.md}, \file{C-D7\_referee.md}. Check
\file{docs/implementation/slot24/cd8\_check.py} (output \file{cd8\_check.out}).
\begin{itemize}\setlength{\itemsep}{0pt}
\item \PNEW{} \S\ref{sec:treepair}. \cl{def:treepair} (untagged): the cusp model $\mathsf B^{(\ell)}_{\mathcal T}$ (the cone of $r_{(1-T)e_{\mathfrak e}}$ on the tree's edge column; ranks $1,2,2,\dots$ at $\ell=2$,
  $1,1$ at $\ell\ge3$), the axis map $\mathsf a^{(\ell)}\colon\mathsf B^{(\ell)}_{\mathcal T}\to\Iso^{(\ell)}_{\mathcal T}$ in closed form ($\mathfrak e_\infty\mapsto(S^{-1}-T)\mathfrak e$ plus periodic column terms), the tree cone
  $\Cone{\mathsf a^{(\ell)\ast}}$, and the geometric tables $\mathsf i_\partial$, $\mathsf r_\partial$, $\mathsf G_\partial$, $\mathsf f_{\mathsf a}$, $\mathsf g_\square$, $\mathsf g'$ on the line complex $C_{\bullet}(\mathcal L)$.
\item \PNEW{} \cl{lem:treepaircyl} (\PROVED, \BSTD): the mapping cylinder $(\mathrm{Cyl}(\mathsf a),\mathsf B)$ is a pair model in the sense of \cl{def:descentpair}; its cone deforms onto
  $\Cone{\mathsf a^{\ast}}$ by explicit $\varpi$, $\varsigma$, $\varsigma'$, $\mathsf h_{\mathrm{Cyl}}$, strictly compatible with $\coneinc$, $\coneproj$; so \cl{prop:descentpair}(1) holds on the tree cone with its signs.
\item \PNEW{} \cl{prop:treepair} (\PROVED, \BSTD): what it computes in every degree over every $\mathsf O_\ell$-algebra; the cochains in closed form at $\ell=2$, $3$, $\ge5$ (at $\ell\ge5$ the
  modular-symbol complex, $\Hcpt{2}$ the coinvariants); effective comparisons with $(\Iso,\Isob)$ --- $\Phi^{\Z}_\partial$, $\Psi^{\Z}_\partial$ by (L), the square homotopies $\mathsf y_\ell$, $\mathsf y'_\ell$ by
  (L$'$), the cone maps $\mathsf j_\Phi$, $\mathsf j_\Psi$, mutually inverse on cohomology through comparisons of pairs; $T_p$ by transport,
  $T^{\mathrm c}_{\mathcal T}=\mathsf j_\Psi\Phi^{(N)}_p\mathsf j_\Phi$, with a \textbf{non-zero, determined off-diagonal $\eta_{\mathcal T}$}; the precision rule; refusal (RSS, precision, the pair lift's codes,
  $q_{\max}$, the cap) and failure (a failed certificate of a constructed map) and scope.
\item \PNEW{} \cl{rmk:treepairscope} (\PROVED): where the plan's \S5 is corrected --- the minimal cusp resolution, the tree cone in place of the cylinder's, $\eta_{\mathcal T}\ne0$ (``no $\eta$ on the tree
  side'' is withdrawn), the comparisons of pairs as proof devices, cost a prediction.
\item \S\ref{sec:descenthecke}'s code map: D3 \textbf{audited and merged} at \texttt{dc8261a3}; the $\ell\ge5$ cokernel torsion read refuses \code{[RssCeiling]} before building;
  \code{HeldOutMismatch} and \code{[ResidueRank]} throw \code{hecke::HeckeInternalError} (failures, not refusals). Pointers to \S\ref{sec:treepair} in \S\ref{sec:isomodeltree} (scope, code map),
  \S\ref{sec:descentpair} (introduction) and \cl{rmk:integralscope}(5); the code-map table gains a row and records the D3 merge.
\item Four statements added, the last of their section; no label is removed or renamed; no existing number or tag set moves; no \code{\textbackslash LEAN} marker is added ($222$ on $86$ labels).
\end{itemize}

\subsection*{v1.23.0-pre.32 --- 2026-10-26 --- hecke-extensions slot 24, the C-D3h referee's fixes}

\noindent\emph{This version.} \textbf{A pre-release, not owner-released, not re-refereed.} It carries v1.23.0-pre.31, refereed ACCEPT after fixes (0 M / 3 m / 6 N,
\file{docs/implementation/slot24/C-D3h\_referee.md}, check \file{pre31referee/pre31\_check.py}), and changes \S\ref{sec:descenthecke} only. No statement is weakened.
\begin{itemize}\setlength{\itemsep}{0pt}
\item \textbf{m1, N4} Proposition~\ref{prop:strictliftN}(2), independence of the lift: the coboundary of $\CI$ is unsigned (Lemma~\ref{lem:cosettransport}(2)), so $\mathsf H=\Theta_{h_\Iso}$ in every degree;
  ``$\pm$'' and ``re-signed degree by degree'' are gone; $\mathsf H$ is ``constructible'', not ``explicit''.
\item \textbf{m2} Theorem~\ref{thm:descentcert}(4) at $\ell\ge5$: with $k=K'$, $K>\sigma_{\max}+v_\ell(\chi^{(N)}_{p,1}(c))$ suffices ($\sigma_{\max}$ of $\Delta$'s certified form, $0$ without $\ell$-torsion), with
  its proof; a short read is refused, never wrong.
\item \textbf{m3} Proposition~\ref{prop:torsionhecke}(3$'$): any exact lifts $x_j\equiv\tilde x_j\pmod{\ell^{K}}$ form a basis and give the same $X$ modulo $\ell^{K}$ ($P=I+\ell^{K}M$), proved.
\item \textbf{N1--N3, N5, N6} (3$'$)'s independence per certifying read; $U=ST$ in Theorem~\ref{thm:descentcert}'s proof; the free part's $K'$ is sharp; the code map notes that
  \code{certify()} at $q\le2$ is all the clause requires and that $\mathrm{Rel}_N$-stability is a theorem; ``What is used'' completed.
\item No label is removed or renamed; no number, tag set or marker moves ($222$ on $86$ labels).
\end{itemize}

\subsection*{v1.23.0-pre.31 --- 2026-10-26 --- hecke-extensions slot 24, D3: the cokernel presentation, obligation (b) through the lift, any admissible lift}

\noindent\emph{This version.} \textbf{A pre-release, not owner-released, not refereed} (an adversarial referee follows). It carries v1.23.0-pre.30 and changes \S\ref{sec:descenthecke} only. Source:
\S6 of \file{docs/implementation/slot24/DESCENT3\_report.md} (D3, branch \texttt{s24-d3impl}, tip \texttt{a1375159}, not merged), and the D3 code (\file{Integral/Descent/DescentHecke*}). No statement is weakened;
three are widened to what the code reads, each with its proof.
\begin{itemize}\setlength{\itemsep}{0pt}
\item \textbf{The cokernel presentation} (D3's note~1). New sub-item Proposition~\ref{prop:torsionhecke}(3$'$): the read of~(3) on $Q=\coker\Delta$ with an endomorphism $F$ preserving
  $\operatorname{im}\Delta$, verbatim (proof: the two-term complex, with $F^{q-1}$ chosen by freeness and never read); the free part $Q/\mathrm{tors}$ is read as well, modulo $\ell^{K'}$
  (the ``moreover'' of Proposition~\ref{prop:descentmaps}(2)(g)(i)), as the input of~(4); the invariants do not depend on the presentation. Proposition~\ref{prop:crtcharpolyN}(2) admits
  $[H_{N,p}]$ on $M_{N,\F_{\ell'}}/\mathrm{Rel}_N$ (Proposition~\ref{prop:heilbronn}(1), (2) through $\varkappa_N$), Theorem~\ref{thm:descentcert}(3), (4) admit it on $M_{N,\Z_{(\ell)}}/\mathrm{Rel}_N$ at
  $\ell\ge5$, with the certifying rank $m_N(n+1)-b_1(M_N)$; the proofs say why ($6$ a unit, $\varkappa_N$ an isomorphism intertwining $[H_{N,p}]$ with $T_p$, flat base change).
\item \textbf{Obligation (b) through the lift} (D3's note~2). Proposition~\ref{prop:strictliftN}(2): an integral table and a $\tau_{\Iso}$ certified by (T1)--(T3) and the strict identity on
  \emph{every} basis element ($\cmpI\tau_{\Iso}(s_je_c)=0$ at $t\ge1$ included) make the table the table of a $g$-twisted chain datum over $\Z$; proof: $\cmpI\tau_{\Iso}$ kills $\ker\cmpI$ (its
  $t=0$ part is $d^{(0)}$ of $t=1$, by flatness of $\Z[\Ga]$ over $\Z[\Ga_\sigma]$), so it factors through $\cmpI$, and the factor inherits (T1)--(T3). A separate $\tau^{s}$ check on $\Xbulk$ is
  therefore not needed; the code obligation's ``before the lift'' reads ``before the lift is used''.
\item \textbf{Any admissible lift} (D3's note~3). Proposition~\ref{prop:strictliftN}(2): ``It is produced \dots\ with no boundary step'' describes one construction; \emph{every} additive $\tau_{\Iso}$ with
  (T1)--(T3) and the strict identity (fixed independently of $r_t$, $N$ where (3)(c), (7) need it) is admissible, the bulk component of D4's pair lift among them; (3)(c) records the pair lift as fixed.
  \emph{Independence of the lift}: two admissible lifts are homotopic through (T1)-maps (the fundamental lemma over $\Z[\Ga''_g]$), and the operators differ by $d_{\mathrm I}\mathsf H+\mathsf Hd_{\mathrm I}$ with $\mathsf H$
  the formula \eqref{eq:levelcorr} applied to the homotopy; so the record's reads coincide, and for a fixed Smith form even $\tilde\Phi$ where it is exact.
\item \textbf{Code map}: a status paragraph --- D3 \textbf{built on branch \texttt{s24-d3impl}, pending merge and audit}: the files, the two validations, the transport at $2$ and $3$ for every $N$,
  the presentation route at $\ell\ge5$, obligations (a)--(d); \textbf{not built}: Theorem~\ref{thm:descentcert}(4) at $\ell\in\{2,3\}$, and (5). The code-map table row and the ``What is used'' list follow.
\item No label is removed or renamed; the new sub-item (3$'$) has no label of its own. No number, tag set or marker moves ($222$ on $86$ labels).
\end{itemize}

\subsection*{v1.23.0-pre.30 --- 2026-10-26 --- hecke-extensions slot 24, D5 merged: the code map of the tree models; \texttt{HeckeInternalError}}

\noindent\emph{This version.} \textbf{A pre-release, not owner-released, not re-refereed}; non-mathematical. It carries v1.23.0-pre.29 and changes only code-map text; no statement, label, number, tag set or marker moves
($222$ on $86$ labels). Sources: \file{docs/implementation/slot24/DESCENT5\_report.md} (with its \S8) and the code; D5 merged at \texttt{5eb7e9de}.
\begin{itemize}\setlength{\itemsep}{0pt}
\item \S\ref{sec:isomodeltree}'s code map (and the code-map table row): ``Not built'' replaced by \textbf{built and merged} --- the files, the two validations, the tree as the default route for the dense $\ell$-local bulk
  pieces of D0, D1, D2 with rollback \code{HECKE\_DESCENT\_TREE=0}, Hecke by transport as a hook whose real arm waits for D3, the tree cusp pair (D5c) not built. The other descent code maps were checked: D0--D2 and D4 already
  record what is built; the remaining ``not built'' items (D3's level-$N$ operators, (7)(a)--(b), $\ell\ge5$ ``raise $K$'') are still true.
\item Code map (f) of \S\ref{sec:descentpair} and the pre.29 changelog's ``known width'': \code{hecke::HeckeInternalError} (\texttt{bc128cb7}) --- genuine internal errors propagate as failures; only a real overflow refuses as \code{IntegerCap}.
\end{itemize}

\subsection*{v1.23.0-pre.29 --- 2026-10-25 --- hecke-extensions slot 24, D4: the C-D7p referee's fixes M1, m1--m4, N1; the code map for the merged D4}

\noindent\emph{This version.} \textbf{A pre-release, not owner-released}; it carries v1.23.0-pre.28 (refereed ACCEPT-with-fixes, \file{docs/implementation/slot24/C-D7p\_referee.md}: one
binding clause wrong, four minors, three notes; the referee's independent check \file{pre28referee/pre28\_pairlift\_check.py}, $132$ checks, $0$ failing, rerun at this version) with every
finding applied; \textbf{not re-refereed}. Only \S\ref{sec:descentpair} and two historical parentheses change; no label, number, tag set or marker moves ($222$ on $86$ labels). The code side
of the report (\texttt{6daf1e03}) and the D4 code audit's fixes (\texttt{46fc0c73}) are merged with D4 at \texttt{3991a2f1}; source of the code's refusal list:
\file{docs/implementation/slot24/DESCENT4\_report.md} \S11, checked against \file{Integral/Descent/DescentPair.h}, \file{DescentPairHecke.cpp}.
\begin{itemize}\setlength{\itemsep}{0pt}
\item \textbf{M1} (``refused exactly when'' did not match the code): \cl{prop:descentpairhecke}(1) now defines the cap as the code's checked bounds --- $64$-bit signed exact products,
  $32$-bit signed $\tau$-table coefficients, the coset tables' $32$-bit unsigned domain, a failed self-certification of an input table --- plus the RSS ceiling ($4$\,GiB, rule~19, before
  building), and says that every other exception is a failure, not a refusal; (5) lists the eleven refusal conditions (a)--(k), the code's \code{PairRefusal} codes \code{NotPrime} \dots\
  \code{StrictIdentity}, the strict identity and $\Isob$-closure on the layout included, and the cap during the operator's assembly as well; ``in the last two cases'' became ``in cases
  (d)--(k)''. The code map's (f) names the codes and the classification (\code{ClassifyPairException}); \cl{rmk:descentpairscope}(iii) and the two historical parentheses
  (\cl{rmk:strictliftbuilt}, \cl{rmk:integralscope}(2)) point to (5). Recorded as a known width, not hidden: the exact-integer layer's overflow type is also thrown by some internal guards
  of the integer geometry on the lift's path, which the code therefore refuses as \code{IntegerCap} instead of failing.
\item \textbf{m1}: the proof of (1) carries the sufficiency argument --- (i) a twisted map on $H_0(\Isob)$ by (T1), (T2) and $\Isob$-preservation, (ii) the parabolic $g^{-1}P_c^{\,k}g$ forces
  a multiple of $[g^{-1}c]$, (iii) \cl{def:descentpair}'s augmentation compatibility and (T3) give $1$; ``Homotopy of pairs'' cites it; the route through $\cmpI$ and \cl{lem:bdrycomponent}
  (which uses (B-$\varepsilon$)) kept as a second proof.
\item \textbf{m2}: the pair is served for $q_{\max}\le5$; $q_{\mathrm{top}}>6$ (and $q_{\max}>5$) is refused; higher degrees come from (7), never from extending the lift; the code
  map's (g) reads $4\le q_{\max}\le5$. \textbf{m3}: an overflow during the operator's assembly refuses (now true in the code). \textbf{m4}: the code map's new (i) lists the
  negative controls, the homotopy faults $\tau+dK+Kd$ (\code{StrictLift::with\_homotopy}) refused by (S) and the boundary items alone among them.
\item \textbf{N1}: (1) says the measured sizes belong to the shipped model (on the referee's model, $\approx\times3$ per two degrees), which is why (7), not a longer lift, serves the degrees
  above $6$. N2, N3: no change needed (N3's separate assembly of $\widetilde T_{\Isob}$ is merged and named in (f)).
\item The code map: D4 \textbf{merged} at \texttt{3991a2f1} (no longer ``not merged''), \file{Integral/Model/StrictLift.\{h,cpp\}}'s validation hook, \file{cmake/slot24/DESCENT4.cmake},
  \file{Registry/Slot24/DESCENT4.cpp} added; the code-map table row likewise.
\end{itemize}

\subsection*{v1.23.0-pre.28 --- 2026-10-25 --- hecke-extensions slot 24, D4: $T_p$ on the boundary and the cone for every prime $p\nmid N$ (owner scope addition)}

\noindent\emph{This version.} \textbf{A pre-release, not owner-released, not refereed} (C-D7 refereed v1.23.0-pre.26, its fixes at pre.27); it carries v1.23.0-pre.27 and amends
\S\ref{sec:descentpair} and two historical remarks to the owner's scope addition of 2026-10-25: D4 computes $T_p$ on the boundary and the cone for \emph{general} $p\nmid N$, at level $1$
and level $N$, from the strict pair lift built on demand for each requested $p$, certified and cached. No label, number or marker moves ($222$ on $86$ labels); no tag set moves. Source:
\file{docs/implementation/slot24/DESCENT4\_report.md} \S\S4--6.
\begin{itemize}\setlength{\itemsep}{0pt}
\item \cl{prop:descentpairhecke}(1): ``Built and certified: $p\in\{2,3,5\}$ \dots; for every other $p$ \dots\ not built'' replaced: built on demand for every prime $p$, degrees $\le6$,
  certified by \cl{rmk:strictliftbuilt}(2)'s checks with the boundary items before first use, cached; refused on a failed certificate or at the cap ($64$-bit signed products); why the
  certificate suffices on its own (the component map on $H_0(\Isob)$ forced by (T1), (T3)); measured sizes and costs to $p=127$ as evidence, no bound. Proof of (1): the boundary step on $\Isob$
  spelled out; termination and integrality carry no restriction on $p$ --- no step was verified only at $p\in\{2,3,5\}$; only the implementation was, and the run-time certificate covers it.
\item \cl{prop:descentpairhecke}(5): certified at every $\ell$ and every $N$ for every prime $p\nmid N$ for which a pair lift has been built and its certificate passed (at run time), the strict
  identity $\rres\widetilde T_{\Iso}=\widetilde T_{\Isob}\rres$ checked entrywise in every degree used; refused exactly for $p$ not prime, $p\mid N$, a failed certificate or the cap (\emph{incomplete}: corrected at v1.23.0-pre.29, referee M1, to the code's eleven refusal codes); pre.27's
  ``in every other case (level $N$ at $p\notin\{2,3,5\}$) \dots\ refused'' restricted to the last two.
\item \cl{rmk:descentpairscope}(iii) updated; the code map: \textbf{built} on branch \texttt{s24-d4impl} (\texttt{d82c7761}, not merged) --- \file{Integral/Descent/DescentPair.\{h,cpp\}},
  \file{DescentPairHecke.cpp}, \code{IntegralDescentPairValidation} (+\code{Slow}) --- and (f) restated for every $p\nmid N$ (no longer ``not a D4 obligation''); the code-map table row likewise.
\item Historical, kept: \cl{rmk:strictliftbuilt} (slot~10's $p\in\{2,3,5\}$) gains a ``later builds'' sentence; \cl{rmk:integralscope}(2) a v1.23.0-pre.28 parenthesis.
\end{itemize}

\subsection*{v1.23.0-pre.27 --- 2026-10-25 --- hecke-extensions slot 24, D4: the C-D7 referee's fixes M1--M3, m1--m5}

\noindent\emph{This version.} \textbf{A pre-release, not owner-released}; it carries v1.23.0-pre.26 (refereed ACCEPT-with-fixes,
\file{docs/implementation/slot24/C-D7\_referee.md}: three fixes that matter, five minors, two notes) with every finding applied; \textbf{not re-refereed}. Only
\S\ref{sec:descentpair} changes; no label, number or marker moves ($222$ on $86$ labels); one item, \cl{prop:descentpair}(5)(e), is added inside an existing statement. No code exists yet
for D4; the code map's obligations (c) and (f) are amended before it.
\begin{itemize}\setlength{\itemsep}{0pt}
\item \textbf{M1} the check paragraph: ``on these cells $\Hcpt{2}$ has no torsion at $\ell\ge5$'' was false by the script's own output ($\Z/4\oplus\Z/60$ at $(1,8)$,
  $\Z/2\oplus\Z/12\oplus\Z/1260$ at $(1,12)$); replaced by (5)(e) and the cells' $\ell\ge5$ parts; the Eisenstein reading withdrawn.
\item \textbf{M2} \cl{prop:descentpair} new (5)(e): $\Hcpt{2}\otimes\Z_{(\ell)}\cong H_0(\Ga;M)\otimes\Z_{(\ell)}$ for $\ell\ge5$, with proof (Manin's resolution, the norms of the stabilisers,
  generation by $S$, $U$); for $M_N$ it is $\Vcoef_A$, zero for $\ell>n+1$ (Borel coinvariants, Nakayama). (6)(c): $\Dcone^{0}$ by the Betti recursion \emph{or} a prime $\ell'\ge5$ with
  $\ell'\nmid\mathrm{tors}\,H_0$ --- for $M_N$ every $\ell'\ge5$, $\ell'>n+1$, the bulk's rule; the a-priori cap on $\Dcone^{0}$ at $\ell\ge5$; \cl{rmk:descentpairscope}(ii) and the code map (c).
\item \textbf{M3} \cl{prop:descentpairhecke}(5) \emph{weakened} (the referee's first option): $T_p$ on the boundary and the cone is certified, at every $\ell$, exactly at the $p$ with a built
  pair lift ($p\in\{2,3,5\}$), and at $\ell\ge5$, level $1$, where the shipped cone operator is built; at level $N$ and $p\notin\{2,3,5\}$, $\ell\ge5$ included, it exists and is determined but is not
  computed, and the code refuses it; the code map (f) says it is not a D4 obligation; \cl{rmk:descentpairscope}(iii).
\item \textbf{m1} locators: Brown VII.6 withdrawn; (4)'s periodicity by the tensor resolution $F\otimes F'$ and the shift of the periodic resolution; (5)(d) by the coefficient sequence
  $0\to M\to\Hom(\Z[\mathbb{P}^1],M)\to\Hom(\mathrm{St},M)\to0$ (note n1), with the Shapiro--cup compatibility for $\mathrm{Coind}$ stated once; pre.26's dangling
  ``\cl{lem:bulkperiodic}(4)'' removed; III.9 read III.8--9.
\item \textbf{m2} the cone is unbounded at $3$ exactly when $H^{q}(A;\Vcoef)_{(3)}\ne0$ for some $q\ge2$ (not in general). \textbf{m3} (4): order $n!$ for $n$ even;
  $2^{\lfloor n/2\rfloor+1}$ for $n$ odd. \textbf{m4} (1)(c)'s signs pinned: $\coneinc{}_\ast=j_{\mathrm{St}}^{\ast}$, $\coneproj{}_\ast=-\delta$ under the stated identification, with proof.
  \textbf{m5} (5)(c): exponent $\le12$ observed, not asserted.
\item \file{cd7\_pair\_check.py} gains arm R (M2; $308$ checks): $2528$ checks, $0$ failing; the referee's \file{cd7referee/cd7referee\_check.py} rerun: $2457$ checks, $0$ failing.
\end{itemize}

\subsection*{v1.23.0-pre.26 --- 2026-10-25 --- hecke-extensions slot 24, D4: the boundary pair and the cone over $\Z$ (clause C-D7)}

\noindent\emph{This version.} \textbf{A pre-release draft, not owner-released, not refereed}; it carries v1.23.0-pre.25 plus one subsection, stated \emph{before} D4's code
(contract first, rule~21): clause C-D7 of \file{docs/implementation/slot24/DESCENT\_plan.md} v0.4 (its D4, the \S2.2 row ``Boundary LES'', D-$\delta$6), under the owner's rulings
[R-Q2] (no big integer), [R-Q3] (invariants only) and [R-Q4] ($T_p$, $p\nmid N$; $U_p$ deferred). The tree cusp pair (D5c, C-D8) is not part of it. No marker moves ($222$ on $86$
labels); the new statements are the last of their section, so no statement number moves.
\begin{itemize}\setlength{\itemsep}{0pt}
\item \PNEW{} \S\ref{sec:descentpair}, with a paragraph naming what it uses.
\item \PNEW{} \cl{def:descentpair}: $S$-pair models --- $(\Iso,\Isob)$ over every ring, $(\Iso^{(\ell)},{\Isob}^{(\ell)})$ over $\mathsf O_\ell$, the cellular pair with $6$ inverted --- their
  cochains, cones and $\Hcpt{\bullet}$; comparisons and homotopies of pairs.
\item \PNEW{} \cl{lem:pairlifting} (\PROVED, \STD{} \cite[Thm.~2.2.6]{weibel}): comparisons of pairs exist, boundary first, are unique up to homotopy of pairs, and induce
  quasi-isomorphisms of cones with the cone homotopy $\Upsilon_k$.
\item \PNEW{} \cl{prop:descentpair} (\PROVED, \STD{} Brown, Weibel). (1) What a pair model computes over every ring; $\Hcpt{q}=\Ext^{q-1}(\mathrm{St},M)$, the sequence
  \eqref{eq:descentpairles} the $\Ext$-sequence of $0\to\mathrm{St}\to\Z[\mathbb{P}^1]\to\Z\to0$; $\Hcpt{1}$ torsion-free over $\Z$; ${\Isob}^{(2)}$ free, ${\Isob}^{(\ell)}$ for $\ell\ge3$ the circle model; the
  shipped E8 pair as the $\Z[1/6]$-piece at level $1$. (2) The glued sequence over $\Z$: the pair map $\cmpI$ the only glueing datum, exactness local, maps by the precision rule. (3)
  Level $N$ by Shapiro, block-diagonal by cusps, widths from $\mathbb{P}^1(\Z/N)$. (4) The boundary in closed form ($X^{t}$, $X_t$, $(X/2)^{t}$, $(X/2)_t$), $\abs{F_{n,w}}=n!\,w^{n}$, bounded away
  from $2$, periodic. (5) $\Hcpt{0}=0$, $\dim\Hcpt{2}=\dim(M^{\vee})^{\Ga}$, $\dim\Hcpt{1}=b_1$ (for $M_N(\chi)$), $48\Hcpt{q}=0$ ($q\ge3$), the shapes at $\ell\ge5$, $3$, $2$; periodicity, $q_{\max}\ge4$ for the
  pair. (6) The certified exponent on boundary and cone: admissible rank sources ($\Dcone^{0}$ by the Betti recursion only), a-priori exponents, \ref{open:ell2}'s truncation --- the
  plan's D-$\delta$6.
\item \PNEW{} \cl{prop:descentpairhecke} (\PROVED, \STD{} Brown, Weibel). (1) The pair lift for every $p$: produced, certified by the existing checks, homotopy-unique as pairs;
  built for $p\in\{2,3,5\}$. (2) The level-$N$ operator on $\Isob$, strict with $\rres$, $\hder=0$. (3) What it induces --- $T_p$ of \S\ref{sec:descenthecke} on the bulk, the double-coset
  operator on $\Ext(\mathrm{St},\VR)$ and on the boundary --- choice-free over $\Z$; the sequence $T_p$-equivariant over $\Z$, $\Z_{(2)}$, $\Z_{(3)}$. (4) The pieces and the glue. (5) What is
  certified at which $p$: boundary and cone at $\ell\ge5$ for every $p$, at $\ell\in\{2,3\}$ only with a built pair lift; never the bulk-only lift on the cone. (6) D2's maps on the pair.
  (7) Periodicity, Hecke-equivariantly.
\item \PNEW{} \cl{rmk:descentpairscope} (\PROVED): where the plan's D4 wording is corrected --- C-D7 a clause, D-$\delta$6 sharpened (boundary unbounded only at $2$, cone at $2$ and
  $3$, $q_{\max}\ge4$), ``built for $p\in\{2,3,5\}$'' versus produced for every $p$, arm (Q) at level $1$ only, coherence for the Hecke operator only, level $N$ by cusps.
\item \cl{rmk:integralscope}(4), (5) and \cl{thm:descentcert}(7)(c): pointers to \S\ref{sec:descentpair} (no statement changes).
\item A check outside the build, \file{docs/implementation/slot24/cd7\_pair\_check.py} (output \file{cd7\_pair\_check.out}): $24$ cells, $2220$ checks, $0$ failing; three independent
  complexes (the tree pair with a computed comparison, the Manin resolution of $\mathrm{St}$, the Bredon pair); planted faults each detected.
\item The status appendix ($402\to407$ lines), the label index, the code-map row and the referee pointer.
\end{itemize}

\subsection*{v1.23.0-pre.25 --- 2026-10-25 --- hecke-extensions slot 24, D3: the C-D4 referee's fixes M1, M2, m1--m5}

\noindent\emph{This version.} \textbf{A pre-release, not owner-released}; it carries v1.23.0-pre.24 (refereed ACCEPT-with-fixes,
\file{docs/implementation/slot24/C-D4\_referee.md}: two fixes that matter, five minors, one note, three trivial) with every finding applied; \textbf{not re-refereed}. Only
\S\ref{sec:descenthecke} changes, with one sentence each in \cl{prop:descentmaps}(2)(g)(ii) and \cl{rmk:bredonhecke}; no statement, label, number or marker moves ($222$ on
$86$ labels). No code exists yet for D3; the code map's obligations are stated before it.
\begin{itemize}\setlength{\itemsep}{0pt}
\item \textbf{M1} \cl{prop:torsionhecke}: new (1)(e) --- the layer-$j$ block's characteristic polynomial is choice-free modulo $\ell^{e_j}$, $e_j=\min(j,d_j)$, $d_j$ the distance from
  $j$ to the nearest other layer (a closed path through a layer $a\ne j$ is divisible by $\ell^{\abs{a-j}}$); (2) no longer says ``only modulo $\ell$'': sharpness pairs on
  $\Z/\ell^{j\mp d}\oplus\Z/\ell^{j}$ for every $1\le d<j$ (e.g.\ $\Z/2\oplus\Z/8$: modulo $4$, not $8$); \cl{rmk:descentcertscope}(iii) and \cl{thm:descentcert}(3) follow.
\item \textbf{M2}, resolution (a): the $3$-local piece of \cl{thm:descentcert}(3) is the transport $T'^{(3)}_p$ (or $T'_{\mathcal T,p}$) for \emph{every} $N$;
  \cl{prop:strictliftN}(6)(b) is restated for an integral datum certified by (T1)--(T3) over $\Z$ and is read by no statement; so the subsection's ``Not used'' list and
  \cl{rmk:bredonhecke} (one sentence added) hold as written. Code obligation: no branch on $\nu_3(N)$, the shipped operator read only with $6$ inverted.
\item \textbf{m1} \cl{prop:torsionhecke}(3): $2\sigma_{\max}\le K$ is necessary for the module (a minimal certified example; $33$ real cases in the referee's check), and is not
  \cl{prop:descentmaps}(2)(g)(ii)'s conservative policy (that one reads exact vectors; a sentence there says so); the $P_j$, the Fitting types and the block polynomials modulo
  $\ell^{\min(e_j,K')}$ need only a certified form; \cl{rmk:descentcertscope}(ii).
\item \textbf{m2} \cl{prop:crtcharpolyN}(4): residue primes $<2^{32}$; a word mulmod for Horner modulo $Q>2^{32}$; the negative case $c^{+}\bmod Q-P\bmod Q$.
\item \textbf{m3} \cl{prop:strictliftN}(2): $\tau_X$ is any $g$-twisted datum over $\Z$ --- the canonical $\tau$ once certified, with no boundary datum at all; the procedure is
  uniform in $p$, its input is not. \textbf{m4} (3)(c): $\tau_{\Iso}$ is built once, independent of $r_t$ and $N$. \textbf{m5} \cl{thm:descentcert}(7)(a),
  \cl{rmk:descentcertscope}(i): ``not asserted'' (finite determinacy expected for a squarefree $\chi$, no bound; repeated roots at old forms) in place of ``open mathematics''.
\item \textbf{n1} (2): $\ell'=p$ admissible. \textbf{t1} (3): attained iff $r_N>0$. \textbf{t2} proof of (1): \cl{prop:symconv}(2). \textbf{t3} the ``What is used'' list:
  \cl{prop:levelconehecke}(4) with $\tau_{\Iso}$.
\item The code map's four obligations (M2, m2, m3, m4); the check paragraph names the referee's check and its new arm \code{M}; the provenance lines, the Appendix~\ref{app:status}
  note (no line, no count moves) and its referee pointer.
\end{itemize}

\subsection*{v1.23.0-pre.24 --- 2026-10-25 --- hecke-extensions slot 24, D3: Hecke $T_p$, $p\nmid N$, on the descent (clauses C-D4, C-D5)}

\noindent\emph{This version.} \textbf{A pre-release draft, not owner-released, not refereed}; it carries v1.23.0-pre.23 plus one subsection, stated \emph{before} D3's code
(contract first, rule~21): clauses C-D4 and C-D5 of \file{docs/implementation/slot24/DESCENT\_plan.md} v0.4, under the owner's rulings [R-Q2] (residues and a certified bound),
[R-Q3] (invariants only; integer Hecke matrices are D8), [R-Q4] ($T_p$, $p\nmid N$; every $U_p$ deferred) and [R-GT] (the tree is the $2$-adic route). No marker moves
($222$ on $86$ labels); the new statements are the last of their section, so no statement number moves.
\begin{itemize}\setlength{\itemsep}{0pt}
\item \PNEW{} \S\ref{sec:descenthecke}, with a paragraph naming exactly the Hecke clauses it uses (and those it does not: the Bredon operator at $\ell\in\{2,3\}$,
  \cl{cor:stablecert}, any $\tau$-table on the tree).
\item \PNEW{} \cl{prop:strictliftN} (\PROVED, \STD{} \cite[Ch.~III.8--10]{brown}, \cite[Thm.~2.2.6]{weibel}). (1) $AgA=A\smat p001A$ for $p\nmid N$ (explicit $\gamma_1,\gamma_2$): no
  adjoint question for $T_p$. (2) A bulk strict lift $\tau_{\Iso}$ for every $p$, \emph{constructed} by (L) of \cl{prop:isomodeltree}(4)(b) without a boundary step; no size
  bound. (3) The induced-fibre form \eqref{eq:levelcorr} on $\Iso_{\bullet}$: integral, a cochain map over every ring, inducing $T_p$ in every degree with no hypothesis on~$6$.
  (4) $\cmpI^{\ast}$ intertwines it with the shipped operator on the nose; the glue. (5) Transports to $\Iso^{(\ell)}$ and to the tree models (the route at $2$). (6) The
  cheap routes at $\ell\ge5$ and at $3$ when $\nu_3(N)=0$. (7) The strict square with $\alpha$ for $p\nmid MN$: the level sequence \eqref{eq:descentmapsles} is
  $T_p$-equivariant over $\Z$ and over $\Z_{(2)}$, $\Z_{(3)}$ (what \cl{prop:descentmaps}(3)(c) had left open). (8) $T_p$, $T_{p'}$ commute.
\item \PNEW{} \cl{prop:crtcharpolyN} (\PROVED, \STD{} Diamond--Im, Deligne, Diamond--Shurman for the bound, which is the input). $\chi^{(N)}_{p,q}$ is $1$ for $q\ge2$,
  $x-(p+1)$ or $1$ for $q=0$; $\chi^{(N)}_{p,1}$ from residues on the $\mathsf O$-piece alone (no $2$-, $3$-adic input); $\abs\lambda\le1+p^{\,n+1}$, coefficient bits
  $\mathsf b_{p,i}=r_N+i(n+1)\mathrm{bitlen}(p)$; recovery in $64$-bit words (Garner, sign by digits, residues modulo any word); the output rule [R-Q2]: integer if certified
  and below $2^{63}$, digits and bound if certified and larger, residues and bound if not certified --- never a lifted integer.
\item \PNEW{} \cl{prop:torsionhecke} (\PROVED). For an endomorphism of a finite $\ell$-group: the layer polynomials $P_j\in\F_\ell[x]$, the Fitting components (whose types are
  read off the $P_j$), $\ker$/$\coker$ of $\phi-c$ and the joint module are choice-free; the characteristic polynomial over $\Z/\ell^{j}$ on Smith summands is not (an
  example). The exact read from a certified local Smith form needs $2\sigma_{\max}\le K$. On a free quotient, the invariant factors of $\psi-c$ are certified by the count
  $r_k=r$, reached for $k>v_\ell(\chi(c))$, never at a root.
\item \PNEW{} \cl{thm:descentcert} (\PROVED): the record of C-D5 --- groups, $\chi^{(N)}_{p,q}$, the torsion Hecke modules, the free-quotient invariants at non-roots,
  D2's maps with $\chi^{(N)}_{p,1}\mid\chi^{(MN)}_{p,1}$, every degree by periodicity --- and what is not asserted (the lattice at the roots and the extension class: the
  plan's open D-$\delta$7; tier B; $U_p$, nebentypus, the pair).
\item \PNEW{} \cl{rmk:descentcertscope} (\PROVED): where the plan's wording is corrected --- torsion is not open; the bound is $2\sigma_{\max}\le K$ on torsion and
  $v_\ell(\chi(c))$ on the lattice; ``char poly on Smith summands'' is not defined; $\chi$ needs no $2$-, $3$-adic input; C-D4$'$ not drafted.
\item \cl{rmk:integralscope}(2), (4) and \cl{prop:descentmaps}(3)(c): pointers to \S\ref{sec:descenthecke} (no statement changes). Bibliography: \cite{deligne_weil1}
  (locator not verified).
\item A check outside the build, \file{docs/implementation/slot24/cd4\_hecke\_check.py} (output \file{cd4\_hecke\_check.out}): $21$ cells, $931$ checks, $0$ failing;
  independent oracles ($q$-expansions, point counts on $X_0(11)$, genus-$0$ Eisenstein); planted faults each detected.
\item The status appendix ($397\to402$ lines), the label index, the code-map row and the referee pointer.
\end{itemize}

\subsection*{v1.23.0-pre.23 --- 2026-10-25 --- hecke-extensions slot 24, D2: the C-D3g referee's fixes M1, m1--m5, t1, t2}

\noindent\emph{This version.} \textbf{A pre-release, not owner-released}; it carries v1.23.0-pre.22 (refereed ACCEPT-with-fixes,
\file{docs/implementation/slot24/C-D3g\_referee.md}: no blocker, one medium, five minors, two trivial) with all eight findings applied; \textbf{not re-refereed}. Only
\cl{prop:descentmaps} and, editorially, \cl{prop:isomodeltree} change; no statement, label, number or marker moves ($222$ on $86$ labels). No code value changes: the code
was already doing what the sharpened statement allows (comments in \file{Integral/Descent/DescentMaps} now cite it).
\begin{itemize}\setlength{\itemsep}{0pt}
\item \textbf{M1} (2)(g)(i) ``moreover'': the first $s$ rows of $\tilde L$, the last $r-s$ columns of $\tilde L^{-1}$ and the last $c-s$ rows of $\tilde R^{-1}$ agree with $L$,
  $L^{-1}$, $R^{-1}$ modulo $\ell^{K}$ (read off the factorisations of the Schur-complement proof); (ii) lattice coordinates: for $y\in A\Z_{(\ell)}^{c}$ the coordinates in the
  basis $(\tilde L^{-1}De_i)$ are $D^{-1}(Ly)_{\le s}$, known modulo $\ell^{K-\sigma_i}$, and $(Ly)_{>s}\equiv0\pmod{\ell^{K}}$ --- what \code{LatticeQuotient} reads.
\item \textbf{m1} (ii): torsion coordinates are exact for every $\sigma_i<K$ (not ``iff $\sigma_i\le K'$''); \code{BuildPres}'s refusal at $K\le2\sigma_{\max}$ is a conservative
  policy, sufficient and not necessary.
\item \textbf{m2} (iii): exact division by $\ell^{\sigma_i}$ (costs $\sigma_i$ digits), reduction of torsion rows (free), mixing precisions (the minimum).
\item \textbf{m3} (iii): the admissible certifying ranks of derived forms ($[F\mid\Delta_Y]$, $L_f$, the sublattice coordinates $C$, and $\rank_{\Q}$ of $\res$, $\cor$, $j$, the
  connecting map, $c_\sigma$, the invariants map), with the rational reason ($\cor\circ\res=[A:B]$).
\item \textbf{m4} (g): the pivots are $u_i\ell^{\sigma_i}$, $u_i$ units; the proof is verbatim.
\item \textbf{m5} (2)(b): $\delta[x]=[d\tilde x]$ stated; the sign does not affect Hecke equivariance, only entrywise comparisons with an independent $\delta$ and compositions
  through $\delta$.
\item \textbf{t1} \cl{prop:isomodeltree}(6)(b): ``$\Phi_\ell$ of \cl{prop:isomodelell}(4), $\Psi_{\mathcal T}$ of (4)(d)''. \textbf{t2} (4)(b): the five files calling
  \code{decompose()}.
\item Pointers: the provenance line, (2)(g)'s heading, the Appendix~\ref{app:status} note (no line, no count moves) and its referee pointer.
\end{itemize}

\subsection*{v1.23.0-pre.22 --- 2026-10-25 --- hecke-extensions slot 24, D2: the code audit's precision rule and signs}

\noindent\emph{This version.} \textbf{A pre-release, not owner-released, not refereed}; it carries v1.23.0-pre.21 (re-refereed ACCEPT) plus two findings of the D2 code audit
(\file{docs/implementation/slot24/D2\_code\_audit.md}, ACCEPT-with-fixes, no value-changing defect) and the C-D6 re-referee's four editorial findings. Only
\cl{prop:descentmaps} and, editorially, \cl{prop:isomodeltree} change; no statement, label, number or marker
moves ($222$ on $86$ labels).
\begin{itemize}\setlength{\itemsep}{0pt}
\item \textbf{F1} New item (2)(g): for a \emph{certified} Smith form $LAR=\Sigma$ over $\Z/\ell^{K}$ ($s=\rank_{\Q}A$ pivots, all of valuation $<K$), an exact Smith form over
  $\Z_{(\ell)}$ exists whose transforms and their inverses agree with $L$, $R$, $L^{-1}$, $R^{-1}$ modulo $\ell^{K-\sigma_{\max}}$; so kernel bases, generators, coordinates
  ($\sigma_i\le K-\sigma_{\max}$) and induced maps read from the transforms are exact to that precision; precisions decrease additively along a chain of certified forms, and an
  exhausted precision is a refusal. Proof: the audit's Appendix~B (a Schur complement, forced to vanish by the certification), made integral over $\Z_{(\ell)}$. Without the
  certification the rule fails ($A=(\ell^{K})$). Until now (d) and \cl{prop:certifiedeell}(6) certified only group types; D2's maps rest on (2)(g).
\item \textbf{F8} (2)(b): the maps $H^{q}(B)\to H^{q+1}(A,B)$, $H^{q}(A,B)\to H^{q}(A)$ are induced by $c'\mapsto(0,c')$, $(c,c')\mapsto c$; the second is $-\delta$ through $\pi$.
\item \textbf{The C-D6 re-referee of pre.21} (\file{C-D6\_referee.md}, ``Re-referee of pre.21 (M1)'': \textbf{ACCEPT}, four editorial findings), in \cl{prop:isomodeltree}:
  \textbf{N1} (4)(b) and the code map cite \code{decompose()} of \file{PeriodicColumn.h} (signed lexicographic order, $j$ relative to the fixed generator) as the representative
  convention; \textbf{N2} (4)(b): (L), (L$'$) finite per degree, the comparisons not periodic (coefficients grow linearly in $q$), built to the degree in use; \textbf{N3} (L$'$):
  $p'$ any chain map, its surjectivity unused; \textbf{N4} (6)(b) names the effective comparison $\vartheta=\Phi_\ell\circ\Psi_{\mathcal T}$.
\item Pointers: the provenance line, the code map (D2 built, \file{DescentMaps}), the Appendix~\ref{app:status} note (no line, no count moves) and its referee pointer.
\end{itemize}

\subsection*{v1.23.0-pre.21 --- 2026-10-25 --- hecke-extensions slot 24, D5: the C-D6 referee's fixes M1, m1--m5}

\noindent\emph{This version.} \textbf{A pre-release, not owner-released}; it carries v1.23.0-pre.20 (refereed ACCEPT-with-fixes,
\file{docs/implementation/slot24/C-D6\_referee.md}: no blocker, one major, five minors) with all six findings applied; \textbf{not re-refereed}. Only
\cl{prop:isomodeltree} changes; no statement, label, number or marker moves ($222$ on $86$ labels).
\begin{itemize}\setlength{\itemsep}{0pt}
\item \textbf{M1} (4) is rewritten as (a)--(d) and made \textbf{effective}. (a) The tree is the $\Ga$-subcomplex of $\Xbulk$ on the orbits of $v_6$, $v_4$, $\{v_4,v_6\}$ ($\mathsf i_{\mathcal T}$);
  an explicit retraction $\mathsf r_{\mathcal T}$ and homotopy $\mathsf G_{\mathcal T}$ on the $19$ representatives, $\mathsf r_{\mathcal T}\mathsf i_{\mathcal T}=\id$,
  $\id-\mathsf i_{\mathcal T}\mathsf r_{\mathcal T}=\partial\mathsf G_{\mathcal T}+\mathsf G_{\mathcal T}\partial$ (\cl{prop:bredon}(2) in closed form). (b) The columns' $\Z$-linear contractions with fixed
  coset representatives, the solvers $\mathsf h_p$ of $\ker\cmpI$ and $\ker\varepsilon_{\mathcal T}$ by descending column index, the lifting formula (L) (\cl{lem:strictlifting} made
  effective) and its homotopy form (L$'$). (c) $\Phi^{\Z}_{\mathcal T}$ is (L) of $\mathsf r_{\mathcal T}\circ\cmpI$, $\Psi^{\Z}_{\mathcal T}$ is (L) of $\mathsf i_{\mathcal T}\circ\varepsilon_{\mathcal T}$; explicit
  homotopies $\mathsf H_{\mathcal T}$, $\mathsf H_X$, $\mathsf H_\ell$. (d) $\Phi_{\mathcal T}$, $\Psi_{\mathcal T}$ with explicit homotopies $\mathsf K$, $\mathsf K'$. (5)'s (T3) and (6) now rest on
  constructed maps; the code map names the route on the code's $\Iso$.
\item \textbf{m1} \cl{prop:isomodelell}(4), (5) are used over every $\mathsf O_\ell$-algebra; (4)(d) and (6)(b) record that their proofs need only $\abs{Q_\sigma}\in\cring^{\times}$.
\item \textbf{m2} (3): ``strict deformation retract'' replaced by ``a retract with $\Psi^{\ast}\Phi^{\ast}=\id$ strictly and $\Phi^{\ast}\Psi^{\ast}\simeq\id$'', the homotopy $\mathsf H_\ell^{\ast}$ of (4)(c); no
  side conditions asserted.
\item \textbf{m3} The proof of (1)(d)(iv): $a_{yx}=e_ya_{yx}e_x$ for the column differentials and the edge's $-\varepsilon_t$ by centrality of $e_x$ in $\cring[\Ga_x]$.
\item \textbf{m4} (5): its last sentence is about groups and fibre-map-induced maps; other maps between the glued groups enter through $\Phi_{\mathcal T}$, $\Psi_{\mathcal T}$.
\item \textbf{m5} (6) is stated for every cochain endomorphism $T$, instantiated at level $N$ by $T_p$, $p\nmid N$ ([R-Q4]); pre.20's changelog line ``(6) \dots\ stated now for
  $p\nmid N$'' is to be read so. The proof of (1)(b): the sign reversal is relative to \cl{def:isomodelell}'s $\bar g-1$, not to \cl{rmk:treeconeres}.
\item A companion check, \file{docs/implementation/slot24/cd6\_iso\_check.py}, builds an $\Iso_{\bullet}$ and checks (4) entirely: $12\,870$ checks, $0$ mismatches; the pre.20 check and
  the referee's scripts rerun unchanged, $0$ mismatches.
\item Pointers: the provenance line, the \S\ref{sec:codemap} row, the Appendix~\ref{app:status} note (no line, no count moves) and its referee pointer.
\end{itemize}

\subsection*{v1.23.0-pre.20 --- 2026-10-25 --- hecke-extensions slot 24, D5: the tree models, bulk only (clause C-D6)}

\noindent\emph{This version.} \textbf{A pre-release draft, not owner-released, not refereed}; it carries v1.23.0-pre.19 plus one new subsection with two statements,
stated \emph{before} D5's code (contract first, rule~21): clause C-D6 of \file{docs/implementation/slot24/DESCENT\_plan.md} v0.4, whose milestone D5 the owner brought
forward at gate G-T on 2026-10-25, ahead of D3 and D4 (the evidence: \file{DESCENT1\_report.md}, the dense $\Iso^{(2)}$ memory wall at $D\approx3\,080$, the tree at $2$
at $D=11\,110$). No marker moves ($222$ on $86$ labels); the new statements are the last of their section, so no statement number moves.
\begin{itemize}\setlength{\itemsep}{0pt}
\item \PNEW{} \S\ref{sec:isomodeltree} and \cl{def:isomodeltree}: for $x$ the two vertices and the edge of the Bass--Serre tree, $\ell$-local columns $\cring[\Ga]e_x$ (periodic
  where the $\ell$-Sylow $P_x\ne1$, degree $0$ otherwise) over $\mathsf O_\ell$-algebras, and the model as the cone of a lift $\varphi'$ of the tree boundary; canonical lifts
  at $\ell=2$ ($(1,-e_U)$, $(1+S,-e_U)$, $e_U=\frac13(1+U^{2}+U^{4})$), $\ell=3$ and $\ell\ge5$; the free model is \cl{rmk:treeconeres}'s.
\item \PNEW{} \cl{prop:isomodeltree} (\PROVED, \STD{} \cite[Thm.~2.2.6]{weibel}, \cite[Ch.~III.6.2]{brown}). (1) Projective resolutions computing $H^{\bullet}(\Ga;M)$ in every
  degree for every $M$ ($\Ind_{\Ga_0(N)}^{\Ga}$ included), cochains $\bigoplus M^{Q_x}$ with ranks by traces, one integral form over $\mathsf O_\ell$, and
  \cl{prop:certifiedeell}(2)--(6) on it with its own rank sources. (2) The cochains in closed form at $\ell=2$, $3$, $\ge5$; periodicity from degree $1$, resp.\ $2$;
  \cl{prop:ellthreeclosed}(2) read off the cochains. (3) Closed-form comparisons $\Phi^{\mathcal T}_\ell$, $\Psi^{\mathcal T}_\ell$ with the free tree (column maps
  $\abs{Q_x}^{\pm\lfloor t/2\rfloor}$ and explicit corrections), $\Phi^{\mathcal T}_\ell\Psi^{\mathcal T}_\ell=\id$ strictly: the $\ell$-local cochains are a strict deformation retract of the
  free tree's. (4) Comparisons $\Phi_{\mathcal T}$, $\Psi_{\mathcal T}$ with $\Iso$ over the augmentation to $\Z$ (\cl{lem:strictlifting}), not over $C_{\bullet}(\Xbulk)$, mutually
  inverse on cohomology. (5) (T1)--(T3) of \cl{prop:descentmaps}(5) hold, so its fibre maps, level cone, long exact sequence and certified rule hold on the tree at
  $\ell\in\{2,3\}$. (6) Hecke by transport, $T'_{\mathcal T}=\Psi^{\ast}_{\mathcal T}\circ T\circ\Phi^{\ast}_{\mathcal T}$: it induces the operator, equals the transport on $\Iso^{(\ell)}$ on
  cohomology, respects relations up to homotopy and fibre maps strictly; stated now for $p\nmid N$ (the plan's [R-Q4]), implemented with D3. Bulk only: nothing on the
  boundary, the pair or the cusp (the plan's C-D8). A code-map paragraph (not built) and a formalization note (no marker).
\item A check outside the build, \file{docs/implementation/slot24/cd6\_tree\_check.py} (own group ring, own induced module, the models built as chain data): the chain
  identities exactly in $\Q[\Ga]$, and on $145$ cells and $18$ level pairs at $\ell=2$, $3$ (and $5$, $7$ on the Bredon tree) the closed forms, the groups against the oracle
  tables, the comparisons and the level cone: $49\,303$ checks, $0$ mismatches.
\item Pointers: the \S\ref{sec:codemap} row; Appendix~\ref{app:status}: two lines, $395\to397$, \code{def} $66\to67$, \code{prop} $94\to95$, \PROVED{} $230\to231$, \BSTD{} $99\to100$,
  tagged $265\to266$; its referee pointer; Appendix~\ref{app:labelindex}: three labels.
\end{itemize}

\subsection*{v1.23.0-pre.19 --- 2026-10-25 --- hecke-extensions slot 24, D2: the C-D3 referee's fixes m1--m5}

\noindent\emph{This version.} \textbf{A pre-release, not owner-released}; it carries v1.23.0-pre.18 (refereed ACCEPT-with-fixes,
\file{docs/implementation/slot24/C-D3\_referee.md}: no blocker, no major, five minors) with the five minors applied; \textbf{not re-refereed}. No statement, label, number
or marker moves ($222$ on $86$ labels).
\begin{itemize}\setlength{\itemsep}{0pt}
\item \textbf{m1} \cl{prop:descentmaps}(3): (a) and (b) say that $[AgA]$, $g=\smat100p$, is \emph{this document's} $U_p$, the adjoint $U_p^{\ast}$ on $p$-old spaces
  (\cl{prop:levelconehecke}(3)(d)); (b)'s failure of the square with $\res$ is that convention's, for $p\mid N$ --- the classical $U_p=[A\smat p001A]$ has a strict square at
  $\chi=1$ --- and for $p\nmid N$, $p\mid M$ both squares fail. \PNEW{} item \textbf{(3)(b$'$)}: if $p\mid N$, $\cor\circ[BgB]=[AgA]\circ\cor$ on $H^{q}$ for every $q$,
  every $M$, over every ring and over $\Z$, with its proof (the double coset formula with $A'=A\cap g^{-1}Ag$: $B\cap A'=B'$ and $A=A'B$ by the transversal $T^{j}\in B$);
  on the Hida tower $\Ga_0(p^{r+1})\subset\Ga_0(p^{r})$ corestriction is $U_p$-equivariant over $\Z$ in every degree. The proof's check paragraph records the referee's
  $32$ cells plus $10$ more run at this version (the square with $\cor$: $15$ cells with $p\mid N$ commute, $5$ with $p\nmid N$ fail).
\item \textbf{m2} (3)(a) cites the criterion \cl{prop:levelconehecke}(1) alone; its instance (2) for $U_p$ needs equal $\Ga_0$-levels.
\item \textbf{m3} (3)(c) assumes $p\nmid M$.
\item \textbf{m4} (2)(d)(i): the rank is read on the reduction mod $\ell'$ of the model's integral form ($\Iso$ over $\Z$, $\Iso^{(\ell)}$ over $\mathsf O_\ell$ with
  $\ell'\nmid\mathsf q_\ell$), as in \cl{prop:certifiedeell}(4)(c).
\item \textbf{m5} (4)(c): $H^{q}(\Ga_0(p^{r});\Sym^{n}\Z^{2})\cong H^{q-1}((\Z/p^{r})^{\times};H^{1}(\Ga_1(p^{r})))$ for $q\ge2$, all of it $2$- and $3$-primary; no $\ell$-torsion
  for $\ell\ge5$ even when $\ell\mid\abs{\mathsf G}$.
\item The statement's provenance line, the \S\ref{app:status} note (no line, no count moves) and the referee pointer of the status appendix.
\end{itemize}

\subsection*{v1.23.0-pre.18 --- 2026-10-25 --- hecke-extensions slot 24, D2: maps and the level cone over $\Z$ (clause C-D3)}

\noindent\emph{This version.} \textbf{A pre-release, not owner-released, not refereed}; it carries v1.23.0-pre.17 (refereed ACCEPT,
\file{docs/implementation/slot24/C-D1\_three\_referee.md}, four minor remarks, folded in below) plus one statement, stated \emph{before} D2's code (contract first,
rule~21): clause C-D3 of \file{docs/implementation/slot24/DESCENT\_plan.md} v0.3, and the part of the owner's slot-22 LES-0 request (the level-change sequence for
$\Ga_0(MN)\subset\Ga_0(N)$, inflation--restriction in the normal case, the Hida $p$-tower first) that \S\ref{sec:nebles} had left at $6\in\cring^{\times}$. No marker moves
($222$ on $86$ labels); the new statement is the last of its section, so no statement number moves.
\begin{itemize}\setlength{\itemsep}{0pt}
\item \PNEW{} \cl{prop:descentmaps} (\PROVED, \STD{} \cite[Ch.~III.6.2, III.8--10, VII.6]{brown}, \cite[Thm.~2.2.6]{weibel}). \emph{Models}: an $S$-model is any finitely
  generated projective resolution of $S$ over $S[\Ga]$; comparisons are chain maps over $\id_S$. (1) A $\Z[\Ga]$-map of fibres commutes strictly with every comparison
  ($\cmpI^{\ast}$, $\Phi_\ell^{\ast}$, $\Psi_\ell^{\ast}$, a tree model's), acts block by block on permutation-type models, and glues to the map over $\Z$; restriction, trace
  and conjugation are induced in every degree over every ring (no $6$ inverted); the split case in all degrees. (2) The level cone on every model is exact and
  computes $H^{\bullet-1}(A;Q_{A/B})$; the long exact sequence \eqref{eq:descentmapsles} over $\Z$ in every degree (not six-term at $2$, $3$: $H^{3}(\Ga_0(4),\Ga_0(8))=\Z/2$ at
  $n=0$), periodic from degree~$3$, glued by \cl{prop:cohglue}; the certified-$K$ rule on the cone, $\rank_{\Q}$ from the Betti recursion $\beta_j=b_j(M_B)-b_j(M_A)$ or from
  the cone's own reduction at a qualifying prime; bad primes of the cone (degree $2$ only at $\ell\ge5$); the tree and sparse-cokernel routes at $\ell\ge5$ and, through
  $S[A]$-models, at $\ell=3$ when $\nu_3(N)=0$. (3) $\res$ and $\cor$ commute with $[\,\cdot\,\smat100p\,\cdot\,]$ in every degree when $p\nmid M$ (double coset formula): $T_p$
  for $p\nmid MN$, $U_p$ for $p\mid N$; level raising fails (\cl{prop:levelconehecke}(3)(c), lifted to $\Z$); the Hecke action on the relative term at $2$, $3$ is
  deferred to C-D4, with what C-D4 must supply. (4) Normal pairs: the $\mathsf G$-action, the norm identities and the five-term sequence over $\Z$; for torsion-free $B$
  two rows over every ring, and for $n>0$ $H^{q}(A)\cong H^{q-1}(\mathsf G;H^{1}(B))$ ($\Ga_1(N)\trianglelefteq\Ga_0(N)$, $N\ge4$: the bottom of the Hida tower; $\abs{\mathsf G}$
  kills $H^{\ge2}(A)$). (5) What an $\ell$-local tree model (D5, C-D6) must supply: (T1)--(T3). A closing list of what is deferred and why. A check outside the build,
  \file{docs/implementation/slot24/cd3\_descentmaps\_check.py} (own cosets, modules, free resolution and exact Smith forms): $74$ cone cells, $12$ two-row cells, $888$ oracle
  rows, $0$ mismatches.
\item The C-D1 three-shortcut referee's minor remarks: \textbf{m1} \cl{prop:descentNthree}(1), why ``no element of order $3$'' is ``no $U$-fixed point''; \textbf{m2}
  \cl{prop:modeltorsion}'s last line points to the full characterisation; \textbf{m3} the header of \file{cd1b3\_three\_shortcut\_check.out} records one command per figure
  and the proof says $84=80+4$ in two runs, and ``$22$ of $32$''; \textbf{m4} the proof records the shared construction and the referee's independent closure.
\item Pointers: \S\ref{sec:nebles}'s scope sentence, \cl{rmk:integralscope}(4). \S\ref{sec:descentN}'s code-map paragraph (D2 not built; D1 merged at \texttt{1cc9a6d3}),
  Formalization note, the \S\ref{sec:codemap} row.
\item Appendix~\ref{app:status}: one line, $394\to395$, \code{prop} $93\to94$, \PROVED{} $229\to230$, \BSTD{} $98\to99$, tagged $264\to265$; Appendix~\ref{app:labelindex}: two labels.
\end{itemize}

\subsection*{v1.23.0-pre.17 --- 2026-10-25 --- hecke-extensions slot 24, D1: the $3$-adic shortcut at $\nu_3(N)=0$}

\noindent\emph{This version.} \textbf{A pre-release, not owner-released, not refereed}; it carries v1.23.0-pre.16 plus one statement, written \emph{after} D1's code
(branch \texttt{s24-d1impl}, \texttt{726c3a67}, \texttt{87e8cf32}) to answer finding M1 of \file{docs/implementation/slot24/D1\_code\_audit.md} (commit \texttt{ffb59130}):
the code's $\ell=3$ route at $\nu_3(N)=0$ was sound and agreed with $\Iso^{(3)}$ on every test cell, but \cl{prop:descentN}(7) stated the route for $\ell\ge5$ only.
No marker moves ($222$ on $86$ labels); the new statement is the last of its section, so no statement number moves.
\begin{itemize}\setlength{\itemsep}{0pt}
\item \PNEW{} \cl{prop:descentNthree} (\PROVED, \STD{} \cite[Ch.~III.6.2]{brown}, \cite[Thm.~2.2.6]{weibel}). (1) $\nu_3(N)=0$ iff $\Ga_0(N)$ has no element of order~$3$ iff
  $9\mid N$ or a prime $p\equiv2\pmod3$ divides $N$; then every finite subgroup of $\Ga_0(N)$ has order $1,2,4$, and $M_N$ is induced over $\langle U^{2}\rangle$.
  (2) $\mathrm T(M_N\otimes\Z_{(3)})$ is free, base-changes to every $\Z_{(3)}$-algebra, and computes $H^{\bullet}(\Ga_0(N);\Vcoef)\otimes\Z_{(3)}$, as does $\Ceq$
  ($\cmpI^{\ast}\otimes\Z_{(3)}$ a quasi-isomorphism): \textbf{no $3$-torsion in $H^{\ge2}$}; checked against \cl{prop:modeltorsion} and \cl{prop:ellthreeclosed}(2),(3).
  (3) \cl{prop:descentN}(7)(b) verbatim at $\ell=3$: the certified-$K$ rule $s=r=\rank\mathrm T^{0}-b_0$ with closed forms, the stopping rule of
  \cl{prop:certifiedeell}(6), the generators (orbit sums; idempotent columns at the $S$-fixed points, order $4$ prime to $3$) and the F4 certificate; there are no
  $U$-fixed points. (4) $n$ odd: nothing $3$-primary. (5) The hypothesis is needed (oracle cells $(1,0)$, $(7,0)$, $(3,4)$ have $3$-torsion in $H^{\ge2}$); $\nu_2$ plays
  no role; $9\mid N$ is covered. A check outside the build, \file{docs/implementation/slot24/cd1b3\_three\_shortcut\_check.py}: $84$ cells, $0$ mismatches.
\item \cl{prop:descentN}(7): a pointer to the new statement. \S\ref{sec:descentN}'s code-map paragraph and Formalization note; a \S\ref{sec:codemap} row for \S\ref{sec:descentN}.
\item Appendix~\ref{app:status}: one line, $393\to394$, \code{prop} $92\to93$, \PROVED{} $228\to229$, \BSTD{} $97\to98$, tagged $263\to264$; Appendix~\ref{app:labelindex}: one label.
\end{itemize}

\subsection*{v1.23.0-pre.16 --- 2026-10-25 --- hecke-extensions slot 24, D1: the Dickson-proof referee's fixes}

\noindent\emph{This version.} \textbf{A pre-release, not owner-released}; it carries v1.23.0-pre.15 plus the fixes of the Dickson-proof referee
(\file{docs/implementation/slot24/C-D1\_dickson\_referee.md}, commit \texttt{509d67b5}: ACCEPT-with-fixes, $1$ Minor, $4$ nits; the mathematics was found correct at every prime);
\textbf{not re-refereed}. No marker moves ($222$ on $86$ labels); one statement is added, so the numbers after Claim~\ref{cl:dimH1} in its section shift by one (labels do not).
\begin{itemize}\setlength{\itemsep}{0pt}
\item \textbf{D-1} (Minor). The proof of Claim~\ref{cl:dimH1}(2) is a standalone \PNEW{} \cl{lem:dicksonSL2} (\PROVED), stated for every prime $\ell$ and every $n$, cited by
  \cl{cl:dimH1}(2), \cl{prop:oddnltwo}(2) (at $\ell=2$) and \cl{prop:dicksonH0}. ``First non-zero at $n=\ell+1$'' is restricted to $\ell\ge3$ ($\delta_2(2)=1$, $Q=d_2$).
\item \textbf{D-2}. \cl{prop:oddnltwo}(2): the \STD{} tag on \cite{wilkerson} removed; a plain pointer to the lemma.
\item \textbf{D-3}. \cl{prop:dicksonH0}(2) rests on the $\SL_2$ half of (1) only; the $\GL_2$ half, $\Fl[L^{\ell-1},Q]$, proved in the lemma ($g_AL=\det A\cdot L$, $g_AQ=Q$).
\item \textbf{D-4}. Step 3: $\theta$ has a minimal polynomial over $\mathcal K^{G}$ because the upper bound made $\mathcal K$ algebraic over $k$; the five coefficients compared are named.
\item \textbf{D-5}. \cl{prop:descentN}(4): $b_1=1$ \emph{exactly} at $N\in\{2,3,5,7,13\}$.
\item Appendix~\ref{app:status}: one line, $392\to393$, \code{lem} $60\to61$, \PROVED{} $227\to228$, tagged $262\to263$; Appendix~\ref{app:labelindex}: one label.
\end{itemize}

\subsection*{v1.23.0-pre.15 --- 2026-10-25 --- hecke-extensions slot 24, D1: the C-D1 / C-D2 referee's fixes}

\noindent\emph{This version.} \textbf{A pre-release, not owner-released}; it carries v1.23.0-pre.14 plus the fixes of the C-D1 / C-D2 referee
(\file{docs/implementation/slot24/C-D1\_referee.md}, commit \texttt{e28dcb09}: ACCEPT-with-fixes, $0$ Blocker, $1$ Major, $4$ minor, $2$ nits), applied to
\S\ref{sec:descentN} and to Claim~\ref{cl:dimH1}(2); \textbf{not re-refereed}. Every new claim is proved in place. No label, statement number or marker moves ($222$
markers on $86$ labels); one tag set moves (Claim~\ref{cl:dimH1}, below).
\begin{itemize}\setlength{\itemsep}{0pt}
\item \textbf{F1} (Major). (a) \cl{prop:badprimesN}(2)'s proof: for $\ell>n>0$, $\ell\nmid N$, Claim~A and the lower unipotent give $s_\ell=0$, so the good-prime rule
  $\ell>n+1$ of \cl{prop:descentN}(6), D1's production path, rests on Claims~A and~B alone. (b) \cl{cl:dimH1}(2), $\Fl[x,y]^{\SL_2(\Fl)}=\Fl[L,Q]$, is
  \PROVED{} in place, for every prime $\ell$, along the referee's sketch with two simplifications: the Moore determinant gives
  $\prod_{v}(T-v)=T^{\ell^{2}}-QT^{\ell}+L^{\ell-1}T$ (integrality); the degree $[\Fl(x,y):\Fl(L,Q)]\le\abs{\SL_2(\Fl)}$ by a tower $k\subseteq k(x)\subseteq k(x,y)$, against
  $[\Fl(x,y):\Fl(x,y)^{G}]\ge\abs{G}$ from the trivial stabiliser of $x+y^{2}$, in place of freeness and Artin's theorem; a factorial ring is integrally closed. So
  the count $\delta_\ell(n)$ at $\ell\nmid N$, $5\le\ell\le n$ (and at level~$1$, \cl{prop:badprimes}(1)) no longer rests on the unverified \cite{wilkerson}, and
  ``not by omission'' holds as stated. Pointers updated: \cl{prop:oddnltwo}(2) (the $\ell=2$ case, $L=d_3$, $Q=d_2$), the proof of \cl{prop:badprimes},
  \cl{prop:dicksonH0}(1) (its $\SL_2$ half; the $\GL_2$ half keeps its citation).
\item \textbf{F2}. \cl{prop:descentN}(4): $6b_1=6+m_N-3\nu_2-4\nu_3\ge6$ for every $N>1$ at $n=0$, $\chi=1$, from \cl{prop:levelorbit}(3) and the local counts of the
  fixed points of $S$, $ST$ (CRT, Hensel), replacing ``by the genus formula, which is not proved here''; equality at $N=2,3,5,7,13$ (the referee's ``only at $N=2$''
  is corrected; the minimum $6$ stands). Nothing else uses it.
\item \textbf{F3}. The code map: \code{SparseCokernel}'s generators are the averaging-idempotent columns and transported blocks, admissible for (7)(b); the blocker is the
  prime-field gate and the RREF elimination.
\item \textbf{F4}. \cl{prop:descentN}(7)(b): the certificate that a family generates $(V_{\Z/\ell^{K}})^{F_j}$ (fixed by a generator, spanning mod~$\ell$; Nakayama), and
  that the comparison value $r$ is the closed form, never the number of rows.
\item \textbf{F5}. \cl{prop:descentN}(5): $\Hom(E,M)$ induces $e\cup-$ for every $M$ (dimension shifting), and $U_p$'s defining label (\cl{prop:levelcorrespondence},
  \cl{prop:nebcorr} at $A=B=\Ga_0(N)$).
\item \textbf{F6, F7} (nits). (7)(a): the source by any generating family, the lifts taken in $e_\sigma M_{\Z_{(\ell)}}$. The subsection's preamble says why it sits under
  \S\ref{sec:isomodelell}.
\item Appendix~\ref{app:status}: \cl{cl:dimH1}'s line gains \PROVED{} and drops \cite{wilkerson}; \cl{prop:descentN}'s gains \cite[Ch.~V, \S3]{brown}; $392$ lines,
  \PROVED{} $226\to227$.
\end{itemize}

\subsection*{v1.23.0-pre.14 --- 2026-10-25 --- hecke-extensions slot 24, D1: the descent at level $N$ and its bad primes}

\noindent\emph{This version.} \textbf{A pre-release draft, not owner-released, not refereed}; it carries v1.23.0-pre.13 plus clauses C-D1 and C-D2 of
\file{docs/implementation/slot24/DESCENT\_plan.md} v0.3 (milestone D1; audit findings M8, m3, m4, m7 of \file{DESCENT\_plan\_audit.md}). Contract first (rule~21).
\textbf{Additive}: one subsection after the last statement of \S\ref{sec:isomodelell}, so that no statement number moves; no existing statement, tag set, marker or
value moves ($222$ markers on $86$ labels; the new statements carry none).
\begin{itemize}\setlength{\itemsep}{0pt}
\item \PNEW{} \cl{prop:descentN} (\PROVED{}), for every $\Z[\Ga]$-module $M$ finitely generated and free over $\Z$, in particular $M_N=\Ind_{\Ga_0(N)}^{\Ga}\Vcoef$ and
  $M_N(\chi)$ over $\mathsf O_\chi=\Z[\zeta_{m_\chi}]$: (1) $\CI(M)$, $\Ceq(M)$, $\mathrm T(M)$ and their base changes; (2) Proposition~\ref{prop:descent} for $M$; (3) maps
  and Hecke operators determined locally, by Proposition~\ref{prop:cohglue}'s morphism clause alone; (4) the coherence term $\Hom^{-1}\cong\Qsix^{b_0b_1}$, so
  Proposition~\ref{prop:homvanish}(2) extends iff $b_0b_1=0$ --- for $M_N(\chi)$ iff $n>0$ or $\chi\ne1$; it fails at $n=0$, $\chi=1$, $N=2$ (finding M8, the
  restriction the plan offered, made exact); (5) Remark~\ref{rmk:truncdevice} for $M$, and periodicity above the vcd by $\res^{\Ga}_{A}e\cup-$, equivariant for every
  $[AgA]$, $T_p$ ($p\nmid N$) included (finding m4); (6) the good prime: for $M_N$ every $\ell'\ge5$, $\ell'>n+1$ ($\ell'\nmid N$ not needed); for $\chi\ne1$,
  $\ell'\nmid m_\chi N$ needed (finding m3); (7) the certified exponent at $\ell\ge5$ on the direct Smith form of $\Ceq(M_{\Z/\ell^{K}})$ (the D0 report's
  recommendation) and on the level-$N$ sparse cokernel over $\Z/\ell^{K}$ (finding m7), with $\rank_{\Q}d^{0}=\rank C^{0}-b_0$, and
  Proposition~\ref{prop:certifiedeell}(4) at level $N$; (8) $M_N(\chi)\otimes\Z/\ell^{K}$ as a sum of \S\ref{sec:nebentypus}'s character models over $W_K(\F_q)$.
\item \PNEW{} \cl{prop:badprimesN} (\PROVED{}): (1) for every such $M$ and $\ell\ge5$ the torsion lies in $H^{1}$, with $\dim H^{0}(\Ga;M/\ell)-b_0$ cyclic summands;
  (2) at level $N$ the count is the level-$1$ Dickson count $\delta_\ell(n)$ for $\ell\nmid N$ and the Borel count $\beta_\ell(n)$ for $\ell\mid N$, which reaches
  $\ell=n+1$; the closed-form list $\mathcal B(n,N)$, exact at $\ell\ge5$, with every $\ell>n+1$ torsion-free by proof; $N=1$ is
  Proposition~\ref{prop:badprimes}; (3) nebentypus at $\ell\nmid6m_\chi$: torsion only at $\ell\mid N$ with $\chi$ of $\ell$-power conductor, count
  $\beta_\ell(n,t_\lambda)$; the primes $\ell\mid m_\chi$ not treated.
\item \S\ref{sec:descentN}'s code map (not built: the direct Smith route, the sparse cokernel over $\Z/\ell^{K}$, the level-$N$ assembly, $E$ on $\CI(M_N)$, $M_N(\chi)$
  over $\mathsf O_\chi$) and formalization note; an exact check outside the build ($1376$ cells, $0$ mismatches), not a tag.
\item Remark~\ref{rmk:integralscope}(4): one parenthesis (the level half of the plan's C-D7). \S\ref{sec:codemap}'s row for \S\ref{sec:isomodelell}: one sentence.
\item Appendix~\ref{app:status}: two lines, counts $390\to392$; Appendix~\ref{app:labelindex}: three labels.
\end{itemize}

\subsection*{v1.23.0-pre.13 --- 2026-10-25 --- hecke-extensions slot 24, D0: the C-D0 referee's fixes}

\noindent\emph{This version.} \textbf{A pre-release, not owner-released}; it carries v1.23.0-pre.12 plus the required fixes of the C-D0 referee
(\file{docs/implementation/slot24/C-D0\_referee.md}, commit \texttt{4089cdf8}: ACCEPT-with-fixes, $0$ Blocker, $2$ Major, $6$ minor), applied to
Proposition~\ref{prop:certifiedeell}; \textbf{not re-refereed}. Every new claim is proved in place. No label, statement number, tag set, marker or value moves ($222$
markers on $86$ labels).
\begin{itemize}\setlength{\itemsep}{0pt}
\item \textbf{F1} (Major). (4) no longer leaves the qualifying primes an input for $M\ne\Vcoef$: (4)(a) $24\cdot H^{j}(\Ga;M)=0$ for $j\ge2$ and every $\Z[\Ga]$-module $M$
  (Mayer--Vietoris from \eqref{eq:treeses} with Shapiro, the argument of Claim~\ref{cl:amalgam}(4)), so every $\ell'\ge5$ qualifies for $d^{q}$, $q\ge1$ ($q=1$, $H^{2}$,
  checked); (4)(b) at $q=0$, $\ell'$ qualifies iff $\dim H^{0}(\Ga;M\otimes\F_{\ell'})=b_0$, measured on the same complex, and $\rank_{\Q}d^{0}=\rank C^{0}_\ell-b_0$; $b_0=0$
  ($n>0$), $1$ ($n=0$) for $\Vcoef$ and $\Ind_{\Ga_0(N)}^{\Ga}\Vcoef$, proved.
\item \textbf{F2} (Major). (4)(c): the two admissible sources of $\rank_{\Q}d^{q}$ (this $\mathsf O_\ell$-complex mod $\ell'$; certified Betti numbers through
  $\rank C^{q}_\ell-b_q-\rank_{\Q}d^{q-1}$); census, shipped-$C_7$ and $\Iso^{(\ell')}$ ranks stated inadmissible, with the reason. Code map: the \code{ell()} guard
  as a D0 requirement (audit note F-5).
\item \textbf{F3}. (6): the Smith valuations a list with multiplicity, their number $s$.
\item \textbf{F4}. (6): $a_t\le v_\ell(24)$ for $q\ge1$ ($3$ at $2$, $1$ at $3$, $0$ at $\ell\ge5$), the stopping rule (raise $K$ to the ring's cap, then refuse); the caps bind
  only at $d^{0}$.
\item \textbf{F5}. Code map: the denominator check (every denominator divides a power of $\mathsf q_\ell$) as a D0 obligation.
\item \textbf{F6}. Code map (ii): $W\le61$ reachable now through \code{ModCtx<uint64\_t>} and \code{Ring2k}'s masking; $W\in[62,64]$ needs the new reduction;
  \S\ref{sec:codemap}'s row, which said ``the reduction to $\Z/2^{64}$'', reworded to agree.
\item \textbf{F7}. Code map (i$'$): the \code{DenseBuf} $\to$ \code{Ring2k} / \code{Ring3k} hand-off; $q_{\max}\ge5$ for the (P) arm at degree~$6$.
\item \textbf{F8}. \file{docs/implementation/slot24/DESCENT\_plan.md} v0.3: row C-D0 ($\abs{Q_\sigma}\in\{2,4\}$ at $\ell=3$, $\mathsf q_3=4$; ``(1)--(3), (4) restated as
  \ref{prop:certifiedeell}(4)'') and the good-prime wording of \S2.1 and D0.
\item Appendix~\ref{app:status}: the line of \ref{prop:certifiedeell} names the proof's new \STD{} input \cite[Ch.~III.10]{brown}; no count moves.
\end{itemize}

\subsection*{v1.23.0-pre.12 --- 2026-10-25 --- hecke-extensions slot 24, D0: the certified-exponent rule on $\Iso^{(\ell)}$}

\noindent\emph{This version.} \textbf{A pre-release draft, not owner-released, not refereed}; it carries v1.23.0-pre.11 plus clause C-D0 of
\file{docs/implementation/slot24/DESCENT\_plan.md} v0.2 (milestone D0; audit finding M3 of \file{DESCENT\_plan\_audit.md}: D0 certifies $K$ on $\Iso^{(2)}$ and
$\Iso^{(3)}$, and Proposition~\ref{prop:certifiede} is stated for complexes of free $\Z$-modules only). Contract first (rule~21): the plan requires a referee
\textbf{before} D0's validation lands. \textbf{Additive}: one statement at the end of \S\ref{sec:isomodelell}'s statements, so that no statement number moves; no
existing statement, tag set, marker or value moves ($222$ markers on $86$ labels; the new statement carries none).
\begin{itemize}\setlength{\itemsep}{0pt}
\item \PNEW{} \cl{prop:certifiedeell} (\PROVED{}): (1) Definition~\ref{def:isomodelell} and Proposition~\ref{prop:isomodelell}(1)--(3) over any ring in which
  $\mathsf q_\ell=\operatorname{lcm}\abs{Q_\sigma}$ is a unit ($\mathsf q_2=3$, $\mathsf q_3=4$, $\mathsf q_{\ell\ge5}=12$); one model over $\mathsf O_\ell=\Z[1/\mathsf q_\ell]$
  whose base changes to $\Z/\ell^{K}$ and $\F_{\ell'}$ ($\ell'\nmid\mathsf q_\ell$) are the complexes built; (2) its cochains free over $\mathsf O_\ell$, compatible with
  base change, computing $H^{q}(\Ga;M)\otimes\mathsf O_\ell$; (3) Proposition~\ref{prop:certifiede}(1)--(3) and Remark~\ref{rmk:certifiedeperdiff} over $\Z_{(\ell)}$;
  (4) $\rank_{\Q}$ at a prime $\ell'\nmid\mathsf q_\ell$ without $\ell'$-torsion in $H^{q+1}$; (5) the code's full layout $\delta^{q}=P_{q+1}\delta^{q}P_q$ is
  $\mathrm{diag}(d^{q},0)$ over $\mathsf O_\ell$; (6) the Smith valuations over $\Z/\ell^{K}$ are the exponents $a_t<K$, their number $r_K$.
\item Its code map (built: $d$ once over $\Q[\Ga]$, denominators $\abs{Q_\sigma}$ only, reduced entrywise; not built: the \code{Ring2k} / \code{Ring3k} read-out on
  $\Iso^{(\ell)}$, a reduction to $\Z/2^{W}$ for $W\ge62$; \code{Ring3k}'s $K\le19$; the plan's ``$\abs{Q}\in\{1,2\}$ at $\ell=3$'' corrected to $\{2,4\}$).
\item \S\ref{sec:isomodelell}: ``What this chapter fixes'' item (v); the notation sentence names the new symbols; the scope sentence (``at best over
  $\Z[1/\mathsf q_\ell]$''); the formalization note (no \code{\textbackslash LEAN} marker). \S\ref{sec:codemap}'s row for \S\ref{sec:isomodelell}: one sentence.
\item Appendix~\ref{app:status}: one line, counts $389\to390$, with the note that the two pre.3 statements still lack lines; Appendix~\ref{app:labelindex}: one label.
\end{itemize}

\subsection*{v1.23.0-pre.11 --- 2026-10-21 --- hecke-extensions slot 24: the G2-H and G6 code-map sentences}

\noindent\emph{This version.} \textbf{A pre-release, not owner-released, not refereed}; it carries v1.23.0-pre.10 plus \textbf{code-map text only}:
\begin{itemize}
\item \S\ref{sec:obligations}, its code map: one sentence on G2-H (\file{docs/implementation/slot24/G2H\_report.md}~\S9; audit \file{G2H\_audit.md}, MERGE after
  fixes F1, F2; its note F6 asks for this sentence) --- on the stream route under \code{RunModuleT}'s lean retention, Obligation~\ref{obl:welldef}(iii) certified by
  \code{unit\_pivot\_residual\_product} without forming $\Th^{(q)}_{\mathrm{eq}}$, the predicate exactly the read-off's residual, so the obligation is discharged as
  written; route~2 compared in place for \ref{obl:agreement}(i); the early release of $X_q$, \code{Zc}; \code{howell\_rows}' adopted store; rollback
  \code{HECKE\_MODULE\_LEAN\_PEAK=0}. The report's draft is corrected in two places against \file{DenseElim.h}: the store is adopted only when also every row is a
  pivot row, and the adoption runs on the stream route in every caller, not only under the lean retention.
\item \S\ref{sec:codemap}, the row for \S\ref{sec:cone}: one sentence on G6 (\file{G6\_report.md}; audit \file{G6\_audit.md}, MERGE after fixes F1--F3) --- the
  span inclusions C2 / C3 / C5 by the unit-row route, the shared $Z$, the census cone release, rollback \code{HECKE\_FIELD\_DIET=0}, and the audit's note N1:
  \code{tape::factor\_basis}'s two re-blockings have no lever.
\item Audit (\file{docs/implementation/slot24/C1230\_pre11\_audit.md}): the G6 sentence now names the tape route's operand threshold
  \code{kSpanTapeMinArea} (below it C2 / C3 / C5 take the dense route of before, unit rows or not) and the heap trim among the levers
  \code{HECKE\_FIELD\_DIET=0} rolls back.
\end{itemize}
No statement, label, number, tag set or \code{\textbackslash LEAN} marker moves ($222$ on $86$ labels); every lever selects a route, never a value; the golden
reference and the digests are unchanged.

\subsection*{v1.23.0-pre.10 --- 2026-10-20 --- hecke-extensions slot 24: the FNTT and MT-1 code-map sentences}

\noindent\emph{This version.} \textbf{A pre-release, not owner-released, not refereed}; it carries v1.23.0-pre.9 plus \textbf{code-map text only}:
\begin{itemize}
\item \S\ref{sec:saturation}, its code map: one sentence on FNTT (\file{docs/implementation/slot24/FNTT\_report.md}; audit \file{FNTT\_audit.md}, fixes
  applied) --- \code{BuildCuspSpace} / \code{BuildModularSpace} in the NTT's transform domain and \code{HeckeOnCusp}'s delayed reduction, rollback
  \code{HECKE\_FORMS\_NTT\_CACHE=0}; and the forms $\ell$-range restated: the G2-X sentence's $\ell\notin[2,2^{16})$ stays a statement about
  \code{MulSeriesFast}; a forms space is built for every $5\le\ell<2^{16}$ and, for $2^{16}\le\ell<2^{27}$, only at $N<1024$ or where the two-prime product is exact
  ($N(\ell-1)^{2}<(5\cdot2^{25}+1)(7\cdot2^{26}+1)$, $2N-1\le2^{25}$), every other cell refused before the module run (the admitted range at the sat route's
  precision from \file{FNTT\_audit.md}~F5).
\item \S\ref{sec:obligations}, its code map: one sentence on MT-1 (\file{MT1\_report.md}; audit \file{MT1\_audit.md}, MERGE after fixes F1--F3) --- the pivot-search
  exit of \code{cyclic\_decomposition} at every thread count and the further splits over the consumer pool, Obligations~\ref{obl:welldef} and
  \ref{obl:agreement}(i) discharged as written, rollback \code{HECKE\_MODULE\_MT=0}.
\item The check \file{docs/implementation/slot23/C1230\_pre10\_check.md} amended these two sentences in place (wording only, against the code): the
  word range is $5\le\ell<2^{16}$; \code{HeckeOnCusp}'s delayed reduction is gated by the two levers and $(m+1)\ell^{2}<2^{64}$ alone, not by the
  transform-domain route's conditions; the range figures $114868$, $77702$ are the largest admitted integers (primes $114859$, $77699$); under
  \code{ENGINE\_HOST\_ELIM\_HOOK=0} only MT-1's assembly splits; the gate runs the read-off and the solves at pools $1$, $2$, $7$, $16$ (only
  \code{cyclic\_decomposition} also at $3$, $13$).
\end{itemize}
The slot-24 record-format event L23-H1 (\code{orbifold\_cell/v2}) concerns no text of this document. No statement, label, number, tag set or
\code{\textbackslash LEAN} marker moves ($222$ on $86$ labels); every lever selects a route, never a value; the golden reference and the digests are unchanged.

\subsection*{v1.23.0-pre.9 --- 2026-10-20 --- hecke-extensions slot 24: ELL-W, the wide tier of the $\tau$-certification (draft)}

\noindent\emph{This version.} \textbf{A pre-release draft, not owner-released, not refereed} (an adversarial audit of the clause
follows; plan question Q5: a PROVED lemma is amended, release at v1.24); it carries v1.23.0-pre.8 plus:
\begin{itemize}
\item \PFIX{} \cl{lem:bounds}: clause~(7) appended (the wide tier: under $2NB\le2^{63}-1$, storage in $64$ bits, products of two
  stored-class operands and determinants exact in $128$ bits, a fixed datum times a stored matrix in $64$ bits, the reach $N\le2^{20}-1$
  when $B=4N^{2}$), with its proof; clauses (1)--(6) and their proofs unchanged.
\item \PFIX{} \cl{thm:backend}: consequence~(4) (the stored wide tier; why a stored $128$-bit type would not extend it) and one proof
  sentence (the order of evaluation); (i)--(iv), (1)--(3) unchanged.
\item Audit repairs (\file{docs/implementation/slot24/ELLW\_audit.md}, F1): (4) states a sufficient condition (``under'', not
  ``iff''), says that (iii) is bounded through (iv), and names the checked, measured-only sites --- the (P2) input
  $\gamma_a^{-1}\gamma_b$ ($\le2B^{2}$) and the edge prefix, the (T1) conjugate ($\le2N^{2}B$), the level-structure consumers;
  Lemma~\ref{lem:bounds}(7)(c) is restricted to the level-one consumers and (7) gains the same list. The statement stays
  \PROVED{} with an inline \COMP{}, as Theorem~\ref{thm:backend}(iv) already does.
\item \cl{obl:backend}(ii): one sentence naming the wide tier's assertion (code~$52$) and the rollback (code~$51$ as before).
\item \cl{rmk:btabsmall}: the measured $B_{\mathrm{tab}}\approx0.88\,N$ past $23170$ and the zero escalations; \S\ref{sec:heckedatum}'s
  reach paragraph and the backend code map: the lift, the helpers, the lever \code{HECKE\_TAU\_WIDE=0}, the gates
  \code{HeckeWideTauValidation}, \code{HeckeWideTauSlowValidation}.
\end{itemize}
No label is added, renamed or removed; no existing clause number moves; no tag set or \code{\textbackslash LEAN} marker moves
($222$ on $86$ labels); the golden reference and the digests are unchanged (every value at $\ell\le23170$ is computed by the same
$64$-bit operations).

\subsection*{v1.23.0-pre.8 --- 2026-10-18 --- hecke-extensions slot 23: the G2-T2 and G2-P code-map sentences}

\noindent\emph{This version.} \textbf{A pre-release, not owner-released, not refereed}; it carries v1.23.0-pre.7 plus \textbf{code-map text only}:
\begin{itemize}
\item \S\ref{sec:obligations}, its code map: one sentence on G2-T2 (the SIMD $\rhon$ fill, rollback \code{HECKE\_RHO\_SIMD=0}; route~2's memo shared across
  chunks, rollback \code{HECKE\_ROUTE2\_SHARED\_MEMO=0}; audit \file{docs/implementation/slot23/G2T2\_audit.md}, PASS) and one on G2-P (the module route on the
  consumer pool, rollback \code{HECKE\_MODULE\_PARALLEL=0}, values independent of the thread count; audit \file{G2P\_audit.md}, PASS, fixes F1--F4 applied).
\end{itemize}
No statement, label, number, tag set or \code{\textbackslash LEAN} marker moves ($222$ on $86$ labels); every lever selects a route, never a value; the golden
reference and the digests are unchanged.

\subsection*{v1.23.0-pre.6 --- 2026-10-18 --- hecke-extensions slot 23: the repairs of C1230-pre3}

\noindent\emph{This version.} \textbf{A repair-only pre-release, not owner-released}; it carries v1.23.0-pre.5 (below) plus the fixes F1--F6 of
\file{docs/implementation/slot23/C1230\_pre3\_referee.md} (Opus, adversarial; the referee of pre.1--pre.3: ACCEPT-WITH-FIXES, BLOCKER 0, MAJOR 0, MINOR 2, NIT 4).
It sits on top of v1.22.0, \textbf{signed by the owner 2026-10-16} (fix F6: v1.22.0's own date line and changelog entry are left as released, ``owner sign-off
pending''; this entry and the \verb|\date| line record the signature, so that the chain reads consistently). \textbf{No label is added, renamed or removed; no
statement number, status tag or value moves; the marker count stays $222$ on $86$ labels} (none in \S\ref{sec:satprecisionsec}).
\begin{itemize}\setlength{\itemsep}{0pt}
\item F1 (MINOR), the proof of Lemma~\ref{lem:precisionreduction}, paragraph ``The integral complex'': the homotopy between the cochain maps of two integral
  $\tau$-tables is now stated over $C=\Ceq(\Z_\ell)$ with $\Upsilon=\widetilde H$ $\Z_\ell$-linear --- Proposition~\ref{prop:coneindep}(2) with
  Lemma~\ref{lem:tauexists}(2) over $\cring=\Z_\ell$, $6\in\Z_\ell^{\times}$, and Proposition~\ref{prop:coneindep}(3),(4) when the transporters $s_j$ or $g$ are
  rechosen --- in place of a homotopy ``of $\Ceq(\Z)$ with $\Upsilon$ integral'' (Proposition~\ref{prop:coneintegral}(4)), which the cited results do not
  supply over $\Z$. The conclusion only needs the homotopy over $C$ and its reduction modulo $\ell^{K}$, so the lemma is unchanged.
  Proposition~\ref{prop:coneintegral}(4)'s own wording of the ring of $\Upsilon$ in its $\tau$-half (older text) is not edited here.
\item F2 (MINOR), Lemma~\ref{lem:precisionreduction}(1), its last sentence, and the proof of (2): (2), (3) and the module structure over $\Z[X_p:p\in S]$ in
  (1) use that the $T_p$ commute on $M_K$, Theorem~\ref{thm:relations}(1), whose proof uses Theorem~\ref{thm:classical}; the sentence ``(2)--(5) use neither''
  now says so. The formalization note of \S\ref{sec:saturation} adds Theorem~\ref{thm:relations}(1) (its marker a \code{\textbackslash LEANlit} with hypothesis
  \code{hcomm}) and Proposition~\ref{prop:coneindep}(2)--(4) with Lemma~\ref{lem:tauexists}(2) in place of Proposition~\ref{prop:coneintegral}(4).
\item F3 (NIT), ``Why this section exists'': $K_{\mathrm{cls}}:=\max(2,e_{\max}+1)$ with $e_{\max}$ of Proposition~\ref{prop:torsionprepass} (which defines
  $e_{\max}$ only); $\mathrm{KMAX}$ defined as the route's level cap, a run setting (\code{LoewyOptions::kmax}, \code{SatOptions::kmax}, default $6$), and
  $K_{\mathrm{stab}}$ as the first level at which a block reads stable in the sense of Corollary~\ref{cor:stablecert}(1) (\code{LoewyBlock::K\_stab}).
\item F4 (NIT), Lemma~\ref{lem:precisionreduction}(3) and its proof: readings of the tower $(M_K,r_{K',K})$ that use the maps between levels are covered, by
  (1)'s naturality; Obligation~\ref{obl:precisionreuse}(iv) names the \code{LOEWY\_RING} fields computed from the listed readings (\code{D}, \code{rank\_T},
  \code{free\_level}, \code{free\_hilbert}, \code{hidden}, \code{cf9}, \code{cert}).
\item F5 (NIT, code, not [C]): \code{HeckeLevelConeValidation}'s two-transversal cell is retagged \code{(hecke\_two\_transversals)} and its two checks
  \code{(gated)}, no longer ``NEG'', matching Obligation~\ref{obl:levelconeexact}(ix) (a message change only; no golden, record or transcript).
\item Header and \verb|\date|; this entry.
\item The NITs N1--N3 of the check \file{docs/implementation/slot23/C1230\_pre6\_check.md} (ACCEPT-WITH-FIXES, NIT 3), applied within pre.6 (no version
  bump): N1, in F1's proof, $H$ of Lemma~\ref{lem:tauexists}(2) (additive) is replaced by the $\Z_\ell$-linear extension of its restriction to
  $C_{\bullet}(\Ku)\otimes1$, which keeps (T1) and the homotopy identity; N2, $K_{\mathrm{stab}}$ is described as the code computes it (the stopping rule read from
  $K=2$, on a torsion block from $K_{\mathrm{cls}}$; the field holds the last level read on a not-stable block and $0$ on \code{reduced\_rank1});
  N3 (catalogue, not [C]) the two-transversal cell of \code{HeckeLevelConeValidation} is no longer listed under NEG. No statement, label, number, tag or marker moves.
\end{itemize}

\subsection*{v1.23.0-pre.5 --- 2026-10-18 --- hecke-extensions slot 23: the G2-F code-map sentence}

\noindent\emph{This version.} \textbf{A pre-release, not owner-released, not refereed}; it carries v1.23.0-pre.4 (below) plus \textbf{code-map text only}, fix F1
of \file{docs/implementation/slot23/G2F\_audit.md} (Opus, adversarial; PASS-WITH-FIXES). \textbf{No statement's mathematics changes; no label is added,
renamed or removed; no statement number, status tag or value moves; the marker count stays $222$ on $86$ labels.}
\begin{itemize}\setlength{\itemsep}{0pt}
\item \S\ref{sec:obligations}, its code map: one sentence on Obligation~\ref{obl:tau} on a \code{TauTableContents::ChainsOnly} table (slot-23 G2-F, the module
  route's table): it carries no consumer table, so (vi) has no object there and its certification is vacuous; (i)--(v), (vii) run unchanged in the constructor;
  the module route evaluates $\rhon$ mod $\ell^{k}$ from each term's exact $\gamma_t$.
\item The status row of Obligation~\ref{obl:tau} in the roster: the same, in a few words, and \code{HeckeModuleStreamValidation} (its (T) arm) as a gate.
\item Header and \verb|\date|; this entry.
\end{itemize}

\subsection*{v1.23.0-pre.4 --- 2026-10-18 --- hecke-extensions slot 23: the G1 / G2 code maps}

\noindent\emph{This version.} \textbf{A pre-release, not owner-released}, in the 1.23.0 series (mid-slot practice: the X.Y.0 release candidate is cut at the
slot close); refereed by \file{docs/implementation/slot23/C1230\_codemap\_referee.md} (Opus, adversarial; ACCEPT-WITH-FIXES, MAJOR 1, MINOR 3, NIT 6, every fix applied). It carries the content of v1.23.0-pre.1 \dots{} pre.3 as described in their entries below, each with the referee
status stated there (pre.1 refereed by F-C123pre1; pre.2 its repairs; pre.3 audited, not refereed), plus \textbf{code-map text only} (the follow-up
``[C] v1.23.0 code-map sentences for G1 / G2'' of \file{docs/implementation/slot23/HANDOVER\_1018.md}; refereed as above). A release referee's coverage record, as
\file{docs/implementation/slot22/release\_referee\_coverage.md} was for v1.22.0, is due with the release candidate at the slot close. \textbf{No statement's mathematics changes; no label is added,
renamed or removed; no statement number, status tag or value moves; the marker count stays $222$ on $86$ labels.}
\begin{itemize}\setlength{\itemsep}{0pt}
\item \S\ref{sec:satprecisionsec}, its code map: \textbf{planned $\to$ built} (G1; \file{docs/implementation/slot23/G1\_impl\_audit.md} finding F8), with the
  clause-by-clause map of Obligation~\ref{obl:precisionreuse} onto \code{ringalg::LevelCache}, \code{ReduceToLevel}, \code{CheckDerivedModule},
  \code{CheckRunModule}, \code{RunLevelPolicy} and the rollback \code{HECKE\_PRECISION\_REUSE=0} (value-neutral).
\item \S\ref{sec:obligations}, its code map: one paragraph, the slot-23 G2 re-implementation of the module route over $\Z/\ell^{k}$ (G2-A, -B, -C, -D, -E, -G),
  which discharges Obligations~\ref{obl:welldef} and \ref{obl:agreement}(i) as written, with the rollback \code{HECKE\_MODULE\_STREAM=0} (value-neutral);
  \file{G2\_plan.md} \S2.7's sentence, widened as \file{G2AB\_report.md} \S6 item~5 and \file{G2CDE\_report.md} \S6 item~5 asked.
\item \S\ref{sec:saturation}, its code map: one sentence on G1 / G2 and one on the G2-X forms product (NTT and CRT) with the rollback
  \code{HECKE\_FORMS\_FASTMUL=0} (value-neutral).
\item \S\ref{sec:codemap}, the row of \S\ref{sec:saturation}: one pointer sentence.
\item Header and \verb|\date|; this entry.
\end{itemize}

\subsection*{v1.23.0-pre.3 --- 2026-10-16 --- hecke-extensions slot 23, G1: one module run per cell}

\noindent\emph{This version.} \textbf{A pre-release, not owner-released}; audited by \file{docs/implementation/slot23/G1\_clause\_audit.md} (Opus, adversarial; ACCEPT-WITH-FIXES, minor 7, nit 5, every fix applied in this entry's text), in the 1.23.0 series on top of v1.22.0 (signed by the owner 2026-10-16; that entry's text is left as released). Licenses lever G1 of \file{docs/implementation/slot23/ELL\_scout.md} \S3
(owner decision 2026-10-16; note \file{docs/implementation/slot23/G1\_clause\_note.md}). \textbf{Additive}: one subsection and two statements at the end of
\S\ref{sec:saturation}, so that no statement number moves; no existing statement, tag set, marker or value moves ($222$ markers on $86$ labels; the new
statements carry none).
\begin{itemize}\setlength{\itemsep}{0pt}
\item \PNEW{} \cl{sec:satprecisionsec}.
\item \PNEW{} \cl{lem:precisionreduction} (\PROVED{}): at $\ell\ge5$, level~$1$, $M_K=\Heq{1}(\Ceq(\Z/\ell^{K}))$ is $M_{K'}/\ell^{K}$ through the reduction map,
  Hecke-equivariantly, for $K\le K'$ (because $H^{2}(\Ceq(\Z_\ell))=0$); the presentation with exponents $\min(a_i,K)$; blocks and every isomorphism-invariant
  reading of the route agree; downward only; the hypothesis is needed (the long exact sequence with $H^{q+1}[\ell^{K}]$; $\Heq{0}$ on the Dickson locus
  is not a reduction).
\item \PNEW{} \cl{obl:precisionreuse}: scope, the restart rule, the derived module exactly, invariance of what is printed (recorded), the per-$K$ route as
  oracle and rollback arm with a byte-identity gate and two negative controls, record honesty for \code{K\_run} (recorded), certificates on derived
  modules.
\item The chapter's ``What this chapter fixes'' gains item (ix); its formalization note one sentence.
\item The audit's fixes (F1--F12, resolution in the audit file): (1)'s second sentence rests on the naturality of Remark~\ref{rmk:orbifold} in the ring and, for
  the classical $T_p$ on $\mathbb L$, on Theorem~\ref{thm:classical} (F1); the hypothesis is ``at each level an integral $\tau$-table'', through
  Proposition~\ref{prop:coneintegral}(4) applied to the bulk complex alone (F2, F11); (i) sends every failure of a run (refusal, exception, cap) to the per-$K$
  route (F3); (v) gains a digit-route cell (F4), negative controls with named checks (F5) and a tier (F7); (vi) the meaning of \code{K\_run} and the
  places to update (F6); (3), (iv) list every printed field (F8); the cochain map of (5) is $\varphi$ (F9); the pre.3 date line (F10); ``a new run'' per event (F12).
\item \textbf{Referee request}: a referee should read Lemma~\ref{lem:precisionreduction}(5)'s converse and its proof's identification of $\Th_g$ at level $K$
  with the base change of one integral cochain map first, then Obligation~\ref{obl:precisionreuse}(iv) against \code{ReadAlgebra} and \code{ReadDefect}.
\end{itemize}

\subsection*{v1.23.0-pre.2 --- 2026-10-12 --- hecke-extensions slot 23, LES-2: repairs of the referee report F-C123pre1}

\noindent\emph{This version.} \textbf{A repair-only pre-release, not owner-released.} Applies \file{docs/implementation/slot23/F-C123pre1\_referee.md} (Fable; verdict ACCEPT; BLOCKER 0, MAJOR 0, MINOR 2, NIT 5) to v1.23.0-pre.1.
\textbf{No label, statement number, tag set, marker or value moves} ($222$ markers on $86$ labels).
\begin{itemize}\setlength{\itemsep}{0pt}
\item MINOR-1: the ``Not edited here'' bullet of v1.23.0-pre.1 no longer quotes the (LES4) sentence as existing text; it is the audit's proposed addition.
\item MINOR-2: Obligation~\ref{obl:levelconeexact}(ix) now states what \code{HeckeLevelConeValidation} does: the two-transversal comparison on $\Ceq$ is \emph{gated} (cell tag \code{(NEG hecke\_fixed\_tables)}), the flat-cochain comparison is recorded, and neither is a negative control. (The tag's ``NEG'' is a naming matter for the LES-3 code pass; no code changed.)
\item NIT-1 (scalar step of the proof of (4) cites $\Delta_0(N_A)$), NIT-2 (class maps named by the indices of \eqref{eq:levelcorr}; both \code{BuildLevelCorrespondence} and \code{BuildLevelCorrespondenceChi}), NIT-3 (rule cross-reference; ``stored equivariant bases, $q=0,1,2$''), NIT-4 (chapter comment before \S\ref{sec:nebles}), NIT-5 (referee pointer): wording only.
\end{itemize}

\subsection*{v1.23.0-pre.1 --- 2026-10-12 --- hecke-extensions slot 23, LES-2: independence of the transversal}

\noindent\emph{This version.} \textbf{A pre-release, not owner-released; refereed by F-C123pre1 (ACCEPT), repaired at pre.2.} Applies the proposed edit of \file{docs/implementation/slot23/LES2\_audit.md} \S3.2 (Fable, MERGE after repairs) with one correction:
the audit's sentence ``on the flat cochains $C^{q}(K;L^{\chi}_A)$ the operator of \eqref{eq:levelcorr} does depend on the transversal'' is \textbf{not} carried over --- repair R1 of the package refuted it
(\file{LES2\_predictions.md} \S7, post-run note; \file{LES2\_report.md} \S3 item~2): the implemented operator uses $r_t$ only through its coset $A''r_t$, so it is transversal-independent on the whole flat
cochain space. \textbf{No label is added, renamed or removed; no statement number, existing tag set, marker or value moves} ($222$ markers on $86$ labels).
\begin{itemize}\setlength{\itemsep}{0pt}
\item Proposition~\ref{prop:levelconehecke}(4) \emph{replaced} (still \PROVED{}): ``Independently fixed tables'' becomes ``Independence of the transversal'', by (T1) of Definition~\ref{def:tau}; the strict square holds for every pair of transversals under (1);
  proof of (4) written out. The old text (the square commutes on cohomology only and a homotopy $\hder$ is needed) is withdrawn, together with its rule ``never (4)''.
\item Obligation~\ref{obl:levelconeexact}(ix): the control ``operators on independently fixed transversals'' (it can never fire) is replaced by the criterion-fails control at $(\Ga_0(25)\subset\Ga_0(5),\ \ell=7,\ n=0)$;
  the equal-on-$\Ceq$ record at $m_A>1$ is stated as a record, not a control.
\item \textbf{Not edited here:} the audit's \emph{proposed} addition to \code{rmk:lescheckable} (LES4) (``\dots{} with any transversals under the criterion, (4)''; not present text of (LES4)), and the erratum pointers in \file{LES0\_spec.md}~\S3.3 and \file{LES\_predictions.md}~\S2, are left to the LES-3 text pass.
\item \textbf{Refereed} (at pre.2): \file{docs/implementation/slot23/F-C123pre1\_referee.md}, ACCEPT, MINOR-1 / MINOR-2 applied. (Original request: a referee should read Proposition~\ref{prop:levelconehecke}(4) and its proof against the audit \S3.2 and \code{BuildLevelCorrespondence}, and check that nothing else in the document still relies on the withdrawn statement.)
\end{itemize}

\subsection*{v1.22.0 --- 2026-10-12 --- hecke-extensions slot 22, CLOSE-T: the slot-22 release candidate, from v1.22.0-pre.17}

\noindent\emph{This version.} \textbf{RELEASE CANDIDATE from v1.22.0-pre.17; the owner's sign-off is pending (CLOSE-T step T8) --- until it is recorded this text is not owner-released.} It releases the content of v1.22.0-pre.1 \dots{} pre.17 as described in their entries below, each with the referee status stated there; the gaps in referee coverage were read by the Fable release referee \file{docs/implementation/slot22/F-T1\_release\_referee.md} (ACCEPT on every gap; its repairs are v1.22.0-pre.17). \textbf{No statement's content changes; no label is added, renamed or removed; no status tag moves; the marker count stays $222$ on $86$ labels ($208$ \code{\textbackslash LEAN}, $14$ \code{\textbackslash LEANlit}).}
\begin{itemize}\setlength{\itemsep}{0pt}
\item Header and \verb|\date|; this entry.
\end{itemize}
Companion pointers: [E] 1.22.0, [U] 0.6.0, [L37] 0.5.2, [RC] 0.2.0, [MPT] 0.3.1, [CLI] 1.4.0, [P] 3.12.0, [PG] 1.1.0.

\subsection*{v1.22.0-pre.17 --- 2026-10-12 --- hecke-extensions slot 22: CLOSE-T T1 repairs (the F-T1 release referee)}

\noindent\emph{This version.} \textbf{A pre-release, not owner-released.} Repairs of \file{docs/implementation/slot22/F-T1\_release\_referee.md} (Fable, ACCEPT on every gap; BLOCKER~0, MAJOR~0, MINOR~2, NIT~15), the [C] items only. \textbf{No statement, number, label, marker, tag set or value moves.}
\begin{itemize}\setlength{\itemsep}{0pt}
\item \textbf{MINOR-1:} the subsection of v1.22.0-pre.9 (below), whose \verb|\date| line pointed at a missing changelog entry.
\item \textbf{Wording / punctuation (NIT-1, -2, -3, -4, -8, -9, -10):} the pre.2 changelog bullet on LN23F \S6 item~3; ``membership in $\mathfrak m^{2}$ is read modulo $p$'' and the reason for excluding $\ell=2$ in Remark~\ref{rmk:levelNcurvefamily}(1); the grammar after the \eqref{eq:negI} deletion; the order of two header comments of \S\ref{sec:mpmodel}; the negative controls of Example~\ref{ex:mpteardrop} (which break $d\circ d=0$, which the dimension gate); the full stop before ``(This generalises \dots'' in Lemma~\ref{lem:centralminusone}.
\item \textbf{Not applied:} NIT-5 (evidence wording at $\ell=5=p$), NIT-6 (``$M>1$ prime to $N$'' versus $M\nmid N$ in Proposition~\ref{prop:hidanormalsteps}) and NIT-7 (the preamble of Proposition~\ref{prop:levelconehecke}) change the scope of a statement and are left to slot 23; NIT-11 to NIT-13, NIT-15 concern [E] / [U].
\end{itemize}

\subsection*{v1.22.0-pre.16 --- 2026-10-12 --- hecke-extensions slot 22: the LR-Y phase-5 markers (the tree sequence in $\mathrm{Rep}\,k\,\Ga$, the Mayer--Vietoris sequence in group cohomology, naturality)}

\noindent\emph{This version.} \textbf{A pre-release, not owner-released, not refereed} as [C] text; the text applied is \S D of
\file{docs/implementation/slot22/LRY\_phase5\_report.md} as repaired by the Fable audit \file{LRY5\_audit.md} (ACCEPT, BLOCKER~0, MAJOR~0, MINOR~1, NIT~6; MINOR-1
names the identifications of [GT] \texttt{thm:mp-mv-terms} in two scopes), applied on the owner's approval of step LR-3 (2026-10-12). Lean side: the LR-Y package, phase 5
(\file{lean/GL2/MP/TreeSES.lean}, \file{lean/GL2/MP/TreeMV.lean}, [GT] rows \texttt{thm:mp-treeses-rep}, \texttt{thm:mp-mv-res}, \texttt{thm:mp-mv-les},
\texttt{thm:mp-mv-natural}; census \code{decls=2776 \dots\ verdict=OK}). \textbf{No new label; no statement, number, existing tag set or value moves.} No shipped value,
default, golden object, setting or console verb changes. Report: \file{docs/implementation/slot22/cpre16\_report.md}.
\begin{itemize}\setlength{\itemsep}{0pt}
\item \begin{sloppypar}\textbf{The LR-Y phase-5 markers}, $4$ \code{\textbackslash LEAN} on $1$ label that already carried markers, each with its own scope parenthesis (every
  [GT] row \code{proved}): \cl{cl:amalgam} ($1$ on (1): \texttt{thm:mp-treeses-rep}, \eqref{eq:treeses} as one short exact sequence in $\mathrm{Rep}\,k\,\Ga$ over every
  commutative ring; $3$ on (2): \texttt{thm:mp-mv-res}, the first map is $(\res_S,\res_U)$ through the bar-built comparison; \texttt{thm:mp-mv-les}, the whole sequence
  exact at every term, $q=0$ included, with a connecting map defined by transport and characterised on cocycles, not identified with mathlib's connecting map of a
  coefficient sequence; \texttt{thm:mp-mv-natural}, each of the three maps natural in $V$). So $222$ markers on $86$ labels in all ($208$ \code{\textbackslash LEAN},
  $14$ \code{\textbackslash LEANlit}).\end{sloppypar}
\item \textbf{Two existing marker texts shortened} (no statement text moves): the scope of \texttt{thm:mp-model7-tree-orbits} in \cl{cl:amalgam}(1) loses ``the sequence is
  not assembled as one object of $\mathrm{Rep}\,k\,\Ga$'' (now \texttt{thm:mp-treeses-rep}), and the ``Not in Lean'' sentence of \cl{cl:amalgam}(2) is removed (its three
  items --- restriction, the connecting map in \code{groupCohomology} terms, naturality in $V$ --- are the three new markers). The \texttt{thm:mp-tree-mv} scope (built on
  $\mathrm{Ext}$) and the \texttt{thm:mp-mv-low} scope (proved directly) stay as they are.
\item The marker-history paragraph before \S\ref{sec:labelpolicy} is brought to pre.16.
\paragraph{Other.} The header comment, the date line. Appendix~\ref{app:status} and the new-statement list are unchanged (no new label, no tag set moves).
\end{itemize}

\subsection*{v1.22.0-pre.14 --- 2026-10-11 --- hecke-extensions slot 22: the Lean dimension results at $\ell\in\{2,3\}$ (LV)}

\noindent\emph{This version.} \textbf{A pre-release, not owner-released, not refereed} as [C] text; worded within \S8 of the Fable audit
\file{docs/implementation/slot22/LV2\_audit.md} (\texttt{3d7b868d}, ACCEPT-WITH-REPAIRS) and the provenance audit \file{LVprov\_audit.md} (\texttt{a2ea84b5}). Lean side:
the LV-0--LV-2 packages ([GT] Part MP, sections LV-1 and LV-2, rows \texttt{thm:lv-f1} \dots\ \texttt{chk:lv3-cells}; census \code{decls=2683 \dots\ verdict=OK}).
\textbf{One new label; no existing statement text, number before it, tag set or value moves}; statement numbers after Remark~\ref{rmk:lvdimsltwothree} within
\S\ref{sec:six} shift (labels do not). No shipped value, default, golden object, setting or console verb changes. Report:
\file{docs/implementation/slot22/cpre14\_report.md}.
\begin{itemize}\setlength{\itemsep}{0pt}
\item \begin{sloppypar}\textbf{The LV markers}, $12$ \code{\textbackslash LEAN} on $4$ labels that carried none, each with its own scope parenthesis (every [GT] row
  \code{proved}, every one in the \code{GL2} library under [GT]'s census): \cl{lem:centralminusone} ($1$, at its last clause: \texttt{thm:lv3-tree-negone}, \emph{for a
  field} $K$ with $2\ne0$ and a finite-dimensional module, not for a ring with $2\in\cring^{\times}$), \cl{prop:oddnltwo} ($1$, at (2): \texttt{chk:lv-cells}, the value of
  $h^{0}$ at the odd $n\le63$ and $n=101,129,255$; not the formula, not the integral groups), \cl{prop:ellthreeclosed} ($1$, at (4): \texttt{thm:lv3-tree-negone}, for
  fields) and the new \cl{rmk:lvdimsltwothree} ($9$). So $218$ markers on $86$ labels in all ($204$ \code{\textbackslash LEAN}, $14$
  \code{\textbackslash LEANlit}). The cells $n>64$ of the library \texttt{L23Cells} are outside [GT]'s census and are cited as data
  (\file{docs/implementation/slot22/lv/}), never by a marker. \cl{prop:ltwothreescalar} carries none.\end{sloppypar}
\item New \cl{rmk:lvdimsltwothree} (\PROVED): $\dim_{\F_\ell}H^{q}(\Ga;\Sym^{n}\F_\ell^{2})$, $\ell\in\{2,3\}$, as ranks of explicit matrices (at $\ell=3$, $n$ even:
  $h^{q}=d-r_3-r_U^{-}$, $q\ge2$, and $h^{1}=2d-s-r_S^{+}-r_3$), formalised on mathlib's \code{groupCohomology} with the module of \eqref{eq:rhondef}, for every
  $n\le700$ and every $q$; a different model from the census's $\CI$, agreeing wherever both exist. It does not say ``Lean-verified record''.
\item The marker-history paragraph before \S\ref{sec:labelpolicy}, which stopped at pre.8, is brought to pre.14.
\paragraph{Other.} Appendix~\ref{app:status}: one line, counts $388\to389$; the header comment, the date line; the new-statement list of this appendix.
\end{itemize}

\subsection*{v1.22.0-pre.13 --- 2026-10-11 --- hecke-extensions slot 22: the LR-Y phase-4 markers (the Mayer--Vietoris terms, $\res-\res$, the natural periodicity)}

\noindent\emph{This version.} \textbf{A pre-release, not owner-released, not refereed} as [C] text; the text applied is \S6 of the Fable audit
\file{docs/implementation/slot22/LRY4\_audit.md} (\texttt{9feb57e2}, ACCEPT-WITH-REPAIRS, BLOCKER~0, MAJOR~1 (process), MINOR~5, NIT~6) verbatim, which corrects the proposal of
\file{LRY\_phase4\_report.md} \S D by MINOR-1, -2, -3 and NIT-1 (MINOR-4, NIT-2--4 concern [GT] and Lean, MINOR-5, NIT-6 the report, MAJOR-1 the process). Lean side: the LR-Y
package, phase 4 (\file{lean/GL2/MP/TreeCyclic.lean}, [GT] rows \texttt{thm:mp-cyc-coh} \dots\ \texttt{thm:mp-periodic-natural}; census \code{decls=2477 \dots\ verdict=OK}).
\textbf{No new label; no statement, number, existing tag set or value moves.} No shipped value, default, golden object, setting or console verb changes. Report:
\file{docs/implementation/slot22/cpre13\_report.md}.
\begin{itemize}\setlength{\itemsep}{0pt}
\item \begin{sloppypar}\textbf{The LR-Y phase-4 markers}, $4$ \code{\textbackslash LEAN} on $2$ labels that already carried markers, each with its own scope parenthesis (every [GT]
  row \code{proved}): \cl{cl:amalgam} ($3$, on (2): \texttt{thm:mp-mv-terms}, the periodic complexes compute mathlib's \code{groupCohomology} of $\langle S\rangle$,
  $\langle U\rangle$, $\langle-I\rangle$; \texttt{thm:mp-mv-resres}, under those identifications the middle map is $\res_S-\res_U$ with the sign as displayed;
  \texttt{thm:mp-mv-low}, the beginning $0\to V^{\Ga}\to V^{S}\oplus V^{U}\to V^{-I}$, proved directly), \cl{lem:bulkperiodic} ($1$, on (2):
  \texttt{thm:mp-periodic-natural}, \emph{a} periodicity isomorphism natural on \code{groupCohomology}, not shown to coincide with that of \texttt{thm:mp-tree-periodic} ---
  not $e\cup-$, not restriction-compatibility, not the $q=1$ clause). So $206$ markers on $82$ labels in all ($192$ \code{\textbackslash LEAN}, $14$
  \code{\textbackslash LEANlit}).\end{sloppypar}
\item \textbf{Two existing marker scopes rewritten} (no statement text moves): \cl{cl:amalgam}(2)'s \texttt{thm:mp-tree-mv} (the cyclic-group terms are no longer
  ``replaced by'' the periodic complexes but the sequence is ``built on'' them; the comparison is named, mathlib's \code{groupCohomologyIso} through $\mathrm{Ext}$; the
  ``Not in Lean'' list now holds only the map $H^{q}(\Ga;V)\to H^{q}(\langle S\rangle;V)\oplus H^{q}(\langle U\rangle;V)$ as restriction, the connecting map in
  \code{groupCohomology} terms and naturality in $V$), and \cl{lem:bulkperiodic}(2)'s \texttt{thm:mp-tree-periodic} (its ``not naturality on $H^{q}(\Ga;V)$ itself'' is
  dropped, since the new marker covers naturality --- for a periodicity isomorphism not shown to be the same one; LRY4 audit MINOR-1, -2).
\paragraph{Other.} The header comment, the date line. Appendix~\ref{app:status} and the new-statement list are unchanged (no new label, no tag set moves).
\end{itemize}

\subsection*{v1.22.0-pre.12 --- 2026-10-11 --- hecke-extensions slot 22: the LR-Y markers on the amalgam, the bulk periodicity and the tree resolution}

\noindent\emph{This version.} \textbf{A pre-release, not owner-released, not refereed} as [C] text; the text applied is the proposal of
\file{docs/implementation/slot22/LRY\_phase3\_report.md} \S C as audited by Fable (\file{LRY3\_audit.md}, \texttt{c3bbe552}, ACCEPT-WITH-REPAIRS, BLOCKER~0, MAJOR~0, MINOR~6,
NIT~6), with MINOR-1--6 and the [C] NITs~1--3, 5 applied (NIT-4, NIT-6 and MINOR-6 concern [GT] and plan files, not [C]). Lean side: the LR-Y package (phases 1--3,
[GT] Part~G12, \file{lean/GL2/MP/TreeRetract.lean} \dots\ \file{TreeCochains.lean}; census \code{decls=2344 \dots\ verdict=OK}). \textbf{One new label; no statement, number of an
earlier statement, existing tag set or value moves}; statement numbers after Remark~\ref{rmk:treeconeres} within \S\ref{sec:six} shift (labels do not). No shipped value,
default, golden object, setting or console verb changes. Report: \file{docs/implementation/slot22/cpre12\_report.md}.
\begin{itemize}\setlength{\itemsep}{0pt}
\item \begin{sloppypar}\textbf{The LR-Y markers}, $7$ \code{\textbackslash LEAN} on $3$ labels, each with its own scope parenthesis (every [GT] row \code{proved}):
  \cl{cl:amalgam} ($3$: (1) \texttt{thm:mp-treeses-exact}, \texttt{thm:mp-model7-tree-orbits}; (2) \texttt{thm:mp-tree-mv}, the long exact sequence with the cyclic-group terms
  replaced by their periodic cochain complexes and the middle map induced by $\varphi^{*}$ --- not $\res-\res$ for $q\ge1$, not the Shapiro identifications, not naturality),
  \cl{lem:bulkperiodic} ($1$: (2) \texttt{thm:mp-tree-periodic}, an isomorphism $H^{q}\cong H^{q+2}$, $q\ge2$, natural on the explicit complex --- not $e\cup-$, not the $q=1$
  clause, not restriction-compatibility; clause (1) carries none), \cl{rmk:treeconeres} ($3$: \texttt{thm:mp-model7-cone-res}, \texttt{thm:mp-tree-mp-comparison},
  \texttt{thm:mp-tree-cochains}). So $202$ markers on $82$ labels in all ($188$ \code{\textbackslash LEAN}, $14$ \code{\textbackslash LEANlit}).\end{sloppypar}
\item New \cl{rmk:treeconeres} (\PROVED): the mapping cone of the lift $\varphi$ of the tree boundary between the $2$-periodic resolutions is a free resolution
  $P_\bullet\colon\Z[\Ga]^{2}\leftarrow\Z[\Ga]^{3}\leftarrow\Z[\Ga]^{3}\leftarrow\cdots$ over every ring, its terms $2$-periodic from degree~$1$ and its differentials from
  degree~$2$; $\Hom_\Ga(P_\bullet,V)$ written out as $V^{2}\to V^{3}\to V^{3}\to\cdots$; the orientation $\mathfrak e\mapsto\mathfrak v_S-\mathfrak v_U$, the negative of
  \S\ref{sec:levelN}'s convention for $\mathcal Y$ (which moves no statement); its relation to Lemma~\ref{lem:treeres}.
\item \textbf{\cl{cl:amalgam}(1), the one change to an existing clause's wording (owner-visible; LRY3 audit MINOR-4).} The last sentence of its standard pointer
  (\cite[Ch.~I, \S4]{serre_trees}), ``A reader who prefers not to rely on it may take \eqref{eq:treeses} itself as the hypothesis; everything below uses the amalgam only through it'', is replaced by
  ``\eqref{eq:treeses} is proved independently of this pointer (the markers below), so nothing in this document rests on it; the amalgam decomposition itself is used nowhere
  else'': the trust basis of \eqref{eq:treeses} moves from the pointer to the Lean proof. The clause also gains, after \eqref{eq:treeses}, the description of the tree inside
  $Y$ (vertex set $\Ga/\langle S\rangle\sqcup\Ga/\langle U\rangle$, edge set $\Ga/\langle-I\rangle$, as Lemma~\ref{lem:bulkperiodic} cites it).
\paragraph{Other.} Appendix~\ref{app:status}: one line, counts $387\to388$; the header comment, the date line; the new-statement list of this appendix.
\end{itemize}

\subsection*{v1.22.0-pre.11 --- 2026-10-10 --- hecke-extensions slot 22: the F-L23 and F-MPT repairs, and the $Q$-reduced nerve (Gap~G1)}

\noindent\emph{This version.} \textbf{A pre-release, not owner-released, not refereed.} Inputs: the Fable referee F-L23 of \S\ref{sec:oddnltwo}
(\file{docs/implementation/slot22/F-L23\_referee.md}, \texttt{d3ef8b2d}, ACCEPT with two required and three recommended repairs) and the Fable referee F-MPT of [MPT]
(\file{docs/implementation/slot22/mp\_triangulations\_referee.md}, \texttt{54a55ec0}), whose \S3 proves the comparison theorem of the $Q$-reduced nerve. \textbf{Two new labels;
no existing tag set, \code{\textbackslash LEAN} marker or value moves}; statement numbers after Remark~\ref{rmk:mphybrid} within \S\ref{sec:mpmodel} shift (labels do not). No shipped
value, default, golden object, setting or console verb changes. Report: \file{docs/implementation/slot22/pre11\_repairs\_report.md}.
\begin{itemize}\setlength{\itemsep}{0pt}
\item F-L23 R1 (required): \cl{prop:oddnhecke}(4) --- ``equivalently'' replaced by the one-way implication: the \CONJ{} is the induced rank-one idempotent on $H^{2}(\Ga;\Vcoef)$;
  the stronger form on $H^{1}(\Ga;\Mn)$ implies it and is not asserted equivalent. F-L23 R2 (required): \cl{prop:ltwothreescalar} --- ``the whole prime-to-$\ell$ Hecke algebra''
  replaced by ``the algebra generated by all $T_p$''.
\item F-L23 R3: (I1)'s independence of the representatives proved in place as a cochain identity, the pointers aligned with \cl{prop:barformula}'s; R4: (F1)'s module
  identification $H^{1}(\Ga(3);\F_3)\cong\Sym^{2}\F_3^{2}$ derived from the cusp count; R5: Step~0(b) says $d_p=1$ and the pair is the identity.
\item New \cl{def:mpqd} ($Q$-domination (QD), the dominant vertex $\mathsf p(x)$; holds on $\mathsf M_7$ and $\mathsf R_{\Z}$ at $\ell\in\{2,3\}$) and \cl{prop:mpqreduced}
  (\PROVED, \COMP\textsuperscript{a}): the choice-free quotient $\mathfrak D=\mathfrak F'/{\sim}$ is a projective resolution computing $H^{\bullet}(\Ga;M)$ with the coefficient system
  $M^{Q}$ and the twisted differential, without adaptedness and without lifts; \cl{rmk:mphybrid}(3) rewritten from ``Gap G1, recorded'' to ``closed''.
\paragraph{Other.} Appendix~\ref{app:status}: two lines, counts $385\to387$; the header comment, the date line; the new-statement list of this appendix.
\end{itemize}

\subsection*{v1.22.0-pre.10 --- 2026-10-10 --- hecke-extensions slot 22: odd $n$ at $\ell=2$, and $T_p=1+p$ at $\ell\in\{2,3\}$ (L23-0 extended)}

\noindent\emph{This version.} \textbf{A pre-release, not owner-released, not refereed} (the Fable referee F-L23 is pending). Inputs: the L23 master plan
\file{docs/implementation/slot22/L23\_master\_plan.md} v0.2 and its Fable review \file{L23\_master\_plan\_review.md} (\texttt{dd3254f3}; only what it marks PROVED enters
as a theorem); the census reviews \file{census\_l23\_classification\_review.md} (with its erratum \S7) and \file{census\_l23\_plan\_v02\_review.md}; the owner's sign-off of
the plan's seven decisions (2026-10-10). \textbf{New subsection \S\ref{sec:oddnltwo}, eleven new labels; no existing statement, tag, \code{\textbackslash LEAN} marker or value
moves}; statement numbers after \S\ref{sec:oddnltwo} within \S\ref{sec:six} shift (labels do not). No shipped value, default, golden object, setting or console verb changes.
Report: \file{docs/implementation/slot22/L23\_0\_report.md}.
\begin{itemize}\setlength{\itemsep}{0pt}
\item \cl{lem:centralminusone} (\PROVED): $-I\mapsto-1$ forces $2H^{q}=0$ for every module and degree, $\Z/2^{k}$ coefficients included (generalises
  \cl{prop:badprimes}(2)); \cl{lem:bocksteinallell} (\PROVED): \cl{cl:bockstein}(2)'s Bockstein isomorphism at every prime, and its first isomorphism exactly where
  $H^{2}(\Ga;\Vcoef)[\ell]=0$ (false at $\ell=2$, $n$ odd); \cl{lem:cyclicres} (\PROVED): cyclic cohomology and restriction.
\item \cl{prop:oddnltwo} (\PROVED; the review's Theorem~O(2)--(3) with its self-contained proof): Mayer--Vietoris collapse, $c_U=2h^{0}$,
  $h^{0}=\lfloor(\lfloor n/3\rfloor+1)/2\rfloor$, $H^{2j+1}\cong\F_2^{h^{0}}$, $H^{2j+2}\cong\F_2^{(n+1)/2-h^{0}}$, periodicity from degree $1$, the odd-degree classes as
  $e^{j}\cup\beta$ of Dickson invariants; \cl{lem:oddnsplit} (\PROVED): $H^{q}(\Ga;\Vcoef/2^{k})\cong H^{q}\oplus H^{q+1}$ Hecke-equivariantly for $k\ge2$.
\item \cl{lem:tellrank} (\PROVED): $\rank T_\ell\le2$ through $H^{q}(\Ga_0(\ell);\Fl)$, of dimensions $1,2,2,\dots$ ($\ell=2$) and $1,2,1,\dots$ ($\ell=3$);
  \cl{prop:ltwothreescalar} (\PROVED{} given (F1) and (I1), both \BSTD{} with section-level pointers \textbf{not} verified): $T_p=(1+p)\id$ on $H^{q}(\Ga;\Sym^{n}\Fl^{2})$,
  every $q$, $\ell\in\{2,3\}$, $p\ne\ell$ --- the review's input (F2) (a Shapiro--Hecke compatibility) is not needed: the proof uses the level-$\Ga(\ell)$ coset
  operator directly; \cl{prop:oddnhecke}: $T_p=0$ ($p$ odd) on $H^{q}(\Vcoef)$ and $H^{q}(\Vcoef/2^{k})$, $T_2=0$ in odd degrees (\PROVED), $T_2$ a rank-one idempotent in
  even degrees (\CONJ).
\item \cl{lem:cprule} (\PROVED, the census review's $c_p=\min(v_\ell(p^{n}-1),\min\mathrm{ed}_{\mathrm{eis},p})$); \cl{rmk:standardtwist} (\PROVED, with the $p=2$ clause);
  \cl{rmk:oddnltwoerr} (the [R]~B(6) erratum and the list of what is deliberately not stated: the Gouv\^ea bound, the $\ell=3$ $\pm$ law, $e_1=v_\ell(n)+1$, the
  $T_\ell$ idempotent for even $n$, Hodge--Newton).
\end{itemize}
\paragraph{Other.} Appendix~\ref{app:status}: eleven lines, counts $374\to385$; the header comment, the date line; the new-statement list of this appendix.

\subsection*{v1.22.0-pre.9 --- 2026-10-10 --- hecke-extensions slot 22: \S\ref{sec:mpmodel}, code-vs-chapter wording (a)--(e)}

\noindent\emph{This version.} \textbf{A pre-release, not owner-released.} Wording only, five sites: Remark~\ref{rmk:mprulesread} (the goodness bullet: \code{good\_fails\_at} R$_{\mathbb Z^{+}}$ at $m$, R$_{\mathbb Z}$ at $v_1,v_2,m$; the finite-group and strong forms agree on the four oracle rows, \code{good\_forms\_agree}; the \COMP$^{r}$ sentence: $4\times3=12$ gated cells, the R$_\zeta$ row reported by the ungated \code{OUTPUT P1.12} line), Example~\ref{ex:mpteardrop} (the three negative controls named as built), \ref{MPC3} (only the $H^{2}$ change is gated; the $H^{0}$ / $H^{1}$ / $\mathbb Z/\ell^{k}$ invariance is an ungated observation). No statement, number, label, tag, marker or value moves; not refereed at pre.9 (read by F-T1 at the close, \file{docs/implementation/slot22/F-T1\_release\_referee.md}, \S3).

\subsection*{v1.22.0-pre.8 --- 2026-10-10 --- hecke-extensions slot 22: the pre.7 referee's repairs, the MP-4 Hecke remark, the MP-3b records of \S\ref{sec:mpmodel}}

\noindent\emph{This version.} \textbf{A pre-release, not owner-released, not refereed.} Inputs: the Fable referee of pre.7 (\file{docs/implementation/slot22/C\_pre7\_referee.md},
ACCEPT, four minor findings F1--F4); the slot-22 code package MP-4 (\file{MP4\_report.md}, merged at \texttt{eac3f3b3}) and its audit F-MP4 (\file{F-MP4\_audit.md}, ACCEPT WITH
REPAIRS, \S3 the contract sentences); the package MP-3b (\file{MP3\_report.md} \S6, merged at \texttt{ec982f93}) and its audit F-MP3b (\file{F-MP3b\_audit.md}, ACCEPT WITH REPAIRS,
R1--R3). \textbf{No statement number, label of an existing statement, \code{\textbackslash LEAN} marker or value moves}; one tag set grows (\COMP\textsuperscript{r}); one item label
is added (\ref{MPC12}); no shipped value, default, golden object, setting or console verb changes. Map of every edit: \file{docs/implementation/slot22/C\_pre8\_report.md}.
\begin{itemize}\setlength{\itemsep}{0pt}
\item \textbf{The pre.7 referee, F1--F4} (all minor, the referee's exact text): the proof of Proposition~\ref{prop:mpbredon}, ``clause (2) below holds for $\Phi^{\ast}$ as well'' (F1, three
  consecutive ``(2)''); Proposition~\ref{prop:mprescor}(1), ``any free $\Ga$-set resolving $\Z$ as in Theorem~\ref{thm:mpcomparison}(2) (free over every subgroup)'' in place of the empty
  qualifier ``free $A$-set that is also free over $B$'' (F2); Definition~\ref{def:mprules}(3), the stabiliser of a point lies in that of a vertex of its carrier ``under (IO)'' (F3); the
  proof of Theorem~\ref{thm:mpnatural}(4), ``($Y$ being connected by (AC))'' (F4, optional, applied). The referee's optional sentence of its \S6 is not adopted (the counts carry it).
\item \textbf{MP-4: Remark~\ref{rmk:mphecke} rewritten} (F-MP4 \S3 verbatim, F-MP4 F11 and R3): the route described at pre.7 (lifting the shipped $\tau$-tables through
  $\mathrm{Tot}\to C_\bullet(\Xbulk)$) was not taken, and its plan (slot 23, $T_2,T_3,T_5$ on $H^{\le3}$ against \code{IntegralHecke}) is superseded. Now: a $\conj_g$-twisted chain map
  on the MP resolution $N\Z[\mathfrak F'_\bullet]$ itself, built with the MP model's own filling, \emph{unique up to twisted homotopy over every commutative ring because
  $N\Z[\mathfrak F']$ is free over $\Z[\Ga'']$} --- no hypothesis $6\in\cring^{\times}$, unlike Lemma~\ref{lem:tauexists}(2) (F-MP4 F1, ``the new feature, unstated''); the label factor
  $\rho_n(g)$; the comparison $T_p\Phi^{\ast}=\Phi^{\ast}\hat\Theta_g$ over every ring; the $g$-twist not a labelled morphism, G3 not needed to compute. A notation sentence ties the audit's
  symbols to [C]'s ($\Phi^{\mathsf{MP}}$, $\Th_{\Iso,g}$, $\Ga''=\Ga_0(p)$) and replaces its undefined ``Section~sec:hecke'' by \S\ref{sec:chainlevel}. A \COMP\textsuperscript{r} record
  follows the remark (\code{OrbifoldMPHeckeValidation}, $2059$ checks) with the non-vacuous counts of F-MP4 R1 ($354/592$, $254/464$, $66/108$), the degree-$2$ caveat of F-MP4
  F7 and the not-gated list (F9: $H^{3}$, $T_{p^{\nu}}$, $\Z/\ell^{k}$ charpolys, level $N$, $j^{\ast}$, the $\iota$-sign, a second filling); the remark's tag: commentary
  $\to$ \COMP\textsuperscript{r}; \ref{MPC12} added. Definition~\ref{def:mpmorphism}(3): the $g$-twist's label factor is $\rho_n(g)$ on a twisted chain map of one $\mathfrak F'$, not a
  morphism of MP data. The scope paragraph, the code map and the status table of \S\ref{sec:implementation}: MP-4 built (``builds no Hecke datum'' $\to$ no \emph{canonical geometric}
  one).
\item \textbf{MP-3b} (\file{MP3\_report.md} \S6.4 with F-MP3b R1, R2): \ref{MPC11} $537\to723$ checks, with the pair sequence over $\Z/4$, $\Z/9$, $\Z/25$ against an integral oracle, the
  pairs on $\mathsf R_{\Z}$ and $\Ga_1(8)<\Ga_1(4)$ with $-I\notin A$; \ref{MPC10} $331\to332$ (the oracle-file digest); the ``not gated'' clauses for $\Z/\ell^{k}$, $\mathsf R_{\Z}$ and
  $-I\notin A$ removed from \ref{MPC11}, from the \COMP\textsuperscript{r} record of Proposition~\ref{prop:mprescor} (``no label $-I$ occurs'' $\to$ ``labels with $u=-I$ occur
  \emph{only} when $-I\notin A$, and do occur, $60$ of $120$'', R2) and from the sentence after the list (the stale ``Remark~\ref{rmk:mphecke} (MP-4)'' remainder dropped, R2); the
  negative control for the $u=-I$ branch called structurally guaranteed (R1). R3 is a docstring of \file{mp3b\_oracle.py} (output unchanged), not a [C] matter.
\end{itemize}
\paragraph{Other.} Appendix~\ref{app:status}: \ref{rmk:mphecke} gains \COMP\textsuperscript{r}, \COMP{} $56\to57$ (registered $28\to29$), tagged $245\to246$. \S\ref{sec:status} (no
marker added), the header comment, the date line, the chapter's header comment, the label index (\ref{MPC12}); pre.7 is recorded as refereed at pre.8.

\subsection*{v1.22.0-pre.7 --- 2026-10-10 --- hecke-extensions slot 22: the MP-L3 Lean markers, the unnormalised clauses, the MP-3 records of \S\ref{sec:mpmodel}}

\noindent\emph{This version.} \textbf{A pre-release, not owner-released, not refereed at pre.7} (refereed at pre.8 by Fable, \file{C\_pre7\_referee.md}, ACCEPT; the four minor
repairs F1--F4 applied there). Inputs: the slot-22 Lean package MP-L3 (\file{docs/implementation/slot22/MPL3\_report.md}
v1.0, merged at \texttt{c5617bc7}; [GT] Part~G12), the slot-22 code package MP-3 (\file{MP3\_report.md}) and its audit F-MP3 (\file{F-MP3code\_audit.md}, ACCEPT WITH
REPAIRS, \S10 the contract sentences). \textbf{No statement number, label of an existing statement, or value moves}; two tag sets grow (\COMP\textsuperscript{r}); no shipped value,
default, golden object, setting or console verb changes. Map of every edit: \file{docs/implementation/slot22/C\_pre7\_report.md}.
\begin{itemize}\setlength{\itemsep}{0pt}
\item \begin{sloppypar}\textbf{The MP-L3 markers}, per MPL3 report \S5 (every [GT] row \code{proved}): $25$ \code{\textbackslash LEAN} --- \cl{def:mpmorphism} ($4$: (1),(2) \texttt{def:mp-morph},
  \texttt{def:mp-orbmap}; (3) \texttt{lem:mp-morph-simplicial}, \texttt{def:mp-pull}), \cl{thm:mpnatural} ($10$: (1) \texttt{thm:mp-pull-cochain}, \texttt{lem:mp-pull-comp},
  \texttt{def:mp-zu}; (2) \texttt{thm:mp-pull-natural}, \texttt{def:mp-barmap}, \texttt{thm:mp-homotopy-exact}; (3) \texttt{thm:mp-refine-invariance}, \texttt{thm:mp-label-indep};
  (4) \texttt{thm:mp-untwisted}, \texttt{thm:mp-rz-untwisted}) and its \COMP\textsuperscript{r} record ($3$: \texttt{def:mp-rz-refine}, \texttt{thm:mp-rz-refine},
  \texttt{thm:mp-rz-untwisted}), \cl{prop:mpbredon} ($3$: (1) \texttt{def:mp-epsmp}, \texttt{thm:mp-epsmp-canonical}; (2), the ``if'' only, \texttt{thm:mp-epsmp-quasiiso}),
  \cl{def:mpset} ($1$: \texttt{thm:mp-norm-dgn}), \cl{def:mpfree} ($2$: \texttt{def:mp-simpK}, \texttt{thm:mp-simpK-alt}), \cl{thm:mpcomparison} ($2$: (2)
  \texttt{thm:mp-norm-heq}, (4) \texttt{thm:mp-comparison-normalised}). \texttt{thm:mp-dd}'s scope note ``not the normalised complex'' is dropped. So $195$ markers on
  $79$ labels in all ($181$ \code{\textbackslash LEAN}, $14$ \code{\textbackslash LEANlit}).\end{sloppypar}
\item \textbf{The unnormalised clauses} (MPL3 report \S4.1). Lean proves Theorem~\ref{thm:mpnatural}(1)--(3) for $C^{\bullet}_{\mathrm{un}}$; [C] now states them there too:
  Definition~\ref{def:mpmorphism}(3) defines $u^{\ast}$ on all strings, the normalised map being its restriction; Theorem~\ref{thm:mpnatural}(1), (2), (3) gain ``likewise for
  $C^{\bullet}_{\mathrm{un}}$'', (2) also that the inclusions $C^{\bullet}\subseteq C^{\bullet}_{\mathrm{un}}$ are quasi-isomorphisms commuting with $u^{\ast}$; the proofs gain one sentence
  each ($u$ commutes with the degeneracies; the projection $\Z[\mathfrak F']\to N\Z[\mathfrak F']$ is a map of free resolutions over $\id_\Z$). The markers are scoped to the
  unnormalised clauses; the normalised $u^{\ast}$ as a cochain map is not marked.
\item \textbf{Refinement certified by a vertex map} (\S4.2). Definition~\ref{def:mpmorphism}(2): the table need not come from a simplicial map; an equivariant map of vertex sets with the
  containments certifies it (how Lean certifies $\mathsf R_{\Z}\to\mathsf M_7$). Proof of Theorem~\ref{thm:mpnatural}(1): ``non-empty to non-empty'' is all that is used.
\item \textbf{Label independence} (\S4.3): [C]'s chain-homotopy claim in Theorem~\ref{thm:mpnatural}(3) stays (its proof, the comparison theorem, covers it); the marker is scoped
  to cohomology.
\item \textbf{$\varepsilon^{\mathsf{MP}}$} (\S4.4). Proposition~\ref{prop:mpbredon}(1) gains: for any lift $\Phi$ of $\id_\cring$ from $\cring[\mathfrak F'_\bullet]$ into
  $C_\bullet(\Xbulk)\otimes\cring$, $\Phi^{\ast}$ is homotopic to $\varepsilon^{\mathsf{MP}}$ followed by $C^{\bullet}\subseteq C^{\bullet}_{\mathrm{un}}$ (the proof: the same uniqueness up to
  homotopy, over $\cring$); so the Lean $\varepsilon^{\mathsf{MP}}=\Phi^{\ast}$ and the written $(\Phi^{\mathsf{MP}})^{\ast}\circ\cmpI^{\ast}$ agree up to homotopy, and (2) holds for both.
\item \textbf{MP-3} (F-MP3 \S10). Proposition~\ref{prop:mprescor}(1): the identity holds on any free $A$-set free over $B$, in particular on $\mathfrak F'_{\Ga}$ (the code's
  choice), with $\res$, $\cor$ written there and one proof sentence; Proposition~\ref{prop:mpboundary}(2): the relative indexing of the cone and its identification with
  $\Cone{\rres}[-1]$; \COMP\textsuperscript{r} records after both propositions (\code{OrbifoldMPBoundaryValidation}, \code{OrbifoldMPLevelConeValidation}), their tags
  \PROVED{} $\to$ \PROVED{} \COMP\textsuperscript{r}; Remark~\ref{rmk:mpcheckable} gains \ref{MPC10}, \ref{MPC11} (two new item labels) and says what is not gated: the level pairs
  over $\Z/\ell^{k}$, on $\mathsf R_{\Z}$, and with $-I\notin A$. The audit's optional sentence for Theorem~\ref{thm:mpnatural}(4) (``non-central for every pair with
  $[\Ga:B]>2$'') is not adopted: $\Ga^{2}<\Ga$ has index $2$ and non-central labels; the record of \ref{MPC11} states the counts instead.
\end{itemize}
\paragraph{Other.} Appendix~\ref{app:status}: two lines gain \COMP\textsuperscript{r}, \COMP{} $54\to56$ (registered $26\to28$). The code map and the status table of
\S\ref{sec:mpmodel}: MP-3 built. The formalization note of \S\ref{sec:mpmodel}, \S\ref{sec:status}, the header comment, the date line, the chapter's header comment, the label index
(\ref{MPC10}, \ref{MPC11}).

\subsection*{v1.22.0-pre.6 --- 2026-10-10 --- hecke-extensions slot 22: the F-MP2 referee repairs and the first Lean markers of \S\ref{sec:mpmodel}}

\noindent\emph{This version.} \textbf{A pre-release, not owner-released.} \S\ref{sec:mpmodel} is refereed by F-MP2 (contract part, Fable,
\file{docs/implementation/slot22/F-MP2\_contract\_referee.md}: ACCEPT WITH REPAIRS, one blocking finding), and its repairs R1--R8 are applied as given there;
the slot-22 Lean packages MP-L1 / MP-L2 (\file{docs/implementation/slot22/MPL1\_report.md} v1.1) supply one wording repair and the markers. \textbf{No statement number,
label, tag set, status or value moves}; no shipped value, default, golden object, setting or console verb changes. Map of every edit:
\file{docs/implementation/slot22/C\_pre6\_report.md}.
\begin{itemize}\setlength{\itemsep}{0pt}
\item \textbf{R1 (blocking), Theorem~\ref{thm:mpnatural}(4).} The first sentence was false: labels acting by $\pm1$ do \emph{not} make $f\mapsto f\circ u$ a cochain map (mixed labels
  $\pm I$ fail at odd $n$, the \code{REFINE\_TWIST} control); the sufficient condition is that all labels act by one and the same central automorphism. Proof: one clause
  (for labels $\pm I$, iff the labels are constant or $n$ is even). The clause's last sentence and the theorem's tag are unchanged.
\item \textbf{R2, Proposition~\ref{prop:mpboundary}(2).} The proof cited Theorem~\ref{thm:mpnatural}(2), which compares resolutions of $\Z$ over $\id_{\Z}$; it now lifts the
  augmentation $\Z[\Ga/\Gbdry]\to\Z$, takes $\Ext_{\Z[\Ga]}(\varepsilon,M)$ and identifies it with $\res^{\Ga}_{\Gbdry}$ by Shapiro, \STD{} Weibel Thm.~2.2.6 and
  Brown Ch.~III.9 (\VERIFYloc{}, added to the list of the pre.5 item); its line in Appendix~\ref{app:status} reads \cite[Ch.~III.6.2, III.9]{brown}.
\item \textbf{R3} Definition~\ref{def:mptri}(3): (PD) is now pure dimension (as written it held in every finite-dimensional complex). \textbf{R4} Definition~\ref{def:mpphi}: the choices of
  $h$ and $c$ are named, the homotopy class is choice-free, the code's choice is \code{PhiLift}'s. \textbf{R5} Remark~\ref{rmk:mphybrid}(1): the hybrid ``is an $\ell$-local isotropy
  model in the sense of'' Definition~\ref{def:isomodelell} with any projective column (was ``is the''). \textbf{R6} Remark~\ref{rmk:mpwhat}(3): the coset sentence
  (adaptedness nests the shared stabilisers in $\Ga_{v(\tau)}$); its \COMP\textsuperscript{r} paragraph says which count is $\mathsf S$'s ($132,2800$) and which the
  $\mu$-form's ($124,2704$). \textbf{R7} the symbol $D$ of Proposition~\ref{prop:mpalldegrees}(2) declared in the notation paragraph.
\item \textbf{R8, the eight hygiene items.} Proposition~\ref{prop:mpboundary}(1) ``the proof of Theorem~\ref{thm:mpcomparison}, parts (1)--(2)'' and (2) ``computes the relative
  object of Definition~\ref{def:borelmodel}''; ``Lemma~'' in the proof of Theorem~\ref{thm:mpcomparison}(5); the cyclicity of the finite subgroups of $\SL_2(\Z)$ argued in
  Remark~\ref{rmk:mprulesread}(3); the label coset $\{h,h\cdot(-I)\}$ in the \COMP\textsuperscript{r} paragraph after Theorem~\ref{thm:mpnatural}; Definition~\ref{def:mpmorphism}(1)
  ``recorded on the descriptors by a label table \dots, which is all that the induced maps of (3) use''; the bound group element $u$ renamed $w$ in the Reading paragraph after
  Theorem~\ref{thm:mpcomparison} and in Definition~\ref{def:isomodelell}'s $e_\sigma$ (a bound variable; no statement changes); the header comment and the pre.5 item now
  say ``refereed at pre.6''.
\item \textbf{Definition~\ref{def:mprules}(3), the strong form of goodness} (MPL1 report \S4.1). Pre.5 wrote $\mathsf T(v)\subseteq\Ga_v$; the
  C++ (\code{Good}, \code{good\_strong}), \file{mp\_rules.py}, Table~\ref{tab:mprules}'s ``fails at'' column and Lean all use ``$\mathsf T(v)$ in the stabiliser $\Ga_p$ of one
  vertex $p$ of $\mathrm{St}(v)$''; the written form fails at every vertex orbit of $\mathsf M_7$ and at $v_0,v_1,v_2,m$ on $\mathsf R_{\Z}$ (Lean: \code{goodVtx\_discrepancy}). The definition now states the
  form in use; Remark~\ref{rmk:mprulesread}(3)'s ``the finite-group and the strong form agree'' is true for it and now says so. A wording repair of a definition: no value,
  table entry, tag or validation changes.
\item \textbf{The unnormalised MP complex $C^{\bullet}_{\mathrm{un}}$} (MPL1 report \S4.2) is named in Definition~\ref{def:mpset}, and Theorem~\ref{thm:mpcomparison}(3), (4) state it
  beside the normalised one (the proof's argument verbatim on the resolution $\Z[\mathfrak F']$ of (2)), so that the markers are clause-scoped.
\item \textbf{The first \code{\textbackslash LEAN} markers of \S\ref{sec:mpmodel}}, per MPL1 report \S5 ([GT] Part~G12, every row \code{proved}): $30$ on $11$ labels ---
  \cl{def:mprules} ($2$: \texttt{def:mp-datum}, \texttt{thm:mp-goodvtx}), \cl{def:mpell} ($1$), \cl{def:mpRZ} ($2$: \texttt{def:mp-models}, \texttt{def:mp-rz-tau}),
  \cl{tab:mprules} ($3$, rows $\mathsf M_7$, $\mathsf R_{\Z}$), \cl{def:mpset} ($4$), \cl{rmk:mpwhat} ($2$, item (1)), \cl{def:mpfree} ($1$), \cl{lem:mprows}, \cl{lem:mpcols},
  \cl{lem:mpassembly} ($1$ each; the last in its acyclic form) and \cl{thm:mpcomparison} ($12$: (1) \texttt{lem:mp-free}, \texttt{lem:mp-set-finite}; (2)
  \texttt{thm:mp-resolution}, \texttt{thm:mp-projres}; (3) \texttt{thm:mp-homequiv}; (4) \texttt{thm:mp-comparison}, \texttt{thm:mp-comparison-natural}; (5)
  \texttt{thm:mp-models-cover}, \texttt{thm:mp-ac-bridge}, \texttt{thm:mp-model7-ac}, \texttt{thm:mp-retr}, \texttt{thm:mp-rz-ac}, i.e.\ (AC) for $\mathsf M_7$ and $\mathsf R_{\Z}$,
  \ref{H4}(ii) of the subdivision included, and the comparison on both with no hypothesis left). So $170$ markers on $76$ labels in all ($156$ \code{\textbackslash LEAN},
  $14$ \code{\textbackslash LEANlit}); \S\ref{sec:status} names the list, the formalization note of \S\ref{sec:mpmodel} says what is not marked, \ref{MPC3} that (AC) is now proved in Lean.
\end{itemize}
\paragraph{Other.} Appendix~\ref{app:status}: no line or count moves (one locator, R2). Non-normative: the header comment, the date line, the chapter's header comment.

\subsection*{v1.22.0-pre.5 --- 2026-10-09 --- hecke-extensions slot 22: the Moerdijk--Pronk model (MP-CT)}

\noindent\emph{This version.} \textbf{A pre-release, not owner-released, not refereed at pre.5} (refereed at pre.6 by F-MP2, contract part, ACCEPT WITH REPAIRS; repairs R1--R8 applied there). Sources:
\file{docs/implementation/slot22/MP\_math\_plan.md} v0.2.1 (audited, \file{MP\_math\_plan\_audit.md}, repairs R1--R8), \file{MP\_impl\_plan.md} v0.2.1 \S4 and its review
(R1, R2, R5, R6), the owner rulings Q1--Q7 of 2026-10-09 (at the recommendations: Q1 the choice-free set is normative, Q2 $\mathsf R_{\Z}$ MP-conformant and $\mathsf M_7$
economical, Q3 the $\ell$-relevant structure, Q4 no shipped value or default changes, Q5 Lean L0--L10 as the M6 ``orbifold'' item, Q6 the \ref{rmk:negI} erratum of pre.4, Q7 no
teardrop in Lean), the F-MP2 code audit \file{F-MP2\_code\_audit.md} \S1.3 (the label factor $\rho(h_{\tilde\sigma_0})$) and the built code and validations
(\file{MP1a\_report.md}, \file{MP1b\_report.md}, \file{MPT\_report.md}, \file{MP2\_report.md}). \textbf{No existing statement number, label, tag set or status moves; the
marker count stays $140$ on $65$ labels} (the new labels carry none; no Lean module of the lane exists yet). \textbf{No shipped value, default, golden object, setting or console verb
changes}; $\F_2$ and $\F_3$ keep the Bredon semantics of Definition~\ref{def:ringsemantics}.
\begin{itemize}\setlength{\itemsep}{0pt}
\item \textbf{New: \S\ref{sec:mpmodel}}, a chapter after \S\ref{sec:isomodelell}, before the appendices (the implementation plan named the place after
  \S\ref{sec:orbifoldmodel}; placed last, as \S\ref{sec:isomodelell} was, so that no chapter number moves): Definitions~\ref{def:mptri}, \ref{def:mprules}, \ref{def:mpell},
  \ref{def:mpRZ} and Table~\ref{tab:mprules} --- orbifold triangulations as descriptors, Moerdijk--Pronk's rules as decidable predicates (subgroup form), the $\ell$-relevant
  structure, $\mathsf M_7$ and $\mathsf R_{\Z}$; Remark~\ref{rmk:mprulesread} (\PROVED, \COMP\textsuperscript{r}); Definition~\ref{def:mpset}, the choice-free MP set, and
  Remark~\ref{rmk:mpwhat} (\PROVED, \COMP\textsuperscript{r}: $\pi_0$ of the restricted action groupoid; Moerdijk--Pronk's own set exactly in the effective, adapted, simplicial
  case; a non-effective generalisation for $\SL_2(\Z)$); Definition~\ref{def:mpfree} and Lemmas~\ref{lem:mprows}, \ref{lem:mpcols}, \ref{lem:mpassembly} (\PROVED, the last with
  a \STD{} alternative); \textbf{Theorem~\ref{thm:mpcomparison}} (\PROVED, \STD{} Weibel) --- the MP complex computes $H^{q}(\Ga;M)$ for every $q$, ring and module, $-I$ included,
  via the free resolution $\Z[\mathfrak F']$; Remark~\ref{rmk:mphybrid} (the hybrid is Proposition~\ref{prop:isomodelell}, its periodic-column instance; gap G1); Definition~\ref{def:mpphi}
  and Proposition~\ref{prop:mpalldegrees} (\PROVED, \STD, \COMP\textsuperscript{r}: $D=3$, the $24\cdot H^{\ge2}=0$ completeness of $k\le3$, what is certified with its $n$-range, scope
  bulk and $\Ga=\SL_2(\Z)$); Definition~\ref{def:mpmorphism} and Theorem~\ref{thm:mpnatural} (\PROVED, \STD, \COMP\textsuperscript{r}: naturality, refinement invariance, the label factor
  $\mathsf h_{\sigma_0}$, which [MP] v0.2 had dropped at the quotient level and the code carried --- F-MP2 \S1.3); Proposition~\ref{prop:mpbredon} (\PROVED, \STD,
  \COMP\textsuperscript{r}: quasi-isomorphism if $6\in\cring^{\times}$, the $\Z[1/2]$ odd-$n$ witness against ``only if'', injective on $H^{\le1}$); Propositions~\ref{prop:mpboundary},
  \ref{prop:mprescor} (\PROVED, \STD{} Brown; code planned, slot-23 MP-3); Remark~\ref{rmk:mphecke} (commentary, MP-4; G3 open); Example~\ref{ex:mpteardrop}
  (\COMP\textsuperscript{r}); Remark~\ref{rmk:mpcheckable} (\ref{MPC1}--\ref{MPC9}, mapping each clause to its registered validation); scope, code map, formalization note.
  The new \STD{} locators Weibel \S8.3, Weibel Lemma~2.7.3 (an alternative to the written proof of Lemma~\ref{lem:mpassembly}) and Brown Ch.~III.8 are
  \VERIFYloc{}; Weibel Thm.~2.2.6 and Brown Ch.~III.6.2, III.9--10 are the document's existing pointers (pre.6: Brown Ch.~III.9 as used in the repaired proof of
  Proposition~\ref{prop:mpboundary}(2), restriction as the map induced by $\Z[\Ga/H]\to\Z$, is added to the \VERIFYloc{} list).
\item \textbf{The forward-dated labels resolved.} \code{thm:mpcomparison} (the \code{via=} of the \code{MP\_ALL\_DEGREES} lines), \code{def:mpmorphism}, \code{rmk:mphybrid},
  \code{prop:mpalldegrees}, \code{sec:mpmodel} and \code{def:mpset}, cited ahead of this version by the catalogue, the MP reports and a C++ output string, now exist.
\item \textbf{Addenda} (no tag moves): Remark~\ref{rmk:mpcaveat} (items (2), (3) resolved by \S\ref{sec:mpmodel}, not withdrawn), Remark~\ref{rmk:orbifoldrefine} (owner ruling Q2),
  Remark~\ref{rmk:isomodellmp}(3) ($\mathsf R_{\zeta}$ against $\mathsf R_{\Z}$); a row of \S\ref{sec:codemap}.
\end{itemize}
\paragraph{Other.} Appendix~\ref{app:status}: twenty-three lines added (its counts paragraph). Appendix~\ref{app:labelindex}: \cl{sec:mpmodel} and the thirty-nine further new
labels. The referee reading list gains one item. Non-normative: the header comment, the date line.

\subsection*{v1.22.0-pre.4 --- 2026-10-09 --- hecke-extensions slot 22: Remark~\ref{rmk:negI}, the superfluous hypothesis $6\in R^{\times}$ (MP-0, erratum)}
A pre-release, not owner-released; the erratum is owner-approved (question Q6 of the Moerdijk--Pronk mathematical plan). In
Remark~\ref{rmk:negI}, second bullet ($n$ odd, $\ell\neq2$), the words ``under the additional hypothesis $6\in R^{\times}$'' are deleted: the
group-side argument ($-I$ is central, so it acts trivially on $H^{\bullet}(\Ga;-)$ and by $(-1)^n=-1$ on the coefficients, whence $2H^{\bullet}=0$ and, $2$ being
invertible, $H^{\bullet}=0$) uses $2\in R^{\times}$ only, which is the bullet's own hypothesis. Wording only: no statement, label, tag set, marker
(140 on 65) or value moves. Predictions for the Moerdijk--Pronk validations: \texttt{docs/implementation/slot22/MP\_predictions.md}.

\subsection*{v1.22.0-pre.3 --- 2026-10-09 --- hecke-extensions slot 22: the level-change cone and inflation--restriction (LES-C)}

\noindent\emph{This version.} \textbf{A pre-release, not owner-released}; refereed by F-LES (Fable, \file{docs/implementation/slot22/LES-C\_referee.md}: ACCEPT WITH
REPAIRS), repairs R1--R6 applied (below). Source:
\file{docs/implementation/slot22/LES0\_spec.md} v0.1 \S\S1--2, written out with proofs; its owner questions Q1--Q7 taken at the spec's recommendations by the lead pending
owner objection (the corrected normality picture; $\ell\in\{2,3\}$ labelled Bredon, the group-cohomology sequence deferred; $\F_\ell$ and $\Z/\ell^{K}$ only; the old-space
$U_p$ cone deferred; no console verb; the text in [C] as \S\ref{sec:nebles}). \textbf{No existing statement number, label, tag set or status moves; the marker count
stays $140$ on $65$ labels} (the new labels carry none).
\begin{itemize}\setlength{\itemsep}{0pt}
\item \textbf{New: \S\ref{sec:nebles}}, the last subsection of \S\ref{sec:nebentypus}: Definitions~\ref{def:levelpair}, \ref{def:levelres}, \ref{def:levelcone},
  \ref{def:normalpair}; Lemma~\ref{lem:nebcompat} (\PROVED); Lemma~\ref{lem:levelresmorph} (\PROVED, \STD{} Brown) --- restriction and trace as fibre maps;
  Lemma~\ref{lem:levelconeexact} (\PROVED) --- the six-term sequence at every ring and every $\ell$, unconditionally; Propositions~\ref{prop:levelconeshapiro},
  \ref{prop:modeltorsion}, \ref{prop:levelsplit} (\PROVED); Proposition~\ref{prop:levelconehecke} (\PROVED, \STD{} Brown, \COMP\textsuperscript{a}; level-raising $U_p$ \PROVED{}
  negative after R1);
  Propositions~\ref{prop:levelinfres} (\PROVED, \STD{} Brown, Hochschild--Serre) and \ref{prop:cycliccone} (\PROVED, \STD{} Brown); Lemmas~\ref{lem:levelfreeeuler},
  \ref{lem:unipotentinv} (\PROVED); Propositions~\ref{prop:gamma0normal}, \ref{prop:hidanormalsteps} (\PROVED, \COMP\textsuperscript{a}); Remark~\ref{rmk:normalN1}
  (\PROVED); Obligation~\ref{obl:levelconeexact}; Remark~\ref{rmk:lescheckable} (\ref{LES1}--\ref{LES9}); Remark~\ref{rmk:oldspacecone}; the code map (planned,
  slot-22 LES-1..3). All new Brown locators are \VERIFYloc{} or marked as read at an earlier R20-L pass.
\item \textbf{Departures from the spec}, each argued in the text: the same-level $U_p$ square is \PROVED{} strict (the spec had $U_p$ for $p\mid N_A$ in general); the
  level-raising $U_p$ square does \emph{not} commute in this document's convention (the spec asserted it on cohomology from the classical right cosets, which belong to
  the classical $U_p=[A\,\mathrm{diag}(p,1)A]$, not to this document's $U^{(N)}_p=U_p^{\ast}$); the Adamson / Bieri--Eckmann identification is a parenthesis, not a claim; $\chi_4$ modulo $5$ is odd (its model vanishes for even $n$) and has no values in $\F_7$; every cell of the
  spec's arm~(i) is split ($\ell\nmid p+1$).
\item \textbf{F-LES repairs} (2026-10-09, same pre-release). \textbf{R1}: Proposition~\ref{prop:levelconehecke}(3) \OPENSTAT{} $\to$ \PROVED{} \COMP\textsuperscript{a}
  with negative content --- (a) $U^{B}_p\circ(\alpha_{B/A})_{\ast}=p\cdot[A\,g\,B]$ as cochain maps (via (T1) of Definition~\ref{def:tau}, as in
  Proposition~\ref{prop:oldspace}(1); $=0$ at $\ell=p$), (b) $U^{A}_p=\cor\circ[A\,g\,A'']$ at cochain level, (c) for $B=A''$, $6,p\in\cring^{\times}$ the square commutes on
  $H^{q}$ iff $\mathrm{dg}_g(H^{q}(A))\subseteq\res H^{q}(A)$, and fails at $(\Ga_0(25)\subset\Ga_0(5),\ \ell=7,\ n=0)$ by an explicit Eisenstein computation, (d) this
  document's $U^{(N)}_p$ is $U_p^{\ast}$ on old spaces, the classical $U_p=[A\,\mathrm{diag}(p,1)A]$ with $\tau'=S\tau S^{-1}$ is strict by (1) (deferred); the
  ``not settled here'' parenthesis and the double-coset count with its Brown~III.9.5(ii) locator are gone (so R2 is void). Propagated to
  Obligation~\ref{obl:levelconeexact}(vi) (the strict identity gated, the square recorded), \ref{LES9} (\PROVED), Remark~\ref{rmk:oldspacecone}, Appendix~\ref{app:status}
  (two \OPENSTAT{} tags go, $8\to6$; \COMP{} $45\to46$), the reading list; evidence \file{les0/levelraise\_up\_check.py}, \file{les0/out\_levelraise\_up.txt}. \textbf{R3}:
  Proposition~\ref{prop:levelconehecke}: $N_A\mid N'_A$, $N_B\mid N'_B$ stated, the coset invariant is $a_\gamma c_\gamma/N'\bmod p$. \textbf{R4}:
  Proposition~\ref{prop:hidanormalsteps}: the non-normality witness inside $\Ga_0(N)$ is $\smat10{NN'}1$. \textbf{R5}: Lemma~\ref{lem:levelresmorph}(3)'s \STD{} marked as
  read at R20-L; \cite{adamson1954}, \cite{bieri_eckmann1978} added to the bibliography (\VERIFYloc). \textbf{R6}: the $T_p$ instance of
  Proposition~\ref{prop:levelconehecke}(2) extended to $\Ga(M)$, $\Ga^{2}$, $\Ga^{3}$; $\Ga(3)$ added to Proposition~\ref{prop:modeltorsion}'s torsion-free list
  (Remark~\ref{rmk:normalN1}: $M\ge3$). No label, number or marker moves.
\end{itemize}
\paragraph{Other.} Appendix~\ref{app:status}: twenty-one lines added (its counts paragraph). Appendix~\ref{app:labelindex}: \cl{sec:nebles} and the twenty-four new
labels. The referee reading list of Appendix~\ref{app:changelog}'s provenance section gains one item. Two bibliography entries (R5). Non-normative: the header comment,
the date line.

\subsection*{v1.22.0-pre.2 --- 2026-10-08 --- hecke-extensions slot 22: Remark~\ref{rmk:levelNcurvefamily}(1), sufficient, not necessary (DOC22)}

\noindent\emph{This version.} \textbf{A pre-release, not owner-released.} One rewording, from \file{docs/implementation/slot22/LN23F\_report.md}~\S1, \S6 item~1 (a finding of LN23-F against the F4 referee's ``iff''): in
Remark~\ref{rmk:levelNcurvefamily}(1), ``By Mazur this holds iff $\mathcal F$ contains a good prime'' becomes ``if'', with the $p$-odd, $\ell=2$ sentence and the $p=2$ counter-example $(2473,2)$ (\COMP), and the next sentence reads
``may read'' for ``reads''. The remark's tags (\BSTD, \CONJ, \COMP\textsuperscript{a}) are unchanged. Not applied from the same report: its \S6 items~2 and~3 (the family-certificate figures; item~3 not added as a data line: $e_3(1747)=3$ appears only through ``Loewy $4=e_3+1$'' in item~1's sentence). No statement number, label, tag set or marker moves.

\subsection*{v1.22.0-pre.1 --- 2026-10-06 --- hecke-extensions slot 22: wild characters, \texorpdfstring{$\Ga_{\mathsf H}(N)$}{Gamma\_H(N)} and the tower at \texorpdfstring{$\ell$}{l} (P10-C)}

\noindent\emph{This version.} \textbf{A pre-release, not owner-released}; decisions taken as recommended. The text of
\file{docs/implementation/slot22/P10C\_contract\_text.tex.txt} v0.2, part~(A), inserted verbatim (``v1.22.0-pre.N'' numbered~$1$), as refereed by the slot-22 Fable referee F-U5
(\file{docs/implementation/slot22/P10C\_F4\_referee\_report.md}, ACCEPT WITH REPAIRS, $0$ blocking / $3$ major / $13$ minor, all applied in that text; the annotations
``v0.6.0-pre.2, F-U5 \dots'' inside it name the referee's revision v0.2 of the text, released here and in [U] 0.6.0-pre.1). Implementation spec \file{docs/implementation/slot22/P10\_spec.md} v0.2.
\textbf{No existing statement number, label, tag set or status moves; the marker count stays $140$ on $65$ labels} (the new labels carry none; no marked label is touched).
\begin{itemize}\setlength{\itemsep}{0pt}
\item \textbf{New: \S\ref{sec:nebwildsec}}, the last subsection of \S\ref{sec:nebentypus}, so that no statement number of v1.21.0 moves: Definition~\ref{def:nebwild} and
  Lemma~\ref{lem:nebwild} (\PROVED, \STD{} Washington, \VERIFYloc) --- the chain rings $\mathcal O_{j,e}$ and the wild characters $\chi_{i,t}$; Definition~\ref{def:nebgammaH},
  Propositions~\ref{prop:nebgammaH}, \ref{prop:nebgammaHheilbronn} (\PROVED; F-U5 m6, m7) --- $\Ga_{\mathsf H}(N)$, its table, diamonds, operators and Heilbronn action;
  Lemma~\ref{lem:nebshapiro} (\PROVED; F-U5 M1: $\Ga'$ acts on $P$ through the inverse of the quotient map) --- specialisation and coinvariants without division;
  Lemmas~\ref{lem:nebwildtors} (\PROVED, \STD{} Ash--Stevens) and \ref{lem:nebchainrref} (\PROVED); Remark~\ref{rmk:nebwildcheckable} (\ref{NB7}--\ref{NB10}, \PROVED; \ref{NB8}'s rank
  clause \PROVED{} \emph{given} [U,~\texttt{prop:u-towerorient}], \COMP{} at $r=1$, gated; F-U5 M3); Remark~\ref{rmk:nebcosetlevel} (the code contract); the code map (planned, slot-22 P10).
\item \cl{rmk:nebfield}(2): one sentence (wild characters admitted over $\mathcal O_{j,e}$; over $\F_q$ and $W_K(\F_q)$ the refusal stands).
\item The refusal box after \cl{rmk:nebcost}: item~(ii) narrowed, not reversed (\code{character\_wild} over $\F_q$, $W_K(\F_q)$; new names \code{wild\_ring\_order}, \code{wild\_precision},
  \code{wild\_tame\_level}).
\item After \cl{thm:hidacontrol}: an unnumbered addendum to its~(2) (\file{docs/implementation/slot22/P10C\_addenda.md}~(1), drafted by F-U5 and verified on Hida pp.~589, 593) ---
  ``no twist intervenes'' is relative to Hida's own $\Lambda$-structure; the index $a$ of $\mathcal V^{0}(a)$ is the natural one. An addendum, not an erratum; the theorem's statement is unchanged.
\end{itemize}
\paragraph{Other.} Appendix~\ref{app:status}: ten lines added (its counts paragraph). Appendix~\ref{app:labelindex}: \cl{sec:nebwildsec} and the ten new labels.
Macros \verb|\Zl|, \verb|\Ql|, \verb|\VERIFYloc| provided (as in [U]). The one departure from the verbatim text: in \ref{NB8} the two cross-document references to [U]'s \texttt{thm:u-towercontrol} and \texttt{prop:u-towerorient} are
printed as label names (\verb|\texttt|), not \verb|\ref|, as [C] cites [U] elsewhere (a \verb|\ref| into another document is undefined here). Non-normative: the header comment, the date line. Companion: [U] 0.6.0-pre.1 (part~5, [U,~\texttt{sec:u-tower}]).

\subsection*{v1.21.0 --- 2026-10-06 --- hecke-extensions slot 21, the close: the slot-21 release, from v1.21.0-pre.4}

\noindent\emph{This version.} \textbf{RELEASE from v1.21.0-pre.4 under the owner's decision D16 of 2026-10-06} (release of the pre-release documents at the
slot-21 close, accepted as recommended; \file{docs/implementation/slot21/slot21\_status\_and\_close\_plan.md}~\S5). It releases the content of
v1.21.0-pre.1 \dots{} pre.4 as described in their entries below, each with the referee status stated there.
\textbf{No statement's content changes; no label is added, renamed or removed; no status tag moves; the marker count stays $140$ on $65$ labels.}
\begin{itemize}\setlength{\itemsep}{0pt}
\item Header and \verb|\date|; this entry.
\end{itemize}
Companion pointers: [E] 1.21.0, [U] 0.5.0, [L37] 0.5.1, [RC] 0.1.0, [LG] 0.1.5, [EX] 1.4.1, [CLI] 1.3.0.

\subsection*{v1.21.0-pre.4 --- 2026-10-06 --- hecke-extensions slot 21: the weight-two curve object (LN23)}

\noindent\emph{This version.} \textbf{A pre-release, not owner-released}; the owner's decision O6 on slot-22 LN23, taken as the F4 referee recommended
(\file{docs/implementation/slot22/LN23\_F4\_referee\_report.md}~\S11, ACCEPT WITH REPAIRS, 0~B / 2~M / 9~m, the repairs in the delivered text).
\textbf{No existing statement number, label, tag set or status moves; the marker count stays $140$ on $65$ labels} (the new labels carry none; no marked label is touched).
\begin{itemize}\setlength{\itemsep}{0pt}
\item \textbf{New: Proposition~\ref{prop:levelNcurve}}, in a new last subsection \S\ref{sec:levelNcurve} of \S\ref{sec:levelN}, so that no statement number of
  v1.21.0-pre.3 moves (\PROVED; text as delivered in \file{docs/implementation/slot22/LN23\_C\_proposition.tex.txt}): at $n=0$ the cokernel presentation with
  orbit-sum fixed spaces is the cellular $H^{1}$ of the quotient graph $\Ga_0(N)\backslash\mathcal T$, i.e.\ $H^{1}(Y_0(N);\Z)=\Hom(\Ga_0(N),\Z)$, free of
  rank $2g+c-1$, and its reduction $H^{1}(Y_0(N);\F_p)$ for \emph{every} prime $p$ --- at $p\in\{2,3\}$ not group cohomology but the Bredon object of \ref{H3};
  $H_{N,q}$ induces $T_q$ over $\Z$ (Proposition~\ref{prop:heilbronn} over $\Z[1/6]$ and torsion-freeness). One \textsf{[standard]} step (the spine, \cite[Ch.~I, \S4.2]{serre_trees},
  locator not verified) with Proposition~\ref{prop:bredon}(2) as the alternative; \COMP\textsuperscript{a} for the torsion of the two-/three-term lattice and the
  unimodular elimination (\file{docs/implementation/slot22/ln23ref/}).
\item \textbf{New: Remark~\ref{rmk:levelNcurvefamily}} (\STD{} Mazur II \S\S9, 16--17 read on the scan, Calegari--Emerton held, Sturm held; \CONJ;
  \COMP\textsuperscript{a}): Mazur's Eisenstein ideal at $p\in\{2,3\}$ on the curve object, and the two ways a finite Hecke family fails --- generation (no
  Mazur good prime) and separation (another maximal ideal).
\item \cl{def:levelNmodel}(c): admits $p\in\{2,3\}$ at $n=0$ on the curve object (a definition; no tag).
\item \ref{LN6} (in \cl{rmk:levelNcheckable}): an addendum carrying the F2 caveat to the extension of \ref{LN4}, \ref{LN6} to $p\in\{2,3\}$; no tag of the remark moves.
\end{itemize}
\paragraph{Other.} \S\ref{sec:levelN}'s scope paragraph names the exception. Appendix~\ref{app:status}: two lines added (its counts paragraph).
Appendix~\ref{app:labelindex}: \cl{sec:levelNcurve}, \cl{prop:levelNcurve}, \cl{rmk:levelNcurvefamily}. \file{gl2\_refs.bib}: \code{calegari\_emerton}.
Non-normative: the header comment, the date line. Companions unchanged.

\subsection*{v1.21.0-pre.3 --- 2026-10-06 --- hecke-extensions slot 21: the (U17)(a) addendum}

\noindent\emph{This version.} \textbf{A pre-release, not owner-released.} One addendum sentence.
\textbf{No clause's mathematics changes; no label, tag or status moves; the marker count stays $140$ on $65$ labels} (no marked label is touched).
\begin{itemize}\setlength{\itemsep}{0pt}
\item \ref{Hi5} (in \cl{rmk:hidacheckable}): for $n<\ell-1$, [U]'s (U17)(a) (the ordinary $\mathfrak m_{\mathrm{Eis}}$-block of $H^{+}$ free of rank one over its Hecke algebra) is
  equivalent to \Espec{cl:lveis}(2) on that block, as read in [E] after the proof of \Espec{cl:lveis} (\file{docs/implementation/slot21/P8\_F4\_referee\_report.md}, M3 and \S4 item~3;
  [RC]~\S9.4); for $n\ge\ell-1$ both stay conjectural. This is the row v1.21.0-pre.2 listed as ``not done here''.
\end{itemize}
Companions: [E] 1.21.0-pre.5, [U] 0.5.0-pre.6, [L37] 0.5.1-pre.6, [RC] 0.1.0-pre.8.

\subsection*{v1.21.0-pre.2 --- 2026-10-05 --- hecke-extensions slot 21: addenda from [U] part~4 and the literature pass}

\noindent\emph{This version.} \textbf{A pre-release, not owner-released.} Addenda and locator repairs only.
\textbf{No clause's mathematics changes; no label, tag or status moves; the marker count stays $140$ on $65$ labels} (no marked label is touched).
\begin{itemize}\setlength{\itemsep}{0pt}
\item \cl{rmk:stabconvention}(3): an \emph{addendum}, not an erratum --- since [U] v0.5.0-pre.2 (\texttt{lem:u-latticeord}, \PROVED{} there; the slot-21 referee F4 of [U]
  part~4, \file{docs/implementation/slot21/P8\_F4\_referee\_report.md}, \S4 item~1) the ordinary part does not see the lattice: $D_\ast$ is an isomorphism on $e_1H^{q}$ for every $q$,
  on $\partial$ and at level~$\ell$; \ref{Hi8}'s arms (a) and (b) coincide and its report (c) ($45/45$) is a theorem. \ref{Hi8}'s own text and tag are unchanged.
\item \cl{thm:hidacontrol}(2): ``its normalisation is not restated here'' now continues with the normalisation read on the page (Hida, Invent.\ p.~545, F4 M2): Hida's
  $P_n=(1+X)-(1+\ell)^{n+2}$, the $P_k$ of item~(3) with $k=n+2$ (F4 \S4 item~2).
\item Literature pass of 2026-10-05 (\file{docs/implementation/slot21/missing\_literature.md}~\S7; LIT-15, LIT-20, LIT-23 received): \cl{rmk:thetaclassical}(1) (Katz,
  Theorem~(2), p.~55, for $N\ge3$; the ``iff'' from Jochnowitz's local-components paper, Lemma~4.1(b), p.~257), (4) (that paper's \S7, pp.~261--264); \cl{rmk:nebcheckable}
  \ref{NB1} (Cohen--Oesterl\'e, Th\'eor\`eme~1, pp.~71--72; Shimura Thm.~8.4 cuspidal only, R20-L). The \textsf{[standard]} tags stay; the bib notes of \code{katz1977},
  \code{jochnowitz\_local}, \code{cohen\_oesterle} record the reading.
\item Not done here (owner's call, listed): a [C] row for [U]'s (U17)(a) $=$ \Espec{cl:lveis}(2) on the ordinary block is in [E] 1.21.0-pre.1, not in [C].
\end{itemize}
Companions: [E] 1.21.0-pre.1, [U] 0.5.0-pre.3, [L37] 0.5.1-pre.2, [EX] 1.4.1-pre.1.

\subsection*{v1.21.0-pre.1 --- 2026-10-05 --- hecke-extensions slot 21, R20-L: the Shimura and Brown locators verified}

\noindent\emph{This version.} \textbf{A pre-release, not owner-released.} Step R20-L (slot-21 plan \S0 item 7; condition (C1) of the owner's Q5
ruling, \file{docs/implementation/slot21/owner\_rulings\_Q5\_1005.md}) read every locator of the nine Shimura \code{\textbackslash LEANlit}
markers on the owner's scan of \cite{shimura} Ch.~3 and Ch.~8 (page images), \cite{brown} III.8--10 and \cite{ash_stevens} \S1.1 (GDZ page images);
report \file{docs/implementation/slot21/R20-L\_report.md}. Every marked clause is \textbf{verified}; the locators, not the clauses, are restated.
\textbf{No clause's mathematics changes; no label, tag or status moves; the marker count stays $140$ on $65$ labels.}
\begin{itemize}\setlength{\itemsep}{0pt}
\item The nine markers: ``locator unverified (R20-L)'' replaced by the verified locators --- \texttt{hcomm}: Prop.~3.8 (p.~54) as applied
  in \S3.2 (p.~56), for $\det<0$ the same argument with a diagonal representative (Thm.~3.24 dropped: it does not state commutativity);
  \texttt{hTp}: Thm.~3.24(5) with Prop.~3.17 (Prop.~3.30 dropped: it is the level-$N$-to-level-$1$ homomorphism); \texttt{hcop}: Prop.~3.16
  (Thm.~3.24(3) is the full-operator form); \texttt{hTpchi0}, \texttt{hTpchi}: Thm.~3.24(5) and Prop.~3.17 carried to
  $R(\Ga_0(N),\De_0(N))$ by Prop.~3.31 (Shimura's ring carries no character; $\chi$ enters through $T(p,p)$); \texttt{hUchi}: Thm.~3.34(3) with
  Prop.~3.33. The transport is \cite[\S1.1, p.~194]{ash_stevens}, where the action of the Hecke algebra on cohomology is stated as well known
  (with a reference to Kuga--Parry--Sah), not proved.
\item \textsf{[standard]} notes: \cl{prop:cosets} (\S3.2, Lemmas~3.11--3.13, not \S3.1; also its status row); \eqref{eq:fullop} (Thm.~3.24(1), (2),
  (7): the notation confirmed); \cl{thm:relations} (its note; the proof's ``$T(p,p)^{\pm1}$'' corrected to Thm.~3.20's polynomial ring on
  $T(1,p)$, $T(p,p)$; the coprime clause Prop.~3.16 / Thm.~3.24(3), not ``\S3.3''; also its status row); \cl{obl:algebra}'s note;
  \cl{def:nebdata}(2) ((3.3.3), p.~68); the $\Ga_1(N)$-double-coset statement in the proof of \cl{prop:nebcorr} (printed Lemma~3.29 and
  Prop.~3.32--3.33, not ``Prop.~3.29''); \cl{thm:nebrelations}'s note; \ref{err:e1}'s evidence (Thm.~3.24(1), (2), not Thm.~3.34);
  \cl{lem:heckerankone}'s secondary pointers (Thm.~3.41 / 3.42 and Prop.~8.5 verified; Diamond--Shurman still not); the two Thm.~8.4 sites
  (the parabolic part); one \cite{brown} note (Cor.~III.10.2 with Prop.~III.9.5(ii)).
\item \S\ref{sec:status}'s paragraph on the slot-20 \code{\textbackslash LEANlit} markers; header and \verb|\date|.
\end{itemize}
Left to the [GT] owner (not edited here, \file{lean/} untouched): the \verb|\VERIFYloc| texts of \texttt{lit:hecke-Tp-prim}, \texttt{lit:hecke-coprime},
\texttt{lit:hecke-commute}, the nebentypus paragraph, \texttt{conv:hecke}, and \texttt{lit:hecke-classical}, whose Ash--Stevens half is
\textbf{not as cited} (Ash--Stevens \S1.1 gives the cocycle formula of $T_g$, not $\cor\circ c_g\circ\res$; no [C] marker rests on it).
Companion pointers unchanged.

\subsection*{v1.20.0 --- 2026-10-05 --- hecke-extensions slot 20, the close: the release, from v1.20.0-pre.2, with the milestone-M4 \texorpdfstring{\code{\textbackslash LEAN}}{LEAN} markers}

\noindent\emph{This version.} \textbf{RELEASE from v1.20.0-pre.2 under the owner's ruling O6 of 2026-10-04 (releases after the Fable windows F1 and F1$'$); the milestone-M4
markers clause by clause per L20-6 \S8 as amended by L20-6\_repairs} (\file{docs/implementation/slot20/L20-6\_F3\_audit.md}, \file{docs/implementation/slot20/L20-6\_repairs.md}).
Referees: F1 (\file{docs/implementation/slot20/WP-3\_F1\_review.md}) and F1$'$ (\file{docs/implementation/slot20/F1prime\_review.md}) ACCEPT for [C] (code-map text, no finding);
the Lean audit F3 ACCEPT WITH REPAIRS on M4 (one MAJOR, \cl{thm:classical} formalised in a weak form only; repaired), Lean census
\code{decls=1303 proved=1238 proved\_mod\_lit=37 stated=0 planned=0 sorries=0 global\_axioms=0 \dots{} verdict=OK}. \textbf{No statement's content changes; no label is added, renamed or removed; no status tag moves.}
\begin{itemize}\setlength{\itemsep}{0pt}
\item \textbf{The milestone-M4 markers} ($42$ new: $33$ \code{\textbackslash LEAN}, $9$ \code{\textbackslash LEANlit}; $140$ markers on $65$ labels in all, $126$ \code{\textbackslash LEAN} and $14$ \code{\textbackslash LEANlit}), clause by clause per L20-6 \S8 as amended by L20-6\_repairs: \cl{lem:tauexists}(2) last sentence --- \texttt{lem:m4-tauR-exists}; \cl{def:hecke}, \cl{prop:cosets}(5) --- \texttt{lem:m4-cosetdata}; \cl{thm:welldef}(a) --- \texttt{prop:fullassembly-model} ($\tau$ over $\cring$); \cl{prop:diamond} --- \texttt{prop:diamond-model}; \cl{prop:fullassembly} --- the summand only; \cl{def:operators}(1), (2), (3) --- \texttt{def:m4-hecke}, (3) only in the assembled form; \cl{thm:relations}(1), its caveat, (2) and \cl{obl:algebra}'s commutativity, (ii), (iii) --- \code{\textbackslash LEANlit} modulo \texttt{hTp} / \texttt{hcomm} / \texttt{hcop} (\cite{shimura} Prop.~3.8, Thm.~3.24, Prop.~3.30, locators unverified, R20-L); \cl{prop:lhcompliant}(4), \cl{def:inducedmodule}, \cl{prop:blockmatrices}, \cl{def:levelellmodel} --- \texttt{def:levelellmodel-model}; \cl{prop:levelcorrespondence}(2), (3), (4)(b)(c) --- \texttt{prop:levelcorrespondence-model}; \cl{def:nebdata}(2)(3)(5), \cl{lem:nebdata}(1)(2)(4) --- \texttt{def:nebdata-model}; \cl{prop:nebcorr} --- \texttt{prop:nebcorr-model}; \cl{thm:nebrelations}(2), (3) --- \code{\textbackslash LEANlit} modulo \texttt{hTpchi0}, \texttt{hTpchi}, \texttt{hUchi} (Thm.~3.34, 3.35, \S3.3, unverified); \cl{lem:merelpath} --- \texttt{lem:merelpath-model}; \cl{prop:heilbronn} --- \texttt{prop:heilbronn-model}, \texttt{chk:merel-n5}; \cl{lem:cosetrepsN}(4) --- \texttt{prop:heilbronn-model}. Not tagged: \cl{thm:classical}, \eqref{eq:corconjres} (not even its maps, although \texttt{def:m4-classical} formalises them), \eqref{eq:fullop}, the ring homomorphism of \cl{obl:algebra}, choice-freeness of the level-$N$ / nebentypus operators. \S\ref{sec:status} names the v1.20.0 list.
\item Header and \verb|\date|.
\end{itemize}
Companion pointers: [E] 1.20.0, [U] 0.4.0, [L37] 0.5.0, [EX] 1.4.0 (unchanged).

\subsection*{v1.20.0-pre.2 --- 2026-10-04 --- hecke-extensions slot 20, the close delta: code-map text}

\noindent\emph{This version.} A pre-release, \textbf{not owner-released}, for the slot-20 Fable referee F1$'$, made with [E] 1.20.0-pre.6, [U] 0.4.0-pre.3 and [L37] 0.5.0-pre.5 (list:
\file{docs/implementation/slot20/close\_delta.md}). \textbf{No statement's content changes; no label is added, renamed or removed; no status tag moves}.
\begin{itemize}\setlength{\itemsep}{0pt}
\item \S\ref{sec:notimplemented}: the v1.20.0-pre.1 box gains a ``built'' paragraph --- WP-MEAS, S1\,/\,S2, NB-4, WP-LV3 merged, with their commits.
\item \S\ref{sec:codemap}, the \S\ref{sec:involution} row: $\inv^{\ast}$ over $\Z/\ell^{k}$ built (WP-LV3(d)).
\item Header and \verb|\date|.
\end{itemize}
Companion pointers: [E] 1.20.0-pre.6, [U] 0.4.0-pre.3, [L37] 0.5.0-pre.5, [EX] 1.4.0 (unchanged).

\subsection*{v1.20.0-pre.1 --- 2026-10-04 --- hecke-extensions slot 20, WP-2: code-map text for slot 20}

\noindent\emph{This version.} A pre-release, \textbf{not owner-released}, for the slot-20 Fable referee F1 (WP-3), made with [E] 1.20.0-pre.1, [U] 0.4.0-pre.1 and [L37] 0.5.0-pre.1
(list of new and changed labels: \file{docs/implementation/slot20/WP-2\_drafts.md}). \textbf{No statement's content changes; no label is added, renamed or removed; no status tag moves};
Appendices~\ref{app:status} and \ref{app:labelindex} are unchanged.
\begin{itemize}\setlength{\itemsep}{0pt}
\item \S\ref{sec:notimplemented}: a v1.20.0-pre.1 box --- WP-LV3, WP-MEAS, NB-4 planned, WP-JET value-neutral; their mathematics in [E] 1.20.0-pre.1 (\Espec{sec:lvmeasure}, \Espec{sec:lvdepth2}).
\item \S\ref{sec:codemap}, the \S\ref{sec:involution} row: $\inv^{\ast}$ over $\Z/\ell^{k}$ planned (WP-LV3(d)); Propositions~\ref{prop:iotasplit}, \ref{prop:iotacommutes} cover it.
\item Not in this pre-release: the milestone-M4 \code{\textbackslash LEAN} markers (slot-20 L20-6 and its repairs); the Shimura locators stay marked unverified until R20-L reads the scan.
\item Header and \verb|\date|.
\end{itemize}
Companion pointers: [E] 1.20.0-pre.1, [U] 0.4.0-pre.1, [L37] 0.5.0-pre.1, [EX] 1.4.0 (unchanged).

\subsection*{v1.19.0 --- 2026-10-04 --- hecke-extensions slot 19, the close: the release, from v1.19.0-pre.6, with the milestone-M3 \texorpdfstring{\code{\textbackslash LEAN}}{LEAN} markers}

\noindent\emph{This version.} \textbf{RELEASE from v1.19.0-pre.6 under the owner's ruling O3, on F1 / F1-delta ACCEPT WITH REPAIRS; the milestone-M3
markers clause by clause per L19-6 \S8 as amended by L19-6\_repairs \S3 and L19-4 \S3} (\file{docs/implementation/slot19/L19-6\_F3\_audit.md},
\file{docs/implementation/slot19/L19-6\_repairs.md}, \file{docs/implementation/slot19/L19-4\_report.md}). Referees: F1 (\file{docs/implementation/slot19/WP-3\_F1\_referee\_report.md})
and F1-delta (\file{docs/implementation/slot19/F1delta\_referee\_report.md}) ACCEPT WITH REPAIRS, repairs at pre.2 and pre.6; the Lean audit F3 ACCEPT WITH REPAIRS
on M3, its eight MINORs repaired in Lean (Lean census \code{decls=1136 proved=1077 proved\_mod\_lit=19 stated=0 planned=16 sorries=0 global\_axioms=0
verdict=OK}). \textbf{No statement's mathematics changes; no label is added, renamed or removed; no status tag moves}; Appendices~\ref{app:status} and
\ref{app:labelindex} are unchanged (markers are not counted). The Shimura locators of v1.19.0-pre.5 remain marked unverified (from memory; a scan is requested).
\begin{itemize}\setlength{\itemsep}{0pt}
\item \textbf{The milestone-M3 markers} ($17$ new, all \code{\textbackslash LEAN}; one \code{\textbackslash LEANlit} becomes \code{\textbackslash LEAN};
  $98$ markers on $48$ labels in all, $93$ \code{\textbackslash LEAN} and $5$ \code{\textbackslash LEANlit}):
  \cl{prop:model7h4}(ii) --- the marker \texttt{prop:model7h4-ii-dim} (``the dimension only'') replaced by \texttt{prop:model7h4-ii}: the acyclicity in the
  three forms of \eqref{eq:contractibleuse} and $\dim\Xbulk=2$; topological contractibility not formalised;
  \eqref{eq:contractibleuse} --- \texttt{prop:model7h4-ii} (over $\Z$, on the model), \texttt{thm:xbulk-repaug} and \texttt{lem:xbulk-contractz} (after $\otimes\cring$);
  \cl{rmk:orbifold} --- \texttt{lem:xbulk-orient} (the direct-sum decomposition), \texttt{thm:xbulk-repaug} (the resolution), \texttt{thm:model7-rmkorbifold}
  (the isomorphism under $6\in\cring^{\times}$, every acyclic descriptor with \ref{H4}(iv), $L$ in the $(G,\Ga,\Delta,e,\rho)$ form; the statement, not the
  argument; $K$ outside the descriptor framework not claimed);
  \cl{lem:tauexists}(2) --- \texttt{lem:tauexists-2-model} ((1) not);
  \cl{thm:welldef}(b) --- \texttt{thm:welldef-b} ($\tau$ over $\cring$, on $\Heq{n}$ itself); (c) --- \texttt{thm:welldef-c} now at the whole of (c); (a) ---
  any two choices of representatives; \cl{def:hecke} --- the caveat names $\tau$ over $\cring$ on the model;
  \cl{cor:stablecert}(1) --- \code{\textbackslash LEANlit} modulo Eichler--Shimura at level one (\texttt{hES}) in place of ``modulo
  Lemma~\ref{lem:heckerankone}''; (2)'s first sentence --- \texttt{cor:stablecert-uncond} (\code{\textbackslash LEAN});
  the cone chapter on the model (L19-4) --- \S\ref{sec:cone}'s standing hypotheses \texttt{thm:model7-pairhyp}, \cl{def:pair} \texttt{def:xbdry} and
  \texttt{thm:xbdry-pair}, \ref{B1}--\ref{B3} \texttt{def:model7bd} (the component map of \ref{B2} not formalised), \cl{lem:etaexists}, \cl{lem:etaunique} and
  \cl{thm:conewelldef}(b) \texttt{cor:model7-etaexists}, \cl{obl:eta}'s link to Definition~\ref{def:pairhomotopy} \texttt{thm:etaspec-pairhomotopy} and
  \texttt{cor:etap2-pairhomotopy} (that realisations of the tables exist is not formalised).
  Unchanged: \cl{lem:conerelative}(5) ``given (4)''; \cl{lem:heckerankone}, \cl{prop:cohforms}. Not tagged: \ref{H4}(ii) as contractibility of a space,
  \cl{lem:tauexists}(1), \cl{rmk:orbifold}'s general $K$ and its proof route.
\item \cl{prop:model7h4}: the v1.19.0-pre.5 parenthesis at the \ref{H4}(ii) marker (``this marker's scope moves only at the slot-19 close's tagging pass'')
  is replaced by the marker's new scope.
\item \S\ref{sec:status}: the v1.19.0 list named. Header and \verb|\date|.
\end{itemize}
Companion pointers: [E] 1.19.0, [U] 0.3.0, [L37] 0.4.0, [EX] 1.4.0 (unchanged).

\subsection*{v1.19.0-pre.6 --- 2026-10-04 --- F1-delta m-C1}

\noindent A pre-release, \textbf{not owner-released}: the slot-19 F1-delta's m-C1 --- one sentence in Appendix entry~\ref{err:e6} (Shimura's naming of
$T(1,q^{2})$), textual only. No statement, tag, label or count change.

\subsection*{v1.19.0-pre.5 --- 2026-10-04 --- hecke-extensions slot 19, WP-2 continued: Shimura's notation, H4(ii), WP-BIL and WP-LIN}

\noindent A pre-release, \textbf{not owner-released}, made with [E] 1.19.0-pre.6, [U] 0.3.0-pre.4 and [L37] 0.4.0-pre.3.
\begin{itemize}\setlength{\itemsep}{0pt}
\item \textbf{Shimura's notation} (\file{docs/implementation/slot19/shimura\_convention\_check.md}): \S\ref{sec:primfull}, Definition~\ref{def:operators}(3), Theorem~\ref{thm:relations}'s caveat
  and proof, the proof in \S\ref{sec:relationsmodel} and Appendix entry~\ref{err:e1}'s evidence called the full operator \eqref{eq:fullop} ``Shimura's
  $T(1,p^{\nu})$''. In Shimura's notation $T(1,p^{\nu})$ is the single double coset $\Deprim_{p^{\nu}}$ and $T(p^{\nu})$ the sum of all double cosets of determinant $p^{\nu}$;
  the text now says $T(p^{\nu})$ for the full operator (and $T(p)T(p^{\nu})=T(p^{\nu+1})+p\,T(p,p)T(p^{\nu-1})$, $T(M)T(N)=T(MN)$). The mathematics --- the definition, the
  recursion, $T_p^{2}=T_{p^{2}}+p^{\,n+1}$ --- is unchanged; the later uses of $T(1,p^{2})$ for the primitive coset (\S\ref{sec:levelN}) were already in Shimura's notation.
  The locators (\S3.1--3.3, Thm.~3.24, Thm.~3.34) are from memory: the book is not in the library and is to be checked when a scan is available.
\item Proposition~\ref{prop:model7h4}, the \ref{H4}(ii) clause of its marker: ``contractibility is stated in [GT], not proved'' was superseded by slot-19 L19-3
  (\texttt{model7\_acyclic}: the acyclicity used through \eqref{eq:contractibleuse}, proved combinatorially). A statement of fact; the marker itself moves only at the close.
\item \S\ref{sec:notimplemented}: a v1.19.0-pre.5 box --- built: WP-BIL, WP-LIN.
\item Header and \verb|\date|.
\end{itemize}
\textbf{No statement's mathematics changes; no label is added, renamed or removed; no tag moves.} Companion pointers: [E] 1.19.0-pre.6, [U] 0.3.0-pre.4, [L37] 0.4.0-pre.3,
[EX] 1.4.0 (unchanged).

\subsection*{v1.19.0-pre.4 --- 2026-10-04 --- hecke-extensions slot 19, WP-2 continued: WP-LV2 and WP-HL built}

\noindent A pre-release, \textbf{not owner-released}, made with [E] 1.19.0-pre.5, [U] 0.3.0-pre.4 and [L37] 0.4.0-pre.2 (change table
\file{docs/implementation/slot19/WP-2\_pre5\_changes.md}). \textbf{Code-map text only.}
\begin{itemize}\setlength{\itemsep}{0pt}
\item \S\ref{sec:codemap}: ``At v1.19.0-pre.4'' in the rows of \S\ref{sec:nebentypus} (WP-LV2: the twisted route built), \S\ref{sec:hidafull} (WP-HL built,
  \code{HeckeLambdaValidation}) and \S\ref{sec:saturation} (WP-LV2's free part and deep rungs).
\item \S\ref{sec:notimplemented}: a v1.19.0-pre.4 box --- built: WP-LV twisted (WP-LV2), WP-HL.
\item Header and \verb|\date|.
\end{itemize}
\textbf{No statement's content changes; no label is added, renamed or removed; no tag moves}; Appendices~\ref{app:status} and \ref{app:labelindex} are unchanged; the
milestone-M3 \code{\textbackslash LEAN} markers still wait for the slot-19 close (F3's list). Companion pointers: [E] 1.19.0-pre.5, [U] 0.3.0-pre.4, [L37] 0.4.0-pre.2,
[EX] 1.4.0 (unchanged).

\subsection*{v1.19.0-pre.3 --- 2026-10-04 --- hecke-extensions slot 19, WP-2 continued: the slot-19 code maps as built}

\noindent A pre-release, \textbf{not owner-released}, made with [E] 1.19.0-pre.3, [U] 0.3.0-pre.3 and [L37] 0.4.0-pre.1 after the slot-19 Fable referee F4
(\file{docs/implementation/slot19/R19-6\_F4\_referee\_report.md}; change table \file{docs/implementation/slot19/WP-2\_pre3\_changes.md}). \textbf{Code-map text only.}
\begin{itemize}\setlength{\itemsep}{0pt}
\item \S\ref{sec:codemap}: ``At v1.19.0-pre.3'' in the rows of \S\ref{sec:saturation} (WP-LC and WP-LV untwisted: built and merged), \S\ref{sec:nebentypus} (WP-LV
  twisted: not built) and \S\ref{sec:hidafull} (WP-HL: not built; R19-2's prototype).
\item \S\ref{sec:notimplemented}: a v1.19.0-pre.3 box --- built: WP-LC, WP-LV untwisted, WP-NF, NB-2; not built: WP-HL, WP-LV twisted.
\item Header and \verb|\date|.
\end{itemize}
\textbf{No statement's content changes; no label is added, renamed or removed; no tag moves}; Appendices~\ref{app:status} and \ref{app:labelindex} are unchanged; the
milestone-M3 \code{\textbackslash LEAN} markers still wait for the slot-19 close (F3's list). Companion pointers: [E] 1.19.0-pre.3, [U] 0.3.0-pre.3, [L37] 0.4.0-pre.1,
[EX] 1.4.0 (unchanged).

\subsection*{v1.19.0-pre.2 --- 2026-10-03 --- hecke-extensions slot 19, WP-3 repairs: typography}

\noindent A pre-release, \textbf{not owner-released}. The slot-19 referee F1 (\file{docs/implementation/slot19/WP-3\_F1\_referee\_report.md}) ACCEPTED v1.19.0-pre.1
with one typographic note (m-C1): in the \S\ref{sec:codemap} rows of \S\ref{sec:hidafull}, \S\ref{sec:nebentypus} and \S\ref{sec:saturation} the space before the full stop
preceding ``Planned at v1.19.0-pre.1'' is removed and each planned sentence ends with a period. The \S\ref{sec:saturation} row's WP-LV pointer now reads ``on a saturated
block (\code{stable} at $K=1$)'', following [E] 1.19.0-pre.2's repair M-1 of \Espec{cl:lveis}(3). \textbf{No statement, tag, label or count changes.} Companion
pointers: [E] 1.19.0-pre.2, [U] 0.3.0-pre.2, [EX] 1.4.0 (unchanged).

\subsection*{v1.19.0-pre.1 --- 2026-10-03 --- hecke-extensions slot 19, WP-2: code-map text for slot 19}

\noindent\emph{This version.} A pre-release, \textbf{not owner-released}, for the slot-19 Fable referee F1 (WP-3). \textbf{No statement's content changes; no label is
added, renamed or removed; no tag moves}; Appendices~\ref{app:status} and \ref{app:labelindex} are unchanged.
\begin{itemize}\setlength{\itemsep}{0pt}
\item \S\ref{sec:notimplemented}: a v1.19.0-pre.1 box --- WP-LC, WP-LV, WP-HL, WP-NF and NB-2, what each reads of this document and that none is built.
\item \S\ref{sec:codemap}: ``Planned at v1.19.0-pre.1'' in the rows of \S\ref{sec:hidafull} (WP-HL), \S\ref{sec:nebentypus} (WP-LV, twisted) and \S\ref{sec:saturation}
  (WP-LC; WP-LV, untwisted).
\item Not in this pre-release: the milestone-M3 \code{\textbackslash LEAN} markers (they follow slot-19 L19-6 and the Lean audit F3), and the measurements of
  Remark~\ref{rmk:levelNscale}(4)'s window (the open repair of the slot-18 F1, pending slot-19 WP-RC's runs).
\item Header and \verb|\date|.
\end{itemize}
Companion pointers: [E] 1.19.0-pre.1, [U] 0.3.0-pre.1, [EX] 1.4.0 (unchanged).

\subsection*{v1.18.0 --- 2026-10-03 --- hecke-extensions slot 18, the close: the release, from v1.18.0-pre.2, with the milestone-M2 \texorpdfstring{\code{\textbackslash LEAN}}{LEAN} markers and an erratum}

\noindent\emph{This version.} Released at the slot-18 close (WP-18) under the owner's ruling O3 of 2026-10-03 (``release [C] 1.18.0, [E] 1.18.0,
[EX] 1.4.0''), on the slot-18 referee F1's ACCEPT WITH REPAIRS (repairs at pre.2). \textbf{No statement's content changes; no label is added, renamed or
removed; no tag moves} relative to pre.2; Appendices~\ref{app:status} and \ref{app:labelindex} are unchanged (markers are not counted).
\begin{itemize}\setlength{\itemsep}{0pt}
\item \textbf{Erratum} (\textbf{confirmed by the owner}, 2026-10-03): the proof of \cl{thm:welldef}(c) --- representatives $\{\gamma_1s_j\}$
  of $\Ga''_1\backslash\Ga$, not $\{s_j\gamma_1^{-1}\}$ --- recorded at pre.2 and \textbf{confirmed} by the slot-18 Lean audit F3 with a counterexample
  ($\Ga''=\Ga_0(p)$, $\gamma_1=S$: the $p$ elements $\gamma_1^{-1}s_j$, $j<p$, lie in one coset), now written into the proof's parenthesis. The statement,
  its tag and its conclusion are unchanged.
\item \textbf{The milestone-M2 markers}, clause by clause per F3 \S8 and its repairs ($39$ new markers, $36$ \code{\textbackslash LEAN} and $3$
  \code{\textbackslash LEANlit}; $22$ newly tagged labels; Lean census \code{decls=823 proved=764 proved\_mod\_lit=34 stated=1 sorries=1 global\_axioms=0
  hyp\_bad=0 verdict=OK}): G2 --- \cl{def:model7} (the simplex list), \cl{def:model7rel} (the eight rows), \cl{lem:model7coincidences}(1),(2),
  \cl{prop:model7h4} ((i) without $\Ku\cong\Xbulk$; (ii) the dimension only; (iii) \ref{FD1}, \ref{FD2}; (iv)), \cl{def:eq_cochains}, \cl{lem:eqgroupoid},
  \cl{lemma:d_maps_C_eq}, \cl{thm:explicit_complex_equiv}(1), \cl{thm:welldef} (the first bullet, (a) at cochain level, (c)'s transported half), \cl{def:hecke},
  \cl{prop:cosets}(3) now with ``group isomorphism'' and (5) for any transversal, \cl{prop:mu}(1), \cl{prop:bdrycosets}(2) with the cusp characterisation,
  \cl{prop:bdrychar}(1)'s first sentence; G3 --- \cl{obl:taubdry}(i)--(iv) and \cl{obl:eta}(i)--(iv) (verified at the listed $(p,\nu)$),
  \cl{def:pairhomotopy}, \cl{lem:conerelative}(2),(3),(5) (given (4)), \cl{cl:les}(2)'s injectivity and Euler identity, \cl{prop:coneambiguity}(2) over a
  field, \cl{lem:etaexists}, \cl{lem:etaunique}, \cl{thm:conewelldef}(a),(b),(d), \cl{prop:coneindep}(1); G10 --- \cl{lem:heckerankone} and
  \cl{prop:cohforms}(3), (4)'s second surjection (\code{\textbackslash LEANlit}, modulo \texttt{hES} / \texttt{hF}), \cl{lem:freepart}(2) in full.
  Not tagged: \ref{H4}(ii)'s contractibility, \ref{FD3}, \ref{FD4}, \cl{thm:welldef}(b) and (c)'s general case, \cl{thm:explicit_complex_equiv}(2),
  \cl{obl:taubdry}(v), \cl{obl:eta}(v),(vi), \cl{lem:tauexists}, \cl{rmk:orbifold} (G4: only the abstract projective reduction is in [GT], slot-18 L18-7).
\item \S\ref{sec:status}: the v1.18.0 list named. Header and \verb|\date|.
\end{itemize}
Companion pointers: [E] 1.18.0, [EX] 1.4.0, [U] 0.2.1 (unchanged).

\subsection*{v1.18.0-pre.2 --- 2026-10-03 --- hecke-extensions slot 18, WP-3 repairs: the referee F1's repairs, the slot's code maps, and an erratum}

\noindent\emph{This version.} A pre-release, \textbf{not owner-released}; the release at the slot-18 close on the owner's yes (ruling~O3). The slot-18 Fable
referee F1 (\file{docs/implementation/slot18/WP-3\_F1\_referee\_report.md}) read v1.18.0-pre.1 as \textbf{ACCEPT WITH REPAIRS}, $0$ blocking, $1$ major
(M-1), $2$ minor; all applied here (\file{docs/implementation/slot18/WP-3\_repairs.md}). No statement's content changes, no label is added, renamed or
removed; two tags move; Appendix~\ref{app:status}'s counts are unchanged (the \cl{IL3}, \cl{IL4} items are not rows of it).
\begin{itemize}\setlength{\itemsep}{0pt}
\item \textbf{M-1} \cl{IL5}: the proof route is \cl{prop:strictlift}(3)'s argument with \cl{prop:isomodelell}(4),(5) in place of \cl{rmk:borelfree} (the
  model over $Y'$ is projective, not free); the hypothesis ``$\Iso^{(\ell)}(Y')$ is an $\ell$-local isotropy model over $Y'$'' is stated, (T1) is over
  $\Ga''$, and the strict identity is named a certificate of the built map, not a hypothesis. The \PROVED{} tag is unchanged; the item gains its
  \COMP\textsuperscript{a} line (slot-18 WP-7, arm (IL5)). \textbf{m-C2}: at $\ell=2$ the item's checks are a code-path check.
\item \textbf{m-C1} \cl{rmk:levelNscale}(4): the figures given per cell ($99330$ at $k=2$: $m_N=D_N=304\,128$, $10\,991$~MiB; $90090$ at $k=4$:
  $m_N=290\,304$, $D_N=870\,912$, $9622$~MiB). And the \emph{outcome} (slot-18 WP-5): the structure-aware fill model \textbf{failed} its factor-$1.5$ rule,
  so \code{hook\_rss/v1} and \code{CokernelCostModel}'s law stay (a finding); W5-F1's root cause is the generator order (the fill at $99330$, $k=2$, is
  $54.45$~M); the sorted order is built, value-free, gated on $21$ rungs; the two window measurements are pending.
\item \textbf{Code maps as built} (slot-18 WP-5, WP-6, WP-7, merged): \S\ref{sec:levelN} (``Built at v1.18.0-pre.2''), \S\ref{sec:nebentypus} (items
  (iv)--(vi) built; \cl{rmk:nebcost}(4)), \S\ref{sec:isomodelell} (``Run at v1.18.0-pre.2'': (v), (vi); the Scope paragraph), \S\ref{sec:codemap}'s three rows,
  \S\ref{sec:notimplemented}'s box; each pre.1 plan kept for the record.
\item \textbf{Tag move} (the referee F1's C-4: ``move both at the close on WP-7's \cl{IL5} + (vi) pass''; WP-7 passed, \code{--all} $156/156$): \cl{IL3},
  \cl{IL4} from \COMP\textsuperscript{a} \emph{partial} to \COMP\textsuperscript{a}.
\item \textbf{Erratum, for the owner's confirmation} (slot-18 L18-2, the Lean proof \texttt{GL2.Model.thetaG\_transported}): the proof of
  \cl{thm:welldef}(c) took $\{s_j\gamma_1^{-1}\}$ as representatives of $\Ga''_1\backslash\Ga$ for $g_1=\gamma_1g\gamma_2$; that family is not a transversal in
  general. The correct representatives are $\{\gamma_1s_j\}$, with $(\gamma_1s_j)^{-1}g_1=\delta_j\gamma_2$; with them the operator is reproduced verbatim. The
  statement (c), its tag and its conclusion are unchanged; only the proof text is corrected (and the sentence on which factors cancel).
\item \textbf{Made explicit} (slot-18 L18-4, \texttt{GL2.Cone.eta\_exists}'s hypothesis \texttt{hjaug}): \cl{def:pair} states $\aug\circ j=\aug_\partial$, true
  because $j$ is the inclusion of a subcomplex, and \cl{lem:etaexists}(V) cites it; no statement changes.
\end{itemize}
Companion pointers: [E] 1.18.0-pre.2 (the referee F1's M-2 and minors; Prediction \texttt{pr:depth}'s $(2,3)$ datum), [EX] 1.4.0-pre.2, [U] 0.2.1 (unchanged).

\subsection*{v1.18.0-pre.1 --- 2026-10-03 --- hecke-extensions slot 18, WP-2: contract first for slot 18, and the Eichler--Shimura locator}

\noindent\emph{This version.} A pre-release, \textbf{not owner-released}, for the slot-18 Fable referee F1 (WP-3); the release at the slot-18 close only on its
acceptance and the owner's yes (ruling~O3 of 2026-10-03; the \code{\textbackslash LEAN} tags of slot~18's F3 list enter only at that release, WP-18). \textbf{No
statement's content changes, no label is renamed or removed, no tag moves}; one checkable item is added; Appendix~\ref{app:status}'s counts are unchanged.
\begin{itemize}\setlength{\itemsep}{0pt}
\item \textbf{New item} \cl{IL5} in \cl{rmk:isomodellcheckable} (\PROVED{} given (T1)--(T3) of the built map; committed before slot-18 WP-7's code, rule~20): the
  strict action on a synthetic refined descriptor agrees on cohomology with the transported one. \cl{app:labelindex} gains it.
\item \textbf{Contract first} (rule~21; nothing built): \S\ref{sec:levelN}'s code map (slot-18 WP-5) and \cl{rmk:levelNscale}(4)'s paragraph \emph{the
  structure-aware fill model} (a cost statement, with its replacement rule; the two window measurements; the order change gated by bit-identity);
  \S\ref{sec:nebentypus}'s code map (slot-18 WP-6: items (iv) $W_K(\F_q)$ at $d(\chi)\ge2$, (v) the model-route $T_{p^{2}}$ through
  $[\mathrm{diag}(1,p^{2})]_\chi+\chi(p)p^{\,n}$ against \cl{thm:nebrelations}(2), (vi) the WP-6 audit minors) and \cl{rmk:nebcost}(4)'s pointer;
  \S\ref{sec:isomodelell}'s code map (slot-18 WP-7: (v) the strict lift on a refined descriptor and its MP requirements, (vi) \cl{IL3} over $\F_9$ and
  $\Z/9$); \S\ref{sec:codemap}'s three rows; a new box in \S\ref{sec:notimplemented}.
\item \textbf{Locator} (the slot-17 referee F1's m-C4): \cl{lem:heckerankone}'s Eichler--Shimura input is \cite[Thm.~12.2.2, Prop.~12.4.10]{diamond_im},
  verified on page images (slot-17 L17-5 repairs); \cite[Ch.~8]{shimura} kept as an unverified secondary; \cl{app:status}'s row follows.
\end{itemize}
Companion pointers: [E] 1.18.0-pre.1 (new: Claims \texttt{cl:seedtwist}, \texttt{cl:selfmirror}, \texttt{cl:dimformula}, Prediction \texttt{pr:depth};
\texttt{pr:kfonset} generalised), [EX] 1.4.0-pre.1, [U] 0.2.1 (unchanged).

\subsection*{v1.17.0 --- 2026-10-02 --- hecke-extensions slot 17, the close: the release, from v1.17.0-pre.5, with the first \texorpdfstring{\code{\textbackslash LEAN}}{LEAN} markers}

\noindent\emph{This version.} The slot's $X$ bump, whose text v1.17.0-pre.1 \dots\ pre.5 carried, is released at the slot-17 close (WP-18) under the
owner's ruling O3 of 2026-10-01 (``release [C] v1.17.0, [E] 1.17.0, [U] 0.2.1 and [EX] 1.3.0 on F1's accept, with the \code{\textbackslash LEAN} tags'').
Its referee (slot-17 WP-3, Fable, \file{docs/implementation/slot17/WP-3\_referee\_F1.md}) was \textbf{ACCEPT WITH REPAIRS} for this document, $0$ blocking,
$1$ major (M-1), $5$ minor; the repairs were applied at pre.3, and M-1 was settled by the owner's ruling of 2026-10-02 at pre.4. \textbf{No statement is added,
no label is renamed or removed, and no tag moves} relative to pre.5; Appendices~\ref{app:status} and \ref{app:labelindex} are unchanged.

\paragraph{Over pre.5: the Lean milestone M1.}
\begin{itemize}\setlength{\itemsep}{0pt}
\item \textbf{The first markers}, exactly per \file{docs/implementation/slot17/lean\_taggable\_labels.md} (the F3 audit's clause-scoped list after its
  repairs; Lean census \code{decls=440 proved=383 proved\_mod\_lit=20 stated=13 sorries=13 global\_axioms=0 hyp\_bad=0 verdict=OK}): $39$
  \code{\textbackslash LEAN} and $3$ \code{\textbackslash LEANlit} on $20$ labels --- Part~0: \cl{prop:symconv}(2),(3), \cl{lem:detsym}, \cl{def:contragredient}
  (\eqref{eq:dualaction}), \cl{prop:pairing}, \cl{lem:symdual}; G1: \cl{def:cosetreps}, \cl{prop:cosets}(1)--(5) with the list's restrictions, \cl{lem:cocycle},
  \cl{lem:bounds}(1)--(4) restricted, \cl{prop:tworoutes}, \cl{prop:bdrycosets}(1)--(4), \cl{prop:bdrychar}(2); G10: \cl{def:satdata}, \cl{prop:saturation}(1)--(4)
  and (5)'s first sentence, \cl{prop:satdefect}(1)--(3), \cl{lem:freepart}(1)--(3), \cl{cor:stablecert}(1) and (2)'s first sentence (\code{\textbackslash LEANlit},
  modulo Lemma~\ref{lem:heckerankone}), \cl{prop:cohforms}(1),(2),(4) restricted, \cl{prop:hidaalgconst}(1) (\code{\textbackslash LEANlit}, modulo
  Theorem~\ref{thm:hidacontrol}(3)), \cl{rmk:thetaforms}(1)'s two power-series facts. Not tagged: \cl{lem:heckerankone}, \cl{prop:cohforms}(3),(5),(6),
  \cl{prop:mu}, \cl{prop:bdrychar}(1),(3), \cl{cor:stablecert}(3), \cl{lem:bounds}(5),(6), and the G3 cone statements (milestone M3).
\item \S\ref{sec:status}: the marker is clause-scoped, a definition may carry it, and the v1.17.0 list is named.
\end{itemize}
Companion pointers: [E] 1.17.0, [U] 0.2.1, [EX] 1.3.0.

\subsection*{v1.17.0-pre.5 --- 2026-10-02 --- hecke-extensions slot 17, WP-2: verified locators and the WP-7 code map}

\noindent\emph{This version.} A pre-release, \textbf{not owner-released}. Pointers proposed for WP-2 by other slot-17 packages; no statement, tag, label or count
changes (\file{docs/implementation/slot17/WP-2\_report.md}~\S8).
\begin{itemize}\setlength{\itemsep}{0pt}
\item \textbf{R17-4} (\file{docs/implementation/slot17/R17-4\_report.md}, locators read on page images). \ref{NB4}(c): the \STD{} input is
  \cite[Lemma~1.1.1, p.~194]{ash_stevens} (the long exact sequence; the referee's recollection ``Lemma~1.1.3'' is restriction / corestriction) and
  \cite[Def.~1.1.2, p.~194; Lemma~1.1.4, p.~195]{ash_stevens} (Shapiro for a compatible pair), \textbf{verified}; the ``not verified'' flag is dropped, the
  \PROVED{} tag of v1.17.0-pre.4 stands and the \COMP\textsuperscript{r} fallback is not triggered; compatibility~(b) is stated in Ash--Stevens' orientation
  ($\Ga\De_0(\ell)$), with~(c). Proposition~\ref{prop:barformula} and \eqref{eq:corconjres}: ``in exactly this form'' / ``in this exact form'' dropped --- Ash--Stevens
  give the right-module form on homogeneous cochains (p.~194). Remark~\ref{rmk:thetaclassical}: item~(3) cites \cite[Lemma~2.1, p.~204]{ash_stevens} for the
  inclusion only, the equality and the twist bookkeeping being \Espec{cl:glover}'s and Proposition~\ref{prop:thetatwist}'s; item~(2) cites
  \cite[Thm.~4.1, p.~276]{jochnowitz}. Lemma~\ref{lem:heckerankone}: the eigenbasis from \cite[\S\S3.5--3.6, pp.~54--55]{diamond_im} (in the library,
  verified; new bibliography key), Shimura and Diamond--Shurman kept as unverified secondary pointers. \ref{LN5}: \cite[Th\'eor\`eme~2, p.~72]{merel96},
  verified, with its hypotheses ($p>3$, $p\mid(N-1)/\gcd(N-1,12)$, in this document's letters). The \cite{ash_stevens} sites of
  Proposition~\ref{prop:relationsmodel}'s \STD{} gain the page and lose ``not verified''. \cl{app:status}: the rows of \cl{prop:barformula},
  \cl{rmk:thetaclassical}, \cl{rmk:levelNcheckable}, \cl{rmk:nebcheckable}, \cl{lem:heckerankone} name the locators; no count moves.
\item \textbf{WP-7} (\file{docs/implementation/slot17/WP-7\_report.md}~\S5, checked against \file{docs/implementation/slot17/WP-7\_audit.md}): \ref{IL3}'s and
  \ref{IL4}'s status sentences record the runs (\ref{IL3} at level $N$ against the pipeline over $\F_3$ only, MINOR-3; \ref{IL4} over $\Z/3^{e}$ and the pair and
  cone); the code map records the descriptor as data, the synthetic refined descriptor, and arm (REF) as a descriptor round trip (MINOR-1); the
  \S\ref{sec:codemap} row and the \S\ref{sec:notimplemented} box follow. The tags stay \COMP\textsuperscript{a} partial; the move is for the next referee.
\item \textbf{Lean} (\file{docs/implementation/slot17/lean\_taggable\_labels.md}): \S\ref{sec:status} notes that Corollary~\ref{cor:stablecert} will take
  \code{\textbackslash LEANlit} (its [GT] proof carries Lemma~\ref{lem:heckerankone} as a hypothesis); no marker is used at this version.
\item The [L37] referee F4's pointers: \Espec{pr:kfonset}~(c) records the window's last cell off~$37$ in [E] 1.17.0-pre.5; no sentence of this document
  states ``$e=2$ at the eleven cells''.
\end{itemize}
Companion pointer: [E] 1.17.0-pre.5; [EX] 1.3.0-pre.4 and [U] 0.2.1-pre.2 unchanged.

\subsection*{v1.17.0-pre.4 --- 2026-10-02 --- hecke-extensions slot 17, WP-2: the owner's rulings on the referee F1's M-1 and B-1}

\noindent\emph{This version.} A pre-release, \textbf{not owner-released}. Two owner rulings of 2026-10-02 on the slot-17 Fable referee F1
(\file{docs/implementation/slot17/WP-3\_referee\_F1.md}):
\begin{itemize}\setlength{\itemsep}{0pt}
\item \textbf{M-1} \ref{NB4}(c): the sharp form at $n=0$ is tagged \PROVED{} (given the one \STD{} input named), as the referee proposed, with the locator
  \cite[\S1.1]{ash_stevens} flagged \textbf{not verified} --- the pattern of \cl{lem:heckerankone}. At pre.3 the referee's alternative (\COMP\textsuperscript{r}
  only, \PROVED{} on verification) had been applied; the owner chose the proposed tag. Reading the locator stays on the owner's scan list
  (\file{docs/implementation/slot17/L17-5\_report.md}~\S3). \cl{app:status}: \cl{rmk:nebcheckable}'s tag set already holds \PROVED{}, so it is unchanged and no count
  moves; the counts paragraph records it. The \S\ref{sec:nebentypus} scope paragraph follows.
\item \textbf{B-1} (the companion's \Espec{pr:latticesplit}~(R), restated on \emph{congruence-free} classes in [E] 1.17.0-pre.4): the two dependent sentences
  repaired at pre.3 (m-C3) take their final wording --- \ref{LN6}'s route-R3 sentence (level~$1$ shows Loewy length $\ge3$ only on classes that are not
  congruence-free) and \S\ref{sec:saturation}'s scope paragraph (the pre.2 Kummer-free form contradicted on the forms side, the regular $\ell=139$ among the
  primes; (R) now on congruence-free classes).
\end{itemize}
No label added, renamed or removed; no count moves. Companion pointer: [E] 1.17.0-pre.4, [EX] 1.3.0-pre.4 in the tree.

\subsection*{v1.17.0-pre.3 --- 2026-10-02 --- hecke-extensions slot 17, WP-2 resumed: the referee F1's repairs}

\noindent\emph{This version.} A pre-release, \textbf{not owner-released}. The slot-17 Fable referee F1 (\file{docs/implementation/slot17/WP-3\_referee\_F1.md})
read v1.17.0-pre.2 as \textbf{ACCEPT WITH REPAIRS}, $0$ blocking, $1$ major (M-1), $5$ minor, and WP-6's proposed code-map text as acceptable with repairs (M-2,
m-W1--m-W3). Every repair is applied as worded except where stated; the referee file's ``Repairs'' table lists each.
\begin{itemize}\setlength{\itemsep}{0pt}
\item \textbf{M-1} \ref{NB4}(c): the orientation step is proved in this document's conventions (Proposition~\ref{prop:symconv}, Definition~\ref{def:inducedmodule},
  Definition~\ref{def:nebdata}(5)): $0\to\Sym^{\ell-1-i}\otimes{\det}^{i}\to I\to\Sym^{i}\to0$ is exact, so $\Sym^{i}$ is the untwisted constituent, and the second
  \STD{} (the composition factors) is no longer needed. The remaining \STD{} input carries the locator \cite[\S1.1]{ash_stevens}, \textbf{not verified} (no copy in
  the library); the referee's alternative therefore applies --- the sharp form at $n=0$ is re-tagged \CONJ{} $\to$ \COMP\textsuperscript{r}, and \CONJ{} $\to$
  \PROVED{} is made when the locator is verified. \cl{app:status}: the \texttt{ash\_stevens} cite gains $[\S1.1]$; the tag set already holds \COMP\textsuperscript{r},
  so no count moves.
\item \textbf{M-2, m-W1--m-W3} \S\ref{sec:nebentypus}'s code map: slot-17 WP-6 recorded as built at this version (merged, audited ACCEPT): the sparse cokernel with
  a character over $\F_q$, and $H_{N,p^{\nu},\chi}$ from $\mathrm{Mer}_{p^{\nu}}$ --- \emph{tested, not proved} to be $T_{p^{\nu}}$ (Merel at composite determinant is
  not a statement of this document); the model-route double coset of Proposition~\ref{prop:nebcorr} not built. The \S\ref{sec:codemap} row and
  \S\ref{sec:notimplemented}'s v1.17.0-pre.1 box follow.
\item \textbf{WP-5} \S\ref{sec:levelN}'s code map: slot-17 WP-5 recorded as built (merged, audited ACCEPT WITH REPAIRS): the coset products in \code{int64} then
  \code{\_\_int128}, bit-identical below the old reach; the ladder to $N=90090$ at $k=2$, $60060$ at $k=4$; the predictions' verdicts held and their numbers did
  not all hold (WP-5 audit MINOR-1). \cl{rmk:levelNscale}(4): the $N=99330$ point (an open measurement and a failed prediction, WP-5 audit MINOR-3).
\item \textbf{m-C1} \cl{rmk:sattor}(4): ``every order is saturated on its $\Z_\ell$-dual''. \textbf{m-C3} two sites cited [E] \Espec{pr:latticesplit} as ``the census at
  level~$1$ never shows Loewy length $\ge3$'': \ref{LN6}'s route-R3 sentence and \S\ref{sec:saturation}'s scope paragraph, both repaired; pre.2's changelog sentence
  ``nothing in this document cites (R)'' is corrected in place. \textbf{m-C4} \cl{lem:heckerankone}: \cite[Thm.~5.5.4]{diamond_shurman} named, still not verified.
  \textbf{m-C5} \S\ref{sec:isomodelell}'s code map, item~(iv): a control cell with $-I\mapsto-1$. \textbf{m-C2} (gate \code{n12\_held == 240} in arm (X4)) is a
  carry for a code package; no text change.
\end{itemize}
No label added, renamed or removed; no tag set of a numbered statement changes; no count moves. Companion pointer: [E] 1.17.0-pre.3 in the tree, whose
\Espec{pr:latticesplit}~(R) awaits the owner's ruling on the referee's B-1 (\cl{LN6}'s and \S\ref{sec:saturation}'s repaired sentences follow its final form).

\subsection*{v1.17.0-pre.2 --- 2026-10-02 --- hecke-extensions slot 17, WP-2, second pass: the eigenbasis citation of Lemma~\ref{lem:heckerankone}}

\noindent\emph{This version.} A pre-release, \textbf{not owner-released}, for the same referee as pre.1 (slot-17 WP-3, Fable, window F1). One citation
repair, the one proposed in the addendum of \file{docs/implementation/slot17/WP-3\_referee\_brief.md} (from slot-17 L17-5's report \S7, finding~2, and its
audit's (a)):
\begin{itemize}
\item \cl{lem:heckerankone}'s proof cites its two \STD{} inputs separately. \emph{The eigenbasis} of $S_k(\Cx)$ (commuting Hecke operators, self-adjoint for
  the Petersson product): \cite[Thm.~3.41]{shimura} or \cite[\S5.5]{diamond_shurman}, in place of pre.1's Shimura Thms.~3.24, 3.29, 3.34, 3.35. \textbf{Neither
  locator is verified} (rule~\ref{A2}): neither book is in the project's library, and both pointers were relayed by L17-5 unread. \emph{Multiplicity one}
  (the distinctness of the systems): \cite[Ch.~VII, \S5, Thm.~7 and its Cor.~1, p.~102; Prop.~11, p.~101]{serre_cours}, verified on page images by L17-5.
  Serre's \S5 gives only this half.
\item \cl{app:status}'s row for \cl{lem:heckerankone} names \cite[Thm.~3.41]{shimura} and \cite[\S5.5]{diamond_shurman}; the bibliography gains
  \texttt{diamond\_shurman}.
\end{itemize}
No statement, tag, label or count changes; no number moves. The companion's 1.17.0-pre.2 (in the tree) restates \Espec{pr:latticesplit}~(R) again (owner ruling
O-R17-2a of 2026-10-02); nothing in this document cites (R). \emph{(Corrected at v1.17.0-pre.3: two sites do --- \S\ref{sec:saturation}'s scope paragraph and
\ref{LN6}'s route-R3 sentence; both repaired there.)}

\subsection*{v1.17.0-pre.1 --- 2026-10-01 --- hecke-extensions slot 17, WP-2 (drafted during the slot-16 close): the deferred items, for one referee}

\noindent\emph{This version.} A pre-release, \textbf{not owner-released}. It drafts, for a single referee (slot-17 WP-3, Fable, window F1), the items that
the slot-16 release left out because they are new mathematics (\file{docs/implementation/slot16/contract\_release\_report.md}~\S4), and writes slot~17's
code-map texts before the code (rule~21). Release at the slot-17 close only on the referee's acceptance and the owner's yes (ruling O3 of 2026-10-01). Report
\file{docs/implementation/slot17/WP-2\_report.md}; referee brief \file{docs/implementation/slot17/WP-3\_referee\_brief.md}.

\paragraph{Statements.}
\begin{itemize}\setlength{\itemsep}{0pt}
\item \textbf{New: Remark~\ref{rmk:sattor}}, in a new last subsection \S\ref{sec:sattorsec} of \S\ref{sec:saturation}, so that no statement number of v1.16.0 moves (\PROVED; slot-16 WP-5, \file{WP-5\_tor\_note.md}~\S1, \S5): the saturation defect's $\mathrm{Tor}$ term
  $I=\mathrm{Tor}_1^{\Z_\ell}(\Talg^{\mathrm{sat}}/\Talg,\Fl)=(\Talg\cap\ell\Talg^{\mathrm{sat}})/\ell\Talg=\ker(\Talg/\ell\to\End(L/\ell))$, its dimension $\varepsilon^{\mathrm{sat}}_1$,
  its reading at a stable $K$, and its distinction from $\mathrm{Tor}_1^{\Talg}(L,\Talg/\midl)$. Over the note: item~(4)'s ``the converse fails'' is proved here
  by an example ($\Talg^{\vee}$ over the non-Gorenstein order of Corollary~\ref{cor:stablecert}(2)), which the note asserted without one.
\item \textbf{\ref{NB4}} (slot-16 WP-11, \file{WP-11\_report.md}~\S1, \S4): (a) the sharp form at $n=0$ recorded as \COMP\textsuperscript{r} ($132/132$, $\ell\le43$),
  with the orientation it pins; (b) the $n>0$ reading (Clebsch--Gordan, the $H^{0}$ correction at $i=n$; $240/240$, reported); (c) a proof of (a) by the
  Jordan--H\"older long exact sequence, \textbf{proposed}: the referee accepts it (\CONJ{} $\to$ \PROVED{} for the sharp form at $n=0$) or downgrades the re-tag
  to \CONJ{} $\to$ \COMP\textsuperscript{r}. The general statement stays \CONJ.
\item \cl{rmk:levelNscale}(4): slot-16 WP-13's F-3 --- the fitted fill law underestimates \code{final\_nnz} at composite $N$, $k>2$, by $1.47$--$6.80$
  (\COMP\textsuperscript{a}); a statement about the cost model only.
\item \cl{lem:heckerankone}'s proof: the eigenbasis of $M_k(\Cx)$ is cited from Petersson's adjointness (\cite{shimura} Ch.~3, \textbf{VERIFY}), the
  multiplicity-one part from Serre's \emph{Cours} Ch.~VII \S5 with verified page locators (slot-17 L17-5, relayed by the lead): Serre's \S5 alone gives the
  distinctness of the systems, not the existence of the basis. No statement or tag changes. (Its locators were replaced at v1.17.0-pre.2, below.)
\end{itemize}

\paragraph{Code maps, contract first (no statement changes).} \S\ref{sec:levelN}: the table's reach beyond composite $N\approx5.5\cdot10^{4}$ (slot-17
WP-5, F-1). \S\ref{sec:nebentypus}: the sparse cokernel with a character and $T^{(N,\chi)}_{p^{\nu}}$, $\nu\ge2$ (WP-6). \S\ref{sec:isomodelell}: \ref{IL3} at
level $N$, \ref{IL4} over $\Z/3^{e}$, the pair and cone half of \cl{prop:isomodelell}(5), and what the MP constraint requires of a synthetic refined descriptor
(WP-7). \S\ref{sec:saturation}: slot-16 WP-5's \file{HeckeExt/HeckeSaturationDefect} recorded as built. \S\ref{sec:codemap}'s row and a
\S\ref{sec:notimplemented} box follow.

\paragraph{Other.} \S\ref{sec:status}: the \code{\textbackslash LEAN} / \code{\textbackslash LEANlit} marker (the Lean programme's Q9, owner-ruled
2026-10-01), defined and \textbf{unused}; it is not a status tag and is not counted. Appendix~\ref{app:status}: one line added, two tag sets grown (its
counts paragraph). Appendix~\ref{app:labelindex}: \cl{sec:sattorsec}, \cl{rmk:sattor}. Not in this version: slot-17 R17-1's double-divisibility observation (not merged
when this draft was written; [E]'s owner may add it in a later patch). Non-normative: the header comment, the date line; the companion pointer stays
at [E] v1.16.3 (its 1.17.0-pre.1 is in the tree, unreleased).

\subsection*{v1.16.0 --- 2026-10-01 --- hecke-extensions slot 16, the close: the release, from v1.16.0-pre.3}

\noindent\emph{This version.} The slot's $X$ bump, whose text v1.16.0-pre.1, pre.2 and pre.3 carried, is released at the slot-16 close under the owner's
ruling O2 of 2026-10-01 (``release [C] v1.16.0 and [U] 0.2.0 since WP-14 accepted and the repairs are confirmed''). Its referee (slot-16 WP-14, Fable,
\file{docs/implementation/slot16/WP-14\_referee.md}) was \textbf{ACCEPT WITH REPAIRS}, $0$ blocking, $1$ major (C-1), $8$ minor; every repair was applied at pre.3
and \textbf{confirmed}, $16$ of $16$ findings with [U]'s, by \file{docs/implementation/slot16/WP-14\_repair\_confirmation.md}. \textbf{No statement is added, no label is
renamed or removed, and no tag moves} relative to pre.3; Appendices~\ref{app:status} and \ref{app:labelindex} are unchanged. The edits over pre.3 are locators,
one sentence of proof commentary, and code-map and status sentences; each is listed with its source in
\file{docs/implementation/slot16/contract\_release\_report.md}.

\paragraph{Over pre.3.}
\begin{itemize}\setlength{\itemsep}{0pt}
\item \textbf{Locators} (slot-16 WP-8, \file{docs/implementation/slot16/WP-8\_citations.md}~\S7): \ref{LN5}'s source restated as \cite[Thm.~1.5.1]{wwe}, verified on the arXiv
  text (Merel's own text not read); \ref{LN6}'s \cite[p.~8, Cor.~7.2.2]{wwe}, verified on the arXiv text. The \CONJ{} tags stay: the chain to
  $H^{1}(\Ga_0(N);\F_p)$ is not carried out. \file{gl2\_refs.bib}: \code{wwe}'s note; \code{haberland} and \code{kiming\_verrill} added for [EX].
\item \textbf{\S\ref{sec:isomodelell}} (the slot-16 WP-16 report \S3 and its audit's F-1, F-7): the proof of \cl{prop:isomodelell}(1) keeps the average and says why it is
  a no-op on $d$ (the lifts start at simplices of dimension $\ge1$, stabiliser $\{\pm I\}$, central) and that it is load-bearing in $\Psi_\ell$; \ref{IL3} and \ref{IL4}
  carry \COMP\textsuperscript{a} \emph{partial} sentences (\ref{IL3}'s level-$N$ clause, \ref{IL4} over $\Z/3^{e}$ and the pair and cone half of
  \cl{prop:isomodelell}(5), not run); the code map names $\Psi_\ell$'s well-definedness check as where the average is tested and records the model as built.
\item \textbf{Code maps} (\S\ref{sec:codemap}, the chapters' own, \S\ref{sec:notimplemented}'s boxes): \S\ref{sec:levelNdegsec} built (slot-16 WP-13,
  \file{Hecke/HeckeLevelNDegen}); \S\ref{sec:nebentypus} built (slot-16 WP-10, WP-11, \file{Hecke/HeckeNebentypus}: \ref{NB4}'s sharp form at $n=0$ gated; $m=2$ and
  $H^{1}_{\mathfrak m}\cong A_W^{2}$ on the dihedral blocks; $H^{1}_{\mathfrak m}\not\cong\Talg_{\mathfrak m}^{2}$ on the F1 hidden blocks over the full Hecke algebra);
  \S\ref{sec:isomodelell} built at $\ell=3$ (slot-16 WP-16, \file{Integral/Model/IsotropyModelEll}).
\item \textbf{Verified, no edit}: the L16-2 audit's note~5 (referee C-9) is applied at pre.3 --- \ref{NB5}'s dictionary is \PROVED{} modulo (L1)--(L5) of
  \textsf{[GT,~\texttt{bridge:nebentypus-dictionary}]}, \CONJ{} only for the literature inputs; the note's conditional second clause (an eigenbasis in place of
  ``the newforms'' if a cell with oldforms is added) has no cell to act on.
\item \textbf{Not in this release} (new mathematics, for a referee): slot-16 WP-5's proposed remark on saturation and torsion, WP-11's proposed statements on
  \ref{NB4} and the hidden blocks, the [L37] report's proposed window bound for \Espec{pr:kfonset}; and, outside this release's brief, WP-13's composite-$N$ caveat
  (its F-3) to the cost figure of \cl{rmk:levelNscale}(4).
\item Non-normative: the header comment, the date line, the companion pointer (now [E] v1.16.2).
\end{itemize}

\subsection*{v1.16.0-pre.3 --- 2026-10-01 --- hecke-extensions slot 16, WP-12 resumed: the WP-14 referee's repairs}

\noindent\emph{This version.} The repairs answering the Fable referee \file{docs/implementation/slot16/WP-14\_referee.md}: [C] \textbf{ACCEPT WITH REPAIRS},
$0$ blocking, $1$ major, $8$ minor; every finding applied as worded. No label added, renamed or removed; no tag set of a numbered statement changes (\cl{rmk:nebcheckable}
already carries \PROVED{}, \BSTD{}, \CONJ{}), so Appendix~\ref{app:status}'s counts are unchanged.

\paragraph{Repairs.}
\begin{itemize}\setlength{\itemsep}{0pt}
\item \textbf{C-1} (major) \cl{prop:isomodelell}(1),(4): every lift produced by the contracting homotopies is averaged, $y\mapsto e_\sigma y$, since $1\otimes e_\sigma$ is
  $Q_\sigma$-fixed; constructively $y_0:=e_\sigma y_0'$, $k:=e_\sigma k'$; \S\ref{sec:isomodelell}'s code map and \ref{IL1} test $d\circ d=0$, \ref{IM1}, \ref{IM2} on the
  built complex.
\item \textbf{C-2} \cl{prop:kappamodel}: on a refinement, $\lift{e_{46}}$ is replaced by $\mathrm{path}(\lift{v_4}\to\lift{v_6})$. \textbf{C-3} \cl{prop:kappamodel}(2) and
  \cl{prop:lhcompliant}(4): the identification is comparison of the projective resolution $C_{\bullet}(\Xbulk)\otimes\cring$, then Shapiro.
\item \textbf{C-4} \ref{LN8}: ``the analogue of \cl{prop:oldspace}(1),(3) with $\Ga_0(M)\to\Ga_0(Mp)$'' (\BSTD{}, \COMP{}), and one cell with $d=p\mid M$ for WP-13;
  \ref{LN9}: the torsion-freeness of $H^{1}(\Ga_0(N);\Sym^{k-2}\Z_p^{2})$ behind $p>k+1$. Same wording in \cl{rmk:degenpresconv}(3).
\item \textbf{C-5} \cl{prop:degenpres}(7), \cl{rmk:levelNscale}(5): at most $2m_N/m_d$ admissible $i$ per $\epsilon$; cost $O(\abs{\mathrm{Mer}_d}m_Md_V^{2})$ with the classes
  enumerated per $\epsilon$. \textbf{C-6} \cl{rmk:levelNscale}(4): \code{SparseCokernel}'s \code{final\_nnz}. \textbf{C-7} \cl{rmk:isomodellmp}(3): the simplices with
  $\abs{P_\sigma}>2$.
\item \textbf{C-8} \ref{NB5}: $H\oplus H^{\mathrm{anti}}$ (not $H\oplus\bar H$ for non-real $\chi$); the characteristic-$0$-to-$\F_q$ clause ($H^{0}=0$ on the cells); the
  reduction rule presumes $d(\chi)=1$. \textbf{C-9} (the owner's call, ruled): \ref{NB5}'s eigenvalue dictionary is \PROVED{} modulo (L1)--(L5) of
  \textsf{[GT,~\texttt{bridge:nebentypus-dictionary}]}, \CONJ{} only for those literature inputs.
\end{itemize}

\subsection*{v1.16.0-pre.2 --- 2026-10-01 --- hecke-extensions slot 16, WP-12 follow-up: (NB5)'s product form squared}

\noindent\emph{This version.} A repair of \ref{NB5} in \cl{rmk:nebcheckable}, ruled by the slot-16 WP-9 / WP-9b audit
(\file{docs/implementation/slot16/WP-9\_audit.md} \S3): on the $\chi$-part the cuspidal factor of the characteristic polynomial of $T_p$, $U_p$ is
$\prod_{f\in S_k(\Ga_0(N),\chi)}(X-a_p(f))^{2}$, not $\prod_f(X-a_p(f))(X-\overline{a_p(f)})$ (the two agree for real $\chi$; the written form failed in
$20$ of $24$ non-real cases, the squared one held in $24$ of $24$). Added: the proof of the squaring by the involution of
$\varepsilon=\mathrm{diag}(1,-1)$, which swaps the Eichler--Shimura halves (\PROVED{}, with \STD{} Eichler--Shimura), and the reduction rule
$\psi(g)=s$ (least primitive root) that makes the mod-$\ell$ reading unambiguous. The eigenvalue dictionary ($\chi$ versus $\bar\chi$) stays \CONJ{} and
L16-2's. No label added, renamed or removed; the tag set of \cl{rmk:nebcheckable} already contains \PROVED{}, \BSTD{} and \CONJ{}, so no count moves.
To be read by the WP-14 referee with v1.16.0-pre.1.

\subsection*{v1.16.0-pre.1 --- 2026-10-01 --- hecke-extensions slot 16, WP-12: degeneracy maps on the presentation, level $N$ at scale, the $\ell$-local isotropy model --- the pre-release of an \texorpdfstring{$X$}{X} bump}

\noindent\emph{This version.} Slot-16 WP-12 (plan \file{docs/implementation/slot16/00-slot16-plan.md} \S2 ``WP-12'', items (a)--(c); owner rulings of
2026-10-01). \textbf{A pre-release, not owner-released}; the Fable referee is slot-16 WP-14, and the release at the slot close needs its ACCEPT and the
owner's sign-off (ruling O2). One subsection added at the end of \S\ref{sec:levelN} and one chapter after \S\ref{sec:saturation}, so that no chapter or
subsection number moves; no existing statement changes, no label is renamed or removed, no tag moves.

\paragraph{New statements.}
\begin{itemize}\setlength{\itemsep}{0pt}
\item \textbf{\S\ref{sec:levelNdegsec}} \PNEW{} (contract first for slot-16 WP-13): \cl{prop:kappamodel} ($\kappa^{\mathrm{mod}}_N(f)=-f(e_{46})$, the pinned
  model-to-presentation map; naturality for fibre maps --- v1.15.0-pre.3's unrefereed scope sentence, now proved; the entrywise comparison of the model's and the
  presentation's operators), \cl{prop:degenpres} (\eqref{eq:degenpres}: $\alpha_d$ on the cokernel presentation by $\mathrm{Mer}_d$ through $\adjg$, for every $d$ with
  $dM\mid N$, proved equal to Proposition~\ref{prop:levelcorrespondence}(4)(c) on $H^{1}$; $d=1$ is $\alpha^{0}$; composite $d$; the trace $\omega\circ\alpha_d=[A\,g\,A]$;
  with a character), \cl{rmk:degenpresconv} (two mis-transcriptions caught; \COMP\textsuperscript{a}), \cl{rmk:levelNscale} (item (c): the chapter's scope at
  $k>2$ and composite $N$, the memory model of the normal form), \cl{rmk:levelNcheckabletwo} (\ref{LN7}--\ref{LN9}).
\item \textbf{\S\ref{sec:isomodelell}} \PNEW{} (item (b)): \cl{def:isomodelell}, \cl{prop:isomodelell}, \cl{prop:ellthreeclosed}, \cl{rmk:isomodellmp},
  \cl{rmk:isomodellcheckable} (\ref{IL1}--\ref{IL4}).
\end{itemize}

\paragraph{Non-normative.}
\begin{itemize}\setlength{\itemsep}{0pt}
\item \S\ref{sec:levelN}'s scope paragraph and code map; \S\ref{sec:codemap}: one row amended, one added; \S\ref{sec:notimplemented}: one box.
  Appendix~\ref{app:status}: $307\to317$ lines; Appendix~\ref{app:labelindex}: the new labels; the provenance counts: \PNEW{} $226\to236$.
\item \file{gl2\_refs.bib}: \code{billerey\_menares} added (not verified). The header comment and the date line. The companion pointer stays at [E] v1.16.1.
\item Evidence scripts (tracked, read-only, not registered validations): \file{docs/implementation/slot16/wp12/degen.py}, \file{negctl.py}. Report:
  \file{docs/implementation/slot16/WP-12\_report.md}, with the referee brief for WP-14.
\end{itemize}

\subsection*{v1.15.0 --- 2026-10-01 --- hecke-extensions slot 15, WP-18: the release at the close, from v1.15.0-pre.3}

\noindent\emph{This version.} The slot's $X$ bump, whose text v1.15.0-pre.1, pre.2 and pre.3 carried, is released at the slot-15 close
under the owner's ruling Q2 of 2026-09-30 (``release [C] v1.15.0 at the close if WP-4 accepts''). Its referee (slot-15 WP-4, Fable,
\file{docs/implementation/slot15/WP-4\_referee.md}) was \textbf{ACCEPT WITH REPAIRS} for \S\ref{sec:nebentypus} and \cl{cor:stablecert}, and
\textbf{ACCEPT} for \cl{prop:saturation}, \cl{prop:satdefect}, \cl{prop:cohforms}, \cl{prop:hidaalgconst} and \cl{rmk:thetaforms}, with $0$
blocking findings; every repair of its \S9 was applied at pre.2. \S\ref{sec:nebentypus} is released \emph{unbuilt} (contract first;
the code is slot~16's, owner ruling Q3). \S\ref{sec:saturation}'s \cl{cor:stablecert} is cited by the
certified \code{stable} verdict of slot-15 WP-5 (Fable audit ACCEPT WITH REPAIRS, repaired). \textbf{No statement is added, no label is
renamed or removed, and no tag moves} relative to pre.3; Appendices~\ref{app:status} and \ref{app:labelindex} are unchanged.

\paragraph{Not refereed (listed for the slot-16 referee).} pre.3's one sentence in \S\ref{sec:levelN}'s scope paragraph (slot-15 WP-17, proposed
by WP-8): a $\Ga$-module map between induced modules induces the map on the cokernel presentations of Proposition~\ref{prop:levelorbit}(2), so the
presentation-level fibre maps $\alpha^{0}$, $\omega$ rest on a statement of this document, while $\alpha_d$, $d>1$, is asserted on the model only.
It adds no numbered statement and moves no tag; it was written after WP-4's referee and is released \emph{unrefereed}.

\paragraph{Non-normative.}
\begin{itemize}\setlength{\itemsep}{0pt}
\item The header comment and the date line. The companion pointer stays at [E] v1.15.0: [E]'s tree text 1.15.1-pre.6 and [EX]'s
  1.2.0-pre.13 are \emph{not} released at this close; their releases (Q13: [E] 1.16.0, [EX] 1.2.0) wait for the owner's ruling on the
  \Espec{pr:latticesplit} edit (refuted as stated at $\ell=37$; \S\ref{sec:saturation}'s scope sentence, pre.2's B4).
\item Release note: \file{docs/implementation/slot15/WP-18\_close\_report.md}.
\end{itemize}

\subsection*{v1.15.0-pre.2 --- 2026-09-30 --- hecke-extensions slot 15, WP-3 resumed: the referee's repairs}

\noindent\emph{This version.} The repairs answering the Fable referee \file{docs/implementation/slot15/WP-4\_referee.md} (commit
\texttt{6215ccd}): \S\ref{sec:nebentypus} ACCEPT WITH REPAIRS; \cl{prop:saturation}, \cl{prop:satdefect}, \cl{prop:hidaalgconst},
\cl{rmk:thetaforms}, \cl{prop:cohforms} ACCEPT; \cl{cor:stablecert} ACCEPT WITH REPAIRS; \textbf{0 blocking}, every repair of its \S9 applied.
\textbf{No statement added or removed, no label renamed, no tag moved}; Appendices~\ref{app:status} and \ref{app:labelindex} unchanged.

\paragraph{Repairs.}
\begin{itemize}\setlength{\itemsep}{0pt}
\item \textbf{A1} \cl{prop:nebcorr}(4): the $\Ga_1(N)$/$\Ga_0(N)$ compatibility proved in-document (the coset bijection
  $A''_1\backslash\Ga_1(N)\to A''_0\backslash\Ga_0(N)$, $r_t\in\Ga_1(N)$ for both models); \STD{} only for the general $\langle d\rangle$-commutation.
\item \textbf{A2} \cl{rmk:nebheilbronn}: the parity clause. \textbf{A4} (NB4): the sharp form with the two constituents $\Sym^{i}$ and
  $\Sym^{\ell-1-i}\otimes\det^{i}$ (still \CONJ). \textbf{A5} \cl{lem:gammaonetable}(3): $L_{\Ga_1(N)}:=L^{\rhon}_{\Ga_1(N)}$. \textbf{A9}
  \cl{lem:teichmuller}: the Galois ring does not depend on the lift. \textbf{N1} \cl{rmk:nebcost}(1): one solve per distinct scalar.
\item \textbf{B1} \cl{lem:heckerankone} for every idempotent $e$ of $\Talg^{\mathbb L}$ (at most one Eisenstein system in any block);
  \cl{cor:stablecert} for a maximal ideal $\midl'$ of $\Talg'=\Z_\ell[T_p:p\in S]$, the code's joint system --- what slot-15 WP-5's certificate
  cites. \textbf{B4} \S\ref{sec:saturation}'s scope: faithfulness on $M_k(\Fl)$ is \Espec{cl:mkfaithful}; \Espec{pr:latticesplit} refuted as
  stated at $\ell=37$.
\end{itemize}

\paragraph{Non-normative.}
\begin{itemize}\setlength{\itemsep}{0pt}
\item Code maps: $H_{N,p}$ of \eqref{eq:heilbronn} \emph{built} (slot-15 WP-7, \file{Hecke/HeckeLevelNHeilbronn},
  \code{HeckeLevelNHeilbronnValidation}; \ref{LN3}'s Heilbronn half gains a \COMP\textsuperscript{a} sentence); \file{HeckeExt/HeckeHidaComponents}
  \emph{built} (slot-15 WP-6) in \S\ref{sec:codemap}'s \S\ref{sec:hidafull} row. The header comment and the date line.
\end{itemize}

\subsection*{v1.15.0-pre.1 --- 2026-09-30 --- hecke-extensions slot 15, WP-3: nebentypus, and the \texorpdfstring{$\Z/\ell^{K}$}{Z/l^K} Hecke algebra --- the pre-release of an \texorpdfstring{$X$}{X} bump}

\noindent\emph{This version.} Slot-15 WP-3 (plan \file{docs/implementation/slot15/00-slot15-plan.md} \S2 ``WP-3''; owner rulings of
2026-09-30, Q2: \emph{general $N$, characters into $\F_q^{\times}$, $q=\ell^{d}$ general, no layering}; Q3: the nebentypus code is slot~16's;
Q4: the certified \code{stable} verdict on the referee's ACCEPT). \textbf{A pre-release, not owner-released}; the Fable referee is
slot-15 WP-4. Two chapters added after \S\ref{sec:levelN}, before the appendices, so that no chapter number moves; no existing statement
changes, no label is renamed or removed, no tag moves.

\paragraph{New statements.}
\begin{itemize}\setlength{\itemsep}{0pt}
\item \textbf{Part~VII, \S\ref{sec:nebentypus}} \PNEW{} (contract first, rule~21; not built): \cl{def:nebdata} and \cl{lem:nebdata}
  ($\Ga_1(N)$, $\De_0(N)$, characters $\chi\colon(\Z/N)^{\times}\to\F_q^{\times}$, the module $\VR(\chi)$ with $\rho_{n,\chi}=\chi(a)\rhon$),
  \cl{lem:teichmuller} (the Galois ring $W_K(\F_q)$ on $C_{\ell,d}$, the Teichm\"uller lift $\tilde x^{q^{K-1}}$), \cl{rmk:nebfield},
  \cl{lem:inducedchar} (the character model is monomial: \eqref{eq:MHchi}), \cl{lem:gammaonetable}, \cl{def:nebmodel},
  \cl{prop:nebdecomp} (diamonds as fibre maps; the character model as the $\chi$-eigenspace; (NB2), (NB3)), \cl{prop:nebcorr}
  (\eqref{eq:nebcorr}: $T_p$, $U_p$, $T(p,p)=\chi(p)p^{n}$ on the level-$1$ $\tau$-table), \cl{thm:nebrelations}, \cl{prop:neborbit}
  (orbit-wise, parity, cokernel, Euler characteristic), \cl{rmk:nebheilbronn}, \cl{rmk:nebcost}, the refusal box, \cl{rmk:nebcheckable}
  ((NB1)--(NB6)).
\item \textbf{Part~VII, \S\ref{sec:saturation}} \PNEW{}: \cl{def:satdata}, \cl{prop:saturation} (P15 $=$ P-WP6-1: $\mathsf A_K=\Talg/(\Talg\cap
  \ell^{K}\Talg^{\mathrm{sat}})$, the stopping-rule equivalence), \cl{prop:satdefect} (P17: $\varepsilon^{\mathrm{sat}}_K=\mathrm{length}((\Talg^{\mathrm{sat}}/\Talg)
  [\ell^{K}])$), \cl{lem:heckerankone}, \cl{lem:freepart} (P-WP6-2), \cl{cor:stablecert} (the certified reading and the generation caveat),
  \cl{prop:cohforms} (the torsion-aware comparison, plan WP-19~(c) as corrected in \texttt{c1b45cd}), \cl{prop:hidaalgconst} (P18),
  \cl{rmk:thetaforms} (P20).
\end{itemize}

\paragraph{Non-normative.}
\begin{itemize}\setlength{\itemsep}{0pt}
\item \ref{thm:hidaweight}: one pointer sentence to \cl{prop:hidaalgconst} (P18); no tag moves. \S\ref{sec:levelN}'s scope paragraph: a
  pointer to \S\ref{sec:nebentypus}.
\item \S\ref{sec:codemap}: two rows; \S\ref{sec:notimplemented}: one box. Appendix~\ref{app:status}: $284\to307$ lines;
  Appendix~\ref{app:labelindex}: the new labels; the provenance counts: \PNEW{} $203\to226$.
\item \file{gl2\_refs.bib}: \code{cohen\_oesterle}, \code{ribet77} added (not verified). The header comment and the date line. The
  companion pointer stays at [E] v1.15.0.
\end{itemize}


\subsection*{v1.14.2 --- 2026-09-29 --- hecke-extensions slot 14: \S\ref{sec:levelN} built (WP-12)}

\noindent\emph{This version.} A $Z$-bump by the owner's ruling \textbf{Q-F} (2026-09-29): slot-14 WP-12 (commit \texttt{7a5158d}, level $N$
on the model by Shapiro, first build) is in the tree, refereed by the Fable code review of commit \texttt{64e6b9c}
(\file{docs/implementation/slot14/WP-12\_code\_review.md} \S4, whose texts are used verbatim; verdicts \S5.1--\S5.3 accept with
change, \S5.4 accept). \textbf{No statement changes, no label is added, renamed or removed, and no tag moves}; Appendix~\ref{app:status}
and Appendix~\ref{app:labelindex} are unchanged.

\paragraph{Non-normative.}
\begin{itemize}\setlength{\itemsep}{0pt}
\item \S\ref{sec:levelN}'s code map, \S\ref{sec:codemap}'s row: ``not built'' $\to$ \emph{built at v1.14.2}, naming \file{Hecke/HeckeLevelN}
  (\code{Gamma0LevelN}, \code{Gamma0SubTransversalN}, \code{InducedRelationsLevel}, \code{OrbitWiseFixedSpace},
  \code{CokernelFromFixedSpaces}), \emph{not} \file{Hecke/InducedLocalSystem} (review S-4). Not built: $H_{N,p}$, the sparse
  cokernel elimination, degeneracy maps at composite $N$.
\item \ref{levelO3} \textbf{closed}; \ref{levelO2} stays open (its $\tau$-route at composite $N$ is built). The two consequential sites
  of \S\ref{sec:hidafull} (code map and row; the scope paragraph) gain their v1.14.2 sentence.
\item \ref{LN3}: the $\tau$-route half's \COMP\textsuperscript{a} sentence ($20$ cells). \ref{LN6}: the second-route
  \COMP\textsuperscript{a} sentence with the $q\le7$ and $e$-read-off caveats and ``the Loewy length is the subalgebra's''.
\item The $r_t$ alternative (slot~13's \code{Gamma0SubTransversal}) is stated in the box after Lemma~\ref{lem:cosetrepsN}, not in
  Definition~\ref{def:cosetrepsN}(6). Remark~\ref{rmk:levelcost}(3): $\nu_2$, $\nu_3$ checked against the table (one parenthesis).
\item The header comment and the date line. The companion pointer stays at [E] v1.15.0.
\end{itemize}


\subsection*{v1.14.1 --- 2026-09-29 --- hecke-extensions slot 14: the second cluster run's data sentences, and the \texorpdfstring{$U_\ell$}{U\_l} lever's rank lemma}

\noindent\emph{This version.} The owner's rulings of 2026-09-29 on that day's Fable audit, the file
\file{AUDIT\_contract\_currency\_2026-09-29.md} in the directory \file{docs/implementation/slot14/}: \textbf{Q-B} (the second-cluster-run data, job \code{3467349}, signed off; proposals
L1--L4 of \file{WP-L1\_binding\_\allowbreak doc\_proposals.md} in the same directory, with the audit's
corrections: the probe's peaks and the sampler's stated apart, and the admission predictor's refit named) and \textbf{Q-H} (audit item
C4, the lever lemma, into this document). A $Y$-level release by the owner's ruling: one new statement, proved, about the
implementation's own matrices; \textbf{no existing statement changes, no label is renamed or removed, and no tag moves}. \textbf{Q-F}:
slot-14 WP-12's ``not built $\to$ built'' flip of \S\ref{sec:levelN}'s three code-map sites is \emph{not} made here (WP-12 is in review);
it will be a later, separate bump.

\paragraph{New statement.}
\begin{itemize}\setlength{\itemsep}{0pt}
\item \cl{lem:rhorank} \PNEW{} --- over $\Z/\ell^{k}$, $\rhon(g)$ with $\ell\mid\det g$ has at most $\min(k,n+1)$ non-zero invariant factors
  (Smith form over $\Z_\ell$, $\rhon$ of a diagonal matrix); over $\Fl$ it is $w\,v^{\mathsf T}$ in closed form. The claim
  \file{Hecke/HeckeLowRank.h} \code{ENGINE\_CHECK}s at every block (slot-13 WP-4; slot-14 G0's closed form). Placed after
  \S\ref{sec:hidafull}'s code map; Appendix~\ref{app:status} gains its line (counts $283\to284$), Appendix~\ref{app:labelindex} its label.
\end{itemize}

\paragraph{Non-normative (Q-B: L1--L4).}
\begin{itemize}\setlength{\itemsep}{0pt}
\item \cl{rmk:hidacheckable2}: \ref{Hi11} gains a \COMP\textsuperscript{a} sentence (L1: held at $n=4150$, $711.8$~s, $49.9$~GiB, the probe's
  gated \code{HI11} line; $n=3460$ refused by name); \ref{Hi12} gains one (L2: held at $(4830,2411)$, $68.4$~GiB probe / $65.3$~GiB sampler;
  $(2420,2411)$ recomputed with the levers in $19.2$~GiB; $(1210,2411)$ reproduced) and its last clause is repaired (the build
  \texttt{e85e25a} prints a (Hi12) verdict). No item's tag moves.
\item \S\ref{sec:hidafull}'s code map (L3): the second run's provenance line, the refit of the campaign's admission predictor
  (\texttt{549200e}), and the warning that \ref{levelO5}'s fit predates the slot-14 levers; the lever sentence points to \cl{lem:rhorank}.
\item \S\ref{sec:codemap}'s v1.13.0 scope box (L4): ``at v1.14.1 a second cluster run added \ref{Hi11} at $n=4150$ and \ref{Hi12} at
  $(4830,2411)$, both held''.
\item The header comment and the date line. The companion pointer stays at [E] v1.15.0.
\end{itemize}


\subsection*{v1.14.0 --- 2026-09-28 --- hecke-extensions slot 14, WP-13: the owner-signed release, from v1.14.0-pre.2}

\noindent\emph{This version.} The owner releases the slot's $X$ bump, whose text v1.14.0-pre.1 and pre.2 carried: \S\ref{sec:levelN}
(slot-14 WP-4). Its referee (slot-14 WP-5, \file{docs/implementation/slot14/WP-5\_referee.md}) was \textbf{ACCEPT WITH
REPAIRS}, with $0$ blocking, $3$ should-fix and $8$ minor findings, all applied at pre.2. The owner's decision of 2026-09-28, ``Go for publication and
release'', answers [S14] Q2 and Q3 with the recommended options. \S\ref{sec:levelN} is released \emph{unbuilt} (contract first,
rule~21). Its code maps say ``not built at v1.14.0''. \textbf{No statement is added, no label is renamed or removed, and no tag
moves}; Appendix~\ref{app:status} is unchanged. Release note: \file{docs/implementation/slot14/WP-13\_release\_note.md}.

\paragraph{Non-normative (the [S14] D1 routing, Q3).}
\begin{itemize}\setlength{\itemsep}{0pt}
\item \textbf{P1} (\ref{Hi12} run once at $(2420,2411)$ and held), \textbf{P2} and \textbf{P8} (\ref{Hi11} at $n=1390$, $2080$, $2770$, and
  $(700,691)$ at $m=2$ on a second machine) and \textbf{P3} (Remark~\ref{rmk:bocksteinexp}(3), the $\Z/\ell^{2}$ depth window, $17$ cells,
  exactly the loci of (2)) are re-inserted \emph{verbatim} from pre.1 (commit \texttt{cc97dd7}). pre.2 had taken them out under Q2's
  default, which no longer holds. Only their provenance clauses change: ``staged, publication pending'' becomes ``published 2026-09-28 on
  the owner's sign-off'', because the extract \file{docs/matrix-engine/results/noctua2\_run1/} is now a published set. Each remains a
  \COMP\textsuperscript{a} sentence. No item's tag moves, and the converse of Remark~\ref{rmk:bocksteinexp}(2) is still not asserted.
\item \textbf{P7} (the $493.3$~GiB provenance) is restored to its pre.1 form in \S\ref{sec:hidafull}'s code map. The
  \S\ref{sec:codemap} row and v1.13.0's scope box record that \ref{Hi12} ran and held.
\item \emph{Deferred to slot~15}, as new statements that the rule for post-referee proposals excludes from a release: P14 (the
  ordinary part of section~M as a \COMP{} sentence; its cluster data are also unpublished), P15 (= slot-14 WP-6's P-WP6-1 and
  P-WP6-2, two propositions), P17 (the saturation identity for $\varepsilon_K$ and the generation caveat) and P18 (an attribution of
  algebra-level consequences to Theorem~\ref{thm:hidacontrol} and Proposition~\ref{prop:stabilisation}). The release note gives the reasons.
\item The header comment, the date line, the companion pointer ($\to$ [E] v1.15.0), and the two code maps of \S\ref{sec:levelN}.
\end{itemize}


\subsection*{v1.14.0-pre.2 --- 2026-09-28 --- hecke-extensions slot 14, WP-5: the referee's repairs}

\noindent\emph{This version.} The repairs answering the Fable referee of v1.14.0-pre.1
(\file{docs/implementation/slot14/WP-5\_referee.md}: \textbf{ACCEPT WITH REPAIRS}, $0$ blocking, $3$ should-fix, $8$ minor; every
proof of \S\ref{sec:levelN} correct as written). All eleven findings are applied (F10 is a check, no change). \textbf{No statement
is added, no label is renamed or removed, and no tag moves}; no \PROVED{} statement's text changes. Still a pre-release ([S14] Q3).
Report: \file{docs/implementation/slot14/WP-4\_report.md}, section ``pre.2 repairs''.

\paragraph{Normative.}
\begin{enumerate}\setlength{\itemsep}{0pt}
\item \textbf{F1} \cl{rmk:heilbronncost}(3): a Hecke \emph{matrix} on $H^{1}$ costs $O(\abs{\mathrm{Mer}_p}\cdot d\cdot(\dim H^{1})^{2})$
  through the $D_N\times\dim H^{1}$ normal-form table of $\mathrm{Rel}_N$ (or one sparse solve per column) --- quadratic in $m_N$;
  obstacle \ref{levelO1} loses the $N_q^{2}$ factor and the dense projector, not the quadratic order. (4) names the normal form as the
  route's one $D_N\times\dim H^{1}$ object, (2)'s ``no $D_N\times D_N$ object'' being about one application; (5) says ``quadratic per
  matrix''.
\item \textbf{F2} \ref{LN4}, \ref{LN6}: the scan's $W_{N,p}$ is joint for $T_2$, $T_3$, $T_5$, $T_7$ only and contains \ref{LN4}'s, equality
  being part of the prediction; $e$ is read off $\dim W_{N,p}$, so the cells are tabulated by $\dim W_{N,p}\in\{3,5,7\}$ and $e=1,2,3$
  only under \ref{LN4}.
\item \textbf{F3} \cl{def:levelellmodel}: the composite-$N$ table is fixed by \cl{def:cosetrepsN}.
\item \textbf{F4} \cl{def:cosetrepsN}(3): the case $N=1$ ($(c_1,d_1)=(0,1)$, $w_1=I$).
\end{enumerate}

\paragraph{Non-normative.}
\begin{itemize}\setlength{\itemsep}{0pt}
\item \textbf{F5} the contractbox after \cl{lem:cosetrepsN}: case (iii) is a bounded loop plus one inverse, table-building time.
  \textbf{F6} the notation paragraph lists $\mathrm{Mer}_n$, $W_A$, $W_{N,p}$, $\mu_N$, $\zeta_A$, $H_{N,p}$, $\mathrm{Div}$, $D_0$ as new
  symbols (rule~\ref{A5}), with $W_A\ne$ \cl{def:rhon}'s $W$ and $\mu_N\ne\mu_{\bullet}$. \textbf{F7} \cl{lem:levelduality}(5): the Brown
  pointer is definitional. \textbf{F8} \ref{LN3}: the independent route tests \cl{prop:heilbronn}(2), the $\tau$-route half is WP-12's;
  the referee's seven re-run cells added. \textbf{F9} \cl{rmk:heilbronncost}(1): $\sum_Lx_0(L)=p^{2}+1$ spelled out,
  $\abs{\mathrm{Mer}_p}/(p\log p)$ increasing. \textbf{F10} (\cl{rmk:heilbronnconv}(2)) verified by the referee, no change.
\item \textbf{F11}: [S14] Q2 is unanswered, so its default (``P1--P4, P8 left out of the binding drafts'') stands. P1 (\ref{Hi12}), P2
  and P8 (\ref{Hi11}) and P3 (\cl{rmk:bocksteinexp}(3)) are taken back out, with their traces in \S\ref{sec:codemap}'s
  \S\ref{sec:hidafull} row and \S\ref{sec:hidafull}'s code map; P7 (the $493.3$~GiB provenance line) is kept, reworded to stand alone.
  Re-insertion notes: \file{docs/implementation/slot14/WP-D1\_binding\_doc\_proposals.md}. No label was cited by the removed text alone.
\item The header comment, the date line, the two code maps of \S\ref{sec:levelN} (``not built through v1.14.0-pre.2''); v1.14.0-pre.1's
  commit hash filled in below.
\end{itemize}


\subsection*{v1.14.0-pre.1 --- 2026-09-28 --- hecke-extensions slot 14, WP-4: level \texorpdfstring{$N$}{N}, part 2 --- the pre-release of an \texorpdfstring{$X$}{X} bump}

\noindent\emph{This version.} The slot's ``contract first'' package for level $N$ (\file{docs/implementation/slot14/00-slot14-plan.md}
\S2 ``WP-4'', rules~20 and~21; the owner's questions of its \S9 at their defaults, recorded as lead rulings in [X] v1.18.2 --- in
particular Q3: drafts only, no release; Q6: literature locators unverified). The Fable referee (slot-14 WP-5) follows and its
repairs are v1.14.0-pre.2. \textbf{Twelve statements are added, \PNEW{}; no label is renamed or removed and no existing tag
moves}; Appendix~\ref{app:status} gains twelve lines. Base \texttt{9269d25}. Report:
\file{docs/implementation/slot14/WP-4\_report.md}. Commit: \texttt{cc97dd7}.

\paragraph{Normative.}
\begin{enumerate}\setlength{\itemsep}{0pt}
\item \textbf{Part~VII, \S\ref{sec:levelN}} \PNEW{}: \cl{def:cosetrepsN}, \cl{lem:cosetrepsN} (obstacle O3: the binding coset table
  for every $N$), \cl{def:levelNmodel}, \cl{lem:treeres}, \cl{lem:manin}, \cl{def:kappaN} (the isomorphism of
  \cl{prop:levelorbit}(2), pinned), \cl{lem:levelduality}, \cl{lem:merelpath}, \cl{prop:heilbronn} (obstacle O2: $T_p$, $p\nmid N$, and
  $U_p$, $p\mid N$, on the cokernel presentation, proved equal to \cl{prop:levelcorrespondence}(4)(d) on $H^{1}$),
  \cl{rmk:heilbronnconv}, \cl{rmk:heilbronncost}, \cl{rmk:levelNcheckable} (the predictions (LN1)--(LN6), committed before
  slot-14 WP-12's code).
\item \textbf{Where this departs from the plan's statements.} (a) The Heilbronn--Merel set enters \emph{through the adjugate}: the
  value on block $j\mapsto j\cdot\epsilon$ is $\rhon(\adjg h_j(\epsilon))$, and $\mathrm{Mer}_p$ itself does not satisfy the right-coset
  condition the proof needs (\cl{rmk:heilbronnconv}(2), with two mis-transcriptions shown to fail). (b) $U_p$, $p\mid N$, is not left
  open: the same formula with the non-primitive terms dropped. (c) The plan's ``$\abs{X_p}=O(p\log p)$'' is not what Merel's set
  gives: $\le p^{2}+1$ proved, about $0.4\,p\log^{2}p$ counted; the $O(p\log p)$ Heilbronn set from continued fractions is named as
  an admissible replacement, not specified (\cl{rmk:heilbronncost}(1)). (d) (N6) is sharpened to mod-$p$ Loewy length $e+1$ on the
  Eisenstein part of the full $H^{1}$ (the cuspidal part $e$, the Eisenstein class one more), from the scan (\ref{LN6}). (e) The
  prediction labels are (LN1)--(LN6), not (N1)--(N6): \texttt{N3} is a notation rule of \S\ref{sec:notation}.
\item \cl{rmk:levelscaling}: \ref{levelO2}, \ref{levelO3} gain a status sentence (specified and proved, resp.\ pinned; not built).
\end{enumerate}

\paragraph{Non-normative.}
\begin{itemize}\setlength{\itemsep}{0pt}
\item \textbf{Data remarks} of \file{docs/implementation/slot14/WP-D1\_binding\_doc\_proposals.md}, each marked as the first
  cluster run's, its publication pending the owner's sign-off ([S14] Q2): P1 (\ref{Hi12} held at $(2420,2411)$), P2 and P8
  (\ref{Hi11} at $n=1390$, $2080$, $2770$; $(700,691)$ at $m=2$ on a second machine) as \COMP\textsuperscript{a} sentences of the
  items; P3 (the $\Z/\ell^{2}$ depth window, $17$ cells, exactly the loci of \cl{rmk:bocksteinexp}(2)) in its clause~(3); P7 (the
  $493.3$~GiB of $(2420,2411)$, $0.912$ of the lever fit) in \S\ref{sec:hidafull}'s code map rather than its Scope paragraph (a
  provenance line, not a scope statement). No tag set moves.
\item Bibliography: \cite{merel96}, \cite{wwe} added (not verified); the header comment, the date line, the abstract, ``How to read
  this document'' (Part~VII), \S\ref{sec:codemap} (a row), \S\ref{sec:notimplemented} (a v1.14.0-pre.1 box), \S\ref{sec:hidafull}'s
  Scope and code map, Appendix~\ref{app:status} (twelve lines, the counts), Appendix~\ref{app:labelindex}, this appendix's
  provenance table and roster. The \textbf{companion pointer moves to [E] v1.14.0} (released at the slot-13 close after this
  document's v1.13.0, which left the pointer at v1.13.1).
\end{itemize}

\paragraph{Not done here, deliberately.}
\begin{enumerate}\setlength{\itemsep}{0pt}
\item \textbf{No code, no build}: WP-12 (gated on WP-5's verdict and the owner's Q4) builds the table, the orbit-wise route and
  $H_{N,p}$; the sparse elimination at large $N$ is slot~15's.
\item No literature locator is verified ([S14] Q6's default); copies of \cite{manin}, \cite{merel}, \cite{mazur} and \cite{cremona} exist
  in the owner's read-only library (papers/) and were \emph{not} read, the ruling being
  unverified.
\item No release: [C] stays at v1.13.0 as released; this text is a pre-release for the owner (Q3).
\end{enumerate}


\subsection*{v1.13.0 --- 2026-09-27 --- hecke-extensions slot 13, the close (WP-15): the owner-signed release, from v1.13.0-pre.2}

\noindent\emph{This version.} The owner's release of the slot's $X$ bump whose text v1.13.0-pre.1 and -pre.2 carried:
\S\ref{sec:hidafull} (slot-13 WP-2, refereed by WP-3, \file{docs/implementation/slot13/WP-3\_referee.md}: \textbf{ACCEPT WITH
REPAIRS}, $0$ blocking, $2$ major, $10$ minor, all applied, none outstanding). Owner decision Q14 (2026-09-26) made the release
conditional on the referee's verdict being ACCEPT with $0$ blocking and on slot-13 WP-6's arms holding; both hold
(\file{docs/implementation/slot13/WP-6\_report.md}: \eqref{eq:Ulevell} held against the independent group-cocycle route,
Proposition~\ref{prop:oldspace}(1), (2) as cochain maps, Proposition~\ref{prop:stabilisation} modulo $\ell$, $\ell^{2}$, $\ell^{3}$,
\ref{Hi10} on $92$ cells). Proposition~\ref{prop:stabilisation} is released with its per-item tags as at pre.2. \textbf{No statement
is added, no label is renamed or removed, and no tag moves.}

\paragraph{Normative.}
\begin{enumerate}\setlength{\itemsep}{0pt}
\item \cl{def:hidalattice}, variant (H): the sentence ``nothing else changes'' is amended, as the slot-13 WP-13 code audit (A3)
  required: on the module route the (free) cochain module's Howell form may carry a non-unit pivot on Hida's lattice, and the route
  then takes a Nakayama basis, checked free by length (slot-13 WP-5); value-preserving wherever the route computed before.
\end{enumerate}

\paragraph{Non-normative.} The code maps of \S\ref{sec:hidafull} and \S\ref{sec:codemap}, and \S\ref{sec:notimplemented}'s box, read
``built'' (\file{Hecke/HeckeLevelEll}, \file{HeckeExt/HeckeHidaTower} part~2, \code{HeckeHidaHeavyCellProbe}); \ref{Hi7}--\ref{Hi11}
held, \ref{Hi12} not run. Companion pointer unchanged ([E] v1.13.1; its v1.14.0-pre.1 \ldots{} pre.3 are slot-13 pre-releases).
Report \file{docs/implementation/slot13/WP-15\_report.md}.


\subsection*{v1.13.0-pre.2 --- 2026-09-26 --- hecke-extensions slot 13, WP-3: the referee's repairs}

\noindent\emph{This version.} The repairs answering the Fable referee of v1.13.0-pre.1
(\file{docs/implementation/slot13/WP-3\_referee.md}: \textbf{ACCEPT WITH REPAIRS}, $0$ blocking, $2$ major, $10$ minor);
\textbf{all twelve findings are applied}, none declined; the finding-by-finding table with locations is
\file{docs/implementation/slot13/WP-3\_response.md}. \textbf{One statement is added, \PNEW{}}
(\cl{rmk:levelscaling}); no label is renamed or removed; two tag sets move (\cl{rmk:hidacheckable2} gains \BSTD{};
\cl{prop:stabilisation} is tagged per item, its tag set unchanged); Appendix~\ref{app:status} gains one line. No code.
Base \texttt{6be7502}. Commit: \texttt{[hash]}.

\paragraph{Normative.}
\begin{enumerate}\setlength{\itemsep}{0pt}
\item \textbf{M1 --- Hida's lattice.} \cl{rmk:stabconvention}(3) replaces pre.1's false ``up to the order of the basis'': Hida's
  $L_n$ carries $\rho^{\mathrm H}_n(A)=\rhon(A^{T})^{T}=D^{-1}\rhon(A)D$, $D=\mathrm{diag}(\binom nj)$, a different $\Z_\ell$-lattice
  from $\Vcoef$ unless $n+1=c\ell^{s}$, $1\le c\le\ell-1$ (so always for $n<\ell$); what is lattice-independent (ranks,
  characteristic polynomials) and what is not (the finite-level module class). Glosses in \cl{thm:hidacontrol},
  \cl{thm:hidaweight}; \cl{def:hidalattice} names its lattice and gains the variant (H) on Hida's lattice (one
  \code{FibreRepresentation} variant); \ref{Hi8} gains arms (a) $\Vcoef$, (b) $L_n$, (c) the $D_\ast$ comparison (reported), and its
  source is restated: (b) extends Hida's parabolic control to the full $H^{1}$, (a) is this document's extrapolation to a
  different lattice. \ref{Hi8} stays \CONJ{}.
\item \textbf{M2 --- the cost model.} \cl{rmk:levelcost}(1), (5): a Hecke \emph{matrix} on $H^{1}$ on the built dense route costs
  the $\Theta((N_qD_N)^{2})$ projector per column, cubic in $m_N(k-1)$; the count of (5) is for the cochain map. New
  \cl{rmk:levelscaling}: the two regimes (large $N$ at small $k$; large $k$ at level $1$ or $\ell$), the obstacles
  \ref{levelO1}--\ref{levelO6} (the owner's scaling requirement, tracked here), and the Hecke action on the cokernel presentation
  named as the Manin--Merel Heilbronn-matrix formula, \textbf{specified by pointer, not transcribed, not built}
  (\cite{manin}, \cite[\S2]{merel}, \cite{stein}, \cite{cremona}; not verified).
\item \textbf{Minors.} m1: \cl{prop:stabilisation} per-item tags, and its proof of (4) cites \cl{rmk:stabconvention}(1). m2:
  \cl{def:levelellmodel} (the count for composite $N$), \cl{prop:levelcorrespondence}(4)(c) (prime-power $d$; composite $d$ by
  composition). m3: \cl{rmk:levelcost}(3), the exact $\nu_2$, $\nu_3$. m4: \ref{Hi8}'s test is a decision procedure, no
  ``inconclusive''. m5: \cl{thm:hidacontrol} records Hida's $p\ge5$ (page image; the text layer says $p>5$) and $N=1$. m7: the
  Scope paragraph at $\ell\in\{2,3\}$ (cochain level and the Bredon reading hold). m8: \cl{rmk:stabconvention}(4), the
  $U^{\ast}$ old-space matrix and $\omega\,e(U^{\ast})\,\mathrm{dg}_1\in\ell\cdot\Z_\ell^{\times}$. m9: \ref{Hi10}'s commutativity \BSTD{}.
  m6, m10: \cl{prop:oldspace}'s \COMP\textsuperscript{a} scope (cohomology only; the block formula untested until WP-6).
\end{enumerate}

\paragraph{Non-normative.}
\begin{itemize}\setlength{\itemsep}{0pt}
\item The chapter's code map, \S\ref{sec:codemap}'s row and \S\ref{sec:notimplemented}'s box (the variant (H), \ref{levelO2}); the
  header comment, the date line, the companion pointer's wording (unchanged: [E] v1.13.1); Appendix~\ref{app:status} (the
  counts), Appendix~\ref{app:labelindex}, the provenance table and roster (and the stale ``$106$ assertions'' left by pre.1,
  now $116$); pre.1's commit back-filled. \cl{prop:stabilisation}'s \COMP\textsuperscript{a} gains the referee's $\ell=13$ values.
\end{itemize}

\paragraph{What WP-5 and WP-6 must test differently.} WP-5: \ref{Hi8} arms (a) and (b) (the latter on the $\rho^{\mathrm H}_n$ fibre),
(c) reported, by the decision procedure. WP-6: the pair $(e_\ell\circ\mathrm{dg}_\ell,\ c_{\mathrm{st}}^{-1}\omega)$ and the old-space
relations, never a trace inverse in the $U^{\ast}$ convention (\cl{rmk:stabconvention}(4)); \eqref{eq:Ulevell} is evaluated on the
model for the first time there.

\subsection*{v1.13.0-pre.1 --- 2026-09-26 --- hecke-extensions slot 13, WP-2: control, level \texorpdfstring{$\ell$}{l}, the lattice --- the pre-release of an \texorpdfstring{$X$}{X} bump}

\noindent\emph{This version.} The slot's ``contract first'' package (\file{docs/implementation/slot13/00-slot13-plan.md}
\S2 ``WP-2'' and \S9, the owner decisions of 2026-09-26, rules~20 and~21). It carries its own identifier so that
no version number is published twice under different text (\ref{V1}, \ref{V2}); the Fable referee (slot-13 WP-3)
follows and its repairs are v1.13.0-pre.2. \textbf{Thirteen statements are added, \PNEW{}; no label is renamed or
removed and no existing tag moves}; Appendix~\ref{app:status} gains thirteen lines. Base \texttt{26f3f45}.
Report: \file{docs/implementation/slot13/WP-2\_report.md}. Commit: \texttt{83376c5} (back-filled at pre.2, \ref{CL6}).

\paragraph{Normative.}
\begin{enumerate}\setlength{\itemsep}{0pt}
\item \textbf{Part~VII, \S\ref{sec:hidafull}} \PNEW{}: \cl{def:levelellmodel}, \cl{prop:levelcorrespondence},
  \cl{prop:levelellU}, \cl{prop:oldspace}, \cl{prop:levelorbit}, \cl{rmk:levelcost}, \cl{lem:levelelltorsion},
  \cl{prop:stabilisation}, \cl{rmk:stabconvention}, \cl{thm:hidacontrol}, \cl{thm:hidaweight}, \cl{def:hidalattice},
  \cl{rmk:hidacheckable2} (the predictions (Hi7)--(Hi12), committed before slot-13 WP-5, WP-6 and WP-10's code).
\item \textbf{Where this departs from the plan's statements.} (a) The level-$\ell$ operator needs \emph{no} new coset or
  $\tau$-table: \cl{prop:levelellU} reads the level-$1$ $U_\ell$ table, so the plan's fallback (``degeneracy-map
  comparison only'') is not needed; the general form \cl{prop:levelcorrespondence} covers every level-$N$ Hecke
  operator, degeneracy and level-change map (the owner's level-$\Ga_0(N)$ amendment). (b) The plan's
  ``$e_{\mathrm{ord}}H(\Ga_0(\ell))\cong e_{\mathrm{ord}}H(\Ga)$'' is made precise: in this document's convention the
  integral isomorphism is $e_\ell\circ\mathrm{dg}_\ell$, and $e_\ell\circ\mathrm{dg}_1$ is divisible by $\ell$
  (\cl{prop:stabilisation}(1),(3)); Hida's ``$e\circ\res$'' is the same statement for the Atkin--Lehner conjugate
  $U^{\ast}$ (\cl{rmk:stabconvention}). (c) The plan's old-space matrix is $\mathsf M=\bigl(\begin{smallmatrix}0&-\ell^{n}\\
  \ell&T_\ell\end{smallmatrix}\bigr)$ in the basis $(\mathrm{dg}_1,\mathrm{dg}_\ell)$, with two of its relations exact at
  cochain level (\cl{prop:oldspace}). (d) The plan's labels \texttt{thm:hidacontrol} / \texttt{thm:hidaweight} are
  literature (\BSTD{}), with locators \textbf{verified} against the owner's read-only PDFs (the Q2 amendment); the
  book \cite{hida_book} is not among them and stays unverified. (e) The plan's $\alpha_1$, $\alpha_\ell$ are written
  $\mathrm{dg}_1$, $\mathrm{dg}_\ell$ (rule~\ref{A5}: $\alpha$ is \cl{prop:alpha}'s). (f) (Hi10) is \PROVED{} as the plan
  says, but its $k=2$ half rests on \BSTD{} pointers not verified (Shimura, Atkin--Lehner).
\end{enumerate}

\paragraph{Non-normative.}
\begin{itemize}\setlength{\itemsep}{0pt}
\item \cl{rmk:hidacheckable}'s \ref{Hi3}, \ref{Hi4}: a v1.13.0-pre.1 locator note each; tags unchanged.
\item Bibliography: \cite{atkin_lehner} added (not verified); the notes of \cite{hida_ens} and \cite{hida_invent} record
  the verified locators. Hida's papers are read from the owner's local library (papers/h/Hida/), read-only;
  nothing is copied into \file{docs/literature/} (the owner's Q2 amendment).
\item The abstract, ``How to read this document'' (Part~VII), \S\ref{sec:codemap} (a row),
  \S\ref{sec:notimplemented} (a v1.13.0-pre.1 box), Appendix~\ref{app:status} (thirteen lines, the counts),
  Appendix~\ref{app:labelindex}, this appendix's provenance table and roster, the notation table of
  \S\ref{sec:notation-dickson}, the header comment and the date line. The companion pointer is unchanged
  ([E] v1.13.1).
\end{itemize}

\paragraph{Not done here, deliberately.}
\begin{enumerate}\setlength{\itemsep}{0pt}
\item \textbf{No Hida theorem is proved}: \ref{Hi7}--\ref{Hi9}, \ref{Hi11}, \ref{Hi12} are \CONJ{}; \ref{Hi3}, \ref{Hi4}
  keep \CONJ{} although their locators are now verified (a candidate for the release, the owner's call).
\item \textbf{No code}: the chapter's code map says what slot-13 WP-5, WP-6 and WP-10 are to build; the monomial and
  orbit-wise routes of \S\ref{sec:levelcost} are an interface for general $N$.
\item No Hecke action on the cokernel presentation (\cl{prop:levelorbit}(2)); no level $\ell^{r}$, $r\ge2$, beyond the
  general formula.
\end{enumerate}


\subsection*{v1.12.1 --- 2026-09-26 --- hecke-extensions slot 12, post-close, S12-census2: the code maps catch up}

\noindent\emph{This version.} A \textbf{non-normative} $Y$ bump (\ref{V2}): code maps and the companion pointer
only. \textbf{No statement, hypothesis, label or tag changes.} Base \texttt{bc96fae} (the finality landing
S12-FIN, [E] v1.13.1). Report: \file{docs/implementation/slot12/census2\_report.md}. Commit: \texttt{[hash]}.

\paragraph{Non-normative.}
\begin{itemize}\setlength{\itemsep}{0pt}
\item \S\ref{sec:hida}'s code map: ``Not built at v1.12.0-pre.1 \ldots\ What is built will be recorded here''
  is replaced by the record of what slot-12 WP-4 and WP-5 built --- \file{Hecke/HeckeOrdinaryPart}
  (\cl{prop:torsionprepass}, \cl{lem:ordexponent}; \code{HeckeOrdinaryPrecisionValidation}) and
  \file{HeckeExt/HeckeHidaTower} (\cl{def:hidatower}, \cl{Hi2}--\cl{Hi6}; \code{HeckeHidaTowerValidation},
  \code{HeckeHidaTowerSlowValidation}); \S\ref{sec:codemap}'s \S\ref{sec:hida} row likewise.
\item \S\ref{sec:codemap}'s \S\ref{sec:boundarydatum} and \S\ref{sec:dicksontheta} rows name
  \file{HeckeExt/HeckeFinalityRoutes} (slot-12 S12-FIN): (d1) ``boundary'' resting on
  \cl{prop:bdrytranscription} and \cl{lem:bdrytorus}, (d0) ``dickson'' on \cl{prop:heckeH0}(1) --- the
  certificates themselves are [E]'s (\texttt{cert:boundary}, \texttt{cert:dickson}, [E] v1.13.0--v1.13.1).
\item The companion pointer: [E] v1.12.1 $\to$ v1.13.1. The header comment and the date line.
\end{itemize}


\subsection*{v1.12.0 --- 2026-09-26 --- hecke-extensions slot 12, post-close: the owner-signed release, from v1.12.0-pre.3}

\noindent\emph{This version.} The owner's release of the slot's $X$ bump whose text v1.12.0-pre.1, -pre.2
and -pre.3 carried: \S\ref{sec:hida}, the foundation of Hida theory at level~$1$ (slot-12 WP-2, refereed by
WP-3, \file{docs/implementation/slot12/WP-3\_referee.md}: $0$ blocking, $0$ major, $10$ minor, all
applied), and the Shapiro transcription in degree~$1$ (Proposition~\ref{prop:bdrytranscription},
Lemma~\ref{lem:bdrytorus}; slot-12 WP-16, refereed by \file{docs/implementation/slot12/WP-16\_referee.md}:
all three results ACCEPT, status notes ACCEPT WITH MINOR REPAIRS, eight minor findings, all applied) ---
both accepted after repairs, all applied, none outstanding. \textbf{The one status change this version
makes}: Question~\ref{open:ordprecision} is answered at $\ell\ge5$, $n\ge1$ by
Corollary~\ref{cor:ordprecision} (with Proposition~\ref{prop:torsionprepass} and
Remark~\ref{rmk:ordtorsionfree}), so its tag moves \OPENSTAT{} $\to$ \PROVED{} and the item is retitled
``--- ANSWERED 2026-09-26'', exactly as \ref{open:ell2} lost its \code{[open]} tag at v1.3.0. \textbf{No
statement is added, no label is renamed or removed, and one tag moves}; the question text and the boxed
statement of \ref{open:ordprecision} are unchanged (rules~\ref{L1}, \ref{L2}). Base \texttt{7d71fdc}: every
slot-12 mathematics package (WP-2, WP-3, WP-16 and its referee repairs) merged. Report: this release note
(no separate WP report; owner sign-off 2026-09-26, on \file{docs/implementation/slot12/WP-2\_report.md},
\file{WP-3\_referee.md}, \file{WP-16\_report.md} and \file{WP-16\_referee.md}). Commit: \texttt{[hash]}.

\paragraph{Normative.}
\begin{enumerate}\setlength{\itemsep}{0pt}
\item \cl{open:ordprecision}: retitled ``Certified precision for the ordinary part --- ANSWERED
  2026-09-26''; the tag \OPENSTAT{} is replaced by \PROVED{}, discharged by \cl{cor:ordprecision} at
  $\ell\ge5$, $n\ge1$ ($e=e'$ for the ordinary part, $K=e'+e_{\max}$ for the free quotient, $e_{\max}$
  computable before the run from $\rhon(S)$, $\rhon(T)$ alone). At $\ell\in\{2,3\}$ nothing is asserted;
  the question text and the boxed statement are unchanged.
\end{enumerate}

\paragraph{Non-normative.}
\begin{itemize}\setlength{\itemsep}{0pt}
\item Appendix~\ref{app:status}: the \code{open} row's status column (\PROVED{} in place of \OPENSTAT{})
  and the counts ($\PROVED{}$ $129\to130$, $\OPENSTAT{}$ $7\to6$; tagged/untagged unchanged).
\item The abstract's parenthesis on \ref{open:ordprecision}, the companion-pointer sentence (unchanged
  version, [E] v1.12.1), the header comment and the date line.
\item \textbf{Outside this document}: \file{CLAUDE.md}'s and \file{Hecke/README.md}'s mentions of this
  document's version move to v1.12.0; \file{docs/implementation/hecke\_extensions\_implementation\_plan.md}
  [X]'s deferred list marks ``[C] v1.12.0 release'' done (the release and the status change; the
  \S\ref{sec:hida} code map's own ``what is built will be recorded here'' is untouched by this release ---
  it is slot-12 WP-4/WP-5's code, not asked for here).
\end{itemize}

\paragraph{Not done here, deliberately.}
\begin{enumerate}\setlength{\itemsep}{0pt}
\item \textbf{No mathematics changes}: every normative statement, hypothesis, displayed formula and label
  of v1.12.0-pre.1--pre.3 is carried verbatim; this version only lifts the pre-release status and answers
  \ref{open:ordprecision}.
\item \textbf{No finality certificate}: (d0)/(d1) are [E] matters and land with their code, as
  v1.12.0-pre.3 recorded.
\item \textbf{No code-map update}: \S\ref{sec:hida}'s ``Not built at v1.12.0-pre.1 \ldots\ What is built
  will be recorded here'' is left for the package that lands that record.
\end{enumerate}


\subsection*{v1.12.0-pre.3 --- 2026-09-26 --- hecke-extensions slot 12, WP-16 (post-close, owner-approved): the Shapiro transcription in degree~\texorpdfstring{$1$}{1}, proved}

\noindent\emph{This version.} The one blocking finding (F1) of the slot-12 WP-13 referee
(\file{docs/implementation/slot12/WP-13\_referee.md}) asked for a proof of ``Proposition~T'', the identity
\cl{prop:bdrychar}(3) stated and left to Obligation~\ref{obl:taubdry}(v.b); the finality proposal v2
(\file{docs/implementation/slot12/WP-12\_finality\_proposal.md} \S3.3) sketched one. This version proves it,
in a form \textbf{stronger than the one asked for}, and adds the lemma it was wanted for. \textbf{Two
statements are added, \PNEW{}; no label is renamed or removed and no existing tag moves.} The finality
certificates (d0)/(d1) are \emph{not} added --- they are [E] matters, owner-gated, and land with their code.
Base \texttt{fb3fb84}. Report: \file{docs/implementation/slot12/WP-16\_report.md}. Commit: \texttt{4ca3555}
(the referee repairs: \texttt{[hash]}); pre.2's commit is \texttt{b73a292}.
\textbf{Referee repairs P1--P8} (the Fable referee, \file{docs/implementation/slot12/WP-16\_referee.md}:
all three results ACCEPT, status notes ACCEPT WITH MINOR REPAIRS, eight minor findings, all applied within
this version): P1 the OP1-box sentence scoped to (v.a)/(v.b), (v.c) a checked route; P2 ``for $n$ even and
$p\neq\ell$'' in \cl{prop:bdrychar}(4)'s pointer; P3 \cl{cl:chi}(2)'s ``class of $y^{n}$ carrying $\chi_0$''
recorded as proved for every even $n$; P4 \cl{prop:bdrytranscription}(2)'s right-hand side well defined over
every $\cring$ for positive-diagonal representatives; P5 the corroboration paragraph: the orientation control
is a sign control, the referee's extended grid and (T3$\partial$)-breaking control cited; P6, P7 one clause each
in the proof (why $T$ is a shift; \ref{B2} for the component of $s_iv_0$); P8 report only.

\paragraph{Normative.}
\begin{enumerate}\setlength{\itemsep}{0pt}
\item \cl{prop:bdrytranscription} \PNEW{} (with the display \cl{eq:bdrygroupformnu}): (1) the coordinate
  $\Phi$ is well defined over every $\cring$, independent of the base vertex, and an isomorphism whenever
  $\rhon(-I)=\id$ --- which \cl{prop:bdrychar}(3) had shown only inside a clause assuming $n<\ell$ and a
  field; (2) at every $\nu$ and every $p$, $\Phi\circ\Thb_g$ is an explicit sum over the left cosets
  $B_N/\Gbdry$; (3) at $\nu=1$ it is \cl{eq:bdrygroupform}. \textbf{Where this departs from the requested
  statement} (WP-13 F1, proposal \S3.3): the hypothesis $p\in\cring^{\times}$ is \emph{not needed} --- the
  proof evaluates on one period and the correction terms telescope under $j\mapsto j\cdot T$, so no division
  by $p$ occurs and $p=\ell$ is included; ``$n$ even'' (or $2=0$) is needed only for $\Phi$ to be an
  isomorphism; (T3$\partial$) is used only at the level of $H_{0}$; and the sketch's period argument (which
  needs $p\in\cring^{\times}$) is replaced. At $\nu\ge2$ the proposition gives an $H^{1}$ formula without
  the double-coset structure \cl{rmk:bdrynugeneral} leaves unstated.
\item \cl{lem:bdrytorus} \PNEW{} (the proposal's ``Lemma~B$'$'', WP-13 \S2): over $\Fl$, $n$ even, $p\neq\ell$,
  the boundary operator in degrees $0$ and $1$ lies in the image of $\Fl[\mathbf T(\Fl)]$, is diagonalisable,
  the operators commute and depend on $p\bmod\ell$ only; with the eigenvalue on a torus weight vector
  --- $\chi_{k}(p)$ in degree $0$ and $\chi_{n-k}(p)$ in degree $1$ (the proposal's ``$\chi_k$ up to
  $b\leftrightarrow n+1-b$'' in degree $1$ is corrected). All parts \PROVED{}; part (1) rests on
  \cl{prop:bdrytranscription}.
\end{enumerate}

\paragraph{Non-normative.}
\begin{itemize}\setlength{\itemsep}{0pt}
\item \cl{prop:bdrychar}(3): the sentence ``which this document does not prove'' replaced by a pointer to
  \cl{prop:bdrytranscription}(3), the old wording quoted; (4): a pointer to \cl{lem:bdrytorus}.
\item \cl{obl:taubdry}(v): the sentence on what a disagreement is; (v.b) a status note --- it is now
  corroboration of a theorem, testing the implementation; the $\nu\ge2$ sentence notes that
  \cl{eq:bdrygroupformnu} is not made a clause (owner confirmation 2 of the slot-2 plan stands). Predicates,
  rings and actions on failure unchanged.
\item \cl{cl:chi}(2): a status note --- the semisimplicity half is proved (\cl{lem:bdrytorus}), the multiset
  half stays \CONJ{}; the slot-12 evidence (WP-13's $2080$ group-form cases, WP-14's $814/814$ recorded
  operators) cited. The tag is unchanged.
\item \cl{rmk:bdrynugeneral}: one sentence (the new $H^{1}$ formula at every $\nu$; the double-coset shape
  still not asserted). The decision box after \cl{def:bdryop}: one sentence (OP1 unchanged).
\item Appendix~\ref{app:status} (two lines, the counts), Appendix~\ref{app:labelindex} (three labels), this
  appendix's provenance table and Part~IV roster, the header comment and the date line.
\end{itemize}

\paragraph{Not done here, deliberately.}
\begin{enumerate}\setlength{\itemsep}{0pt}
\item \textbf{No finality certificate}: (d0) ``dickson'' and (d1) ``boundary'' (proposal \S3.1, \S3.4) are
  [E] \code{sec:certgen} and schema matters and land with their code (owner-gated); this version only
  removes (d1)'s mathematical condition.
\item \textbf{No obligation clause added}: \cl{eq:bdrygroupformnu} at $\nu\ge2$ is not made a (v.b).
\item \textbf{No tag moves}: \cl{cl:chi} keeps its tags; its (2) is not split into two labels.
\end{enumerate}


\subsection*{v1.12.0-pre.2 --- 2026-09-26 --- hecke-extensions slot 12, WP-3: the Fable referee's repairs}

\noindent\emph{This version.} The Fable referee of v1.12.0-pre.1
(\file{docs/implementation/slot12/WP-3\_referee.md}; \textbf{0 blocking, 0 major, 10 minor}, R1--R10, R2 and R5
optional) is disposed of: \textbf{every finding is applied}. No label is renamed or removed and no statement
is added; one \COMP\textsuperscript{a} tag is added (R5). Every change is marked ``v1.12.0-pre.2'' at its
site. Base \texttt{dce4ab0}. Report: \file{docs/implementation/slot12/WP-2\_report.md} \S6. Commit:
\texttt{[hash]}; pre.1's commit is \texttt{30e5b69}.

\paragraph{Normative.}
\begin{enumerate}\setlength{\itemsep}{0pt}
\item \cl{prop:torsionprepass}(6) (R1): $H^{1}(\Ceq(\Z))$ agrees with $H^{1}(\Ga;\Vcoef)$ only after inverting
  $6$ (\cl{prop:isomodel}(4)); the conclusion (the $\Z_{(\ell)}$-types agree at $\ell\ge5$) is unchanged.
\item \cl{cor:ordprecision}(1) (R3): $\epsilon$ is Hida's projector $\lim_jU^{j!}$ on $H$, pinning
  ``the ordinary part of $H^{1}\otimes\Z_\ell$'' as an object; the exponent is written $e=e'$ (was
  $\max(1,e')$, vacuous under $e'\ge1$); \cl{open:ordprecision}'s pre.1 note likewise.
\item \cl{rmk:hidacheckable} (Hi2) (R7): the two equalities stated on named modules ($F'/\ell^{e'}F'$ and
  $H_{e'}$).
\end{enumerate}

\paragraph{Non-normative.}
\begin{itemize}\setlength{\itemsep}{0pt}
\item \cl{prop:torsionprepass}(2) (R2): why $\ell\ge5$ is needed for $e_{\max}=0$ below $\ell+1$ ($(2,2)$ has $\Z/2$).
\item \cl{cor:ordprecision} (R4): $h=2d_k-1$ (even $n$), $0$ (odd $n$), known before the run.
\item \cl{rmk:ordexponentuse}(1) (R5): \COMP\textsuperscript{a} --- the exponent of
  $\mathrm{GL}_{\mathfrak g}(\Z/\ell^{K})$ equals $N^{\mathrm{lcm}}_{\mathrm{ord}}$ in eleven tested cases (general
  equality not claimed); (2) (R6): the $F'$-block sentence reworded --- $\eord(U_{F'})$ exists, and is canonical
  only modulo $\ell^{e'}$.
\item \cl{rmk:hidacheckable}: (Hi3) (R8) the weight-$2$ step explicit, $\Delta\equiv f_{11a}\pmod{11}$ named,
  the referee's classical check recorded under the item's \COMP\textsuperscript{a}; (Hi5) (R9) the Eisenstein
  line a simultaneous eigenline; (Hi6) (R10) the Bezout step named and $\tau(2411)$'s independent re-derivation
  recorded.
\item Appendix~\ref{app:status}: \cl{rmk:ordexponentuse}'s row gains \COMP\textsuperscript{a}, the counts
  (\COMP{} $25\to26$). Header comment, date line.
\end{itemize}


\subsection*{v1.12.0-pre.1 --- 2026-09-26 --- hecke-extensions slot 12, WP-2: the foundation of Hida theory at level one --- the pre-release of an \texorpdfstring{$X$}{X} bump}

\noindent\emph{This version.} The slot's ``contract first'' package
(\file{docs/implementation/slot12/00-slot12-plan.md} \S2 ``WP-2'', rules~20 and~21), written with the owner
away for the slot: the plan's ruling L3 keeps \cl{open:ordprecision}'s status macro until the owner-signed
release. It carries its own identifier so that no version number is published twice under different text
(\ref{V1}, \ref{V2}); the Fable referee (slot-12 WP-3) follows and its repairs are v1.12.0-pre.2.
\textbf{Six statements are added, \PNEW{}; no label is renamed or removed and no existing tag moves};
Appendix~\ref{app:status} gains six lines. Base \texttt{e2d137a}. Report:
\file{docs/implementation/slot12/WP-2\_report.md}. Commit: \texttt{[hash]}.

\paragraph{Normative.}
\begin{enumerate}\setlength{\itemsep}{0pt}
\item \textbf{Part~VII, \S\ref{sec:hida}} \PNEW{}: \cl{prop:torsionprepass}, \cl{cor:ordprecision},
  \cl{lem:ordexponent}, \cl{rmk:ordexponentuse}, \cl{def:hidatower}, \cl{rmk:hidacheckable} (the
  predictions (Hi1)--(Hi6), committed before slot-12 WP-4 and WP-5's code).
\item \textbf{Where this departs from the plan's statements.} (a) The plan states the pre-pass at
  $\ell\ge5$; \cl{prop:torsionprepass}(1)--(5) hold at \emph{every} prime for group cohomology (only
  $H^{0}(\Ga;\Vcoef)=0$ is used), and what needs $\ell\ge5$ is its use as a precision for the model
  (clause~(6)). (b) The plan's working precision $K=k+e_{\max}$ for ``the charpolys of every $T_p$ on the
  ordinary part'' is sufficient but not needed: \cl{cor:ordprecision}(1) shows the ordinary part is read at
  the target precision itself, the torsion being non-ordinary; $K=e'+e_{\max}$ is what the \emph{free
  quotient} needs (clauses~(2)--(4), which also justify slot-11 WP-18c's free-block reading). (c) The
  plan's (H1)--(H6) and $r$ are written (Hi1)--(Hi6) and $\mathfrak g$ (rule~\ref{A5}: \ref{H1}--\ref{H4}
  and the restriction map $r$ are taken); the exponent is $N_{\mathrm{ord}}$ ($N$ is the Hecke level) and
  the target precision $e'$ ($k$ is the weight). (d) \cl{lem:ordexponent} also gives the lcm variant and
  every multiple; \cl{rmk:ordexponentuse} shows no term of the exponent can be dropped.
\item \cl{open:ordprecision}: a v1.12.0-pre.1 status note --- answered at $\ell\ge5$, $n\ge1$ by
  \cl{cor:ordprecision}, pending release; \OPENSTAT{} kept. The question text is kept.
\end{enumerate}

\paragraph{Non-normative.}
\begin{itemize}\setlength{\itemsep}{0pt}
\item \cl{rmk:ordtorsionfree}: one parenthesis pointing at \cl{cor:ordprecision}; its tag and statement are
  unchanged.
\item Three bibliography entries (\cite{hida_ens}, \cite{hida_invent}, \cite{hida_book}), none of them in
  \file{docs/literature/} and every locator into them flagged ``not verified''.
\item The abstract, ``How to read this document'' (Part~VII), \S\ref{sec:codemap} (a row),
  \S\ref{sec:notimplemented} (a v1.12.0-pre.1 box), Appendix~\ref{app:status} (six lines, the counts),
  Appendix~\ref{app:labelindex}, this appendix's provenance table and rosters, the notation table of
  \S\ref{sec:notation-dickson}, the header comment and the date line. The companion pointer is unchanged
  ([E] v1.12.1).
\end{itemize}

\paragraph{Not done here, deliberately.}
\begin{enumerate}\setlength{\itemsep}{0pt}
\item \textbf{No Hida theorem}: control and weight independence stay predictions (\ref{Hi3}, \ref{Hi4},
  \CONJ); they, the $\Ga_0(\ell)$ model, the lattice and the large cells are slot~13's
  (\S\ref{sec:hida}'s scope paragraph).
\item \textbf{No tag moves}; \cl{open:ordprecision} keeps \OPENSTAT{} (ruling L3).
\item \textbf{No code}: the chapter's code map says what slot-12 WP-4 and WP-5 are to build.
\end{enumerate}


\subsection*{v1.11.1 --- 2026-09-25 --- hecke-extensions slot 11, WP-17 (the slot close): the WP-15 referee's repairs}

\noindent\emph{This version.} A $Y$ bump (\ref{V2}): the Fable referee of v1.11.0
(\file{docs/implementation/slot11/WP-15\_referee.md}; \textbf{accept with repairs, 0 B / 2 M / 6 m}) is
disposed of. \textbf{No normative statement changes}; one \PROVED{} remark is added (new non-normative
material, a corollary of statements already here); no label is renamed or removed. Every change is marked
``v1.11.1'' at its site. Base: \texttt{hecke-extensions} at \texttt{9833321}. Report:
\file{docs/implementation/slot11/WP-17\_report.md}. Commit: \texttt{[hash]}.

\paragraph{Non-normative.}
\begin{itemize}\setlength{\itemsep}{0pt}
\item \cl{rmk:thetacheckable} (M-1): the \COMP\textsuperscript{r} paragraph's provenance parenthesis
  restated to the reports --- WP-6's predictions committed before their runs; WP-5's one rule-20 deviation
  (two closed forms checked on their first run); WP-6 \S1a's re-alignment of (C5) and (C7) after the first
  runs --- with the v1.11.0 wording quoted. The item's evidence and tags are unchanged.
\item \cl{rmk:ordtorsionfree} \textbf{new}, \PROVED{} (m-2, the ruling on the author's F-4): at $\ell\ge5$,
  $n\ge1$ the torsion of $H^{1}(\Ga;\Sym^{n}\Z^{2})\otimes\Z_\ell$ is non-ordinary, so the ordinary part is
  torsion-free at level~$1$ and is the ordinary part of the free quotient; proof from
  \cl{cor:thetanonordinary}(3) and \cl{prop:ordinary}(1)--(4). Appendix~\ref{app:status}: one row, the
  counts ($248\to249$).
\item \cl{open:ordprecision}: a v1.11.1 status note --- the torsion half answered by
  \cl{rmk:ordtorsionfree}; the per-cell bound computable before the run stays \OPENSTAT. The question text
  is kept.
\item Part~VII's code map (m-3): the pre.1 sentence on \cl{def:ordinary} marked ``(through
  v1.11.0-pre.2)''; and the $\ell\in\{2,3\}$ refusal recorded as \Espec{sec:thetaobj}~(h)'s precondition
  (the ruling on F-2; the code changed at slot-11 WP-17).
\item The companion pointer (m-4) names the companion's v1.11.1 (\texttt{4999a21}), which preceded this
  document's v1.11.0-pre.1 and which the pointer had skipped; the pointer reads [E] v1.12.1.
\item \S\ref{sec:notimplemented}: a v1.11.1 box with the counts re-derived from \code{--list} ($141$
  registered, $82$ \code{Hecke*}). The abstract's parenthesis. Header comment, date line.
\item v1.11.0's commit (\texttt{a8e0afb}) back-filled (\ref{CL6}).
\item \textbf{Outside this document} (the same package): the companion's v1.12.1 (the referee's M-2, m-1,
  m-6); the console plan's m-5 (wording of a schedule row); the examples document [EX] v1.0.0 (owner
  sign-off); the implementation plan [X] v1.15.
\end{itemize}

\paragraph{Not done here, deliberately.}
\begin{enumerate}\setlength{\itemsep}{0pt}
\item \textbf{No Hida section}: it is slot~12's (or slot~13's, by the owner's leave), with the per-cell bound
  of \cl{open:ordprecision}.
\item \textbf{No tag moves}; \cl{open:ordprecision} keeps \OPENSTAT.
\end{enumerate}


\subsection*{v1.11.0 --- 2026-09-25 --- hecke-extensions slot 11, WP-15: the slot's \texorpdfstring{$X$}{X} bump published, with what slot 11 built}

\noindent\emph{This version.} The slot's documents package
(\file{docs/implementation/hecke\_extensions\_slot11\_plan.md} \S3 ``WP-15'', \S4.3) publishes the $X$ bump
whose text v1.11.0-pre.1 and v1.11.0-pre.2 carried, with the status of what slot~11 built recorded at the
sites that promised it (``recorded at v1.11.0''). Sources: the reports of slot-11 WP-5, WP-6, WP-11,
WP-18a, WP-18b and WP-18c and the slot's independent review
(\file{docs/implementation/slot11/WP-14\_review.md}, \file{WP-14\_review\_pass2.md}), all in
\file{docs/implementation/slot11/}. Base \texttt{1106200}: every slot-11 code package merged. One version
is one text: pre.1 and pre.2 keep theirs. \textbf{No statement is added, no label is renamed or removed,
and one tag is added} (\COMP\textsuperscript{r} on \cl{rmk:thetacheckable}); Appendix~\ref{app:status}
changes in that row and in the counts. Every change is marked ``v1.11.0'' at its site. A short Fable
referee of the delta since pre.2 follows (slot-11 WP-15). Report:
\file{docs/implementation/slot11/WP-15\_report.md}. Commit: \texttt{a8e0afb} (back-filled at v1.11.1, \ref{CL6}).

\paragraph{Normative.}
\begin{enumerate}\setlength{\itemsep}{0pt}
\item \cl{obl:backend}(iv), the scope of the exclusion, worded to the digit route as built
  (\file{ENGINE\_API.md} \S13.11.2, \S13.11.8): a digit primitive is a single-word operation or one product
  into a fixed double word (pre.2 said ``a native $w$-bit operation of the mod-$\ell^{m}$ backend of (v)'',
  and (v) is the mod-$\ell$ backend); $K$ is fixed per cell before the digit run from the certified bound
  $K=k+e_{\max}$ and capped by [E]'s grammar; the integer-tier ladder governs \cl{thm:backend}'s tables and
  is unchanged, and the digit route is not a rung of it (\file{ENGINE\_API.md} \S13.11.8 had called it
  ``the fixed-width wide rung''; it now reads as here). The ban and its owner ruling are unchanged.
\item \cl{rmk:thetacheckable}(C9): the indexing of $\rank\delta$ is said --- at the source $n'=n-6$
  ($n'\in\{30,38\}$, the targets $36$, $44$), which is what the question needs, and at the cell itself
  ($38$, $42$), which is what pre.2's parenthesis stated (slot-11 WP-6's finding~1, confirmed by the review).
  Both are true; the item's claim is unchanged.
\end{enumerate}

\paragraph{Non-normative (status).}
\begin{itemize}\setlength{\itemsep}{0pt}
\item \cl{rmk:thetacheckable}: a last paragraph, \COMP\textsuperscript{r}, with the registered checks
  (\code{HeckeDicksonThetaValidation}, \code{HeckeDicksonThetaHeckeValidation} and their \code{slow}
  companions) and their outcome: every checkable item held, (C7)'s open alternative decided at $n=14$, (C9)
  decided cell by cell. The re-derivation note after \cl{thm:thetales} points to it.
\item \S\ref{sec:notimplemented}: a v1.11.0 box. Part~VII's code map: ``Built at v1.11.0''.
  \S\ref{sec:codemap}: a Part~VII row, and (g) in \S\ref{sec:integral}'s row. The abstract's parenthesis.
\item \cl{prop:relationsmodel}: the note after its proof records slot-11 WP-11 --- the lifts for
  $\mathrm{diag}(1,p^{2})$, $\mathrm{diag}(1,p^{3})$ and $pI$ built and certified, clause~(2) holding for
  $\nu\le2$ on $H^{q}(\CI(\cring))$, $q\le5$, at $\ell\in\{2,3\}$ with $p=\ell$. \cl{rmk:integralevidence}:
  a ninth bullet (\code{HeckeStrictLiftRelationsValidation}) and the title.
\item \cl{open:ordprecision}: a status note --- slot-11 WP-18c's per-cell certified bound for the free part
  ($K=k+e_{\max}$) and the ordinary idempotent on it are built; neither answers the question, which stays
  \OPENSTAT{}; its pointer into \file{ENGINE\_API.md} corrected (\S13.12 $\to$ \S13.11.11).
\item The commits of v1.11.0-pre.2 (\texttt{fb9dcbf}) and v1.10.1 (\texttt{bee8afb}) back-filled
  (\ref{CL6}). Header comment, date line, companion pointer ([E] v1.12.0, published in the same package).
  The registered-validation counts of \S\ref{sec:codemap}'s obligations row stay left to the slot close.
\item \textbf{Outside this document} (the same package): the companion's v1.12.0; the console plan's
  v3.9; the integral-line design brief's erratum items 17--19 (X7 in full, arbitrary $k$ by the digit
  induction, the ordinary projector as built); \file{ENGINE\_API.md} \S13.11.4, \S13.11.6, \S13.11.8 and
  \S13.12.9 (the review's m-3 and m-4, and \S13.11.8 worded to agree with \cl{obl:backend}(iv)); the
  slot-11 plan's Hida-exponent line corrected.
\end{itemize}

\paragraph{Not done here, deliberately.}
\begin{enumerate}\setlength{\itemsep}{0pt}
\item \textbf{No Hida section} (decision~2 (A)); \cl{open:ordprecision} stays open.
\item \textbf{No tag upgrade}: the registered checks of Part~VII and of \cl{prop:relationsmodel} are
  evidence; no \PROVED{}, \BSTD{} or \CONJ{} tag moves, and the literature locators of
  \cl{rmk:thetaclassical} stay ``not verified''.
\item \textbf{Nothing of [EX]} (\file{docs/latex/hecke\_examples.tex}): it cites this document and is not
  cited by it.
\end{enumerate}


\subsection*{v1.11.0-pre.2 --- 2026-09-25 --- hecke-extensions slot 11, WP-2: the WP-3 referee's repairs}

\noindent\emph{This version.} The Fable referee of v1.11.0-pre.1 (slot-11 WP-3,
\file{docs/implementation/slot11/WP-3\_referee.md}) accepted it after repairs: no blocking, one major
(M-1), nine minor findings (m-1 \ldots{} m-9), and \textbf{confirmed \cl{prop:heckeH0} independently}
(the plan's ``$U_\ell=0$ on all of $H^{0}$'' is false; $U_5\DkD^{6}=-\DkL^{20}$). One version is one text:
pre.1 keeps its text under its number. \textbf{No label is renamed or removed and no tag moves};
Appendix~\ref{app:status} is unchanged except one added citation in \cl{rmk:thetaclassical}'s row. Every
change is marked ``v1.11.0-pre.2'' at its site. Base: \texttt{hecke-extensions} at \texttt{9c2ee40}.
Report: \file{docs/implementation/slot11/WP-2\_report.md}. Commit: \texttt{fb9dcbf} (back-filled at
v1.11.0).

\paragraph{Normative.}
\begin{enumerate}\setlength{\itemsep}{0pt}
\item \textbf{M-1} --- \cl{obl:backend}(iv) gains the \emph{scope of the exclusion} (owner ruling (ii) of
  2026-09-25, hecke-extensions slot~11): the $\ell$-adic digit induction of \file{ENGINE\_API.md}
  \S13.11 with a digit count fixed per cell from the certified precision and native-word digit
  operations is not arbitrary precision; a run-time-growing digit count remains excluded. The referee's
  draft, verbatim. The ban itself is unchanged.
\item \textbf{m-1} --- \cl{prop:heckeH0}: the hypotheses are stated per clause ((3) for $\ell$ odd), and
  the proof of (4) gains the case $\ell=2$ ($U_2\DkD^{3}=\DkL^{2}$ at $m=6=\ell^{3}-\ell$), so (4) holds at
  every prime as stated.
\end{enumerate}

\paragraph{Non-normative.}
\begin{itemize}\setlength{\itemsep}{0pt}
\item \textbf{m-2} --- \cl{rmk:thetacheckable}(C5): $H^{1}(M_{22})$ at $\ell=11$ is $\thL H^{1}(M_{10})$ plus
  the Eisenstein class plus a second $p+p^{2}$ class; that system has multiplicity $2$ there.
\item \textbf{m-3} --- the paragraph after \cl{thm:thetales} is a re-derivation note and no longer carries
  an ad hoc tag, so \cl{app:status}'s row and counts are consistent as they stand.
\item \textbf{m-4} --- \cl{rmk:thetasmallell}: which stabilisers fail at $\ell=3$ ($\langle U\rangle$) and at
  $\ell=2$ (all).
\item \textbf{m-5} --- \cl{rmk:thetaclassical}(5): \cite{coleman_voloch} beside \cite{gross} (not
  verified), and the endpoint $k=\ell-1$ named as the case (C8) uses.
\item \textbf{m-7} --- \cl{prop:heckeH0}(3)'s proof: an unused clause deleted.
\item \textbf{m-8} --- pre.1's commit \texttt{3b223ef} back-filled.
\item \textbf{m-9} --- \cl{rmk:thetacheckable}(C7): the pipeline's answer ($p^{2}\tau(p)$ first at $n=14$,
  $\ell=13$, multiplicity $2$) recorded beside the unchanged disjunction.
\item \textbf{m-6} concerns the WP-2 report only (the Ash--Stevens paper named).
\item Header comment, date line, one bibliography entry.
\end{itemize}


\subsection*{v1.11.0-pre.1 --- 2026-09-25 --- hecke-extensions slot 11, WP-2: the Dickson invariant and $\thL$, the ordinary projector, and the Hecke relations on the free model --- the pre-release of an \texorpdfstring{$X$}{X} bump}

\noindent\emph{This version.} The slot's ``contract first'' package
(\file{docs/implementation/hecke\_extensions\_slot11\_plan.md} \S3 ``WP-2'', rule~21), written on the
owner's rulings of 2026-09-25 (\S10 there: every decision at recommendation; decision~2 (A), so no Hida
section), with the plan's audit (\file{docs/implementation/hecke\_extensions\_slot11\_plan\_audit.md}) M-1,
m-1, m-2, m-11, m-12, m-13 taken up. It carries its own identifier so that no version number is published
twice under different text (\ref{V1}, \ref{V2}); the Fable referee (slot-11 WP-3) follows and its repairs
are v1.11.0-pre.2, the slot's documents package publishes v1.11.0. \textbf{Twenty statements are added,
\PNEW{}; no label is renamed or removed and no existing tag moves}; Appendix~\ref{app:status} gains twenty
lines. Base \texttt{54bc2b5} (R0). Report: \file{docs/implementation/slot11/WP-2\_report.md}. Commit:
\texttt{3b223ef} (back-filled at v1.11.0-pre.2).

\paragraph{Normative.}
\begin{enumerate}\setlength{\itemsep}{0pt}
\item \textbf{Part~VII, \S\ref{sec:dicksontheta}} \PNEW{}: \cl{def:dicksoninv}, \cl{lem:dicksontwist}
  (\eqref{eq:dicksontwist} for every $A\in\Mat_2(\Z)$, singular included), \cl{prop:dicksonH0},
  \cl{prop:heckeH0}, \cl{rmk:heckeH0plan}, \cl{lem:heckefibre}, \cl{def:Qn} (the complement basis
  $(\bar u_0,\ldots,\bar u_{\ell-1},\bar u_{\nplus})$ \emph{fixed}), \cl{lem:Qnfunctions}, \cl{def:theta},
  \cl{prop:thetatwist} (the twist is $p$, derived from \cl{def:rhon}'s covariance and \ref{H2}; the plan's
  audit m-1), \cl{thm:thetales}, \cl{rmk:thetasmallell} (audit m-12), \cl{rmk:thetaclassical} (audit
  m-11), \cl{rmk:thetacheckable} (WP-6's prediction list).
\item \textbf{The plan's ``$U_\ell=0$ on all of $H^{0}$'' is not adopted}: \cl{prop:heckeH0}(2)--(4) proves
  $U_\ell v=v(1,0)\,G_m$, rank $\le1$ and $U_\ell^{2}=0$ on $H^{0}$, $U_\ell=0$ there for $n<\ell^{3}-\ell$,
  and $U_\ell\neq0$ at $n=\ell^{3}-\ell$; $U_5\DkD^{6}=-\DkL^{20}$ at $(120,5)$, inside the slot's grid.
  The audit's M-1 had corrected an earlier formula and checked $\DkD$, $\DkD^{2}$ only. ``Every Dickson class
  is non-ordinary'' survives, as \cl{cor:thetanonordinary}(2)--(3).
\item \textbf{Part~VII, \S\ref{sec:ordinary}} \PNEW{}: \cl{def:ordinary}, \cl{prop:ordinary},
  \cl{cor:thetanonordinary}, \cl{rmk:ordinaryeis} (audit m-13), \cl{open:ordprecision} (\OPENSTAT{}: the
  certified precision, slot~12's).
\item \textbf{\S\ref{sec:relationsmodel}} \PNEW{}: \cl{prop:relationsmodel} --- \cl{thm:relations}'s
  recursion on $H^{q}(\CI(\cring))$ for every commutative $\cring$ (slot-10 review finding~6; line~2's
  M-3). \cl{thm:relations} itself is unchanged and keeps $6\in\cring^{\times}$.
\end{enumerate}

\paragraph{Non-normative.}
\begin{itemize}\setlength{\itemsep}{0pt}
\item Notation \S\ref{sec:notation-dickson} (sans-serif $\DkL$, $\DkD$ and $\thL$, since $L$, $\Dn$ and
  $\theta_g$ are taken, rule~\ref{A5}; $\nplus$, since $N$ is the Hecke level only), five preamble macros,
  eight bibliography entries (\cite{dickson1911}, \cite{swinnerton_dyer}, \cite{katz1977},
  \cite{jochnowitz_local}, \cite{edixhoven}, \cite{gross}, \cite{rhie_whaples}, \cite{lucas1878}), none of
  them in \file{docs/literature/} and every locator into them flagged ``not verified''.
\item The abstract, ``How to read this document'', \S\ref{sec:notimplemented} (a v1.11.0-pre.1 box),
  Appendix~\ref{app:status} (twenty lines, the counts), Appendix~\ref{app:labelindex}, this appendix's
  provenance table and rosters, the header comment and the date line.
\item Companion pointers unchanged ([E] v1.11.0); [E]'s \texttt{cl:Tlnil} is upgraded to a pointer at
  \cl{prop:thetatwist} by slot-11 WP-4, not here.
\end{itemize}

\paragraph{Not done here, deliberately.}
\begin{enumerate}\setlength{\itemsep}{0pt}
\item \textbf{No Hida section} (decision~2 (A)): the ordinary-part predictions at $(700,691)$ and
  $(10,2411)$, control and the lattice are slot~12's.
\item \textbf{No \COMP{} for Part~VII beyond the ad hoc script} of \cl{prop:heckeH0}'s proof and the
  paragraph after \cl{thm:thetales}: the registered checks are WP-5 and WP-6's, recorded at v1.11.0.
\item \textbf{No claim about $\thL$'s long exact sequence at $\ell\in\{2,3\}$} (\cl{rmk:thetasmallell}).
\end{enumerate}


\subsection*{v1.10.1 --- 2026-09-25 --- hecke-extensions slot 10, WP-15, phase 2: the referee's repairs --- a \texorpdfstring{$Y$}{Y} bump}

\noindent\emph{This version.} The second phase of the package that wrote v1.10.0
(\file{docs/implementation/hecke\_extensions\_slot10\_plan.md} \S3 ``WP-15''). The Fable referee of
v1.10.0 and of the companion's v1.11.0-pre.2 (\file{docs/implementation/slot10/WP-15\_referee\_phase1.md})
accepted them after repairs: two major findings (M-1, M-2), each a one-line text repair, and six
minor ones, of which m-1 \ldots{} m-4 concern this document (m-5 the companion, m-6 the slot plan).
One version is one text: v1.10.0 keeps its text under its number, and these repairs carry their
own. \textbf{No normative predicate, no hypothesis, no displayed formula, no tag and no label
moves}; Appendix~\ref{app:status} is unchanged. Every change is marked ``v1.10.1'' at its site.
Base \texttt{47e3ff9}: every slot-10 code package merged. Report:
\file{docs/implementation/slot10/WP-15\_report.md}. Commit: \texttt{bee8afb} (back-filled at v1.11.0).

\paragraph{Non-normative.}
\begin{itemize}\setlength{\itemsep}{0pt}
\item \textbf{M-1} --- \cl{rmk:bredonhecke} said that at $e=1$ the shipped $\Fl$ pipeline computes
  (a) and (b) for the canonical $\tau$ ``unchanged since v1.0.0''. The canonical $\tau$ was
  re-defined at v1.1.0 and the golden file re-baselined (\cl{rmk:iotatausym}, ``Two consequences''
  (b)), so neither $\tau$ nor (a) is unchanged since v1.0.0. The sentence now reads: the canonical
  $\tau$ of \cl{alg:tau} since v1.1.0; its operator (b) did not move in that re-baseline and equals
  the legacy $\tau$'s at every $\ell\in\{2,3\}$ golden point, \cl{rmk:tauexistsscope}(1) --- an
  observation, not a theorem. ``Compares nothing there'' becomes ``compares none of these operators
  there'' ($\Th^{s}_g$, the canonical $\Th_g$, $\Th_{\Iso,g}$), with a pointer to that observation.
\item \textbf{M-2} --- \cl{rmk:strictliftbuilt}(3): $349\to359$ cells, an arithmetic slip inherited
  from a report's summary line (per ring $34,34,42,55,42,42,55,55$).
\item \textbf{m-1} --- v1.10.0's header comment said \cl{rmk:integralscope}(2) was ``narrowed to what
  is not built''; it gains a pointer, and the comment now says so.
\item \textbf{m-2} --- \cl{rmk:strictliftbuilt}(3): ``on the rings where a $\Th^{s}_g$ is shipped''
  becomes ``on eight rings among those with a shipped $\Th^{s}_g$'', since $\Th^{s}_g$ is shipped over
  $\Fl$ at every $\ell$.
\item \textbf{m-3} --- \cl{rmk:bredonhecke}(a): why the canonical $\tau$ is integral --- in practice,
  because since v1.1.0 it vanishes in degree~$2$ (\cl{rmk:iotatausym}, consequence (a)), so
  \cl{alg:tau}'s mod-$\ell$ triangle-fill fallback is never used.
\item \textbf{m-4} --- \cl{rmk:integralevidence}, the \code{HeckePeriodicityClassValidation} bullet:
  $\kappa_C(\res_C e)=\chi_C$ is read on the stabilisers; the $224$-element check is $\chi=\zeta_+$.
\item \textbf{Status} --- \cl{rmk:integralevidence}'s check counts after slot-10 WP-8's validation
  carries: \code{HeckeFreeModelValidation} $803\to801$, \code{HeckePeriodicityClassValidation}
  $2482\to2483$. Slot-10 WP-7's levers (the $\iota$ sign split and route~2 of
  \code{HeckeOperatorOnK} as host tasks) and WP-8's ninth fingerprint bit move no value and touch no
  site of this document. The commit of v1.10.0 (\texttt{ca6c9b9}) is back-filled (\ref{CL6}).
  Header comment, date line, companion pointer ([E] v1.11.0, published at this version's commit).
  The registered-validation counts stay left to the slot close.
\end{itemize}


\subsection*{v1.10.0 --- 2026-09-25 --- hecke-extensions slot 10, WP-15: the strict lift as built, the Bredon-object operator not compared, the sign of the Bockstein, and (vii) over the ring --- an \texorpdfstring{$X$}{X} bump}

\noindent\emph{This version.} The slot's one $X$ revision of this document
(\file{docs/implementation/hecke\_extensions\_slot10\_plan.md} \S3 ``WP-15'', \S4.3), made on the
owner's rulings of 2026-09-25 on the slot-10 decisions (\S10 there, every one at
recommendation~(A)): 12~(A), this $X$ with a Fable referee; 3~(A), \cl{obl:coneexact}(vii) over
$\Z/\ell^{e}$; 4~(A), the Hecke action on the Bredon object at $\ell\in\{2,3\}$ declined, with a
remark here. Sources: the reports of slot-10 WP-3, WP-10, WP-11 and WP-13
(\file{docs/implementation/slot10/}) and the slot's independent review
(\file{docs/implementation/slot10/WP-14\_review.md}, passes 1 and 2), whose findings M-1, M-3, m-1,
m-6, m-7, m-15 and m-17 this version takes up. Base \texttt{f0640d0}: WP-2, WP-3, WP-6, WP-9,
WP-10, WP-11, WP-12 and WP-13 merged; slot-10 WP-7 and WP-8 (lane~B) touch no site of this
document. \textbf{Two statements are added, \PNEW{}; one proof gains a sentence fixing a sign
convention}; the rest is status. \textbf{No label is renamed or removed and no existing tag moves};
Appendix~\ref{app:status} gains two lines. Report:
\file{docs/implementation/slot10/WP-15\_report.md}. Commit: \texttt{ca6c9b9} (back-filled at
v1.10.1).

\paragraph{Normative.}
\begin{enumerate}\setlength{\itemsep}{0pt}
\item \cl{lem:bulkperiodic}, proof, Step~0 --- \textbf{the sign of $\delta$ is fixed}: the
  snake-lemma connecting map $\delta[c]=[d\tilde c]$, the cochain differential carrying no sign.
  Through v1.9.1 the step named the Bockstein without its sign, so clause~(1)'s ``unique class $e$''
  was determined only up to $\pm$; it is now the class the built one is (the slot's review, m-6).
  \textbf{Pins a convention}; no clause's truth moves, since (2) and (3) are insensitive to the
  sign.
\item \cl{rmk:strictliftbuilt} \PNEW{} (\COMP\textsuperscript{r}) --- \cl{prop:strictlift}(1), (3)
  and (4) as built for $g=\mathrm{diag}(1,p)$, $p\in\{2,3,5\}$, in degrees $\le q_{\max}+1=6$: the
  construction (the preimage read off $\tau^{s}$, $\ker\cmpI$ contracted by the columns' closed-form
  homotopies, no solve), the five certificates, (3) and (4) in use, and the integer sizes measured
  (the first evidence on [IB]'s K1). Its last paragraph keeps ``existence only'' for every other
  $g$ and degree. \cl{prop:strictlift}'s closing paragraph and \cl{rmk:integralscope}(2) gain
  pointers to it; neither statement changes.
\item \cl{rmk:bredonhecke} \PNEW{} (commentary) --- beside \cl{prop:strictlift}(4): at
  $\ell\in\{2,3\}$ the operator on cohomology of the Bredon object is \textbf{not shown} independent
  of $\tau$ (homotopy uniqueness, \cl{lem:tauexists}(2), needs $6\in\cring^{\times}$), this document
  compares nothing there, and that is the absence of a proof, not a proof of dependence (owner
  ruling 4~(A); plan \S9 audit finding M-2's wording). No line-2 package exhibited two $\tau$ with
  different operators, so the ``not shown'' reading stands.
\end{enumerate}

\paragraph{Non-normative.}
\begin{itemize}\setlength{\itemsep}{0pt}
\item \cl{obl:coneexact}(vii) --- a v1.10.0 status paragraph over $\Z/\ell^{e}$: \textbf{sentences
  1--2 discharged at $\ell^{e}<2^{32}$, every E8 point} (slot-10 WP-3), the cross-routes refusing
  above $2^{32}$ by the \texttt{uint32\_t} regime rule; \textbf{sentence~3 recorded as unavailable
  over the ring} --- the pairing and the contragredient refuse a composite ring. The slot plan's and
  WP-3's hand-over read ``(vii) over the ring discharged''; the review's M-1 scoped it to the two
  sentences built, and this version follows the review. The v1.9.0 sentence is kept, in the past
  tense, with its staleness window (\texttt{c9684e9} to this version). The obligation roster's row
  says the same.
\item \cl{sec:codemap} --- the \S\ref{sec:cone} row gains the cross-routes over the ring; the
  \S\ref{sec:integral} row gains (d) $\tau_{\Iso}$ and $\Th_{\Iso,g}$ (WP-10), (e) $e$ and $E_q$
  (WP-13), (f) the basis-free Hecke action (WP-11). \S\ref{sec:integral}'s closing ``Code map''
  paragraph, stale from v1.9.0, is pointed at the row.
\item \cl{sec:notimplemented} --- a v1.10.0 paragraph: (vii) over the ring as above; E8 complete;
  at $\ell\in\{2,3\}$ the refusals stand for \cl{rmk:bredonhecke}'s reason; \S\ref{sec:integral}
  with the lifted datum, $e$ and the action built, $T_pT_{p'}=T_{p'}T_p$ tested and
  \cl{thm:relations}(1)'s recursion not (a determinant-$p^{2}$ lift is not built; the review's M-3);
  the descent, the value $\Z$ and M1(b) not built.
\item \cl{rmk:integralevidence} --- title and count (``slots 8, 9 and 10'', eight validations) and
  three bullets: \code{HeckeStrictLiftValidation}, \code{HeckePeriodicityClassValidation} (with the
  identification of the realised connecting map, the review's m-7) and
  \code{HeckeIntegralActionValidation} (with the diagonalisability the congruence-module readings
  assume, m-15, and the $(16,3)$ row recorded without a $T_3$ oracle, m-17); the
  \code{HeckeFreeModelValidation} bullet names where its tests (3)--(5) are hosted. Tag unchanged.
\item \cl{def:ringsemantics}(3)(d), \cl{rmk:mpcaveat}(4), \S\ref{sec:integral}'s ``What is not
  here'', the abstract --- pointers.
\item Appendix~\ref{app:status} (two lines, the counts), Appendix~\ref{app:labelindex} (two
  labels), this appendix's provenance table and Part~VI roster (whose ``nothing tests the rest'',
  stale since v1.9.0, is pointed), the header comment, the date line and the companion pointer
  ([E] v1.11.0). The commit of v1.9.1 (\texttt{1f190e5}) is back-filled (\ref{CL6}). Containers.
\item \textbf{Outside this document} (\ref{CL2}): the companion specification's v1.11.0 (whose
  pre-release v1.11.0-pre.2 is at this version's commit), the console plan's v3.8 and the design
  brief's slot-10 erratum items, by the same package.
\end{itemize}

\paragraph{Not done here, deliberately.}
\begin{enumerate}\setlength{\itemsep}{0pt}
\item \textbf{\cl{obl:coneexact}(vii) is not narrowed}: its third sentence stays as stated; over the
  ring it is recorded unavailable by its own ``until then'' rule, and the E7 pairing over a
  composite ring is code for a later slot.
\item \textbf{No effectivity statement for $\tau_{\Iso}$ is added to \cl{prop:strictlift}}: the
  evidence is a remark, and ``existence only'' stays a statement about every $g$.
\item \textbf{The registered-validation counts} of \S\ref{sec:codemap}'s \S\ref{sec:obligations}
  row are left at v1.9.1's: slot-10 packages still running when this version was written may
  register validations, and the slot close re-derives them from \code{--list}.
\item \textbf{No statement about the Bredon-object operator at $\ell\in\{2,3\}$ beyond
  \cl{rmk:bredonhecke}}: in particular no claim that it depends on $\tau$.
\end{enumerate}


\subsection*{v1.9.1 --- 2026-09-25 --- hecke-extensions slot 9, WP-17, the slot close: the referee's repairs and status --- a \texorpdfstring{$Y$}{Y} bump}

\noindent\emph{This version.} The slot close (\file{docs/implementation/hecke\_extensions\_slot9\_plan.md}
\S3 ``WP-17''). The Fable referee of v1.9.0 and of the companion's v1.10.0 (plan WP-15) accepted
them after repairs: two major findings (M-1, M-2), each a one-sentence or one-file repair, and
eleven minor ones, of which m-1 \ldots{} m-4 concern this document (the others the companion,
the console plan, the design brief and the slot close). \textbf{No normative predicate, no
hypothesis, no displayed formula, no tag and no label moves}; Appendix~\ref{app:status} is
unchanged. Report: \file{docs/implementation/slot9/WP-17\_report.md}. Commit: \texttt{1f190e5}
(back-filled at v1.10.0).

\paragraph{Non-normative.}
\begin{itemize}\setlength{\itemsep}{0pt}
\item \textbf{M-1} --- v1.9.0 cited ``the slot's independent review (WP-14, pass~1)'' as evidence,
  and no such file was in the tree: the reviewer's file was lost when the host ran out of disk
  on 2026-09-25. It is reconstructed from the reviewer's two returned summaries and committed as
  \file{docs/implementation/slot9/WP-14\_review.md}, with a provenance note; \cl{rmk:certifiedeperdiff}
  cites it. v1.9.0's own changelog text is left as written, and this entry resolves its citation.
\item \textbf{M-2} --- \cl{obl:coneexact}(iv)'s as-built sentence named \code{HeckeReductionValidation}
  arm (4b-C) at the seven points. Since slot-9 WP-4t (owner-approved 2026-09-25, merged before
  v1.9.0 was written) the point $(60,5)$ at both levels runs in
  \code{HeckeReductionCoverageSlowValidation}; arm (4b-C) covers the other five points and the two
  word-regime cells. The obligation roster's joint attribution was already true and is unchanged.
\item \textbf{m-1} --- the commits of v1.9.0 (\texttt{1d7599a}) and v1.8.1 (\texttt{b7391f7}) are
  back-filled, each found by \texttt{git log -S} on its changelog header (\ref{CL6}).
\item \textbf{m-2} --- \cl{def:ringsemantics}(2) and the \S\ref{sec:integral} code-map row cited
  \cl{prop:ringinforce}(2),(3) for the non-freeness of $\Ceq(\Z/\ell^{e})$ at $\ell\in\{2,3\}$;
  neither clause states it. The citation is now the invariants computation of (2)'s proof read
  over $\Z/2^{e}$ ($(V_R)^{\langle S\rangle}\cong\Z/2^{e}\oplus\Z/2^{e-1}$ at $n=2$), which the
  companion's \Espec{sec:ring}~(c1$'$) carries out.
\item \textbf{m-3} --- the grammar of \cl{obl:coneexact}(vii)'s as-built sentence (``\ldots{} and
  the relative sub-complex constructor that route~3 needs refuse \ldots'').
\item \textbf{m-4} --- the \S\ref{sec:integral} code-map row opened ``Partly tested, not built
  (since v1.8.0)'' and later said ``Since v1.9.0: partly built''; its lead now reads ``Partly built
  since v1.9.0 (partly tested, not built, v1.8.0--v1.8.1)''.
\item \textbf{Status} --- the \S\ref{sec:obligations} code-map row's count is re-derived from
  \code{--list}: $71$ \code{Hecke*} entries of $125$ registered ($52$ \code{fast}, $63$
  \code{standard}, $10$ \code{slow}), the slot-9 names registered through the per-package
  fragments \file{Registry/Slot9/}; \S\ref{sec:notimplemented} gains a v1.9.1 sentence;
  \ref{CL5}'s list of $Y$ bumps gains v1.9.1; the header comment, the date line and the companion
  pointer ([E] v1.10.1).
\item \textbf{Outside this document} (\ref{CL2}): the companion's v1.10.1, the console plan's
  v3.7.1, the design brief's erratum item~8 and the implementation plan's v1.13, by the same
  package.
\end{itemize}


\subsection*{v1.9.0 --- 2026-09-24 --- hecke-extensions slot 9, WP-15: the F-A remark at every prime, the certified-exponent rule per differential, the coefficient ring at $2$ and $3$ as built, and phase E8 item~4b --- an \texorpdfstring{$X$}{X} bump}

\noindent\emph{This version.} The slot's one $X$ revision of this document
(\file{docs/implementation/hecke\_extensions\_slot9\_plan.md} \S3 ``WP-15'', scope~3, \S4.3),
made on the owner's rulings of 2026-09-24 on the slot-9 decisions (\S10 there, every one at
recommendation~(A)): 10~(A), this $X$ with a Fable referee; 3~(iii)~(A), \cl{obl:coneexact}(vii)
over $\Z/\ell^{e}$ by route~1 only, its composite half recorded as open; 4~(A), the lift at
$\ell\in\{2,3\}$ for the cohomology object only. Sources: the reports of slot-9 WP-2, WP-3,
WP-4, WP-10, WP-11, WP-12 and WP-13 (\file{docs/implementation/slot9/}) and the slot's
independent review (WP-14, pass~1), whose findings on this document are the per-differential
remark and the two stale sentences. \textbf{One statement is added, \PNEW{}; one remark's
clause is stated at every prime; one definition's clause is replaced by the as-built statement};
the rest is status. \textbf{No label is renamed or removed and no existing tag moves};
Appendix~\ref{app:status} gains one line. Report: \file{docs/implementation/slot9/WP-15\_report.md}.
Commit: \texttt{1d7599a}.

\paragraph{Normative.}
\begin{enumerate}\setlength{\itemsep}{0pt}
\item \cl{rmk:bocksteinexp}(2) --- \textbf{stated at every prime $\ell$}, $2$ and $3$ included;
  clauses (1), (3), (4) keep $\ell\ge5$. The proof is unchanged except for one parenthesis: its
  only use of \cl{cl:bockstein} is $H^{0}(\Ga;\Vcoef)=0$, whose proof is about $\Z$-coefficients.
  The clause names the design brief's instances (F-4): $\Z/2^{j}$ at $n=2^{j}$ from $m=2$ at
  $\ell=2$, and $\Z/9$ at $n=12,18$ from $m=4,6$ at $\ell=3$; clause~(3) records, outside its own
  scope, that the probe's table shows exactly these exponents at $n\le20$. \textbf{Widens a
  proved statement to what its proof already gave} (slot-8 WP-9 finding~6; owner ruling 10~(A)).
  Tags unchanged (\PROVED, \COMP\textsuperscript{a} and a standard-citation pointer).
\item \cl{rmk:certifiedeperdiff} \PNEW{} (\PROVED{}) --- \cl{prop:certifiede}(3)'s hypothesis may be
  checked differential by differential, each $d^{q}$ at its own level $e_q$; which levels the type
  of $H^{q}$ and the universal-coefficient display need; clause~(3) is the uniform case. A
  clarification, not a correction: no statement or proof of \cl{prop:certifiede} changes. It is
  the reading line~2's milestone M0 uses (\code{HeckeFreeModelValidation}); slot-9 WP-11
  finding~2, confirmed by the slot's review.
\item \cl{def:ringsemantics}(2) --- \textbf{the clause ``and the implementation of this build
  refuses it there'' is replaced} by the as-built statement: at $\ell\in\{2,3\}$ this build
  computes the Bredon object's lengths and invariant factors through a flat-coordinate module
  route, and refuses the $K$-mode (undefined there, \cl{prop:certifiede}(5)), the Hecke operator
  and the reductions by name. The clause was stale from slot-9 WP-3's merge (\texttt{b53836f})
  until this version, a window the slot-9 plan recorded (\S4.3) rather than closed by a second
  editor. \textbf{A status repair in a definition}: what the ring names is unchanged. Measured by
  \code{HeckeBredonTreeRingValidation}. (3)(a) and (3)(b) gain pointers: the ranks may come from
  $\CI(\Z)$ as \cl{prop:certifiede}(4) already allows, and ``per degree'' is read per
  differential.
\item \cl{obl:coneexact}(iv) --- \textbf{the parenthesis ``until it exists the non-field case is
  recorded as unchecked'' is restated}: a ring over which no module solve is built records the
  case unchecked; over $\Z/\ell^{e}$, $\ell\ge5$, the check is made (the as-built paragraph names
  \code{ConeModuleCertificate} item~14 and \code{HeckeReductionValidation} arm (4b-C)), and the
  length identity, with the composite-zero check, implies exactness independently of the
  witnesses. The same change in the clause's last sentence. \textbf{No predicate moves}: the
  clause always required the solve over a non-field ring; what changes is that the conditional
  ``until it exists'' is discharged.
\end{enumerate}

\paragraph{Non-normative.}
\begin{itemize}\setlength{\itemsep}{0pt}
\item \cl{obl:coneexact}(vi) and (vii) --- one status paragraph each, over $\Z/\ell^{e}$: (vi) is
  partial there (the commutator asserted, arm (4b-H); the $\nu=2$ recursion not computed over the
  ring); (vii) is route~1 only, its composite half open (owner ruling 3~(iii)~(A)). The obligation
  roster's \cl{obl:coneexact} row says the same.
\item \cl{sec:codemap} --- the \S\ref{sec:cone} row gains the dense module route (slot-9 WP-4:
  \code{ModuleConeT}, \code{ConeModuleCertificate}, \code{ModuleConeActionT},
  \code{HeckeBoundaryOperatorT}; the \code{uint32\_t} classes below $2^{32}$, the templated route in
  every word regime); the \S\ref{sec:integral} row is \textbf{partly built}: the scoped lift and the
  flat-coordinate route at $\ell\in\{2,3\}$ (WP-3), $\Iso_{\bullet}$, $\Isob_{\bullet}$ and $\cmpI$
  (WP-10, L2-1a) and the integral groups for $n\le20$ (WP-11, M0).
\item \cl{sec:notimplemented} --- a v1.9.0 paragraph: phase E8 complete as specified (items 1--4b;
  4b over $\Z/\ell^{e}$ by route~1; \cl{obl:coneexact}(vii)'s composite half open, (vi) partial
  over the ring); \S\ref{sec:integral} partly built. The registered-validation counts are the
  slot close's and are not re-derived here.
\item \cl{rmk:integralevidence} --- its title and count (``slots 8 and 9'', five validations) and
  three bullets: \code{HeckeBredonTreeRingValidation}, \code{HeckeIsotropyModelValidation},
  \code{HeckeFreeModelValidation}. Its tag is unchanged.
\item \S\ref{sec:integral}'s opening ``What is not here'', \cl{rmk:mpcaveat}'s last parenthesis
  and the abstract --- pointers: the free model, its comparison map and the integral groups at
  small $n$ are built; the lifted datum and the descent are not.
\item Appendix~\ref{app:status} (one line, the counts), Appendix~\ref{app:labelindex} (one
  label), this appendix's provenance table and Part~VI roster, the header comment, the date line
  and the companion pointer ([E] v1.10.0). Containers.
\item \textbf{Outside this document} (\ref{CL2}): the companion specification's v1.10.0 and the
  design brief's erratum block, by the same package.
\end{itemize}

\paragraph{Not done here, deliberately.}
\begin{enumerate}\setlength{\itemsep}{0pt}
\item \textbf{\cl{obl:coneexact}(vi) and (vii) are not narrowed.} Both stay as stated, over every
  ring with $6$ invertible; this build discharges them over $\Z/\ell^{e}$ in part, and says so. The
  recursion over the ring and the two cross-routes over the ring are code, for a later slot.
\item \textbf{The registered-validation counts} of \S\ref{sec:codemap}'s \S\ref{sec:obligations}
  row are left at v1.8.1's: slot-9 packages still running when this version was written register
  or re-count validations, and the slot close re-derives them from \code{--list}.
\item \textbf{No obligation form of \cl{prop:certifiede}}, and no record-schema rule for the
  integral line: line~2's objects carry no golden entry and no record key in slot~9.
\end{enumerate}


\subsection*{v1.8.1 --- 2026-09-24 --- hecke-extensions slot 8, WP-14, the slot close: the referee's wording repairs, the changelog's gaps, and status --- a \texorpdfstring{$Y$}{Y} bump}

\noindent\emph{This version.} The slot close (\file{docs/implementation/hecke\_extensions\_slot8\_plan.md}
\S3 ``WP-14'', scope~4). The Fable referee of v1.8.0 (plan WP-9) accepted it with six minor
findings, all wording, ``none needs a $Y$ bump before merge, all can go into the slot close's
status pass''; m-1 to m-5 are applied here and m-6 concerns the plan's wording, not this
document. \textbf{No normative statement, hypothesis, displayed formula, asserted predicate, tag
or label moves}, so Appendix~\ref{app:status}'s counts are unchanged. Close report:
\file{docs/implementation/slot8/WP-14\_report.md}. Commit: \texttt{b7391f7}.

\paragraph{Repairs (non-normative; the referee's m-1 to m-5).}
\begin{enumerate}\setlength{\itemsep}{0pt}
\item \cl{lem:strictlifting}, proof --- ``(for $q=0$, $z=0$)'' assumed $E_{-1}=0$, which the lemma
  does not; now ``for $q=0$ the same step applies; in the applications below $E_{-1}=0$ and
  $z=0$''. The proof's general step already covered $q=0$ (m-1).
\item \cl{sec:notation-integral}, the row of $F^{\sigma}_{\bullet}$ --- ``an instance of the
  $F_{\bullet}$ of \cl{def:borelmodel}'' becomes ``the analogue, for the stabiliser'': the two are
  resolutions over $\Z[\Ga_\sigma]$ and over $\Z[\Ga]$, the same kind of object, not one object
  (\ref{N2}; m-2). A description, not a rule: the symbol and its meaning are unchanged.
\item \cl{thm:specialisation}, the v1.8.0 sentence of its closing paragraph --- now says what
  v1.8.0 adds for $(\Xbulk,\Xbdry)$ (a finite-rank free model with its comparison map and a Hecke
  datum of the pair, as existence) instead of reading as if the unrefined case had been deferred;
  clauses (1)--(4) and their tags untouched (m-3).
\item \cl{prop:bredon} --- one parenthesis after the definition of $\mathcal D$: it is the family
  of all finite subgroups of $\Ga$, which this document does not use (m-4).
\item \cl{rmk:truncdevice} --- the Hecke-equivariance clause cites \cl{lem:bulkperiodic}(2),(3),
  not (2) alone (m-5).
\end{enumerate}

\paragraph{Pre-existing gaps closed (non-normative).}
\begin{itemize}\setlength{\itemsep}{0pt}
\item The abstract --- ``Two chapters, \S\ref{sec:cone} and \S\ref{sec:orbifoldmodel}, state
  objects that nothing builds'' had been stale since v1.7.3 (\S\ref{sec:cone} implemented in
  full); it now names \S\ref{sec:orbifoldmodel} and \S\ref{sec:integral}.
\item This appendix --- subsections for v1.7.2, v1.7.3 and v1.7.4 (the three status-only slot
  closes, which v1.8.0's entry found missing, ``Not done here'' item~4), reconstructed from their
  date lines, header comments and commits and adding nothing those do not say; the commit
  hashes of fifteen entries back-filled (\ref{CL6}; v0.13.0's commit could not be identified from
  the history and keeps its \texttt{[hash]}); \ref{CL5}'s list of $Y$ bumps completed from v1.6.1
  on; \S\ref{app:changelog-gaps}'s sentence that the citation census is ``specified and not
  written'', stale since hecke-extensions slot~3, pointed at \file{label\_census.sh}.
\end{itemize}

\paragraph{Status (non-normative).}
\begin{itemize}\setlength{\itemsep}{0pt}
\item \cl{sec:codemap} --- the \S\ref{sec:symconv} row gains $\rho_n$ over $\Z/\ell^{k}$
  (\code{ComputeSymmetricPowerMatrixRing}, slot-8 WP-3); the \S\ref{sec:six} row gains the
  reductions and the Bockstein over the ring in force (\file{Hecke/HeckeReduction}, WP-4) and the
  count identity of \cl{rmk:bocksteinexp}(1) as it is asserted; the \S\ref{sec:integral} row says
  the third complex of \cl{def:ringinforce} at $\Z/\ell^{k}$, $\ell\ge5$, is built (phase E8 items
  2 and 3) and refused at $\ell\in\{2,3\}$, as \cl{prop:certifiede}'s scope clause requires; the
  \S\ref{sec:obligations} row's count is re-derived ($65$ \code{Hecke*}; $114$ registered).
\item \cl{sec:notimplemented} --- a v1.8.1 paragraph: the free models of
  \S\ref{sec:orbifoldmodel} and \S\ref{sec:integral} are unbuilt; the boundary and the cone over
  $\Z/\ell^{k}$ (phase E8 item~4b) are the next slot's, so the non-field half of
  \cl{obl:coneexact}(iv) stays unchecked.
\item The header comment, the date line and the companion pointer ([E] v1.9.1). Containers.
\item \textbf{Outside this document} (\ref{CL2}): the extension plan's v1.12 and the engine
  contract's \S13 pointer repairs, by the same package.
\end{itemize}

\paragraph{Not done here, deliberately.}
\begin{enumerate}\setlength{\itemsep}{0pt}
\item \textbf{\cl{rmk:bocksteinexp}(2) at every $\ell$.} The referee recommends stating clause~(2)
  for every $\ell$ at the next $X$ (its proof uses no hypothesis on $\ell$ and at $\ell\in\{2,3\}$
  proves the observed $\Z/2^{j}$ at $n=2^{j}$ and $\Z/9$ at $n=12,18$); clause~(1) and the summand
  count stay at $\ell\ge5$. It changes a statement, so it is an $X$ bump's, and it is recorded
  for the next one (owner default of 2026-09-24).
\item \textbf{Part~VI's placement} is kept, as the referee recommends: moving it would renumber
  chapters 18--22 for no gain.
\end{enumerate}


\subsection*{v1.8.0 --- 2026-09-24 --- hecke-extensions slot 8, WP-9: the integral line (L2-C), C-A5 re-scoped onto the unrefined tessellation, and the F-A remark --- an \texorpdfstring{$X$}{X} bump}

\noindent\emph{This version.} The slot's one $X$ revision of this document
(\file{docs/implementation/hecke\_extensions\_slot8\_plan.md} \S3 ``WP-9'', \S4.3), made on the
owner's rulings of 2026-09-24 (\S10 there): 6b~(A), the chapter must-accept; 6c~(A), the
extension plan's contract item C-A5 folded into it on the \emph{unrefined} tessellation, OD-T1
answered ``not needed''; 6d, the integral-line decisions OD-3~(c), OD-4~(a), OD-6~(a),
OD-7~(a), OD-9~(a), OD-10~(a) (record keys only under the value naming $\Z$) and
OD-11~(b)/(a) confirmed; 10~(A), the F-A remark beside \cl{cl:bockstein}. Sources: the
integral-line design brief, revision~2
(\file{docs/implementation/integral\_line\_design\_brief.md}) and its audit, and the
$\ell$-adic investigation (\file{docs/implementation/e8\_ladic\_investigation.md}) ---
project notes, not sources of record; evidence: slot-8 WP-7 and WP-8. \textbf{Twenty
numbered statements are added, all \PNEW{}}; one hypothesis (\ref{H3}) and one theorem's
closing paragraph (\cl{thm:specialisation}) are amended; the notation chapter gains a
subsection; \textbf{no label is renamed or removed and no existing tag moves}.
Appendix~\ref{app:status}'s counts are re-derived mechanically from its table. Report:
\file{docs/implementation/slot8/WP-9\_report.md}. Commit: \texttt{1760a4a}.

\paragraph{Normative.}
\begin{enumerate}\setlength{\itemsep}{0pt}
\item \cl{sec:integral} --- \textbf{new chapter}, Part~VI, placed after Part~V so that no
  chapter number of v1.7.4 moves; containers \cl{sec:integralobjects}, \cl{sec:isomodel},
  \cl{sec:descent}, \cl{sec:strictlift}, \cl{sec:badprimes}, \cl{sec:ringsemantics}, the
  display \cl{eq:treemodel} and the rule tags \cl{IM1}, \cl{IM2}.
\item \cl{def:ringinforce} \PNEW{} --- the three complexes on the model (integral, its base
  change, the complex over the ring in force) and the notation rule that
  $\Ceq(\cring)$ is always the third.
\item \cl{prop:ringinforce} \PNEW{} --- the reduction map injective, an isomorphism at
  $6\in\cring^{\times}$ (from \cl{prop:coneintegral}(2)); \emph{not} a quasi-isomorphism at
  $(n,\cring)=(2,\F_2)$, with a proof (the design brief's audit had shown it on the golden
  $\dim H^{0}$ values); the odd-$n$ divergence of the four objects, including
  $H^{2}(\Ga;\Sym^{1}\Z^{2})\neq0$.
\item \cl{prop:bredon} \PNEW{} --- the shipped complex is Bredon cohomology over the ring
  in force: the fixed-set check on $\Xbulk$ and on the Bass--Serre tree, the $\Z[\Ga]$-chain
  homotopy equivalence $C_{\bullet}(\Xbulk)\simeq C_{\bullet}(\mathcal T)$ through the orbit
  category of the family $\mathcal D$, the tree complex \eqref{eq:treemodel}, independence
  of the triangulation. The design brief rested this on the classifying-space theory of
  proper actions with a pointer it had to drop as unverified; the proof here uses only the
  comparison theorem in a functor category and \cl{cl:amalgam}(1).
\item \cl{def:isomodel}, \cl{prop:isomodel} \PNEW{} --- the isotropy-resolution model on
  the unrefined $(\Xbulk,\Xbdry)$, its existence by a terminating procedure (boundary
  summands first), what it computes over every ring in every degree, the one integral
  comparison map $\cmpI^{\ast}$ (cone annihilated by $48$, a quasi-isomorphism at
  $6\in\cring^{\times}$), and its comparison with \cl{def:borelmodel}. \textbf{The ranks are
  $5$, $14$, $19$, $19$, \ldots}: the design brief (\S2.5(b), risk K10) and its audit (m-9)
  counted the $7$, $11$, $5$ simplices of $K$ instead of the $5$, $9$, $5$ $\Ga$-orbits of
  \cl{lem:model7coincidences}(2), and so gave $7(n+1)$, $18(n+1)$, $23(n+1)$ where the
  cochain ranks are $5(n+1)$, $14(n+1)$, $19(n+1)$.
\item \cl{rmk:cafive} \PNEW{} --- C-A5 discharged item by item on the unrefined
  tessellation; OD-T1 answered ``not needed'', with the reasons already in this document;
  a refinement, if ever wanted, compares integrally.
\item \cl{prop:fracture}, \cl{prop:cohglue}, \cl{prop:descent}, \cl{prop:homvanish}
  \PNEW{} --- the arithmetic fracture square as an exact sequence of complexes; glueing at
  the cohomology level exact over $\Z$, \textbf{with the design brief's finite-generation
  hypothesis removed} (the proof uses flatness only); the owner's descent realised by the one
  map $\cmpI^{\ast}$; and $\Hom^{-1}$ over $\Q_2\times\Q_3$ zero for the bulk, the coherence
  datum real only for the pair (\cl{prop:coneambiguity}).
\item \cl{rmk:truncdevice} \PNEW{} --- the untruncated complex is what is glued, the brutal
  truncation is \cl{open:ell2}'s box on the free model, and the object is not perfect
  ($n=0$).
\item \cl{lem:strictlifting}, \cl{prop:strictlift} \PNEW{} --- the strict lift of the
  strict pair datum of \cl{prop:strictpair} along $\cmpI$, boundary first, with vanishing
  pair homotopy; its coset transport intertwined with that of $\tau^{s}$ on the nose over
  every ring; agreement with the shipped operator on cohomology at $6\in\cring^{\times}$ only.
  \textbf{Existence only}: effectivity and certification are code (risk K1).
\item \cl{prop:badprimes} \PNEW{} --- the bad primes as a closed-form list, and $2$-torsion
  only at odd $n$.
\item \cl{prop:certifiede} \PNEW{} --- the certified-exponent rule for any complex of
  finitely generated free $\Z$-modules, hence at $2$ and $3$ on the free model; where the
  rank over $\Q$ comes from; and the \textbf{scope clause}: not the model complex over
  $\Z/\ell^{e}$ at $\ell\in\{2,3\}$, which is no reduction of an integral complex.
\item \cl{prop:crtcharpoly} \PNEW{} --- the residue identity behind $\chi_p$ over $\Z$ and its
  determination by residues against a bound, with no integer beyond a word; the bound is an
  input.
\item \cl{def:ringsemantics} \PNEW{} --- what each coefficient ring names; the value
  naming $\Z$ is group cohomology, never $\Ceq(\Z)$, basis-free (OD-11~(b)); the setting
  itself is \Espec{sec:ring}'s and is not restated (single-editor rule).
\item \cl{rmk:integralevidence} \PNEW{} (\COMP\textsuperscript{r}) and
  \cl{rmk:integralscope} \PNEW{} (commentary) --- what the two slot-8 validations test, and
  what the chapter does not say (no arbitrary precision: OD-3~(c), \cl{thm:backend}(3),
  \cl{obl:backend}(iv) unchanged).
\item \cl{rmk:bocksteinexp} \PNEW{} --- \textbf{the F-A remark} beside \cl{cl:bockstein}: the
  Bockstein counts the exponent-one summands, and the Frobenius lower bound (an element of
  order $\ell^{j+1}$ at $n=\ell^{j}m$, $m$ a Dickson degree), both \PROVED{}; the probe's exact
  exponents at $(30,5)$, $(60,5)$, $(150,5)$ and the rest of its sweep,
  \COMP\textsuperscript{a}. \textbf{A remark, not a claim change}: \cl{cl:bockstein} is
  untouched and nothing beyond ``$\ge$'' is asserted.
\item \ref{H3} --- \textbf{amended}: ``no comparison map between the two is asserted here, in
  either direction'' becomes ``no isomorphism \ldots\ is asserted there'', with the pointer to
  \cl{prop:bredon} and $\cmpI^{\ast}$. \textbf{Narrows a negative statement to what stays
  true}; every hypothesis the comparison statements carry, $6\in\cring^{\times}$, is unchanged.
\item \cl{thm:specialisation} --- its closing ``Not asserted'' paragraph gains the pointer
  that the deferred construction is no longer deferred for $(\Xbulk,\Xbdry)$ itself. Clauses
  (1)--(4) and their tags are unchanged.
\item \cl{sec:notation-integral} --- new subsection of the notation chapter (normative under
  \ref{CL5}): the symbols of \S\ref{sec:integral} and how the design notes' symbols read here;
  \file{gl2\_preamble.tex} gains the seven macros behind them.
\end{enumerate}

\paragraph{Non-normative.}
\begin{itemize}\setlength{\itemsep}{0pt}
\item \cl{rmk:orbifold} --- its design-direction paragraph records the answer (a
  direction nothing depends on, \ref{CL5}); \cl{rmk:mpcaveat}(4), \cl{rmk:orbifoldrefine}
  and \cl{sec:orbifoldmodel}'s opening paragraph gain pointers.
\item \S\ref{sec:scope}'s reading guide names Part~VI.
\item \cl{sec:codemap} --- a row for \cl{sec:integral}, and the \cl{sec:six} /
  \cl{sec:orbifoldmodel} row's last sentence, which this version falsifies (C-A5 no longer
  fixes a refinement); \cl{sec:notimplemented} --- a v1.8.0 paragraph, for the same reason.
  The fuller status is the slot close's.
\item Appendix~\ref{app:status} (twenty lines, the counts), Appendix~\ref{app:formalization}
  (one item), Appendix~\ref{app:labelindex} (the new labels), this appendix's provenance
  table and roster, the header comment and the date line. Containers.
\item \textbf{Outside this document} (\ref{CL2}): none. The extension plan records the
  re-scoping of C-A5, E9 and R20 at its next revision (the slot close); the companion
  specification is not touched.
\end{itemize}

\paragraph{Not done here, deliberately.}
\begin{enumerate}\setlength{\itemsep}{0pt}
\item \textbf{The arbitrary-precision ban is not re-scoped} --- \cl{thm:backend}(3), the box
  of \S\ref{sec:fpfilter}, \cl{obl:backend}(iv): owner decision OD-3~(c).
\item \textbf{No record-schema rule for the value naming $\Z$} --- the keys are the companion
  specification's; \cl{sec:recordschema} is not amended.
\item \textbf{The pair and the cone over $\Z$} beyond \cl{prop:strictlift}(1)'s datum, and
  any obligation form of \cl{prop:certifiede} --- a later revision's, with the code that
  needs them.
\item \textbf{Versions 1.7.2--1.7.4} (the three status-only slot closes) carry no subsection
  of this appendix; their record is the header comment and the date line. Found while this
  entry was written; not repaired here, since reconstructing another version's entry is the
  slot close's call.
\end{enumerate}


\subsection*{v1.7.4 --- 2026-09-23 --- hecke-extensions slot 7, WP-14, the slot close: status only --- a \texorpdfstring{$Y$}{Y} bump}

\noindent\emph{This version} (entry reconstructed at v1.8.1 from the date line, the header comment
and the commit; nothing is added that those do not say). \textbf{No normative predicate, no
hypothesis, no displayed formula, no tag and no label moves.} Close report:
\file{docs/implementation/slot7/WP-14\_report.md}. Commit: \texttt{cc166d4}.

\paragraph{Status (non-normative).}
\begin{itemize}\setlength{\itemsep}{0pt}
\item \cl{sec:codemap} and \cl{sec:notimplemented} --- phase E7 is \textbf{accepted}:
  \cl{cl:pd}(1)--(2) are \emph{tested} on all $54$ cells of the extension stage's case list at
  $\ell\ge5$ (slot-7 WP-1) and fail on none; \cl{cl:pd} stays open. The extension record of
  \cl{sec:recordschema} item~6 is \code{hecke\_ext/v3} (slot-7 stage~3), its new keys being [E]'s
  field inventory, so no rule of item~6 moves. Ten mathematical keys are registered.
\item The registered-validation counts re-derived: $106\to108$ registered, $60\to61$
  \code{Hecke*}. Companion pointer [E] v1.8.2. The header comment and the date line.
\end{itemize}


\subsection*{v1.7.3 --- 2026-09-22 --- hecke-extensions slot 6, WP-14, the slot close: status only --- a \texorpdfstring{$Y$}{Y} bump}

\noindent\emph{This version} (entry reconstructed at v1.8.1, as above). \textbf{No normative
predicate, no hypothesis, no displayed formula, no tag and no label moves.} Close report:
\file{docs/implementation/slot6/WP-14\_report.md}. Commit: \texttt{41792ac}.

\paragraph{Status (non-normative).}
\begin{itemize}\setlength{\itemsep}{0pt}
\item \cl{sec:codemap} --- \cl{sec:contragredient} is \textbf{implemented} (slot-6 WP-1: the
  contragredient, $\Phi_n$, the coefficient pairing, $H^{\ast}(D_n)$ and $A(D_n)(T_p)$);
  \cl{obl:coneexact}(vii) is \textbf{implemented in all three sentences}: the strict-pair route
  (WP-4), route~3 (WP-5) and the Poincar\'e-duality gate for $X_3$ at $\ell\ge5$ (WP-3), which
  \cl{cl:pd} still only \emph{tests}. The presentation and atlas units and identity-certificate
  items 6 and 7 are named.
\item \cl{sec:notimplemented} --- one chapter unimplemented (\S\ref{sec:orbifoldmodel}), none
  partial except \cl{obl:indep}'s clause (iii).
\item The registered-validation counts re-derived by \code{--list}: $100\to106$ registered,
  $55\to60$ \code{Hecke*}. Companion pointer [E] v1.7.2. The header comment and the date line.
\end{itemize}


\subsection*{v1.7.2 --- 2026-09-22 --- hecke-extensions slot 5, WP-13, the slot close: status only --- a \texorpdfstring{$Y$}{Y} bump}

\noindent\emph{This version} (entry reconstructed at v1.8.1, as above). \textbf{No normative
predicate, no hypothesis, no displayed formula, no tag and no label moves.} Close report:
\file{docs/implementation/slot5/WP-13\_report.md}. Commit: \texttt{53563c3}.

\paragraph{Status (non-normative).}
\begin{itemize}\setlength{\itemsep}{0pt}
\item \cl{sec:codemap} records what slot~5 shipped: the record is \code{hecke\_ext/v2}, and the
  boundary, pair and cone complexes are emitted objects (WP-1); the cone's $Z/B/H/P$ and the cone
  certificate's span inclusions run on the tape, and its item~11 evaluates \eqref{eq:conedims}(1)
  on its own guard (WP-4, WP-5), both value-preserving; \cl{obl:coneexact}(iii)'s $r$-half is
  asserted at every $\ell$ (WP-4), and the cone action's certificate gains item~12, record
  honesty (WP-12a, owner ruling Q2); \cl{obl:iota}(ii)'s rider gets an emitted record (WP-12a,
  owner ruling Q3); \cl{sec:recordschema} item~5's refusal is a \code{verify()} code (WP-11a);
  the $\rho_n$ memo, the Dickson multiplication and the canonical field are named.
\item \textbf{Decided}: the owner's reading of \cl{obl:taubdry}(vii) left open by v1.7.1 ---
  Reading~A, a distinct weaker clause, so v1.7.1's text stands as drafted.
\item The registered-validation counts re-derived by \code{--list}: $94\to100$ registered,
  $49\to55$ \code{Hecke*}. The header comment and the date line.
\end{itemize}


\subsection*{v1.7.1 --- 2026-09-22 --- hecke-extensions slot 5, WP-12: the independent referee pass on v1.7.0 --- pointers and self-contradictions only, a \texorpdfstring{$Y$}{Y} bump}

\noindent\emph{This version.} The slot's independent referee pass on v1.7.0 and on [E] v1.6.0
(\file{docs/implementation/hecke\_extensions\_slot5\_plan.md} \S3 WP-12), written by an agent
that did not write v1.7.0, under the operating rule of \code{[OD]} slot-4 close-out~4: a
referee may repair a self-contradiction the version under review itself introduced, including
in a status sentence, provided the repair is pointers-only; anything wider is a finding for the
close. \textbf{No normative statement, hypothesis, displayed formula, asserted predicate, tag
or label moves}, so Appendix~\ref{app:status}'s counts are unchanged. Report:
\file{docs/implementation/referee\_report\_2026-09-22\_contract\_v1\_7\_0.md}, with the full
numbered finding list and the questions routed to the owner; package report
\file{docs/implementation/slot5/WP-12\_report.md}. Commit: \texttt{5d1c083}.

\paragraph{Repairs (non-normative; five of the six are of v1.7.0's own text).}
\begin{enumerate}\setlength{\itemsep}{0pt}
\item \cl{sec:recordschema} item~6(f) --- \textbf{rule (i) was bound and unbound at once}. The
  item's preamble binds ``the nine below --- (a)--(f) since v1.6.0, (g)--(i) with \code{v2}'',
  and [E] \S7 says item~6(i) ``now makes'' the matching rule a rule; (f), which scopes what
  this document binds, still read ``(a)--(e) and (g)--(h)''. It now reads ``(a)--(e) and
  (g)--(i)''. (At v1.6.2 the list ended at (f) and the sentence was correct; v1.7.0 appended
  (g), (h) and (i) and updated (f) for two of the three.)
\item \cl{rmk:conerestrictsmallell} --- \textbf{``load-bearing for no clause'' two lines before
  the clause it is load-bearing for}. The remark's new opening said it is load-bearing for no
  clause and then, in the same sentence, that it is ``still the record that discharges
  Obligation~\ref{obl:taubdry}(vii) at $\ell\in\{2,3\}$''. The v1.7.0 changelog entry has the
  qualifier the body lost; the body now reads ``load-bearing for no clause of
  \S\ref{sec:cone}''.
\item \cl{obl:taubdry}(vii), the asymmetry paragraph --- \textbf{a premise this clause's own
  statement falsifies}. The paragraph rested on ``the data \emph{this} chapter stores are
  $\taub$ and $\Thb_g$ alone''; but (vii) reads the bulk $\Tdc{g}$ and the restriction
  $\rres$, and says so itself (``the only assertion that ties \eqref{eq:bdrychain}'s choice of
  coset system to the bulk''). The distinguishing datum is $\etapair$ and $\hder$, not the
  bulk operator, and the paragraph now says that. \textbf{The conclusion is unchanged} --- the
  proof of \cl{obl:coneexact}(iii) is unavailable without $\hder$ --- so no predicate moves.
\item \cl{obl:taubdry}(vii), the opening sentence --- \textbf{``belongs to the pair and is
  stated in full at \cl{obl:coneexact}''}, which was exactly true while the two clauses had
  the same content and which v1.7.0 falsified by strengthening (iii) alone. Aligned with
  v1.7.0's own new wording in (iii) (``the pair's form of Obligation~\ref{obl:taubdry}(vii)''):
  (vii) is the boundary datum's form of a statement expounded in full at
  \cl{obl:coneexact}(iii). \textbf{Open for the owner}, below.
\item \cl{obl:coneexact}(iii), the proof --- \textbf{$z$ is named}. The proof read ``let $z$ be
  a stored equivariant cocycle''; \emph{this document uses ``stored'' of the flat matrices as
  well} (clause (i)'s own v1.6.0 paragraph: ``the stored matrices are honest linear maps of the
  flat cochain spaces''), where clause (i) is recorded and false in $395$ of $424$ degrees, so
  the adjective carrying the argument was ``equivariant'' alone. It now reads $z\in\Zq{q}$,
  ``an element of the span of the stored $\Ceq$ basis''. This is WP-11's own residual
  finding~5, which asked for exactly one more word if a referee found the reading.
\item \cl{sec:recordschema} item~5, the implementation note --- \textbf{the refusal landed}.
  v1.7.0 wrote the rule and recorded that the tree did not yet enforce it; hecke-extensions
  slot-5 WP-11a added \code{HeckeExtRecord::verify()} code~22 and inverted
  \code{HeckeExtRecordValidation}'s arm~13 to assert it, with a positive control on a source
  complex. The note records that, and the clause is no longer ahead of the code. (The v1.7.0
  changelog entry above is left as written: it was true when written.)
\end{enumerate}

\paragraph{Open for the owner --- not repaired here.}
\begin{enumerate}\setlength{\itemsep}{0pt}
\item \textbf{Is \cl{obl:taubdry}(vii) a distinct clause or a pointer to
  \cl{obl:coneexact}(iii)?} Repair~4 takes the first reading, because v1.7.0's own (iii) calls
  itself ``the pair's form of'' (vii) and its asymmetry paragraph is built on the two binding
  different implementations. The second reading --- that (vii) always was (iii), named in the
  boundary chapter for signposting --- is equally available from the pre-v1.7.0 text, and under
  it v1.7.0 strengthened (vii) too, so (vii)'s carve-out sentence and (iii)'s whole asymmetry
  paragraph come out. \emph{Which reading is normative is an owner choice}; this version makes
  the document consistent under one of them and records the other.
\item \textbf{\code{hecke\_cone\_action/v1}'s \code{restriction\_equivariance\_asserted} is
  \code{false} at $\ell\in\{2,3\}$}, being the stored \code{six\_invertible}, although
  slot-5 WP-4 made the $\rres$-half assertion unconditional in $\ell$ and this document asserts
  it at every $\ell$ since v1.7.0. A record that says ``not asserted'' where the run asserted is
  the mirror of what rule~\ref{A6} forbids. A code change, so not this pass's.
\item \textbf{\cl{obl:iota}(ii)'s rider asks every implementation to ``record which rows it ran
  the exact half on''}; the shipped \code{IotaCertificate} carries it (per-row
  \code{has\_exact}, plus \code{exact\_relation\_data}), but no \emph{emitted} record does ---
  $\inv$ is in neither \code{hecke\_report/v2} nor \code{hecke\_ext/v2}. Either the rider means
  the certificate, or it is an open implementation obligation of the same kind item~5's was.
\end{enumerate}


\subsection*{v1.7.0 --- 2026-09-22 --- hecke-extensions slot 5, WP-11: \texorpdfstring{\cl{obl:coneexact}}{obl:coneexact}(iii) at every \texorpdfstring{$\ell$}{ell} with its proof, the \texorpdfstring{\code{hecke\_ext/v2}}{hecke ext v2} schema rules, and four deferred items --- an \texorpdfstring{$X$}{X} bump}

\noindent\emph{This version.} The slot's one $X$ revision
(\file{docs/implementation/hecke\_extensions\_slot5\_plan.md} \S3 WP-11), written on the
findings of the packages accepted when it started --- slot-5 WP-4 (the cone contract
tightenings), WP-1 (the \code{hecke\_ext/v2} schema) and WP-10 (the console) --- and on the
owner's rulings of 2026-09-22. \textbf{The one clause that moves an asserted predicate is
\cl{obl:coneexact}(iii)}: its $\ell\in\{2,3\}$ carve-out on the $\rres$-half is dropped and
the identity is asserted at every $\ell$. \textbf{No label is added, renamed or removed, and
no tag moves}, so Appendix~\ref{app:status}'s counts are unchanged. Status sentences --- code
maps, the validation counts, \cl{sec:notimplemented} --- are \textbf{not} this version's and
stay with the slot close, the exceptions being the five places this version's own text
falsifies: \cl{rmk:conerestrictsmallell}'s opening and closing sentences,
\cl{obl:taubdry}(vii)'s last sentence, \cl{obl:coneexact}(v)'s division of labour between
certificate and validation arm, the obligation roster's \cl{obl:coneexact} row, and
\cl{sec:recordschema} item~6's ``implemented at'' paragraph. Report:
\file{docs/implementation/slot5/WP-11\_report.md}. Commit: \texttt{2390850}.

\paragraph{Normative.}
\begin{enumerate}\setlength{\itemsep}{0pt}
\item \cl{obl:coneexact}(iii) --- \textbf{the $\rres$-half is asserted at every $\ell$, and
  the proof is written in}. $\rres^{q}\Amat{q}=\Abdry{q}\rres^{q}$ aborts on failure with no
  hypothesis on $\ell$; clause~(i), itself unconditional in $\ell$ and aborting, gives
  $\Thb_g\rres z-\rres\Th_g z=(d\hder+\hder d)z=d(\hder z)$ on a stored equivariant cocycle,
  a coboundary, whence the identity on cohomology. \textbf{Widens an asserted predicate} ---
  owner ruling \code{[OD]} slot-4 close-out 3, Option A, which also asks that the implication
  stop being re-derived and become contract text. Two further paragraphs are part of the
  clause and not commentary: \emph{what it does not say} --- at $\ell\in\{2,3\}$ the assertion
  is about $\Heq{\bullet}(K,\Rel;L)$, $H^{\bullet}(C^{\bullet}_{\partial})$ and the stored
  maps, and \cl{rmk:conetorsion} keeps those separate from $H^{\bullet}(\Ga;\Mn)$, so the
  clause neither asserts nor licenses the group-cohomological reading there --- and
  \emph{the asymmetry} with \cl{obl:taubdry}(vii), below. Discharged as written by
  \code{HeckeConeActionValidation}, whose A6 arm re-derives the identity independently of the
  certificate and, since slot-5 WP-4, \emph{asserts} at every $\ell$ what it previously only
  measured: $424/424$ $(\text{cell},q)$ pairs, the $74$ operators at $\ell\in\{2,3\}$
  included; and by \code{ConeActionCertificate}, whose $\rres$-half guard lost its
  $6\in\cring^{\times}$ conjunct in the same package.
\item \cl{obl:taubdry}(vii) --- \textbf{keeps its carve-out, and gains the reason}. The
  boundary chapter's stored data are $\taub$ and $\Thb_g$ alone: there is no $\etapair$ and
  no $\hder$, so item~1's proof is unavailable to an implementation of the boundary datum,
  which is what that obligation binds. The two clauses bind different implementations and
  therefore do not contradict each other; for one that discharges both, the identity is
  asserted at every $\ell$ under (iii) and this carve-out is vacuous. \textbf{No predicate
  moves here.} Discharged as before by \code{HeckeBoundaryOperatorValidation} ($216/216$
  recorded at $\ell\in\{2,3\}$).
\item \cl{rmk:conerestrictsmallell} --- \textbf{demoted from the clause's authority to the
  measurement that preceded the proof}. It is \COMP\textsuperscript{r}, which never upgrades
  and never disappears (\ref{CL4}, \S\ref{sec:status}), so the remark stays, its title now
  saying what it is; it remains the record that discharges \cl{obl:taubdry}(vii) at
  $\ell\in\{2,3\}$, and is load-bearing for no clause of \S\ref{sec:cone}. Its closing
  sentence, which said the $424/424$ were ``recorded there and not asserted'', is corrected:
  they are asserted, on the strength of clause (i) and not of this record, which licenses
  nothing (rule~\ref{A6}). \textbf{Every recorded pair is unchanged.}
\item \cl{obl:coneexact}(v) --- \textbf{the unconditional first identity of
  \eqref{eq:conedims} is the certificate's too}. The clause's closing sentence assigned the
  second identity to the stored object's certificate and the first to ``the validation arm
  that carries it''; it now names \emph{both}, each identity under its own minimal guard.
  \textbf{No predicate moves} --- the clause already required the first identity
  unconditionally --- but the division of labour it states is now the shipped one:
  \code{ConeCertificate}'s item C11 is split in place, keeping the item code, and the split
  is shown non-vacuous on a constructed cell (owner ruling \code{[OD]} slot-4 close-out 2,
  Option A; slot-5 WP-4 \S3). Discharged by \code{ConeCertificate} item C11 and by
  \code{HeckeConeValidation}'s G5 arm and its new negative control.
\item \cl{sec:recordschema} item~5 --- \textbf{the names of the complexes a complex is
  derived from are the record's own derived flag}. Non-empty exactly on a derived complex;
  consequently \code{q\_max\_source = "set"} on a complex with a non-empty list is
  \emph{non-conforming} and a verifier that checks internal consistency refuses it.
  \textbf{Adds a refusal, adds no field} (v1.5.0 already required the names). This closes
  slot-5 WP-1 finding~3, which that package deliberately left open because the choice changes
  what a conforming record may say; the alternative --- that the rule binds only the builder
  --- is rejected because it leaves the violation undetectable from the record, which is
  what this section exists to prevent. \code{HeckeExtRecordValidation} pins the gap today,
  in its negative form, and is where the refusal becomes an assertion.
\item \cl{sec:recordschema} item~6 --- \textbf{\code{hecke\_ext/v2}}, with nine binding rules
  in place of six. (b) is amended: in the extension family a verdict that is a property of
  the whole run is carried \textbf{at record level, once, with a single \code{resumable}},
  and the incomplete-item count is then a count of \textbf{decisions} --- v1.6.0's rule
  stands for a schema that does replicate such a verdict, and the change of an existing key's
  \emph{value} is why this is a version bump and not an amendment in place, made before the
  first emission (owner ruling \code{[OD]} slot-4 close-out 5). (c) gains the rule one level
  down: a key the schema requires but this record has not computed is \code{null}, never a
  default and never dropped --- dropping makes a closed list a function of the cell, and a
  \code{false} on an unevaluated diagnostic destroys rule~\ref{A6}'s distinction (slot-5 WP-1
  findings 2 and 4). (g): a record about several complexes reports the \emph{maps} between
  them, per spot, with \cl{obl:coneexact}(iv)'s notion of ``reported'' governing which spots
  exist and (c)'s \code{null} governing the rest. (h): a $\pm$-split record is two further
  records of the same cell, each ideal carrying its record's sign, and the per-sign
  dimensions \textbf{sum} to the unsplit ones --- \emph{asserted}, being
  \cl{prop:iotasplit}(1) itself. (i): a matching a record performed is a matching it reports,
  with its witness; a count is not a report. Discharged by \code{HeckeExtRecordValidation}
  and \code{HeckeExtRamanujanValidation}.
\item \cl{sec:recordschema} item~6(a) and (d) --- \textbf{the two halves of v1.6.0's deferred
  item~3, answered}. (a): the settings block does \emph{not} name the engine settings it
  excludes; the criterion is that a setting belongs iff it may change a published value
  (owner decision R1) and the \code{ENGINE\_*} family is excluded \emph{by that rule}, since
  enumerating it would make a record a function of the build. (d): the closed list of
  certificate names lives in [E, \S7] with the rest of the field inventory; this document
  binds that it is closed, that the two payload forms are distinguishable by a reader holding
  only the record, and that a payload outside it is non-conforming.
\item \cl{obl:iota}(ii) --- \textbf{the rider for a complex without exact group data}. The
  \emph{fibre} half is asserted unconditionally, always. The \emph{exact} half is asserted on
  every row whose $\gamma$ the complex stores, an implementation whose complex stores them
  may not skip it, and every implementation \textbf{records} which rows it ran it on, so that
  a complex carrying no exact rows discharges the fibre half and records the exact half as
  not evaluated rather than reporting the clause passed. \textbf{Names the data the clause
  reads; excuses nothing.} Deferred item~2 of the v1.6.0 entry (slot-4 WP-3 finding~3).
  Discharged by \code{HeckeInvolution}'s \code{IotaCertificate} (the recorded flag, and the
  \code{verify()} clause that re-derives the exact comparison where the rows are present)
  and by \code{HeckeInvolutionValidation}.
\end{enumerate}

\paragraph{Non-normative.}
\begin{itemize}\setlength{\itemsep}{0pt}
\item The obligation roster's \cl{obl:coneexact} row --- the (iii) cell rewritten to the new
  clause, with the assertion status of the validation's own re-derivation. Falsified by
  item~1.
\item \cl{sec:recordschema} item~6's ``implemented at'' paragraph and its status sentence ---
  \code{hecke\_ext/v2}, slot-5 WP-1, \file{HeckeExt/HeckeExtComplexRecords} named. Falsified
  by item~6. The fuller status is the slot close's.
\item The title page's companion pointer --- \file{hecke\_extensions.tex} v1.5.1 $\to$
  \textbf{v1.6.0}, that companion having been revised in this same package. Container.
\item The header comment, the date line and this appendix. Containers.
\item \textbf{Outside this document} (\ref{CL2}): [E] v1.6.0 --- the \code{det}-twist
  convention of \code{cl:glover}'s preamble recorded as \textbf{the gate on phase E7}, which
  [E] had not said; \S7's two booked schema repairs marked taken at \code{hecke\_ext/v2}; and
  the three schema findings of slot-5 WP-1 that are inventory rather than rule. [P] v3.3 ---
  the two console verbs bound, the C5-early scripts, the counts. [X] and \file{CLAUDE.md} are
  the slot close's.
\end{itemize}

\paragraph{Deferred to the next \texorpdfstring{$X$}{X} revision.}
\begin{enumerate}\setlength{\itemsep}{0pt}
\item \textbf{The $\inv$-equivariance of the derived homotopy} --- a clause of
  \cl{obl:eta} or a counterexample cell. Deferred at v1.6.0 because no shipped validation
  asserts it and one cell is not evidence, and deferred again here: of the packages accepted
  when this revision started, none widened the measurement. Slot-5 WP-1 widened the per-sign
  \emph{dimension} assertion, which is \cl{prop:iotasplit}(1) and not this; $[\inv_{c},
  \Acone{q}]=0$ is still recorded at $(10,691)$ alone (\S\ref{sec:coneaction}'s closing
  paragraph, corrected at v1.6.1), and $\inv\hder=\hder\inv$ is measured nowhere.
\item \textbf{Any normative finding of a package accepted after this revision started.} The
  slot rule routes it to the close's open-point list and to the next $X$, which is what keeps
  the document chain off the critical path.
\end{enumerate}


\subsection*{v1.6.2 --- 2026-09-22 --- hecke-extensions slot 4, WP-11: the slot close --- status only, a \texorpdfstring{$Y$}{Y} bump}

\noindent\emph{This version.} The slot close. \textbf{No normative statement, hypothesis,
displayed formula, asserted predicate, tag or label moves}; every edit below reports what
the tree now contains. The slot itself moved \textbf{no golden value and emitted no golden
entry}: \code{git diff 9ae1d7a} is empty for \file{Golden/GoldenReference.inc}, for
\file{docs/implementation/tau\_table\_p2.txt} and for \file{Hecke/HeckeReport.\{h,cpp\}};
\code{GOLDEN\_STRICT=1} reports tier~A $9446$ failing $0$ and tier~B $4603$ mismatching
$0$; \code{EngineDeterminismValidation} gives digest \code{0xe446b3bd2d28bd6e} on all
eight arms.

\paragraph{Status recorded.}
\begin{itemize}
\item \textbf{The production $\inv^{\ast}$ exists.} \cl{sec:involution}'s 2026-09-16
  census paragraph, its specification paragraph, \cl{sec:codemap}'s row and
  \cl{sec:notimplemented} are rewritten from the minimal pointers v1.6.1 left (the
  referee's B1), naming \file{Hecke/HeckeInvolution} (the morphism, with
  \code{compute()} discharging Obligation~\ref{obl:iota}(i),(ii),(iii) blockwise and
  exactly over $\Z$ and aborting on failure), \file{HeckeExt/HeckeSignSplit} (clause (v),
  $A^{(q)}_{\pm}$, the per-sign decomposition, the Eisenstein line) and
  \code{HeckeInvolutionValidation}, whose acceptance is \emph{entrywise against the two
  validations' independent construction}. Slot-4 WP-3, phase E3.
\item \textbf{The extension record is written, bulk complex.} \cl{sec:recordschema}
  item~6 gains that status and names \file{HeckeExt/HeckeExtDriver}; \cl{sec:codemap}'s
  level (b),(c) row gains \file{HeckeExt/HeckeExtRecord}, \file{HeckeExtReport},
  \file{HeckeExtFinal}, the settings keys \code{commutant\_budget} and
  \code{final\_mode} (six keys, defaults and fingerprint $0$ unchanged), and the first
  two E6 units \file{HeckeExt/HeckePredictions} and \file{HeckeExt/GloverFactors}, which
  realise no statement of this document. Slot-4 WP-1, WP-2, WP-7.
\item \textbf{Chapter~\ref{sec:cone}'s cost work.} The row records WP-4's three
  re-implementations and states explicitly that \emph{no clause was sampled, relaxed or
  moved out of \code{compute()}}: both certificates still run every item unconditionally,
  $A_c^{(q)}$ is bit-identical and both validations' printed output is byte-identical to
  the pre-slot-4 binary's. The cost \emph{gate} is [X]'s E5 acceptance (d), not this
  document's --- [C] states no cost bound anywhere --- and WP-4's finding that (d) is not
  decidable as written is recorded there.
\item \textbf{The citation census.} Remark~\ref{rmk:labelcensus} records the
  content-keyed allow-list (\code{relpath|<fragment>}, a stale entry failing, the retired
  \code{relpath:line} form rejected by name) and the slot-4 run: $429$ labels defined,
  $1575$ sites in $128$ of $470$ files, $4$ allow-listed, $0$ unresolved, $0$ stale.
  Slot-4 WP-6.
\item \textbf{Counts}, taken from \code{--list} and not asserted: $89\to94$ registered
  validations, $45\to49$ \code{Hecke*}. Appendix~\ref{app:status}'s own row counts are
  unchanged, no statement having moved.
\end{itemize}

\paragraph{Left for the owner, decided nowhere here.} The three questions v1.6.1's
referee raised and did not answer stand open, verbatim: \textbf{(1)} \code{ConeCertificate}
item C11 gates \emph{both} identities of \eqref{eq:conedims} on a stored $\Heq{2}=0$ while
clause (v) says the first may not be conditional --- the clause is not false (it assigns
the unconditional identity to the validation arm, and \code{HeckeConeValidation}'s G5 arm
does assert it on all $74$ cells), but splitting C11 so the certificate itself carries it
is a \emph{code} change and is slot 5's; \textbf{(2)} Obligation~\ref{obl:coneexact}(iii)'s
$\ell\in\{2,3\}$ carve-out is weaker than clause (i) already proves, so asserting the
$\rres$-half at every $\ell$ is a candidate for the next $X$; \textbf{(3)} whether the B1
pointer repair belonged to the referee pass at all --- \emph{decided by the slot lead: B1
stays}, minimal and pointers-only, and this version completes it.


\subsection*{v1.6.1 --- 2026-09-22 --- hecke-extensions slot 4, WP-9: referee pass on v1.6.0 --- two recorded counts corrected, three self-contradicting status sentences pointed, a \texorpdfstring{$Z$}{Z} bump}

\noindent\emph{This version.} The independent referee pass the slot-4 plan puts after the
document package (\file{hecke\_extensions\_slot4\_plan.md} \S3 WP-9), run on \cl{sec:cone},
\cl{sec:coneaction}, \cl{sec:involution}, \cl{sec:recordschema} and
Appendix~\ref{app:status} from the diff, the source tree and the validations.
\textbf{Changes no normative statement} (\ref{V2}, \ref{CL5}): no hypothesis, displayed
formula, asserted predicate, tag or label moves, and Appendix~\ref{app:status}'s counts are
unchanged ($117/88$, \COMP{} $19$) --- they were recounted mechanically from the table and
every one of v1.6.0's numbers is right. \textbf{The referee could not fault} the rewritten
proof of \cl{lem:conerelative}(5), the equivariant reading of \cl{prop:coneambiguity}(2),
\cl{lem:etaexists}(V)'s repair, \cl{obl:coneexact}(i)--(vi) as rewritten or
\cl{sec:recordschema} item~6, and confirmed each of them discharged as written by the
validation that runs it. Report:
\file{docs/implementation/referee\_report\_2026-09-22\_contract\_v1\_6\_0.md}.
Commit: \texttt{8994c35}.

\paragraph{Normative.}
\textbf{None.} \ref{CL5}'s test is met vacuously: the two items below correct \emph{recorded
measurements}, which rule~\ref{A6} makes records and not claims, and the third is a status
pointer.

\paragraph{Corrections to recorded measurements and to status.}
\begin{enumerate}\setlength{\itemsep}{0pt}
\item \cl{rmk:iotapairing} --- \textbf{``thirteen cells'' $\to$ \emph{eight}}. The recorded
  grid is \code{HeckeInvolutionValidation}'s $15$ cells
  $\ell\in\{2,3,5,7,691\}\times n\in\{0,2,10\}$; $(10,691)$ is one of the twelve with
  $\ell\neq2$, so ``$(10,691)$ and the twelve cells with $\ell\in\{3,5,7,691\}$'' counts it
  twice, and the four at $n=0$ --- where $\Hcpt{1}=0$ and \cl{obl:iota}(v)'s $d_k$ is not
  even defined, being stated for even $k\ge4$ --- were folded into the count while the
  sentence described only $n\in\{2,10\}$. The eight cells it describes are the eight it now
  claims, and the remark records that at $\ell=3$ \cl{obl:iota}(v)'s own hypothesis
  $\ell\ge5$ does not hold, so the comparison there is measured against measured. \COMP
  \textsuperscript{r} kept; every recorded pair is unchanged.
\item \S\ref{sec:coneaction}'s closing paragraph --- \textbf{the commutator is recorded at
  $(10,691)$ alone}, not on the $15$-cell grid, whose arm builds $\inv_{c}$, $\Dcone$ and the
  $\pm$ dimensions and no cone Hecke action. The paragraph's own next sentence, ``One cell
  is not evidence'', already said so.
\item \S\ref{sec:involution} and \S\ref{sec:notimplemented} --- \textbf{three status
  sentences that v1.6.0's own new text falsifies} (``no \code{HeckeInvolution} translation
  unit''; the specification paragraph's ``stands unchanged''; \S\ref{sec:codemap}'s ``Both
  are validations: no \emph{production} $\inv^{\ast}$ exists''; and
  \S\ref{sec:notimplemented}'s ``what that chapter still lacks'') each gain a
  \textbf{pointer} to \file{Hecke/HeckeInvolution} and \cl{rmk:iotapairing}. This is the
  exception v1.6.0 made for \cl{sec:recordschema} item~5 and owed here; the rewrite of the
  code maps and of the row itself stays with the slot close
  (\file{hecke\_extensions\_slot4\_plan.md} \S4.3), which is why the pointers say so.
\end{enumerate}

\paragraph{Non-normative.}
\begin{itemize}\setlength{\itemsep}{0pt}
\item The title page --- \textbf{``Companion specification: \file{hecke\_extensions.tex}
  v1.3.3'' $\to$ v1.5.1}, three revisions stale (since [E] v1.4.0, slot 2) and not
  introduced by v1.6.0, but named because v1.6.0 moved [E]'s pointer to this document while
  leaving this document's pointer to [E]. Container.
\item The header comment and the date line. Containers.
\item \textbf{Outside this document} (\ref{CL2}): [E] v1.5.1, the companion repairs ---
  [E]'s \code{cl:harder}, its new predicate paragraph (\emph{false} at its own edge cell
  $(n,\ell)=(0,2)$ and silent on every even $k$ with $(\ell-1)\mid k$, i.e.\ on all of
  $\ell\in\{2,3\}$); \code{cl:glover}(4)'s raw-factor bound (\emph{false}: two factors against
  one at $n=\ell$); and two proof defects in \code{cl:sharpbound}. [P] v3.2's one off-by-one
  (``thirteen keys'' against ``the twelve keys'' at its \S10 C2 (g)) is reported and left,
  [P] being outside this pass's editing scope; [X] and \file{CLAUDE.md} are the slot
  close's.
\end{itemize}

\paragraph{Left for the owner, decided nowhere here.}
\begin{enumerate}\setlength{\itemsep}{0pt}
\item \code{ConeCertificate} item C11 gates \emph{both} identities of \eqref{eq:conedims} on
  a stored $\Heq{2}=0$, while \cl{obl:coneexact}(v) now says an implementation ``may not
  make the assertion conditional'' on it for the first. The clause is not false --- it
  assigns the unconditional identity to the validation arm, and
  \code{HeckeConeValidation}'s does carry it on all $74$ cells --- but splitting C11 is a
  code change, and is not this pass's to make.
\item \cl{obl:coneexact}(iii)'s new $\ell\in\{2,3\}$ carve-out on the $\rres$-half is
  weaker than clause~(i) now proves: on an equivariant cocycle, (i) gives
  $\rres\Th_gz-\Thb_g\rres z=-d(\hder z)$, a coboundary, at \emph{every} $\ell$. The
  carve-out is what v1.5.2's deferred list proposed and the owner pre-approved, so it
  stands; asserting the $\rres$-half at every $\ell$ is a candidate for the next $X$.
\item Whether item~3 above belongs to this pass at all or wholly to the slot close.
\end{enumerate}


\subsection*{v1.6.0 --- 2026-09-21 --- hecke-extensions slot 4, WP-8: the eight items deferred at v1.5.2, and the normative findings of the accepted slot-4 packages --- an \texorpdfstring{$X$}{X} bump}

\noindent\emph{This version.} The $X$ revision that v1.5.2's \emph{Deferred to the next
$X$} list asks for, run in slot~4 rather than slot~5 by owner confirmation
(\file{docs/implementation/hecke\_extensions\_owner\_decisions.md}, ``Slot 4
confirmations'' of 2026-09-21, item~2; \file{hecke\_extensions\_slot4\_plan.md} \S3 WP-8).
\textbf{All eight items are resolved}, the three that needed a real choice as the owner
answered them --- item~1(a) \emph{yes}: \eqref{eq:conedims}'s first identity is asserted
on every computed cell, without consulting $\Heq{2}$, so that a cell with $\Heq{2}\neq0$
would be covered too (none occurs on the shipped grid);
item~1(b) \emph{yes}: the explicit $\im=\ker$ check is required in every degree; item~1(c)
\emph{per mode} (an [E] item, not one of this document's). Together with them, the
normative findings of the slot-4 packages accepted while this revision ran: WP-1's
record-schema findings (\file{docs/implementation/slot4/WP-1\_report.md} findings 1--3,
5), WP-3's involution findings (\file{WP-3\_report.md} findings 1 and 5) and WP-2's two
record-schema rules (\file{WP-2\_report.md} findings 4 and 5). \textbf{WP-4's findings are
not this document's}: E5 acceptance (d) is a clause of the implementation plan [X], not of
this contract, which states no cost bound anywhere; its two defects --- (d) pins neither
the denominator's measurement protocol nor which E5 object the numerator is --- go to the
slot close with the measured numbers.
\textbf{No label is added, renamed or removed}; one tag changes (\cl{rmk:iotapairing}
gains \COMP\textsuperscript{r}) and Appendix~\ref{app:status}'s counts move with it, to
$117/88$ and \COMP{} $19$ ($14$ registered). Status sentences --- code maps, the
obligation roster, \cl{sec:notimplemented}, the validation counts --- are \textbf{not}
this version's and stay with the slot close (\file{hecke\_extensions\_slot4\_plan.md}
\S4.3); the one exception is the half-sentence of \cl{sec:recordschema} item~5 that said
the extension record ``is not yet written'', which the new item~6 replaces and which
would otherwise contradict this version's own text. Report:
\file{docs/implementation/slot4/WP-8\_report.md}. Commit: \texttt{ec92e3a} (\ref{CL6},
back-filled at v1.6.1).

\paragraph{Normative.}
\begin{enumerate}\setlength{\itemsep}{0pt}
\item \cl{obl:coneexact}(i) --- \textbf{``as matrices'' is named: on the stored $\Ceq$
  basis}. The cone identity $d\hder+\hder d=\Thb_g\rres-\rres\Th_g$ and
  $\Dcone\Phig=\Phig\Dcone$ are asserted there, $\Dcone^{2}=0$ everywhere; the same two
  identities on the full flat cochain space are \textbf{recorded, never asserted}, the
  coset transport of \cl{def:cosettransport} being defined on $\Hom_{\Z[\Ga]}(-,\VR)$
  ($424/424$ computed degrees on $\Ceq$, $29$ of $424$ flat; $\Dcone\Phig=\Phig\Dcone$
  flat on $1$ of $212$ operators). The qualifier is carried onto clause (viii)(b)'s
  tampered $\hder'$. Deferred item~1 (WP-6 finding~1); discharged as written by
  \code{HeckeConeActionValidation}.
\item \cl{obl:coneexact}(iii) --- \textbf{the $\ell\in\{2,3\}$ carve-out on the
  $\rres$-half}. $\rres^{q}\Amat{q}=\Abdry{q}\rres^{q}$ is asserted for
  $6\in\cring^{\times}$ and \emph{recorded} otherwise, which is the carve-out
  \cl{obl:taubdry}(vii) and \cl{rmk:conerestrictsmallell} already carried and which (iii)
  contradicted by aborting with no exception; the other maps of \eqref{eq:coneles} are
  asserted at every $\ell$. Deferred item~2 (WP-6 finding~2); discharged as written by
  \code{HeckeConeActionValidation} ($424/424$ recorded).
\item \cl{obl:coneexact}(iv) --- \textbf{the explicit $\im=\ker$ check is
  degree-generic}. ``$\im(\Hcpt{1}\to\Heq{1})=\ker\rres^{1}$ is checked explicitly'' is
  replaced by $\im(\coneproj{}_{\ast}^{\,q})=\ker\rres^{q}$, both inclusions by solves, at
  \emph{every} bulk spot all three of whose terms are reported. \textbf{Widens the
  asserted set} --- owner item~1(b), answered \emph{yes}. Deferred item~6 (WP-5
  finding~5); discharged as written by \code{ConeCertificate} item C9, which already ran
  in every fully reported degree, and by \code{HeckeConeValidation}'s own inclusion
  checks.
\item \cl{obl:coneexact}(v) and \eqref{eq:conedims} --- \textbf{the first identity is
  asserted without consulting $\Heq{2}$ at all}, the second where $\Heq{2}=0$ holds in the
  stored data. \textbf{Weakens a hypothesis} --- owner item~1(a), answered \emph{yes}.
  Deferred item~5 (WP-5 finding~3); discharged as written by \code{HeckeConeValidation},
  whose G5 arm asserts the first identity on all $74$ cells of its grid, $\ell\in\{2,3\}$
  included and with no reference to $\dim\Heq{2}$, and by \code{ConeCertificate} item C11
  for the second. \emph{Stated as a dependence, not as a counterexample}: on the shipped
  grid $\Heq{2}=0$ everywhere, at $\ell\in\{2,3\}$ recorded rather than asserted, so no
  computed cell separates the two readings today; the clause says so in as many words
  rather than implying, as a first draft of it did, that the $\ell\in\{2,3\}$ cells are
  cells with $\Heq{2}\neq0$.
\item \cl{lem:conerelative}(5) --- restructured to match: \eqref{eq:conedims} is stated
  for a field, its \textbf{first identity with no further hypothesis} (exactness of
  \eqref{eq:coneles} and $\Hcpt{0}=0$ from clause (4) suffice) and its second under
  $\Heq{2}=0$, which is also the hypothesis of the six-term window. \textbf{The proof is
  rewritten} in two halves accordingly, the first exhibiting the four-term exact sequence
  $0\to\Heq{0}\to H^{0}(\partial)\to\Hcpt{1}\to\ker\rres^{1}\to0$ and the second saying
  where $\Heq{2}=0$ is used. Joint with item~4 above; \PROVED{} kept.
\item \cl{obl:coneexact}(vi) --- \textbf{where the \emph{full} $T_{p^{\nu}}$ comes
  from}: the cone pipeline at $g=\mathrm{diag}(1,p^{\nu})$ emits the primitive operator,
  and the full one is assembled by \eqref{eq:fullassembly}, at $\nu=2$ by adding
  $p^{\,n}I$ (\cl{sec:primfull}, \cl{err:e1}). Deferred item~4 (WP-6 finding~4);
  discharged as written by \code{ConeRelationsCertificate}.
\item \cl{obl:coneexact}'s preamble --- ``the six maps of \eqref{eq:coneles}'' $\to$
  \textbf{the three map families $\coneproj{}_{\ast}^{\,q}$, $\rres^{q}$,
  $\coneinc{}_{\ast}^{\,q}$, stored per degree}, with the eight non-trivial instances on
  the model complexes named and the two zero-dimensional ones marked. Deferred item~7
  (WP-5 finding~11).
\item \cl{prop:coneambiguity}(2) --- \textbf{the realisation is a construction inside
  $\Ceq$}: $C^{1}:=\Cq{1}(K,\Rel;L)$, $B^{1}:=d\,\Cq{0}$, $Z^{1}:=\Zq{1}$ stated
  explicitly in the clause and in its proof, with the reason the flat reading is vacuous
  (on the contractible $\Xbulk$ the flat complex is acyclic in degree~$1$,
  \eqref{eq:contractibleuse}). Deferred item~3 (WP-6 finding~3, joint with item~1).
  \PROVED{} kept.
\item \cl{lem:etaexists}(V) --- ``reduces to the single edge path'' $\to$ ``reduces to
  \emph{the one path joining the two vertices}'', with the reason (a stored path template
  may be several edges long). Deferred item~8 (WP-9 \S7 item~4). The statement of (V) is
  untouched, as it was at v1.5.0.
\item \cl{sec:recordschema} --- \textbf{new item~6, the extension record}
  \code{hecke\_ext/v1}: a second writer and not an extension of the per-degree one
  (whose text is a golden object); and its six binding rules --- (a) the settings block
  carries value, \code{source} and the fingerprint of the non-default entries, which
  enters the record's hashes and the writing job's identity, so that \emph{how} a value
  was set is itself value-bearing; (b) a budgeted item's \code{resumable} object, and the
  sidecar, whose name must be \textbf{injective} across several states of one kind and
  whose reference is \textbf{a path relative to the directory holding the record}; (c) an
  unavailable complex is the placeholder \code{\{"status": "unavailable"\}} and its
  degree fields are not emitted; (d) a non-final verdict's reason list, with
  \code{certificate\_failed:<code or certificate name>} --- a \emph{closed name list}
  beside the clause index, because a certificate route may answer ``no'' with no
  \code{verify()} having failed --- and \code{phase\_reached} per item beside the run's
  high-water mark under separate keys; (e) the two recorded cone flags
  \code{phi\_flat\_identity} / \code{residual\_flat\_zero} of item~1 above, named here so
  that they have one spelling from the start; (f) what this document binds and what is
  [E, \S7]'s. WP-1 findings 1--3 and 5 (\file{docs/implementation/slot4/WP-1\_report.md});
  discharged as written by \code{HeckeExtRecordValidation}, whose structural
  \code{verify()} runs (a)--(e). \textbf{Two further rules from the E2b driver} (slot-4
  WP-2, \file{docs/implementation/slot4/WP-2\_report.md} findings 4 and 5), each normative
  because each decides a published number: in (a), \emph{a resumed run writes the settings
  block a fresh run at the same effective values would write}, so that a resume to the
  default budget is byte-identical to a fresh default run, fingerprint included --- without
  the rule two conforming implementations differ in a published hash (discharged as written
  by \code{HeckeExtRamanujanValidation} arm~4, which asserts that byte identity); and in
  (b), \emph{the count of incomplete resumable items counts record entries}, which may
  exceed the number of distinct incomplete decisions where a schema attaches a run-level
  verdict to each of several sub-objects --- legal, and what the shipped record does with
  its per-ideal $U_{\ell}$ verdict, but a schema that does it must say so (discharged as
  written by \code{HeckeExtRecordValidation}, whose \code{verify()} asserts the count equals
  the number of incomplete entries).
\item \cl{rmk:iotapairing} --- \textbf{the mirror-split prediction becomes a record.} The
  remark gains \COMP\textsuperscript{r} and a closing paragraph: $\inv^{\ast}_{\partial}$
  and $\inv_{c}$ now exist, $\rres^{q}\inv^{\ast}=\inv^{\ast}_{\partial}\rres^{q}$ and
  $\Dcone\inv_{c}=\inv_{c}\Dcone$ are asserted, and on thirteen cells $\Hcpt{1}$ carries
  exactly the mirror of the split \cl{obl:iota}(v) puts on $\Heq{1}$. The \emph{pairing}
  statement stays contingent on \cl{cl:pd}; $\ell=2$'s $(k,k)$ pairs are correctly not
  summed. WP-3 finding~5; discharged as written by \code{HeckeInvolutionValidation}.
\item \S\ref{sec:coneaction}, a new closing paragraph --- \textbf{what this document does
  \emph{not} say about $\inv$ and the cone action.} $[\inv_{c},\Acone{q}]=0$ follows from
  nothing here: beyond \cl{prop:iotacommutes} and \cl{obl:taubdry} it needs
  $\inv^{\ast}_{\partial}\hder=\hder\,\inv^{\ast}$ on the stored $\Ceq$ basis, and neither
  \cl{lem:etaexists}'s construction nor \cl{alg:taubdry}'s vertex policy is built to
  respect $\inv$. The commutator is \textbf{recorded} zero at $(10,691)$ for $T_2$ and is
  \textbf{asserted nowhere}; the statement, or a counterexample, is left for the next $X$.
  WP-3 finding~1, which proposed the recorded form rather than a proved one.
\end{enumerate}

\paragraph{Non-normative.}
\begin{itemize}\setlength{\itemsep}{0pt}
\item Appendix~\ref{app:status} --- \cl{rmk:iotapairing}'s row \Brmk{} $\to$
  \COMP\textsuperscript{r}; the \COMP{} split paragraph ``thirteen of the eighteen'' $\to$
  ``fourteen of the nineteen'', with the v1.6.0 clause; the counts paragraph
  \COMP{} $18\to19$ ($13\to14$ registered), tagged/untagged $116/89\to117/88$, untagged
  remarks $36\to35$. Recounted mechanically from the table: $205$ rows, \PROVED{} $93$,
  \BSTD{} $34$, \COMP{} $19$, \CONJ{} $7$, \OPENSTAT{} $6$. No row added or removed ---
  v1.6.0 adds no numbered statement.
\item \cl{sec:recordschema} item~5 --- ``the extension record of E2 \dots{} is not yet
  written'' $\to$ a pointer to item~6. Container; the item's rule is unchanged.
\item This appendix --- the v1.5.2 entry's \texttt{[hash]} back-filled with
  \texttt{cfc6f0b} (\ref{CL6}); the list of $Y$ bumps under \ref{CL5} is unchanged,
  v1.6.0 being an $X$. The header comment and the date line. Containers.
\item \textbf{Outside this document} (\ref{CL2}): [E] v1.5.0 (its own items (a)--(e), and
  the companion pointer $\to$ v1.6.0) and the console plan [P] v3.2, both revised by the
  same package; [X] and \file{CLAUDE.md} are the slot close's.
\end{itemize}

\paragraph{Deferred to the next $X$ revision (slot 5, with E5 acceptance --- findings no
status bump may act on).}
\begin{enumerate}\setlength{\itemsep}{0pt}
\item $\inv$-equivariance of the derived homotopy: either a clause of \cl{obl:eta}
  requiring $\inv^{\ast}_{\partial}\hder=\hder\,\inv^{\ast}$ where an $\inv$-equivariant
  $\etapair$ has been constructed, or a counterexample cell. One measured cell is not
  evidence (WP-3 finding~1; the recorded form is item~12 above).
\item \cl{obl:iota}(ii)'s exact half asks for the image triple ``as exact $\Mat_2(\Z)$
  data together with stored simplex indices and their signs'', which a local system
  carrying only $\rhon(g_r)$ cannot supply; a rider saying what an implementation on a
  complex without exact group data must do instead --- the fibre half
  $\rhon(\inv)=M_r\rhon(\inv)M_{r'}$, which is what makes $\inv_C$ preserve $\Cq{q}$ ---
  would help a re-implementer (WP-3 finding~3).
\item Whether \cl{sec:recordschema} item~6(a)'s settings block should also name the
  \emph{engine} settings it deliberately excludes, and whether (d)'s closed name list
  belongs in this document or in [E, \S7] alone. Raised by WP-1 finding~5's observation
  that [E, \S7]'s \code{groups} object mixes groups with scalars indistinguishably --- an
  [E] schema-hygiene point that item~6(f) leaves with [E].
\item Any normative finding of a slot-4 package accepted after this revision started
  (\file{hecke\_extensions\_slot4\_plan.md} \S4.3).
\end{enumerate}

\subsection*{v1.5.2 --- 2026-09-17 --- hecke-extensions slot 3, WP-11: the slot close --- status only: strict locality implemented, the cone chapter implemented except \texorpdfstring{\cl{obl:coneexact}}{obl:coneexact}(vii), the \texorpdfstring{$U_\ell$}{U ell} repair and the census in the code maps}

\noindent\emph{This version.} The close of hecke-extensions slot 3
(\file{docs/implementation/hecke\_extensions\_slot3\_plan.md} \S3 WP-11, \S4.3: ``WP-11 acts
only on status''). Every sentence changed here reports what the tree does on 2026-09-17 at
the slot's close (commit \texttt{d65e96f} before this version) and nothing else: the
$U_\ell$ reach repair (WP-1), the Sturm sweep on real blocks (WP-2), strict locality mode
(WP-3), the equivariant cone and its long exact sequence (WP-5), the cone Hecke action
(WP-6), console C1a (WP-7, which realises no statement of this document) and the citation
census (WP-10); the two document packages WP-8 / WP-9 are v1.5.0 / v1.5.1. The slot made
\textbf{no golden re-baseline and no golden emission}: \file{GoldenReference.inc} and
\file{tau\_table\_p2.txt} are byte-identical to the slot's base, and the determinism digest
is unchanged. \textbf{Changes no normative statement} (\ref{V2}, \ref{CL5}): no hypothesis,
displayed formula, asserted predicate, tag or label moves; Appendix~\ref{app:status}'s
counts are unchanged ($116/89$). Report:
\file{docs/implementation/slot3/WP-11\_report.md}. Commit: \texttt{cfc6f0b}.

\paragraph{Normative.}
\textbf{None.} \ref{CL5}'s test is met vacuously.

\paragraph{Non-normative.}
\begin{itemize}\setlength{\itemsep}{0pt}
\item \S\ref{sec:codemap}, the level (b), (c) row --- ``\code{strict} accepted and throwing
  a named not-implemented error until slot 3'', \textbf{false since WP-3 merged}, the one
  outright-false status sentence v1.5.0 and v1.5.1 carried (WP-8 \S6 item~1, WP-9 \S7
  item~3): strict mode implemented (the Frobenius radical, the Berlekamp subalgebra, the
  residue witness over $\bar A$, terminating, byte-identical to budget mode wherever budget
  mode certifies), \code{NotImplementedInThisBuild} surviving as a named error no settings
  value raises; the keys \code{u\_ell}, \code{sturm}, \code{sturm\_unavailable}; the
  extension layer's \file{HeckeExt/HeckeSturmSweep}. Container.
\item \cl{sec:cone}'s code map --- ``\textbf{Not implemented}: the cone \dots\ --- slot 3''
  $\to$ \textbf{implemented except \cl{obl:coneexact}(vii)}: \file{EquivariantMappingCone},
  \file{Hecke/SL2ZBoundaryPair}, \file{Hecke/HeckeConeAction}, \code{EtaVertexPolicy} /
  \code{TauBoundaryVertexPolicy} and \code{table\_fnv64()}, the two validations, the
  measured flat-versus-$\Ceq$ reading of clause (i), the collectors defined and never
  emitted; the v1.4.2 paragraph kept below it as the record, its census sentence in the past
  tense, and a note that the class its item (c) anticipated now exists. Container.
\item \S\ref{sec:obligationroster} --- \cl{obl:eta}: \textbf{partial} $\to$
  \textbf{discharged} ((v) by \code{HeckeConeActionValidation}, the tampered $\hder'$
  failing exactly (v) at $(10,691)$); \cl{obl:coneexact}: \textbf{unimplemented} $\to$
  \textbf{partial} (every clause but (vii), with the measured figures of WP-5 / WP-6); the
  preamble's count $80\to89$; ``Reading the table'': eleven $\to$ twelve rows discharged,
  two unimplemented $\to$ one partial. Container; the obligations' text unchanged.
\item \S\ref{sec:codemap}, the rows \cl{sec:heckedatum} (\code{QuadPointT<I>},
  \code{escalate\_on\_overflow}, \file{HeckeExt/HeckeUEllBudget}, the three $U_\ell$
  validations), \cl{sec:cone} (as above), \cl{sec:six} / \cl{sec:orbifoldmodel} (the
  restriction lifted as \file{Hecke/SL2ZBoundaryPair}) and \cl{sec:obligations} ($45$
  \code{Hecke*} entries); \cl{lem:bounds} / \cl{thm:backend}'s code-map paragraph in
  \S\ref{sec:reduction} (WP-8 \S6 item~4). Containers.
\item \S\ref{sec:notimplemented} --- \cl{sec:cone} ``partial since v1.4.2'' $\to$
  implemented except clause (vii), PD adjointness and the record integration; and
  Appendix~\ref{app:formalization}'s oracle item (``two chapters'' $\to$ one chapter and one
  clause; the count $80\to89$). Containers.
\item \S\ref{sec:machinery} item~2 and \cl{rmk:labelcensus} --- \code{label\_census.sh} is
  written and registered as \code{ContractCensusValidation} (WP-10); the remark gains the
  registered run of the slot close ($429$ labels, $1408$ sites in $121$ files, $0$
  unresolved). Rule container; remark, tag \COMP{} kept.
\item \cl{rmk:conerestrictsmallell} --- the record gains the cone's own run ($424/424$,
  WP-6 finding~2). Remark, tag \COMP{} kept; nothing is asserted that was recorded.
\item \cl{rmk:coneambiguity691} --- the closing status sentence: [X]'s stale $4$ is
  corrected at [X] v1.7, and the measured $\dim\Hcpt{1}=3$ is \code{HeckeConeValidation}'s.
  Remark, tags unchanged.
\item \cl{sec:involution}'s census sentence --- the count $80\to89$; there is still no
  \code{HeckeInvolution} (phase E3 moved to slot 4 by owner confirmation). Container.
\item This appendix --- the v1.5.1 entry's \texttt{[hash]} back-filled with
  \texttt{de697e6} (\ref{CL6}; the commit as the referee made it, cherry-picked onto
  \code{hecke-extensions} as \texttt{ef8d1b4}, exactly as v1.5.0's \texttt{1f19cf5} landed
  as \texttt{821c621}); the list of $Y$ bumps under \ref{CL5} gains v1.5.2. The header
  comment and the date line. Containers.
\item \textbf{Outside this document} (\ref{CL2}): \file{CLAUDE.md} (the suite table
  re-measured, the slot-3 roster and architecture), \file{docs/KNOWN\_ISSUES.md} \S3 (the
  RSS figure reconciled with the frozen model), the implementation plan [X] v1.7, [E]
  v1.4.2 (status: strict mode, E1b, the cone; the [C] pointer $\to$ v1.5.2) and the console
  plan [P] v3.1 (status: C1a implemented).
\end{itemize}

\paragraph{Deferred to the next $X$ revision (slot 5, with E5 acceptance --- findings no
status bump may act on).}
\begin{enumerate}\setlength{\itemsep}{0pt}
\item \cl{obl:coneexact}(i): ``$d\hder+\hder d=\Thb_g\rres-\rres\Th_g$ as matrices'' and
  ``$\Dcone\Phig=\Phig\Dcone$'' read as statements on the stored $\Ceq$ basis; the same
  identities on the full flat cochain space \textbf{recorded}, not asserted (they hold in
  $29$ of $424$ degrees and on $1$ of $212$ operators), the coset transport being defined on
  $\Hom_{\Z[\Ga]}(-,V)$; \S\ref{sec:recordschema} to carry the two recorded flags
  (\code{phi\_flat\_identity}, \code{residual\_flat\_zero}). Normative (WP-6 finding~1).
\item \cl{obl:coneexact}(iii) against \cl{rmk:conerestrictsmallell}: (iii) aborts with no
  $\ell$ carve-out while the remark records the $\rres$-half at $\ell\in\{2,3\}$; add the
  carve-out to (iii). Normative (WP-6 finding~2).
\item \cl{prop:coneambiguity}(2): state $C^{1}:=C^{1}_{\mathrm{eq}}(K,\Rel;L)$ and $B^{1}:=d\,C^{0}_{\mathrm{eq}}$
  explicitly (on the contractible $\Xbulk$ the flat reading is vacuous), and in
  \cl{obl:coneexact}(viii)(b) that $\hder'$ satisfies the cone identity on $\Ceq$. Normative
  (WP-6 finding~3, joint with item~1).
\item \cl{obl:coneexact}(vi) / \S\ref{sec:primfull}: one sentence on where the \emph{full}
  $T_{p^{2}}$ comes from when the cone pipeline runs at $g=\mathrm{diag}(1,p^{\nu})$
  (\eqref{eq:fullassembly}; at $\nu=2$ adding $p^{n}I$). Wording (WP-6 finding~4).
\item \eqref{eq:conedims}: the first identity needs no $\Heq{2}=0$ (exactness and
  $\Hcpt{0}=0$ suffice), while \cl{obl:coneexact}(v) gates both on it --- a weakened
  hypothesis, i.e.\ a strengthening, for the owner (WP-5 finding~3; WP-8 \S4 item~1).
\item \cl{obl:coneexact}(iv): ``$\im(\Hcpt{1}\to\Heq{1})=\ker\rres^{1}$ is checked
  explicitly'' read degree-generically, as implemented --- widens the asserted set (WP-5
  finding~5; WP-8 \S4 item~2).
\item \cl{obl:coneexact}'s preamble, ``the six maps of \eqref{eq:coneles}'': on the model
  eight map instances, two zero-dimensional; a wording nit (WP-5 finding~11).
\item \cl{lem:etaexists}(V): ``the single edge path'' $\to$ ``the one path joining the two
  vertices'' (WP-9 \S7 item~4; not touched by a status bump).
\end{enumerate}
\noindent Not for [C] and listed here only so that none is lost: [E]
\code{sec:localcert}'s hypothesis on a \emph{reported} matrix witness ($w\in A_S$, or
commuting with the family), the sharpness identity $\dim_{\Fl}\bar A=\operatorname{lcm}_p
\deg\pi_p$ on a local block, the mode-dependence of the sub-block order after a split and
\code{phase\_reached} per block versus high-water mark (WP-3 findings 2, 3, 4, 6 --- the
next [E] revision); [P]'s phase column in \S3.5, \code{golden()}'s home, the extension
layer's console bindings and C1b's obsolete \code{strict} \code{expect\_error} line (WP-7
findings 1--4, WP-3 finding~7 --- the next [P] revision, with C1b).

\subsection*{v1.5.1 --- 2026-09-17 --- hecke-extensions slot 3, WP-9: referee pass on v1.5.0 --- the ten resolutions accepted, four arithmetic and two proof defects in the \texorpdfstring{$U_\ell$}{U ell} item repaired}

\noindent\emph{This version.} The independent read of v1.5.0 that
\file{docs/implementation/hecke\_extensions\_slot3\_plan.md} \S3 WP-9 asks for, on the
model of v1.4.1 (\texttt{9c54b91}) and v1.3.1 (\texttt{f69cda0}). Every one of the ten
resolutions was re-derived: \cl{rmk:coneambiguity691}'s arithmetic line by line against
the measured row for $(10,691)$ of \file{docs/implementation/slot3/WP-5\_report.md} \S3
($\dim\Hcpt{\bullet}=[0,3,0]$, $\dim\Heq{\bullet}=[0,3,0]$, $\dim
H^{\bullet}(\partial)=[1,1]$, $\rank\rres^{0},\rres^{1}=0,1$, $\dim H^{1}_{!}=2$) and
against \cl{eq:conedims}, \cl{cl:dimH1}(1)--(2) and
\file{docs/implementation/slot2/WP-7\_report.md} \S3.2--3.3; \cl{obl:coneexact}(iv)'s
per-complex ``reported'' against \ref{open:ell2}'s box, its rider (iii) and the slot-2
owner confirmation~1 (``per complex, never as a maximum'') --- it is the box's own
per-term sentence applied literally, it never asserts a spot with an unreported term,
and where the three $q_{\max}$ agree (the only case the spelled-out clause defines) it
reproduces that clause verbatim on a truncated model and adds, on an untruncated one,
only spots all of whose terms are $0$; the three (T1) helpers, the one conjugation primitive and
the three \code{table\_fnv64()} accessors against the tree (\file{Hecke/HeckeTauTable},
\file{Hecke/HeckeBoundaryDatum}, \file{Hecke/HeckePairHomotopy.h:247,252,253}); the
``MAY omit'' branch of \cl{obl:eta}(iv) against the shipped solve regime (OP2's abort
becomes \emph{unconditional} there, hence is not weakened); \cl{lem:etaexists}(V)'s
routing and the $\etapair$-record's four fields against
\file{Hecke/HeckePairHomotopy.h}; and Appendix~\ref{app:status}'s counts by a mechanical
recount of its table ($205$ rows; $\COMP{}=18=13+5$, \PROVED{} $93$, \BSTD{} $34$,
\CONJ{} $7$, \OPENSTAT{} $6$, tagged/untagged $116/89$ with $10$ untagged obligations ---
every figure as v1.5.0 states it). \ref{L1}: no label was added, renamed or removed by
v1.5.0 or by this version. \textbf{Changes no normative statement} (\ref{V2},
\ref{CL5}): no hypothesis, displayed formula, asserted predicate, tag or label moves; the
tier criteria of \cl{thm:backend} are the same criteria in the same quantities. What was
wrong lay in the $U_\ell$ item --- two arithmetic slips of \cl{thm:backend}(1), one
unproved inference in its item~(iv), two steps of \cl{lem:bounds}(6)'s proof --- with
one over-stated sentence at \cl{obl:oracle} and one stale quotation at \ref{A4} beside
them; each is repaired under its own statement's label (\ref{L3}), and the seven repairs
are listed below. Report:
\file{docs/implementation/slot3/WP-9\_referee\_report.md}. Commit: \texttt{de697e6}.

\paragraph{Normative.}
\textbf{None.} \ref{CL5}'s test is met vacuously.

\paragraph{Non-normative.}
\begin{itemize}\setlength{\itemsep}{0pt}
\item \cl{thm:backend} --- clause~(1), \textbf{two arithmetic corrections}, tag
  \PROVED{} kept, both criteria unchanged. \emph{(a)} ``for every $N\le2^{15}$ with such
  a table'' does not follow from $B\le2^{31}$ and $B=\max(B_{\mathrm{tab}},4N^{2})$: at
  $N=2^{15}$ already $4N^{2}=2^{32}$. The condition is $4N^{2}\le2^{31}-1$, i.e.\
  $N\le23170$ --- which is exactly the reach \S\ref{sec:heckedatum} records in the same
  version, clause~(e) of \cl{obl:tau}'s certification being this criterion, so the two
  sites now agree. \emph{(b)} ``$7B_{\mathrm{red}}^{2}<2^{63}$ \dots\ means
  $N\lesssim4580$'' mixed a criterion with a measurement: the criterion gives
  $N\lesssim2700$, and $4580$ is the \emph{measured} first overflow, which
  \cl{rmk:btabsmall} derives from the intermediate actually formed, $\approx
  B_{\mathrm{red}}^{2}$. Both figures are kept, each with its own provenance;
  clause~(2)'s $1.8\cdot10^{8}$, which \emph{is} derived from the criterion, was what
  exposed the mismatch.
\item \cl{thm:backend} --- item~(iv), \textbf{an unproved inference removed}.
  ``$B_{\mathrm{red}}\le7C^{2}N^{2}$ by \cl{lem:bounds}(6)'' is not what that lemma
  proves: (6b) bounds the \emph{entry} point $z_{0}$, and the reduction's own steps are
  not bounded by their entry norm by any argument in this document (a single $S$ step on
  a point with $\abs z<1$ can multiply $\lVert z\rVert$ by a factor). The item now says
  so and cites \cl{rmk:btabsmall}, where the constancy of $\max\lVert z\rVert$ along the
  reduction is \COMP{} \emph{measured} over $\ell\in\{11,\dots,3617\}$ --- which is all
  the clause needs, its criteria being stated in $B_{\mathrm{red}}$ itself and clause~(3)
  requiring overflow-checked arithmetic regardless. \cl{lem:bounds}(6c)'s ``with
  $\max_z\lVert z\rVert$ of order $C^{2}N^{2}$'' is struck for the same reason and the
  measurement left where it is recorded. Nothing an implementation must satisfy changes.
\item \cl{lem:bounds} --- proof of (6), \textbf{two steps repaired}, statement and tag
  unchanged. \emph{(a)} (6a) bounds one M\"obius step by $7\lVert\gamma\rVert^{2}\lVert
  z\rVert^{2}$, but \cl{thm:backend}(iv) quotes $7\lVert z\rVert^{2}$, with no
  $\lVert\gamma\rVert$ --- and $\lVert T^{-q}\rVert=\abs q$ is not $1$ for a shift by a
  large quotient, so the step needs the two shapes \cl{alg:tau} step~1 actually applies:
  $S$, where $\lVert\gamma\rVert=1$; and the \emph{nearest}-integer shift, where $c=0$,
  $\Delta=1$ and $\abs{u-qw}\le w/2$ collapse the formulas to magnitudes
  $\le2\lVert z\rVert^{2}$. The proof now carries that specialisation. \emph{(b)} the
  attainment argument of (6b) cited $v'=vw\Delta$, which for $\adjg(g)s_j$ is of order
  $C^{2}N$, not $C^{2}N^{2}$; the coordinate that attains the order is $w'$, through its
  summand $D\,c^{2}v^{2}$ with $c$ of order $N$. Also $144.8\,N^{2}=C^{2}N^{2}$ (here
  and at \cl{rmk:btabsmall}) is an equality between $144.8$ and $144$; it is now
  $\approx$ / $\asymp$, and \cl{rmk:btabsmall}'s int64 threshold is derived from the
  measured $144.8$, which is where its $4580$ comes from.
\item \ref{A4} --- its worked example quoted \cl{obl:eta}(ii) in the form v1.5.0
  corrected (``both sides evaluated through the shared (T1) helper''), which was the one
  site of that wording left over: the other five --- \S\ref{sec:codemap},
  \cl{obl:tau}(iii), \cl{alg:tau}(0), \S\ref{sec:heckedatum} and the v0.18.0 entry ---
  are about the $\tau$-table alone, where there \emph{is} one helper
  (\code{eval\_tau\_on\_translate}), and are correct as they stand. Rule, no tag.
\item \cl{obl:oracle} --- one sentence of the v1.5.0 paragraph. ``Off the Kummer loci
  $E$ is a line'' was asserted flatly; it is measured, on the $43$ off-Dickson
  non-Kummer cells with $k\le26$ of slot-2 WP-7, and is [E]'s \code{cl:harder} in
  general. The sentence says so, and the reason no assertion of the obligation moves is
  now stated positively --- (i)--(iii) assert dimensions, \cl{eq:esoracle} and the
  full/primitive distinction, none of which names a class inside $E$ --- instead of
  resting on the line case. Obligation, no tag.
\item Appendix~\ref{app:changelog} --- the v1.5.0 \cl{lem:bounds} line gains \PNEW{}
  under \ref{CL4}'s rule for a preserved label whose statement is strictly generalised
  (clause~(6) is new content with its own proof), with the line saying that
  Appendix~\ref{app:status}'s row stays \BFIX{} and no provenance count moves. That is
  the referee's answer to the open question WP-8's report raised under that head.
  \ref{CL5}'s list of $Y$ bumps gains v1.5.1. The v1.5.0 entry's \texttt{[hash]}
  back-filled with \texttt{1f19cf5} (\ref{CL6}). The header comment and the date line.
  Containers, no tag.
\item \textbf{Outside this document} (\ref{CL2}): nothing. [E] v1.4.1 was read and
  accepted as it stands (the canonical generator exists by Nakayama, is a function of
  the reported data, and no claim of [E] depends on it, so the $Y$ bump is right);
  \file{CLAUDE.md}, \file{docs/KNOWN\_ISSUES.md}, [X] and [P] are the slot close's.
\end{itemize}

\paragraph{The ten resolutions of v1.5.0: accepted.}
All ten stand, eight of them unaltered. The two that moved are
\cl{thm:backend}/\cl{lem:bounds} --- item~7 of the deferred list, whose \emph{mathematics}
is accepted (the gap is real, clause~(6) is the right place to close it, and the site,
the bracket $3617<N<5003$ and the new reach $\ell\le23170$ are re-derived here) and
whose arithmetic and proof carried the six defects above. \cl{obl:coneexact}(iv), which
v1.5.0's entry names as ``the one entry a referee should attack'', is accepted as
written: it is \ref{open:ell2}'s own per-term formulation (``exact at every spot all
three of whose terms are reported''), it is what the owner's slot-2 confirmation~1
decided (``per complex, never as a maximum''), and where the three $q_{\max}$ differ
the single-$q_{\max}$ sentence has no reading at all --- which is the defect the rule
repairs, not a set it widens.

\subsection*{v1.5.0 --- 2026-09-17 --- hecke-extensions slot 3, WP-8: the seven deferred items of v1.4.2, and the normative findings of the accepted slot-3 packages --- an \texorpdfstring{$X$}{X} bump}

\noindent\emph{This version.} The one $X$ revision of hecke-extensions slot 3
(\file{docs/implementation/hecke\_extensions\_slot3\_plan.md} \S3 WP-8, \S2.2; the owner
pre-approved the seven ``as the slot-2 reports propose'' on 2026-09-17,
\file{docs/implementation/hecke\_extensions\_owner\_decisions.md}, ``Slot 3
confirmations'' item~1). It does two things and nothing else. \textbf{(1)} It resolves
the deferred list the v1.4.2 entry below hands it --- all seven, none carried forward.
\textbf{(2)} It applies the \emph{normative} findings of the two slot-3 packages
accepted when it started: WP-1, the $U_\ell$ exact-integer reach repair
(\file{docs/implementation/slot3/WP-1\_report.md}), and WP-5, the equivariant cone
complex and its long exact sequence
(\file{docs/implementation/slot3/WP-5\_report.md}). \textbf{No label is added, renamed or
removed} (\ref{L1}, \ref{L6}): every change is a clause, a paragraph or a sentence
inside a statement this document already had, so \ref{app:status}'s line count is
unchanged and no provenance tag is created. \textbf{Status sentences are not this
version's} (\S3 WP-8 scope): the obligation roster, the code maps,
\S\ref{sec:notimplemented} and \S\ref{sec:codemap} are v1.5.2's, so that a status never
precedes the merge it describes --- with the single exception the deferred list itself
demands, item~7's $U_\ell$ reach record, which \emph{is} the finding. The normative
findings of WP-6 (the cone Hecke action), and any finding that arrives after this
version started, go to the next $X$ revision (slot~5), exactly as slot~2 handed its
seven here. Report: \file{docs/implementation/slot3/WP-8\_report.md}. Commit:
\texttt{1f19cf5}.

\paragraph{Normative.}
\begin{itemize}\setlength{\itemsep}{0pt}
\item \cl{rmk:coneambiguity691} --- \textbf{deferred item~1}, the rewrite. The displayed
  arithmetic was computed under $\rres^{1}=0$ at $(10,691)$, which is false. Rewritten
  from the measured values: $\rank\rres^{0}=0$ and $\coker\rres^{0}$ one-dimensional;
  $\rank\rres^{1}=1$, $\dim\ker\rres^{1}=\dim H^{1}_{!}=2$; hence
  $\dim X_3=1-0+2=\mathbf3$ by \cl{eq:conedims}, confirmed by a direct
  $\dim\Hcpt{1}=3$ (slot-2 WP-7 \code{HeckeBoundaryParabolicValidation}; slot-3 WP-5
  \code{HeckeConeValidation}, which finds the same on all ten Kummer cells), and
  $\dim\Hcpt{2}=0$; the ambiguity
  $\Hom(\ker\rres^{1},\coker\rres^{0})\cong\F_{691}^{2}\neq0$. \textbf{The remark's
  conclusion is unchanged} --- a free cochain-level $\hder$ leaves the $T_p$-action on
  $X_3$ undetermined and the derived one does not --- and it is why the remark exists;
  only two dimensions moved, and the v1.4.2 numbers are quoted at the end of the remark
  so that a reader of the old PDF is not left guessing. The new paragraph ``Why
  $\rres^{1}\neq0$'' states the mod-$\ell$ reading of ``the Eisenstein class'' as the
  joint $1+p^{\,n+1}$-eigenspace $E$ ([E] v1.4.0 \code{cl:harder}) and records that it
  is the one input this document does not prove. Remark; \OPENSTAT{} kept, with its
  reason corrected (it had said ``\S\ref{sec:six} not yet written at this version'',
  which has not been true since v0.19.0), \PROVED{} kept, and \COMP\textsuperscript{r}
  \textbf{added} --- the dimensions are now measured and re-runnable by name. The one
  tag change of this version.
\item \cl{obl:oracle} --- \textbf{deferred item~1}, second half. A new paragraph after
  \cl{eq:esoracle}: the \emph{factor} $X-(1+p^{\,n+1})$ is well defined wherever the
  obligation applies; the \emph{class} the second standard pointer names
  (\cite[Ch.~2--3]{harder3}, a characteristic-$0$ statement) need not survive reduction
  as a line, and at a Kummer cell does not --- read ``the Eisenstein class modulo
  $\ell$'' as the joint eigenspace $E$ throughout this document, \cl{obl:iota}(v) giving
  the two conditions under which a class rather than a space is locatable. \textbf{No
  assertion of the obligation moves}: off the Kummer loci $E$ is a line and the two
  readings coincide. Obligation; no tag change (the measurement is cited to
  \cl{rmk:coneambiguity691}, which carries it).
\item \cl{obl:eta}(ii) --- \textbf{deferred item~2}. ``Both sides evaluated only through
  the shared (T1) helper --- the same helper the construction uses'' named a helper that
  does not exist. Corrected to ``through each table's own (T1) helper, all built on the
  one (P1) split'', with the three helpers and the one conjugation primitive
  (\code{HeckeCosetTables::twisted\_conjugate}) named, and what the clause is really
  about --- construction and certification sharing the derivation --- said in place of
  what it had said. The asserted predicate is unchanged. Obligation clause, no tag.
\item \cl{obl:eta}(iv) --- \textbf{deferred item~3}. ``Behind the explicit setting that
  permits it'' presupposed a setting no implementation need have. Now: a mod-$m$ path,
  \emph{where one exists}, is behind an explicit setting (unchanged), \textbf{and an
  implementation MAY omit the mod-$m$ regime altogether}, in which case a tier failure
  aborts unconditionally and the recorded regime is $\Z$ by construction. The shipped
  implementation is of that kind, and the absence of the setting is conformance rather
  than a gap. \textbf{OP2's abort is not weakened} --- it is made unconditional on the
  omitting branch. Obligation clause, no tag.
\item \cl{lem:etaexists}(V) --- \textbf{deferred item~4}. A sentence after the procedure
  naming the constructive choice the implementation makes for the free vertex chain: the
  pairwise path routing of \cl{alg:tau} step~2, which needs no linear solve and reduces
  to the single edge path when both values are single vertices. \textbf{The statement is
  unchanged} --- (V) still admits \emph{any} $1$-chain with the stated boundary, which
  \cl{lem:etaunique} is what makes the naming harmless. Lemma, \PROVED{} kept.
\item \cl{sec:recordschema} item~5 --- \textbf{deferred item~5}, and WP-5 finding~2. A
  new block, ``Maps, and complexes built from complexes'': a record about a \emph{map}
  ($\rres$, $\etapair$, $\hder$) carries the four fields of its \textbf{source} complex
  and the degree shift is a schema constant fixed by the map's definition, never a
  field; on a truncated free model the map is built and certified in the source degrees
  the box of \ref{open:ell2} builds its source complex in, and reported in the source
  degrees $\le q_{\max}(\text{source})$. A record about a complex \emph{derived} from
  others (the cone) carries \code{q\_max\_source = "derived"} unconditionally with the
  names of its sources, and its $q_{\max}$ / \code{truncated} are the stated functions of
  theirs; \code{"set"} is legal on a source complex only, which is rider~(i) of
  \ref{open:ell2} one level up. Container (\ref{CL4}: no tag), but a normative schema
  rule, hence this list.
\item \cl{obl:eta} preamble --- \textbf{deferred item~6}. ``Together with the stored
  $\tau$- and $\taub$-tables it was built against'' is now ``together with the
  \emph{table hashes} of'' them, with a paragraph fixing what a table hash is (one
  function, \code{table\_fnv64()}, on each of the three tables, the FNV-1a-64 of
  \file{ModArithmetic/Fnv.h}) and what it is for: \emph{identifying and never
  certifying}. Two admissible $\etapair$ differ in the first hash while $\Acone{q}$ does
  not (\cl{lem:etaunique}, \cl{prop:coneindep}(1)), so a hash may be recorded and
  compared but not branched on (\ref{A6}), and clause~(i) --- not a hash --- is what
  ties the table to its certified inputs. The three accessors exist in the tree.
  Obligation preamble, no tag.
\item \cl{lem:bounds} \PNEW{} --- \textbf{deferred item~7}, and the substance of WP-1's finding:
  a \textbf{new clause (6)}, with proof. Clauses (1)--(5) bound the matrices the pipeline
  \emph{stores}; \cl{alg:tau} step~1 also carries a quadratic point $z$ on cleared
  denominators that none of them mentions. (6a) one M\"obius step forms integers of
  magnitude $\le7\lVert\gamma\rVert^{2}\lVert z\rVert^{2}$; (6b) the point step~1 starts
  from satisfies $\lVert z_{0}\rVert\le7C^{2}N^{2}$ with $C=\max_v\lVert
  \mathrm{pos}(v)\rVert$ ($C=12$ here), and the order $C^{2}N^{2}$ is attained; (6c) so
  the tier must be chosen against $\max(B,7C^{2}N^{2})$ and \emph{its} square, not
  against $B$. Lemma, \PROVED{}; the tag is unchanged and the proof carries the new
  clause. \textbf{The bound was not wrong --- it did not exist}, which is why this is a
  new clause and not a repair. (\PNEW{} added at v1.5.1 under \ref{CL4}'s rule for a
  preserved label whose statement is strictly generalised: the tag on this \emph{line}
  is the new clause's. Appendix~\ref{app:status}'s row for \cl{lem:bounds} stays
  \BFIX{} --- the statement as a whole is still a transported-and-repaired one --- so
  no provenance count moves.)
\item \cl{thm:backend} --- \textbf{deferred item~7} continued. A new item~(iv) in the
  closure list (step~1's M\"obius steps, of magnitude $\le7B_{\mathrm{red}}^{2}$ with
  $B_{\mathrm{red}}\le7C^{2}N^{2}$ by \cl{lem:bounds}(6), and \emph{not} bounded by
  $B$); and clause~(1) \textbf{corrected}: it had said that $2B^{2}<2^{63}$ ``covers
  every $N<2^{15}$ outright'', which is false --- at $N=5003<2^{15}$ the certified $B$
  leaves $2B^{2}$ four orders of magnitude inside the tier while step~1 leaves it.
  Clause~(1) now states the two criteria separately, and clause~(2) likewise for the
  $128$-bit tier. The wording ``the ladder is entered per operation'' is item~(iv)'s
  consequence and is what the implementation does. Theorem, \PROVED{} kept; clause~(3)
  (arbitrary precision never required) is unaffected --- the $192$- and $320$-bit widths
  bound $7B_{\mathrm{red}}^{2}$ a priori as well.
\item \cl{obl:coneexact}(iv) --- WP-5 finding~1, the per-complex reading of ``reported''.
  A term $H^{q}(F)$ is reported iff $q<0$, or $q\le q_{\max}(F)$, or $F$ is not
  truncated; the clause's existing sentence is this rule in the case
  $q_{\max}(\text{bulk})=q_{\max}(\text{boundary})$, which is what \ref{open:ell2}'s
  spelled-out clause states. \textbf{Nothing on the shipped models changes} --- all
  three complexes are finite and untruncated there --- and on a truncated model the
  rule is the one \code{HeckeConeValidation} implements and exercises. Obligation
  clause, no tag. \textbf{This is the one entry a referee should attack}: on a
  hypothetical model with $q_{\max}(\text{bulk})\neq q_{\max}(\text{boundary})$ it
  asserts more spots than the single-$q_{\max}$ sentence did.
\end{itemize}

\paragraph{Non-normative.}
\begin{itemize}\setlength{\itemsep}{0pt}
\item \cl{rmk:btabsmall} --- a second \COMP{} record beside the $B\le196$ one: the
  measurement that found \cl{lem:bounds}(6)'s gap ($\lVert z\rVert=144.8\,N^{2}$ over
  $\ell\in\{11,\dots,3617\}$, constant along the reduction, independent of $n$), the
  int64 threshold $N\ge4580$ it predicts --- exactly the slot-1 bracket $3617<N<5003$ ---
  and the fact that the slot-1 record's attribution of that bracket to the coset
  construction was wrong. Remark, \COMP\textsuperscript{r} kept.
\item \S\ref{sec:heckedatum}'s $U_\ell$ reach record --- \textbf{deferred item~7}'s
  status half, which the item explicitly asks this version for. The v1.4.2 sentence
  (``the repair did not run'') is replaced by the measured outcome: the site was
  step~1, not the coset construction; \cl{lem:bounds} did not fail; and after the
  per-operation escalation the binding constraint is again \cl{lem:bounds}(5), giving
  $\ell\le23170$, past which the table is \emph{refused by \cl{obl:tau}'s certification}
  rather than lost to an anonymous overflow. With the measured $(16,5003)$ and
  $(16,10007)$ points. Container.
\item Appendix~\ref{app:status} --- \cl{rmk:coneambiguity691}'s status column and the two
  counts that follow from it: \COMP{} $17\to18$, its registered half $12\to13$. The
  tagged/untagged split $116/89$ is unchanged (the remark was already tagged), and no
  line was added or removed. Container.
\item This appendix --- the v1.4.2 entry's \texttt{[hash]} back-filled with
  \texttt{833c98f} (\ref{CL6}); the deferred list of that entry marked resolved. The
  header comment and the date line. Containers.
\item \textbf{Outside this document} (\ref{CL2}): \file{hecke\_extensions.tex} goes to
  [E] v1.4.1, stating \code{cert:cyclic}'s canonical generator as part of the
  certificate rather than only as an implementation note, and re-pointing its companion
  citation at [C] v1.5.0. \file{CLAUDE.md}, \file{docs/KNOWN\_ISSUES.md} \S3, the
  implementation plan [X] and the console plan [P] are slot~3's close (WP-11), not this
  version's.
\end{itemize}

\paragraph{The v1.4.2 deferred list: resolved.}
All seven items below are closed by this version --- items 1--6 as the slot-2 reports
proposed and the owner pre-approved, item~7 by WP-1's report, which found more than the
item expected: not a status sentence alone but a missing clause of \cl{lem:bounds} and a
false clause of \cl{thm:backend}(1). \textbf{Nothing is carried forward.} What is
\emph{not} here, and why: WP-5's findings 3 (\cl{eq:conedims}'s first identity needs no
$\Heq{2}=0$) and 5 (\cl{obl:coneexact}(iv)'s ``$\im(\Hcpt{1}\to\Heq{1})=\ker\rres^{1}$ is
checked explicitly'' read degree-generically) would each \emph{strengthen} a statement or
widen an obligation's asserted set, which \S3 WP-8 of the slot plan reserves for the
owner; they are listed in WP-8's report for the next $X$. So are the findings of slot-3
WP-3 (strict locality) that bear on [E] \code{sec:localcert} rather than on this
document.

\subsection*{v1.4.2 --- 2026-09-16 --- hecke-extensions slot 2, WP-13: the slot close --- status only: \texorpdfstring{\cl{sec:recordschema}}{sec:recordschema} item 5 implemented, the code maps for the slot-2 units, \texorpdfstring{\cl{obl:eta}}{obl:eta}(i)--(iv) discharged}

\noindent\emph{This version.} The close of hecke-extensions slot 2
(\file{docs/implementation/hecke\_extensions\_slot2\_plan.md} \S3 WP-13, \S4.3: ``status only'').
Every sentence changed here reports what the tree does on 2026-09-16 at the slot's close
(commit \texttt{45f7c4e} before this version), and nothing else: the slot's one golden
re-baseline (WP-2, \code{hecke\_report/v2}, $384$ report lines, $0$ elsewhere), the locality
certificate inside the refinement loop and the first mathematical setting (WP-1), the local
algebra and both generation certificates (WP-4, WP-5), $H^{1}_{!}$, Harder and the two
recorded diagnostics (WP-7), and the certified $\etapair$-table (WP-8). The slot's planned
$U_\ell$ reach repair (WP-3) \textbf{did not run}, and the record of \cl{sec:heckedatum}
says so. \textbf{Changes no normative statement} (\ref{V2}, \ref{CL5}): no hypothesis,
displayed formula, asserted predicate, tag or label moves; Appendix~\ref{app:status}'s
counts are unchanged ($116/89$, WP-10's recount). Report:
\file{docs/implementation/slot2/WP-13\_report.md}. Commit: \texttt{833c98f}.

\paragraph{Normative.}
\textbf{None.} \ref{CL5}'s test is met vacuously.

\paragraph{Non-normative.}
\begin{itemize}\setlength{\itemsep}{0pt}
\item \cl{sec:recordschema} --- the writer's schema string (\code{v2} since v1.4.2), and
  item~5's status: ``scheduled'' $\to$ \textbf{implemented at \code{hecke\_report/v2}}, per
  complex, derived from \code{HeckeReportInput::num\_degrees}, which the writer requires;
  \code{"set"} reserved for E9; the campaign writer's three un-hashed fields; the
  re-baseline's count and note. Container; site~(6) of \cl{sec:oc1sites} unchanged.
\item \cl{rmk:conestrict}(1) --- the \COMP{} record's counts re-read on the canonical
  $\tau$ (v1.1.0): $10$ of $12$ vertex slots and $5$ of $9$ edge slots leave $\Xbdry$
  (\code{hecke::cusp\_support\_walk}, WP-7, confirmed by hand on
  \file{tau\_table\_p2.txt}); the ``$j=1$ slots all stay'' sentence was the pre-v1.1.0
  table's and is withdrawn with the old counts kept in a parenthesis; the second statistic,
  $11$ of $12$ and $6$ of $9$ slots with a non-zero $\etapair$ (WP-8), added and
  distinguished from the first. The conclusion of (1) is unchanged and holds a fortiori.
  Remark, tag \PROVED{} kept (it tags (2), not the record).
\item \cl{rmk:coneambiguity691} --- a status sentence: $\rres^{1}$ is now computed and has
  rank $1$ at $(10,691)$ (WP-7), so $\dim\ker\rres^{1}=2$ and $\dim X_3=3$; the stale
  arithmetic is left standing under its \OPENSTAT{} tag and listed below for the next $X$.
\item \cl{obl:eta} --- the roster row of \S\ref{sec:obligationroster}: ``unimplemented''
  $\to$ \textbf{partial}, (i)--(iv) discharged and (vi) recorded by
  \code{HeckePairHomotopyValidation}, (v) awaiting the cone. Obligation text unchanged.
\item \cl{sec:cone}'s code map --- ``nothing in this chapter is implemented'' $\to$
  implemented in part: \file{Hecke/HeckePairHomotopy}, \file{Hecke/HeckeIntegerSolve}, the
  validation, what is still missing; the census sentence past tense. Container.
\item \cl{sec:heckedatum}'s $U_\ell$ reach record --- a status sentence: the repair
  (slot-2 WP-3) did not run; slot 3. Container.
\item \cl{sec:codemap} --- rows \cl{sec:heckedatum} (\file{HeckeIntegerSolve}),
  \cl{sec:boundarydatum} (\file{HeckeExt/HeckeBoundary},
  \file{Hecke/HeckeBoundaryDiagnostics}), \cl{sec:cone} split from \cl{sec:six} /
  \cl{sec:orbifoldmodel} (the former in part), \cl{sec:obligations} ($39$ \code{Hecke*}
  entries), the level (b),(c) row (\file{HeckeModuleDecomposition}'s certified loop,
  \file{HeckeSettings}, \file{HeckeExt/HeckeLocalAlgebra},
  \file{HeckeExt/HeckeGenerationCertificates}), and the plan pointer v1.5 $\to$ v1.6 in the
  \cl{sec:symconv} and \cl{sec:orbifoldmodel} rows. Container.
\item \cl{sec:notimplemented} --- ``two chapters have no implementation'' $\to$ one
  (\cl{sec:orbifoldmodel}), \cl{sec:cone} partial; chapter~\ref{sec:boundarydatum}'s
  sentence gains $H^{1}_{!}$ and Harder. Container.
\item \cl{sec:oraclesnew} and \cl{sec:example}'s closing paragraph --- the family (A)/(B)
  status sentences (the second had said ``no boundary Hecke operator'' since before
  v1.3.3). Containers.
\item \cl{sec:involution}'s census sentence, \S\ref{sec:obligationroster}'s preamble and
  Appendix~\ref{app:formalization}'s oracle item --- the count $77\to80$ registered
  validations, and the oracle item's two chapters $\to$ one and a half. Containers.
\item This appendix --- the v1.4.1 entry's \texttt{[hash]} back-filled with
  \texttt{9c54b91} (\ref{CL6}); the list of $Y$ bumps under \ref{CL5} gains v1.4.2. The
  header comment and the date line. Containers.
\item \textbf{Outside this document} (\ref{CL2}): \file{CLAUDE.md} (the ``[C] v1.3.3''
  pointer, the suite table, the slot-2 roster), \file{docs/KNOWN\_ISSUES.md} \S3 (the E1b
  booking, WP-3 not run) and the implementation plan [X] v1.6 --- all WP-13's;
  \file{hecke\_extensions.tex} stays at [E] v1.4.0 citing [C] v1.4.1, which is unaffected
  by a status bump (no label moved).
\end{itemize}

\paragraph{Deferred to the next $X$ revision (findings the slot-2 packages handed to
this version, none of which a status bump may act on).}
\begin{enumerate}\setlength{\itemsep}{0pt}
\item \cl{rmk:coneambiguity691}: the displayed arithmetic assumes $\rres^{1}=0$ at
  $(10,691)$; measured $\operatorname{rank}\rres^{1}=1$, $\dim\ker\rres^{1}=2$,
  $\dim X_3=3$ (WP-7 \S3.3). Rewrite from the measured values, with the mod-$\ell$ reading
  of ``the Eisenstein class'' at a Kummer cell as the joint eigenspace $E$ ([E] v1.4.0
  \code{cl:harder}); the same reading applies wherever \cl{obl:oracle}(b)'s \STD{}
  citation speaks of the Eisenstein class mod $\ell$ (WP-7 finding 3, WP-11 item 4).
\item \cl{obl:eta}(ii), ``the shared (T1) helper --- the same helper the construction
  uses'': there is no one helper shared by three tables; each table has its own (T1)
  helper, all on the one (P1)/conjugation primitive
  \code{HeckeCosetTables::twisted\_conjugate}. Proposed: ``through each table's (T1)
  helper, all built on the one (P1) split'' (WP-8 finding 1).
\item \cl{obl:eta}(iv), ``the explicit setting that permits'' a mod-$m$ fill: the
  shipped implementation has no mod-$m$ path and no such setting; a tier failure aborts.
  Proposed: an implementation MAY omit the mod-$m$ regime, and then the record's regime is
  $\Z$ by construction (WP-8 finding 2).
\item \cl{lem:etaexists}(V): a sentence naming the constructive choice actually made ---
  the pairwise path routing of \cl{alg:tau} step~2, which needs no linear solve and
  reduces to the edge path when both values are single vertices (WP-8 finding 3).
\item \cl{sec:recordschema} item~5 / \cl{open:ell2} for a chain \emph{map}: the
  $\etapair$ record carries the \emph{source} (boundary) degrees $[0,1]$ with $q_{\max}=1$,
  $\etapair(q)$ landing in bulk degree $q+1$; the rule for a truncated free model (E9) is
  to be stated (WP-8 finding 4).
\item \cl{obl:eta}'s preamble, ``together with the stored $\tau$- and $\taub$-tables it
  was built against'': neither base table exposes a table hash; the $\etapair$ record
  computes both (\code{built\_against.tau\_fnv64} / \code{tau\_bd\_fnv64}). Proposed: the
  base tables expose \code{table\_fnv64()} and E2's record carries the three hashes as
  tier-B objects (WP-8 finding 5).
\item \cl{sec:heckedatum}'s $U_\ell$ record and \cl{lem:bounds} / \cl{thm:backend}: no
  finding --- WP-3 did not run, so the reach limit stands as recorded and nothing is
  flagged; the repair's report (slot 3) may flag \cl{lem:bounds} then.
\end{enumerate}

\subsection*{v1.4.1 --- 2026-09-16 --- hecke-extensions slot 2, WP-10: referee pass on v1.4.0 --- all nine accepted, tags kept, \texorpdfstring{\cl{rmk:taubdrynoiota}}{rmk:taubdrynoiota}(3)'s exhibit moved to the right slots, prose repaired}

\noindent\emph{This version.} The independent read of v1.4.0 that
\file{docs/implementation/hecke\_extensions\_slot2\_plan.md} \S3 WP-10 asks for, on the
model of v1.3.1 (\texttt{f69cda0}): every one of the nine resolutions was re-derived ---
the arm map of \cl{obl:taubdry}(v.a)--(v.c) against \file{Hecke/HeckeBoundaryOperator.cpp};
the general-$\nu$ coset shape; \ref{FD2}, \ref{FD4} and the coincidence groupoid of
$(\Kbdry,\Xbdry)$ for \cl{prop:circlemodel}(4); the bijectivity of $\Phi$ and the
orientation sign; $\chi_0-\chi_1=(1-p)(1-p^{n})$ and the Dickson-power congruence
$b\equiv n-b\pmod{\ell-1}$; $\VR^{T^{p}}=\VR^{T}$; \cl{alg:taubdry} step by step against
\file{Hecke/HeckeBoundaryDatum.cpp} and \code{boundary\_line\_nearest}; the
$216/216$ record against \code{run\_restriction\_arm}; the per-complex schema rule against
owner confirmation 1 and \cl{open:ell2} rider (i). \textbf{Changes no normative statement}
(\ref{V2}, \ref{CL5}): no hypothesis, displayed formula, asserted predicate or tag moves.
One proof was wrong as written and is repaired under its own label (\ref{L3}):
\cl{rmk:taubdrynoiota}(3) computed the $(2,1)$ exhibit with $g$ in place of $g^{-1}$
($g^{-1}$ acts on $\Hyp$ as $z\mapsto Nz$, \cl{alg:tau} step~1), placing it on the slots
$Sv_1$, $Sv_2$, where the transported real parts are $\mp\tfrac1{12}$ and no tie occurs;
the asymmetry is real and sits on $v_1,v_2$ and $v_0,v_3$ of $\Kbdry$ --- confirmed on
the shipped table. \ref{L1}: the label set is unchanged (v1.3.3 $\to$ v1.4.0 added exactly
the four labels the v1.4.0 entry names, none removed). Appendix~\ref{app:status} was
recounted mechanically from its table and agrees with ``The counts'' in every figure
except the tagged/untagged split, corrected below. Report:
\file{docs/implementation/slot2/WP-10\_referee\_report.md}. Commit: \texttt{9c54b91}
(the report, committed before any repair, is \texttt{2fed619}).

\paragraph{Normative.}
\textbf{None.} \ref{CL5}'s test is met vacuously.

\paragraph{Non-normative.}
\begin{itemize}\setlength{\itemsep}{0pt}
\item \cl{rmk:taubdrynoiota} --- proof and prose, tag \PROVED{} kept. (2): the sentence
  ``consequently \ldots cannot be arranged by the choice of tie rule for all $(p,\nu)$ at
  once'' was asserted without the step that the mirror point is ever a transported real
  part; it now says that an integer transported real part occurs at every $(2,\nu)$
  ($x(v_0)=-N/2$) and obstructs every $T$-equivariant vertex rule there, and that for odd
  $p$ it obstructs nothing; the proof gains the reduction of the identity at an integer
  to the identity at $0$. (3): the exhibit moved from $Sv_1,Sv_2$ to $v_1,v_2$ (transported
  real parts $N\,\mathrm{Re}\,v=\mp\tfrac13$, ties to $v_1$ and $v_3$) and $v_0,v_3$
  ($x=\mp1$, ties to $T^{-1}v_2$ and $Tv_2$, $\inv_{\Xbulk}(T^{-1}v_2)=Tv_1$), with the
  general criterion for an integer transported real part
  ($x=(p^{a}\mathrm{Re}\,v+B)/p^{\nu-a}$; integral only for $p=2$, $v\in\{v_0,v_3\}$) and a
  parenthetical recording what v1.4.0 had computed. (4): made consistent with (2)+(3) ---
  the table-level identity fails at every $(2,\nu)$ under every $T$-equivariant rule; at odd
  $p$ the mirror-consistent-tie-rule question stays not settled and not claimed. What is
  asserted after the repair is what (2) already asserted; it is now proved.
\item \cl{obl:taubdry} --- prose, two clauses. (v.c): ``at cochain level the two data are
  homotopic, not equal (\cl{lem:taubdryexists}(3))'' was unproved --- the $\Gbdry$-level
  datum of route A is not a boundary $g$-twisted chain datum, and the decision box after
  \cl{def:bdryop} calls the transcription a validation route only; the clause now says the
  two operators are not asserted equal and no chain homotopy is claimed, the asserted
  predicate (cohomology-level, entrywise, aborting) unchanged. (vii): ``with the
  right-coset sum $M_{\beta}$ of \cl{prop:bdrychar}(2) in place of $\Abdry{1}$'' was a type
  error, $M_{\beta}$ being an endomorphism of $\VR^{T}=H^{0}$; the $H^{1}$ control is the
  scalar $\chi_1(p)$ in place of $\Abdry{1}=\chi_0(p)$ on the one-dimensional $H^{1}$
  ($n<\ell$), which is what \code{HeckeBoundaryOperatorValidation} substitutes; the
  condition $\chi_0(p)\neq\chi_1(p)$, $\rres^{1}\neq0$ is unchanged. Obligation, no tag.
\item \cl{prop:circlemodel} --- proof of (4), one parenthetical: ``\ref{FD1} fails,
  $\Xbdry$ being disconnected'' was false (\ref{FD1} concerns $\abs{K}$ and $\widetilde D$,
  both fine for the path); what fails is \cl{ass:fund_domain}'s standing hypothesis that
  $\abs{\Ku}$ be connected, and the sentence now says so and names the face closure
  \cl{ass:consistent_face_id} of $\Rel_{\partial}$, which the theorem's hypotheses require
  and the proof had not mentioned. Statement, hypotheses and tag unchanged.
\item \cl{alg:taubdry} --- step~3, wording: the two vertex values are ``already stored (a
  derived face's value being the one the step-0 helper produced)'', which is what
  \code{construct\_edge} does, in place of ``fetched through the step-0 helper where a
  face is derived''. Procedure, no tag.
\item Appendix~\ref{app:status} --- ``The counts'': $115/90$ tagged/untagged with $11$
  untagged obligations $\to$ $116/89$ with $10$, since \cl{obl:taubdry} has carried
  \COMP\textsuperscript{r} since v1.3.3 and is counted in the $17$; every other figure
  confirmed by a mechanical recount of the table ($205$ rows; $63$ \code{rmk}, $4$
  \code{alg}; \PROVED{} $93$, \COMP{} $17=12+5$; \BNEW{} $125$, \BRH{} $45$, \BFIX{} $34+1$).
  This appendix --- the v1.4.0 entry's \texttt{[hash]} back-filled with \texttt{75c0c27}
  (\ref{CL6}); the list of $Y$ bumps under \ref{CL5} gains v1.4.1. The header comment and
  the date line. Containers, no tag.
\item \textbf{Outside this document} (\ref{CL2}): nothing --- [E]'s pointer remains WP-11's,
  CLAUDE.md the close package's.
\end{itemize}

\subsection*{v1.4.0 --- 2026-09-16 --- hecke-extensions slot 2, WP-9: the nine open points of v1.3.3 resolved --- \texorpdfstring{\cl{obl:taubdry}}{obl:taubdry}(v) split and (vii) qualified, \texorpdfstring{\cl{prop:circlemodel}}{prop:circlemodel}(4), \texorpdfstring{\cl{alg:taubdry}}{alg:taubdry}, the per-complex record fields}

\noindent\emph{This version.} The slot's normative revision of the boundary chapter
(\file{docs/implementation/hecke\_extensions\_slot2\_plan.md} \S2.2 and \S3 WP-9;
report \file{docs/implementation/slot2/WP-9\_report.md}). It resolves \textbf{exactly}
the nine points that v1.3.3's entry below recorded as ``open points recorded, not
resolved'', each as the slot-1 work-package report that found it proposed
(\file{docs/implementation/slot1/WP-4a\_report.md} findings 2, 3, 5;
\file{WP-4b\_report.md} findings 9--15, 17), under the owner's close-out decision (a) of
slot 1 (``as the work-package reports propose, in the next normative ($X$) revision'')
and the two slot-2 confirmations that the reports had left open (confirmation 1: the
record fields are per complex, never a maximum; confirmation 2: clause (v) is narrowed to
degree $0$ at $\nu\ge2$, the general-$\nu$ lemma recorded as not stated). \textbf{This is
an $X$ bump} (\ref{V2}, \ref{CL5}): an obligation's clauses change, a proposition gains a
clause, a proposition's statement is made precise, an algorithm and a proved remark are
added, and the record schema's rule changes. Nothing outside chapter~\ref{sec:boundarydatum},
the two schema sites, one remark in \S\ref{sec:cone} and the indices is touched; chapter
\ref{sec:involution}'s $\pm$ statements, \cl{lem:bulkperiodic} and everything the
concurrent E4b package may discharge are deliberately left for later bumps. No erratum is
filed (\ref{CL2}); \ref{L2} is not engaged --- every label keeps its subject, and the one
statement that was false as written, \cl{prop:circlemodel}'s ``single relation'', is
corrected under its own label, as \ref{L3} directs. Every validation that discharges
\cl{obl:taubdry} evaluates the rewritten clauses \emph{as written}, with no arm changed
(the roster row says which arm is which sub-clause). Commit: \texttt{75c0c27}.

\paragraph{Normative.}
\begin{itemize}\setlength{\itemsep}{0pt}
\item \cl{obl:taubdry} \PNEW{} --- \textbf{clause (v)} (open point 1; [WP-4b] 10):
  \emph{before}, ``the operator obtained from \cl{eq:bdrychain} and the operator
  $\sum_{\delta}\rhon(\delta)(\cdot)$ written in the coordinates of
  \cl{prop:circlemodel}(1) agree \ldots in every computed degree'' --- one named
  comparison, where the chapter offers three and the implementation runs all three;
  \emph{after}, three sub-clauses, each naming its degrees and $\nu$: (v.a) degree $0$,
  every $\nu$, against \cl{prop:bdryop}(2) through \cl{prop:circlemodel}(1); (v.b) degree
  $1$, $\nu=1$, against \cl{eq:bdrygroupform} in the coordinate $\Phi$ of
  \cl{prop:bdrychar}(3); (v.c) every computed degree, $\nu=1$, against the $\Gbdry$-level
  datum of the decision box after \cl{def:bdryop}, asserted on cohomology only.
  \textbf{Clause (v) at $\nu\ge2$} (open point 2; [WP-4b] 11; owner confirmation 2):
  \emph{before}, silent --- (v.b)'s and (v.c)'s inputs are $\nu=1$ statements, so at
  $\nu\ge2$ only the $H^{0}$ half was available; \emph{after}, ``at $\nu\ge2$ clause (v)
  is (v.a) alone'', $H^{1}$ there being covered by (vii) and \cl{prop:bdryop}(1) as run.
  A narrowing to what is proved, not an extension. \textbf{Clause (vii)} (open point 5;
  [WP-4b] 13): \emph{before}, ``it holds for $1+p^{n+1}$ with the left-coset operator and
  fails with the right-coset one'' --- over-broad, the failure being invisible where
  $\chi_0(p)=\chi_1(p)$ ($62$ of $232$ $n<\ell$ cells) and never on $H^{0}$ of the model;
  \emph{after}, ``fails whenever $\chi_0(p)\neq\chi_1(p)$ in $\cring$ and
  $\rres^{1}\neq0$'', with the $H^{0}$ statement made exact: the control cannot bite on
  any power of the first Dickson class ($b\equiv n-b\pmod{\ell-1}$, \cl{cl:chi}(3)), hence
  not on the model complexes below degree $\ell(\ell-1)$; plus the pointer to
  \cl{rmk:conerestrictsmallell}. Underspecified made precise; over-broad corrected.
\item \cl{prop:circlemodel} \PFIX{} --- (open point 3; [WP-4b] 9) \emph{before}, the
  preamble called $\Ceq(\Kbdry,\Rel_{\partial};L)$ ``the circle model with its single
  relation $f(v_3)=\rhon(T)f(v_0)$'' and identified it with $C^{\bullet}_{\partial}$,
  which is \textbf{false for $n$ odd, $2\in\cring^{\times}$}: $\Hom_{\Z[\Ga]}$ carries the
  central $-I$ constraint and the single-relation model does not (the defect was found by
  the phase that implemented the chapter, and the implementation had already added the
  self-relations; a second repair of a \PFIX{} statement, \ref{CL3}); \emph{after},
  $\Rel_{\partial}$ is \emph{defined} as the cusp relation together with the $-I$
  self-relations $(\sigma,\sigma,-I)$ on every simplex, and a new \textbf{clause (4)}
  proves $\Ceq(\Kbdry,\Rel_{\partial};L)=C^{\bullet}_{\partial}$ for every $n$ and every
  $\cring$ --- via \ref{FD2} and \ref{FD4} for $(\Kbdry,\Rel_{\partial})$ and the proof of
  \cl{thm:explicit_complex_equiv}(1), which uses exactly those two --- with the
  single-relation form kept as the corollary ``literally when $\rhon(-I)=\id$, i.e.\ $n$
  even or $2=0$ in $\cring$'', the zero complex for $n$ odd at $2\in\cring^{\times}$, and
  the explicit statement that the single-relation model on $\VR$ is not
  $C^{\bullet}_{\partial}$ there. Clauses (1)--(3) unchanged, numbering kept. Tag
  \PROVED{} kept, the new clause proved.
\item \cl{prop:bdrychar} \PFIX{} --- \textbf{clause (3)} (open point 4; [WP-4b] 12):
  \emph{before}, the identification $H^{1}(C^{\bullet}_{\partial})\cong(\VR)_{T}$ that
  \cl{eq:bdrygroupform} depends on was used silently in the proof; \emph{after}, the
  statement carries it: $\Phi\colon[\varphi]\mapsto[\varphi(\langle v_0,Tv_0\rangle)]$,
  the signed sum over the oriented fundamental path, equivalently $[f(T)]$ for the group
  cocycle $f(\gamma)=\varphi(\langle v_0,\gamma v_0\rangle)$, with its inverse; the proof
  gains the verification that $\Phi$ is a well-defined bijection. \textbf{Resolved with
  one correction to the report's remark}: reversing the orientation gives $-\Phi$, which
  leaves the transported operator \emph{unchanged}; what must agree is the orientation
  used by a representative map and by a lift, and with the two on opposite orientations
  the transport is off by $-1$ --- which is the disagreement [WP-4b] observed. The
  statement says so. \textbf{Clause (2)} (open point 6; [WP-4b] 14): \emph{before},
  nothing said what object the right-coset substitution is; \emph{after}, a clarifying
  clause: $M_{\beta}=\sum_b\rhon(\beta_b)+\rhon(\beta_{\infty})$ is a well-defined
  endomorphism of $\VR^{T}$ (proved: $T\beta_b=\beta_bT^{p}$,
  $T^{p}\beta_{\infty}=\beta_{\infty}T$, $\VR^{T^{p}}=\VR^{T}$ for $p\neq\ell$), the
  monomials are its eigenbasis, it is asserted as an operator on $\VR^{T}$ only, it is
  not asserted to be induced by any cochain map, and the implementation builds it as a
  fibre-level negative control. Underspecified made precise, both clauses.
\item \cl{alg:taubdry} \PNEW{} --- \textbf{new} (open point 7, first half; [WP-4a] 2,
  5): the canonical $\taub$ as a terminating procedure beside \cl{alg:tau}, per
  rule~\ref{A3}: the slot/derivation step through the single (T1$\partial$) helper, the
  transported real part $x=\mathrm{Re}(\gamma_{\mathrm{red}}\,g^{-1}s_j\mathrm{pos}(v))$,
  the vertex rule ``nearest line vertex, ties to the larger real part'', the edge walk,
  constructor-time certification, and the normalisation
  $\gamma_j=\mathrm{red}(g^{-1}s_j\infty)^{-1}$ of \cl{prop:bdrycosets}(2). Four
  properties \PROVED: well defined and independent of the derivation order and of the
  cusp reduction (the rule is a $\Ga''$-twisted-equivariant function of the vertex), a
  boundary datum over $\Z$ with $\{0,\pm1\}$ coefficients, homotopic to every other
  admissible datum, and free of $\ell$ and $n$. It matches
  \file{Hecke/HeckeBoundaryDatum}'s \code{construct\_vertex} / \code{construct\_edge} /
  \code{try\_derive} and \file{Hecke/SL2ZTessellationTables}'
  \code{boundary\_line\_nearest} step by step. \cl{lem:taubdryexists}(1) gains a pointer
  to it (prose).
\item \cl{rmk:taubdrynoiota} \PNEW{} --- \textbf{new}, \PROVED{} (open point 7, second
  half; [WP-4a] 3): (1) $\inv_{\Xbulk}\taub\inv_{\Xbulk}$ is an admissible boundary datum,
  homotopic to $\taub$; (2) no $T$-equivariant vertex rule --- any tie rule, or none ---
  is $\inv$-equivariant, because at the mirror point $x=0$ it would need an
  $\inv_{\Xbulk}$-fixed vertex of $\Xbdry_{\infty}$ and there is none (the $\Gbdry$-orbit
  structure of \cl{rmk:bdrymodel}); (3) the canonical table at $(p,\nu)=(2,1)$ is not
  $\inv$-symmetric, by an explicit computation on the slots $Sv_1$, $Sv_2$ (transported
  real parts $\mp\tfrac13$, ties resolved to $v_1$ and $v_3$, $\invK v_3=v_0\neq v_1$);
  (4) what is and is not asserted --- the table-level identity is asserted nowhere and
  recorded nowhere, and whether a mirror-consistent tie rule would give it at every
  $(p,\nu)$ is \emph{not settled}. \textbf{Resolved with a scope correction to the
  report}: [WP-4a] 3's argument (``the tied set at the mirror is a full $\invK$-orbit'')
  proves the impossibility for the \emph{rule}, not for the emitted table, since the
  mirror point need not be a transported real part; the remark says exactly what each
  argument proves. \cl{obl:iota} unchanged.
\item \cl{rmk:bdrynugeneral} \PNEW{} --- \textbf{new}, commentary, no tag (open point 2's
  companion; owner confirmation 2): records that no double-coset decomposition of $B_N$
  for $\nu\ge2$, no general-$\nu$ \cl{eq:bdrygroupform} and no $\Gbdry$-level datum for
  (v.c) at $\nu\ge2$ is stated, and records the shape a later addition would have to
  prove (indexing by $u$ and $b\bmod p^{\min(u,\nu-u)}$, $p^{\max(0,2u-\nu)}$ left cosets
  per double coset), explicitly as not asserted.
\item \cl{sec:recordschema} item~5 --- (open point 9; [WP-4b] 17; owner confirmation 1)
  \emph{before}, ``the record carries three fields'' --- one $q_{\max}$ per record, the
  reported-degree list as the third field, and the instruction addressed to E5 ``not
  implemented now''; \emph{after}, the four fields \code{q\_max}, \code{truncated},
  \code{q\_max\_source}, \code{degrees} are \textbf{per complex}, never a maximum over a
  record; a single-complex record carries them at record level, a multi-complex record
  inside each complex's object; the campaign's reported-degree selection is a distinct
  job-level field; the job identity includes each complex's tuple (the campaign job-id
  fix is done, \code{CampaignSmokeValidation}); the model values are bulk
  $2/\code{false}/\code{"derived"}/[0,1,2]$ and boundary $1/\code{false}/\code{"derived"}/[0,1]$;
  the instruction is re-addressed to E0b (slot-2 WP-2) as \emph{scheduled}, with the
  ``implemented'' sentence left to the version after the re-baseline. \cl{open:ell2}'s
  boxed record sentence gains ``per complex it is about'' (its ``verbatim the v0.14.0
  working assumption'' note says so); \cl{sec:oc1sites} items (3) (the split clause) and
  (6) (the four fields) follow. Schema rule; container labels, no tag.
\end{itemize}

\paragraph{Non-normative.}
\begin{itemize}\setlength{\itemsep}{0pt}
\item \cl{rmk:conerestrictsmallell} \PNEW{} --- \textbf{new} in \S\ref{sec:cone},
  \COMP\textsuperscript{r} (open point 8; [WP-4b] 15): \cl{obl:taubdry}(vii) $=$
  \cl{obl:coneexact}(iii) holds on all $216/216$ $(\text{cell},q)$ pairs of
  \code{HeckeBoundaryOperatorValidation}'s grid at $\ell\in\{2,3\}$ (run 2026-09-16),
  recorded as a computational fact with the validation named, consistent with a pair
  homotopy over $\Z$, claimed nowhere, and licensing nothing (\ref{A6}). No statement
  moves.
\item \cl{rmk:bdryoddn} --- one sentence: since \cl{prop:circlemodel}(4) the odd-$n$ zero
  complex is literally what the model computes. Remark unchanged in claim; prose.
\item \cl{sec:boundarydatum} --- code map: the (v) paragraph names the three sub-clauses
  and which arm discharges each; the ``choices the implementation fixed'' bullets point
  at \cl{alg:taubdry}, \cl{rmk:taubdrynoiota} and \cl{prop:circlemodel}(4); the
  ``what it records'' paragraph's per-complex question is closed by \cl{sec:recordschema}
  item~5; the (vii) bullet points at \cl{rmk:conerestrictsmallell} and at the
  qualified clause. Container.
\item \cl{obl:taubdry} --- roster row (\cl{sec:obligationroster}): (v) per sub-clause
  with its arm, (vii)'s control failures as the qualified clause predicts, the
  $\ell\in\{2,3\}$ pointer. Evidence.
\item Appendix~\ref{app:status} --- four new rows (\cl{rmk:bdrynugeneral} commentary,
  \cl{alg:taubdry} procedure, \cl{rmk:taubdrynoiota} \PROVED, \cl{rmk:conerestrictsmallell}
  \COMP\textsuperscript{r}); the \COMP{} split sentence (eleven of sixteen $\to$ twelve of
  seventeen); the counts ($201\to205$ statements, $60\to63$ remarks, $3\to4$ algorithms,
  \PROVED{} $92\to93$, tagged $113\to115$ / untagged $88\to90$, \BNEW{} $121\to125$) ---
  and the \COMP{} figure of ``The counts'', which v1.3.3 had left at $13$ ($8$
  registered) against a table of $16$ ($11$), corrected to $17$ ($12$) with a note.
  Appendix~\ref{app:labelindex} --- the four labels dereferenced. This appendix ---
  the provenance table (\PNEW{} $121\to125$: $70$ assertions or definitions, $55$
  remarks; $190\to194$ labels on \PNEW{} lines) and the Part~IV boundary roster ($10\to11$,
  \cl{alg:taubdry}). The header comment and the date line. Containers, no tag.
\item \textbf{Open points recorded at v1.3.3: resolved, all nine} (the v1.3.3 entry
  below keeps its wording, as a published version's text does). (1) $\to$ \cl{obl:taubdry}(v.a)--(v.c);
  (2) $\to$ \cl{obl:taubdry}(v)'s $\nu\ge2$ sentence and \cl{rmk:bdrynugeneral};
  (3) $\to$ \cl{prop:circlemodel}(4); (4) $\to$ \cl{prop:bdrychar}(3); (5) $\to$
  \cl{obl:taubdry}(vii); (6) $\to$ \cl{prop:bdrychar}(2); (7) $\to$ \cl{alg:taubdry} and
  \cl{rmk:taubdrynoiota}; (8) $\to$ \cl{rmk:conerestrictsmallell}; (9) $\to$
  \cl{sec:recordschema} item~5. None remains open; the referee pass on this version is
  the slot's next package (WP-10, a $Y$ bump), and its status sentence for
  \code{hecke\_report/v2} the one after (WP-13).
\item \textbf{Outside this document} (\ref{CL2}): nothing --- the companion
  \file{hecke\_extensions.tex} pointer is left at v1.3.3 for the package that owns it
  (WP-11); CLAUDE.md is the close package's.
\end{itemize}

\subsection*{v1.3.3 --- 2026-09-16 --- hecke-extensions slot 1 close: chapter \texorpdfstring{\cl{sec:boundarydatum}}{sec:boundarydatum} implemented, \texorpdfstring{\cl{obl:taubdry}}{obl:taubdry} discharged, status sites and counts}

\noindent\emph{This version.} The close of slot 1 of the hecke-extensions stage
(\file{docs/implementation/hecke\_extensions\_slot1\_plan.md} \S3 WP-7;
\file{docs/implementation/slot1/WP-7\_report.md}). \textbf{Changes no normative statement}
(\ref{V2}, \ref{CL5}): no definition, obligation, hypothesis, displayed formula or tag on a
statement moves; what it records is \emph{evidence and status}, exactly as v1.1.1 did for
\cl{sec:involution}. The seven clauses of \cl{obl:taubdry}, written at v1.0.0 as a
specification for the phase that would build the boundary datum, were implemented and run
by that phase --- slot-1 packages WP-4a (\file{Hecke/HeckeBoundaryDatum}, the boundary
subcomplex of \file{Hecke/SL2ZTessellationTables}, \code{HeckeCosetTables::twisted\_conjugate})
and WP-4b (\file{Hecke/HeckeBoundaryOperator}, the circle model as a library call in
\file{LocalSystemSL2Z}) --- and two registered validations,
\code{HeckeBoundaryTauValidation} ((i)--(iv), $138$ tables, $52\,686$ checks, tier
\code{standard}) and \code{HeckeBoundaryOperatorValidation} ((v)--(vii), $546$ cells,
$19\,669$ checks, \code{standard}), evaluate every aborting predicate and record what the
clauses mark as recorded; nothing failed. The registered-validation count moves from $68$
to $77$ (slot 1 added \code{HeckeLocalityValidation}, \code{HeckeLocalityProbe},
\code{HeckeIdentityValidation}, \code{HeckeSturmBoundValidation},
\code{HeckeExtArchitectureValidation}, \code{ConsoleC0Validation} and the two above,
besides \code{GoldenDump}); the \code{Hecke*} count from $29$ to $36$. No golden value
changed in the slot (\code{GOLDEN\_STRICT=1}: tier~A $9446$ checked, $0$ failing; tier~B
$4603$ checked, $0$ mismatching, against the unchanged reference). Commit: \texttt{13760f7}.

\paragraph{Normative.}
\textbf{None.} \ref{CL5}'s test is met vacuously.

\paragraph{Non-normative.}
\begin{itemize}\setlength{\itemsep}{0pt}
\item \cl{sec:boundarydatum} --- the ``What this chapter fixes'' paragraph gains one
  status sentence, and the chapter's code map is rewritten from a census that found
  nothing to the record of what was built: the four units, the two validations with their
  grids, the clause-by-clause \COMP{} record of \cl{obl:taubdry}(v)--(vii) (both routes of
  (v) asserted in \code{compute()}; (vi)'s asserted parts and the recorded $330/330$
  agreement of \cl{cl:chi}(2); (vii) asserted on $864$ pairs at $6\in\cring^{\times}$ and
  recorded, $216/216$, at $\ell\in\{2,3\}$; the right-coset negative control of
  \cl{prop:bdrycosets}(4) returning $\chi_{b+1}$ and breaking (vii) on $170$ cells), the
  choices the implementation fixed where the statements leave them open (the
  nearest-line-vertex rule with ties to the larger real part in \cl{lem:taubdryexists}(1);
  the normalisation of $\gamma_j$ in \cl{prop:bdrycosets}(2); the reading of clause (iii)'s
  transporter set; $\Rel_{\partial}$ including the $-I$ self-relations, so that the circle
  model is $\Hom_{\Z[\Gbdry]}$ literally for every $n$), and the recorded fields
  (\code{boundary\_model}, \code{boundary\_is\_group\_cohomology}, \code{q\_max = 1}
  derived, \code{truncated = false}). The census's ``simplices not those of $\Ku$'' is
  corrected to ``slot set not that of $\Ku$''. Container and code map, no tag.
\item \cl{obl:taubdry} --- roster row (\cl{sec:obligationroster}): unimplemented $\to$
  \textbf{discharged}, clause by clause, with the two validation names; the roster's
  count sentence (``ten of the fifteen rows'' $\to$ eleven) and its prose (``three are
  unimplemented'' $\to$ two, naming \cl{obl:eta} and \cl{obl:coneexact}). Obligation
  unchanged; evidence.
\item \cl{sec:oraclesnew} --- opening paragraph: family (A) is a run, not a specification;
  family (A) gains its \COMP{} run record (the numbers above, the $(2,5,2)$ instance
  reading $4/4/4$). Container.
\item \cl{sec:codemap} --- the chapter-14 row now names the four units and the
  statements each realises; the row for \cl{sec:cone}, \cl{sec:six},
  \cl{sec:orbifoldmodel} stays ``---'' and points at phase E9 / C-A5; the
  \cl{sec:symconv} row records that $\Dn$ and the pairing are \emph{scheduled} (phase E7);
  the \cl{sec:obligations} row's \code{Hecke*} count; the \cl{rmk:dickson} roster row
  names \code{hecke::IsDicksonExceptionalCount} (\file{Hecke/DicksonCount.h}), the one
  Dickson count that replaced three file-local copies (slot-1 WP-2; the golden
  \code{dickson} entry unchanged). Container.
\item \cl{sec:notimplemented} --- ``three chapters'' $\to$ two, with a sentence on
  \cl{sec:boundarydatum} in the style of the v1.1.1 sentence on \cl{sec:involution}.
  Container.
\item \cl{sec:buildsurface} --- five targets $\to$ six (\code{ConsoleCore}, the console's
  STL-only language core, slot-1 WP-5; it realises no statement of this document).
  Container.
\item \cl{sec:cone}'s code map --- the census sentence no longer lists a missing
  $\taub$-table; the pair homotopy is what is still missing. Container.
\item \cl{thm:backend}'s consumer paragraph in \cl{sec:heckedatum} (the $U_\ell$ entry-point
  paragraph) --- gains a \COMP{} record of a reach limit found by slot-1 WP-2: for
  $U_\ell$ at $\nu=1$ the \code{int64} tier is exhausted between $\ell=3617$ (runs) and
  $\ell=5003$ (\code{HeckeOverflow} in the coset construction), independently of $n$,
  while $\ell=43867$ with a good prime runs; recorded, not repaired
  (\file{docs/KNOWN\_ISSUES.md} \S3). Container.
\item Appendix~\ref{app:status} --- \cl{prop:bdrycosets}, \cl{prop:bdrychar} and
  \cl{obl:taubdry} gain \COMP\textsuperscript{r}; the count sentence (eight of thirteen
  $\to$ eleven of sixteen). Appendix~\ref{app:formalization} --- ``three chapters'' $\to$
  two, the validation count, ``eight'' $\to$ ``eleven'' \COMP\textsuperscript{r} claims.
  \S\ref{app:changelog-gaps} --- ``three chapters'' $\to$ two. The header comment, the
  abstract, \ref{V5}'s closing sentence, the involution chapter's census count, the date
  line and the companion pointer (\file{hecke\_extensions.tex} v1.3.3). Containers, no tag.
\item \textbf{Open points recorded, not resolved} (each would be an $X$ bump by its own
  editor, or an owner decision; none blocks the discharge above, which holds under the
  statements as written): (1) \cl{obl:taubdry}(v) names one comparison where the chapter
  offers three (\cl{prop:bdryop}(2) on $H^{0}$, \cl{eq:bdrygroupform} on $H^{1}$, and the
  $\Gbdry$-level datum of the decision box after \cl{def:bdryop}); the implementation runs
  all of them --- the clause should say which are meant. (2) At $\nu\ge2$ only the
  $H^{0}$ half of (v) is available as stated: \cl{prop:bdrycosets}(3) and
  \cl{eq:bdrygroupform} are $\nu=1$ statements, and the double-coset decomposition of
  $B_N$ for general $\nu$ is not in the chapter. (3) \cl{prop:circlemodel}'s
  identification with the circle model ``with its single relation'' holds for $n$ even
  (or $\ell=2$) only; $\Hom_{\Z[\Gbdry]}$ carries the $-I$ constraint, and the
  implementation adds it (code map). (4) The $H^{1}\cong(\VR)_{T}$ coordinate that
  \cl{eq:bdrygroupform} depends on ($\varphi\mapsto\varphi([v_0,Tv_0])$) is used silently
  in the proof of \cl{prop:bdrychar}(3) and should be stated. (5) \cl{obl:taubdry}(vii)'s
  ``fails with the right-coset one'' holds only where $\chi_0(p)\neq\chi_1(p)$
  (\cl{cl:chi}(3)); on $62$ of the $232$ $n<\ell$ cells the two coincide and the control
  cannot bite, and on $H^{0}$ it cannot bite at all on this model. (6)
  \cl{prop:bdrychar}(2)'s right-coset substitution is an operator on $\VR^{T}$ only
  because the monomials are an eigenbasis; it has no cochain-level meaning, which a
  reader may try to give it. (7) Whether the nearest-line-vertex rule should become an
  \code{alg:taubdry} beside \cl{alg:tau}; and $\taub$ is not $\inv$-symmetric, which no
  tie rule can repair (code map) --- if \S\ref{sec:cone} or the involution ever need
  $\invK\taub\invK=\taub$, that is a contract addition. (8) \cl{obl:taubdry}(vii) holds
  at $\ell\in\{2,3\}$ on the model complexes ($216/216$), consistent with a pair homotopy
  over $\Z$; recorded for \S\ref{sec:cone}, claimed nowhere. (9) \cl{sec:recordschema}
  item 5 speaks of one $q_{\max}$; the boundary object derives $1$ and the bulk $2$ ---
  E0b's record integration decides.
\item \textbf{Outside this document} (\ref{CL2}): \file{hecke\_extensions.tex} v1.3.2
  $\to$ v1.3.3 (the $(160,691)$ probe verdict, the Sturm-bound values and implementation
  status, the reduced-character wording, the stale OQ2-table status sentence; PDF
  rebuilt); \file{hecke\_extensions\_implementation\_plan.md} v1.5 $\to$ v1.5.1 (\S9.3's
  verdicts and constants, \S1.3.1's resume-invariance wording); \file{CLAUDE.md};
  \file{docs/KNOWN\_ISSUES.md}; the P8 comment of \file{HeckeIotaSignValidation.cpp}
  that v1.3.2 left.
\end{itemize}

\subsection*{v1.3.2 --- 2026-09-16 --- coherence pass after OQ3 and OC-1: stale status sites, an over-broad $\pm$ rule, the direction of \texorpdfstring{\cl{rmk:iotapairing}}{rmk:iotapairing}'s prediction, and the provenance counts}

\noindent\emph{This version.} Applies the required fixes of the read-only coherence
review \file{docs/implementation/review\_2026-09-16\_oq3\_oc1\_coherence.md} ($89$ sites
checked after v1.2.0 and v1.3.0, $10$ inconsistent), less the one that v1.3.1 had already
made. \textbf{Changes no normative statement} (\ref{V2}, \ref{CL5}): no definition,
obligation, hypothesis, displayed formula or tag moves. One remark's \emph{prediction}
changes direction --- \cl{rmk:iotapairing}, indexed by Appendix~\ref{app:status} as
commentary, tagged nothing, consumed by no obligation, oracle or code map, and stating of
itself that it is ``asserted here as nothing more'' --- and that is recorded with its own
line below (\ref{CL1}, \ref{CL3}) rather than by an $X$ bump, which \ref{V2} reserves for
normative statements. Commit: \texttt{c2c6b44}.

\paragraph{Normative.}
\textbf{None.} \ref{CL5}'s test is met vacuously.

\paragraph{Non-normative.}
\begin{itemize}\setlength{\itemsep}{0pt}
\item \cl{def:iota} --- clause (iv), prose: ``not the pullback along any lift of
  $z\mapsto-\bar z$'' was false for $n$ odd, where $\rhon(-\inv)=-\rhon(\inv)$ and
  $\inv^{\ast}$ \emph{is} the pullback along $-\inv$; now qualified ``for $n$ even'', with
  the odd case and why it is harmless ($\Heq{\bullet}=0$ for $\ell\neq2$, the sign
  invisible at $\ell=2$) in a parenthesis (review R2). Two polish items: the sentence
  ``with $\det|_{\Ga}=1$'' now says what that buys (the relations of $\Rel$ are preserved,
  P2), and the flat-coordinate description is $-\varepsilon_{\sigma}\rhon(\inv)$ at block
  position $(\sigma,\abs{\invK\sigma})$, the form of the code map and of the validation
  (P3); the matrix is unchanged, being its own transpose. Definition; the formula
  \eqref{eq:iotacochain} untouched.
\item The $\pm$ subsection's opening paragraph (container) --- ``a $\pm$ read from a
  source written before v1.2.0 \ldots\ must be negated'' was over-broad: only the
  \emph{computed} signs of the v1.1.1 run record are geometric; the $\pm$
  \emph{predicted} from OQ3 (\cl{obl:iota}(v)'s split, family (C) of
  \cl{sec:oraclesnew}, the $H^{1}_{\pm}$ statements of \file{hecke\_extensions.tex}) were
  written in the convention v1.2.0 adopted and hold as printed, as the v1.2.0 entry
  already said (R3).
\item \cl{rmk:iotapairing} --- \textbf{the contingent prediction changes direction.}
  Until v1.3.2 the remark predicted that duality would pair each $\inv^{\ast}$-sign with
  itself, from an expected equivariance $(-1)^{n}$ read off the coefficient level alone.
  That omitted the orientation factor: $z\mapsto-\bar z$ reverses the orientation of
  $\Hyp$, so any pairing through a fundamental class of the pair --- which is what
  \cl{cl:pd}(1) asserts --- picks up a further $-1$, the expected equivariance is
  $-(-1)^{n}$, equal signs annihilate for $n$ even, and duality pairs $+$ with $-$
  (the modular-symbol picture, \cite[Ch.~2]{cremona}). The remark now says so, and adds
  the dimension cross-check from this document's own predictions: $(d_{\,n+2},d_{\,n+2}-1)$
  on $H^{1}_{\mathrm{eq},\pm}$ (\cl{obl:iota}(v)) against the mirror split on $H^{1}_{c}$
  forced by the odd one-dimensional $H^{0}(\partial;\Mn)$, which self-pairing could not
  match. Its status is unchanged: contingent, ``asserted here as nothing more'', no tag.
  (Review R4, adjudicated by the referee: the reviewer's sign is right, and the
  contract's own dimension predictions independently exclude the old direction.)
\item \cl{rmk:iotaelltwo} --- closing sentence: ``open construction item \ldots\ a tagged
  working assumption'' (stale since v1.3.0) $\to$ ``decided by the owner on 2026-09-16
  (\cl{oc1a})''. Remark, commentary (R5).
\item The remark before \cl{obl:coneexact} that names \S\ref{sec:cone}'s two reading
  sites (container) --- ``open construction item'' $\to$ decided, with the date and
  \cl{oc1a} (R6).
\item \cl{sec:oc1sites} --- item (7) gains the two status sites above, marked as missed
  at v1.3.0 (R5, R6). Container.
\item \cl{sec:oraclesnew} --- family (C): ``on every cell'' $\to$ ``on every cell on
  which it can be located (the Dickson cells carry no sign)'' (P4). Container.
\item \cl{err:e12} --- ``what is not settled at $\ell=2$'' now also names the policy
  decided for it (P9). Erratum entry, no tag.
\item Appendix~\ref{app:changelog}, the provenance table and prose --- \textbf{recounted
  from Appendix~\ref{app:status}} ($121$ \BNEW{} rows: $52$ \code{rmk} and $69$ others,
  namely $21$ \code{def}, $19$ \code{prop}, $15$ \code{lem}, $4$ \code{cl}, $4$
  \code{obl}, $2$ \code{thm}, $1$ each of \code{alg}, \code{cor}, \code{ex}, \code{open}):
  the table's ``$70$ assertions or definitions, $51$ remarks'' becomes $69$ / $52$,
  agreeing with the roster below it (which lists $69$ and was right); the prose
  ``$118$ / $49$ / excess over $111$ and $35$'' and the heading ``The $68$ statements''
  --- stale since before v1.3.0 --- become $121$ / $52$ / $121$ and $45$ / $69$ (R10).
  \ref{CL6} no longer says ``ten commits'' and says what \code{[hash]} means; the
  \code{[hash]} placeholders of v1.1.1, v1.2.0, v1.3.0 and v1.3.1 are back-filled
  (\texttt{a047147}, \texttt{c823f3a}, \texttt{b1c0aaa}/\texttt{4f84a32},
  \texttt{f69cda0}; P11); the ``$Y$ field'' sentence under the rules lists v1.0.0,
  v1.1.1, v1.3.1 and v1.3.2 (R9, completed).
\item The date line (companion pointer \file{hecke\_extensions.tex} v1.3.2, P1) and the
  header comment. Container, no tag.
\item \textbf{Outside this document} (\ref{CL2}): \file{hecke\_extensions.tex} v1.3.1
  $\to$ v1.3.2 --- its (A2) closing sentence no longer calls the truncation policy
  ``open \ldots\ not decided here'' but points at [C] \cl{open:ell2} / \cl{oc1a} (R8),
  its item \code{add:iota}'s ``not a pullback along any lift'' gains ``for $n$ even''
  (P5), and
  ``$\iota$ is not implemented'' becomes ``no production implementation; built and
  \cl{obl:iota} discharged by two validations'' (P6); PDF rebuilt.
  \file{hecke\_extensions\_implementation\_plan.md} risk R4(ii) and the Q2 row (R7, P10);
  \file{hecke\_extensions\_owner\_decisions.md} summary row 3 (P7). Not applied: P8, a
  comment in \file{HeckeIotaSignValidation.cpp} --- no C++ change in a documentation pass.
\end{itemize}

\subsection*{v1.3.1 --- 2026-09-16 --- referee pass on \texorpdfstring{\cl{lem:bulkperiodic}}{lem:bulkperiodic}: verified, tag kept, prose repaired}

\noindent\emph{This version.} The audit the owner decision of 2026-09-16 asked for
(``write the periodicity argument and then audit / referee the result, including Hecke
actions''), executed on \cl{lem:bulkperiodic} as written at v1.3.0. \textbf{Changes no
normative statement} (\ref{V2}, \ref{CL5}): every clause of the lemma and every step of
its proof was re-derived independently --- the five lemma and its lower edge, the cyclic
periodicity with its degree-$0$ kernel $N\,M$, the well-posedness and
conjugation-invariance of $\zeta_+$ (in particular that $S$ and $S^{-1}$, with
$\zeta_+=\pm i$, are not conjugate in $\GL_2(\Q)^{+}$, so the definition is consistent),
$H^{2}(\Ga;\Z)\cong\Z/12$, the detection of $H^{2}(\Ga'';\Z)$ on finite cyclic
subgroups, the projection formula and the compatible-pair bookkeeping of
\eqref{eq:corconjres} --- and checked against $H^{\bullet}(\Ga;\F_2)$,
$H^{\bullet}(\Ga;\F_3)$ and the $6\in\cring^{\times}$ case. No error was found; the tag
\PROVED{} stands, the \BSTD{} inputs being cited at the level \S\ref{sec:status}
permits. Four prose defects were repaired, none touching a hypothesis, a displayed
formula or a conclusion. Report:
\file{docs/implementation/referee\_report\_2026-09-16\_bulkperiodic.md}. Commit:
\texttt{f69cda0}.

\paragraph{Normative.}
\textbf{None.} \ref{CL5}'s test is met vacuously.

\paragraph{Non-normative.}
\begin{itemize}\setlength{\itemsep}{0pt}
\item \cl{lem:bulkperiodic} --- prose only, three places. (i) The ``External inputs''
  preamble claimed that every unverified pointer is ``a direct check on the cochain
  formula''; that is true of (E1) and the first two items of (E2) and false of (E2)'s
  last item, (E3) and (E4), which are now said to be used as stated, (E4) as the meaning
  of $\cor$ in \eqref{eq:corconjres}; the preamble also now names
  \cite[Ch.~III.10.2]{brown} (Steps 0, 5) and the in-place proof of the fixed-vertex
  statement (Step 3), which the count ``four standard facts'' had left out. (ii) Clause
  (3): the operator is now said to be \eqref{eq:corconjres} on ordinary group
  cohomology, defined for every $\cring$, $n$ and $\ell$ with no model, the pointers to
  Theorem~\ref{thm:classical} and Remark~\ref{rmk:borelfree} being identifications the
  clause does not depend on --- where v1.3.0 called it ``the Hecke operator of the free
  $\Ga$-model'' and cited Theorem~\ref{thm:specialisation}(3), a $6\in\cring^{\times}$
  statement (rule~\ref{A1}: the ring and $n$ are now in the clause). (iii) The closing
  sentence of (3): $\Amat{q}(\Tdc{g})$ is glossed as the matrix on $H^{q}(\Ga;\VR)$, not
  the shipped model's at $\ell\in\{2,3\}$. Lemma, tag unchanged.
\item \cl{rmk:ell2window} --- clause (2): ``proved from Claim~\ref{cl:amalgam}(2)'' now
  says from the tree sequence \eqref{eq:treeses} of \cl{cl:amalgam}(1), the
  Mayer--Vietoris sequence re-derived for an arbitrary coefficient module
  (\cl{cl:amalgam}(2) is stated for $\VR$ only). Remark, commentary.
\item \cl{rmk:bulkperiodicscope} --- (3)(b): the identification
  $\Hcpt{\bullet}=H^{\bullet}_{\Ga}(Y,\partial Y;\VR)$ now points at
  Definition~\ref{def:borelmodel}, Theorem~\ref{thm:specialisation}(2) being a
  $6\in\cring^{\times}$ statement; the closing sentence of (3) said a run with
  $q_{\max}<4$ ``cannot exhibit one full period of the cone'', which misstates what
  $q_{\max}=4$ shows --- it exhibits the first bulk isomorphism $H^{2}\cong H^{4}$, while
  the cone's first, $\Hcpt{3}\cong\Hcpt{5}$, needs $q_{\max}\ge5$; the sentence now says
  so. Remark, commentary.
\item \cl{open:ell2} --- the \emph{non-normative recommendation} paragraph: ``one full
  period of the bulk needs degrees $2$ to $5$ built, i.e.\ $q_{\max}\ge4$'' was
  internally inconsistent (degree $5$ is built as scaffolding and never reported, so
  $H^{3}\cong H^{5}$ is not observed at $q_{\max}=4$); it now says which isomorphism
  $q_{\max}\ge4$ exhibits and that the second and the cone's first need $q_{\max}\ge5$.
  The recommended value $q_{\max}\ge4$, the box, the riders and the decision are
  untouched. Policy, no tag.
\item Appendix~\ref{app:changelog} --- the sentence ``the $Y$ field has never yet been
  used'' under the rules, stale since v1.1.1, corrected. Container, no tag.
\item The date line and the header comment (version $1.3.1$). Container, no tag.
\item \textbf{Outside this document} (\ref{CL2}): the referee report
  \file{docs/implementation/referee\_report\_2026-09-16\_bulkperiodic.md} --- verdict per
  claim, the sanity computations $H^{\bullet}(\Ga;\F_2)$, $H^{\bullet}(\Ga;\F_3)$,
  $H^{2}(\Ga;\Z)$ written out, the external facts actually needed with precise
  statements, and what a phase-E5 computation on the free model could test.
\end{itemize}

\subsection*{v1.3.0 --- 2026-09-16 --- OC-1 decided: \texorpdfstring{\cl{open:ell2}}{open:ell2} becomes the binding truncation policy, and the bulk periodicity above the vcd is proved, Hecke-equivariantly}

\noindent\emph{This version.} Executes the owner decision of 2026-09-16 on the
$\ell=2$ truncation policy --- ``OC-1: adoption per brief, write the periodicity
argument and then audit / referee the result, including Hecke actions'' --- on the
decision brief \file{docs/implementation/oc1\_decision\_brief.md}: \cl{oc1a} is adopted
with the brief's three riders, and the argument the brief sketched in its \S4.2 is
written into the contract as \cl{lem:bulkperiodic}, in a stronger form than sketched
(a canonical class, and the Hecke equivariance the brief listed as the open half). The
audit/referee pass the decision asks for is the \emph{next} slot, not this one; the
lemma is written to be refereed and says at each step which standard fact it uses.
\textbf{This is an $X$ bump} (\ref{V2}, \ref{CL5}): \cl{open:ell2} loses its
\code{[open]} tag, and \S\ref{sec:status} makes the tag part of the statement; its
boxed content becomes binding rather than assumed; \cl{obl:coneexact}(iv), an
obligation, changes wording; the record schema gains its adopted fields; and a new
normative statement, \cl{lem:bulkperiodic}, is added. \textbf{Ordering note}: v1.2.0 ---
OQ3, the $\inv$ sign and the $\pm$ decomposition of Chapter~\ref{sec:involution} ---
was prepared \emph{in parallel} with this version, from the same v1.1.1 base, and
integrated first; neither version edits a statement of the other, the two entries are
independent, and the date line and this appendix list them in integration order
(v1.3.0 above v1.2.0). No erratum is filed:
\ref{CL2} reserves Appendix~\ref{app:errata} for the superseded corpus, and \ref{L2}
is not engaged, every label keeping its subject. Commit: \texttt{b1c0aaa} (integrated
with v1.2.0, PDFs rebuilt, at \texttt{4f84a32}).

\paragraph{Normative.}
\begin{itemize}\setlength{\itemsep}{0pt}
\item \cl{open:ell2} --- retitled ``OC-1: the $\ell=2$ truncation policy --- DECIDED
  2026-09-16''; the tag \code{[open]} and the header ``reserved by the owner; a tagged
  working assumption, not a decision'' are replaced by the decision sentence and its
  date; the boxed statement is adopted \textbf{verbatim} and followed by the three
  riders ($q_{\max}$ derived on a finite model and settable only where
  \code{truncated = true}, distinct from the reported-degree list; the periodicity is
  a theorem; \cl{obl:coneexact}(iv) carries the box's qualification) and by the
  non-normative recommendation $q_{\max}\ge4$ at $\ell\in\{2,3\}$; ``The alternatives,
  kept reachable'' becomes ``\cl{oc1a} taken, \cl{oc1c} kept reachable''; \cl{oc1c}'s
  ``the Farrell--Tate route named, and deliberately not asserted'' and ``Neither
  hypothesis is proved anywhere in this document'' become a pointer to
  \cl{lem:bulkperiodic} and the precise residue of \cl{rmk:bulkperiodicscope}(3);
  ``Recommendation, and why'' becomes ``Decision, and why''. The labels \cl{oc1a}--\cl{oc1d}
  are kept (\ref{L1}). Open construction item; status tag \code{[open]} removed, the
  truncation lemma's \PROVED{} kept; \PNEW{} content unchanged in its box (\ref{CL4}).
\item \cl{lem:bulkperiodic} \PNEW{} --- \textbf{new}: for
  $\Ga=\langle S\rangle\ast_{\langle-I\rangle}\langle U\rangle$ and \emph{any}
  $\Z[\Ga]$-module $V$, cup product with an explicitly constructed generator
  $e\in H^{2}(\Ga;\Z)\cong\Z/12$ --- the unique class whose restriction to each finite
  cyclic subgroup is the eigenvalue character $\zeta_+$ on the fixed point in $\Hyp$ ---
  is an isomorphism $H^{q}(\Ga;V)\to H^{q+2}(\Ga;V)$ for every $q\ge2$, a surjection with
  kernel $\partial(N_{-I}V)$ for $q=1$; $e$ is $\De$-invariant, and every
  $\Tdc{g}=\cor\circ(\conj_g)_{\ast}\circ\res$ commutes with $e\cup-$, so the Hecke module
  $H^{q}(\Ga;\VR)$ is $2$-periodic in $q\ge2$ for every $\cring$ and $n$. Proof by the
  five lemma on \cl{cl:amalgam}(2), a dimension-shifting proof of the periodicity of
  finite cyclic groups, detection of $H^{2}$ of a finite-index subgroup on its finite
  cyclic subgroups via the Bass--Serre tree, and the projection formula
  $\cor(\res u\cup y)=u\cup\cor y$ derived from Shapiro's lemma. Four standard inputs
  are cited at section level (Brown Ch.~V.3, III.5--6, III.8--9; Serre, \emph{Trees}
  I.4.3) with rule~\ref{A2} caveats where the pointer is unverified. Ordinary
  cohomology throughout; Farrell cohomology is neither used nor asserted about.
\item \cl{rmk:ell2window} --- clause (2): the parenthesis ``the standard route would be
  Farrell--Tate cohomology \ldots\ \code{[conjectural]} as a citation: this document has
  not verified the precise form'' is replaced by the statement of what
  \cl{lem:bulkperiodic} proves and by the note that Farrell theory is not used. The
  remark loses its only tag and becomes commentary; \ref{L2} is not engaged, the clause
  now asserting \emph{more}, not something different. Remark, tag \code{[conjectural]}
  removed.
\item \cl{obl:coneexact} --- clause (iv): ``Each composite of consecutive maps of
  \eqref{eq:coneles} is zero in every computed degree'' becomes ``each composite both of
  whose end terms are reported is zero'', with the truncated-model reading spelled out
  (only the incoming composite at $H^{q_{\max}}(\partial)$), and two \textbf{recorded}
  quantities at that spot: $\rank\rres^{q_{\max}}$ and the injection
  $\coker\rres^{q_{\max}}\hookrightarrow\Hcpt{q_{\max}+1}$, the latter at $\ell=2$ a
  recorded lower bound and never a reported group (the brief's \S4.3 under-claim).
  Obligation, predicate narrowed (\ref{CL5}: normative).
\item \S\ref{sec:implementation}, record schema item~5 (container, no tag) --- ``on a
  truncated free model they are the parameter \cl{open:ell2} fixes \ldots\ which is why
  the shipped writer does not yet carry them'' becomes the adopted fields
  (\code{q\_max}, \code{truncated}, the reported-degree list, the derived/set marker) and
  a concrete instruction for phase E5 naming \file{Hecke/HeckeReport.\{h,cpp\}} and
  \file{Campaign/CampaignDriver.cpp}, and requiring the degree list and \code{q\_max}
  in a record's identity.
\end{itemize}

\paragraph{Non-normative.}
\begin{itemize}\setlength{\itemsep}{0pt}
\item \cl{rmk:bulkperiodicscope} \PNEW{} --- new commentary: ordinary versus Farrell
  cohomology; the (unverified, unused) identification of $e$ with the Euler class of
  $B_3\to\SL_2(\Z)$; what \cl{lem:bulkperiodic} makes of \cl{oc1c} and the three things
  that remain open in it (a boundary periodicity isomorphism --- candidate
  $\res^{\Ga}_{\Gbdry}e\cup-$, automatically compatible with $\rres$ ---, the
  $H^{\bullet}(\Ga;\Z)$-module structure of \eqref{eq:coneles} on the cone, and the Hecke
  equivariance of both); why $q_{\max}\ge4$ is the observing bound.
\item \cl{sec:oc1sites} --- item (1) reworded for a decided item; a seventh item added
  listing the sites that assert the item's \emph{status} (the brief's finding: the list
  enumerated content sites only). Container, no tag.
\item The status-asserting sites themselves (item (7) of \cl{sec:oc1sites}): the
  abstract; \cl{sec:scope}'s closing paragraph; \ref{V5}; \cl{sec:model}'s $\ell=2$
  paragraph; \cl{sec:boundarydatum}'s ``What this chapter fixes''; the opening of
  \cl{sec:ell2}; Appendix~\ref{app:formalization}'s ``A decision on \cl{open:ell2}'';
  \S\ref{app:changelog-gaps}. Each now says ``decided 2026-09-16'' where it said
  ``open''. Containers, no tag.
\item Appendix~\ref{app:status} --- rows for \cl{lem:bulkperiodic} and
  \cl{rmk:bulkperiodicscope} added; \cl{rmk:ell2window}'s row \code{[conjectural]}
  $\to$ commentary; \cl{open:ell2}'s row loses \code{[open]}; the counts paragraph
  ($201$ statements, $23$ \code{lem}, $60$ \code{rmk}; \PROVED{} $92$, \BSTD{} $34$,
  \CONJ{} $7$, \OPENSTAT{} $6$; $88$ untagged) and the provenance box ($121$ \BNEW).
\item \ref{CL4} --- the list of statement kinds gains \code{open}, as the v0.14.0 entry
  asked; the provenance table ($121$ \PNEW, $70$ assertions or definitions, $51$
  remarks) and the roster of new statements (Part~IV, the boundary: $10$) follow.
\item The date line and the header comment (version $1.3.0$, with the concurrency note
  on v1.2.0). Container, no tag.
\item \textbf{Outside this document} (\ref{CL2}; recorded where they happened): the
  OC-1 entry in \file{docs/implementation/hecke\_extensions\_owner\_decisions.md}, the
  ``DECIDED'' note at the head of the brief, and the pointers in
  \file{hecke\_extensions\_implementation\_plan.md}, \file{hecke\_extensions\_open\_questions.md},
  \file{docs/console/00-console-interface-plan.md} and
  \file{lean\_formalization\_assessment.md}.
\end{itemize}

\subsection*{v1.2.0 --- 2026-09-16 --- OQ3 resolved: \texorpdfstring{\cl{def:iota}}{def:iota}(iv) negated, ``$+$'' is by definition the Eisenstein side}

\noindent\emph{This version.} Executes the owner decision of 2026-09-16 on OQ3, taken on
the v1.1.1 evidence: ``negate \cl{def:iota}(iv) in the following sense: we \emph{define}
`$+$' to be the sign where Eisenstein classes live; this means that we have to switch
$+/-$ notationally everywhere.'' Of the three realizations weighed in
\file{docs/implementation/oq3\_resolution.md} --- global negation, notation-only, and a
degree-parity twist --- the owner's own (global negation) is the cheapest and is taken:
$\inv^{\ast}:=-\inv^{\ast}_{\mathrm{geo}}$ in every degree, the geometric pullback kept under
its own name. Its honest cost, stated in the definition: $\inv^{\ast}$ is the geometric
involution twisted by the sign character $\det$ of $\GL_2(\Z)/\SL_2(\Z)$, not the pullback
along any lift of $z\mapsto-\bar z$; constants are odd ($\inv^{\ast}=-\id$ on
$H^{0}(\Ga;\Fl)$); on the boundary $H^{0}(\partial)$ is odd and $H^{1}(\partial)$ ($\ell>n$)
even; at $\ell=2$ nothing changes. What it buys: every $\pm$ prediction the document had
written from OQ3 --- $(d_{\,n+2},d_{\,n+2}-1)$, Eisenstein in $H^{1}_{\mathrm{eq},+}$ ---
holds as printed, the former recorded half of \cl{obl:iota}(v) becomes a gating
definition-consistency oracle, and no PROVED statement of chapter~\ref{sec:involution}
needs more than ``the sign cancels''. No pipeline value moved: $\inv^{\ast}$ occurs in no
golden object, and \code{GoldenRegressionValidation} was run in the same process as the
changed validation. The validation is
\code{HeckeIotaSignValidation} (tier \code{standard}, $1.4$\,s, $4\,104\,267$ exact
assertions, $0$ gating failures; it now builds both involutions and prints
\code{IOTA\_SIGN eisenstein=plus (geometric=minus) ... h0\_anchor=-1 (geometric=+1)}).
$X$ increments (\ref{V2}, \ref{CL5}): a definition, an obligation and a remark change.
Commit: \texttt{c823f3a}.

\paragraph{Normative.}
\begin{itemize}\setlength{\itemsep}{0pt}
\item \cl{def:iota} \PNEW{} --- clause (iv) restated: \cl{eq:iotacochain} now defines
  $\inv^{\ast}_{\mathrm{geo}}$ (the former $\inv^{\ast}$) \emph{and} $\inv^{\ast}:=-\inv^{\ast}_{\mathrm{geo}}$,
  for every degree and every coefficient module ($\Dn$ included); $\inv^{\ast}$ is declared
  the involution of the document; the sign's nature ($(\det\otimes\rhon)(\inv)$; an
  automorphism of $(K,\Rel,L)$ but no lift of $z\mapsto-\bar z$), its consequences
  ($-\id$ on $H^{0}(\Ga;\Fl)$; invisible at $\ell=2$) and the flat-coordinate block
  $\mp\rhon(\inv)$ are stated. Clauses (i)--(iii) unchanged.
\item \cl{obl:iota} \PNEW{} --- clause (v) rewritten: the sign statement (``the Eisenstein
  class lies in $H^{1}_{\mathrm{eq},+}$'', by both routes, on every signed cell, with
  $\inv^{\ast}=-\inv^{\ast}_{\mathrm{geo}}$ entrywise) moves from \emph{recorded} to
  \emph{aborting} as a definition-consistency oracle, and what a failure means is stated
  (a non-uniform geometric sign, i.e.\ an ill-posed definition); the degree-$0$ sentence
  becomes $-\rhon(\inv)$ on the invariants, $-\id$ at $n=0$; the $\pm$ dimension
  statement stays recorded, with the reason (it is the even cuspidal split, the \CONJ{}
  Eichler--Shimura prediction). The \COMP{} record is rewritten for the v1.2.0 run, the
  geometric numbers kept as the recorded reading; \code{[OQ3-PENDING]} removed.
\item \cl{rmk:iotasign} \PNEW{} (\ref{CL3}) --- rewritten: the sign was an assertion about
  the geometric involution and was checked (the mechanism named); the resolution negated
  the definition and not the record; \code{iota\_sign} now \emph{names the
  normalisation} $\inv^{\ast}=-\inv^{\ast}_{\mathrm{geo}}$, by which a third party reproduces
  a record; what is not chosen (uniformity of the geometric sign, now gated by
  \cl{obl:iota}(v)). Its prohibition on call-site flips is unchanged.
\end{itemize}

\paragraph{Non-normative.}
\begin{itemize}\setlength{\itemsep}{0pt}
\item \cl{prop:iotacommutes} --- statement (2) names $\inv^{\ast}_{\mathrm{geo}}$ alongside
  $\inv^{\ast}$ and transports through $\inv^{\#}:=-\inv^{\#}_{\mathrm{geo}}$; the proofs of
  (2) and (3a) are written for the geometric pair with the one-line reduction (the
  sign cancels; $(-\id)^2=\id$). Content of the proposition unchanged, \PROVED{} tag
  unchanged, so no normative line (the statement's predicates are the same).
\item \cl{rmk:iotaelltwo} --- one sentence: at $\ell=2$, $-1=+1$, so
  $\inv^{\ast}=\inv^{\ast}_{\mathrm{geo}}$ and the normalisation is invisible. Remark, no tag
  change.
\item \cl{rmk:iotapairing} --- one sentence: the sign applies to both coefficient
  modules, the two factors $-1$ cancel in the pairing, and every $\pm$ consequence is
  invariant under negating both involutions at once. Remark, no tag change.
\item \cl{sec:involution} --- the ``What this chapter fixes'' paragraph (the sign fixed,
  its cost named); the $\pm$ subsection's opening paragraph (the eigenspaces are the
  geometric ones with labels exchanged; how to read a pre-v1.2.0 $\pm$); the
  Eichler--Shimura paragraph after \cl{obl:iota}; the code map (both involutions built,
  the sign gating, the production block $\mp\rhon(\inv)$). Container, no tag
  (\ref{CL4}).
\item \cl{sec:lateroblig}, \cl{sec:obligationroster} (the \cl{obl:iota} row and the
  ``Reading the table'' paragraph), \cl{sec:oraclesnew} (opening and family (C), now
  carrying both runs), \cl{sec:codemap} (the \cl{sec:involution} row),
  \cl{sec:notimplemented} (chapter~\ref{sec:involution} has no open item) --- each
  updated from ``pending'' to the resolution. Containers, no tag.
\item The date line and Appendix~\ref{app:changelog}'s closing paragraph. Container,
  no tag.
\item Appendix~\ref{app:status} is \textbf{unchanged}: \cl{def:iota} stays a definition,
  \cl{obl:iota}'s line already reads \CONJ{} \COMP\textsuperscript{r} \BSTD{} (the
  \CONJ{} is the dimension prediction, still recorded), and \cl{rmk:iotasign} stays
  commentary.
\item Outside this document, in the same commit: \file{hecke\_extensions.tex} Stage~0
  item (A3) and the \code{iota\_sign} paragraph of its \S7;
  \file{hecke\_extensions\_owner\_decisions.md} \S OQ3 (the decision of 2026-09-16 with
  the verdict that motivated it), \file{hecke\_extensions\_open\_questions.md} items~3
  and~18, the implementation plan's C-A1 and Q3 rows;
  \code{HeckeIotaSignValidation.\{h,cpp\}}, its registry comment and
  \file{Hecke/README.md}; and the banner on
  \file{docs/implementation/oq3\_iota\_sign\_verification.md} saying its $\pm$ are
  geometric.
\end{itemize}

\subsection*{v1.1.1 --- 2026-09-16 --- \texorpdfstring{\cl{obl:iota}}{obl:iota}(ii),(iii),(v) run: status only, and the OQ3 sign comes out \texorpdfstring{$-$}{minus}}

\noindent\emph{This version.} \textbf{Changes no normative statement of this
document}, which is why $Y$ and not $X$ increments (\ref{V2}, \ref{CL5}): no clause was
added, removed or amended, no predicate, hypothesis, ring or action-on-failure moved,
and no displayed formula changed. What it records is \emph{evidence}. The three clauses
of \cl{obl:iota} that Chapter~\ref{sec:involution}'s code map had listed as
unimplemented --- (ii), (iii) and (v) --- were implemented and run by the new registered
validation \code{HeckeIotaSignValidation}
(\file{src/hacs/Hecke/HeckeIotaSignValidation.\{h,cpp\}}, tier
\code{standard}, $1.4$\,s, $4\,103\,488$ exact assertions, $0$ gating failures), which
builds $\inv^{\ast}$ from \cl{def:iota} alone. \textbf{Every aborting predicate of the
obligation passes.} The half of clause (v) that the clause itself marks
\emph{recorded rather than aborting} --- owner decision OQ3's ``the Eisenstein class
lies in $H^{1}_{\mathrm{eq},+}$'' --- is \textbf{false}: on all $74$ cells of a $77$-cell
grid ($n\in\{2,\dots,30\}$ even, $\ell\in\{5,7,11,13,101,134217757,2^{31}-1\}$) on which
a sign is defined, and by two independent routes that agree wherever both are defined,
the class lies in $H^{1}_{\mathrm{eq},-}$, and on all $66$ off-Dickson cells the $\pm$
split is the exact mirror $(d_{\,n+2}-1,\,d_{\,n+2})$ of the printed prediction. Per
\cl{rmk:iotasign} that is resolved by an owner decision on a normative statement ---
negate \cl{def:iota}(iv), or amend OQ3 --- and \textbf{the decision is pending}: this
version takes neither, tags the place \code{[OQ3-PENDING]}, and leaves the predicted
numbers standing next to the observed ones so that the two can be compared. Nothing was
flipped at a call site, which is the failure mode \cl{rmk:iotasign} exists to forbid.
The evidence, the full $77$-cell table and both resolutions in the owner's own terms:
\file{docs/implementation/oq3\_iota\_sign\_verification.md}. Commit: \texttt{a047147}.

\paragraph{Normative.}
\textbf{None.} \ref{CL5}'s test is met vacuously: no entry below is a change to a
statement's content.

\paragraph{Non-normative.}
\begin{itemize}\setlength{\itemsep}{0pt}
\item \cl{obl:iota} --- a \COMP{} evidence record added to clause (v), reporting the run
  above: what passes, what is recorded false, the two dimensions as observed, and the
  \code{[OQ3-PENDING]} tag. The clause's own text --- its predicate, its
  $6\in\cring^{\times}$ hypothesis, its displayed prediction, its action on failure and
  its degree-$0$ sentence --- is untouched, so this is evidence about a statement and
  not a change to one. Obligation, no tag change (\ref{CL4}).
\item \cl{rmk:iotasign} --- one bullet added: the check it specifies has been run and
  fails, so the choice between its two named resolutions is now live. The remark's
  content is unchanged, because it is precisely the statement that says what such a
  failure means and what may not be done about it. Remark, no tag change.
\item \cl{sec:involution} --- the chapter's ``What this chapter fixes'' paragraph (all
  five clauses now implemented and run, with the one recorded failure named) and its
  code map (the second validation, what it builds and which clauses it discharges; the
  census re-run; the \code{EquivariantCochainComplexMorphismOnK} paragraph re-labelled
  as the specification for a \emph{production} $\inv^{\ast}$, which still does not
  exist). Container, no tag (\ref{CL4}).
\item \cl{sec:lateroblig} --- the \cl{obl:iota} bullet gains the (ii),(iii),(v)
  paragraph. Container, no tag.
\item \cl{sec:obligationroster} --- the \cl{obl:iota} row's status and validation
  columns, and the ``Reading the table'' paragraph, which now describes that row as
  evaluated in full with two negative records rather than as partly unimplemented.
  Container, no tag.
\item \cl{sec:oraclesnew} --- the section's opening (family (C) is a run, not a
  specification) and family (C) itself, which gains a \COMP{} paragraph with the
  observed split and loses the stale sentence calling the \cl{obl:iota}(iv) diagnostic
  false without saying that v1.1.0 repaired it. Container, no tag.
\item \cl{sec:codemap} --- the \cl{sec:involution} row of the chapter-to-file table
  names both validations; the \cl{sec:obligations} row's \code{Hecke*Validation} count
  moves from $27$ to $29$. Container, no tag.
\item \cl{sec:notimplemented} --- Chapter~\ref{sec:involution} is removed from the list
  of chapters whose obligations are unimplemented; what it still lacks is named
  precisely (a production $\inv^{\ast}$, and an owner decision). Container, no tag.
\item The count of \textbf{registered validations} moves from $67$ to $68$ at its six
  occurrences (\cl{sec:involution}'s code map, \cl{sec:boundarydatum}'s and
  \cl{sec:cone}'s census paragraphs, \cl{sec:obligationroster},
  \cl{sec:notimplemented}, and Appendix~\ref{app:formalization}'s
  ``second, independent witness''). Counts, no tag.
\item The date line and Appendix~\ref{app:changelog}'s closing paragraph (``not yet
  written''), which now names \cl{obl:iota}(v)'s pending decision alongside
  \cl{obl:iota}(iv)'s two-version record. Container, no tag.
\item Appendix~\ref{app:status} is \textbf{unchanged}: it is mechanically derived from
  the tags a statement carries, \cl{obl:iota}'s line already reads
  \CONJ{} \COMP\textsuperscript{r} \BSTD{}, and this version adds no tag to any
  statement --- the new \COMP{} record is a second \COMP\textsuperscript{r} on a line
  that already has one, and \cl{rmk:iotasign} stays commentary. The
  ``eight of the thirteen'' \COMP{} count is therefore also unchanged.
\end{itemize}

\subsection*{v1.1.0 --- 2026-09-16 --- the canonical $\tau$ made \texorpdfstring{$\inv$}{iota}-symmetric: owner decision, re-baseline}

\noindent\emph{This version.} Executes the owner decision of 2026-09-16 that
\cl{rmk:iotatausym} had recorded as available and not taken: the canonical $\tau$ of
\cl{alg:tau} is rebuilt so that $\inv_{\Xbulk}\tau\inv_{\Xbulk}=\tau$, which makes
\cl{prop:iotacommutes}(3b) --- the only route to $[\inv^{\ast},\Tdc{g}]=0$ at
$\ell\in\{2,3\}$ --- available for it. The feasibility pass found that the remark's
proposed repair could not work as stated (no choice \emph{within} an $\invK$-swapped
tied pair commutes with $\invK$; the witness point $4i$ lies on the mirror), and that
step~2 of \cl{alg:tau} was a second obstruction the degree-$0$ witness could not show;
both are repaired and both are recorded. Every frozen cochain-level reference value
(\file{Golden/GoldenReference.inc}: $2097$ of $14049$ lines --- every
$\Th^{(q)}_{\mathrm{flat}}$, its rank, every $\Th^{(q)}_{\mathrm{eq}}$, the JSON records)
was re-baselined and \emph{no} cohomology-level value moved, $\ell\in\{2,3\}$ included;
the object-by-object record is
\file{docs/matrix-engine/rebaseline-2026-09-16-iota-symmetric-tau.md}. The pre-v1.1.0
$\tau$ is kept in the tree as \code{TauVertexPolicy::GeometricNearestLegacy} and re-run
by the same validation, so the v0.13.0 numbers stay reproducible. $X$ increments
(\ref{V2}): \cl{alg:tau} and \cl{obl:iota} change. Commit: \texttt{a047147}.

\paragraph{Normative.}
\begin{itemize}\setlength{\itemsep}{0pt}
\item \cl{alg:tau} \PNEW{} --- step~(1b): ties among nearest vertices are broken
  $\invK$-equivariantly (a tied set containing a full $\invK$-orbit goes to $v_6$, the
  unique $\invK$-fixed vertex; otherwise to its member in the orbit of smallest index),
  in place of ascending vertex index. Step~2: the path templates are tree-routed ---
  spoke to the Bass--Serre tree $Y$ (the $\Ga$-orbit of $\{v_4v_6\},\{v_5v_6\}$), the
  unique $Y$-path, spoke back --- so that concatenation along any word telescopes to a
  chain of the two endpoints alone, $\Ga$- and $\inv_{\Xbulk}$-equivariant and
  antisymmetric, in place of breadth-first paths in ascending neighbour order. Neither
  change touches \cl{def:tau}, \cl{obl:tau} or step~4; the rule is new to this document
  and to every source.
\item \cl{obl:iota} \PNEW{} --- clause (iv) amended: still recorded and still licensing
  nothing beyond (3b), but on the canonical $\tau$ of \cl{alg:tau} a false value is now a
  \emph{failure} (a regression gate on the construction, not a branch on the diagnostic);
  on any other admissible $\tau$ it is not. Its \COMP{} record moves from ``discharged,
  negatively'' to ``negatively at v0.13.0 on the then-canonical $\tau$, positively since
  v1.1.0 on the canonical $\tau$'': $0$ of $21528$ slots over $144$ tables, exact over
  $\Z$; the legacy $\tau$ reproduces its $3600$ of $10764$.
\end{itemize}

\paragraph{Non-normative.}
\begin{itemize}\setlength{\itemsep}{0pt}
\item \cl{rmk:iotatausym} \PNEW{} (\ref{CL3}) --- rewritten a second time. The v0.13.0
  finding is kept verbatim as history (grid, counts by degree, the witness, the
  mechanism); added: why no tie-break of the tied set can repair it (the mirror
  argument, with $i$ the only vertex of $\Xbulk$ on the mirror), the three tied
  configurations that actually occur, the repair taken, the measurement that step~1
  alone left $1848/5148$ and $1068/2340$ differing slots in degrees $1$ and $2$ (the
  second obstruction, in step~2), the verdict since v1.1.0 with its two stated limits
  (verified per table by a run, not proved for every $(p,\nu)$; $\inv^{\ast}$ on cochains
  still not implemented), the consequence $[\inv^{\ast},\Tdc{g}]=0$ at $\ell\in\{2,3\}$
  for the canonical $\tau$ on the grid, the vanishing of the canonical $\tau$ in degree
  $2$, and the re-baseline's counts. The sentence ``a repair exists and is not taken
  here'' is withdrawn: the repair as proposed there did not exist, and the one that does
  is taken.
\item \cl{rmk:certoverGeometry} \PNEW{} --- item (1) revised: the (1b) tie-break and the
  step-2 templates are now $\invK$-equivariant and the canonical $\tau$ is
  $\inv$-symmetric; the pre-v1.1.0 choices were not, each on its own.
\item \cl{rmk:tauexistsscope} --- one sentence added to the \CONJ{} paragraph: the two
  certified $\tau$ the validation holds induce the \emph{same} $\Amat{q}$ at every
  $\ell\in\{2,3\}$ point of the golden grid, so no counterexample to homotopy uniqueness
  at $\ell\in\{2,3\}$ is known there. Remark, no tag change.
\item \cl{sec:involution} --- the chapter's ``What this chapter fixes'' paragraph
  (clause (iv) now holds on the canonical $\tau$ and gates it) and its code map
  (three policies, $432$ tables, $64699$ assertions; the recorded verdict positive on
  the canonical $\tau$, negative on the legacy one). Container, no tag (\ref{CL4}).
\item \cl{sec:lateroblig} --- the clause-(iv) paragraph of the obligations summary,
  from ``discharged negatively'' to the two-version record and the route now open at
  $\ell\in\{2,3\}$. Container, no tag.
\item \cl{sec:codemap} --- the \cl{alg:tau} line names both template families; the
  \cl{sec:involution} row of the chapter-to-file table records that (iv) gates.
  Container, no tag.
\item Appendix~\ref{app:changelog}'s closing paragraph (``not yet written''): (iv) is no
  longer listed as discharged negatively. Appendix, no tag.
\end{itemize}

\subsection*{v1.0.0 --- 2026-09-16 --- the leading version field becomes $1$: \texorpdfstring{\ref{V4}}{(V4)}'s freezes, and the companion's half of the scope cut}

\noindent\emph{This version.} \textbf{Changes no normative statement of this document.}
It records the discharge of \ref{V5}\ref{V5ii} --- the criterion v0.19.0 added and did
not meet --- by the two acts that criterion names, both of which are edits to
\emph{other} files:
\begin{enumerate}[label=(\arabic*)]
\item \textbf{\ref{V4}'s two supersession freezes, carried out.}
  \file{hecke\_action.tex} is frozen at \emph{Version 1.0 --- 2026-07-02} and
  \file{monodromy\_v5.tex} at \emph{Version 5.0 --- 2026-07-03}, each date being that of
  the last commit that changed the file's content, each replacing a
  \code{\textbackslash date\{\textbackslash today\}} (entry~\ref{err:e17}), and each
  followed by a banner naming this document as the successor and Appendix~\ref{app:errata}
  as required reading. Neither document's content was otherwise touched. [v5]'s banner
  says where \cl{thm:explicit_complex_equiv} now lives --- its five C++ citing sites use
  the \code{[C]} key and resolve here (entry~\ref{err:e7}) --- and that \S12's $\beta$ is
  not re-homed (entry~\ref{err:e19}).
\item \textbf{\file{hecke\_extensions.tex} v1.3 cites the nine labels instead of defining
  them.} Its former \S1.1--\S1.2 and the boundary-eigenvalue claim of \S1.4 are replaced
  by a pointer section; the four of the nine this document \emph{corrected} are named
  there as superseded, with the correction and its site: \cl{cl:chi}'s Dickson-locus
  sentence (entry~\ref{err:e21}, \cl{rmk:chicoincide}), \cl{cl:boundary}'s
  $H^{i}(\partial)=0$ for $i\ge2$ without $2\in\cring^{\times}$ (\cl{rmk:bdryrepairs}),
  \cl{cl:les}'s hypothesis pair (\cl{rmk:sixscope}), \cl{cl:pd}'s self-adjointness
  wording (\cl{rmk:pdscope}). Its Appendix~A loses the seven STATUS rows for the moved
  claims, and its \code{[C]} key now means this document.
\end{enumerate}
The evidence that the criterion is met is \cl{rmk:labelcensus}'s second run: $0$
dangling \code{[C, \ldots]} citations, and an \emph{empty} intersection between this
document's label set and that of the only live companion. Commit: \texttt{b456448}.
Page count unchanged apart from the paragraphs below.

\paragraph{Normative.}
\begin{itemize}\setlength{\itemsep}{0pt}
\item \cl{V4} \PFIX{} --- from ``each receive a frozen version marker and a
  superseded-by banner when this document ships. Freezing them is a separate, later
  step'' to a record that the step is taken, with both frozen version numbers and what
  the banners say. The rule is unchanged; what changed is that it is now discharged, and
  \ref{V4}'s last sentence no longer defers.
\item \cl{V5} \PFIX{} --- \ref{V5ii} recorded as \textbf{met}. The rule itself is
  untouched; the sentence ``\ref{V5ii} is not [met]'' and the list of what blocked it are
  replaced by the discharge and a pointer to its evidence, and the paragraph now says in
  terms what a leading $1$ does \emph{not} assert.
\end{itemize}

\paragraph{Non-normative.}
\begin{itemize}\setlength{\itemsep}{0pt}
\item \cl{err:e17} --- ``Corrected in \ldots the frozen banners that both superseded
  documents receive when this document ships'' becomes a record of the two markers as
  written, with their dates and the reason those dates were chosen.
\item \cl{err:e18} --- gains a \emph{Completed at v1.0.0} paragraph: what the companion's
  v1.3 removed, which four of the nine it names as superseded, and the one surviving
  \code{[E, cl:les]} citation in the corpus --- a verbatim quotation inside the dated
  review record \file{docs/implementation/hecke\_extensions\_OQ2\_verdict.md}, left as
  written because rewriting a quotation would falsify the record.
\item \cl{rmk:labelcensus} --- gains the v1.0.0 re-run ($63$ labels, $778$ sites, $79$
  files, $0$ dangling) \emph{and} the two label-set checks the census does not perform.
  Those found one collision outside the nine, which the companion's v1.3 removed: it used
  \code{A1}--\code{A3} for its three required additions while this document uses them for
  authoring rules \ref{A1}--\ref{A3}. No citation used either reading, so no citation
  census could have found it --- which is now the stated reason for \S\ref{sec:machinery}'s
  guard to compare label sets and not only citations.
\item The abstract and the header comment are re-written for the same reason, and
  \S\ref{app:changelog-gaps}'s ``Not yet written'' and ``two supersession steps''
  paragraphs are updated. \textbf{No chapter of Parts~I--V was touched.}
\item \textbf{What this version does not claim.} Not that the three unimplemented
  chapters are implemented, not that \ref{open:ell2} is settled, not that
  \ref{obl:indep} is complete, and not that \S\ref{sec:machinery}'s census guard is
  written. \S\ref{app:changelog-gaps} keeps all four.
\end{itemize}

\subsection*{v0.19.0 --- 2026-09-16 --- chapters 21--22 and Appendices B--C: the worked example, the oracles, the machinery, the STATUS index and the formalization notes}

\noindent\emph{This version.} Writes \S\ref{sec:example}, \S\ref{sec:implementation}
with \S\ref{sec:machinery}, and Appendices \ref{app:status} and
\ref{app:formalization}, which were the document's last five declared stubs, adds
one errata entry, and \textbf{amends \ref{V5}}. \ref{V5} as written would have taken
the leading version field to $1$ on the strength of the last stub alone. It now requires
two criteria, and the second --- an unambiguous label namespace --- is \textbf{not} met:
the nine labels that entry~\ref{err:e18}'s scope cut moved into \S\ref{sec:six} are
still defined in the live \file{hecke\_extensions.tex} as well, so a citation
\code{[C,~cl:pd]} does not resolve to a unique statement and \ref{L6} cannot be honoured
across the corpus. The version is therefore \textbf{0.19.0} and the leading $0$ now
carries that information as well as the stub count. Criterion \ref{V5i} \emph{is} met,
and \ref{V5} says so. None of this asserts that the three unimplemented chapters have
been implemented (they have not) or that \ref{open:ell2} is settled (it is not);
\S\ref{app:changelog-gaps} says so in those words. Commit: \texttt{0a3358b}. Pages $142$--$147$, $148$--$151$,
$161$--$166$ and $167$--$170$, i.e.\ $6+4+6+4=20$~pp against an architecture budget of
$4+2+5+2.5=13.5$ (the overrun is $1.5\times$ overall, and is worst in
\S\ref{sec:implementation} at $4$~pp against $2$ --- its code map is a whole-document table, and the configuration surface and record
schema are new material the architecture budgeted as ``file$\to$section map''). Sources:
\file{hecke\_action.tex}:1512--1612 (\S\ref{sec:example}), :1613--1690
(\S\ref{sec:implementation}); this document itself, read mechanically, for
Appendix~\ref{app:status}; the eighteen per-chapter \emph{Formalization note} paragraphs
for Appendix~\ref{app:formalization}; and the shipped tree, read and run on 2026-09-16,
for every recorded number in all four.

\paragraph{Normative.}
\begin{itemize}\setlength{\itemsep}{0pt}
\item \cl{V5} \PFIX{} --- \textbf{amended}. Its criterion for the leading version field
  was ``the last stub is written''; it is now that criterion, \ref{V5i}, \emph{and}
  \ref{V5ii}, an unambiguous label namespace: \ref{V4}'s two freezes done, and no label
  of this document also defined by a live companion. The rule as it stood would have
  certified $1.X.Y$ with nine labels of \S\ref{sec:six} duplicated in
  \file{hecke\_extensions.tex} --- which \ref{L1}--\ref{L4} bind and \ref{L6} is meant
  to machine-check, and which entry~\ref{err:e18} recorded as a relocation rather than a
  duplication. The nine are named in \ref{V5} so that the criterion can be discharged by
  inspection. Owner decision of 2026-09-16.
\item \S\ref{sec:example}, \eqref{eq:t2datum}, \S\ref{sec:t2datum} --- the $T_2$ datum.
  \cl{prop:t2cocycle} \PRH{}, \eqref{eq:t2cocycle}, \eqref{eq:t2relcheck} --- the twelve
  cocycles and the three permutations, transported unchanged from the source table
  (which is correct: it was recomputed from $p$ alone in this pass and agrees entry for
  entry) but promoted from a display to a proposition with a proof, so that
  \S\ref{sec:chainlevel}'s forward reference has a numbered target.
\item \cl{ex:t2tau} \PNEW{}, \COMP{} --- the $\tau$-table at $p=2$ as \emph{emitted}:
  $69$ chains, $15$ triangle fills all over $\Z$, $B_{\mathrm{tab}}=2$, $B=16$,
  coefficients in $\{0,\pm1\}$. The superseded text gave the count as ``at most $69$''
  and derived it from a bijection that \cl{alg:tau}'s repair showed to be false; this
  example gives the count as a measurement and names step~0 as the reason the $69$ slots
  are not $69$ constructions. \emph{Not} added to the roster of statements nothing else
  states: its content is a recorded run, not new mathematics.
\item \S\ref{sec:workedinstance} \PNEW{} --- a fully worked $(n,\ell,p)=(2,5,2)$, with
  $\rhon$ of the six blocks of $\alpha$ and $\omega$ and of the three $\delta_j$ printed
  over $\Z$, and the emitted $\Amat{0}(T_2)=(3)$, $\Amat{1}(T_2)=(4)$. The superseded
  chapter worked the datum symbolically in $n$ and never instantiated it.
\item \S\ref{sec:oraclelist} --- the nine oracles, re-homed with the hypotheses written
  onto each clause (\ref{A1}) and the deciding of abort-versus-record handed to
  \cl{obl:oracle}(i),(ii) rather than restated. Two carry errata.
  \textbf{Oracle~6} carries \cl{err:e1}: it now states both identities --- full
  $T_4=T_2^2-2^{n+1}I$ and primitive $=T_2^2-2^{n+1}I-2^nI$ --- with the emitting entry
  point of each, which is what \cl{obl:oracle}(iii) binds it to. \textbf{Oracle~7}
  carries \cl{err:e23}, new in this version.
\item \cl{err:e23} \PNEW{} --- \emph{found in this pass}, in the oracle list being
  re-homed: the superseded oracle~7 asserts $U_4=U_2^{2}$ at $p=\ell=2$ citing
  \cl{cor:ulrecursion}, whose clause~(1) carries $\ell\ge5$ and whose clause~(2) says
  something else. The implementation already records rather than asserts it at
  $\ell\in\{2,3\}$ (amendment A-F1), so the defect is in the contract text alone.
  Oracle~7 now splits by hypothesis. \cl{cor:ulrecursion} itself is unchanged.
\item \S\ref{sec:oraclesnew} \PNEW{} --- the boundary, cone and involution oracle
  families, as \emph{specifications}: each names the obligation it serves, and each is
  marked as not runnable, because the objects do not exist. Two numerical corrections
  were made while writing them and are the reason they are numbers and not prose: the
  boundary $H^{0}$ has rank $\lfloor n/\ell\rfloor+1$ and not $1$
  (\cl{prop:bdrychar}(1)), so the cone dimension book is instantiated at
  $(n,\ell,p)=(2,5,2)$, where $n<\ell$ and every term is independently known, and not at
  $n=10$, $\ell=5$, where it is not.
\item \cl{rmk:dicksontwo} \PNEW{}, \COMP{} --- two implementations of the gating
  predicate of \cl{rmk:dickson} exist in the tree under one name; the weaker one
  (first Dickson generator only) is sound on the grid it is used on, checked
  exhaustively, and unsound in general.
\item \S\ref{sec:implementation}, \S\ref{sec:codemap}, \S\ref{sec:notimplemented},
  \S\ref{sec:buildsurface}, \S\ref{sec:recordschema} \PNEW{} --- the whole-document code
  map; the boxed rule that \textbf{no switch may change a computed value}, with the two
  rollbacks, the two shadow builds and the execution switches classified by which of
  those they are; and the record schema, whose fifth item is site~(6) of
  \S\ref{sec:oc1sites} and carries \ref{open:ell2}'s $q_{\max}$, \code{truncated} and
  reported-degree list. \S\ref{sec:notimplemented} is a contract box and states in one
  place what four chapters had each stated in their own code maps.
\item \S\ref{sec:machinery} \PNEW{} --- this document's four guards, of which three
  exist. \cl{rmk:labelcensus} \PNEW{}, \COMP{}: the fourth, the citation census that
  would have caught \cl{err:e7}, is \emph{not written}, and was run by hand instead ---
  $64$ distinct labels cited from $78$ files, $0$ cited labels this document does not
  define.
\item Appendix~\ref{app:status} \PNEW{} --- $199$ statements indexed with status and
  provenance, mechanically derived and build-gated: the number column is a real
  \code{\textbackslash ref}, so an index line naming an undefined label fails the build,
  which is what \ref{L6} requires and which was verified by deliberately mistyping one
  label. \COMP{} is split into \emph{registered} ($8$) and \emph{ad hoc} ($5$), because
  a check the repository re-runs and a sweep done once are not the same evidence.
\item Appendix~\ref{app:formalization} \PNEW{} --- the eighteen per-chapter
  formalization notes reconciled into: what must be built, what must be imported, the
  five hypotheses that must not be lost, the shapes to factor once, and a cheapest-target
  table read against Appendix~\ref{app:changelog}'s roster of the statements nothing
  else states. It performs \emph{no} cost assessment --- that is the commission --- and
  it boxes the warning that every claim in this document about what \code{mathlib4}
  contains is unverified.
\end{itemize}

\paragraph{Non-normative.}
\begin{itemize}\setlength{\itemsep}{0pt}
\item \textbf{Two gaps in this document were found by compiling
  Appendix~\ref{app:status} and are recorded, not repaired}, because repairing another
  chapter's statement from an index is the silent edit \ref{V3} exists to prevent, and
  because \ref{CL2} reserves Appendix~\ref{app:errata} for the superseded corpus.
  \textbf{(i)} \cl{cor:ulrecursion} carries no status tag although
  \S\ref{sec:status} requires one; it is the only untagged statement in the document
  that is not a definition, a hypothesis, a procedure, a predicate or commentary, and
  the tag it should carry is \PROVED. \textbf{(ii)} \cl{rmk:taurings} has no line in
  this appendix: it entered in v0.5.0 alongside \cl{def:tau}'s widening, which names it,
  but got no line of its own and is in no roster. Its line in
  Appendix~\ref{app:status} is starred for that reason.
\item The provenance counts above are updated for this version: \PNEW{} $119$
  ($69$ assertions or definitions, $50$ remarks), \PRH{} $45$, \PFIX{} $38$.
  Appendix~\ref{app:status}'s column totals are \emph{statement} counts at current
  content and are deliberately different; the contract box at the head of that appendix
  says why neither refutes the other.
\item The roster of statements that no superseded document and no project note states is
  \textbf{unchanged} by this version. Of its four new statements one is \PRH{}, two are
  remarks, and \cl{ex:t2tau} is a recorded run.
\end{itemize}

\subsection*{v0.18.0 --- 2026-09-16 --- chapters 18--20: reduction, bounds, the floating-point filter, and the obligations}

\noindent\emph{This version.} Writes \S\ref{sec:reduction}, \S\ref{sec:fpfilter} and
\S\ref{sec:obligations}, which between them were eighteen declared stubs, and adds one
errata entry. All three are re-homings of \file{hecke\_action.tex}:1162--1511, and
almost all of the new writing is in the two places where re-homing was not enough:
\ref{A3} (a definition is the terminating procedure where the code's is), which turned
\cl{alg:tau} from a four-step sketch into a five-step algorithm whose emitted table is
well defined, and \ref{A4} (an obligation is a checkable predicate on emitted data),
which turned eleven prose restatements of theorems into clause lists naming data,
predicate, ring and failure action. Commit: \texttt{3d67bf5}. Pages $123$--$128$,
$129$--$131$ and $132$--$140$ (architecture budget $4+2+3$~pp; the overrun is
$2.0\times$ and is concentrated in \S\ref{sec:obligations} at $9$~pp against $3$ --- the
roster of \S\ref{sec:obligationroster} has fifteen rows, not eleven, because the
obligations of \S\S\ref{sec:involution}, \ref{sec:boundarydatum} and \ref{sec:cone} are
integrated here, and \ref{A4}'s four fields are written out for every clause of every
one of them). Sources: \file{hecke\_action.tex}:1162--1338 (\S\ref{sec:reduction}),
:1339--1407 (\S\ref{sec:fpfilter}), :1408--1511 (\S\ref{sec:obligations}); and, for the
constructive content of \cl{alg:tau} and for every status claim in
\S\ref{sec:obligationroster}, the shipped tree read on 2026-09-16.

\paragraph{Normative.}
\begin{itemize}\setlength{\itemsep}{0pt}
\item \cl{def:primitives} \PRH{} --- (P1)--(P4) with their output specifications.
  Two changes of content, both recorded in the statement: (P2)'s parenthetical
  ``all intermediate matrices have entries bounded by the entries of the input'' is
  dropped, because it contradicts \cl{lem:bounds}(4) of the same source, which says
  \emph{twice} the maximum and is what the implementation checks --- the true statement
  about the input's own entries is that $\abs{c}$ strictly decreases, and that is what
  is written; and (P3) is recorded as \textbf{specified with no live call site}, the
  closed forms of \cl{def:cosetreps} producing $\delta_j$ directly. (P4) now specifies
  the canonical key (lexicographic minimum over the transporter set, sign from the
  vertex-slot permutation, quotient by $-I$) where the source said ``a fixed choice of
  orbit section''.
\item \cl{rmk:primitiveclosure} \PNEW{} --- the closure statement: $\Z$ is touched in
  exactly four places, which is what makes \cl{thm:backend} a statement about a finite
  workload.
\item \cl{alg:tau} \PFIX{}, \eqref{eq:taurecord} --- the $\tau$-table construction.
  \textbf{The source's step list is not enough to determine its own output.} It stores
  ``one representative per orbit, since $\{(i,\sigma)\}\leftrightarrow$ simplices of
  $\Ga''\backslash\Ku$'' --- and that bijection is false: $\orient{q}(K)$ is not a
  section of the orbits, already because $\lift{v_3}=T\cdot\lift{v_0}$. The repaired
  algorithm has a \textbf{step~0}: one slot per $\Ga''$-orbit is constructed and every
  other is derived through a single shared (T1) helper used by \emph{both} construction
  and certification. Step~1 additionally carries the two branches that later chapters
  depend on --- the exact-position hit, and the ascending-index tie-break that
  \cl{rmk:iotatausym} shows is structurally incompatible with $\invK$ --- and the
  $\adjg(g)$ projective device that keeps the geometry rational-free. Step~3 carries
  owner decision OP2's default-forbidden mod-$\ell$ fallback. Step~4 is delegated to
  \cl{obl:tau} rather than restated (\ref{A4}: one home).
\item \cl{rmk:certoverGeometry} \PNEW{} --- certification over geometry, with the two
  consequences the rest of the document uses: the choices of steps~1--3 are free and
  changing one changes the table but not $\Amat{q}$; elliptic slots are forced by (T1),
  which explains the construction and proves nothing.
\item \cl{lem:bounds} \PFIX{} --- the entry bounds. Clause (4) was asserted
  (``never exceed twice the maximum of their input entries'') and is now
  \textbf{proved}, in the sharper form $\max(\lVert\gamma\rVert,2)$, from
  $\abs{a'}\le\abs{c}/2$ and $b'=(a'd-1)/c$; the implementation checks the weaker
  constant at every step. Clauses (2) and (3) are stated in the level $N=p^{\nu}$, not
  in $p$: the source wrote $\lVert s_j\rVert\le p$ and $4p^2$, which are the $\nu=1$
  values of bounds that are really about $\mathbb{P}^1(\Z/N)$, and the implementation's
  certified $B$ is $\max(B_{\mathrm{tab}},4N^2)$.
\item \cl{rmk:btabsmall} \PNEW{} --- \COMP{}: $C_K$ is never evaluated and
  $B\le196$ over the $66$ shipped points, so the first tier of \cl{thm:backend} is never
  approached. (Re-reads the run already recorded at \cl{rmk:tauintegral}; nothing new is
  executed.)
\item \cl{thm:backend} \PFIX{} --- the backend. The source's clause (iii) bounds the
  predicates by $2B^2$ with $B$ the bound on the \emph{stored matrices}, which does not
  cover the one predicate whose inputs are not stored matrices: the nearest-vertex
  comparison of \cl{alg:tau}(1b) forms a difference of squared distances on cleared
  denominators, needs the $128$-bit tier and a $\gcd$ reduction before the $64$-bit
  predicate can take it, and aborts if that is not enough. That clause is written out.
  The early-reduction scheme and the mod-$\ell$ box are transported.
\item \cl{lem:predshapes} \PNEW{}, \eqref{eq:dot2} --- every branch decision is a
  two-term integer dot product, and there are exactly four: (D), (Q1), (Q2) and
  \textbf{(Q3)}, the cross-$D$ nearest-vertex comparison $\sgn(A+B\sqrt3)$, which the
  source did not list at all. Proved, including the closure of the exact quadratic
  points under the M\"obius action of a positive-determinant matrix. This is the lemma
  that makes the rest of \S\ref{sec:fpfilter} a self-contained statement about
  IEEE-754 rather than a claim about the pipeline.
\item \cl{thm:fp64} \PFIX{}, \eqref{eq:fperr} --- the filter. Its floating-point
  hypotheses were unstated and its proof used all three: (FP1) exact input conversion,
  (FP2) the round-to-nearest standard model, and \textbf{(FP3) the evaluation order
  $\fl(\fl(a_1b_1)\pm\fl(a_2b_2))$ with no fused multiply-add}, which the implementation
  enforces with \code{volatile} intermediates and which \ref{A1} requires in the
  statement. The source's displayed error chain does not parse as written; the bound is
  re-derived, and the sharper $\tfrac12E(B)(1+2^{-53})$ is recorded because
  \cl{rmk:thresholdrounding} needs it.
\item \cl{rmk:thresholdrounding} \PNEW{} --- the threshold is itself computed in
  floating point, and the factor $2$ of slack between the proved bound and the normative
  constant $E(B)=2^{-50}B^2$ is what makes that sound. This is why the binding constant
  is not the sharper $2^{-51}B^2$.
\item \cl{cor:fptotal} \PFIX{}, \eqref{eq:fptotalcrit} --- totality. The criterion is
  now the non-strict $E(B)\le\tfrac12$, with a second and independent sufficient
  condition (exact representability, $2B^2\le2^{53}$), and ``total \emph{iff}'' is
  corrected to a sufficient condition. Erratum \cl{err:e22}.
\item \cl{cor:fp32} \PFIX{} --- the binary32 analogue. Threshold $B\le2^{10}$
  \emph{unchanged}, and now proved: it fails the source's own strict criterion with
  equality. Erratum \cl{err:e22}.
\item \cl{obl:welldef} \PRH{} --- recast as the three predicates that
  \cl{alg:assembly}'s contract box delegates here: $y_j$ satisfies every row of $\Rel$;
  $d\,y_j=0$; the $\Th^{(q)}_{\mathrm{eq}}$ solve succeeds with identically zero
  residual. All over $\cring$, all aborting.
\item \cl{obl:indep} \PFIX{} --- independence. \textbf{The source asserted as a runtime
  test something the construction cannot test.} Clause (iii), independence of the
  double-coset representative $g$, is not reachable through an interface that fixes
  $g=\mathrm{diag}(1,N)$ and $\Ga''=\Ga_0(N)$; what a run can do is the constant-$h_j$
  case of clause (i), and the obligation now says so and says what would make (iii)
  checkable. Clause (i) also records, rather than asserts, the cochain-level agreement
  under a non-trivial relabelling, because rebuilding the chain rebuilds $\tau$
  (\cl{rmk:indeplevel}); clause (ii) is gated on $6\in\cring^{\times}$, which is where
  \cl{lem:tauexists}(2) is used.
\item \cl{obl:integral} \PRH{} --- four clauses: integrality of every $\delta$, no
  inversion at determinant $N$, the integral normalisation $\Amat{0}(T_p)=1+p$ with
  coefficient $p^{\,n+1}$, and the bad prime through the same pipeline.
\item \cl{obl:algebra} \PFIX{} --- \textbf{erratum E1, the second of its two landing
  sites} (\cl{err:e1}). The source's runtime test compared the \emph{primitive}
  $p^2+p$-coset operator with $T_p^2-p^{\,n+1}$, which is the \emph{full} operator's
  identity, and is therefore wrong by $p^{\,n}I$. The obligation now asserts \emph{both}
  identities, names which entry point emits which operator, and requires the full
  $T_{p^2}$ of clause (ii) to be obtained by the double-coset assembly
  \eqref{eq:fullassembly} rather than by the recursion being tested. Clauses (iv) and
  (v), the diamond and coprime multiplicativity, are new as clauses; (iv) carries no
  hypothesis on $\ell$, by \cl{prop:diamond}.
\item \cl{obl:agreement} \PRH{} --- four clauses, each with its level of comparison
  fixed by \cl{rmk:indeplevel}: cochain-level bit-identity only between constructions on
  the same datum, cohomology-level for the bar-cochain route of \cl{prop:barformula}.
  The clause requiring agreement through the public evaluation interface is written out,
  because the internal assertion can be discharged against private state.
\item \cl{obl:shapiro} \PRH{} --- the dimension comparison, with the observation that
  the \emph{equality of the two computed dimensions} carries no hypothesis on $\ell$
  while identifying either side with $H^{\bullet}(\Ga_0(N);\VR)$ does.
\item \cl{obl:oracle} \PFIX{}, \eqref{eq:esoracle} --- the Eichler--Shimura oracle. The
  source wrote the cuspidal factor as $(X-a_p(f))(X-a_p(f))$, a doubled factor whose
  multiplicity is the point; it is now $(X-a_p(f))^2$ with the multiplicity attributed
  to the two summands of \cl{cl:dimH1}(1). What is \emph{asserted} and what is
  \emph{recorded} is now part of the obligation --- only independently verified $a_p(f)$
  constants gate --- and clause (iii) binds \S\ref{sec:example}'s oracle~6 to the
  full-versus-primitive distinction of \cl{err:e1}.
\item \cl{rmk:dickson} \PRH{} --- the exceptional loci. The invariant-theoretic input
  and its citation are \emph{deferred} to \cl{cl:dimH1}(2) rather than restated with a
  second pointer; what is added is that both semigroup generators matter, with the
  witness $(n,\ell)=(20,5)$ where $\ell+1\nmid n$, and \COMP{} that the implemented
  predicate is semigroup membership and that $(20,5)$ is in the shipped grid.
\item \cl{obl:fp} \PRH{} --- four clauses. The differential test must not share its
  exact routine with the production fallback; the counters witnessing dead code are a
  diagnostic (\ref{A6}); and clause (iv), the pure-integer bit-identity, is recorded as
  a \textbf{two-build protocol rather than a predicate a single run can evaluate}.
\item \cl{obl:backend} \PRH{} --- five clauses: the certified bound recounted, the tier
  inequality asserted, overflow aborting rather than wrapping, the absence of arbitrary
  precision recorded rather than asserted at runtime, and the mod-$\ell$ backend exact
  for the $\ell$ in force.
\item \cl{obl:tau} \PRH{} --- seven clauses, each named at the point
  \cl{alg:tau}(4) defers to them, with the ring per clause ($\Z$, or $\Fl$ for an
  identity involving a mod-$\ell$-solved chain) and the (T2) right-hand side required to
  go through the \emph{same} (T1) helper the construction uses.
\item \S\ref{sec:obligationroster} --- the roster. Fifteen obligations with, for each,
  the ring, the failure action, the status and the registered validations behind it;
  ten are discharged by a registered validation, one (\cl{obl:indep}) is partial for a
  stated structural reason, one clause of \cl{obl:iota} is discharged
  \emph{negatively}, and three (\cl{obl:taubdry}, \cl{obl:eta}, \cl{obl:coneexact}) are
  unimplemented. \COMP{} The one clause that is a protocol rather than a predicate,
  \cl{obl:fp}(iv), was executed on 2026-09-16: two out-of-source builds,
  \code{HECKE\_PURE\_INTEGER} off and on, both reporting tier~A $9446$ failing $0$ and
  tier~B $4603$ mismatching $0$ against the same frozen reference under
  \code{GOLDEN\_STRICT=1}, hence bit-identical on all $14049$ hashed objects, with five
  named Hecke validations exiting $0$ in both
  (\file{docs/matrix-engine/tools/cross\_config.sh}). That is a recorded run of an
  existing script, not a registered validation, and the roster says so.
\end{itemize}

\paragraph{Non-normative.}
\begin{itemize}\setlength{\itemsep}{0pt}
\item \cl{err:e22} --- a new errata entry: \code{cor:fp32}'s totality threshold
  $B\le2^{10}$ does not follow from \code{cor:fptotal}'s strict criterion
  $E(B)<\tfrac12$, which it meets with equality. The threshold and every shipped
  constant are unchanged; what was missing is the proof, which the implementation's own
  validation had already written down in a comment and the contract had not. Adding an
  erratum entry is itself a change to this document (\ref{CL2}), and the summary table
  of \cl{app:errata} gains its row.
\item \textbf{One symbol was deliberately not introduced.} \S\ref{sec:fpfilter} needs the
  unit roundoff, whose customary letter is $\varepsilon$ --- which is
  \cl{app:labelindex}'s $\aug$, the augmentation, four pages earlier in \cl{alg:tau}(4).
  Rather than add a twentieth collision row, the two values are written as the explicit
  constants $2^{-53}$ and $2^{-24}$ throughout and the quantity is given no letter; the
  chapter header records the decision. \file{gl2\_notation.tex} is unchanged by this
  version. The four symbols \S\ref{sec:reduction} does need
  ($\lVert\cdot\rVert$, $\mathcal{F}$, $\mathrm{pos}$, $\mathcal{N}$) are declared
  chapter-locally, in the manner of the proof of \cl{prop:fullassembly}, and occur
  nowhere else.
\item \textbf{No new statement of these three chapters carries a \CONJ{} or an
  \OPENSTAT{} tag}, and no \STD{} pointer is newly asserted: \S\ref{sec:reduction} and
  \S\ref{sec:fpfilter} are self-contained, and \S\ref{sec:obligations} cites
  \cite{shimura}, \cite{ash_stevens}, \cite{brown}, \cite{harder3} and \cite{deng} only
  at pointers \cl{thm:relations}, \cl{thm:classical} and \cl{cl:dimH1} already carry, and
  \cl{obl:oracle} says so in its own statement rather than implying the texts were
  re-checked here.
  \file{gl2\_refs.bib} is unchanged by this version.
\item \cl{app:changelog} --- provenance counts updated from $111/35/29$ to $116/44/38$
  (\PNEW{}/\PRH{}/\PFIX{}) and the \PNEW{} split from $67/44$ to $68/48$; this version
  contributes $5$ \PNEW{}, $9$ \PRH{} and $9$ \PFIX{} statements. The tagged-line
  totals go from $182/64$ to $190/80$, this version contributing $8$ labels on \PNEW{}
  lines and $16$ on \PRH{} lines, counted as every \code{\textbackslash cl} on an item
  whose \emph{subject} carries that tag. A third such figure is \emph{not} introduced:
  this version's \PFIX{} items name $21$ labels, but earlier versions did not record
  theirs, and a column that starts at v0.18.0 would read as a total.
\item \cl{app:changelog-gaps} --- \S\ref{sec:reduction}, \S\ref{sec:fpfilter} and
  \S\ref{sec:obligations} removed from the stub roster, with their eighteen stub labels.
  \cl{app:labelindex} --- the nineteen new labels of the three chapters dereferenced,
  and \cl{err:e22} added to the errata list, per \ref{L6}.
\item The abstract, the file header and the date line: twenty of the twenty-two
  chapters are written.
\end{itemize}

\subsection*{v0.17.0 --- 2026-09-16 --- chapters 16--17: the six groups, the long exact sequence, and the general orbifold model}

\noindent\emph{This version.} Writes \S\ref{sec:six} and \S\ref{sec:orbifoldmodel},
which between them were eleven declared stubs. \S\ref{sec:six} is the
\textbf{move out of [E]} that the scope cut requires (\cl{err:e18}): nine labels that
were defined in \file{hecke\_extensions.tex} are defined here, each now carrying its own
hypotheses (\ref{A1}) and its own citation (\ref{A2}) instead of a blanket
``$\ell\ge5$, $n$ even'' declared once. \S\ref{sec:orbifoldmodel} is
\textbf{definitions now, construction deferred} (owner decision OP4), written
refinement-parametrically per owner direction OD-T1. Commit: \texttt{44c355e}. Pages
$111$--$118$ and $119$--$122$ (architecture budget $4+3$~pp; the overrun is $1.7\times$,
and it is the apparatus \ref{A1} and \ref{A2} force --- [E] proved none of these nine
statements and cited three of them to a working note). Sources:
\file{hecke\_extensions.tex}:186--338 for \S\ref{sec:six}; the OQ2 verdict \S5, the
implementation plan's C-A3(ii.7) and the Moerdijk--Pronk source in
\file{docs/literature/orbifolds/} for \S\ref{sec:orbifoldmodel}; owner decisions OP4,
OP5 and direction OD-T1. \textbf{One erratum lands here}, \cl{err:e21}, found in this
pass.

\paragraph{The three things a reader should check first.}
\textbf{(i)} \cl{cl:pd} is stated in the \emph{adjointness} form and in no other. [E]'s
``$T_q$ self-adjoint once $\Dn$ carries the adjugate $\De$-structure'' is true of
$\langle\rhon(g)v,\varphi\rangle=\langle v,\rhodual(\adjg g)\varphi\rangle$
(\cl{prop:pairing}(3), twist-free) and \emph{false} of
$\langle\rhon(g)v,\rhodual(g)\varphi\rangle=(\det g)^{n}\langle v,\varphi\rangle$
(\cl{prop:pairing}(4), which carries $p^{\,n}$). The v0.3.0 entry below recorded that it
``must be re-homed that way''; it now is. \textbf{(ii)} \cl{cl:les} drops the hypothesis
$\Hcpt{0}=0$, which \cl{lem:conerelative}(4) makes unconditional, and drops ``$n$ even'',
which exactness never used; it \emph{keeps} $H^{2}(\Ga;\VR)=0$, and
\cl{lem:h2torsion}(2) shows why --- at $(\ell,n)=(2,0)$, $H^{q}(\Ga;\F_2)\neq0$ for
every $q\ge0$. \textbf{(iii)} \cl{err:e21}: one sentence of [E]'s \cl{cl:chi} is false.

\paragraph{Normative.}
\begin{itemize}\setlength{\itemsep}{0pt}
\item \S\ref{sec:six}, \cl{eq:sixgroups} --- the chapter and the contract box fixing
  $X_1,\dots,X_6$, $\Htot$ and $d$ in the notation of
  \S\ref{sec:notation-cone}, with the three recurring hypotheses named once and repeated
  in every statement.
\item \cl{cl:amalgam} \PNEW{} --- \textbf{strictly generalised} (\ref{CL4}: the tag is on
  the content). [E] asserted the Mayer--Vietoris sequence and the vanishing
  $H^{i}(\Ga;M)=0$, $i\ge2$, for $\ell\ge5$, as \STD{}. Here: (1) the amalgam and its
  Bass--Serre tree sequence \cl{eq:treeses}, the only external input; (2) the long exact
  sequence \cl{eq:mvamalgam} in \emph{every} degree over \emph{every} commutative ring,
  proved from (1) by Shapiro; (3) the $6\in\cring^{\times}$ consequences, with
  $\VR^{S}=(1+\rhon(S))\VR$ and $\VR^{U}=(1+\rhon(U)+\rhon(U)^{2})\VR$ derived rather
  than asserted; (4) a torsion bound of $24$ proved from (2) with \emph{no} citation at
  all, beside the sharp classical $12$. The Euler-characteristic clause no longer names
  the quantity $\chi$, which is the boundary character.
\item \cl{lem:h2torsion} \PNEW{} --- (1) the vanishing for $6\in\cring^{\times}$;
  \textbf{(2) the proof that it is not removable}: $H^{q}(\SL_2(\Z);\F_2)\neq0$ for every
  $q\ge0$, by the dimension count in \cl{eq:mvamalgam} with $\F_2$-coefficients. This
  also settles, for $n=0$, the half of \cl{rmk:ell2window}(2) that was left open there;
  that remark is \textbf{not} edited (it is the v0.14.0 slot's record), and the new
  statement says so at its own site.
\item \cl{cl:dimH1} \PRH{} --- content unchanged; restated in the closed form
  $2d_k-1+\dim H^{0}$ with the \emph{same} $d_k$ that \cl{obl:iota}(v) defines, rather
  than in the symbols $M_k$, $S_k$, which would collide with $\Mn=\Sym^{n}\Fl^{2}$ and
  with the operator set of [E] \S7 (\S\ref{sec:notation-collisions}, C12, C14). No
  modular-form notation was added: see the non-normative note below.
\item \cl{cl:bockstein} \PRH{} --- content unchanged; its \emph{status} moves from
  \STD{} to \PROVED{}, the $2,3$-torsion input now being \cl{cl:amalgam}(4) and the rest
  the long exact cohomology sequence of
  $0\to\Vcoef\xrightarrow{\ell}\Vcoef\to\Mn\to0$, whose connecting map \emph{is} the
  Bockstein. Clause~(3), the Hecke-equivariance, is given its reason: that sequence is one
  of $\De$-modules, and $\res$, $(\conj_g)_{\ast}$, $\cor$ are morphisms of
  $\delta$-functors. The vanishing $H^{0}(\Ga;\Vcoef)=0$ is proved in two lines rather
  than cited.
\item \cl{cl:boundary} \PNEW{} --- \textbf{hypothesis repaired}. [E]'s
  ``$H^{i}(\partial;M)=0$ for $i\ge2$'' is false at $\ell=2$
  (\cl{lem:ell2unbounded}); the statement now carries $2\in\cring^{\times}$, and is
  proved from \cl{prop:circlemodel} and \cl{prop:bdrychar}(1) rather than asserted.
  \cl{rmk:bdryrepairs} \PNEW{} records that repair and the C18 spelling $\mathrm{Nm}$.
\item \cl{cl:les}, \cl{eq:les}, \cl{eq:dimbook}, \cl{eq:r1formula} --- \PNEW{} for the
  generalisation. Four clauses: (1) the sequence in all degrees with \emph{no} hypothesis
  on $\ell$ or $n$, Hecke-equivariant by \cl{thm:conewelldef}(a) --- where [E] made
  equivariance conditional on ``$\Phi_g$ a cochain map'', which is now a theorem;
  (2) the collapse, from \cl{lem:conerelative}(4),(5); (3) exactly when \cl{eq:les} is a
  statement about \cl{eq:sixgroups}; (4) the formula for $\rres^{1}$, proved
  \emph{well defined} here and tagged \CONJ{} as a citation for the identification,
  its only source being the working note \cl{err:e14} declares not a source of record.
\item \cl{cl:chi} --- \PNEW{} for clause~(3), which is now \PROVED{} (Vandermonde on a
  generator of $\Fl^{\times}$, plus Dirichlet only for the passage from primes to
  residues). Clause~(1) cites \cl{prop:bdrychar}; clause~(2), the general-$n$ $H^{1}$
  multiset, stays \CONJ{} --- it is what \cl{obl:taubdry}(vi) records rather than gates
  --- with a \COMP{} corroboration on $468$ cases.
\item \cl{cl:pd} \PFIX{} --- restated in the adjointness form (see (i) above), with
  $\adjg(\Ga g\Ga)=\Ga g\Ga$ from \cl{prop:pairing}(5) supplying the
  \emph{self}-adjointness and no determinant factor; \cl{err:e14}'s socle-versus-Loewy
  correction carried; status \OPENSTAT{}, as \cl{prop:coneroutes}(3) and
  \cl{thm:conewelldef}(c$'$) already read it. \cl{rmk:pdscope} \PNEW{} states the
  $p^{\,n}$ twist, what the claim depends on, and what has no certificate yet.
\item \cl{rmk:sixscope} \PNEW{} --- which two of [E]'s hypotheses are dropped and which
  one is not.
\item \S\ref{sec:orbifoldmodel}, \cl{def:borelmodel}, \cl{eq:borelfree},
  \ref{G1}--\ref{G4} \PNEW{} --- the free $\Ga$-model $F_{\bullet}\otimes C_{\bullet}(Y)$
  of an \emph{arbitrary} pair satisfying \ref{G1}--\ref{G4}, its orbifold cochain
  complexes in a degree range $\le q_{\max}+1$, and twisted chain data on it, with
  (T1)--(T3) and the pair identity carried over verbatim. Stated so that a
  $\Ga$-equivariant refinement of \cl{def:model7} is admissible input, per OD-T1.
\item \cl{rmk:borelfree} \PNEW{} --- why the free model is wanted: both clauses of
  \cl{lem:tauexists} hold on it over $\Z$ for every $\ell$, so
  \cl{thm:conewelldef}(c) holds at every $\ell$ there, where on the $7$-vertex model it
  is asserted only for $6\in\cring^{\times}$ (\cl{thm:conewelldef}(e)).
\item \cl{thm:specialisation} \PNEW{} --- ``the $7$-vertex complex is \emph{a}
  specialisation at $6\in\cring^{\times}$'', in a degree range and for an arbitrary
  input: (1) the isotropy spectral sequence collapses to $\Ceq(K_Y,\Rel_Y;L)$;
  (2) the relative version is $\Hcpt{\bullet}$; (3) the two Hecke actions agree, both
  being $\cor\circ(\conj_g)_{\ast}\circ\res$; (4) the shipped model is one admissible
  input, a refinement another.
\item \cl{rmk:orbifoldrefine} \PNEW{} --- what a refinement preserves
  (\cl{prop:model7h4}(i)--(iv), \cl{def:cosetreps}), what must be recomputed
  ($N_0,N_1,N_2$, \cl{def:model7rel}, the stored simplex order, \cl{alg:tau},
  \cl{lem:bounds}, the golden references), that OD-T1 revises nothing, and that a
  refinement does not by itself make the complexes finite.
\item \cl{err:e21} --- a new errata entry: [E]'s ``$\chi_1\equiv\chi_2$ at every Dickson
  locus'' is false; the criterion is $(\ell-1)\mid(n-2)$, first counterexample
  $(n,\ell)=(12,5)$. Adding an erratum entry is itself a change to this document
  (\ref{CL2}), and the summary table of \cl{app:errata} gains its row.
\end{itemize}

\paragraph{Non-normative.}
\begin{itemize}\setlength{\itemsep}{0pt}
\item \cl{rmk:mpcaveat} \PNEW{} --- Moerdijk--Pronk, cited for the Borel-construction
  identification and for $S(\mathcal{T})$ as one computational model, and \emph{not} for
  the normative definition: their charts are effective and
  $[\SL_2(\Z)\backslash\Hyp]$ is not, and they twice omit the independence proof for
  $\mu$. \textbf{Every quotation and line pointer in that remark was checked against
  \file{docs/literature/orbifolds/moerdijk\_pronk\_\ldots/simplicial.tex} on
  2026-09-16}; the published theorem \emph{number} was not checked against the journal
  text and the remark says so.
\item \file{gl2\_refs.bib} --- two entries added, \code{serre\_trees} (the amalgam, its
  tree, and free subgroups of free products) and \code{wilkerson} (the Dickson
  invariants). \textbf{Neither is in \file{docs/literature/} and neither section pointer
  is verified}; both notes say so, and both statements carry either an independent proof
  from \cl{eq:treeses} or a \COMP{} corroboration. \code{moerdijk\_pronk}'s existing note
  already said what \cl{rmk:mpcaveat} now spells out.
\item \textbf{No modular-form notation was added}, and \file{gl2\_notation.tex} is
  unchanged by this version. Slot~7 recorded that the contract has no symbol for $M_k$ or
  $S_k$ and that chapter~\ref{sec:six} might need them. It does not: the only statement
  that would use them, \cl{cl:dimH1}, states the same numbers as $d_k$ and $d_k-1$, which
  \cl{obl:iota}(v) already defines with its citation. Adding $\mathcal{M}_k$ and
  $\mathcal{S}_k$ would have cost a twentieth collision row --- $\mathcal{S}$ is [E]'s
  operator set (C14) and $M$ already has four live readings (C12) --- for no statement
  this document makes.
\item \COMP{} \textbf{corroborations, both \emph{ad hoc} and neither a registered
  validation.} \cl{cl:dimH1}: the shipped campaign table
  \file{docs/matrix-engine/results/results\_table.tsv} re-read on 2026-09-16 --- $2118$
  rows, $710$ distinct $(n,\ell)$ with $\ell\ge5$, $n>0$, $203$ on the Dickson locus;
  $0$ disagreements on $\dim\Heq{0}$, on $\dim\Heq{1}=2d_k-1+\dim\Heq{0}$, and on
  $\dim\Heq{2}=0$. \cl{cl:chi}(2),(3) and \cl{rmk:chicoincide}: direct computation over
  $\Fl$ --- $468$ operator spectra and $190\,560$ instances of the coincidence criterion,
  $0$ disagreements. There is nothing in the tree to re-run for either; \cl{obl:oracle}
  and \cl{obl:taubdry}(vi) are where they would become registered.
\item \cl{app:changelog} --- provenance counts updated from $97/33/28$ to $111/35/29$
  (\PNEW{}/\PRH{}/\PFIX{}), the \PNEW{} split from $60/37$ to $67/44$, the tagged-line
  totals from $144/60$ to $182/64$ --- this version contributes $38$ labels on \PNEW{}
  lines and $4$ on \PRH{} lines, counted as every \code{\textbackslash cl} on an item
  whose \emph{subject} carries that tag, which is the convention the entries above use
  and which is stated here because it is not self-evident from them --- and a Part~IV
  roster bullet of $7$ new statements added.
\item \cl{sec:oc1sites} --- site~(4) updated: \cl{def:borelmodel} and
  \cl{thm:specialisation} now \emph{are} stated in a degree range rather than being
  scheduled to be, and \cl{rmk:orbifoldrefine}(4) is added as that chapter's third
  \ref{open:ell2}-dependent site. This is the only edit this version makes to
  \S\ref{sec:boundarydatum}, and it is the one \cl{sec:oc1sites} asks for by existing.
\item \cl{app:changelog-gaps} --- \S\ref{sec:six}
  and \S\ref{sec:orbifoldmodel} removed from the stub roster. \cl{app:labelindex} ---
  the seventeen new labels of the two chapters dereferenced, and \cl{err:e21} added to the
  errata list, per \ref{L6}.
\item The abstract, the file header and the date line: seventeen of the twenty-two
  chapters are written.
\end{itemize}


\subsection*{v0.16.0 --- 2026-09-16 --- \texorpdfstring{\ref{lem:tauexists}}{lem:tauexists} apportioned; notation C18, C19; a wrong bibliography entry}

\noindent\emph{This version.} Three things, none of them a new chapter. (i) The
hypothesis $6\in\cring^{\times}$ on Lemma~\ref{lem:tauexists} is apportioned per clause:
\textbf{existence holds over $\Z$ for every $\ell$}, by the construction
Proposition~\ref{prop:strictpair} already contained; \textbf{homotopy uniqueness keeps
the hypothesis}. (ii) The two notation collisions that Chapters~14 and~15 flagged and
could not close are closed, as C18 and C19. (iii) \file{gl2\_refs.bib}'s
\code{harder\_eisenstein} carried the wrong title. Commit: \texttt{3c3366c}. Sources: this
document only --- Proposition~\ref{prop:strictpair} (v0.15.0), the two flagged notation
items in the v0.15.0 and v0.14.0 entries below, and \cite{serre_cours} read against the
text. No superseded document and no project note is an input, so nothing here is
\PRH{}.

\paragraph{The verdict, stated plainly.} The strict-pair construction \emph{does} establish existence of a datum
satisfying (T1)--(T3) of \cl{def:tau} over $\Z$ for every $\ell$: steps (V), (E), (F) of
\cl{prop:strictpair} never invert a stabiliser order --- (V) chooses the vertex value
\emph{inside} the fixed-point set of the twisted stabiliser, which exists by Serre's
classification of the elliptic points, and (E), (F) fill freely because edge and triangle
stabilisers are central --- and step (B), the only part that uses \ref{B1}--\ref{B3} and
$\taub$, is what makes the result a map of \emph{pairs} and is not needed for bare
existence. It does \textbf{not} establish uniqueness up to twisted homotopy: a homotopy
$H$ must have a \emph{prescribed} boundary \emph{and} be stabiliser-invariant, and only
the averaging idempotent supplies that. \textbf{No downstream hypothesis is therefore
weakened.} The statements checked, one by one, against their own proofs, and all
unchanged: \cl{thm:welldef}(b) and, through it, (c) in the general case;
\cl{prop:iotacommutes}(3c); \cl{thm:classical} (both of its two uses);
\cl{prop:barformula}'s agreement clause; \cl{prop:coneindep}(2); \cl{thm:conewelldef}(c),
(c$'$) and (e); \cl{rmk:conesteinberg}(i),(ii). \textbf{Not} re-examined, and therefore
not weakened either: everything whose $6\in\cring^{\times}$ comes from
\cl{rmk:orbifold}/\ref{H3} rather than from \cl{lem:tauexists} --- \cl{prop:bdrycohom},
\cl{thm:relations}, \cl{obl:oracle}, \cl{prop:iotasplit}, the $\Heq{2}=0$ citation of
\S\ref{sec:degreewise} --- none of which cites this lemma at all.

\paragraph{Normative.}
\begin{itemize}\setlength{\itemsep}{0pt}
\item \cl{lem:tauexists} --- restated as two clauses with their hypotheses apportioned
  (\ref{A1}). \textbf{(1) existence over $\Z$, no hypothesis on $\ell$}, new here, with
  the terminating construction (V), (E), (F) written out and the missing step supplied
  (a finite subgroup of order $4$ or $6$ is cyclic and a generator is elliptic, hence has
  a fixed point); \textbf{(2) homotopy uniqueness}, the v0.5.0 statement verbatim with
  its hypothesis and its proof unchanged. The contract box after it is restated
  accordingly: at $\ell\in\{2,3\}$ it is \emph{(2) and not the lemma} that is
  unavailable. The label already stands in the \PFIX{} column from v0.5.0 and is not
  re-counted; clause (1) is new mathematics and is what this line records.
\item \cl{rmk:tauexistsscope} \PNEW{} --- the apportionment: why the construction does
  not give uniqueness and what would settle it (\CONJ{} either way); the list of
  downstream statements checked and left alone; and the warning that existence of
  \emph{an} integral datum is not integrality of the \emph{canonical} table, which is
  \cl{rmk:tauintegral}'s subject and owner decision OP2's.
\item \cl{rmk:tauintegral} --- the ``what is \emph{not} claimed'' bullet corrected: it
  said Lemma~\ref{lem:tauexists} gives existence ``\emph{not} over $\Z$'', which
  clause~(1) makes false. What survives, and is now the bullet's point, is the
  distinction between existence of some integral datum and integrality of
  Algorithm~\ref{alg:tau}'s fills.
\item \cl{rmk:taubdryintegral} --- the parallel bullet corrected: ``that
  Lemma~\ref{lem:tauexists} holds over $\Z$ \ldots\ it does not'' narrowed to the
  uniqueness half, which is what that sentence was about.
\item \cl{prop:strictpair} --- statement unchanged; its \STD{} pointer corrected from
  Serre's Thm.~2 (the generation by $S$ and $T$) to Thm.~1(1),(3) (the fundamental domain
  and the stabilisers), the target list extended from $\{i,\zeta\}$ to
  $\{i,\zeta^{2},\zeta\}$ to match the source exactly, and the fixed-point step made
  explicit. Verified against \cite{serre_cours} Ch.~VII \S1.2 itself on 2026-09-16.
\item \S\ref{sec:notation-collisions}, row \textbf{C18} --- $N$ the Hecke level versus
  $N$ the degree-$\ell$ boundary invariant $y^{\ell}-x^{\ell-1}y$; resolved to
  $\mathrm{Nm}$, with its crosswalk row. Flagged by the v0.14.0 slot and resolved locally
  in \cl{prop:bdrychar}; that local note now names C18 instead of promising a revision.
\item \S\ref{sec:notation-collisions}, row \textbf{C19} --- the tilde, with its four live
  readings (the lift \code{\textbackslash lift}; coset transport $\widetilde F$; the
  cochain-level $\widetilde T$; Part~I's $\widetilde K$, $\widetilde X$). Resolved by a
  reading rule and \emph{no} rename --- the decoration is read from what it decorates ---
  plus the standing rule that no new tilde may be introduced without a row. Flagged by
  the v0.15.0 slot; \cl{def:cosettransport}'s notation note now names C19 and lists all
  four readings. The collision count is nineteen; \S\ref{sec:notation-collisions}'s
  heading, its opening sentence and \S\ref{sec:notation-framework}'s
  \code{\textbackslash lift} row are updated.
\end{itemize}

\paragraph{Non-normative.}
\begin{itemize}\setlength{\itemsep}{0pt}
\item \file{gl2\_refs.bib}, \code{harder\_eisenstein} --- the title was
  ``Eisenstein cohomology for $\SL_2(\Z)$ and special values of $L$-functions'', which is
  not a paper of Harder's in Invent.~Math.~89. The correct title is \emph{Eisenstein
  cohomology of arithmetic groups. The case $\GL_2$}; author, journal, volume, year and
  pages $37$--$118$ were already right and are now verified. Under \ref{CL2} the
  bibliography is part of this document, so this is an entry.
\item \file{gl2\_refs.bib}, \code{deng} --- page range $87$--$156$ added from the
  article. \code{serre\_cours} --- the note records what Thms.~1 and~2 of Ch.~VII \S1.2
  actually say, so the next pointer cannot repeat the slip fixed above. The remaining
  $29$ entries were spot-checked for author, title, journal, volume and year; no further
  error was found, and no section- or theorem-level pointer other than
  \code{serre\_cours}'s was verified against a text.
\item \cl{rmk:conesteinberg}(iii) and \cl{thm:conewelldef}(c$'$) --- the \ref{A2} caveat
  on unverified section numbers extended from \cite{borel_serre} to
  \cite{harder_eisenstein}, which it should always have covered; and
  \cite{harder_eisenstein}'s subject stated as $\GL_2$ over a number field rather than
  $\SL_2(\Z)$, matching the corrected title.
\item Formalization notes of \S\ref{sec:chainlevel} and \S\ref{sec:boundarydatum}, and
  the standing-hypotheses paragraph of \S\ref{sec:chainlevel} --- all rewritten to name
  the clause, since ``the comparison theorem is what a formalisation must not lose'' is
  now true of (2) and false of (1); (1) is the cheaper target and should be done first.
  Every reference to \cl{lem:tauexists} in the document now names the clause it means,
  except the handful that mean the lemma as a whole or its contract box.
\item The paragraph after \cl{prop:strictpair} --- rewritten: that proposition no longer
  makes a bulk-existence statement ``\cl{lem:tauexists} does not make'', because it now
  does; what it adds is strictness. Neither proof uses the other.
\item \cl{app:changelog} --- provenance counts updated from $96/33/28$ to $97/33/28$
  (\PNEW{}/\PRH{}/\PFIX{}), the \PNEW{} split from $60/36$ to $60/37$ (the one addition
  is a remark), and the tagged-line totals from $143/60$ to $144/60$. \cl{lem:tauexists},
  \cl{prop:strictpair}, \cl{rmk:tauintegral} and \cl{rmk:taubdryintegral} are labels
  already in the table and are not re-counted (\ref{CL3}, \ref{CL4}: the table counts
  labels, this appendix counts revisions).
\item \cl{app:labelindex} --- \cl{rmk:tauexistsscope} dereferenced, per \ref{L6}.
\item The date line. The v0.14.0 and v0.15.0 entries below are \textbf{not} edited: they
  are the append-only record of what those slots found and flagged, and this entry is
  where the flags are discharged.
\end{itemize}


\subsection*{v0.15.0 --- 2026-09-16 --- the pair homotopy $\etapair$ and the cone action $\Phig$}

\noindent\emph{This version.} Writes \S\ref{sec:cone}, which was a stub declaring twelve
labels, and with it the Hecke action on compactly supported cohomology as a
\emph{definition}: the pair homotopy $\etapair$ at chain level, the derived $\hder$,
the cone map $\Phig$, and the theorem that the induced operator is independent of
$\etapair$ over $\Z$ for every $\ell$. Commit: \texttt{0ae11d3}. Sources: the OQ2
verdict \S\S0, 2.1--2.10, 4, 6 (\file{hecke\_extensions\_OQ2\_verdict.md}); the
implementation plan's amendment C-A3(ii)--(iii) and its decision record and errata
(\file{hecke\_extensions\_implementation\_plan.md}:412--566, 1994--2098); the OQ2
revision audit (F5, B.1--B.4); the specification's (A2)
(\file{hecke\_extensions.tex}:789--880); \file{monodromy\_v5.tex}:1148--1265 for the cone
convention, the only superseded-document input; owner decisions OP2, OP3, OP6. \textbf{No
superseded document states any of this}: the sources are project notes, so under
\ref{CL4} every statement below is \PNEW{} or --- where a note's argument was defective
--- \PFIX{}, and the \PRH{} roster for this version is \textbf{empty}. \textbf{Five
errata land here}: \cl{err:e2} and \cl{err:e3} (\cl{prop:coneambiguity}), \cl{err:e5}
(\cl{lem:etaexists}), \cl{err:e13} (\cl{lem:conesign}, \cl{lem:derivedhomotopy}) and
\cl{err:e15} (the chapter as a whole); \cl{err:e14}'s cone half is applied by citing
\cl{thm:relations} and never restating the $T_p^2$ relation. Pages $91$--$106$
(architecture budget: $8$~pp; the overrun is $2.0\times$, the same rate as chapter~14's
mathematics, and it is proofs: every statement is proved in full rather than by
reference to the verdict, because the verdict is not a source of record). The sign
of the cone identity was \textbf{re-derived for this version} from the cone
differential alone and confirms the plan v1.2 / verdict \S2.1 / revision audit B.1
orientation (\cl{lem:conesign}); the claim that the shipped $\tau$ is not a map of
pairs was \textbf{re-verified} on \file{docs/implementation/tau\_table\_p2.txt}
(\cl{rmk:conestrict}(1): $10$ of the $21$ stored boundary slots leave $\Xbdry$).

\paragraph{Normative.}
\begin{itemize}\setlength{\itemsep}{0pt}
\item \cl{eq:contractibleuse} --- the three consequences of \ref{H4}(ii) ($\aug$ an
  isomorphism on $H_0$, $Z_1=B_1$, $Z_2=0$, over $\Z$) written once and cited at
  every use; the OQ2 revision audit's F5 found them load-bearing and unwritten in
  every source. Display label, no tag (\ref{CL4}).
\item \cl{def:pair} \PNEW{} (\cl{eq:conediff}) --- the pair, $\rres=j^{\ast}$, the
  relative complex $\ker\rres$, the cone in \file{monodromy\_v5.tex} \S8's convention
  with the notation of \S\ref{sec:notation-cone} ($q$, $(c,c_\partial)$,
  $\coneinc$/$\coneproj$), and $\Hcpt{q}:=H^{q-1}(\Cone{\rres})$. The convention is
  re-homed; the pair and $\Hcpt{\bullet}$ as objects of the contract are new.
\item \cl{lem:conerelative} \PNEW{} (\cl{eq:coneles}, \cl{eq:conedims}) --- $\rres$ is
  onto (extension by zero, equivariant because $\Xbdry$ is $\Ga$-stable, and
  \emph{not} a cochain map); the cone is a complex with the long exact sequence
  \cl{eq:coneles} whose connecting map is $+\rres$; $(\ker\rres)[1]\to\Cone{\rres}$ is a
  quasi-isomorphism, so $\Hcpt{\bullet}$ \emph{is} relative equivariant cohomology;
  $\Hcpt{q}=0$ for $q\le0$ and $q\ge3$ --- in particular $\Hcpt{0}=0$ \textbf{with no
  hypothesis on $\ell$ or $n$}, which the verdict and the plan stated only for
  $6\in\cring^{\times}$, $n$ even; and the six-term window closes under the single
  hypothesis $\Heq{2}=0$. New: no source proves any of these.
\item \cl{def:cosettransport} \PNEW{} (\cl{eq:cosettransport}), \cl{lem:cosettransport}
  \PNEW{} --- the device of \cl{thm:welldef}(b) named and generalised to any additive
  (T1)-map of any degree between the chain complexes of the pair, with its
  equivariance and \emph{four identities} ($\widetilde{\partial F}=\widetilde F d$,
  $\widetilde{F\partial}=d\widetilde F$, $\widetilde{jG}=\widetilde G\rres$,
  $\widetilde{Gj}=\rres\widetilde G$) that make every later proof a two-line
  computation. The tilde on a map (coset transport) is distinguished from the tilde on
  a simplex (\code{\textbackslash lift}, C16) at the point of definition.
\item \cl{def:pairhomotopy} \PNEW{} (\cl{eq:pairhomotopy}) --- the pair homotopy and
  the Hecke datum of the pair $(\tau,\taub,\etapair)$; with the owner-decision box
  recording OP6, the sign $\partial\etapair+\etapair\partial=j\taub-\tau j$ and
  $\hder=\widetilde{\etapair}$ with no sign, and stating what that decision does
  \emph{not} move.
\item \cl{lem:etaexists} \PFIX{} (erratum \cl{err:e5}) --- existence over $\Z$ for every
  $\ell$ by a terminating construction (\ref{A3}): the vertex step is the augmentation
  argument ($\aug(j\taub x-\tau jx)=1-1=0$), not ``a path between two vertices''; the
  edge fills exist by $Z_1=B_1$ and are \emph{unique} by $Z_2=0$; \ref{B3} makes the
  (T1)-extension well defined without inverting a stabiliser order, and the proof
  says at each step which hypothesis is doing the work.
\item \cl{lem:etaunique} \PNEW{} --- two pair homotopies differ by
  $\partial\Lam-\Lam\partial$ for a \textbf{unique} $\Lam\colon C_0(\Xbdry)\to
  C_2(\Xbulk)$, with (T1) \emph{automatic} from uniqueness; the degree-$1$ identity is
  forced. Stated with \ref{H4}(ii) as the hypothesis (audit F5) and with the
  observation, new here, that \ref{B3} is \emph{not} needed for this lemma.
\item \cl{lem:conesign} \PFIX{} (erratum \cl{err:e13}; \cl{eq:conesign}) --- for an
  \emph{arbitrary} degree-$(-1)$ map $\hder$,
  $\Dcone\Phi-\Phi\Dcone=(0,(d\hder+\hder d+\rres\Th-\Thb\rres)c)$, so $\Phi$ is a
  cochain map iff $d\hder+\hder d=\Thb\rres-\rres\Th$; and
  $\Phi\coneinc=\coneinc\Thb$, $\coneproj\Phi=\Th\coneproj$ for every $\hder$. The sign
  is asserted here and nowhere else.
\item \cl{def:coneaction} \PNEW{} (\cl{eq:derivedh}, \cl{eq:conemap}) --- the derived
  $\hder:=\widetilde{\etapair}$, the cone map $\Phig$, $T^{c}_{\Ga g\Ga}$ and
  $\Acone{q}$ by \cl{alg:assembly} on $\Cone{\rres}$; the degree bookkeeping
  ($\hder^{1}$ from vertices, $\hder^{2}$ from edges, $\hder^{0}=\hder^{3}=0$); and
  the rule that $\hder$ is never chosen at cochain level.
\item \cl{lem:derivedhomotopy} \PNEW{} --- $d\hder+\hder d=\Thb_g\rres-\rres\Th_g$ for
  the derived $\hder$, with no hypothesis on $\ell$, as a substitution into
  \cl{lem:cosettransport}: the cochain identity is \emph{derived} from the chain-level
  convention (\cl{err:e13}), which is why it has only one home.
\item \cl{thm:conewelldef} \PNEW{} --- (a) $\Phig$ is a cochain map and every map of
  \cl{eq:coneles} intertwines, no hypothesis on $\ell$ or $n$; (b) independence of
  $\etapair$ over $\Z$ for every $\ell$; (c) for $6\in\cring^{\times}$, independence of
  $\tau,\taub,s_j,g$, identification with $\cor\circ(\conj_g)_{\ast}\circ\res$ on the
  relative complex; (c$'$) the Hecke relations on the nose (\PROVED{} via the Steinberg
  form), the geometric identification (\STD{}, section numbers of Borel--Serre
  unverified) and the PD adjointness (\OPENSTAT{}), each with its own tag; (d) $\etapair\equiv0$
  iff $\tau j=j\taub$, and the entrywise identity does not license $\hder=0$; (e) the
  $6\notin\cring^{\times}$ reading, as for $\Tdc{g}$.
\item \cl{prop:coneindep} \PNEW{} (\cl{eq:conehomotopyeta}, \cl{eq:conehomotopytau}) ---
  the four independence proofs by explicit cone homotopies: $\Upsilon=(0,-\widetilde\Lam c)$
  for $\etapair$ (over $\Z$, integral, filtration-preserving, non-zero only on
  $\Cone{\rres}^{1}$); $\Upsilon=(-\widetilde Hc,\widetilde{N_\partial}c_\partial)$ for
  $(\tau,\taub)$, with $6\in\cring^{\times}$ entering only through the existence of $H$
  and the $\taub$-half unconditional; the $s_j$ summand by summand with the tables
  fixed; $g$ by transport.
\item \cl{prop:strictpair} \PNEW{} --- strict pairs exist over $\Z$ for every $\ell$, by a
  construction that pre-assigns the boundary slots and needs, beyond
  \cl{eq:contractibleuse}, only that a finite subgroup of $\Ga$ of order $4$ or $6$
  fixes a vertex of $\Xbulk$ (Serre, cited). The verdict's Prop.~2.6.2 argued through
  \cl{alg:tau}; this is self-contained and, as a by-product, an \emph{existence}
  statement for the bulk $\tau$ over $\Z$ with no hypothesis on $\ell$, which
  \cl{lem:tauexists} does not make and which does not repair that lemma's uniqueness
  half.
\item \cl{prop:coneambiguity} \PFIX{} (errata \cl{err:e2}, \cl{err:e3}) --- the
  $\Hom(\ker\rres^{q+1},\coker\rres^{q})$ ambiguity of a free cochain-level $\hder$,
  over any $\cring$; its \emph{realisation in full}, stated \textbf{over a field} with
  $Z^{0}(C^{\bullet}_{\partial})$ and never $B^{0}$, and with $\Z/\ell^{e}$, $\Z/m$
  explicitly not claimed; and the consequence that for the derived $\hder$ the
  ambiguity is a cone homotopy and vanishes.
\item \cl{prop:coneintegral} \PNEW{} --- tables over $\Z$ once, matrices per $\cring$;
  the cone over $\Z$ is a complex of free $\Z$-modules (saturated sublattices) and
  $\Cone{\rres}_{\Z}\otimes\cring=\Cone{\rres}_{\cring}$ for $6\in\cring^{\times}$
  (averaging); reductions and Bockstein connecting maps are $\Phi_{g,\Z}$-equivariant
  \emph{on the nose}; two data give the same maps on every page of the $\ell$-adic
  spectral sequence, by an integral filtration-preserving homotopy; and the
  observation, new here, that the $\etapair$ fills are unique over every $\Z/m$, so a
  mod-$m$ regime on the $\etapair$-table can only ever be a tier failure (OP2 applied).
\item \cl{prop:coneroutes} \PNEW{} (\cl{eq:route3}) --- the strict-pair route (needs
  $6\in\cring^{\times}$, via \cl{prop:coneindep}(2)) and route 3,
  $\Th'_g=\Th_g+dE+Ed$ with $E=\mathrm{ext}_0\circ\hder$, which satisfies
  $\rres\Th'_g=\Thb_g\rres$ on the nose and is cone-homotopic to $\Phig$ via
  $(Ec,0)$ --- \emph{for every $\ell$}, a distinction the sources do not draw; and the
  PD gate, stated conditionally on \cl{cl:pd} (a stub) and tagged \OPENSTAT{}.
  Owner decision OP3 (both routes committed; PD the gate) recorded.
\item \cl{obl:eta} \PNEW{} --- six clauses on the $\etapair$-table and the derived
  $\hder$ per \ref{A4}: structural; the homotopy identity through the \emph{shared}
  (T1) helper; (T1) bookkeeping and the \cl{lem:bounds} recount; the solve regime
  under OP2 (a mod-$m$ fill aborts by default, is a recorded non-integral datum behind
  an explicit setting, and is never mathematically necessary); \textbf{clause (v), that
  the stored $\hder$ is the transport of the stored $\etapair$, which is what pins
  $\Acone{q}$}; and the diagnostics \code{eta\_zero}, \code{cochain\_commutes}, recorded
  and never branched on (\ref{A6}), with $\code{eta\_zero}\Rightarrow
  \code{cochain\_commutes}$ the only asserted relation between them.
\item \cl{obl:coneexact} \PNEW{} --- eight clauses over $\cring$: the cone identity and
  $\Dcone\Phig=\Phig\Dcone$; assembly; intertwining; exactness \emph{where asserted}, with
  the non-field case recorded until a non-field elimination exists and with
  \cl{open:ell2} narrowing the asserted spots; bookkeeping with $\Hcpt{0}=0$
  unconditional and $\Heq{2}=0$ asserted only for $6\in\cring^{\times}$; the Hecke
  relations re-checked as theorems for $6\in\cring^{\times}$, recorded otherwise; the
  two cross-routes and the PD gate; and the two mandatory validation arms
  (independence; the tampered-$\hder$ negative control), the second being the
  executable form of the statement that clauses (i)--(vi) certify the sequence and
  \emph{not} the module --- any $\Phi$ of \cl{prop:coneambiguity} passes them.
\end{itemize}

\paragraph{Non-normative.}
\begin{itemize}\setlength{\itemsep}{0pt}
\item \cl{rmk:conesteinberg} \PNEW{} --- the status of the two sentences of
  \cl{thm:conewelldef}(c) that reach outside the document. (i) The operator on the
  relative complex is canonical (\PROVED{}: two strict pairs give homotopic
  $\tau_{\mathrm{rel}}$, by extending the $\Lam$ of \cl{lem:etaunique} by zero to a
  homotopy of pairs). (ii) \textbf{The Steinberg form}: for $6\in\cring^{\times}$,
  $\Hcpt{q}\cong H^{q-1}(\Ga;\Hom_{\Z}(\mathrm{St},\VR))$ with $\mathrm{St}$ the Steinberg
  module, the cone operator is the double-coset operator on that $\De$-module, and the
  Hecke relations therefore hold on $\Hcpt{\bullet}$ \emph{on the same external inputs as}
  \cl{thm:relations} --- a theorem, not a citation of geometry; \PROVED{} as a sketch
  whose derived-category bookkeeping is named as the part a formalisation must supply.
  This route appears in no source. (iii) The identification with the Borel--Serre pair
  and $H^{\bullet}_c$ of the open orbifold: \STD{} on Harder \S\S1--2 and Borel--Serre,
  with the caveat, per \ref{A2}, that the section numbers of the latter are not
  verified here and the half is to be read as \CONJ{} \emph{as a citation}; nothing in
  the document depends on it. (iv) PD is \cl{cl:pd}'s business and \OPENSTAT{} here.
\item \cl{rmk:conestrict} \PNEW{} --- (1) \COMP{} the shipped $\tau$ is not a map of
  pairs: on \file{tau\_table\_p2.txt}, $8$ of $12$ stored boundary-vertex slots and $2$
  of $9$ boundary-edge slots leave $\Xbdry$, with the specific slots named; a property of
  \cl{alg:tau}'s nearest-vertex step, not of the model. (2) The entrywise identity
  $\rres\Th_g=\Thb_g\rres$ does not license $\hder=0$; whether the residual class
  vanishes for some actual $(\tau,\taub)$ is \CONJ{} either way; the diagnostics gate
  nothing.
\item \cl{rmk:coneambiguity691} \PNEW{} --- the $(10,691)$ witness:
  $\Hom(\ker\rres^{1},\coker\rres^{0})\cong\F_{691}^{3}\neq0$ and $\dim X_3=4$, tagged
  \OPENSTAT{} \emph{in this document} because two of its three inputs are stubs of
  \cl{sec:six} and the third is [E]'s; $X_2\cong\F_{691}$ is \PROVED{} here via
  \cl{prop:bdrychar}(1).
\item \cl{rmk:conetorsion} \PNEW{} --- what holds at $6\notin\cring^{\times}$ on the model
  complexes (everything that used no hypothesis on $\ell$), what is only \emph{defined}
  there, and the two sites at which this chapter reads \cl{open:ell2}.
\item \S\cl{sec:pair}, \S\cl{sec:pairhomotopy}, \S\cl{sec:coneaction},
  \S\cl{sec:coneintegral}, \S\cl{sec:coneroutes}, \S\cl{sec:coneobligations} --- the
  chapter's subsections. Containers, no tag (\ref{CL4}).
\item \S\cl{sec:oc1sites} --- item~3 gains the two sites in \S\ref{sec:cone} that read
  \cl{open:ell2} (\cl{def:coneaction}, \cl{rmk:conetorsion}). Container, no tag.
\item The code map records that \textbf{nothing in this chapter is implemented} --- a
  census of the tree finds no $\etapair$-table, no equivariant cone, no $\Phig$, no
  $\Acone{q}$ and no cone validation --- and that the one existing cone,
  \code{CochainComplexMorphism::ConeComplex}, is the non-equivariant model of
  \file{monodromy\_v5.tex} \S8; and it fixes the shape of the phase that builds the rest
  (E5), including the golden tiers of owner decision OP7.
\item A notation point found while writing and \emph{not} resolvable here: the tilde
  denotes both the lift of a simplex (\code{\textbackslash lift}, C16) and the coset transport of a map
  (\cl{thm:welldef}(b), and now \cl{def:cosettransport}). They never meet in one formula
  and are typographically distinct (a decorated symbol versus a decorated operator), but
  \S\ref{sec:notation-collisions} should record the pair at its next revision.
  \file{gl2\_notation.tex} was not edited.
\item \cl{app:changelog} --- provenance counts updated from $77/33/25$ to $96/33/28$
  (\PNEW{}/\PRH{}/\PFIX{}), the \PNEW{} split from $45/32$ to $60/36$, the tagged-line
  totals from $114/60$ to $143/60$, and a Part~IV cone roster of $15$ new statements
  added; the sentence ``over all eleven versions'' corrected to ``over all versions to
  date''.
\item \cl{app:labelindex} --- the forty-one labels of \S\ref{sec:cone} dereferenced,
  per \ref{L6} (thirteen declared by the stub, \S\ref{sec:cone} included, twenty-eight new).
\item \cl{app:changelog-gaps} --- \S\ref{sec:cone} removed from the stub roster.
\item The abstract, the file header and the date line: fifteen of the twenty-two
  chapters are written.
\end{itemize}

\subsection*{v0.14.0 --- 2026-09-16 --- the boundary datum, and \texorpdfstring{\ref{open:ell2}}{OC-1}}

\noindent\emph{This version.} Writes \S\ref{sec:boundarydatum} and \S\ref{sec:ell2},
which were a stub declaring nine labels, and states the $\ell=2$ truncation policy
\ref{open:ell2} (\textbf{OC-1}) as a \emph{tagged working assumption}, the owner having
reserved the decision. Commit: \texttt{01b4b75}. Sources: the implementation plan's
amendment C-A2 (\file{hecke\_extensions\_implementation\_plan.md}:312--382) and
C-A3(ii.2) (:412--429); the OQ2 verdict \S\S2.6, 5.1--5.5; the specification's Stage~0
item (A1) (\file{hecke\_extensions.tex}:761--788); owner decisions OP1, OP2, OP5. All
four sources are prose or Markdown --- no superseded document states the $\Ga$-level
$\taub$, its existence proof, its integrality or the $\ell=2$ degree behaviour --- so
all but one statement below is written rather than transported. \textbf{Three errata
land here}: \cl{err:e12} (a definition and its own normative construction disagreeing
inside one document), \cl{err:e4} ($\taub$-uniqueness derived from a theorem whose
hypothesis the construction exists to avoid) and the boundary half of \cl{err:e14} (the
characters in the right-coset convention). The \PRH{} roster for this version is the
single entry \cl{def:bdrygroup}, transported unchanged from \cl{sec:six}'s definition of
$X_2$ and $X_5$. Pages
$77$--$90$ (architecture budget: $5$~pp). \textbf{The overrun is $2.8\times$ and it is
\ref{open:ell2}}: the boxed assumption, the four enumerated alternatives, the
spelled-out exactness clause, the site list of \S\ref{sec:oc1sites} and the
recommendation are about $2.5$~pp that the architecture costed as one label. The
mathematics --- \S\S\ref{sec:bdrygroup}--\ref{sec:bdryop} --- came in at about
$8.5$~pp against the same $5$~pp, i.e.\ $1.7\times$, which is the ordinary
\emph{RC+NEW} rate.

\paragraph{Normative.}
\begin{itemize}\setlength{\itemsep}{0pt}
\item \cl{B1}, \cl{B2}, \cl{B3} --- the chapter's own standing hypotheses: the boundary
  subcomplex, its components (each a contractible $1$-complex), and the centrality of
  its stabilisers. Stated for \emph{the boundary triangulation in force} and not for the
  four horocycle vertices of \cl{def:model7}, per owner direction OD-T1 as recorded in
  \cl{rmk:orbifold}; \cl{rmk:bdrymodel} proves that the shipped triangulation and every
  equivariant refinement of it satisfy them. Container labels, no tag (\ref{CL4}).
\item \cl{def:bdrygroup} \PRH{} --- $\Gbdry=\langle T\rangle\times\langle-I\rangle$
  \cl{eq:bdrysplit} and the \emph{definition} $H^{\bullet}(\partial;\VR):=
  H^{\bullet}(\Gbdry;\VR)$ (owner decision OP5, the Borel option), transported unchanged
  from \cl{sec:six}. The label previously named \code{[to be written]}. It is the half of
  \cl{err:e12} that was already right; the other half is \cl{prop:circlemodel}.
\item \cl{prop:bdrycohom} \PNEW{} --- the cohomology of $\Gbdry$ in \emph{every} degree,
  as the two-column Lyndon--Hochschild--Serre sequence \cl{eq:bdrylhs}. New: no source
  computes $H^{q}(\Gbdry;\VR)$ for $q\ge2$ at all.
\item \cl{prop:circlemodel} \PFIX{} (erratum \cl{err:e12}) --- what the circle model
  computes ($H^{\bullet}(\langle T\rangle;\VR)$, unconditionally, by Shapiro), and that
  it is $H^{\bullet}(\partial;\VR)$ \textbf{exactly when} $2\in\cring^{\times}$, with the
  failure at $\cring=\Z/2^{e}$ exhibited in degree $2$. [E] \emph{defines} the boundary
  cohomology as $H^{\bullet}(\Gbdry;M)$ at \code{:189} and then makes the circle model
  its normative construction at \code{:755--768}; the hypothesis separating the two was
  in neither statement and is now in this one.
\item \cl{prop:bdrycosets} \PNEW{} --- $B_N$, its $m$ left $\Gbdry$-cosets \textbf{for
  every $\nu$}, the bijection $B_N/\Gbdry\to\Ga g\Ga/\Ga$, the decomposition of $B_p$
  into \emph{two} $\Gbdry$-double cosets, and the right cosets $\beta_b$ as the wrong
  system. [E]~(A1) states the $\nu=1$ table and the bijection without proof; the
  general-$\nu$ count, the two-double-coset structure and the proofs are new.
\item \cl{def:taubdry} \PNEW{} --- the boundary $g$-twisted chain datum, with
  (T3$\partial$) written as ``the induced map on $H_0(\Xbdry)$ is $\pi_{\ast}$ for a map
  of sets'', which is the correct analogue of (T3) on a \emph{disconnected} complex; the
  sources say ``(T3) on every component'' without saying what that means.
\item \cl{lem:bdrycomponent} \PNEW{} --- (T1$\partial$) and (T3$\partial$) force the
  component map, and it is $c\mapsto g^{-1}c$ on $\mathbb{P}^1(\Q)$. The plan asserts
  only the two cases $\infty\mapsto\infty$ and $0\mapsto0$, parenthetically and without
  proof; the closed form and the proof are new.
\item \cl{lem:taubdryexists} \PFIX{} (erratum \cl{err:e4}) --- existence over $\Z$ \emph{for every $\ell$} by
  a terminating construction (\ref{A3}), coefficients in $\{0,\pm1\}$, and homotopy
  uniqueness over $\Z$ for every $\ell$ with $N_{\partial}=0$ in degree $1$. The verdict
  states the conclusion at \S2.6.1 but derives the uniqueness half from ``the comparison
  theorem'', whose relative-projectivity hypothesis is precisely what forces $\ell\ge5$
  in \cl{lem:tauexists} --- circular, and \cl{err:e4}. The proof here is explicit and
  never cites that theorem; it also separates the two hypotheses it uses (\ref{B3} for
  the (T1$\partial$)-extension, \ref{B2} for the fills) and says at each step which one
  is doing the work.
\item \cl{def:bdryop} \PNEW{} --- $\Thb_g$ by \cl{eq:bdrychain}, i.e.\
  \cl{eq:heckechain} with $\taub$ and \textbf{the same $s_j$, $\delta_j$ as the bulk};
  with the owner-decision box recording OP1, that the $\Ga$-level datum is normative and
  the left-coset $(\delta^{\partial}_{\infty},\delta^{\partial}_{t})$ operator is its
  Shapiro transcription and a validation route only.
\item \cl{prop:bdryop} \PNEW{} --- $\Thb_g$ is well defined with no hypothesis on
  $\ell$, and its independence of $\taub$ holds \emph{over $\Z$ for every $\ell$}, which
  the bulk statement \cl{thm:welldef}(b) does not; and in degree $0$ it is the classical
  double-coset action $\sum_{B_N/\Gbdry}\rhon(\delta)$. The degree-$0$ identification is
  new and is what licenses \cl{prop:bdrychar}(1) for the normative operator.
\item \cl{prop:bdrychar} \PFIX{} (erratum \cl{err:e14}, boundary half) --- $\VR^{T}=\cring[x,\mathrm{Nm}]_n$ with a proof (no
  citation: the invariant-ring argument is four lines and a formalisation target), the
  characters $\chi_b$ on $H^{0}$, the $\chi_{b+1}$ of the right-coset system, and the
  $n<\ell$ case of $H^{1}$ with the group-cohomology form \cl{eq:bdrygroupform}. [E]'s
  \cl{cl:chi} is \emph{not} superseded: its general-$n$ $H^{1}$ multiset is left to
  \S\ref{sec:six} and named as such. What is repaired is
  \file{hecke\_extensions\_mathematics.md}:67, which presents the characters in the
  $\beta_b$ convention as though it were \emph{the} answer with only a ``convention
  warning'' as a hedge: here the normative convention is fixed, the $\det$-twist is part
  of the statement, and \cl{prop:bdrycosets}(4) says why the two conventions split $B_p$
  in opposite ways.
\item \cl{obl:taubdry} \PNEW{} --- seven clauses, each naming its data, predicate, ring
  and failure action per \ref{A4}. Clause (iv) is where OP2 lands: on the boundary the
  mod-$\ell$ fallback \emph{cannot arise}, so a $\taub$-table with
  \code{integral = false} is an implementation defect and \textbf{aborts}, where the
  bulk clause of \cl{rmk:tauintegral} records a weaker datum instead.
\item \cl{lem:ell2unbounded} \PNEW{} --- at $\ell=2$, $H^{q}(\partial;\VR)\neq0$ for
  \emph{every} $q\ge0$, with the exact sequence and, over a field, the dimensions
  $\dim\VR^{T}$ and $2\dim\VR^{T}$. The plan states the non-vanishing; the sequence, the
  dimensions and the proof are new.
\item \cl{open:ell2} \PNEW{} --- \textbf{OC-1}, stated as a tagged working assumption
  and \emph{not} as a decision: $q_{\max}$ an explicit recorded run parameter defaulting
  to $2$, complexes built to $q_{\max}+1$, cohomology reported only to $q_{\max}$,
  exactness asserted only where all three terms are reported. The four alternatives
  \cl{oc1a}--\cl{oc1d} are enumerated and kept reachable, the recommendation is
  \cl{oc1a} with its reason (it is the only one that keeps \cl{oc1c}, a stabilisation
  theorem, free), and \S\cl{sec:oc1sites} lists every site that reads the tag so the
  owner can flip it mechanically. \ref{CL4}'s list of statement kinds does not name
  \code{open}; this is a numbered statement and is tagged as one, and that list should
  gain the kind at its next revision.
\end{itemize}

\paragraph{Non-normative.}
\begin{itemize}\setlength{\itemsep}{0pt}
\item \cl{rmk:bdrymodel} \PNEW{} --- the shipped triangulation satisfies
  \ref{B1}--\ref{B3}, and so does every $\Ga$-equivariant refinement of it; what is
  \emph{not} inherited by a refinement is named (orbit counts, table size,
  \cl{lem:bounds}). This is the chapter's discharge of OD-T1.
\item \cl{rmk:bdryoddn} \PNEW{} --- for $n$ odd and $2\in\cring^{\times}$ the boundary
  complex is zero, for the same reason as \cl{rmk:negI} in the bulk, and both fail at
  $\ell=2$.
\item \cl{rmk:taubdryintegral} \PNEW{} --- the integrality claim, its two independent
  reasons, and \textbf{what it is not}: it does not repair \cl{lem:tauexists}, and it
  does not use the comparison theorem. Written in the four-field shape because this is
  the area in which an over-claim would be hardest to see (cf.\ \cl{err:e3}).
\item \cl{rmk:ell2window} \PNEW{} --- the three complexes that coincide for
  $6\in\cring^{\times}$ and come apart at $\ell=2$ (model / group-cohomological / free
  $\Ga$-model), which of them the ``unbounded'' statement is about, and that
  \cl{open:ell2} governs only the third. It also states explicitly that whether
  $H^{q}(\Ga;\VR)$ is non-zero in every degree is \textbf{not} asserted, and that
  nothing here needs it --- the Farrell--Tate route is named with a \CONJ{} \emph{as a
  citation}, per \ref{A2}.
\item \S\cl{sec:bdrygroup}, \S\cl{sec:bdrycosets}, \S\cl{sec:taubdry},
  \S\cl{sec:bdryop}, \S\cl{sec:ell2}, \S\cl{sec:oc1sites} --- the chapter's
  subsections. Containers, no tag (\ref{CL4}).
\item A notation collision found while writing and \emph{not} resolvable here:
  \cl{cl:boundary} and \cl{cl:chi} write the degree-$\ell$ boundary invariant
  $y^{\ell}-x^{\ell-1}y$ as $N$, while $N=p^{\nu}$ is this document's Hecke level. It is
  written $\mathrm{Nm}$ in \cl{prop:bdrychar}, with the clash stated at the point of
  use; \S\ref{sec:notation-collisions} lists seventeen collisions and should gain this
  as an eighteenth at its next revision. \file{gl2\_notation.tex} was not edited.
\item \cl{prop:bdrychar}'s corroboration paragraph --- an \emph{ad hoc} polynomial
  check of all three parts over $\ell\in\{2,3,5,7,11,13\}$, $n\le24$,
  $p\in\{2,3,5,7,11,13\}$, $p\neq\ell$, $0$ disagreements, carrying \textbf{no} \COMP{}
  tag because nothing in the tree re-runs it. \cl{obl:taubdry}(vi) is where it becomes a
  registered check.
\item \cl{app:changelog} --- provenance counts updated from $64/32/22$ to $77/33/25$
  (\PNEW{}/\PRH{}/\PFIX{}), the \PNEW{} split from $36/28$ to $45/32$, the tagged-line
  totals from $94/58$ to $114/60$, and a Part~IV roster of $9$ new statements added.
\item \cl{app:labelindex} --- the thirty-five labels of \S\ref{sec:boundarydatum} and
  \S\ref{sec:ell2} dereferenced, per \ref{L6}.
\item \cl{app:changelog-gaps} --- \S\ref{sec:boundarydatum} and \S\ref{sec:ell2} removed
  from the stub roster.
\item The abstract, the file header and the date line: fourteen of the twenty-two
  chapters are written.
\end{itemize}

\subsection*{v0.13.0 --- 2026-09-16 --- \texorpdfstring{$\inv_{\Xbulk}\tau\inv_{\Xbulk}\neq\tau$}{iota tau iota != tau}, computed}

\noindent\emph{This version.} Records a computation that \S\ref{sec:involution} asked
for and that had never been performed. No statement is generalised, restricted or
withdrawn: an obligation is \emph{discharged}, and a remark that recorded a gap now
records the gap's cause. Commit: \texttt{[hash]}; the evidence is the run of
\code{HeckeIotaTauSymmetryValidation}, added to the code base as \texttt{e0b64c4} and
re-run on 2026-09-16 for this entry. Pages $43$--$50$ (was $43$--$48$).

\paragraph{Normative.}
\begin{itemize}\setlength{\itemsep}{0pt}
\item \cl{obl:iota} \PNEW{} (\ref{CL3}: the statement was written in v0.12.0 and changes
  here) --- clause (iv), the $\tau$-symmetry diagnostic, gains its \COMP{} discharge: the
  recorded value is \emph{false} on the canonical $\tau$ at every $\ell\in\{2,3\}$ of the
  shipped grid. The predicate, the ring and the failure action are unchanged, and the
  clause explicitly stays in force for any other $\tau$, vertex policy or refined
  triangulation. Clause (i) is recorded as passing. The verdict remains unbranchable
  (\ref{A6}).
\end{itemize}

\paragraph{Non-normative.}
\begin{itemize}\setlength{\itemsep}{0pt}
\item \cl{rmk:iotatausym} \PNEW{} (\ref{CL3}) --- rewritten. v0.12.0 said the hypothesis
  of \cl{prop:iotacommutes}(3b) was \emph{not known}, and attributed the doubt to the
  stored edge-path templates of \cl{alg:tau} step~2. Both halves are superseded by
  measurement: the hypothesis is \textbf{false}, with $3600$ of $10764$ slots differing,
  every comparison exact over $\Z$, and a minimal sign-free witness printed; and the
  mechanism is one level earlier --- step~1's nearest-vertex tie-break, which is by
  ascending vertex index and is therefore structurally incompatible with $\inv$-symmetry,
  since $\invK$ exchanges $v_1$ and $v_2$. The remark also records that the negative is
  \emph{certified}: $\tau^{\inv}$ satisfies (T1)--(T3) of \cl{def:tau} on every stored
  slot, so this is two admissible data for one $g$, exactly as
  \cl{prop:iotacommutes}(3a) predicts. Consequence for the document: at $\ell\in\{2,3\}$
  there is \textbf{still} no route to $[\inv^{\ast},\Tdc{g}]=0$, and the gap is now open
  for a known reason. An $\inv$-symmetric tie-break is recorded as an available repair
  that changes no normative statement but re-baselines every cochain-level reference
  value, and is \emph{not} taken --- an owner decision.
\item \cl{sec:involution} --- the chapter's ``What this chapter fixes'' paragraph and
  its code map: the chapter is no longer wholly unimplemented. The code map now names
  the validation, says which clauses of \cl{obl:iota} it discharges and which three it
  does not, and records that it re-verifies \cl{lem:iotamodel}'s unstated hypothesis
  $\mathrm{Re}\,v_1=-\mathrm{Re}\,v_2$ and certifies
  \cl{prop:iotacommutes}(3a) on the shipped tables. The validation count in that
  paragraph moves from $66$ to $67$. Container, no tag (\ref{CL4}).
\item The date line.
\end{itemize}

\subsection*{v0.12.0 --- 2026-09-15 --- the involution $\inv$}

\noindent\emph{This version.} Writes \S\ref{sec:involution}, which was a stub declaring
three labels. Commit: \texttt{7beceb2}. Sources: the implementation plan's amendment
C-A1 (\file{hecke\_extensions\_implementation\_plan.md}:287--305), the specification's
Stage~0 item (A3) (\file{hecke\_extensions.tex}:882--889) and owner decision OQ3 ---
all three prose, none of them LaTeX, so every statement below is written rather than
transported and carries \PNEW{}. There is no \PRH{} roster for this version and no
\PFIX{} line: no superseded document states anything about $\inv$, and no erratum lands
here. Pages $43$--$48$ (architecture budget: $3$~pp).

\paragraph{Normative.}
\begin{itemize}\setlength{\itemsep}{0pt}
\item \cl{I1} --- the chapter's own standing hypothesis, that the tessellation
  $\Xbulk$ and the fundamental domain $K$ are stable under $z\mapsto-\bar z$. Stated
  for an arbitrary $\Ga$-tessellation, not for the seven vertices, per owner direction
  OD-T1 as recorded in \cl{rmk:orbifold}. Container label, no tag (\ref{CL4}).
\item \cl{def:iota} \PNEW{} --- $\inv=\mathrm{diag}(-1,1)$, the conjugation table
  \cl{eq:iotaconj}, the symmetric power \cl{eq:rhoniota}, the realisations $\invK$ and
  $\inv_{\Xbulk}$ on chains, and the cochain involution \cl{eq:iotacochain}. The label
  previously named \code{[to be written]}. The definition is given on all of
  $\Mat_2(\Z)$ rather than on $\De$, which is available only because
  \cl{def:rhon} was stated there: $\det\inv=-1$, so $\inv\notin\De$.
\item \cl{lem:iotamodel} \PNEW{} --- the seven-vertex model satisfies \ref{I1}, with
  the vertex permutation and the orientation sign of every simplex. \textbf{It carries
  a hypothesis that no source states}: $\mathrm{Re}\,v_1=-\mathrm{Re}\,v_2$.
  \cl{def:model7} fixes $\mathrm{Re}\,v_0=-\tfrac12$ and $\mathrm{Re}\,v_3=+\tfrac12$
  and leaves the two middle horocycle vertices free, so ``the model is symmetric under
  $z\mapsto-\bar z$'' is true of the model only under that placement. Nothing else in
  the document uses it, and it is compatible with \cl{lem:model7coincidences}, whose
  proof needs only $\abs{\mathrm{Re}}<\tfrac12$ there.
\item \cl{prop:iotacommutes} \PNEW{} --- (1) $(\invK,\rhon(\inv))$ carries every
  relation of $\Rel$ to the inverse of a relation, hence is an automorphism of
  $\langle\Rel\rangle$ and not of the list; (2) $\inv^{\ast}$ is an involution of
  $\Ceq$ commuting with $d$, and corresponds to $\inv^{\#}$ under $\Psi$ of
  \cl{thm:explicit_complex_equiv}; (3) the commutation with $\Tdc{g}$,
  \textbf{split into three statements with three different hypotheses} --- (3a)
  unconditional, at cochain level, against the $\inv$-transported datum; (3b)
  unconditional and on the nose \emph{if} $\inv_{\Xbulk}\tau\inv_{\Xbulk}=\tau$;
  (3c) on $\Heq{q}$ under $6\in\cring^{\times}$, via \cl{thm:welldef}(a),(b). The
  stub's title claimed only the unqualified ``$\inv$ commutes with $\Tdc{g}$'', which
  is what \ref{A1} forbids; the label keeps its subject and gains its hypotheses.
\item \cl{prop:iotasplit} \PNEW{} --- the $\pm$ decomposition of $\Heq{q}$ holds when
  $2\in\cring^{\times}$ and is \emph{false in general} otherwise; at $\ell=2$ what
  holds is the Hecke-stable filtration $0\subseteq\im\mathfrak{n}\subseteq
  \ker\mathfrak{n}\subseteq\Heq{q}$ with $\mathfrak{n}=\inv^{\ast}-\id$ square-zero,
  and $\Heq{q}$ is a sum of trivial and free $\F_2[x]/(x^2)$-modules. The sources
  assert the splitting with no hypothesis at all
  (\file{hecke\_extensions\_mathematics.md}:14, ``it splits every cohomology group into
  $\pm$ eigenspaces''; \file{hecke\_extensions.tex}:888, ``$\Htot=H^+\oplus H^-$ halves
  every matrix''), which is exactly the class of unstated invertibility hypothesis
  \ref{A1} exists for.
\item \cl{obl:iota} \PNEW{} --- five clauses, each naming its data, predicate, ring and
  failure action per \ref{A4}: (i) simplicial automorphism, (ii) local-system
  automorphism, (iii) cochain involution, (iv) the $\tau$-symmetry
  \emph{diagnostic} --- recorded, never branched on (\ref{A6}) --- and (v) commutation
  plus the $\dim H^{1}_{\mathrm{eq},\pm}$ oracle, which records rather than aborts
  because its mod-$\ell$ form is a prediction. The stub's title promised a ``sign
  convention''; see \cl{rmk:iotasign} for why there is none to fix.
\end{itemize}

\paragraph{Non-normative.}
\begin{itemize}\setlength{\itemsep}{0pt}
\item \cl{rmk:iotasigns} \PNEW{} --- the six minus signs are forced by the stored
  orientation and are not a convention.
\item \cl{rmk:iotatausym} \PNEW{} --- at $\ell\in\{2,3\}$, \cl{prop:iotacommutes}(3b)
  is the only available route to the commutation, and whether the canonical $\tau$ of
  \cl{alg:tau} satisfies its hypothesis is \textbf{not known}. Recorded as a gap.
\item \cl{rmk:iotaelltwo} \PNEW{} --- the three degeneracies at $\ell=2$
  ($\rhon(\inv)=I$, no splitting, and the failure of \cl{rmk:negI}'s odd-$n$ vanishing)
  are distinct and coincide only in that regime.
\item \cl{rmk:iotapairing} \PNEW{} --- $\rhodual(\inv)=(-1)^{n}\rhon(\inv)$ and
  $\langle\rhon(\inv)v,\rhodual(\inv)\varphi\rangle=(-1)^{n}\langle v,\varphi\rangle$
  from \cl{prop:pairing}(4), with an explicit statement that the cohomological half is
  \cl{cl:pd}'s and is not available here, and that the actions of $\inv$ on
  $H^{\bullet}(\partial)$ and $\Hcpt{\bullet}$ belong to \cl{sec:boundarydatum} and
  \cl{sec:cone}. The duality-restricts-to-each-sign consequence is marked a prediction.
\item \cl{rmk:iotasign} \PNEW{} --- \code{iota\_sign}: the definition leaves no sign
  free (the two lifts $\inv$, $-\inv$ of $z\mapsto-\bar z$ differ by $(-1)^{n}$ and
  agree wherever $\Heq{\bullet}\neq0$ with $6\in\cring^{\times}$), so owner decision
  OQ3 is an \emph{assertion to be checked}, not a normalisation, and the record field
  says which convention a record was produced under.
\item \cl{app:changelog} --- the provenance counts of this appendix updated from
  $54/31/23$ to $64/36/28$ \PNEW{}, the Part~II roster of new statements extended by
  five, the retrospective-numbering box scoped to the first eleven versions, and the
  \ref{CL5} tally to ten writing slots.
\item \cl{app:labelindex} --- the fourteen labels of \S\ref{sec:involution}
  dereferenced, per \ref{L6}.
\item \cl{app:changelog-gaps} --- \S\ref{sec:involution} removed from the stub roster.
\item The abstract and the date line: thirteen of the twenty-two chapters are written.
\end{itemize}

\subsection*{v0.11.0 --- 2026-09-15 --- the changelog}

\noindent\emph{This version.} Adds the record and the rule that makes the record's own
version number well-formed. No mathematical statement changes.

\paragraph{Normative.}
\begin{itemize}\setlength{\itemsep}{0pt}
\item \cl{V5} --- the pre-$1.0$ numbering rule: while any chapter is still a
  stub the version is $0.X.Y$, the leading $0$ marking incompleteness and $X.Y$
  carrying \ref{V2}'s bump rules unchanged; the leading field becomes $1$ when the last
  stub is written. Without it the date line below would violate \ref{V1}.
\end{itemize}

\paragraph{Non-normative.}
\begin{itemize}\setlength{\itemsep}{0pt}
\item \cl{app:changelog} --- this appendix; versions 0.1.0--0.10.0 back-filled.
\item \cl{V3} --- amended to name where the record lives.
\item \cl{sec:machinery} --- the stub now points at this appendix for the record and at
  \S\ref{sec:versioning} for the convention, so the chapter that will describe the
  machinery does not have to re-derive where the changelog went.
\item \cl{app:labelindex} --- \cl{app:changelog} dereferenced, per \ref{L6}.
\item The abstract and the date line stop calling the document a skeleton: twelve of
  the twenty-two chapters and two of the four appendices are written.
\end{itemize}

\subsection*{v0.10.0 --- 2026-09-15 --- \texorpdfstring{\ref{FD4}}{FD4} proved; \texorpdfstring{\ref{FD1}}{FD1} and \texorpdfstring{\ref{FD3}}{FD3} repaired for a non-effective action}

\noindent\emph{Commit \code{3e5ffbc}.} $78\to82$ pages. Two connected repairs to
Part~I, arising from an audit of how \cl{prop:model7h4} discharges \ref{H4}(iii). No
label altered; no source file touched.

\paragraph{Normative.}
\begin{itemize}\setlength{\itemsep}{0pt}
\item \cl{lem:fd4criterion}, \cl{FD4a}, \cl{FD4b} \PNEW{} --- $\Rel$ generates the
  coincidence groupoid \emph{iff} (a) at every simplex the $\Rel$-loops reach all of
  $\mathrm{Stab}_G(\lift\sigma)$ and (b) every pair of distinct simplices with
  non-empty transporter is joined by $\Rel$. Completeness is split into local isotropy
  and global orbit content, each finitely decidable, with no enumeration of $\Ku$ or
  of $G$.
\item \cl{rmk:fd4criterion} \PNEW{} --- why the criterion has that shape.
\item \cl{lem:model7coincidences}, \cl{serreA}, \cl{serreB} \PNEW{} --- exactly four
  pairs of distinct simplices of the model are $\Ga$-equivalent, derived from Serre's
  fundamental-domain theorems alone: a moved simplex's relative interior lies in
  $\mathrm{Re}\,z=\pm\tfrac12$ or $\abs{z}=1$, hence has dimension $\leq1$.
\item \cl{prop:model7h4} \PFIX{} --- \ref{FD4} is now discharged by
  \cl{lem:fd4criterion} rather than ``by inspection'' of the relation table: the three
  elliptic rows give (a) at $v_4,v_5,v_6$, the $23$ central rows give (a) at the other
  $20$ simplices, and four seed rows give (b); the eighth seed row is shown redundant
  for \ref{FD4}. Also corrected: the $31$ rows are $8$ seeds $+\,23$ central, not
  $8+3+23$ --- the three elliptic stabiliser rows are inside the eight.
\item \cl{def:kernelN} \PNEW{} --- the ineffective kernel $N:=\ker(G\to\mathrm{Aut}(\Ku))$,
  normal, contained in every stabiliser, with every transporter a union of $N$-cosets;
  $N=1$ in the free instance and $N=\{\pm I\}$ for the model. A non-effective action
  becomes an explicit supported case of Part~I instead of an unstated exception.
\item \cl{ass:fund_domain}, \ref{FD1} \PFIX{} --- false as written. $-I$ acts trivially
  on the tessellation while being a non-identity group element, so every simplex admits
  the coincidence $(\sigma,\sigma,-I)$; the old $K_{\mathrm{cl}}$ was all of $K$,
  $\widetilde D$ was empty, and \ref{FD1} failed outright. $K_{\mathrm{cl}}$ is now
  generated by the simplices a coincidence \emph{moves}, and \ref{FD1} additionally
  asserts that $\widetilde D$ is dense in $\abs{K}$ --- which
  \cl{thm:fund_group_generators}'s proof already used and \ref{FD3} only pretended to
  supply.
\item \cl{ass:fund_domain}, \ref{FD3} \PFIX{} --- disjointness is now asserted over
  $G/N$: $\gamma\widetilde D\cap\widetilde D=\emptyset$ for every $\gamma\notin N$, and
  nothing else. As written it failed at $\gamma=-I$ once $\widetilde D$ was repaired.
\item \cl{rmk:fdconsequences}, \cl{eq:fddecomp} \PNEW{} --- the former conclusion
  ``hence $\abs{\Ku}=\bigsqcup_{\gamma}\gamma\widetilde D$'' was a non-sequitur, wrong
  even for an effective action, since the union omits the translates of
  $\abs{K_{\mathrm{cl}}}$; it holds only for a half-open domain, which this document
  neither needs nor builds. Removed and replaced by a proved decomposition, disjoint in
  its first half and not in its second.
\item \cl{rmk:fd13centre} \PNEW{} --- why $N$ is quantified over and never divided out,
  stated once and in Part~I: the coefficient action does not factor through $G/N$ ---
  $\rhon(-I)=(-1)^n\,\mathrm{id}$ is the odd-$n$ vanishing mechanism --- so quotienting
  would delete exactly the constraint \ref{FD4} imposes. Owner decision, 2026-09-15,
  taken against the alternative of restating Part~I over the effective quotient.
\item \cl{rmk:twosettings} \PFIX{} --- both bullets updated: the free instance has
  $N=1$, the model is non-effective with $N=\{\pm I\}$.
\item \cl{rmk:fd4computed} \PNEW{}, \COMP{} --- \ref{FD4} for the model re-verified as
  a finite computation: \code{HeckeTessellationValidation} computes $\langle\Rel\rangle$
  and asserts exactly conditions (a) and (b); re-run, PASS. Recorded with the caveat
  that the computation supplies $\langle\Rel\rangle$ but not the true stabilisers and
  orbits it is compared against.
\end{itemize}

\subsection*{v0.9.0 --- 2026-09-15 --- Part I re-homed from [v5]; erratum E7 closed}

\noindent\emph{Commit \code{d5a2b2d}, writing slots 5--6.} Chapters
\ref{sec:framework}--\ref{sec:systematiccomputation}, $14$ pp; $65\to78$ pages. All $26$
assigned labels carried verbatim, each defined exactly once, $0$ removed, $32$ new
under \ref{L5}.

\paragraph{Normative.}
\begin{itemize}\setlength{\itemsep}{0pt}
\item \cl{thm:explicit_complex_equiv} \PRH{} --- \cl{err:e7} closed: the label is now
  defined in this document under exactly the string six source files cite, and its
  part~(1) is what those sites assert. The erratum entry records five files; there are
  six sites in six files (it lists all six paths but counts a validation
  \code{.h}/\code{.cpp} pair as one).
\item \cl{thm:explicit_complex_equiv} \PFIX{} --- freeness of the action is a
  hypothesis of part~(2) \emph{only}, not of part~(1). [v5] never marked the use, and
  the marking is what licenses this pipeline --- whose action is not free --- to cite
  part~(1).
\item \cl{def:relationdatum} \PNEW{} --- relations, transporters and the coincidence
  groupoid. [v5]'s relation datum is indexed by ordered pairs $(u,v)$ with $u<v$ and
  therefore structurally cannot express a self-relation $(\sigma,\sigma,\gamma)$ with
  $\gamma\neq1$, while the model's $\Rel$ consists of seed rows \emph{and} $23$ central
  $-I$ rows. Transcribed as written, the re-homed theorem would not have covered the
  model it is cited for.
\item \cl{ass:fund_domain}, \ref{FD1}--\ref{FD3} \PRH{} --- [v5]'s fundamental-domain
  hypothesis, written out as three separately citable and separately repairable
  conditions. All three are repaired in v0.10.0.
\item \cl{ass:fund_domain}, \ref{FD4} \PNEW{} --- the completeness condition that makes
  $\Psi$ surjective. \ref{H4}(iii)'s forward reference ``presents the quotient in the
  sense of the framework'' had, until this version, resolved to nothing.
\item \cl{rmk:twosettings} \PNEW{} --- the two instances of the framework this document
  uses: the free one, and the model. Repaired in v0.10.0, once the model's action was
  recognised as non-effective.
\item \cl{def:orderedrelation}, \cl{def:quotientcomplex}, \cl{def:saturated},
  \cl{def:bdryint} \PNEW{} --- the framework's definitional layer, stated rather than
  assumed.
\item \cl{lem:eqgroupoid} \PNEW{} --- one groupoid argument replacing [v5]'s
  well-definedness case analysis, which offered ``interior $\Rightarrow h=e$'' or
  ``boundary $\Rightarrow h=g_r^{-1}$'' and had no case for
  $\gamma_1^{-1}\gamma_2\in\mathrm{Stab}(\lift\sigma)$ --- the stabiliser case the model
  needs. It also makes visible that dropping the central rows breaks surjectivity of
  $\Psi$, which is the mathematical content of the \code{FIX-negI} change and until now
  lived only in an implementation-plan bullet.
\item \cl{prop:Bqimage} \PNEW{} --- $\Bq{q}$ computed as the image basis of
  $d\,\Cq{q-1}$ equals [v5]'s $d\,W^{q-1}$, proved column for column. The shipped code
  takes the former and [v5] \S7.2 computes the latter; the equality was justified only
  in a source comment.
\item \cl{prop:saturated_complexes} \PFIX{} --- [v5] proved it ``via explicit
  construction'' and gave no construction; supplied.
\item \cl{prop:injectivity} \PFIX{} --- its decomposition was asserted; now derived.
\item \cl{thm:fund_group_generators} \PFIX{} --- [v5]'s $70$-line path induction
  replaced by a $15$-line simplicial argument, which also yields the non-free case.
\item \cl{prop:fund_group_action} \PRH{} --- relation elements. Its part~(1), that a
  non-empty transporter is a left coset of the stabiliser, is what \cl{lem:fd4criterion}
  runs on in v0.10.0.
\item \cl{def:storedbases}, \cl{prop:ACq}, \cl{prop:AZq}, \cl{prop:projector},
  \cl{eq:projcert}, \cl{alg:cohomology} \PNEW{} --- chapter
  \ref{sec:systematiccomputation}. [v5] \S7.2 is an informal recipe; the stored objects,
  the two cochain systems, the projector and its certificate are stated and proved here
  for the first time.
\item \cl{eq:psiinverse}, \cl{eq:equivariance}, \cl{eq:naivecoboundary},
  \cl{eq:quotientcoboundary}, \cl{def:naivecochains}, \cl{def:eqcohomology},
  \cl{cor:threecategories} \PNEW{} --- the displayed forms the later chapters transport
  formulas through.
\item Nine hypotheses [v5] left implicit are now stated at the point of use
  (\ref{A1}): action without inversions, so that
  $\mathrm{Tr}(\sigma,\sigma)=\mathrm{Stab}$; completeness \ref{FD4}; injectivity of
  $\sigma\mapsto\lift\sigma$; contractibility of $\abs{\Ku}$; and ring versus field ---
  chapters \ref{sec:localsystems}--\ref{sec:explicitcochaincomplex} hold over any
  commutative ring, chapter \ref{sec:systematiccomputation} over a field. [v5] wrote
  $F$ once and never marked its use, and $\cring=\Z/\ell^{e}$ is in scope.
\item \emph{Re-homed unchanged} (content as in [v5], notation crosswalk applied):
  \cl{sec:framework}, \cl{ass:K_contractible}, \cl{ass:edge_non_ident},
  \cl{ass:open_interior}, \cl{lemma:pR_simplicial}, \cl{def:universal_cover},
  \cl{prop:universalcover}, \cl{rmk:injectivityofi}, \cl{ex:S1},
  \cl{sec:localsystems}, \cl{def:matrix_ls}, \cl{prop:matrix_rep_equiv},
  \cl{thm:lcs_rep_equiv}, \cl{thm:rep_ls_equiv},
  \cl{sec:explicitcochaincomplex}, \cl{rmk:universalcovercohomology},
  \cl{thm:cohom_equiv}, \cl{def:eq_cochains}, \cl{ass:consistent_face_id},
  \cl{lemma:d_maps_C_eq}, \cl{sec:systematiccomputation}.
\end{itemize}

\paragraph{Non-normative.}
\begin{itemize}\setlength{\itemsep}{0pt}
\item \cl{err:e11} applied --- the $c/z/b$ basis-name scheme, $w^{(q)}$ for $W^q$, and
  [v5]'s $\beta^{(q)}_j$ renamed $z^{(q)}_{1,j}$, since C8 reserves $\beta$ for the
  right-coset representative.
\item \cl{rmk:derivedmatrices}, \cl{rmk:lcsnotfree}, \cl{rmk:psiuse},
  \cl{rmk:psiinjectivelift}, \cl{rmk:facesautomatic}, \cl{rmk:cost} \PNEW{} --- what
  the equivalence does not say, what the theorem is used for, and the cost of the
  algorithm.
\item \cl{err:e8} and \cl{err:e10} referenced, not restated; \cl{err:e19} and
  \cl{err:e20} deliberately not duplicated, chapters
  \ref{sec:framework}--\ref{sec:systematiccomputation} touching neither $\mu$ nor the
  bar formula.
\item Recorded, and repaired in v0.10.0: \cl{prop:model7h4}'s ``\ref{H4}(iii) follows by
  inspection'' now discharges a stronger condition than before --- which it does, but
  by inspection of a $31$-row table rather than by an argument.
\end{itemize}

\subsection*{v0.8.0 --- 2026-09-15 --- C17, and two stale uniqueness claims}

\noindent\emph{Commit \code{24c235e}, slot 4 follow-ups.} $65$ pages, unchanged. Three
corrections, all of them claims that were true when written in v0.7.0 and stopped being
true later in v0.7.0 when \cl{prop:mu}(3) was added (\ref{CL3}).

\paragraph{Normative.}
\begin{itemize}\setlength{\itemsep}{0pt}
\item \cl{sec:notation-collisions}, row C17, with its \cl{sec:notation-crosswalk} row
  --- the relation element of $\Rel$ versus the Hecke element
  $g=\smat{1}{0}{0}{N}$; the two meet in one formula, \cl{thm:xi}(1). Relation elements
  are $\gamma\in\Ga$; $g$ is the Hecke element only. Checked against what the sources
  write (subscripted $g_r$, not bare $g$).
\item \cl{rmk:indeplevel} \PFIX{} --- the enumeration of licensed cochain-level
  bit-identities omitted $\omega_{\bullet}\circ\mu_{\bullet}\circ\alpha_{\bullet}$
  against $\Th_g$. Added with its hypothesis: both sides consume the same $g$, the same
  $s_j$ and the same $\tau$-table, per \cl{obl:agreement}. The bar-cocycle route is
  correctly excluded and was left unchanged.
\item \cl{prop:tworoutes} \PFIX{} --- its contract box called the proposition ``the one
  place in this document where bit-identity of cochain-level matrices is asserted''. It
  is one of two; the other, \cl{prop:mu}(3), is now named.
\end{itemize}

\paragraph{Non-normative.}
\begin{itemize}\setlength{\itemsep}{0pt}
\item \cl{def:model7rel}, \cl{def:contragredient}, \cl{sec:induced} --- C17
  conformance, $g\to\gamma$. The collision was wider than slot 4 reported:
  \cl{def:model7rel} still wrote $g$ in four places (table column, prose, contract box,
  and the contragredient local-system definition), and chapter \ref{sec:equivalence}
  carried a paragraph describing the discrepancy as open.
\item \cl{sec:degreewise}, Formalization note --- it called \cl{prop:tworoutes} ``the
  only claim in the document that a bit-level agreement is mathematically forced rather
  than empirically observed''. Corrected to two, with the other named and its proof
  route given. The forced-versus-observed distinction is intact; only the count was
  wrong, and a false uniqueness claim in a formalization note misleads precisely the
  reader it is addressed to.
\end{itemize}

\subsection*{v0.7.0 --- 2026-09-15 --- induced systems, \texorpdfstring{$\Xig$}{Xi}, \texorpdfstring{$\Th_g$}{Theta}, and the equivalent formulations}

\noindent\emph{Commit \code{592c6e8}, writing slot 4.} Chapters \ref{sec:induced} and
\ref{sec:equivalence}, $13$ pp; $53\to65$ pages. All $17$ assigned labels carried
verbatim, $0$ unplaceable, $18$ new.

\paragraph{Normative.}
\begin{itemize}\setlength{\itemsep}{0pt}
\item \cl{prop:mu}, \cl{sec:threemap} \PNEW{} --- the middle cochain map, in closed
  form, derived and proved. The source's ``Recovering the three-map factorization''
  remark is \emph{circular}: $\beta^{\bullet}$ is ``determined by
  $\beta^{\bullet}\alpha^{\bullet}=\Xig$ on $\mathrm{im}\,\alpha^{\bullet}$ and by the
  analogous formula in general'', which is not a definition. A map the contract claimed
  to specify was in fact specified only by the code, and \code{HeckeCrossChecks}
  implements what is now proved.
\item \cl{eq:evalone} --- the same remark's composite is \emph{ill-typed}:
  $\rhon(g)$ maps $V\to V$ while $\phi^{L'}$ lands in $\mathrm{Ind}_{\Ga'}V$. Fixed by
  ordering the operations --- evaluate at $1$, then apply $\rhon(g)$.
\item \cl{prop:lhcompliant} \PFIX{} --- part~(3) added: $M_H(-I)=(-1)^n I$, the odd-$n$
  vanishing mechanism. \cl{def:inducedls} and \cl{prop:lhcompliant} predate the central
  relations and gave no $M_H(-I)$ at all.
\item \cl{rmk:nogo}, \cl{sec:nogo} \PFIX{} --- \cl{err:e19} applied: the map $\beta$ of
  [v5] \S12 is a no-go, not a repair; the middle map is $\mu$.
\item \cl{rmk:barformula}, \cl{prop:barformula}, \cl{eq:barcorrected} \PFIX{} ---
  \cl{err:e20} applied: the corrected degree-$1$ group-cochain formula carries the
  coset permutation $j\mapsto\gamma\ast j$ explicitly. Both keep their labels
  (\ref{L2}).
\item \cl{rmk:kappaexact} \PNEW{} --- the corrected formula would otherwise invert a
  matrix of determinant $N$, which \cl{obl:integral} forbids; the implementation uses
  the exact division $N\kappa=\adjg(g)\,s_k\gamma\,s_j^{-1}g$.
\item \cl{prop:tworoutes}, \cl{sec:chainpicture} \PNEW{} --- the two bracketings agree
  bit for bit, proved term by term rather than observed.
\item \cl{eq:psiinv} --- well-definedness of the inverse transport, from the
  strict fundamental domain \emph{plus} the stabiliser and central relations.
\item \cl{thm:xi} \PFIX{} --- $\lift{\partial\sigma}=\partial\lift\sigma$, from
  \ref{H4}(iii), is stated as the hypothesis part~(2) uses; and the proof's own
  collision (an oriented simplex called $\tau$, while $\tau$ is the chain datum
  everywhere else) is resolved by C4.
\item \cl{alg:assembly} \PFIX{} --- $P_q H_{\mathrm{eq}}=I$ and $P_q B_{\mathrm{eq}}=0$
  named as the defining property the algorithm consumes, rather than left to the
  reader.
\item \cl{rmk:omegaconvention} \PNEW{} --- \cl{err:e9}: the corestriction convention
  stated once, with the $\gamma\to\gamma^{-1}$ substitution shown.
\item \cl{rmk:thetachoices}, \cl{rmk:shapirouse}, \cl{sec:xitheta}, \cl{sec:alphaomega},
  \cl{sec:assembly}, \cl{sec:inducedmodule}, \cl{sec:barformula}, \cl{sec:degreewise}
  \PNEW{} --- which choices $\Th_g$ depends on, what the Shapiro statement is and is
  not for, and the degree-wise specialisations.
\item $6\in\cring^{\times}$ is isolated to exactly two statements of the two chapters;
  every other statement is over an arbitrary commutative $\cring$, and every chapter-%
  \ref{sec:induced} statement is parametric in $(K,\Rel)$, per owner direction OD-T1.
\item \emph{Re-homed unchanged}: \cl{sec:induced}, \cl{def:inducedmodule},
  \cl{prop:blockmatrices}, \cl{eq:MH}, \cl{def:inducedls}, \cl{prop:alpha},
  \cl{prop:omega}, \cl{def:xi}, \cl{eq:xidef}, \cl{eq:xiexplicit},
  \cl{sec:equivalence}.
\end{itemize}

\paragraph{Non-normative.}
\begin{itemize}\setlength{\itemsep}{0pt}
\item Three route comparisons recorded at the level each actually holds, none
  overclaimed: fused versus route~2 bit-identical (\cl{prop:tworoutes});
  $\omega_{\bullet}\mu_{\bullet}\alpha_{\bullet}$ versus $\Th_g$ bit-identical on the
  same data (\cl{prop:mu}(3)); and the bar-cocycle route versus $A^{(1)}(T_p)$
  \emph{spectral only} --- characteristic and minimal polynomial, never entrywise ---
  because that route produces its matrix in an unrelated basis of $H^1$.
\end{itemize}

\subsection*{v0.6.0 --- 2026-09-15 --- C16 and the lift macro}

\noindent\emph{Commit \code{3560525}, slot 3 follow-ups.} $52\to53$ pages. Two notation
gaps closed before slot 4, which carries the heaviest notation load, and one owner
direction recorded.

\paragraph{Normative.}
\begin{itemize}\setlength{\itemsep}{0pt}
\item \cl{sec:notation-framework}, row C16 --- the contract had \emph{no}
  symbol for the inclusion $K\to\Ku$: both sources write $i(\sigma)$, which C15
  retired. The macro \code{\textbackslash lift} is registered, with a collision row and
  two crosswalk rows ($i(\sigma)\to\lift{\sigma}$; $s_i,\delta_i,n^{(i,\sigma)}_k\to
  s_j,\delta_j,n^{(j,\sigma)}_t$), so the mechanical crosswalk later slots run cannot
  reintroduce either collision.
\item \cl{sec:notation-hecke} --- the notation chapter's own Hecke table was
  stale against the rule it declares. It now writes $\Ga=\bigsqcup_j\Ga''s_j$ and
  $\delta_j=s_j^{-1}g$, matching chapter \ref{sec:chainlevel}, with C7, C10, C14 and
  the $k$-bullet re-indexed. Historical quotations of the superseded documents keep
  their $i$.
\end{itemize}

\paragraph{Non-normative.}
\begin{itemize}\setlength{\itemsep}{0pt}
\item \cl{sec:operatordef} --- the distinguished lift is now written $\lift\sigma$ via
  the macro (C15, C16).
\item \cl{rmk:tauintegral} --- $\tau(\partial\tilde t)\to\tau(\partial\lift\sigma)$: $t$
  was naming a triangle while \ref{N4} reserves it as a running index.
\item \cl{err:e11} --- $\delta_i\to\delta_j$. Checked to be the contract's own voice and
  not a quotation of the defective text: an erratum's quoted statement must keep the old
  spelling.
\item \cl{rmk:orbifold} --- owner direction OD-T1 (2026-09-15) recorded: a correct
  orbifold treatment at $\ell\in\{2,3\}$ and over $\Z$ is expected to require a
  \emph{refinement} of the $7$-vertex triangulation. It is marked explicitly as a design
  direction and not a theorem; nothing stated in this document depends on it, and its
  force is on the authors of \S\ref{sec:orbifoldmodel} and \S\ref{sec:boundarydatum},
  who must phrase their definitions to admit a refinement.
\end{itemize}

\subsection*{v0.5.0 --- 2026-09-15 --- the Hecke datum, \texorpdfstring{$\tau$}{tau}, and the definition of the operator}

\noindent\emph{Commit \code{ef87af6}, writing slot 3.} Chapters \ref{sec:heckedatum}
and \ref{sec:chainlevel}, $13$ pp; $39\to52$ pages. All $17$ assigned labels carried
verbatim --- including \cl{alg:tau} ($36$ code citation sites) and \cl{obl:tau} ($28$)
--- $0$ unplaceable, $19$ new.

\paragraph{Normative.}
\begin{itemize}\setlength{\itemsep}{0pt}
\item \cl{prop:cosets} \PNEW{} (extension of a re-homed statement) --- stated and proved
  for every $\nu\geq1$, via the bijection with $\mathbf{P}^1(\Z/N)$ and the count
  $p^{\nu}+p^{\nu-1}$. The source states it only for $g=\smat{1}{0}{0}{p}$, yet the
  pipeline runs at $g=\smat{1}{0}{0}{p^{\nu}}$, and the code said so outright, calling
  its own $\nu\geq2$ case an undocumented generalisation. The Hermite-normal-form
  citation the source leaned on is replaced by a proof.
\item \cl{def:cosetreps}, \cl{eq:cocycle}, \cl{lem:cocycle} \PNEW{} --- the closed-form
  lifts are pinned \emph{binding as written}, with the class map and its cocycle
  property. A generic extended-Euclid B\'ezout solve returns a different, equally valid
  lift, and the cochain-level matrices depend on the choice, so an implementation that
  ``reimplements the obvious thing'' would produce different matrices with no test to
  catch it.
\item \cl{prop:fullassembly}, \cl{eq:fullassembly} \PNEW{} --- the primitive-to-full
  assembly, promoted from a parenthetical to a proposition proved at cochain level with
  no hypothesis on $\ell$. \cl{err:e1}.
\item \cl{sec:primfull}, \cl{rmk:whichoperator}, \cl{rmk:rawprimitive} \PNEW{} ---
  \cl{err:e1} made unrepeatable: the two objects are named in a contract box before
  anything can refer to them, unqualified $T_{p^{\nu}}$ \emph{always} means the full
  operator, and the entry-point table sits in the code map.
\item \cl{thm:relations} \PFIX{} --- now carries $6\in\cring^{\times}$, the hypothesis
  its own proof needs, since it routes through \cl{obl:algebra} and \cl{thm:classical}.
  The source omitted it, and so did a source-code comment, which was corrected
  separately in v0.6.0's commit.
\item \cl{cor:ulrecursion} \PFIX{} --- split: part~1 needs $\ell\geq5$, part~2 does not.
\item \cl{thm:welldef} \PFIX{} --- split into what holds unconditionally
  (equivariance, part~(a)), what needs $6\in\cring^{\times}$ (part~(b)), and what is
  unconditional for the transported choices but conditional in general (part~(c)). The
  source blurred the three.
\item \cl{lem:tauexists} \PFIX{} --- $6\in\cring^{\times}$ made explicit \emph{and}
  derived, from the elliptic stabiliser orders $4$ and $6$ of \cl{prop:model7h4}(iv)
  rather than re-assumed; acyclicity from \cl{prop:model7h4}(ii).
\item \cl{def:tau} \PFIX{} --- widened from $\Fl$ to a general commutative $\cring$,
  recorded in \cl{rmk:taurings} as a widening and not a retargeting (\ref{L2}), because
  the integrality requirement and ``tables once, matrices per $\cring$'' both need it.
\item \cl{prop:diamond} \PRH{} --- proved with no hypothesis on $\ell$.
\item \cl{rmk:tauintegral} \PNEW{}, \COMP{} --- $\tau$'s integrality stated narrowly:
  steps 1--2 of \cl{alg:tau} are integral unconditionally, only the step-3 triangle fill
  can lose it, and there is \emph{no} theorem that the fills solve over $\Z$ ---
  \cl{lem:tauexists} gives existence over $\cring$ with $6$ invertible. Verified per
  table, with the measured grid in the statement: $0$ of $66$ grid points used the
  mod-$\ell$ triangle solve, maximum certified bound $B=196$, $69/69$ passed.
\item \cl{rmk:relcoeff} \PNEW{} --- the coefficient of $T_{p^{\nu-1}}$ is $p^{\,n+1}$,
  not $(p+1)p^{\,n}$; \cl{err:e6}, settled by the degree count.
\item \cl{rmk:T3}, \cl{rmk:badprimesingular} \PNEW{} --- what (T3) is for, and why the
  singularity at $p=\ell$ is not a defect of the construction.
\item \emph{Re-homed unchanged}: \cl{sec:heckedatum}, \cl{sec:correspondencedatum},
  \cl{sec:operatordefs}, \cl{def:operators}, \cl{eq:fullop}, \cl{sec:relations},
  \cl{sec:chainlevel}, \cl{sec:correspondence}, \cl{eq:corconjres}, \cl{sec:taudatum},
  \cl{sec:operatordef}, \cl{def:hecke}, \cl{eq:heckechain},
  \cl{sec:classicalcomparison}, \cl{thm:classical}, \cl{rmk:indeplevel}.
\end{itemize}

\paragraph{Non-normative.}
\begin{itemize}\setlength{\itemsep}{0pt}
\item Two citations that could not be checked were removed under \ref{A2}: the Hermite
  normal form is now proved inside \cl{prop:cosets}, and Shimura Ch.~3 replaces a pair
  of unverifiable theorem numbers.
\end{itemize}

\subsection*{v0.4.0 --- 2026-09-15 --- C15: \texorpdfstring{$i$}{i} is \texorpdfstring{$\sqrt{-1}$}{sqrt(-1)} only}

\noindent\emph{Commit \code{6d8abae}.} $39$ pages, unchanged; the change is entirely in
chapter \ref{sec:notation}. Slot 2 had found a fifteenth symbol collision the notation
chapter did not list and resolved it locally, but nothing told the remaining chapter
authors to do the same.

\paragraph{Normative.}
\begin{itemize}\setlength{\itemsep}{0pt}
\item \cl{N4} --- amended from ``a running index is $i$, $j$, $t$ or $s$'' to
  ``a running index is $t$, $j$ or $s$; $i$ is never an index''. Owner decision,
  2026-09-15; the rejected alternative, renaming the elliptic vertex value, churns a
  normative object for no gain.
\item \cl{sec:notation-collisions}, row C15 --- $i$ as a running index versus
  $i=\sqrt{-1}=v_6$, the elliptic point of the model; collision count updated.
\item \cl{sec:notation-crosswalk} --- crosswalk row ($i$ as an index $\to$
  $t/j/s$). It earns its place: all three source documents use $i$ heavily as a running
  index, so the mechanical crosswalk later slots run over re-homed text would otherwise
  reintroduce the collision chapter by chapter.
\end{itemize}

\paragraph{Non-normative.}
\begin{itemize}\setlength{\itemsep}{0pt}
\item \cl{sec:notation-open} --- C15 recorded as decided, not as an open item.
\end{itemize}

\subsection*{v0.3.0 --- 2026-09-15 --- conventions, the 7-vertex model, \texorpdfstring{$\Sym^n$}{Sym\^{}n} and the contragredient}

\noindent\emph{Commit \code{1c84f7c}, writing slot 2.} Chapters \ref{sec:conventions}
and \ref{sec:symconv}, $9$ pp against a $7$ pp budget; $30\to39$ pages. All $12$
assigned labels carried verbatim, $17$ new. This is the first re-homing slot, and it
established the notation-crosswalk workflow for the rest.

\paragraph{Normative.}
\begin{itemize}\setlength{\itemsep}{0pt}
\item \cl{H4}, \cl{prop:model7h4} \PNEW{} --- that $\Ku$ is the $\Ga$-tessellation of
  the truncated half-plane is used by the source's \cl{rmk:orbifold} but the source
  never states that it is contractible and two-dimensional, nor the stabiliser
  inventory. Both are load-bearing here and in chapter \ref{sec:cone}. Written as
  \ref{H4}, with \cl{prop:model7h4} proving (i)--(iv). Contractibility goes by the
  horoball retraction, which needs $Y>1$ --- a condition the sources do not state
  either, and which is now in the proof.
\item \cl{H1}, \cl{H2}, \cl{H3} --- the standing hypotheses, kept under those
  exact strings because they are cited by name from four source files and four plan
  documents.
\item \cl{sec:contragredient}, \cl{def:contragredient}, \cl{eq:adjrules},
  \cl{eq:dualaction}, \cl{prop:pairing}, \cl{eq:pairadjoint}, \cl{lem:symdual},
  \cl{eq:BnGram} \PNEW{} --- the contragredient $\Dn$, the perfect pairing, and the
  integral invariant form, per owner decision OQ4. The adjointness identity
  $\langle\rhon(g)v,\phi\rangle=\langle v,\rhodual(\adjg g)\phi\rangle$ is exact,
  integral and carries no determinant factor; the \emph{diagonal} identity does carry a
  $p^{\,n}$ twist, so [E]'s ``$T_q$ self-adjoint once $\Dn$ carries the adjugate
  structure'' is true of the adjointness form only and must be re-homed that way.
\item \cl{rmk:selfadjoint} \PNEW{} --- states that \cl{prop:pairing} asserts nothing
  about cohomology. The architecture's assigned caption for it, ``self-adjointness of
  $T_p$'', over-promises: the operators are not defined until chapters
  \ref{sec:heckedatum}--\ref{sec:induced}. Only the coefficient-level content is proved
  here, and chapter \ref{sec:six} must not assume more.
\item \cl{rmk:divided} \PNEW{} --- divided powers, and what is not claimed.
\item \cl{rmk:orbifold} \PFIX{} --- \cl{err:e8} at its first landing site: the torsion
  is stated as a hypothesis and the place where ignoring it changes an answer is named,
  replacing ``issues with torsion, which we ignore''.
\item \cl{rmk:negI} \PFIX{} --- \cl{err:e8} at its second: the odd-$n$ vanishing
  argument replaces the same phrase.
\item \cl{sec:conventions} --- \cl{err:e16} applied: the composite-index
  operator is $T_N$, never $T_n$, $n$ being the $\Sym$ index.
\item \cl{def:model7}, \cl{def:model7rel}, \cl{eq:negI}, \cl{eq:flatindex},
  \cl{eq:rhondef} \PNEW{} --- the model, its relation table with a code tag per row,
  the central relation, the flat index, and the displayed $\rhon$.
\item \emph{Re-homed unchanged}: \cl{sec:model}, \cl{sec:symconv}, \cl{def:rhon},
  \cl{prop:symconv}, \cl{def:semigroup}, \cl{lem:detsym}.
\end{itemize}

\paragraph{Non-normative.}
\begin{itemize}\setlength{\itemsep}{0pt}
\item The model was verified against the code rather than the prose: $7$ vertices, $11$
  edges, $5$ faces ($\chi=1$), hence $\abs{\Rel}=8+23=31$.
\item \file{gl2\_refs.bib} gains \code{serre\_cours}, the fundamental-domain,
  generation and stabiliser-order citation \ref{A2} requires.
\item A fifteenth symbol collision ($i$) was found and resolved locally only; closing it
  needed an owner call and became C15 in v0.4.0.
\end{itemize}

\subsection*{v0.2.0 --- 2026-09-15 --- scope, the errata appendix, and the \texorpdfstring{$T_{p^{\nu}}$}{T\_\{p\^{}nu\}} semantics}

\noindent\emph{Commit \code{7f1f5f2}, writing slot 1.} Chapter \ref{sec:scope} ($3.5$
pp) and Appendix~\ref{app:errata} ($7.5$ pp, $20$ entries); $18\to30$ pages. $38$ labels
added, $0$ removed.

\paragraph{Normative.}
\begin{itemize}\setlength{\itemsep}{0pt}
\item \cl{sec:scope} --- the [C]/[E] cut, moved: predictions and plain facts
  about the objects had been on the same side of it. Nine statement groups --- the six
  groups $X_1,\dots,X_6$ and their long exact sequence, the boundary groups and their
  characters, Poincar\'e duality, the amalgam presentation and $\dim H^1$, and the
  Dickson/Bockstein statement --- move into this document, recorded as \cl{err:e18}.
\item \cl{sec:status} --- the STATUS vocabulary. \code{\textbackslash STD}
  takes a \emph{mandatory} precise citation and will not compile without one, so
  ``[standard]'' with no pointer is a build error; \COMP{} never upgrades to \PROVED{},
  and carries the range it covers and the validation that produces it.
\item \cl{sec:labelpolicy}, \cl{L1}--\cl{L6} --- a label is an API: never
  renamed, never deleted, never retargeted to a statement asserting something different;
  a retired statement keeps its label, attached to the erratum that retires it.
\item \cl{sec:authoring}, \cl{A1}--\cl{A6} --- every hypothesis at the point of
  use; no claim rests on ``standard''; definitions constructive where the code is;
  obligations as checkable predicates on emitted data; one symbol one meaning;
  diagnostics never branched on.
\item \cl{sec:versioning}, \cl{V1}--\cl{V4} --- this document's version
  convention, prompted by \cl{err:e17}. Honoured for the first time in v0.11.0.
\end{itemize}

\paragraph{Non-normative.}
\begin{itemize}\setlength{\itemsep}{0pt}
\item \cl{app:errata} --- the errata appendix, $20$ entries
  \cl{err:e1}--\cl{err:e20}, each giving the defective statement with its
  \code{path:line}, the correct statement, the evidence, and the landing site.
\item \cl{err:e18}, \cl{err:e19}, \cl{err:e20} --- three entries the
  architecture's list did not have: the [E]$\to$[C] label move (a consequence of the
  scope cut, not a defect), [v5] \S12's $\beta$ no-go, and [v5] \S12.3's bar-cochain
  formula omitting the coset permutation.
\item \cl{err:e1} --- the primitive-versus-full $T_{p^{\nu}}$ semantics settled
  empirically rather than assumed, by running the pipeline and not by reading it: the
  entry point returns the \emph{full} operator, the raw route the \emph{primitive}
  double coset, and both entry points are named. The hand-checkable instance at
  $(p,\ell,n)=(2,5,0)$ is recorded with the entry.
\item Erratum numbering unified to E1--E20 across the document and the architecture,
  which had been written from two different drafts; later chapters cite E-numbers in
  their template lines, so a chapter author reading the wrong list would have applied
  the wrong correction.
\end{itemize}

\subsection*{v0.1.0 --- 2026-09-15 --- the skeleton}

\noindent\emph{Commit \code{91cf200}.} $18$ pages. The architecture pass: every chapter
and appendix of the plan present, $144$ labels declared, and every statement other than
chapter \ref{sec:notation} a stub reading \code{[to be written in step 3]}. The purpose
was to prove that the merged label namespace compiles and is collision-free.

\paragraph{Normative.}
\begin{itemize}\setlength{\itemsep}{0pt}
\item \cl{sec:notation} --- the notation chapter, written in full because it
  binds every later chapter, and declared normative and exhaustive. No document of the
  superseded corpus had one. \cl{N1}--\cl{N5}; \cl{sec:notation-primes},
  \cl{sec:notation-framework}, \cl{sec:notation-cochains}, \cl{sec:notation-hecke},
  \cl{sec:notation-cone}; \cl{sec:notation-collisions} with $14$ resolved collisions
  C1--C14 (including \cl{coll:p}, and the three the corpus survey missed: $M$ with five
  meanings, $\Rel$ versus the coefficient ring, and $s/S$);
  \cl{sec:notation-crosswalk}; \cl{sec:notation-open}, \cl{O1}--\cl{O7}.
\end{itemize}

\paragraph{Non-normative.}
\begin{itemize}\setlength{\itemsep}{0pt}
\item \cl{app:labelindex} --- every externally cited label dereferenced, so a
  missing, renamed or duplicated label is a build error and not a silent breakage
  (\ref{L6}). The $68$ preserved labels --- $67$ from \file{hecke\_action.tex} plus
  \cl{thm:explicit_complex_equiv} from \file{monodromy\_v5.tex} --- are placed, $0$
  unplaceable, $0$ collisions; the three namespaces are pairwise disjoint.
\item Every other label of the document is declared here as a stub; each is written by
  the version recorded above, and the ones still unwritten are listed in
  \S\ref{app:changelog-gaps}.
\end{itemize}

\subsection{What is not recorded here}\label{app:changelog-gaps}

\paragraph{Not recorded, for want of evidence.}
Nothing. Every entry above is carried by a commit that names the statement it changed,
by the label census of that commit's revision, or by the text itself. Two classes of
change were deliberately \emph{excluded} rather than lost:

\begin{itemize}\setlength{\itemsep}{0pt}
\item \textbf{Changes outside this document} (in the sense of \ref{CL2}). Three of the
  ten commits also changed something else: v0.1.0's added the corpus survey, the
  architecture note and the build script; v0.2.0's reconciled the architecture note
  against what slot 1 had actually written; and v0.6.0's corrected a
  \code{HeckeOperatorFamily.h} comment that claimed the $T_{p^{\nu}}$ recursion is
  contract-backed unconditionally, while \cl{thm:relations} requires
  $6\in\cring^{\times}$, and filed owner direction OD-T1 in the owner-decision record.
  Those are not versions of this document and are recorded where they happened; the
  contract-side half of each is above.
\item \textbf{Page counts and build numbers} are given per version as context, not as
  entries.
\end{itemize}

\paragraph{Not yet written.}
\textbf{Nothing.} v0.19.0 wrote the last five stubs --- \S\ref{sec:example},
\S\ref{sec:implementation}, \S\ref{sec:machinery} and
Appendices~\ref{app:status} and \ref{app:formalization} --- so \ref{V5}\ref{V5i} is
met; v1.0.0 discharged \ref{V5}\ref{V5ii} (\cl{rmk:labelcensus}), so the leading version
field is $1$. \textbf{A written document with an unambiguous namespace is still not a
finished programme.} Two chapters state objects no code builds
(\S\ref{sec:notimplemented}; three until v1.3.3, when \S\ref{sec:boundarydatum} was
implemented); the construction item \ref{open:ell2} is \emph{decided}
(v1.3.0, \ref{oc1a}) but governs a truncated free model that nothing builds, and the
stabilisation theorem \ref{oc1c} that would supersede it is half proved
(\ref{lem:bulkperiodic}) and half open (\ref{rmk:bulkperiodicscope}(3)); \ref{obl:indep}
is partial; \ref{obl:iota}(iv), discharged
negatively at v0.13.0, is discharged positively on the canonical $\tau$ since v1.1.0,
and \ref{obl:iota}(v)'s sign prediction, discharged \emph{negatively} for the geometric
involution at v1.1.1, holds by definition and gates since the owner's v1.2.0 negation of
\cl{def:iota}(iv); and
the citation-census guard of \S\ref{sec:machinery} was specified and not
written until the slot-3 tooling (v1.8.1: \file{docs/matrix-engine/tools/label\_census.sh},
run by \code{ContractCensusValidation} since hecke-extensions slot~3; this sentence had not
been updated) --- v1.0.0's census was run by hand, as \cl{rmk:labelcensus} says, and found one
collision that a citation census could not have found. \ref{V5} does not claim otherwise:
its leading field speaks about this document, not about the implementation.

\paragraph{The two supersession steps \texorpdfstring{\ref{V4}}{(V4)} requires} ---
freezing \file{hecke\_action.tex} and \file{monodromy\_v5.tex} with a version marker and
a superseded-by banner --- \textbf{were taken at v1.0.0}: the first is frozen at
Version~1.0 (2026-07-02), the second at Version~5.0 (2026-07-03), and each banner directs
its reader to Appendix~\ref{app:errata} and says that a \code{[C,~<label>]} citation
resolves in this document. Their content is otherwise unchanged; they remain the
historical record, and a reader of either must still apply
Appendix~\ref{app:errata}.

\nocite{*}
\bibliographystyle{alpha}
\bibliography{gl2_refs}

\end{document}
